\section{Mean Field Conservation Laws}
\label{sec:laws}\label{sec:conservation-hyperbolic}

\paragraph{Gradient flow and covariance closure}
Although training evolves all neuron weights, its effect on the
predictor is governed by a fixed-size covariance matrix.
We use the width-normalized gradient flow
$(\dot U_t,\dot V_t)=-n\nabla_{(U,V)}\cL(U_t,V_t)$, namely
\begin{equation}\label{eq:mean-field-gradient-flow}
 \dot U_t=-E_t^\top V_t,\qquad
 \dot V_t=-E_tU_t,\qquad E_t:=\nabla\ell(V_tU_t^\top/n).
\end{equation}
The factor $n$ fixes a common time scale across widths; equivalently
this is Euclidean gradient flow of $n\cL$.
With $W_t:=\binom{U_t}{V_t}$, define the uncentered covariance
\[
 \Sigma_t:=\tfrac1nW_tW_t^\top
 =\begin{psmallmatrix}\Suu_t&\Suv_t\\\Svu_t&\Svv_t\end{psmallmatrix},
 \qquad
 \Suu_t:=\tfrac{U_tU_t^\top}{n},\quad
 \Svv_t:=\tfrac{V_tV_t^\top}{n},\quad
 \Svu_t:=\tfrac{V_tU_t^\top}{n}.
\]
The hidden blocks $\Suu_t,\Svv_t$ do not enter the loss, but govern
the predictor's motion. We use the same notation $\Sigma$ whether the
covariance is specified directly or realized by finite-width factors:
the closed equation depends only on this matrix.
Set
$G_E:=\left(\begin{smallmatrix}0&E^\top\\E&0\end{smallmatrix}\right)$,
$G_t:=G_{E_t}$, and $G(S):=G_{\nabla\ell(S)}$, so $\dot W_t=-G_tW_t$.
The split structure of the gradient is encoded by the signature matrix
\[
 J:=\diag(\Id_{d_u},-\Id_{d_v}),\qquad G_tJ+JG_t=0.
\]
Our first result combines covariance closure with width-independent spectral conservation laws.
\begin{proposition}[Covariance closure and conservation laws]
\label{prop:covariance-closure-laws}
Let $\mathcal H_k(\Sigma):=\tr\!\left((J\Sigma)^k\right)$.
On its interval of existence, the factor flow
\eqref{eq:mean-field-gradient-flow} closes on the covariance $\Sigma_t$ as a Bures equation, for which the characteristic polynomial of $J\Sigma$ is preserved:
\begin{equation}\label{eq:bures-G}
 \dot\Sigma_t=-(G_t\Sigma_t+\Sigma_tG_t), \qquad
 \tfrac{\dd}{\dd t}\mathcal H_k(\Sigma_t)=0, \quad k=1,\ldots,m,
\end{equation}
\end{proposition}

\begin{comment}
\emph{Proof sketch.}
Differentiate $W_tW_t^\top/n$ to obtain the closed equation.
The identity $JG_t=-G_tJ$ shows that $J\Sigma_t$ evolves following a Lax equation: $\frac{\dd}{\dd t}(J\Sigma_t)=[G_t,J\Sigma_t] = G_t J \Sigma_t - J \Sigma_t G_t$, whose trace powers
are constant by cyclicity of the trace. Newton identities give the
characteristic polynomial; see Appendix~\ref{app:covariance-laws} for a full proof.\qed
\end{comment}

The covariance equation \eqref{eq:bures-G} is the Bures gradient
flow~\citep{Takatsu2011,BhatiaJainLim2019} of
$F(\Sigma):=\ell(\Svu)$, since $\nabla_\Sigma F=G/2$.
Its fundamental matrix $\Phi_t$, defined by $\dot\Phi_t=-G_t\Phi_t$ and $\Phi_0=\Id$, satisfies
\begin{equation}\label{eq:orbit}
 \Phi_t^\top J\Phi_t=  \Phi_tJ\Phi_t^\top=J,\qquad
 \Sigma_t=\Phi_t\Sigma_0\Phi_t^\top.
\end{equation}
The functions $\mathcal H_k$ are conserved at every width $n$, defining mean-field laws.
Appendix~\ref{app:mean-field} gives their measure formulation, the corresponding Wasserstein PDE, and describes a larger family of laws.

\paragraph{Conserved spectral margins}
The spectral conservation laws \eqref{eq:bures-G} give uniform lower bounds on the spectrum of the hidden blocks $(\Suu, \Svv)$ through two margins
of the full covariance.
\begin{definition}[Covariance support margins and gap]\label{def:support-gap}
For $\Sigma\succeq0$, the support margins and their sum are
% Compact spacing local to this definition.
\setlength{\abovedisplayskip}{2pt}
\setlength{\belowdisplayskip}{2pt}
\setlength{\abovedisplayshortskip}{0pt}
\setlength{\belowdisplayshortskip}{2pt}
\begin{equation}\label{eq:margins}
 a_\pm(\Sigma):=\max\{a\ge0:\Sigma\mp aJ\succeq0\},\qquad
 \gap{\Sigma}:=\aplus{\Sigma}+\aminus{\Sigma}.
\end{equation}
\end{definition}
These maxima are finite and nonnegative. If $\lambda_1\le\cdots\le\lambda_m$
are the eigenvalues of $\Sigma^{1/2}J\Sigma^{1/2}$, \emph{including zeros}, then
\begin{equation}\label{eq:gap-spectrum}
 \aplus{\Sigma}=\lambda_{d_v+1},\qquad
 \aminus{\Sigma}=-\lambda_{d_v},\qquad
 \gap{\Sigma}=\lambda_{d_v+1}-\lambda_{d_v}.
\end{equation}
Thus $\gap{\Sigma}$ is an indexed spectral gap. Equivalently, it is the
length of the positive-semidefinite interval
$[-\aminus{\Sigma},\aplus{\Sigma}]$ of the pencil $\Sigma-tJ$;
such intervals are standard in matrix-pencil theory
\citep{NguyenNguyen2023}. Appendix~\ref{app:technical-details} proves
\eqref{eq:gap-spectrum}, including at singular covariances.
For $\Sigma\succ0$, $\aplus{\Sigma}$ and $-\aminus{\Sigma}$ are the
signed eigenvalues of $\Sigma^{1/2} J \Sigma^{1/2}$ nearest zero on their respective sides, and
$\gap{\Sigma}\ge2\lambda_{\min}(\Sigma)>0$.

\begin{proposition}[Transported lower bounds]\label{lem:margins}
On its interval of existence, every solution of \eqref{eq:bures-G} from
$\Sigma_0\succeq0$ preserves $a_\pm(\Sigma_t)=a_\pm(\Sigma_0)$ and satisfies
\[
  \Sigma_t \mp a_{\pm}(\Sigma_0) J \succeq 0, \qquad
  \Suu_t\succeq\aplus{\Sigma_0}\Id_{d_u},\qquad
  \Svv_t\succeq\aminus{\Sigma_0}\Id_{d_v}.
\]
\end{proposition}

\begin{comment}
\emph{Proof sketch.}
The key is to track matrix inequalities, rather than eigenvalues.
The supporting inequalities hold at $t = 0$ by definition. Congruence by
$\Phi_t$ preserves them because $\Phi_tJ\Phi_t^\top=J$; congruence by
$\Phi_t^{-1}$ proves that the maximal margins are unchanged.
Taking principal blocks then gives the hidden bounds. The complete proof is in
Appendix~\ref{app:technical-details}.\qed
\end{comment}

The margins $a_\pm(\Sigma)$ need not equal the minimum eigenvalues of the initial hidden blocks: they satisfy
$\aplus{\Sigma}\le\lambda_{\min}(\Suu)$ and
$\aminus{\Sigma}\le\lambda_{\min}(\Svv)$, with equality when $\Svu=0$.
For instance, letting $\Sigma=\left(\begin{smallmatrix}1&r\\r&1\end{smallmatrix}\right)$,
$|r|\le1$, both margins are $\sqrt{1-r^2}$, whereas both hidden variances are one.
The support gap $g(\Sigma)$ controls predictor dissipation and can remain positive even when the full covariance is degenerate.
\begin{wrapfigure}[22]{r}{0.32\textwidth}
  \centering
  \resizebox{\linewidth}{!}{%% Creator: Matplotlib, PGF backend
%%
%% To include the figure in your LaTeX document, write
%%   \input{<filename>.pgf}
%%
%% Make sure the required packages are loaded in your preamble
%%   \usepackage{pgf}
%%
%% Also ensure that all the required font packages are loaded; for instance,
%% the lmodern package is sometimes necessary when using math font.
%%   \usepackage{lmodern}
%%
%% Figures using additional raster images can only be included by \input if
%% they are in the same directory as the main LaTeX file. For loading figures
%% from other directories you can use the `import` package
%%   \usepackage{import}
%%
%% and then include the figures with
%%   \import{<path to file>}{<filename>.pgf}
%%
%% Matplotlib used the following preamble
%%   \def\mathdefault#1{#1}
%%   \everymath=\expandafter{\the\everymath\displaystyle}
%%   \IfFileExists{scrextend.sty}{
%%     \usepackage[fontsize=6.000000pt]{scrextend}
%%   }{
%%     \renewcommand{\normalsize}{\fontsize{6.000000}{7.200000}\selectfont}
%%     \normalsize
%%   }
%%   \usepackage{amsfonts}
%%   \usepackage{amssymb}
%%   \ifdefined\pdftexversion\else  % non-pdftex case.
%%     \usepackage{fontspec}
%%   \fi
%%   \makeatletter\@ifpackageloaded{underscore}{}{\usepackage[strings]{underscore}}\makeatother
%%
\begingroup%
\makeatletter%
\begin{pgfpicture}%
\pgfpathrectangle{\pgfpointorigin}{\pgfqpoint{2.150000in}{3.350000in}}%
\pgfusepath{use as bounding box, clip}%
\begin{pgfscope}%
\pgfsetbuttcap%
\pgfsetmiterjoin%
\definecolor{currentfill}{rgb}{1.000000,1.000000,1.000000}%
\pgfsetfillcolor{currentfill}%
\pgfsetlinewidth{0.000000pt}%
\definecolor{currentstroke}{rgb}{1.000000,1.000000,1.000000}%
\pgfsetstrokecolor{currentstroke}%
\pgfsetdash{}{0pt}%
\pgfpathmoveto{\pgfqpoint{0.000000in}{0.000000in}}%
\pgfpathlineto{\pgfqpoint{2.150000in}{0.000000in}}%
\pgfpathlineto{\pgfqpoint{2.150000in}{3.350000in}}%
\pgfpathlineto{\pgfqpoint{0.000000in}{3.350000in}}%
\pgfpathlineto{\pgfqpoint{0.000000in}{0.000000in}}%
\pgfpathclose%
\pgfusepath{fill}%
\end{pgfscope}%
\begin{pgfscope}%
\pgfsetbuttcap%
\pgfsetmiterjoin%
\definecolor{currentfill}{rgb}{1.000000,1.000000,1.000000}%
\pgfsetfillcolor{currentfill}%
\pgfsetlinewidth{0.000000pt}%
\definecolor{currentstroke}{rgb}{0.000000,0.000000,0.000000}%
\pgfsetstrokecolor{currentstroke}%
\pgfsetstrokeopacity{0.000000}%
\pgfsetdash{}{0pt}%
\pgfpathmoveto{\pgfqpoint{0.537500in}{2.059688in}}%
\pgfpathlineto{\pgfqpoint{2.042500in}{2.059688in}}%
\pgfpathlineto{\pgfqpoint{2.042500in}{3.216000in}}%
\pgfpathlineto{\pgfqpoint{0.537500in}{3.216000in}}%
\pgfpathlineto{\pgfqpoint{0.537500in}{2.059688in}}%
\pgfpathclose%
\pgfusepath{fill}%
\end{pgfscope}%
\begin{pgfscope}%
\pgfpathrectangle{\pgfqpoint{0.537500in}{2.059688in}}{\pgfqpoint{1.505000in}{1.156312in}}%
\pgfusepath{clip}%
\pgfsetbuttcap%
\pgfsetroundjoin%
\definecolor{currentfill}{rgb}{0.933333,0.941176,0.949020}%
\pgfsetfillcolor{currentfill}%
\pgfsetlinewidth{1.003750pt}%
\definecolor{currentstroke}{rgb}{0.933333,0.941176,0.949020}%
\pgfsetstrokecolor{currentstroke}%
\pgfsetdash{}{0pt}%
\pgfsys@defobject{currentmarker}{\pgfqpoint{0.537500in}{2.060627in}}{\pgfqpoint{2.042500in}{3.216000in}}{%
\pgfpathmoveto{\pgfqpoint{0.537500in}{3.216000in}}%
\pgfpathlineto{\pgfqpoint{0.537500in}{2.840980in}}%
\pgfpathlineto{\pgfqpoint{0.539653in}{2.838566in}}%
\pgfpathlineto{\pgfqpoint{0.541806in}{2.836151in}}%
\pgfpathlineto{\pgfqpoint{0.543959in}{2.833737in}}%
\pgfpathlineto{\pgfqpoint{0.546112in}{2.831323in}}%
\pgfpathlineto{\pgfqpoint{0.548265in}{2.828909in}}%
\pgfpathlineto{\pgfqpoint{0.550418in}{2.826494in}}%
\pgfpathlineto{\pgfqpoint{0.552572in}{2.824080in}}%
\pgfpathlineto{\pgfqpoint{0.554725in}{2.821666in}}%
\pgfpathlineto{\pgfqpoint{0.556878in}{2.819251in}}%
\pgfpathlineto{\pgfqpoint{0.559031in}{2.816837in}}%
\pgfpathlineto{\pgfqpoint{0.561184in}{2.814423in}}%
\pgfpathlineto{\pgfqpoint{0.563337in}{2.812009in}}%
\pgfpathlineto{\pgfqpoint{0.565490in}{2.809594in}}%
\pgfpathlineto{\pgfqpoint{0.567643in}{2.807180in}}%
\pgfpathlineto{\pgfqpoint{0.569796in}{2.804766in}}%
\pgfpathlineto{\pgfqpoint{0.571949in}{2.802351in}}%
\pgfpathlineto{\pgfqpoint{0.574102in}{2.799937in}}%
\pgfpathlineto{\pgfqpoint{0.576255in}{2.797523in}}%
\pgfpathlineto{\pgfqpoint{0.578408in}{2.795108in}}%
\pgfpathlineto{\pgfqpoint{0.580562in}{2.792694in}}%
\pgfpathlineto{\pgfqpoint{0.582715in}{2.790280in}}%
\pgfpathlineto{\pgfqpoint{0.584868in}{2.787866in}}%
\pgfpathlineto{\pgfqpoint{0.587021in}{2.785451in}}%
\pgfpathlineto{\pgfqpoint{0.589174in}{2.783037in}}%
\pgfpathlineto{\pgfqpoint{0.591327in}{2.780623in}}%
\pgfpathlineto{\pgfqpoint{0.593480in}{2.778208in}}%
\pgfpathlineto{\pgfqpoint{0.595633in}{2.775794in}}%
\pgfpathlineto{\pgfqpoint{0.597786in}{2.773380in}}%
\pgfpathlineto{\pgfqpoint{0.599939in}{2.770966in}}%
\pgfpathlineto{\pgfqpoint{0.602092in}{2.768551in}}%
\pgfpathlineto{\pgfqpoint{0.604245in}{2.766137in}}%
\pgfpathlineto{\pgfqpoint{0.606398in}{2.763723in}}%
\pgfpathlineto{\pgfqpoint{0.608552in}{2.761308in}}%
\pgfpathlineto{\pgfqpoint{0.610705in}{2.758894in}}%
\pgfpathlineto{\pgfqpoint{0.612858in}{2.756480in}}%
\pgfpathlineto{\pgfqpoint{0.615011in}{2.754066in}}%
\pgfpathlineto{\pgfqpoint{0.617164in}{2.751651in}}%
\pgfpathlineto{\pgfqpoint{0.619317in}{2.749237in}}%
\pgfpathlineto{\pgfqpoint{0.621470in}{2.746823in}}%
\pgfpathlineto{\pgfqpoint{0.623623in}{2.744408in}}%
\pgfpathlineto{\pgfqpoint{0.625776in}{2.741994in}}%
\pgfpathlineto{\pgfqpoint{0.627929in}{2.739580in}}%
\pgfpathlineto{\pgfqpoint{0.630082in}{2.737165in}}%
\pgfpathlineto{\pgfqpoint{0.632235in}{2.734751in}}%
\pgfpathlineto{\pgfqpoint{0.634388in}{2.732337in}}%
\pgfpathlineto{\pgfqpoint{0.636541in}{2.729923in}}%
\pgfpathlineto{\pgfqpoint{0.638695in}{2.727508in}}%
\pgfpathlineto{\pgfqpoint{0.640848in}{2.725094in}}%
\pgfpathlineto{\pgfqpoint{0.643001in}{2.722680in}}%
\pgfpathlineto{\pgfqpoint{0.645154in}{2.720265in}}%
\pgfpathlineto{\pgfqpoint{0.647307in}{2.717851in}}%
\pgfpathlineto{\pgfqpoint{0.649460in}{2.715437in}}%
\pgfpathlineto{\pgfqpoint{0.651613in}{2.713023in}}%
\pgfpathlineto{\pgfqpoint{0.653766in}{2.710608in}}%
\pgfpathlineto{\pgfqpoint{0.655919in}{2.708194in}}%
\pgfpathlineto{\pgfqpoint{0.658072in}{2.705780in}}%
\pgfpathlineto{\pgfqpoint{0.660225in}{2.703365in}}%
\pgfpathlineto{\pgfqpoint{0.662378in}{2.700951in}}%
\pgfpathlineto{\pgfqpoint{0.664531in}{2.698537in}}%
\pgfpathlineto{\pgfqpoint{0.666685in}{2.696123in}}%
\pgfpathlineto{\pgfqpoint{0.668838in}{2.693708in}}%
\pgfpathlineto{\pgfqpoint{0.670991in}{2.691294in}}%
\pgfpathlineto{\pgfqpoint{0.673144in}{2.688880in}}%
\pgfpathlineto{\pgfqpoint{0.675297in}{2.686465in}}%
\pgfpathlineto{\pgfqpoint{0.677450in}{2.684051in}}%
\pgfpathlineto{\pgfqpoint{0.679603in}{2.681637in}}%
\pgfpathlineto{\pgfqpoint{0.681756in}{2.679222in}}%
\pgfpathlineto{\pgfqpoint{0.683909in}{2.676808in}}%
\pgfpathlineto{\pgfqpoint{0.686062in}{2.674394in}}%
\pgfpathlineto{\pgfqpoint{0.688215in}{2.671980in}}%
\pgfpathlineto{\pgfqpoint{0.690368in}{2.669565in}}%
\pgfpathlineto{\pgfqpoint{0.692521in}{2.667151in}}%
\pgfpathlineto{\pgfqpoint{0.694675in}{2.664737in}}%
\pgfpathlineto{\pgfqpoint{0.696828in}{2.662322in}}%
\pgfpathlineto{\pgfqpoint{0.698981in}{2.659908in}}%
\pgfpathlineto{\pgfqpoint{0.701134in}{2.657494in}}%
\pgfpathlineto{\pgfqpoint{0.703287in}{2.655080in}}%
\pgfpathlineto{\pgfqpoint{0.705440in}{2.652665in}}%
\pgfpathlineto{\pgfqpoint{0.707593in}{2.650251in}}%
\pgfpathlineto{\pgfqpoint{0.709746in}{2.647837in}}%
\pgfpathlineto{\pgfqpoint{0.711899in}{2.645422in}}%
\pgfpathlineto{\pgfqpoint{0.714052in}{2.643008in}}%
\pgfpathlineto{\pgfqpoint{0.716205in}{2.640594in}}%
\pgfpathlineto{\pgfqpoint{0.718358in}{2.638179in}}%
\pgfpathlineto{\pgfqpoint{0.720511in}{2.635765in}}%
\pgfpathlineto{\pgfqpoint{0.722665in}{2.633351in}}%
\pgfpathlineto{\pgfqpoint{0.724818in}{2.630937in}}%
\pgfpathlineto{\pgfqpoint{0.726971in}{2.628522in}}%
\pgfpathlineto{\pgfqpoint{0.729124in}{2.626108in}}%
\pgfpathlineto{\pgfqpoint{0.731277in}{2.623694in}}%
\pgfpathlineto{\pgfqpoint{0.733430in}{2.621279in}}%
\pgfpathlineto{\pgfqpoint{0.735583in}{2.618865in}}%
\pgfpathlineto{\pgfqpoint{0.737736in}{2.616451in}}%
\pgfpathlineto{\pgfqpoint{0.739889in}{2.614037in}}%
\pgfpathlineto{\pgfqpoint{0.742042in}{2.611622in}}%
\pgfpathlineto{\pgfqpoint{0.744195in}{2.609208in}}%
\pgfpathlineto{\pgfqpoint{0.746348in}{2.606794in}}%
\pgfpathlineto{\pgfqpoint{0.748501in}{2.604379in}}%
\pgfpathlineto{\pgfqpoint{0.750655in}{2.601965in}}%
\pgfpathlineto{\pgfqpoint{0.752808in}{2.599551in}}%
\pgfpathlineto{\pgfqpoint{0.754961in}{2.597137in}}%
\pgfpathlineto{\pgfqpoint{0.757114in}{2.594722in}}%
\pgfpathlineto{\pgfqpoint{0.759267in}{2.592308in}}%
\pgfpathlineto{\pgfqpoint{0.761420in}{2.589894in}}%
\pgfpathlineto{\pgfqpoint{0.763573in}{2.587479in}}%
\pgfpathlineto{\pgfqpoint{0.765726in}{2.585065in}}%
\pgfpathlineto{\pgfqpoint{0.767879in}{2.582651in}}%
\pgfpathlineto{\pgfqpoint{0.770032in}{2.580236in}}%
\pgfpathlineto{\pgfqpoint{0.772185in}{2.577822in}}%
\pgfpathlineto{\pgfqpoint{0.774338in}{2.575408in}}%
\pgfpathlineto{\pgfqpoint{0.776491in}{2.572994in}}%
\pgfpathlineto{\pgfqpoint{0.778644in}{2.570579in}}%
\pgfpathlineto{\pgfqpoint{0.780798in}{2.568165in}}%
\pgfpathlineto{\pgfqpoint{0.782951in}{2.565751in}}%
\pgfpathlineto{\pgfqpoint{0.785104in}{2.563336in}}%
\pgfpathlineto{\pgfqpoint{0.787257in}{2.560922in}}%
\pgfpathlineto{\pgfqpoint{0.789410in}{2.558508in}}%
\pgfpathlineto{\pgfqpoint{0.791563in}{2.556094in}}%
\pgfpathlineto{\pgfqpoint{0.793716in}{2.553679in}}%
\pgfpathlineto{\pgfqpoint{0.795869in}{2.551265in}}%
\pgfpathlineto{\pgfqpoint{0.798022in}{2.548851in}}%
\pgfpathlineto{\pgfqpoint{0.800175in}{2.546436in}}%
\pgfpathlineto{\pgfqpoint{0.802328in}{2.544022in}}%
\pgfpathlineto{\pgfqpoint{0.804481in}{2.541608in}}%
\pgfpathlineto{\pgfqpoint{0.806634in}{2.539194in}}%
\pgfpathlineto{\pgfqpoint{0.808788in}{2.536779in}}%
\pgfpathlineto{\pgfqpoint{0.810941in}{2.534365in}}%
\pgfpathlineto{\pgfqpoint{0.813094in}{2.531951in}}%
\pgfpathlineto{\pgfqpoint{0.815247in}{2.529536in}}%
\pgfpathlineto{\pgfqpoint{0.817400in}{2.527122in}}%
\pgfpathlineto{\pgfqpoint{0.819553in}{2.524708in}}%
\pgfpathlineto{\pgfqpoint{0.821706in}{2.522293in}}%
\pgfpathlineto{\pgfqpoint{0.823859in}{2.519879in}}%
\pgfpathlineto{\pgfqpoint{0.826012in}{2.517465in}}%
\pgfpathlineto{\pgfqpoint{0.828165in}{2.515051in}}%
\pgfpathlineto{\pgfqpoint{0.830318in}{2.512636in}}%
\pgfpathlineto{\pgfqpoint{0.832471in}{2.510222in}}%
\pgfpathlineto{\pgfqpoint{0.834624in}{2.507808in}}%
\pgfpathlineto{\pgfqpoint{0.836778in}{2.505393in}}%
\pgfpathlineto{\pgfqpoint{0.838931in}{2.502979in}}%
\pgfpathlineto{\pgfqpoint{0.841084in}{2.500565in}}%
\pgfpathlineto{\pgfqpoint{0.843237in}{2.498151in}}%
\pgfpathlineto{\pgfqpoint{0.845390in}{2.495736in}}%
\pgfpathlineto{\pgfqpoint{0.847543in}{2.493322in}}%
\pgfpathlineto{\pgfqpoint{0.849696in}{2.490908in}}%
\pgfpathlineto{\pgfqpoint{0.851849in}{2.488493in}}%
\pgfpathlineto{\pgfqpoint{0.854002in}{2.486079in}}%
\pgfpathlineto{\pgfqpoint{0.856155in}{2.483665in}}%
\pgfpathlineto{\pgfqpoint{0.858308in}{2.481251in}}%
\pgfpathlineto{\pgfqpoint{0.860461in}{2.478836in}}%
\pgfpathlineto{\pgfqpoint{0.862614in}{2.476422in}}%
\pgfpathlineto{\pgfqpoint{0.864768in}{2.474008in}}%
\pgfpathlineto{\pgfqpoint{0.866921in}{2.471593in}}%
\pgfpathlineto{\pgfqpoint{0.869074in}{2.469179in}}%
\pgfpathlineto{\pgfqpoint{0.871227in}{2.466765in}}%
\pgfpathlineto{\pgfqpoint{0.873380in}{2.464350in}}%
\pgfpathlineto{\pgfqpoint{0.875533in}{2.461936in}}%
\pgfpathlineto{\pgfqpoint{0.877686in}{2.459522in}}%
\pgfpathlineto{\pgfqpoint{0.879839in}{2.457108in}}%
\pgfpathlineto{\pgfqpoint{0.881992in}{2.454693in}}%
\pgfpathlineto{\pgfqpoint{0.884145in}{2.452279in}}%
\pgfpathlineto{\pgfqpoint{0.886298in}{2.449865in}}%
\pgfpathlineto{\pgfqpoint{0.888451in}{2.447450in}}%
\pgfpathlineto{\pgfqpoint{0.890604in}{2.445036in}}%
\pgfpathlineto{\pgfqpoint{0.892758in}{2.442622in}}%
\pgfpathlineto{\pgfqpoint{0.894911in}{2.440208in}}%
\pgfpathlineto{\pgfqpoint{0.897064in}{2.437793in}}%
\pgfpathlineto{\pgfqpoint{0.899217in}{2.435379in}}%
\pgfpathlineto{\pgfqpoint{0.901370in}{2.432965in}}%
\pgfpathlineto{\pgfqpoint{0.903523in}{2.430550in}}%
\pgfpathlineto{\pgfqpoint{0.905676in}{2.428136in}}%
\pgfpathlineto{\pgfqpoint{0.907829in}{2.425722in}}%
\pgfpathlineto{\pgfqpoint{0.909982in}{2.423308in}}%
\pgfpathlineto{\pgfqpoint{0.912135in}{2.420893in}}%
\pgfpathlineto{\pgfqpoint{0.914288in}{2.418479in}}%
\pgfpathlineto{\pgfqpoint{0.916441in}{2.416065in}}%
\pgfpathlineto{\pgfqpoint{0.918594in}{2.413650in}}%
\pgfpathlineto{\pgfqpoint{0.920747in}{2.411236in}}%
\pgfpathlineto{\pgfqpoint{0.922901in}{2.408822in}}%
\pgfpathlineto{\pgfqpoint{0.925054in}{2.406407in}}%
\pgfpathlineto{\pgfqpoint{0.927207in}{2.403993in}}%
\pgfpathlineto{\pgfqpoint{0.929360in}{2.401579in}}%
\pgfpathlineto{\pgfqpoint{0.931513in}{2.399165in}}%
\pgfpathlineto{\pgfqpoint{0.933666in}{2.396750in}}%
\pgfpathlineto{\pgfqpoint{0.935819in}{2.394336in}}%
\pgfpathlineto{\pgfqpoint{0.937972in}{2.391922in}}%
\pgfpathlineto{\pgfqpoint{0.940125in}{2.389507in}}%
\pgfpathlineto{\pgfqpoint{0.942278in}{2.387093in}}%
\pgfpathlineto{\pgfqpoint{0.944431in}{2.384679in}}%
\pgfpathlineto{\pgfqpoint{0.946584in}{2.382265in}}%
\pgfpathlineto{\pgfqpoint{0.948737in}{2.379850in}}%
\pgfpathlineto{\pgfqpoint{0.950891in}{2.377436in}}%
\pgfpathlineto{\pgfqpoint{0.953044in}{2.375022in}}%
\pgfpathlineto{\pgfqpoint{0.955197in}{2.372607in}}%
\pgfpathlineto{\pgfqpoint{0.957350in}{2.370193in}}%
\pgfpathlineto{\pgfqpoint{0.959503in}{2.367779in}}%
\pgfpathlineto{\pgfqpoint{0.961656in}{2.365365in}}%
\pgfpathlineto{\pgfqpoint{0.963809in}{2.362950in}}%
\pgfpathlineto{\pgfqpoint{0.965962in}{2.360536in}}%
\pgfpathlineto{\pgfqpoint{0.968115in}{2.358122in}}%
\pgfpathlineto{\pgfqpoint{0.970268in}{2.355707in}}%
\pgfpathlineto{\pgfqpoint{0.972421in}{2.353293in}}%
\pgfpathlineto{\pgfqpoint{0.974574in}{2.350879in}}%
\pgfpathlineto{\pgfqpoint{0.976727in}{2.348464in}}%
\pgfpathlineto{\pgfqpoint{0.978881in}{2.346050in}}%
\pgfpathlineto{\pgfqpoint{0.981034in}{2.343636in}}%
\pgfpathlineto{\pgfqpoint{0.983187in}{2.341222in}}%
\pgfpathlineto{\pgfqpoint{0.985340in}{2.338807in}}%
\pgfpathlineto{\pgfqpoint{0.987493in}{2.336393in}}%
\pgfpathlineto{\pgfqpoint{0.989646in}{2.333979in}}%
\pgfpathlineto{\pgfqpoint{0.991799in}{2.331564in}}%
\pgfpathlineto{\pgfqpoint{0.993952in}{2.329150in}}%
\pgfpathlineto{\pgfqpoint{0.996105in}{2.326736in}}%
\pgfpathlineto{\pgfqpoint{0.998258in}{2.324322in}}%
\pgfpathlineto{\pgfqpoint{1.000411in}{2.321907in}}%
\pgfpathlineto{\pgfqpoint{1.002564in}{2.319493in}}%
\pgfpathlineto{\pgfqpoint{1.004717in}{2.317079in}}%
\pgfpathlineto{\pgfqpoint{1.006871in}{2.314664in}}%
\pgfpathlineto{\pgfqpoint{1.009024in}{2.312250in}}%
\pgfpathlineto{\pgfqpoint{1.011177in}{2.309836in}}%
\pgfpathlineto{\pgfqpoint{1.013330in}{2.307422in}}%
\pgfpathlineto{\pgfqpoint{1.015483in}{2.305007in}}%
\pgfpathlineto{\pgfqpoint{1.017636in}{2.302593in}}%
\pgfpathlineto{\pgfqpoint{1.019789in}{2.300179in}}%
\pgfpathlineto{\pgfqpoint{1.021942in}{2.297764in}}%
\pgfpathlineto{\pgfqpoint{1.024095in}{2.295350in}}%
\pgfpathlineto{\pgfqpoint{1.026248in}{2.292936in}}%
\pgfpathlineto{\pgfqpoint{1.028401in}{2.290521in}}%
\pgfpathlineto{\pgfqpoint{1.030554in}{2.288107in}}%
\pgfpathlineto{\pgfqpoint{1.032707in}{2.285693in}}%
\pgfpathlineto{\pgfqpoint{1.034861in}{2.283279in}}%
\pgfpathlineto{\pgfqpoint{1.037014in}{2.280864in}}%
\pgfpathlineto{\pgfqpoint{1.039167in}{2.278450in}}%
\pgfpathlineto{\pgfqpoint{1.041320in}{2.276036in}}%
\pgfpathlineto{\pgfqpoint{1.043473in}{2.273621in}}%
\pgfpathlineto{\pgfqpoint{1.045626in}{2.271207in}}%
\pgfpathlineto{\pgfqpoint{1.047779in}{2.268793in}}%
\pgfpathlineto{\pgfqpoint{1.049932in}{2.266379in}}%
\pgfpathlineto{\pgfqpoint{1.052085in}{2.263964in}}%
\pgfpathlineto{\pgfqpoint{1.054238in}{2.261550in}}%
\pgfpathlineto{\pgfqpoint{1.056391in}{2.259136in}}%
\pgfpathlineto{\pgfqpoint{1.058544in}{2.256721in}}%
\pgfpathlineto{\pgfqpoint{1.060697in}{2.254307in}}%
\pgfpathlineto{\pgfqpoint{1.062851in}{2.251893in}}%
\pgfpathlineto{\pgfqpoint{1.065004in}{2.249479in}}%
\pgfpathlineto{\pgfqpoint{1.067157in}{2.247064in}}%
\pgfpathlineto{\pgfqpoint{1.069310in}{2.244650in}}%
\pgfpathlineto{\pgfqpoint{1.071463in}{2.242236in}}%
\pgfpathlineto{\pgfqpoint{1.073616in}{2.239821in}}%
\pgfpathlineto{\pgfqpoint{1.075769in}{2.237407in}}%
\pgfpathlineto{\pgfqpoint{1.077922in}{2.234993in}}%
\pgfpathlineto{\pgfqpoint{1.080075in}{2.232578in}}%
\pgfpathlineto{\pgfqpoint{1.082228in}{2.230164in}}%
\pgfpathlineto{\pgfqpoint{1.084381in}{2.227750in}}%
\pgfpathlineto{\pgfqpoint{1.086534in}{2.225336in}}%
\pgfpathlineto{\pgfqpoint{1.088687in}{2.222921in}}%
\pgfpathlineto{\pgfqpoint{1.090840in}{2.220507in}}%
\pgfpathlineto{\pgfqpoint{1.092994in}{2.218093in}}%
\pgfpathlineto{\pgfqpoint{1.095147in}{2.215678in}}%
\pgfpathlineto{\pgfqpoint{1.097300in}{2.213264in}}%
\pgfpathlineto{\pgfqpoint{1.099453in}{2.210850in}}%
\pgfpathlineto{\pgfqpoint{1.101606in}{2.208436in}}%
\pgfpathlineto{\pgfqpoint{1.103759in}{2.206021in}}%
\pgfpathlineto{\pgfqpoint{1.105912in}{2.203607in}}%
\pgfpathlineto{\pgfqpoint{1.108065in}{2.201193in}}%
\pgfpathlineto{\pgfqpoint{1.110218in}{2.198778in}}%
\pgfpathlineto{\pgfqpoint{1.112371in}{2.196364in}}%
\pgfpathlineto{\pgfqpoint{1.114524in}{2.193950in}}%
\pgfpathlineto{\pgfqpoint{1.116677in}{2.191536in}}%
\pgfpathlineto{\pgfqpoint{1.118830in}{2.189121in}}%
\pgfpathlineto{\pgfqpoint{1.120984in}{2.186707in}}%
\pgfpathlineto{\pgfqpoint{1.123137in}{2.184293in}}%
\pgfpathlineto{\pgfqpoint{1.125290in}{2.181878in}}%
\pgfpathlineto{\pgfqpoint{1.127443in}{2.179464in}}%
\pgfpathlineto{\pgfqpoint{1.129596in}{2.177050in}}%
\pgfpathlineto{\pgfqpoint{1.131749in}{2.174635in}}%
\pgfpathlineto{\pgfqpoint{1.133902in}{2.172221in}}%
\pgfpathlineto{\pgfqpoint{1.136055in}{2.169807in}}%
\pgfpathlineto{\pgfqpoint{1.138208in}{2.167393in}}%
\pgfpathlineto{\pgfqpoint{1.140361in}{2.164978in}}%
\pgfpathlineto{\pgfqpoint{1.142514in}{2.162564in}}%
\pgfpathlineto{\pgfqpoint{1.144667in}{2.160150in}}%
\pgfpathlineto{\pgfqpoint{1.146820in}{2.157735in}}%
\pgfpathlineto{\pgfqpoint{1.148974in}{2.155321in}}%
\pgfpathlineto{\pgfqpoint{1.151127in}{2.152907in}}%
\pgfpathlineto{\pgfqpoint{1.153280in}{2.150493in}}%
\pgfpathlineto{\pgfqpoint{1.155433in}{2.148078in}}%
\pgfpathlineto{\pgfqpoint{1.157586in}{2.145664in}}%
\pgfpathlineto{\pgfqpoint{1.159739in}{2.143250in}}%
\pgfpathlineto{\pgfqpoint{1.161892in}{2.140835in}}%
\pgfpathlineto{\pgfqpoint{1.164045in}{2.138421in}}%
\pgfpathlineto{\pgfqpoint{1.166198in}{2.136007in}}%
\pgfpathlineto{\pgfqpoint{1.168351in}{2.133593in}}%
\pgfpathlineto{\pgfqpoint{1.170504in}{2.131178in}}%
\pgfpathlineto{\pgfqpoint{1.172657in}{2.128764in}}%
\pgfpathlineto{\pgfqpoint{1.174810in}{2.126350in}}%
\pgfpathlineto{\pgfqpoint{1.176964in}{2.123935in}}%
\pgfpathlineto{\pgfqpoint{1.179117in}{2.121521in}}%
\pgfpathlineto{\pgfqpoint{1.181270in}{2.119107in}}%
\pgfpathlineto{\pgfqpoint{1.183423in}{2.116692in}}%
\pgfpathlineto{\pgfqpoint{1.185576in}{2.114278in}}%
\pgfpathlineto{\pgfqpoint{1.187729in}{2.111864in}}%
\pgfpathlineto{\pgfqpoint{1.189882in}{2.109450in}}%
\pgfpathlineto{\pgfqpoint{1.192035in}{2.107035in}}%
\pgfpathlineto{\pgfqpoint{1.194188in}{2.104621in}}%
\pgfpathlineto{\pgfqpoint{1.196341in}{2.102207in}}%
\pgfpathlineto{\pgfqpoint{1.198494in}{2.099792in}}%
\pgfpathlineto{\pgfqpoint{1.200647in}{2.097378in}}%
\pgfpathlineto{\pgfqpoint{1.202800in}{2.094964in}}%
\pgfpathlineto{\pgfqpoint{1.204954in}{2.092550in}}%
\pgfpathlineto{\pgfqpoint{1.207107in}{2.090135in}}%
\pgfpathlineto{\pgfqpoint{1.209260in}{2.087721in}}%
\pgfpathlineto{\pgfqpoint{1.211413in}{2.085307in}}%
\pgfpathlineto{\pgfqpoint{1.213566in}{2.082892in}}%
\pgfpathlineto{\pgfqpoint{1.215719in}{2.080478in}}%
\pgfpathlineto{\pgfqpoint{1.217872in}{2.078064in}}%
\pgfpathlineto{\pgfqpoint{1.220025in}{2.075650in}}%
\pgfpathlineto{\pgfqpoint{1.222178in}{2.073235in}}%
\pgfpathlineto{\pgfqpoint{1.224331in}{2.070821in}}%
\pgfpathlineto{\pgfqpoint{1.226484in}{2.068407in}}%
\pgfpathlineto{\pgfqpoint{1.228637in}{2.065992in}}%
\pgfpathlineto{\pgfqpoint{1.230790in}{2.063578in}}%
\pgfpathlineto{\pgfqpoint{1.232943in}{2.061164in}}%
\pgfpathlineto{\pgfqpoint{1.235097in}{2.060627in}}%
\pgfpathlineto{\pgfqpoint{1.237250in}{2.063042in}}%
\pgfpathlineto{\pgfqpoint{1.239403in}{2.065456in}}%
\pgfpathlineto{\pgfqpoint{1.241556in}{2.067870in}}%
\pgfpathlineto{\pgfqpoint{1.243709in}{2.070284in}}%
\pgfpathlineto{\pgfqpoint{1.245862in}{2.072699in}}%
\pgfpathlineto{\pgfqpoint{1.248015in}{2.075113in}}%
\pgfpathlineto{\pgfqpoint{1.250168in}{2.077527in}}%
\pgfpathlineto{\pgfqpoint{1.252321in}{2.079942in}}%
\pgfpathlineto{\pgfqpoint{1.254474in}{2.082356in}}%
\pgfpathlineto{\pgfqpoint{1.256627in}{2.084770in}}%
\pgfpathlineto{\pgfqpoint{1.258780in}{2.087184in}}%
\pgfpathlineto{\pgfqpoint{1.260933in}{2.089599in}}%
\pgfpathlineto{\pgfqpoint{1.263087in}{2.092013in}}%
\pgfpathlineto{\pgfqpoint{1.265240in}{2.094427in}}%
\pgfpathlineto{\pgfqpoint{1.267393in}{2.096842in}}%
\pgfpathlineto{\pgfqpoint{1.269546in}{2.099256in}}%
\pgfpathlineto{\pgfqpoint{1.271699in}{2.101670in}}%
\pgfpathlineto{\pgfqpoint{1.273852in}{2.104085in}}%
\pgfpathlineto{\pgfqpoint{1.276005in}{2.106499in}}%
\pgfpathlineto{\pgfqpoint{1.278158in}{2.108913in}}%
\pgfpathlineto{\pgfqpoint{1.280311in}{2.111327in}}%
\pgfpathlineto{\pgfqpoint{1.282464in}{2.113742in}}%
\pgfpathlineto{\pgfqpoint{1.284617in}{2.116156in}}%
\pgfpathlineto{\pgfqpoint{1.286770in}{2.118570in}}%
\pgfpathlineto{\pgfqpoint{1.288923in}{2.120985in}}%
\pgfpathlineto{\pgfqpoint{1.291077in}{2.123399in}}%
\pgfpathlineto{\pgfqpoint{1.293230in}{2.125813in}}%
\pgfpathlineto{\pgfqpoint{1.295383in}{2.128227in}}%
\pgfpathlineto{\pgfqpoint{1.297536in}{2.130642in}}%
\pgfpathlineto{\pgfqpoint{1.299689in}{2.133056in}}%
\pgfpathlineto{\pgfqpoint{1.301842in}{2.135470in}}%
\pgfpathlineto{\pgfqpoint{1.303995in}{2.137885in}}%
\pgfpathlineto{\pgfqpoint{1.306148in}{2.140299in}}%
\pgfpathlineto{\pgfqpoint{1.308301in}{2.142713in}}%
\pgfpathlineto{\pgfqpoint{1.310454in}{2.145127in}}%
\pgfpathlineto{\pgfqpoint{1.312607in}{2.147542in}}%
\pgfpathlineto{\pgfqpoint{1.314760in}{2.149956in}}%
\pgfpathlineto{\pgfqpoint{1.316913in}{2.152370in}}%
\pgfpathlineto{\pgfqpoint{1.319067in}{2.154785in}}%
\pgfpathlineto{\pgfqpoint{1.321220in}{2.157199in}}%
\pgfpathlineto{\pgfqpoint{1.323373in}{2.159613in}}%
\pgfpathlineto{\pgfqpoint{1.325526in}{2.162028in}}%
\pgfpathlineto{\pgfqpoint{1.327679in}{2.164442in}}%
\pgfpathlineto{\pgfqpoint{1.329832in}{2.166856in}}%
\pgfpathlineto{\pgfqpoint{1.331985in}{2.169270in}}%
\pgfpathlineto{\pgfqpoint{1.334138in}{2.171685in}}%
\pgfpathlineto{\pgfqpoint{1.336291in}{2.174099in}}%
\pgfpathlineto{\pgfqpoint{1.338444in}{2.176513in}}%
\pgfpathlineto{\pgfqpoint{1.340597in}{2.178928in}}%
\pgfpathlineto{\pgfqpoint{1.342750in}{2.181342in}}%
\pgfpathlineto{\pgfqpoint{1.344903in}{2.183756in}}%
\pgfpathlineto{\pgfqpoint{1.347057in}{2.186170in}}%
\pgfpathlineto{\pgfqpoint{1.349210in}{2.188585in}}%
\pgfpathlineto{\pgfqpoint{1.351363in}{2.190999in}}%
\pgfpathlineto{\pgfqpoint{1.353516in}{2.193413in}}%
\pgfpathlineto{\pgfqpoint{1.355669in}{2.195828in}}%
\pgfpathlineto{\pgfqpoint{1.357822in}{2.198242in}}%
\pgfpathlineto{\pgfqpoint{1.359975in}{2.200656in}}%
\pgfpathlineto{\pgfqpoint{1.362128in}{2.203070in}}%
\pgfpathlineto{\pgfqpoint{1.364281in}{2.205485in}}%
\pgfpathlineto{\pgfqpoint{1.366434in}{2.207899in}}%
\pgfpathlineto{\pgfqpoint{1.368587in}{2.210313in}}%
\pgfpathlineto{\pgfqpoint{1.370740in}{2.212728in}}%
\pgfpathlineto{\pgfqpoint{1.372893in}{2.215142in}}%
\pgfpathlineto{\pgfqpoint{1.375046in}{2.217556in}}%
\pgfpathlineto{\pgfqpoint{1.377200in}{2.219971in}}%
\pgfpathlineto{\pgfqpoint{1.379353in}{2.222385in}}%
\pgfpathlineto{\pgfqpoint{1.381506in}{2.224799in}}%
\pgfpathlineto{\pgfqpoint{1.383659in}{2.227213in}}%
\pgfpathlineto{\pgfqpoint{1.385812in}{2.229628in}}%
\pgfpathlineto{\pgfqpoint{1.387965in}{2.232042in}}%
\pgfpathlineto{\pgfqpoint{1.390118in}{2.234456in}}%
\pgfpathlineto{\pgfqpoint{1.392271in}{2.236871in}}%
\pgfpathlineto{\pgfqpoint{1.394424in}{2.239285in}}%
\pgfpathlineto{\pgfqpoint{1.396577in}{2.241699in}}%
\pgfpathlineto{\pgfqpoint{1.398730in}{2.244113in}}%
\pgfpathlineto{\pgfqpoint{1.400883in}{2.246528in}}%
\pgfpathlineto{\pgfqpoint{1.403036in}{2.248942in}}%
\pgfpathlineto{\pgfqpoint{1.405190in}{2.251356in}}%
\pgfpathlineto{\pgfqpoint{1.407343in}{2.253771in}}%
\pgfpathlineto{\pgfqpoint{1.409496in}{2.256185in}}%
\pgfpathlineto{\pgfqpoint{1.411649in}{2.258599in}}%
\pgfpathlineto{\pgfqpoint{1.413802in}{2.261013in}}%
\pgfpathlineto{\pgfqpoint{1.415955in}{2.263428in}}%
\pgfpathlineto{\pgfqpoint{1.418108in}{2.265842in}}%
\pgfpathlineto{\pgfqpoint{1.420261in}{2.268256in}}%
\pgfpathlineto{\pgfqpoint{1.422414in}{2.270671in}}%
\pgfpathlineto{\pgfqpoint{1.424567in}{2.273085in}}%
\pgfpathlineto{\pgfqpoint{1.426720in}{2.275499in}}%
\pgfpathlineto{\pgfqpoint{1.428873in}{2.277914in}}%
\pgfpathlineto{\pgfqpoint{1.431026in}{2.280328in}}%
\pgfpathlineto{\pgfqpoint{1.433180in}{2.282742in}}%
\pgfpathlineto{\pgfqpoint{1.435333in}{2.285156in}}%
\pgfpathlineto{\pgfqpoint{1.437486in}{2.287571in}}%
\pgfpathlineto{\pgfqpoint{1.439639in}{2.289985in}}%
\pgfpathlineto{\pgfqpoint{1.441792in}{2.292399in}}%
\pgfpathlineto{\pgfqpoint{1.443945in}{2.294814in}}%
\pgfpathlineto{\pgfqpoint{1.446098in}{2.297228in}}%
\pgfpathlineto{\pgfqpoint{1.448251in}{2.299642in}}%
\pgfpathlineto{\pgfqpoint{1.450404in}{2.302056in}}%
\pgfpathlineto{\pgfqpoint{1.452557in}{2.304471in}}%
\pgfpathlineto{\pgfqpoint{1.454710in}{2.306885in}}%
\pgfpathlineto{\pgfqpoint{1.456863in}{2.309299in}}%
\pgfpathlineto{\pgfqpoint{1.459016in}{2.311714in}}%
\pgfpathlineto{\pgfqpoint{1.461170in}{2.314128in}}%
\pgfpathlineto{\pgfqpoint{1.463323in}{2.316542in}}%
\pgfpathlineto{\pgfqpoint{1.465476in}{2.318956in}}%
\pgfpathlineto{\pgfqpoint{1.467629in}{2.321371in}}%
\pgfpathlineto{\pgfqpoint{1.469782in}{2.323785in}}%
\pgfpathlineto{\pgfqpoint{1.471935in}{2.326199in}}%
\pgfpathlineto{\pgfqpoint{1.474088in}{2.328614in}}%
\pgfpathlineto{\pgfqpoint{1.476241in}{2.331028in}}%
\pgfpathlineto{\pgfqpoint{1.478394in}{2.333442in}}%
\pgfpathlineto{\pgfqpoint{1.480547in}{2.335857in}}%
\pgfpathlineto{\pgfqpoint{1.482700in}{2.338271in}}%
\pgfpathlineto{\pgfqpoint{1.484853in}{2.340685in}}%
\pgfpathlineto{\pgfqpoint{1.487006in}{2.343099in}}%
\pgfpathlineto{\pgfqpoint{1.489160in}{2.345514in}}%
\pgfpathlineto{\pgfqpoint{1.491313in}{2.347928in}}%
\pgfpathlineto{\pgfqpoint{1.493466in}{2.350342in}}%
\pgfpathlineto{\pgfqpoint{1.495619in}{2.352757in}}%
\pgfpathlineto{\pgfqpoint{1.497772in}{2.355171in}}%
\pgfpathlineto{\pgfqpoint{1.499925in}{2.357585in}}%
\pgfpathlineto{\pgfqpoint{1.502078in}{2.359999in}}%
\pgfpathlineto{\pgfqpoint{1.504231in}{2.362414in}}%
\pgfpathlineto{\pgfqpoint{1.506384in}{2.364828in}}%
\pgfpathlineto{\pgfqpoint{1.508537in}{2.367242in}}%
\pgfpathlineto{\pgfqpoint{1.510690in}{2.369657in}}%
\pgfpathlineto{\pgfqpoint{1.512843in}{2.372071in}}%
\pgfpathlineto{\pgfqpoint{1.514996in}{2.374485in}}%
\pgfpathlineto{\pgfqpoint{1.517149in}{2.376899in}}%
\pgfpathlineto{\pgfqpoint{1.519303in}{2.379314in}}%
\pgfpathlineto{\pgfqpoint{1.521456in}{2.381728in}}%
\pgfpathlineto{\pgfqpoint{1.523609in}{2.384142in}}%
\pgfpathlineto{\pgfqpoint{1.525762in}{2.386557in}}%
\pgfpathlineto{\pgfqpoint{1.527915in}{2.388971in}}%
\pgfpathlineto{\pgfqpoint{1.530068in}{2.391385in}}%
\pgfpathlineto{\pgfqpoint{1.532221in}{2.393800in}}%
\pgfpathlineto{\pgfqpoint{1.534374in}{2.396214in}}%
\pgfpathlineto{\pgfqpoint{1.536527in}{2.398628in}}%
\pgfpathlineto{\pgfqpoint{1.538680in}{2.401042in}}%
\pgfpathlineto{\pgfqpoint{1.540833in}{2.403457in}}%
\pgfpathlineto{\pgfqpoint{1.542986in}{2.405871in}}%
\pgfpathlineto{\pgfqpoint{1.545139in}{2.408285in}}%
\pgfpathlineto{\pgfqpoint{1.547293in}{2.410700in}}%
\pgfpathlineto{\pgfqpoint{1.549446in}{2.413114in}}%
\pgfpathlineto{\pgfqpoint{1.551599in}{2.415528in}}%
\pgfpathlineto{\pgfqpoint{1.553752in}{2.417942in}}%
\pgfpathlineto{\pgfqpoint{1.555905in}{2.420357in}}%
\pgfpathlineto{\pgfqpoint{1.558058in}{2.422771in}}%
\pgfpathlineto{\pgfqpoint{1.560211in}{2.425185in}}%
\pgfpathlineto{\pgfqpoint{1.562364in}{2.427600in}}%
\pgfpathlineto{\pgfqpoint{1.564517in}{2.430014in}}%
\pgfpathlineto{\pgfqpoint{1.566670in}{2.432428in}}%
\pgfpathlineto{\pgfqpoint{1.568823in}{2.434842in}}%
\pgfpathlineto{\pgfqpoint{1.570976in}{2.437257in}}%
\pgfpathlineto{\pgfqpoint{1.573129in}{2.439671in}}%
\pgfpathlineto{\pgfqpoint{1.575283in}{2.442085in}}%
\pgfpathlineto{\pgfqpoint{1.577436in}{2.444500in}}%
\pgfpathlineto{\pgfqpoint{1.579589in}{2.446914in}}%
\pgfpathlineto{\pgfqpoint{1.581742in}{2.449328in}}%
\pgfpathlineto{\pgfqpoint{1.583895in}{2.451743in}}%
\pgfpathlineto{\pgfqpoint{1.586048in}{2.454157in}}%
\pgfpathlineto{\pgfqpoint{1.588201in}{2.456571in}}%
\pgfpathlineto{\pgfqpoint{1.590354in}{2.458985in}}%
\pgfpathlineto{\pgfqpoint{1.592507in}{2.461400in}}%
\pgfpathlineto{\pgfqpoint{1.594660in}{2.463814in}}%
\pgfpathlineto{\pgfqpoint{1.596813in}{2.466228in}}%
\pgfpathlineto{\pgfqpoint{1.598966in}{2.468643in}}%
\pgfpathlineto{\pgfqpoint{1.601119in}{2.471057in}}%
\pgfpathlineto{\pgfqpoint{1.603273in}{2.473471in}}%
\pgfpathlineto{\pgfqpoint{1.605426in}{2.475885in}}%
\pgfpathlineto{\pgfqpoint{1.607579in}{2.478300in}}%
\pgfpathlineto{\pgfqpoint{1.609732in}{2.480714in}}%
\pgfpathlineto{\pgfqpoint{1.611885in}{2.483128in}}%
\pgfpathlineto{\pgfqpoint{1.614038in}{2.485543in}}%
\pgfpathlineto{\pgfqpoint{1.616191in}{2.487957in}}%
\pgfpathlineto{\pgfqpoint{1.618344in}{2.490371in}}%
\pgfpathlineto{\pgfqpoint{1.620497in}{2.492785in}}%
\pgfpathlineto{\pgfqpoint{1.622650in}{2.495200in}}%
\pgfpathlineto{\pgfqpoint{1.624803in}{2.497614in}}%
\pgfpathlineto{\pgfqpoint{1.626956in}{2.500028in}}%
\pgfpathlineto{\pgfqpoint{1.629109in}{2.502443in}}%
\pgfpathlineto{\pgfqpoint{1.631263in}{2.504857in}}%
\pgfpathlineto{\pgfqpoint{1.633416in}{2.507271in}}%
\pgfpathlineto{\pgfqpoint{1.635569in}{2.509686in}}%
\pgfpathlineto{\pgfqpoint{1.637722in}{2.512100in}}%
\pgfpathlineto{\pgfqpoint{1.639875in}{2.514514in}}%
\pgfpathlineto{\pgfqpoint{1.642028in}{2.516928in}}%
\pgfpathlineto{\pgfqpoint{1.644181in}{2.519343in}}%
\pgfpathlineto{\pgfqpoint{1.646334in}{2.521757in}}%
\pgfpathlineto{\pgfqpoint{1.648487in}{2.524171in}}%
\pgfpathlineto{\pgfqpoint{1.650640in}{2.526586in}}%
\pgfpathlineto{\pgfqpoint{1.652793in}{2.529000in}}%
\pgfpathlineto{\pgfqpoint{1.654946in}{2.531414in}}%
\pgfpathlineto{\pgfqpoint{1.657099in}{2.533828in}}%
\pgfpathlineto{\pgfqpoint{1.659253in}{2.536243in}}%
\pgfpathlineto{\pgfqpoint{1.661406in}{2.538657in}}%
\pgfpathlineto{\pgfqpoint{1.663559in}{2.541071in}}%
\pgfpathlineto{\pgfqpoint{1.665712in}{2.543486in}}%
\pgfpathlineto{\pgfqpoint{1.667865in}{2.545900in}}%
\pgfpathlineto{\pgfqpoint{1.670018in}{2.548314in}}%
\pgfpathlineto{\pgfqpoint{1.672171in}{2.550728in}}%
\pgfpathlineto{\pgfqpoint{1.674324in}{2.553143in}}%
\pgfpathlineto{\pgfqpoint{1.676477in}{2.555557in}}%
\pgfpathlineto{\pgfqpoint{1.678630in}{2.557971in}}%
\pgfpathlineto{\pgfqpoint{1.680783in}{2.560386in}}%
\pgfpathlineto{\pgfqpoint{1.682936in}{2.562800in}}%
\pgfpathlineto{\pgfqpoint{1.685089in}{2.565214in}}%
\pgfpathlineto{\pgfqpoint{1.687242in}{2.567629in}}%
\pgfpathlineto{\pgfqpoint{1.689396in}{2.570043in}}%
\pgfpathlineto{\pgfqpoint{1.691549in}{2.572457in}}%
\pgfpathlineto{\pgfqpoint{1.693702in}{2.574871in}}%
\pgfpathlineto{\pgfqpoint{1.695855in}{2.577286in}}%
\pgfpathlineto{\pgfqpoint{1.698008in}{2.579700in}}%
\pgfpathlineto{\pgfqpoint{1.700161in}{2.582114in}}%
\pgfpathlineto{\pgfqpoint{1.702314in}{2.584529in}}%
\pgfpathlineto{\pgfqpoint{1.704467in}{2.586943in}}%
\pgfpathlineto{\pgfqpoint{1.706620in}{2.589357in}}%
\pgfpathlineto{\pgfqpoint{1.708773in}{2.591771in}}%
\pgfpathlineto{\pgfqpoint{1.710926in}{2.594186in}}%
\pgfpathlineto{\pgfqpoint{1.713079in}{2.596600in}}%
\pgfpathlineto{\pgfqpoint{1.715232in}{2.599014in}}%
\pgfpathlineto{\pgfqpoint{1.717386in}{2.601429in}}%
\pgfpathlineto{\pgfqpoint{1.719539in}{2.603843in}}%
\pgfpathlineto{\pgfqpoint{1.721692in}{2.606257in}}%
\pgfpathlineto{\pgfqpoint{1.723845in}{2.608671in}}%
\pgfpathlineto{\pgfqpoint{1.725998in}{2.611086in}}%
\pgfpathlineto{\pgfqpoint{1.728151in}{2.613500in}}%
\pgfpathlineto{\pgfqpoint{1.730304in}{2.615914in}}%
\pgfpathlineto{\pgfqpoint{1.732457in}{2.618329in}}%
\pgfpathlineto{\pgfqpoint{1.734610in}{2.620743in}}%
\pgfpathlineto{\pgfqpoint{1.736763in}{2.623157in}}%
\pgfpathlineto{\pgfqpoint{1.738916in}{2.625572in}}%
\pgfpathlineto{\pgfqpoint{1.741069in}{2.627986in}}%
\pgfpathlineto{\pgfqpoint{1.743222in}{2.630400in}}%
\pgfpathlineto{\pgfqpoint{1.745376in}{2.632814in}}%
\pgfpathlineto{\pgfqpoint{1.747529in}{2.635229in}}%
\pgfpathlineto{\pgfqpoint{1.749682in}{2.637643in}}%
\pgfpathlineto{\pgfqpoint{1.751835in}{2.640057in}}%
\pgfpathlineto{\pgfqpoint{1.753988in}{2.642472in}}%
\pgfpathlineto{\pgfqpoint{1.756141in}{2.644886in}}%
\pgfpathlineto{\pgfqpoint{1.758294in}{2.647300in}}%
\pgfpathlineto{\pgfqpoint{1.760447in}{2.649714in}}%
\pgfpathlineto{\pgfqpoint{1.762600in}{2.652129in}}%
\pgfpathlineto{\pgfqpoint{1.764753in}{2.654543in}}%
\pgfpathlineto{\pgfqpoint{1.766906in}{2.656957in}}%
\pgfpathlineto{\pgfqpoint{1.769059in}{2.659372in}}%
\pgfpathlineto{\pgfqpoint{1.771212in}{2.661786in}}%
\pgfpathlineto{\pgfqpoint{1.773366in}{2.664200in}}%
\pgfpathlineto{\pgfqpoint{1.775519in}{2.666614in}}%
\pgfpathlineto{\pgfqpoint{1.777672in}{2.669029in}}%
\pgfpathlineto{\pgfqpoint{1.779825in}{2.671443in}}%
\pgfpathlineto{\pgfqpoint{1.781978in}{2.673857in}}%
\pgfpathlineto{\pgfqpoint{1.784131in}{2.676272in}}%
\pgfpathlineto{\pgfqpoint{1.786284in}{2.678686in}}%
\pgfpathlineto{\pgfqpoint{1.788437in}{2.681100in}}%
\pgfpathlineto{\pgfqpoint{1.790590in}{2.683515in}}%
\pgfpathlineto{\pgfqpoint{1.792743in}{2.685929in}}%
\pgfpathlineto{\pgfqpoint{1.794896in}{2.688343in}}%
\pgfpathlineto{\pgfqpoint{1.797049in}{2.690757in}}%
\pgfpathlineto{\pgfqpoint{1.799202in}{2.693172in}}%
\pgfpathlineto{\pgfqpoint{1.801356in}{2.695586in}}%
\pgfpathlineto{\pgfqpoint{1.803509in}{2.698000in}}%
\pgfpathlineto{\pgfqpoint{1.805662in}{2.700415in}}%
\pgfpathlineto{\pgfqpoint{1.807815in}{2.702829in}}%
\pgfpathlineto{\pgfqpoint{1.809968in}{2.705243in}}%
\pgfpathlineto{\pgfqpoint{1.812121in}{2.707657in}}%
\pgfpathlineto{\pgfqpoint{1.814274in}{2.710072in}}%
\pgfpathlineto{\pgfqpoint{1.816427in}{2.712486in}}%
\pgfpathlineto{\pgfqpoint{1.818580in}{2.714900in}}%
\pgfpathlineto{\pgfqpoint{1.820733in}{2.717315in}}%
\pgfpathlineto{\pgfqpoint{1.822886in}{2.719729in}}%
\pgfpathlineto{\pgfqpoint{1.825039in}{2.722143in}}%
\pgfpathlineto{\pgfqpoint{1.827192in}{2.724557in}}%
\pgfpathlineto{\pgfqpoint{1.829345in}{2.726972in}}%
\pgfpathlineto{\pgfqpoint{1.831499in}{2.729386in}}%
\pgfpathlineto{\pgfqpoint{1.833652in}{2.731800in}}%
\pgfpathlineto{\pgfqpoint{1.835805in}{2.734215in}}%
\pgfpathlineto{\pgfqpoint{1.837958in}{2.736629in}}%
\pgfpathlineto{\pgfqpoint{1.840111in}{2.739043in}}%
\pgfpathlineto{\pgfqpoint{1.842264in}{2.741458in}}%
\pgfpathlineto{\pgfqpoint{1.844417in}{2.743872in}}%
\pgfpathlineto{\pgfqpoint{1.846570in}{2.746286in}}%
\pgfpathlineto{\pgfqpoint{1.848723in}{2.748700in}}%
\pgfpathlineto{\pgfqpoint{1.850876in}{2.751115in}}%
\pgfpathlineto{\pgfqpoint{1.853029in}{2.753529in}}%
\pgfpathlineto{\pgfqpoint{1.855182in}{2.755943in}}%
\pgfpathlineto{\pgfqpoint{1.857335in}{2.758358in}}%
\pgfpathlineto{\pgfqpoint{1.859489in}{2.760772in}}%
\pgfpathlineto{\pgfqpoint{1.861642in}{2.763186in}}%
\pgfpathlineto{\pgfqpoint{1.863795in}{2.765600in}}%
\pgfpathlineto{\pgfqpoint{1.865948in}{2.768015in}}%
\pgfpathlineto{\pgfqpoint{1.868101in}{2.770429in}}%
\pgfpathlineto{\pgfqpoint{1.870254in}{2.772843in}}%
\pgfpathlineto{\pgfqpoint{1.872407in}{2.775258in}}%
\pgfpathlineto{\pgfqpoint{1.874560in}{2.777672in}}%
\pgfpathlineto{\pgfqpoint{1.876713in}{2.780086in}}%
\pgfpathlineto{\pgfqpoint{1.878866in}{2.782500in}}%
\pgfpathlineto{\pgfqpoint{1.881019in}{2.784915in}}%
\pgfpathlineto{\pgfqpoint{1.883172in}{2.787329in}}%
\pgfpathlineto{\pgfqpoint{1.885325in}{2.789743in}}%
\pgfpathlineto{\pgfqpoint{1.887479in}{2.792158in}}%
\pgfpathlineto{\pgfqpoint{1.889632in}{2.794572in}}%
\pgfpathlineto{\pgfqpoint{1.891785in}{2.796986in}}%
\pgfpathlineto{\pgfqpoint{1.893938in}{2.799401in}}%
\pgfpathlineto{\pgfqpoint{1.896091in}{2.801815in}}%
\pgfpathlineto{\pgfqpoint{1.898244in}{2.804229in}}%
\pgfpathlineto{\pgfqpoint{1.900397in}{2.806643in}}%
\pgfpathlineto{\pgfqpoint{1.902550in}{2.809058in}}%
\pgfpathlineto{\pgfqpoint{1.904703in}{2.811472in}}%
\pgfpathlineto{\pgfqpoint{1.906856in}{2.813886in}}%
\pgfpathlineto{\pgfqpoint{1.909009in}{2.816301in}}%
\pgfpathlineto{\pgfqpoint{1.911162in}{2.818715in}}%
\pgfpathlineto{\pgfqpoint{1.913315in}{2.821129in}}%
\pgfpathlineto{\pgfqpoint{1.915469in}{2.823543in}}%
\pgfpathlineto{\pgfqpoint{1.917622in}{2.825958in}}%
\pgfpathlineto{\pgfqpoint{1.919775in}{2.828372in}}%
\pgfpathlineto{\pgfqpoint{1.921928in}{2.830786in}}%
\pgfpathlineto{\pgfqpoint{1.924081in}{2.833201in}}%
\pgfpathlineto{\pgfqpoint{1.926234in}{2.835615in}}%
\pgfpathlineto{\pgfqpoint{1.928387in}{2.838029in}}%
\pgfpathlineto{\pgfqpoint{1.930540in}{2.840444in}}%
\pgfpathlineto{\pgfqpoint{1.932693in}{2.842858in}}%
\pgfpathlineto{\pgfqpoint{1.934846in}{2.845272in}}%
\pgfpathlineto{\pgfqpoint{1.936999in}{2.847686in}}%
\pgfpathlineto{\pgfqpoint{1.939152in}{2.850101in}}%
\pgfpathlineto{\pgfqpoint{1.941305in}{2.852515in}}%
\pgfpathlineto{\pgfqpoint{1.943459in}{2.854929in}}%
\pgfpathlineto{\pgfqpoint{1.945612in}{2.857344in}}%
\pgfpathlineto{\pgfqpoint{1.947765in}{2.859758in}}%
\pgfpathlineto{\pgfqpoint{1.949918in}{2.862172in}}%
\pgfpathlineto{\pgfqpoint{1.952071in}{2.864586in}}%
\pgfpathlineto{\pgfqpoint{1.954224in}{2.867001in}}%
\pgfpathlineto{\pgfqpoint{1.956377in}{2.869415in}}%
\pgfpathlineto{\pgfqpoint{1.958530in}{2.871829in}}%
\pgfpathlineto{\pgfqpoint{1.960683in}{2.874244in}}%
\pgfpathlineto{\pgfqpoint{1.962836in}{2.876658in}}%
\pgfpathlineto{\pgfqpoint{1.964989in}{2.879072in}}%
\pgfpathlineto{\pgfqpoint{1.967142in}{2.881486in}}%
\pgfpathlineto{\pgfqpoint{1.969295in}{2.883901in}}%
\pgfpathlineto{\pgfqpoint{1.971448in}{2.886315in}}%
\pgfpathlineto{\pgfqpoint{1.973602in}{2.888729in}}%
\pgfpathlineto{\pgfqpoint{1.975755in}{2.891144in}}%
\pgfpathlineto{\pgfqpoint{1.977908in}{2.893558in}}%
\pgfpathlineto{\pgfqpoint{1.980061in}{2.895972in}}%
\pgfpathlineto{\pgfqpoint{1.982214in}{2.898387in}}%
\pgfpathlineto{\pgfqpoint{1.984367in}{2.900801in}}%
\pgfpathlineto{\pgfqpoint{1.986520in}{2.903215in}}%
\pgfpathlineto{\pgfqpoint{1.988673in}{2.905629in}}%
\pgfpathlineto{\pgfqpoint{1.990826in}{2.908044in}}%
\pgfpathlineto{\pgfqpoint{1.992979in}{2.910458in}}%
\pgfpathlineto{\pgfqpoint{1.995132in}{2.912872in}}%
\pgfpathlineto{\pgfqpoint{1.997285in}{2.915287in}}%
\pgfpathlineto{\pgfqpoint{1.999438in}{2.917701in}}%
\pgfpathlineto{\pgfqpoint{2.001592in}{2.920115in}}%
\pgfpathlineto{\pgfqpoint{2.003745in}{2.922529in}}%
\pgfpathlineto{\pgfqpoint{2.005898in}{2.924944in}}%
\pgfpathlineto{\pgfqpoint{2.008051in}{2.927358in}}%
\pgfpathlineto{\pgfqpoint{2.010204in}{2.929772in}}%
\pgfpathlineto{\pgfqpoint{2.012357in}{2.932187in}}%
\pgfpathlineto{\pgfqpoint{2.014510in}{2.934601in}}%
\pgfpathlineto{\pgfqpoint{2.016663in}{2.937015in}}%
\pgfpathlineto{\pgfqpoint{2.018816in}{2.939429in}}%
\pgfpathlineto{\pgfqpoint{2.020969in}{2.941844in}}%
\pgfpathlineto{\pgfqpoint{2.023122in}{2.944258in}}%
\pgfpathlineto{\pgfqpoint{2.025275in}{2.946672in}}%
\pgfpathlineto{\pgfqpoint{2.027428in}{2.949087in}}%
\pgfpathlineto{\pgfqpoint{2.029582in}{2.951501in}}%
\pgfpathlineto{\pgfqpoint{2.031735in}{2.953915in}}%
\pgfpathlineto{\pgfqpoint{2.033888in}{2.956330in}}%
\pgfpathlineto{\pgfqpoint{2.036041in}{2.958744in}}%
\pgfpathlineto{\pgfqpoint{2.038194in}{2.961158in}}%
\pgfpathlineto{\pgfqpoint{2.040347in}{2.963572in}}%
\pgfpathlineto{\pgfqpoint{2.042500in}{2.965987in}}%
\pgfpathlineto{\pgfqpoint{2.042500in}{3.216000in}}%
\pgfpathlineto{\pgfqpoint{2.042500in}{3.216000in}}%
\pgfpathlineto{\pgfqpoint{2.040347in}{3.216000in}}%
\pgfpathlineto{\pgfqpoint{2.038194in}{3.216000in}}%
\pgfpathlineto{\pgfqpoint{2.036041in}{3.216000in}}%
\pgfpathlineto{\pgfqpoint{2.033888in}{3.216000in}}%
\pgfpathlineto{\pgfqpoint{2.031735in}{3.216000in}}%
\pgfpathlineto{\pgfqpoint{2.029582in}{3.216000in}}%
\pgfpathlineto{\pgfqpoint{2.027428in}{3.216000in}}%
\pgfpathlineto{\pgfqpoint{2.025275in}{3.216000in}}%
\pgfpathlineto{\pgfqpoint{2.023122in}{3.216000in}}%
\pgfpathlineto{\pgfqpoint{2.020969in}{3.216000in}}%
\pgfpathlineto{\pgfqpoint{2.018816in}{3.216000in}}%
\pgfpathlineto{\pgfqpoint{2.016663in}{3.216000in}}%
\pgfpathlineto{\pgfqpoint{2.014510in}{3.216000in}}%
\pgfpathlineto{\pgfqpoint{2.012357in}{3.216000in}}%
\pgfpathlineto{\pgfqpoint{2.010204in}{3.216000in}}%
\pgfpathlineto{\pgfqpoint{2.008051in}{3.216000in}}%
\pgfpathlineto{\pgfqpoint{2.005898in}{3.216000in}}%
\pgfpathlineto{\pgfqpoint{2.003745in}{3.216000in}}%
\pgfpathlineto{\pgfqpoint{2.001592in}{3.216000in}}%
\pgfpathlineto{\pgfqpoint{1.999438in}{3.216000in}}%
\pgfpathlineto{\pgfqpoint{1.997285in}{3.216000in}}%
\pgfpathlineto{\pgfqpoint{1.995132in}{3.216000in}}%
\pgfpathlineto{\pgfqpoint{1.992979in}{3.216000in}}%
\pgfpathlineto{\pgfqpoint{1.990826in}{3.216000in}}%
\pgfpathlineto{\pgfqpoint{1.988673in}{3.216000in}}%
\pgfpathlineto{\pgfqpoint{1.986520in}{3.216000in}}%
\pgfpathlineto{\pgfqpoint{1.984367in}{3.216000in}}%
\pgfpathlineto{\pgfqpoint{1.982214in}{3.216000in}}%
\pgfpathlineto{\pgfqpoint{1.980061in}{3.216000in}}%
\pgfpathlineto{\pgfqpoint{1.977908in}{3.216000in}}%
\pgfpathlineto{\pgfqpoint{1.975755in}{3.216000in}}%
\pgfpathlineto{\pgfqpoint{1.973602in}{3.216000in}}%
\pgfpathlineto{\pgfqpoint{1.971448in}{3.216000in}}%
\pgfpathlineto{\pgfqpoint{1.969295in}{3.216000in}}%
\pgfpathlineto{\pgfqpoint{1.967142in}{3.216000in}}%
\pgfpathlineto{\pgfqpoint{1.964989in}{3.216000in}}%
\pgfpathlineto{\pgfqpoint{1.962836in}{3.216000in}}%
\pgfpathlineto{\pgfqpoint{1.960683in}{3.216000in}}%
\pgfpathlineto{\pgfqpoint{1.958530in}{3.216000in}}%
\pgfpathlineto{\pgfqpoint{1.956377in}{3.216000in}}%
\pgfpathlineto{\pgfqpoint{1.954224in}{3.216000in}}%
\pgfpathlineto{\pgfqpoint{1.952071in}{3.216000in}}%
\pgfpathlineto{\pgfqpoint{1.949918in}{3.216000in}}%
\pgfpathlineto{\pgfqpoint{1.947765in}{3.216000in}}%
\pgfpathlineto{\pgfqpoint{1.945612in}{3.216000in}}%
\pgfpathlineto{\pgfqpoint{1.943459in}{3.216000in}}%
\pgfpathlineto{\pgfqpoint{1.941305in}{3.216000in}}%
\pgfpathlineto{\pgfqpoint{1.939152in}{3.216000in}}%
\pgfpathlineto{\pgfqpoint{1.936999in}{3.216000in}}%
\pgfpathlineto{\pgfqpoint{1.934846in}{3.216000in}}%
\pgfpathlineto{\pgfqpoint{1.932693in}{3.216000in}}%
\pgfpathlineto{\pgfqpoint{1.930540in}{3.216000in}}%
\pgfpathlineto{\pgfqpoint{1.928387in}{3.216000in}}%
\pgfpathlineto{\pgfqpoint{1.926234in}{3.216000in}}%
\pgfpathlineto{\pgfqpoint{1.924081in}{3.216000in}}%
\pgfpathlineto{\pgfqpoint{1.921928in}{3.216000in}}%
\pgfpathlineto{\pgfqpoint{1.919775in}{3.216000in}}%
\pgfpathlineto{\pgfqpoint{1.917622in}{3.216000in}}%
\pgfpathlineto{\pgfqpoint{1.915469in}{3.216000in}}%
\pgfpathlineto{\pgfqpoint{1.913315in}{3.216000in}}%
\pgfpathlineto{\pgfqpoint{1.911162in}{3.216000in}}%
\pgfpathlineto{\pgfqpoint{1.909009in}{3.216000in}}%
\pgfpathlineto{\pgfqpoint{1.906856in}{3.216000in}}%
\pgfpathlineto{\pgfqpoint{1.904703in}{3.216000in}}%
\pgfpathlineto{\pgfqpoint{1.902550in}{3.216000in}}%
\pgfpathlineto{\pgfqpoint{1.900397in}{3.216000in}}%
\pgfpathlineto{\pgfqpoint{1.898244in}{3.216000in}}%
\pgfpathlineto{\pgfqpoint{1.896091in}{3.216000in}}%
\pgfpathlineto{\pgfqpoint{1.893938in}{3.216000in}}%
\pgfpathlineto{\pgfqpoint{1.891785in}{3.216000in}}%
\pgfpathlineto{\pgfqpoint{1.889632in}{3.216000in}}%
\pgfpathlineto{\pgfqpoint{1.887479in}{3.216000in}}%
\pgfpathlineto{\pgfqpoint{1.885325in}{3.216000in}}%
\pgfpathlineto{\pgfqpoint{1.883172in}{3.216000in}}%
\pgfpathlineto{\pgfqpoint{1.881019in}{3.216000in}}%
\pgfpathlineto{\pgfqpoint{1.878866in}{3.216000in}}%
\pgfpathlineto{\pgfqpoint{1.876713in}{3.216000in}}%
\pgfpathlineto{\pgfqpoint{1.874560in}{3.216000in}}%
\pgfpathlineto{\pgfqpoint{1.872407in}{3.216000in}}%
\pgfpathlineto{\pgfqpoint{1.870254in}{3.216000in}}%
\pgfpathlineto{\pgfqpoint{1.868101in}{3.216000in}}%
\pgfpathlineto{\pgfqpoint{1.865948in}{3.216000in}}%
\pgfpathlineto{\pgfqpoint{1.863795in}{3.216000in}}%
\pgfpathlineto{\pgfqpoint{1.861642in}{3.216000in}}%
\pgfpathlineto{\pgfqpoint{1.859489in}{3.216000in}}%
\pgfpathlineto{\pgfqpoint{1.857335in}{3.216000in}}%
\pgfpathlineto{\pgfqpoint{1.855182in}{3.216000in}}%
\pgfpathlineto{\pgfqpoint{1.853029in}{3.216000in}}%
\pgfpathlineto{\pgfqpoint{1.850876in}{3.216000in}}%
\pgfpathlineto{\pgfqpoint{1.848723in}{3.216000in}}%
\pgfpathlineto{\pgfqpoint{1.846570in}{3.216000in}}%
\pgfpathlineto{\pgfqpoint{1.844417in}{3.216000in}}%
\pgfpathlineto{\pgfqpoint{1.842264in}{3.216000in}}%
\pgfpathlineto{\pgfqpoint{1.840111in}{3.216000in}}%
\pgfpathlineto{\pgfqpoint{1.837958in}{3.216000in}}%
\pgfpathlineto{\pgfqpoint{1.835805in}{3.216000in}}%
\pgfpathlineto{\pgfqpoint{1.833652in}{3.216000in}}%
\pgfpathlineto{\pgfqpoint{1.831499in}{3.216000in}}%
\pgfpathlineto{\pgfqpoint{1.829345in}{3.216000in}}%
\pgfpathlineto{\pgfqpoint{1.827192in}{3.216000in}}%
\pgfpathlineto{\pgfqpoint{1.825039in}{3.216000in}}%
\pgfpathlineto{\pgfqpoint{1.822886in}{3.216000in}}%
\pgfpathlineto{\pgfqpoint{1.820733in}{3.216000in}}%
\pgfpathlineto{\pgfqpoint{1.818580in}{3.216000in}}%
\pgfpathlineto{\pgfqpoint{1.816427in}{3.216000in}}%
\pgfpathlineto{\pgfqpoint{1.814274in}{3.216000in}}%
\pgfpathlineto{\pgfqpoint{1.812121in}{3.216000in}}%
\pgfpathlineto{\pgfqpoint{1.809968in}{3.216000in}}%
\pgfpathlineto{\pgfqpoint{1.807815in}{3.216000in}}%
\pgfpathlineto{\pgfqpoint{1.805662in}{3.216000in}}%
\pgfpathlineto{\pgfqpoint{1.803509in}{3.216000in}}%
\pgfpathlineto{\pgfqpoint{1.801356in}{3.216000in}}%
\pgfpathlineto{\pgfqpoint{1.799202in}{3.216000in}}%
\pgfpathlineto{\pgfqpoint{1.797049in}{3.216000in}}%
\pgfpathlineto{\pgfqpoint{1.794896in}{3.216000in}}%
\pgfpathlineto{\pgfqpoint{1.792743in}{3.216000in}}%
\pgfpathlineto{\pgfqpoint{1.790590in}{3.216000in}}%
\pgfpathlineto{\pgfqpoint{1.788437in}{3.216000in}}%
\pgfpathlineto{\pgfqpoint{1.786284in}{3.216000in}}%
\pgfpathlineto{\pgfqpoint{1.784131in}{3.216000in}}%
\pgfpathlineto{\pgfqpoint{1.781978in}{3.216000in}}%
\pgfpathlineto{\pgfqpoint{1.779825in}{3.216000in}}%
\pgfpathlineto{\pgfqpoint{1.777672in}{3.216000in}}%
\pgfpathlineto{\pgfqpoint{1.775519in}{3.216000in}}%
\pgfpathlineto{\pgfqpoint{1.773366in}{3.216000in}}%
\pgfpathlineto{\pgfqpoint{1.771212in}{3.216000in}}%
\pgfpathlineto{\pgfqpoint{1.769059in}{3.216000in}}%
\pgfpathlineto{\pgfqpoint{1.766906in}{3.216000in}}%
\pgfpathlineto{\pgfqpoint{1.764753in}{3.216000in}}%
\pgfpathlineto{\pgfqpoint{1.762600in}{3.216000in}}%
\pgfpathlineto{\pgfqpoint{1.760447in}{3.216000in}}%
\pgfpathlineto{\pgfqpoint{1.758294in}{3.216000in}}%
\pgfpathlineto{\pgfqpoint{1.756141in}{3.216000in}}%
\pgfpathlineto{\pgfqpoint{1.753988in}{3.216000in}}%
\pgfpathlineto{\pgfqpoint{1.751835in}{3.216000in}}%
\pgfpathlineto{\pgfqpoint{1.749682in}{3.216000in}}%
\pgfpathlineto{\pgfqpoint{1.747529in}{3.216000in}}%
\pgfpathlineto{\pgfqpoint{1.745376in}{3.216000in}}%
\pgfpathlineto{\pgfqpoint{1.743222in}{3.216000in}}%
\pgfpathlineto{\pgfqpoint{1.741069in}{3.216000in}}%
\pgfpathlineto{\pgfqpoint{1.738916in}{3.216000in}}%
\pgfpathlineto{\pgfqpoint{1.736763in}{3.216000in}}%
\pgfpathlineto{\pgfqpoint{1.734610in}{3.216000in}}%
\pgfpathlineto{\pgfqpoint{1.732457in}{3.216000in}}%
\pgfpathlineto{\pgfqpoint{1.730304in}{3.216000in}}%
\pgfpathlineto{\pgfqpoint{1.728151in}{3.216000in}}%
\pgfpathlineto{\pgfqpoint{1.725998in}{3.216000in}}%
\pgfpathlineto{\pgfqpoint{1.723845in}{3.216000in}}%
\pgfpathlineto{\pgfqpoint{1.721692in}{3.216000in}}%
\pgfpathlineto{\pgfqpoint{1.719539in}{3.216000in}}%
\pgfpathlineto{\pgfqpoint{1.717386in}{3.216000in}}%
\pgfpathlineto{\pgfqpoint{1.715232in}{3.216000in}}%
\pgfpathlineto{\pgfqpoint{1.713079in}{3.216000in}}%
\pgfpathlineto{\pgfqpoint{1.710926in}{3.216000in}}%
\pgfpathlineto{\pgfqpoint{1.708773in}{3.216000in}}%
\pgfpathlineto{\pgfqpoint{1.706620in}{3.216000in}}%
\pgfpathlineto{\pgfqpoint{1.704467in}{3.216000in}}%
\pgfpathlineto{\pgfqpoint{1.702314in}{3.216000in}}%
\pgfpathlineto{\pgfqpoint{1.700161in}{3.216000in}}%
\pgfpathlineto{\pgfqpoint{1.698008in}{3.216000in}}%
\pgfpathlineto{\pgfqpoint{1.695855in}{3.216000in}}%
\pgfpathlineto{\pgfqpoint{1.693702in}{3.216000in}}%
\pgfpathlineto{\pgfqpoint{1.691549in}{3.216000in}}%
\pgfpathlineto{\pgfqpoint{1.689396in}{3.216000in}}%
\pgfpathlineto{\pgfqpoint{1.687242in}{3.216000in}}%
\pgfpathlineto{\pgfqpoint{1.685089in}{3.216000in}}%
\pgfpathlineto{\pgfqpoint{1.682936in}{3.216000in}}%
\pgfpathlineto{\pgfqpoint{1.680783in}{3.216000in}}%
\pgfpathlineto{\pgfqpoint{1.678630in}{3.216000in}}%
\pgfpathlineto{\pgfqpoint{1.676477in}{3.216000in}}%
\pgfpathlineto{\pgfqpoint{1.674324in}{3.216000in}}%
\pgfpathlineto{\pgfqpoint{1.672171in}{3.216000in}}%
\pgfpathlineto{\pgfqpoint{1.670018in}{3.216000in}}%
\pgfpathlineto{\pgfqpoint{1.667865in}{3.216000in}}%
\pgfpathlineto{\pgfqpoint{1.665712in}{3.216000in}}%
\pgfpathlineto{\pgfqpoint{1.663559in}{3.216000in}}%
\pgfpathlineto{\pgfqpoint{1.661406in}{3.216000in}}%
\pgfpathlineto{\pgfqpoint{1.659253in}{3.216000in}}%
\pgfpathlineto{\pgfqpoint{1.657099in}{3.216000in}}%
\pgfpathlineto{\pgfqpoint{1.654946in}{3.216000in}}%
\pgfpathlineto{\pgfqpoint{1.652793in}{3.216000in}}%
\pgfpathlineto{\pgfqpoint{1.650640in}{3.216000in}}%
\pgfpathlineto{\pgfqpoint{1.648487in}{3.216000in}}%
\pgfpathlineto{\pgfqpoint{1.646334in}{3.216000in}}%
\pgfpathlineto{\pgfqpoint{1.644181in}{3.216000in}}%
\pgfpathlineto{\pgfqpoint{1.642028in}{3.216000in}}%
\pgfpathlineto{\pgfqpoint{1.639875in}{3.216000in}}%
\pgfpathlineto{\pgfqpoint{1.637722in}{3.216000in}}%
\pgfpathlineto{\pgfqpoint{1.635569in}{3.216000in}}%
\pgfpathlineto{\pgfqpoint{1.633416in}{3.216000in}}%
\pgfpathlineto{\pgfqpoint{1.631263in}{3.216000in}}%
\pgfpathlineto{\pgfqpoint{1.629109in}{3.216000in}}%
\pgfpathlineto{\pgfqpoint{1.626956in}{3.216000in}}%
\pgfpathlineto{\pgfqpoint{1.624803in}{3.216000in}}%
\pgfpathlineto{\pgfqpoint{1.622650in}{3.216000in}}%
\pgfpathlineto{\pgfqpoint{1.620497in}{3.216000in}}%
\pgfpathlineto{\pgfqpoint{1.618344in}{3.216000in}}%
\pgfpathlineto{\pgfqpoint{1.616191in}{3.216000in}}%
\pgfpathlineto{\pgfqpoint{1.614038in}{3.216000in}}%
\pgfpathlineto{\pgfqpoint{1.611885in}{3.216000in}}%
\pgfpathlineto{\pgfqpoint{1.609732in}{3.216000in}}%
\pgfpathlineto{\pgfqpoint{1.607579in}{3.216000in}}%
\pgfpathlineto{\pgfqpoint{1.605426in}{3.216000in}}%
\pgfpathlineto{\pgfqpoint{1.603273in}{3.216000in}}%
\pgfpathlineto{\pgfqpoint{1.601119in}{3.216000in}}%
\pgfpathlineto{\pgfqpoint{1.598966in}{3.216000in}}%
\pgfpathlineto{\pgfqpoint{1.596813in}{3.216000in}}%
\pgfpathlineto{\pgfqpoint{1.594660in}{3.216000in}}%
\pgfpathlineto{\pgfqpoint{1.592507in}{3.216000in}}%
\pgfpathlineto{\pgfqpoint{1.590354in}{3.216000in}}%
\pgfpathlineto{\pgfqpoint{1.588201in}{3.216000in}}%
\pgfpathlineto{\pgfqpoint{1.586048in}{3.216000in}}%
\pgfpathlineto{\pgfqpoint{1.583895in}{3.216000in}}%
\pgfpathlineto{\pgfqpoint{1.581742in}{3.216000in}}%
\pgfpathlineto{\pgfqpoint{1.579589in}{3.216000in}}%
\pgfpathlineto{\pgfqpoint{1.577436in}{3.216000in}}%
\pgfpathlineto{\pgfqpoint{1.575283in}{3.216000in}}%
\pgfpathlineto{\pgfqpoint{1.573129in}{3.216000in}}%
\pgfpathlineto{\pgfqpoint{1.570976in}{3.216000in}}%
\pgfpathlineto{\pgfqpoint{1.568823in}{3.216000in}}%
\pgfpathlineto{\pgfqpoint{1.566670in}{3.216000in}}%
\pgfpathlineto{\pgfqpoint{1.564517in}{3.216000in}}%
\pgfpathlineto{\pgfqpoint{1.562364in}{3.216000in}}%
\pgfpathlineto{\pgfqpoint{1.560211in}{3.216000in}}%
\pgfpathlineto{\pgfqpoint{1.558058in}{3.216000in}}%
\pgfpathlineto{\pgfqpoint{1.555905in}{3.216000in}}%
\pgfpathlineto{\pgfqpoint{1.553752in}{3.216000in}}%
\pgfpathlineto{\pgfqpoint{1.551599in}{3.216000in}}%
\pgfpathlineto{\pgfqpoint{1.549446in}{3.216000in}}%
\pgfpathlineto{\pgfqpoint{1.547293in}{3.216000in}}%
\pgfpathlineto{\pgfqpoint{1.545139in}{3.216000in}}%
\pgfpathlineto{\pgfqpoint{1.542986in}{3.216000in}}%
\pgfpathlineto{\pgfqpoint{1.540833in}{3.216000in}}%
\pgfpathlineto{\pgfqpoint{1.538680in}{3.216000in}}%
\pgfpathlineto{\pgfqpoint{1.536527in}{3.216000in}}%
\pgfpathlineto{\pgfqpoint{1.534374in}{3.216000in}}%
\pgfpathlineto{\pgfqpoint{1.532221in}{3.216000in}}%
\pgfpathlineto{\pgfqpoint{1.530068in}{3.216000in}}%
\pgfpathlineto{\pgfqpoint{1.527915in}{3.216000in}}%
\pgfpathlineto{\pgfqpoint{1.525762in}{3.216000in}}%
\pgfpathlineto{\pgfqpoint{1.523609in}{3.216000in}}%
\pgfpathlineto{\pgfqpoint{1.521456in}{3.216000in}}%
\pgfpathlineto{\pgfqpoint{1.519303in}{3.216000in}}%
\pgfpathlineto{\pgfqpoint{1.517149in}{3.216000in}}%
\pgfpathlineto{\pgfqpoint{1.514996in}{3.216000in}}%
\pgfpathlineto{\pgfqpoint{1.512843in}{3.216000in}}%
\pgfpathlineto{\pgfqpoint{1.510690in}{3.216000in}}%
\pgfpathlineto{\pgfqpoint{1.508537in}{3.216000in}}%
\pgfpathlineto{\pgfqpoint{1.506384in}{3.216000in}}%
\pgfpathlineto{\pgfqpoint{1.504231in}{3.216000in}}%
\pgfpathlineto{\pgfqpoint{1.502078in}{3.216000in}}%
\pgfpathlineto{\pgfqpoint{1.499925in}{3.216000in}}%
\pgfpathlineto{\pgfqpoint{1.497772in}{3.216000in}}%
\pgfpathlineto{\pgfqpoint{1.495619in}{3.216000in}}%
\pgfpathlineto{\pgfqpoint{1.493466in}{3.216000in}}%
\pgfpathlineto{\pgfqpoint{1.491313in}{3.216000in}}%
\pgfpathlineto{\pgfqpoint{1.489160in}{3.216000in}}%
\pgfpathlineto{\pgfqpoint{1.487006in}{3.216000in}}%
\pgfpathlineto{\pgfqpoint{1.484853in}{3.216000in}}%
\pgfpathlineto{\pgfqpoint{1.482700in}{3.216000in}}%
\pgfpathlineto{\pgfqpoint{1.480547in}{3.216000in}}%
\pgfpathlineto{\pgfqpoint{1.478394in}{3.216000in}}%
\pgfpathlineto{\pgfqpoint{1.476241in}{3.216000in}}%
\pgfpathlineto{\pgfqpoint{1.474088in}{3.216000in}}%
\pgfpathlineto{\pgfqpoint{1.471935in}{3.216000in}}%
\pgfpathlineto{\pgfqpoint{1.469782in}{3.216000in}}%
\pgfpathlineto{\pgfqpoint{1.467629in}{3.216000in}}%
\pgfpathlineto{\pgfqpoint{1.465476in}{3.216000in}}%
\pgfpathlineto{\pgfqpoint{1.463323in}{3.216000in}}%
\pgfpathlineto{\pgfqpoint{1.461170in}{3.216000in}}%
\pgfpathlineto{\pgfqpoint{1.459016in}{3.216000in}}%
\pgfpathlineto{\pgfqpoint{1.456863in}{3.216000in}}%
\pgfpathlineto{\pgfqpoint{1.454710in}{3.216000in}}%
\pgfpathlineto{\pgfqpoint{1.452557in}{3.216000in}}%
\pgfpathlineto{\pgfqpoint{1.450404in}{3.216000in}}%
\pgfpathlineto{\pgfqpoint{1.448251in}{3.216000in}}%
\pgfpathlineto{\pgfqpoint{1.446098in}{3.216000in}}%
\pgfpathlineto{\pgfqpoint{1.443945in}{3.216000in}}%
\pgfpathlineto{\pgfqpoint{1.441792in}{3.216000in}}%
\pgfpathlineto{\pgfqpoint{1.439639in}{3.216000in}}%
\pgfpathlineto{\pgfqpoint{1.437486in}{3.216000in}}%
\pgfpathlineto{\pgfqpoint{1.435333in}{3.216000in}}%
\pgfpathlineto{\pgfqpoint{1.433180in}{3.216000in}}%
\pgfpathlineto{\pgfqpoint{1.431026in}{3.216000in}}%
\pgfpathlineto{\pgfqpoint{1.428873in}{3.216000in}}%
\pgfpathlineto{\pgfqpoint{1.426720in}{3.216000in}}%
\pgfpathlineto{\pgfqpoint{1.424567in}{3.216000in}}%
\pgfpathlineto{\pgfqpoint{1.422414in}{3.216000in}}%
\pgfpathlineto{\pgfqpoint{1.420261in}{3.216000in}}%
\pgfpathlineto{\pgfqpoint{1.418108in}{3.216000in}}%
\pgfpathlineto{\pgfqpoint{1.415955in}{3.216000in}}%
\pgfpathlineto{\pgfqpoint{1.413802in}{3.216000in}}%
\pgfpathlineto{\pgfqpoint{1.411649in}{3.216000in}}%
\pgfpathlineto{\pgfqpoint{1.409496in}{3.216000in}}%
\pgfpathlineto{\pgfqpoint{1.407343in}{3.216000in}}%
\pgfpathlineto{\pgfqpoint{1.405190in}{3.216000in}}%
\pgfpathlineto{\pgfqpoint{1.403036in}{3.216000in}}%
\pgfpathlineto{\pgfqpoint{1.400883in}{3.216000in}}%
\pgfpathlineto{\pgfqpoint{1.398730in}{3.216000in}}%
\pgfpathlineto{\pgfqpoint{1.396577in}{3.216000in}}%
\pgfpathlineto{\pgfqpoint{1.394424in}{3.216000in}}%
\pgfpathlineto{\pgfqpoint{1.392271in}{3.216000in}}%
\pgfpathlineto{\pgfqpoint{1.390118in}{3.216000in}}%
\pgfpathlineto{\pgfqpoint{1.387965in}{3.216000in}}%
\pgfpathlineto{\pgfqpoint{1.385812in}{3.216000in}}%
\pgfpathlineto{\pgfqpoint{1.383659in}{3.216000in}}%
\pgfpathlineto{\pgfqpoint{1.381506in}{3.216000in}}%
\pgfpathlineto{\pgfqpoint{1.379353in}{3.216000in}}%
\pgfpathlineto{\pgfqpoint{1.377200in}{3.216000in}}%
\pgfpathlineto{\pgfqpoint{1.375046in}{3.216000in}}%
\pgfpathlineto{\pgfqpoint{1.372893in}{3.216000in}}%
\pgfpathlineto{\pgfqpoint{1.370740in}{3.216000in}}%
\pgfpathlineto{\pgfqpoint{1.368587in}{3.216000in}}%
\pgfpathlineto{\pgfqpoint{1.366434in}{3.216000in}}%
\pgfpathlineto{\pgfqpoint{1.364281in}{3.216000in}}%
\pgfpathlineto{\pgfqpoint{1.362128in}{3.216000in}}%
\pgfpathlineto{\pgfqpoint{1.359975in}{3.216000in}}%
\pgfpathlineto{\pgfqpoint{1.357822in}{3.216000in}}%
\pgfpathlineto{\pgfqpoint{1.355669in}{3.216000in}}%
\pgfpathlineto{\pgfqpoint{1.353516in}{3.216000in}}%
\pgfpathlineto{\pgfqpoint{1.351363in}{3.216000in}}%
\pgfpathlineto{\pgfqpoint{1.349210in}{3.216000in}}%
\pgfpathlineto{\pgfqpoint{1.347057in}{3.216000in}}%
\pgfpathlineto{\pgfqpoint{1.344903in}{3.216000in}}%
\pgfpathlineto{\pgfqpoint{1.342750in}{3.216000in}}%
\pgfpathlineto{\pgfqpoint{1.340597in}{3.216000in}}%
\pgfpathlineto{\pgfqpoint{1.338444in}{3.216000in}}%
\pgfpathlineto{\pgfqpoint{1.336291in}{3.216000in}}%
\pgfpathlineto{\pgfqpoint{1.334138in}{3.216000in}}%
\pgfpathlineto{\pgfqpoint{1.331985in}{3.216000in}}%
\pgfpathlineto{\pgfqpoint{1.329832in}{3.216000in}}%
\pgfpathlineto{\pgfqpoint{1.327679in}{3.216000in}}%
\pgfpathlineto{\pgfqpoint{1.325526in}{3.216000in}}%
\pgfpathlineto{\pgfqpoint{1.323373in}{3.216000in}}%
\pgfpathlineto{\pgfqpoint{1.321220in}{3.216000in}}%
\pgfpathlineto{\pgfqpoint{1.319067in}{3.216000in}}%
\pgfpathlineto{\pgfqpoint{1.316913in}{3.216000in}}%
\pgfpathlineto{\pgfqpoint{1.314760in}{3.216000in}}%
\pgfpathlineto{\pgfqpoint{1.312607in}{3.216000in}}%
\pgfpathlineto{\pgfqpoint{1.310454in}{3.216000in}}%
\pgfpathlineto{\pgfqpoint{1.308301in}{3.216000in}}%
\pgfpathlineto{\pgfqpoint{1.306148in}{3.216000in}}%
\pgfpathlineto{\pgfqpoint{1.303995in}{3.216000in}}%
\pgfpathlineto{\pgfqpoint{1.301842in}{3.216000in}}%
\pgfpathlineto{\pgfqpoint{1.299689in}{3.216000in}}%
\pgfpathlineto{\pgfqpoint{1.297536in}{3.216000in}}%
\pgfpathlineto{\pgfqpoint{1.295383in}{3.216000in}}%
\pgfpathlineto{\pgfqpoint{1.293230in}{3.216000in}}%
\pgfpathlineto{\pgfqpoint{1.291077in}{3.216000in}}%
\pgfpathlineto{\pgfqpoint{1.288923in}{3.216000in}}%
\pgfpathlineto{\pgfqpoint{1.286770in}{3.216000in}}%
\pgfpathlineto{\pgfqpoint{1.284617in}{3.216000in}}%
\pgfpathlineto{\pgfqpoint{1.282464in}{3.216000in}}%
\pgfpathlineto{\pgfqpoint{1.280311in}{3.216000in}}%
\pgfpathlineto{\pgfqpoint{1.278158in}{3.216000in}}%
\pgfpathlineto{\pgfqpoint{1.276005in}{3.216000in}}%
\pgfpathlineto{\pgfqpoint{1.273852in}{3.216000in}}%
\pgfpathlineto{\pgfqpoint{1.271699in}{3.216000in}}%
\pgfpathlineto{\pgfqpoint{1.269546in}{3.216000in}}%
\pgfpathlineto{\pgfqpoint{1.267393in}{3.216000in}}%
\pgfpathlineto{\pgfqpoint{1.265240in}{3.216000in}}%
\pgfpathlineto{\pgfqpoint{1.263087in}{3.216000in}}%
\pgfpathlineto{\pgfqpoint{1.260933in}{3.216000in}}%
\pgfpathlineto{\pgfqpoint{1.258780in}{3.216000in}}%
\pgfpathlineto{\pgfqpoint{1.256627in}{3.216000in}}%
\pgfpathlineto{\pgfqpoint{1.254474in}{3.216000in}}%
\pgfpathlineto{\pgfqpoint{1.252321in}{3.216000in}}%
\pgfpathlineto{\pgfqpoint{1.250168in}{3.216000in}}%
\pgfpathlineto{\pgfqpoint{1.248015in}{3.216000in}}%
\pgfpathlineto{\pgfqpoint{1.245862in}{3.216000in}}%
\pgfpathlineto{\pgfqpoint{1.243709in}{3.216000in}}%
\pgfpathlineto{\pgfqpoint{1.241556in}{3.216000in}}%
\pgfpathlineto{\pgfqpoint{1.239403in}{3.216000in}}%
\pgfpathlineto{\pgfqpoint{1.237250in}{3.216000in}}%
\pgfpathlineto{\pgfqpoint{1.235097in}{3.216000in}}%
\pgfpathlineto{\pgfqpoint{1.232943in}{3.216000in}}%
\pgfpathlineto{\pgfqpoint{1.230790in}{3.216000in}}%
\pgfpathlineto{\pgfqpoint{1.228637in}{3.216000in}}%
\pgfpathlineto{\pgfqpoint{1.226484in}{3.216000in}}%
\pgfpathlineto{\pgfqpoint{1.224331in}{3.216000in}}%
\pgfpathlineto{\pgfqpoint{1.222178in}{3.216000in}}%
\pgfpathlineto{\pgfqpoint{1.220025in}{3.216000in}}%
\pgfpathlineto{\pgfqpoint{1.217872in}{3.216000in}}%
\pgfpathlineto{\pgfqpoint{1.215719in}{3.216000in}}%
\pgfpathlineto{\pgfqpoint{1.213566in}{3.216000in}}%
\pgfpathlineto{\pgfqpoint{1.211413in}{3.216000in}}%
\pgfpathlineto{\pgfqpoint{1.209260in}{3.216000in}}%
\pgfpathlineto{\pgfqpoint{1.207107in}{3.216000in}}%
\pgfpathlineto{\pgfqpoint{1.204954in}{3.216000in}}%
\pgfpathlineto{\pgfqpoint{1.202800in}{3.216000in}}%
\pgfpathlineto{\pgfqpoint{1.200647in}{3.216000in}}%
\pgfpathlineto{\pgfqpoint{1.198494in}{3.216000in}}%
\pgfpathlineto{\pgfqpoint{1.196341in}{3.216000in}}%
\pgfpathlineto{\pgfqpoint{1.194188in}{3.216000in}}%
\pgfpathlineto{\pgfqpoint{1.192035in}{3.216000in}}%
\pgfpathlineto{\pgfqpoint{1.189882in}{3.216000in}}%
\pgfpathlineto{\pgfqpoint{1.187729in}{3.216000in}}%
\pgfpathlineto{\pgfqpoint{1.185576in}{3.216000in}}%
\pgfpathlineto{\pgfqpoint{1.183423in}{3.216000in}}%
\pgfpathlineto{\pgfqpoint{1.181270in}{3.216000in}}%
\pgfpathlineto{\pgfqpoint{1.179117in}{3.216000in}}%
\pgfpathlineto{\pgfqpoint{1.176964in}{3.216000in}}%
\pgfpathlineto{\pgfqpoint{1.174810in}{3.216000in}}%
\pgfpathlineto{\pgfqpoint{1.172657in}{3.216000in}}%
\pgfpathlineto{\pgfqpoint{1.170504in}{3.216000in}}%
\pgfpathlineto{\pgfqpoint{1.168351in}{3.216000in}}%
\pgfpathlineto{\pgfqpoint{1.166198in}{3.216000in}}%
\pgfpathlineto{\pgfqpoint{1.164045in}{3.216000in}}%
\pgfpathlineto{\pgfqpoint{1.161892in}{3.216000in}}%
\pgfpathlineto{\pgfqpoint{1.159739in}{3.216000in}}%
\pgfpathlineto{\pgfqpoint{1.157586in}{3.216000in}}%
\pgfpathlineto{\pgfqpoint{1.155433in}{3.216000in}}%
\pgfpathlineto{\pgfqpoint{1.153280in}{3.216000in}}%
\pgfpathlineto{\pgfqpoint{1.151127in}{3.216000in}}%
\pgfpathlineto{\pgfqpoint{1.148974in}{3.216000in}}%
\pgfpathlineto{\pgfqpoint{1.146820in}{3.216000in}}%
\pgfpathlineto{\pgfqpoint{1.144667in}{3.216000in}}%
\pgfpathlineto{\pgfqpoint{1.142514in}{3.216000in}}%
\pgfpathlineto{\pgfqpoint{1.140361in}{3.216000in}}%
\pgfpathlineto{\pgfqpoint{1.138208in}{3.216000in}}%
\pgfpathlineto{\pgfqpoint{1.136055in}{3.216000in}}%
\pgfpathlineto{\pgfqpoint{1.133902in}{3.216000in}}%
\pgfpathlineto{\pgfqpoint{1.131749in}{3.216000in}}%
\pgfpathlineto{\pgfqpoint{1.129596in}{3.216000in}}%
\pgfpathlineto{\pgfqpoint{1.127443in}{3.216000in}}%
\pgfpathlineto{\pgfqpoint{1.125290in}{3.216000in}}%
\pgfpathlineto{\pgfqpoint{1.123137in}{3.216000in}}%
\pgfpathlineto{\pgfqpoint{1.120984in}{3.216000in}}%
\pgfpathlineto{\pgfqpoint{1.118830in}{3.216000in}}%
\pgfpathlineto{\pgfqpoint{1.116677in}{3.216000in}}%
\pgfpathlineto{\pgfqpoint{1.114524in}{3.216000in}}%
\pgfpathlineto{\pgfqpoint{1.112371in}{3.216000in}}%
\pgfpathlineto{\pgfqpoint{1.110218in}{3.216000in}}%
\pgfpathlineto{\pgfqpoint{1.108065in}{3.216000in}}%
\pgfpathlineto{\pgfqpoint{1.105912in}{3.216000in}}%
\pgfpathlineto{\pgfqpoint{1.103759in}{3.216000in}}%
\pgfpathlineto{\pgfqpoint{1.101606in}{3.216000in}}%
\pgfpathlineto{\pgfqpoint{1.099453in}{3.216000in}}%
\pgfpathlineto{\pgfqpoint{1.097300in}{3.216000in}}%
\pgfpathlineto{\pgfqpoint{1.095147in}{3.216000in}}%
\pgfpathlineto{\pgfqpoint{1.092994in}{3.216000in}}%
\pgfpathlineto{\pgfqpoint{1.090840in}{3.216000in}}%
\pgfpathlineto{\pgfqpoint{1.088687in}{3.216000in}}%
\pgfpathlineto{\pgfqpoint{1.086534in}{3.216000in}}%
\pgfpathlineto{\pgfqpoint{1.084381in}{3.216000in}}%
\pgfpathlineto{\pgfqpoint{1.082228in}{3.216000in}}%
\pgfpathlineto{\pgfqpoint{1.080075in}{3.216000in}}%
\pgfpathlineto{\pgfqpoint{1.077922in}{3.216000in}}%
\pgfpathlineto{\pgfqpoint{1.075769in}{3.216000in}}%
\pgfpathlineto{\pgfqpoint{1.073616in}{3.216000in}}%
\pgfpathlineto{\pgfqpoint{1.071463in}{3.216000in}}%
\pgfpathlineto{\pgfqpoint{1.069310in}{3.216000in}}%
\pgfpathlineto{\pgfqpoint{1.067157in}{3.216000in}}%
\pgfpathlineto{\pgfqpoint{1.065004in}{3.216000in}}%
\pgfpathlineto{\pgfqpoint{1.062851in}{3.216000in}}%
\pgfpathlineto{\pgfqpoint{1.060697in}{3.216000in}}%
\pgfpathlineto{\pgfqpoint{1.058544in}{3.216000in}}%
\pgfpathlineto{\pgfqpoint{1.056391in}{3.216000in}}%
\pgfpathlineto{\pgfqpoint{1.054238in}{3.216000in}}%
\pgfpathlineto{\pgfqpoint{1.052085in}{3.216000in}}%
\pgfpathlineto{\pgfqpoint{1.049932in}{3.216000in}}%
\pgfpathlineto{\pgfqpoint{1.047779in}{3.216000in}}%
\pgfpathlineto{\pgfqpoint{1.045626in}{3.216000in}}%
\pgfpathlineto{\pgfqpoint{1.043473in}{3.216000in}}%
\pgfpathlineto{\pgfqpoint{1.041320in}{3.216000in}}%
\pgfpathlineto{\pgfqpoint{1.039167in}{3.216000in}}%
\pgfpathlineto{\pgfqpoint{1.037014in}{3.216000in}}%
\pgfpathlineto{\pgfqpoint{1.034861in}{3.216000in}}%
\pgfpathlineto{\pgfqpoint{1.032707in}{3.216000in}}%
\pgfpathlineto{\pgfqpoint{1.030554in}{3.216000in}}%
\pgfpathlineto{\pgfqpoint{1.028401in}{3.216000in}}%
\pgfpathlineto{\pgfqpoint{1.026248in}{3.216000in}}%
\pgfpathlineto{\pgfqpoint{1.024095in}{3.216000in}}%
\pgfpathlineto{\pgfqpoint{1.021942in}{3.216000in}}%
\pgfpathlineto{\pgfqpoint{1.019789in}{3.216000in}}%
\pgfpathlineto{\pgfqpoint{1.017636in}{3.216000in}}%
\pgfpathlineto{\pgfqpoint{1.015483in}{3.216000in}}%
\pgfpathlineto{\pgfqpoint{1.013330in}{3.216000in}}%
\pgfpathlineto{\pgfqpoint{1.011177in}{3.216000in}}%
\pgfpathlineto{\pgfqpoint{1.009024in}{3.216000in}}%
\pgfpathlineto{\pgfqpoint{1.006871in}{3.216000in}}%
\pgfpathlineto{\pgfqpoint{1.004717in}{3.216000in}}%
\pgfpathlineto{\pgfqpoint{1.002564in}{3.216000in}}%
\pgfpathlineto{\pgfqpoint{1.000411in}{3.216000in}}%
\pgfpathlineto{\pgfqpoint{0.998258in}{3.216000in}}%
\pgfpathlineto{\pgfqpoint{0.996105in}{3.216000in}}%
\pgfpathlineto{\pgfqpoint{0.993952in}{3.216000in}}%
\pgfpathlineto{\pgfqpoint{0.991799in}{3.216000in}}%
\pgfpathlineto{\pgfqpoint{0.989646in}{3.216000in}}%
\pgfpathlineto{\pgfqpoint{0.987493in}{3.216000in}}%
\pgfpathlineto{\pgfqpoint{0.985340in}{3.216000in}}%
\pgfpathlineto{\pgfqpoint{0.983187in}{3.216000in}}%
\pgfpathlineto{\pgfqpoint{0.981034in}{3.216000in}}%
\pgfpathlineto{\pgfqpoint{0.978881in}{3.216000in}}%
\pgfpathlineto{\pgfqpoint{0.976727in}{3.216000in}}%
\pgfpathlineto{\pgfqpoint{0.974574in}{3.216000in}}%
\pgfpathlineto{\pgfqpoint{0.972421in}{3.216000in}}%
\pgfpathlineto{\pgfqpoint{0.970268in}{3.216000in}}%
\pgfpathlineto{\pgfqpoint{0.968115in}{3.216000in}}%
\pgfpathlineto{\pgfqpoint{0.965962in}{3.216000in}}%
\pgfpathlineto{\pgfqpoint{0.963809in}{3.216000in}}%
\pgfpathlineto{\pgfqpoint{0.961656in}{3.216000in}}%
\pgfpathlineto{\pgfqpoint{0.959503in}{3.216000in}}%
\pgfpathlineto{\pgfqpoint{0.957350in}{3.216000in}}%
\pgfpathlineto{\pgfqpoint{0.955197in}{3.216000in}}%
\pgfpathlineto{\pgfqpoint{0.953044in}{3.216000in}}%
\pgfpathlineto{\pgfqpoint{0.950891in}{3.216000in}}%
\pgfpathlineto{\pgfqpoint{0.948737in}{3.216000in}}%
\pgfpathlineto{\pgfqpoint{0.946584in}{3.216000in}}%
\pgfpathlineto{\pgfqpoint{0.944431in}{3.216000in}}%
\pgfpathlineto{\pgfqpoint{0.942278in}{3.216000in}}%
\pgfpathlineto{\pgfqpoint{0.940125in}{3.216000in}}%
\pgfpathlineto{\pgfqpoint{0.937972in}{3.216000in}}%
\pgfpathlineto{\pgfqpoint{0.935819in}{3.216000in}}%
\pgfpathlineto{\pgfqpoint{0.933666in}{3.216000in}}%
\pgfpathlineto{\pgfqpoint{0.931513in}{3.216000in}}%
\pgfpathlineto{\pgfqpoint{0.929360in}{3.216000in}}%
\pgfpathlineto{\pgfqpoint{0.927207in}{3.216000in}}%
\pgfpathlineto{\pgfqpoint{0.925054in}{3.216000in}}%
\pgfpathlineto{\pgfqpoint{0.922901in}{3.216000in}}%
\pgfpathlineto{\pgfqpoint{0.920747in}{3.216000in}}%
\pgfpathlineto{\pgfqpoint{0.918594in}{3.216000in}}%
\pgfpathlineto{\pgfqpoint{0.916441in}{3.216000in}}%
\pgfpathlineto{\pgfqpoint{0.914288in}{3.216000in}}%
\pgfpathlineto{\pgfqpoint{0.912135in}{3.216000in}}%
\pgfpathlineto{\pgfqpoint{0.909982in}{3.216000in}}%
\pgfpathlineto{\pgfqpoint{0.907829in}{3.216000in}}%
\pgfpathlineto{\pgfqpoint{0.905676in}{3.216000in}}%
\pgfpathlineto{\pgfqpoint{0.903523in}{3.216000in}}%
\pgfpathlineto{\pgfqpoint{0.901370in}{3.216000in}}%
\pgfpathlineto{\pgfqpoint{0.899217in}{3.216000in}}%
\pgfpathlineto{\pgfqpoint{0.897064in}{3.216000in}}%
\pgfpathlineto{\pgfqpoint{0.894911in}{3.216000in}}%
\pgfpathlineto{\pgfqpoint{0.892758in}{3.216000in}}%
\pgfpathlineto{\pgfqpoint{0.890604in}{3.216000in}}%
\pgfpathlineto{\pgfqpoint{0.888451in}{3.216000in}}%
\pgfpathlineto{\pgfqpoint{0.886298in}{3.216000in}}%
\pgfpathlineto{\pgfqpoint{0.884145in}{3.216000in}}%
\pgfpathlineto{\pgfqpoint{0.881992in}{3.216000in}}%
\pgfpathlineto{\pgfqpoint{0.879839in}{3.216000in}}%
\pgfpathlineto{\pgfqpoint{0.877686in}{3.216000in}}%
\pgfpathlineto{\pgfqpoint{0.875533in}{3.216000in}}%
\pgfpathlineto{\pgfqpoint{0.873380in}{3.216000in}}%
\pgfpathlineto{\pgfqpoint{0.871227in}{3.216000in}}%
\pgfpathlineto{\pgfqpoint{0.869074in}{3.216000in}}%
\pgfpathlineto{\pgfqpoint{0.866921in}{3.216000in}}%
\pgfpathlineto{\pgfqpoint{0.864768in}{3.216000in}}%
\pgfpathlineto{\pgfqpoint{0.862614in}{3.216000in}}%
\pgfpathlineto{\pgfqpoint{0.860461in}{3.216000in}}%
\pgfpathlineto{\pgfqpoint{0.858308in}{3.216000in}}%
\pgfpathlineto{\pgfqpoint{0.856155in}{3.216000in}}%
\pgfpathlineto{\pgfqpoint{0.854002in}{3.216000in}}%
\pgfpathlineto{\pgfqpoint{0.851849in}{3.216000in}}%
\pgfpathlineto{\pgfqpoint{0.849696in}{3.216000in}}%
\pgfpathlineto{\pgfqpoint{0.847543in}{3.216000in}}%
\pgfpathlineto{\pgfqpoint{0.845390in}{3.216000in}}%
\pgfpathlineto{\pgfqpoint{0.843237in}{3.216000in}}%
\pgfpathlineto{\pgfqpoint{0.841084in}{3.216000in}}%
\pgfpathlineto{\pgfqpoint{0.838931in}{3.216000in}}%
\pgfpathlineto{\pgfqpoint{0.836778in}{3.216000in}}%
\pgfpathlineto{\pgfqpoint{0.834624in}{3.216000in}}%
\pgfpathlineto{\pgfqpoint{0.832471in}{3.216000in}}%
\pgfpathlineto{\pgfqpoint{0.830318in}{3.216000in}}%
\pgfpathlineto{\pgfqpoint{0.828165in}{3.216000in}}%
\pgfpathlineto{\pgfqpoint{0.826012in}{3.216000in}}%
\pgfpathlineto{\pgfqpoint{0.823859in}{3.216000in}}%
\pgfpathlineto{\pgfqpoint{0.821706in}{3.216000in}}%
\pgfpathlineto{\pgfqpoint{0.819553in}{3.216000in}}%
\pgfpathlineto{\pgfqpoint{0.817400in}{3.216000in}}%
\pgfpathlineto{\pgfqpoint{0.815247in}{3.216000in}}%
\pgfpathlineto{\pgfqpoint{0.813094in}{3.216000in}}%
\pgfpathlineto{\pgfqpoint{0.810941in}{3.216000in}}%
\pgfpathlineto{\pgfqpoint{0.808788in}{3.216000in}}%
\pgfpathlineto{\pgfqpoint{0.806634in}{3.216000in}}%
\pgfpathlineto{\pgfqpoint{0.804481in}{3.216000in}}%
\pgfpathlineto{\pgfqpoint{0.802328in}{3.216000in}}%
\pgfpathlineto{\pgfqpoint{0.800175in}{3.216000in}}%
\pgfpathlineto{\pgfqpoint{0.798022in}{3.216000in}}%
\pgfpathlineto{\pgfqpoint{0.795869in}{3.216000in}}%
\pgfpathlineto{\pgfqpoint{0.793716in}{3.216000in}}%
\pgfpathlineto{\pgfqpoint{0.791563in}{3.216000in}}%
\pgfpathlineto{\pgfqpoint{0.789410in}{3.216000in}}%
\pgfpathlineto{\pgfqpoint{0.787257in}{3.216000in}}%
\pgfpathlineto{\pgfqpoint{0.785104in}{3.216000in}}%
\pgfpathlineto{\pgfqpoint{0.782951in}{3.216000in}}%
\pgfpathlineto{\pgfqpoint{0.780798in}{3.216000in}}%
\pgfpathlineto{\pgfqpoint{0.778644in}{3.216000in}}%
\pgfpathlineto{\pgfqpoint{0.776491in}{3.216000in}}%
\pgfpathlineto{\pgfqpoint{0.774338in}{3.216000in}}%
\pgfpathlineto{\pgfqpoint{0.772185in}{3.216000in}}%
\pgfpathlineto{\pgfqpoint{0.770032in}{3.216000in}}%
\pgfpathlineto{\pgfqpoint{0.767879in}{3.216000in}}%
\pgfpathlineto{\pgfqpoint{0.765726in}{3.216000in}}%
\pgfpathlineto{\pgfqpoint{0.763573in}{3.216000in}}%
\pgfpathlineto{\pgfqpoint{0.761420in}{3.216000in}}%
\pgfpathlineto{\pgfqpoint{0.759267in}{3.216000in}}%
\pgfpathlineto{\pgfqpoint{0.757114in}{3.216000in}}%
\pgfpathlineto{\pgfqpoint{0.754961in}{3.216000in}}%
\pgfpathlineto{\pgfqpoint{0.752808in}{3.216000in}}%
\pgfpathlineto{\pgfqpoint{0.750655in}{3.216000in}}%
\pgfpathlineto{\pgfqpoint{0.748501in}{3.216000in}}%
\pgfpathlineto{\pgfqpoint{0.746348in}{3.216000in}}%
\pgfpathlineto{\pgfqpoint{0.744195in}{3.216000in}}%
\pgfpathlineto{\pgfqpoint{0.742042in}{3.216000in}}%
\pgfpathlineto{\pgfqpoint{0.739889in}{3.216000in}}%
\pgfpathlineto{\pgfqpoint{0.737736in}{3.216000in}}%
\pgfpathlineto{\pgfqpoint{0.735583in}{3.216000in}}%
\pgfpathlineto{\pgfqpoint{0.733430in}{3.216000in}}%
\pgfpathlineto{\pgfqpoint{0.731277in}{3.216000in}}%
\pgfpathlineto{\pgfqpoint{0.729124in}{3.216000in}}%
\pgfpathlineto{\pgfqpoint{0.726971in}{3.216000in}}%
\pgfpathlineto{\pgfqpoint{0.724818in}{3.216000in}}%
\pgfpathlineto{\pgfqpoint{0.722665in}{3.216000in}}%
\pgfpathlineto{\pgfqpoint{0.720511in}{3.216000in}}%
\pgfpathlineto{\pgfqpoint{0.718358in}{3.216000in}}%
\pgfpathlineto{\pgfqpoint{0.716205in}{3.216000in}}%
\pgfpathlineto{\pgfqpoint{0.714052in}{3.216000in}}%
\pgfpathlineto{\pgfqpoint{0.711899in}{3.216000in}}%
\pgfpathlineto{\pgfqpoint{0.709746in}{3.216000in}}%
\pgfpathlineto{\pgfqpoint{0.707593in}{3.216000in}}%
\pgfpathlineto{\pgfqpoint{0.705440in}{3.216000in}}%
\pgfpathlineto{\pgfqpoint{0.703287in}{3.216000in}}%
\pgfpathlineto{\pgfqpoint{0.701134in}{3.216000in}}%
\pgfpathlineto{\pgfqpoint{0.698981in}{3.216000in}}%
\pgfpathlineto{\pgfqpoint{0.696828in}{3.216000in}}%
\pgfpathlineto{\pgfqpoint{0.694675in}{3.216000in}}%
\pgfpathlineto{\pgfqpoint{0.692521in}{3.216000in}}%
\pgfpathlineto{\pgfqpoint{0.690368in}{3.216000in}}%
\pgfpathlineto{\pgfqpoint{0.688215in}{3.216000in}}%
\pgfpathlineto{\pgfqpoint{0.686062in}{3.216000in}}%
\pgfpathlineto{\pgfqpoint{0.683909in}{3.216000in}}%
\pgfpathlineto{\pgfqpoint{0.681756in}{3.216000in}}%
\pgfpathlineto{\pgfqpoint{0.679603in}{3.216000in}}%
\pgfpathlineto{\pgfqpoint{0.677450in}{3.216000in}}%
\pgfpathlineto{\pgfqpoint{0.675297in}{3.216000in}}%
\pgfpathlineto{\pgfqpoint{0.673144in}{3.216000in}}%
\pgfpathlineto{\pgfqpoint{0.670991in}{3.216000in}}%
\pgfpathlineto{\pgfqpoint{0.668838in}{3.216000in}}%
\pgfpathlineto{\pgfqpoint{0.666685in}{3.216000in}}%
\pgfpathlineto{\pgfqpoint{0.664531in}{3.216000in}}%
\pgfpathlineto{\pgfqpoint{0.662378in}{3.216000in}}%
\pgfpathlineto{\pgfqpoint{0.660225in}{3.216000in}}%
\pgfpathlineto{\pgfqpoint{0.658072in}{3.216000in}}%
\pgfpathlineto{\pgfqpoint{0.655919in}{3.216000in}}%
\pgfpathlineto{\pgfqpoint{0.653766in}{3.216000in}}%
\pgfpathlineto{\pgfqpoint{0.651613in}{3.216000in}}%
\pgfpathlineto{\pgfqpoint{0.649460in}{3.216000in}}%
\pgfpathlineto{\pgfqpoint{0.647307in}{3.216000in}}%
\pgfpathlineto{\pgfqpoint{0.645154in}{3.216000in}}%
\pgfpathlineto{\pgfqpoint{0.643001in}{3.216000in}}%
\pgfpathlineto{\pgfqpoint{0.640848in}{3.216000in}}%
\pgfpathlineto{\pgfqpoint{0.638695in}{3.216000in}}%
\pgfpathlineto{\pgfqpoint{0.636541in}{3.216000in}}%
\pgfpathlineto{\pgfqpoint{0.634388in}{3.216000in}}%
\pgfpathlineto{\pgfqpoint{0.632235in}{3.216000in}}%
\pgfpathlineto{\pgfqpoint{0.630082in}{3.216000in}}%
\pgfpathlineto{\pgfqpoint{0.627929in}{3.216000in}}%
\pgfpathlineto{\pgfqpoint{0.625776in}{3.216000in}}%
\pgfpathlineto{\pgfqpoint{0.623623in}{3.216000in}}%
\pgfpathlineto{\pgfqpoint{0.621470in}{3.216000in}}%
\pgfpathlineto{\pgfqpoint{0.619317in}{3.216000in}}%
\pgfpathlineto{\pgfqpoint{0.617164in}{3.216000in}}%
\pgfpathlineto{\pgfqpoint{0.615011in}{3.216000in}}%
\pgfpathlineto{\pgfqpoint{0.612858in}{3.216000in}}%
\pgfpathlineto{\pgfqpoint{0.610705in}{3.216000in}}%
\pgfpathlineto{\pgfqpoint{0.608552in}{3.216000in}}%
\pgfpathlineto{\pgfqpoint{0.606398in}{3.216000in}}%
\pgfpathlineto{\pgfqpoint{0.604245in}{3.216000in}}%
\pgfpathlineto{\pgfqpoint{0.602092in}{3.216000in}}%
\pgfpathlineto{\pgfqpoint{0.599939in}{3.216000in}}%
\pgfpathlineto{\pgfqpoint{0.597786in}{3.216000in}}%
\pgfpathlineto{\pgfqpoint{0.595633in}{3.216000in}}%
\pgfpathlineto{\pgfqpoint{0.593480in}{3.216000in}}%
\pgfpathlineto{\pgfqpoint{0.591327in}{3.216000in}}%
\pgfpathlineto{\pgfqpoint{0.589174in}{3.216000in}}%
\pgfpathlineto{\pgfqpoint{0.587021in}{3.216000in}}%
\pgfpathlineto{\pgfqpoint{0.584868in}{3.216000in}}%
\pgfpathlineto{\pgfqpoint{0.582715in}{3.216000in}}%
\pgfpathlineto{\pgfqpoint{0.580562in}{3.216000in}}%
\pgfpathlineto{\pgfqpoint{0.578408in}{3.216000in}}%
\pgfpathlineto{\pgfqpoint{0.576255in}{3.216000in}}%
\pgfpathlineto{\pgfqpoint{0.574102in}{3.216000in}}%
\pgfpathlineto{\pgfqpoint{0.571949in}{3.216000in}}%
\pgfpathlineto{\pgfqpoint{0.569796in}{3.216000in}}%
\pgfpathlineto{\pgfqpoint{0.567643in}{3.216000in}}%
\pgfpathlineto{\pgfqpoint{0.565490in}{3.216000in}}%
\pgfpathlineto{\pgfqpoint{0.563337in}{3.216000in}}%
\pgfpathlineto{\pgfqpoint{0.561184in}{3.216000in}}%
\pgfpathlineto{\pgfqpoint{0.559031in}{3.216000in}}%
\pgfpathlineto{\pgfqpoint{0.556878in}{3.216000in}}%
\pgfpathlineto{\pgfqpoint{0.554725in}{3.216000in}}%
\pgfpathlineto{\pgfqpoint{0.552572in}{3.216000in}}%
\pgfpathlineto{\pgfqpoint{0.550418in}{3.216000in}}%
\pgfpathlineto{\pgfqpoint{0.548265in}{3.216000in}}%
\pgfpathlineto{\pgfqpoint{0.546112in}{3.216000in}}%
\pgfpathlineto{\pgfqpoint{0.543959in}{3.216000in}}%
\pgfpathlineto{\pgfqpoint{0.541806in}{3.216000in}}%
\pgfpathlineto{\pgfqpoint{0.539653in}{3.216000in}}%
\pgfpathlineto{\pgfqpoint{0.537500in}{3.216000in}}%
\pgfpathlineto{\pgfqpoint{0.537500in}{3.216000in}}%
\pgfpathclose%
\pgfusepath{stroke,fill}%
}%
\begin{pgfscope}%
\pgfsys@transformshift{0.000000in}{0.000000in}%
\pgfsys@useobject{currentmarker}{}%
\end{pgfscope}%
\end{pgfscope}%
\begin{pgfscope}%
\pgfpathrectangle{\pgfqpoint{0.537500in}{2.059688in}}{\pgfqpoint{1.505000in}{1.156312in}}%
\pgfusepath{clip}%
\pgfsetrectcap%
\pgfsetroundjoin%
\pgfsetlinewidth{0.903375pt}%
\definecolor{currentstroke}{rgb}{0.717647,0.717647,0.717647}%
\pgfsetstrokecolor{currentstroke}%
\pgfsetstrokeopacity{0.900000}%
\pgfsetdash{}{0pt}%
\pgfpathmoveto{\pgfqpoint{0.537500in}{2.850917in}}%
\pgfpathlineto{\pgfqpoint{0.686062in}{2.686976in}}%
\pgfpathlineto{\pgfqpoint{0.791563in}{2.571591in}}%
\pgfpathlineto{\pgfqpoint{0.871227in}{2.485526in}}%
\pgfpathlineto{\pgfqpoint{0.931513in}{2.421449in}}%
\pgfpathlineto{\pgfqpoint{0.978881in}{2.372146in}}%
\pgfpathlineto{\pgfqpoint{1.017636in}{2.332872in}}%
\pgfpathlineto{\pgfqpoint{1.047779in}{2.303310in}}%
\pgfpathlineto{\pgfqpoint{1.073616in}{2.278947in}}%
\pgfpathlineto{\pgfqpoint{1.095147in}{2.259587in}}%
\pgfpathlineto{\pgfqpoint{1.114524in}{2.243135in}}%
\pgfpathlineto{\pgfqpoint{1.131749in}{2.229510in}}%
\pgfpathlineto{\pgfqpoint{1.146820in}{2.218559in}}%
\pgfpathlineto{\pgfqpoint{1.159739in}{2.210052in}}%
\pgfpathlineto{\pgfqpoint{1.172657in}{2.202510in}}%
\pgfpathlineto{\pgfqpoint{1.183423in}{2.197079in}}%
\pgfpathlineto{\pgfqpoint{1.194188in}{2.192525in}}%
\pgfpathlineto{\pgfqpoint{1.202800in}{2.189577in}}%
\pgfpathlineto{\pgfqpoint{1.211413in}{2.187293in}}%
\pgfpathlineto{\pgfqpoint{1.220025in}{2.185710in}}%
\pgfpathlineto{\pgfqpoint{1.228637in}{2.184854in}}%
\pgfpathlineto{\pgfqpoint{1.237250in}{2.184740in}}%
\pgfpathlineto{\pgfqpoint{1.245862in}{2.185370in}}%
\pgfpathlineto{\pgfqpoint{1.254474in}{2.186734in}}%
\pgfpathlineto{\pgfqpoint{1.263087in}{2.188807in}}%
\pgfpathlineto{\pgfqpoint{1.271699in}{2.191556in}}%
\pgfpathlineto{\pgfqpoint{1.282464in}{2.195881in}}%
\pgfpathlineto{\pgfqpoint{1.293230in}{2.201107in}}%
\pgfpathlineto{\pgfqpoint{1.303995in}{2.207138in}}%
\pgfpathlineto{\pgfqpoint{1.316913in}{2.215305in}}%
\pgfpathlineto{\pgfqpoint{1.329832in}{2.224344in}}%
\pgfpathlineto{\pgfqpoint{1.344903in}{2.235812in}}%
\pgfpathlineto{\pgfqpoint{1.362128in}{2.249912in}}%
\pgfpathlineto{\pgfqpoint{1.381506in}{2.266783in}}%
\pgfpathlineto{\pgfqpoint{1.403036in}{2.286500in}}%
\pgfpathlineto{\pgfqpoint{1.428873in}{2.311182in}}%
\pgfpathlineto{\pgfqpoint{1.459016in}{2.341013in}}%
\pgfpathlineto{\pgfqpoint{1.495619in}{2.378304in}}%
\pgfpathlineto{\pgfqpoint{1.540833in}{2.425480in}}%
\pgfpathlineto{\pgfqpoint{1.596813in}{2.485013in}}%
\pgfpathlineto{\pgfqpoint{1.667865in}{2.561713in}}%
\pgfpathlineto{\pgfqpoint{1.760447in}{2.662811in}}%
\pgfpathlineto{\pgfqpoint{1.883172in}{2.797989in}}%
\pgfpathlineto{\pgfqpoint{2.042500in}{2.974567in}}%
\pgfpathlineto{\pgfqpoint{2.042500in}{2.974567in}}%
\pgfusepath{stroke}%
\end{pgfscope}%
\begin{pgfscope}%
\pgfpathrectangle{\pgfqpoint{0.537500in}{2.059688in}}{\pgfqpoint{1.505000in}{1.156312in}}%
\pgfusepath{clip}%
\pgfsetrectcap%
\pgfsetroundjoin%
\pgfsetlinewidth{0.903375pt}%
\definecolor{currentstroke}{rgb}{0.717647,0.717647,0.717647}%
\pgfsetstrokecolor{currentstroke}%
\pgfsetstrokeopacity{0.900000}%
\pgfsetdash{}{0pt}%
\pgfpathmoveto{\pgfqpoint{0.537500in}{2.871029in}}%
\pgfpathlineto{\pgfqpoint{0.630082in}{2.771610in}}%
\pgfpathlineto{\pgfqpoint{0.705440in}{2.691731in}}%
\pgfpathlineto{\pgfqpoint{0.767879in}{2.626563in}}%
\pgfpathlineto{\pgfqpoint{0.821706in}{2.571411in}}%
\pgfpathlineto{\pgfqpoint{0.866921in}{2.526081in}}%
\pgfpathlineto{\pgfqpoint{0.907829in}{2.486112in}}%
\pgfpathlineto{\pgfqpoint{0.942278in}{2.453453in}}%
\pgfpathlineto{\pgfqpoint{0.972421in}{2.425831in}}%
\pgfpathlineto{\pgfqpoint{1.000411in}{2.401178in}}%
\pgfpathlineto{\pgfqpoint{1.024095in}{2.381236in}}%
\pgfpathlineto{\pgfqpoint{1.045626in}{2.363986in}}%
\pgfpathlineto{\pgfqpoint{1.065004in}{2.349304in}}%
\pgfpathlineto{\pgfqpoint{1.084381in}{2.335553in}}%
\pgfpathlineto{\pgfqpoint{1.101606in}{2.324230in}}%
\pgfpathlineto{\pgfqpoint{1.116677in}{2.315110in}}%
\pgfpathlineto{\pgfqpoint{1.131749in}{2.306811in}}%
\pgfpathlineto{\pgfqpoint{1.146820in}{2.299417in}}%
\pgfpathlineto{\pgfqpoint{1.159739in}{2.293866in}}%
\pgfpathlineto{\pgfqpoint{1.172657in}{2.289097in}}%
\pgfpathlineto{\pgfqpoint{1.185576in}{2.285158in}}%
\pgfpathlineto{\pgfqpoint{1.198494in}{2.282096in}}%
\pgfpathlineto{\pgfqpoint{1.209260in}{2.280239in}}%
\pgfpathlineto{\pgfqpoint{1.220025in}{2.279032in}}%
\pgfpathlineto{\pgfqpoint{1.230790in}{2.278485in}}%
\pgfpathlineto{\pgfqpoint{1.241556in}{2.278603in}}%
\pgfpathlineto{\pgfqpoint{1.252321in}{2.279386in}}%
\pgfpathlineto{\pgfqpoint{1.263087in}{2.280825in}}%
\pgfpathlineto{\pgfqpoint{1.273852in}{2.282910in}}%
\pgfpathlineto{\pgfqpoint{1.286770in}{2.286236in}}%
\pgfpathlineto{\pgfqpoint{1.299689in}{2.290425in}}%
\pgfpathlineto{\pgfqpoint{1.312607in}{2.295432in}}%
\pgfpathlineto{\pgfqpoint{1.325526in}{2.301204in}}%
\pgfpathlineto{\pgfqpoint{1.340597in}{2.308836in}}%
\pgfpathlineto{\pgfqpoint{1.355669in}{2.317352in}}%
\pgfpathlineto{\pgfqpoint{1.370740in}{2.326667in}}%
\pgfpathlineto{\pgfqpoint{1.387965in}{2.338189in}}%
\pgfpathlineto{\pgfqpoint{1.407343in}{2.352134in}}%
\pgfpathlineto{\pgfqpoint{1.426720in}{2.366985in}}%
\pgfpathlineto{\pgfqpoint{1.448251in}{2.384395in}}%
\pgfpathlineto{\pgfqpoint{1.471935in}{2.404485in}}%
\pgfpathlineto{\pgfqpoint{1.497772in}{2.427338in}}%
\pgfpathlineto{\pgfqpoint{1.527915in}{2.455016in}}%
\pgfpathlineto{\pgfqpoint{1.562364in}{2.487725in}}%
\pgfpathlineto{\pgfqpoint{1.601119in}{2.525608in}}%
\pgfpathlineto{\pgfqpoint{1.644181in}{2.568745in}}%
\pgfpathlineto{\pgfqpoint{1.695855in}{2.621617in}}%
\pgfpathlineto{\pgfqpoint{1.753988in}{2.682178in}}%
\pgfpathlineto{\pgfqpoint{1.822886in}{2.755037in}}%
\pgfpathlineto{\pgfqpoint{1.904703in}{2.842654in}}%
\pgfpathlineto{\pgfqpoint{2.003745in}{2.949830in}}%
\pgfpathlineto{\pgfqpoint{2.042500in}{2.992015in}}%
\pgfpathlineto{\pgfqpoint{2.042500in}{2.992015in}}%
\pgfusepath{stroke}%
\end{pgfscope}%
\begin{pgfscope}%
\pgfpathrectangle{\pgfqpoint{0.537500in}{2.059688in}}{\pgfqpoint{1.505000in}{1.156312in}}%
\pgfusepath{clip}%
\pgfsetrectcap%
\pgfsetroundjoin%
\pgfsetlinewidth{0.903375pt}%
\definecolor{currentstroke}{rgb}{0.717647,0.717647,0.717647}%
\pgfsetstrokecolor{currentstroke}%
\pgfsetstrokeopacity{0.900000}%
\pgfsetdash{}{0pt}%
\pgfpathmoveto{\pgfqpoint{0.537500in}{2.913265in}}%
\pgfpathlineto{\pgfqpoint{0.602092in}{2.847510in}}%
\pgfpathlineto{\pgfqpoint{0.658072in}{2.791543in}}%
\pgfpathlineto{\pgfqpoint{0.707593in}{2.743019in}}%
\pgfpathlineto{\pgfqpoint{0.752808in}{2.699710in}}%
\pgfpathlineto{\pgfqpoint{0.793716in}{2.661522in}}%
\pgfpathlineto{\pgfqpoint{0.830318in}{2.628317in}}%
\pgfpathlineto{\pgfqpoint{0.864768in}{2.598054in}}%
\pgfpathlineto{\pgfqpoint{0.894911in}{2.572496in}}%
\pgfpathlineto{\pgfqpoint{0.922901in}{2.549659in}}%
\pgfpathlineto{\pgfqpoint{0.948737in}{2.529455in}}%
\pgfpathlineto{\pgfqpoint{0.974574in}{2.510209in}}%
\pgfpathlineto{\pgfqpoint{0.998258in}{2.493517in}}%
\pgfpathlineto{\pgfqpoint{1.019789in}{2.479226in}}%
\pgfpathlineto{\pgfqpoint{1.041320in}{2.465869in}}%
\pgfpathlineto{\pgfqpoint{1.060697in}{2.454724in}}%
\pgfpathlineto{\pgfqpoint{1.080075in}{2.444484in}}%
\pgfpathlineto{\pgfqpoint{1.097300in}{2.436202in}}%
\pgfpathlineto{\pgfqpoint{1.114524in}{2.428745in}}%
\pgfpathlineto{\pgfqpoint{1.131749in}{2.422165in}}%
\pgfpathlineto{\pgfqpoint{1.146820in}{2.417165in}}%
\pgfpathlineto{\pgfqpoint{1.161892in}{2.412904in}}%
\pgfpathlineto{\pgfqpoint{1.176964in}{2.409409in}}%
\pgfpathlineto{\pgfqpoint{1.192035in}{2.406702in}}%
\pgfpathlineto{\pgfqpoint{1.207107in}{2.404802in}}%
\pgfpathlineto{\pgfqpoint{1.222178in}{2.403724in}}%
\pgfpathlineto{\pgfqpoint{1.237250in}{2.403473in}}%
\pgfpathlineto{\pgfqpoint{1.252321in}{2.404053in}}%
\pgfpathlineto{\pgfqpoint{1.267393in}{2.405459in}}%
\pgfpathlineto{\pgfqpoint{1.282464in}{2.407680in}}%
\pgfpathlineto{\pgfqpoint{1.297536in}{2.410703in}}%
\pgfpathlineto{\pgfqpoint{1.312607in}{2.414505in}}%
\pgfpathlineto{\pgfqpoint{1.327679in}{2.419063in}}%
\pgfpathlineto{\pgfqpoint{1.342750in}{2.424347in}}%
\pgfpathlineto{\pgfqpoint{1.359975in}{2.431237in}}%
\pgfpathlineto{\pgfqpoint{1.377200in}{2.438986in}}%
\pgfpathlineto{\pgfqpoint{1.394424in}{2.447543in}}%
\pgfpathlineto{\pgfqpoint{1.413802in}{2.458071in}}%
\pgfpathlineto{\pgfqpoint{1.433180in}{2.469481in}}%
\pgfpathlineto{\pgfqpoint{1.452557in}{2.481701in}}%
\pgfpathlineto{\pgfqpoint{1.474088in}{2.496149in}}%
\pgfpathlineto{\pgfqpoint{1.497772in}{2.512995in}}%
\pgfpathlineto{\pgfqpoint{1.521456in}{2.530736in}}%
\pgfpathlineto{\pgfqpoint{1.547293in}{2.550998in}}%
\pgfpathlineto{\pgfqpoint{1.575283in}{2.573891in}}%
\pgfpathlineto{\pgfqpoint{1.605426in}{2.599500in}}%
\pgfpathlineto{\pgfqpoint{1.637722in}{2.627890in}}%
\pgfpathlineto{\pgfqpoint{1.674324in}{2.661082in}}%
\pgfpathlineto{\pgfqpoint{1.713079in}{2.697223in}}%
\pgfpathlineto{\pgfqpoint{1.756141in}{2.738388in}}%
\pgfpathlineto{\pgfqpoint{1.803509in}{2.784684in}}%
\pgfpathlineto{\pgfqpoint{1.857335in}{2.838351in}}%
\pgfpathlineto{\pgfqpoint{1.917622in}{2.899537in}}%
\pgfpathlineto{\pgfqpoint{1.984367in}{2.968339in}}%
\pgfpathlineto{\pgfqpoint{2.042500in}{3.028994in}}%
\pgfpathlineto{\pgfqpoint{2.042500in}{3.028994in}}%
\pgfusepath{stroke}%
\end{pgfscope}%
\begin{pgfscope}%
\pgfpathrectangle{\pgfqpoint{0.537500in}{2.059688in}}{\pgfqpoint{1.505000in}{1.156312in}}%
\pgfusepath{clip}%
\pgfsetrectcap%
\pgfsetroundjoin%
\pgfsetlinewidth{0.903375pt}%
\definecolor{currentstroke}{rgb}{0.717647,0.717647,0.717647}%
\pgfsetstrokecolor{currentstroke}%
\pgfsetstrokeopacity{0.900000}%
\pgfsetdash{}{0pt}%
\pgfpathmoveto{\pgfqpoint{0.537500in}{2.987289in}}%
\pgfpathlineto{\pgfqpoint{0.587021in}{2.941027in}}%
\pgfpathlineto{\pgfqpoint{0.632235in}{2.899770in}}%
\pgfpathlineto{\pgfqpoint{0.675297in}{2.861484in}}%
\pgfpathlineto{\pgfqpoint{0.714052in}{2.827991in}}%
\pgfpathlineto{\pgfqpoint{0.750655in}{2.797313in}}%
\pgfpathlineto{\pgfqpoint{0.785104in}{2.769398in}}%
\pgfpathlineto{\pgfqpoint{0.817400in}{2.744174in}}%
\pgfpathlineto{\pgfqpoint{0.847543in}{2.721553in}}%
\pgfpathlineto{\pgfqpoint{0.875533in}{2.701428in}}%
\pgfpathlineto{\pgfqpoint{0.903523in}{2.682236in}}%
\pgfpathlineto{\pgfqpoint{0.929360in}{2.665424in}}%
\pgfpathlineto{\pgfqpoint{0.953044in}{2.650841in}}%
\pgfpathlineto{\pgfqpoint{0.976727in}{2.637112in}}%
\pgfpathlineto{\pgfqpoint{0.998258in}{2.625425in}}%
\pgfpathlineto{\pgfqpoint{1.019789in}{2.614542in}}%
\pgfpathlineto{\pgfqpoint{1.041320in}{2.604512in}}%
\pgfpathlineto{\pgfqpoint{1.060697in}{2.596254in}}%
\pgfpathlineto{\pgfqpoint{1.080075in}{2.588761in}}%
\pgfpathlineto{\pgfqpoint{1.099453in}{2.582064in}}%
\pgfpathlineto{\pgfqpoint{1.118830in}{2.576195in}}%
\pgfpathlineto{\pgfqpoint{1.136055in}{2.571697in}}%
\pgfpathlineto{\pgfqpoint{1.153280in}{2.567893in}}%
\pgfpathlineto{\pgfqpoint{1.170504in}{2.564800in}}%
\pgfpathlineto{\pgfqpoint{1.187729in}{2.562430in}}%
\pgfpathlineto{\pgfqpoint{1.204954in}{2.560794in}}%
\pgfpathlineto{\pgfqpoint{1.222178in}{2.559898in}}%
\pgfpathlineto{\pgfqpoint{1.239403in}{2.559748in}}%
\pgfpathlineto{\pgfqpoint{1.256627in}{2.560344in}}%
\pgfpathlineto{\pgfqpoint{1.273852in}{2.561682in}}%
\pgfpathlineto{\pgfqpoint{1.291077in}{2.563757in}}%
\pgfpathlineto{\pgfqpoint{1.308301in}{2.566561in}}%
\pgfpathlineto{\pgfqpoint{1.325526in}{2.570080in}}%
\pgfpathlineto{\pgfqpoint{1.342750in}{2.574301in}}%
\pgfpathlineto{\pgfqpoint{1.362128in}{2.579866in}}%
\pgfpathlineto{\pgfqpoint{1.381506in}{2.586270in}}%
\pgfpathlineto{\pgfqpoint{1.400883in}{2.593482in}}%
\pgfpathlineto{\pgfqpoint{1.420261in}{2.601470in}}%
\pgfpathlineto{\pgfqpoint{1.439639in}{2.610200in}}%
\pgfpathlineto{\pgfqpoint{1.461170in}{2.620729in}}%
\pgfpathlineto{\pgfqpoint{1.482700in}{2.632082in}}%
\pgfpathlineto{\pgfqpoint{1.506384in}{2.645467in}}%
\pgfpathlineto{\pgfqpoint{1.530068in}{2.659730in}}%
\pgfpathlineto{\pgfqpoint{1.555905in}{2.676218in}}%
\pgfpathlineto{\pgfqpoint{1.581742in}{2.693602in}}%
\pgfpathlineto{\pgfqpoint{1.609732in}{2.713362in}}%
\pgfpathlineto{\pgfqpoint{1.637722in}{2.734005in}}%
\pgfpathlineto{\pgfqpoint{1.667865in}{2.757132in}}%
\pgfpathlineto{\pgfqpoint{1.700161in}{2.782845in}}%
\pgfpathlineto{\pgfqpoint{1.734610in}{2.811226in}}%
\pgfpathlineto{\pgfqpoint{1.771212in}{2.842343in}}%
\pgfpathlineto{\pgfqpoint{1.809968in}{2.876247in}}%
\pgfpathlineto{\pgfqpoint{1.853029in}{2.914932in}}%
\pgfpathlineto{\pgfqpoint{1.898244in}{2.956554in}}%
\pgfpathlineto{\pgfqpoint{1.947765in}{3.003160in}}%
\pgfpathlineto{\pgfqpoint{2.001592in}{3.054857in}}%
\pgfpathlineto{\pgfqpoint{2.042500in}{3.094774in}}%
\pgfpathlineto{\pgfqpoint{2.042500in}{3.094774in}}%
\pgfusepath{stroke}%
\end{pgfscope}%
\begin{pgfscope}%
\pgfpathrectangle{\pgfqpoint{0.537500in}{2.059688in}}{\pgfqpoint{1.505000in}{1.156312in}}%
\pgfusepath{clip}%
\pgfsetrectcap%
\pgfsetroundjoin%
\pgfsetlinewidth{0.903375pt}%
\definecolor{currentstroke}{rgb}{0.717647,0.717647,0.717647}%
\pgfsetstrokecolor{currentstroke}%
\pgfsetstrokeopacity{0.900000}%
\pgfsetdash{}{0pt}%
\pgfpathmoveto{\pgfqpoint{0.537500in}{3.100420in}}%
\pgfpathlineto{\pgfqpoint{0.580562in}{3.064678in}}%
\pgfpathlineto{\pgfqpoint{0.619317in}{3.033436in}}%
\pgfpathlineto{\pgfqpoint{0.658072in}{3.003161in}}%
\pgfpathlineto{\pgfqpoint{0.694675in}{2.975543in}}%
\pgfpathlineto{\pgfqpoint{0.729124in}{2.950496in}}%
\pgfpathlineto{\pgfqpoint{0.761420in}{2.927919in}}%
\pgfpathlineto{\pgfqpoint{0.791563in}{2.907701in}}%
\pgfpathlineto{\pgfqpoint{0.821706in}{2.888368in}}%
\pgfpathlineto{\pgfqpoint{0.849696in}{2.871265in}}%
\pgfpathlineto{\pgfqpoint{0.877686in}{2.855033in}}%
\pgfpathlineto{\pgfqpoint{0.903523in}{2.840870in}}%
\pgfpathlineto{\pgfqpoint{0.929360in}{2.827540in}}%
\pgfpathlineto{\pgfqpoint{0.955197in}{2.815086in}}%
\pgfpathlineto{\pgfqpoint{0.978881in}{2.804477in}}%
\pgfpathlineto{\pgfqpoint{1.002564in}{2.794675in}}%
\pgfpathlineto{\pgfqpoint{1.026248in}{2.785712in}}%
\pgfpathlineto{\pgfqpoint{1.047779in}{2.778320in}}%
\pgfpathlineto{\pgfqpoint{1.069310in}{2.771670in}}%
\pgfpathlineto{\pgfqpoint{1.090840in}{2.765783in}}%
\pgfpathlineto{\pgfqpoint{1.112371in}{2.760678in}}%
\pgfpathlineto{\pgfqpoint{1.133902in}{2.756374in}}%
\pgfpathlineto{\pgfqpoint{1.155433in}{2.752883in}}%
\pgfpathlineto{\pgfqpoint{1.176964in}{2.750220in}}%
\pgfpathlineto{\pgfqpoint{1.196341in}{2.748538in}}%
\pgfpathlineto{\pgfqpoint{1.215719in}{2.747539in}}%
\pgfpathlineto{\pgfqpoint{1.235097in}{2.747226in}}%
\pgfpathlineto{\pgfqpoint{1.254474in}{2.747599in}}%
\pgfpathlineto{\pgfqpoint{1.273852in}{2.748657in}}%
\pgfpathlineto{\pgfqpoint{1.293230in}{2.750398in}}%
\pgfpathlineto{\pgfqpoint{1.312607in}{2.752815in}}%
\pgfpathlineto{\pgfqpoint{1.334138in}{2.756287in}}%
\pgfpathlineto{\pgfqpoint{1.355669in}{2.760574in}}%
\pgfpathlineto{\pgfqpoint{1.377200in}{2.765661in}}%
\pgfpathlineto{\pgfqpoint{1.398730in}{2.771531in}}%
\pgfpathlineto{\pgfqpoint{1.420261in}{2.778164in}}%
\pgfpathlineto{\pgfqpoint{1.441792in}{2.785540in}}%
\pgfpathlineto{\pgfqpoint{1.465476in}{2.794485in}}%
\pgfpathlineto{\pgfqpoint{1.489160in}{2.804271in}}%
\pgfpathlineto{\pgfqpoint{1.512843in}{2.814864in}}%
\pgfpathlineto{\pgfqpoint{1.536527in}{2.826230in}}%
\pgfpathlineto{\pgfqpoint{1.562364in}{2.839474in}}%
\pgfpathlineto{\pgfqpoint{1.588201in}{2.853554in}}%
\pgfpathlineto{\pgfqpoint{1.616191in}{2.869701in}}%
\pgfpathlineto{\pgfqpoint{1.644181in}{2.886724in}}%
\pgfpathlineto{\pgfqpoint{1.674324in}{2.905977in}}%
\pgfpathlineto{\pgfqpoint{1.704467in}{2.926120in}}%
\pgfpathlineto{\pgfqpoint{1.736763in}{2.948622in}}%
\pgfpathlineto{\pgfqpoint{1.769059in}{2.972008in}}%
\pgfpathlineto{\pgfqpoint{1.803509in}{2.997851in}}%
\pgfpathlineto{\pgfqpoint{1.840111in}{3.026244in}}%
\pgfpathlineto{\pgfqpoint{1.878866in}{3.057267in}}%
\pgfpathlineto{\pgfqpoint{1.919775in}{3.090989in}}%
\pgfpathlineto{\pgfqpoint{1.962836in}{3.127464in}}%
\pgfpathlineto{\pgfqpoint{2.008051in}{3.166737in}}%
\pgfpathlineto{\pgfqpoint{2.042500in}{3.197266in}}%
\pgfpathlineto{\pgfqpoint{2.042500in}{3.197266in}}%
\pgfusepath{stroke}%
\end{pgfscope}%
\begin{pgfscope}%
\pgfpathrectangle{\pgfqpoint{0.537500in}{2.059688in}}{\pgfqpoint{1.505000in}{1.156312in}}%
\pgfusepath{clip}%
\pgfsetrectcap%
\pgfsetroundjoin%
\pgfsetlinewidth{0.903375pt}%
\definecolor{currentstroke}{rgb}{0.717647,0.717647,0.717647}%
\pgfsetstrokecolor{currentstroke}%
\pgfsetstrokeopacity{0.900000}%
\pgfsetdash{}{0pt}%
\pgfpathmoveto{\pgfqpoint{0.555627in}{3.219333in}}%
\pgfpathlineto{\pgfqpoint{0.593480in}{3.191933in}}%
\pgfpathlineto{\pgfqpoint{0.630082in}{3.166342in}}%
\pgfpathlineto{\pgfqpoint{0.664531in}{3.143125in}}%
\pgfpathlineto{\pgfqpoint{0.698981in}{3.120806in}}%
\pgfpathlineto{\pgfqpoint{0.731277in}{3.100750in}}%
\pgfpathlineto{\pgfqpoint{0.763573in}{3.081584in}}%
\pgfpathlineto{\pgfqpoint{0.793716in}{3.064543in}}%
\pgfpathlineto{\pgfqpoint{0.823859in}{3.048365in}}%
\pgfpathlineto{\pgfqpoint{0.851849in}{3.034153in}}%
\pgfpathlineto{\pgfqpoint{0.879839in}{3.020755in}}%
\pgfpathlineto{\pgfqpoint{0.907829in}{3.008207in}}%
\pgfpathlineto{\pgfqpoint{0.933666in}{2.997408in}}%
\pgfpathlineto{\pgfqpoint{0.959503in}{2.987388in}}%
\pgfpathlineto{\pgfqpoint{0.985340in}{2.978173in}}%
\pgfpathlineto{\pgfqpoint{1.011177in}{2.969787in}}%
\pgfpathlineto{\pgfqpoint{1.034861in}{2.962849in}}%
\pgfpathlineto{\pgfqpoint{1.058544in}{2.956644in}}%
\pgfpathlineto{\pgfqpoint{1.082228in}{2.951186in}}%
\pgfpathlineto{\pgfqpoint{1.105912in}{2.946491in}}%
\pgfpathlineto{\pgfqpoint{1.129596in}{2.942570in}}%
\pgfpathlineto{\pgfqpoint{1.153280in}{2.939434in}}%
\pgfpathlineto{\pgfqpoint{1.176964in}{2.937090in}}%
\pgfpathlineto{\pgfqpoint{1.200647in}{2.935546in}}%
\pgfpathlineto{\pgfqpoint{1.224331in}{2.934806in}}%
\pgfpathlineto{\pgfqpoint{1.248015in}{2.934871in}}%
\pgfpathlineto{\pgfqpoint{1.271699in}{2.935742in}}%
\pgfpathlineto{\pgfqpoint{1.295383in}{2.937415in}}%
\pgfpathlineto{\pgfqpoint{1.319067in}{2.939887in}}%
\pgfpathlineto{\pgfqpoint{1.342750in}{2.943151in}}%
\pgfpathlineto{\pgfqpoint{1.366434in}{2.947198in}}%
\pgfpathlineto{\pgfqpoint{1.390118in}{2.952017in}}%
\pgfpathlineto{\pgfqpoint{1.413802in}{2.957596in}}%
\pgfpathlineto{\pgfqpoint{1.437486in}{2.963921in}}%
\pgfpathlineto{\pgfqpoint{1.461170in}{2.970977in}}%
\pgfpathlineto{\pgfqpoint{1.487006in}{2.979486in}}%
\pgfpathlineto{\pgfqpoint{1.512843in}{2.988822in}}%
\pgfpathlineto{\pgfqpoint{1.538680in}{2.998959in}}%
\pgfpathlineto{\pgfqpoint{1.564517in}{3.009871in}}%
\pgfpathlineto{\pgfqpoint{1.592507in}{3.022538in}}%
\pgfpathlineto{\pgfqpoint{1.620497in}{3.036049in}}%
\pgfpathlineto{\pgfqpoint{1.648487in}{3.050370in}}%
\pgfpathlineto{\pgfqpoint{1.678630in}{3.066660in}}%
\pgfpathlineto{\pgfqpoint{1.708773in}{3.083807in}}%
\pgfpathlineto{\pgfqpoint{1.741069in}{3.103081in}}%
\pgfpathlineto{\pgfqpoint{1.773366in}{3.123240in}}%
\pgfpathlineto{\pgfqpoint{1.807815in}{3.145661in}}%
\pgfpathlineto{\pgfqpoint{1.842264in}{3.168975in}}%
\pgfpathlineto{\pgfqpoint{1.878866in}{3.194661in}}%
\pgfpathlineto{\pgfqpoint{1.912892in}{3.219333in}}%
\pgfpathlineto{\pgfqpoint{1.912892in}{3.219333in}}%
\pgfusepath{stroke}%
\end{pgfscope}%
\begin{pgfscope}%
\pgfpathrectangle{\pgfqpoint{0.537500in}{2.059688in}}{\pgfqpoint{1.505000in}{1.156312in}}%
\pgfusepath{clip}%
\pgfsetrectcap%
\pgfsetroundjoin%
\pgfsetlinewidth{1.003750pt}%
\definecolor{currentstroke}{rgb}{1.000000,1.000000,1.000000}%
\pgfsetstrokecolor{currentstroke}%
\pgfsetdash{}{0pt}%
\pgfpathmoveto{\pgfqpoint{0.676852in}{2.059688in}}%
\pgfpathlineto{\pgfqpoint{0.676852in}{3.216000in}}%
\pgfusepath{stroke}%
\end{pgfscope}%
\begin{pgfscope}%
\pgfsetbuttcap%
\pgfsetroundjoin%
\definecolor{currentfill}{rgb}{0.000000,0.000000,0.000000}%
\pgfsetfillcolor{currentfill}%
\pgfsetlinewidth{0.803000pt}%
\definecolor{currentstroke}{rgb}{0.000000,0.000000,0.000000}%
\pgfsetstrokecolor{currentstroke}%
\pgfsetdash{}{0pt}%
\pgfsys@defobject{currentmarker}{\pgfqpoint{0.000000in}{-0.048611in}}{\pgfqpoint{0.000000in}{0.000000in}}{%
\pgfpathmoveto{\pgfqpoint{0.000000in}{0.000000in}}%
\pgfpathlineto{\pgfqpoint{0.000000in}{-0.048611in}}%
\pgfusepath{stroke,fill}%
}%
\begin{pgfscope}%
\pgfsys@transformshift{0.676852in}{2.059688in}%
\pgfsys@useobject{currentmarker}{}%
\end{pgfscope}%
\end{pgfscope}%
\begin{pgfscope}%
\definecolor{textcolor}{rgb}{0.000000,0.000000,0.000000}%
\pgfsetstrokecolor{textcolor}%
\pgfsetfillcolor{textcolor}%
\pgftext[x=0.676852in,y=1.962466in,,top]{\color{textcolor}{\rmfamily\fontsize{6.000000}{7.200000}\selectfont\catcode`\^=\active\def^{\ifmmode\sp\else\^{}\fi}\catcode`\%=\active\def%{\%}\ensuremath{-}1.0}}%
\end{pgfscope}%
\begin{pgfscope}%
\pgfpathrectangle{\pgfqpoint{0.537500in}{2.059688in}}{\pgfqpoint{1.505000in}{1.156312in}}%
\pgfusepath{clip}%
\pgfsetrectcap%
\pgfsetroundjoin%
\pgfsetlinewidth{1.003750pt}%
\definecolor{currentstroke}{rgb}{1.000000,1.000000,1.000000}%
\pgfsetstrokecolor{currentstroke}%
\pgfsetdash{}{0pt}%
\pgfpathmoveto{\pgfqpoint{0.955556in}{2.059688in}}%
\pgfpathlineto{\pgfqpoint{0.955556in}{3.216000in}}%
\pgfusepath{stroke}%
\end{pgfscope}%
\begin{pgfscope}%
\pgfsetbuttcap%
\pgfsetroundjoin%
\definecolor{currentfill}{rgb}{0.000000,0.000000,0.000000}%
\pgfsetfillcolor{currentfill}%
\pgfsetlinewidth{0.803000pt}%
\definecolor{currentstroke}{rgb}{0.000000,0.000000,0.000000}%
\pgfsetstrokecolor{currentstroke}%
\pgfsetdash{}{0pt}%
\pgfsys@defobject{currentmarker}{\pgfqpoint{0.000000in}{-0.048611in}}{\pgfqpoint{0.000000in}{0.000000in}}{%
\pgfpathmoveto{\pgfqpoint{0.000000in}{0.000000in}}%
\pgfpathlineto{\pgfqpoint{0.000000in}{-0.048611in}}%
\pgfusepath{stroke,fill}%
}%
\begin{pgfscope}%
\pgfsys@transformshift{0.955556in}{2.059688in}%
\pgfsys@useobject{currentmarker}{}%
\end{pgfscope}%
\end{pgfscope}%
\begin{pgfscope}%
\definecolor{textcolor}{rgb}{0.000000,0.000000,0.000000}%
\pgfsetstrokecolor{textcolor}%
\pgfsetfillcolor{textcolor}%
\pgftext[x=0.955556in,y=1.962466in,,top]{\color{textcolor}{\rmfamily\fontsize{6.000000}{7.200000}\selectfont\catcode`\^=\active\def^{\ifmmode\sp\else\^{}\fi}\catcode`\%=\active\def%{\%}\ensuremath{-}0.5}}%
\end{pgfscope}%
\begin{pgfscope}%
\pgfpathrectangle{\pgfqpoint{0.537500in}{2.059688in}}{\pgfqpoint{1.505000in}{1.156312in}}%
\pgfusepath{clip}%
\pgfsetrectcap%
\pgfsetroundjoin%
\pgfsetlinewidth{1.003750pt}%
\definecolor{currentstroke}{rgb}{1.000000,1.000000,1.000000}%
\pgfsetstrokecolor{currentstroke}%
\pgfsetdash{}{0pt}%
\pgfpathmoveto{\pgfqpoint{1.234259in}{2.059688in}}%
\pgfpathlineto{\pgfqpoint{1.234259in}{3.216000in}}%
\pgfusepath{stroke}%
\end{pgfscope}%
\begin{pgfscope}%
\pgfsetbuttcap%
\pgfsetroundjoin%
\definecolor{currentfill}{rgb}{0.000000,0.000000,0.000000}%
\pgfsetfillcolor{currentfill}%
\pgfsetlinewidth{0.803000pt}%
\definecolor{currentstroke}{rgb}{0.000000,0.000000,0.000000}%
\pgfsetstrokecolor{currentstroke}%
\pgfsetdash{}{0pt}%
\pgfsys@defobject{currentmarker}{\pgfqpoint{0.000000in}{-0.048611in}}{\pgfqpoint{0.000000in}{0.000000in}}{%
\pgfpathmoveto{\pgfqpoint{0.000000in}{0.000000in}}%
\pgfpathlineto{\pgfqpoint{0.000000in}{-0.048611in}}%
\pgfusepath{stroke,fill}%
}%
\begin{pgfscope}%
\pgfsys@transformshift{1.234259in}{2.059688in}%
\pgfsys@useobject{currentmarker}{}%
\end{pgfscope}%
\end{pgfscope}%
\begin{pgfscope}%
\definecolor{textcolor}{rgb}{0.000000,0.000000,0.000000}%
\pgfsetstrokecolor{textcolor}%
\pgfsetfillcolor{textcolor}%
\pgftext[x=1.234259in,y=1.962466in,,top]{\color{textcolor}{\rmfamily\fontsize{6.000000}{7.200000}\selectfont\catcode`\^=\active\def^{\ifmmode\sp\else\^{}\fi}\catcode`\%=\active\def%{\%}0.0}}%
\end{pgfscope}%
\begin{pgfscope}%
\pgfpathrectangle{\pgfqpoint{0.537500in}{2.059688in}}{\pgfqpoint{1.505000in}{1.156312in}}%
\pgfusepath{clip}%
\pgfsetrectcap%
\pgfsetroundjoin%
\pgfsetlinewidth{1.003750pt}%
\definecolor{currentstroke}{rgb}{1.000000,1.000000,1.000000}%
\pgfsetstrokecolor{currentstroke}%
\pgfsetdash{}{0pt}%
\pgfpathmoveto{\pgfqpoint{1.512963in}{2.059688in}}%
\pgfpathlineto{\pgfqpoint{1.512963in}{3.216000in}}%
\pgfusepath{stroke}%
\end{pgfscope}%
\begin{pgfscope}%
\pgfsetbuttcap%
\pgfsetroundjoin%
\definecolor{currentfill}{rgb}{0.000000,0.000000,0.000000}%
\pgfsetfillcolor{currentfill}%
\pgfsetlinewidth{0.803000pt}%
\definecolor{currentstroke}{rgb}{0.000000,0.000000,0.000000}%
\pgfsetstrokecolor{currentstroke}%
\pgfsetdash{}{0pt}%
\pgfsys@defobject{currentmarker}{\pgfqpoint{0.000000in}{-0.048611in}}{\pgfqpoint{0.000000in}{0.000000in}}{%
\pgfpathmoveto{\pgfqpoint{0.000000in}{0.000000in}}%
\pgfpathlineto{\pgfqpoint{0.000000in}{-0.048611in}}%
\pgfusepath{stroke,fill}%
}%
\begin{pgfscope}%
\pgfsys@transformshift{1.512963in}{2.059688in}%
\pgfsys@useobject{currentmarker}{}%
\end{pgfscope}%
\end{pgfscope}%
\begin{pgfscope}%
\definecolor{textcolor}{rgb}{0.000000,0.000000,0.000000}%
\pgfsetstrokecolor{textcolor}%
\pgfsetfillcolor{textcolor}%
\pgftext[x=1.512963in,y=1.962466in,,top]{\color{textcolor}{\rmfamily\fontsize{6.000000}{7.200000}\selectfont\catcode`\^=\active\def^{\ifmmode\sp\else\^{}\fi}\catcode`\%=\active\def%{\%}0.5}}%
\end{pgfscope}%
\begin{pgfscope}%
\pgfpathrectangle{\pgfqpoint{0.537500in}{2.059688in}}{\pgfqpoint{1.505000in}{1.156312in}}%
\pgfusepath{clip}%
\pgfsetrectcap%
\pgfsetroundjoin%
\pgfsetlinewidth{1.003750pt}%
\definecolor{currentstroke}{rgb}{1.000000,1.000000,1.000000}%
\pgfsetstrokecolor{currentstroke}%
\pgfsetdash{}{0pt}%
\pgfpathmoveto{\pgfqpoint{1.791667in}{2.059688in}}%
\pgfpathlineto{\pgfqpoint{1.791667in}{3.216000in}}%
\pgfusepath{stroke}%
\end{pgfscope}%
\begin{pgfscope}%
\pgfsetbuttcap%
\pgfsetroundjoin%
\definecolor{currentfill}{rgb}{0.000000,0.000000,0.000000}%
\pgfsetfillcolor{currentfill}%
\pgfsetlinewidth{0.803000pt}%
\definecolor{currentstroke}{rgb}{0.000000,0.000000,0.000000}%
\pgfsetstrokecolor{currentstroke}%
\pgfsetdash{}{0pt}%
\pgfsys@defobject{currentmarker}{\pgfqpoint{0.000000in}{-0.048611in}}{\pgfqpoint{0.000000in}{0.000000in}}{%
\pgfpathmoveto{\pgfqpoint{0.000000in}{0.000000in}}%
\pgfpathlineto{\pgfqpoint{0.000000in}{-0.048611in}}%
\pgfusepath{stroke,fill}%
}%
\begin{pgfscope}%
\pgfsys@transformshift{1.791667in}{2.059688in}%
\pgfsys@useobject{currentmarker}{}%
\end{pgfscope}%
\end{pgfscope}%
\begin{pgfscope}%
\definecolor{textcolor}{rgb}{0.000000,0.000000,0.000000}%
\pgfsetstrokecolor{textcolor}%
\pgfsetfillcolor{textcolor}%
\pgftext[x=1.791667in,y=1.962466in,,top]{\color{textcolor}{\rmfamily\fontsize{6.000000}{7.200000}\selectfont\catcode`\^=\active\def^{\ifmmode\sp\else\^{}\fi}\catcode`\%=\active\def%{\%}1.0}}%
\end{pgfscope}%
\begin{pgfscope}%
\definecolor{textcolor}{rgb}{0.000000,0.000000,0.000000}%
\pgfsetstrokecolor{textcolor}%
\pgfsetfillcolor{textcolor}%
\pgftext[x=1.290000in,y=1.823577in,,top]{\color{textcolor}{\rmfamily\fontsize{6.000000}{7.200000}\selectfont\catcode`\^=\active\def^{\ifmmode\sp\else\^{}\fi}\catcode`\%=\active\def%{\%}predictor $r$}}%
\end{pgfscope}%
\begin{pgfscope}%
\pgfpathrectangle{\pgfqpoint{0.537500in}{2.059688in}}{\pgfqpoint{1.505000in}{1.156312in}}%
\pgfusepath{clip}%
\pgfsetrectcap%
\pgfsetroundjoin%
\pgfsetlinewidth{1.003750pt}%
\definecolor{currentstroke}{rgb}{1.000000,1.000000,1.000000}%
\pgfsetstrokecolor{currentstroke}%
\pgfsetdash{}{0pt}%
\pgfpathmoveto{\pgfqpoint{0.537500in}{2.059688in}}%
\pgfpathlineto{\pgfqpoint{2.042500in}{2.059688in}}%
\pgfusepath{stroke}%
\end{pgfscope}%
\begin{pgfscope}%
\pgfsetbuttcap%
\pgfsetroundjoin%
\definecolor{currentfill}{rgb}{0.000000,0.000000,0.000000}%
\pgfsetfillcolor{currentfill}%
\pgfsetlinewidth{0.803000pt}%
\definecolor{currentstroke}{rgb}{0.000000,0.000000,0.000000}%
\pgfsetstrokecolor{currentstroke}%
\pgfsetdash{}{0pt}%
\pgfsys@defobject{currentmarker}{\pgfqpoint{-0.048611in}{0.000000in}}{\pgfqpoint{-0.000000in}{0.000000in}}{%
\pgfpathmoveto{\pgfqpoint{-0.000000in}{0.000000in}}%
\pgfpathlineto{\pgfqpoint{-0.048611in}{0.000000in}}%
\pgfusepath{stroke,fill}%
}%
\begin{pgfscope}%
\pgfsys@transformshift{0.537500in}{2.059688in}%
\pgfsys@useobject{currentmarker}{}%
\end{pgfscope}%
\end{pgfscope}%
\begin{pgfscope}%
\definecolor{textcolor}{rgb}{0.000000,0.000000,0.000000}%
\pgfsetstrokecolor{textcolor}%
\pgfsetfillcolor{textcolor}%
\pgftext[x=0.307750in, y=2.028031in, left, base]{\color{textcolor}{\rmfamily\fontsize{6.000000}{7.200000}\selectfont\catcode`\^=\active\def^{\ifmmode\sp\else\^{}\fi}\catcode`\%=\active\def%{\%}0.0}}%
\end{pgfscope}%
\begin{pgfscope}%
\pgfpathrectangle{\pgfqpoint{0.537500in}{2.059688in}}{\pgfqpoint{1.505000in}{1.156312in}}%
\pgfusepath{clip}%
\pgfsetrectcap%
\pgfsetroundjoin%
\pgfsetlinewidth{1.003750pt}%
\definecolor{currentstroke}{rgb}{1.000000,1.000000,1.000000}%
\pgfsetstrokecolor{currentstroke}%
\pgfsetdash{}{0pt}%
\pgfpathmoveto{\pgfqpoint{0.537500in}{2.372205in}}%
\pgfpathlineto{\pgfqpoint{2.042500in}{2.372205in}}%
\pgfusepath{stroke}%
\end{pgfscope}%
\begin{pgfscope}%
\pgfsetbuttcap%
\pgfsetroundjoin%
\definecolor{currentfill}{rgb}{0.000000,0.000000,0.000000}%
\pgfsetfillcolor{currentfill}%
\pgfsetlinewidth{0.803000pt}%
\definecolor{currentstroke}{rgb}{0.000000,0.000000,0.000000}%
\pgfsetstrokecolor{currentstroke}%
\pgfsetdash{}{0pt}%
\pgfsys@defobject{currentmarker}{\pgfqpoint{-0.048611in}{0.000000in}}{\pgfqpoint{-0.000000in}{0.000000in}}{%
\pgfpathmoveto{\pgfqpoint{-0.000000in}{0.000000in}}%
\pgfpathlineto{\pgfqpoint{-0.048611in}{0.000000in}}%
\pgfusepath{stroke,fill}%
}%
\begin{pgfscope}%
\pgfsys@transformshift{0.537500in}{2.372205in}%
\pgfsys@useobject{currentmarker}{}%
\end{pgfscope}%
\end{pgfscope}%
\begin{pgfscope}%
\definecolor{textcolor}{rgb}{0.000000,0.000000,0.000000}%
\pgfsetstrokecolor{textcolor}%
\pgfsetfillcolor{textcolor}%
\pgftext[x=0.307750in, y=2.340548in, left, base]{\color{textcolor}{\rmfamily\fontsize{6.000000}{7.200000}\selectfont\catcode`\^=\active\def^{\ifmmode\sp\else\^{}\fi}\catcode`\%=\active\def%{\%}0.5}}%
\end{pgfscope}%
\begin{pgfscope}%
\pgfpathrectangle{\pgfqpoint{0.537500in}{2.059688in}}{\pgfqpoint{1.505000in}{1.156312in}}%
\pgfusepath{clip}%
\pgfsetrectcap%
\pgfsetroundjoin%
\pgfsetlinewidth{1.003750pt}%
\definecolor{currentstroke}{rgb}{1.000000,1.000000,1.000000}%
\pgfsetstrokecolor{currentstroke}%
\pgfsetdash{}{0pt}%
\pgfpathmoveto{\pgfqpoint{0.537500in}{2.684722in}}%
\pgfpathlineto{\pgfqpoint{2.042500in}{2.684722in}}%
\pgfusepath{stroke}%
\end{pgfscope}%
\begin{pgfscope}%
\pgfsetbuttcap%
\pgfsetroundjoin%
\definecolor{currentfill}{rgb}{0.000000,0.000000,0.000000}%
\pgfsetfillcolor{currentfill}%
\pgfsetlinewidth{0.803000pt}%
\definecolor{currentstroke}{rgb}{0.000000,0.000000,0.000000}%
\pgfsetstrokecolor{currentstroke}%
\pgfsetdash{}{0pt}%
\pgfsys@defobject{currentmarker}{\pgfqpoint{-0.048611in}{0.000000in}}{\pgfqpoint{-0.000000in}{0.000000in}}{%
\pgfpathmoveto{\pgfqpoint{-0.000000in}{0.000000in}}%
\pgfpathlineto{\pgfqpoint{-0.048611in}{0.000000in}}%
\pgfusepath{stroke,fill}%
}%
\begin{pgfscope}%
\pgfsys@transformshift{0.537500in}{2.684722in}%
\pgfsys@useobject{currentmarker}{}%
\end{pgfscope}%
\end{pgfscope}%
\begin{pgfscope}%
\definecolor{textcolor}{rgb}{0.000000,0.000000,0.000000}%
\pgfsetstrokecolor{textcolor}%
\pgfsetfillcolor{textcolor}%
\pgftext[x=0.307750in, y=2.653065in, left, base]{\color{textcolor}{\rmfamily\fontsize{6.000000}{7.200000}\selectfont\catcode`\^=\active\def^{\ifmmode\sp\else\^{}\fi}\catcode`\%=\active\def%{\%}1.0}}%
\end{pgfscope}%
\begin{pgfscope}%
\pgfpathrectangle{\pgfqpoint{0.537500in}{2.059688in}}{\pgfqpoint{1.505000in}{1.156312in}}%
\pgfusepath{clip}%
\pgfsetrectcap%
\pgfsetroundjoin%
\pgfsetlinewidth{1.003750pt}%
\definecolor{currentstroke}{rgb}{1.000000,1.000000,1.000000}%
\pgfsetstrokecolor{currentstroke}%
\pgfsetdash{}{0pt}%
\pgfpathmoveto{\pgfqpoint{0.537500in}{2.997238in}}%
\pgfpathlineto{\pgfqpoint{2.042500in}{2.997238in}}%
\pgfusepath{stroke}%
\end{pgfscope}%
\begin{pgfscope}%
\pgfsetbuttcap%
\pgfsetroundjoin%
\definecolor{currentfill}{rgb}{0.000000,0.000000,0.000000}%
\pgfsetfillcolor{currentfill}%
\pgfsetlinewidth{0.803000pt}%
\definecolor{currentstroke}{rgb}{0.000000,0.000000,0.000000}%
\pgfsetstrokecolor{currentstroke}%
\pgfsetdash{}{0pt}%
\pgfsys@defobject{currentmarker}{\pgfqpoint{-0.048611in}{0.000000in}}{\pgfqpoint{-0.000000in}{0.000000in}}{%
\pgfpathmoveto{\pgfqpoint{-0.000000in}{0.000000in}}%
\pgfpathlineto{\pgfqpoint{-0.048611in}{0.000000in}}%
\pgfusepath{stroke,fill}%
}%
\begin{pgfscope}%
\pgfsys@transformshift{0.537500in}{2.997238in}%
\pgfsys@useobject{currentmarker}{}%
\end{pgfscope}%
\end{pgfscope}%
\begin{pgfscope}%
\definecolor{textcolor}{rgb}{0.000000,0.000000,0.000000}%
\pgfsetstrokecolor{textcolor}%
\pgfsetfillcolor{textcolor}%
\pgftext[x=0.307750in, y=2.965581in, left, base]{\color{textcolor}{\rmfamily\fontsize{6.000000}{7.200000}\selectfont\catcode`\^=\active\def^{\ifmmode\sp\else\^{}\fi}\catcode`\%=\active\def%{\%}1.5}}%
\end{pgfscope}%
\begin{pgfscope}%
\definecolor{textcolor}{rgb}{0.000000,0.000000,0.000000}%
\pgfsetstrokecolor{textcolor}%
\pgfsetfillcolor{textcolor}%
\pgftext[x=0.252195in,y=2.637844in,,bottom,rotate=90.000000]{\color{textcolor}{\rmfamily\fontsize{6.000000}{7.200000}\selectfont\catcode`\^=\active\def^{\ifmmode\sp\else\^{}\fi}\catcode`\%=\active\def%{\%}shared variance $s$}}%
\end{pgfscope}%
\begin{pgfscope}%
\pgfpathrectangle{\pgfqpoint{0.537500in}{2.059688in}}{\pgfqpoint{1.505000in}{1.156312in}}%
\pgfusepath{clip}%
\pgfsetrectcap%
\pgfsetroundjoin%
\pgfsetlinewidth{1.003750pt}%
\definecolor{currentstroke}{rgb}{0.447059,0.447059,0.447059}%
\pgfsetstrokecolor{currentstroke}%
\pgfsetdash{}{0pt}%
\pgfpathmoveto{\pgfqpoint{0.537500in}{2.840980in}}%
\pgfpathlineto{\pgfqpoint{1.232943in}{2.061164in}}%
\pgfpathlineto{\pgfqpoint{1.235097in}{2.060627in}}%
\pgfpathlineto{\pgfqpoint{2.042500in}{2.965987in}}%
\pgfpathlineto{\pgfqpoint{2.042500in}{2.965987in}}%
\pgfusepath{stroke}%
\end{pgfscope}%
\begin{pgfscope}%
\pgfpathrectangle{\pgfqpoint{0.537500in}{2.059688in}}{\pgfqpoint{1.505000in}{1.156312in}}%
\pgfusepath{clip}%
\pgfsetrectcap%
\pgfsetroundjoin%
\pgfsetlinewidth{2.007500pt}%
\definecolor{currentstroke}{rgb}{0.168627,0.611765,0.411765}%
\pgfsetstrokecolor{currentstroke}%
\pgfsetdash{}{0pt}%
\pgfpathmoveto{\pgfqpoint{1.791667in}{2.684722in}}%
\pgfpathlineto{\pgfqpoint{1.791667in}{3.216000in}}%
\pgfusepath{stroke}%
\end{pgfscope}%
\begin{pgfscope}%
\pgfsetrectcap%
\pgfsetmiterjoin%
\pgfsetlinewidth{0.803000pt}%
\definecolor{currentstroke}{rgb}{0.000000,0.000000,0.000000}%
\pgfsetstrokecolor{currentstroke}%
\pgfsetdash{}{0pt}%
\pgfpathmoveto{\pgfqpoint{0.537500in}{2.059688in}}%
\pgfpathlineto{\pgfqpoint{0.537500in}{3.216000in}}%
\pgfusepath{stroke}%
\end{pgfscope}%
\begin{pgfscope}%
\pgfsetrectcap%
\pgfsetmiterjoin%
\pgfsetlinewidth{0.803000pt}%
\definecolor{currentstroke}{rgb}{0.000000,0.000000,0.000000}%
\pgfsetstrokecolor{currentstroke}%
\pgfsetdash{}{0pt}%
\pgfpathmoveto{\pgfqpoint{0.537500in}{2.059688in}}%
\pgfpathlineto{\pgfqpoint{2.042500in}{2.059688in}}%
\pgfusepath{stroke}%
\end{pgfscope}%
\begin{pgfscope}%
\pgfpathrectangle{\pgfqpoint{0.537500in}{2.059688in}}{\pgfqpoint{1.505000in}{1.156312in}}%
\pgfusepath{clip}%
\pgfsetrectcap%
\pgfsetroundjoin%
\pgfsetlinewidth{2.158062pt}%
\definecolor{currentstroke}{rgb}{0.141176,0.345098,0.650980}%
\pgfsetstrokecolor{currentstroke}%
\pgfsetdash{}{0pt}%
\pgfpathmoveto{\pgfqpoint{0.648981in}{2.732899in}}%
\pgfpathlineto{\pgfqpoint{0.756547in}{2.615965in}}%
\pgfpathlineto{\pgfqpoint{0.826223in}{2.541192in}}%
\pgfpathlineto{\pgfqpoint{0.883040in}{2.481120in}}%
\pgfpathlineto{\pgfqpoint{0.930240in}{2.432137in}}%
\pgfpathlineto{\pgfqpoint{0.970077in}{2.391737in}}%
\pgfpathlineto{\pgfqpoint{1.004172in}{2.358130in}}%
\pgfpathlineto{\pgfqpoint{1.033717in}{2.330003in}}%
\pgfpathlineto{\pgfqpoint{1.059614in}{2.306373in}}%
\pgfpathlineto{\pgfqpoint{1.082556in}{2.286490in}}%
\pgfpathlineto{\pgfqpoint{1.103089in}{2.269774in}}%
\pgfpathlineto{\pgfqpoint{1.121646in}{2.255770in}}%
\pgfpathlineto{\pgfqpoint{1.135307in}{2.246273in}}%
\pgfpathlineto{\pgfqpoint{1.148079in}{2.238129in}}%
\pgfpathlineto{\pgfqpoint{1.160091in}{2.231208in}}%
\pgfpathlineto{\pgfqpoint{1.171454in}{2.225405in}}%
\pgfpathlineto{\pgfqpoint{1.182266in}{2.220628in}}%
\pgfpathlineto{\pgfqpoint{1.192612in}{2.216798in}}%
\pgfpathlineto{\pgfqpoint{1.202568in}{2.213848in}}%
\pgfpathlineto{\pgfqpoint{1.212200in}{2.211722in}}%
\pgfpathlineto{\pgfqpoint{1.221568in}{2.210370in}}%
\pgfpathlineto{\pgfqpoint{1.230725in}{2.209749in}}%
\pgfpathlineto{\pgfqpoint{1.239720in}{2.209821in}}%
\pgfpathlineto{\pgfqpoint{1.248593in}{2.210555in}}%
\pgfpathlineto{\pgfqpoint{1.257385in}{2.211921in}}%
\pgfpathlineto{\pgfqpoint{1.268311in}{2.214480in}}%
\pgfpathlineto{\pgfqpoint{1.279221in}{2.217942in}}%
\pgfpathlineto{\pgfqpoint{1.290168in}{2.222269in}}%
\pgfpathlineto{\pgfqpoint{1.301194in}{2.227425in}}%
\pgfpathlineto{\pgfqpoint{1.312336in}{2.233376in}}%
\pgfpathlineto{\pgfqpoint{1.325902in}{2.241519in}}%
\pgfpathlineto{\pgfqpoint{1.339718in}{2.250702in}}%
\pgfpathlineto{\pgfqpoint{1.356186in}{2.262652in}}%
\pgfpathlineto{\pgfqpoint{1.373047in}{2.275841in}}%
\pgfpathlineto{\pgfqpoint{1.392794in}{2.292291in}}%
\pgfpathlineto{\pgfqpoint{1.415564in}{2.312342in}}%
\pgfpathlineto{\pgfqpoint{1.441420in}{2.336208in}}%
\pgfpathlineto{\pgfqpoint{1.470280in}{2.363901in}}%
\pgfpathlineto{\pgfqpoint{1.504453in}{2.397765in}}%
\pgfpathlineto{\pgfqpoint{1.545628in}{2.439694in}}%
\pgfpathlineto{\pgfqpoint{1.593898in}{2.489956in}}%
\pgfpathlineto{\pgfqpoint{1.653933in}{2.553608in}}%
\pgfpathlineto{\pgfqpoint{1.727725in}{2.632995in}}%
\pgfpathlineto{\pgfqpoint{1.787975in}{2.698446in}}%
\pgfpathlineto{\pgfqpoint{1.787975in}{2.698446in}}%
\pgfusepath{stroke}%
\end{pgfscope}%
\begin{pgfscope}%
\pgfsetroundcap%
\pgfsetroundjoin%
\pgfsetlinewidth{1.254687pt}%
\definecolor{currentstroke}{rgb}{0.141176,0.345098,0.650980}%
\pgfsetstrokecolor{currentstroke}%
\pgfsetdash{}{0pt}%
\pgfpathmoveto{\pgfqpoint{1.577532in}{2.472805in}}%
\pgfpathquadraticcurveto{\pgfqpoint{1.578720in}{2.474047in}}{\pgfqpoint{1.566490in}{2.461263in}}%
\pgfusepath{stroke}%
\end{pgfscope}%
\begin{pgfscope}%
\pgfsetroundcap%
\pgfsetroundjoin%
\definecolor{currentfill}{rgb}{0.141176,0.345098,0.650980}%
\pgfsetfillcolor{currentfill}%
\pgfsetlinewidth{1.254687pt}%
\definecolor{currentstroke}{rgb}{0.141176,0.345098,0.650980}%
\pgfsetstrokecolor{currentstroke}%
\pgfsetdash{}{0pt}%
\pgfpathmoveto{\pgfqpoint{1.531404in}{2.448699in}}%
\pgfpathlineto{\pgfqpoint{1.566490in}{2.461263in}}%
\pgfpathlineto{\pgfqpoint{1.555490in}{2.425656in}}%
\pgfpathlineto{\pgfqpoint{1.531404in}{2.448699in}}%
\pgfpathclose%
\pgfusepath{stroke,fill}%
\end{pgfscope}%
\begin{pgfscope}%
\pgfpathrectangle{\pgfqpoint{0.537500in}{2.059688in}}{\pgfqpoint{1.505000in}{1.156312in}}%
\pgfusepath{clip}%
\pgfsetrectcap%
\pgfsetroundjoin%
\pgfsetlinewidth{2.158062pt}%
\definecolor{currentstroke}{rgb}{0.819608,0.286275,0.356863}%
\pgfsetstrokecolor{currentstroke}%
\pgfsetdash{}{0pt}%
\pgfpathmoveto{\pgfqpoint{0.721444in}{2.681840in}}%
\pgfpathlineto{\pgfqpoint{0.771731in}{2.630130in}}%
\pgfpathlineto{\pgfqpoint{0.815874in}{2.585530in}}%
\pgfpathlineto{\pgfqpoint{0.854966in}{2.546823in}}%
\pgfpathlineto{\pgfqpoint{0.889861in}{2.513063in}}%
\pgfpathlineto{\pgfqpoint{0.921236in}{2.483496in}}%
\pgfpathlineto{\pgfqpoint{0.949635in}{2.457523in}}%
\pgfpathlineto{\pgfqpoint{0.975502in}{2.434654in}}%
\pgfpathlineto{\pgfqpoint{0.999199in}{2.414492in}}%
\pgfpathlineto{\pgfqpoint{1.021029in}{2.396706in}}%
\pgfpathlineto{\pgfqpoint{1.041245in}{2.381021in}}%
\pgfpathlineto{\pgfqpoint{1.060061in}{2.367206in}}%
\pgfpathlineto{\pgfqpoint{1.077657in}{2.355066in}}%
\pgfpathlineto{\pgfqpoint{1.094189in}{2.344436in}}%
\pgfpathlineto{\pgfqpoint{1.109792in}{2.335173in}}%
\pgfpathlineto{\pgfqpoint{1.124580in}{2.327154in}}%
\pgfpathlineto{\pgfqpoint{1.138655in}{2.320274in}}%
\pgfpathlineto{\pgfqpoint{1.152105in}{2.314440in}}%
\pgfpathlineto{\pgfqpoint{1.165008in}{2.309573in}}%
\pgfpathlineto{\pgfqpoint{1.177431in}{2.305601in}}%
\pgfpathlineto{\pgfqpoint{1.189435in}{2.302461in}}%
\pgfpathlineto{\pgfqpoint{1.201074in}{2.300098in}}%
\pgfpathlineto{\pgfqpoint{1.212394in}{2.298463in}}%
\pgfpathlineto{\pgfqpoint{1.223438in}{2.297511in}}%
\pgfpathlineto{\pgfqpoint{1.234243in}{2.297201in}}%
\pgfpathlineto{\pgfqpoint{1.244843in}{2.297497in}}%
\pgfpathlineto{\pgfqpoint{1.255266in}{2.298366in}}%
\pgfpathlineto{\pgfqpoint{1.268935in}{2.300363in}}%
\pgfpathlineto{\pgfqpoint{1.282389in}{2.303255in}}%
\pgfpathlineto{\pgfqpoint{1.295675in}{2.306983in}}%
\pgfpathlineto{\pgfqpoint{1.308828in}{2.311489in}}%
\pgfpathlineto{\pgfqpoint{1.321879in}{2.316720in}}%
\pgfpathlineto{\pgfqpoint{1.334851in}{2.322624in}}%
\pgfpathlineto{\pgfqpoint{1.350983in}{2.330877in}}%
\pgfpathlineto{\pgfqpoint{1.367044in}{2.340013in}}%
\pgfpathlineto{\pgfqpoint{1.383044in}{2.349941in}}%
\pgfpathlineto{\pgfqpoint{1.402166in}{2.362774in}}%
\pgfpathlineto{\pgfqpoint{1.421189in}{2.376466in}}%
\pgfpathlineto{\pgfqpoint{1.443216in}{2.393324in}}%
\pgfpathlineto{\pgfqpoint{1.464997in}{2.410907in}}%
\pgfpathlineto{\pgfqpoint{1.489477in}{2.431592in}}%
\pgfpathlineto{\pgfqpoint{1.516315in}{2.455216in}}%
\pgfpathlineto{\pgfqpoint{1.547806in}{2.483982in}}%
\pgfpathlineto{\pgfqpoint{1.582688in}{2.516919in}}%
\pgfpathlineto{\pgfqpoint{1.621471in}{2.554595in}}%
\pgfpathlineto{\pgfqpoint{1.666455in}{2.599391in}}%
\pgfpathlineto{\pgfqpoint{1.717184in}{2.651000in}}%
\pgfpathlineto{\pgfqpoint{1.776258in}{2.712206in}}%
\pgfpathlineto{\pgfqpoint{1.790609in}{2.727219in}}%
\pgfpathlineto{\pgfqpoint{1.790609in}{2.727219in}}%
\pgfusepath{stroke}%
\end{pgfscope}%
\begin{pgfscope}%
\pgfsetroundcap%
\pgfsetroundjoin%
\pgfsetlinewidth{1.254687pt}%
\definecolor{currentstroke}{rgb}{0.819608,0.286275,0.356863}%
\pgfsetstrokecolor{currentstroke}%
\pgfsetdash{}{0pt}%
\pgfpathmoveto{\pgfqpoint{1.694524in}{2.627822in}}%
\pgfpathquadraticcurveto{\pgfqpoint{1.695262in}{2.628575in}}{\pgfqpoint{1.682405in}{2.615474in}}%
\pgfusepath{stroke}%
\end{pgfscope}%
\begin{pgfscope}%
\pgfsetroundcap%
\pgfsetroundjoin%
\definecolor{currentfill}{rgb}{0.819608,0.286275,0.356863}%
\pgfsetfillcolor{currentfill}%
\pgfsetlinewidth{1.254687pt}%
\definecolor{currentstroke}{rgb}{0.819608,0.286275,0.356863}%
\pgfsetstrokecolor{currentstroke}%
\pgfsetdash{}{0pt}%
\pgfpathmoveto{\pgfqpoint{1.647161in}{2.603358in}}%
\pgfpathlineto{\pgfqpoint{1.682405in}{2.615474in}}%
\pgfpathlineto{\pgfqpoint{1.670951in}{2.580010in}}%
\pgfpathlineto{\pgfqpoint{1.647161in}{2.603358in}}%
\pgfpathclose%
\pgfusepath{stroke,fill}%
\end{pgfscope}%
\begin{pgfscope}%
\pgfpathrectangle{\pgfqpoint{0.537500in}{2.059688in}}{\pgfqpoint{1.505000in}{1.156312in}}%
\pgfusepath{clip}%
\pgfsetrectcap%
\pgfsetroundjoin%
\pgfsetlinewidth{2.158062pt}%
\definecolor{currentstroke}{rgb}{0.000000,0.529412,0.423529}%
\pgfsetstrokecolor{currentstroke}%
\pgfsetdash{}{0pt}%
\pgfpathmoveto{\pgfqpoint{0.810630in}{2.657241in}}%
\pgfpathlineto{\pgfqpoint{0.841008in}{2.630536in}}%
\pgfpathlineto{\pgfqpoint{0.869187in}{2.606496in}}%
\pgfpathlineto{\pgfqpoint{0.895430in}{2.584829in}}%
\pgfpathlineto{\pgfqpoint{0.919959in}{2.565285in}}%
\pgfpathlineto{\pgfqpoint{0.942969in}{2.547651in}}%
\pgfpathlineto{\pgfqpoint{0.964625in}{2.531741in}}%
\pgfpathlineto{\pgfqpoint{0.985072in}{2.517395in}}%
\pgfpathlineto{\pgfqpoint{1.004437in}{2.504472in}}%
\pgfpathlineto{\pgfqpoint{1.022828in}{2.492849in}}%
\pgfpathlineto{\pgfqpoint{1.040344in}{2.482419in}}%
\pgfpathlineto{\pgfqpoint{1.057071in}{2.473084in}}%
\pgfpathlineto{\pgfqpoint{1.073083in}{2.464761in}}%
\pgfpathlineto{\pgfqpoint{1.088450in}{2.457372in}}%
\pgfpathlineto{\pgfqpoint{1.103230in}{2.450850in}}%
\pgfpathlineto{\pgfqpoint{1.117477in}{2.445134in}}%
\pgfpathlineto{\pgfqpoint{1.131240in}{2.440168in}}%
\pgfpathlineto{\pgfqpoint{1.144561in}{2.435902in}}%
\pgfpathlineto{\pgfqpoint{1.157479in}{2.432291in}}%
\pgfpathlineto{\pgfqpoint{1.176175in}{2.428012in}}%
\pgfpathlineto{\pgfqpoint{1.194144in}{2.424988in}}%
\pgfpathlineto{\pgfqpoint{1.211474in}{2.423107in}}%
\pgfpathlineto{\pgfqpoint{1.228240in}{2.422271in}}%
\pgfpathlineto{\pgfqpoint{1.244505in}{2.422390in}}%
\pgfpathlineto{\pgfqpoint{1.260326in}{2.423384in}}%
\pgfpathlineto{\pgfqpoint{1.275747in}{2.425180in}}%
\pgfpathlineto{\pgfqpoint{1.290809in}{2.427712in}}%
\pgfpathlineto{\pgfqpoint{1.305544in}{2.430915in}}%
\pgfpathlineto{\pgfqpoint{1.319979in}{2.434734in}}%
\pgfpathlineto{\pgfqpoint{1.334137in}{2.439113in}}%
\pgfpathlineto{\pgfqpoint{1.352615in}{2.445737in}}%
\pgfpathlineto{\pgfqpoint{1.370666in}{2.453155in}}%
\pgfpathlineto{\pgfqpoint{1.388312in}{2.461261in}}%
\pgfpathlineto{\pgfqpoint{1.405569in}{2.469957in}}%
\pgfpathlineto{\pgfqpoint{1.426602in}{2.481515in}}%
\pgfpathlineto{\pgfqpoint{1.447042in}{2.493681in}}%
\pgfpathlineto{\pgfqpoint{1.466880in}{2.506297in}}%
\pgfpathlineto{\pgfqpoint{1.489872in}{2.521828in}}%
\pgfpathlineto{\pgfqpoint{1.515531in}{2.540204in}}%
\pgfpathlineto{\pgfqpoint{1.539883in}{2.558553in}}%
\pgfpathlineto{\pgfqpoint{1.566046in}{2.579144in}}%
\pgfpathlineto{\pgfqpoint{1.596159in}{2.603838in}}%
\pgfpathlineto{\pgfqpoint{1.628427in}{2.631330in}}%
\pgfpathlineto{\pgfqpoint{1.662901in}{2.661719in}}%
\pgfpathlineto{\pgfqpoint{1.700993in}{2.696339in}}%
\pgfpathlineto{\pgfqpoint{1.742005in}{2.734652in}}%
\pgfpathlineto{\pgfqpoint{1.787371in}{2.778081in}}%
\pgfpathlineto{\pgfqpoint{1.791433in}{2.782018in}}%
\pgfpathlineto{\pgfqpoint{1.791433in}{2.782018in}}%
\pgfusepath{stroke}%
\end{pgfscope}%
\begin{pgfscope}%
\pgfsetroundcap%
\pgfsetroundjoin%
\pgfsetlinewidth{1.254687pt}%
\definecolor{currentstroke}{rgb}{0.000000,0.529412,0.423529}%
\pgfsetstrokecolor{currentstroke}%
\pgfsetdash{}{0pt}%
\pgfpathmoveto{\pgfqpoint{1.756640in}{2.748551in}}%
\pgfpathquadraticcurveto{\pgfqpoint{1.756962in}{2.748858in}}{\pgfqpoint{1.743237in}{2.735770in}}%
\pgfusepath{stroke}%
\end{pgfscope}%
\begin{pgfscope}%
\pgfsetroundcap%
\pgfsetroundjoin%
\definecolor{currentfill}{rgb}{0.000000,0.529412,0.423529}%
\pgfsetfillcolor{currentfill}%
\pgfsetlinewidth{1.254687pt}%
\definecolor{currentstroke}{rgb}{0.000000,0.529412,0.423529}%
\pgfsetstrokecolor{currentstroke}%
\pgfsetdash{}{0pt}%
\pgfpathmoveto{\pgfqpoint{1.707613in}{2.724826in}}%
\pgfpathlineto{\pgfqpoint{1.743237in}{2.735770in}}%
\pgfpathlineto{\pgfqpoint{1.730617in}{2.700704in}}%
\pgfpathlineto{\pgfqpoint{1.707613in}{2.724826in}}%
\pgfpathclose%
\pgfusepath{stroke,fill}%
\end{pgfscope}%
\begin{pgfscope}%
\pgfpathrectangle{\pgfqpoint{0.537500in}{2.059688in}}{\pgfqpoint{1.505000in}{1.156312in}}%
\pgfusepath{clip}%
\pgfsetrectcap%
\pgfsetroundjoin%
\pgfsetlinewidth{2.158062pt}%
\definecolor{currentstroke}{rgb}{0.478431,0.243137,0.615686}%
\pgfsetstrokecolor{currentstroke}%
\pgfsetdash{}{0pt}%
\pgfpathmoveto{\pgfqpoint{0.910963in}{2.687466in}}%
\pgfpathlineto{\pgfqpoint{0.939822in}{2.669352in}}%
\pgfpathlineto{\pgfqpoint{0.966980in}{2.653412in}}%
\pgfpathlineto{\pgfqpoint{0.992629in}{2.639426in}}%
\pgfpathlineto{\pgfqpoint{1.016935in}{2.627200in}}%
\pgfpathlineto{\pgfqpoint{1.040039in}{2.616567in}}%
\pgfpathlineto{\pgfqpoint{1.062065in}{2.607380in}}%
\pgfpathlineto{\pgfqpoint{1.083121in}{2.599509in}}%
\pgfpathlineto{\pgfqpoint{1.103300in}{2.592838in}}%
\pgfpathlineto{\pgfqpoint{1.122683in}{2.587265in}}%
\pgfpathlineto{\pgfqpoint{1.141344in}{2.582698in}}%
\pgfpathlineto{\pgfqpoint{1.159345in}{2.579054in}}%
\pgfpathlineto{\pgfqpoint{1.176741in}{2.576258in}}%
\pgfpathlineto{\pgfqpoint{1.193581in}{2.574241in}}%
\pgfpathlineto{\pgfqpoint{1.209907in}{2.572943in}}%
\pgfpathlineto{\pgfqpoint{1.225759in}{2.572304in}}%
\pgfpathlineto{\pgfqpoint{1.241170in}{2.572274in}}%
\pgfpathlineto{\pgfqpoint{1.256167in}{2.572804in}}%
\pgfpathlineto{\pgfqpoint{1.270779in}{2.573849in}}%
\pgfpathlineto{\pgfqpoint{1.292021in}{2.576292in}}%
\pgfpathlineto{\pgfqpoint{1.312509in}{2.579672in}}%
\pgfpathlineto{\pgfqpoint{1.332299in}{2.583873in}}%
\pgfpathlineto{\pgfqpoint{1.351434in}{2.588789in}}%
\pgfpathlineto{\pgfqpoint{1.369950in}{2.594323in}}%
\pgfpathlineto{\pgfqpoint{1.387876in}{2.600388in}}%
\pgfpathlineto{\pgfqpoint{1.405238in}{2.606901in}}%
\pgfpathlineto{\pgfqpoint{1.427539in}{2.616156in}}%
\pgfpathlineto{\pgfqpoint{1.448905in}{2.625917in}}%
\pgfpathlineto{\pgfqpoint{1.469364in}{2.636042in}}%
\pgfpathlineto{\pgfqpoint{1.493692in}{2.649020in}}%
\pgfpathlineto{\pgfqpoint{1.516669in}{2.662154in}}%
\pgfpathlineto{\pgfqpoint{1.542496in}{2.677868in}}%
\pgfpathlineto{\pgfqpoint{1.566466in}{2.693288in}}%
\pgfpathlineto{\pgfqpoint{1.592151in}{2.710638in}}%
\pgfpathlineto{\pgfqpoint{1.621712in}{2.731580in}}%
\pgfpathlineto{\pgfqpoint{1.650254in}{2.752706in}}%
\pgfpathlineto{\pgfqpoint{1.681196in}{2.776519in}}%
\pgfpathlineto{\pgfqpoint{1.714144in}{2.802818in}}%
\pgfpathlineto{\pgfqpoint{1.749686in}{2.832166in}}%
\pgfpathlineto{\pgfqpoint{1.787588in}{2.864459in}}%
\pgfpathlineto{\pgfqpoint{1.791626in}{2.867954in}}%
\pgfpathlineto{\pgfqpoint{1.791626in}{2.867954in}}%
\pgfusepath{stroke}%
\end{pgfscope}%
\begin{pgfscope}%
\pgfsetroundcap%
\pgfsetroundjoin%
\pgfsetlinewidth{1.254687pt}%
\definecolor{currentstroke}{rgb}{0.478431,0.243137,0.615686}%
\pgfsetstrokecolor{currentstroke}%
\pgfsetdash{}{0pt}%
\pgfpathmoveto{\pgfqpoint{1.780164in}{2.858058in}}%
\pgfpathquadraticcurveto{\pgfqpoint{1.780286in}{2.858163in}}{\pgfqpoint{1.765691in}{2.845613in}}%
\pgfusepath{stroke}%
\end{pgfscope}%
\begin{pgfscope}%
\pgfsetroundcap%
\pgfsetroundjoin%
\definecolor{currentfill}{rgb}{0.478431,0.243137,0.615686}%
\pgfsetfillcolor{currentfill}%
\pgfsetlinewidth{1.254687pt}%
\definecolor{currentstroke}{rgb}{0.478431,0.243137,0.615686}%
\pgfsetstrokecolor{currentstroke}%
\pgfsetdash{}{0pt}%
\pgfpathmoveto{\pgfqpoint{1.729550in}{2.836518in}}%
\pgfpathlineto{\pgfqpoint{1.765691in}{2.845613in}}%
\pgfpathlineto{\pgfqpoint{1.751282in}{2.811243in}}%
\pgfpathlineto{\pgfqpoint{1.729550in}{2.836518in}}%
\pgfpathclose%
\pgfusepath{stroke,fill}%
\end{pgfscope}%
\begin{pgfscope}%
\pgfpathrectangle{\pgfqpoint{0.537500in}{2.059688in}}{\pgfqpoint{1.505000in}{1.156312in}}%
\pgfusepath{clip}%
\pgfsetrectcap%
\pgfsetroundjoin%
\pgfsetlinewidth{2.158062pt}%
\definecolor{currentstroke}{rgb}{0.878431,0.552941,0.176471}%
\pgfsetstrokecolor{currentstroke}%
\pgfsetdash{}{0pt}%
\pgfpathmoveto{\pgfqpoint{1.044741in}{2.791271in}}%
\pgfpathlineto{\pgfqpoint{1.073428in}{2.782582in}}%
\pgfpathlineto{\pgfqpoint{1.100723in}{2.775561in}}%
\pgfpathlineto{\pgfqpoint{1.126755in}{2.770029in}}%
\pgfpathlineto{\pgfqpoint{1.151640in}{2.765829in}}%
\pgfpathlineto{\pgfqpoint{1.175474in}{2.762822in}}%
\pgfpathlineto{\pgfqpoint{1.198341in}{2.760883in}}%
\pgfpathlineto{\pgfqpoint{1.220316in}{2.759900in}}%
\pgfpathlineto{\pgfqpoint{1.241460in}{2.759772in}}%
\pgfpathlineto{\pgfqpoint{1.261829in}{2.760408in}}%
\pgfpathlineto{\pgfqpoint{1.281472in}{2.761725in}}%
\pgfpathlineto{\pgfqpoint{1.300429in}{2.763647in}}%
\pgfpathlineto{\pgfqpoint{1.318737in}{2.766106in}}%
\pgfpathlineto{\pgfqpoint{1.336430in}{2.769038in}}%
\pgfpathlineto{\pgfqpoint{1.361875in}{2.774202in}}%
\pgfpathlineto{\pgfqpoint{1.386079in}{2.780128in}}%
\pgfpathlineto{\pgfqpoint{1.409114in}{2.786665in}}%
\pgfpathlineto{\pgfqpoint{1.431038in}{2.793677in}}%
\pgfpathlineto{\pgfqpoint{1.451904in}{2.801047in}}%
\pgfpathlineto{\pgfqpoint{1.478156in}{2.811251in}}%
\pgfpathlineto{\pgfqpoint{1.502701in}{2.821698in}}%
\pgfpathlineto{\pgfqpoint{1.525626in}{2.832214in}}%
\pgfpathlineto{\pgfqpoint{1.552122in}{2.845240in}}%
\pgfpathlineto{\pgfqpoint{1.576353in}{2.857936in}}%
\pgfpathlineto{\pgfqpoint{1.602635in}{2.872509in}}%
\pgfpathlineto{\pgfqpoint{1.629737in}{2.888366in}}%
\pgfpathlineto{\pgfqpoint{1.659505in}{2.906698in}}%
\pgfpathlineto{\pgfqpoint{1.689135in}{2.925838in}}%
\pgfpathlineto{\pgfqpoint{1.719723in}{2.946470in}}%
\pgfpathlineto{\pgfqpoint{1.752399in}{2.969422in}}%
\pgfpathlineto{\pgfqpoint{1.786587in}{2.994371in}}%
\pgfpathlineto{\pgfqpoint{1.791662in}{2.998151in}}%
\pgfpathlineto{\pgfqpoint{1.791662in}{2.998151in}}%
\pgfusepath{stroke}%
\end{pgfscope}%
\begin{pgfscope}%
\pgfsetroundcap%
\pgfsetroundjoin%
\pgfsetlinewidth{1.254687pt}%
\definecolor{currentstroke}{rgb}{0.878431,0.552941,0.176471}%
\pgfsetstrokecolor{currentstroke}%
\pgfsetdash{}{0pt}%
\pgfpathmoveto{\pgfqpoint{1.788674in}{2.995923in}}%
\pgfpathquadraticcurveto{\pgfqpoint{1.788711in}{2.995951in}}{\pgfqpoint{1.773180in}{2.984386in}}%
\pgfusepath{stroke}%
\end{pgfscope}%
\begin{pgfscope}%
\pgfsetroundcap%
\pgfsetroundjoin%
\definecolor{currentfill}{rgb}{0.878431,0.552941,0.176471}%
\pgfsetfillcolor{currentfill}%
\pgfsetlinewidth{1.254687pt}%
\definecolor{currentstroke}{rgb}{0.878431,0.552941,0.176471}%
\pgfsetstrokecolor{currentstroke}%
\pgfsetdash{}{0pt}%
\pgfpathmoveto{\pgfqpoint{1.736491in}{2.977846in}}%
\pgfpathlineto{\pgfqpoint{1.773180in}{2.984386in}}%
\pgfpathlineto{\pgfqpoint{1.756398in}{2.951111in}}%
\pgfpathlineto{\pgfqpoint{1.736491in}{2.977846in}}%
\pgfpathclose%
\pgfusepath{stroke,fill}%
\end{pgfscope}%
\begin{pgfscope}%
\definecolor{textcolor}{rgb}{0.372549,0.372549,0.372549}%
\pgfsetstrokecolor{textcolor}%
\pgfsetfillcolor{textcolor}%
\pgftext[x=0.576519in,y=3.122245in,left,base]{\color{textcolor}{\rmfamily\fontsize{6.000000}{7.200000}\selectfont\catcode`\^=\active\def^{\ifmmode\sp\else\^{}\fi}\catcode`\%=\active\def%{\%}leaves $s^2-r^2=\alpha^2$}}%
\end{pgfscope}%
\begin{pgfscope}%
\definecolor{textcolor}{rgb}{0.125490,0.498039,0.321569}%
\pgfsetstrokecolor{textcolor}%
\pgfsetfillcolor{textcolor}%
\pgftext[x=1.811176in,y=3.028490in,left,base]{\color{textcolor}{\rmfamily\fontsize{6.000000}{7.200000}\selectfont\catcode`\^=\active\def^{\ifmmode\sp\else\^{}\fi}\catcode`\%=\active\def%{\%}$r=\beta$}}%
\end{pgfscope}%
\begin{pgfscope}%
\definecolor{textcolor}{rgb}{0.000000,0.000000,0.000000}%
\pgfsetstrokecolor{textcolor}%
\pgfsetfillcolor{textcolor}%
\pgftext[x=1.290000in,y=3.243778in,,base]{\color{textcolor}{\rmfamily\fontsize{8.000000}{9.600000}\selectfont\catcode`\^=\active\def^{\ifmmode\sp\else\^{}\fi}\catcode`\%=\active\def%{\%}(a) Equal variances}}%
\end{pgfscope}%
\begin{pgfscope}%
\pgfpathrectangle{\pgfqpoint{0.537500in}{2.059688in}}{\pgfqpoint{1.505000in}{1.156312in}}%
\pgfusepath{clip}%
\pgfsetbuttcap%
\pgfsetroundjoin%
\definecolor{currentfill}{rgb}{0.141176,0.345098,0.650980}%
\pgfsetfillcolor{currentfill}%
\pgfsetlinewidth{1.003750pt}%
\definecolor{currentstroke}{rgb}{0.141176,0.345098,0.650980}%
\pgfsetstrokecolor{currentstroke}%
\pgfsetdash{}{0pt}%
\pgfsys@defobject{currentmarker}{\pgfqpoint{-0.034021in}{-0.034021in}}{\pgfqpoint{0.034021in}{0.034021in}}{%
\pgfpathmoveto{\pgfqpoint{0.000000in}{-0.034021in}}%
\pgfpathcurveto{\pgfqpoint{0.009022in}{-0.034021in}}{\pgfqpoint{0.017676in}{-0.030436in}}{\pgfqpoint{0.024056in}{-0.024056in}}%
\pgfpathcurveto{\pgfqpoint{0.030436in}{-0.017676in}}{\pgfqpoint{0.034021in}{-0.009022in}}{\pgfqpoint{0.034021in}{0.000000in}}%
\pgfpathcurveto{\pgfqpoint{0.034021in}{0.009022in}}{\pgfqpoint{0.030436in}{0.017676in}}{\pgfqpoint{0.024056in}{0.024056in}}%
\pgfpathcurveto{\pgfqpoint{0.017676in}{0.030436in}}{\pgfqpoint{0.009022in}{0.034021in}}{\pgfqpoint{0.000000in}{0.034021in}}%
\pgfpathcurveto{\pgfqpoint{-0.009022in}{0.034021in}}{\pgfqpoint{-0.017676in}{0.030436in}}{\pgfqpoint{-0.024056in}{0.024056in}}%
\pgfpathcurveto{\pgfqpoint{-0.030436in}{0.017676in}}{\pgfqpoint{-0.034021in}{0.009022in}}{\pgfqpoint{-0.034021in}{0.000000in}}%
\pgfpathcurveto{\pgfqpoint{-0.034021in}{-0.009022in}}{\pgfqpoint{-0.030436in}{-0.017676in}}{\pgfqpoint{-0.024056in}{-0.024056in}}%
\pgfpathcurveto{\pgfqpoint{-0.017676in}{-0.030436in}}{\pgfqpoint{-0.009022in}{-0.034021in}}{\pgfqpoint{0.000000in}{-0.034021in}}%
\pgfpathlineto{\pgfqpoint{0.000000in}{-0.034021in}}%
\pgfpathclose%
\pgfusepath{stroke,fill}%
}%
\begin{pgfscope}%
\pgfsys@transformshift{0.648981in}{2.732899in}%
\pgfsys@useobject{currentmarker}{}%
\end{pgfscope}%
\end{pgfscope}%
\begin{pgfscope}%
\pgfpathrectangle{\pgfqpoint{0.537500in}{2.059688in}}{\pgfqpoint{1.505000in}{1.156312in}}%
\pgfusepath{clip}%
\pgfsetbuttcap%
\pgfsetroundjoin%
\definecolor{currentfill}{rgb}{0.141176,0.345098,0.650980}%
\pgfsetfillcolor{currentfill}%
\pgfsetlinewidth{1.003750pt}%
\definecolor{currentstroke}{rgb}{0.141176,0.345098,0.650980}%
\pgfsetstrokecolor{currentstroke}%
\pgfsetdash{}{0pt}%
\pgfsys@defobject{currentmarker}{\pgfqpoint{-0.036747in}{-0.036747in}}{\pgfqpoint{0.036747in}{0.036747in}}{%
\pgfpathmoveto{\pgfqpoint{-0.036747in}{-0.036747in}}%
\pgfpathlineto{\pgfqpoint{0.036747in}{-0.036747in}}%
\pgfpathlineto{\pgfqpoint{0.036747in}{0.036747in}}%
\pgfpathlineto{\pgfqpoint{-0.036747in}{0.036747in}}%
\pgfpathlineto{\pgfqpoint{-0.036747in}{-0.036747in}}%
\pgfpathclose%
\pgfusepath{stroke,fill}%
}%
\begin{pgfscope}%
\pgfsys@transformshift{1.787975in}{2.698446in}%
\pgfsys@useobject{currentmarker}{}%
\end{pgfscope}%
\end{pgfscope}%
\begin{pgfscope}%
\pgfpathrectangle{\pgfqpoint{0.537500in}{2.059688in}}{\pgfqpoint{1.505000in}{1.156312in}}%
\pgfusepath{clip}%
\pgfsetbuttcap%
\pgfsetroundjoin%
\definecolor{currentfill}{rgb}{0.819608,0.286275,0.356863}%
\pgfsetfillcolor{currentfill}%
\pgfsetlinewidth{1.003750pt}%
\definecolor{currentstroke}{rgb}{0.819608,0.286275,0.356863}%
\pgfsetstrokecolor{currentstroke}%
\pgfsetdash{}{0pt}%
\pgfsys@defobject{currentmarker}{\pgfqpoint{-0.034021in}{-0.034021in}}{\pgfqpoint{0.034021in}{0.034021in}}{%
\pgfpathmoveto{\pgfqpoint{0.000000in}{-0.034021in}}%
\pgfpathcurveto{\pgfqpoint{0.009022in}{-0.034021in}}{\pgfqpoint{0.017676in}{-0.030436in}}{\pgfqpoint{0.024056in}{-0.024056in}}%
\pgfpathcurveto{\pgfqpoint{0.030436in}{-0.017676in}}{\pgfqpoint{0.034021in}{-0.009022in}}{\pgfqpoint{0.034021in}{0.000000in}}%
\pgfpathcurveto{\pgfqpoint{0.034021in}{0.009022in}}{\pgfqpoint{0.030436in}{0.017676in}}{\pgfqpoint{0.024056in}{0.024056in}}%
\pgfpathcurveto{\pgfqpoint{0.017676in}{0.030436in}}{\pgfqpoint{0.009022in}{0.034021in}}{\pgfqpoint{0.000000in}{0.034021in}}%
\pgfpathcurveto{\pgfqpoint{-0.009022in}{0.034021in}}{\pgfqpoint{-0.017676in}{0.030436in}}{\pgfqpoint{-0.024056in}{0.024056in}}%
\pgfpathcurveto{\pgfqpoint{-0.030436in}{0.017676in}}{\pgfqpoint{-0.034021in}{0.009022in}}{\pgfqpoint{-0.034021in}{0.000000in}}%
\pgfpathcurveto{\pgfqpoint{-0.034021in}{-0.009022in}}{\pgfqpoint{-0.030436in}{-0.017676in}}{\pgfqpoint{-0.024056in}{-0.024056in}}%
\pgfpathcurveto{\pgfqpoint{-0.017676in}{-0.030436in}}{\pgfqpoint{-0.009022in}{-0.034021in}}{\pgfqpoint{0.000000in}{-0.034021in}}%
\pgfpathlineto{\pgfqpoint{0.000000in}{-0.034021in}}%
\pgfpathclose%
\pgfusepath{stroke,fill}%
}%
\begin{pgfscope}%
\pgfsys@transformshift{0.721444in}{2.681840in}%
\pgfsys@useobject{currentmarker}{}%
\end{pgfscope}%
\end{pgfscope}%
\begin{pgfscope}%
\pgfpathrectangle{\pgfqpoint{0.537500in}{2.059688in}}{\pgfqpoint{1.505000in}{1.156312in}}%
\pgfusepath{clip}%
\pgfsetbuttcap%
\pgfsetroundjoin%
\definecolor{currentfill}{rgb}{0.819608,0.286275,0.356863}%
\pgfsetfillcolor{currentfill}%
\pgfsetlinewidth{1.003750pt}%
\definecolor{currentstroke}{rgb}{0.819608,0.286275,0.356863}%
\pgfsetstrokecolor{currentstroke}%
\pgfsetdash{}{0pt}%
\pgfsys@defobject{currentmarker}{\pgfqpoint{-0.036747in}{-0.036747in}}{\pgfqpoint{0.036747in}{0.036747in}}{%
\pgfpathmoveto{\pgfqpoint{-0.036747in}{-0.036747in}}%
\pgfpathlineto{\pgfqpoint{0.036747in}{-0.036747in}}%
\pgfpathlineto{\pgfqpoint{0.036747in}{0.036747in}}%
\pgfpathlineto{\pgfqpoint{-0.036747in}{0.036747in}}%
\pgfpathlineto{\pgfqpoint{-0.036747in}{-0.036747in}}%
\pgfpathclose%
\pgfusepath{stroke,fill}%
}%
\begin{pgfscope}%
\pgfsys@transformshift{1.790609in}{2.727219in}%
\pgfsys@useobject{currentmarker}{}%
\end{pgfscope}%
\end{pgfscope}%
\begin{pgfscope}%
\pgfpathrectangle{\pgfqpoint{0.537500in}{2.059688in}}{\pgfqpoint{1.505000in}{1.156312in}}%
\pgfusepath{clip}%
\pgfsetbuttcap%
\pgfsetroundjoin%
\definecolor{currentfill}{rgb}{0.000000,0.529412,0.423529}%
\pgfsetfillcolor{currentfill}%
\pgfsetlinewidth{1.003750pt}%
\definecolor{currentstroke}{rgb}{0.000000,0.529412,0.423529}%
\pgfsetstrokecolor{currentstroke}%
\pgfsetdash{}{0pt}%
\pgfsys@defobject{currentmarker}{\pgfqpoint{-0.034021in}{-0.034021in}}{\pgfqpoint{0.034021in}{0.034021in}}{%
\pgfpathmoveto{\pgfqpoint{0.000000in}{-0.034021in}}%
\pgfpathcurveto{\pgfqpoint{0.009022in}{-0.034021in}}{\pgfqpoint{0.017676in}{-0.030436in}}{\pgfqpoint{0.024056in}{-0.024056in}}%
\pgfpathcurveto{\pgfqpoint{0.030436in}{-0.017676in}}{\pgfqpoint{0.034021in}{-0.009022in}}{\pgfqpoint{0.034021in}{0.000000in}}%
\pgfpathcurveto{\pgfqpoint{0.034021in}{0.009022in}}{\pgfqpoint{0.030436in}{0.017676in}}{\pgfqpoint{0.024056in}{0.024056in}}%
\pgfpathcurveto{\pgfqpoint{0.017676in}{0.030436in}}{\pgfqpoint{0.009022in}{0.034021in}}{\pgfqpoint{0.000000in}{0.034021in}}%
\pgfpathcurveto{\pgfqpoint{-0.009022in}{0.034021in}}{\pgfqpoint{-0.017676in}{0.030436in}}{\pgfqpoint{-0.024056in}{0.024056in}}%
\pgfpathcurveto{\pgfqpoint{-0.030436in}{0.017676in}}{\pgfqpoint{-0.034021in}{0.009022in}}{\pgfqpoint{-0.034021in}{0.000000in}}%
\pgfpathcurveto{\pgfqpoint{-0.034021in}{-0.009022in}}{\pgfqpoint{-0.030436in}{-0.017676in}}{\pgfqpoint{-0.024056in}{-0.024056in}}%
\pgfpathcurveto{\pgfqpoint{-0.017676in}{-0.030436in}}{\pgfqpoint{-0.009022in}{-0.034021in}}{\pgfqpoint{0.000000in}{-0.034021in}}%
\pgfpathlineto{\pgfqpoint{0.000000in}{-0.034021in}}%
\pgfpathclose%
\pgfusepath{stroke,fill}%
}%
\begin{pgfscope}%
\pgfsys@transformshift{0.810630in}{2.657241in}%
\pgfsys@useobject{currentmarker}{}%
\end{pgfscope}%
\end{pgfscope}%
\begin{pgfscope}%
\pgfpathrectangle{\pgfqpoint{0.537500in}{2.059688in}}{\pgfqpoint{1.505000in}{1.156312in}}%
\pgfusepath{clip}%
\pgfsetbuttcap%
\pgfsetroundjoin%
\definecolor{currentfill}{rgb}{0.000000,0.529412,0.423529}%
\pgfsetfillcolor{currentfill}%
\pgfsetlinewidth{1.003750pt}%
\definecolor{currentstroke}{rgb}{0.000000,0.529412,0.423529}%
\pgfsetstrokecolor{currentstroke}%
\pgfsetdash{}{0pt}%
\pgfsys@defobject{currentmarker}{\pgfqpoint{-0.036747in}{-0.036747in}}{\pgfqpoint{0.036747in}{0.036747in}}{%
\pgfpathmoveto{\pgfqpoint{-0.036747in}{-0.036747in}}%
\pgfpathlineto{\pgfqpoint{0.036747in}{-0.036747in}}%
\pgfpathlineto{\pgfqpoint{0.036747in}{0.036747in}}%
\pgfpathlineto{\pgfqpoint{-0.036747in}{0.036747in}}%
\pgfpathlineto{\pgfqpoint{-0.036747in}{-0.036747in}}%
\pgfpathclose%
\pgfusepath{stroke,fill}%
}%
\begin{pgfscope}%
\pgfsys@transformshift{1.791433in}{2.782018in}%
\pgfsys@useobject{currentmarker}{}%
\end{pgfscope}%
\end{pgfscope}%
\begin{pgfscope}%
\pgfpathrectangle{\pgfqpoint{0.537500in}{2.059688in}}{\pgfqpoint{1.505000in}{1.156312in}}%
\pgfusepath{clip}%
\pgfsetbuttcap%
\pgfsetroundjoin%
\definecolor{currentfill}{rgb}{0.478431,0.243137,0.615686}%
\pgfsetfillcolor{currentfill}%
\pgfsetlinewidth{1.003750pt}%
\definecolor{currentstroke}{rgb}{0.478431,0.243137,0.615686}%
\pgfsetstrokecolor{currentstroke}%
\pgfsetdash{}{0pt}%
\pgfsys@defobject{currentmarker}{\pgfqpoint{-0.034021in}{-0.034021in}}{\pgfqpoint{0.034021in}{0.034021in}}{%
\pgfpathmoveto{\pgfqpoint{0.000000in}{-0.034021in}}%
\pgfpathcurveto{\pgfqpoint{0.009022in}{-0.034021in}}{\pgfqpoint{0.017676in}{-0.030436in}}{\pgfqpoint{0.024056in}{-0.024056in}}%
\pgfpathcurveto{\pgfqpoint{0.030436in}{-0.017676in}}{\pgfqpoint{0.034021in}{-0.009022in}}{\pgfqpoint{0.034021in}{0.000000in}}%
\pgfpathcurveto{\pgfqpoint{0.034021in}{0.009022in}}{\pgfqpoint{0.030436in}{0.017676in}}{\pgfqpoint{0.024056in}{0.024056in}}%
\pgfpathcurveto{\pgfqpoint{0.017676in}{0.030436in}}{\pgfqpoint{0.009022in}{0.034021in}}{\pgfqpoint{0.000000in}{0.034021in}}%
\pgfpathcurveto{\pgfqpoint{-0.009022in}{0.034021in}}{\pgfqpoint{-0.017676in}{0.030436in}}{\pgfqpoint{-0.024056in}{0.024056in}}%
\pgfpathcurveto{\pgfqpoint{-0.030436in}{0.017676in}}{\pgfqpoint{-0.034021in}{0.009022in}}{\pgfqpoint{-0.034021in}{0.000000in}}%
\pgfpathcurveto{\pgfqpoint{-0.034021in}{-0.009022in}}{\pgfqpoint{-0.030436in}{-0.017676in}}{\pgfqpoint{-0.024056in}{-0.024056in}}%
\pgfpathcurveto{\pgfqpoint{-0.017676in}{-0.030436in}}{\pgfqpoint{-0.009022in}{-0.034021in}}{\pgfqpoint{0.000000in}{-0.034021in}}%
\pgfpathlineto{\pgfqpoint{0.000000in}{-0.034021in}}%
\pgfpathclose%
\pgfusepath{stroke,fill}%
}%
\begin{pgfscope}%
\pgfsys@transformshift{0.910963in}{2.687466in}%
\pgfsys@useobject{currentmarker}{}%
\end{pgfscope}%
\end{pgfscope}%
\begin{pgfscope}%
\pgfpathrectangle{\pgfqpoint{0.537500in}{2.059688in}}{\pgfqpoint{1.505000in}{1.156312in}}%
\pgfusepath{clip}%
\pgfsetbuttcap%
\pgfsetroundjoin%
\definecolor{currentfill}{rgb}{0.478431,0.243137,0.615686}%
\pgfsetfillcolor{currentfill}%
\pgfsetlinewidth{1.003750pt}%
\definecolor{currentstroke}{rgb}{0.478431,0.243137,0.615686}%
\pgfsetstrokecolor{currentstroke}%
\pgfsetdash{}{0pt}%
\pgfsys@defobject{currentmarker}{\pgfqpoint{-0.036747in}{-0.036747in}}{\pgfqpoint{0.036747in}{0.036747in}}{%
\pgfpathmoveto{\pgfqpoint{-0.036747in}{-0.036747in}}%
\pgfpathlineto{\pgfqpoint{0.036747in}{-0.036747in}}%
\pgfpathlineto{\pgfqpoint{0.036747in}{0.036747in}}%
\pgfpathlineto{\pgfqpoint{-0.036747in}{0.036747in}}%
\pgfpathlineto{\pgfqpoint{-0.036747in}{-0.036747in}}%
\pgfpathclose%
\pgfusepath{stroke,fill}%
}%
\begin{pgfscope}%
\pgfsys@transformshift{1.791626in}{2.867954in}%
\pgfsys@useobject{currentmarker}{}%
\end{pgfscope}%
\end{pgfscope}%
\begin{pgfscope}%
\pgfpathrectangle{\pgfqpoint{0.537500in}{2.059688in}}{\pgfqpoint{1.505000in}{1.156312in}}%
\pgfusepath{clip}%
\pgfsetbuttcap%
\pgfsetroundjoin%
\definecolor{currentfill}{rgb}{0.878431,0.552941,0.176471}%
\pgfsetfillcolor{currentfill}%
\pgfsetlinewidth{1.003750pt}%
\definecolor{currentstroke}{rgb}{0.878431,0.552941,0.176471}%
\pgfsetstrokecolor{currentstroke}%
\pgfsetdash{}{0pt}%
\pgfsys@defobject{currentmarker}{\pgfqpoint{-0.034021in}{-0.034021in}}{\pgfqpoint{0.034021in}{0.034021in}}{%
\pgfpathmoveto{\pgfqpoint{0.000000in}{-0.034021in}}%
\pgfpathcurveto{\pgfqpoint{0.009022in}{-0.034021in}}{\pgfqpoint{0.017676in}{-0.030436in}}{\pgfqpoint{0.024056in}{-0.024056in}}%
\pgfpathcurveto{\pgfqpoint{0.030436in}{-0.017676in}}{\pgfqpoint{0.034021in}{-0.009022in}}{\pgfqpoint{0.034021in}{0.000000in}}%
\pgfpathcurveto{\pgfqpoint{0.034021in}{0.009022in}}{\pgfqpoint{0.030436in}{0.017676in}}{\pgfqpoint{0.024056in}{0.024056in}}%
\pgfpathcurveto{\pgfqpoint{0.017676in}{0.030436in}}{\pgfqpoint{0.009022in}{0.034021in}}{\pgfqpoint{0.000000in}{0.034021in}}%
\pgfpathcurveto{\pgfqpoint{-0.009022in}{0.034021in}}{\pgfqpoint{-0.017676in}{0.030436in}}{\pgfqpoint{-0.024056in}{0.024056in}}%
\pgfpathcurveto{\pgfqpoint{-0.030436in}{0.017676in}}{\pgfqpoint{-0.034021in}{0.009022in}}{\pgfqpoint{-0.034021in}{0.000000in}}%
\pgfpathcurveto{\pgfqpoint{-0.034021in}{-0.009022in}}{\pgfqpoint{-0.030436in}{-0.017676in}}{\pgfqpoint{-0.024056in}{-0.024056in}}%
\pgfpathcurveto{\pgfqpoint{-0.017676in}{-0.030436in}}{\pgfqpoint{-0.009022in}{-0.034021in}}{\pgfqpoint{0.000000in}{-0.034021in}}%
\pgfpathlineto{\pgfqpoint{0.000000in}{-0.034021in}}%
\pgfpathclose%
\pgfusepath{stroke,fill}%
}%
\begin{pgfscope}%
\pgfsys@transformshift{1.044741in}{2.791271in}%
\pgfsys@useobject{currentmarker}{}%
\end{pgfscope}%
\end{pgfscope}%
\begin{pgfscope}%
\pgfpathrectangle{\pgfqpoint{0.537500in}{2.059688in}}{\pgfqpoint{1.505000in}{1.156312in}}%
\pgfusepath{clip}%
\pgfsetbuttcap%
\pgfsetroundjoin%
\definecolor{currentfill}{rgb}{0.878431,0.552941,0.176471}%
\pgfsetfillcolor{currentfill}%
\pgfsetlinewidth{1.003750pt}%
\definecolor{currentstroke}{rgb}{0.878431,0.552941,0.176471}%
\pgfsetstrokecolor{currentstroke}%
\pgfsetdash{}{0pt}%
\pgfsys@defobject{currentmarker}{\pgfqpoint{-0.036747in}{-0.036747in}}{\pgfqpoint{0.036747in}{0.036747in}}{%
\pgfpathmoveto{\pgfqpoint{-0.036747in}{-0.036747in}}%
\pgfpathlineto{\pgfqpoint{0.036747in}{-0.036747in}}%
\pgfpathlineto{\pgfqpoint{0.036747in}{0.036747in}}%
\pgfpathlineto{\pgfqpoint{-0.036747in}{0.036747in}}%
\pgfpathlineto{\pgfqpoint{-0.036747in}{-0.036747in}}%
\pgfpathclose%
\pgfusepath{stroke,fill}%
}%
\begin{pgfscope}%
\pgfsys@transformshift{1.791662in}{2.998151in}%
\pgfsys@useobject{currentmarker}{}%
\end{pgfscope}%
\end{pgfscope}%
\begin{pgfscope}%
\pgfpathrectangle{\pgfqpoint{0.537500in}{2.059688in}}{\pgfqpoint{1.505000in}{1.156312in}}%
\pgfusepath{clip}%
\pgfsetbuttcap%
\pgfsetroundjoin%
\definecolor{currentfill}{rgb}{0.168627,0.611765,0.411765}%
\pgfsetfillcolor{currentfill}%
\pgfsetlinewidth{1.003750pt}%
\definecolor{currentstroke}{rgb}{0.168627,0.611765,0.411765}%
\pgfsetstrokecolor{currentstroke}%
\pgfsetdash{}{0pt}%
\pgfsys@defobject{currentmarker}{\pgfqpoint{-0.034021in}{-0.034021in}}{\pgfqpoint{0.034021in}{0.034021in}}{%
\pgfpathmoveto{\pgfqpoint{0.000000in}{-0.034021in}}%
\pgfpathcurveto{\pgfqpoint{0.009022in}{-0.034021in}}{\pgfqpoint{0.017676in}{-0.030436in}}{\pgfqpoint{0.024056in}{-0.024056in}}%
\pgfpathcurveto{\pgfqpoint{0.030436in}{-0.017676in}}{\pgfqpoint{0.034021in}{-0.009022in}}{\pgfqpoint{0.034021in}{0.000000in}}%
\pgfpathcurveto{\pgfqpoint{0.034021in}{0.009022in}}{\pgfqpoint{0.030436in}{0.017676in}}{\pgfqpoint{0.024056in}{0.024056in}}%
\pgfpathcurveto{\pgfqpoint{0.017676in}{0.030436in}}{\pgfqpoint{0.009022in}{0.034021in}}{\pgfqpoint{0.000000in}{0.034021in}}%
\pgfpathcurveto{\pgfqpoint{-0.009022in}{0.034021in}}{\pgfqpoint{-0.017676in}{0.030436in}}{\pgfqpoint{-0.024056in}{0.024056in}}%
\pgfpathcurveto{\pgfqpoint{-0.030436in}{0.017676in}}{\pgfqpoint{-0.034021in}{0.009022in}}{\pgfqpoint{-0.034021in}{0.000000in}}%
\pgfpathcurveto{\pgfqpoint{-0.034021in}{-0.009022in}}{\pgfqpoint{-0.030436in}{-0.017676in}}{\pgfqpoint{-0.024056in}{-0.024056in}}%
\pgfpathcurveto{\pgfqpoint{-0.017676in}{-0.030436in}}{\pgfqpoint{-0.009022in}{-0.034021in}}{\pgfqpoint{0.000000in}{-0.034021in}}%
\pgfpathlineto{\pgfqpoint{0.000000in}{-0.034021in}}%
\pgfpathclose%
\pgfusepath{stroke,fill}%
}%
\begin{pgfscope}%
\pgfsys@transformshift{1.791667in}{2.684722in}%
\pgfsys@useobject{currentmarker}{}%
\end{pgfscope}%
\end{pgfscope}%
\begin{pgfscope}%
\pgfsetbuttcap%
\pgfsetmiterjoin%
\definecolor{currentfill}{rgb}{1.000000,1.000000,1.000000}%
\pgfsetfillcolor{currentfill}%
\pgfsetlinewidth{0.000000pt}%
\definecolor{currentstroke}{rgb}{0.000000,0.000000,0.000000}%
\pgfsetstrokecolor{currentstroke}%
\pgfsetstrokeopacity{0.000000}%
\pgfsetdash{}{0pt}%
\pgfpathmoveto{\pgfqpoint{0.407592in}{0.222904in}}%
\pgfpathlineto{\pgfqpoint{1.818754in}{0.222904in}}%
\pgfpathlineto{\pgfqpoint{1.818754in}{1.634067in}}%
\pgfpathlineto{\pgfqpoint{0.407592in}{1.634067in}}%
\pgfpathlineto{\pgfqpoint{0.407592in}{0.222904in}}%
\pgfpathclose%
\pgfusepath{fill}%
\end{pgfscope}%
\begin{pgfscope}%
\pgfsetbuttcap%
\pgfsetmiterjoin%
\pgfsetlinewidth{0.000000pt}%
\definecolor{currentstroke}{rgb}{0.950000,0.950000,0.950000}%
\pgfsetstrokecolor{currentstroke}%
\pgfsetstrokeopacity{0.000000}%
\pgfsetdash{}{0pt}%
\pgfpathmoveto{\pgfqpoint{0.418841in}{0.578914in}}%
\pgfpathlineto{\pgfqpoint{0.882898in}{0.862989in}}%
\pgfpathlineto{\pgfqpoint{0.873607in}{1.529400in}}%
\pgfpathlineto{\pgfqpoint{0.389895in}{1.293091in}}%
\pgfusepath{}%
\end{pgfscope}%
\begin{pgfscope}%
\pgfsetbuttcap%
\pgfsetmiterjoin%
\pgfsetlinewidth{0.000000pt}%
\definecolor{currentstroke}{rgb}{0.900000,0.900000,0.900000}%
\pgfsetstrokecolor{currentstroke}%
\pgfsetstrokeopacity{0.000000}%
\pgfsetdash{}{0pt}%
\pgfpathmoveto{\pgfqpoint{0.882898in}{0.862989in}}%
\pgfpathlineto{\pgfqpoint{1.837675in}{0.619202in}}%
\pgfpathlineto{\pgfqpoint{1.865965in}{1.326696in}}%
\pgfpathlineto{\pgfqpoint{0.873607in}{1.529400in}}%
\pgfusepath{}%
\end{pgfscope}%
\begin{pgfscope}%
\pgfsetbuttcap%
\pgfsetmiterjoin%
\pgfsetlinewidth{0.000000pt}%
\definecolor{currentstroke}{rgb}{0.925000,0.925000,0.925000}%
\pgfsetstrokecolor{currentstroke}%
\pgfsetstrokeopacity{0.000000}%
\pgfsetdash{}{0pt}%
\pgfpathmoveto{\pgfqpoint{0.418841in}{0.578914in}}%
\pgfpathlineto{\pgfqpoint{1.424602in}{0.289843in}}%
\pgfpathlineto{\pgfqpoint{1.837675in}{0.619202in}}%
\pgfpathlineto{\pgfqpoint{0.882898in}{0.862989in}}%
\pgfusepath{}%
\end{pgfscope}%
\begin{pgfscope}%
\pgfsetrectcap%
\pgfsetroundjoin%
\pgfsetlinewidth{0.803000pt}%
\definecolor{currentstroke}{rgb}{0.000000,0.000000,0.000000}%
\pgfsetstrokecolor{currentstroke}%
\pgfsetdash{}{0pt}%
\pgfpathmoveto{\pgfqpoint{0.418841in}{0.578914in}}%
\pgfpathlineto{\pgfqpoint{1.424602in}{0.289843in}}%
\pgfusepath{stroke}%
\end{pgfscope}%
\begin{pgfscope}%
\pgfsetrectcap%
\pgfsetroundjoin%
\pgfsetlinewidth{0.803000pt}%
\definecolor{currentstroke}{rgb}{0.000000,0.000000,0.000000}%
\pgfsetstrokecolor{currentstroke}%
\pgfsetdash{}{0pt}%
\pgfpathmoveto{\pgfqpoint{0.500034in}{0.559223in}}%
\pgfpathlineto{\pgfqpoint{0.488021in}{0.551726in}}%
\pgfusepath{stroke}%
\end{pgfscope}%
\begin{pgfscope}%
\definecolor{textcolor}{rgb}{0.000000,0.000000,0.000000}%
\pgfsetstrokecolor{textcolor}%
\pgfsetfillcolor{textcolor}%
\pgftext[x=0.432780in,y=0.423379in,,top]{\color{textcolor}{\rmfamily\fontsize{5.000000}{6.000000}\selectfont\catcode`\^=\active\def^{\ifmmode\sp\else\^{}\fi}\catcode`\%=\active\def%{\%}\ensuremath{-}1.0}}%
\end{pgfscope}%
\begin{pgfscope}%
\pgfsetrectcap%
\pgfsetroundjoin%
\pgfsetlinewidth{0.803000pt}%
\definecolor{currentstroke}{rgb}{0.000000,0.000000,0.000000}%
\pgfsetstrokecolor{currentstroke}%
\pgfsetdash{}{0pt}%
\pgfpathmoveto{\pgfqpoint{0.678989in}{0.507840in}}%
\pgfpathlineto{\pgfqpoint{0.667181in}{0.500129in}}%
\pgfusepath{stroke}%
\end{pgfscope}%
\begin{pgfscope}%
\definecolor{textcolor}{rgb}{0.000000,0.000000,0.000000}%
\pgfsetstrokecolor{textcolor}%
\pgfsetfillcolor{textcolor}%
\pgftext[x=0.612032in,y=0.369697in,,top]{\color{textcolor}{\rmfamily\fontsize{5.000000}{6.000000}\selectfont\catcode`\^=\active\def^{\ifmmode\sp\else\^{}\fi}\catcode`\%=\active\def%{\%}\ensuremath{-}0.5}}%
\end{pgfscope}%
\begin{pgfscope}%
\pgfsetrectcap%
\pgfsetroundjoin%
\pgfsetlinewidth{0.803000pt}%
\definecolor{currentstroke}{rgb}{0.000000,0.000000,0.000000}%
\pgfsetstrokecolor{currentstroke}%
\pgfsetdash{}{0pt}%
\pgfpathmoveto{\pgfqpoint{0.863089in}{0.454980in}}%
\pgfpathlineto{\pgfqpoint{0.851503in}{0.447045in}}%
\pgfusepath{stroke}%
\end{pgfscope}%
\begin{pgfscope}%
\definecolor{textcolor}{rgb}{0.000000,0.000000,0.000000}%
\pgfsetstrokecolor{textcolor}%
\pgfsetfillcolor{textcolor}%
\pgftext[x=0.796474in,y=0.314461in,,top]{\color{textcolor}{\rmfamily\fontsize{5.000000}{6.000000}\selectfont\catcode`\^=\active\def^{\ifmmode\sp\else\^{}\fi}\catcode`\%=\active\def%{\%}0.0}}%
\end{pgfscope}%
\begin{pgfscope}%
\pgfsetrectcap%
\pgfsetroundjoin%
\pgfsetlinewidth{0.803000pt}%
\definecolor{currentstroke}{rgb}{0.000000,0.000000,0.000000}%
\pgfsetstrokecolor{currentstroke}%
\pgfsetdash{}{0pt}%
\pgfpathmoveto{\pgfqpoint{1.052559in}{0.400579in}}%
\pgfpathlineto{\pgfqpoint{1.041212in}{0.392410in}}%
\pgfusepath{stroke}%
\end{pgfscope}%
\begin{pgfscope}%
\definecolor{textcolor}{rgb}{0.000000,0.000000,0.000000}%
\pgfsetstrokecolor{textcolor}%
\pgfsetfillcolor{textcolor}%
\pgftext[x=0.986334in,y=0.257602in,,top]{\color{textcolor}{\rmfamily\fontsize{5.000000}{6.000000}\selectfont\catcode`\^=\active\def^{\ifmmode\sp\else\^{}\fi}\catcode`\%=\active\def%{\%}0.5}}%
\end{pgfscope}%
\begin{pgfscope}%
\pgfsetrectcap%
\pgfsetroundjoin%
\pgfsetlinewidth{0.803000pt}%
\definecolor{currentstroke}{rgb}{0.000000,0.000000,0.000000}%
\pgfsetstrokecolor{currentstroke}%
\pgfsetdash{}{0pt}%
\pgfpathmoveto{\pgfqpoint{1.247636in}{0.344567in}}%
\pgfpathlineto{\pgfqpoint{1.236548in}{0.336154in}}%
\pgfusepath{stroke}%
\end{pgfscope}%
\begin{pgfscope}%
\definecolor{textcolor}{rgb}{0.000000,0.000000,0.000000}%
\pgfsetstrokecolor{textcolor}%
\pgfsetfillcolor{textcolor}%
\pgftext[x=1.181855in,y=0.199048in,,top]{\color{textcolor}{\rmfamily\fontsize{5.000000}{6.000000}\selectfont\catcode`\^=\active\def^{\ifmmode\sp\else\^{}\fi}\catcode`\%=\active\def%{\%}1.0}}%
\end{pgfscope}%
\begin{pgfscope}%
\definecolor{textcolor}{rgb}{0.000000,0.000000,0.000000}%
\pgfsetstrokecolor{textcolor}%
\pgfsetfillcolor{textcolor}%
\pgftext[x=0.757857in,y=0.120925in,,]{\color{textcolor}{\rmfamily\fontsize{6.000000}{7.200000}\selectfont\catcode`\^=\active\def^{\ifmmode\sp\else\^{}\fi}\catcode`\%=\active\def%{\%}$r$}}%
\end{pgfscope}%
\begin{pgfscope}%
\pgfsetrectcap%
\pgfsetroundjoin%
\pgfsetlinewidth{0.803000pt}%
\definecolor{currentstroke}{rgb}{0.000000,0.000000,0.000000}%
\pgfsetstrokecolor{currentstroke}%
\pgfsetdash{}{0pt}%
\pgfpathmoveto{\pgfqpoint{1.837675in}{0.619202in}}%
\pgfpathlineto{\pgfqpoint{1.424602in}{0.289843in}}%
\pgfusepath{stroke}%
\end{pgfscope}%
\begin{pgfscope}%
\pgfsetrectcap%
\pgfsetroundjoin%
\pgfsetlinewidth{0.803000pt}%
\definecolor{currentstroke}{rgb}{0.000000,0.000000,0.000000}%
\pgfsetstrokecolor{currentstroke}%
\pgfsetdash{}{0pt}%
\pgfpathmoveto{\pgfqpoint{1.516903in}{0.372665in}}%
\pgfpathlineto{\pgfqpoint{1.542643in}{0.365473in}}%
\pgfusepath{stroke}%
\end{pgfscope}%
\begin{pgfscope}%
\definecolor{textcolor}{rgb}{0.000000,0.000000,0.000000}%
\pgfsetstrokecolor{textcolor}%
\pgfsetfillcolor{textcolor}%
\pgftext[x=1.665803in,y=0.235089in,,top]{\color{textcolor}{\rmfamily\fontsize{5.000000}{6.000000}\selectfont\catcode`\^=\active\def^{\ifmmode\sp\else\^{}\fi}\catcode`\%=\active\def%{\%}\ensuremath{-}0.5}}%
\end{pgfscope}%
\begin{pgfscope}%
\pgfsetrectcap%
\pgfsetroundjoin%
\pgfsetlinewidth{0.803000pt}%
\definecolor{currentstroke}{rgb}{0.000000,0.000000,0.000000}%
\pgfsetstrokecolor{currentstroke}%
\pgfsetdash{}{0pt}%
\pgfpathmoveto{\pgfqpoint{1.631821in}{0.464072in}}%
\pgfpathlineto{\pgfqpoint{1.657155in}{0.457221in}}%
\pgfusepath{stroke}%
\end{pgfscope}%
\begin{pgfscope}%
\definecolor{textcolor}{rgb}{0.000000,0.000000,0.000000}%
\pgfsetstrokecolor{textcolor}%
\pgfsetfillcolor{textcolor}%
\pgftext[x=1.777758in,y=0.330311in,,top]{\color{textcolor}{\rmfamily\fontsize{5.000000}{6.000000}\selectfont\catcode`\^=\active\def^{\ifmmode\sp\else\^{}\fi}\catcode`\%=\active\def%{\%}0.0}}%
\end{pgfscope}%
\begin{pgfscope}%
\pgfsetrectcap%
\pgfsetroundjoin%
\pgfsetlinewidth{0.803000pt}%
\definecolor{currentstroke}{rgb}{0.000000,0.000000,0.000000}%
\pgfsetstrokecolor{currentstroke}%
\pgfsetdash{}{0pt}%
\pgfpathmoveto{\pgfqpoint{1.741358in}{0.551200in}}%
\pgfpathlineto{\pgfqpoint{1.766295in}{0.544666in}}%
\pgfusepath{stroke}%
\end{pgfscope}%
\begin{pgfscope}%
\definecolor{textcolor}{rgb}{0.000000,0.000000,0.000000}%
\pgfsetstrokecolor{textcolor}%
\pgfsetfillcolor{textcolor}%
\pgftext[x=1.884443in,y=0.421050in,,top]{\color{textcolor}{\rmfamily\fontsize{5.000000}{6.000000}\selectfont\catcode`\^=\active\def^{\ifmmode\sp\else\^{}\fi}\catcode`\%=\active\def%{\%}0.5}}%
\end{pgfscope}%
\begin{pgfscope}%
\definecolor{textcolor}{rgb}{0.000000,0.000000,0.000000}%
\pgfsetstrokecolor{textcolor}%
\pgfsetfillcolor{textcolor}%
\pgftext[x=1.976512in,y=0.140262in,,]{\color{textcolor}{\rmfamily\fontsize{6.000000}{7.200000}\selectfont\catcode`\^=\active\def^{\ifmmode\sp\else\^{}\fi}\catcode`\%=\active\def%{\%}$b$}}%
\end{pgfscope}%
\begin{pgfscope}%
\pgfsetrectcap%
\pgfsetroundjoin%
\pgfsetlinewidth{0.803000pt}%
\definecolor{currentstroke}{rgb}{0.000000,0.000000,0.000000}%
\pgfsetstrokecolor{currentstroke}%
\pgfsetdash{}{0pt}%
\pgfpathmoveto{\pgfqpoint{1.837675in}{0.619202in}}%
\pgfpathlineto{\pgfqpoint{1.865965in}{1.326696in}}%
\pgfusepath{stroke}%
\end{pgfscope}%
\begin{pgfscope}%
\pgfsetrectcap%
\pgfsetroundjoin%
\pgfsetlinewidth{0.803000pt}%
\definecolor{currentstroke}{rgb}{0.000000,0.000000,0.000000}%
\pgfsetstrokecolor{currentstroke}%
\pgfsetdash{}{0pt}%
\pgfpathmoveto{\pgfqpoint{1.829482in}{0.621294in}}%
\pgfpathlineto{\pgfqpoint{1.854092in}{0.615010in}}%
\pgfusepath{stroke}%
\end{pgfscope}%
\begin{pgfscope}%
\definecolor{textcolor}{rgb}{0.000000,0.000000,0.000000}%
\pgfsetstrokecolor{textcolor}%
\pgfsetfillcolor{textcolor}%
\pgftext[x=2.022987in,y=0.624464in,,top]{\color{textcolor}{\rmfamily\fontsize{5.000000}{6.000000}\selectfont\catcode`\^=\active\def^{\ifmmode\sp\else\^{}\fi}\catcode`\%=\active\def%{\%}0.0}}%
\end{pgfscope}%
\begin{pgfscope}%
\pgfsetrectcap%
\pgfsetroundjoin%
\pgfsetlinewidth{0.803000pt}%
\definecolor{currentstroke}{rgb}{0.000000,0.000000,0.000000}%
\pgfsetstrokecolor{currentstroke}%
\pgfsetdash{}{0pt}%
\pgfpathmoveto{\pgfqpoint{1.836942in}{0.810078in}}%
\pgfpathlineto{\pgfqpoint{1.861829in}{0.804063in}}%
\pgfusepath{stroke}%
\end{pgfscope}%
\begin{pgfscope}%
\definecolor{textcolor}{rgb}{0.000000,0.000000,0.000000}%
\pgfsetstrokecolor{textcolor}%
\pgfsetfillcolor{textcolor}%
\pgftext[x=2.032509in,y=0.813112in,,top]{\color{textcolor}{\rmfamily\fontsize{5.000000}{6.000000}\selectfont\catcode`\^=\active\def^{\ifmmode\sp\else\^{}\fi}\catcode`\%=\active\def%{\%}0.5}}%
\end{pgfscope}%
\begin{pgfscope}%
\pgfsetrectcap%
\pgfsetroundjoin%
\pgfsetlinewidth{0.803000pt}%
\definecolor{currentstroke}{rgb}{0.000000,0.000000,0.000000}%
\pgfsetstrokecolor{currentstroke}%
\pgfsetdash{}{0pt}%
\pgfpathmoveto{\pgfqpoint{1.844563in}{1.002944in}}%
\pgfpathlineto{\pgfqpoint{1.869734in}{0.997212in}}%
\pgfusepath{stroke}%
\end{pgfscope}%
\begin{pgfscope}%
\definecolor{textcolor}{rgb}{0.000000,0.000000,0.000000}%
\pgfsetstrokecolor{textcolor}%
\pgfsetfillcolor{textcolor}%
\pgftext[x=2.042236in,y=1.005836in,,top]{\color{textcolor}{\rmfamily\fontsize{5.000000}{6.000000}\selectfont\catcode`\^=\active\def^{\ifmmode\sp\else\^{}\fi}\catcode`\%=\active\def%{\%}1.0}}%
\end{pgfscope}%
\begin{pgfscope}%
\pgfsetrectcap%
\pgfsetroundjoin%
\pgfsetlinewidth{0.803000pt}%
\definecolor{currentstroke}{rgb}{0.000000,0.000000,0.000000}%
\pgfsetstrokecolor{currentstroke}%
\pgfsetdash{}{0pt}%
\pgfpathmoveto{\pgfqpoint{1.852351in}{1.200029in}}%
\pgfpathlineto{\pgfqpoint{1.877812in}{1.194592in}}%
\pgfusepath{stroke}%
\end{pgfscope}%
\begin{pgfscope}%
\definecolor{textcolor}{rgb}{0.000000,0.000000,0.000000}%
\pgfsetstrokecolor{textcolor}%
\pgfsetfillcolor{textcolor}%
\pgftext[x=2.052177in,y=1.202771in,,top]{\color{textcolor}{\rmfamily\fontsize{5.000000}{6.000000}\selectfont\catcode`\^=\active\def^{\ifmmode\sp\else\^{}\fi}\catcode`\%=\active\def%{\%}1.5}}%
\end{pgfscope}%
\begin{pgfscope}%
\definecolor{textcolor}{rgb}{0.000000,0.000000,0.000000}%
\pgfsetstrokecolor{textcolor}%
\pgfsetfillcolor{textcolor}%
\pgftext[x=2.228830in,y=0.975668in,,]{\color{textcolor}{\rmfamily\fontsize{6.000000}{7.200000}\selectfont\catcode`\^=\active\def^{\ifmmode\sp\else\^{}\fi}\catcode`\%=\active\def%{\%}$a$}}%
\end{pgfscope}%
\begin{pgfscope}%
\pgfpathrectangle{\pgfqpoint{0.407592in}{0.222904in}}{\pgfqpoint{1.411163in}{1.411163in}}%
\pgfusepath{clip}%
\pgfsetrectcap%
\pgfsetroundjoin%
\pgfsetlinewidth{0.451687pt}%
\definecolor{currentstroke}{rgb}{0.560784,0.588235,0.611765}%
\pgfsetstrokecolor{currentstroke}%
\pgfsetstrokeopacity{0.450000}%
\pgfsetdash{}{0pt}%
\pgfpathmoveto{\pgfqpoint{1.257376in}{0.737514in}}%
\pgfpathlineto{\pgfqpoint{1.265273in}{0.743397in}}%
\pgfpathlineto{\pgfqpoint{1.272040in}{0.749493in}}%
\pgfpathlineto{\pgfqpoint{1.277643in}{0.755763in}}%
\pgfpathlineto{\pgfqpoint{1.282050in}{0.762165in}}%
\pgfpathlineto{\pgfqpoint{1.285242in}{0.768659in}}%
\pgfpathlineto{\pgfqpoint{1.286688in}{0.773018in}}%
\pgfpathlineto{\pgfqpoint{1.287585in}{0.777387in}}%
\pgfpathlineto{\pgfqpoint{1.287934in}{0.781754in}}%
\pgfpathlineto{\pgfqpoint{1.287735in}{0.786108in}}%
\pgfpathlineto{\pgfqpoint{1.286991in}{0.790435in}}%
\pgfpathlineto{\pgfqpoint{1.285706in}{0.794724in}}%
\pgfpathlineto{\pgfqpoint{1.283886in}{0.798964in}}%
\pgfpathlineto{\pgfqpoint{1.281538in}{0.803143in}}%
\pgfpathlineto{\pgfqpoint{1.277043in}{0.809273in}}%
\pgfpathlineto{\pgfqpoint{1.271412in}{0.815205in}}%
\pgfpathlineto{\pgfqpoint{1.264686in}{0.820905in}}%
\pgfpathlineto{\pgfqpoint{1.256908in}{0.826340in}}%
\pgfpathlineto{\pgfqpoint{1.248132in}{0.831477in}}%
\pgfpathlineto{\pgfqpoint{1.238414in}{0.836289in}}%
\pgfpathlineto{\pgfqpoint{1.227816in}{0.840747in}}%
\pgfpathlineto{\pgfqpoint{1.216402in}{0.844828in}}%
\pgfpathlineto{\pgfqpoint{1.204244in}{0.848508in}}%
\pgfpathlineto{\pgfqpoint{1.191415in}{0.851768in}}%
\pgfpathlineto{\pgfqpoint{1.177991in}{0.854589in}}%
\pgfpathlineto{\pgfqpoint{1.159305in}{0.857643in}}%
\pgfpathlineto{\pgfqpoint{1.139899in}{0.859862in}}%
\pgfpathlineto{\pgfqpoint{1.119974in}{0.861224in}}%
\pgfpathlineto{\pgfqpoint{1.099735in}{0.861716in}}%
\pgfpathlineto{\pgfqpoint{1.079390in}{0.861334in}}%
\pgfpathlineto{\pgfqpoint{1.059147in}{0.860082in}}%
\pgfpathlineto{\pgfqpoint{1.039215in}{0.857971in}}%
\pgfpathlineto{\pgfqpoint{1.019800in}{0.855021in}}%
\pgfpathlineto{\pgfqpoint{1.001103in}{0.851261in}}%
\pgfpathlineto{\pgfqpoint{0.983321in}{0.846727in}}%
\pgfpathlineto{\pgfqpoint{0.970698in}{0.842844in}}%
\pgfpathlineto{\pgfqpoint{0.958773in}{0.838572in}}%
\pgfpathlineto{\pgfqpoint{0.947617in}{0.833935in}}%
\pgfpathlineto{\pgfqpoint{0.937302in}{0.828957in}}%
\pgfpathlineto{\pgfqpoint{0.927891in}{0.823668in}}%
\pgfpathlineto{\pgfqpoint{0.919445in}{0.818098in}}%
\pgfpathlineto{\pgfqpoint{0.912020in}{0.812278in}}%
\pgfpathlineto{\pgfqpoint{0.905666in}{0.806242in}}%
\pgfpathlineto{\pgfqpoint{0.900427in}{0.800027in}}%
\pgfpathlineto{\pgfqpoint{0.896342in}{0.793668in}}%
\pgfpathlineto{\pgfqpoint{0.893441in}{0.787205in}}%
\pgfpathlineto{\pgfqpoint{0.892179in}{0.782857in}}%
\pgfpathlineto{\pgfqpoint{0.891460in}{0.778492in}}%
\pgfpathlineto{\pgfqpoint{0.891288in}{0.774122in}}%
\pgfpathlineto{\pgfqpoint{0.891666in}{0.769759in}}%
\pgfpathlineto{\pgfqpoint{0.892596in}{0.765416in}}%
\pgfpathlineto{\pgfqpoint{0.894076in}{0.761103in}}%
\pgfpathlineto{\pgfqpoint{0.896104in}{0.756834in}}%
\pgfpathlineto{\pgfqpoint{0.898678in}{0.752621in}}%
\pgfpathlineto{\pgfqpoint{0.903548in}{0.746433in}}%
\pgfpathlineto{\pgfqpoint{0.909607in}{0.740438in}}%
\pgfpathlineto{\pgfqpoint{0.916822in}{0.734676in}}%
\pgfpathlineto{\pgfqpoint{0.925152in}{0.729186in}}%
\pgfpathlineto{\pgfqpoint{0.934549in}{0.724004in}}%
\pgfpathlineto{\pgfqpoint{0.944956in}{0.719165in}}%
\pgfpathlineto{\pgfqpoint{0.956310in}{0.714704in}}%
\pgfpathlineto{\pgfqpoint{0.968538in}{0.710650in}}%
\pgfpathlineto{\pgfqpoint{0.981563in}{0.707032in}}%
\pgfpathlineto{\pgfqpoint{0.995303in}{0.703875in}}%
\pgfpathlineto{\pgfqpoint{1.009666in}{0.701201in}}%
\pgfpathlineto{\pgfqpoint{1.024560in}{0.699029in}}%
\pgfpathlineto{\pgfqpoint{1.039886in}{0.697375in}}%
\pgfpathlineto{\pgfqpoint{1.060817in}{0.695995in}}%
\pgfpathlineto{\pgfqpoint{1.082086in}{0.695572in}}%
\pgfpathlineto{\pgfqpoint{1.103443in}{0.696114in}}%
\pgfpathlineto{\pgfqpoint{1.124633in}{0.697611in}}%
\pgfpathlineto{\pgfqpoint{1.145406in}{0.700047in}}%
\pgfpathlineto{\pgfqpoint{1.160563in}{0.702471in}}%
\pgfpathlineto{\pgfqpoint{1.175248in}{0.705388in}}%
\pgfpathlineto{\pgfqpoint{1.189364in}{0.708778in}}%
\pgfpathlineto{\pgfqpoint{1.202820in}{0.712616in}}%
\pgfpathlineto{\pgfqpoint{1.215529in}{0.716877in}}%
\pgfpathlineto{\pgfqpoint{1.227410in}{0.721530in}}%
\pgfpathlineto{\pgfqpoint{1.238388in}{0.726544in}}%
\pgfpathlineto{\pgfqpoint{1.248397in}{0.731884in}}%
\pgfpathlineto{\pgfqpoint{1.257376in}{0.737514in}}%
\pgfpathlineto{\pgfqpoint{1.257376in}{0.737514in}}%
\pgfusepath{stroke}%
\end{pgfscope}%
\begin{pgfscope}%
\pgfpathrectangle{\pgfqpoint{0.407592in}{0.222904in}}{\pgfqpoint{1.411163in}{1.411163in}}%
\pgfusepath{clip}%
\pgfsetrectcap%
\pgfsetroundjoin%
\pgfsetlinewidth{0.451687pt}%
\definecolor{currentstroke}{rgb}{0.560784,0.588235,0.611765}%
\pgfsetstrokecolor{currentstroke}%
\pgfsetstrokeopacity{0.450000}%
\pgfsetdash{}{0pt}%
\pgfpathmoveto{\pgfqpoint{1.432199in}{0.870327in}}%
\pgfpathlineto{\pgfqpoint{1.442805in}{0.877983in}}%
\pgfpathlineto{\pgfqpoint{1.452363in}{0.885817in}}%
\pgfpathlineto{\pgfqpoint{1.460854in}{0.893806in}}%
\pgfpathlineto{\pgfqpoint{1.468260in}{0.901926in}}%
\pgfpathlineto{\pgfqpoint{1.474568in}{0.910151in}}%
\pgfpathlineto{\pgfqpoint{1.479767in}{0.918458in}}%
\pgfpathlineto{\pgfqpoint{1.483853in}{0.926822in}}%
\pgfpathlineto{\pgfqpoint{1.486821in}{0.935219in}}%
\pgfpathlineto{\pgfqpoint{1.488671in}{0.943625in}}%
\pgfpathlineto{\pgfqpoint{1.489408in}{0.952017in}}%
\pgfpathlineto{\pgfqpoint{1.489038in}{0.960371in}}%
\pgfpathlineto{\pgfqpoint{1.487569in}{0.968665in}}%
\pgfpathlineto{\pgfqpoint{1.485016in}{0.976877in}}%
\pgfpathlineto{\pgfqpoint{1.481393in}{0.984984in}}%
\pgfpathlineto{\pgfqpoint{1.476717in}{0.992967in}}%
\pgfpathlineto{\pgfqpoint{1.471010in}{1.000805in}}%
\pgfpathlineto{\pgfqpoint{1.464293in}{1.008477in}}%
\pgfpathlineto{\pgfqpoint{1.456592in}{1.015966in}}%
\pgfpathlineto{\pgfqpoint{1.447934in}{1.023252in}}%
\pgfpathlineto{\pgfqpoint{1.438347in}{1.030319in}}%
\pgfpathlineto{\pgfqpoint{1.427863in}{1.037150in}}%
\pgfpathlineto{\pgfqpoint{1.416514in}{1.043729in}}%
\pgfpathlineto{\pgfqpoint{1.404334in}{1.050040in}}%
\pgfpathlineto{\pgfqpoint{1.391359in}{1.056070in}}%
\pgfpathlineto{\pgfqpoint{1.377625in}{1.061805in}}%
\pgfpathlineto{\pgfqpoint{1.363173in}{1.067232in}}%
\pgfpathlineto{\pgfqpoint{1.348040in}{1.072339in}}%
\pgfpathlineto{\pgfqpoint{1.332267in}{1.077116in}}%
\pgfpathlineto{\pgfqpoint{1.307500in}{1.083640in}}%
\pgfpathlineto{\pgfqpoint{1.281531in}{1.089367in}}%
\pgfpathlineto{\pgfqpoint{1.254509in}{1.094270in}}%
\pgfpathlineto{\pgfqpoint{1.226584in}{1.098326in}}%
\pgfpathlineto{\pgfqpoint{1.197911in}{1.101517in}}%
\pgfpathlineto{\pgfqpoint{1.168645in}{1.103829in}}%
\pgfpathlineto{\pgfqpoint{1.138946in}{1.105250in}}%
\pgfpathlineto{\pgfqpoint{1.108971in}{1.105775in}}%
\pgfpathlineto{\pgfqpoint{1.078883in}{1.105402in}}%
\pgfpathlineto{\pgfqpoint{1.048840in}{1.104131in}}%
\pgfpathlineto{\pgfqpoint{1.019005in}{1.101969in}}%
\pgfpathlineto{\pgfqpoint{0.989538in}{1.098925in}}%
\pgfpathlineto{\pgfqpoint{0.960598in}{1.095014in}}%
\pgfpathlineto{\pgfqpoint{0.932343in}{1.090252in}}%
\pgfpathlineto{\pgfqpoint{0.904932in}{1.084663in}}%
\pgfpathlineto{\pgfqpoint{0.878516in}{1.078271in}}%
\pgfpathlineto{\pgfqpoint{0.853249in}{1.071108in}}%
\pgfpathlineto{\pgfqpoint{0.829279in}{1.063207in}}%
\pgfpathlineto{\pgfqpoint{0.814089in}{1.057548in}}%
\pgfpathlineto{\pgfqpoint{0.799582in}{1.051592in}}%
\pgfpathlineto{\pgfqpoint{0.785796in}{1.045351in}}%
\pgfpathlineto{\pgfqpoint{0.772771in}{1.038838in}}%
\pgfpathlineto{\pgfqpoint{0.760545in}{1.032070in}}%
\pgfpathlineto{\pgfqpoint{0.749155in}{1.025061in}}%
\pgfpathlineto{\pgfqpoint{0.738635in}{1.017828in}}%
\pgfpathlineto{\pgfqpoint{0.729018in}{1.010389in}}%
\pgfpathlineto{\pgfqpoint{0.720338in}{1.002761in}}%
\pgfpathlineto{\pgfqpoint{0.712622in}{0.994964in}}%
\pgfpathlineto{\pgfqpoint{0.705899in}{0.987016in}}%
\pgfpathlineto{\pgfqpoint{0.700194in}{0.978938in}}%
\pgfpathlineto{\pgfqpoint{0.695530in}{0.970750in}}%
\pgfpathlineto{\pgfqpoint{0.691928in}{0.962475in}}%
\pgfpathlineto{\pgfqpoint{0.689406in}{0.954134in}}%
\pgfpathlineto{\pgfqpoint{0.687979in}{0.945750in}}%
\pgfpathlineto{\pgfqpoint{0.687661in}{0.937345in}}%
\pgfpathlineto{\pgfqpoint{0.688459in}{0.928943in}}%
\pgfpathlineto{\pgfqpoint{0.690382in}{0.920569in}}%
\pgfpathlineto{\pgfqpoint{0.693431in}{0.912245in}}%
\pgfpathlineto{\pgfqpoint{0.697608in}{0.903997in}}%
\pgfpathlineto{\pgfqpoint{0.702907in}{0.895848in}}%
\pgfpathlineto{\pgfqpoint{0.709324in}{0.887824in}}%
\pgfpathlineto{\pgfqpoint{0.716845in}{0.879948in}}%
\pgfpathlineto{\pgfqpoint{0.725458in}{0.872245in}}%
\pgfpathlineto{\pgfqpoint{0.735145in}{0.864739in}}%
\pgfpathlineto{\pgfqpoint{0.745883in}{0.857455in}}%
\pgfpathlineto{\pgfqpoint{0.757648in}{0.850414in}}%
\pgfpathlineto{\pgfqpoint{0.770409in}{0.843641in}}%
\pgfpathlineto{\pgfqpoint{0.784135in}{0.837157in}}%
\pgfpathlineto{\pgfqpoint{0.798789in}{0.830985in}}%
\pgfpathlineto{\pgfqpoint{0.814331in}{0.825144in}}%
\pgfpathlineto{\pgfqpoint{0.830717in}{0.819655in}}%
\pgfpathlineto{\pgfqpoint{0.847901in}{0.814536in}}%
\pgfpathlineto{\pgfqpoint{0.865832in}{0.809805in}}%
\pgfpathlineto{\pgfqpoint{0.884458in}{0.805478in}}%
\pgfpathlineto{\pgfqpoint{0.903723in}{0.801572in}}%
\pgfpathlineto{\pgfqpoint{0.923567in}{0.798099in}}%
\pgfpathlineto{\pgfqpoint{0.943931in}{0.795071in}}%
\pgfpathlineto{\pgfqpoint{0.964751in}{0.792500in}}%
\pgfpathlineto{\pgfqpoint{0.985961in}{0.790395in}}%
\pgfpathlineto{\pgfqpoint{1.007495in}{0.788764in}}%
\pgfpathlineto{\pgfqpoint{1.040254in}{0.787216in}}%
\pgfpathlineto{\pgfqpoint{1.073354in}{0.786759in}}%
\pgfpathlineto{\pgfqpoint{1.106558in}{0.787396in}}%
\pgfpathlineto{\pgfqpoint{1.139626in}{0.789122in}}%
\pgfpathlineto{\pgfqpoint{1.172322in}{0.791924in}}%
\pgfpathlineto{\pgfqpoint{1.204413in}{0.795777in}}%
\pgfpathlineto{\pgfqpoint{1.225359in}{0.798916in}}%
\pgfpathlineto{\pgfqpoint{1.245870in}{0.802497in}}%
\pgfpathlineto{\pgfqpoint{1.265883in}{0.806509in}}%
\pgfpathlineto{\pgfqpoint{1.285337in}{0.810937in}}%
\pgfpathlineto{\pgfqpoint{1.304173in}{0.815765in}}%
\pgfpathlineto{\pgfqpoint{1.322334in}{0.820977in}}%
\pgfpathlineto{\pgfqpoint{1.339768in}{0.826555in}}%
\pgfpathlineto{\pgfqpoint{1.356425in}{0.832480in}}%
\pgfpathlineto{\pgfqpoint{1.372257in}{0.838731in}}%
\pgfpathlineto{\pgfqpoint{1.387220in}{0.845289in}}%
\pgfpathlineto{\pgfqpoint{1.401274in}{0.852130in}}%
\pgfpathlineto{\pgfqpoint{1.414383in}{0.859233in}}%
\pgfpathlineto{\pgfqpoint{1.426512in}{0.866574in}}%
\pgfpathlineto{\pgfqpoint{1.432199in}{0.870327in}}%
\pgfpathlineto{\pgfqpoint{1.432199in}{0.870327in}}%
\pgfusepath{stroke}%
\end{pgfscope}%
\begin{pgfscope}%
\pgfpathrectangle{\pgfqpoint{0.407592in}{0.222904in}}{\pgfqpoint{1.411163in}{1.411163in}}%
\pgfusepath{clip}%
\pgfsetrectcap%
\pgfsetroundjoin%
\pgfsetlinewidth{0.451687pt}%
\definecolor{currentstroke}{rgb}{0.560784,0.588235,0.611765}%
\pgfsetstrokecolor{currentstroke}%
\pgfsetstrokeopacity{0.450000}%
\pgfsetdash{}{0pt}%
\pgfpathmoveto{\pgfqpoint{1.615239in}{1.009383in}}%
\pgfpathlineto{\pgfqpoint{1.630853in}{1.020628in}}%
\pgfpathlineto{\pgfqpoint{1.644808in}{1.032109in}}%
\pgfpathlineto{\pgfqpoint{1.657082in}{1.043786in}}%
\pgfpathlineto{\pgfqpoint{1.667655in}{1.055623in}}%
\pgfpathlineto{\pgfqpoint{1.676515in}{1.067583in}}%
\pgfpathlineto{\pgfqpoint{1.680301in}{1.073598in}}%
\pgfpathlineto{\pgfqpoint{1.683657in}{1.079630in}}%
\pgfpathlineto{\pgfqpoint{1.686582in}{1.085674in}}%
\pgfpathlineto{\pgfqpoint{1.689078in}{1.091727in}}%
\pgfpathlineto{\pgfqpoint{1.691145in}{1.097783in}}%
\pgfpathlineto{\pgfqpoint{1.692784in}{1.103839in}}%
\pgfpathlineto{\pgfqpoint{1.693996in}{1.109890in}}%
\pgfpathlineto{\pgfqpoint{1.694783in}{1.115932in}}%
\pgfpathlineto{\pgfqpoint{1.695148in}{1.121961in}}%
\pgfpathlineto{\pgfqpoint{1.695091in}{1.127972in}}%
\pgfpathlineto{\pgfqpoint{1.694616in}{1.133962in}}%
\pgfpathlineto{\pgfqpoint{1.693726in}{1.139926in}}%
\pgfpathlineto{\pgfqpoint{1.692424in}{1.145861in}}%
\pgfpathlineto{\pgfqpoint{1.690712in}{1.151763in}}%
\pgfpathlineto{\pgfqpoint{1.688595in}{1.157627in}}%
\pgfpathlineto{\pgfqpoint{1.686076in}{1.163451in}}%
\pgfpathlineto{\pgfqpoint{1.683159in}{1.169230in}}%
\pgfpathlineto{\pgfqpoint{1.676148in}{1.180641in}}%
\pgfpathlineto{\pgfqpoint{1.667600in}{1.191832in}}%
\pgfpathlineto{\pgfqpoint{1.657553in}{1.202776in}}%
\pgfpathlineto{\pgfqpoint{1.646050in}{1.213448in}}%
\pgfpathlineto{\pgfqpoint{1.633135in}{1.223823in}}%
\pgfpathlineto{\pgfqpoint{1.618855in}{1.233880in}}%
\pgfpathlineto{\pgfqpoint{1.603260in}{1.243595in}}%
\pgfpathlineto{\pgfqpoint{1.586402in}{1.252949in}}%
\pgfpathlineto{\pgfqpoint{1.568332in}{1.261922in}}%
\pgfpathlineto{\pgfqpoint{1.549106in}{1.270496in}}%
\pgfpathlineto{\pgfqpoint{1.528781in}{1.278654in}}%
\pgfpathlineto{\pgfqpoint{1.507412in}{1.286380in}}%
\pgfpathlineto{\pgfqpoint{1.485058in}{1.293659in}}%
\pgfpathlineto{\pgfqpoint{1.461779in}{1.300479in}}%
\pgfpathlineto{\pgfqpoint{1.437636in}{1.306826in}}%
\pgfpathlineto{\pgfqpoint{1.412688in}{1.312689in}}%
\pgfpathlineto{\pgfqpoint{1.386997in}{1.318057in}}%
\pgfpathlineto{\pgfqpoint{1.360627in}{1.322922in}}%
\pgfpathlineto{\pgfqpoint{1.333638in}{1.327275in}}%
\pgfpathlineto{\pgfqpoint{1.306095in}{1.331109in}}%
\pgfpathlineto{\pgfqpoint{1.278060in}{1.334416in}}%
\pgfpathlineto{\pgfqpoint{1.249598in}{1.337193in}}%
\pgfpathlineto{\pgfqpoint{1.220773in}{1.339433in}}%
\pgfpathlineto{\pgfqpoint{1.191648in}{1.341134in}}%
\pgfpathlineto{\pgfqpoint{1.162288in}{1.342292in}}%
\pgfpathlineto{\pgfqpoint{1.132758in}{1.342906in}}%
\pgfpathlineto{\pgfqpoint{1.103122in}{1.342975in}}%
\pgfpathlineto{\pgfqpoint{1.073446in}{1.342499in}}%
\pgfpathlineto{\pgfqpoint{1.043795in}{1.341478in}}%
\pgfpathlineto{\pgfqpoint{1.014232in}{1.339914in}}%
\pgfpathlineto{\pgfqpoint{0.984825in}{1.337810in}}%
\pgfpathlineto{\pgfqpoint{0.955636in}{1.335169in}}%
\pgfpathlineto{\pgfqpoint{0.926733in}{1.331995in}}%
\pgfpathlineto{\pgfqpoint{0.898179in}{1.328294in}}%
\pgfpathlineto{\pgfqpoint{0.870040in}{1.324072in}}%
\pgfpathlineto{\pgfqpoint{0.842381in}{1.319336in}}%
\pgfpathlineto{\pgfqpoint{0.815266in}{1.314094in}}%
\pgfpathlineto{\pgfqpoint{0.788760in}{1.308355in}}%
\pgfpathlineto{\pgfqpoint{0.762926in}{1.302129in}}%
\pgfpathlineto{\pgfqpoint{0.737829in}{1.295428in}}%
\pgfpathlineto{\pgfqpoint{0.713532in}{1.288264in}}%
\pgfpathlineto{\pgfqpoint{0.690098in}{1.280649in}}%
\pgfpathlineto{\pgfqpoint{0.667589in}{1.272599in}}%
\pgfpathlineto{\pgfqpoint{0.646065in}{1.264129in}}%
\pgfpathlineto{\pgfqpoint{0.625588in}{1.255255in}}%
\pgfpathlineto{\pgfqpoint{0.606216in}{1.245996in}}%
\pgfpathlineto{\pgfqpoint{0.588007in}{1.236370in}}%
\pgfpathlineto{\pgfqpoint{0.571019in}{1.226398in}}%
\pgfpathlineto{\pgfqpoint{0.555305in}{1.216101in}}%
\pgfpathlineto{\pgfqpoint{0.540920in}{1.205501in}}%
\pgfpathlineto{\pgfqpoint{0.527913in}{1.194624in}}%
\pgfpathlineto{\pgfqpoint{0.516335in}{1.183493in}}%
\pgfpathlineto{\pgfqpoint{0.506231in}{1.172135in}}%
\pgfpathlineto{\pgfqpoint{0.497646in}{1.160578in}}%
\pgfpathlineto{\pgfqpoint{0.490618in}{1.148850in}}%
\pgfpathlineto{\pgfqpoint{0.487701in}{1.142932in}}%
\pgfpathlineto{\pgfqpoint{0.485187in}{1.136982in}}%
\pgfpathlineto{\pgfqpoint{0.483080in}{1.131005in}}%
\pgfpathlineto{\pgfqpoint{0.481385in}{1.125004in}}%
\pgfpathlineto{\pgfqpoint{0.480104in}{1.118983in}}%
\pgfpathlineto{\pgfqpoint{0.479242in}{1.112948in}}%
\pgfpathlineto{\pgfqpoint{0.478801in}{1.106901in}}%
\pgfpathlineto{\pgfqpoint{0.478784in}{1.100847in}}%
\pgfpathlineto{\pgfqpoint{0.479194in}{1.094790in}}%
\pgfpathlineto{\pgfqpoint{0.480033in}{1.088735in}}%
\pgfpathlineto{\pgfqpoint{0.481302in}{1.082686in}}%
\pgfpathlineto{\pgfqpoint{0.483004in}{1.076647in}}%
\pgfpathlineto{\pgfqpoint{0.485140in}{1.070623in}}%
\pgfpathlineto{\pgfqpoint{0.487710in}{1.064619in}}%
\pgfpathlineto{\pgfqpoint{0.490715in}{1.058638in}}%
\pgfpathlineto{\pgfqpoint{0.494156in}{1.052685in}}%
\pgfpathlineto{\pgfqpoint{0.498033in}{1.046766in}}%
\pgfpathlineto{\pgfqpoint{0.502344in}{1.040884in}}%
\pgfpathlineto{\pgfqpoint{0.512269in}{1.029252in}}%
\pgfpathlineto{\pgfqpoint{0.523920in}{1.017826in}}%
\pgfpathlineto{\pgfqpoint{0.537279in}{1.006644in}}%
\pgfpathlineto{\pgfqpoint{0.552324in}{0.995743in}}%
\pgfpathlineto{\pgfqpoint{0.569025in}{0.985160in}}%
\pgfpathlineto{\pgfqpoint{0.587346in}{0.974931in}}%
\pgfpathlineto{\pgfqpoint{0.607243in}{0.965093in}}%
\pgfpathlineto{\pgfqpoint{0.628667in}{0.955681in}}%
\pgfpathlineto{\pgfqpoint{0.651562in}{0.946729in}}%
\pgfpathlineto{\pgfqpoint{0.675867in}{0.938270in}}%
\pgfpathlineto{\pgfqpoint{0.701511in}{0.930336in}}%
\pgfpathlineto{\pgfqpoint{0.728419in}{0.922957in}}%
\pgfpathlineto{\pgfqpoint{0.756512in}{0.916161in}}%
\pgfpathlineto{\pgfqpoint{0.785702in}{0.909975in}}%
\pgfpathlineto{\pgfqpoint{0.815897in}{0.904423in}}%
\pgfpathlineto{\pgfqpoint{0.847001in}{0.899526in}}%
\pgfpathlineto{\pgfqpoint{0.878911in}{0.895306in}}%
\pgfpathlineto{\pgfqpoint{0.911523in}{0.891777in}}%
\pgfpathlineto{\pgfqpoint{0.944727in}{0.888956in}}%
\pgfpathlineto{\pgfqpoint{0.978411in}{0.886852in}}%
\pgfpathlineto{\pgfqpoint{1.012462in}{0.885476in}}%
\pgfpathlineto{\pgfqpoint{1.046761in}{0.884831in}}%
\pgfpathlineto{\pgfqpoint{1.081192in}{0.884922in}}%
\pgfpathlineto{\pgfqpoint{1.115637in}{0.885748in}}%
\pgfpathlineto{\pgfqpoint{1.149978in}{0.887305in}}%
\pgfpathlineto{\pgfqpoint{1.184097in}{0.889586in}}%
\pgfpathlineto{\pgfqpoint{1.217878in}{0.892584in}}%
\pgfpathlineto{\pgfqpoint{1.251206in}{0.896285in}}%
\pgfpathlineto{\pgfqpoint{1.283971in}{0.900674in}}%
\pgfpathlineto{\pgfqpoint{1.316063in}{0.905734in}}%
\pgfpathlineto{\pgfqpoint{1.347377in}{0.911445in}}%
\pgfpathlineto{\pgfqpoint{1.377811in}{0.917785in}}%
\pgfpathlineto{\pgfqpoint{1.407270in}{0.924727in}}%
\pgfpathlineto{\pgfqpoint{1.435662in}{0.932246in}}%
\pgfpathlineto{\pgfqpoint{1.462900in}{0.940313in}}%
\pgfpathlineto{\pgfqpoint{1.488904in}{0.948897in}}%
\pgfpathlineto{\pgfqpoint{1.513598in}{0.957965in}}%
\pgfpathlineto{\pgfqpoint{1.536914in}{0.967486in}}%
\pgfpathlineto{\pgfqpoint{1.558789in}{0.977424in}}%
\pgfpathlineto{\pgfqpoint{1.579167in}{0.987743in}}%
\pgfpathlineto{\pgfqpoint{1.597998in}{0.998409in}}%
\pgfpathlineto{\pgfqpoint{1.615239in}{1.009383in}}%
\pgfpathlineto{\pgfqpoint{1.615239in}{1.009383in}}%
\pgfusepath{stroke}%
\end{pgfscope}%
\begin{pgfscope}%
\pgfpathrectangle{\pgfqpoint{0.407592in}{0.222904in}}{\pgfqpoint{1.411163in}{1.411163in}}%
\pgfusepath{clip}%
\pgfsetrectcap%
\pgfsetroundjoin%
\pgfsetlinewidth{0.451687pt}%
\definecolor{currentstroke}{rgb}{0.560784,0.588235,0.611765}%
\pgfsetstrokecolor{currentstroke}%
\pgfsetstrokeopacity{0.450000}%
\pgfsetdash{}{0pt}%
\pgfpathmoveto{\pgfqpoint{1.763662in}{1.122140in}}%
\pgfpathlineto{\pgfqpoint{1.782996in}{1.136039in}}%
\pgfpathlineto{\pgfqpoint{1.791849in}{1.143089in}}%
\pgfpathlineto{\pgfqpoint{1.800155in}{1.150199in}}%
\pgfpathlineto{\pgfqpoint{1.807911in}{1.157362in}}%
\pgfpathlineto{\pgfqpoint{1.815114in}{1.164572in}}%
\pgfpathlineto{\pgfqpoint{1.822088in}{1.172211in}}%
\pgfpathmoveto{\pgfqpoint{1.822088in}{1.340804in}}%
\pgfpathlineto{\pgfqpoint{1.818991in}{1.344252in}}%
\pgfpathlineto{\pgfqpoint{1.812535in}{1.350843in}}%
\pgfpathlineto{\pgfqpoint{1.805624in}{1.357350in}}%
\pgfpathlineto{\pgfqpoint{1.790464in}{1.370098in}}%
\pgfpathlineto{\pgfqpoint{1.773573in}{1.382469in}}%
\pgfpathlineto{\pgfqpoint{1.755016in}{1.394439in}}%
\pgfpathlineto{\pgfqpoint{1.734858in}{1.405982in}}%
\pgfpathlineto{\pgfqpoint{1.713167in}{1.417078in}}%
\pgfpathlineto{\pgfqpoint{1.690012in}{1.427704in}}%
\pgfpathlineto{\pgfqpoint{1.665464in}{1.437843in}}%
\pgfpathlineto{\pgfqpoint{1.639594in}{1.447475in}}%
\pgfpathlineto{\pgfqpoint{1.612475in}{1.456585in}}%
\pgfpathlineto{\pgfqpoint{1.584179in}{1.465157in}}%
\pgfpathlineto{\pgfqpoint{1.554782in}{1.473178in}}%
\pgfpathlineto{\pgfqpoint{1.524355in}{1.480635in}}%
\pgfpathlineto{\pgfqpoint{1.492975in}{1.487515in}}%
\pgfpathlineto{\pgfqpoint{1.460716in}{1.493810in}}%
\pgfpathlineto{\pgfqpoint{1.427652in}{1.499508in}}%
\pgfpathlineto{\pgfqpoint{1.393859in}{1.504603in}}%
\pgfpathlineto{\pgfqpoint{1.359412in}{1.509087in}}%
\pgfpathlineto{\pgfqpoint{1.324386in}{1.512953in}}%
\pgfpathlineto{\pgfqpoint{1.288857in}{1.516197in}}%
\pgfpathlineto{\pgfqpoint{1.252899in}{1.518813in}}%
\pgfpathlineto{\pgfqpoint{1.216588in}{1.520799in}}%
\pgfpathlineto{\pgfqpoint{1.179999in}{1.522151in}}%
\pgfpathlineto{\pgfqpoint{1.143208in}{1.522867in}}%
\pgfpathlineto{\pgfqpoint{1.106290in}{1.522948in}}%
\pgfpathlineto{\pgfqpoint{1.069322in}{1.522392in}}%
\pgfpathlineto{\pgfqpoint{1.032377in}{1.521200in}}%
\pgfpathlineto{\pgfqpoint{0.995534in}{1.519375in}}%
\pgfpathlineto{\pgfqpoint{0.958866in}{1.516918in}}%
\pgfpathlineto{\pgfqpoint{0.922452in}{1.513833in}}%
\pgfpathlineto{\pgfqpoint{0.886366in}{1.510123in}}%
\pgfpathlineto{\pgfqpoint{0.850686in}{1.505795in}}%
\pgfpathlineto{\pgfqpoint{0.815487in}{1.500854in}}%
\pgfpathlineto{\pgfqpoint{0.780849in}{1.495307in}}%
\pgfpathlineto{\pgfqpoint{0.746846in}{1.489163in}}%
\pgfpathlineto{\pgfqpoint{0.713557in}{1.482430in}}%
\pgfpathlineto{\pgfqpoint{0.681058in}{1.475118in}}%
\pgfpathlineto{\pgfqpoint{0.649427in}{1.467239in}}%
\pgfpathlineto{\pgfqpoint{0.618742in}{1.458804in}}%
\pgfpathlineto{\pgfqpoint{0.589080in}{1.449829in}}%
\pgfpathlineto{\pgfqpoint{0.560517in}{1.440327in}}%
\pgfpathlineto{\pgfqpoint{0.533130in}{1.430315in}}%
\pgfpathlineto{\pgfqpoint{0.506995in}{1.419810in}}%
\pgfpathlineto{\pgfqpoint{0.482188in}{1.408831in}}%
\pgfpathlineto{\pgfqpoint{0.458784in}{1.397399in}}%
\pgfpathlineto{\pgfqpoint{0.436856in}{1.385535in}}%
\pgfpathlineto{\pgfqpoint{0.416476in}{1.373263in}}%
\pgfpathlineto{\pgfqpoint{0.404258in}{1.365154in}}%
\pgfpathmoveto{\pgfqpoint{0.404258in}{1.102084in}}%
\pgfpathlineto{\pgfqpoint{0.409555in}{1.098627in}}%
\pgfpathlineto{\pgfqpoint{0.431448in}{1.085692in}}%
\pgfpathlineto{\pgfqpoint{0.455462in}{1.073197in}}%
\pgfpathlineto{\pgfqpoint{0.481542in}{1.061190in}}%
\pgfpathlineto{\pgfqpoint{0.509624in}{1.049715in}}%
\pgfpathlineto{\pgfqpoint{0.539635in}{1.038819in}}%
\pgfpathlineto{\pgfqpoint{0.571492in}{1.028544in}}%
\pgfpathlineto{\pgfqpoint{0.605103in}{1.018932in}}%
\pgfpathlineto{\pgfqpoint{0.640368in}{1.010023in}}%
\pgfpathlineto{\pgfqpoint{0.677177in}{1.001854in}}%
\pgfpathlineto{\pgfqpoint{0.715411in}{0.994459in}}%
\pgfpathlineto{\pgfqpoint{0.754945in}{0.987870in}}%
\pgfpathlineto{\pgfqpoint{0.795645in}{0.982116in}}%
\pgfpathlineto{\pgfqpoint{0.837373in}{0.977221in}}%
\pgfpathlineto{\pgfqpoint{0.879982in}{0.973208in}}%
\pgfpathlineto{\pgfqpoint{0.923322in}{0.970093in}}%
\pgfpathlineto{\pgfqpoint{0.967239in}{0.967892in}}%
\pgfpathlineto{\pgfqpoint{1.011573in}{0.966613in}}%
\pgfpathlineto{\pgfqpoint{1.056165in}{0.966263in}}%
\pgfpathlineto{\pgfqpoint{1.100854in}{0.966843in}}%
\pgfpathlineto{\pgfqpoint{1.145476in}{0.968350in}}%
\pgfpathlineto{\pgfqpoint{1.189870in}{0.970778in}}%
\pgfpathlineto{\pgfqpoint{1.233876in}{0.974116in}}%
\pgfpathlineto{\pgfqpoint{1.277338in}{0.978349in}}%
\pgfpathlineto{\pgfqpoint{1.320099in}{0.983459in}}%
\pgfpathlineto{\pgfqpoint{1.362012in}{0.989422in}}%
\pgfpathlineto{\pgfqpoint{1.402931in}{0.996213in}}%
\pgfpathlineto{\pgfqpoint{1.442717in}{1.003802in}}%
\pgfpathlineto{\pgfqpoint{1.481240in}{1.012157in}}%
\pgfpathlineto{\pgfqpoint{1.518373in}{1.021244in}}%
\pgfpathlineto{\pgfqpoint{1.553999in}{1.031023in}}%
\pgfpathlineto{\pgfqpoint{1.588011in}{1.041455in}}%
\pgfpathlineto{\pgfqpoint{1.620307in}{1.052498in}}%
\pgfpathlineto{\pgfqpoint{1.650797in}{1.064109in}}%
\pgfpathlineto{\pgfqpoint{1.679399in}{1.076241in}}%
\pgfpathlineto{\pgfqpoint{1.706039in}{1.088849in}}%
\pgfpathlineto{\pgfqpoint{1.730654in}{1.101886in}}%
\pgfpathlineto{\pgfqpoint{1.753189in}{1.115303in}}%
\pgfpathlineto{\pgfqpoint{1.763662in}{1.122140in}}%
\pgfpathlineto{\pgfqpoint{1.763662in}{1.122140in}}%
\pgfusepath{stroke}%
\end{pgfscope}%
\begin{pgfscope}%
\pgfpathrectangle{\pgfqpoint{0.407592in}{0.222904in}}{\pgfqpoint{1.411163in}{1.411163in}}%
\pgfusepath{clip}%
\pgfsetrectcap%
\pgfsetroundjoin%
\pgfsetlinewidth{0.401500pt}%
\definecolor{currentstroke}{rgb}{0.560784,0.588235,0.611765}%
\pgfsetstrokecolor{currentstroke}%
\pgfsetstrokeopacity{0.350000}%
\pgfsetdash{}{0pt}%
\pgfpathmoveto{\pgfqpoint{1.090230in}{0.610533in}}%
\pgfpathlineto{\pgfqpoint{1.098147in}{0.616548in}}%
\pgfpathlineto{\pgfqpoint{1.106082in}{0.622576in}}%
\pgfpathlineto{\pgfqpoint{1.114034in}{0.628617in}}%
\pgfpathlineto{\pgfqpoint{1.122002in}{0.634671in}}%
\pgfpathlineto{\pgfqpoint{1.129988in}{0.640737in}}%
\pgfpathlineto{\pgfqpoint{1.137991in}{0.646817in}}%
\pgfpathlineto{\pgfqpoint{1.146011in}{0.652910in}}%
\pgfpathlineto{\pgfqpoint{1.154048in}{0.659015in}}%
\pgfpathlineto{\pgfqpoint{1.162102in}{0.665134in}}%
\pgfpathlineto{\pgfqpoint{1.170173in}{0.671266in}}%
\pgfpathlineto{\pgfqpoint{1.178262in}{0.677411in}}%
\pgfpathlineto{\pgfqpoint{1.186368in}{0.683570in}}%
\pgfpathlineto{\pgfqpoint{1.194492in}{0.689741in}}%
\pgfpathlineto{\pgfqpoint{1.202633in}{0.695926in}}%
\pgfpathlineto{\pgfqpoint{1.210792in}{0.702124in}}%
\pgfpathlineto{\pgfqpoint{1.218969in}{0.708336in}}%
\pgfpathlineto{\pgfqpoint{1.227163in}{0.714561in}}%
\pgfpathlineto{\pgfqpoint{1.235375in}{0.720800in}}%
\pgfpathlineto{\pgfqpoint{1.243605in}{0.727052in}}%
\pgfpathlineto{\pgfqpoint{1.251852in}{0.733318in}}%
\pgfpathlineto{\pgfqpoint{1.260118in}{0.739597in}}%
\pgfpathlineto{\pgfqpoint{1.268401in}{0.745890in}}%
\pgfpathlineto{\pgfqpoint{1.276703in}{0.752196in}}%
\pgfpathlineto{\pgfqpoint{1.285022in}{0.758517in}}%
\pgfpathlineto{\pgfqpoint{1.293360in}{0.764851in}}%
\pgfpathlineto{\pgfqpoint{1.301716in}{0.771199in}}%
\pgfpathlineto{\pgfqpoint{1.310090in}{0.777561in}}%
\pgfpathlineto{\pgfqpoint{1.318483in}{0.783937in}}%
\pgfpathlineto{\pgfqpoint{1.326894in}{0.790327in}}%
\pgfpathlineto{\pgfqpoint{1.335323in}{0.796731in}}%
\pgfpathlineto{\pgfqpoint{1.343771in}{0.803148in}}%
\pgfpathlineto{\pgfqpoint{1.352238in}{0.809580in}}%
\pgfpathlineto{\pgfqpoint{1.360723in}{0.816027in}}%
\pgfpathlineto{\pgfqpoint{1.369227in}{0.822487in}}%
\pgfpathlineto{\pgfqpoint{1.377749in}{0.828962in}}%
\pgfpathlineto{\pgfqpoint{1.386291in}{0.835451in}}%
\pgfpathlineto{\pgfqpoint{1.394851in}{0.841954in}}%
\pgfpathlineto{\pgfqpoint{1.403430in}{0.848472in}}%
\pgfpathlineto{\pgfqpoint{1.412029in}{0.855004in}}%
\pgfpathlineto{\pgfqpoint{1.420646in}{0.861550in}}%
\pgfpathlineto{\pgfqpoint{1.429283in}{0.868112in}}%
\pgfpathlineto{\pgfqpoint{1.437938in}{0.874687in}}%
\pgfpathlineto{\pgfqpoint{1.446613in}{0.881278in}}%
\pgfpathlineto{\pgfqpoint{1.455308in}{0.887883in}}%
\pgfpathlineto{\pgfqpoint{1.464021in}{0.894503in}}%
\pgfpathlineto{\pgfqpoint{1.472755in}{0.901137in}}%
\pgfpathlineto{\pgfqpoint{1.481507in}{0.907787in}}%
\pgfpathlineto{\pgfqpoint{1.490280in}{0.914451in}}%
\pgfpathlineto{\pgfqpoint{1.499072in}{0.921130in}}%
\pgfpathlineto{\pgfqpoint{1.507883in}{0.927824in}}%
\pgfpathlineto{\pgfqpoint{1.516715in}{0.934534in}}%
\pgfpathlineto{\pgfqpoint{1.525566in}{0.941258in}}%
\pgfpathlineto{\pgfqpoint{1.534437in}{0.947998in}}%
\pgfpathlineto{\pgfqpoint{1.543329in}{0.954752in}}%
\pgfpathlineto{\pgfqpoint{1.552240in}{0.961522in}}%
\pgfpathlineto{\pgfqpoint{1.561171in}{0.968307in}}%
\pgfpathlineto{\pgfqpoint{1.570123in}{0.975108in}}%
\pgfpathlineto{\pgfqpoint{1.579095in}{0.981924in}}%
\pgfpathlineto{\pgfqpoint{1.588087in}{0.988755in}}%
\pgfpathlineto{\pgfqpoint{1.597100in}{0.995602in}}%
\pgfpathlineto{\pgfqpoint{1.606133in}{1.002465in}}%
\pgfpathlineto{\pgfqpoint{1.615186in}{1.009343in}}%
\pgfpathlineto{\pgfqpoint{1.624261in}{1.016236in}}%
\pgfpathlineto{\pgfqpoint{1.633356in}{1.023146in}}%
\pgfpathlineto{\pgfqpoint{1.642471in}{1.030071in}}%
\pgfpathlineto{\pgfqpoint{1.651608in}{1.037012in}}%
\pgfpathlineto{\pgfqpoint{1.660765in}{1.043969in}}%
\pgfpathlineto{\pgfqpoint{1.669944in}{1.050942in}}%
\pgfpathlineto{\pgfqpoint{1.679143in}{1.057931in}}%
\pgfpathlineto{\pgfqpoint{1.688364in}{1.064935in}}%
\pgfpathlineto{\pgfqpoint{1.697605in}{1.071956in}}%
\pgfpathlineto{\pgfqpoint{1.706868in}{1.078993in}}%
\pgfpathlineto{\pgfqpoint{1.716153in}{1.086047in}}%
\pgfpathlineto{\pgfqpoint{1.725458in}{1.093116in}}%
\pgfpathlineto{\pgfqpoint{1.734786in}{1.100202in}}%
\pgfpathlineto{\pgfqpoint{1.744134in}{1.107304in}}%
\pgfpathlineto{\pgfqpoint{1.753505in}{1.114423in}}%
\pgfpathlineto{\pgfqpoint{1.762897in}{1.121558in}}%
\pgfpathlineto{\pgfqpoint{1.772311in}{1.128710in}}%
\pgfusepath{stroke}%
\end{pgfscope}%
\begin{pgfscope}%
\pgfpathrectangle{\pgfqpoint{0.407592in}{0.222904in}}{\pgfqpoint{1.411163in}{1.411163in}}%
\pgfusepath{clip}%
\pgfsetrectcap%
\pgfsetroundjoin%
\pgfsetlinewidth{0.401500pt}%
\definecolor{currentstroke}{rgb}{0.560784,0.588235,0.611765}%
\pgfsetstrokecolor{currentstroke}%
\pgfsetstrokeopacity{0.350000}%
\pgfsetdash{}{0pt}%
\pgfpathmoveto{\pgfqpoint{1.090230in}{0.610533in}}%
\pgfpathlineto{\pgfqpoint{1.099520in}{0.619706in}}%
\pgfpathlineto{\pgfqpoint{1.108814in}{0.628883in}}%
\pgfpathlineto{\pgfqpoint{1.118111in}{0.638063in}}%
\pgfpathlineto{\pgfqpoint{1.127411in}{0.647246in}}%
\pgfpathlineto{\pgfqpoint{1.136715in}{0.656433in}}%
\pgfpathlineto{\pgfqpoint{1.146022in}{0.665623in}}%
\pgfpathlineto{\pgfqpoint{1.155332in}{0.674816in}}%
\pgfpathlineto{\pgfqpoint{1.164646in}{0.684012in}}%
\pgfpathlineto{\pgfqpoint{1.173963in}{0.693212in}}%
\pgfpathlineto{\pgfqpoint{1.183283in}{0.702414in}}%
\pgfpathlineto{\pgfqpoint{1.192606in}{0.711621in}}%
\pgfpathlineto{\pgfqpoint{1.201933in}{0.720830in}}%
\pgfpathlineto{\pgfqpoint{1.211263in}{0.730043in}}%
\pgfpathlineto{\pgfqpoint{1.220597in}{0.739259in}}%
\pgfpathlineto{\pgfqpoint{1.229934in}{0.748478in}}%
\pgfpathlineto{\pgfqpoint{1.239274in}{0.757700in}}%
\pgfpathlineto{\pgfqpoint{1.248617in}{0.766926in}}%
\pgfpathlineto{\pgfqpoint{1.257964in}{0.776155in}}%
\pgfpathlineto{\pgfqpoint{1.267314in}{0.785388in}}%
\pgfpathlineto{\pgfqpoint{1.276668in}{0.794623in}}%
\pgfpathlineto{\pgfqpoint{1.286024in}{0.803862in}}%
\pgfpathlineto{\pgfqpoint{1.295384in}{0.813104in}}%
\pgfpathlineto{\pgfqpoint{1.304748in}{0.822350in}}%
\pgfpathlineto{\pgfqpoint{1.314115in}{0.831599in}}%
\pgfpathlineto{\pgfqpoint{1.323485in}{0.840851in}}%
\pgfpathlineto{\pgfqpoint{1.332858in}{0.850106in}}%
\pgfpathlineto{\pgfqpoint{1.342235in}{0.859365in}}%
\pgfpathlineto{\pgfqpoint{1.351615in}{0.868627in}}%
\pgfpathlineto{\pgfqpoint{1.360998in}{0.877892in}}%
\pgfpathlineto{\pgfqpoint{1.370385in}{0.887161in}}%
\pgfpathlineto{\pgfqpoint{1.379775in}{0.896433in}}%
\pgfpathlineto{\pgfqpoint{1.389169in}{0.905708in}}%
\pgfpathlineto{\pgfqpoint{1.398566in}{0.914986in}}%
\pgfpathlineto{\pgfqpoint{1.407966in}{0.924268in}}%
\pgfpathlineto{\pgfqpoint{1.417369in}{0.933553in}}%
\pgfpathlineto{\pgfqpoint{1.426776in}{0.942842in}}%
\pgfpathlineto{\pgfqpoint{1.436187in}{0.952133in}}%
\pgfpathlineto{\pgfqpoint{1.445600in}{0.961429in}}%
\pgfpathlineto{\pgfqpoint{1.455017in}{0.970727in}}%
\pgfpathlineto{\pgfqpoint{1.464438in}{0.980029in}}%
\pgfpathlineto{\pgfqpoint{1.473861in}{0.989334in}}%
\pgfpathlineto{\pgfqpoint{1.483288in}{0.998642in}}%
\pgfpathlineto{\pgfqpoint{1.492719in}{1.007954in}}%
\pgfpathlineto{\pgfqpoint{1.502153in}{1.017269in}}%
\pgfpathlineto{\pgfqpoint{1.511590in}{1.026587in}}%
\pgfpathlineto{\pgfqpoint{1.521031in}{1.035909in}}%
\pgfpathlineto{\pgfqpoint{1.530475in}{1.045234in}}%
\pgfpathlineto{\pgfqpoint{1.539922in}{1.054562in}}%
\pgfpathlineto{\pgfqpoint{1.549373in}{1.063894in}}%
\pgfpathlineto{\pgfqpoint{1.558827in}{1.073229in}}%
\pgfpathlineto{\pgfqpoint{1.568284in}{1.082568in}}%
\pgfpathlineto{\pgfqpoint{1.577745in}{1.091909in}}%
\pgfpathlineto{\pgfqpoint{1.587209in}{1.101254in}}%
\pgfpathlineto{\pgfqpoint{1.596677in}{1.110603in}}%
\pgfpathlineto{\pgfqpoint{1.606148in}{1.119955in}}%
\pgfpathlineto{\pgfqpoint{1.615622in}{1.129310in}}%
\pgfpathlineto{\pgfqpoint{1.625100in}{1.138668in}}%
\pgfpathlineto{\pgfqpoint{1.634581in}{1.148030in}}%
\pgfpathlineto{\pgfqpoint{1.644066in}{1.157395in}}%
\pgfpathlineto{\pgfqpoint{1.653554in}{1.166764in}}%
\pgfpathlineto{\pgfqpoint{1.663046in}{1.176136in}}%
\pgfpathlineto{\pgfqpoint{1.672540in}{1.185511in}}%
\pgfpathlineto{\pgfqpoint{1.682039in}{1.194890in}}%
\pgfpathlineto{\pgfqpoint{1.691540in}{1.204272in}}%
\pgfpathlineto{\pgfqpoint{1.701045in}{1.213657in}}%
\pgfpathlineto{\pgfqpoint{1.710554in}{1.223046in}}%
\pgfpathlineto{\pgfqpoint{1.720066in}{1.232438in}}%
\pgfpathlineto{\pgfqpoint{1.729581in}{1.241833in}}%
\pgfpathlineto{\pgfqpoint{1.739100in}{1.251232in}}%
\pgfpathlineto{\pgfqpoint{1.748622in}{1.260634in}}%
\pgfpathlineto{\pgfqpoint{1.758147in}{1.270040in}}%
\pgfpathlineto{\pgfqpoint{1.767676in}{1.279449in}}%
\pgfpathlineto{\pgfqpoint{1.777209in}{1.288861in}}%
\pgfpathlineto{\pgfqpoint{1.786745in}{1.298277in}}%
\pgfpathlineto{\pgfqpoint{1.796284in}{1.307696in}}%
\pgfpathlineto{\pgfqpoint{1.805827in}{1.317119in}}%
\pgfpathlineto{\pgfqpoint{1.815373in}{1.326545in}}%
\pgfpathlineto{\pgfqpoint{1.822088in}{1.333175in}}%
\pgfusepath{stroke}%
\end{pgfscope}%
\begin{pgfscope}%
\pgfpathrectangle{\pgfqpoint{0.407592in}{0.222904in}}{\pgfqpoint{1.411163in}{1.411163in}}%
\pgfusepath{clip}%
\pgfsetrectcap%
\pgfsetroundjoin%
\pgfsetlinewidth{0.401500pt}%
\definecolor{currentstroke}{rgb}{0.560784,0.588235,0.611765}%
\pgfsetstrokecolor{currentstroke}%
\pgfsetstrokeopacity{0.350000}%
\pgfsetdash{}{0pt}%
\pgfpathmoveto{\pgfqpoint{1.090230in}{0.610533in}}%
\pgfpathlineto{\pgfqpoint{1.095442in}{0.622265in}}%
\pgfpathlineto{\pgfqpoint{1.100649in}{0.633984in}}%
\pgfpathlineto{\pgfqpoint{1.105851in}{0.645690in}}%
\pgfpathlineto{\pgfqpoint{1.111047in}{0.657384in}}%
\pgfpathlineto{\pgfqpoint{1.116237in}{0.669065in}}%
\pgfpathlineto{\pgfqpoint{1.121422in}{0.680734in}}%
\pgfpathlineto{\pgfqpoint{1.126601in}{0.692390in}}%
\pgfpathlineto{\pgfqpoint{1.131774in}{0.704033in}}%
\pgfpathlineto{\pgfqpoint{1.136942in}{0.715664in}}%
\pgfpathlineto{\pgfqpoint{1.142105in}{0.727283in}}%
\pgfpathlineto{\pgfqpoint{1.147262in}{0.738889in}}%
\pgfpathlineto{\pgfqpoint{1.152413in}{0.750482in}}%
\pgfpathlineto{\pgfqpoint{1.157559in}{0.762063in}}%
\pgfpathlineto{\pgfqpoint{1.162699in}{0.773632in}}%
\pgfpathlineto{\pgfqpoint{1.167834in}{0.785188in}}%
\pgfpathlineto{\pgfqpoint{1.172963in}{0.796732in}}%
\pgfpathlineto{\pgfqpoint{1.178087in}{0.808263in}}%
\pgfpathlineto{\pgfqpoint{1.183205in}{0.819782in}}%
\pgfpathlineto{\pgfqpoint{1.188318in}{0.831289in}}%
\pgfpathlineto{\pgfqpoint{1.193425in}{0.842784in}}%
\pgfpathlineto{\pgfqpoint{1.198527in}{0.854266in}}%
\pgfpathlineto{\pgfqpoint{1.203623in}{0.865736in}}%
\pgfpathlineto{\pgfqpoint{1.208714in}{0.877194in}}%
\pgfpathlineto{\pgfqpoint{1.213800in}{0.888639in}}%
\pgfpathlineto{\pgfqpoint{1.218880in}{0.900072in}}%
\pgfpathlineto{\pgfqpoint{1.223955in}{0.911493in}}%
\pgfpathlineto{\pgfqpoint{1.229024in}{0.922902in}}%
\pgfpathlineto{\pgfqpoint{1.234088in}{0.934299in}}%
\pgfpathlineto{\pgfqpoint{1.239146in}{0.945684in}}%
\pgfpathlineto{\pgfqpoint{1.244200in}{0.957056in}}%
\pgfpathlineto{\pgfqpoint{1.249247in}{0.968417in}}%
\pgfpathlineto{\pgfqpoint{1.254290in}{0.979765in}}%
\pgfpathlineto{\pgfqpoint{1.259327in}{0.991101in}}%
\pgfpathlineto{\pgfqpoint{1.264358in}{1.002426in}}%
\pgfpathlineto{\pgfqpoint{1.269385in}{1.013738in}}%
\pgfpathlineto{\pgfqpoint{1.274406in}{1.025038in}}%
\pgfpathlineto{\pgfqpoint{1.279422in}{1.036326in}}%
\pgfpathlineto{\pgfqpoint{1.284432in}{1.047603in}}%
\pgfpathlineto{\pgfqpoint{1.289437in}{1.058867in}}%
\pgfpathlineto{\pgfqpoint{1.294437in}{1.070120in}}%
\pgfpathlineto{\pgfqpoint{1.299431in}{1.081360in}}%
\pgfpathlineto{\pgfqpoint{1.304421in}{1.092589in}}%
\pgfpathlineto{\pgfqpoint{1.309405in}{1.103806in}}%
\pgfpathlineto{\pgfqpoint{1.314383in}{1.115011in}}%
\pgfpathlineto{\pgfqpoint{1.319357in}{1.126204in}}%
\pgfpathlineto{\pgfqpoint{1.324325in}{1.137385in}}%
\pgfpathlineto{\pgfqpoint{1.329288in}{1.148555in}}%
\pgfpathlineto{\pgfqpoint{1.334246in}{1.159713in}}%
\pgfpathlineto{\pgfqpoint{1.339198in}{1.170859in}}%
\pgfpathlineto{\pgfqpoint{1.344145in}{1.181993in}}%
\pgfpathlineto{\pgfqpoint{1.349087in}{1.193116in}}%
\pgfpathlineto{\pgfqpoint{1.354024in}{1.204227in}}%
\pgfpathlineto{\pgfqpoint{1.358956in}{1.215326in}}%
\pgfpathlineto{\pgfqpoint{1.363883in}{1.226413in}}%
\pgfpathlineto{\pgfqpoint{1.368804in}{1.237489in}}%
\pgfpathlineto{\pgfqpoint{1.373720in}{1.248554in}}%
\pgfpathlineto{\pgfqpoint{1.378631in}{1.259606in}}%
\pgfpathlineto{\pgfqpoint{1.383537in}{1.270647in}}%
\pgfpathlineto{\pgfqpoint{1.388438in}{1.281677in}}%
\pgfpathlineto{\pgfqpoint{1.393333in}{1.292695in}}%
\pgfpathlineto{\pgfqpoint{1.398224in}{1.303702in}}%
\pgfpathlineto{\pgfqpoint{1.403109in}{1.314697in}}%
\pgfpathlineto{\pgfqpoint{1.407989in}{1.325680in}}%
\pgfpathlineto{\pgfqpoint{1.412865in}{1.336652in}}%
\pgfpathlineto{\pgfqpoint{1.417735in}{1.347613in}}%
\pgfpathlineto{\pgfqpoint{1.422600in}{1.358562in}}%
\pgfpathlineto{\pgfqpoint{1.427460in}{1.369499in}}%
\pgfpathlineto{\pgfqpoint{1.432314in}{1.380426in}}%
\pgfpathlineto{\pgfqpoint{1.437164in}{1.391341in}}%
\pgfpathlineto{\pgfqpoint{1.442009in}{1.402244in}}%
\pgfpathlineto{\pgfqpoint{1.446849in}{1.413136in}}%
\pgfpathlineto{\pgfqpoint{1.451683in}{1.424017in}}%
\pgfpathlineto{\pgfqpoint{1.456513in}{1.434887in}}%
\pgfpathlineto{\pgfqpoint{1.461338in}{1.445745in}}%
\pgfpathlineto{\pgfqpoint{1.466157in}{1.456592in}}%
\pgfpathlineto{\pgfqpoint{1.470972in}{1.467428in}}%
\pgfpathlineto{\pgfqpoint{1.475781in}{1.478252in}}%
\pgfpathlineto{\pgfqpoint{1.480586in}{1.489065in}}%
\pgfpathlineto{\pgfqpoint{1.485386in}{1.499867in}}%
\pgfusepath{stroke}%
\end{pgfscope}%
\begin{pgfscope}%
\pgfpathrectangle{\pgfqpoint{0.407592in}{0.222904in}}{\pgfqpoint{1.411163in}{1.411163in}}%
\pgfusepath{clip}%
\pgfsetrectcap%
\pgfsetroundjoin%
\pgfsetlinewidth{0.401500pt}%
\definecolor{currentstroke}{rgb}{0.560784,0.588235,0.611765}%
\pgfsetstrokecolor{currentstroke}%
\pgfsetstrokeopacity{0.350000}%
\pgfsetdash{}{0pt}%
\pgfpathmoveto{\pgfqpoint{1.090230in}{0.610533in}}%
\pgfpathlineto{\pgfqpoint{1.088307in}{0.622729in}}%
\pgfpathlineto{\pgfqpoint{1.086386in}{0.634909in}}%
\pgfpathlineto{\pgfqpoint{1.084468in}{0.647072in}}%
\pgfpathlineto{\pgfqpoint{1.082553in}{0.659219in}}%
\pgfpathlineto{\pgfqpoint{1.080640in}{0.671350in}}%
\pgfpathlineto{\pgfqpoint{1.078730in}{0.683464in}}%
\pgfpathlineto{\pgfqpoint{1.076822in}{0.695563in}}%
\pgfpathlineto{\pgfqpoint{1.074917in}{0.707645in}}%
\pgfpathlineto{\pgfqpoint{1.073015in}{0.719711in}}%
\pgfpathlineto{\pgfqpoint{1.071115in}{0.731762in}}%
\pgfpathlineto{\pgfqpoint{1.069217in}{0.743796in}}%
\pgfpathlineto{\pgfqpoint{1.067322in}{0.755814in}}%
\pgfpathlineto{\pgfqpoint{1.065430in}{0.767816in}}%
\pgfpathlineto{\pgfqpoint{1.063540in}{0.779802in}}%
\pgfpathlineto{\pgfqpoint{1.061653in}{0.791772in}}%
\pgfpathlineto{\pgfqpoint{1.059768in}{0.803726in}}%
\pgfpathlineto{\pgfqpoint{1.057885in}{0.815665in}}%
\pgfpathlineto{\pgfqpoint{1.056005in}{0.827587in}}%
\pgfpathlineto{\pgfqpoint{1.054128in}{0.839494in}}%
\pgfpathlineto{\pgfqpoint{1.052253in}{0.851385in}}%
\pgfpathlineto{\pgfqpoint{1.050381in}{0.863260in}}%
\pgfpathlineto{\pgfqpoint{1.048511in}{0.875120in}}%
\pgfpathlineto{\pgfqpoint{1.046643in}{0.886964in}}%
\pgfpathlineto{\pgfqpoint{1.044778in}{0.898792in}}%
\pgfpathlineto{\pgfqpoint{1.042915in}{0.910605in}}%
\pgfpathlineto{\pgfqpoint{1.041055in}{0.922402in}}%
\pgfpathlineto{\pgfqpoint{1.039198in}{0.934184in}}%
\pgfpathlineto{\pgfqpoint{1.037342in}{0.945950in}}%
\pgfpathlineto{\pgfqpoint{1.035490in}{0.957700in}}%
\pgfpathlineto{\pgfqpoint{1.033639in}{0.969435in}}%
\pgfpathlineto{\pgfqpoint{1.031791in}{0.981155in}}%
\pgfpathlineto{\pgfqpoint{1.029946in}{0.992859in}}%
\pgfpathlineto{\pgfqpoint{1.028103in}{1.004548in}}%
\pgfpathlineto{\pgfqpoint{1.026262in}{1.016222in}}%
\pgfpathlineto{\pgfqpoint{1.024424in}{1.027880in}}%
\pgfpathlineto{\pgfqpoint{1.022588in}{1.039523in}}%
\pgfpathlineto{\pgfqpoint{1.020755in}{1.051151in}}%
\pgfpathlineto{\pgfqpoint{1.018924in}{1.062763in}}%
\pgfpathlineto{\pgfqpoint{1.017095in}{1.074361in}}%
\pgfpathlineto{\pgfqpoint{1.015269in}{1.085943in}}%
\pgfpathlineto{\pgfqpoint{1.013445in}{1.097510in}}%
\pgfpathlineto{\pgfqpoint{1.011624in}{1.109062in}}%
\pgfpathlineto{\pgfqpoint{1.009804in}{1.120599in}}%
\pgfpathlineto{\pgfqpoint{1.007988in}{1.132121in}}%
\pgfpathlineto{\pgfqpoint{1.006173in}{1.143628in}}%
\pgfpathlineto{\pgfqpoint{1.004361in}{1.155120in}}%
\pgfpathlineto{\pgfqpoint{1.002552in}{1.166597in}}%
\pgfpathlineto{\pgfqpoint{1.000744in}{1.178059in}}%
\pgfpathlineto{\pgfqpoint{0.998940in}{1.189506in}}%
\pgfpathlineto{\pgfqpoint{0.997137in}{1.200938in}}%
\pgfpathlineto{\pgfqpoint{0.995337in}{1.212355in}}%
\pgfpathlineto{\pgfqpoint{0.993539in}{1.223758in}}%
\pgfpathlineto{\pgfqpoint{0.991743in}{1.235146in}}%
\pgfpathlineto{\pgfqpoint{0.989950in}{1.246519in}}%
\pgfpathlineto{\pgfqpoint{0.988159in}{1.257878in}}%
\pgfpathlineto{\pgfqpoint{0.986370in}{1.269221in}}%
\pgfpathlineto{\pgfqpoint{0.984584in}{1.280550in}}%
\pgfpathlineto{\pgfqpoint{0.982800in}{1.291865in}}%
\pgfpathlineto{\pgfqpoint{0.981018in}{1.303165in}}%
\pgfpathlineto{\pgfqpoint{0.979239in}{1.314450in}}%
\pgfpathlineto{\pgfqpoint{0.977462in}{1.325721in}}%
\pgfpathlineto{\pgfqpoint{0.975687in}{1.336977in}}%
\pgfpathlineto{\pgfqpoint{0.973914in}{1.348219in}}%
\pgfpathlineto{\pgfqpoint{0.972144in}{1.359446in}}%
\pgfpathlineto{\pgfqpoint{0.970376in}{1.370659in}}%
\pgfpathlineto{\pgfqpoint{0.968610in}{1.381857in}}%
\pgfpathlineto{\pgfqpoint{0.966847in}{1.393041in}}%
\pgfpathlineto{\pgfqpoint{0.965086in}{1.404211in}}%
\pgfpathlineto{\pgfqpoint{0.963327in}{1.415366in}}%
\pgfpathlineto{\pgfqpoint{0.961570in}{1.426507in}}%
\pgfpathlineto{\pgfqpoint{0.959816in}{1.437634in}}%
\pgfpathlineto{\pgfqpoint{0.958064in}{1.448746in}}%
\pgfpathlineto{\pgfqpoint{0.956314in}{1.459845in}}%
\pgfpathlineto{\pgfqpoint{0.954566in}{1.470929in}}%
\pgfpathlineto{\pgfqpoint{0.952820in}{1.481999in}}%
\pgfpathlineto{\pgfqpoint{0.951077in}{1.493055in}}%
\pgfpathlineto{\pgfqpoint{0.949336in}{1.504096in}}%
\pgfpathlineto{\pgfqpoint{0.947597in}{1.515124in}}%
\pgfpathlineto{\pgfqpoint{0.945861in}{1.526138in}}%
\pgfusepath{stroke}%
\end{pgfscope}%
\begin{pgfscope}%
\pgfpathrectangle{\pgfqpoint{0.407592in}{0.222904in}}{\pgfqpoint{1.411163in}{1.411163in}}%
\pgfusepath{clip}%
\pgfsetrectcap%
\pgfsetroundjoin%
\pgfsetlinewidth{0.401500pt}%
\definecolor{currentstroke}{rgb}{0.560784,0.588235,0.611765}%
\pgfsetstrokecolor{currentstroke}%
\pgfsetstrokeopacity{0.350000}%
\pgfsetdash{}{0pt}%
\pgfpathmoveto{\pgfqpoint{1.090230in}{0.610533in}}%
\pgfpathlineto{\pgfqpoint{1.082282in}{0.620829in}}%
\pgfpathlineto{\pgfqpoint{1.074337in}{0.631122in}}%
\pgfpathlineto{\pgfqpoint{1.066394in}{0.641412in}}%
\pgfpathlineto{\pgfqpoint{1.058454in}{0.651699in}}%
\pgfpathlineto{\pgfqpoint{1.050515in}{0.661983in}}%
\pgfpathlineto{\pgfqpoint{1.042579in}{0.672265in}}%
\pgfpathlineto{\pgfqpoint{1.034645in}{0.682544in}}%
\pgfpathlineto{\pgfqpoint{1.026713in}{0.692820in}}%
\pgfpathlineto{\pgfqpoint{1.018783in}{0.703093in}}%
\pgfpathlineto{\pgfqpoint{1.010856in}{0.713363in}}%
\pgfpathlineto{\pgfqpoint{1.002930in}{0.723631in}}%
\pgfpathlineto{\pgfqpoint{0.995007in}{0.733896in}}%
\pgfpathlineto{\pgfqpoint{0.987086in}{0.744157in}}%
\pgfpathlineto{\pgfqpoint{0.979167in}{0.754417in}}%
\pgfpathlineto{\pgfqpoint{0.971250in}{0.764673in}}%
\pgfpathlineto{\pgfqpoint{0.963336in}{0.774926in}}%
\pgfpathlineto{\pgfqpoint{0.955423in}{0.785177in}}%
\pgfpathlineto{\pgfqpoint{0.947513in}{0.795425in}}%
\pgfpathlineto{\pgfqpoint{0.939605in}{0.805670in}}%
\pgfpathlineto{\pgfqpoint{0.931699in}{0.815912in}}%
\pgfpathlineto{\pgfqpoint{0.923795in}{0.826151in}}%
\pgfpathlineto{\pgfqpoint{0.915894in}{0.836388in}}%
\pgfpathlineto{\pgfqpoint{0.907994in}{0.846622in}}%
\pgfpathlineto{\pgfqpoint{0.900097in}{0.856853in}}%
\pgfpathlineto{\pgfqpoint{0.892202in}{0.867081in}}%
\pgfpathlineto{\pgfqpoint{0.884309in}{0.877306in}}%
\pgfpathlineto{\pgfqpoint{0.876418in}{0.887529in}}%
\pgfpathlineto{\pgfqpoint{0.868530in}{0.897749in}}%
\pgfpathlineto{\pgfqpoint{0.860643in}{0.907966in}}%
\pgfpathlineto{\pgfqpoint{0.852759in}{0.918180in}}%
\pgfpathlineto{\pgfqpoint{0.844877in}{0.928391in}}%
\pgfpathlineto{\pgfqpoint{0.836997in}{0.938600in}}%
\pgfpathlineto{\pgfqpoint{0.829119in}{0.948806in}}%
\pgfpathlineto{\pgfqpoint{0.821244in}{0.959009in}}%
\pgfpathlineto{\pgfqpoint{0.813370in}{0.969209in}}%
\pgfpathlineto{\pgfqpoint{0.805499in}{0.979406in}}%
\pgfpathlineto{\pgfqpoint{0.797630in}{0.989601in}}%
\pgfpathlineto{\pgfqpoint{0.789763in}{0.999793in}}%
\pgfpathlineto{\pgfqpoint{0.781898in}{1.009982in}}%
\pgfpathlineto{\pgfqpoint{0.774035in}{1.020168in}}%
\pgfpathlineto{\pgfqpoint{0.766175in}{1.030352in}}%
\pgfpathlineto{\pgfqpoint{0.758316in}{1.040532in}}%
\pgfpathlineto{\pgfqpoint{0.750460in}{1.050710in}}%
\pgfpathlineto{\pgfqpoint{0.742606in}{1.060886in}}%
\pgfpathlineto{\pgfqpoint{0.734754in}{1.071058in}}%
\pgfpathlineto{\pgfqpoint{0.726904in}{1.081227in}}%
\pgfpathlineto{\pgfqpoint{0.719056in}{1.091394in}}%
\pgfpathlineto{\pgfqpoint{0.711211in}{1.101558in}}%
\pgfpathlineto{\pgfqpoint{0.703367in}{1.111720in}}%
\pgfpathlineto{\pgfqpoint{0.695526in}{1.121878in}}%
\pgfpathlineto{\pgfqpoint{0.687687in}{1.132034in}}%
\pgfpathlineto{\pgfqpoint{0.679850in}{1.142187in}}%
\pgfpathlineto{\pgfqpoint{0.672015in}{1.152337in}}%
\pgfpathlineto{\pgfqpoint{0.664182in}{1.162484in}}%
\pgfpathlineto{\pgfqpoint{0.656352in}{1.172629in}}%
\pgfpathlineto{\pgfqpoint{0.648523in}{1.182771in}}%
\pgfpathlineto{\pgfqpoint{0.640697in}{1.192910in}}%
\pgfpathlineto{\pgfqpoint{0.632873in}{1.203046in}}%
\pgfpathlineto{\pgfqpoint{0.625051in}{1.213180in}}%
\pgfpathlineto{\pgfqpoint{0.617231in}{1.223311in}}%
\pgfpathlineto{\pgfqpoint{0.609413in}{1.233439in}}%
\pgfpathlineto{\pgfqpoint{0.601598in}{1.243564in}}%
\pgfpathlineto{\pgfqpoint{0.593784in}{1.253687in}}%
\pgfpathlineto{\pgfqpoint{0.585973in}{1.263807in}}%
\pgfpathlineto{\pgfqpoint{0.578163in}{1.273924in}}%
\pgfpathlineto{\pgfqpoint{0.570356in}{1.284038in}}%
\pgfpathlineto{\pgfqpoint{0.562551in}{1.294149in}}%
\pgfpathlineto{\pgfqpoint{0.554748in}{1.304258in}}%
\pgfpathlineto{\pgfqpoint{0.546948in}{1.314364in}}%
\pgfpathlineto{\pgfqpoint{0.539149in}{1.324467in}}%
\pgfpathlineto{\pgfqpoint{0.531352in}{1.334568in}}%
\pgfpathlineto{\pgfqpoint{0.523558in}{1.344666in}}%
\pgfpathlineto{\pgfqpoint{0.515766in}{1.354761in}}%
\pgfpathlineto{\pgfqpoint{0.507976in}{1.364853in}}%
\pgfpathlineto{\pgfqpoint{0.500188in}{1.374942in}}%
\pgfpathlineto{\pgfqpoint{0.492402in}{1.385029in}}%
\pgfpathlineto{\pgfqpoint{0.484618in}{1.395113in}}%
\pgfpathlineto{\pgfqpoint{0.476836in}{1.405195in}}%
\pgfpathlineto{\pgfqpoint{0.469057in}{1.415273in}}%
\pgfusepath{stroke}%
\end{pgfscope}%
\begin{pgfscope}%
\pgfpathrectangle{\pgfqpoint{0.407592in}{0.222904in}}{\pgfqpoint{1.411163in}{1.411163in}}%
\pgfusepath{clip}%
\pgfsetrectcap%
\pgfsetroundjoin%
\pgfsetlinewidth{0.401500pt}%
\definecolor{currentstroke}{rgb}{0.560784,0.588235,0.611765}%
\pgfsetstrokecolor{currentstroke}%
\pgfsetstrokeopacity{0.350000}%
\pgfsetdash{}{0pt}%
\pgfpathmoveto{\pgfqpoint{1.090230in}{0.610533in}}%
\pgfpathlineto{\pgfqpoint{1.080895in}{0.617672in}}%
\pgfpathlineto{\pgfqpoint{1.071546in}{0.624823in}}%
\pgfpathlineto{\pgfqpoint{1.062184in}{0.631983in}}%
\pgfpathlineto{\pgfqpoint{1.052807in}{0.639155in}}%
\pgfpathlineto{\pgfqpoint{1.043416in}{0.646337in}}%
\pgfpathlineto{\pgfqpoint{1.034011in}{0.653531in}}%
\pgfpathlineto{\pgfqpoint{1.024592in}{0.660734in}}%
\pgfpathlineto{\pgfqpoint{1.015159in}{0.667949in}}%
\pgfpathlineto{\pgfqpoint{1.005712in}{0.675175in}}%
\pgfpathlineto{\pgfqpoint{0.996250in}{0.682411in}}%
\pgfpathlineto{\pgfqpoint{0.986775in}{0.689659in}}%
\pgfpathlineto{\pgfqpoint{0.977284in}{0.696917in}}%
\pgfpathlineto{\pgfqpoint{0.967780in}{0.704187in}}%
\pgfpathlineto{\pgfqpoint{0.958261in}{0.711467in}}%
\pgfpathlineto{\pgfqpoint{0.948728in}{0.718758in}}%
\pgfpathlineto{\pgfqpoint{0.939180in}{0.726060in}}%
\pgfpathlineto{\pgfqpoint{0.929618in}{0.733374in}}%
\pgfpathlineto{\pgfqpoint{0.920042in}{0.740698in}}%
\pgfpathlineto{\pgfqpoint{0.910451in}{0.748034in}}%
\pgfpathlineto{\pgfqpoint{0.900845in}{0.755381in}}%
\pgfpathlineto{\pgfqpoint{0.891225in}{0.762739in}}%
\pgfpathlineto{\pgfqpoint{0.881590in}{0.770108in}}%
\pgfpathlineto{\pgfqpoint{0.871941in}{0.777488in}}%
\pgfpathlineto{\pgfqpoint{0.862276in}{0.784879in}}%
\pgfpathlineto{\pgfqpoint{0.852597in}{0.792282in}}%
\pgfpathlineto{\pgfqpoint{0.842904in}{0.799696in}}%
\pgfpathlineto{\pgfqpoint{0.833195in}{0.807121in}}%
\pgfpathlineto{\pgfqpoint{0.823472in}{0.814558in}}%
\pgfpathlineto{\pgfqpoint{0.813734in}{0.822006in}}%
\pgfpathlineto{\pgfqpoint{0.803981in}{0.829466in}}%
\pgfpathlineto{\pgfqpoint{0.794213in}{0.836936in}}%
\pgfpathlineto{\pgfqpoint{0.784430in}{0.844419in}}%
\pgfpathlineto{\pgfqpoint{0.774632in}{0.851912in}}%
\pgfpathlineto{\pgfqpoint{0.764819in}{0.859418in}}%
\pgfpathlineto{\pgfqpoint{0.754991in}{0.866934in}}%
\pgfpathlineto{\pgfqpoint{0.745148in}{0.874463in}}%
\pgfpathlineto{\pgfqpoint{0.735290in}{0.882003in}}%
\pgfpathlineto{\pgfqpoint{0.725417in}{0.889554in}}%
\pgfpathlineto{\pgfqpoint{0.715528in}{0.897117in}}%
\pgfpathlineto{\pgfqpoint{0.705624in}{0.904692in}}%
\pgfpathlineto{\pgfqpoint{0.695705in}{0.912278in}}%
\pgfpathlineto{\pgfqpoint{0.685771in}{0.919877in}}%
\pgfpathlineto{\pgfqpoint{0.675821in}{0.927487in}}%
\pgfpathlineto{\pgfqpoint{0.665856in}{0.935108in}}%
\pgfpathlineto{\pgfqpoint{0.655876in}{0.942742in}}%
\pgfpathlineto{\pgfqpoint{0.645880in}{0.950387in}}%
\pgfpathlineto{\pgfqpoint{0.635868in}{0.958044in}}%
\pgfpathlineto{\pgfqpoint{0.625841in}{0.965713in}}%
\pgfpathlineto{\pgfqpoint{0.615799in}{0.973394in}}%
\pgfpathlineto{\pgfqpoint{0.605741in}{0.981087in}}%
\pgfpathlineto{\pgfqpoint{0.595667in}{0.988791in}}%
\pgfpathlineto{\pgfqpoint{0.585578in}{0.996508in}}%
\pgfpathlineto{\pgfqpoint{0.575472in}{1.004237in}}%
\pgfpathlineto{\pgfqpoint{0.565352in}{1.011978in}}%
\pgfpathlineto{\pgfqpoint{0.555215in}{1.019730in}}%
\pgfpathlineto{\pgfqpoint{0.545062in}{1.027495in}}%
\pgfpathlineto{\pgfqpoint{0.534894in}{1.035272in}}%
\pgfpathlineto{\pgfqpoint{0.524710in}{1.043062in}}%
\pgfpathlineto{\pgfqpoint{0.514510in}{1.050863in}}%
\pgfpathlineto{\pgfqpoint{0.504294in}{1.058677in}}%
\pgfpathlineto{\pgfqpoint{0.494062in}{1.066502in}}%
\pgfpathlineto{\pgfqpoint{0.483814in}{1.074341in}}%
\pgfpathlineto{\pgfqpoint{0.473549in}{1.082191in}}%
\pgfpathlineto{\pgfqpoint{0.463269in}{1.090054in}}%
\pgfpathlineto{\pgfqpoint{0.452973in}{1.097929in}}%
\pgfpathlineto{\pgfqpoint{0.442660in}{1.105816in}}%
\pgfpathlineto{\pgfqpoint{0.432331in}{1.113716in}}%
\pgfpathlineto{\pgfqpoint{0.421986in}{1.121628in}}%
\pgfpathlineto{\pgfqpoint{0.411625in}{1.129553in}}%
\pgfpathlineto{\pgfqpoint{0.404258in}{1.135187in}}%
\pgfusepath{stroke}%
\end{pgfscope}%
\begin{pgfscope}%
\pgfpathrectangle{\pgfqpoint{0.407592in}{0.222904in}}{\pgfqpoint{1.411163in}{1.411163in}}%
\pgfusepath{clip}%
\pgfsetrectcap%
\pgfsetroundjoin%
\pgfsetlinewidth{0.401500pt}%
\definecolor{currentstroke}{rgb}{0.560784,0.588235,0.611765}%
\pgfsetstrokecolor{currentstroke}%
\pgfsetstrokeopacity{0.350000}%
\pgfsetdash{}{0pt}%
\pgfpathmoveto{\pgfqpoint{1.090230in}{0.610533in}}%
\pgfpathlineto{\pgfqpoint{1.084967in}{0.615108in}}%
\pgfpathlineto{\pgfqpoint{1.079690in}{0.619696in}}%
\pgfpathlineto{\pgfqpoint{1.074397in}{0.624297in}}%
\pgfpathlineto{\pgfqpoint{1.069088in}{0.628912in}}%
\pgfpathlineto{\pgfqpoint{1.063764in}{0.633541in}}%
\pgfpathlineto{\pgfqpoint{1.058424in}{0.638184in}}%
\pgfpathlineto{\pgfqpoint{1.053068in}{0.642840in}}%
\pgfpathlineto{\pgfqpoint{1.047696in}{0.647510in}}%
\pgfpathlineto{\pgfqpoint{1.042308in}{0.652194in}}%
\pgfpathlineto{\pgfqpoint{1.036904in}{0.656891in}}%
\pgfpathlineto{\pgfqpoint{1.031484in}{0.661603in}}%
\pgfpathlineto{\pgfqpoint{1.026048in}{0.666329in}}%
\pgfpathlineto{\pgfqpoint{1.020596in}{0.671069in}}%
\pgfpathlineto{\pgfqpoint{1.015127in}{0.675823in}}%
\pgfpathlineto{\pgfqpoint{1.009642in}{0.680592in}}%
\pgfpathlineto{\pgfqpoint{1.004141in}{0.685374in}}%
\pgfpathlineto{\pgfqpoint{0.998623in}{0.690171in}}%
\pgfpathlineto{\pgfqpoint{0.993088in}{0.694983in}}%
\pgfpathlineto{\pgfqpoint{0.987537in}{0.699809in}}%
\pgfpathlineto{\pgfqpoint{0.981969in}{0.704649in}}%
\pgfpathlineto{\pgfqpoint{0.976384in}{0.709505in}}%
\pgfpathlineto{\pgfqpoint{0.970782in}{0.714375in}}%
\pgfpathlineto{\pgfqpoint{0.965164in}{0.719259in}}%
\pgfpathlineto{\pgfqpoint{0.959528in}{0.724159in}}%
\pgfpathlineto{\pgfqpoint{0.953875in}{0.729073in}}%
\pgfpathlineto{\pgfqpoint{0.948205in}{0.734003in}}%
\pgfpathlineto{\pgfqpoint{0.942517in}{0.738947in}}%
\pgfpathlineto{\pgfqpoint{0.936812in}{0.743907in}}%
\pgfpathlineto{\pgfqpoint{0.931090in}{0.748881in}}%
\pgfpathlineto{\pgfqpoint{0.925350in}{0.753871in}}%
\pgfpathlineto{\pgfqpoint{0.919592in}{0.758877in}}%
\pgfpathlineto{\pgfqpoint{0.913817in}{0.763897in}}%
\pgfpathlineto{\pgfqpoint{0.908024in}{0.768933in}}%
\pgfpathlineto{\pgfqpoint{0.902213in}{0.773985in}}%
\pgfpathlineto{\pgfqpoint{0.896385in}{0.779052in}}%
\pgfpathlineto{\pgfqpoint{0.890538in}{0.784135in}}%
\pgfpathlineto{\pgfqpoint{0.884673in}{0.789234in}}%
\pgfpathlineto{\pgfqpoint{0.878790in}{0.794349in}}%
\pgfpathlineto{\pgfqpoint{0.872888in}{0.799479in}}%
\pgfpathlineto{\pgfqpoint{0.866968in}{0.804625in}}%
\pgfpathlineto{\pgfqpoint{0.861030in}{0.809788in}}%
\pgfpathlineto{\pgfqpoint{0.855073in}{0.814966in}}%
\pgfpathlineto{\pgfqpoint{0.849098in}{0.820161in}}%
\pgfpathlineto{\pgfqpoint{0.843104in}{0.825372in}}%
\pgfpathlineto{\pgfqpoint{0.837091in}{0.830600in}}%
\pgfpathlineto{\pgfqpoint{0.831059in}{0.835843in}}%
\pgfpathlineto{\pgfqpoint{0.825008in}{0.841104in}}%
\pgfpathlineto{\pgfqpoint{0.818938in}{0.846381in}}%
\pgfpathlineto{\pgfqpoint{0.812849in}{0.851674in}}%
\pgfpathlineto{\pgfqpoint{0.806741in}{0.856984in}}%
\pgfpathlineto{\pgfqpoint{0.800613in}{0.862311in}}%
\pgfpathlineto{\pgfqpoint{0.794466in}{0.867655in}}%
\pgfpathlineto{\pgfqpoint{0.788299in}{0.873016in}}%
\pgfpathlineto{\pgfqpoint{0.782113in}{0.878394in}}%
\pgfpathlineto{\pgfqpoint{0.775907in}{0.883790in}}%
\pgfpathlineto{\pgfqpoint{0.769681in}{0.889202in}}%
\pgfpathlineto{\pgfqpoint{0.763436in}{0.894632in}}%
\pgfpathlineto{\pgfqpoint{0.757170in}{0.900079in}}%
\pgfpathlineto{\pgfqpoint{0.750884in}{0.905543in}}%
\pgfpathlineto{\pgfqpoint{0.744578in}{0.911026in}}%
\pgfpathlineto{\pgfqpoint{0.738252in}{0.916525in}}%
\pgfpathlineto{\pgfqpoint{0.731905in}{0.922043in}}%
\pgfpathlineto{\pgfqpoint{0.725538in}{0.927578in}}%
\pgfpathlineto{\pgfqpoint{0.719150in}{0.933131in}}%
\pgfpathlineto{\pgfqpoint{0.712742in}{0.938703in}}%
\pgfpathlineto{\pgfqpoint{0.706312in}{0.944292in}}%
\pgfpathlineto{\pgfqpoint{0.699862in}{0.949899in}}%
\pgfpathlineto{\pgfqpoint{0.693391in}{0.955525in}}%
\pgfpathlineto{\pgfqpoint{0.686899in}{0.961169in}}%
\pgfpathlineto{\pgfqpoint{0.680385in}{0.966831in}}%
\pgfpathlineto{\pgfqpoint{0.673851in}{0.972512in}}%
\pgfpathlineto{\pgfqpoint{0.667295in}{0.978212in}}%
\pgfpathlineto{\pgfqpoint{0.660717in}{0.983930in}}%
\pgfpathlineto{\pgfqpoint{0.654118in}{0.989667in}}%
\pgfpathlineto{\pgfqpoint{0.647497in}{0.995423in}}%
\pgfpathlineto{\pgfqpoint{0.640854in}{1.001198in}}%
\pgfpathlineto{\pgfqpoint{0.634189in}{1.006992in}}%
\pgfpathlineto{\pgfqpoint{0.627502in}{1.012806in}}%
\pgfpathlineto{\pgfqpoint{0.620793in}{1.018638in}}%
\pgfusepath{stroke}%
\end{pgfscope}%
\begin{pgfscope}%
\pgfpathrectangle{\pgfqpoint{0.407592in}{0.222904in}}{\pgfqpoint{1.411163in}{1.411163in}}%
\pgfusepath{clip}%
\pgfsetrectcap%
\pgfsetroundjoin%
\pgfsetlinewidth{0.401500pt}%
\definecolor{currentstroke}{rgb}{0.560784,0.588235,0.611765}%
\pgfsetstrokecolor{currentstroke}%
\pgfsetstrokeopacity{0.350000}%
\pgfsetdash{}{0pt}%
\pgfpathmoveto{\pgfqpoint{1.090230in}{0.610533in}}%
\pgfpathlineto{\pgfqpoint{1.092118in}{0.614641in}}%
\pgfpathlineto{\pgfqpoint{1.094013in}{0.618763in}}%
\pgfpathlineto{\pgfqpoint{1.095913in}{0.622898in}}%
\pgfpathlineto{\pgfqpoint{1.097820in}{0.627046in}}%
\pgfpathlineto{\pgfqpoint{1.099733in}{0.631207in}}%
\pgfpathlineto{\pgfqpoint{1.101652in}{0.635382in}}%
\pgfpathlineto{\pgfqpoint{1.103577in}{0.639570in}}%
\pgfpathlineto{\pgfqpoint{1.105508in}{0.643772in}}%
\pgfpathlineto{\pgfqpoint{1.107446in}{0.647987in}}%
\pgfpathlineto{\pgfqpoint{1.109390in}{0.652216in}}%
\pgfpathlineto{\pgfqpoint{1.111340in}{0.656459in}}%
\pgfpathlineto{\pgfqpoint{1.113297in}{0.660715in}}%
\pgfpathlineto{\pgfqpoint{1.115260in}{0.664986in}}%
\pgfpathlineto{\pgfqpoint{1.117229in}{0.669270in}}%
\pgfpathlineto{\pgfqpoint{1.119205in}{0.673569in}}%
\pgfpathlineto{\pgfqpoint{1.121187in}{0.677881in}}%
\pgfpathlineto{\pgfqpoint{1.123176in}{0.682208in}}%
\pgfpathlineto{\pgfqpoint{1.125171in}{0.686549in}}%
\pgfpathlineto{\pgfqpoint{1.127173in}{0.690904in}}%
\pgfpathlineto{\pgfqpoint{1.129182in}{0.695273in}}%
\pgfpathlineto{\pgfqpoint{1.131197in}{0.699657in}}%
\pgfpathlineto{\pgfqpoint{1.133219in}{0.704056in}}%
\pgfpathlineto{\pgfqpoint{1.135247in}{0.708469in}}%
\pgfpathlineto{\pgfqpoint{1.137282in}{0.712896in}}%
\pgfpathlineto{\pgfqpoint{1.139324in}{0.717338in}}%
\pgfpathlineto{\pgfqpoint{1.141373in}{0.721796in}}%
\pgfpathlineto{\pgfqpoint{1.143429in}{0.726268in}}%
\pgfpathlineto{\pgfqpoint{1.145491in}{0.730755in}}%
\pgfpathlineto{\pgfqpoint{1.147561in}{0.735256in}}%
\pgfpathlineto{\pgfqpoint{1.149637in}{0.739774in}}%
\pgfpathlineto{\pgfqpoint{1.151720in}{0.744306in}}%
\pgfpathlineto{\pgfqpoint{1.153811in}{0.748853in}}%
\pgfpathlineto{\pgfqpoint{1.155908in}{0.753416in}}%
\pgfpathlineto{\pgfqpoint{1.158012in}{0.757994in}}%
\pgfpathlineto{\pgfqpoint{1.160124in}{0.762588in}}%
\pgfpathlineto{\pgfqpoint{1.162242in}{0.767197in}}%
\pgfpathlineto{\pgfqpoint{1.164368in}{0.771822in}}%
\pgfpathlineto{\pgfqpoint{1.166501in}{0.776462in}}%
\pgfpathlineto{\pgfqpoint{1.168642in}{0.781119in}}%
\pgfpathlineto{\pgfqpoint{1.170789in}{0.785791in}}%
\pgfpathlineto{\pgfqpoint{1.172944in}{0.790479in}}%
\pgfpathlineto{\pgfqpoint{1.175107in}{0.795183in}}%
\pgfpathlineto{\pgfqpoint{1.177276in}{0.799903in}}%
\pgfpathlineto{\pgfqpoint{1.179454in}{0.804640in}}%
\pgfpathlineto{\pgfqpoint{1.181638in}{0.809393in}}%
\pgfpathlineto{\pgfqpoint{1.183831in}{0.814162in}}%
\pgfpathlineto{\pgfqpoint{1.186030in}{0.818947in}}%
\pgfpathlineto{\pgfqpoint{1.188238in}{0.823750in}}%
\pgfpathlineto{\pgfqpoint{1.190453in}{0.828568in}}%
\pgfpathlineto{\pgfqpoint{1.192675in}{0.833404in}}%
\pgfpathlineto{\pgfqpoint{1.194906in}{0.838256in}}%
\pgfpathlineto{\pgfqpoint{1.197144in}{0.843125in}}%
\pgfpathlineto{\pgfqpoint{1.199390in}{0.848011in}}%
\pgfpathlineto{\pgfqpoint{1.201644in}{0.852915in}}%
\pgfpathlineto{\pgfqpoint{1.203906in}{0.857835in}}%
\pgfpathlineto{\pgfqpoint{1.206175in}{0.862773in}}%
\pgfpathlineto{\pgfqpoint{1.208453in}{0.867727in}}%
\pgfpathlineto{\pgfqpoint{1.210738in}{0.872700in}}%
\pgfpathlineto{\pgfqpoint{1.213032in}{0.877690in}}%
\pgfpathlineto{\pgfqpoint{1.215334in}{0.882697in}}%
\pgfpathlineto{\pgfqpoint{1.217644in}{0.887722in}}%
\pgfpathlineto{\pgfqpoint{1.219962in}{0.892765in}}%
\pgfpathlineto{\pgfqpoint{1.222288in}{0.897826in}}%
\pgfpathlineto{\pgfqpoint{1.224623in}{0.902905in}}%
\pgfpathlineto{\pgfqpoint{1.226966in}{0.908002in}}%
\pgfpathlineto{\pgfqpoint{1.229317in}{0.913117in}}%
\pgfpathlineto{\pgfqpoint{1.231676in}{0.918250in}}%
\pgfpathlineto{\pgfqpoint{1.234045in}{0.923402in}}%
\pgfpathlineto{\pgfqpoint{1.236421in}{0.928572in}}%
\pgfpathlineto{\pgfqpoint{1.238806in}{0.933761in}}%
\pgfpathlineto{\pgfqpoint{1.241200in}{0.938969in}}%
\pgfpathlineto{\pgfqpoint{1.243602in}{0.944195in}}%
\pgfpathlineto{\pgfqpoint{1.246013in}{0.949440in}}%
\pgfpathlineto{\pgfqpoint{1.248433in}{0.954704in}}%
\pgfpathlineto{\pgfqpoint{1.250862in}{0.959988in}}%
\pgfpathlineto{\pgfqpoint{1.253299in}{0.965290in}}%
\pgfpathlineto{\pgfqpoint{1.255745in}{0.970612in}}%
\pgfpathlineto{\pgfqpoint{1.258200in}{0.975953in}}%
\pgfpathlineto{\pgfqpoint{1.260664in}{0.981313in}}%
\pgfusepath{stroke}%
\end{pgfscope}%
\begin{pgfscope}%
\pgfpathrectangle{\pgfqpoint{0.407592in}{0.222904in}}{\pgfqpoint{1.411163in}{1.411163in}}%
\pgfusepath{clip}%
\pgfsetrectcap%
\pgfsetroundjoin%
\pgfsetlinewidth{1.806750pt}%
\definecolor{currentstroke}{rgb}{0.168627,0.611765,0.411765}%
\pgfsetstrokecolor{currentstroke}%
\pgfsetdash{}{0pt}%
\pgfpathmoveto{\pgfqpoint{1.247522in}{0.897771in}}%
\pgfpathlineto{\pgfqpoint{1.264082in}{0.890537in}}%
\pgfpathlineto{\pgfqpoint{1.280506in}{0.884135in}}%
\pgfpathlineto{\pgfqpoint{1.296799in}{0.878629in}}%
\pgfpathlineto{\pgfqpoint{1.310664in}{0.874677in}}%
\pgfpathlineto{\pgfqpoint{1.324439in}{0.871482in}}%
\pgfpathlineto{\pgfqpoint{1.338128in}{0.869093in}}%
\pgfpathlineto{\pgfqpoint{1.351734in}{0.867560in}}%
\pgfpathlineto{\pgfqpoint{1.363012in}{0.866973in}}%
\pgfpathlineto{\pgfqpoint{1.374236in}{0.867045in}}%
\pgfpathlineto{\pgfqpoint{1.385410in}{0.867806in}}%
\pgfpathlineto{\pgfqpoint{1.396535in}{0.869281in}}%
\pgfpathlineto{\pgfqpoint{1.407613in}{0.871498in}}%
\pgfpathlineto{\pgfqpoint{1.418647in}{0.874479in}}%
\pgfpathlineto{\pgfqpoint{1.429639in}{0.878244in}}%
\pgfpathlineto{\pgfqpoint{1.440591in}{0.882810in}}%
\pgfpathlineto{\pgfqpoint{1.451505in}{0.888189in}}%
\pgfpathlineto{\pgfqpoint{1.462383in}{0.894389in}}%
\pgfpathlineto{\pgfqpoint{1.473228in}{0.901413in}}%
\pgfpathlineto{\pgfqpoint{1.484040in}{0.909259in}}%
\pgfpathlineto{\pgfqpoint{1.494821in}{0.917918in}}%
\pgfpathlineto{\pgfqpoint{1.505574in}{0.927379in}}%
\pgfpathlineto{\pgfqpoint{1.518442in}{0.939765in}}%
\pgfpathlineto{\pgfqpoint{1.531271in}{0.953241in}}%
\pgfpathlineto{\pgfqpoint{1.544065in}{0.967763in}}%
\pgfpathlineto{\pgfqpoint{1.556825in}{0.983277in}}%
\pgfpathlineto{\pgfqpoint{1.571670in}{1.002559in}}%
\pgfpathlineto{\pgfqpoint{1.586472in}{1.023023in}}%
\pgfpathlineto{\pgfqpoint{1.603337in}{1.047737in}}%
\pgfpathlineto{\pgfqpoint{1.620147in}{1.073730in}}%
\pgfpathlineto{\pgfqpoint{1.638995in}{1.104332in}}%
\pgfpathlineto{\pgfqpoint{1.657776in}{1.136199in}}%
\pgfpathlineto{\pgfqpoint{1.678565in}{1.172889in}}%
\pgfpathlineto{\pgfqpoint{1.682713in}{1.180372in}}%
\pgfpathlineto{\pgfqpoint{1.682713in}{1.180372in}}%
\pgfusepath{stroke}%
\end{pgfscope}%
\begin{pgfscope}%
\pgfpathrectangle{\pgfqpoint{0.407592in}{0.222904in}}{\pgfqpoint{1.411163in}{1.411163in}}%
\pgfusepath{clip}%
\pgfsetrectcap%
\pgfsetroundjoin%
\pgfsetlinewidth{2.308625pt}%
\definecolor{currentstroke}{rgb}{0.141176,0.345098,0.650980}%
\pgfsetstrokecolor{currentstroke}%
\pgfsetdash{}{0pt}%
\pgfpathmoveto{\pgfqpoint{0.537673in}{1.095850in}}%
\pgfpathlineto{\pgfqpoint{0.583281in}{1.045325in}}%
\pgfpathlineto{\pgfqpoint{0.623058in}{1.002745in}}%
\pgfpathlineto{\pgfqpoint{0.658171in}{0.966597in}}%
\pgfpathlineto{\pgfqpoint{0.689504in}{0.935734in}}%
\pgfpathlineto{\pgfqpoint{0.717737in}{0.909269in}}%
\pgfpathlineto{\pgfqpoint{0.743404in}{0.886509in}}%
\pgfpathlineto{\pgfqpoint{0.766930in}{0.866900in}}%
\pgfpathlineto{\pgfqpoint{0.788654in}{0.849996in}}%
\pgfpathlineto{\pgfqpoint{0.808853in}{0.835433in}}%
\pgfpathlineto{\pgfqpoint{0.827752in}{0.822909in}}%
\pgfpathlineto{\pgfqpoint{0.845537in}{0.812174in}}%
\pgfpathlineto{\pgfqpoint{0.862362in}{0.803018in}}%
\pgfpathlineto{\pgfqpoint{0.878355in}{0.795260in}}%
\pgfpathlineto{\pgfqpoint{0.893624in}{0.788747in}}%
\pgfpathlineto{\pgfqpoint{0.908259in}{0.783345in}}%
\pgfpathlineto{\pgfqpoint{0.922333in}{0.778940in}}%
\pgfpathlineto{\pgfqpoint{0.935911in}{0.775428in}}%
\pgfpathlineto{\pgfqpoint{0.949044in}{0.772720in}}%
\pgfpathlineto{\pgfqpoint{0.961776in}{0.770736in}}%
\pgfpathlineto{\pgfqpoint{0.974145in}{0.769403in}}%
\pgfpathlineto{\pgfqpoint{0.990120in}{0.768527in}}%
\pgfpathlineto{\pgfqpoint{1.005557in}{0.768557in}}%
\pgfpathlineto{\pgfqpoint{1.020497in}{0.769369in}}%
\pgfpathlineto{\pgfqpoint{1.034971in}{0.770855in}}%
\pgfpathlineto{\pgfqpoint{1.052439in}{0.773510in}}%
\pgfpathlineto{\pgfqpoint{1.069237in}{0.776894in}}%
\pgfpathlineto{\pgfqpoint{1.085379in}{0.780854in}}%
\pgfpathlineto{\pgfqpoint{1.103888in}{0.786181in}}%
\pgfpathlineto{\pgfqpoint{1.124289in}{0.792937in}}%
\pgfpathlineto{\pgfqpoint{1.146015in}{0.801046in}}%
\pgfpathlineto{\pgfqpoint{1.168412in}{0.810280in}}%
\pgfpathlineto{\pgfqpoint{1.192850in}{0.821246in}}%
\pgfpathlineto{\pgfqpoint{1.217673in}{0.833215in}}%
\pgfpathlineto{\pgfqpoint{1.245306in}{0.847389in}}%
\pgfpathlineto{\pgfqpoint{1.275238in}{0.863608in}}%
\pgfpathlineto{\pgfqpoint{1.307959in}{0.882211in}}%
\pgfpathlineto{\pgfqpoint{1.318296in}{0.888253in}}%
\pgfpathlineto{\pgfqpoint{1.318296in}{0.888253in}}%
\pgfusepath{stroke}%
\end{pgfscope}%
\begin{pgfscope}%
\pgfpathrectangle{\pgfqpoint{0.407592in}{0.222904in}}{\pgfqpoint{1.411163in}{1.411163in}}%
\pgfusepath{clip}%
\pgfsetrectcap%
\pgfsetroundjoin%
\pgfsetlinewidth{2.308625pt}%
\definecolor{currentstroke}{rgb}{0.819608,0.286275,0.356863}%
\pgfsetstrokecolor{currentstroke}%
\pgfsetdash{}{0pt}%
\pgfpathmoveto{\pgfqpoint{0.672155in}{1.057312in}}%
\pgfpathlineto{\pgfqpoint{0.708311in}{1.016342in}}%
\pgfpathlineto{\pgfqpoint{0.740163in}{0.981221in}}%
\pgfpathlineto{\pgfqpoint{0.768505in}{0.950926in}}%
\pgfpathlineto{\pgfqpoint{0.793952in}{0.924665in}}%
\pgfpathlineto{\pgfqpoint{0.816990in}{0.901810in}}%
\pgfpathlineto{\pgfqpoint{0.838008in}{0.881864in}}%
\pgfpathlineto{\pgfqpoint{0.857320in}{0.864421in}}%
\pgfpathlineto{\pgfqpoint{0.875184in}{0.849151in}}%
\pgfpathlineto{\pgfqpoint{0.891812in}{0.835782in}}%
\pgfpathlineto{\pgfqpoint{0.907379in}{0.824087in}}%
\pgfpathlineto{\pgfqpoint{0.922035in}{0.813875in}}%
\pgfpathlineto{\pgfqpoint{0.935904in}{0.804984in}}%
\pgfpathlineto{\pgfqpoint{0.949090in}{0.797275in}}%
\pgfpathlineto{\pgfqpoint{0.961686in}{0.790631in}}%
\pgfpathlineto{\pgfqpoint{0.973766in}{0.784947in}}%
\pgfpathlineto{\pgfqpoint{0.985398in}{0.780135in}}%
\pgfpathlineto{\pgfqpoint{0.996636in}{0.776115in}}%
\pgfpathlineto{\pgfqpoint{1.007531in}{0.772818in}}%
\pgfpathlineto{\pgfqpoint{1.021591in}{0.769438in}}%
\pgfpathlineto{\pgfqpoint{1.035194in}{0.767104in}}%
\pgfpathlineto{\pgfqpoint{1.048406in}{0.765700in}}%
\pgfpathlineto{\pgfqpoint{1.061278in}{0.765128in}}%
\pgfpathlineto{\pgfqpoint{1.073852in}{0.765297in}}%
\pgfpathlineto{\pgfqpoint{1.086159in}{0.766126in}}%
\pgfpathlineto{\pgfqpoint{1.101203in}{0.767980in}}%
\pgfpathlineto{\pgfqpoint{1.115895in}{0.770622in}}%
\pgfpathlineto{\pgfqpoint{1.130252in}{0.773931in}}%
\pgfpathlineto{\pgfqpoint{1.147043in}{0.778635in}}%
\pgfpathlineto{\pgfqpoint{1.163349in}{0.783976in}}%
\pgfpathlineto{\pgfqpoint{1.181730in}{0.790809in}}%
\pgfpathlineto{\pgfqpoint{1.201832in}{0.799150in}}%
\pgfpathlineto{\pgfqpoint{1.223199in}{0.808874in}}%
\pgfpathlineto{\pgfqpoint{1.245264in}{0.819713in}}%
\pgfpathlineto{\pgfqpoint{1.271139in}{0.833297in}}%
\pgfpathlineto{\pgfqpoint{1.298486in}{0.848510in}}%
\pgfpathlineto{\pgfqpoint{1.329398in}{0.866578in}}%
\pgfpathlineto{\pgfqpoint{1.364170in}{0.887801in}}%
\pgfpathlineto{\pgfqpoint{1.396682in}{0.908341in}}%
\pgfpathlineto{\pgfqpoint{1.396682in}{0.908341in}}%
\pgfusepath{stroke}%
\end{pgfscope}%
\begin{pgfscope}%
\pgfpathrectangle{\pgfqpoint{0.407592in}{0.222904in}}{\pgfqpoint{1.411163in}{1.411163in}}%
\pgfusepath{clip}%
\pgfsetrectcap%
\pgfsetroundjoin%
\pgfsetlinewidth{2.308625pt}%
\definecolor{currentstroke}{rgb}{0.000000,0.529412,0.423529}%
\pgfsetstrokecolor{currentstroke}%
\pgfsetdash{}{0pt}%
\pgfpathmoveto{\pgfqpoint{0.812524in}{1.057517in}}%
\pgfpathlineto{\pgfqpoint{0.843065in}{1.026044in}}%
\pgfpathlineto{\pgfqpoint{0.870461in}{0.998966in}}%
\pgfpathlineto{\pgfqpoint{0.895259in}{0.975582in}}%
\pgfpathlineto{\pgfqpoint{0.917889in}{0.955335in}}%
\pgfpathlineto{\pgfqpoint{0.938697in}{0.937777in}}%
\pgfpathlineto{\pgfqpoint{0.957966in}{0.922543in}}%
\pgfpathlineto{\pgfqpoint{0.975926in}{0.909335in}}%
\pgfpathlineto{\pgfqpoint{0.992767in}{0.897903in}}%
\pgfpathlineto{\pgfqpoint{1.008646in}{0.888038in}}%
\pgfpathlineto{\pgfqpoint{1.023697in}{0.879564in}}%
\pgfpathlineto{\pgfqpoint{1.038029in}{0.872331in}}%
\pgfpathlineto{\pgfqpoint{1.051738in}{0.866208in}}%
\pgfpathlineto{\pgfqpoint{1.064901in}{0.861083in}}%
\pgfpathlineto{\pgfqpoint{1.077585in}{0.856858in}}%
\pgfpathlineto{\pgfqpoint{1.089847in}{0.853446in}}%
\pgfpathlineto{\pgfqpoint{1.101736in}{0.850771in}}%
\pgfpathlineto{\pgfqpoint{1.113290in}{0.848764in}}%
\pgfpathlineto{\pgfqpoint{1.128233in}{0.847022in}}%
\pgfpathlineto{\pgfqpoint{1.142706in}{0.846229in}}%
\pgfpathlineto{\pgfqpoint{1.156759in}{0.846271in}}%
\pgfpathlineto{\pgfqpoint{1.170429in}{0.847044in}}%
\pgfpathlineto{\pgfqpoint{1.183745in}{0.848456in}}%
\pgfpathlineto{\pgfqpoint{1.199918in}{0.850992in}}%
\pgfpathlineto{\pgfqpoint{1.215588in}{0.854249in}}%
\pgfpathlineto{\pgfqpoint{1.230765in}{0.858093in}}%
\pgfpathlineto{\pgfqpoint{1.248325in}{0.863316in}}%
\pgfpathlineto{\pgfqpoint{1.267896in}{0.870019in}}%
\pgfpathlineto{\pgfqpoint{1.286450in}{0.877142in}}%
\pgfpathlineto{\pgfqpoint{1.306365in}{0.885522in}}%
\pgfpathlineto{\pgfqpoint{1.329259in}{0.895981in}}%
\pgfpathlineto{\pgfqpoint{1.353738in}{0.908011in}}%
\pgfpathlineto{\pgfqpoint{1.379794in}{0.921636in}}%
\pgfpathlineto{\pgfqpoint{1.409527in}{0.938060in}}%
\pgfpathlineto{\pgfqpoint{1.441853in}{0.956797in}}%
\pgfpathlineto{\pgfqpoint{1.473472in}{0.975862in}}%
\pgfpathlineto{\pgfqpoint{1.473472in}{0.975862in}}%
\pgfusepath{stroke}%
\end{pgfscope}%
\begin{pgfscope}%
\pgfpathrectangle{\pgfqpoint{0.407592in}{0.222904in}}{\pgfqpoint{1.411163in}{1.411163in}}%
\pgfusepath{clip}%
\pgfsetrectcap%
\pgfsetroundjoin%
\pgfsetlinewidth{2.308625pt}%
\definecolor{currentstroke}{rgb}{0.478431,0.243137,0.615686}%
\pgfsetstrokecolor{currentstroke}%
\pgfsetdash{}{0pt}%
\pgfpathmoveto{\pgfqpoint{0.835551in}{1.147966in}}%
\pgfpathlineto{\pgfqpoint{0.870813in}{1.108269in}}%
\pgfpathlineto{\pgfqpoint{0.901885in}{1.074249in}}%
\pgfpathlineto{\pgfqpoint{0.929538in}{1.044913in}}%
\pgfpathlineto{\pgfqpoint{0.954375in}{1.019491in}}%
\pgfpathlineto{\pgfqpoint{0.976868in}{0.997376in}}%
\pgfpathlineto{\pgfqpoint{0.997395in}{0.978083in}}%
\pgfpathlineto{\pgfqpoint{1.016264in}{0.961221in}}%
\pgfpathlineto{\pgfqpoint{1.033725in}{0.946468in}}%
\pgfpathlineto{\pgfqpoint{1.049983in}{0.933561in}}%
\pgfpathlineto{\pgfqpoint{1.065212in}{0.922279in}}%
\pgfpathlineto{\pgfqpoint{1.079555in}{0.912438in}}%
\pgfpathlineto{\pgfqpoint{1.093134in}{0.903879in}}%
\pgfpathlineto{\pgfqpoint{1.106052in}{0.896470in}}%
\pgfpathlineto{\pgfqpoint{1.118395in}{0.890094in}}%
\pgfpathlineto{\pgfqpoint{1.130240in}{0.884652in}}%
\pgfpathlineto{\pgfqpoint{1.141649in}{0.880057in}}%
\pgfpathlineto{\pgfqpoint{1.152678in}{0.876233in}}%
\pgfpathlineto{\pgfqpoint{1.166872in}{0.872215in}}%
\pgfpathlineto{\pgfqpoint{1.180566in}{0.869306in}}%
\pgfpathlineto{\pgfqpoint{1.193835in}{0.867384in}}%
\pgfpathlineto{\pgfqpoint{1.206739in}{0.866342in}}%
\pgfpathlineto{\pgfqpoint{1.219327in}{0.866085in}}%
\pgfpathlineto{\pgfqpoint{1.231635in}{0.866529in}}%
\pgfpathlineto{\pgfqpoint{1.243692in}{0.867597in}}%
\pgfpathlineto{\pgfqpoint{1.258439in}{0.869702in}}%
\pgfpathlineto{\pgfqpoint{1.272850in}{0.872548in}}%
\pgfpathlineto{\pgfqpoint{1.286937in}{0.876022in}}%
\pgfpathlineto{\pgfqpoint{1.303416in}{0.880874in}}%
\pgfpathlineto{\pgfqpoint{1.319417in}{0.886317in}}%
\pgfpathlineto{\pgfqpoint{1.337451in}{0.893219in}}%
\pgfpathlineto{\pgfqpoint{1.357165in}{0.901581in}}%
\pgfpathlineto{\pgfqpoint{1.378102in}{0.911272in}}%
\pgfpathlineto{\pgfqpoint{1.401762in}{0.923080in}}%
\pgfpathlineto{\pgfqpoint{1.426813in}{0.936417in}}%
\pgfpathlineto{\pgfqpoint{1.454691in}{0.952099in}}%
\pgfpathlineto{\pgfqpoint{1.486132in}{0.970656in}}%
\pgfpathlineto{\pgfqpoint{1.521413in}{0.992369in}}%
\pgfpathlineto{\pgfqpoint{1.547118in}{1.008671in}}%
\pgfpathlineto{\pgfqpoint{1.547118in}{1.008671in}}%
\pgfusepath{stroke}%
\end{pgfscope}%
\begin{pgfscope}%
\pgfpathrectangle{\pgfqpoint{0.407592in}{0.222904in}}{\pgfqpoint{1.411163in}{1.411163in}}%
\pgfusepath{clip}%
\pgfsetrectcap%
\pgfsetroundjoin%
\pgfsetlinewidth{2.308625pt}%
\definecolor{currentstroke}{rgb}{0.878431,0.552941,0.176471}%
\pgfsetstrokecolor{currentstroke}%
\pgfsetdash{}{0pt}%
\pgfpathmoveto{\pgfqpoint{0.867618in}{1.272438in}}%
\pgfpathlineto{\pgfqpoint{0.911021in}{1.225022in}}%
\pgfpathlineto{\pgfqpoint{0.948898in}{1.185089in}}%
\pgfpathlineto{\pgfqpoint{0.982354in}{1.151214in}}%
\pgfpathlineto{\pgfqpoint{1.012228in}{1.122317in}}%
\pgfpathlineto{\pgfqpoint{1.039168in}{1.097564in}}%
\pgfpathlineto{\pgfqpoint{1.063680in}{1.076302in}}%
\pgfpathlineto{\pgfqpoint{1.086165in}{1.058008in}}%
\pgfpathlineto{\pgfqpoint{1.106947in}{1.042264in}}%
\pgfpathlineto{\pgfqpoint{1.126288in}{1.028726in}}%
\pgfpathlineto{\pgfqpoint{1.144400in}{1.017112in}}%
\pgfpathlineto{\pgfqpoint{1.161461in}{1.007185in}}%
\pgfpathlineto{\pgfqpoint{1.177617in}{0.998746in}}%
\pgfpathlineto{\pgfqpoint{1.192989in}{0.991626in}}%
\pgfpathlineto{\pgfqpoint{1.207677in}{0.985681in}}%
\pgfpathlineto{\pgfqpoint{1.221768in}{0.980784in}}%
\pgfpathlineto{\pgfqpoint{1.235331in}{0.976828in}}%
\pgfpathlineto{\pgfqpoint{1.248426in}{0.973714in}}%
\pgfpathlineto{\pgfqpoint{1.261103in}{0.971359in}}%
\pgfpathlineto{\pgfqpoint{1.273402in}{0.969687in}}%
\pgfpathlineto{\pgfqpoint{1.289272in}{0.968402in}}%
\pgfpathlineto{\pgfqpoint{1.304595in}{0.968066in}}%
\pgfpathlineto{\pgfqpoint{1.319420in}{0.968553in}}%
\pgfpathlineto{\pgfqpoint{1.333783in}{0.969750in}}%
\pgfpathlineto{\pgfqpoint{1.347708in}{0.971555in}}%
\pgfpathlineto{\pgfqpoint{1.364527in}{0.974530in}}%
\pgfpathlineto{\pgfqpoint{1.380708in}{0.978156in}}%
\pgfpathlineto{\pgfqpoint{1.399297in}{0.983169in}}%
\pgfpathlineto{\pgfqpoint{1.416982in}{0.988707in}}%
\pgfpathlineto{\pgfqpoint{1.436461in}{0.995588in}}%
\pgfpathlineto{\pgfqpoint{1.457181in}{1.003713in}}%
\pgfpathlineto{\pgfqpoint{1.478506in}{1.012845in}}%
\pgfpathlineto{\pgfqpoint{1.501727in}{1.023571in}}%
\pgfpathlineto{\pgfqpoint{1.526894in}{1.035997in}}%
\pgfpathlineto{\pgfqpoint{1.555052in}{1.050749in}}%
\pgfpathlineto{\pgfqpoint{1.586043in}{1.067868in}}%
\pgfpathlineto{\pgfqpoint{1.619859in}{1.087432in}}%
\pgfpathlineto{\pgfqpoint{1.619876in}{1.087442in}}%
\pgfpathlineto{\pgfqpoint{1.619876in}{1.087442in}}%
\pgfusepath{stroke}%
\end{pgfscope}%
\begin{pgfscope}%
\pgfpathrectangle{\pgfqpoint{0.407592in}{0.222904in}}{\pgfqpoint{1.411163in}{1.411163in}}%
\pgfusepath{clip}%
\pgfsetbuttcap%
\pgfsetroundjoin%
\definecolor{currentfill}{rgb}{0.878431,0.552941,0.176471}%
\pgfsetfillcolor{currentfill}%
\pgfsetlinewidth{1.003750pt}%
\definecolor{currentstroke}{rgb}{0.878431,0.552941,0.176471}%
\pgfsetstrokecolor{currentstroke}%
\pgfsetdash{}{0pt}%
\pgfsys@defobject{currentmarker}{\pgfqpoint{-0.034021in}{-0.034021in}}{\pgfqpoint{0.034021in}{0.034021in}}{%
\pgfpathmoveto{\pgfqpoint{0.000000in}{-0.034021in}}%
\pgfpathcurveto{\pgfqpoint{0.009022in}{-0.034021in}}{\pgfqpoint{0.017676in}{-0.030436in}}{\pgfqpoint{0.024056in}{-0.024056in}}%
\pgfpathcurveto{\pgfqpoint{0.030436in}{-0.017676in}}{\pgfqpoint{0.034021in}{-0.009022in}}{\pgfqpoint{0.034021in}{0.000000in}}%
\pgfpathcurveto{\pgfqpoint{0.034021in}{0.009022in}}{\pgfqpoint{0.030436in}{0.017676in}}{\pgfqpoint{0.024056in}{0.024056in}}%
\pgfpathcurveto{\pgfqpoint{0.017676in}{0.030436in}}{\pgfqpoint{0.009022in}{0.034021in}}{\pgfqpoint{0.000000in}{0.034021in}}%
\pgfpathcurveto{\pgfqpoint{-0.009022in}{0.034021in}}{\pgfqpoint{-0.017676in}{0.030436in}}{\pgfqpoint{-0.024056in}{0.024056in}}%
\pgfpathcurveto{\pgfqpoint{-0.030436in}{0.017676in}}{\pgfqpoint{-0.034021in}{0.009022in}}{\pgfqpoint{-0.034021in}{0.000000in}}%
\pgfpathcurveto{\pgfqpoint{-0.034021in}{-0.009022in}}{\pgfqpoint{-0.030436in}{-0.017676in}}{\pgfqpoint{-0.024056in}{-0.024056in}}%
\pgfpathcurveto{\pgfqpoint{-0.017676in}{-0.030436in}}{\pgfqpoint{-0.009022in}{-0.034021in}}{\pgfqpoint{0.000000in}{-0.034021in}}%
\pgfpathlineto{\pgfqpoint{0.000000in}{-0.034021in}}%
\pgfpathclose%
\pgfusepath{stroke,fill}%
}%
\begin{pgfscope}%
\pgfsys@transformshift{0.867618in}{1.272438in}%
\pgfsys@useobject{currentmarker}{}%
\end{pgfscope}%
\end{pgfscope}%
\begin{pgfscope}%
\definecolor{textcolor}{rgb}{0.349020,0.376471,0.403922}%
\pgfsetstrokecolor{textcolor}%
\pgfsetfillcolor{textcolor}%
\pgftext[x=0.450955in,y=1.358414in,left,base]{\color{textcolor}{\rmfamily\fontsize{6.000000}{7.200000}\selectfont\catcode`\^=\active\def^{\ifmmode\sp\else\^{}\fi}\catcode`\%=\active\def%{\%}cone $a^2\geq r^2+b^2$}}%
\end{pgfscope}%
\begin{pgfscope}%
\definecolor{textcolor}{rgb}{0.125490,0.498039,0.321569}%
\pgfsetstrokecolor{textcolor}%
\pgfsetfillcolor{textcolor}%
\pgftext[x=1.299191in,y=1.332664in,left,base]{\color{textcolor}{\rmfamily\fontsize{6.000000}{7.200000}\selectfont\catcode`\^=\active\def^{\ifmmode\sp\else\^{}\fi}\catcode`\%=\active\def%{\%}minima $r=\beta$}}%
\end{pgfscope}%
\begin{pgfscope}%
\definecolor{textcolor}{rgb}{0.000000,0.000000,0.000000}%
\pgfsetstrokecolor{textcolor}%
\pgfsetfillcolor{textcolor}%
\pgftext[x=1.113173in,y=1.660759in,,base]{\color{textcolor}{\rmfamily\fontsize{8.000000}{9.600000}\selectfont\catcode`\^=\active\def^{\ifmmode\sp\else\^{}\fi}\catcode`\%=\active\def%{\%}(b) Unequal variances}}%
\end{pgfscope}%
\begin{pgfscope}%
\pgfpathrectangle{\pgfqpoint{0.407592in}{0.222904in}}{\pgfqpoint{1.411163in}{1.411163in}}%
\pgfusepath{clip}%
\pgfsetbuttcap%
\pgfsetroundjoin%
\definecolor{currentfill}{rgb}{0.478431,0.243137,0.615686}%
\pgfsetfillcolor{currentfill}%
\pgfsetlinewidth{1.003750pt}%
\definecolor{currentstroke}{rgb}{0.478431,0.243137,0.615686}%
\pgfsetstrokecolor{currentstroke}%
\pgfsetdash{}{0pt}%
\pgfsys@defobject{currentmarker}{\pgfqpoint{-0.034021in}{-0.034021in}}{\pgfqpoint{0.034021in}{0.034021in}}{%
\pgfpathmoveto{\pgfqpoint{0.000000in}{-0.034021in}}%
\pgfpathcurveto{\pgfqpoint{0.009022in}{-0.034021in}}{\pgfqpoint{0.017676in}{-0.030436in}}{\pgfqpoint{0.024056in}{-0.024056in}}%
\pgfpathcurveto{\pgfqpoint{0.030436in}{-0.017676in}}{\pgfqpoint{0.034021in}{-0.009022in}}{\pgfqpoint{0.034021in}{0.000000in}}%
\pgfpathcurveto{\pgfqpoint{0.034021in}{0.009022in}}{\pgfqpoint{0.030436in}{0.017676in}}{\pgfqpoint{0.024056in}{0.024056in}}%
\pgfpathcurveto{\pgfqpoint{0.017676in}{0.030436in}}{\pgfqpoint{0.009022in}{0.034021in}}{\pgfqpoint{0.000000in}{0.034021in}}%
\pgfpathcurveto{\pgfqpoint{-0.009022in}{0.034021in}}{\pgfqpoint{-0.017676in}{0.030436in}}{\pgfqpoint{-0.024056in}{0.024056in}}%
\pgfpathcurveto{\pgfqpoint{-0.030436in}{0.017676in}}{\pgfqpoint{-0.034021in}{0.009022in}}{\pgfqpoint{-0.034021in}{0.000000in}}%
\pgfpathcurveto{\pgfqpoint{-0.034021in}{-0.009022in}}{\pgfqpoint{-0.030436in}{-0.017676in}}{\pgfqpoint{-0.024056in}{-0.024056in}}%
\pgfpathcurveto{\pgfqpoint{-0.017676in}{-0.030436in}}{\pgfqpoint{-0.009022in}{-0.034021in}}{\pgfqpoint{0.000000in}{-0.034021in}}%
\pgfpathlineto{\pgfqpoint{0.000000in}{-0.034021in}}%
\pgfpathclose%
\pgfusepath{stroke,fill}%
}%
\begin{pgfscope}%
\pgfsys@transformshift{0.835551in}{1.147966in}%
\pgfsys@useobject{currentmarker}{}%
\end{pgfscope}%
\end{pgfscope}%
\begin{pgfscope}%
\pgfpathrectangle{\pgfqpoint{0.407592in}{0.222904in}}{\pgfqpoint{1.411163in}{1.411163in}}%
\pgfusepath{clip}%
\pgfsetbuttcap%
\pgfsetroundjoin%
\definecolor{currentfill}{rgb}{0.000000,0.529412,0.423529}%
\pgfsetfillcolor{currentfill}%
\pgfsetlinewidth{1.003750pt}%
\definecolor{currentstroke}{rgb}{0.000000,0.529412,0.423529}%
\pgfsetstrokecolor{currentstroke}%
\pgfsetdash{}{0pt}%
\pgfsys@defobject{currentmarker}{\pgfqpoint{-0.034021in}{-0.034021in}}{\pgfqpoint{0.034021in}{0.034021in}}{%
\pgfpathmoveto{\pgfqpoint{0.000000in}{-0.034021in}}%
\pgfpathcurveto{\pgfqpoint{0.009022in}{-0.034021in}}{\pgfqpoint{0.017676in}{-0.030436in}}{\pgfqpoint{0.024056in}{-0.024056in}}%
\pgfpathcurveto{\pgfqpoint{0.030436in}{-0.017676in}}{\pgfqpoint{0.034021in}{-0.009022in}}{\pgfqpoint{0.034021in}{0.000000in}}%
\pgfpathcurveto{\pgfqpoint{0.034021in}{0.009022in}}{\pgfqpoint{0.030436in}{0.017676in}}{\pgfqpoint{0.024056in}{0.024056in}}%
\pgfpathcurveto{\pgfqpoint{0.017676in}{0.030436in}}{\pgfqpoint{0.009022in}{0.034021in}}{\pgfqpoint{0.000000in}{0.034021in}}%
\pgfpathcurveto{\pgfqpoint{-0.009022in}{0.034021in}}{\pgfqpoint{-0.017676in}{0.030436in}}{\pgfqpoint{-0.024056in}{0.024056in}}%
\pgfpathcurveto{\pgfqpoint{-0.030436in}{0.017676in}}{\pgfqpoint{-0.034021in}{0.009022in}}{\pgfqpoint{-0.034021in}{0.000000in}}%
\pgfpathcurveto{\pgfqpoint{-0.034021in}{-0.009022in}}{\pgfqpoint{-0.030436in}{-0.017676in}}{\pgfqpoint{-0.024056in}{-0.024056in}}%
\pgfpathcurveto{\pgfqpoint{-0.017676in}{-0.030436in}}{\pgfqpoint{-0.009022in}{-0.034021in}}{\pgfqpoint{0.000000in}{-0.034021in}}%
\pgfpathlineto{\pgfqpoint{0.000000in}{-0.034021in}}%
\pgfpathclose%
\pgfusepath{stroke,fill}%
}%
\begin{pgfscope}%
\pgfsys@transformshift{0.812524in}{1.057517in}%
\pgfsys@useobject{currentmarker}{}%
\end{pgfscope}%
\end{pgfscope}%
\begin{pgfscope}%
\pgfpathrectangle{\pgfqpoint{0.407592in}{0.222904in}}{\pgfqpoint{1.411163in}{1.411163in}}%
\pgfusepath{clip}%
\pgfsetbuttcap%
\pgfsetroundjoin%
\definecolor{currentfill}{rgb}{0.819608,0.286275,0.356863}%
\pgfsetfillcolor{currentfill}%
\pgfsetlinewidth{1.003750pt}%
\definecolor{currentstroke}{rgb}{0.819608,0.286275,0.356863}%
\pgfsetstrokecolor{currentstroke}%
\pgfsetdash{}{0pt}%
\pgfsys@defobject{currentmarker}{\pgfqpoint{-0.034021in}{-0.034021in}}{\pgfqpoint{0.034021in}{0.034021in}}{%
\pgfpathmoveto{\pgfqpoint{0.000000in}{-0.034021in}}%
\pgfpathcurveto{\pgfqpoint{0.009022in}{-0.034021in}}{\pgfqpoint{0.017676in}{-0.030436in}}{\pgfqpoint{0.024056in}{-0.024056in}}%
\pgfpathcurveto{\pgfqpoint{0.030436in}{-0.017676in}}{\pgfqpoint{0.034021in}{-0.009022in}}{\pgfqpoint{0.034021in}{0.000000in}}%
\pgfpathcurveto{\pgfqpoint{0.034021in}{0.009022in}}{\pgfqpoint{0.030436in}{0.017676in}}{\pgfqpoint{0.024056in}{0.024056in}}%
\pgfpathcurveto{\pgfqpoint{0.017676in}{0.030436in}}{\pgfqpoint{0.009022in}{0.034021in}}{\pgfqpoint{0.000000in}{0.034021in}}%
\pgfpathcurveto{\pgfqpoint{-0.009022in}{0.034021in}}{\pgfqpoint{-0.017676in}{0.030436in}}{\pgfqpoint{-0.024056in}{0.024056in}}%
\pgfpathcurveto{\pgfqpoint{-0.030436in}{0.017676in}}{\pgfqpoint{-0.034021in}{0.009022in}}{\pgfqpoint{-0.034021in}{0.000000in}}%
\pgfpathcurveto{\pgfqpoint{-0.034021in}{-0.009022in}}{\pgfqpoint{-0.030436in}{-0.017676in}}{\pgfqpoint{-0.024056in}{-0.024056in}}%
\pgfpathcurveto{\pgfqpoint{-0.017676in}{-0.030436in}}{\pgfqpoint{-0.009022in}{-0.034021in}}{\pgfqpoint{0.000000in}{-0.034021in}}%
\pgfpathlineto{\pgfqpoint{0.000000in}{-0.034021in}}%
\pgfpathclose%
\pgfusepath{stroke,fill}%
}%
\begin{pgfscope}%
\pgfsys@transformshift{0.672155in}{1.057312in}%
\pgfsys@useobject{currentmarker}{}%
\end{pgfscope}%
\end{pgfscope}%
\begin{pgfscope}%
\pgfpathrectangle{\pgfqpoint{0.407592in}{0.222904in}}{\pgfqpoint{1.411163in}{1.411163in}}%
\pgfusepath{clip}%
\pgfsetbuttcap%
\pgfsetroundjoin%
\definecolor{currentfill}{rgb}{0.878431,0.552941,0.176471}%
\pgfsetfillcolor{currentfill}%
\pgfsetlinewidth{1.003750pt}%
\definecolor{currentstroke}{rgb}{0.878431,0.552941,0.176471}%
\pgfsetstrokecolor{currentstroke}%
\pgfsetdash{}{0pt}%
\pgfsys@defobject{currentmarker}{\pgfqpoint{-0.038036in}{-0.038036in}}{\pgfqpoint{0.038036in}{0.038036in}}{%
\pgfpathmoveto{\pgfqpoint{-0.038036in}{-0.038036in}}%
\pgfpathlineto{\pgfqpoint{0.038036in}{-0.038036in}}%
\pgfpathlineto{\pgfqpoint{0.038036in}{0.038036in}}%
\pgfpathlineto{\pgfqpoint{-0.038036in}{0.038036in}}%
\pgfpathlineto{\pgfqpoint{-0.038036in}{-0.038036in}}%
\pgfpathclose%
\pgfusepath{stroke,fill}%
}%
\begin{pgfscope}%
\pgfsys@transformshift{1.619876in}{1.087442in}%
\pgfsys@useobject{currentmarker}{}%
\end{pgfscope}%
\end{pgfscope}%
\begin{pgfscope}%
\pgfpathrectangle{\pgfqpoint{0.407592in}{0.222904in}}{\pgfqpoint{1.411163in}{1.411163in}}%
\pgfusepath{clip}%
\pgfsetbuttcap%
\pgfsetroundjoin%
\definecolor{currentfill}{rgb}{0.141176,0.345098,0.650980}%
\pgfsetfillcolor{currentfill}%
\pgfsetlinewidth{1.003750pt}%
\definecolor{currentstroke}{rgb}{0.141176,0.345098,0.650980}%
\pgfsetstrokecolor{currentstroke}%
\pgfsetdash{}{0pt}%
\pgfsys@defobject{currentmarker}{\pgfqpoint{-0.034021in}{-0.034021in}}{\pgfqpoint{0.034021in}{0.034021in}}{%
\pgfpathmoveto{\pgfqpoint{0.000000in}{-0.034021in}}%
\pgfpathcurveto{\pgfqpoint{0.009022in}{-0.034021in}}{\pgfqpoint{0.017676in}{-0.030436in}}{\pgfqpoint{0.024056in}{-0.024056in}}%
\pgfpathcurveto{\pgfqpoint{0.030436in}{-0.017676in}}{\pgfqpoint{0.034021in}{-0.009022in}}{\pgfqpoint{0.034021in}{0.000000in}}%
\pgfpathcurveto{\pgfqpoint{0.034021in}{0.009022in}}{\pgfqpoint{0.030436in}{0.017676in}}{\pgfqpoint{0.024056in}{0.024056in}}%
\pgfpathcurveto{\pgfqpoint{0.017676in}{0.030436in}}{\pgfqpoint{0.009022in}{0.034021in}}{\pgfqpoint{0.000000in}{0.034021in}}%
\pgfpathcurveto{\pgfqpoint{-0.009022in}{0.034021in}}{\pgfqpoint{-0.017676in}{0.030436in}}{\pgfqpoint{-0.024056in}{0.024056in}}%
\pgfpathcurveto{\pgfqpoint{-0.030436in}{0.017676in}}{\pgfqpoint{-0.034021in}{0.009022in}}{\pgfqpoint{-0.034021in}{0.000000in}}%
\pgfpathcurveto{\pgfqpoint{-0.034021in}{-0.009022in}}{\pgfqpoint{-0.030436in}{-0.017676in}}{\pgfqpoint{-0.024056in}{-0.024056in}}%
\pgfpathcurveto{\pgfqpoint{-0.017676in}{-0.030436in}}{\pgfqpoint{-0.009022in}{-0.034021in}}{\pgfqpoint{0.000000in}{-0.034021in}}%
\pgfpathlineto{\pgfqpoint{0.000000in}{-0.034021in}}%
\pgfpathclose%
\pgfusepath{stroke,fill}%
}%
\begin{pgfscope}%
\pgfsys@transformshift{0.537673in}{1.095850in}%
\pgfsys@useobject{currentmarker}{}%
\end{pgfscope}%
\end{pgfscope}%
\begin{pgfscope}%
\pgfpathrectangle{\pgfqpoint{0.407592in}{0.222904in}}{\pgfqpoint{1.411163in}{1.411163in}}%
\pgfusepath{clip}%
\pgfsetbuttcap%
\pgfsetroundjoin%
\definecolor{currentfill}{rgb}{0.478431,0.243137,0.615686}%
\pgfsetfillcolor{currentfill}%
\pgfsetlinewidth{1.003750pt}%
\definecolor{currentstroke}{rgb}{0.478431,0.243137,0.615686}%
\pgfsetstrokecolor{currentstroke}%
\pgfsetdash{}{0pt}%
\pgfsys@defobject{currentmarker}{\pgfqpoint{-0.038036in}{-0.038036in}}{\pgfqpoint{0.038036in}{0.038036in}}{%
\pgfpathmoveto{\pgfqpoint{-0.038036in}{-0.038036in}}%
\pgfpathlineto{\pgfqpoint{0.038036in}{-0.038036in}}%
\pgfpathlineto{\pgfqpoint{0.038036in}{0.038036in}}%
\pgfpathlineto{\pgfqpoint{-0.038036in}{0.038036in}}%
\pgfpathlineto{\pgfqpoint{-0.038036in}{-0.038036in}}%
\pgfpathclose%
\pgfusepath{stroke,fill}%
}%
\begin{pgfscope}%
\pgfsys@transformshift{1.547118in}{1.008671in}%
\pgfsys@useobject{currentmarker}{}%
\end{pgfscope}%
\end{pgfscope}%
\begin{pgfscope}%
\pgfpathrectangle{\pgfqpoint{0.407592in}{0.222904in}}{\pgfqpoint{1.411163in}{1.411163in}}%
\pgfusepath{clip}%
\pgfsetbuttcap%
\pgfsetroundjoin%
\definecolor{currentfill}{rgb}{0.000000,0.529412,0.423529}%
\pgfsetfillcolor{currentfill}%
\pgfsetlinewidth{1.003750pt}%
\definecolor{currentstroke}{rgb}{0.000000,0.529412,0.423529}%
\pgfsetstrokecolor{currentstroke}%
\pgfsetdash{}{0pt}%
\pgfsys@defobject{currentmarker}{\pgfqpoint{-0.038036in}{-0.038036in}}{\pgfqpoint{0.038036in}{0.038036in}}{%
\pgfpathmoveto{\pgfqpoint{-0.038036in}{-0.038036in}}%
\pgfpathlineto{\pgfqpoint{0.038036in}{-0.038036in}}%
\pgfpathlineto{\pgfqpoint{0.038036in}{0.038036in}}%
\pgfpathlineto{\pgfqpoint{-0.038036in}{0.038036in}}%
\pgfpathlineto{\pgfqpoint{-0.038036in}{-0.038036in}}%
\pgfpathclose%
\pgfusepath{stroke,fill}%
}%
\begin{pgfscope}%
\pgfsys@transformshift{1.473472in}{0.975862in}%
\pgfsys@useobject{currentmarker}{}%
\end{pgfscope}%
\end{pgfscope}%
\begin{pgfscope}%
\pgfpathrectangle{\pgfqpoint{0.407592in}{0.222904in}}{\pgfqpoint{1.411163in}{1.411163in}}%
\pgfusepath{clip}%
\pgfsetbuttcap%
\pgfsetroundjoin%
\definecolor{currentfill}{rgb}{0.819608,0.286275,0.356863}%
\pgfsetfillcolor{currentfill}%
\pgfsetlinewidth{1.003750pt}%
\definecolor{currentstroke}{rgb}{0.819608,0.286275,0.356863}%
\pgfsetstrokecolor{currentstroke}%
\pgfsetdash{}{0pt}%
\pgfsys@defobject{currentmarker}{\pgfqpoint{-0.038036in}{-0.038036in}}{\pgfqpoint{0.038036in}{0.038036in}}{%
\pgfpathmoveto{\pgfqpoint{-0.038036in}{-0.038036in}}%
\pgfpathlineto{\pgfqpoint{0.038036in}{-0.038036in}}%
\pgfpathlineto{\pgfqpoint{0.038036in}{0.038036in}}%
\pgfpathlineto{\pgfqpoint{-0.038036in}{0.038036in}}%
\pgfpathlineto{\pgfqpoint{-0.038036in}{-0.038036in}}%
\pgfpathclose%
\pgfusepath{stroke,fill}%
}%
\begin{pgfscope}%
\pgfsys@transformshift{1.396682in}{0.908341in}%
\pgfsys@useobject{currentmarker}{}%
\end{pgfscope}%
\end{pgfscope}%
\begin{pgfscope}%
\pgfpathrectangle{\pgfqpoint{0.407592in}{0.222904in}}{\pgfqpoint{1.411163in}{1.411163in}}%
\pgfusepath{clip}%
\pgfsetbuttcap%
\pgfsetroundjoin%
\definecolor{currentfill}{rgb}{0.141176,0.345098,0.650980}%
\pgfsetfillcolor{currentfill}%
\pgfsetlinewidth{1.003750pt}%
\definecolor{currentstroke}{rgb}{0.141176,0.345098,0.650980}%
\pgfsetstrokecolor{currentstroke}%
\pgfsetdash{}{0pt}%
\pgfsys@defobject{currentmarker}{\pgfqpoint{-0.038036in}{-0.038036in}}{\pgfqpoint{0.038036in}{0.038036in}}{%
\pgfpathmoveto{\pgfqpoint{-0.038036in}{-0.038036in}}%
\pgfpathlineto{\pgfqpoint{0.038036in}{-0.038036in}}%
\pgfpathlineto{\pgfqpoint{0.038036in}{0.038036in}}%
\pgfpathlineto{\pgfqpoint{-0.038036in}{0.038036in}}%
\pgfpathlineto{\pgfqpoint{-0.038036in}{-0.038036in}}%
\pgfpathclose%
\pgfusepath{stroke,fill}%
}%
\begin{pgfscope}%
\pgfsys@transformshift{1.318296in}{0.888253in}%
\pgfsys@useobject{currentmarker}{}%
\end{pgfscope}%
\end{pgfscope}%
\begin{pgfscope}%
\pgfpathrectangle{\pgfqpoint{0.407592in}{0.222904in}}{\pgfqpoint{1.411163in}{1.411163in}}%
\pgfusepath{clip}%
\pgfsetbuttcap%
\pgfsetroundjoin%
\definecolor{currentfill}{rgb}{0.141176,0.345098,0.650980}%
\pgfsetfillcolor{currentfill}%
\pgfsetfillopacity{0.080000}%
\pgfsetlinewidth{0.000000pt}%
\definecolor{currentstroke}{rgb}{0.000000,0.000000,0.000000}%
\pgfsetstrokecolor{currentstroke}%
\pgfsetdash{}{0pt}%
\pgfpathmoveto{\pgfqpoint{0.554113in}{1.064739in}}%
\pgfusepath{fill}%
\end{pgfscope}%
\begin{pgfscope}%
\pgfpathrectangle{\pgfqpoint{0.407592in}{0.222904in}}{\pgfqpoint{1.411163in}{1.411163in}}%
\pgfusepath{clip}%
\pgfsetbuttcap%
\pgfsetroundjoin%
\definecolor{currentfill}{rgb}{0.141176,0.345098,0.650980}%
\pgfsetfillcolor{currentfill}%
\pgfsetfillopacity{0.080000}%
\pgfsetlinewidth{0.000000pt}%
\definecolor{currentstroke}{rgb}{0.000000,0.000000,0.000000}%
\pgfsetstrokecolor{currentstroke}%
\pgfsetdash{}{0pt}%
\pgfpathmoveto{\pgfqpoint{0.554113in}{1.064739in}}%
\pgfpathlineto{\pgfqpoint{0.552790in}{1.106160in}}%
\pgfusepath{fill}%
\end{pgfscope}%
\begin{pgfscope}%
\pgfpathrectangle{\pgfqpoint{0.407592in}{0.222904in}}{\pgfqpoint{1.411163in}{1.411163in}}%
\pgfusepath{clip}%
\pgfsetbuttcap%
\pgfsetroundjoin%
\definecolor{currentfill}{rgb}{0.141176,0.345098,0.650980}%
\pgfsetfillcolor{currentfill}%
\pgfsetfillopacity{0.080000}%
\pgfsetlinewidth{0.000000pt}%
\definecolor{currentstroke}{rgb}{0.000000,0.000000,0.000000}%
\pgfsetstrokecolor{currentstroke}%
\pgfsetdash{}{0pt}%
\pgfpathmoveto{\pgfqpoint{0.552790in}{1.106160in}}%
\pgfpathlineto{\pgfqpoint{0.551459in}{1.147771in}}%
\pgfpathlineto{\pgfqpoint{0.489239in}{1.162676in}}%
\pgfusepath{fill}%
\end{pgfscope}%
\begin{pgfscope}%
\pgfpathrectangle{\pgfqpoint{0.407592in}{0.222904in}}{\pgfqpoint{1.411163in}{1.411163in}}%
\pgfusepath{clip}%
\pgfsetbuttcap%
\pgfsetroundjoin%
\definecolor{currentfill}{rgb}{0.141176,0.345098,0.650980}%
\pgfsetfillcolor{currentfill}%
\pgfsetfillopacity{0.080000}%
\pgfsetlinewidth{0.000000pt}%
\definecolor{currentstroke}{rgb}{0.000000,0.000000,0.000000}%
\pgfsetstrokecolor{currentstroke}%
\pgfsetdash{}{0pt}%
\pgfpathmoveto{\pgfqpoint{0.617856in}{1.007946in}}%
\pgfpathlineto{\pgfqpoint{0.616675in}{1.049346in}}%
\pgfpathlineto{\pgfqpoint{0.554113in}{1.064739in}}%
\pgfusepath{fill}%
\end{pgfscope}%
\begin{pgfscope}%
\pgfpathrectangle{\pgfqpoint{0.407592in}{0.222904in}}{\pgfqpoint{1.411163in}{1.411163in}}%
\pgfusepath{clip}%
\pgfsetbuttcap%
\pgfsetroundjoin%
\definecolor{currentfill}{rgb}{0.141176,0.345098,0.650980}%
\pgfsetfillcolor{currentfill}%
\pgfsetfillopacity{0.080000}%
\pgfsetlinewidth{0.000000pt}%
\definecolor{currentstroke}{rgb}{0.000000,0.000000,0.000000}%
\pgfsetstrokecolor{currentstroke}%
\pgfsetdash{}{0pt}%
\pgfpathmoveto{\pgfqpoint{0.617856in}{1.007946in}}%
\pgfusepath{fill}%
\end{pgfscope}%
\begin{pgfscope}%
\pgfpathrectangle{\pgfqpoint{0.407592in}{0.222904in}}{\pgfqpoint{1.411163in}{1.411163in}}%
\pgfusepath{clip}%
\pgfsetbuttcap%
\pgfsetroundjoin%
\definecolor{currentfill}{rgb}{0.141176,0.345098,0.650980}%
\pgfsetfillcolor{currentfill}%
\pgfsetfillopacity{0.080000}%
\pgfsetlinewidth{0.000000pt}%
\definecolor{currentstroke}{rgb}{0.000000,0.000000,0.000000}%
\pgfsetstrokecolor{currentstroke}%
\pgfsetdash{}{0pt}%
\pgfpathmoveto{\pgfqpoint{0.489239in}{1.162676in}}%
\pgfpathlineto{\pgfqpoint{0.551459in}{1.147771in}}%
\pgfpathlineto{\pgfqpoint{0.550123in}{1.189574in}}%
\pgfpathlineto{\pgfqpoint{0.487767in}{1.204307in}}%
\pgfpathlineto{\pgfqpoint{0.489239in}{1.162676in}}%
\pgfpathclose%
\pgfusepath{fill}%
\end{pgfscope}%
\begin{pgfscope}%
\pgfpathrectangle{\pgfqpoint{0.407592in}{0.222904in}}{\pgfqpoint{1.411163in}{1.411163in}}%
\pgfusepath{clip}%
\pgfsetbuttcap%
\pgfsetroundjoin%
\definecolor{currentfill}{rgb}{0.141176,0.345098,0.650980}%
\pgfsetfillcolor{currentfill}%
\pgfsetfillopacity{0.080000}%
\pgfsetlinewidth{0.000000pt}%
\definecolor{currentstroke}{rgb}{0.000000,0.000000,0.000000}%
\pgfsetstrokecolor{currentstroke}%
\pgfsetdash{}{0pt}%
\pgfpathmoveto{\pgfqpoint{0.554113in}{1.064739in}}%
\pgfpathlineto{\pgfqpoint{0.616675in}{1.049346in}}%
\pgfpathlineto{\pgfqpoint{0.615489in}{1.090936in}}%
\pgfpathlineto{\pgfqpoint{0.552790in}{1.106160in}}%
\pgfpathlineto{\pgfqpoint{0.554113in}{1.064739in}}%
\pgfpathclose%
\pgfusepath{fill}%
\end{pgfscope}%
\begin{pgfscope}%
\pgfpathrectangle{\pgfqpoint{0.407592in}{0.222904in}}{\pgfqpoint{1.411163in}{1.411163in}}%
\pgfusepath{clip}%
\pgfsetbuttcap%
\pgfsetroundjoin%
\definecolor{currentfill}{rgb}{0.141176,0.345098,0.650980}%
\pgfsetfillcolor{currentfill}%
\pgfsetfillopacity{0.080000}%
\pgfsetlinewidth{0.000000pt}%
\definecolor{currentstroke}{rgb}{0.000000,0.000000,0.000000}%
\pgfsetstrokecolor{currentstroke}%
\pgfsetdash{}{0pt}%
\pgfpathmoveto{\pgfqpoint{0.681935in}{0.950853in}}%
\pgfpathlineto{\pgfqpoint{0.680899in}{0.992231in}}%
\pgfpathlineto{\pgfqpoint{0.617856in}{1.007946in}}%
\pgfusepath{fill}%
\end{pgfscope}%
\begin{pgfscope}%
\pgfpathrectangle{\pgfqpoint{0.407592in}{0.222904in}}{\pgfqpoint{1.411163in}{1.411163in}}%
\pgfusepath{clip}%
\pgfsetbuttcap%
\pgfsetroundjoin%
\definecolor{currentfill}{rgb}{0.141176,0.345098,0.650980}%
\pgfsetfillcolor{currentfill}%
\pgfsetfillopacity{0.080000}%
\pgfsetlinewidth{0.000000pt}%
\definecolor{currentstroke}{rgb}{0.000000,0.000000,0.000000}%
\pgfsetstrokecolor{currentstroke}%
\pgfsetdash{}{0pt}%
\pgfpathmoveto{\pgfqpoint{0.681935in}{0.950853in}}%
\pgfusepath{fill}%
\end{pgfscope}%
\begin{pgfscope}%
\pgfpathrectangle{\pgfqpoint{0.407592in}{0.222904in}}{\pgfqpoint{1.411163in}{1.411163in}}%
\pgfusepath{clip}%
\pgfsetbuttcap%
\pgfsetroundjoin%
\definecolor{currentfill}{rgb}{0.141176,0.345098,0.650980}%
\pgfsetfillcolor{currentfill}%
\pgfsetfillopacity{0.080000}%
\pgfsetlinewidth{0.000000pt}%
\definecolor{currentstroke}{rgb}{0.000000,0.000000,0.000000}%
\pgfsetstrokecolor{currentstroke}%
\pgfsetdash{}{0pt}%
\pgfpathmoveto{\pgfqpoint{0.487767in}{1.204307in}}%
\pgfpathlineto{\pgfqpoint{0.550123in}{1.189574in}}%
\pgfpathlineto{\pgfqpoint{0.548781in}{1.231569in}}%
\pgfpathlineto{\pgfqpoint{0.486288in}{1.246129in}}%
\pgfpathlineto{\pgfqpoint{0.487767in}{1.204307in}}%
\pgfpathclose%
\pgfusepath{fill}%
\end{pgfscope}%
\begin{pgfscope}%
\pgfpathrectangle{\pgfqpoint{0.407592in}{0.222904in}}{\pgfqpoint{1.411163in}{1.411163in}}%
\pgfusepath{clip}%
\pgfsetbuttcap%
\pgfsetroundjoin%
\definecolor{currentfill}{rgb}{0.141176,0.345098,0.650980}%
\pgfsetfillcolor{currentfill}%
\pgfsetfillopacity{0.080000}%
\pgfsetlinewidth{0.000000pt}%
\definecolor{currentstroke}{rgb}{0.000000,0.000000,0.000000}%
\pgfsetstrokecolor{currentstroke}%
\pgfsetdash{}{0pt}%
\pgfpathmoveto{\pgfqpoint{0.552790in}{1.106160in}}%
\pgfpathlineto{\pgfqpoint{0.615489in}{1.090936in}}%
\pgfpathlineto{\pgfqpoint{0.614297in}{1.132718in}}%
\pgfpathlineto{\pgfqpoint{0.551459in}{1.147771in}}%
\pgfpathlineto{\pgfqpoint{0.552790in}{1.106160in}}%
\pgfpathclose%
\pgfusepath{fill}%
\end{pgfscope}%
\begin{pgfscope}%
\pgfpathrectangle{\pgfqpoint{0.407592in}{0.222904in}}{\pgfqpoint{1.411163in}{1.411163in}}%
\pgfusepath{clip}%
\pgfsetbuttcap%
\pgfsetroundjoin%
\definecolor{currentfill}{rgb}{0.141176,0.345098,0.650980}%
\pgfsetfillcolor{currentfill}%
\pgfsetfillopacity{0.080000}%
\pgfsetlinewidth{0.000000pt}%
\definecolor{currentstroke}{rgb}{0.000000,0.000000,0.000000}%
\pgfsetstrokecolor{currentstroke}%
\pgfsetdash{}{0pt}%
\pgfpathmoveto{\pgfqpoint{0.617856in}{1.007946in}}%
\pgfpathlineto{\pgfqpoint{0.680899in}{0.992231in}}%
\pgfpathlineto{\pgfqpoint{0.679857in}{1.033800in}}%
\pgfpathlineto{\pgfqpoint{0.616675in}{1.049346in}}%
\pgfpathlineto{\pgfqpoint{0.617856in}{1.007946in}}%
\pgfpathclose%
\pgfusepath{fill}%
\end{pgfscope}%
\begin{pgfscope}%
\pgfpathrectangle{\pgfqpoint{0.407592in}{0.222904in}}{\pgfqpoint{1.411163in}{1.411163in}}%
\pgfusepath{clip}%
\pgfsetbuttcap%
\pgfsetroundjoin%
\definecolor{currentfill}{rgb}{0.141176,0.345098,0.650980}%
\pgfsetfillcolor{currentfill}%
\pgfsetfillopacity{0.080000}%
\pgfsetlinewidth{0.000000pt}%
\definecolor{currentstroke}{rgb}{0.000000,0.000000,0.000000}%
\pgfsetstrokecolor{currentstroke}%
\pgfsetdash{}{0pt}%
\pgfpathmoveto{\pgfqpoint{0.746353in}{0.893459in}}%
\pgfpathlineto{\pgfqpoint{0.745462in}{0.934814in}}%
\pgfpathlineto{\pgfqpoint{0.681935in}{0.950853in}}%
\pgfusepath{fill}%
\end{pgfscope}%
\begin{pgfscope}%
\pgfpathrectangle{\pgfqpoint{0.407592in}{0.222904in}}{\pgfqpoint{1.411163in}{1.411163in}}%
\pgfusepath{clip}%
\pgfsetbuttcap%
\pgfsetroundjoin%
\definecolor{currentfill}{rgb}{0.141176,0.345098,0.650980}%
\pgfsetfillcolor{currentfill}%
\pgfsetfillopacity{0.080000}%
\pgfsetlinewidth{0.000000pt}%
\definecolor{currentstroke}{rgb}{0.000000,0.000000,0.000000}%
\pgfsetstrokecolor{currentstroke}%
\pgfsetdash{}{0pt}%
\pgfpathmoveto{\pgfqpoint{0.486288in}{1.246129in}}%
\pgfpathlineto{\pgfqpoint{0.548781in}{1.231569in}}%
\pgfpathlineto{\pgfqpoint{0.547433in}{1.273759in}}%
\pgfpathlineto{\pgfqpoint{0.484803in}{1.288144in}}%
\pgfpathlineto{\pgfqpoint{0.486288in}{1.246129in}}%
\pgfpathclose%
\pgfusepath{fill}%
\end{pgfscope}%
\begin{pgfscope}%
\pgfpathrectangle{\pgfqpoint{0.407592in}{0.222904in}}{\pgfqpoint{1.411163in}{1.411163in}}%
\pgfusepath{clip}%
\pgfsetbuttcap%
\pgfsetroundjoin%
\definecolor{currentfill}{rgb}{0.141176,0.345098,0.650980}%
\pgfsetfillcolor{currentfill}%
\pgfsetfillopacity{0.080000}%
\pgfsetlinewidth{0.000000pt}%
\definecolor{currentstroke}{rgb}{0.000000,0.000000,0.000000}%
\pgfsetstrokecolor{currentstroke}%
\pgfsetdash{}{0pt}%
\pgfpathmoveto{\pgfqpoint{0.746353in}{0.893459in}}%
\pgfusepath{fill}%
\end{pgfscope}%
\begin{pgfscope}%
\pgfpathrectangle{\pgfqpoint{0.407592in}{0.222904in}}{\pgfqpoint{1.411163in}{1.411163in}}%
\pgfusepath{clip}%
\pgfsetbuttcap%
\pgfsetroundjoin%
\definecolor{currentfill}{rgb}{0.141176,0.345098,0.650980}%
\pgfsetfillcolor{currentfill}%
\pgfsetfillopacity{0.080000}%
\pgfsetlinewidth{0.000000pt}%
\definecolor{currentstroke}{rgb}{0.000000,0.000000,0.000000}%
\pgfsetstrokecolor{currentstroke}%
\pgfsetdash{}{0pt}%
\pgfpathmoveto{\pgfqpoint{0.551459in}{1.147771in}}%
\pgfpathlineto{\pgfqpoint{0.614297in}{1.132718in}}%
\pgfpathlineto{\pgfqpoint{0.613099in}{1.174694in}}%
\pgfpathlineto{\pgfqpoint{0.550123in}{1.189574in}}%
\pgfpathlineto{\pgfqpoint{0.551459in}{1.147771in}}%
\pgfpathclose%
\pgfusepath{fill}%
\end{pgfscope}%
\begin{pgfscope}%
\pgfpathrectangle{\pgfqpoint{0.407592in}{0.222904in}}{\pgfqpoint{1.411163in}{1.411163in}}%
\pgfusepath{clip}%
\pgfsetbuttcap%
\pgfsetroundjoin%
\definecolor{currentfill}{rgb}{0.141176,0.345098,0.650980}%
\pgfsetfillcolor{currentfill}%
\pgfsetfillopacity{0.080000}%
\pgfsetlinewidth{0.000000pt}%
\definecolor{currentstroke}{rgb}{0.000000,0.000000,0.000000}%
\pgfsetstrokecolor{currentstroke}%
\pgfsetdash{}{0pt}%
\pgfpathmoveto{\pgfqpoint{0.616675in}{1.049346in}}%
\pgfpathlineto{\pgfqpoint{0.679857in}{1.033800in}}%
\pgfpathlineto{\pgfqpoint{0.678811in}{1.075561in}}%
\pgfpathlineto{\pgfqpoint{0.615489in}{1.090936in}}%
\pgfpathlineto{\pgfqpoint{0.616675in}{1.049346in}}%
\pgfpathclose%
\pgfusepath{fill}%
\end{pgfscope}%
\begin{pgfscope}%
\pgfpathrectangle{\pgfqpoint{0.407592in}{0.222904in}}{\pgfqpoint{1.411163in}{1.411163in}}%
\pgfusepath{clip}%
\pgfsetbuttcap%
\pgfsetroundjoin%
\definecolor{currentfill}{rgb}{0.141176,0.345098,0.650980}%
\pgfsetfillcolor{currentfill}%
\pgfsetfillopacity{0.080000}%
\pgfsetlinewidth{0.000000pt}%
\definecolor{currentstroke}{rgb}{0.000000,0.000000,0.000000}%
\pgfsetstrokecolor{currentstroke}%
\pgfsetdash{}{0pt}%
\pgfpathmoveto{\pgfqpoint{0.681935in}{0.950853in}}%
\pgfpathlineto{\pgfqpoint{0.745462in}{0.934814in}}%
\pgfpathlineto{\pgfqpoint{0.744568in}{0.976360in}}%
\pgfpathlineto{\pgfqpoint{0.680899in}{0.992231in}}%
\pgfpathlineto{\pgfqpoint{0.681935in}{0.950853in}}%
\pgfpathclose%
\pgfusepath{fill}%
\end{pgfscope}%
\begin{pgfscope}%
\pgfpathrectangle{\pgfqpoint{0.407592in}{0.222904in}}{\pgfqpoint{1.411163in}{1.411163in}}%
\pgfusepath{clip}%
\pgfsetbuttcap%
\pgfsetroundjoin%
\definecolor{currentfill}{rgb}{0.141176,0.345098,0.650980}%
\pgfsetfillcolor{currentfill}%
\pgfsetfillopacity{0.080000}%
\pgfsetlinewidth{0.000000pt}%
\definecolor{currentstroke}{rgb}{0.000000,0.000000,0.000000}%
\pgfsetstrokecolor{currentstroke}%
\pgfsetdash{}{0pt}%
\pgfpathmoveto{\pgfqpoint{0.811113in}{0.835760in}}%
\pgfpathlineto{\pgfqpoint{0.810369in}{0.877091in}}%
\pgfpathlineto{\pgfqpoint{0.746353in}{0.893459in}}%
\pgfusepath{fill}%
\end{pgfscope}%
\begin{pgfscope}%
\pgfpathrectangle{\pgfqpoint{0.407592in}{0.222904in}}{\pgfqpoint{1.411163in}{1.411163in}}%
\pgfusepath{clip}%
\pgfsetbuttcap%
\pgfsetroundjoin%
\definecolor{currentfill}{rgb}{0.141176,0.345098,0.650980}%
\pgfsetfillcolor{currentfill}%
\pgfsetfillopacity{0.080000}%
\pgfsetlinewidth{0.000000pt}%
\definecolor{currentstroke}{rgb}{0.000000,0.000000,0.000000}%
\pgfsetstrokecolor{currentstroke}%
\pgfsetdash{}{0pt}%
\pgfpathmoveto{\pgfqpoint{0.550123in}{1.189574in}}%
\pgfpathlineto{\pgfqpoint{0.613099in}{1.174694in}}%
\pgfpathlineto{\pgfqpoint{0.611896in}{1.216864in}}%
\pgfpathlineto{\pgfqpoint{0.548781in}{1.231569in}}%
\pgfpathlineto{\pgfqpoint{0.550123in}{1.189574in}}%
\pgfpathclose%
\pgfusepath{fill}%
\end{pgfscope}%
\begin{pgfscope}%
\pgfpathrectangle{\pgfqpoint{0.407592in}{0.222904in}}{\pgfqpoint{1.411163in}{1.411163in}}%
\pgfusepath{clip}%
\pgfsetbuttcap%
\pgfsetroundjoin%
\definecolor{currentfill}{rgb}{0.141176,0.345098,0.650980}%
\pgfsetfillcolor{currentfill}%
\pgfsetfillopacity{0.080000}%
\pgfsetlinewidth{0.000000pt}%
\definecolor{currentstroke}{rgb}{0.000000,0.000000,0.000000}%
\pgfsetstrokecolor{currentstroke}%
\pgfsetdash{}{0pt}%
\pgfpathmoveto{\pgfqpoint{0.811113in}{0.835760in}}%
\pgfusepath{fill}%
\end{pgfscope}%
\begin{pgfscope}%
\pgfpathrectangle{\pgfqpoint{0.407592in}{0.222904in}}{\pgfqpoint{1.411163in}{1.411163in}}%
\pgfusepath{clip}%
\pgfsetbuttcap%
\pgfsetroundjoin%
\definecolor{currentfill}{rgb}{0.141176,0.345098,0.650980}%
\pgfsetfillcolor{currentfill}%
\pgfsetfillopacity{0.080000}%
\pgfsetlinewidth{0.000000pt}%
\definecolor{currentstroke}{rgb}{0.000000,0.000000,0.000000}%
\pgfsetstrokecolor{currentstroke}%
\pgfsetdash{}{0pt}%
\pgfpathmoveto{\pgfqpoint{0.615489in}{1.090936in}}%
\pgfpathlineto{\pgfqpoint{0.678811in}{1.075561in}}%
\pgfpathlineto{\pgfqpoint{0.677760in}{1.117516in}}%
\pgfpathlineto{\pgfqpoint{0.614297in}{1.132718in}}%
\pgfpathlineto{\pgfqpoint{0.615489in}{1.090936in}}%
\pgfpathclose%
\pgfusepath{fill}%
\end{pgfscope}%
\begin{pgfscope}%
\pgfpathrectangle{\pgfqpoint{0.407592in}{0.222904in}}{\pgfqpoint{1.411163in}{1.411163in}}%
\pgfusepath{clip}%
\pgfsetbuttcap%
\pgfsetroundjoin%
\definecolor{currentfill}{rgb}{0.141176,0.345098,0.650980}%
\pgfsetfillcolor{currentfill}%
\pgfsetfillopacity{0.080000}%
\pgfsetlinewidth{0.000000pt}%
\definecolor{currentstroke}{rgb}{0.000000,0.000000,0.000000}%
\pgfsetstrokecolor{currentstroke}%
\pgfsetdash{}{0pt}%
\pgfpathmoveto{\pgfqpoint{0.680899in}{0.992231in}}%
\pgfpathlineto{\pgfqpoint{0.744568in}{0.976360in}}%
\pgfpathlineto{\pgfqpoint{0.743669in}{1.018099in}}%
\pgfpathlineto{\pgfqpoint{0.679857in}{1.033800in}}%
\pgfpathlineto{\pgfqpoint{0.680899in}{0.992231in}}%
\pgfpathclose%
\pgfusepath{fill}%
\end{pgfscope}%
\begin{pgfscope}%
\pgfpathrectangle{\pgfqpoint{0.407592in}{0.222904in}}{\pgfqpoint{1.411163in}{1.411163in}}%
\pgfusepath{clip}%
\pgfsetbuttcap%
\pgfsetroundjoin%
\definecolor{currentfill}{rgb}{0.141176,0.345098,0.650980}%
\pgfsetfillcolor{currentfill}%
\pgfsetfillopacity{0.080000}%
\pgfsetlinewidth{0.000000pt}%
\definecolor{currentstroke}{rgb}{0.000000,0.000000,0.000000}%
\pgfsetstrokecolor{currentstroke}%
\pgfsetdash{}{0pt}%
\pgfpathmoveto{\pgfqpoint{0.746353in}{0.893459in}}%
\pgfpathlineto{\pgfqpoint{0.810369in}{0.877091in}}%
\pgfpathlineto{\pgfqpoint{0.809623in}{0.918615in}}%
\pgfpathlineto{\pgfqpoint{0.745462in}{0.934814in}}%
\pgfpathlineto{\pgfqpoint{0.746353in}{0.893459in}}%
\pgfpathclose%
\pgfusepath{fill}%
\end{pgfscope}%
\begin{pgfscope}%
\pgfpathrectangle{\pgfqpoint{0.407592in}{0.222904in}}{\pgfqpoint{1.411163in}{1.411163in}}%
\pgfusepath{clip}%
\pgfsetbuttcap%
\pgfsetroundjoin%
\definecolor{currentfill}{rgb}{0.141176,0.345098,0.650980}%
\pgfsetfillcolor{currentfill}%
\pgfsetfillopacity{0.080000}%
\pgfsetlinewidth{0.000000pt}%
\definecolor{currentstroke}{rgb}{0.000000,0.000000,0.000000}%
\pgfsetstrokecolor{currentstroke}%
\pgfsetdash{}{0pt}%
\pgfpathmoveto{\pgfqpoint{0.548781in}{1.231569in}}%
\pgfpathlineto{\pgfqpoint{0.611896in}{1.216864in}}%
\pgfpathlineto{\pgfqpoint{0.610687in}{1.259230in}}%
\pgfpathlineto{\pgfqpoint{0.547433in}{1.273759in}}%
\pgfpathlineto{\pgfqpoint{0.548781in}{1.231569in}}%
\pgfpathclose%
\pgfusepath{fill}%
\end{pgfscope}%
\begin{pgfscope}%
\pgfpathrectangle{\pgfqpoint{0.407592in}{0.222904in}}{\pgfqpoint{1.411163in}{1.411163in}}%
\pgfusepath{clip}%
\pgfsetbuttcap%
\pgfsetroundjoin%
\definecolor{currentfill}{rgb}{0.141176,0.345098,0.650980}%
\pgfsetfillcolor{currentfill}%
\pgfsetfillopacity{0.080000}%
\pgfsetlinewidth{0.000000pt}%
\definecolor{currentstroke}{rgb}{0.000000,0.000000,0.000000}%
\pgfsetstrokecolor{currentstroke}%
\pgfsetdash{}{0pt}%
\pgfpathmoveto{\pgfqpoint{0.876216in}{0.777754in}}%
\pgfpathlineto{\pgfqpoint{0.875622in}{0.819061in}}%
\pgfpathlineto{\pgfqpoint{0.811113in}{0.835760in}}%
\pgfusepath{fill}%
\end{pgfscope}%
\begin{pgfscope}%
\pgfpathrectangle{\pgfqpoint{0.407592in}{0.222904in}}{\pgfqpoint{1.411163in}{1.411163in}}%
\pgfusepath{clip}%
\pgfsetbuttcap%
\pgfsetroundjoin%
\definecolor{currentfill}{rgb}{0.141176,0.345098,0.650980}%
\pgfsetfillcolor{currentfill}%
\pgfsetfillopacity{0.080000}%
\pgfsetlinewidth{0.000000pt}%
\definecolor{currentstroke}{rgb}{0.000000,0.000000,0.000000}%
\pgfsetstrokecolor{currentstroke}%
\pgfsetdash{}{0pt}%
\pgfpathmoveto{\pgfqpoint{0.614297in}{1.132718in}}%
\pgfpathlineto{\pgfqpoint{0.677760in}{1.117516in}}%
\pgfpathlineto{\pgfqpoint{0.676704in}{1.159665in}}%
\pgfpathlineto{\pgfqpoint{0.613099in}{1.174694in}}%
\pgfpathlineto{\pgfqpoint{0.614297in}{1.132718in}}%
\pgfpathclose%
\pgfusepath{fill}%
\end{pgfscope}%
\begin{pgfscope}%
\pgfpathrectangle{\pgfqpoint{0.407592in}{0.222904in}}{\pgfqpoint{1.411163in}{1.411163in}}%
\pgfusepath{clip}%
\pgfsetbuttcap%
\pgfsetroundjoin%
\definecolor{currentfill}{rgb}{0.141176,0.345098,0.650980}%
\pgfsetfillcolor{currentfill}%
\pgfsetfillopacity{0.080000}%
\pgfsetlinewidth{0.000000pt}%
\definecolor{currentstroke}{rgb}{0.000000,0.000000,0.000000}%
\pgfsetstrokecolor{currentstroke}%
\pgfsetdash{}{0pt}%
\pgfpathmoveto{\pgfqpoint{0.876216in}{0.777754in}}%
\pgfusepath{fill}%
\end{pgfscope}%
\begin{pgfscope}%
\pgfpathrectangle{\pgfqpoint{0.407592in}{0.222904in}}{\pgfqpoint{1.411163in}{1.411163in}}%
\pgfusepath{clip}%
\pgfsetbuttcap%
\pgfsetroundjoin%
\definecolor{currentfill}{rgb}{0.141176,0.345098,0.650980}%
\pgfsetfillcolor{currentfill}%
\pgfsetfillopacity{0.080000}%
\pgfsetlinewidth{0.000000pt}%
\definecolor{currentstroke}{rgb}{0.000000,0.000000,0.000000}%
\pgfsetstrokecolor{currentstroke}%
\pgfsetdash{}{0pt}%
\pgfpathmoveto{\pgfqpoint{0.679857in}{1.033800in}}%
\pgfpathlineto{\pgfqpoint{0.743669in}{1.018099in}}%
\pgfpathlineto{\pgfqpoint{0.742765in}{1.060032in}}%
\pgfpathlineto{\pgfqpoint{0.678811in}{1.075561in}}%
\pgfpathlineto{\pgfqpoint{0.679857in}{1.033800in}}%
\pgfpathclose%
\pgfusepath{fill}%
\end{pgfscope}%
\begin{pgfscope}%
\pgfpathrectangle{\pgfqpoint{0.407592in}{0.222904in}}{\pgfqpoint{1.411163in}{1.411163in}}%
\pgfusepath{clip}%
\pgfsetbuttcap%
\pgfsetroundjoin%
\definecolor{currentfill}{rgb}{0.141176,0.345098,0.650980}%
\pgfsetfillcolor{currentfill}%
\pgfsetfillopacity{0.080000}%
\pgfsetlinewidth{0.000000pt}%
\definecolor{currentstroke}{rgb}{0.000000,0.000000,0.000000}%
\pgfsetstrokecolor{currentstroke}%
\pgfsetdash{}{0pt}%
\pgfpathmoveto{\pgfqpoint{0.745462in}{0.934814in}}%
\pgfpathlineto{\pgfqpoint{0.809623in}{0.918615in}}%
\pgfpathlineto{\pgfqpoint{0.808873in}{0.960331in}}%
\pgfpathlineto{\pgfqpoint{0.744568in}{0.976360in}}%
\pgfpathlineto{\pgfqpoint{0.745462in}{0.934814in}}%
\pgfpathclose%
\pgfusepath{fill}%
\end{pgfscope}%
\begin{pgfscope}%
\pgfpathrectangle{\pgfqpoint{0.407592in}{0.222904in}}{\pgfqpoint{1.411163in}{1.411163in}}%
\pgfusepath{clip}%
\pgfsetbuttcap%
\pgfsetroundjoin%
\definecolor{currentfill}{rgb}{0.141176,0.345098,0.650980}%
\pgfsetfillcolor{currentfill}%
\pgfsetfillopacity{0.080000}%
\pgfsetlinewidth{0.000000pt}%
\definecolor{currentstroke}{rgb}{0.000000,0.000000,0.000000}%
\pgfsetstrokecolor{currentstroke}%
\pgfsetdash{}{0pt}%
\pgfpathmoveto{\pgfqpoint{0.811113in}{0.835760in}}%
\pgfpathlineto{\pgfqpoint{0.875622in}{0.819061in}}%
\pgfpathlineto{\pgfqpoint{0.875025in}{0.860560in}}%
\pgfpathlineto{\pgfqpoint{0.810369in}{0.877091in}}%
\pgfpathlineto{\pgfqpoint{0.811113in}{0.835760in}}%
\pgfpathclose%
\pgfusepath{fill}%
\end{pgfscope}%
\begin{pgfscope}%
\pgfpathrectangle{\pgfqpoint{0.407592in}{0.222904in}}{\pgfqpoint{1.411163in}{1.411163in}}%
\pgfusepath{clip}%
\pgfsetbuttcap%
\pgfsetroundjoin%
\definecolor{currentfill}{rgb}{0.141176,0.345098,0.650980}%
\pgfsetfillcolor{currentfill}%
\pgfsetfillopacity{0.080000}%
\pgfsetlinewidth{0.000000pt}%
\definecolor{currentstroke}{rgb}{0.000000,0.000000,0.000000}%
\pgfsetstrokecolor{currentstroke}%
\pgfsetdash{}{0pt}%
\pgfpathmoveto{\pgfqpoint{0.613099in}{1.174694in}}%
\pgfpathlineto{\pgfqpoint{0.676704in}{1.159665in}}%
\pgfpathlineto{\pgfqpoint{0.675643in}{1.202012in}}%
\pgfpathlineto{\pgfqpoint{0.611896in}{1.216864in}}%
\pgfpathlineto{\pgfqpoint{0.613099in}{1.174694in}}%
\pgfpathclose%
\pgfusepath{fill}%
\end{pgfscope}%
\begin{pgfscope}%
\pgfpathrectangle{\pgfqpoint{0.407592in}{0.222904in}}{\pgfqpoint{1.411163in}{1.411163in}}%
\pgfusepath{clip}%
\pgfsetbuttcap%
\pgfsetroundjoin%
\definecolor{currentfill}{rgb}{0.141176,0.345098,0.650980}%
\pgfsetfillcolor{currentfill}%
\pgfsetfillopacity{0.080000}%
\pgfsetlinewidth{0.000000pt}%
\definecolor{currentstroke}{rgb}{0.000000,0.000000,0.000000}%
\pgfsetstrokecolor{currentstroke}%
\pgfsetdash{}{0pt}%
\pgfpathmoveto{\pgfqpoint{0.941224in}{0.760721in}}%
\pgfpathlineto{\pgfqpoint{0.876216in}{0.777754in}}%
\pgfusepath{fill}%
\end{pgfscope}%
\begin{pgfscope}%
\pgfpathrectangle{\pgfqpoint{0.407592in}{0.222904in}}{\pgfqpoint{1.411163in}{1.411163in}}%
\pgfusepath{clip}%
\pgfsetbuttcap%
\pgfsetroundjoin%
\definecolor{currentfill}{rgb}{0.141176,0.345098,0.650980}%
\pgfsetfillcolor{currentfill}%
\pgfsetfillopacity{0.080000}%
\pgfsetlinewidth{0.000000pt}%
\definecolor{currentstroke}{rgb}{0.000000,0.000000,0.000000}%
\pgfsetstrokecolor{currentstroke}%
\pgfsetdash{}{0pt}%
\pgfpathmoveto{\pgfqpoint{0.678811in}{1.075561in}}%
\pgfpathlineto{\pgfqpoint{0.742765in}{1.060032in}}%
\pgfpathlineto{\pgfqpoint{0.741858in}{1.102161in}}%
\pgfpathlineto{\pgfqpoint{0.677760in}{1.117516in}}%
\pgfpathlineto{\pgfqpoint{0.678811in}{1.075561in}}%
\pgfpathclose%
\pgfusepath{fill}%
\end{pgfscope}%
\begin{pgfscope}%
\pgfpathrectangle{\pgfqpoint{0.407592in}{0.222904in}}{\pgfqpoint{1.411163in}{1.411163in}}%
\pgfusepath{clip}%
\pgfsetbuttcap%
\pgfsetroundjoin%
\definecolor{currentfill}{rgb}{0.141176,0.345098,0.650980}%
\pgfsetfillcolor{currentfill}%
\pgfsetfillopacity{0.080000}%
\pgfsetlinewidth{0.000000pt}%
\definecolor{currentstroke}{rgb}{0.000000,0.000000,0.000000}%
\pgfsetstrokecolor{currentstroke}%
\pgfsetdash{}{0pt}%
\pgfpathmoveto{\pgfqpoint{0.744568in}{0.976360in}}%
\pgfpathlineto{\pgfqpoint{0.808873in}{0.960331in}}%
\pgfpathlineto{\pgfqpoint{0.808119in}{1.002242in}}%
\pgfpathlineto{\pgfqpoint{0.743669in}{1.018099in}}%
\pgfpathlineto{\pgfqpoint{0.744568in}{0.976360in}}%
\pgfpathclose%
\pgfusepath{fill}%
\end{pgfscope}%
\begin{pgfscope}%
\pgfpathrectangle{\pgfqpoint{0.407592in}{0.222904in}}{\pgfqpoint{1.411163in}{1.411163in}}%
\pgfusepath{clip}%
\pgfsetbuttcap%
\pgfsetroundjoin%
\definecolor{currentfill}{rgb}{0.141176,0.345098,0.650980}%
\pgfsetfillcolor{currentfill}%
\pgfsetfillopacity{0.080000}%
\pgfsetlinewidth{0.000000pt}%
\definecolor{currentstroke}{rgb}{0.000000,0.000000,0.000000}%
\pgfsetstrokecolor{currentstroke}%
\pgfsetdash{}{0pt}%
\pgfpathmoveto{\pgfqpoint{0.810369in}{0.877091in}}%
\pgfpathlineto{\pgfqpoint{0.875025in}{0.860560in}}%
\pgfpathlineto{\pgfqpoint{0.874426in}{0.902253in}}%
\pgfpathlineto{\pgfqpoint{0.809623in}{0.918615in}}%
\pgfpathlineto{\pgfqpoint{0.810369in}{0.877091in}}%
\pgfpathclose%
\pgfusepath{fill}%
\end{pgfscope}%
\begin{pgfscope}%
\pgfpathrectangle{\pgfqpoint{0.407592in}{0.222904in}}{\pgfqpoint{1.411163in}{1.411163in}}%
\pgfusepath{clip}%
\pgfsetbuttcap%
\pgfsetroundjoin%
\definecolor{currentfill}{rgb}{0.141176,0.345098,0.650980}%
\pgfsetfillcolor{currentfill}%
\pgfsetfillopacity{0.080000}%
\pgfsetlinewidth{0.000000pt}%
\definecolor{currentstroke}{rgb}{0.000000,0.000000,0.000000}%
\pgfsetstrokecolor{currentstroke}%
\pgfsetdash{}{0pt}%
\pgfpathmoveto{\pgfqpoint{0.876216in}{0.777754in}}%
\pgfpathlineto{\pgfqpoint{0.941224in}{0.760721in}}%
\pgfpathlineto{\pgfqpoint{0.940778in}{0.802195in}}%
\pgfpathlineto{\pgfqpoint{0.875622in}{0.819061in}}%
\pgfpathlineto{\pgfqpoint{0.876216in}{0.777754in}}%
\pgfpathclose%
\pgfusepath{fill}%
\end{pgfscope}%
\begin{pgfscope}%
\pgfpathrectangle{\pgfqpoint{0.407592in}{0.222904in}}{\pgfqpoint{1.411163in}{1.411163in}}%
\pgfusepath{clip}%
\pgfsetbuttcap%
\pgfsetroundjoin%
\definecolor{currentfill}{rgb}{0.141176,0.345098,0.650980}%
\pgfsetfillcolor{currentfill}%
\pgfsetfillopacity{0.080000}%
\pgfsetlinewidth{0.000000pt}%
\definecolor{currentstroke}{rgb}{0.000000,0.000000,0.000000}%
\pgfsetstrokecolor{currentstroke}%
\pgfsetdash{}{0pt}%
\pgfpathmoveto{\pgfqpoint{0.611896in}{1.216864in}}%
\pgfpathlineto{\pgfqpoint{0.675643in}{1.202012in}}%
\pgfpathlineto{\pgfqpoint{0.674577in}{1.244555in}}%
\pgfpathlineto{\pgfqpoint{0.610687in}{1.259230in}}%
\pgfpathlineto{\pgfqpoint{0.611896in}{1.216864in}}%
\pgfpathclose%
\pgfusepath{fill}%
\end{pgfscope}%
\begin{pgfscope}%
\pgfpathrectangle{\pgfqpoint{0.407592in}{0.222904in}}{\pgfqpoint{1.411163in}{1.411163in}}%
\pgfusepath{clip}%
\pgfsetbuttcap%
\pgfsetroundjoin%
\definecolor{currentfill}{rgb}{0.141176,0.345098,0.650980}%
\pgfsetfillcolor{currentfill}%
\pgfsetfillopacity{0.080000}%
\pgfsetlinewidth{0.000000pt}%
\definecolor{currentstroke}{rgb}{0.000000,0.000000,0.000000}%
\pgfsetstrokecolor{currentstroke}%
\pgfsetdash{}{0pt}%
\pgfpathmoveto{\pgfqpoint{0.677760in}{1.117516in}}%
\pgfpathlineto{\pgfqpoint{0.741858in}{1.102161in}}%
\pgfpathlineto{\pgfqpoint{0.740946in}{1.144486in}}%
\pgfpathlineto{\pgfqpoint{0.676704in}{1.159665in}}%
\pgfpathlineto{\pgfqpoint{0.677760in}{1.117516in}}%
\pgfpathclose%
\pgfusepath{fill}%
\end{pgfscope}%
\begin{pgfscope}%
\pgfpathrectangle{\pgfqpoint{0.407592in}{0.222904in}}{\pgfqpoint{1.411163in}{1.411163in}}%
\pgfusepath{clip}%
\pgfsetbuttcap%
\pgfsetroundjoin%
\definecolor{currentfill}{rgb}{0.141176,0.345098,0.650980}%
\pgfsetfillcolor{currentfill}%
\pgfsetfillopacity{0.080000}%
\pgfsetlinewidth{0.000000pt}%
\definecolor{currentstroke}{rgb}{0.000000,0.000000,0.000000}%
\pgfsetstrokecolor{currentstroke}%
\pgfsetdash{}{0pt}%
\pgfpathmoveto{\pgfqpoint{0.743669in}{1.018099in}}%
\pgfpathlineto{\pgfqpoint{0.808119in}{1.002242in}}%
\pgfpathlineto{\pgfqpoint{0.807362in}{1.044348in}}%
\pgfpathlineto{\pgfqpoint{0.742765in}{1.060032in}}%
\pgfpathlineto{\pgfqpoint{0.743669in}{1.018099in}}%
\pgfpathclose%
\pgfusepath{fill}%
\end{pgfscope}%
\begin{pgfscope}%
\pgfpathrectangle{\pgfqpoint{0.407592in}{0.222904in}}{\pgfqpoint{1.411163in}{1.411163in}}%
\pgfusepath{clip}%
\pgfsetbuttcap%
\pgfsetroundjoin%
\definecolor{currentfill}{rgb}{0.141176,0.345098,0.650980}%
\pgfsetfillcolor{currentfill}%
\pgfsetfillopacity{0.080000}%
\pgfsetlinewidth{0.000000pt}%
\definecolor{currentstroke}{rgb}{0.000000,0.000000,0.000000}%
\pgfsetstrokecolor{currentstroke}%
\pgfsetdash{}{0pt}%
\pgfpathmoveto{\pgfqpoint{0.809623in}{0.918615in}}%
\pgfpathlineto{\pgfqpoint{0.874426in}{0.902253in}}%
\pgfpathlineto{\pgfqpoint{0.873823in}{0.944141in}}%
\pgfpathlineto{\pgfqpoint{0.808873in}{0.960331in}}%
\pgfpathlineto{\pgfqpoint{0.809623in}{0.918615in}}%
\pgfpathclose%
\pgfusepath{fill}%
\end{pgfscope}%
\begin{pgfscope}%
\pgfpathrectangle{\pgfqpoint{0.407592in}{0.222904in}}{\pgfqpoint{1.411163in}{1.411163in}}%
\pgfusepath{clip}%
\pgfsetbuttcap%
\pgfsetroundjoin%
\definecolor{currentfill}{rgb}{0.141176,0.345098,0.650980}%
\pgfsetfillcolor{currentfill}%
\pgfsetfillopacity{0.080000}%
\pgfsetlinewidth{0.000000pt}%
\definecolor{currentstroke}{rgb}{0.000000,0.000000,0.000000}%
\pgfsetstrokecolor{currentstroke}%
\pgfsetdash{}{0pt}%
\pgfpathmoveto{\pgfqpoint{0.875622in}{0.819061in}}%
\pgfpathlineto{\pgfqpoint{0.940778in}{0.802195in}}%
\pgfpathlineto{\pgfqpoint{0.940331in}{0.843863in}}%
\pgfpathlineto{\pgfqpoint{0.875025in}{0.860560in}}%
\pgfpathlineto{\pgfqpoint{0.875622in}{0.819061in}}%
\pgfpathclose%
\pgfusepath{fill}%
\end{pgfscope}%
\begin{pgfscope}%
\pgfpathrectangle{\pgfqpoint{0.407592in}{0.222904in}}{\pgfqpoint{1.411163in}{1.411163in}}%
\pgfusepath{clip}%
\pgfsetbuttcap%
\pgfsetroundjoin%
\definecolor{currentfill}{rgb}{0.141176,0.345098,0.650980}%
\pgfsetfillcolor{currentfill}%
\pgfsetfillopacity{0.080000}%
\pgfsetlinewidth{0.000000pt}%
\definecolor{currentstroke}{rgb}{0.000000,0.000000,0.000000}%
\pgfsetstrokecolor{currentstroke}%
\pgfsetdash{}{0pt}%
\pgfpathmoveto{\pgfqpoint{0.676704in}{1.159665in}}%
\pgfpathlineto{\pgfqpoint{0.740946in}{1.144486in}}%
\pgfpathlineto{\pgfqpoint{0.740030in}{1.187010in}}%
\pgfpathlineto{\pgfqpoint{0.675643in}{1.202012in}}%
\pgfpathlineto{\pgfqpoint{0.676704in}{1.159665in}}%
\pgfpathclose%
\pgfusepath{fill}%
\end{pgfscope}%
\begin{pgfscope}%
\pgfpathrectangle{\pgfqpoint{0.407592in}{0.222904in}}{\pgfqpoint{1.411163in}{1.411163in}}%
\pgfusepath{clip}%
\pgfsetbuttcap%
\pgfsetroundjoin%
\definecolor{currentfill}{rgb}{0.141176,0.345098,0.650980}%
\pgfsetfillcolor{currentfill}%
\pgfsetfillopacity{0.080000}%
\pgfsetlinewidth{0.000000pt}%
\definecolor{currentstroke}{rgb}{0.000000,0.000000,0.000000}%
\pgfsetstrokecolor{currentstroke}%
\pgfsetdash{}{0pt}%
\pgfpathmoveto{\pgfqpoint{1.006884in}{0.743517in}}%
\pgfpathlineto{\pgfqpoint{0.941224in}{0.760721in}}%
\pgfusepath{fill}%
\end{pgfscope}%
\begin{pgfscope}%
\pgfpathrectangle{\pgfqpoint{0.407592in}{0.222904in}}{\pgfqpoint{1.411163in}{1.411163in}}%
\pgfusepath{clip}%
\pgfsetbuttcap%
\pgfsetroundjoin%
\definecolor{currentfill}{rgb}{0.141176,0.345098,0.650980}%
\pgfsetfillcolor{currentfill}%
\pgfsetfillopacity{0.080000}%
\pgfsetlinewidth{0.000000pt}%
\definecolor{currentstroke}{rgb}{0.000000,0.000000,0.000000}%
\pgfsetstrokecolor{currentstroke}%
\pgfsetdash{}{0pt}%
\pgfpathmoveto{\pgfqpoint{0.742765in}{1.060032in}}%
\pgfpathlineto{\pgfqpoint{0.807362in}{1.044348in}}%
\pgfpathlineto{\pgfqpoint{0.806601in}{1.086652in}}%
\pgfpathlineto{\pgfqpoint{0.741858in}{1.102161in}}%
\pgfpathlineto{\pgfqpoint{0.742765in}{1.060032in}}%
\pgfpathclose%
\pgfusepath{fill}%
\end{pgfscope}%
\begin{pgfscope}%
\pgfpathrectangle{\pgfqpoint{0.407592in}{0.222904in}}{\pgfqpoint{1.411163in}{1.411163in}}%
\pgfusepath{clip}%
\pgfsetbuttcap%
\pgfsetroundjoin%
\definecolor{currentfill}{rgb}{0.141176,0.345098,0.650980}%
\pgfsetfillcolor{currentfill}%
\pgfsetfillopacity{0.080000}%
\pgfsetlinewidth{0.000000pt}%
\definecolor{currentstroke}{rgb}{0.000000,0.000000,0.000000}%
\pgfsetstrokecolor{currentstroke}%
\pgfsetdash{}{0pt}%
\pgfpathmoveto{\pgfqpoint{0.808873in}{0.960331in}}%
\pgfpathlineto{\pgfqpoint{0.873823in}{0.944141in}}%
\pgfpathlineto{\pgfqpoint{0.873218in}{0.986224in}}%
\pgfpathlineto{\pgfqpoint{0.808119in}{1.002242in}}%
\pgfpathlineto{\pgfqpoint{0.808873in}{0.960331in}}%
\pgfpathclose%
\pgfusepath{fill}%
\end{pgfscope}%
\begin{pgfscope}%
\pgfpathrectangle{\pgfqpoint{0.407592in}{0.222904in}}{\pgfqpoint{1.411163in}{1.411163in}}%
\pgfusepath{clip}%
\pgfsetbuttcap%
\pgfsetroundjoin%
\definecolor{currentfill}{rgb}{0.141176,0.345098,0.650980}%
\pgfsetfillcolor{currentfill}%
\pgfsetfillopacity{0.080000}%
\pgfsetlinewidth{0.000000pt}%
\definecolor{currentstroke}{rgb}{0.000000,0.000000,0.000000}%
\pgfsetstrokecolor{currentstroke}%
\pgfsetdash{}{0pt}%
\pgfpathmoveto{\pgfqpoint{0.875025in}{0.860560in}}%
\pgfpathlineto{\pgfqpoint{0.940331in}{0.843863in}}%
\pgfpathlineto{\pgfqpoint{0.939881in}{0.885727in}}%
\pgfpathlineto{\pgfqpoint{0.874426in}{0.902253in}}%
\pgfpathlineto{\pgfqpoint{0.875025in}{0.860560in}}%
\pgfpathclose%
\pgfusepath{fill}%
\end{pgfscope}%
\begin{pgfscope}%
\pgfpathrectangle{\pgfqpoint{0.407592in}{0.222904in}}{\pgfqpoint{1.411163in}{1.411163in}}%
\pgfusepath{clip}%
\pgfsetbuttcap%
\pgfsetroundjoin%
\definecolor{currentfill}{rgb}{0.141176,0.345098,0.650980}%
\pgfsetfillcolor{currentfill}%
\pgfsetfillopacity{0.080000}%
\pgfsetlinewidth{0.000000pt}%
\definecolor{currentstroke}{rgb}{0.000000,0.000000,0.000000}%
\pgfsetstrokecolor{currentstroke}%
\pgfsetdash{}{0pt}%
\pgfpathmoveto{\pgfqpoint{0.941224in}{0.760721in}}%
\pgfpathlineto{\pgfqpoint{1.006884in}{0.743517in}}%
\pgfpathlineto{\pgfqpoint{1.006591in}{0.785160in}}%
\pgfpathlineto{\pgfqpoint{0.940778in}{0.802195in}}%
\pgfpathlineto{\pgfqpoint{0.941224in}{0.760721in}}%
\pgfpathclose%
\pgfusepath{fill}%
\end{pgfscope}%
\begin{pgfscope}%
\pgfpathrectangle{\pgfqpoint{0.407592in}{0.222904in}}{\pgfqpoint{1.411163in}{1.411163in}}%
\pgfusepath{clip}%
\pgfsetbuttcap%
\pgfsetroundjoin%
\definecolor{currentfill}{rgb}{0.141176,0.345098,0.650980}%
\pgfsetfillcolor{currentfill}%
\pgfsetfillopacity{0.080000}%
\pgfsetlinewidth{0.000000pt}%
\definecolor{currentstroke}{rgb}{0.000000,0.000000,0.000000}%
\pgfsetstrokecolor{currentstroke}%
\pgfsetdash{}{0pt}%
\pgfpathmoveto{\pgfqpoint{0.675643in}{1.202012in}}%
\pgfpathlineto{\pgfqpoint{0.740030in}{1.187010in}}%
\pgfpathlineto{\pgfqpoint{0.739110in}{1.229733in}}%
\pgfpathlineto{\pgfqpoint{0.674577in}{1.244555in}}%
\pgfpathlineto{\pgfqpoint{0.675643in}{1.202012in}}%
\pgfpathclose%
\pgfusepath{fill}%
\end{pgfscope}%
\begin{pgfscope}%
\pgfpathrectangle{\pgfqpoint{0.407592in}{0.222904in}}{\pgfqpoint{1.411163in}{1.411163in}}%
\pgfusepath{clip}%
\pgfsetbuttcap%
\pgfsetroundjoin%
\definecolor{currentfill}{rgb}{0.141176,0.345098,0.650980}%
\pgfsetfillcolor{currentfill}%
\pgfsetfillopacity{0.080000}%
\pgfsetlinewidth{0.000000pt}%
\definecolor{currentstroke}{rgb}{0.000000,0.000000,0.000000}%
\pgfsetstrokecolor{currentstroke}%
\pgfsetdash{}{0pt}%
\pgfpathmoveto{\pgfqpoint{0.741858in}{1.102161in}}%
\pgfpathlineto{\pgfqpoint{0.806601in}{1.086652in}}%
\pgfpathlineto{\pgfqpoint{0.805837in}{1.129154in}}%
\pgfpathlineto{\pgfqpoint{0.740946in}{1.144486in}}%
\pgfpathlineto{\pgfqpoint{0.741858in}{1.102161in}}%
\pgfpathclose%
\pgfusepath{fill}%
\end{pgfscope}%
\begin{pgfscope}%
\pgfpathrectangle{\pgfqpoint{0.407592in}{0.222904in}}{\pgfqpoint{1.411163in}{1.411163in}}%
\pgfusepath{clip}%
\pgfsetbuttcap%
\pgfsetroundjoin%
\definecolor{currentfill}{rgb}{0.141176,0.345098,0.650980}%
\pgfsetfillcolor{currentfill}%
\pgfsetfillopacity{0.080000}%
\pgfsetlinewidth{0.000000pt}%
\definecolor{currentstroke}{rgb}{0.000000,0.000000,0.000000}%
\pgfsetstrokecolor{currentstroke}%
\pgfsetdash{}{0pt}%
\pgfpathmoveto{\pgfqpoint{1.006884in}{0.743517in}}%
\pgfusepath{fill}%
\end{pgfscope}%
\begin{pgfscope}%
\pgfpathrectangle{\pgfqpoint{0.407592in}{0.222904in}}{\pgfqpoint{1.411163in}{1.411163in}}%
\pgfusepath{clip}%
\pgfsetbuttcap%
\pgfsetroundjoin%
\definecolor{currentfill}{rgb}{0.141176,0.345098,0.650980}%
\pgfsetfillcolor{currentfill}%
\pgfsetfillopacity{0.080000}%
\pgfsetlinewidth{0.000000pt}%
\definecolor{currentstroke}{rgb}{0.000000,0.000000,0.000000}%
\pgfsetstrokecolor{currentstroke}%
\pgfsetdash{}{0pt}%
\pgfpathmoveto{\pgfqpoint{0.808119in}{1.002242in}}%
\pgfpathlineto{\pgfqpoint{0.873218in}{0.986224in}}%
\pgfpathlineto{\pgfqpoint{0.872610in}{1.028506in}}%
\pgfpathlineto{\pgfqpoint{0.807362in}{1.044348in}}%
\pgfpathlineto{\pgfqpoint{0.808119in}{1.002242in}}%
\pgfpathclose%
\pgfusepath{fill}%
\end{pgfscope}%
\begin{pgfscope}%
\pgfpathrectangle{\pgfqpoint{0.407592in}{0.222904in}}{\pgfqpoint{1.411163in}{1.411163in}}%
\pgfusepath{clip}%
\pgfsetbuttcap%
\pgfsetroundjoin%
\definecolor{currentfill}{rgb}{0.141176,0.345098,0.650980}%
\pgfsetfillcolor{currentfill}%
\pgfsetfillopacity{0.080000}%
\pgfsetlinewidth{0.000000pt}%
\definecolor{currentstroke}{rgb}{0.000000,0.000000,0.000000}%
\pgfsetstrokecolor{currentstroke}%
\pgfsetdash{}{0pt}%
\pgfpathmoveto{\pgfqpoint{0.874426in}{0.902253in}}%
\pgfpathlineto{\pgfqpoint{0.939881in}{0.885727in}}%
\pgfpathlineto{\pgfqpoint{0.939429in}{0.927787in}}%
\pgfpathlineto{\pgfqpoint{0.873823in}{0.944141in}}%
\pgfpathlineto{\pgfqpoint{0.874426in}{0.902253in}}%
\pgfpathclose%
\pgfusepath{fill}%
\end{pgfscope}%
\begin{pgfscope}%
\pgfpathrectangle{\pgfqpoint{0.407592in}{0.222904in}}{\pgfqpoint{1.411163in}{1.411163in}}%
\pgfusepath{clip}%
\pgfsetbuttcap%
\pgfsetroundjoin%
\definecolor{currentfill}{rgb}{0.141176,0.345098,0.650980}%
\pgfsetfillcolor{currentfill}%
\pgfsetfillopacity{0.080000}%
\pgfsetlinewidth{0.000000pt}%
\definecolor{currentstroke}{rgb}{0.000000,0.000000,0.000000}%
\pgfsetstrokecolor{currentstroke}%
\pgfsetdash{}{0pt}%
\pgfpathmoveto{\pgfqpoint{0.940778in}{0.802195in}}%
\pgfpathlineto{\pgfqpoint{1.006591in}{0.785160in}}%
\pgfpathlineto{\pgfqpoint{1.006295in}{0.826998in}}%
\pgfpathlineto{\pgfqpoint{0.940331in}{0.843863in}}%
\pgfpathlineto{\pgfqpoint{0.940778in}{0.802195in}}%
\pgfpathclose%
\pgfusepath{fill}%
\end{pgfscope}%
\begin{pgfscope}%
\pgfpathrectangle{\pgfqpoint{0.407592in}{0.222904in}}{\pgfqpoint{1.411163in}{1.411163in}}%
\pgfusepath{clip}%
\pgfsetbuttcap%
\pgfsetroundjoin%
\definecolor{currentfill}{rgb}{0.141176,0.345098,0.650980}%
\pgfsetfillcolor{currentfill}%
\pgfsetfillopacity{0.080000}%
\pgfsetlinewidth{0.000000pt}%
\definecolor{currentstroke}{rgb}{0.000000,0.000000,0.000000}%
\pgfsetstrokecolor{currentstroke}%
\pgfsetdash{}{0pt}%
\pgfpathmoveto{\pgfqpoint{0.740946in}{1.144486in}}%
\pgfpathlineto{\pgfqpoint{0.805837in}{1.129154in}}%
\pgfpathlineto{\pgfqpoint{0.805069in}{1.171856in}}%
\pgfpathlineto{\pgfqpoint{0.740030in}{1.187010in}}%
\pgfpathlineto{\pgfqpoint{0.740946in}{1.144486in}}%
\pgfpathclose%
\pgfusepath{fill}%
\end{pgfscope}%
\begin{pgfscope}%
\pgfpathrectangle{\pgfqpoint{0.407592in}{0.222904in}}{\pgfqpoint{1.411163in}{1.411163in}}%
\pgfusepath{clip}%
\pgfsetbuttcap%
\pgfsetroundjoin%
\definecolor{currentfill}{rgb}{0.141176,0.345098,0.650980}%
\pgfsetfillcolor{currentfill}%
\pgfsetfillopacity{0.080000}%
\pgfsetlinewidth{0.000000pt}%
\definecolor{currentstroke}{rgb}{0.000000,0.000000,0.000000}%
\pgfsetstrokecolor{currentstroke}%
\pgfsetdash{}{0pt}%
\pgfpathmoveto{\pgfqpoint{0.807362in}{1.044348in}}%
\pgfpathlineto{\pgfqpoint{0.872610in}{1.028506in}}%
\pgfpathlineto{\pgfqpoint{0.871999in}{1.070986in}}%
\pgfpathlineto{\pgfqpoint{0.806601in}{1.086652in}}%
\pgfpathlineto{\pgfqpoint{0.807362in}{1.044348in}}%
\pgfpathclose%
\pgfusepath{fill}%
\end{pgfscope}%
\begin{pgfscope}%
\pgfpathrectangle{\pgfqpoint{0.407592in}{0.222904in}}{\pgfqpoint{1.411163in}{1.411163in}}%
\pgfusepath{clip}%
\pgfsetbuttcap%
\pgfsetroundjoin%
\definecolor{currentfill}{rgb}{0.141176,0.345098,0.650980}%
\pgfsetfillcolor{currentfill}%
\pgfsetfillopacity{0.080000}%
\pgfsetlinewidth{0.000000pt}%
\definecolor{currentstroke}{rgb}{0.000000,0.000000,0.000000}%
\pgfsetstrokecolor{currentstroke}%
\pgfsetdash{}{0pt}%
\pgfpathmoveto{\pgfqpoint{0.873823in}{0.944141in}}%
\pgfpathlineto{\pgfqpoint{0.939429in}{0.927787in}}%
\pgfpathlineto{\pgfqpoint{0.938975in}{0.970045in}}%
\pgfpathlineto{\pgfqpoint{0.873218in}{0.986224in}}%
\pgfpathlineto{\pgfqpoint{0.873823in}{0.944141in}}%
\pgfpathclose%
\pgfusepath{fill}%
\end{pgfscope}%
\begin{pgfscope}%
\pgfpathrectangle{\pgfqpoint{0.407592in}{0.222904in}}{\pgfqpoint{1.411163in}{1.411163in}}%
\pgfusepath{clip}%
\pgfsetbuttcap%
\pgfsetroundjoin%
\definecolor{currentfill}{rgb}{0.141176,0.345098,0.650980}%
\pgfsetfillcolor{currentfill}%
\pgfsetfillopacity{0.080000}%
\pgfsetlinewidth{0.000000pt}%
\definecolor{currentstroke}{rgb}{0.000000,0.000000,0.000000}%
\pgfsetstrokecolor{currentstroke}%
\pgfsetdash{}{0pt}%
\pgfpathmoveto{\pgfqpoint{1.006884in}{0.743517in}}%
\pgfpathlineto{\pgfqpoint{1.073069in}{0.767951in}}%
\pgfpathlineto{\pgfqpoint{1.006591in}{0.785160in}}%
\pgfusepath{fill}%
\end{pgfscope}%
\begin{pgfscope}%
\pgfpathrectangle{\pgfqpoint{0.407592in}{0.222904in}}{\pgfqpoint{1.411163in}{1.411163in}}%
\pgfusepath{clip}%
\pgfsetbuttcap%
\pgfsetroundjoin%
\definecolor{currentfill}{rgb}{0.141176,0.345098,0.650980}%
\pgfsetfillcolor{currentfill}%
\pgfsetfillopacity{0.080000}%
\pgfsetlinewidth{0.000000pt}%
\definecolor{currentstroke}{rgb}{0.000000,0.000000,0.000000}%
\pgfsetstrokecolor{currentstroke}%
\pgfsetdash{}{0pt}%
\pgfpathmoveto{\pgfqpoint{0.940331in}{0.843863in}}%
\pgfpathlineto{\pgfqpoint{1.006295in}{0.826998in}}%
\pgfpathlineto{\pgfqpoint{1.005999in}{0.869033in}}%
\pgfpathlineto{\pgfqpoint{0.939881in}{0.885727in}}%
\pgfpathlineto{\pgfqpoint{0.940331in}{0.843863in}}%
\pgfpathclose%
\pgfusepath{fill}%
\end{pgfscope}%
\begin{pgfscope}%
\pgfpathrectangle{\pgfqpoint{0.407592in}{0.222904in}}{\pgfqpoint{1.411163in}{1.411163in}}%
\pgfusepath{clip}%
\pgfsetbuttcap%
\pgfsetroundjoin%
\definecolor{currentfill}{rgb}{0.141176,0.345098,0.650980}%
\pgfsetfillcolor{currentfill}%
\pgfsetfillopacity{0.080000}%
\pgfsetlinewidth{0.000000pt}%
\definecolor{currentstroke}{rgb}{0.000000,0.000000,0.000000}%
\pgfsetstrokecolor{currentstroke}%
\pgfsetdash{}{0pt}%
\pgfpathmoveto{\pgfqpoint{0.740030in}{1.187010in}}%
\pgfpathlineto{\pgfqpoint{0.805069in}{1.171856in}}%
\pgfpathlineto{\pgfqpoint{0.804297in}{1.214760in}}%
\pgfpathlineto{\pgfqpoint{0.739110in}{1.229733in}}%
\pgfpathlineto{\pgfqpoint{0.740030in}{1.187010in}}%
\pgfpathclose%
\pgfusepath{fill}%
\end{pgfscope}%
\begin{pgfscope}%
\pgfpathrectangle{\pgfqpoint{0.407592in}{0.222904in}}{\pgfqpoint{1.411163in}{1.411163in}}%
\pgfusepath{clip}%
\pgfsetbuttcap%
\pgfsetroundjoin%
\definecolor{currentfill}{rgb}{0.141176,0.345098,0.650980}%
\pgfsetfillcolor{currentfill}%
\pgfsetfillopacity{0.080000}%
\pgfsetlinewidth{0.000000pt}%
\definecolor{currentstroke}{rgb}{0.000000,0.000000,0.000000}%
\pgfsetstrokecolor{currentstroke}%
\pgfsetdash{}{0pt}%
\pgfpathmoveto{\pgfqpoint{0.806601in}{1.086652in}}%
\pgfpathlineto{\pgfqpoint{0.871999in}{1.070986in}}%
\pgfpathlineto{\pgfqpoint{0.871385in}{1.113667in}}%
\pgfpathlineto{\pgfqpoint{0.805837in}{1.129154in}}%
\pgfpathlineto{\pgfqpoint{0.806601in}{1.086652in}}%
\pgfpathclose%
\pgfusepath{fill}%
\end{pgfscope}%
\begin{pgfscope}%
\pgfpathrectangle{\pgfqpoint{0.407592in}{0.222904in}}{\pgfqpoint{1.411163in}{1.411163in}}%
\pgfusepath{clip}%
\pgfsetbuttcap%
\pgfsetroundjoin%
\definecolor{currentfill}{rgb}{0.141176,0.345098,0.650980}%
\pgfsetfillcolor{currentfill}%
\pgfsetfillopacity{0.080000}%
\pgfsetlinewidth{0.000000pt}%
\definecolor{currentstroke}{rgb}{0.000000,0.000000,0.000000}%
\pgfsetstrokecolor{currentstroke}%
\pgfsetdash{}{0pt}%
\pgfpathmoveto{\pgfqpoint{0.873218in}{0.986224in}}%
\pgfpathlineto{\pgfqpoint{0.938975in}{0.970045in}}%
\pgfpathlineto{\pgfqpoint{0.938519in}{1.012502in}}%
\pgfpathlineto{\pgfqpoint{0.872610in}{1.028506in}}%
\pgfpathlineto{\pgfqpoint{0.873218in}{0.986224in}}%
\pgfpathclose%
\pgfusepath{fill}%
\end{pgfscope}%
\begin{pgfscope}%
\pgfpathrectangle{\pgfqpoint{0.407592in}{0.222904in}}{\pgfqpoint{1.411163in}{1.411163in}}%
\pgfusepath{clip}%
\pgfsetbuttcap%
\pgfsetroundjoin%
\definecolor{currentfill}{rgb}{0.141176,0.345098,0.650980}%
\pgfsetfillcolor{currentfill}%
\pgfsetfillopacity{0.080000}%
\pgfsetlinewidth{0.000000pt}%
\definecolor{currentstroke}{rgb}{0.000000,0.000000,0.000000}%
\pgfsetstrokecolor{currentstroke}%
\pgfsetdash{}{0pt}%
\pgfpathmoveto{\pgfqpoint{0.939881in}{0.885727in}}%
\pgfpathlineto{\pgfqpoint{1.005999in}{0.869033in}}%
\pgfpathlineto{\pgfqpoint{1.005701in}{0.911267in}}%
\pgfpathlineto{\pgfqpoint{0.939429in}{0.927787in}}%
\pgfpathlineto{\pgfqpoint{0.939881in}{0.885727in}}%
\pgfpathclose%
\pgfusepath{fill}%
\end{pgfscope}%
\begin{pgfscope}%
\pgfpathrectangle{\pgfqpoint{0.407592in}{0.222904in}}{\pgfqpoint{1.411163in}{1.411163in}}%
\pgfusepath{clip}%
\pgfsetbuttcap%
\pgfsetroundjoin%
\definecolor{currentfill}{rgb}{0.141176,0.345098,0.650980}%
\pgfsetfillcolor{currentfill}%
\pgfsetfillopacity{0.080000}%
\pgfsetlinewidth{0.000000pt}%
\definecolor{currentstroke}{rgb}{0.000000,0.000000,0.000000}%
\pgfsetstrokecolor{currentstroke}%
\pgfsetdash{}{0pt}%
\pgfpathmoveto{\pgfqpoint{1.006591in}{0.785160in}}%
\pgfpathlineto{\pgfqpoint{1.073069in}{0.767951in}}%
\pgfpathlineto{\pgfqpoint{1.072929in}{0.809961in}}%
\pgfpathlineto{\pgfqpoint{1.006295in}{0.826998in}}%
\pgfpathlineto{\pgfqpoint{1.006591in}{0.785160in}}%
\pgfpathclose%
\pgfusepath{fill}%
\end{pgfscope}%
\begin{pgfscope}%
\pgfpathrectangle{\pgfqpoint{0.407592in}{0.222904in}}{\pgfqpoint{1.411163in}{1.411163in}}%
\pgfusepath{clip}%
\pgfsetbuttcap%
\pgfsetroundjoin%
\definecolor{currentfill}{rgb}{0.141176,0.345098,0.650980}%
\pgfsetfillcolor{currentfill}%
\pgfsetfillopacity{0.080000}%
\pgfsetlinewidth{0.000000pt}%
\definecolor{currentstroke}{rgb}{0.000000,0.000000,0.000000}%
\pgfsetstrokecolor{currentstroke}%
\pgfsetdash{}{0pt}%
\pgfpathmoveto{\pgfqpoint{0.805837in}{1.129154in}}%
\pgfpathlineto{\pgfqpoint{0.871385in}{1.113667in}}%
\pgfpathlineto{\pgfqpoint{0.870768in}{1.156549in}}%
\pgfpathlineto{\pgfqpoint{0.805069in}{1.171856in}}%
\pgfpathlineto{\pgfqpoint{0.805837in}{1.129154in}}%
\pgfpathclose%
\pgfusepath{fill}%
\end{pgfscope}%
\begin{pgfscope}%
\pgfpathrectangle{\pgfqpoint{0.407592in}{0.222904in}}{\pgfqpoint{1.411163in}{1.411163in}}%
\pgfusepath{clip}%
\pgfsetbuttcap%
\pgfsetroundjoin%
\definecolor{currentfill}{rgb}{0.141176,0.345098,0.650980}%
\pgfsetfillcolor{currentfill}%
\pgfsetfillopacity{0.080000}%
\pgfsetlinewidth{0.000000pt}%
\definecolor{currentstroke}{rgb}{0.000000,0.000000,0.000000}%
\pgfsetstrokecolor{currentstroke}%
\pgfsetdash{}{0pt}%
\pgfpathmoveto{\pgfqpoint{1.073069in}{0.767951in}}%
\pgfusepath{fill}%
\end{pgfscope}%
\begin{pgfscope}%
\pgfpathrectangle{\pgfqpoint{0.407592in}{0.222904in}}{\pgfqpoint{1.411163in}{1.411163in}}%
\pgfusepath{clip}%
\pgfsetbuttcap%
\pgfsetroundjoin%
\definecolor{currentfill}{rgb}{0.141176,0.345098,0.650980}%
\pgfsetfillcolor{currentfill}%
\pgfsetfillopacity{0.080000}%
\pgfsetlinewidth{0.000000pt}%
\definecolor{currentstroke}{rgb}{0.000000,0.000000,0.000000}%
\pgfsetstrokecolor{currentstroke}%
\pgfsetdash{}{0pt}%
\pgfpathmoveto{\pgfqpoint{0.872610in}{1.028506in}}%
\pgfpathlineto{\pgfqpoint{0.938519in}{1.012502in}}%
\pgfpathlineto{\pgfqpoint{0.938061in}{1.055161in}}%
\pgfpathlineto{\pgfqpoint{0.871999in}{1.070986in}}%
\pgfpathlineto{\pgfqpoint{0.872610in}{1.028506in}}%
\pgfpathclose%
\pgfusepath{fill}%
\end{pgfscope}%
\begin{pgfscope}%
\pgfpathrectangle{\pgfqpoint{0.407592in}{0.222904in}}{\pgfqpoint{1.411163in}{1.411163in}}%
\pgfusepath{clip}%
\pgfsetbuttcap%
\pgfsetroundjoin%
\definecolor{currentfill}{rgb}{0.141176,0.345098,0.650980}%
\pgfsetfillcolor{currentfill}%
\pgfsetfillopacity{0.080000}%
\pgfsetlinewidth{0.000000pt}%
\definecolor{currentstroke}{rgb}{0.000000,0.000000,0.000000}%
\pgfsetstrokecolor{currentstroke}%
\pgfsetdash{}{0pt}%
\pgfpathmoveto{\pgfqpoint{0.939429in}{0.927787in}}%
\pgfpathlineto{\pgfqpoint{1.005701in}{0.911267in}}%
\pgfpathlineto{\pgfqpoint{1.005401in}{0.953701in}}%
\pgfpathlineto{\pgfqpoint{0.938975in}{0.970045in}}%
\pgfpathlineto{\pgfqpoint{0.939429in}{0.927787in}}%
\pgfpathclose%
\pgfusepath{fill}%
\end{pgfscope}%
\begin{pgfscope}%
\pgfpathrectangle{\pgfqpoint{0.407592in}{0.222904in}}{\pgfqpoint{1.411163in}{1.411163in}}%
\pgfusepath{clip}%
\pgfsetbuttcap%
\pgfsetroundjoin%
\definecolor{currentfill}{rgb}{0.141176,0.345098,0.650980}%
\pgfsetfillcolor{currentfill}%
\pgfsetfillopacity{0.080000}%
\pgfsetlinewidth{0.000000pt}%
\definecolor{currentstroke}{rgb}{0.000000,0.000000,0.000000}%
\pgfsetstrokecolor{currentstroke}%
\pgfsetdash{}{0pt}%
\pgfpathmoveto{\pgfqpoint{1.006295in}{0.826998in}}%
\pgfpathlineto{\pgfqpoint{1.072929in}{0.809961in}}%
\pgfpathlineto{\pgfqpoint{1.072789in}{0.852170in}}%
\pgfpathlineto{\pgfqpoint{1.005999in}{0.869033in}}%
\pgfpathlineto{\pgfqpoint{1.006295in}{0.826998in}}%
\pgfpathclose%
\pgfusepath{fill}%
\end{pgfscope}%
\begin{pgfscope}%
\pgfpathrectangle{\pgfqpoint{0.407592in}{0.222904in}}{\pgfqpoint{1.411163in}{1.411163in}}%
\pgfusepath{clip}%
\pgfsetbuttcap%
\pgfsetroundjoin%
\definecolor{currentfill}{rgb}{0.141176,0.345098,0.650980}%
\pgfsetfillcolor{currentfill}%
\pgfsetfillopacity{0.080000}%
\pgfsetlinewidth{0.000000pt}%
\definecolor{currentstroke}{rgb}{0.000000,0.000000,0.000000}%
\pgfsetstrokecolor{currentstroke}%
\pgfsetdash{}{0pt}%
\pgfpathmoveto{\pgfqpoint{0.805069in}{1.171856in}}%
\pgfpathlineto{\pgfqpoint{0.870768in}{1.156549in}}%
\pgfpathlineto{\pgfqpoint{0.870148in}{1.199635in}}%
\pgfpathlineto{\pgfqpoint{0.804297in}{1.214760in}}%
\pgfpathlineto{\pgfqpoint{0.805069in}{1.171856in}}%
\pgfpathclose%
\pgfusepath{fill}%
\end{pgfscope}%
\begin{pgfscope}%
\pgfpathrectangle{\pgfqpoint{0.407592in}{0.222904in}}{\pgfqpoint{1.411163in}{1.411163in}}%
\pgfusepath{clip}%
\pgfsetbuttcap%
\pgfsetroundjoin%
\definecolor{currentfill}{rgb}{0.141176,0.345098,0.650980}%
\pgfsetfillcolor{currentfill}%
\pgfsetfillopacity{0.080000}%
\pgfsetlinewidth{0.000000pt}%
\definecolor{currentstroke}{rgb}{0.000000,0.000000,0.000000}%
\pgfsetstrokecolor{currentstroke}%
\pgfsetdash{}{0pt}%
\pgfpathmoveto{\pgfqpoint{0.871999in}{1.070986in}}%
\pgfpathlineto{\pgfqpoint{0.938061in}{1.055161in}}%
\pgfpathlineto{\pgfqpoint{0.937601in}{1.098021in}}%
\pgfpathlineto{\pgfqpoint{0.871385in}{1.113667in}}%
\pgfpathlineto{\pgfqpoint{0.871999in}{1.070986in}}%
\pgfpathclose%
\pgfusepath{fill}%
\end{pgfscope}%
\begin{pgfscope}%
\pgfpathrectangle{\pgfqpoint{0.407592in}{0.222904in}}{\pgfqpoint{1.411163in}{1.411163in}}%
\pgfusepath{clip}%
\pgfsetbuttcap%
\pgfsetroundjoin%
\definecolor{currentfill}{rgb}{0.141176,0.345098,0.650980}%
\pgfsetfillcolor{currentfill}%
\pgfsetfillopacity{0.080000}%
\pgfsetlinewidth{0.000000pt}%
\definecolor{currentstroke}{rgb}{0.000000,0.000000,0.000000}%
\pgfsetstrokecolor{currentstroke}%
\pgfsetdash{}{0pt}%
\pgfpathmoveto{\pgfqpoint{0.938975in}{0.970045in}}%
\pgfpathlineto{\pgfqpoint{1.005401in}{0.953701in}}%
\pgfpathlineto{\pgfqpoint{1.005101in}{0.996336in}}%
\pgfpathlineto{\pgfqpoint{0.938519in}{1.012502in}}%
\pgfpathlineto{\pgfqpoint{0.938975in}{0.970045in}}%
\pgfpathclose%
\pgfusepath{fill}%
\end{pgfscope}%
\begin{pgfscope}%
\pgfpathrectangle{\pgfqpoint{0.407592in}{0.222904in}}{\pgfqpoint{1.411163in}{1.411163in}}%
\pgfusepath{clip}%
\pgfsetbuttcap%
\pgfsetroundjoin%
\definecolor{currentfill}{rgb}{0.141176,0.345098,0.650980}%
\pgfsetfillcolor{currentfill}%
\pgfsetfillopacity{0.080000}%
\pgfsetlinewidth{0.000000pt}%
\definecolor{currentstroke}{rgb}{0.000000,0.000000,0.000000}%
\pgfsetstrokecolor{currentstroke}%
\pgfsetdash{}{0pt}%
\pgfpathmoveto{\pgfqpoint{1.073069in}{0.767951in}}%
\pgfpathlineto{\pgfqpoint{1.140243in}{0.792751in}}%
\pgfpathlineto{\pgfqpoint{1.072929in}{0.809961in}}%
\pgfusepath{fill}%
\end{pgfscope}%
\begin{pgfscope}%
\pgfpathrectangle{\pgfqpoint{0.407592in}{0.222904in}}{\pgfqpoint{1.411163in}{1.411163in}}%
\pgfusepath{clip}%
\pgfsetbuttcap%
\pgfsetroundjoin%
\definecolor{currentfill}{rgb}{0.141176,0.345098,0.650980}%
\pgfsetfillcolor{currentfill}%
\pgfsetfillopacity{0.080000}%
\pgfsetlinewidth{0.000000pt}%
\definecolor{currentstroke}{rgb}{0.000000,0.000000,0.000000}%
\pgfsetstrokecolor{currentstroke}%
\pgfsetdash{}{0pt}%
\pgfpathmoveto{\pgfqpoint{1.005999in}{0.869033in}}%
\pgfpathlineto{\pgfqpoint{1.072789in}{0.852170in}}%
\pgfpathlineto{\pgfqpoint{1.072648in}{0.894579in}}%
\pgfpathlineto{\pgfqpoint{1.005701in}{0.911267in}}%
\pgfpathlineto{\pgfqpoint{1.005999in}{0.869033in}}%
\pgfpathclose%
\pgfusepath{fill}%
\end{pgfscope}%
\begin{pgfscope}%
\pgfpathrectangle{\pgfqpoint{0.407592in}{0.222904in}}{\pgfqpoint{1.411163in}{1.411163in}}%
\pgfusepath{clip}%
\pgfsetbuttcap%
\pgfsetroundjoin%
\definecolor{currentfill}{rgb}{0.141176,0.345098,0.650980}%
\pgfsetfillcolor{currentfill}%
\pgfsetfillopacity{0.080000}%
\pgfsetlinewidth{0.000000pt}%
\definecolor{currentstroke}{rgb}{0.000000,0.000000,0.000000}%
\pgfsetstrokecolor{currentstroke}%
\pgfsetdash{}{0pt}%
\pgfpathmoveto{\pgfqpoint{0.871385in}{1.113667in}}%
\pgfpathlineto{\pgfqpoint{0.937601in}{1.098021in}}%
\pgfpathlineto{\pgfqpoint{0.937138in}{1.141086in}}%
\pgfpathlineto{\pgfqpoint{0.870768in}{1.156549in}}%
\pgfpathlineto{\pgfqpoint{0.871385in}{1.113667in}}%
\pgfpathclose%
\pgfusepath{fill}%
\end{pgfscope}%
\begin{pgfscope}%
\pgfpathrectangle{\pgfqpoint{0.407592in}{0.222904in}}{\pgfqpoint{1.411163in}{1.411163in}}%
\pgfusepath{clip}%
\pgfsetbuttcap%
\pgfsetroundjoin%
\definecolor{currentfill}{rgb}{0.141176,0.345098,0.650980}%
\pgfsetfillcolor{currentfill}%
\pgfsetfillopacity{0.080000}%
\pgfsetlinewidth{0.000000pt}%
\definecolor{currentstroke}{rgb}{0.000000,0.000000,0.000000}%
\pgfsetstrokecolor{currentstroke}%
\pgfsetdash{}{0pt}%
\pgfpathmoveto{\pgfqpoint{0.938519in}{1.012502in}}%
\pgfpathlineto{\pgfqpoint{1.005101in}{0.996336in}}%
\pgfpathlineto{\pgfqpoint{1.004798in}{1.039174in}}%
\pgfpathlineto{\pgfqpoint{0.938061in}{1.055161in}}%
\pgfpathlineto{\pgfqpoint{0.938519in}{1.012502in}}%
\pgfpathclose%
\pgfusepath{fill}%
\end{pgfscope}%
\begin{pgfscope}%
\pgfpathrectangle{\pgfqpoint{0.407592in}{0.222904in}}{\pgfqpoint{1.411163in}{1.411163in}}%
\pgfusepath{clip}%
\pgfsetbuttcap%
\pgfsetroundjoin%
\definecolor{currentfill}{rgb}{0.141176,0.345098,0.650980}%
\pgfsetfillcolor{currentfill}%
\pgfsetfillopacity{0.080000}%
\pgfsetlinewidth{0.000000pt}%
\definecolor{currentstroke}{rgb}{0.000000,0.000000,0.000000}%
\pgfsetstrokecolor{currentstroke}%
\pgfsetdash{}{0pt}%
\pgfpathmoveto{\pgfqpoint{1.005701in}{0.911267in}}%
\pgfpathlineto{\pgfqpoint{1.072648in}{0.894579in}}%
\pgfpathlineto{\pgfqpoint{1.072506in}{0.937190in}}%
\pgfpathlineto{\pgfqpoint{1.005401in}{0.953701in}}%
\pgfpathlineto{\pgfqpoint{1.005701in}{0.911267in}}%
\pgfpathclose%
\pgfusepath{fill}%
\end{pgfscope}%
\begin{pgfscope}%
\pgfpathrectangle{\pgfqpoint{0.407592in}{0.222904in}}{\pgfqpoint{1.411163in}{1.411163in}}%
\pgfusepath{clip}%
\pgfsetbuttcap%
\pgfsetroundjoin%
\definecolor{currentfill}{rgb}{0.141176,0.345098,0.650980}%
\pgfsetfillcolor{currentfill}%
\pgfsetfillopacity{0.080000}%
\pgfsetlinewidth{0.000000pt}%
\definecolor{currentstroke}{rgb}{0.000000,0.000000,0.000000}%
\pgfsetstrokecolor{currentstroke}%
\pgfsetdash{}{0pt}%
\pgfpathmoveto{\pgfqpoint{1.072929in}{0.809961in}}%
\pgfpathlineto{\pgfqpoint{1.140243in}{0.792751in}}%
\pgfpathlineto{\pgfqpoint{1.140262in}{0.835135in}}%
\pgfpathlineto{\pgfqpoint{1.072789in}{0.852170in}}%
\pgfpathlineto{\pgfqpoint{1.072929in}{0.809961in}}%
\pgfpathclose%
\pgfusepath{fill}%
\end{pgfscope}%
\begin{pgfscope}%
\pgfpathrectangle{\pgfqpoint{0.407592in}{0.222904in}}{\pgfqpoint{1.411163in}{1.411163in}}%
\pgfusepath{clip}%
\pgfsetbuttcap%
\pgfsetroundjoin%
\definecolor{currentfill}{rgb}{0.141176,0.345098,0.650980}%
\pgfsetfillcolor{currentfill}%
\pgfsetfillopacity{0.080000}%
\pgfsetlinewidth{0.000000pt}%
\definecolor{currentstroke}{rgb}{0.000000,0.000000,0.000000}%
\pgfsetstrokecolor{currentstroke}%
\pgfsetdash{}{0pt}%
\pgfpathmoveto{\pgfqpoint{0.870768in}{1.156549in}}%
\pgfpathlineto{\pgfqpoint{0.937138in}{1.141086in}}%
\pgfpathlineto{\pgfqpoint{0.936673in}{1.184355in}}%
\pgfpathlineto{\pgfqpoint{0.870148in}{1.199635in}}%
\pgfpathlineto{\pgfqpoint{0.870768in}{1.156549in}}%
\pgfpathclose%
\pgfusepath{fill}%
\end{pgfscope}%
\begin{pgfscope}%
\pgfpathrectangle{\pgfqpoint{0.407592in}{0.222904in}}{\pgfqpoint{1.411163in}{1.411163in}}%
\pgfusepath{clip}%
\pgfsetbuttcap%
\pgfsetroundjoin%
\definecolor{currentfill}{rgb}{0.141176,0.345098,0.650980}%
\pgfsetfillcolor{currentfill}%
\pgfsetfillopacity{0.080000}%
\pgfsetlinewidth{0.000000pt}%
\definecolor{currentstroke}{rgb}{0.000000,0.000000,0.000000}%
\pgfsetstrokecolor{currentstroke}%
\pgfsetdash{}{0pt}%
\pgfpathmoveto{\pgfqpoint{1.140243in}{0.792751in}}%
\pgfusepath{fill}%
\end{pgfscope}%
\begin{pgfscope}%
\pgfpathrectangle{\pgfqpoint{0.407592in}{0.222904in}}{\pgfqpoint{1.411163in}{1.411163in}}%
\pgfusepath{clip}%
\pgfsetbuttcap%
\pgfsetroundjoin%
\definecolor{currentfill}{rgb}{0.141176,0.345098,0.650980}%
\pgfsetfillcolor{currentfill}%
\pgfsetfillopacity{0.080000}%
\pgfsetlinewidth{0.000000pt}%
\definecolor{currentstroke}{rgb}{0.000000,0.000000,0.000000}%
\pgfsetstrokecolor{currentstroke}%
\pgfsetdash{}{0pt}%
\pgfpathmoveto{\pgfqpoint{0.938061in}{1.055161in}}%
\pgfpathlineto{\pgfqpoint{1.004798in}{1.039174in}}%
\pgfpathlineto{\pgfqpoint{1.004495in}{1.082216in}}%
\pgfpathlineto{\pgfqpoint{0.937601in}{1.098021in}}%
\pgfpathlineto{\pgfqpoint{0.938061in}{1.055161in}}%
\pgfpathclose%
\pgfusepath{fill}%
\end{pgfscope}%
\begin{pgfscope}%
\pgfpathrectangle{\pgfqpoint{0.407592in}{0.222904in}}{\pgfqpoint{1.411163in}{1.411163in}}%
\pgfusepath{clip}%
\pgfsetbuttcap%
\pgfsetroundjoin%
\definecolor{currentfill}{rgb}{0.141176,0.345098,0.650980}%
\pgfsetfillcolor{currentfill}%
\pgfsetfillopacity{0.080000}%
\pgfsetlinewidth{0.000000pt}%
\definecolor{currentstroke}{rgb}{0.000000,0.000000,0.000000}%
\pgfsetstrokecolor{currentstroke}%
\pgfsetdash{}{0pt}%
\pgfpathmoveto{\pgfqpoint{1.005401in}{0.953701in}}%
\pgfpathlineto{\pgfqpoint{1.072506in}{0.937190in}}%
\pgfpathlineto{\pgfqpoint{1.072364in}{0.980004in}}%
\pgfpathlineto{\pgfqpoint{1.005101in}{0.996336in}}%
\pgfpathlineto{\pgfqpoint{1.005401in}{0.953701in}}%
\pgfpathclose%
\pgfusepath{fill}%
\end{pgfscope}%
\begin{pgfscope}%
\pgfpathrectangle{\pgfqpoint{0.407592in}{0.222904in}}{\pgfqpoint{1.411163in}{1.411163in}}%
\pgfusepath{clip}%
\pgfsetbuttcap%
\pgfsetroundjoin%
\definecolor{currentfill}{rgb}{0.141176,0.345098,0.650980}%
\pgfsetfillcolor{currentfill}%
\pgfsetfillopacity{0.080000}%
\pgfsetlinewidth{0.000000pt}%
\definecolor{currentstroke}{rgb}{0.000000,0.000000,0.000000}%
\pgfsetstrokecolor{currentstroke}%
\pgfsetdash{}{0pt}%
\pgfpathmoveto{\pgfqpoint{1.072789in}{0.852170in}}%
\pgfpathlineto{\pgfqpoint{1.140262in}{0.835135in}}%
\pgfpathlineto{\pgfqpoint{1.140281in}{0.877720in}}%
\pgfpathlineto{\pgfqpoint{1.072648in}{0.894579in}}%
\pgfpathlineto{\pgfqpoint{1.072789in}{0.852170in}}%
\pgfpathclose%
\pgfusepath{fill}%
\end{pgfscope}%
\begin{pgfscope}%
\pgfpathrectangle{\pgfqpoint{0.407592in}{0.222904in}}{\pgfqpoint{1.411163in}{1.411163in}}%
\pgfusepath{clip}%
\pgfsetbuttcap%
\pgfsetroundjoin%
\definecolor{currentfill}{rgb}{0.141176,0.345098,0.650980}%
\pgfsetfillcolor{currentfill}%
\pgfsetfillopacity{0.080000}%
\pgfsetlinewidth{0.000000pt}%
\definecolor{currentstroke}{rgb}{0.000000,0.000000,0.000000}%
\pgfsetstrokecolor{currentstroke}%
\pgfsetdash{}{0pt}%
\pgfpathmoveto{\pgfqpoint{0.937601in}{1.098021in}}%
\pgfpathlineto{\pgfqpoint{1.004495in}{1.082216in}}%
\pgfpathlineto{\pgfqpoint{1.004189in}{1.125463in}}%
\pgfpathlineto{\pgfqpoint{0.937138in}{1.141086in}}%
\pgfpathlineto{\pgfqpoint{0.937601in}{1.098021in}}%
\pgfpathclose%
\pgfusepath{fill}%
\end{pgfscope}%
\begin{pgfscope}%
\pgfpathrectangle{\pgfqpoint{0.407592in}{0.222904in}}{\pgfqpoint{1.411163in}{1.411163in}}%
\pgfusepath{clip}%
\pgfsetbuttcap%
\pgfsetroundjoin%
\definecolor{currentfill}{rgb}{0.141176,0.345098,0.650980}%
\pgfsetfillcolor{currentfill}%
\pgfsetfillopacity{0.080000}%
\pgfsetlinewidth{0.000000pt}%
\definecolor{currentstroke}{rgb}{0.000000,0.000000,0.000000}%
\pgfsetstrokecolor{currentstroke}%
\pgfsetdash{}{0pt}%
\pgfpathmoveto{\pgfqpoint{1.005101in}{0.996336in}}%
\pgfpathlineto{\pgfqpoint{1.072364in}{0.980004in}}%
\pgfpathlineto{\pgfqpoint{1.072221in}{1.023023in}}%
\pgfpathlineto{\pgfqpoint{1.004798in}{1.039174in}}%
\pgfpathlineto{\pgfqpoint{1.005101in}{0.996336in}}%
\pgfpathclose%
\pgfusepath{fill}%
\end{pgfscope}%
\begin{pgfscope}%
\pgfpathrectangle{\pgfqpoint{0.407592in}{0.222904in}}{\pgfqpoint{1.411163in}{1.411163in}}%
\pgfusepath{clip}%
\pgfsetbuttcap%
\pgfsetroundjoin%
\definecolor{currentfill}{rgb}{0.141176,0.345098,0.650980}%
\pgfsetfillcolor{currentfill}%
\pgfsetfillopacity{0.080000}%
\pgfsetlinewidth{0.000000pt}%
\definecolor{currentstroke}{rgb}{0.000000,0.000000,0.000000}%
\pgfsetstrokecolor{currentstroke}%
\pgfsetdash{}{0pt}%
\pgfpathmoveto{\pgfqpoint{1.140243in}{0.792751in}}%
\pgfpathlineto{\pgfqpoint{1.208428in}{0.817924in}}%
\pgfpathlineto{\pgfqpoint{1.140262in}{0.835135in}}%
\pgfusepath{fill}%
\end{pgfscope}%
\begin{pgfscope}%
\pgfpathrectangle{\pgfqpoint{0.407592in}{0.222904in}}{\pgfqpoint{1.411163in}{1.411163in}}%
\pgfusepath{clip}%
\pgfsetbuttcap%
\pgfsetroundjoin%
\definecolor{currentfill}{rgb}{0.141176,0.345098,0.650980}%
\pgfsetfillcolor{currentfill}%
\pgfsetfillopacity{0.080000}%
\pgfsetlinewidth{0.000000pt}%
\definecolor{currentstroke}{rgb}{0.000000,0.000000,0.000000}%
\pgfsetstrokecolor{currentstroke}%
\pgfsetdash{}{0pt}%
\pgfpathmoveto{\pgfqpoint{1.072648in}{0.894579in}}%
\pgfpathlineto{\pgfqpoint{1.140281in}{0.877720in}}%
\pgfpathlineto{\pgfqpoint{1.140300in}{0.920510in}}%
\pgfpathlineto{\pgfqpoint{1.072506in}{0.937190in}}%
\pgfpathlineto{\pgfqpoint{1.072648in}{0.894579in}}%
\pgfpathclose%
\pgfusepath{fill}%
\end{pgfscope}%
\begin{pgfscope}%
\pgfpathrectangle{\pgfqpoint{0.407592in}{0.222904in}}{\pgfqpoint{1.411163in}{1.411163in}}%
\pgfusepath{clip}%
\pgfsetbuttcap%
\pgfsetroundjoin%
\definecolor{currentfill}{rgb}{0.141176,0.345098,0.650980}%
\pgfsetfillcolor{currentfill}%
\pgfsetfillopacity{0.080000}%
\pgfsetlinewidth{0.000000pt}%
\definecolor{currentstroke}{rgb}{0.000000,0.000000,0.000000}%
\pgfsetstrokecolor{currentstroke}%
\pgfsetdash{}{0pt}%
\pgfpathmoveto{\pgfqpoint{0.937138in}{1.141086in}}%
\pgfpathlineto{\pgfqpoint{1.004189in}{1.125463in}}%
\pgfpathlineto{\pgfqpoint{1.003883in}{1.168918in}}%
\pgfpathlineto{\pgfqpoint{0.936673in}{1.184355in}}%
\pgfpathlineto{\pgfqpoint{0.937138in}{1.141086in}}%
\pgfpathclose%
\pgfusepath{fill}%
\end{pgfscope}%
\begin{pgfscope}%
\pgfpathrectangle{\pgfqpoint{0.407592in}{0.222904in}}{\pgfqpoint{1.411163in}{1.411163in}}%
\pgfusepath{clip}%
\pgfsetbuttcap%
\pgfsetroundjoin%
\definecolor{currentfill}{rgb}{0.141176,0.345098,0.650980}%
\pgfsetfillcolor{currentfill}%
\pgfsetfillopacity{0.080000}%
\pgfsetlinewidth{0.000000pt}%
\definecolor{currentstroke}{rgb}{0.000000,0.000000,0.000000}%
\pgfsetstrokecolor{currentstroke}%
\pgfsetdash{}{0pt}%
\pgfpathmoveto{\pgfqpoint{1.004798in}{1.039174in}}%
\pgfpathlineto{\pgfqpoint{1.072221in}{1.023023in}}%
\pgfpathlineto{\pgfqpoint{1.072077in}{1.066248in}}%
\pgfpathlineto{\pgfqpoint{1.004495in}{1.082216in}}%
\pgfpathlineto{\pgfqpoint{1.004798in}{1.039174in}}%
\pgfpathclose%
\pgfusepath{fill}%
\end{pgfscope}%
\begin{pgfscope}%
\pgfpathrectangle{\pgfqpoint{0.407592in}{0.222904in}}{\pgfqpoint{1.411163in}{1.411163in}}%
\pgfusepath{clip}%
\pgfsetbuttcap%
\pgfsetroundjoin%
\definecolor{currentfill}{rgb}{0.141176,0.345098,0.650980}%
\pgfsetfillcolor{currentfill}%
\pgfsetfillopacity{0.080000}%
\pgfsetlinewidth{0.000000pt}%
\definecolor{currentstroke}{rgb}{0.000000,0.000000,0.000000}%
\pgfsetstrokecolor{currentstroke}%
\pgfsetdash{}{0pt}%
\pgfpathmoveto{\pgfqpoint{1.072506in}{0.937190in}}%
\pgfpathlineto{\pgfqpoint{1.140300in}{0.920510in}}%
\pgfpathlineto{\pgfqpoint{1.140319in}{0.963504in}}%
\pgfpathlineto{\pgfqpoint{1.072364in}{0.980004in}}%
\pgfpathlineto{\pgfqpoint{1.072506in}{0.937190in}}%
\pgfpathclose%
\pgfusepath{fill}%
\end{pgfscope}%
\begin{pgfscope}%
\pgfpathrectangle{\pgfqpoint{0.407592in}{0.222904in}}{\pgfqpoint{1.411163in}{1.411163in}}%
\pgfusepath{clip}%
\pgfsetbuttcap%
\pgfsetroundjoin%
\definecolor{currentfill}{rgb}{0.141176,0.345098,0.650980}%
\pgfsetfillcolor{currentfill}%
\pgfsetfillopacity{0.080000}%
\pgfsetlinewidth{0.000000pt}%
\definecolor{currentstroke}{rgb}{0.000000,0.000000,0.000000}%
\pgfsetstrokecolor{currentstroke}%
\pgfsetdash{}{0pt}%
\pgfpathmoveto{\pgfqpoint{1.140262in}{0.835135in}}%
\pgfpathlineto{\pgfqpoint{1.208428in}{0.817924in}}%
\pgfpathlineto{\pgfqpoint{1.208610in}{0.860688in}}%
\pgfpathlineto{\pgfqpoint{1.140281in}{0.877720in}}%
\pgfpathlineto{\pgfqpoint{1.140262in}{0.835135in}}%
\pgfpathclose%
\pgfusepath{fill}%
\end{pgfscope}%
\begin{pgfscope}%
\pgfpathrectangle{\pgfqpoint{0.407592in}{0.222904in}}{\pgfqpoint{1.411163in}{1.411163in}}%
\pgfusepath{clip}%
\pgfsetbuttcap%
\pgfsetroundjoin%
\definecolor{currentfill}{rgb}{0.141176,0.345098,0.650980}%
\pgfsetfillcolor{currentfill}%
\pgfsetfillopacity{0.080000}%
\pgfsetlinewidth{0.000000pt}%
\definecolor{currentstroke}{rgb}{0.000000,0.000000,0.000000}%
\pgfsetstrokecolor{currentstroke}%
\pgfsetdash{}{0pt}%
\pgfpathmoveto{\pgfqpoint{1.208428in}{0.817924in}}%
\pgfusepath{fill}%
\end{pgfscope}%
\begin{pgfscope}%
\pgfpathrectangle{\pgfqpoint{0.407592in}{0.222904in}}{\pgfqpoint{1.411163in}{1.411163in}}%
\pgfusepath{clip}%
\pgfsetbuttcap%
\pgfsetroundjoin%
\definecolor{currentfill}{rgb}{0.141176,0.345098,0.650980}%
\pgfsetfillcolor{currentfill}%
\pgfsetfillopacity{0.080000}%
\pgfsetlinewidth{0.000000pt}%
\definecolor{currentstroke}{rgb}{0.000000,0.000000,0.000000}%
\pgfsetstrokecolor{currentstroke}%
\pgfsetdash{}{0pt}%
\pgfpathmoveto{\pgfqpoint{1.004495in}{1.082216in}}%
\pgfpathlineto{\pgfqpoint{1.072077in}{1.066248in}}%
\pgfpathlineto{\pgfqpoint{1.071932in}{1.109680in}}%
\pgfpathlineto{\pgfqpoint{1.004189in}{1.125463in}}%
\pgfpathlineto{\pgfqpoint{1.004495in}{1.082216in}}%
\pgfpathclose%
\pgfusepath{fill}%
\end{pgfscope}%
\begin{pgfscope}%
\pgfpathrectangle{\pgfqpoint{0.407592in}{0.222904in}}{\pgfqpoint{1.411163in}{1.411163in}}%
\pgfusepath{clip}%
\pgfsetbuttcap%
\pgfsetroundjoin%
\definecolor{currentfill}{rgb}{0.141176,0.345098,0.650980}%
\pgfsetfillcolor{currentfill}%
\pgfsetfillopacity{0.080000}%
\pgfsetlinewidth{0.000000pt}%
\definecolor{currentstroke}{rgb}{0.000000,0.000000,0.000000}%
\pgfsetstrokecolor{currentstroke}%
\pgfsetdash{}{0pt}%
\pgfpathmoveto{\pgfqpoint{1.072364in}{0.980004in}}%
\pgfpathlineto{\pgfqpoint{1.140319in}{0.963504in}}%
\pgfpathlineto{\pgfqpoint{1.140339in}{1.006705in}}%
\pgfpathlineto{\pgfqpoint{1.072221in}{1.023023in}}%
\pgfpathlineto{\pgfqpoint{1.072364in}{0.980004in}}%
\pgfpathclose%
\pgfusepath{fill}%
\end{pgfscope}%
\begin{pgfscope}%
\pgfpathrectangle{\pgfqpoint{0.407592in}{0.222904in}}{\pgfqpoint{1.411163in}{1.411163in}}%
\pgfusepath{clip}%
\pgfsetbuttcap%
\pgfsetroundjoin%
\definecolor{currentfill}{rgb}{0.141176,0.345098,0.650980}%
\pgfsetfillcolor{currentfill}%
\pgfsetfillopacity{0.080000}%
\pgfsetlinewidth{0.000000pt}%
\definecolor{currentstroke}{rgb}{0.000000,0.000000,0.000000}%
\pgfsetstrokecolor{currentstroke}%
\pgfsetdash{}{0pt}%
\pgfpathmoveto{\pgfqpoint{1.140281in}{0.877720in}}%
\pgfpathlineto{\pgfqpoint{1.208610in}{0.860688in}}%
\pgfpathlineto{\pgfqpoint{1.208794in}{0.903657in}}%
\pgfpathlineto{\pgfqpoint{1.140300in}{0.920510in}}%
\pgfpathlineto{\pgfqpoint{1.140281in}{0.877720in}}%
\pgfpathclose%
\pgfusepath{fill}%
\end{pgfscope}%
\begin{pgfscope}%
\pgfpathrectangle{\pgfqpoint{0.407592in}{0.222904in}}{\pgfqpoint{1.411163in}{1.411163in}}%
\pgfusepath{clip}%
\pgfsetbuttcap%
\pgfsetroundjoin%
\definecolor{currentfill}{rgb}{0.141176,0.345098,0.650980}%
\pgfsetfillcolor{currentfill}%
\pgfsetfillopacity{0.080000}%
\pgfsetlinewidth{0.000000pt}%
\definecolor{currentstroke}{rgb}{0.000000,0.000000,0.000000}%
\pgfsetstrokecolor{currentstroke}%
\pgfsetdash{}{0pt}%
\pgfpathmoveto{\pgfqpoint{1.208428in}{0.817924in}}%
\pgfpathlineto{\pgfqpoint{1.208610in}{0.860688in}}%
\pgfusepath{fill}%
\end{pgfscope}%
\begin{pgfscope}%
\pgfpathrectangle{\pgfqpoint{0.407592in}{0.222904in}}{\pgfqpoint{1.411163in}{1.411163in}}%
\pgfusepath{clip}%
\pgfsetbuttcap%
\pgfsetroundjoin%
\definecolor{currentfill}{rgb}{0.141176,0.345098,0.650980}%
\pgfsetfillcolor{currentfill}%
\pgfsetfillopacity{0.080000}%
\pgfsetlinewidth{0.000000pt}%
\definecolor{currentstroke}{rgb}{0.000000,0.000000,0.000000}%
\pgfsetstrokecolor{currentstroke}%
\pgfsetdash{}{0pt}%
\pgfpathmoveto{\pgfqpoint{1.004189in}{1.125463in}}%
\pgfpathlineto{\pgfqpoint{1.071932in}{1.109680in}}%
\pgfpathlineto{\pgfqpoint{1.071787in}{1.153321in}}%
\pgfpathlineto{\pgfqpoint{1.003883in}{1.168918in}}%
\pgfpathlineto{\pgfqpoint{1.004189in}{1.125463in}}%
\pgfpathclose%
\pgfusepath{fill}%
\end{pgfscope}%
\begin{pgfscope}%
\pgfpathrectangle{\pgfqpoint{0.407592in}{0.222904in}}{\pgfqpoint{1.411163in}{1.411163in}}%
\pgfusepath{clip}%
\pgfsetbuttcap%
\pgfsetroundjoin%
\definecolor{currentfill}{rgb}{0.141176,0.345098,0.650980}%
\pgfsetfillcolor{currentfill}%
\pgfsetfillopacity{0.080000}%
\pgfsetlinewidth{0.000000pt}%
\definecolor{currentstroke}{rgb}{0.000000,0.000000,0.000000}%
\pgfsetstrokecolor{currentstroke}%
\pgfsetdash{}{0pt}%
\pgfpathmoveto{\pgfqpoint{1.072221in}{1.023023in}}%
\pgfpathlineto{\pgfqpoint{1.140339in}{1.006705in}}%
\pgfpathlineto{\pgfqpoint{1.140358in}{1.050114in}}%
\pgfpathlineto{\pgfqpoint{1.072077in}{1.066248in}}%
\pgfpathlineto{\pgfqpoint{1.072221in}{1.023023in}}%
\pgfpathclose%
\pgfusepath{fill}%
\end{pgfscope}%
\begin{pgfscope}%
\pgfpathrectangle{\pgfqpoint{0.407592in}{0.222904in}}{\pgfqpoint{1.411163in}{1.411163in}}%
\pgfusepath{clip}%
\pgfsetbuttcap%
\pgfsetroundjoin%
\definecolor{currentfill}{rgb}{0.141176,0.345098,0.650980}%
\pgfsetfillcolor{currentfill}%
\pgfsetfillopacity{0.080000}%
\pgfsetlinewidth{0.000000pt}%
\definecolor{currentstroke}{rgb}{0.000000,0.000000,0.000000}%
\pgfsetstrokecolor{currentstroke}%
\pgfsetdash{}{0pt}%
\pgfpathmoveto{\pgfqpoint{1.140300in}{0.920510in}}%
\pgfpathlineto{\pgfqpoint{1.208794in}{0.903657in}}%
\pgfpathlineto{\pgfqpoint{1.208978in}{0.946833in}}%
\pgfpathlineto{\pgfqpoint{1.140319in}{0.963504in}}%
\pgfpathlineto{\pgfqpoint{1.140300in}{0.920510in}}%
\pgfpathclose%
\pgfusepath{fill}%
\end{pgfscope}%
\begin{pgfscope}%
\pgfpathrectangle{\pgfqpoint{0.407592in}{0.222904in}}{\pgfqpoint{1.411163in}{1.411163in}}%
\pgfusepath{clip}%
\pgfsetbuttcap%
\pgfsetroundjoin%
\definecolor{currentfill}{rgb}{0.141176,0.345098,0.650980}%
\pgfsetfillcolor{currentfill}%
\pgfsetfillopacity{0.080000}%
\pgfsetlinewidth{0.000000pt}%
\definecolor{currentstroke}{rgb}{0.000000,0.000000,0.000000}%
\pgfsetstrokecolor{currentstroke}%
\pgfsetdash{}{0pt}%
\pgfpathmoveto{\pgfqpoint{1.072077in}{1.066248in}}%
\pgfpathlineto{\pgfqpoint{1.140358in}{1.050114in}}%
\pgfpathlineto{\pgfqpoint{1.140378in}{1.093733in}}%
\pgfpathlineto{\pgfqpoint{1.071932in}{1.109680in}}%
\pgfpathlineto{\pgfqpoint{1.072077in}{1.066248in}}%
\pgfpathclose%
\pgfusepath{fill}%
\end{pgfscope}%
\begin{pgfscope}%
\pgfpathrectangle{\pgfqpoint{0.407592in}{0.222904in}}{\pgfqpoint{1.411163in}{1.411163in}}%
\pgfusepath{clip}%
\pgfsetbuttcap%
\pgfsetroundjoin%
\definecolor{currentfill}{rgb}{0.141176,0.345098,0.650980}%
\pgfsetfillcolor{currentfill}%
\pgfsetfillopacity{0.080000}%
\pgfsetlinewidth{0.000000pt}%
\definecolor{currentstroke}{rgb}{0.000000,0.000000,0.000000}%
\pgfsetstrokecolor{currentstroke}%
\pgfsetdash{}{0pt}%
\pgfpathmoveto{\pgfqpoint{1.208610in}{0.860688in}}%
\pgfpathlineto{\pgfqpoint{1.277998in}{0.886630in}}%
\pgfpathlineto{\pgfqpoint{1.208794in}{0.903657in}}%
\pgfusepath{fill}%
\end{pgfscope}%
\begin{pgfscope}%
\pgfpathrectangle{\pgfqpoint{0.407592in}{0.222904in}}{\pgfqpoint{1.411163in}{1.411163in}}%
\pgfusepath{clip}%
\pgfsetbuttcap%
\pgfsetroundjoin%
\definecolor{currentfill}{rgb}{0.141176,0.345098,0.650980}%
\pgfsetfillcolor{currentfill}%
\pgfsetfillopacity{0.080000}%
\pgfsetlinewidth{0.000000pt}%
\definecolor{currentstroke}{rgb}{0.000000,0.000000,0.000000}%
\pgfsetstrokecolor{currentstroke}%
\pgfsetdash{}{0pt}%
\pgfpathmoveto{\pgfqpoint{1.140319in}{0.963504in}}%
\pgfpathlineto{\pgfqpoint{1.208978in}{0.946833in}}%
\pgfpathlineto{\pgfqpoint{1.209163in}{0.990218in}}%
\pgfpathlineto{\pgfqpoint{1.140339in}{1.006705in}}%
\pgfpathlineto{\pgfqpoint{1.140319in}{0.963504in}}%
\pgfpathclose%
\pgfusepath{fill}%
\end{pgfscope}%
\begin{pgfscope}%
\pgfpathrectangle{\pgfqpoint{0.407592in}{0.222904in}}{\pgfqpoint{1.411163in}{1.411163in}}%
\pgfusepath{clip}%
\pgfsetbuttcap%
\pgfsetroundjoin%
\definecolor{currentfill}{rgb}{0.141176,0.345098,0.650980}%
\pgfsetfillcolor{currentfill}%
\pgfsetfillopacity{0.080000}%
\pgfsetlinewidth{0.000000pt}%
\definecolor{currentstroke}{rgb}{0.000000,0.000000,0.000000}%
\pgfsetstrokecolor{currentstroke}%
\pgfsetdash{}{0pt}%
\pgfpathmoveto{\pgfqpoint{1.071932in}{1.109680in}}%
\pgfpathlineto{\pgfqpoint{1.140378in}{1.093733in}}%
\pgfpathlineto{\pgfqpoint{1.140398in}{1.137562in}}%
\pgfpathlineto{\pgfqpoint{1.071787in}{1.153321in}}%
\pgfpathlineto{\pgfqpoint{1.071932in}{1.109680in}}%
\pgfpathclose%
\pgfusepath{fill}%
\end{pgfscope}%
\begin{pgfscope}%
\pgfpathrectangle{\pgfqpoint{0.407592in}{0.222904in}}{\pgfqpoint{1.411163in}{1.411163in}}%
\pgfusepath{clip}%
\pgfsetbuttcap%
\pgfsetroundjoin%
\definecolor{currentfill}{rgb}{0.141176,0.345098,0.650980}%
\pgfsetfillcolor{currentfill}%
\pgfsetfillopacity{0.080000}%
\pgfsetlinewidth{0.000000pt}%
\definecolor{currentstroke}{rgb}{0.000000,0.000000,0.000000}%
\pgfsetstrokecolor{currentstroke}%
\pgfsetdash{}{0pt}%
\pgfpathmoveto{\pgfqpoint{1.140339in}{1.006705in}}%
\pgfpathlineto{\pgfqpoint{1.209163in}{0.990218in}}%
\pgfpathlineto{\pgfqpoint{1.209350in}{1.033813in}}%
\pgfpathlineto{\pgfqpoint{1.140358in}{1.050114in}}%
\pgfpathlineto{\pgfqpoint{1.140339in}{1.006705in}}%
\pgfpathclose%
\pgfusepath{fill}%
\end{pgfscope}%
\begin{pgfscope}%
\pgfpathrectangle{\pgfqpoint{0.407592in}{0.222904in}}{\pgfqpoint{1.411163in}{1.411163in}}%
\pgfusepath{clip}%
\pgfsetbuttcap%
\pgfsetroundjoin%
\definecolor{currentfill}{rgb}{0.141176,0.345098,0.650980}%
\pgfsetfillcolor{currentfill}%
\pgfsetfillopacity{0.080000}%
\pgfsetlinewidth{0.000000pt}%
\definecolor{currentstroke}{rgb}{0.000000,0.000000,0.000000}%
\pgfsetstrokecolor{currentstroke}%
\pgfsetdash{}{0pt}%
\pgfpathmoveto{\pgfqpoint{1.208794in}{0.903657in}}%
\pgfpathlineto{\pgfqpoint{1.277998in}{0.886630in}}%
\pgfpathlineto{\pgfqpoint{1.278351in}{0.929989in}}%
\pgfpathlineto{\pgfqpoint{1.208978in}{0.946833in}}%
\pgfpathlineto{\pgfqpoint{1.208794in}{0.903657in}}%
\pgfpathclose%
\pgfusepath{fill}%
\end{pgfscope}%
\begin{pgfscope}%
\pgfpathrectangle{\pgfqpoint{0.407592in}{0.222904in}}{\pgfqpoint{1.411163in}{1.411163in}}%
\pgfusepath{clip}%
\pgfsetbuttcap%
\pgfsetroundjoin%
\definecolor{currentfill}{rgb}{0.141176,0.345098,0.650980}%
\pgfsetfillcolor{currentfill}%
\pgfsetfillopacity{0.080000}%
\pgfsetlinewidth{0.000000pt}%
\definecolor{currentstroke}{rgb}{0.000000,0.000000,0.000000}%
\pgfsetstrokecolor{currentstroke}%
\pgfsetdash{}{0pt}%
\pgfpathmoveto{\pgfqpoint{1.277998in}{0.886630in}}%
\pgfusepath{fill}%
\end{pgfscope}%
\begin{pgfscope}%
\pgfpathrectangle{\pgfqpoint{0.407592in}{0.222904in}}{\pgfqpoint{1.411163in}{1.411163in}}%
\pgfusepath{clip}%
\pgfsetbuttcap%
\pgfsetroundjoin%
\definecolor{currentfill}{rgb}{0.141176,0.345098,0.650980}%
\pgfsetfillcolor{currentfill}%
\pgfsetfillopacity{0.080000}%
\pgfsetlinewidth{0.000000pt}%
\definecolor{currentstroke}{rgb}{0.000000,0.000000,0.000000}%
\pgfsetstrokecolor{currentstroke}%
\pgfsetdash{}{0pt}%
\pgfpathmoveto{\pgfqpoint{1.140358in}{1.050114in}}%
\pgfpathlineto{\pgfqpoint{1.209350in}{1.033813in}}%
\pgfpathlineto{\pgfqpoint{1.209537in}{1.077619in}}%
\pgfpathlineto{\pgfqpoint{1.140378in}{1.093733in}}%
\pgfpathlineto{\pgfqpoint{1.140358in}{1.050114in}}%
\pgfpathclose%
\pgfusepath{fill}%
\end{pgfscope}%
\begin{pgfscope}%
\pgfpathrectangle{\pgfqpoint{0.407592in}{0.222904in}}{\pgfqpoint{1.411163in}{1.411163in}}%
\pgfusepath{clip}%
\pgfsetbuttcap%
\pgfsetroundjoin%
\definecolor{currentfill}{rgb}{0.141176,0.345098,0.650980}%
\pgfsetfillcolor{currentfill}%
\pgfsetfillopacity{0.080000}%
\pgfsetlinewidth{0.000000pt}%
\definecolor{currentstroke}{rgb}{0.000000,0.000000,0.000000}%
\pgfsetstrokecolor{currentstroke}%
\pgfsetdash{}{0pt}%
\pgfpathmoveto{\pgfqpoint{1.208978in}{0.946833in}}%
\pgfpathlineto{\pgfqpoint{1.278351in}{0.929989in}}%
\pgfpathlineto{\pgfqpoint{1.278706in}{0.973559in}}%
\pgfpathlineto{\pgfqpoint{1.209163in}{0.990218in}}%
\pgfpathlineto{\pgfqpoint{1.208978in}{0.946833in}}%
\pgfpathclose%
\pgfusepath{fill}%
\end{pgfscope}%
\begin{pgfscope}%
\pgfpathrectangle{\pgfqpoint{0.407592in}{0.222904in}}{\pgfqpoint{1.411163in}{1.411163in}}%
\pgfusepath{clip}%
\pgfsetbuttcap%
\pgfsetroundjoin%
\definecolor{currentfill}{rgb}{0.141176,0.345098,0.650980}%
\pgfsetfillcolor{currentfill}%
\pgfsetfillopacity{0.080000}%
\pgfsetlinewidth{0.000000pt}%
\definecolor{currentstroke}{rgb}{0.000000,0.000000,0.000000}%
\pgfsetstrokecolor{currentstroke}%
\pgfsetdash{}{0pt}%
\pgfpathmoveto{\pgfqpoint{1.140378in}{1.093733in}}%
\pgfpathlineto{\pgfqpoint{1.209537in}{1.077619in}}%
\pgfpathlineto{\pgfqpoint{1.209725in}{1.121639in}}%
\pgfpathlineto{\pgfqpoint{1.140398in}{1.137562in}}%
\pgfpathlineto{\pgfqpoint{1.140378in}{1.093733in}}%
\pgfpathclose%
\pgfusepath{fill}%
\end{pgfscope}%
\begin{pgfscope}%
\pgfpathrectangle{\pgfqpoint{0.407592in}{0.222904in}}{\pgfqpoint{1.411163in}{1.411163in}}%
\pgfusepath{clip}%
\pgfsetbuttcap%
\pgfsetroundjoin%
\definecolor{currentfill}{rgb}{0.141176,0.345098,0.650980}%
\pgfsetfillcolor{currentfill}%
\pgfsetfillopacity{0.080000}%
\pgfsetlinewidth{0.000000pt}%
\definecolor{currentstroke}{rgb}{0.000000,0.000000,0.000000}%
\pgfsetstrokecolor{currentstroke}%
\pgfsetdash{}{0pt}%
\pgfpathmoveto{\pgfqpoint{1.277998in}{0.886630in}}%
\pgfpathlineto{\pgfqpoint{1.348450in}{0.912969in}}%
\pgfpathlineto{\pgfqpoint{1.278351in}{0.929989in}}%
\pgfusepath{fill}%
\end{pgfscope}%
\begin{pgfscope}%
\pgfpathrectangle{\pgfqpoint{0.407592in}{0.222904in}}{\pgfqpoint{1.411163in}{1.411163in}}%
\pgfusepath{clip}%
\pgfsetbuttcap%
\pgfsetroundjoin%
\definecolor{currentfill}{rgb}{0.141176,0.345098,0.650980}%
\pgfsetfillcolor{currentfill}%
\pgfsetfillopacity{0.080000}%
\pgfsetlinewidth{0.000000pt}%
\definecolor{currentstroke}{rgb}{0.000000,0.000000,0.000000}%
\pgfsetstrokecolor{currentstroke}%
\pgfsetdash{}{0pt}%
\pgfpathmoveto{\pgfqpoint{1.209163in}{0.990218in}}%
\pgfpathlineto{\pgfqpoint{1.278706in}{0.973559in}}%
\pgfpathlineto{\pgfqpoint{1.279062in}{1.017342in}}%
\pgfpathlineto{\pgfqpoint{1.209350in}{1.033813in}}%
\pgfpathlineto{\pgfqpoint{1.209163in}{0.990218in}}%
\pgfpathclose%
\pgfusepath{fill}%
\end{pgfscope}%
\begin{pgfscope}%
\pgfpathrectangle{\pgfqpoint{0.407592in}{0.222904in}}{\pgfqpoint{1.411163in}{1.411163in}}%
\pgfusepath{clip}%
\pgfsetbuttcap%
\pgfsetroundjoin%
\definecolor{currentfill}{rgb}{0.141176,0.345098,0.650980}%
\pgfsetfillcolor{currentfill}%
\pgfsetfillopacity{0.080000}%
\pgfsetlinewidth{0.000000pt}%
\definecolor{currentstroke}{rgb}{0.000000,0.000000,0.000000}%
\pgfsetstrokecolor{currentstroke}%
\pgfsetdash{}{0pt}%
\pgfpathmoveto{\pgfqpoint{1.209350in}{1.033813in}}%
\pgfpathlineto{\pgfqpoint{1.279062in}{1.017342in}}%
\pgfpathlineto{\pgfqpoint{1.279420in}{1.061337in}}%
\pgfpathlineto{\pgfqpoint{1.209537in}{1.077619in}}%
\pgfpathlineto{\pgfqpoint{1.209350in}{1.033813in}}%
\pgfpathclose%
\pgfusepath{fill}%
\end{pgfscope}%
\begin{pgfscope}%
\pgfpathrectangle{\pgfqpoint{0.407592in}{0.222904in}}{\pgfqpoint{1.411163in}{1.411163in}}%
\pgfusepath{clip}%
\pgfsetbuttcap%
\pgfsetroundjoin%
\definecolor{currentfill}{rgb}{0.141176,0.345098,0.650980}%
\pgfsetfillcolor{currentfill}%
\pgfsetfillopacity{0.080000}%
\pgfsetlinewidth{0.000000pt}%
\definecolor{currentstroke}{rgb}{0.000000,0.000000,0.000000}%
\pgfsetstrokecolor{currentstroke}%
\pgfsetdash{}{0pt}%
\pgfpathmoveto{\pgfqpoint{1.278351in}{0.929989in}}%
\pgfpathlineto{\pgfqpoint{1.348450in}{0.912969in}}%
\pgfpathlineto{\pgfqpoint{1.348977in}{0.956726in}}%
\pgfpathlineto{\pgfqpoint{1.278706in}{0.973559in}}%
\pgfpathlineto{\pgfqpoint{1.278351in}{0.929989in}}%
\pgfpathclose%
\pgfusepath{fill}%
\end{pgfscope}%
\begin{pgfscope}%
\pgfpathrectangle{\pgfqpoint{0.407592in}{0.222904in}}{\pgfqpoint{1.411163in}{1.411163in}}%
\pgfusepath{clip}%
\pgfsetbuttcap%
\pgfsetroundjoin%
\definecolor{currentfill}{rgb}{0.141176,0.345098,0.650980}%
\pgfsetfillcolor{currentfill}%
\pgfsetfillopacity{0.080000}%
\pgfsetlinewidth{0.000000pt}%
\definecolor{currentstroke}{rgb}{0.000000,0.000000,0.000000}%
\pgfsetstrokecolor{currentstroke}%
\pgfsetdash{}{0pt}%
\pgfpathmoveto{\pgfqpoint{1.348450in}{0.912969in}}%
\pgfusepath{fill}%
\end{pgfscope}%
\begin{pgfscope}%
\pgfpathrectangle{\pgfqpoint{0.407592in}{0.222904in}}{\pgfqpoint{1.411163in}{1.411163in}}%
\pgfusepath{clip}%
\pgfsetbuttcap%
\pgfsetroundjoin%
\definecolor{currentfill}{rgb}{0.141176,0.345098,0.650980}%
\pgfsetfillcolor{currentfill}%
\pgfsetfillopacity{0.080000}%
\pgfsetlinewidth{0.000000pt}%
\definecolor{currentstroke}{rgb}{0.000000,0.000000,0.000000}%
\pgfsetstrokecolor{currentstroke}%
\pgfsetdash{}{0pt}%
\pgfpathmoveto{\pgfqpoint{1.209537in}{1.077619in}}%
\pgfpathlineto{\pgfqpoint{1.279420in}{1.061337in}}%
\pgfpathlineto{\pgfqpoint{1.279780in}{1.105548in}}%
\pgfpathlineto{\pgfqpoint{1.209725in}{1.121639in}}%
\pgfpathlineto{\pgfqpoint{1.209537in}{1.077619in}}%
\pgfpathclose%
\pgfusepath{fill}%
\end{pgfscope}%
\begin{pgfscope}%
\pgfpathrectangle{\pgfqpoint{0.407592in}{0.222904in}}{\pgfqpoint{1.411163in}{1.411163in}}%
\pgfusepath{clip}%
\pgfsetbuttcap%
\pgfsetroundjoin%
\definecolor{currentfill}{rgb}{0.141176,0.345098,0.650980}%
\pgfsetfillcolor{currentfill}%
\pgfsetfillopacity{0.080000}%
\pgfsetlinewidth{0.000000pt}%
\definecolor{currentstroke}{rgb}{0.000000,0.000000,0.000000}%
\pgfsetstrokecolor{currentstroke}%
\pgfsetdash{}{0pt}%
\pgfpathmoveto{\pgfqpoint{1.278706in}{0.973559in}}%
\pgfpathlineto{\pgfqpoint{1.348977in}{0.956726in}}%
\pgfpathlineto{\pgfqpoint{1.349507in}{1.000697in}}%
\pgfpathlineto{\pgfqpoint{1.279062in}{1.017342in}}%
\pgfpathlineto{\pgfqpoint{1.278706in}{0.973559in}}%
\pgfpathclose%
\pgfusepath{fill}%
\end{pgfscope}%
\begin{pgfscope}%
\pgfpathrectangle{\pgfqpoint{0.407592in}{0.222904in}}{\pgfqpoint{1.411163in}{1.411163in}}%
\pgfusepath{clip}%
\pgfsetbuttcap%
\pgfsetroundjoin%
\definecolor{currentfill}{rgb}{0.141176,0.345098,0.650980}%
\pgfsetfillcolor{currentfill}%
\pgfsetfillopacity{0.080000}%
\pgfsetlinewidth{0.000000pt}%
\definecolor{currentstroke}{rgb}{0.000000,0.000000,0.000000}%
\pgfsetstrokecolor{currentstroke}%
\pgfsetdash{}{0pt}%
\pgfpathmoveto{\pgfqpoint{1.348450in}{0.912969in}}%
\pgfpathlineto{\pgfqpoint{1.419989in}{0.939715in}}%
\pgfpathlineto{\pgfqpoint{1.348977in}{0.956726in}}%
\pgfusepath{fill}%
\end{pgfscope}%
\begin{pgfscope}%
\pgfpathrectangle{\pgfqpoint{0.407592in}{0.222904in}}{\pgfqpoint{1.411163in}{1.411163in}}%
\pgfusepath{clip}%
\pgfsetbuttcap%
\pgfsetroundjoin%
\definecolor{currentfill}{rgb}{0.141176,0.345098,0.650980}%
\pgfsetfillcolor{currentfill}%
\pgfsetfillopacity{0.080000}%
\pgfsetlinewidth{0.000000pt}%
\definecolor{currentstroke}{rgb}{0.000000,0.000000,0.000000}%
\pgfsetstrokecolor{currentstroke}%
\pgfsetdash{}{0pt}%
\pgfpathmoveto{\pgfqpoint{1.279062in}{1.017342in}}%
\pgfpathlineto{\pgfqpoint{1.349507in}{1.000697in}}%
\pgfpathlineto{\pgfqpoint{1.350040in}{1.044884in}}%
\pgfpathlineto{\pgfqpoint{1.279420in}{1.061337in}}%
\pgfpathlineto{\pgfqpoint{1.279062in}{1.017342in}}%
\pgfpathclose%
\pgfusepath{fill}%
\end{pgfscope}%
\begin{pgfscope}%
\pgfpathrectangle{\pgfqpoint{0.407592in}{0.222904in}}{\pgfqpoint{1.411163in}{1.411163in}}%
\pgfusepath{clip}%
\pgfsetbuttcap%
\pgfsetroundjoin%
\definecolor{currentfill}{rgb}{0.141176,0.345098,0.650980}%
\pgfsetfillcolor{currentfill}%
\pgfsetfillopacity{0.080000}%
\pgfsetlinewidth{0.000000pt}%
\definecolor{currentstroke}{rgb}{0.000000,0.000000,0.000000}%
\pgfsetstrokecolor{currentstroke}%
\pgfsetdash{}{0pt}%
\pgfpathmoveto{\pgfqpoint{1.279420in}{1.061337in}}%
\pgfpathlineto{\pgfqpoint{1.350040in}{1.044884in}}%
\pgfpathlineto{\pgfqpoint{1.350575in}{1.089287in}}%
\pgfpathlineto{\pgfqpoint{1.279780in}{1.105548in}}%
\pgfpathlineto{\pgfqpoint{1.279420in}{1.061337in}}%
\pgfpathclose%
\pgfusepath{fill}%
\end{pgfscope}%
\begin{pgfscope}%
\pgfpathrectangle{\pgfqpoint{0.407592in}{0.222904in}}{\pgfqpoint{1.411163in}{1.411163in}}%
\pgfusepath{clip}%
\pgfsetbuttcap%
\pgfsetroundjoin%
\definecolor{currentfill}{rgb}{0.141176,0.345098,0.650980}%
\pgfsetfillcolor{currentfill}%
\pgfsetfillopacity{0.080000}%
\pgfsetlinewidth{0.000000pt}%
\definecolor{currentstroke}{rgb}{0.000000,0.000000,0.000000}%
\pgfsetstrokecolor{currentstroke}%
\pgfsetdash{}{0pt}%
\pgfpathmoveto{\pgfqpoint{1.348977in}{0.956726in}}%
\pgfpathlineto{\pgfqpoint{1.419989in}{0.939715in}}%
\pgfpathlineto{\pgfqpoint{1.420696in}{0.983877in}}%
\pgfpathlineto{\pgfqpoint{1.349507in}{1.000697in}}%
\pgfpathlineto{\pgfqpoint{1.348977in}{0.956726in}}%
\pgfpathclose%
\pgfusepath{fill}%
\end{pgfscope}%
\begin{pgfscope}%
\pgfpathrectangle{\pgfqpoint{0.407592in}{0.222904in}}{\pgfqpoint{1.411163in}{1.411163in}}%
\pgfusepath{clip}%
\pgfsetbuttcap%
\pgfsetroundjoin%
\definecolor{currentfill}{rgb}{0.141176,0.345098,0.650980}%
\pgfsetfillcolor{currentfill}%
\pgfsetfillopacity{0.080000}%
\pgfsetlinewidth{0.000000pt}%
\definecolor{currentstroke}{rgb}{0.000000,0.000000,0.000000}%
\pgfsetstrokecolor{currentstroke}%
\pgfsetdash{}{0pt}%
\pgfpathmoveto{\pgfqpoint{1.419989in}{0.939715in}}%
\pgfusepath{fill}%
\end{pgfscope}%
\begin{pgfscope}%
\pgfpathrectangle{\pgfqpoint{0.407592in}{0.222904in}}{\pgfqpoint{1.411163in}{1.411163in}}%
\pgfusepath{clip}%
\pgfsetbuttcap%
\pgfsetroundjoin%
\definecolor{currentfill}{rgb}{0.141176,0.345098,0.650980}%
\pgfsetfillcolor{currentfill}%
\pgfsetfillopacity{0.080000}%
\pgfsetlinewidth{0.000000pt}%
\definecolor{currentstroke}{rgb}{0.000000,0.000000,0.000000}%
\pgfsetstrokecolor{currentstroke}%
\pgfsetdash{}{0pt}%
\pgfpathmoveto{\pgfqpoint{1.349507in}{1.000697in}}%
\pgfpathlineto{\pgfqpoint{1.420696in}{0.983877in}}%
\pgfpathlineto{\pgfqpoint{1.421407in}{1.028256in}}%
\pgfpathlineto{\pgfqpoint{1.350040in}{1.044884in}}%
\pgfpathlineto{\pgfqpoint{1.349507in}{1.000697in}}%
\pgfpathclose%
\pgfusepath{fill}%
\end{pgfscope}%
\begin{pgfscope}%
\pgfpathrectangle{\pgfqpoint{0.407592in}{0.222904in}}{\pgfqpoint{1.411163in}{1.411163in}}%
\pgfusepath{clip}%
\pgfsetbuttcap%
\pgfsetroundjoin%
\definecolor{currentfill}{rgb}{0.141176,0.345098,0.650980}%
\pgfsetfillcolor{currentfill}%
\pgfsetfillopacity{0.080000}%
\pgfsetlinewidth{0.000000pt}%
\definecolor{currentstroke}{rgb}{0.000000,0.000000,0.000000}%
\pgfsetstrokecolor{currentstroke}%
\pgfsetdash{}{0pt}%
\pgfpathmoveto{\pgfqpoint{1.419989in}{0.939715in}}%
\pgfpathlineto{\pgfqpoint{1.420696in}{0.983877in}}%
\pgfusepath{fill}%
\end{pgfscope}%
\begin{pgfscope}%
\pgfpathrectangle{\pgfqpoint{0.407592in}{0.222904in}}{\pgfqpoint{1.411163in}{1.411163in}}%
\pgfusepath{clip}%
\pgfsetbuttcap%
\pgfsetroundjoin%
\definecolor{currentfill}{rgb}{0.141176,0.345098,0.650980}%
\pgfsetfillcolor{currentfill}%
\pgfsetfillopacity{0.080000}%
\pgfsetlinewidth{0.000000pt}%
\definecolor{currentstroke}{rgb}{0.000000,0.000000,0.000000}%
\pgfsetstrokecolor{currentstroke}%
\pgfsetdash{}{0pt}%
\pgfpathmoveto{\pgfqpoint{1.350040in}{1.044884in}}%
\pgfpathlineto{\pgfqpoint{1.421407in}{1.028256in}}%
\pgfpathlineto{\pgfqpoint{1.422122in}{1.072854in}}%
\pgfpathlineto{\pgfqpoint{1.350575in}{1.089287in}}%
\pgfpathlineto{\pgfqpoint{1.350040in}{1.044884in}}%
\pgfpathclose%
\pgfusepath{fill}%
\end{pgfscope}%
\begin{pgfscope}%
\pgfpathrectangle{\pgfqpoint{0.407592in}{0.222904in}}{\pgfqpoint{1.411163in}{1.411163in}}%
\pgfusepath{clip}%
\pgfsetbuttcap%
\pgfsetroundjoin%
\definecolor{currentfill}{rgb}{0.141176,0.345098,0.650980}%
\pgfsetfillcolor{currentfill}%
\pgfsetfillopacity{0.080000}%
\pgfsetlinewidth{0.000000pt}%
\definecolor{currentstroke}{rgb}{0.000000,0.000000,0.000000}%
\pgfsetstrokecolor{currentstroke}%
\pgfsetdash{}{0pt}%
\pgfpathmoveto{\pgfqpoint{1.420696in}{0.983877in}}%
\pgfpathlineto{\pgfqpoint{1.493535in}{1.011451in}}%
\pgfpathlineto{\pgfqpoint{1.421407in}{1.028256in}}%
\pgfusepath{fill}%
\end{pgfscope}%
\begin{pgfscope}%
\pgfpathrectangle{\pgfqpoint{0.407592in}{0.222904in}}{\pgfqpoint{1.411163in}{1.411163in}}%
\pgfusepath{clip}%
\pgfsetbuttcap%
\pgfsetroundjoin%
\definecolor{currentfill}{rgb}{0.141176,0.345098,0.650980}%
\pgfsetfillcolor{currentfill}%
\pgfsetfillopacity{0.080000}%
\pgfsetlinewidth{0.000000pt}%
\definecolor{currentstroke}{rgb}{0.000000,0.000000,0.000000}%
\pgfsetstrokecolor{currentstroke}%
\pgfsetdash{}{0pt}%
\pgfpathmoveto{\pgfqpoint{1.421407in}{1.028256in}}%
\pgfpathlineto{\pgfqpoint{1.493535in}{1.011451in}}%
\pgfpathlineto{\pgfqpoint{1.494432in}{1.056245in}}%
\pgfpathlineto{\pgfqpoint{1.422122in}{1.072854in}}%
\pgfpathlineto{\pgfqpoint{1.421407in}{1.028256in}}%
\pgfpathclose%
\pgfusepath{fill}%
\end{pgfscope}%
\begin{pgfscope}%
\pgfpathrectangle{\pgfqpoint{0.407592in}{0.222904in}}{\pgfqpoint{1.411163in}{1.411163in}}%
\pgfusepath{clip}%
\pgfsetbuttcap%
\pgfsetroundjoin%
\definecolor{currentfill}{rgb}{0.168627,0.611765,0.411765}%
\pgfsetfillcolor{currentfill}%
\pgfsetfillopacity{0.180000}%
\pgfsetlinewidth{0.000000pt}%
\definecolor{currentstroke}{rgb}{0.000000,0.000000,0.000000}%
\pgfsetstrokecolor{currentstroke}%
\pgfsetdash{}{0pt}%
\pgfpathmoveto{\pgfqpoint{1.668468in}{1.154894in}}%
\pgfpathlineto{\pgfqpoint{1.682713in}{1.180372in}}%
\pgfpathlineto{\pgfqpoint{1.683370in}{1.202112in}}%
\pgfpathlineto{\pgfqpoint{1.669109in}{1.176649in}}%
\pgfpathlineto{\pgfqpoint{1.668468in}{1.154894in}}%
\pgfpathclose%
\pgfusepath{fill}%
\end{pgfscope}%
\begin{pgfscope}%
\pgfpathrectangle{\pgfqpoint{0.407592in}{0.222904in}}{\pgfqpoint{1.411163in}{1.411163in}}%
\pgfusepath{clip}%
\pgfsetbuttcap%
\pgfsetroundjoin%
\definecolor{currentfill}{rgb}{0.168627,0.611765,0.411765}%
\pgfsetfillcolor{currentfill}%
\pgfsetfillopacity{0.180000}%
\pgfsetlinewidth{0.000000pt}%
\definecolor{currentstroke}{rgb}{0.000000,0.000000,0.000000}%
\pgfsetstrokecolor{currentstroke}%
\pgfsetdash{}{0pt}%
\pgfpathmoveto{\pgfqpoint{1.669109in}{1.176649in}}%
\pgfpathlineto{\pgfqpoint{1.683370in}{1.202112in}}%
\pgfpathlineto{\pgfqpoint{1.684029in}{1.223905in}}%
\pgfpathlineto{\pgfqpoint{1.669753in}{1.198456in}}%
\pgfpathlineto{\pgfqpoint{1.669109in}{1.176649in}}%
\pgfpathclose%
\pgfusepath{fill}%
\end{pgfscope}%
\begin{pgfscope}%
\pgfpathrectangle{\pgfqpoint{0.407592in}{0.222904in}}{\pgfqpoint{1.411163in}{1.411163in}}%
\pgfusepath{clip}%
\pgfsetbuttcap%
\pgfsetroundjoin%
\definecolor{currentfill}{rgb}{0.168627,0.611765,0.411765}%
\pgfsetfillcolor{currentfill}%
\pgfsetfillopacity{0.180000}%
\pgfsetlinewidth{0.000000pt}%
\definecolor{currentstroke}{rgb}{0.000000,0.000000,0.000000}%
\pgfsetstrokecolor{currentstroke}%
\pgfsetdash{}{0pt}%
\pgfpathmoveto{\pgfqpoint{1.654183in}{1.130004in}}%
\pgfpathlineto{\pgfqpoint{1.668468in}{1.154894in}}%
\pgfpathlineto{\pgfqpoint{1.669109in}{1.176649in}}%
\pgfpathlineto{\pgfqpoint{1.654809in}{1.151775in}}%
\pgfpathlineto{\pgfqpoint{1.654183in}{1.130004in}}%
\pgfpathclose%
\pgfusepath{fill}%
\end{pgfscope}%
\begin{pgfscope}%
\pgfpathrectangle{\pgfqpoint{0.407592in}{0.222904in}}{\pgfqpoint{1.411163in}{1.411163in}}%
\pgfusepath{clip}%
\pgfsetbuttcap%
\pgfsetroundjoin%
\definecolor{currentfill}{rgb}{0.168627,0.611765,0.411765}%
\pgfsetfillcolor{currentfill}%
\pgfsetfillopacity{0.180000}%
\pgfsetlinewidth{0.000000pt}%
\definecolor{currentstroke}{rgb}{0.000000,0.000000,0.000000}%
\pgfsetstrokecolor{currentstroke}%
\pgfsetdash{}{0pt}%
\pgfpathmoveto{\pgfqpoint{1.669753in}{1.198456in}}%
\pgfpathlineto{\pgfqpoint{1.684029in}{1.223905in}}%
\pgfpathlineto{\pgfqpoint{1.684690in}{1.245750in}}%
\pgfpathlineto{\pgfqpoint{1.670397in}{1.220316in}}%
\pgfpathlineto{\pgfqpoint{1.669753in}{1.198456in}}%
\pgfpathclose%
\pgfusepath{fill}%
\end{pgfscope}%
\begin{pgfscope}%
\pgfpathrectangle{\pgfqpoint{0.407592in}{0.222904in}}{\pgfqpoint{1.411163in}{1.411163in}}%
\pgfusepath{clip}%
\pgfsetbuttcap%
\pgfsetroundjoin%
\definecolor{currentfill}{rgb}{0.168627,0.611765,0.411765}%
\pgfsetfillcolor{currentfill}%
\pgfsetfillopacity{0.180000}%
\pgfsetlinewidth{0.000000pt}%
\definecolor{currentstroke}{rgb}{0.000000,0.000000,0.000000}%
\pgfsetstrokecolor{currentstroke}%
\pgfsetdash{}{0pt}%
\pgfpathmoveto{\pgfqpoint{1.654809in}{1.151775in}}%
\pgfpathlineto{\pgfqpoint{1.669109in}{1.176649in}}%
\pgfpathlineto{\pgfqpoint{1.669753in}{1.198456in}}%
\pgfpathlineto{\pgfqpoint{1.655437in}{1.173599in}}%
\pgfpathlineto{\pgfqpoint{1.654809in}{1.151775in}}%
\pgfpathclose%
\pgfusepath{fill}%
\end{pgfscope}%
\begin{pgfscope}%
\pgfpathrectangle{\pgfqpoint{0.407592in}{0.222904in}}{\pgfqpoint{1.411163in}{1.411163in}}%
\pgfusepath{clip}%
\pgfsetbuttcap%
\pgfsetroundjoin%
\definecolor{currentfill}{rgb}{0.168627,0.611765,0.411765}%
\pgfsetfillcolor{currentfill}%
\pgfsetfillopacity{0.180000}%
\pgfsetlinewidth{0.000000pt}%
\definecolor{currentstroke}{rgb}{0.000000,0.000000,0.000000}%
\pgfsetstrokecolor{currentstroke}%
\pgfsetdash{}{0pt}%
\pgfpathmoveto{\pgfqpoint{1.670397in}{1.220316in}}%
\pgfpathlineto{\pgfqpoint{1.684690in}{1.245750in}}%
\pgfpathlineto{\pgfqpoint{1.685352in}{1.267647in}}%
\pgfpathlineto{\pgfqpoint{1.671044in}{1.242228in}}%
\pgfpathlineto{\pgfqpoint{1.670397in}{1.220316in}}%
\pgfpathclose%
\pgfusepath{fill}%
\end{pgfscope}%
\begin{pgfscope}%
\pgfpathrectangle{\pgfqpoint{0.407592in}{0.222904in}}{\pgfqpoint{1.411163in}{1.411163in}}%
\pgfusepath{clip}%
\pgfsetbuttcap%
\pgfsetroundjoin%
\definecolor{currentfill}{rgb}{0.168627,0.611765,0.411765}%
\pgfsetfillcolor{currentfill}%
\pgfsetfillopacity{0.180000}%
\pgfsetlinewidth{0.000000pt}%
\definecolor{currentstroke}{rgb}{0.000000,0.000000,0.000000}%
\pgfsetstrokecolor{currentstroke}%
\pgfsetdash{}{0pt}%
\pgfpathmoveto{\pgfqpoint{1.639860in}{1.105771in}}%
\pgfpathlineto{\pgfqpoint{1.654183in}{1.130004in}}%
\pgfpathlineto{\pgfqpoint{1.654809in}{1.151775in}}%
\pgfpathlineto{\pgfqpoint{1.640470in}{1.127559in}}%
\pgfpathlineto{\pgfqpoint{1.639860in}{1.105771in}}%
\pgfpathclose%
\pgfusepath{fill}%
\end{pgfscope}%
\begin{pgfscope}%
\pgfpathrectangle{\pgfqpoint{0.407592in}{0.222904in}}{\pgfqpoint{1.411163in}{1.411163in}}%
\pgfusepath{clip}%
\pgfsetbuttcap%
\pgfsetroundjoin%
\definecolor{currentfill}{rgb}{0.168627,0.611765,0.411765}%
\pgfsetfillcolor{currentfill}%
\pgfsetfillopacity{0.180000}%
\pgfsetlinewidth{0.000000pt}%
\definecolor{currentstroke}{rgb}{0.000000,0.000000,0.000000}%
\pgfsetstrokecolor{currentstroke}%
\pgfsetdash{}{0pt}%
\pgfpathmoveto{\pgfqpoint{1.655437in}{1.173599in}}%
\pgfpathlineto{\pgfqpoint{1.669753in}{1.198456in}}%
\pgfpathlineto{\pgfqpoint{1.670397in}{1.220316in}}%
\pgfpathlineto{\pgfqpoint{1.656066in}{1.195474in}}%
\pgfpathlineto{\pgfqpoint{1.655437in}{1.173599in}}%
\pgfpathclose%
\pgfusepath{fill}%
\end{pgfscope}%
\begin{pgfscope}%
\pgfpathrectangle{\pgfqpoint{0.407592in}{0.222904in}}{\pgfqpoint{1.411163in}{1.411163in}}%
\pgfusepath{clip}%
\pgfsetbuttcap%
\pgfsetroundjoin%
\definecolor{currentfill}{rgb}{0.168627,0.611765,0.411765}%
\pgfsetfillcolor{currentfill}%
\pgfsetfillopacity{0.180000}%
\pgfsetlinewidth{0.000000pt}%
\definecolor{currentstroke}{rgb}{0.000000,0.000000,0.000000}%
\pgfsetstrokecolor{currentstroke}%
\pgfsetdash{}{0pt}%
\pgfpathmoveto{\pgfqpoint{1.671044in}{1.242228in}}%
\pgfpathlineto{\pgfqpoint{1.685352in}{1.267647in}}%
\pgfpathlineto{\pgfqpoint{1.686016in}{1.289597in}}%
\pgfpathlineto{\pgfqpoint{1.671692in}{1.264192in}}%
\pgfpathlineto{\pgfqpoint{1.671044in}{1.242228in}}%
\pgfpathclose%
\pgfusepath{fill}%
\end{pgfscope}%
\begin{pgfscope}%
\pgfpathrectangle{\pgfqpoint{0.407592in}{0.222904in}}{\pgfqpoint{1.411163in}{1.411163in}}%
\pgfusepath{clip}%
\pgfsetbuttcap%
\pgfsetroundjoin%
\definecolor{currentfill}{rgb}{0.168627,0.611765,0.411765}%
\pgfsetfillcolor{currentfill}%
\pgfsetfillopacity{0.180000}%
\pgfsetlinewidth{0.000000pt}%
\definecolor{currentstroke}{rgb}{0.000000,0.000000,0.000000}%
\pgfsetstrokecolor{currentstroke}%
\pgfsetdash{}{0pt}%
\pgfpathmoveto{\pgfqpoint{1.640470in}{1.127559in}}%
\pgfpathlineto{\pgfqpoint{1.654809in}{1.151775in}}%
\pgfpathlineto{\pgfqpoint{1.655437in}{1.173599in}}%
\pgfpathlineto{\pgfqpoint{1.641082in}{1.149401in}}%
\pgfpathlineto{\pgfqpoint{1.640470in}{1.127559in}}%
\pgfpathclose%
\pgfusepath{fill}%
\end{pgfscope}%
\begin{pgfscope}%
\pgfpathrectangle{\pgfqpoint{0.407592in}{0.222904in}}{\pgfqpoint{1.411163in}{1.411163in}}%
\pgfusepath{clip}%
\pgfsetbuttcap%
\pgfsetroundjoin%
\definecolor{currentfill}{rgb}{0.168627,0.611765,0.411765}%
\pgfsetfillcolor{currentfill}%
\pgfsetfillopacity{0.180000}%
\pgfsetlinewidth{0.000000pt}%
\definecolor{currentstroke}{rgb}{0.000000,0.000000,0.000000}%
\pgfsetstrokecolor{currentstroke}%
\pgfsetdash{}{0pt}%
\pgfpathmoveto{\pgfqpoint{1.656066in}{1.195474in}}%
\pgfpathlineto{\pgfqpoint{1.670397in}{1.220316in}}%
\pgfpathlineto{\pgfqpoint{1.671044in}{1.242228in}}%
\pgfpathlineto{\pgfqpoint{1.656696in}{1.217403in}}%
\pgfpathlineto{\pgfqpoint{1.656066in}{1.195474in}}%
\pgfpathclose%
\pgfusepath{fill}%
\end{pgfscope}%
\begin{pgfscope}%
\pgfpathrectangle{\pgfqpoint{0.407592in}{0.222904in}}{\pgfqpoint{1.411163in}{1.411163in}}%
\pgfusepath{clip}%
\pgfsetbuttcap%
\pgfsetroundjoin%
\definecolor{currentfill}{rgb}{0.168627,0.611765,0.411765}%
\pgfsetfillcolor{currentfill}%
\pgfsetfillopacity{0.180000}%
\pgfsetlinewidth{0.000000pt}%
\definecolor{currentstroke}{rgb}{0.000000,0.000000,0.000000}%
\pgfsetstrokecolor{currentstroke}%
\pgfsetdash{}{0pt}%
\pgfpathmoveto{\pgfqpoint{1.671692in}{1.264192in}}%
\pgfpathlineto{\pgfqpoint{1.686016in}{1.289597in}}%
\pgfpathlineto{\pgfqpoint{1.686681in}{1.311599in}}%
\pgfpathlineto{\pgfqpoint{1.672341in}{1.286210in}}%
\pgfpathlineto{\pgfqpoint{1.671692in}{1.264192in}}%
\pgfpathclose%
\pgfusepath{fill}%
\end{pgfscope}%
\begin{pgfscope}%
\pgfpathrectangle{\pgfqpoint{0.407592in}{0.222904in}}{\pgfqpoint{1.411163in}{1.411163in}}%
\pgfusepath{clip}%
\pgfsetbuttcap%
\pgfsetroundjoin%
\definecolor{currentfill}{rgb}{0.168627,0.611765,0.411765}%
\pgfsetfillcolor{currentfill}%
\pgfsetfillopacity{0.180000}%
\pgfsetlinewidth{0.000000pt}%
\definecolor{currentstroke}{rgb}{0.000000,0.000000,0.000000}%
\pgfsetstrokecolor{currentstroke}%
\pgfsetdash{}{0pt}%
\pgfpathmoveto{\pgfqpoint{1.625498in}{1.082268in}}%
\pgfpathlineto{\pgfqpoint{1.639860in}{1.105771in}}%
\pgfpathlineto{\pgfqpoint{1.640470in}{1.127559in}}%
\pgfpathlineto{\pgfqpoint{1.626092in}{1.104076in}}%
\pgfpathlineto{\pgfqpoint{1.625498in}{1.082268in}}%
\pgfpathclose%
\pgfusepath{fill}%
\end{pgfscope}%
\begin{pgfscope}%
\pgfpathrectangle{\pgfqpoint{0.407592in}{0.222904in}}{\pgfqpoint{1.411163in}{1.411163in}}%
\pgfusepath{clip}%
\pgfsetbuttcap%
\pgfsetroundjoin%
\definecolor{currentfill}{rgb}{0.168627,0.611765,0.411765}%
\pgfsetfillcolor{currentfill}%
\pgfsetfillopacity{0.180000}%
\pgfsetlinewidth{0.000000pt}%
\definecolor{currentstroke}{rgb}{0.000000,0.000000,0.000000}%
\pgfsetstrokecolor{currentstroke}%
\pgfsetdash{}{0pt}%
\pgfpathmoveto{\pgfqpoint{1.641082in}{1.149401in}}%
\pgfpathlineto{\pgfqpoint{1.655437in}{1.173599in}}%
\pgfpathlineto{\pgfqpoint{1.656066in}{1.195474in}}%
\pgfpathlineto{\pgfqpoint{1.641695in}{1.171294in}}%
\pgfpathlineto{\pgfqpoint{1.641082in}{1.149401in}}%
\pgfpathclose%
\pgfusepath{fill}%
\end{pgfscope}%
\begin{pgfscope}%
\pgfpathrectangle{\pgfqpoint{0.407592in}{0.222904in}}{\pgfqpoint{1.411163in}{1.411163in}}%
\pgfusepath{clip}%
\pgfsetbuttcap%
\pgfsetroundjoin%
\definecolor{currentfill}{rgb}{0.168627,0.611765,0.411765}%
\pgfsetfillcolor{currentfill}%
\pgfsetfillopacity{0.180000}%
\pgfsetlinewidth{0.000000pt}%
\definecolor{currentstroke}{rgb}{0.000000,0.000000,0.000000}%
\pgfsetstrokecolor{currentstroke}%
\pgfsetdash{}{0pt}%
\pgfpathmoveto{\pgfqpoint{1.656696in}{1.217403in}}%
\pgfpathlineto{\pgfqpoint{1.671044in}{1.242228in}}%
\pgfpathlineto{\pgfqpoint{1.671692in}{1.264192in}}%
\pgfpathlineto{\pgfqpoint{1.657328in}{1.239384in}}%
\pgfpathlineto{\pgfqpoint{1.656696in}{1.217403in}}%
\pgfpathclose%
\pgfusepath{fill}%
\end{pgfscope}%
\begin{pgfscope}%
\pgfpathrectangle{\pgfqpoint{0.407592in}{0.222904in}}{\pgfqpoint{1.411163in}{1.411163in}}%
\pgfusepath{clip}%
\pgfsetbuttcap%
\pgfsetroundjoin%
\definecolor{currentfill}{rgb}{0.168627,0.611765,0.411765}%
\pgfsetfillcolor{currentfill}%
\pgfsetfillopacity{0.180000}%
\pgfsetlinewidth{0.000000pt}%
\definecolor{currentstroke}{rgb}{0.000000,0.000000,0.000000}%
\pgfsetstrokecolor{currentstroke}%
\pgfsetdash{}{0pt}%
\pgfpathmoveto{\pgfqpoint{1.672341in}{1.286210in}}%
\pgfpathlineto{\pgfqpoint{1.686681in}{1.311599in}}%
\pgfpathlineto{\pgfqpoint{1.687348in}{1.333655in}}%
\pgfpathlineto{\pgfqpoint{1.672992in}{1.308281in}}%
\pgfpathlineto{\pgfqpoint{1.672341in}{1.286210in}}%
\pgfpathclose%
\pgfusepath{fill}%
\end{pgfscope}%
\begin{pgfscope}%
\pgfpathrectangle{\pgfqpoint{0.407592in}{0.222904in}}{\pgfqpoint{1.411163in}{1.411163in}}%
\pgfusepath{clip}%
\pgfsetbuttcap%
\pgfsetroundjoin%
\definecolor{currentfill}{rgb}{0.168627,0.611765,0.411765}%
\pgfsetfillcolor{currentfill}%
\pgfsetfillopacity{0.180000}%
\pgfsetlinewidth{0.000000pt}%
\definecolor{currentstroke}{rgb}{0.000000,0.000000,0.000000}%
\pgfsetstrokecolor{currentstroke}%
\pgfsetdash{}{0pt}%
\pgfpathmoveto{\pgfqpoint{1.626092in}{1.104076in}}%
\pgfpathlineto{\pgfqpoint{1.640470in}{1.127559in}}%
\pgfpathlineto{\pgfqpoint{1.641082in}{1.149401in}}%
\pgfpathlineto{\pgfqpoint{1.626687in}{1.125937in}}%
\pgfpathlineto{\pgfqpoint{1.626092in}{1.104076in}}%
\pgfpathclose%
\pgfusepath{fill}%
\end{pgfscope}%
\begin{pgfscope}%
\pgfpathrectangle{\pgfqpoint{0.407592in}{0.222904in}}{\pgfqpoint{1.411163in}{1.411163in}}%
\pgfusepath{clip}%
\pgfsetbuttcap%
\pgfsetroundjoin%
\definecolor{currentfill}{rgb}{0.168627,0.611765,0.411765}%
\pgfsetfillcolor{currentfill}%
\pgfsetfillopacity{0.180000}%
\pgfsetlinewidth{0.000000pt}%
\definecolor{currentstroke}{rgb}{0.000000,0.000000,0.000000}%
\pgfsetstrokecolor{currentstroke}%
\pgfsetdash{}{0pt}%
\pgfpathmoveto{\pgfqpoint{1.641695in}{1.171294in}}%
\pgfpathlineto{\pgfqpoint{1.656066in}{1.195474in}}%
\pgfpathlineto{\pgfqpoint{1.656696in}{1.217403in}}%
\pgfpathlineto{\pgfqpoint{1.642309in}{1.193241in}}%
\pgfpathlineto{\pgfqpoint{1.641695in}{1.171294in}}%
\pgfpathclose%
\pgfusepath{fill}%
\end{pgfscope}%
\begin{pgfscope}%
\pgfpathrectangle{\pgfqpoint{0.407592in}{0.222904in}}{\pgfqpoint{1.411163in}{1.411163in}}%
\pgfusepath{clip}%
\pgfsetbuttcap%
\pgfsetroundjoin%
\definecolor{currentfill}{rgb}{0.168627,0.611765,0.411765}%
\pgfsetfillcolor{currentfill}%
\pgfsetfillopacity{0.180000}%
\pgfsetlinewidth{0.000000pt}%
\definecolor{currentstroke}{rgb}{0.000000,0.000000,0.000000}%
\pgfsetstrokecolor{currentstroke}%
\pgfsetdash{}{0pt}%
\pgfpathmoveto{\pgfqpoint{1.657328in}{1.239384in}}%
\pgfpathlineto{\pgfqpoint{1.671692in}{1.264192in}}%
\pgfpathlineto{\pgfqpoint{1.672341in}{1.286210in}}%
\pgfpathlineto{\pgfqpoint{1.657962in}{1.261418in}}%
\pgfpathlineto{\pgfqpoint{1.657328in}{1.239384in}}%
\pgfpathclose%
\pgfusepath{fill}%
\end{pgfscope}%
\begin{pgfscope}%
\pgfpathrectangle{\pgfqpoint{0.407592in}{0.222904in}}{\pgfqpoint{1.411163in}{1.411163in}}%
\pgfusepath{clip}%
\pgfsetbuttcap%
\pgfsetroundjoin%
\definecolor{currentfill}{rgb}{0.168627,0.611765,0.411765}%
\pgfsetfillcolor{currentfill}%
\pgfsetfillopacity{0.180000}%
\pgfsetlinewidth{0.000000pt}%
\definecolor{currentstroke}{rgb}{0.000000,0.000000,0.000000}%
\pgfsetstrokecolor{currentstroke}%
\pgfsetdash{}{0pt}%
\pgfpathmoveto{\pgfqpoint{1.672992in}{1.308281in}}%
\pgfpathlineto{\pgfqpoint{1.687348in}{1.333655in}}%
\pgfpathlineto{\pgfqpoint{1.688017in}{1.355763in}}%
\pgfpathlineto{\pgfqpoint{1.673645in}{1.330405in}}%
\pgfpathlineto{\pgfqpoint{1.672992in}{1.308281in}}%
\pgfpathclose%
\pgfusepath{fill}%
\end{pgfscope}%
\begin{pgfscope}%
\pgfpathrectangle{\pgfqpoint{0.407592in}{0.222904in}}{\pgfqpoint{1.411163in}{1.411163in}}%
\pgfusepath{clip}%
\pgfsetbuttcap%
\pgfsetroundjoin%
\definecolor{currentfill}{rgb}{0.168627,0.611765,0.411765}%
\pgfsetfillcolor{currentfill}%
\pgfsetfillopacity{0.180000}%
\pgfsetlinewidth{0.000000pt}%
\definecolor{currentstroke}{rgb}{0.000000,0.000000,0.000000}%
\pgfsetstrokecolor{currentstroke}%
\pgfsetdash{}{0pt}%
\pgfpathmoveto{\pgfqpoint{1.611096in}{1.059575in}}%
\pgfpathlineto{\pgfqpoint{1.625498in}{1.082268in}}%
\pgfpathlineto{\pgfqpoint{1.626092in}{1.104076in}}%
\pgfpathlineto{\pgfqpoint{1.611674in}{1.081405in}}%
\pgfpathlineto{\pgfqpoint{1.611096in}{1.059575in}}%
\pgfpathclose%
\pgfusepath{fill}%
\end{pgfscope}%
\begin{pgfscope}%
\pgfpathrectangle{\pgfqpoint{0.407592in}{0.222904in}}{\pgfqpoint{1.411163in}{1.411163in}}%
\pgfusepath{clip}%
\pgfsetbuttcap%
\pgfsetroundjoin%
\definecolor{currentfill}{rgb}{0.168627,0.611765,0.411765}%
\pgfsetfillcolor{currentfill}%
\pgfsetfillopacity{0.180000}%
\pgfsetlinewidth{0.000000pt}%
\definecolor{currentstroke}{rgb}{0.000000,0.000000,0.000000}%
\pgfsetstrokecolor{currentstroke}%
\pgfsetdash{}{0pt}%
\pgfpathmoveto{\pgfqpoint{1.626687in}{1.125937in}}%
\pgfpathlineto{\pgfqpoint{1.641082in}{1.149401in}}%
\pgfpathlineto{\pgfqpoint{1.641695in}{1.171294in}}%
\pgfpathlineto{\pgfqpoint{1.627284in}{1.147851in}}%
\pgfpathlineto{\pgfqpoint{1.626687in}{1.125937in}}%
\pgfpathclose%
\pgfusepath{fill}%
\end{pgfscope}%
\begin{pgfscope}%
\pgfpathrectangle{\pgfqpoint{0.407592in}{0.222904in}}{\pgfqpoint{1.411163in}{1.411163in}}%
\pgfusepath{clip}%
\pgfsetbuttcap%
\pgfsetroundjoin%
\definecolor{currentfill}{rgb}{0.168627,0.611765,0.411765}%
\pgfsetfillcolor{currentfill}%
\pgfsetfillopacity{0.180000}%
\pgfsetlinewidth{0.000000pt}%
\definecolor{currentstroke}{rgb}{0.000000,0.000000,0.000000}%
\pgfsetstrokecolor{currentstroke}%
\pgfsetdash{}{0pt}%
\pgfpathmoveto{\pgfqpoint{1.642309in}{1.193241in}}%
\pgfpathlineto{\pgfqpoint{1.656696in}{1.217403in}}%
\pgfpathlineto{\pgfqpoint{1.657328in}{1.239384in}}%
\pgfpathlineto{\pgfqpoint{1.642925in}{1.215241in}}%
\pgfpathlineto{\pgfqpoint{1.642309in}{1.193241in}}%
\pgfpathclose%
\pgfusepath{fill}%
\end{pgfscope}%
\begin{pgfscope}%
\pgfpathrectangle{\pgfqpoint{0.407592in}{0.222904in}}{\pgfqpoint{1.411163in}{1.411163in}}%
\pgfusepath{clip}%
\pgfsetbuttcap%
\pgfsetroundjoin%
\definecolor{currentfill}{rgb}{0.168627,0.611765,0.411765}%
\pgfsetfillcolor{currentfill}%
\pgfsetfillopacity{0.180000}%
\pgfsetlinewidth{0.000000pt}%
\definecolor{currentstroke}{rgb}{0.000000,0.000000,0.000000}%
\pgfsetstrokecolor{currentstroke}%
\pgfsetdash{}{0pt}%
\pgfpathmoveto{\pgfqpoint{1.657962in}{1.261418in}}%
\pgfpathlineto{\pgfqpoint{1.672341in}{1.286210in}}%
\pgfpathlineto{\pgfqpoint{1.672992in}{1.308281in}}%
\pgfpathlineto{\pgfqpoint{1.658597in}{1.283506in}}%
\pgfpathlineto{\pgfqpoint{1.657962in}{1.261418in}}%
\pgfpathclose%
\pgfusepath{fill}%
\end{pgfscope}%
\begin{pgfscope}%
\pgfpathrectangle{\pgfqpoint{0.407592in}{0.222904in}}{\pgfqpoint{1.411163in}{1.411163in}}%
\pgfusepath{clip}%
\pgfsetbuttcap%
\pgfsetroundjoin%
\definecolor{currentfill}{rgb}{0.168627,0.611765,0.411765}%
\pgfsetfillcolor{currentfill}%
\pgfsetfillopacity{0.180000}%
\pgfsetlinewidth{0.000000pt}%
\definecolor{currentstroke}{rgb}{0.000000,0.000000,0.000000}%
\pgfsetstrokecolor{currentstroke}%
\pgfsetdash{}{0pt}%
\pgfpathmoveto{\pgfqpoint{1.673645in}{1.330405in}}%
\pgfpathlineto{\pgfqpoint{1.688017in}{1.355763in}}%
\pgfpathlineto{\pgfqpoint{1.688687in}{1.377925in}}%
\pgfpathlineto{\pgfqpoint{1.674299in}{1.352582in}}%
\pgfpathlineto{\pgfqpoint{1.673645in}{1.330405in}}%
\pgfpathclose%
\pgfusepath{fill}%
\end{pgfscope}%
\begin{pgfscope}%
\pgfpathrectangle{\pgfqpoint{0.407592in}{0.222904in}}{\pgfqpoint{1.411163in}{1.411163in}}%
\pgfusepath{clip}%
\pgfsetbuttcap%
\pgfsetroundjoin%
\definecolor{currentfill}{rgb}{0.168627,0.611765,0.411765}%
\pgfsetfillcolor{currentfill}%
\pgfsetfillopacity{0.180000}%
\pgfsetlinewidth{0.000000pt}%
\definecolor{currentstroke}{rgb}{0.000000,0.000000,0.000000}%
\pgfsetstrokecolor{currentstroke}%
\pgfsetdash{}{0pt}%
\pgfpathmoveto{\pgfqpoint{1.611674in}{1.081405in}}%
\pgfpathlineto{\pgfqpoint{1.626092in}{1.104076in}}%
\pgfpathlineto{\pgfqpoint{1.626687in}{1.125937in}}%
\pgfpathlineto{\pgfqpoint{1.612254in}{1.103288in}}%
\pgfpathlineto{\pgfqpoint{1.611674in}{1.081405in}}%
\pgfpathclose%
\pgfusepath{fill}%
\end{pgfscope}%
\begin{pgfscope}%
\pgfpathrectangle{\pgfqpoint{0.407592in}{0.222904in}}{\pgfqpoint{1.411163in}{1.411163in}}%
\pgfusepath{clip}%
\pgfsetbuttcap%
\pgfsetroundjoin%
\definecolor{currentfill}{rgb}{0.168627,0.611765,0.411765}%
\pgfsetfillcolor{currentfill}%
\pgfsetfillopacity{0.180000}%
\pgfsetlinewidth{0.000000pt}%
\definecolor{currentstroke}{rgb}{0.000000,0.000000,0.000000}%
\pgfsetstrokecolor{currentstroke}%
\pgfsetdash{}{0pt}%
\pgfpathmoveto{\pgfqpoint{1.627284in}{1.147851in}}%
\pgfpathlineto{\pgfqpoint{1.641695in}{1.171294in}}%
\pgfpathlineto{\pgfqpoint{1.642309in}{1.193241in}}%
\pgfpathlineto{\pgfqpoint{1.627883in}{1.169818in}}%
\pgfpathlineto{\pgfqpoint{1.627284in}{1.147851in}}%
\pgfpathclose%
\pgfusepath{fill}%
\end{pgfscope}%
\begin{pgfscope}%
\pgfpathrectangle{\pgfqpoint{0.407592in}{0.222904in}}{\pgfqpoint{1.411163in}{1.411163in}}%
\pgfusepath{clip}%
\pgfsetbuttcap%
\pgfsetroundjoin%
\definecolor{currentfill}{rgb}{0.168627,0.611765,0.411765}%
\pgfsetfillcolor{currentfill}%
\pgfsetfillopacity{0.180000}%
\pgfsetlinewidth{0.000000pt}%
\definecolor{currentstroke}{rgb}{0.000000,0.000000,0.000000}%
\pgfsetstrokecolor{currentstroke}%
\pgfsetdash{}{0pt}%
\pgfpathmoveto{\pgfqpoint{1.642925in}{1.215241in}}%
\pgfpathlineto{\pgfqpoint{1.657328in}{1.239384in}}%
\pgfpathlineto{\pgfqpoint{1.657962in}{1.261418in}}%
\pgfpathlineto{\pgfqpoint{1.643543in}{1.237293in}}%
\pgfpathlineto{\pgfqpoint{1.642925in}{1.215241in}}%
\pgfpathclose%
\pgfusepath{fill}%
\end{pgfscope}%
\begin{pgfscope}%
\pgfpathrectangle{\pgfqpoint{0.407592in}{0.222904in}}{\pgfqpoint{1.411163in}{1.411163in}}%
\pgfusepath{clip}%
\pgfsetbuttcap%
\pgfsetroundjoin%
\definecolor{currentfill}{rgb}{0.168627,0.611765,0.411765}%
\pgfsetfillcolor{currentfill}%
\pgfsetfillopacity{0.180000}%
\pgfsetlinewidth{0.000000pt}%
\definecolor{currentstroke}{rgb}{0.000000,0.000000,0.000000}%
\pgfsetstrokecolor{currentstroke}%
\pgfsetdash{}{0pt}%
\pgfpathmoveto{\pgfqpoint{1.658597in}{1.283506in}}%
\pgfpathlineto{\pgfqpoint{1.672992in}{1.308281in}}%
\pgfpathlineto{\pgfqpoint{1.673645in}{1.330405in}}%
\pgfpathlineto{\pgfqpoint{1.659233in}{1.305647in}}%
\pgfpathlineto{\pgfqpoint{1.658597in}{1.283506in}}%
\pgfpathclose%
\pgfusepath{fill}%
\end{pgfscope}%
\begin{pgfscope}%
\pgfpathrectangle{\pgfqpoint{0.407592in}{0.222904in}}{\pgfqpoint{1.411163in}{1.411163in}}%
\pgfusepath{clip}%
\pgfsetbuttcap%
\pgfsetroundjoin%
\definecolor{currentfill}{rgb}{0.168627,0.611765,0.411765}%
\pgfsetfillcolor{currentfill}%
\pgfsetfillopacity{0.180000}%
\pgfsetlinewidth{0.000000pt}%
\definecolor{currentstroke}{rgb}{0.000000,0.000000,0.000000}%
\pgfsetstrokecolor{currentstroke}%
\pgfsetdash{}{0pt}%
\pgfpathmoveto{\pgfqpoint{1.674299in}{1.352582in}}%
\pgfpathlineto{\pgfqpoint{1.688687in}{1.377925in}}%
\pgfpathlineto{\pgfqpoint{1.689359in}{1.400140in}}%
\pgfpathlineto{\pgfqpoint{1.674955in}{1.374813in}}%
\pgfpathlineto{\pgfqpoint{1.674299in}{1.352582in}}%
\pgfpathclose%
\pgfusepath{fill}%
\end{pgfscope}%
\begin{pgfscope}%
\pgfpathrectangle{\pgfqpoint{0.407592in}{0.222904in}}{\pgfqpoint{1.411163in}{1.411163in}}%
\pgfusepath{clip}%
\pgfsetbuttcap%
\pgfsetroundjoin%
\definecolor{currentfill}{rgb}{0.168627,0.611765,0.411765}%
\pgfsetfillcolor{currentfill}%
\pgfsetfillopacity{0.180000}%
\pgfsetlinewidth{0.000000pt}%
\definecolor{currentstroke}{rgb}{0.000000,0.000000,0.000000}%
\pgfsetstrokecolor{currentstroke}%
\pgfsetdash{}{0pt}%
\pgfpathmoveto{\pgfqpoint{1.596655in}{1.037776in}}%
\pgfpathlineto{\pgfqpoint{1.611096in}{1.059575in}}%
\pgfpathlineto{\pgfqpoint{1.611674in}{1.081405in}}%
\pgfpathlineto{\pgfqpoint{1.597217in}{1.059631in}}%
\pgfpathlineto{\pgfqpoint{1.596655in}{1.037776in}}%
\pgfpathclose%
\pgfusepath{fill}%
\end{pgfscope}%
\begin{pgfscope}%
\pgfpathrectangle{\pgfqpoint{0.407592in}{0.222904in}}{\pgfqpoint{1.411163in}{1.411163in}}%
\pgfusepath{clip}%
\pgfsetbuttcap%
\pgfsetroundjoin%
\definecolor{currentfill}{rgb}{0.168627,0.611765,0.411765}%
\pgfsetfillcolor{currentfill}%
\pgfsetfillopacity{0.180000}%
\pgfsetlinewidth{0.000000pt}%
\definecolor{currentstroke}{rgb}{0.000000,0.000000,0.000000}%
\pgfsetstrokecolor{currentstroke}%
\pgfsetdash{}{0pt}%
\pgfpathmoveto{\pgfqpoint{1.612254in}{1.103288in}}%
\pgfpathlineto{\pgfqpoint{1.626687in}{1.125937in}}%
\pgfpathlineto{\pgfqpoint{1.627284in}{1.147851in}}%
\pgfpathlineto{\pgfqpoint{1.612835in}{1.125224in}}%
\pgfpathlineto{\pgfqpoint{1.612254in}{1.103288in}}%
\pgfpathclose%
\pgfusepath{fill}%
\end{pgfscope}%
\begin{pgfscope}%
\pgfpathrectangle{\pgfqpoint{0.407592in}{0.222904in}}{\pgfqpoint{1.411163in}{1.411163in}}%
\pgfusepath{clip}%
\pgfsetbuttcap%
\pgfsetroundjoin%
\definecolor{currentfill}{rgb}{0.168627,0.611765,0.411765}%
\pgfsetfillcolor{currentfill}%
\pgfsetfillopacity{0.180000}%
\pgfsetlinewidth{0.000000pt}%
\definecolor{currentstroke}{rgb}{0.000000,0.000000,0.000000}%
\pgfsetstrokecolor{currentstroke}%
\pgfsetdash{}{0pt}%
\pgfpathmoveto{\pgfqpoint{1.627883in}{1.169818in}}%
\pgfpathlineto{\pgfqpoint{1.642309in}{1.193241in}}%
\pgfpathlineto{\pgfqpoint{1.642925in}{1.215241in}}%
\pgfpathlineto{\pgfqpoint{1.628483in}{1.191838in}}%
\pgfpathlineto{\pgfqpoint{1.627883in}{1.169818in}}%
\pgfpathclose%
\pgfusepath{fill}%
\end{pgfscope}%
\begin{pgfscope}%
\pgfpathrectangle{\pgfqpoint{0.407592in}{0.222904in}}{\pgfqpoint{1.411163in}{1.411163in}}%
\pgfusepath{clip}%
\pgfsetbuttcap%
\pgfsetroundjoin%
\definecolor{currentfill}{rgb}{0.168627,0.611765,0.411765}%
\pgfsetfillcolor{currentfill}%
\pgfsetfillopacity{0.180000}%
\pgfsetlinewidth{0.000000pt}%
\definecolor{currentstroke}{rgb}{0.000000,0.000000,0.000000}%
\pgfsetstrokecolor{currentstroke}%
\pgfsetdash{}{0pt}%
\pgfpathmoveto{\pgfqpoint{1.643543in}{1.237293in}}%
\pgfpathlineto{\pgfqpoint{1.657962in}{1.261418in}}%
\pgfpathlineto{\pgfqpoint{1.658597in}{1.283506in}}%
\pgfpathlineto{\pgfqpoint{1.644162in}{1.259399in}}%
\pgfpathlineto{\pgfqpoint{1.643543in}{1.237293in}}%
\pgfpathclose%
\pgfusepath{fill}%
\end{pgfscope}%
\begin{pgfscope}%
\pgfpathrectangle{\pgfqpoint{0.407592in}{0.222904in}}{\pgfqpoint{1.411163in}{1.411163in}}%
\pgfusepath{clip}%
\pgfsetbuttcap%
\pgfsetroundjoin%
\definecolor{currentfill}{rgb}{0.168627,0.611765,0.411765}%
\pgfsetfillcolor{currentfill}%
\pgfsetfillopacity{0.180000}%
\pgfsetlinewidth{0.000000pt}%
\definecolor{currentstroke}{rgb}{0.000000,0.000000,0.000000}%
\pgfsetstrokecolor{currentstroke}%
\pgfsetdash{}{0pt}%
\pgfpathmoveto{\pgfqpoint{1.659233in}{1.305647in}}%
\pgfpathlineto{\pgfqpoint{1.673645in}{1.330405in}}%
\pgfpathlineto{\pgfqpoint{1.674299in}{1.352582in}}%
\pgfpathlineto{\pgfqpoint{1.659872in}{1.327841in}}%
\pgfpathlineto{\pgfqpoint{1.659233in}{1.305647in}}%
\pgfpathclose%
\pgfusepath{fill}%
\end{pgfscope}%
\begin{pgfscope}%
\pgfpathrectangle{\pgfqpoint{0.407592in}{0.222904in}}{\pgfqpoint{1.411163in}{1.411163in}}%
\pgfusepath{clip}%
\pgfsetbuttcap%
\pgfsetroundjoin%
\definecolor{currentfill}{rgb}{0.168627,0.611765,0.411765}%
\pgfsetfillcolor{currentfill}%
\pgfsetfillopacity{0.180000}%
\pgfsetlinewidth{0.000000pt}%
\definecolor{currentstroke}{rgb}{0.000000,0.000000,0.000000}%
\pgfsetstrokecolor{currentstroke}%
\pgfsetdash{}{0pt}%
\pgfpathmoveto{\pgfqpoint{1.674955in}{1.374813in}}%
\pgfpathlineto{\pgfqpoint{1.689359in}{1.400140in}}%
\pgfpathlineto{\pgfqpoint{1.690032in}{1.422409in}}%
\pgfpathlineto{\pgfqpoint{1.675612in}{1.397098in}}%
\pgfpathlineto{\pgfqpoint{1.674955in}{1.374813in}}%
\pgfpathclose%
\pgfusepath{fill}%
\end{pgfscope}%
\begin{pgfscope}%
\pgfpathrectangle{\pgfqpoint{0.407592in}{0.222904in}}{\pgfqpoint{1.411163in}{1.411163in}}%
\pgfusepath{clip}%
\pgfsetbuttcap%
\pgfsetroundjoin%
\definecolor{currentfill}{rgb}{0.168627,0.611765,0.411765}%
\pgfsetfillcolor{currentfill}%
\pgfsetfillopacity{0.180000}%
\pgfsetlinewidth{0.000000pt}%
\definecolor{currentstroke}{rgb}{0.000000,0.000000,0.000000}%
\pgfsetstrokecolor{currentstroke}%
\pgfsetdash{}{0pt}%
\pgfpathmoveto{\pgfqpoint{1.597217in}{1.059631in}}%
\pgfpathlineto{\pgfqpoint{1.611674in}{1.081405in}}%
\pgfpathlineto{\pgfqpoint{1.612254in}{1.103288in}}%
\pgfpathlineto{\pgfqpoint{1.597781in}{1.081538in}}%
\pgfpathlineto{\pgfqpoint{1.597217in}{1.059631in}}%
\pgfpathclose%
\pgfusepath{fill}%
\end{pgfscope}%
\begin{pgfscope}%
\pgfpathrectangle{\pgfqpoint{0.407592in}{0.222904in}}{\pgfqpoint{1.411163in}{1.411163in}}%
\pgfusepath{clip}%
\pgfsetbuttcap%
\pgfsetroundjoin%
\definecolor{currentfill}{rgb}{0.168627,0.611765,0.411765}%
\pgfsetfillcolor{currentfill}%
\pgfsetfillopacity{0.180000}%
\pgfsetlinewidth{0.000000pt}%
\definecolor{currentstroke}{rgb}{0.000000,0.000000,0.000000}%
\pgfsetstrokecolor{currentstroke}%
\pgfsetdash{}{0pt}%
\pgfpathmoveto{\pgfqpoint{1.612835in}{1.125224in}}%
\pgfpathlineto{\pgfqpoint{1.627284in}{1.147851in}}%
\pgfpathlineto{\pgfqpoint{1.627883in}{1.169818in}}%
\pgfpathlineto{\pgfqpoint{1.613417in}{1.147213in}}%
\pgfpathlineto{\pgfqpoint{1.612835in}{1.125224in}}%
\pgfpathclose%
\pgfusepath{fill}%
\end{pgfscope}%
\begin{pgfscope}%
\pgfpathrectangle{\pgfqpoint{0.407592in}{0.222904in}}{\pgfqpoint{1.411163in}{1.411163in}}%
\pgfusepath{clip}%
\pgfsetbuttcap%
\pgfsetroundjoin%
\definecolor{currentfill}{rgb}{0.168627,0.611765,0.411765}%
\pgfsetfillcolor{currentfill}%
\pgfsetfillopacity{0.180000}%
\pgfsetlinewidth{0.000000pt}%
\definecolor{currentstroke}{rgb}{0.000000,0.000000,0.000000}%
\pgfsetstrokecolor{currentstroke}%
\pgfsetdash{}{0pt}%
\pgfpathmoveto{\pgfqpoint{1.628483in}{1.191838in}}%
\pgfpathlineto{\pgfqpoint{1.642925in}{1.215241in}}%
\pgfpathlineto{\pgfqpoint{1.643543in}{1.237293in}}%
\pgfpathlineto{\pgfqpoint{1.629084in}{1.213911in}}%
\pgfpathlineto{\pgfqpoint{1.628483in}{1.191838in}}%
\pgfpathclose%
\pgfusepath{fill}%
\end{pgfscope}%
\begin{pgfscope}%
\pgfpathrectangle{\pgfqpoint{0.407592in}{0.222904in}}{\pgfqpoint{1.411163in}{1.411163in}}%
\pgfusepath{clip}%
\pgfsetbuttcap%
\pgfsetroundjoin%
\definecolor{currentfill}{rgb}{0.168627,0.611765,0.411765}%
\pgfsetfillcolor{currentfill}%
\pgfsetfillopacity{0.180000}%
\pgfsetlinewidth{0.000000pt}%
\definecolor{currentstroke}{rgb}{0.000000,0.000000,0.000000}%
\pgfsetstrokecolor{currentstroke}%
\pgfsetdash{}{0pt}%
\pgfpathmoveto{\pgfqpoint{1.644162in}{1.259399in}}%
\pgfpathlineto{\pgfqpoint{1.658597in}{1.283506in}}%
\pgfpathlineto{\pgfqpoint{1.659233in}{1.305647in}}%
\pgfpathlineto{\pgfqpoint{1.644782in}{1.281559in}}%
\pgfpathlineto{\pgfqpoint{1.644162in}{1.259399in}}%
\pgfpathclose%
\pgfusepath{fill}%
\end{pgfscope}%
\begin{pgfscope}%
\pgfpathrectangle{\pgfqpoint{0.407592in}{0.222904in}}{\pgfqpoint{1.411163in}{1.411163in}}%
\pgfusepath{clip}%
\pgfsetbuttcap%
\pgfsetroundjoin%
\definecolor{currentfill}{rgb}{0.168627,0.611765,0.411765}%
\pgfsetfillcolor{currentfill}%
\pgfsetfillopacity{0.180000}%
\pgfsetlinewidth{0.000000pt}%
\definecolor{currentstroke}{rgb}{0.000000,0.000000,0.000000}%
\pgfsetstrokecolor{currentstroke}%
\pgfsetdash{}{0pt}%
\pgfpathmoveto{\pgfqpoint{1.659872in}{1.327841in}}%
\pgfpathlineto{\pgfqpoint{1.674299in}{1.352582in}}%
\pgfpathlineto{\pgfqpoint{1.674955in}{1.374813in}}%
\pgfpathlineto{\pgfqpoint{1.660511in}{1.350090in}}%
\pgfpathlineto{\pgfqpoint{1.659872in}{1.327841in}}%
\pgfpathclose%
\pgfusepath{fill}%
\end{pgfscope}%
\begin{pgfscope}%
\pgfpathrectangle{\pgfqpoint{0.407592in}{0.222904in}}{\pgfqpoint{1.411163in}{1.411163in}}%
\pgfusepath{clip}%
\pgfsetbuttcap%
\pgfsetroundjoin%
\definecolor{currentfill}{rgb}{0.168627,0.611765,0.411765}%
\pgfsetfillcolor{currentfill}%
\pgfsetfillopacity{0.180000}%
\pgfsetlinewidth{0.000000pt}%
\definecolor{currentstroke}{rgb}{0.000000,0.000000,0.000000}%
\pgfsetstrokecolor{currentstroke}%
\pgfsetdash{}{0pt}%
\pgfpathmoveto{\pgfqpoint{1.675612in}{1.397098in}}%
\pgfpathlineto{\pgfqpoint{1.690032in}{1.422409in}}%
\pgfpathlineto{\pgfqpoint{1.690707in}{1.444732in}}%
\pgfpathlineto{\pgfqpoint{1.676271in}{1.419437in}}%
\pgfpathlineto{\pgfqpoint{1.675612in}{1.397098in}}%
\pgfpathclose%
\pgfusepath{fill}%
\end{pgfscope}%
\begin{pgfscope}%
\pgfpathrectangle{\pgfqpoint{0.407592in}{0.222904in}}{\pgfqpoint{1.411163in}{1.411163in}}%
\pgfusepath{clip}%
\pgfsetbuttcap%
\pgfsetroundjoin%
\definecolor{currentfill}{rgb}{0.168627,0.611765,0.411765}%
\pgfsetfillcolor{currentfill}%
\pgfsetfillopacity{0.180000}%
\pgfsetlinewidth{0.000000pt}%
\definecolor{currentstroke}{rgb}{0.000000,0.000000,0.000000}%
\pgfsetstrokecolor{currentstroke}%
\pgfsetdash{}{0pt}%
\pgfpathmoveto{\pgfqpoint{1.582174in}{1.016959in}}%
\pgfpathlineto{\pgfqpoint{1.596655in}{1.037776in}}%
\pgfpathlineto{\pgfqpoint{1.597217in}{1.059631in}}%
\pgfpathlineto{\pgfqpoint{1.582720in}{1.038841in}}%
\pgfpathlineto{\pgfqpoint{1.582174in}{1.016959in}}%
\pgfpathclose%
\pgfusepath{fill}%
\end{pgfscope}%
\begin{pgfscope}%
\pgfpathrectangle{\pgfqpoint{0.407592in}{0.222904in}}{\pgfqpoint{1.411163in}{1.411163in}}%
\pgfusepath{clip}%
\pgfsetbuttcap%
\pgfsetroundjoin%
\definecolor{currentfill}{rgb}{0.168627,0.611765,0.411765}%
\pgfsetfillcolor{currentfill}%
\pgfsetfillopacity{0.180000}%
\pgfsetlinewidth{0.000000pt}%
\definecolor{currentstroke}{rgb}{0.000000,0.000000,0.000000}%
\pgfsetstrokecolor{currentstroke}%
\pgfsetdash{}{0pt}%
\pgfpathmoveto{\pgfqpoint{1.597781in}{1.081538in}}%
\pgfpathlineto{\pgfqpoint{1.612254in}{1.103288in}}%
\pgfpathlineto{\pgfqpoint{1.612835in}{1.125224in}}%
\pgfpathlineto{\pgfqpoint{1.598345in}{1.103499in}}%
\pgfpathlineto{\pgfqpoint{1.597781in}{1.081538in}}%
\pgfpathclose%
\pgfusepath{fill}%
\end{pgfscope}%
\begin{pgfscope}%
\pgfpathrectangle{\pgfqpoint{0.407592in}{0.222904in}}{\pgfqpoint{1.411163in}{1.411163in}}%
\pgfusepath{clip}%
\pgfsetbuttcap%
\pgfsetroundjoin%
\definecolor{currentfill}{rgb}{0.168627,0.611765,0.411765}%
\pgfsetfillcolor{currentfill}%
\pgfsetfillopacity{0.180000}%
\pgfsetlinewidth{0.000000pt}%
\definecolor{currentstroke}{rgb}{0.000000,0.000000,0.000000}%
\pgfsetstrokecolor{currentstroke}%
\pgfsetdash{}{0pt}%
\pgfpathmoveto{\pgfqpoint{1.613417in}{1.147213in}}%
\pgfpathlineto{\pgfqpoint{1.627883in}{1.169818in}}%
\pgfpathlineto{\pgfqpoint{1.628483in}{1.191838in}}%
\pgfpathlineto{\pgfqpoint{1.614001in}{1.169256in}}%
\pgfpathlineto{\pgfqpoint{1.613417in}{1.147213in}}%
\pgfpathclose%
\pgfusepath{fill}%
\end{pgfscope}%
\begin{pgfscope}%
\pgfpathrectangle{\pgfqpoint{0.407592in}{0.222904in}}{\pgfqpoint{1.411163in}{1.411163in}}%
\pgfusepath{clip}%
\pgfsetbuttcap%
\pgfsetroundjoin%
\definecolor{currentfill}{rgb}{0.168627,0.611765,0.411765}%
\pgfsetfillcolor{currentfill}%
\pgfsetfillopacity{0.180000}%
\pgfsetlinewidth{0.000000pt}%
\definecolor{currentstroke}{rgb}{0.000000,0.000000,0.000000}%
\pgfsetstrokecolor{currentstroke}%
\pgfsetdash{}{0pt}%
\pgfpathmoveto{\pgfqpoint{1.629084in}{1.213911in}}%
\pgfpathlineto{\pgfqpoint{1.643543in}{1.237293in}}%
\pgfpathlineto{\pgfqpoint{1.644162in}{1.259399in}}%
\pgfpathlineto{\pgfqpoint{1.629687in}{1.236038in}}%
\pgfpathlineto{\pgfqpoint{1.629084in}{1.213911in}}%
\pgfpathclose%
\pgfusepath{fill}%
\end{pgfscope}%
\begin{pgfscope}%
\pgfpathrectangle{\pgfqpoint{0.407592in}{0.222904in}}{\pgfqpoint{1.411163in}{1.411163in}}%
\pgfusepath{clip}%
\pgfsetbuttcap%
\pgfsetroundjoin%
\definecolor{currentfill}{rgb}{0.168627,0.611765,0.411765}%
\pgfsetfillcolor{currentfill}%
\pgfsetfillopacity{0.180000}%
\pgfsetlinewidth{0.000000pt}%
\definecolor{currentstroke}{rgb}{0.000000,0.000000,0.000000}%
\pgfsetstrokecolor{currentstroke}%
\pgfsetdash{}{0pt}%
\pgfpathmoveto{\pgfqpoint{1.644782in}{1.281559in}}%
\pgfpathlineto{\pgfqpoint{1.659233in}{1.305647in}}%
\pgfpathlineto{\pgfqpoint{1.659872in}{1.327841in}}%
\pgfpathlineto{\pgfqpoint{1.645404in}{1.303773in}}%
\pgfpathlineto{\pgfqpoint{1.644782in}{1.281559in}}%
\pgfpathclose%
\pgfusepath{fill}%
\end{pgfscope}%
\begin{pgfscope}%
\pgfpathrectangle{\pgfqpoint{0.407592in}{0.222904in}}{\pgfqpoint{1.411163in}{1.411163in}}%
\pgfusepath{clip}%
\pgfsetbuttcap%
\pgfsetroundjoin%
\definecolor{currentfill}{rgb}{0.168627,0.611765,0.411765}%
\pgfsetfillcolor{currentfill}%
\pgfsetfillopacity{0.180000}%
\pgfsetlinewidth{0.000000pt}%
\definecolor{currentstroke}{rgb}{0.000000,0.000000,0.000000}%
\pgfsetstrokecolor{currentstroke}%
\pgfsetdash{}{0pt}%
\pgfpathmoveto{\pgfqpoint{1.660511in}{1.350090in}}%
\pgfpathlineto{\pgfqpoint{1.674955in}{1.374813in}}%
\pgfpathlineto{\pgfqpoint{1.675612in}{1.397098in}}%
\pgfpathlineto{\pgfqpoint{1.661153in}{1.372392in}}%
\pgfpathlineto{\pgfqpoint{1.660511in}{1.350090in}}%
\pgfpathclose%
\pgfusepath{fill}%
\end{pgfscope}%
\begin{pgfscope}%
\pgfpathrectangle{\pgfqpoint{0.407592in}{0.222904in}}{\pgfqpoint{1.411163in}{1.411163in}}%
\pgfusepath{clip}%
\pgfsetbuttcap%
\pgfsetroundjoin%
\definecolor{currentfill}{rgb}{0.168627,0.611765,0.411765}%
\pgfsetfillcolor{currentfill}%
\pgfsetfillopacity{0.180000}%
\pgfsetlinewidth{0.000000pt}%
\definecolor{currentstroke}{rgb}{0.000000,0.000000,0.000000}%
\pgfsetstrokecolor{currentstroke}%
\pgfsetdash{}{0pt}%
\pgfpathmoveto{\pgfqpoint{1.676271in}{1.419437in}}%
\pgfpathlineto{\pgfqpoint{1.690707in}{1.444732in}}%
\pgfpathlineto{\pgfqpoint{1.691384in}{1.467109in}}%
\pgfpathlineto{\pgfqpoint{1.676932in}{1.441830in}}%
\pgfpathlineto{\pgfqpoint{1.676271in}{1.419437in}}%
\pgfpathclose%
\pgfusepath{fill}%
\end{pgfscope}%
\begin{pgfscope}%
\pgfpathrectangle{\pgfqpoint{0.407592in}{0.222904in}}{\pgfqpoint{1.411163in}{1.411163in}}%
\pgfusepath{clip}%
\pgfsetbuttcap%
\pgfsetroundjoin%
\definecolor{currentfill}{rgb}{0.168627,0.611765,0.411765}%
\pgfsetfillcolor{currentfill}%
\pgfsetfillopacity{0.180000}%
\pgfsetlinewidth{0.000000pt}%
\definecolor{currentstroke}{rgb}{0.000000,0.000000,0.000000}%
\pgfsetstrokecolor{currentstroke}%
\pgfsetdash{}{0pt}%
\pgfpathmoveto{\pgfqpoint{1.582720in}{1.038841in}}%
\pgfpathlineto{\pgfqpoint{1.597217in}{1.059631in}}%
\pgfpathlineto{\pgfqpoint{1.597781in}{1.081538in}}%
\pgfpathlineto{\pgfqpoint{1.583267in}{1.060775in}}%
\pgfpathlineto{\pgfqpoint{1.582720in}{1.038841in}}%
\pgfpathclose%
\pgfusepath{fill}%
\end{pgfscope}%
\begin{pgfscope}%
\pgfpathrectangle{\pgfqpoint{0.407592in}{0.222904in}}{\pgfqpoint{1.411163in}{1.411163in}}%
\pgfusepath{clip}%
\pgfsetbuttcap%
\pgfsetroundjoin%
\definecolor{currentfill}{rgb}{0.168627,0.611765,0.411765}%
\pgfsetfillcolor{currentfill}%
\pgfsetfillopacity{0.180000}%
\pgfsetlinewidth{0.000000pt}%
\definecolor{currentstroke}{rgb}{0.000000,0.000000,0.000000}%
\pgfsetstrokecolor{currentstroke}%
\pgfsetdash{}{0pt}%
\pgfpathmoveto{\pgfqpoint{1.598345in}{1.103499in}}%
\pgfpathlineto{\pgfqpoint{1.612835in}{1.125224in}}%
\pgfpathlineto{\pgfqpoint{1.613417in}{1.147213in}}%
\pgfpathlineto{\pgfqpoint{1.598911in}{1.125513in}}%
\pgfpathlineto{\pgfqpoint{1.598345in}{1.103499in}}%
\pgfpathclose%
\pgfusepath{fill}%
\end{pgfscope}%
\begin{pgfscope}%
\pgfpathrectangle{\pgfqpoint{0.407592in}{0.222904in}}{\pgfqpoint{1.411163in}{1.411163in}}%
\pgfusepath{clip}%
\pgfsetbuttcap%
\pgfsetroundjoin%
\definecolor{currentfill}{rgb}{0.168627,0.611765,0.411765}%
\pgfsetfillcolor{currentfill}%
\pgfsetfillopacity{0.180000}%
\pgfsetlinewidth{0.000000pt}%
\definecolor{currentstroke}{rgb}{0.000000,0.000000,0.000000}%
\pgfsetstrokecolor{currentstroke}%
\pgfsetdash{}{0pt}%
\pgfpathmoveto{\pgfqpoint{1.614001in}{1.169256in}}%
\pgfpathlineto{\pgfqpoint{1.628483in}{1.191838in}}%
\pgfpathlineto{\pgfqpoint{1.629084in}{1.213911in}}%
\pgfpathlineto{\pgfqpoint{1.614586in}{1.191352in}}%
\pgfpathlineto{\pgfqpoint{1.614001in}{1.169256in}}%
\pgfpathclose%
\pgfusepath{fill}%
\end{pgfscope}%
\begin{pgfscope}%
\pgfpathrectangle{\pgfqpoint{0.407592in}{0.222904in}}{\pgfqpoint{1.411163in}{1.411163in}}%
\pgfusepath{clip}%
\pgfsetbuttcap%
\pgfsetroundjoin%
\definecolor{currentfill}{rgb}{0.168627,0.611765,0.411765}%
\pgfsetfillcolor{currentfill}%
\pgfsetfillopacity{0.180000}%
\pgfsetlinewidth{0.000000pt}%
\definecolor{currentstroke}{rgb}{0.000000,0.000000,0.000000}%
\pgfsetstrokecolor{currentstroke}%
\pgfsetdash{}{0pt}%
\pgfpathmoveto{\pgfqpoint{1.629687in}{1.236038in}}%
\pgfpathlineto{\pgfqpoint{1.644162in}{1.259399in}}%
\pgfpathlineto{\pgfqpoint{1.644782in}{1.281559in}}%
\pgfpathlineto{\pgfqpoint{1.630291in}{1.258218in}}%
\pgfpathlineto{\pgfqpoint{1.629687in}{1.236038in}}%
\pgfpathclose%
\pgfusepath{fill}%
\end{pgfscope}%
\begin{pgfscope}%
\pgfpathrectangle{\pgfqpoint{0.407592in}{0.222904in}}{\pgfqpoint{1.411163in}{1.411163in}}%
\pgfusepath{clip}%
\pgfsetbuttcap%
\pgfsetroundjoin%
\definecolor{currentfill}{rgb}{0.168627,0.611765,0.411765}%
\pgfsetfillcolor{currentfill}%
\pgfsetfillopacity{0.180000}%
\pgfsetlinewidth{0.000000pt}%
\definecolor{currentstroke}{rgb}{0.000000,0.000000,0.000000}%
\pgfsetstrokecolor{currentstroke}%
\pgfsetdash{}{0pt}%
\pgfpathmoveto{\pgfqpoint{1.645404in}{1.303773in}}%
\pgfpathlineto{\pgfqpoint{1.659872in}{1.327841in}}%
\pgfpathlineto{\pgfqpoint{1.660511in}{1.350090in}}%
\pgfpathlineto{\pgfqpoint{1.646028in}{1.326040in}}%
\pgfpathlineto{\pgfqpoint{1.645404in}{1.303773in}}%
\pgfpathclose%
\pgfusepath{fill}%
\end{pgfscope}%
\begin{pgfscope}%
\pgfpathrectangle{\pgfqpoint{0.407592in}{0.222904in}}{\pgfqpoint{1.411163in}{1.411163in}}%
\pgfusepath{clip}%
\pgfsetbuttcap%
\pgfsetroundjoin%
\definecolor{currentfill}{rgb}{0.168627,0.611765,0.411765}%
\pgfsetfillcolor{currentfill}%
\pgfsetfillopacity{0.180000}%
\pgfsetlinewidth{0.000000pt}%
\definecolor{currentstroke}{rgb}{0.000000,0.000000,0.000000}%
\pgfsetstrokecolor{currentstroke}%
\pgfsetdash{}{0pt}%
\pgfpathmoveto{\pgfqpoint{1.661153in}{1.372392in}}%
\pgfpathlineto{\pgfqpoint{1.675612in}{1.397098in}}%
\pgfpathlineto{\pgfqpoint{1.676271in}{1.419437in}}%
\pgfpathlineto{\pgfqpoint{1.661795in}{1.394749in}}%
\pgfpathlineto{\pgfqpoint{1.661153in}{1.372392in}}%
\pgfpathclose%
\pgfusepath{fill}%
\end{pgfscope}%
\begin{pgfscope}%
\pgfpathrectangle{\pgfqpoint{0.407592in}{0.222904in}}{\pgfqpoint{1.411163in}{1.411163in}}%
\pgfusepath{clip}%
\pgfsetbuttcap%
\pgfsetroundjoin%
\definecolor{currentfill}{rgb}{0.168627,0.611765,0.411765}%
\pgfsetfillcolor{currentfill}%
\pgfsetfillopacity{0.180000}%
\pgfsetlinewidth{0.000000pt}%
\definecolor{currentstroke}{rgb}{0.000000,0.000000,0.000000}%
\pgfsetstrokecolor{currentstroke}%
\pgfsetdash{}{0pt}%
\pgfpathmoveto{\pgfqpoint{1.567653in}{0.997214in}}%
\pgfpathlineto{\pgfqpoint{1.582174in}{1.016959in}}%
\pgfpathlineto{\pgfqpoint{1.582720in}{1.038841in}}%
\pgfpathlineto{\pgfqpoint{1.568182in}{1.019125in}}%
\pgfpathlineto{\pgfqpoint{1.567653in}{0.997214in}}%
\pgfpathclose%
\pgfusepath{fill}%
\end{pgfscope}%
\begin{pgfscope}%
\pgfpathrectangle{\pgfqpoint{0.407592in}{0.222904in}}{\pgfqpoint{1.411163in}{1.411163in}}%
\pgfusepath{clip}%
\pgfsetbuttcap%
\pgfsetroundjoin%
\definecolor{currentfill}{rgb}{0.168627,0.611765,0.411765}%
\pgfsetfillcolor{currentfill}%
\pgfsetfillopacity{0.180000}%
\pgfsetlinewidth{0.000000pt}%
\definecolor{currentstroke}{rgb}{0.000000,0.000000,0.000000}%
\pgfsetstrokecolor{currentstroke}%
\pgfsetdash{}{0pt}%
\pgfpathmoveto{\pgfqpoint{1.583267in}{1.060775in}}%
\pgfpathlineto{\pgfqpoint{1.597781in}{1.081538in}}%
\pgfpathlineto{\pgfqpoint{1.598345in}{1.103499in}}%
\pgfpathlineto{\pgfqpoint{1.583816in}{1.082763in}}%
\pgfpathlineto{\pgfqpoint{1.583267in}{1.060775in}}%
\pgfpathclose%
\pgfusepath{fill}%
\end{pgfscope}%
\begin{pgfscope}%
\pgfpathrectangle{\pgfqpoint{0.407592in}{0.222904in}}{\pgfqpoint{1.411163in}{1.411163in}}%
\pgfusepath{clip}%
\pgfsetbuttcap%
\pgfsetroundjoin%
\definecolor{currentfill}{rgb}{0.168627,0.611765,0.411765}%
\pgfsetfillcolor{currentfill}%
\pgfsetfillopacity{0.180000}%
\pgfsetlinewidth{0.000000pt}%
\definecolor{currentstroke}{rgb}{0.000000,0.000000,0.000000}%
\pgfsetstrokecolor{currentstroke}%
\pgfsetdash{}{0pt}%
\pgfpathmoveto{\pgfqpoint{1.598911in}{1.125513in}}%
\pgfpathlineto{\pgfqpoint{1.613417in}{1.147213in}}%
\pgfpathlineto{\pgfqpoint{1.614001in}{1.169256in}}%
\pgfpathlineto{\pgfqpoint{1.599479in}{1.147580in}}%
\pgfpathlineto{\pgfqpoint{1.598911in}{1.125513in}}%
\pgfpathclose%
\pgfusepath{fill}%
\end{pgfscope}%
\begin{pgfscope}%
\pgfpathrectangle{\pgfqpoint{0.407592in}{0.222904in}}{\pgfqpoint{1.411163in}{1.411163in}}%
\pgfusepath{clip}%
\pgfsetbuttcap%
\pgfsetroundjoin%
\definecolor{currentfill}{rgb}{0.168627,0.611765,0.411765}%
\pgfsetfillcolor{currentfill}%
\pgfsetfillopacity{0.180000}%
\pgfsetlinewidth{0.000000pt}%
\definecolor{currentstroke}{rgb}{0.000000,0.000000,0.000000}%
\pgfsetstrokecolor{currentstroke}%
\pgfsetdash{}{0pt}%
\pgfpathmoveto{\pgfqpoint{1.614586in}{1.191352in}}%
\pgfpathlineto{\pgfqpoint{1.629084in}{1.213911in}}%
\pgfpathlineto{\pgfqpoint{1.629687in}{1.236038in}}%
\pgfpathlineto{\pgfqpoint{1.615173in}{1.213502in}}%
\pgfpathlineto{\pgfqpoint{1.614586in}{1.191352in}}%
\pgfpathclose%
\pgfusepath{fill}%
\end{pgfscope}%
\begin{pgfscope}%
\pgfpathrectangle{\pgfqpoint{0.407592in}{0.222904in}}{\pgfqpoint{1.411163in}{1.411163in}}%
\pgfusepath{clip}%
\pgfsetbuttcap%
\pgfsetroundjoin%
\definecolor{currentfill}{rgb}{0.168627,0.611765,0.411765}%
\pgfsetfillcolor{currentfill}%
\pgfsetfillopacity{0.180000}%
\pgfsetlinewidth{0.000000pt}%
\definecolor{currentstroke}{rgb}{0.000000,0.000000,0.000000}%
\pgfsetstrokecolor{currentstroke}%
\pgfsetdash{}{0pt}%
\pgfpathmoveto{\pgfqpoint{1.630291in}{1.258218in}}%
\pgfpathlineto{\pgfqpoint{1.644782in}{1.281559in}}%
\pgfpathlineto{\pgfqpoint{1.645404in}{1.303773in}}%
\pgfpathlineto{\pgfqpoint{1.630897in}{1.280453in}}%
\pgfpathlineto{\pgfqpoint{1.630291in}{1.258218in}}%
\pgfpathclose%
\pgfusepath{fill}%
\end{pgfscope}%
\begin{pgfscope}%
\pgfpathrectangle{\pgfqpoint{0.407592in}{0.222904in}}{\pgfqpoint{1.411163in}{1.411163in}}%
\pgfusepath{clip}%
\pgfsetbuttcap%
\pgfsetroundjoin%
\definecolor{currentfill}{rgb}{0.168627,0.611765,0.411765}%
\pgfsetfillcolor{currentfill}%
\pgfsetfillopacity{0.180000}%
\pgfsetlinewidth{0.000000pt}%
\definecolor{currentstroke}{rgb}{0.000000,0.000000,0.000000}%
\pgfsetstrokecolor{currentstroke}%
\pgfsetdash{}{0pt}%
\pgfpathmoveto{\pgfqpoint{1.646028in}{1.326040in}}%
\pgfpathlineto{\pgfqpoint{1.660511in}{1.350090in}}%
\pgfpathlineto{\pgfqpoint{1.661153in}{1.372392in}}%
\pgfpathlineto{\pgfqpoint{1.646653in}{1.348362in}}%
\pgfpathlineto{\pgfqpoint{1.646028in}{1.326040in}}%
\pgfpathclose%
\pgfusepath{fill}%
\end{pgfscope}%
\begin{pgfscope}%
\pgfpathrectangle{\pgfqpoint{0.407592in}{0.222904in}}{\pgfqpoint{1.411163in}{1.411163in}}%
\pgfusepath{clip}%
\pgfsetbuttcap%
\pgfsetroundjoin%
\definecolor{currentfill}{rgb}{0.168627,0.611765,0.411765}%
\pgfsetfillcolor{currentfill}%
\pgfsetfillopacity{0.180000}%
\pgfsetlinewidth{0.000000pt}%
\definecolor{currentstroke}{rgb}{0.000000,0.000000,0.000000}%
\pgfsetstrokecolor{currentstroke}%
\pgfsetdash{}{0pt}%
\pgfpathmoveto{\pgfqpoint{1.661795in}{1.394749in}}%
\pgfpathlineto{\pgfqpoint{1.676271in}{1.419437in}}%
\pgfpathlineto{\pgfqpoint{1.676932in}{1.441830in}}%
\pgfpathlineto{\pgfqpoint{1.662440in}{1.417160in}}%
\pgfpathlineto{\pgfqpoint{1.661795in}{1.394749in}}%
\pgfpathclose%
\pgfusepath{fill}%
\end{pgfscope}%
\begin{pgfscope}%
\pgfpathrectangle{\pgfqpoint{0.407592in}{0.222904in}}{\pgfqpoint{1.411163in}{1.411163in}}%
\pgfusepath{clip}%
\pgfsetbuttcap%
\pgfsetroundjoin%
\definecolor{currentfill}{rgb}{0.168627,0.611765,0.411765}%
\pgfsetfillcolor{currentfill}%
\pgfsetfillopacity{0.180000}%
\pgfsetlinewidth{0.000000pt}%
\definecolor{currentstroke}{rgb}{0.000000,0.000000,0.000000}%
\pgfsetstrokecolor{currentstroke}%
\pgfsetdash{}{0pt}%
\pgfpathmoveto{\pgfqpoint{1.568182in}{1.019125in}}%
\pgfpathlineto{\pgfqpoint{1.582720in}{1.038841in}}%
\pgfpathlineto{\pgfqpoint{1.583267in}{1.060775in}}%
\pgfpathlineto{\pgfqpoint{1.568713in}{1.041090in}}%
\pgfpathlineto{\pgfqpoint{1.568182in}{1.019125in}}%
\pgfpathclose%
\pgfusepath{fill}%
\end{pgfscope}%
\begin{pgfscope}%
\pgfpathrectangle{\pgfqpoint{0.407592in}{0.222904in}}{\pgfqpoint{1.411163in}{1.411163in}}%
\pgfusepath{clip}%
\pgfsetbuttcap%
\pgfsetroundjoin%
\definecolor{currentfill}{rgb}{0.168627,0.611765,0.411765}%
\pgfsetfillcolor{currentfill}%
\pgfsetfillopacity{0.180000}%
\pgfsetlinewidth{0.000000pt}%
\definecolor{currentstroke}{rgb}{0.000000,0.000000,0.000000}%
\pgfsetstrokecolor{currentstroke}%
\pgfsetdash{}{0pt}%
\pgfpathmoveto{\pgfqpoint{1.583816in}{1.082763in}}%
\pgfpathlineto{\pgfqpoint{1.598345in}{1.103499in}}%
\pgfpathlineto{\pgfqpoint{1.598911in}{1.125513in}}%
\pgfpathlineto{\pgfqpoint{1.584365in}{1.104805in}}%
\pgfpathlineto{\pgfqpoint{1.583816in}{1.082763in}}%
\pgfpathclose%
\pgfusepath{fill}%
\end{pgfscope}%
\begin{pgfscope}%
\pgfpathrectangle{\pgfqpoint{0.407592in}{0.222904in}}{\pgfqpoint{1.411163in}{1.411163in}}%
\pgfusepath{clip}%
\pgfsetbuttcap%
\pgfsetroundjoin%
\definecolor{currentfill}{rgb}{0.168627,0.611765,0.411765}%
\pgfsetfillcolor{currentfill}%
\pgfsetfillopacity{0.180000}%
\pgfsetlinewidth{0.000000pt}%
\definecolor{currentstroke}{rgb}{0.000000,0.000000,0.000000}%
\pgfsetstrokecolor{currentstroke}%
\pgfsetdash{}{0pt}%
\pgfpathmoveto{\pgfqpoint{1.599479in}{1.147580in}}%
\pgfpathlineto{\pgfqpoint{1.614001in}{1.169256in}}%
\pgfpathlineto{\pgfqpoint{1.614586in}{1.191352in}}%
\pgfpathlineto{\pgfqpoint{1.600048in}{1.169701in}}%
\pgfpathlineto{\pgfqpoint{1.599479in}{1.147580in}}%
\pgfpathclose%
\pgfusepath{fill}%
\end{pgfscope}%
\begin{pgfscope}%
\pgfpathrectangle{\pgfqpoint{0.407592in}{0.222904in}}{\pgfqpoint{1.411163in}{1.411163in}}%
\pgfusepath{clip}%
\pgfsetbuttcap%
\pgfsetroundjoin%
\definecolor{currentfill}{rgb}{0.168627,0.611765,0.411765}%
\pgfsetfillcolor{currentfill}%
\pgfsetfillopacity{0.180000}%
\pgfsetlinewidth{0.000000pt}%
\definecolor{currentstroke}{rgb}{0.000000,0.000000,0.000000}%
\pgfsetstrokecolor{currentstroke}%
\pgfsetdash{}{0pt}%
\pgfpathmoveto{\pgfqpoint{1.615173in}{1.213502in}}%
\pgfpathlineto{\pgfqpoint{1.629687in}{1.236038in}}%
\pgfpathlineto{\pgfqpoint{1.630291in}{1.258218in}}%
\pgfpathlineto{\pgfqpoint{1.615761in}{1.235705in}}%
\pgfpathlineto{\pgfqpoint{1.615173in}{1.213502in}}%
\pgfpathclose%
\pgfusepath{fill}%
\end{pgfscope}%
\begin{pgfscope}%
\pgfpathrectangle{\pgfqpoint{0.407592in}{0.222904in}}{\pgfqpoint{1.411163in}{1.411163in}}%
\pgfusepath{clip}%
\pgfsetbuttcap%
\pgfsetroundjoin%
\definecolor{currentfill}{rgb}{0.168627,0.611765,0.411765}%
\pgfsetfillcolor{currentfill}%
\pgfsetfillopacity{0.180000}%
\pgfsetlinewidth{0.000000pt}%
\definecolor{currentstroke}{rgb}{0.000000,0.000000,0.000000}%
\pgfsetstrokecolor{currentstroke}%
\pgfsetdash{}{0pt}%
\pgfpathmoveto{\pgfqpoint{1.630897in}{1.280453in}}%
\pgfpathlineto{\pgfqpoint{1.645404in}{1.303773in}}%
\pgfpathlineto{\pgfqpoint{1.646028in}{1.326040in}}%
\pgfpathlineto{\pgfqpoint{1.631504in}{1.302742in}}%
\pgfpathlineto{\pgfqpoint{1.630897in}{1.280453in}}%
\pgfpathclose%
\pgfusepath{fill}%
\end{pgfscope}%
\begin{pgfscope}%
\pgfpathrectangle{\pgfqpoint{0.407592in}{0.222904in}}{\pgfqpoint{1.411163in}{1.411163in}}%
\pgfusepath{clip}%
\pgfsetbuttcap%
\pgfsetroundjoin%
\definecolor{currentfill}{rgb}{0.168627,0.611765,0.411765}%
\pgfsetfillcolor{currentfill}%
\pgfsetfillopacity{0.180000}%
\pgfsetlinewidth{0.000000pt}%
\definecolor{currentstroke}{rgb}{0.000000,0.000000,0.000000}%
\pgfsetstrokecolor{currentstroke}%
\pgfsetdash{}{0pt}%
\pgfpathmoveto{\pgfqpoint{1.646653in}{1.348362in}}%
\pgfpathlineto{\pgfqpoint{1.661153in}{1.372392in}}%
\pgfpathlineto{\pgfqpoint{1.661795in}{1.394749in}}%
\pgfpathlineto{\pgfqpoint{1.647279in}{1.370738in}}%
\pgfpathlineto{\pgfqpoint{1.646653in}{1.348362in}}%
\pgfpathclose%
\pgfusepath{fill}%
\end{pgfscope}%
\begin{pgfscope}%
\pgfpathrectangle{\pgfqpoint{0.407592in}{0.222904in}}{\pgfqpoint{1.411163in}{1.411163in}}%
\pgfusepath{clip}%
\pgfsetbuttcap%
\pgfsetroundjoin%
\definecolor{currentfill}{rgb}{0.168627,0.611765,0.411765}%
\pgfsetfillcolor{currentfill}%
\pgfsetfillopacity{0.180000}%
\pgfsetlinewidth{0.000000pt}%
\definecolor{currentstroke}{rgb}{0.000000,0.000000,0.000000}%
\pgfsetstrokecolor{currentstroke}%
\pgfsetdash{}{0pt}%
\pgfpathmoveto{\pgfqpoint{1.568713in}{1.041090in}}%
\pgfpathlineto{\pgfqpoint{1.583267in}{1.060775in}}%
\pgfpathlineto{\pgfqpoint{1.583816in}{1.082763in}}%
\pgfpathlineto{\pgfqpoint{1.569245in}{1.063108in}}%
\pgfpathlineto{\pgfqpoint{1.568713in}{1.041090in}}%
\pgfpathclose%
\pgfusepath{fill}%
\end{pgfscope}%
\begin{pgfscope}%
\pgfpathrectangle{\pgfqpoint{0.407592in}{0.222904in}}{\pgfqpoint{1.411163in}{1.411163in}}%
\pgfusepath{clip}%
\pgfsetbuttcap%
\pgfsetroundjoin%
\definecolor{currentfill}{rgb}{0.168627,0.611765,0.411765}%
\pgfsetfillcolor{currentfill}%
\pgfsetfillopacity{0.180000}%
\pgfsetlinewidth{0.000000pt}%
\definecolor{currentstroke}{rgb}{0.000000,0.000000,0.000000}%
\pgfsetstrokecolor{currentstroke}%
\pgfsetdash{}{0pt}%
\pgfpathmoveto{\pgfqpoint{1.584365in}{1.104805in}}%
\pgfpathlineto{\pgfqpoint{1.598911in}{1.125513in}}%
\pgfpathlineto{\pgfqpoint{1.599479in}{1.147580in}}%
\pgfpathlineto{\pgfqpoint{1.584916in}{1.126900in}}%
\pgfpathlineto{\pgfqpoint{1.584365in}{1.104805in}}%
\pgfpathclose%
\pgfusepath{fill}%
\end{pgfscope}%
\begin{pgfscope}%
\pgfpathrectangle{\pgfqpoint{0.407592in}{0.222904in}}{\pgfqpoint{1.411163in}{1.411163in}}%
\pgfusepath{clip}%
\pgfsetbuttcap%
\pgfsetroundjoin%
\definecolor{currentfill}{rgb}{0.168627,0.611765,0.411765}%
\pgfsetfillcolor{currentfill}%
\pgfsetfillopacity{0.180000}%
\pgfsetlinewidth{0.000000pt}%
\definecolor{currentstroke}{rgb}{0.000000,0.000000,0.000000}%
\pgfsetstrokecolor{currentstroke}%
\pgfsetdash{}{0pt}%
\pgfpathmoveto{\pgfqpoint{1.553089in}{0.978630in}}%
\pgfpathlineto{\pgfqpoint{1.567653in}{0.997214in}}%
\pgfpathlineto{\pgfqpoint{1.568182in}{1.019125in}}%
\pgfpathlineto{\pgfqpoint{1.553602in}{1.000574in}}%
\pgfpathlineto{\pgfqpoint{1.553089in}{0.978630in}}%
\pgfpathclose%
\pgfusepath{fill}%
\end{pgfscope}%
\begin{pgfscope}%
\pgfpathrectangle{\pgfqpoint{0.407592in}{0.222904in}}{\pgfqpoint{1.411163in}{1.411163in}}%
\pgfusepath{clip}%
\pgfsetbuttcap%
\pgfsetroundjoin%
\definecolor{currentfill}{rgb}{0.168627,0.611765,0.411765}%
\pgfsetfillcolor{currentfill}%
\pgfsetfillopacity{0.180000}%
\pgfsetlinewidth{0.000000pt}%
\definecolor{currentstroke}{rgb}{0.000000,0.000000,0.000000}%
\pgfsetstrokecolor{currentstroke}%
\pgfsetdash{}{0pt}%
\pgfpathmoveto{\pgfqpoint{1.600048in}{1.169701in}}%
\pgfpathlineto{\pgfqpoint{1.614586in}{1.191352in}}%
\pgfpathlineto{\pgfqpoint{1.615173in}{1.213502in}}%
\pgfpathlineto{\pgfqpoint{1.600618in}{1.191877in}}%
\pgfpathlineto{\pgfqpoint{1.600048in}{1.169701in}}%
\pgfpathclose%
\pgfusepath{fill}%
\end{pgfscope}%
\begin{pgfscope}%
\pgfpathrectangle{\pgfqpoint{0.407592in}{0.222904in}}{\pgfqpoint{1.411163in}{1.411163in}}%
\pgfusepath{clip}%
\pgfsetbuttcap%
\pgfsetroundjoin%
\definecolor{currentfill}{rgb}{0.168627,0.611765,0.411765}%
\pgfsetfillcolor{currentfill}%
\pgfsetfillopacity{0.180000}%
\pgfsetlinewidth{0.000000pt}%
\definecolor{currentstroke}{rgb}{0.000000,0.000000,0.000000}%
\pgfsetstrokecolor{currentstroke}%
\pgfsetdash{}{0pt}%
\pgfpathmoveto{\pgfqpoint{1.615761in}{1.235705in}}%
\pgfpathlineto{\pgfqpoint{1.630291in}{1.258218in}}%
\pgfpathlineto{\pgfqpoint{1.630897in}{1.280453in}}%
\pgfpathlineto{\pgfqpoint{1.616350in}{1.257963in}}%
\pgfpathlineto{\pgfqpoint{1.615761in}{1.235705in}}%
\pgfpathclose%
\pgfusepath{fill}%
\end{pgfscope}%
\begin{pgfscope}%
\pgfpathrectangle{\pgfqpoint{0.407592in}{0.222904in}}{\pgfqpoint{1.411163in}{1.411163in}}%
\pgfusepath{clip}%
\pgfsetbuttcap%
\pgfsetroundjoin%
\definecolor{currentfill}{rgb}{0.168627,0.611765,0.411765}%
\pgfsetfillcolor{currentfill}%
\pgfsetfillopacity{0.180000}%
\pgfsetlinewidth{0.000000pt}%
\definecolor{currentstroke}{rgb}{0.000000,0.000000,0.000000}%
\pgfsetstrokecolor{currentstroke}%
\pgfsetdash{}{0pt}%
\pgfpathmoveto{\pgfqpoint{1.631504in}{1.302742in}}%
\pgfpathlineto{\pgfqpoint{1.646028in}{1.326040in}}%
\pgfpathlineto{\pgfqpoint{1.646653in}{1.348362in}}%
\pgfpathlineto{\pgfqpoint{1.632113in}{1.325085in}}%
\pgfpathlineto{\pgfqpoint{1.631504in}{1.302742in}}%
\pgfpathclose%
\pgfusepath{fill}%
\end{pgfscope}%
\begin{pgfscope}%
\pgfpathrectangle{\pgfqpoint{0.407592in}{0.222904in}}{\pgfqpoint{1.411163in}{1.411163in}}%
\pgfusepath{clip}%
\pgfsetbuttcap%
\pgfsetroundjoin%
\definecolor{currentfill}{rgb}{0.168627,0.611765,0.411765}%
\pgfsetfillcolor{currentfill}%
\pgfsetfillopacity{0.180000}%
\pgfsetlinewidth{0.000000pt}%
\definecolor{currentstroke}{rgb}{0.000000,0.000000,0.000000}%
\pgfsetstrokecolor{currentstroke}%
\pgfsetdash{}{0pt}%
\pgfpathmoveto{\pgfqpoint{1.647279in}{1.370738in}}%
\pgfpathlineto{\pgfqpoint{1.661795in}{1.394749in}}%
\pgfpathlineto{\pgfqpoint{1.662440in}{1.417160in}}%
\pgfpathlineto{\pgfqpoint{1.647907in}{1.393169in}}%
\pgfpathlineto{\pgfqpoint{1.647279in}{1.370738in}}%
\pgfpathclose%
\pgfusepath{fill}%
\end{pgfscope}%
\begin{pgfscope}%
\pgfpathrectangle{\pgfqpoint{0.407592in}{0.222904in}}{\pgfqpoint{1.411163in}{1.411163in}}%
\pgfusepath{clip}%
\pgfsetbuttcap%
\pgfsetroundjoin%
\definecolor{currentfill}{rgb}{0.168627,0.611765,0.411765}%
\pgfsetfillcolor{currentfill}%
\pgfsetfillopacity{0.180000}%
\pgfsetlinewidth{0.000000pt}%
\definecolor{currentstroke}{rgb}{0.000000,0.000000,0.000000}%
\pgfsetstrokecolor{currentstroke}%
\pgfsetdash{}{0pt}%
\pgfpathmoveto{\pgfqpoint{1.569245in}{1.063108in}}%
\pgfpathlineto{\pgfqpoint{1.583816in}{1.082763in}}%
\pgfpathlineto{\pgfqpoint{1.584365in}{1.104805in}}%
\pgfpathlineto{\pgfqpoint{1.569778in}{1.085180in}}%
\pgfpathlineto{\pgfqpoint{1.569245in}{1.063108in}}%
\pgfpathclose%
\pgfusepath{fill}%
\end{pgfscope}%
\begin{pgfscope}%
\pgfpathrectangle{\pgfqpoint{0.407592in}{0.222904in}}{\pgfqpoint{1.411163in}{1.411163in}}%
\pgfusepath{clip}%
\pgfsetbuttcap%
\pgfsetroundjoin%
\definecolor{currentfill}{rgb}{0.168627,0.611765,0.411765}%
\pgfsetfillcolor{currentfill}%
\pgfsetfillopacity{0.180000}%
\pgfsetlinewidth{0.000000pt}%
\definecolor{currentstroke}{rgb}{0.000000,0.000000,0.000000}%
\pgfsetstrokecolor{currentstroke}%
\pgfsetdash{}{0pt}%
\pgfpathmoveto{\pgfqpoint{1.584916in}{1.126900in}}%
\pgfpathlineto{\pgfqpoint{1.599479in}{1.147580in}}%
\pgfpathlineto{\pgfqpoint{1.600048in}{1.169701in}}%
\pgfpathlineto{\pgfqpoint{1.585469in}{1.149049in}}%
\pgfpathlineto{\pgfqpoint{1.584916in}{1.126900in}}%
\pgfpathclose%
\pgfusepath{fill}%
\end{pgfscope}%
\begin{pgfscope}%
\pgfpathrectangle{\pgfqpoint{0.407592in}{0.222904in}}{\pgfqpoint{1.411163in}{1.411163in}}%
\pgfusepath{clip}%
\pgfsetbuttcap%
\pgfsetroundjoin%
\definecolor{currentfill}{rgb}{0.168627,0.611765,0.411765}%
\pgfsetfillcolor{currentfill}%
\pgfsetfillopacity{0.180000}%
\pgfsetlinewidth{0.000000pt}%
\definecolor{currentstroke}{rgb}{0.000000,0.000000,0.000000}%
\pgfsetstrokecolor{currentstroke}%
\pgfsetdash{}{0pt}%
\pgfpathmoveto{\pgfqpoint{1.553602in}{1.000574in}}%
\pgfpathlineto{\pgfqpoint{1.568182in}{1.019125in}}%
\pgfpathlineto{\pgfqpoint{1.568713in}{1.041090in}}%
\pgfpathlineto{\pgfqpoint{1.554116in}{1.022572in}}%
\pgfpathlineto{\pgfqpoint{1.553602in}{1.000574in}}%
\pgfpathclose%
\pgfusepath{fill}%
\end{pgfscope}%
\begin{pgfscope}%
\pgfpathrectangle{\pgfqpoint{0.407592in}{0.222904in}}{\pgfqpoint{1.411163in}{1.411163in}}%
\pgfusepath{clip}%
\pgfsetbuttcap%
\pgfsetroundjoin%
\definecolor{currentfill}{rgb}{0.168627,0.611765,0.411765}%
\pgfsetfillcolor{currentfill}%
\pgfsetfillopacity{0.180000}%
\pgfsetlinewidth{0.000000pt}%
\definecolor{currentstroke}{rgb}{0.000000,0.000000,0.000000}%
\pgfsetstrokecolor{currentstroke}%
\pgfsetdash{}{0pt}%
\pgfpathmoveto{\pgfqpoint{1.600618in}{1.191877in}}%
\pgfpathlineto{\pgfqpoint{1.615173in}{1.213502in}}%
\pgfpathlineto{\pgfqpoint{1.615761in}{1.235705in}}%
\pgfpathlineto{\pgfqpoint{1.601190in}{1.214106in}}%
\pgfpathlineto{\pgfqpoint{1.600618in}{1.191877in}}%
\pgfpathclose%
\pgfusepath{fill}%
\end{pgfscope}%
\begin{pgfscope}%
\pgfpathrectangle{\pgfqpoint{0.407592in}{0.222904in}}{\pgfqpoint{1.411163in}{1.411163in}}%
\pgfusepath{clip}%
\pgfsetbuttcap%
\pgfsetroundjoin%
\definecolor{currentfill}{rgb}{0.168627,0.611765,0.411765}%
\pgfsetfillcolor{currentfill}%
\pgfsetfillopacity{0.180000}%
\pgfsetlinewidth{0.000000pt}%
\definecolor{currentstroke}{rgb}{0.000000,0.000000,0.000000}%
\pgfsetstrokecolor{currentstroke}%
\pgfsetdash{}{0pt}%
\pgfpathmoveto{\pgfqpoint{1.616350in}{1.257963in}}%
\pgfpathlineto{\pgfqpoint{1.630897in}{1.280453in}}%
\pgfpathlineto{\pgfqpoint{1.631504in}{1.302742in}}%
\pgfpathlineto{\pgfqpoint{1.616941in}{1.280275in}}%
\pgfpathlineto{\pgfqpoint{1.616350in}{1.257963in}}%
\pgfpathclose%
\pgfusepath{fill}%
\end{pgfscope}%
\begin{pgfscope}%
\pgfpathrectangle{\pgfqpoint{0.407592in}{0.222904in}}{\pgfqpoint{1.411163in}{1.411163in}}%
\pgfusepath{clip}%
\pgfsetbuttcap%
\pgfsetroundjoin%
\definecolor{currentfill}{rgb}{0.168627,0.611765,0.411765}%
\pgfsetfillcolor{currentfill}%
\pgfsetfillopacity{0.180000}%
\pgfsetlinewidth{0.000000pt}%
\definecolor{currentstroke}{rgb}{0.000000,0.000000,0.000000}%
\pgfsetstrokecolor{currentstroke}%
\pgfsetdash{}{0pt}%
\pgfpathmoveto{\pgfqpoint{1.632113in}{1.325085in}}%
\pgfpathlineto{\pgfqpoint{1.646653in}{1.348362in}}%
\pgfpathlineto{\pgfqpoint{1.647279in}{1.370738in}}%
\pgfpathlineto{\pgfqpoint{1.632723in}{1.347483in}}%
\pgfpathlineto{\pgfqpoint{1.632113in}{1.325085in}}%
\pgfpathclose%
\pgfusepath{fill}%
\end{pgfscope}%
\begin{pgfscope}%
\pgfpathrectangle{\pgfqpoint{0.407592in}{0.222904in}}{\pgfqpoint{1.411163in}{1.411163in}}%
\pgfusepath{clip}%
\pgfsetbuttcap%
\pgfsetroundjoin%
\definecolor{currentfill}{rgb}{0.168627,0.611765,0.411765}%
\pgfsetfillcolor{currentfill}%
\pgfsetfillopacity{0.180000}%
\pgfsetlinewidth{0.000000pt}%
\definecolor{currentstroke}{rgb}{0.000000,0.000000,0.000000}%
\pgfsetstrokecolor{currentstroke}%
\pgfsetdash{}{0pt}%
\pgfpathmoveto{\pgfqpoint{1.569778in}{1.085180in}}%
\pgfpathlineto{\pgfqpoint{1.584365in}{1.104805in}}%
\pgfpathlineto{\pgfqpoint{1.584916in}{1.126900in}}%
\pgfpathlineto{\pgfqpoint{1.570313in}{1.107305in}}%
\pgfpathlineto{\pgfqpoint{1.569778in}{1.085180in}}%
\pgfpathclose%
\pgfusepath{fill}%
\end{pgfscope}%
\begin{pgfscope}%
\pgfpathrectangle{\pgfqpoint{0.407592in}{0.222904in}}{\pgfqpoint{1.411163in}{1.411163in}}%
\pgfusepath{clip}%
\pgfsetbuttcap%
\pgfsetroundjoin%
\definecolor{currentfill}{rgb}{0.168627,0.611765,0.411765}%
\pgfsetfillcolor{currentfill}%
\pgfsetfillopacity{0.180000}%
\pgfsetlinewidth{0.000000pt}%
\definecolor{currentstroke}{rgb}{0.000000,0.000000,0.000000}%
\pgfsetstrokecolor{currentstroke}%
\pgfsetdash{}{0pt}%
\pgfpathmoveto{\pgfqpoint{1.585469in}{1.149049in}}%
\pgfpathlineto{\pgfqpoint{1.600048in}{1.169701in}}%
\pgfpathlineto{\pgfqpoint{1.600618in}{1.191877in}}%
\pgfpathlineto{\pgfqpoint{1.586023in}{1.171252in}}%
\pgfpathlineto{\pgfqpoint{1.585469in}{1.149049in}}%
\pgfpathclose%
\pgfusepath{fill}%
\end{pgfscope}%
\begin{pgfscope}%
\pgfpathrectangle{\pgfqpoint{0.407592in}{0.222904in}}{\pgfqpoint{1.411163in}{1.411163in}}%
\pgfusepath{clip}%
\pgfsetbuttcap%
\pgfsetroundjoin%
\definecolor{currentfill}{rgb}{0.168627,0.611765,0.411765}%
\pgfsetfillcolor{currentfill}%
\pgfsetfillopacity{0.180000}%
\pgfsetlinewidth{0.000000pt}%
\definecolor{currentstroke}{rgb}{0.000000,0.000000,0.000000}%
\pgfsetstrokecolor{currentstroke}%
\pgfsetdash{}{0pt}%
\pgfpathmoveto{\pgfqpoint{1.554116in}{1.022572in}}%
\pgfpathlineto{\pgfqpoint{1.568713in}{1.041090in}}%
\pgfpathlineto{\pgfqpoint{1.569245in}{1.063108in}}%
\pgfpathlineto{\pgfqpoint{1.554632in}{1.044623in}}%
\pgfpathlineto{\pgfqpoint{1.554116in}{1.022572in}}%
\pgfpathclose%
\pgfusepath{fill}%
\end{pgfscope}%
\begin{pgfscope}%
\pgfpathrectangle{\pgfqpoint{0.407592in}{0.222904in}}{\pgfqpoint{1.411163in}{1.411163in}}%
\pgfusepath{clip}%
\pgfsetbuttcap%
\pgfsetroundjoin%
\definecolor{currentfill}{rgb}{0.168627,0.611765,0.411765}%
\pgfsetfillcolor{currentfill}%
\pgfsetfillopacity{0.180000}%
\pgfsetlinewidth{0.000000pt}%
\definecolor{currentstroke}{rgb}{0.000000,0.000000,0.000000}%
\pgfsetstrokecolor{currentstroke}%
\pgfsetdash{}{0pt}%
\pgfpathmoveto{\pgfqpoint{1.601190in}{1.214106in}}%
\pgfpathlineto{\pgfqpoint{1.615761in}{1.235705in}}%
\pgfpathlineto{\pgfqpoint{1.616350in}{1.257963in}}%
\pgfpathlineto{\pgfqpoint{1.601763in}{1.236389in}}%
\pgfpathlineto{\pgfqpoint{1.601190in}{1.214106in}}%
\pgfpathclose%
\pgfusepath{fill}%
\end{pgfscope}%
\begin{pgfscope}%
\pgfpathrectangle{\pgfqpoint{0.407592in}{0.222904in}}{\pgfqpoint{1.411163in}{1.411163in}}%
\pgfusepath{clip}%
\pgfsetbuttcap%
\pgfsetroundjoin%
\definecolor{currentfill}{rgb}{0.168627,0.611765,0.411765}%
\pgfsetfillcolor{currentfill}%
\pgfsetfillopacity{0.180000}%
\pgfsetlinewidth{0.000000pt}%
\definecolor{currentstroke}{rgb}{0.000000,0.000000,0.000000}%
\pgfsetstrokecolor{currentstroke}%
\pgfsetdash{}{0pt}%
\pgfpathmoveto{\pgfqpoint{1.538481in}{0.961294in}}%
\pgfpathlineto{\pgfqpoint{1.553089in}{0.978630in}}%
\pgfpathlineto{\pgfqpoint{1.553602in}{1.000574in}}%
\pgfpathlineto{\pgfqpoint{1.538978in}{0.983275in}}%
\pgfpathlineto{\pgfqpoint{1.538481in}{0.961294in}}%
\pgfpathclose%
\pgfusepath{fill}%
\end{pgfscope}%
\begin{pgfscope}%
\pgfpathrectangle{\pgfqpoint{0.407592in}{0.222904in}}{\pgfqpoint{1.411163in}{1.411163in}}%
\pgfusepath{clip}%
\pgfsetbuttcap%
\pgfsetroundjoin%
\definecolor{currentfill}{rgb}{0.168627,0.611765,0.411765}%
\pgfsetfillcolor{currentfill}%
\pgfsetfillopacity{0.180000}%
\pgfsetlinewidth{0.000000pt}%
\definecolor{currentstroke}{rgb}{0.000000,0.000000,0.000000}%
\pgfsetstrokecolor{currentstroke}%
\pgfsetdash{}{0pt}%
\pgfpathmoveto{\pgfqpoint{1.616941in}{1.280275in}}%
\pgfpathlineto{\pgfqpoint{1.631504in}{1.302742in}}%
\pgfpathlineto{\pgfqpoint{1.632113in}{1.325085in}}%
\pgfpathlineto{\pgfqpoint{1.617533in}{1.302642in}}%
\pgfpathlineto{\pgfqpoint{1.616941in}{1.280275in}}%
\pgfpathclose%
\pgfusepath{fill}%
\end{pgfscope}%
\begin{pgfscope}%
\pgfpathrectangle{\pgfqpoint{0.407592in}{0.222904in}}{\pgfqpoint{1.411163in}{1.411163in}}%
\pgfusepath{clip}%
\pgfsetbuttcap%
\pgfsetroundjoin%
\definecolor{currentfill}{rgb}{0.168627,0.611765,0.411765}%
\pgfsetfillcolor{currentfill}%
\pgfsetfillopacity{0.180000}%
\pgfsetlinewidth{0.000000pt}%
\definecolor{currentstroke}{rgb}{0.000000,0.000000,0.000000}%
\pgfsetstrokecolor{currentstroke}%
\pgfsetdash{}{0pt}%
\pgfpathmoveto{\pgfqpoint{1.632723in}{1.347483in}}%
\pgfpathlineto{\pgfqpoint{1.647279in}{1.370738in}}%
\pgfpathlineto{\pgfqpoint{1.647907in}{1.393169in}}%
\pgfpathlineto{\pgfqpoint{1.633335in}{1.369935in}}%
\pgfpathlineto{\pgfqpoint{1.632723in}{1.347483in}}%
\pgfpathclose%
\pgfusepath{fill}%
\end{pgfscope}%
\begin{pgfscope}%
\pgfpathrectangle{\pgfqpoint{0.407592in}{0.222904in}}{\pgfqpoint{1.411163in}{1.411163in}}%
\pgfusepath{clip}%
\pgfsetbuttcap%
\pgfsetroundjoin%
\definecolor{currentfill}{rgb}{0.168627,0.611765,0.411765}%
\pgfsetfillcolor{currentfill}%
\pgfsetfillopacity{0.180000}%
\pgfsetlinewidth{0.000000pt}%
\definecolor{currentstroke}{rgb}{0.000000,0.000000,0.000000}%
\pgfsetstrokecolor{currentstroke}%
\pgfsetdash{}{0pt}%
\pgfpathmoveto{\pgfqpoint{1.570313in}{1.107305in}}%
\pgfpathlineto{\pgfqpoint{1.584916in}{1.126900in}}%
\pgfpathlineto{\pgfqpoint{1.585469in}{1.149049in}}%
\pgfpathlineto{\pgfqpoint{1.570849in}{1.129485in}}%
\pgfpathlineto{\pgfqpoint{1.570313in}{1.107305in}}%
\pgfpathclose%
\pgfusepath{fill}%
\end{pgfscope}%
\begin{pgfscope}%
\pgfpathrectangle{\pgfqpoint{0.407592in}{0.222904in}}{\pgfqpoint{1.411163in}{1.411163in}}%
\pgfusepath{clip}%
\pgfsetbuttcap%
\pgfsetroundjoin%
\definecolor{currentfill}{rgb}{0.168627,0.611765,0.411765}%
\pgfsetfillcolor{currentfill}%
\pgfsetfillopacity{0.180000}%
\pgfsetlinewidth{0.000000pt}%
\definecolor{currentstroke}{rgb}{0.000000,0.000000,0.000000}%
\pgfsetstrokecolor{currentstroke}%
\pgfsetdash{}{0pt}%
\pgfpathmoveto{\pgfqpoint{1.586023in}{1.171252in}}%
\pgfpathlineto{\pgfqpoint{1.600618in}{1.191877in}}%
\pgfpathlineto{\pgfqpoint{1.601190in}{1.214106in}}%
\pgfpathlineto{\pgfqpoint{1.586578in}{1.193510in}}%
\pgfpathlineto{\pgfqpoint{1.586023in}{1.171252in}}%
\pgfpathclose%
\pgfusepath{fill}%
\end{pgfscope}%
\begin{pgfscope}%
\pgfpathrectangle{\pgfqpoint{0.407592in}{0.222904in}}{\pgfqpoint{1.411163in}{1.411163in}}%
\pgfusepath{clip}%
\pgfsetbuttcap%
\pgfsetroundjoin%
\definecolor{currentfill}{rgb}{0.168627,0.611765,0.411765}%
\pgfsetfillcolor{currentfill}%
\pgfsetfillopacity{0.180000}%
\pgfsetlinewidth{0.000000pt}%
\definecolor{currentstroke}{rgb}{0.000000,0.000000,0.000000}%
\pgfsetstrokecolor{currentstroke}%
\pgfsetdash{}{0pt}%
\pgfpathmoveto{\pgfqpoint{1.554632in}{1.044623in}}%
\pgfpathlineto{\pgfqpoint{1.569245in}{1.063108in}}%
\pgfpathlineto{\pgfqpoint{1.569778in}{1.085180in}}%
\pgfpathlineto{\pgfqpoint{1.555148in}{1.066728in}}%
\pgfpathlineto{\pgfqpoint{1.554632in}{1.044623in}}%
\pgfpathclose%
\pgfusepath{fill}%
\end{pgfscope}%
\begin{pgfscope}%
\pgfpathrectangle{\pgfqpoint{0.407592in}{0.222904in}}{\pgfqpoint{1.411163in}{1.411163in}}%
\pgfusepath{clip}%
\pgfsetbuttcap%
\pgfsetroundjoin%
\definecolor{currentfill}{rgb}{0.168627,0.611765,0.411765}%
\pgfsetfillcolor{currentfill}%
\pgfsetfillopacity{0.180000}%
\pgfsetlinewidth{0.000000pt}%
\definecolor{currentstroke}{rgb}{0.000000,0.000000,0.000000}%
\pgfsetstrokecolor{currentstroke}%
\pgfsetdash{}{0pt}%
\pgfpathmoveto{\pgfqpoint{1.601763in}{1.236389in}}%
\pgfpathlineto{\pgfqpoint{1.616350in}{1.257963in}}%
\pgfpathlineto{\pgfqpoint{1.616941in}{1.280275in}}%
\pgfpathlineto{\pgfqpoint{1.602337in}{1.258728in}}%
\pgfpathlineto{\pgfqpoint{1.601763in}{1.236389in}}%
\pgfpathclose%
\pgfusepath{fill}%
\end{pgfscope}%
\begin{pgfscope}%
\pgfpathrectangle{\pgfqpoint{0.407592in}{0.222904in}}{\pgfqpoint{1.411163in}{1.411163in}}%
\pgfusepath{clip}%
\pgfsetbuttcap%
\pgfsetroundjoin%
\definecolor{currentfill}{rgb}{0.168627,0.611765,0.411765}%
\pgfsetfillcolor{currentfill}%
\pgfsetfillopacity{0.180000}%
\pgfsetlinewidth{0.000000pt}%
\definecolor{currentstroke}{rgb}{0.000000,0.000000,0.000000}%
\pgfsetstrokecolor{currentstroke}%
\pgfsetdash{}{0pt}%
\pgfpathmoveto{\pgfqpoint{1.538978in}{0.983275in}}%
\pgfpathlineto{\pgfqpoint{1.553602in}{1.000574in}}%
\pgfpathlineto{\pgfqpoint{1.554116in}{1.022572in}}%
\pgfpathlineto{\pgfqpoint{1.539476in}{1.005309in}}%
\pgfpathlineto{\pgfqpoint{1.538978in}{0.983275in}}%
\pgfpathclose%
\pgfusepath{fill}%
\end{pgfscope}%
\begin{pgfscope}%
\pgfpathrectangle{\pgfqpoint{0.407592in}{0.222904in}}{\pgfqpoint{1.411163in}{1.411163in}}%
\pgfusepath{clip}%
\pgfsetbuttcap%
\pgfsetroundjoin%
\definecolor{currentfill}{rgb}{0.168627,0.611765,0.411765}%
\pgfsetfillcolor{currentfill}%
\pgfsetfillopacity{0.180000}%
\pgfsetlinewidth{0.000000pt}%
\definecolor{currentstroke}{rgb}{0.000000,0.000000,0.000000}%
\pgfsetstrokecolor{currentstroke}%
\pgfsetdash{}{0pt}%
\pgfpathmoveto{\pgfqpoint{1.617533in}{1.302642in}}%
\pgfpathlineto{\pgfqpoint{1.632113in}{1.325085in}}%
\pgfpathlineto{\pgfqpoint{1.632723in}{1.347483in}}%
\pgfpathlineto{\pgfqpoint{1.618127in}{1.325064in}}%
\pgfpathlineto{\pgfqpoint{1.617533in}{1.302642in}}%
\pgfpathclose%
\pgfusepath{fill}%
\end{pgfscope}%
\begin{pgfscope}%
\pgfpathrectangle{\pgfqpoint{0.407592in}{0.222904in}}{\pgfqpoint{1.411163in}{1.411163in}}%
\pgfusepath{clip}%
\pgfsetbuttcap%
\pgfsetroundjoin%
\definecolor{currentfill}{rgb}{0.168627,0.611765,0.411765}%
\pgfsetfillcolor{currentfill}%
\pgfsetfillopacity{0.180000}%
\pgfsetlinewidth{0.000000pt}%
\definecolor{currentstroke}{rgb}{0.000000,0.000000,0.000000}%
\pgfsetstrokecolor{currentstroke}%
\pgfsetdash{}{0pt}%
\pgfpathmoveto{\pgfqpoint{1.570849in}{1.129485in}}%
\pgfpathlineto{\pgfqpoint{1.585469in}{1.149049in}}%
\pgfpathlineto{\pgfqpoint{1.586023in}{1.171252in}}%
\pgfpathlineto{\pgfqpoint{1.571386in}{1.151719in}}%
\pgfpathlineto{\pgfqpoint{1.570849in}{1.129485in}}%
\pgfpathclose%
\pgfusepath{fill}%
\end{pgfscope}%
\begin{pgfscope}%
\pgfpathrectangle{\pgfqpoint{0.407592in}{0.222904in}}{\pgfqpoint{1.411163in}{1.411163in}}%
\pgfusepath{clip}%
\pgfsetbuttcap%
\pgfsetroundjoin%
\definecolor{currentfill}{rgb}{0.168627,0.611765,0.411765}%
\pgfsetfillcolor{currentfill}%
\pgfsetfillopacity{0.180000}%
\pgfsetlinewidth{0.000000pt}%
\definecolor{currentstroke}{rgb}{0.000000,0.000000,0.000000}%
\pgfsetstrokecolor{currentstroke}%
\pgfsetdash{}{0pt}%
\pgfpathmoveto{\pgfqpoint{1.586578in}{1.193510in}}%
\pgfpathlineto{\pgfqpoint{1.601190in}{1.214106in}}%
\pgfpathlineto{\pgfqpoint{1.601763in}{1.236389in}}%
\pgfpathlineto{\pgfqpoint{1.587134in}{1.215822in}}%
\pgfpathlineto{\pgfqpoint{1.586578in}{1.193510in}}%
\pgfpathclose%
\pgfusepath{fill}%
\end{pgfscope}%
\begin{pgfscope}%
\pgfpathrectangle{\pgfqpoint{0.407592in}{0.222904in}}{\pgfqpoint{1.411163in}{1.411163in}}%
\pgfusepath{clip}%
\pgfsetbuttcap%
\pgfsetroundjoin%
\definecolor{currentfill}{rgb}{0.168627,0.611765,0.411765}%
\pgfsetfillcolor{currentfill}%
\pgfsetfillopacity{0.180000}%
\pgfsetlinewidth{0.000000pt}%
\definecolor{currentstroke}{rgb}{0.000000,0.000000,0.000000}%
\pgfsetstrokecolor{currentstroke}%
\pgfsetdash{}{0pt}%
\pgfpathmoveto{\pgfqpoint{1.555148in}{1.066728in}}%
\pgfpathlineto{\pgfqpoint{1.569778in}{1.085180in}}%
\pgfpathlineto{\pgfqpoint{1.570313in}{1.107305in}}%
\pgfpathlineto{\pgfqpoint{1.555666in}{1.088888in}}%
\pgfpathlineto{\pgfqpoint{1.555148in}{1.066728in}}%
\pgfpathclose%
\pgfusepath{fill}%
\end{pgfscope}%
\begin{pgfscope}%
\pgfpathrectangle{\pgfqpoint{0.407592in}{0.222904in}}{\pgfqpoint{1.411163in}{1.411163in}}%
\pgfusepath{clip}%
\pgfsetbuttcap%
\pgfsetroundjoin%
\definecolor{currentfill}{rgb}{0.168627,0.611765,0.411765}%
\pgfsetfillcolor{currentfill}%
\pgfsetfillopacity{0.180000}%
\pgfsetlinewidth{0.000000pt}%
\definecolor{currentstroke}{rgb}{0.000000,0.000000,0.000000}%
\pgfsetstrokecolor{currentstroke}%
\pgfsetdash{}{0pt}%
\pgfpathmoveto{\pgfqpoint{1.602337in}{1.258728in}}%
\pgfpathlineto{\pgfqpoint{1.616941in}{1.280275in}}%
\pgfpathlineto{\pgfqpoint{1.617533in}{1.302642in}}%
\pgfpathlineto{\pgfqpoint{1.602913in}{1.281120in}}%
\pgfpathlineto{\pgfqpoint{1.602337in}{1.258728in}}%
\pgfpathclose%
\pgfusepath{fill}%
\end{pgfscope}%
\begin{pgfscope}%
\pgfpathrectangle{\pgfqpoint{0.407592in}{0.222904in}}{\pgfqpoint{1.411163in}{1.411163in}}%
\pgfusepath{clip}%
\pgfsetbuttcap%
\pgfsetroundjoin%
\definecolor{currentfill}{rgb}{0.168627,0.611765,0.411765}%
\pgfsetfillcolor{currentfill}%
\pgfsetfillopacity{0.180000}%
\pgfsetlinewidth{0.000000pt}%
\definecolor{currentstroke}{rgb}{0.000000,0.000000,0.000000}%
\pgfsetstrokecolor{currentstroke}%
\pgfsetdash{}{0pt}%
\pgfpathmoveto{\pgfqpoint{1.539476in}{1.005309in}}%
\pgfpathlineto{\pgfqpoint{1.554116in}{1.022572in}}%
\pgfpathlineto{\pgfqpoint{1.554632in}{1.044623in}}%
\pgfpathlineto{\pgfqpoint{1.539975in}{1.027396in}}%
\pgfpathlineto{\pgfqpoint{1.539476in}{1.005309in}}%
\pgfpathclose%
\pgfusepath{fill}%
\end{pgfscope}%
\begin{pgfscope}%
\pgfpathrectangle{\pgfqpoint{0.407592in}{0.222904in}}{\pgfqpoint{1.411163in}{1.411163in}}%
\pgfusepath{clip}%
\pgfsetbuttcap%
\pgfsetroundjoin%
\definecolor{currentfill}{rgb}{0.168627,0.611765,0.411765}%
\pgfsetfillcolor{currentfill}%
\pgfsetfillopacity{0.180000}%
\pgfsetlinewidth{0.000000pt}%
\definecolor{currentstroke}{rgb}{0.000000,0.000000,0.000000}%
\pgfsetstrokecolor{currentstroke}%
\pgfsetdash{}{0pt}%
\pgfpathmoveto{\pgfqpoint{1.618127in}{1.325064in}}%
\pgfpathlineto{\pgfqpoint{1.632723in}{1.347483in}}%
\pgfpathlineto{\pgfqpoint{1.633335in}{1.369935in}}%
\pgfpathlineto{\pgfqpoint{1.618722in}{1.347540in}}%
\pgfpathlineto{\pgfqpoint{1.618127in}{1.325064in}}%
\pgfpathclose%
\pgfusepath{fill}%
\end{pgfscope}%
\begin{pgfscope}%
\pgfpathrectangle{\pgfqpoint{0.407592in}{0.222904in}}{\pgfqpoint{1.411163in}{1.411163in}}%
\pgfusepath{clip}%
\pgfsetbuttcap%
\pgfsetroundjoin%
\definecolor{currentfill}{rgb}{0.168627,0.611765,0.411765}%
\pgfsetfillcolor{currentfill}%
\pgfsetfillopacity{0.180000}%
\pgfsetlinewidth{0.000000pt}%
\definecolor{currentstroke}{rgb}{0.000000,0.000000,0.000000}%
\pgfsetstrokecolor{currentstroke}%
\pgfsetdash{}{0pt}%
\pgfpathmoveto{\pgfqpoint{1.523829in}{0.945288in}}%
\pgfpathlineto{\pgfqpoint{1.538481in}{0.961294in}}%
\pgfpathlineto{\pgfqpoint{1.538978in}{0.983275in}}%
\pgfpathlineto{\pgfqpoint{1.524308in}{0.967308in}}%
\pgfpathlineto{\pgfqpoint{1.523829in}{0.945288in}}%
\pgfpathclose%
\pgfusepath{fill}%
\end{pgfscope}%
\begin{pgfscope}%
\pgfpathrectangle{\pgfqpoint{0.407592in}{0.222904in}}{\pgfqpoint{1.411163in}{1.411163in}}%
\pgfusepath{clip}%
\pgfsetbuttcap%
\pgfsetroundjoin%
\definecolor{currentfill}{rgb}{0.168627,0.611765,0.411765}%
\pgfsetfillcolor{currentfill}%
\pgfsetfillopacity{0.180000}%
\pgfsetlinewidth{0.000000pt}%
\definecolor{currentstroke}{rgb}{0.000000,0.000000,0.000000}%
\pgfsetstrokecolor{currentstroke}%
\pgfsetdash{}{0pt}%
\pgfpathmoveto{\pgfqpoint{1.571386in}{1.151719in}}%
\pgfpathlineto{\pgfqpoint{1.586023in}{1.171252in}}%
\pgfpathlineto{\pgfqpoint{1.586578in}{1.193510in}}%
\pgfpathlineto{\pgfqpoint{1.571924in}{1.174008in}}%
\pgfpathlineto{\pgfqpoint{1.571386in}{1.151719in}}%
\pgfpathclose%
\pgfusepath{fill}%
\end{pgfscope}%
\begin{pgfscope}%
\pgfpathrectangle{\pgfqpoint{0.407592in}{0.222904in}}{\pgfqpoint{1.411163in}{1.411163in}}%
\pgfusepath{clip}%
\pgfsetbuttcap%
\pgfsetroundjoin%
\definecolor{currentfill}{rgb}{0.168627,0.611765,0.411765}%
\pgfsetfillcolor{currentfill}%
\pgfsetfillopacity{0.180000}%
\pgfsetlinewidth{0.000000pt}%
\definecolor{currentstroke}{rgb}{0.000000,0.000000,0.000000}%
\pgfsetstrokecolor{currentstroke}%
\pgfsetdash{}{0pt}%
\pgfpathmoveto{\pgfqpoint{1.587134in}{1.215822in}}%
\pgfpathlineto{\pgfqpoint{1.601763in}{1.236389in}}%
\pgfpathlineto{\pgfqpoint{1.602337in}{1.258728in}}%
\pgfpathlineto{\pgfqpoint{1.587692in}{1.238188in}}%
\pgfpathlineto{\pgfqpoint{1.587134in}{1.215822in}}%
\pgfpathclose%
\pgfusepath{fill}%
\end{pgfscope}%
\begin{pgfscope}%
\pgfpathrectangle{\pgfqpoint{0.407592in}{0.222904in}}{\pgfqpoint{1.411163in}{1.411163in}}%
\pgfusepath{clip}%
\pgfsetbuttcap%
\pgfsetroundjoin%
\definecolor{currentfill}{rgb}{0.168627,0.611765,0.411765}%
\pgfsetfillcolor{currentfill}%
\pgfsetfillopacity{0.180000}%
\pgfsetlinewidth{0.000000pt}%
\definecolor{currentstroke}{rgb}{0.000000,0.000000,0.000000}%
\pgfsetstrokecolor{currentstroke}%
\pgfsetdash{}{0pt}%
\pgfpathmoveto{\pgfqpoint{1.555666in}{1.088888in}}%
\pgfpathlineto{\pgfqpoint{1.570313in}{1.107305in}}%
\pgfpathlineto{\pgfqpoint{1.570849in}{1.129485in}}%
\pgfpathlineto{\pgfqpoint{1.556186in}{1.111101in}}%
\pgfpathlineto{\pgfqpoint{1.555666in}{1.088888in}}%
\pgfpathclose%
\pgfusepath{fill}%
\end{pgfscope}%
\begin{pgfscope}%
\pgfpathrectangle{\pgfqpoint{0.407592in}{0.222904in}}{\pgfqpoint{1.411163in}{1.411163in}}%
\pgfusepath{clip}%
\pgfsetbuttcap%
\pgfsetroundjoin%
\definecolor{currentfill}{rgb}{0.168627,0.611765,0.411765}%
\pgfsetfillcolor{currentfill}%
\pgfsetfillopacity{0.180000}%
\pgfsetlinewidth{0.000000pt}%
\definecolor{currentstroke}{rgb}{0.000000,0.000000,0.000000}%
\pgfsetstrokecolor{currentstroke}%
\pgfsetdash{}{0pt}%
\pgfpathmoveto{\pgfqpoint{1.602913in}{1.281120in}}%
\pgfpathlineto{\pgfqpoint{1.617533in}{1.302642in}}%
\pgfpathlineto{\pgfqpoint{1.618127in}{1.325064in}}%
\pgfpathlineto{\pgfqpoint{1.603490in}{1.303568in}}%
\pgfpathlineto{\pgfqpoint{1.602913in}{1.281120in}}%
\pgfpathclose%
\pgfusepath{fill}%
\end{pgfscope}%
\begin{pgfscope}%
\pgfpathrectangle{\pgfqpoint{0.407592in}{0.222904in}}{\pgfqpoint{1.411163in}{1.411163in}}%
\pgfusepath{clip}%
\pgfsetbuttcap%
\pgfsetroundjoin%
\definecolor{currentfill}{rgb}{0.168627,0.611765,0.411765}%
\pgfsetfillcolor{currentfill}%
\pgfsetfillopacity{0.180000}%
\pgfsetlinewidth{0.000000pt}%
\definecolor{currentstroke}{rgb}{0.000000,0.000000,0.000000}%
\pgfsetstrokecolor{currentstroke}%
\pgfsetdash{}{0pt}%
\pgfpathmoveto{\pgfqpoint{1.539975in}{1.027396in}}%
\pgfpathlineto{\pgfqpoint{1.554632in}{1.044623in}}%
\pgfpathlineto{\pgfqpoint{1.555148in}{1.066728in}}%
\pgfpathlineto{\pgfqpoint{1.540475in}{1.049539in}}%
\pgfpathlineto{\pgfqpoint{1.539975in}{1.027396in}}%
\pgfpathclose%
\pgfusepath{fill}%
\end{pgfscope}%
\begin{pgfscope}%
\pgfpathrectangle{\pgfqpoint{0.407592in}{0.222904in}}{\pgfqpoint{1.411163in}{1.411163in}}%
\pgfusepath{clip}%
\pgfsetbuttcap%
\pgfsetroundjoin%
\definecolor{currentfill}{rgb}{0.168627,0.611765,0.411765}%
\pgfsetfillcolor{currentfill}%
\pgfsetfillopacity{0.180000}%
\pgfsetlinewidth{0.000000pt}%
\definecolor{currentstroke}{rgb}{0.000000,0.000000,0.000000}%
\pgfsetstrokecolor{currentstroke}%
\pgfsetdash{}{0pt}%
\pgfpathmoveto{\pgfqpoint{1.524308in}{0.967308in}}%
\pgfpathlineto{\pgfqpoint{1.538978in}{0.983275in}}%
\pgfpathlineto{\pgfqpoint{1.539476in}{1.005309in}}%
\pgfpathlineto{\pgfqpoint{1.524790in}{0.989382in}}%
\pgfpathlineto{\pgfqpoint{1.524308in}{0.967308in}}%
\pgfpathclose%
\pgfusepath{fill}%
\end{pgfscope}%
\begin{pgfscope}%
\pgfpathrectangle{\pgfqpoint{0.407592in}{0.222904in}}{\pgfqpoint{1.411163in}{1.411163in}}%
\pgfusepath{clip}%
\pgfsetbuttcap%
\pgfsetroundjoin%
\definecolor{currentfill}{rgb}{0.168627,0.611765,0.411765}%
\pgfsetfillcolor{currentfill}%
\pgfsetfillopacity{0.180000}%
\pgfsetlinewidth{0.000000pt}%
\definecolor{currentstroke}{rgb}{0.000000,0.000000,0.000000}%
\pgfsetstrokecolor{currentstroke}%
\pgfsetdash{}{0pt}%
\pgfpathmoveto{\pgfqpoint{1.571924in}{1.174008in}}%
\pgfpathlineto{\pgfqpoint{1.586578in}{1.193510in}}%
\pgfpathlineto{\pgfqpoint{1.587134in}{1.215822in}}%
\pgfpathlineto{\pgfqpoint{1.572464in}{1.196352in}}%
\pgfpathlineto{\pgfqpoint{1.571924in}{1.174008in}}%
\pgfpathclose%
\pgfusepath{fill}%
\end{pgfscope}%
\begin{pgfscope}%
\pgfpathrectangle{\pgfqpoint{0.407592in}{0.222904in}}{\pgfqpoint{1.411163in}{1.411163in}}%
\pgfusepath{clip}%
\pgfsetbuttcap%
\pgfsetroundjoin%
\definecolor{currentfill}{rgb}{0.168627,0.611765,0.411765}%
\pgfsetfillcolor{currentfill}%
\pgfsetfillopacity{0.180000}%
\pgfsetlinewidth{0.000000pt}%
\definecolor{currentstroke}{rgb}{0.000000,0.000000,0.000000}%
\pgfsetstrokecolor{currentstroke}%
\pgfsetdash{}{0pt}%
\pgfpathmoveto{\pgfqpoint{1.587692in}{1.238188in}}%
\pgfpathlineto{\pgfqpoint{1.602337in}{1.258728in}}%
\pgfpathlineto{\pgfqpoint{1.602913in}{1.281120in}}%
\pgfpathlineto{\pgfqpoint{1.588251in}{1.260610in}}%
\pgfpathlineto{\pgfqpoint{1.587692in}{1.238188in}}%
\pgfpathclose%
\pgfusepath{fill}%
\end{pgfscope}%
\begin{pgfscope}%
\pgfpathrectangle{\pgfqpoint{0.407592in}{0.222904in}}{\pgfqpoint{1.411163in}{1.411163in}}%
\pgfusepath{clip}%
\pgfsetbuttcap%
\pgfsetroundjoin%
\definecolor{currentfill}{rgb}{0.168627,0.611765,0.411765}%
\pgfsetfillcolor{currentfill}%
\pgfsetfillopacity{0.180000}%
\pgfsetlinewidth{0.000000pt}%
\definecolor{currentstroke}{rgb}{0.000000,0.000000,0.000000}%
\pgfsetstrokecolor{currentstroke}%
\pgfsetdash{}{0pt}%
\pgfpathmoveto{\pgfqpoint{1.556186in}{1.111101in}}%
\pgfpathlineto{\pgfqpoint{1.570849in}{1.129485in}}%
\pgfpathlineto{\pgfqpoint{1.571386in}{1.151719in}}%
\pgfpathlineto{\pgfqpoint{1.556706in}{1.133370in}}%
\pgfpathlineto{\pgfqpoint{1.556186in}{1.111101in}}%
\pgfpathclose%
\pgfusepath{fill}%
\end{pgfscope}%
\begin{pgfscope}%
\pgfpathrectangle{\pgfqpoint{0.407592in}{0.222904in}}{\pgfqpoint{1.411163in}{1.411163in}}%
\pgfusepath{clip}%
\pgfsetbuttcap%
\pgfsetroundjoin%
\definecolor{currentfill}{rgb}{0.168627,0.611765,0.411765}%
\pgfsetfillcolor{currentfill}%
\pgfsetfillopacity{0.180000}%
\pgfsetlinewidth{0.000000pt}%
\definecolor{currentstroke}{rgb}{0.000000,0.000000,0.000000}%
\pgfsetstrokecolor{currentstroke}%
\pgfsetdash{}{0pt}%
\pgfpathmoveto{\pgfqpoint{1.603490in}{1.303568in}}%
\pgfpathlineto{\pgfqpoint{1.618127in}{1.325064in}}%
\pgfpathlineto{\pgfqpoint{1.618722in}{1.347540in}}%
\pgfpathlineto{\pgfqpoint{1.604069in}{1.326071in}}%
\pgfpathlineto{\pgfqpoint{1.603490in}{1.303568in}}%
\pgfpathclose%
\pgfusepath{fill}%
\end{pgfscope}%
\begin{pgfscope}%
\pgfpathrectangle{\pgfqpoint{0.407592in}{0.222904in}}{\pgfqpoint{1.411163in}{1.411163in}}%
\pgfusepath{clip}%
\pgfsetbuttcap%
\pgfsetroundjoin%
\definecolor{currentfill}{rgb}{0.168627,0.611765,0.411765}%
\pgfsetfillcolor{currentfill}%
\pgfsetfillopacity{0.180000}%
\pgfsetlinewidth{0.000000pt}%
\definecolor{currentstroke}{rgb}{0.000000,0.000000,0.000000}%
\pgfsetstrokecolor{currentstroke}%
\pgfsetdash{}{0pt}%
\pgfpathmoveto{\pgfqpoint{1.540475in}{1.049539in}}%
\pgfpathlineto{\pgfqpoint{1.555148in}{1.066728in}}%
\pgfpathlineto{\pgfqpoint{1.555666in}{1.088888in}}%
\pgfpathlineto{\pgfqpoint{1.540976in}{1.071735in}}%
\pgfpathlineto{\pgfqpoint{1.540475in}{1.049539in}}%
\pgfpathclose%
\pgfusepath{fill}%
\end{pgfscope}%
\begin{pgfscope}%
\pgfpathrectangle{\pgfqpoint{0.407592in}{0.222904in}}{\pgfqpoint{1.411163in}{1.411163in}}%
\pgfusepath{clip}%
\pgfsetbuttcap%
\pgfsetroundjoin%
\definecolor{currentfill}{rgb}{0.168627,0.611765,0.411765}%
\pgfsetfillcolor{currentfill}%
\pgfsetfillopacity{0.180000}%
\pgfsetlinewidth{0.000000pt}%
\definecolor{currentstroke}{rgb}{0.000000,0.000000,0.000000}%
\pgfsetstrokecolor{currentstroke}%
\pgfsetdash{}{0pt}%
\pgfpathmoveto{\pgfqpoint{1.524790in}{0.989382in}}%
\pgfpathlineto{\pgfqpoint{1.539476in}{1.005309in}}%
\pgfpathlineto{\pgfqpoint{1.539975in}{1.027396in}}%
\pgfpathlineto{\pgfqpoint{1.525272in}{1.011510in}}%
\pgfpathlineto{\pgfqpoint{1.524790in}{0.989382in}}%
\pgfpathclose%
\pgfusepath{fill}%
\end{pgfscope}%
\begin{pgfscope}%
\pgfpathrectangle{\pgfqpoint{0.407592in}{0.222904in}}{\pgfqpoint{1.411163in}{1.411163in}}%
\pgfusepath{clip}%
\pgfsetbuttcap%
\pgfsetroundjoin%
\definecolor{currentfill}{rgb}{0.168627,0.611765,0.411765}%
\pgfsetfillcolor{currentfill}%
\pgfsetfillopacity{0.180000}%
\pgfsetlinewidth{0.000000pt}%
\definecolor{currentstroke}{rgb}{0.000000,0.000000,0.000000}%
\pgfsetstrokecolor{currentstroke}%
\pgfsetdash{}{0pt}%
\pgfpathmoveto{\pgfqpoint{1.509128in}{0.930684in}}%
\pgfpathlineto{\pgfqpoint{1.523829in}{0.945288in}}%
\pgfpathlineto{\pgfqpoint{1.524308in}{0.967308in}}%
\pgfpathlineto{\pgfqpoint{1.509591in}{0.952748in}}%
\pgfpathlineto{\pgfqpoint{1.509128in}{0.930684in}}%
\pgfpathclose%
\pgfusepath{fill}%
\end{pgfscope}%
\begin{pgfscope}%
\pgfpathrectangle{\pgfqpoint{0.407592in}{0.222904in}}{\pgfqpoint{1.411163in}{1.411163in}}%
\pgfusepath{clip}%
\pgfsetbuttcap%
\pgfsetroundjoin%
\definecolor{currentfill}{rgb}{0.168627,0.611765,0.411765}%
\pgfsetfillcolor{currentfill}%
\pgfsetfillopacity{0.180000}%
\pgfsetlinewidth{0.000000pt}%
\definecolor{currentstroke}{rgb}{0.000000,0.000000,0.000000}%
\pgfsetstrokecolor{currentstroke}%
\pgfsetdash{}{0pt}%
\pgfpathmoveto{\pgfqpoint{1.572464in}{1.196352in}}%
\pgfpathlineto{\pgfqpoint{1.587134in}{1.215822in}}%
\pgfpathlineto{\pgfqpoint{1.587692in}{1.238188in}}%
\pgfpathlineto{\pgfqpoint{1.573006in}{1.218750in}}%
\pgfpathlineto{\pgfqpoint{1.572464in}{1.196352in}}%
\pgfpathclose%
\pgfusepath{fill}%
\end{pgfscope}%
\begin{pgfscope}%
\pgfpathrectangle{\pgfqpoint{0.407592in}{0.222904in}}{\pgfqpoint{1.411163in}{1.411163in}}%
\pgfusepath{clip}%
\pgfsetbuttcap%
\pgfsetroundjoin%
\definecolor{currentfill}{rgb}{0.168627,0.611765,0.411765}%
\pgfsetfillcolor{currentfill}%
\pgfsetfillopacity{0.180000}%
\pgfsetlinewidth{0.000000pt}%
\definecolor{currentstroke}{rgb}{0.000000,0.000000,0.000000}%
\pgfsetstrokecolor{currentstroke}%
\pgfsetdash{}{0pt}%
\pgfpathmoveto{\pgfqpoint{1.588251in}{1.260610in}}%
\pgfpathlineto{\pgfqpoint{1.602913in}{1.281120in}}%
\pgfpathlineto{\pgfqpoint{1.603490in}{1.303568in}}%
\pgfpathlineto{\pgfqpoint{1.588812in}{1.283087in}}%
\pgfpathlineto{\pgfqpoint{1.588251in}{1.260610in}}%
\pgfpathclose%
\pgfusepath{fill}%
\end{pgfscope}%
\begin{pgfscope}%
\pgfpathrectangle{\pgfqpoint{0.407592in}{0.222904in}}{\pgfqpoint{1.411163in}{1.411163in}}%
\pgfusepath{clip}%
\pgfsetbuttcap%
\pgfsetroundjoin%
\definecolor{currentfill}{rgb}{0.168627,0.611765,0.411765}%
\pgfsetfillcolor{currentfill}%
\pgfsetfillopacity{0.180000}%
\pgfsetlinewidth{0.000000pt}%
\definecolor{currentstroke}{rgb}{0.000000,0.000000,0.000000}%
\pgfsetstrokecolor{currentstroke}%
\pgfsetdash{}{0pt}%
\pgfpathmoveto{\pgfqpoint{1.556706in}{1.133370in}}%
\pgfpathlineto{\pgfqpoint{1.571386in}{1.151719in}}%
\pgfpathlineto{\pgfqpoint{1.571924in}{1.174008in}}%
\pgfpathlineto{\pgfqpoint{1.557228in}{1.155693in}}%
\pgfpathlineto{\pgfqpoint{1.556706in}{1.133370in}}%
\pgfpathclose%
\pgfusepath{fill}%
\end{pgfscope}%
\begin{pgfscope}%
\pgfpathrectangle{\pgfqpoint{0.407592in}{0.222904in}}{\pgfqpoint{1.411163in}{1.411163in}}%
\pgfusepath{clip}%
\pgfsetbuttcap%
\pgfsetroundjoin%
\definecolor{currentfill}{rgb}{0.168627,0.611765,0.411765}%
\pgfsetfillcolor{currentfill}%
\pgfsetfillopacity{0.180000}%
\pgfsetlinewidth{0.000000pt}%
\definecolor{currentstroke}{rgb}{0.000000,0.000000,0.000000}%
\pgfsetstrokecolor{currentstroke}%
\pgfsetdash{}{0pt}%
\pgfpathmoveto{\pgfqpoint{1.540976in}{1.071735in}}%
\pgfpathlineto{\pgfqpoint{1.555666in}{1.088888in}}%
\pgfpathlineto{\pgfqpoint{1.556186in}{1.111101in}}%
\pgfpathlineto{\pgfqpoint{1.541479in}{1.093986in}}%
\pgfpathlineto{\pgfqpoint{1.540976in}{1.071735in}}%
\pgfpathclose%
\pgfusepath{fill}%
\end{pgfscope}%
\begin{pgfscope}%
\pgfpathrectangle{\pgfqpoint{0.407592in}{0.222904in}}{\pgfqpoint{1.411163in}{1.411163in}}%
\pgfusepath{clip}%
\pgfsetbuttcap%
\pgfsetroundjoin%
\definecolor{currentfill}{rgb}{0.168627,0.611765,0.411765}%
\pgfsetfillcolor{currentfill}%
\pgfsetfillopacity{0.180000}%
\pgfsetlinewidth{0.000000pt}%
\definecolor{currentstroke}{rgb}{0.000000,0.000000,0.000000}%
\pgfsetstrokecolor{currentstroke}%
\pgfsetdash{}{0pt}%
\pgfpathmoveto{\pgfqpoint{1.525272in}{1.011510in}}%
\pgfpathlineto{\pgfqpoint{1.539975in}{1.027396in}}%
\pgfpathlineto{\pgfqpoint{1.540475in}{1.049539in}}%
\pgfpathlineto{\pgfqpoint{1.525755in}{1.033693in}}%
\pgfpathlineto{\pgfqpoint{1.525272in}{1.011510in}}%
\pgfpathclose%
\pgfusepath{fill}%
\end{pgfscope}%
\begin{pgfscope}%
\pgfpathrectangle{\pgfqpoint{0.407592in}{0.222904in}}{\pgfqpoint{1.411163in}{1.411163in}}%
\pgfusepath{clip}%
\pgfsetbuttcap%
\pgfsetroundjoin%
\definecolor{currentfill}{rgb}{0.168627,0.611765,0.411765}%
\pgfsetfillcolor{currentfill}%
\pgfsetfillopacity{0.180000}%
\pgfsetlinewidth{0.000000pt}%
\definecolor{currentstroke}{rgb}{0.000000,0.000000,0.000000}%
\pgfsetstrokecolor{currentstroke}%
\pgfsetdash{}{0pt}%
\pgfpathmoveto{\pgfqpoint{1.509591in}{0.952748in}}%
\pgfpathlineto{\pgfqpoint{1.524308in}{0.967308in}}%
\pgfpathlineto{\pgfqpoint{1.524790in}{0.989382in}}%
\pgfpathlineto{\pgfqpoint{1.510055in}{0.974865in}}%
\pgfpathlineto{\pgfqpoint{1.509591in}{0.952748in}}%
\pgfpathclose%
\pgfusepath{fill}%
\end{pgfscope}%
\begin{pgfscope}%
\pgfpathrectangle{\pgfqpoint{0.407592in}{0.222904in}}{\pgfqpoint{1.411163in}{1.411163in}}%
\pgfusepath{clip}%
\pgfsetbuttcap%
\pgfsetroundjoin%
\definecolor{currentfill}{rgb}{0.168627,0.611765,0.411765}%
\pgfsetfillcolor{currentfill}%
\pgfsetfillopacity{0.180000}%
\pgfsetlinewidth{0.000000pt}%
\definecolor{currentstroke}{rgb}{0.000000,0.000000,0.000000}%
\pgfsetstrokecolor{currentstroke}%
\pgfsetdash{}{0pt}%
\pgfpathmoveto{\pgfqpoint{1.573006in}{1.218750in}}%
\pgfpathlineto{\pgfqpoint{1.587692in}{1.238188in}}%
\pgfpathlineto{\pgfqpoint{1.588251in}{1.260610in}}%
\pgfpathlineto{\pgfqpoint{1.573548in}{1.241203in}}%
\pgfpathlineto{\pgfqpoint{1.573006in}{1.218750in}}%
\pgfpathclose%
\pgfusepath{fill}%
\end{pgfscope}%
\begin{pgfscope}%
\pgfpathrectangle{\pgfqpoint{0.407592in}{0.222904in}}{\pgfqpoint{1.411163in}{1.411163in}}%
\pgfusepath{clip}%
\pgfsetbuttcap%
\pgfsetroundjoin%
\definecolor{currentfill}{rgb}{0.168627,0.611765,0.411765}%
\pgfsetfillcolor{currentfill}%
\pgfsetfillopacity{0.180000}%
\pgfsetlinewidth{0.000000pt}%
\definecolor{currentstroke}{rgb}{0.000000,0.000000,0.000000}%
\pgfsetstrokecolor{currentstroke}%
\pgfsetdash{}{0pt}%
\pgfpathmoveto{\pgfqpoint{1.588812in}{1.283087in}}%
\pgfpathlineto{\pgfqpoint{1.603490in}{1.303568in}}%
\pgfpathlineto{\pgfqpoint{1.604069in}{1.326071in}}%
\pgfpathlineto{\pgfqpoint{1.589374in}{1.305619in}}%
\pgfpathlineto{\pgfqpoint{1.588812in}{1.283087in}}%
\pgfpathclose%
\pgfusepath{fill}%
\end{pgfscope}%
\begin{pgfscope}%
\pgfpathrectangle{\pgfqpoint{0.407592in}{0.222904in}}{\pgfqpoint{1.411163in}{1.411163in}}%
\pgfusepath{clip}%
\pgfsetbuttcap%
\pgfsetroundjoin%
\definecolor{currentfill}{rgb}{0.168627,0.611765,0.411765}%
\pgfsetfillcolor{currentfill}%
\pgfsetfillopacity{0.180000}%
\pgfsetlinewidth{0.000000pt}%
\definecolor{currentstroke}{rgb}{0.000000,0.000000,0.000000}%
\pgfsetstrokecolor{currentstroke}%
\pgfsetdash{}{0pt}%
\pgfpathmoveto{\pgfqpoint{1.557228in}{1.155693in}}%
\pgfpathlineto{\pgfqpoint{1.571924in}{1.174008in}}%
\pgfpathlineto{\pgfqpoint{1.572464in}{1.196352in}}%
\pgfpathlineto{\pgfqpoint{1.557751in}{1.178071in}}%
\pgfpathlineto{\pgfqpoint{1.557228in}{1.155693in}}%
\pgfpathclose%
\pgfusepath{fill}%
\end{pgfscope}%
\begin{pgfscope}%
\pgfpathrectangle{\pgfqpoint{0.407592in}{0.222904in}}{\pgfqpoint{1.411163in}{1.411163in}}%
\pgfusepath{clip}%
\pgfsetbuttcap%
\pgfsetroundjoin%
\definecolor{currentfill}{rgb}{0.168627,0.611765,0.411765}%
\pgfsetfillcolor{currentfill}%
\pgfsetfillopacity{0.180000}%
\pgfsetlinewidth{0.000000pt}%
\definecolor{currentstroke}{rgb}{0.000000,0.000000,0.000000}%
\pgfsetstrokecolor{currentstroke}%
\pgfsetdash{}{0pt}%
\pgfpathmoveto{\pgfqpoint{1.541479in}{1.093986in}}%
\pgfpathlineto{\pgfqpoint{1.556186in}{1.111101in}}%
\pgfpathlineto{\pgfqpoint{1.556706in}{1.133370in}}%
\pgfpathlineto{\pgfqpoint{1.541983in}{1.116292in}}%
\pgfpathlineto{\pgfqpoint{1.541479in}{1.093986in}}%
\pgfpathclose%
\pgfusepath{fill}%
\end{pgfscope}%
\begin{pgfscope}%
\pgfpathrectangle{\pgfqpoint{0.407592in}{0.222904in}}{\pgfqpoint{1.411163in}{1.411163in}}%
\pgfusepath{clip}%
\pgfsetbuttcap%
\pgfsetroundjoin%
\definecolor{currentfill}{rgb}{0.168627,0.611765,0.411765}%
\pgfsetfillcolor{currentfill}%
\pgfsetfillopacity{0.180000}%
\pgfsetlinewidth{0.000000pt}%
\definecolor{currentstroke}{rgb}{0.000000,0.000000,0.000000}%
\pgfsetstrokecolor{currentstroke}%
\pgfsetdash{}{0pt}%
\pgfpathmoveto{\pgfqpoint{1.525755in}{1.033693in}}%
\pgfpathlineto{\pgfqpoint{1.540475in}{1.049539in}}%
\pgfpathlineto{\pgfqpoint{1.540976in}{1.071735in}}%
\pgfpathlineto{\pgfqpoint{1.526240in}{1.055930in}}%
\pgfpathlineto{\pgfqpoint{1.525755in}{1.033693in}}%
\pgfpathclose%
\pgfusepath{fill}%
\end{pgfscope}%
\begin{pgfscope}%
\pgfpathrectangle{\pgfqpoint{0.407592in}{0.222904in}}{\pgfqpoint{1.411163in}{1.411163in}}%
\pgfusepath{clip}%
\pgfsetbuttcap%
\pgfsetroundjoin%
\definecolor{currentfill}{rgb}{0.168627,0.611765,0.411765}%
\pgfsetfillcolor{currentfill}%
\pgfsetfillopacity{0.180000}%
\pgfsetlinewidth{0.000000pt}%
\definecolor{currentstroke}{rgb}{0.000000,0.000000,0.000000}%
\pgfsetstrokecolor{currentstroke}%
\pgfsetdash{}{0pt}%
\pgfpathmoveto{\pgfqpoint{1.510055in}{0.974865in}}%
\pgfpathlineto{\pgfqpoint{1.524790in}{0.989382in}}%
\pgfpathlineto{\pgfqpoint{1.525272in}{1.011510in}}%
\pgfpathlineto{\pgfqpoint{1.510521in}{0.997038in}}%
\pgfpathlineto{\pgfqpoint{1.510055in}{0.974865in}}%
\pgfpathclose%
\pgfusepath{fill}%
\end{pgfscope}%
\begin{pgfscope}%
\pgfpathrectangle{\pgfqpoint{0.407592in}{0.222904in}}{\pgfqpoint{1.411163in}{1.411163in}}%
\pgfusepath{clip}%
\pgfsetbuttcap%
\pgfsetroundjoin%
\definecolor{currentfill}{rgb}{0.168627,0.611765,0.411765}%
\pgfsetfillcolor{currentfill}%
\pgfsetfillopacity{0.180000}%
\pgfsetlinewidth{0.000000pt}%
\definecolor{currentstroke}{rgb}{0.000000,0.000000,0.000000}%
\pgfsetstrokecolor{currentstroke}%
\pgfsetdash{}{0pt}%
\pgfpathmoveto{\pgfqpoint{1.494376in}{0.917544in}}%
\pgfpathlineto{\pgfqpoint{1.509128in}{0.930684in}}%
\pgfpathlineto{\pgfqpoint{1.509591in}{0.952748in}}%
\pgfpathlineto{\pgfqpoint{1.494822in}{0.939655in}}%
\pgfpathlineto{\pgfqpoint{1.494376in}{0.917544in}}%
\pgfpathclose%
\pgfusepath{fill}%
\end{pgfscope}%
\begin{pgfscope}%
\pgfpathrectangle{\pgfqpoint{0.407592in}{0.222904in}}{\pgfqpoint{1.411163in}{1.411163in}}%
\pgfusepath{clip}%
\pgfsetbuttcap%
\pgfsetroundjoin%
\definecolor{currentfill}{rgb}{0.168627,0.611765,0.411765}%
\pgfsetfillcolor{currentfill}%
\pgfsetfillopacity{0.180000}%
\pgfsetlinewidth{0.000000pt}%
\definecolor{currentstroke}{rgb}{0.000000,0.000000,0.000000}%
\pgfsetstrokecolor{currentstroke}%
\pgfsetdash{}{0pt}%
\pgfpathmoveto{\pgfqpoint{1.573548in}{1.241203in}}%
\pgfpathlineto{\pgfqpoint{1.588251in}{1.260610in}}%
\pgfpathlineto{\pgfqpoint{1.588812in}{1.283087in}}%
\pgfpathlineto{\pgfqpoint{1.574092in}{1.263712in}}%
\pgfpathlineto{\pgfqpoint{1.573548in}{1.241203in}}%
\pgfpathclose%
\pgfusepath{fill}%
\end{pgfscope}%
\begin{pgfscope}%
\pgfpathrectangle{\pgfqpoint{0.407592in}{0.222904in}}{\pgfqpoint{1.411163in}{1.411163in}}%
\pgfusepath{clip}%
\pgfsetbuttcap%
\pgfsetroundjoin%
\definecolor{currentfill}{rgb}{0.168627,0.611765,0.411765}%
\pgfsetfillcolor{currentfill}%
\pgfsetfillopacity{0.180000}%
\pgfsetlinewidth{0.000000pt}%
\definecolor{currentstroke}{rgb}{0.000000,0.000000,0.000000}%
\pgfsetstrokecolor{currentstroke}%
\pgfsetdash{}{0pt}%
\pgfpathmoveto{\pgfqpoint{1.557751in}{1.178071in}}%
\pgfpathlineto{\pgfqpoint{1.572464in}{1.196352in}}%
\pgfpathlineto{\pgfqpoint{1.573006in}{1.218750in}}%
\pgfpathlineto{\pgfqpoint{1.558276in}{1.200504in}}%
\pgfpathlineto{\pgfqpoint{1.557751in}{1.178071in}}%
\pgfpathclose%
\pgfusepath{fill}%
\end{pgfscope}%
\begin{pgfscope}%
\pgfpathrectangle{\pgfqpoint{0.407592in}{0.222904in}}{\pgfqpoint{1.411163in}{1.411163in}}%
\pgfusepath{clip}%
\pgfsetbuttcap%
\pgfsetroundjoin%
\definecolor{currentfill}{rgb}{0.168627,0.611765,0.411765}%
\pgfsetfillcolor{currentfill}%
\pgfsetfillopacity{0.180000}%
\pgfsetlinewidth{0.000000pt}%
\definecolor{currentstroke}{rgb}{0.000000,0.000000,0.000000}%
\pgfsetstrokecolor{currentstroke}%
\pgfsetdash{}{0pt}%
\pgfpathmoveto{\pgfqpoint{1.541983in}{1.116292in}}%
\pgfpathlineto{\pgfqpoint{1.556706in}{1.133370in}}%
\pgfpathlineto{\pgfqpoint{1.557228in}{1.155693in}}%
\pgfpathlineto{\pgfqpoint{1.542488in}{1.138653in}}%
\pgfpathlineto{\pgfqpoint{1.541983in}{1.116292in}}%
\pgfpathclose%
\pgfusepath{fill}%
\end{pgfscope}%
\begin{pgfscope}%
\pgfpathrectangle{\pgfqpoint{0.407592in}{0.222904in}}{\pgfqpoint{1.411163in}{1.411163in}}%
\pgfusepath{clip}%
\pgfsetbuttcap%
\pgfsetroundjoin%
\definecolor{currentfill}{rgb}{0.168627,0.611765,0.411765}%
\pgfsetfillcolor{currentfill}%
\pgfsetfillopacity{0.180000}%
\pgfsetlinewidth{0.000000pt}%
\definecolor{currentstroke}{rgb}{0.000000,0.000000,0.000000}%
\pgfsetstrokecolor{currentstroke}%
\pgfsetdash{}{0pt}%
\pgfpathmoveto{\pgfqpoint{1.526240in}{1.055930in}}%
\pgfpathlineto{\pgfqpoint{1.540976in}{1.071735in}}%
\pgfpathlineto{\pgfqpoint{1.541479in}{1.093986in}}%
\pgfpathlineto{\pgfqpoint{1.526726in}{1.078222in}}%
\pgfpathlineto{\pgfqpoint{1.526240in}{1.055930in}}%
\pgfpathclose%
\pgfusepath{fill}%
\end{pgfscope}%
\begin{pgfscope}%
\pgfpathrectangle{\pgfqpoint{0.407592in}{0.222904in}}{\pgfqpoint{1.411163in}{1.411163in}}%
\pgfusepath{clip}%
\pgfsetbuttcap%
\pgfsetroundjoin%
\definecolor{currentfill}{rgb}{0.168627,0.611765,0.411765}%
\pgfsetfillcolor{currentfill}%
\pgfsetfillopacity{0.180000}%
\pgfsetlinewidth{0.000000pt}%
\definecolor{currentstroke}{rgb}{0.000000,0.000000,0.000000}%
\pgfsetstrokecolor{currentstroke}%
\pgfsetdash{}{0pt}%
\pgfpathmoveto{\pgfqpoint{1.510521in}{0.997038in}}%
\pgfpathlineto{\pgfqpoint{1.525272in}{1.011510in}}%
\pgfpathlineto{\pgfqpoint{1.525755in}{1.033693in}}%
\pgfpathlineto{\pgfqpoint{1.510987in}{1.019264in}}%
\pgfpathlineto{\pgfqpoint{1.510521in}{0.997038in}}%
\pgfpathclose%
\pgfusepath{fill}%
\end{pgfscope}%
\begin{pgfscope}%
\pgfpathrectangle{\pgfqpoint{0.407592in}{0.222904in}}{\pgfqpoint{1.411163in}{1.411163in}}%
\pgfusepath{clip}%
\pgfsetbuttcap%
\pgfsetroundjoin%
\definecolor{currentfill}{rgb}{0.168627,0.611765,0.411765}%
\pgfsetfillcolor{currentfill}%
\pgfsetfillopacity{0.180000}%
\pgfsetlinewidth{0.000000pt}%
\definecolor{currentstroke}{rgb}{0.000000,0.000000,0.000000}%
\pgfsetstrokecolor{currentstroke}%
\pgfsetdash{}{0pt}%
\pgfpathmoveto{\pgfqpoint{1.494822in}{0.939655in}}%
\pgfpathlineto{\pgfqpoint{1.509591in}{0.952748in}}%
\pgfpathlineto{\pgfqpoint{1.510055in}{0.974865in}}%
\pgfpathlineto{\pgfqpoint{1.495270in}{0.961820in}}%
\pgfpathlineto{\pgfqpoint{1.494822in}{0.939655in}}%
\pgfpathclose%
\pgfusepath{fill}%
\end{pgfscope}%
\begin{pgfscope}%
\pgfpathrectangle{\pgfqpoint{0.407592in}{0.222904in}}{\pgfqpoint{1.411163in}{1.411163in}}%
\pgfusepath{clip}%
\pgfsetbuttcap%
\pgfsetroundjoin%
\definecolor{currentfill}{rgb}{0.168627,0.611765,0.411765}%
\pgfsetfillcolor{currentfill}%
\pgfsetfillopacity{0.180000}%
\pgfsetlinewidth{0.000000pt}%
\definecolor{currentstroke}{rgb}{0.000000,0.000000,0.000000}%
\pgfsetstrokecolor{currentstroke}%
\pgfsetdash{}{0pt}%
\pgfpathmoveto{\pgfqpoint{1.574092in}{1.263712in}}%
\pgfpathlineto{\pgfqpoint{1.588812in}{1.283087in}}%
\pgfpathlineto{\pgfqpoint{1.589374in}{1.305619in}}%
\pgfpathlineto{\pgfqpoint{1.574637in}{1.286277in}}%
\pgfpathlineto{\pgfqpoint{1.574092in}{1.263712in}}%
\pgfpathclose%
\pgfusepath{fill}%
\end{pgfscope}%
\begin{pgfscope}%
\pgfpathrectangle{\pgfqpoint{0.407592in}{0.222904in}}{\pgfqpoint{1.411163in}{1.411163in}}%
\pgfusepath{clip}%
\pgfsetbuttcap%
\pgfsetroundjoin%
\definecolor{currentfill}{rgb}{0.168627,0.611765,0.411765}%
\pgfsetfillcolor{currentfill}%
\pgfsetfillopacity{0.180000}%
\pgfsetlinewidth{0.000000pt}%
\definecolor{currentstroke}{rgb}{0.000000,0.000000,0.000000}%
\pgfsetstrokecolor{currentstroke}%
\pgfsetdash{}{0pt}%
\pgfpathmoveto{\pgfqpoint{1.558276in}{1.200504in}}%
\pgfpathlineto{\pgfqpoint{1.573006in}{1.218750in}}%
\pgfpathlineto{\pgfqpoint{1.573548in}{1.241203in}}%
\pgfpathlineto{\pgfqpoint{1.558802in}{1.222993in}}%
\pgfpathlineto{\pgfqpoint{1.558276in}{1.200504in}}%
\pgfpathclose%
\pgfusepath{fill}%
\end{pgfscope}%
\begin{pgfscope}%
\pgfpathrectangle{\pgfqpoint{0.407592in}{0.222904in}}{\pgfqpoint{1.411163in}{1.411163in}}%
\pgfusepath{clip}%
\pgfsetbuttcap%
\pgfsetroundjoin%
\definecolor{currentfill}{rgb}{0.168627,0.611765,0.411765}%
\pgfsetfillcolor{currentfill}%
\pgfsetfillopacity{0.180000}%
\pgfsetlinewidth{0.000000pt}%
\definecolor{currentstroke}{rgb}{0.000000,0.000000,0.000000}%
\pgfsetstrokecolor{currentstroke}%
\pgfsetdash{}{0pt}%
\pgfpathmoveto{\pgfqpoint{1.542488in}{1.138653in}}%
\pgfpathlineto{\pgfqpoint{1.557228in}{1.155693in}}%
\pgfpathlineto{\pgfqpoint{1.557751in}{1.178071in}}%
\pgfpathlineto{\pgfqpoint{1.542994in}{1.161069in}}%
\pgfpathlineto{\pgfqpoint{1.542488in}{1.138653in}}%
\pgfpathclose%
\pgfusepath{fill}%
\end{pgfscope}%
\begin{pgfscope}%
\pgfpathrectangle{\pgfqpoint{0.407592in}{0.222904in}}{\pgfqpoint{1.411163in}{1.411163in}}%
\pgfusepath{clip}%
\pgfsetbuttcap%
\pgfsetroundjoin%
\definecolor{currentfill}{rgb}{0.168627,0.611765,0.411765}%
\pgfsetfillcolor{currentfill}%
\pgfsetfillopacity{0.180000}%
\pgfsetlinewidth{0.000000pt}%
\definecolor{currentstroke}{rgb}{0.000000,0.000000,0.000000}%
\pgfsetstrokecolor{currentstroke}%
\pgfsetdash{}{0pt}%
\pgfpathmoveto{\pgfqpoint{1.526726in}{1.078222in}}%
\pgfpathlineto{\pgfqpoint{1.541479in}{1.093986in}}%
\pgfpathlineto{\pgfqpoint{1.541983in}{1.116292in}}%
\pgfpathlineto{\pgfqpoint{1.527213in}{1.100569in}}%
\pgfpathlineto{\pgfqpoint{1.526726in}{1.078222in}}%
\pgfpathclose%
\pgfusepath{fill}%
\end{pgfscope}%
\begin{pgfscope}%
\pgfpathrectangle{\pgfqpoint{0.407592in}{0.222904in}}{\pgfqpoint{1.411163in}{1.411163in}}%
\pgfusepath{clip}%
\pgfsetbuttcap%
\pgfsetroundjoin%
\definecolor{currentfill}{rgb}{0.168627,0.611765,0.411765}%
\pgfsetfillcolor{currentfill}%
\pgfsetfillopacity{0.180000}%
\pgfsetlinewidth{0.000000pt}%
\definecolor{currentstroke}{rgb}{0.000000,0.000000,0.000000}%
\pgfsetstrokecolor{currentstroke}%
\pgfsetdash{}{0pt}%
\pgfpathmoveto{\pgfqpoint{1.510987in}{1.019264in}}%
\pgfpathlineto{\pgfqpoint{1.525755in}{1.033693in}}%
\pgfpathlineto{\pgfqpoint{1.526240in}{1.055930in}}%
\pgfpathlineto{\pgfqpoint{1.511455in}{1.041546in}}%
\pgfpathlineto{\pgfqpoint{1.510987in}{1.019264in}}%
\pgfpathclose%
\pgfusepath{fill}%
\end{pgfscope}%
\begin{pgfscope}%
\pgfpathrectangle{\pgfqpoint{0.407592in}{0.222904in}}{\pgfqpoint{1.411163in}{1.411163in}}%
\pgfusepath{clip}%
\pgfsetbuttcap%
\pgfsetroundjoin%
\definecolor{currentfill}{rgb}{0.168627,0.611765,0.411765}%
\pgfsetfillcolor{currentfill}%
\pgfsetfillopacity{0.180000}%
\pgfsetlinewidth{0.000000pt}%
\definecolor{currentstroke}{rgb}{0.000000,0.000000,0.000000}%
\pgfsetstrokecolor{currentstroke}%
\pgfsetdash{}{0pt}%
\pgfpathmoveto{\pgfqpoint{1.495270in}{0.961820in}}%
\pgfpathlineto{\pgfqpoint{1.510055in}{0.974865in}}%
\pgfpathlineto{\pgfqpoint{1.510521in}{0.997038in}}%
\pgfpathlineto{\pgfqpoint{1.495718in}{0.984040in}}%
\pgfpathlineto{\pgfqpoint{1.495270in}{0.961820in}}%
\pgfpathclose%
\pgfusepath{fill}%
\end{pgfscope}%
\begin{pgfscope}%
\pgfpathrectangle{\pgfqpoint{0.407592in}{0.222904in}}{\pgfqpoint{1.411163in}{1.411163in}}%
\pgfusepath{clip}%
\pgfsetbuttcap%
\pgfsetroundjoin%
\definecolor{currentfill}{rgb}{0.168627,0.611765,0.411765}%
\pgfsetfillcolor{currentfill}%
\pgfsetfillopacity{0.180000}%
\pgfsetlinewidth{0.000000pt}%
\definecolor{currentstroke}{rgb}{0.000000,0.000000,0.000000}%
\pgfsetstrokecolor{currentstroke}%
\pgfsetdash{}{0pt}%
\pgfpathmoveto{\pgfqpoint{1.558802in}{1.222993in}}%
\pgfpathlineto{\pgfqpoint{1.573548in}{1.241203in}}%
\pgfpathlineto{\pgfqpoint{1.574092in}{1.263712in}}%
\pgfpathlineto{\pgfqpoint{1.559329in}{1.245537in}}%
\pgfpathlineto{\pgfqpoint{1.558802in}{1.222993in}}%
\pgfpathclose%
\pgfusepath{fill}%
\end{pgfscope}%
\begin{pgfscope}%
\pgfpathrectangle{\pgfqpoint{0.407592in}{0.222904in}}{\pgfqpoint{1.411163in}{1.411163in}}%
\pgfusepath{clip}%
\pgfsetbuttcap%
\pgfsetroundjoin%
\definecolor{currentfill}{rgb}{0.168627,0.611765,0.411765}%
\pgfsetfillcolor{currentfill}%
\pgfsetfillopacity{0.180000}%
\pgfsetlinewidth{0.000000pt}%
\definecolor{currentstroke}{rgb}{0.000000,0.000000,0.000000}%
\pgfsetstrokecolor{currentstroke}%
\pgfsetdash{}{0pt}%
\pgfpathmoveto{\pgfqpoint{1.542994in}{1.161069in}}%
\pgfpathlineto{\pgfqpoint{1.557751in}{1.178071in}}%
\pgfpathlineto{\pgfqpoint{1.558276in}{1.200504in}}%
\pgfpathlineto{\pgfqpoint{1.543502in}{1.183540in}}%
\pgfpathlineto{\pgfqpoint{1.542994in}{1.161069in}}%
\pgfpathclose%
\pgfusepath{fill}%
\end{pgfscope}%
\begin{pgfscope}%
\pgfpathrectangle{\pgfqpoint{0.407592in}{0.222904in}}{\pgfqpoint{1.411163in}{1.411163in}}%
\pgfusepath{clip}%
\pgfsetbuttcap%
\pgfsetroundjoin%
\definecolor{currentfill}{rgb}{0.168627,0.611765,0.411765}%
\pgfsetfillcolor{currentfill}%
\pgfsetfillopacity{0.180000}%
\pgfsetlinewidth{0.000000pt}%
\definecolor{currentstroke}{rgb}{0.000000,0.000000,0.000000}%
\pgfsetstrokecolor{currentstroke}%
\pgfsetdash{}{0pt}%
\pgfpathmoveto{\pgfqpoint{1.479570in}{0.905913in}}%
\pgfpathlineto{\pgfqpoint{1.494376in}{0.917544in}}%
\pgfpathlineto{\pgfqpoint{1.494822in}{0.939655in}}%
\pgfpathlineto{\pgfqpoint{1.479999in}{0.928075in}}%
\pgfpathlineto{\pgfqpoint{1.479570in}{0.905913in}}%
\pgfpathclose%
\pgfusepath{fill}%
\end{pgfscope}%
\begin{pgfscope}%
\pgfpathrectangle{\pgfqpoint{0.407592in}{0.222904in}}{\pgfqpoint{1.411163in}{1.411163in}}%
\pgfusepath{clip}%
\pgfsetbuttcap%
\pgfsetroundjoin%
\definecolor{currentfill}{rgb}{0.168627,0.611765,0.411765}%
\pgfsetfillcolor{currentfill}%
\pgfsetfillopacity{0.180000}%
\pgfsetlinewidth{0.000000pt}%
\definecolor{currentstroke}{rgb}{0.000000,0.000000,0.000000}%
\pgfsetstrokecolor{currentstroke}%
\pgfsetdash{}{0pt}%
\pgfpathmoveto{\pgfqpoint{1.527213in}{1.100569in}}%
\pgfpathlineto{\pgfqpoint{1.541983in}{1.116292in}}%
\pgfpathlineto{\pgfqpoint{1.542488in}{1.138653in}}%
\pgfpathlineto{\pgfqpoint{1.527701in}{1.122971in}}%
\pgfpathlineto{\pgfqpoint{1.527213in}{1.100569in}}%
\pgfpathclose%
\pgfusepath{fill}%
\end{pgfscope}%
\begin{pgfscope}%
\pgfpathrectangle{\pgfqpoint{0.407592in}{0.222904in}}{\pgfqpoint{1.411163in}{1.411163in}}%
\pgfusepath{clip}%
\pgfsetbuttcap%
\pgfsetroundjoin%
\definecolor{currentfill}{rgb}{0.168627,0.611765,0.411765}%
\pgfsetfillcolor{currentfill}%
\pgfsetfillopacity{0.180000}%
\pgfsetlinewidth{0.000000pt}%
\definecolor{currentstroke}{rgb}{0.000000,0.000000,0.000000}%
\pgfsetstrokecolor{currentstroke}%
\pgfsetdash{}{0pt}%
\pgfpathmoveto{\pgfqpoint{1.511455in}{1.041546in}}%
\pgfpathlineto{\pgfqpoint{1.526240in}{1.055930in}}%
\pgfpathlineto{\pgfqpoint{1.526726in}{1.078222in}}%
\pgfpathlineto{\pgfqpoint{1.511924in}{1.063883in}}%
\pgfpathlineto{\pgfqpoint{1.511455in}{1.041546in}}%
\pgfpathclose%
\pgfusepath{fill}%
\end{pgfscope}%
\begin{pgfscope}%
\pgfpathrectangle{\pgfqpoint{0.407592in}{0.222904in}}{\pgfqpoint{1.411163in}{1.411163in}}%
\pgfusepath{clip}%
\pgfsetbuttcap%
\pgfsetroundjoin%
\definecolor{currentfill}{rgb}{0.168627,0.611765,0.411765}%
\pgfsetfillcolor{currentfill}%
\pgfsetfillopacity{0.180000}%
\pgfsetlinewidth{0.000000pt}%
\definecolor{currentstroke}{rgb}{0.000000,0.000000,0.000000}%
\pgfsetstrokecolor{currentstroke}%
\pgfsetdash{}{0pt}%
\pgfpathmoveto{\pgfqpoint{1.495718in}{0.984040in}}%
\pgfpathlineto{\pgfqpoint{1.510521in}{0.997038in}}%
\pgfpathlineto{\pgfqpoint{1.510987in}{1.019264in}}%
\pgfpathlineto{\pgfqpoint{1.496168in}{1.006315in}}%
\pgfpathlineto{\pgfqpoint{1.495718in}{0.984040in}}%
\pgfpathclose%
\pgfusepath{fill}%
\end{pgfscope}%
\begin{pgfscope}%
\pgfpathrectangle{\pgfqpoint{0.407592in}{0.222904in}}{\pgfqpoint{1.411163in}{1.411163in}}%
\pgfusepath{clip}%
\pgfsetbuttcap%
\pgfsetroundjoin%
\definecolor{currentfill}{rgb}{0.168627,0.611765,0.411765}%
\pgfsetfillcolor{currentfill}%
\pgfsetfillopacity{0.180000}%
\pgfsetlinewidth{0.000000pt}%
\definecolor{currentstroke}{rgb}{0.000000,0.000000,0.000000}%
\pgfsetstrokecolor{currentstroke}%
\pgfsetdash{}{0pt}%
\pgfpathmoveto{\pgfqpoint{1.559329in}{1.245537in}}%
\pgfpathlineto{\pgfqpoint{1.574092in}{1.263712in}}%
\pgfpathlineto{\pgfqpoint{1.574637in}{1.286277in}}%
\pgfpathlineto{\pgfqpoint{1.559857in}{1.268137in}}%
\pgfpathlineto{\pgfqpoint{1.559329in}{1.245537in}}%
\pgfpathclose%
\pgfusepath{fill}%
\end{pgfscope}%
\begin{pgfscope}%
\pgfpathrectangle{\pgfqpoint{0.407592in}{0.222904in}}{\pgfqpoint{1.411163in}{1.411163in}}%
\pgfusepath{clip}%
\pgfsetbuttcap%
\pgfsetroundjoin%
\definecolor{currentfill}{rgb}{0.168627,0.611765,0.411765}%
\pgfsetfillcolor{currentfill}%
\pgfsetfillopacity{0.180000}%
\pgfsetlinewidth{0.000000pt}%
\definecolor{currentstroke}{rgb}{0.000000,0.000000,0.000000}%
\pgfsetstrokecolor{currentstroke}%
\pgfsetdash{}{0pt}%
\pgfpathmoveto{\pgfqpoint{1.543502in}{1.183540in}}%
\pgfpathlineto{\pgfqpoint{1.558276in}{1.200504in}}%
\pgfpathlineto{\pgfqpoint{1.558802in}{1.222993in}}%
\pgfpathlineto{\pgfqpoint{1.544011in}{1.206067in}}%
\pgfpathlineto{\pgfqpoint{1.543502in}{1.183540in}}%
\pgfpathclose%
\pgfusepath{fill}%
\end{pgfscope}%
\begin{pgfscope}%
\pgfpathrectangle{\pgfqpoint{0.407592in}{0.222904in}}{\pgfqpoint{1.411163in}{1.411163in}}%
\pgfusepath{clip}%
\pgfsetbuttcap%
\pgfsetroundjoin%
\definecolor{currentfill}{rgb}{0.168627,0.611765,0.411765}%
\pgfsetfillcolor{currentfill}%
\pgfsetfillopacity{0.180000}%
\pgfsetlinewidth{0.000000pt}%
\definecolor{currentstroke}{rgb}{0.000000,0.000000,0.000000}%
\pgfsetstrokecolor{currentstroke}%
\pgfsetdash{}{0pt}%
\pgfpathmoveto{\pgfqpoint{1.479999in}{0.928075in}}%
\pgfpathlineto{\pgfqpoint{1.494822in}{0.939655in}}%
\pgfpathlineto{\pgfqpoint{1.495270in}{0.961820in}}%
\pgfpathlineto{\pgfqpoint{1.480429in}{0.950293in}}%
\pgfpathlineto{\pgfqpoint{1.479999in}{0.928075in}}%
\pgfpathclose%
\pgfusepath{fill}%
\end{pgfscope}%
\begin{pgfscope}%
\pgfpathrectangle{\pgfqpoint{0.407592in}{0.222904in}}{\pgfqpoint{1.411163in}{1.411163in}}%
\pgfusepath{clip}%
\pgfsetbuttcap%
\pgfsetroundjoin%
\definecolor{currentfill}{rgb}{0.168627,0.611765,0.411765}%
\pgfsetfillcolor{currentfill}%
\pgfsetfillopacity{0.180000}%
\pgfsetlinewidth{0.000000pt}%
\definecolor{currentstroke}{rgb}{0.000000,0.000000,0.000000}%
\pgfsetstrokecolor{currentstroke}%
\pgfsetdash{}{0pt}%
\pgfpathmoveto{\pgfqpoint{1.527701in}{1.122971in}}%
\pgfpathlineto{\pgfqpoint{1.542488in}{1.138653in}}%
\pgfpathlineto{\pgfqpoint{1.542994in}{1.161069in}}%
\pgfpathlineto{\pgfqpoint{1.528190in}{1.145429in}}%
\pgfpathlineto{\pgfqpoint{1.527701in}{1.122971in}}%
\pgfpathclose%
\pgfusepath{fill}%
\end{pgfscope}%
\begin{pgfscope}%
\pgfpathrectangle{\pgfqpoint{0.407592in}{0.222904in}}{\pgfqpoint{1.411163in}{1.411163in}}%
\pgfusepath{clip}%
\pgfsetbuttcap%
\pgfsetroundjoin%
\definecolor{currentfill}{rgb}{0.168627,0.611765,0.411765}%
\pgfsetfillcolor{currentfill}%
\pgfsetfillopacity{0.180000}%
\pgfsetlinewidth{0.000000pt}%
\definecolor{currentstroke}{rgb}{0.000000,0.000000,0.000000}%
\pgfsetstrokecolor{currentstroke}%
\pgfsetdash{}{0pt}%
\pgfpathmoveto{\pgfqpoint{1.511924in}{1.063883in}}%
\pgfpathlineto{\pgfqpoint{1.526726in}{1.078222in}}%
\pgfpathlineto{\pgfqpoint{1.527213in}{1.100569in}}%
\pgfpathlineto{\pgfqpoint{1.512394in}{1.086275in}}%
\pgfpathlineto{\pgfqpoint{1.511924in}{1.063883in}}%
\pgfpathclose%
\pgfusepath{fill}%
\end{pgfscope}%
\begin{pgfscope}%
\pgfpathrectangle{\pgfqpoint{0.407592in}{0.222904in}}{\pgfqpoint{1.411163in}{1.411163in}}%
\pgfusepath{clip}%
\pgfsetbuttcap%
\pgfsetroundjoin%
\definecolor{currentfill}{rgb}{0.168627,0.611765,0.411765}%
\pgfsetfillcolor{currentfill}%
\pgfsetfillopacity{0.180000}%
\pgfsetlinewidth{0.000000pt}%
\definecolor{currentstroke}{rgb}{0.000000,0.000000,0.000000}%
\pgfsetstrokecolor{currentstroke}%
\pgfsetdash{}{0pt}%
\pgfpathmoveto{\pgfqpoint{1.496168in}{1.006315in}}%
\pgfpathlineto{\pgfqpoint{1.510987in}{1.019264in}}%
\pgfpathlineto{\pgfqpoint{1.511455in}{1.041546in}}%
\pgfpathlineto{\pgfqpoint{1.496619in}{1.028646in}}%
\pgfpathlineto{\pgfqpoint{1.496168in}{1.006315in}}%
\pgfpathclose%
\pgfusepath{fill}%
\end{pgfscope}%
\begin{pgfscope}%
\pgfpathrectangle{\pgfqpoint{0.407592in}{0.222904in}}{\pgfqpoint{1.411163in}{1.411163in}}%
\pgfusepath{clip}%
\pgfsetbuttcap%
\pgfsetroundjoin%
\definecolor{currentfill}{rgb}{0.168627,0.611765,0.411765}%
\pgfsetfillcolor{currentfill}%
\pgfsetfillopacity{0.180000}%
\pgfsetlinewidth{0.000000pt}%
\definecolor{currentstroke}{rgb}{0.000000,0.000000,0.000000}%
\pgfsetstrokecolor{currentstroke}%
\pgfsetdash{}{0pt}%
\pgfpathmoveto{\pgfqpoint{1.544011in}{1.206067in}}%
\pgfpathlineto{\pgfqpoint{1.558802in}{1.222993in}}%
\pgfpathlineto{\pgfqpoint{1.559329in}{1.245537in}}%
\pgfpathlineto{\pgfqpoint{1.544521in}{1.228650in}}%
\pgfpathlineto{\pgfqpoint{1.544011in}{1.206067in}}%
\pgfpathclose%
\pgfusepath{fill}%
\end{pgfscope}%
\begin{pgfscope}%
\pgfpathrectangle{\pgfqpoint{0.407592in}{0.222904in}}{\pgfqpoint{1.411163in}{1.411163in}}%
\pgfusepath{clip}%
\pgfsetbuttcap%
\pgfsetroundjoin%
\definecolor{currentfill}{rgb}{0.168627,0.611765,0.411765}%
\pgfsetfillcolor{currentfill}%
\pgfsetfillopacity{0.180000}%
\pgfsetlinewidth{0.000000pt}%
\definecolor{currentstroke}{rgb}{0.000000,0.000000,0.000000}%
\pgfsetstrokecolor{currentstroke}%
\pgfsetdash{}{0pt}%
\pgfpathmoveto{\pgfqpoint{1.480429in}{0.950293in}}%
\pgfpathlineto{\pgfqpoint{1.495270in}{0.961820in}}%
\pgfpathlineto{\pgfqpoint{1.495718in}{0.984040in}}%
\pgfpathlineto{\pgfqpoint{1.480861in}{0.972565in}}%
\pgfpathlineto{\pgfqpoint{1.480429in}{0.950293in}}%
\pgfpathclose%
\pgfusepath{fill}%
\end{pgfscope}%
\begin{pgfscope}%
\pgfpathrectangle{\pgfqpoint{0.407592in}{0.222904in}}{\pgfqpoint{1.411163in}{1.411163in}}%
\pgfusepath{clip}%
\pgfsetbuttcap%
\pgfsetroundjoin%
\definecolor{currentfill}{rgb}{0.168627,0.611765,0.411765}%
\pgfsetfillcolor{currentfill}%
\pgfsetfillopacity{0.180000}%
\pgfsetlinewidth{0.000000pt}%
\definecolor{currentstroke}{rgb}{0.000000,0.000000,0.000000}%
\pgfsetstrokecolor{currentstroke}%
\pgfsetdash{}{0pt}%
\pgfpathmoveto{\pgfqpoint{1.528190in}{1.145429in}}%
\pgfpathlineto{\pgfqpoint{1.542994in}{1.161069in}}%
\pgfpathlineto{\pgfqpoint{1.543502in}{1.183540in}}%
\pgfpathlineto{\pgfqpoint{1.528681in}{1.167943in}}%
\pgfpathlineto{\pgfqpoint{1.528190in}{1.145429in}}%
\pgfpathclose%
\pgfusepath{fill}%
\end{pgfscope}%
\begin{pgfscope}%
\pgfpathrectangle{\pgfqpoint{0.407592in}{0.222904in}}{\pgfqpoint{1.411163in}{1.411163in}}%
\pgfusepath{clip}%
\pgfsetbuttcap%
\pgfsetroundjoin%
\definecolor{currentfill}{rgb}{0.168627,0.611765,0.411765}%
\pgfsetfillcolor{currentfill}%
\pgfsetfillopacity{0.180000}%
\pgfsetlinewidth{0.000000pt}%
\definecolor{currentstroke}{rgb}{0.000000,0.000000,0.000000}%
\pgfsetstrokecolor{currentstroke}%
\pgfsetdash{}{0pt}%
\pgfpathmoveto{\pgfqpoint{1.512394in}{1.086275in}}%
\pgfpathlineto{\pgfqpoint{1.527213in}{1.100569in}}%
\pgfpathlineto{\pgfqpoint{1.527701in}{1.122971in}}%
\pgfpathlineto{\pgfqpoint{1.512865in}{1.108723in}}%
\pgfpathlineto{\pgfqpoint{1.512394in}{1.086275in}}%
\pgfpathclose%
\pgfusepath{fill}%
\end{pgfscope}%
\begin{pgfscope}%
\pgfpathrectangle{\pgfqpoint{0.407592in}{0.222904in}}{\pgfqpoint{1.411163in}{1.411163in}}%
\pgfusepath{clip}%
\pgfsetbuttcap%
\pgfsetroundjoin%
\definecolor{currentfill}{rgb}{0.168627,0.611765,0.411765}%
\pgfsetfillcolor{currentfill}%
\pgfsetfillopacity{0.180000}%
\pgfsetlinewidth{0.000000pt}%
\definecolor{currentstroke}{rgb}{0.000000,0.000000,0.000000}%
\pgfsetstrokecolor{currentstroke}%
\pgfsetdash{}{0pt}%
\pgfpathmoveto{\pgfqpoint{1.464705in}{0.895822in}}%
\pgfpathlineto{\pgfqpoint{1.479570in}{0.905913in}}%
\pgfpathlineto{\pgfqpoint{1.479999in}{0.928075in}}%
\pgfpathlineto{\pgfqpoint{1.465117in}{0.918040in}}%
\pgfpathlineto{\pgfqpoint{1.464705in}{0.895822in}}%
\pgfpathclose%
\pgfusepath{fill}%
\end{pgfscope}%
\begin{pgfscope}%
\pgfpathrectangle{\pgfqpoint{0.407592in}{0.222904in}}{\pgfqpoint{1.411163in}{1.411163in}}%
\pgfusepath{clip}%
\pgfsetbuttcap%
\pgfsetroundjoin%
\definecolor{currentfill}{rgb}{0.168627,0.611765,0.411765}%
\pgfsetfillcolor{currentfill}%
\pgfsetfillopacity{0.180000}%
\pgfsetlinewidth{0.000000pt}%
\definecolor{currentstroke}{rgb}{0.000000,0.000000,0.000000}%
\pgfsetstrokecolor{currentstroke}%
\pgfsetdash{}{0pt}%
\pgfpathmoveto{\pgfqpoint{1.496619in}{1.028646in}}%
\pgfpathlineto{\pgfqpoint{1.511455in}{1.041546in}}%
\pgfpathlineto{\pgfqpoint{1.511924in}{1.063883in}}%
\pgfpathlineto{\pgfqpoint{1.497070in}{1.051031in}}%
\pgfpathlineto{\pgfqpoint{1.496619in}{1.028646in}}%
\pgfpathclose%
\pgfusepath{fill}%
\end{pgfscope}%
\begin{pgfscope}%
\pgfpathrectangle{\pgfqpoint{0.407592in}{0.222904in}}{\pgfqpoint{1.411163in}{1.411163in}}%
\pgfusepath{clip}%
\pgfsetbuttcap%
\pgfsetroundjoin%
\definecolor{currentfill}{rgb}{0.168627,0.611765,0.411765}%
\pgfsetfillcolor{currentfill}%
\pgfsetfillopacity{0.180000}%
\pgfsetlinewidth{0.000000pt}%
\definecolor{currentstroke}{rgb}{0.000000,0.000000,0.000000}%
\pgfsetstrokecolor{currentstroke}%
\pgfsetdash{}{0pt}%
\pgfpathmoveto{\pgfqpoint{1.544521in}{1.228650in}}%
\pgfpathlineto{\pgfqpoint{1.559329in}{1.245537in}}%
\pgfpathlineto{\pgfqpoint{1.559857in}{1.268137in}}%
\pgfpathlineto{\pgfqpoint{1.545032in}{1.251289in}}%
\pgfpathlineto{\pgfqpoint{1.544521in}{1.228650in}}%
\pgfpathclose%
\pgfusepath{fill}%
\end{pgfscope}%
\begin{pgfscope}%
\pgfpathrectangle{\pgfqpoint{0.407592in}{0.222904in}}{\pgfqpoint{1.411163in}{1.411163in}}%
\pgfusepath{clip}%
\pgfsetbuttcap%
\pgfsetroundjoin%
\definecolor{currentfill}{rgb}{0.168627,0.611765,0.411765}%
\pgfsetfillcolor{currentfill}%
\pgfsetfillopacity{0.180000}%
\pgfsetlinewidth{0.000000pt}%
\definecolor{currentstroke}{rgb}{0.000000,0.000000,0.000000}%
\pgfsetstrokecolor{currentstroke}%
\pgfsetdash{}{0pt}%
\pgfpathmoveto{\pgfqpoint{1.480861in}{0.972565in}}%
\pgfpathlineto{\pgfqpoint{1.495718in}{0.984040in}}%
\pgfpathlineto{\pgfqpoint{1.496168in}{1.006315in}}%
\pgfpathlineto{\pgfqpoint{1.481293in}{0.994893in}}%
\pgfpathlineto{\pgfqpoint{1.480861in}{0.972565in}}%
\pgfpathclose%
\pgfusepath{fill}%
\end{pgfscope}%
\begin{pgfscope}%
\pgfpathrectangle{\pgfqpoint{0.407592in}{0.222904in}}{\pgfqpoint{1.411163in}{1.411163in}}%
\pgfusepath{clip}%
\pgfsetbuttcap%
\pgfsetroundjoin%
\definecolor{currentfill}{rgb}{0.168627,0.611765,0.411765}%
\pgfsetfillcolor{currentfill}%
\pgfsetfillopacity{0.180000}%
\pgfsetlinewidth{0.000000pt}%
\definecolor{currentstroke}{rgb}{0.000000,0.000000,0.000000}%
\pgfsetstrokecolor{currentstroke}%
\pgfsetdash{}{0pt}%
\pgfpathmoveto{\pgfqpoint{1.528681in}{1.167943in}}%
\pgfpathlineto{\pgfqpoint{1.543502in}{1.183540in}}%
\pgfpathlineto{\pgfqpoint{1.544011in}{1.206067in}}%
\pgfpathlineto{\pgfqpoint{1.529173in}{1.190512in}}%
\pgfpathlineto{\pgfqpoint{1.528681in}{1.167943in}}%
\pgfpathclose%
\pgfusepath{fill}%
\end{pgfscope}%
\begin{pgfscope}%
\pgfpathrectangle{\pgfqpoint{0.407592in}{0.222904in}}{\pgfqpoint{1.411163in}{1.411163in}}%
\pgfusepath{clip}%
\pgfsetbuttcap%
\pgfsetroundjoin%
\definecolor{currentfill}{rgb}{0.168627,0.611765,0.411765}%
\pgfsetfillcolor{currentfill}%
\pgfsetfillopacity{0.180000}%
\pgfsetlinewidth{0.000000pt}%
\definecolor{currentstroke}{rgb}{0.000000,0.000000,0.000000}%
\pgfsetstrokecolor{currentstroke}%
\pgfsetdash{}{0pt}%
\pgfpathmoveto{\pgfqpoint{1.512865in}{1.108723in}}%
\pgfpathlineto{\pgfqpoint{1.527701in}{1.122971in}}%
\pgfpathlineto{\pgfqpoint{1.528190in}{1.145429in}}%
\pgfpathlineto{\pgfqpoint{1.513338in}{1.131226in}}%
\pgfpathlineto{\pgfqpoint{1.512865in}{1.108723in}}%
\pgfpathclose%
\pgfusepath{fill}%
\end{pgfscope}%
\begin{pgfscope}%
\pgfpathrectangle{\pgfqpoint{0.407592in}{0.222904in}}{\pgfqpoint{1.411163in}{1.411163in}}%
\pgfusepath{clip}%
\pgfsetbuttcap%
\pgfsetroundjoin%
\definecolor{currentfill}{rgb}{0.168627,0.611765,0.411765}%
\pgfsetfillcolor{currentfill}%
\pgfsetfillopacity{0.180000}%
\pgfsetlinewidth{0.000000pt}%
\definecolor{currentstroke}{rgb}{0.000000,0.000000,0.000000}%
\pgfsetstrokecolor{currentstroke}%
\pgfsetdash{}{0pt}%
\pgfpathmoveto{\pgfqpoint{1.465117in}{0.918040in}}%
\pgfpathlineto{\pgfqpoint{1.479999in}{0.928075in}}%
\pgfpathlineto{\pgfqpoint{1.480429in}{0.950293in}}%
\pgfpathlineto{\pgfqpoint{1.465530in}{0.940313in}}%
\pgfpathlineto{\pgfqpoint{1.465117in}{0.918040in}}%
\pgfpathclose%
\pgfusepath{fill}%
\end{pgfscope}%
\begin{pgfscope}%
\pgfpathrectangle{\pgfqpoint{0.407592in}{0.222904in}}{\pgfqpoint{1.411163in}{1.411163in}}%
\pgfusepath{clip}%
\pgfsetbuttcap%
\pgfsetroundjoin%
\definecolor{currentfill}{rgb}{0.168627,0.611765,0.411765}%
\pgfsetfillcolor{currentfill}%
\pgfsetfillopacity{0.180000}%
\pgfsetlinewidth{0.000000pt}%
\definecolor{currentstroke}{rgb}{0.000000,0.000000,0.000000}%
\pgfsetstrokecolor{currentstroke}%
\pgfsetdash{}{0pt}%
\pgfpathmoveto{\pgfqpoint{1.497070in}{1.051031in}}%
\pgfpathlineto{\pgfqpoint{1.511924in}{1.063883in}}%
\pgfpathlineto{\pgfqpoint{1.512394in}{1.086275in}}%
\pgfpathlineto{\pgfqpoint{1.497523in}{1.073473in}}%
\pgfpathlineto{\pgfqpoint{1.497070in}{1.051031in}}%
\pgfpathclose%
\pgfusepath{fill}%
\end{pgfscope}%
\begin{pgfscope}%
\pgfpathrectangle{\pgfqpoint{0.407592in}{0.222904in}}{\pgfqpoint{1.411163in}{1.411163in}}%
\pgfusepath{clip}%
\pgfsetbuttcap%
\pgfsetroundjoin%
\definecolor{currentfill}{rgb}{0.168627,0.611765,0.411765}%
\pgfsetfillcolor{currentfill}%
\pgfsetfillopacity{0.180000}%
\pgfsetlinewidth{0.000000pt}%
\definecolor{currentstroke}{rgb}{0.000000,0.000000,0.000000}%
\pgfsetstrokecolor{currentstroke}%
\pgfsetdash{}{0pt}%
\pgfpathmoveto{\pgfqpoint{1.481293in}{0.994893in}}%
\pgfpathlineto{\pgfqpoint{1.496168in}{1.006315in}}%
\pgfpathlineto{\pgfqpoint{1.496619in}{1.028646in}}%
\pgfpathlineto{\pgfqpoint{1.481727in}{1.017276in}}%
\pgfpathlineto{\pgfqpoint{1.481293in}{0.994893in}}%
\pgfpathclose%
\pgfusepath{fill}%
\end{pgfscope}%
\begin{pgfscope}%
\pgfpathrectangle{\pgfqpoint{0.407592in}{0.222904in}}{\pgfqpoint{1.411163in}{1.411163in}}%
\pgfusepath{clip}%
\pgfsetbuttcap%
\pgfsetroundjoin%
\definecolor{currentfill}{rgb}{0.168627,0.611765,0.411765}%
\pgfsetfillcolor{currentfill}%
\pgfsetfillopacity{0.180000}%
\pgfsetlinewidth{0.000000pt}%
\definecolor{currentstroke}{rgb}{0.000000,0.000000,0.000000}%
\pgfsetstrokecolor{currentstroke}%
\pgfsetdash{}{0pt}%
\pgfpathmoveto{\pgfqpoint{1.529173in}{1.190512in}}%
\pgfpathlineto{\pgfqpoint{1.544011in}{1.206067in}}%
\pgfpathlineto{\pgfqpoint{1.544521in}{1.228650in}}%
\pgfpathlineto{\pgfqpoint{1.529666in}{1.213137in}}%
\pgfpathlineto{\pgfqpoint{1.529173in}{1.190512in}}%
\pgfpathclose%
\pgfusepath{fill}%
\end{pgfscope}%
\begin{pgfscope}%
\pgfpathrectangle{\pgfqpoint{0.407592in}{0.222904in}}{\pgfqpoint{1.411163in}{1.411163in}}%
\pgfusepath{clip}%
\pgfsetbuttcap%
\pgfsetroundjoin%
\definecolor{currentfill}{rgb}{0.168627,0.611765,0.411765}%
\pgfsetfillcolor{currentfill}%
\pgfsetfillopacity{0.180000}%
\pgfsetlinewidth{0.000000pt}%
\definecolor{currentstroke}{rgb}{0.000000,0.000000,0.000000}%
\pgfsetstrokecolor{currentstroke}%
\pgfsetdash{}{0pt}%
\pgfpathmoveto{\pgfqpoint{1.513338in}{1.131226in}}%
\pgfpathlineto{\pgfqpoint{1.528190in}{1.145429in}}%
\pgfpathlineto{\pgfqpoint{1.528681in}{1.167943in}}%
\pgfpathlineto{\pgfqpoint{1.513811in}{1.153785in}}%
\pgfpathlineto{\pgfqpoint{1.513338in}{1.131226in}}%
\pgfpathclose%
\pgfusepath{fill}%
\end{pgfscope}%
\begin{pgfscope}%
\pgfpathrectangle{\pgfqpoint{0.407592in}{0.222904in}}{\pgfqpoint{1.411163in}{1.411163in}}%
\pgfusepath{clip}%
\pgfsetbuttcap%
\pgfsetroundjoin%
\definecolor{currentfill}{rgb}{0.168627,0.611765,0.411765}%
\pgfsetfillcolor{currentfill}%
\pgfsetfillopacity{0.180000}%
\pgfsetlinewidth{0.000000pt}%
\definecolor{currentstroke}{rgb}{0.000000,0.000000,0.000000}%
\pgfsetstrokecolor{currentstroke}%
\pgfsetdash{}{0pt}%
\pgfpathmoveto{\pgfqpoint{1.465530in}{0.940313in}}%
\pgfpathlineto{\pgfqpoint{1.480429in}{0.950293in}}%
\pgfpathlineto{\pgfqpoint{1.480861in}{0.972565in}}%
\pgfpathlineto{\pgfqpoint{1.465944in}{0.962642in}}%
\pgfpathlineto{\pgfqpoint{1.465530in}{0.940313in}}%
\pgfpathclose%
\pgfusepath{fill}%
\end{pgfscope}%
\begin{pgfscope}%
\pgfpathrectangle{\pgfqpoint{0.407592in}{0.222904in}}{\pgfqpoint{1.411163in}{1.411163in}}%
\pgfusepath{clip}%
\pgfsetbuttcap%
\pgfsetroundjoin%
\definecolor{currentfill}{rgb}{0.168627,0.611765,0.411765}%
\pgfsetfillcolor{currentfill}%
\pgfsetfillopacity{0.180000}%
\pgfsetlinewidth{0.000000pt}%
\definecolor{currentstroke}{rgb}{0.000000,0.000000,0.000000}%
\pgfsetstrokecolor{currentstroke}%
\pgfsetdash{}{0pt}%
\pgfpathmoveto{\pgfqpoint{1.497523in}{1.073473in}}%
\pgfpathlineto{\pgfqpoint{1.512394in}{1.086275in}}%
\pgfpathlineto{\pgfqpoint{1.512865in}{1.108723in}}%
\pgfpathlineto{\pgfqpoint{1.497978in}{1.095970in}}%
\pgfpathlineto{\pgfqpoint{1.497523in}{1.073473in}}%
\pgfpathclose%
\pgfusepath{fill}%
\end{pgfscope}%
\begin{pgfscope}%
\pgfpathrectangle{\pgfqpoint{0.407592in}{0.222904in}}{\pgfqpoint{1.411163in}{1.411163in}}%
\pgfusepath{clip}%
\pgfsetbuttcap%
\pgfsetroundjoin%
\definecolor{currentfill}{rgb}{0.168627,0.611765,0.411765}%
\pgfsetfillcolor{currentfill}%
\pgfsetfillopacity{0.180000}%
\pgfsetlinewidth{0.000000pt}%
\definecolor{currentstroke}{rgb}{0.000000,0.000000,0.000000}%
\pgfsetstrokecolor{currentstroke}%
\pgfsetdash{}{0pt}%
\pgfpathmoveto{\pgfqpoint{1.449776in}{0.887281in}}%
\pgfpathlineto{\pgfqpoint{1.464705in}{0.895822in}}%
\pgfpathlineto{\pgfqpoint{1.465117in}{0.918040in}}%
\pgfpathlineto{\pgfqpoint{1.450171in}{0.909559in}}%
\pgfpathlineto{\pgfqpoint{1.449776in}{0.887281in}}%
\pgfpathclose%
\pgfusepath{fill}%
\end{pgfscope}%
\begin{pgfscope}%
\pgfpathrectangle{\pgfqpoint{0.407592in}{0.222904in}}{\pgfqpoint{1.411163in}{1.411163in}}%
\pgfusepath{clip}%
\pgfsetbuttcap%
\pgfsetroundjoin%
\definecolor{currentfill}{rgb}{0.168627,0.611765,0.411765}%
\pgfsetfillcolor{currentfill}%
\pgfsetfillopacity{0.180000}%
\pgfsetlinewidth{0.000000pt}%
\definecolor{currentstroke}{rgb}{0.000000,0.000000,0.000000}%
\pgfsetstrokecolor{currentstroke}%
\pgfsetdash{}{0pt}%
\pgfpathmoveto{\pgfqpoint{1.481727in}{1.017276in}}%
\pgfpathlineto{\pgfqpoint{1.496619in}{1.028646in}}%
\pgfpathlineto{\pgfqpoint{1.497070in}{1.051031in}}%
\pgfpathlineto{\pgfqpoint{1.482162in}{1.039714in}}%
\pgfpathlineto{\pgfqpoint{1.481727in}{1.017276in}}%
\pgfpathclose%
\pgfusepath{fill}%
\end{pgfscope}%
\begin{pgfscope}%
\pgfpathrectangle{\pgfqpoint{0.407592in}{0.222904in}}{\pgfqpoint{1.411163in}{1.411163in}}%
\pgfusepath{clip}%
\pgfsetbuttcap%
\pgfsetroundjoin%
\definecolor{currentfill}{rgb}{0.168627,0.611765,0.411765}%
\pgfsetfillcolor{currentfill}%
\pgfsetfillopacity{0.180000}%
\pgfsetlinewidth{0.000000pt}%
\definecolor{currentstroke}{rgb}{0.000000,0.000000,0.000000}%
\pgfsetstrokecolor{currentstroke}%
\pgfsetdash{}{0pt}%
\pgfpathmoveto{\pgfqpoint{1.529666in}{1.213137in}}%
\pgfpathlineto{\pgfqpoint{1.544521in}{1.228650in}}%
\pgfpathlineto{\pgfqpoint{1.545032in}{1.251289in}}%
\pgfpathlineto{\pgfqpoint{1.530160in}{1.235819in}}%
\pgfpathlineto{\pgfqpoint{1.529666in}{1.213137in}}%
\pgfpathclose%
\pgfusepath{fill}%
\end{pgfscope}%
\begin{pgfscope}%
\pgfpathrectangle{\pgfqpoint{0.407592in}{0.222904in}}{\pgfqpoint{1.411163in}{1.411163in}}%
\pgfusepath{clip}%
\pgfsetbuttcap%
\pgfsetroundjoin%
\definecolor{currentfill}{rgb}{0.168627,0.611765,0.411765}%
\pgfsetfillcolor{currentfill}%
\pgfsetfillopacity{0.180000}%
\pgfsetlinewidth{0.000000pt}%
\definecolor{currentstroke}{rgb}{0.000000,0.000000,0.000000}%
\pgfsetstrokecolor{currentstroke}%
\pgfsetdash{}{0pt}%
\pgfpathmoveto{\pgfqpoint{1.513811in}{1.153785in}}%
\pgfpathlineto{\pgfqpoint{1.528681in}{1.167943in}}%
\pgfpathlineto{\pgfqpoint{1.529173in}{1.190512in}}%
\pgfpathlineto{\pgfqpoint{1.514286in}{1.176401in}}%
\pgfpathlineto{\pgfqpoint{1.513811in}{1.153785in}}%
\pgfpathclose%
\pgfusepath{fill}%
\end{pgfscope}%
\begin{pgfscope}%
\pgfpathrectangle{\pgfqpoint{0.407592in}{0.222904in}}{\pgfqpoint{1.411163in}{1.411163in}}%
\pgfusepath{clip}%
\pgfsetbuttcap%
\pgfsetroundjoin%
\definecolor{currentfill}{rgb}{0.168627,0.611765,0.411765}%
\pgfsetfillcolor{currentfill}%
\pgfsetfillopacity{0.180000}%
\pgfsetlinewidth{0.000000pt}%
\definecolor{currentstroke}{rgb}{0.000000,0.000000,0.000000}%
\pgfsetstrokecolor{currentstroke}%
\pgfsetdash{}{0pt}%
\pgfpathmoveto{\pgfqpoint{1.465944in}{0.962642in}}%
\pgfpathlineto{\pgfqpoint{1.480861in}{0.972565in}}%
\pgfpathlineto{\pgfqpoint{1.481293in}{0.994893in}}%
\pgfpathlineto{\pgfqpoint{1.466360in}{0.985026in}}%
\pgfpathlineto{\pgfqpoint{1.465944in}{0.962642in}}%
\pgfpathclose%
\pgfusepath{fill}%
\end{pgfscope}%
\begin{pgfscope}%
\pgfpathrectangle{\pgfqpoint{0.407592in}{0.222904in}}{\pgfqpoint{1.411163in}{1.411163in}}%
\pgfusepath{clip}%
\pgfsetbuttcap%
\pgfsetroundjoin%
\definecolor{currentfill}{rgb}{0.168627,0.611765,0.411765}%
\pgfsetfillcolor{currentfill}%
\pgfsetfillopacity{0.180000}%
\pgfsetlinewidth{0.000000pt}%
\definecolor{currentstroke}{rgb}{0.000000,0.000000,0.000000}%
\pgfsetstrokecolor{currentstroke}%
\pgfsetdash{}{0pt}%
\pgfpathmoveto{\pgfqpoint{1.497978in}{1.095970in}}%
\pgfpathlineto{\pgfqpoint{1.512865in}{1.108723in}}%
\pgfpathlineto{\pgfqpoint{1.513338in}{1.131226in}}%
\pgfpathlineto{\pgfqpoint{1.498433in}{1.118523in}}%
\pgfpathlineto{\pgfqpoint{1.497978in}{1.095970in}}%
\pgfpathclose%
\pgfusepath{fill}%
\end{pgfscope}%
\begin{pgfscope}%
\pgfpathrectangle{\pgfqpoint{0.407592in}{0.222904in}}{\pgfqpoint{1.411163in}{1.411163in}}%
\pgfusepath{clip}%
\pgfsetbuttcap%
\pgfsetroundjoin%
\definecolor{currentfill}{rgb}{0.168627,0.611765,0.411765}%
\pgfsetfillcolor{currentfill}%
\pgfsetfillopacity{0.180000}%
\pgfsetlinewidth{0.000000pt}%
\definecolor{currentstroke}{rgb}{0.000000,0.000000,0.000000}%
\pgfsetstrokecolor{currentstroke}%
\pgfsetdash{}{0pt}%
\pgfpathmoveto{\pgfqpoint{1.450171in}{0.909559in}}%
\pgfpathlineto{\pgfqpoint{1.465117in}{0.918040in}}%
\pgfpathlineto{\pgfqpoint{1.465530in}{0.940313in}}%
\pgfpathlineto{\pgfqpoint{1.450567in}{0.931893in}}%
\pgfpathlineto{\pgfqpoint{1.450171in}{0.909559in}}%
\pgfpathclose%
\pgfusepath{fill}%
\end{pgfscope}%
\begin{pgfscope}%
\pgfpathrectangle{\pgfqpoint{0.407592in}{0.222904in}}{\pgfqpoint{1.411163in}{1.411163in}}%
\pgfusepath{clip}%
\pgfsetbuttcap%
\pgfsetroundjoin%
\definecolor{currentfill}{rgb}{0.168627,0.611765,0.411765}%
\pgfsetfillcolor{currentfill}%
\pgfsetfillopacity{0.180000}%
\pgfsetlinewidth{0.000000pt}%
\definecolor{currentstroke}{rgb}{0.000000,0.000000,0.000000}%
\pgfsetstrokecolor{currentstroke}%
\pgfsetdash{}{0pt}%
\pgfpathmoveto{\pgfqpoint{1.482162in}{1.039714in}}%
\pgfpathlineto{\pgfqpoint{1.497070in}{1.051031in}}%
\pgfpathlineto{\pgfqpoint{1.497523in}{1.073473in}}%
\pgfpathlineto{\pgfqpoint{1.482597in}{1.062209in}}%
\pgfpathlineto{\pgfqpoint{1.482162in}{1.039714in}}%
\pgfpathclose%
\pgfusepath{fill}%
\end{pgfscope}%
\begin{pgfscope}%
\pgfpathrectangle{\pgfqpoint{0.407592in}{0.222904in}}{\pgfqpoint{1.411163in}{1.411163in}}%
\pgfusepath{clip}%
\pgfsetbuttcap%
\pgfsetroundjoin%
\definecolor{currentfill}{rgb}{0.168627,0.611765,0.411765}%
\pgfsetfillcolor{currentfill}%
\pgfsetfillopacity{0.180000}%
\pgfsetlinewidth{0.000000pt}%
\definecolor{currentstroke}{rgb}{0.000000,0.000000,0.000000}%
\pgfsetstrokecolor{currentstroke}%
\pgfsetdash{}{0pt}%
\pgfpathmoveto{\pgfqpoint{1.514286in}{1.176401in}}%
\pgfpathlineto{\pgfqpoint{1.529173in}{1.190512in}}%
\pgfpathlineto{\pgfqpoint{1.529666in}{1.213137in}}%
\pgfpathlineto{\pgfqpoint{1.514762in}{1.199073in}}%
\pgfpathlineto{\pgfqpoint{1.514286in}{1.176401in}}%
\pgfpathclose%
\pgfusepath{fill}%
\end{pgfscope}%
\begin{pgfscope}%
\pgfpathrectangle{\pgfqpoint{0.407592in}{0.222904in}}{\pgfqpoint{1.411163in}{1.411163in}}%
\pgfusepath{clip}%
\pgfsetbuttcap%
\pgfsetroundjoin%
\definecolor{currentfill}{rgb}{0.168627,0.611765,0.411765}%
\pgfsetfillcolor{currentfill}%
\pgfsetfillopacity{0.180000}%
\pgfsetlinewidth{0.000000pt}%
\definecolor{currentstroke}{rgb}{0.000000,0.000000,0.000000}%
\pgfsetstrokecolor{currentstroke}%
\pgfsetdash{}{0pt}%
\pgfpathmoveto{\pgfqpoint{1.466360in}{0.985026in}}%
\pgfpathlineto{\pgfqpoint{1.481293in}{0.994893in}}%
\pgfpathlineto{\pgfqpoint{1.481727in}{1.017276in}}%
\pgfpathlineto{\pgfqpoint{1.466776in}{1.007466in}}%
\pgfpathlineto{\pgfqpoint{1.466360in}{0.985026in}}%
\pgfpathclose%
\pgfusepath{fill}%
\end{pgfscope}%
\begin{pgfscope}%
\pgfpathrectangle{\pgfqpoint{0.407592in}{0.222904in}}{\pgfqpoint{1.411163in}{1.411163in}}%
\pgfusepath{clip}%
\pgfsetbuttcap%
\pgfsetroundjoin%
\definecolor{currentfill}{rgb}{0.168627,0.611765,0.411765}%
\pgfsetfillcolor{currentfill}%
\pgfsetfillopacity{0.180000}%
\pgfsetlinewidth{0.000000pt}%
\definecolor{currentstroke}{rgb}{0.000000,0.000000,0.000000}%
\pgfsetstrokecolor{currentstroke}%
\pgfsetdash{}{0pt}%
\pgfpathmoveto{\pgfqpoint{1.498433in}{1.118523in}}%
\pgfpathlineto{\pgfqpoint{1.513338in}{1.131226in}}%
\pgfpathlineto{\pgfqpoint{1.513811in}{1.153785in}}%
\pgfpathlineto{\pgfqpoint{1.498889in}{1.141132in}}%
\pgfpathlineto{\pgfqpoint{1.498433in}{1.118523in}}%
\pgfpathclose%
\pgfusepath{fill}%
\end{pgfscope}%
\begin{pgfscope}%
\pgfpathrectangle{\pgfqpoint{0.407592in}{0.222904in}}{\pgfqpoint{1.411163in}{1.411163in}}%
\pgfusepath{clip}%
\pgfsetbuttcap%
\pgfsetroundjoin%
\definecolor{currentfill}{rgb}{0.168627,0.611765,0.411765}%
\pgfsetfillcolor{currentfill}%
\pgfsetfillopacity{0.180000}%
\pgfsetlinewidth{0.000000pt}%
\definecolor{currentstroke}{rgb}{0.000000,0.000000,0.000000}%
\pgfsetstrokecolor{currentstroke}%
\pgfsetdash{}{0pt}%
\pgfpathmoveto{\pgfqpoint{1.450567in}{0.931893in}}%
\pgfpathlineto{\pgfqpoint{1.465530in}{0.940313in}}%
\pgfpathlineto{\pgfqpoint{1.465944in}{0.962642in}}%
\pgfpathlineto{\pgfqpoint{1.450964in}{0.954282in}}%
\pgfpathlineto{\pgfqpoint{1.450567in}{0.931893in}}%
\pgfpathclose%
\pgfusepath{fill}%
\end{pgfscope}%
\begin{pgfscope}%
\pgfpathrectangle{\pgfqpoint{0.407592in}{0.222904in}}{\pgfqpoint{1.411163in}{1.411163in}}%
\pgfusepath{clip}%
\pgfsetbuttcap%
\pgfsetroundjoin%
\definecolor{currentfill}{rgb}{0.168627,0.611765,0.411765}%
\pgfsetfillcolor{currentfill}%
\pgfsetfillopacity{0.180000}%
\pgfsetlinewidth{0.000000pt}%
\definecolor{currentstroke}{rgb}{0.000000,0.000000,0.000000}%
\pgfsetstrokecolor{currentstroke}%
\pgfsetdash{}{0pt}%
\pgfpathmoveto{\pgfqpoint{1.482597in}{1.062209in}}%
\pgfpathlineto{\pgfqpoint{1.497523in}{1.073473in}}%
\pgfpathlineto{\pgfqpoint{1.497978in}{1.095970in}}%
\pgfpathlineto{\pgfqpoint{1.483034in}{1.084760in}}%
\pgfpathlineto{\pgfqpoint{1.482597in}{1.062209in}}%
\pgfpathclose%
\pgfusepath{fill}%
\end{pgfscope}%
\begin{pgfscope}%
\pgfpathrectangle{\pgfqpoint{0.407592in}{0.222904in}}{\pgfqpoint{1.411163in}{1.411163in}}%
\pgfusepath{clip}%
\pgfsetbuttcap%
\pgfsetroundjoin%
\definecolor{currentfill}{rgb}{0.168627,0.611765,0.411765}%
\pgfsetfillcolor{currentfill}%
\pgfsetfillopacity{0.180000}%
\pgfsetlinewidth{0.000000pt}%
\definecolor{currentstroke}{rgb}{0.000000,0.000000,0.000000}%
\pgfsetstrokecolor{currentstroke}%
\pgfsetdash{}{0pt}%
\pgfpathmoveto{\pgfqpoint{1.514762in}{1.199073in}}%
\pgfpathlineto{\pgfqpoint{1.529666in}{1.213137in}}%
\pgfpathlineto{\pgfqpoint{1.530160in}{1.235819in}}%
\pgfpathlineto{\pgfqpoint{1.515239in}{1.221801in}}%
\pgfpathlineto{\pgfqpoint{1.514762in}{1.199073in}}%
\pgfpathclose%
\pgfusepath{fill}%
\end{pgfscope}%
\begin{pgfscope}%
\pgfpathrectangle{\pgfqpoint{0.407592in}{0.222904in}}{\pgfqpoint{1.411163in}{1.411163in}}%
\pgfusepath{clip}%
\pgfsetbuttcap%
\pgfsetroundjoin%
\definecolor{currentfill}{rgb}{0.168627,0.611765,0.411765}%
\pgfsetfillcolor{currentfill}%
\pgfsetfillopacity{0.180000}%
\pgfsetlinewidth{0.000000pt}%
\definecolor{currentstroke}{rgb}{0.000000,0.000000,0.000000}%
\pgfsetstrokecolor{currentstroke}%
\pgfsetdash{}{0pt}%
\pgfpathmoveto{\pgfqpoint{1.466776in}{1.007466in}}%
\pgfpathlineto{\pgfqpoint{1.481727in}{1.017276in}}%
\pgfpathlineto{\pgfqpoint{1.482162in}{1.039714in}}%
\pgfpathlineto{\pgfqpoint{1.467193in}{1.029962in}}%
\pgfpathlineto{\pgfqpoint{1.466776in}{1.007466in}}%
\pgfpathclose%
\pgfusepath{fill}%
\end{pgfscope}%
\begin{pgfscope}%
\pgfpathrectangle{\pgfqpoint{0.407592in}{0.222904in}}{\pgfqpoint{1.411163in}{1.411163in}}%
\pgfusepath{clip}%
\pgfsetbuttcap%
\pgfsetroundjoin%
\definecolor{currentfill}{rgb}{0.168627,0.611765,0.411765}%
\pgfsetfillcolor{currentfill}%
\pgfsetfillopacity{0.180000}%
\pgfsetlinewidth{0.000000pt}%
\definecolor{currentstroke}{rgb}{0.000000,0.000000,0.000000}%
\pgfsetstrokecolor{currentstroke}%
\pgfsetdash{}{0pt}%
\pgfpathmoveto{\pgfqpoint{1.434780in}{0.880284in}}%
\pgfpathlineto{\pgfqpoint{1.449776in}{0.887281in}}%
\pgfpathlineto{\pgfqpoint{1.450171in}{0.909559in}}%
\pgfpathlineto{\pgfqpoint{1.435158in}{0.902627in}}%
\pgfpathlineto{\pgfqpoint{1.434780in}{0.880284in}}%
\pgfpathclose%
\pgfusepath{fill}%
\end{pgfscope}%
\begin{pgfscope}%
\pgfpathrectangle{\pgfqpoint{0.407592in}{0.222904in}}{\pgfqpoint{1.411163in}{1.411163in}}%
\pgfusepath{clip}%
\pgfsetbuttcap%
\pgfsetroundjoin%
\definecolor{currentfill}{rgb}{0.168627,0.611765,0.411765}%
\pgfsetfillcolor{currentfill}%
\pgfsetfillopacity{0.180000}%
\pgfsetlinewidth{0.000000pt}%
\definecolor{currentstroke}{rgb}{0.000000,0.000000,0.000000}%
\pgfsetstrokecolor{currentstroke}%
\pgfsetdash{}{0pt}%
\pgfpathmoveto{\pgfqpoint{1.498889in}{1.141132in}}%
\pgfpathlineto{\pgfqpoint{1.513811in}{1.153785in}}%
\pgfpathlineto{\pgfqpoint{1.514286in}{1.176401in}}%
\pgfpathlineto{\pgfqpoint{1.499347in}{1.163798in}}%
\pgfpathlineto{\pgfqpoint{1.498889in}{1.141132in}}%
\pgfpathclose%
\pgfusepath{fill}%
\end{pgfscope}%
\begin{pgfscope}%
\pgfpathrectangle{\pgfqpoint{0.407592in}{0.222904in}}{\pgfqpoint{1.411163in}{1.411163in}}%
\pgfusepath{clip}%
\pgfsetbuttcap%
\pgfsetroundjoin%
\definecolor{currentfill}{rgb}{0.168627,0.611765,0.411765}%
\pgfsetfillcolor{currentfill}%
\pgfsetfillopacity{0.180000}%
\pgfsetlinewidth{0.000000pt}%
\definecolor{currentstroke}{rgb}{0.000000,0.000000,0.000000}%
\pgfsetstrokecolor{currentstroke}%
\pgfsetdash{}{0pt}%
\pgfpathmoveto{\pgfqpoint{1.450964in}{0.954282in}}%
\pgfpathlineto{\pgfqpoint{1.465944in}{0.962642in}}%
\pgfpathlineto{\pgfqpoint{1.466360in}{0.985026in}}%
\pgfpathlineto{\pgfqpoint{1.451362in}{0.976727in}}%
\pgfpathlineto{\pgfqpoint{1.450964in}{0.954282in}}%
\pgfpathclose%
\pgfusepath{fill}%
\end{pgfscope}%
\begin{pgfscope}%
\pgfpathrectangle{\pgfqpoint{0.407592in}{0.222904in}}{\pgfqpoint{1.411163in}{1.411163in}}%
\pgfusepath{clip}%
\pgfsetbuttcap%
\pgfsetroundjoin%
\definecolor{currentfill}{rgb}{0.168627,0.611765,0.411765}%
\pgfsetfillcolor{currentfill}%
\pgfsetfillopacity{0.180000}%
\pgfsetlinewidth{0.000000pt}%
\definecolor{currentstroke}{rgb}{0.000000,0.000000,0.000000}%
\pgfsetstrokecolor{currentstroke}%
\pgfsetdash{}{0pt}%
\pgfpathmoveto{\pgfqpoint{1.483034in}{1.084760in}}%
\pgfpathlineto{\pgfqpoint{1.497978in}{1.095970in}}%
\pgfpathlineto{\pgfqpoint{1.498433in}{1.118523in}}%
\pgfpathlineto{\pgfqpoint{1.483472in}{1.107367in}}%
\pgfpathlineto{\pgfqpoint{1.483034in}{1.084760in}}%
\pgfpathclose%
\pgfusepath{fill}%
\end{pgfscope}%
\begin{pgfscope}%
\pgfpathrectangle{\pgfqpoint{0.407592in}{0.222904in}}{\pgfqpoint{1.411163in}{1.411163in}}%
\pgfusepath{clip}%
\pgfsetbuttcap%
\pgfsetroundjoin%
\definecolor{currentfill}{rgb}{0.168627,0.611765,0.411765}%
\pgfsetfillcolor{currentfill}%
\pgfsetfillopacity{0.180000}%
\pgfsetlinewidth{0.000000pt}%
\definecolor{currentstroke}{rgb}{0.000000,0.000000,0.000000}%
\pgfsetstrokecolor{currentstroke}%
\pgfsetdash{}{0pt}%
\pgfpathmoveto{\pgfqpoint{1.467193in}{1.029962in}}%
\pgfpathlineto{\pgfqpoint{1.482162in}{1.039714in}}%
\pgfpathlineto{\pgfqpoint{1.482597in}{1.062209in}}%
\pgfpathlineto{\pgfqpoint{1.467612in}{1.052514in}}%
\pgfpathlineto{\pgfqpoint{1.467193in}{1.029962in}}%
\pgfpathclose%
\pgfusepath{fill}%
\end{pgfscope}%
\begin{pgfscope}%
\pgfpathrectangle{\pgfqpoint{0.407592in}{0.222904in}}{\pgfqpoint{1.411163in}{1.411163in}}%
\pgfusepath{clip}%
\pgfsetbuttcap%
\pgfsetroundjoin%
\definecolor{currentfill}{rgb}{0.168627,0.611765,0.411765}%
\pgfsetfillcolor{currentfill}%
\pgfsetfillopacity{0.180000}%
\pgfsetlinewidth{0.000000pt}%
\definecolor{currentstroke}{rgb}{0.000000,0.000000,0.000000}%
\pgfsetstrokecolor{currentstroke}%
\pgfsetdash{}{0pt}%
\pgfpathmoveto{\pgfqpoint{1.435158in}{0.902627in}}%
\pgfpathlineto{\pgfqpoint{1.450171in}{0.909559in}}%
\pgfpathlineto{\pgfqpoint{1.450567in}{0.931893in}}%
\pgfpathlineto{\pgfqpoint{1.435536in}{0.925025in}}%
\pgfpathlineto{\pgfqpoint{1.435158in}{0.902627in}}%
\pgfpathclose%
\pgfusepath{fill}%
\end{pgfscope}%
\begin{pgfscope}%
\pgfpathrectangle{\pgfqpoint{0.407592in}{0.222904in}}{\pgfqpoint{1.411163in}{1.411163in}}%
\pgfusepath{clip}%
\pgfsetbuttcap%
\pgfsetroundjoin%
\definecolor{currentfill}{rgb}{0.168627,0.611765,0.411765}%
\pgfsetfillcolor{currentfill}%
\pgfsetfillopacity{0.180000}%
\pgfsetlinewidth{0.000000pt}%
\definecolor{currentstroke}{rgb}{0.000000,0.000000,0.000000}%
\pgfsetstrokecolor{currentstroke}%
\pgfsetdash{}{0pt}%
\pgfpathmoveto{\pgfqpoint{1.499347in}{1.163798in}}%
\pgfpathlineto{\pgfqpoint{1.514286in}{1.176401in}}%
\pgfpathlineto{\pgfqpoint{1.514762in}{1.199073in}}%
\pgfpathlineto{\pgfqpoint{1.499805in}{1.186520in}}%
\pgfpathlineto{\pgfqpoint{1.499347in}{1.163798in}}%
\pgfpathclose%
\pgfusepath{fill}%
\end{pgfscope}%
\begin{pgfscope}%
\pgfpathrectangle{\pgfqpoint{0.407592in}{0.222904in}}{\pgfqpoint{1.411163in}{1.411163in}}%
\pgfusepath{clip}%
\pgfsetbuttcap%
\pgfsetroundjoin%
\definecolor{currentfill}{rgb}{0.168627,0.611765,0.411765}%
\pgfsetfillcolor{currentfill}%
\pgfsetfillopacity{0.180000}%
\pgfsetlinewidth{0.000000pt}%
\definecolor{currentstroke}{rgb}{0.000000,0.000000,0.000000}%
\pgfsetstrokecolor{currentstroke}%
\pgfsetdash{}{0pt}%
\pgfpathmoveto{\pgfqpoint{1.451362in}{0.976727in}}%
\pgfpathlineto{\pgfqpoint{1.466360in}{0.985026in}}%
\pgfpathlineto{\pgfqpoint{1.466776in}{1.007466in}}%
\pgfpathlineto{\pgfqpoint{1.451761in}{0.999229in}}%
\pgfpathlineto{\pgfqpoint{1.451362in}{0.976727in}}%
\pgfpathclose%
\pgfusepath{fill}%
\end{pgfscope}%
\begin{pgfscope}%
\pgfpathrectangle{\pgfqpoint{0.407592in}{0.222904in}}{\pgfqpoint{1.411163in}{1.411163in}}%
\pgfusepath{clip}%
\pgfsetbuttcap%
\pgfsetroundjoin%
\definecolor{currentfill}{rgb}{0.168627,0.611765,0.411765}%
\pgfsetfillcolor{currentfill}%
\pgfsetfillopacity{0.180000}%
\pgfsetlinewidth{0.000000pt}%
\definecolor{currentstroke}{rgb}{0.000000,0.000000,0.000000}%
\pgfsetstrokecolor{currentstroke}%
\pgfsetdash{}{0pt}%
\pgfpathmoveto{\pgfqpoint{1.483472in}{1.107367in}}%
\pgfpathlineto{\pgfqpoint{1.498433in}{1.118523in}}%
\pgfpathlineto{\pgfqpoint{1.498889in}{1.141132in}}%
\pgfpathlineto{\pgfqpoint{1.483911in}{1.130030in}}%
\pgfpathlineto{\pgfqpoint{1.483472in}{1.107367in}}%
\pgfpathclose%
\pgfusepath{fill}%
\end{pgfscope}%
\begin{pgfscope}%
\pgfpathrectangle{\pgfqpoint{0.407592in}{0.222904in}}{\pgfqpoint{1.411163in}{1.411163in}}%
\pgfusepath{clip}%
\pgfsetbuttcap%
\pgfsetroundjoin%
\definecolor{currentfill}{rgb}{0.168627,0.611765,0.411765}%
\pgfsetfillcolor{currentfill}%
\pgfsetfillopacity{0.180000}%
\pgfsetlinewidth{0.000000pt}%
\definecolor{currentstroke}{rgb}{0.000000,0.000000,0.000000}%
\pgfsetstrokecolor{currentstroke}%
\pgfsetdash{}{0pt}%
\pgfpathmoveto{\pgfqpoint{1.467612in}{1.052514in}}%
\pgfpathlineto{\pgfqpoint{1.482597in}{1.062209in}}%
\pgfpathlineto{\pgfqpoint{1.483034in}{1.084760in}}%
\pgfpathlineto{\pgfqpoint{1.468031in}{1.075123in}}%
\pgfpathlineto{\pgfqpoint{1.467612in}{1.052514in}}%
\pgfpathclose%
\pgfusepath{fill}%
\end{pgfscope}%
\begin{pgfscope}%
\pgfpathrectangle{\pgfqpoint{0.407592in}{0.222904in}}{\pgfqpoint{1.411163in}{1.411163in}}%
\pgfusepath{clip}%
\pgfsetbuttcap%
\pgfsetroundjoin%
\definecolor{currentfill}{rgb}{0.168627,0.611765,0.411765}%
\pgfsetfillcolor{currentfill}%
\pgfsetfillopacity{0.180000}%
\pgfsetlinewidth{0.000000pt}%
\definecolor{currentstroke}{rgb}{0.000000,0.000000,0.000000}%
\pgfsetstrokecolor{currentstroke}%
\pgfsetdash{}{0pt}%
\pgfpathmoveto{\pgfqpoint{1.435536in}{0.925025in}}%
\pgfpathlineto{\pgfqpoint{1.450567in}{0.931893in}}%
\pgfpathlineto{\pgfqpoint{1.450964in}{0.954282in}}%
\pgfpathlineto{\pgfqpoint{1.435916in}{0.947480in}}%
\pgfpathlineto{\pgfqpoint{1.435536in}{0.925025in}}%
\pgfpathclose%
\pgfusepath{fill}%
\end{pgfscope}%
\begin{pgfscope}%
\pgfpathrectangle{\pgfqpoint{0.407592in}{0.222904in}}{\pgfqpoint{1.411163in}{1.411163in}}%
\pgfusepath{clip}%
\pgfsetbuttcap%
\pgfsetroundjoin%
\definecolor{currentfill}{rgb}{0.168627,0.611765,0.411765}%
\pgfsetfillcolor{currentfill}%
\pgfsetfillopacity{0.180000}%
\pgfsetlinewidth{0.000000pt}%
\definecolor{currentstroke}{rgb}{0.000000,0.000000,0.000000}%
\pgfsetstrokecolor{currentstroke}%
\pgfsetdash{}{0pt}%
\pgfpathmoveto{\pgfqpoint{1.499805in}{1.186520in}}%
\pgfpathlineto{\pgfqpoint{1.514762in}{1.199073in}}%
\pgfpathlineto{\pgfqpoint{1.515239in}{1.221801in}}%
\pgfpathlineto{\pgfqpoint{1.500265in}{1.209299in}}%
\pgfpathlineto{\pgfqpoint{1.499805in}{1.186520in}}%
\pgfpathclose%
\pgfusepath{fill}%
\end{pgfscope}%
\begin{pgfscope}%
\pgfpathrectangle{\pgfqpoint{0.407592in}{0.222904in}}{\pgfqpoint{1.411163in}{1.411163in}}%
\pgfusepath{clip}%
\pgfsetbuttcap%
\pgfsetroundjoin%
\definecolor{currentfill}{rgb}{0.168627,0.611765,0.411765}%
\pgfsetfillcolor{currentfill}%
\pgfsetfillopacity{0.180000}%
\pgfsetlinewidth{0.000000pt}%
\definecolor{currentstroke}{rgb}{0.000000,0.000000,0.000000}%
\pgfsetstrokecolor{currentstroke}%
\pgfsetdash{}{0pt}%
\pgfpathmoveto{\pgfqpoint{1.451761in}{0.999229in}}%
\pgfpathlineto{\pgfqpoint{1.466776in}{1.007466in}}%
\pgfpathlineto{\pgfqpoint{1.467193in}{1.029962in}}%
\pgfpathlineto{\pgfqpoint{1.452161in}{1.021787in}}%
\pgfpathlineto{\pgfqpoint{1.451761in}{0.999229in}}%
\pgfpathclose%
\pgfusepath{fill}%
\end{pgfscope}%
\begin{pgfscope}%
\pgfpathrectangle{\pgfqpoint{0.407592in}{0.222904in}}{\pgfqpoint{1.411163in}{1.411163in}}%
\pgfusepath{clip}%
\pgfsetbuttcap%
\pgfsetroundjoin%
\definecolor{currentfill}{rgb}{0.168627,0.611765,0.411765}%
\pgfsetfillcolor{currentfill}%
\pgfsetfillopacity{0.180000}%
\pgfsetlinewidth{0.000000pt}%
\definecolor{currentstroke}{rgb}{0.000000,0.000000,0.000000}%
\pgfsetstrokecolor{currentstroke}%
\pgfsetdash{}{0pt}%
\pgfpathmoveto{\pgfqpoint{1.483911in}{1.130030in}}%
\pgfpathlineto{\pgfqpoint{1.498889in}{1.141132in}}%
\pgfpathlineto{\pgfqpoint{1.499347in}{1.163798in}}%
\pgfpathlineto{\pgfqpoint{1.484352in}{1.152750in}}%
\pgfpathlineto{\pgfqpoint{1.483911in}{1.130030in}}%
\pgfpathclose%
\pgfusepath{fill}%
\end{pgfscope}%
\begin{pgfscope}%
\pgfpathrectangle{\pgfqpoint{0.407592in}{0.222904in}}{\pgfqpoint{1.411163in}{1.411163in}}%
\pgfusepath{clip}%
\pgfsetbuttcap%
\pgfsetroundjoin%
\definecolor{currentfill}{rgb}{0.168627,0.611765,0.411765}%
\pgfsetfillcolor{currentfill}%
\pgfsetfillopacity{0.180000}%
\pgfsetlinewidth{0.000000pt}%
\definecolor{currentstroke}{rgb}{0.000000,0.000000,0.000000}%
\pgfsetstrokecolor{currentstroke}%
\pgfsetdash{}{0pt}%
\pgfpathmoveto{\pgfqpoint{1.419710in}{0.874808in}}%
\pgfpathlineto{\pgfqpoint{1.434780in}{0.880284in}}%
\pgfpathlineto{\pgfqpoint{1.435158in}{0.902627in}}%
\pgfpathlineto{\pgfqpoint{1.420070in}{0.897219in}}%
\pgfpathlineto{\pgfqpoint{1.419710in}{0.874808in}}%
\pgfpathclose%
\pgfusepath{fill}%
\end{pgfscope}%
\begin{pgfscope}%
\pgfpathrectangle{\pgfqpoint{0.407592in}{0.222904in}}{\pgfqpoint{1.411163in}{1.411163in}}%
\pgfusepath{clip}%
\pgfsetbuttcap%
\pgfsetroundjoin%
\definecolor{currentfill}{rgb}{0.168627,0.611765,0.411765}%
\pgfsetfillcolor{currentfill}%
\pgfsetfillopacity{0.180000}%
\pgfsetlinewidth{0.000000pt}%
\definecolor{currentstroke}{rgb}{0.000000,0.000000,0.000000}%
\pgfsetstrokecolor{currentstroke}%
\pgfsetdash{}{0pt}%
\pgfpathmoveto{\pgfqpoint{1.468031in}{1.075123in}}%
\pgfpathlineto{\pgfqpoint{1.483034in}{1.084760in}}%
\pgfpathlineto{\pgfqpoint{1.483472in}{1.107367in}}%
\pgfpathlineto{\pgfqpoint{1.468452in}{1.097788in}}%
\pgfpathlineto{\pgfqpoint{1.468031in}{1.075123in}}%
\pgfpathclose%
\pgfusepath{fill}%
\end{pgfscope}%
\begin{pgfscope}%
\pgfpathrectangle{\pgfqpoint{0.407592in}{0.222904in}}{\pgfqpoint{1.411163in}{1.411163in}}%
\pgfusepath{clip}%
\pgfsetbuttcap%
\pgfsetroundjoin%
\definecolor{currentfill}{rgb}{0.168627,0.611765,0.411765}%
\pgfsetfillcolor{currentfill}%
\pgfsetfillopacity{0.180000}%
\pgfsetlinewidth{0.000000pt}%
\definecolor{currentstroke}{rgb}{0.000000,0.000000,0.000000}%
\pgfsetstrokecolor{currentstroke}%
\pgfsetdash{}{0pt}%
\pgfpathmoveto{\pgfqpoint{1.435916in}{0.947480in}}%
\pgfpathlineto{\pgfqpoint{1.450964in}{0.954282in}}%
\pgfpathlineto{\pgfqpoint{1.451362in}{0.976727in}}%
\pgfpathlineto{\pgfqpoint{1.436296in}{0.969990in}}%
\pgfpathlineto{\pgfqpoint{1.435916in}{0.947480in}}%
\pgfpathclose%
\pgfusepath{fill}%
\end{pgfscope}%
\begin{pgfscope}%
\pgfpathrectangle{\pgfqpoint{0.407592in}{0.222904in}}{\pgfqpoint{1.411163in}{1.411163in}}%
\pgfusepath{clip}%
\pgfsetbuttcap%
\pgfsetroundjoin%
\definecolor{currentfill}{rgb}{0.168627,0.611765,0.411765}%
\pgfsetfillcolor{currentfill}%
\pgfsetfillopacity{0.180000}%
\pgfsetlinewidth{0.000000pt}%
\definecolor{currentstroke}{rgb}{0.000000,0.000000,0.000000}%
\pgfsetstrokecolor{currentstroke}%
\pgfsetdash{}{0pt}%
\pgfpathmoveto{\pgfqpoint{1.452161in}{1.021787in}}%
\pgfpathlineto{\pgfqpoint{1.467193in}{1.029962in}}%
\pgfpathlineto{\pgfqpoint{1.467612in}{1.052514in}}%
\pgfpathlineto{\pgfqpoint{1.452562in}{1.044401in}}%
\pgfpathlineto{\pgfqpoint{1.452161in}{1.021787in}}%
\pgfpathclose%
\pgfusepath{fill}%
\end{pgfscope}%
\begin{pgfscope}%
\pgfpathrectangle{\pgfqpoint{0.407592in}{0.222904in}}{\pgfqpoint{1.411163in}{1.411163in}}%
\pgfusepath{clip}%
\pgfsetbuttcap%
\pgfsetroundjoin%
\definecolor{currentfill}{rgb}{0.168627,0.611765,0.411765}%
\pgfsetfillcolor{currentfill}%
\pgfsetfillopacity{0.180000}%
\pgfsetlinewidth{0.000000pt}%
\definecolor{currentstroke}{rgb}{0.000000,0.000000,0.000000}%
\pgfsetstrokecolor{currentstroke}%
\pgfsetdash{}{0pt}%
\pgfpathmoveto{\pgfqpoint{1.484352in}{1.152750in}}%
\pgfpathlineto{\pgfqpoint{1.499347in}{1.163798in}}%
\pgfpathlineto{\pgfqpoint{1.499805in}{1.186520in}}%
\pgfpathlineto{\pgfqpoint{1.484793in}{1.175528in}}%
\pgfpathlineto{\pgfqpoint{1.484352in}{1.152750in}}%
\pgfpathclose%
\pgfusepath{fill}%
\end{pgfscope}%
\begin{pgfscope}%
\pgfpathrectangle{\pgfqpoint{0.407592in}{0.222904in}}{\pgfqpoint{1.411163in}{1.411163in}}%
\pgfusepath{clip}%
\pgfsetbuttcap%
\pgfsetroundjoin%
\definecolor{currentfill}{rgb}{0.168627,0.611765,0.411765}%
\pgfsetfillcolor{currentfill}%
\pgfsetfillopacity{0.180000}%
\pgfsetlinewidth{0.000000pt}%
\definecolor{currentstroke}{rgb}{0.000000,0.000000,0.000000}%
\pgfsetstrokecolor{currentstroke}%
\pgfsetdash{}{0pt}%
\pgfpathmoveto{\pgfqpoint{1.420070in}{0.897219in}}%
\pgfpathlineto{\pgfqpoint{1.435158in}{0.902627in}}%
\pgfpathlineto{\pgfqpoint{1.435536in}{0.925025in}}%
\pgfpathlineto{\pgfqpoint{1.420431in}{0.919686in}}%
\pgfpathlineto{\pgfqpoint{1.420070in}{0.897219in}}%
\pgfpathclose%
\pgfusepath{fill}%
\end{pgfscope}%
\begin{pgfscope}%
\pgfpathrectangle{\pgfqpoint{0.407592in}{0.222904in}}{\pgfqpoint{1.411163in}{1.411163in}}%
\pgfusepath{clip}%
\pgfsetbuttcap%
\pgfsetroundjoin%
\definecolor{currentfill}{rgb}{0.168627,0.611765,0.411765}%
\pgfsetfillcolor{currentfill}%
\pgfsetfillopacity{0.180000}%
\pgfsetlinewidth{0.000000pt}%
\definecolor{currentstroke}{rgb}{0.000000,0.000000,0.000000}%
\pgfsetstrokecolor{currentstroke}%
\pgfsetdash{}{0pt}%
\pgfpathmoveto{\pgfqpoint{1.468452in}{1.097788in}}%
\pgfpathlineto{\pgfqpoint{1.483472in}{1.107367in}}%
\pgfpathlineto{\pgfqpoint{1.483911in}{1.130030in}}%
\pgfpathlineto{\pgfqpoint{1.468874in}{1.120510in}}%
\pgfpathlineto{\pgfqpoint{1.468452in}{1.097788in}}%
\pgfpathclose%
\pgfusepath{fill}%
\end{pgfscope}%
\begin{pgfscope}%
\pgfpathrectangle{\pgfqpoint{0.407592in}{0.222904in}}{\pgfqpoint{1.411163in}{1.411163in}}%
\pgfusepath{clip}%
\pgfsetbuttcap%
\pgfsetroundjoin%
\definecolor{currentfill}{rgb}{0.168627,0.611765,0.411765}%
\pgfsetfillcolor{currentfill}%
\pgfsetfillopacity{0.180000}%
\pgfsetlinewidth{0.000000pt}%
\definecolor{currentstroke}{rgb}{0.000000,0.000000,0.000000}%
\pgfsetstrokecolor{currentstroke}%
\pgfsetdash{}{0pt}%
\pgfpathmoveto{\pgfqpoint{1.436296in}{0.969990in}}%
\pgfpathlineto{\pgfqpoint{1.451362in}{0.976727in}}%
\pgfpathlineto{\pgfqpoint{1.451761in}{0.999229in}}%
\pgfpathlineto{\pgfqpoint{1.436677in}{0.992557in}}%
\pgfpathlineto{\pgfqpoint{1.436296in}{0.969990in}}%
\pgfpathclose%
\pgfusepath{fill}%
\end{pgfscope}%
\begin{pgfscope}%
\pgfpathrectangle{\pgfqpoint{0.407592in}{0.222904in}}{\pgfqpoint{1.411163in}{1.411163in}}%
\pgfusepath{clip}%
\pgfsetbuttcap%
\pgfsetroundjoin%
\definecolor{currentfill}{rgb}{0.168627,0.611765,0.411765}%
\pgfsetfillcolor{currentfill}%
\pgfsetfillopacity{0.180000}%
\pgfsetlinewidth{0.000000pt}%
\definecolor{currentstroke}{rgb}{0.000000,0.000000,0.000000}%
\pgfsetstrokecolor{currentstroke}%
\pgfsetdash{}{0pt}%
\pgfpathmoveto{\pgfqpoint{1.452562in}{1.044401in}}%
\pgfpathlineto{\pgfqpoint{1.467612in}{1.052514in}}%
\pgfpathlineto{\pgfqpoint{1.468031in}{1.075123in}}%
\pgfpathlineto{\pgfqpoint{1.452964in}{1.067072in}}%
\pgfpathlineto{\pgfqpoint{1.452562in}{1.044401in}}%
\pgfpathclose%
\pgfusepath{fill}%
\end{pgfscope}%
\begin{pgfscope}%
\pgfpathrectangle{\pgfqpoint{0.407592in}{0.222904in}}{\pgfqpoint{1.411163in}{1.411163in}}%
\pgfusepath{clip}%
\pgfsetbuttcap%
\pgfsetroundjoin%
\definecolor{currentfill}{rgb}{0.168627,0.611765,0.411765}%
\pgfsetfillcolor{currentfill}%
\pgfsetfillopacity{0.180000}%
\pgfsetlinewidth{0.000000pt}%
\definecolor{currentstroke}{rgb}{0.000000,0.000000,0.000000}%
\pgfsetstrokecolor{currentstroke}%
\pgfsetdash{}{0pt}%
\pgfpathmoveto{\pgfqpoint{1.484793in}{1.175528in}}%
\pgfpathlineto{\pgfqpoint{1.499805in}{1.186520in}}%
\pgfpathlineto{\pgfqpoint{1.500265in}{1.209299in}}%
\pgfpathlineto{\pgfqpoint{1.485235in}{1.198362in}}%
\pgfpathlineto{\pgfqpoint{1.484793in}{1.175528in}}%
\pgfpathclose%
\pgfusepath{fill}%
\end{pgfscope}%
\begin{pgfscope}%
\pgfpathrectangle{\pgfqpoint{0.407592in}{0.222904in}}{\pgfqpoint{1.411163in}{1.411163in}}%
\pgfusepath{clip}%
\pgfsetbuttcap%
\pgfsetroundjoin%
\definecolor{currentfill}{rgb}{0.168627,0.611765,0.411765}%
\pgfsetfillcolor{currentfill}%
\pgfsetfillopacity{0.180000}%
\pgfsetlinewidth{0.000000pt}%
\definecolor{currentstroke}{rgb}{0.000000,0.000000,0.000000}%
\pgfsetstrokecolor{currentstroke}%
\pgfsetdash{}{0pt}%
\pgfpathmoveto{\pgfqpoint{1.420431in}{0.919686in}}%
\pgfpathlineto{\pgfqpoint{1.435536in}{0.925025in}}%
\pgfpathlineto{\pgfqpoint{1.435916in}{0.947480in}}%
\pgfpathlineto{\pgfqpoint{1.420793in}{0.942210in}}%
\pgfpathlineto{\pgfqpoint{1.420431in}{0.919686in}}%
\pgfpathclose%
\pgfusepath{fill}%
\end{pgfscope}%
\begin{pgfscope}%
\pgfpathrectangle{\pgfqpoint{0.407592in}{0.222904in}}{\pgfqpoint{1.411163in}{1.411163in}}%
\pgfusepath{clip}%
\pgfsetbuttcap%
\pgfsetroundjoin%
\definecolor{currentfill}{rgb}{0.168627,0.611765,0.411765}%
\pgfsetfillcolor{currentfill}%
\pgfsetfillopacity{0.180000}%
\pgfsetlinewidth{0.000000pt}%
\definecolor{currentstroke}{rgb}{0.000000,0.000000,0.000000}%
\pgfsetstrokecolor{currentstroke}%
\pgfsetdash{}{0pt}%
\pgfpathmoveto{\pgfqpoint{1.468874in}{1.120510in}}%
\pgfpathlineto{\pgfqpoint{1.483911in}{1.130030in}}%
\pgfpathlineto{\pgfqpoint{1.484352in}{1.152750in}}%
\pgfpathlineto{\pgfqpoint{1.469296in}{1.143289in}}%
\pgfpathlineto{\pgfqpoint{1.468874in}{1.120510in}}%
\pgfpathclose%
\pgfusepath{fill}%
\end{pgfscope}%
\begin{pgfscope}%
\pgfpathrectangle{\pgfqpoint{0.407592in}{0.222904in}}{\pgfqpoint{1.411163in}{1.411163in}}%
\pgfusepath{clip}%
\pgfsetbuttcap%
\pgfsetroundjoin%
\definecolor{currentfill}{rgb}{0.168627,0.611765,0.411765}%
\pgfsetfillcolor{currentfill}%
\pgfsetfillopacity{0.180000}%
\pgfsetlinewidth{0.000000pt}%
\definecolor{currentstroke}{rgb}{0.000000,0.000000,0.000000}%
\pgfsetstrokecolor{currentstroke}%
\pgfsetdash{}{0pt}%
\pgfpathmoveto{\pgfqpoint{1.436677in}{0.992557in}}%
\pgfpathlineto{\pgfqpoint{1.451761in}{0.999229in}}%
\pgfpathlineto{\pgfqpoint{1.452161in}{1.021787in}}%
\pgfpathlineto{\pgfqpoint{1.437060in}{1.015181in}}%
\pgfpathlineto{\pgfqpoint{1.436677in}{0.992557in}}%
\pgfpathclose%
\pgfusepath{fill}%
\end{pgfscope}%
\begin{pgfscope}%
\pgfpathrectangle{\pgfqpoint{0.407592in}{0.222904in}}{\pgfqpoint{1.411163in}{1.411163in}}%
\pgfusepath{clip}%
\pgfsetbuttcap%
\pgfsetroundjoin%
\definecolor{currentfill}{rgb}{0.168627,0.611765,0.411765}%
\pgfsetfillcolor{currentfill}%
\pgfsetfillopacity{0.180000}%
\pgfsetlinewidth{0.000000pt}%
\definecolor{currentstroke}{rgb}{0.000000,0.000000,0.000000}%
\pgfsetstrokecolor{currentstroke}%
\pgfsetdash{}{0pt}%
\pgfpathmoveto{\pgfqpoint{1.404562in}{0.870811in}}%
\pgfpathlineto{\pgfqpoint{1.419710in}{0.874808in}}%
\pgfpathlineto{\pgfqpoint{1.420070in}{0.897219in}}%
\pgfpathlineto{\pgfqpoint{1.404904in}{0.893295in}}%
\pgfpathlineto{\pgfqpoint{1.404562in}{0.870811in}}%
\pgfpathclose%
\pgfusepath{fill}%
\end{pgfscope}%
\begin{pgfscope}%
\pgfpathrectangle{\pgfqpoint{0.407592in}{0.222904in}}{\pgfqpoint{1.411163in}{1.411163in}}%
\pgfusepath{clip}%
\pgfsetbuttcap%
\pgfsetroundjoin%
\definecolor{currentfill}{rgb}{0.168627,0.611765,0.411765}%
\pgfsetfillcolor{currentfill}%
\pgfsetfillopacity{0.180000}%
\pgfsetlinewidth{0.000000pt}%
\definecolor{currentstroke}{rgb}{0.000000,0.000000,0.000000}%
\pgfsetstrokecolor{currentstroke}%
\pgfsetdash{}{0pt}%
\pgfpathmoveto{\pgfqpoint{1.452964in}{1.067072in}}%
\pgfpathlineto{\pgfqpoint{1.468031in}{1.075123in}}%
\pgfpathlineto{\pgfqpoint{1.468452in}{1.097788in}}%
\pgfpathlineto{\pgfqpoint{1.453367in}{1.089800in}}%
\pgfpathlineto{\pgfqpoint{1.452964in}{1.067072in}}%
\pgfpathclose%
\pgfusepath{fill}%
\end{pgfscope}%
\begin{pgfscope}%
\pgfpathrectangle{\pgfqpoint{0.407592in}{0.222904in}}{\pgfqpoint{1.411163in}{1.411163in}}%
\pgfusepath{clip}%
\pgfsetbuttcap%
\pgfsetroundjoin%
\definecolor{currentfill}{rgb}{0.168627,0.611765,0.411765}%
\pgfsetfillcolor{currentfill}%
\pgfsetfillopacity{0.180000}%
\pgfsetlinewidth{0.000000pt}%
\definecolor{currentstroke}{rgb}{0.000000,0.000000,0.000000}%
\pgfsetstrokecolor{currentstroke}%
\pgfsetdash{}{0pt}%
\pgfpathmoveto{\pgfqpoint{1.420793in}{0.942210in}}%
\pgfpathlineto{\pgfqpoint{1.435916in}{0.947480in}}%
\pgfpathlineto{\pgfqpoint{1.436296in}{0.969990in}}%
\pgfpathlineto{\pgfqpoint{1.421156in}{0.964790in}}%
\pgfpathlineto{\pgfqpoint{1.420793in}{0.942210in}}%
\pgfpathclose%
\pgfusepath{fill}%
\end{pgfscope}%
\begin{pgfscope}%
\pgfpathrectangle{\pgfqpoint{0.407592in}{0.222904in}}{\pgfqpoint{1.411163in}{1.411163in}}%
\pgfusepath{clip}%
\pgfsetbuttcap%
\pgfsetroundjoin%
\definecolor{currentfill}{rgb}{0.168627,0.611765,0.411765}%
\pgfsetfillcolor{currentfill}%
\pgfsetfillopacity{0.180000}%
\pgfsetlinewidth{0.000000pt}%
\definecolor{currentstroke}{rgb}{0.000000,0.000000,0.000000}%
\pgfsetstrokecolor{currentstroke}%
\pgfsetdash{}{0pt}%
\pgfpathmoveto{\pgfqpoint{1.469296in}{1.143289in}}%
\pgfpathlineto{\pgfqpoint{1.484352in}{1.152750in}}%
\pgfpathlineto{\pgfqpoint{1.484793in}{1.175528in}}%
\pgfpathlineto{\pgfqpoint{1.469720in}{1.166126in}}%
\pgfpathlineto{\pgfqpoint{1.469296in}{1.143289in}}%
\pgfpathclose%
\pgfusepath{fill}%
\end{pgfscope}%
\begin{pgfscope}%
\pgfpathrectangle{\pgfqpoint{0.407592in}{0.222904in}}{\pgfqpoint{1.411163in}{1.411163in}}%
\pgfusepath{clip}%
\pgfsetbuttcap%
\pgfsetroundjoin%
\definecolor{currentfill}{rgb}{0.168627,0.611765,0.411765}%
\pgfsetfillcolor{currentfill}%
\pgfsetfillopacity{0.180000}%
\pgfsetlinewidth{0.000000pt}%
\definecolor{currentstroke}{rgb}{0.000000,0.000000,0.000000}%
\pgfsetstrokecolor{currentstroke}%
\pgfsetdash{}{0pt}%
\pgfpathmoveto{\pgfqpoint{1.437060in}{1.015181in}}%
\pgfpathlineto{\pgfqpoint{1.452161in}{1.021787in}}%
\pgfpathlineto{\pgfqpoint{1.452562in}{1.044401in}}%
\pgfpathlineto{\pgfqpoint{1.437443in}{1.037862in}}%
\pgfpathlineto{\pgfqpoint{1.437060in}{1.015181in}}%
\pgfpathclose%
\pgfusepath{fill}%
\end{pgfscope}%
\begin{pgfscope}%
\pgfpathrectangle{\pgfqpoint{0.407592in}{0.222904in}}{\pgfqpoint{1.411163in}{1.411163in}}%
\pgfusepath{clip}%
\pgfsetbuttcap%
\pgfsetroundjoin%
\definecolor{currentfill}{rgb}{0.168627,0.611765,0.411765}%
\pgfsetfillcolor{currentfill}%
\pgfsetfillopacity{0.180000}%
\pgfsetlinewidth{0.000000pt}%
\definecolor{currentstroke}{rgb}{0.000000,0.000000,0.000000}%
\pgfsetstrokecolor{currentstroke}%
\pgfsetdash{}{0pt}%
\pgfpathmoveto{\pgfqpoint{1.404904in}{0.893295in}}%
\pgfpathlineto{\pgfqpoint{1.420070in}{0.897219in}}%
\pgfpathlineto{\pgfqpoint{1.420431in}{0.919686in}}%
\pgfpathlineto{\pgfqpoint{1.405247in}{0.915836in}}%
\pgfpathlineto{\pgfqpoint{1.404904in}{0.893295in}}%
\pgfpathclose%
\pgfusepath{fill}%
\end{pgfscope}%
\begin{pgfscope}%
\pgfpathrectangle{\pgfqpoint{0.407592in}{0.222904in}}{\pgfqpoint{1.411163in}{1.411163in}}%
\pgfusepath{clip}%
\pgfsetbuttcap%
\pgfsetroundjoin%
\definecolor{currentfill}{rgb}{0.168627,0.611765,0.411765}%
\pgfsetfillcolor{currentfill}%
\pgfsetfillopacity{0.180000}%
\pgfsetlinewidth{0.000000pt}%
\definecolor{currentstroke}{rgb}{0.000000,0.000000,0.000000}%
\pgfsetstrokecolor{currentstroke}%
\pgfsetdash{}{0pt}%
\pgfpathmoveto{\pgfqpoint{1.453367in}{1.089800in}}%
\pgfpathlineto{\pgfqpoint{1.468452in}{1.097788in}}%
\pgfpathlineto{\pgfqpoint{1.468874in}{1.120510in}}%
\pgfpathlineto{\pgfqpoint{1.453771in}{1.112585in}}%
\pgfpathlineto{\pgfqpoint{1.453367in}{1.089800in}}%
\pgfpathclose%
\pgfusepath{fill}%
\end{pgfscope}%
\begin{pgfscope}%
\pgfpathrectangle{\pgfqpoint{0.407592in}{0.222904in}}{\pgfqpoint{1.411163in}{1.411163in}}%
\pgfusepath{clip}%
\pgfsetbuttcap%
\pgfsetroundjoin%
\definecolor{currentfill}{rgb}{0.168627,0.611765,0.411765}%
\pgfsetfillcolor{currentfill}%
\pgfsetfillopacity{0.180000}%
\pgfsetlinewidth{0.000000pt}%
\definecolor{currentstroke}{rgb}{0.000000,0.000000,0.000000}%
\pgfsetstrokecolor{currentstroke}%
\pgfsetdash{}{0pt}%
\pgfpathmoveto{\pgfqpoint{1.421156in}{0.964790in}}%
\pgfpathlineto{\pgfqpoint{1.436296in}{0.969990in}}%
\pgfpathlineto{\pgfqpoint{1.436677in}{0.992557in}}%
\pgfpathlineto{\pgfqpoint{1.421519in}{0.987428in}}%
\pgfpathlineto{\pgfqpoint{1.421156in}{0.964790in}}%
\pgfpathclose%
\pgfusepath{fill}%
\end{pgfscope}%
\begin{pgfscope}%
\pgfpathrectangle{\pgfqpoint{0.407592in}{0.222904in}}{\pgfqpoint{1.411163in}{1.411163in}}%
\pgfusepath{clip}%
\pgfsetbuttcap%
\pgfsetroundjoin%
\definecolor{currentfill}{rgb}{0.168627,0.611765,0.411765}%
\pgfsetfillcolor{currentfill}%
\pgfsetfillopacity{0.180000}%
\pgfsetlinewidth{0.000000pt}%
\definecolor{currentstroke}{rgb}{0.000000,0.000000,0.000000}%
\pgfsetstrokecolor{currentstroke}%
\pgfsetdash{}{0pt}%
\pgfpathmoveto{\pgfqpoint{1.469720in}{1.166126in}}%
\pgfpathlineto{\pgfqpoint{1.484793in}{1.175528in}}%
\pgfpathlineto{\pgfqpoint{1.485235in}{1.198362in}}%
\pgfpathlineto{\pgfqpoint{1.470145in}{1.189020in}}%
\pgfpathlineto{\pgfqpoint{1.469720in}{1.166126in}}%
\pgfpathclose%
\pgfusepath{fill}%
\end{pgfscope}%
\begin{pgfscope}%
\pgfpathrectangle{\pgfqpoint{0.407592in}{0.222904in}}{\pgfqpoint{1.411163in}{1.411163in}}%
\pgfusepath{clip}%
\pgfsetbuttcap%
\pgfsetroundjoin%
\definecolor{currentfill}{rgb}{0.168627,0.611765,0.411765}%
\pgfsetfillcolor{currentfill}%
\pgfsetfillopacity{0.180000}%
\pgfsetlinewidth{0.000000pt}%
\definecolor{currentstroke}{rgb}{0.000000,0.000000,0.000000}%
\pgfsetstrokecolor{currentstroke}%
\pgfsetdash{}{0pt}%
\pgfpathmoveto{\pgfqpoint{1.437443in}{1.037862in}}%
\pgfpathlineto{\pgfqpoint{1.452562in}{1.044401in}}%
\pgfpathlineto{\pgfqpoint{1.452964in}{1.067072in}}%
\pgfpathlineto{\pgfqpoint{1.437827in}{1.060600in}}%
\pgfpathlineto{\pgfqpoint{1.437443in}{1.037862in}}%
\pgfpathclose%
\pgfusepath{fill}%
\end{pgfscope}%
\begin{pgfscope}%
\pgfpathrectangle{\pgfqpoint{0.407592in}{0.222904in}}{\pgfqpoint{1.411163in}{1.411163in}}%
\pgfusepath{clip}%
\pgfsetbuttcap%
\pgfsetroundjoin%
\definecolor{currentfill}{rgb}{0.168627,0.611765,0.411765}%
\pgfsetfillcolor{currentfill}%
\pgfsetfillopacity{0.180000}%
\pgfsetlinewidth{0.000000pt}%
\definecolor{currentstroke}{rgb}{0.000000,0.000000,0.000000}%
\pgfsetstrokecolor{currentstroke}%
\pgfsetdash{}{0pt}%
\pgfpathmoveto{\pgfqpoint{1.405247in}{0.915836in}}%
\pgfpathlineto{\pgfqpoint{1.420431in}{0.919686in}}%
\pgfpathlineto{\pgfqpoint{1.420793in}{0.942210in}}%
\pgfpathlineto{\pgfqpoint{1.405591in}{0.938433in}}%
\pgfpathlineto{\pgfqpoint{1.405247in}{0.915836in}}%
\pgfpathclose%
\pgfusepath{fill}%
\end{pgfscope}%
\begin{pgfscope}%
\pgfpathrectangle{\pgfqpoint{0.407592in}{0.222904in}}{\pgfqpoint{1.411163in}{1.411163in}}%
\pgfusepath{clip}%
\pgfsetbuttcap%
\pgfsetroundjoin%
\definecolor{currentfill}{rgb}{0.168627,0.611765,0.411765}%
\pgfsetfillcolor{currentfill}%
\pgfsetfillopacity{0.180000}%
\pgfsetlinewidth{0.000000pt}%
\definecolor{currentstroke}{rgb}{0.000000,0.000000,0.000000}%
\pgfsetstrokecolor{currentstroke}%
\pgfsetdash{}{0pt}%
\pgfpathmoveto{\pgfqpoint{1.453771in}{1.112585in}}%
\pgfpathlineto{\pgfqpoint{1.468874in}{1.120510in}}%
\pgfpathlineto{\pgfqpoint{1.469296in}{1.143289in}}%
\pgfpathlineto{\pgfqpoint{1.454176in}{1.135428in}}%
\pgfpathlineto{\pgfqpoint{1.453771in}{1.112585in}}%
\pgfpathclose%
\pgfusepath{fill}%
\end{pgfscope}%
\begin{pgfscope}%
\pgfpathrectangle{\pgfqpoint{0.407592in}{0.222904in}}{\pgfqpoint{1.411163in}{1.411163in}}%
\pgfusepath{clip}%
\pgfsetbuttcap%
\pgfsetroundjoin%
\definecolor{currentfill}{rgb}{0.168627,0.611765,0.411765}%
\pgfsetfillcolor{currentfill}%
\pgfsetfillopacity{0.180000}%
\pgfsetlinewidth{0.000000pt}%
\definecolor{currentstroke}{rgb}{0.000000,0.000000,0.000000}%
\pgfsetstrokecolor{currentstroke}%
\pgfsetdash{}{0pt}%
\pgfpathmoveto{\pgfqpoint{1.421519in}{0.987428in}}%
\pgfpathlineto{\pgfqpoint{1.436677in}{0.992557in}}%
\pgfpathlineto{\pgfqpoint{1.437060in}{1.015181in}}%
\pgfpathlineto{\pgfqpoint{1.421884in}{1.010122in}}%
\pgfpathlineto{\pgfqpoint{1.421519in}{0.987428in}}%
\pgfpathclose%
\pgfusepath{fill}%
\end{pgfscope}%
\begin{pgfscope}%
\pgfpathrectangle{\pgfqpoint{0.407592in}{0.222904in}}{\pgfqpoint{1.411163in}{1.411163in}}%
\pgfusepath{clip}%
\pgfsetbuttcap%
\pgfsetroundjoin%
\definecolor{currentfill}{rgb}{0.168627,0.611765,0.411765}%
\pgfsetfillcolor{currentfill}%
\pgfsetfillopacity{0.180000}%
\pgfsetlinewidth{0.000000pt}%
\definecolor{currentstroke}{rgb}{0.000000,0.000000,0.000000}%
\pgfsetstrokecolor{currentstroke}%
\pgfsetdash{}{0pt}%
\pgfpathmoveto{\pgfqpoint{1.437827in}{1.060600in}}%
\pgfpathlineto{\pgfqpoint{1.452964in}{1.067072in}}%
\pgfpathlineto{\pgfqpoint{1.453367in}{1.089800in}}%
\pgfpathlineto{\pgfqpoint{1.438213in}{1.083395in}}%
\pgfpathlineto{\pgfqpoint{1.437827in}{1.060600in}}%
\pgfpathclose%
\pgfusepath{fill}%
\end{pgfscope}%
\begin{pgfscope}%
\pgfpathrectangle{\pgfqpoint{0.407592in}{0.222904in}}{\pgfqpoint{1.411163in}{1.411163in}}%
\pgfusepath{clip}%
\pgfsetbuttcap%
\pgfsetroundjoin%
\definecolor{currentfill}{rgb}{0.168627,0.611765,0.411765}%
\pgfsetfillcolor{currentfill}%
\pgfsetfillopacity{0.180000}%
\pgfsetlinewidth{0.000000pt}%
\definecolor{currentstroke}{rgb}{0.000000,0.000000,0.000000}%
\pgfsetstrokecolor{currentstroke}%
\pgfsetdash{}{0pt}%
\pgfpathmoveto{\pgfqpoint{1.389328in}{0.868242in}}%
\pgfpathlineto{\pgfqpoint{1.404562in}{0.870811in}}%
\pgfpathlineto{\pgfqpoint{1.404904in}{0.893295in}}%
\pgfpathlineto{\pgfqpoint{1.389653in}{0.890803in}}%
\pgfpathlineto{\pgfqpoint{1.389328in}{0.868242in}}%
\pgfpathclose%
\pgfusepath{fill}%
\end{pgfscope}%
\begin{pgfscope}%
\pgfpathrectangle{\pgfqpoint{0.407592in}{0.222904in}}{\pgfqpoint{1.411163in}{1.411163in}}%
\pgfusepath{clip}%
\pgfsetbuttcap%
\pgfsetroundjoin%
\definecolor{currentfill}{rgb}{0.168627,0.611765,0.411765}%
\pgfsetfillcolor{currentfill}%
\pgfsetfillopacity{0.180000}%
\pgfsetlinewidth{0.000000pt}%
\definecolor{currentstroke}{rgb}{0.000000,0.000000,0.000000}%
\pgfsetstrokecolor{currentstroke}%
\pgfsetdash{}{0pt}%
\pgfpathmoveto{\pgfqpoint{1.405591in}{0.938433in}}%
\pgfpathlineto{\pgfqpoint{1.420793in}{0.942210in}}%
\pgfpathlineto{\pgfqpoint{1.421156in}{0.964790in}}%
\pgfpathlineto{\pgfqpoint{1.405936in}{0.961088in}}%
\pgfpathlineto{\pgfqpoint{1.405591in}{0.938433in}}%
\pgfpathclose%
\pgfusepath{fill}%
\end{pgfscope}%
\begin{pgfscope}%
\pgfpathrectangle{\pgfqpoint{0.407592in}{0.222904in}}{\pgfqpoint{1.411163in}{1.411163in}}%
\pgfusepath{clip}%
\pgfsetbuttcap%
\pgfsetroundjoin%
\definecolor{currentfill}{rgb}{0.168627,0.611765,0.411765}%
\pgfsetfillcolor{currentfill}%
\pgfsetfillopacity{0.180000}%
\pgfsetlinewidth{0.000000pt}%
\definecolor{currentstroke}{rgb}{0.000000,0.000000,0.000000}%
\pgfsetstrokecolor{currentstroke}%
\pgfsetdash{}{0pt}%
\pgfpathmoveto{\pgfqpoint{1.454176in}{1.135428in}}%
\pgfpathlineto{\pgfqpoint{1.469296in}{1.143289in}}%
\pgfpathlineto{\pgfqpoint{1.469720in}{1.166126in}}%
\pgfpathlineto{\pgfqpoint{1.454582in}{1.158328in}}%
\pgfpathlineto{\pgfqpoint{1.454176in}{1.135428in}}%
\pgfpathclose%
\pgfusepath{fill}%
\end{pgfscope}%
\begin{pgfscope}%
\pgfpathrectangle{\pgfqpoint{0.407592in}{0.222904in}}{\pgfqpoint{1.411163in}{1.411163in}}%
\pgfusepath{clip}%
\pgfsetbuttcap%
\pgfsetroundjoin%
\definecolor{currentfill}{rgb}{0.168627,0.611765,0.411765}%
\pgfsetfillcolor{currentfill}%
\pgfsetfillopacity{0.180000}%
\pgfsetlinewidth{0.000000pt}%
\definecolor{currentstroke}{rgb}{0.000000,0.000000,0.000000}%
\pgfsetstrokecolor{currentstroke}%
\pgfsetdash{}{0pt}%
\pgfpathmoveto{\pgfqpoint{1.421884in}{1.010122in}}%
\pgfpathlineto{\pgfqpoint{1.437060in}{1.015181in}}%
\pgfpathlineto{\pgfqpoint{1.437443in}{1.037862in}}%
\pgfpathlineto{\pgfqpoint{1.422249in}{1.032873in}}%
\pgfpathlineto{\pgfqpoint{1.421884in}{1.010122in}}%
\pgfpathclose%
\pgfusepath{fill}%
\end{pgfscope}%
\begin{pgfscope}%
\pgfpathrectangle{\pgfqpoint{0.407592in}{0.222904in}}{\pgfqpoint{1.411163in}{1.411163in}}%
\pgfusepath{clip}%
\pgfsetbuttcap%
\pgfsetroundjoin%
\definecolor{currentfill}{rgb}{0.168627,0.611765,0.411765}%
\pgfsetfillcolor{currentfill}%
\pgfsetfillopacity{0.180000}%
\pgfsetlinewidth{0.000000pt}%
\definecolor{currentstroke}{rgb}{0.000000,0.000000,0.000000}%
\pgfsetstrokecolor{currentstroke}%
\pgfsetdash{}{0pt}%
\pgfpathmoveto{\pgfqpoint{1.438213in}{1.083395in}}%
\pgfpathlineto{\pgfqpoint{1.453367in}{1.089800in}}%
\pgfpathlineto{\pgfqpoint{1.453771in}{1.112585in}}%
\pgfpathlineto{\pgfqpoint{1.438599in}{1.106248in}}%
\pgfpathlineto{\pgfqpoint{1.438213in}{1.083395in}}%
\pgfpathclose%
\pgfusepath{fill}%
\end{pgfscope}%
\begin{pgfscope}%
\pgfpathrectangle{\pgfqpoint{0.407592in}{0.222904in}}{\pgfqpoint{1.411163in}{1.411163in}}%
\pgfusepath{clip}%
\pgfsetbuttcap%
\pgfsetroundjoin%
\definecolor{currentfill}{rgb}{0.168627,0.611765,0.411765}%
\pgfsetfillcolor{currentfill}%
\pgfsetfillopacity{0.180000}%
\pgfsetlinewidth{0.000000pt}%
\definecolor{currentstroke}{rgb}{0.000000,0.000000,0.000000}%
\pgfsetstrokecolor{currentstroke}%
\pgfsetdash{}{0pt}%
\pgfpathmoveto{\pgfqpoint{1.389653in}{0.890803in}}%
\pgfpathlineto{\pgfqpoint{1.404904in}{0.893295in}}%
\pgfpathlineto{\pgfqpoint{1.405247in}{0.915836in}}%
\pgfpathlineto{\pgfqpoint{1.389978in}{0.913421in}}%
\pgfpathlineto{\pgfqpoint{1.389653in}{0.890803in}}%
\pgfpathclose%
\pgfusepath{fill}%
\end{pgfscope}%
\begin{pgfscope}%
\pgfpathrectangle{\pgfqpoint{0.407592in}{0.222904in}}{\pgfqpoint{1.411163in}{1.411163in}}%
\pgfusepath{clip}%
\pgfsetbuttcap%
\pgfsetroundjoin%
\definecolor{currentfill}{rgb}{0.168627,0.611765,0.411765}%
\pgfsetfillcolor{currentfill}%
\pgfsetfillopacity{0.180000}%
\pgfsetlinewidth{0.000000pt}%
\definecolor{currentstroke}{rgb}{0.000000,0.000000,0.000000}%
\pgfsetstrokecolor{currentstroke}%
\pgfsetdash{}{0pt}%
\pgfpathmoveto{\pgfqpoint{1.405936in}{0.961088in}}%
\pgfpathlineto{\pgfqpoint{1.421156in}{0.964790in}}%
\pgfpathlineto{\pgfqpoint{1.421519in}{0.987428in}}%
\pgfpathlineto{\pgfqpoint{1.406281in}{0.983799in}}%
\pgfpathlineto{\pgfqpoint{1.405936in}{0.961088in}}%
\pgfpathclose%
\pgfusepath{fill}%
\end{pgfscope}%
\begin{pgfscope}%
\pgfpathrectangle{\pgfqpoint{0.407592in}{0.222904in}}{\pgfqpoint{1.411163in}{1.411163in}}%
\pgfusepath{clip}%
\pgfsetbuttcap%
\pgfsetroundjoin%
\definecolor{currentfill}{rgb}{0.168627,0.611765,0.411765}%
\pgfsetfillcolor{currentfill}%
\pgfsetfillopacity{0.180000}%
\pgfsetlinewidth{0.000000pt}%
\definecolor{currentstroke}{rgb}{0.000000,0.000000,0.000000}%
\pgfsetstrokecolor{currentstroke}%
\pgfsetdash{}{0pt}%
\pgfpathmoveto{\pgfqpoint{1.454582in}{1.158328in}}%
\pgfpathlineto{\pgfqpoint{1.469720in}{1.166126in}}%
\pgfpathlineto{\pgfqpoint{1.470145in}{1.189020in}}%
\pgfpathlineto{\pgfqpoint{1.454989in}{1.181286in}}%
\pgfpathlineto{\pgfqpoint{1.454582in}{1.158328in}}%
\pgfpathclose%
\pgfusepath{fill}%
\end{pgfscope}%
\begin{pgfscope}%
\pgfpathrectangle{\pgfqpoint{0.407592in}{0.222904in}}{\pgfqpoint{1.411163in}{1.411163in}}%
\pgfusepath{clip}%
\pgfsetbuttcap%
\pgfsetroundjoin%
\definecolor{currentfill}{rgb}{0.168627,0.611765,0.411765}%
\pgfsetfillcolor{currentfill}%
\pgfsetfillopacity{0.180000}%
\pgfsetlinewidth{0.000000pt}%
\definecolor{currentstroke}{rgb}{0.000000,0.000000,0.000000}%
\pgfsetstrokecolor{currentstroke}%
\pgfsetdash{}{0pt}%
\pgfpathmoveto{\pgfqpoint{1.422249in}{1.032873in}}%
\pgfpathlineto{\pgfqpoint{1.437443in}{1.037862in}}%
\pgfpathlineto{\pgfqpoint{1.437827in}{1.060600in}}%
\pgfpathlineto{\pgfqpoint{1.422616in}{1.055682in}}%
\pgfpathlineto{\pgfqpoint{1.422249in}{1.032873in}}%
\pgfpathclose%
\pgfusepath{fill}%
\end{pgfscope}%
\begin{pgfscope}%
\pgfpathrectangle{\pgfqpoint{0.407592in}{0.222904in}}{\pgfqpoint{1.411163in}{1.411163in}}%
\pgfusepath{clip}%
\pgfsetbuttcap%
\pgfsetroundjoin%
\definecolor{currentfill}{rgb}{0.168627,0.611765,0.411765}%
\pgfsetfillcolor{currentfill}%
\pgfsetfillopacity{0.180000}%
\pgfsetlinewidth{0.000000pt}%
\definecolor{currentstroke}{rgb}{0.000000,0.000000,0.000000}%
\pgfsetstrokecolor{currentstroke}%
\pgfsetdash{}{0pt}%
\pgfpathmoveto{\pgfqpoint{1.438599in}{1.106248in}}%
\pgfpathlineto{\pgfqpoint{1.453771in}{1.112585in}}%
\pgfpathlineto{\pgfqpoint{1.454176in}{1.135428in}}%
\pgfpathlineto{\pgfqpoint{1.438986in}{1.129158in}}%
\pgfpathlineto{\pgfqpoint{1.438599in}{1.106248in}}%
\pgfpathclose%
\pgfusepath{fill}%
\end{pgfscope}%
\begin{pgfscope}%
\pgfpathrectangle{\pgfqpoint{0.407592in}{0.222904in}}{\pgfqpoint{1.411163in}{1.411163in}}%
\pgfusepath{clip}%
\pgfsetbuttcap%
\pgfsetroundjoin%
\definecolor{currentfill}{rgb}{0.168627,0.611765,0.411765}%
\pgfsetfillcolor{currentfill}%
\pgfsetfillopacity{0.180000}%
\pgfsetlinewidth{0.000000pt}%
\definecolor{currentstroke}{rgb}{0.000000,0.000000,0.000000}%
\pgfsetstrokecolor{currentstroke}%
\pgfsetdash{}{0pt}%
\pgfpathmoveto{\pgfqpoint{1.389978in}{0.913421in}}%
\pgfpathlineto{\pgfqpoint{1.405247in}{0.915836in}}%
\pgfpathlineto{\pgfqpoint{1.405591in}{0.938433in}}%
\pgfpathlineto{\pgfqpoint{1.390303in}{0.936096in}}%
\pgfpathlineto{\pgfqpoint{1.389978in}{0.913421in}}%
\pgfpathclose%
\pgfusepath{fill}%
\end{pgfscope}%
\begin{pgfscope}%
\pgfpathrectangle{\pgfqpoint{0.407592in}{0.222904in}}{\pgfqpoint{1.411163in}{1.411163in}}%
\pgfusepath{clip}%
\pgfsetbuttcap%
\pgfsetroundjoin%
\definecolor{currentfill}{rgb}{0.168627,0.611765,0.411765}%
\pgfsetfillcolor{currentfill}%
\pgfsetfillopacity{0.180000}%
\pgfsetlinewidth{0.000000pt}%
\definecolor{currentstroke}{rgb}{0.000000,0.000000,0.000000}%
\pgfsetstrokecolor{currentstroke}%
\pgfsetdash{}{0pt}%
\pgfpathmoveto{\pgfqpoint{1.406281in}{0.983799in}}%
\pgfpathlineto{\pgfqpoint{1.421519in}{0.987428in}}%
\pgfpathlineto{\pgfqpoint{1.421884in}{1.010122in}}%
\pgfpathlineto{\pgfqpoint{1.406628in}{1.006568in}}%
\pgfpathlineto{\pgfqpoint{1.406281in}{0.983799in}}%
\pgfpathclose%
\pgfusepath{fill}%
\end{pgfscope}%
\begin{pgfscope}%
\pgfpathrectangle{\pgfqpoint{0.407592in}{0.222904in}}{\pgfqpoint{1.411163in}{1.411163in}}%
\pgfusepath{clip}%
\pgfsetbuttcap%
\pgfsetroundjoin%
\definecolor{currentfill}{rgb}{0.168627,0.611765,0.411765}%
\pgfsetfillcolor{currentfill}%
\pgfsetfillopacity{0.180000}%
\pgfsetlinewidth{0.000000pt}%
\definecolor{currentstroke}{rgb}{0.000000,0.000000,0.000000}%
\pgfsetstrokecolor{currentstroke}%
\pgfsetdash{}{0pt}%
\pgfpathmoveto{\pgfqpoint{1.422616in}{1.055682in}}%
\pgfpathlineto{\pgfqpoint{1.437827in}{1.060600in}}%
\pgfpathlineto{\pgfqpoint{1.438213in}{1.083395in}}%
\pgfpathlineto{\pgfqpoint{1.422983in}{1.078549in}}%
\pgfpathlineto{\pgfqpoint{1.422616in}{1.055682in}}%
\pgfpathclose%
\pgfusepath{fill}%
\end{pgfscope}%
\begin{pgfscope}%
\pgfpathrectangle{\pgfqpoint{0.407592in}{0.222904in}}{\pgfqpoint{1.411163in}{1.411163in}}%
\pgfusepath{clip}%
\pgfsetbuttcap%
\pgfsetroundjoin%
\definecolor{currentfill}{rgb}{0.168627,0.611765,0.411765}%
\pgfsetfillcolor{currentfill}%
\pgfsetfillopacity{0.180000}%
\pgfsetlinewidth{0.000000pt}%
\definecolor{currentstroke}{rgb}{0.000000,0.000000,0.000000}%
\pgfsetstrokecolor{currentstroke}%
\pgfsetdash{}{0pt}%
\pgfpathmoveto{\pgfqpoint{1.438986in}{1.129158in}}%
\pgfpathlineto{\pgfqpoint{1.454176in}{1.135428in}}%
\pgfpathlineto{\pgfqpoint{1.454582in}{1.158328in}}%
\pgfpathlineto{\pgfqpoint{1.439374in}{1.152126in}}%
\pgfpathlineto{\pgfqpoint{1.438986in}{1.129158in}}%
\pgfpathclose%
\pgfusepath{fill}%
\end{pgfscope}%
\begin{pgfscope}%
\pgfpathrectangle{\pgfqpoint{0.407592in}{0.222904in}}{\pgfqpoint{1.411163in}{1.411163in}}%
\pgfusepath{clip}%
\pgfsetbuttcap%
\pgfsetroundjoin%
\definecolor{currentfill}{rgb}{0.168627,0.611765,0.411765}%
\pgfsetfillcolor{currentfill}%
\pgfsetfillopacity{0.180000}%
\pgfsetlinewidth{0.000000pt}%
\definecolor{currentstroke}{rgb}{0.000000,0.000000,0.000000}%
\pgfsetstrokecolor{currentstroke}%
\pgfsetdash{}{0pt}%
\pgfpathmoveto{\pgfqpoint{1.374005in}{0.867037in}}%
\pgfpathlineto{\pgfqpoint{1.389328in}{0.868242in}}%
\pgfpathlineto{\pgfqpoint{1.389653in}{0.890803in}}%
\pgfpathlineto{\pgfqpoint{1.374311in}{0.889679in}}%
\pgfpathlineto{\pgfqpoint{1.374005in}{0.867037in}}%
\pgfpathclose%
\pgfusepath{fill}%
\end{pgfscope}%
\begin{pgfscope}%
\pgfpathrectangle{\pgfqpoint{0.407592in}{0.222904in}}{\pgfqpoint{1.411163in}{1.411163in}}%
\pgfusepath{clip}%
\pgfsetbuttcap%
\pgfsetroundjoin%
\definecolor{currentfill}{rgb}{0.168627,0.611765,0.411765}%
\pgfsetfillcolor{currentfill}%
\pgfsetfillopacity{0.180000}%
\pgfsetlinewidth{0.000000pt}%
\definecolor{currentstroke}{rgb}{0.000000,0.000000,0.000000}%
\pgfsetstrokecolor{currentstroke}%
\pgfsetdash{}{0pt}%
\pgfpathmoveto{\pgfqpoint{1.390303in}{0.936096in}}%
\pgfpathlineto{\pgfqpoint{1.405591in}{0.938433in}}%
\pgfpathlineto{\pgfqpoint{1.405936in}{0.961088in}}%
\pgfpathlineto{\pgfqpoint{1.390630in}{0.958829in}}%
\pgfpathlineto{\pgfqpoint{1.390303in}{0.936096in}}%
\pgfpathclose%
\pgfusepath{fill}%
\end{pgfscope}%
\begin{pgfscope}%
\pgfpathrectangle{\pgfqpoint{0.407592in}{0.222904in}}{\pgfqpoint{1.411163in}{1.411163in}}%
\pgfusepath{clip}%
\pgfsetbuttcap%
\pgfsetroundjoin%
\definecolor{currentfill}{rgb}{0.168627,0.611765,0.411765}%
\pgfsetfillcolor{currentfill}%
\pgfsetfillopacity{0.180000}%
\pgfsetlinewidth{0.000000pt}%
\definecolor{currentstroke}{rgb}{0.000000,0.000000,0.000000}%
\pgfsetstrokecolor{currentstroke}%
\pgfsetdash{}{0pt}%
\pgfpathmoveto{\pgfqpoint{1.406628in}{1.006568in}}%
\pgfpathlineto{\pgfqpoint{1.421884in}{1.010122in}}%
\pgfpathlineto{\pgfqpoint{1.422249in}{1.032873in}}%
\pgfpathlineto{\pgfqpoint{1.406975in}{1.029395in}}%
\pgfpathlineto{\pgfqpoint{1.406628in}{1.006568in}}%
\pgfpathclose%
\pgfusepath{fill}%
\end{pgfscope}%
\begin{pgfscope}%
\pgfpathrectangle{\pgfqpoint{0.407592in}{0.222904in}}{\pgfqpoint{1.411163in}{1.411163in}}%
\pgfusepath{clip}%
\pgfsetbuttcap%
\pgfsetroundjoin%
\definecolor{currentfill}{rgb}{0.168627,0.611765,0.411765}%
\pgfsetfillcolor{currentfill}%
\pgfsetfillopacity{0.180000}%
\pgfsetlinewidth{0.000000pt}%
\definecolor{currentstroke}{rgb}{0.000000,0.000000,0.000000}%
\pgfsetstrokecolor{currentstroke}%
\pgfsetdash{}{0pt}%
\pgfpathmoveto{\pgfqpoint{1.422983in}{1.078549in}}%
\pgfpathlineto{\pgfqpoint{1.438213in}{1.083395in}}%
\pgfpathlineto{\pgfqpoint{1.438599in}{1.106248in}}%
\pgfpathlineto{\pgfqpoint{1.423351in}{1.101474in}}%
\pgfpathlineto{\pgfqpoint{1.422983in}{1.078549in}}%
\pgfpathclose%
\pgfusepath{fill}%
\end{pgfscope}%
\begin{pgfscope}%
\pgfpathrectangle{\pgfqpoint{0.407592in}{0.222904in}}{\pgfqpoint{1.411163in}{1.411163in}}%
\pgfusepath{clip}%
\pgfsetbuttcap%
\pgfsetroundjoin%
\definecolor{currentfill}{rgb}{0.168627,0.611765,0.411765}%
\pgfsetfillcolor{currentfill}%
\pgfsetfillopacity{0.180000}%
\pgfsetlinewidth{0.000000pt}%
\definecolor{currentstroke}{rgb}{0.000000,0.000000,0.000000}%
\pgfsetstrokecolor{currentstroke}%
\pgfsetdash{}{0pt}%
\pgfpathmoveto{\pgfqpoint{1.439374in}{1.152126in}}%
\pgfpathlineto{\pgfqpoint{1.454582in}{1.158328in}}%
\pgfpathlineto{\pgfqpoint{1.454989in}{1.181286in}}%
\pgfpathlineto{\pgfqpoint{1.439763in}{1.175153in}}%
\pgfpathlineto{\pgfqpoint{1.439374in}{1.152126in}}%
\pgfpathclose%
\pgfusepath{fill}%
\end{pgfscope}%
\begin{pgfscope}%
\pgfpathrectangle{\pgfqpoint{0.407592in}{0.222904in}}{\pgfqpoint{1.411163in}{1.411163in}}%
\pgfusepath{clip}%
\pgfsetbuttcap%
\pgfsetroundjoin%
\definecolor{currentfill}{rgb}{0.168627,0.611765,0.411765}%
\pgfsetfillcolor{currentfill}%
\pgfsetfillopacity{0.180000}%
\pgfsetlinewidth{0.000000pt}%
\definecolor{currentstroke}{rgb}{0.000000,0.000000,0.000000}%
\pgfsetstrokecolor{currentstroke}%
\pgfsetdash{}{0pt}%
\pgfpathmoveto{\pgfqpoint{1.374311in}{0.889679in}}%
\pgfpathlineto{\pgfqpoint{1.389653in}{0.890803in}}%
\pgfpathlineto{\pgfqpoint{1.389978in}{0.913421in}}%
\pgfpathlineto{\pgfqpoint{1.374618in}{0.912379in}}%
\pgfpathlineto{\pgfqpoint{1.374311in}{0.889679in}}%
\pgfpathclose%
\pgfusepath{fill}%
\end{pgfscope}%
\begin{pgfscope}%
\pgfpathrectangle{\pgfqpoint{0.407592in}{0.222904in}}{\pgfqpoint{1.411163in}{1.411163in}}%
\pgfusepath{clip}%
\pgfsetbuttcap%
\pgfsetroundjoin%
\definecolor{currentfill}{rgb}{0.168627,0.611765,0.411765}%
\pgfsetfillcolor{currentfill}%
\pgfsetfillopacity{0.180000}%
\pgfsetlinewidth{0.000000pt}%
\definecolor{currentstroke}{rgb}{0.000000,0.000000,0.000000}%
\pgfsetstrokecolor{currentstroke}%
\pgfsetdash{}{0pt}%
\pgfpathmoveto{\pgfqpoint{1.390630in}{0.958829in}}%
\pgfpathlineto{\pgfqpoint{1.405936in}{0.961088in}}%
\pgfpathlineto{\pgfqpoint{1.406281in}{0.983799in}}%
\pgfpathlineto{\pgfqpoint{1.390958in}{0.981619in}}%
\pgfpathlineto{\pgfqpoint{1.390630in}{0.958829in}}%
\pgfpathclose%
\pgfusepath{fill}%
\end{pgfscope}%
\begin{pgfscope}%
\pgfpathrectangle{\pgfqpoint{0.407592in}{0.222904in}}{\pgfqpoint{1.411163in}{1.411163in}}%
\pgfusepath{clip}%
\pgfsetbuttcap%
\pgfsetroundjoin%
\definecolor{currentfill}{rgb}{0.168627,0.611765,0.411765}%
\pgfsetfillcolor{currentfill}%
\pgfsetfillopacity{0.180000}%
\pgfsetlinewidth{0.000000pt}%
\definecolor{currentstroke}{rgb}{0.000000,0.000000,0.000000}%
\pgfsetstrokecolor{currentstroke}%
\pgfsetdash{}{0pt}%
\pgfpathmoveto{\pgfqpoint{1.406975in}{1.029395in}}%
\pgfpathlineto{\pgfqpoint{1.422249in}{1.032873in}}%
\pgfpathlineto{\pgfqpoint{1.422616in}{1.055682in}}%
\pgfpathlineto{\pgfqpoint{1.407324in}{1.052280in}}%
\pgfpathlineto{\pgfqpoint{1.406975in}{1.029395in}}%
\pgfpathclose%
\pgfusepath{fill}%
\end{pgfscope}%
\begin{pgfscope}%
\pgfpathrectangle{\pgfqpoint{0.407592in}{0.222904in}}{\pgfqpoint{1.411163in}{1.411163in}}%
\pgfusepath{clip}%
\pgfsetbuttcap%
\pgfsetroundjoin%
\definecolor{currentfill}{rgb}{0.168627,0.611765,0.411765}%
\pgfsetfillcolor{currentfill}%
\pgfsetfillopacity{0.180000}%
\pgfsetlinewidth{0.000000pt}%
\definecolor{currentstroke}{rgb}{0.000000,0.000000,0.000000}%
\pgfsetstrokecolor{currentstroke}%
\pgfsetdash{}{0pt}%
\pgfpathmoveto{\pgfqpoint{1.423351in}{1.101474in}}%
\pgfpathlineto{\pgfqpoint{1.438599in}{1.106248in}}%
\pgfpathlineto{\pgfqpoint{1.438986in}{1.129158in}}%
\pgfpathlineto{\pgfqpoint{1.423720in}{1.124457in}}%
\pgfpathlineto{\pgfqpoint{1.423351in}{1.101474in}}%
\pgfpathclose%
\pgfusepath{fill}%
\end{pgfscope}%
\begin{pgfscope}%
\pgfpathrectangle{\pgfqpoint{0.407592in}{0.222904in}}{\pgfqpoint{1.411163in}{1.411163in}}%
\pgfusepath{clip}%
\pgfsetbuttcap%
\pgfsetroundjoin%
\definecolor{currentfill}{rgb}{0.168627,0.611765,0.411765}%
\pgfsetfillcolor{currentfill}%
\pgfsetfillopacity{0.180000}%
\pgfsetlinewidth{0.000000pt}%
\definecolor{currentstroke}{rgb}{0.000000,0.000000,0.000000}%
\pgfsetstrokecolor{currentstroke}%
\pgfsetdash{}{0pt}%
\pgfpathmoveto{\pgfqpoint{1.374618in}{0.912379in}}%
\pgfpathlineto{\pgfqpoint{1.389978in}{0.913421in}}%
\pgfpathlineto{\pgfqpoint{1.390303in}{0.936096in}}%
\pgfpathlineto{\pgfqpoint{1.374925in}{0.935136in}}%
\pgfpathlineto{\pgfqpoint{1.374618in}{0.912379in}}%
\pgfpathclose%
\pgfusepath{fill}%
\end{pgfscope}%
\begin{pgfscope}%
\pgfpathrectangle{\pgfqpoint{0.407592in}{0.222904in}}{\pgfqpoint{1.411163in}{1.411163in}}%
\pgfusepath{clip}%
\pgfsetbuttcap%
\pgfsetroundjoin%
\definecolor{currentfill}{rgb}{0.168627,0.611765,0.411765}%
\pgfsetfillcolor{currentfill}%
\pgfsetfillopacity{0.180000}%
\pgfsetlinewidth{0.000000pt}%
\definecolor{currentstroke}{rgb}{0.000000,0.000000,0.000000}%
\pgfsetstrokecolor{currentstroke}%
\pgfsetdash{}{0pt}%
\pgfpathmoveto{\pgfqpoint{1.390958in}{0.981619in}}%
\pgfpathlineto{\pgfqpoint{1.406281in}{0.983799in}}%
\pgfpathlineto{\pgfqpoint{1.406628in}{1.006568in}}%
\pgfpathlineto{\pgfqpoint{1.391286in}{1.004467in}}%
\pgfpathlineto{\pgfqpoint{1.390958in}{0.981619in}}%
\pgfpathclose%
\pgfusepath{fill}%
\end{pgfscope}%
\begin{pgfscope}%
\pgfpathrectangle{\pgfqpoint{0.407592in}{0.222904in}}{\pgfqpoint{1.411163in}{1.411163in}}%
\pgfusepath{clip}%
\pgfsetbuttcap%
\pgfsetroundjoin%
\definecolor{currentfill}{rgb}{0.168627,0.611765,0.411765}%
\pgfsetfillcolor{currentfill}%
\pgfsetfillopacity{0.180000}%
\pgfsetlinewidth{0.000000pt}%
\definecolor{currentstroke}{rgb}{0.000000,0.000000,0.000000}%
\pgfsetstrokecolor{currentstroke}%
\pgfsetdash{}{0pt}%
\pgfpathmoveto{\pgfqpoint{1.407324in}{1.052280in}}%
\pgfpathlineto{\pgfqpoint{1.422616in}{1.055682in}}%
\pgfpathlineto{\pgfqpoint{1.422983in}{1.078549in}}%
\pgfpathlineto{\pgfqpoint{1.407673in}{1.075222in}}%
\pgfpathlineto{\pgfqpoint{1.407324in}{1.052280in}}%
\pgfpathclose%
\pgfusepath{fill}%
\end{pgfscope}%
\begin{pgfscope}%
\pgfpathrectangle{\pgfqpoint{0.407592in}{0.222904in}}{\pgfqpoint{1.411163in}{1.411163in}}%
\pgfusepath{clip}%
\pgfsetbuttcap%
\pgfsetroundjoin%
\definecolor{currentfill}{rgb}{0.168627,0.611765,0.411765}%
\pgfsetfillcolor{currentfill}%
\pgfsetfillopacity{0.180000}%
\pgfsetlinewidth{0.000000pt}%
\definecolor{currentstroke}{rgb}{0.000000,0.000000,0.000000}%
\pgfsetstrokecolor{currentstroke}%
\pgfsetdash{}{0pt}%
\pgfpathmoveto{\pgfqpoint{1.423720in}{1.124457in}}%
\pgfpathlineto{\pgfqpoint{1.438986in}{1.129158in}}%
\pgfpathlineto{\pgfqpoint{1.439374in}{1.152126in}}%
\pgfpathlineto{\pgfqpoint{1.424090in}{1.147498in}}%
\pgfpathlineto{\pgfqpoint{1.423720in}{1.124457in}}%
\pgfpathclose%
\pgfusepath{fill}%
\end{pgfscope}%
\begin{pgfscope}%
\pgfpathrectangle{\pgfqpoint{0.407592in}{0.222904in}}{\pgfqpoint{1.411163in}{1.411163in}}%
\pgfusepath{clip}%
\pgfsetbuttcap%
\pgfsetroundjoin%
\definecolor{currentfill}{rgb}{0.168627,0.611765,0.411765}%
\pgfsetfillcolor{currentfill}%
\pgfsetfillopacity{0.180000}%
\pgfsetlinewidth{0.000000pt}%
\definecolor{currentstroke}{rgb}{0.000000,0.000000,0.000000}%
\pgfsetstrokecolor{currentstroke}%
\pgfsetdash{}{0pt}%
\pgfpathmoveto{\pgfqpoint{1.358585in}{0.867127in}}%
\pgfpathlineto{\pgfqpoint{1.374005in}{0.867037in}}%
\pgfpathlineto{\pgfqpoint{1.374311in}{0.889679in}}%
\pgfpathlineto{\pgfqpoint{1.358873in}{0.889854in}}%
\pgfpathlineto{\pgfqpoint{1.358585in}{0.867127in}}%
\pgfpathclose%
\pgfusepath{fill}%
\end{pgfscope}%
\begin{pgfscope}%
\pgfpathrectangle{\pgfqpoint{0.407592in}{0.222904in}}{\pgfqpoint{1.411163in}{1.411163in}}%
\pgfusepath{clip}%
\pgfsetbuttcap%
\pgfsetroundjoin%
\definecolor{currentfill}{rgb}{0.168627,0.611765,0.411765}%
\pgfsetfillcolor{currentfill}%
\pgfsetfillopacity{0.180000}%
\pgfsetlinewidth{0.000000pt}%
\definecolor{currentstroke}{rgb}{0.000000,0.000000,0.000000}%
\pgfsetstrokecolor{currentstroke}%
\pgfsetdash{}{0pt}%
\pgfpathmoveto{\pgfqpoint{1.374925in}{0.935136in}}%
\pgfpathlineto{\pgfqpoint{1.390303in}{0.936096in}}%
\pgfpathlineto{\pgfqpoint{1.390630in}{0.958829in}}%
\pgfpathlineto{\pgfqpoint{1.375233in}{0.957951in}}%
\pgfpathlineto{\pgfqpoint{1.374925in}{0.935136in}}%
\pgfpathclose%
\pgfusepath{fill}%
\end{pgfscope}%
\begin{pgfscope}%
\pgfpathrectangle{\pgfqpoint{0.407592in}{0.222904in}}{\pgfqpoint{1.411163in}{1.411163in}}%
\pgfusepath{clip}%
\pgfsetbuttcap%
\pgfsetroundjoin%
\definecolor{currentfill}{rgb}{0.168627,0.611765,0.411765}%
\pgfsetfillcolor{currentfill}%
\pgfsetfillopacity{0.180000}%
\pgfsetlinewidth{0.000000pt}%
\definecolor{currentstroke}{rgb}{0.000000,0.000000,0.000000}%
\pgfsetstrokecolor{currentstroke}%
\pgfsetdash{}{0pt}%
\pgfpathmoveto{\pgfqpoint{1.391286in}{1.004467in}}%
\pgfpathlineto{\pgfqpoint{1.406628in}{1.006568in}}%
\pgfpathlineto{\pgfqpoint{1.406975in}{1.029395in}}%
\pgfpathlineto{\pgfqpoint{1.391615in}{1.027374in}}%
\pgfpathlineto{\pgfqpoint{1.391286in}{1.004467in}}%
\pgfpathclose%
\pgfusepath{fill}%
\end{pgfscope}%
\begin{pgfscope}%
\pgfpathrectangle{\pgfqpoint{0.407592in}{0.222904in}}{\pgfqpoint{1.411163in}{1.411163in}}%
\pgfusepath{clip}%
\pgfsetbuttcap%
\pgfsetroundjoin%
\definecolor{currentfill}{rgb}{0.168627,0.611765,0.411765}%
\pgfsetfillcolor{currentfill}%
\pgfsetfillopacity{0.180000}%
\pgfsetlinewidth{0.000000pt}%
\definecolor{currentstroke}{rgb}{0.000000,0.000000,0.000000}%
\pgfsetstrokecolor{currentstroke}%
\pgfsetdash{}{0pt}%
\pgfpathmoveto{\pgfqpoint{1.407673in}{1.075222in}}%
\pgfpathlineto{\pgfqpoint{1.422983in}{1.078549in}}%
\pgfpathlineto{\pgfqpoint{1.423351in}{1.101474in}}%
\pgfpathlineto{\pgfqpoint{1.408023in}{1.098223in}}%
\pgfpathlineto{\pgfqpoint{1.407673in}{1.075222in}}%
\pgfpathclose%
\pgfusepath{fill}%
\end{pgfscope}%
\begin{pgfscope}%
\pgfpathrectangle{\pgfqpoint{0.407592in}{0.222904in}}{\pgfqpoint{1.411163in}{1.411163in}}%
\pgfusepath{clip}%
\pgfsetbuttcap%
\pgfsetroundjoin%
\definecolor{currentfill}{rgb}{0.168627,0.611765,0.411765}%
\pgfsetfillcolor{currentfill}%
\pgfsetfillopacity{0.180000}%
\pgfsetlinewidth{0.000000pt}%
\definecolor{currentstroke}{rgb}{0.000000,0.000000,0.000000}%
\pgfsetstrokecolor{currentstroke}%
\pgfsetdash{}{0pt}%
\pgfpathmoveto{\pgfqpoint{1.424090in}{1.147498in}}%
\pgfpathlineto{\pgfqpoint{1.439374in}{1.152126in}}%
\pgfpathlineto{\pgfqpoint{1.439763in}{1.175153in}}%
\pgfpathlineto{\pgfqpoint{1.424461in}{1.170598in}}%
\pgfpathlineto{\pgfqpoint{1.424090in}{1.147498in}}%
\pgfpathclose%
\pgfusepath{fill}%
\end{pgfscope}%
\begin{pgfscope}%
\pgfpathrectangle{\pgfqpoint{0.407592in}{0.222904in}}{\pgfqpoint{1.411163in}{1.411163in}}%
\pgfusepath{clip}%
\pgfsetbuttcap%
\pgfsetroundjoin%
\definecolor{currentfill}{rgb}{0.168627,0.611765,0.411765}%
\pgfsetfillcolor{currentfill}%
\pgfsetfillopacity{0.180000}%
\pgfsetlinewidth{0.000000pt}%
\definecolor{currentstroke}{rgb}{0.000000,0.000000,0.000000}%
\pgfsetstrokecolor{currentstroke}%
\pgfsetdash{}{0pt}%
\pgfpathmoveto{\pgfqpoint{1.358873in}{0.889854in}}%
\pgfpathlineto{\pgfqpoint{1.374311in}{0.889679in}}%
\pgfpathlineto{\pgfqpoint{1.374618in}{0.912379in}}%
\pgfpathlineto{\pgfqpoint{1.359161in}{0.912639in}}%
\pgfpathlineto{\pgfqpoint{1.358873in}{0.889854in}}%
\pgfpathclose%
\pgfusepath{fill}%
\end{pgfscope}%
\begin{pgfscope}%
\pgfpathrectangle{\pgfqpoint{0.407592in}{0.222904in}}{\pgfqpoint{1.411163in}{1.411163in}}%
\pgfusepath{clip}%
\pgfsetbuttcap%
\pgfsetroundjoin%
\definecolor{currentfill}{rgb}{0.168627,0.611765,0.411765}%
\pgfsetfillcolor{currentfill}%
\pgfsetfillopacity{0.180000}%
\pgfsetlinewidth{0.000000pt}%
\definecolor{currentstroke}{rgb}{0.000000,0.000000,0.000000}%
\pgfsetstrokecolor{currentstroke}%
\pgfsetdash{}{0pt}%
\pgfpathmoveto{\pgfqpoint{1.375233in}{0.957951in}}%
\pgfpathlineto{\pgfqpoint{1.390630in}{0.958829in}}%
\pgfpathlineto{\pgfqpoint{1.390958in}{0.981619in}}%
\pgfpathlineto{\pgfqpoint{1.375542in}{0.980824in}}%
\pgfpathlineto{\pgfqpoint{1.375233in}{0.957951in}}%
\pgfpathclose%
\pgfusepath{fill}%
\end{pgfscope}%
\begin{pgfscope}%
\pgfpathrectangle{\pgfqpoint{0.407592in}{0.222904in}}{\pgfqpoint{1.411163in}{1.411163in}}%
\pgfusepath{clip}%
\pgfsetbuttcap%
\pgfsetroundjoin%
\definecolor{currentfill}{rgb}{0.168627,0.611765,0.411765}%
\pgfsetfillcolor{currentfill}%
\pgfsetfillopacity{0.180000}%
\pgfsetlinewidth{0.000000pt}%
\definecolor{currentstroke}{rgb}{0.000000,0.000000,0.000000}%
\pgfsetstrokecolor{currentstroke}%
\pgfsetdash{}{0pt}%
\pgfpathmoveto{\pgfqpoint{1.391615in}{1.027374in}}%
\pgfpathlineto{\pgfqpoint{1.406975in}{1.029395in}}%
\pgfpathlineto{\pgfqpoint{1.407324in}{1.052280in}}%
\pgfpathlineto{\pgfqpoint{1.391945in}{1.050338in}}%
\pgfpathlineto{\pgfqpoint{1.391615in}{1.027374in}}%
\pgfpathclose%
\pgfusepath{fill}%
\end{pgfscope}%
\begin{pgfscope}%
\pgfpathrectangle{\pgfqpoint{0.407592in}{0.222904in}}{\pgfqpoint{1.411163in}{1.411163in}}%
\pgfusepath{clip}%
\pgfsetbuttcap%
\pgfsetroundjoin%
\definecolor{currentfill}{rgb}{0.168627,0.611765,0.411765}%
\pgfsetfillcolor{currentfill}%
\pgfsetfillopacity{0.180000}%
\pgfsetlinewidth{0.000000pt}%
\definecolor{currentstroke}{rgb}{0.000000,0.000000,0.000000}%
\pgfsetstrokecolor{currentstroke}%
\pgfsetdash{}{0pt}%
\pgfpathmoveto{\pgfqpoint{1.408023in}{1.098223in}}%
\pgfpathlineto{\pgfqpoint{1.423351in}{1.101474in}}%
\pgfpathlineto{\pgfqpoint{1.423720in}{1.124457in}}%
\pgfpathlineto{\pgfqpoint{1.408374in}{1.121283in}}%
\pgfpathlineto{\pgfqpoint{1.408023in}{1.098223in}}%
\pgfpathclose%
\pgfusepath{fill}%
\end{pgfscope}%
\begin{pgfscope}%
\pgfpathrectangle{\pgfqpoint{0.407592in}{0.222904in}}{\pgfqpoint{1.411163in}{1.411163in}}%
\pgfusepath{clip}%
\pgfsetbuttcap%
\pgfsetroundjoin%
\definecolor{currentfill}{rgb}{0.168627,0.611765,0.411765}%
\pgfsetfillcolor{currentfill}%
\pgfsetfillopacity{0.180000}%
\pgfsetlinewidth{0.000000pt}%
\definecolor{currentstroke}{rgb}{0.000000,0.000000,0.000000}%
\pgfsetstrokecolor{currentstroke}%
\pgfsetdash{}{0pt}%
\pgfpathmoveto{\pgfqpoint{1.359161in}{0.912639in}}%
\pgfpathlineto{\pgfqpoint{1.374618in}{0.912379in}}%
\pgfpathlineto{\pgfqpoint{1.374925in}{0.935136in}}%
\pgfpathlineto{\pgfqpoint{1.359450in}{0.935482in}}%
\pgfpathlineto{\pgfqpoint{1.359161in}{0.912639in}}%
\pgfpathclose%
\pgfusepath{fill}%
\end{pgfscope}%
\begin{pgfscope}%
\pgfpathrectangle{\pgfqpoint{0.407592in}{0.222904in}}{\pgfqpoint{1.411163in}{1.411163in}}%
\pgfusepath{clip}%
\pgfsetbuttcap%
\pgfsetroundjoin%
\definecolor{currentfill}{rgb}{0.168627,0.611765,0.411765}%
\pgfsetfillcolor{currentfill}%
\pgfsetfillopacity{0.180000}%
\pgfsetlinewidth{0.000000pt}%
\definecolor{currentstroke}{rgb}{0.000000,0.000000,0.000000}%
\pgfsetstrokecolor{currentstroke}%
\pgfsetdash{}{0pt}%
\pgfpathmoveto{\pgfqpoint{1.375542in}{0.980824in}}%
\pgfpathlineto{\pgfqpoint{1.390958in}{0.981619in}}%
\pgfpathlineto{\pgfqpoint{1.391286in}{1.004467in}}%
\pgfpathlineto{\pgfqpoint{1.375852in}{1.003755in}}%
\pgfpathlineto{\pgfqpoint{1.375542in}{0.980824in}}%
\pgfpathclose%
\pgfusepath{fill}%
\end{pgfscope}%
\begin{pgfscope}%
\pgfpathrectangle{\pgfqpoint{0.407592in}{0.222904in}}{\pgfqpoint{1.411163in}{1.411163in}}%
\pgfusepath{clip}%
\pgfsetbuttcap%
\pgfsetroundjoin%
\definecolor{currentfill}{rgb}{0.168627,0.611765,0.411765}%
\pgfsetfillcolor{currentfill}%
\pgfsetfillopacity{0.180000}%
\pgfsetlinewidth{0.000000pt}%
\definecolor{currentstroke}{rgb}{0.000000,0.000000,0.000000}%
\pgfsetstrokecolor{currentstroke}%
\pgfsetdash{}{0pt}%
\pgfpathmoveto{\pgfqpoint{1.391945in}{1.050338in}}%
\pgfpathlineto{\pgfqpoint{1.407324in}{1.052280in}}%
\pgfpathlineto{\pgfqpoint{1.407673in}{1.075222in}}%
\pgfpathlineto{\pgfqpoint{1.392276in}{1.073361in}}%
\pgfpathlineto{\pgfqpoint{1.391945in}{1.050338in}}%
\pgfpathclose%
\pgfusepath{fill}%
\end{pgfscope}%
\begin{pgfscope}%
\pgfpathrectangle{\pgfqpoint{0.407592in}{0.222904in}}{\pgfqpoint{1.411163in}{1.411163in}}%
\pgfusepath{clip}%
\pgfsetbuttcap%
\pgfsetroundjoin%
\definecolor{currentfill}{rgb}{0.168627,0.611765,0.411765}%
\pgfsetfillcolor{currentfill}%
\pgfsetfillopacity{0.180000}%
\pgfsetlinewidth{0.000000pt}%
\definecolor{currentstroke}{rgb}{0.000000,0.000000,0.000000}%
\pgfsetstrokecolor{currentstroke}%
\pgfsetdash{}{0pt}%
\pgfpathmoveto{\pgfqpoint{1.408374in}{1.121283in}}%
\pgfpathlineto{\pgfqpoint{1.423720in}{1.124457in}}%
\pgfpathlineto{\pgfqpoint{1.424090in}{1.147498in}}%
\pgfpathlineto{\pgfqpoint{1.408726in}{1.144401in}}%
\pgfpathlineto{\pgfqpoint{1.408374in}{1.121283in}}%
\pgfpathclose%
\pgfusepath{fill}%
\end{pgfscope}%
\begin{pgfscope}%
\pgfpathrectangle{\pgfqpoint{0.407592in}{0.222904in}}{\pgfqpoint{1.411163in}{1.411163in}}%
\pgfusepath{clip}%
\pgfsetbuttcap%
\pgfsetroundjoin%
\definecolor{currentfill}{rgb}{0.168627,0.611765,0.411765}%
\pgfsetfillcolor{currentfill}%
\pgfsetfillopacity{0.180000}%
\pgfsetlinewidth{0.000000pt}%
\definecolor{currentstroke}{rgb}{0.000000,0.000000,0.000000}%
\pgfsetstrokecolor{currentstroke}%
\pgfsetdash{}{0pt}%
\pgfpathmoveto{\pgfqpoint{1.343064in}{0.868436in}}%
\pgfpathlineto{\pgfqpoint{1.358585in}{0.867127in}}%
\pgfpathlineto{\pgfqpoint{1.358873in}{0.889854in}}%
\pgfpathlineto{\pgfqpoint{1.343333in}{0.891252in}}%
\pgfpathlineto{\pgfqpoint{1.343064in}{0.868436in}}%
\pgfpathclose%
\pgfusepath{fill}%
\end{pgfscope}%
\begin{pgfscope}%
\pgfpathrectangle{\pgfqpoint{0.407592in}{0.222904in}}{\pgfqpoint{1.411163in}{1.411163in}}%
\pgfusepath{clip}%
\pgfsetbuttcap%
\pgfsetroundjoin%
\definecolor{currentfill}{rgb}{0.168627,0.611765,0.411765}%
\pgfsetfillcolor{currentfill}%
\pgfsetfillopacity{0.180000}%
\pgfsetlinewidth{0.000000pt}%
\definecolor{currentstroke}{rgb}{0.000000,0.000000,0.000000}%
\pgfsetstrokecolor{currentstroke}%
\pgfsetdash{}{0pt}%
\pgfpathmoveto{\pgfqpoint{1.359450in}{0.935482in}}%
\pgfpathlineto{\pgfqpoint{1.374925in}{0.935136in}}%
\pgfpathlineto{\pgfqpoint{1.375233in}{0.957951in}}%
\pgfpathlineto{\pgfqpoint{1.359740in}{0.958383in}}%
\pgfpathlineto{\pgfqpoint{1.359450in}{0.935482in}}%
\pgfpathclose%
\pgfusepath{fill}%
\end{pgfscope}%
\begin{pgfscope}%
\pgfpathrectangle{\pgfqpoint{0.407592in}{0.222904in}}{\pgfqpoint{1.411163in}{1.411163in}}%
\pgfusepath{clip}%
\pgfsetbuttcap%
\pgfsetroundjoin%
\definecolor{currentfill}{rgb}{0.168627,0.611765,0.411765}%
\pgfsetfillcolor{currentfill}%
\pgfsetfillopacity{0.180000}%
\pgfsetlinewidth{0.000000pt}%
\definecolor{currentstroke}{rgb}{0.000000,0.000000,0.000000}%
\pgfsetstrokecolor{currentstroke}%
\pgfsetdash{}{0pt}%
\pgfpathmoveto{\pgfqpoint{1.375852in}{1.003755in}}%
\pgfpathlineto{\pgfqpoint{1.391286in}{1.004467in}}%
\pgfpathlineto{\pgfqpoint{1.391615in}{1.027374in}}%
\pgfpathlineto{\pgfqpoint{1.376163in}{1.026745in}}%
\pgfpathlineto{\pgfqpoint{1.375852in}{1.003755in}}%
\pgfpathclose%
\pgfusepath{fill}%
\end{pgfscope}%
\begin{pgfscope}%
\pgfpathrectangle{\pgfqpoint{0.407592in}{0.222904in}}{\pgfqpoint{1.411163in}{1.411163in}}%
\pgfusepath{clip}%
\pgfsetbuttcap%
\pgfsetroundjoin%
\definecolor{currentfill}{rgb}{0.168627,0.611765,0.411765}%
\pgfsetfillcolor{currentfill}%
\pgfsetfillopacity{0.180000}%
\pgfsetlinewidth{0.000000pt}%
\definecolor{currentstroke}{rgb}{0.000000,0.000000,0.000000}%
\pgfsetstrokecolor{currentstroke}%
\pgfsetdash{}{0pt}%
\pgfpathmoveto{\pgfqpoint{1.392276in}{1.073361in}}%
\pgfpathlineto{\pgfqpoint{1.407673in}{1.075222in}}%
\pgfpathlineto{\pgfqpoint{1.408023in}{1.098223in}}%
\pgfpathlineto{\pgfqpoint{1.392608in}{1.096443in}}%
\pgfpathlineto{\pgfqpoint{1.392276in}{1.073361in}}%
\pgfpathclose%
\pgfusepath{fill}%
\end{pgfscope}%
\begin{pgfscope}%
\pgfpathrectangle{\pgfqpoint{0.407592in}{0.222904in}}{\pgfqpoint{1.411163in}{1.411163in}}%
\pgfusepath{clip}%
\pgfsetbuttcap%
\pgfsetroundjoin%
\definecolor{currentfill}{rgb}{0.168627,0.611765,0.411765}%
\pgfsetfillcolor{currentfill}%
\pgfsetfillopacity{0.180000}%
\pgfsetlinewidth{0.000000pt}%
\definecolor{currentstroke}{rgb}{0.000000,0.000000,0.000000}%
\pgfsetstrokecolor{currentstroke}%
\pgfsetdash{}{0pt}%
\pgfpathmoveto{\pgfqpoint{1.408726in}{1.144401in}}%
\pgfpathlineto{\pgfqpoint{1.424090in}{1.147498in}}%
\pgfpathlineto{\pgfqpoint{1.424461in}{1.170598in}}%
\pgfpathlineto{\pgfqpoint{1.409078in}{1.167579in}}%
\pgfpathlineto{\pgfqpoint{1.408726in}{1.144401in}}%
\pgfpathclose%
\pgfusepath{fill}%
\end{pgfscope}%
\begin{pgfscope}%
\pgfpathrectangle{\pgfqpoint{0.407592in}{0.222904in}}{\pgfqpoint{1.411163in}{1.411163in}}%
\pgfusepath{clip}%
\pgfsetbuttcap%
\pgfsetroundjoin%
\definecolor{currentfill}{rgb}{0.168627,0.611765,0.411765}%
\pgfsetfillcolor{currentfill}%
\pgfsetfillopacity{0.180000}%
\pgfsetlinewidth{0.000000pt}%
\definecolor{currentstroke}{rgb}{0.000000,0.000000,0.000000}%
\pgfsetstrokecolor{currentstroke}%
\pgfsetdash{}{0pt}%
\pgfpathmoveto{\pgfqpoint{1.343333in}{0.891252in}}%
\pgfpathlineto{\pgfqpoint{1.358873in}{0.889854in}}%
\pgfpathlineto{\pgfqpoint{1.359161in}{0.912639in}}%
\pgfpathlineto{\pgfqpoint{1.343602in}{0.914126in}}%
\pgfpathlineto{\pgfqpoint{1.343333in}{0.891252in}}%
\pgfpathclose%
\pgfusepath{fill}%
\end{pgfscope}%
\begin{pgfscope}%
\pgfpathrectangle{\pgfqpoint{0.407592in}{0.222904in}}{\pgfqpoint{1.411163in}{1.411163in}}%
\pgfusepath{clip}%
\pgfsetbuttcap%
\pgfsetroundjoin%
\definecolor{currentfill}{rgb}{0.168627,0.611765,0.411765}%
\pgfsetfillcolor{currentfill}%
\pgfsetfillopacity{0.180000}%
\pgfsetlinewidth{0.000000pt}%
\definecolor{currentstroke}{rgb}{0.000000,0.000000,0.000000}%
\pgfsetstrokecolor{currentstroke}%
\pgfsetdash{}{0pt}%
\pgfpathmoveto{\pgfqpoint{1.359740in}{0.958383in}}%
\pgfpathlineto{\pgfqpoint{1.375233in}{0.957951in}}%
\pgfpathlineto{\pgfqpoint{1.375542in}{0.980824in}}%
\pgfpathlineto{\pgfqpoint{1.360030in}{0.981342in}}%
\pgfpathlineto{\pgfqpoint{1.359740in}{0.958383in}}%
\pgfpathclose%
\pgfusepath{fill}%
\end{pgfscope}%
\begin{pgfscope}%
\pgfpathrectangle{\pgfqpoint{0.407592in}{0.222904in}}{\pgfqpoint{1.411163in}{1.411163in}}%
\pgfusepath{clip}%
\pgfsetbuttcap%
\pgfsetroundjoin%
\definecolor{currentfill}{rgb}{0.168627,0.611765,0.411765}%
\pgfsetfillcolor{currentfill}%
\pgfsetfillopacity{0.180000}%
\pgfsetlinewidth{0.000000pt}%
\definecolor{currentstroke}{rgb}{0.000000,0.000000,0.000000}%
\pgfsetstrokecolor{currentstroke}%
\pgfsetdash{}{0pt}%
\pgfpathmoveto{\pgfqpoint{1.376163in}{1.026745in}}%
\pgfpathlineto{\pgfqpoint{1.391615in}{1.027374in}}%
\pgfpathlineto{\pgfqpoint{1.391945in}{1.050338in}}%
\pgfpathlineto{\pgfqpoint{1.376475in}{1.049793in}}%
\pgfpathlineto{\pgfqpoint{1.376163in}{1.026745in}}%
\pgfpathclose%
\pgfusepath{fill}%
\end{pgfscope}%
\begin{pgfscope}%
\pgfpathrectangle{\pgfqpoint{0.407592in}{0.222904in}}{\pgfqpoint{1.411163in}{1.411163in}}%
\pgfusepath{clip}%
\pgfsetbuttcap%
\pgfsetroundjoin%
\definecolor{currentfill}{rgb}{0.168627,0.611765,0.411765}%
\pgfsetfillcolor{currentfill}%
\pgfsetfillopacity{0.180000}%
\pgfsetlinewidth{0.000000pt}%
\definecolor{currentstroke}{rgb}{0.000000,0.000000,0.000000}%
\pgfsetstrokecolor{currentstroke}%
\pgfsetdash{}{0pt}%
\pgfpathmoveto{\pgfqpoint{1.392608in}{1.096443in}}%
\pgfpathlineto{\pgfqpoint{1.408023in}{1.098223in}}%
\pgfpathlineto{\pgfqpoint{1.408374in}{1.121283in}}%
\pgfpathlineto{\pgfqpoint{1.392940in}{1.119583in}}%
\pgfpathlineto{\pgfqpoint{1.392608in}{1.096443in}}%
\pgfpathclose%
\pgfusepath{fill}%
\end{pgfscope}%
\begin{pgfscope}%
\pgfpathrectangle{\pgfqpoint{0.407592in}{0.222904in}}{\pgfqpoint{1.411163in}{1.411163in}}%
\pgfusepath{clip}%
\pgfsetbuttcap%
\pgfsetroundjoin%
\definecolor{currentfill}{rgb}{0.168627,0.611765,0.411765}%
\pgfsetfillcolor{currentfill}%
\pgfsetfillopacity{0.180000}%
\pgfsetlinewidth{0.000000pt}%
\definecolor{currentstroke}{rgb}{0.000000,0.000000,0.000000}%
\pgfsetstrokecolor{currentstroke}%
\pgfsetdash{}{0pt}%
\pgfpathmoveto{\pgfqpoint{1.343602in}{0.914126in}}%
\pgfpathlineto{\pgfqpoint{1.359161in}{0.912639in}}%
\pgfpathlineto{\pgfqpoint{1.359450in}{0.935482in}}%
\pgfpathlineto{\pgfqpoint{1.343873in}{0.937058in}}%
\pgfpathlineto{\pgfqpoint{1.343602in}{0.914126in}}%
\pgfpathclose%
\pgfusepath{fill}%
\end{pgfscope}%
\begin{pgfscope}%
\pgfpathrectangle{\pgfqpoint{0.407592in}{0.222904in}}{\pgfqpoint{1.411163in}{1.411163in}}%
\pgfusepath{clip}%
\pgfsetbuttcap%
\pgfsetroundjoin%
\definecolor{currentfill}{rgb}{0.168627,0.611765,0.411765}%
\pgfsetfillcolor{currentfill}%
\pgfsetfillopacity{0.180000}%
\pgfsetlinewidth{0.000000pt}%
\definecolor{currentstroke}{rgb}{0.000000,0.000000,0.000000}%
\pgfsetstrokecolor{currentstroke}%
\pgfsetdash{}{0pt}%
\pgfpathmoveto{\pgfqpoint{1.360030in}{0.981342in}}%
\pgfpathlineto{\pgfqpoint{1.375542in}{0.980824in}}%
\pgfpathlineto{\pgfqpoint{1.375852in}{1.003755in}}%
\pgfpathlineto{\pgfqpoint{1.360321in}{1.004361in}}%
\pgfpathlineto{\pgfqpoint{1.360030in}{0.981342in}}%
\pgfpathclose%
\pgfusepath{fill}%
\end{pgfscope}%
\begin{pgfscope}%
\pgfpathrectangle{\pgfqpoint{0.407592in}{0.222904in}}{\pgfqpoint{1.411163in}{1.411163in}}%
\pgfusepath{clip}%
\pgfsetbuttcap%
\pgfsetroundjoin%
\definecolor{currentfill}{rgb}{0.168627,0.611765,0.411765}%
\pgfsetfillcolor{currentfill}%
\pgfsetfillopacity{0.180000}%
\pgfsetlinewidth{0.000000pt}%
\definecolor{currentstroke}{rgb}{0.000000,0.000000,0.000000}%
\pgfsetstrokecolor{currentstroke}%
\pgfsetdash{}{0pt}%
\pgfpathmoveto{\pgfqpoint{1.376475in}{1.049793in}}%
\pgfpathlineto{\pgfqpoint{1.391945in}{1.050338in}}%
\pgfpathlineto{\pgfqpoint{1.392276in}{1.073361in}}%
\pgfpathlineto{\pgfqpoint{1.376787in}{1.072901in}}%
\pgfpathlineto{\pgfqpoint{1.376475in}{1.049793in}}%
\pgfpathclose%
\pgfusepath{fill}%
\end{pgfscope}%
\begin{pgfscope}%
\pgfpathrectangle{\pgfqpoint{0.407592in}{0.222904in}}{\pgfqpoint{1.411163in}{1.411163in}}%
\pgfusepath{clip}%
\pgfsetbuttcap%
\pgfsetroundjoin%
\definecolor{currentfill}{rgb}{0.168627,0.611765,0.411765}%
\pgfsetfillcolor{currentfill}%
\pgfsetfillopacity{0.180000}%
\pgfsetlinewidth{0.000000pt}%
\definecolor{currentstroke}{rgb}{0.000000,0.000000,0.000000}%
\pgfsetstrokecolor{currentstroke}%
\pgfsetdash{}{0pt}%
\pgfpathmoveto{\pgfqpoint{1.392940in}{1.119583in}}%
\pgfpathlineto{\pgfqpoint{1.408374in}{1.121283in}}%
\pgfpathlineto{\pgfqpoint{1.408726in}{1.144401in}}%
\pgfpathlineto{\pgfqpoint{1.393273in}{1.142783in}}%
\pgfpathlineto{\pgfqpoint{1.392940in}{1.119583in}}%
\pgfpathclose%
\pgfusepath{fill}%
\end{pgfscope}%
\begin{pgfscope}%
\pgfpathrectangle{\pgfqpoint{0.407592in}{0.222904in}}{\pgfqpoint{1.411163in}{1.411163in}}%
\pgfusepath{clip}%
\pgfsetbuttcap%
\pgfsetroundjoin%
\definecolor{currentfill}{rgb}{0.168627,0.611765,0.411765}%
\pgfsetfillcolor{currentfill}%
\pgfsetfillopacity{0.180000}%
\pgfsetlinewidth{0.000000pt}%
\definecolor{currentstroke}{rgb}{0.000000,0.000000,0.000000}%
\pgfsetstrokecolor{currentstroke}%
\pgfsetdash{}{0pt}%
\pgfpathmoveto{\pgfqpoint{1.343873in}{0.937058in}}%
\pgfpathlineto{\pgfqpoint{1.359450in}{0.935482in}}%
\pgfpathlineto{\pgfqpoint{1.359740in}{0.958383in}}%
\pgfpathlineto{\pgfqpoint{1.344144in}{0.960050in}}%
\pgfpathlineto{\pgfqpoint{1.343873in}{0.937058in}}%
\pgfpathclose%
\pgfusepath{fill}%
\end{pgfscope}%
\begin{pgfscope}%
\pgfpathrectangle{\pgfqpoint{0.407592in}{0.222904in}}{\pgfqpoint{1.411163in}{1.411163in}}%
\pgfusepath{clip}%
\pgfsetbuttcap%
\pgfsetroundjoin%
\definecolor{currentfill}{rgb}{0.168627,0.611765,0.411765}%
\pgfsetfillcolor{currentfill}%
\pgfsetfillopacity{0.180000}%
\pgfsetlinewidth{0.000000pt}%
\definecolor{currentstroke}{rgb}{0.000000,0.000000,0.000000}%
\pgfsetstrokecolor{currentstroke}%
\pgfsetdash{}{0pt}%
\pgfpathmoveto{\pgfqpoint{1.360321in}{1.004361in}}%
\pgfpathlineto{\pgfqpoint{1.375852in}{1.003755in}}%
\pgfpathlineto{\pgfqpoint{1.376163in}{1.026745in}}%
\pgfpathlineto{\pgfqpoint{1.360613in}{1.027438in}}%
\pgfpathlineto{\pgfqpoint{1.360321in}{1.004361in}}%
\pgfpathclose%
\pgfusepath{fill}%
\end{pgfscope}%
\begin{pgfscope}%
\pgfpathrectangle{\pgfqpoint{0.407592in}{0.222904in}}{\pgfqpoint{1.411163in}{1.411163in}}%
\pgfusepath{clip}%
\pgfsetbuttcap%
\pgfsetroundjoin%
\definecolor{currentfill}{rgb}{0.168627,0.611765,0.411765}%
\pgfsetfillcolor{currentfill}%
\pgfsetfillopacity{0.180000}%
\pgfsetlinewidth{0.000000pt}%
\definecolor{currentstroke}{rgb}{0.000000,0.000000,0.000000}%
\pgfsetstrokecolor{currentstroke}%
\pgfsetdash{}{0pt}%
\pgfpathmoveto{\pgfqpoint{1.327436in}{0.870889in}}%
\pgfpathlineto{\pgfqpoint{1.343064in}{0.868436in}}%
\pgfpathlineto{\pgfqpoint{1.343333in}{0.891252in}}%
\pgfpathlineto{\pgfqpoint{1.327686in}{0.893797in}}%
\pgfpathlineto{\pgfqpoint{1.327436in}{0.870889in}}%
\pgfpathclose%
\pgfusepath{fill}%
\end{pgfscope}%
\begin{pgfscope}%
\pgfpathrectangle{\pgfqpoint{0.407592in}{0.222904in}}{\pgfqpoint{1.411163in}{1.411163in}}%
\pgfusepath{clip}%
\pgfsetbuttcap%
\pgfsetroundjoin%
\definecolor{currentfill}{rgb}{0.168627,0.611765,0.411765}%
\pgfsetfillcolor{currentfill}%
\pgfsetfillopacity{0.180000}%
\pgfsetlinewidth{0.000000pt}%
\definecolor{currentstroke}{rgb}{0.000000,0.000000,0.000000}%
\pgfsetstrokecolor{currentstroke}%
\pgfsetdash{}{0pt}%
\pgfpathmoveto{\pgfqpoint{1.376787in}{1.072901in}}%
\pgfpathlineto{\pgfqpoint{1.392276in}{1.073361in}}%
\pgfpathlineto{\pgfqpoint{1.392608in}{1.096443in}}%
\pgfpathlineto{\pgfqpoint{1.377100in}{1.096067in}}%
\pgfpathlineto{\pgfqpoint{1.376787in}{1.072901in}}%
\pgfpathclose%
\pgfusepath{fill}%
\end{pgfscope}%
\begin{pgfscope}%
\pgfpathrectangle{\pgfqpoint{0.407592in}{0.222904in}}{\pgfqpoint{1.411163in}{1.411163in}}%
\pgfusepath{clip}%
\pgfsetbuttcap%
\pgfsetroundjoin%
\definecolor{currentfill}{rgb}{0.168627,0.611765,0.411765}%
\pgfsetfillcolor{currentfill}%
\pgfsetfillopacity{0.180000}%
\pgfsetlinewidth{0.000000pt}%
\definecolor{currentstroke}{rgb}{0.000000,0.000000,0.000000}%
\pgfsetstrokecolor{currentstroke}%
\pgfsetdash{}{0pt}%
\pgfpathmoveto{\pgfqpoint{1.393273in}{1.142783in}}%
\pgfpathlineto{\pgfqpoint{1.408726in}{1.144401in}}%
\pgfpathlineto{\pgfqpoint{1.409078in}{1.167579in}}%
\pgfpathlineto{\pgfqpoint{1.393608in}{1.166043in}}%
\pgfpathlineto{\pgfqpoint{1.393273in}{1.142783in}}%
\pgfpathclose%
\pgfusepath{fill}%
\end{pgfscope}%
\begin{pgfscope}%
\pgfpathrectangle{\pgfqpoint{0.407592in}{0.222904in}}{\pgfqpoint{1.411163in}{1.411163in}}%
\pgfusepath{clip}%
\pgfsetbuttcap%
\pgfsetroundjoin%
\definecolor{currentfill}{rgb}{0.168627,0.611765,0.411765}%
\pgfsetfillcolor{currentfill}%
\pgfsetfillopacity{0.180000}%
\pgfsetlinewidth{0.000000pt}%
\definecolor{currentstroke}{rgb}{0.000000,0.000000,0.000000}%
\pgfsetstrokecolor{currentstroke}%
\pgfsetdash{}{0pt}%
\pgfpathmoveto{\pgfqpoint{1.344144in}{0.960050in}}%
\pgfpathlineto{\pgfqpoint{1.359740in}{0.958383in}}%
\pgfpathlineto{\pgfqpoint{1.360030in}{0.981342in}}%
\pgfpathlineto{\pgfqpoint{1.344415in}{0.983100in}}%
\pgfpathlineto{\pgfqpoint{1.344144in}{0.960050in}}%
\pgfpathclose%
\pgfusepath{fill}%
\end{pgfscope}%
\begin{pgfscope}%
\pgfpathrectangle{\pgfqpoint{0.407592in}{0.222904in}}{\pgfqpoint{1.411163in}{1.411163in}}%
\pgfusepath{clip}%
\pgfsetbuttcap%
\pgfsetroundjoin%
\definecolor{currentfill}{rgb}{0.168627,0.611765,0.411765}%
\pgfsetfillcolor{currentfill}%
\pgfsetfillopacity{0.180000}%
\pgfsetlinewidth{0.000000pt}%
\definecolor{currentstroke}{rgb}{0.000000,0.000000,0.000000}%
\pgfsetstrokecolor{currentstroke}%
\pgfsetdash{}{0pt}%
\pgfpathmoveto{\pgfqpoint{1.360613in}{1.027438in}}%
\pgfpathlineto{\pgfqpoint{1.376163in}{1.026745in}}%
\pgfpathlineto{\pgfqpoint{1.376475in}{1.049793in}}%
\pgfpathlineto{\pgfqpoint{1.360906in}{1.050575in}}%
\pgfpathlineto{\pgfqpoint{1.360613in}{1.027438in}}%
\pgfpathclose%
\pgfusepath{fill}%
\end{pgfscope}%
\begin{pgfscope}%
\pgfpathrectangle{\pgfqpoint{0.407592in}{0.222904in}}{\pgfqpoint{1.411163in}{1.411163in}}%
\pgfusepath{clip}%
\pgfsetbuttcap%
\pgfsetroundjoin%
\definecolor{currentfill}{rgb}{0.168627,0.611765,0.411765}%
\pgfsetfillcolor{currentfill}%
\pgfsetfillopacity{0.180000}%
\pgfsetlinewidth{0.000000pt}%
\definecolor{currentstroke}{rgb}{0.000000,0.000000,0.000000}%
\pgfsetstrokecolor{currentstroke}%
\pgfsetdash{}{0pt}%
\pgfpathmoveto{\pgfqpoint{1.327686in}{0.893797in}}%
\pgfpathlineto{\pgfqpoint{1.343333in}{0.891252in}}%
\pgfpathlineto{\pgfqpoint{1.343602in}{0.914126in}}%
\pgfpathlineto{\pgfqpoint{1.327936in}{0.916764in}}%
\pgfpathlineto{\pgfqpoint{1.327686in}{0.893797in}}%
\pgfpathclose%
\pgfusepath{fill}%
\end{pgfscope}%
\begin{pgfscope}%
\pgfpathrectangle{\pgfqpoint{0.407592in}{0.222904in}}{\pgfqpoint{1.411163in}{1.411163in}}%
\pgfusepath{clip}%
\pgfsetbuttcap%
\pgfsetroundjoin%
\definecolor{currentfill}{rgb}{0.168627,0.611765,0.411765}%
\pgfsetfillcolor{currentfill}%
\pgfsetfillopacity{0.180000}%
\pgfsetlinewidth{0.000000pt}%
\definecolor{currentstroke}{rgb}{0.000000,0.000000,0.000000}%
\pgfsetstrokecolor{currentstroke}%
\pgfsetdash{}{0pt}%
\pgfpathmoveto{\pgfqpoint{1.377100in}{1.096067in}}%
\pgfpathlineto{\pgfqpoint{1.392608in}{1.096443in}}%
\pgfpathlineto{\pgfqpoint{1.392940in}{1.119583in}}%
\pgfpathlineto{\pgfqpoint{1.377414in}{1.119293in}}%
\pgfpathlineto{\pgfqpoint{1.377100in}{1.096067in}}%
\pgfpathclose%
\pgfusepath{fill}%
\end{pgfscope}%
\begin{pgfscope}%
\pgfpathrectangle{\pgfqpoint{0.407592in}{0.222904in}}{\pgfqpoint{1.411163in}{1.411163in}}%
\pgfusepath{clip}%
\pgfsetbuttcap%
\pgfsetroundjoin%
\definecolor{currentfill}{rgb}{0.168627,0.611765,0.411765}%
\pgfsetfillcolor{currentfill}%
\pgfsetfillopacity{0.180000}%
\pgfsetlinewidth{0.000000pt}%
\definecolor{currentstroke}{rgb}{0.000000,0.000000,0.000000}%
\pgfsetstrokecolor{currentstroke}%
\pgfsetdash{}{0pt}%
\pgfpathmoveto{\pgfqpoint{1.344415in}{0.983100in}}%
\pgfpathlineto{\pgfqpoint{1.360030in}{0.981342in}}%
\pgfpathlineto{\pgfqpoint{1.360321in}{1.004361in}}%
\pgfpathlineto{\pgfqpoint{1.344687in}{1.006209in}}%
\pgfpathlineto{\pgfqpoint{1.344415in}{0.983100in}}%
\pgfpathclose%
\pgfusepath{fill}%
\end{pgfscope}%
\begin{pgfscope}%
\pgfpathrectangle{\pgfqpoint{0.407592in}{0.222904in}}{\pgfqpoint{1.411163in}{1.411163in}}%
\pgfusepath{clip}%
\pgfsetbuttcap%
\pgfsetroundjoin%
\definecolor{currentfill}{rgb}{0.168627,0.611765,0.411765}%
\pgfsetfillcolor{currentfill}%
\pgfsetfillopacity{0.180000}%
\pgfsetlinewidth{0.000000pt}%
\definecolor{currentstroke}{rgb}{0.000000,0.000000,0.000000}%
\pgfsetstrokecolor{currentstroke}%
\pgfsetdash{}{0pt}%
\pgfpathmoveto{\pgfqpoint{1.360906in}{1.050575in}}%
\pgfpathlineto{\pgfqpoint{1.376475in}{1.049793in}}%
\pgfpathlineto{\pgfqpoint{1.376787in}{1.072901in}}%
\pgfpathlineto{\pgfqpoint{1.361200in}{1.073771in}}%
\pgfpathlineto{\pgfqpoint{1.360906in}{1.050575in}}%
\pgfpathclose%
\pgfusepath{fill}%
\end{pgfscope}%
\begin{pgfscope}%
\pgfpathrectangle{\pgfqpoint{0.407592in}{0.222904in}}{\pgfqpoint{1.411163in}{1.411163in}}%
\pgfusepath{clip}%
\pgfsetbuttcap%
\pgfsetroundjoin%
\definecolor{currentfill}{rgb}{0.168627,0.611765,0.411765}%
\pgfsetfillcolor{currentfill}%
\pgfsetfillopacity{0.180000}%
\pgfsetlinewidth{0.000000pt}%
\definecolor{currentstroke}{rgb}{0.000000,0.000000,0.000000}%
\pgfsetstrokecolor{currentstroke}%
\pgfsetdash{}{0pt}%
\pgfpathmoveto{\pgfqpoint{1.327936in}{0.916764in}}%
\pgfpathlineto{\pgfqpoint{1.343602in}{0.914126in}}%
\pgfpathlineto{\pgfqpoint{1.343873in}{0.937058in}}%
\pgfpathlineto{\pgfqpoint{1.328188in}{0.939790in}}%
\pgfpathlineto{\pgfqpoint{1.327936in}{0.916764in}}%
\pgfpathclose%
\pgfusepath{fill}%
\end{pgfscope}%
\begin{pgfscope}%
\pgfpathrectangle{\pgfqpoint{0.407592in}{0.222904in}}{\pgfqpoint{1.411163in}{1.411163in}}%
\pgfusepath{clip}%
\pgfsetbuttcap%
\pgfsetroundjoin%
\definecolor{currentfill}{rgb}{0.168627,0.611765,0.411765}%
\pgfsetfillcolor{currentfill}%
\pgfsetfillopacity{0.180000}%
\pgfsetlinewidth{0.000000pt}%
\definecolor{currentstroke}{rgb}{0.000000,0.000000,0.000000}%
\pgfsetstrokecolor{currentstroke}%
\pgfsetdash{}{0pt}%
\pgfpathmoveto{\pgfqpoint{1.377414in}{1.119293in}}%
\pgfpathlineto{\pgfqpoint{1.392940in}{1.119583in}}%
\pgfpathlineto{\pgfqpoint{1.393273in}{1.142783in}}%
\pgfpathlineto{\pgfqpoint{1.377728in}{1.142579in}}%
\pgfpathlineto{\pgfqpoint{1.377414in}{1.119293in}}%
\pgfpathclose%
\pgfusepath{fill}%
\end{pgfscope}%
\begin{pgfscope}%
\pgfpathrectangle{\pgfqpoint{0.407592in}{0.222904in}}{\pgfqpoint{1.411163in}{1.411163in}}%
\pgfusepath{clip}%
\pgfsetbuttcap%
\pgfsetroundjoin%
\definecolor{currentfill}{rgb}{0.168627,0.611765,0.411765}%
\pgfsetfillcolor{currentfill}%
\pgfsetfillopacity{0.180000}%
\pgfsetlinewidth{0.000000pt}%
\definecolor{currentstroke}{rgb}{0.000000,0.000000,0.000000}%
\pgfsetstrokecolor{currentstroke}%
\pgfsetdash{}{0pt}%
\pgfpathmoveto{\pgfqpoint{1.344687in}{1.006209in}}%
\pgfpathlineto{\pgfqpoint{1.360321in}{1.004361in}}%
\pgfpathlineto{\pgfqpoint{1.360613in}{1.027438in}}%
\pgfpathlineto{\pgfqpoint{1.344960in}{1.029378in}}%
\pgfpathlineto{\pgfqpoint{1.344687in}{1.006209in}}%
\pgfpathclose%
\pgfusepath{fill}%
\end{pgfscope}%
\begin{pgfscope}%
\pgfpathrectangle{\pgfqpoint{0.407592in}{0.222904in}}{\pgfqpoint{1.411163in}{1.411163in}}%
\pgfusepath{clip}%
\pgfsetbuttcap%
\pgfsetroundjoin%
\definecolor{currentfill}{rgb}{0.168627,0.611765,0.411765}%
\pgfsetfillcolor{currentfill}%
\pgfsetfillopacity{0.180000}%
\pgfsetlinewidth{0.000000pt}%
\definecolor{currentstroke}{rgb}{0.000000,0.000000,0.000000}%
\pgfsetstrokecolor{currentstroke}%
\pgfsetdash{}{0pt}%
\pgfpathmoveto{\pgfqpoint{1.361200in}{1.073771in}}%
\pgfpathlineto{\pgfqpoint{1.376787in}{1.072901in}}%
\pgfpathlineto{\pgfqpoint{1.377100in}{1.096067in}}%
\pgfpathlineto{\pgfqpoint{1.361494in}{1.097026in}}%
\pgfpathlineto{\pgfqpoint{1.361200in}{1.073771in}}%
\pgfpathclose%
\pgfusepath{fill}%
\end{pgfscope}%
\begin{pgfscope}%
\pgfpathrectangle{\pgfqpoint{0.407592in}{0.222904in}}{\pgfqpoint{1.411163in}{1.411163in}}%
\pgfusepath{clip}%
\pgfsetbuttcap%
\pgfsetroundjoin%
\definecolor{currentfill}{rgb}{0.168627,0.611765,0.411765}%
\pgfsetfillcolor{currentfill}%
\pgfsetfillopacity{0.180000}%
\pgfsetlinewidth{0.000000pt}%
\definecolor{currentstroke}{rgb}{0.000000,0.000000,0.000000}%
\pgfsetstrokecolor{currentstroke}%
\pgfsetdash{}{0pt}%
\pgfpathmoveto{\pgfqpoint{1.328188in}{0.939790in}}%
\pgfpathlineto{\pgfqpoint{1.343873in}{0.937058in}}%
\pgfpathlineto{\pgfqpoint{1.344144in}{0.960050in}}%
\pgfpathlineto{\pgfqpoint{1.328439in}{0.962875in}}%
\pgfpathlineto{\pgfqpoint{1.328188in}{0.939790in}}%
\pgfpathclose%
\pgfusepath{fill}%
\end{pgfscope}%
\begin{pgfscope}%
\pgfpathrectangle{\pgfqpoint{0.407592in}{0.222904in}}{\pgfqpoint{1.411163in}{1.411163in}}%
\pgfusepath{clip}%
\pgfsetbuttcap%
\pgfsetroundjoin%
\definecolor{currentfill}{rgb}{0.168627,0.611765,0.411765}%
\pgfsetfillcolor{currentfill}%
\pgfsetfillopacity{0.180000}%
\pgfsetlinewidth{0.000000pt}%
\definecolor{currentstroke}{rgb}{0.000000,0.000000,0.000000}%
\pgfsetstrokecolor{currentstroke}%
\pgfsetdash{}{0pt}%
\pgfpathmoveto{\pgfqpoint{1.377728in}{1.142579in}}%
\pgfpathlineto{\pgfqpoint{1.393273in}{1.142783in}}%
\pgfpathlineto{\pgfqpoint{1.393608in}{1.166043in}}%
\pgfpathlineto{\pgfqpoint{1.378044in}{1.165924in}}%
\pgfpathlineto{\pgfqpoint{1.377728in}{1.142579in}}%
\pgfpathclose%
\pgfusepath{fill}%
\end{pgfscope}%
\begin{pgfscope}%
\pgfpathrectangle{\pgfqpoint{0.407592in}{0.222904in}}{\pgfqpoint{1.411163in}{1.411163in}}%
\pgfusepath{clip}%
\pgfsetbuttcap%
\pgfsetroundjoin%
\definecolor{currentfill}{rgb}{0.168627,0.611765,0.411765}%
\pgfsetfillcolor{currentfill}%
\pgfsetfillopacity{0.180000}%
\pgfsetlinewidth{0.000000pt}%
\definecolor{currentstroke}{rgb}{0.000000,0.000000,0.000000}%
\pgfsetstrokecolor{currentstroke}%
\pgfsetdash{}{0pt}%
\pgfpathmoveto{\pgfqpoint{1.311696in}{0.874411in}}%
\pgfpathlineto{\pgfqpoint{1.327436in}{0.870889in}}%
\pgfpathlineto{\pgfqpoint{1.327686in}{0.893797in}}%
\pgfpathlineto{\pgfqpoint{1.311927in}{0.897415in}}%
\pgfpathlineto{\pgfqpoint{1.311696in}{0.874411in}}%
\pgfpathclose%
\pgfusepath{fill}%
\end{pgfscope}%
\begin{pgfscope}%
\pgfpathrectangle{\pgfqpoint{0.407592in}{0.222904in}}{\pgfqpoint{1.411163in}{1.411163in}}%
\pgfusepath{clip}%
\pgfsetbuttcap%
\pgfsetroundjoin%
\definecolor{currentfill}{rgb}{0.168627,0.611765,0.411765}%
\pgfsetfillcolor{currentfill}%
\pgfsetfillopacity{0.180000}%
\pgfsetlinewidth{0.000000pt}%
\definecolor{currentstroke}{rgb}{0.000000,0.000000,0.000000}%
\pgfsetstrokecolor{currentstroke}%
\pgfsetdash{}{0pt}%
\pgfpathmoveto{\pgfqpoint{1.344960in}{1.029378in}}%
\pgfpathlineto{\pgfqpoint{1.360613in}{1.027438in}}%
\pgfpathlineto{\pgfqpoint{1.360906in}{1.050575in}}%
\pgfpathlineto{\pgfqpoint{1.345234in}{1.052607in}}%
\pgfpathlineto{\pgfqpoint{1.344960in}{1.029378in}}%
\pgfpathclose%
\pgfusepath{fill}%
\end{pgfscope}%
\begin{pgfscope}%
\pgfpathrectangle{\pgfqpoint{0.407592in}{0.222904in}}{\pgfqpoint{1.411163in}{1.411163in}}%
\pgfusepath{clip}%
\pgfsetbuttcap%
\pgfsetroundjoin%
\definecolor{currentfill}{rgb}{0.168627,0.611765,0.411765}%
\pgfsetfillcolor{currentfill}%
\pgfsetfillopacity{0.180000}%
\pgfsetlinewidth{0.000000pt}%
\definecolor{currentstroke}{rgb}{0.000000,0.000000,0.000000}%
\pgfsetstrokecolor{currentstroke}%
\pgfsetdash{}{0pt}%
\pgfpathmoveto{\pgfqpoint{1.361494in}{1.097026in}}%
\pgfpathlineto{\pgfqpoint{1.377100in}{1.096067in}}%
\pgfpathlineto{\pgfqpoint{1.377414in}{1.119293in}}%
\pgfpathlineto{\pgfqpoint{1.361789in}{1.120342in}}%
\pgfpathlineto{\pgfqpoint{1.361494in}{1.097026in}}%
\pgfpathclose%
\pgfusepath{fill}%
\end{pgfscope}%
\begin{pgfscope}%
\pgfpathrectangle{\pgfqpoint{0.407592in}{0.222904in}}{\pgfqpoint{1.411163in}{1.411163in}}%
\pgfusepath{clip}%
\pgfsetbuttcap%
\pgfsetroundjoin%
\definecolor{currentfill}{rgb}{0.168627,0.611765,0.411765}%
\pgfsetfillcolor{currentfill}%
\pgfsetfillopacity{0.180000}%
\pgfsetlinewidth{0.000000pt}%
\definecolor{currentstroke}{rgb}{0.000000,0.000000,0.000000}%
\pgfsetstrokecolor{currentstroke}%
\pgfsetdash{}{0pt}%
\pgfpathmoveto{\pgfqpoint{1.328439in}{0.962875in}}%
\pgfpathlineto{\pgfqpoint{1.344144in}{0.960050in}}%
\pgfpathlineto{\pgfqpoint{1.344415in}{0.983100in}}%
\pgfpathlineto{\pgfqpoint{1.328692in}{0.986020in}}%
\pgfpathlineto{\pgfqpoint{1.328439in}{0.962875in}}%
\pgfpathclose%
\pgfusepath{fill}%
\end{pgfscope}%
\begin{pgfscope}%
\pgfpathrectangle{\pgfqpoint{0.407592in}{0.222904in}}{\pgfqpoint{1.411163in}{1.411163in}}%
\pgfusepath{clip}%
\pgfsetbuttcap%
\pgfsetroundjoin%
\definecolor{currentfill}{rgb}{0.168627,0.611765,0.411765}%
\pgfsetfillcolor{currentfill}%
\pgfsetfillopacity{0.180000}%
\pgfsetlinewidth{0.000000pt}%
\definecolor{currentstroke}{rgb}{0.000000,0.000000,0.000000}%
\pgfsetstrokecolor{currentstroke}%
\pgfsetdash{}{0pt}%
\pgfpathmoveto{\pgfqpoint{1.311927in}{0.897415in}}%
\pgfpathlineto{\pgfqpoint{1.327686in}{0.893797in}}%
\pgfpathlineto{\pgfqpoint{1.327936in}{0.916764in}}%
\pgfpathlineto{\pgfqpoint{1.312158in}{0.920478in}}%
\pgfpathlineto{\pgfqpoint{1.311927in}{0.897415in}}%
\pgfpathclose%
\pgfusepath{fill}%
\end{pgfscope}%
\begin{pgfscope}%
\pgfpathrectangle{\pgfqpoint{0.407592in}{0.222904in}}{\pgfqpoint{1.411163in}{1.411163in}}%
\pgfusepath{clip}%
\pgfsetbuttcap%
\pgfsetroundjoin%
\definecolor{currentfill}{rgb}{0.168627,0.611765,0.411765}%
\pgfsetfillcolor{currentfill}%
\pgfsetfillopacity{0.180000}%
\pgfsetlinewidth{0.000000pt}%
\definecolor{currentstroke}{rgb}{0.000000,0.000000,0.000000}%
\pgfsetstrokecolor{currentstroke}%
\pgfsetdash{}{0pt}%
\pgfpathmoveto{\pgfqpoint{1.345234in}{1.052607in}}%
\pgfpathlineto{\pgfqpoint{1.360906in}{1.050575in}}%
\pgfpathlineto{\pgfqpoint{1.361200in}{1.073771in}}%
\pgfpathlineto{\pgfqpoint{1.345509in}{1.075895in}}%
\pgfpathlineto{\pgfqpoint{1.345234in}{1.052607in}}%
\pgfpathclose%
\pgfusepath{fill}%
\end{pgfscope}%
\begin{pgfscope}%
\pgfpathrectangle{\pgfqpoint{0.407592in}{0.222904in}}{\pgfqpoint{1.411163in}{1.411163in}}%
\pgfusepath{clip}%
\pgfsetbuttcap%
\pgfsetroundjoin%
\definecolor{currentfill}{rgb}{0.168627,0.611765,0.411765}%
\pgfsetfillcolor{currentfill}%
\pgfsetfillopacity{0.180000}%
\pgfsetlinewidth{0.000000pt}%
\definecolor{currentstroke}{rgb}{0.000000,0.000000,0.000000}%
\pgfsetstrokecolor{currentstroke}%
\pgfsetdash{}{0pt}%
\pgfpathmoveto{\pgfqpoint{1.361789in}{1.120342in}}%
\pgfpathlineto{\pgfqpoint{1.377414in}{1.119293in}}%
\pgfpathlineto{\pgfqpoint{1.377728in}{1.142579in}}%
\pgfpathlineto{\pgfqpoint{1.362085in}{1.143717in}}%
\pgfpathlineto{\pgfqpoint{1.361789in}{1.120342in}}%
\pgfpathclose%
\pgfusepath{fill}%
\end{pgfscope}%
\begin{pgfscope}%
\pgfpathrectangle{\pgfqpoint{0.407592in}{0.222904in}}{\pgfqpoint{1.411163in}{1.411163in}}%
\pgfusepath{clip}%
\pgfsetbuttcap%
\pgfsetroundjoin%
\definecolor{currentfill}{rgb}{0.168627,0.611765,0.411765}%
\pgfsetfillcolor{currentfill}%
\pgfsetfillopacity{0.180000}%
\pgfsetlinewidth{0.000000pt}%
\definecolor{currentstroke}{rgb}{0.000000,0.000000,0.000000}%
\pgfsetstrokecolor{currentstroke}%
\pgfsetdash{}{0pt}%
\pgfpathmoveto{\pgfqpoint{1.328692in}{0.986020in}}%
\pgfpathlineto{\pgfqpoint{1.344415in}{0.983100in}}%
\pgfpathlineto{\pgfqpoint{1.344687in}{1.006209in}}%
\pgfpathlineto{\pgfqpoint{1.328945in}{1.009224in}}%
\pgfpathlineto{\pgfqpoint{1.328692in}{0.986020in}}%
\pgfpathclose%
\pgfusepath{fill}%
\end{pgfscope}%
\begin{pgfscope}%
\pgfpathrectangle{\pgfqpoint{0.407592in}{0.222904in}}{\pgfqpoint{1.411163in}{1.411163in}}%
\pgfusepath{clip}%
\pgfsetbuttcap%
\pgfsetroundjoin%
\definecolor{currentfill}{rgb}{0.168627,0.611765,0.411765}%
\pgfsetfillcolor{currentfill}%
\pgfsetfillopacity{0.180000}%
\pgfsetlinewidth{0.000000pt}%
\definecolor{currentstroke}{rgb}{0.000000,0.000000,0.000000}%
\pgfsetstrokecolor{currentstroke}%
\pgfsetdash{}{0pt}%
\pgfpathmoveto{\pgfqpoint{1.312158in}{0.920478in}}%
\pgfpathlineto{\pgfqpoint{1.327936in}{0.916764in}}%
\pgfpathlineto{\pgfqpoint{1.328188in}{0.939790in}}%
\pgfpathlineto{\pgfqpoint{1.312390in}{0.943601in}}%
\pgfpathlineto{\pgfqpoint{1.312158in}{0.920478in}}%
\pgfpathclose%
\pgfusepath{fill}%
\end{pgfscope}%
\begin{pgfscope}%
\pgfpathrectangle{\pgfqpoint{0.407592in}{0.222904in}}{\pgfqpoint{1.411163in}{1.411163in}}%
\pgfusepath{clip}%
\pgfsetbuttcap%
\pgfsetroundjoin%
\definecolor{currentfill}{rgb}{0.168627,0.611765,0.411765}%
\pgfsetfillcolor{currentfill}%
\pgfsetfillopacity{0.180000}%
\pgfsetlinewidth{0.000000pt}%
\definecolor{currentstroke}{rgb}{0.000000,0.000000,0.000000}%
\pgfsetstrokecolor{currentstroke}%
\pgfsetdash{}{0pt}%
\pgfpathmoveto{\pgfqpoint{1.345509in}{1.075895in}}%
\pgfpathlineto{\pgfqpoint{1.361200in}{1.073771in}}%
\pgfpathlineto{\pgfqpoint{1.361494in}{1.097026in}}%
\pgfpathlineto{\pgfqpoint{1.345784in}{1.099243in}}%
\pgfpathlineto{\pgfqpoint{1.345509in}{1.075895in}}%
\pgfpathclose%
\pgfusepath{fill}%
\end{pgfscope}%
\begin{pgfscope}%
\pgfpathrectangle{\pgfqpoint{0.407592in}{0.222904in}}{\pgfqpoint{1.411163in}{1.411163in}}%
\pgfusepath{clip}%
\pgfsetbuttcap%
\pgfsetroundjoin%
\definecolor{currentfill}{rgb}{0.168627,0.611765,0.411765}%
\pgfsetfillcolor{currentfill}%
\pgfsetfillopacity{0.180000}%
\pgfsetlinewidth{0.000000pt}%
\definecolor{currentstroke}{rgb}{0.000000,0.000000,0.000000}%
\pgfsetstrokecolor{currentstroke}%
\pgfsetdash{}{0pt}%
\pgfpathmoveto{\pgfqpoint{1.362085in}{1.143717in}}%
\pgfpathlineto{\pgfqpoint{1.377728in}{1.142579in}}%
\pgfpathlineto{\pgfqpoint{1.378044in}{1.165924in}}%
\pgfpathlineto{\pgfqpoint{1.362381in}{1.167153in}}%
\pgfpathlineto{\pgfqpoint{1.362085in}{1.143717in}}%
\pgfpathclose%
\pgfusepath{fill}%
\end{pgfscope}%
\begin{pgfscope}%
\pgfpathrectangle{\pgfqpoint{0.407592in}{0.222904in}}{\pgfqpoint{1.411163in}{1.411163in}}%
\pgfusepath{clip}%
\pgfsetbuttcap%
\pgfsetroundjoin%
\definecolor{currentfill}{rgb}{0.168627,0.611765,0.411765}%
\pgfsetfillcolor{currentfill}%
\pgfsetfillopacity{0.180000}%
\pgfsetlinewidth{0.000000pt}%
\definecolor{currentstroke}{rgb}{0.000000,0.000000,0.000000}%
\pgfsetstrokecolor{currentstroke}%
\pgfsetdash{}{0pt}%
\pgfpathmoveto{\pgfqpoint{1.328945in}{1.009224in}}%
\pgfpathlineto{\pgfqpoint{1.344687in}{1.006209in}}%
\pgfpathlineto{\pgfqpoint{1.344960in}{1.029378in}}%
\pgfpathlineto{\pgfqpoint{1.329199in}{1.032488in}}%
\pgfpathlineto{\pgfqpoint{1.328945in}{1.009224in}}%
\pgfpathclose%
\pgfusepath{fill}%
\end{pgfscope}%
\begin{pgfscope}%
\pgfpathrectangle{\pgfqpoint{0.407592in}{0.222904in}}{\pgfqpoint{1.411163in}{1.411163in}}%
\pgfusepath{clip}%
\pgfsetbuttcap%
\pgfsetroundjoin%
\definecolor{currentfill}{rgb}{0.168627,0.611765,0.411765}%
\pgfsetfillcolor{currentfill}%
\pgfsetfillopacity{0.180000}%
\pgfsetlinewidth{0.000000pt}%
\definecolor{currentstroke}{rgb}{0.000000,0.000000,0.000000}%
\pgfsetstrokecolor{currentstroke}%
\pgfsetdash{}{0pt}%
\pgfpathmoveto{\pgfqpoint{1.312390in}{0.943601in}}%
\pgfpathlineto{\pgfqpoint{1.328188in}{0.939790in}}%
\pgfpathlineto{\pgfqpoint{1.328439in}{0.962875in}}%
\pgfpathlineto{\pgfqpoint{1.312622in}{0.966783in}}%
\pgfpathlineto{\pgfqpoint{1.312390in}{0.943601in}}%
\pgfpathclose%
\pgfusepath{fill}%
\end{pgfscope}%
\begin{pgfscope}%
\pgfpathrectangle{\pgfqpoint{0.407592in}{0.222904in}}{\pgfqpoint{1.411163in}{1.411163in}}%
\pgfusepath{clip}%
\pgfsetbuttcap%
\pgfsetroundjoin%
\definecolor{currentfill}{rgb}{0.168627,0.611765,0.411765}%
\pgfsetfillcolor{currentfill}%
\pgfsetfillopacity{0.180000}%
\pgfsetlinewidth{0.000000pt}%
\definecolor{currentstroke}{rgb}{0.000000,0.000000,0.000000}%
\pgfsetstrokecolor{currentstroke}%
\pgfsetdash{}{0pt}%
\pgfpathmoveto{\pgfqpoint{1.295839in}{0.878928in}}%
\pgfpathlineto{\pgfqpoint{1.311696in}{0.874411in}}%
\pgfpathlineto{\pgfqpoint{1.311927in}{0.897415in}}%
\pgfpathlineto{\pgfqpoint{1.296051in}{0.902031in}}%
\pgfpathlineto{\pgfqpoint{1.295839in}{0.878928in}}%
\pgfpathclose%
\pgfusepath{fill}%
\end{pgfscope}%
\begin{pgfscope}%
\pgfpathrectangle{\pgfqpoint{0.407592in}{0.222904in}}{\pgfqpoint{1.411163in}{1.411163in}}%
\pgfusepath{clip}%
\pgfsetbuttcap%
\pgfsetroundjoin%
\definecolor{currentfill}{rgb}{0.168627,0.611765,0.411765}%
\pgfsetfillcolor{currentfill}%
\pgfsetfillopacity{0.180000}%
\pgfsetlinewidth{0.000000pt}%
\definecolor{currentstroke}{rgb}{0.000000,0.000000,0.000000}%
\pgfsetstrokecolor{currentstroke}%
\pgfsetdash{}{0pt}%
\pgfpathmoveto{\pgfqpoint{1.345784in}{1.099243in}}%
\pgfpathlineto{\pgfqpoint{1.361494in}{1.097026in}}%
\pgfpathlineto{\pgfqpoint{1.361789in}{1.120342in}}%
\pgfpathlineto{\pgfqpoint{1.346060in}{1.122652in}}%
\pgfpathlineto{\pgfqpoint{1.345784in}{1.099243in}}%
\pgfpathclose%
\pgfusepath{fill}%
\end{pgfscope}%
\begin{pgfscope}%
\pgfpathrectangle{\pgfqpoint{0.407592in}{0.222904in}}{\pgfqpoint{1.411163in}{1.411163in}}%
\pgfusepath{clip}%
\pgfsetbuttcap%
\pgfsetroundjoin%
\definecolor{currentfill}{rgb}{0.168627,0.611765,0.411765}%
\pgfsetfillcolor{currentfill}%
\pgfsetfillopacity{0.180000}%
\pgfsetlinewidth{0.000000pt}%
\definecolor{currentstroke}{rgb}{0.000000,0.000000,0.000000}%
\pgfsetstrokecolor{currentstroke}%
\pgfsetdash{}{0pt}%
\pgfpathmoveto{\pgfqpoint{1.329199in}{1.032488in}}%
\pgfpathlineto{\pgfqpoint{1.344960in}{1.029378in}}%
\pgfpathlineto{\pgfqpoint{1.345234in}{1.052607in}}%
\pgfpathlineto{\pgfqpoint{1.329453in}{1.055812in}}%
\pgfpathlineto{\pgfqpoint{1.329199in}{1.032488in}}%
\pgfpathclose%
\pgfusepath{fill}%
\end{pgfscope}%
\begin{pgfscope}%
\pgfpathrectangle{\pgfqpoint{0.407592in}{0.222904in}}{\pgfqpoint{1.411163in}{1.411163in}}%
\pgfusepath{clip}%
\pgfsetbuttcap%
\pgfsetroundjoin%
\definecolor{currentfill}{rgb}{0.168627,0.611765,0.411765}%
\pgfsetfillcolor{currentfill}%
\pgfsetfillopacity{0.180000}%
\pgfsetlinewidth{0.000000pt}%
\definecolor{currentstroke}{rgb}{0.000000,0.000000,0.000000}%
\pgfsetstrokecolor{currentstroke}%
\pgfsetdash{}{0pt}%
\pgfpathmoveto{\pgfqpoint{1.312622in}{0.966783in}}%
\pgfpathlineto{\pgfqpoint{1.328439in}{0.962875in}}%
\pgfpathlineto{\pgfqpoint{1.328692in}{0.986020in}}%
\pgfpathlineto{\pgfqpoint{1.312855in}{0.990025in}}%
\pgfpathlineto{\pgfqpoint{1.312622in}{0.966783in}}%
\pgfpathclose%
\pgfusepath{fill}%
\end{pgfscope}%
\begin{pgfscope}%
\pgfpathrectangle{\pgfqpoint{0.407592in}{0.222904in}}{\pgfqpoint{1.411163in}{1.411163in}}%
\pgfusepath{clip}%
\pgfsetbuttcap%
\pgfsetroundjoin%
\definecolor{currentfill}{rgb}{0.168627,0.611765,0.411765}%
\pgfsetfillcolor{currentfill}%
\pgfsetfillopacity{0.180000}%
\pgfsetlinewidth{0.000000pt}%
\definecolor{currentstroke}{rgb}{0.000000,0.000000,0.000000}%
\pgfsetstrokecolor{currentstroke}%
\pgfsetdash{}{0pt}%
\pgfpathmoveto{\pgfqpoint{1.296051in}{0.902031in}}%
\pgfpathlineto{\pgfqpoint{1.311927in}{0.897415in}}%
\pgfpathlineto{\pgfqpoint{1.312158in}{0.920478in}}%
\pgfpathlineto{\pgfqpoint{1.296262in}{0.925194in}}%
\pgfpathlineto{\pgfqpoint{1.296051in}{0.902031in}}%
\pgfpathclose%
\pgfusepath{fill}%
\end{pgfscope}%
\begin{pgfscope}%
\pgfpathrectangle{\pgfqpoint{0.407592in}{0.222904in}}{\pgfqpoint{1.411163in}{1.411163in}}%
\pgfusepath{clip}%
\pgfsetbuttcap%
\pgfsetroundjoin%
\definecolor{currentfill}{rgb}{0.168627,0.611765,0.411765}%
\pgfsetfillcolor{currentfill}%
\pgfsetfillopacity{0.180000}%
\pgfsetlinewidth{0.000000pt}%
\definecolor{currentstroke}{rgb}{0.000000,0.000000,0.000000}%
\pgfsetstrokecolor{currentstroke}%
\pgfsetdash{}{0pt}%
\pgfpathmoveto{\pgfqpoint{1.346060in}{1.122652in}}%
\pgfpathlineto{\pgfqpoint{1.361789in}{1.120342in}}%
\pgfpathlineto{\pgfqpoint{1.362085in}{1.143717in}}%
\pgfpathlineto{\pgfqpoint{1.346336in}{1.146121in}}%
\pgfpathlineto{\pgfqpoint{1.346060in}{1.122652in}}%
\pgfpathclose%
\pgfusepath{fill}%
\end{pgfscope}%
\begin{pgfscope}%
\pgfpathrectangle{\pgfqpoint{0.407592in}{0.222904in}}{\pgfqpoint{1.411163in}{1.411163in}}%
\pgfusepath{clip}%
\pgfsetbuttcap%
\pgfsetroundjoin%
\definecolor{currentfill}{rgb}{0.168627,0.611765,0.411765}%
\pgfsetfillcolor{currentfill}%
\pgfsetfillopacity{0.180000}%
\pgfsetlinewidth{0.000000pt}%
\definecolor{currentstroke}{rgb}{0.000000,0.000000,0.000000}%
\pgfsetstrokecolor{currentstroke}%
\pgfsetdash{}{0pt}%
\pgfpathmoveto{\pgfqpoint{1.329453in}{1.055812in}}%
\pgfpathlineto{\pgfqpoint{1.345234in}{1.052607in}}%
\pgfpathlineto{\pgfqpoint{1.345509in}{1.075895in}}%
\pgfpathlineto{\pgfqpoint{1.329708in}{1.079196in}}%
\pgfpathlineto{\pgfqpoint{1.329453in}{1.055812in}}%
\pgfpathclose%
\pgfusepath{fill}%
\end{pgfscope}%
\begin{pgfscope}%
\pgfpathrectangle{\pgfqpoint{0.407592in}{0.222904in}}{\pgfqpoint{1.411163in}{1.411163in}}%
\pgfusepath{clip}%
\pgfsetbuttcap%
\pgfsetroundjoin%
\definecolor{currentfill}{rgb}{0.168627,0.611765,0.411765}%
\pgfsetfillcolor{currentfill}%
\pgfsetfillopacity{0.180000}%
\pgfsetlinewidth{0.000000pt}%
\definecolor{currentstroke}{rgb}{0.000000,0.000000,0.000000}%
\pgfsetstrokecolor{currentstroke}%
\pgfsetdash{}{0pt}%
\pgfpathmoveto{\pgfqpoint{1.312855in}{0.990025in}}%
\pgfpathlineto{\pgfqpoint{1.328692in}{0.986020in}}%
\pgfpathlineto{\pgfqpoint{1.328945in}{1.009224in}}%
\pgfpathlineto{\pgfqpoint{1.313089in}{1.013327in}}%
\pgfpathlineto{\pgfqpoint{1.312855in}{0.990025in}}%
\pgfpathclose%
\pgfusepath{fill}%
\end{pgfscope}%
\begin{pgfscope}%
\pgfpathrectangle{\pgfqpoint{0.407592in}{0.222904in}}{\pgfqpoint{1.411163in}{1.411163in}}%
\pgfusepath{clip}%
\pgfsetbuttcap%
\pgfsetroundjoin%
\definecolor{currentfill}{rgb}{0.168627,0.611765,0.411765}%
\pgfsetfillcolor{currentfill}%
\pgfsetfillopacity{0.180000}%
\pgfsetlinewidth{0.000000pt}%
\definecolor{currentstroke}{rgb}{0.000000,0.000000,0.000000}%
\pgfsetstrokecolor{currentstroke}%
\pgfsetdash{}{0pt}%
\pgfpathmoveto{\pgfqpoint{1.296262in}{0.925194in}}%
\pgfpathlineto{\pgfqpoint{1.312158in}{0.920478in}}%
\pgfpathlineto{\pgfqpoint{1.312390in}{0.943601in}}%
\pgfpathlineto{\pgfqpoint{1.296474in}{0.948416in}}%
\pgfpathlineto{\pgfqpoint{1.296262in}{0.925194in}}%
\pgfpathclose%
\pgfusepath{fill}%
\end{pgfscope}%
\begin{pgfscope}%
\pgfpathrectangle{\pgfqpoint{0.407592in}{0.222904in}}{\pgfqpoint{1.411163in}{1.411163in}}%
\pgfusepath{clip}%
\pgfsetbuttcap%
\pgfsetroundjoin%
\definecolor{currentfill}{rgb}{0.168627,0.611765,0.411765}%
\pgfsetfillcolor{currentfill}%
\pgfsetfillopacity{0.180000}%
\pgfsetlinewidth{0.000000pt}%
\definecolor{currentstroke}{rgb}{0.000000,0.000000,0.000000}%
\pgfsetstrokecolor{currentstroke}%
\pgfsetdash{}{0pt}%
\pgfpathmoveto{\pgfqpoint{1.346336in}{1.146121in}}%
\pgfpathlineto{\pgfqpoint{1.362085in}{1.143717in}}%
\pgfpathlineto{\pgfqpoint{1.362381in}{1.167153in}}%
\pgfpathlineto{\pgfqpoint{1.346613in}{1.169651in}}%
\pgfpathlineto{\pgfqpoint{1.346336in}{1.146121in}}%
\pgfpathclose%
\pgfusepath{fill}%
\end{pgfscope}%
\begin{pgfscope}%
\pgfpathrectangle{\pgfqpoint{0.407592in}{0.222904in}}{\pgfqpoint{1.411163in}{1.411163in}}%
\pgfusepath{clip}%
\pgfsetbuttcap%
\pgfsetroundjoin%
\definecolor{currentfill}{rgb}{0.168627,0.611765,0.411765}%
\pgfsetfillcolor{currentfill}%
\pgfsetfillopacity{0.180000}%
\pgfsetlinewidth{0.000000pt}%
\definecolor{currentstroke}{rgb}{0.000000,0.000000,0.000000}%
\pgfsetstrokecolor{currentstroke}%
\pgfsetdash{}{0pt}%
\pgfpathmoveto{\pgfqpoint{1.329708in}{1.079196in}}%
\pgfpathlineto{\pgfqpoint{1.345509in}{1.075895in}}%
\pgfpathlineto{\pgfqpoint{1.345784in}{1.099243in}}%
\pgfpathlineto{\pgfqpoint{1.329964in}{1.102641in}}%
\pgfpathlineto{\pgfqpoint{1.329708in}{1.079196in}}%
\pgfpathclose%
\pgfusepath{fill}%
\end{pgfscope}%
\begin{pgfscope}%
\pgfpathrectangle{\pgfqpoint{0.407592in}{0.222904in}}{\pgfqpoint{1.411163in}{1.411163in}}%
\pgfusepath{clip}%
\pgfsetbuttcap%
\pgfsetroundjoin%
\definecolor{currentfill}{rgb}{0.168627,0.611765,0.411765}%
\pgfsetfillcolor{currentfill}%
\pgfsetfillopacity{0.180000}%
\pgfsetlinewidth{0.000000pt}%
\definecolor{currentstroke}{rgb}{0.000000,0.000000,0.000000}%
\pgfsetstrokecolor{currentstroke}%
\pgfsetdash{}{0pt}%
\pgfpathmoveto{\pgfqpoint{1.313089in}{1.013327in}}%
\pgfpathlineto{\pgfqpoint{1.328945in}{1.009224in}}%
\pgfpathlineto{\pgfqpoint{1.329199in}{1.032488in}}%
\pgfpathlineto{\pgfqpoint{1.313323in}{1.036690in}}%
\pgfpathlineto{\pgfqpoint{1.313089in}{1.013327in}}%
\pgfpathclose%
\pgfusepath{fill}%
\end{pgfscope}%
\begin{pgfscope}%
\pgfpathrectangle{\pgfqpoint{0.407592in}{0.222904in}}{\pgfqpoint{1.411163in}{1.411163in}}%
\pgfusepath{clip}%
\pgfsetbuttcap%
\pgfsetroundjoin%
\definecolor{currentfill}{rgb}{0.168627,0.611765,0.411765}%
\pgfsetfillcolor{currentfill}%
\pgfsetfillopacity{0.180000}%
\pgfsetlinewidth{0.000000pt}%
\definecolor{currentstroke}{rgb}{0.000000,0.000000,0.000000}%
\pgfsetstrokecolor{currentstroke}%
\pgfsetdash{}{0pt}%
\pgfpathmoveto{\pgfqpoint{1.296474in}{0.948416in}}%
\pgfpathlineto{\pgfqpoint{1.312390in}{0.943601in}}%
\pgfpathlineto{\pgfqpoint{1.312622in}{0.966783in}}%
\pgfpathlineto{\pgfqpoint{1.296687in}{0.971699in}}%
\pgfpathlineto{\pgfqpoint{1.296474in}{0.948416in}}%
\pgfpathclose%
\pgfusepath{fill}%
\end{pgfscope}%
\begin{pgfscope}%
\pgfpathrectangle{\pgfqpoint{0.407592in}{0.222904in}}{\pgfqpoint{1.411163in}{1.411163in}}%
\pgfusepath{clip}%
\pgfsetbuttcap%
\pgfsetroundjoin%
\definecolor{currentfill}{rgb}{0.168627,0.611765,0.411765}%
\pgfsetfillcolor{currentfill}%
\pgfsetfillopacity{0.180000}%
\pgfsetlinewidth{0.000000pt}%
\definecolor{currentstroke}{rgb}{0.000000,0.000000,0.000000}%
\pgfsetstrokecolor{currentstroke}%
\pgfsetdash{}{0pt}%
\pgfpathmoveto{\pgfqpoint{1.279861in}{0.884370in}}%
\pgfpathlineto{\pgfqpoint{1.295839in}{0.878928in}}%
\pgfpathlineto{\pgfqpoint{1.296051in}{0.902031in}}%
\pgfpathlineto{\pgfqpoint{1.280052in}{0.907576in}}%
\pgfpathlineto{\pgfqpoint{1.279861in}{0.884370in}}%
\pgfpathclose%
\pgfusepath{fill}%
\end{pgfscope}%
\begin{pgfscope}%
\pgfpathrectangle{\pgfqpoint{0.407592in}{0.222904in}}{\pgfqpoint{1.411163in}{1.411163in}}%
\pgfusepath{clip}%
\pgfsetbuttcap%
\pgfsetroundjoin%
\definecolor{currentfill}{rgb}{0.168627,0.611765,0.411765}%
\pgfsetfillcolor{currentfill}%
\pgfsetfillopacity{0.180000}%
\pgfsetlinewidth{0.000000pt}%
\definecolor{currentstroke}{rgb}{0.000000,0.000000,0.000000}%
\pgfsetstrokecolor{currentstroke}%
\pgfsetdash{}{0pt}%
\pgfpathmoveto{\pgfqpoint{1.329964in}{1.102641in}}%
\pgfpathlineto{\pgfqpoint{1.345784in}{1.099243in}}%
\pgfpathlineto{\pgfqpoint{1.346060in}{1.122652in}}%
\pgfpathlineto{\pgfqpoint{1.330221in}{1.126147in}}%
\pgfpathlineto{\pgfqpoint{1.329964in}{1.102641in}}%
\pgfpathclose%
\pgfusepath{fill}%
\end{pgfscope}%
\begin{pgfscope}%
\pgfpathrectangle{\pgfqpoint{0.407592in}{0.222904in}}{\pgfqpoint{1.411163in}{1.411163in}}%
\pgfusepath{clip}%
\pgfsetbuttcap%
\pgfsetroundjoin%
\definecolor{currentfill}{rgb}{0.168627,0.611765,0.411765}%
\pgfsetfillcolor{currentfill}%
\pgfsetfillopacity{0.180000}%
\pgfsetlinewidth{0.000000pt}%
\definecolor{currentstroke}{rgb}{0.000000,0.000000,0.000000}%
\pgfsetstrokecolor{currentstroke}%
\pgfsetdash{}{0pt}%
\pgfpathmoveto{\pgfqpoint{1.313323in}{1.036690in}}%
\pgfpathlineto{\pgfqpoint{1.329199in}{1.032488in}}%
\pgfpathlineto{\pgfqpoint{1.329453in}{1.055812in}}%
\pgfpathlineto{\pgfqpoint{1.313558in}{1.060113in}}%
\pgfpathlineto{\pgfqpoint{1.313323in}{1.036690in}}%
\pgfpathclose%
\pgfusepath{fill}%
\end{pgfscope}%
\begin{pgfscope}%
\pgfpathrectangle{\pgfqpoint{0.407592in}{0.222904in}}{\pgfqpoint{1.411163in}{1.411163in}}%
\pgfusepath{clip}%
\pgfsetbuttcap%
\pgfsetroundjoin%
\definecolor{currentfill}{rgb}{0.168627,0.611765,0.411765}%
\pgfsetfillcolor{currentfill}%
\pgfsetfillopacity{0.180000}%
\pgfsetlinewidth{0.000000pt}%
\definecolor{currentstroke}{rgb}{0.000000,0.000000,0.000000}%
\pgfsetstrokecolor{currentstroke}%
\pgfsetdash{}{0pt}%
\pgfpathmoveto{\pgfqpoint{1.296687in}{0.971699in}}%
\pgfpathlineto{\pgfqpoint{1.312622in}{0.966783in}}%
\pgfpathlineto{\pgfqpoint{1.312855in}{0.990025in}}%
\pgfpathlineto{\pgfqpoint{1.296901in}{0.995042in}}%
\pgfpathlineto{\pgfqpoint{1.296687in}{0.971699in}}%
\pgfpathclose%
\pgfusepath{fill}%
\end{pgfscope}%
\begin{pgfscope}%
\pgfpathrectangle{\pgfqpoint{0.407592in}{0.222904in}}{\pgfqpoint{1.411163in}{1.411163in}}%
\pgfusepath{clip}%
\pgfsetbuttcap%
\pgfsetroundjoin%
\definecolor{currentfill}{rgb}{0.168627,0.611765,0.411765}%
\pgfsetfillcolor{currentfill}%
\pgfsetfillopacity{0.180000}%
\pgfsetlinewidth{0.000000pt}%
\definecolor{currentstroke}{rgb}{0.000000,0.000000,0.000000}%
\pgfsetstrokecolor{currentstroke}%
\pgfsetdash{}{0pt}%
\pgfpathmoveto{\pgfqpoint{1.280052in}{0.907576in}}%
\pgfpathlineto{\pgfqpoint{1.296051in}{0.902031in}}%
\pgfpathlineto{\pgfqpoint{1.296262in}{0.925194in}}%
\pgfpathlineto{\pgfqpoint{1.280244in}{0.930841in}}%
\pgfpathlineto{\pgfqpoint{1.280052in}{0.907576in}}%
\pgfpathclose%
\pgfusepath{fill}%
\end{pgfscope}%
\begin{pgfscope}%
\pgfpathrectangle{\pgfqpoint{0.407592in}{0.222904in}}{\pgfqpoint{1.411163in}{1.411163in}}%
\pgfusepath{clip}%
\pgfsetbuttcap%
\pgfsetroundjoin%
\definecolor{currentfill}{rgb}{0.168627,0.611765,0.411765}%
\pgfsetfillcolor{currentfill}%
\pgfsetfillopacity{0.180000}%
\pgfsetlinewidth{0.000000pt}%
\definecolor{currentstroke}{rgb}{0.000000,0.000000,0.000000}%
\pgfsetstrokecolor{currentstroke}%
\pgfsetdash{}{0pt}%
\pgfpathmoveto{\pgfqpoint{1.330221in}{1.126147in}}%
\pgfpathlineto{\pgfqpoint{1.346060in}{1.122652in}}%
\pgfpathlineto{\pgfqpoint{1.346336in}{1.146121in}}%
\pgfpathlineto{\pgfqpoint{1.330478in}{1.149713in}}%
\pgfpathlineto{\pgfqpoint{1.330221in}{1.126147in}}%
\pgfpathclose%
\pgfusepath{fill}%
\end{pgfscope}%
\begin{pgfscope}%
\pgfpathrectangle{\pgfqpoint{0.407592in}{0.222904in}}{\pgfqpoint{1.411163in}{1.411163in}}%
\pgfusepath{clip}%
\pgfsetbuttcap%
\pgfsetroundjoin%
\definecolor{currentfill}{rgb}{0.168627,0.611765,0.411765}%
\pgfsetfillcolor{currentfill}%
\pgfsetfillopacity{0.180000}%
\pgfsetlinewidth{0.000000pt}%
\definecolor{currentstroke}{rgb}{0.000000,0.000000,0.000000}%
\pgfsetstrokecolor{currentstroke}%
\pgfsetdash{}{0pt}%
\pgfpathmoveto{\pgfqpoint{1.313558in}{1.060113in}}%
\pgfpathlineto{\pgfqpoint{1.329453in}{1.055812in}}%
\pgfpathlineto{\pgfqpoint{1.329708in}{1.079196in}}%
\pgfpathlineto{\pgfqpoint{1.313794in}{1.083597in}}%
\pgfpathlineto{\pgfqpoint{1.313558in}{1.060113in}}%
\pgfpathclose%
\pgfusepath{fill}%
\end{pgfscope}%
\begin{pgfscope}%
\pgfpathrectangle{\pgfqpoint{0.407592in}{0.222904in}}{\pgfqpoint{1.411163in}{1.411163in}}%
\pgfusepath{clip}%
\pgfsetbuttcap%
\pgfsetroundjoin%
\definecolor{currentfill}{rgb}{0.168627,0.611765,0.411765}%
\pgfsetfillcolor{currentfill}%
\pgfsetfillopacity{0.180000}%
\pgfsetlinewidth{0.000000pt}%
\definecolor{currentstroke}{rgb}{0.000000,0.000000,0.000000}%
\pgfsetstrokecolor{currentstroke}%
\pgfsetdash{}{0pt}%
\pgfpathmoveto{\pgfqpoint{1.296901in}{0.995042in}}%
\pgfpathlineto{\pgfqpoint{1.312855in}{0.990025in}}%
\pgfpathlineto{\pgfqpoint{1.313089in}{1.013327in}}%
\pgfpathlineto{\pgfqpoint{1.297114in}{1.018446in}}%
\pgfpathlineto{\pgfqpoint{1.296901in}{0.995042in}}%
\pgfpathclose%
\pgfusepath{fill}%
\end{pgfscope}%
\begin{pgfscope}%
\pgfpathrectangle{\pgfqpoint{0.407592in}{0.222904in}}{\pgfqpoint{1.411163in}{1.411163in}}%
\pgfusepath{clip}%
\pgfsetbuttcap%
\pgfsetroundjoin%
\definecolor{currentfill}{rgb}{0.168627,0.611765,0.411765}%
\pgfsetfillcolor{currentfill}%
\pgfsetfillopacity{0.180000}%
\pgfsetlinewidth{0.000000pt}%
\definecolor{currentstroke}{rgb}{0.000000,0.000000,0.000000}%
\pgfsetstrokecolor{currentstroke}%
\pgfsetdash{}{0pt}%
\pgfpathmoveto{\pgfqpoint{1.280244in}{0.930841in}}%
\pgfpathlineto{\pgfqpoint{1.296262in}{0.925194in}}%
\pgfpathlineto{\pgfqpoint{1.296474in}{0.948416in}}%
\pgfpathlineto{\pgfqpoint{1.280436in}{0.954167in}}%
\pgfpathlineto{\pgfqpoint{1.280244in}{0.930841in}}%
\pgfpathclose%
\pgfusepath{fill}%
\end{pgfscope}%
\begin{pgfscope}%
\pgfpathrectangle{\pgfqpoint{0.407592in}{0.222904in}}{\pgfqpoint{1.411163in}{1.411163in}}%
\pgfusepath{clip}%
\pgfsetbuttcap%
\pgfsetroundjoin%
\definecolor{currentfill}{rgb}{0.168627,0.611765,0.411765}%
\pgfsetfillcolor{currentfill}%
\pgfsetfillopacity{0.180000}%
\pgfsetlinewidth{0.000000pt}%
\definecolor{currentstroke}{rgb}{0.000000,0.000000,0.000000}%
\pgfsetstrokecolor{currentstroke}%
\pgfsetdash{}{0pt}%
\pgfpathmoveto{\pgfqpoint{1.330478in}{1.149713in}}%
\pgfpathlineto{\pgfqpoint{1.346336in}{1.146121in}}%
\pgfpathlineto{\pgfqpoint{1.346613in}{1.169651in}}%
\pgfpathlineto{\pgfqpoint{1.330735in}{1.173342in}}%
\pgfpathlineto{\pgfqpoint{1.330478in}{1.149713in}}%
\pgfpathclose%
\pgfusepath{fill}%
\end{pgfscope}%
\begin{pgfscope}%
\pgfpathrectangle{\pgfqpoint{0.407592in}{0.222904in}}{\pgfqpoint{1.411163in}{1.411163in}}%
\pgfusepath{clip}%
\pgfsetbuttcap%
\pgfsetroundjoin%
\definecolor{currentfill}{rgb}{0.168627,0.611765,0.411765}%
\pgfsetfillcolor{currentfill}%
\pgfsetfillopacity{0.180000}%
\pgfsetlinewidth{0.000000pt}%
\definecolor{currentstroke}{rgb}{0.000000,0.000000,0.000000}%
\pgfsetstrokecolor{currentstroke}%
\pgfsetdash{}{0pt}%
\pgfpathmoveto{\pgfqpoint{1.313794in}{1.083597in}}%
\pgfpathlineto{\pgfqpoint{1.329708in}{1.079196in}}%
\pgfpathlineto{\pgfqpoint{1.329964in}{1.102641in}}%
\pgfpathlineto{\pgfqpoint{1.314030in}{1.107142in}}%
\pgfpathlineto{\pgfqpoint{1.313794in}{1.083597in}}%
\pgfpathclose%
\pgfusepath{fill}%
\end{pgfscope}%
\begin{pgfscope}%
\pgfpathrectangle{\pgfqpoint{0.407592in}{0.222904in}}{\pgfqpoint{1.411163in}{1.411163in}}%
\pgfusepath{clip}%
\pgfsetbuttcap%
\pgfsetroundjoin%
\definecolor{currentfill}{rgb}{0.168627,0.611765,0.411765}%
\pgfsetfillcolor{currentfill}%
\pgfsetfillopacity{0.180000}%
\pgfsetlinewidth{0.000000pt}%
\definecolor{currentstroke}{rgb}{0.000000,0.000000,0.000000}%
\pgfsetstrokecolor{currentstroke}%
\pgfsetdash{}{0pt}%
\pgfpathmoveto{\pgfqpoint{1.297114in}{1.018446in}}%
\pgfpathlineto{\pgfqpoint{1.313089in}{1.013327in}}%
\pgfpathlineto{\pgfqpoint{1.313323in}{1.036690in}}%
\pgfpathlineto{\pgfqpoint{1.297329in}{1.041911in}}%
\pgfpathlineto{\pgfqpoint{1.297114in}{1.018446in}}%
\pgfpathclose%
\pgfusepath{fill}%
\end{pgfscope}%
\begin{pgfscope}%
\pgfpathrectangle{\pgfqpoint{0.407592in}{0.222904in}}{\pgfqpoint{1.411163in}{1.411163in}}%
\pgfusepath{clip}%
\pgfsetbuttcap%
\pgfsetroundjoin%
\definecolor{currentfill}{rgb}{0.168627,0.611765,0.411765}%
\pgfsetfillcolor{currentfill}%
\pgfsetfillopacity{0.180000}%
\pgfsetlinewidth{0.000000pt}%
\definecolor{currentstroke}{rgb}{0.000000,0.000000,0.000000}%
\pgfsetstrokecolor{currentstroke}%
\pgfsetdash{}{0pt}%
\pgfpathmoveto{\pgfqpoint{1.280436in}{0.954167in}}%
\pgfpathlineto{\pgfqpoint{1.296474in}{0.948416in}}%
\pgfpathlineto{\pgfqpoint{1.296687in}{0.971699in}}%
\pgfpathlineto{\pgfqpoint{1.280629in}{0.977553in}}%
\pgfpathlineto{\pgfqpoint{1.280436in}{0.954167in}}%
\pgfpathclose%
\pgfusepath{fill}%
\end{pgfscope}%
\begin{pgfscope}%
\pgfpathrectangle{\pgfqpoint{0.407592in}{0.222904in}}{\pgfqpoint{1.411163in}{1.411163in}}%
\pgfusepath{clip}%
\pgfsetbuttcap%
\pgfsetroundjoin%
\definecolor{currentfill}{rgb}{0.168627,0.611765,0.411765}%
\pgfsetfillcolor{currentfill}%
\pgfsetfillopacity{0.180000}%
\pgfsetlinewidth{0.000000pt}%
\definecolor{currentstroke}{rgb}{0.000000,0.000000,0.000000}%
\pgfsetstrokecolor{currentstroke}%
\pgfsetdash{}{0pt}%
\pgfpathmoveto{\pgfqpoint{1.263757in}{0.890672in}}%
\pgfpathlineto{\pgfqpoint{1.279861in}{0.884370in}}%
\pgfpathlineto{\pgfqpoint{1.280052in}{0.907576in}}%
\pgfpathlineto{\pgfqpoint{1.263928in}{0.913982in}}%
\pgfpathlineto{\pgfqpoint{1.263757in}{0.890672in}}%
\pgfpathclose%
\pgfusepath{fill}%
\end{pgfscope}%
\begin{pgfscope}%
\pgfpathrectangle{\pgfqpoint{0.407592in}{0.222904in}}{\pgfqpoint{1.411163in}{1.411163in}}%
\pgfusepath{clip}%
\pgfsetbuttcap%
\pgfsetroundjoin%
\definecolor{currentfill}{rgb}{0.168627,0.611765,0.411765}%
\pgfsetfillcolor{currentfill}%
\pgfsetfillopacity{0.180000}%
\pgfsetlinewidth{0.000000pt}%
\definecolor{currentstroke}{rgb}{0.000000,0.000000,0.000000}%
\pgfsetstrokecolor{currentstroke}%
\pgfsetdash{}{0pt}%
\pgfpathmoveto{\pgfqpoint{1.314030in}{1.107142in}}%
\pgfpathlineto{\pgfqpoint{1.329964in}{1.102641in}}%
\pgfpathlineto{\pgfqpoint{1.330221in}{1.126147in}}%
\pgfpathlineto{\pgfqpoint{1.314266in}{1.130749in}}%
\pgfpathlineto{\pgfqpoint{1.314030in}{1.107142in}}%
\pgfpathclose%
\pgfusepath{fill}%
\end{pgfscope}%
\begin{pgfscope}%
\pgfpathrectangle{\pgfqpoint{0.407592in}{0.222904in}}{\pgfqpoint{1.411163in}{1.411163in}}%
\pgfusepath{clip}%
\pgfsetbuttcap%
\pgfsetroundjoin%
\definecolor{currentfill}{rgb}{0.168627,0.611765,0.411765}%
\pgfsetfillcolor{currentfill}%
\pgfsetfillopacity{0.180000}%
\pgfsetlinewidth{0.000000pt}%
\definecolor{currentstroke}{rgb}{0.000000,0.000000,0.000000}%
\pgfsetstrokecolor{currentstroke}%
\pgfsetdash{}{0pt}%
\pgfpathmoveto{\pgfqpoint{1.297329in}{1.041911in}}%
\pgfpathlineto{\pgfqpoint{1.313323in}{1.036690in}}%
\pgfpathlineto{\pgfqpoint{1.313558in}{1.060113in}}%
\pgfpathlineto{\pgfqpoint{1.297544in}{1.065437in}}%
\pgfpathlineto{\pgfqpoint{1.297329in}{1.041911in}}%
\pgfpathclose%
\pgfusepath{fill}%
\end{pgfscope}%
\begin{pgfscope}%
\pgfpathrectangle{\pgfqpoint{0.407592in}{0.222904in}}{\pgfqpoint{1.411163in}{1.411163in}}%
\pgfusepath{clip}%
\pgfsetbuttcap%
\pgfsetroundjoin%
\definecolor{currentfill}{rgb}{0.168627,0.611765,0.411765}%
\pgfsetfillcolor{currentfill}%
\pgfsetfillopacity{0.180000}%
\pgfsetlinewidth{0.000000pt}%
\definecolor{currentstroke}{rgb}{0.000000,0.000000,0.000000}%
\pgfsetstrokecolor{currentstroke}%
\pgfsetdash{}{0pt}%
\pgfpathmoveto{\pgfqpoint{1.280629in}{0.977553in}}%
\pgfpathlineto{\pgfqpoint{1.296687in}{0.971699in}}%
\pgfpathlineto{\pgfqpoint{1.296901in}{0.995042in}}%
\pgfpathlineto{\pgfqpoint{1.280823in}{1.001001in}}%
\pgfpathlineto{\pgfqpoint{1.280629in}{0.977553in}}%
\pgfpathclose%
\pgfusepath{fill}%
\end{pgfscope}%
\begin{pgfscope}%
\pgfpathrectangle{\pgfqpoint{0.407592in}{0.222904in}}{\pgfqpoint{1.411163in}{1.411163in}}%
\pgfusepath{clip}%
\pgfsetbuttcap%
\pgfsetroundjoin%
\definecolor{currentfill}{rgb}{0.168627,0.611765,0.411765}%
\pgfsetfillcolor{currentfill}%
\pgfsetfillopacity{0.180000}%
\pgfsetlinewidth{0.000000pt}%
\definecolor{currentstroke}{rgb}{0.000000,0.000000,0.000000}%
\pgfsetstrokecolor{currentstroke}%
\pgfsetdash{}{0pt}%
\pgfpathmoveto{\pgfqpoint{1.263928in}{0.913982in}}%
\pgfpathlineto{\pgfqpoint{1.280052in}{0.907576in}}%
\pgfpathlineto{\pgfqpoint{1.280244in}{0.930841in}}%
\pgfpathlineto{\pgfqpoint{1.264099in}{0.937353in}}%
\pgfpathlineto{\pgfqpoint{1.263928in}{0.913982in}}%
\pgfpathclose%
\pgfusepath{fill}%
\end{pgfscope}%
\begin{pgfscope}%
\pgfpathrectangle{\pgfqpoint{0.407592in}{0.222904in}}{\pgfqpoint{1.411163in}{1.411163in}}%
\pgfusepath{clip}%
\pgfsetbuttcap%
\pgfsetroundjoin%
\definecolor{currentfill}{rgb}{0.168627,0.611765,0.411765}%
\pgfsetfillcolor{currentfill}%
\pgfsetfillopacity{0.180000}%
\pgfsetlinewidth{0.000000pt}%
\definecolor{currentstroke}{rgb}{0.000000,0.000000,0.000000}%
\pgfsetstrokecolor{currentstroke}%
\pgfsetdash{}{0pt}%
\pgfpathmoveto{\pgfqpoint{1.314266in}{1.130749in}}%
\pgfpathlineto{\pgfqpoint{1.330221in}{1.126147in}}%
\pgfpathlineto{\pgfqpoint{1.330478in}{1.149713in}}%
\pgfpathlineto{\pgfqpoint{1.314504in}{1.154417in}}%
\pgfpathlineto{\pgfqpoint{1.314266in}{1.130749in}}%
\pgfpathclose%
\pgfusepath{fill}%
\end{pgfscope}%
\begin{pgfscope}%
\pgfpathrectangle{\pgfqpoint{0.407592in}{0.222904in}}{\pgfqpoint{1.411163in}{1.411163in}}%
\pgfusepath{clip}%
\pgfsetbuttcap%
\pgfsetroundjoin%
\definecolor{currentfill}{rgb}{0.168627,0.611765,0.411765}%
\pgfsetfillcolor{currentfill}%
\pgfsetfillopacity{0.180000}%
\pgfsetlinewidth{0.000000pt}%
\definecolor{currentstroke}{rgb}{0.000000,0.000000,0.000000}%
\pgfsetstrokecolor{currentstroke}%
\pgfsetdash{}{0pt}%
\pgfpathmoveto{\pgfqpoint{1.297544in}{1.065437in}}%
\pgfpathlineto{\pgfqpoint{1.313558in}{1.060113in}}%
\pgfpathlineto{\pgfqpoint{1.313794in}{1.083597in}}%
\pgfpathlineto{\pgfqpoint{1.297759in}{1.089024in}}%
\pgfpathlineto{\pgfqpoint{1.297544in}{1.065437in}}%
\pgfpathclose%
\pgfusepath{fill}%
\end{pgfscope}%
\begin{pgfscope}%
\pgfpathrectangle{\pgfqpoint{0.407592in}{0.222904in}}{\pgfqpoint{1.411163in}{1.411163in}}%
\pgfusepath{clip}%
\pgfsetbuttcap%
\pgfsetroundjoin%
\definecolor{currentfill}{rgb}{0.168627,0.611765,0.411765}%
\pgfsetfillcolor{currentfill}%
\pgfsetfillopacity{0.180000}%
\pgfsetlinewidth{0.000000pt}%
\definecolor{currentstroke}{rgb}{0.000000,0.000000,0.000000}%
\pgfsetstrokecolor{currentstroke}%
\pgfsetdash{}{0pt}%
\pgfpathmoveto{\pgfqpoint{1.280823in}{1.001001in}}%
\pgfpathlineto{\pgfqpoint{1.296901in}{0.995042in}}%
\pgfpathlineto{\pgfqpoint{1.297114in}{1.018446in}}%
\pgfpathlineto{\pgfqpoint{1.281016in}{1.024509in}}%
\pgfpathlineto{\pgfqpoint{1.280823in}{1.001001in}}%
\pgfpathclose%
\pgfusepath{fill}%
\end{pgfscope}%
\begin{pgfscope}%
\pgfpathrectangle{\pgfqpoint{0.407592in}{0.222904in}}{\pgfqpoint{1.411163in}{1.411163in}}%
\pgfusepath{clip}%
\pgfsetbuttcap%
\pgfsetroundjoin%
\definecolor{currentfill}{rgb}{0.168627,0.611765,0.411765}%
\pgfsetfillcolor{currentfill}%
\pgfsetfillopacity{0.180000}%
\pgfsetlinewidth{0.000000pt}%
\definecolor{currentstroke}{rgb}{0.000000,0.000000,0.000000}%
\pgfsetstrokecolor{currentstroke}%
\pgfsetdash{}{0pt}%
\pgfpathmoveto{\pgfqpoint{1.264099in}{0.937353in}}%
\pgfpathlineto{\pgfqpoint{1.280244in}{0.930841in}}%
\pgfpathlineto{\pgfqpoint{1.280436in}{0.954167in}}%
\pgfpathlineto{\pgfqpoint{1.264271in}{0.960785in}}%
\pgfpathlineto{\pgfqpoint{1.264099in}{0.937353in}}%
\pgfpathclose%
\pgfusepath{fill}%
\end{pgfscope}%
\begin{pgfscope}%
\pgfpathrectangle{\pgfqpoint{0.407592in}{0.222904in}}{\pgfqpoint{1.411163in}{1.411163in}}%
\pgfusepath{clip}%
\pgfsetbuttcap%
\pgfsetroundjoin%
\definecolor{currentfill}{rgb}{0.168627,0.611765,0.411765}%
\pgfsetfillcolor{currentfill}%
\pgfsetfillopacity{0.180000}%
\pgfsetlinewidth{0.000000pt}%
\definecolor{currentstroke}{rgb}{0.000000,0.000000,0.000000}%
\pgfsetstrokecolor{currentstroke}%
\pgfsetdash{}{0pt}%
\pgfpathmoveto{\pgfqpoint{1.314504in}{1.154417in}}%
\pgfpathlineto{\pgfqpoint{1.330478in}{1.149713in}}%
\pgfpathlineto{\pgfqpoint{1.330735in}{1.173342in}}%
\pgfpathlineto{\pgfqpoint{1.314742in}{1.178147in}}%
\pgfpathlineto{\pgfqpoint{1.314504in}{1.154417in}}%
\pgfpathclose%
\pgfusepath{fill}%
\end{pgfscope}%
\begin{pgfscope}%
\pgfpathrectangle{\pgfqpoint{0.407592in}{0.222904in}}{\pgfqpoint{1.411163in}{1.411163in}}%
\pgfusepath{clip}%
\pgfsetbuttcap%
\pgfsetroundjoin%
\definecolor{currentfill}{rgb}{0.168627,0.611765,0.411765}%
\pgfsetfillcolor{currentfill}%
\pgfsetfillopacity{0.180000}%
\pgfsetlinewidth{0.000000pt}%
\definecolor{currentstroke}{rgb}{0.000000,0.000000,0.000000}%
\pgfsetstrokecolor{currentstroke}%
\pgfsetdash{}{0pt}%
\pgfpathmoveto{\pgfqpoint{1.297759in}{1.089024in}}%
\pgfpathlineto{\pgfqpoint{1.313794in}{1.083597in}}%
\pgfpathlineto{\pgfqpoint{1.314030in}{1.107142in}}%
\pgfpathlineto{\pgfqpoint{1.297975in}{1.112673in}}%
\pgfpathlineto{\pgfqpoint{1.297759in}{1.089024in}}%
\pgfpathclose%
\pgfusepath{fill}%
\end{pgfscope}%
\begin{pgfscope}%
\pgfpathrectangle{\pgfqpoint{0.407592in}{0.222904in}}{\pgfqpoint{1.411163in}{1.411163in}}%
\pgfusepath{clip}%
\pgfsetbuttcap%
\pgfsetroundjoin%
\definecolor{currentfill}{rgb}{0.168627,0.611765,0.411765}%
\pgfsetfillcolor{currentfill}%
\pgfsetfillopacity{0.180000}%
\pgfsetlinewidth{0.000000pt}%
\definecolor{currentstroke}{rgb}{0.000000,0.000000,0.000000}%
\pgfsetstrokecolor{currentstroke}%
\pgfsetdash{}{0pt}%
\pgfpathmoveto{\pgfqpoint{1.281016in}{1.024509in}}%
\pgfpathlineto{\pgfqpoint{1.297114in}{1.018446in}}%
\pgfpathlineto{\pgfqpoint{1.297329in}{1.041911in}}%
\pgfpathlineto{\pgfqpoint{1.281211in}{1.048079in}}%
\pgfpathlineto{\pgfqpoint{1.281016in}{1.024509in}}%
\pgfpathclose%
\pgfusepath{fill}%
\end{pgfscope}%
\begin{pgfscope}%
\pgfpathrectangle{\pgfqpoint{0.407592in}{0.222904in}}{\pgfqpoint{1.411163in}{1.411163in}}%
\pgfusepath{clip}%
\pgfsetbuttcap%
\pgfsetroundjoin%
\definecolor{currentfill}{rgb}{0.168627,0.611765,0.411765}%
\pgfsetfillcolor{currentfill}%
\pgfsetfillopacity{0.180000}%
\pgfsetlinewidth{0.000000pt}%
\definecolor{currentstroke}{rgb}{0.000000,0.000000,0.000000}%
\pgfsetstrokecolor{currentstroke}%
\pgfsetdash{}{0pt}%
\pgfpathmoveto{\pgfqpoint{1.264271in}{0.960785in}}%
\pgfpathlineto{\pgfqpoint{1.280436in}{0.954167in}}%
\pgfpathlineto{\pgfqpoint{1.280629in}{0.977553in}}%
\pgfpathlineto{\pgfqpoint{1.264444in}{0.984279in}}%
\pgfpathlineto{\pgfqpoint{1.264271in}{0.960785in}}%
\pgfpathclose%
\pgfusepath{fill}%
\end{pgfscope}%
\begin{pgfscope}%
\pgfpathrectangle{\pgfqpoint{0.407592in}{0.222904in}}{\pgfqpoint{1.411163in}{1.411163in}}%
\pgfusepath{clip}%
\pgfsetbuttcap%
\pgfsetroundjoin%
\definecolor{currentfill}{rgb}{0.168627,0.611765,0.411765}%
\pgfsetfillcolor{currentfill}%
\pgfsetfillopacity{0.180000}%
\pgfsetlinewidth{0.000000pt}%
\definecolor{currentstroke}{rgb}{0.000000,0.000000,0.000000}%
\pgfsetstrokecolor{currentstroke}%
\pgfsetdash{}{0pt}%
\pgfpathmoveto{\pgfqpoint{1.247522in}{0.897771in}}%
\pgfpathlineto{\pgfqpoint{1.263757in}{0.890672in}}%
\pgfpathlineto{\pgfqpoint{1.263928in}{0.913982in}}%
\pgfpathlineto{\pgfqpoint{1.247673in}{0.921190in}}%
\pgfpathlineto{\pgfqpoint{1.247522in}{0.897771in}}%
\pgfpathclose%
\pgfusepath{fill}%
\end{pgfscope}%
\begin{pgfscope}%
\pgfpathrectangle{\pgfqpoint{0.407592in}{0.222904in}}{\pgfqpoint{1.411163in}{1.411163in}}%
\pgfusepath{clip}%
\pgfsetbuttcap%
\pgfsetroundjoin%
\definecolor{currentfill}{rgb}{0.168627,0.611765,0.411765}%
\pgfsetfillcolor{currentfill}%
\pgfsetfillopacity{0.180000}%
\pgfsetlinewidth{0.000000pt}%
\definecolor{currentstroke}{rgb}{0.000000,0.000000,0.000000}%
\pgfsetstrokecolor{currentstroke}%
\pgfsetdash{}{0pt}%
\pgfpathmoveto{\pgfqpoint{1.297975in}{1.112673in}}%
\pgfpathlineto{\pgfqpoint{1.314030in}{1.107142in}}%
\pgfpathlineto{\pgfqpoint{1.314266in}{1.130749in}}%
\pgfpathlineto{\pgfqpoint{1.298192in}{1.136383in}}%
\pgfpathlineto{\pgfqpoint{1.297975in}{1.112673in}}%
\pgfpathclose%
\pgfusepath{fill}%
\end{pgfscope}%
\begin{pgfscope}%
\pgfpathrectangle{\pgfqpoint{0.407592in}{0.222904in}}{\pgfqpoint{1.411163in}{1.411163in}}%
\pgfusepath{clip}%
\pgfsetbuttcap%
\pgfsetroundjoin%
\definecolor{currentfill}{rgb}{0.168627,0.611765,0.411765}%
\pgfsetfillcolor{currentfill}%
\pgfsetfillopacity{0.180000}%
\pgfsetlinewidth{0.000000pt}%
\definecolor{currentstroke}{rgb}{0.000000,0.000000,0.000000}%
\pgfsetstrokecolor{currentstroke}%
\pgfsetdash{}{0pt}%
\pgfpathmoveto{\pgfqpoint{1.281211in}{1.048079in}}%
\pgfpathlineto{\pgfqpoint{1.297329in}{1.041911in}}%
\pgfpathlineto{\pgfqpoint{1.297544in}{1.065437in}}%
\pgfpathlineto{\pgfqpoint{1.281405in}{1.071711in}}%
\pgfpathlineto{\pgfqpoint{1.281211in}{1.048079in}}%
\pgfpathclose%
\pgfusepath{fill}%
\end{pgfscope}%
\begin{pgfscope}%
\pgfpathrectangle{\pgfqpoint{0.407592in}{0.222904in}}{\pgfqpoint{1.411163in}{1.411163in}}%
\pgfusepath{clip}%
\pgfsetbuttcap%
\pgfsetroundjoin%
\definecolor{currentfill}{rgb}{0.168627,0.611765,0.411765}%
\pgfsetfillcolor{currentfill}%
\pgfsetfillopacity{0.180000}%
\pgfsetlinewidth{0.000000pt}%
\definecolor{currentstroke}{rgb}{0.000000,0.000000,0.000000}%
\pgfsetstrokecolor{currentstroke}%
\pgfsetdash{}{0pt}%
\pgfpathmoveto{\pgfqpoint{1.264444in}{0.984279in}}%
\pgfpathlineto{\pgfqpoint{1.280629in}{0.977553in}}%
\pgfpathlineto{\pgfqpoint{1.280823in}{1.001001in}}%
\pgfpathlineto{\pgfqpoint{1.264617in}{1.007834in}}%
\pgfpathlineto{\pgfqpoint{1.264444in}{0.984279in}}%
\pgfpathclose%
\pgfusepath{fill}%
\end{pgfscope}%
\begin{pgfscope}%
\pgfpathrectangle{\pgfqpoint{0.407592in}{0.222904in}}{\pgfqpoint{1.411163in}{1.411163in}}%
\pgfusepath{clip}%
\pgfsetbuttcap%
\pgfsetroundjoin%
\definecolor{currentfill}{rgb}{0.168627,0.611765,0.411765}%
\pgfsetfillcolor{currentfill}%
\pgfsetfillopacity{0.180000}%
\pgfsetlinewidth{0.000000pt}%
\definecolor{currentstroke}{rgb}{0.000000,0.000000,0.000000}%
\pgfsetstrokecolor{currentstroke}%
\pgfsetdash{}{0pt}%
\pgfpathmoveto{\pgfqpoint{1.247673in}{0.921190in}}%
\pgfpathlineto{\pgfqpoint{1.263928in}{0.913982in}}%
\pgfpathlineto{\pgfqpoint{1.264099in}{0.937353in}}%
\pgfpathlineto{\pgfqpoint{1.247824in}{0.944669in}}%
\pgfpathlineto{\pgfqpoint{1.247673in}{0.921190in}}%
\pgfpathclose%
\pgfusepath{fill}%
\end{pgfscope}%
\begin{pgfscope}%
\pgfpathrectangle{\pgfqpoint{0.407592in}{0.222904in}}{\pgfqpoint{1.411163in}{1.411163in}}%
\pgfusepath{clip}%
\pgfsetbuttcap%
\pgfsetroundjoin%
\definecolor{currentfill}{rgb}{0.168627,0.611765,0.411765}%
\pgfsetfillcolor{currentfill}%
\pgfsetfillopacity{0.180000}%
\pgfsetlinewidth{0.000000pt}%
\definecolor{currentstroke}{rgb}{0.000000,0.000000,0.000000}%
\pgfsetstrokecolor{currentstroke}%
\pgfsetdash{}{0pt}%
\pgfpathmoveto{\pgfqpoint{1.298192in}{1.136383in}}%
\pgfpathlineto{\pgfqpoint{1.314266in}{1.130749in}}%
\pgfpathlineto{\pgfqpoint{1.314504in}{1.154417in}}%
\pgfpathlineto{\pgfqpoint{1.298409in}{1.160156in}}%
\pgfpathlineto{\pgfqpoint{1.298192in}{1.136383in}}%
\pgfpathclose%
\pgfusepath{fill}%
\end{pgfscope}%
\begin{pgfscope}%
\pgfpathrectangle{\pgfqpoint{0.407592in}{0.222904in}}{\pgfqpoint{1.411163in}{1.411163in}}%
\pgfusepath{clip}%
\pgfsetbuttcap%
\pgfsetroundjoin%
\definecolor{currentfill}{rgb}{0.168627,0.611765,0.411765}%
\pgfsetfillcolor{currentfill}%
\pgfsetfillopacity{0.180000}%
\pgfsetlinewidth{0.000000pt}%
\definecolor{currentstroke}{rgb}{0.000000,0.000000,0.000000}%
\pgfsetstrokecolor{currentstroke}%
\pgfsetdash{}{0pt}%
\pgfpathmoveto{\pgfqpoint{1.281405in}{1.071711in}}%
\pgfpathlineto{\pgfqpoint{1.297544in}{1.065437in}}%
\pgfpathlineto{\pgfqpoint{1.297759in}{1.089024in}}%
\pgfpathlineto{\pgfqpoint{1.281601in}{1.095405in}}%
\pgfpathlineto{\pgfqpoint{1.281405in}{1.071711in}}%
\pgfpathclose%
\pgfusepath{fill}%
\end{pgfscope}%
\begin{pgfscope}%
\pgfpathrectangle{\pgfqpoint{0.407592in}{0.222904in}}{\pgfqpoint{1.411163in}{1.411163in}}%
\pgfusepath{clip}%
\pgfsetbuttcap%
\pgfsetroundjoin%
\definecolor{currentfill}{rgb}{0.168627,0.611765,0.411765}%
\pgfsetfillcolor{currentfill}%
\pgfsetfillopacity{0.180000}%
\pgfsetlinewidth{0.000000pt}%
\definecolor{currentstroke}{rgb}{0.000000,0.000000,0.000000}%
\pgfsetstrokecolor{currentstroke}%
\pgfsetdash{}{0pt}%
\pgfpathmoveto{\pgfqpoint{1.264617in}{1.007834in}}%
\pgfpathlineto{\pgfqpoint{1.280823in}{1.001001in}}%
\pgfpathlineto{\pgfqpoint{1.281016in}{1.024509in}}%
\pgfpathlineto{\pgfqpoint{1.264790in}{1.031450in}}%
\pgfpathlineto{\pgfqpoint{1.264617in}{1.007834in}}%
\pgfpathclose%
\pgfusepath{fill}%
\end{pgfscope}%
\begin{pgfscope}%
\pgfpathrectangle{\pgfqpoint{0.407592in}{0.222904in}}{\pgfqpoint{1.411163in}{1.411163in}}%
\pgfusepath{clip}%
\pgfsetbuttcap%
\pgfsetroundjoin%
\definecolor{currentfill}{rgb}{0.168627,0.611765,0.411765}%
\pgfsetfillcolor{currentfill}%
\pgfsetfillopacity{0.180000}%
\pgfsetlinewidth{0.000000pt}%
\definecolor{currentstroke}{rgb}{0.000000,0.000000,0.000000}%
\pgfsetstrokecolor{currentstroke}%
\pgfsetdash{}{0pt}%
\pgfpathmoveto{\pgfqpoint{1.247824in}{0.944669in}}%
\pgfpathlineto{\pgfqpoint{1.264099in}{0.937353in}}%
\pgfpathlineto{\pgfqpoint{1.264271in}{0.960785in}}%
\pgfpathlineto{\pgfqpoint{1.247975in}{0.968211in}}%
\pgfpathlineto{\pgfqpoint{1.247824in}{0.944669in}}%
\pgfpathclose%
\pgfusepath{fill}%
\end{pgfscope}%
\begin{pgfscope}%
\pgfpathrectangle{\pgfqpoint{0.407592in}{0.222904in}}{\pgfqpoint{1.411163in}{1.411163in}}%
\pgfusepath{clip}%
\pgfsetbuttcap%
\pgfsetroundjoin%
\definecolor{currentfill}{rgb}{0.168627,0.611765,0.411765}%
\pgfsetfillcolor{currentfill}%
\pgfsetfillopacity{0.180000}%
\pgfsetlinewidth{0.000000pt}%
\definecolor{currentstroke}{rgb}{0.000000,0.000000,0.000000}%
\pgfsetstrokecolor{currentstroke}%
\pgfsetdash{}{0pt}%
\pgfpathmoveto{\pgfqpoint{1.298409in}{1.160156in}}%
\pgfpathlineto{\pgfqpoint{1.314504in}{1.154417in}}%
\pgfpathlineto{\pgfqpoint{1.314742in}{1.178147in}}%
\pgfpathlineto{\pgfqpoint{1.298627in}{1.183991in}}%
\pgfpathlineto{\pgfqpoint{1.298409in}{1.160156in}}%
\pgfpathclose%
\pgfusepath{fill}%
\end{pgfscope}%
\begin{pgfscope}%
\pgfpathrectangle{\pgfqpoint{0.407592in}{0.222904in}}{\pgfqpoint{1.411163in}{1.411163in}}%
\pgfusepath{clip}%
\pgfsetbuttcap%
\pgfsetroundjoin%
\definecolor{currentfill}{rgb}{0.168627,0.611765,0.411765}%
\pgfsetfillcolor{currentfill}%
\pgfsetfillopacity{0.180000}%
\pgfsetlinewidth{0.000000pt}%
\definecolor{currentstroke}{rgb}{0.000000,0.000000,0.000000}%
\pgfsetstrokecolor{currentstroke}%
\pgfsetdash{}{0pt}%
\pgfpathmoveto{\pgfqpoint{1.281601in}{1.095405in}}%
\pgfpathlineto{\pgfqpoint{1.297759in}{1.089024in}}%
\pgfpathlineto{\pgfqpoint{1.297975in}{1.112673in}}%
\pgfpathlineto{\pgfqpoint{1.281797in}{1.119160in}}%
\pgfpathlineto{\pgfqpoint{1.281601in}{1.095405in}}%
\pgfpathclose%
\pgfusepath{fill}%
\end{pgfscope}%
\begin{pgfscope}%
\pgfpathrectangle{\pgfqpoint{0.407592in}{0.222904in}}{\pgfqpoint{1.411163in}{1.411163in}}%
\pgfusepath{clip}%
\pgfsetbuttcap%
\pgfsetroundjoin%
\definecolor{currentfill}{rgb}{0.168627,0.611765,0.411765}%
\pgfsetfillcolor{currentfill}%
\pgfsetfillopacity{0.180000}%
\pgfsetlinewidth{0.000000pt}%
\definecolor{currentstroke}{rgb}{0.000000,0.000000,0.000000}%
\pgfsetstrokecolor{currentstroke}%
\pgfsetdash{}{0pt}%
\pgfpathmoveto{\pgfqpoint{1.264790in}{1.031450in}}%
\pgfpathlineto{\pgfqpoint{1.281016in}{1.024509in}}%
\pgfpathlineto{\pgfqpoint{1.281211in}{1.048079in}}%
\pgfpathlineto{\pgfqpoint{1.264964in}{1.055128in}}%
\pgfpathlineto{\pgfqpoint{1.264790in}{1.031450in}}%
\pgfpathclose%
\pgfusepath{fill}%
\end{pgfscope}%
\begin{pgfscope}%
\pgfpathrectangle{\pgfqpoint{0.407592in}{0.222904in}}{\pgfqpoint{1.411163in}{1.411163in}}%
\pgfusepath{clip}%
\pgfsetbuttcap%
\pgfsetroundjoin%
\definecolor{currentfill}{rgb}{0.168627,0.611765,0.411765}%
\pgfsetfillcolor{currentfill}%
\pgfsetfillopacity{0.180000}%
\pgfsetlinewidth{0.000000pt}%
\definecolor{currentstroke}{rgb}{0.000000,0.000000,0.000000}%
\pgfsetstrokecolor{currentstroke}%
\pgfsetdash{}{0pt}%
\pgfpathmoveto{\pgfqpoint{1.247975in}{0.968211in}}%
\pgfpathlineto{\pgfqpoint{1.264271in}{0.960785in}}%
\pgfpathlineto{\pgfqpoint{1.264444in}{0.984279in}}%
\pgfpathlineto{\pgfqpoint{1.248127in}{0.991814in}}%
\pgfpathlineto{\pgfqpoint{1.247975in}{0.968211in}}%
\pgfpathclose%
\pgfusepath{fill}%
\end{pgfscope}%
\begin{pgfscope}%
\pgfpathrectangle{\pgfqpoint{0.407592in}{0.222904in}}{\pgfqpoint{1.411163in}{1.411163in}}%
\pgfusepath{clip}%
\pgfsetbuttcap%
\pgfsetroundjoin%
\definecolor{currentfill}{rgb}{0.168627,0.611765,0.411765}%
\pgfsetfillcolor{currentfill}%
\pgfsetfillopacity{0.180000}%
\pgfsetlinewidth{0.000000pt}%
\definecolor{currentstroke}{rgb}{0.000000,0.000000,0.000000}%
\pgfsetstrokecolor{currentstroke}%
\pgfsetdash{}{0pt}%
\pgfpathmoveto{\pgfqpoint{1.281797in}{1.119160in}}%
\pgfpathlineto{\pgfqpoint{1.297975in}{1.112673in}}%
\pgfpathlineto{\pgfqpoint{1.298192in}{1.136383in}}%
\pgfpathlineto{\pgfqpoint{1.281993in}{1.142978in}}%
\pgfpathlineto{\pgfqpoint{1.281797in}{1.119160in}}%
\pgfpathclose%
\pgfusepath{fill}%
\end{pgfscope}%
\begin{pgfscope}%
\pgfpathrectangle{\pgfqpoint{0.407592in}{0.222904in}}{\pgfqpoint{1.411163in}{1.411163in}}%
\pgfusepath{clip}%
\pgfsetbuttcap%
\pgfsetroundjoin%
\definecolor{currentfill}{rgb}{0.168627,0.611765,0.411765}%
\pgfsetfillcolor{currentfill}%
\pgfsetfillopacity{0.180000}%
\pgfsetlinewidth{0.000000pt}%
\definecolor{currentstroke}{rgb}{0.000000,0.000000,0.000000}%
\pgfsetstrokecolor{currentstroke}%
\pgfsetdash{}{0pt}%
\pgfpathmoveto{\pgfqpoint{1.264964in}{1.055128in}}%
\pgfpathlineto{\pgfqpoint{1.281211in}{1.048079in}}%
\pgfpathlineto{\pgfqpoint{1.281405in}{1.071711in}}%
\pgfpathlineto{\pgfqpoint{1.265138in}{1.078869in}}%
\pgfpathlineto{\pgfqpoint{1.264964in}{1.055128in}}%
\pgfpathclose%
\pgfusepath{fill}%
\end{pgfscope}%
\begin{pgfscope}%
\pgfpathrectangle{\pgfqpoint{0.407592in}{0.222904in}}{\pgfqpoint{1.411163in}{1.411163in}}%
\pgfusepath{clip}%
\pgfsetbuttcap%
\pgfsetroundjoin%
\definecolor{currentfill}{rgb}{0.168627,0.611765,0.411765}%
\pgfsetfillcolor{currentfill}%
\pgfsetfillopacity{0.180000}%
\pgfsetlinewidth{0.000000pt}%
\definecolor{currentstroke}{rgb}{0.000000,0.000000,0.000000}%
\pgfsetstrokecolor{currentstroke}%
\pgfsetdash{}{0pt}%
\pgfpathmoveto{\pgfqpoint{1.248127in}{0.991814in}}%
\pgfpathlineto{\pgfqpoint{1.264444in}{0.984279in}}%
\pgfpathlineto{\pgfqpoint{1.264617in}{1.007834in}}%
\pgfpathlineto{\pgfqpoint{1.248279in}{1.015479in}}%
\pgfpathlineto{\pgfqpoint{1.248127in}{0.991814in}}%
\pgfpathclose%
\pgfusepath{fill}%
\end{pgfscope}%
\begin{pgfscope}%
\pgfpathrectangle{\pgfqpoint{0.407592in}{0.222904in}}{\pgfqpoint{1.411163in}{1.411163in}}%
\pgfusepath{clip}%
\pgfsetbuttcap%
\pgfsetroundjoin%
\definecolor{currentfill}{rgb}{0.168627,0.611765,0.411765}%
\pgfsetfillcolor{currentfill}%
\pgfsetfillopacity{0.180000}%
\pgfsetlinewidth{0.000000pt}%
\definecolor{currentstroke}{rgb}{0.000000,0.000000,0.000000}%
\pgfsetstrokecolor{currentstroke}%
\pgfsetdash{}{0pt}%
\pgfpathmoveto{\pgfqpoint{1.281993in}{1.142978in}}%
\pgfpathlineto{\pgfqpoint{1.298192in}{1.136383in}}%
\pgfpathlineto{\pgfqpoint{1.298409in}{1.160156in}}%
\pgfpathlineto{\pgfqpoint{1.282190in}{1.166859in}}%
\pgfpathlineto{\pgfqpoint{1.281993in}{1.142978in}}%
\pgfpathclose%
\pgfusepath{fill}%
\end{pgfscope}%
\begin{pgfscope}%
\pgfpathrectangle{\pgfqpoint{0.407592in}{0.222904in}}{\pgfqpoint{1.411163in}{1.411163in}}%
\pgfusepath{clip}%
\pgfsetbuttcap%
\pgfsetroundjoin%
\definecolor{currentfill}{rgb}{0.168627,0.611765,0.411765}%
\pgfsetfillcolor{currentfill}%
\pgfsetfillopacity{0.180000}%
\pgfsetlinewidth{0.000000pt}%
\definecolor{currentstroke}{rgb}{0.000000,0.000000,0.000000}%
\pgfsetstrokecolor{currentstroke}%
\pgfsetdash{}{0pt}%
\pgfpathmoveto{\pgfqpoint{1.265138in}{1.078869in}}%
\pgfpathlineto{\pgfqpoint{1.281405in}{1.071711in}}%
\pgfpathlineto{\pgfqpoint{1.281601in}{1.095405in}}%
\pgfpathlineto{\pgfqpoint{1.265313in}{1.102672in}}%
\pgfpathlineto{\pgfqpoint{1.265138in}{1.078869in}}%
\pgfpathclose%
\pgfusepath{fill}%
\end{pgfscope}%
\begin{pgfscope}%
\pgfpathrectangle{\pgfqpoint{0.407592in}{0.222904in}}{\pgfqpoint{1.411163in}{1.411163in}}%
\pgfusepath{clip}%
\pgfsetbuttcap%
\pgfsetroundjoin%
\definecolor{currentfill}{rgb}{0.168627,0.611765,0.411765}%
\pgfsetfillcolor{currentfill}%
\pgfsetfillopacity{0.180000}%
\pgfsetlinewidth{0.000000pt}%
\definecolor{currentstroke}{rgb}{0.000000,0.000000,0.000000}%
\pgfsetstrokecolor{currentstroke}%
\pgfsetdash{}{0pt}%
\pgfpathmoveto{\pgfqpoint{1.248279in}{1.015479in}}%
\pgfpathlineto{\pgfqpoint{1.264617in}{1.007834in}}%
\pgfpathlineto{\pgfqpoint{1.264790in}{1.031450in}}%
\pgfpathlineto{\pgfqpoint{1.248432in}{1.039206in}}%
\pgfpathlineto{\pgfqpoint{1.248279in}{1.015479in}}%
\pgfpathclose%
\pgfusepath{fill}%
\end{pgfscope}%
\begin{pgfscope}%
\pgfpathrectangle{\pgfqpoint{0.407592in}{0.222904in}}{\pgfqpoint{1.411163in}{1.411163in}}%
\pgfusepath{clip}%
\pgfsetbuttcap%
\pgfsetroundjoin%
\definecolor{currentfill}{rgb}{0.168627,0.611765,0.411765}%
\pgfsetfillcolor{currentfill}%
\pgfsetfillopacity{0.180000}%
\pgfsetlinewidth{0.000000pt}%
\definecolor{currentstroke}{rgb}{0.000000,0.000000,0.000000}%
\pgfsetstrokecolor{currentstroke}%
\pgfsetdash{}{0pt}%
\pgfpathmoveto{\pgfqpoint{1.282190in}{1.166859in}}%
\pgfpathlineto{\pgfqpoint{1.298409in}{1.160156in}}%
\pgfpathlineto{\pgfqpoint{1.298627in}{1.183991in}}%
\pgfpathlineto{\pgfqpoint{1.282387in}{1.190802in}}%
\pgfpathlineto{\pgfqpoint{1.282190in}{1.166859in}}%
\pgfpathclose%
\pgfusepath{fill}%
\end{pgfscope}%
\begin{pgfscope}%
\pgfpathrectangle{\pgfqpoint{0.407592in}{0.222904in}}{\pgfqpoint{1.411163in}{1.411163in}}%
\pgfusepath{clip}%
\pgfsetbuttcap%
\pgfsetroundjoin%
\definecolor{currentfill}{rgb}{0.168627,0.611765,0.411765}%
\pgfsetfillcolor{currentfill}%
\pgfsetfillopacity{0.180000}%
\pgfsetlinewidth{0.000000pt}%
\definecolor{currentstroke}{rgb}{0.000000,0.000000,0.000000}%
\pgfsetstrokecolor{currentstroke}%
\pgfsetdash{}{0pt}%
\pgfpathmoveto{\pgfqpoint{1.265313in}{1.102672in}}%
\pgfpathlineto{\pgfqpoint{1.281601in}{1.095405in}}%
\pgfpathlineto{\pgfqpoint{1.281797in}{1.119160in}}%
\pgfpathlineto{\pgfqpoint{1.265488in}{1.126538in}}%
\pgfpathlineto{\pgfqpoint{1.265313in}{1.102672in}}%
\pgfpathclose%
\pgfusepath{fill}%
\end{pgfscope}%
\begin{pgfscope}%
\pgfpathrectangle{\pgfqpoint{0.407592in}{0.222904in}}{\pgfqpoint{1.411163in}{1.411163in}}%
\pgfusepath{clip}%
\pgfsetbuttcap%
\pgfsetroundjoin%
\definecolor{currentfill}{rgb}{0.168627,0.611765,0.411765}%
\pgfsetfillcolor{currentfill}%
\pgfsetfillopacity{0.180000}%
\pgfsetlinewidth{0.000000pt}%
\definecolor{currentstroke}{rgb}{0.000000,0.000000,0.000000}%
\pgfsetstrokecolor{currentstroke}%
\pgfsetdash{}{0pt}%
\pgfpathmoveto{\pgfqpoint{1.248432in}{1.039206in}}%
\pgfpathlineto{\pgfqpoint{1.264790in}{1.031450in}}%
\pgfpathlineto{\pgfqpoint{1.264964in}{1.055128in}}%
\pgfpathlineto{\pgfqpoint{1.248585in}{1.062996in}}%
\pgfpathlineto{\pgfqpoint{1.248432in}{1.039206in}}%
\pgfpathclose%
\pgfusepath{fill}%
\end{pgfscope}%
\begin{pgfscope}%
\pgfpathrectangle{\pgfqpoint{0.407592in}{0.222904in}}{\pgfqpoint{1.411163in}{1.411163in}}%
\pgfusepath{clip}%
\pgfsetbuttcap%
\pgfsetroundjoin%
\definecolor{currentfill}{rgb}{0.168627,0.611765,0.411765}%
\pgfsetfillcolor{currentfill}%
\pgfsetfillopacity{0.180000}%
\pgfsetlinewidth{0.000000pt}%
\definecolor{currentstroke}{rgb}{0.000000,0.000000,0.000000}%
\pgfsetstrokecolor{currentstroke}%
\pgfsetdash{}{0pt}%
\pgfpathmoveto{\pgfqpoint{1.265488in}{1.126538in}}%
\pgfpathlineto{\pgfqpoint{1.281797in}{1.119160in}}%
\pgfpathlineto{\pgfqpoint{1.281993in}{1.142978in}}%
\pgfpathlineto{\pgfqpoint{1.265664in}{1.150467in}}%
\pgfpathlineto{\pgfqpoint{1.265488in}{1.126538in}}%
\pgfpathclose%
\pgfusepath{fill}%
\end{pgfscope}%
\begin{pgfscope}%
\pgfpathrectangle{\pgfqpoint{0.407592in}{0.222904in}}{\pgfqpoint{1.411163in}{1.411163in}}%
\pgfusepath{clip}%
\pgfsetbuttcap%
\pgfsetroundjoin%
\definecolor{currentfill}{rgb}{0.168627,0.611765,0.411765}%
\pgfsetfillcolor{currentfill}%
\pgfsetfillopacity{0.180000}%
\pgfsetlinewidth{0.000000pt}%
\definecolor{currentstroke}{rgb}{0.000000,0.000000,0.000000}%
\pgfsetstrokecolor{currentstroke}%
\pgfsetdash{}{0pt}%
\pgfpathmoveto{\pgfqpoint{1.248585in}{1.062996in}}%
\pgfpathlineto{\pgfqpoint{1.264964in}{1.055128in}}%
\pgfpathlineto{\pgfqpoint{1.265138in}{1.078869in}}%
\pgfpathlineto{\pgfqpoint{1.248738in}{1.086849in}}%
\pgfpathlineto{\pgfqpoint{1.248585in}{1.062996in}}%
\pgfpathclose%
\pgfusepath{fill}%
\end{pgfscope}%
\begin{pgfscope}%
\pgfpathrectangle{\pgfqpoint{0.407592in}{0.222904in}}{\pgfqpoint{1.411163in}{1.411163in}}%
\pgfusepath{clip}%
\pgfsetbuttcap%
\pgfsetroundjoin%
\definecolor{currentfill}{rgb}{0.168627,0.611765,0.411765}%
\pgfsetfillcolor{currentfill}%
\pgfsetfillopacity{0.180000}%
\pgfsetlinewidth{0.000000pt}%
\definecolor{currentstroke}{rgb}{0.000000,0.000000,0.000000}%
\pgfsetstrokecolor{currentstroke}%
\pgfsetdash{}{0pt}%
\pgfpathmoveto{\pgfqpoint{1.265664in}{1.150467in}}%
\pgfpathlineto{\pgfqpoint{1.281993in}{1.142978in}}%
\pgfpathlineto{\pgfqpoint{1.282190in}{1.166859in}}%
\pgfpathlineto{\pgfqpoint{1.265840in}{1.174459in}}%
\pgfpathlineto{\pgfqpoint{1.265664in}{1.150467in}}%
\pgfpathclose%
\pgfusepath{fill}%
\end{pgfscope}%
\begin{pgfscope}%
\pgfpathrectangle{\pgfqpoint{0.407592in}{0.222904in}}{\pgfqpoint{1.411163in}{1.411163in}}%
\pgfusepath{clip}%
\pgfsetbuttcap%
\pgfsetroundjoin%
\definecolor{currentfill}{rgb}{0.168627,0.611765,0.411765}%
\pgfsetfillcolor{currentfill}%
\pgfsetfillopacity{0.180000}%
\pgfsetlinewidth{0.000000pt}%
\definecolor{currentstroke}{rgb}{0.000000,0.000000,0.000000}%
\pgfsetstrokecolor{currentstroke}%
\pgfsetdash{}{0pt}%
\pgfpathmoveto{\pgfqpoint{1.248738in}{1.086849in}}%
\pgfpathlineto{\pgfqpoint{1.265138in}{1.078869in}}%
\pgfpathlineto{\pgfqpoint{1.265313in}{1.102672in}}%
\pgfpathlineto{\pgfqpoint{1.248892in}{1.110765in}}%
\pgfpathlineto{\pgfqpoint{1.248738in}{1.086849in}}%
\pgfpathclose%
\pgfusepath{fill}%
\end{pgfscope}%
\begin{pgfscope}%
\pgfpathrectangle{\pgfqpoint{0.407592in}{0.222904in}}{\pgfqpoint{1.411163in}{1.411163in}}%
\pgfusepath{clip}%
\pgfsetbuttcap%
\pgfsetroundjoin%
\definecolor{currentfill}{rgb}{0.168627,0.611765,0.411765}%
\pgfsetfillcolor{currentfill}%
\pgfsetfillopacity{0.180000}%
\pgfsetlinewidth{0.000000pt}%
\definecolor{currentstroke}{rgb}{0.000000,0.000000,0.000000}%
\pgfsetstrokecolor{currentstroke}%
\pgfsetdash{}{0pt}%
\pgfpathmoveto{\pgfqpoint{1.265840in}{1.174459in}}%
\pgfpathlineto{\pgfqpoint{1.282190in}{1.166859in}}%
\pgfpathlineto{\pgfqpoint{1.282387in}{1.190802in}}%
\pgfpathlineto{\pgfqpoint{1.266017in}{1.198514in}}%
\pgfpathlineto{\pgfqpoint{1.265840in}{1.174459in}}%
\pgfpathclose%
\pgfusepath{fill}%
\end{pgfscope}%
\begin{pgfscope}%
\pgfpathrectangle{\pgfqpoint{0.407592in}{0.222904in}}{\pgfqpoint{1.411163in}{1.411163in}}%
\pgfusepath{clip}%
\pgfsetbuttcap%
\pgfsetroundjoin%
\definecolor{currentfill}{rgb}{0.168627,0.611765,0.411765}%
\pgfsetfillcolor{currentfill}%
\pgfsetfillopacity{0.180000}%
\pgfsetlinewidth{0.000000pt}%
\definecolor{currentstroke}{rgb}{0.000000,0.000000,0.000000}%
\pgfsetstrokecolor{currentstroke}%
\pgfsetdash{}{0pt}%
\pgfpathmoveto{\pgfqpoint{1.248892in}{1.110765in}}%
\pgfpathlineto{\pgfqpoint{1.265313in}{1.102672in}}%
\pgfpathlineto{\pgfqpoint{1.265488in}{1.126538in}}%
\pgfpathlineto{\pgfqpoint{1.249047in}{1.134743in}}%
\pgfpathlineto{\pgfqpoint{1.248892in}{1.110765in}}%
\pgfpathclose%
\pgfusepath{fill}%
\end{pgfscope}%
\begin{pgfscope}%
\pgfpathrectangle{\pgfqpoint{0.407592in}{0.222904in}}{\pgfqpoint{1.411163in}{1.411163in}}%
\pgfusepath{clip}%
\pgfsetbuttcap%
\pgfsetroundjoin%
\definecolor{currentfill}{rgb}{0.168627,0.611765,0.411765}%
\pgfsetfillcolor{currentfill}%
\pgfsetfillopacity{0.180000}%
\pgfsetlinewidth{0.000000pt}%
\definecolor{currentstroke}{rgb}{0.000000,0.000000,0.000000}%
\pgfsetstrokecolor{currentstroke}%
\pgfsetdash{}{0pt}%
\pgfpathmoveto{\pgfqpoint{1.249047in}{1.134743in}}%
\pgfpathlineto{\pgfqpoint{1.265488in}{1.126538in}}%
\pgfpathlineto{\pgfqpoint{1.265664in}{1.150467in}}%
\pgfpathlineto{\pgfqpoint{1.249201in}{1.158786in}}%
\pgfpathlineto{\pgfqpoint{1.249047in}{1.134743in}}%
\pgfpathclose%
\pgfusepath{fill}%
\end{pgfscope}%
\begin{pgfscope}%
\pgfpathrectangle{\pgfqpoint{0.407592in}{0.222904in}}{\pgfqpoint{1.411163in}{1.411163in}}%
\pgfusepath{clip}%
\pgfsetbuttcap%
\pgfsetroundjoin%
\definecolor{currentfill}{rgb}{0.168627,0.611765,0.411765}%
\pgfsetfillcolor{currentfill}%
\pgfsetfillopacity{0.180000}%
\pgfsetlinewidth{0.000000pt}%
\definecolor{currentstroke}{rgb}{0.000000,0.000000,0.000000}%
\pgfsetstrokecolor{currentstroke}%
\pgfsetdash{}{0pt}%
\pgfpathmoveto{\pgfqpoint{1.249201in}{1.158786in}}%
\pgfpathlineto{\pgfqpoint{1.265664in}{1.150467in}}%
\pgfpathlineto{\pgfqpoint{1.265840in}{1.174459in}}%
\pgfpathlineto{\pgfqpoint{1.249356in}{1.182892in}}%
\pgfpathlineto{\pgfqpoint{1.249201in}{1.158786in}}%
\pgfpathclose%
\pgfusepath{fill}%
\end{pgfscope}%
\begin{pgfscope}%
\pgfpathrectangle{\pgfqpoint{0.407592in}{0.222904in}}{\pgfqpoint{1.411163in}{1.411163in}}%
\pgfusepath{clip}%
\pgfsetbuttcap%
\pgfsetroundjoin%
\definecolor{currentfill}{rgb}{0.168627,0.611765,0.411765}%
\pgfsetfillcolor{currentfill}%
\pgfsetfillopacity{0.180000}%
\pgfsetlinewidth{0.000000pt}%
\definecolor{currentstroke}{rgb}{0.000000,0.000000,0.000000}%
\pgfsetstrokecolor{currentstroke}%
\pgfsetdash{}{0pt}%
\pgfpathmoveto{\pgfqpoint{1.249356in}{1.182892in}}%
\pgfpathlineto{\pgfqpoint{1.265840in}{1.174459in}}%
\pgfpathlineto{\pgfqpoint{1.266017in}{1.198514in}}%
\pgfpathlineto{\pgfqpoint{1.249512in}{1.207062in}}%
\pgfpathlineto{\pgfqpoint{1.249356in}{1.182892in}}%
\pgfpathclose%
\pgfusepath{fill}%
\end{pgfscope}%
\begin{pgfscope}%
\pgfpathrectangle{\pgfqpoint{0.407592in}{0.222904in}}{\pgfqpoint{1.411163in}{1.411163in}}%
\pgfusepath{clip}%
\pgfsetbuttcap%
\pgfsetroundjoin%
\definecolor{currentfill}{rgb}{0.556863,0.654902,0.847059}%
\pgfsetfillcolor{currentfill}%
\pgfsetfillopacity{0.130000}%
\pgfsetlinewidth{0.000000pt}%
\definecolor{currentstroke}{rgb}{0.000000,0.000000,0.000000}%
\pgfsetstrokecolor{currentstroke}%
\pgfsetdash{}{0pt}%
\pgfpathmoveto{\pgfqpoint{0.871478in}{1.394437in}}%
\pgfpathlineto{\pgfqpoint{0.886848in}{1.378173in}}%
\pgfpathlineto{\pgfqpoint{0.902230in}{1.362091in}}%
\pgfpathlineto{\pgfqpoint{0.910091in}{1.373743in}}%
\pgfpathlineto{\pgfqpoint{0.894724in}{1.389630in}}%
\pgfpathlineto{\pgfqpoint{0.879369in}{1.405704in}}%
\pgfusepath{fill}%
\end{pgfscope}%
\begin{pgfscope}%
\pgfpathrectangle{\pgfqpoint{0.407592in}{0.222904in}}{\pgfqpoint{1.411163in}{1.411163in}}%
\pgfusepath{clip}%
\pgfsetbuttcap%
\pgfsetroundjoin%
\definecolor{currentfill}{rgb}{0.556863,0.654902,0.847059}%
\pgfsetfillcolor{currentfill}%
\pgfsetfillopacity{0.130000}%
\pgfsetlinewidth{0.000000pt}%
\definecolor{currentstroke}{rgb}{0.000000,0.000000,0.000000}%
\pgfsetstrokecolor{currentstroke}%
\pgfsetdash{}{0pt}%
\pgfpathmoveto{\pgfqpoint{0.902230in}{1.362091in}}%
\pgfpathlineto{\pgfqpoint{0.917624in}{1.346205in}}%
\pgfpathlineto{\pgfqpoint{0.933031in}{1.330531in}}%
\pgfpathlineto{\pgfqpoint{0.940863in}{1.342588in}}%
\pgfpathlineto{\pgfqpoint{0.925470in}{1.358057in}}%
\pgfpathlineto{\pgfqpoint{0.910091in}{1.373743in}}%
\pgfusepath{fill}%
\end{pgfscope}%
\begin{pgfscope}%
\pgfpathrectangle{\pgfqpoint{0.407592in}{0.222904in}}{\pgfqpoint{1.411163in}{1.411163in}}%
\pgfusepath{clip}%
\pgfsetbuttcap%
\pgfsetroundjoin%
\definecolor{currentfill}{rgb}{0.556863,0.654902,0.847059}%
\pgfsetfillcolor{currentfill}%
\pgfsetfillopacity{0.130000}%
\pgfsetlinewidth{0.000000pt}%
\definecolor{currentstroke}{rgb}{0.000000,0.000000,0.000000}%
\pgfsetstrokecolor{currentstroke}%
\pgfsetdash{}{0pt}%
\pgfpathmoveto{\pgfqpoint{0.855637in}{1.372404in}}%
\pgfpathlineto{\pgfqpoint{0.871038in}{1.355759in}}%
\pgfpathlineto{\pgfqpoint{0.886451in}{1.339285in}}%
\pgfpathlineto{\pgfqpoint{0.894350in}{1.350602in}}%
\pgfpathlineto{\pgfqpoint{0.902230in}{1.362091in}}%
\pgfpathlineto{\pgfqpoint{0.886848in}{1.378173in}}%
\pgfpathlineto{\pgfqpoint{0.871478in}{1.394437in}}%
\pgfpathlineto{\pgfqpoint{0.863568in}{1.383335in}}%
\pgfpathlineto{\pgfqpoint{0.855637in}{1.372404in}}%
\pgfpathclose%
\pgfusepath{fill}%
\end{pgfscope}%
\begin{pgfscope}%
\pgfpathrectangle{\pgfqpoint{0.407592in}{0.222904in}}{\pgfqpoint{1.411163in}{1.411163in}}%
\pgfusepath{clip}%
\pgfsetbuttcap%
\pgfsetroundjoin%
\definecolor{currentfill}{rgb}{0.556863,0.654902,0.847059}%
\pgfsetfillcolor{currentfill}%
\pgfsetfillopacity{0.130000}%
\pgfsetlinewidth{0.000000pt}%
\definecolor{currentstroke}{rgb}{0.000000,0.000000,0.000000}%
\pgfsetstrokecolor{currentstroke}%
\pgfsetdash{}{0pt}%
\pgfpathmoveto{\pgfqpoint{0.933031in}{1.330531in}}%
\pgfpathlineto{\pgfqpoint{0.948450in}{1.315085in}}%
\pgfpathlineto{\pgfqpoint{0.963882in}{1.299885in}}%
\pgfpathlineto{\pgfqpoint{0.971687in}{1.312368in}}%
\pgfpathlineto{\pgfqpoint{0.956268in}{1.327353in}}%
\pgfpathlineto{\pgfqpoint{0.940863in}{1.342588in}}%
\pgfusepath{fill}%
\end{pgfscope}%
\begin{pgfscope}%
\pgfpathrectangle{\pgfqpoint{0.407592in}{0.222904in}}{\pgfqpoint{1.411163in}{1.411163in}}%
\pgfusepath{clip}%
\pgfsetbuttcap%
\pgfsetroundjoin%
\definecolor{currentfill}{rgb}{0.556863,0.654902,0.847059}%
\pgfsetfillcolor{currentfill}%
\pgfsetfillopacity{0.130000}%
\pgfsetlinewidth{0.000000pt}%
\definecolor{currentstroke}{rgb}{0.000000,0.000000,0.000000}%
\pgfsetstrokecolor{currentstroke}%
\pgfsetdash{}{0pt}%
\pgfpathmoveto{\pgfqpoint{0.886451in}{1.339285in}}%
\pgfpathlineto{\pgfqpoint{0.901875in}{1.322996in}}%
\pgfpathlineto{\pgfqpoint{0.917311in}{1.306908in}}%
\pgfpathlineto{\pgfqpoint{0.925180in}{1.318634in}}%
\pgfpathlineto{\pgfqpoint{0.933031in}{1.330531in}}%
\pgfpathlineto{\pgfqpoint{0.917624in}{1.346205in}}%
\pgfpathlineto{\pgfqpoint{0.902230in}{1.362091in}}%
\pgfpathlineto{\pgfqpoint{0.894350in}{1.350602in}}%
\pgfpathlineto{\pgfqpoint{0.886451in}{1.339285in}}%
\pgfpathclose%
\pgfusepath{fill}%
\end{pgfscope}%
\begin{pgfscope}%
\pgfpathrectangle{\pgfqpoint{0.407592in}{0.222904in}}{\pgfqpoint{1.411163in}{1.411163in}}%
\pgfusepath{clip}%
\pgfsetbuttcap%
\pgfsetroundjoin%
\definecolor{currentfill}{rgb}{0.556863,0.654902,0.847059}%
\pgfsetfillcolor{currentfill}%
\pgfsetfillopacity{0.130000}%
\pgfsetlinewidth{0.000000pt}%
\definecolor{currentstroke}{rgb}{0.000000,0.000000,0.000000}%
\pgfsetstrokecolor{currentstroke}%
\pgfsetdash{}{0pt}%
\pgfpathmoveto{\pgfqpoint{0.963882in}{1.299885in}}%
\pgfpathlineto{\pgfqpoint{0.979329in}{1.284951in}}%
\pgfpathlineto{\pgfqpoint{0.994790in}{1.270304in}}%
\pgfpathlineto{\pgfqpoint{1.002568in}{1.283228in}}%
\pgfpathlineto{\pgfqpoint{0.987120in}{1.297653in}}%
\pgfpathlineto{\pgfqpoint{0.971687in}{1.312368in}}%
\pgfusepath{fill}%
\end{pgfscope}%
\begin{pgfscope}%
\pgfpathrectangle{\pgfqpoint{0.407592in}{0.222904in}}{\pgfqpoint{1.411163in}{1.411163in}}%
\pgfusepath{clip}%
\pgfsetbuttcap%
\pgfsetroundjoin%
\definecolor{currentfill}{rgb}{0.556863,0.654902,0.847059}%
\pgfsetfillcolor{currentfill}%
\pgfsetfillopacity{0.130000}%
\pgfsetlinewidth{0.000000pt}%
\definecolor{currentstroke}{rgb}{0.000000,0.000000,0.000000}%
\pgfsetstrokecolor{currentstroke}%
\pgfsetdash{}{0pt}%
\pgfpathmoveto{\pgfqpoint{0.917311in}{1.306908in}}%
\pgfpathlineto{\pgfqpoint{0.932759in}{1.291037in}}%
\pgfpathlineto{\pgfqpoint{0.948220in}{1.275402in}}%
\pgfpathlineto{\pgfqpoint{0.956060in}{1.287560in}}%
\pgfpathlineto{\pgfqpoint{0.963882in}{1.299885in}}%
\pgfpathlineto{\pgfqpoint{0.948450in}{1.315085in}}%
\pgfpathlineto{\pgfqpoint{0.933031in}{1.330531in}}%
\pgfpathlineto{\pgfqpoint{0.925180in}{1.318634in}}%
\pgfpathlineto{\pgfqpoint{0.917311in}{1.306908in}}%
\pgfpathclose%
\pgfusepath{fill}%
\end{pgfscope}%
\begin{pgfscope}%
\pgfpathrectangle{\pgfqpoint{0.407592in}{0.222904in}}{\pgfqpoint{1.411163in}{1.411163in}}%
\pgfusepath{clip}%
\pgfsetbuttcap%
\pgfsetroundjoin%
\definecolor{currentfill}{rgb}{0.556863,0.654902,0.847059}%
\pgfsetfillcolor{currentfill}%
\pgfsetfillopacity{0.130000}%
\pgfsetlinewidth{0.000000pt}%
\definecolor{currentstroke}{rgb}{0.000000,0.000000,0.000000}%
\pgfsetstrokecolor{currentstroke}%
\pgfsetdash{}{0pt}%
\pgfpathmoveto{\pgfqpoint{0.839713in}{1.351088in}}%
\pgfpathlineto{\pgfqpoint{0.855147in}{1.334064in}}%
\pgfpathlineto{\pgfqpoint{0.870591in}{1.317199in}}%
\pgfpathlineto{\pgfqpoint{0.878531in}{1.328148in}}%
\pgfpathlineto{\pgfqpoint{0.886451in}{1.339285in}}%
\pgfpathlineto{\pgfqpoint{0.871038in}{1.355759in}}%
\pgfpathlineto{\pgfqpoint{0.855637in}{1.372404in}}%
\pgfpathlineto{\pgfqpoint{0.847685in}{1.361653in}}%
\pgfpathlineto{\pgfqpoint{0.839713in}{1.351088in}}%
\pgfpathclose%
\pgfusepath{fill}%
\end{pgfscope}%
\begin{pgfscope}%
\pgfpathrectangle{\pgfqpoint{0.407592in}{0.222904in}}{\pgfqpoint{1.411163in}{1.411163in}}%
\pgfusepath{clip}%
\pgfsetbuttcap%
\pgfsetroundjoin%
\definecolor{currentfill}{rgb}{0.556863,0.654902,0.847059}%
\pgfsetfillcolor{currentfill}%
\pgfsetfillopacity{0.130000}%
\pgfsetlinewidth{0.000000pt}%
\definecolor{currentstroke}{rgb}{0.000000,0.000000,0.000000}%
\pgfsetstrokecolor{currentstroke}%
\pgfsetdash{}{0pt}%
\pgfpathmoveto{\pgfqpoint{0.994790in}{1.270304in}}%
\pgfpathlineto{\pgfqpoint{1.010266in}{1.255965in}}%
\pgfpathlineto{\pgfqpoint{1.025758in}{1.241959in}}%
\pgfpathlineto{\pgfqpoint{1.033512in}{1.255338in}}%
\pgfpathlineto{\pgfqpoint{1.018032in}{1.269115in}}%
\pgfpathlineto{\pgfqpoint{1.002568in}{1.283228in}}%
\pgfusepath{fill}%
\end{pgfscope}%
\begin{pgfscope}%
\pgfpathrectangle{\pgfqpoint{0.407592in}{0.222904in}}{\pgfqpoint{1.411163in}{1.411163in}}%
\pgfusepath{clip}%
\pgfsetbuttcap%
\pgfsetroundjoin%
\definecolor{currentfill}{rgb}{0.556863,0.654902,0.847059}%
\pgfsetfillcolor{currentfill}%
\pgfsetfillopacity{0.130000}%
\pgfsetlinewidth{0.000000pt}%
\definecolor{currentstroke}{rgb}{0.000000,0.000000,0.000000}%
\pgfsetstrokecolor{currentstroke}%
\pgfsetdash{}{0pt}%
\pgfpathmoveto{\pgfqpoint{0.948220in}{1.275402in}}%
\pgfpathlineto{\pgfqpoint{0.963694in}{1.260023in}}%
\pgfpathlineto{\pgfqpoint{0.979181in}{1.244921in}}%
\pgfpathlineto{\pgfqpoint{0.986994in}{1.257531in}}%
\pgfpathlineto{\pgfqpoint{0.994790in}{1.270304in}}%
\pgfpathlineto{\pgfqpoint{0.979329in}{1.284951in}}%
\pgfpathlineto{\pgfqpoint{0.963882in}{1.299885in}}%
\pgfpathlineto{\pgfqpoint{0.956060in}{1.287560in}}%
\pgfpathlineto{\pgfqpoint{0.948220in}{1.275402in}}%
\pgfpathclose%
\pgfusepath{fill}%
\end{pgfscope}%
\begin{pgfscope}%
\pgfpathrectangle{\pgfqpoint{0.407592in}{0.222904in}}{\pgfqpoint{1.411163in}{1.411163in}}%
\pgfusepath{clip}%
\pgfsetbuttcap%
\pgfsetroundjoin%
\definecolor{currentfill}{rgb}{0.556863,0.654902,0.847059}%
\pgfsetfillcolor{currentfill}%
\pgfsetfillopacity{0.130000}%
\pgfsetlinewidth{0.000000pt}%
\definecolor{currentstroke}{rgb}{0.000000,0.000000,0.000000}%
\pgfsetstrokecolor{currentstroke}%
\pgfsetdash{}{0pt}%
\pgfpathmoveto{\pgfqpoint{0.870591in}{1.317199in}}%
\pgfpathlineto{\pgfqpoint{0.886047in}{1.300507in}}%
\pgfpathlineto{\pgfqpoint{0.901514in}{1.284003in}}%
\pgfpathlineto{\pgfqpoint{0.909422in}{1.295361in}}%
\pgfpathlineto{\pgfqpoint{0.917311in}{1.306908in}}%
\pgfpathlineto{\pgfqpoint{0.901875in}{1.322996in}}%
\pgfpathlineto{\pgfqpoint{0.886451in}{1.339285in}}%
\pgfpathlineto{\pgfqpoint{0.878531in}{1.328148in}}%
\pgfpathlineto{\pgfqpoint{0.870591in}{1.317199in}}%
\pgfpathclose%
\pgfusepath{fill}%
\end{pgfscope}%
\begin{pgfscope}%
\pgfpathrectangle{\pgfqpoint{0.407592in}{0.222904in}}{\pgfqpoint{1.411163in}{1.411163in}}%
\pgfusepath{clip}%
\pgfsetbuttcap%
\pgfsetroundjoin%
\definecolor{currentfill}{rgb}{0.556863,0.654902,0.847059}%
\pgfsetfillcolor{currentfill}%
\pgfsetfillopacity{0.130000}%
\pgfsetlinewidth{0.000000pt}%
\definecolor{currentstroke}{rgb}{0.000000,0.000000,0.000000}%
\pgfsetstrokecolor{currentstroke}%
\pgfsetdash{}{0pt}%
\pgfpathmoveto{\pgfqpoint{1.025758in}{1.241959in}}%
\pgfpathlineto{\pgfqpoint{1.041266in}{1.228311in}}%
\pgfpathlineto{\pgfqpoint{1.056792in}{1.215048in}}%
\pgfpathlineto{\pgfqpoint{1.064524in}{1.228888in}}%
\pgfpathlineto{\pgfqpoint{1.049009in}{1.241920in}}%
\pgfpathlineto{\pgfqpoint{1.033512in}{1.255338in}}%
\pgfusepath{fill}%
\end{pgfscope}%
\begin{pgfscope}%
\pgfpathrectangle{\pgfqpoint{0.407592in}{0.222904in}}{\pgfqpoint{1.411163in}{1.411163in}}%
\pgfusepath{clip}%
\pgfsetbuttcap%
\pgfsetroundjoin%
\definecolor{currentfill}{rgb}{0.556863,0.654902,0.847059}%
\pgfsetfillcolor{currentfill}%
\pgfsetfillopacity{0.130000}%
\pgfsetlinewidth{0.000000pt}%
\definecolor{currentstroke}{rgb}{0.000000,0.000000,0.000000}%
\pgfsetstrokecolor{currentstroke}%
\pgfsetdash{}{0pt}%
\pgfpathmoveto{\pgfqpoint{0.901514in}{1.284003in}}%
\pgfpathlineto{\pgfqpoint{0.916992in}{1.267704in}}%
\pgfpathlineto{\pgfqpoint{0.932483in}{1.251628in}}%
\pgfpathlineto{\pgfqpoint{0.940361in}{1.263421in}}%
\pgfpathlineto{\pgfqpoint{0.948220in}{1.275402in}}%
\pgfpathlineto{\pgfqpoint{0.932759in}{1.291037in}}%
\pgfpathlineto{\pgfqpoint{0.917311in}{1.306908in}}%
\pgfpathlineto{\pgfqpoint{0.909422in}{1.295361in}}%
\pgfpathlineto{\pgfqpoint{0.901514in}{1.284003in}}%
\pgfpathclose%
\pgfusepath{fill}%
\end{pgfscope}%
\begin{pgfscope}%
\pgfpathrectangle{\pgfqpoint{0.407592in}{0.222904in}}{\pgfqpoint{1.411163in}{1.411163in}}%
\pgfusepath{clip}%
\pgfsetbuttcap%
\pgfsetroundjoin%
\definecolor{currentfill}{rgb}{0.556863,0.654902,0.847059}%
\pgfsetfillcolor{currentfill}%
\pgfsetfillopacity{0.130000}%
\pgfsetlinewidth{0.000000pt}%
\definecolor{currentstroke}{rgb}{0.000000,0.000000,0.000000}%
\pgfsetstrokecolor{currentstroke}%
\pgfsetdash{}{0pt}%
\pgfpathmoveto{\pgfqpoint{0.979181in}{1.244921in}}%
\pgfpathlineto{\pgfqpoint{0.994683in}{1.230119in}}%
\pgfpathlineto{\pgfqpoint{1.010200in}{1.215642in}}%
\pgfpathlineto{\pgfqpoint{1.017987in}{1.228723in}}%
\pgfpathlineto{\pgfqpoint{1.025758in}{1.241959in}}%
\pgfpathlineto{\pgfqpoint{1.010266in}{1.255965in}}%
\pgfpathlineto{\pgfqpoint{0.994790in}{1.270304in}}%
\pgfpathlineto{\pgfqpoint{0.986994in}{1.257531in}}%
\pgfpathlineto{\pgfqpoint{0.979181in}{1.244921in}}%
\pgfpathclose%
\pgfusepath{fill}%
\end{pgfscope}%
\begin{pgfscope}%
\pgfpathrectangle{\pgfqpoint{0.407592in}{0.222904in}}{\pgfqpoint{1.411163in}{1.411163in}}%
\pgfusepath{clip}%
\pgfsetbuttcap%
\pgfsetroundjoin%
\definecolor{currentfill}{rgb}{0.556863,0.654902,0.847059}%
\pgfsetfillcolor{currentfill}%
\pgfsetfillopacity{0.130000}%
\pgfsetlinewidth{0.000000pt}%
\definecolor{currentstroke}{rgb}{0.000000,0.000000,0.000000}%
\pgfsetstrokecolor{currentstroke}%
\pgfsetdash{}{0pt}%
\pgfpathmoveto{\pgfqpoint{0.823702in}{1.330551in}}%
\pgfpathlineto{\pgfqpoint{0.839170in}{1.313154in}}%
\pgfpathlineto{\pgfqpoint{0.854648in}{1.295902in}}%
\pgfpathlineto{\pgfqpoint{0.862630in}{1.306447in}}%
\pgfpathlineto{\pgfqpoint{0.870591in}{1.317199in}}%
\pgfpathlineto{\pgfqpoint{0.855147in}{1.334064in}}%
\pgfpathlineto{\pgfqpoint{0.839713in}{1.351088in}}%
\pgfpathlineto{\pgfqpoint{0.831718in}{1.340718in}}%
\pgfpathlineto{\pgfqpoint{0.823702in}{1.330551in}}%
\pgfpathclose%
\pgfusepath{fill}%
\end{pgfscope}%
\begin{pgfscope}%
\pgfpathrectangle{\pgfqpoint{0.407592in}{0.222904in}}{\pgfqpoint{1.411163in}{1.411163in}}%
\pgfusepath{clip}%
\pgfsetbuttcap%
\pgfsetroundjoin%
\definecolor{currentfill}{rgb}{0.556863,0.654902,0.847059}%
\pgfsetfillcolor{currentfill}%
\pgfsetfillopacity{0.130000}%
\pgfsetlinewidth{0.000000pt}%
\definecolor{currentstroke}{rgb}{0.000000,0.000000,0.000000}%
\pgfsetstrokecolor{currentstroke}%
\pgfsetdash{}{0pt}%
\pgfpathmoveto{\pgfqpoint{0.932483in}{1.251628in}}%
\pgfpathlineto{\pgfqpoint{0.947986in}{1.235795in}}%
\pgfpathlineto{\pgfqpoint{0.963502in}{1.220228in}}%
\pgfpathlineto{\pgfqpoint{0.971351in}{1.232482in}}%
\pgfpathlineto{\pgfqpoint{0.979181in}{1.244921in}}%
\pgfpathlineto{\pgfqpoint{0.963694in}{1.260023in}}%
\pgfpathlineto{\pgfqpoint{0.948220in}{1.275402in}}%
\pgfpathlineto{\pgfqpoint{0.940361in}{1.263421in}}%
\pgfpathlineto{\pgfqpoint{0.932483in}{1.251628in}}%
\pgfpathclose%
\pgfusepath{fill}%
\end{pgfscope}%
\begin{pgfscope}%
\pgfpathrectangle{\pgfqpoint{0.407592in}{0.222904in}}{\pgfqpoint{1.411163in}{1.411163in}}%
\pgfusepath{clip}%
\pgfsetbuttcap%
\pgfsetroundjoin%
\definecolor{currentfill}{rgb}{0.556863,0.654902,0.847059}%
\pgfsetfillcolor{currentfill}%
\pgfsetfillopacity{0.130000}%
\pgfsetlinewidth{0.000000pt}%
\definecolor{currentstroke}{rgb}{0.000000,0.000000,0.000000}%
\pgfsetstrokecolor{currentstroke}%
\pgfsetdash{}{0pt}%
\pgfpathmoveto{\pgfqpoint{1.056792in}{1.215048in}}%
\pgfpathlineto{\pgfqpoint{1.072337in}{1.202198in}}%
\pgfpathlineto{\pgfqpoint{1.087900in}{1.189791in}}%
\pgfpathlineto{\pgfqpoint{1.095612in}{1.204093in}}%
\pgfpathlineto{\pgfqpoint{1.080058in}{1.216270in}}%
\pgfpathlineto{\pgfqpoint{1.064524in}{1.228888in}}%
\pgfusepath{fill}%
\end{pgfscope}%
\begin{pgfscope}%
\pgfpathrectangle{\pgfqpoint{0.407592in}{0.222904in}}{\pgfqpoint{1.411163in}{1.411163in}}%
\pgfusepath{clip}%
\pgfsetbuttcap%
\pgfsetroundjoin%
\definecolor{currentfill}{rgb}{0.556863,0.654902,0.847059}%
\pgfsetfillcolor{currentfill}%
\pgfsetfillopacity{0.130000}%
\pgfsetlinewidth{0.000000pt}%
\definecolor{currentstroke}{rgb}{0.000000,0.000000,0.000000}%
\pgfsetstrokecolor{currentstroke}%
\pgfsetdash{}{0pt}%
\pgfpathmoveto{\pgfqpoint{1.010200in}{1.215642in}}%
\pgfpathlineto{\pgfqpoint{1.025732in}{1.201517in}}%
\pgfpathlineto{\pgfqpoint{1.041281in}{1.187774in}}%
\pgfpathlineto{\pgfqpoint{1.049045in}{1.201339in}}%
\pgfpathlineto{\pgfqpoint{1.056792in}{1.215048in}}%
\pgfpathlineto{\pgfqpoint{1.041266in}{1.228311in}}%
\pgfpathlineto{\pgfqpoint{1.025758in}{1.241959in}}%
\pgfpathlineto{\pgfqpoint{1.017987in}{1.228723in}}%
\pgfpathlineto{\pgfqpoint{1.010200in}{1.215642in}}%
\pgfpathclose%
\pgfusepath{fill}%
\end{pgfscope}%
\begin{pgfscope}%
\pgfpathrectangle{\pgfqpoint{0.407592in}{0.222904in}}{\pgfqpoint{1.411163in}{1.411163in}}%
\pgfusepath{clip}%
\pgfsetbuttcap%
\pgfsetroundjoin%
\definecolor{currentfill}{rgb}{0.556863,0.654902,0.847059}%
\pgfsetfillcolor{currentfill}%
\pgfsetfillopacity{0.130000}%
\pgfsetlinewidth{0.000000pt}%
\definecolor{currentstroke}{rgb}{0.000000,0.000000,0.000000}%
\pgfsetstrokecolor{currentstroke}%
\pgfsetdash{}{0pt}%
\pgfpathmoveto{\pgfqpoint{0.854648in}{1.295902in}}%
\pgfpathlineto{\pgfqpoint{0.870136in}{1.278810in}}%
\pgfpathlineto{\pgfqpoint{0.885636in}{1.261892in}}%
\pgfpathlineto{\pgfqpoint{0.893585in}{1.272843in}}%
\pgfpathlineto{\pgfqpoint{0.901514in}{1.284003in}}%
\pgfpathlineto{\pgfqpoint{0.886047in}{1.300507in}}%
\pgfpathlineto{\pgfqpoint{0.870591in}{1.317199in}}%
\pgfpathlineto{\pgfqpoint{0.862630in}{1.306447in}}%
\pgfpathlineto{\pgfqpoint{0.854648in}{1.295902in}}%
\pgfpathclose%
\pgfusepath{fill}%
\end{pgfscope}%
\begin{pgfscope}%
\pgfpathrectangle{\pgfqpoint{0.407592in}{0.222904in}}{\pgfqpoint{1.411163in}{1.411163in}}%
\pgfusepath{clip}%
\pgfsetbuttcap%
\pgfsetroundjoin%
\definecolor{currentfill}{rgb}{0.556863,0.654902,0.847059}%
\pgfsetfillcolor{currentfill}%
\pgfsetfillopacity{0.130000}%
\pgfsetlinewidth{0.000000pt}%
\definecolor{currentstroke}{rgb}{0.000000,0.000000,0.000000}%
\pgfsetstrokecolor{currentstroke}%
\pgfsetdash{}{0pt}%
\pgfpathmoveto{\pgfqpoint{0.963502in}{1.220228in}}%
\pgfpathlineto{\pgfqpoint{0.979031in}{1.204950in}}%
\pgfpathlineto{\pgfqpoint{0.994574in}{1.189987in}}%
\pgfpathlineto{\pgfqpoint{1.002396in}{1.202726in}}%
\pgfpathlineto{\pgfqpoint{1.010200in}{1.215642in}}%
\pgfpathlineto{\pgfqpoint{0.994683in}{1.230119in}}%
\pgfpathlineto{\pgfqpoint{0.979181in}{1.244921in}}%
\pgfpathlineto{\pgfqpoint{0.971351in}{1.232482in}}%
\pgfpathlineto{\pgfqpoint{0.963502in}{1.220228in}}%
\pgfpathclose%
\pgfusepath{fill}%
\end{pgfscope}%
\begin{pgfscope}%
\pgfpathrectangle{\pgfqpoint{0.407592in}{0.222904in}}{\pgfqpoint{1.411163in}{1.411163in}}%
\pgfusepath{clip}%
\pgfsetbuttcap%
\pgfsetroundjoin%
\definecolor{currentfill}{rgb}{0.556863,0.654902,0.847059}%
\pgfsetfillcolor{currentfill}%
\pgfsetfillopacity{0.130000}%
\pgfsetlinewidth{0.000000pt}%
\definecolor{currentstroke}{rgb}{0.000000,0.000000,0.000000}%
\pgfsetstrokecolor{currentstroke}%
\pgfsetdash{}{0pt}%
\pgfpathmoveto{\pgfqpoint{0.885636in}{1.261892in}}%
\pgfpathlineto{\pgfqpoint{0.901146in}{1.245164in}}%
\pgfpathlineto{\pgfqpoint{0.916668in}{1.228646in}}%
\pgfpathlineto{\pgfqpoint{0.924585in}{1.240032in}}%
\pgfpathlineto{\pgfqpoint{0.932483in}{1.251628in}}%
\pgfpathlineto{\pgfqpoint{0.916992in}{1.267704in}}%
\pgfpathlineto{\pgfqpoint{0.901514in}{1.284003in}}%
\pgfpathlineto{\pgfqpoint{0.893585in}{1.272843in}}%
\pgfpathlineto{\pgfqpoint{0.885636in}{1.261892in}}%
\pgfpathclose%
\pgfusepath{fill}%
\end{pgfscope}%
\begin{pgfscope}%
\pgfpathrectangle{\pgfqpoint{0.407592in}{0.222904in}}{\pgfqpoint{1.411163in}{1.411163in}}%
\pgfusepath{clip}%
\pgfsetbuttcap%
\pgfsetroundjoin%
\definecolor{currentfill}{rgb}{0.556863,0.654902,0.847059}%
\pgfsetfillcolor{currentfill}%
\pgfsetfillopacity{0.130000}%
\pgfsetlinewidth{0.000000pt}%
\definecolor{currentstroke}{rgb}{0.000000,0.000000,0.000000}%
\pgfsetstrokecolor{currentstroke}%
\pgfsetdash{}{0pt}%
\pgfpathmoveto{\pgfqpoint{1.087900in}{1.189791in}}%
\pgfpathlineto{\pgfqpoint{1.103485in}{1.177859in}}%
\pgfpathlineto{\pgfqpoint{1.119091in}{1.166434in}}%
\pgfpathlineto{\pgfqpoint{1.126785in}{1.181183in}}%
\pgfpathlineto{\pgfqpoint{1.111187in}{1.192387in}}%
\pgfpathlineto{\pgfqpoint{1.095612in}{1.204093in}}%
\pgfusepath{fill}%
\end{pgfscope}%
\begin{pgfscope}%
\pgfpathrectangle{\pgfqpoint{0.407592in}{0.222904in}}{\pgfqpoint{1.411163in}{1.411163in}}%
\pgfusepath{clip}%
\pgfsetbuttcap%
\pgfsetroundjoin%
\definecolor{currentfill}{rgb}{0.556863,0.654902,0.847059}%
\pgfsetfillcolor{currentfill}%
\pgfsetfillopacity{0.130000}%
\pgfsetlinewidth{0.000000pt}%
\definecolor{currentstroke}{rgb}{0.000000,0.000000,0.000000}%
\pgfsetstrokecolor{currentstroke}%
\pgfsetdash{}{0pt}%
\pgfpathmoveto{\pgfqpoint{1.041281in}{1.187774in}}%
\pgfpathlineto{\pgfqpoint{1.056848in}{1.174443in}}%
\pgfpathlineto{\pgfqpoint{1.072432in}{1.161557in}}%
\pgfpathlineto{\pgfqpoint{1.080174in}{1.175609in}}%
\pgfpathlineto{\pgfqpoint{1.087900in}{1.189791in}}%
\pgfpathlineto{\pgfqpoint{1.072337in}{1.202198in}}%
\pgfpathlineto{\pgfqpoint{1.056792in}{1.215048in}}%
\pgfpathlineto{\pgfqpoint{1.049045in}{1.201339in}}%
\pgfpathlineto{\pgfqpoint{1.041281in}{1.187774in}}%
\pgfpathclose%
\pgfusepath{fill}%
\end{pgfscope}%
\begin{pgfscope}%
\pgfpathrectangle{\pgfqpoint{0.407592in}{0.222904in}}{\pgfqpoint{1.411163in}{1.411163in}}%
\pgfusepath{clip}%
\pgfsetbuttcap%
\pgfsetroundjoin%
\definecolor{currentfill}{rgb}{0.556863,0.654902,0.847059}%
\pgfsetfillcolor{currentfill}%
\pgfsetfillopacity{0.130000}%
\pgfsetlinewidth{0.000000pt}%
\definecolor{currentstroke}{rgb}{0.000000,0.000000,0.000000}%
\pgfsetstrokecolor{currentstroke}%
\pgfsetdash{}{0pt}%
\pgfpathmoveto{\pgfqpoint{0.807601in}{1.310858in}}%
\pgfpathlineto{\pgfqpoint{0.823104in}{1.293097in}}%
\pgfpathlineto{\pgfqpoint{0.838617in}{1.275467in}}%
\pgfpathlineto{\pgfqpoint{0.846643in}{1.285572in}}%
\pgfpathlineto{\pgfqpoint{0.854648in}{1.295902in}}%
\pgfpathlineto{\pgfqpoint{0.839170in}{1.313154in}}%
\pgfpathlineto{\pgfqpoint{0.823702in}{1.330551in}}%
\pgfpathlineto{\pgfqpoint{0.815663in}{1.320595in}}%
\pgfpathlineto{\pgfqpoint{0.807601in}{1.310858in}}%
\pgfpathclose%
\pgfusepath{fill}%
\end{pgfscope}%
\begin{pgfscope}%
\pgfpathrectangle{\pgfqpoint{0.407592in}{0.222904in}}{\pgfqpoint{1.411163in}{1.411163in}}%
\pgfusepath{clip}%
\pgfsetbuttcap%
\pgfsetroundjoin%
\definecolor{currentfill}{rgb}{0.556863,0.654902,0.847059}%
\pgfsetfillcolor{currentfill}%
\pgfsetfillopacity{0.130000}%
\pgfsetlinewidth{0.000000pt}%
\definecolor{currentstroke}{rgb}{0.000000,0.000000,0.000000}%
\pgfsetstrokecolor{currentstroke}%
\pgfsetdash{}{0pt}%
\pgfpathmoveto{\pgfqpoint{0.916668in}{1.228646in}}%
\pgfpathlineto{\pgfqpoint{0.932201in}{1.212356in}}%
\pgfpathlineto{\pgfqpoint{0.947747in}{1.196318in}}%
\pgfpathlineto{\pgfqpoint{0.955634in}{1.208169in}}%
\pgfpathlineto{\pgfqpoint{0.963502in}{1.220228in}}%
\pgfpathlineto{\pgfqpoint{0.947986in}{1.235795in}}%
\pgfpathlineto{\pgfqpoint{0.932483in}{1.251628in}}%
\pgfpathlineto{\pgfqpoint{0.924585in}{1.240032in}}%
\pgfpathlineto{\pgfqpoint{0.916668in}{1.228646in}}%
\pgfpathclose%
\pgfusepath{fill}%
\end{pgfscope}%
\begin{pgfscope}%
\pgfpathrectangle{\pgfqpoint{0.407592in}{0.222904in}}{\pgfqpoint{1.411163in}{1.411163in}}%
\pgfusepath{clip}%
\pgfsetbuttcap%
\pgfsetroundjoin%
\definecolor{currentfill}{rgb}{0.556863,0.654902,0.847059}%
\pgfsetfillcolor{currentfill}%
\pgfsetfillopacity{0.130000}%
\pgfsetlinewidth{0.000000pt}%
\definecolor{currentstroke}{rgb}{0.000000,0.000000,0.000000}%
\pgfsetstrokecolor{currentstroke}%
\pgfsetdash{}{0pt}%
\pgfpathmoveto{\pgfqpoint{0.994574in}{1.189987in}}%
\pgfpathlineto{\pgfqpoint{1.010132in}{1.175368in}}%
\pgfpathlineto{\pgfqpoint{1.025705in}{1.161124in}}%
\pgfpathlineto{\pgfqpoint{1.033502in}{1.174365in}}%
\pgfpathlineto{\pgfqpoint{1.041281in}{1.187774in}}%
\pgfpathlineto{\pgfqpoint{1.025732in}{1.201517in}}%
\pgfpathlineto{\pgfqpoint{1.010200in}{1.215642in}}%
\pgfpathlineto{\pgfqpoint{1.002396in}{1.202726in}}%
\pgfpathlineto{\pgfqpoint{0.994574in}{1.189987in}}%
\pgfpathclose%
\pgfusepath{fill}%
\end{pgfscope}%
\begin{pgfscope}%
\pgfpathrectangle{\pgfqpoint{0.407592in}{0.222904in}}{\pgfqpoint{1.411163in}{1.411163in}}%
\pgfusepath{clip}%
\pgfsetbuttcap%
\pgfsetroundjoin%
\definecolor{currentfill}{rgb}{0.556863,0.654902,0.847059}%
\pgfsetfillcolor{currentfill}%
\pgfsetfillopacity{0.130000}%
\pgfsetlinewidth{0.000000pt}%
\definecolor{currentstroke}{rgb}{0.000000,0.000000,0.000000}%
\pgfsetstrokecolor{currentstroke}%
\pgfsetdash{}{0pt}%
\pgfpathmoveto{\pgfqpoint{0.838617in}{1.275467in}}%
\pgfpathlineto{\pgfqpoint{0.854140in}{1.257983in}}%
\pgfpathlineto{\pgfqpoint{0.869673in}{1.240657in}}%
\pgfpathlineto{\pgfqpoint{0.877665in}{1.251160in}}%
\pgfpathlineto{\pgfqpoint{0.885636in}{1.261892in}}%
\pgfpathlineto{\pgfqpoint{0.870136in}{1.278810in}}%
\pgfpathlineto{\pgfqpoint{0.854648in}{1.295902in}}%
\pgfpathlineto{\pgfqpoint{0.846643in}{1.285572in}}%
\pgfpathlineto{\pgfqpoint{0.838617in}{1.275467in}}%
\pgfpathclose%
\pgfusepath{fill}%
\end{pgfscope}%
\begin{pgfscope}%
\pgfpathrectangle{\pgfqpoint{0.407592in}{0.222904in}}{\pgfqpoint{1.411163in}{1.411163in}}%
\pgfusepath{clip}%
\pgfsetbuttcap%
\pgfsetroundjoin%
\definecolor{currentfill}{rgb}{0.556863,0.654902,0.847059}%
\pgfsetfillcolor{currentfill}%
\pgfsetfillopacity{0.130000}%
\pgfsetlinewidth{0.000000pt}%
\definecolor{currentstroke}{rgb}{0.000000,0.000000,0.000000}%
\pgfsetstrokecolor{currentstroke}%
\pgfsetdash{}{0pt}%
\pgfpathmoveto{\pgfqpoint{1.072432in}{1.161557in}}%
\pgfpathlineto{\pgfqpoint{1.088036in}{1.149150in}}%
\pgfpathlineto{\pgfqpoint{1.103662in}{1.137258in}}%
\pgfpathlineto{\pgfqpoint{1.111384in}{1.151788in}}%
\pgfpathlineto{\pgfqpoint{1.119091in}{1.166434in}}%
\pgfpathlineto{\pgfqpoint{1.103485in}{1.177859in}}%
\pgfpathlineto{\pgfqpoint{1.087900in}{1.189791in}}%
\pgfpathlineto{\pgfqpoint{1.080174in}{1.175609in}}%
\pgfpathlineto{\pgfqpoint{1.072432in}{1.161557in}}%
\pgfpathclose%
\pgfusepath{fill}%
\end{pgfscope}%
\begin{pgfscope}%
\pgfpathrectangle{\pgfqpoint{0.407592in}{0.222904in}}{\pgfqpoint{1.411163in}{1.411163in}}%
\pgfusepath{clip}%
\pgfsetbuttcap%
\pgfsetroundjoin%
\definecolor{currentfill}{rgb}{0.556863,0.654902,0.847059}%
\pgfsetfillcolor{currentfill}%
\pgfsetfillopacity{0.130000}%
\pgfsetlinewidth{0.000000pt}%
\definecolor{currentstroke}{rgb}{0.000000,0.000000,0.000000}%
\pgfsetstrokecolor{currentstroke}%
\pgfsetdash{}{0pt}%
\pgfpathmoveto{\pgfqpoint{1.119091in}{1.166434in}}%
\pgfpathlineto{\pgfqpoint{1.134721in}{1.155547in}}%
\pgfpathlineto{\pgfqpoint{1.150376in}{1.145233in}}%
\pgfpathlineto{\pgfqpoint{1.158053in}{1.160402in}}%
\pgfpathlineto{\pgfqpoint{1.142406in}{1.170511in}}%
\pgfpathlineto{\pgfqpoint{1.126785in}{1.181183in}}%
\pgfusepath{fill}%
\end{pgfscope}%
\begin{pgfscope}%
\pgfpathrectangle{\pgfqpoint{0.407592in}{0.222904in}}{\pgfqpoint{1.411163in}{1.411163in}}%
\pgfusepath{clip}%
\pgfsetbuttcap%
\pgfsetroundjoin%
\definecolor{currentfill}{rgb}{0.556863,0.654902,0.847059}%
\pgfsetfillcolor{currentfill}%
\pgfsetfillopacity{0.130000}%
\pgfsetlinewidth{0.000000pt}%
\definecolor{currentstroke}{rgb}{0.000000,0.000000,0.000000}%
\pgfsetstrokecolor{currentstroke}%
\pgfsetdash{}{0pt}%
\pgfpathmoveto{\pgfqpoint{0.947747in}{1.196318in}}%
\pgfpathlineto{\pgfqpoint{0.963305in}{1.180555in}}%
\pgfpathlineto{\pgfqpoint{0.978877in}{1.165094in}}%
\pgfpathlineto{\pgfqpoint{0.986734in}{1.177439in}}%
\pgfpathlineto{\pgfqpoint{0.994574in}{1.189987in}}%
\pgfpathlineto{\pgfqpoint{0.979031in}{1.204950in}}%
\pgfpathlineto{\pgfqpoint{0.963502in}{1.220228in}}%
\pgfpathlineto{\pgfqpoint{0.955634in}{1.208169in}}%
\pgfpathlineto{\pgfqpoint{0.947747in}{1.196318in}}%
\pgfpathclose%
\pgfusepath{fill}%
\end{pgfscope}%
\begin{pgfscope}%
\pgfpathrectangle{\pgfqpoint{0.407592in}{0.222904in}}{\pgfqpoint{1.411163in}{1.411163in}}%
\pgfusepath{clip}%
\pgfsetbuttcap%
\pgfsetroundjoin%
\definecolor{currentfill}{rgb}{0.556863,0.654902,0.847059}%
\pgfsetfillcolor{currentfill}%
\pgfsetfillopacity{0.130000}%
\pgfsetlinewidth{0.000000pt}%
\definecolor{currentstroke}{rgb}{0.000000,0.000000,0.000000}%
\pgfsetstrokecolor{currentstroke}%
\pgfsetdash{}{0pt}%
\pgfpathmoveto{\pgfqpoint{0.869673in}{1.240657in}}%
\pgfpathlineto{\pgfqpoint{0.885216in}{1.223506in}}%
\pgfpathlineto{\pgfqpoint{0.900771in}{1.206548in}}%
\pgfpathlineto{\pgfqpoint{0.908730in}{1.217481in}}%
\pgfpathlineto{\pgfqpoint{0.916668in}{1.228646in}}%
\pgfpathlineto{\pgfqpoint{0.901146in}{1.245164in}}%
\pgfpathlineto{\pgfqpoint{0.885636in}{1.261892in}}%
\pgfpathlineto{\pgfqpoint{0.877665in}{1.251160in}}%
\pgfpathlineto{\pgfqpoint{0.869673in}{1.240657in}}%
\pgfpathclose%
\pgfusepath{fill}%
\end{pgfscope}%
\begin{pgfscope}%
\pgfpathrectangle{\pgfqpoint{0.407592in}{0.222904in}}{\pgfqpoint{1.411163in}{1.411163in}}%
\pgfusepath{clip}%
\pgfsetbuttcap%
\pgfsetroundjoin%
\definecolor{currentfill}{rgb}{0.556863,0.654902,0.847059}%
\pgfsetfillcolor{currentfill}%
\pgfsetfillopacity{0.130000}%
\pgfsetlinewidth{0.000000pt}%
\definecolor{currentstroke}{rgb}{0.000000,0.000000,0.000000}%
\pgfsetstrokecolor{currentstroke}%
\pgfsetdash{}{0pt}%
\pgfpathmoveto{\pgfqpoint{1.025705in}{1.161124in}}%
\pgfpathlineto{\pgfqpoint{1.041295in}{1.147287in}}%
\pgfpathlineto{\pgfqpoint{1.056902in}{1.133893in}}%
\pgfpathlineto{\pgfqpoint{1.064675in}{1.147647in}}%
\pgfpathlineto{\pgfqpoint{1.072432in}{1.161557in}}%
\pgfpathlineto{\pgfqpoint{1.056848in}{1.174443in}}%
\pgfpathlineto{\pgfqpoint{1.041281in}{1.187774in}}%
\pgfpathlineto{\pgfqpoint{1.033502in}{1.174365in}}%
\pgfpathlineto{\pgfqpoint{1.025705in}{1.161124in}}%
\pgfpathclose%
\pgfusepath{fill}%
\end{pgfscope}%
\begin{pgfscope}%
\pgfpathrectangle{\pgfqpoint{0.407592in}{0.222904in}}{\pgfqpoint{1.411163in}{1.411163in}}%
\pgfusepath{clip}%
\pgfsetbuttcap%
\pgfsetroundjoin%
\definecolor{currentfill}{rgb}{0.556863,0.654902,0.847059}%
\pgfsetfillcolor{currentfill}%
\pgfsetfillopacity{0.130000}%
\pgfsetlinewidth{0.000000pt}%
\definecolor{currentstroke}{rgb}{0.000000,0.000000,0.000000}%
\pgfsetstrokecolor{currentstroke}%
\pgfsetdash{}{0pt}%
\pgfpathmoveto{\pgfqpoint{0.791405in}{1.292074in}}%
\pgfpathlineto{\pgfqpoint{0.806944in}{1.273963in}}%
\pgfpathlineto{\pgfqpoint{0.822493in}{1.255969in}}%
\pgfpathlineto{\pgfqpoint{0.830567in}{1.265596in}}%
\pgfpathlineto{\pgfqpoint{0.838617in}{1.275467in}}%
\pgfpathlineto{\pgfqpoint{0.823104in}{1.293097in}}%
\pgfpathlineto{\pgfqpoint{0.807601in}{1.310858in}}%
\pgfpathlineto{\pgfqpoint{0.799515in}{1.301348in}}%
\pgfpathlineto{\pgfqpoint{0.791405in}{1.292074in}}%
\pgfpathclose%
\pgfusepath{fill}%
\end{pgfscope}%
\begin{pgfscope}%
\pgfpathrectangle{\pgfqpoint{0.407592in}{0.222904in}}{\pgfqpoint{1.411163in}{1.411163in}}%
\pgfusepath{clip}%
\pgfsetbuttcap%
\pgfsetroundjoin%
\definecolor{currentfill}{rgb}{0.556863,0.654902,0.847059}%
\pgfsetfillcolor{currentfill}%
\pgfsetfillopacity{0.130000}%
\pgfsetlinewidth{0.000000pt}%
\definecolor{currentstroke}{rgb}{0.000000,0.000000,0.000000}%
\pgfsetstrokecolor{currentstroke}%
\pgfsetdash{}{0pt}%
\pgfpathmoveto{\pgfqpoint{0.978877in}{1.165094in}}%
\pgfpathlineto{\pgfqpoint{0.994462in}{1.149965in}}%
\pgfpathlineto{\pgfqpoint{1.010061in}{1.135201in}}%
\pgfpathlineto{\pgfqpoint{1.017892in}{1.148064in}}%
\pgfpathlineto{\pgfqpoint{1.025705in}{1.161124in}}%
\pgfpathlineto{\pgfqpoint{1.010132in}{1.175368in}}%
\pgfpathlineto{\pgfqpoint{0.994574in}{1.189987in}}%
\pgfpathlineto{\pgfqpoint{0.986734in}{1.177439in}}%
\pgfpathlineto{\pgfqpoint{0.978877in}{1.165094in}}%
\pgfpathclose%
\pgfusepath{fill}%
\end{pgfscope}%
\begin{pgfscope}%
\pgfpathrectangle{\pgfqpoint{0.407592in}{0.222904in}}{\pgfqpoint{1.411163in}{1.411163in}}%
\pgfusepath{clip}%
\pgfsetbuttcap%
\pgfsetroundjoin%
\definecolor{currentfill}{rgb}{0.556863,0.654902,0.847059}%
\pgfsetfillcolor{currentfill}%
\pgfsetfillopacity{0.130000}%
\pgfsetlinewidth{0.000000pt}%
\definecolor{currentstroke}{rgb}{0.000000,0.000000,0.000000}%
\pgfsetstrokecolor{currentstroke}%
\pgfsetdash{}{0pt}%
\pgfpathmoveto{\pgfqpoint{0.900771in}{1.206548in}}%
\pgfpathlineto{\pgfqpoint{0.916337in}{1.189803in}}%
\pgfpathlineto{\pgfqpoint{0.931914in}{1.173292in}}%
\pgfpathlineto{\pgfqpoint{0.939840in}{1.184688in}}%
\pgfpathlineto{\pgfqpoint{0.947747in}{1.196318in}}%
\pgfpathlineto{\pgfqpoint{0.932201in}{1.212356in}}%
\pgfpathlineto{\pgfqpoint{0.916668in}{1.228646in}}%
\pgfpathlineto{\pgfqpoint{0.908730in}{1.217481in}}%
\pgfpathlineto{\pgfqpoint{0.900771in}{1.206548in}}%
\pgfpathclose%
\pgfusepath{fill}%
\end{pgfscope}%
\begin{pgfscope}%
\pgfpathrectangle{\pgfqpoint{0.407592in}{0.222904in}}{\pgfqpoint{1.411163in}{1.411163in}}%
\pgfusepath{clip}%
\pgfsetbuttcap%
\pgfsetroundjoin%
\definecolor{currentfill}{rgb}{0.556863,0.654902,0.847059}%
\pgfsetfillcolor{currentfill}%
\pgfsetfillopacity{0.130000}%
\pgfsetlinewidth{0.000000pt}%
\definecolor{currentstroke}{rgb}{0.000000,0.000000,0.000000}%
\pgfsetstrokecolor{currentstroke}%
\pgfsetdash{}{0pt}%
\pgfpathmoveto{\pgfqpoint{1.103662in}{1.137258in}}%
\pgfpathlineto{\pgfqpoint{1.119309in}{1.125917in}}%
\pgfpathlineto{\pgfqpoint{1.134980in}{1.115167in}}%
\pgfpathlineto{\pgfqpoint{1.142685in}{1.130151in}}%
\pgfpathlineto{\pgfqpoint{1.150376in}{1.145233in}}%
\pgfpathlineto{\pgfqpoint{1.134721in}{1.155547in}}%
\pgfpathlineto{\pgfqpoint{1.119091in}{1.166434in}}%
\pgfpathlineto{\pgfqpoint{1.111384in}{1.151788in}}%
\pgfpathlineto{\pgfqpoint{1.103662in}{1.137258in}}%
\pgfpathclose%
\pgfusepath{fill}%
\end{pgfscope}%
\begin{pgfscope}%
\pgfpathrectangle{\pgfqpoint{0.407592in}{0.222904in}}{\pgfqpoint{1.411163in}{1.411163in}}%
\pgfusepath{clip}%
\pgfsetbuttcap%
\pgfsetroundjoin%
\definecolor{currentfill}{rgb}{0.556863,0.654902,0.847059}%
\pgfsetfillcolor{currentfill}%
\pgfsetfillopacity{0.130000}%
\pgfsetlinewidth{0.000000pt}%
\definecolor{currentstroke}{rgb}{0.000000,0.000000,0.000000}%
\pgfsetstrokecolor{currentstroke}%
\pgfsetdash{}{0pt}%
\pgfpathmoveto{\pgfqpoint{1.150376in}{1.145233in}}%
\pgfpathlineto{\pgfqpoint{1.166057in}{1.135524in}}%
\pgfpathlineto{\pgfqpoint{1.181767in}{1.126452in}}%
\pgfpathlineto{\pgfqpoint{1.189430in}{1.141996in}}%
\pgfpathlineto{\pgfqpoint{1.173727in}{1.150887in}}%
\pgfpathlineto{\pgfqpoint{1.158053in}{1.160402in}}%
\pgfusepath{fill}%
\end{pgfscope}%
\begin{pgfscope}%
\pgfpathrectangle{\pgfqpoint{0.407592in}{0.222904in}}{\pgfqpoint{1.411163in}{1.411163in}}%
\pgfusepath{clip}%
\pgfsetbuttcap%
\pgfsetroundjoin%
\definecolor{currentfill}{rgb}{0.556863,0.654902,0.847059}%
\pgfsetfillcolor{currentfill}%
\pgfsetfillopacity{0.130000}%
\pgfsetlinewidth{0.000000pt}%
\definecolor{currentstroke}{rgb}{0.000000,0.000000,0.000000}%
\pgfsetstrokecolor{currentstroke}%
\pgfsetdash{}{0pt}%
\pgfpathmoveto{\pgfqpoint{1.056902in}{1.133893in}}%
\pgfpathlineto{\pgfqpoint{1.072527in}{1.120981in}}%
\pgfpathlineto{\pgfqpoint{1.088173in}{1.108589in}}%
\pgfpathlineto{\pgfqpoint{1.095925in}{1.122853in}}%
\pgfpathlineto{\pgfqpoint{1.103662in}{1.137258in}}%
\pgfpathlineto{\pgfqpoint{1.088036in}{1.149150in}}%
\pgfpathlineto{\pgfqpoint{1.072432in}{1.161557in}}%
\pgfpathlineto{\pgfqpoint{1.064675in}{1.147647in}}%
\pgfpathlineto{\pgfqpoint{1.056902in}{1.133893in}}%
\pgfpathclose%
\pgfusepath{fill}%
\end{pgfscope}%
\begin{pgfscope}%
\pgfpathrectangle{\pgfqpoint{0.407592in}{0.222904in}}{\pgfqpoint{1.411163in}{1.411163in}}%
\pgfusepath{clip}%
\pgfsetbuttcap%
\pgfsetroundjoin%
\definecolor{currentfill}{rgb}{0.556863,0.654902,0.847059}%
\pgfsetfillcolor{currentfill}%
\pgfsetfillopacity{0.130000}%
\pgfsetlinewidth{0.000000pt}%
\definecolor{currentstroke}{rgb}{0.000000,0.000000,0.000000}%
\pgfsetstrokecolor{currentstroke}%
\pgfsetdash{}{0pt}%
\pgfpathmoveto{\pgfqpoint{0.822493in}{1.255969in}}%
\pgfpathlineto{\pgfqpoint{0.838052in}{1.238104in}}%
\pgfpathlineto{\pgfqpoint{0.853620in}{1.220383in}}%
\pgfpathlineto{\pgfqpoint{0.861658in}{1.230394in}}%
\pgfpathlineto{\pgfqpoint{0.869673in}{1.240657in}}%
\pgfpathlineto{\pgfqpoint{0.854140in}{1.257983in}}%
\pgfpathlineto{\pgfqpoint{0.838617in}{1.275467in}}%
\pgfpathlineto{\pgfqpoint{0.830567in}{1.265596in}}%
\pgfpathlineto{\pgfqpoint{0.822493in}{1.255969in}}%
\pgfpathclose%
\pgfusepath{fill}%
\end{pgfscope}%
\begin{pgfscope}%
\pgfpathrectangle{\pgfqpoint{0.407592in}{0.222904in}}{\pgfqpoint{1.411163in}{1.411163in}}%
\pgfusepath{clip}%
\pgfsetbuttcap%
\pgfsetroundjoin%
\definecolor{currentfill}{rgb}{0.556863,0.654902,0.847059}%
\pgfsetfillcolor{currentfill}%
\pgfsetfillopacity{0.130000}%
\pgfsetlinewidth{0.000000pt}%
\definecolor{currentstroke}{rgb}{0.000000,0.000000,0.000000}%
\pgfsetstrokecolor{currentstroke}%
\pgfsetdash{}{0pt}%
\pgfpathmoveto{\pgfqpoint{0.931914in}{1.173292in}}%
\pgfpathlineto{\pgfqpoint{0.947503in}{1.157040in}}%
\pgfpathlineto{\pgfqpoint{0.963104in}{1.141075in}}%
\pgfpathlineto{\pgfqpoint{0.971000in}{1.152968in}}%
\pgfpathlineto{\pgfqpoint{0.978877in}{1.165094in}}%
\pgfpathlineto{\pgfqpoint{0.963305in}{1.180555in}}%
\pgfpathlineto{\pgfqpoint{0.947747in}{1.196318in}}%
\pgfpathlineto{\pgfqpoint{0.939840in}{1.184688in}}%
\pgfpathlineto{\pgfqpoint{0.931914in}{1.173292in}}%
\pgfpathclose%
\pgfusepath{fill}%
\end{pgfscope}%
\begin{pgfscope}%
\pgfpathrectangle{\pgfqpoint{0.407592in}{0.222904in}}{\pgfqpoint{1.411163in}{1.411163in}}%
\pgfusepath{clip}%
\pgfsetbuttcap%
\pgfsetroundjoin%
\definecolor{currentfill}{rgb}{0.556863,0.654902,0.847059}%
\pgfsetfillcolor{currentfill}%
\pgfsetfillopacity{0.130000}%
\pgfsetlinewidth{0.000000pt}%
\definecolor{currentstroke}{rgb}{0.000000,0.000000,0.000000}%
\pgfsetstrokecolor{currentstroke}%
\pgfsetdash{}{0pt}%
\pgfpathmoveto{\pgfqpoint{1.010061in}{1.135201in}}%
\pgfpathlineto{\pgfqpoint{1.025676in}{1.120837in}}%
\pgfpathlineto{\pgfqpoint{1.041307in}{1.106910in}}%
\pgfpathlineto{\pgfqpoint{1.049113in}{1.120309in}}%
\pgfpathlineto{\pgfqpoint{1.056902in}{1.133893in}}%
\pgfpathlineto{\pgfqpoint{1.041295in}{1.147287in}}%
\pgfpathlineto{\pgfqpoint{1.025705in}{1.161124in}}%
\pgfpathlineto{\pgfqpoint{1.017892in}{1.148064in}}%
\pgfpathlineto{\pgfqpoint{1.010061in}{1.135201in}}%
\pgfpathclose%
\pgfusepath{fill}%
\end{pgfscope}%
\begin{pgfscope}%
\pgfpathrectangle{\pgfqpoint{0.407592in}{0.222904in}}{\pgfqpoint{1.411163in}{1.411163in}}%
\pgfusepath{clip}%
\pgfsetbuttcap%
\pgfsetroundjoin%
\definecolor{currentfill}{rgb}{0.556863,0.654902,0.847059}%
\pgfsetfillcolor{currentfill}%
\pgfsetfillopacity{0.130000}%
\pgfsetlinewidth{0.000000pt}%
\definecolor{currentstroke}{rgb}{0.000000,0.000000,0.000000}%
\pgfsetstrokecolor{currentstroke}%
\pgfsetdash{}{0pt}%
\pgfpathmoveto{\pgfqpoint{0.853620in}{1.220383in}}%
\pgfpathlineto{\pgfqpoint{0.869199in}{1.202819in}}%
\pgfpathlineto{\pgfqpoint{0.884788in}{1.185431in}}%
\pgfpathlineto{\pgfqpoint{0.892791in}{1.195861in}}%
\pgfpathlineto{\pgfqpoint{0.900771in}{1.206548in}}%
\pgfpathlineto{\pgfqpoint{0.885216in}{1.223506in}}%
\pgfpathlineto{\pgfqpoint{0.869673in}{1.240657in}}%
\pgfpathlineto{\pgfqpoint{0.861658in}{1.230394in}}%
\pgfpathlineto{\pgfqpoint{0.853620in}{1.220383in}}%
\pgfpathclose%
\pgfusepath{fill}%
\end{pgfscope}%
\begin{pgfscope}%
\pgfpathrectangle{\pgfqpoint{0.407592in}{0.222904in}}{\pgfqpoint{1.411163in}{1.411163in}}%
\pgfusepath{clip}%
\pgfsetbuttcap%
\pgfsetroundjoin%
\definecolor{currentfill}{rgb}{0.556863,0.654902,0.847059}%
\pgfsetfillcolor{currentfill}%
\pgfsetfillopacity{0.130000}%
\pgfsetlinewidth{0.000000pt}%
\definecolor{currentstroke}{rgb}{0.000000,0.000000,0.000000}%
\pgfsetstrokecolor{currentstroke}%
\pgfsetdash{}{0pt}%
\pgfpathmoveto{\pgfqpoint{1.134980in}{1.115167in}}%
\pgfpathlineto{\pgfqpoint{1.150676in}{1.105044in}}%
\pgfpathlineto{\pgfqpoint{1.166400in}{1.095585in}}%
\pgfpathlineto{\pgfqpoint{1.174090in}{1.110979in}}%
\pgfpathlineto{\pgfqpoint{1.181767in}{1.126452in}}%
\pgfpathlineto{\pgfqpoint{1.166057in}{1.135524in}}%
\pgfpathlineto{\pgfqpoint{1.150376in}{1.145233in}}%
\pgfpathlineto{\pgfqpoint{1.142685in}{1.130151in}}%
\pgfpathlineto{\pgfqpoint{1.134980in}{1.115167in}}%
\pgfpathclose%
\pgfusepath{fill}%
\end{pgfscope}%
\begin{pgfscope}%
\pgfpathrectangle{\pgfqpoint{0.407592in}{0.222904in}}{\pgfqpoint{1.411163in}{1.411163in}}%
\pgfusepath{clip}%
\pgfsetbuttcap%
\pgfsetroundjoin%
\definecolor{currentfill}{rgb}{0.556863,0.654902,0.847059}%
\pgfsetfillcolor{currentfill}%
\pgfsetfillopacity{0.130000}%
\pgfsetlinewidth{0.000000pt}%
\definecolor{currentstroke}{rgb}{0.000000,0.000000,0.000000}%
\pgfsetstrokecolor{currentstroke}%
\pgfsetdash{}{0pt}%
\pgfpathmoveto{\pgfqpoint{1.088173in}{1.108589in}}%
\pgfpathlineto{\pgfqpoint{1.103839in}{1.096760in}}%
\pgfpathlineto{\pgfqpoint{1.119527in}{1.085537in}}%
\pgfpathlineto{\pgfqpoint{1.127261in}{1.100291in}}%
\pgfpathlineto{\pgfqpoint{1.134980in}{1.115167in}}%
\pgfpathlineto{\pgfqpoint{1.119309in}{1.125917in}}%
\pgfpathlineto{\pgfqpoint{1.103662in}{1.137258in}}%
\pgfpathlineto{\pgfqpoint{1.095925in}{1.122853in}}%
\pgfpathlineto{\pgfqpoint{1.088173in}{1.108589in}}%
\pgfpathclose%
\pgfusepath{fill}%
\end{pgfscope}%
\begin{pgfscope}%
\pgfpathrectangle{\pgfqpoint{0.407592in}{0.222904in}}{\pgfqpoint{1.411163in}{1.411163in}}%
\pgfusepath{clip}%
\pgfsetbuttcap%
\pgfsetroundjoin%
\definecolor{currentfill}{rgb}{0.556863,0.654902,0.847059}%
\pgfsetfillcolor{currentfill}%
\pgfsetfillopacity{0.130000}%
\pgfsetlinewidth{0.000000pt}%
\definecolor{currentstroke}{rgb}{0.000000,0.000000,0.000000}%
\pgfsetstrokecolor{currentstroke}%
\pgfsetdash{}{0pt}%
\pgfpathmoveto{\pgfqpoint{0.963104in}{1.141075in}}%
\pgfpathlineto{\pgfqpoint{0.978718in}{1.125426in}}%
\pgfpathlineto{\pgfqpoint{0.994346in}{1.110128in}}%
\pgfpathlineto{\pgfqpoint{1.002213in}{1.122550in}}%
\pgfpathlineto{\pgfqpoint{1.010061in}{1.135201in}}%
\pgfpathlineto{\pgfqpoint{0.994462in}{1.149965in}}%
\pgfpathlineto{\pgfqpoint{0.978877in}{1.165094in}}%
\pgfpathlineto{\pgfqpoint{0.971000in}{1.152968in}}%
\pgfpathlineto{\pgfqpoint{0.963104in}{1.141075in}}%
\pgfpathclose%
\pgfusepath{fill}%
\end{pgfscope}%
\begin{pgfscope}%
\pgfpathrectangle{\pgfqpoint{0.407592in}{0.222904in}}{\pgfqpoint{1.411163in}{1.411163in}}%
\pgfusepath{clip}%
\pgfsetbuttcap%
\pgfsetroundjoin%
\definecolor{currentfill}{rgb}{0.556863,0.654902,0.847059}%
\pgfsetfillcolor{currentfill}%
\pgfsetfillopacity{0.130000}%
\pgfsetlinewidth{0.000000pt}%
\definecolor{currentstroke}{rgb}{0.000000,0.000000,0.000000}%
\pgfsetstrokecolor{currentstroke}%
\pgfsetdash{}{0pt}%
\pgfpathmoveto{\pgfqpoint{1.181767in}{1.126452in}}%
\pgfpathlineto{\pgfqpoint{1.197506in}{1.118048in}}%
\pgfpathlineto{\pgfqpoint{1.213278in}{1.110341in}}%
\pgfpathlineto{\pgfqpoint{1.220930in}{1.126195in}}%
\pgfpathlineto{\pgfqpoint{1.205163in}{1.133756in}}%
\pgfpathlineto{\pgfqpoint{1.189430in}{1.141996in}}%
\pgfusepath{fill}%
\end{pgfscope}%
\begin{pgfscope}%
\pgfpathrectangle{\pgfqpoint{0.407592in}{0.222904in}}{\pgfqpoint{1.411163in}{1.411163in}}%
\pgfusepath{clip}%
\pgfsetbuttcap%
\pgfsetroundjoin%
\definecolor{currentfill}{rgb}{0.556863,0.654902,0.847059}%
\pgfsetfillcolor{currentfill}%
\pgfsetfillopacity{0.130000}%
\pgfsetlinewidth{0.000000pt}%
\definecolor{currentstroke}{rgb}{0.000000,0.000000,0.000000}%
\pgfsetstrokecolor{currentstroke}%
\pgfsetdash{}{0pt}%
\pgfpathmoveto{\pgfqpoint{0.775110in}{1.274265in}}%
\pgfpathlineto{\pgfqpoint{0.790687in}{1.255823in}}%
\pgfpathlineto{\pgfqpoint{0.806273in}{1.237483in}}%
\pgfpathlineto{\pgfqpoint{0.814396in}{1.246595in}}%
\pgfpathlineto{\pgfqpoint{0.822493in}{1.255969in}}%
\pgfpathlineto{\pgfqpoint{0.806944in}{1.273963in}}%
\pgfpathlineto{\pgfqpoint{0.791405in}{1.292074in}}%
\pgfpathlineto{\pgfqpoint{0.783270in}{1.283044in}}%
\pgfpathlineto{\pgfqpoint{0.775110in}{1.274265in}}%
\pgfpathclose%
\pgfusepath{fill}%
\end{pgfscope}%
\begin{pgfscope}%
\pgfpathrectangle{\pgfqpoint{0.407592in}{0.222904in}}{\pgfqpoint{1.411163in}{1.411163in}}%
\pgfusepath{clip}%
\pgfsetbuttcap%
\pgfsetroundjoin%
\definecolor{currentfill}{rgb}{0.556863,0.654902,0.847059}%
\pgfsetfillcolor{currentfill}%
\pgfsetfillopacity{0.130000}%
\pgfsetlinewidth{0.000000pt}%
\definecolor{currentstroke}{rgb}{0.000000,0.000000,0.000000}%
\pgfsetstrokecolor{currentstroke}%
\pgfsetdash{}{0pt}%
\pgfpathmoveto{\pgfqpoint{0.884788in}{1.185431in}}%
\pgfpathlineto{\pgfqpoint{0.900387in}{1.168237in}}%
\pgfpathlineto{\pgfqpoint{0.915998in}{1.151260in}}%
\pgfpathlineto{\pgfqpoint{0.923966in}{1.162144in}}%
\pgfpathlineto{\pgfqpoint{0.931914in}{1.173292in}}%
\pgfpathlineto{\pgfqpoint{0.916337in}{1.189803in}}%
\pgfpathlineto{\pgfqpoint{0.900771in}{1.206548in}}%
\pgfpathlineto{\pgfqpoint{0.892791in}{1.195861in}}%
\pgfpathlineto{\pgfqpoint{0.884788in}{1.185431in}}%
\pgfpathclose%
\pgfusepath{fill}%
\end{pgfscope}%
\begin{pgfscope}%
\pgfpathrectangle{\pgfqpoint{0.407592in}{0.222904in}}{\pgfqpoint{1.411163in}{1.411163in}}%
\pgfusepath{clip}%
\pgfsetbuttcap%
\pgfsetroundjoin%
\definecolor{currentfill}{rgb}{0.556863,0.654902,0.847059}%
\pgfsetfillcolor{currentfill}%
\pgfsetfillopacity{0.130000}%
\pgfsetlinewidth{0.000000pt}%
\definecolor{currentstroke}{rgb}{0.000000,0.000000,0.000000}%
\pgfsetstrokecolor{currentstroke}%
\pgfsetdash{}{0pt}%
\pgfpathmoveto{\pgfqpoint{1.041307in}{1.106910in}}%
\pgfpathlineto{\pgfqpoint{1.056955in}{1.093462in}}%
\pgfpathlineto{\pgfqpoint{1.072622in}{1.080538in}}%
\pgfpathlineto{\pgfqpoint{1.080405in}{1.094479in}}%
\pgfpathlineto{\pgfqpoint{1.088173in}{1.108589in}}%
\pgfpathlineto{\pgfqpoint{1.072527in}{1.120981in}}%
\pgfpathlineto{\pgfqpoint{1.056902in}{1.133893in}}%
\pgfpathlineto{\pgfqpoint{1.049113in}{1.120309in}}%
\pgfpathlineto{\pgfqpoint{1.041307in}{1.106910in}}%
\pgfpathclose%
\pgfusepath{fill}%
\end{pgfscope}%
\begin{pgfscope}%
\pgfpathrectangle{\pgfqpoint{0.407592in}{0.222904in}}{\pgfqpoint{1.411163in}{1.411163in}}%
\pgfusepath{clip}%
\pgfsetbuttcap%
\pgfsetroundjoin%
\definecolor{currentfill}{rgb}{0.556863,0.654902,0.847059}%
\pgfsetfillcolor{currentfill}%
\pgfsetfillopacity{0.130000}%
\pgfsetlinewidth{0.000000pt}%
\definecolor{currentstroke}{rgb}{0.000000,0.000000,0.000000}%
\pgfsetstrokecolor{currentstroke}%
\pgfsetdash{}{0pt}%
\pgfpathmoveto{\pgfqpoint{0.806273in}{1.237483in}}%
\pgfpathlineto{\pgfqpoint{0.821869in}{1.219256in}}%
\pgfpathlineto{\pgfqpoint{0.837474in}{1.201157in}}%
\pgfpathlineto{\pgfqpoint{0.845559in}{1.210633in}}%
\pgfpathlineto{\pgfqpoint{0.853620in}{1.220383in}}%
\pgfpathlineto{\pgfqpoint{0.838052in}{1.238104in}}%
\pgfpathlineto{\pgfqpoint{0.822493in}{1.255969in}}%
\pgfpathlineto{\pgfqpoint{0.814396in}{1.246595in}}%
\pgfpathlineto{\pgfqpoint{0.806273in}{1.237483in}}%
\pgfpathclose%
\pgfusepath{fill}%
\end{pgfscope}%
\begin{pgfscope}%
\pgfpathrectangle{\pgfqpoint{0.407592in}{0.222904in}}{\pgfqpoint{1.411163in}{1.411163in}}%
\pgfusepath{clip}%
\pgfsetbuttcap%
\pgfsetroundjoin%
\definecolor{currentfill}{rgb}{0.556863,0.654902,0.847059}%
\pgfsetfillcolor{currentfill}%
\pgfsetfillopacity{0.130000}%
\pgfsetlinewidth{0.000000pt}%
\definecolor{currentstroke}{rgb}{0.000000,0.000000,0.000000}%
\pgfsetstrokecolor{currentstroke}%
\pgfsetdash{}{0pt}%
\pgfpathmoveto{\pgfqpoint{0.994346in}{1.110128in}}%
\pgfpathlineto{\pgfqpoint{1.009988in}{1.095217in}}%
\pgfpathlineto{\pgfqpoint{1.025645in}{1.080735in}}%
\pgfpathlineto{\pgfqpoint{1.033485in}{1.093713in}}%
\pgfpathlineto{\pgfqpoint{1.041307in}{1.106910in}}%
\pgfpathlineto{\pgfqpoint{1.025676in}{1.120837in}}%
\pgfpathlineto{\pgfqpoint{1.010061in}{1.135201in}}%
\pgfpathlineto{\pgfqpoint{1.002213in}{1.122550in}}%
\pgfpathlineto{\pgfqpoint{0.994346in}{1.110128in}}%
\pgfpathclose%
\pgfusepath{fill}%
\end{pgfscope}%
\begin{pgfscope}%
\pgfpathrectangle{\pgfqpoint{0.407592in}{0.222904in}}{\pgfqpoint{1.411163in}{1.411163in}}%
\pgfusepath{clip}%
\pgfsetbuttcap%
\pgfsetroundjoin%
\definecolor{currentfill}{rgb}{0.556863,0.654902,0.847059}%
\pgfsetfillcolor{currentfill}%
\pgfsetfillopacity{0.130000}%
\pgfsetlinewidth{0.000000pt}%
\definecolor{currentstroke}{rgb}{0.000000,0.000000,0.000000}%
\pgfsetstrokecolor{currentstroke}%
\pgfsetdash{}{0pt}%
\pgfpathmoveto{\pgfqpoint{0.915998in}{1.151260in}}%
\pgfpathlineto{\pgfqpoint{0.931619in}{1.134522in}}%
\pgfpathlineto{\pgfqpoint{0.947252in}{1.118052in}}%
\pgfpathlineto{\pgfqpoint{0.955188in}{1.129431in}}%
\pgfpathlineto{\pgfqpoint{0.963104in}{1.141075in}}%
\pgfpathlineto{\pgfqpoint{0.947503in}{1.157040in}}%
\pgfpathlineto{\pgfqpoint{0.931914in}{1.173292in}}%
\pgfpathlineto{\pgfqpoint{0.923966in}{1.162144in}}%
\pgfpathlineto{\pgfqpoint{0.915998in}{1.151260in}}%
\pgfpathclose%
\pgfusepath{fill}%
\end{pgfscope}%
\begin{pgfscope}%
\pgfpathrectangle{\pgfqpoint{0.407592in}{0.222904in}}{\pgfqpoint{1.411163in}{1.411163in}}%
\pgfusepath{clip}%
\pgfsetbuttcap%
\pgfsetroundjoin%
\definecolor{currentfill}{rgb}{0.556863,0.654902,0.847059}%
\pgfsetfillcolor{currentfill}%
\pgfsetfillopacity{0.130000}%
\pgfsetlinewidth{0.000000pt}%
\definecolor{currentstroke}{rgb}{0.000000,0.000000,0.000000}%
\pgfsetstrokecolor{currentstroke}%
\pgfsetdash{}{0pt}%
\pgfpathmoveto{\pgfqpoint{1.166400in}{1.095585in}}%
\pgfpathlineto{\pgfqpoint{1.182153in}{1.086828in}}%
\pgfpathlineto{\pgfqpoint{1.197937in}{1.078806in}}%
\pgfpathlineto{\pgfqpoint{1.205614in}{1.094541in}}%
\pgfpathlineto{\pgfqpoint{1.213278in}{1.110341in}}%
\pgfpathlineto{\pgfqpoint{1.197506in}{1.118048in}}%
\pgfpathlineto{\pgfqpoint{1.181767in}{1.126452in}}%
\pgfpathlineto{\pgfqpoint{1.174090in}{1.110979in}}%
\pgfpathlineto{\pgfqpoint{1.166400in}{1.095585in}}%
\pgfpathclose%
\pgfusepath{fill}%
\end{pgfscope}%
\begin{pgfscope}%
\pgfpathrectangle{\pgfqpoint{0.407592in}{0.222904in}}{\pgfqpoint{1.411163in}{1.411163in}}%
\pgfusepath{clip}%
\pgfsetbuttcap%
\pgfsetroundjoin%
\definecolor{currentfill}{rgb}{0.556863,0.654902,0.847059}%
\pgfsetfillcolor{currentfill}%
\pgfsetfillopacity{0.130000}%
\pgfsetlinewidth{0.000000pt}%
\definecolor{currentstroke}{rgb}{0.000000,0.000000,0.000000}%
\pgfsetstrokecolor{currentstroke}%
\pgfsetdash{}{0pt}%
\pgfpathmoveto{\pgfqpoint{1.119527in}{1.085537in}}%
\pgfpathlineto{\pgfqpoint{1.135240in}{1.074963in}}%
\pgfpathlineto{\pgfqpoint{1.150979in}{1.065082in}}%
\pgfpathlineto{\pgfqpoint{1.158696in}{1.080282in}}%
\pgfpathlineto{\pgfqpoint{1.166400in}{1.095585in}}%
\pgfpathlineto{\pgfqpoint{1.150676in}{1.105044in}}%
\pgfpathlineto{\pgfqpoint{1.134980in}{1.115167in}}%
\pgfpathlineto{\pgfqpoint{1.127261in}{1.100291in}}%
\pgfpathlineto{\pgfqpoint{1.119527in}{1.085537in}}%
\pgfpathclose%
\pgfusepath{fill}%
\end{pgfscope}%
\begin{pgfscope}%
\pgfpathrectangle{\pgfqpoint{0.407592in}{0.222904in}}{\pgfqpoint{1.411163in}{1.411163in}}%
\pgfusepath{clip}%
\pgfsetbuttcap%
\pgfsetroundjoin%
\definecolor{currentfill}{rgb}{0.556863,0.654902,0.847059}%
\pgfsetfillcolor{currentfill}%
\pgfsetfillopacity{0.130000}%
\pgfsetlinewidth{0.000000pt}%
\definecolor{currentstroke}{rgb}{0.000000,0.000000,0.000000}%
\pgfsetstrokecolor{currentstroke}%
\pgfsetdash{}{0pt}%
\pgfpathmoveto{\pgfqpoint{0.837474in}{1.201157in}}%
\pgfpathlineto{\pgfqpoint{0.853089in}{1.183198in}}%
\pgfpathlineto{\pgfqpoint{0.868714in}{1.165395in}}%
\pgfpathlineto{\pgfqpoint{0.876762in}{1.175272in}}%
\pgfpathlineto{\pgfqpoint{0.884788in}{1.185431in}}%
\pgfpathlineto{\pgfqpoint{0.869199in}{1.202819in}}%
\pgfpathlineto{\pgfqpoint{0.853620in}{1.220383in}}%
\pgfpathlineto{\pgfqpoint{0.845559in}{1.210633in}}%
\pgfpathlineto{\pgfqpoint{0.837474in}{1.201157in}}%
\pgfpathclose%
\pgfusepath{fill}%
\end{pgfscope}%
\begin{pgfscope}%
\pgfpathrectangle{\pgfqpoint{0.407592in}{0.222904in}}{\pgfqpoint{1.411163in}{1.411163in}}%
\pgfusepath{clip}%
\pgfsetbuttcap%
\pgfsetroundjoin%
\definecolor{currentfill}{rgb}{0.556863,0.654902,0.847059}%
\pgfsetfillcolor{currentfill}%
\pgfsetfillopacity{0.130000}%
\pgfsetlinewidth{0.000000pt}%
\definecolor{currentstroke}{rgb}{0.000000,0.000000,0.000000}%
\pgfsetstrokecolor{currentstroke}%
\pgfsetdash{}{0pt}%
\pgfpathmoveto{\pgfqpoint{1.072622in}{1.080538in}}%
\pgfpathlineto{\pgfqpoint{1.088308in}{1.068185in}}%
\pgfpathlineto{\pgfqpoint{1.104016in}{1.056450in}}%
\pgfpathlineto{\pgfqpoint{1.111779in}{1.070918in}}%
\pgfpathlineto{\pgfqpoint{1.119527in}{1.085537in}}%
\pgfpathlineto{\pgfqpoint{1.103839in}{1.096760in}}%
\pgfpathlineto{\pgfqpoint{1.088173in}{1.108589in}}%
\pgfpathlineto{\pgfqpoint{1.080405in}{1.094479in}}%
\pgfpathlineto{\pgfqpoint{1.072622in}{1.080538in}}%
\pgfpathclose%
\pgfusepath{fill}%
\end{pgfscope}%
\begin{pgfscope}%
\pgfpathrectangle{\pgfqpoint{0.407592in}{0.222904in}}{\pgfqpoint{1.411163in}{1.411163in}}%
\pgfusepath{clip}%
\pgfsetbuttcap%
\pgfsetroundjoin%
\definecolor{currentfill}{rgb}{0.556863,0.654902,0.847059}%
\pgfsetfillcolor{currentfill}%
\pgfsetfillopacity{0.130000}%
\pgfsetlinewidth{0.000000pt}%
\definecolor{currentstroke}{rgb}{0.000000,0.000000,0.000000}%
\pgfsetstrokecolor{currentstroke}%
\pgfsetdash{}{0pt}%
\pgfpathmoveto{\pgfqpoint{1.213278in}{1.110341in}}%
\pgfpathlineto{\pgfqpoint{1.229085in}{1.103356in}}%
\pgfpathlineto{\pgfqpoint{1.244927in}{1.097116in}}%
\pgfpathlineto{\pgfqpoint{1.252568in}{1.113201in}}%
\pgfpathlineto{\pgfqpoint{1.236730in}{1.119337in}}%
\pgfpathlineto{\pgfqpoint{1.220930in}{1.126195in}}%
\pgfusepath{fill}%
\end{pgfscope}%
\begin{pgfscope}%
\pgfpathrectangle{\pgfqpoint{0.407592in}{0.222904in}}{\pgfqpoint{1.411163in}{1.411163in}}%
\pgfusepath{clip}%
\pgfsetbuttcap%
\pgfsetroundjoin%
\definecolor{currentfill}{rgb}{0.556863,0.654902,0.847059}%
\pgfsetfillcolor{currentfill}%
\pgfsetfillopacity{0.130000}%
\pgfsetlinewidth{0.000000pt}%
\definecolor{currentstroke}{rgb}{0.000000,0.000000,0.000000}%
\pgfsetstrokecolor{currentstroke}%
\pgfsetdash{}{0pt}%
\pgfpathmoveto{\pgfqpoint{0.947252in}{1.118052in}}%
\pgfpathlineto{\pgfqpoint{0.962897in}{1.101880in}}%
\pgfpathlineto{\pgfqpoint{0.978555in}{1.086041in}}%
\pgfpathlineto{\pgfqpoint{0.986460in}{1.097952in}}%
\pgfpathlineto{\pgfqpoint{0.994346in}{1.110128in}}%
\pgfpathlineto{\pgfqpoint{0.978718in}{1.125426in}}%
\pgfpathlineto{\pgfqpoint{0.963104in}{1.141075in}}%
\pgfpathlineto{\pgfqpoint{0.955188in}{1.129431in}}%
\pgfpathlineto{\pgfqpoint{0.947252in}{1.118052in}}%
\pgfpathclose%
\pgfusepath{fill}%
\end{pgfscope}%
\begin{pgfscope}%
\pgfpathrectangle{\pgfqpoint{0.407592in}{0.222904in}}{\pgfqpoint{1.411163in}{1.411163in}}%
\pgfusepath{clip}%
\pgfsetbuttcap%
\pgfsetroundjoin%
\definecolor{currentfill}{rgb}{0.556863,0.654902,0.847059}%
\pgfsetfillcolor{currentfill}%
\pgfsetfillopacity{0.130000}%
\pgfsetlinewidth{0.000000pt}%
\definecolor{currentstroke}{rgb}{0.000000,0.000000,0.000000}%
\pgfsetstrokecolor{currentstroke}%
\pgfsetdash{}{0pt}%
\pgfpathmoveto{\pgfqpoint{1.025645in}{1.080735in}}%
\pgfpathlineto{\pgfqpoint{1.041317in}{1.066726in}}%
\pgfpathlineto{\pgfqpoint{1.057007in}{1.053238in}}%
\pgfpathlineto{\pgfqpoint{1.064823in}{1.066785in}}%
\pgfpathlineto{\pgfqpoint{1.072622in}{1.080538in}}%
\pgfpathlineto{\pgfqpoint{1.056955in}{1.093462in}}%
\pgfpathlineto{\pgfqpoint{1.041307in}{1.106910in}}%
\pgfpathlineto{\pgfqpoint{1.033485in}{1.093713in}}%
\pgfpathlineto{\pgfqpoint{1.025645in}{1.080735in}}%
\pgfpathclose%
\pgfusepath{fill}%
\end{pgfscope}%
\begin{pgfscope}%
\pgfpathrectangle{\pgfqpoint{0.407592in}{0.222904in}}{\pgfqpoint{1.411163in}{1.411163in}}%
\pgfusepath{clip}%
\pgfsetbuttcap%
\pgfsetroundjoin%
\definecolor{currentfill}{rgb}{0.556863,0.654902,0.847059}%
\pgfsetfillcolor{currentfill}%
\pgfsetfillopacity{0.130000}%
\pgfsetlinewidth{0.000000pt}%
\definecolor{currentstroke}{rgb}{0.000000,0.000000,0.000000}%
\pgfsetstrokecolor{currentstroke}%
\pgfsetdash{}{0pt}%
\pgfpathmoveto{\pgfqpoint{0.868714in}{1.165395in}}%
\pgfpathlineto{\pgfqpoint{0.884348in}{1.147768in}}%
\pgfpathlineto{\pgfqpoint{0.899994in}{1.130337in}}%
\pgfpathlineto{\pgfqpoint{0.908007in}{1.140652in}}%
\pgfpathlineto{\pgfqpoint{0.915998in}{1.151260in}}%
\pgfpathlineto{\pgfqpoint{0.900387in}{1.168237in}}%
\pgfpathlineto{\pgfqpoint{0.884788in}{1.185431in}}%
\pgfpathlineto{\pgfqpoint{0.876762in}{1.175272in}}%
\pgfpathlineto{\pgfqpoint{0.868714in}{1.165395in}}%
\pgfpathclose%
\pgfusepath{fill}%
\end{pgfscope}%
\begin{pgfscope}%
\pgfpathrectangle{\pgfqpoint{0.407592in}{0.222904in}}{\pgfqpoint{1.411163in}{1.411163in}}%
\pgfusepath{clip}%
\pgfsetbuttcap%
\pgfsetroundjoin%
\definecolor{currentfill}{rgb}{0.556863,0.654902,0.847059}%
\pgfsetfillcolor{currentfill}%
\pgfsetfillopacity{0.130000}%
\pgfsetlinewidth{0.000000pt}%
\definecolor{currentstroke}{rgb}{0.000000,0.000000,0.000000}%
\pgfsetstrokecolor{currentstroke}%
\pgfsetdash{}{0pt}%
\pgfpathmoveto{\pgfqpoint{0.758710in}{1.257496in}}%
\pgfpathlineto{\pgfqpoint{0.774326in}{1.238746in}}%
\pgfpathlineto{\pgfqpoint{0.789951in}{1.220083in}}%
\pgfpathlineto{\pgfqpoint{0.798125in}{1.228643in}}%
\pgfpathlineto{\pgfqpoint{0.806273in}{1.237483in}}%
\pgfpathlineto{\pgfqpoint{0.790687in}{1.255823in}}%
\pgfpathlineto{\pgfqpoint{0.775110in}{1.274265in}}%
\pgfpathlineto{\pgfqpoint{0.766923in}{1.265747in}}%
\pgfpathlineto{\pgfqpoint{0.758710in}{1.257496in}}%
\pgfpathclose%
\pgfusepath{fill}%
\end{pgfscope}%
\begin{pgfscope}%
\pgfpathrectangle{\pgfqpoint{0.407592in}{0.222904in}}{\pgfqpoint{1.411163in}{1.411163in}}%
\pgfusepath{clip}%
\pgfsetbuttcap%
\pgfsetroundjoin%
\definecolor{currentfill}{rgb}{0.556863,0.654902,0.847059}%
\pgfsetfillcolor{currentfill}%
\pgfsetfillopacity{0.130000}%
\pgfsetlinewidth{0.000000pt}%
\definecolor{currentstroke}{rgb}{0.000000,0.000000,0.000000}%
\pgfsetstrokecolor{currentstroke}%
\pgfsetdash{}{0pt}%
\pgfpathmoveto{\pgfqpoint{1.150979in}{1.065082in}}%
\pgfpathlineto{\pgfqpoint{1.166746in}{1.055936in}}%
\pgfpathlineto{\pgfqpoint{1.182543in}{1.047567in}}%
\pgfpathlineto{\pgfqpoint{1.190247in}{1.063144in}}%
\pgfpathlineto{\pgfqpoint{1.197937in}{1.078806in}}%
\pgfpathlineto{\pgfqpoint{1.182153in}{1.086828in}}%
\pgfpathlineto{\pgfqpoint{1.166400in}{1.095585in}}%
\pgfpathlineto{\pgfqpoint{1.158696in}{1.080282in}}%
\pgfpathlineto{\pgfqpoint{1.150979in}{1.065082in}}%
\pgfpathclose%
\pgfusepath{fill}%
\end{pgfscope}%
\begin{pgfscope}%
\pgfpathrectangle{\pgfqpoint{0.407592in}{0.222904in}}{\pgfqpoint{1.411163in}{1.411163in}}%
\pgfusepath{clip}%
\pgfsetbuttcap%
\pgfsetroundjoin%
\definecolor{currentfill}{rgb}{0.556863,0.654902,0.847059}%
\pgfsetfillcolor{currentfill}%
\pgfsetfillopacity{0.130000}%
\pgfsetlinewidth{0.000000pt}%
\definecolor{currentstroke}{rgb}{0.000000,0.000000,0.000000}%
\pgfsetstrokecolor{currentstroke}%
\pgfsetdash{}{0pt}%
\pgfpathmoveto{\pgfqpoint{0.978555in}{1.086041in}}%
\pgfpathlineto{\pgfqpoint{0.994226in}{1.070572in}}%
\pgfpathlineto{\pgfqpoint{1.009910in}{1.055518in}}%
\pgfpathlineto{\pgfqpoint{1.017787in}{1.067997in}}%
\pgfpathlineto{\pgfqpoint{1.025645in}{1.080735in}}%
\pgfpathlineto{\pgfqpoint{1.009988in}{1.095217in}}%
\pgfpathlineto{\pgfqpoint{0.994346in}{1.110128in}}%
\pgfpathlineto{\pgfqpoint{0.986460in}{1.097952in}}%
\pgfpathlineto{\pgfqpoint{0.978555in}{1.086041in}}%
\pgfpathclose%
\pgfusepath{fill}%
\end{pgfscope}%
\begin{pgfscope}%
\pgfpathrectangle{\pgfqpoint{0.407592in}{0.222904in}}{\pgfqpoint{1.411163in}{1.411163in}}%
\pgfusepath{clip}%
\pgfsetbuttcap%
\pgfsetroundjoin%
\definecolor{currentfill}{rgb}{0.556863,0.654902,0.847059}%
\pgfsetfillcolor{currentfill}%
\pgfsetfillopacity{0.130000}%
\pgfsetlinewidth{0.000000pt}%
\definecolor{currentstroke}{rgb}{0.000000,0.000000,0.000000}%
\pgfsetstrokecolor{currentstroke}%
\pgfsetdash{}{0pt}%
\pgfpathmoveto{\pgfqpoint{1.104016in}{1.056450in}}%
\pgfpathlineto{\pgfqpoint{1.119746in}{1.045385in}}%
\pgfpathlineto{\pgfqpoint{1.135502in}{1.035041in}}%
\pgfpathlineto{\pgfqpoint{1.143247in}{1.049996in}}%
\pgfpathlineto{\pgfqpoint{1.150979in}{1.065082in}}%
\pgfpathlineto{\pgfqpoint{1.135240in}{1.074963in}}%
\pgfpathlineto{\pgfqpoint{1.119527in}{1.085537in}}%
\pgfpathlineto{\pgfqpoint{1.111779in}{1.070918in}}%
\pgfpathlineto{\pgfqpoint{1.104016in}{1.056450in}}%
\pgfpathclose%
\pgfusepath{fill}%
\end{pgfscope}%
\begin{pgfscope}%
\pgfpathrectangle{\pgfqpoint{0.407592in}{0.222904in}}{\pgfqpoint{1.411163in}{1.411163in}}%
\pgfusepath{clip}%
\pgfsetbuttcap%
\pgfsetroundjoin%
\definecolor{currentfill}{rgb}{0.556863,0.654902,0.847059}%
\pgfsetfillcolor{currentfill}%
\pgfsetfillopacity{0.130000}%
\pgfsetlinewidth{0.000000pt}%
\definecolor{currentstroke}{rgb}{0.000000,0.000000,0.000000}%
\pgfsetstrokecolor{currentstroke}%
\pgfsetdash{}{0pt}%
\pgfpathmoveto{\pgfqpoint{0.789951in}{1.220083in}}%
\pgfpathlineto{\pgfqpoint{0.805585in}{1.201520in}}%
\pgfpathlineto{\pgfqpoint{0.821228in}{1.183066in}}%
\pgfpathlineto{\pgfqpoint{0.829364in}{1.191964in}}%
\pgfpathlineto{\pgfqpoint{0.837474in}{1.201157in}}%
\pgfpathlineto{\pgfqpoint{0.821869in}{1.219256in}}%
\pgfpathlineto{\pgfqpoint{0.806273in}{1.237483in}}%
\pgfpathlineto{\pgfqpoint{0.798125in}{1.228643in}}%
\pgfpathlineto{\pgfqpoint{0.789951in}{1.220083in}}%
\pgfpathclose%
\pgfusepath{fill}%
\end{pgfscope}%
\begin{pgfscope}%
\pgfpathrectangle{\pgfqpoint{0.407592in}{0.222904in}}{\pgfqpoint{1.411163in}{1.411163in}}%
\pgfusepath{clip}%
\pgfsetbuttcap%
\pgfsetroundjoin%
\definecolor{currentfill}{rgb}{0.556863,0.654902,0.847059}%
\pgfsetfillcolor{currentfill}%
\pgfsetfillopacity{0.130000}%
\pgfsetlinewidth{0.000000pt}%
\definecolor{currentstroke}{rgb}{0.000000,0.000000,0.000000}%
\pgfsetstrokecolor{currentstroke}%
\pgfsetdash{}{0pt}%
\pgfpathmoveto{\pgfqpoint{1.197937in}{1.078806in}}%
\pgfpathlineto{\pgfqpoint{1.213755in}{1.071550in}}%
\pgfpathlineto{\pgfqpoint{1.229608in}{1.065087in}}%
\pgfpathlineto{\pgfqpoint{1.237274in}{1.081075in}}%
\pgfpathlineto{\pgfqpoint{1.244927in}{1.097116in}}%
\pgfpathlineto{\pgfqpoint{1.229085in}{1.103356in}}%
\pgfpathlineto{\pgfqpoint{1.213278in}{1.110341in}}%
\pgfpathlineto{\pgfqpoint{1.205614in}{1.094541in}}%
\pgfpathlineto{\pgfqpoint{1.197937in}{1.078806in}}%
\pgfpathclose%
\pgfusepath{fill}%
\end{pgfscope}%
\begin{pgfscope}%
\pgfpathrectangle{\pgfqpoint{0.407592in}{0.222904in}}{\pgfqpoint{1.411163in}{1.411163in}}%
\pgfusepath{clip}%
\pgfsetbuttcap%
\pgfsetroundjoin%
\definecolor{currentfill}{rgb}{0.556863,0.654902,0.847059}%
\pgfsetfillcolor{currentfill}%
\pgfsetfillopacity{0.130000}%
\pgfsetlinewidth{0.000000pt}%
\definecolor{currentstroke}{rgb}{0.000000,0.000000,0.000000}%
\pgfsetstrokecolor{currentstroke}%
\pgfsetdash{}{0pt}%
\pgfpathmoveto{\pgfqpoint{0.899994in}{1.130337in}}%
\pgfpathlineto{\pgfqpoint{0.915649in}{1.113125in}}%
\pgfpathlineto{\pgfqpoint{0.931316in}{1.096159in}}%
\pgfpathlineto{\pgfqpoint{0.939295in}{1.106956in}}%
\pgfpathlineto{\pgfqpoint{0.947252in}{1.118052in}}%
\pgfpathlineto{\pgfqpoint{0.931619in}{1.134522in}}%
\pgfpathlineto{\pgfqpoint{0.915998in}{1.151260in}}%
\pgfpathlineto{\pgfqpoint{0.908007in}{1.140652in}}%
\pgfpathlineto{\pgfqpoint{0.899994in}{1.130337in}}%
\pgfpathclose%
\pgfusepath{fill}%
\end{pgfscope}%
\begin{pgfscope}%
\pgfpathrectangle{\pgfqpoint{0.407592in}{0.222904in}}{\pgfqpoint{1.411163in}{1.411163in}}%
\pgfusepath{clip}%
\pgfsetbuttcap%
\pgfsetroundjoin%
\definecolor{currentfill}{rgb}{0.556863,0.654902,0.847059}%
\pgfsetfillcolor{currentfill}%
\pgfsetfillopacity{0.130000}%
\pgfsetlinewidth{0.000000pt}%
\definecolor{currentstroke}{rgb}{0.000000,0.000000,0.000000}%
\pgfsetstrokecolor{currentstroke}%
\pgfsetdash{}{0pt}%
\pgfpathmoveto{\pgfqpoint{1.057007in}{1.053238in}}%
\pgfpathlineto{\pgfqpoint{1.072716in}{1.040324in}}%
\pgfpathlineto{\pgfqpoint{1.088444in}{1.028039in}}%
\pgfpathlineto{\pgfqpoint{1.096238in}{1.042151in}}%
\pgfpathlineto{\pgfqpoint{1.104016in}{1.056450in}}%
\pgfpathlineto{\pgfqpoint{1.088308in}{1.068185in}}%
\pgfpathlineto{\pgfqpoint{1.072622in}{1.080538in}}%
\pgfpathlineto{\pgfqpoint{1.064823in}{1.066785in}}%
\pgfpathlineto{\pgfqpoint{1.057007in}{1.053238in}}%
\pgfpathclose%
\pgfusepath{fill}%
\end{pgfscope}%
\begin{pgfscope}%
\pgfpathrectangle{\pgfqpoint{0.407592in}{0.222904in}}{\pgfqpoint{1.411163in}{1.411163in}}%
\pgfusepath{clip}%
\pgfsetbuttcap%
\pgfsetroundjoin%
\definecolor{currentfill}{rgb}{0.556863,0.654902,0.847059}%
\pgfsetfillcolor{currentfill}%
\pgfsetfillopacity{0.130000}%
\pgfsetlinewidth{0.000000pt}%
\definecolor{currentstroke}{rgb}{0.000000,0.000000,0.000000}%
\pgfsetstrokecolor{currentstroke}%
\pgfsetdash{}{0pt}%
\pgfpathmoveto{\pgfqpoint{0.821228in}{1.183066in}}%
\pgfpathlineto{\pgfqpoint{0.836881in}{1.164736in}}%
\pgfpathlineto{\pgfqpoint{0.852543in}{1.146543in}}%
\pgfpathlineto{\pgfqpoint{0.860641in}{1.155815in}}%
\pgfpathlineto{\pgfqpoint{0.868714in}{1.165395in}}%
\pgfpathlineto{\pgfqpoint{0.853089in}{1.183198in}}%
\pgfpathlineto{\pgfqpoint{0.837474in}{1.201157in}}%
\pgfpathlineto{\pgfqpoint{0.829364in}{1.191964in}}%
\pgfpathlineto{\pgfqpoint{0.821228in}{1.183066in}}%
\pgfpathclose%
\pgfusepath{fill}%
\end{pgfscope}%
\begin{pgfscope}%
\pgfpathrectangle{\pgfqpoint{0.407592in}{0.222904in}}{\pgfqpoint{1.411163in}{1.411163in}}%
\pgfusepath{clip}%
\pgfsetbuttcap%
\pgfsetroundjoin%
\definecolor{currentfill}{rgb}{0.556863,0.654902,0.847059}%
\pgfsetfillcolor{currentfill}%
\pgfsetfillopacity{0.130000}%
\pgfsetlinewidth{0.000000pt}%
\definecolor{currentstroke}{rgb}{0.000000,0.000000,0.000000}%
\pgfsetstrokecolor{currentstroke}%
\pgfsetdash{}{0pt}%
\pgfpathmoveto{\pgfqpoint{1.244927in}{1.097116in}}%
\pgfpathlineto{\pgfqpoint{1.260808in}{1.091640in}}%
\pgfpathlineto{\pgfqpoint{1.276730in}{1.086941in}}%
\pgfpathlineto{\pgfqpoint{1.284361in}{1.103162in}}%
\pgfpathlineto{\pgfqpoint{1.268444in}{1.107805in}}%
\pgfpathlineto{\pgfqpoint{1.252568in}{1.113201in}}%
\pgfusepath{fill}%
\end{pgfscope}%
\begin{pgfscope}%
\pgfpathrectangle{\pgfqpoint{0.407592in}{0.222904in}}{\pgfqpoint{1.411163in}{1.411163in}}%
\pgfusepath{clip}%
\pgfsetbuttcap%
\pgfsetroundjoin%
\definecolor{currentfill}{rgb}{0.556863,0.654902,0.847059}%
\pgfsetfillcolor{currentfill}%
\pgfsetfillopacity{0.130000}%
\pgfsetlinewidth{0.000000pt}%
\definecolor{currentstroke}{rgb}{0.000000,0.000000,0.000000}%
\pgfsetstrokecolor{currentstroke}%
\pgfsetdash{}{0pt}%
\pgfpathmoveto{\pgfqpoint{0.931316in}{1.096159in}}%
\pgfpathlineto{\pgfqpoint{0.946994in}{1.079469in}}%
\pgfpathlineto{\pgfqpoint{0.962683in}{1.063091in}}%
\pgfpathlineto{\pgfqpoint{0.970629in}{1.074413in}}%
\pgfpathlineto{\pgfqpoint{0.978555in}{1.086041in}}%
\pgfpathlineto{\pgfqpoint{0.962897in}{1.101880in}}%
\pgfpathlineto{\pgfqpoint{0.947252in}{1.118052in}}%
\pgfpathlineto{\pgfqpoint{0.939295in}{1.106956in}}%
\pgfpathlineto{\pgfqpoint{0.931316in}{1.096159in}}%
\pgfpathclose%
\pgfusepath{fill}%
\end{pgfscope}%
\begin{pgfscope}%
\pgfpathrectangle{\pgfqpoint{0.407592in}{0.222904in}}{\pgfqpoint{1.411163in}{1.411163in}}%
\pgfusepath{clip}%
\pgfsetbuttcap%
\pgfsetroundjoin%
\definecolor{currentfill}{rgb}{0.556863,0.654902,0.847059}%
\pgfsetfillcolor{currentfill}%
\pgfsetfillopacity{0.130000}%
\pgfsetlinewidth{0.000000pt}%
\definecolor{currentstroke}{rgb}{0.000000,0.000000,0.000000}%
\pgfsetstrokecolor{currentstroke}%
\pgfsetdash{}{0pt}%
\pgfpathmoveto{\pgfqpoint{1.009910in}{1.055518in}}%
\pgfpathlineto{\pgfqpoint{1.025610in}{1.040925in}}%
\pgfpathlineto{\pgfqpoint{1.041325in}{1.026847in}}%
\pgfpathlineto{\pgfqpoint{1.049175in}{1.039918in}}%
\pgfpathlineto{\pgfqpoint{1.057007in}{1.053238in}}%
\pgfpathlineto{\pgfqpoint{1.041317in}{1.066726in}}%
\pgfpathlineto{\pgfqpoint{1.025645in}{1.080735in}}%
\pgfpathlineto{\pgfqpoint{1.017787in}{1.067997in}}%
\pgfpathlineto{\pgfqpoint{1.009910in}{1.055518in}}%
\pgfpathclose%
\pgfusepath{fill}%
\end{pgfscope}%
\begin{pgfscope}%
\pgfpathrectangle{\pgfqpoint{0.407592in}{0.222904in}}{\pgfqpoint{1.411163in}{1.411163in}}%
\pgfusepath{clip}%
\pgfsetbuttcap%
\pgfsetroundjoin%
\definecolor{currentfill}{rgb}{0.556863,0.654902,0.847059}%
\pgfsetfillcolor{currentfill}%
\pgfsetfillopacity{0.130000}%
\pgfsetlinewidth{0.000000pt}%
\definecolor{currentstroke}{rgb}{0.000000,0.000000,0.000000}%
\pgfsetstrokecolor{currentstroke}%
\pgfsetdash{}{0pt}%
\pgfpathmoveto{\pgfqpoint{0.852543in}{1.146543in}}%
\pgfpathlineto{\pgfqpoint{0.868215in}{1.128507in}}%
\pgfpathlineto{\pgfqpoint{0.883896in}{1.110644in}}%
\pgfpathlineto{\pgfqpoint{0.891957in}{1.120329in}}%
\pgfpathlineto{\pgfqpoint{0.899994in}{1.130337in}}%
\pgfpathlineto{\pgfqpoint{0.884348in}{1.147768in}}%
\pgfpathlineto{\pgfqpoint{0.868714in}{1.165395in}}%
\pgfpathlineto{\pgfqpoint{0.860641in}{1.155815in}}%
\pgfpathlineto{\pgfqpoint{0.852543in}{1.146543in}}%
\pgfpathclose%
\pgfusepath{fill}%
\end{pgfscope}%
\begin{pgfscope}%
\pgfpathrectangle{\pgfqpoint{0.407592in}{0.222904in}}{\pgfqpoint{1.411163in}{1.411163in}}%
\pgfusepath{clip}%
\pgfsetbuttcap%
\pgfsetroundjoin%
\definecolor{currentfill}{rgb}{0.556863,0.654902,0.847059}%
\pgfsetfillcolor{currentfill}%
\pgfsetfillopacity{0.130000}%
\pgfsetlinewidth{0.000000pt}%
\definecolor{currentstroke}{rgb}{0.000000,0.000000,0.000000}%
\pgfsetstrokecolor{currentstroke}%
\pgfsetdash{}{0pt}%
\pgfpathmoveto{\pgfqpoint{1.135502in}{1.035041in}}%
\pgfpathlineto{\pgfqpoint{1.151283in}{1.025468in}}%
\pgfpathlineto{\pgfqpoint{1.167094in}{1.016715in}}%
\pgfpathlineto{\pgfqpoint{1.174825in}{1.032086in}}%
\pgfpathlineto{\pgfqpoint{1.182543in}{1.047567in}}%
\pgfpathlineto{\pgfqpoint{1.166746in}{1.055936in}}%
\pgfpathlineto{\pgfqpoint{1.150979in}{1.065082in}}%
\pgfpathlineto{\pgfqpoint{1.143247in}{1.049996in}}%
\pgfpathlineto{\pgfqpoint{1.135502in}{1.035041in}}%
\pgfpathclose%
\pgfusepath{fill}%
\end{pgfscope}%
\begin{pgfscope}%
\pgfpathrectangle{\pgfqpoint{0.407592in}{0.222904in}}{\pgfqpoint{1.411163in}{1.411163in}}%
\pgfusepath{clip}%
\pgfsetbuttcap%
\pgfsetroundjoin%
\definecolor{currentfill}{rgb}{0.556863,0.654902,0.847059}%
\pgfsetfillcolor{currentfill}%
\pgfsetfillopacity{0.130000}%
\pgfsetlinewidth{0.000000pt}%
\definecolor{currentstroke}{rgb}{0.000000,0.000000,0.000000}%
\pgfsetstrokecolor{currentstroke}%
\pgfsetdash{}{0pt}%
\pgfpathmoveto{\pgfqpoint{1.182543in}{1.047567in}}%
\pgfpathlineto{\pgfqpoint{1.198372in}{1.040011in}}%
\pgfpathlineto{\pgfqpoint{1.214237in}{1.033302in}}%
\pgfpathlineto{\pgfqpoint{1.221929in}{1.049159in}}%
\pgfpathlineto{\pgfqpoint{1.229608in}{1.065087in}}%
\pgfpathlineto{\pgfqpoint{1.213755in}{1.071550in}}%
\pgfpathlineto{\pgfqpoint{1.197937in}{1.078806in}}%
\pgfpathlineto{\pgfqpoint{1.190247in}{1.063144in}}%
\pgfpathlineto{\pgfqpoint{1.182543in}{1.047567in}}%
\pgfpathclose%
\pgfusepath{fill}%
\end{pgfscope}%
\begin{pgfscope}%
\pgfpathrectangle{\pgfqpoint{0.407592in}{0.222904in}}{\pgfqpoint{1.411163in}{1.411163in}}%
\pgfusepath{clip}%
\pgfsetbuttcap%
\pgfsetroundjoin%
\definecolor{currentfill}{rgb}{0.556863,0.654902,0.847059}%
\pgfsetfillcolor{currentfill}%
\pgfsetfillopacity{0.130000}%
\pgfsetlinewidth{0.000000pt}%
\definecolor{currentstroke}{rgb}{0.000000,0.000000,0.000000}%
\pgfsetstrokecolor{currentstroke}%
\pgfsetdash{}{0pt}%
\pgfpathmoveto{\pgfqpoint{0.742202in}{1.241828in}}%
\pgfpathlineto{\pgfqpoint{0.757857in}{1.222798in}}%
\pgfpathlineto{\pgfqpoint{0.773522in}{1.203843in}}%
\pgfpathlineto{\pgfqpoint{0.781750in}{1.211814in}}%
\pgfpathlineto{\pgfqpoint{0.789951in}{1.220083in}}%
\pgfpathlineto{\pgfqpoint{0.774326in}{1.238746in}}%
\pgfpathlineto{\pgfqpoint{0.758710in}{1.257496in}}%
\pgfpathlineto{\pgfqpoint{0.750470in}{1.249521in}}%
\pgfpathlineto{\pgfqpoint{0.742202in}{1.241828in}}%
\pgfpathclose%
\pgfusepath{fill}%
\end{pgfscope}%
\begin{pgfscope}%
\pgfpathrectangle{\pgfqpoint{0.407592in}{0.222904in}}{\pgfqpoint{1.411163in}{1.411163in}}%
\pgfusepath{clip}%
\pgfsetbuttcap%
\pgfsetroundjoin%
\definecolor{currentfill}{rgb}{0.556863,0.654902,0.847059}%
\pgfsetfillcolor{currentfill}%
\pgfsetfillopacity{0.130000}%
\pgfsetlinewidth{0.000000pt}%
\definecolor{currentstroke}{rgb}{0.000000,0.000000,0.000000}%
\pgfsetstrokecolor{currentstroke}%
\pgfsetdash{}{0pt}%
\pgfpathmoveto{\pgfqpoint{1.088444in}{1.028039in}}%
\pgfpathlineto{\pgfqpoint{1.104193in}{1.016442in}}%
\pgfpathlineto{\pgfqpoint{1.119966in}{1.005592in}}%
\pgfpathlineto{\pgfqpoint{1.127741in}{1.020233in}}%
\pgfpathlineto{\pgfqpoint{1.135502in}{1.035041in}}%
\pgfpathlineto{\pgfqpoint{1.119746in}{1.045385in}}%
\pgfpathlineto{\pgfqpoint{1.104016in}{1.056450in}}%
\pgfpathlineto{\pgfqpoint{1.096238in}{1.042151in}}%
\pgfpathlineto{\pgfqpoint{1.088444in}{1.028039in}}%
\pgfpathclose%
\pgfusepath{fill}%
\end{pgfscope}%
\begin{pgfscope}%
\pgfpathrectangle{\pgfqpoint{0.407592in}{0.222904in}}{\pgfqpoint{1.411163in}{1.411163in}}%
\pgfusepath{clip}%
\pgfsetbuttcap%
\pgfsetroundjoin%
\definecolor{currentfill}{rgb}{0.556863,0.654902,0.847059}%
\pgfsetfillcolor{currentfill}%
\pgfsetfillopacity{0.130000}%
\pgfsetlinewidth{0.000000pt}%
\definecolor{currentstroke}{rgb}{0.000000,0.000000,0.000000}%
\pgfsetstrokecolor{currentstroke}%
\pgfsetdash{}{0pt}%
\pgfpathmoveto{\pgfqpoint{0.962683in}{1.063091in}}%
\pgfpathlineto{\pgfqpoint{0.978385in}{1.047062in}}%
\pgfpathlineto{\pgfqpoint{0.994100in}{1.031429in}}%
\pgfpathlineto{\pgfqpoint{1.002015in}{1.043321in}}%
\pgfpathlineto{\pgfqpoint{1.009910in}{1.055518in}}%
\pgfpathlineto{\pgfqpoint{0.994226in}{1.070572in}}%
\pgfpathlineto{\pgfqpoint{0.978555in}{1.086041in}}%
\pgfpathlineto{\pgfqpoint{0.970629in}{1.074413in}}%
\pgfpathlineto{\pgfqpoint{0.962683in}{1.063091in}}%
\pgfpathclose%
\pgfusepath{fill}%
\end{pgfscope}%
\begin{pgfscope}%
\pgfpathrectangle{\pgfqpoint{0.407592in}{0.222904in}}{\pgfqpoint{1.411163in}{1.411163in}}%
\pgfusepath{clip}%
\pgfsetbuttcap%
\pgfsetroundjoin%
\definecolor{currentfill}{rgb}{0.556863,0.654902,0.847059}%
\pgfsetfillcolor{currentfill}%
\pgfsetfillopacity{0.130000}%
\pgfsetlinewidth{0.000000pt}%
\definecolor{currentstroke}{rgb}{0.000000,0.000000,0.000000}%
\pgfsetstrokecolor{currentstroke}%
\pgfsetdash{}{0pt}%
\pgfpathmoveto{\pgfqpoint{1.229608in}{1.065087in}}%
\pgfpathlineto{\pgfqpoint{1.245499in}{1.059439in}}%
\pgfpathlineto{\pgfqpoint{1.261431in}{1.054624in}}%
\pgfpathlineto{\pgfqpoint{1.269087in}{1.070760in}}%
\pgfpathlineto{\pgfqpoint{1.276730in}{1.086941in}}%
\pgfpathlineto{\pgfqpoint{1.260808in}{1.091640in}}%
\pgfpathlineto{\pgfqpoint{1.244927in}{1.097116in}}%
\pgfpathlineto{\pgfqpoint{1.237274in}{1.081075in}}%
\pgfpathlineto{\pgfqpoint{1.229608in}{1.065087in}}%
\pgfpathclose%
\pgfusepath{fill}%
\end{pgfscope}%
\begin{pgfscope}%
\pgfpathrectangle{\pgfqpoint{0.407592in}{0.222904in}}{\pgfqpoint{1.411163in}{1.411163in}}%
\pgfusepath{clip}%
\pgfsetbuttcap%
\pgfsetroundjoin%
\definecolor{currentfill}{rgb}{0.556863,0.654902,0.847059}%
\pgfsetfillcolor{currentfill}%
\pgfsetfillopacity{0.130000}%
\pgfsetlinewidth{0.000000pt}%
\definecolor{currentstroke}{rgb}{0.000000,0.000000,0.000000}%
\pgfsetstrokecolor{currentstroke}%
\pgfsetdash{}{0pt}%
\pgfpathmoveto{\pgfqpoint{1.041325in}{1.026847in}}%
\pgfpathlineto{\pgfqpoint{1.057057in}{1.013340in}}%
\pgfpathlineto{\pgfqpoint{1.072808in}{1.000468in}}%
\pgfpathlineto{\pgfqpoint{1.080634in}{1.014137in}}%
\pgfpathlineto{\pgfqpoint{1.088444in}{1.028039in}}%
\pgfpathlineto{\pgfqpoint{1.072716in}{1.040324in}}%
\pgfpathlineto{\pgfqpoint{1.057007in}{1.053238in}}%
\pgfpathlineto{\pgfqpoint{1.049175in}{1.039918in}}%
\pgfpathlineto{\pgfqpoint{1.041325in}{1.026847in}}%
\pgfpathclose%
\pgfusepath{fill}%
\end{pgfscope}%
\begin{pgfscope}%
\pgfpathrectangle{\pgfqpoint{0.407592in}{0.222904in}}{\pgfqpoint{1.411163in}{1.411163in}}%
\pgfusepath{clip}%
\pgfsetbuttcap%
\pgfsetroundjoin%
\definecolor{currentfill}{rgb}{0.556863,0.654902,0.847059}%
\pgfsetfillcolor{currentfill}%
\pgfsetfillopacity{0.130000}%
\pgfsetlinewidth{0.000000pt}%
\definecolor{currentstroke}{rgb}{0.000000,0.000000,0.000000}%
\pgfsetstrokecolor{currentstroke}%
\pgfsetdash{}{0pt}%
\pgfpathmoveto{\pgfqpoint{0.883896in}{1.110644in}}%
\pgfpathlineto{\pgfqpoint{0.899588in}{1.092980in}}%
\pgfpathlineto{\pgfqpoint{0.915290in}{1.075537in}}%
\pgfpathlineto{\pgfqpoint{0.923314in}{1.085680in}}%
\pgfpathlineto{\pgfqpoint{0.931316in}{1.096159in}}%
\pgfpathlineto{\pgfqpoint{0.915649in}{1.113125in}}%
\pgfpathlineto{\pgfqpoint{0.899994in}{1.130337in}}%
\pgfpathlineto{\pgfqpoint{0.891957in}{1.120329in}}%
\pgfpathlineto{\pgfqpoint{0.883896in}{1.110644in}}%
\pgfpathclose%
\pgfusepath{fill}%
\end{pgfscope}%
\begin{pgfscope}%
\pgfpathrectangle{\pgfqpoint{0.407592in}{0.222904in}}{\pgfqpoint{1.411163in}{1.411163in}}%
\pgfusepath{clip}%
\pgfsetbuttcap%
\pgfsetroundjoin%
\definecolor{currentfill}{rgb}{0.556863,0.654902,0.847059}%
\pgfsetfillcolor{currentfill}%
\pgfsetfillopacity{0.130000}%
\pgfsetlinewidth{0.000000pt}%
\definecolor{currentstroke}{rgb}{0.000000,0.000000,0.000000}%
\pgfsetstrokecolor{currentstroke}%
\pgfsetdash{}{0pt}%
\pgfpathmoveto{\pgfqpoint{0.773522in}{1.203843in}}%
\pgfpathlineto{\pgfqpoint{0.789195in}{1.184972in}}%
\pgfpathlineto{\pgfqpoint{0.804877in}{1.166196in}}%
\pgfpathlineto{\pgfqpoint{0.813066in}{1.174473in}}%
\pgfpathlineto{\pgfqpoint{0.821228in}{1.183066in}}%
\pgfpathlineto{\pgfqpoint{0.805585in}{1.201520in}}%
\pgfpathlineto{\pgfqpoint{0.789951in}{1.220083in}}%
\pgfpathlineto{\pgfqpoint{0.781750in}{1.211814in}}%
\pgfpathlineto{\pgfqpoint{0.773522in}{1.203843in}}%
\pgfpathclose%
\pgfusepath{fill}%
\end{pgfscope}%
\begin{pgfscope}%
\pgfpathrectangle{\pgfqpoint{0.407592in}{0.222904in}}{\pgfqpoint{1.411163in}{1.411163in}}%
\pgfusepath{clip}%
\pgfsetbuttcap%
\pgfsetroundjoin%
\definecolor{currentfill}{rgb}{0.556863,0.654902,0.847059}%
\pgfsetfillcolor{currentfill}%
\pgfsetfillopacity{0.130000}%
\pgfsetlinewidth{0.000000pt}%
\definecolor{currentstroke}{rgb}{0.000000,0.000000,0.000000}%
\pgfsetstrokecolor{currentstroke}%
\pgfsetdash{}{0pt}%
\pgfpathmoveto{\pgfqpoint{1.276730in}{1.086941in}}%
\pgfpathlineto{\pgfqpoint{1.292696in}{1.083029in}}%
\pgfpathlineto{\pgfqpoint{1.308706in}{1.079908in}}%
\pgfpathlineto{\pgfqpoint{1.316328in}{1.096162in}}%
\pgfpathlineto{\pgfqpoint{1.300322in}{1.099279in}}%
\pgfpathlineto{\pgfqpoint{1.284361in}{1.103162in}}%
\pgfusepath{fill}%
\end{pgfscope}%
\begin{pgfscope}%
\pgfpathrectangle{\pgfqpoint{0.407592in}{0.222904in}}{\pgfqpoint{1.411163in}{1.411163in}}%
\pgfusepath{clip}%
\pgfsetbuttcap%
\pgfsetroundjoin%
\definecolor{currentfill}{rgb}{0.556863,0.654902,0.847059}%
\pgfsetfillcolor{currentfill}%
\pgfsetfillopacity{0.130000}%
\pgfsetlinewidth{0.000000pt}%
\definecolor{currentstroke}{rgb}{0.000000,0.000000,0.000000}%
\pgfsetstrokecolor{currentstroke}%
\pgfsetdash{}{0pt}%
\pgfpathmoveto{\pgfqpoint{0.994100in}{1.031429in}}%
\pgfpathlineto{\pgfqpoint{1.009828in}{1.016239in}}%
\pgfpathlineto{\pgfqpoint{1.025571in}{1.001551in}}%
\pgfpathlineto{\pgfqpoint{1.033457in}{1.014049in}}%
\pgfpathlineto{\pgfqpoint{1.041325in}{1.026847in}}%
\pgfpathlineto{\pgfqpoint{1.025610in}{1.040925in}}%
\pgfpathlineto{\pgfqpoint{1.009910in}{1.055518in}}%
\pgfpathlineto{\pgfqpoint{1.002015in}{1.043321in}}%
\pgfpathlineto{\pgfqpoint{0.994100in}{1.031429in}}%
\pgfpathclose%
\pgfusepath{fill}%
\end{pgfscope}%
\begin{pgfscope}%
\pgfpathrectangle{\pgfqpoint{0.407592in}{0.222904in}}{\pgfqpoint{1.411163in}{1.411163in}}%
\pgfusepath{clip}%
\pgfsetbuttcap%
\pgfsetroundjoin%
\definecolor{currentfill}{rgb}{0.556863,0.654902,0.847059}%
\pgfsetfillcolor{currentfill}%
\pgfsetfillopacity{0.130000}%
\pgfsetlinewidth{0.000000pt}%
\definecolor{currentstroke}{rgb}{0.000000,0.000000,0.000000}%
\pgfsetstrokecolor{currentstroke}%
\pgfsetdash{}{0pt}%
\pgfpathmoveto{\pgfqpoint{0.804877in}{1.166196in}}%
\pgfpathlineto{\pgfqpoint{0.820569in}{1.147526in}}%
\pgfpathlineto{\pgfqpoint{0.836270in}{1.128977in}}%
\pgfpathlineto{\pgfqpoint{0.844419in}{1.137593in}}%
\pgfpathlineto{\pgfqpoint{0.852543in}{1.146543in}}%
\pgfpathlineto{\pgfqpoint{0.836881in}{1.164736in}}%
\pgfpathlineto{\pgfqpoint{0.821228in}{1.183066in}}%
\pgfpathlineto{\pgfqpoint{0.813066in}{1.174473in}}%
\pgfpathlineto{\pgfqpoint{0.804877in}{1.166196in}}%
\pgfpathclose%
\pgfusepath{fill}%
\end{pgfscope}%
\begin{pgfscope}%
\pgfpathrectangle{\pgfqpoint{0.407592in}{0.222904in}}{\pgfqpoint{1.411163in}{1.411163in}}%
\pgfusepath{clip}%
\pgfsetbuttcap%
\pgfsetroundjoin%
\definecolor{currentfill}{rgb}{0.556863,0.654902,0.847059}%
\pgfsetfillcolor{currentfill}%
\pgfsetfillopacity{0.130000}%
\pgfsetlinewidth{0.000000pt}%
\definecolor{currentstroke}{rgb}{0.000000,0.000000,0.000000}%
\pgfsetstrokecolor{currentstroke}%
\pgfsetdash{}{0pt}%
\pgfpathmoveto{\pgfqpoint{0.915290in}{1.075537in}}%
\pgfpathlineto{\pgfqpoint{0.931002in}{1.058348in}}%
\pgfpathlineto{\pgfqpoint{0.946726in}{1.041444in}}%
\pgfpathlineto{\pgfqpoint{0.954716in}{1.052094in}}%
\pgfpathlineto{\pgfqpoint{0.962683in}{1.063091in}}%
\pgfpathlineto{\pgfqpoint{0.946994in}{1.079469in}}%
\pgfpathlineto{\pgfqpoint{0.931316in}{1.096159in}}%
\pgfpathlineto{\pgfqpoint{0.923314in}{1.085680in}}%
\pgfpathlineto{\pgfqpoint{0.915290in}{1.075537in}}%
\pgfpathclose%
\pgfusepath{fill}%
\end{pgfscope}%
\begin{pgfscope}%
\pgfpathrectangle{\pgfqpoint{0.407592in}{0.222904in}}{\pgfqpoint{1.411163in}{1.411163in}}%
\pgfusepath{clip}%
\pgfsetbuttcap%
\pgfsetroundjoin%
\definecolor{currentfill}{rgb}{0.556863,0.654902,0.847059}%
\pgfsetfillcolor{currentfill}%
\pgfsetfillopacity{0.130000}%
\pgfsetlinewidth{0.000000pt}%
\definecolor{currentstroke}{rgb}{0.000000,0.000000,0.000000}%
\pgfsetstrokecolor{currentstroke}%
\pgfsetdash{}{0pt}%
\pgfpathmoveto{\pgfqpoint{1.119966in}{1.005592in}}%
\pgfpathlineto{\pgfqpoint{1.135765in}{0.995548in}}%
\pgfpathlineto{\pgfqpoint{1.151591in}{0.986371in}}%
\pgfpathlineto{\pgfqpoint{1.159350in}{1.001471in}}%
\pgfpathlineto{\pgfqpoint{1.167094in}{1.016715in}}%
\pgfpathlineto{\pgfqpoint{1.151283in}{1.025468in}}%
\pgfpathlineto{\pgfqpoint{1.135502in}{1.035041in}}%
\pgfpathlineto{\pgfqpoint{1.127741in}{1.020233in}}%
\pgfpathlineto{\pgfqpoint{1.119966in}{1.005592in}}%
\pgfpathclose%
\pgfusepath{fill}%
\end{pgfscope}%
\begin{pgfscope}%
\pgfpathrectangle{\pgfqpoint{0.407592in}{0.222904in}}{\pgfqpoint{1.411163in}{1.411163in}}%
\pgfusepath{clip}%
\pgfsetbuttcap%
\pgfsetroundjoin%
\definecolor{currentfill}{rgb}{0.556863,0.654902,0.847059}%
\pgfsetfillcolor{currentfill}%
\pgfsetfillopacity{0.130000}%
\pgfsetlinewidth{0.000000pt}%
\definecolor{currentstroke}{rgb}{0.000000,0.000000,0.000000}%
\pgfsetstrokecolor{currentstroke}%
\pgfsetdash{}{0pt}%
\pgfpathmoveto{\pgfqpoint{1.167094in}{1.016715in}}%
\pgfpathlineto{\pgfqpoint{1.182937in}{1.008828in}}%
\pgfpathlineto{\pgfqpoint{1.198813in}{1.001845in}}%
\pgfpathlineto{\pgfqpoint{1.206532in}{1.017526in}}%
\pgfpathlineto{\pgfqpoint{1.214237in}{1.033302in}}%
\pgfpathlineto{\pgfqpoint{1.198372in}{1.040011in}}%
\pgfpathlineto{\pgfqpoint{1.182543in}{1.047567in}}%
\pgfpathlineto{\pgfqpoint{1.174825in}{1.032086in}}%
\pgfpathlineto{\pgfqpoint{1.167094in}{1.016715in}}%
\pgfpathclose%
\pgfusepath{fill}%
\end{pgfscope}%
\begin{pgfscope}%
\pgfpathrectangle{\pgfqpoint{0.407592in}{0.222904in}}{\pgfqpoint{1.411163in}{1.411163in}}%
\pgfusepath{clip}%
\pgfsetbuttcap%
\pgfsetroundjoin%
\definecolor{currentfill}{rgb}{0.556863,0.654902,0.847059}%
\pgfsetfillcolor{currentfill}%
\pgfsetfillopacity{0.130000}%
\pgfsetlinewidth{0.000000pt}%
\definecolor{currentstroke}{rgb}{0.000000,0.000000,0.000000}%
\pgfsetstrokecolor{currentstroke}%
\pgfsetdash{}{0pt}%
\pgfpathmoveto{\pgfqpoint{1.072808in}{1.000468in}}%
\pgfpathlineto{\pgfqpoint{1.088579in}{0.988297in}}%
\pgfpathlineto{\pgfqpoint{1.104371in}{0.976896in}}%
\pgfpathlineto{\pgfqpoint{1.112176in}{0.991138in}}%
\pgfpathlineto{\pgfqpoint{1.119966in}{1.005592in}}%
\pgfpathlineto{\pgfqpoint{1.104193in}{1.016442in}}%
\pgfpathlineto{\pgfqpoint{1.088444in}{1.028039in}}%
\pgfpathlineto{\pgfqpoint{1.080634in}{1.014137in}}%
\pgfpathlineto{\pgfqpoint{1.072808in}{1.000468in}}%
\pgfpathclose%
\pgfusepath{fill}%
\end{pgfscope}%
\begin{pgfscope}%
\pgfpathrectangle{\pgfqpoint{0.407592in}{0.222904in}}{\pgfqpoint{1.411163in}{1.411163in}}%
\pgfusepath{clip}%
\pgfsetbuttcap%
\pgfsetroundjoin%
\definecolor{currentfill}{rgb}{0.556863,0.654902,0.847059}%
\pgfsetfillcolor{currentfill}%
\pgfsetfillopacity{0.130000}%
\pgfsetlinewidth{0.000000pt}%
\definecolor{currentstroke}{rgb}{0.000000,0.000000,0.000000}%
\pgfsetstrokecolor{currentstroke}%
\pgfsetdash{}{0pt}%
\pgfpathmoveto{\pgfqpoint{1.214237in}{1.033302in}}%
\pgfpathlineto{\pgfqpoint{1.230139in}{1.027466in}}%
\pgfpathlineto{\pgfqpoint{1.246081in}{1.022524in}}%
\pgfpathlineto{\pgfqpoint{1.253762in}{1.038542in}}%
\pgfpathlineto{\pgfqpoint{1.261431in}{1.054624in}}%
\pgfpathlineto{\pgfqpoint{1.245499in}{1.059439in}}%
\pgfpathlineto{\pgfqpoint{1.229608in}{1.065087in}}%
\pgfpathlineto{\pgfqpoint{1.221929in}{1.049159in}}%
\pgfpathlineto{\pgfqpoint{1.214237in}{1.033302in}}%
\pgfpathclose%
\pgfusepath{fill}%
\end{pgfscope}%
\begin{pgfscope}%
\pgfpathrectangle{\pgfqpoint{0.407592in}{0.222904in}}{\pgfqpoint{1.411163in}{1.411163in}}%
\pgfusepath{clip}%
\pgfsetbuttcap%
\pgfsetroundjoin%
\definecolor{currentfill}{rgb}{0.556863,0.654902,0.847059}%
\pgfsetfillcolor{currentfill}%
\pgfsetfillopacity{0.130000}%
\pgfsetlinewidth{0.000000pt}%
\definecolor{currentstroke}{rgb}{0.000000,0.000000,0.000000}%
\pgfsetstrokecolor{currentstroke}%
\pgfsetdash{}{0pt}%
\pgfpathmoveto{\pgfqpoint{0.836270in}{1.128977in}}%
\pgfpathlineto{\pgfqpoint{0.851980in}{1.110563in}}%
\pgfpathlineto{\pgfqpoint{0.867700in}{1.092303in}}%
\pgfpathlineto{\pgfqpoint{0.875811in}{1.101297in}}%
\pgfpathlineto{\pgfqpoint{0.883896in}{1.110644in}}%
\pgfpathlineto{\pgfqpoint{0.868215in}{1.128507in}}%
\pgfpathlineto{\pgfqpoint{0.852543in}{1.146543in}}%
\pgfpathlineto{\pgfqpoint{0.844419in}{1.137593in}}%
\pgfpathlineto{\pgfqpoint{0.836270in}{1.128977in}}%
\pgfpathclose%
\pgfusepath{fill}%
\end{pgfscope}%
\begin{pgfscope}%
\pgfpathrectangle{\pgfqpoint{0.407592in}{0.222904in}}{\pgfqpoint{1.411163in}{1.411163in}}%
\pgfusepath{clip}%
\pgfsetbuttcap%
\pgfsetroundjoin%
\definecolor{currentfill}{rgb}{0.556863,0.654902,0.847059}%
\pgfsetfillcolor{currentfill}%
\pgfsetfillopacity{0.130000}%
\pgfsetlinewidth{0.000000pt}%
\definecolor{currentstroke}{rgb}{0.000000,0.000000,0.000000}%
\pgfsetstrokecolor{currentstroke}%
\pgfsetdash{}{0pt}%
\pgfpathmoveto{\pgfqpoint{0.946726in}{1.041444in}}%
\pgfpathlineto{\pgfqpoint{0.962461in}{1.024867in}}%
\pgfpathlineto{\pgfqpoint{0.978208in}{1.008659in}}%
\pgfpathlineto{\pgfqpoint{0.986164in}{1.019866in}}%
\pgfpathlineto{\pgfqpoint{0.994100in}{1.031429in}}%
\pgfpathlineto{\pgfqpoint{0.978385in}{1.047062in}}%
\pgfpathlineto{\pgfqpoint{0.962683in}{1.063091in}}%
\pgfpathlineto{\pgfqpoint{0.954716in}{1.052094in}}%
\pgfpathlineto{\pgfqpoint{0.946726in}{1.041444in}}%
\pgfpathclose%
\pgfusepath{fill}%
\end{pgfscope}%
\begin{pgfscope}%
\pgfpathrectangle{\pgfqpoint{0.407592in}{0.222904in}}{\pgfqpoint{1.411163in}{1.411163in}}%
\pgfusepath{clip}%
\pgfsetbuttcap%
\pgfsetroundjoin%
\definecolor{currentfill}{rgb}{0.556863,0.654902,0.847059}%
\pgfsetfillcolor{currentfill}%
\pgfsetfillopacity{0.130000}%
\pgfsetlinewidth{0.000000pt}%
\definecolor{currentstroke}{rgb}{0.000000,0.000000,0.000000}%
\pgfsetstrokecolor{currentstroke}%
\pgfsetdash{}{0pt}%
\pgfpathmoveto{\pgfqpoint{1.025571in}{1.001551in}}%
\pgfpathlineto{\pgfqpoint{1.041329in}{0.987427in}}%
\pgfpathlineto{\pgfqpoint{1.057105in}{0.973936in}}%
\pgfpathlineto{\pgfqpoint{1.064965in}{0.987058in}}%
\pgfpathlineto{\pgfqpoint{1.072808in}{1.000468in}}%
\pgfpathlineto{\pgfqpoint{1.057057in}{1.013340in}}%
\pgfpathlineto{\pgfqpoint{1.041325in}{1.026847in}}%
\pgfpathlineto{\pgfqpoint{1.033457in}{1.014049in}}%
\pgfpathlineto{\pgfqpoint{1.025571in}{1.001551in}}%
\pgfpathclose%
\pgfusepath{fill}%
\end{pgfscope}%
\begin{pgfscope}%
\pgfpathrectangle{\pgfqpoint{0.407592in}{0.222904in}}{\pgfqpoint{1.411163in}{1.411163in}}%
\pgfusepath{clip}%
\pgfsetbuttcap%
\pgfsetroundjoin%
\definecolor{currentfill}{rgb}{0.556863,0.654902,0.847059}%
\pgfsetfillcolor{currentfill}%
\pgfsetfillopacity{0.130000}%
\pgfsetlinewidth{0.000000pt}%
\definecolor{currentstroke}{rgb}{0.000000,0.000000,0.000000}%
\pgfsetstrokecolor{currentstroke}%
\pgfsetdash{}{0pt}%
\pgfpathmoveto{\pgfqpoint{1.261431in}{1.054624in}}%
\pgfpathlineto{\pgfqpoint{1.277406in}{1.050651in}}%
\pgfpathlineto{\pgfqpoint{1.293426in}{1.047526in}}%
\pgfpathlineto{\pgfqpoint{1.301073in}{1.063694in}}%
\pgfpathlineto{\pgfqpoint{1.308706in}{1.079908in}}%
\pgfpathlineto{\pgfqpoint{1.292696in}{1.083029in}}%
\pgfpathlineto{\pgfqpoint{1.276730in}{1.086941in}}%
\pgfpathlineto{\pgfqpoint{1.269087in}{1.070760in}}%
\pgfpathlineto{\pgfqpoint{1.261431in}{1.054624in}}%
\pgfpathclose%
\pgfusepath{fill}%
\end{pgfscope}%
\begin{pgfscope}%
\pgfpathrectangle{\pgfqpoint{0.407592in}{0.222904in}}{\pgfqpoint{1.411163in}{1.411163in}}%
\pgfusepath{clip}%
\pgfsetbuttcap%
\pgfsetroundjoin%
\definecolor{currentfill}{rgb}{0.556863,0.654902,0.847059}%
\pgfsetfillcolor{currentfill}%
\pgfsetfillopacity{0.130000}%
\pgfsetlinewidth{0.000000pt}%
\definecolor{currentstroke}{rgb}{0.000000,0.000000,0.000000}%
\pgfsetstrokecolor{currentstroke}%
\pgfsetdash{}{0pt}%
\pgfpathmoveto{\pgfqpoint{0.725579in}{1.227318in}}%
\pgfpathlineto{\pgfqpoint{0.741275in}{1.208042in}}%
\pgfpathlineto{\pgfqpoint{0.756980in}{1.188829in}}%
\pgfpathlineto{\pgfqpoint{0.765265in}{1.196178in}}%
\pgfpathlineto{\pgfqpoint{0.773522in}{1.203843in}}%
\pgfpathlineto{\pgfqpoint{0.757857in}{1.222798in}}%
\pgfpathlineto{\pgfqpoint{0.742202in}{1.241828in}}%
\pgfpathlineto{\pgfqpoint{0.733905in}{1.234425in}}%
\pgfpathlineto{\pgfqpoint{0.725579in}{1.227318in}}%
\pgfpathclose%
\pgfusepath{fill}%
\end{pgfscope}%
\begin{pgfscope}%
\pgfpathrectangle{\pgfqpoint{0.407592in}{0.222904in}}{\pgfqpoint{1.411163in}{1.411163in}}%
\pgfusepath{clip}%
\pgfsetbuttcap%
\pgfsetroundjoin%
\definecolor{currentfill}{rgb}{0.556863,0.654902,0.847059}%
\pgfsetfillcolor{currentfill}%
\pgfsetfillopacity{0.130000}%
\pgfsetlinewidth{0.000000pt}%
\definecolor{currentstroke}{rgb}{0.000000,0.000000,0.000000}%
\pgfsetstrokecolor{currentstroke}%
\pgfsetdash{}{0pt}%
\pgfpathmoveto{\pgfqpoint{0.867700in}{1.092303in}}%
\pgfpathlineto{\pgfqpoint{0.883429in}{1.074219in}}%
\pgfpathlineto{\pgfqpoint{0.899168in}{1.056333in}}%
\pgfpathlineto{\pgfqpoint{0.907241in}{1.065749in}}%
\pgfpathlineto{\pgfqpoint{0.915290in}{1.075537in}}%
\pgfpathlineto{\pgfqpoint{0.899588in}{1.092980in}}%
\pgfpathlineto{\pgfqpoint{0.883896in}{1.110644in}}%
\pgfpathlineto{\pgfqpoint{0.875811in}{1.101297in}}%
\pgfpathlineto{\pgfqpoint{0.867700in}{1.092303in}}%
\pgfpathclose%
\pgfusepath{fill}%
\end{pgfscope}%
\begin{pgfscope}%
\pgfpathrectangle{\pgfqpoint{0.407592in}{0.222904in}}{\pgfqpoint{1.411163in}{1.411163in}}%
\pgfusepath{clip}%
\pgfsetbuttcap%
\pgfsetroundjoin%
\definecolor{currentfill}{rgb}{0.556863,0.654902,0.847059}%
\pgfsetfillcolor{currentfill}%
\pgfsetfillopacity{0.130000}%
\pgfsetlinewidth{0.000000pt}%
\definecolor{currentstroke}{rgb}{0.000000,0.000000,0.000000}%
\pgfsetstrokecolor{currentstroke}%
\pgfsetdash{}{0pt}%
\pgfpathmoveto{\pgfqpoint{0.756980in}{1.188829in}}%
\pgfpathlineto{\pgfqpoint{0.772693in}{1.169687in}}%
\pgfpathlineto{\pgfqpoint{0.788416in}{1.150626in}}%
\pgfpathlineto{\pgfqpoint{0.796661in}{1.158244in}}%
\pgfpathlineto{\pgfqpoint{0.804877in}{1.166196in}}%
\pgfpathlineto{\pgfqpoint{0.789195in}{1.184972in}}%
\pgfpathlineto{\pgfqpoint{0.773522in}{1.203843in}}%
\pgfpathlineto{\pgfqpoint{0.765265in}{1.196178in}}%
\pgfpathlineto{\pgfqpoint{0.756980in}{1.188829in}}%
\pgfpathclose%
\pgfusepath{fill}%
\end{pgfscope}%
\begin{pgfscope}%
\pgfpathrectangle{\pgfqpoint{0.407592in}{0.222904in}}{\pgfqpoint{1.411163in}{1.411163in}}%
\pgfusepath{clip}%
\pgfsetbuttcap%
\pgfsetroundjoin%
\definecolor{currentfill}{rgb}{0.556863,0.654902,0.847059}%
\pgfsetfillcolor{currentfill}%
\pgfsetfillopacity{0.130000}%
\pgfsetlinewidth{0.000000pt}%
\definecolor{currentstroke}{rgb}{0.000000,0.000000,0.000000}%
\pgfsetstrokecolor{currentstroke}%
\pgfsetdash{}{0pt}%
\pgfpathmoveto{\pgfqpoint{0.978208in}{1.008659in}}%
\pgfpathlineto{\pgfqpoint{0.993967in}{0.992874in}}%
\pgfpathlineto{\pgfqpoint{1.009740in}{0.977571in}}%
\pgfpathlineto{\pgfqpoint{1.017665in}{0.989382in}}%
\pgfpathlineto{\pgfqpoint{1.025571in}{1.001551in}}%
\pgfpathlineto{\pgfqpoint{1.009828in}{1.016239in}}%
\pgfpathlineto{\pgfqpoint{0.994100in}{1.031429in}}%
\pgfpathlineto{\pgfqpoint{0.986164in}{1.019866in}}%
\pgfpathlineto{\pgfqpoint{0.978208in}{1.008659in}}%
\pgfpathclose%
\pgfusepath{fill}%
\end{pgfscope}%
\begin{pgfscope}%
\pgfpathrectangle{\pgfqpoint{0.407592in}{0.222904in}}{\pgfqpoint{1.411163in}{1.411163in}}%
\pgfusepath{clip}%
\pgfsetbuttcap%
\pgfsetroundjoin%
\definecolor{currentfill}{rgb}{0.556863,0.654902,0.847059}%
\pgfsetfillcolor{currentfill}%
\pgfsetfillopacity{0.130000}%
\pgfsetlinewidth{0.000000pt}%
\definecolor{currentstroke}{rgb}{0.000000,0.000000,0.000000}%
\pgfsetstrokecolor{currentstroke}%
\pgfsetdash{}{0pt}%
\pgfpathmoveto{\pgfqpoint{1.151591in}{0.986371in}}%
\pgfpathlineto{\pgfqpoint{1.167447in}{0.978114in}}%
\pgfpathlineto{\pgfqpoint{1.183336in}{0.970828in}}%
\pgfpathlineto{\pgfqpoint{1.191081in}{0.986273in}}%
\pgfpathlineto{\pgfqpoint{1.198813in}{1.001845in}}%
\pgfpathlineto{\pgfqpoint{1.182937in}{1.008828in}}%
\pgfpathlineto{\pgfqpoint{1.167094in}{1.016715in}}%
\pgfpathlineto{\pgfqpoint{1.159350in}{1.001471in}}%
\pgfpathlineto{\pgfqpoint{1.151591in}{0.986371in}}%
\pgfpathclose%
\pgfusepath{fill}%
\end{pgfscope}%
\begin{pgfscope}%
\pgfpathrectangle{\pgfqpoint{0.407592in}{0.222904in}}{\pgfqpoint{1.411163in}{1.411163in}}%
\pgfusepath{clip}%
\pgfsetbuttcap%
\pgfsetroundjoin%
\definecolor{currentfill}{rgb}{0.556863,0.654902,0.847059}%
\pgfsetfillcolor{currentfill}%
\pgfsetfillopacity{0.130000}%
\pgfsetlinewidth{0.000000pt}%
\definecolor{currentstroke}{rgb}{0.000000,0.000000,0.000000}%
\pgfsetstrokecolor{currentstroke}%
\pgfsetdash{}{0pt}%
\pgfpathmoveto{\pgfqpoint{1.104371in}{0.976896in}}%
\pgfpathlineto{\pgfqpoint{1.120187in}{0.966336in}}%
\pgfpathlineto{\pgfqpoint{1.136029in}{0.956689in}}%
\pgfpathlineto{\pgfqpoint{1.143817in}{0.971436in}}%
\pgfpathlineto{\pgfqpoint{1.151591in}{0.986371in}}%
\pgfpathlineto{\pgfqpoint{1.135765in}{0.995548in}}%
\pgfpathlineto{\pgfqpoint{1.119966in}{1.005592in}}%
\pgfpathlineto{\pgfqpoint{1.112176in}{0.991138in}}%
\pgfpathlineto{\pgfqpoint{1.104371in}{0.976896in}}%
\pgfpathclose%
\pgfusepath{fill}%
\end{pgfscope}%
\begin{pgfscope}%
\pgfpathrectangle{\pgfqpoint{0.407592in}{0.222904in}}{\pgfqpoint{1.411163in}{1.411163in}}%
\pgfusepath{clip}%
\pgfsetbuttcap%
\pgfsetroundjoin%
\definecolor{currentfill}{rgb}{0.556863,0.654902,0.847059}%
\pgfsetfillcolor{currentfill}%
\pgfsetfillopacity{0.130000}%
\pgfsetlinewidth{0.000000pt}%
\definecolor{currentstroke}{rgb}{0.000000,0.000000,0.000000}%
\pgfsetstrokecolor{currentstroke}%
\pgfsetdash{}{0pt}%
\pgfpathmoveto{\pgfqpoint{1.308706in}{1.079908in}}%
\pgfpathlineto{\pgfqpoint{1.324765in}{1.077576in}}%
\pgfpathlineto{\pgfqpoint{1.340873in}{1.076026in}}%
\pgfpathlineto{\pgfqpoint{1.348485in}{1.092211in}}%
\pgfpathlineto{\pgfqpoint{1.332382in}{1.093808in}}%
\pgfpathlineto{\pgfqpoint{1.316328in}{1.096162in}}%
\pgfusepath{fill}%
\end{pgfscope}%
\begin{pgfscope}%
\pgfpathrectangle{\pgfqpoint{0.407592in}{0.222904in}}{\pgfqpoint{1.411163in}{1.411163in}}%
\pgfusepath{clip}%
\pgfsetbuttcap%
\pgfsetroundjoin%
\definecolor{currentfill}{rgb}{0.556863,0.654902,0.847059}%
\pgfsetfillcolor{currentfill}%
\pgfsetfillopacity{0.130000}%
\pgfsetlinewidth{0.000000pt}%
\definecolor{currentstroke}{rgb}{0.000000,0.000000,0.000000}%
\pgfsetstrokecolor{currentstroke}%
\pgfsetdash{}{0pt}%
\pgfpathmoveto{\pgfqpoint{1.198813in}{1.001845in}}%
\pgfpathlineto{\pgfqpoint{1.214726in}{0.995801in}}%
\pgfpathlineto{\pgfqpoint{1.230679in}{0.990721in}}%
\pgfpathlineto{\pgfqpoint{1.238386in}{1.006579in}}%
\pgfpathlineto{\pgfqpoint{1.246081in}{1.022524in}}%
\pgfpathlineto{\pgfqpoint{1.230139in}{1.027466in}}%
\pgfpathlineto{\pgfqpoint{1.214237in}{1.033302in}}%
\pgfpathlineto{\pgfqpoint{1.206532in}{1.017526in}}%
\pgfpathlineto{\pgfqpoint{1.198813in}{1.001845in}}%
\pgfpathclose%
\pgfusepath{fill}%
\end{pgfscope}%
\begin{pgfscope}%
\pgfpathrectangle{\pgfqpoint{0.407592in}{0.222904in}}{\pgfqpoint{1.411163in}{1.411163in}}%
\pgfusepath{clip}%
\pgfsetbuttcap%
\pgfsetroundjoin%
\definecolor{currentfill}{rgb}{0.556863,0.654902,0.847059}%
\pgfsetfillcolor{currentfill}%
\pgfsetfillopacity{0.130000}%
\pgfsetlinewidth{0.000000pt}%
\definecolor{currentstroke}{rgb}{0.000000,0.000000,0.000000}%
\pgfsetstrokecolor{currentstroke}%
\pgfsetdash{}{0pt}%
\pgfpathmoveto{\pgfqpoint{0.899168in}{1.056333in}}%
\pgfpathlineto{\pgfqpoint{0.914917in}{1.038674in}}%
\pgfpathlineto{\pgfqpoint{0.930677in}{1.021276in}}%
\pgfpathlineto{\pgfqpoint{0.938713in}{1.031164in}}%
\pgfpathlineto{\pgfqpoint{0.946726in}{1.041444in}}%
\pgfpathlineto{\pgfqpoint{0.931002in}{1.058348in}}%
\pgfpathlineto{\pgfqpoint{0.915290in}{1.075537in}}%
\pgfpathlineto{\pgfqpoint{0.907241in}{1.065749in}}%
\pgfpathlineto{\pgfqpoint{0.899168in}{1.056333in}}%
\pgfpathclose%
\pgfusepath{fill}%
\end{pgfscope}%
\begin{pgfscope}%
\pgfpathrectangle{\pgfqpoint{0.407592in}{0.222904in}}{\pgfqpoint{1.411163in}{1.411163in}}%
\pgfusepath{clip}%
\pgfsetbuttcap%
\pgfsetroundjoin%
\definecolor{currentfill}{rgb}{0.556863,0.654902,0.847059}%
\pgfsetfillcolor{currentfill}%
\pgfsetfillopacity{0.130000}%
\pgfsetlinewidth{0.000000pt}%
\definecolor{currentstroke}{rgb}{0.000000,0.000000,0.000000}%
\pgfsetstrokecolor{currentstroke}%
\pgfsetdash{}{0pt}%
\pgfpathmoveto{\pgfqpoint{1.057105in}{0.973936in}}%
\pgfpathlineto{\pgfqpoint{1.072898in}{0.961153in}}%
\pgfpathlineto{\pgfqpoint{1.088712in}{0.949159in}}%
\pgfpathlineto{\pgfqpoint{1.096550in}{0.962892in}}%
\pgfpathlineto{\pgfqpoint{1.104371in}{0.976896in}}%
\pgfpathlineto{\pgfqpoint{1.088579in}{0.988297in}}%
\pgfpathlineto{\pgfqpoint{1.072808in}{1.000468in}}%
\pgfpathlineto{\pgfqpoint{1.064965in}{0.987058in}}%
\pgfpathlineto{\pgfqpoint{1.057105in}{0.973936in}}%
\pgfpathclose%
\pgfusepath{fill}%
\end{pgfscope}%
\begin{pgfscope}%
\pgfpathrectangle{\pgfqpoint{0.407592in}{0.222904in}}{\pgfqpoint{1.411163in}{1.411163in}}%
\pgfusepath{clip}%
\pgfsetbuttcap%
\pgfsetroundjoin%
\definecolor{currentfill}{rgb}{0.556863,0.654902,0.847059}%
\pgfsetfillcolor{currentfill}%
\pgfsetfillopacity{0.130000}%
\pgfsetlinewidth{0.000000pt}%
\definecolor{currentstroke}{rgb}{0.000000,0.000000,0.000000}%
\pgfsetstrokecolor{currentstroke}%
\pgfsetdash{}{0pt}%
\pgfpathmoveto{\pgfqpoint{0.788416in}{1.150626in}}%
\pgfpathlineto{\pgfqpoint{0.804148in}{1.131657in}}%
\pgfpathlineto{\pgfqpoint{0.819889in}{1.112792in}}%
\pgfpathlineto{\pgfqpoint{0.828093in}{1.120706in}}%
\pgfpathlineto{\pgfqpoint{0.836270in}{1.128977in}}%
\pgfpathlineto{\pgfqpoint{0.820569in}{1.147526in}}%
\pgfpathlineto{\pgfqpoint{0.804877in}{1.166196in}}%
\pgfpathlineto{\pgfqpoint{0.796661in}{1.158244in}}%
\pgfpathlineto{\pgfqpoint{0.788416in}{1.150626in}}%
\pgfpathclose%
\pgfusepath{fill}%
\end{pgfscope}%
\begin{pgfscope}%
\pgfpathrectangle{\pgfqpoint{0.407592in}{0.222904in}}{\pgfqpoint{1.411163in}{1.411163in}}%
\pgfusepath{clip}%
\pgfsetbuttcap%
\pgfsetroundjoin%
\definecolor{currentfill}{rgb}{0.556863,0.654902,0.847059}%
\pgfsetfillcolor{currentfill}%
\pgfsetfillopacity{0.130000}%
\pgfsetlinewidth{0.000000pt}%
\definecolor{currentstroke}{rgb}{0.000000,0.000000,0.000000}%
\pgfsetstrokecolor{currentstroke}%
\pgfsetdash{}{0pt}%
\pgfpathmoveto{\pgfqpoint{1.246081in}{1.022524in}}%
\pgfpathlineto{\pgfqpoint{1.262065in}{1.018489in}}%
\pgfpathlineto{\pgfqpoint{1.278095in}{1.015365in}}%
\pgfpathlineto{\pgfqpoint{1.285767in}{1.031413in}}%
\pgfpathlineto{\pgfqpoint{1.293426in}{1.047526in}}%
\pgfpathlineto{\pgfqpoint{1.277406in}{1.050651in}}%
\pgfpathlineto{\pgfqpoint{1.261431in}{1.054624in}}%
\pgfpathlineto{\pgfqpoint{1.253762in}{1.038542in}}%
\pgfpathlineto{\pgfqpoint{1.246081in}{1.022524in}}%
\pgfpathclose%
\pgfusepath{fill}%
\end{pgfscope}%
\begin{pgfscope}%
\pgfpathrectangle{\pgfqpoint{0.407592in}{0.222904in}}{\pgfqpoint{1.411163in}{1.411163in}}%
\pgfusepath{clip}%
\pgfsetbuttcap%
\pgfsetroundjoin%
\definecolor{currentfill}{rgb}{0.556863,0.654902,0.847059}%
\pgfsetfillcolor{currentfill}%
\pgfsetfillopacity{0.130000}%
\pgfsetlinewidth{0.000000pt}%
\definecolor{currentstroke}{rgb}{0.000000,0.000000,0.000000}%
\pgfsetstrokecolor{currentstroke}%
\pgfsetdash{}{0pt}%
\pgfpathmoveto{\pgfqpoint{1.009740in}{0.977571in}}%
\pgfpathlineto{\pgfqpoint{1.025527in}{0.962816in}}%
\pgfpathlineto{\pgfqpoint{1.041329in}{0.948685in}}%
\pgfpathlineto{\pgfqpoint{1.049226in}{0.961133in}}%
\pgfpathlineto{\pgfqpoint{1.057105in}{0.973936in}}%
\pgfpathlineto{\pgfqpoint{1.041329in}{0.987427in}}%
\pgfpathlineto{\pgfqpoint{1.025571in}{1.001551in}}%
\pgfpathlineto{\pgfqpoint{1.017665in}{0.989382in}}%
\pgfpathlineto{\pgfqpoint{1.009740in}{0.977571in}}%
\pgfpathclose%
\pgfusepath{fill}%
\end{pgfscope}%
\begin{pgfscope}%
\pgfpathrectangle{\pgfqpoint{0.407592in}{0.222904in}}{\pgfqpoint{1.411163in}{1.411163in}}%
\pgfusepath{clip}%
\pgfsetbuttcap%
\pgfsetroundjoin%
\definecolor{currentfill}{rgb}{0.556863,0.654902,0.847059}%
\pgfsetfillcolor{currentfill}%
\pgfsetfillopacity{0.130000}%
\pgfsetlinewidth{0.000000pt}%
\definecolor{currentstroke}{rgb}{0.000000,0.000000,0.000000}%
\pgfsetstrokecolor{currentstroke}%
\pgfsetdash{}{0pt}%
\pgfpathmoveto{\pgfqpoint{0.819889in}{1.112792in}}%
\pgfpathlineto{\pgfqpoint{0.835638in}{1.094044in}}%
\pgfpathlineto{\pgfqpoint{0.851397in}{1.075431in}}%
\pgfpathlineto{\pgfqpoint{0.859562in}{1.083677in}}%
\pgfpathlineto{\pgfqpoint{0.867700in}{1.092303in}}%
\pgfpathlineto{\pgfqpoint{0.851980in}{1.110563in}}%
\pgfpathlineto{\pgfqpoint{0.836270in}{1.128977in}}%
\pgfpathlineto{\pgfqpoint{0.828093in}{1.120706in}}%
\pgfpathlineto{\pgfqpoint{0.819889in}{1.112792in}}%
\pgfpathclose%
\pgfusepath{fill}%
\end{pgfscope}%
\begin{pgfscope}%
\pgfpathrectangle{\pgfqpoint{0.407592in}{0.222904in}}{\pgfqpoint{1.411163in}{1.411163in}}%
\pgfusepath{clip}%
\pgfsetbuttcap%
\pgfsetroundjoin%
\definecolor{currentfill}{rgb}{0.556863,0.654902,0.847059}%
\pgfsetfillcolor{currentfill}%
\pgfsetfillopacity{0.130000}%
\pgfsetlinewidth{0.000000pt}%
\definecolor{currentstroke}{rgb}{0.000000,0.000000,0.000000}%
\pgfsetstrokecolor{currentstroke}%
\pgfsetdash{}{0pt}%
\pgfpathmoveto{\pgfqpoint{0.930677in}{1.021276in}}%
\pgfpathlineto{\pgfqpoint{0.946447in}{1.004176in}}%
\pgfpathlineto{\pgfqpoint{0.962228in}{0.987419in}}%
\pgfpathlineto{\pgfqpoint{0.970229in}{0.997834in}}%
\pgfpathlineto{\pgfqpoint{0.978208in}{1.008659in}}%
\pgfpathlineto{\pgfqpoint{0.962461in}{1.024867in}}%
\pgfpathlineto{\pgfqpoint{0.946726in}{1.041444in}}%
\pgfpathlineto{\pgfqpoint{0.938713in}{1.031164in}}%
\pgfpathlineto{\pgfqpoint{0.930677in}{1.021276in}}%
\pgfpathclose%
\pgfusepath{fill}%
\end{pgfscope}%
\begin{pgfscope}%
\pgfpathrectangle{\pgfqpoint{0.407592in}{0.222904in}}{\pgfqpoint{1.411163in}{1.411163in}}%
\pgfusepath{clip}%
\pgfsetbuttcap%
\pgfsetroundjoin%
\definecolor{currentfill}{rgb}{0.556863,0.654902,0.847059}%
\pgfsetfillcolor{currentfill}%
\pgfsetfillopacity{0.130000}%
\pgfsetlinewidth{0.000000pt}%
\definecolor{currentstroke}{rgb}{0.000000,0.000000,0.000000}%
\pgfsetstrokecolor{currentstroke}%
\pgfsetdash{}{0pt}%
\pgfpathmoveto{\pgfqpoint{1.293426in}{1.047526in}}%
\pgfpathlineto{\pgfqpoint{1.309494in}{1.045246in}}%
\pgfpathlineto{\pgfqpoint{1.325612in}{1.043802in}}%
\pgfpathlineto{\pgfqpoint{1.333249in}{1.059887in}}%
\pgfpathlineto{\pgfqpoint{1.340873in}{1.076026in}}%
\pgfpathlineto{\pgfqpoint{1.324765in}{1.077576in}}%
\pgfpathlineto{\pgfqpoint{1.308706in}{1.079908in}}%
\pgfpathlineto{\pgfqpoint{1.301073in}{1.063694in}}%
\pgfpathlineto{\pgfqpoint{1.293426in}{1.047526in}}%
\pgfpathclose%
\pgfusepath{fill}%
\end{pgfscope}%
\begin{pgfscope}%
\pgfpathrectangle{\pgfqpoint{0.407592in}{0.222904in}}{\pgfqpoint{1.411163in}{1.411163in}}%
\pgfusepath{clip}%
\pgfsetbuttcap%
\pgfsetroundjoin%
\definecolor{currentfill}{rgb}{0.556863,0.654902,0.847059}%
\pgfsetfillcolor{currentfill}%
\pgfsetfillopacity{0.130000}%
\pgfsetlinewidth{0.000000pt}%
\definecolor{currentstroke}{rgb}{0.000000,0.000000,0.000000}%
\pgfsetstrokecolor{currentstroke}%
\pgfsetdash{}{0pt}%
\pgfpathmoveto{\pgfqpoint{0.708836in}{1.214017in}}%
\pgfpathlineto{\pgfqpoint{0.724573in}{1.194532in}}%
\pgfpathlineto{\pgfqpoint{0.740319in}{1.175101in}}%
\pgfpathlineto{\pgfqpoint{0.748665in}{1.181801in}}%
\pgfpathlineto{\pgfqpoint{0.756980in}{1.188829in}}%
\pgfpathlineto{\pgfqpoint{0.741275in}{1.208042in}}%
\pgfpathlineto{\pgfqpoint{0.725579in}{1.227318in}}%
\pgfpathlineto{\pgfqpoint{0.717222in}{1.220514in}}%
\pgfpathlineto{\pgfqpoint{0.708836in}{1.214017in}}%
\pgfpathclose%
\pgfusepath{fill}%
\end{pgfscope}%
\begin{pgfscope}%
\pgfpathrectangle{\pgfqpoint{0.407592in}{0.222904in}}{\pgfqpoint{1.411163in}{1.411163in}}%
\pgfusepath{clip}%
\pgfsetbuttcap%
\pgfsetroundjoin%
\definecolor{currentfill}{rgb}{0.556863,0.654902,0.847059}%
\pgfsetfillcolor{currentfill}%
\pgfsetfillopacity{0.130000}%
\pgfsetlinewidth{0.000000pt}%
\definecolor{currentstroke}{rgb}{0.000000,0.000000,0.000000}%
\pgfsetstrokecolor{currentstroke}%
\pgfsetdash{}{0pt}%
\pgfpathmoveto{\pgfqpoint{1.136029in}{0.956689in}}%
\pgfpathlineto{\pgfqpoint{1.151900in}{0.948023in}}%
\pgfpathlineto{\pgfqpoint{1.167803in}{0.940398in}}%
\pgfpathlineto{\pgfqpoint{1.175576in}{0.955528in}}%
\pgfpathlineto{\pgfqpoint{1.183336in}{0.970828in}}%
\pgfpathlineto{\pgfqpoint{1.167447in}{0.978114in}}%
\pgfpathlineto{\pgfqpoint{1.151591in}{0.986371in}}%
\pgfpathlineto{\pgfqpoint{1.143817in}{0.971436in}}%
\pgfpathlineto{\pgfqpoint{1.136029in}{0.956689in}}%
\pgfpathclose%
\pgfusepath{fill}%
\end{pgfscope}%
\begin{pgfscope}%
\pgfpathrectangle{\pgfqpoint{0.407592in}{0.222904in}}{\pgfqpoint{1.411163in}{1.411163in}}%
\pgfusepath{clip}%
\pgfsetbuttcap%
\pgfsetroundjoin%
\definecolor{currentfill}{rgb}{0.556863,0.654902,0.847059}%
\pgfsetfillcolor{currentfill}%
\pgfsetfillopacity{0.130000}%
\pgfsetlinewidth{0.000000pt}%
\definecolor{currentstroke}{rgb}{0.000000,0.000000,0.000000}%
\pgfsetstrokecolor{currentstroke}%
\pgfsetdash{}{0pt}%
\pgfpathmoveto{\pgfqpoint{1.088712in}{0.949159in}}%
\pgfpathlineto{\pgfqpoint{1.104548in}{0.938038in}}%
\pgfpathlineto{\pgfqpoint{1.120408in}{0.927876in}}%
\pgfpathlineto{\pgfqpoint{1.128227in}{0.942159in}}%
\pgfpathlineto{\pgfqpoint{1.136029in}{0.956689in}}%
\pgfpathlineto{\pgfqpoint{1.120187in}{0.966336in}}%
\pgfpathlineto{\pgfqpoint{1.104371in}{0.976896in}}%
\pgfpathlineto{\pgfqpoint{1.096550in}{0.962892in}}%
\pgfpathlineto{\pgfqpoint{1.088712in}{0.949159in}}%
\pgfpathclose%
\pgfusepath{fill}%
\end{pgfscope}%
\begin{pgfscope}%
\pgfpathrectangle{\pgfqpoint{0.407592in}{0.222904in}}{\pgfqpoint{1.411163in}{1.411163in}}%
\pgfusepath{clip}%
\pgfsetbuttcap%
\pgfsetroundjoin%
\definecolor{currentfill}{rgb}{0.556863,0.654902,0.847059}%
\pgfsetfillcolor{currentfill}%
\pgfsetfillopacity{0.130000}%
\pgfsetlinewidth{0.000000pt}%
\definecolor{currentstroke}{rgb}{0.000000,0.000000,0.000000}%
\pgfsetstrokecolor{currentstroke}%
\pgfsetdash{}{0pt}%
\pgfpathmoveto{\pgfqpoint{0.851397in}{1.075431in}}%
\pgfpathlineto{\pgfqpoint{0.867166in}{1.056972in}}%
\pgfpathlineto{\pgfqpoint{0.882944in}{1.038690in}}%
\pgfpathlineto{\pgfqpoint{0.891069in}{1.047307in}}%
\pgfpathlineto{\pgfqpoint{0.899168in}{1.056333in}}%
\pgfpathlineto{\pgfqpoint{0.883429in}{1.074219in}}%
\pgfpathlineto{\pgfqpoint{0.867700in}{1.092303in}}%
\pgfpathlineto{\pgfqpoint{0.859562in}{1.083677in}}%
\pgfpathlineto{\pgfqpoint{0.851397in}{1.075431in}}%
\pgfpathclose%
\pgfusepath{fill}%
\end{pgfscope}%
\begin{pgfscope}%
\pgfpathrectangle{\pgfqpoint{0.407592in}{0.222904in}}{\pgfqpoint{1.411163in}{1.411163in}}%
\pgfusepath{clip}%
\pgfsetbuttcap%
\pgfsetroundjoin%
\definecolor{currentfill}{rgb}{0.556863,0.654902,0.847059}%
\pgfsetfillcolor{currentfill}%
\pgfsetfillopacity{0.130000}%
\pgfsetlinewidth{0.000000pt}%
\definecolor{currentstroke}{rgb}{0.000000,0.000000,0.000000}%
\pgfsetstrokecolor{currentstroke}%
\pgfsetdash{}{0pt}%
\pgfpathmoveto{\pgfqpoint{1.183336in}{0.970828in}}%
\pgfpathlineto{\pgfqpoint{1.199261in}{0.964553in}}%
\pgfpathlineto{\pgfqpoint{1.215224in}{0.959322in}}%
\pgfpathlineto{\pgfqpoint{1.222958in}{0.974963in}}%
\pgfpathlineto{\pgfqpoint{1.230679in}{0.990721in}}%
\pgfpathlineto{\pgfqpoint{1.214726in}{0.995801in}}%
\pgfpathlineto{\pgfqpoint{1.198813in}{1.001845in}}%
\pgfpathlineto{\pgfqpoint{1.191081in}{0.986273in}}%
\pgfpathlineto{\pgfqpoint{1.183336in}{0.970828in}}%
\pgfpathclose%
\pgfusepath{fill}%
\end{pgfscope}%
\begin{pgfscope}%
\pgfpathrectangle{\pgfqpoint{0.407592in}{0.222904in}}{\pgfqpoint{1.411163in}{1.411163in}}%
\pgfusepath{clip}%
\pgfsetbuttcap%
\pgfsetroundjoin%
\definecolor{currentfill}{rgb}{0.556863,0.654902,0.847059}%
\pgfsetfillcolor{currentfill}%
\pgfsetfillopacity{0.130000}%
\pgfsetlinewidth{0.000000pt}%
\definecolor{currentstroke}{rgb}{0.000000,0.000000,0.000000}%
\pgfsetstrokecolor{currentstroke}%
\pgfsetdash{}{0pt}%
\pgfpathmoveto{\pgfqpoint{0.962228in}{0.987419in}}%
\pgfpathlineto{\pgfqpoint{0.978021in}{0.971057in}}%
\pgfpathlineto{\pgfqpoint{0.993826in}{0.955152in}}%
\pgfpathlineto{\pgfqpoint{1.001794in}{0.966150in}}%
\pgfpathlineto{\pgfqpoint{1.009740in}{0.977571in}}%
\pgfpathlineto{\pgfqpoint{0.993967in}{0.992874in}}%
\pgfpathlineto{\pgfqpoint{0.978208in}{1.008659in}}%
\pgfpathlineto{\pgfqpoint{0.970229in}{0.997834in}}%
\pgfpathlineto{\pgfqpoint{0.962228in}{0.987419in}}%
\pgfpathclose%
\pgfusepath{fill}%
\end{pgfscope}%
\begin{pgfscope}%
\pgfpathrectangle{\pgfqpoint{0.407592in}{0.222904in}}{\pgfqpoint{1.411163in}{1.411163in}}%
\pgfusepath{clip}%
\pgfsetbuttcap%
\pgfsetroundjoin%
\definecolor{currentfill}{rgb}{0.556863,0.654902,0.847059}%
\pgfsetfillcolor{currentfill}%
\pgfsetfillopacity{0.130000}%
\pgfsetlinewidth{0.000000pt}%
\definecolor{currentstroke}{rgb}{0.000000,0.000000,0.000000}%
\pgfsetstrokecolor{currentstroke}%
\pgfsetdash{}{0pt}%
\pgfpathmoveto{\pgfqpoint{1.041329in}{0.948685in}}%
\pgfpathlineto{\pgfqpoint{1.057148in}{0.935262in}}%
\pgfpathlineto{\pgfqpoint{1.072986in}{0.922640in}}%
\pgfpathlineto{\pgfqpoint{1.080858in}{0.935729in}}%
\pgfpathlineto{\pgfqpoint{1.088712in}{0.949159in}}%
\pgfpathlineto{\pgfqpoint{1.072898in}{0.961153in}}%
\pgfpathlineto{\pgfqpoint{1.057105in}{0.973936in}}%
\pgfpathlineto{\pgfqpoint{1.049226in}{0.961133in}}%
\pgfpathlineto{\pgfqpoint{1.041329in}{0.948685in}}%
\pgfpathclose%
\pgfusepath{fill}%
\end{pgfscope}%
\begin{pgfscope}%
\pgfpathrectangle{\pgfqpoint{0.407592in}{0.222904in}}{\pgfqpoint{1.411163in}{1.411163in}}%
\pgfusepath{clip}%
\pgfsetbuttcap%
\pgfsetroundjoin%
\definecolor{currentfill}{rgb}{0.556863,0.654902,0.847059}%
\pgfsetfillcolor{currentfill}%
\pgfsetfillopacity{0.130000}%
\pgfsetlinewidth{0.000000pt}%
\definecolor{currentstroke}{rgb}{0.000000,0.000000,0.000000}%
\pgfsetstrokecolor{currentstroke}%
\pgfsetdash{}{0pt}%
\pgfpathmoveto{\pgfqpoint{0.740319in}{1.175101in}}%
\pgfpathlineto{\pgfqpoint{0.756074in}{1.155731in}}%
\pgfpathlineto{\pgfqpoint{0.771838in}{1.136431in}}%
\pgfpathlineto{\pgfqpoint{0.780142in}{1.143353in}}%
\pgfpathlineto{\pgfqpoint{0.788416in}{1.150626in}}%
\pgfpathlineto{\pgfqpoint{0.772693in}{1.169687in}}%
\pgfpathlineto{\pgfqpoint{0.756980in}{1.188829in}}%
\pgfpathlineto{\pgfqpoint{0.748665in}{1.181801in}}%
\pgfpathlineto{\pgfqpoint{0.740319in}{1.175101in}}%
\pgfpathclose%
\pgfusepath{fill}%
\end{pgfscope}%
\begin{pgfscope}%
\pgfpathrectangle{\pgfqpoint{0.407592in}{0.222904in}}{\pgfqpoint{1.411163in}{1.411163in}}%
\pgfusepath{clip}%
\pgfsetbuttcap%
\pgfsetroundjoin%
\definecolor{currentfill}{rgb}{0.556863,0.654902,0.847059}%
\pgfsetfillcolor{currentfill}%
\pgfsetfillopacity{0.130000}%
\pgfsetlinewidth{0.000000pt}%
\definecolor{currentstroke}{rgb}{0.000000,0.000000,0.000000}%
\pgfsetstrokecolor{currentstroke}%
\pgfsetdash{}{0pt}%
\pgfpathmoveto{\pgfqpoint{1.230679in}{0.990721in}}%
\pgfpathlineto{\pgfqpoint{1.246673in}{0.986620in}}%
\pgfpathlineto{\pgfqpoint{1.262713in}{0.983505in}}%
\pgfpathlineto{\pgfqpoint{1.270411in}{0.999391in}}%
\pgfpathlineto{\pgfqpoint{1.278095in}{1.015365in}}%
\pgfpathlineto{\pgfqpoint{1.262065in}{1.018489in}}%
\pgfpathlineto{\pgfqpoint{1.246081in}{1.022524in}}%
\pgfpathlineto{\pgfqpoint{1.238386in}{1.006579in}}%
\pgfpathlineto{\pgfqpoint{1.230679in}{0.990721in}}%
\pgfpathclose%
\pgfusepath{fill}%
\end{pgfscope}%
\begin{pgfscope}%
\pgfpathrectangle{\pgfqpoint{0.407592in}{0.222904in}}{\pgfqpoint{1.411163in}{1.411163in}}%
\pgfusepath{clip}%
\pgfsetbuttcap%
\pgfsetroundjoin%
\definecolor{currentfill}{rgb}{0.556863,0.654902,0.847059}%
\pgfsetfillcolor{currentfill}%
\pgfsetfillopacity{0.130000}%
\pgfsetlinewidth{0.000000pt}%
\definecolor{currentstroke}{rgb}{0.000000,0.000000,0.000000}%
\pgfsetstrokecolor{currentstroke}%
\pgfsetdash{}{0pt}%
\pgfpathmoveto{\pgfqpoint{0.882944in}{1.038690in}}%
\pgfpathlineto{\pgfqpoint{0.898731in}{1.020610in}}%
\pgfpathlineto{\pgfqpoint{0.914528in}{1.002763in}}%
\pgfpathlineto{\pgfqpoint{0.922615in}{1.011802in}}%
\pgfpathlineto{\pgfqpoint{0.930677in}{1.021276in}}%
\pgfpathlineto{\pgfqpoint{0.914917in}{1.038674in}}%
\pgfpathlineto{\pgfqpoint{0.899168in}{1.056333in}}%
\pgfpathlineto{\pgfqpoint{0.891069in}{1.047307in}}%
\pgfpathlineto{\pgfqpoint{0.882944in}{1.038690in}}%
\pgfpathclose%
\pgfusepath{fill}%
\end{pgfscope}%
\begin{pgfscope}%
\pgfpathrectangle{\pgfqpoint{0.407592in}{0.222904in}}{\pgfqpoint{1.411163in}{1.411163in}}%
\pgfusepath{clip}%
\pgfsetbuttcap%
\pgfsetroundjoin%
\definecolor{currentfill}{rgb}{0.556863,0.654902,0.847059}%
\pgfsetfillcolor{currentfill}%
\pgfsetfillopacity{0.130000}%
\pgfsetlinewidth{0.000000pt}%
\definecolor{currentstroke}{rgb}{0.000000,0.000000,0.000000}%
\pgfsetstrokecolor{currentstroke}%
\pgfsetdash{}{0pt}%
\pgfpathmoveto{\pgfqpoint{1.340873in}{1.076026in}}%
\pgfpathlineto{\pgfqpoint{1.357034in}{1.075246in}}%
\pgfpathlineto{\pgfqpoint{1.373249in}{1.075222in}}%
\pgfpathlineto{\pgfqpoint{1.380851in}{1.091245in}}%
\pgfpathlineto{\pgfqpoint{1.364641in}{1.091362in}}%
\pgfpathlineto{\pgfqpoint{1.348485in}{1.092211in}}%
\pgfusepath{fill}%
\end{pgfscope}%
\begin{pgfscope}%
\pgfpathrectangle{\pgfqpoint{0.407592in}{0.222904in}}{\pgfqpoint{1.411163in}{1.411163in}}%
\pgfusepath{clip}%
\pgfsetbuttcap%
\pgfsetroundjoin%
\definecolor{currentfill}{rgb}{0.556863,0.654902,0.847059}%
\pgfsetfillcolor{currentfill}%
\pgfsetfillopacity{0.130000}%
\pgfsetlinewidth{0.000000pt}%
\definecolor{currentstroke}{rgb}{0.000000,0.000000,0.000000}%
\pgfsetstrokecolor{currentstroke}%
\pgfsetdash{}{0pt}%
\pgfpathmoveto{\pgfqpoint{0.771838in}{1.136431in}}%
\pgfpathlineto{\pgfqpoint{0.787611in}{1.117209in}}%
\pgfpathlineto{\pgfqpoint{0.803393in}{1.098077in}}%
\pgfpathlineto{\pgfqpoint{0.811655in}{1.105246in}}%
\pgfpathlineto{\pgfqpoint{0.819889in}{1.112792in}}%
\pgfpathlineto{\pgfqpoint{0.804148in}{1.131657in}}%
\pgfpathlineto{\pgfqpoint{0.788416in}{1.150626in}}%
\pgfpathlineto{\pgfqpoint{0.780142in}{1.143353in}}%
\pgfpathlineto{\pgfqpoint{0.771838in}{1.136431in}}%
\pgfpathclose%
\pgfusepath{fill}%
\end{pgfscope}%
\begin{pgfscope}%
\pgfpathrectangle{\pgfqpoint{0.407592in}{0.222904in}}{\pgfqpoint{1.411163in}{1.411163in}}%
\pgfusepath{clip}%
\pgfsetbuttcap%
\pgfsetroundjoin%
\definecolor{currentfill}{rgb}{0.556863,0.654902,0.847059}%
\pgfsetfillcolor{currentfill}%
\pgfsetfillopacity{0.130000}%
\pgfsetlinewidth{0.000000pt}%
\definecolor{currentstroke}{rgb}{0.000000,0.000000,0.000000}%
\pgfsetstrokecolor{currentstroke}%
\pgfsetdash{}{0pt}%
\pgfpathmoveto{\pgfqpoint{1.278095in}{1.015365in}}%
\pgfpathlineto{\pgfqpoint{1.294173in}{1.013150in}}%
\pgfpathlineto{\pgfqpoint{1.310301in}{1.011832in}}%
\pgfpathlineto{\pgfqpoint{1.317963in}{1.027780in}}%
\pgfpathlineto{\pgfqpoint{1.325612in}{1.043802in}}%
\pgfpathlineto{\pgfqpoint{1.309494in}{1.045246in}}%
\pgfpathlineto{\pgfqpoint{1.293426in}{1.047526in}}%
\pgfpathlineto{\pgfqpoint{1.285767in}{1.031413in}}%
\pgfpathlineto{\pgfqpoint{1.278095in}{1.015365in}}%
\pgfpathclose%
\pgfusepath{fill}%
\end{pgfscope}%
\begin{pgfscope}%
\pgfpathrectangle{\pgfqpoint{0.407592in}{0.222904in}}{\pgfqpoint{1.411163in}{1.411163in}}%
\pgfusepath{clip}%
\pgfsetbuttcap%
\pgfsetroundjoin%
\definecolor{currentfill}{rgb}{0.556863,0.654902,0.847059}%
\pgfsetfillcolor{currentfill}%
\pgfsetfillopacity{0.130000}%
\pgfsetlinewidth{0.000000pt}%
\definecolor{currentstroke}{rgb}{0.000000,0.000000,0.000000}%
\pgfsetstrokecolor{currentstroke}%
\pgfsetdash{}{0pt}%
\pgfpathmoveto{\pgfqpoint{0.993826in}{0.955152in}}%
\pgfpathlineto{\pgfqpoint{1.009644in}{0.939772in}}%
\pgfpathlineto{\pgfqpoint{1.025476in}{0.925000in}}%
\pgfpathlineto{\pgfqpoint{1.033413in}{0.936627in}}%
\pgfpathlineto{\pgfqpoint{1.041329in}{0.948685in}}%
\pgfpathlineto{\pgfqpoint{1.025527in}{0.962816in}}%
\pgfpathlineto{\pgfqpoint{1.009740in}{0.977571in}}%
\pgfpathlineto{\pgfqpoint{1.001794in}{0.966150in}}%
\pgfpathlineto{\pgfqpoint{0.993826in}{0.955152in}}%
\pgfpathclose%
\pgfusepath{fill}%
\end{pgfscope}%
\begin{pgfscope}%
\pgfpathrectangle{\pgfqpoint{0.407592in}{0.222904in}}{\pgfqpoint{1.411163in}{1.411163in}}%
\pgfusepath{clip}%
\pgfsetbuttcap%
\pgfsetroundjoin%
\definecolor{currentfill}{rgb}{0.556863,0.654902,0.847059}%
\pgfsetfillcolor{currentfill}%
\pgfsetfillopacity{0.130000}%
\pgfsetlinewidth{0.000000pt}%
\definecolor{currentstroke}{rgb}{0.000000,0.000000,0.000000}%
\pgfsetstrokecolor{currentstroke}%
\pgfsetdash{}{0pt}%
\pgfpathmoveto{\pgfqpoint{0.914528in}{1.002763in}}%
\pgfpathlineto{\pgfqpoint{0.930335in}{0.985188in}}%
\pgfpathlineto{\pgfqpoint{0.946153in}{0.967927in}}%
\pgfpathlineto{\pgfqpoint{0.954203in}{0.977441in}}%
\pgfpathlineto{\pgfqpoint{0.962228in}{0.987419in}}%
\pgfpathlineto{\pgfqpoint{0.946447in}{1.004176in}}%
\pgfpathlineto{\pgfqpoint{0.930677in}{1.021276in}}%
\pgfpathlineto{\pgfqpoint{0.922615in}{1.011802in}}%
\pgfpathlineto{\pgfqpoint{0.914528in}{1.002763in}}%
\pgfpathclose%
\pgfusepath{fill}%
\end{pgfscope}%
\begin{pgfscope}%
\pgfpathrectangle{\pgfqpoint{0.407592in}{0.222904in}}{\pgfqpoint{1.411163in}{1.411163in}}%
\pgfusepath{clip}%
\pgfsetbuttcap%
\pgfsetroundjoin%
\definecolor{currentfill}{rgb}{0.556863,0.654902,0.847059}%
\pgfsetfillcolor{currentfill}%
\pgfsetfillopacity{0.130000}%
\pgfsetlinewidth{0.000000pt}%
\definecolor{currentstroke}{rgb}{0.000000,0.000000,0.000000}%
\pgfsetstrokecolor{currentstroke}%
\pgfsetdash{}{0pt}%
\pgfpathmoveto{\pgfqpoint{1.120408in}{0.927876in}}%
\pgfpathlineto{\pgfqpoint{1.136296in}{0.918755in}}%
\pgfpathlineto{\pgfqpoint{1.152214in}{0.910754in}}%
\pgfpathlineto{\pgfqpoint{1.160016in}{0.925463in}}%
\pgfpathlineto{\pgfqpoint{1.167803in}{0.940398in}}%
\pgfpathlineto{\pgfqpoint{1.151900in}{0.948023in}}%
\pgfpathlineto{\pgfqpoint{1.136029in}{0.956689in}}%
\pgfpathlineto{\pgfqpoint{1.128227in}{0.942159in}}%
\pgfpathlineto{\pgfqpoint{1.120408in}{0.927876in}}%
\pgfpathclose%
\pgfusepath{fill}%
\end{pgfscope}%
\begin{pgfscope}%
\pgfpathrectangle{\pgfqpoint{0.407592in}{0.222904in}}{\pgfqpoint{1.411163in}{1.411163in}}%
\pgfusepath{clip}%
\pgfsetbuttcap%
\pgfsetroundjoin%
\definecolor{currentfill}{rgb}{0.556863,0.654902,0.847059}%
\pgfsetfillcolor{currentfill}%
\pgfsetfillopacity{0.130000}%
\pgfsetlinewidth{0.000000pt}%
\definecolor{currentstroke}{rgb}{0.000000,0.000000,0.000000}%
\pgfsetstrokecolor{currentstroke}%
\pgfsetdash{}{0pt}%
\pgfpathmoveto{\pgfqpoint{0.803393in}{1.098077in}}%
\pgfpathlineto{\pgfqpoint{0.819183in}{1.079049in}}%
\pgfpathlineto{\pgfqpoint{0.834983in}{1.060138in}}%
\pgfpathlineto{\pgfqpoint{0.843205in}{1.067581in}}%
\pgfpathlineto{\pgfqpoint{0.851397in}{1.075431in}}%
\pgfpathlineto{\pgfqpoint{0.835638in}{1.094044in}}%
\pgfpathlineto{\pgfqpoint{0.819889in}{1.112792in}}%
\pgfpathlineto{\pgfqpoint{0.811655in}{1.105246in}}%
\pgfpathlineto{\pgfqpoint{0.803393in}{1.098077in}}%
\pgfpathclose%
\pgfusepath{fill}%
\end{pgfscope}%
\begin{pgfscope}%
\pgfpathrectangle{\pgfqpoint{0.407592in}{0.222904in}}{\pgfqpoint{1.411163in}{1.411163in}}%
\pgfusepath{clip}%
\pgfsetbuttcap%
\pgfsetroundjoin%
\definecolor{currentfill}{rgb}{0.556863,0.654902,0.847059}%
\pgfsetfillcolor{currentfill}%
\pgfsetfillopacity{0.130000}%
\pgfsetlinewidth{0.000000pt}%
\definecolor{currentstroke}{rgb}{0.000000,0.000000,0.000000}%
\pgfsetstrokecolor{currentstroke}%
\pgfsetdash{}{0pt}%
\pgfpathmoveto{\pgfqpoint{1.167803in}{0.940398in}}%
\pgfpathlineto{\pgfqpoint{1.183741in}{0.933866in}}%
\pgfpathlineto{\pgfqpoint{1.199716in}{0.928468in}}%
\pgfpathlineto{\pgfqpoint{1.207477in}{0.943816in}}%
\pgfpathlineto{\pgfqpoint{1.215224in}{0.959322in}}%
\pgfpathlineto{\pgfqpoint{1.199261in}{0.964553in}}%
\pgfpathlineto{\pgfqpoint{1.183336in}{0.970828in}}%
\pgfpathlineto{\pgfqpoint{1.175576in}{0.955528in}}%
\pgfpathlineto{\pgfqpoint{1.167803in}{0.940398in}}%
\pgfpathclose%
\pgfusepath{fill}%
\end{pgfscope}%
\begin{pgfscope}%
\pgfpathrectangle{\pgfqpoint{0.407592in}{0.222904in}}{\pgfqpoint{1.411163in}{1.411163in}}%
\pgfusepath{clip}%
\pgfsetbuttcap%
\pgfsetroundjoin%
\definecolor{currentfill}{rgb}{0.556863,0.654902,0.847059}%
\pgfsetfillcolor{currentfill}%
\pgfsetfillopacity{0.130000}%
\pgfsetlinewidth{0.000000pt}%
\definecolor{currentstroke}{rgb}{0.000000,0.000000,0.000000}%
\pgfsetstrokecolor{currentstroke}%
\pgfsetdash{}{0pt}%
\pgfpathmoveto{\pgfqpoint{1.072986in}{0.922640in}}%
\pgfpathlineto{\pgfqpoint{1.088844in}{0.910918in}}%
\pgfpathlineto{\pgfqpoint{1.104724in}{0.900197in}}%
\pgfpathlineto{\pgfqpoint{1.112574in}{0.913875in}}%
\pgfpathlineto{\pgfqpoint{1.120408in}{0.927876in}}%
\pgfpathlineto{\pgfqpoint{1.104548in}{0.938038in}}%
\pgfpathlineto{\pgfqpoint{1.088712in}{0.949159in}}%
\pgfpathlineto{\pgfqpoint{1.080858in}{0.935729in}}%
\pgfpathlineto{\pgfqpoint{1.072986in}{0.922640in}}%
\pgfpathclose%
\pgfusepath{fill}%
\end{pgfscope}%
\begin{pgfscope}%
\pgfpathrectangle{\pgfqpoint{0.407592in}{0.222904in}}{\pgfqpoint{1.411163in}{1.411163in}}%
\pgfusepath{clip}%
\pgfsetbuttcap%
\pgfsetroundjoin%
\definecolor{currentfill}{rgb}{0.556863,0.654902,0.847059}%
\pgfsetfillcolor{currentfill}%
\pgfsetfillopacity{0.130000}%
\pgfsetlinewidth{0.000000pt}%
\definecolor{currentstroke}{rgb}{0.000000,0.000000,0.000000}%
\pgfsetstrokecolor{currentstroke}%
\pgfsetdash{}{0pt}%
\pgfpathmoveto{\pgfqpoint{1.325612in}{1.043802in}}%
\pgfpathlineto{\pgfqpoint{1.341783in}{1.043180in}}%
\pgfpathlineto{\pgfqpoint{1.358008in}{1.043360in}}%
\pgfpathlineto{\pgfqpoint{1.365634in}{1.059257in}}%
\pgfpathlineto{\pgfqpoint{1.373249in}{1.075222in}}%
\pgfpathlineto{\pgfqpoint{1.357034in}{1.075246in}}%
\pgfpathlineto{\pgfqpoint{1.340873in}{1.076026in}}%
\pgfpathlineto{\pgfqpoint{1.333249in}{1.059887in}}%
\pgfpathlineto{\pgfqpoint{1.325612in}{1.043802in}}%
\pgfpathclose%
\pgfusepath{fill}%
\end{pgfscope}%
\begin{pgfscope}%
\pgfpathrectangle{\pgfqpoint{0.407592in}{0.222904in}}{\pgfqpoint{1.411163in}{1.411163in}}%
\pgfusepath{clip}%
\pgfsetbuttcap%
\pgfsetroundjoin%
\definecolor{currentfill}{rgb}{0.556863,0.654902,0.847059}%
\pgfsetfillcolor{currentfill}%
\pgfsetfillopacity{0.130000}%
\pgfsetlinewidth{0.000000pt}%
\definecolor{currentstroke}{rgb}{0.000000,0.000000,0.000000}%
\pgfsetstrokecolor{currentstroke}%
\pgfsetdash{}{0pt}%
\pgfpathmoveto{\pgfqpoint{1.215224in}{0.959322in}}%
\pgfpathlineto{\pgfqpoint{1.231229in}{0.955153in}}%
\pgfpathlineto{\pgfqpoint{1.247279in}{0.952054in}}%
\pgfpathlineto{\pgfqpoint{1.255003in}{0.967721in}}%
\pgfpathlineto{\pgfqpoint{1.262713in}{0.983505in}}%
\pgfpathlineto{\pgfqpoint{1.246673in}{0.986620in}}%
\pgfpathlineto{\pgfqpoint{1.230679in}{0.990721in}}%
\pgfpathlineto{\pgfqpoint{1.222958in}{0.974963in}}%
\pgfpathlineto{\pgfqpoint{1.215224in}{0.959322in}}%
\pgfpathclose%
\pgfusepath{fill}%
\end{pgfscope}%
\begin{pgfscope}%
\pgfpathrectangle{\pgfqpoint{0.407592in}{0.222904in}}{\pgfqpoint{1.411163in}{1.411163in}}%
\pgfusepath{clip}%
\pgfsetbuttcap%
\pgfsetroundjoin%
\definecolor{currentfill}{rgb}{0.556863,0.654902,0.847059}%
\pgfsetfillcolor{currentfill}%
\pgfsetfillopacity{0.130000}%
\pgfsetlinewidth{0.000000pt}%
\definecolor{currentstroke}{rgb}{0.000000,0.000000,0.000000}%
\pgfsetstrokecolor{currentstroke}%
\pgfsetdash{}{0pt}%
\pgfpathmoveto{\pgfqpoint{0.691968in}{1.201968in}}%
\pgfpathlineto{\pgfqpoint{0.707747in}{1.182317in}}%
\pgfpathlineto{\pgfqpoint{0.723535in}{1.162712in}}%
\pgfpathlineto{\pgfqpoint{0.731943in}{1.168737in}}%
\pgfpathlineto{\pgfqpoint{0.740319in}{1.175101in}}%
\pgfpathlineto{\pgfqpoint{0.724573in}{1.194532in}}%
\pgfpathlineto{\pgfqpoint{0.708836in}{1.214017in}}%
\pgfpathlineto{\pgfqpoint{0.700418in}{1.207834in}}%
\pgfpathlineto{\pgfqpoint{0.691968in}{1.201968in}}%
\pgfpathclose%
\pgfusepath{fill}%
\end{pgfscope}%
\begin{pgfscope}%
\pgfpathrectangle{\pgfqpoint{0.407592in}{0.222904in}}{\pgfqpoint{1.411163in}{1.411163in}}%
\pgfusepath{clip}%
\pgfsetbuttcap%
\pgfsetroundjoin%
\definecolor{currentfill}{rgb}{0.556863,0.654902,0.847059}%
\pgfsetfillcolor{currentfill}%
\pgfsetfillopacity{0.130000}%
\pgfsetlinewidth{0.000000pt}%
\definecolor{currentstroke}{rgb}{0.000000,0.000000,0.000000}%
\pgfsetstrokecolor{currentstroke}%
\pgfsetdash{}{0pt}%
\pgfpathmoveto{\pgfqpoint{0.834983in}{1.060138in}}%
\pgfpathlineto{\pgfqpoint{0.850791in}{1.041362in}}%
\pgfpathlineto{\pgfqpoint{0.866609in}{1.022743in}}%
\pgfpathlineto{\pgfqpoint{0.874791in}{1.030496in}}%
\pgfpathlineto{\pgfqpoint{0.882944in}{1.038690in}}%
\pgfpathlineto{\pgfqpoint{0.867166in}{1.056972in}}%
\pgfpathlineto{\pgfqpoint{0.851397in}{1.075431in}}%
\pgfpathlineto{\pgfqpoint{0.843205in}{1.067581in}}%
\pgfpathlineto{\pgfqpoint{0.834983in}{1.060138in}}%
\pgfpathclose%
\pgfusepath{fill}%
\end{pgfscope}%
\begin{pgfscope}%
\pgfpathrectangle{\pgfqpoint{0.407592in}{0.222904in}}{\pgfqpoint{1.411163in}{1.411163in}}%
\pgfusepath{clip}%
\pgfsetbuttcap%
\pgfsetroundjoin%
\definecolor{currentfill}{rgb}{0.556863,0.654902,0.847059}%
\pgfsetfillcolor{currentfill}%
\pgfsetfillopacity{0.130000}%
\pgfsetlinewidth{0.000000pt}%
\definecolor{currentstroke}{rgb}{0.000000,0.000000,0.000000}%
\pgfsetstrokecolor{currentstroke}%
\pgfsetdash{}{0pt}%
\pgfpathmoveto{\pgfqpoint{1.025476in}{0.925000in}}%
\pgfpathlineto{\pgfqpoint{1.041323in}{0.910928in}}%
\pgfpathlineto{\pgfqpoint{1.057187in}{0.897660in}}%
\pgfpathlineto{\pgfqpoint{1.065096in}{0.909935in}}%
\pgfpathlineto{\pgfqpoint{1.072986in}{0.922640in}}%
\pgfpathlineto{\pgfqpoint{1.057148in}{0.935262in}}%
\pgfpathlineto{\pgfqpoint{1.041329in}{0.948685in}}%
\pgfpathlineto{\pgfqpoint{1.033413in}{0.936627in}}%
\pgfpathlineto{\pgfqpoint{1.025476in}{0.925000in}}%
\pgfpathclose%
\pgfusepath{fill}%
\end{pgfscope}%
\begin{pgfscope}%
\pgfpathrectangle{\pgfqpoint{0.407592in}{0.222904in}}{\pgfqpoint{1.411163in}{1.411163in}}%
\pgfusepath{clip}%
\pgfsetbuttcap%
\pgfsetroundjoin%
\definecolor{currentfill}{rgb}{0.556863,0.654902,0.847059}%
\pgfsetfillcolor{currentfill}%
\pgfsetfillopacity{0.130000}%
\pgfsetlinewidth{0.000000pt}%
\definecolor{currentstroke}{rgb}{0.000000,0.000000,0.000000}%
\pgfsetstrokecolor{currentstroke}%
\pgfsetdash{}{0pt}%
\pgfpathmoveto{\pgfqpoint{0.946153in}{0.967927in}}%
\pgfpathlineto{\pgfqpoint{0.961982in}{0.951031in}}%
\pgfpathlineto{\pgfqpoint{0.977821in}{0.934563in}}%
\pgfpathlineto{\pgfqpoint{0.985835in}{0.944611in}}%
\pgfpathlineto{\pgfqpoint{0.993826in}{0.955152in}}%
\pgfpathlineto{\pgfqpoint{0.978021in}{0.971057in}}%
\pgfpathlineto{\pgfqpoint{0.962228in}{0.987419in}}%
\pgfpathlineto{\pgfqpoint{0.954203in}{0.977441in}}%
\pgfpathlineto{\pgfqpoint{0.946153in}{0.967927in}}%
\pgfpathclose%
\pgfusepath{fill}%
\end{pgfscope}%
\begin{pgfscope}%
\pgfpathrectangle{\pgfqpoint{0.407592in}{0.222904in}}{\pgfqpoint{1.411163in}{1.411163in}}%
\pgfusepath{clip}%
\pgfsetbuttcap%
\pgfsetroundjoin%
\definecolor{currentfill}{rgb}{0.556863,0.654902,0.847059}%
\pgfsetfillcolor{currentfill}%
\pgfsetfillopacity{0.130000}%
\pgfsetlinewidth{0.000000pt}%
\definecolor{currentstroke}{rgb}{0.000000,0.000000,0.000000}%
\pgfsetstrokecolor{currentstroke}%
\pgfsetdash{}{0pt}%
\pgfpathmoveto{\pgfqpoint{1.262713in}{0.983505in}}%
\pgfpathlineto{\pgfqpoint{1.278801in}{0.981371in}}%
\pgfpathlineto{\pgfqpoint{1.294939in}{0.980202in}}%
\pgfpathlineto{\pgfqpoint{1.302626in}{0.995968in}}%
\pgfpathlineto{\pgfqpoint{1.310301in}{1.011832in}}%
\pgfpathlineto{\pgfqpoint{1.294173in}{1.013150in}}%
\pgfpathlineto{\pgfqpoint{1.278095in}{1.015365in}}%
\pgfpathlineto{\pgfqpoint{1.270411in}{0.999391in}}%
\pgfpathlineto{\pgfqpoint{1.262713in}{0.983505in}}%
\pgfpathclose%
\pgfusepath{fill}%
\end{pgfscope}%
\begin{pgfscope}%
\pgfpathrectangle{\pgfqpoint{0.407592in}{0.222904in}}{\pgfqpoint{1.411163in}{1.411163in}}%
\pgfusepath{clip}%
\pgfsetbuttcap%
\pgfsetroundjoin%
\definecolor{currentfill}{rgb}{0.556863,0.654902,0.847059}%
\pgfsetfillcolor{currentfill}%
\pgfsetfillopacity{0.130000}%
\pgfsetlinewidth{0.000000pt}%
\definecolor{currentstroke}{rgb}{0.000000,0.000000,0.000000}%
\pgfsetstrokecolor{currentstroke}%
\pgfsetdash{}{0pt}%
\pgfpathmoveto{\pgfqpoint{0.723535in}{1.162712in}}%
\pgfpathlineto{\pgfqpoint{0.739332in}{1.143161in}}%
\pgfpathlineto{\pgfqpoint{0.755138in}{1.123671in}}%
\pgfpathlineto{\pgfqpoint{0.763504in}{1.129868in}}%
\pgfpathlineto{\pgfqpoint{0.771838in}{1.136431in}}%
\pgfpathlineto{\pgfqpoint{0.756074in}{1.155731in}}%
\pgfpathlineto{\pgfqpoint{0.740319in}{1.175101in}}%
\pgfpathlineto{\pgfqpoint{0.731943in}{1.168737in}}%
\pgfpathlineto{\pgfqpoint{0.723535in}{1.162712in}}%
\pgfpathclose%
\pgfusepath{fill}%
\end{pgfscope}%
\begin{pgfscope}%
\pgfpathrectangle{\pgfqpoint{0.407592in}{0.222904in}}{\pgfqpoint{1.411163in}{1.411163in}}%
\pgfusepath{clip}%
\pgfsetbuttcap%
\pgfsetroundjoin%
\definecolor{currentfill}{rgb}{0.556863,0.654902,0.847059}%
\pgfsetfillcolor{currentfill}%
\pgfsetfillopacity{0.130000}%
\pgfsetlinewidth{0.000000pt}%
\definecolor{currentstroke}{rgb}{0.000000,0.000000,0.000000}%
\pgfsetstrokecolor{currentstroke}%
\pgfsetdash{}{0pt}%
\pgfpathmoveto{\pgfqpoint{0.866609in}{1.022743in}}%
\pgfpathlineto{\pgfqpoint{0.882436in}{1.004305in}}%
\pgfpathlineto{\pgfqpoint{0.898273in}{0.986078in}}%
\pgfpathlineto{\pgfqpoint{0.906414in}{0.994182in}}%
\pgfpathlineto{\pgfqpoint{0.914528in}{1.002763in}}%
\pgfpathlineto{\pgfqpoint{0.898731in}{1.020610in}}%
\pgfpathlineto{\pgfqpoint{0.882944in}{1.038690in}}%
\pgfpathlineto{\pgfqpoint{0.874791in}{1.030496in}}%
\pgfpathlineto{\pgfqpoint{0.866609in}{1.022743in}}%
\pgfpathclose%
\pgfusepath{fill}%
\end{pgfscope}%
\begin{pgfscope}%
\pgfpathrectangle{\pgfqpoint{0.407592in}{0.222904in}}{\pgfqpoint{1.411163in}{1.411163in}}%
\pgfusepath{clip}%
\pgfsetbuttcap%
\pgfsetroundjoin%
\definecolor{currentfill}{rgb}{0.556863,0.654902,0.847059}%
\pgfsetfillcolor{currentfill}%
\pgfsetfillopacity{0.130000}%
\pgfsetlinewidth{0.000000pt}%
\definecolor{currentstroke}{rgb}{0.000000,0.000000,0.000000}%
\pgfsetstrokecolor{currentstroke}%
\pgfsetdash{}{0pt}%
\pgfpathmoveto{\pgfqpoint{0.977821in}{0.934563in}}%
\pgfpathlineto{\pgfqpoint{0.993673in}{0.918594in}}%
\pgfpathlineto{\pgfqpoint{1.009537in}{0.903208in}}%
\pgfpathlineto{\pgfqpoint{1.017518in}{0.913846in}}%
\pgfpathlineto{\pgfqpoint{1.025476in}{0.925000in}}%
\pgfpathlineto{\pgfqpoint{1.009644in}{0.939772in}}%
\pgfpathlineto{\pgfqpoint{0.993826in}{0.955152in}}%
\pgfpathlineto{\pgfqpoint{0.985835in}{0.944611in}}%
\pgfpathlineto{\pgfqpoint{0.977821in}{0.934563in}}%
\pgfpathclose%
\pgfusepath{fill}%
\end{pgfscope}%
\begin{pgfscope}%
\pgfpathrectangle{\pgfqpoint{0.407592in}{0.222904in}}{\pgfqpoint{1.411163in}{1.411163in}}%
\pgfusepath{clip}%
\pgfsetbuttcap%
\pgfsetroundjoin%
\definecolor{currentfill}{rgb}{0.556863,0.654902,0.847059}%
\pgfsetfillcolor{currentfill}%
\pgfsetfillopacity{0.130000}%
\pgfsetlinewidth{0.000000pt}%
\definecolor{currentstroke}{rgb}{0.000000,0.000000,0.000000}%
\pgfsetstrokecolor{currentstroke}%
\pgfsetdash{}{0pt}%
\pgfpathmoveto{\pgfqpoint{1.104724in}{0.900197in}}%
\pgfpathlineto{\pgfqpoint{1.120630in}{0.890580in}}%
\pgfpathlineto{\pgfqpoint{1.136564in}{0.882163in}}%
\pgfpathlineto{\pgfqpoint{1.144397in}{0.896307in}}%
\pgfpathlineto{\pgfqpoint{1.152214in}{0.910754in}}%
\pgfpathlineto{\pgfqpoint{1.136296in}{0.918755in}}%
\pgfpathlineto{\pgfqpoint{1.120408in}{0.927876in}}%
\pgfpathlineto{\pgfqpoint{1.112574in}{0.913875in}}%
\pgfpathlineto{\pgfqpoint{1.104724in}{0.900197in}}%
\pgfpathclose%
\pgfusepath{fill}%
\end{pgfscope}%
\begin{pgfscope}%
\pgfpathrectangle{\pgfqpoint{0.407592in}{0.222904in}}{\pgfqpoint{1.411163in}{1.411163in}}%
\pgfusepath{clip}%
\pgfsetbuttcap%
\pgfsetroundjoin%
\definecolor{currentfill}{rgb}{0.556863,0.654902,0.847059}%
\pgfsetfillcolor{currentfill}%
\pgfsetfillopacity{0.130000}%
\pgfsetlinewidth{0.000000pt}%
\definecolor{currentstroke}{rgb}{0.000000,0.000000,0.000000}%
\pgfsetstrokecolor{currentstroke}%
\pgfsetdash{}{0pt}%
\pgfpathmoveto{\pgfqpoint{1.373249in}{1.075222in}}%
\pgfpathlineto{\pgfqpoint{1.389519in}{1.075932in}}%
\pgfpathlineto{\pgfqpoint{1.405848in}{1.077353in}}%
\pgfpathlineto{\pgfqpoint{1.413439in}{1.093134in}}%
\pgfpathlineto{\pgfqpoint{1.397116in}{1.091843in}}%
\pgfpathlineto{\pgfqpoint{1.380851in}{1.091245in}}%
\pgfusepath{fill}%
\end{pgfscope}%
\begin{pgfscope}%
\pgfpathrectangle{\pgfqpoint{0.407592in}{0.222904in}}{\pgfqpoint{1.411163in}{1.411163in}}%
\pgfusepath{clip}%
\pgfsetbuttcap%
\pgfsetroundjoin%
\definecolor{currentfill}{rgb}{0.556863,0.654902,0.847059}%
\pgfsetfillcolor{currentfill}%
\pgfsetfillopacity{0.130000}%
\pgfsetlinewidth{0.000000pt}%
\definecolor{currentstroke}{rgb}{0.000000,0.000000,0.000000}%
\pgfsetstrokecolor{currentstroke}%
\pgfsetdash{}{0pt}%
\pgfpathmoveto{\pgfqpoint{1.152214in}{0.910754in}}%
\pgfpathlineto{\pgfqpoint{1.168165in}{0.903937in}}%
\pgfpathlineto{\pgfqpoint{1.184153in}{0.898356in}}%
\pgfpathlineto{\pgfqpoint{1.191941in}{0.913305in}}%
\pgfpathlineto{\pgfqpoint{1.199716in}{0.928468in}}%
\pgfpathlineto{\pgfqpoint{1.183741in}{0.933866in}}%
\pgfpathlineto{\pgfqpoint{1.167803in}{0.940398in}}%
\pgfpathlineto{\pgfqpoint{1.160016in}{0.925463in}}%
\pgfpathlineto{\pgfqpoint{1.152214in}{0.910754in}}%
\pgfpathclose%
\pgfusepath{fill}%
\end{pgfscope}%
\begin{pgfscope}%
\pgfpathrectangle{\pgfqpoint{0.407592in}{0.222904in}}{\pgfqpoint{1.411163in}{1.411163in}}%
\pgfusepath{clip}%
\pgfsetbuttcap%
\pgfsetroundjoin%
\definecolor{currentfill}{rgb}{0.556863,0.654902,0.847059}%
\pgfsetfillcolor{currentfill}%
\pgfsetfillopacity{0.130000}%
\pgfsetlinewidth{0.000000pt}%
\definecolor{currentstroke}{rgb}{0.000000,0.000000,0.000000}%
\pgfsetstrokecolor{currentstroke}%
\pgfsetdash{}{0pt}%
\pgfpathmoveto{\pgfqpoint{0.755138in}{1.123671in}}%
\pgfpathlineto{\pgfqpoint{0.770952in}{1.104251in}}%
\pgfpathlineto{\pgfqpoint{0.786776in}{1.084910in}}%
\pgfpathlineto{\pgfqpoint{0.795100in}{1.091296in}}%
\pgfpathlineto{\pgfqpoint{0.803393in}{1.098077in}}%
\pgfpathlineto{\pgfqpoint{0.787611in}{1.117209in}}%
\pgfpathlineto{\pgfqpoint{0.771838in}{1.136431in}}%
\pgfpathlineto{\pgfqpoint{0.763504in}{1.129868in}}%
\pgfpathlineto{\pgfqpoint{0.755138in}{1.123671in}}%
\pgfpathclose%
\pgfusepath{fill}%
\end{pgfscope}%
\begin{pgfscope}%
\pgfpathrectangle{\pgfqpoint{0.407592in}{0.222904in}}{\pgfqpoint{1.411163in}{1.411163in}}%
\pgfusepath{clip}%
\pgfsetbuttcap%
\pgfsetroundjoin%
\definecolor{currentfill}{rgb}{0.556863,0.654902,0.847059}%
\pgfsetfillcolor{currentfill}%
\pgfsetfillopacity{0.130000}%
\pgfsetlinewidth{0.000000pt}%
\definecolor{currentstroke}{rgb}{0.000000,0.000000,0.000000}%
\pgfsetstrokecolor{currentstroke}%
\pgfsetdash{}{0pt}%
\pgfpathmoveto{\pgfqpoint{1.310301in}{1.011832in}}%
\pgfpathlineto{\pgfqpoint{1.326481in}{1.011393in}}%
\pgfpathlineto{\pgfqpoint{1.342717in}{1.011809in}}%
\pgfpathlineto{\pgfqpoint{1.350369in}{1.027540in}}%
\pgfpathlineto{\pgfqpoint{1.358008in}{1.043360in}}%
\pgfpathlineto{\pgfqpoint{1.341783in}{1.043180in}}%
\pgfpathlineto{\pgfqpoint{1.325612in}{1.043802in}}%
\pgfpathlineto{\pgfqpoint{1.317963in}{1.027780in}}%
\pgfpathlineto{\pgfqpoint{1.310301in}{1.011832in}}%
\pgfpathclose%
\pgfusepath{fill}%
\end{pgfscope}%
\begin{pgfscope}%
\pgfpathrectangle{\pgfqpoint{0.407592in}{0.222904in}}{\pgfqpoint{1.411163in}{1.411163in}}%
\pgfusepath{clip}%
\pgfsetbuttcap%
\pgfsetroundjoin%
\definecolor{currentfill}{rgb}{0.556863,0.654902,0.847059}%
\pgfsetfillcolor{currentfill}%
\pgfsetfillopacity{0.130000}%
\pgfsetlinewidth{0.000000pt}%
\definecolor{currentstroke}{rgb}{0.000000,0.000000,0.000000}%
\pgfsetstrokecolor{currentstroke}%
\pgfsetdash{}{0pt}%
\pgfpathmoveto{\pgfqpoint{1.057187in}{0.897660in}}%
\pgfpathlineto{\pgfqpoint{1.073069in}{0.885309in}}%
\pgfpathlineto{\pgfqpoint{1.088972in}{0.873997in}}%
\pgfpathlineto{\pgfqpoint{1.096857in}{0.886887in}}%
\pgfpathlineto{\pgfqpoint{1.104724in}{0.900197in}}%
\pgfpathlineto{\pgfqpoint{1.088844in}{0.910918in}}%
\pgfpathlineto{\pgfqpoint{1.072986in}{0.922640in}}%
\pgfpathlineto{\pgfqpoint{1.065096in}{0.909935in}}%
\pgfpathlineto{\pgfqpoint{1.057187in}{0.897660in}}%
\pgfpathclose%
\pgfusepath{fill}%
\end{pgfscope}%
\begin{pgfscope}%
\pgfpathrectangle{\pgfqpoint{0.407592in}{0.222904in}}{\pgfqpoint{1.411163in}{1.411163in}}%
\pgfusepath{clip}%
\pgfsetbuttcap%
\pgfsetroundjoin%
\definecolor{currentfill}{rgb}{0.556863,0.654902,0.847059}%
\pgfsetfillcolor{currentfill}%
\pgfsetfillopacity{0.130000}%
\pgfsetlinewidth{0.000000pt}%
\definecolor{currentstroke}{rgb}{0.000000,0.000000,0.000000}%
\pgfsetstrokecolor{currentstroke}%
\pgfsetdash{}{0pt}%
\pgfpathmoveto{\pgfqpoint{1.199716in}{0.928468in}}%
\pgfpathlineto{\pgfqpoint{1.215732in}{0.924229in}}%
\pgfpathlineto{\pgfqpoint{1.231793in}{0.921155in}}%
\pgfpathlineto{\pgfqpoint{1.239543in}{0.936525in}}%
\pgfpathlineto{\pgfqpoint{1.247279in}{0.952054in}}%
\pgfpathlineto{\pgfqpoint{1.231229in}{0.955153in}}%
\pgfpathlineto{\pgfqpoint{1.215224in}{0.959322in}}%
\pgfpathlineto{\pgfqpoint{1.207477in}{0.943816in}}%
\pgfpathlineto{\pgfqpoint{1.199716in}{0.928468in}}%
\pgfpathclose%
\pgfusepath{fill}%
\end{pgfscope}%
\begin{pgfscope}%
\pgfpathrectangle{\pgfqpoint{0.407592in}{0.222904in}}{\pgfqpoint{1.411163in}{1.411163in}}%
\pgfusepath{clip}%
\pgfsetbuttcap%
\pgfsetroundjoin%
\definecolor{currentfill}{rgb}{0.556863,0.654902,0.847059}%
\pgfsetfillcolor{currentfill}%
\pgfsetfillopacity{0.130000}%
\pgfsetlinewidth{0.000000pt}%
\definecolor{currentstroke}{rgb}{0.000000,0.000000,0.000000}%
\pgfsetstrokecolor{currentstroke}%
\pgfsetdash{}{0pt}%
\pgfpathmoveto{\pgfqpoint{0.898273in}{0.986078in}}%
\pgfpathlineto{\pgfqpoint{0.914119in}{0.968095in}}%
\pgfpathlineto{\pgfqpoint{0.929975in}{0.950400in}}%
\pgfpathlineto{\pgfqpoint{0.938078in}{0.958905in}}%
\pgfpathlineto{\pgfqpoint{0.946153in}{0.967927in}}%
\pgfpathlineto{\pgfqpoint{0.930335in}{0.985188in}}%
\pgfpathlineto{\pgfqpoint{0.914528in}{1.002763in}}%
\pgfpathlineto{\pgfqpoint{0.906414in}{0.994182in}}%
\pgfpathlineto{\pgfqpoint{0.898273in}{0.986078in}}%
\pgfpathclose%
\pgfusepath{fill}%
\end{pgfscope}%
\begin{pgfscope}%
\pgfpathrectangle{\pgfqpoint{0.407592in}{0.222904in}}{\pgfqpoint{1.411163in}{1.411163in}}%
\pgfusepath{clip}%
\pgfsetbuttcap%
\pgfsetroundjoin%
\definecolor{currentfill}{rgb}{0.556863,0.654902,0.847059}%
\pgfsetfillcolor{currentfill}%
\pgfsetfillopacity{0.130000}%
\pgfsetlinewidth{0.000000pt}%
\definecolor{currentstroke}{rgb}{0.000000,0.000000,0.000000}%
\pgfsetstrokecolor{currentstroke}%
\pgfsetdash{}{0pt}%
\pgfpathmoveto{\pgfqpoint{0.786776in}{1.084910in}}%
\pgfpathlineto{\pgfqpoint{0.802608in}{1.065660in}}%
\pgfpathlineto{\pgfqpoint{0.818449in}{1.046516in}}%
\pgfpathlineto{\pgfqpoint{0.826731in}{1.053113in}}%
\pgfpathlineto{\pgfqpoint{0.834983in}{1.060138in}}%
\pgfpathlineto{\pgfqpoint{0.819183in}{1.079049in}}%
\pgfpathlineto{\pgfqpoint{0.803393in}{1.098077in}}%
\pgfpathlineto{\pgfqpoint{0.795100in}{1.091296in}}%
\pgfpathlineto{\pgfqpoint{0.786776in}{1.084910in}}%
\pgfpathclose%
\pgfusepath{fill}%
\end{pgfscope}%
\begin{pgfscope}%
\pgfpathrectangle{\pgfqpoint{0.407592in}{0.222904in}}{\pgfqpoint{1.411163in}{1.411163in}}%
\pgfusepath{clip}%
\pgfsetbuttcap%
\pgfsetroundjoin%
\definecolor{currentfill}{rgb}{0.556863,0.654902,0.847059}%
\pgfsetfillcolor{currentfill}%
\pgfsetfillopacity{0.130000}%
\pgfsetlinewidth{0.000000pt}%
\definecolor{currentstroke}{rgb}{0.000000,0.000000,0.000000}%
\pgfsetstrokecolor{currentstroke}%
\pgfsetdash{}{0pt}%
\pgfpathmoveto{\pgfqpoint{1.247279in}{0.952054in}}%
\pgfpathlineto{\pgfqpoint{1.263377in}{0.950018in}}%
\pgfpathlineto{\pgfqpoint{1.279526in}{0.949027in}}%
\pgfpathlineto{\pgfqpoint{1.287239in}{0.964549in}}%
\pgfpathlineto{\pgfqpoint{1.294939in}{0.980202in}}%
\pgfpathlineto{\pgfqpoint{1.278801in}{0.981371in}}%
\pgfpathlineto{\pgfqpoint{1.262713in}{0.983505in}}%
\pgfpathlineto{\pgfqpoint{1.255003in}{0.967721in}}%
\pgfpathlineto{\pgfqpoint{1.247279in}{0.952054in}}%
\pgfpathclose%
\pgfusepath{fill}%
\end{pgfscope}%
\begin{pgfscope}%
\pgfpathrectangle{\pgfqpoint{0.407592in}{0.222904in}}{\pgfqpoint{1.411163in}{1.411163in}}%
\pgfusepath{clip}%
\pgfsetbuttcap%
\pgfsetroundjoin%
\definecolor{currentfill}{rgb}{0.556863,0.654902,0.847059}%
\pgfsetfillcolor{currentfill}%
\pgfsetfillopacity{0.130000}%
\pgfsetlinewidth{0.000000pt}%
\definecolor{currentstroke}{rgb}{0.000000,0.000000,0.000000}%
\pgfsetstrokecolor{currentstroke}%
\pgfsetdash{}{0pt}%
\pgfpathmoveto{\pgfqpoint{1.009537in}{0.903208in}}%
\pgfpathlineto{\pgfqpoint{1.025415in}{0.888506in}}%
\pgfpathlineto{\pgfqpoint{1.041308in}{0.874601in}}%
\pgfpathlineto{\pgfqpoint{1.049258in}{0.885864in}}%
\pgfpathlineto{\pgfqpoint{1.057187in}{0.897660in}}%
\pgfpathlineto{\pgfqpoint{1.041323in}{0.910928in}}%
\pgfpathlineto{\pgfqpoint{1.025476in}{0.925000in}}%
\pgfpathlineto{\pgfqpoint{1.017518in}{0.913846in}}%
\pgfpathlineto{\pgfqpoint{1.009537in}{0.903208in}}%
\pgfpathclose%
\pgfusepath{fill}%
\end{pgfscope}%
\begin{pgfscope}%
\pgfpathrectangle{\pgfqpoint{0.407592in}{0.222904in}}{\pgfqpoint{1.411163in}{1.411163in}}%
\pgfusepath{clip}%
\pgfsetbuttcap%
\pgfsetroundjoin%
\definecolor{currentfill}{rgb}{0.556863,0.654902,0.847059}%
\pgfsetfillcolor{currentfill}%
\pgfsetfillopacity{0.130000}%
\pgfsetlinewidth{0.000000pt}%
\definecolor{currentstroke}{rgb}{0.000000,0.000000,0.000000}%
\pgfsetstrokecolor{currentstroke}%
\pgfsetdash{}{0pt}%
\pgfpathmoveto{\pgfqpoint{1.358008in}{1.043360in}}%
\pgfpathlineto{\pgfqpoint{1.374289in}{1.044317in}}%
\pgfpathlineto{\pgfqpoint{1.390629in}{1.046024in}}%
\pgfpathlineto{\pgfqpoint{1.398245in}{1.061646in}}%
\pgfpathlineto{\pgfqpoint{1.405848in}{1.077353in}}%
\pgfpathlineto{\pgfqpoint{1.389519in}{1.075932in}}%
\pgfpathlineto{\pgfqpoint{1.373249in}{1.075222in}}%
\pgfpathlineto{\pgfqpoint{1.365634in}{1.059257in}}%
\pgfpathlineto{\pgfqpoint{1.358008in}{1.043360in}}%
\pgfpathclose%
\pgfusepath{fill}%
\end{pgfscope}%
\begin{pgfscope}%
\pgfpathrectangle{\pgfqpoint{0.407592in}{0.222904in}}{\pgfqpoint{1.411163in}{1.411163in}}%
\pgfusepath{clip}%
\pgfsetbuttcap%
\pgfsetroundjoin%
\definecolor{currentfill}{rgb}{0.556863,0.654902,0.847059}%
\pgfsetfillcolor{currentfill}%
\pgfsetfillopacity{0.130000}%
\pgfsetlinewidth{0.000000pt}%
\definecolor{currentstroke}{rgb}{0.000000,0.000000,0.000000}%
\pgfsetstrokecolor{currentstroke}%
\pgfsetdash{}{0pt}%
\pgfpathmoveto{\pgfqpoint{0.929975in}{0.950400in}}%
\pgfpathlineto{\pgfqpoint{0.945841in}{0.933042in}}%
\pgfpathlineto{\pgfqpoint{0.961718in}{0.916081in}}%
\pgfpathlineto{\pgfqpoint{0.969783in}{0.925042in}}%
\pgfpathlineto{\pgfqpoint{0.977821in}{0.934563in}}%
\pgfpathlineto{\pgfqpoint{0.961982in}{0.951031in}}%
\pgfpathlineto{\pgfqpoint{0.946153in}{0.967927in}}%
\pgfpathlineto{\pgfqpoint{0.938078in}{0.958905in}}%
\pgfpathlineto{\pgfqpoint{0.929975in}{0.950400in}}%
\pgfpathclose%
\pgfusepath{fill}%
\end{pgfscope}%
\begin{pgfscope}%
\pgfpathrectangle{\pgfqpoint{0.407592in}{0.222904in}}{\pgfqpoint{1.411163in}{1.411163in}}%
\pgfusepath{clip}%
\pgfsetbuttcap%
\pgfsetroundjoin%
\definecolor{currentfill}{rgb}{0.556863,0.654902,0.847059}%
\pgfsetfillcolor{currentfill}%
\pgfsetfillopacity{0.130000}%
\pgfsetlinewidth{0.000000pt}%
\definecolor{currentstroke}{rgb}{0.000000,0.000000,0.000000}%
\pgfsetstrokecolor{currentstroke}%
\pgfsetdash{}{0pt}%
\pgfpathmoveto{\pgfqpoint{0.818449in}{1.046516in}}%
\pgfpathlineto{\pgfqpoint{0.834299in}{1.027494in}}%
\pgfpathlineto{\pgfqpoint{0.850158in}{1.008613in}}%
\pgfpathlineto{\pgfqpoint{0.858399in}{1.015445in}}%
\pgfpathlineto{\pgfqpoint{0.866609in}{1.022743in}}%
\pgfpathlineto{\pgfqpoint{0.850791in}{1.041362in}}%
\pgfpathlineto{\pgfqpoint{0.834983in}{1.060138in}}%
\pgfpathlineto{\pgfqpoint{0.826731in}{1.053113in}}%
\pgfpathlineto{\pgfqpoint{0.818449in}{1.046516in}}%
\pgfpathclose%
\pgfusepath{fill}%
\end{pgfscope}%
\begin{pgfscope}%
\pgfpathrectangle{\pgfqpoint{0.407592in}{0.222904in}}{\pgfqpoint{1.411163in}{1.411163in}}%
\pgfusepath{clip}%
\pgfsetbuttcap%
\pgfsetroundjoin%
\definecolor{currentfill}{rgb}{0.556863,0.654902,0.847059}%
\pgfsetfillcolor{currentfill}%
\pgfsetfillopacity{0.130000}%
\pgfsetlinewidth{0.000000pt}%
\definecolor{currentstroke}{rgb}{0.000000,0.000000,0.000000}%
\pgfsetstrokecolor{currentstroke}%
\pgfsetdash{}{0pt}%
\pgfpathmoveto{\pgfqpoint{0.674969in}{1.191205in}}%
\pgfpathlineto{\pgfqpoint{0.690791in}{1.171432in}}%
\pgfpathlineto{\pgfqpoint{0.706621in}{1.151702in}}%
\pgfpathlineto{\pgfqpoint{0.715095in}{1.157033in}}%
\pgfpathlineto{\pgfqpoint{0.723535in}{1.162712in}}%
\pgfpathlineto{\pgfqpoint{0.707747in}{1.182317in}}%
\pgfpathlineto{\pgfqpoint{0.691968in}{1.201968in}}%
\pgfpathlineto{\pgfqpoint{0.683485in}{1.196424in}}%
\pgfpathlineto{\pgfqpoint{0.674969in}{1.191205in}}%
\pgfpathclose%
\pgfusepath{fill}%
\end{pgfscope}%
\begin{pgfscope}%
\pgfpathrectangle{\pgfqpoint{0.407592in}{0.222904in}}{\pgfqpoint{1.411163in}{1.411163in}}%
\pgfusepath{clip}%
\pgfsetbuttcap%
\pgfsetroundjoin%
\definecolor{currentfill}{rgb}{0.556863,0.654902,0.847059}%
\pgfsetfillcolor{currentfill}%
\pgfsetfillopacity{0.130000}%
\pgfsetlinewidth{0.000000pt}%
\definecolor{currentstroke}{rgb}{0.000000,0.000000,0.000000}%
\pgfsetstrokecolor{currentstroke}%
\pgfsetdash{}{0pt}%
\pgfpathmoveto{\pgfqpoint{1.136564in}{0.882163in}}%
\pgfpathlineto{\pgfqpoint{1.152530in}{0.875031in}}%
\pgfpathlineto{\pgfqpoint{1.168532in}{0.869248in}}%
\pgfpathlineto{\pgfqpoint{1.176350in}{0.883656in}}%
\pgfpathlineto{\pgfqpoint{1.184153in}{0.898356in}}%
\pgfpathlineto{\pgfqpoint{1.168165in}{0.903937in}}%
\pgfpathlineto{\pgfqpoint{1.152214in}{0.910754in}}%
\pgfpathlineto{\pgfqpoint{1.144397in}{0.896307in}}%
\pgfpathlineto{\pgfqpoint{1.136564in}{0.882163in}}%
\pgfpathclose%
\pgfusepath{fill}%
\end{pgfscope}%
\begin{pgfscope}%
\pgfpathrectangle{\pgfqpoint{0.407592in}{0.222904in}}{\pgfqpoint{1.411163in}{1.411163in}}%
\pgfusepath{clip}%
\pgfsetbuttcap%
\pgfsetroundjoin%
\definecolor{currentfill}{rgb}{0.556863,0.654902,0.847059}%
\pgfsetfillcolor{currentfill}%
\pgfsetfillopacity{0.130000}%
\pgfsetlinewidth{0.000000pt}%
\definecolor{currentstroke}{rgb}{0.000000,0.000000,0.000000}%
\pgfsetstrokecolor{currentstroke}%
\pgfsetdash{}{0pt}%
\pgfpathmoveto{\pgfqpoint{1.088972in}{0.873997in}}%
\pgfpathlineto{\pgfqpoint{1.104899in}{0.863848in}}%
\pgfpathlineto{\pgfqpoint{1.120852in}{0.854983in}}%
\pgfpathlineto{\pgfqpoint{1.128716in}{0.868370in}}%
\pgfpathlineto{\pgfqpoint{1.136564in}{0.882163in}}%
\pgfpathlineto{\pgfqpoint{1.120630in}{0.890580in}}%
\pgfpathlineto{\pgfqpoint{1.104724in}{0.900197in}}%
\pgfpathlineto{\pgfqpoint{1.096857in}{0.886887in}}%
\pgfpathlineto{\pgfqpoint{1.088972in}{0.873997in}}%
\pgfpathclose%
\pgfusepath{fill}%
\end{pgfscope}%
\begin{pgfscope}%
\pgfpathrectangle{\pgfqpoint{0.407592in}{0.222904in}}{\pgfqpoint{1.411163in}{1.411163in}}%
\pgfusepath{clip}%
\pgfsetbuttcap%
\pgfsetroundjoin%
\definecolor{currentfill}{rgb}{0.556863,0.654902,0.847059}%
\pgfsetfillcolor{currentfill}%
\pgfsetfillopacity{0.130000}%
\pgfsetlinewidth{0.000000pt}%
\definecolor{currentstroke}{rgb}{0.000000,0.000000,0.000000}%
\pgfsetstrokecolor{currentstroke}%
\pgfsetdash{}{0pt}%
\pgfpathmoveto{\pgfqpoint{1.294939in}{0.980202in}}%
\pgfpathlineto{\pgfqpoint{1.311130in}{0.979977in}}%
\pgfpathlineto{\pgfqpoint{1.327377in}{0.980663in}}%
\pgfpathlineto{\pgfqpoint{1.335053in}{0.996178in}}%
\pgfpathlineto{\pgfqpoint{1.342717in}{1.011809in}}%
\pgfpathlineto{\pgfqpoint{1.326481in}{1.011393in}}%
\pgfpathlineto{\pgfqpoint{1.310301in}{1.011832in}}%
\pgfpathlineto{\pgfqpoint{1.302626in}{0.995968in}}%
\pgfpathlineto{\pgfqpoint{1.294939in}{0.980202in}}%
\pgfpathclose%
\pgfusepath{fill}%
\end{pgfscope}%
\begin{pgfscope}%
\pgfpathrectangle{\pgfqpoint{0.407592in}{0.222904in}}{\pgfqpoint{1.411163in}{1.411163in}}%
\pgfusepath{clip}%
\pgfsetbuttcap%
\pgfsetroundjoin%
\definecolor{currentfill}{rgb}{0.556863,0.654902,0.847059}%
\pgfsetfillcolor{currentfill}%
\pgfsetfillopacity{0.130000}%
\pgfsetlinewidth{0.000000pt}%
\definecolor{currentstroke}{rgb}{0.000000,0.000000,0.000000}%
\pgfsetstrokecolor{currentstroke}%
\pgfsetdash{}{0pt}%
\pgfpathmoveto{\pgfqpoint{1.184153in}{0.898356in}}%
\pgfpathlineto{\pgfqpoint{1.200181in}{0.894041in}}%
\pgfpathlineto{\pgfqpoint{1.216253in}{0.891004in}}%
\pgfpathlineto{\pgfqpoint{1.224030in}{0.905972in}}%
\pgfpathlineto{\pgfqpoint{1.231793in}{0.921155in}}%
\pgfpathlineto{\pgfqpoint{1.215732in}{0.924229in}}%
\pgfpathlineto{\pgfqpoint{1.199716in}{0.928468in}}%
\pgfpathlineto{\pgfqpoint{1.191941in}{0.913305in}}%
\pgfpathlineto{\pgfqpoint{1.184153in}{0.898356in}}%
\pgfpathclose%
\pgfusepath{fill}%
\end{pgfscope}%
\begin{pgfscope}%
\pgfpathrectangle{\pgfqpoint{0.407592in}{0.222904in}}{\pgfqpoint{1.411163in}{1.411163in}}%
\pgfusepath{clip}%
\pgfsetbuttcap%
\pgfsetroundjoin%
\definecolor{currentfill}{rgb}{0.556863,0.654902,0.847059}%
\pgfsetfillcolor{currentfill}%
\pgfsetfillopacity{0.130000}%
\pgfsetlinewidth{0.000000pt}%
\definecolor{currentstroke}{rgb}{0.000000,0.000000,0.000000}%
\pgfsetstrokecolor{currentstroke}%
\pgfsetdash{}{0pt}%
\pgfpathmoveto{\pgfqpoint{0.850158in}{1.008613in}}%
\pgfpathlineto{\pgfqpoint{0.866026in}{0.989897in}}%
\pgfpathlineto{\pgfqpoint{0.881903in}{0.971372in}}%
\pgfpathlineto{\pgfqpoint{0.890103in}{0.978469in}}%
\pgfpathlineto{\pgfqpoint{0.898273in}{0.986078in}}%
\pgfpathlineto{\pgfqpoint{0.882436in}{1.004305in}}%
\pgfpathlineto{\pgfqpoint{0.866609in}{1.022743in}}%
\pgfpathlineto{\pgfqpoint{0.858399in}{1.015445in}}%
\pgfpathlineto{\pgfqpoint{0.850158in}{1.008613in}}%
\pgfpathclose%
\pgfusepath{fill}%
\end{pgfscope}%
\begin{pgfscope}%
\pgfpathrectangle{\pgfqpoint{0.407592in}{0.222904in}}{\pgfqpoint{1.411163in}{1.411163in}}%
\pgfusepath{clip}%
\pgfsetbuttcap%
\pgfsetroundjoin%
\definecolor{currentfill}{rgb}{0.556863,0.654902,0.847059}%
\pgfsetfillcolor{currentfill}%
\pgfsetfillopacity{0.130000}%
\pgfsetlinewidth{0.000000pt}%
\definecolor{currentstroke}{rgb}{0.000000,0.000000,0.000000}%
\pgfsetstrokecolor{currentstroke}%
\pgfsetdash{}{0pt}%
\pgfpathmoveto{\pgfqpoint{0.706621in}{1.151702in}}%
\pgfpathlineto{\pgfqpoint{0.722461in}{1.132021in}}%
\pgfpathlineto{\pgfqpoint{0.738309in}{1.112396in}}%
\pgfpathlineto{\pgfqpoint{0.746740in}{1.117846in}}%
\pgfpathlineto{\pgfqpoint{0.755138in}{1.123671in}}%
\pgfpathlineto{\pgfqpoint{0.739332in}{1.143161in}}%
\pgfpathlineto{\pgfqpoint{0.723535in}{1.162712in}}%
\pgfpathlineto{\pgfqpoint{0.715095in}{1.157033in}}%
\pgfpathlineto{\pgfqpoint{0.706621in}{1.151702in}}%
\pgfpathclose%
\pgfusepath{fill}%
\end{pgfscope}%
\begin{pgfscope}%
\pgfpathrectangle{\pgfqpoint{0.407592in}{0.222904in}}{\pgfqpoint{1.411163in}{1.411163in}}%
\pgfusepath{clip}%
\pgfsetbuttcap%
\pgfsetroundjoin%
\definecolor{currentfill}{rgb}{0.556863,0.654902,0.847059}%
\pgfsetfillcolor{currentfill}%
\pgfsetfillopacity{0.130000}%
\pgfsetlinewidth{0.000000pt}%
\definecolor{currentstroke}{rgb}{0.000000,0.000000,0.000000}%
\pgfsetstrokecolor{currentstroke}%
\pgfsetdash{}{0pt}%
\pgfpathmoveto{\pgfqpoint{0.961718in}{0.916081in}}%
\pgfpathlineto{\pgfqpoint{0.977606in}{0.899591in}}%
\pgfpathlineto{\pgfqpoint{0.993505in}{0.883658in}}%
\pgfpathlineto{\pgfqpoint{1.001533in}{0.893131in}}%
\pgfpathlineto{\pgfqpoint{1.009537in}{0.903208in}}%
\pgfpathlineto{\pgfqpoint{0.993673in}{0.918594in}}%
\pgfpathlineto{\pgfqpoint{0.977821in}{0.934563in}}%
\pgfpathlineto{\pgfqpoint{0.969783in}{0.925042in}}%
\pgfpathlineto{\pgfqpoint{0.961718in}{0.916081in}}%
\pgfpathclose%
\pgfusepath{fill}%
\end{pgfscope}%
\begin{pgfscope}%
\pgfpathrectangle{\pgfqpoint{0.407592in}{0.222904in}}{\pgfqpoint{1.411163in}{1.411163in}}%
\pgfusepath{clip}%
\pgfsetbuttcap%
\pgfsetroundjoin%
\definecolor{currentfill}{rgb}{0.556863,0.654902,0.847059}%
\pgfsetfillcolor{currentfill}%
\pgfsetfillopacity{0.130000}%
\pgfsetlinewidth{0.000000pt}%
\definecolor{currentstroke}{rgb}{0.000000,0.000000,0.000000}%
\pgfsetstrokecolor{currentstroke}%
\pgfsetdash{}{0pt}%
\pgfpathmoveto{\pgfqpoint{1.041308in}{0.874601in}}%
\pgfpathlineto{\pgfqpoint{1.057218in}{0.861622in}}%
\pgfpathlineto{\pgfqpoint{1.073146in}{0.849711in}}%
\pgfpathlineto{\pgfqpoint{1.081069in}{0.861584in}}%
\pgfpathlineto{\pgfqpoint{1.088972in}{0.873997in}}%
\pgfpathlineto{\pgfqpoint{1.073069in}{0.885309in}}%
\pgfpathlineto{\pgfqpoint{1.057187in}{0.897660in}}%
\pgfpathlineto{\pgfqpoint{1.049258in}{0.885864in}}%
\pgfpathlineto{\pgfqpoint{1.041308in}{0.874601in}}%
\pgfpathclose%
\pgfusepath{fill}%
\end{pgfscope}%
\begin{pgfscope}%
\pgfpathrectangle{\pgfqpoint{0.407592in}{0.222904in}}{\pgfqpoint{1.411163in}{1.411163in}}%
\pgfusepath{clip}%
\pgfsetbuttcap%
\pgfsetroundjoin%
\definecolor{currentfill}{rgb}{0.556863,0.654902,0.847059}%
\pgfsetfillcolor{currentfill}%
\pgfsetfillopacity{0.130000}%
\pgfsetlinewidth{0.000000pt}%
\definecolor{currentstroke}{rgb}{0.000000,0.000000,0.000000}%
\pgfsetstrokecolor{currentstroke}%
\pgfsetdash{}{0pt}%
\pgfpathmoveto{\pgfqpoint{1.231793in}{0.921155in}}%
\pgfpathlineto{\pgfqpoint{1.247902in}{0.919239in}}%
\pgfpathlineto{\pgfqpoint{1.264061in}{0.918456in}}%
\pgfpathlineto{\pgfqpoint{1.271800in}{0.933655in}}%
\pgfpathlineto{\pgfqpoint{1.279526in}{0.949027in}}%
\pgfpathlineto{\pgfqpoint{1.263377in}{0.950018in}}%
\pgfpathlineto{\pgfqpoint{1.247279in}{0.952054in}}%
\pgfpathlineto{\pgfqpoint{1.239543in}{0.936525in}}%
\pgfpathlineto{\pgfqpoint{1.231793in}{0.921155in}}%
\pgfpathclose%
\pgfusepath{fill}%
\end{pgfscope}%
\begin{pgfscope}%
\pgfpathrectangle{\pgfqpoint{0.407592in}{0.222904in}}{\pgfqpoint{1.411163in}{1.411163in}}%
\pgfusepath{clip}%
\pgfsetbuttcap%
\pgfsetroundjoin%
\definecolor{currentfill}{rgb}{0.556863,0.654902,0.847059}%
\pgfsetfillcolor{currentfill}%
\pgfsetfillopacity{0.130000}%
\pgfsetlinewidth{0.000000pt}%
\definecolor{currentstroke}{rgb}{0.000000,0.000000,0.000000}%
\pgfsetstrokecolor{currentstroke}%
\pgfsetdash{}{0pt}%
\pgfpathmoveto{\pgfqpoint{1.405848in}{1.077353in}}%
\pgfpathlineto{\pgfqpoint{1.422236in}{1.079459in}}%
\pgfpathlineto{\pgfqpoint{1.438685in}{1.082223in}}%
\pgfpathlineto{\pgfqpoint{1.446264in}{1.097699in}}%
\pgfpathlineto{\pgfqpoint{1.429821in}{1.095095in}}%
\pgfpathlineto{\pgfqpoint{1.413439in}{1.093134in}}%
\pgfusepath{fill}%
\end{pgfscope}%
\begin{pgfscope}%
\pgfpathrectangle{\pgfqpoint{0.407592in}{0.222904in}}{\pgfqpoint{1.411163in}{1.411163in}}%
\pgfusepath{clip}%
\pgfsetbuttcap%
\pgfsetroundjoin%
\definecolor{currentfill}{rgb}{0.556863,0.654902,0.847059}%
\pgfsetfillcolor{currentfill}%
\pgfsetfillopacity{0.130000}%
\pgfsetlinewidth{0.000000pt}%
\definecolor{currentstroke}{rgb}{0.000000,0.000000,0.000000}%
\pgfsetstrokecolor{currentstroke}%
\pgfsetdash{}{0pt}%
\pgfpathmoveto{\pgfqpoint{0.738309in}{1.112396in}}%
\pgfpathlineto{\pgfqpoint{0.754166in}{1.092835in}}%
\pgfpathlineto{\pgfqpoint{0.770031in}{1.073348in}}%
\pgfpathlineto{\pgfqpoint{0.778420in}{1.078925in}}%
\pgfpathlineto{\pgfqpoint{0.786776in}{1.084910in}}%
\pgfpathlineto{\pgfqpoint{0.770952in}{1.104251in}}%
\pgfpathlineto{\pgfqpoint{0.755138in}{1.123671in}}%
\pgfpathlineto{\pgfqpoint{0.746740in}{1.117846in}}%
\pgfpathlineto{\pgfqpoint{0.738309in}{1.112396in}}%
\pgfpathclose%
\pgfusepath{fill}%
\end{pgfscope}%
\begin{pgfscope}%
\pgfpathrectangle{\pgfqpoint{0.407592in}{0.222904in}}{\pgfqpoint{1.411163in}{1.411163in}}%
\pgfusepath{clip}%
\pgfsetbuttcap%
\pgfsetroundjoin%
\definecolor{currentfill}{rgb}{0.556863,0.654902,0.847059}%
\pgfsetfillcolor{currentfill}%
\pgfsetfillopacity{0.130000}%
\pgfsetlinewidth{0.000000pt}%
\definecolor{currentstroke}{rgb}{0.000000,0.000000,0.000000}%
\pgfsetstrokecolor{currentstroke}%
\pgfsetdash{}{0pt}%
\pgfpathmoveto{\pgfqpoint{1.342717in}{1.011809in}}%
\pgfpathlineto{\pgfqpoint{1.359010in}{1.013048in}}%
\pgfpathlineto{\pgfqpoint{1.375362in}{1.015076in}}%
\pgfpathlineto{\pgfqpoint{1.383002in}{1.030496in}}%
\pgfpathlineto{\pgfqpoint{1.390629in}{1.046024in}}%
\pgfpathlineto{\pgfqpoint{1.374289in}{1.044317in}}%
\pgfpathlineto{\pgfqpoint{1.358008in}{1.043360in}}%
\pgfpathlineto{\pgfqpoint{1.350369in}{1.027540in}}%
\pgfpathlineto{\pgfqpoint{1.342717in}{1.011809in}}%
\pgfpathclose%
\pgfusepath{fill}%
\end{pgfscope}%
\begin{pgfscope}%
\pgfpathrectangle{\pgfqpoint{0.407592in}{0.222904in}}{\pgfqpoint{1.411163in}{1.411163in}}%
\pgfusepath{clip}%
\pgfsetbuttcap%
\pgfsetroundjoin%
\definecolor{currentfill}{rgb}{0.556863,0.654902,0.847059}%
\pgfsetfillcolor{currentfill}%
\pgfsetfillopacity{0.130000}%
\pgfsetlinewidth{0.000000pt}%
\definecolor{currentstroke}{rgb}{0.000000,0.000000,0.000000}%
\pgfsetstrokecolor{currentstroke}%
\pgfsetdash{}{0pt}%
\pgfpathmoveto{\pgfqpoint{0.881903in}{0.971372in}}%
\pgfpathlineto{\pgfqpoint{0.897790in}{0.953072in}}%
\pgfpathlineto{\pgfqpoint{0.913686in}{0.935037in}}%
\pgfpathlineto{\pgfqpoint{0.921845in}{0.942437in}}%
\pgfpathlineto{\pgfqpoint{0.929975in}{0.950400in}}%
\pgfpathlineto{\pgfqpoint{0.914119in}{0.968095in}}%
\pgfpathlineto{\pgfqpoint{0.898273in}{0.986078in}}%
\pgfpathlineto{\pgfqpoint{0.890103in}{0.978469in}}%
\pgfpathlineto{\pgfqpoint{0.881903in}{0.971372in}}%
\pgfpathclose%
\pgfusepath{fill}%
\end{pgfscope}%
\begin{pgfscope}%
\pgfpathrectangle{\pgfqpoint{0.407592in}{0.222904in}}{\pgfqpoint{1.411163in}{1.411163in}}%
\pgfusepath{clip}%
\pgfsetbuttcap%
\pgfsetroundjoin%
\definecolor{currentfill}{rgb}{0.556863,0.654902,0.847059}%
\pgfsetfillcolor{currentfill}%
\pgfsetfillopacity{0.130000}%
\pgfsetlinewidth{0.000000pt}%
\definecolor{currentstroke}{rgb}{0.000000,0.000000,0.000000}%
\pgfsetstrokecolor{currentstroke}%
\pgfsetdash{}{0pt}%
\pgfpathmoveto{\pgfqpoint{0.993505in}{0.883658in}}%
\pgfpathlineto{\pgfqpoint{1.009416in}{0.868386in}}%
\pgfpathlineto{\pgfqpoint{1.025342in}{0.853898in}}%
\pgfpathlineto{\pgfqpoint{1.033336in}{0.863927in}}%
\pgfpathlineto{\pgfqpoint{1.041308in}{0.874601in}}%
\pgfpathlineto{\pgfqpoint{1.025415in}{0.888506in}}%
\pgfpathlineto{\pgfqpoint{1.009537in}{0.903208in}}%
\pgfpathlineto{\pgfqpoint{1.001533in}{0.893131in}}%
\pgfpathlineto{\pgfqpoint{0.993505in}{0.883658in}}%
\pgfpathclose%
\pgfusepath{fill}%
\end{pgfscope}%
\begin{pgfscope}%
\pgfpathrectangle{\pgfqpoint{0.407592in}{0.222904in}}{\pgfqpoint{1.411163in}{1.411163in}}%
\pgfusepath{clip}%
\pgfsetbuttcap%
\pgfsetroundjoin%
\definecolor{currentfill}{rgb}{0.556863,0.654902,0.847059}%
\pgfsetfillcolor{currentfill}%
\pgfsetfillopacity{0.130000}%
\pgfsetlinewidth{0.000000pt}%
\definecolor{currentstroke}{rgb}{0.000000,0.000000,0.000000}%
\pgfsetstrokecolor{currentstroke}%
\pgfsetdash{}{0pt}%
\pgfpathmoveto{\pgfqpoint{0.770031in}{1.073348in}}%
\pgfpathlineto{\pgfqpoint{0.785906in}{1.053946in}}%
\pgfpathlineto{\pgfqpoint{0.801789in}{1.034642in}}%
\pgfpathlineto{\pgfqpoint{0.810135in}{1.040357in}}%
\pgfpathlineto{\pgfqpoint{0.818449in}{1.046516in}}%
\pgfpathlineto{\pgfqpoint{0.802608in}{1.065660in}}%
\pgfpathlineto{\pgfqpoint{0.786776in}{1.084910in}}%
\pgfpathlineto{\pgfqpoint{0.778420in}{1.078925in}}%
\pgfpathlineto{\pgfqpoint{0.770031in}{1.073348in}}%
\pgfpathclose%
\pgfusepath{fill}%
\end{pgfscope}%
\begin{pgfscope}%
\pgfpathrectangle{\pgfqpoint{0.407592in}{0.222904in}}{\pgfqpoint{1.411163in}{1.411163in}}%
\pgfusepath{clip}%
\pgfsetbuttcap%
\pgfsetroundjoin%
\definecolor{currentfill}{rgb}{0.556863,0.654902,0.847059}%
\pgfsetfillcolor{currentfill}%
\pgfsetfillopacity{0.130000}%
\pgfsetlinewidth{0.000000pt}%
\definecolor{currentstroke}{rgb}{0.000000,0.000000,0.000000}%
\pgfsetstrokecolor{currentstroke}%
\pgfsetdash{}{0pt}%
\pgfpathmoveto{\pgfqpoint{1.279526in}{0.949027in}}%
\pgfpathlineto{\pgfqpoint{1.295728in}{0.949049in}}%
\pgfpathlineto{\pgfqpoint{1.311986in}{0.950046in}}%
\pgfpathlineto{\pgfqpoint{1.319688in}{0.965280in}}%
\pgfpathlineto{\pgfqpoint{1.327377in}{0.980663in}}%
\pgfpathlineto{\pgfqpoint{1.311130in}{0.979977in}}%
\pgfpathlineto{\pgfqpoint{1.294939in}{0.980202in}}%
\pgfpathlineto{\pgfqpoint{1.287239in}{0.964549in}}%
\pgfpathlineto{\pgfqpoint{1.279526in}{0.949027in}}%
\pgfpathclose%
\pgfusepath{fill}%
\end{pgfscope}%
\begin{pgfscope}%
\pgfpathrectangle{\pgfqpoint{0.407592in}{0.222904in}}{\pgfqpoint{1.411163in}{1.411163in}}%
\pgfusepath{clip}%
\pgfsetbuttcap%
\pgfsetroundjoin%
\definecolor{currentfill}{rgb}{0.556863,0.654902,0.847059}%
\pgfsetfillcolor{currentfill}%
\pgfsetfillopacity{0.130000}%
\pgfsetlinewidth{0.000000pt}%
\definecolor{currentstroke}{rgb}{0.000000,0.000000,0.000000}%
\pgfsetstrokecolor{currentstroke}%
\pgfsetdash{}{0pt}%
\pgfpathmoveto{\pgfqpoint{1.120852in}{0.854983in}}%
\pgfpathlineto{\pgfqpoint{1.136835in}{0.847507in}}%
\pgfpathlineto{\pgfqpoint{1.152852in}{0.841507in}}%
\pgfpathlineto{\pgfqpoint{1.160700in}{0.855179in}}%
\pgfpathlineto{\pgfqpoint{1.168532in}{0.869248in}}%
\pgfpathlineto{\pgfqpoint{1.152530in}{0.875031in}}%
\pgfpathlineto{\pgfqpoint{1.136564in}{0.882163in}}%
\pgfpathlineto{\pgfqpoint{1.128716in}{0.868370in}}%
\pgfpathlineto{\pgfqpoint{1.120852in}{0.854983in}}%
\pgfpathclose%
\pgfusepath{fill}%
\end{pgfscope}%
\begin{pgfscope}%
\pgfpathrectangle{\pgfqpoint{0.407592in}{0.222904in}}{\pgfqpoint{1.411163in}{1.411163in}}%
\pgfusepath{clip}%
\pgfsetbuttcap%
\pgfsetroundjoin%
\definecolor{currentfill}{rgb}{0.556863,0.654902,0.847059}%
\pgfsetfillcolor{currentfill}%
\pgfsetfillopacity{0.130000}%
\pgfsetlinewidth{0.000000pt}%
\definecolor{currentstroke}{rgb}{0.000000,0.000000,0.000000}%
\pgfsetstrokecolor{currentstroke}%
\pgfsetdash{}{0pt}%
\pgfpathmoveto{\pgfqpoint{0.913686in}{0.935037in}}%
\pgfpathlineto{\pgfqpoint{0.929591in}{0.917316in}}%
\pgfpathlineto{\pgfqpoint{0.945507in}{0.899967in}}%
\pgfpathlineto{\pgfqpoint{0.953626in}{0.907713in}}%
\pgfpathlineto{\pgfqpoint{0.961718in}{0.916081in}}%
\pgfpathlineto{\pgfqpoint{0.945841in}{0.933042in}}%
\pgfpathlineto{\pgfqpoint{0.929975in}{0.950400in}}%
\pgfpathlineto{\pgfqpoint{0.921845in}{0.942437in}}%
\pgfpathlineto{\pgfqpoint{0.913686in}{0.935037in}}%
\pgfpathclose%
\pgfusepath{fill}%
\end{pgfscope}%
\begin{pgfscope}%
\pgfpathrectangle{\pgfqpoint{0.407592in}{0.222904in}}{\pgfqpoint{1.411163in}{1.411163in}}%
\pgfusepath{clip}%
\pgfsetbuttcap%
\pgfsetroundjoin%
\definecolor{currentfill}{rgb}{0.556863,0.654902,0.847059}%
\pgfsetfillcolor{currentfill}%
\pgfsetfillopacity{0.130000}%
\pgfsetlinewidth{0.000000pt}%
\definecolor{currentstroke}{rgb}{0.000000,0.000000,0.000000}%
\pgfsetstrokecolor{currentstroke}%
\pgfsetdash{}{0pt}%
\pgfpathmoveto{\pgfqpoint{1.073146in}{0.849711in}}%
\pgfpathlineto{\pgfqpoint{1.089096in}{0.839014in}}%
\pgfpathlineto{\pgfqpoint{1.105070in}{0.829681in}}%
\pgfpathlineto{\pgfqpoint{1.112970in}{0.842063in}}%
\pgfpathlineto{\pgfqpoint{1.120852in}{0.854983in}}%
\pgfpathlineto{\pgfqpoint{1.104899in}{0.863848in}}%
\pgfpathlineto{\pgfqpoint{1.088972in}{0.873997in}}%
\pgfpathlineto{\pgfqpoint{1.081069in}{0.861584in}}%
\pgfpathlineto{\pgfqpoint{1.073146in}{0.849711in}}%
\pgfpathclose%
\pgfusepath{fill}%
\end{pgfscope}%
\begin{pgfscope}%
\pgfpathrectangle{\pgfqpoint{0.407592in}{0.222904in}}{\pgfqpoint{1.411163in}{1.411163in}}%
\pgfusepath{clip}%
\pgfsetbuttcap%
\pgfsetroundjoin%
\definecolor{currentfill}{rgb}{0.556863,0.654902,0.847059}%
\pgfsetfillcolor{currentfill}%
\pgfsetfillopacity{0.130000}%
\pgfsetlinewidth{0.000000pt}%
\definecolor{currentstroke}{rgb}{0.000000,0.000000,0.000000}%
\pgfsetstrokecolor{currentstroke}%
\pgfsetdash{}{0pt}%
\pgfpathmoveto{\pgfqpoint{1.168532in}{0.869248in}}%
\pgfpathlineto{\pgfqpoint{1.184574in}{0.864856in}}%
\pgfpathlineto{\pgfqpoint{1.200658in}{0.861867in}}%
\pgfpathlineto{\pgfqpoint{1.208463in}{0.876289in}}%
\pgfpathlineto{\pgfqpoint{1.216253in}{0.891004in}}%
\pgfpathlineto{\pgfqpoint{1.200181in}{0.894041in}}%
\pgfpathlineto{\pgfqpoint{1.184153in}{0.898356in}}%
\pgfpathlineto{\pgfqpoint{1.176350in}{0.883656in}}%
\pgfpathlineto{\pgfqpoint{1.168532in}{0.869248in}}%
\pgfpathclose%
\pgfusepath{fill}%
\end{pgfscope}%
\begin{pgfscope}%
\pgfpathrectangle{\pgfqpoint{0.407592in}{0.222904in}}{\pgfqpoint{1.411163in}{1.411163in}}%
\pgfusepath{clip}%
\pgfsetbuttcap%
\pgfsetroundjoin%
\definecolor{currentfill}{rgb}{0.556863,0.654902,0.847059}%
\pgfsetfillcolor{currentfill}%
\pgfsetfillopacity{0.130000}%
\pgfsetlinewidth{0.000000pt}%
\definecolor{currentstroke}{rgb}{0.000000,0.000000,0.000000}%
\pgfsetstrokecolor{currentstroke}%
\pgfsetdash{}{0pt}%
\pgfpathmoveto{\pgfqpoint{1.390629in}{1.046024in}}%
\pgfpathlineto{\pgfqpoint{1.407029in}{1.048448in}}%
\pgfpathlineto{\pgfqpoint{1.423491in}{1.051556in}}%
\pgfpathlineto{\pgfqpoint{1.431094in}{1.066839in}}%
\pgfpathlineto{\pgfqpoint{1.438685in}{1.082223in}}%
\pgfpathlineto{\pgfqpoint{1.422236in}{1.079459in}}%
\pgfpathlineto{\pgfqpoint{1.405848in}{1.077353in}}%
\pgfpathlineto{\pgfqpoint{1.398245in}{1.061646in}}%
\pgfpathlineto{\pgfqpoint{1.390629in}{1.046024in}}%
\pgfpathclose%
\pgfusepath{fill}%
\end{pgfscope}%
\begin{pgfscope}%
\pgfpathrectangle{\pgfqpoint{0.407592in}{0.222904in}}{\pgfqpoint{1.411163in}{1.411163in}}%
\pgfusepath{clip}%
\pgfsetbuttcap%
\pgfsetroundjoin%
\definecolor{currentfill}{rgb}{0.556863,0.654902,0.847059}%
\pgfsetfillcolor{currentfill}%
\pgfsetfillopacity{0.130000}%
\pgfsetlinewidth{0.000000pt}%
\definecolor{currentstroke}{rgb}{0.000000,0.000000,0.000000}%
\pgfsetstrokecolor{currentstroke}%
\pgfsetdash{}{0pt}%
\pgfpathmoveto{\pgfqpoint{0.801789in}{1.034642in}}%
\pgfpathlineto{\pgfqpoint{0.817681in}{1.015451in}}%
\pgfpathlineto{\pgfqpoint{0.833582in}{0.996393in}}%
\pgfpathlineto{\pgfqpoint{0.841886in}{1.002260in}}%
\pgfpathlineto{\pgfqpoint{0.850158in}{1.008613in}}%
\pgfpathlineto{\pgfqpoint{0.834299in}{1.027494in}}%
\pgfpathlineto{\pgfqpoint{0.818449in}{1.046516in}}%
\pgfpathlineto{\pgfqpoint{0.810135in}{1.040357in}}%
\pgfpathlineto{\pgfqpoint{0.801789in}{1.034642in}}%
\pgfpathclose%
\pgfusepath{fill}%
\end{pgfscope}%
\begin{pgfscope}%
\pgfpathrectangle{\pgfqpoint{0.407592in}{0.222904in}}{\pgfqpoint{1.411163in}{1.411163in}}%
\pgfusepath{clip}%
\pgfsetbuttcap%
\pgfsetroundjoin%
\definecolor{currentfill}{rgb}{0.556863,0.654902,0.847059}%
\pgfsetfillcolor{currentfill}%
\pgfsetfillopacity{0.130000}%
\pgfsetlinewidth{0.000000pt}%
\definecolor{currentstroke}{rgb}{0.000000,0.000000,0.000000}%
\pgfsetstrokecolor{currentstroke}%
\pgfsetdash{}{0pt}%
\pgfpathmoveto{\pgfqpoint{0.657834in}{1.181752in}}%
\pgfpathlineto{\pgfqpoint{0.673699in}{1.161904in}}%
\pgfpathlineto{\pgfqpoint{0.689572in}{1.142098in}}%
\pgfpathlineto{\pgfqpoint{0.698114in}{1.146723in}}%
\pgfpathlineto{\pgfqpoint{0.706621in}{1.151702in}}%
\pgfpathlineto{\pgfqpoint{0.690791in}{1.171432in}}%
\pgfpathlineto{\pgfqpoint{0.674969in}{1.191205in}}%
\pgfpathlineto{\pgfqpoint{0.666419in}{1.186314in}}%
\pgfpathlineto{\pgfqpoint{0.657834in}{1.181752in}}%
\pgfpathclose%
\pgfusepath{fill}%
\end{pgfscope}%
\begin{pgfscope}%
\pgfpathrectangle{\pgfqpoint{0.407592in}{0.222904in}}{\pgfqpoint{1.411163in}{1.411163in}}%
\pgfusepath{clip}%
\pgfsetbuttcap%
\pgfsetroundjoin%
\definecolor{currentfill}{rgb}{0.556863,0.654902,0.847059}%
\pgfsetfillcolor{currentfill}%
\pgfsetfillopacity{0.130000}%
\pgfsetlinewidth{0.000000pt}%
\definecolor{currentstroke}{rgb}{0.000000,0.000000,0.000000}%
\pgfsetstrokecolor{currentstroke}%
\pgfsetdash{}{0pt}%
\pgfpathmoveto{\pgfqpoint{1.025342in}{0.853898in}}%
\pgfpathlineto{\pgfqpoint{1.041282in}{0.840334in}}%
\pgfpathlineto{\pgfqpoint{1.057239in}{0.827855in}}%
\pgfpathlineto{\pgfqpoint{1.065203in}{0.838444in}}%
\pgfpathlineto{\pgfqpoint{1.073146in}{0.849711in}}%
\pgfpathlineto{\pgfqpoint{1.057218in}{0.861622in}}%
\pgfpathlineto{\pgfqpoint{1.041308in}{0.874601in}}%
\pgfpathlineto{\pgfqpoint{1.033336in}{0.863927in}}%
\pgfpathlineto{\pgfqpoint{1.025342in}{0.853898in}}%
\pgfpathclose%
\pgfusepath{fill}%
\end{pgfscope}%
\begin{pgfscope}%
\pgfpathrectangle{\pgfqpoint{0.407592in}{0.222904in}}{\pgfqpoint{1.411163in}{1.411163in}}%
\pgfusepath{clip}%
\pgfsetbuttcap%
\pgfsetroundjoin%
\definecolor{currentfill}{rgb}{0.556863,0.654902,0.847059}%
\pgfsetfillcolor{currentfill}%
\pgfsetfillopacity{0.130000}%
\pgfsetlinewidth{0.000000pt}%
\definecolor{currentstroke}{rgb}{0.000000,0.000000,0.000000}%
\pgfsetstrokecolor{currentstroke}%
\pgfsetdash{}{0pt}%
\pgfpathmoveto{\pgfqpoint{1.216253in}{0.891004in}}%
\pgfpathlineto{\pgfqpoint{1.232373in}{0.889232in}}%
\pgfpathlineto{\pgfqpoint{1.248544in}{0.888691in}}%
\pgfpathlineto{\pgfqpoint{1.256310in}{0.903458in}}%
\pgfpathlineto{\pgfqpoint{1.264061in}{0.918456in}}%
\pgfpathlineto{\pgfqpoint{1.247902in}{0.919239in}}%
\pgfpathlineto{\pgfqpoint{1.231793in}{0.921155in}}%
\pgfpathlineto{\pgfqpoint{1.224030in}{0.905972in}}%
\pgfpathlineto{\pgfqpoint{1.216253in}{0.891004in}}%
\pgfpathclose%
\pgfusepath{fill}%
\end{pgfscope}%
\begin{pgfscope}%
\pgfpathrectangle{\pgfqpoint{0.407592in}{0.222904in}}{\pgfqpoint{1.411163in}{1.411163in}}%
\pgfusepath{clip}%
\pgfsetbuttcap%
\pgfsetroundjoin%
\definecolor{currentfill}{rgb}{0.556863,0.654902,0.847059}%
\pgfsetfillcolor{currentfill}%
\pgfsetfillopacity{0.130000}%
\pgfsetlinewidth{0.000000pt}%
\definecolor{currentstroke}{rgb}{0.000000,0.000000,0.000000}%
\pgfsetstrokecolor{currentstroke}%
\pgfsetdash{}{0pt}%
\pgfpathmoveto{\pgfqpoint{0.945507in}{0.899967in}}%
\pgfpathlineto{\pgfqpoint{0.961432in}{0.883064in}}%
\pgfpathlineto{\pgfqpoint{0.977369in}{0.866693in}}%
\pgfpathlineto{\pgfqpoint{0.985450in}{0.874832in}}%
\pgfpathlineto{\pgfqpoint{0.993505in}{0.883658in}}%
\pgfpathlineto{\pgfqpoint{0.977606in}{0.899591in}}%
\pgfpathlineto{\pgfqpoint{0.961718in}{0.916081in}}%
\pgfpathlineto{\pgfqpoint{0.953626in}{0.907713in}}%
\pgfpathlineto{\pgfqpoint{0.945507in}{0.899967in}}%
\pgfpathclose%
\pgfusepath{fill}%
\end{pgfscope}%
\begin{pgfscope}%
\pgfpathrectangle{\pgfqpoint{0.407592in}{0.222904in}}{\pgfqpoint{1.411163in}{1.411163in}}%
\pgfusepath{clip}%
\pgfsetbuttcap%
\pgfsetroundjoin%
\definecolor{currentfill}{rgb}{0.556863,0.654902,0.847059}%
\pgfsetfillcolor{currentfill}%
\pgfsetfillopacity{0.130000}%
\pgfsetlinewidth{0.000000pt}%
\definecolor{currentstroke}{rgb}{0.000000,0.000000,0.000000}%
\pgfsetstrokecolor{currentstroke}%
\pgfsetdash{}{0pt}%
\pgfpathmoveto{\pgfqpoint{1.327377in}{0.980663in}}%
\pgfpathlineto{\pgfqpoint{1.343681in}{0.982223in}}%
\pgfpathlineto{\pgfqpoint{1.360045in}{0.984614in}}%
\pgfpathlineto{\pgfqpoint{1.367710in}{0.999777in}}%
\pgfpathlineto{\pgfqpoint{1.375362in}{1.015076in}}%
\pgfpathlineto{\pgfqpoint{1.359010in}{1.013048in}}%
\pgfpathlineto{\pgfqpoint{1.342717in}{1.011809in}}%
\pgfpathlineto{\pgfqpoint{1.335053in}{0.996178in}}%
\pgfpathlineto{\pgfqpoint{1.327377in}{0.980663in}}%
\pgfpathclose%
\pgfusepath{fill}%
\end{pgfscope}%
\begin{pgfscope}%
\pgfpathrectangle{\pgfqpoint{0.407592in}{0.222904in}}{\pgfqpoint{1.411163in}{1.411163in}}%
\pgfusepath{clip}%
\pgfsetbuttcap%
\pgfsetroundjoin%
\definecolor{currentfill}{rgb}{0.556863,0.654902,0.847059}%
\pgfsetfillcolor{currentfill}%
\pgfsetfillopacity{0.130000}%
\pgfsetlinewidth{0.000000pt}%
\definecolor{currentstroke}{rgb}{0.000000,0.000000,0.000000}%
\pgfsetstrokecolor{currentstroke}%
\pgfsetdash{}{0pt}%
\pgfpathmoveto{\pgfqpoint{0.833582in}{0.996393in}}%
\pgfpathlineto{\pgfqpoint{0.849492in}{0.977490in}}%
\pgfpathlineto{\pgfqpoint{0.865410in}{0.958767in}}%
\pgfpathlineto{\pgfqpoint{0.873673in}{0.964800in}}%
\pgfpathlineto{\pgfqpoint{0.881903in}{0.971372in}}%
\pgfpathlineto{\pgfqpoint{0.866026in}{0.989897in}}%
\pgfpathlineto{\pgfqpoint{0.850158in}{1.008613in}}%
\pgfpathlineto{\pgfqpoint{0.841886in}{1.002260in}}%
\pgfpathlineto{\pgfqpoint{0.833582in}{0.996393in}}%
\pgfpathclose%
\pgfusepath{fill}%
\end{pgfscope}%
\begin{pgfscope}%
\pgfpathrectangle{\pgfqpoint{0.407592in}{0.222904in}}{\pgfqpoint{1.411163in}{1.411163in}}%
\pgfusepath{clip}%
\pgfsetbuttcap%
\pgfsetroundjoin%
\definecolor{currentfill}{rgb}{0.556863,0.654902,0.847059}%
\pgfsetfillcolor{currentfill}%
\pgfsetfillopacity{0.130000}%
\pgfsetlinewidth{0.000000pt}%
\definecolor{currentstroke}{rgb}{0.000000,0.000000,0.000000}%
\pgfsetstrokecolor{currentstroke}%
\pgfsetdash{}{0pt}%
\pgfpathmoveto{\pgfqpoint{0.689572in}{1.142098in}}%
\pgfpathlineto{\pgfqpoint{0.705454in}{1.122340in}}%
\pgfpathlineto{\pgfqpoint{0.721345in}{1.102638in}}%
\pgfpathlineto{\pgfqpoint{0.729844in}{1.107326in}}%
\pgfpathlineto{\pgfqpoint{0.738309in}{1.112396in}}%
\pgfpathlineto{\pgfqpoint{0.722461in}{1.132021in}}%
\pgfpathlineto{\pgfqpoint{0.706621in}{1.151702in}}%
\pgfpathlineto{\pgfqpoint{0.698114in}{1.146723in}}%
\pgfpathlineto{\pgfqpoint{0.689572in}{1.142098in}}%
\pgfpathclose%
\pgfusepath{fill}%
\end{pgfscope}%
\begin{pgfscope}%
\pgfpathrectangle{\pgfqpoint{0.407592in}{0.222904in}}{\pgfqpoint{1.411163in}{1.411163in}}%
\pgfusepath{clip}%
\pgfsetbuttcap%
\pgfsetroundjoin%
\definecolor{currentfill}{rgb}{0.556863,0.654902,0.847059}%
\pgfsetfillcolor{currentfill}%
\pgfsetfillopacity{0.130000}%
\pgfsetlinewidth{0.000000pt}%
\definecolor{currentstroke}{rgb}{0.000000,0.000000,0.000000}%
\pgfsetstrokecolor{currentstroke}%
\pgfsetdash{}{0pt}%
\pgfpathmoveto{\pgfqpoint{1.264061in}{0.918456in}}%
\pgfpathlineto{\pgfqpoint{1.280275in}{0.918765in}}%
\pgfpathlineto{\pgfqpoint{1.296545in}{0.920117in}}%
\pgfpathlineto{\pgfqpoint{1.304272in}{0.934984in}}%
\pgfpathlineto{\pgfqpoint{1.311986in}{0.950046in}}%
\pgfpathlineto{\pgfqpoint{1.295728in}{0.949049in}}%
\pgfpathlineto{\pgfqpoint{1.279526in}{0.949027in}}%
\pgfpathlineto{\pgfqpoint{1.271800in}{0.933655in}}%
\pgfpathlineto{\pgfqpoint{1.264061in}{0.918456in}}%
\pgfpathclose%
\pgfusepath{fill}%
\end{pgfscope}%
\begin{pgfscope}%
\pgfpathrectangle{\pgfqpoint{0.407592in}{0.222904in}}{\pgfqpoint{1.411163in}{1.411163in}}%
\pgfusepath{clip}%
\pgfsetbuttcap%
\pgfsetroundjoin%
\definecolor{currentfill}{rgb}{0.556863,0.654902,0.847059}%
\pgfsetfillcolor{currentfill}%
\pgfsetfillopacity{0.130000}%
\pgfsetlinewidth{0.000000pt}%
\definecolor{currentstroke}{rgb}{0.000000,0.000000,0.000000}%
\pgfsetstrokecolor{currentstroke}%
\pgfsetdash{}{0pt}%
\pgfpathmoveto{\pgfqpoint{0.865410in}{0.958767in}}%
\pgfpathlineto{\pgfqpoint{0.881338in}{0.940256in}}%
\pgfpathlineto{\pgfqpoint{0.897275in}{0.921998in}}%
\pgfpathlineto{\pgfqpoint{0.905496in}{0.928219in}}%
\pgfpathlineto{\pgfqpoint{0.913686in}{0.935037in}}%
\pgfpathlineto{\pgfqpoint{0.897790in}{0.953072in}}%
\pgfpathlineto{\pgfqpoint{0.881903in}{0.971372in}}%
\pgfpathlineto{\pgfqpoint{0.873673in}{0.964800in}}%
\pgfpathlineto{\pgfqpoint{0.865410in}{0.958767in}}%
\pgfpathclose%
\pgfusepath{fill}%
\end{pgfscope}%
\begin{pgfscope}%
\pgfpathrectangle{\pgfqpoint{0.407592in}{0.222904in}}{\pgfqpoint{1.411163in}{1.411163in}}%
\pgfusepath{clip}%
\pgfsetbuttcap%
\pgfsetroundjoin%
\definecolor{currentfill}{rgb}{0.556863,0.654902,0.847059}%
\pgfsetfillcolor{currentfill}%
\pgfsetfillopacity{0.130000}%
\pgfsetlinewidth{0.000000pt}%
\definecolor{currentstroke}{rgb}{0.000000,0.000000,0.000000}%
\pgfsetstrokecolor{currentstroke}%
\pgfsetdash{}{0pt}%
\pgfpathmoveto{\pgfqpoint{0.721345in}{1.102638in}}%
\pgfpathlineto{\pgfqpoint{0.737244in}{1.082998in}}%
\pgfpathlineto{\pgfqpoint{0.753153in}{1.063432in}}%
\pgfpathlineto{\pgfqpoint{0.761609in}{1.068182in}}%
\pgfpathlineto{\pgfqpoint{0.770031in}{1.073348in}}%
\pgfpathlineto{\pgfqpoint{0.754166in}{1.092835in}}%
\pgfpathlineto{\pgfqpoint{0.738309in}{1.112396in}}%
\pgfpathlineto{\pgfqpoint{0.729844in}{1.107326in}}%
\pgfpathlineto{\pgfqpoint{0.721345in}{1.102638in}}%
\pgfpathclose%
\pgfusepath{fill}%
\end{pgfscope}%
\begin{pgfscope}%
\pgfpathrectangle{\pgfqpoint{0.407592in}{0.222904in}}{\pgfqpoint{1.411163in}{1.411163in}}%
\pgfusepath{clip}%
\pgfsetbuttcap%
\pgfsetroundjoin%
\definecolor{currentfill}{rgb}{0.556863,0.654902,0.847059}%
\pgfsetfillcolor{currentfill}%
\pgfsetfillopacity{0.130000}%
\pgfsetlinewidth{0.000000pt}%
\definecolor{currentstroke}{rgb}{0.000000,0.000000,0.000000}%
\pgfsetstrokecolor{currentstroke}%
\pgfsetdash{}{0pt}%
\pgfpathmoveto{\pgfqpoint{1.105070in}{0.829681in}}%
\pgfpathlineto{\pgfqpoint{1.121073in}{0.821845in}}%
\pgfpathlineto{\pgfqpoint{1.137108in}{0.815616in}}%
\pgfpathlineto{\pgfqpoint{1.144988in}{0.828294in}}%
\pgfpathlineto{\pgfqpoint{1.152852in}{0.841507in}}%
\pgfpathlineto{\pgfqpoint{1.136835in}{0.847507in}}%
\pgfpathlineto{\pgfqpoint{1.120852in}{0.854983in}}%
\pgfpathlineto{\pgfqpoint{1.112970in}{0.842063in}}%
\pgfpathlineto{\pgfqpoint{1.105070in}{0.829681in}}%
\pgfpathclose%
\pgfusepath{fill}%
\end{pgfscope}%
\begin{pgfscope}%
\pgfpathrectangle{\pgfqpoint{0.407592in}{0.222904in}}{\pgfqpoint{1.411163in}{1.411163in}}%
\pgfusepath{clip}%
\pgfsetbuttcap%
\pgfsetroundjoin%
\definecolor{currentfill}{rgb}{0.556863,0.654902,0.847059}%
\pgfsetfillcolor{currentfill}%
\pgfsetfillopacity{0.130000}%
\pgfsetlinewidth{0.000000pt}%
\definecolor{currentstroke}{rgb}{0.000000,0.000000,0.000000}%
\pgfsetstrokecolor{currentstroke}%
\pgfsetdash{}{0pt}%
\pgfpathmoveto{\pgfqpoint{0.977369in}{0.866693in}}%
\pgfpathlineto{\pgfqpoint{0.993317in}{0.850961in}}%
\pgfpathlineto{\pgfqpoint{1.009277in}{0.835997in}}%
\pgfpathlineto{\pgfqpoint{1.017322in}{0.844570in}}%
\pgfpathlineto{\pgfqpoint{1.025342in}{0.853898in}}%
\pgfpathlineto{\pgfqpoint{1.009416in}{0.868386in}}%
\pgfpathlineto{\pgfqpoint{0.993505in}{0.883658in}}%
\pgfpathlineto{\pgfqpoint{0.985450in}{0.874832in}}%
\pgfpathlineto{\pgfqpoint{0.977369in}{0.866693in}}%
\pgfpathclose%
\pgfusepath{fill}%
\end{pgfscope}%
\begin{pgfscope}%
\pgfpathrectangle{\pgfqpoint{0.407592in}{0.222904in}}{\pgfqpoint{1.411163in}{1.411163in}}%
\pgfusepath{clip}%
\pgfsetbuttcap%
\pgfsetroundjoin%
\definecolor{currentfill}{rgb}{0.556863,0.654902,0.847059}%
\pgfsetfillcolor{currentfill}%
\pgfsetfillopacity{0.130000}%
\pgfsetlinewidth{0.000000pt}%
\definecolor{currentstroke}{rgb}{0.000000,0.000000,0.000000}%
\pgfsetstrokecolor{currentstroke}%
\pgfsetdash{}{0pt}%
\pgfpathmoveto{\pgfqpoint{1.438685in}{1.082223in}}%
\pgfpathlineto{\pgfqpoint{1.455197in}{1.085615in}}%
\pgfpathlineto{\pgfqpoint{1.471773in}{1.089605in}}%
\pgfpathlineto{\pgfqpoint{1.479338in}{1.104731in}}%
\pgfpathlineto{\pgfqpoint{1.462769in}{1.100921in}}%
\pgfpathlineto{\pgfqpoint{1.446264in}{1.097699in}}%
\pgfusepath{fill}%
\end{pgfscope}%
\begin{pgfscope}%
\pgfpathrectangle{\pgfqpoint{0.407592in}{0.222904in}}{\pgfqpoint{1.411163in}{1.411163in}}%
\pgfusepath{clip}%
\pgfsetbuttcap%
\pgfsetroundjoin%
\definecolor{currentfill}{rgb}{0.556863,0.654902,0.847059}%
\pgfsetfillcolor{currentfill}%
\pgfsetfillopacity{0.130000}%
\pgfsetlinewidth{0.000000pt}%
\definecolor{currentstroke}{rgb}{0.000000,0.000000,0.000000}%
\pgfsetstrokecolor{currentstroke}%
\pgfsetdash{}{0pt}%
\pgfpathmoveto{\pgfqpoint{1.375362in}{1.015076in}}%
\pgfpathlineto{\pgfqpoint{1.391774in}{1.017856in}}%
\pgfpathlineto{\pgfqpoint{1.408250in}{1.021348in}}%
\pgfpathlineto{\pgfqpoint{1.415877in}{1.036388in}}%
\pgfpathlineto{\pgfqpoint{1.423491in}{1.051556in}}%
\pgfpathlineto{\pgfqpoint{1.407029in}{1.048448in}}%
\pgfpathlineto{\pgfqpoint{1.390629in}{1.046024in}}%
\pgfpathlineto{\pgfqpoint{1.383002in}{1.030496in}}%
\pgfpathlineto{\pgfqpoint{1.375362in}{1.015076in}}%
\pgfpathclose%
\pgfusepath{fill}%
\end{pgfscope}%
\begin{pgfscope}%
\pgfpathrectangle{\pgfqpoint{0.407592in}{0.222904in}}{\pgfqpoint{1.411163in}{1.411163in}}%
\pgfusepath{clip}%
\pgfsetbuttcap%
\pgfsetroundjoin%
\definecolor{currentfill}{rgb}{0.556863,0.654902,0.847059}%
\pgfsetfillcolor{currentfill}%
\pgfsetfillopacity{0.130000}%
\pgfsetlinewidth{0.000000pt}%
\definecolor{currentstroke}{rgb}{0.000000,0.000000,0.000000}%
\pgfsetstrokecolor{currentstroke}%
\pgfsetdash{}{0pt}%
\pgfpathmoveto{\pgfqpoint{1.152852in}{0.841507in}}%
\pgfpathlineto{\pgfqpoint{1.168908in}{0.837034in}}%
\pgfpathlineto{\pgfqpoint{1.185006in}{0.834104in}}%
\pgfpathlineto{\pgfqpoint{1.192840in}{0.847786in}}%
\pgfpathlineto{\pgfqpoint{1.200658in}{0.861867in}}%
\pgfpathlineto{\pgfqpoint{1.184574in}{0.864856in}}%
\pgfpathlineto{\pgfqpoint{1.168532in}{0.869248in}}%
\pgfpathlineto{\pgfqpoint{1.160700in}{0.855179in}}%
\pgfpathlineto{\pgfqpoint{1.152852in}{0.841507in}}%
\pgfpathclose%
\pgfusepath{fill}%
\end{pgfscope}%
\begin{pgfscope}%
\pgfpathrectangle{\pgfqpoint{0.407592in}{0.222904in}}{\pgfqpoint{1.411163in}{1.411163in}}%
\pgfusepath{clip}%
\pgfsetbuttcap%
\pgfsetroundjoin%
\definecolor{currentfill}{rgb}{0.556863,0.654902,0.847059}%
\pgfsetfillcolor{currentfill}%
\pgfsetfillopacity{0.130000}%
\pgfsetlinewidth{0.000000pt}%
\definecolor{currentstroke}{rgb}{0.000000,0.000000,0.000000}%
\pgfsetstrokecolor{currentstroke}%
\pgfsetdash{}{0pt}%
\pgfpathmoveto{\pgfqpoint{1.057239in}{0.827855in}}%
\pgfpathlineto{\pgfqpoint{1.073215in}{0.816632in}}%
\pgfpathlineto{\pgfqpoint{1.089214in}{0.806844in}}%
\pgfpathlineto{\pgfqpoint{1.097152in}{0.817914in}}%
\pgfpathlineto{\pgfqpoint{1.105070in}{0.829681in}}%
\pgfpathlineto{\pgfqpoint{1.089096in}{0.839014in}}%
\pgfpathlineto{\pgfqpoint{1.073146in}{0.849711in}}%
\pgfpathlineto{\pgfqpoint{1.065203in}{0.838444in}}%
\pgfpathlineto{\pgfqpoint{1.057239in}{0.827855in}}%
\pgfpathclose%
\pgfusepath{fill}%
\end{pgfscope}%
\begin{pgfscope}%
\pgfpathrectangle{\pgfqpoint{0.407592in}{0.222904in}}{\pgfqpoint{1.411163in}{1.411163in}}%
\pgfusepath{clip}%
\pgfsetbuttcap%
\pgfsetroundjoin%
\definecolor{currentfill}{rgb}{0.556863,0.654902,0.847059}%
\pgfsetfillcolor{currentfill}%
\pgfsetfillopacity{0.130000}%
\pgfsetlinewidth{0.000000pt}%
\definecolor{currentstroke}{rgb}{0.000000,0.000000,0.000000}%
\pgfsetstrokecolor{currentstroke}%
\pgfsetdash{}{0pt}%
\pgfpathmoveto{\pgfqpoint{0.753153in}{1.063432in}}%
\pgfpathlineto{\pgfqpoint{0.769070in}{1.043949in}}%
\pgfpathlineto{\pgfqpoint{0.784995in}{1.024562in}}%
\pgfpathlineto{\pgfqpoint{0.793409in}{1.029376in}}%
\pgfpathlineto{\pgfqpoint{0.801789in}{1.034642in}}%
\pgfpathlineto{\pgfqpoint{0.785906in}{1.053946in}}%
\pgfpathlineto{\pgfqpoint{0.770031in}{1.073348in}}%
\pgfpathlineto{\pgfqpoint{0.761609in}{1.068182in}}%
\pgfpathlineto{\pgfqpoint{0.753153in}{1.063432in}}%
\pgfpathclose%
\pgfusepath{fill}%
\end{pgfscope}%
\begin{pgfscope}%
\pgfpathrectangle{\pgfqpoint{0.407592in}{0.222904in}}{\pgfqpoint{1.411163in}{1.411163in}}%
\pgfusepath{clip}%
\pgfsetbuttcap%
\pgfsetroundjoin%
\definecolor{currentfill}{rgb}{0.556863,0.654902,0.847059}%
\pgfsetfillcolor{currentfill}%
\pgfsetfillopacity{0.130000}%
\pgfsetlinewidth{0.000000pt}%
\definecolor{currentstroke}{rgb}{0.000000,0.000000,0.000000}%
\pgfsetstrokecolor{currentstroke}%
\pgfsetdash{}{0pt}%
\pgfpathmoveto{\pgfqpoint{0.897275in}{0.921998in}}%
\pgfpathlineto{\pgfqpoint{0.913221in}{0.904039in}}%
\pgfpathlineto{\pgfqpoint{0.929177in}{0.886437in}}%
\pgfpathlineto{\pgfqpoint{0.937357in}{0.892868in}}%
\pgfpathlineto{\pgfqpoint{0.945507in}{0.899967in}}%
\pgfpathlineto{\pgfqpoint{0.929591in}{0.917316in}}%
\pgfpathlineto{\pgfqpoint{0.913686in}{0.935037in}}%
\pgfpathlineto{\pgfqpoint{0.905496in}{0.928219in}}%
\pgfpathlineto{\pgfqpoint{0.897275in}{0.921998in}}%
\pgfpathclose%
\pgfusepath{fill}%
\end{pgfscope}%
\begin{pgfscope}%
\pgfpathrectangle{\pgfqpoint{0.407592in}{0.222904in}}{\pgfqpoint{1.411163in}{1.411163in}}%
\pgfusepath{clip}%
\pgfsetbuttcap%
\pgfsetroundjoin%
\definecolor{currentfill}{rgb}{0.556863,0.654902,0.847059}%
\pgfsetfillcolor{currentfill}%
\pgfsetfillopacity{0.130000}%
\pgfsetlinewidth{0.000000pt}%
\definecolor{currentstroke}{rgb}{0.000000,0.000000,0.000000}%
\pgfsetstrokecolor{currentstroke}%
\pgfsetdash{}{0pt}%
\pgfpathmoveto{\pgfqpoint{1.200658in}{0.861867in}}%
\pgfpathlineto{\pgfqpoint{1.216791in}{0.860263in}}%
\pgfpathlineto{\pgfqpoint{1.232974in}{0.860001in}}%
\pgfpathlineto{\pgfqpoint{1.240766in}{0.874191in}}%
\pgfpathlineto{\pgfqpoint{1.248544in}{0.888691in}}%
\pgfpathlineto{\pgfqpoint{1.232373in}{0.889232in}}%
\pgfpathlineto{\pgfqpoint{1.216253in}{0.891004in}}%
\pgfpathlineto{\pgfqpoint{1.208463in}{0.876289in}}%
\pgfpathlineto{\pgfqpoint{1.200658in}{0.861867in}}%
\pgfpathclose%
\pgfusepath{fill}%
\end{pgfscope}%
\begin{pgfscope}%
\pgfpathrectangle{\pgfqpoint{0.407592in}{0.222904in}}{\pgfqpoint{1.411163in}{1.411163in}}%
\pgfusepath{clip}%
\pgfsetbuttcap%
\pgfsetroundjoin%
\definecolor{currentfill}{rgb}{0.556863,0.654902,0.847059}%
\pgfsetfillcolor{currentfill}%
\pgfsetfillopacity{0.130000}%
\pgfsetlinewidth{0.000000pt}%
\definecolor{currentstroke}{rgb}{0.000000,0.000000,0.000000}%
\pgfsetstrokecolor{currentstroke}%
\pgfsetdash{}{0pt}%
\pgfpathmoveto{\pgfqpoint{1.311986in}{0.950046in}}%
\pgfpathlineto{\pgfqpoint{1.328302in}{0.951970in}}%
\pgfpathlineto{\pgfqpoint{1.344679in}{0.954769in}}%
\pgfpathlineto{\pgfqpoint{1.352368in}{0.969605in}}%
\pgfpathlineto{\pgfqpoint{1.360045in}{0.984614in}}%
\pgfpathlineto{\pgfqpoint{1.343681in}{0.982223in}}%
\pgfpathlineto{\pgfqpoint{1.327377in}{0.980663in}}%
\pgfpathlineto{\pgfqpoint{1.319688in}{0.965280in}}%
\pgfpathlineto{\pgfqpoint{1.311986in}{0.950046in}}%
\pgfpathclose%
\pgfusepath{fill}%
\end{pgfscope}%
\begin{pgfscope}%
\pgfpathrectangle{\pgfqpoint{0.407592in}{0.222904in}}{\pgfqpoint{1.411163in}{1.411163in}}%
\pgfusepath{clip}%
\pgfsetbuttcap%
\pgfsetroundjoin%
\definecolor{currentfill}{rgb}{0.556863,0.654902,0.847059}%
\pgfsetfillcolor{currentfill}%
\pgfsetfillopacity{0.130000}%
\pgfsetlinewidth{0.000000pt}%
\definecolor{currentstroke}{rgb}{0.000000,0.000000,0.000000}%
\pgfsetstrokecolor{currentstroke}%
\pgfsetdash{}{0pt}%
\pgfpathmoveto{\pgfqpoint{1.009277in}{0.835997in}}%
\pgfpathlineto{\pgfqpoint{1.025250in}{0.821950in}}%
\pgfpathlineto{\pgfqpoint{1.041239in}{0.808995in}}%
\pgfpathlineto{\pgfqpoint{1.049251in}{0.818015in}}%
\pgfpathlineto{\pgfqpoint{1.057239in}{0.827855in}}%
\pgfpathlineto{\pgfqpoint{1.041282in}{0.840334in}}%
\pgfpathlineto{\pgfqpoint{1.025342in}{0.853898in}}%
\pgfpathlineto{\pgfqpoint{1.017322in}{0.844570in}}%
\pgfpathlineto{\pgfqpoint{1.009277in}{0.835997in}}%
\pgfpathclose%
\pgfusepath{fill}%
\end{pgfscope}%
\begin{pgfscope}%
\pgfpathrectangle{\pgfqpoint{0.407592in}{0.222904in}}{\pgfqpoint{1.411163in}{1.411163in}}%
\pgfusepath{clip}%
\pgfsetbuttcap%
\pgfsetroundjoin%
\definecolor{currentfill}{rgb}{0.556863,0.654902,0.847059}%
\pgfsetfillcolor{currentfill}%
\pgfsetfillopacity{0.130000}%
\pgfsetlinewidth{0.000000pt}%
\definecolor{currentstroke}{rgb}{0.000000,0.000000,0.000000}%
\pgfsetstrokecolor{currentstroke}%
\pgfsetdash{}{0pt}%
\pgfpathmoveto{\pgfqpoint{0.784995in}{1.024562in}}%
\pgfpathlineto{\pgfqpoint{0.800930in}{1.005288in}}%
\pgfpathlineto{\pgfqpoint{0.816873in}{0.986145in}}%
\pgfpathlineto{\pgfqpoint{0.825244in}{0.991020in}}%
\pgfpathlineto{\pgfqpoint{0.833582in}{0.996393in}}%
\pgfpathlineto{\pgfqpoint{0.817681in}{1.015451in}}%
\pgfpathlineto{\pgfqpoint{0.801789in}{1.034642in}}%
\pgfpathlineto{\pgfqpoint{0.793409in}{1.029376in}}%
\pgfpathlineto{\pgfqpoint{0.784995in}{1.024562in}}%
\pgfpathclose%
\pgfusepath{fill}%
\end{pgfscope}%
\begin{pgfscope}%
\pgfpathrectangle{\pgfqpoint{0.407592in}{0.222904in}}{\pgfqpoint{1.411163in}{1.411163in}}%
\pgfusepath{clip}%
\pgfsetbuttcap%
\pgfsetroundjoin%
\definecolor{currentfill}{rgb}{0.556863,0.654902,0.847059}%
\pgfsetfillcolor{currentfill}%
\pgfsetfillopacity{0.130000}%
\pgfsetlinewidth{0.000000pt}%
\definecolor{currentstroke}{rgb}{0.000000,0.000000,0.000000}%
\pgfsetstrokecolor{currentstroke}%
\pgfsetdash{}{0pt}%
\pgfpathmoveto{\pgfqpoint{1.423491in}{1.051556in}}%
\pgfpathlineto{\pgfqpoint{1.440017in}{1.055315in}}%
\pgfpathlineto{\pgfqpoint{1.456607in}{1.059690in}}%
\pgfpathlineto{\pgfqpoint{1.464196in}{1.074587in}}%
\pgfpathlineto{\pgfqpoint{1.471773in}{1.089605in}}%
\pgfpathlineto{\pgfqpoint{1.455197in}{1.085615in}}%
\pgfpathlineto{\pgfqpoint{1.438685in}{1.082223in}}%
\pgfpathlineto{\pgfqpoint{1.431094in}{1.066839in}}%
\pgfpathlineto{\pgfqpoint{1.423491in}{1.051556in}}%
\pgfpathclose%
\pgfusepath{fill}%
\end{pgfscope}%
\begin{pgfscope}%
\pgfpathrectangle{\pgfqpoint{0.407592in}{0.222904in}}{\pgfqpoint{1.411163in}{1.411163in}}%
\pgfusepath{clip}%
\pgfsetbuttcap%
\pgfsetroundjoin%
\definecolor{currentfill}{rgb}{0.556863,0.654902,0.847059}%
\pgfsetfillcolor{currentfill}%
\pgfsetfillopacity{0.130000}%
\pgfsetlinewidth{0.000000pt}%
\definecolor{currentstroke}{rgb}{0.000000,0.000000,0.000000}%
\pgfsetstrokecolor{currentstroke}%
\pgfsetdash{}{0pt}%
\pgfpathmoveto{\pgfqpoint{1.248544in}{0.888691in}}%
\pgfpathlineto{\pgfqpoint{1.264770in}{0.889330in}}%
\pgfpathlineto{\pgfqpoint{1.281052in}{0.891084in}}%
\pgfpathlineto{\pgfqpoint{1.288805in}{0.905473in}}%
\pgfpathlineto{\pgfqpoint{1.296545in}{0.920117in}}%
\pgfpathlineto{\pgfqpoint{1.280275in}{0.918765in}}%
\pgfpathlineto{\pgfqpoint{1.264061in}{0.918456in}}%
\pgfpathlineto{\pgfqpoint{1.256310in}{0.903458in}}%
\pgfpathlineto{\pgfqpoint{1.248544in}{0.888691in}}%
\pgfpathclose%
\pgfusepath{fill}%
\end{pgfscope}%
\begin{pgfscope}%
\pgfpathrectangle{\pgfqpoint{0.407592in}{0.222904in}}{\pgfqpoint{1.411163in}{1.411163in}}%
\pgfusepath{clip}%
\pgfsetbuttcap%
\pgfsetroundjoin%
\definecolor{currentfill}{rgb}{0.556863,0.654902,0.847059}%
\pgfsetfillcolor{currentfill}%
\pgfsetfillopacity{0.130000}%
\pgfsetlinewidth{0.000000pt}%
\definecolor{currentstroke}{rgb}{0.000000,0.000000,0.000000}%
\pgfsetstrokecolor{currentstroke}%
\pgfsetdash{}{0pt}%
\pgfpathmoveto{\pgfqpoint{0.929177in}{0.886437in}}%
\pgfpathlineto{\pgfqpoint{0.945143in}{0.869264in}}%
\pgfpathlineto{\pgfqpoint{0.961118in}{0.852608in}}%
\pgfpathlineto{\pgfqpoint{0.969259in}{0.859276in}}%
\pgfpathlineto{\pgfqpoint{0.977369in}{0.866693in}}%
\pgfpathlineto{\pgfqpoint{0.961432in}{0.883064in}}%
\pgfpathlineto{\pgfqpoint{0.945507in}{0.899967in}}%
\pgfpathlineto{\pgfqpoint{0.937357in}{0.892868in}}%
\pgfpathlineto{\pgfqpoint{0.929177in}{0.886437in}}%
\pgfpathclose%
\pgfusepath{fill}%
\end{pgfscope}%
\begin{pgfscope}%
\pgfpathrectangle{\pgfqpoint{0.407592in}{0.222904in}}{\pgfqpoint{1.411163in}{1.411163in}}%
\pgfusepath{clip}%
\pgfsetbuttcap%
\pgfsetroundjoin%
\definecolor{currentfill}{rgb}{0.556863,0.654902,0.847059}%
\pgfsetfillcolor{currentfill}%
\pgfsetfillopacity{0.130000}%
\pgfsetlinewidth{0.000000pt}%
\definecolor{currentstroke}{rgb}{0.000000,0.000000,0.000000}%
\pgfsetstrokecolor{currentstroke}%
\pgfsetdash{}{0pt}%
\pgfpathmoveto{\pgfqpoint{0.640558in}{1.173621in}}%
\pgfpathlineto{\pgfqpoint{0.656466in}{1.153746in}}%
\pgfpathlineto{\pgfqpoint{0.672382in}{1.133915in}}%
\pgfpathlineto{\pgfqpoint{0.680995in}{1.137829in}}%
\pgfpathlineto{\pgfqpoint{0.689572in}{1.142098in}}%
\pgfpathlineto{\pgfqpoint{0.673699in}{1.161904in}}%
\pgfpathlineto{\pgfqpoint{0.657834in}{1.181752in}}%
\pgfpathlineto{\pgfqpoint{0.649214in}{1.177521in}}%
\pgfpathlineto{\pgfqpoint{0.640558in}{1.173621in}}%
\pgfpathclose%
\pgfusepath{fill}%
\end{pgfscope}%
\begin{pgfscope}%
\pgfpathrectangle{\pgfqpoint{0.407592in}{0.222904in}}{\pgfqpoint{1.411163in}{1.411163in}}%
\pgfusepath{clip}%
\pgfsetbuttcap%
\pgfsetroundjoin%
\definecolor{currentfill}{rgb}{0.556863,0.654902,0.847059}%
\pgfsetfillcolor{currentfill}%
\pgfsetfillopacity{0.130000}%
\pgfsetlinewidth{0.000000pt}%
\definecolor{currentstroke}{rgb}{0.000000,0.000000,0.000000}%
\pgfsetstrokecolor{currentstroke}%
\pgfsetdash{}{0pt}%
\pgfpathmoveto{\pgfqpoint{1.360045in}{0.984614in}}%
\pgfpathlineto{\pgfqpoint{1.376471in}{0.987791in}}%
\pgfpathlineto{\pgfqpoint{1.392960in}{0.991706in}}%
\pgfpathlineto{\pgfqpoint{1.400611in}{1.006448in}}%
\pgfpathlineto{\pgfqpoint{1.408250in}{1.021348in}}%
\pgfpathlineto{\pgfqpoint{1.391774in}{1.017856in}}%
\pgfpathlineto{\pgfqpoint{1.375362in}{1.015076in}}%
\pgfpathlineto{\pgfqpoint{1.367710in}{0.999777in}}%
\pgfpathlineto{\pgfqpoint{1.360045in}{0.984614in}}%
\pgfpathclose%
\pgfusepath{fill}%
\end{pgfscope}%
\begin{pgfscope}%
\pgfpathrectangle{\pgfqpoint{0.407592in}{0.222904in}}{\pgfqpoint{1.411163in}{1.411163in}}%
\pgfusepath{clip}%
\pgfsetbuttcap%
\pgfsetroundjoin%
\definecolor{currentfill}{rgb}{0.556863,0.654902,0.847059}%
\pgfsetfillcolor{currentfill}%
\pgfsetfillopacity{0.130000}%
\pgfsetlinewidth{0.000000pt}%
\definecolor{currentstroke}{rgb}{0.000000,0.000000,0.000000}%
\pgfsetstrokecolor{currentstroke}%
\pgfsetdash{}{0pt}%
\pgfpathmoveto{\pgfqpoint{0.816873in}{0.986145in}}%
\pgfpathlineto{\pgfqpoint{0.832825in}{0.967154in}}%
\pgfpathlineto{\pgfqpoint{0.848786in}{0.948342in}}%
\pgfpathlineto{\pgfqpoint{0.857115in}{0.953279in}}%
\pgfpathlineto{\pgfqpoint{0.865410in}{0.958767in}}%
\pgfpathlineto{\pgfqpoint{0.849492in}{0.977490in}}%
\pgfpathlineto{\pgfqpoint{0.833582in}{0.996393in}}%
\pgfpathlineto{\pgfqpoint{0.825244in}{0.991020in}}%
\pgfpathlineto{\pgfqpoint{0.816873in}{0.986145in}}%
\pgfpathclose%
\pgfusepath{fill}%
\end{pgfscope}%
\begin{pgfscope}%
\pgfpathrectangle{\pgfqpoint{0.407592in}{0.222904in}}{\pgfqpoint{1.411163in}{1.411163in}}%
\pgfusepath{clip}%
\pgfsetbuttcap%
\pgfsetroundjoin%
\definecolor{currentfill}{rgb}{0.556863,0.654902,0.847059}%
\pgfsetfillcolor{currentfill}%
\pgfsetfillopacity{0.130000}%
\pgfsetlinewidth{0.000000pt}%
\definecolor{currentstroke}{rgb}{0.000000,0.000000,0.000000}%
\pgfsetstrokecolor{currentstroke}%
\pgfsetdash{}{0pt}%
\pgfpathmoveto{\pgfqpoint{1.089214in}{0.806844in}}%
\pgfpathlineto{\pgfqpoint{1.105238in}{0.798654in}}%
\pgfpathlineto{\pgfqpoint{1.121293in}{0.792199in}}%
\pgfpathlineto{\pgfqpoint{1.129210in}{0.803554in}}%
\pgfpathlineto{\pgfqpoint{1.137108in}{0.815616in}}%
\pgfpathlineto{\pgfqpoint{1.121073in}{0.821845in}}%
\pgfpathlineto{\pgfqpoint{1.105070in}{0.829681in}}%
\pgfpathlineto{\pgfqpoint{1.097152in}{0.817914in}}%
\pgfpathlineto{\pgfqpoint{1.089214in}{0.806844in}}%
\pgfpathclose%
\pgfusepath{fill}%
\end{pgfscope}%
\begin{pgfscope}%
\pgfpathrectangle{\pgfqpoint{0.407592in}{0.222904in}}{\pgfqpoint{1.411163in}{1.411163in}}%
\pgfusepath{clip}%
\pgfsetbuttcap%
\pgfsetroundjoin%
\definecolor{currentfill}{rgb}{0.556863,0.654902,0.847059}%
\pgfsetfillcolor{currentfill}%
\pgfsetfillopacity{0.130000}%
\pgfsetlinewidth{0.000000pt}%
\definecolor{currentstroke}{rgb}{0.000000,0.000000,0.000000}%
\pgfsetstrokecolor{currentstroke}%
\pgfsetdash{}{0pt}%
\pgfpathmoveto{\pgfqpoint{1.137108in}{0.815616in}}%
\pgfpathlineto{\pgfqpoint{1.153180in}{0.811061in}}%
\pgfpathlineto{\pgfqpoint{1.169293in}{0.808200in}}%
\pgfpathlineto{\pgfqpoint{1.177157in}{0.820884in}}%
\pgfpathlineto{\pgfqpoint{1.185006in}{0.834104in}}%
\pgfpathlineto{\pgfqpoint{1.168908in}{0.837034in}}%
\pgfpathlineto{\pgfqpoint{1.152852in}{0.841507in}}%
\pgfpathlineto{\pgfqpoint{1.144988in}{0.828294in}}%
\pgfpathlineto{\pgfqpoint{1.137108in}{0.815616in}}%
\pgfpathclose%
\pgfusepath{fill}%
\end{pgfscope}%
\begin{pgfscope}%
\pgfpathrectangle{\pgfqpoint{0.407592in}{0.222904in}}{\pgfqpoint{1.411163in}{1.411163in}}%
\pgfusepath{clip}%
\pgfsetbuttcap%
\pgfsetroundjoin%
\definecolor{currentfill}{rgb}{0.556863,0.654902,0.847059}%
\pgfsetfillcolor{currentfill}%
\pgfsetfillopacity{0.130000}%
\pgfsetlinewidth{0.000000pt}%
\definecolor{currentstroke}{rgb}{0.000000,0.000000,0.000000}%
\pgfsetstrokecolor{currentstroke}%
\pgfsetdash{}{0pt}%
\pgfpathmoveto{\pgfqpoint{0.672382in}{1.133915in}}%
\pgfpathlineto{\pgfqpoint{0.688307in}{1.114135in}}%
\pgfpathlineto{\pgfqpoint{0.704240in}{1.094413in}}%
\pgfpathlineto{\pgfqpoint{0.712811in}{1.098333in}}%
\pgfpathlineto{\pgfqpoint{0.721345in}{1.102638in}}%
\pgfpathlineto{\pgfqpoint{0.705454in}{1.122340in}}%
\pgfpathlineto{\pgfqpoint{0.689572in}{1.142098in}}%
\pgfpathlineto{\pgfqpoint{0.680995in}{1.137829in}}%
\pgfpathlineto{\pgfqpoint{0.672382in}{1.133915in}}%
\pgfpathclose%
\pgfusepath{fill}%
\end{pgfscope}%
\begin{pgfscope}%
\pgfpathrectangle{\pgfqpoint{0.407592in}{0.222904in}}{\pgfqpoint{1.411163in}{1.411163in}}%
\pgfusepath{clip}%
\pgfsetbuttcap%
\pgfsetroundjoin%
\definecolor{currentfill}{rgb}{0.556863,0.654902,0.847059}%
\pgfsetfillcolor{currentfill}%
\pgfsetfillopacity{0.130000}%
\pgfsetlinewidth{0.000000pt}%
\definecolor{currentstroke}{rgb}{0.000000,0.000000,0.000000}%
\pgfsetstrokecolor{currentstroke}%
\pgfsetdash{}{0pt}%
\pgfpathmoveto{\pgfqpoint{0.961118in}{0.852608in}}%
\pgfpathlineto{\pgfqpoint{0.977104in}{0.836576in}}%
\pgfpathlineto{\pgfqpoint{0.993102in}{0.821299in}}%
\pgfpathlineto{\pgfqpoint{1.001204in}{0.828226in}}%
\pgfpathlineto{\pgfqpoint{1.009277in}{0.835997in}}%
\pgfpathlineto{\pgfqpoint{0.993317in}{0.850961in}}%
\pgfpathlineto{\pgfqpoint{0.977369in}{0.866693in}}%
\pgfpathlineto{\pgfqpoint{0.969259in}{0.859276in}}%
\pgfpathlineto{\pgfqpoint{0.961118in}{0.852608in}}%
\pgfpathclose%
\pgfusepath{fill}%
\end{pgfscope}%
\begin{pgfscope}%
\pgfpathrectangle{\pgfqpoint{0.407592in}{0.222904in}}{\pgfqpoint{1.411163in}{1.411163in}}%
\pgfusepath{clip}%
\pgfsetbuttcap%
\pgfsetroundjoin%
\definecolor{currentfill}{rgb}{0.556863,0.654902,0.847059}%
\pgfsetfillcolor{currentfill}%
\pgfsetfillopacity{0.130000}%
\pgfsetlinewidth{0.000000pt}%
\definecolor{currentstroke}{rgb}{0.000000,0.000000,0.000000}%
\pgfsetstrokecolor{currentstroke}%
\pgfsetdash{}{0pt}%
\pgfpathmoveto{\pgfqpoint{1.296545in}{0.920117in}}%
\pgfpathlineto{\pgfqpoint{1.312873in}{0.922452in}}%
\pgfpathlineto{\pgfqpoint{1.329263in}{0.925705in}}%
\pgfpathlineto{\pgfqpoint{1.336977in}{0.940127in}}%
\pgfpathlineto{\pgfqpoint{1.344679in}{0.954769in}}%
\pgfpathlineto{\pgfqpoint{1.328302in}{0.951970in}}%
\pgfpathlineto{\pgfqpoint{1.311986in}{0.950046in}}%
\pgfpathlineto{\pgfqpoint{1.304272in}{0.934984in}}%
\pgfpathlineto{\pgfqpoint{1.296545in}{0.920117in}}%
\pgfpathclose%
\pgfusepath{fill}%
\end{pgfscope}%
\begin{pgfscope}%
\pgfpathrectangle{\pgfqpoint{0.407592in}{0.222904in}}{\pgfqpoint{1.411163in}{1.411163in}}%
\pgfusepath{clip}%
\pgfsetbuttcap%
\pgfsetroundjoin%
\definecolor{currentfill}{rgb}{0.556863,0.654902,0.847059}%
\pgfsetfillcolor{currentfill}%
\pgfsetfillopacity{0.130000}%
\pgfsetlinewidth{0.000000pt}%
\definecolor{currentstroke}{rgb}{0.000000,0.000000,0.000000}%
\pgfsetstrokecolor{currentstroke}%
\pgfsetdash{}{0pt}%
\pgfpathmoveto{\pgfqpoint{1.185006in}{0.834104in}}%
\pgfpathlineto{\pgfqpoint{1.201151in}{0.832694in}}%
\pgfpathlineto{\pgfqpoint{1.217348in}{0.832744in}}%
\pgfpathlineto{\pgfqpoint{1.225168in}{0.846167in}}%
\pgfpathlineto{\pgfqpoint{1.232974in}{0.860001in}}%
\pgfpathlineto{\pgfqpoint{1.216791in}{0.860263in}}%
\pgfpathlineto{\pgfqpoint{1.200658in}{0.861867in}}%
\pgfpathlineto{\pgfqpoint{1.192840in}{0.847786in}}%
\pgfpathlineto{\pgfqpoint{1.185006in}{0.834104in}}%
\pgfpathclose%
\pgfusepath{fill}%
\end{pgfscope}%
\begin{pgfscope}%
\pgfpathrectangle{\pgfqpoint{0.407592in}{0.222904in}}{\pgfqpoint{1.411163in}{1.411163in}}%
\pgfusepath{clip}%
\pgfsetbuttcap%
\pgfsetroundjoin%
\definecolor{currentfill}{rgb}{0.556863,0.654902,0.847059}%
\pgfsetfillcolor{currentfill}%
\pgfsetfillopacity{0.130000}%
\pgfsetlinewidth{0.000000pt}%
\definecolor{currentstroke}{rgb}{0.000000,0.000000,0.000000}%
\pgfsetstrokecolor{currentstroke}%
\pgfsetdash{}{0pt}%
\pgfpathmoveto{\pgfqpoint{1.041239in}{0.808995in}}%
\pgfpathlineto{\pgfqpoint{1.057245in}{0.797326in}}%
\pgfpathlineto{\pgfqpoint{1.073272in}{0.787145in}}%
\pgfpathlineto{\pgfqpoint{1.081254in}{0.796559in}}%
\pgfpathlineto{\pgfqpoint{1.089214in}{0.806844in}}%
\pgfpathlineto{\pgfqpoint{1.073215in}{0.816632in}}%
\pgfpathlineto{\pgfqpoint{1.057239in}{0.827855in}}%
\pgfpathlineto{\pgfqpoint{1.049251in}{0.818015in}}%
\pgfpathlineto{\pgfqpoint{1.041239in}{0.808995in}}%
\pgfpathclose%
\pgfusepath{fill}%
\end{pgfscope}%
\begin{pgfscope}%
\pgfpathrectangle{\pgfqpoint{0.407592in}{0.222904in}}{\pgfqpoint{1.411163in}{1.411163in}}%
\pgfusepath{clip}%
\pgfsetbuttcap%
\pgfsetroundjoin%
\definecolor{currentfill}{rgb}{0.556863,0.654902,0.847059}%
\pgfsetfillcolor{currentfill}%
\pgfsetfillopacity{0.130000}%
\pgfsetlinewidth{0.000000pt}%
\definecolor{currentstroke}{rgb}{0.000000,0.000000,0.000000}%
\pgfsetstrokecolor{currentstroke}%
\pgfsetdash{}{0pt}%
\pgfpathmoveto{\pgfqpoint{0.848786in}{0.948342in}}%
\pgfpathlineto{\pgfqpoint{0.864756in}{0.929741in}}%
\pgfpathlineto{\pgfqpoint{0.880735in}{0.911389in}}%
\pgfpathlineto{\pgfqpoint{0.889022in}{0.916386in}}%
\pgfpathlineto{\pgfqpoint{0.897275in}{0.921998in}}%
\pgfpathlineto{\pgfqpoint{0.881338in}{0.940256in}}%
\pgfpathlineto{\pgfqpoint{0.865410in}{0.958767in}}%
\pgfpathlineto{\pgfqpoint{0.857115in}{0.953279in}}%
\pgfpathlineto{\pgfqpoint{0.848786in}{0.948342in}}%
\pgfpathclose%
\pgfusepath{fill}%
\end{pgfscope}%
\begin{pgfscope}%
\pgfpathrectangle{\pgfqpoint{0.407592in}{0.222904in}}{\pgfqpoint{1.411163in}{1.411163in}}%
\pgfusepath{clip}%
\pgfsetbuttcap%
\pgfsetroundjoin%
\definecolor{currentfill}{rgb}{0.556863,0.654902,0.847059}%
\pgfsetfillcolor{currentfill}%
\pgfsetfillopacity{0.130000}%
\pgfsetlinewidth{0.000000pt}%
\definecolor{currentstroke}{rgb}{0.000000,0.000000,0.000000}%
\pgfsetstrokecolor{currentstroke}%
\pgfsetdash{}{0pt}%
\pgfpathmoveto{\pgfqpoint{0.704240in}{1.094413in}}%
\pgfpathlineto{\pgfqpoint{0.720183in}{1.074758in}}%
\pgfpathlineto{\pgfqpoint{0.736134in}{1.055179in}}%
\pgfpathlineto{\pgfqpoint{0.744661in}{1.059097in}}%
\pgfpathlineto{\pgfqpoint{0.753153in}{1.063432in}}%
\pgfpathlineto{\pgfqpoint{0.737244in}{1.082998in}}%
\pgfpathlineto{\pgfqpoint{0.721345in}{1.102638in}}%
\pgfpathlineto{\pgfqpoint{0.712811in}{1.098333in}}%
\pgfpathlineto{\pgfqpoint{0.704240in}{1.094413in}}%
\pgfpathclose%
\pgfusepath{fill}%
\end{pgfscope}%
\begin{pgfscope}%
\pgfpathrectangle{\pgfqpoint{0.407592in}{0.222904in}}{\pgfqpoint{1.411163in}{1.411163in}}%
\pgfusepath{clip}%
\pgfsetbuttcap%
\pgfsetroundjoin%
\definecolor{currentfill}{rgb}{0.556863,0.654902,0.847059}%
\pgfsetfillcolor{currentfill}%
\pgfsetfillopacity{0.130000}%
\pgfsetlinewidth{0.000000pt}%
\definecolor{currentstroke}{rgb}{0.000000,0.000000,0.000000}%
\pgfsetstrokecolor{currentstroke}%
\pgfsetdash{}{0pt}%
\pgfpathmoveto{\pgfqpoint{1.471773in}{1.089605in}}%
\pgfpathlineto{\pgfqpoint{1.488414in}{1.094163in}}%
\pgfpathlineto{\pgfqpoint{1.505121in}{1.099259in}}%
\pgfpathlineto{\pgfqpoint{1.512671in}{1.114007in}}%
\pgfpathlineto{\pgfqpoint{1.495971in}{1.109103in}}%
\pgfpathlineto{\pgfqpoint{1.479338in}{1.104731in}}%
\pgfusepath{fill}%
\end{pgfscope}%
\begin{pgfscope}%
\pgfpathrectangle{\pgfqpoint{0.407592in}{0.222904in}}{\pgfqpoint{1.411163in}{1.411163in}}%
\pgfusepath{clip}%
\pgfsetbuttcap%
\pgfsetroundjoin%
\definecolor{currentfill}{rgb}{0.556863,0.654902,0.847059}%
\pgfsetfillcolor{currentfill}%
\pgfsetfillopacity{0.130000}%
\pgfsetlinewidth{0.000000pt}%
\definecolor{currentstroke}{rgb}{0.000000,0.000000,0.000000}%
\pgfsetstrokecolor{currentstroke}%
\pgfsetdash{}{0pt}%
\pgfpathmoveto{\pgfqpoint{1.408250in}{1.021348in}}%
\pgfpathlineto{\pgfqpoint{1.424789in}{1.025511in}}%
\pgfpathlineto{\pgfqpoint{1.441394in}{1.030305in}}%
\pgfpathlineto{\pgfqpoint{1.449006in}{1.044924in}}%
\pgfpathlineto{\pgfqpoint{1.456607in}{1.059690in}}%
\pgfpathlineto{\pgfqpoint{1.440017in}{1.055315in}}%
\pgfpathlineto{\pgfqpoint{1.423491in}{1.051556in}}%
\pgfpathlineto{\pgfqpoint{1.415877in}{1.036388in}}%
\pgfpathlineto{\pgfqpoint{1.408250in}{1.021348in}}%
\pgfpathclose%
\pgfusepath{fill}%
\end{pgfscope}%
\begin{pgfscope}%
\pgfpathrectangle{\pgfqpoint{0.407592in}{0.222904in}}{\pgfqpoint{1.411163in}{1.411163in}}%
\pgfusepath{clip}%
\pgfsetbuttcap%
\pgfsetroundjoin%
\definecolor{currentfill}{rgb}{0.556863,0.654902,0.847059}%
\pgfsetfillcolor{currentfill}%
\pgfsetfillopacity{0.130000}%
\pgfsetlinewidth{0.000000pt}%
\definecolor{currentstroke}{rgb}{0.000000,0.000000,0.000000}%
\pgfsetstrokecolor{currentstroke}%
\pgfsetdash{}{0pt}%
\pgfpathmoveto{\pgfqpoint{1.232974in}{0.860001in}}%
\pgfpathlineto{\pgfqpoint{1.249211in}{0.861012in}}%
\pgfpathlineto{\pgfqpoint{1.265506in}{0.863214in}}%
\pgfpathlineto{\pgfqpoint{1.273286in}{0.876984in}}%
\pgfpathlineto{\pgfqpoint{1.281052in}{0.891084in}}%
\pgfpathlineto{\pgfqpoint{1.264770in}{0.889330in}}%
\pgfpathlineto{\pgfqpoint{1.248544in}{0.888691in}}%
\pgfpathlineto{\pgfqpoint{1.240766in}{0.874191in}}%
\pgfpathlineto{\pgfqpoint{1.232974in}{0.860001in}}%
\pgfpathclose%
\pgfusepath{fill}%
\end{pgfscope}%
\begin{pgfscope}%
\pgfpathrectangle{\pgfqpoint{0.407592in}{0.222904in}}{\pgfqpoint{1.411163in}{1.411163in}}%
\pgfusepath{clip}%
\pgfsetbuttcap%
\pgfsetroundjoin%
\definecolor{currentfill}{rgb}{0.556863,0.654902,0.847059}%
\pgfsetfillcolor{currentfill}%
\pgfsetfillopacity{0.130000}%
\pgfsetlinewidth{0.000000pt}%
\definecolor{currentstroke}{rgb}{0.000000,0.000000,0.000000}%
\pgfsetstrokecolor{currentstroke}%
\pgfsetdash{}{0pt}%
\pgfpathmoveto{\pgfqpoint{0.880735in}{0.911389in}}%
\pgfpathlineto{\pgfqpoint{0.896723in}{0.893334in}}%
\pgfpathlineto{\pgfqpoint{0.912720in}{0.875633in}}%
\pgfpathlineto{\pgfqpoint{0.920965in}{0.880689in}}%
\pgfpathlineto{\pgfqpoint{0.929177in}{0.886437in}}%
\pgfpathlineto{\pgfqpoint{0.913221in}{0.904039in}}%
\pgfpathlineto{\pgfqpoint{0.897275in}{0.921998in}}%
\pgfpathlineto{\pgfqpoint{0.889022in}{0.916386in}}%
\pgfpathlineto{\pgfqpoint{0.880735in}{0.911389in}}%
\pgfpathclose%
\pgfusepath{fill}%
\end{pgfscope}%
\begin{pgfscope}%
\pgfpathrectangle{\pgfqpoint{0.407592in}{0.222904in}}{\pgfqpoint{1.411163in}{1.411163in}}%
\pgfusepath{clip}%
\pgfsetbuttcap%
\pgfsetroundjoin%
\definecolor{currentfill}{rgb}{0.556863,0.654902,0.847059}%
\pgfsetfillcolor{currentfill}%
\pgfsetfillopacity{0.130000}%
\pgfsetlinewidth{0.000000pt}%
\definecolor{currentstroke}{rgb}{0.000000,0.000000,0.000000}%
\pgfsetstrokecolor{currentstroke}%
\pgfsetdash{}{0pt}%
\pgfpathmoveto{\pgfqpoint{0.736134in}{1.055179in}}%
\pgfpathlineto{\pgfqpoint{0.752094in}{1.035688in}}%
\pgfpathlineto{\pgfqpoint{0.768062in}{1.016299in}}%
\pgfpathlineto{\pgfqpoint{0.776547in}{1.020204in}}%
\pgfpathlineto{\pgfqpoint{0.784995in}{1.024562in}}%
\pgfpathlineto{\pgfqpoint{0.769070in}{1.043949in}}%
\pgfpathlineto{\pgfqpoint{0.753153in}{1.063432in}}%
\pgfpathlineto{\pgfqpoint{0.744661in}{1.059097in}}%
\pgfpathlineto{\pgfqpoint{0.736134in}{1.055179in}}%
\pgfpathclose%
\pgfusepath{fill}%
\end{pgfscope}%
\begin{pgfscope}%
\pgfpathrectangle{\pgfqpoint{0.407592in}{0.222904in}}{\pgfqpoint{1.411163in}{1.411163in}}%
\pgfusepath{clip}%
\pgfsetbuttcap%
\pgfsetroundjoin%
\definecolor{currentfill}{rgb}{0.556863,0.654902,0.847059}%
\pgfsetfillcolor{currentfill}%
\pgfsetfillopacity{0.130000}%
\pgfsetlinewidth{0.000000pt}%
\definecolor{currentstroke}{rgb}{0.000000,0.000000,0.000000}%
\pgfsetstrokecolor{currentstroke}%
\pgfsetdash{}{0pt}%
\pgfpathmoveto{\pgfqpoint{0.993102in}{0.821299in}}%
\pgfpathlineto{\pgfqpoint{1.009112in}{0.806933in}}%
\pgfpathlineto{\pgfqpoint{1.025135in}{0.793661in}}%
\pgfpathlineto{\pgfqpoint{1.033201in}{0.800859in}}%
\pgfpathlineto{\pgfqpoint{1.041239in}{0.808995in}}%
\pgfpathlineto{\pgfqpoint{1.025250in}{0.821950in}}%
\pgfpathlineto{\pgfqpoint{1.009277in}{0.835997in}}%
\pgfpathlineto{\pgfqpoint{1.001204in}{0.828226in}}%
\pgfpathlineto{\pgfqpoint{0.993102in}{0.821299in}}%
\pgfpathclose%
\pgfusepath{fill}%
\end{pgfscope}%
\begin{pgfscope}%
\pgfpathrectangle{\pgfqpoint{0.407592in}{0.222904in}}{\pgfqpoint{1.411163in}{1.411163in}}%
\pgfusepath{clip}%
\pgfsetbuttcap%
\pgfsetroundjoin%
\definecolor{currentfill}{rgb}{0.556863,0.654902,0.847059}%
\pgfsetfillcolor{currentfill}%
\pgfsetfillopacity{0.130000}%
\pgfsetlinewidth{0.000000pt}%
\definecolor{currentstroke}{rgb}{0.000000,0.000000,0.000000}%
\pgfsetstrokecolor{currentstroke}%
\pgfsetdash{}{0pt}%
\pgfpathmoveto{\pgfqpoint{1.344679in}{0.954769in}}%
\pgfpathlineto{\pgfqpoint{1.361118in}{0.958387in}}%
\pgfpathlineto{\pgfqpoint{1.377621in}{0.962768in}}%
\pgfpathlineto{\pgfqpoint{1.385296in}{0.977140in}}%
\pgfpathlineto{\pgfqpoint{1.392960in}{0.991706in}}%
\pgfpathlineto{\pgfqpoint{1.376471in}{0.987791in}}%
\pgfpathlineto{\pgfqpoint{1.360045in}{0.984614in}}%
\pgfpathlineto{\pgfqpoint{1.352368in}{0.969605in}}%
\pgfpathlineto{\pgfqpoint{1.344679in}{0.954769in}}%
\pgfpathclose%
\pgfusepath{fill}%
\end{pgfscope}%
\begin{pgfscope}%
\pgfpathrectangle{\pgfqpoint{0.407592in}{0.222904in}}{\pgfqpoint{1.411163in}{1.411163in}}%
\pgfusepath{clip}%
\pgfsetbuttcap%
\pgfsetroundjoin%
\definecolor{currentfill}{rgb}{0.556863,0.654902,0.847059}%
\pgfsetfillcolor{currentfill}%
\pgfsetfillopacity{0.130000}%
\pgfsetlinewidth{0.000000pt}%
\definecolor{currentstroke}{rgb}{0.000000,0.000000,0.000000}%
\pgfsetstrokecolor{currentstroke}%
\pgfsetdash{}{0pt}%
\pgfpathmoveto{\pgfqpoint{0.768062in}{1.016299in}}%
\pgfpathlineto{\pgfqpoint{0.784039in}{0.997028in}}%
\pgfpathlineto{\pgfqpoint{0.800025in}{0.977893in}}%
\pgfpathlineto{\pgfqpoint{0.808467in}{0.981770in}}%
\pgfpathlineto{\pgfqpoint{0.816873in}{0.986145in}}%
\pgfpathlineto{\pgfqpoint{0.800930in}{1.005288in}}%
\pgfpathlineto{\pgfqpoint{0.784995in}{1.024562in}}%
\pgfpathlineto{\pgfqpoint{0.776547in}{1.020204in}}%
\pgfpathlineto{\pgfqpoint{0.768062in}{1.016299in}}%
\pgfpathclose%
\pgfusepath{fill}%
\end{pgfscope}%
\begin{pgfscope}%
\pgfpathrectangle{\pgfqpoint{0.407592in}{0.222904in}}{\pgfqpoint{1.411163in}{1.411163in}}%
\pgfusepath{clip}%
\pgfsetbuttcap%
\pgfsetroundjoin%
\definecolor{currentfill}{rgb}{0.556863,0.654902,0.847059}%
\pgfsetfillcolor{currentfill}%
\pgfsetfillopacity{0.130000}%
\pgfsetlinewidth{0.000000pt}%
\definecolor{currentstroke}{rgb}{0.000000,0.000000,0.000000}%
\pgfsetstrokecolor{currentstroke}%
\pgfsetdash{}{0pt}%
\pgfpathmoveto{\pgfqpoint{1.121293in}{0.792199in}}%
\pgfpathlineto{\pgfqpoint{1.137384in}{0.787564in}}%
\pgfpathlineto{\pgfqpoint{1.153515in}{0.784774in}}%
\pgfpathlineto{\pgfqpoint{1.161412in}{0.796133in}}%
\pgfpathlineto{\pgfqpoint{1.169293in}{0.808200in}}%
\pgfpathlineto{\pgfqpoint{1.153180in}{0.811061in}}%
\pgfpathlineto{\pgfqpoint{1.137108in}{0.815616in}}%
\pgfpathlineto{\pgfqpoint{1.129210in}{0.803554in}}%
\pgfpathlineto{\pgfqpoint{1.121293in}{0.792199in}}%
\pgfpathclose%
\pgfusepath{fill}%
\end{pgfscope}%
\begin{pgfscope}%
\pgfpathrectangle{\pgfqpoint{0.407592in}{0.222904in}}{\pgfqpoint{1.411163in}{1.411163in}}%
\pgfusepath{clip}%
\pgfsetbuttcap%
\pgfsetroundjoin%
\definecolor{currentfill}{rgb}{0.556863,0.654902,0.847059}%
\pgfsetfillcolor{currentfill}%
\pgfsetfillopacity{0.130000}%
\pgfsetlinewidth{0.000000pt}%
\definecolor{currentstroke}{rgb}{0.000000,0.000000,0.000000}%
\pgfsetstrokecolor{currentstroke}%
\pgfsetdash{}{0pt}%
\pgfpathmoveto{\pgfqpoint{0.912720in}{0.875633in}}%
\pgfpathlineto{\pgfqpoint{0.928726in}{0.858359in}}%
\pgfpathlineto{\pgfqpoint{0.944742in}{0.841599in}}%
\pgfpathlineto{\pgfqpoint{0.952947in}{0.846712in}}%
\pgfpathlineto{\pgfqpoint{0.961118in}{0.852608in}}%
\pgfpathlineto{\pgfqpoint{0.945143in}{0.869264in}}%
\pgfpathlineto{\pgfqpoint{0.929177in}{0.886437in}}%
\pgfpathlineto{\pgfqpoint{0.920965in}{0.880689in}}%
\pgfpathlineto{\pgfqpoint{0.912720in}{0.875633in}}%
\pgfpathclose%
\pgfusepath{fill}%
\end{pgfscope}%
\begin{pgfscope}%
\pgfpathrectangle{\pgfqpoint{0.407592in}{0.222904in}}{\pgfqpoint{1.411163in}{1.411163in}}%
\pgfusepath{clip}%
\pgfsetbuttcap%
\pgfsetroundjoin%
\definecolor{currentfill}{rgb}{0.556863,0.654902,0.847059}%
\pgfsetfillcolor{currentfill}%
\pgfsetfillopacity{0.130000}%
\pgfsetlinewidth{0.000000pt}%
\definecolor{currentstroke}{rgb}{0.000000,0.000000,0.000000}%
\pgfsetstrokecolor{currentstroke}%
\pgfsetdash{}{0pt}%
\pgfpathmoveto{\pgfqpoint{1.073272in}{0.787145in}}%
\pgfpathlineto{\pgfqpoint{1.089322in}{0.778647in}}%
\pgfpathlineto{\pgfqpoint{1.105400in}{0.771992in}}%
\pgfpathlineto{\pgfqpoint{1.113357in}{0.781646in}}%
\pgfpathlineto{\pgfqpoint{1.121293in}{0.792199in}}%
\pgfpathlineto{\pgfqpoint{1.105238in}{0.798654in}}%
\pgfpathlineto{\pgfqpoint{1.089214in}{0.806844in}}%
\pgfpathlineto{\pgfqpoint{1.081254in}{0.796559in}}%
\pgfpathlineto{\pgfqpoint{1.073272in}{0.787145in}}%
\pgfpathclose%
\pgfusepath{fill}%
\end{pgfscope}%
\begin{pgfscope}%
\pgfpathrectangle{\pgfqpoint{0.407592in}{0.222904in}}{\pgfqpoint{1.411163in}{1.411163in}}%
\pgfusepath{clip}%
\pgfsetbuttcap%
\pgfsetroundjoin%
\definecolor{currentfill}{rgb}{0.556863,0.654902,0.847059}%
\pgfsetfillcolor{currentfill}%
\pgfsetfillopacity{0.130000}%
\pgfsetlinewidth{0.000000pt}%
\definecolor{currentstroke}{rgb}{0.000000,0.000000,0.000000}%
\pgfsetstrokecolor{currentstroke}%
\pgfsetdash{}{0pt}%
\pgfpathmoveto{\pgfqpoint{1.281052in}{0.891084in}}%
\pgfpathlineto{\pgfqpoint{1.297393in}{0.893877in}}%
\pgfpathlineto{\pgfqpoint{1.313797in}{0.897632in}}%
\pgfpathlineto{\pgfqpoint{1.321536in}{0.911530in}}%
\pgfpathlineto{\pgfqpoint{1.329263in}{0.925705in}}%
\pgfpathlineto{\pgfqpoint{1.312873in}{0.922452in}}%
\pgfpathlineto{\pgfqpoint{1.296545in}{0.920117in}}%
\pgfpathlineto{\pgfqpoint{1.288805in}{0.905473in}}%
\pgfpathlineto{\pgfqpoint{1.281052in}{0.891084in}}%
\pgfpathclose%
\pgfusepath{fill}%
\end{pgfscope}%
\begin{pgfscope}%
\pgfpathrectangle{\pgfqpoint{0.407592in}{0.222904in}}{\pgfqpoint{1.411163in}{1.411163in}}%
\pgfusepath{clip}%
\pgfsetbuttcap%
\pgfsetroundjoin%
\definecolor{currentfill}{rgb}{0.556863,0.654902,0.847059}%
\pgfsetfillcolor{currentfill}%
\pgfsetfillopacity{0.130000}%
\pgfsetlinewidth{0.000000pt}%
\definecolor{currentstroke}{rgb}{0.000000,0.000000,0.000000}%
\pgfsetstrokecolor{currentstroke}%
\pgfsetdash{}{0pt}%
\pgfpathmoveto{\pgfqpoint{1.456607in}{1.059690in}}%
\pgfpathlineto{\pgfqpoint{1.473263in}{1.064645in}}%
\pgfpathlineto{\pgfqpoint{1.489986in}{1.070146in}}%
\pgfpathlineto{\pgfqpoint{1.497560in}{1.084635in}}%
\pgfpathlineto{\pgfqpoint{1.505121in}{1.099259in}}%
\pgfpathlineto{\pgfqpoint{1.488414in}{1.094163in}}%
\pgfpathlineto{\pgfqpoint{1.471773in}{1.089605in}}%
\pgfpathlineto{\pgfqpoint{1.464196in}{1.074587in}}%
\pgfpathlineto{\pgfqpoint{1.456607in}{1.059690in}}%
\pgfpathclose%
\pgfusepath{fill}%
\end{pgfscope}%
\begin{pgfscope}%
\pgfpathrectangle{\pgfqpoint{0.407592in}{0.222904in}}{\pgfqpoint{1.411163in}{1.411163in}}%
\pgfusepath{clip}%
\pgfsetbuttcap%
\pgfsetroundjoin%
\definecolor{currentfill}{rgb}{0.556863,0.654902,0.847059}%
\pgfsetfillcolor{currentfill}%
\pgfsetfillopacity{0.130000}%
\pgfsetlinewidth{0.000000pt}%
\definecolor{currentstroke}{rgb}{0.000000,0.000000,0.000000}%
\pgfsetstrokecolor{currentstroke}%
\pgfsetdash{}{0pt}%
\pgfpathmoveto{\pgfqpoint{0.623136in}{1.166815in}}%
\pgfpathlineto{\pgfqpoint{0.639086in}{1.146958in}}%
\pgfpathlineto{\pgfqpoint{0.655045in}{1.127152in}}%
\pgfpathlineto{\pgfqpoint{0.663732in}{1.130356in}}%
\pgfpathlineto{\pgfqpoint{0.672382in}{1.133915in}}%
\pgfpathlineto{\pgfqpoint{0.656466in}{1.153746in}}%
\pgfpathlineto{\pgfqpoint{0.640558in}{1.173621in}}%
\pgfpathlineto{\pgfqpoint{0.631866in}{1.170053in}}%
\pgfpathlineto{\pgfqpoint{0.623136in}{1.166815in}}%
\pgfpathclose%
\pgfusepath{fill}%
\end{pgfscope}%
\begin{pgfscope}%
\pgfpathrectangle{\pgfqpoint{0.407592in}{0.222904in}}{\pgfqpoint{1.411163in}{1.411163in}}%
\pgfusepath{clip}%
\pgfsetbuttcap%
\pgfsetroundjoin%
\definecolor{currentfill}{rgb}{0.556863,0.654902,0.847059}%
\pgfsetfillcolor{currentfill}%
\pgfsetfillopacity{0.130000}%
\pgfsetlinewidth{0.000000pt}%
\definecolor{currentstroke}{rgb}{0.000000,0.000000,0.000000}%
\pgfsetstrokecolor{currentstroke}%
\pgfsetdash{}{0pt}%
\pgfpathmoveto{\pgfqpoint{1.169293in}{0.808200in}}%
\pgfpathlineto{\pgfqpoint{1.185453in}{0.807000in}}%
\pgfpathlineto{\pgfqpoint{1.201664in}{0.807384in}}%
\pgfpathlineto{\pgfqpoint{1.209513in}{0.819793in}}%
\pgfpathlineto{\pgfqpoint{1.217348in}{0.832744in}}%
\pgfpathlineto{\pgfqpoint{1.201151in}{0.832694in}}%
\pgfpathlineto{\pgfqpoint{1.185006in}{0.834104in}}%
\pgfpathlineto{\pgfqpoint{1.177157in}{0.820884in}}%
\pgfpathlineto{\pgfqpoint{1.169293in}{0.808200in}}%
\pgfpathclose%
\pgfusepath{fill}%
\end{pgfscope}%
\begin{pgfscope}%
\pgfpathrectangle{\pgfqpoint{0.407592in}{0.222904in}}{\pgfqpoint{1.411163in}{1.411163in}}%
\pgfusepath{clip}%
\pgfsetbuttcap%
\pgfsetroundjoin%
\definecolor{currentfill}{rgb}{0.556863,0.654902,0.847059}%
\pgfsetfillcolor{currentfill}%
\pgfsetfillopacity{0.130000}%
\pgfsetlinewidth{0.000000pt}%
\definecolor{currentstroke}{rgb}{0.000000,0.000000,0.000000}%
\pgfsetstrokecolor{currentstroke}%
\pgfsetdash{}{0pt}%
\pgfpathmoveto{\pgfqpoint{0.800025in}{0.977893in}}%
\pgfpathlineto{\pgfqpoint{0.816020in}{0.958917in}}%
\pgfpathlineto{\pgfqpoint{0.832024in}{0.940128in}}%
\pgfpathlineto{\pgfqpoint{0.840423in}{0.943959in}}%
\pgfpathlineto{\pgfqpoint{0.848786in}{0.948342in}}%
\pgfpathlineto{\pgfqpoint{0.832825in}{0.967154in}}%
\pgfpathlineto{\pgfqpoint{0.816873in}{0.986145in}}%
\pgfpathlineto{\pgfqpoint{0.808467in}{0.981770in}}%
\pgfpathlineto{\pgfqpoint{0.800025in}{0.977893in}}%
\pgfpathclose%
\pgfusepath{fill}%
\end{pgfscope}%
\begin{pgfscope}%
\pgfpathrectangle{\pgfqpoint{0.407592in}{0.222904in}}{\pgfqpoint{1.411163in}{1.411163in}}%
\pgfusepath{clip}%
\pgfsetbuttcap%
\pgfsetroundjoin%
\definecolor{currentfill}{rgb}{0.556863,0.654902,0.847059}%
\pgfsetfillcolor{currentfill}%
\pgfsetfillopacity{0.130000}%
\pgfsetlinewidth{0.000000pt}%
\definecolor{currentstroke}{rgb}{0.000000,0.000000,0.000000}%
\pgfsetstrokecolor{currentstroke}%
\pgfsetdash{}{0pt}%
\pgfpathmoveto{\pgfqpoint{1.392960in}{0.991706in}}%
\pgfpathlineto{\pgfqpoint{1.409513in}{0.996312in}}%
\pgfpathlineto{\pgfqpoint{1.426133in}{1.001563in}}%
\pgfpathlineto{\pgfqpoint{1.433770in}{1.015845in}}%
\pgfpathlineto{\pgfqpoint{1.441394in}{1.030305in}}%
\pgfpathlineto{\pgfqpoint{1.424789in}{1.025511in}}%
\pgfpathlineto{\pgfqpoint{1.408250in}{1.021348in}}%
\pgfpathlineto{\pgfqpoint{1.400611in}{1.006448in}}%
\pgfpathlineto{\pgfqpoint{1.392960in}{0.991706in}}%
\pgfpathclose%
\pgfusepath{fill}%
\end{pgfscope}%
\begin{pgfscope}%
\pgfpathrectangle{\pgfqpoint{0.407592in}{0.222904in}}{\pgfqpoint{1.411163in}{1.411163in}}%
\pgfusepath{clip}%
\pgfsetbuttcap%
\pgfsetroundjoin%
\definecolor{currentfill}{rgb}{0.556863,0.654902,0.847059}%
\pgfsetfillcolor{currentfill}%
\pgfsetfillopacity{0.130000}%
\pgfsetlinewidth{0.000000pt}%
\definecolor{currentstroke}{rgb}{0.000000,0.000000,0.000000}%
\pgfsetstrokecolor{currentstroke}%
\pgfsetdash{}{0pt}%
\pgfpathmoveto{\pgfqpoint{1.025135in}{0.793661in}}%
\pgfpathlineto{\pgfqpoint{1.041175in}{0.781692in}}%
\pgfpathlineto{\pgfqpoint{1.057232in}{0.771245in}}%
\pgfpathlineto{\pgfqpoint{1.065265in}{0.778684in}}%
\pgfpathlineto{\pgfqpoint{1.073272in}{0.787145in}}%
\pgfpathlineto{\pgfqpoint{1.057245in}{0.797326in}}%
\pgfpathlineto{\pgfqpoint{1.041239in}{0.808995in}}%
\pgfpathlineto{\pgfqpoint{1.033201in}{0.800859in}}%
\pgfpathlineto{\pgfqpoint{1.025135in}{0.793661in}}%
\pgfpathclose%
\pgfusepath{fill}%
\end{pgfscope}%
\begin{pgfscope}%
\pgfpathrectangle{\pgfqpoint{0.407592in}{0.222904in}}{\pgfqpoint{1.411163in}{1.411163in}}%
\pgfusepath{clip}%
\pgfsetbuttcap%
\pgfsetroundjoin%
\definecolor{currentfill}{rgb}{0.556863,0.654902,0.847059}%
\pgfsetfillcolor{currentfill}%
\pgfsetfillopacity{0.130000}%
\pgfsetlinewidth{0.000000pt}%
\definecolor{currentstroke}{rgb}{0.000000,0.000000,0.000000}%
\pgfsetstrokecolor{currentstroke}%
\pgfsetdash{}{0pt}%
\pgfpathmoveto{\pgfqpoint{0.655045in}{1.127152in}}%
\pgfpathlineto{\pgfqpoint{0.671013in}{1.107403in}}%
\pgfpathlineto{\pgfqpoint{0.686990in}{1.087719in}}%
\pgfpathlineto{\pgfqpoint{0.695634in}{1.090875in}}%
\pgfpathlineto{\pgfqpoint{0.704240in}{1.094413in}}%
\pgfpathlineto{\pgfqpoint{0.688307in}{1.114135in}}%
\pgfpathlineto{\pgfqpoint{0.672382in}{1.133915in}}%
\pgfpathlineto{\pgfqpoint{0.663732in}{1.130356in}}%
\pgfpathlineto{\pgfqpoint{0.655045in}{1.127152in}}%
\pgfpathclose%
\pgfusepath{fill}%
\end{pgfscope}%
\begin{pgfscope}%
\pgfpathrectangle{\pgfqpoint{0.407592in}{0.222904in}}{\pgfqpoint{1.411163in}{1.411163in}}%
\pgfusepath{clip}%
\pgfsetbuttcap%
\pgfsetroundjoin%
\definecolor{currentfill}{rgb}{0.556863,0.654902,0.847059}%
\pgfsetfillcolor{currentfill}%
\pgfsetfillopacity{0.130000}%
\pgfsetlinewidth{0.000000pt}%
\definecolor{currentstroke}{rgb}{0.000000,0.000000,0.000000}%
\pgfsetstrokecolor{currentstroke}%
\pgfsetdash{}{0pt}%
\pgfpathmoveto{\pgfqpoint{0.944742in}{0.841599in}}%
\pgfpathlineto{\pgfqpoint{0.960768in}{0.825461in}}%
\pgfpathlineto{\pgfqpoint{0.976805in}{0.810076in}}%
\pgfpathlineto{\pgfqpoint{0.984970in}{0.815243in}}%
\pgfpathlineto{\pgfqpoint{0.993102in}{0.821299in}}%
\pgfpathlineto{\pgfqpoint{0.977104in}{0.836576in}}%
\pgfpathlineto{\pgfqpoint{0.961118in}{0.852608in}}%
\pgfpathlineto{\pgfqpoint{0.952947in}{0.846712in}}%
\pgfpathlineto{\pgfqpoint{0.944742in}{0.841599in}}%
\pgfpathclose%
\pgfusepath{fill}%
\end{pgfscope}%
\begin{pgfscope}%
\pgfpathrectangle{\pgfqpoint{0.407592in}{0.222904in}}{\pgfqpoint{1.411163in}{1.411163in}}%
\pgfusepath{clip}%
\pgfsetbuttcap%
\pgfsetroundjoin%
\definecolor{currentfill}{rgb}{0.556863,0.654902,0.847059}%
\pgfsetfillcolor{currentfill}%
\pgfsetfillopacity{0.130000}%
\pgfsetlinewidth{0.000000pt}%
\definecolor{currentstroke}{rgb}{0.000000,0.000000,0.000000}%
\pgfsetstrokecolor{currentstroke}%
\pgfsetdash{}{0pt}%
\pgfpathmoveto{\pgfqpoint{1.217348in}{0.832744in}}%
\pgfpathlineto{\pgfqpoint{1.233599in}{0.834165in}}%
\pgfpathlineto{\pgfqpoint{1.249907in}{0.836852in}}%
\pgfpathlineto{\pgfqpoint{1.257714in}{0.849820in}}%
\pgfpathlineto{\pgfqpoint{1.265506in}{0.863214in}}%
\pgfpathlineto{\pgfqpoint{1.249211in}{0.861012in}}%
\pgfpathlineto{\pgfqpoint{1.232974in}{0.860001in}}%
\pgfpathlineto{\pgfqpoint{1.225168in}{0.846167in}}%
\pgfpathlineto{\pgfqpoint{1.217348in}{0.832744in}}%
\pgfpathclose%
\pgfusepath{fill}%
\end{pgfscope}%
\begin{pgfscope}%
\pgfpathrectangle{\pgfqpoint{0.407592in}{0.222904in}}{\pgfqpoint{1.411163in}{1.411163in}}%
\pgfusepath{clip}%
\pgfsetbuttcap%
\pgfsetroundjoin%
\definecolor{currentfill}{rgb}{0.556863,0.654902,0.847059}%
\pgfsetfillcolor{currentfill}%
\pgfsetfillopacity{0.130000}%
\pgfsetlinewidth{0.000000pt}%
\definecolor{currentstroke}{rgb}{0.000000,0.000000,0.000000}%
\pgfsetstrokecolor{currentstroke}%
\pgfsetdash{}{0pt}%
\pgfpathmoveto{\pgfqpoint{0.832024in}{0.940128in}}%
\pgfpathlineto{\pgfqpoint{0.848036in}{0.921557in}}%
\pgfpathlineto{\pgfqpoint{0.864057in}{0.903246in}}%
\pgfpathlineto{\pgfqpoint{0.872414in}{0.907010in}}%
\pgfpathlineto{\pgfqpoint{0.880735in}{0.911389in}}%
\pgfpathlineto{\pgfqpoint{0.864756in}{0.929741in}}%
\pgfpathlineto{\pgfqpoint{0.848786in}{0.948342in}}%
\pgfpathlineto{\pgfqpoint{0.840423in}{0.943959in}}%
\pgfpathlineto{\pgfqpoint{0.832024in}{0.940128in}}%
\pgfpathclose%
\pgfusepath{fill}%
\end{pgfscope}%
\begin{pgfscope}%
\pgfpathrectangle{\pgfqpoint{0.407592in}{0.222904in}}{\pgfqpoint{1.411163in}{1.411163in}}%
\pgfusepath{clip}%
\pgfsetbuttcap%
\pgfsetroundjoin%
\definecolor{currentfill}{rgb}{0.556863,0.654902,0.847059}%
\pgfsetfillcolor{currentfill}%
\pgfsetfillopacity{0.130000}%
\pgfsetlinewidth{0.000000pt}%
\definecolor{currentstroke}{rgb}{0.000000,0.000000,0.000000}%
\pgfsetstrokecolor{currentstroke}%
\pgfsetdash{}{0pt}%
\pgfpathmoveto{\pgfqpoint{1.329263in}{0.925705in}}%
\pgfpathlineto{\pgfqpoint{1.345716in}{0.929810in}}%
\pgfpathlineto{\pgfqpoint{1.362233in}{0.934699in}}%
\pgfpathlineto{\pgfqpoint{1.369933in}{0.948613in}}%
\pgfpathlineto{\pgfqpoint{1.377621in}{0.962768in}}%
\pgfpathlineto{\pgfqpoint{1.361118in}{0.958387in}}%
\pgfpathlineto{\pgfqpoint{1.344679in}{0.954769in}}%
\pgfpathlineto{\pgfqpoint{1.336977in}{0.940127in}}%
\pgfpathlineto{\pgfqpoint{1.329263in}{0.925705in}}%
\pgfpathclose%
\pgfusepath{fill}%
\end{pgfscope}%
\begin{pgfscope}%
\pgfpathrectangle{\pgfqpoint{0.407592in}{0.222904in}}{\pgfqpoint{1.411163in}{1.411163in}}%
\pgfusepath{clip}%
\pgfsetbuttcap%
\pgfsetroundjoin%
\definecolor{currentfill}{rgb}{0.556863,0.654902,0.847059}%
\pgfsetfillcolor{currentfill}%
\pgfsetfillopacity{0.130000}%
\pgfsetlinewidth{0.000000pt}%
\definecolor{currentstroke}{rgb}{0.000000,0.000000,0.000000}%
\pgfsetstrokecolor{currentstroke}%
\pgfsetdash{}{0pt}%
\pgfpathmoveto{\pgfqpoint{0.686990in}{1.087719in}}%
\pgfpathlineto{\pgfqpoint{0.702975in}{1.068110in}}%
\pgfpathlineto{\pgfqpoint{0.718969in}{1.048586in}}%
\pgfpathlineto{\pgfqpoint{0.727570in}{1.051676in}}%
\pgfpathlineto{\pgfqpoint{0.736134in}{1.055179in}}%
\pgfpathlineto{\pgfqpoint{0.720183in}{1.074758in}}%
\pgfpathlineto{\pgfqpoint{0.704240in}{1.094413in}}%
\pgfpathlineto{\pgfqpoint{0.695634in}{1.090875in}}%
\pgfpathlineto{\pgfqpoint{0.686990in}{1.087719in}}%
\pgfpathclose%
\pgfusepath{fill}%
\end{pgfscope}%
\begin{pgfscope}%
\pgfpathrectangle{\pgfqpoint{0.407592in}{0.222904in}}{\pgfqpoint{1.411163in}{1.411163in}}%
\pgfusepath{clip}%
\pgfsetbuttcap%
\pgfsetroundjoin%
\definecolor{currentfill}{rgb}{0.556863,0.654902,0.847059}%
\pgfsetfillcolor{currentfill}%
\pgfsetfillopacity{0.130000}%
\pgfsetlinewidth{0.000000pt}%
\definecolor{currentstroke}{rgb}{0.000000,0.000000,0.000000}%
\pgfsetstrokecolor{currentstroke}%
\pgfsetdash{}{0pt}%
\pgfpathmoveto{\pgfqpoint{1.505121in}{1.099259in}}%
\pgfpathlineto{\pgfqpoint{1.521896in}{1.104863in}}%
\pgfpathlineto{\pgfqpoint{1.538740in}{1.110948in}}%
\pgfpathlineto{\pgfqpoint{1.546274in}{1.125305in}}%
\pgfpathlineto{\pgfqpoint{1.529438in}{1.119417in}}%
\pgfpathlineto{\pgfqpoint{1.512671in}{1.114007in}}%
\pgfusepath{fill}%
\end{pgfscope}%
\begin{pgfscope}%
\pgfpathrectangle{\pgfqpoint{0.407592in}{0.222904in}}{\pgfqpoint{1.411163in}{1.411163in}}%
\pgfusepath{clip}%
\pgfsetbuttcap%
\pgfsetroundjoin%
\definecolor{currentfill}{rgb}{0.556863,0.654902,0.847059}%
\pgfsetfillcolor{currentfill}%
\pgfsetfillopacity{0.130000}%
\pgfsetlinewidth{0.000000pt}%
\definecolor{currentstroke}{rgb}{0.000000,0.000000,0.000000}%
\pgfsetstrokecolor{currentstroke}%
\pgfsetdash{}{0pt}%
\pgfpathmoveto{\pgfqpoint{0.864057in}{0.903246in}}%
\pgfpathlineto{\pgfqpoint{0.880087in}{0.885241in}}%
\pgfpathlineto{\pgfqpoint{0.896126in}{0.867602in}}%
\pgfpathlineto{\pgfqpoint{0.904440in}{0.871272in}}%
\pgfpathlineto{\pgfqpoint{0.912720in}{0.875633in}}%
\pgfpathlineto{\pgfqpoint{0.896723in}{0.893334in}}%
\pgfpathlineto{\pgfqpoint{0.880735in}{0.911389in}}%
\pgfpathlineto{\pgfqpoint{0.872414in}{0.907010in}}%
\pgfpathlineto{\pgfqpoint{0.864057in}{0.903246in}}%
\pgfpathclose%
\pgfusepath{fill}%
\end{pgfscope}%
\begin{pgfscope}%
\pgfpathrectangle{\pgfqpoint{0.407592in}{0.222904in}}{\pgfqpoint{1.411163in}{1.411163in}}%
\pgfusepath{clip}%
\pgfsetbuttcap%
\pgfsetroundjoin%
\definecolor{currentfill}{rgb}{0.556863,0.654902,0.847059}%
\pgfsetfillcolor{currentfill}%
\pgfsetfillopacity{0.130000}%
\pgfsetlinewidth{0.000000pt}%
\definecolor{currentstroke}{rgb}{0.000000,0.000000,0.000000}%
\pgfsetstrokecolor{currentstroke}%
\pgfsetdash{}{0pt}%
\pgfpathmoveto{\pgfqpoint{1.441394in}{1.030305in}}%
\pgfpathlineto{\pgfqpoint{1.458065in}{1.035690in}}%
\pgfpathlineto{\pgfqpoint{1.474805in}{1.041628in}}%
\pgfpathlineto{\pgfqpoint{1.482401in}{1.055806in}}%
\pgfpathlineto{\pgfqpoint{1.489986in}{1.070146in}}%
\pgfpathlineto{\pgfqpoint{1.473263in}{1.064645in}}%
\pgfpathlineto{\pgfqpoint{1.456607in}{1.059690in}}%
\pgfpathlineto{\pgfqpoint{1.449006in}{1.044924in}}%
\pgfpathlineto{\pgfqpoint{1.441394in}{1.030305in}}%
\pgfpathclose%
\pgfusepath{fill}%
\end{pgfscope}%
\begin{pgfscope}%
\pgfpathrectangle{\pgfqpoint{0.407592in}{0.222904in}}{\pgfqpoint{1.411163in}{1.411163in}}%
\pgfusepath{clip}%
\pgfsetbuttcap%
\pgfsetroundjoin%
\definecolor{currentfill}{rgb}{0.556863,0.654902,0.847059}%
\pgfsetfillcolor{currentfill}%
\pgfsetfillopacity{0.130000}%
\pgfsetlinewidth{0.000000pt}%
\definecolor{currentstroke}{rgb}{0.000000,0.000000,0.000000}%
\pgfsetstrokecolor{currentstroke}%
\pgfsetdash{}{0pt}%
\pgfpathmoveto{\pgfqpoint{1.265506in}{0.863214in}}%
\pgfpathlineto{\pgfqpoint{1.281862in}{0.866511in}}%
\pgfpathlineto{\pgfqpoint{1.298279in}{0.870807in}}%
\pgfpathlineto{\pgfqpoint{1.306044in}{0.884045in}}%
\pgfpathlineto{\pgfqpoint{1.313797in}{0.897632in}}%
\pgfpathlineto{\pgfqpoint{1.297393in}{0.893877in}}%
\pgfpathlineto{\pgfqpoint{1.281052in}{0.891084in}}%
\pgfpathlineto{\pgfqpoint{1.273286in}{0.876984in}}%
\pgfpathlineto{\pgfqpoint{1.265506in}{0.863214in}}%
\pgfpathclose%
\pgfusepath{fill}%
\end{pgfscope}%
\begin{pgfscope}%
\pgfpathrectangle{\pgfqpoint{0.407592in}{0.222904in}}{\pgfqpoint{1.411163in}{1.411163in}}%
\pgfusepath{clip}%
\pgfsetbuttcap%
\pgfsetroundjoin%
\definecolor{currentfill}{rgb}{0.556863,0.654902,0.847059}%
\pgfsetfillcolor{currentfill}%
\pgfsetfillopacity{0.130000}%
\pgfsetlinewidth{0.000000pt}%
\definecolor{currentstroke}{rgb}{0.000000,0.000000,0.000000}%
\pgfsetstrokecolor{currentstroke}%
\pgfsetdash{}{0pt}%
\pgfpathmoveto{\pgfqpoint{0.718969in}{1.048586in}}%
\pgfpathlineto{\pgfqpoint{0.734972in}{1.029159in}}%
\pgfpathlineto{\pgfqpoint{0.750983in}{1.009845in}}%
\pgfpathlineto{\pgfqpoint{0.759541in}{1.012848in}}%
\pgfpathlineto{\pgfqpoint{0.768062in}{1.016299in}}%
\pgfpathlineto{\pgfqpoint{0.752094in}{1.035688in}}%
\pgfpathlineto{\pgfqpoint{0.736134in}{1.055179in}}%
\pgfpathlineto{\pgfqpoint{0.727570in}{1.051676in}}%
\pgfpathlineto{\pgfqpoint{0.718969in}{1.048586in}}%
\pgfpathclose%
\pgfusepath{fill}%
\end{pgfscope}%
\begin{pgfscope}%
\pgfpathrectangle{\pgfqpoint{0.407592in}{0.222904in}}{\pgfqpoint{1.411163in}{1.411163in}}%
\pgfusepath{clip}%
\pgfsetbuttcap%
\pgfsetroundjoin%
\definecolor{currentfill}{rgb}{0.556863,0.654902,0.847059}%
\pgfsetfillcolor{currentfill}%
\pgfsetfillopacity{0.130000}%
\pgfsetlinewidth{0.000000pt}%
\definecolor{currentstroke}{rgb}{0.000000,0.000000,0.000000}%
\pgfsetstrokecolor{currentstroke}%
\pgfsetdash{}{0pt}%
\pgfpathmoveto{\pgfqpoint{0.976805in}{0.810076in}}%
\pgfpathlineto{\pgfqpoint{0.992853in}{0.795603in}}%
\pgfpathlineto{\pgfqpoint{1.008914in}{0.782227in}}%
\pgfpathlineto{\pgfqpoint{1.017040in}{0.787443in}}%
\pgfpathlineto{\pgfqpoint{1.025135in}{0.793661in}}%
\pgfpathlineto{\pgfqpoint{1.009112in}{0.806933in}}%
\pgfpathlineto{\pgfqpoint{0.993102in}{0.821299in}}%
\pgfpathlineto{\pgfqpoint{0.984970in}{0.815243in}}%
\pgfpathlineto{\pgfqpoint{0.976805in}{0.810076in}}%
\pgfpathclose%
\pgfusepath{fill}%
\end{pgfscope}%
\begin{pgfscope}%
\pgfpathrectangle{\pgfqpoint{0.407592in}{0.222904in}}{\pgfqpoint{1.411163in}{1.411163in}}%
\pgfusepath{clip}%
\pgfsetbuttcap%
\pgfsetroundjoin%
\definecolor{currentfill}{rgb}{0.556863,0.654902,0.847059}%
\pgfsetfillcolor{currentfill}%
\pgfsetfillopacity{0.130000}%
\pgfsetlinewidth{0.000000pt}%
\definecolor{currentstroke}{rgb}{0.000000,0.000000,0.000000}%
\pgfsetstrokecolor{currentstroke}%
\pgfsetdash{}{0pt}%
\pgfpathmoveto{\pgfqpoint{1.105400in}{0.771992in}}%
\pgfpathlineto{\pgfqpoint{1.121512in}{0.767285in}}%
\pgfpathlineto{\pgfqpoint{1.137663in}{0.764555in}}%
\pgfpathlineto{\pgfqpoint{1.145599in}{0.774216in}}%
\pgfpathlineto{\pgfqpoint{1.153515in}{0.784774in}}%
\pgfpathlineto{\pgfqpoint{1.137384in}{0.787564in}}%
\pgfpathlineto{\pgfqpoint{1.121293in}{0.792199in}}%
\pgfpathlineto{\pgfqpoint{1.113357in}{0.781646in}}%
\pgfpathlineto{\pgfqpoint{1.105400in}{0.771992in}}%
\pgfpathclose%
\pgfusepath{fill}%
\end{pgfscope}%
\begin{pgfscope}%
\pgfpathrectangle{\pgfqpoint{0.407592in}{0.222904in}}{\pgfqpoint{1.411163in}{1.411163in}}%
\pgfusepath{clip}%
\pgfsetbuttcap%
\pgfsetroundjoin%
\definecolor{currentfill}{rgb}{0.556863,0.654902,0.847059}%
\pgfsetfillcolor{currentfill}%
\pgfsetfillopacity{0.130000}%
\pgfsetlinewidth{0.000000pt}%
\definecolor{currentstroke}{rgb}{0.000000,0.000000,0.000000}%
\pgfsetstrokecolor{currentstroke}%
\pgfsetdash{}{0pt}%
\pgfpathmoveto{\pgfqpoint{1.153515in}{0.784774in}}%
\pgfpathlineto{\pgfqpoint{1.169691in}{0.783786in}}%
\pgfpathlineto{\pgfqpoint{1.185918in}{0.784499in}}%
\pgfpathlineto{\pgfqpoint{1.193799in}{0.795592in}}%
\pgfpathlineto{\pgfqpoint{1.201664in}{0.807384in}}%
\pgfpathlineto{\pgfqpoint{1.185453in}{0.807000in}}%
\pgfpathlineto{\pgfqpoint{1.169293in}{0.808200in}}%
\pgfpathlineto{\pgfqpoint{1.161412in}{0.796133in}}%
\pgfpathlineto{\pgfqpoint{1.153515in}{0.784774in}}%
\pgfpathclose%
\pgfusepath{fill}%
\end{pgfscope}%
\begin{pgfscope}%
\pgfpathrectangle{\pgfqpoint{0.407592in}{0.222904in}}{\pgfqpoint{1.411163in}{1.411163in}}%
\pgfusepath{clip}%
\pgfsetbuttcap%
\pgfsetroundjoin%
\definecolor{currentfill}{rgb}{0.556863,0.654902,0.847059}%
\pgfsetfillcolor{currentfill}%
\pgfsetfillopacity{0.130000}%
\pgfsetlinewidth{0.000000pt}%
\definecolor{currentstroke}{rgb}{0.000000,0.000000,0.000000}%
\pgfsetstrokecolor{currentstroke}%
\pgfsetdash{}{0pt}%
\pgfpathmoveto{\pgfqpoint{1.057232in}{0.771245in}}%
\pgfpathlineto{\pgfqpoint{1.073312in}{0.762532in}}%
\pgfpathlineto{\pgfqpoint{1.089418in}{0.755736in}}%
\pgfpathlineto{\pgfqpoint{1.097421in}{0.763328in}}%
\pgfpathlineto{\pgfqpoint{1.105400in}{0.771992in}}%
\pgfpathlineto{\pgfqpoint{1.089322in}{0.778647in}}%
\pgfpathlineto{\pgfqpoint{1.073272in}{0.787145in}}%
\pgfpathlineto{\pgfqpoint{1.065265in}{0.778684in}}%
\pgfpathlineto{\pgfqpoint{1.057232in}{0.771245in}}%
\pgfpathclose%
\pgfusepath{fill}%
\end{pgfscope}%
\begin{pgfscope}%
\pgfpathrectangle{\pgfqpoint{0.407592in}{0.222904in}}{\pgfqpoint{1.411163in}{1.411163in}}%
\pgfusepath{clip}%
\pgfsetbuttcap%
\pgfsetroundjoin%
\definecolor{currentfill}{rgb}{0.556863,0.654902,0.847059}%
\pgfsetfillcolor{currentfill}%
\pgfsetfillopacity{0.130000}%
\pgfsetlinewidth{0.000000pt}%
\definecolor{currentstroke}{rgb}{0.000000,0.000000,0.000000}%
\pgfsetstrokecolor{currentstroke}%
\pgfsetdash{}{0pt}%
\pgfpathmoveto{\pgfqpoint{1.377621in}{0.962768in}}%
\pgfpathlineto{\pgfqpoint{1.394190in}{0.967857in}}%
\pgfpathlineto{\pgfqpoint{1.410825in}{0.973598in}}%
\pgfpathlineto{\pgfqpoint{1.418485in}{0.987474in}}%
\pgfpathlineto{\pgfqpoint{1.426133in}{1.001563in}}%
\pgfpathlineto{\pgfqpoint{1.409513in}{0.996312in}}%
\pgfpathlineto{\pgfqpoint{1.392960in}{0.991706in}}%
\pgfpathlineto{\pgfqpoint{1.385296in}{0.977140in}}%
\pgfpathlineto{\pgfqpoint{1.377621in}{0.962768in}}%
\pgfpathclose%
\pgfusepath{fill}%
\end{pgfscope}%
\begin{pgfscope}%
\pgfpathrectangle{\pgfqpoint{0.407592in}{0.222904in}}{\pgfqpoint{1.411163in}{1.411163in}}%
\pgfusepath{clip}%
\pgfsetbuttcap%
\pgfsetroundjoin%
\definecolor{currentfill}{rgb}{0.556863,0.654902,0.847059}%
\pgfsetfillcolor{currentfill}%
\pgfsetfillopacity{0.130000}%
\pgfsetlinewidth{0.000000pt}%
\definecolor{currentstroke}{rgb}{0.000000,0.000000,0.000000}%
\pgfsetstrokecolor{currentstroke}%
\pgfsetdash{}{0pt}%
\pgfpathmoveto{\pgfqpoint{0.750983in}{1.009845in}}%
\pgfpathlineto{\pgfqpoint{0.767003in}{0.990659in}}%
\pgfpathlineto{\pgfqpoint{0.783032in}{0.971623in}}%
\pgfpathlineto{\pgfqpoint{0.791547in}{0.974513in}}%
\pgfpathlineto{\pgfqpoint{0.800025in}{0.977893in}}%
\pgfpathlineto{\pgfqpoint{0.784039in}{0.997028in}}%
\pgfpathlineto{\pgfqpoint{0.768062in}{1.016299in}}%
\pgfpathlineto{\pgfqpoint{0.759541in}{1.012848in}}%
\pgfpathlineto{\pgfqpoint{0.750983in}{1.009845in}}%
\pgfpathclose%
\pgfusepath{fill}%
\end{pgfscope}%
\begin{pgfscope}%
\pgfpathrectangle{\pgfqpoint{0.407592in}{0.222904in}}{\pgfqpoint{1.411163in}{1.411163in}}%
\pgfusepath{clip}%
\pgfsetbuttcap%
\pgfsetroundjoin%
\definecolor{currentfill}{rgb}{0.556863,0.654902,0.847059}%
\pgfsetfillcolor{currentfill}%
\pgfsetfillopacity{0.130000}%
\pgfsetlinewidth{0.000000pt}%
\definecolor{currentstroke}{rgb}{0.000000,0.000000,0.000000}%
\pgfsetstrokecolor{currentstroke}%
\pgfsetdash{}{0pt}%
\pgfpathmoveto{\pgfqpoint{0.896126in}{0.867602in}}%
\pgfpathlineto{\pgfqpoint{0.912173in}{0.850400in}}%
\pgfpathlineto{\pgfqpoint{0.928231in}{0.833726in}}%
\pgfpathlineto{\pgfqpoint{0.936504in}{0.837272in}}%
\pgfpathlineto{\pgfqpoint{0.944742in}{0.841599in}}%
\pgfpathlineto{\pgfqpoint{0.928726in}{0.858359in}}%
\pgfpathlineto{\pgfqpoint{0.912720in}{0.875633in}}%
\pgfpathlineto{\pgfqpoint{0.904440in}{0.871272in}}%
\pgfpathlineto{\pgfqpoint{0.896126in}{0.867602in}}%
\pgfpathclose%
\pgfusepath{fill}%
\end{pgfscope}%
\begin{pgfscope}%
\pgfpathrectangle{\pgfqpoint{0.407592in}{0.222904in}}{\pgfqpoint{1.411163in}{1.411163in}}%
\pgfusepath{clip}%
\pgfsetbuttcap%
\pgfsetroundjoin%
\definecolor{currentfill}{rgb}{0.556863,0.654902,0.847059}%
\pgfsetfillcolor{currentfill}%
\pgfsetfillopacity{0.130000}%
\pgfsetlinewidth{0.000000pt}%
\definecolor{currentstroke}{rgb}{0.000000,0.000000,0.000000}%
\pgfsetstrokecolor{currentstroke}%
\pgfsetdash{}{0pt}%
\pgfpathmoveto{\pgfqpoint{1.201664in}{0.807384in}}%
\pgfpathlineto{\pgfqpoint{1.217929in}{0.809237in}}%
\pgfpathlineto{\pgfqpoint{1.234253in}{0.812426in}}%
\pgfpathlineto{\pgfqpoint{1.242087in}{0.824367in}}%
\pgfpathlineto{\pgfqpoint{1.249907in}{0.836852in}}%
\pgfpathlineto{\pgfqpoint{1.233599in}{0.834165in}}%
\pgfpathlineto{\pgfqpoint{1.217348in}{0.832744in}}%
\pgfpathlineto{\pgfqpoint{1.209513in}{0.819793in}}%
\pgfpathlineto{\pgfqpoint{1.201664in}{0.807384in}}%
\pgfpathclose%
\pgfusepath{fill}%
\end{pgfscope}%
\begin{pgfscope}%
\pgfpathrectangle{\pgfqpoint{0.407592in}{0.222904in}}{\pgfqpoint{1.411163in}{1.411163in}}%
\pgfusepath{clip}%
\pgfsetbuttcap%
\pgfsetroundjoin%
\definecolor{currentfill}{rgb}{0.556863,0.654902,0.847059}%
\pgfsetfillcolor{currentfill}%
\pgfsetfillopacity{0.130000}%
\pgfsetlinewidth{0.000000pt}%
\definecolor{currentstroke}{rgb}{0.000000,0.000000,0.000000}%
\pgfsetstrokecolor{currentstroke}%
\pgfsetdash{}{0pt}%
\pgfpathmoveto{\pgfqpoint{1.313797in}{0.897632in}}%
\pgfpathlineto{\pgfqpoint{1.330264in}{0.902265in}}%
\pgfpathlineto{\pgfqpoint{1.346797in}{0.907699in}}%
\pgfpathlineto{\pgfqpoint{1.354521in}{0.921051in}}%
\pgfpathlineto{\pgfqpoint{1.362233in}{0.934699in}}%
\pgfpathlineto{\pgfqpoint{1.345716in}{0.929810in}}%
\pgfpathlineto{\pgfqpoint{1.329263in}{0.925705in}}%
\pgfpathlineto{\pgfqpoint{1.321536in}{0.911530in}}%
\pgfpathlineto{\pgfqpoint{1.313797in}{0.897632in}}%
\pgfpathclose%
\pgfusepath{fill}%
\end{pgfscope}%
\begin{pgfscope}%
\pgfpathrectangle{\pgfqpoint{0.407592in}{0.222904in}}{\pgfqpoint{1.411163in}{1.411163in}}%
\pgfusepath{clip}%
\pgfsetbuttcap%
\pgfsetroundjoin%
\definecolor{currentfill}{rgb}{0.556863,0.654902,0.847059}%
\pgfsetfillcolor{currentfill}%
\pgfsetfillopacity{0.130000}%
\pgfsetlinewidth{0.000000pt}%
\definecolor{currentstroke}{rgb}{0.000000,0.000000,0.000000}%
\pgfsetstrokecolor{currentstroke}%
\pgfsetdash{}{0pt}%
\pgfpathmoveto{\pgfqpoint{0.605563in}{1.161323in}}%
\pgfpathlineto{\pgfqpoint{0.621556in}{1.141530in}}%
\pgfpathlineto{\pgfqpoint{0.637558in}{1.121797in}}%
\pgfpathlineto{\pgfqpoint{0.646321in}{1.124300in}}%
\pgfpathlineto{\pgfqpoint{0.655045in}{1.127152in}}%
\pgfpathlineto{\pgfqpoint{0.639086in}{1.146958in}}%
\pgfpathlineto{\pgfqpoint{0.623136in}{1.166815in}}%
\pgfpathlineto{\pgfqpoint{0.614368in}{1.163906in}}%
\pgfpathlineto{\pgfqpoint{0.605563in}{1.161323in}}%
\pgfpathclose%
\pgfusepath{fill}%
\end{pgfscope}%
\begin{pgfscope}%
\pgfpathrectangle{\pgfqpoint{0.407592in}{0.222904in}}{\pgfqpoint{1.411163in}{1.411163in}}%
\pgfusepath{clip}%
\pgfsetbuttcap%
\pgfsetroundjoin%
\definecolor{currentfill}{rgb}{0.556863,0.654902,0.847059}%
\pgfsetfillcolor{currentfill}%
\pgfsetfillopacity{0.130000}%
\pgfsetlinewidth{0.000000pt}%
\definecolor{currentstroke}{rgb}{0.000000,0.000000,0.000000}%
\pgfsetstrokecolor{currentstroke}%
\pgfsetdash{}{0pt}%
\pgfpathmoveto{\pgfqpoint{1.489986in}{1.070146in}}%
\pgfpathlineto{\pgfqpoint{1.506778in}{1.076162in}}%
\pgfpathlineto{\pgfqpoint{1.523638in}{1.082659in}}%
\pgfpathlineto{\pgfqpoint{1.531195in}{1.096729in}}%
\pgfpathlineto{\pgfqpoint{1.538740in}{1.110948in}}%
\pgfpathlineto{\pgfqpoint{1.521896in}{1.104863in}}%
\pgfpathlineto{\pgfqpoint{1.505121in}{1.099259in}}%
\pgfpathlineto{\pgfqpoint{1.497560in}{1.084635in}}%
\pgfpathlineto{\pgfqpoint{1.489986in}{1.070146in}}%
\pgfpathclose%
\pgfusepath{fill}%
\end{pgfscope}%
\begin{pgfscope}%
\pgfpathrectangle{\pgfqpoint{0.407592in}{0.222904in}}{\pgfqpoint{1.411163in}{1.411163in}}%
\pgfusepath{clip}%
\pgfsetbuttcap%
\pgfsetroundjoin%
\definecolor{currentfill}{rgb}{0.556863,0.654902,0.847059}%
\pgfsetfillcolor{currentfill}%
\pgfsetfillopacity{0.130000}%
\pgfsetlinewidth{0.000000pt}%
\definecolor{currentstroke}{rgb}{0.000000,0.000000,0.000000}%
\pgfsetstrokecolor{currentstroke}%
\pgfsetdash{}{0pt}%
\pgfpathmoveto{\pgfqpoint{0.783032in}{0.971623in}}%
\pgfpathlineto{\pgfqpoint{0.799069in}{0.952761in}}%
\pgfpathlineto{\pgfqpoint{0.815115in}{0.934101in}}%
\pgfpathlineto{\pgfqpoint{0.823588in}{0.936844in}}%
\pgfpathlineto{\pgfqpoint{0.832024in}{0.940128in}}%
\pgfpathlineto{\pgfqpoint{0.816020in}{0.958917in}}%
\pgfpathlineto{\pgfqpoint{0.800025in}{0.977893in}}%
\pgfpathlineto{\pgfqpoint{0.791547in}{0.974513in}}%
\pgfpathlineto{\pgfqpoint{0.783032in}{0.971623in}}%
\pgfpathclose%
\pgfusepath{fill}%
\end{pgfscope}%
\begin{pgfscope}%
\pgfpathrectangle{\pgfqpoint{0.407592in}{0.222904in}}{\pgfqpoint{1.411163in}{1.411163in}}%
\pgfusepath{clip}%
\pgfsetbuttcap%
\pgfsetroundjoin%
\definecolor{currentfill}{rgb}{0.556863,0.654902,0.847059}%
\pgfsetfillcolor{currentfill}%
\pgfsetfillopacity{0.130000}%
\pgfsetlinewidth{0.000000pt}%
\definecolor{currentstroke}{rgb}{0.000000,0.000000,0.000000}%
\pgfsetstrokecolor{currentstroke}%
\pgfsetdash{}{0pt}%
\pgfpathmoveto{\pgfqpoint{1.008914in}{0.782227in}}%
\pgfpathlineto{\pgfqpoint{1.024989in}{0.770158in}}%
\pgfpathlineto{\pgfqpoint{1.041082in}{0.759620in}}%
\pgfpathlineto{\pgfqpoint{1.049172in}{0.764880in}}%
\pgfpathlineto{\pgfqpoint{1.057232in}{0.771245in}}%
\pgfpathlineto{\pgfqpoint{1.041175in}{0.781692in}}%
\pgfpathlineto{\pgfqpoint{1.025135in}{0.793661in}}%
\pgfpathlineto{\pgfqpoint{1.017040in}{0.787443in}}%
\pgfpathlineto{\pgfqpoint{1.008914in}{0.782227in}}%
\pgfpathclose%
\pgfusepath{fill}%
\end{pgfscope}%
\begin{pgfscope}%
\pgfpathrectangle{\pgfqpoint{0.407592in}{0.222904in}}{\pgfqpoint{1.411163in}{1.411163in}}%
\pgfusepath{clip}%
\pgfsetbuttcap%
\pgfsetroundjoin%
\definecolor{currentfill}{rgb}{0.556863,0.654902,0.847059}%
\pgfsetfillcolor{currentfill}%
\pgfsetfillopacity{0.130000}%
\pgfsetlinewidth{0.000000pt}%
\definecolor{currentstroke}{rgb}{0.000000,0.000000,0.000000}%
\pgfsetstrokecolor{currentstroke}%
\pgfsetdash{}{0pt}%
\pgfpathmoveto{\pgfqpoint{0.637558in}{1.121797in}}%
\pgfpathlineto{\pgfqpoint{0.653569in}{1.102130in}}%
\pgfpathlineto{\pgfqpoint{0.669588in}{1.082539in}}%
\pgfpathlineto{\pgfqpoint{0.678308in}{1.084942in}}%
\pgfpathlineto{\pgfqpoint{0.686990in}{1.087719in}}%
\pgfpathlineto{\pgfqpoint{0.671013in}{1.107403in}}%
\pgfpathlineto{\pgfqpoint{0.655045in}{1.127152in}}%
\pgfpathlineto{\pgfqpoint{0.646321in}{1.124300in}}%
\pgfpathlineto{\pgfqpoint{0.637558in}{1.121797in}}%
\pgfpathclose%
\pgfusepath{fill}%
\end{pgfscope}%
\begin{pgfscope}%
\pgfpathrectangle{\pgfqpoint{0.407592in}{0.222904in}}{\pgfqpoint{1.411163in}{1.411163in}}%
\pgfusepath{clip}%
\pgfsetbuttcap%
\pgfsetroundjoin%
\definecolor{currentfill}{rgb}{0.556863,0.654902,0.847059}%
\pgfsetfillcolor{currentfill}%
\pgfsetfillopacity{0.130000}%
\pgfsetlinewidth{0.000000pt}%
\definecolor{currentstroke}{rgb}{0.000000,0.000000,0.000000}%
\pgfsetstrokecolor{currentstroke}%
\pgfsetdash{}{0pt}%
\pgfpathmoveto{\pgfqpoint{0.928231in}{0.833726in}}%
\pgfpathlineto{\pgfqpoint{0.944298in}{0.817685in}}%
\pgfpathlineto{\pgfqpoint{0.960375in}{0.802408in}}%
\pgfpathlineto{\pgfqpoint{0.968607in}{0.805801in}}%
\pgfpathlineto{\pgfqpoint{0.976805in}{0.810076in}}%
\pgfpathlineto{\pgfqpoint{0.960768in}{0.825461in}}%
\pgfpathlineto{\pgfqpoint{0.944742in}{0.841599in}}%
\pgfpathlineto{\pgfqpoint{0.936504in}{0.837272in}}%
\pgfpathlineto{\pgfqpoint{0.928231in}{0.833726in}}%
\pgfpathclose%
\pgfusepath{fill}%
\end{pgfscope}%
\begin{pgfscope}%
\pgfpathrectangle{\pgfqpoint{0.407592in}{0.222904in}}{\pgfqpoint{1.411163in}{1.411163in}}%
\pgfusepath{clip}%
\pgfsetbuttcap%
\pgfsetroundjoin%
\definecolor{currentfill}{rgb}{0.556863,0.654902,0.847059}%
\pgfsetfillcolor{currentfill}%
\pgfsetfillopacity{0.130000}%
\pgfsetlinewidth{0.000000pt}%
\definecolor{currentstroke}{rgb}{0.000000,0.000000,0.000000}%
\pgfsetstrokecolor{currentstroke}%
\pgfsetdash{}{0pt}%
\pgfpathmoveto{\pgfqpoint{1.426133in}{1.001563in}}%
\pgfpathlineto{\pgfqpoint{1.442821in}{1.007411in}}%
\pgfpathlineto{\pgfqpoint{1.459576in}{1.013815in}}%
\pgfpathlineto{\pgfqpoint{1.467196in}{1.027626in}}%
\pgfpathlineto{\pgfqpoint{1.474805in}{1.041628in}}%
\pgfpathlineto{\pgfqpoint{1.458065in}{1.035690in}}%
\pgfpathlineto{\pgfqpoint{1.441394in}{1.030305in}}%
\pgfpathlineto{\pgfqpoint{1.433770in}{1.015845in}}%
\pgfpathlineto{\pgfqpoint{1.426133in}{1.001563in}}%
\pgfpathclose%
\pgfusepath{fill}%
\end{pgfscope}%
\begin{pgfscope}%
\pgfpathrectangle{\pgfqpoint{0.407592in}{0.222904in}}{\pgfqpoint{1.411163in}{1.411163in}}%
\pgfusepath{clip}%
\pgfsetbuttcap%
\pgfsetroundjoin%
\definecolor{currentfill}{rgb}{0.556863,0.654902,0.847059}%
\pgfsetfillcolor{currentfill}%
\pgfsetfillopacity{0.130000}%
\pgfsetlinewidth{0.000000pt}%
\definecolor{currentstroke}{rgb}{0.000000,0.000000,0.000000}%
\pgfsetstrokecolor{currentstroke}%
\pgfsetdash{}{0pt}%
\pgfpathmoveto{\pgfqpoint{1.249907in}{0.836852in}}%
\pgfpathlineto{\pgfqpoint{1.266277in}{0.840686in}}%
\pgfpathlineto{\pgfqpoint{1.282710in}{0.845550in}}%
\pgfpathlineto{\pgfqpoint{1.290501in}{0.857961in}}%
\pgfpathlineto{\pgfqpoint{1.298279in}{0.870807in}}%
\pgfpathlineto{\pgfqpoint{1.281862in}{0.866511in}}%
\pgfpathlineto{\pgfqpoint{1.265506in}{0.863214in}}%
\pgfpathlineto{\pgfqpoint{1.257714in}{0.849820in}}%
\pgfpathlineto{\pgfqpoint{1.249907in}{0.836852in}}%
\pgfpathclose%
\pgfusepath{fill}%
\end{pgfscope}%
\begin{pgfscope}%
\pgfpathrectangle{\pgfqpoint{0.407592in}{0.222904in}}{\pgfqpoint{1.411163in}{1.411163in}}%
\pgfusepath{clip}%
\pgfsetbuttcap%
\pgfsetroundjoin%
\definecolor{currentfill}{rgb}{0.556863,0.654902,0.847059}%
\pgfsetfillcolor{currentfill}%
\pgfsetfillopacity{0.130000}%
\pgfsetlinewidth{0.000000pt}%
\definecolor{currentstroke}{rgb}{0.000000,0.000000,0.000000}%
\pgfsetstrokecolor{currentstroke}%
\pgfsetdash{}{0pt}%
\pgfpathmoveto{\pgfqpoint{0.815115in}{0.934101in}}%
\pgfpathlineto{\pgfqpoint{0.831170in}{0.915678in}}%
\pgfpathlineto{\pgfqpoint{0.847233in}{0.897533in}}%
\pgfpathlineto{\pgfqpoint{0.855664in}{0.900090in}}%
\pgfpathlineto{\pgfqpoint{0.864057in}{0.903246in}}%
\pgfpathlineto{\pgfqpoint{0.848036in}{0.921557in}}%
\pgfpathlineto{\pgfqpoint{0.832024in}{0.940128in}}%
\pgfpathlineto{\pgfqpoint{0.823588in}{0.936844in}}%
\pgfpathlineto{\pgfqpoint{0.815115in}{0.934101in}}%
\pgfpathclose%
\pgfusepath{fill}%
\end{pgfscope}%
\begin{pgfscope}%
\pgfpathrectangle{\pgfqpoint{0.407592in}{0.222904in}}{\pgfqpoint{1.411163in}{1.411163in}}%
\pgfusepath{clip}%
\pgfsetbuttcap%
\pgfsetroundjoin%
\definecolor{currentfill}{rgb}{0.556863,0.654902,0.847059}%
\pgfsetfillcolor{currentfill}%
\pgfsetfillopacity{0.130000}%
\pgfsetlinewidth{0.000000pt}%
\definecolor{currentstroke}{rgb}{0.000000,0.000000,0.000000}%
\pgfsetstrokecolor{currentstroke}%
\pgfsetdash{}{0pt}%
\pgfpathmoveto{\pgfqpoint{1.362233in}{0.934699in}}%
\pgfpathlineto{\pgfqpoint{1.378817in}{0.940307in}}%
\pgfpathlineto{\pgfqpoint{1.395469in}{0.946572in}}%
\pgfpathlineto{\pgfqpoint{1.403153in}{0.959957in}}%
\pgfpathlineto{\pgfqpoint{1.410825in}{0.973598in}}%
\pgfpathlineto{\pgfqpoint{1.394190in}{0.967857in}}%
\pgfpathlineto{\pgfqpoint{1.377621in}{0.962768in}}%
\pgfpathlineto{\pgfqpoint{1.369933in}{0.948613in}}%
\pgfpathlineto{\pgfqpoint{1.362233in}{0.934699in}}%
\pgfpathclose%
\pgfusepath{fill}%
\end{pgfscope}%
\begin{pgfscope}%
\pgfpathrectangle{\pgfqpoint{0.407592in}{0.222904in}}{\pgfqpoint{1.411163in}{1.411163in}}%
\pgfusepath{clip}%
\pgfsetbuttcap%
\pgfsetroundjoin%
\definecolor{currentfill}{rgb}{0.556863,0.654902,0.847059}%
\pgfsetfillcolor{currentfill}%
\pgfsetfillopacity{0.130000}%
\pgfsetlinewidth{0.000000pt}%
\definecolor{currentstroke}{rgb}{0.000000,0.000000,0.000000}%
\pgfsetstrokecolor{currentstroke}%
\pgfsetdash{}{0pt}%
\pgfpathmoveto{\pgfqpoint{0.669588in}{1.082539in}}%
\pgfpathlineto{\pgfqpoint{0.685616in}{1.063034in}}%
\pgfpathlineto{\pgfqpoint{0.701653in}{1.043626in}}%
\pgfpathlineto{\pgfqpoint{0.710330in}{1.045904in}}%
\pgfpathlineto{\pgfqpoint{0.718969in}{1.048586in}}%
\pgfpathlineto{\pgfqpoint{0.702975in}{1.068110in}}%
\pgfpathlineto{\pgfqpoint{0.686990in}{1.087719in}}%
\pgfpathlineto{\pgfqpoint{0.678308in}{1.084942in}}%
\pgfpathlineto{\pgfqpoint{0.669588in}{1.082539in}}%
\pgfpathclose%
\pgfusepath{fill}%
\end{pgfscope}%
\begin{pgfscope}%
\pgfpathrectangle{\pgfqpoint{0.407592in}{0.222904in}}{\pgfqpoint{1.411163in}{1.411163in}}%
\pgfusepath{clip}%
\pgfsetbuttcap%
\pgfsetroundjoin%
\definecolor{currentfill}{rgb}{0.556863,0.654902,0.847059}%
\pgfsetfillcolor{currentfill}%
\pgfsetfillopacity{0.130000}%
\pgfsetlinewidth{0.000000pt}%
\definecolor{currentstroke}{rgb}{0.000000,0.000000,0.000000}%
\pgfsetstrokecolor{currentstroke}%
\pgfsetdash{}{0pt}%
\pgfpathmoveto{\pgfqpoint{1.089418in}{0.755736in}}%
\pgfpathlineto{\pgfqpoint{1.105555in}{0.750973in}}%
\pgfpathlineto{\pgfqpoint{1.121729in}{0.748276in}}%
\pgfpathlineto{\pgfqpoint{1.129707in}{0.755882in}}%
\pgfpathlineto{\pgfqpoint{1.137663in}{0.764555in}}%
\pgfpathlineto{\pgfqpoint{1.121512in}{0.767285in}}%
\pgfpathlineto{\pgfqpoint{1.105400in}{0.771992in}}%
\pgfpathlineto{\pgfqpoint{1.097421in}{0.763328in}}%
\pgfpathlineto{\pgfqpoint{1.089418in}{0.755736in}}%
\pgfpathclose%
\pgfusepath{fill}%
\end{pgfscope}%
\begin{pgfscope}%
\pgfpathrectangle{\pgfqpoint{0.407592in}{0.222904in}}{\pgfqpoint{1.411163in}{1.411163in}}%
\pgfusepath{clip}%
\pgfsetbuttcap%
\pgfsetroundjoin%
\definecolor{currentfill}{rgb}{0.556863,0.654902,0.847059}%
\pgfsetfillcolor{currentfill}%
\pgfsetfillopacity{0.130000}%
\pgfsetlinewidth{0.000000pt}%
\definecolor{currentstroke}{rgb}{0.000000,0.000000,0.000000}%
\pgfsetstrokecolor{currentstroke}%
\pgfsetdash{}{0pt}%
\pgfpathmoveto{\pgfqpoint{1.137663in}{0.764555in}}%
\pgfpathlineto{\pgfqpoint{1.153859in}{0.763749in}}%
\pgfpathlineto{\pgfqpoint{1.170103in}{0.764745in}}%
\pgfpathlineto{\pgfqpoint{1.178019in}{0.774189in}}%
\pgfpathlineto{\pgfqpoint{1.185918in}{0.784499in}}%
\pgfpathlineto{\pgfqpoint{1.169691in}{0.783786in}}%
\pgfpathlineto{\pgfqpoint{1.153515in}{0.784774in}}%
\pgfpathlineto{\pgfqpoint{1.145599in}{0.774216in}}%
\pgfpathlineto{\pgfqpoint{1.137663in}{0.764555in}}%
\pgfpathclose%
\pgfusepath{fill}%
\end{pgfscope}%
\begin{pgfscope}%
\pgfpathrectangle{\pgfqpoint{0.407592in}{0.222904in}}{\pgfqpoint{1.411163in}{1.411163in}}%
\pgfusepath{clip}%
\pgfsetbuttcap%
\pgfsetroundjoin%
\definecolor{currentfill}{rgb}{0.556863,0.654902,0.847059}%
\pgfsetfillcolor{currentfill}%
\pgfsetfillopacity{0.130000}%
\pgfsetlinewidth{0.000000pt}%
\definecolor{currentstroke}{rgb}{0.000000,0.000000,0.000000}%
\pgfsetstrokecolor{currentstroke}%
\pgfsetdash{}{0pt}%
\pgfpathmoveto{\pgfqpoint{0.960375in}{0.802408in}}%
\pgfpathlineto{\pgfqpoint{0.976463in}{0.788051in}}%
\pgfpathlineto{\pgfqpoint{0.992563in}{0.774795in}}%
\pgfpathlineto{\pgfqpoint{1.000755in}{0.778016in}}%
\pgfpathlineto{\pgfqpoint{1.008914in}{0.782227in}}%
\pgfpathlineto{\pgfqpoint{0.992853in}{0.795603in}}%
\pgfpathlineto{\pgfqpoint{0.976805in}{0.810076in}}%
\pgfpathlineto{\pgfqpoint{0.968607in}{0.805801in}}%
\pgfpathlineto{\pgfqpoint{0.960375in}{0.802408in}}%
\pgfpathclose%
\pgfusepath{fill}%
\end{pgfscope}%
\begin{pgfscope}%
\pgfpathrectangle{\pgfqpoint{0.407592in}{0.222904in}}{\pgfqpoint{1.411163in}{1.411163in}}%
\pgfusepath{clip}%
\pgfsetbuttcap%
\pgfsetroundjoin%
\definecolor{currentfill}{rgb}{0.556863,0.654902,0.847059}%
\pgfsetfillcolor{currentfill}%
\pgfsetfillopacity{0.130000}%
\pgfsetlinewidth{0.000000pt}%
\definecolor{currentstroke}{rgb}{0.000000,0.000000,0.000000}%
\pgfsetstrokecolor{currentstroke}%
\pgfsetdash{}{0pt}%
\pgfpathmoveto{\pgfqpoint{0.847233in}{0.897533in}}%
\pgfpathlineto{\pgfqpoint{0.863305in}{0.879715in}}%
\pgfpathlineto{\pgfqpoint{0.879386in}{0.862286in}}%
\pgfpathlineto{\pgfqpoint{0.887774in}{0.864612in}}%
\pgfpathlineto{\pgfqpoint{0.896126in}{0.867602in}}%
\pgfpathlineto{\pgfqpoint{0.880087in}{0.885241in}}%
\pgfpathlineto{\pgfqpoint{0.864057in}{0.903246in}}%
\pgfpathlineto{\pgfqpoint{0.855664in}{0.900090in}}%
\pgfpathlineto{\pgfqpoint{0.847233in}{0.897533in}}%
\pgfpathclose%
\pgfusepath{fill}%
\end{pgfscope}%
\begin{pgfscope}%
\pgfpathrectangle{\pgfqpoint{0.407592in}{0.222904in}}{\pgfqpoint{1.411163in}{1.411163in}}%
\pgfusepath{clip}%
\pgfsetbuttcap%
\pgfsetroundjoin%
\definecolor{currentfill}{rgb}{0.556863,0.654902,0.847059}%
\pgfsetfillcolor{currentfill}%
\pgfsetfillopacity{0.130000}%
\pgfsetlinewidth{0.000000pt}%
\definecolor{currentstroke}{rgb}{0.000000,0.000000,0.000000}%
\pgfsetstrokecolor{currentstroke}%
\pgfsetdash{}{0pt}%
\pgfpathmoveto{\pgfqpoint{0.701653in}{1.043626in}}%
\pgfpathlineto{\pgfqpoint{0.717698in}{1.024330in}}%
\pgfpathlineto{\pgfqpoint{0.733752in}{1.005162in}}%
\pgfpathlineto{\pgfqpoint{0.742387in}{1.007285in}}%
\pgfpathlineto{\pgfqpoint{0.750983in}{1.009845in}}%
\pgfpathlineto{\pgfqpoint{0.734972in}{1.029159in}}%
\pgfpathlineto{\pgfqpoint{0.718969in}{1.048586in}}%
\pgfpathlineto{\pgfqpoint{0.710330in}{1.045904in}}%
\pgfpathlineto{\pgfqpoint{0.701653in}{1.043626in}}%
\pgfpathclose%
\pgfusepath{fill}%
\end{pgfscope}%
\begin{pgfscope}%
\pgfpathrectangle{\pgfqpoint{0.407592in}{0.222904in}}{\pgfqpoint{1.411163in}{1.411163in}}%
\pgfusepath{clip}%
\pgfsetbuttcap%
\pgfsetroundjoin%
\definecolor{currentfill}{rgb}{0.556863,0.654902,0.847059}%
\pgfsetfillcolor{currentfill}%
\pgfsetfillopacity{0.130000}%
\pgfsetlinewidth{0.000000pt}%
\definecolor{currentstroke}{rgb}{0.000000,0.000000,0.000000}%
\pgfsetstrokecolor{currentstroke}%
\pgfsetdash{}{0pt}%
\pgfpathmoveto{\pgfqpoint{1.298279in}{0.870807in}}%
\pgfpathlineto{\pgfqpoint{1.314761in}{0.876005in}}%
\pgfpathlineto{\pgfqpoint{1.331310in}{0.882011in}}%
\pgfpathlineto{\pgfqpoint{1.339060in}{0.894674in}}%
\pgfpathlineto{\pgfqpoint{1.346797in}{0.907699in}}%
\pgfpathlineto{\pgfqpoint{1.330264in}{0.902265in}}%
\pgfpathlineto{\pgfqpoint{1.313797in}{0.897632in}}%
\pgfpathlineto{\pgfqpoint{1.306044in}{0.884045in}}%
\pgfpathlineto{\pgfqpoint{1.298279in}{0.870807in}}%
\pgfpathclose%
\pgfusepath{fill}%
\end{pgfscope}%
\begin{pgfscope}%
\pgfpathrectangle{\pgfqpoint{0.407592in}{0.222904in}}{\pgfqpoint{1.411163in}{1.411163in}}%
\pgfusepath{clip}%
\pgfsetbuttcap%
\pgfsetroundjoin%
\definecolor{currentfill}{rgb}{0.556863,0.654902,0.847059}%
\pgfsetfillcolor{currentfill}%
\pgfsetfillopacity{0.130000}%
\pgfsetlinewidth{0.000000pt}%
\definecolor{currentstroke}{rgb}{0.000000,0.000000,0.000000}%
\pgfsetstrokecolor{currentstroke}%
\pgfsetdash{}{0pt}%
\pgfpathmoveto{\pgfqpoint{1.041082in}{0.759620in}}%
\pgfpathlineto{\pgfqpoint{1.057194in}{0.750833in}}%
\pgfpathlineto{\pgfqpoint{1.073331in}{0.743981in}}%
\pgfpathlineto{\pgfqpoint{1.081388in}{0.749276in}}%
\pgfpathlineto{\pgfqpoint{1.089418in}{0.755736in}}%
\pgfpathlineto{\pgfqpoint{1.073312in}{0.762532in}}%
\pgfpathlineto{\pgfqpoint{1.057232in}{0.771245in}}%
\pgfpathlineto{\pgfqpoint{1.049172in}{0.764880in}}%
\pgfpathlineto{\pgfqpoint{1.041082in}{0.759620in}}%
\pgfpathclose%
\pgfusepath{fill}%
\end{pgfscope}%
\begin{pgfscope}%
\pgfpathrectangle{\pgfqpoint{0.407592in}{0.222904in}}{\pgfqpoint{1.411163in}{1.411163in}}%
\pgfusepath{clip}%
\pgfsetbuttcap%
\pgfsetroundjoin%
\definecolor{currentfill}{rgb}{0.556863,0.654902,0.847059}%
\pgfsetfillcolor{currentfill}%
\pgfsetfillopacity{0.130000}%
\pgfsetlinewidth{0.000000pt}%
\definecolor{currentstroke}{rgb}{0.000000,0.000000,0.000000}%
\pgfsetstrokecolor{currentstroke}%
\pgfsetdash{}{0pt}%
\pgfpathmoveto{\pgfqpoint{1.185918in}{0.784499in}}%
\pgfpathlineto{\pgfqpoint{1.202199in}{0.786772in}}%
\pgfpathlineto{\pgfqpoint{1.218538in}{0.790440in}}%
\pgfpathlineto{\pgfqpoint{1.226403in}{0.801094in}}%
\pgfpathlineto{\pgfqpoint{1.234253in}{0.812426in}}%
\pgfpathlineto{\pgfqpoint{1.217929in}{0.809237in}}%
\pgfpathlineto{\pgfqpoint{1.201664in}{0.807384in}}%
\pgfpathlineto{\pgfqpoint{1.193799in}{0.795592in}}%
\pgfpathlineto{\pgfqpoint{1.185918in}{0.784499in}}%
\pgfpathclose%
\pgfusepath{fill}%
\end{pgfscope}%
\begin{pgfscope}%
\pgfpathrectangle{\pgfqpoint{0.407592in}{0.222904in}}{\pgfqpoint{1.411163in}{1.411163in}}%
\pgfusepath{clip}%
\pgfsetbuttcap%
\pgfsetroundjoin%
\definecolor{currentfill}{rgb}{0.556863,0.654902,0.847059}%
\pgfsetfillcolor{currentfill}%
\pgfsetfillopacity{0.130000}%
\pgfsetlinewidth{0.000000pt}%
\definecolor{currentstroke}{rgb}{0.000000,0.000000,0.000000}%
\pgfsetstrokecolor{currentstroke}%
\pgfsetdash{}{0pt}%
\pgfpathmoveto{\pgfqpoint{1.474805in}{1.041628in}}%
\pgfpathlineto{\pgfqpoint{1.491613in}{1.048082in}}%
\pgfpathlineto{\pgfqpoint{1.508491in}{1.055016in}}%
\pgfpathlineto{\pgfqpoint{1.516070in}{1.068750in}}%
\pgfpathlineto{\pgfqpoint{1.523638in}{1.082659in}}%
\pgfpathlineto{\pgfqpoint{1.506778in}{1.076162in}}%
\pgfpathlineto{\pgfqpoint{1.489986in}{1.070146in}}%
\pgfpathlineto{\pgfqpoint{1.482401in}{1.055806in}}%
\pgfpathlineto{\pgfqpoint{1.474805in}{1.041628in}}%
\pgfpathclose%
\pgfusepath{fill}%
\end{pgfscope}%
\begin{pgfscope}%
\pgfpathrectangle{\pgfqpoint{0.407592in}{0.222904in}}{\pgfqpoint{1.411163in}{1.411163in}}%
\pgfusepath{clip}%
\pgfsetbuttcap%
\pgfsetroundjoin%
\definecolor{currentfill}{rgb}{0.556863,0.654902,0.847059}%
\pgfsetfillcolor{currentfill}%
\pgfsetfillopacity{0.130000}%
\pgfsetlinewidth{0.000000pt}%
\definecolor{currentstroke}{rgb}{0.000000,0.000000,0.000000}%
\pgfsetstrokecolor{currentstroke}%
\pgfsetdash{}{0pt}%
\pgfpathmoveto{\pgfqpoint{1.538740in}{1.110948in}}%
\pgfpathlineto{\pgfqpoint{1.555653in}{1.117485in}}%
\pgfpathlineto{\pgfqpoint{1.572637in}{1.124449in}}%
\pgfpathlineto{\pgfqpoint{1.580153in}{1.138411in}}%
\pgfpathlineto{\pgfqpoint{1.563178in}{1.131644in}}%
\pgfpathlineto{\pgfqpoint{1.546274in}{1.125305in}}%
\pgfusepath{fill}%
\end{pgfscope}%
\begin{pgfscope}%
\pgfpathrectangle{\pgfqpoint{0.407592in}{0.222904in}}{\pgfqpoint{1.411163in}{1.411163in}}%
\pgfusepath{clip}%
\pgfsetbuttcap%
\pgfsetroundjoin%
\definecolor{currentfill}{rgb}{0.556863,0.654902,0.847059}%
\pgfsetfillcolor{currentfill}%
\pgfsetfillopacity{0.130000}%
\pgfsetlinewidth{0.000000pt}%
\definecolor{currentstroke}{rgb}{0.000000,0.000000,0.000000}%
\pgfsetstrokecolor{currentstroke}%
\pgfsetdash{}{0pt}%
\pgfpathmoveto{\pgfqpoint{1.410825in}{0.973598in}}%
\pgfpathlineto{\pgfqpoint{1.427529in}{0.979942in}}%
\pgfpathlineto{\pgfqpoint{1.444301in}{0.986838in}}%
\pgfpathlineto{\pgfqpoint{1.451945in}{1.000213in}}%
\pgfpathlineto{\pgfqpoint{1.459576in}{1.013815in}}%
\pgfpathlineto{\pgfqpoint{1.442821in}{1.007411in}}%
\pgfpathlineto{\pgfqpoint{1.426133in}{1.001563in}}%
\pgfpathlineto{\pgfqpoint{1.418485in}{0.987474in}}%
\pgfpathlineto{\pgfqpoint{1.410825in}{0.973598in}}%
\pgfpathclose%
\pgfusepath{fill}%
\end{pgfscope}%
\begin{pgfscope}%
\pgfpathrectangle{\pgfqpoint{0.407592in}{0.222904in}}{\pgfqpoint{1.411163in}{1.411163in}}%
\pgfusepath{clip}%
\pgfsetbuttcap%
\pgfsetroundjoin%
\definecolor{currentfill}{rgb}{0.556863,0.654902,0.847059}%
\pgfsetfillcolor{currentfill}%
\pgfsetfillopacity{0.130000}%
\pgfsetlinewidth{0.000000pt}%
\definecolor{currentstroke}{rgb}{0.000000,0.000000,0.000000}%
\pgfsetstrokecolor{currentstroke}%
\pgfsetdash{}{0pt}%
\pgfpathmoveto{\pgfqpoint{0.733752in}{1.005162in}}%
\pgfpathlineto{\pgfqpoint{0.749815in}{0.986140in}}%
\pgfpathlineto{\pgfqpoint{0.765886in}{0.967286in}}%
\pgfpathlineto{\pgfqpoint{0.774478in}{0.969217in}}%
\pgfpathlineto{\pgfqpoint{0.783032in}{0.971623in}}%
\pgfpathlineto{\pgfqpoint{0.767003in}{0.990659in}}%
\pgfpathlineto{\pgfqpoint{0.750983in}{1.009845in}}%
\pgfpathlineto{\pgfqpoint{0.742387in}{1.007285in}}%
\pgfpathlineto{\pgfqpoint{0.733752in}{1.005162in}}%
\pgfpathclose%
\pgfusepath{fill}%
\end{pgfscope}%
\begin{pgfscope}%
\pgfpathrectangle{\pgfqpoint{0.407592in}{0.222904in}}{\pgfqpoint{1.411163in}{1.411163in}}%
\pgfusepath{clip}%
\pgfsetbuttcap%
\pgfsetroundjoin%
\definecolor{currentfill}{rgb}{0.556863,0.654902,0.847059}%
\pgfsetfillcolor{currentfill}%
\pgfsetfillopacity{0.130000}%
\pgfsetlinewidth{0.000000pt}%
\definecolor{currentstroke}{rgb}{0.000000,0.000000,0.000000}%
\pgfsetstrokecolor{currentstroke}%
\pgfsetdash{}{0pt}%
\pgfpathmoveto{\pgfqpoint{0.879386in}{0.862286in}}%
\pgfpathlineto{\pgfqpoint{0.895476in}{0.845319in}}%
\pgfpathlineto{\pgfqpoint{0.911574in}{0.828902in}}%
\pgfpathlineto{\pgfqpoint{0.919921in}{0.830943in}}%
\pgfpathlineto{\pgfqpoint{0.928231in}{0.833726in}}%
\pgfpathlineto{\pgfqpoint{0.912173in}{0.850400in}}%
\pgfpathlineto{\pgfqpoint{0.896126in}{0.867602in}}%
\pgfpathlineto{\pgfqpoint{0.887774in}{0.864612in}}%
\pgfpathlineto{\pgfqpoint{0.879386in}{0.862286in}}%
\pgfpathclose%
\pgfusepath{fill}%
\end{pgfscope}%
\begin{pgfscope}%
\pgfpathrectangle{\pgfqpoint{0.407592in}{0.222904in}}{\pgfqpoint{1.411163in}{1.411163in}}%
\pgfusepath{clip}%
\pgfsetbuttcap%
\pgfsetroundjoin%
\definecolor{currentfill}{rgb}{0.556863,0.654902,0.847059}%
\pgfsetfillcolor{currentfill}%
\pgfsetfillopacity{0.130000}%
\pgfsetlinewidth{0.000000pt}%
\definecolor{currentstroke}{rgb}{0.000000,0.000000,0.000000}%
\pgfsetstrokecolor{currentstroke}%
\pgfsetdash{}{0pt}%
\pgfpathmoveto{\pgfqpoint{1.234253in}{0.812426in}}%
\pgfpathlineto{\pgfqpoint{1.250637in}{0.816806in}}%
\pgfpathlineto{\pgfqpoint{1.267086in}{0.822236in}}%
\pgfpathlineto{\pgfqpoint{1.274905in}{0.833625in}}%
\pgfpathlineto{\pgfqpoint{1.282710in}{0.845550in}}%
\pgfpathlineto{\pgfqpoint{1.266277in}{0.840686in}}%
\pgfpathlineto{\pgfqpoint{1.249907in}{0.836852in}}%
\pgfpathlineto{\pgfqpoint{1.242087in}{0.824367in}}%
\pgfpathlineto{\pgfqpoint{1.234253in}{0.812426in}}%
\pgfpathclose%
\pgfusepath{fill}%
\end{pgfscope}%
\begin{pgfscope}%
\pgfpathrectangle{\pgfqpoint{0.407592in}{0.222904in}}{\pgfqpoint{1.411163in}{1.411163in}}%
\pgfusepath{clip}%
\pgfsetbuttcap%
\pgfsetroundjoin%
\definecolor{currentfill}{rgb}{0.556863,0.654902,0.847059}%
\pgfsetfillcolor{currentfill}%
\pgfsetfillopacity{0.130000}%
\pgfsetlinewidth{0.000000pt}%
\definecolor{currentstroke}{rgb}{0.000000,0.000000,0.000000}%
\pgfsetstrokecolor{currentstroke}%
\pgfsetdash{}{0pt}%
\pgfpathmoveto{\pgfqpoint{0.587834in}{1.157126in}}%
\pgfpathlineto{\pgfqpoint{0.603870in}{1.137440in}}%
\pgfpathlineto{\pgfqpoint{0.619915in}{1.117823in}}%
\pgfpathlineto{\pgfqpoint{0.628756in}{1.119639in}}%
\pgfpathlineto{\pgfqpoint{0.637558in}{1.121797in}}%
\pgfpathlineto{\pgfqpoint{0.621556in}{1.141530in}}%
\pgfpathlineto{\pgfqpoint{0.605563in}{1.161323in}}%
\pgfpathlineto{\pgfqpoint{0.596718in}{1.159064in}}%
\pgfpathlineto{\pgfqpoint{0.587834in}{1.157126in}}%
\pgfpathclose%
\pgfusepath{fill}%
\end{pgfscope}%
\begin{pgfscope}%
\pgfpathrectangle{\pgfqpoint{0.407592in}{0.222904in}}{\pgfqpoint{1.411163in}{1.411163in}}%
\pgfusepath{clip}%
\pgfsetbuttcap%
\pgfsetroundjoin%
\definecolor{currentfill}{rgb}{0.556863,0.654902,0.847059}%
\pgfsetfillcolor{currentfill}%
\pgfsetfillopacity{0.130000}%
\pgfsetlinewidth{0.000000pt}%
\definecolor{currentstroke}{rgb}{0.000000,0.000000,0.000000}%
\pgfsetstrokecolor{currentstroke}%
\pgfsetdash{}{0pt}%
\pgfpathmoveto{\pgfqpoint{0.992563in}{0.774795in}}%
\pgfpathlineto{\pgfqpoint{1.008677in}{0.762842in}}%
\pgfpathlineto{\pgfqpoint{1.024806in}{0.752405in}}%
\pgfpathlineto{\pgfqpoint{1.032960in}{0.755469in}}%
\pgfpathlineto{\pgfqpoint{1.041082in}{0.759620in}}%
\pgfpathlineto{\pgfqpoint{1.024989in}{0.770158in}}%
\pgfpathlineto{\pgfqpoint{1.008914in}{0.782227in}}%
\pgfpathlineto{\pgfqpoint{1.000755in}{0.778016in}}%
\pgfpathlineto{\pgfqpoint{0.992563in}{0.774795in}}%
\pgfpathclose%
\pgfusepath{fill}%
\end{pgfscope}%
\begin{pgfscope}%
\pgfpathrectangle{\pgfqpoint{0.407592in}{0.222904in}}{\pgfqpoint{1.411163in}{1.411163in}}%
\pgfusepath{clip}%
\pgfsetbuttcap%
\pgfsetroundjoin%
\definecolor{currentfill}{rgb}{0.556863,0.654902,0.847059}%
\pgfsetfillcolor{currentfill}%
\pgfsetfillopacity{0.130000}%
\pgfsetlinewidth{0.000000pt}%
\definecolor{currentstroke}{rgb}{0.000000,0.000000,0.000000}%
\pgfsetstrokecolor{currentstroke}%
\pgfsetdash{}{0pt}%
\pgfpathmoveto{\pgfqpoint{1.346797in}{0.907699in}}%
\pgfpathlineto{\pgfqpoint{1.363396in}{0.913858in}}%
\pgfpathlineto{\pgfqpoint{1.380064in}{0.920672in}}%
\pgfpathlineto{\pgfqpoint{1.387773in}{0.933468in}}%
\pgfpathlineto{\pgfqpoint{1.395469in}{0.946572in}}%
\pgfpathlineto{\pgfqpoint{1.378817in}{0.940307in}}%
\pgfpathlineto{\pgfqpoint{1.362233in}{0.934699in}}%
\pgfpathlineto{\pgfqpoint{1.354521in}{0.921051in}}%
\pgfpathlineto{\pgfqpoint{1.346797in}{0.907699in}}%
\pgfpathclose%
\pgfusepath{fill}%
\end{pgfscope}%
\begin{pgfscope}%
\pgfpathrectangle{\pgfqpoint{0.407592in}{0.222904in}}{\pgfqpoint{1.411163in}{1.411163in}}%
\pgfusepath{clip}%
\pgfsetbuttcap%
\pgfsetroundjoin%
\definecolor{currentfill}{rgb}{0.556863,0.654902,0.847059}%
\pgfsetfillcolor{currentfill}%
\pgfsetfillopacity{0.130000}%
\pgfsetlinewidth{0.000000pt}%
\definecolor{currentstroke}{rgb}{0.000000,0.000000,0.000000}%
\pgfsetstrokecolor{currentstroke}%
\pgfsetdash{}{0pt}%
\pgfpathmoveto{\pgfqpoint{0.765886in}{0.967286in}}%
\pgfpathlineto{\pgfqpoint{0.781966in}{0.948626in}}%
\pgfpathlineto{\pgfqpoint{0.798055in}{0.930191in}}%
\pgfpathlineto{\pgfqpoint{0.806604in}{0.931888in}}%
\pgfpathlineto{\pgfqpoint{0.815115in}{0.934101in}}%
\pgfpathlineto{\pgfqpoint{0.799069in}{0.952761in}}%
\pgfpathlineto{\pgfqpoint{0.783032in}{0.971623in}}%
\pgfpathlineto{\pgfqpoint{0.774478in}{0.969217in}}%
\pgfpathlineto{\pgfqpoint{0.765886in}{0.967286in}}%
\pgfpathclose%
\pgfusepath{fill}%
\end{pgfscope}%
\begin{pgfscope}%
\pgfpathrectangle{\pgfqpoint{0.407592in}{0.222904in}}{\pgfqpoint{1.411163in}{1.411163in}}%
\pgfusepath{clip}%
\pgfsetbuttcap%
\pgfsetroundjoin%
\definecolor{currentfill}{rgb}{0.556863,0.654902,0.847059}%
\pgfsetfillcolor{currentfill}%
\pgfsetfillopacity{0.130000}%
\pgfsetlinewidth{0.000000pt}%
\definecolor{currentstroke}{rgb}{0.000000,0.000000,0.000000}%
\pgfsetstrokecolor{currentstroke}%
\pgfsetdash{}{0pt}%
\pgfpathmoveto{\pgfqpoint{1.523638in}{1.082659in}}%
\pgfpathlineto{\pgfqpoint{1.540569in}{1.089608in}}%
\pgfpathlineto{\pgfqpoint{1.557571in}{1.096981in}}%
\pgfpathlineto{\pgfqpoint{1.565110in}{1.110635in}}%
\pgfpathlineto{\pgfqpoint{1.572637in}{1.124449in}}%
\pgfpathlineto{\pgfqpoint{1.555653in}{1.117485in}}%
\pgfpathlineto{\pgfqpoint{1.538740in}{1.110948in}}%
\pgfpathlineto{\pgfqpoint{1.531195in}{1.096729in}}%
\pgfpathlineto{\pgfqpoint{1.523638in}{1.082659in}}%
\pgfpathclose%
\pgfusepath{fill}%
\end{pgfscope}%
\begin{pgfscope}%
\pgfpathrectangle{\pgfqpoint{0.407592in}{0.222904in}}{\pgfqpoint{1.411163in}{1.411163in}}%
\pgfusepath{clip}%
\pgfsetbuttcap%
\pgfsetroundjoin%
\definecolor{currentfill}{rgb}{0.556863,0.654902,0.847059}%
\pgfsetfillcolor{currentfill}%
\pgfsetfillopacity{0.130000}%
\pgfsetlinewidth{0.000000pt}%
\definecolor{currentstroke}{rgb}{0.000000,0.000000,0.000000}%
\pgfsetstrokecolor{currentstroke}%
\pgfsetdash{}{0pt}%
\pgfpathmoveto{\pgfqpoint{0.911574in}{0.828902in}}%
\pgfpathlineto{\pgfqpoint{0.927683in}{0.813142in}}%
\pgfpathlineto{\pgfqpoint{0.943801in}{0.798165in}}%
\pgfpathlineto{\pgfqpoint{0.952106in}{0.799874in}}%
\pgfpathlineto{\pgfqpoint{0.960375in}{0.802408in}}%
\pgfpathlineto{\pgfqpoint{0.944298in}{0.817685in}}%
\pgfpathlineto{\pgfqpoint{0.928231in}{0.833726in}}%
\pgfpathlineto{\pgfqpoint{0.919921in}{0.830943in}}%
\pgfpathlineto{\pgfqpoint{0.911574in}{0.828902in}}%
\pgfpathclose%
\pgfusepath{fill}%
\end{pgfscope}%
\begin{pgfscope}%
\pgfpathrectangle{\pgfqpoint{0.407592in}{0.222904in}}{\pgfqpoint{1.411163in}{1.411163in}}%
\pgfusepath{clip}%
\pgfsetbuttcap%
\pgfsetroundjoin%
\definecolor{currentfill}{rgb}{0.556863,0.654902,0.847059}%
\pgfsetfillcolor{currentfill}%
\pgfsetfillopacity{0.130000}%
\pgfsetlinewidth{0.000000pt}%
\definecolor{currentstroke}{rgb}{0.000000,0.000000,0.000000}%
\pgfsetstrokecolor{currentstroke}%
\pgfsetdash{}{0pt}%
\pgfpathmoveto{\pgfqpoint{0.619915in}{1.117823in}}%
\pgfpathlineto{\pgfqpoint{0.635969in}{1.098286in}}%
\pgfpathlineto{\pgfqpoint{0.652031in}{1.078837in}}%
\pgfpathlineto{\pgfqpoint{0.660829in}{1.080506in}}%
\pgfpathlineto{\pgfqpoint{0.669588in}{1.082539in}}%
\pgfpathlineto{\pgfqpoint{0.653569in}{1.102130in}}%
\pgfpathlineto{\pgfqpoint{0.637558in}{1.121797in}}%
\pgfpathlineto{\pgfqpoint{0.628756in}{1.119639in}}%
\pgfpathlineto{\pgfqpoint{0.619915in}{1.117823in}}%
\pgfpathclose%
\pgfusepath{fill}%
\end{pgfscope}%
\begin{pgfscope}%
\pgfpathrectangle{\pgfqpoint{0.407592in}{0.222904in}}{\pgfqpoint{1.411163in}{1.411163in}}%
\pgfusepath{clip}%
\pgfsetbuttcap%
\pgfsetroundjoin%
\definecolor{currentfill}{rgb}{0.556863,0.654902,0.847059}%
\pgfsetfillcolor{currentfill}%
\pgfsetfillopacity{0.130000}%
\pgfsetlinewidth{0.000000pt}%
\definecolor{currentstroke}{rgb}{0.000000,0.000000,0.000000}%
\pgfsetstrokecolor{currentstroke}%
\pgfsetdash{}{0pt}%
\pgfpathmoveto{\pgfqpoint{1.121729in}{0.748276in}}%
\pgfpathlineto{\pgfqpoint{1.137947in}{0.747584in}}%
\pgfpathlineto{\pgfqpoint{1.154213in}{0.748759in}}%
\pgfpathlineto{\pgfqpoint{1.162168in}{0.756246in}}%
\pgfpathlineto{\pgfqpoint{1.170103in}{0.764745in}}%
\pgfpathlineto{\pgfqpoint{1.153859in}{0.763749in}}%
\pgfpathlineto{\pgfqpoint{1.137663in}{0.764555in}}%
\pgfpathlineto{\pgfqpoint{1.129707in}{0.755882in}}%
\pgfpathlineto{\pgfqpoint{1.121729in}{0.748276in}}%
\pgfpathclose%
\pgfusepath{fill}%
\end{pgfscope}%
\begin{pgfscope}%
\pgfpathrectangle{\pgfqpoint{0.407592in}{0.222904in}}{\pgfqpoint{1.411163in}{1.411163in}}%
\pgfusepath{clip}%
\pgfsetbuttcap%
\pgfsetroundjoin%
\definecolor{currentfill}{rgb}{0.556863,0.654902,0.847059}%
\pgfsetfillcolor{currentfill}%
\pgfsetfillopacity{0.130000}%
\pgfsetlinewidth{0.000000pt}%
\definecolor{currentstroke}{rgb}{0.000000,0.000000,0.000000}%
\pgfsetstrokecolor{currentstroke}%
\pgfsetdash{}{0pt}%
\pgfpathmoveto{\pgfqpoint{1.459576in}{1.013815in}}%
\pgfpathlineto{\pgfqpoint{1.476402in}{1.020733in}}%
\pgfpathlineto{\pgfqpoint{1.493298in}{1.028128in}}%
\pgfpathlineto{\pgfqpoint{1.500900in}{1.041471in}}%
\pgfpathlineto{\pgfqpoint{1.508491in}{1.055016in}}%
\pgfpathlineto{\pgfqpoint{1.491613in}{1.048082in}}%
\pgfpathlineto{\pgfqpoint{1.474805in}{1.041628in}}%
\pgfpathlineto{\pgfqpoint{1.467196in}{1.027626in}}%
\pgfpathlineto{\pgfqpoint{1.459576in}{1.013815in}}%
\pgfpathclose%
\pgfusepath{fill}%
\end{pgfscope}%
\begin{pgfscope}%
\pgfpathrectangle{\pgfqpoint{0.407592in}{0.222904in}}{\pgfqpoint{1.411163in}{1.411163in}}%
\pgfusepath{clip}%
\pgfsetbuttcap%
\pgfsetroundjoin%
\definecolor{currentfill}{rgb}{0.556863,0.654902,0.847059}%
\pgfsetfillcolor{currentfill}%
\pgfsetfillopacity{0.130000}%
\pgfsetlinewidth{0.000000pt}%
\definecolor{currentstroke}{rgb}{0.000000,0.000000,0.000000}%
\pgfsetstrokecolor{currentstroke}%
\pgfsetdash{}{0pt}%
\pgfpathmoveto{\pgfqpoint{1.282710in}{0.845550in}}%
\pgfpathlineto{\pgfqpoint{1.299207in}{0.851329in}}%
\pgfpathlineto{\pgfqpoint{1.315772in}{0.857916in}}%
\pgfpathlineto{\pgfqpoint{1.323547in}{0.869745in}}%
\pgfpathlineto{\pgfqpoint{1.331310in}{0.882011in}}%
\pgfpathlineto{\pgfqpoint{1.314761in}{0.876005in}}%
\pgfpathlineto{\pgfqpoint{1.298279in}{0.870807in}}%
\pgfpathlineto{\pgfqpoint{1.290501in}{0.857961in}}%
\pgfpathlineto{\pgfqpoint{1.282710in}{0.845550in}}%
\pgfpathclose%
\pgfusepath{fill}%
\end{pgfscope}%
\begin{pgfscope}%
\pgfpathrectangle{\pgfqpoint{0.407592in}{0.222904in}}{\pgfqpoint{1.411163in}{1.411163in}}%
\pgfusepath{clip}%
\pgfsetbuttcap%
\pgfsetroundjoin%
\definecolor{currentfill}{rgb}{0.556863,0.654902,0.847059}%
\pgfsetfillcolor{currentfill}%
\pgfsetfillopacity{0.130000}%
\pgfsetlinewidth{0.000000pt}%
\definecolor{currentstroke}{rgb}{0.000000,0.000000,0.000000}%
\pgfsetstrokecolor{currentstroke}%
\pgfsetdash{}{0pt}%
\pgfpathmoveto{\pgfqpoint{1.073331in}{0.743981in}}%
\pgfpathlineto{\pgfqpoint{1.089498in}{0.739185in}}%
\pgfpathlineto{\pgfqpoint{1.105701in}{0.736477in}}%
\pgfpathlineto{\pgfqpoint{1.113728in}{0.741796in}}%
\pgfpathlineto{\pgfqpoint{1.121729in}{0.748276in}}%
\pgfpathlineto{\pgfqpoint{1.105555in}{0.750973in}}%
\pgfpathlineto{\pgfqpoint{1.089418in}{0.755736in}}%
\pgfpathlineto{\pgfqpoint{1.081388in}{0.749276in}}%
\pgfpathlineto{\pgfqpoint{1.073331in}{0.743981in}}%
\pgfpathclose%
\pgfusepath{fill}%
\end{pgfscope}%
\begin{pgfscope}%
\pgfpathrectangle{\pgfqpoint{0.407592in}{0.222904in}}{\pgfqpoint{1.411163in}{1.411163in}}%
\pgfusepath{clip}%
\pgfsetbuttcap%
\pgfsetroundjoin%
\definecolor{currentfill}{rgb}{0.556863,0.654902,0.847059}%
\pgfsetfillcolor{currentfill}%
\pgfsetfillopacity{0.130000}%
\pgfsetlinewidth{0.000000pt}%
\definecolor{currentstroke}{rgb}{0.000000,0.000000,0.000000}%
\pgfsetstrokecolor{currentstroke}%
\pgfsetdash{}{0pt}%
\pgfpathmoveto{\pgfqpoint{0.798055in}{0.930191in}}%
\pgfpathlineto{\pgfqpoint{0.814151in}{0.912018in}}%
\pgfpathlineto{\pgfqpoint{0.830257in}{0.894149in}}%
\pgfpathlineto{\pgfqpoint{0.838764in}{0.895558in}}%
\pgfpathlineto{\pgfqpoint{0.847233in}{0.897533in}}%
\pgfpathlineto{\pgfqpoint{0.831170in}{0.915678in}}%
\pgfpathlineto{\pgfqpoint{0.815115in}{0.934101in}}%
\pgfpathlineto{\pgfqpoint{0.806604in}{0.931888in}}%
\pgfpathlineto{\pgfqpoint{0.798055in}{0.930191in}}%
\pgfpathclose%
\pgfusepath{fill}%
\end{pgfscope}%
\begin{pgfscope}%
\pgfpathrectangle{\pgfqpoint{0.407592in}{0.222904in}}{\pgfqpoint{1.411163in}{1.411163in}}%
\pgfusepath{clip}%
\pgfsetbuttcap%
\pgfsetroundjoin%
\definecolor{currentfill}{rgb}{0.556863,0.654902,0.847059}%
\pgfsetfillcolor{currentfill}%
\pgfsetfillopacity{0.130000}%
\pgfsetlinewidth{0.000000pt}%
\definecolor{currentstroke}{rgb}{0.000000,0.000000,0.000000}%
\pgfsetstrokecolor{currentstroke}%
\pgfsetdash{}{0pt}%
\pgfpathmoveto{\pgfqpoint{1.170103in}{0.764745in}}%
\pgfpathlineto{\pgfqpoint{1.186402in}{0.767373in}}%
\pgfpathlineto{\pgfqpoint{1.202760in}{0.771439in}}%
\pgfpathlineto{\pgfqpoint{1.210657in}{0.780533in}}%
\pgfpathlineto{\pgfqpoint{1.218538in}{0.790440in}}%
\pgfpathlineto{\pgfqpoint{1.202199in}{0.786772in}}%
\pgfpathlineto{\pgfqpoint{1.185918in}{0.784499in}}%
\pgfpathlineto{\pgfqpoint{1.178019in}{0.774189in}}%
\pgfpathlineto{\pgfqpoint{1.170103in}{0.764745in}}%
\pgfpathclose%
\pgfusepath{fill}%
\end{pgfscope}%
\begin{pgfscope}%
\pgfpathrectangle{\pgfqpoint{0.407592in}{0.222904in}}{\pgfqpoint{1.411163in}{1.411163in}}%
\pgfusepath{clip}%
\pgfsetbuttcap%
\pgfsetroundjoin%
\definecolor{currentfill}{rgb}{0.556863,0.654902,0.847059}%
\pgfsetfillcolor{currentfill}%
\pgfsetfillopacity{0.130000}%
\pgfsetlinewidth{0.000000pt}%
\definecolor{currentstroke}{rgb}{0.000000,0.000000,0.000000}%
\pgfsetstrokecolor{currentstroke}%
\pgfsetdash{}{0pt}%
\pgfpathmoveto{\pgfqpoint{0.652031in}{1.078837in}}%
\pgfpathlineto{\pgfqpoint{0.668101in}{1.059489in}}%
\pgfpathlineto{\pgfqpoint{0.684181in}{1.040255in}}%
\pgfpathlineto{\pgfqpoint{0.692937in}{1.041746in}}%
\pgfpathlineto{\pgfqpoint{0.701653in}{1.043626in}}%
\pgfpathlineto{\pgfqpoint{0.685616in}{1.063034in}}%
\pgfpathlineto{\pgfqpoint{0.669588in}{1.082539in}}%
\pgfpathlineto{\pgfqpoint{0.660829in}{1.080506in}}%
\pgfpathlineto{\pgfqpoint{0.652031in}{1.078837in}}%
\pgfpathclose%
\pgfusepath{fill}%
\end{pgfscope}%
\begin{pgfscope}%
\pgfpathrectangle{\pgfqpoint{0.407592in}{0.222904in}}{\pgfqpoint{1.411163in}{1.411163in}}%
\pgfusepath{clip}%
\pgfsetbuttcap%
\pgfsetroundjoin%
\definecolor{currentfill}{rgb}{0.556863,0.654902,0.847059}%
\pgfsetfillcolor{currentfill}%
\pgfsetfillopacity{0.130000}%
\pgfsetlinewidth{0.000000pt}%
\definecolor{currentstroke}{rgb}{0.000000,0.000000,0.000000}%
\pgfsetstrokecolor{currentstroke}%
\pgfsetdash{}{0pt}%
\pgfpathmoveto{\pgfqpoint{1.395469in}{0.946572in}}%
\pgfpathlineto{\pgfqpoint{1.412189in}{0.953437in}}%
\pgfpathlineto{\pgfqpoint{1.428979in}{0.960848in}}%
\pgfpathlineto{\pgfqpoint{1.436646in}{0.973709in}}%
\pgfpathlineto{\pgfqpoint{1.444301in}{0.986838in}}%
\pgfpathlineto{\pgfqpoint{1.427529in}{0.979942in}}%
\pgfpathlineto{\pgfqpoint{1.410825in}{0.973598in}}%
\pgfpathlineto{\pgfqpoint{1.403153in}{0.959957in}}%
\pgfpathlineto{\pgfqpoint{1.395469in}{0.946572in}}%
\pgfpathclose%
\pgfusepath{fill}%
\end{pgfscope}%
\begin{pgfscope}%
\pgfpathrectangle{\pgfqpoint{0.407592in}{0.222904in}}{\pgfqpoint{1.411163in}{1.411163in}}%
\pgfusepath{clip}%
\pgfsetbuttcap%
\pgfsetroundjoin%
\definecolor{currentfill}{rgb}{0.556863,0.654902,0.847059}%
\pgfsetfillcolor{currentfill}%
\pgfsetfillopacity{0.130000}%
\pgfsetlinewidth{0.000000pt}%
\definecolor{currentstroke}{rgb}{0.000000,0.000000,0.000000}%
\pgfsetstrokecolor{currentstroke}%
\pgfsetdash{}{0pt}%
\pgfpathmoveto{\pgfqpoint{1.024806in}{0.752405in}}%
\pgfpathlineto{\pgfqpoint{1.040956in}{0.743693in}}%
\pgfpathlineto{\pgfqpoint{1.057128in}{0.736877in}}%
\pgfpathlineto{\pgfqpoint{1.065245in}{0.739857in}}%
\pgfpathlineto{\pgfqpoint{1.073331in}{0.743981in}}%
\pgfpathlineto{\pgfqpoint{1.057194in}{0.750833in}}%
\pgfpathlineto{\pgfqpoint{1.041082in}{0.759620in}}%
\pgfpathlineto{\pgfqpoint{1.032960in}{0.755469in}}%
\pgfpathlineto{\pgfqpoint{1.024806in}{0.752405in}}%
\pgfpathclose%
\pgfusepath{fill}%
\end{pgfscope}%
\begin{pgfscope}%
\pgfpathrectangle{\pgfqpoint{0.407592in}{0.222904in}}{\pgfqpoint{1.411163in}{1.411163in}}%
\pgfusepath{clip}%
\pgfsetbuttcap%
\pgfsetroundjoin%
\definecolor{currentfill}{rgb}{0.556863,0.654902,0.847059}%
\pgfsetfillcolor{currentfill}%
\pgfsetfillopacity{0.130000}%
\pgfsetlinewidth{0.000000pt}%
\definecolor{currentstroke}{rgb}{0.000000,0.000000,0.000000}%
\pgfsetstrokecolor{currentstroke}%
\pgfsetdash{}{0pt}%
\pgfpathmoveto{\pgfqpoint{0.943801in}{0.798165in}}%
\pgfpathlineto{\pgfqpoint{0.959930in}{0.784119in}}%
\pgfpathlineto{\pgfqpoint{0.976071in}{0.771175in}}%
\pgfpathlineto{\pgfqpoint{0.984335in}{0.772531in}}%
\pgfpathlineto{\pgfqpoint{0.992563in}{0.774795in}}%
\pgfpathlineto{\pgfqpoint{0.976463in}{0.788051in}}%
\pgfpathlineto{\pgfqpoint{0.960375in}{0.802408in}}%
\pgfpathlineto{\pgfqpoint{0.952106in}{0.799874in}}%
\pgfpathlineto{\pgfqpoint{0.943801in}{0.798165in}}%
\pgfpathclose%
\pgfusepath{fill}%
\end{pgfscope}%
\begin{pgfscope}%
\pgfpathrectangle{\pgfqpoint{0.407592in}{0.222904in}}{\pgfqpoint{1.411163in}{1.411163in}}%
\pgfusepath{clip}%
\pgfsetbuttcap%
\pgfsetroundjoin%
\definecolor{currentfill}{rgb}{0.556863,0.654902,0.847059}%
\pgfsetfillcolor{currentfill}%
\pgfsetfillopacity{0.130000}%
\pgfsetlinewidth{0.000000pt}%
\definecolor{currentstroke}{rgb}{0.000000,0.000000,0.000000}%
\pgfsetstrokecolor{currentstroke}%
\pgfsetdash{}{0pt}%
\pgfpathmoveto{\pgfqpoint{0.830257in}{0.894149in}}%
\pgfpathlineto{\pgfqpoint{0.846371in}{0.876637in}}%
\pgfpathlineto{\pgfqpoint{0.862494in}{0.859543in}}%
\pgfpathlineto{\pgfqpoint{0.870959in}{0.860605in}}%
\pgfpathlineto{\pgfqpoint{0.879386in}{0.862286in}}%
\pgfpathlineto{\pgfqpoint{0.863305in}{0.879715in}}%
\pgfpathlineto{\pgfqpoint{0.847233in}{0.897533in}}%
\pgfpathlineto{\pgfqpoint{0.838764in}{0.895558in}}%
\pgfpathlineto{\pgfqpoint{0.830257in}{0.894149in}}%
\pgfpathclose%
\pgfusepath{fill}%
\end{pgfscope}%
\begin{pgfscope}%
\pgfpathrectangle{\pgfqpoint{0.407592in}{0.222904in}}{\pgfqpoint{1.411163in}{1.411163in}}%
\pgfusepath{clip}%
\pgfsetbuttcap%
\pgfsetroundjoin%
\definecolor{currentfill}{rgb}{0.556863,0.654902,0.847059}%
\pgfsetfillcolor{currentfill}%
\pgfsetfillopacity{0.130000}%
\pgfsetlinewidth{0.000000pt}%
\definecolor{currentstroke}{rgb}{0.000000,0.000000,0.000000}%
\pgfsetstrokecolor{currentstroke}%
\pgfsetdash{}{0pt}%
\pgfpathmoveto{\pgfqpoint{0.684181in}{1.040255in}}%
\pgfpathlineto{\pgfqpoint{0.700269in}{1.021150in}}%
\pgfpathlineto{\pgfqpoint{0.716365in}{1.002191in}}%
\pgfpathlineto{\pgfqpoint{0.725079in}{1.003467in}}%
\pgfpathlineto{\pgfqpoint{0.733752in}{1.005162in}}%
\pgfpathlineto{\pgfqpoint{0.717698in}{1.024330in}}%
\pgfpathlineto{\pgfqpoint{0.701653in}{1.043626in}}%
\pgfpathlineto{\pgfqpoint{0.692937in}{1.041746in}}%
\pgfpathlineto{\pgfqpoint{0.684181in}{1.040255in}}%
\pgfpathclose%
\pgfusepath{fill}%
\end{pgfscope}%
\begin{pgfscope}%
\pgfpathrectangle{\pgfqpoint{0.407592in}{0.222904in}}{\pgfqpoint{1.411163in}{1.411163in}}%
\pgfusepath{clip}%
\pgfsetbuttcap%
\pgfsetroundjoin%
\definecolor{currentfill}{rgb}{0.556863,0.654902,0.847059}%
\pgfsetfillcolor{currentfill}%
\pgfsetfillopacity{0.130000}%
\pgfsetlinewidth{0.000000pt}%
\definecolor{currentstroke}{rgb}{0.000000,0.000000,0.000000}%
\pgfsetstrokecolor{currentstroke}%
\pgfsetdash{}{0pt}%
\pgfpathmoveto{\pgfqpoint{1.218538in}{0.790440in}}%
\pgfpathlineto{\pgfqpoint{1.234939in}{0.795332in}}%
\pgfpathlineto{\pgfqpoint{1.251404in}{0.801285in}}%
\pgfpathlineto{\pgfqpoint{1.259252in}{0.811438in}}%
\pgfpathlineto{\pgfqpoint{1.267086in}{0.822236in}}%
\pgfpathlineto{\pgfqpoint{1.250637in}{0.816806in}}%
\pgfpathlineto{\pgfqpoint{1.234253in}{0.812426in}}%
\pgfpathlineto{\pgfqpoint{1.226403in}{0.801094in}}%
\pgfpathlineto{\pgfqpoint{1.218538in}{0.790440in}}%
\pgfpathclose%
\pgfusepath{fill}%
\end{pgfscope}%
\begin{pgfscope}%
\pgfpathrectangle{\pgfqpoint{0.407592in}{0.222904in}}{\pgfqpoint{1.411163in}{1.411163in}}%
\pgfusepath{clip}%
\pgfsetbuttcap%
\pgfsetroundjoin%
\definecolor{currentfill}{rgb}{0.556863,0.654902,0.847059}%
\pgfsetfillcolor{currentfill}%
\pgfsetfillopacity{0.130000}%
\pgfsetlinewidth{0.000000pt}%
\definecolor{currentstroke}{rgb}{0.000000,0.000000,0.000000}%
\pgfsetstrokecolor{currentstroke}%
\pgfsetdash{}{0pt}%
\pgfpathmoveto{\pgfqpoint{1.331310in}{0.882011in}}%
\pgfpathlineto{\pgfqpoint{1.347926in}{0.888740in}}%
\pgfpathlineto{\pgfqpoint{1.364611in}{0.896115in}}%
\pgfpathlineto{\pgfqpoint{1.372344in}{0.908210in}}%
\pgfpathlineto{\pgfqpoint{1.380064in}{0.920672in}}%
\pgfpathlineto{\pgfqpoint{1.363396in}{0.913858in}}%
\pgfpathlineto{\pgfqpoint{1.346797in}{0.907699in}}%
\pgfpathlineto{\pgfqpoint{1.339060in}{0.894674in}}%
\pgfpathlineto{\pgfqpoint{1.331310in}{0.882011in}}%
\pgfpathclose%
\pgfusepath{fill}%
\end{pgfscope}%
\begin{pgfscope}%
\pgfpathrectangle{\pgfqpoint{0.407592in}{0.222904in}}{\pgfqpoint{1.411163in}{1.411163in}}%
\pgfusepath{clip}%
\pgfsetbuttcap%
\pgfsetroundjoin%
\definecolor{currentfill}{rgb}{0.556863,0.654902,0.847059}%
\pgfsetfillcolor{currentfill}%
\pgfsetfillopacity{0.130000}%
\pgfsetlinewidth{0.000000pt}%
\definecolor{currentstroke}{rgb}{0.000000,0.000000,0.000000}%
\pgfsetstrokecolor{currentstroke}%
\pgfsetdash{}{0pt}%
\pgfpathmoveto{\pgfqpoint{1.508491in}{1.055016in}}%
\pgfpathlineto{\pgfqpoint{1.525440in}{1.062400in}}%
\pgfpathlineto{\pgfqpoint{1.542460in}{1.070200in}}%
\pgfpathlineto{\pgfqpoint{1.550021in}{1.083499in}}%
\pgfpathlineto{\pgfqpoint{1.557571in}{1.096981in}}%
\pgfpathlineto{\pgfqpoint{1.540569in}{1.089608in}}%
\pgfpathlineto{\pgfqpoint{1.523638in}{1.082659in}}%
\pgfpathlineto{\pgfqpoint{1.516070in}{1.068750in}}%
\pgfpathlineto{\pgfqpoint{1.508491in}{1.055016in}}%
\pgfpathclose%
\pgfusepath{fill}%
\end{pgfscope}%
\begin{pgfscope}%
\pgfpathrectangle{\pgfqpoint{0.407592in}{0.222904in}}{\pgfqpoint{1.411163in}{1.411163in}}%
\pgfusepath{clip}%
\pgfsetbuttcap%
\pgfsetroundjoin%
\definecolor{currentfill}{rgb}{0.556863,0.654902,0.847059}%
\pgfsetfillcolor{currentfill}%
\pgfsetfillopacity{0.130000}%
\pgfsetlinewidth{0.000000pt}%
\definecolor{currentstroke}{rgb}{0.000000,0.000000,0.000000}%
\pgfsetstrokecolor{currentstroke}%
\pgfsetdash{}{0pt}%
\pgfpathmoveto{\pgfqpoint{1.572637in}{1.124449in}}%
\pgfpathlineto{\pgfqpoint{1.589692in}{1.131815in}}%
\pgfpathlineto{\pgfqpoint{1.606818in}{1.139560in}}%
\pgfpathlineto{\pgfqpoint{1.614316in}{1.153133in}}%
\pgfpathlineto{\pgfqpoint{1.597198in}{1.145582in}}%
\pgfpathlineto{\pgfqpoint{1.580153in}{1.138411in}}%
\pgfusepath{fill}%
\end{pgfscope}%
\begin{pgfscope}%
\pgfpathrectangle{\pgfqpoint{0.407592in}{0.222904in}}{\pgfqpoint{1.411163in}{1.411163in}}%
\pgfusepath{clip}%
\pgfsetbuttcap%
\pgfsetroundjoin%
\definecolor{currentfill}{rgb}{0.556863,0.654902,0.847059}%
\pgfsetfillcolor{currentfill}%
\pgfsetfillopacity{0.130000}%
\pgfsetlinewidth{0.000000pt}%
\definecolor{currentstroke}{rgb}{0.000000,0.000000,0.000000}%
\pgfsetstrokecolor{currentstroke}%
\pgfsetdash{}{0pt}%
\pgfpathmoveto{\pgfqpoint{0.716365in}{1.002191in}}%
\pgfpathlineto{\pgfqpoint{0.732470in}{0.983400in}}%
\pgfpathlineto{\pgfqpoint{0.748584in}{0.964799in}}%
\pgfpathlineto{\pgfqpoint{0.757255in}{0.965817in}}%
\pgfpathlineto{\pgfqpoint{0.765886in}{0.967286in}}%
\pgfpathlineto{\pgfqpoint{0.749815in}{0.986140in}}%
\pgfpathlineto{\pgfqpoint{0.733752in}{1.005162in}}%
\pgfpathlineto{\pgfqpoint{0.725079in}{1.003467in}}%
\pgfpathlineto{\pgfqpoint{0.716365in}{1.002191in}}%
\pgfpathclose%
\pgfusepath{fill}%
\end{pgfscope}%
\begin{pgfscope}%
\pgfpathrectangle{\pgfqpoint{0.407592in}{0.222904in}}{\pgfqpoint{1.411163in}{1.411163in}}%
\pgfusepath{clip}%
\pgfsetbuttcap%
\pgfsetroundjoin%
\definecolor{currentfill}{rgb}{0.556863,0.654902,0.847059}%
\pgfsetfillcolor{currentfill}%
\pgfsetfillopacity{0.130000}%
\pgfsetlinewidth{0.000000pt}%
\definecolor{currentstroke}{rgb}{0.000000,0.000000,0.000000}%
\pgfsetstrokecolor{currentstroke}%
\pgfsetdash{}{0pt}%
\pgfpathmoveto{\pgfqpoint{0.862494in}{0.859543in}}%
\pgfpathlineto{\pgfqpoint{0.878626in}{0.842941in}}%
\pgfpathlineto{\pgfqpoint{0.894767in}{0.826919in}}%
\pgfpathlineto{\pgfqpoint{0.903190in}{0.827572in}}%
\pgfpathlineto{\pgfqpoint{0.911574in}{0.828902in}}%
\pgfpathlineto{\pgfqpoint{0.895476in}{0.845319in}}%
\pgfpathlineto{\pgfqpoint{0.879386in}{0.862286in}}%
\pgfpathlineto{\pgfqpoint{0.870959in}{0.860605in}}%
\pgfpathlineto{\pgfqpoint{0.862494in}{0.859543in}}%
\pgfpathclose%
\pgfusepath{fill}%
\end{pgfscope}%
\begin{pgfscope}%
\pgfpathrectangle{\pgfqpoint{0.407592in}{0.222904in}}{\pgfqpoint{1.411163in}{1.411163in}}%
\pgfusepath{clip}%
\pgfsetbuttcap%
\pgfsetroundjoin%
\definecolor{currentfill}{rgb}{0.556863,0.654902,0.847059}%
\pgfsetfillcolor{currentfill}%
\pgfsetfillopacity{0.130000}%
\pgfsetlinewidth{0.000000pt}%
\definecolor{currentstroke}{rgb}{0.000000,0.000000,0.000000}%
\pgfsetstrokecolor{currentstroke}%
\pgfsetdash{}{0pt}%
\pgfpathmoveto{\pgfqpoint{1.444301in}{0.986838in}}%
\pgfpathlineto{\pgfqpoint{1.461144in}{0.994243in}}%
\pgfpathlineto{\pgfqpoint{1.478058in}{1.002116in}}%
\pgfpathlineto{\pgfqpoint{1.485684in}{1.015004in}}%
\pgfpathlineto{\pgfqpoint{1.493298in}{1.028128in}}%
\pgfpathlineto{\pgfqpoint{1.476402in}{1.020733in}}%
\pgfpathlineto{\pgfqpoint{1.459576in}{1.013815in}}%
\pgfpathlineto{\pgfqpoint{1.451945in}{1.000213in}}%
\pgfpathlineto{\pgfqpoint{1.444301in}{0.986838in}}%
\pgfpathclose%
\pgfusepath{fill}%
\end{pgfscope}%
\begin{pgfscope}%
\pgfpathrectangle{\pgfqpoint{0.407592in}{0.222904in}}{\pgfqpoint{1.411163in}{1.411163in}}%
\pgfusepath{clip}%
\pgfsetbuttcap%
\pgfsetroundjoin%
\definecolor{currentfill}{rgb}{0.556863,0.654902,0.847059}%
\pgfsetfillcolor{currentfill}%
\pgfsetfillopacity{0.130000}%
\pgfsetlinewidth{0.000000pt}%
\definecolor{currentstroke}{rgb}{0.000000,0.000000,0.000000}%
\pgfsetstrokecolor{currentstroke}%
\pgfsetdash{}{0pt}%
\pgfpathmoveto{\pgfqpoint{1.267086in}{0.822236in}}%
\pgfpathlineto{\pgfqpoint{1.283600in}{0.828585in}}%
\pgfpathlineto{\pgfqpoint{1.300181in}{0.835734in}}%
\pgfpathlineto{\pgfqpoint{1.307983in}{0.846565in}}%
\pgfpathlineto{\pgfqpoint{1.315772in}{0.857916in}}%
\pgfpathlineto{\pgfqpoint{1.299207in}{0.851329in}}%
\pgfpathlineto{\pgfqpoint{1.282710in}{0.845550in}}%
\pgfpathlineto{\pgfqpoint{1.274905in}{0.833625in}}%
\pgfpathlineto{\pgfqpoint{1.267086in}{0.822236in}}%
\pgfpathclose%
\pgfusepath{fill}%
\end{pgfscope}%
\begin{pgfscope}%
\pgfpathrectangle{\pgfqpoint{0.407592in}{0.222904in}}{\pgfqpoint{1.411163in}{1.411163in}}%
\pgfusepath{clip}%
\pgfsetbuttcap%
\pgfsetroundjoin%
\definecolor{currentfill}{rgb}{0.556863,0.654902,0.847059}%
\pgfsetfillcolor{currentfill}%
\pgfsetfillopacity{0.130000}%
\pgfsetlinewidth{0.000000pt}%
\definecolor{currentstroke}{rgb}{0.000000,0.000000,0.000000}%
\pgfsetstrokecolor{currentstroke}%
\pgfsetdash{}{0pt}%
\pgfpathmoveto{\pgfqpoint{0.569946in}{1.154195in}}%
\pgfpathlineto{\pgfqpoint{0.586025in}{1.134653in}}%
\pgfpathlineto{\pgfqpoint{0.602112in}{1.115194in}}%
\pgfpathlineto{\pgfqpoint{0.611034in}{1.116343in}}%
\pgfpathlineto{\pgfqpoint{0.619915in}{1.117823in}}%
\pgfpathlineto{\pgfqpoint{0.603870in}{1.137440in}}%
\pgfpathlineto{\pgfqpoint{0.587834in}{1.157126in}}%
\pgfpathlineto{\pgfqpoint{0.578910in}{1.155504in}}%
\pgfpathlineto{\pgfqpoint{0.569946in}{1.154195in}}%
\pgfpathclose%
\pgfusepath{fill}%
\end{pgfscope}%
\begin{pgfscope}%
\pgfpathrectangle{\pgfqpoint{0.407592in}{0.222904in}}{\pgfqpoint{1.411163in}{1.411163in}}%
\pgfusepath{clip}%
\pgfsetbuttcap%
\pgfsetroundjoin%
\definecolor{currentfill}{rgb}{0.556863,0.654902,0.847059}%
\pgfsetfillcolor{currentfill}%
\pgfsetfillopacity{0.130000}%
\pgfsetlinewidth{0.000000pt}%
\definecolor{currentstroke}{rgb}{0.000000,0.000000,0.000000}%
\pgfsetstrokecolor{currentstroke}%
\pgfsetdash{}{0pt}%
\pgfpathmoveto{\pgfqpoint{0.976071in}{0.771175in}}%
\pgfpathlineto{\pgfqpoint{0.992225in}{0.759518in}}%
\pgfpathlineto{\pgfqpoint{1.008395in}{0.749341in}}%
\pgfpathlineto{\pgfqpoint{1.016618in}{0.750384in}}%
\pgfpathlineto{\pgfqpoint{1.024806in}{0.752405in}}%
\pgfpathlineto{\pgfqpoint{1.008677in}{0.762842in}}%
\pgfpathlineto{\pgfqpoint{0.992563in}{0.774795in}}%
\pgfpathlineto{\pgfqpoint{0.984335in}{0.772531in}}%
\pgfpathlineto{\pgfqpoint{0.976071in}{0.771175in}}%
\pgfpathclose%
\pgfusepath{fill}%
\end{pgfscope}%
\begin{pgfscope}%
\pgfpathrectangle{\pgfqpoint{0.407592in}{0.222904in}}{\pgfqpoint{1.411163in}{1.411163in}}%
\pgfusepath{clip}%
\pgfsetbuttcap%
\pgfsetroundjoin%
\definecolor{currentfill}{rgb}{0.556863,0.654902,0.847059}%
\pgfsetfillcolor{currentfill}%
\pgfsetfillopacity{0.130000}%
\pgfsetlinewidth{0.000000pt}%
\definecolor{currentstroke}{rgb}{0.000000,0.000000,0.000000}%
\pgfsetstrokecolor{currentstroke}%
\pgfsetdash{}{0pt}%
\pgfpathmoveto{\pgfqpoint{0.748584in}{0.964799in}}%
\pgfpathlineto{\pgfqpoint{0.764706in}{0.946418in}}%
\pgfpathlineto{\pgfqpoint{0.780837in}{0.928288in}}%
\pgfpathlineto{\pgfqpoint{0.789466in}{0.928997in}}%
\pgfpathlineto{\pgfqpoint{0.798055in}{0.930191in}}%
\pgfpathlineto{\pgfqpoint{0.781966in}{0.948626in}}%
\pgfpathlineto{\pgfqpoint{0.765886in}{0.967286in}}%
\pgfpathlineto{\pgfqpoint{0.757255in}{0.965817in}}%
\pgfpathlineto{\pgfqpoint{0.748584in}{0.964799in}}%
\pgfpathclose%
\pgfusepath{fill}%
\end{pgfscope}%
\begin{pgfscope}%
\pgfpathrectangle{\pgfqpoint{0.407592in}{0.222904in}}{\pgfqpoint{1.411163in}{1.411163in}}%
\pgfusepath{clip}%
\pgfsetbuttcap%
\pgfsetroundjoin%
\definecolor{currentfill}{rgb}{0.556863,0.654902,0.847059}%
\pgfsetfillcolor{currentfill}%
\pgfsetfillopacity{0.130000}%
\pgfsetlinewidth{0.000000pt}%
\definecolor{currentstroke}{rgb}{0.000000,0.000000,0.000000}%
\pgfsetstrokecolor{currentstroke}%
\pgfsetdash{}{0pt}%
\pgfpathmoveto{\pgfqpoint{1.105701in}{0.736477in}}%
\pgfpathlineto{\pgfqpoint{1.121944in}{0.735796in}}%
\pgfpathlineto{\pgfqpoint{1.138235in}{0.737000in}}%
\pgfpathlineto{\pgfqpoint{1.146236in}{0.742334in}}%
\pgfpathlineto{\pgfqpoint{1.154213in}{0.748759in}}%
\pgfpathlineto{\pgfqpoint{1.137947in}{0.747584in}}%
\pgfpathlineto{\pgfqpoint{1.121729in}{0.748276in}}%
\pgfpathlineto{\pgfqpoint{1.113728in}{0.741796in}}%
\pgfpathlineto{\pgfqpoint{1.105701in}{0.736477in}}%
\pgfpathclose%
\pgfusepath{fill}%
\end{pgfscope}%
\begin{pgfscope}%
\pgfpathrectangle{\pgfqpoint{0.407592in}{0.222904in}}{\pgfqpoint{1.411163in}{1.411163in}}%
\pgfusepath{clip}%
\pgfsetbuttcap%
\pgfsetroundjoin%
\definecolor{currentfill}{rgb}{0.556863,0.654902,0.847059}%
\pgfsetfillcolor{currentfill}%
\pgfsetfillopacity{0.130000}%
\pgfsetlinewidth{0.000000pt}%
\definecolor{currentstroke}{rgb}{0.000000,0.000000,0.000000}%
\pgfsetstrokecolor{currentstroke}%
\pgfsetdash{}{0pt}%
\pgfpathmoveto{\pgfqpoint{1.380064in}{0.920672in}}%
\pgfpathlineto{\pgfqpoint{1.396802in}{0.928076in}}%
\pgfpathlineto{\pgfqpoint{1.413609in}{0.936014in}}%
\pgfpathlineto{\pgfqpoint{1.421300in}{0.948275in}}%
\pgfpathlineto{\pgfqpoint{1.428979in}{0.960848in}}%
\pgfpathlineto{\pgfqpoint{1.412189in}{0.953437in}}%
\pgfpathlineto{\pgfqpoint{1.395469in}{0.946572in}}%
\pgfpathlineto{\pgfqpoint{1.387773in}{0.933468in}}%
\pgfpathlineto{\pgfqpoint{1.380064in}{0.920672in}}%
\pgfpathclose%
\pgfusepath{fill}%
\end{pgfscope}%
\begin{pgfscope}%
\pgfpathrectangle{\pgfqpoint{0.407592in}{0.222904in}}{\pgfqpoint{1.411163in}{1.411163in}}%
\pgfusepath{clip}%
\pgfsetbuttcap%
\pgfsetroundjoin%
\definecolor{currentfill}{rgb}{0.556863,0.654902,0.847059}%
\pgfsetfillcolor{currentfill}%
\pgfsetfillopacity{0.130000}%
\pgfsetlinewidth{0.000000pt}%
\definecolor{currentstroke}{rgb}{0.000000,0.000000,0.000000}%
\pgfsetstrokecolor{currentstroke}%
\pgfsetdash{}{0pt}%
\pgfpathmoveto{\pgfqpoint{1.154213in}{0.748759in}}%
\pgfpathlineto{\pgfqpoint{1.170532in}{0.751611in}}%
\pgfpathlineto{\pgfqpoint{1.186909in}{0.755929in}}%
\pgfpathlineto{\pgfqpoint{1.194844in}{0.763221in}}%
\pgfpathlineto{\pgfqpoint{1.202760in}{0.771439in}}%
\pgfpathlineto{\pgfqpoint{1.186402in}{0.767373in}}%
\pgfpathlineto{\pgfqpoint{1.170103in}{0.764745in}}%
\pgfpathlineto{\pgfqpoint{1.162168in}{0.756246in}}%
\pgfpathlineto{\pgfqpoint{1.154213in}{0.748759in}}%
\pgfpathclose%
\pgfusepath{fill}%
\end{pgfscope}%
\begin{pgfscope}%
\pgfpathrectangle{\pgfqpoint{0.407592in}{0.222904in}}{\pgfqpoint{1.411163in}{1.411163in}}%
\pgfusepath{clip}%
\pgfsetbuttcap%
\pgfsetroundjoin%
\definecolor{currentfill}{rgb}{0.556863,0.654902,0.847059}%
\pgfsetfillcolor{currentfill}%
\pgfsetfillopacity{0.130000}%
\pgfsetlinewidth{0.000000pt}%
\definecolor{currentstroke}{rgb}{0.000000,0.000000,0.000000}%
\pgfsetstrokecolor{currentstroke}%
\pgfsetdash{}{0pt}%
\pgfpathmoveto{\pgfqpoint{1.557571in}{1.096981in}}%
\pgfpathlineto{\pgfqpoint{1.574644in}{1.104752in}}%
\pgfpathlineto{\pgfqpoint{1.591791in}{1.112895in}}%
\pgfpathlineto{\pgfqpoint{1.599310in}{1.126143in}}%
\pgfpathlineto{\pgfqpoint{1.606818in}{1.139560in}}%
\pgfpathlineto{\pgfqpoint{1.589692in}{1.131815in}}%
\pgfpathlineto{\pgfqpoint{1.572637in}{1.124449in}}%
\pgfpathlineto{\pgfqpoint{1.565110in}{1.110635in}}%
\pgfpathlineto{\pgfqpoint{1.557571in}{1.096981in}}%
\pgfpathclose%
\pgfusepath{fill}%
\end{pgfscope}%
\begin{pgfscope}%
\pgfpathrectangle{\pgfqpoint{0.407592in}{0.222904in}}{\pgfqpoint{1.411163in}{1.411163in}}%
\pgfusepath{clip}%
\pgfsetbuttcap%
\pgfsetroundjoin%
\definecolor{currentfill}{rgb}{0.556863,0.654902,0.847059}%
\pgfsetfillcolor{currentfill}%
\pgfsetfillopacity{0.130000}%
\pgfsetlinewidth{0.000000pt}%
\definecolor{currentstroke}{rgb}{0.000000,0.000000,0.000000}%
\pgfsetstrokecolor{currentstroke}%
\pgfsetdash{}{0pt}%
\pgfpathmoveto{\pgfqpoint{0.602112in}{1.115194in}}%
\pgfpathlineto{\pgfqpoint{0.618208in}{1.095828in}}%
\pgfpathlineto{\pgfqpoint{0.634313in}{1.076567in}}%
\pgfpathlineto{\pgfqpoint{0.643192in}{1.077527in}}%
\pgfpathlineto{\pgfqpoint{0.652031in}{1.078837in}}%
\pgfpathlineto{\pgfqpoint{0.635969in}{1.098286in}}%
\pgfpathlineto{\pgfqpoint{0.619915in}{1.117823in}}%
\pgfpathlineto{\pgfqpoint{0.611034in}{1.116343in}}%
\pgfpathlineto{\pgfqpoint{0.602112in}{1.115194in}}%
\pgfpathclose%
\pgfusepath{fill}%
\end{pgfscope}%
\begin{pgfscope}%
\pgfpathrectangle{\pgfqpoint{0.407592in}{0.222904in}}{\pgfqpoint{1.411163in}{1.411163in}}%
\pgfusepath{clip}%
\pgfsetbuttcap%
\pgfsetroundjoin%
\definecolor{currentfill}{rgb}{0.556863,0.654902,0.847059}%
\pgfsetfillcolor{currentfill}%
\pgfsetfillopacity{0.130000}%
\pgfsetlinewidth{0.000000pt}%
\definecolor{currentstroke}{rgb}{0.000000,0.000000,0.000000}%
\pgfsetstrokecolor{currentstroke}%
\pgfsetdash{}{0pt}%
\pgfpathmoveto{\pgfqpoint{0.894767in}{0.826919in}}%
\pgfpathlineto{\pgfqpoint{0.910917in}{0.811580in}}%
\pgfpathlineto{\pgfqpoint{0.927077in}{0.797044in}}%
\pgfpathlineto{\pgfqpoint{0.935458in}{0.797237in}}%
\pgfpathlineto{\pgfqpoint{0.943801in}{0.798165in}}%
\pgfpathlineto{\pgfqpoint{0.927683in}{0.813142in}}%
\pgfpathlineto{\pgfqpoint{0.911574in}{0.828902in}}%
\pgfpathlineto{\pgfqpoint{0.903190in}{0.827572in}}%
\pgfpathlineto{\pgfqpoint{0.894767in}{0.826919in}}%
\pgfpathclose%
\pgfusepath{fill}%
\end{pgfscope}%
\begin{pgfscope}%
\pgfpathrectangle{\pgfqpoint{0.407592in}{0.222904in}}{\pgfqpoint{1.411163in}{1.411163in}}%
\pgfusepath{clip}%
\pgfsetbuttcap%
\pgfsetroundjoin%
\definecolor{currentfill}{rgb}{0.556863,0.654902,0.847059}%
\pgfsetfillcolor{currentfill}%
\pgfsetfillopacity{0.130000}%
\pgfsetlinewidth{0.000000pt}%
\definecolor{currentstroke}{rgb}{0.000000,0.000000,0.000000}%
\pgfsetstrokecolor{currentstroke}%
\pgfsetdash{}{0pt}%
\pgfpathmoveto{\pgfqpoint{1.057128in}{0.736877in}}%
\pgfpathlineto{\pgfqpoint{1.073329in}{0.732072in}}%
\pgfpathlineto{\pgfqpoint{1.089563in}{0.729307in}}%
\pgfpathlineto{\pgfqpoint{1.097646in}{0.732322in}}%
\pgfpathlineto{\pgfqpoint{1.105701in}{0.736477in}}%
\pgfpathlineto{\pgfqpoint{1.089498in}{0.739185in}}%
\pgfpathlineto{\pgfqpoint{1.073331in}{0.743981in}}%
\pgfpathlineto{\pgfqpoint{1.065245in}{0.739857in}}%
\pgfpathlineto{\pgfqpoint{1.057128in}{0.736877in}}%
\pgfpathclose%
\pgfusepath{fill}%
\end{pgfscope}%
\begin{pgfscope}%
\pgfpathrectangle{\pgfqpoint{0.407592in}{0.222904in}}{\pgfqpoint{1.411163in}{1.411163in}}%
\pgfusepath{clip}%
\pgfsetbuttcap%
\pgfsetroundjoin%
\definecolor{currentfill}{rgb}{0.556863,0.654902,0.847059}%
\pgfsetfillcolor{currentfill}%
\pgfsetfillopacity{0.130000}%
\pgfsetlinewidth{0.000000pt}%
\definecolor{currentstroke}{rgb}{0.000000,0.000000,0.000000}%
\pgfsetstrokecolor{currentstroke}%
\pgfsetdash{}{0pt}%
\pgfpathmoveto{\pgfqpoint{0.780837in}{0.928288in}}%
\pgfpathlineto{\pgfqpoint{0.796976in}{0.910448in}}%
\pgfpathlineto{\pgfqpoint{0.813124in}{0.892943in}}%
\pgfpathlineto{\pgfqpoint{0.821710in}{0.893285in}}%
\pgfpathlineto{\pgfqpoint{0.830257in}{0.894149in}}%
\pgfpathlineto{\pgfqpoint{0.814151in}{0.912018in}}%
\pgfpathlineto{\pgfqpoint{0.798055in}{0.930191in}}%
\pgfpathlineto{\pgfqpoint{0.789466in}{0.928997in}}%
\pgfpathlineto{\pgfqpoint{0.780837in}{0.928288in}}%
\pgfpathclose%
\pgfusepath{fill}%
\end{pgfscope}%
\begin{pgfscope}%
\pgfpathrectangle{\pgfqpoint{0.407592in}{0.222904in}}{\pgfqpoint{1.411163in}{1.411163in}}%
\pgfusepath{clip}%
\pgfsetbuttcap%
\pgfsetroundjoin%
\definecolor{currentfill}{rgb}{0.556863,0.654902,0.847059}%
\pgfsetfillcolor{currentfill}%
\pgfsetfillopacity{0.130000}%
\pgfsetlinewidth{0.000000pt}%
\definecolor{currentstroke}{rgb}{0.000000,0.000000,0.000000}%
\pgfsetstrokecolor{currentstroke}%
\pgfsetdash{}{0pt}%
\pgfpathmoveto{\pgfqpoint{1.315772in}{0.857916in}}%
\pgfpathlineto{\pgfqpoint{1.332405in}{0.865216in}}%
\pgfpathlineto{\pgfqpoint{1.349107in}{0.873144in}}%
\pgfpathlineto{\pgfqpoint{1.356865in}{0.884415in}}%
\pgfpathlineto{\pgfqpoint{1.364611in}{0.896115in}}%
\pgfpathlineto{\pgfqpoint{1.347926in}{0.888740in}}%
\pgfpathlineto{\pgfqpoint{1.331310in}{0.882011in}}%
\pgfpathlineto{\pgfqpoint{1.323547in}{0.869745in}}%
\pgfpathlineto{\pgfqpoint{1.315772in}{0.857916in}}%
\pgfpathclose%
\pgfusepath{fill}%
\end{pgfscope}%
\begin{pgfscope}%
\pgfpathrectangle{\pgfqpoint{0.407592in}{0.222904in}}{\pgfqpoint{1.411163in}{1.411163in}}%
\pgfusepath{clip}%
\pgfsetbuttcap%
\pgfsetroundjoin%
\definecolor{currentfill}{rgb}{0.556863,0.654902,0.847059}%
\pgfsetfillcolor{currentfill}%
\pgfsetfillopacity{0.130000}%
\pgfsetlinewidth{0.000000pt}%
\definecolor{currentstroke}{rgb}{0.000000,0.000000,0.000000}%
\pgfsetstrokecolor{currentstroke}%
\pgfsetdash{}{0pt}%
\pgfpathmoveto{\pgfqpoint{1.493298in}{1.028128in}}%
\pgfpathlineto{\pgfqpoint{1.510265in}{1.035964in}}%
\pgfpathlineto{\pgfqpoint{1.527304in}{1.044208in}}%
\pgfpathlineto{\pgfqpoint{1.534888in}{1.057099in}}%
\pgfpathlineto{\pgfqpoint{1.542460in}{1.070200in}}%
\pgfpathlineto{\pgfqpoint{1.525440in}{1.062400in}}%
\pgfpathlineto{\pgfqpoint{1.508491in}{1.055016in}}%
\pgfpathlineto{\pgfqpoint{1.500900in}{1.041471in}}%
\pgfpathlineto{\pgfqpoint{1.493298in}{1.028128in}}%
\pgfpathclose%
\pgfusepath{fill}%
\end{pgfscope}%
\begin{pgfscope}%
\pgfpathrectangle{\pgfqpoint{0.407592in}{0.222904in}}{\pgfqpoint{1.411163in}{1.411163in}}%
\pgfusepath{clip}%
\pgfsetbuttcap%
\pgfsetroundjoin%
\definecolor{currentfill}{rgb}{0.556863,0.654902,0.847059}%
\pgfsetfillcolor{currentfill}%
\pgfsetfillopacity{0.130000}%
\pgfsetlinewidth{0.000000pt}%
\definecolor{currentstroke}{rgb}{0.000000,0.000000,0.000000}%
\pgfsetstrokecolor{currentstroke}%
\pgfsetdash{}{0pt}%
\pgfpathmoveto{\pgfqpoint{1.202760in}{0.771439in}}%
\pgfpathlineto{\pgfqpoint{1.219178in}{0.776751in}}%
\pgfpathlineto{\pgfqpoint{1.235661in}{0.783126in}}%
\pgfpathlineto{\pgfqpoint{1.243541in}{0.791830in}}%
\pgfpathlineto{\pgfqpoint{1.251404in}{0.801285in}}%
\pgfpathlineto{\pgfqpoint{1.234939in}{0.795332in}}%
\pgfpathlineto{\pgfqpoint{1.218538in}{0.790440in}}%
\pgfpathlineto{\pgfqpoint{1.210657in}{0.780533in}}%
\pgfpathlineto{\pgfqpoint{1.202760in}{0.771439in}}%
\pgfpathclose%
\pgfusepath{fill}%
\end{pgfscope}%
\begin{pgfscope}%
\pgfpathrectangle{\pgfqpoint{0.407592in}{0.222904in}}{\pgfqpoint{1.411163in}{1.411163in}}%
\pgfusepath{clip}%
\pgfsetbuttcap%
\pgfsetroundjoin%
\definecolor{currentfill}{rgb}{0.556863,0.654902,0.847059}%
\pgfsetfillcolor{currentfill}%
\pgfsetfillopacity{0.130000}%
\pgfsetlinewidth{0.000000pt}%
\definecolor{currentstroke}{rgb}{0.000000,0.000000,0.000000}%
\pgfsetstrokecolor{currentstroke}%
\pgfsetdash{}{0pt}%
\pgfpathmoveto{\pgfqpoint{0.634313in}{1.076567in}}%
\pgfpathlineto{\pgfqpoint{0.650426in}{1.057423in}}%
\pgfpathlineto{\pgfqpoint{0.666548in}{1.038411in}}%
\pgfpathlineto{\pgfqpoint{0.675385in}{1.039146in}}%
\pgfpathlineto{\pgfqpoint{0.684181in}{1.040255in}}%
\pgfpathlineto{\pgfqpoint{0.668101in}{1.059489in}}%
\pgfpathlineto{\pgfqpoint{0.652031in}{1.078837in}}%
\pgfpathlineto{\pgfqpoint{0.643192in}{1.077527in}}%
\pgfpathlineto{\pgfqpoint{0.634313in}{1.076567in}}%
\pgfpathclose%
\pgfusepath{fill}%
\end{pgfscope}%
\begin{pgfscope}%
\pgfpathrectangle{\pgfqpoint{0.407592in}{0.222904in}}{\pgfqpoint{1.411163in}{1.411163in}}%
\pgfusepath{clip}%
\pgfsetbuttcap%
\pgfsetroundjoin%
\definecolor{currentfill}{rgb}{0.556863,0.654902,0.847059}%
\pgfsetfillcolor{currentfill}%
\pgfsetfillopacity{0.130000}%
\pgfsetlinewidth{0.000000pt}%
\definecolor{currentstroke}{rgb}{0.000000,0.000000,0.000000}%
\pgfsetstrokecolor{currentstroke}%
\pgfsetdash{}{0pt}%
\pgfpathmoveto{\pgfqpoint{1.428979in}{0.960848in}}%
\pgfpathlineto{\pgfqpoint{1.445840in}{0.968756in}}%
\pgfpathlineto{\pgfqpoint{1.462772in}{0.977118in}}%
\pgfpathlineto{\pgfqpoint{1.470421in}{0.989481in}}%
\pgfpathlineto{\pgfqpoint{1.478058in}{1.002116in}}%
\pgfpathlineto{\pgfqpoint{1.461144in}{0.994243in}}%
\pgfpathlineto{\pgfqpoint{1.444301in}{0.986838in}}%
\pgfpathlineto{\pgfqpoint{1.436646in}{0.973709in}}%
\pgfpathlineto{\pgfqpoint{1.428979in}{0.960848in}}%
\pgfpathclose%
\pgfusepath{fill}%
\end{pgfscope}%
\begin{pgfscope}%
\pgfpathrectangle{\pgfqpoint{0.407592in}{0.222904in}}{\pgfqpoint{1.411163in}{1.411163in}}%
\pgfusepath{clip}%
\pgfsetbuttcap%
\pgfsetroundjoin%
\definecolor{currentfill}{rgb}{0.556863,0.654902,0.847059}%
\pgfsetfillcolor{currentfill}%
\pgfsetfillopacity{0.130000}%
\pgfsetlinewidth{0.000000pt}%
\definecolor{currentstroke}{rgb}{0.000000,0.000000,0.000000}%
\pgfsetstrokecolor{currentstroke}%
\pgfsetdash{}{0pt}%
\pgfpathmoveto{\pgfqpoint{1.008395in}{0.749341in}}%
\pgfpathlineto{\pgfqpoint{1.024584in}{0.740825in}}%
\pgfpathlineto{\pgfqpoint{1.040795in}{0.734121in}}%
\pgfpathlineto{\pgfqpoint{1.048978in}{0.734990in}}%
\pgfpathlineto{\pgfqpoint{1.057128in}{0.736877in}}%
\pgfpathlineto{\pgfqpoint{1.040956in}{0.743693in}}%
\pgfpathlineto{\pgfqpoint{1.024806in}{0.752405in}}%
\pgfpathlineto{\pgfqpoint{1.016618in}{0.750384in}}%
\pgfpathlineto{\pgfqpoint{1.008395in}{0.749341in}}%
\pgfpathclose%
\pgfusepath{fill}%
\end{pgfscope}%
\begin{pgfscope}%
\pgfpathrectangle{\pgfqpoint{0.407592in}{0.222904in}}{\pgfqpoint{1.411163in}{1.411163in}}%
\pgfusepath{clip}%
\pgfsetbuttcap%
\pgfsetroundjoin%
\definecolor{currentfill}{rgb}{0.556863,0.654902,0.847059}%
\pgfsetfillcolor{currentfill}%
\pgfsetfillopacity{0.130000}%
\pgfsetlinewidth{0.000000pt}%
\definecolor{currentstroke}{rgb}{0.000000,0.000000,0.000000}%
\pgfsetstrokecolor{currentstroke}%
\pgfsetdash{}{0pt}%
\pgfpathmoveto{\pgfqpoint{0.927077in}{0.797044in}}%
\pgfpathlineto{\pgfqpoint{0.943248in}{0.783450in}}%
\pgfpathlineto{\pgfqpoint{0.959431in}{0.770949in}}%
\pgfpathlineto{\pgfqpoint{0.967770in}{0.770670in}}%
\pgfpathlineto{\pgfqpoint{0.976071in}{0.771175in}}%
\pgfpathlineto{\pgfqpoint{0.959930in}{0.784119in}}%
\pgfpathlineto{\pgfqpoint{0.943801in}{0.798165in}}%
\pgfpathlineto{\pgfqpoint{0.935458in}{0.797237in}}%
\pgfpathlineto{\pgfqpoint{0.927077in}{0.797044in}}%
\pgfpathclose%
\pgfusepath{fill}%
\end{pgfscope}%
\begin{pgfscope}%
\pgfpathrectangle{\pgfqpoint{0.407592in}{0.222904in}}{\pgfqpoint{1.411163in}{1.411163in}}%
\pgfusepath{clip}%
\pgfsetbuttcap%
\pgfsetroundjoin%
\definecolor{currentfill}{rgb}{0.556863,0.654902,0.847059}%
\pgfsetfillcolor{currentfill}%
\pgfsetfillopacity{0.130000}%
\pgfsetlinewidth{0.000000pt}%
\definecolor{currentstroke}{rgb}{0.000000,0.000000,0.000000}%
\pgfsetstrokecolor{currentstroke}%
\pgfsetdash{}{0pt}%
\pgfpathmoveto{\pgfqpoint{0.813124in}{0.892943in}}%
\pgfpathlineto{\pgfqpoint{0.829280in}{0.875827in}}%
\pgfpathlineto{\pgfqpoint{0.845445in}{0.859161in}}%
\pgfpathlineto{\pgfqpoint{0.853990in}{0.859072in}}%
\pgfpathlineto{\pgfqpoint{0.862494in}{0.859543in}}%
\pgfpathlineto{\pgfqpoint{0.846371in}{0.876637in}}%
\pgfpathlineto{\pgfqpoint{0.830257in}{0.894149in}}%
\pgfpathlineto{\pgfqpoint{0.821710in}{0.893285in}}%
\pgfpathlineto{\pgfqpoint{0.813124in}{0.892943in}}%
\pgfpathclose%
\pgfusepath{fill}%
\end{pgfscope}%
\begin{pgfscope}%
\pgfpathrectangle{\pgfqpoint{0.407592in}{0.222904in}}{\pgfqpoint{1.411163in}{1.411163in}}%
\pgfusepath{clip}%
\pgfsetbuttcap%
\pgfsetroundjoin%
\definecolor{currentfill}{rgb}{0.556863,0.654902,0.847059}%
\pgfsetfillcolor{currentfill}%
\pgfsetfillopacity{0.130000}%
\pgfsetlinewidth{0.000000pt}%
\definecolor{currentstroke}{rgb}{0.000000,0.000000,0.000000}%
\pgfsetstrokecolor{currentstroke}%
\pgfsetdash{}{0pt}%
\pgfpathmoveto{\pgfqpoint{0.666548in}{1.038411in}}%
\pgfpathlineto{\pgfqpoint{0.682678in}{1.019548in}}%
\pgfpathlineto{\pgfqpoint{0.698817in}{1.000853in}}%
\pgfpathlineto{\pgfqpoint{0.707612in}{1.001324in}}%
\pgfpathlineto{\pgfqpoint{0.716365in}{1.002191in}}%
\pgfpathlineto{\pgfqpoint{0.700269in}{1.021150in}}%
\pgfpathlineto{\pgfqpoint{0.684181in}{1.040255in}}%
\pgfpathlineto{\pgfqpoint{0.675385in}{1.039146in}}%
\pgfpathlineto{\pgfqpoint{0.666548in}{1.038411in}}%
\pgfpathclose%
\pgfusepath{fill}%
\end{pgfscope}%
\begin{pgfscope}%
\pgfpathrectangle{\pgfqpoint{0.407592in}{0.222904in}}{\pgfqpoint{1.411163in}{1.411163in}}%
\pgfusepath{clip}%
\pgfsetbuttcap%
\pgfsetroundjoin%
\definecolor{currentfill}{rgb}{0.556863,0.654902,0.847059}%
\pgfsetfillcolor{currentfill}%
\pgfsetfillopacity{0.130000}%
\pgfsetlinewidth{0.000000pt}%
\definecolor{currentstroke}{rgb}{0.000000,0.000000,0.000000}%
\pgfsetstrokecolor{currentstroke}%
\pgfsetdash{}{0pt}%
\pgfpathmoveto{\pgfqpoint{1.251404in}{0.801285in}}%
\pgfpathlineto{\pgfqpoint{1.267935in}{0.808151in}}%
\pgfpathlineto{\pgfqpoint{1.284534in}{0.815801in}}%
\pgfpathlineto{\pgfqpoint{1.292365in}{0.825465in}}%
\pgfpathlineto{\pgfqpoint{1.300181in}{0.835734in}}%
\pgfpathlineto{\pgfqpoint{1.283600in}{0.828585in}}%
\pgfpathlineto{\pgfqpoint{1.267086in}{0.822236in}}%
\pgfpathlineto{\pgfqpoint{1.259252in}{0.811438in}}%
\pgfpathlineto{\pgfqpoint{1.251404in}{0.801285in}}%
\pgfpathclose%
\pgfusepath{fill}%
\end{pgfscope}%
\begin{pgfscope}%
\pgfpathrectangle{\pgfqpoint{0.407592in}{0.222904in}}{\pgfqpoint{1.411163in}{1.411163in}}%
\pgfusepath{clip}%
\pgfsetbuttcap%
\pgfsetroundjoin%
\definecolor{currentfill}{rgb}{0.556863,0.654902,0.847059}%
\pgfsetfillcolor{currentfill}%
\pgfsetfillopacity{0.130000}%
\pgfsetlinewidth{0.000000pt}%
\definecolor{currentstroke}{rgb}{0.000000,0.000000,0.000000}%
\pgfsetstrokecolor{currentstroke}%
\pgfsetdash{}{0pt}%
\pgfpathmoveto{\pgfqpoint{1.364611in}{0.896115in}}%
\pgfpathlineto{\pgfqpoint{1.381366in}{0.904064in}}%
\pgfpathlineto{\pgfqpoint{1.398191in}{0.912528in}}%
\pgfpathlineto{\pgfqpoint{1.405906in}{0.924090in}}%
\pgfpathlineto{\pgfqpoint{1.413609in}{0.936014in}}%
\pgfpathlineto{\pgfqpoint{1.396802in}{0.928076in}}%
\pgfpathlineto{\pgfqpoint{1.380064in}{0.920672in}}%
\pgfpathlineto{\pgfqpoint{1.372344in}{0.908210in}}%
\pgfpathlineto{\pgfqpoint{1.364611in}{0.896115in}}%
\pgfpathclose%
\pgfusepath{fill}%
\end{pgfscope}%
\begin{pgfscope}%
\pgfpathrectangle{\pgfqpoint{0.407592in}{0.222904in}}{\pgfqpoint{1.411163in}{1.411163in}}%
\pgfusepath{clip}%
\pgfsetbuttcap%
\pgfsetroundjoin%
\definecolor{currentfill}{rgb}{0.556863,0.654902,0.847059}%
\pgfsetfillcolor{currentfill}%
\pgfsetfillopacity{0.130000}%
\pgfsetlinewidth{0.000000pt}%
\definecolor{currentstroke}{rgb}{0.000000,0.000000,0.000000}%
\pgfsetstrokecolor{currentstroke}%
\pgfsetdash{}{0pt}%
\pgfpathmoveto{\pgfqpoint{1.606818in}{1.139560in}}%
\pgfpathlineto{\pgfqpoint{1.624018in}{1.147661in}}%
\pgfpathlineto{\pgfqpoint{1.641291in}{1.156100in}}%
\pgfpathlineto{\pgfqpoint{1.648768in}{1.169296in}}%
\pgfpathlineto{\pgfqpoint{1.631505in}{1.161045in}}%
\pgfpathlineto{\pgfqpoint{1.614316in}{1.153133in}}%
\pgfusepath{fill}%
\end{pgfscope}%
\begin{pgfscope}%
\pgfpathrectangle{\pgfqpoint{0.407592in}{0.222904in}}{\pgfqpoint{1.411163in}{1.411163in}}%
\pgfusepath{clip}%
\pgfsetbuttcap%
\pgfsetroundjoin%
\definecolor{currentfill}{rgb}{0.556863,0.654902,0.847059}%
\pgfsetfillcolor{currentfill}%
\pgfsetfillopacity{0.130000}%
\pgfsetlinewidth{0.000000pt}%
\definecolor{currentstroke}{rgb}{0.000000,0.000000,0.000000}%
\pgfsetstrokecolor{currentstroke}%
\pgfsetdash{}{0pt}%
\pgfpathmoveto{\pgfqpoint{1.542460in}{1.070200in}}%
\pgfpathlineto{\pgfqpoint{1.559553in}{1.078391in}}%
\pgfpathlineto{\pgfqpoint{1.576718in}{1.086945in}}%
\pgfpathlineto{\pgfqpoint{1.584260in}{1.099825in}}%
\pgfpathlineto{\pgfqpoint{1.591791in}{1.112895in}}%
\pgfpathlineto{\pgfqpoint{1.574644in}{1.104752in}}%
\pgfpathlineto{\pgfqpoint{1.557571in}{1.096981in}}%
\pgfpathlineto{\pgfqpoint{1.550021in}{1.083499in}}%
\pgfpathlineto{\pgfqpoint{1.542460in}{1.070200in}}%
\pgfpathclose%
\pgfusepath{fill}%
\end{pgfscope}%
\begin{pgfscope}%
\pgfpathrectangle{\pgfqpoint{0.407592in}{0.222904in}}{\pgfqpoint{1.411163in}{1.411163in}}%
\pgfusepath{clip}%
\pgfsetbuttcap%
\pgfsetroundjoin%
\definecolor{currentfill}{rgb}{0.556863,0.654902,0.847059}%
\pgfsetfillcolor{currentfill}%
\pgfsetfillopacity{0.130000}%
\pgfsetlinewidth{0.000000pt}%
\definecolor{currentstroke}{rgb}{0.000000,0.000000,0.000000}%
\pgfsetstrokecolor{currentstroke}%
\pgfsetdash{}{0pt}%
\pgfpathmoveto{\pgfqpoint{0.698817in}{1.000853in}}%
\pgfpathlineto{\pgfqpoint{0.714965in}{0.982348in}}%
\pgfpathlineto{\pgfqpoint{0.731121in}{0.964059in}}%
\pgfpathlineto{\pgfqpoint{0.739873in}{0.964218in}}%
\pgfpathlineto{\pgfqpoint{0.748584in}{0.964799in}}%
\pgfpathlineto{\pgfqpoint{0.732470in}{0.983400in}}%
\pgfpathlineto{\pgfqpoint{0.716365in}{1.002191in}}%
\pgfpathlineto{\pgfqpoint{0.707612in}{1.001324in}}%
\pgfpathlineto{\pgfqpoint{0.698817in}{1.000853in}}%
\pgfpathclose%
\pgfusepath{fill}%
\end{pgfscope}%
\begin{pgfscope}%
\pgfpathrectangle{\pgfqpoint{0.407592in}{0.222904in}}{\pgfqpoint{1.411163in}{1.411163in}}%
\pgfusepath{clip}%
\pgfsetbuttcap%
\pgfsetroundjoin%
\definecolor{currentfill}{rgb}{0.556863,0.654902,0.847059}%
\pgfsetfillcolor{currentfill}%
\pgfsetfillopacity{0.130000}%
\pgfsetlinewidth{0.000000pt}%
\definecolor{currentstroke}{rgb}{0.000000,0.000000,0.000000}%
\pgfsetstrokecolor{currentstroke}%
\pgfsetdash{}{0pt}%
\pgfpathmoveto{\pgfqpoint{0.845445in}{0.859161in}}%
\pgfpathlineto{\pgfqpoint{0.861619in}{0.843019in}}%
\pgfpathlineto{\pgfqpoint{0.877802in}{0.827486in}}%
\pgfpathlineto{\pgfqpoint{0.886304in}{0.826904in}}%
\pgfpathlineto{\pgfqpoint{0.894767in}{0.826919in}}%
\pgfpathlineto{\pgfqpoint{0.878626in}{0.842941in}}%
\pgfpathlineto{\pgfqpoint{0.862494in}{0.859543in}}%
\pgfpathlineto{\pgfqpoint{0.853990in}{0.859072in}}%
\pgfpathlineto{\pgfqpoint{0.845445in}{0.859161in}}%
\pgfpathclose%
\pgfusepath{fill}%
\end{pgfscope}%
\begin{pgfscope}%
\pgfpathrectangle{\pgfqpoint{0.407592in}{0.222904in}}{\pgfqpoint{1.411163in}{1.411163in}}%
\pgfusepath{clip}%
\pgfsetbuttcap%
\pgfsetroundjoin%
\definecolor{currentfill}{rgb}{0.556863,0.654902,0.847059}%
\pgfsetfillcolor{currentfill}%
\pgfsetfillopacity{0.130000}%
\pgfsetlinewidth{0.000000pt}%
\definecolor{currentstroke}{rgb}{0.000000,0.000000,0.000000}%
\pgfsetstrokecolor{currentstroke}%
\pgfsetdash{}{0pt}%
\pgfpathmoveto{\pgfqpoint{1.478058in}{1.002116in}}%
\pgfpathlineto{\pgfqpoint{1.495044in}{1.010418in}}%
\pgfpathlineto{\pgfqpoint{1.512102in}{1.019116in}}%
\pgfpathlineto{\pgfqpoint{1.519709in}{1.031542in}}%
\pgfpathlineto{\pgfqpoint{1.527304in}{1.044208in}}%
\pgfpathlineto{\pgfqpoint{1.510265in}{1.035964in}}%
\pgfpathlineto{\pgfqpoint{1.493298in}{1.028128in}}%
\pgfpathlineto{\pgfqpoint{1.485684in}{1.015004in}}%
\pgfpathlineto{\pgfqpoint{1.478058in}{1.002116in}}%
\pgfpathclose%
\pgfusepath{fill}%
\end{pgfscope}%
\begin{pgfscope}%
\pgfpathrectangle{\pgfqpoint{0.407592in}{0.222904in}}{\pgfqpoint{1.411163in}{1.411163in}}%
\pgfusepath{clip}%
\pgfsetbuttcap%
\pgfsetroundjoin%
\definecolor{currentfill}{rgb}{0.556863,0.654902,0.847059}%
\pgfsetfillcolor{currentfill}%
\pgfsetfillopacity{0.130000}%
\pgfsetlinewidth{0.000000pt}%
\definecolor{currentstroke}{rgb}{0.000000,0.000000,0.000000}%
\pgfsetstrokecolor{currentstroke}%
\pgfsetdash{}{0pt}%
\pgfpathmoveto{\pgfqpoint{0.551894in}{1.152493in}}%
\pgfpathlineto{\pgfqpoint{0.568015in}{1.133129in}}%
\pgfpathlineto{\pgfqpoint{0.584145in}{1.113863in}}%
\pgfpathlineto{\pgfqpoint{0.593149in}{1.114369in}}%
\pgfpathlineto{\pgfqpoint{0.602112in}{1.115194in}}%
\pgfpathlineto{\pgfqpoint{0.586025in}{1.134653in}}%
\pgfpathlineto{\pgfqpoint{0.569946in}{1.154195in}}%
\pgfpathlineto{\pgfqpoint{0.560940in}{1.153193in}}%
\pgfpathlineto{\pgfqpoint{0.551894in}{1.152493in}}%
\pgfpathclose%
\pgfusepath{fill}%
\end{pgfscope}%
\begin{pgfscope}%
\pgfpathrectangle{\pgfqpoint{0.407592in}{0.222904in}}{\pgfqpoint{1.411163in}{1.411163in}}%
\pgfusepath{clip}%
\pgfsetbuttcap%
\pgfsetroundjoin%
\definecolor{currentfill}{rgb}{0.556863,0.654902,0.847059}%
\pgfsetfillcolor{currentfill}%
\pgfsetfillopacity{0.130000}%
\pgfsetlinewidth{0.000000pt}%
\definecolor{currentstroke}{rgb}{0.000000,0.000000,0.000000}%
\pgfsetstrokecolor{currentstroke}%
\pgfsetdash{}{0pt}%
\pgfpathmoveto{\pgfqpoint{1.138235in}{0.737000in}}%
\pgfpathlineto{\pgfqpoint{1.154578in}{0.739893in}}%
\pgfpathlineto{\pgfqpoint{1.170978in}{0.744261in}}%
\pgfpathlineto{\pgfqpoint{1.178955in}{0.749601in}}%
\pgfpathlineto{\pgfqpoint{1.186909in}{0.755929in}}%
\pgfpathlineto{\pgfqpoint{1.170532in}{0.751611in}}%
\pgfpathlineto{\pgfqpoint{1.154213in}{0.748759in}}%
\pgfpathlineto{\pgfqpoint{1.146236in}{0.742334in}}%
\pgfpathlineto{\pgfqpoint{1.138235in}{0.737000in}}%
\pgfpathclose%
\pgfusepath{fill}%
\end{pgfscope}%
\begin{pgfscope}%
\pgfpathrectangle{\pgfqpoint{0.407592in}{0.222904in}}{\pgfqpoint{1.411163in}{1.411163in}}%
\pgfusepath{clip}%
\pgfsetbuttcap%
\pgfsetroundjoin%
\definecolor{currentfill}{rgb}{0.556863,0.654902,0.847059}%
\pgfsetfillcolor{currentfill}%
\pgfsetfillopacity{0.130000}%
\pgfsetlinewidth{0.000000pt}%
\definecolor{currentstroke}{rgb}{0.000000,0.000000,0.000000}%
\pgfsetstrokecolor{currentstroke}%
\pgfsetdash{}{0pt}%
\pgfpathmoveto{\pgfqpoint{1.089563in}{0.729307in}}%
\pgfpathlineto{\pgfqpoint{1.105838in}{0.728526in}}%
\pgfpathlineto{\pgfqpoint{1.122158in}{0.729595in}}%
\pgfpathlineto{\pgfqpoint{1.130210in}{0.732760in}}%
\pgfpathlineto{\pgfqpoint{1.138235in}{0.737000in}}%
\pgfpathlineto{\pgfqpoint{1.121944in}{0.735796in}}%
\pgfpathlineto{\pgfqpoint{1.105701in}{0.736477in}}%
\pgfpathlineto{\pgfqpoint{1.097646in}{0.732322in}}%
\pgfpathlineto{\pgfqpoint{1.089563in}{0.729307in}}%
\pgfpathclose%
\pgfusepath{fill}%
\end{pgfscope}%
\begin{pgfscope}%
\pgfpathrectangle{\pgfqpoint{0.407592in}{0.222904in}}{\pgfqpoint{1.411163in}{1.411163in}}%
\pgfusepath{clip}%
\pgfsetbuttcap%
\pgfsetroundjoin%
\definecolor{currentfill}{rgb}{0.556863,0.654902,0.847059}%
\pgfsetfillcolor{currentfill}%
\pgfsetfillopacity{0.130000}%
\pgfsetlinewidth{0.000000pt}%
\definecolor{currentstroke}{rgb}{0.000000,0.000000,0.000000}%
\pgfsetstrokecolor{currentstroke}%
\pgfsetdash{}{0pt}%
\pgfpathmoveto{\pgfqpoint{1.300181in}{0.835734in}}%
\pgfpathlineto{\pgfqpoint{1.316831in}{0.843577in}}%
\pgfpathlineto{\pgfqpoint{1.333551in}{0.852024in}}%
\pgfpathlineto{\pgfqpoint{1.341335in}{0.862336in}}%
\pgfpathlineto{\pgfqpoint{1.349107in}{0.873144in}}%
\pgfpathlineto{\pgfqpoint{1.332405in}{0.865216in}}%
\pgfpathlineto{\pgfqpoint{1.315772in}{0.857916in}}%
\pgfpathlineto{\pgfqpoint{1.307983in}{0.846565in}}%
\pgfpathlineto{\pgfqpoint{1.300181in}{0.835734in}}%
\pgfpathclose%
\pgfusepath{fill}%
\end{pgfscope}%
\begin{pgfscope}%
\pgfpathrectangle{\pgfqpoint{0.407592in}{0.222904in}}{\pgfqpoint{1.411163in}{1.411163in}}%
\pgfusepath{clip}%
\pgfsetbuttcap%
\pgfsetroundjoin%
\definecolor{currentfill}{rgb}{0.556863,0.654902,0.847059}%
\pgfsetfillcolor{currentfill}%
\pgfsetfillopacity{0.130000}%
\pgfsetlinewidth{0.000000pt}%
\definecolor{currentstroke}{rgb}{0.000000,0.000000,0.000000}%
\pgfsetstrokecolor{currentstroke}%
\pgfsetdash{}{0pt}%
\pgfpathmoveto{\pgfqpoint{0.959431in}{0.770949in}}%
\pgfpathlineto{\pgfqpoint{0.975627in}{0.759706in}}%
\pgfpathlineto{\pgfqpoint{0.991839in}{0.749886in}}%
\pgfpathlineto{\pgfqpoint{1.000136in}{0.749202in}}%
\pgfpathlineto{\pgfqpoint{1.008395in}{0.749341in}}%
\pgfpathlineto{\pgfqpoint{0.992225in}{0.759518in}}%
\pgfpathlineto{\pgfqpoint{0.976071in}{0.771175in}}%
\pgfpathlineto{\pgfqpoint{0.967770in}{0.770670in}}%
\pgfpathlineto{\pgfqpoint{0.959431in}{0.770949in}}%
\pgfpathclose%
\pgfusepath{fill}%
\end{pgfscope}%
\begin{pgfscope}%
\pgfpathrectangle{\pgfqpoint{0.407592in}{0.222904in}}{\pgfqpoint{1.411163in}{1.411163in}}%
\pgfusepath{clip}%
\pgfsetbuttcap%
\pgfsetroundjoin%
\definecolor{currentfill}{rgb}{0.556863,0.654902,0.847059}%
\pgfsetfillcolor{currentfill}%
\pgfsetfillopacity{0.130000}%
\pgfsetlinewidth{0.000000pt}%
\definecolor{currentstroke}{rgb}{0.000000,0.000000,0.000000}%
\pgfsetstrokecolor{currentstroke}%
\pgfsetdash{}{0pt}%
\pgfpathmoveto{\pgfqpoint{0.731121in}{0.964059in}}%
\pgfpathlineto{\pgfqpoint{0.747285in}{0.946016in}}%
\pgfpathlineto{\pgfqpoint{0.763458in}{0.928253in}}%
\pgfpathlineto{\pgfqpoint{0.772168in}{0.928046in}}%
\pgfpathlineto{\pgfqpoint{0.780837in}{0.928288in}}%
\pgfpathlineto{\pgfqpoint{0.764706in}{0.946418in}}%
\pgfpathlineto{\pgfqpoint{0.748584in}{0.964799in}}%
\pgfpathlineto{\pgfqpoint{0.739873in}{0.964218in}}%
\pgfpathlineto{\pgfqpoint{0.731121in}{0.964059in}}%
\pgfpathclose%
\pgfusepath{fill}%
\end{pgfscope}%
\begin{pgfscope}%
\pgfpathrectangle{\pgfqpoint{0.407592in}{0.222904in}}{\pgfqpoint{1.411163in}{1.411163in}}%
\pgfusepath{clip}%
\pgfsetbuttcap%
\pgfsetroundjoin%
\definecolor{currentfill}{rgb}{0.556863,0.654902,0.847059}%
\pgfsetfillcolor{currentfill}%
\pgfsetfillopacity{0.130000}%
\pgfsetlinewidth{0.000000pt}%
\definecolor{currentstroke}{rgb}{0.000000,0.000000,0.000000}%
\pgfsetstrokecolor{currentstroke}%
\pgfsetdash{}{0pt}%
\pgfpathmoveto{\pgfqpoint{1.186909in}{0.755929in}}%
\pgfpathlineto{\pgfqpoint{1.203348in}{0.761502in}}%
\pgfpathlineto{\pgfqpoint{1.219850in}{0.768138in}}%
\pgfpathlineto{\pgfqpoint{1.227764in}{0.775216in}}%
\pgfpathlineto{\pgfqpoint{1.235661in}{0.783126in}}%
\pgfpathlineto{\pgfqpoint{1.219178in}{0.776751in}}%
\pgfpathlineto{\pgfqpoint{1.202760in}{0.771439in}}%
\pgfpathlineto{\pgfqpoint{1.194844in}{0.763221in}}%
\pgfpathlineto{\pgfqpoint{1.186909in}{0.755929in}}%
\pgfpathclose%
\pgfusepath{fill}%
\end{pgfscope}%
\begin{pgfscope}%
\pgfpathrectangle{\pgfqpoint{0.407592in}{0.222904in}}{\pgfqpoint{1.411163in}{1.411163in}}%
\pgfusepath{clip}%
\pgfsetbuttcap%
\pgfsetroundjoin%
\definecolor{currentfill}{rgb}{0.556863,0.654902,0.847059}%
\pgfsetfillcolor{currentfill}%
\pgfsetfillopacity{0.130000}%
\pgfsetlinewidth{0.000000pt}%
\definecolor{currentstroke}{rgb}{0.000000,0.000000,0.000000}%
\pgfsetstrokecolor{currentstroke}%
\pgfsetdash{}{0pt}%
\pgfpathmoveto{\pgfqpoint{1.413609in}{0.936014in}}%
\pgfpathlineto{\pgfqpoint{1.430489in}{0.944433in}}%
\pgfpathlineto{\pgfqpoint{1.447440in}{0.953287in}}%
\pgfpathlineto{\pgfqpoint{1.455112in}{0.965047in}}%
\pgfpathlineto{\pgfqpoint{1.462772in}{0.977118in}}%
\pgfpathlineto{\pgfqpoint{1.445840in}{0.968756in}}%
\pgfpathlineto{\pgfqpoint{1.428979in}{0.960848in}}%
\pgfpathlineto{\pgfqpoint{1.421300in}{0.948275in}}%
\pgfpathlineto{\pgfqpoint{1.413609in}{0.936014in}}%
\pgfpathclose%
\pgfusepath{fill}%
\end{pgfscope}%
\begin{pgfscope}%
\pgfpathrectangle{\pgfqpoint{0.407592in}{0.222904in}}{\pgfqpoint{1.411163in}{1.411163in}}%
\pgfusepath{clip}%
\pgfsetbuttcap%
\pgfsetroundjoin%
\definecolor{currentfill}{rgb}{0.556863,0.654902,0.847059}%
\pgfsetfillcolor{currentfill}%
\pgfsetfillopacity{0.130000}%
\pgfsetlinewidth{0.000000pt}%
\definecolor{currentstroke}{rgb}{0.000000,0.000000,0.000000}%
\pgfsetstrokecolor{currentstroke}%
\pgfsetdash{}{0pt}%
\pgfpathmoveto{\pgfqpoint{1.040795in}{0.734121in}}%
\pgfpathlineto{\pgfqpoint{1.057034in}{0.729326in}}%
\pgfpathlineto{\pgfqpoint{1.073305in}{0.726467in}}%
\pgfpathlineto{\pgfqpoint{1.081450in}{0.727379in}}%
\pgfpathlineto{\pgfqpoint{1.089563in}{0.729307in}}%
\pgfpathlineto{\pgfqpoint{1.073329in}{0.732072in}}%
\pgfpathlineto{\pgfqpoint{1.057128in}{0.736877in}}%
\pgfpathlineto{\pgfqpoint{1.048978in}{0.734990in}}%
\pgfpathlineto{\pgfqpoint{1.040795in}{0.734121in}}%
\pgfpathclose%
\pgfusepath{fill}%
\end{pgfscope}%
\begin{pgfscope}%
\pgfpathrectangle{\pgfqpoint{0.407592in}{0.222904in}}{\pgfqpoint{1.411163in}{1.411163in}}%
\pgfusepath{clip}%
\pgfsetbuttcap%
\pgfsetroundjoin%
\definecolor{currentfill}{rgb}{0.556863,0.654902,0.847059}%
\pgfsetfillcolor{currentfill}%
\pgfsetfillopacity{0.130000}%
\pgfsetlinewidth{0.000000pt}%
\definecolor{currentstroke}{rgb}{0.000000,0.000000,0.000000}%
\pgfsetstrokecolor{currentstroke}%
\pgfsetdash{}{0pt}%
\pgfpathmoveto{\pgfqpoint{0.584145in}{1.113863in}}%
\pgfpathlineto{\pgfqpoint{0.600284in}{1.094706in}}%
\pgfpathlineto{\pgfqpoint{0.616431in}{1.075670in}}%
\pgfpathlineto{\pgfqpoint{0.625393in}{1.075951in}}%
\pgfpathlineto{\pgfqpoint{0.634313in}{1.076567in}}%
\pgfpathlineto{\pgfqpoint{0.618208in}{1.095828in}}%
\pgfpathlineto{\pgfqpoint{0.602112in}{1.115194in}}%
\pgfpathlineto{\pgfqpoint{0.593149in}{1.114369in}}%
\pgfpathlineto{\pgfqpoint{0.584145in}{1.113863in}}%
\pgfpathclose%
\pgfusepath{fill}%
\end{pgfscope}%
\begin{pgfscope}%
\pgfpathrectangle{\pgfqpoint{0.407592in}{0.222904in}}{\pgfqpoint{1.411163in}{1.411163in}}%
\pgfusepath{clip}%
\pgfsetbuttcap%
\pgfsetroundjoin%
\definecolor{currentfill}{rgb}{0.556863,0.654902,0.847059}%
\pgfsetfillcolor{currentfill}%
\pgfsetfillopacity{0.130000}%
\pgfsetlinewidth{0.000000pt}%
\definecolor{currentstroke}{rgb}{0.000000,0.000000,0.000000}%
\pgfsetstrokecolor{currentstroke}%
\pgfsetdash{}{0pt}%
\pgfpathmoveto{\pgfqpoint{0.877802in}{0.827486in}}%
\pgfpathlineto{\pgfqpoint{0.893995in}{0.812659in}}%
\pgfpathlineto{\pgfqpoint{0.910197in}{0.798651in}}%
\pgfpathlineto{\pgfqpoint{0.918657in}{0.797533in}}%
\pgfpathlineto{\pgfqpoint{0.927077in}{0.797044in}}%
\pgfpathlineto{\pgfqpoint{0.910917in}{0.811580in}}%
\pgfpathlineto{\pgfqpoint{0.894767in}{0.826919in}}%
\pgfpathlineto{\pgfqpoint{0.886304in}{0.826904in}}%
\pgfpathlineto{\pgfqpoint{0.877802in}{0.827486in}}%
\pgfpathclose%
\pgfusepath{fill}%
\end{pgfscope}%
\begin{pgfscope}%
\pgfpathrectangle{\pgfqpoint{0.407592in}{0.222904in}}{\pgfqpoint{1.411163in}{1.411163in}}%
\pgfusepath{clip}%
\pgfsetbuttcap%
\pgfsetroundjoin%
\definecolor{currentfill}{rgb}{0.556863,0.654902,0.847059}%
\pgfsetfillcolor{currentfill}%
\pgfsetfillopacity{0.130000}%
\pgfsetlinewidth{0.000000pt}%
\definecolor{currentstroke}{rgb}{0.000000,0.000000,0.000000}%
\pgfsetstrokecolor{currentstroke}%
\pgfsetdash{}{0pt}%
\pgfpathmoveto{\pgfqpoint{1.591791in}{1.112895in}}%
\pgfpathlineto{\pgfqpoint{1.609010in}{1.121387in}}%
\pgfpathlineto{\pgfqpoint{1.626303in}{1.130208in}}%
\pgfpathlineto{\pgfqpoint{1.633802in}{1.143067in}}%
\pgfpathlineto{\pgfqpoint{1.641291in}{1.156100in}}%
\pgfpathlineto{\pgfqpoint{1.624018in}{1.147661in}}%
\pgfpathlineto{\pgfqpoint{1.606818in}{1.139560in}}%
\pgfpathlineto{\pgfqpoint{1.599310in}{1.126143in}}%
\pgfpathlineto{\pgfqpoint{1.591791in}{1.112895in}}%
\pgfpathclose%
\pgfusepath{fill}%
\end{pgfscope}%
\begin{pgfscope}%
\pgfpathrectangle{\pgfqpoint{0.407592in}{0.222904in}}{\pgfqpoint{1.411163in}{1.411163in}}%
\pgfusepath{clip}%
\pgfsetbuttcap%
\pgfsetroundjoin%
\definecolor{currentfill}{rgb}{0.556863,0.654902,0.847059}%
\pgfsetfillcolor{currentfill}%
\pgfsetfillopacity{0.130000}%
\pgfsetlinewidth{0.000000pt}%
\definecolor{currentstroke}{rgb}{0.000000,0.000000,0.000000}%
\pgfsetstrokecolor{currentstroke}%
\pgfsetdash{}{0pt}%
\pgfpathmoveto{\pgfqpoint{0.763458in}{0.928253in}}%
\pgfpathlineto{\pgfqpoint{0.779639in}{0.910810in}}%
\pgfpathlineto{\pgfqpoint{0.795829in}{0.893733in}}%
\pgfpathlineto{\pgfqpoint{0.804497in}{0.893101in}}%
\pgfpathlineto{\pgfqpoint{0.813124in}{0.892943in}}%
\pgfpathlineto{\pgfqpoint{0.796976in}{0.910448in}}%
\pgfpathlineto{\pgfqpoint{0.780837in}{0.928288in}}%
\pgfpathlineto{\pgfqpoint{0.772168in}{0.928046in}}%
\pgfpathlineto{\pgfqpoint{0.763458in}{0.928253in}}%
\pgfpathclose%
\pgfusepath{fill}%
\end{pgfscope}%
\begin{pgfscope}%
\pgfpathrectangle{\pgfqpoint{0.407592in}{0.222904in}}{\pgfqpoint{1.411163in}{1.411163in}}%
\pgfusepath{clip}%
\pgfsetbuttcap%
\pgfsetroundjoin%
\definecolor{currentfill}{rgb}{0.556863,0.654902,0.847059}%
\pgfsetfillcolor{currentfill}%
\pgfsetfillopacity{0.130000}%
\pgfsetlinewidth{0.000000pt}%
\definecolor{currentstroke}{rgb}{0.000000,0.000000,0.000000}%
\pgfsetstrokecolor{currentstroke}%
\pgfsetdash{}{0pt}%
\pgfpathmoveto{\pgfqpoint{1.235661in}{0.783126in}}%
\pgfpathlineto{\pgfqpoint{1.252210in}{0.790404in}}%
\pgfpathlineto{\pgfqpoint{1.268827in}{0.798448in}}%
\pgfpathlineto{\pgfqpoint{1.276688in}{0.806783in}}%
\pgfpathlineto{\pgfqpoint{1.284534in}{0.815801in}}%
\pgfpathlineto{\pgfqpoint{1.267935in}{0.808151in}}%
\pgfpathlineto{\pgfqpoint{1.251404in}{0.801285in}}%
\pgfpathlineto{\pgfqpoint{1.243541in}{0.791830in}}%
\pgfpathlineto{\pgfqpoint{1.235661in}{0.783126in}}%
\pgfpathclose%
\pgfusepath{fill}%
\end{pgfscope}%
\begin{pgfscope}%
\pgfpathrectangle{\pgfqpoint{0.407592in}{0.222904in}}{\pgfqpoint{1.411163in}{1.411163in}}%
\pgfusepath{clip}%
\pgfsetbuttcap%
\pgfsetroundjoin%
\definecolor{currentfill}{rgb}{0.556863,0.654902,0.847059}%
\pgfsetfillcolor{currentfill}%
\pgfsetfillopacity{0.130000}%
\pgfsetlinewidth{0.000000pt}%
\definecolor{currentstroke}{rgb}{0.000000,0.000000,0.000000}%
\pgfsetstrokecolor{currentstroke}%
\pgfsetdash{}{0pt}%
\pgfpathmoveto{\pgfqpoint{1.349107in}{0.873144in}}%
\pgfpathlineto{\pgfqpoint{1.365879in}{0.881626in}}%
\pgfpathlineto{\pgfqpoint{1.382724in}{0.890597in}}%
\pgfpathlineto{\pgfqpoint{1.390464in}{0.901355in}}%
\pgfpathlineto{\pgfqpoint{1.398191in}{0.912528in}}%
\pgfpathlineto{\pgfqpoint{1.381366in}{0.904064in}}%
\pgfpathlineto{\pgfqpoint{1.364611in}{0.896115in}}%
\pgfpathlineto{\pgfqpoint{1.356865in}{0.884415in}}%
\pgfpathlineto{\pgfqpoint{1.349107in}{0.873144in}}%
\pgfpathclose%
\pgfusepath{fill}%
\end{pgfscope}%
\begin{pgfscope}%
\pgfpathrectangle{\pgfqpoint{0.407592in}{0.222904in}}{\pgfqpoint{1.411163in}{1.411163in}}%
\pgfusepath{clip}%
\pgfsetbuttcap%
\pgfsetroundjoin%
\definecolor{currentfill}{rgb}{0.556863,0.654902,0.847059}%
\pgfsetfillcolor{currentfill}%
\pgfsetfillopacity{0.130000}%
\pgfsetlinewidth{0.000000pt}%
\definecolor{currentstroke}{rgb}{0.000000,0.000000,0.000000}%
\pgfsetstrokecolor{currentstroke}%
\pgfsetdash{}{0pt}%
\pgfpathmoveto{\pgfqpoint{0.616431in}{1.075670in}}%
\pgfpathlineto{\pgfqpoint{0.632587in}{1.056770in}}%
\pgfpathlineto{\pgfqpoint{0.648751in}{1.038021in}}%
\pgfpathlineto{\pgfqpoint{0.657670in}{1.038039in}}%
\pgfpathlineto{\pgfqpoint{0.666548in}{1.038411in}}%
\pgfpathlineto{\pgfqpoint{0.650426in}{1.057423in}}%
\pgfpathlineto{\pgfqpoint{0.634313in}{1.076567in}}%
\pgfpathlineto{\pgfqpoint{0.625393in}{1.075951in}}%
\pgfpathlineto{\pgfqpoint{0.616431in}{1.075670in}}%
\pgfpathclose%
\pgfusepath{fill}%
\end{pgfscope}%
\begin{pgfscope}%
\pgfpathrectangle{\pgfqpoint{0.407592in}{0.222904in}}{\pgfqpoint{1.411163in}{1.411163in}}%
\pgfusepath{clip}%
\pgfsetbuttcap%
\pgfsetroundjoin%
\definecolor{currentfill}{rgb}{0.556863,0.654902,0.847059}%
\pgfsetfillcolor{currentfill}%
\pgfsetfillopacity{0.130000}%
\pgfsetlinewidth{0.000000pt}%
\definecolor{currentstroke}{rgb}{0.000000,0.000000,0.000000}%
\pgfsetstrokecolor{currentstroke}%
\pgfsetdash{}{0pt}%
\pgfpathmoveto{\pgfqpoint{1.527304in}{1.044208in}}%
\pgfpathlineto{\pgfqpoint{1.544416in}{1.052830in}}%
\pgfpathlineto{\pgfqpoint{1.561602in}{1.061805in}}%
\pgfpathlineto{\pgfqpoint{1.569166in}{1.074267in}}%
\pgfpathlineto{\pgfqpoint{1.576718in}{1.086945in}}%
\pgfpathlineto{\pgfqpoint{1.559553in}{1.078391in}}%
\pgfpathlineto{\pgfqpoint{1.542460in}{1.070200in}}%
\pgfpathlineto{\pgfqpoint{1.534888in}{1.057099in}}%
\pgfpathlineto{\pgfqpoint{1.527304in}{1.044208in}}%
\pgfpathclose%
\pgfusepath{fill}%
\end{pgfscope}%
\begin{pgfscope}%
\pgfpathrectangle{\pgfqpoint{0.407592in}{0.222904in}}{\pgfqpoint{1.411163in}{1.411163in}}%
\pgfusepath{clip}%
\pgfsetbuttcap%
\pgfsetroundjoin%
\definecolor{currentfill}{rgb}{0.556863,0.654902,0.847059}%
\pgfsetfillcolor{currentfill}%
\pgfsetfillopacity{0.130000}%
\pgfsetlinewidth{0.000000pt}%
\definecolor{currentstroke}{rgb}{0.000000,0.000000,0.000000}%
\pgfsetstrokecolor{currentstroke}%
\pgfsetdash{}{0pt}%
\pgfpathmoveto{\pgfqpoint{0.991839in}{0.749886in}}%
\pgfpathlineto{\pgfqpoint{1.008070in}{0.741640in}}%
\pgfpathlineto{\pgfqpoint{1.024323in}{0.735092in}}%
\pgfpathlineto{\pgfqpoint{1.032577in}{0.734186in}}%
\pgfpathlineto{\pgfqpoint{1.040795in}{0.734121in}}%
\pgfpathlineto{\pgfqpoint{1.024584in}{0.740825in}}%
\pgfpathlineto{\pgfqpoint{1.008395in}{0.749341in}}%
\pgfpathlineto{\pgfqpoint{1.000136in}{0.749202in}}%
\pgfpathlineto{\pgfqpoint{0.991839in}{0.749886in}}%
\pgfpathclose%
\pgfusepath{fill}%
\end{pgfscope}%
\begin{pgfscope}%
\pgfpathrectangle{\pgfqpoint{0.407592in}{0.222904in}}{\pgfqpoint{1.411163in}{1.411163in}}%
\pgfusepath{clip}%
\pgfsetbuttcap%
\pgfsetroundjoin%
\definecolor{currentfill}{rgb}{0.556863,0.654902,0.847059}%
\pgfsetfillcolor{currentfill}%
\pgfsetfillopacity{0.130000}%
\pgfsetlinewidth{0.000000pt}%
\definecolor{currentstroke}{rgb}{0.000000,0.000000,0.000000}%
\pgfsetstrokecolor{currentstroke}%
\pgfsetdash{}{0pt}%
\pgfpathmoveto{\pgfqpoint{1.462772in}{0.977118in}}%
\pgfpathlineto{\pgfqpoint{1.479777in}{0.985894in}}%
\pgfpathlineto{\pgfqpoint{1.496855in}{0.995049in}}%
\pgfpathlineto{\pgfqpoint{1.504485in}{1.006946in}}%
\pgfpathlineto{\pgfqpoint{1.512102in}{1.019116in}}%
\pgfpathlineto{\pgfqpoint{1.495044in}{1.010418in}}%
\pgfpathlineto{\pgfqpoint{1.478058in}{1.002116in}}%
\pgfpathlineto{\pgfqpoint{1.470421in}{0.989481in}}%
\pgfpathlineto{\pgfqpoint{1.462772in}{0.977118in}}%
\pgfpathclose%
\pgfusepath{fill}%
\end{pgfscope}%
\begin{pgfscope}%
\pgfpathrectangle{\pgfqpoint{0.407592in}{0.222904in}}{\pgfqpoint{1.411163in}{1.411163in}}%
\pgfusepath{clip}%
\pgfsetbuttcap%
\pgfsetroundjoin%
\definecolor{currentfill}{rgb}{0.556863,0.654902,0.847059}%
\pgfsetfillcolor{currentfill}%
\pgfsetfillopacity{0.130000}%
\pgfsetlinewidth{0.000000pt}%
\definecolor{currentstroke}{rgb}{0.000000,0.000000,0.000000}%
\pgfsetstrokecolor{currentstroke}%
\pgfsetdash{}{0pt}%
\pgfpathmoveto{\pgfqpoint{0.910197in}{0.798651in}}%
\pgfpathlineto{\pgfqpoint{0.926411in}{0.785584in}}%
\pgfpathlineto{\pgfqpoint{0.942637in}{0.773594in}}%
\pgfpathlineto{\pgfqpoint{0.951053in}{0.771946in}}%
\pgfpathlineto{\pgfqpoint{0.959431in}{0.770949in}}%
\pgfpathlineto{\pgfqpoint{0.943248in}{0.783450in}}%
\pgfpathlineto{\pgfqpoint{0.927077in}{0.797044in}}%
\pgfpathlineto{\pgfqpoint{0.918657in}{0.797533in}}%
\pgfpathlineto{\pgfqpoint{0.910197in}{0.798651in}}%
\pgfpathclose%
\pgfusepath{fill}%
\end{pgfscope}%
\begin{pgfscope}%
\pgfpathrectangle{\pgfqpoint{0.407592in}{0.222904in}}{\pgfqpoint{1.411163in}{1.411163in}}%
\pgfusepath{clip}%
\pgfsetbuttcap%
\pgfsetroundjoin%
\definecolor{currentfill}{rgb}{0.556863,0.654902,0.847059}%
\pgfsetfillcolor{currentfill}%
\pgfsetfillopacity{0.130000}%
\pgfsetlinewidth{0.000000pt}%
\definecolor{currentstroke}{rgb}{0.000000,0.000000,0.000000}%
\pgfsetstrokecolor{currentstroke}%
\pgfsetdash{}{0pt}%
\pgfpathmoveto{\pgfqpoint{0.648751in}{1.038021in}}%
\pgfpathlineto{\pgfqpoint{0.664924in}{1.019442in}}%
\pgfpathlineto{\pgfqpoint{0.681105in}{1.001054in}}%
\pgfpathlineto{\pgfqpoint{0.689982in}{1.000767in}}%
\pgfpathlineto{\pgfqpoint{0.698817in}{1.000853in}}%
\pgfpathlineto{\pgfqpoint{0.682678in}{1.019548in}}%
\pgfpathlineto{\pgfqpoint{0.666548in}{1.038411in}}%
\pgfpathlineto{\pgfqpoint{0.657670in}{1.038039in}}%
\pgfpathlineto{\pgfqpoint{0.648751in}{1.038021in}}%
\pgfpathclose%
\pgfusepath{fill}%
\end{pgfscope}%
\begin{pgfscope}%
\pgfpathrectangle{\pgfqpoint{0.407592in}{0.222904in}}{\pgfqpoint{1.411163in}{1.411163in}}%
\pgfusepath{clip}%
\pgfsetbuttcap%
\pgfsetroundjoin%
\definecolor{currentfill}{rgb}{0.556863,0.654902,0.847059}%
\pgfsetfillcolor{currentfill}%
\pgfsetfillopacity{0.130000}%
\pgfsetlinewidth{0.000000pt}%
\definecolor{currentstroke}{rgb}{0.000000,0.000000,0.000000}%
\pgfsetstrokecolor{currentstroke}%
\pgfsetdash{}{0pt}%
\pgfpathmoveto{\pgfqpoint{0.795829in}{0.893733in}}%
\pgfpathlineto{\pgfqpoint{0.812028in}{0.877075in}}%
\pgfpathlineto{\pgfqpoint{0.828235in}{0.860898in}}%
\pgfpathlineto{\pgfqpoint{0.836860in}{0.859781in}}%
\pgfpathlineto{\pgfqpoint{0.845445in}{0.859161in}}%
\pgfpathlineto{\pgfqpoint{0.829280in}{0.875827in}}%
\pgfpathlineto{\pgfqpoint{0.813124in}{0.892943in}}%
\pgfpathlineto{\pgfqpoint{0.804497in}{0.893101in}}%
\pgfpathlineto{\pgfqpoint{0.795829in}{0.893733in}}%
\pgfpathclose%
\pgfusepath{fill}%
\end{pgfscope}%
\begin{pgfscope}%
\pgfpathrectangle{\pgfqpoint{0.407592in}{0.222904in}}{\pgfqpoint{1.411163in}{1.411163in}}%
\pgfusepath{clip}%
\pgfsetbuttcap%
\pgfsetroundjoin%
\definecolor{currentfill}{rgb}{0.556863,0.654902,0.847059}%
\pgfsetfillcolor{currentfill}%
\pgfsetfillopacity{0.130000}%
\pgfsetlinewidth{0.000000pt}%
\definecolor{currentstroke}{rgb}{0.000000,0.000000,0.000000}%
\pgfsetstrokecolor{currentstroke}%
\pgfsetdash{}{0pt}%
\pgfpathmoveto{\pgfqpoint{1.284534in}{0.815801in}}%
\pgfpathlineto{\pgfqpoint{1.301202in}{0.824123in}}%
\pgfpathlineto{\pgfqpoint{1.317940in}{0.833022in}}%
\pgfpathlineto{\pgfqpoint{1.325752in}{0.842242in}}%
\pgfpathlineto{\pgfqpoint{1.333551in}{0.852024in}}%
\pgfpathlineto{\pgfqpoint{1.316831in}{0.843577in}}%
\pgfpathlineto{\pgfqpoint{1.300181in}{0.835734in}}%
\pgfpathlineto{\pgfqpoint{1.292365in}{0.825465in}}%
\pgfpathlineto{\pgfqpoint{1.284534in}{0.815801in}}%
\pgfpathclose%
\pgfusepath{fill}%
\end{pgfscope}%
\begin{pgfscope}%
\pgfpathrectangle{\pgfqpoint{0.407592in}{0.222904in}}{\pgfqpoint{1.411163in}{1.411163in}}%
\pgfusepath{clip}%
\pgfsetbuttcap%
\pgfsetroundjoin%
\definecolor{currentfill}{rgb}{0.556863,0.654902,0.847059}%
\pgfsetfillcolor{currentfill}%
\pgfsetfillopacity{0.130000}%
\pgfsetlinewidth{0.000000pt}%
\definecolor{currentstroke}{rgb}{0.000000,0.000000,0.000000}%
\pgfsetstrokecolor{currentstroke}%
\pgfsetdash{}{0pt}%
\pgfpathmoveto{\pgfqpoint{1.398191in}{0.912528in}}%
\pgfpathlineto{\pgfqpoint{1.415089in}{0.921451in}}%
\pgfpathlineto{\pgfqpoint{1.432059in}{0.930787in}}%
\pgfpathlineto{\pgfqpoint{1.439755in}{0.941860in}}%
\pgfpathlineto{\pgfqpoint{1.447440in}{0.953287in}}%
\pgfpathlineto{\pgfqpoint{1.430489in}{0.944433in}}%
\pgfpathlineto{\pgfqpoint{1.413609in}{0.936014in}}%
\pgfpathlineto{\pgfqpoint{1.405906in}{0.924090in}}%
\pgfpathlineto{\pgfqpoint{1.398191in}{0.912528in}}%
\pgfpathclose%
\pgfusepath{fill}%
\end{pgfscope}%
\begin{pgfscope}%
\pgfpathrectangle{\pgfqpoint{0.407592in}{0.222904in}}{\pgfqpoint{1.411163in}{1.411163in}}%
\pgfusepath{clip}%
\pgfsetbuttcap%
\pgfsetroundjoin%
\definecolor{currentfill}{rgb}{0.556863,0.654902,0.847059}%
\pgfsetfillcolor{currentfill}%
\pgfsetfillopacity{0.130000}%
\pgfsetlinewidth{0.000000pt}%
\definecolor{currentstroke}{rgb}{0.000000,0.000000,0.000000}%
\pgfsetstrokecolor{currentstroke}%
\pgfsetdash{}{0pt}%
\pgfpathmoveto{\pgfqpoint{1.641291in}{1.156100in}}%
\pgfpathlineto{\pgfqpoint{1.658638in}{1.164855in}}%
\pgfpathlineto{\pgfqpoint{1.676059in}{1.173911in}}%
\pgfpathlineto{\pgfqpoint{1.683516in}{1.186745in}}%
\pgfpathlineto{\pgfqpoint{1.666105in}{1.177868in}}%
\pgfpathlineto{\pgfqpoint{1.648768in}{1.169296in}}%
\pgfusepath{fill}%
\end{pgfscope}%
\begin{pgfscope}%
\pgfpathrectangle{\pgfqpoint{0.407592in}{0.222904in}}{\pgfqpoint{1.411163in}{1.411163in}}%
\pgfusepath{clip}%
\pgfsetbuttcap%
\pgfsetroundjoin%
\definecolor{currentfill}{rgb}{0.556863,0.654902,0.847059}%
\pgfsetfillcolor{currentfill}%
\pgfsetfillopacity{0.130000}%
\pgfsetlinewidth{0.000000pt}%
\definecolor{currentstroke}{rgb}{0.000000,0.000000,0.000000}%
\pgfsetstrokecolor{currentstroke}%
\pgfsetdash{}{0pt}%
\pgfpathmoveto{\pgfqpoint{1.576718in}{1.086945in}}%
\pgfpathlineto{\pgfqpoint{1.593958in}{1.095838in}}%
\pgfpathlineto{\pgfqpoint{1.611272in}{1.105049in}}%
\pgfpathlineto{\pgfqpoint{1.618793in}{1.117531in}}%
\pgfpathlineto{\pgfqpoint{1.626303in}{1.130208in}}%
\pgfpathlineto{\pgfqpoint{1.609010in}{1.121387in}}%
\pgfpathlineto{\pgfqpoint{1.591791in}{1.112895in}}%
\pgfpathlineto{\pgfqpoint{1.584260in}{1.099825in}}%
\pgfpathlineto{\pgfqpoint{1.576718in}{1.086945in}}%
\pgfpathclose%
\pgfusepath{fill}%
\end{pgfscope}%
\begin{pgfscope}%
\pgfpathrectangle{\pgfqpoint{0.407592in}{0.222904in}}{\pgfqpoint{1.411163in}{1.411163in}}%
\pgfusepath{clip}%
\pgfsetbuttcap%
\pgfsetroundjoin%
\definecolor{currentfill}{rgb}{0.556863,0.654902,0.847059}%
\pgfsetfillcolor{currentfill}%
\pgfsetfillopacity{0.130000}%
\pgfsetlinewidth{0.000000pt}%
\definecolor{currentstroke}{rgb}{0.000000,0.000000,0.000000}%
\pgfsetstrokecolor{currentstroke}%
\pgfsetdash{}{0pt}%
\pgfpathmoveto{\pgfqpoint{1.122158in}{0.729595in}}%
\pgfpathlineto{\pgfqpoint{1.138528in}{0.732333in}}%
\pgfpathlineto{\pgfqpoint{1.154954in}{0.736535in}}%
\pgfpathlineto{\pgfqpoint{1.162978in}{0.739910in}}%
\pgfpathlineto{\pgfqpoint{1.170978in}{0.744261in}}%
\pgfpathlineto{\pgfqpoint{1.154578in}{0.739893in}}%
\pgfpathlineto{\pgfqpoint{1.138235in}{0.737000in}}%
\pgfpathlineto{\pgfqpoint{1.130210in}{0.732760in}}%
\pgfpathlineto{\pgfqpoint{1.122158in}{0.729595in}}%
\pgfpathclose%
\pgfusepath{fill}%
\end{pgfscope}%
\begin{pgfscope}%
\pgfpathrectangle{\pgfqpoint{0.407592in}{0.222904in}}{\pgfqpoint{1.411163in}{1.411163in}}%
\pgfusepath{clip}%
\pgfsetbuttcap%
\pgfsetroundjoin%
\definecolor{currentfill}{rgb}{0.556863,0.654902,0.847059}%
\pgfsetfillcolor{currentfill}%
\pgfsetfillopacity{0.130000}%
\pgfsetlinewidth{0.000000pt}%
\definecolor{currentstroke}{rgb}{0.000000,0.000000,0.000000}%
\pgfsetstrokecolor{currentstroke}%
\pgfsetdash{}{0pt}%
\pgfpathmoveto{\pgfqpoint{0.681105in}{1.001054in}}%
\pgfpathlineto{\pgfqpoint{0.697294in}{0.982880in}}%
\pgfpathlineto{\pgfqpoint{0.713492in}{0.964947in}}%
\pgfpathlineto{\pgfqpoint{0.722327in}{0.964308in}}%
\pgfpathlineto{\pgfqpoint{0.731121in}{0.964059in}}%
\pgfpathlineto{\pgfqpoint{0.714965in}{0.982348in}}%
\pgfpathlineto{\pgfqpoint{0.698817in}{1.000853in}}%
\pgfpathlineto{\pgfqpoint{0.689982in}{1.000767in}}%
\pgfpathlineto{\pgfqpoint{0.681105in}{1.001054in}}%
\pgfpathclose%
\pgfusepath{fill}%
\end{pgfscope}%
\begin{pgfscope}%
\pgfpathrectangle{\pgfqpoint{0.407592in}{0.222904in}}{\pgfqpoint{1.411163in}{1.411163in}}%
\pgfusepath{clip}%
\pgfsetbuttcap%
\pgfsetroundjoin%
\definecolor{currentfill}{rgb}{0.556863,0.654902,0.847059}%
\pgfsetfillcolor{currentfill}%
\pgfsetfillopacity{0.130000}%
\pgfsetlinewidth{0.000000pt}%
\definecolor{currentstroke}{rgb}{0.000000,0.000000,0.000000}%
\pgfsetstrokecolor{currentstroke}%
\pgfsetdash{}{0pt}%
\pgfpathmoveto{\pgfqpoint{1.073305in}{0.726467in}}%
\pgfpathlineto{\pgfqpoint{1.089614in}{0.725496in}}%
\pgfpathlineto{\pgfqpoint{1.105968in}{0.726299in}}%
\pgfpathlineto{\pgfqpoint{1.114077in}{0.727461in}}%
\pgfpathlineto{\pgfqpoint{1.122158in}{0.729595in}}%
\pgfpathlineto{\pgfqpoint{1.105838in}{0.728526in}}%
\pgfpathlineto{\pgfqpoint{1.089563in}{0.729307in}}%
\pgfpathlineto{\pgfqpoint{1.081450in}{0.727379in}}%
\pgfpathlineto{\pgfqpoint{1.073305in}{0.726467in}}%
\pgfpathclose%
\pgfusepath{fill}%
\end{pgfscope}%
\begin{pgfscope}%
\pgfpathrectangle{\pgfqpoint{0.407592in}{0.222904in}}{\pgfqpoint{1.411163in}{1.411163in}}%
\pgfusepath{clip}%
\pgfsetbuttcap%
\pgfsetroundjoin%
\definecolor{currentfill}{rgb}{0.556863,0.654902,0.847059}%
\pgfsetfillcolor{currentfill}%
\pgfsetfillopacity{0.130000}%
\pgfsetlinewidth{0.000000pt}%
\definecolor{currentstroke}{rgb}{0.000000,0.000000,0.000000}%
\pgfsetstrokecolor{currentstroke}%
\pgfsetdash{}{0pt}%
\pgfpathmoveto{\pgfqpoint{0.828235in}{0.860898in}}%
\pgfpathlineto{\pgfqpoint{0.844451in}{0.845273in}}%
\pgfpathlineto{\pgfqpoint{0.860677in}{0.830281in}}%
\pgfpathlineto{\pgfqpoint{0.869260in}{0.828625in}}%
\pgfpathlineto{\pgfqpoint{0.877802in}{0.827486in}}%
\pgfpathlineto{\pgfqpoint{0.861619in}{0.843019in}}%
\pgfpathlineto{\pgfqpoint{0.845445in}{0.859161in}}%
\pgfpathlineto{\pgfqpoint{0.836860in}{0.859781in}}%
\pgfpathlineto{\pgfqpoint{0.828235in}{0.860898in}}%
\pgfpathclose%
\pgfusepath{fill}%
\end{pgfscope}%
\begin{pgfscope}%
\pgfpathrectangle{\pgfqpoint{0.407592in}{0.222904in}}{\pgfqpoint{1.411163in}{1.411163in}}%
\pgfusepath{clip}%
\pgfsetbuttcap%
\pgfsetroundjoin%
\definecolor{currentfill}{rgb}{0.556863,0.654902,0.847059}%
\pgfsetfillcolor{currentfill}%
\pgfsetfillopacity{0.130000}%
\pgfsetlinewidth{0.000000pt}%
\definecolor{currentstroke}{rgb}{0.000000,0.000000,0.000000}%
\pgfsetstrokecolor{currentstroke}%
\pgfsetdash{}{0pt}%
\pgfpathmoveto{\pgfqpoint{1.170978in}{0.744261in}}%
\pgfpathlineto{\pgfqpoint{1.187438in}{0.749888in}}%
\pgfpathlineto{\pgfqpoint{1.203961in}{0.756580in}}%
\pgfpathlineto{\pgfqpoint{1.211916in}{0.761921in}}%
\pgfpathlineto{\pgfqpoint{1.219850in}{0.768138in}}%
\pgfpathlineto{\pgfqpoint{1.203348in}{0.761502in}}%
\pgfpathlineto{\pgfqpoint{1.186909in}{0.755929in}}%
\pgfpathlineto{\pgfqpoint{1.178955in}{0.749601in}}%
\pgfpathlineto{\pgfqpoint{1.170978in}{0.744261in}}%
\pgfpathclose%
\pgfusepath{fill}%
\end{pgfscope}%
\begin{pgfscope}%
\pgfpathrectangle{\pgfqpoint{0.407592in}{0.222904in}}{\pgfqpoint{1.411163in}{1.411163in}}%
\pgfusepath{clip}%
\pgfsetbuttcap%
\pgfsetroundjoin%
\definecolor{currentfill}{rgb}{0.556863,0.654902,0.847059}%
\pgfsetfillcolor{currentfill}%
\pgfsetfillopacity{0.130000}%
\pgfsetlinewidth{0.000000pt}%
\definecolor{currentstroke}{rgb}{0.000000,0.000000,0.000000}%
\pgfsetstrokecolor{currentstroke}%
\pgfsetdash{}{0pt}%
\pgfpathmoveto{\pgfqpoint{0.533675in}{1.151976in}}%
\pgfpathlineto{\pgfqpoint{0.549838in}{1.132820in}}%
\pgfpathlineto{\pgfqpoint{0.566011in}{1.113777in}}%
\pgfpathlineto{\pgfqpoint{0.575099in}{1.113668in}}%
\pgfpathlineto{\pgfqpoint{0.584145in}{1.113863in}}%
\pgfpathlineto{\pgfqpoint{0.568015in}{1.133129in}}%
\pgfpathlineto{\pgfqpoint{0.551894in}{1.152493in}}%
\pgfpathlineto{\pgfqpoint{0.542805in}{1.152089in}}%
\pgfpathlineto{\pgfqpoint{0.533675in}{1.151976in}}%
\pgfpathclose%
\pgfusepath{fill}%
\end{pgfscope}%
\begin{pgfscope}%
\pgfpathrectangle{\pgfqpoint{0.407592in}{0.222904in}}{\pgfqpoint{1.411163in}{1.411163in}}%
\pgfusepath{clip}%
\pgfsetbuttcap%
\pgfsetroundjoin%
\definecolor{currentfill}{rgb}{0.556863,0.654902,0.847059}%
\pgfsetfillcolor{currentfill}%
\pgfsetfillopacity{0.130000}%
\pgfsetlinewidth{0.000000pt}%
\definecolor{currentstroke}{rgb}{0.000000,0.000000,0.000000}%
\pgfsetstrokecolor{currentstroke}%
\pgfsetdash{}{0pt}%
\pgfpathmoveto{\pgfqpoint{1.512102in}{1.019116in}}%
\pgfpathlineto{\pgfqpoint{1.529234in}{1.028178in}}%
\pgfpathlineto{\pgfqpoint{1.546440in}{1.037577in}}%
\pgfpathlineto{\pgfqpoint{1.554027in}{1.049570in}}%
\pgfpathlineto{\pgfqpoint{1.561602in}{1.061805in}}%
\pgfpathlineto{\pgfqpoint{1.544416in}{1.052830in}}%
\pgfpathlineto{\pgfqpoint{1.527304in}{1.044208in}}%
\pgfpathlineto{\pgfqpoint{1.519709in}{1.031542in}}%
\pgfpathlineto{\pgfqpoint{1.512102in}{1.019116in}}%
\pgfpathclose%
\pgfusepath{fill}%
\end{pgfscope}%
\begin{pgfscope}%
\pgfpathrectangle{\pgfqpoint{0.407592in}{0.222904in}}{\pgfqpoint{1.411163in}{1.411163in}}%
\pgfusepath{clip}%
\pgfsetbuttcap%
\pgfsetroundjoin%
\definecolor{currentfill}{rgb}{0.556863,0.654902,0.847059}%
\pgfsetfillcolor{currentfill}%
\pgfsetfillopacity{0.130000}%
\pgfsetlinewidth{0.000000pt}%
\definecolor{currentstroke}{rgb}{0.000000,0.000000,0.000000}%
\pgfsetstrokecolor{currentstroke}%
\pgfsetdash{}{0pt}%
\pgfpathmoveto{\pgfqpoint{0.942637in}{0.773594in}}%
\pgfpathlineto{\pgfqpoint{0.958877in}{0.762819in}}%
\pgfpathlineto{\pgfqpoint{0.975132in}{0.753395in}}%
\pgfpathlineto{\pgfqpoint{0.983505in}{0.751310in}}%
\pgfpathlineto{\pgfqpoint{0.991839in}{0.749886in}}%
\pgfpathlineto{\pgfqpoint{0.975627in}{0.759706in}}%
\pgfpathlineto{\pgfqpoint{0.959431in}{0.770949in}}%
\pgfpathlineto{\pgfqpoint{0.951053in}{0.771946in}}%
\pgfpathlineto{\pgfqpoint{0.942637in}{0.773594in}}%
\pgfpathclose%
\pgfusepath{fill}%
\end{pgfscope}%
\begin{pgfscope}%
\pgfpathrectangle{\pgfqpoint{0.407592in}{0.222904in}}{\pgfqpoint{1.411163in}{1.411163in}}%
\pgfusepath{clip}%
\pgfsetbuttcap%
\pgfsetroundjoin%
\definecolor{currentfill}{rgb}{0.556863,0.654902,0.847059}%
\pgfsetfillcolor{currentfill}%
\pgfsetfillopacity{0.130000}%
\pgfsetlinewidth{0.000000pt}%
\definecolor{currentstroke}{rgb}{0.000000,0.000000,0.000000}%
\pgfsetstrokecolor{currentstroke}%
\pgfsetdash{}{0pt}%
\pgfpathmoveto{\pgfqpoint{1.333551in}{0.852024in}}%
\pgfpathlineto{\pgfqpoint{1.350342in}{0.860999in}}%
\pgfpathlineto{\pgfqpoint{1.367204in}{0.870434in}}%
\pgfpathlineto{\pgfqpoint{1.374970in}{0.880281in}}%
\pgfpathlineto{\pgfqpoint{1.382724in}{0.890597in}}%
\pgfpathlineto{\pgfqpoint{1.365879in}{0.881626in}}%
\pgfpathlineto{\pgfqpoint{1.349107in}{0.873144in}}%
\pgfpathlineto{\pgfqpoint{1.341335in}{0.862336in}}%
\pgfpathlineto{\pgfqpoint{1.333551in}{0.852024in}}%
\pgfpathclose%
\pgfusepath{fill}%
\end{pgfscope}%
\begin{pgfscope}%
\pgfpathrectangle{\pgfqpoint{0.407592in}{0.222904in}}{\pgfqpoint{1.411163in}{1.411163in}}%
\pgfusepath{clip}%
\pgfsetbuttcap%
\pgfsetroundjoin%
\definecolor{currentfill}{rgb}{0.556863,0.654902,0.847059}%
\pgfsetfillcolor{currentfill}%
\pgfsetfillopacity{0.130000}%
\pgfsetlinewidth{0.000000pt}%
\definecolor{currentstroke}{rgb}{0.000000,0.000000,0.000000}%
\pgfsetstrokecolor{currentstroke}%
\pgfsetdash{}{0pt}%
\pgfpathmoveto{\pgfqpoint{0.713492in}{0.964947in}}%
\pgfpathlineto{\pgfqpoint{0.729699in}{0.947286in}}%
\pgfpathlineto{\pgfqpoint{0.745914in}{0.929934in}}%
\pgfpathlineto{\pgfqpoint{0.754707in}{0.928889in}}%
\pgfpathlineto{\pgfqpoint{0.763458in}{0.928253in}}%
\pgfpathlineto{\pgfqpoint{0.747285in}{0.946016in}}%
\pgfpathlineto{\pgfqpoint{0.731121in}{0.964059in}}%
\pgfpathlineto{\pgfqpoint{0.722327in}{0.964308in}}%
\pgfpathlineto{\pgfqpoint{0.713492in}{0.964947in}}%
\pgfpathclose%
\pgfusepath{fill}%
\end{pgfscope}%
\begin{pgfscope}%
\pgfpathrectangle{\pgfqpoint{0.407592in}{0.222904in}}{\pgfqpoint{1.411163in}{1.411163in}}%
\pgfusepath{clip}%
\pgfsetbuttcap%
\pgfsetroundjoin%
\definecolor{currentfill}{rgb}{0.556863,0.654902,0.847059}%
\pgfsetfillcolor{currentfill}%
\pgfsetfillopacity{0.130000}%
\pgfsetlinewidth{0.000000pt}%
\definecolor{currentstroke}{rgb}{0.000000,0.000000,0.000000}%
\pgfsetstrokecolor{currentstroke}%
\pgfsetdash{}{0pt}%
\pgfpathmoveto{\pgfqpoint{1.219850in}{0.768138in}}%
\pgfpathlineto{\pgfqpoint{1.236418in}{0.775670in}}%
\pgfpathlineto{\pgfqpoint{1.253054in}{0.783954in}}%
\pgfpathlineto{\pgfqpoint{1.260949in}{0.790829in}}%
\pgfpathlineto{\pgfqpoint{1.268827in}{0.798448in}}%
\pgfpathlineto{\pgfqpoint{1.252210in}{0.790404in}}%
\pgfpathlineto{\pgfqpoint{1.235661in}{0.783126in}}%
\pgfpathlineto{\pgfqpoint{1.227764in}{0.775216in}}%
\pgfpathlineto{\pgfqpoint{1.219850in}{0.768138in}}%
\pgfpathclose%
\pgfusepath{fill}%
\end{pgfscope}%
\begin{pgfscope}%
\pgfpathrectangle{\pgfqpoint{0.407592in}{0.222904in}}{\pgfqpoint{1.411163in}{1.411163in}}%
\pgfusepath{clip}%
\pgfsetbuttcap%
\pgfsetroundjoin%
\definecolor{currentfill}{rgb}{0.556863,0.654902,0.847059}%
\pgfsetfillcolor{currentfill}%
\pgfsetfillopacity{0.130000}%
\pgfsetlinewidth{0.000000pt}%
\definecolor{currentstroke}{rgb}{0.000000,0.000000,0.000000}%
\pgfsetstrokecolor{currentstroke}%
\pgfsetdash{}{0pt}%
\pgfpathmoveto{\pgfqpoint{0.566011in}{1.113777in}}%
\pgfpathlineto{\pgfqpoint{0.582192in}{1.094859in}}%
\pgfpathlineto{\pgfqpoint{0.598381in}{1.076080in}}%
\pgfpathlineto{\pgfqpoint{0.607427in}{1.075716in}}%
\pgfpathlineto{\pgfqpoint{0.616431in}{1.075670in}}%
\pgfpathlineto{\pgfqpoint{0.600284in}{1.094706in}}%
\pgfpathlineto{\pgfqpoint{0.584145in}{1.113863in}}%
\pgfpathlineto{\pgfqpoint{0.575099in}{1.113668in}}%
\pgfpathlineto{\pgfqpoint{0.566011in}{1.113777in}}%
\pgfpathclose%
\pgfusepath{fill}%
\end{pgfscope}%
\begin{pgfscope}%
\pgfpathrectangle{\pgfqpoint{0.407592in}{0.222904in}}{\pgfqpoint{1.411163in}{1.411163in}}%
\pgfusepath{clip}%
\pgfsetbuttcap%
\pgfsetroundjoin%
\definecolor{currentfill}{rgb}{0.556863,0.654902,0.847059}%
\pgfsetfillcolor{currentfill}%
\pgfsetfillopacity{0.130000}%
\pgfsetlinewidth{0.000000pt}%
\definecolor{currentstroke}{rgb}{0.000000,0.000000,0.000000}%
\pgfsetstrokecolor{currentstroke}%
\pgfsetdash{}{0pt}%
\pgfpathmoveto{\pgfqpoint{1.024323in}{0.735092in}}%
\pgfpathlineto{\pgfqpoint{1.040603in}{0.730320in}}%
\pgfpathlineto{\pgfqpoint{1.056914in}{0.727347in}}%
\pgfpathlineto{\pgfqpoint{1.065127in}{0.726487in}}%
\pgfpathlineto{\pgfqpoint{1.073305in}{0.726467in}}%
\pgfpathlineto{\pgfqpoint{1.057034in}{0.729326in}}%
\pgfpathlineto{\pgfqpoint{1.040795in}{0.734121in}}%
\pgfpathlineto{\pgfqpoint{1.032577in}{0.734186in}}%
\pgfpathlineto{\pgfqpoint{1.024323in}{0.735092in}}%
\pgfpathclose%
\pgfusepath{fill}%
\end{pgfscope}%
\begin{pgfscope}%
\pgfpathrectangle{\pgfqpoint{0.407592in}{0.222904in}}{\pgfqpoint{1.411163in}{1.411163in}}%
\pgfusepath{clip}%
\pgfsetbuttcap%
\pgfsetroundjoin%
\definecolor{currentfill}{rgb}{0.556863,0.654902,0.847059}%
\pgfsetfillcolor{currentfill}%
\pgfsetfillopacity{0.130000}%
\pgfsetlinewidth{0.000000pt}%
\definecolor{currentstroke}{rgb}{0.000000,0.000000,0.000000}%
\pgfsetstrokecolor{currentstroke}%
\pgfsetdash{}{0pt}%
\pgfpathmoveto{\pgfqpoint{1.447440in}{0.953287in}}%
\pgfpathlineto{\pgfqpoint{1.464464in}{0.962535in}}%
\pgfpathlineto{\pgfqpoint{1.481561in}{0.972142in}}%
\pgfpathlineto{\pgfqpoint{1.489214in}{0.983442in}}%
\pgfpathlineto{\pgfqpoint{1.496855in}{0.995049in}}%
\pgfpathlineto{\pgfqpoint{1.479777in}{0.985894in}}%
\pgfpathlineto{\pgfqpoint{1.462772in}{0.977118in}}%
\pgfpathlineto{\pgfqpoint{1.455112in}{0.965047in}}%
\pgfpathlineto{\pgfqpoint{1.447440in}{0.953287in}}%
\pgfpathclose%
\pgfusepath{fill}%
\end{pgfscope}%
\begin{pgfscope}%
\pgfpathrectangle{\pgfqpoint{0.407592in}{0.222904in}}{\pgfqpoint{1.411163in}{1.411163in}}%
\pgfusepath{clip}%
\pgfsetbuttcap%
\pgfsetroundjoin%
\definecolor{currentfill}{rgb}{0.556863,0.654902,0.847059}%
\pgfsetfillcolor{currentfill}%
\pgfsetfillopacity{0.130000}%
\pgfsetlinewidth{0.000000pt}%
\definecolor{currentstroke}{rgb}{0.000000,0.000000,0.000000}%
\pgfsetstrokecolor{currentstroke}%
\pgfsetdash{}{0pt}%
\pgfpathmoveto{\pgfqpoint{0.860677in}{0.830281in}}%
\pgfpathlineto{\pgfqpoint{0.876912in}{0.816011in}}%
\pgfpathlineto{\pgfqpoint{0.893158in}{0.802566in}}%
\pgfpathlineto{\pgfqpoint{0.901698in}{0.800345in}}%
\pgfpathlineto{\pgfqpoint{0.910197in}{0.798651in}}%
\pgfpathlineto{\pgfqpoint{0.893995in}{0.812659in}}%
\pgfpathlineto{\pgfqpoint{0.877802in}{0.827486in}}%
\pgfpathlineto{\pgfqpoint{0.869260in}{0.828625in}}%
\pgfpathlineto{\pgfqpoint{0.860677in}{0.830281in}}%
\pgfpathclose%
\pgfusepath{fill}%
\end{pgfscope}%
\begin{pgfscope}%
\pgfpathrectangle{\pgfqpoint{0.407592in}{0.222904in}}{\pgfqpoint{1.411163in}{1.411163in}}%
\pgfusepath{clip}%
\pgfsetbuttcap%
\pgfsetroundjoin%
\definecolor{currentfill}{rgb}{0.556863,0.654902,0.847059}%
\pgfsetfillcolor{currentfill}%
\pgfsetfillopacity{0.130000}%
\pgfsetlinewidth{0.000000pt}%
\definecolor{currentstroke}{rgb}{0.000000,0.000000,0.000000}%
\pgfsetstrokecolor{currentstroke}%
\pgfsetdash{}{0pt}%
\pgfpathmoveto{\pgfqpoint{1.626303in}{1.130208in}}%
\pgfpathlineto{\pgfqpoint{1.643671in}{1.139337in}}%
\pgfpathlineto{\pgfqpoint{1.661114in}{1.148756in}}%
\pgfpathlineto{\pgfqpoint{1.668592in}{1.161245in}}%
\pgfpathlineto{\pgfqpoint{1.676059in}{1.173911in}}%
\pgfpathlineto{\pgfqpoint{1.658638in}{1.164855in}}%
\pgfpathlineto{\pgfqpoint{1.641291in}{1.156100in}}%
\pgfpathlineto{\pgfqpoint{1.633802in}{1.143067in}}%
\pgfpathlineto{\pgfqpoint{1.626303in}{1.130208in}}%
\pgfpathclose%
\pgfusepath{fill}%
\end{pgfscope}%
\begin{pgfscope}%
\pgfpathrectangle{\pgfqpoint{0.407592in}{0.222904in}}{\pgfqpoint{1.411163in}{1.411163in}}%
\pgfusepath{clip}%
\pgfsetbuttcap%
\pgfsetroundjoin%
\definecolor{currentfill}{rgb}{0.556863,0.654902,0.847059}%
\pgfsetfillcolor{currentfill}%
\pgfsetfillopacity{0.130000}%
\pgfsetlinewidth{0.000000pt}%
\definecolor{currentstroke}{rgb}{0.000000,0.000000,0.000000}%
\pgfsetstrokecolor{currentstroke}%
\pgfsetdash{}{0pt}%
\pgfpathmoveto{\pgfqpoint{0.745914in}{0.929934in}}%
\pgfpathlineto{\pgfqpoint{0.762138in}{0.912929in}}%
\pgfpathlineto{\pgfqpoint{0.778370in}{0.896319in}}%
\pgfpathlineto{\pgfqpoint{0.787120in}{0.894814in}}%
\pgfpathlineto{\pgfqpoint{0.795829in}{0.893733in}}%
\pgfpathlineto{\pgfqpoint{0.779639in}{0.910810in}}%
\pgfpathlineto{\pgfqpoint{0.763458in}{0.928253in}}%
\pgfpathlineto{\pgfqpoint{0.754707in}{0.928889in}}%
\pgfpathlineto{\pgfqpoint{0.745914in}{0.929934in}}%
\pgfpathclose%
\pgfusepath{fill}%
\end{pgfscope}%
\begin{pgfscope}%
\pgfpathrectangle{\pgfqpoint{0.407592in}{0.222904in}}{\pgfqpoint{1.411163in}{1.411163in}}%
\pgfusepath{clip}%
\pgfsetbuttcap%
\pgfsetroundjoin%
\definecolor{currentfill}{rgb}{0.556863,0.654902,0.847059}%
\pgfsetfillcolor{currentfill}%
\pgfsetfillopacity{0.130000}%
\pgfsetlinewidth{0.000000pt}%
\definecolor{currentstroke}{rgb}{0.000000,0.000000,0.000000}%
\pgfsetstrokecolor{currentstroke}%
\pgfsetdash{}{0pt}%
\pgfpathmoveto{\pgfqpoint{0.598381in}{1.076080in}}%
\pgfpathlineto{\pgfqpoint{0.614579in}{1.057456in}}%
\pgfpathlineto{\pgfqpoint{0.630786in}{1.039003in}}%
\pgfpathlineto{\pgfqpoint{0.639789in}{1.038346in}}%
\pgfpathlineto{\pgfqpoint{0.648751in}{1.038021in}}%
\pgfpathlineto{\pgfqpoint{0.632587in}{1.056770in}}%
\pgfpathlineto{\pgfqpoint{0.616431in}{1.075670in}}%
\pgfpathlineto{\pgfqpoint{0.607427in}{1.075716in}}%
\pgfpathlineto{\pgfqpoint{0.598381in}{1.076080in}}%
\pgfpathclose%
\pgfusepath{fill}%
\end{pgfscope}%
\begin{pgfscope}%
\pgfpathrectangle{\pgfqpoint{0.407592in}{0.222904in}}{\pgfqpoint{1.411163in}{1.411163in}}%
\pgfusepath{clip}%
\pgfsetbuttcap%
\pgfsetroundjoin%
\definecolor{currentfill}{rgb}{0.556863,0.654902,0.847059}%
\pgfsetfillcolor{currentfill}%
\pgfsetfillopacity{0.130000}%
\pgfsetlinewidth{0.000000pt}%
\definecolor{currentstroke}{rgb}{0.000000,0.000000,0.000000}%
\pgfsetstrokecolor{currentstroke}%
\pgfsetdash{}{0pt}%
\pgfpathmoveto{\pgfqpoint{1.268827in}{0.798448in}}%
\pgfpathlineto{\pgfqpoint{1.285513in}{0.807141in}}%
\pgfpathlineto{\pgfqpoint{1.302270in}{0.816388in}}%
\pgfpathlineto{\pgfqpoint{1.310112in}{0.824395in}}%
\pgfpathlineto{\pgfqpoint{1.317940in}{0.833022in}}%
\pgfpathlineto{\pgfqpoint{1.301202in}{0.824123in}}%
\pgfpathlineto{\pgfqpoint{1.284534in}{0.815801in}}%
\pgfpathlineto{\pgfqpoint{1.276688in}{0.806783in}}%
\pgfpathlineto{\pgfqpoint{1.268827in}{0.798448in}}%
\pgfpathclose%
\pgfusepath{fill}%
\end{pgfscope}%
\begin{pgfscope}%
\pgfpathrectangle{\pgfqpoint{0.407592in}{0.222904in}}{\pgfqpoint{1.411163in}{1.411163in}}%
\pgfusepath{clip}%
\pgfsetbuttcap%
\pgfsetroundjoin%
\definecolor{currentfill}{rgb}{0.556863,0.654902,0.847059}%
\pgfsetfillcolor{currentfill}%
\pgfsetfillopacity{0.130000}%
\pgfsetlinewidth{0.000000pt}%
\definecolor{currentstroke}{rgb}{0.000000,0.000000,0.000000}%
\pgfsetstrokecolor{currentstroke}%
\pgfsetdash{}{0pt}%
\pgfpathmoveto{\pgfqpoint{1.561602in}{1.061805in}}%
\pgfpathlineto{\pgfqpoint{1.578862in}{1.071106in}}%
\pgfpathlineto{\pgfqpoint{1.596196in}{1.080712in}}%
\pgfpathlineto{\pgfqpoint{1.603740in}{1.092772in}}%
\pgfpathlineto{\pgfqpoint{1.611272in}{1.105049in}}%
\pgfpathlineto{\pgfqpoint{1.593958in}{1.095838in}}%
\pgfpathlineto{\pgfqpoint{1.576718in}{1.086945in}}%
\pgfpathlineto{\pgfqpoint{1.569166in}{1.074267in}}%
\pgfpathlineto{\pgfqpoint{1.561602in}{1.061805in}}%
\pgfpathclose%
\pgfusepath{fill}%
\end{pgfscope}%
\begin{pgfscope}%
\pgfpathrectangle{\pgfqpoint{0.407592in}{0.222904in}}{\pgfqpoint{1.411163in}{1.411163in}}%
\pgfusepath{clip}%
\pgfsetbuttcap%
\pgfsetroundjoin%
\definecolor{currentfill}{rgb}{0.556863,0.654902,0.847059}%
\pgfsetfillcolor{currentfill}%
\pgfsetfillopacity{0.130000}%
\pgfsetlinewidth{0.000000pt}%
\definecolor{currentstroke}{rgb}{0.000000,0.000000,0.000000}%
\pgfsetstrokecolor{currentstroke}%
\pgfsetdash{}{0pt}%
\pgfpathmoveto{\pgfqpoint{1.382724in}{0.890597in}}%
\pgfpathlineto{\pgfqpoint{1.399640in}{0.900001in}}%
\pgfpathlineto{\pgfqpoint{1.416629in}{0.909791in}}%
\pgfpathlineto{\pgfqpoint{1.424350in}{0.920090in}}%
\pgfpathlineto{\pgfqpoint{1.432059in}{0.930787in}}%
\pgfpathlineto{\pgfqpoint{1.415089in}{0.921451in}}%
\pgfpathlineto{\pgfqpoint{1.398191in}{0.912528in}}%
\pgfpathlineto{\pgfqpoint{1.390464in}{0.901355in}}%
\pgfpathlineto{\pgfqpoint{1.382724in}{0.890597in}}%
\pgfpathclose%
\pgfusepath{fill}%
\end{pgfscope}%
\begin{pgfscope}%
\pgfpathrectangle{\pgfqpoint{0.407592in}{0.222904in}}{\pgfqpoint{1.411163in}{1.411163in}}%
\pgfusepath{clip}%
\pgfsetbuttcap%
\pgfsetroundjoin%
\definecolor{currentfill}{rgb}{0.556863,0.654902,0.847059}%
\pgfsetfillcolor{currentfill}%
\pgfsetfillopacity{0.130000}%
\pgfsetlinewidth{0.000000pt}%
\definecolor{currentstroke}{rgb}{0.000000,0.000000,0.000000}%
\pgfsetstrokecolor{currentstroke}%
\pgfsetdash{}{0pt}%
\pgfpathmoveto{\pgfqpoint{0.975132in}{0.753395in}}%
\pgfpathlineto{\pgfqpoint{0.991407in}{0.745446in}}%
\pgfpathlineto{\pgfqpoint{1.007705in}{0.739069in}}%
\pgfpathlineto{\pgfqpoint{1.016032in}{0.736748in}}%
\pgfpathlineto{\pgfqpoint{1.024323in}{0.735092in}}%
\pgfpathlineto{\pgfqpoint{1.008070in}{0.741640in}}%
\pgfpathlineto{\pgfqpoint{0.991839in}{0.749886in}}%
\pgfpathlineto{\pgfqpoint{0.983505in}{0.751310in}}%
\pgfpathlineto{\pgfqpoint{0.975132in}{0.753395in}}%
\pgfpathclose%
\pgfusepath{fill}%
\end{pgfscope}%
\begin{pgfscope}%
\pgfpathrectangle{\pgfqpoint{0.407592in}{0.222904in}}{\pgfqpoint{1.411163in}{1.411163in}}%
\pgfusepath{clip}%
\pgfsetbuttcap%
\pgfsetroundjoin%
\definecolor{currentfill}{rgb}{0.556863,0.654902,0.847059}%
\pgfsetfillcolor{currentfill}%
\pgfsetfillopacity{0.130000}%
\pgfsetlinewidth{0.000000pt}%
\definecolor{currentstroke}{rgb}{0.000000,0.000000,0.000000}%
\pgfsetstrokecolor{currentstroke}%
\pgfsetdash{}{0pt}%
\pgfpathmoveto{\pgfqpoint{0.630786in}{1.039003in}}%
\pgfpathlineto{\pgfqpoint{0.647001in}{1.020741in}}%
\pgfpathlineto{\pgfqpoint{0.663224in}{1.002692in}}%
\pgfpathlineto{\pgfqpoint{0.672186in}{1.001700in}}%
\pgfpathlineto{\pgfqpoint{0.681105in}{1.001054in}}%
\pgfpathlineto{\pgfqpoint{0.664924in}{1.019442in}}%
\pgfpathlineto{\pgfqpoint{0.648751in}{1.038021in}}%
\pgfpathlineto{\pgfqpoint{0.639789in}{1.038346in}}%
\pgfpathlineto{\pgfqpoint{0.630786in}{1.039003in}}%
\pgfpathclose%
\pgfusepath{fill}%
\end{pgfscope}%
\begin{pgfscope}%
\pgfpathrectangle{\pgfqpoint{0.407592in}{0.222904in}}{\pgfqpoint{1.411163in}{1.411163in}}%
\pgfusepath{clip}%
\pgfsetbuttcap%
\pgfsetroundjoin%
\definecolor{currentfill}{rgb}{0.556863,0.654902,0.847059}%
\pgfsetfillcolor{currentfill}%
\pgfsetfillopacity{0.130000}%
\pgfsetlinewidth{0.000000pt}%
\definecolor{currentstroke}{rgb}{0.000000,0.000000,0.000000}%
\pgfsetstrokecolor{currentstroke}%
\pgfsetdash{}{0pt}%
\pgfpathmoveto{\pgfqpoint{1.496855in}{0.995049in}}%
\pgfpathlineto{\pgfqpoint{1.514007in}{1.004551in}}%
\pgfpathlineto{\pgfqpoint{1.531233in}{1.014372in}}%
\pgfpathlineto{\pgfqpoint{1.538842in}{1.025840in}}%
\pgfpathlineto{\pgfqpoint{1.546440in}{1.037577in}}%
\pgfpathlineto{\pgfqpoint{1.529234in}{1.028178in}}%
\pgfpathlineto{\pgfqpoint{1.512102in}{1.019116in}}%
\pgfpathlineto{\pgfqpoint{1.504485in}{1.006946in}}%
\pgfpathlineto{\pgfqpoint{1.496855in}{0.995049in}}%
\pgfpathclose%
\pgfusepath{fill}%
\end{pgfscope}%
\begin{pgfscope}%
\pgfpathrectangle{\pgfqpoint{0.407592in}{0.222904in}}{\pgfqpoint{1.411163in}{1.411163in}}%
\pgfusepath{clip}%
\pgfsetbuttcap%
\pgfsetroundjoin%
\definecolor{currentfill}{rgb}{0.556863,0.654902,0.847059}%
\pgfsetfillcolor{currentfill}%
\pgfsetfillopacity{0.130000}%
\pgfsetlinewidth{0.000000pt}%
\definecolor{currentstroke}{rgb}{0.000000,0.000000,0.000000}%
\pgfsetstrokecolor{currentstroke}%
\pgfsetdash{}{0pt}%
\pgfpathmoveto{\pgfqpoint{0.778370in}{0.896319in}}%
\pgfpathlineto{\pgfqpoint{0.794611in}{0.880157in}}%
\pgfpathlineto{\pgfqpoint{0.810860in}{0.864500in}}%
\pgfpathlineto{\pgfqpoint{0.819568in}{0.862481in}}%
\pgfpathlineto{\pgfqpoint{0.828235in}{0.860898in}}%
\pgfpathlineto{\pgfqpoint{0.812028in}{0.877075in}}%
\pgfpathlineto{\pgfqpoint{0.795829in}{0.893733in}}%
\pgfpathlineto{\pgfqpoint{0.787120in}{0.894814in}}%
\pgfpathlineto{\pgfqpoint{0.778370in}{0.896319in}}%
\pgfpathclose%
\pgfusepath{fill}%
\end{pgfscope}%
\begin{pgfscope}%
\pgfpathrectangle{\pgfqpoint{0.407592in}{0.222904in}}{\pgfqpoint{1.411163in}{1.411163in}}%
\pgfusepath{clip}%
\pgfsetbuttcap%
\pgfsetroundjoin%
\definecolor{currentfill}{rgb}{0.556863,0.654902,0.847059}%
\pgfsetfillcolor{currentfill}%
\pgfsetfillopacity{0.130000}%
\pgfsetlinewidth{0.000000pt}%
\definecolor{currentstroke}{rgb}{0.000000,0.000000,0.000000}%
\pgfsetstrokecolor{currentstroke}%
\pgfsetdash{}{0pt}%
\pgfpathmoveto{\pgfqpoint{0.893158in}{0.802566in}}%
\pgfpathlineto{\pgfqpoint{0.909415in}{0.790053in}}%
\pgfpathlineto{\pgfqpoint{0.925686in}{0.778588in}}%
\pgfpathlineto{\pgfqpoint{0.934181in}{0.775828in}}%
\pgfpathlineto{\pgfqpoint{0.942637in}{0.773594in}}%
\pgfpathlineto{\pgfqpoint{0.926411in}{0.785584in}}%
\pgfpathlineto{\pgfqpoint{0.910197in}{0.798651in}}%
\pgfpathlineto{\pgfqpoint{0.901698in}{0.800345in}}%
\pgfpathlineto{\pgfqpoint{0.893158in}{0.802566in}}%
\pgfpathclose%
\pgfusepath{fill}%
\end{pgfscope}%
\begin{pgfscope}%
\pgfpathrectangle{\pgfqpoint{0.407592in}{0.222904in}}{\pgfqpoint{1.411163in}{1.411163in}}%
\pgfusepath{clip}%
\pgfsetbuttcap%
\pgfsetroundjoin%
\definecolor{currentfill}{rgb}{0.556863,0.654902,0.847059}%
\pgfsetfillcolor{currentfill}%
\pgfsetfillopacity{0.130000}%
\pgfsetlinewidth{0.000000pt}%
\definecolor{currentstroke}{rgb}{0.000000,0.000000,0.000000}%
\pgfsetstrokecolor{currentstroke}%
\pgfsetdash{}{0pt}%
\pgfpathmoveto{\pgfqpoint{1.105968in}{0.726299in}}%
\pgfpathlineto{\pgfqpoint{1.122370in}{0.728719in}}%
\pgfpathlineto{\pgfqpoint{1.138826in}{0.732573in}}%
\pgfpathlineto{\pgfqpoint{1.146904in}{0.734104in}}%
\pgfpathlineto{\pgfqpoint{1.154954in}{0.736535in}}%
\pgfpathlineto{\pgfqpoint{1.138528in}{0.732333in}}%
\pgfpathlineto{\pgfqpoint{1.122158in}{0.729595in}}%
\pgfpathlineto{\pgfqpoint{1.114077in}{0.727461in}}%
\pgfpathlineto{\pgfqpoint{1.105968in}{0.726299in}}%
\pgfpathclose%
\pgfusepath{fill}%
\end{pgfscope}%
\begin{pgfscope}%
\pgfpathrectangle{\pgfqpoint{0.407592in}{0.222904in}}{\pgfqpoint{1.411163in}{1.411163in}}%
\pgfusepath{clip}%
\pgfsetbuttcap%
\pgfsetroundjoin%
\definecolor{currentfill}{rgb}{0.556863,0.654902,0.847059}%
\pgfsetfillcolor{currentfill}%
\pgfsetfillopacity{0.130000}%
\pgfsetlinewidth{0.000000pt}%
\definecolor{currentstroke}{rgb}{0.000000,0.000000,0.000000}%
\pgfsetstrokecolor{currentstroke}%
\pgfsetdash{}{0pt}%
\pgfpathmoveto{\pgfqpoint{1.317940in}{0.833022in}}%
\pgfpathlineto{\pgfqpoint{1.334749in}{0.842420in}}%
\pgfpathlineto{\pgfqpoint{1.351630in}{0.852252in}}%
\pgfpathlineto{\pgfqpoint{1.359424in}{0.861083in}}%
\pgfpathlineto{\pgfqpoint{1.367204in}{0.870434in}}%
\pgfpathlineto{\pgfqpoint{1.350342in}{0.860999in}}%
\pgfpathlineto{\pgfqpoint{1.333551in}{0.852024in}}%
\pgfpathlineto{\pgfqpoint{1.325752in}{0.842242in}}%
\pgfpathlineto{\pgfqpoint{1.317940in}{0.833022in}}%
\pgfpathclose%
\pgfusepath{fill}%
\end{pgfscope}%
\begin{pgfscope}%
\pgfpathrectangle{\pgfqpoint{0.407592in}{0.222904in}}{\pgfqpoint{1.411163in}{1.411163in}}%
\pgfusepath{clip}%
\pgfsetbuttcap%
\pgfsetroundjoin%
\definecolor{currentfill}{rgb}{0.556863,0.654902,0.847059}%
\pgfsetfillcolor{currentfill}%
\pgfsetfillopacity{0.130000}%
\pgfsetlinewidth{0.000000pt}%
\definecolor{currentstroke}{rgb}{0.000000,0.000000,0.000000}%
\pgfsetstrokecolor{currentstroke}%
\pgfsetdash{}{0pt}%
\pgfpathmoveto{\pgfqpoint{1.154954in}{0.736535in}}%
\pgfpathlineto{\pgfqpoint{1.171439in}{0.741996in}}%
\pgfpathlineto{\pgfqpoint{1.187987in}{0.748527in}}%
\pgfpathlineto{\pgfqpoint{1.195986in}{0.752118in}}%
\pgfpathlineto{\pgfqpoint{1.203961in}{0.756580in}}%
\pgfpathlineto{\pgfqpoint{1.187438in}{0.749888in}}%
\pgfpathlineto{\pgfqpoint{1.170978in}{0.744261in}}%
\pgfpathlineto{\pgfqpoint{1.162978in}{0.739910in}}%
\pgfpathlineto{\pgfqpoint{1.154954in}{0.736535in}}%
\pgfpathclose%
\pgfusepath{fill}%
\end{pgfscope}%
\begin{pgfscope}%
\pgfpathrectangle{\pgfqpoint{0.407592in}{0.222904in}}{\pgfqpoint{1.411163in}{1.411163in}}%
\pgfusepath{clip}%
\pgfsetbuttcap%
\pgfsetroundjoin%
\definecolor{currentfill}{rgb}{0.556863,0.654902,0.847059}%
\pgfsetfillcolor{currentfill}%
\pgfsetfillopacity{0.130000}%
\pgfsetlinewidth{0.000000pt}%
\definecolor{currentstroke}{rgb}{0.000000,0.000000,0.000000}%
\pgfsetstrokecolor{currentstroke}%
\pgfsetdash{}{0pt}%
\pgfpathmoveto{\pgfqpoint{0.663224in}{1.002692in}}%
\pgfpathlineto{\pgfqpoint{0.679456in}{0.984880in}}%
\pgfpathlineto{\pgfqpoint{0.695696in}{0.967335in}}%
\pgfpathlineto{\pgfqpoint{0.704615in}{0.965961in}}%
\pgfpathlineto{\pgfqpoint{0.713492in}{0.964947in}}%
\pgfpathlineto{\pgfqpoint{0.697294in}{0.982880in}}%
\pgfpathlineto{\pgfqpoint{0.681105in}{1.001054in}}%
\pgfpathlineto{\pgfqpoint{0.672186in}{1.001700in}}%
\pgfpathlineto{\pgfqpoint{0.663224in}{1.002692in}}%
\pgfpathclose%
\pgfusepath{fill}%
\end{pgfscope}%
\begin{pgfscope}%
\pgfpathrectangle{\pgfqpoint{0.407592in}{0.222904in}}{\pgfqpoint{1.411163in}{1.411163in}}%
\pgfusepath{clip}%
\pgfsetbuttcap%
\pgfsetroundjoin%
\definecolor{currentfill}{rgb}{0.556863,0.654902,0.847059}%
\pgfsetfillcolor{currentfill}%
\pgfsetfillopacity{0.130000}%
\pgfsetlinewidth{0.000000pt}%
\definecolor{currentstroke}{rgb}{0.000000,0.000000,0.000000}%
\pgfsetstrokecolor{currentstroke}%
\pgfsetdash{}{0pt}%
\pgfpathmoveto{\pgfqpoint{1.432059in}{0.930787in}}%
\pgfpathlineto{\pgfqpoint{1.449103in}{0.940493in}}%
\pgfpathlineto{\pgfqpoint{1.466220in}{0.950535in}}%
\pgfpathlineto{\pgfqpoint{1.473897in}{0.961167in}}%
\pgfpathlineto{\pgfqpoint{1.481561in}{0.972142in}}%
\pgfpathlineto{\pgfqpoint{1.464464in}{0.962535in}}%
\pgfpathlineto{\pgfqpoint{1.447440in}{0.953287in}}%
\pgfpathlineto{\pgfqpoint{1.439755in}{0.941860in}}%
\pgfpathlineto{\pgfqpoint{1.432059in}{0.930787in}}%
\pgfpathclose%
\pgfusepath{fill}%
\end{pgfscope}%
\begin{pgfscope}%
\pgfpathrectangle{\pgfqpoint{0.407592in}{0.222904in}}{\pgfqpoint{1.411163in}{1.411163in}}%
\pgfusepath{clip}%
\pgfsetbuttcap%
\pgfsetroundjoin%
\definecolor{currentfill}{rgb}{0.556863,0.654902,0.847059}%
\pgfsetfillcolor{currentfill}%
\pgfsetfillopacity{0.130000}%
\pgfsetlinewidth{0.000000pt}%
\definecolor{currentstroke}{rgb}{0.000000,0.000000,0.000000}%
\pgfsetstrokecolor{currentstroke}%
\pgfsetdash{}{0pt}%
\pgfpathmoveto{\pgfqpoint{1.676059in}{1.173911in}}%
\pgfpathlineto{\pgfqpoint{1.693557in}{1.183250in}}%
\pgfpathlineto{\pgfqpoint{1.711130in}{1.192858in}}%
\pgfpathlineto{\pgfqpoint{1.718564in}{1.205346in}}%
\pgfpathlineto{\pgfqpoint{1.701002in}{1.195909in}}%
\pgfpathlineto{\pgfqpoint{1.683516in}{1.186745in}}%
\pgfusepath{fill}%
\end{pgfscope}%
\begin{pgfscope}%
\pgfpathrectangle{\pgfqpoint{0.407592in}{0.222904in}}{\pgfqpoint{1.411163in}{1.411163in}}%
\pgfusepath{clip}%
\pgfsetbuttcap%
\pgfsetroundjoin%
\definecolor{currentfill}{rgb}{0.556863,0.654902,0.847059}%
\pgfsetfillcolor{currentfill}%
\pgfsetfillopacity{0.130000}%
\pgfsetlinewidth{0.000000pt}%
\definecolor{currentstroke}{rgb}{0.000000,0.000000,0.000000}%
\pgfsetstrokecolor{currentstroke}%
\pgfsetdash{}{0pt}%
\pgfpathmoveto{\pgfqpoint{1.611272in}{1.105049in}}%
\pgfpathlineto{\pgfqpoint{1.628661in}{1.114558in}}%
\pgfpathlineto{\pgfqpoint{1.646125in}{1.124345in}}%
\pgfpathlineto{\pgfqpoint{1.653625in}{1.136453in}}%
\pgfpathlineto{\pgfqpoint{1.661114in}{1.148756in}}%
\pgfpathlineto{\pgfqpoint{1.643671in}{1.139337in}}%
\pgfpathlineto{\pgfqpoint{1.626303in}{1.130208in}}%
\pgfpathlineto{\pgfqpoint{1.618793in}{1.117531in}}%
\pgfpathlineto{\pgfqpoint{1.611272in}{1.105049in}}%
\pgfpathclose%
\pgfusepath{fill}%
\end{pgfscope}%
\begin{pgfscope}%
\pgfpathrectangle{\pgfqpoint{0.407592in}{0.222904in}}{\pgfqpoint{1.411163in}{1.411163in}}%
\pgfusepath{clip}%
\pgfsetbuttcap%
\pgfsetroundjoin%
\definecolor{currentfill}{rgb}{0.556863,0.654902,0.847059}%
\pgfsetfillcolor{currentfill}%
\pgfsetfillopacity{0.130000}%
\pgfsetlinewidth{0.000000pt}%
\definecolor{currentstroke}{rgb}{0.000000,0.000000,0.000000}%
\pgfsetstrokecolor{currentstroke}%
\pgfsetdash{}{0pt}%
\pgfpathmoveto{\pgfqpoint{1.056914in}{0.727347in}}%
\pgfpathlineto{\pgfqpoint{1.073263in}{0.726135in}}%
\pgfpathlineto{\pgfqpoint{1.089654in}{0.726594in}}%
\pgfpathlineto{\pgfqpoint{1.097827in}{0.726036in}}%
\pgfpathlineto{\pgfqpoint{1.105968in}{0.726299in}}%
\pgfpathlineto{\pgfqpoint{1.089614in}{0.725496in}}%
\pgfpathlineto{\pgfqpoint{1.073305in}{0.726467in}}%
\pgfpathlineto{\pgfqpoint{1.065127in}{0.726487in}}%
\pgfpathlineto{\pgfqpoint{1.056914in}{0.727347in}}%
\pgfpathclose%
\pgfusepath{fill}%
\end{pgfscope}%
\begin{pgfscope}%
\pgfpathrectangle{\pgfqpoint{0.407592in}{0.222904in}}{\pgfqpoint{1.411163in}{1.411163in}}%
\pgfusepath{clip}%
\pgfsetbuttcap%
\pgfsetroundjoin%
\definecolor{currentfill}{rgb}{0.556863,0.654902,0.847059}%
\pgfsetfillcolor{currentfill}%
\pgfsetfillopacity{0.130000}%
\pgfsetlinewidth{0.000000pt}%
\definecolor{currentstroke}{rgb}{0.000000,0.000000,0.000000}%
\pgfsetstrokecolor{currentstroke}%
\pgfsetdash{}{0pt}%
\pgfpathmoveto{\pgfqpoint{0.515285in}{1.152597in}}%
\pgfpathlineto{\pgfqpoint{0.531491in}{1.133673in}}%
\pgfpathlineto{\pgfqpoint{0.547706in}{1.114878in}}%
\pgfpathlineto{\pgfqpoint{0.556880in}{1.114183in}}%
\pgfpathlineto{\pgfqpoint{0.566011in}{1.113777in}}%
\pgfpathlineto{\pgfqpoint{0.549838in}{1.132820in}}%
\pgfpathlineto{\pgfqpoint{0.533675in}{1.151976in}}%
\pgfpathlineto{\pgfqpoint{0.524501in}{1.152148in}}%
\pgfpathlineto{\pgfqpoint{0.515285in}{1.152597in}}%
\pgfpathclose%
\pgfusepath{fill}%
\end{pgfscope}%
\begin{pgfscope}%
\pgfpathrectangle{\pgfqpoint{0.407592in}{0.222904in}}{\pgfqpoint{1.411163in}{1.411163in}}%
\pgfusepath{clip}%
\pgfsetbuttcap%
\pgfsetroundjoin%
\definecolor{currentfill}{rgb}{0.556863,0.654902,0.847059}%
\pgfsetfillcolor{currentfill}%
\pgfsetfillopacity{0.130000}%
\pgfsetlinewidth{0.000000pt}%
\definecolor{currentstroke}{rgb}{0.000000,0.000000,0.000000}%
\pgfsetstrokecolor{currentstroke}%
\pgfsetdash{}{0pt}%
\pgfpathmoveto{\pgfqpoint{1.203961in}{0.756580in}}%
\pgfpathlineto{\pgfqpoint{1.220551in}{0.764167in}}%
\pgfpathlineto{\pgfqpoint{1.237207in}{0.772505in}}%
\pgfpathlineto{\pgfqpoint{1.245140in}{0.777842in}}%
\pgfpathlineto{\pgfqpoint{1.253054in}{0.783954in}}%
\pgfpathlineto{\pgfqpoint{1.236418in}{0.775670in}}%
\pgfpathlineto{\pgfqpoint{1.219850in}{0.768138in}}%
\pgfpathlineto{\pgfqpoint{1.211916in}{0.761921in}}%
\pgfpathlineto{\pgfqpoint{1.203961in}{0.756580in}}%
\pgfpathclose%
\pgfusepath{fill}%
\end{pgfscope}%
\begin{pgfscope}%
\pgfpathrectangle{\pgfqpoint{0.407592in}{0.222904in}}{\pgfqpoint{1.411163in}{1.411163in}}%
\pgfusepath{clip}%
\pgfsetbuttcap%
\pgfsetroundjoin%
\definecolor{currentfill}{rgb}{0.556863,0.654902,0.847059}%
\pgfsetfillcolor{currentfill}%
\pgfsetfillopacity{0.130000}%
\pgfsetlinewidth{0.000000pt}%
\definecolor{currentstroke}{rgb}{0.000000,0.000000,0.000000}%
\pgfsetstrokecolor{currentstroke}%
\pgfsetdash{}{0pt}%
\pgfpathmoveto{\pgfqpoint{0.810860in}{0.864500in}}%
\pgfpathlineto{\pgfqpoint{0.827120in}{0.849418in}}%
\pgfpathlineto{\pgfqpoint{0.843388in}{0.834985in}}%
\pgfpathlineto{\pgfqpoint{0.852053in}{0.832413in}}%
\pgfpathlineto{\pgfqpoint{0.860677in}{0.830281in}}%
\pgfpathlineto{\pgfqpoint{0.844451in}{0.845273in}}%
\pgfpathlineto{\pgfqpoint{0.828235in}{0.860898in}}%
\pgfpathlineto{\pgfqpoint{0.819568in}{0.862481in}}%
\pgfpathlineto{\pgfqpoint{0.810860in}{0.864500in}}%
\pgfpathclose%
\pgfusepath{fill}%
\end{pgfscope}%
\begin{pgfscope}%
\pgfpathrectangle{\pgfqpoint{0.407592in}{0.222904in}}{\pgfqpoint{1.411163in}{1.411163in}}%
\pgfusepath{clip}%
\pgfsetbuttcap%
\pgfsetroundjoin%
\definecolor{currentfill}{rgb}{0.556863,0.654902,0.847059}%
\pgfsetfillcolor{currentfill}%
\pgfsetfillopacity{0.130000}%
\pgfsetlinewidth{0.000000pt}%
\definecolor{currentstroke}{rgb}{0.000000,0.000000,0.000000}%
\pgfsetstrokecolor{currentstroke}%
\pgfsetdash{}{0pt}%
\pgfpathmoveto{\pgfqpoint{1.546440in}{1.037577in}}%
\pgfpathlineto{\pgfqpoint{1.563721in}{1.047288in}}%
\pgfpathlineto{\pgfqpoint{1.581076in}{1.057287in}}%
\pgfpathlineto{\pgfqpoint{1.588642in}{1.068879in}}%
\pgfpathlineto{\pgfqpoint{1.596196in}{1.080712in}}%
\pgfpathlineto{\pgfqpoint{1.578862in}{1.071106in}}%
\pgfpathlineto{\pgfqpoint{1.561602in}{1.061805in}}%
\pgfpathlineto{\pgfqpoint{1.554027in}{1.049570in}}%
\pgfpathlineto{\pgfqpoint{1.546440in}{1.037577in}}%
\pgfpathclose%
\pgfusepath{fill}%
\end{pgfscope}%
\begin{pgfscope}%
\pgfpathrectangle{\pgfqpoint{0.407592in}{0.222904in}}{\pgfqpoint{1.411163in}{1.411163in}}%
\pgfusepath{clip}%
\pgfsetbuttcap%
\pgfsetroundjoin%
\definecolor{currentfill}{rgb}{0.556863,0.654902,0.847059}%
\pgfsetfillcolor{currentfill}%
\pgfsetfillopacity{0.130000}%
\pgfsetlinewidth{0.000000pt}%
\definecolor{currentstroke}{rgb}{0.000000,0.000000,0.000000}%
\pgfsetstrokecolor{currentstroke}%
\pgfsetdash{}{0pt}%
\pgfpathmoveto{\pgfqpoint{0.695696in}{0.967335in}}%
\pgfpathlineto{\pgfqpoint{0.711945in}{0.950086in}}%
\pgfpathlineto{\pgfqpoint{0.728202in}{0.933171in}}%
\pgfpathlineto{\pgfqpoint{0.737079in}{0.931368in}}%
\pgfpathlineto{\pgfqpoint{0.745914in}{0.929934in}}%
\pgfpathlineto{\pgfqpoint{0.729699in}{0.947286in}}%
\pgfpathlineto{\pgfqpoint{0.713492in}{0.964947in}}%
\pgfpathlineto{\pgfqpoint{0.704615in}{0.965961in}}%
\pgfpathlineto{\pgfqpoint{0.695696in}{0.967335in}}%
\pgfpathclose%
\pgfusepath{fill}%
\end{pgfscope}%
\begin{pgfscope}%
\pgfpathrectangle{\pgfqpoint{0.407592in}{0.222904in}}{\pgfqpoint{1.411163in}{1.411163in}}%
\pgfusepath{clip}%
\pgfsetbuttcap%
\pgfsetroundjoin%
\definecolor{currentfill}{rgb}{0.556863,0.654902,0.847059}%
\pgfsetfillcolor{currentfill}%
\pgfsetfillopacity{0.130000}%
\pgfsetlinewidth{0.000000pt}%
\definecolor{currentstroke}{rgb}{0.000000,0.000000,0.000000}%
\pgfsetstrokecolor{currentstroke}%
\pgfsetdash{}{0pt}%
\pgfpathmoveto{\pgfqpoint{0.925686in}{0.778588in}}%
\pgfpathlineto{\pgfqpoint{0.941970in}{0.768287in}}%
\pgfpathlineto{\pgfqpoint{0.958272in}{0.759259in}}%
\pgfpathlineto{\pgfqpoint{0.966722in}{0.756067in}}%
\pgfpathlineto{\pgfqpoint{0.975132in}{0.753395in}}%
\pgfpathlineto{\pgfqpoint{0.958877in}{0.762819in}}%
\pgfpathlineto{\pgfqpoint{0.942637in}{0.773594in}}%
\pgfpathlineto{\pgfqpoint{0.934181in}{0.775828in}}%
\pgfpathlineto{\pgfqpoint{0.925686in}{0.778588in}}%
\pgfpathclose%
\pgfusepath{fill}%
\end{pgfscope}%
\begin{pgfscope}%
\pgfpathrectangle{\pgfqpoint{0.407592in}{0.222904in}}{\pgfqpoint{1.411163in}{1.411163in}}%
\pgfusepath{clip}%
\pgfsetbuttcap%
\pgfsetroundjoin%
\definecolor{currentfill}{rgb}{0.556863,0.654902,0.847059}%
\pgfsetfillcolor{currentfill}%
\pgfsetfillopacity{0.130000}%
\pgfsetlinewidth{0.000000pt}%
\definecolor{currentstroke}{rgb}{0.000000,0.000000,0.000000}%
\pgfsetstrokecolor{currentstroke}%
\pgfsetdash{}{0pt}%
\pgfpathmoveto{\pgfqpoint{1.367204in}{0.870434in}}%
\pgfpathlineto{\pgfqpoint{1.384140in}{0.880275in}}%
\pgfpathlineto{\pgfqpoint{1.401148in}{0.890475in}}%
\pgfpathlineto{\pgfqpoint{1.408895in}{0.899912in}}%
\pgfpathlineto{\pgfqpoint{1.416629in}{0.909791in}}%
\pgfpathlineto{\pgfqpoint{1.399640in}{0.900001in}}%
\pgfpathlineto{\pgfqpoint{1.382724in}{0.890597in}}%
\pgfpathlineto{\pgfqpoint{1.374970in}{0.880281in}}%
\pgfpathlineto{\pgfqpoint{1.367204in}{0.870434in}}%
\pgfpathclose%
\pgfusepath{fill}%
\end{pgfscope}%
\begin{pgfscope}%
\pgfpathrectangle{\pgfqpoint{0.407592in}{0.222904in}}{\pgfqpoint{1.411163in}{1.411163in}}%
\pgfusepath{clip}%
\pgfsetbuttcap%
\pgfsetroundjoin%
\definecolor{currentfill}{rgb}{0.556863,0.654902,0.847059}%
\pgfsetfillcolor{currentfill}%
\pgfsetfillopacity{0.130000}%
\pgfsetlinewidth{0.000000pt}%
\definecolor{currentstroke}{rgb}{0.000000,0.000000,0.000000}%
\pgfsetstrokecolor{currentstroke}%
\pgfsetdash{}{0pt}%
\pgfpathmoveto{\pgfqpoint{0.547706in}{1.114878in}}%
\pgfpathlineto{\pgfqpoint{0.563929in}{1.096224in}}%
\pgfpathlineto{\pgfqpoint{0.580161in}{1.077728in}}%
\pgfpathlineto{\pgfqpoint{0.589292in}{1.076754in}}%
\pgfpathlineto{\pgfqpoint{0.598381in}{1.076080in}}%
\pgfpathlineto{\pgfqpoint{0.582192in}{1.094859in}}%
\pgfpathlineto{\pgfqpoint{0.566011in}{1.113777in}}%
\pgfpathlineto{\pgfqpoint{0.556880in}{1.114183in}}%
\pgfpathlineto{\pgfqpoint{0.547706in}{1.114878in}}%
\pgfpathclose%
\pgfusepath{fill}%
\end{pgfscope}%
\begin{pgfscope}%
\pgfpathrectangle{\pgfqpoint{0.407592in}{0.222904in}}{\pgfqpoint{1.411163in}{1.411163in}}%
\pgfusepath{clip}%
\pgfsetbuttcap%
\pgfsetroundjoin%
\definecolor{currentfill}{rgb}{0.556863,0.654902,0.847059}%
\pgfsetfillcolor{currentfill}%
\pgfsetfillopacity{0.130000}%
\pgfsetlinewidth{0.000000pt}%
\definecolor{currentstroke}{rgb}{0.000000,0.000000,0.000000}%
\pgfsetstrokecolor{currentstroke}%
\pgfsetdash{}{0pt}%
\pgfpathmoveto{\pgfqpoint{1.253054in}{0.783954in}}%
\pgfpathlineto{\pgfqpoint{1.269759in}{0.792873in}}%
\pgfpathlineto{\pgfqpoint{1.286535in}{0.802331in}}%
\pgfpathlineto{\pgfqpoint{1.294411in}{0.809027in}}%
\pgfpathlineto{\pgfqpoint{1.302270in}{0.816388in}}%
\pgfpathlineto{\pgfqpoint{1.285513in}{0.807141in}}%
\pgfpathlineto{\pgfqpoint{1.268827in}{0.798448in}}%
\pgfpathlineto{\pgfqpoint{1.260949in}{0.790829in}}%
\pgfpathlineto{\pgfqpoint{1.253054in}{0.783954in}}%
\pgfpathclose%
\pgfusepath{fill}%
\end{pgfscope}%
\begin{pgfscope}%
\pgfpathrectangle{\pgfqpoint{0.407592in}{0.222904in}}{\pgfqpoint{1.411163in}{1.411163in}}%
\pgfusepath{clip}%
\pgfsetbuttcap%
\pgfsetroundjoin%
\definecolor{currentfill}{rgb}{0.556863,0.654902,0.847059}%
\pgfsetfillcolor{currentfill}%
\pgfsetfillopacity{0.130000}%
\pgfsetlinewidth{0.000000pt}%
\definecolor{currentstroke}{rgb}{0.000000,0.000000,0.000000}%
\pgfsetstrokecolor{currentstroke}%
\pgfsetdash{}{0pt}%
\pgfpathmoveto{\pgfqpoint{1.007705in}{0.739069in}}%
\pgfpathlineto{\pgfqpoint{1.024029in}{0.734327in}}%
\pgfpathlineto{\pgfqpoint{1.040384in}{0.731236in}}%
\pgfpathlineto{\pgfqpoint{1.048667in}{0.728959in}}%
\pgfpathlineto{\pgfqpoint{1.056914in}{0.727347in}}%
\pgfpathlineto{\pgfqpoint{1.040603in}{0.730320in}}%
\pgfpathlineto{\pgfqpoint{1.024323in}{0.735092in}}%
\pgfpathlineto{\pgfqpoint{1.016032in}{0.736748in}}%
\pgfpathlineto{\pgfqpoint{1.007705in}{0.739069in}}%
\pgfpathclose%
\pgfusepath{fill}%
\end{pgfscope}%
\begin{pgfscope}%
\pgfpathrectangle{\pgfqpoint{0.407592in}{0.222904in}}{\pgfqpoint{1.411163in}{1.411163in}}%
\pgfusepath{clip}%
\pgfsetbuttcap%
\pgfsetroundjoin%
\definecolor{currentfill}{rgb}{0.556863,0.654902,0.847059}%
\pgfsetfillcolor{currentfill}%
\pgfsetfillopacity{0.130000}%
\pgfsetlinewidth{0.000000pt}%
\definecolor{currentstroke}{rgb}{0.000000,0.000000,0.000000}%
\pgfsetstrokecolor{currentstroke}%
\pgfsetdash{}{0pt}%
\pgfpathmoveto{\pgfqpoint{1.481561in}{0.972142in}}%
\pgfpathlineto{\pgfqpoint{1.498733in}{0.982075in}}%
\pgfpathlineto{\pgfqpoint{1.515980in}{0.992308in}}%
\pgfpathlineto{\pgfqpoint{1.523612in}{1.003190in}}%
\pgfpathlineto{\pgfqpoint{1.531233in}{1.014372in}}%
\pgfpathlineto{\pgfqpoint{1.514007in}{1.004551in}}%
\pgfpathlineto{\pgfqpoint{1.496855in}{0.995049in}}%
\pgfpathlineto{\pgfqpoint{1.489214in}{0.983442in}}%
\pgfpathlineto{\pgfqpoint{1.481561in}{0.972142in}}%
\pgfpathclose%
\pgfusepath{fill}%
\end{pgfscope}%
\begin{pgfscope}%
\pgfpathrectangle{\pgfqpoint{0.407592in}{0.222904in}}{\pgfqpoint{1.411163in}{1.411163in}}%
\pgfusepath{clip}%
\pgfsetbuttcap%
\pgfsetroundjoin%
\definecolor{currentfill}{rgb}{0.556863,0.654902,0.847059}%
\pgfsetfillcolor{currentfill}%
\pgfsetfillopacity{0.130000}%
\pgfsetlinewidth{0.000000pt}%
\definecolor{currentstroke}{rgb}{0.000000,0.000000,0.000000}%
\pgfsetstrokecolor{currentstroke}%
\pgfsetdash{}{0pt}%
\pgfpathmoveto{\pgfqpoint{0.843388in}{0.834985in}}%
\pgfpathlineto{\pgfqpoint{0.859667in}{0.821282in}}%
\pgfpathlineto{\pgfqpoint{0.875957in}{0.808400in}}%
\pgfpathlineto{\pgfqpoint{0.884578in}{0.805265in}}%
\pgfpathlineto{\pgfqpoint{0.893158in}{0.802566in}}%
\pgfpathlineto{\pgfqpoint{0.876912in}{0.816011in}}%
\pgfpathlineto{\pgfqpoint{0.860677in}{0.830281in}}%
\pgfpathlineto{\pgfqpoint{0.852053in}{0.832413in}}%
\pgfpathlineto{\pgfqpoint{0.843388in}{0.834985in}}%
\pgfpathclose%
\pgfusepath{fill}%
\end{pgfscope}%
\begin{pgfscope}%
\pgfpathrectangle{\pgfqpoint{0.407592in}{0.222904in}}{\pgfqpoint{1.411163in}{1.411163in}}%
\pgfusepath{clip}%
\pgfsetbuttcap%
\pgfsetroundjoin%
\definecolor{currentfill}{rgb}{0.556863,0.654902,0.847059}%
\pgfsetfillcolor{currentfill}%
\pgfsetfillopacity{0.130000}%
\pgfsetlinewidth{0.000000pt}%
\definecolor{currentstroke}{rgb}{0.000000,0.000000,0.000000}%
\pgfsetstrokecolor{currentstroke}%
\pgfsetdash{}{0pt}%
\pgfpathmoveto{\pgfqpoint{1.661114in}{1.148756in}}%
\pgfpathlineto{\pgfqpoint{1.678633in}{1.158449in}}%
\pgfpathlineto{\pgfqpoint{1.696228in}{1.168401in}}%
\pgfpathlineto{\pgfqpoint{1.703684in}{1.180540in}}%
\pgfpathlineto{\pgfqpoint{1.711130in}{1.192858in}}%
\pgfpathlineto{\pgfqpoint{1.693557in}{1.183250in}}%
\pgfpathlineto{\pgfqpoint{1.676059in}{1.173911in}}%
\pgfpathlineto{\pgfqpoint{1.668592in}{1.161245in}}%
\pgfpathlineto{\pgfqpoint{1.661114in}{1.148756in}}%
\pgfpathclose%
\pgfusepath{fill}%
\end{pgfscope}%
\begin{pgfscope}%
\pgfpathrectangle{\pgfqpoint{0.407592in}{0.222904in}}{\pgfqpoint{1.411163in}{1.411163in}}%
\pgfusepath{clip}%
\pgfsetbuttcap%
\pgfsetroundjoin%
\definecolor{currentfill}{rgb}{0.556863,0.654902,0.847059}%
\pgfsetfillcolor{currentfill}%
\pgfsetfillopacity{0.130000}%
\pgfsetlinewidth{0.000000pt}%
\definecolor{currentstroke}{rgb}{0.000000,0.000000,0.000000}%
\pgfsetstrokecolor{currentstroke}%
\pgfsetdash{}{0pt}%
\pgfpathmoveto{\pgfqpoint{0.580161in}{1.077728in}}%
\pgfpathlineto{\pgfqpoint{0.596401in}{1.059404in}}%
\pgfpathlineto{\pgfqpoint{0.612650in}{1.041271in}}%
\pgfpathlineto{\pgfqpoint{0.621739in}{1.039982in}}%
\pgfpathlineto{\pgfqpoint{0.630786in}{1.039003in}}%
\pgfpathlineto{\pgfqpoint{0.614579in}{1.057456in}}%
\pgfpathlineto{\pgfqpoint{0.598381in}{1.076080in}}%
\pgfpathlineto{\pgfqpoint{0.589292in}{1.076754in}}%
\pgfpathlineto{\pgfqpoint{0.580161in}{1.077728in}}%
\pgfpathclose%
\pgfusepath{fill}%
\end{pgfscope}%
\begin{pgfscope}%
\pgfpathrectangle{\pgfqpoint{0.407592in}{0.222904in}}{\pgfqpoint{1.411163in}{1.411163in}}%
\pgfusepath{clip}%
\pgfsetbuttcap%
\pgfsetroundjoin%
\definecolor{currentfill}{rgb}{0.556863,0.654902,0.847059}%
\pgfsetfillcolor{currentfill}%
\pgfsetfillopacity{0.130000}%
\pgfsetlinewidth{0.000000pt}%
\definecolor{currentstroke}{rgb}{0.000000,0.000000,0.000000}%
\pgfsetstrokecolor{currentstroke}%
\pgfsetdash{}{0pt}%
\pgfpathmoveto{\pgfqpoint{0.728202in}{0.933171in}}%
\pgfpathlineto{\pgfqpoint{0.744468in}{0.916630in}}%
\pgfpathlineto{\pgfqpoint{0.760743in}{0.900507in}}%
\pgfpathlineto{\pgfqpoint{0.769577in}{0.898225in}}%
\pgfpathlineto{\pgfqpoint{0.778370in}{0.896319in}}%
\pgfpathlineto{\pgfqpoint{0.762138in}{0.912929in}}%
\pgfpathlineto{\pgfqpoint{0.745914in}{0.929934in}}%
\pgfpathlineto{\pgfqpoint{0.737079in}{0.931368in}}%
\pgfpathlineto{\pgfqpoint{0.728202in}{0.933171in}}%
\pgfpathclose%
\pgfusepath{fill}%
\end{pgfscope}%
\begin{pgfscope}%
\pgfpathrectangle{\pgfqpoint{0.407592in}{0.222904in}}{\pgfqpoint{1.411163in}{1.411163in}}%
\pgfusepath{clip}%
\pgfsetbuttcap%
\pgfsetroundjoin%
\definecolor{currentfill}{rgb}{0.556863,0.654902,0.847059}%
\pgfsetfillcolor{currentfill}%
\pgfsetfillopacity{0.130000}%
\pgfsetlinewidth{0.000000pt}%
\definecolor{currentstroke}{rgb}{0.000000,0.000000,0.000000}%
\pgfsetstrokecolor{currentstroke}%
\pgfsetdash{}{0pt}%
\pgfpathmoveto{\pgfqpoint{1.302270in}{0.816388in}}%
\pgfpathlineto{\pgfqpoint{1.319098in}{0.826110in}}%
\pgfpathlineto{\pgfqpoint{1.335998in}{0.836242in}}%
\pgfpathlineto{\pgfqpoint{1.343822in}{0.843965in}}%
\pgfpathlineto{\pgfqpoint{1.351630in}{0.852252in}}%
\pgfpathlineto{\pgfqpoint{1.334749in}{0.842420in}}%
\pgfpathlineto{\pgfqpoint{1.317940in}{0.833022in}}%
\pgfpathlineto{\pgfqpoint{1.310112in}{0.824395in}}%
\pgfpathlineto{\pgfqpoint{1.302270in}{0.816388in}}%
\pgfpathclose%
\pgfusepath{fill}%
\end{pgfscope}%
\begin{pgfscope}%
\pgfpathrectangle{\pgfqpoint{0.407592in}{0.222904in}}{\pgfqpoint{1.411163in}{1.411163in}}%
\pgfusepath{clip}%
\pgfsetbuttcap%
\pgfsetroundjoin%
\definecolor{currentfill}{rgb}{0.556863,0.654902,0.847059}%
\pgfsetfillcolor{currentfill}%
\pgfsetfillopacity{0.130000}%
\pgfsetlinewidth{0.000000pt}%
\definecolor{currentstroke}{rgb}{0.000000,0.000000,0.000000}%
\pgfsetstrokecolor{currentstroke}%
\pgfsetdash{}{0pt}%
\pgfpathmoveto{\pgfqpoint{1.596196in}{1.080712in}}%
\pgfpathlineto{\pgfqpoint{1.613607in}{1.090602in}}%
\pgfpathlineto{\pgfqpoint{1.631093in}{1.100757in}}%
\pgfpathlineto{\pgfqpoint{1.638615in}{1.112443in}}%
\pgfpathlineto{\pgfqpoint{1.646125in}{1.124345in}}%
\pgfpathlineto{\pgfqpoint{1.628661in}{1.114558in}}%
\pgfpathlineto{\pgfqpoint{1.611272in}{1.105049in}}%
\pgfpathlineto{\pgfqpoint{1.603740in}{1.092772in}}%
\pgfpathlineto{\pgfqpoint{1.596196in}{1.080712in}}%
\pgfpathclose%
\pgfusepath{fill}%
\end{pgfscope}%
\begin{pgfscope}%
\pgfpathrectangle{\pgfqpoint{0.407592in}{0.222904in}}{\pgfqpoint{1.411163in}{1.411163in}}%
\pgfusepath{clip}%
\pgfsetbuttcap%
\pgfsetroundjoin%
\definecolor{currentfill}{rgb}{0.556863,0.654902,0.847059}%
\pgfsetfillcolor{currentfill}%
\pgfsetfillopacity{0.130000}%
\pgfsetlinewidth{0.000000pt}%
\definecolor{currentstroke}{rgb}{0.000000,0.000000,0.000000}%
\pgfsetstrokecolor{currentstroke}%
\pgfsetdash{}{0pt}%
\pgfpathmoveto{\pgfqpoint{1.416629in}{0.909791in}}%
\pgfpathlineto{\pgfqpoint{1.433692in}{0.919927in}}%
\pgfpathlineto{\pgfqpoint{1.450830in}{0.930373in}}%
\pgfpathlineto{\pgfqpoint{1.458531in}{0.940264in}}%
\pgfpathlineto{\pgfqpoint{1.466220in}{0.950535in}}%
\pgfpathlineto{\pgfqpoint{1.449103in}{0.940493in}}%
\pgfpathlineto{\pgfqpoint{1.432059in}{0.930787in}}%
\pgfpathlineto{\pgfqpoint{1.424350in}{0.920090in}}%
\pgfpathlineto{\pgfqpoint{1.416629in}{0.909791in}}%
\pgfpathclose%
\pgfusepath{fill}%
\end{pgfscope}%
\begin{pgfscope}%
\pgfpathrectangle{\pgfqpoint{0.407592in}{0.222904in}}{\pgfqpoint{1.411163in}{1.411163in}}%
\pgfusepath{clip}%
\pgfsetbuttcap%
\pgfsetroundjoin%
\definecolor{currentfill}{rgb}{0.556863,0.654902,0.847059}%
\pgfsetfillcolor{currentfill}%
\pgfsetfillopacity{0.130000}%
\pgfsetlinewidth{0.000000pt}%
\definecolor{currentstroke}{rgb}{0.000000,0.000000,0.000000}%
\pgfsetstrokecolor{currentstroke}%
\pgfsetdash{}{0pt}%
\pgfpathmoveto{\pgfqpoint{0.612650in}{1.041271in}}%
\pgfpathlineto{\pgfqpoint{0.628907in}{1.023350in}}%
\pgfpathlineto{\pgfqpoint{0.645172in}{1.005662in}}%
\pgfpathlineto{\pgfqpoint{0.654220in}{1.004017in}}%
\pgfpathlineto{\pgfqpoint{0.663224in}{1.002692in}}%
\pgfpathlineto{\pgfqpoint{0.647001in}{1.020741in}}%
\pgfpathlineto{\pgfqpoint{0.630786in}{1.039003in}}%
\pgfpathlineto{\pgfqpoint{0.621739in}{1.039982in}}%
\pgfpathlineto{\pgfqpoint{0.612650in}{1.041271in}}%
\pgfpathclose%
\pgfusepath{fill}%
\end{pgfscope}%
\begin{pgfscope}%
\pgfpathrectangle{\pgfqpoint{0.407592in}{0.222904in}}{\pgfqpoint{1.411163in}{1.411163in}}%
\pgfusepath{clip}%
\pgfsetbuttcap%
\pgfsetroundjoin%
\definecolor{currentfill}{rgb}{0.556863,0.654902,0.847059}%
\pgfsetfillcolor{currentfill}%
\pgfsetfillopacity{0.130000}%
\pgfsetlinewidth{0.000000pt}%
\definecolor{currentstroke}{rgb}{0.000000,0.000000,0.000000}%
\pgfsetstrokecolor{currentstroke}%
\pgfsetdash{}{0pt}%
\pgfpathmoveto{\pgfqpoint{0.958272in}{0.759259in}}%
\pgfpathlineto{\pgfqpoint{0.974593in}{0.751603in}}%
\pgfpathlineto{\pgfqpoint{0.990936in}{0.745394in}}%
\pgfpathlineto{\pgfqpoint{0.999340in}{0.741975in}}%
\pgfpathlineto{\pgfqpoint{1.007705in}{0.739069in}}%
\pgfpathlineto{\pgfqpoint{0.991407in}{0.745446in}}%
\pgfpathlineto{\pgfqpoint{0.975132in}{0.753395in}}%
\pgfpathlineto{\pgfqpoint{0.966722in}{0.756067in}}%
\pgfpathlineto{\pgfqpoint{0.958272in}{0.759259in}}%
\pgfpathclose%
\pgfusepath{fill}%
\end{pgfscope}%
\begin{pgfscope}%
\pgfpathrectangle{\pgfqpoint{0.407592in}{0.222904in}}{\pgfqpoint{1.411163in}{1.411163in}}%
\pgfusepath{clip}%
\pgfsetbuttcap%
\pgfsetroundjoin%
\definecolor{currentfill}{rgb}{0.556863,0.654902,0.847059}%
\pgfsetfillcolor{currentfill}%
\pgfsetfillopacity{0.130000}%
\pgfsetlinewidth{0.000000pt}%
\definecolor{currentstroke}{rgb}{0.000000,0.000000,0.000000}%
\pgfsetstrokecolor{currentstroke}%
\pgfsetdash{}{0pt}%
\pgfpathmoveto{\pgfqpoint{1.138826in}{0.732573in}}%
\pgfpathlineto{\pgfqpoint{1.155340in}{0.737677in}}%
\pgfpathlineto{\pgfqpoint{1.171915in}{0.743857in}}%
\pgfpathlineto{\pgfqpoint{1.179964in}{0.745784in}}%
\pgfpathlineto{\pgfqpoint{1.187987in}{0.748527in}}%
\pgfpathlineto{\pgfqpoint{1.171439in}{0.741996in}}%
\pgfpathlineto{\pgfqpoint{1.154954in}{0.736535in}}%
\pgfpathlineto{\pgfqpoint{1.146904in}{0.734104in}}%
\pgfpathlineto{\pgfqpoint{1.138826in}{0.732573in}}%
\pgfpathclose%
\pgfusepath{fill}%
\end{pgfscope}%
\begin{pgfscope}%
\pgfpathrectangle{\pgfqpoint{0.407592in}{0.222904in}}{\pgfqpoint{1.411163in}{1.411163in}}%
\pgfusepath{clip}%
\pgfsetbuttcap%
\pgfsetroundjoin%
\definecolor{currentfill}{rgb}{0.556863,0.654902,0.847059}%
\pgfsetfillcolor{currentfill}%
\pgfsetfillopacity{0.130000}%
\pgfsetlinewidth{0.000000pt}%
\definecolor{currentstroke}{rgb}{0.000000,0.000000,0.000000}%
\pgfsetstrokecolor{currentstroke}%
\pgfsetdash{}{0pt}%
\pgfpathmoveto{\pgfqpoint{0.760743in}{0.900507in}}%
\pgfpathlineto{\pgfqpoint{0.777027in}{0.884854in}}%
\pgfpathlineto{\pgfqpoint{0.793319in}{0.869728in}}%
\pgfpathlineto{\pgfqpoint{0.802111in}{0.866925in}}%
\pgfpathlineto{\pgfqpoint{0.810860in}{0.864500in}}%
\pgfpathlineto{\pgfqpoint{0.794611in}{0.880157in}}%
\pgfpathlineto{\pgfqpoint{0.778370in}{0.896319in}}%
\pgfpathlineto{\pgfqpoint{0.769577in}{0.898225in}}%
\pgfpathlineto{\pgfqpoint{0.760743in}{0.900507in}}%
\pgfpathclose%
\pgfusepath{fill}%
\end{pgfscope}%
\begin{pgfscope}%
\pgfpathrectangle{\pgfqpoint{0.407592in}{0.222904in}}{\pgfqpoint{1.411163in}{1.411163in}}%
\pgfusepath{clip}%
\pgfsetbuttcap%
\pgfsetroundjoin%
\definecolor{currentfill}{rgb}{0.556863,0.654902,0.847059}%
\pgfsetfillcolor{currentfill}%
\pgfsetfillopacity{0.130000}%
\pgfsetlinewidth{0.000000pt}%
\definecolor{currentstroke}{rgb}{0.000000,0.000000,0.000000}%
\pgfsetstrokecolor{currentstroke}%
\pgfsetdash{}{0pt}%
\pgfpathmoveto{\pgfqpoint{1.531233in}{1.014372in}}%
\pgfpathlineto{\pgfqpoint{1.548535in}{1.024487in}}%
\pgfpathlineto{\pgfqpoint{1.565912in}{1.034874in}}%
\pgfpathlineto{\pgfqpoint{1.573500in}{1.045948in}}%
\pgfpathlineto{\pgfqpoint{1.581076in}{1.057287in}}%
\pgfpathlineto{\pgfqpoint{1.563721in}{1.047288in}}%
\pgfpathlineto{\pgfqpoint{1.546440in}{1.037577in}}%
\pgfpathlineto{\pgfqpoint{1.538842in}{1.025840in}}%
\pgfpathlineto{\pgfqpoint{1.531233in}{1.014372in}}%
\pgfpathclose%
\pgfusepath{fill}%
\end{pgfscope}%
\begin{pgfscope}%
\pgfpathrectangle{\pgfqpoint{0.407592in}{0.222904in}}{\pgfqpoint{1.411163in}{1.411163in}}%
\pgfusepath{clip}%
\pgfsetbuttcap%
\pgfsetroundjoin%
\definecolor{currentfill}{rgb}{0.556863,0.654902,0.847059}%
\pgfsetfillcolor{currentfill}%
\pgfsetfillopacity{0.130000}%
\pgfsetlinewidth{0.000000pt}%
\definecolor{currentstroke}{rgb}{0.000000,0.000000,0.000000}%
\pgfsetstrokecolor{currentstroke}%
\pgfsetdash{}{0pt}%
\pgfpathmoveto{\pgfqpoint{1.089654in}{0.726594in}}%
\pgfpathlineto{\pgfqpoint{1.106092in}{0.728593in}}%
\pgfpathlineto{\pgfqpoint{1.122582in}{0.731981in}}%
\pgfpathlineto{\pgfqpoint{1.130719in}{0.731886in}}%
\pgfpathlineto{\pgfqpoint{1.138826in}{0.732573in}}%
\pgfpathlineto{\pgfqpoint{1.122370in}{0.728719in}}%
\pgfpathlineto{\pgfqpoint{1.105968in}{0.726299in}}%
\pgfpathlineto{\pgfqpoint{1.097827in}{0.726036in}}%
\pgfpathlineto{\pgfqpoint{1.089654in}{0.726594in}}%
\pgfpathclose%
\pgfusepath{fill}%
\end{pgfscope}%
\begin{pgfscope}%
\pgfpathrectangle{\pgfqpoint{0.407592in}{0.222904in}}{\pgfqpoint{1.411163in}{1.411163in}}%
\pgfusepath{clip}%
\pgfsetbuttcap%
\pgfsetroundjoin%
\definecolor{currentfill}{rgb}{0.556863,0.654902,0.847059}%
\pgfsetfillcolor{currentfill}%
\pgfsetfillopacity{0.130000}%
\pgfsetlinewidth{0.000000pt}%
\definecolor{currentstroke}{rgb}{0.000000,0.000000,0.000000}%
\pgfsetstrokecolor{currentstroke}%
\pgfsetdash{}{0pt}%
\pgfpathmoveto{\pgfqpoint{0.875957in}{0.808400in}}%
\pgfpathlineto{\pgfqpoint{0.892259in}{0.796432in}}%
\pgfpathlineto{\pgfqpoint{0.908575in}{0.785475in}}%
\pgfpathlineto{\pgfqpoint{0.917150in}{0.781820in}}%
\pgfpathlineto{\pgfqpoint{0.925686in}{0.778588in}}%
\pgfpathlineto{\pgfqpoint{0.909415in}{0.790053in}}%
\pgfpathlineto{\pgfqpoint{0.893158in}{0.802566in}}%
\pgfpathlineto{\pgfqpoint{0.884578in}{0.805265in}}%
\pgfpathlineto{\pgfqpoint{0.875957in}{0.808400in}}%
\pgfpathclose%
\pgfusepath{fill}%
\end{pgfscope}%
\begin{pgfscope}%
\pgfpathrectangle{\pgfqpoint{0.407592in}{0.222904in}}{\pgfqpoint{1.411163in}{1.411163in}}%
\pgfusepath{clip}%
\pgfsetbuttcap%
\pgfsetroundjoin%
\definecolor{currentfill}{rgb}{0.556863,0.654902,0.847059}%
\pgfsetfillcolor{currentfill}%
\pgfsetfillopacity{0.130000}%
\pgfsetlinewidth{0.000000pt}%
\definecolor{currentstroke}{rgb}{0.000000,0.000000,0.000000}%
\pgfsetstrokecolor{currentstroke}%
\pgfsetdash{}{0pt}%
\pgfpathmoveto{\pgfqpoint{1.187987in}{0.748527in}}%
\pgfpathlineto{\pgfqpoint{1.204599in}{0.755962in}}%
\pgfpathlineto{\pgfqpoint{1.221279in}{0.764160in}}%
\pgfpathlineto{\pgfqpoint{1.229254in}{0.767947in}}%
\pgfpathlineto{\pgfqpoint{1.237207in}{0.772505in}}%
\pgfpathlineto{\pgfqpoint{1.220551in}{0.764167in}}%
\pgfpathlineto{\pgfqpoint{1.203961in}{0.756580in}}%
\pgfpathlineto{\pgfqpoint{1.195986in}{0.752118in}}%
\pgfpathlineto{\pgfqpoint{1.187987in}{0.748527in}}%
\pgfpathclose%
\pgfusepath{fill}%
\end{pgfscope}%
\begin{pgfscope}%
\pgfpathrectangle{\pgfqpoint{0.407592in}{0.222904in}}{\pgfqpoint{1.411163in}{1.411163in}}%
\pgfusepath{clip}%
\pgfsetbuttcap%
\pgfsetroundjoin%
\definecolor{currentfill}{rgb}{0.556863,0.654902,0.847059}%
\pgfsetfillcolor{currentfill}%
\pgfsetfillopacity{0.130000}%
\pgfsetlinewidth{0.000000pt}%
\definecolor{currentstroke}{rgb}{0.000000,0.000000,0.000000}%
\pgfsetstrokecolor{currentstroke}%
\pgfsetdash{}{0pt}%
\pgfpathmoveto{\pgfqpoint{1.351630in}{0.852252in}}%
\pgfpathlineto{\pgfqpoint{1.368585in}{0.862462in}}%
\pgfpathlineto{\pgfqpoint{1.385613in}{0.873006in}}%
\pgfpathlineto{\pgfqpoint{1.393388in}{0.881500in}}%
\pgfpathlineto{\pgfqpoint{1.401148in}{0.890475in}}%
\pgfpathlineto{\pgfqpoint{1.384140in}{0.880275in}}%
\pgfpathlineto{\pgfqpoint{1.367204in}{0.870434in}}%
\pgfpathlineto{\pgfqpoint{1.359424in}{0.861083in}}%
\pgfpathlineto{\pgfqpoint{1.351630in}{0.852252in}}%
\pgfpathclose%
\pgfusepath{fill}%
\end{pgfscope}%
\begin{pgfscope}%
\pgfpathrectangle{\pgfqpoint{0.407592in}{0.222904in}}{\pgfqpoint{1.411163in}{1.411163in}}%
\pgfusepath{clip}%
\pgfsetbuttcap%
\pgfsetroundjoin%
\definecolor{currentfill}{rgb}{0.556863,0.654902,0.847059}%
\pgfsetfillcolor{currentfill}%
\pgfsetfillopacity{0.130000}%
\pgfsetlinewidth{0.000000pt}%
\definecolor{currentstroke}{rgb}{0.000000,0.000000,0.000000}%
\pgfsetstrokecolor{currentstroke}%
\pgfsetdash{}{0pt}%
\pgfpathmoveto{\pgfqpoint{0.645172in}{1.005662in}}%
\pgfpathlineto{\pgfqpoint{0.661447in}{0.988235in}}%
\pgfpathlineto{\pgfqpoint{0.677729in}{0.971095in}}%
\pgfpathlineto{\pgfqpoint{0.686734in}{0.969051in}}%
\pgfpathlineto{\pgfqpoint{0.695696in}{0.967335in}}%
\pgfpathlineto{\pgfqpoint{0.679456in}{0.984880in}}%
\pgfpathlineto{\pgfqpoint{0.663224in}{1.002692in}}%
\pgfpathlineto{\pgfqpoint{0.654220in}{1.004017in}}%
\pgfpathlineto{\pgfqpoint{0.645172in}{1.005662in}}%
\pgfpathclose%
\pgfusepath{fill}%
\end{pgfscope}%
\begin{pgfscope}%
\pgfpathrectangle{\pgfqpoint{0.407592in}{0.222904in}}{\pgfqpoint{1.411163in}{1.411163in}}%
\pgfusepath{clip}%
\pgfsetbuttcap%
\pgfsetroundjoin%
\definecolor{currentfill}{rgb}{0.556863,0.654902,0.847059}%
\pgfsetfillcolor{currentfill}%
\pgfsetfillopacity{0.130000}%
\pgfsetlinewidth{0.000000pt}%
\definecolor{currentstroke}{rgb}{0.000000,0.000000,0.000000}%
\pgfsetstrokecolor{currentstroke}%
\pgfsetdash{}{0pt}%
\pgfpathmoveto{\pgfqpoint{1.040384in}{0.731236in}}%
\pgfpathlineto{\pgfqpoint{1.056775in}{0.729769in}}%
\pgfpathlineto{\pgfqpoint{1.073207in}{0.729855in}}%
\pgfpathlineto{\pgfqpoint{1.081448in}{0.727893in}}%
\pgfpathlineto{\pgfqpoint{1.089654in}{0.726594in}}%
\pgfpathlineto{\pgfqpoint{1.073263in}{0.726135in}}%
\pgfpathlineto{\pgfqpoint{1.056914in}{0.727347in}}%
\pgfpathlineto{\pgfqpoint{1.048667in}{0.728959in}}%
\pgfpathlineto{\pgfqpoint{1.040384in}{0.731236in}}%
\pgfpathclose%
\pgfusepath{fill}%
\end{pgfscope}%
\begin{pgfscope}%
\pgfpathrectangle{\pgfqpoint{0.407592in}{0.222904in}}{\pgfqpoint{1.411163in}{1.411163in}}%
\pgfusepath{clip}%
\pgfsetbuttcap%
\pgfsetroundjoin%
\definecolor{currentfill}{rgb}{0.556863,0.654902,0.847059}%
\pgfsetfillcolor{currentfill}%
\pgfsetfillopacity{0.130000}%
\pgfsetlinewidth{0.000000pt}%
\definecolor{currentstroke}{rgb}{0.000000,0.000000,0.000000}%
\pgfsetstrokecolor{currentstroke}%
\pgfsetdash{}{0pt}%
\pgfpathmoveto{\pgfqpoint{0.496721in}{1.154304in}}%
\pgfpathlineto{\pgfqpoint{0.512970in}{1.135632in}}%
\pgfpathlineto{\pgfqpoint{0.529227in}{1.117104in}}%
\pgfpathlineto{\pgfqpoint{0.538488in}{1.115854in}}%
\pgfpathlineto{\pgfqpoint{0.547706in}{1.114878in}}%
\pgfpathlineto{\pgfqpoint{0.531491in}{1.133673in}}%
\pgfpathlineto{\pgfqpoint{0.515285in}{1.152597in}}%
\pgfpathlineto{\pgfqpoint{0.506025in}{1.153318in}}%
\pgfpathlineto{\pgfqpoint{0.496721in}{1.154304in}}%
\pgfpathclose%
\pgfusepath{fill}%
\end{pgfscope}%
\begin{pgfscope}%
\pgfpathrectangle{\pgfqpoint{0.407592in}{0.222904in}}{\pgfqpoint{1.411163in}{1.411163in}}%
\pgfusepath{clip}%
\pgfsetbuttcap%
\pgfsetroundjoin%
\definecolor{currentfill}{rgb}{0.556863,0.654902,0.847059}%
\pgfsetfillcolor{currentfill}%
\pgfsetfillopacity{0.130000}%
\pgfsetlinewidth{0.000000pt}%
\definecolor{currentstroke}{rgb}{0.000000,0.000000,0.000000}%
\pgfsetstrokecolor{currentstroke}%
\pgfsetdash{}{0pt}%
\pgfpathmoveto{\pgfqpoint{1.711130in}{1.192858in}}%
\pgfpathlineto{\pgfqpoint{1.728779in}{1.202720in}}%
\pgfpathlineto{\pgfqpoint{1.746506in}{1.212824in}}%
\pgfpathlineto{\pgfqpoint{1.753918in}{1.224984in}}%
\pgfpathlineto{\pgfqpoint{1.736203in}{1.215042in}}%
\pgfpathlineto{\pgfqpoint{1.718564in}{1.205346in}}%
\pgfusepath{fill}%
\end{pgfscope}%
\begin{pgfscope}%
\pgfpathrectangle{\pgfqpoint{0.407592in}{0.222904in}}{\pgfqpoint{1.411163in}{1.411163in}}%
\pgfusepath{clip}%
\pgfsetbuttcap%
\pgfsetroundjoin%
\definecolor{currentfill}{rgb}{0.556863,0.654902,0.847059}%
\pgfsetfillcolor{currentfill}%
\pgfsetfillopacity{0.130000}%
\pgfsetlinewidth{0.000000pt}%
\definecolor{currentstroke}{rgb}{0.000000,0.000000,0.000000}%
\pgfsetstrokecolor{currentstroke}%
\pgfsetdash{}{0pt}%
\pgfpathmoveto{\pgfqpoint{1.466220in}{0.950535in}}%
\pgfpathlineto{\pgfqpoint{1.483412in}{0.960881in}}%
\pgfpathlineto{\pgfqpoint{1.500680in}{0.971504in}}%
\pgfpathlineto{\pgfqpoint{1.508336in}{0.981741in}}%
\pgfpathlineto{\pgfqpoint{1.515980in}{0.992308in}}%
\pgfpathlineto{\pgfqpoint{1.498733in}{0.982075in}}%
\pgfpathlineto{\pgfqpoint{1.481561in}{0.972142in}}%
\pgfpathlineto{\pgfqpoint{1.473897in}{0.961167in}}%
\pgfpathlineto{\pgfqpoint{1.466220in}{0.950535in}}%
\pgfpathclose%
\pgfusepath{fill}%
\end{pgfscope}%
\begin{pgfscope}%
\pgfpathrectangle{\pgfqpoint{0.407592in}{0.222904in}}{\pgfqpoint{1.411163in}{1.411163in}}%
\pgfusepath{clip}%
\pgfsetbuttcap%
\pgfsetroundjoin%
\definecolor{currentfill}{rgb}{0.556863,0.654902,0.847059}%
\pgfsetfillcolor{currentfill}%
\pgfsetfillopacity{0.130000}%
\pgfsetlinewidth{0.000000pt}%
\definecolor{currentstroke}{rgb}{0.000000,0.000000,0.000000}%
\pgfsetstrokecolor{currentstroke}%
\pgfsetdash{}{0pt}%
\pgfpathmoveto{\pgfqpoint{1.646125in}{1.124345in}}%
\pgfpathlineto{\pgfqpoint{1.663666in}{1.134395in}}%
\pgfpathlineto{\pgfqpoint{1.681284in}{1.144691in}}%
\pgfpathlineto{\pgfqpoint{1.688761in}{1.156448in}}%
\pgfpathlineto{\pgfqpoint{1.696228in}{1.168401in}}%
\pgfpathlineto{\pgfqpoint{1.678633in}{1.158449in}}%
\pgfpathlineto{\pgfqpoint{1.661114in}{1.148756in}}%
\pgfpathlineto{\pgfqpoint{1.653625in}{1.136453in}}%
\pgfpathlineto{\pgfqpoint{1.646125in}{1.124345in}}%
\pgfpathclose%
\pgfusepath{fill}%
\end{pgfscope}%
\begin{pgfscope}%
\pgfpathrectangle{\pgfqpoint{0.407592in}{0.222904in}}{\pgfqpoint{1.411163in}{1.411163in}}%
\pgfusepath{clip}%
\pgfsetbuttcap%
\pgfsetroundjoin%
\definecolor{currentfill}{rgb}{0.556863,0.654902,0.847059}%
\pgfsetfillcolor{currentfill}%
\pgfsetfillopacity{0.130000}%
\pgfsetlinewidth{0.000000pt}%
\definecolor{currentstroke}{rgb}{0.000000,0.000000,0.000000}%
\pgfsetstrokecolor{currentstroke}%
\pgfsetdash{}{0pt}%
\pgfpathmoveto{\pgfqpoint{1.237207in}{0.772505in}}%
\pgfpathlineto{\pgfqpoint{1.253933in}{0.781477in}}%
\pgfpathlineto{\pgfqpoint{1.270730in}{0.790984in}}%
\pgfpathlineto{\pgfqpoint{1.278642in}{0.796313in}}%
\pgfpathlineto{\pgfqpoint{1.286535in}{0.802331in}}%
\pgfpathlineto{\pgfqpoint{1.269759in}{0.792873in}}%
\pgfpathlineto{\pgfqpoint{1.253054in}{0.783954in}}%
\pgfpathlineto{\pgfqpoint{1.245140in}{0.777842in}}%
\pgfpathlineto{\pgfqpoint{1.237207in}{0.772505in}}%
\pgfpathclose%
\pgfusepath{fill}%
\end{pgfscope}%
\begin{pgfscope}%
\pgfpathrectangle{\pgfqpoint{0.407592in}{0.222904in}}{\pgfqpoint{1.411163in}{1.411163in}}%
\pgfusepath{clip}%
\pgfsetbuttcap%
\pgfsetroundjoin%
\definecolor{currentfill}{rgb}{0.556863,0.654902,0.847059}%
\pgfsetfillcolor{currentfill}%
\pgfsetfillopacity{0.130000}%
\pgfsetlinewidth{0.000000pt}%
\definecolor{currentstroke}{rgb}{0.000000,0.000000,0.000000}%
\pgfsetstrokecolor{currentstroke}%
\pgfsetdash{}{0pt}%
\pgfpathmoveto{\pgfqpoint{0.793319in}{0.869728in}}%
\pgfpathlineto{\pgfqpoint{0.809622in}{0.855190in}}%
\pgfpathlineto{\pgfqpoint{0.825934in}{0.841309in}}%
\pgfpathlineto{\pgfqpoint{0.834682in}{0.837961in}}%
\pgfpathlineto{\pgfqpoint{0.843388in}{0.834985in}}%
\pgfpathlineto{\pgfqpoint{0.827120in}{0.849418in}}%
\pgfpathlineto{\pgfqpoint{0.810860in}{0.864500in}}%
\pgfpathlineto{\pgfqpoint{0.802111in}{0.866925in}}%
\pgfpathlineto{\pgfqpoint{0.793319in}{0.869728in}}%
\pgfpathclose%
\pgfusepath{fill}%
\end{pgfscope}%
\begin{pgfscope}%
\pgfpathrectangle{\pgfqpoint{0.407592in}{0.222904in}}{\pgfqpoint{1.411163in}{1.411163in}}%
\pgfusepath{clip}%
\pgfsetbuttcap%
\pgfsetroundjoin%
\definecolor{currentfill}{rgb}{0.556863,0.654902,0.847059}%
\pgfsetfillcolor{currentfill}%
\pgfsetfillopacity{0.130000}%
\pgfsetlinewidth{0.000000pt}%
\definecolor{currentstroke}{rgb}{0.000000,0.000000,0.000000}%
\pgfsetstrokecolor{currentstroke}%
\pgfsetdash{}{0pt}%
\pgfpathmoveto{\pgfqpoint{1.581076in}{1.057287in}}%
\pgfpathlineto{\pgfqpoint{1.598508in}{1.067556in}}%
\pgfpathlineto{\pgfqpoint{1.616016in}{1.078076in}}%
\pgfpathlineto{\pgfqpoint{1.623560in}{1.089298in}}%
\pgfpathlineto{\pgfqpoint{1.631093in}{1.100757in}}%
\pgfpathlineto{\pgfqpoint{1.613607in}{1.090602in}}%
\pgfpathlineto{\pgfqpoint{1.596196in}{1.080712in}}%
\pgfpathlineto{\pgfqpoint{1.588642in}{1.068879in}}%
\pgfpathlineto{\pgfqpoint{1.581076in}{1.057287in}}%
\pgfpathclose%
\pgfusepath{fill}%
\end{pgfscope}%
\begin{pgfscope}%
\pgfpathrectangle{\pgfqpoint{0.407592in}{0.222904in}}{\pgfqpoint{1.411163in}{1.411163in}}%
\pgfusepath{clip}%
\pgfsetbuttcap%
\pgfsetroundjoin%
\definecolor{currentfill}{rgb}{0.556863,0.654902,0.847059}%
\pgfsetfillcolor{currentfill}%
\pgfsetfillopacity{0.130000}%
\pgfsetlinewidth{0.000000pt}%
\definecolor{currentstroke}{rgb}{0.000000,0.000000,0.000000}%
\pgfsetstrokecolor{currentstroke}%
\pgfsetdash{}{0pt}%
\pgfpathmoveto{\pgfqpoint{0.677729in}{0.971095in}}%
\pgfpathlineto{\pgfqpoint{0.694020in}{0.954275in}}%
\pgfpathlineto{\pgfqpoint{0.710320in}{0.937811in}}%
\pgfpathlineto{\pgfqpoint{0.719283in}{0.935325in}}%
\pgfpathlineto{\pgfqpoint{0.728202in}{0.933171in}}%
\pgfpathlineto{\pgfqpoint{0.711945in}{0.950086in}}%
\pgfpathlineto{\pgfqpoint{0.695696in}{0.967335in}}%
\pgfpathlineto{\pgfqpoint{0.686734in}{0.969051in}}%
\pgfpathlineto{\pgfqpoint{0.677729in}{0.971095in}}%
\pgfpathclose%
\pgfusepath{fill}%
\end{pgfscope}%
\begin{pgfscope}%
\pgfpathrectangle{\pgfqpoint{0.407592in}{0.222904in}}{\pgfqpoint{1.411163in}{1.411163in}}%
\pgfusepath{clip}%
\pgfsetbuttcap%
\pgfsetroundjoin%
\definecolor{currentfill}{rgb}{0.556863,0.654902,0.847059}%
\pgfsetfillcolor{currentfill}%
\pgfsetfillopacity{0.130000}%
\pgfsetlinewidth{0.000000pt}%
\definecolor{currentstroke}{rgb}{0.000000,0.000000,0.000000}%
\pgfsetstrokecolor{currentstroke}%
\pgfsetdash{}{0pt}%
\pgfpathmoveto{\pgfqpoint{0.529227in}{1.117104in}}%
\pgfpathlineto{\pgfqpoint{0.545492in}{1.098734in}}%
\pgfpathlineto{\pgfqpoint{0.561766in}{1.080538in}}%
\pgfpathlineto{\pgfqpoint{0.570985in}{1.078992in}}%
\pgfpathlineto{\pgfqpoint{0.580161in}{1.077728in}}%
\pgfpathlineto{\pgfqpoint{0.563929in}{1.096224in}}%
\pgfpathlineto{\pgfqpoint{0.547706in}{1.114878in}}%
\pgfpathlineto{\pgfqpoint{0.538488in}{1.115854in}}%
\pgfpathlineto{\pgfqpoint{0.529227in}{1.117104in}}%
\pgfpathclose%
\pgfusepath{fill}%
\end{pgfscope}%
\begin{pgfscope}%
\pgfpathrectangle{\pgfqpoint{0.407592in}{0.222904in}}{\pgfqpoint{1.411163in}{1.411163in}}%
\pgfusepath{clip}%
\pgfsetbuttcap%
\pgfsetroundjoin%
\definecolor{currentfill}{rgb}{0.556863,0.654902,0.847059}%
\pgfsetfillcolor{currentfill}%
\pgfsetfillopacity{0.130000}%
\pgfsetlinewidth{0.000000pt}%
\definecolor{currentstroke}{rgb}{0.000000,0.000000,0.000000}%
\pgfsetstrokecolor{currentstroke}%
\pgfsetdash{}{0pt}%
\pgfpathmoveto{\pgfqpoint{0.908575in}{0.785475in}}%
\pgfpathlineto{\pgfqpoint{0.924906in}{0.775624in}}%
\pgfpathlineto{\pgfqpoint{0.941254in}{0.766967in}}%
\pgfpathlineto{\pgfqpoint{0.949783in}{0.762910in}}%
\pgfpathlineto{\pgfqpoint{0.958272in}{0.759259in}}%
\pgfpathlineto{\pgfqpoint{0.941970in}{0.768287in}}%
\pgfpathlineto{\pgfqpoint{0.925686in}{0.778588in}}%
\pgfpathlineto{\pgfqpoint{0.917150in}{0.781820in}}%
\pgfpathlineto{\pgfqpoint{0.908575in}{0.785475in}}%
\pgfpathclose%
\pgfusepath{fill}%
\end{pgfscope}%
\begin{pgfscope}%
\pgfpathrectangle{\pgfqpoint{0.407592in}{0.222904in}}{\pgfqpoint{1.411163in}{1.411163in}}%
\pgfusepath{clip}%
\pgfsetbuttcap%
\pgfsetroundjoin%
\definecolor{currentfill}{rgb}{0.556863,0.654902,0.847059}%
\pgfsetfillcolor{currentfill}%
\pgfsetfillopacity{0.130000}%
\pgfsetlinewidth{0.000000pt}%
\definecolor{currentstroke}{rgb}{0.000000,0.000000,0.000000}%
\pgfsetstrokecolor{currentstroke}%
\pgfsetdash{}{0pt}%
\pgfpathmoveto{\pgfqpoint{1.401148in}{0.890475in}}%
\pgfpathlineto{\pgfqpoint{1.418231in}{0.900994in}}%
\pgfpathlineto{\pgfqpoint{1.435389in}{0.911799in}}%
\pgfpathlineto{\pgfqpoint{1.443116in}{0.920879in}}%
\pgfpathlineto{\pgfqpoint{1.450830in}{0.930373in}}%
\pgfpathlineto{\pgfqpoint{1.433692in}{0.919927in}}%
\pgfpathlineto{\pgfqpoint{1.416629in}{0.909791in}}%
\pgfpathlineto{\pgfqpoint{1.408895in}{0.899912in}}%
\pgfpathlineto{\pgfqpoint{1.401148in}{0.890475in}}%
\pgfpathclose%
\pgfusepath{fill}%
\end{pgfscope}%
\begin{pgfscope}%
\pgfpathrectangle{\pgfqpoint{0.407592in}{0.222904in}}{\pgfqpoint{1.411163in}{1.411163in}}%
\pgfusepath{clip}%
\pgfsetbuttcap%
\pgfsetroundjoin%
\definecolor{currentfill}{rgb}{0.556863,0.654902,0.847059}%
\pgfsetfillcolor{currentfill}%
\pgfsetfillopacity{0.130000}%
\pgfsetlinewidth{0.000000pt}%
\definecolor{currentstroke}{rgb}{0.000000,0.000000,0.000000}%
\pgfsetstrokecolor{currentstroke}%
\pgfsetdash{}{0pt}%
\pgfpathmoveto{\pgfqpoint{1.286535in}{0.802331in}}%
\pgfpathlineto{\pgfqpoint{1.303383in}{0.812248in}}%
\pgfpathlineto{\pgfqpoint{1.320303in}{0.822560in}}%
\pgfpathlineto{\pgfqpoint{1.328159in}{0.829102in}}%
\pgfpathlineto{\pgfqpoint{1.335998in}{0.836242in}}%
\pgfpathlineto{\pgfqpoint{1.319098in}{0.826110in}}%
\pgfpathlineto{\pgfqpoint{1.302270in}{0.816388in}}%
\pgfpathlineto{\pgfqpoint{1.294411in}{0.809027in}}%
\pgfpathlineto{\pgfqpoint{1.286535in}{0.802331in}}%
\pgfpathclose%
\pgfusepath{fill}%
\end{pgfscope}%
\begin{pgfscope}%
\pgfpathrectangle{\pgfqpoint{0.407592in}{0.222904in}}{\pgfqpoint{1.411163in}{1.411163in}}%
\pgfusepath{clip}%
\pgfsetbuttcap%
\pgfsetroundjoin%
\definecolor{currentfill}{rgb}{0.556863,0.654902,0.847059}%
\pgfsetfillcolor{currentfill}%
\pgfsetfillopacity{0.130000}%
\pgfsetlinewidth{0.000000pt}%
\definecolor{currentstroke}{rgb}{0.000000,0.000000,0.000000}%
\pgfsetstrokecolor{currentstroke}%
\pgfsetdash{}{0pt}%
\pgfpathmoveto{\pgfqpoint{0.990936in}{0.745394in}}%
\pgfpathlineto{\pgfqpoint{1.007307in}{0.740681in}}%
\pgfpathlineto{\pgfqpoint{1.023708in}{0.737478in}}%
\pgfpathlineto{\pgfqpoint{1.032064in}{0.734099in}}%
\pgfpathlineto{\pgfqpoint{1.040384in}{0.731236in}}%
\pgfpathlineto{\pgfqpoint{1.024029in}{0.734327in}}%
\pgfpathlineto{\pgfqpoint{1.007705in}{0.739069in}}%
\pgfpathlineto{\pgfqpoint{0.999340in}{0.741975in}}%
\pgfpathlineto{\pgfqpoint{0.990936in}{0.745394in}}%
\pgfpathclose%
\pgfusepath{fill}%
\end{pgfscope}%
\begin{pgfscope}%
\pgfpathrectangle{\pgfqpoint{0.407592in}{0.222904in}}{\pgfqpoint{1.411163in}{1.411163in}}%
\pgfusepath{clip}%
\pgfsetbuttcap%
\pgfsetroundjoin%
\definecolor{currentfill}{rgb}{0.556863,0.654902,0.847059}%
\pgfsetfillcolor{currentfill}%
\pgfsetfillopacity{0.130000}%
\pgfsetlinewidth{0.000000pt}%
\definecolor{currentstroke}{rgb}{0.000000,0.000000,0.000000}%
\pgfsetstrokecolor{currentstroke}%
\pgfsetdash{}{0pt}%
\pgfpathmoveto{\pgfqpoint{1.515980in}{0.992308in}}%
\pgfpathlineto{\pgfqpoint{1.533302in}{1.002815in}}%
\pgfpathlineto{\pgfqpoint{1.550701in}{1.013574in}}%
\pgfpathlineto{\pgfqpoint{1.558312in}{1.024079in}}%
\pgfpathlineto{\pgfqpoint{1.565912in}{1.034874in}}%
\pgfpathlineto{\pgfqpoint{1.548535in}{1.024487in}}%
\pgfpathlineto{\pgfqpoint{1.531233in}{1.014372in}}%
\pgfpathlineto{\pgfqpoint{1.523612in}{1.003190in}}%
\pgfpathlineto{\pgfqpoint{1.515980in}{0.992308in}}%
\pgfpathclose%
\pgfusepath{fill}%
\end{pgfscope}%
\begin{pgfscope}%
\pgfpathrectangle{\pgfqpoint{0.407592in}{0.222904in}}{\pgfqpoint{1.411163in}{1.411163in}}%
\pgfusepath{clip}%
\pgfsetbuttcap%
\pgfsetroundjoin%
\definecolor{currentfill}{rgb}{0.556863,0.654902,0.847059}%
\pgfsetfillcolor{currentfill}%
\pgfsetfillopacity{0.130000}%
\pgfsetlinewidth{0.000000pt}%
\definecolor{currentstroke}{rgb}{0.000000,0.000000,0.000000}%
\pgfsetstrokecolor{currentstroke}%
\pgfsetdash{}{0pt}%
\pgfpathmoveto{\pgfqpoint{1.696228in}{1.168401in}}%
\pgfpathlineto{\pgfqpoint{1.713901in}{1.178597in}}%
\pgfpathlineto{\pgfqpoint{1.731650in}{1.189024in}}%
\pgfpathlineto{\pgfqpoint{1.739084in}{1.200835in}}%
\pgfpathlineto{\pgfqpoint{1.746506in}{1.212824in}}%
\pgfpathlineto{\pgfqpoint{1.728779in}{1.202720in}}%
\pgfpathlineto{\pgfqpoint{1.711130in}{1.192858in}}%
\pgfpathlineto{\pgfqpoint{1.703684in}{1.180540in}}%
\pgfpathlineto{\pgfqpoint{1.696228in}{1.168401in}}%
\pgfpathclose%
\pgfusepath{fill}%
\end{pgfscope}%
\begin{pgfscope}%
\pgfpathrectangle{\pgfqpoint{0.407592in}{0.222904in}}{\pgfqpoint{1.411163in}{1.411163in}}%
\pgfusepath{clip}%
\pgfsetbuttcap%
\pgfsetroundjoin%
\definecolor{currentfill}{rgb}{0.556863,0.654902,0.847059}%
\pgfsetfillcolor{currentfill}%
\pgfsetfillopacity{0.130000}%
\pgfsetlinewidth{0.000000pt}%
\definecolor{currentstroke}{rgb}{0.000000,0.000000,0.000000}%
\pgfsetstrokecolor{currentstroke}%
\pgfsetdash{}{0pt}%
\pgfpathmoveto{\pgfqpoint{0.561766in}{1.080538in}}%
\pgfpathlineto{\pgfqpoint{0.578049in}{1.062533in}}%
\pgfpathlineto{\pgfqpoint{0.594340in}{1.044738in}}%
\pgfpathlineto{\pgfqpoint{0.603517in}{1.042860in}}%
\pgfpathlineto{\pgfqpoint{0.612650in}{1.041271in}}%
\pgfpathlineto{\pgfqpoint{0.596401in}{1.059404in}}%
\pgfpathlineto{\pgfqpoint{0.580161in}{1.077728in}}%
\pgfpathlineto{\pgfqpoint{0.570985in}{1.078992in}}%
\pgfpathlineto{\pgfqpoint{0.561766in}{1.080538in}}%
\pgfpathclose%
\pgfusepath{fill}%
\end{pgfscope}%
\begin{pgfscope}%
\pgfpathrectangle{\pgfqpoint{0.407592in}{0.222904in}}{\pgfqpoint{1.411163in}{1.411163in}}%
\pgfusepath{clip}%
\pgfsetbuttcap%
\pgfsetroundjoin%
\definecolor{currentfill}{rgb}{0.556863,0.654902,0.847059}%
\pgfsetfillcolor{currentfill}%
\pgfsetfillopacity{0.130000}%
\pgfsetlinewidth{0.000000pt}%
\definecolor{currentstroke}{rgb}{0.000000,0.000000,0.000000}%
\pgfsetstrokecolor{currentstroke}%
\pgfsetdash{}{0pt}%
\pgfpathmoveto{\pgfqpoint{0.825934in}{0.841309in}}%
\pgfpathlineto{\pgfqpoint{0.842258in}{0.828159in}}%
\pgfpathlineto{\pgfqpoint{0.858593in}{0.815816in}}%
\pgfpathlineto{\pgfqpoint{0.867296in}{0.811929in}}%
\pgfpathlineto{\pgfqpoint{0.875957in}{0.808400in}}%
\pgfpathlineto{\pgfqpoint{0.859667in}{0.821282in}}%
\pgfpathlineto{\pgfqpoint{0.843388in}{0.834985in}}%
\pgfpathlineto{\pgfqpoint{0.834682in}{0.837961in}}%
\pgfpathlineto{\pgfqpoint{0.825934in}{0.841309in}}%
\pgfpathclose%
\pgfusepath{fill}%
\end{pgfscope}%
\begin{pgfscope}%
\pgfpathrectangle{\pgfqpoint{0.407592in}{0.222904in}}{\pgfqpoint{1.411163in}{1.411163in}}%
\pgfusepath{clip}%
\pgfsetbuttcap%
\pgfsetroundjoin%
\definecolor{currentfill}{rgb}{0.556863,0.654902,0.847059}%
\pgfsetfillcolor{currentfill}%
\pgfsetfillopacity{0.130000}%
\pgfsetlinewidth{0.000000pt}%
\definecolor{currentstroke}{rgb}{0.000000,0.000000,0.000000}%
\pgfsetstrokecolor{currentstroke}%
\pgfsetdash{}{0pt}%
\pgfpathmoveto{\pgfqpoint{0.710320in}{0.937811in}}%
\pgfpathlineto{\pgfqpoint{0.726629in}{0.921741in}}%
\pgfpathlineto{\pgfqpoint{0.742947in}{0.906110in}}%
\pgfpathlineto{\pgfqpoint{0.751866in}{0.903143in}}%
\pgfpathlineto{\pgfqpoint{0.760743in}{0.900507in}}%
\pgfpathlineto{\pgfqpoint{0.744468in}{0.916630in}}%
\pgfpathlineto{\pgfqpoint{0.728202in}{0.933171in}}%
\pgfpathlineto{\pgfqpoint{0.719283in}{0.935325in}}%
\pgfpathlineto{\pgfqpoint{0.710320in}{0.937811in}}%
\pgfpathclose%
\pgfusepath{fill}%
\end{pgfscope}%
\begin{pgfscope}%
\pgfpathrectangle{\pgfqpoint{0.407592in}{0.222904in}}{\pgfqpoint{1.411163in}{1.411163in}}%
\pgfusepath{clip}%
\pgfsetbuttcap%
\pgfsetroundjoin%
\definecolor{currentfill}{rgb}{0.556863,0.654902,0.847059}%
\pgfsetfillcolor{currentfill}%
\pgfsetfillopacity{0.130000}%
\pgfsetlinewidth{0.000000pt}%
\definecolor{currentstroke}{rgb}{0.000000,0.000000,0.000000}%
\pgfsetstrokecolor{currentstroke}%
\pgfsetdash{}{0pt}%
\pgfpathmoveto{\pgfqpoint{1.631093in}{1.100757in}}%
\pgfpathlineto{\pgfqpoint{1.648656in}{1.111162in}}%
\pgfpathlineto{\pgfqpoint{1.666296in}{1.121801in}}%
\pgfpathlineto{\pgfqpoint{1.673795in}{1.133139in}}%
\pgfpathlineto{\pgfqpoint{1.681284in}{1.144691in}}%
\pgfpathlineto{\pgfqpoint{1.663666in}{1.134395in}}%
\pgfpathlineto{\pgfqpoint{1.646125in}{1.124345in}}%
\pgfpathlineto{\pgfqpoint{1.638615in}{1.112443in}}%
\pgfpathlineto{\pgfqpoint{1.631093in}{1.100757in}}%
\pgfpathclose%
\pgfusepath{fill}%
\end{pgfscope}%
\begin{pgfscope}%
\pgfpathrectangle{\pgfqpoint{0.407592in}{0.222904in}}{\pgfqpoint{1.411163in}{1.411163in}}%
\pgfusepath{clip}%
\pgfsetbuttcap%
\pgfsetroundjoin%
\definecolor{currentfill}{rgb}{0.556863,0.654902,0.847059}%
\pgfsetfillcolor{currentfill}%
\pgfsetfillopacity{0.130000}%
\pgfsetlinewidth{0.000000pt}%
\definecolor{currentstroke}{rgb}{0.000000,0.000000,0.000000}%
\pgfsetstrokecolor{currentstroke}%
\pgfsetdash{}{0pt}%
\pgfpathmoveto{\pgfqpoint{1.335998in}{0.836242in}}%
\pgfpathlineto{\pgfqpoint{1.352972in}{0.846730in}}%
\pgfpathlineto{\pgfqpoint{1.370020in}{0.857532in}}%
\pgfpathlineto{\pgfqpoint{1.377824in}{0.865011in}}%
\pgfpathlineto{\pgfqpoint{1.385613in}{0.873006in}}%
\pgfpathlineto{\pgfqpoint{1.368585in}{0.862462in}}%
\pgfpathlineto{\pgfqpoint{1.351630in}{0.852252in}}%
\pgfpathlineto{\pgfqpoint{1.343822in}{0.843965in}}%
\pgfpathlineto{\pgfqpoint{1.335998in}{0.836242in}}%
\pgfpathclose%
\pgfusepath{fill}%
\end{pgfscope}%
\begin{pgfscope}%
\pgfpathrectangle{\pgfqpoint{0.407592in}{0.222904in}}{\pgfqpoint{1.411163in}{1.411163in}}%
\pgfusepath{clip}%
\pgfsetbuttcap%
\pgfsetroundjoin%
\definecolor{currentfill}{rgb}{0.556863,0.654902,0.847059}%
\pgfsetfillcolor{currentfill}%
\pgfsetfillopacity{0.130000}%
\pgfsetlinewidth{0.000000pt}%
\definecolor{currentstroke}{rgb}{0.000000,0.000000,0.000000}%
\pgfsetstrokecolor{currentstroke}%
\pgfsetdash{}{0pt}%
\pgfpathmoveto{\pgfqpoint{1.450830in}{0.930373in}}%
\pgfpathlineto{\pgfqpoint{1.468042in}{0.941100in}}%
\pgfpathlineto{\pgfqpoint{1.485330in}{0.952082in}}%
\pgfpathlineto{\pgfqpoint{1.493011in}{0.961613in}}%
\pgfpathlineto{\pgfqpoint{1.500680in}{0.971504in}}%
\pgfpathlineto{\pgfqpoint{1.483412in}{0.960881in}}%
\pgfpathlineto{\pgfqpoint{1.466220in}{0.950535in}}%
\pgfpathlineto{\pgfqpoint{1.458531in}{0.940264in}}%
\pgfpathlineto{\pgfqpoint{1.450830in}{0.930373in}}%
\pgfpathclose%
\pgfusepath{fill}%
\end{pgfscope}%
\begin{pgfscope}%
\pgfpathrectangle{\pgfqpoint{0.407592in}{0.222904in}}{\pgfqpoint{1.411163in}{1.411163in}}%
\pgfusepath{clip}%
\pgfsetbuttcap%
\pgfsetroundjoin%
\definecolor{currentfill}{rgb}{0.556863,0.654902,0.847059}%
\pgfsetfillcolor{currentfill}%
\pgfsetfillopacity{0.130000}%
\pgfsetlinewidth{0.000000pt}%
\definecolor{currentstroke}{rgb}{0.000000,0.000000,0.000000}%
\pgfsetstrokecolor{currentstroke}%
\pgfsetdash{}{0pt}%
\pgfpathmoveto{\pgfqpoint{1.122582in}{0.731981in}}%
\pgfpathlineto{\pgfqpoint{1.139128in}{0.736598in}}%
\pgfpathlineto{\pgfqpoint{1.155734in}{0.742290in}}%
\pgfpathlineto{\pgfqpoint{1.163838in}{0.742708in}}%
\pgfpathlineto{\pgfqpoint{1.171915in}{0.743857in}}%
\pgfpathlineto{\pgfqpoint{1.155340in}{0.737677in}}%
\pgfpathlineto{\pgfqpoint{1.138826in}{0.732573in}}%
\pgfpathlineto{\pgfqpoint{1.130719in}{0.731886in}}%
\pgfpathlineto{\pgfqpoint{1.122582in}{0.731981in}}%
\pgfpathclose%
\pgfusepath{fill}%
\end{pgfscope}%
\begin{pgfscope}%
\pgfpathrectangle{\pgfqpoint{0.407592in}{0.222904in}}{\pgfqpoint{1.411163in}{1.411163in}}%
\pgfusepath{clip}%
\pgfsetbuttcap%
\pgfsetroundjoin%
\definecolor{currentfill}{rgb}{0.556863,0.654902,0.847059}%
\pgfsetfillcolor{currentfill}%
\pgfsetfillopacity{0.130000}%
\pgfsetlinewidth{0.000000pt}%
\definecolor{currentstroke}{rgb}{0.000000,0.000000,0.000000}%
\pgfsetstrokecolor{currentstroke}%
\pgfsetdash{}{0pt}%
\pgfpathmoveto{\pgfqpoint{0.594340in}{1.044738in}}%
\pgfpathlineto{\pgfqpoint{0.610639in}{1.027172in}}%
\pgfpathlineto{\pgfqpoint{0.626947in}{1.009861in}}%
\pgfpathlineto{\pgfqpoint{0.636082in}{1.007614in}}%
\pgfpathlineto{\pgfqpoint{0.645172in}{1.005662in}}%
\pgfpathlineto{\pgfqpoint{0.628907in}{1.023350in}}%
\pgfpathlineto{\pgfqpoint{0.612650in}{1.041271in}}%
\pgfpathlineto{\pgfqpoint{0.603517in}{1.042860in}}%
\pgfpathlineto{\pgfqpoint{0.594340in}{1.044738in}}%
\pgfpathclose%
\pgfusepath{fill}%
\end{pgfscope}%
\begin{pgfscope}%
\pgfpathrectangle{\pgfqpoint{0.407592in}{0.222904in}}{\pgfqpoint{1.411163in}{1.411163in}}%
\pgfusepath{clip}%
\pgfsetbuttcap%
\pgfsetroundjoin%
\definecolor{currentfill}{rgb}{0.556863,0.654902,0.847059}%
\pgfsetfillcolor{currentfill}%
\pgfsetfillopacity{0.130000}%
\pgfsetlinewidth{0.000000pt}%
\definecolor{currentstroke}{rgb}{0.000000,0.000000,0.000000}%
\pgfsetstrokecolor{currentstroke}%
\pgfsetdash{}{0pt}%
\pgfpathmoveto{\pgfqpoint{1.171915in}{0.743857in}}%
\pgfpathlineto{\pgfqpoint{1.188553in}{0.750956in}}%
\pgfpathlineto{\pgfqpoint{1.205258in}{0.758838in}}%
\pgfpathlineto{\pgfqpoint{1.213280in}{0.761131in}}%
\pgfpathlineto{\pgfqpoint{1.221279in}{0.764160in}}%
\pgfpathlineto{\pgfqpoint{1.204599in}{0.755962in}}%
\pgfpathlineto{\pgfqpoint{1.187987in}{0.748527in}}%
\pgfpathlineto{\pgfqpoint{1.179964in}{0.745784in}}%
\pgfpathlineto{\pgfqpoint{1.171915in}{0.743857in}}%
\pgfpathclose%
\pgfusepath{fill}%
\end{pgfscope}%
\begin{pgfscope}%
\pgfpathrectangle{\pgfqpoint{0.407592in}{0.222904in}}{\pgfqpoint{1.411163in}{1.411163in}}%
\pgfusepath{clip}%
\pgfsetbuttcap%
\pgfsetroundjoin%
\definecolor{currentfill}{rgb}{0.556863,0.654902,0.847059}%
\pgfsetfillcolor{currentfill}%
\pgfsetfillopacity{0.130000}%
\pgfsetlinewidth{0.000000pt}%
\definecolor{currentstroke}{rgb}{0.000000,0.000000,0.000000}%
\pgfsetstrokecolor{currentstroke}%
\pgfsetdash{}{0pt}%
\pgfpathmoveto{\pgfqpoint{1.073207in}{0.729855in}}%
\pgfpathlineto{\pgfqpoint{1.089685in}{0.731394in}}%
\pgfpathlineto{\pgfqpoint{1.106213in}{0.734261in}}%
\pgfpathlineto{\pgfqpoint{1.114414in}{0.732794in}}%
\pgfpathlineto{\pgfqpoint{1.122582in}{0.731981in}}%
\pgfpathlineto{\pgfqpoint{1.106092in}{0.728593in}}%
\pgfpathlineto{\pgfqpoint{1.089654in}{0.726594in}}%
\pgfpathlineto{\pgfqpoint{1.081448in}{0.727893in}}%
\pgfpathlineto{\pgfqpoint{1.073207in}{0.729855in}}%
\pgfpathclose%
\pgfusepath{fill}%
\end{pgfscope}%
\begin{pgfscope}%
\pgfpathrectangle{\pgfqpoint{0.407592in}{0.222904in}}{\pgfqpoint{1.411163in}{1.411163in}}%
\pgfusepath{clip}%
\pgfsetbuttcap%
\pgfsetroundjoin%
\definecolor{currentfill}{rgb}{0.556863,0.654902,0.847059}%
\pgfsetfillcolor{currentfill}%
\pgfsetfillopacity{0.130000}%
\pgfsetlinewidth{0.000000pt}%
\definecolor{currentstroke}{rgb}{0.000000,0.000000,0.000000}%
\pgfsetstrokecolor{currentstroke}%
\pgfsetdash{}{0pt}%
\pgfpathmoveto{\pgfqpoint{0.941254in}{0.766967in}}%
\pgfpathlineto{\pgfqpoint{0.957622in}{0.759583in}}%
\pgfpathlineto{\pgfqpoint{0.974014in}{0.753530in}}%
\pgfpathlineto{\pgfqpoint{0.982494in}{0.749264in}}%
\pgfpathlineto{\pgfqpoint{0.990936in}{0.745394in}}%
\pgfpathlineto{\pgfqpoint{0.974593in}{0.751603in}}%
\pgfpathlineto{\pgfqpoint{0.958272in}{0.759259in}}%
\pgfpathlineto{\pgfqpoint{0.949783in}{0.762910in}}%
\pgfpathlineto{\pgfqpoint{0.941254in}{0.766967in}}%
\pgfpathclose%
\pgfusepath{fill}%
\end{pgfscope}%
\begin{pgfscope}%
\pgfpathrectangle{\pgfqpoint{0.407592in}{0.222904in}}{\pgfqpoint{1.411163in}{1.411163in}}%
\pgfusepath{clip}%
\pgfsetbuttcap%
\pgfsetroundjoin%
\definecolor{currentfill}{rgb}{0.556863,0.654902,0.847059}%
\pgfsetfillcolor{currentfill}%
\pgfsetfillopacity{0.130000}%
\pgfsetlinewidth{0.000000pt}%
\definecolor{currentstroke}{rgb}{0.000000,0.000000,0.000000}%
\pgfsetstrokecolor{currentstroke}%
\pgfsetdash{}{0pt}%
\pgfpathmoveto{\pgfqpoint{0.742947in}{0.906110in}}%
\pgfpathlineto{\pgfqpoint{0.759273in}{0.890966in}}%
\pgfpathlineto{\pgfqpoint{0.775610in}{0.876361in}}%
\pgfpathlineto{\pgfqpoint{0.784486in}{0.872881in}}%
\pgfpathlineto{\pgfqpoint{0.793319in}{0.869728in}}%
\pgfpathlineto{\pgfqpoint{0.777027in}{0.884854in}}%
\pgfpathlineto{\pgfqpoint{0.760743in}{0.900507in}}%
\pgfpathlineto{\pgfqpoint{0.751866in}{0.903143in}}%
\pgfpathlineto{\pgfqpoint{0.742947in}{0.906110in}}%
\pgfpathclose%
\pgfusepath{fill}%
\end{pgfscope}%
\begin{pgfscope}%
\pgfpathrectangle{\pgfqpoint{0.407592in}{0.222904in}}{\pgfqpoint{1.411163in}{1.411163in}}%
\pgfusepath{clip}%
\pgfsetbuttcap%
\pgfsetroundjoin%
\definecolor{currentfill}{rgb}{0.556863,0.654902,0.847059}%
\pgfsetfillcolor{currentfill}%
\pgfsetfillopacity{0.130000}%
\pgfsetlinewidth{0.000000pt}%
\definecolor{currentstroke}{rgb}{0.000000,0.000000,0.000000}%
\pgfsetstrokecolor{currentstroke}%
\pgfsetdash{}{0pt}%
\pgfpathmoveto{\pgfqpoint{1.565912in}{1.034874in}}%
\pgfpathlineto{\pgfqpoint{1.583365in}{1.045513in}}%
\pgfpathlineto{\pgfqpoint{1.600895in}{1.056387in}}%
\pgfpathlineto{\pgfqpoint{1.608461in}{1.067102in}}%
\pgfpathlineto{\pgfqpoint{1.616016in}{1.078076in}}%
\pgfpathlineto{\pgfqpoint{1.598508in}{1.067556in}}%
\pgfpathlineto{\pgfqpoint{1.581076in}{1.057287in}}%
\pgfpathlineto{\pgfqpoint{1.573500in}{1.045948in}}%
\pgfpathlineto{\pgfqpoint{1.565912in}{1.034874in}}%
\pgfpathclose%
\pgfusepath{fill}%
\end{pgfscope}%
\begin{pgfscope}%
\pgfpathrectangle{\pgfqpoint{0.407592in}{0.222904in}}{\pgfqpoint{1.411163in}{1.411163in}}%
\pgfusepath{clip}%
\pgfsetbuttcap%
\pgfsetroundjoin%
\definecolor{currentfill}{rgb}{0.556863,0.654902,0.847059}%
\pgfsetfillcolor{currentfill}%
\pgfsetfillopacity{0.130000}%
\pgfsetlinewidth{0.000000pt}%
\definecolor{currentstroke}{rgb}{0.000000,0.000000,0.000000}%
\pgfsetstrokecolor{currentstroke}%
\pgfsetdash{}{0pt}%
\pgfpathmoveto{\pgfqpoint{1.221279in}{0.764160in}}%
\pgfpathlineto{\pgfqpoint{1.238027in}{0.773003in}}%
\pgfpathlineto{\pgfqpoint{1.254845in}{0.782394in}}%
\pgfpathlineto{\pgfqpoint{1.262798in}{0.786346in}}%
\pgfpathlineto{\pgfqpoint{1.270730in}{0.790984in}}%
\pgfpathlineto{\pgfqpoint{1.253933in}{0.781477in}}%
\pgfpathlineto{\pgfqpoint{1.237207in}{0.772505in}}%
\pgfpathlineto{\pgfqpoint{1.229254in}{0.767947in}}%
\pgfpathlineto{\pgfqpoint{1.221279in}{0.764160in}}%
\pgfpathclose%
\pgfusepath{fill}%
\end{pgfscope}%
\begin{pgfscope}%
\pgfpathrectangle{\pgfqpoint{0.407592in}{0.222904in}}{\pgfqpoint{1.411163in}{1.411163in}}%
\pgfusepath{clip}%
\pgfsetbuttcap%
\pgfsetroundjoin%
\definecolor{currentfill}{rgb}{0.556863,0.654902,0.847059}%
\pgfsetfillcolor{currentfill}%
\pgfsetfillopacity{0.130000}%
\pgfsetlinewidth{0.000000pt}%
\definecolor{currentstroke}{rgb}{0.000000,0.000000,0.000000}%
\pgfsetstrokecolor{currentstroke}%
\pgfsetdash{}{0pt}%
\pgfpathmoveto{\pgfqpoint{0.858593in}{0.815816in}}%
\pgfpathlineto{\pgfqpoint{0.874940in}{0.804363in}}%
\pgfpathlineto{\pgfqpoint{0.891302in}{0.793879in}}%
\pgfpathlineto{\pgfqpoint{0.899959in}{0.789507in}}%
\pgfpathlineto{\pgfqpoint{0.908575in}{0.785475in}}%
\pgfpathlineto{\pgfqpoint{0.892259in}{0.796432in}}%
\pgfpathlineto{\pgfqpoint{0.875957in}{0.808400in}}%
\pgfpathlineto{\pgfqpoint{0.867296in}{0.811929in}}%
\pgfpathlineto{\pgfqpoint{0.858593in}{0.815816in}}%
\pgfpathclose%
\pgfusepath{fill}%
\end{pgfscope}%
\begin{pgfscope}%
\pgfpathrectangle{\pgfqpoint{0.407592in}{0.222904in}}{\pgfqpoint{1.411163in}{1.411163in}}%
\pgfusepath{clip}%
\pgfsetbuttcap%
\pgfsetroundjoin%
\definecolor{currentfill}{rgb}{0.556863,0.654902,0.847059}%
\pgfsetfillcolor{currentfill}%
\pgfsetfillopacity{0.130000}%
\pgfsetlinewidth{0.000000pt}%
\definecolor{currentstroke}{rgb}{0.000000,0.000000,0.000000}%
\pgfsetstrokecolor{currentstroke}%
\pgfsetdash{}{0pt}%
\pgfpathmoveto{\pgfqpoint{0.626947in}{1.009861in}}%
\pgfpathlineto{\pgfqpoint{0.643264in}{0.992828in}}%
\pgfpathlineto{\pgfqpoint{0.659589in}{0.976103in}}%
\pgfpathlineto{\pgfqpoint{0.668681in}{0.973451in}}%
\pgfpathlineto{\pgfqpoint{0.677729in}{0.971095in}}%
\pgfpathlineto{\pgfqpoint{0.661447in}{0.988235in}}%
\pgfpathlineto{\pgfqpoint{0.645172in}{1.005662in}}%
\pgfpathlineto{\pgfqpoint{0.636082in}{1.007614in}}%
\pgfpathlineto{\pgfqpoint{0.626947in}{1.009861in}}%
\pgfpathclose%
\pgfusepath{fill}%
\end{pgfscope}%
\begin{pgfscope}%
\pgfpathrectangle{\pgfqpoint{0.407592in}{0.222904in}}{\pgfqpoint{1.411163in}{1.411163in}}%
\pgfusepath{clip}%
\pgfsetbuttcap%
\pgfsetroundjoin%
\definecolor{currentfill}{rgb}{0.556863,0.654902,0.847059}%
\pgfsetfillcolor{currentfill}%
\pgfsetfillopacity{0.130000}%
\pgfsetlinewidth{0.000000pt}%
\definecolor{currentstroke}{rgb}{0.000000,0.000000,0.000000}%
\pgfsetstrokecolor{currentstroke}%
\pgfsetdash{}{0pt}%
\pgfpathmoveto{\pgfqpoint{1.385613in}{0.873006in}}%
\pgfpathlineto{\pgfqpoint{1.402716in}{0.883845in}}%
\pgfpathlineto{\pgfqpoint{1.419893in}{0.894947in}}%
\pgfpathlineto{\pgfqpoint{1.427648in}{0.903150in}}%
\pgfpathlineto{\pgfqpoint{1.435389in}{0.911799in}}%
\pgfpathlineto{\pgfqpoint{1.418231in}{0.900994in}}%
\pgfpathlineto{\pgfqpoint{1.401148in}{0.890475in}}%
\pgfpathlineto{\pgfqpoint{1.393388in}{0.881500in}}%
\pgfpathlineto{\pgfqpoint{1.385613in}{0.873006in}}%
\pgfpathclose%
\pgfusepath{fill}%
\end{pgfscope}%
\begin{pgfscope}%
\pgfpathrectangle{\pgfqpoint{0.407592in}{0.222904in}}{\pgfqpoint{1.411163in}{1.411163in}}%
\pgfusepath{clip}%
\pgfsetbuttcap%
\pgfsetroundjoin%
\definecolor{currentfill}{rgb}{0.556863,0.654902,0.847059}%
\pgfsetfillcolor{currentfill}%
\pgfsetfillopacity{0.130000}%
\pgfsetlinewidth{0.000000pt}%
\definecolor{currentstroke}{rgb}{0.000000,0.000000,0.000000}%
\pgfsetstrokecolor{currentstroke}%
\pgfsetdash{}{0pt}%
\pgfpathmoveto{\pgfqpoint{0.477981in}{1.157044in}}%
\pgfpathlineto{\pgfqpoint{0.494272in}{1.138639in}}%
\pgfpathlineto{\pgfqpoint{0.510571in}{1.120393in}}%
\pgfpathlineto{\pgfqpoint{0.519921in}{1.118620in}}%
\pgfpathlineto{\pgfqpoint{0.529227in}{1.117104in}}%
\pgfpathlineto{\pgfqpoint{0.512970in}{1.135632in}}%
\pgfpathlineto{\pgfqpoint{0.496721in}{1.154304in}}%
\pgfpathlineto{\pgfqpoint{0.487373in}{1.155548in}}%
\pgfpathlineto{\pgfqpoint{0.477981in}{1.157044in}}%
\pgfpathclose%
\pgfusepath{fill}%
\end{pgfscope}%
\begin{pgfscope}%
\pgfpathrectangle{\pgfqpoint{0.407592in}{0.222904in}}{\pgfqpoint{1.411163in}{1.411163in}}%
\pgfusepath{clip}%
\pgfsetbuttcap%
\pgfsetroundjoin%
\definecolor{currentfill}{rgb}{0.556863,0.654902,0.847059}%
\pgfsetfillcolor{currentfill}%
\pgfsetfillopacity{0.130000}%
\pgfsetlinewidth{0.000000pt}%
\definecolor{currentstroke}{rgb}{0.000000,0.000000,0.000000}%
\pgfsetstrokecolor{currentstroke}%
\pgfsetdash{}{0pt}%
\pgfpathmoveto{\pgfqpoint{1.023708in}{0.737478in}}%
\pgfpathlineto{\pgfqpoint{1.040144in}{0.735765in}}%
\pgfpathlineto{\pgfqpoint{1.056620in}{0.735487in}}%
\pgfpathlineto{\pgfqpoint{1.064932in}{0.732409in}}%
\pgfpathlineto{\pgfqpoint{1.073207in}{0.729855in}}%
\pgfpathlineto{\pgfqpoint{1.056775in}{0.729769in}}%
\pgfpathlineto{\pgfqpoint{1.040384in}{0.731236in}}%
\pgfpathlineto{\pgfqpoint{1.032064in}{0.734099in}}%
\pgfpathlineto{\pgfqpoint{1.023708in}{0.737478in}}%
\pgfpathclose%
\pgfusepath{fill}%
\end{pgfscope}%
\begin{pgfscope}%
\pgfpathrectangle{\pgfqpoint{0.407592in}{0.222904in}}{\pgfqpoint{1.411163in}{1.411163in}}%
\pgfusepath{clip}%
\pgfsetbuttcap%
\pgfsetroundjoin%
\definecolor{currentfill}{rgb}{0.556863,0.654902,0.847059}%
\pgfsetfillcolor{currentfill}%
\pgfsetfillopacity{0.130000}%
\pgfsetlinewidth{0.000000pt}%
\definecolor{currentstroke}{rgb}{0.000000,0.000000,0.000000}%
\pgfsetstrokecolor{currentstroke}%
\pgfsetdash{}{0pt}%
\pgfpathmoveto{\pgfqpoint{1.746506in}{1.212824in}}%
\pgfpathlineto{\pgfqpoint{1.764311in}{1.223157in}}%
\pgfpathlineto{\pgfqpoint{1.782194in}{1.233709in}}%
\pgfpathlineto{\pgfqpoint{1.789583in}{1.245560in}}%
\pgfpathlineto{\pgfqpoint{1.771711in}{1.235160in}}%
\pgfpathlineto{\pgfqpoint{1.753918in}{1.224984in}}%
\pgfusepath{fill}%
\end{pgfscope}%
\begin{pgfscope}%
\pgfpathrectangle{\pgfqpoint{0.407592in}{0.222904in}}{\pgfqpoint{1.411163in}{1.411163in}}%
\pgfusepath{clip}%
\pgfsetbuttcap%
\pgfsetroundjoin%
\definecolor{currentfill}{rgb}{0.556863,0.654902,0.847059}%
\pgfsetfillcolor{currentfill}%
\pgfsetfillopacity{0.130000}%
\pgfsetlinewidth{0.000000pt}%
\definecolor{currentstroke}{rgb}{0.000000,0.000000,0.000000}%
\pgfsetstrokecolor{currentstroke}%
\pgfsetdash{}{0pt}%
\pgfpathmoveto{\pgfqpoint{1.681284in}{1.144691in}}%
\pgfpathlineto{\pgfqpoint{1.698979in}{1.155221in}}%
\pgfpathlineto{\pgfqpoint{1.716752in}{1.165971in}}%
\pgfpathlineto{\pgfqpoint{1.724207in}{1.177400in}}%
\pgfpathlineto{\pgfqpoint{1.731650in}{1.189024in}}%
\pgfpathlineto{\pgfqpoint{1.713901in}{1.178597in}}%
\pgfpathlineto{\pgfqpoint{1.696228in}{1.168401in}}%
\pgfpathlineto{\pgfqpoint{1.688761in}{1.156448in}}%
\pgfpathlineto{\pgfqpoint{1.681284in}{1.144691in}}%
\pgfpathclose%
\pgfusepath{fill}%
\end{pgfscope}%
\begin{pgfscope}%
\pgfpathrectangle{\pgfqpoint{0.407592in}{0.222904in}}{\pgfqpoint{1.411163in}{1.411163in}}%
\pgfusepath{clip}%
\pgfsetbuttcap%
\pgfsetroundjoin%
\definecolor{currentfill}{rgb}{0.556863,0.654902,0.847059}%
\pgfsetfillcolor{currentfill}%
\pgfsetfillopacity{0.130000}%
\pgfsetlinewidth{0.000000pt}%
\definecolor{currentstroke}{rgb}{0.000000,0.000000,0.000000}%
\pgfsetstrokecolor{currentstroke}%
\pgfsetdash{}{0pt}%
\pgfpathmoveto{\pgfqpoint{1.500680in}{0.971504in}}%
\pgfpathlineto{\pgfqpoint{1.518023in}{0.982381in}}%
\pgfpathlineto{\pgfqpoint{1.535443in}{0.993491in}}%
\pgfpathlineto{\pgfqpoint{1.543078in}{1.003374in}}%
\pgfpathlineto{\pgfqpoint{1.550701in}{1.013574in}}%
\pgfpathlineto{\pgfqpoint{1.533302in}{1.002815in}}%
\pgfpathlineto{\pgfqpoint{1.515980in}{0.992308in}}%
\pgfpathlineto{\pgfqpoint{1.508336in}{0.981741in}}%
\pgfpathlineto{\pgfqpoint{1.500680in}{0.971504in}}%
\pgfpathclose%
\pgfusepath{fill}%
\end{pgfscope}%
\begin{pgfscope}%
\pgfpathrectangle{\pgfqpoint{0.407592in}{0.222904in}}{\pgfqpoint{1.411163in}{1.411163in}}%
\pgfusepath{clip}%
\pgfsetbuttcap%
\pgfsetroundjoin%
\definecolor{currentfill}{rgb}{0.556863,0.654902,0.847059}%
\pgfsetfillcolor{currentfill}%
\pgfsetfillopacity{0.130000}%
\pgfsetlinewidth{0.000000pt}%
\definecolor{currentstroke}{rgb}{0.000000,0.000000,0.000000}%
\pgfsetstrokecolor{currentstroke}%
\pgfsetdash{}{0pt}%
\pgfpathmoveto{\pgfqpoint{1.270730in}{0.790984in}}%
\pgfpathlineto{\pgfqpoint{1.287598in}{0.800949in}}%
\pgfpathlineto{\pgfqpoint{1.304539in}{0.811307in}}%
\pgfpathlineto{\pgfqpoint{1.312430in}{0.816626in}}%
\pgfpathlineto{\pgfqpoint{1.320303in}{0.822560in}}%
\pgfpathlineto{\pgfqpoint{1.303383in}{0.812248in}}%
\pgfpathlineto{\pgfqpoint{1.286535in}{0.802331in}}%
\pgfpathlineto{\pgfqpoint{1.278642in}{0.796313in}}%
\pgfpathlineto{\pgfqpoint{1.270730in}{0.790984in}}%
\pgfpathclose%
\pgfusepath{fill}%
\end{pgfscope}%
\begin{pgfscope}%
\pgfpathrectangle{\pgfqpoint{0.407592in}{0.222904in}}{\pgfqpoint{1.411163in}{1.411163in}}%
\pgfusepath{clip}%
\pgfsetbuttcap%
\pgfsetroundjoin%
\definecolor{currentfill}{rgb}{0.556863,0.654902,0.847059}%
\pgfsetfillcolor{currentfill}%
\pgfsetfillopacity{0.130000}%
\pgfsetlinewidth{0.000000pt}%
\definecolor{currentstroke}{rgb}{0.000000,0.000000,0.000000}%
\pgfsetstrokecolor{currentstroke}%
\pgfsetdash{}{0pt}%
\pgfpathmoveto{\pgfqpoint{0.775610in}{0.876361in}}%
\pgfpathlineto{\pgfqpoint{0.791956in}{0.862353in}}%
\pgfpathlineto{\pgfqpoint{0.808313in}{0.849002in}}%
\pgfpathlineto{\pgfqpoint{0.817145in}{0.844999in}}%
\pgfpathlineto{\pgfqpoint{0.825934in}{0.841309in}}%
\pgfpathlineto{\pgfqpoint{0.809622in}{0.855190in}}%
\pgfpathlineto{\pgfqpoint{0.793319in}{0.869728in}}%
\pgfpathlineto{\pgfqpoint{0.784486in}{0.872881in}}%
\pgfpathlineto{\pgfqpoint{0.775610in}{0.876361in}}%
\pgfpathclose%
\pgfusepath{fill}%
\end{pgfscope}%
\begin{pgfscope}%
\pgfpathrectangle{\pgfqpoint{0.407592in}{0.222904in}}{\pgfqpoint{1.411163in}{1.411163in}}%
\pgfusepath{clip}%
\pgfsetbuttcap%
\pgfsetroundjoin%
\definecolor{currentfill}{rgb}{0.556863,0.654902,0.847059}%
\pgfsetfillcolor{currentfill}%
\pgfsetfillopacity{0.130000}%
\pgfsetlinewidth{0.000000pt}%
\definecolor{currentstroke}{rgb}{0.000000,0.000000,0.000000}%
\pgfsetstrokecolor{currentstroke}%
\pgfsetdash{}{0pt}%
\pgfpathmoveto{\pgfqpoint{0.510571in}{1.120393in}}%
\pgfpathlineto{\pgfqpoint{0.526879in}{1.102321in}}%
\pgfpathlineto{\pgfqpoint{0.543195in}{1.084439in}}%
\pgfpathlineto{\pgfqpoint{0.552503in}{1.082357in}}%
\pgfpathlineto{\pgfqpoint{0.561766in}{1.080538in}}%
\pgfpathlineto{\pgfqpoint{0.545492in}{1.098734in}}%
\pgfpathlineto{\pgfqpoint{0.529227in}{1.117104in}}%
\pgfpathlineto{\pgfqpoint{0.519921in}{1.118620in}}%
\pgfpathlineto{\pgfqpoint{0.510571in}{1.120393in}}%
\pgfpathclose%
\pgfusepath{fill}%
\end{pgfscope}%
\begin{pgfscope}%
\pgfpathrectangle{\pgfqpoint{0.407592in}{0.222904in}}{\pgfqpoint{1.411163in}{1.411163in}}%
\pgfusepath{clip}%
\pgfsetbuttcap%
\pgfsetroundjoin%
\definecolor{currentfill}{rgb}{0.556863,0.654902,0.847059}%
\pgfsetfillcolor{currentfill}%
\pgfsetfillopacity{0.130000}%
\pgfsetlinewidth{0.000000pt}%
\definecolor{currentstroke}{rgb}{0.000000,0.000000,0.000000}%
\pgfsetstrokecolor{currentstroke}%
\pgfsetdash{}{0pt}%
\pgfpathmoveto{\pgfqpoint{0.659589in}{0.976103in}}%
\pgfpathlineto{\pgfqpoint{0.675923in}{0.959718in}}%
\pgfpathlineto{\pgfqpoint{0.692266in}{0.943706in}}%
\pgfpathlineto{\pgfqpoint{0.701315in}{0.940610in}}%
\pgfpathlineto{\pgfqpoint{0.710320in}{0.937811in}}%
\pgfpathlineto{\pgfqpoint{0.694020in}{0.954275in}}%
\pgfpathlineto{\pgfqpoint{0.677729in}{0.971095in}}%
\pgfpathlineto{\pgfqpoint{0.668681in}{0.973451in}}%
\pgfpathlineto{\pgfqpoint{0.659589in}{0.976103in}}%
\pgfpathclose%
\pgfusepath{fill}%
\end{pgfscope}%
\begin{pgfscope}%
\pgfpathrectangle{\pgfqpoint{0.407592in}{0.222904in}}{\pgfqpoint{1.411163in}{1.411163in}}%
\pgfusepath{clip}%
\pgfsetbuttcap%
\pgfsetroundjoin%
\definecolor{currentfill}{rgb}{0.556863,0.654902,0.847059}%
\pgfsetfillcolor{currentfill}%
\pgfsetfillopacity{0.130000}%
\pgfsetlinewidth{0.000000pt}%
\definecolor{currentstroke}{rgb}{0.000000,0.000000,0.000000}%
\pgfsetstrokecolor{currentstroke}%
\pgfsetdash{}{0pt}%
\pgfpathmoveto{\pgfqpoint{1.616016in}{1.078076in}}%
\pgfpathlineto{\pgfqpoint{1.633602in}{1.088830in}}%
\pgfpathlineto{\pgfqpoint{1.651264in}{1.099804in}}%
\pgfpathlineto{\pgfqpoint{1.658786in}{1.110686in}}%
\pgfpathlineto{\pgfqpoint{1.666296in}{1.121801in}}%
\pgfpathlineto{\pgfqpoint{1.648656in}{1.111162in}}%
\pgfpathlineto{\pgfqpoint{1.631093in}{1.100757in}}%
\pgfpathlineto{\pgfqpoint{1.623560in}{1.089298in}}%
\pgfpathlineto{\pgfqpoint{1.616016in}{1.078076in}}%
\pgfpathclose%
\pgfusepath{fill}%
\end{pgfscope}%
\begin{pgfscope}%
\pgfpathrectangle{\pgfqpoint{0.407592in}{0.222904in}}{\pgfqpoint{1.411163in}{1.411163in}}%
\pgfusepath{clip}%
\pgfsetbuttcap%
\pgfsetroundjoin%
\definecolor{currentfill}{rgb}{0.556863,0.654902,0.847059}%
\pgfsetfillcolor{currentfill}%
\pgfsetfillopacity{0.130000}%
\pgfsetlinewidth{0.000000pt}%
\definecolor{currentstroke}{rgb}{0.000000,0.000000,0.000000}%
\pgfsetstrokecolor{currentstroke}%
\pgfsetdash{}{0pt}%
\pgfpathmoveto{\pgfqpoint{0.891302in}{0.793879in}}%
\pgfpathlineto{\pgfqpoint{0.907681in}{0.784441in}}%
\pgfpathlineto{\pgfqpoint{0.924077in}{0.776122in}}%
\pgfpathlineto{\pgfqpoint{0.932685in}{0.771385in}}%
\pgfpathlineto{\pgfqpoint{0.941254in}{0.766967in}}%
\pgfpathlineto{\pgfqpoint{0.924906in}{0.775624in}}%
\pgfpathlineto{\pgfqpoint{0.908575in}{0.785475in}}%
\pgfpathlineto{\pgfqpoint{0.899959in}{0.789507in}}%
\pgfpathlineto{\pgfqpoint{0.891302in}{0.793879in}}%
\pgfpathclose%
\pgfusepath{fill}%
\end{pgfscope}%
\begin{pgfscope}%
\pgfpathrectangle{\pgfqpoint{0.407592in}{0.222904in}}{\pgfqpoint{1.411163in}{1.411163in}}%
\pgfusepath{clip}%
\pgfsetbuttcap%
\pgfsetroundjoin%
\definecolor{currentfill}{rgb}{0.556863,0.654902,0.847059}%
\pgfsetfillcolor{currentfill}%
\pgfsetfillopacity{0.130000}%
\pgfsetlinewidth{0.000000pt}%
\definecolor{currentstroke}{rgb}{0.000000,0.000000,0.000000}%
\pgfsetstrokecolor{currentstroke}%
\pgfsetdash{}{0pt}%
\pgfpathmoveto{\pgfqpoint{1.435389in}{0.911799in}}%
\pgfpathlineto{\pgfqpoint{1.452621in}{0.922862in}}%
\pgfpathlineto{\pgfqpoint{1.469930in}{0.934159in}}%
\pgfpathlineto{\pgfqpoint{1.477637in}{0.942926in}}%
\pgfpathlineto{\pgfqpoint{1.485330in}{0.952082in}}%
\pgfpathlineto{\pgfqpoint{1.468042in}{0.941100in}}%
\pgfpathlineto{\pgfqpoint{1.450830in}{0.930373in}}%
\pgfpathlineto{\pgfqpoint{1.443116in}{0.920879in}}%
\pgfpathlineto{\pgfqpoint{1.435389in}{0.911799in}}%
\pgfpathclose%
\pgfusepath{fill}%
\end{pgfscope}%
\begin{pgfscope}%
\pgfpathrectangle{\pgfqpoint{0.407592in}{0.222904in}}{\pgfqpoint{1.411163in}{1.411163in}}%
\pgfusepath{clip}%
\pgfsetbuttcap%
\pgfsetroundjoin%
\definecolor{currentfill}{rgb}{0.556863,0.654902,0.847059}%
\pgfsetfillcolor{currentfill}%
\pgfsetfillopacity{0.130000}%
\pgfsetlinewidth{0.000000pt}%
\definecolor{currentstroke}{rgb}{0.000000,0.000000,0.000000}%
\pgfsetstrokecolor{currentstroke}%
\pgfsetdash{}{0pt}%
\pgfpathmoveto{\pgfqpoint{1.320303in}{0.822560in}}%
\pgfpathlineto{\pgfqpoint{1.337297in}{0.833216in}}%
\pgfpathlineto{\pgfqpoint{1.354365in}{0.844172in}}%
\pgfpathlineto{\pgfqpoint{1.362201in}{0.850582in}}%
\pgfpathlineto{\pgfqpoint{1.370020in}{0.857532in}}%
\pgfpathlineto{\pgfqpoint{1.352972in}{0.846730in}}%
\pgfpathlineto{\pgfqpoint{1.335998in}{0.836242in}}%
\pgfpathlineto{\pgfqpoint{1.328159in}{0.829102in}}%
\pgfpathlineto{\pgfqpoint{1.320303in}{0.822560in}}%
\pgfpathclose%
\pgfusepath{fill}%
\end{pgfscope}%
\begin{pgfscope}%
\pgfpathrectangle{\pgfqpoint{0.407592in}{0.222904in}}{\pgfqpoint{1.411163in}{1.411163in}}%
\pgfusepath{clip}%
\pgfsetbuttcap%
\pgfsetroundjoin%
\definecolor{currentfill}{rgb}{0.556863,0.654902,0.847059}%
\pgfsetfillcolor{currentfill}%
\pgfsetfillopacity{0.130000}%
\pgfsetlinewidth{0.000000pt}%
\definecolor{currentstroke}{rgb}{0.000000,0.000000,0.000000}%
\pgfsetstrokecolor{currentstroke}%
\pgfsetdash{}{0pt}%
\pgfpathmoveto{\pgfqpoint{0.974014in}{0.753530in}}%
\pgfpathlineto{\pgfqpoint{0.990432in}{0.748845in}}%
\pgfpathlineto{\pgfqpoint{1.006881in}{0.745540in}}%
\pgfpathlineto{\pgfqpoint{1.015313in}{0.741310in}}%
\pgfpathlineto{\pgfqpoint{1.023708in}{0.737478in}}%
\pgfpathlineto{\pgfqpoint{1.007307in}{0.740681in}}%
\pgfpathlineto{\pgfqpoint{0.990936in}{0.745394in}}%
\pgfpathlineto{\pgfqpoint{0.982494in}{0.749264in}}%
\pgfpathlineto{\pgfqpoint{0.974014in}{0.753530in}}%
\pgfpathclose%
\pgfusepath{fill}%
\end{pgfscope}%
\begin{pgfscope}%
\pgfpathrectangle{\pgfqpoint{0.407592in}{0.222904in}}{\pgfqpoint{1.411163in}{1.411163in}}%
\pgfusepath{clip}%
\pgfsetbuttcap%
\pgfsetroundjoin%
\definecolor{currentfill}{rgb}{0.556863,0.654902,0.847059}%
\pgfsetfillcolor{currentfill}%
\pgfsetfillopacity{0.130000}%
\pgfsetlinewidth{0.000000pt}%
\definecolor{currentstroke}{rgb}{0.000000,0.000000,0.000000}%
\pgfsetstrokecolor{currentstroke}%
\pgfsetdash{}{0pt}%
\pgfpathmoveto{\pgfqpoint{0.543195in}{1.084439in}}%
\pgfpathlineto{\pgfqpoint{0.559520in}{1.066764in}}%
\pgfpathlineto{\pgfqpoint{0.575854in}{1.049316in}}%
\pgfpathlineto{\pgfqpoint{0.585119in}{1.046893in}}%
\pgfpathlineto{\pgfqpoint{0.594340in}{1.044738in}}%
\pgfpathlineto{\pgfqpoint{0.578049in}{1.062533in}}%
\pgfpathlineto{\pgfqpoint{0.561766in}{1.080538in}}%
\pgfpathlineto{\pgfqpoint{0.552503in}{1.082357in}}%
\pgfpathlineto{\pgfqpoint{0.543195in}{1.084439in}}%
\pgfpathclose%
\pgfusepath{fill}%
\end{pgfscope}%
\begin{pgfscope}%
\pgfpathrectangle{\pgfqpoint{0.407592in}{0.222904in}}{\pgfqpoint{1.411163in}{1.411163in}}%
\pgfusepath{clip}%
\pgfsetbuttcap%
\pgfsetroundjoin%
\definecolor{currentfill}{rgb}{0.556863,0.654902,0.847059}%
\pgfsetfillcolor{currentfill}%
\pgfsetfillopacity{0.130000}%
\pgfsetlinewidth{0.000000pt}%
\definecolor{currentstroke}{rgb}{0.000000,0.000000,0.000000}%
\pgfsetstrokecolor{currentstroke}%
\pgfsetdash{}{0pt}%
\pgfpathmoveto{\pgfqpoint{1.550701in}{1.013574in}}%
\pgfpathlineto{\pgfqpoint{1.568176in}{1.024568in}}%
\pgfpathlineto{\pgfqpoint{1.585728in}{1.035780in}}%
\pgfpathlineto{\pgfqpoint{1.593317in}{1.045943in}}%
\pgfpathlineto{\pgfqpoint{1.600895in}{1.056387in}}%
\pgfpathlineto{\pgfqpoint{1.583365in}{1.045513in}}%
\pgfpathlineto{\pgfqpoint{1.565912in}{1.034874in}}%
\pgfpathlineto{\pgfqpoint{1.558312in}{1.024079in}}%
\pgfpathlineto{\pgfqpoint{1.550701in}{1.013574in}}%
\pgfpathclose%
\pgfusepath{fill}%
\end{pgfscope}%
\begin{pgfscope}%
\pgfpathrectangle{\pgfqpoint{0.407592in}{0.222904in}}{\pgfqpoint{1.411163in}{1.411163in}}%
\pgfusepath{clip}%
\pgfsetbuttcap%
\pgfsetroundjoin%
\definecolor{currentfill}{rgb}{0.556863,0.654902,0.847059}%
\pgfsetfillcolor{currentfill}%
\pgfsetfillopacity{0.130000}%
\pgfsetlinewidth{0.000000pt}%
\definecolor{currentstroke}{rgb}{0.000000,0.000000,0.000000}%
\pgfsetstrokecolor{currentstroke}%
\pgfsetdash{}{0pt}%
\pgfpathmoveto{\pgfqpoint{1.731650in}{1.189024in}}%
\pgfpathlineto{\pgfqpoint{1.749479in}{1.199672in}}%
\pgfpathlineto{\pgfqpoint{1.767385in}{1.210528in}}%
\pgfpathlineto{\pgfqpoint{1.774795in}{1.222030in}}%
\pgfpathlineto{\pgfqpoint{1.782194in}{1.233709in}}%
\pgfpathlineto{\pgfqpoint{1.764311in}{1.223157in}}%
\pgfpathlineto{\pgfqpoint{1.746506in}{1.212824in}}%
\pgfpathlineto{\pgfqpoint{1.739084in}{1.200835in}}%
\pgfpathlineto{\pgfqpoint{1.731650in}{1.189024in}}%
\pgfpathclose%
\pgfusepath{fill}%
\end{pgfscope}%
\begin{pgfscope}%
\pgfpathrectangle{\pgfqpoint{0.407592in}{0.222904in}}{\pgfqpoint{1.411163in}{1.411163in}}%
\pgfusepath{clip}%
\pgfsetbuttcap%
\pgfsetroundjoin%
\definecolor{currentfill}{rgb}{0.556863,0.654902,0.847059}%
\pgfsetfillcolor{currentfill}%
\pgfsetfillopacity{0.130000}%
\pgfsetlinewidth{0.000000pt}%
\definecolor{currentstroke}{rgb}{0.000000,0.000000,0.000000}%
\pgfsetstrokecolor{currentstroke}%
\pgfsetdash{}{0pt}%
\pgfpathmoveto{\pgfqpoint{0.808313in}{0.849002in}}%
\pgfpathlineto{\pgfqpoint{0.824682in}{0.836375in}}%
\pgfpathlineto{\pgfqpoint{0.841063in}{0.824538in}}%
\pgfpathlineto{\pgfqpoint{0.849849in}{0.820029in}}%
\pgfpathlineto{\pgfqpoint{0.858593in}{0.815816in}}%
\pgfpathlineto{\pgfqpoint{0.842258in}{0.828159in}}%
\pgfpathlineto{\pgfqpoint{0.825934in}{0.841309in}}%
\pgfpathlineto{\pgfqpoint{0.817145in}{0.844999in}}%
\pgfpathlineto{\pgfqpoint{0.808313in}{0.849002in}}%
\pgfpathclose%
\pgfusepath{fill}%
\end{pgfscope}%
\begin{pgfscope}%
\pgfpathrectangle{\pgfqpoint{0.407592in}{0.222904in}}{\pgfqpoint{1.411163in}{1.411163in}}%
\pgfusepath{clip}%
\pgfsetbuttcap%
\pgfsetroundjoin%
\definecolor{currentfill}{rgb}{0.556863,0.654902,0.847059}%
\pgfsetfillcolor{currentfill}%
\pgfsetfillopacity{0.130000}%
\pgfsetlinewidth{0.000000pt}%
\definecolor{currentstroke}{rgb}{0.000000,0.000000,0.000000}%
\pgfsetstrokecolor{currentstroke}%
\pgfsetdash{}{0pt}%
\pgfpathmoveto{\pgfqpoint{0.692266in}{0.943706in}}%
\pgfpathlineto{\pgfqpoint{0.708618in}{0.928106in}}%
\pgfpathlineto{\pgfqpoint{0.724979in}{0.912959in}}%
\pgfpathlineto{\pgfqpoint{0.733984in}{0.909389in}}%
\pgfpathlineto{\pgfqpoint{0.742947in}{0.906110in}}%
\pgfpathlineto{\pgfqpoint{0.726629in}{0.921741in}}%
\pgfpathlineto{\pgfqpoint{0.710320in}{0.937811in}}%
\pgfpathlineto{\pgfqpoint{0.701315in}{0.940610in}}%
\pgfpathlineto{\pgfqpoint{0.692266in}{0.943706in}}%
\pgfpathclose%
\pgfusepath{fill}%
\end{pgfscope}%
\begin{pgfscope}%
\pgfpathrectangle{\pgfqpoint{0.407592in}{0.222904in}}{\pgfqpoint{1.411163in}{1.411163in}}%
\pgfusepath{clip}%
\pgfsetbuttcap%
\pgfsetroundjoin%
\definecolor{currentfill}{rgb}{0.556863,0.654902,0.847059}%
\pgfsetfillcolor{currentfill}%
\pgfsetfillopacity{0.130000}%
\pgfsetlinewidth{0.000000pt}%
\definecolor{currentstroke}{rgb}{0.000000,0.000000,0.000000}%
\pgfsetstrokecolor{currentstroke}%
\pgfsetdash{}{0pt}%
\pgfpathmoveto{\pgfqpoint{1.155734in}{0.742290in}}%
\pgfpathlineto{\pgfqpoint{1.172402in}{0.748915in}}%
\pgfpathlineto{\pgfqpoint{1.189135in}{0.756344in}}%
\pgfpathlineto{\pgfqpoint{1.197210in}{0.757253in}}%
\pgfpathlineto{\pgfqpoint{1.205258in}{0.758838in}}%
\pgfpathlineto{\pgfqpoint{1.188553in}{0.750956in}}%
\pgfpathlineto{\pgfqpoint{1.171915in}{0.743857in}}%
\pgfpathlineto{\pgfqpoint{1.163838in}{0.742708in}}%
\pgfpathlineto{\pgfqpoint{1.155734in}{0.742290in}}%
\pgfpathclose%
\pgfusepath{fill}%
\end{pgfscope}%
\begin{pgfscope}%
\pgfpathrectangle{\pgfqpoint{0.407592in}{0.222904in}}{\pgfqpoint{1.411163in}{1.411163in}}%
\pgfusepath{clip}%
\pgfsetbuttcap%
\pgfsetroundjoin%
\definecolor{currentfill}{rgb}{0.556863,0.654902,0.847059}%
\pgfsetfillcolor{currentfill}%
\pgfsetfillopacity{0.130000}%
\pgfsetlinewidth{0.000000pt}%
\definecolor{currentstroke}{rgb}{0.000000,0.000000,0.000000}%
\pgfsetstrokecolor{currentstroke}%
\pgfsetdash{}{0pt}%
\pgfpathmoveto{\pgfqpoint{1.106213in}{0.734261in}}%
\pgfpathlineto{\pgfqpoint{1.122795in}{0.738324in}}%
\pgfpathlineto{\pgfqpoint{1.139435in}{0.743451in}}%
\pgfpathlineto{\pgfqpoint{1.147600in}{0.742555in}}%
\pgfpathlineto{\pgfqpoint{1.155734in}{0.742290in}}%
\pgfpathlineto{\pgfqpoint{1.139128in}{0.736598in}}%
\pgfpathlineto{\pgfqpoint{1.122582in}{0.731981in}}%
\pgfpathlineto{\pgfqpoint{1.114414in}{0.732794in}}%
\pgfpathlineto{\pgfqpoint{1.106213in}{0.734261in}}%
\pgfpathclose%
\pgfusepath{fill}%
\end{pgfscope}%
\begin{pgfscope}%
\pgfpathrectangle{\pgfqpoint{0.407592in}{0.222904in}}{\pgfqpoint{1.411163in}{1.411163in}}%
\pgfusepath{clip}%
\pgfsetbuttcap%
\pgfsetroundjoin%
\definecolor{currentfill}{rgb}{0.556863,0.654902,0.847059}%
\pgfsetfillcolor{currentfill}%
\pgfsetfillopacity{0.130000}%
\pgfsetlinewidth{0.000000pt}%
\definecolor{currentstroke}{rgb}{0.000000,0.000000,0.000000}%
\pgfsetstrokecolor{currentstroke}%
\pgfsetdash{}{0pt}%
\pgfpathmoveto{\pgfqpoint{1.666296in}{1.121801in}}%
\pgfpathlineto{\pgfqpoint{1.684014in}{1.132660in}}%
\pgfpathlineto{\pgfqpoint{1.701811in}{1.143728in}}%
\pgfpathlineto{\pgfqpoint{1.709287in}{1.154744in}}%
\pgfpathlineto{\pgfqpoint{1.716752in}{1.165971in}}%
\pgfpathlineto{\pgfqpoint{1.698979in}{1.155221in}}%
\pgfpathlineto{\pgfqpoint{1.681284in}{1.144691in}}%
\pgfpathlineto{\pgfqpoint{1.673795in}{1.133139in}}%
\pgfpathlineto{\pgfqpoint{1.666296in}{1.121801in}}%
\pgfpathclose%
\pgfusepath{fill}%
\end{pgfscope}%
\begin{pgfscope}%
\pgfpathrectangle{\pgfqpoint{0.407592in}{0.222904in}}{\pgfqpoint{1.411163in}{1.411163in}}%
\pgfusepath{clip}%
\pgfsetbuttcap%
\pgfsetroundjoin%
\definecolor{currentfill}{rgb}{0.556863,0.654902,0.847059}%
\pgfsetfillcolor{currentfill}%
\pgfsetfillopacity{0.130000}%
\pgfsetlinewidth{0.000000pt}%
\definecolor{currentstroke}{rgb}{0.000000,0.000000,0.000000}%
\pgfsetstrokecolor{currentstroke}%
\pgfsetdash{}{0pt}%
\pgfpathmoveto{\pgfqpoint{1.370020in}{0.857532in}}%
\pgfpathlineto{\pgfqpoint{1.387143in}{0.868610in}}%
\pgfpathlineto{\pgfqpoint{1.404341in}{0.879934in}}%
\pgfpathlineto{\pgfqpoint{1.412124in}{0.887205in}}%
\pgfpathlineto{\pgfqpoint{1.419893in}{0.894947in}}%
\pgfpathlineto{\pgfqpoint{1.402716in}{0.883845in}}%
\pgfpathlineto{\pgfqpoint{1.385613in}{0.873006in}}%
\pgfpathlineto{\pgfqpoint{1.377824in}{0.865011in}}%
\pgfpathlineto{\pgfqpoint{1.370020in}{0.857532in}}%
\pgfpathclose%
\pgfusepath{fill}%
\end{pgfscope}%
\begin{pgfscope}%
\pgfpathrectangle{\pgfqpoint{0.407592in}{0.222904in}}{\pgfqpoint{1.411163in}{1.411163in}}%
\pgfusepath{clip}%
\pgfsetbuttcap%
\pgfsetroundjoin%
\definecolor{currentfill}{rgb}{0.556863,0.654902,0.847059}%
\pgfsetfillcolor{currentfill}%
\pgfsetfillopacity{0.130000}%
\pgfsetlinewidth{0.000000pt}%
\definecolor{currentstroke}{rgb}{0.000000,0.000000,0.000000}%
\pgfsetstrokecolor{currentstroke}%
\pgfsetdash{}{0pt}%
\pgfpathmoveto{\pgfqpoint{0.575854in}{1.049316in}}%
\pgfpathlineto{\pgfqpoint{0.592196in}{1.032116in}}%
\pgfpathlineto{\pgfqpoint{0.608546in}{1.015186in}}%
\pgfpathlineto{\pgfqpoint{0.617769in}{1.012388in}}%
\pgfpathlineto{\pgfqpoint{0.626947in}{1.009861in}}%
\pgfpathlineto{\pgfqpoint{0.610639in}{1.027172in}}%
\pgfpathlineto{\pgfqpoint{0.594340in}{1.044738in}}%
\pgfpathlineto{\pgfqpoint{0.585119in}{1.046893in}}%
\pgfpathlineto{\pgfqpoint{0.575854in}{1.049316in}}%
\pgfpathclose%
\pgfusepath{fill}%
\end{pgfscope}%
\begin{pgfscope}%
\pgfpathrectangle{\pgfqpoint{0.407592in}{0.222904in}}{\pgfqpoint{1.411163in}{1.411163in}}%
\pgfusepath{clip}%
\pgfsetbuttcap%
\pgfsetroundjoin%
\definecolor{currentfill}{rgb}{0.556863,0.654902,0.847059}%
\pgfsetfillcolor{currentfill}%
\pgfsetfillopacity{0.130000}%
\pgfsetlinewidth{0.000000pt}%
\definecolor{currentstroke}{rgb}{0.000000,0.000000,0.000000}%
\pgfsetstrokecolor{currentstroke}%
\pgfsetdash{}{0pt}%
\pgfpathmoveto{\pgfqpoint{1.485330in}{0.952082in}}%
\pgfpathlineto{\pgfqpoint{1.502695in}{0.963296in}}%
\pgfpathlineto{\pgfqpoint{1.520136in}{0.974725in}}%
\pgfpathlineto{\pgfqpoint{1.527795in}{0.983937in}}%
\pgfpathlineto{\pgfqpoint{1.535443in}{0.993491in}}%
\pgfpathlineto{\pgfqpoint{1.518023in}{0.982381in}}%
\pgfpathlineto{\pgfqpoint{1.500680in}{0.971504in}}%
\pgfpathlineto{\pgfqpoint{1.493011in}{0.961613in}}%
\pgfpathlineto{\pgfqpoint{1.485330in}{0.952082in}}%
\pgfpathclose%
\pgfusepath{fill}%
\end{pgfscope}%
\begin{pgfscope}%
\pgfpathrectangle{\pgfqpoint{0.407592in}{0.222904in}}{\pgfqpoint{1.411163in}{1.411163in}}%
\pgfusepath{clip}%
\pgfsetbuttcap%
\pgfsetroundjoin%
\definecolor{currentfill}{rgb}{0.556863,0.654902,0.847059}%
\pgfsetfillcolor{currentfill}%
\pgfsetfillopacity{0.130000}%
\pgfsetlinewidth{0.000000pt}%
\definecolor{currentstroke}{rgb}{0.000000,0.000000,0.000000}%
\pgfsetstrokecolor{currentstroke}%
\pgfsetdash{}{0pt}%
\pgfpathmoveto{\pgfqpoint{1.205258in}{0.758838in}}%
\pgfpathlineto{\pgfqpoint{1.222031in}{0.767388in}}%
\pgfpathlineto{\pgfqpoint{1.238873in}{0.776510in}}%
\pgfpathlineto{\pgfqpoint{1.246871in}{0.779121in}}%
\pgfpathlineto{\pgfqpoint{1.254845in}{0.782394in}}%
\pgfpathlineto{\pgfqpoint{1.238027in}{0.773003in}}%
\pgfpathlineto{\pgfqpoint{1.221279in}{0.764160in}}%
\pgfpathlineto{\pgfqpoint{1.213280in}{0.761131in}}%
\pgfpathlineto{\pgfqpoint{1.205258in}{0.758838in}}%
\pgfpathclose%
\pgfusepath{fill}%
\end{pgfscope}%
\begin{pgfscope}%
\pgfpathrectangle{\pgfqpoint{0.407592in}{0.222904in}}{\pgfqpoint{1.411163in}{1.411163in}}%
\pgfusepath{clip}%
\pgfsetbuttcap%
\pgfsetroundjoin%
\definecolor{currentfill}{rgb}{0.556863,0.654902,0.847059}%
\pgfsetfillcolor{currentfill}%
\pgfsetfillopacity{0.130000}%
\pgfsetlinewidth{0.000000pt}%
\definecolor{currentstroke}{rgb}{0.000000,0.000000,0.000000}%
\pgfsetstrokecolor{currentstroke}%
\pgfsetdash{}{0pt}%
\pgfpathmoveto{\pgfqpoint{1.056620in}{0.735487in}}%
\pgfpathlineto{\pgfqpoint{1.073141in}{0.736568in}}%
\pgfpathlineto{\pgfqpoint{1.089710in}{0.738908in}}%
\pgfpathlineto{\pgfqpoint{1.097979in}{0.736318in}}%
\pgfpathlineto{\pgfqpoint{1.106213in}{0.734261in}}%
\pgfpathlineto{\pgfqpoint{1.089685in}{0.731394in}}%
\pgfpathlineto{\pgfqpoint{1.073207in}{0.729855in}}%
\pgfpathlineto{\pgfqpoint{1.064932in}{0.732409in}}%
\pgfpathlineto{\pgfqpoint{1.056620in}{0.735487in}}%
\pgfpathclose%
\pgfusepath{fill}%
\end{pgfscope}%
\begin{pgfscope}%
\pgfpathrectangle{\pgfqpoint{0.407592in}{0.222904in}}{\pgfqpoint{1.411163in}{1.411163in}}%
\pgfusepath{clip}%
\pgfsetbuttcap%
\pgfsetroundjoin%
\definecolor{currentfill}{rgb}{0.556863,0.654902,0.847059}%
\pgfsetfillcolor{currentfill}%
\pgfsetfillopacity{0.130000}%
\pgfsetlinewidth{0.000000pt}%
\definecolor{currentstroke}{rgb}{0.000000,0.000000,0.000000}%
\pgfsetstrokecolor{currentstroke}%
\pgfsetdash{}{0pt}%
\pgfpathmoveto{\pgfqpoint{0.724979in}{0.912959in}}%
\pgfpathlineto{\pgfqpoint{0.741349in}{0.898310in}}%
\pgfpathlineto{\pgfqpoint{0.757730in}{0.884208in}}%
\pgfpathlineto{\pgfqpoint{0.766691in}{0.880144in}}%
\pgfpathlineto{\pgfqpoint{0.775610in}{0.876361in}}%
\pgfpathlineto{\pgfqpoint{0.759273in}{0.890966in}}%
\pgfpathlineto{\pgfqpoint{0.742947in}{0.906110in}}%
\pgfpathlineto{\pgfqpoint{0.733984in}{0.909389in}}%
\pgfpathlineto{\pgfqpoint{0.724979in}{0.912959in}}%
\pgfpathclose%
\pgfusepath{fill}%
\end{pgfscope}%
\begin{pgfscope}%
\pgfpathrectangle{\pgfqpoint{0.407592in}{0.222904in}}{\pgfqpoint{1.411163in}{1.411163in}}%
\pgfusepath{clip}%
\pgfsetbuttcap%
\pgfsetroundjoin%
\definecolor{currentfill}{rgb}{0.556863,0.654902,0.847059}%
\pgfsetfillcolor{currentfill}%
\pgfsetfillopacity{0.130000}%
\pgfsetlinewidth{0.000000pt}%
\definecolor{currentstroke}{rgb}{0.000000,0.000000,0.000000}%
\pgfsetstrokecolor{currentstroke}%
\pgfsetdash{}{0pt}%
\pgfpathmoveto{\pgfqpoint{0.924077in}{0.776122in}}%
\pgfpathlineto{\pgfqpoint{0.940494in}{0.768983in}}%
\pgfpathlineto{\pgfqpoint{0.956935in}{0.763069in}}%
\pgfpathlineto{\pgfqpoint{0.965494in}{0.758145in}}%
\pgfpathlineto{\pgfqpoint{0.974014in}{0.753530in}}%
\pgfpathlineto{\pgfqpoint{0.957622in}{0.759583in}}%
\pgfpathlineto{\pgfqpoint{0.941254in}{0.766967in}}%
\pgfpathlineto{\pgfqpoint{0.932685in}{0.771385in}}%
\pgfpathlineto{\pgfqpoint{0.924077in}{0.776122in}}%
\pgfpathclose%
\pgfusepath{fill}%
\end{pgfscope}%
\begin{pgfscope}%
\pgfpathrectangle{\pgfqpoint{0.407592in}{0.222904in}}{\pgfqpoint{1.411163in}{1.411163in}}%
\pgfusepath{clip}%
\pgfsetbuttcap%
\pgfsetroundjoin%
\definecolor{currentfill}{rgb}{0.556863,0.654902,0.847059}%
\pgfsetfillcolor{currentfill}%
\pgfsetfillopacity{0.130000}%
\pgfsetlinewidth{0.000000pt}%
\definecolor{currentstroke}{rgb}{0.000000,0.000000,0.000000}%
\pgfsetstrokecolor{currentstroke}%
\pgfsetdash{}{0pt}%
\pgfpathmoveto{\pgfqpoint{1.600895in}{1.056387in}}%
\pgfpathlineto{\pgfqpoint{1.618503in}{1.067480in}}%
\pgfpathlineto{\pgfqpoint{1.636188in}{1.078777in}}%
\pgfpathlineto{\pgfqpoint{1.643732in}{1.089165in}}%
\pgfpathlineto{\pgfqpoint{1.651264in}{1.099804in}}%
\pgfpathlineto{\pgfqpoint{1.633602in}{1.088830in}}%
\pgfpathlineto{\pgfqpoint{1.616016in}{1.078076in}}%
\pgfpathlineto{\pgfqpoint{1.608461in}{1.067102in}}%
\pgfpathlineto{\pgfqpoint{1.600895in}{1.056387in}}%
\pgfpathclose%
\pgfusepath{fill}%
\end{pgfscope}%
\begin{pgfscope}%
\pgfpathrectangle{\pgfqpoint{0.407592in}{0.222904in}}{\pgfqpoint{1.411163in}{1.411163in}}%
\pgfusepath{clip}%
\pgfsetbuttcap%
\pgfsetroundjoin%
\definecolor{currentfill}{rgb}{0.556863,0.654902,0.847059}%
\pgfsetfillcolor{currentfill}%
\pgfsetfillopacity{0.130000}%
\pgfsetlinewidth{0.000000pt}%
\definecolor{currentstroke}{rgb}{0.000000,0.000000,0.000000}%
\pgfsetstrokecolor{currentstroke}%
\pgfsetdash{}{0pt}%
\pgfpathmoveto{\pgfqpoint{1.254845in}{0.782394in}}%
\pgfpathlineto{\pgfqpoint{1.271735in}{0.792254in}}%
\pgfpathlineto{\pgfqpoint{1.288698in}{0.802519in}}%
\pgfpathlineto{\pgfqpoint{1.296628in}{0.806605in}}%
\pgfpathlineto{\pgfqpoint{1.304539in}{0.811307in}}%
\pgfpathlineto{\pgfqpoint{1.287598in}{0.800949in}}%
\pgfpathlineto{\pgfqpoint{1.270730in}{0.790984in}}%
\pgfpathlineto{\pgfqpoint{1.262798in}{0.786346in}}%
\pgfpathlineto{\pgfqpoint{1.254845in}{0.782394in}}%
\pgfpathclose%
\pgfusepath{fill}%
\end{pgfscope}%
\begin{pgfscope}%
\pgfpathrectangle{\pgfqpoint{0.407592in}{0.222904in}}{\pgfqpoint{1.411163in}{1.411163in}}%
\pgfusepath{clip}%
\pgfsetbuttcap%
\pgfsetroundjoin%
\definecolor{currentfill}{rgb}{0.556863,0.654902,0.847059}%
\pgfsetfillcolor{currentfill}%
\pgfsetfillopacity{0.130000}%
\pgfsetlinewidth{0.000000pt}%
\definecolor{currentstroke}{rgb}{0.000000,0.000000,0.000000}%
\pgfsetstrokecolor{currentstroke}%
\pgfsetdash{}{0pt}%
\pgfpathmoveto{\pgfqpoint{0.841063in}{0.824538in}}%
\pgfpathlineto{\pgfqpoint{0.857457in}{0.813559in}}%
\pgfpathlineto{\pgfqpoint{0.873867in}{0.803505in}}%
\pgfpathlineto{\pgfqpoint{0.882605in}{0.798555in}}%
\pgfpathlineto{\pgfqpoint{0.891302in}{0.793879in}}%
\pgfpathlineto{\pgfqpoint{0.874940in}{0.804363in}}%
\pgfpathlineto{\pgfqpoint{0.858593in}{0.815816in}}%
\pgfpathlineto{\pgfqpoint{0.849849in}{0.820029in}}%
\pgfpathlineto{\pgfqpoint{0.841063in}{0.824538in}}%
\pgfpathclose%
\pgfusepath{fill}%
\end{pgfscope}%
\begin{pgfscope}%
\pgfpathrectangle{\pgfqpoint{0.407592in}{0.222904in}}{\pgfqpoint{1.411163in}{1.411163in}}%
\pgfusepath{clip}%
\pgfsetbuttcap%
\pgfsetroundjoin%
\definecolor{currentfill}{rgb}{0.556863,0.654902,0.847059}%
\pgfsetfillcolor{currentfill}%
\pgfsetfillopacity{0.130000}%
\pgfsetlinewidth{0.000000pt}%
\definecolor{currentstroke}{rgb}{0.000000,0.000000,0.000000}%
\pgfsetstrokecolor{currentstroke}%
\pgfsetdash{}{0pt}%
\pgfpathmoveto{\pgfqpoint{0.459061in}{1.160764in}}%
\pgfpathlineto{\pgfqpoint{0.475395in}{1.142637in}}%
\pgfpathlineto{\pgfqpoint{0.491736in}{1.124683in}}%
\pgfpathlineto{\pgfqpoint{0.501176in}{1.122417in}}%
\pgfpathlineto{\pgfqpoint{0.510571in}{1.120393in}}%
\pgfpathlineto{\pgfqpoint{0.494272in}{1.138639in}}%
\pgfpathlineto{\pgfqpoint{0.477981in}{1.157044in}}%
\pgfpathlineto{\pgfqpoint{0.468544in}{1.158785in}}%
\pgfpathlineto{\pgfqpoint{0.459061in}{1.160764in}}%
\pgfpathclose%
\pgfusepath{fill}%
\end{pgfscope}%
\begin{pgfscope}%
\pgfpathrectangle{\pgfqpoint{0.407592in}{0.222904in}}{\pgfqpoint{1.411163in}{1.411163in}}%
\pgfusepath{clip}%
\pgfsetbuttcap%
\pgfsetroundjoin%
\definecolor{currentfill}{rgb}{0.556863,0.654902,0.847059}%
\pgfsetfillcolor{currentfill}%
\pgfsetfillopacity{0.130000}%
\pgfsetlinewidth{0.000000pt}%
\definecolor{currentstroke}{rgb}{0.000000,0.000000,0.000000}%
\pgfsetstrokecolor{currentstroke}%
\pgfsetdash{}{0pt}%
\pgfpathmoveto{\pgfqpoint{0.608546in}{1.015186in}}%
\pgfpathlineto{\pgfqpoint{0.624906in}{0.998553in}}%
\pgfpathlineto{\pgfqpoint{0.641274in}{0.982244in}}%
\pgfpathlineto{\pgfqpoint{0.650453in}{0.979039in}}%
\pgfpathlineto{\pgfqpoint{0.659589in}{0.976103in}}%
\pgfpathlineto{\pgfqpoint{0.643264in}{0.992828in}}%
\pgfpathlineto{\pgfqpoint{0.626947in}{1.009861in}}%
\pgfpathlineto{\pgfqpoint{0.617769in}{1.012388in}}%
\pgfpathlineto{\pgfqpoint{0.608546in}{1.015186in}}%
\pgfpathclose%
\pgfusepath{fill}%
\end{pgfscope}%
\begin{pgfscope}%
\pgfpathrectangle{\pgfqpoint{0.407592in}{0.222904in}}{\pgfqpoint{1.411163in}{1.411163in}}%
\pgfusepath{clip}%
\pgfsetbuttcap%
\pgfsetroundjoin%
\definecolor{currentfill}{rgb}{0.556863,0.654902,0.847059}%
\pgfsetfillcolor{currentfill}%
\pgfsetfillopacity{0.130000}%
\pgfsetlinewidth{0.000000pt}%
\definecolor{currentstroke}{rgb}{0.000000,0.000000,0.000000}%
\pgfsetstrokecolor{currentstroke}%
\pgfsetdash{}{0pt}%
\pgfpathmoveto{\pgfqpoint{1.419893in}{0.894947in}}%
\pgfpathlineto{\pgfqpoint{1.437147in}{0.906288in}}%
\pgfpathlineto{\pgfqpoint{1.454476in}{0.917843in}}%
\pgfpathlineto{\pgfqpoint{1.462210in}{0.925794in}}%
\pgfpathlineto{\pgfqpoint{1.469930in}{0.934159in}}%
\pgfpathlineto{\pgfqpoint{1.452621in}{0.922862in}}%
\pgfpathlineto{\pgfqpoint{1.435389in}{0.911799in}}%
\pgfpathlineto{\pgfqpoint{1.427648in}{0.903150in}}%
\pgfpathlineto{\pgfqpoint{1.419893in}{0.894947in}}%
\pgfpathclose%
\pgfusepath{fill}%
\end{pgfscope}%
\begin{pgfscope}%
\pgfpathrectangle{\pgfqpoint{0.407592in}{0.222904in}}{\pgfqpoint{1.411163in}{1.411163in}}%
\pgfusepath{clip}%
\pgfsetbuttcap%
\pgfsetroundjoin%
\definecolor{currentfill}{rgb}{0.556863,0.654902,0.847059}%
\pgfsetfillcolor{currentfill}%
\pgfsetfillopacity{0.130000}%
\pgfsetlinewidth{0.000000pt}%
\definecolor{currentstroke}{rgb}{0.000000,0.000000,0.000000}%
\pgfsetstrokecolor{currentstroke}%
\pgfsetdash{}{0pt}%
\pgfpathmoveto{\pgfqpoint{1.782194in}{1.233709in}}%
\pgfpathlineto{\pgfqpoint{1.800156in}{1.244470in}}%
\pgfpathlineto{\pgfqpoint{1.818197in}{1.255429in}}%
\pgfpathlineto{\pgfqpoint{1.825562in}{1.266988in}}%
\pgfpathlineto{\pgfqpoint{1.807533in}{1.256172in}}%
\pgfpathlineto{\pgfqpoint{1.789583in}{1.245560in}}%
\pgfusepath{fill}%
\end{pgfscope}%
\begin{pgfscope}%
\pgfpathrectangle{\pgfqpoint{0.407592in}{0.222904in}}{\pgfqpoint{1.411163in}{1.411163in}}%
\pgfusepath{clip}%
\pgfsetbuttcap%
\pgfsetroundjoin%
\definecolor{currentfill}{rgb}{0.556863,0.654902,0.847059}%
\pgfsetfillcolor{currentfill}%
\pgfsetfillopacity{0.130000}%
\pgfsetlinewidth{0.000000pt}%
\definecolor{currentstroke}{rgb}{0.000000,0.000000,0.000000}%
\pgfsetstrokecolor{currentstroke}%
\pgfsetdash{}{0pt}%
\pgfpathmoveto{\pgfqpoint{1.006881in}{0.745540in}}%
\pgfpathlineto{\pgfqpoint{1.023365in}{0.743601in}}%
\pgfpathlineto{\pgfqpoint{1.039888in}{0.742986in}}%
\pgfpathlineto{\pgfqpoint{1.048272in}{0.739031in}}%
\pgfpathlineto{\pgfqpoint{1.056620in}{0.735487in}}%
\pgfpathlineto{\pgfqpoint{1.040144in}{0.735765in}}%
\pgfpathlineto{\pgfqpoint{1.023708in}{0.737478in}}%
\pgfpathlineto{\pgfqpoint{1.015313in}{0.741310in}}%
\pgfpathlineto{\pgfqpoint{1.006881in}{0.745540in}}%
\pgfpathclose%
\pgfusepath{fill}%
\end{pgfscope}%
\begin{pgfscope}%
\pgfpathrectangle{\pgfqpoint{0.407592in}{0.222904in}}{\pgfqpoint{1.411163in}{1.411163in}}%
\pgfusepath{clip}%
\pgfsetbuttcap%
\pgfsetroundjoin%
\definecolor{currentfill}{rgb}{0.556863,0.654902,0.847059}%
\pgfsetfillcolor{currentfill}%
\pgfsetfillopacity{0.130000}%
\pgfsetlinewidth{0.000000pt}%
\definecolor{currentstroke}{rgb}{0.000000,0.000000,0.000000}%
\pgfsetstrokecolor{currentstroke}%
\pgfsetdash{}{0pt}%
\pgfpathmoveto{\pgfqpoint{1.716752in}{1.165971in}}%
\pgfpathlineto{\pgfqpoint{1.734604in}{1.176930in}}%
\pgfpathlineto{\pgfqpoint{1.752534in}{1.188087in}}%
\pgfpathlineto{\pgfqpoint{1.759965in}{1.199211in}}%
\pgfpathlineto{\pgfqpoint{1.767385in}{1.210528in}}%
\pgfpathlineto{\pgfqpoint{1.749479in}{1.199672in}}%
\pgfpathlineto{\pgfqpoint{1.731650in}{1.189024in}}%
\pgfpathlineto{\pgfqpoint{1.724207in}{1.177400in}}%
\pgfpathlineto{\pgfqpoint{1.716752in}{1.165971in}}%
\pgfpathclose%
\pgfusepath{fill}%
\end{pgfscope}%
\begin{pgfscope}%
\pgfpathrectangle{\pgfqpoint{0.407592in}{0.222904in}}{\pgfqpoint{1.411163in}{1.411163in}}%
\pgfusepath{clip}%
\pgfsetbuttcap%
\pgfsetroundjoin%
\definecolor{currentfill}{rgb}{0.556863,0.654902,0.847059}%
\pgfsetfillcolor{currentfill}%
\pgfsetfillopacity{0.130000}%
\pgfsetlinewidth{0.000000pt}%
\definecolor{currentstroke}{rgb}{0.000000,0.000000,0.000000}%
\pgfsetstrokecolor{currentstroke}%
\pgfsetdash{}{0pt}%
\pgfpathmoveto{\pgfqpoint{1.535443in}{0.993491in}}%
\pgfpathlineto{\pgfqpoint{1.552939in}{1.004816in}}%
\pgfpathlineto{\pgfqpoint{1.570514in}{1.016342in}}%
\pgfpathlineto{\pgfqpoint{1.578127in}{1.025909in}}%
\pgfpathlineto{\pgfqpoint{1.585728in}{1.035780in}}%
\pgfpathlineto{\pgfqpoint{1.568176in}{1.024568in}}%
\pgfpathlineto{\pgfqpoint{1.550701in}{1.013574in}}%
\pgfpathlineto{\pgfqpoint{1.543078in}{1.003374in}}%
\pgfpathlineto{\pgfqpoint{1.535443in}{0.993491in}}%
\pgfpathclose%
\pgfusepath{fill}%
\end{pgfscope}%
\begin{pgfscope}%
\pgfpathrectangle{\pgfqpoint{0.407592in}{0.222904in}}{\pgfqpoint{1.411163in}{1.411163in}}%
\pgfusepath{clip}%
\pgfsetbuttcap%
\pgfsetroundjoin%
\definecolor{currentfill}{rgb}{0.556863,0.654902,0.847059}%
\pgfsetfillcolor{currentfill}%
\pgfsetfillopacity{0.130000}%
\pgfsetlinewidth{0.000000pt}%
\definecolor{currentstroke}{rgb}{0.000000,0.000000,0.000000}%
\pgfsetstrokecolor{currentstroke}%
\pgfsetdash{}{0pt}%
\pgfpathmoveto{\pgfqpoint{1.304539in}{0.811307in}}%
\pgfpathlineto{\pgfqpoint{1.321553in}{0.822006in}}%
\pgfpathlineto{\pgfqpoint{1.338642in}{0.833004in}}%
\pgfpathlineto{\pgfqpoint{1.346512in}{0.838311in}}%
\pgfpathlineto{\pgfqpoint{1.354365in}{0.844172in}}%
\pgfpathlineto{\pgfqpoint{1.337297in}{0.833216in}}%
\pgfpathlineto{\pgfqpoint{1.320303in}{0.822560in}}%
\pgfpathlineto{\pgfqpoint{1.312430in}{0.816626in}}%
\pgfpathlineto{\pgfqpoint{1.304539in}{0.811307in}}%
\pgfpathclose%
\pgfusepath{fill}%
\end{pgfscope}%
\begin{pgfscope}%
\pgfpathrectangle{\pgfqpoint{0.407592in}{0.222904in}}{\pgfqpoint{1.411163in}{1.411163in}}%
\pgfusepath{clip}%
\pgfsetbuttcap%
\pgfsetroundjoin%
\definecolor{currentfill}{rgb}{0.556863,0.654902,0.847059}%
\pgfsetfillcolor{currentfill}%
\pgfsetfillopacity{0.130000}%
\pgfsetlinewidth{0.000000pt}%
\definecolor{currentstroke}{rgb}{0.000000,0.000000,0.000000}%
\pgfsetstrokecolor{currentstroke}%
\pgfsetdash{}{0pt}%
\pgfpathmoveto{\pgfqpoint{0.491736in}{1.124683in}}%
\pgfpathlineto{\pgfqpoint{0.508087in}{1.106918in}}%
\pgfpathlineto{\pgfqpoint{0.524446in}{1.089358in}}%
\pgfpathlineto{\pgfqpoint{0.533843in}{1.086776in}}%
\pgfpathlineto{\pgfqpoint{0.543195in}{1.084439in}}%
\pgfpathlineto{\pgfqpoint{0.526879in}{1.102321in}}%
\pgfpathlineto{\pgfqpoint{0.510571in}{1.120393in}}%
\pgfpathlineto{\pgfqpoint{0.501176in}{1.122417in}}%
\pgfpathlineto{\pgfqpoint{0.491736in}{1.124683in}}%
\pgfpathclose%
\pgfusepath{fill}%
\end{pgfscope}%
\begin{pgfscope}%
\pgfpathrectangle{\pgfqpoint{0.407592in}{0.222904in}}{\pgfqpoint{1.411163in}{1.411163in}}%
\pgfusepath{clip}%
\pgfsetbuttcap%
\pgfsetroundjoin%
\definecolor{currentfill}{rgb}{0.556863,0.654902,0.847059}%
\pgfsetfillcolor{currentfill}%
\pgfsetfillopacity{0.130000}%
\pgfsetlinewidth{0.000000pt}%
\definecolor{currentstroke}{rgb}{0.000000,0.000000,0.000000}%
\pgfsetstrokecolor{currentstroke}%
\pgfsetdash{}{0pt}%
\pgfpathmoveto{\pgfqpoint{0.757730in}{0.884208in}}%
\pgfpathlineto{\pgfqpoint{0.774121in}{0.870705in}}%
\pgfpathlineto{\pgfqpoint{0.790524in}{0.857854in}}%
\pgfpathlineto{\pgfqpoint{0.799440in}{0.853295in}}%
\pgfpathlineto{\pgfqpoint{0.808313in}{0.849002in}}%
\pgfpathlineto{\pgfqpoint{0.791956in}{0.862353in}}%
\pgfpathlineto{\pgfqpoint{0.775610in}{0.876361in}}%
\pgfpathlineto{\pgfqpoint{0.766691in}{0.880144in}}%
\pgfpathlineto{\pgfqpoint{0.757730in}{0.884208in}}%
\pgfpathclose%
\pgfusepath{fill}%
\end{pgfscope}%
\begin{pgfscope}%
\pgfpathrectangle{\pgfqpoint{0.407592in}{0.222904in}}{\pgfqpoint{1.411163in}{1.411163in}}%
\pgfusepath{clip}%
\pgfsetbuttcap%
\pgfsetroundjoin%
\definecolor{currentfill}{rgb}{0.556863,0.654902,0.847059}%
\pgfsetfillcolor{currentfill}%
\pgfsetfillopacity{0.130000}%
\pgfsetlinewidth{0.000000pt}%
\definecolor{currentstroke}{rgb}{0.000000,0.000000,0.000000}%
\pgfsetstrokecolor{currentstroke}%
\pgfsetdash{}{0pt}%
\pgfpathmoveto{\pgfqpoint{0.641274in}{0.982244in}}%
\pgfpathlineto{\pgfqpoint{0.657651in}{0.966290in}}%
\pgfpathlineto{\pgfqpoint{0.674037in}{0.950724in}}%
\pgfpathlineto{\pgfqpoint{0.683173in}{0.947083in}}%
\pgfpathlineto{\pgfqpoint{0.692266in}{0.943706in}}%
\pgfpathlineto{\pgfqpoint{0.675923in}{0.959718in}}%
\pgfpathlineto{\pgfqpoint{0.659589in}{0.976103in}}%
\pgfpathlineto{\pgfqpoint{0.650453in}{0.979039in}}%
\pgfpathlineto{\pgfqpoint{0.641274in}{0.982244in}}%
\pgfpathclose%
\pgfusepath{fill}%
\end{pgfscope}%
\begin{pgfscope}%
\pgfpathrectangle{\pgfqpoint{0.407592in}{0.222904in}}{\pgfqpoint{1.411163in}{1.411163in}}%
\pgfusepath{clip}%
\pgfsetbuttcap%
\pgfsetroundjoin%
\definecolor{currentfill}{rgb}{0.556863,0.654902,0.847059}%
\pgfsetfillcolor{currentfill}%
\pgfsetfillopacity{0.130000}%
\pgfsetlinewidth{0.000000pt}%
\definecolor{currentstroke}{rgb}{0.000000,0.000000,0.000000}%
\pgfsetstrokecolor{currentstroke}%
\pgfsetdash{}{0pt}%
\pgfpathmoveto{\pgfqpoint{1.651264in}{1.099804in}}%
\pgfpathlineto{\pgfqpoint{1.669006in}{1.110985in}}%
\pgfpathlineto{\pgfqpoint{1.686825in}{1.122362in}}%
\pgfpathlineto{\pgfqpoint{1.694323in}{1.132931in}}%
\pgfpathlineto{\pgfqpoint{1.701811in}{1.143728in}}%
\pgfpathlineto{\pgfqpoint{1.684014in}{1.132660in}}%
\pgfpathlineto{\pgfqpoint{1.666296in}{1.121801in}}%
\pgfpathlineto{\pgfqpoint{1.658786in}{1.110686in}}%
\pgfpathlineto{\pgfqpoint{1.651264in}{1.099804in}}%
\pgfpathclose%
\pgfusepath{fill}%
\end{pgfscope}%
\begin{pgfscope}%
\pgfpathrectangle{\pgfqpoint{0.407592in}{0.222904in}}{\pgfqpoint{1.411163in}{1.411163in}}%
\pgfusepath{clip}%
\pgfsetbuttcap%
\pgfsetroundjoin%
\definecolor{currentfill}{rgb}{0.556863,0.654902,0.847059}%
\pgfsetfillcolor{currentfill}%
\pgfsetfillopacity{0.130000}%
\pgfsetlinewidth{0.000000pt}%
\definecolor{currentstroke}{rgb}{0.000000,0.000000,0.000000}%
\pgfsetstrokecolor{currentstroke}%
\pgfsetdash{}{0pt}%
\pgfpathmoveto{\pgfqpoint{1.469930in}{0.934159in}}%
\pgfpathlineto{\pgfqpoint{1.487315in}{0.945669in}}%
\pgfpathlineto{\pgfqpoint{1.504778in}{0.957375in}}%
\pgfpathlineto{\pgfqpoint{1.512463in}{0.965867in}}%
\pgfpathlineto{\pgfqpoint{1.520136in}{0.974725in}}%
\pgfpathlineto{\pgfqpoint{1.502695in}{0.963296in}}%
\pgfpathlineto{\pgfqpoint{1.485330in}{0.952082in}}%
\pgfpathlineto{\pgfqpoint{1.477637in}{0.942926in}}%
\pgfpathlineto{\pgfqpoint{1.469930in}{0.934159in}}%
\pgfpathclose%
\pgfusepath{fill}%
\end{pgfscope}%
\begin{pgfscope}%
\pgfpathrectangle{\pgfqpoint{0.407592in}{0.222904in}}{\pgfqpoint{1.411163in}{1.411163in}}%
\pgfusepath{clip}%
\pgfsetbuttcap%
\pgfsetroundjoin%
\definecolor{currentfill}{rgb}{0.556863,0.654902,0.847059}%
\pgfsetfillcolor{currentfill}%
\pgfsetfillopacity{0.130000}%
\pgfsetlinewidth{0.000000pt}%
\definecolor{currentstroke}{rgb}{0.000000,0.000000,0.000000}%
\pgfsetstrokecolor{currentstroke}%
\pgfsetdash{}{0pt}%
\pgfpathmoveto{\pgfqpoint{0.873867in}{0.803505in}}%
\pgfpathlineto{\pgfqpoint{0.890293in}{0.794441in}}%
\pgfpathlineto{\pgfqpoint{0.906739in}{0.786422in}}%
\pgfpathlineto{\pgfqpoint{0.915428in}{0.781145in}}%
\pgfpathlineto{\pgfqpoint{0.924077in}{0.776122in}}%
\pgfpathlineto{\pgfqpoint{0.907681in}{0.784441in}}%
\pgfpathlineto{\pgfqpoint{0.891302in}{0.793879in}}%
\pgfpathlineto{\pgfqpoint{0.882605in}{0.798555in}}%
\pgfpathlineto{\pgfqpoint{0.873867in}{0.803505in}}%
\pgfpathclose%
\pgfusepath{fill}%
\end{pgfscope}%
\begin{pgfscope}%
\pgfpathrectangle{\pgfqpoint{0.407592in}{0.222904in}}{\pgfqpoint{1.411163in}{1.411163in}}%
\pgfusepath{clip}%
\pgfsetbuttcap%
\pgfsetroundjoin%
\definecolor{currentfill}{rgb}{0.556863,0.654902,0.847059}%
\pgfsetfillcolor{currentfill}%
\pgfsetfillopacity{0.130000}%
\pgfsetlinewidth{0.000000pt}%
\definecolor{currentstroke}{rgb}{0.000000,0.000000,0.000000}%
\pgfsetstrokecolor{currentstroke}%
\pgfsetdash{}{0pt}%
\pgfpathmoveto{\pgfqpoint{1.354365in}{0.844172in}}%
\pgfpathlineto{\pgfqpoint{1.371508in}{0.855394in}}%
\pgfpathlineto{\pgfqpoint{1.388726in}{0.866852in}}%
\pgfpathlineto{\pgfqpoint{1.396542in}{0.873147in}}%
\pgfpathlineto{\pgfqpoint{1.404341in}{0.879934in}}%
\pgfpathlineto{\pgfqpoint{1.387143in}{0.868610in}}%
\pgfpathlineto{\pgfqpoint{1.370020in}{0.857532in}}%
\pgfpathlineto{\pgfqpoint{1.362201in}{0.850582in}}%
\pgfpathlineto{\pgfqpoint{1.354365in}{0.844172in}}%
\pgfpathclose%
\pgfusepath{fill}%
\end{pgfscope}%
\begin{pgfscope}%
\pgfpathrectangle{\pgfqpoint{0.407592in}{0.222904in}}{\pgfqpoint{1.411163in}{1.411163in}}%
\pgfusepath{clip}%
\pgfsetbuttcap%
\pgfsetroundjoin%
\definecolor{currentfill}{rgb}{0.556863,0.654902,0.847059}%
\pgfsetfillcolor{currentfill}%
\pgfsetfillopacity{0.130000}%
\pgfsetlinewidth{0.000000pt}%
\definecolor{currentstroke}{rgb}{0.000000,0.000000,0.000000}%
\pgfsetstrokecolor{currentstroke}%
\pgfsetdash{}{0pt}%
\pgfpathmoveto{\pgfqpoint{0.524446in}{1.089358in}}%
\pgfpathlineto{\pgfqpoint{0.540813in}{1.072020in}}%
\pgfpathlineto{\pgfqpoint{0.557189in}{1.054924in}}%
\pgfpathlineto{\pgfqpoint{0.566544in}{1.051996in}}%
\pgfpathlineto{\pgfqpoint{0.575854in}{1.049316in}}%
\pgfpathlineto{\pgfqpoint{0.559520in}{1.066764in}}%
\pgfpathlineto{\pgfqpoint{0.543195in}{1.084439in}}%
\pgfpathlineto{\pgfqpoint{0.533843in}{1.086776in}}%
\pgfpathlineto{\pgfqpoint{0.524446in}{1.089358in}}%
\pgfpathclose%
\pgfusepath{fill}%
\end{pgfscope}%
\begin{pgfscope}%
\pgfpathrectangle{\pgfqpoint{0.407592in}{0.222904in}}{\pgfqpoint{1.411163in}{1.411163in}}%
\pgfusepath{clip}%
\pgfsetbuttcap%
\pgfsetroundjoin%
\definecolor{currentfill}{rgb}{0.556863,0.654902,0.847059}%
\pgfsetfillcolor{currentfill}%
\pgfsetfillopacity{0.130000}%
\pgfsetlinewidth{0.000000pt}%
\definecolor{currentstroke}{rgb}{0.000000,0.000000,0.000000}%
\pgfsetstrokecolor{currentstroke}%
\pgfsetdash{}{0pt}%
\pgfpathmoveto{\pgfqpoint{0.956935in}{0.763069in}}%
\pgfpathlineto{\pgfqpoint{0.973403in}{0.758409in}}%
\pgfpathlineto{\pgfqpoint{0.989901in}{0.755014in}}%
\pgfpathlineto{\pgfqpoint{0.998411in}{0.750122in}}%
\pgfpathlineto{\pgfqpoint{1.006881in}{0.745540in}}%
\pgfpathlineto{\pgfqpoint{0.990432in}{0.748845in}}%
\pgfpathlineto{\pgfqpoint{0.974014in}{0.753530in}}%
\pgfpathlineto{\pgfqpoint{0.965494in}{0.758145in}}%
\pgfpathlineto{\pgfqpoint{0.956935in}{0.763069in}}%
\pgfpathclose%
\pgfusepath{fill}%
\end{pgfscope}%
\begin{pgfscope}%
\pgfpathrectangle{\pgfqpoint{0.407592in}{0.222904in}}{\pgfqpoint{1.411163in}{1.411163in}}%
\pgfusepath{clip}%
\pgfsetbuttcap%
\pgfsetroundjoin%
\definecolor{currentfill}{rgb}{0.556863,0.654902,0.847059}%
\pgfsetfillcolor{currentfill}%
\pgfsetfillopacity{0.130000}%
\pgfsetlinewidth{0.000000pt}%
\definecolor{currentstroke}{rgb}{0.000000,0.000000,0.000000}%
\pgfsetstrokecolor{currentstroke}%
\pgfsetdash{}{0pt}%
\pgfpathmoveto{\pgfqpoint{1.767385in}{1.210528in}}%
\pgfpathlineto{\pgfqpoint{1.785372in}{1.221583in}}%
\pgfpathlineto{\pgfqpoint{1.803438in}{1.232829in}}%
\pgfpathlineto{\pgfqpoint{1.810823in}{1.244041in}}%
\pgfpathlineto{\pgfqpoint{1.818197in}{1.255429in}}%
\pgfpathlineto{\pgfqpoint{1.800156in}{1.244470in}}%
\pgfpathlineto{\pgfqpoint{1.782194in}{1.233709in}}%
\pgfpathlineto{\pgfqpoint{1.774795in}{1.222030in}}%
\pgfpathlineto{\pgfqpoint{1.767385in}{1.210528in}}%
\pgfpathclose%
\pgfusepath{fill}%
\end{pgfscope}%
\begin{pgfscope}%
\pgfpathrectangle{\pgfqpoint{0.407592in}{0.222904in}}{\pgfqpoint{1.411163in}{1.411163in}}%
\pgfusepath{clip}%
\pgfsetbuttcap%
\pgfsetroundjoin%
\definecolor{currentfill}{rgb}{0.556863,0.654902,0.847059}%
\pgfsetfillcolor{currentfill}%
\pgfsetfillopacity{0.130000}%
\pgfsetlinewidth{0.000000pt}%
\definecolor{currentstroke}{rgb}{0.000000,0.000000,0.000000}%
\pgfsetstrokecolor{currentstroke}%
\pgfsetdash{}{0pt}%
\pgfpathmoveto{\pgfqpoint{1.585728in}{1.035780in}}%
\pgfpathlineto{\pgfqpoint{1.603358in}{1.047194in}}%
\pgfpathlineto{\pgfqpoint{1.621066in}{1.058798in}}%
\pgfpathlineto{\pgfqpoint{1.628633in}{1.068652in}}%
\pgfpathlineto{\pgfqpoint{1.636188in}{1.078777in}}%
\pgfpathlineto{\pgfqpoint{1.618503in}{1.067480in}}%
\pgfpathlineto{\pgfqpoint{1.600895in}{1.056387in}}%
\pgfpathlineto{\pgfqpoint{1.593317in}{1.045943in}}%
\pgfpathlineto{\pgfqpoint{1.585728in}{1.035780in}}%
\pgfpathclose%
\pgfusepath{fill}%
\end{pgfscope}%
\begin{pgfscope}%
\pgfpathrectangle{\pgfqpoint{0.407592in}{0.222904in}}{\pgfqpoint{1.411163in}{1.411163in}}%
\pgfusepath{clip}%
\pgfsetbuttcap%
\pgfsetroundjoin%
\definecolor{currentfill}{rgb}{0.556863,0.654902,0.847059}%
\pgfsetfillcolor{currentfill}%
\pgfsetfillopacity{0.130000}%
\pgfsetlinewidth{0.000000pt}%
\definecolor{currentstroke}{rgb}{0.000000,0.000000,0.000000}%
\pgfsetstrokecolor{currentstroke}%
\pgfsetdash{}{0pt}%
\pgfpathmoveto{\pgfqpoint{0.674037in}{0.950724in}}%
\pgfpathlineto{\pgfqpoint{0.690432in}{0.935583in}}%
\pgfpathlineto{\pgfqpoint{0.706837in}{0.920904in}}%
\pgfpathlineto{\pgfqpoint{0.715930in}{0.916803in}}%
\pgfpathlineto{\pgfqpoint{0.724979in}{0.912959in}}%
\pgfpathlineto{\pgfqpoint{0.708618in}{0.928106in}}%
\pgfpathlineto{\pgfqpoint{0.692266in}{0.943706in}}%
\pgfpathlineto{\pgfqpoint{0.683173in}{0.947083in}}%
\pgfpathlineto{\pgfqpoint{0.674037in}{0.950724in}}%
\pgfpathclose%
\pgfusepath{fill}%
\end{pgfscope}%
\begin{pgfscope}%
\pgfpathrectangle{\pgfqpoint{0.407592in}{0.222904in}}{\pgfqpoint{1.411163in}{1.411163in}}%
\pgfusepath{clip}%
\pgfsetbuttcap%
\pgfsetroundjoin%
\definecolor{currentfill}{rgb}{0.556863,0.654902,0.847059}%
\pgfsetfillcolor{currentfill}%
\pgfsetfillopacity{0.130000}%
\pgfsetlinewidth{0.000000pt}%
\definecolor{currentstroke}{rgb}{0.000000,0.000000,0.000000}%
\pgfsetstrokecolor{currentstroke}%
\pgfsetdash{}{0pt}%
\pgfpathmoveto{\pgfqpoint{1.139435in}{0.743451in}}%
\pgfpathlineto{\pgfqpoint{1.156136in}{0.749515in}}%
\pgfpathlineto{\pgfqpoint{1.172900in}{0.756403in}}%
\pgfpathlineto{\pgfqpoint{1.181032in}{0.756073in}}%
\pgfpathlineto{\pgfqpoint{1.189135in}{0.756344in}}%
\pgfpathlineto{\pgfqpoint{1.172402in}{0.748915in}}%
\pgfpathlineto{\pgfqpoint{1.155734in}{0.742290in}}%
\pgfpathlineto{\pgfqpoint{1.147600in}{0.742555in}}%
\pgfpathlineto{\pgfqpoint{1.139435in}{0.743451in}}%
\pgfpathclose%
\pgfusepath{fill}%
\end{pgfscope}%
\begin{pgfscope}%
\pgfpathrectangle{\pgfqpoint{0.407592in}{0.222904in}}{\pgfqpoint{1.411163in}{1.411163in}}%
\pgfusepath{clip}%
\pgfsetbuttcap%
\pgfsetroundjoin%
\definecolor{currentfill}{rgb}{0.556863,0.654902,0.847059}%
\pgfsetfillcolor{currentfill}%
\pgfsetfillopacity{0.130000}%
\pgfsetlinewidth{0.000000pt}%
\definecolor{currentstroke}{rgb}{0.000000,0.000000,0.000000}%
\pgfsetstrokecolor{currentstroke}%
\pgfsetdash{}{0pt}%
\pgfpathmoveto{\pgfqpoint{0.790524in}{0.857854in}}%
\pgfpathlineto{\pgfqpoint{0.806938in}{0.845714in}}%
\pgfpathlineto{\pgfqpoint{0.823366in}{0.834341in}}%
\pgfpathlineto{\pgfqpoint{0.832235in}{0.829316in}}%
\pgfpathlineto{\pgfqpoint{0.841063in}{0.824538in}}%
\pgfpathlineto{\pgfqpoint{0.824682in}{0.836375in}}%
\pgfpathlineto{\pgfqpoint{0.808313in}{0.849002in}}%
\pgfpathlineto{\pgfqpoint{0.799440in}{0.853295in}}%
\pgfpathlineto{\pgfqpoint{0.790524in}{0.857854in}}%
\pgfpathclose%
\pgfusepath{fill}%
\end{pgfscope}%
\begin{pgfscope}%
\pgfpathrectangle{\pgfqpoint{0.407592in}{0.222904in}}{\pgfqpoint{1.411163in}{1.411163in}}%
\pgfusepath{clip}%
\pgfsetbuttcap%
\pgfsetroundjoin%
\definecolor{currentfill}{rgb}{0.556863,0.654902,0.847059}%
\pgfsetfillcolor{currentfill}%
\pgfsetfillopacity{0.130000}%
\pgfsetlinewidth{0.000000pt}%
\definecolor{currentstroke}{rgb}{0.000000,0.000000,0.000000}%
\pgfsetstrokecolor{currentstroke}%
\pgfsetdash{}{0pt}%
\pgfpathmoveto{\pgfqpoint{1.189135in}{0.756344in}}%
\pgfpathlineto{\pgfqpoint{1.205935in}{0.764469in}}%
\pgfpathlineto{\pgfqpoint{1.222804in}{0.773194in}}%
\pgfpathlineto{\pgfqpoint{1.230851in}{0.774542in}}%
\pgfpathlineto{\pgfqpoint{1.238873in}{0.776510in}}%
\pgfpathlineto{\pgfqpoint{1.222031in}{0.767388in}}%
\pgfpathlineto{\pgfqpoint{1.205258in}{0.758838in}}%
\pgfpathlineto{\pgfqpoint{1.197210in}{0.757253in}}%
\pgfpathlineto{\pgfqpoint{1.189135in}{0.756344in}}%
\pgfpathclose%
\pgfusepath{fill}%
\end{pgfscope}%
\begin{pgfscope}%
\pgfpathrectangle{\pgfqpoint{0.407592in}{0.222904in}}{\pgfqpoint{1.411163in}{1.411163in}}%
\pgfusepath{clip}%
\pgfsetbuttcap%
\pgfsetroundjoin%
\definecolor{currentfill}{rgb}{0.556863,0.654902,0.847059}%
\pgfsetfillcolor{currentfill}%
\pgfsetfillopacity{0.130000}%
\pgfsetlinewidth{0.000000pt}%
\definecolor{currentstroke}{rgb}{0.000000,0.000000,0.000000}%
\pgfsetstrokecolor{currentstroke}%
\pgfsetdash{}{0pt}%
\pgfpathmoveto{\pgfqpoint{1.089710in}{0.738908in}}%
\pgfpathlineto{\pgfqpoint{1.106332in}{0.742401in}}%
\pgfpathlineto{\pgfqpoint{1.123009in}{0.746937in}}%
\pgfpathlineto{\pgfqpoint{1.131239in}{0.744928in}}%
\pgfpathlineto{\pgfqpoint{1.139435in}{0.743451in}}%
\pgfpathlineto{\pgfqpoint{1.122795in}{0.738324in}}%
\pgfpathlineto{\pgfqpoint{1.106213in}{0.734261in}}%
\pgfpathlineto{\pgfqpoint{1.097979in}{0.736318in}}%
\pgfpathlineto{\pgfqpoint{1.089710in}{0.738908in}}%
\pgfpathclose%
\pgfusepath{fill}%
\end{pgfscope}%
\begin{pgfscope}%
\pgfpathrectangle{\pgfqpoint{0.407592in}{0.222904in}}{\pgfqpoint{1.411163in}{1.411163in}}%
\pgfusepath{clip}%
\pgfsetbuttcap%
\pgfsetroundjoin%
\definecolor{currentfill}{rgb}{0.556863,0.654902,0.847059}%
\pgfsetfillcolor{currentfill}%
\pgfsetfillopacity{0.130000}%
\pgfsetlinewidth{0.000000pt}%
\definecolor{currentstroke}{rgb}{0.000000,0.000000,0.000000}%
\pgfsetstrokecolor{currentstroke}%
\pgfsetdash{}{0pt}%
\pgfpathmoveto{\pgfqpoint{1.701811in}{1.143728in}}%
\pgfpathlineto{\pgfqpoint{1.719686in}{1.154992in}}%
\pgfpathlineto{\pgfqpoint{1.737641in}{1.166444in}}%
\pgfpathlineto{\pgfqpoint{1.745093in}{1.177162in}}%
\pgfpathlineto{\pgfqpoint{1.752534in}{1.188087in}}%
\pgfpathlineto{\pgfqpoint{1.734604in}{1.176930in}}%
\pgfpathlineto{\pgfqpoint{1.716752in}{1.165971in}}%
\pgfpathlineto{\pgfqpoint{1.709287in}{1.154744in}}%
\pgfpathlineto{\pgfqpoint{1.701811in}{1.143728in}}%
\pgfpathclose%
\pgfusepath{fill}%
\end{pgfscope}%
\begin{pgfscope}%
\pgfpathrectangle{\pgfqpoint{0.407592in}{0.222904in}}{\pgfqpoint{1.411163in}{1.411163in}}%
\pgfusepath{clip}%
\pgfsetbuttcap%
\pgfsetroundjoin%
\definecolor{currentfill}{rgb}{0.556863,0.654902,0.847059}%
\pgfsetfillcolor{currentfill}%
\pgfsetfillopacity{0.130000}%
\pgfsetlinewidth{0.000000pt}%
\definecolor{currentstroke}{rgb}{0.000000,0.000000,0.000000}%
\pgfsetstrokecolor{currentstroke}%
\pgfsetdash{}{0pt}%
\pgfpathmoveto{\pgfqpoint{0.557189in}{1.054924in}}%
\pgfpathlineto{\pgfqpoint{0.573574in}{1.038091in}}%
\pgfpathlineto{\pgfqpoint{0.589967in}{1.021543in}}%
\pgfpathlineto{\pgfqpoint{0.599279in}{1.018241in}}%
\pgfpathlineto{\pgfqpoint{0.608546in}{1.015186in}}%
\pgfpathlineto{\pgfqpoint{0.592196in}{1.032116in}}%
\pgfpathlineto{\pgfqpoint{0.575854in}{1.049316in}}%
\pgfpathlineto{\pgfqpoint{0.566544in}{1.051996in}}%
\pgfpathlineto{\pgfqpoint{0.557189in}{1.054924in}}%
\pgfpathclose%
\pgfusepath{fill}%
\end{pgfscope}%
\begin{pgfscope}%
\pgfpathrectangle{\pgfqpoint{0.407592in}{0.222904in}}{\pgfqpoint{1.411163in}{1.411163in}}%
\pgfusepath{clip}%
\pgfsetbuttcap%
\pgfsetroundjoin%
\definecolor{currentfill}{rgb}{0.556863,0.654902,0.847059}%
\pgfsetfillcolor{currentfill}%
\pgfsetfillopacity{0.130000}%
\pgfsetlinewidth{0.000000pt}%
\definecolor{currentstroke}{rgb}{0.000000,0.000000,0.000000}%
\pgfsetstrokecolor{currentstroke}%
\pgfsetdash{}{0pt}%
\pgfpathmoveto{\pgfqpoint{1.238873in}{0.776510in}}%
\pgfpathlineto{\pgfqpoint{1.255787in}{0.786123in}}%
\pgfpathlineto{\pgfqpoint{1.272772in}{0.796163in}}%
\pgfpathlineto{\pgfqpoint{1.280746in}{0.799041in}}%
\pgfpathlineto{\pgfqpoint{1.288698in}{0.802519in}}%
\pgfpathlineto{\pgfqpoint{1.271735in}{0.792254in}}%
\pgfpathlineto{\pgfqpoint{1.254845in}{0.782394in}}%
\pgfpathlineto{\pgfqpoint{1.246871in}{0.779121in}}%
\pgfpathlineto{\pgfqpoint{1.238873in}{0.776510in}}%
\pgfpathclose%
\pgfusepath{fill}%
\end{pgfscope}%
\begin{pgfscope}%
\pgfpathrectangle{\pgfqpoint{0.407592in}{0.222904in}}{\pgfqpoint{1.411163in}{1.411163in}}%
\pgfusepath{clip}%
\pgfsetbuttcap%
\pgfsetroundjoin%
\definecolor{currentfill}{rgb}{0.556863,0.654902,0.847059}%
\pgfsetfillcolor{currentfill}%
\pgfsetfillopacity{0.130000}%
\pgfsetlinewidth{0.000000pt}%
\definecolor{currentstroke}{rgb}{0.000000,0.000000,0.000000}%
\pgfsetstrokecolor{currentstroke}%
\pgfsetdash{}{0pt}%
\pgfpathmoveto{\pgfqpoint{1.404341in}{0.879934in}}%
\pgfpathlineto{\pgfqpoint{1.421615in}{0.891481in}}%
\pgfpathlineto{\pgfqpoint{1.438965in}{0.903229in}}%
\pgfpathlineto{\pgfqpoint{1.446728in}{0.910318in}}%
\pgfpathlineto{\pgfqpoint{1.454476in}{0.917843in}}%
\pgfpathlineto{\pgfqpoint{1.437147in}{0.906288in}}%
\pgfpathlineto{\pgfqpoint{1.419893in}{0.894947in}}%
\pgfpathlineto{\pgfqpoint{1.412124in}{0.887205in}}%
\pgfpathlineto{\pgfqpoint{1.404341in}{0.879934in}}%
\pgfpathclose%
\pgfusepath{fill}%
\end{pgfscope}%
\begin{pgfscope}%
\pgfpathrectangle{\pgfqpoint{0.407592in}{0.222904in}}{\pgfqpoint{1.411163in}{1.411163in}}%
\pgfusepath{clip}%
\pgfsetbuttcap%
\pgfsetroundjoin%
\definecolor{currentfill}{rgb}{0.556863,0.654902,0.847059}%
\pgfsetfillcolor{currentfill}%
\pgfsetfillopacity{0.130000}%
\pgfsetlinewidth{0.000000pt}%
\definecolor{currentstroke}{rgb}{0.000000,0.000000,0.000000}%
\pgfsetstrokecolor{currentstroke}%
\pgfsetdash{}{0pt}%
\pgfpathmoveto{\pgfqpoint{1.520136in}{0.974725in}}%
\pgfpathlineto{\pgfqpoint{1.537654in}{0.986351in}}%
\pgfpathlineto{\pgfqpoint{1.555250in}{0.998161in}}%
\pgfpathlineto{\pgfqpoint{1.562888in}{1.007089in}}%
\pgfpathlineto{\pgfqpoint{1.570514in}{1.016342in}}%
\pgfpathlineto{\pgfqpoint{1.552939in}{1.004816in}}%
\pgfpathlineto{\pgfqpoint{1.535443in}{0.993491in}}%
\pgfpathlineto{\pgfqpoint{1.527795in}{0.983937in}}%
\pgfpathlineto{\pgfqpoint{1.520136in}{0.974725in}}%
\pgfpathclose%
\pgfusepath{fill}%
\end{pgfscope}%
\begin{pgfscope}%
\pgfpathrectangle{\pgfqpoint{0.407592in}{0.222904in}}{\pgfqpoint{1.411163in}{1.411163in}}%
\pgfusepath{clip}%
\pgfsetbuttcap%
\pgfsetroundjoin%
\definecolor{currentfill}{rgb}{0.556863,0.654902,0.847059}%
\pgfsetfillcolor{currentfill}%
\pgfsetfillopacity{0.130000}%
\pgfsetlinewidth{0.000000pt}%
\definecolor{currentstroke}{rgb}{0.000000,0.000000,0.000000}%
\pgfsetstrokecolor{currentstroke}%
\pgfsetdash{}{0pt}%
\pgfpathmoveto{\pgfqpoint{1.818197in}{1.255429in}}%
\pgfpathlineto{\pgfqpoint{1.836319in}{1.266580in}}%
\pgfpathlineto{\pgfqpoint{1.843671in}{1.277999in}}%
\pgfpathlineto{\pgfqpoint{1.825562in}{1.266988in}}%
\pgfusepath{fill}%
\end{pgfscope}%
\begin{pgfscope}%
\pgfpathrectangle{\pgfqpoint{0.407592in}{0.222904in}}{\pgfqpoint{1.411163in}{1.411163in}}%
\pgfusepath{clip}%
\pgfsetbuttcap%
\pgfsetroundjoin%
\definecolor{currentfill}{rgb}{0.556863,0.654902,0.847059}%
\pgfsetfillcolor{currentfill}%
\pgfsetfillopacity{0.130000}%
\pgfsetlinewidth{0.000000pt}%
\definecolor{currentstroke}{rgb}{0.000000,0.000000,0.000000}%
\pgfsetstrokecolor{currentstroke}%
\pgfsetdash{}{0pt}%
\pgfpathmoveto{\pgfqpoint{1.039888in}{0.742986in}}%
\pgfpathlineto{\pgfqpoint{1.056454in}{0.743637in}}%
\pgfpathlineto{\pgfqpoint{1.073067in}{0.745476in}}%
\pgfpathlineto{\pgfqpoint{1.081406in}{0.741977in}}%
\pgfpathlineto{\pgfqpoint{1.089710in}{0.738908in}}%
\pgfpathlineto{\pgfqpoint{1.073141in}{0.736568in}}%
\pgfpathlineto{\pgfqpoint{1.056620in}{0.735487in}}%
\pgfpathlineto{\pgfqpoint{1.048272in}{0.739031in}}%
\pgfpathlineto{\pgfqpoint{1.039888in}{0.742986in}}%
\pgfpathclose%
\pgfusepath{fill}%
\end{pgfscope}%
\begin{pgfscope}%
\pgfpathrectangle{\pgfqpoint{0.407592in}{0.222904in}}{\pgfqpoint{1.411163in}{1.411163in}}%
\pgfusepath{clip}%
\pgfsetbuttcap%
\pgfsetroundjoin%
\definecolor{currentfill}{rgb}{0.556863,0.654902,0.847059}%
\pgfsetfillcolor{currentfill}%
\pgfsetfillopacity{0.130000}%
\pgfsetlinewidth{0.000000pt}%
\definecolor{currentstroke}{rgb}{0.000000,0.000000,0.000000}%
\pgfsetstrokecolor{currentstroke}%
\pgfsetdash{}{0pt}%
\pgfpathmoveto{\pgfqpoint{0.706837in}{0.920904in}}%
\pgfpathlineto{\pgfqpoint{0.723253in}{0.906731in}}%
\pgfpathlineto{\pgfqpoint{0.739678in}{0.893106in}}%
\pgfpathlineto{\pgfqpoint{0.748726in}{0.888535in}}%
\pgfpathlineto{\pgfqpoint{0.757730in}{0.884208in}}%
\pgfpathlineto{\pgfqpoint{0.741349in}{0.898310in}}%
\pgfpathlineto{\pgfqpoint{0.724979in}{0.912959in}}%
\pgfpathlineto{\pgfqpoint{0.715930in}{0.916803in}}%
\pgfpathlineto{\pgfqpoint{0.706837in}{0.920904in}}%
\pgfpathclose%
\pgfusepath{fill}%
\end{pgfscope}%
\begin{pgfscope}%
\pgfpathrectangle{\pgfqpoint{0.407592in}{0.222904in}}{\pgfqpoint{1.411163in}{1.411163in}}%
\pgfusepath{clip}%
\pgfsetbuttcap%
\pgfsetroundjoin%
\definecolor{currentfill}{rgb}{0.556863,0.654902,0.847059}%
\pgfsetfillcolor{currentfill}%
\pgfsetfillopacity{0.130000}%
\pgfsetlinewidth{0.000000pt}%
\definecolor{currentstroke}{rgb}{0.000000,0.000000,0.000000}%
\pgfsetstrokecolor{currentstroke}%
\pgfsetdash{}{0pt}%
\pgfpathmoveto{\pgfqpoint{1.636188in}{1.078777in}}%
\pgfpathlineto{\pgfqpoint{1.653952in}{1.090268in}}%
\pgfpathlineto{\pgfqpoint{1.671795in}{1.101941in}}%
\pgfpathlineto{\pgfqpoint{1.679316in}{1.112029in}}%
\pgfpathlineto{\pgfqpoint{1.686825in}{1.122362in}}%
\pgfpathlineto{\pgfqpoint{1.669006in}{1.110985in}}%
\pgfpathlineto{\pgfqpoint{1.651264in}{1.099804in}}%
\pgfpathlineto{\pgfqpoint{1.643732in}{1.089165in}}%
\pgfpathlineto{\pgfqpoint{1.636188in}{1.078777in}}%
\pgfpathclose%
\pgfusepath{fill}%
\end{pgfscope}%
\begin{pgfscope}%
\pgfpathrectangle{\pgfqpoint{0.407592in}{0.222904in}}{\pgfqpoint{1.411163in}{1.411163in}}%
\pgfusepath{clip}%
\pgfsetbuttcap%
\pgfsetroundjoin%
\definecolor{currentfill}{rgb}{0.556863,0.654902,0.847059}%
\pgfsetfillcolor{currentfill}%
\pgfsetfillopacity{0.130000}%
\pgfsetlinewidth{0.000000pt}%
\definecolor{currentstroke}{rgb}{0.000000,0.000000,0.000000}%
\pgfsetstrokecolor{currentstroke}%
\pgfsetdash{}{0pt}%
\pgfpathmoveto{\pgfqpoint{0.439960in}{1.165411in}}%
\pgfpathlineto{\pgfqpoint{0.456335in}{1.147569in}}%
\pgfpathlineto{\pgfqpoint{0.472720in}{1.129914in}}%
\pgfpathlineto{\pgfqpoint{0.482251in}{1.127185in}}%
\pgfpathlineto{\pgfqpoint{0.491736in}{1.124683in}}%
\pgfpathlineto{\pgfqpoint{0.475395in}{1.142637in}}%
\pgfpathlineto{\pgfqpoint{0.459061in}{1.160764in}}%
\pgfpathlineto{\pgfqpoint{0.449533in}{1.162975in}}%
\pgfpathlineto{\pgfqpoint{0.439960in}{1.165411in}}%
\pgfpathclose%
\pgfusepath{fill}%
\end{pgfscope}%
\begin{pgfscope}%
\pgfpathrectangle{\pgfqpoint{0.407592in}{0.222904in}}{\pgfqpoint{1.411163in}{1.411163in}}%
\pgfusepath{clip}%
\pgfsetbuttcap%
\pgfsetroundjoin%
\definecolor{currentfill}{rgb}{0.556863,0.654902,0.847059}%
\pgfsetfillcolor{currentfill}%
\pgfsetfillopacity{0.130000}%
\pgfsetlinewidth{0.000000pt}%
\definecolor{currentstroke}{rgb}{0.000000,0.000000,0.000000}%
\pgfsetstrokecolor{currentstroke}%
\pgfsetdash{}{0pt}%
\pgfpathmoveto{\pgfqpoint{1.288698in}{0.802519in}}%
\pgfpathlineto{\pgfqpoint{1.305734in}{0.813134in}}%
\pgfpathlineto{\pgfqpoint{1.322844in}{0.824057in}}%
\pgfpathlineto{\pgfqpoint{1.330753in}{0.828253in}}%
\pgfpathlineto{\pgfqpoint{1.338642in}{0.833004in}}%
\pgfpathlineto{\pgfqpoint{1.321553in}{0.822006in}}%
\pgfpathlineto{\pgfqpoint{1.304539in}{0.811307in}}%
\pgfpathlineto{\pgfqpoint{1.296628in}{0.806605in}}%
\pgfpathlineto{\pgfqpoint{1.288698in}{0.802519in}}%
\pgfpathclose%
\pgfusepath{fill}%
\end{pgfscope}%
\begin{pgfscope}%
\pgfpathrectangle{\pgfqpoint{0.407592in}{0.222904in}}{\pgfqpoint{1.411163in}{1.411163in}}%
\pgfusepath{clip}%
\pgfsetbuttcap%
\pgfsetroundjoin%
\definecolor{currentfill}{rgb}{0.556863,0.654902,0.847059}%
\pgfsetfillcolor{currentfill}%
\pgfsetfillopacity{0.130000}%
\pgfsetlinewidth{0.000000pt}%
\definecolor{currentstroke}{rgb}{0.000000,0.000000,0.000000}%
\pgfsetstrokecolor{currentstroke}%
\pgfsetdash{}{0pt}%
\pgfpathmoveto{\pgfqpoint{0.906739in}{0.786422in}}%
\pgfpathlineto{\pgfqpoint{0.923206in}{0.779499in}}%
\pgfpathlineto{\pgfqpoint{0.939697in}{0.773708in}}%
\pgfpathlineto{\pgfqpoint{0.948336in}{0.768266in}}%
\pgfpathlineto{\pgfqpoint{0.956935in}{0.763069in}}%
\pgfpathlineto{\pgfqpoint{0.940494in}{0.768983in}}%
\pgfpathlineto{\pgfqpoint{0.924077in}{0.776122in}}%
\pgfpathlineto{\pgfqpoint{0.915428in}{0.781145in}}%
\pgfpathlineto{\pgfqpoint{0.906739in}{0.786422in}}%
\pgfpathclose%
\pgfusepath{fill}%
\end{pgfscope}%
\begin{pgfscope}%
\pgfpathrectangle{\pgfqpoint{0.407592in}{0.222904in}}{\pgfqpoint{1.411163in}{1.411163in}}%
\pgfusepath{clip}%
\pgfsetbuttcap%
\pgfsetroundjoin%
\definecolor{currentfill}{rgb}{0.556863,0.654902,0.847059}%
\pgfsetfillcolor{currentfill}%
\pgfsetfillopacity{0.130000}%
\pgfsetlinewidth{0.000000pt}%
\definecolor{currentstroke}{rgb}{0.000000,0.000000,0.000000}%
\pgfsetstrokecolor{currentstroke}%
\pgfsetdash{}{0pt}%
\pgfpathmoveto{\pgfqpoint{0.589967in}{1.021543in}}%
\pgfpathlineto{\pgfqpoint{0.606370in}{1.005307in}}%
\pgfpathlineto{\pgfqpoint{0.622781in}{0.989409in}}%
\pgfpathlineto{\pgfqpoint{0.632050in}{0.985705in}}%
\pgfpathlineto{\pgfqpoint{0.641274in}{0.982244in}}%
\pgfpathlineto{\pgfqpoint{0.624906in}{0.998553in}}%
\pgfpathlineto{\pgfqpoint{0.608546in}{1.015186in}}%
\pgfpathlineto{\pgfqpoint{0.599279in}{1.018241in}}%
\pgfpathlineto{\pgfqpoint{0.589967in}{1.021543in}}%
\pgfpathclose%
\pgfusepath{fill}%
\end{pgfscope}%
\begin{pgfscope}%
\pgfpathrectangle{\pgfqpoint{0.407592in}{0.222904in}}{\pgfqpoint{1.411163in}{1.411163in}}%
\pgfusepath{clip}%
\pgfsetbuttcap%
\pgfsetroundjoin%
\definecolor{currentfill}{rgb}{0.556863,0.654902,0.847059}%
\pgfsetfillcolor{currentfill}%
\pgfsetfillopacity{0.130000}%
\pgfsetlinewidth{0.000000pt}%
\definecolor{currentstroke}{rgb}{0.000000,0.000000,0.000000}%
\pgfsetstrokecolor{currentstroke}%
\pgfsetdash{}{0pt}%
\pgfpathmoveto{\pgfqpoint{0.823366in}{0.834341in}}%
\pgfpathlineto{\pgfqpoint{0.839808in}{0.823794in}}%
\pgfpathlineto{\pgfqpoint{0.856266in}{0.814128in}}%
\pgfpathlineto{\pgfqpoint{0.865087in}{0.808704in}}%
\pgfpathlineto{\pgfqpoint{0.873867in}{0.803505in}}%
\pgfpathlineto{\pgfqpoint{0.857457in}{0.813559in}}%
\pgfpathlineto{\pgfqpoint{0.841063in}{0.824538in}}%
\pgfpathlineto{\pgfqpoint{0.832235in}{0.829316in}}%
\pgfpathlineto{\pgfqpoint{0.823366in}{0.834341in}}%
\pgfpathclose%
\pgfusepath{fill}%
\end{pgfscope}%
\begin{pgfscope}%
\pgfpathrectangle{\pgfqpoint{0.407592in}{0.222904in}}{\pgfqpoint{1.411163in}{1.411163in}}%
\pgfusepath{clip}%
\pgfsetbuttcap%
\pgfsetroundjoin%
\definecolor{currentfill}{rgb}{0.556863,0.654902,0.847059}%
\pgfsetfillcolor{currentfill}%
\pgfsetfillopacity{0.130000}%
\pgfsetlinewidth{0.000000pt}%
\definecolor{currentstroke}{rgb}{0.000000,0.000000,0.000000}%
\pgfsetstrokecolor{currentstroke}%
\pgfsetdash{}{0pt}%
\pgfpathmoveto{\pgfqpoint{1.454476in}{0.917843in}}%
\pgfpathlineto{\pgfqpoint{1.471883in}{0.929595in}}%
\pgfpathlineto{\pgfqpoint{1.489366in}{0.941528in}}%
\pgfpathlineto{\pgfqpoint{1.497079in}{0.949258in}}%
\pgfpathlineto{\pgfqpoint{1.504778in}{0.957375in}}%
\pgfpathlineto{\pgfqpoint{1.487315in}{0.945669in}}%
\pgfpathlineto{\pgfqpoint{1.469930in}{0.934159in}}%
\pgfpathlineto{\pgfqpoint{1.462210in}{0.925794in}}%
\pgfpathlineto{\pgfqpoint{1.454476in}{0.917843in}}%
\pgfpathclose%
\pgfusepath{fill}%
\end{pgfscope}%
\begin{pgfscope}%
\pgfpathrectangle{\pgfqpoint{0.407592in}{0.222904in}}{\pgfqpoint{1.411163in}{1.411163in}}%
\pgfusepath{clip}%
\pgfsetbuttcap%
\pgfsetroundjoin%
\definecolor{currentfill}{rgb}{0.556863,0.654902,0.847059}%
\pgfsetfillcolor{currentfill}%
\pgfsetfillopacity{0.130000}%
\pgfsetlinewidth{0.000000pt}%
\definecolor{currentstroke}{rgb}{0.000000,0.000000,0.000000}%
\pgfsetstrokecolor{currentstroke}%
\pgfsetdash{}{0pt}%
\pgfpathmoveto{\pgfqpoint{1.752534in}{1.188087in}}%
\pgfpathlineto{\pgfqpoint{1.770545in}{1.199433in}}%
\pgfpathlineto{\pgfqpoint{1.788636in}{1.210959in}}%
\pgfpathlineto{\pgfqpoint{1.796042in}{1.221799in}}%
\pgfpathlineto{\pgfqpoint{1.803438in}{1.232829in}}%
\pgfpathlineto{\pgfqpoint{1.785372in}{1.221583in}}%
\pgfpathlineto{\pgfqpoint{1.767385in}{1.210528in}}%
\pgfpathlineto{\pgfqpoint{1.759965in}{1.199211in}}%
\pgfpathlineto{\pgfqpoint{1.752534in}{1.188087in}}%
\pgfpathclose%
\pgfusepath{fill}%
\end{pgfscope}%
\begin{pgfscope}%
\pgfpathrectangle{\pgfqpoint{0.407592in}{0.222904in}}{\pgfqpoint{1.411163in}{1.411163in}}%
\pgfusepath{clip}%
\pgfsetbuttcap%
\pgfsetroundjoin%
\definecolor{currentfill}{rgb}{0.556863,0.654902,0.847059}%
\pgfsetfillcolor{currentfill}%
\pgfsetfillopacity{0.130000}%
\pgfsetlinewidth{0.000000pt}%
\definecolor{currentstroke}{rgb}{0.000000,0.000000,0.000000}%
\pgfsetstrokecolor{currentstroke}%
\pgfsetdash{}{0pt}%
\pgfpathmoveto{\pgfqpoint{0.989901in}{0.755014in}}%
\pgfpathlineto{\pgfqpoint{1.006434in}{0.752873in}}%
\pgfpathlineto{\pgfqpoint{1.023005in}{0.751955in}}%
\pgfpathlineto{\pgfqpoint{1.031465in}{0.747308in}}%
\pgfpathlineto{\pgfqpoint{1.039888in}{0.742986in}}%
\pgfpathlineto{\pgfqpoint{1.023365in}{0.743601in}}%
\pgfpathlineto{\pgfqpoint{1.006881in}{0.745540in}}%
\pgfpathlineto{\pgfqpoint{0.998411in}{0.750122in}}%
\pgfpathlineto{\pgfqpoint{0.989901in}{0.755014in}}%
\pgfpathclose%
\pgfusepath{fill}%
\end{pgfscope}%
\begin{pgfscope}%
\pgfpathrectangle{\pgfqpoint{0.407592in}{0.222904in}}{\pgfqpoint{1.411163in}{1.411163in}}%
\pgfusepath{clip}%
\pgfsetbuttcap%
\pgfsetroundjoin%
\definecolor{currentfill}{rgb}{0.556863,0.654902,0.847059}%
\pgfsetfillcolor{currentfill}%
\pgfsetfillopacity{0.130000}%
\pgfsetlinewidth{0.000000pt}%
\definecolor{currentstroke}{rgb}{0.000000,0.000000,0.000000}%
\pgfsetstrokecolor{currentstroke}%
\pgfsetdash{}{0pt}%
\pgfpathmoveto{\pgfqpoint{1.570514in}{1.016342in}}%
\pgfpathlineto{\pgfqpoint{1.588166in}{1.028055in}}%
\pgfpathlineto{\pgfqpoint{1.605897in}{1.039942in}}%
\pgfpathlineto{\pgfqpoint{1.613488in}{1.049225in}}%
\pgfpathlineto{\pgfqpoint{1.621066in}{1.058798in}}%
\pgfpathlineto{\pgfqpoint{1.603358in}{1.047194in}}%
\pgfpathlineto{\pgfqpoint{1.585728in}{1.035780in}}%
\pgfpathlineto{\pgfqpoint{1.578127in}{1.025909in}}%
\pgfpathlineto{\pgfqpoint{1.570514in}{1.016342in}}%
\pgfpathclose%
\pgfusepath{fill}%
\end{pgfscope}%
\begin{pgfscope}%
\pgfpathrectangle{\pgfqpoint{0.407592in}{0.222904in}}{\pgfqpoint{1.411163in}{1.411163in}}%
\pgfusepath{clip}%
\pgfsetbuttcap%
\pgfsetroundjoin%
\definecolor{currentfill}{rgb}{0.556863,0.654902,0.847059}%
\pgfsetfillcolor{currentfill}%
\pgfsetfillopacity{0.130000}%
\pgfsetlinewidth{0.000000pt}%
\definecolor{currentstroke}{rgb}{0.000000,0.000000,0.000000}%
\pgfsetstrokecolor{currentstroke}%
\pgfsetdash{}{0pt}%
\pgfpathmoveto{\pgfqpoint{0.472720in}{1.129914in}}%
\pgfpathlineto{\pgfqpoint{0.489113in}{1.112460in}}%
\pgfpathlineto{\pgfqpoint{0.505514in}{1.095225in}}%
\pgfpathlineto{\pgfqpoint{0.515003in}{1.092177in}}%
\pgfpathlineto{\pgfqpoint{0.524446in}{1.089358in}}%
\pgfpathlineto{\pgfqpoint{0.508087in}{1.106918in}}%
\pgfpathlineto{\pgfqpoint{0.491736in}{1.124683in}}%
\pgfpathlineto{\pgfqpoint{0.482251in}{1.127185in}}%
\pgfpathlineto{\pgfqpoint{0.472720in}{1.129914in}}%
\pgfpathclose%
\pgfusepath{fill}%
\end{pgfscope}%
\begin{pgfscope}%
\pgfpathrectangle{\pgfqpoint{0.407592in}{0.222904in}}{\pgfqpoint{1.411163in}{1.411163in}}%
\pgfusepath{clip}%
\pgfsetbuttcap%
\pgfsetroundjoin%
\definecolor{currentfill}{rgb}{0.556863,0.654902,0.847059}%
\pgfsetfillcolor{currentfill}%
\pgfsetfillopacity{0.130000}%
\pgfsetlinewidth{0.000000pt}%
\definecolor{currentstroke}{rgb}{0.000000,0.000000,0.000000}%
\pgfsetstrokecolor{currentstroke}%
\pgfsetdash{}{0pt}%
\pgfpathmoveto{\pgfqpoint{1.338642in}{0.833004in}}%
\pgfpathlineto{\pgfqpoint{1.355806in}{0.844265in}}%
\pgfpathlineto{\pgfqpoint{1.373045in}{0.855762in}}%
\pgfpathlineto{\pgfqpoint{1.380894in}{0.861056in}}%
\pgfpathlineto{\pgfqpoint{1.388726in}{0.866852in}}%
\pgfpathlineto{\pgfqpoint{1.371508in}{0.855394in}}%
\pgfpathlineto{\pgfqpoint{1.354365in}{0.844172in}}%
\pgfpathlineto{\pgfqpoint{1.346512in}{0.838311in}}%
\pgfpathlineto{\pgfqpoint{1.338642in}{0.833004in}}%
\pgfpathclose%
\pgfusepath{fill}%
\end{pgfscope}%
\begin{pgfscope}%
\pgfpathrectangle{\pgfqpoint{0.407592in}{0.222904in}}{\pgfqpoint{1.411163in}{1.411163in}}%
\pgfusepath{clip}%
\pgfsetbuttcap%
\pgfsetroundjoin%
\definecolor{currentfill}{rgb}{0.556863,0.654902,0.847059}%
\pgfsetfillcolor{currentfill}%
\pgfsetfillopacity{0.130000}%
\pgfsetlinewidth{0.000000pt}%
\definecolor{currentstroke}{rgb}{0.000000,0.000000,0.000000}%
\pgfsetstrokecolor{currentstroke}%
\pgfsetdash{}{0pt}%
\pgfpathmoveto{\pgfqpoint{0.739678in}{0.893106in}}%
\pgfpathlineto{\pgfqpoint{0.756115in}{0.880077in}}%
\pgfpathlineto{\pgfqpoint{0.772564in}{0.867691in}}%
\pgfpathlineto{\pgfqpoint{0.781565in}{0.862659in}}%
\pgfpathlineto{\pgfqpoint{0.790524in}{0.857854in}}%
\pgfpathlineto{\pgfqpoint{0.774121in}{0.870705in}}%
\pgfpathlineto{\pgfqpoint{0.757730in}{0.884208in}}%
\pgfpathlineto{\pgfqpoint{0.748726in}{0.888535in}}%
\pgfpathlineto{\pgfqpoint{0.739678in}{0.893106in}}%
\pgfpathclose%
\pgfusepath{fill}%
\end{pgfscope}%
\begin{pgfscope}%
\pgfpathrectangle{\pgfqpoint{0.407592in}{0.222904in}}{\pgfqpoint{1.411163in}{1.411163in}}%
\pgfusepath{clip}%
\pgfsetbuttcap%
\pgfsetroundjoin%
\definecolor{currentfill}{rgb}{0.556863,0.654902,0.847059}%
\pgfsetfillcolor{currentfill}%
\pgfsetfillopacity{0.130000}%
\pgfsetlinewidth{0.000000pt}%
\definecolor{currentstroke}{rgb}{0.000000,0.000000,0.000000}%
\pgfsetstrokecolor{currentstroke}%
\pgfsetdash{}{0pt}%
\pgfpathmoveto{\pgfqpoint{0.622781in}{0.989409in}}%
\pgfpathlineto{\pgfqpoint{0.639202in}{0.973878in}}%
\pgfpathlineto{\pgfqpoint{0.655632in}{0.958746in}}%
\pgfpathlineto{\pgfqpoint{0.664856in}{0.954617in}}%
\pgfpathlineto{\pgfqpoint{0.674037in}{0.950724in}}%
\pgfpathlineto{\pgfqpoint{0.657651in}{0.966290in}}%
\pgfpathlineto{\pgfqpoint{0.641274in}{0.982244in}}%
\pgfpathlineto{\pgfqpoint{0.632050in}{0.985705in}}%
\pgfpathlineto{\pgfqpoint{0.622781in}{0.989409in}}%
\pgfpathclose%
\pgfusepath{fill}%
\end{pgfscope}%
\begin{pgfscope}%
\pgfpathrectangle{\pgfqpoint{0.407592in}{0.222904in}}{\pgfqpoint{1.411163in}{1.411163in}}%
\pgfusepath{clip}%
\pgfsetbuttcap%
\pgfsetroundjoin%
\definecolor{currentfill}{rgb}{0.556863,0.654902,0.847059}%
\pgfsetfillcolor{currentfill}%
\pgfsetfillopacity{0.130000}%
\pgfsetlinewidth{0.000000pt}%
\definecolor{currentstroke}{rgb}{0.000000,0.000000,0.000000}%
\pgfsetstrokecolor{currentstroke}%
\pgfsetdash{}{0pt}%
\pgfpathmoveto{\pgfqpoint{1.803438in}{1.232829in}}%
\pgfpathlineto{\pgfqpoint{1.821584in}{1.244256in}}%
\pgfpathlineto{\pgfqpoint{1.828957in}{1.255330in}}%
\pgfpathlineto{\pgfqpoint{1.836319in}{1.266580in}}%
\pgfpathlineto{\pgfqpoint{1.818197in}{1.255429in}}%
\pgfpathlineto{\pgfqpoint{1.810823in}{1.244041in}}%
\pgfusepath{fill}%
\end{pgfscope}%
\begin{pgfscope}%
\pgfpathrectangle{\pgfqpoint{0.407592in}{0.222904in}}{\pgfqpoint{1.411163in}{1.411163in}}%
\pgfusepath{clip}%
\pgfsetbuttcap%
\pgfsetroundjoin%
\definecolor{currentfill}{rgb}{0.556863,0.654902,0.847059}%
\pgfsetfillcolor{currentfill}%
\pgfsetfillopacity{0.130000}%
\pgfsetlinewidth{0.000000pt}%
\definecolor{currentstroke}{rgb}{0.000000,0.000000,0.000000}%
\pgfsetstrokecolor{currentstroke}%
\pgfsetdash{}{0pt}%
\pgfpathmoveto{\pgfqpoint{1.686825in}{1.122362in}}%
\pgfpathlineto{\pgfqpoint{1.704724in}{1.133923in}}%
\pgfpathlineto{\pgfqpoint{1.722703in}{1.145660in}}%
\pgfpathlineto{\pgfqpoint{1.730178in}{1.155941in}}%
\pgfpathlineto{\pgfqpoint{1.737641in}{1.166444in}}%
\pgfpathlineto{\pgfqpoint{1.719686in}{1.154992in}}%
\pgfpathlineto{\pgfqpoint{1.701811in}{1.143728in}}%
\pgfpathlineto{\pgfqpoint{1.694323in}{1.132931in}}%
\pgfpathlineto{\pgfqpoint{1.686825in}{1.122362in}}%
\pgfpathclose%
\pgfusepath{fill}%
\end{pgfscope}%
\begin{pgfscope}%
\pgfpathrectangle{\pgfqpoint{0.407592in}{0.222904in}}{\pgfqpoint{1.411163in}{1.411163in}}%
\pgfusepath{clip}%
\pgfsetbuttcap%
\pgfsetroundjoin%
\definecolor{currentfill}{rgb}{0.556863,0.654902,0.847059}%
\pgfsetfillcolor{currentfill}%
\pgfsetfillopacity{0.130000}%
\pgfsetlinewidth{0.000000pt}%
\definecolor{currentstroke}{rgb}{0.000000,0.000000,0.000000}%
\pgfsetstrokecolor{currentstroke}%
\pgfsetdash{}{0pt}%
\pgfpathmoveto{\pgfqpoint{0.505514in}{1.095225in}}%
\pgfpathlineto{\pgfqpoint{0.521925in}{1.078226in}}%
\pgfpathlineto{\pgfqpoint{0.538344in}{1.061482in}}%
\pgfpathlineto{\pgfqpoint{0.547789in}{1.058089in}}%
\pgfpathlineto{\pgfqpoint{0.557189in}{1.054924in}}%
\pgfpathlineto{\pgfqpoint{0.540813in}{1.072020in}}%
\pgfpathlineto{\pgfqpoint{0.524446in}{1.089358in}}%
\pgfpathlineto{\pgfqpoint{0.515003in}{1.092177in}}%
\pgfpathlineto{\pgfqpoint{0.505514in}{1.095225in}}%
\pgfpathclose%
\pgfusepath{fill}%
\end{pgfscope}%
\begin{pgfscope}%
\pgfpathrectangle{\pgfqpoint{0.407592in}{0.222904in}}{\pgfqpoint{1.411163in}{1.411163in}}%
\pgfusepath{clip}%
\pgfsetbuttcap%
\pgfsetroundjoin%
\definecolor{currentfill}{rgb}{0.556863,0.654902,0.847059}%
\pgfsetfillcolor{currentfill}%
\pgfsetfillopacity{0.130000}%
\pgfsetlinewidth{0.000000pt}%
\definecolor{currentstroke}{rgb}{0.000000,0.000000,0.000000}%
\pgfsetstrokecolor{currentstroke}%
\pgfsetdash{}{0pt}%
\pgfpathmoveto{\pgfqpoint{1.504778in}{0.957375in}}%
\pgfpathlineto{\pgfqpoint{1.522318in}{0.969261in}}%
\pgfpathlineto{\pgfqpoint{1.539936in}{0.981316in}}%
\pgfpathlineto{\pgfqpoint{1.547600in}{0.989566in}}%
\pgfpathlineto{\pgfqpoint{1.555250in}{0.998161in}}%
\pgfpathlineto{\pgfqpoint{1.537654in}{0.986351in}}%
\pgfpathlineto{\pgfqpoint{1.520136in}{0.974725in}}%
\pgfpathlineto{\pgfqpoint{1.512463in}{0.965867in}}%
\pgfpathlineto{\pgfqpoint{1.504778in}{0.957375in}}%
\pgfpathclose%
\pgfusepath{fill}%
\end{pgfscope}%
\begin{pgfscope}%
\pgfpathrectangle{\pgfqpoint{0.407592in}{0.222904in}}{\pgfqpoint{1.411163in}{1.411163in}}%
\pgfusepath{clip}%
\pgfsetbuttcap%
\pgfsetroundjoin%
\definecolor{currentfill}{rgb}{0.556863,0.654902,0.847059}%
\pgfsetfillcolor{currentfill}%
\pgfsetfillopacity{0.130000}%
\pgfsetlinewidth{0.000000pt}%
\definecolor{currentstroke}{rgb}{0.000000,0.000000,0.000000}%
\pgfsetstrokecolor{currentstroke}%
\pgfsetdash{}{0pt}%
\pgfpathmoveto{\pgfqpoint{1.388726in}{0.866852in}}%
\pgfpathlineto{\pgfqpoint{1.406021in}{0.878524in}}%
\pgfpathlineto{\pgfqpoint{1.423392in}{0.890388in}}%
\pgfpathlineto{\pgfqpoint{1.431187in}{0.896583in}}%
\pgfpathlineto{\pgfqpoint{1.438965in}{0.903229in}}%
\pgfpathlineto{\pgfqpoint{1.421615in}{0.891481in}}%
\pgfpathlineto{\pgfqpoint{1.404341in}{0.879934in}}%
\pgfpathlineto{\pgfqpoint{1.396542in}{0.873147in}}%
\pgfpathlineto{\pgfqpoint{1.388726in}{0.866852in}}%
\pgfpathclose%
\pgfusepath{fill}%
\end{pgfscope}%
\begin{pgfscope}%
\pgfpathrectangle{\pgfqpoint{0.407592in}{0.222904in}}{\pgfqpoint{1.411163in}{1.411163in}}%
\pgfusepath{clip}%
\pgfsetbuttcap%
\pgfsetroundjoin%
\definecolor{currentfill}{rgb}{0.556863,0.654902,0.847059}%
\pgfsetfillcolor{currentfill}%
\pgfsetfillopacity{0.130000}%
\pgfsetlinewidth{0.000000pt}%
\definecolor{currentstroke}{rgb}{0.000000,0.000000,0.000000}%
\pgfsetstrokecolor{currentstroke}%
\pgfsetdash{}{0pt}%
\pgfpathmoveto{\pgfqpoint{0.856266in}{0.814128in}}%
\pgfpathlineto{\pgfqpoint{0.872742in}{0.805395in}}%
\pgfpathlineto{\pgfqpoint{0.889237in}{0.797644in}}%
\pgfpathlineto{\pgfqpoint{0.898008in}{0.791929in}}%
\pgfpathlineto{\pgfqpoint{0.906739in}{0.786422in}}%
\pgfpathlineto{\pgfqpoint{0.890293in}{0.794441in}}%
\pgfpathlineto{\pgfqpoint{0.873867in}{0.803505in}}%
\pgfpathlineto{\pgfqpoint{0.865087in}{0.808704in}}%
\pgfpathlineto{\pgfqpoint{0.856266in}{0.814128in}}%
\pgfpathclose%
\pgfusepath{fill}%
\end{pgfscope}%
\begin{pgfscope}%
\pgfpathrectangle{\pgfqpoint{0.407592in}{0.222904in}}{\pgfqpoint{1.411163in}{1.411163in}}%
\pgfusepath{clip}%
\pgfsetbuttcap%
\pgfsetroundjoin%
\definecolor{currentfill}{rgb}{0.556863,0.654902,0.847059}%
\pgfsetfillcolor{currentfill}%
\pgfsetfillopacity{0.130000}%
\pgfsetlinewidth{0.000000pt}%
\definecolor{currentstroke}{rgb}{0.000000,0.000000,0.000000}%
\pgfsetstrokecolor{currentstroke}%
\pgfsetdash{}{0pt}%
\pgfpathmoveto{\pgfqpoint{1.172900in}{0.756403in}}%
\pgfpathlineto{\pgfqpoint{1.189731in}{0.764011in}}%
\pgfpathlineto{\pgfqpoint{1.206629in}{0.772249in}}%
\pgfpathlineto{\pgfqpoint{1.214730in}{0.772440in}}%
\pgfpathlineto{\pgfqpoint{1.222804in}{0.773194in}}%
\pgfpathlineto{\pgfqpoint{1.205935in}{0.764469in}}%
\pgfpathlineto{\pgfqpoint{1.189135in}{0.756344in}}%
\pgfpathlineto{\pgfqpoint{1.181032in}{0.756073in}}%
\pgfpathlineto{\pgfqpoint{1.172900in}{0.756403in}}%
\pgfpathclose%
\pgfusepath{fill}%
\end{pgfscope}%
\begin{pgfscope}%
\pgfpathrectangle{\pgfqpoint{0.407592in}{0.222904in}}{\pgfqpoint{1.411163in}{1.411163in}}%
\pgfusepath{clip}%
\pgfsetbuttcap%
\pgfsetroundjoin%
\definecolor{currentfill}{rgb}{0.556863,0.654902,0.847059}%
\pgfsetfillcolor{currentfill}%
\pgfsetfillopacity{0.130000}%
\pgfsetlinewidth{0.000000pt}%
\definecolor{currentstroke}{rgb}{0.000000,0.000000,0.000000}%
\pgfsetstrokecolor{currentstroke}%
\pgfsetdash{}{0pt}%
\pgfpathmoveto{\pgfqpoint{1.123009in}{0.746937in}}%
\pgfpathlineto{\pgfqpoint{1.139746in}{0.752406in}}%
\pgfpathlineto{\pgfqpoint{1.156545in}{0.758707in}}%
\pgfpathlineto{\pgfqpoint{1.164738in}{0.757294in}}%
\pgfpathlineto{\pgfqpoint{1.172900in}{0.756403in}}%
\pgfpathlineto{\pgfqpoint{1.156136in}{0.749515in}}%
\pgfpathlineto{\pgfqpoint{1.139435in}{0.743451in}}%
\pgfpathlineto{\pgfqpoint{1.131239in}{0.744928in}}%
\pgfpathlineto{\pgfqpoint{1.123009in}{0.746937in}}%
\pgfpathclose%
\pgfusepath{fill}%
\end{pgfscope}%
\begin{pgfscope}%
\pgfpathrectangle{\pgfqpoint{0.407592in}{0.222904in}}{\pgfqpoint{1.411163in}{1.411163in}}%
\pgfusepath{clip}%
\pgfsetbuttcap%
\pgfsetroundjoin%
\definecolor{currentfill}{rgb}{0.556863,0.654902,0.847059}%
\pgfsetfillcolor{currentfill}%
\pgfsetfillopacity{0.130000}%
\pgfsetlinewidth{0.000000pt}%
\definecolor{currentstroke}{rgb}{0.000000,0.000000,0.000000}%
\pgfsetstrokecolor{currentstroke}%
\pgfsetdash{}{0pt}%
\pgfpathmoveto{\pgfqpoint{0.939697in}{0.773708in}}%
\pgfpathlineto{\pgfqpoint{0.956215in}{0.769072in}}%
\pgfpathlineto{\pgfqpoint{0.972764in}{0.765598in}}%
\pgfpathlineto{\pgfqpoint{0.981353in}{0.760183in}}%
\pgfpathlineto{\pgfqpoint{0.989901in}{0.755014in}}%
\pgfpathlineto{\pgfqpoint{0.973403in}{0.758409in}}%
\pgfpathlineto{\pgfqpoint{0.956935in}{0.763069in}}%
\pgfpathlineto{\pgfqpoint{0.948336in}{0.768266in}}%
\pgfpathlineto{\pgfqpoint{0.939697in}{0.773708in}}%
\pgfpathclose%
\pgfusepath{fill}%
\end{pgfscope}%
\begin{pgfscope}%
\pgfpathrectangle{\pgfqpoint{0.407592in}{0.222904in}}{\pgfqpoint{1.411163in}{1.411163in}}%
\pgfusepath{clip}%
\pgfsetbuttcap%
\pgfsetroundjoin%
\definecolor{currentfill}{rgb}{0.556863,0.654902,0.847059}%
\pgfsetfillcolor{currentfill}%
\pgfsetfillopacity{0.130000}%
\pgfsetlinewidth{0.000000pt}%
\definecolor{currentstroke}{rgb}{0.000000,0.000000,0.000000}%
\pgfsetstrokecolor{currentstroke}%
\pgfsetdash{}{0pt}%
\pgfpathmoveto{\pgfqpoint{1.621066in}{1.058798in}}%
\pgfpathlineto{\pgfqpoint{1.638853in}{1.070580in}}%
\pgfpathlineto{\pgfqpoint{1.656720in}{1.082531in}}%
\pgfpathlineto{\pgfqpoint{1.664264in}{1.092105in}}%
\pgfpathlineto{\pgfqpoint{1.671795in}{1.101941in}}%
\pgfpathlineto{\pgfqpoint{1.653952in}{1.090268in}}%
\pgfpathlineto{\pgfqpoint{1.636188in}{1.078777in}}%
\pgfpathlineto{\pgfqpoint{1.628633in}{1.068652in}}%
\pgfpathlineto{\pgfqpoint{1.621066in}{1.058798in}}%
\pgfpathclose%
\pgfusepath{fill}%
\end{pgfscope}%
\begin{pgfscope}%
\pgfpathrectangle{\pgfqpoint{0.407592in}{0.222904in}}{\pgfqpoint{1.411163in}{1.411163in}}%
\pgfusepath{clip}%
\pgfsetbuttcap%
\pgfsetroundjoin%
\definecolor{currentfill}{rgb}{0.556863,0.654902,0.847059}%
\pgfsetfillcolor{currentfill}%
\pgfsetfillopacity{0.130000}%
\pgfsetlinewidth{0.000000pt}%
\definecolor{currentstroke}{rgb}{0.000000,0.000000,0.000000}%
\pgfsetstrokecolor{currentstroke}%
\pgfsetdash{}{0pt}%
\pgfpathmoveto{\pgfqpoint{0.655632in}{0.958746in}}%
\pgfpathlineto{\pgfqpoint{0.672071in}{0.944046in}}%
\pgfpathlineto{\pgfqpoint{0.688521in}{0.929816in}}%
\pgfpathlineto{\pgfqpoint{0.697701in}{0.925247in}}%
\pgfpathlineto{\pgfqpoint{0.706837in}{0.920904in}}%
\pgfpathlineto{\pgfqpoint{0.690432in}{0.935583in}}%
\pgfpathlineto{\pgfqpoint{0.674037in}{0.950724in}}%
\pgfpathlineto{\pgfqpoint{0.664856in}{0.954617in}}%
\pgfpathlineto{\pgfqpoint{0.655632in}{0.958746in}}%
\pgfpathclose%
\pgfusepath{fill}%
\end{pgfscope}%
\begin{pgfscope}%
\pgfpathrectangle{\pgfqpoint{0.407592in}{0.222904in}}{\pgfqpoint{1.411163in}{1.411163in}}%
\pgfusepath{clip}%
\pgfsetbuttcap%
\pgfsetroundjoin%
\definecolor{currentfill}{rgb}{0.556863,0.654902,0.847059}%
\pgfsetfillcolor{currentfill}%
\pgfsetfillopacity{0.130000}%
\pgfsetlinewidth{0.000000pt}%
\definecolor{currentstroke}{rgb}{0.000000,0.000000,0.000000}%
\pgfsetstrokecolor{currentstroke}%
\pgfsetdash{}{0pt}%
\pgfpathmoveto{\pgfqpoint{1.222804in}{0.773194in}}%
\pgfpathlineto{\pgfqpoint{1.239743in}{0.782442in}}%
\pgfpathlineto{\pgfqpoint{1.256754in}{0.792144in}}%
\pgfpathlineto{\pgfqpoint{1.264775in}{0.793869in}}%
\pgfpathlineto{\pgfqpoint{1.272772in}{0.796163in}}%
\pgfpathlineto{\pgfqpoint{1.255787in}{0.786123in}}%
\pgfpathlineto{\pgfqpoint{1.238873in}{0.776510in}}%
\pgfpathlineto{\pgfqpoint{1.230851in}{0.774542in}}%
\pgfpathlineto{\pgfqpoint{1.222804in}{0.773194in}}%
\pgfpathclose%
\pgfusepath{fill}%
\end{pgfscope}%
\begin{pgfscope}%
\pgfpathrectangle{\pgfqpoint{0.407592in}{0.222904in}}{\pgfqpoint{1.411163in}{1.411163in}}%
\pgfusepath{clip}%
\pgfsetbuttcap%
\pgfsetroundjoin%
\definecolor{currentfill}{rgb}{0.556863,0.654902,0.847059}%
\pgfsetfillcolor{currentfill}%
\pgfsetfillopacity{0.130000}%
\pgfsetlinewidth{0.000000pt}%
\definecolor{currentstroke}{rgb}{0.000000,0.000000,0.000000}%
\pgfsetstrokecolor{currentstroke}%
\pgfsetdash{}{0pt}%
\pgfpathmoveto{\pgfqpoint{1.073067in}{0.745476in}}%
\pgfpathlineto{\pgfqpoint{1.089731in}{0.748417in}}%
\pgfpathlineto{\pgfqpoint{1.106449in}{0.752370in}}%
\pgfpathlineto{\pgfqpoint{1.114747in}{0.749432in}}%
\pgfpathlineto{\pgfqpoint{1.123009in}{0.746937in}}%
\pgfpathlineto{\pgfqpoint{1.106332in}{0.742401in}}%
\pgfpathlineto{\pgfqpoint{1.089710in}{0.738908in}}%
\pgfpathlineto{\pgfqpoint{1.081406in}{0.741977in}}%
\pgfpathlineto{\pgfqpoint{1.073067in}{0.745476in}}%
\pgfpathclose%
\pgfusepath{fill}%
\end{pgfscope}%
\begin{pgfscope}%
\pgfpathrectangle{\pgfqpoint{0.407592in}{0.222904in}}{\pgfqpoint{1.411163in}{1.411163in}}%
\pgfusepath{clip}%
\pgfsetbuttcap%
\pgfsetroundjoin%
\definecolor{currentfill}{rgb}{0.556863,0.654902,0.847059}%
\pgfsetfillcolor{currentfill}%
\pgfsetfillopacity{0.130000}%
\pgfsetlinewidth{0.000000pt}%
\definecolor{currentstroke}{rgb}{0.000000,0.000000,0.000000}%
\pgfsetstrokecolor{currentstroke}%
\pgfsetdash{}{0pt}%
\pgfpathmoveto{\pgfqpoint{0.772564in}{0.867691in}}%
\pgfpathlineto{\pgfqpoint{0.789025in}{0.855999in}}%
\pgfpathlineto{\pgfqpoint{0.805501in}{0.845049in}}%
\pgfpathlineto{\pgfqpoint{0.814454in}{0.839591in}}%
\pgfpathlineto{\pgfqpoint{0.823366in}{0.834341in}}%
\pgfpathlineto{\pgfqpoint{0.806938in}{0.845714in}}%
\pgfpathlineto{\pgfqpoint{0.790524in}{0.857854in}}%
\pgfpathlineto{\pgfqpoint{0.781565in}{0.862659in}}%
\pgfpathlineto{\pgfqpoint{0.772564in}{0.867691in}}%
\pgfpathclose%
\pgfusepath{fill}%
\end{pgfscope}%
\begin{pgfscope}%
\pgfpathrectangle{\pgfqpoint{0.407592in}{0.222904in}}{\pgfqpoint{1.411163in}{1.411163in}}%
\pgfusepath{clip}%
\pgfsetbuttcap%
\pgfsetroundjoin%
\definecolor{currentfill}{rgb}{0.556863,0.654902,0.847059}%
\pgfsetfillcolor{currentfill}%
\pgfsetfillopacity{0.130000}%
\pgfsetlinewidth{0.000000pt}%
\definecolor{currentstroke}{rgb}{0.000000,0.000000,0.000000}%
\pgfsetstrokecolor{currentstroke}%
\pgfsetdash{}{0pt}%
\pgfpathmoveto{\pgfqpoint{0.538344in}{1.061482in}}%
\pgfpathlineto{\pgfqpoint{0.554771in}{1.045015in}}%
\pgfpathlineto{\pgfqpoint{0.571208in}{1.028845in}}%
\pgfpathlineto{\pgfqpoint{0.580610in}{1.025081in}}%
\pgfpathlineto{\pgfqpoint{0.589967in}{1.021543in}}%
\pgfpathlineto{\pgfqpoint{0.573574in}{1.038091in}}%
\pgfpathlineto{\pgfqpoint{0.557189in}{1.054924in}}%
\pgfpathlineto{\pgfqpoint{0.547789in}{1.058089in}}%
\pgfpathlineto{\pgfqpoint{0.538344in}{1.061482in}}%
\pgfpathclose%
\pgfusepath{fill}%
\end{pgfscope}%
\begin{pgfscope}%
\pgfpathrectangle{\pgfqpoint{0.407592in}{0.222904in}}{\pgfqpoint{1.411163in}{1.411163in}}%
\pgfusepath{clip}%
\pgfsetbuttcap%
\pgfsetroundjoin%
\definecolor{currentfill}{rgb}{0.556863,0.654902,0.847059}%
\pgfsetfillcolor{currentfill}%
\pgfsetfillopacity{0.130000}%
\pgfsetlinewidth{0.000000pt}%
\definecolor{currentstroke}{rgb}{0.000000,0.000000,0.000000}%
\pgfsetstrokecolor{currentstroke}%
\pgfsetdash{}{0pt}%
\pgfpathmoveto{\pgfqpoint{1.737641in}{1.166444in}}%
\pgfpathlineto{\pgfqpoint{1.755676in}{1.178074in}}%
\pgfpathlineto{\pgfqpoint{1.773791in}{1.189874in}}%
\pgfpathlineto{\pgfqpoint{1.781219in}{1.200315in}}%
\pgfpathlineto{\pgfqpoint{1.788636in}{1.210959in}}%
\pgfpathlineto{\pgfqpoint{1.770545in}{1.199433in}}%
\pgfpathlineto{\pgfqpoint{1.752534in}{1.188087in}}%
\pgfpathlineto{\pgfqpoint{1.745093in}{1.177162in}}%
\pgfpathlineto{\pgfqpoint{1.737641in}{1.166444in}}%
\pgfpathclose%
\pgfusepath{fill}%
\end{pgfscope}%
\begin{pgfscope}%
\pgfpathrectangle{\pgfqpoint{0.407592in}{0.222904in}}{\pgfqpoint{1.411163in}{1.411163in}}%
\pgfusepath{clip}%
\pgfsetbuttcap%
\pgfsetroundjoin%
\definecolor{currentfill}{rgb}{0.556863,0.654902,0.847059}%
\pgfsetfillcolor{currentfill}%
\pgfsetfillopacity{0.130000}%
\pgfsetlinewidth{0.000000pt}%
\definecolor{currentstroke}{rgb}{0.000000,0.000000,0.000000}%
\pgfsetstrokecolor{currentstroke}%
\pgfsetdash{}{0pt}%
\pgfpathmoveto{\pgfqpoint{1.272772in}{0.796163in}}%
\pgfpathlineto{\pgfqpoint{1.289831in}{0.806573in}}%
\pgfpathlineto{\pgfqpoint{1.306964in}{0.817310in}}%
\pgfpathlineto{\pgfqpoint{1.314915in}{0.820412in}}%
\pgfpathlineto{\pgfqpoint{1.322844in}{0.824057in}}%
\pgfpathlineto{\pgfqpoint{1.305734in}{0.813134in}}%
\pgfpathlineto{\pgfqpoint{1.288698in}{0.802519in}}%
\pgfpathlineto{\pgfqpoint{1.280746in}{0.799041in}}%
\pgfpathlineto{\pgfqpoint{1.272772in}{0.796163in}}%
\pgfpathclose%
\pgfusepath{fill}%
\end{pgfscope}%
\begin{pgfscope}%
\pgfpathrectangle{\pgfqpoint{0.407592in}{0.222904in}}{\pgfqpoint{1.411163in}{1.411163in}}%
\pgfusepath{clip}%
\pgfsetbuttcap%
\pgfsetroundjoin%
\definecolor{currentfill}{rgb}{0.556863,0.654902,0.847059}%
\pgfsetfillcolor{currentfill}%
\pgfsetfillopacity{0.130000}%
\pgfsetlinewidth{0.000000pt}%
\definecolor{currentstroke}{rgb}{0.000000,0.000000,0.000000}%
\pgfsetstrokecolor{currentstroke}%
\pgfsetdash{}{0pt}%
\pgfpathmoveto{\pgfqpoint{1.438965in}{0.903229in}}%
\pgfpathlineto{\pgfqpoint{1.456392in}{0.915160in}}%
\pgfpathlineto{\pgfqpoint{1.473898in}{0.927261in}}%
\pgfpathlineto{\pgfqpoint{1.481639in}{0.934193in}}%
\pgfpathlineto{\pgfqpoint{1.489366in}{0.941528in}}%
\pgfpathlineto{\pgfqpoint{1.471883in}{0.929595in}}%
\pgfpathlineto{\pgfqpoint{1.454476in}{0.917843in}}%
\pgfpathlineto{\pgfqpoint{1.446728in}{0.910318in}}%
\pgfpathlineto{\pgfqpoint{1.438965in}{0.903229in}}%
\pgfpathclose%
\pgfusepath{fill}%
\end{pgfscope}%
\begin{pgfscope}%
\pgfpathrectangle{\pgfqpoint{0.407592in}{0.222904in}}{\pgfqpoint{1.411163in}{1.411163in}}%
\pgfusepath{clip}%
\pgfsetbuttcap%
\pgfsetroundjoin%
\definecolor{currentfill}{rgb}{0.556863,0.654902,0.847059}%
\pgfsetfillcolor{currentfill}%
\pgfsetfillopacity{0.130000}%
\pgfsetlinewidth{0.000000pt}%
\definecolor{currentstroke}{rgb}{0.000000,0.000000,0.000000}%
\pgfsetstrokecolor{currentstroke}%
\pgfsetdash{}{0pt}%
\pgfpathmoveto{\pgfqpoint{1.555250in}{0.998161in}}%
\pgfpathlineto{\pgfqpoint{1.572925in}{1.010141in}}%
\pgfpathlineto{\pgfqpoint{1.590679in}{1.022282in}}%
\pgfpathlineto{\pgfqpoint{1.598294in}{1.030958in}}%
\pgfpathlineto{\pgfqpoint{1.605897in}{1.039942in}}%
\pgfpathlineto{\pgfqpoint{1.588166in}{1.028055in}}%
\pgfpathlineto{\pgfqpoint{1.570514in}{1.016342in}}%
\pgfpathlineto{\pgfqpoint{1.562888in}{1.007089in}}%
\pgfpathlineto{\pgfqpoint{1.555250in}{0.998161in}}%
\pgfpathclose%
\pgfusepath{fill}%
\end{pgfscope}%
\begin{pgfscope}%
\pgfpathrectangle{\pgfqpoint{0.407592in}{0.222904in}}{\pgfqpoint{1.411163in}{1.411163in}}%
\pgfusepath{clip}%
\pgfsetbuttcap%
\pgfsetroundjoin%
\definecolor{currentfill}{rgb}{0.556863,0.654902,0.847059}%
\pgfsetfillcolor{currentfill}%
\pgfsetfillopacity{0.130000}%
\pgfsetlinewidth{0.000000pt}%
\definecolor{currentstroke}{rgb}{0.000000,0.000000,0.000000}%
\pgfsetstrokecolor{currentstroke}%
\pgfsetdash{}{0pt}%
\pgfpathmoveto{\pgfqpoint{1.023005in}{0.751955in}}%
\pgfpathlineto{\pgfqpoint{1.039619in}{0.752215in}}%
\pgfpathlineto{\pgfqpoint{1.056278in}{0.753592in}}%
\pgfpathlineto{\pgfqpoint{1.064691in}{0.749360in}}%
\pgfpathlineto{\pgfqpoint{1.073067in}{0.745476in}}%
\pgfpathlineto{\pgfqpoint{1.056454in}{0.743637in}}%
\pgfpathlineto{\pgfqpoint{1.039888in}{0.742986in}}%
\pgfpathlineto{\pgfqpoint{1.031465in}{0.747308in}}%
\pgfpathlineto{\pgfqpoint{1.023005in}{0.751955in}}%
\pgfpathclose%
\pgfusepath{fill}%
\end{pgfscope}%
\begin{pgfscope}%
\pgfpathrectangle{\pgfqpoint{0.407592in}{0.222904in}}{\pgfqpoint{1.411163in}{1.411163in}}%
\pgfusepath{clip}%
\pgfsetbuttcap%
\pgfsetroundjoin%
\definecolor{currentfill}{rgb}{0.556863,0.654902,0.847059}%
\pgfsetfillcolor{currentfill}%
\pgfsetfillopacity{0.130000}%
\pgfsetlinewidth{0.000000pt}%
\definecolor{currentstroke}{rgb}{0.000000,0.000000,0.000000}%
\pgfsetstrokecolor{currentstroke}%
\pgfsetdash{}{0pt}%
\pgfpathmoveto{\pgfqpoint{1.788636in}{1.210959in}}%
\pgfpathlineto{\pgfqpoint{1.806807in}{1.222657in}}%
\pgfpathlineto{\pgfqpoint{1.814201in}{1.233363in}}%
\pgfpathlineto{\pgfqpoint{1.821584in}{1.244256in}}%
\pgfpathlineto{\pgfqpoint{1.803438in}{1.232829in}}%
\pgfpathlineto{\pgfqpoint{1.796042in}{1.221799in}}%
\pgfusepath{fill}%
\end{pgfscope}%
\begin{pgfscope}%
\pgfpathrectangle{\pgfqpoint{0.407592in}{0.222904in}}{\pgfqpoint{1.411163in}{1.411163in}}%
\pgfusepath{clip}%
\pgfsetbuttcap%
\pgfsetroundjoin%
\definecolor{currentfill}{rgb}{0.556863,0.654902,0.847059}%
\pgfsetfillcolor{currentfill}%
\pgfsetfillopacity{0.130000}%
\pgfsetlinewidth{0.000000pt}%
\definecolor{currentstroke}{rgb}{0.000000,0.000000,0.000000}%
\pgfsetstrokecolor{currentstroke}%
\pgfsetdash{}{0pt}%
\pgfpathmoveto{\pgfqpoint{0.420674in}{1.170931in}}%
\pgfpathlineto{\pgfqpoint{0.437092in}{1.153379in}}%
\pgfpathlineto{\pgfqpoint{0.453519in}{1.136025in}}%
\pgfpathlineto{\pgfqpoint{0.463143in}{1.132863in}}%
\pgfpathlineto{\pgfqpoint{0.472720in}{1.129914in}}%
\pgfpathlineto{\pgfqpoint{0.456335in}{1.147569in}}%
\pgfpathlineto{\pgfqpoint{0.439960in}{1.165411in}}%
\pgfpathlineto{\pgfqpoint{0.430340in}{1.168065in}}%
\pgfpathlineto{\pgfqpoint{0.420674in}{1.170931in}}%
\pgfpathclose%
\pgfusepath{fill}%
\end{pgfscope}%
\begin{pgfscope}%
\pgfpathrectangle{\pgfqpoint{0.407592in}{0.222904in}}{\pgfqpoint{1.411163in}{1.411163in}}%
\pgfusepath{clip}%
\pgfsetbuttcap%
\pgfsetroundjoin%
\definecolor{currentfill}{rgb}{0.556863,0.654902,0.847059}%
\pgfsetfillcolor{currentfill}%
\pgfsetfillopacity{0.130000}%
\pgfsetlinewidth{0.000000pt}%
\definecolor{currentstroke}{rgb}{0.000000,0.000000,0.000000}%
\pgfsetstrokecolor{currentstroke}%
\pgfsetdash{}{0pt}%
\pgfpathmoveto{\pgfqpoint{0.688521in}{0.929816in}}%
\pgfpathlineto{\pgfqpoint{0.704981in}{0.916092in}}%
\pgfpathlineto{\pgfqpoint{0.721453in}{0.902916in}}%
\pgfpathlineto{\pgfqpoint{0.730587in}{0.897905in}}%
\pgfpathlineto{\pgfqpoint{0.739678in}{0.893106in}}%
\pgfpathlineto{\pgfqpoint{0.723253in}{0.906731in}}%
\pgfpathlineto{\pgfqpoint{0.706837in}{0.920904in}}%
\pgfpathlineto{\pgfqpoint{0.697701in}{0.925247in}}%
\pgfpathlineto{\pgfqpoint{0.688521in}{0.929816in}}%
\pgfpathclose%
\pgfusepath{fill}%
\end{pgfscope}%
\begin{pgfscope}%
\pgfpathrectangle{\pgfqpoint{0.407592in}{0.222904in}}{\pgfqpoint{1.411163in}{1.411163in}}%
\pgfusepath{clip}%
\pgfsetbuttcap%
\pgfsetroundjoin%
\definecolor{currentfill}{rgb}{0.556863,0.654902,0.847059}%
\pgfsetfillcolor{currentfill}%
\pgfsetfillopacity{0.130000}%
\pgfsetlinewidth{0.000000pt}%
\definecolor{currentstroke}{rgb}{0.000000,0.000000,0.000000}%
\pgfsetstrokecolor{currentstroke}%
\pgfsetdash{}{0pt}%
\pgfpathmoveto{\pgfqpoint{1.671795in}{1.101941in}}%
\pgfpathlineto{\pgfqpoint{1.689718in}{1.113785in}}%
\pgfpathlineto{\pgfqpoint{1.707721in}{1.125793in}}%
\pgfpathlineto{\pgfqpoint{1.715218in}{1.135608in}}%
\pgfpathlineto{\pgfqpoint{1.722703in}{1.145660in}}%
\pgfpathlineto{\pgfqpoint{1.704724in}{1.133923in}}%
\pgfpathlineto{\pgfqpoint{1.686825in}{1.122362in}}%
\pgfpathlineto{\pgfqpoint{1.679316in}{1.112029in}}%
\pgfpathlineto{\pgfqpoint{1.671795in}{1.101941in}}%
\pgfpathclose%
\pgfusepath{fill}%
\end{pgfscope}%
\begin{pgfscope}%
\pgfpathrectangle{\pgfqpoint{0.407592in}{0.222904in}}{\pgfqpoint{1.411163in}{1.411163in}}%
\pgfusepath{clip}%
\pgfsetbuttcap%
\pgfsetroundjoin%
\definecolor{currentfill}{rgb}{0.556863,0.654902,0.847059}%
\pgfsetfillcolor{currentfill}%
\pgfsetfillopacity{0.130000}%
\pgfsetlinewidth{0.000000pt}%
\definecolor{currentstroke}{rgb}{0.000000,0.000000,0.000000}%
\pgfsetstrokecolor{currentstroke}%
\pgfsetdash{}{0pt}%
\pgfpathmoveto{\pgfqpoint{1.322844in}{0.824057in}}%
\pgfpathlineto{\pgfqpoint{1.340029in}{0.835252in}}%
\pgfpathlineto{\pgfqpoint{1.357291in}{0.846689in}}%
\pgfpathlineto{\pgfqpoint{1.365177in}{0.850974in}}%
\pgfpathlineto{\pgfqpoint{1.373045in}{0.855762in}}%
\pgfpathlineto{\pgfqpoint{1.355806in}{0.844265in}}%
\pgfpathlineto{\pgfqpoint{1.338642in}{0.833004in}}%
\pgfpathlineto{\pgfqpoint{1.330753in}{0.828253in}}%
\pgfpathlineto{\pgfqpoint{1.322844in}{0.824057in}}%
\pgfpathclose%
\pgfusepath{fill}%
\end{pgfscope}%
\begin{pgfscope}%
\pgfpathrectangle{\pgfqpoint{0.407592in}{0.222904in}}{\pgfqpoint{1.411163in}{1.411163in}}%
\pgfusepath{clip}%
\pgfsetbuttcap%
\pgfsetroundjoin%
\definecolor{currentfill}{rgb}{0.556863,0.654902,0.847059}%
\pgfsetfillcolor{currentfill}%
\pgfsetfillopacity{0.130000}%
\pgfsetlinewidth{0.000000pt}%
\definecolor{currentstroke}{rgb}{0.000000,0.000000,0.000000}%
\pgfsetstrokecolor{currentstroke}%
\pgfsetdash{}{0pt}%
\pgfpathmoveto{\pgfqpoint{0.571208in}{1.028845in}}%
\pgfpathlineto{\pgfqpoint{0.587654in}{1.012999in}}%
\pgfpathlineto{\pgfqpoint{0.604109in}{0.997501in}}%
\pgfpathlineto{\pgfqpoint{0.613468in}{0.993344in}}%
\pgfpathlineto{\pgfqpoint{0.622781in}{0.989409in}}%
\pgfpathlineto{\pgfqpoint{0.606370in}{1.005307in}}%
\pgfpathlineto{\pgfqpoint{0.589967in}{1.021543in}}%
\pgfpathlineto{\pgfqpoint{0.580610in}{1.025081in}}%
\pgfpathlineto{\pgfqpoint{0.571208in}{1.028845in}}%
\pgfpathclose%
\pgfusepath{fill}%
\end{pgfscope}%
\begin{pgfscope}%
\pgfpathrectangle{\pgfqpoint{0.407592in}{0.222904in}}{\pgfqpoint{1.411163in}{1.411163in}}%
\pgfusepath{clip}%
\pgfsetbuttcap%
\pgfsetroundjoin%
\definecolor{currentfill}{rgb}{0.556863,0.654902,0.847059}%
\pgfsetfillcolor{currentfill}%
\pgfsetfillopacity{0.130000}%
\pgfsetlinewidth{0.000000pt}%
\definecolor{currentstroke}{rgb}{0.000000,0.000000,0.000000}%
\pgfsetstrokecolor{currentstroke}%
\pgfsetdash{}{0pt}%
\pgfpathmoveto{\pgfqpoint{0.889237in}{0.797644in}}%
\pgfpathlineto{\pgfqpoint{0.905755in}{0.790911in}}%
\pgfpathlineto{\pgfqpoint{0.922297in}{0.785227in}}%
\pgfpathlineto{\pgfqpoint{0.931017in}{0.779369in}}%
\pgfpathlineto{\pgfqpoint{0.939697in}{0.773708in}}%
\pgfpathlineto{\pgfqpoint{0.923206in}{0.779499in}}%
\pgfpathlineto{\pgfqpoint{0.906739in}{0.786422in}}%
\pgfpathlineto{\pgfqpoint{0.898008in}{0.791929in}}%
\pgfpathlineto{\pgfqpoint{0.889237in}{0.797644in}}%
\pgfpathclose%
\pgfusepath{fill}%
\end{pgfscope}%
\begin{pgfscope}%
\pgfpathrectangle{\pgfqpoint{0.407592in}{0.222904in}}{\pgfqpoint{1.411163in}{1.411163in}}%
\pgfusepath{clip}%
\pgfsetbuttcap%
\pgfsetroundjoin%
\definecolor{currentfill}{rgb}{0.556863,0.654902,0.847059}%
\pgfsetfillcolor{currentfill}%
\pgfsetfillopacity{0.130000}%
\pgfsetlinewidth{0.000000pt}%
\definecolor{currentstroke}{rgb}{0.000000,0.000000,0.000000}%
\pgfsetstrokecolor{currentstroke}%
\pgfsetdash{}{0pt}%
\pgfpathmoveto{\pgfqpoint{0.805501in}{0.845049in}}%
\pgfpathlineto{\pgfqpoint{0.821991in}{0.834892in}}%
\pgfpathlineto{\pgfqpoint{0.838498in}{0.825573in}}%
\pgfpathlineto{\pgfqpoint{0.847403in}{0.819756in}}%
\pgfpathlineto{\pgfqpoint{0.856266in}{0.814128in}}%
\pgfpathlineto{\pgfqpoint{0.839808in}{0.823794in}}%
\pgfpathlineto{\pgfqpoint{0.823366in}{0.834341in}}%
\pgfpathlineto{\pgfqpoint{0.814454in}{0.839591in}}%
\pgfpathlineto{\pgfqpoint{0.805501in}{0.845049in}}%
\pgfpathclose%
\pgfusepath{fill}%
\end{pgfscope}%
\begin{pgfscope}%
\pgfpathrectangle{\pgfqpoint{0.407592in}{0.222904in}}{\pgfqpoint{1.411163in}{1.411163in}}%
\pgfusepath{clip}%
\pgfsetbuttcap%
\pgfsetroundjoin%
\definecolor{currentfill}{rgb}{0.556863,0.654902,0.847059}%
\pgfsetfillcolor{currentfill}%
\pgfsetfillopacity{0.130000}%
\pgfsetlinewidth{0.000000pt}%
\definecolor{currentstroke}{rgb}{0.000000,0.000000,0.000000}%
\pgfsetstrokecolor{currentstroke}%
\pgfsetdash{}{0pt}%
\pgfpathmoveto{\pgfqpoint{1.489366in}{0.941528in}}%
\pgfpathlineto{\pgfqpoint{1.506928in}{0.953627in}}%
\pgfpathlineto{\pgfqpoint{1.524568in}{0.965882in}}%
\pgfpathlineto{\pgfqpoint{1.532259in}{0.973418in}}%
\pgfpathlineto{\pgfqpoint{1.539936in}{0.981316in}}%
\pgfpathlineto{\pgfqpoint{1.522318in}{0.969261in}}%
\pgfpathlineto{\pgfqpoint{1.504778in}{0.957375in}}%
\pgfpathlineto{\pgfqpoint{1.497079in}{0.949258in}}%
\pgfpathlineto{\pgfqpoint{1.489366in}{0.941528in}}%
\pgfpathclose%
\pgfusepath{fill}%
\end{pgfscope}%
\begin{pgfscope}%
\pgfpathrectangle{\pgfqpoint{0.407592in}{0.222904in}}{\pgfqpoint{1.411163in}{1.411163in}}%
\pgfusepath{clip}%
\pgfsetbuttcap%
\pgfsetroundjoin%
\definecolor{currentfill}{rgb}{0.556863,0.654902,0.847059}%
\pgfsetfillcolor{currentfill}%
\pgfsetfillopacity{0.130000}%
\pgfsetlinewidth{0.000000pt}%
\definecolor{currentstroke}{rgb}{0.000000,0.000000,0.000000}%
\pgfsetstrokecolor{currentstroke}%
\pgfsetdash{}{0pt}%
\pgfpathmoveto{\pgfqpoint{0.453519in}{1.136025in}}%
\pgfpathlineto{\pgfqpoint{0.469955in}{1.118886in}}%
\pgfpathlineto{\pgfqpoint{0.486399in}{1.101976in}}%
\pgfpathlineto{\pgfqpoint{0.495980in}{1.098494in}}%
\pgfpathlineto{\pgfqpoint{0.505514in}{1.095225in}}%
\pgfpathlineto{\pgfqpoint{0.489113in}{1.112460in}}%
\pgfpathlineto{\pgfqpoint{0.472720in}{1.129914in}}%
\pgfpathlineto{\pgfqpoint{0.463143in}{1.132863in}}%
\pgfpathlineto{\pgfqpoint{0.453519in}{1.136025in}}%
\pgfpathclose%
\pgfusepath{fill}%
\end{pgfscope}%
\begin{pgfscope}%
\pgfpathrectangle{\pgfqpoint{0.407592in}{0.222904in}}{\pgfqpoint{1.411163in}{1.411163in}}%
\pgfusepath{clip}%
\pgfsetbuttcap%
\pgfsetroundjoin%
\definecolor{currentfill}{rgb}{0.556863,0.654902,0.847059}%
\pgfsetfillcolor{currentfill}%
\pgfsetfillopacity{0.130000}%
\pgfsetlinewidth{0.000000pt}%
\definecolor{currentstroke}{rgb}{0.000000,0.000000,0.000000}%
\pgfsetstrokecolor{currentstroke}%
\pgfsetdash{}{0pt}%
\pgfpathmoveto{\pgfqpoint{0.972764in}{0.765598in}}%
\pgfpathlineto{\pgfqpoint{0.989348in}{0.763279in}}%
\pgfpathlineto{\pgfqpoint{1.005969in}{0.762093in}}%
\pgfpathlineto{\pgfqpoint{1.014506in}{0.756894in}}%
\pgfpathlineto{\pgfqpoint{1.023005in}{0.751955in}}%
\pgfpathlineto{\pgfqpoint{1.006434in}{0.752873in}}%
\pgfpathlineto{\pgfqpoint{0.989901in}{0.755014in}}%
\pgfpathlineto{\pgfqpoint{0.981353in}{0.760183in}}%
\pgfpathlineto{\pgfqpoint{0.972764in}{0.765598in}}%
\pgfpathclose%
\pgfusepath{fill}%
\end{pgfscope}%
\begin{pgfscope}%
\pgfpathrectangle{\pgfqpoint{0.407592in}{0.222904in}}{\pgfqpoint{1.411163in}{1.411163in}}%
\pgfusepath{clip}%
\pgfsetbuttcap%
\pgfsetroundjoin%
\definecolor{currentfill}{rgb}{0.556863,0.654902,0.847059}%
\pgfsetfillcolor{currentfill}%
\pgfsetfillopacity{0.130000}%
\pgfsetlinewidth{0.000000pt}%
\definecolor{currentstroke}{rgb}{0.000000,0.000000,0.000000}%
\pgfsetstrokecolor{currentstroke}%
\pgfsetdash{}{0pt}%
\pgfpathmoveto{\pgfqpoint{1.605897in}{1.039942in}}%
\pgfpathlineto{\pgfqpoint{1.623707in}{1.051993in}}%
\pgfpathlineto{\pgfqpoint{1.641597in}{1.064199in}}%
\pgfpathlineto{\pgfqpoint{1.649165in}{1.073227in}}%
\pgfpathlineto{\pgfqpoint{1.656720in}{1.082531in}}%
\pgfpathlineto{\pgfqpoint{1.638853in}{1.070580in}}%
\pgfpathlineto{\pgfqpoint{1.621066in}{1.058798in}}%
\pgfpathlineto{\pgfqpoint{1.613488in}{1.049225in}}%
\pgfpathlineto{\pgfqpoint{1.605897in}{1.039942in}}%
\pgfpathclose%
\pgfusepath{fill}%
\end{pgfscope}%
\begin{pgfscope}%
\pgfpathrectangle{\pgfqpoint{0.407592in}{0.222904in}}{\pgfqpoint{1.411163in}{1.411163in}}%
\pgfusepath{clip}%
\pgfsetbuttcap%
\pgfsetroundjoin%
\definecolor{currentfill}{rgb}{0.556863,0.654902,0.847059}%
\pgfsetfillcolor{currentfill}%
\pgfsetfillopacity{0.130000}%
\pgfsetlinewidth{0.000000pt}%
\definecolor{currentstroke}{rgb}{0.000000,0.000000,0.000000}%
\pgfsetstrokecolor{currentstroke}%
\pgfsetdash{}{0pt}%
\pgfpathmoveto{\pgfqpoint{1.373045in}{0.855762in}}%
\pgfpathlineto{\pgfqpoint{1.390361in}{0.867471in}}%
\pgfpathlineto{\pgfqpoint{1.407754in}{0.879371in}}%
\pgfpathlineto{\pgfqpoint{1.415582in}{0.884650in}}%
\pgfpathlineto{\pgfqpoint{1.423392in}{0.890388in}}%
\pgfpathlineto{\pgfqpoint{1.406021in}{0.878524in}}%
\pgfpathlineto{\pgfqpoint{1.388726in}{0.866852in}}%
\pgfpathlineto{\pgfqpoint{1.380894in}{0.861056in}}%
\pgfpathlineto{\pgfqpoint{1.373045in}{0.855762in}}%
\pgfpathclose%
\pgfusepath{fill}%
\end{pgfscope}%
\begin{pgfscope}%
\pgfpathrectangle{\pgfqpoint{0.407592in}{0.222904in}}{\pgfqpoint{1.411163in}{1.411163in}}%
\pgfusepath{clip}%
\pgfsetbuttcap%
\pgfsetroundjoin%
\definecolor{currentfill}{rgb}{0.556863,0.654902,0.847059}%
\pgfsetfillcolor{currentfill}%
\pgfsetfillopacity{0.130000}%
\pgfsetlinewidth{0.000000pt}%
\definecolor{currentstroke}{rgb}{0.000000,0.000000,0.000000}%
\pgfsetstrokecolor{currentstroke}%
\pgfsetdash{}{0pt}%
\pgfpathmoveto{\pgfqpoint{0.604109in}{0.997501in}}%
\pgfpathlineto{\pgfqpoint{0.620574in}{0.982379in}}%
\pgfpathlineto{\pgfqpoint{0.637048in}{0.967665in}}%
\pgfpathlineto{\pgfqpoint{0.646362in}{0.963099in}}%
\pgfpathlineto{\pgfqpoint{0.655632in}{0.958746in}}%
\pgfpathlineto{\pgfqpoint{0.639202in}{0.973878in}}%
\pgfpathlineto{\pgfqpoint{0.622781in}{0.989409in}}%
\pgfpathlineto{\pgfqpoint{0.613468in}{0.993344in}}%
\pgfpathlineto{\pgfqpoint{0.604109in}{0.997501in}}%
\pgfpathclose%
\pgfusepath{fill}%
\end{pgfscope}%
\begin{pgfscope}%
\pgfpathrectangle{\pgfqpoint{0.407592in}{0.222904in}}{\pgfqpoint{1.411163in}{1.411163in}}%
\pgfusepath{clip}%
\pgfsetbuttcap%
\pgfsetroundjoin%
\definecolor{currentfill}{rgb}{0.556863,0.654902,0.847059}%
\pgfsetfillcolor{currentfill}%
\pgfsetfillopacity{0.130000}%
\pgfsetlinewidth{0.000000pt}%
\definecolor{currentstroke}{rgb}{0.000000,0.000000,0.000000}%
\pgfsetstrokecolor{currentstroke}%
\pgfsetdash{}{0pt}%
\pgfpathmoveto{\pgfqpoint{0.721453in}{0.902916in}}%
\pgfpathlineto{\pgfqpoint{0.737936in}{0.890328in}}%
\pgfpathlineto{\pgfqpoint{0.754432in}{0.878371in}}%
\pgfpathlineto{\pgfqpoint{0.763519in}{0.872933in}}%
\pgfpathlineto{\pgfqpoint{0.772564in}{0.867691in}}%
\pgfpathlineto{\pgfqpoint{0.756115in}{0.880077in}}%
\pgfpathlineto{\pgfqpoint{0.739678in}{0.893106in}}%
\pgfpathlineto{\pgfqpoint{0.730587in}{0.897905in}}%
\pgfpathlineto{\pgfqpoint{0.721453in}{0.902916in}}%
\pgfpathclose%
\pgfusepath{fill}%
\end{pgfscope}%
\begin{pgfscope}%
\pgfpathrectangle{\pgfqpoint{0.407592in}{0.222904in}}{\pgfqpoint{1.411163in}{1.411163in}}%
\pgfusepath{clip}%
\pgfsetbuttcap%
\pgfsetroundjoin%
\definecolor{currentfill}{rgb}{0.556863,0.654902,0.847059}%
\pgfsetfillcolor{currentfill}%
\pgfsetfillopacity{0.130000}%
\pgfsetlinewidth{0.000000pt}%
\definecolor{currentstroke}{rgb}{0.000000,0.000000,0.000000}%
\pgfsetstrokecolor{currentstroke}%
\pgfsetdash{}{0pt}%
\pgfpathmoveto{\pgfqpoint{1.722703in}{1.145660in}}%
\pgfpathlineto{\pgfqpoint{1.740763in}{1.157563in}}%
\pgfpathlineto{\pgfqpoint{1.758903in}{1.169625in}}%
\pgfpathlineto{\pgfqpoint{1.766353in}{1.179642in}}%
\pgfpathlineto{\pgfqpoint{1.773791in}{1.189874in}}%
\pgfpathlineto{\pgfqpoint{1.755676in}{1.178074in}}%
\pgfpathlineto{\pgfqpoint{1.737641in}{1.166444in}}%
\pgfpathlineto{\pgfqpoint{1.730178in}{1.155941in}}%
\pgfpathlineto{\pgfqpoint{1.722703in}{1.145660in}}%
\pgfpathclose%
\pgfusepath{fill}%
\end{pgfscope}%
\begin{pgfscope}%
\pgfpathrectangle{\pgfqpoint{0.407592in}{0.222904in}}{\pgfqpoint{1.411163in}{1.411163in}}%
\pgfusepath{clip}%
\pgfsetbuttcap%
\pgfsetroundjoin%
\definecolor{currentfill}{rgb}{0.556863,0.654902,0.847059}%
\pgfsetfillcolor{currentfill}%
\pgfsetfillopacity{0.130000}%
\pgfsetlinewidth{0.000000pt}%
\definecolor{currentstroke}{rgb}{0.000000,0.000000,0.000000}%
\pgfsetstrokecolor{currentstroke}%
\pgfsetdash{}{0pt}%
\pgfpathmoveto{\pgfqpoint{0.486399in}{1.101976in}}%
\pgfpathlineto{\pgfqpoint{0.502853in}{1.085314in}}%
\pgfpathlineto{\pgfqpoint{0.519315in}{1.068919in}}%
\pgfpathlineto{\pgfqpoint{0.528852in}{1.065095in}}%
\pgfpathlineto{\pgfqpoint{0.538344in}{1.061482in}}%
\pgfpathlineto{\pgfqpoint{0.521925in}{1.078226in}}%
\pgfpathlineto{\pgfqpoint{0.505514in}{1.095225in}}%
\pgfpathlineto{\pgfqpoint{0.495980in}{1.098494in}}%
\pgfpathlineto{\pgfqpoint{0.486399in}{1.101976in}}%
\pgfpathclose%
\pgfusepath{fill}%
\end{pgfscope}%
\begin{pgfscope}%
\pgfpathrectangle{\pgfqpoint{0.407592in}{0.222904in}}{\pgfqpoint{1.411163in}{1.411163in}}%
\pgfusepath{clip}%
\pgfsetbuttcap%
\pgfsetroundjoin%
\definecolor{currentfill}{rgb}{0.556863,0.654902,0.847059}%
\pgfsetfillcolor{currentfill}%
\pgfsetfillopacity{0.130000}%
\pgfsetlinewidth{0.000000pt}%
\definecolor{currentstroke}{rgb}{0.000000,0.000000,0.000000}%
\pgfsetstrokecolor{currentstroke}%
\pgfsetdash{}{0pt}%
\pgfpathmoveto{\pgfqpoint{1.156545in}{0.758707in}}%
\pgfpathlineto{\pgfqpoint{1.173409in}{0.765747in}}%
\pgfpathlineto{\pgfqpoint{1.190339in}{0.773441in}}%
\pgfpathlineto{\pgfqpoint{1.198499in}{0.772594in}}%
\pgfpathlineto{\pgfqpoint{1.206629in}{0.772249in}}%
\pgfpathlineto{\pgfqpoint{1.189731in}{0.764011in}}%
\pgfpathlineto{\pgfqpoint{1.172900in}{0.756403in}}%
\pgfpathlineto{\pgfqpoint{1.164738in}{0.757294in}}%
\pgfpathlineto{\pgfqpoint{1.156545in}{0.758707in}}%
\pgfpathclose%
\pgfusepath{fill}%
\end{pgfscope}%
\begin{pgfscope}%
\pgfpathrectangle{\pgfqpoint{0.407592in}{0.222904in}}{\pgfqpoint{1.411163in}{1.411163in}}%
\pgfusepath{clip}%
\pgfsetbuttcap%
\pgfsetroundjoin%
\definecolor{currentfill}{rgb}{0.556863,0.654902,0.847059}%
\pgfsetfillcolor{currentfill}%
\pgfsetfillopacity{0.130000}%
\pgfsetlinewidth{0.000000pt}%
\definecolor{currentstroke}{rgb}{0.000000,0.000000,0.000000}%
\pgfsetstrokecolor{currentstroke}%
\pgfsetdash{}{0pt}%
\pgfpathmoveto{\pgfqpoint{1.206629in}{0.772249in}}%
\pgfpathlineto{\pgfqpoint{1.223596in}{0.781041in}}%
\pgfpathlineto{\pgfqpoint{1.240634in}{0.790318in}}%
\pgfpathlineto{\pgfqpoint{1.248707in}{0.790967in}}%
\pgfpathlineto{\pgfqpoint{1.256754in}{0.792144in}}%
\pgfpathlineto{\pgfqpoint{1.239743in}{0.782442in}}%
\pgfpathlineto{\pgfqpoint{1.222804in}{0.773194in}}%
\pgfpathlineto{\pgfqpoint{1.214730in}{0.772440in}}%
\pgfpathlineto{\pgfqpoint{1.206629in}{0.772249in}}%
\pgfpathclose%
\pgfusepath{fill}%
\end{pgfscope}%
\begin{pgfscope}%
\pgfpathrectangle{\pgfqpoint{0.407592in}{0.222904in}}{\pgfqpoint{1.411163in}{1.411163in}}%
\pgfusepath{clip}%
\pgfsetbuttcap%
\pgfsetroundjoin%
\definecolor{currentfill}{rgb}{0.556863,0.654902,0.847059}%
\pgfsetfillcolor{currentfill}%
\pgfsetfillopacity{0.130000}%
\pgfsetlinewidth{0.000000pt}%
\definecolor{currentstroke}{rgb}{0.000000,0.000000,0.000000}%
\pgfsetstrokecolor{currentstroke}%
\pgfsetdash{}{0pt}%
\pgfpathmoveto{\pgfqpoint{1.773791in}{1.189874in}}%
\pgfpathlineto{\pgfqpoint{1.791988in}{1.201836in}}%
\pgfpathlineto{\pgfqpoint{1.799403in}{1.212146in}}%
\pgfpathlineto{\pgfqpoint{1.806807in}{1.222657in}}%
\pgfpathlineto{\pgfqpoint{1.788636in}{1.210959in}}%
\pgfpathlineto{\pgfqpoint{1.781219in}{1.200315in}}%
\pgfusepath{fill}%
\end{pgfscope}%
\begin{pgfscope}%
\pgfpathrectangle{\pgfqpoint{0.407592in}{0.222904in}}{\pgfqpoint{1.411163in}{1.411163in}}%
\pgfusepath{clip}%
\pgfsetbuttcap%
\pgfsetroundjoin%
\definecolor{currentfill}{rgb}{0.556863,0.654902,0.847059}%
\pgfsetfillcolor{currentfill}%
\pgfsetfillopacity{0.130000}%
\pgfsetlinewidth{0.000000pt}%
\definecolor{currentstroke}{rgb}{0.000000,0.000000,0.000000}%
\pgfsetstrokecolor{currentstroke}%
\pgfsetdash{}{0pt}%
\pgfpathmoveto{\pgfqpoint{1.539936in}{0.981316in}}%
\pgfpathlineto{\pgfqpoint{1.557633in}{0.993528in}}%
\pgfpathlineto{\pgfqpoint{1.575410in}{1.005887in}}%
\pgfpathlineto{\pgfqpoint{1.583051in}{1.013922in}}%
\pgfpathlineto{\pgfqpoint{1.590679in}{1.022282in}}%
\pgfpathlineto{\pgfqpoint{1.572925in}{1.010141in}}%
\pgfpathlineto{\pgfqpoint{1.555250in}{0.998161in}}%
\pgfpathlineto{\pgfqpoint{1.547600in}{0.989566in}}%
\pgfpathlineto{\pgfqpoint{1.539936in}{0.981316in}}%
\pgfpathclose%
\pgfusepath{fill}%
\end{pgfscope}%
\begin{pgfscope}%
\pgfpathrectangle{\pgfqpoint{0.407592in}{0.222904in}}{\pgfqpoint{1.411163in}{1.411163in}}%
\pgfusepath{clip}%
\pgfsetbuttcap%
\pgfsetroundjoin%
\definecolor{currentfill}{rgb}{0.556863,0.654902,0.847059}%
\pgfsetfillcolor{currentfill}%
\pgfsetfillopacity{0.130000}%
\pgfsetlinewidth{0.000000pt}%
\definecolor{currentstroke}{rgb}{0.000000,0.000000,0.000000}%
\pgfsetstrokecolor{currentstroke}%
\pgfsetdash{}{0pt}%
\pgfpathmoveto{\pgfqpoint{1.106449in}{0.752370in}}%
\pgfpathlineto{\pgfqpoint{1.123225in}{0.757243in}}%
\pgfpathlineto{\pgfqpoint{1.140062in}{0.762947in}}%
\pgfpathlineto{\pgfqpoint{1.148320in}{0.760603in}}%
\pgfpathlineto{\pgfqpoint{1.156545in}{0.758707in}}%
\pgfpathlineto{\pgfqpoint{1.139746in}{0.752406in}}%
\pgfpathlineto{\pgfqpoint{1.123009in}{0.746937in}}%
\pgfpathlineto{\pgfqpoint{1.114747in}{0.749432in}}%
\pgfpathlineto{\pgfqpoint{1.106449in}{0.752370in}}%
\pgfpathclose%
\pgfusepath{fill}%
\end{pgfscope}%
\begin{pgfscope}%
\pgfpathrectangle{\pgfqpoint{0.407592in}{0.222904in}}{\pgfqpoint{1.411163in}{1.411163in}}%
\pgfusepath{clip}%
\pgfsetbuttcap%
\pgfsetroundjoin%
\definecolor{currentfill}{rgb}{0.556863,0.654902,0.847059}%
\pgfsetfillcolor{currentfill}%
\pgfsetfillopacity{0.130000}%
\pgfsetlinewidth{0.000000pt}%
\definecolor{currentstroke}{rgb}{0.000000,0.000000,0.000000}%
\pgfsetstrokecolor{currentstroke}%
\pgfsetdash{}{0pt}%
\pgfpathmoveto{\pgfqpoint{1.423392in}{0.890388in}}%
\pgfpathlineto{\pgfqpoint{1.440841in}{0.902429in}}%
\pgfpathlineto{\pgfqpoint{1.458368in}{0.914632in}}%
\pgfpathlineto{\pgfqpoint{1.466141in}{0.920739in}}%
\pgfpathlineto{\pgfqpoint{1.473898in}{0.927261in}}%
\pgfpathlineto{\pgfqpoint{1.456392in}{0.915160in}}%
\pgfpathlineto{\pgfqpoint{1.438965in}{0.903229in}}%
\pgfpathlineto{\pgfqpoint{1.431187in}{0.896583in}}%
\pgfpathlineto{\pgfqpoint{1.423392in}{0.890388in}}%
\pgfpathclose%
\pgfusepath{fill}%
\end{pgfscope}%
\begin{pgfscope}%
\pgfpathrectangle{\pgfqpoint{0.407592in}{0.222904in}}{\pgfqpoint{1.411163in}{1.411163in}}%
\pgfusepath{clip}%
\pgfsetbuttcap%
\pgfsetroundjoin%
\definecolor{currentfill}{rgb}{0.556863,0.654902,0.847059}%
\pgfsetfillcolor{currentfill}%
\pgfsetfillopacity{0.130000}%
\pgfsetlinewidth{0.000000pt}%
\definecolor{currentstroke}{rgb}{0.000000,0.000000,0.000000}%
\pgfsetstrokecolor{currentstroke}%
\pgfsetdash{}{0pt}%
\pgfpathmoveto{\pgfqpoint{0.838498in}{0.825573in}}%
\pgfpathlineto{\pgfqpoint{0.855024in}{0.817136in}}%
\pgfpathlineto{\pgfqpoint{0.871571in}{0.809619in}}%
\pgfpathlineto{\pgfqpoint{0.880425in}{0.803545in}}%
\pgfpathlineto{\pgfqpoint{0.889237in}{0.797644in}}%
\pgfpathlineto{\pgfqpoint{0.872742in}{0.805395in}}%
\pgfpathlineto{\pgfqpoint{0.856266in}{0.814128in}}%
\pgfpathlineto{\pgfqpoint{0.847403in}{0.819756in}}%
\pgfpathlineto{\pgfqpoint{0.838498in}{0.825573in}}%
\pgfpathclose%
\pgfusepath{fill}%
\end{pgfscope}%
\begin{pgfscope}%
\pgfpathrectangle{\pgfqpoint{0.407592in}{0.222904in}}{\pgfqpoint{1.411163in}{1.411163in}}%
\pgfusepath{clip}%
\pgfsetbuttcap%
\pgfsetroundjoin%
\definecolor{currentfill}{rgb}{0.556863,0.654902,0.847059}%
\pgfsetfillcolor{currentfill}%
\pgfsetfillopacity{0.130000}%
\pgfsetlinewidth{0.000000pt}%
\definecolor{currentstroke}{rgb}{0.000000,0.000000,0.000000}%
\pgfsetstrokecolor{currentstroke}%
\pgfsetdash{}{0pt}%
\pgfpathmoveto{\pgfqpoint{1.656720in}{1.082531in}}%
\pgfpathlineto{\pgfqpoint{1.674667in}{1.094641in}}%
\pgfpathlineto{\pgfqpoint{1.692694in}{1.106903in}}%
\pgfpathlineto{\pgfqpoint{1.700213in}{1.116222in}}%
\pgfpathlineto{\pgfqpoint{1.707721in}{1.125793in}}%
\pgfpathlineto{\pgfqpoint{1.689718in}{1.113785in}}%
\pgfpathlineto{\pgfqpoint{1.671795in}{1.101941in}}%
\pgfpathlineto{\pgfqpoint{1.664264in}{1.092105in}}%
\pgfpathlineto{\pgfqpoint{1.656720in}{1.082531in}}%
\pgfpathclose%
\pgfusepath{fill}%
\end{pgfscope}%
\begin{pgfscope}%
\pgfpathrectangle{\pgfqpoint{0.407592in}{0.222904in}}{\pgfqpoint{1.411163in}{1.411163in}}%
\pgfusepath{clip}%
\pgfsetbuttcap%
\pgfsetroundjoin%
\definecolor{currentfill}{rgb}{0.556863,0.654902,0.847059}%
\pgfsetfillcolor{currentfill}%
\pgfsetfillopacity{0.130000}%
\pgfsetlinewidth{0.000000pt}%
\definecolor{currentstroke}{rgb}{0.000000,0.000000,0.000000}%
\pgfsetstrokecolor{currentstroke}%
\pgfsetdash{}{0pt}%
\pgfpathmoveto{\pgfqpoint{1.256754in}{0.792144in}}%
\pgfpathlineto{\pgfqpoint{1.273838in}{0.802245in}}%
\pgfpathlineto{\pgfqpoint{1.290995in}{0.812698in}}%
\pgfpathlineto{\pgfqpoint{1.298992in}{0.814743in}}%
\pgfpathlineto{\pgfqpoint{1.306964in}{0.817310in}}%
\pgfpathlineto{\pgfqpoint{1.289831in}{0.806573in}}%
\pgfpathlineto{\pgfqpoint{1.272772in}{0.796163in}}%
\pgfpathlineto{\pgfqpoint{1.264775in}{0.793869in}}%
\pgfpathlineto{\pgfqpoint{1.256754in}{0.792144in}}%
\pgfpathclose%
\pgfusepath{fill}%
\end{pgfscope}%
\begin{pgfscope}%
\pgfpathrectangle{\pgfqpoint{0.407592in}{0.222904in}}{\pgfqpoint{1.411163in}{1.411163in}}%
\pgfusepath{clip}%
\pgfsetbuttcap%
\pgfsetroundjoin%
\definecolor{currentfill}{rgb}{0.556863,0.654902,0.847059}%
\pgfsetfillcolor{currentfill}%
\pgfsetfillopacity{0.130000}%
\pgfsetlinewidth{0.000000pt}%
\definecolor{currentstroke}{rgb}{0.000000,0.000000,0.000000}%
\pgfsetstrokecolor{currentstroke}%
\pgfsetdash{}{0pt}%
\pgfpathmoveto{\pgfqpoint{0.922297in}{0.785227in}}%
\pgfpathlineto{\pgfqpoint{0.938867in}{0.780611in}}%
\pgfpathlineto{\pgfqpoint{0.955468in}{0.777068in}}%
\pgfpathlineto{\pgfqpoint{0.964136in}{0.771233in}}%
\pgfpathlineto{\pgfqpoint{0.972764in}{0.765598in}}%
\pgfpathlineto{\pgfqpoint{0.956215in}{0.769072in}}%
\pgfpathlineto{\pgfqpoint{0.939697in}{0.773708in}}%
\pgfpathlineto{\pgfqpoint{0.931017in}{0.779369in}}%
\pgfpathlineto{\pgfqpoint{0.922297in}{0.785227in}}%
\pgfpathclose%
\pgfusepath{fill}%
\end{pgfscope}%
\begin{pgfscope}%
\pgfpathrectangle{\pgfqpoint{0.407592in}{0.222904in}}{\pgfqpoint{1.411163in}{1.411163in}}%
\pgfusepath{clip}%
\pgfsetbuttcap%
\pgfsetroundjoin%
\definecolor{currentfill}{rgb}{0.556863,0.654902,0.847059}%
\pgfsetfillcolor{currentfill}%
\pgfsetfillopacity{0.130000}%
\pgfsetlinewidth{0.000000pt}%
\definecolor{currentstroke}{rgb}{0.000000,0.000000,0.000000}%
\pgfsetstrokecolor{currentstroke}%
\pgfsetdash{}{0pt}%
\pgfpathmoveto{\pgfqpoint{0.637048in}{0.967665in}}%
\pgfpathlineto{\pgfqpoint{0.653533in}{0.953388in}}%
\pgfpathlineto{\pgfqpoint{0.670028in}{0.939582in}}%
\pgfpathlineto{\pgfqpoint{0.679297in}{0.934598in}}%
\pgfpathlineto{\pgfqpoint{0.688521in}{0.929816in}}%
\pgfpathlineto{\pgfqpoint{0.672071in}{0.944046in}}%
\pgfpathlineto{\pgfqpoint{0.655632in}{0.958746in}}%
\pgfpathlineto{\pgfqpoint{0.646362in}{0.963099in}}%
\pgfpathlineto{\pgfqpoint{0.637048in}{0.967665in}}%
\pgfpathclose%
\pgfusepath{fill}%
\end{pgfscope}%
\begin{pgfscope}%
\pgfpathrectangle{\pgfqpoint{0.407592in}{0.222904in}}{\pgfqpoint{1.411163in}{1.411163in}}%
\pgfusepath{clip}%
\pgfsetbuttcap%
\pgfsetroundjoin%
\definecolor{currentfill}{rgb}{0.556863,0.654902,0.847059}%
\pgfsetfillcolor{currentfill}%
\pgfsetfillopacity{0.130000}%
\pgfsetlinewidth{0.000000pt}%
\definecolor{currentstroke}{rgb}{0.000000,0.000000,0.000000}%
\pgfsetstrokecolor{currentstroke}%
\pgfsetdash{}{0pt}%
\pgfpathmoveto{\pgfqpoint{1.056278in}{0.753592in}}%
\pgfpathlineto{\pgfqpoint{1.072987in}{0.756019in}}%
\pgfpathlineto{\pgfqpoint{1.089748in}{0.759419in}}%
\pgfpathlineto{\pgfqpoint{1.098117in}{0.755711in}}%
\pgfpathlineto{\pgfqpoint{1.106449in}{0.752370in}}%
\pgfpathlineto{\pgfqpoint{1.089731in}{0.748417in}}%
\pgfpathlineto{\pgfqpoint{1.073067in}{0.745476in}}%
\pgfpathlineto{\pgfqpoint{1.064691in}{0.749360in}}%
\pgfpathlineto{\pgfqpoint{1.056278in}{0.753592in}}%
\pgfpathclose%
\pgfusepath{fill}%
\end{pgfscope}%
\begin{pgfscope}%
\pgfpathrectangle{\pgfqpoint{0.407592in}{0.222904in}}{\pgfqpoint{1.411163in}{1.411163in}}%
\pgfusepath{clip}%
\pgfsetbuttcap%
\pgfsetroundjoin%
\definecolor{currentfill}{rgb}{0.556863,0.654902,0.847059}%
\pgfsetfillcolor{currentfill}%
\pgfsetfillopacity{0.130000}%
\pgfsetlinewidth{0.000000pt}%
\definecolor{currentstroke}{rgb}{0.000000,0.000000,0.000000}%
\pgfsetstrokecolor{currentstroke}%
\pgfsetdash{}{0pt}%
\pgfpathmoveto{\pgfqpoint{0.754432in}{0.878371in}}%
\pgfpathlineto{\pgfqpoint{0.770941in}{0.867089in}}%
\pgfpathlineto{\pgfqpoint{0.787465in}{0.856524in}}%
\pgfpathlineto{\pgfqpoint{0.796504in}{0.850698in}}%
\pgfpathlineto{\pgfqpoint{0.805501in}{0.845049in}}%
\pgfpathlineto{\pgfqpoint{0.789025in}{0.855999in}}%
\pgfpathlineto{\pgfqpoint{0.772564in}{0.867691in}}%
\pgfpathlineto{\pgfqpoint{0.763519in}{0.872933in}}%
\pgfpathlineto{\pgfqpoint{0.754432in}{0.878371in}}%
\pgfpathclose%
\pgfusepath{fill}%
\end{pgfscope}%
\begin{pgfscope}%
\pgfpathrectangle{\pgfqpoint{0.407592in}{0.222904in}}{\pgfqpoint{1.411163in}{1.411163in}}%
\pgfusepath{clip}%
\pgfsetbuttcap%
\pgfsetroundjoin%
\definecolor{currentfill}{rgb}{0.556863,0.654902,0.847059}%
\pgfsetfillcolor{currentfill}%
\pgfsetfillopacity{0.130000}%
\pgfsetlinewidth{0.000000pt}%
\definecolor{currentstroke}{rgb}{0.000000,0.000000,0.000000}%
\pgfsetstrokecolor{currentstroke}%
\pgfsetdash{}{0pt}%
\pgfpathmoveto{\pgfqpoint{0.519315in}{1.068919in}}%
\pgfpathlineto{\pgfqpoint{0.535786in}{1.052811in}}%
\pgfpathlineto{\pgfqpoint{0.552267in}{1.037011in}}%
\pgfpathlineto{\pgfqpoint{0.561760in}{1.032825in}}%
\pgfpathlineto{\pgfqpoint{0.571208in}{1.028845in}}%
\pgfpathlineto{\pgfqpoint{0.554771in}{1.045015in}}%
\pgfpathlineto{\pgfqpoint{0.538344in}{1.061482in}}%
\pgfpathlineto{\pgfqpoint{0.528852in}{1.065095in}}%
\pgfpathlineto{\pgfqpoint{0.519315in}{1.068919in}}%
\pgfpathclose%
\pgfusepath{fill}%
\end{pgfscope}%
\begin{pgfscope}%
\pgfpathrectangle{\pgfqpoint{0.407592in}{0.222904in}}{\pgfqpoint{1.411163in}{1.411163in}}%
\pgfusepath{clip}%
\pgfsetbuttcap%
\pgfsetroundjoin%
\definecolor{currentfill}{rgb}{0.556863,0.654902,0.847059}%
\pgfsetfillcolor{currentfill}%
\pgfsetfillopacity{0.130000}%
\pgfsetlinewidth{0.000000pt}%
\definecolor{currentstroke}{rgb}{0.000000,0.000000,0.000000}%
\pgfsetstrokecolor{currentstroke}%
\pgfsetdash{}{0pt}%
\pgfpathmoveto{\pgfqpoint{1.306964in}{0.817310in}}%
\pgfpathlineto{\pgfqpoint{1.324173in}{0.828336in}}%
\pgfpathlineto{\pgfqpoint{1.341457in}{0.839620in}}%
\pgfpathlineto{\pgfqpoint{1.349384in}{0.842906in}}%
\pgfpathlineto{\pgfqpoint{1.357291in}{0.846689in}}%
\pgfpathlineto{\pgfqpoint{1.340029in}{0.835252in}}%
\pgfpathlineto{\pgfqpoint{1.322844in}{0.824057in}}%
\pgfpathlineto{\pgfqpoint{1.314915in}{0.820412in}}%
\pgfpathlineto{\pgfqpoint{1.306964in}{0.817310in}}%
\pgfpathclose%
\pgfusepath{fill}%
\end{pgfscope}%
\begin{pgfscope}%
\pgfpathrectangle{\pgfqpoint{0.407592in}{0.222904in}}{\pgfqpoint{1.411163in}{1.411163in}}%
\pgfusepath{clip}%
\pgfsetbuttcap%
\pgfsetroundjoin%
\definecolor{currentfill}{rgb}{0.556863,0.654902,0.847059}%
\pgfsetfillcolor{currentfill}%
\pgfsetfillopacity{0.130000}%
\pgfsetlinewidth{0.000000pt}%
\definecolor{currentstroke}{rgb}{0.000000,0.000000,0.000000}%
\pgfsetstrokecolor{currentstroke}%
\pgfsetdash{}{0pt}%
\pgfpathmoveto{\pgfqpoint{1.473898in}{0.927261in}}%
\pgfpathlineto{\pgfqpoint{1.491481in}{0.939518in}}%
\pgfpathlineto{\pgfqpoint{1.509144in}{0.951921in}}%
\pgfpathlineto{\pgfqpoint{1.516863in}{0.958714in}}%
\pgfpathlineto{\pgfqpoint{1.524568in}{0.965882in}}%
\pgfpathlineto{\pgfqpoint{1.506928in}{0.953627in}}%
\pgfpathlineto{\pgfqpoint{1.489366in}{0.941528in}}%
\pgfpathlineto{\pgfqpoint{1.481639in}{0.934193in}}%
\pgfpathlineto{\pgfqpoint{1.473898in}{0.927261in}}%
\pgfpathclose%
\pgfusepath{fill}%
\end{pgfscope}%
\begin{pgfscope}%
\pgfpathrectangle{\pgfqpoint{0.407592in}{0.222904in}}{\pgfqpoint{1.411163in}{1.411163in}}%
\pgfusepath{clip}%
\pgfsetbuttcap%
\pgfsetroundjoin%
\definecolor{currentfill}{rgb}{0.556863,0.654902,0.847059}%
\pgfsetfillcolor{currentfill}%
\pgfsetfillopacity{0.130000}%
\pgfsetlinewidth{0.000000pt}%
\definecolor{currentstroke}{rgb}{0.000000,0.000000,0.000000}%
\pgfsetstrokecolor{currentstroke}%
\pgfsetdash{}{0pt}%
\pgfpathmoveto{\pgfqpoint{1.590679in}{1.022282in}}%
\pgfpathlineto{\pgfqpoint{1.608512in}{1.034574in}}%
\pgfpathlineto{\pgfqpoint{1.626426in}{1.047008in}}%
\pgfpathlineto{\pgfqpoint{1.634018in}{1.055458in}}%
\pgfpathlineto{\pgfqpoint{1.641597in}{1.064199in}}%
\pgfpathlineto{\pgfqpoint{1.623707in}{1.051993in}}%
\pgfpathlineto{\pgfqpoint{1.605897in}{1.039942in}}%
\pgfpathlineto{\pgfqpoint{1.598294in}{1.030958in}}%
\pgfpathlineto{\pgfqpoint{1.590679in}{1.022282in}}%
\pgfpathclose%
\pgfusepath{fill}%
\end{pgfscope}%
\begin{pgfscope}%
\pgfpathrectangle{\pgfqpoint{0.407592in}{0.222904in}}{\pgfqpoint{1.411163in}{1.411163in}}%
\pgfusepath{clip}%
\pgfsetbuttcap%
\pgfsetroundjoin%
\definecolor{currentfill}{rgb}{0.556863,0.654902,0.847059}%
\pgfsetfillcolor{currentfill}%
\pgfsetfillopacity{0.130000}%
\pgfsetlinewidth{0.000000pt}%
\definecolor{currentstroke}{rgb}{0.000000,0.000000,0.000000}%
\pgfsetstrokecolor{currentstroke}%
\pgfsetdash{}{0pt}%
\pgfpathmoveto{\pgfqpoint{1.005969in}{0.762093in}}%
\pgfpathlineto{\pgfqpoint{1.022631in}{0.762003in}}%
\pgfpathlineto{\pgfqpoint{1.039339in}{0.762963in}}%
\pgfpathlineto{\pgfqpoint{1.047827in}{0.758137in}}%
\pgfpathlineto{\pgfqpoint{1.056278in}{0.753592in}}%
\pgfpathlineto{\pgfqpoint{1.039619in}{0.752215in}}%
\pgfpathlineto{\pgfqpoint{1.023005in}{0.751955in}}%
\pgfpathlineto{\pgfqpoint{1.014506in}{0.756894in}}%
\pgfpathlineto{\pgfqpoint{1.005969in}{0.762093in}}%
\pgfpathclose%
\pgfusepath{fill}%
\end{pgfscope}%
\begin{pgfscope}%
\pgfpathrectangle{\pgfqpoint{0.407592in}{0.222904in}}{\pgfqpoint{1.411163in}{1.411163in}}%
\pgfusepath{clip}%
\pgfsetbuttcap%
\pgfsetroundjoin%
\definecolor{currentfill}{rgb}{0.556863,0.654902,0.847059}%
\pgfsetfillcolor{currentfill}%
\pgfsetfillopacity{0.130000}%
\pgfsetlinewidth{0.000000pt}%
\definecolor{currentstroke}{rgb}{0.000000,0.000000,0.000000}%
\pgfsetstrokecolor{currentstroke}%
\pgfsetdash{}{0pt}%
\pgfpathmoveto{\pgfqpoint{0.670028in}{0.939582in}}%
\pgfpathlineto{\pgfqpoint{0.686534in}{0.926282in}}%
\pgfpathlineto{\pgfqpoint{0.703052in}{0.913523in}}%
\pgfpathlineto{\pgfqpoint{0.712274in}{0.908126in}}%
\pgfpathlineto{\pgfqpoint{0.721453in}{0.902916in}}%
\pgfpathlineto{\pgfqpoint{0.704981in}{0.916092in}}%
\pgfpathlineto{\pgfqpoint{0.688521in}{0.929816in}}%
\pgfpathlineto{\pgfqpoint{0.679297in}{0.934598in}}%
\pgfpathlineto{\pgfqpoint{0.670028in}{0.939582in}}%
\pgfpathclose%
\pgfusepath{fill}%
\end{pgfscope}%
\begin{pgfscope}%
\pgfpathrectangle{\pgfqpoint{0.407592in}{0.222904in}}{\pgfqpoint{1.411163in}{1.411163in}}%
\pgfusepath{clip}%
\pgfsetbuttcap%
\pgfsetroundjoin%
\definecolor{currentfill}{rgb}{0.556863,0.654902,0.847059}%
\pgfsetfillcolor{currentfill}%
\pgfsetfillopacity{0.130000}%
\pgfsetlinewidth{0.000000pt}%
\definecolor{currentstroke}{rgb}{0.000000,0.000000,0.000000}%
\pgfsetstrokecolor{currentstroke}%
\pgfsetdash{}{0pt}%
\pgfpathmoveto{\pgfqpoint{1.707721in}{1.125793in}}%
\pgfpathlineto{\pgfqpoint{1.725805in}{1.137956in}}%
\pgfpathlineto{\pgfqpoint{1.743970in}{1.150268in}}%
\pgfpathlineto{\pgfqpoint{1.751442in}{1.159832in}}%
\pgfpathlineto{\pgfqpoint{1.758903in}{1.169625in}}%
\pgfpathlineto{\pgfqpoint{1.740763in}{1.157563in}}%
\pgfpathlineto{\pgfqpoint{1.722703in}{1.145660in}}%
\pgfpathlineto{\pgfqpoint{1.715218in}{1.135608in}}%
\pgfpathlineto{\pgfqpoint{1.707721in}{1.125793in}}%
\pgfpathclose%
\pgfusepath{fill}%
\end{pgfscope}%
\begin{pgfscope}%
\pgfpathrectangle{\pgfqpoint{0.407592in}{0.222904in}}{\pgfqpoint{1.411163in}{1.411163in}}%
\pgfusepath{clip}%
\pgfsetbuttcap%
\pgfsetroundjoin%
\definecolor{currentfill}{rgb}{0.556863,0.654902,0.847059}%
\pgfsetfillcolor{currentfill}%
\pgfsetfillopacity{0.130000}%
\pgfsetlinewidth{0.000000pt}%
\definecolor{currentstroke}{rgb}{0.000000,0.000000,0.000000}%
\pgfsetstrokecolor{currentstroke}%
\pgfsetdash{}{0pt}%
\pgfpathmoveto{\pgfqpoint{0.552267in}{1.037011in}}%
\pgfpathlineto{\pgfqpoint{0.568757in}{1.021544in}}%
\pgfpathlineto{\pgfqpoint{0.585256in}{1.006433in}}%
\pgfpathlineto{\pgfqpoint{0.594705in}{1.001867in}}%
\pgfpathlineto{\pgfqpoint{0.604109in}{0.997501in}}%
\pgfpathlineto{\pgfqpoint{0.587654in}{1.012999in}}%
\pgfpathlineto{\pgfqpoint{0.571208in}{1.028845in}}%
\pgfpathlineto{\pgfqpoint{0.561760in}{1.032825in}}%
\pgfpathlineto{\pgfqpoint{0.552267in}{1.037011in}}%
\pgfpathclose%
\pgfusepath{fill}%
\end{pgfscope}%
\begin{pgfscope}%
\pgfpathrectangle{\pgfqpoint{0.407592in}{0.222904in}}{\pgfqpoint{1.411163in}{1.411163in}}%
\pgfusepath{clip}%
\pgfsetbuttcap%
\pgfsetroundjoin%
\definecolor{currentfill}{rgb}{0.556863,0.654902,0.847059}%
\pgfsetfillcolor{currentfill}%
\pgfsetfillopacity{0.130000}%
\pgfsetlinewidth{0.000000pt}%
\definecolor{currentstroke}{rgb}{0.000000,0.000000,0.000000}%
\pgfsetstrokecolor{currentstroke}%
\pgfsetdash{}{0pt}%
\pgfpathmoveto{\pgfqpoint{1.357291in}{0.846689in}}%
\pgfpathlineto{\pgfqpoint{1.374628in}{0.858345in}}%
\pgfpathlineto{\pgfqpoint{1.392043in}{0.870198in}}%
\pgfpathlineto{\pgfqpoint{1.399908in}{0.874554in}}%
\pgfpathlineto{\pgfqpoint{1.407754in}{0.879371in}}%
\pgfpathlineto{\pgfqpoint{1.390361in}{0.867471in}}%
\pgfpathlineto{\pgfqpoint{1.373045in}{0.855762in}}%
\pgfpathlineto{\pgfqpoint{1.365177in}{0.850974in}}%
\pgfpathlineto{\pgfqpoint{1.357291in}{0.846689in}}%
\pgfpathclose%
\pgfusepath{fill}%
\end{pgfscope}%
\begin{pgfscope}%
\pgfpathrectangle{\pgfqpoint{0.407592in}{0.222904in}}{\pgfqpoint{1.411163in}{1.411163in}}%
\pgfusepath{clip}%
\pgfsetbuttcap%
\pgfsetroundjoin%
\definecolor{currentfill}{rgb}{0.556863,0.654902,0.847059}%
\pgfsetfillcolor{currentfill}%
\pgfsetfillopacity{0.130000}%
\pgfsetlinewidth{0.000000pt}%
\definecolor{currentstroke}{rgb}{0.000000,0.000000,0.000000}%
\pgfsetstrokecolor{currentstroke}%
\pgfsetdash{}{0pt}%
\pgfpathmoveto{\pgfqpoint{0.871571in}{0.809619in}}%
\pgfpathlineto{\pgfqpoint{0.888140in}{0.803053in}}%
\pgfpathlineto{\pgfqpoint{0.904734in}{0.797463in}}%
\pgfpathlineto{\pgfqpoint{0.913536in}{0.791264in}}%
\pgfpathlineto{\pgfqpoint{0.922297in}{0.785227in}}%
\pgfpathlineto{\pgfqpoint{0.905755in}{0.790911in}}%
\pgfpathlineto{\pgfqpoint{0.889237in}{0.797644in}}%
\pgfpathlineto{\pgfqpoint{0.880425in}{0.803545in}}%
\pgfpathlineto{\pgfqpoint{0.871571in}{0.809619in}}%
\pgfpathclose%
\pgfusepath{fill}%
\end{pgfscope}%
\begin{pgfscope}%
\pgfpathrectangle{\pgfqpoint{0.407592in}{0.222904in}}{\pgfqpoint{1.411163in}{1.411163in}}%
\pgfusepath{clip}%
\pgfsetbuttcap%
\pgfsetroundjoin%
\definecolor{currentfill}{rgb}{0.556863,0.654902,0.847059}%
\pgfsetfillcolor{currentfill}%
\pgfsetfillopacity{0.130000}%
\pgfsetlinewidth{0.000000pt}%
\definecolor{currentstroke}{rgb}{0.000000,0.000000,0.000000}%
\pgfsetstrokecolor{currentstroke}%
\pgfsetdash{}{0pt}%
\pgfpathmoveto{\pgfqpoint{1.758903in}{1.169625in}}%
\pgfpathlineto{\pgfqpoint{1.777125in}{1.181840in}}%
\pgfpathlineto{\pgfqpoint{1.784562in}{1.191731in}}%
\pgfpathlineto{\pgfqpoint{1.791988in}{1.201836in}}%
\pgfpathlineto{\pgfqpoint{1.773791in}{1.189874in}}%
\pgfpathlineto{\pgfqpoint{1.766353in}{1.179642in}}%
\pgfusepath{fill}%
\end{pgfscope}%
\begin{pgfscope}%
\pgfpathrectangle{\pgfqpoint{0.407592in}{0.222904in}}{\pgfqpoint{1.411163in}{1.411163in}}%
\pgfusepath{clip}%
\pgfsetbuttcap%
\pgfsetroundjoin%
\definecolor{currentfill}{rgb}{0.556863,0.654902,0.847059}%
\pgfsetfillcolor{currentfill}%
\pgfsetfillopacity{0.130000}%
\pgfsetlinewidth{0.000000pt}%
\definecolor{currentstroke}{rgb}{0.000000,0.000000,0.000000}%
\pgfsetstrokecolor{currentstroke}%
\pgfsetdash{}{0pt}%
\pgfpathmoveto{\pgfqpoint{0.787465in}{0.856524in}}%
\pgfpathlineto{\pgfqpoint{0.804005in}{0.846717in}}%
\pgfpathlineto{\pgfqpoint{0.820562in}{0.837708in}}%
\pgfpathlineto{\pgfqpoint{0.829551in}{0.831561in}}%
\pgfpathlineto{\pgfqpoint{0.838498in}{0.825573in}}%
\pgfpathlineto{\pgfqpoint{0.821991in}{0.834892in}}%
\pgfpathlineto{\pgfqpoint{0.805501in}{0.845049in}}%
\pgfpathlineto{\pgfqpoint{0.796504in}{0.850698in}}%
\pgfpathlineto{\pgfqpoint{0.787465in}{0.856524in}}%
\pgfpathclose%
\pgfusepath{fill}%
\end{pgfscope}%
\begin{pgfscope}%
\pgfpathrectangle{\pgfqpoint{0.407592in}{0.222904in}}{\pgfqpoint{1.411163in}{1.411163in}}%
\pgfusepath{clip}%
\pgfsetbuttcap%
\pgfsetroundjoin%
\definecolor{currentfill}{rgb}{0.556863,0.654902,0.847059}%
\pgfsetfillcolor{currentfill}%
\pgfsetfillopacity{0.130000}%
\pgfsetlinewidth{0.000000pt}%
\definecolor{currentstroke}{rgb}{0.000000,0.000000,0.000000}%
\pgfsetstrokecolor{currentstroke}%
\pgfsetdash{}{0pt}%
\pgfpathmoveto{\pgfqpoint{1.524568in}{0.965882in}}%
\pgfpathlineto{\pgfqpoint{1.542288in}{0.978282in}}%
\pgfpathlineto{\pgfqpoint{1.560087in}{0.990819in}}%
\pgfpathlineto{\pgfqpoint{1.567755in}{0.998184in}}%
\pgfpathlineto{\pgfqpoint{1.575410in}{1.005887in}}%
\pgfpathlineto{\pgfqpoint{1.557633in}{0.993528in}}%
\pgfpathlineto{\pgfqpoint{1.539936in}{0.981316in}}%
\pgfpathlineto{\pgfqpoint{1.532259in}{0.973418in}}%
\pgfpathlineto{\pgfqpoint{1.524568in}{0.965882in}}%
\pgfpathclose%
\pgfusepath{fill}%
\end{pgfscope}%
\begin{pgfscope}%
\pgfpathrectangle{\pgfqpoint{0.407592in}{0.222904in}}{\pgfqpoint{1.411163in}{1.411163in}}%
\pgfusepath{clip}%
\pgfsetbuttcap%
\pgfsetroundjoin%
\definecolor{currentfill}{rgb}{0.556863,0.654902,0.847059}%
\pgfsetfillcolor{currentfill}%
\pgfsetfillopacity{0.130000}%
\pgfsetlinewidth{0.000000pt}%
\definecolor{currentstroke}{rgb}{0.000000,0.000000,0.000000}%
\pgfsetstrokecolor{currentstroke}%
\pgfsetdash{}{0pt}%
\pgfpathmoveto{\pgfqpoint{1.641597in}{1.064199in}}%
\pgfpathlineto{\pgfqpoint{1.659568in}{1.076553in}}%
\pgfpathlineto{\pgfqpoint{1.677619in}{1.089046in}}%
\pgfpathlineto{\pgfqpoint{1.685162in}{1.097842in}}%
\pgfpathlineto{\pgfqpoint{1.692694in}{1.106903in}}%
\pgfpathlineto{\pgfqpoint{1.674667in}{1.094641in}}%
\pgfpathlineto{\pgfqpoint{1.656720in}{1.082531in}}%
\pgfpathlineto{\pgfqpoint{1.649165in}{1.073227in}}%
\pgfpathlineto{\pgfqpoint{1.641597in}{1.064199in}}%
\pgfpathclose%
\pgfusepath{fill}%
\end{pgfscope}%
\begin{pgfscope}%
\pgfpathrectangle{\pgfqpoint{0.407592in}{0.222904in}}{\pgfqpoint{1.411163in}{1.411163in}}%
\pgfusepath{clip}%
\pgfsetbuttcap%
\pgfsetroundjoin%
\definecolor{currentfill}{rgb}{0.556863,0.654902,0.847059}%
\pgfsetfillcolor{currentfill}%
\pgfsetfillopacity{0.130000}%
\pgfsetlinewidth{0.000000pt}%
\definecolor{currentstroke}{rgb}{0.000000,0.000000,0.000000}%
\pgfsetstrokecolor{currentstroke}%
\pgfsetdash{}{0pt}%
\pgfpathmoveto{\pgfqpoint{0.955468in}{0.777068in}}%
\pgfpathlineto{\pgfqpoint{0.972103in}{0.774594in}}%
\pgfpathlineto{\pgfqpoint{0.988776in}{0.773171in}}%
\pgfpathlineto{\pgfqpoint{0.997392in}{0.767526in}}%
\pgfpathlineto{\pgfqpoint{1.005969in}{0.762093in}}%
\pgfpathlineto{\pgfqpoint{0.989348in}{0.763279in}}%
\pgfpathlineto{\pgfqpoint{0.972764in}{0.765598in}}%
\pgfpathlineto{\pgfqpoint{0.964136in}{0.771233in}}%
\pgfpathlineto{\pgfqpoint{0.955468in}{0.777068in}}%
\pgfpathclose%
\pgfusepath{fill}%
\end{pgfscope}%
\begin{pgfscope}%
\pgfpathrectangle{\pgfqpoint{0.407592in}{0.222904in}}{\pgfqpoint{1.411163in}{1.411163in}}%
\pgfusepath{clip}%
\pgfsetbuttcap%
\pgfsetroundjoin%
\definecolor{currentfill}{rgb}{0.556863,0.654902,0.847059}%
\pgfsetfillcolor{currentfill}%
\pgfsetfillopacity{0.130000}%
\pgfsetlinewidth{0.000000pt}%
\definecolor{currentstroke}{rgb}{0.000000,0.000000,0.000000}%
\pgfsetstrokecolor{currentstroke}%
\pgfsetdash{}{0pt}%
\pgfpathmoveto{\pgfqpoint{1.407754in}{0.879371in}}%
\pgfpathlineto{\pgfqpoint{1.425224in}{0.891446in}}%
\pgfpathlineto{\pgfqpoint{1.442773in}{0.903684in}}%
\pgfpathlineto{\pgfqpoint{1.450579in}{0.908946in}}%
\pgfpathlineto{\pgfqpoint{1.458368in}{0.914632in}}%
\pgfpathlineto{\pgfqpoint{1.440841in}{0.902429in}}%
\pgfpathlineto{\pgfqpoint{1.423392in}{0.890388in}}%
\pgfpathlineto{\pgfqpoint{1.415582in}{0.884650in}}%
\pgfpathlineto{\pgfqpoint{1.407754in}{0.879371in}}%
\pgfpathclose%
\pgfusepath{fill}%
\end{pgfscope}%
\begin{pgfscope}%
\pgfpathrectangle{\pgfqpoint{0.407592in}{0.222904in}}{\pgfqpoint{1.411163in}{1.411163in}}%
\pgfusepath{clip}%
\pgfsetbuttcap%
\pgfsetroundjoin%
\definecolor{currentfill}{rgb}{0.556863,0.654902,0.847059}%
\pgfsetfillcolor{currentfill}%
\pgfsetfillopacity{0.130000}%
\pgfsetlinewidth{0.000000pt}%
\definecolor{currentstroke}{rgb}{0.000000,0.000000,0.000000}%
\pgfsetstrokecolor{currentstroke}%
\pgfsetdash{}{0pt}%
\pgfpathmoveto{\pgfqpoint{0.585256in}{1.006433in}}%
\pgfpathlineto{\pgfqpoint{0.601765in}{0.991706in}}%
\pgfpathlineto{\pgfqpoint{0.618285in}{0.977389in}}%
\pgfpathlineto{\pgfqpoint{0.627689in}{0.972431in}}%
\pgfpathlineto{\pgfqpoint{0.637048in}{0.967665in}}%
\pgfpathlineto{\pgfqpoint{0.620574in}{0.982379in}}%
\pgfpathlineto{\pgfqpoint{0.604109in}{0.997501in}}%
\pgfpathlineto{\pgfqpoint{0.594705in}{1.001867in}}%
\pgfpathlineto{\pgfqpoint{0.585256in}{1.006433in}}%
\pgfpathclose%
\pgfusepath{fill}%
\end{pgfscope}%
\begin{pgfscope}%
\pgfpathrectangle{\pgfqpoint{0.407592in}{0.222904in}}{\pgfqpoint{1.411163in}{1.411163in}}%
\pgfusepath{clip}%
\pgfsetbuttcap%
\pgfsetroundjoin%
\definecolor{currentfill}{rgb}{0.556863,0.654902,0.847059}%
\pgfsetfillcolor{currentfill}%
\pgfsetfillopacity{0.130000}%
\pgfsetlinewidth{0.000000pt}%
\definecolor{currentstroke}{rgb}{0.000000,0.000000,0.000000}%
\pgfsetstrokecolor{currentstroke}%
\pgfsetdash{}{0pt}%
\pgfpathmoveto{\pgfqpoint{0.703052in}{0.913523in}}%
\pgfpathlineto{\pgfqpoint{0.719582in}{0.901344in}}%
\pgfpathlineto{\pgfqpoint{0.736126in}{0.889781in}}%
\pgfpathlineto{\pgfqpoint{0.745300in}{0.883991in}}%
\pgfpathlineto{\pgfqpoint{0.754432in}{0.878371in}}%
\pgfpathlineto{\pgfqpoint{0.737936in}{0.890328in}}%
\pgfpathlineto{\pgfqpoint{0.721453in}{0.902916in}}%
\pgfpathlineto{\pgfqpoint{0.712274in}{0.908126in}}%
\pgfpathlineto{\pgfqpoint{0.703052in}{0.913523in}}%
\pgfpathclose%
\pgfusepath{fill}%
\end{pgfscope}%
\begin{pgfscope}%
\pgfpathrectangle{\pgfqpoint{0.407592in}{0.222904in}}{\pgfqpoint{1.411163in}{1.411163in}}%
\pgfusepath{clip}%
\pgfsetbuttcap%
\pgfsetroundjoin%
\definecolor{currentfill}{rgb}{0.556863,0.654902,0.847059}%
\pgfsetfillcolor{currentfill}%
\pgfsetfillopacity{0.130000}%
\pgfsetlinewidth{0.000000pt}%
\definecolor{currentstroke}{rgb}{0.000000,0.000000,0.000000}%
\pgfsetstrokecolor{currentstroke}%
\pgfsetdash{}{0pt}%
\pgfpathmoveto{\pgfqpoint{1.190339in}{0.773441in}}%
\pgfpathlineto{\pgfqpoint{1.207337in}{0.781717in}}%
\pgfpathlineto{\pgfqpoint{1.224405in}{0.790509in}}%
\pgfpathlineto{\pgfqpoint{1.232534in}{0.790173in}}%
\pgfpathlineto{\pgfqpoint{1.240634in}{0.790318in}}%
\pgfpathlineto{\pgfqpoint{1.223596in}{0.781041in}}%
\pgfpathlineto{\pgfqpoint{1.206629in}{0.772249in}}%
\pgfpathlineto{\pgfqpoint{1.198499in}{0.772594in}}%
\pgfpathlineto{\pgfqpoint{1.190339in}{0.773441in}}%
\pgfpathclose%
\pgfusepath{fill}%
\end{pgfscope}%
\begin{pgfscope}%
\pgfpathrectangle{\pgfqpoint{0.407592in}{0.222904in}}{\pgfqpoint{1.411163in}{1.411163in}}%
\pgfusepath{clip}%
\pgfsetbuttcap%
\pgfsetroundjoin%
\definecolor{currentfill}{rgb}{0.556863,0.654902,0.847059}%
\pgfsetfillcolor{currentfill}%
\pgfsetfillopacity{0.130000}%
\pgfsetlinewidth{0.000000pt}%
\definecolor{currentstroke}{rgb}{0.000000,0.000000,0.000000}%
\pgfsetstrokecolor{currentstroke}%
\pgfsetdash{}{0pt}%
\pgfpathmoveto{\pgfqpoint{1.140062in}{0.762947in}}%
\pgfpathlineto{\pgfqpoint{1.156961in}{0.769400in}}%
\pgfpathlineto{\pgfqpoint{1.173926in}{0.776526in}}%
\pgfpathlineto{\pgfqpoint{1.182148in}{0.774762in}}%
\pgfpathlineto{\pgfqpoint{1.190339in}{0.773441in}}%
\pgfpathlineto{\pgfqpoint{1.173409in}{0.765747in}}%
\pgfpathlineto{\pgfqpoint{1.156545in}{0.758707in}}%
\pgfpathlineto{\pgfqpoint{1.148320in}{0.760603in}}%
\pgfpathlineto{\pgfqpoint{1.140062in}{0.762947in}}%
\pgfpathclose%
\pgfusepath{fill}%
\end{pgfscope}%
\begin{pgfscope}%
\pgfpathrectangle{\pgfqpoint{0.407592in}{0.222904in}}{\pgfqpoint{1.411163in}{1.411163in}}%
\pgfusepath{clip}%
\pgfsetbuttcap%
\pgfsetroundjoin%
\definecolor{currentfill}{rgb}{0.556863,0.654902,0.847059}%
\pgfsetfillcolor{currentfill}%
\pgfsetfillopacity{0.130000}%
\pgfsetlinewidth{0.000000pt}%
\definecolor{currentstroke}{rgb}{0.000000,0.000000,0.000000}%
\pgfsetstrokecolor{currentstroke}%
\pgfsetdash{}{0pt}%
\pgfpathmoveto{\pgfqpoint{1.240634in}{0.790318in}}%
\pgfpathlineto{\pgfqpoint{1.257745in}{0.800025in}}%
\pgfpathlineto{\pgfqpoint{1.274929in}{0.810113in}}%
\pgfpathlineto{\pgfqpoint{1.282975in}{0.811160in}}%
\pgfpathlineto{\pgfqpoint{1.290995in}{0.812698in}}%
\pgfpathlineto{\pgfqpoint{1.273838in}{0.802245in}}%
\pgfpathlineto{\pgfqpoint{1.256754in}{0.792144in}}%
\pgfpathlineto{\pgfqpoint{1.248707in}{0.790967in}}%
\pgfpathlineto{\pgfqpoint{1.240634in}{0.790318in}}%
\pgfpathclose%
\pgfusepath{fill}%
\end{pgfscope}%
\begin{pgfscope}%
\pgfpathrectangle{\pgfqpoint{0.407592in}{0.222904in}}{\pgfqpoint{1.411163in}{1.411163in}}%
\pgfusepath{clip}%
\pgfsetbuttcap%
\pgfsetroundjoin%
\definecolor{currentfill}{rgb}{0.556863,0.654902,0.847059}%
\pgfsetfillcolor{currentfill}%
\pgfsetfillopacity{0.130000}%
\pgfsetlinewidth{0.000000pt}%
\definecolor{currentstroke}{rgb}{0.000000,0.000000,0.000000}%
\pgfsetstrokecolor{currentstroke}%
\pgfsetdash{}{0pt}%
\pgfpathmoveto{\pgfqpoint{1.089748in}{0.759419in}}%
\pgfpathlineto{\pgfqpoint{1.106566in}{0.763719in}}%
\pgfpathlineto{\pgfqpoint{1.123443in}{0.768841in}}%
\pgfpathlineto{\pgfqpoint{1.131770in}{0.765704in}}%
\pgfpathlineto{\pgfqpoint{1.140062in}{0.762947in}}%
\pgfpathlineto{\pgfqpoint{1.123225in}{0.757243in}}%
\pgfpathlineto{\pgfqpoint{1.106449in}{0.752370in}}%
\pgfpathlineto{\pgfqpoint{1.098117in}{0.755711in}}%
\pgfpathlineto{\pgfqpoint{1.089748in}{0.759419in}}%
\pgfpathclose%
\pgfusepath{fill}%
\end{pgfscope}%
\begin{pgfscope}%
\pgfpathrectangle{\pgfqpoint{0.407592in}{0.222904in}}{\pgfqpoint{1.411163in}{1.411163in}}%
\pgfusepath{clip}%
\pgfsetbuttcap%
\pgfsetroundjoin%
\definecolor{currentfill}{rgb}{0.556863,0.654902,0.847059}%
\pgfsetfillcolor{currentfill}%
\pgfsetfillopacity{0.130000}%
\pgfsetlinewidth{0.000000pt}%
\definecolor{currentstroke}{rgb}{0.000000,0.000000,0.000000}%
\pgfsetstrokecolor{currentstroke}%
\pgfsetdash{}{0pt}%
\pgfpathmoveto{\pgfqpoint{1.575410in}{1.005887in}}%
\pgfpathlineto{\pgfqpoint{1.593266in}{1.018386in}}%
\pgfpathlineto{\pgfqpoint{1.611203in}{1.031016in}}%
\pgfpathlineto{\pgfqpoint{1.618821in}{1.038859in}}%
\pgfpathlineto{\pgfqpoint{1.626426in}{1.047008in}}%
\pgfpathlineto{\pgfqpoint{1.608512in}{1.034574in}}%
\pgfpathlineto{\pgfqpoint{1.590679in}{1.022282in}}%
\pgfpathlineto{\pgfqpoint{1.583051in}{1.013922in}}%
\pgfpathlineto{\pgfqpoint{1.575410in}{1.005887in}}%
\pgfpathclose%
\pgfusepath{fill}%
\end{pgfscope}%
\begin{pgfscope}%
\pgfpathrectangle{\pgfqpoint{0.407592in}{0.222904in}}{\pgfqpoint{1.411163in}{1.411163in}}%
\pgfusepath{clip}%
\pgfsetbuttcap%
\pgfsetroundjoin%
\definecolor{currentfill}{rgb}{0.556863,0.654902,0.847059}%
\pgfsetfillcolor{currentfill}%
\pgfsetfillopacity{0.130000}%
\pgfsetlinewidth{0.000000pt}%
\definecolor{currentstroke}{rgb}{0.000000,0.000000,0.000000}%
\pgfsetstrokecolor{currentstroke}%
\pgfsetdash{}{0pt}%
\pgfpathmoveto{\pgfqpoint{1.458368in}{0.914632in}}%
\pgfpathlineto{\pgfqpoint{1.475974in}{0.926987in}}%
\pgfpathlineto{\pgfqpoint{1.493658in}{0.939481in}}%
\pgfpathlineto{\pgfqpoint{1.501409in}{0.945508in}}%
\pgfpathlineto{\pgfqpoint{1.509144in}{0.951921in}}%
\pgfpathlineto{\pgfqpoint{1.491481in}{0.939518in}}%
\pgfpathlineto{\pgfqpoint{1.473898in}{0.927261in}}%
\pgfpathlineto{\pgfqpoint{1.466141in}{0.920739in}}%
\pgfpathlineto{\pgfqpoint{1.458368in}{0.914632in}}%
\pgfpathclose%
\pgfusepath{fill}%
\end{pgfscope}%
\begin{pgfscope}%
\pgfpathrectangle{\pgfqpoint{0.407592in}{0.222904in}}{\pgfqpoint{1.411163in}{1.411163in}}%
\pgfusepath{clip}%
\pgfsetbuttcap%
\pgfsetroundjoin%
\definecolor{currentfill}{rgb}{0.556863,0.654902,0.847059}%
\pgfsetfillcolor{currentfill}%
\pgfsetfillopacity{0.130000}%
\pgfsetlinewidth{0.000000pt}%
\definecolor{currentstroke}{rgb}{0.000000,0.000000,0.000000}%
\pgfsetstrokecolor{currentstroke}%
\pgfsetdash{}{0pt}%
\pgfpathmoveto{\pgfqpoint{1.290995in}{0.812698in}}%
\pgfpathlineto{\pgfqpoint{1.308228in}{0.823463in}}%
\pgfpathlineto{\pgfqpoint{1.325536in}{0.834508in}}%
\pgfpathlineto{\pgfqpoint{1.333508in}{0.836824in}}%
\pgfpathlineto{\pgfqpoint{1.341457in}{0.839620in}}%
\pgfpathlineto{\pgfqpoint{1.324173in}{0.828336in}}%
\pgfpathlineto{\pgfqpoint{1.306964in}{0.817310in}}%
\pgfpathlineto{\pgfqpoint{1.298992in}{0.814743in}}%
\pgfpathlineto{\pgfqpoint{1.290995in}{0.812698in}}%
\pgfpathclose%
\pgfusepath{fill}%
\end{pgfscope}%
\begin{pgfscope}%
\pgfpathrectangle{\pgfqpoint{0.407592in}{0.222904in}}{\pgfqpoint{1.411163in}{1.411163in}}%
\pgfusepath{clip}%
\pgfsetbuttcap%
\pgfsetroundjoin%
\definecolor{currentfill}{rgb}{0.556863,0.654902,0.847059}%
\pgfsetfillcolor{currentfill}%
\pgfsetfillopacity{0.130000}%
\pgfsetlinewidth{0.000000pt}%
\definecolor{currentstroke}{rgb}{0.000000,0.000000,0.000000}%
\pgfsetstrokecolor{currentstroke}%
\pgfsetdash{}{0pt}%
\pgfpathmoveto{\pgfqpoint{1.692694in}{1.106903in}}%
\pgfpathlineto{\pgfqpoint{1.710802in}{1.119308in}}%
\pgfpathlineto{\pgfqpoint{1.728992in}{1.131852in}}%
\pgfpathlineto{\pgfqpoint{1.736487in}{1.140939in}}%
\pgfpathlineto{\pgfqpoint{1.743970in}{1.150268in}}%
\pgfpathlineto{\pgfqpoint{1.725805in}{1.137956in}}%
\pgfpathlineto{\pgfqpoint{1.707721in}{1.125793in}}%
\pgfpathlineto{\pgfqpoint{1.700213in}{1.116222in}}%
\pgfpathlineto{\pgfqpoint{1.692694in}{1.106903in}}%
\pgfpathclose%
\pgfusepath{fill}%
\end{pgfscope}%
\begin{pgfscope}%
\pgfpathrectangle{\pgfqpoint{0.407592in}{0.222904in}}{\pgfqpoint{1.411163in}{1.411163in}}%
\pgfusepath{clip}%
\pgfsetbuttcap%
\pgfsetroundjoin%
\definecolor{currentfill}{rgb}{0.556863,0.654902,0.847059}%
\pgfsetfillcolor{currentfill}%
\pgfsetfillopacity{0.130000}%
\pgfsetlinewidth{0.000000pt}%
\definecolor{currentstroke}{rgb}{0.000000,0.000000,0.000000}%
\pgfsetstrokecolor{currentstroke}%
\pgfsetdash{}{0pt}%
\pgfpathmoveto{\pgfqpoint{0.820562in}{0.837708in}}%
\pgfpathlineto{\pgfqpoint{0.837139in}{0.829532in}}%
\pgfpathlineto{\pgfqpoint{0.853737in}{0.822223in}}%
\pgfpathlineto{\pgfqpoint{0.862675in}{0.815848in}}%
\pgfpathlineto{\pgfqpoint{0.871571in}{0.809619in}}%
\pgfpathlineto{\pgfqpoint{0.855024in}{0.817136in}}%
\pgfpathlineto{\pgfqpoint{0.838498in}{0.825573in}}%
\pgfpathlineto{\pgfqpoint{0.829551in}{0.831561in}}%
\pgfpathlineto{\pgfqpoint{0.820562in}{0.837708in}}%
\pgfpathclose%
\pgfusepath{fill}%
\end{pgfscope}%
\begin{pgfscope}%
\pgfpathrectangle{\pgfqpoint{0.407592in}{0.222904in}}{\pgfqpoint{1.411163in}{1.411163in}}%
\pgfusepath{clip}%
\pgfsetbuttcap%
\pgfsetroundjoin%
\definecolor{currentfill}{rgb}{0.556863,0.654902,0.847059}%
\pgfsetfillcolor{currentfill}%
\pgfsetfillopacity{0.130000}%
\pgfsetlinewidth{0.000000pt}%
\definecolor{currentstroke}{rgb}{0.000000,0.000000,0.000000}%
\pgfsetstrokecolor{currentstroke}%
\pgfsetdash{}{0pt}%
\pgfpathmoveto{\pgfqpoint{0.618285in}{0.977389in}}%
\pgfpathlineto{\pgfqpoint{0.634815in}{0.963512in}}%
\pgfpathlineto{\pgfqpoint{0.651356in}{0.950106in}}%
\pgfpathlineto{\pgfqpoint{0.660714in}{0.944754in}}%
\pgfpathlineto{\pgfqpoint{0.670028in}{0.939582in}}%
\pgfpathlineto{\pgfqpoint{0.653533in}{0.953388in}}%
\pgfpathlineto{\pgfqpoint{0.637048in}{0.967665in}}%
\pgfpathlineto{\pgfqpoint{0.627689in}{0.972431in}}%
\pgfpathlineto{\pgfqpoint{0.618285in}{0.977389in}}%
\pgfpathclose%
\pgfusepath{fill}%
\end{pgfscope}%
\begin{pgfscope}%
\pgfpathrectangle{\pgfqpoint{0.407592in}{0.222904in}}{\pgfqpoint{1.411163in}{1.411163in}}%
\pgfusepath{clip}%
\pgfsetbuttcap%
\pgfsetroundjoin%
\definecolor{currentfill}{rgb}{0.556863,0.654902,0.847059}%
\pgfsetfillcolor{currentfill}%
\pgfsetfillopacity{0.130000}%
\pgfsetlinewidth{0.000000pt}%
\definecolor{currentstroke}{rgb}{0.000000,0.000000,0.000000}%
\pgfsetstrokecolor{currentstroke}%
\pgfsetdash{}{0pt}%
\pgfpathmoveto{\pgfqpoint{1.039339in}{0.762963in}}%
\pgfpathlineto{\pgfqpoint{1.056095in}{0.764919in}}%
\pgfpathlineto{\pgfqpoint{1.072902in}{0.767809in}}%
\pgfpathlineto{\pgfqpoint{1.081344in}{0.763462in}}%
\pgfpathlineto{\pgfqpoint{1.089748in}{0.759419in}}%
\pgfpathlineto{\pgfqpoint{1.072987in}{0.756019in}}%
\pgfpathlineto{\pgfqpoint{1.056278in}{0.753592in}}%
\pgfpathlineto{\pgfqpoint{1.047827in}{0.758137in}}%
\pgfpathlineto{\pgfqpoint{1.039339in}{0.762963in}}%
\pgfpathclose%
\pgfusepath{fill}%
\end{pgfscope}%
\begin{pgfscope}%
\pgfpathrectangle{\pgfqpoint{0.407592in}{0.222904in}}{\pgfqpoint{1.411163in}{1.411163in}}%
\pgfusepath{clip}%
\pgfsetbuttcap%
\pgfsetroundjoin%
\definecolor{currentfill}{rgb}{0.556863,0.654902,0.847059}%
\pgfsetfillcolor{currentfill}%
\pgfsetfillopacity{0.130000}%
\pgfsetlinewidth{0.000000pt}%
\definecolor{currentstroke}{rgb}{0.000000,0.000000,0.000000}%
\pgfsetstrokecolor{currentstroke}%
\pgfsetdash{}{0pt}%
\pgfpathmoveto{\pgfqpoint{0.904734in}{0.797463in}}%
\pgfpathlineto{\pgfqpoint{0.921357in}{0.792865in}}%
\pgfpathlineto{\pgfqpoint{0.938010in}{0.789262in}}%
\pgfpathlineto{\pgfqpoint{0.946760in}{0.783083in}}%
\pgfpathlineto{\pgfqpoint{0.955468in}{0.777068in}}%
\pgfpathlineto{\pgfqpoint{0.938867in}{0.780611in}}%
\pgfpathlineto{\pgfqpoint{0.922297in}{0.785227in}}%
\pgfpathlineto{\pgfqpoint{0.913536in}{0.791264in}}%
\pgfpathlineto{\pgfqpoint{0.904734in}{0.797463in}}%
\pgfpathclose%
\pgfusepath{fill}%
\end{pgfscope}%
\begin{pgfscope}%
\pgfpathrectangle{\pgfqpoint{0.407592in}{0.222904in}}{\pgfqpoint{1.411163in}{1.411163in}}%
\pgfusepath{clip}%
\pgfsetbuttcap%
\pgfsetroundjoin%
\definecolor{currentfill}{rgb}{0.556863,0.654902,0.847059}%
\pgfsetfillcolor{currentfill}%
\pgfsetfillopacity{0.130000}%
\pgfsetlinewidth{0.000000pt}%
\definecolor{currentstroke}{rgb}{0.000000,0.000000,0.000000}%
\pgfsetstrokecolor{currentstroke}%
\pgfsetdash{}{0pt}%
\pgfpathmoveto{\pgfqpoint{1.743970in}{1.150268in}}%
\pgfpathlineto{\pgfqpoint{1.762217in}{1.162721in}}%
\pgfpathlineto{\pgfqpoint{1.769677in}{1.172168in}}%
\pgfpathlineto{\pgfqpoint{1.777125in}{1.181840in}}%
\pgfpathlineto{\pgfqpoint{1.758903in}{1.169625in}}%
\pgfpathlineto{\pgfqpoint{1.751442in}{1.159832in}}%
\pgfusepath{fill}%
\end{pgfscope}%
\begin{pgfscope}%
\pgfpathrectangle{\pgfqpoint{0.407592in}{0.222904in}}{\pgfqpoint{1.411163in}{1.411163in}}%
\pgfusepath{clip}%
\pgfsetbuttcap%
\pgfsetroundjoin%
\definecolor{currentfill}{rgb}{0.556863,0.654902,0.847059}%
\pgfsetfillcolor{currentfill}%
\pgfsetfillopacity{0.130000}%
\pgfsetlinewidth{0.000000pt}%
\definecolor{currentstroke}{rgb}{0.000000,0.000000,0.000000}%
\pgfsetstrokecolor{currentstroke}%
\pgfsetdash{}{0pt}%
\pgfpathmoveto{\pgfqpoint{0.736126in}{0.889781in}}%
\pgfpathlineto{\pgfqpoint{0.752683in}{0.878873in}}%
\pgfpathlineto{\pgfqpoint{0.769257in}{0.868655in}}%
\pgfpathlineto{\pgfqpoint{0.778382in}{0.862513in}}%
\pgfpathlineto{\pgfqpoint{0.787465in}{0.856524in}}%
\pgfpathlineto{\pgfqpoint{0.770941in}{0.867089in}}%
\pgfpathlineto{\pgfqpoint{0.754432in}{0.878371in}}%
\pgfpathlineto{\pgfqpoint{0.745300in}{0.883991in}}%
\pgfpathlineto{\pgfqpoint{0.736126in}{0.889781in}}%
\pgfpathclose%
\pgfusepath{fill}%
\end{pgfscope}%
\begin{pgfscope}%
\pgfpathrectangle{\pgfqpoint{0.407592in}{0.222904in}}{\pgfqpoint{1.411163in}{1.411163in}}%
\pgfusepath{clip}%
\pgfsetbuttcap%
\pgfsetroundjoin%
\definecolor{currentfill}{rgb}{0.556863,0.654902,0.847059}%
\pgfsetfillcolor{currentfill}%
\pgfsetfillopacity{0.130000}%
\pgfsetlinewidth{0.000000pt}%
\definecolor{currentstroke}{rgb}{0.000000,0.000000,0.000000}%
\pgfsetstrokecolor{currentstroke}%
\pgfsetdash{}{0pt}%
\pgfpathmoveto{\pgfqpoint{1.341457in}{0.839620in}}%
\pgfpathlineto{\pgfqpoint{1.358817in}{0.851136in}}%
\pgfpathlineto{\pgfqpoint{1.376255in}{0.862863in}}%
\pgfpathlineto{\pgfqpoint{1.384159in}{0.866302in}}%
\pgfpathlineto{\pgfqpoint{1.392043in}{0.870198in}}%
\pgfpathlineto{\pgfqpoint{1.374628in}{0.858345in}}%
\pgfpathlineto{\pgfqpoint{1.357291in}{0.846689in}}%
\pgfpathlineto{\pgfqpoint{1.349384in}{0.842906in}}%
\pgfpathlineto{\pgfqpoint{1.341457in}{0.839620in}}%
\pgfpathclose%
\pgfusepath{fill}%
\end{pgfscope}%
\begin{pgfscope}%
\pgfpathrectangle{\pgfqpoint{0.407592in}{0.222904in}}{\pgfqpoint{1.411163in}{1.411163in}}%
\pgfusepath{clip}%
\pgfsetbuttcap%
\pgfsetroundjoin%
\definecolor{currentfill}{rgb}{0.556863,0.654902,0.847059}%
\pgfsetfillcolor{currentfill}%
\pgfsetfillopacity{0.130000}%
\pgfsetlinewidth{0.000000pt}%
\definecolor{currentstroke}{rgb}{0.000000,0.000000,0.000000}%
\pgfsetstrokecolor{currentstroke}%
\pgfsetdash{}{0pt}%
\pgfpathmoveto{\pgfqpoint{1.509144in}{0.951921in}}%
\pgfpathlineto{\pgfqpoint{1.526886in}{0.964460in}}%
\pgfpathlineto{\pgfqpoint{1.544708in}{0.977130in}}%
\pgfpathlineto{\pgfqpoint{1.552405in}{0.983799in}}%
\pgfpathlineto{\pgfqpoint{1.560087in}{0.990819in}}%
\pgfpathlineto{\pgfqpoint{1.542288in}{0.978282in}}%
\pgfpathlineto{\pgfqpoint{1.524568in}{0.965882in}}%
\pgfpathlineto{\pgfqpoint{1.516863in}{0.958714in}}%
\pgfpathlineto{\pgfqpoint{1.509144in}{0.951921in}}%
\pgfpathclose%
\pgfusepath{fill}%
\end{pgfscope}%
\begin{pgfscope}%
\pgfpathrectangle{\pgfqpoint{0.407592in}{0.222904in}}{\pgfqpoint{1.411163in}{1.411163in}}%
\pgfusepath{clip}%
\pgfsetbuttcap%
\pgfsetroundjoin%
\definecolor{currentfill}{rgb}{0.556863,0.654902,0.847059}%
\pgfsetfillcolor{currentfill}%
\pgfsetfillopacity{0.130000}%
\pgfsetlinewidth{0.000000pt}%
\definecolor{currentstroke}{rgb}{0.000000,0.000000,0.000000}%
\pgfsetstrokecolor{currentstroke}%
\pgfsetdash{}{0pt}%
\pgfpathmoveto{\pgfqpoint{1.626426in}{1.047008in}}%
\pgfpathlineto{\pgfqpoint{1.644420in}{1.059579in}}%
\pgfpathlineto{\pgfqpoint{1.662495in}{1.072279in}}%
\pgfpathlineto{\pgfqpoint{1.670063in}{1.080523in}}%
\pgfpathlineto{\pgfqpoint{1.677619in}{1.089046in}}%
\pgfpathlineto{\pgfqpoint{1.659568in}{1.076553in}}%
\pgfpathlineto{\pgfqpoint{1.641597in}{1.064199in}}%
\pgfpathlineto{\pgfqpoint{1.634018in}{1.055458in}}%
\pgfpathlineto{\pgfqpoint{1.626426in}{1.047008in}}%
\pgfpathclose%
\pgfusepath{fill}%
\end{pgfscope}%
\begin{pgfscope}%
\pgfpathrectangle{\pgfqpoint{0.407592in}{0.222904in}}{\pgfqpoint{1.411163in}{1.411163in}}%
\pgfusepath{clip}%
\pgfsetbuttcap%
\pgfsetroundjoin%
\definecolor{currentfill}{rgb}{0.556863,0.654902,0.847059}%
\pgfsetfillcolor{currentfill}%
\pgfsetfillopacity{0.130000}%
\pgfsetlinewidth{0.000000pt}%
\definecolor{currentstroke}{rgb}{0.000000,0.000000,0.000000}%
\pgfsetstrokecolor{currentstroke}%
\pgfsetdash{}{0pt}%
\pgfpathmoveto{\pgfqpoint{0.988776in}{0.773171in}}%
\pgfpathlineto{\pgfqpoint{1.005489in}{0.772772in}}%
\pgfpathlineto{\pgfqpoint{1.022246in}{0.773360in}}%
\pgfpathlineto{\pgfqpoint{1.030812in}{0.768046in}}%
\pgfpathlineto{\pgfqpoint{1.039339in}{0.762963in}}%
\pgfpathlineto{\pgfqpoint{1.022631in}{0.762003in}}%
\pgfpathlineto{\pgfqpoint{1.005969in}{0.762093in}}%
\pgfpathlineto{\pgfqpoint{0.997392in}{0.767526in}}%
\pgfpathlineto{\pgfqpoint{0.988776in}{0.773171in}}%
\pgfpathclose%
\pgfusepath{fill}%
\end{pgfscope}%
\begin{pgfscope}%
\pgfpathrectangle{\pgfqpoint{0.407592in}{0.222904in}}{\pgfqpoint{1.411163in}{1.411163in}}%
\pgfusepath{clip}%
\pgfsetbuttcap%
\pgfsetroundjoin%
\definecolor{currentfill}{rgb}{0.556863,0.654902,0.847059}%
\pgfsetfillcolor{currentfill}%
\pgfsetfillopacity{0.130000}%
\pgfsetlinewidth{0.000000pt}%
\definecolor{currentstroke}{rgb}{0.000000,0.000000,0.000000}%
\pgfsetstrokecolor{currentstroke}%
\pgfsetdash{}{0pt}%
\pgfpathmoveto{\pgfqpoint{0.651356in}{0.950106in}}%
\pgfpathlineto{\pgfqpoint{0.667909in}{0.937202in}}%
\pgfpathlineto{\pgfqpoint{0.684474in}{0.924832in}}%
\pgfpathlineto{\pgfqpoint{0.693785in}{0.919095in}}%
\pgfpathlineto{\pgfqpoint{0.703052in}{0.913523in}}%
\pgfpathlineto{\pgfqpoint{0.686534in}{0.926282in}}%
\pgfpathlineto{\pgfqpoint{0.670028in}{0.939582in}}%
\pgfpathlineto{\pgfqpoint{0.660714in}{0.944754in}}%
\pgfpathlineto{\pgfqpoint{0.651356in}{0.950106in}}%
\pgfpathclose%
\pgfusepath{fill}%
\end{pgfscope}%
\begin{pgfscope}%
\pgfpathrectangle{\pgfqpoint{0.407592in}{0.222904in}}{\pgfqpoint{1.411163in}{1.411163in}}%
\pgfusepath{clip}%
\pgfsetbuttcap%
\pgfsetroundjoin%
\definecolor{currentfill}{rgb}{0.556863,0.654902,0.847059}%
\pgfsetfillcolor{currentfill}%
\pgfsetfillopacity{0.130000}%
\pgfsetlinewidth{0.000000pt}%
\definecolor{currentstroke}{rgb}{0.000000,0.000000,0.000000}%
\pgfsetstrokecolor{currentstroke}%
\pgfsetdash{}{0pt}%
\pgfpathmoveto{\pgfqpoint{1.392043in}{0.870198in}}%
\pgfpathlineto{\pgfqpoint{1.409536in}{0.882232in}}%
\pgfpathlineto{\pgfqpoint{1.427107in}{0.894433in}}%
\pgfpathlineto{\pgfqpoint{1.434950in}{0.898846in}}%
\pgfpathlineto{\pgfqpoint{1.442773in}{0.903684in}}%
\pgfpathlineto{\pgfqpoint{1.425224in}{0.891446in}}%
\pgfpathlineto{\pgfqpoint{1.407754in}{0.879371in}}%
\pgfpathlineto{\pgfqpoint{1.399908in}{0.874554in}}%
\pgfpathlineto{\pgfqpoint{1.392043in}{0.870198in}}%
\pgfpathclose%
\pgfusepath{fill}%
\end{pgfscope}%
\begin{pgfscope}%
\pgfpathrectangle{\pgfqpoint{0.407592in}{0.222904in}}{\pgfqpoint{1.411163in}{1.411163in}}%
\pgfusepath{clip}%
\pgfsetbuttcap%
\pgfsetroundjoin%
\definecolor{currentfill}{rgb}{0.556863,0.654902,0.847059}%
\pgfsetfillcolor{currentfill}%
\pgfsetfillopacity{0.130000}%
\pgfsetlinewidth{0.000000pt}%
\definecolor{currentstroke}{rgb}{0.000000,0.000000,0.000000}%
\pgfsetstrokecolor{currentstroke}%
\pgfsetdash{}{0pt}%
\pgfpathmoveto{\pgfqpoint{0.853737in}{0.822223in}}%
\pgfpathlineto{\pgfqpoint{0.870358in}{0.815804in}}%
\pgfpathlineto{\pgfqpoint{0.887006in}{0.810297in}}%
\pgfpathlineto{\pgfqpoint{0.895891in}{0.803812in}}%
\pgfpathlineto{\pgfqpoint{0.904734in}{0.797463in}}%
\pgfpathlineto{\pgfqpoint{0.888140in}{0.803053in}}%
\pgfpathlineto{\pgfqpoint{0.871571in}{0.809619in}}%
\pgfpathlineto{\pgfqpoint{0.862675in}{0.815848in}}%
\pgfpathlineto{\pgfqpoint{0.853737in}{0.822223in}}%
\pgfpathclose%
\pgfusepath{fill}%
\end{pgfscope}%
\begin{pgfscope}%
\pgfpathrectangle{\pgfqpoint{0.407592in}{0.222904in}}{\pgfqpoint{1.411163in}{1.411163in}}%
\pgfusepath{clip}%
\pgfsetbuttcap%
\pgfsetroundjoin%
\definecolor{currentfill}{rgb}{0.556863,0.654902,0.847059}%
\pgfsetfillcolor{currentfill}%
\pgfsetfillopacity{0.130000}%
\pgfsetlinewidth{0.000000pt}%
\definecolor{currentstroke}{rgb}{0.000000,0.000000,0.000000}%
\pgfsetstrokecolor{currentstroke}%
\pgfsetdash{}{0pt}%
\pgfpathmoveto{\pgfqpoint{0.769257in}{0.868655in}}%
\pgfpathlineto{\pgfqpoint{0.785847in}{0.859163in}}%
\pgfpathlineto{\pgfqpoint{0.802455in}{0.850431in}}%
\pgfpathlineto{\pgfqpoint{0.811530in}{0.844001in}}%
\pgfpathlineto{\pgfqpoint{0.820562in}{0.837708in}}%
\pgfpathlineto{\pgfqpoint{0.804005in}{0.846717in}}%
\pgfpathlineto{\pgfqpoint{0.787465in}{0.856524in}}%
\pgfpathlineto{\pgfqpoint{0.778382in}{0.862513in}}%
\pgfpathlineto{\pgfqpoint{0.769257in}{0.868655in}}%
\pgfpathclose%
\pgfusepath{fill}%
\end{pgfscope}%
\begin{pgfscope}%
\pgfpathrectangle{\pgfqpoint{0.407592in}{0.222904in}}{\pgfqpoint{1.411163in}{1.411163in}}%
\pgfusepath{clip}%
\pgfsetbuttcap%
\pgfsetroundjoin%
\definecolor{currentfill}{rgb}{0.556863,0.654902,0.847059}%
\pgfsetfillcolor{currentfill}%
\pgfsetfillopacity{0.130000}%
\pgfsetlinewidth{0.000000pt}%
\definecolor{currentstroke}{rgb}{0.000000,0.000000,0.000000}%
\pgfsetstrokecolor{currentstroke}%
\pgfsetdash{}{0pt}%
\pgfpathmoveto{\pgfqpoint{1.560087in}{0.990819in}}%
\pgfpathlineto{\pgfqpoint{1.577966in}{1.003486in}}%
\pgfpathlineto{\pgfqpoint{1.595927in}{1.016276in}}%
\pgfpathlineto{\pgfqpoint{1.603572in}{1.023487in}}%
\pgfpathlineto{\pgfqpoint{1.611203in}{1.031016in}}%
\pgfpathlineto{\pgfqpoint{1.593266in}{1.018386in}}%
\pgfpathlineto{\pgfqpoint{1.575410in}{1.005887in}}%
\pgfpathlineto{\pgfqpoint{1.567755in}{0.998184in}}%
\pgfpathlineto{\pgfqpoint{1.560087in}{0.990819in}}%
\pgfpathclose%
\pgfusepath{fill}%
\end{pgfscope}%
\begin{pgfscope}%
\pgfpathrectangle{\pgfqpoint{0.407592in}{0.222904in}}{\pgfqpoint{1.411163in}{1.411163in}}%
\pgfusepath{clip}%
\pgfsetbuttcap%
\pgfsetroundjoin%
\definecolor{currentfill}{rgb}{0.556863,0.654902,0.847059}%
\pgfsetfillcolor{currentfill}%
\pgfsetfillopacity{0.130000}%
\pgfsetlinewidth{0.000000pt}%
\definecolor{currentstroke}{rgb}{0.000000,0.000000,0.000000}%
\pgfsetstrokecolor{currentstroke}%
\pgfsetdash{}{0pt}%
\pgfpathmoveto{\pgfqpoint{1.677619in}{1.089046in}}%
\pgfpathlineto{\pgfqpoint{1.695752in}{1.101674in}}%
\pgfpathlineto{\pgfqpoint{1.713967in}{1.114429in}}%
\pgfpathlineto{\pgfqpoint{1.721485in}{1.123013in}}%
\pgfpathlineto{\pgfqpoint{1.728992in}{1.131852in}}%
\pgfpathlineto{\pgfqpoint{1.710802in}{1.119308in}}%
\pgfpathlineto{\pgfqpoint{1.692694in}{1.106903in}}%
\pgfpathlineto{\pgfqpoint{1.685162in}{1.097842in}}%
\pgfpathlineto{\pgfqpoint{1.677619in}{1.089046in}}%
\pgfpathclose%
\pgfusepath{fill}%
\end{pgfscope}%
\begin{pgfscope}%
\pgfpathrectangle{\pgfqpoint{0.407592in}{0.222904in}}{\pgfqpoint{1.411163in}{1.411163in}}%
\pgfusepath{clip}%
\pgfsetbuttcap%
\pgfsetroundjoin%
\definecolor{currentfill}{rgb}{0.556863,0.654902,0.847059}%
\pgfsetfillcolor{currentfill}%
\pgfsetfillopacity{0.130000}%
\pgfsetlinewidth{0.000000pt}%
\definecolor{currentstroke}{rgb}{0.000000,0.000000,0.000000}%
\pgfsetstrokecolor{currentstroke}%
\pgfsetdash{}{0pt}%
\pgfpathmoveto{\pgfqpoint{1.442773in}{0.903684in}}%
\pgfpathlineto{\pgfqpoint{1.460401in}{0.916071in}}%
\pgfpathlineto{\pgfqpoint{1.478108in}{0.928598in}}%
\pgfpathlineto{\pgfqpoint{1.485892in}{0.933844in}}%
\pgfpathlineto{\pgfqpoint{1.493658in}{0.939481in}}%
\pgfpathlineto{\pgfqpoint{1.475974in}{0.926987in}}%
\pgfpathlineto{\pgfqpoint{1.458368in}{0.914632in}}%
\pgfpathlineto{\pgfqpoint{1.450579in}{0.908946in}}%
\pgfpathlineto{\pgfqpoint{1.442773in}{0.903684in}}%
\pgfpathclose%
\pgfusepath{fill}%
\end{pgfscope}%
\begin{pgfscope}%
\pgfpathrectangle{\pgfqpoint{0.407592in}{0.222904in}}{\pgfqpoint{1.411163in}{1.411163in}}%
\pgfusepath{clip}%
\pgfsetbuttcap%
\pgfsetroundjoin%
\definecolor{currentfill}{rgb}{0.556863,0.654902,0.847059}%
\pgfsetfillcolor{currentfill}%
\pgfsetfillopacity{0.130000}%
\pgfsetlinewidth{0.000000pt}%
\definecolor{currentstroke}{rgb}{0.000000,0.000000,0.000000}%
\pgfsetstrokecolor{currentstroke}%
\pgfsetdash{}{0pt}%
\pgfpathmoveto{\pgfqpoint{1.173926in}{0.776526in}}%
\pgfpathlineto{\pgfqpoint{1.190958in}{0.784255in}}%
\pgfpathlineto{\pgfqpoint{1.208059in}{0.792526in}}%
\pgfpathlineto{\pgfqpoint{1.216247in}{0.791301in}}%
\pgfpathlineto{\pgfqpoint{1.224405in}{0.790509in}}%
\pgfpathlineto{\pgfqpoint{1.207337in}{0.781717in}}%
\pgfpathlineto{\pgfqpoint{1.190339in}{0.773441in}}%
\pgfpathlineto{\pgfqpoint{1.182148in}{0.774762in}}%
\pgfpathlineto{\pgfqpoint{1.173926in}{0.776526in}}%
\pgfpathclose%
\pgfusepath{fill}%
\end{pgfscope}%
\begin{pgfscope}%
\pgfpathrectangle{\pgfqpoint{0.407592in}{0.222904in}}{\pgfqpoint{1.411163in}{1.411163in}}%
\pgfusepath{clip}%
\pgfsetbuttcap%
\pgfsetroundjoin%
\definecolor{currentfill}{rgb}{0.556863,0.654902,0.847059}%
\pgfsetfillcolor{currentfill}%
\pgfsetfillopacity{0.130000}%
\pgfsetlinewidth{0.000000pt}%
\definecolor{currentstroke}{rgb}{0.000000,0.000000,0.000000}%
\pgfsetstrokecolor{currentstroke}%
\pgfsetdash{}{0pt}%
\pgfpathmoveto{\pgfqpoint{0.938010in}{0.789262in}}%
\pgfpathlineto{\pgfqpoint{0.954698in}{0.786652in}}%
\pgfpathlineto{\pgfqpoint{0.971422in}{0.785022in}}%
\pgfpathlineto{\pgfqpoint{0.980119in}{0.779009in}}%
\pgfpathlineto{\pgfqpoint{0.988776in}{0.773171in}}%
\pgfpathlineto{\pgfqpoint{0.972103in}{0.774594in}}%
\pgfpathlineto{\pgfqpoint{0.955468in}{0.777068in}}%
\pgfpathlineto{\pgfqpoint{0.946760in}{0.783083in}}%
\pgfpathlineto{\pgfqpoint{0.938010in}{0.789262in}}%
\pgfpathclose%
\pgfusepath{fill}%
\end{pgfscope}%
\begin{pgfscope}%
\pgfpathrectangle{\pgfqpoint{0.407592in}{0.222904in}}{\pgfqpoint{1.411163in}{1.411163in}}%
\pgfusepath{clip}%
\pgfsetbuttcap%
\pgfsetroundjoin%
\definecolor{currentfill}{rgb}{0.556863,0.654902,0.847059}%
\pgfsetfillcolor{currentfill}%
\pgfsetfillopacity{0.130000}%
\pgfsetlinewidth{0.000000pt}%
\definecolor{currentstroke}{rgb}{0.000000,0.000000,0.000000}%
\pgfsetstrokecolor{currentstroke}%
\pgfsetdash{}{0pt}%
\pgfpathmoveto{\pgfqpoint{1.224405in}{0.790509in}}%
\pgfpathlineto{\pgfqpoint{1.241545in}{0.799761in}}%
\pgfpathlineto{\pgfqpoint{1.258758in}{0.809424in}}%
\pgfpathlineto{\pgfqpoint{1.266857in}{0.809540in}}%
\pgfpathlineto{\pgfqpoint{1.274929in}{0.810113in}}%
\pgfpathlineto{\pgfqpoint{1.257745in}{0.800025in}}%
\pgfpathlineto{\pgfqpoint{1.240634in}{0.790318in}}%
\pgfpathlineto{\pgfqpoint{1.232534in}{0.790173in}}%
\pgfpathlineto{\pgfqpoint{1.224405in}{0.790509in}}%
\pgfpathclose%
\pgfusepath{fill}%
\end{pgfscope}%
\begin{pgfscope}%
\pgfpathrectangle{\pgfqpoint{0.407592in}{0.222904in}}{\pgfqpoint{1.411163in}{1.411163in}}%
\pgfusepath{clip}%
\pgfsetbuttcap%
\pgfsetroundjoin%
\definecolor{currentfill}{rgb}{0.556863,0.654902,0.847059}%
\pgfsetfillcolor{currentfill}%
\pgfsetfillopacity{0.130000}%
\pgfsetlinewidth{0.000000pt}%
\definecolor{currentstroke}{rgb}{0.000000,0.000000,0.000000}%
\pgfsetstrokecolor{currentstroke}%
\pgfsetdash{}{0pt}%
\pgfpathmoveto{\pgfqpoint{1.123443in}{0.768841in}}%
\pgfpathlineto{\pgfqpoint{1.140382in}{0.774713in}}%
\pgfpathlineto{\pgfqpoint{1.157384in}{0.781268in}}%
\pgfpathlineto{\pgfqpoint{1.165672in}{0.778704in}}%
\pgfpathlineto{\pgfqpoint{1.173926in}{0.776526in}}%
\pgfpathlineto{\pgfqpoint{1.156961in}{0.769400in}}%
\pgfpathlineto{\pgfqpoint{1.140062in}{0.762947in}}%
\pgfpathlineto{\pgfqpoint{1.131770in}{0.765704in}}%
\pgfpathlineto{\pgfqpoint{1.123443in}{0.768841in}}%
\pgfpathclose%
\pgfusepath{fill}%
\end{pgfscope}%
\begin{pgfscope}%
\pgfpathrectangle{\pgfqpoint{0.407592in}{0.222904in}}{\pgfqpoint{1.411163in}{1.411163in}}%
\pgfusepath{clip}%
\pgfsetbuttcap%
\pgfsetroundjoin%
\definecolor{currentfill}{rgb}{0.556863,0.654902,0.847059}%
\pgfsetfillcolor{currentfill}%
\pgfsetfillopacity{0.130000}%
\pgfsetlinewidth{0.000000pt}%
\definecolor{currentstroke}{rgb}{0.000000,0.000000,0.000000}%
\pgfsetstrokecolor{currentstroke}%
\pgfsetdash{}{0pt}%
\pgfpathmoveto{\pgfqpoint{1.728992in}{1.131852in}}%
\pgfpathlineto{\pgfqpoint{1.747264in}{1.144528in}}%
\pgfpathlineto{\pgfqpoint{1.754747in}{1.153506in}}%
\pgfpathlineto{\pgfqpoint{1.762217in}{1.162721in}}%
\pgfpathlineto{\pgfqpoint{1.743970in}{1.150268in}}%
\pgfpathlineto{\pgfqpoint{1.736487in}{1.140939in}}%
\pgfusepath{fill}%
\end{pgfscope}%
\begin{pgfscope}%
\pgfpathrectangle{\pgfqpoint{0.407592in}{0.222904in}}{\pgfqpoint{1.411163in}{1.411163in}}%
\pgfusepath{clip}%
\pgfsetbuttcap%
\pgfsetroundjoin%
\definecolor{currentfill}{rgb}{0.556863,0.654902,0.847059}%
\pgfsetfillcolor{currentfill}%
\pgfsetfillopacity{0.130000}%
\pgfsetlinewidth{0.000000pt}%
\definecolor{currentstroke}{rgb}{0.000000,0.000000,0.000000}%
\pgfsetstrokecolor{currentstroke}%
\pgfsetdash{}{0pt}%
\pgfpathmoveto{\pgfqpoint{0.684474in}{0.924832in}}%
\pgfpathlineto{\pgfqpoint{0.701052in}{0.913030in}}%
\pgfpathlineto{\pgfqpoint{0.717644in}{0.901828in}}%
\pgfpathlineto{\pgfqpoint{0.726907in}{0.895730in}}%
\pgfpathlineto{\pgfqpoint{0.736126in}{0.889781in}}%
\pgfpathlineto{\pgfqpoint{0.719582in}{0.901344in}}%
\pgfpathlineto{\pgfqpoint{0.703052in}{0.913523in}}%
\pgfpathlineto{\pgfqpoint{0.693785in}{0.919095in}}%
\pgfpathlineto{\pgfqpoint{0.684474in}{0.924832in}}%
\pgfpathclose%
\pgfusepath{fill}%
\end{pgfscope}%
\begin{pgfscope}%
\pgfpathrectangle{\pgfqpoint{0.407592in}{0.222904in}}{\pgfqpoint{1.411163in}{1.411163in}}%
\pgfusepath{clip}%
\pgfsetbuttcap%
\pgfsetroundjoin%
\definecolor{currentfill}{rgb}{0.556863,0.654902,0.847059}%
\pgfsetfillcolor{currentfill}%
\pgfsetfillopacity{0.130000}%
\pgfsetlinewidth{0.000000pt}%
\definecolor{currentstroke}{rgb}{0.000000,0.000000,0.000000}%
\pgfsetstrokecolor{currentstroke}%
\pgfsetdash{}{0pt}%
\pgfpathmoveto{\pgfqpoint{1.274929in}{0.810113in}}%
\pgfpathlineto{\pgfqpoint{1.292188in}{0.820541in}}%
\pgfpathlineto{\pgfqpoint{1.309521in}{0.831275in}}%
\pgfpathlineto{\pgfqpoint{1.317541in}{0.832663in}}%
\pgfpathlineto{\pgfqpoint{1.325536in}{0.834508in}}%
\pgfpathlineto{\pgfqpoint{1.308228in}{0.823463in}}%
\pgfpathlineto{\pgfqpoint{1.290995in}{0.812698in}}%
\pgfpathlineto{\pgfqpoint{1.282975in}{0.811160in}}%
\pgfpathlineto{\pgfqpoint{1.274929in}{0.810113in}}%
\pgfpathclose%
\pgfusepath{fill}%
\end{pgfscope}%
\begin{pgfscope}%
\pgfpathrectangle{\pgfqpoint{0.407592in}{0.222904in}}{\pgfqpoint{1.411163in}{1.411163in}}%
\pgfusepath{clip}%
\pgfsetbuttcap%
\pgfsetroundjoin%
\definecolor{currentfill}{rgb}{0.556863,0.654902,0.847059}%
\pgfsetfillcolor{currentfill}%
\pgfsetfillopacity{0.130000}%
\pgfsetlinewidth{0.000000pt}%
\definecolor{currentstroke}{rgb}{0.000000,0.000000,0.000000}%
\pgfsetstrokecolor{currentstroke}%
\pgfsetdash{}{0pt}%
\pgfpathmoveto{\pgfqpoint{1.072902in}{0.767809in}}%
\pgfpathlineto{\pgfqpoint{1.089764in}{0.771570in}}%
\pgfpathlineto{\pgfqpoint{1.106684in}{0.776140in}}%
\pgfpathlineto{\pgfqpoint{1.115081in}{0.772329in}}%
\pgfpathlineto{\pgfqpoint{1.123443in}{0.768841in}}%
\pgfpathlineto{\pgfqpoint{1.106566in}{0.763719in}}%
\pgfpathlineto{\pgfqpoint{1.089748in}{0.759419in}}%
\pgfpathlineto{\pgfqpoint{1.081344in}{0.763462in}}%
\pgfpathlineto{\pgfqpoint{1.072902in}{0.767809in}}%
\pgfpathclose%
\pgfusepath{fill}%
\end{pgfscope}%
\begin{pgfscope}%
\pgfpathrectangle{\pgfqpoint{0.407592in}{0.222904in}}{\pgfqpoint{1.411163in}{1.411163in}}%
\pgfusepath{clip}%
\pgfsetbuttcap%
\pgfsetroundjoin%
\definecolor{currentfill}{rgb}{0.556863,0.654902,0.847059}%
\pgfsetfillcolor{currentfill}%
\pgfsetfillopacity{0.130000}%
\pgfsetlinewidth{0.000000pt}%
\definecolor{currentstroke}{rgb}{0.000000,0.000000,0.000000}%
\pgfsetstrokecolor{currentstroke}%
\pgfsetdash{}{0pt}%
\pgfpathmoveto{\pgfqpoint{1.611203in}{1.031016in}}%
\pgfpathlineto{\pgfqpoint{1.629221in}{1.043773in}}%
\pgfpathlineto{\pgfqpoint{1.647320in}{1.056650in}}%
\pgfpathlineto{\pgfqpoint{1.654914in}{1.064319in}}%
\pgfpathlineto{\pgfqpoint{1.662495in}{1.072279in}}%
\pgfpathlineto{\pgfqpoint{1.644420in}{1.059579in}}%
\pgfpathlineto{\pgfqpoint{1.626426in}{1.047008in}}%
\pgfpathlineto{\pgfqpoint{1.618821in}{1.038859in}}%
\pgfpathlineto{\pgfqpoint{1.611203in}{1.031016in}}%
\pgfpathclose%
\pgfusepath{fill}%
\end{pgfscope}%
\begin{pgfscope}%
\pgfpathrectangle{\pgfqpoint{0.407592in}{0.222904in}}{\pgfqpoint{1.411163in}{1.411163in}}%
\pgfusepath{clip}%
\pgfsetbuttcap%
\pgfsetroundjoin%
\definecolor{currentfill}{rgb}{0.556863,0.654902,0.847059}%
\pgfsetfillcolor{currentfill}%
\pgfsetfillopacity{0.130000}%
\pgfsetlinewidth{0.000000pt}%
\definecolor{currentstroke}{rgb}{0.000000,0.000000,0.000000}%
\pgfsetstrokecolor{currentstroke}%
\pgfsetdash{}{0pt}%
\pgfpathmoveto{\pgfqpoint{1.493658in}{0.939481in}}%
\pgfpathlineto{\pgfqpoint{1.511423in}{0.952108in}}%
\pgfpathlineto{\pgfqpoint{1.529268in}{0.964860in}}%
\pgfpathlineto{\pgfqpoint{1.536995in}{0.970815in}}%
\pgfpathlineto{\pgfqpoint{1.544708in}{0.977130in}}%
\pgfpathlineto{\pgfqpoint{1.526886in}{0.964460in}}%
\pgfpathlineto{\pgfqpoint{1.509144in}{0.951921in}}%
\pgfpathlineto{\pgfqpoint{1.501409in}{0.945508in}}%
\pgfpathlineto{\pgfqpoint{1.493658in}{0.939481in}}%
\pgfpathclose%
\pgfusepath{fill}%
\end{pgfscope}%
\begin{pgfscope}%
\pgfpathrectangle{\pgfqpoint{0.407592in}{0.222904in}}{\pgfqpoint{1.411163in}{1.411163in}}%
\pgfusepath{clip}%
\pgfsetbuttcap%
\pgfsetroundjoin%
\definecolor{currentfill}{rgb}{0.556863,0.654902,0.847059}%
\pgfsetfillcolor{currentfill}%
\pgfsetfillopacity{0.130000}%
\pgfsetlinewidth{0.000000pt}%
\definecolor{currentstroke}{rgb}{0.000000,0.000000,0.000000}%
\pgfsetstrokecolor{currentstroke}%
\pgfsetdash{}{0pt}%
\pgfpathmoveto{\pgfqpoint{1.325536in}{0.834508in}}%
\pgfpathlineto{\pgfqpoint{1.342921in}{0.845805in}}%
\pgfpathlineto{\pgfqpoint{1.360383in}{0.857330in}}%
\pgfpathlineto{\pgfqpoint{1.368330in}{0.859874in}}%
\pgfpathlineto{\pgfqpoint{1.376255in}{0.862863in}}%
\pgfpathlineto{\pgfqpoint{1.358817in}{0.851136in}}%
\pgfpathlineto{\pgfqpoint{1.341457in}{0.839620in}}%
\pgfpathlineto{\pgfqpoint{1.333508in}{0.836824in}}%
\pgfpathlineto{\pgfqpoint{1.325536in}{0.834508in}}%
\pgfpathclose%
\pgfusepath{fill}%
\end{pgfscope}%
\begin{pgfscope}%
\pgfpathrectangle{\pgfqpoint{0.407592in}{0.222904in}}{\pgfqpoint{1.411163in}{1.411163in}}%
\pgfusepath{clip}%
\pgfsetbuttcap%
\pgfsetroundjoin%
\definecolor{currentfill}{rgb}{0.556863,0.654902,0.847059}%
\pgfsetfillcolor{currentfill}%
\pgfsetfillopacity{0.130000}%
\pgfsetlinewidth{0.000000pt}%
\definecolor{currentstroke}{rgb}{0.000000,0.000000,0.000000}%
\pgfsetstrokecolor{currentstroke}%
\pgfsetdash{}{0pt}%
\pgfpathmoveto{\pgfqpoint{0.802455in}{0.850431in}}%
\pgfpathlineto{\pgfqpoint{0.819083in}{0.842488in}}%
\pgfpathlineto{\pgfqpoint{0.835734in}{0.835362in}}%
\pgfpathlineto{\pgfqpoint{0.844757in}{0.828730in}}%
\pgfpathlineto{\pgfqpoint{0.853737in}{0.822223in}}%
\pgfpathlineto{\pgfqpoint{0.837139in}{0.829532in}}%
\pgfpathlineto{\pgfqpoint{0.820562in}{0.837708in}}%
\pgfpathlineto{\pgfqpoint{0.811530in}{0.844001in}}%
\pgfpathlineto{\pgfqpoint{0.802455in}{0.850431in}}%
\pgfpathclose%
\pgfusepath{fill}%
\end{pgfscope}%
\begin{pgfscope}%
\pgfpathrectangle{\pgfqpoint{0.407592in}{0.222904in}}{\pgfqpoint{1.411163in}{1.411163in}}%
\pgfusepath{clip}%
\pgfsetbuttcap%
\pgfsetroundjoin%
\definecolor{currentfill}{rgb}{0.556863,0.654902,0.847059}%
\pgfsetfillcolor{currentfill}%
\pgfsetfillopacity{0.130000}%
\pgfsetlinewidth{0.000000pt}%
\definecolor{currentstroke}{rgb}{0.000000,0.000000,0.000000}%
\pgfsetstrokecolor{currentstroke}%
\pgfsetdash{}{0pt}%
\pgfpathmoveto{\pgfqpoint{1.022246in}{0.773360in}}%
\pgfpathlineto{\pgfqpoint{1.039050in}{0.774890in}}%
\pgfpathlineto{\pgfqpoint{1.055905in}{0.777314in}}%
\pgfpathlineto{\pgfqpoint{1.064423in}{0.772434in}}%
\pgfpathlineto{\pgfqpoint{1.072902in}{0.767809in}}%
\pgfpathlineto{\pgfqpoint{1.056095in}{0.764919in}}%
\pgfpathlineto{\pgfqpoint{1.039339in}{0.762963in}}%
\pgfpathlineto{\pgfqpoint{1.030812in}{0.768046in}}%
\pgfpathlineto{\pgfqpoint{1.022246in}{0.773360in}}%
\pgfpathclose%
\pgfusepath{fill}%
\end{pgfscope}%
\begin{pgfscope}%
\pgfpathrectangle{\pgfqpoint{0.407592in}{0.222904in}}{\pgfqpoint{1.411163in}{1.411163in}}%
\pgfusepath{clip}%
\pgfsetbuttcap%
\pgfsetroundjoin%
\definecolor{currentfill}{rgb}{0.556863,0.654902,0.847059}%
\pgfsetfillcolor{currentfill}%
\pgfsetfillopacity{0.130000}%
\pgfsetlinewidth{0.000000pt}%
\definecolor{currentstroke}{rgb}{0.000000,0.000000,0.000000}%
\pgfsetstrokecolor{currentstroke}%
\pgfsetdash{}{0pt}%
\pgfpathmoveto{\pgfqpoint{0.887006in}{0.810297in}}%
\pgfpathlineto{\pgfqpoint{0.903681in}{0.805714in}}%
\pgfpathlineto{\pgfqpoint{0.920388in}{0.802059in}}%
\pgfpathlineto{\pgfqpoint{0.929220in}{0.795592in}}%
\pgfpathlineto{\pgfqpoint{0.938010in}{0.789262in}}%
\pgfpathlineto{\pgfqpoint{0.921357in}{0.792865in}}%
\pgfpathlineto{\pgfqpoint{0.904734in}{0.797463in}}%
\pgfpathlineto{\pgfqpoint{0.895891in}{0.803812in}}%
\pgfpathlineto{\pgfqpoint{0.887006in}{0.810297in}}%
\pgfpathclose%
\pgfusepath{fill}%
\end{pgfscope}%
\begin{pgfscope}%
\pgfpathrectangle{\pgfqpoint{0.407592in}{0.222904in}}{\pgfqpoint{1.411163in}{1.411163in}}%
\pgfusepath{clip}%
\pgfsetbuttcap%
\pgfsetroundjoin%
\definecolor{currentfill}{rgb}{0.556863,0.654902,0.847059}%
\pgfsetfillcolor{currentfill}%
\pgfsetfillopacity{0.130000}%
\pgfsetlinewidth{0.000000pt}%
\definecolor{currentstroke}{rgb}{0.000000,0.000000,0.000000}%
\pgfsetstrokecolor{currentstroke}%
\pgfsetdash{}{0pt}%
\pgfpathmoveto{\pgfqpoint{0.717644in}{0.901828in}}%
\pgfpathlineto{\pgfqpoint{0.734251in}{0.891260in}}%
\pgfpathlineto{\pgfqpoint{0.750874in}{0.881356in}}%
\pgfpathlineto{\pgfqpoint{0.760087in}{0.874939in}}%
\pgfpathlineto{\pgfqpoint{0.769257in}{0.868655in}}%
\pgfpathlineto{\pgfqpoint{0.752683in}{0.878873in}}%
\pgfpathlineto{\pgfqpoint{0.736126in}{0.889781in}}%
\pgfpathlineto{\pgfqpoint{0.726907in}{0.895730in}}%
\pgfpathlineto{\pgfqpoint{0.717644in}{0.901828in}}%
\pgfpathclose%
\pgfusepath{fill}%
\end{pgfscope}%
\begin{pgfscope}%
\pgfpathrectangle{\pgfqpoint{0.407592in}{0.222904in}}{\pgfqpoint{1.411163in}{1.411163in}}%
\pgfusepath{clip}%
\pgfsetbuttcap%
\pgfsetroundjoin%
\definecolor{currentfill}{rgb}{0.556863,0.654902,0.847059}%
\pgfsetfillcolor{currentfill}%
\pgfsetfillopacity{0.130000}%
\pgfsetlinewidth{0.000000pt}%
\definecolor{currentstroke}{rgb}{0.000000,0.000000,0.000000}%
\pgfsetstrokecolor{currentstroke}%
\pgfsetdash{}{0pt}%
\pgfpathmoveto{\pgfqpoint{1.376255in}{0.862863in}}%
\pgfpathlineto{\pgfqpoint{1.393771in}{0.874781in}}%
\pgfpathlineto{\pgfqpoint{1.411366in}{0.886876in}}%
\pgfpathlineto{\pgfqpoint{1.419246in}{0.890444in}}%
\pgfpathlineto{\pgfqpoint{1.427107in}{0.894433in}}%
\pgfpathlineto{\pgfqpoint{1.409536in}{0.882232in}}%
\pgfpathlineto{\pgfqpoint{1.392043in}{0.870198in}}%
\pgfpathlineto{\pgfqpoint{1.384159in}{0.866302in}}%
\pgfpathlineto{\pgfqpoint{1.376255in}{0.862863in}}%
\pgfpathclose%
\pgfusepath{fill}%
\end{pgfscope}%
\begin{pgfscope}%
\pgfpathrectangle{\pgfqpoint{0.407592in}{0.222904in}}{\pgfqpoint{1.411163in}{1.411163in}}%
\pgfusepath{clip}%
\pgfsetbuttcap%
\pgfsetroundjoin%
\definecolor{currentfill}{rgb}{0.556863,0.654902,0.847059}%
\pgfsetfillcolor{currentfill}%
\pgfsetfillopacity{0.130000}%
\pgfsetlinewidth{0.000000pt}%
\definecolor{currentstroke}{rgb}{0.000000,0.000000,0.000000}%
\pgfsetstrokecolor{currentstroke}%
\pgfsetdash{}{0pt}%
\pgfpathmoveto{\pgfqpoint{1.544708in}{0.977130in}}%
\pgfpathlineto{\pgfqpoint{1.562610in}{0.989921in}}%
\pgfpathlineto{\pgfqpoint{1.580594in}{1.002830in}}%
\pgfpathlineto{\pgfqpoint{1.588267in}{1.009389in}}%
\pgfpathlineto{\pgfqpoint{1.595927in}{1.016276in}}%
\pgfpathlineto{\pgfqpoint{1.577966in}{1.003486in}}%
\pgfpathlineto{\pgfqpoint{1.560087in}{0.990819in}}%
\pgfpathlineto{\pgfqpoint{1.552405in}{0.983799in}}%
\pgfpathlineto{\pgfqpoint{1.544708in}{0.977130in}}%
\pgfpathclose%
\pgfusepath{fill}%
\end{pgfscope}%
\begin{pgfscope}%
\pgfpathrectangle{\pgfqpoint{0.407592in}{0.222904in}}{\pgfqpoint{1.411163in}{1.411163in}}%
\pgfusepath{clip}%
\pgfsetbuttcap%
\pgfsetroundjoin%
\definecolor{currentfill}{rgb}{0.556863,0.654902,0.847059}%
\pgfsetfillcolor{currentfill}%
\pgfsetfillopacity{0.130000}%
\pgfsetlinewidth{0.000000pt}%
\definecolor{currentstroke}{rgb}{0.000000,0.000000,0.000000}%
\pgfsetstrokecolor{currentstroke}%
\pgfsetdash{}{0pt}%
\pgfpathmoveto{\pgfqpoint{1.662495in}{1.072279in}}%
\pgfpathlineto{\pgfqpoint{1.680652in}{1.085103in}}%
\pgfpathlineto{\pgfqpoint{1.698892in}{1.098045in}}%
\pgfpathlineto{\pgfqpoint{1.706436in}{1.106104in}}%
\pgfpathlineto{\pgfqpoint{1.713967in}{1.114429in}}%
\pgfpathlineto{\pgfqpoint{1.695752in}{1.101674in}}%
\pgfpathlineto{\pgfqpoint{1.677619in}{1.089046in}}%
\pgfpathlineto{\pgfqpoint{1.670063in}{1.080523in}}%
\pgfpathlineto{\pgfqpoint{1.662495in}{1.072279in}}%
\pgfpathclose%
\pgfusepath{fill}%
\end{pgfscope}%
\begin{pgfscope}%
\pgfpathrectangle{\pgfqpoint{0.407592in}{0.222904in}}{\pgfqpoint{1.411163in}{1.411163in}}%
\pgfusepath{clip}%
\pgfsetbuttcap%
\pgfsetroundjoin%
\definecolor{currentfill}{rgb}{0.556863,0.654902,0.847059}%
\pgfsetfillcolor{currentfill}%
\pgfsetfillopacity{0.130000}%
\pgfsetlinewidth{0.000000pt}%
\definecolor{currentstroke}{rgb}{0.000000,0.000000,0.000000}%
\pgfsetstrokecolor{currentstroke}%
\pgfsetdash{}{0pt}%
\pgfpathmoveto{\pgfqpoint{0.971422in}{0.785022in}}%
\pgfpathlineto{\pgfqpoint{0.988187in}{0.784350in}}%
\pgfpathlineto{\pgfqpoint{1.004995in}{0.784606in}}%
\pgfpathlineto{\pgfqpoint{1.013640in}{0.778886in}}%
\pgfpathlineto{\pgfqpoint{1.022246in}{0.773360in}}%
\pgfpathlineto{\pgfqpoint{1.005489in}{0.772772in}}%
\pgfpathlineto{\pgfqpoint{0.988776in}{0.773171in}}%
\pgfpathlineto{\pgfqpoint{0.980119in}{0.779009in}}%
\pgfpathlineto{\pgfqpoint{0.971422in}{0.785022in}}%
\pgfpathclose%
\pgfusepath{fill}%
\end{pgfscope}%
\begin{pgfscope}%
\pgfpathrectangle{\pgfqpoint{0.407592in}{0.222904in}}{\pgfqpoint{1.411163in}{1.411163in}}%
\pgfusepath{clip}%
\pgfsetbuttcap%
\pgfsetroundjoin%
\definecolor{currentfill}{rgb}{0.556863,0.654902,0.847059}%
\pgfsetfillcolor{currentfill}%
\pgfsetfillopacity{0.130000}%
\pgfsetlinewidth{0.000000pt}%
\definecolor{currentstroke}{rgb}{0.000000,0.000000,0.000000}%
\pgfsetstrokecolor{currentstroke}%
\pgfsetdash{}{0pt}%
\pgfpathmoveto{\pgfqpoint{1.713967in}{1.114429in}}%
\pgfpathlineto{\pgfqpoint{1.732264in}{1.127307in}}%
\pgfpathlineto{\pgfqpoint{1.739770in}{1.135793in}}%
\pgfpathlineto{\pgfqpoint{1.747264in}{1.144528in}}%
\pgfpathlineto{\pgfqpoint{1.728992in}{1.131852in}}%
\pgfpathlineto{\pgfqpoint{1.721485in}{1.123013in}}%
\pgfusepath{fill}%
\end{pgfscope}%
\begin{pgfscope}%
\pgfpathrectangle{\pgfqpoint{0.407592in}{0.222904in}}{\pgfqpoint{1.411163in}{1.411163in}}%
\pgfusepath{clip}%
\pgfsetbuttcap%
\pgfsetroundjoin%
\definecolor{currentfill}{rgb}{0.556863,0.654902,0.847059}%
\pgfsetfillcolor{currentfill}%
\pgfsetfillopacity{0.130000}%
\pgfsetlinewidth{0.000000pt}%
\definecolor{currentstroke}{rgb}{0.000000,0.000000,0.000000}%
\pgfsetstrokecolor{currentstroke}%
\pgfsetdash{}{0pt}%
\pgfpathmoveto{\pgfqpoint{1.427107in}{0.894433in}}%
\pgfpathlineto{\pgfqpoint{1.444758in}{0.906787in}}%
\pgfpathlineto{\pgfqpoint{1.462488in}{0.919286in}}%
\pgfpathlineto{\pgfqpoint{1.470307in}{0.923745in}}%
\pgfpathlineto{\pgfqpoint{1.478108in}{0.928598in}}%
\pgfpathlineto{\pgfqpoint{1.460401in}{0.916071in}}%
\pgfpathlineto{\pgfqpoint{1.442773in}{0.903684in}}%
\pgfpathlineto{\pgfqpoint{1.434950in}{0.898846in}}%
\pgfpathlineto{\pgfqpoint{1.427107in}{0.894433in}}%
\pgfpathclose%
\pgfusepath{fill}%
\end{pgfscope}%
\begin{pgfscope}%
\pgfpathrectangle{\pgfqpoint{0.407592in}{0.222904in}}{\pgfqpoint{1.411163in}{1.411163in}}%
\pgfusepath{clip}%
\pgfsetbuttcap%
\pgfsetroundjoin%
\definecolor{currentfill}{rgb}{0.556863,0.654902,0.847059}%
\pgfsetfillcolor{currentfill}%
\pgfsetfillopacity{0.130000}%
\pgfsetlinewidth{0.000000pt}%
\definecolor{currentstroke}{rgb}{0.000000,0.000000,0.000000}%
\pgfsetstrokecolor{currentstroke}%
\pgfsetdash{}{0pt}%
\pgfpathmoveto{\pgfqpoint{0.835734in}{0.835362in}}%
\pgfpathlineto{\pgfqpoint{0.852408in}{0.829074in}}%
\pgfpathlineto{\pgfqpoint{0.869109in}{0.823639in}}%
\pgfpathlineto{\pgfqpoint{0.878078in}{0.816909in}}%
\pgfpathlineto{\pgfqpoint{0.887006in}{0.810297in}}%
\pgfpathlineto{\pgfqpoint{0.870358in}{0.815804in}}%
\pgfpathlineto{\pgfqpoint{0.853737in}{0.822223in}}%
\pgfpathlineto{\pgfqpoint{0.844757in}{0.828730in}}%
\pgfpathlineto{\pgfqpoint{0.835734in}{0.835362in}}%
\pgfpathclose%
\pgfusepath{fill}%
\end{pgfscope}%
\begin{pgfscope}%
\pgfpathrectangle{\pgfqpoint{0.407592in}{0.222904in}}{\pgfqpoint{1.411163in}{1.411163in}}%
\pgfusepath{clip}%
\pgfsetbuttcap%
\pgfsetroundjoin%
\definecolor{currentfill}{rgb}{0.556863,0.654902,0.847059}%
\pgfsetfillcolor{currentfill}%
\pgfsetfillopacity{0.130000}%
\pgfsetlinewidth{0.000000pt}%
\definecolor{currentstroke}{rgb}{0.000000,0.000000,0.000000}%
\pgfsetstrokecolor{currentstroke}%
\pgfsetdash{}{0pt}%
\pgfpathmoveto{\pgfqpoint{1.208059in}{0.792526in}}%
\pgfpathlineto{\pgfqpoint{1.225230in}{0.801284in}}%
\pgfpathlineto{\pgfqpoint{1.242474in}{0.810480in}}%
\pgfpathlineto{\pgfqpoint{1.250630in}{0.809743in}}%
\pgfpathlineto{\pgfqpoint{1.258758in}{0.809424in}}%
\pgfpathlineto{\pgfqpoint{1.241545in}{0.799761in}}%
\pgfpathlineto{\pgfqpoint{1.224405in}{0.790509in}}%
\pgfpathlineto{\pgfqpoint{1.216247in}{0.791301in}}%
\pgfpathlineto{\pgfqpoint{1.208059in}{0.792526in}}%
\pgfpathclose%
\pgfusepath{fill}%
\end{pgfscope}%
\begin{pgfscope}%
\pgfpathrectangle{\pgfqpoint{0.407592in}{0.222904in}}{\pgfqpoint{1.411163in}{1.411163in}}%
\pgfusepath{clip}%
\pgfsetbuttcap%
\pgfsetroundjoin%
\definecolor{currentfill}{rgb}{0.556863,0.654902,0.847059}%
\pgfsetfillcolor{currentfill}%
\pgfsetfillopacity{0.130000}%
\pgfsetlinewidth{0.000000pt}%
\definecolor{currentstroke}{rgb}{0.000000,0.000000,0.000000}%
\pgfsetstrokecolor{currentstroke}%
\pgfsetdash{}{0pt}%
\pgfpathmoveto{\pgfqpoint{1.157384in}{0.781268in}}%
\pgfpathlineto{\pgfqpoint{1.174452in}{0.788442in}}%
\pgfpathlineto{\pgfqpoint{1.191588in}{0.796177in}}%
\pgfpathlineto{\pgfqpoint{1.199839in}{0.794159in}}%
\pgfpathlineto{\pgfqpoint{1.208059in}{0.792526in}}%
\pgfpathlineto{\pgfqpoint{1.190958in}{0.784255in}}%
\pgfpathlineto{\pgfqpoint{1.173926in}{0.776526in}}%
\pgfpathlineto{\pgfqpoint{1.165672in}{0.778704in}}%
\pgfpathlineto{\pgfqpoint{1.157384in}{0.781268in}}%
\pgfpathclose%
\pgfusepath{fill}%
\end{pgfscope}%
\begin{pgfscope}%
\pgfpathrectangle{\pgfqpoint{0.407592in}{0.222904in}}{\pgfqpoint{1.411163in}{1.411163in}}%
\pgfusepath{clip}%
\pgfsetbuttcap%
\pgfsetroundjoin%
\definecolor{currentfill}{rgb}{0.556863,0.654902,0.847059}%
\pgfsetfillcolor{currentfill}%
\pgfsetfillopacity{0.130000}%
\pgfsetlinewidth{0.000000pt}%
\definecolor{currentstroke}{rgb}{0.000000,0.000000,0.000000}%
\pgfsetstrokecolor{currentstroke}%
\pgfsetdash{}{0pt}%
\pgfpathmoveto{\pgfqpoint{1.595927in}{1.016276in}}%
\pgfpathlineto{\pgfqpoint{1.613968in}{1.029183in}}%
\pgfpathlineto{\pgfqpoint{1.632092in}{1.042204in}}%
\pgfpathlineto{\pgfqpoint{1.639713in}{1.049276in}}%
\pgfpathlineto{\pgfqpoint{1.647320in}{1.056650in}}%
\pgfpathlineto{\pgfqpoint{1.629221in}{1.043773in}}%
\pgfpathlineto{\pgfqpoint{1.611203in}{1.031016in}}%
\pgfpathlineto{\pgfqpoint{1.603572in}{1.023487in}}%
\pgfpathlineto{\pgfqpoint{1.595927in}{1.016276in}}%
\pgfpathclose%
\pgfusepath{fill}%
\end{pgfscope}%
\begin{pgfscope}%
\pgfpathrectangle{\pgfqpoint{0.407592in}{0.222904in}}{\pgfqpoint{1.411163in}{1.411163in}}%
\pgfusepath{clip}%
\pgfsetbuttcap%
\pgfsetroundjoin%
\definecolor{currentfill}{rgb}{0.556863,0.654902,0.847059}%
\pgfsetfillcolor{currentfill}%
\pgfsetfillopacity{0.130000}%
\pgfsetlinewidth{0.000000pt}%
\definecolor{currentstroke}{rgb}{0.000000,0.000000,0.000000}%
\pgfsetstrokecolor{currentstroke}%
\pgfsetdash{}{0pt}%
\pgfpathmoveto{\pgfqpoint{0.750874in}{0.881356in}}%
\pgfpathlineto{\pgfqpoint{0.767515in}{0.872148in}}%
\pgfpathlineto{\pgfqpoint{0.784175in}{0.863664in}}%
\pgfpathlineto{\pgfqpoint{0.793337in}{0.856988in}}%
\pgfpathlineto{\pgfqpoint{0.802455in}{0.850431in}}%
\pgfpathlineto{\pgfqpoint{0.785847in}{0.859163in}}%
\pgfpathlineto{\pgfqpoint{0.769257in}{0.868655in}}%
\pgfpathlineto{\pgfqpoint{0.760087in}{0.874939in}}%
\pgfpathlineto{\pgfqpoint{0.750874in}{0.881356in}}%
\pgfpathclose%
\pgfusepath{fill}%
\end{pgfscope}%
\begin{pgfscope}%
\pgfpathrectangle{\pgfqpoint{0.407592in}{0.222904in}}{\pgfqpoint{1.411163in}{1.411163in}}%
\pgfusepath{clip}%
\pgfsetbuttcap%
\pgfsetroundjoin%
\definecolor{currentfill}{rgb}{0.556863,0.654902,0.847059}%
\pgfsetfillcolor{currentfill}%
\pgfsetfillopacity{0.130000}%
\pgfsetlinewidth{0.000000pt}%
\definecolor{currentstroke}{rgb}{0.000000,0.000000,0.000000}%
\pgfsetstrokecolor{currentstroke}%
\pgfsetdash{}{0pt}%
\pgfpathmoveto{\pgfqpoint{1.258758in}{0.809424in}}%
\pgfpathlineto{\pgfqpoint{1.276044in}{0.819455in}}%
\pgfpathlineto{\pgfqpoint{1.293406in}{0.829819in}}%
\pgfpathlineto{\pgfqpoint{1.301477in}{0.830332in}}%
\pgfpathlineto{\pgfqpoint{1.309521in}{0.831275in}}%
\pgfpathlineto{\pgfqpoint{1.292188in}{0.820541in}}%
\pgfpathlineto{\pgfqpoint{1.274929in}{0.810113in}}%
\pgfpathlineto{\pgfqpoint{1.266857in}{0.809540in}}%
\pgfpathlineto{\pgfqpoint{1.258758in}{0.809424in}}%
\pgfpathclose%
\pgfusepath{fill}%
\end{pgfscope}%
\begin{pgfscope}%
\pgfpathrectangle{\pgfqpoint{0.407592in}{0.222904in}}{\pgfqpoint{1.411163in}{1.411163in}}%
\pgfusepath{clip}%
\pgfsetbuttcap%
\pgfsetroundjoin%
\definecolor{currentfill}{rgb}{0.556863,0.654902,0.847059}%
\pgfsetfillcolor{currentfill}%
\pgfsetfillopacity{0.130000}%
\pgfsetlinewidth{0.000000pt}%
\definecolor{currentstroke}{rgb}{0.000000,0.000000,0.000000}%
\pgfsetstrokecolor{currentstroke}%
\pgfsetdash{}{0pt}%
\pgfpathmoveto{\pgfqpoint{1.106684in}{0.776140in}}%
\pgfpathlineto{\pgfqpoint{1.123663in}{0.781453in}}%
\pgfpathlineto{\pgfqpoint{1.140706in}{0.787452in}}%
\pgfpathlineto{\pgfqpoint{1.149062in}{0.784192in}}%
\pgfpathlineto{\pgfqpoint{1.157384in}{0.781268in}}%
\pgfpathlineto{\pgfqpoint{1.140382in}{0.774713in}}%
\pgfpathlineto{\pgfqpoint{1.123443in}{0.768841in}}%
\pgfpathlineto{\pgfqpoint{1.115081in}{0.772329in}}%
\pgfpathlineto{\pgfqpoint{1.106684in}{0.776140in}}%
\pgfpathclose%
\pgfusepath{fill}%
\end{pgfscope}%
\begin{pgfscope}%
\pgfpathrectangle{\pgfqpoint{0.407592in}{0.222904in}}{\pgfqpoint{1.411163in}{1.411163in}}%
\pgfusepath{clip}%
\pgfsetbuttcap%
\pgfsetroundjoin%
\definecolor{currentfill}{rgb}{0.556863,0.654902,0.847059}%
\pgfsetfillcolor{currentfill}%
\pgfsetfillopacity{0.130000}%
\pgfsetlinewidth{0.000000pt}%
\definecolor{currentstroke}{rgb}{0.000000,0.000000,0.000000}%
\pgfsetstrokecolor{currentstroke}%
\pgfsetdash{}{0pt}%
\pgfpathmoveto{\pgfqpoint{1.478108in}{0.928598in}}%
\pgfpathlineto{\pgfqpoint{1.495895in}{0.941256in}}%
\pgfpathlineto{\pgfqpoint{1.513763in}{0.954039in}}%
\pgfpathlineto{\pgfqpoint{1.521524in}{0.959267in}}%
\pgfpathlineto{\pgfqpoint{1.529268in}{0.964860in}}%
\pgfpathlineto{\pgfqpoint{1.511423in}{0.952108in}}%
\pgfpathlineto{\pgfqpoint{1.493658in}{0.939481in}}%
\pgfpathlineto{\pgfqpoint{1.485892in}{0.933844in}}%
\pgfpathlineto{\pgfqpoint{1.478108in}{0.928598in}}%
\pgfpathclose%
\pgfusepath{fill}%
\end{pgfscope}%
\begin{pgfscope}%
\pgfpathrectangle{\pgfqpoint{0.407592in}{0.222904in}}{\pgfqpoint{1.411163in}{1.411163in}}%
\pgfusepath{clip}%
\pgfsetbuttcap%
\pgfsetroundjoin%
\definecolor{currentfill}{rgb}{0.556863,0.654902,0.847059}%
\pgfsetfillcolor{currentfill}%
\pgfsetfillopacity{0.130000}%
\pgfsetlinewidth{0.000000pt}%
\definecolor{currentstroke}{rgb}{0.000000,0.000000,0.000000}%
\pgfsetstrokecolor{currentstroke}%
\pgfsetdash{}{0pt}%
\pgfpathmoveto{\pgfqpoint{0.920388in}{0.802059in}}%
\pgfpathlineto{\pgfqpoint{0.937129in}{0.799331in}}%
\pgfpathlineto{\pgfqpoint{0.953907in}{0.797518in}}%
\pgfpathlineto{\pgfqpoint{0.962685in}{0.791196in}}%
\pgfpathlineto{\pgfqpoint{0.971422in}{0.785022in}}%
\pgfpathlineto{\pgfqpoint{0.954698in}{0.786652in}}%
\pgfpathlineto{\pgfqpoint{0.938010in}{0.789262in}}%
\pgfpathlineto{\pgfqpoint{0.929220in}{0.795592in}}%
\pgfpathlineto{\pgfqpoint{0.920388in}{0.802059in}}%
\pgfpathclose%
\pgfusepath{fill}%
\end{pgfscope}%
\begin{pgfscope}%
\pgfpathrectangle{\pgfqpoint{0.407592in}{0.222904in}}{\pgfqpoint{1.411163in}{1.411163in}}%
\pgfusepath{clip}%
\pgfsetbuttcap%
\pgfsetroundjoin%
\definecolor{currentfill}{rgb}{0.556863,0.654902,0.847059}%
\pgfsetfillcolor{currentfill}%
\pgfsetfillopacity{0.130000}%
\pgfsetlinewidth{0.000000pt}%
\definecolor{currentstroke}{rgb}{0.000000,0.000000,0.000000}%
\pgfsetstrokecolor{currentstroke}%
\pgfsetdash{}{0pt}%
\pgfpathmoveto{\pgfqpoint{1.309521in}{0.831275in}}%
\pgfpathlineto{\pgfqpoint{1.326932in}{0.842284in}}%
\pgfpathlineto{\pgfqpoint{1.344419in}{0.853544in}}%
\pgfpathlineto{\pgfqpoint{1.352413in}{0.855223in}}%
\pgfpathlineto{\pgfqpoint{1.360383in}{0.857330in}}%
\pgfpathlineto{\pgfqpoint{1.342921in}{0.845805in}}%
\pgfpathlineto{\pgfqpoint{1.325536in}{0.834508in}}%
\pgfpathlineto{\pgfqpoint{1.317541in}{0.832663in}}%
\pgfpathlineto{\pgfqpoint{1.309521in}{0.831275in}}%
\pgfpathclose%
\pgfusepath{fill}%
\end{pgfscope}%
\begin{pgfscope}%
\pgfpathrectangle{\pgfqpoint{0.407592in}{0.222904in}}{\pgfqpoint{1.411163in}{1.411163in}}%
\pgfusepath{clip}%
\pgfsetbuttcap%
\pgfsetroundjoin%
\definecolor{currentfill}{rgb}{0.556863,0.654902,0.847059}%
\pgfsetfillcolor{currentfill}%
\pgfsetfillopacity{0.130000}%
\pgfsetlinewidth{0.000000pt}%
\definecolor{currentstroke}{rgb}{0.000000,0.000000,0.000000}%
\pgfsetstrokecolor{currentstroke}%
\pgfsetdash{}{0pt}%
\pgfpathmoveto{\pgfqpoint{1.055905in}{0.777314in}}%
\pgfpathlineto{\pgfqpoint{1.072813in}{0.780580in}}%
\pgfpathlineto{\pgfqpoint{1.089778in}{0.784632in}}%
\pgfpathlineto{\pgfqpoint{1.098249in}{0.780248in}}%
\pgfpathlineto{\pgfqpoint{1.106684in}{0.776140in}}%
\pgfpathlineto{\pgfqpoint{1.089764in}{0.771570in}}%
\pgfpathlineto{\pgfqpoint{1.072902in}{0.767809in}}%
\pgfpathlineto{\pgfqpoint{1.064423in}{0.772434in}}%
\pgfpathlineto{\pgfqpoint{1.055905in}{0.777314in}}%
\pgfpathclose%
\pgfusepath{fill}%
\end{pgfscope}%
\begin{pgfscope}%
\pgfpathrectangle{\pgfqpoint{0.407592in}{0.222904in}}{\pgfqpoint{1.411163in}{1.411163in}}%
\pgfusepath{clip}%
\pgfsetbuttcap%
\pgfsetroundjoin%
\definecolor{currentfill}{rgb}{0.556863,0.654902,0.847059}%
\pgfsetfillcolor{currentfill}%
\pgfsetfillopacity{0.130000}%
\pgfsetlinewidth{0.000000pt}%
\definecolor{currentstroke}{rgb}{0.000000,0.000000,0.000000}%
\pgfsetstrokecolor{currentstroke}%
\pgfsetdash{}{0pt}%
\pgfpathmoveto{\pgfqpoint{1.647320in}{1.056650in}}%
\pgfpathlineto{\pgfqpoint{1.665502in}{1.069642in}}%
\pgfpathlineto{\pgfqpoint{1.683766in}{1.082745in}}%
\pgfpathlineto{\pgfqpoint{1.691336in}{1.090258in}}%
\pgfpathlineto{\pgfqpoint{1.698892in}{1.098045in}}%
\pgfpathlineto{\pgfqpoint{1.680652in}{1.085103in}}%
\pgfpathlineto{\pgfqpoint{1.662495in}{1.072279in}}%
\pgfpathlineto{\pgfqpoint{1.654914in}{1.064319in}}%
\pgfpathlineto{\pgfqpoint{1.647320in}{1.056650in}}%
\pgfpathclose%
\pgfusepath{fill}%
\end{pgfscope}%
\begin{pgfscope}%
\pgfpathrectangle{\pgfqpoint{0.407592in}{0.222904in}}{\pgfqpoint{1.411163in}{1.411163in}}%
\pgfusepath{clip}%
\pgfsetbuttcap%
\pgfsetroundjoin%
\definecolor{currentfill}{rgb}{0.556863,0.654902,0.847059}%
\pgfsetfillcolor{currentfill}%
\pgfsetfillopacity{0.130000}%
\pgfsetlinewidth{0.000000pt}%
\definecolor{currentstroke}{rgb}{0.000000,0.000000,0.000000}%
\pgfsetstrokecolor{currentstroke}%
\pgfsetdash{}{0pt}%
\pgfpathmoveto{\pgfqpoint{1.360383in}{0.857330in}}%
\pgfpathlineto{\pgfqpoint{1.377923in}{0.869064in}}%
\pgfpathlineto{\pgfqpoint{1.395541in}{0.880990in}}%
\pgfpathlineto{\pgfqpoint{1.403464in}{0.883727in}}%
\pgfpathlineto{\pgfqpoint{1.411366in}{0.886876in}}%
\pgfpathlineto{\pgfqpoint{1.393771in}{0.874781in}}%
\pgfpathlineto{\pgfqpoint{1.376255in}{0.862863in}}%
\pgfpathlineto{\pgfqpoint{1.368330in}{0.859874in}}%
\pgfpathlineto{\pgfqpoint{1.360383in}{0.857330in}}%
\pgfpathclose%
\pgfusepath{fill}%
\end{pgfscope}%
\begin{pgfscope}%
\pgfpathrectangle{\pgfqpoint{0.407592in}{0.222904in}}{\pgfqpoint{1.411163in}{1.411163in}}%
\pgfusepath{clip}%
\pgfsetbuttcap%
\pgfsetroundjoin%
\definecolor{currentfill}{rgb}{0.556863,0.654902,0.847059}%
\pgfsetfillcolor{currentfill}%
\pgfsetfillopacity{0.130000}%
\pgfsetlinewidth{0.000000pt}%
\definecolor{currentstroke}{rgb}{0.000000,0.000000,0.000000}%
\pgfsetstrokecolor{currentstroke}%
\pgfsetdash{}{0pt}%
\pgfpathmoveto{\pgfqpoint{1.529268in}{0.964860in}}%
\pgfpathlineto{\pgfqpoint{1.547193in}{0.977731in}}%
\pgfpathlineto{\pgfqpoint{1.565200in}{0.990715in}}%
\pgfpathlineto{\pgfqpoint{1.572905in}{0.996605in}}%
\pgfpathlineto{\pgfqpoint{1.580594in}{1.002830in}}%
\pgfpathlineto{\pgfqpoint{1.562610in}{0.989921in}}%
\pgfpathlineto{\pgfqpoint{1.544708in}{0.977130in}}%
\pgfpathlineto{\pgfqpoint{1.536995in}{0.970815in}}%
\pgfpathlineto{\pgfqpoint{1.529268in}{0.964860in}}%
\pgfpathclose%
\pgfusepath{fill}%
\end{pgfscope}%
\begin{pgfscope}%
\pgfpathrectangle{\pgfqpoint{0.407592in}{0.222904in}}{\pgfqpoint{1.411163in}{1.411163in}}%
\pgfusepath{clip}%
\pgfsetbuttcap%
\pgfsetroundjoin%
\definecolor{currentfill}{rgb}{0.556863,0.654902,0.847059}%
\pgfsetfillcolor{currentfill}%
\pgfsetfillopacity{0.130000}%
\pgfsetlinewidth{0.000000pt}%
\definecolor{currentstroke}{rgb}{0.000000,0.000000,0.000000}%
\pgfsetstrokecolor{currentstroke}%
\pgfsetdash{}{0pt}%
\pgfpathmoveto{\pgfqpoint{1.698892in}{1.098045in}}%
\pgfpathlineto{\pgfqpoint{1.717215in}{1.111103in}}%
\pgfpathlineto{\pgfqpoint{1.724746in}{1.119075in}}%
\pgfpathlineto{\pgfqpoint{1.732264in}{1.127307in}}%
\pgfpathlineto{\pgfqpoint{1.713967in}{1.114429in}}%
\pgfpathlineto{\pgfqpoint{1.706436in}{1.106104in}}%
\pgfusepath{fill}%
\end{pgfscope}%
\begin{pgfscope}%
\pgfpathrectangle{\pgfqpoint{0.407592in}{0.222904in}}{\pgfqpoint{1.411163in}{1.411163in}}%
\pgfusepath{clip}%
\pgfsetbuttcap%
\pgfsetroundjoin%
\definecolor{currentfill}{rgb}{0.556863,0.654902,0.847059}%
\pgfsetfillcolor{currentfill}%
\pgfsetfillopacity{0.130000}%
\pgfsetlinewidth{0.000000pt}%
\definecolor{currentstroke}{rgb}{0.000000,0.000000,0.000000}%
\pgfsetstrokecolor{currentstroke}%
\pgfsetdash{}{0pt}%
\pgfpathmoveto{\pgfqpoint{1.004995in}{0.784606in}}%
\pgfpathlineto{\pgfqpoint{1.021850in}{0.785755in}}%
\pgfpathlineto{\pgfqpoint{1.038754in}{0.787757in}}%
\pgfpathlineto{\pgfqpoint{1.047349in}{0.782428in}}%
\pgfpathlineto{\pgfqpoint{1.055905in}{0.777314in}}%
\pgfpathlineto{\pgfqpoint{1.039050in}{0.774890in}}%
\pgfpathlineto{\pgfqpoint{1.022246in}{0.773360in}}%
\pgfpathlineto{\pgfqpoint{1.013640in}{0.778886in}}%
\pgfpathlineto{\pgfqpoint{1.004995in}{0.784606in}}%
\pgfpathclose%
\pgfusepath{fill}%
\end{pgfscope}%
\begin{pgfscope}%
\pgfpathrectangle{\pgfqpoint{0.407592in}{0.222904in}}{\pgfqpoint{1.411163in}{1.411163in}}%
\pgfusepath{clip}%
\pgfsetbuttcap%
\pgfsetroundjoin%
\definecolor{currentfill}{rgb}{0.556863,0.654902,0.847059}%
\pgfsetfillcolor{currentfill}%
\pgfsetfillopacity{0.130000}%
\pgfsetlinewidth{0.000000pt}%
\definecolor{currentstroke}{rgb}{0.000000,0.000000,0.000000}%
\pgfsetstrokecolor{currentstroke}%
\pgfsetdash{}{0pt}%
\pgfpathmoveto{\pgfqpoint{0.784175in}{0.863664in}}%
\pgfpathlineto{\pgfqpoint{0.800856in}{0.855929in}}%
\pgfpathlineto{\pgfqpoint{0.817559in}{0.848965in}}%
\pgfpathlineto{\pgfqpoint{0.826668in}{0.842110in}}%
\pgfpathlineto{\pgfqpoint{0.835734in}{0.835362in}}%
\pgfpathlineto{\pgfqpoint{0.819083in}{0.842488in}}%
\pgfpathlineto{\pgfqpoint{0.802455in}{0.850431in}}%
\pgfpathlineto{\pgfqpoint{0.793337in}{0.856988in}}%
\pgfpathlineto{\pgfqpoint{0.784175in}{0.863664in}}%
\pgfpathclose%
\pgfusepath{fill}%
\end{pgfscope}%
\begin{pgfscope}%
\pgfpathrectangle{\pgfqpoint{0.407592in}{0.222904in}}{\pgfqpoint{1.411163in}{1.411163in}}%
\pgfusepath{clip}%
\pgfsetbuttcap%
\pgfsetroundjoin%
\definecolor{currentfill}{rgb}{0.556863,0.654902,0.847059}%
\pgfsetfillcolor{currentfill}%
\pgfsetfillopacity{0.130000}%
\pgfsetlinewidth{0.000000pt}%
\definecolor{currentstroke}{rgb}{0.000000,0.000000,0.000000}%
\pgfsetstrokecolor{currentstroke}%
\pgfsetdash{}{0pt}%
\pgfpathmoveto{\pgfqpoint{0.869109in}{0.823639in}}%
\pgfpathlineto{\pgfqpoint{0.885838in}{0.819069in}}%
\pgfpathlineto{\pgfqpoint{0.902599in}{0.815369in}}%
\pgfpathlineto{\pgfqpoint{0.911515in}{0.808655in}}%
\pgfpathlineto{\pgfqpoint{0.920388in}{0.802059in}}%
\pgfpathlineto{\pgfqpoint{0.903681in}{0.805714in}}%
\pgfpathlineto{\pgfqpoint{0.887006in}{0.810297in}}%
\pgfpathlineto{\pgfqpoint{0.878078in}{0.816909in}}%
\pgfpathlineto{\pgfqpoint{0.869109in}{0.823639in}}%
\pgfpathclose%
\pgfusepath{fill}%
\end{pgfscope}%
\begin{pgfscope}%
\pgfpathrectangle{\pgfqpoint{0.407592in}{0.222904in}}{\pgfqpoint{1.411163in}{1.411163in}}%
\pgfusepath{clip}%
\pgfsetbuttcap%
\pgfsetroundjoin%
\definecolor{currentfill}{rgb}{0.556863,0.654902,0.847059}%
\pgfsetfillcolor{currentfill}%
\pgfsetfillopacity{0.130000}%
\pgfsetlinewidth{0.000000pt}%
\definecolor{currentstroke}{rgb}{0.000000,0.000000,0.000000}%
\pgfsetstrokecolor{currentstroke}%
\pgfsetdash{}{0pt}%
\pgfpathmoveto{\pgfqpoint{1.411366in}{0.886876in}}%
\pgfpathlineto{\pgfqpoint{1.429039in}{0.899135in}}%
\pgfpathlineto{\pgfqpoint{1.446793in}{0.911546in}}%
\pgfpathlineto{\pgfqpoint{1.454650in}{0.915220in}}%
\pgfpathlineto{\pgfqpoint{1.462488in}{0.919286in}}%
\pgfpathlineto{\pgfqpoint{1.444758in}{0.906787in}}%
\pgfpathlineto{\pgfqpoint{1.427107in}{0.894433in}}%
\pgfpathlineto{\pgfqpoint{1.419246in}{0.890444in}}%
\pgfpathlineto{\pgfqpoint{1.411366in}{0.886876in}}%
\pgfpathclose%
\pgfusepath{fill}%
\end{pgfscope}%
\begin{pgfscope}%
\pgfpathrectangle{\pgfqpoint{0.407592in}{0.222904in}}{\pgfqpoint{1.411163in}{1.411163in}}%
\pgfusepath{clip}%
\pgfsetbuttcap%
\pgfsetroundjoin%
\definecolor{currentfill}{rgb}{0.556863,0.654902,0.847059}%
\pgfsetfillcolor{currentfill}%
\pgfsetfillopacity{0.130000}%
\pgfsetlinewidth{0.000000pt}%
\definecolor{currentstroke}{rgb}{0.000000,0.000000,0.000000}%
\pgfsetstrokecolor{currentstroke}%
\pgfsetdash{}{0pt}%
\pgfpathmoveto{\pgfqpoint{1.580594in}{1.002830in}}%
\pgfpathlineto{\pgfqpoint{1.598659in}{1.015851in}}%
\pgfpathlineto{\pgfqpoint{1.616807in}{1.028979in}}%
\pgfpathlineto{\pgfqpoint{1.624456in}{1.035437in}}%
\pgfpathlineto{\pgfqpoint{1.632092in}{1.042204in}}%
\pgfpathlineto{\pgfqpoint{1.613968in}{1.029183in}}%
\pgfpathlineto{\pgfqpoint{1.595927in}{1.016276in}}%
\pgfpathlineto{\pgfqpoint{1.588267in}{1.009389in}}%
\pgfpathlineto{\pgfqpoint{1.580594in}{1.002830in}}%
\pgfpathclose%
\pgfusepath{fill}%
\end{pgfscope}%
\begin{pgfscope}%
\pgfpathrectangle{\pgfqpoint{0.407592in}{0.222904in}}{\pgfqpoint{1.411163in}{1.411163in}}%
\pgfusepath{clip}%
\pgfsetbuttcap%
\pgfsetroundjoin%
\definecolor{currentfill}{rgb}{0.556863,0.654902,0.847059}%
\pgfsetfillcolor{currentfill}%
\pgfsetfillopacity{0.130000}%
\pgfsetlinewidth{0.000000pt}%
\definecolor{currentstroke}{rgb}{0.000000,0.000000,0.000000}%
\pgfsetstrokecolor{currentstroke}%
\pgfsetdash{}{0pt}%
\pgfpathmoveto{\pgfqpoint{0.953907in}{0.797518in}}%
\pgfpathlineto{\pgfqpoint{0.970724in}{0.796604in}}%
\pgfpathlineto{\pgfqpoint{0.987584in}{0.796565in}}%
\pgfpathlineto{\pgfqpoint{0.996310in}{0.790503in}}%
\pgfpathlineto{\pgfqpoint{1.004995in}{0.784606in}}%
\pgfpathlineto{\pgfqpoint{0.988187in}{0.784350in}}%
\pgfpathlineto{\pgfqpoint{0.971422in}{0.785022in}}%
\pgfpathlineto{\pgfqpoint{0.962685in}{0.791196in}}%
\pgfpathlineto{\pgfqpoint{0.953907in}{0.797518in}}%
\pgfpathclose%
\pgfusepath{fill}%
\end{pgfscope}%
\begin{pgfscope}%
\pgfpathrectangle{\pgfqpoint{0.407592in}{0.222904in}}{\pgfqpoint{1.411163in}{1.411163in}}%
\pgfusepath{clip}%
\pgfsetbuttcap%
\pgfsetroundjoin%
\definecolor{currentfill}{rgb}{0.556863,0.654902,0.847059}%
\pgfsetfillcolor{currentfill}%
\pgfsetfillopacity{0.130000}%
\pgfsetlinewidth{0.000000pt}%
\definecolor{currentstroke}{rgb}{0.000000,0.000000,0.000000}%
\pgfsetstrokecolor{currentstroke}%
\pgfsetdash{}{0pt}%
\pgfpathmoveto{\pgfqpoint{1.462488in}{0.919286in}}%
\pgfpathlineto{\pgfqpoint{1.480299in}{0.931920in}}%
\pgfpathlineto{\pgfqpoint{1.498190in}{0.944682in}}%
\pgfpathlineto{\pgfqpoint{1.505986in}{0.949177in}}%
\pgfpathlineto{\pgfqpoint{1.513763in}{0.954039in}}%
\pgfpathlineto{\pgfqpoint{1.495895in}{0.941256in}}%
\pgfpathlineto{\pgfqpoint{1.478108in}{0.928598in}}%
\pgfpathlineto{\pgfqpoint{1.470307in}{0.923745in}}%
\pgfpathlineto{\pgfqpoint{1.462488in}{0.919286in}}%
\pgfpathclose%
\pgfusepath{fill}%
\end{pgfscope}%
\begin{pgfscope}%
\pgfpathrectangle{\pgfqpoint{0.407592in}{0.222904in}}{\pgfqpoint{1.411163in}{1.411163in}}%
\pgfusepath{clip}%
\pgfsetbuttcap%
\pgfsetroundjoin%
\definecolor{currentfill}{rgb}{0.556863,0.654902,0.847059}%
\pgfsetfillcolor{currentfill}%
\pgfsetfillopacity{0.130000}%
\pgfsetlinewidth{0.000000pt}%
\definecolor{currentstroke}{rgb}{0.000000,0.000000,0.000000}%
\pgfsetstrokecolor{currentstroke}%
\pgfsetdash{}{0pt}%
\pgfpathmoveto{\pgfqpoint{1.191588in}{0.796177in}}%
\pgfpathlineto{\pgfqpoint{1.208794in}{0.804421in}}%
\pgfpathlineto{\pgfqpoint{1.226070in}{0.813128in}}%
\pgfpathlineto{\pgfqpoint{1.234287in}{0.811615in}}%
\pgfpathlineto{\pgfqpoint{1.242474in}{0.810480in}}%
\pgfpathlineto{\pgfqpoint{1.225230in}{0.801284in}}%
\pgfpathlineto{\pgfqpoint{1.208059in}{0.792526in}}%
\pgfpathlineto{\pgfqpoint{1.199839in}{0.794159in}}%
\pgfpathlineto{\pgfqpoint{1.191588in}{0.796177in}}%
\pgfpathclose%
\pgfusepath{fill}%
\end{pgfscope}%
\begin{pgfscope}%
\pgfpathrectangle{\pgfqpoint{0.407592in}{0.222904in}}{\pgfqpoint{1.411163in}{1.411163in}}%
\pgfusepath{clip}%
\pgfsetbuttcap%
\pgfsetroundjoin%
\definecolor{currentfill}{rgb}{0.556863,0.654902,0.847059}%
\pgfsetfillcolor{currentfill}%
\pgfsetfillopacity{0.130000}%
\pgfsetlinewidth{0.000000pt}%
\definecolor{currentstroke}{rgb}{0.000000,0.000000,0.000000}%
\pgfsetstrokecolor{currentstroke}%
\pgfsetdash{}{0pt}%
\pgfpathmoveto{\pgfqpoint{1.140706in}{0.787452in}}%
\pgfpathlineto{\pgfqpoint{1.157813in}{0.794077in}}%
\pgfpathlineto{\pgfqpoint{1.174986in}{0.801278in}}%
\pgfpathlineto{\pgfqpoint{1.183304in}{0.798557in}}%
\pgfpathlineto{\pgfqpoint{1.191588in}{0.796177in}}%
\pgfpathlineto{\pgfqpoint{1.174452in}{0.788442in}}%
\pgfpathlineto{\pgfqpoint{1.157384in}{0.781268in}}%
\pgfpathlineto{\pgfqpoint{1.149062in}{0.784192in}}%
\pgfpathlineto{\pgfqpoint{1.140706in}{0.787452in}}%
\pgfpathclose%
\pgfusepath{fill}%
\end{pgfscope}%
\begin{pgfscope}%
\pgfpathrectangle{\pgfqpoint{0.407592in}{0.222904in}}{\pgfqpoint{1.411163in}{1.411163in}}%
\pgfusepath{clip}%
\pgfsetbuttcap%
\pgfsetroundjoin%
\definecolor{currentfill}{rgb}{0.556863,0.654902,0.847059}%
\pgfsetfillcolor{currentfill}%
\pgfsetfillopacity{0.130000}%
\pgfsetlinewidth{0.000000pt}%
\definecolor{currentstroke}{rgb}{0.000000,0.000000,0.000000}%
\pgfsetstrokecolor{currentstroke}%
\pgfsetdash{}{0pt}%
\pgfpathmoveto{\pgfqpoint{1.242474in}{0.810480in}}%
\pgfpathlineto{\pgfqpoint{1.259790in}{0.820073in}}%
\pgfpathlineto{\pgfqpoint{1.277182in}{0.830025in}}%
\pgfpathlineto{\pgfqpoint{1.285308in}{0.829722in}}%
\pgfpathlineto{\pgfqpoint{1.293406in}{0.829819in}}%
\pgfpathlineto{\pgfqpoint{1.276044in}{0.819455in}}%
\pgfpathlineto{\pgfqpoint{1.258758in}{0.809424in}}%
\pgfpathlineto{\pgfqpoint{1.250630in}{0.809743in}}%
\pgfpathlineto{\pgfqpoint{1.242474in}{0.810480in}}%
\pgfpathclose%
\pgfusepath{fill}%
\end{pgfscope}%
\begin{pgfscope}%
\pgfpathrectangle{\pgfqpoint{0.407592in}{0.222904in}}{\pgfqpoint{1.411163in}{1.411163in}}%
\pgfusepath{clip}%
\pgfsetbuttcap%
\pgfsetroundjoin%
\definecolor{currentfill}{rgb}{0.556863,0.654902,0.847059}%
\pgfsetfillcolor{currentfill}%
\pgfsetfillopacity{0.130000}%
\pgfsetlinewidth{0.000000pt}%
\definecolor{currentstroke}{rgb}{0.000000,0.000000,0.000000}%
\pgfsetstrokecolor{currentstroke}%
\pgfsetdash{}{0pt}%
\pgfpathmoveto{\pgfqpoint{1.632092in}{1.042204in}}%
\pgfpathlineto{\pgfqpoint{1.650298in}{1.055333in}}%
\pgfpathlineto{\pgfqpoint{1.668587in}{1.068568in}}%
\pgfpathlineto{\pgfqpoint{1.676183in}{1.075514in}}%
\pgfpathlineto{\pgfqpoint{1.683766in}{1.082745in}}%
\pgfpathlineto{\pgfqpoint{1.665502in}{1.069642in}}%
\pgfpathlineto{\pgfqpoint{1.647320in}{1.056650in}}%
\pgfpathlineto{\pgfqpoint{1.639713in}{1.049276in}}%
\pgfpathlineto{\pgfqpoint{1.632092in}{1.042204in}}%
\pgfpathclose%
\pgfusepath{fill}%
\end{pgfscope}%
\begin{pgfscope}%
\pgfpathrectangle{\pgfqpoint{0.407592in}{0.222904in}}{\pgfqpoint{1.411163in}{1.411163in}}%
\pgfusepath{clip}%
\pgfsetbuttcap%
\pgfsetroundjoin%
\definecolor{currentfill}{rgb}{0.556863,0.654902,0.847059}%
\pgfsetfillcolor{currentfill}%
\pgfsetfillopacity{0.130000}%
\pgfsetlinewidth{0.000000pt}%
\definecolor{currentstroke}{rgb}{0.000000,0.000000,0.000000}%
\pgfsetstrokecolor{currentstroke}%
\pgfsetdash{}{0pt}%
\pgfpathmoveto{\pgfqpoint{0.817559in}{0.848965in}}%
\pgfpathlineto{\pgfqpoint{0.834287in}{0.842792in}}%
\pgfpathlineto{\pgfqpoint{0.851042in}{0.837422in}}%
\pgfpathlineto{\pgfqpoint{0.860097in}{0.830479in}}%
\pgfpathlineto{\pgfqpoint{0.869109in}{0.823639in}}%
\pgfpathlineto{\pgfqpoint{0.852408in}{0.829074in}}%
\pgfpathlineto{\pgfqpoint{0.835734in}{0.835362in}}%
\pgfpathlineto{\pgfqpoint{0.826668in}{0.842110in}}%
\pgfpathlineto{\pgfqpoint{0.817559in}{0.848965in}}%
\pgfpathclose%
\pgfusepath{fill}%
\end{pgfscope}%
\begin{pgfscope}%
\pgfpathrectangle{\pgfqpoint{0.407592in}{0.222904in}}{\pgfqpoint{1.411163in}{1.411163in}}%
\pgfusepath{clip}%
\pgfsetbuttcap%
\pgfsetroundjoin%
\definecolor{currentfill}{rgb}{0.556863,0.654902,0.847059}%
\pgfsetfillcolor{currentfill}%
\pgfsetfillopacity{0.130000}%
\pgfsetlinewidth{0.000000pt}%
\definecolor{currentstroke}{rgb}{0.000000,0.000000,0.000000}%
\pgfsetstrokecolor{currentstroke}%
\pgfsetdash{}{0pt}%
\pgfpathmoveto{\pgfqpoint{1.293406in}{0.829819in}}%
\pgfpathlineto{\pgfqpoint{1.310843in}{0.840485in}}%
\pgfpathlineto{\pgfqpoint{1.328358in}{0.851425in}}%
\pgfpathlineto{\pgfqpoint{1.336401in}{0.852281in}}%
\pgfpathlineto{\pgfqpoint{1.344419in}{0.853544in}}%
\pgfpathlineto{\pgfqpoint{1.326932in}{0.842284in}}%
\pgfpathlineto{\pgfqpoint{1.309521in}{0.831275in}}%
\pgfpathlineto{\pgfqpoint{1.301477in}{0.830332in}}%
\pgfpathlineto{\pgfqpoint{1.293406in}{0.829819in}}%
\pgfpathclose%
\pgfusepath{fill}%
\end{pgfscope}%
\begin{pgfscope}%
\pgfpathrectangle{\pgfqpoint{0.407592in}{0.222904in}}{\pgfqpoint{1.411163in}{1.411163in}}%
\pgfusepath{clip}%
\pgfsetbuttcap%
\pgfsetroundjoin%
\definecolor{currentfill}{rgb}{0.556863,0.654902,0.847059}%
\pgfsetfillcolor{currentfill}%
\pgfsetfillopacity{0.130000}%
\pgfsetlinewidth{0.000000pt}%
\definecolor{currentstroke}{rgb}{0.000000,0.000000,0.000000}%
\pgfsetstrokecolor{currentstroke}%
\pgfsetdash{}{0pt}%
\pgfpathmoveto{\pgfqpoint{1.089778in}{0.784632in}}%
\pgfpathlineto{\pgfqpoint{1.106801in}{0.789419in}}%
\pgfpathlineto{\pgfqpoint{1.123886in}{0.794886in}}%
\pgfpathlineto{\pgfqpoint{1.132314in}{0.791023in}}%
\pgfpathlineto{\pgfqpoint{1.140706in}{0.787452in}}%
\pgfpathlineto{\pgfqpoint{1.123663in}{0.781453in}}%
\pgfpathlineto{\pgfqpoint{1.106684in}{0.776140in}}%
\pgfpathlineto{\pgfqpoint{1.098249in}{0.780248in}}%
\pgfpathlineto{\pgfqpoint{1.089778in}{0.784632in}}%
\pgfpathclose%
\pgfusepath{fill}%
\end{pgfscope}%
\begin{pgfscope}%
\pgfpathrectangle{\pgfqpoint{0.407592in}{0.222904in}}{\pgfqpoint{1.411163in}{1.411163in}}%
\pgfusepath{clip}%
\pgfsetbuttcap%
\pgfsetroundjoin%
\definecolor{currentfill}{rgb}{0.556863,0.654902,0.847059}%
\pgfsetfillcolor{currentfill}%
\pgfsetfillopacity{0.130000}%
\pgfsetlinewidth{0.000000pt}%
\definecolor{currentstroke}{rgb}{0.000000,0.000000,0.000000}%
\pgfsetstrokecolor{currentstroke}%
\pgfsetdash{}{0pt}%
\pgfpathmoveto{\pgfqpoint{1.683766in}{1.082745in}}%
\pgfpathlineto{\pgfqpoint{1.702114in}{1.095957in}}%
\pgfpathlineto{\pgfqpoint{1.709671in}{1.103395in}}%
\pgfpathlineto{\pgfqpoint{1.717215in}{1.111103in}}%
\pgfpathlineto{\pgfqpoint{1.698892in}{1.098045in}}%
\pgfpathlineto{\pgfqpoint{1.691336in}{1.090258in}}%
\pgfusepath{fill}%
\end{pgfscope}%
\begin{pgfscope}%
\pgfpathrectangle{\pgfqpoint{0.407592in}{0.222904in}}{\pgfqpoint{1.411163in}{1.411163in}}%
\pgfusepath{clip}%
\pgfsetbuttcap%
\pgfsetroundjoin%
\definecolor{currentfill}{rgb}{0.556863,0.654902,0.847059}%
\pgfsetfillcolor{currentfill}%
\pgfsetfillopacity{0.130000}%
\pgfsetlinewidth{0.000000pt}%
\definecolor{currentstroke}{rgb}{0.000000,0.000000,0.000000}%
\pgfsetstrokecolor{currentstroke}%
\pgfsetdash{}{0pt}%
\pgfpathmoveto{\pgfqpoint{1.513763in}{0.954039in}}%
\pgfpathlineto{\pgfqpoint{1.531712in}{0.966941in}}%
\pgfpathlineto{\pgfqpoint{1.549743in}{0.979955in}}%
\pgfpathlineto{\pgfqpoint{1.557480in}{0.985165in}}%
\pgfpathlineto{\pgfqpoint{1.565200in}{0.990715in}}%
\pgfpathlineto{\pgfqpoint{1.547193in}{0.977731in}}%
\pgfpathlineto{\pgfqpoint{1.529268in}{0.964860in}}%
\pgfpathlineto{\pgfqpoint{1.521524in}{0.959267in}}%
\pgfpathlineto{\pgfqpoint{1.513763in}{0.954039in}}%
\pgfpathclose%
\pgfusepath{fill}%
\end{pgfscope}%
\begin{pgfscope}%
\pgfpathrectangle{\pgfqpoint{0.407592in}{0.222904in}}{\pgfqpoint{1.411163in}{1.411163in}}%
\pgfusepath{clip}%
\pgfsetbuttcap%
\pgfsetroundjoin%
\definecolor{currentfill}{rgb}{0.556863,0.654902,0.847059}%
\pgfsetfillcolor{currentfill}%
\pgfsetfillopacity{0.130000}%
\pgfsetlinewidth{0.000000pt}%
\definecolor{currentstroke}{rgb}{0.000000,0.000000,0.000000}%
\pgfsetstrokecolor{currentstroke}%
\pgfsetdash{}{0pt}%
\pgfpathmoveto{\pgfqpoint{0.902599in}{0.815369in}}%
\pgfpathlineto{\pgfqpoint{0.919394in}{0.812536in}}%
\pgfpathlineto{\pgfqpoint{0.936226in}{0.810562in}}%
\pgfpathlineto{\pgfqpoint{0.945087in}{0.803976in}}%
\pgfpathlineto{\pgfqpoint{0.953907in}{0.797518in}}%
\pgfpathlineto{\pgfqpoint{0.937129in}{0.799331in}}%
\pgfpathlineto{\pgfqpoint{0.920388in}{0.802059in}}%
\pgfpathlineto{\pgfqpoint{0.911515in}{0.808655in}}%
\pgfpathlineto{\pgfqpoint{0.902599in}{0.815369in}}%
\pgfpathclose%
\pgfusepath{fill}%
\end{pgfscope}%
\begin{pgfscope}%
\pgfpathrectangle{\pgfqpoint{0.407592in}{0.222904in}}{\pgfqpoint{1.411163in}{1.411163in}}%
\pgfusepath{clip}%
\pgfsetbuttcap%
\pgfsetroundjoin%
\definecolor{currentfill}{rgb}{0.556863,0.654902,0.847059}%
\pgfsetfillcolor{currentfill}%
\pgfsetfillopacity{0.130000}%
\pgfsetlinewidth{0.000000pt}%
\definecolor{currentstroke}{rgb}{0.000000,0.000000,0.000000}%
\pgfsetstrokecolor{currentstroke}%
\pgfsetdash{}{0pt}%
\pgfpathmoveto{\pgfqpoint{1.344419in}{0.853544in}}%
\pgfpathlineto{\pgfqpoint{1.361984in}{0.865032in}}%
\pgfpathlineto{\pgfqpoint{1.379628in}{0.876730in}}%
\pgfpathlineto{\pgfqpoint{1.387596in}{0.878660in}}%
\pgfpathlineto{\pgfqpoint{1.395541in}{0.880990in}}%
\pgfpathlineto{\pgfqpoint{1.377923in}{0.869064in}}%
\pgfpathlineto{\pgfqpoint{1.360383in}{0.857330in}}%
\pgfpathlineto{\pgfqpoint{1.352413in}{0.855223in}}%
\pgfpathlineto{\pgfqpoint{1.344419in}{0.853544in}}%
\pgfpathclose%
\pgfusepath{fill}%
\end{pgfscope}%
\begin{pgfscope}%
\pgfpathrectangle{\pgfqpoint{0.407592in}{0.222904in}}{\pgfqpoint{1.411163in}{1.411163in}}%
\pgfusepath{clip}%
\pgfsetbuttcap%
\pgfsetroundjoin%
\definecolor{currentfill}{rgb}{0.556863,0.654902,0.847059}%
\pgfsetfillcolor{currentfill}%
\pgfsetfillopacity{0.130000}%
\pgfsetlinewidth{0.000000pt}%
\definecolor{currentstroke}{rgb}{0.000000,0.000000,0.000000}%
\pgfsetstrokecolor{currentstroke}%
\pgfsetdash{}{0pt}%
\pgfpathmoveto{\pgfqpoint{1.038754in}{0.787757in}}%
\pgfpathlineto{\pgfqpoint{1.055710in}{0.790569in}}%
\pgfpathlineto{\pgfqpoint{1.072722in}{0.794145in}}%
\pgfpathlineto{\pgfqpoint{1.081269in}{0.789271in}}%
\pgfpathlineto{\pgfqpoint{1.089778in}{0.784632in}}%
\pgfpathlineto{\pgfqpoint{1.072813in}{0.780580in}}%
\pgfpathlineto{\pgfqpoint{1.055905in}{0.777314in}}%
\pgfpathlineto{\pgfqpoint{1.047349in}{0.782428in}}%
\pgfpathlineto{\pgfqpoint{1.038754in}{0.787757in}}%
\pgfpathclose%
\pgfusepath{fill}%
\end{pgfscope}%
\begin{pgfscope}%
\pgfpathrectangle{\pgfqpoint{0.407592in}{0.222904in}}{\pgfqpoint{1.411163in}{1.411163in}}%
\pgfusepath{clip}%
\pgfsetbuttcap%
\pgfsetroundjoin%
\definecolor{currentfill}{rgb}{0.556863,0.654902,0.847059}%
\pgfsetfillcolor{currentfill}%
\pgfsetfillopacity{0.130000}%
\pgfsetlinewidth{0.000000pt}%
\definecolor{currentstroke}{rgb}{0.000000,0.000000,0.000000}%
\pgfsetstrokecolor{currentstroke}%
\pgfsetdash{}{0pt}%
\pgfpathmoveto{\pgfqpoint{1.395541in}{0.880990in}}%
\pgfpathlineto{\pgfqpoint{1.413239in}{0.893092in}}%
\pgfpathlineto{\pgfqpoint{1.431018in}{0.905360in}}%
\pgfpathlineto{\pgfqpoint{1.438916in}{0.908260in}}%
\pgfpathlineto{\pgfqpoint{1.446793in}{0.911546in}}%
\pgfpathlineto{\pgfqpoint{1.429039in}{0.899135in}}%
\pgfpathlineto{\pgfqpoint{1.411366in}{0.886876in}}%
\pgfpathlineto{\pgfqpoint{1.403464in}{0.883727in}}%
\pgfpathlineto{\pgfqpoint{1.395541in}{0.880990in}}%
\pgfpathclose%
\pgfusepath{fill}%
\end{pgfscope}%
\begin{pgfscope}%
\pgfpathrectangle{\pgfqpoint{0.407592in}{0.222904in}}{\pgfqpoint{1.411163in}{1.411163in}}%
\pgfusepath{clip}%
\pgfsetbuttcap%
\pgfsetroundjoin%
\definecolor{currentfill}{rgb}{0.556863,0.654902,0.847059}%
\pgfsetfillcolor{currentfill}%
\pgfsetfillopacity{0.130000}%
\pgfsetlinewidth{0.000000pt}%
\definecolor{currentstroke}{rgb}{0.000000,0.000000,0.000000}%
\pgfsetstrokecolor{currentstroke}%
\pgfsetdash{}{0pt}%
\pgfpathmoveto{\pgfqpoint{1.565200in}{0.990715in}}%
\pgfpathlineto{\pgfqpoint{1.583290in}{1.003808in}}%
\pgfpathlineto{\pgfqpoint{1.601462in}{1.017006in}}%
\pgfpathlineto{\pgfqpoint{1.609142in}{1.022835in}}%
\pgfpathlineto{\pgfqpoint{1.616807in}{1.028979in}}%
\pgfpathlineto{\pgfqpoint{1.598659in}{1.015851in}}%
\pgfpathlineto{\pgfqpoint{1.580594in}{1.002830in}}%
\pgfpathlineto{\pgfqpoint{1.572905in}{0.996605in}}%
\pgfpathlineto{\pgfqpoint{1.565200in}{0.990715in}}%
\pgfpathclose%
\pgfusepath{fill}%
\end{pgfscope}%
\begin{pgfscope}%
\pgfpathrectangle{\pgfqpoint{0.407592in}{0.222904in}}{\pgfqpoint{1.411163in}{1.411163in}}%
\pgfusepath{clip}%
\pgfsetbuttcap%
\pgfsetroundjoin%
\definecolor{currentfill}{rgb}{0.556863,0.654902,0.847059}%
\pgfsetfillcolor{currentfill}%
\pgfsetfillopacity{0.130000}%
\pgfsetlinewidth{0.000000pt}%
\definecolor{currentstroke}{rgb}{0.000000,0.000000,0.000000}%
\pgfsetstrokecolor{currentstroke}%
\pgfsetdash{}{0pt}%
\pgfpathmoveto{\pgfqpoint{0.987584in}{0.796565in}}%
\pgfpathlineto{\pgfqpoint{1.004490in}{0.797373in}}%
\pgfpathlineto{\pgfqpoint{1.021445in}{0.798995in}}%
\pgfpathlineto{\pgfqpoint{1.030119in}{0.793284in}}%
\pgfpathlineto{\pgfqpoint{1.038754in}{0.787757in}}%
\pgfpathlineto{\pgfqpoint{1.021850in}{0.785755in}}%
\pgfpathlineto{\pgfqpoint{1.004995in}{0.784606in}}%
\pgfpathlineto{\pgfqpoint{0.996310in}{0.790503in}}%
\pgfpathlineto{\pgfqpoint{0.987584in}{0.796565in}}%
\pgfpathclose%
\pgfusepath{fill}%
\end{pgfscope}%
\begin{pgfscope}%
\pgfpathrectangle{\pgfqpoint{0.407592in}{0.222904in}}{\pgfqpoint{1.411163in}{1.411163in}}%
\pgfusepath{clip}%
\pgfsetbuttcap%
\pgfsetroundjoin%
\definecolor{currentfill}{rgb}{0.556863,0.654902,0.847059}%
\pgfsetfillcolor{currentfill}%
\pgfsetfillopacity{0.130000}%
\pgfsetlinewidth{0.000000pt}%
\definecolor{currentstroke}{rgb}{0.000000,0.000000,0.000000}%
\pgfsetstrokecolor{currentstroke}%
\pgfsetdash{}{0pt}%
\pgfpathmoveto{\pgfqpoint{0.851042in}{0.837422in}}%
\pgfpathlineto{\pgfqpoint{0.867826in}{0.832864in}}%
\pgfpathlineto{\pgfqpoint{0.884641in}{0.829124in}}%
\pgfpathlineto{\pgfqpoint{0.893641in}{0.822194in}}%
\pgfpathlineto{\pgfqpoint{0.902599in}{0.815369in}}%
\pgfpathlineto{\pgfqpoint{0.885838in}{0.819069in}}%
\pgfpathlineto{\pgfqpoint{0.869109in}{0.823639in}}%
\pgfpathlineto{\pgfqpoint{0.860097in}{0.830479in}}%
\pgfpathlineto{\pgfqpoint{0.851042in}{0.837422in}}%
\pgfpathclose%
\pgfusepath{fill}%
\end{pgfscope}%
\begin{pgfscope}%
\pgfpathrectangle{\pgfqpoint{0.407592in}{0.222904in}}{\pgfqpoint{1.411163in}{1.411163in}}%
\pgfusepath{clip}%
\pgfsetbuttcap%
\pgfsetroundjoin%
\definecolor{currentfill}{rgb}{0.556863,0.654902,0.847059}%
\pgfsetfillcolor{currentfill}%
\pgfsetfillopacity{0.130000}%
\pgfsetlinewidth{0.000000pt}%
\definecolor{currentstroke}{rgb}{0.000000,0.000000,0.000000}%
\pgfsetstrokecolor{currentstroke}%
\pgfsetdash{}{0pt}%
\pgfpathmoveto{\pgfqpoint{1.446793in}{0.911546in}}%
\pgfpathlineto{\pgfqpoint{1.464627in}{0.924100in}}%
\pgfpathlineto{\pgfqpoint{1.482543in}{0.936789in}}%
\pgfpathlineto{\pgfqpoint{1.490376in}{0.940553in}}%
\pgfpathlineto{\pgfqpoint{1.498190in}{0.944682in}}%
\pgfpathlineto{\pgfqpoint{1.480299in}{0.931920in}}%
\pgfpathlineto{\pgfqpoint{1.462488in}{0.919286in}}%
\pgfpathlineto{\pgfqpoint{1.454650in}{0.915220in}}%
\pgfpathlineto{\pgfqpoint{1.446793in}{0.911546in}}%
\pgfpathclose%
\pgfusepath{fill}%
\end{pgfscope}%
\begin{pgfscope}%
\pgfpathrectangle{\pgfqpoint{0.407592in}{0.222904in}}{\pgfqpoint{1.411163in}{1.411163in}}%
\pgfusepath{clip}%
\pgfsetbuttcap%
\pgfsetroundjoin%
\definecolor{currentfill}{rgb}{0.556863,0.654902,0.847059}%
\pgfsetfillcolor{currentfill}%
\pgfsetfillopacity{0.130000}%
\pgfsetlinewidth{0.000000pt}%
\definecolor{currentstroke}{rgb}{0.000000,0.000000,0.000000}%
\pgfsetstrokecolor{currentstroke}%
\pgfsetdash{}{0pt}%
\pgfpathmoveto{\pgfqpoint{1.616807in}{1.028979in}}%
\pgfpathlineto{\pgfqpoint{1.635037in}{1.042211in}}%
\pgfpathlineto{\pgfqpoint{1.653351in}{1.055544in}}%
\pgfpathlineto{\pgfqpoint{1.660976in}{1.061910in}}%
\pgfpathlineto{\pgfqpoint{1.668587in}{1.068568in}}%
\pgfpathlineto{\pgfqpoint{1.650298in}{1.055333in}}%
\pgfpathlineto{\pgfqpoint{1.632092in}{1.042204in}}%
\pgfpathlineto{\pgfqpoint{1.624456in}{1.035437in}}%
\pgfpathlineto{\pgfqpoint{1.616807in}{1.028979in}}%
\pgfpathclose%
\pgfusepath{fill}%
\end{pgfscope}%
\begin{pgfscope}%
\pgfpathrectangle{\pgfqpoint{0.407592in}{0.222904in}}{\pgfqpoint{1.411163in}{1.411163in}}%
\pgfusepath{clip}%
\pgfsetbuttcap%
\pgfsetroundjoin%
\definecolor{currentfill}{rgb}{0.556863,0.654902,0.847059}%
\pgfsetfillcolor{currentfill}%
\pgfsetfillopacity{0.130000}%
\pgfsetlinewidth{0.000000pt}%
\definecolor{currentstroke}{rgb}{0.000000,0.000000,0.000000}%
\pgfsetstrokecolor{currentstroke}%
\pgfsetdash{}{0pt}%
\pgfpathmoveto{\pgfqpoint{1.668587in}{1.068568in}}%
\pgfpathlineto{\pgfqpoint{1.686960in}{1.081904in}}%
\pgfpathlineto{\pgfqpoint{1.694544in}{1.088792in}}%
\pgfpathlineto{\pgfqpoint{1.702114in}{1.095957in}}%
\pgfpathlineto{\pgfqpoint{1.683766in}{1.082745in}}%
\pgfpathlineto{\pgfqpoint{1.676183in}{1.075514in}}%
\pgfusepath{fill}%
\end{pgfscope}%
\begin{pgfscope}%
\pgfpathrectangle{\pgfqpoint{0.407592in}{0.222904in}}{\pgfqpoint{1.411163in}{1.411163in}}%
\pgfusepath{clip}%
\pgfsetbuttcap%
\pgfsetroundjoin%
\definecolor{currentfill}{rgb}{0.556863,0.654902,0.847059}%
\pgfsetfillcolor{currentfill}%
\pgfsetfillopacity{0.130000}%
\pgfsetlinewidth{0.000000pt}%
\definecolor{currentstroke}{rgb}{0.000000,0.000000,0.000000}%
\pgfsetstrokecolor{currentstroke}%
\pgfsetdash{}{0pt}%
\pgfpathmoveto{\pgfqpoint{0.936226in}{0.810562in}}%
\pgfpathlineto{\pgfqpoint{0.953097in}{0.809434in}}%
\pgfpathlineto{\pgfqpoint{0.970010in}{0.809134in}}%
\pgfpathlineto{\pgfqpoint{0.978818in}{0.802779in}}%
\pgfpathlineto{\pgfqpoint{0.987584in}{0.796565in}}%
\pgfpathlineto{\pgfqpoint{0.970724in}{0.796604in}}%
\pgfpathlineto{\pgfqpoint{0.953907in}{0.797518in}}%
\pgfpathlineto{\pgfqpoint{0.945087in}{0.803976in}}%
\pgfpathlineto{\pgfqpoint{0.936226in}{0.810562in}}%
\pgfpathclose%
\pgfusepath{fill}%
\end{pgfscope}%
\begin{pgfscope}%
\pgfpathrectangle{\pgfqpoint{0.407592in}{0.222904in}}{\pgfqpoint{1.411163in}{1.411163in}}%
\pgfusepath{clip}%
\pgfsetbuttcap%
\pgfsetroundjoin%
\definecolor{currentfill}{rgb}{0.556863,0.654902,0.847059}%
\pgfsetfillcolor{currentfill}%
\pgfsetfillopacity{0.130000}%
\pgfsetlinewidth{0.000000pt}%
\definecolor{currentstroke}{rgb}{0.000000,0.000000,0.000000}%
\pgfsetstrokecolor{currentstroke}%
\pgfsetdash{}{0pt}%
\pgfpathmoveto{\pgfqpoint{1.174986in}{0.801278in}}%
\pgfpathlineto{\pgfqpoint{1.192228in}{0.809005in}}%
\pgfpathlineto{\pgfqpoint{1.209540in}{0.817214in}}%
\pgfpathlineto{\pgfqpoint{1.217822in}{0.815000in}}%
\pgfpathlineto{\pgfqpoint{1.226070in}{0.813128in}}%
\pgfpathlineto{\pgfqpoint{1.208794in}{0.804421in}}%
\pgfpathlineto{\pgfqpoint{1.191588in}{0.796177in}}%
\pgfpathlineto{\pgfqpoint{1.183304in}{0.798557in}}%
\pgfpathlineto{\pgfqpoint{1.174986in}{0.801278in}}%
\pgfpathclose%
\pgfusepath{fill}%
\end{pgfscope}%
\begin{pgfscope}%
\pgfpathrectangle{\pgfqpoint{0.407592in}{0.222904in}}{\pgfqpoint{1.411163in}{1.411163in}}%
\pgfusepath{clip}%
\pgfsetbuttcap%
\pgfsetroundjoin%
\definecolor{currentfill}{rgb}{0.556863,0.654902,0.847059}%
\pgfsetfillcolor{currentfill}%
\pgfsetfillopacity{0.130000}%
\pgfsetlinewidth{0.000000pt}%
\definecolor{currentstroke}{rgb}{0.000000,0.000000,0.000000}%
\pgfsetstrokecolor{currentstroke}%
\pgfsetdash{}{0pt}%
\pgfpathmoveto{\pgfqpoint{1.226070in}{0.813128in}}%
\pgfpathlineto{\pgfqpoint{1.243419in}{0.822256in}}%
\pgfpathlineto{\pgfqpoint{1.260842in}{0.831767in}}%
\pgfpathlineto{\pgfqpoint{1.269027in}{0.830712in}}%
\pgfpathlineto{\pgfqpoint{1.277182in}{0.830025in}}%
\pgfpathlineto{\pgfqpoint{1.259790in}{0.820073in}}%
\pgfpathlineto{\pgfqpoint{1.242474in}{0.810480in}}%
\pgfpathlineto{\pgfqpoint{1.234287in}{0.811615in}}%
\pgfpathlineto{\pgfqpoint{1.226070in}{0.813128in}}%
\pgfpathclose%
\pgfusepath{fill}%
\end{pgfscope}%
\begin{pgfscope}%
\pgfpathrectangle{\pgfqpoint{0.407592in}{0.222904in}}{\pgfqpoint{1.411163in}{1.411163in}}%
\pgfusepath{clip}%
\pgfsetbuttcap%
\pgfsetroundjoin%
\definecolor{currentfill}{rgb}{0.556863,0.654902,0.847059}%
\pgfsetfillcolor{currentfill}%
\pgfsetfillopacity{0.130000}%
\pgfsetlinewidth{0.000000pt}%
\definecolor{currentstroke}{rgb}{0.000000,0.000000,0.000000}%
\pgfsetstrokecolor{currentstroke}%
\pgfsetdash{}{0pt}%
\pgfpathmoveto{\pgfqpoint{1.498190in}{0.944682in}}%
\pgfpathlineto{\pgfqpoint{1.516163in}{0.957565in}}%
\pgfpathlineto{\pgfqpoint{1.534218in}{0.970564in}}%
\pgfpathlineto{\pgfqpoint{1.541989in}{0.975088in}}%
\pgfpathlineto{\pgfqpoint{1.549743in}{0.979955in}}%
\pgfpathlineto{\pgfqpoint{1.531712in}{0.966941in}}%
\pgfpathlineto{\pgfqpoint{1.513763in}{0.954039in}}%
\pgfpathlineto{\pgfqpoint{1.505986in}{0.949177in}}%
\pgfpathlineto{\pgfqpoint{1.498190in}{0.944682in}}%
\pgfpathclose%
\pgfusepath{fill}%
\end{pgfscope}%
\begin{pgfscope}%
\pgfpathrectangle{\pgfqpoint{0.407592in}{0.222904in}}{\pgfqpoint{1.411163in}{1.411163in}}%
\pgfusepath{clip}%
\pgfsetbuttcap%
\pgfsetroundjoin%
\definecolor{currentfill}{rgb}{0.556863,0.654902,0.847059}%
\pgfsetfillcolor{currentfill}%
\pgfsetfillopacity{0.130000}%
\pgfsetlinewidth{0.000000pt}%
\definecolor{currentstroke}{rgb}{0.000000,0.000000,0.000000}%
\pgfsetstrokecolor{currentstroke}%
\pgfsetdash{}{0pt}%
\pgfpathmoveto{\pgfqpoint{1.123886in}{0.794886in}}%
\pgfpathlineto{\pgfqpoint{1.141034in}{0.800983in}}%
\pgfpathlineto{\pgfqpoint{1.158247in}{0.807663in}}%
\pgfpathlineto{\pgfqpoint{1.166634in}{0.804320in}}%
\pgfpathlineto{\pgfqpoint{1.174986in}{0.801278in}}%
\pgfpathlineto{\pgfqpoint{1.157813in}{0.794077in}}%
\pgfpathlineto{\pgfqpoint{1.140706in}{0.787452in}}%
\pgfpathlineto{\pgfqpoint{1.132314in}{0.791023in}}%
\pgfpathlineto{\pgfqpoint{1.123886in}{0.794886in}}%
\pgfpathclose%
\pgfusepath{fill}%
\end{pgfscope}%
\begin{pgfscope}%
\pgfpathrectangle{\pgfqpoint{0.407592in}{0.222904in}}{\pgfqpoint{1.411163in}{1.411163in}}%
\pgfusepath{clip}%
\pgfsetbuttcap%
\pgfsetroundjoin%
\definecolor{currentfill}{rgb}{0.556863,0.654902,0.847059}%
\pgfsetfillcolor{currentfill}%
\pgfsetfillopacity{0.130000}%
\pgfsetlinewidth{0.000000pt}%
\definecolor{currentstroke}{rgb}{0.000000,0.000000,0.000000}%
\pgfsetstrokecolor{currentstroke}%
\pgfsetdash{}{0pt}%
\pgfpathmoveto{\pgfqpoint{1.277182in}{0.830025in}}%
\pgfpathlineto{\pgfqpoint{1.294648in}{0.840304in}}%
\pgfpathlineto{\pgfqpoint{1.312192in}{0.850882in}}%
\pgfpathlineto{\pgfqpoint{1.320288in}{0.850962in}}%
\pgfpathlineto{\pgfqpoint{1.328358in}{0.851425in}}%
\pgfpathlineto{\pgfqpoint{1.310843in}{0.840485in}}%
\pgfpathlineto{\pgfqpoint{1.293406in}{0.829819in}}%
\pgfpathlineto{\pgfqpoint{1.285308in}{0.829722in}}%
\pgfpathlineto{\pgfqpoint{1.277182in}{0.830025in}}%
\pgfpathclose%
\pgfusepath{fill}%
\end{pgfscope}%
\begin{pgfscope}%
\pgfpathrectangle{\pgfqpoint{0.407592in}{0.222904in}}{\pgfqpoint{1.411163in}{1.411163in}}%
\pgfusepath{clip}%
\pgfsetbuttcap%
\pgfsetroundjoin%
\definecolor{currentfill}{rgb}{0.556863,0.654902,0.847059}%
\pgfsetfillcolor{currentfill}%
\pgfsetfillopacity{0.130000}%
\pgfsetlinewidth{0.000000pt}%
\definecolor{currentstroke}{rgb}{0.000000,0.000000,0.000000}%
\pgfsetstrokecolor{currentstroke}%
\pgfsetdash{}{0pt}%
\pgfpathmoveto{\pgfqpoint{1.328358in}{0.851425in}}%
\pgfpathlineto{\pgfqpoint{1.345950in}{0.862615in}}%
\pgfpathlineto{\pgfqpoint{1.363621in}{0.874036in}}%
\pgfpathlineto{\pgfqpoint{1.371637in}{0.875192in}}%
\pgfpathlineto{\pgfqpoint{1.379628in}{0.876730in}}%
\pgfpathlineto{\pgfqpoint{1.361984in}{0.865032in}}%
\pgfpathlineto{\pgfqpoint{1.344419in}{0.853544in}}%
\pgfpathlineto{\pgfqpoint{1.336401in}{0.852281in}}%
\pgfpathlineto{\pgfqpoint{1.328358in}{0.851425in}}%
\pgfpathclose%
\pgfusepath{fill}%
\end{pgfscope}%
\begin{pgfscope}%
\pgfpathrectangle{\pgfqpoint{0.407592in}{0.222904in}}{\pgfqpoint{1.411163in}{1.411163in}}%
\pgfusepath{clip}%
\pgfsetbuttcap%
\pgfsetroundjoin%
\definecolor{currentfill}{rgb}{0.556863,0.654902,0.847059}%
\pgfsetfillcolor{currentfill}%
\pgfsetfillopacity{0.130000}%
\pgfsetlinewidth{0.000000pt}%
\definecolor{currentstroke}{rgb}{0.000000,0.000000,0.000000}%
\pgfsetstrokecolor{currentstroke}%
\pgfsetdash{}{0pt}%
\pgfpathmoveto{\pgfqpoint{1.072722in}{0.794145in}}%
\pgfpathlineto{\pgfqpoint{1.089791in}{0.798439in}}%
\pgfpathlineto{\pgfqpoint{1.106920in}{0.803406in}}%
\pgfpathlineto{\pgfqpoint{1.115421in}{0.799019in}}%
\pgfpathlineto{\pgfqpoint{1.123886in}{0.794886in}}%
\pgfpathlineto{\pgfqpoint{1.106801in}{0.789419in}}%
\pgfpathlineto{\pgfqpoint{1.089778in}{0.784632in}}%
\pgfpathlineto{\pgfqpoint{1.081269in}{0.789271in}}%
\pgfpathlineto{\pgfqpoint{1.072722in}{0.794145in}}%
\pgfpathclose%
\pgfusepath{fill}%
\end{pgfscope}%
\begin{pgfscope}%
\pgfpathrectangle{\pgfqpoint{0.407592in}{0.222904in}}{\pgfqpoint{1.411163in}{1.411163in}}%
\pgfusepath{clip}%
\pgfsetbuttcap%
\pgfsetroundjoin%
\definecolor{currentfill}{rgb}{0.556863,0.654902,0.847059}%
\pgfsetfillcolor{currentfill}%
\pgfsetfillopacity{0.130000}%
\pgfsetlinewidth{0.000000pt}%
\definecolor{currentstroke}{rgb}{0.000000,0.000000,0.000000}%
\pgfsetstrokecolor{currentstroke}%
\pgfsetdash{}{0pt}%
\pgfpathmoveto{\pgfqpoint{1.549743in}{0.979955in}}%
\pgfpathlineto{\pgfqpoint{1.567857in}{0.993078in}}%
\pgfpathlineto{\pgfqpoint{1.586053in}{1.006306in}}%
\pgfpathlineto{\pgfqpoint{1.593765in}{1.011496in}}%
\pgfpathlineto{\pgfqpoint{1.601462in}{1.017006in}}%
\pgfpathlineto{\pgfqpoint{1.583290in}{1.003808in}}%
\pgfpathlineto{\pgfqpoint{1.565200in}{0.990715in}}%
\pgfpathlineto{\pgfqpoint{1.557480in}{0.985165in}}%
\pgfpathlineto{\pgfqpoint{1.549743in}{0.979955in}}%
\pgfpathclose%
\pgfusepath{fill}%
\end{pgfscope}%
\begin{pgfscope}%
\pgfpathrectangle{\pgfqpoint{0.407592in}{0.222904in}}{\pgfqpoint{1.411163in}{1.411163in}}%
\pgfusepath{clip}%
\pgfsetbuttcap%
\pgfsetroundjoin%
\definecolor{currentfill}{rgb}{0.556863,0.654902,0.847059}%
\pgfsetfillcolor{currentfill}%
\pgfsetfillopacity{0.130000}%
\pgfsetlinewidth{0.000000pt}%
\definecolor{currentstroke}{rgb}{0.000000,0.000000,0.000000}%
\pgfsetstrokecolor{currentstroke}%
\pgfsetdash{}{0pt}%
\pgfpathmoveto{\pgfqpoint{1.379628in}{0.876730in}}%
\pgfpathlineto{\pgfqpoint{1.397352in}{0.888622in}}%
\pgfpathlineto{\pgfqpoint{1.415156in}{0.900695in}}%
\pgfpathlineto{\pgfqpoint{1.423098in}{0.902840in}}%
\pgfpathlineto{\pgfqpoint{1.431018in}{0.905360in}}%
\pgfpathlineto{\pgfqpoint{1.413239in}{0.893092in}}%
\pgfpathlineto{\pgfqpoint{1.395541in}{0.880990in}}%
\pgfpathlineto{\pgfqpoint{1.387596in}{0.878660in}}%
\pgfpathlineto{\pgfqpoint{1.379628in}{0.876730in}}%
\pgfpathclose%
\pgfusepath{fill}%
\end{pgfscope}%
\begin{pgfscope}%
\pgfpathrectangle{\pgfqpoint{0.407592in}{0.222904in}}{\pgfqpoint{1.411163in}{1.411163in}}%
\pgfusepath{clip}%
\pgfsetbuttcap%
\pgfsetroundjoin%
\definecolor{currentfill}{rgb}{0.556863,0.654902,0.847059}%
\pgfsetfillcolor{currentfill}%
\pgfsetfillopacity{0.130000}%
\pgfsetlinewidth{0.000000pt}%
\definecolor{currentstroke}{rgb}{0.000000,0.000000,0.000000}%
\pgfsetstrokecolor{currentstroke}%
\pgfsetdash{}{0pt}%
\pgfpathmoveto{\pgfqpoint{0.884641in}{0.829124in}}%
\pgfpathlineto{\pgfqpoint{0.901491in}{0.826198in}}%
\pgfpathlineto{\pgfqpoint{0.918377in}{0.824082in}}%
\pgfpathlineto{\pgfqpoint{0.927322in}{0.817267in}}%
\pgfpathlineto{\pgfqpoint{0.936226in}{0.810562in}}%
\pgfpathlineto{\pgfqpoint{0.919394in}{0.812536in}}%
\pgfpathlineto{\pgfqpoint{0.902599in}{0.815369in}}%
\pgfpathlineto{\pgfqpoint{0.893641in}{0.822194in}}%
\pgfpathlineto{\pgfqpoint{0.884641in}{0.829124in}}%
\pgfpathclose%
\pgfusepath{fill}%
\end{pgfscope}%
\begin{pgfscope}%
\pgfpathrectangle{\pgfqpoint{0.407592in}{0.222904in}}{\pgfqpoint{1.411163in}{1.411163in}}%
\pgfusepath{clip}%
\pgfsetbuttcap%
\pgfsetroundjoin%
\definecolor{currentfill}{rgb}{0.556863,0.654902,0.847059}%
\pgfsetfillcolor{currentfill}%
\pgfsetfillopacity{0.130000}%
\pgfsetlinewidth{0.000000pt}%
\definecolor{currentstroke}{rgb}{0.000000,0.000000,0.000000}%
\pgfsetstrokecolor{currentstroke}%
\pgfsetdash{}{0pt}%
\pgfpathmoveto{\pgfqpoint{1.021445in}{0.798995in}}%
\pgfpathlineto{\pgfqpoint{1.038451in}{0.801395in}}%
\pgfpathlineto{\pgfqpoint{1.055511in}{0.804534in}}%
\pgfpathlineto{\pgfqpoint{1.064136in}{0.799238in}}%
\pgfpathlineto{\pgfqpoint{1.072722in}{0.794145in}}%
\pgfpathlineto{\pgfqpoint{1.055710in}{0.790569in}}%
\pgfpathlineto{\pgfqpoint{1.038754in}{0.787757in}}%
\pgfpathlineto{\pgfqpoint{1.030119in}{0.793284in}}%
\pgfpathlineto{\pgfqpoint{1.021445in}{0.798995in}}%
\pgfpathclose%
\pgfusepath{fill}%
\end{pgfscope}%
\begin{pgfscope}%
\pgfpathrectangle{\pgfqpoint{0.407592in}{0.222904in}}{\pgfqpoint{1.411163in}{1.411163in}}%
\pgfusepath{clip}%
\pgfsetbuttcap%
\pgfsetroundjoin%
\definecolor{currentfill}{rgb}{0.556863,0.654902,0.847059}%
\pgfsetfillcolor{currentfill}%
\pgfsetfillopacity{0.130000}%
\pgfsetlinewidth{0.000000pt}%
\definecolor{currentstroke}{rgb}{0.000000,0.000000,0.000000}%
\pgfsetstrokecolor{currentstroke}%
\pgfsetdash{}{0pt}%
\pgfpathmoveto{\pgfqpoint{1.431018in}{0.905360in}}%
\pgfpathlineto{\pgfqpoint{1.448876in}{0.917782in}}%
\pgfpathlineto{\pgfqpoint{1.466817in}{0.930348in}}%
\pgfpathlineto{\pgfqpoint{1.474690in}{0.933389in}}%
\pgfpathlineto{\pgfqpoint{1.482543in}{0.936789in}}%
\pgfpathlineto{\pgfqpoint{1.464627in}{0.924100in}}%
\pgfpathlineto{\pgfqpoint{1.446793in}{0.911546in}}%
\pgfpathlineto{\pgfqpoint{1.438916in}{0.908260in}}%
\pgfpathlineto{\pgfqpoint{1.431018in}{0.905360in}}%
\pgfpathclose%
\pgfusepath{fill}%
\end{pgfscope}%
\begin{pgfscope}%
\pgfpathrectangle{\pgfqpoint{0.407592in}{0.222904in}}{\pgfqpoint{1.411163in}{1.411163in}}%
\pgfusepath{clip}%
\pgfsetbuttcap%
\pgfsetroundjoin%
\definecolor{currentfill}{rgb}{0.556863,0.654902,0.847059}%
\pgfsetfillcolor{currentfill}%
\pgfsetfillopacity{0.130000}%
\pgfsetlinewidth{0.000000pt}%
\definecolor{currentstroke}{rgb}{0.000000,0.000000,0.000000}%
\pgfsetstrokecolor{currentstroke}%
\pgfsetdash{}{0pt}%
\pgfpathmoveto{\pgfqpoint{1.601462in}{1.017006in}}%
\pgfpathlineto{\pgfqpoint{1.619717in}{1.030305in}}%
\pgfpathlineto{\pgfqpoint{1.638056in}{1.043703in}}%
\pgfpathlineto{\pgfqpoint{1.645711in}{1.049475in}}%
\pgfpathlineto{\pgfqpoint{1.653351in}{1.055544in}}%
\pgfpathlineto{\pgfqpoint{1.635037in}{1.042211in}}%
\pgfpathlineto{\pgfqpoint{1.616807in}{1.028979in}}%
\pgfpathlineto{\pgfqpoint{1.609142in}{1.022835in}}%
\pgfpathlineto{\pgfqpoint{1.601462in}{1.017006in}}%
\pgfpathclose%
\pgfusepath{fill}%
\end{pgfscope}%
\begin{pgfscope}%
\pgfpathrectangle{\pgfqpoint{0.407592in}{0.222904in}}{\pgfqpoint{1.411163in}{1.411163in}}%
\pgfusepath{clip}%
\pgfsetbuttcap%
\pgfsetroundjoin%
\definecolor{currentfill}{rgb}{0.556863,0.654902,0.847059}%
\pgfsetfillcolor{currentfill}%
\pgfsetfillopacity{0.130000}%
\pgfsetlinewidth{0.000000pt}%
\definecolor{currentstroke}{rgb}{0.000000,0.000000,0.000000}%
\pgfsetstrokecolor{currentstroke}%
\pgfsetdash{}{0pt}%
\pgfpathmoveto{\pgfqpoint{1.653351in}{1.055544in}}%
\pgfpathlineto{\pgfqpoint{1.671749in}{1.068975in}}%
\pgfpathlineto{\pgfqpoint{1.679362in}{1.075297in}}%
\pgfpathlineto{\pgfqpoint{1.686960in}{1.081904in}}%
\pgfpathlineto{\pgfqpoint{1.668587in}{1.068568in}}%
\pgfpathlineto{\pgfqpoint{1.660976in}{1.061910in}}%
\pgfusepath{fill}%
\end{pgfscope}%
\begin{pgfscope}%
\pgfpathrectangle{\pgfqpoint{0.407592in}{0.222904in}}{\pgfqpoint{1.411163in}{1.411163in}}%
\pgfusepath{clip}%
\pgfsetbuttcap%
\pgfsetroundjoin%
\definecolor{currentfill}{rgb}{0.556863,0.654902,0.847059}%
\pgfsetfillcolor{currentfill}%
\pgfsetfillopacity{0.130000}%
\pgfsetlinewidth{0.000000pt}%
\definecolor{currentstroke}{rgb}{0.000000,0.000000,0.000000}%
\pgfsetstrokecolor{currentstroke}%
\pgfsetdash{}{0pt}%
\pgfpathmoveto{\pgfqpoint{0.970010in}{0.809134in}}%
\pgfpathlineto{\pgfqpoint{0.986968in}{0.809636in}}%
\pgfpathlineto{\pgfqpoint{1.003975in}{0.810916in}}%
\pgfpathlineto{\pgfqpoint{1.012730in}{0.804876in}}%
\pgfpathlineto{\pgfqpoint{1.021445in}{0.798995in}}%
\pgfpathlineto{\pgfqpoint{1.004490in}{0.797373in}}%
\pgfpathlineto{\pgfqpoint{0.987584in}{0.796565in}}%
\pgfpathlineto{\pgfqpoint{0.978818in}{0.802779in}}%
\pgfpathlineto{\pgfqpoint{0.970010in}{0.809134in}}%
\pgfpathclose%
\pgfusepath{fill}%
\end{pgfscope}%
\begin{pgfscope}%
\pgfpathrectangle{\pgfqpoint{0.407592in}{0.222904in}}{\pgfqpoint{1.411163in}{1.411163in}}%
\pgfusepath{clip}%
\pgfsetbuttcap%
\pgfsetroundjoin%
\definecolor{currentfill}{rgb}{0.556863,0.654902,0.847059}%
\pgfsetfillcolor{currentfill}%
\pgfsetfillopacity{0.130000}%
\pgfsetlinewidth{0.000000pt}%
\definecolor{currentstroke}{rgb}{0.000000,0.000000,0.000000}%
\pgfsetstrokecolor{currentstroke}%
\pgfsetdash{}{0pt}%
\pgfpathmoveto{\pgfqpoint{1.482543in}{0.936789in}}%
\pgfpathlineto{\pgfqpoint{1.500540in}{0.949606in}}%
\pgfpathlineto{\pgfqpoint{1.518620in}{0.962544in}}%
\pgfpathlineto{\pgfqpoint{1.526428in}{0.966383in}}%
\pgfpathlineto{\pgfqpoint{1.534218in}{0.970564in}}%
\pgfpathlineto{\pgfqpoint{1.516163in}{0.957565in}}%
\pgfpathlineto{\pgfqpoint{1.498190in}{0.944682in}}%
\pgfpathlineto{\pgfqpoint{1.490376in}{0.940553in}}%
\pgfpathlineto{\pgfqpoint{1.482543in}{0.936789in}}%
\pgfpathclose%
\pgfusepath{fill}%
\end{pgfscope}%
\begin{pgfscope}%
\pgfpathrectangle{\pgfqpoint{0.407592in}{0.222904in}}{\pgfqpoint{1.411163in}{1.411163in}}%
\pgfusepath{clip}%
\pgfsetbuttcap%
\pgfsetroundjoin%
\definecolor{currentfill}{rgb}{0.556863,0.654902,0.847059}%
\pgfsetfillcolor{currentfill}%
\pgfsetfillopacity{0.130000}%
\pgfsetlinewidth{0.000000pt}%
\definecolor{currentstroke}{rgb}{0.000000,0.000000,0.000000}%
\pgfsetstrokecolor{currentstroke}%
\pgfsetdash{}{0pt}%
\pgfpathmoveto{\pgfqpoint{1.209540in}{0.817214in}}%
\pgfpathlineto{\pgfqpoint{1.226924in}{0.825864in}}%
\pgfpathlineto{\pgfqpoint{1.244381in}{0.834920in}}%
\pgfpathlineto{\pgfqpoint{1.252627in}{0.833175in}}%
\pgfpathlineto{\pgfqpoint{1.260842in}{0.831767in}}%
\pgfpathlineto{\pgfqpoint{1.243419in}{0.822256in}}%
\pgfpathlineto{\pgfqpoint{1.226070in}{0.813128in}}%
\pgfpathlineto{\pgfqpoint{1.217822in}{0.815000in}}%
\pgfpathlineto{\pgfqpoint{1.209540in}{0.817214in}}%
\pgfpathclose%
\pgfusepath{fill}%
\end{pgfscope}%
\begin{pgfscope}%
\pgfpathrectangle{\pgfqpoint{0.407592in}{0.222904in}}{\pgfqpoint{1.411163in}{1.411163in}}%
\pgfusepath{clip}%
\pgfsetbuttcap%
\pgfsetroundjoin%
\definecolor{currentfill}{rgb}{0.556863,0.654902,0.847059}%
\pgfsetfillcolor{currentfill}%
\pgfsetfillopacity{0.130000}%
\pgfsetlinewidth{0.000000pt}%
\definecolor{currentstroke}{rgb}{0.000000,0.000000,0.000000}%
\pgfsetstrokecolor{currentstroke}%
\pgfsetdash{}{0pt}%
\pgfpathmoveto{\pgfqpoint{1.158247in}{0.807663in}}%
\pgfpathlineto{\pgfqpoint{1.175528in}{0.814880in}}%
\pgfpathlineto{\pgfqpoint{1.192878in}{0.822594in}}%
\pgfpathlineto{\pgfqpoint{1.201226in}{0.819751in}}%
\pgfpathlineto{\pgfqpoint{1.209540in}{0.817214in}}%
\pgfpathlineto{\pgfqpoint{1.192228in}{0.809005in}}%
\pgfpathlineto{\pgfqpoint{1.174986in}{0.801278in}}%
\pgfpathlineto{\pgfqpoint{1.166634in}{0.804320in}}%
\pgfpathlineto{\pgfqpoint{1.158247in}{0.807663in}}%
\pgfpathclose%
\pgfusepath{fill}%
\end{pgfscope}%
\begin{pgfscope}%
\pgfpathrectangle{\pgfqpoint{0.407592in}{0.222904in}}{\pgfqpoint{1.411163in}{1.411163in}}%
\pgfusepath{clip}%
\pgfsetbuttcap%
\pgfsetroundjoin%
\definecolor{currentfill}{rgb}{0.556863,0.654902,0.847059}%
\pgfsetfillcolor{currentfill}%
\pgfsetfillopacity{0.130000}%
\pgfsetlinewidth{0.000000pt}%
\definecolor{currentstroke}{rgb}{0.000000,0.000000,0.000000}%
\pgfsetstrokecolor{currentstroke}%
\pgfsetdash{}{0pt}%
\pgfpathmoveto{\pgfqpoint{1.260842in}{0.831767in}}%
\pgfpathlineto{\pgfqpoint{1.278340in}{0.841630in}}%
\pgfpathlineto{\pgfqpoint{1.295914in}{0.851815in}}%
\pgfpathlineto{\pgfqpoint{1.304067in}{0.851171in}}%
\pgfpathlineto{\pgfqpoint{1.312192in}{0.850882in}}%
\pgfpathlineto{\pgfqpoint{1.294648in}{0.840304in}}%
\pgfpathlineto{\pgfqpoint{1.277182in}{0.830025in}}%
\pgfpathlineto{\pgfqpoint{1.269027in}{0.830712in}}%
\pgfpathlineto{\pgfqpoint{1.260842in}{0.831767in}}%
\pgfpathclose%
\pgfusepath{fill}%
\end{pgfscope}%
\begin{pgfscope}%
\pgfpathrectangle{\pgfqpoint{0.407592in}{0.222904in}}{\pgfqpoint{1.411163in}{1.411163in}}%
\pgfusepath{clip}%
\pgfsetbuttcap%
\pgfsetroundjoin%
\definecolor{currentfill}{rgb}{0.556863,0.654902,0.847059}%
\pgfsetfillcolor{currentfill}%
\pgfsetfillopacity{0.130000}%
\pgfsetlinewidth{0.000000pt}%
\definecolor{currentstroke}{rgb}{0.000000,0.000000,0.000000}%
\pgfsetstrokecolor{currentstroke}%
\pgfsetdash{}{0pt}%
\pgfpathmoveto{\pgfqpoint{0.918377in}{0.824082in}}%
\pgfpathlineto{\pgfqpoint{0.935302in}{0.822765in}}%
\pgfpathlineto{\pgfqpoint{0.952269in}{0.822230in}}%
\pgfpathlineto{\pgfqpoint{0.961161in}{0.815620in}}%
\pgfpathlineto{\pgfqpoint{0.970010in}{0.809134in}}%
\pgfpathlineto{\pgfqpoint{0.953097in}{0.809434in}}%
\pgfpathlineto{\pgfqpoint{0.936226in}{0.810562in}}%
\pgfpathlineto{\pgfqpoint{0.927322in}{0.817267in}}%
\pgfpathlineto{\pgfqpoint{0.918377in}{0.824082in}}%
\pgfpathclose%
\pgfusepath{fill}%
\end{pgfscope}%
\begin{pgfscope}%
\pgfpathrectangle{\pgfqpoint{0.407592in}{0.222904in}}{\pgfqpoint{1.411163in}{1.411163in}}%
\pgfusepath{clip}%
\pgfsetbuttcap%
\pgfsetroundjoin%
\definecolor{currentfill}{rgb}{0.556863,0.654902,0.847059}%
\pgfsetfillcolor{currentfill}%
\pgfsetfillopacity{0.130000}%
\pgfsetlinewidth{0.000000pt}%
\definecolor{currentstroke}{rgb}{0.000000,0.000000,0.000000}%
\pgfsetstrokecolor{currentstroke}%
\pgfsetdash{}{0pt}%
\pgfpathmoveto{\pgfqpoint{1.534218in}{0.970564in}}%
\pgfpathlineto{\pgfqpoint{1.552356in}{0.983674in}}%
\pgfpathlineto{\pgfqpoint{1.570577in}{0.996891in}}%
\pgfpathlineto{\pgfqpoint{1.578324in}{1.001437in}}%
\pgfpathlineto{\pgfqpoint{1.586053in}{1.006306in}}%
\pgfpathlineto{\pgfqpoint{1.567857in}{0.993078in}}%
\pgfpathlineto{\pgfqpoint{1.549743in}{0.979955in}}%
\pgfpathlineto{\pgfqpoint{1.541989in}{0.975088in}}%
\pgfpathlineto{\pgfqpoint{1.534218in}{0.970564in}}%
\pgfpathclose%
\pgfusepath{fill}%
\end{pgfscope}%
\begin{pgfscope}%
\pgfpathrectangle{\pgfqpoint{0.407592in}{0.222904in}}{\pgfqpoint{1.411163in}{1.411163in}}%
\pgfusepath{clip}%
\pgfsetbuttcap%
\pgfsetroundjoin%
\definecolor{currentfill}{rgb}{0.556863,0.654902,0.847059}%
\pgfsetfillcolor{currentfill}%
\pgfsetfillopacity{0.130000}%
\pgfsetlinewidth{0.000000pt}%
\definecolor{currentstroke}{rgb}{0.000000,0.000000,0.000000}%
\pgfsetstrokecolor{currentstroke}%
\pgfsetdash{}{0pt}%
\pgfpathmoveto{\pgfqpoint{1.312192in}{0.850882in}}%
\pgfpathlineto{\pgfqpoint{1.329812in}{0.861734in}}%
\pgfpathlineto{\pgfqpoint{1.347511in}{0.872838in}}%
\pgfpathlineto{\pgfqpoint{1.355579in}{0.873255in}}%
\pgfpathlineto{\pgfqpoint{1.363621in}{0.874036in}}%
\pgfpathlineto{\pgfqpoint{1.345950in}{0.862615in}}%
\pgfpathlineto{\pgfqpoint{1.328358in}{0.851425in}}%
\pgfpathlineto{\pgfqpoint{1.320288in}{0.850962in}}%
\pgfpathlineto{\pgfqpoint{1.312192in}{0.850882in}}%
\pgfpathclose%
\pgfusepath{fill}%
\end{pgfscope}%
\begin{pgfscope}%
\pgfpathrectangle{\pgfqpoint{0.407592in}{0.222904in}}{\pgfqpoint{1.411163in}{1.411163in}}%
\pgfusepath{clip}%
\pgfsetbuttcap%
\pgfsetroundjoin%
\definecolor{currentfill}{rgb}{0.556863,0.654902,0.847059}%
\pgfsetfillcolor{currentfill}%
\pgfsetfillopacity{0.130000}%
\pgfsetlinewidth{0.000000pt}%
\definecolor{currentstroke}{rgb}{0.000000,0.000000,0.000000}%
\pgfsetstrokecolor{currentstroke}%
\pgfsetdash{}{0pt}%
\pgfpathmoveto{\pgfqpoint{1.106920in}{0.803406in}}%
\pgfpathlineto{\pgfqpoint{1.124111in}{0.809000in}}%
\pgfpathlineto{\pgfqpoint{1.141367in}{0.815180in}}%
\pgfpathlineto{\pgfqpoint{1.149825in}{0.811288in}}%
\pgfpathlineto{\pgfqpoint{1.158247in}{0.807663in}}%
\pgfpathlineto{\pgfqpoint{1.141034in}{0.800983in}}%
\pgfpathlineto{\pgfqpoint{1.123886in}{0.794886in}}%
\pgfpathlineto{\pgfqpoint{1.115421in}{0.799019in}}%
\pgfpathlineto{\pgfqpoint{1.106920in}{0.803406in}}%
\pgfpathclose%
\pgfusepath{fill}%
\end{pgfscope}%
\begin{pgfscope}%
\pgfpathrectangle{\pgfqpoint{0.407592in}{0.222904in}}{\pgfqpoint{1.411163in}{1.411163in}}%
\pgfusepath{clip}%
\pgfsetbuttcap%
\pgfsetroundjoin%
\definecolor{currentfill}{rgb}{0.556863,0.654902,0.847059}%
\pgfsetfillcolor{currentfill}%
\pgfsetfillopacity{0.130000}%
\pgfsetlinewidth{0.000000pt}%
\definecolor{currentstroke}{rgb}{0.000000,0.000000,0.000000}%
\pgfsetstrokecolor{currentstroke}%
\pgfsetdash{}{0pt}%
\pgfpathmoveto{\pgfqpoint{1.363621in}{0.874036in}}%
\pgfpathlineto{\pgfqpoint{1.381371in}{0.885671in}}%
\pgfpathlineto{\pgfqpoint{1.399201in}{0.897504in}}%
\pgfpathlineto{\pgfqpoint{1.407190in}{0.898918in}}%
\pgfpathlineto{\pgfqpoint{1.415156in}{0.900695in}}%
\pgfpathlineto{\pgfqpoint{1.397352in}{0.888622in}}%
\pgfpathlineto{\pgfqpoint{1.379628in}{0.876730in}}%
\pgfpathlineto{\pgfqpoint{1.371637in}{0.875192in}}%
\pgfpathlineto{\pgfqpoint{1.363621in}{0.874036in}}%
\pgfpathclose%
\pgfusepath{fill}%
\end{pgfscope}%
\begin{pgfscope}%
\pgfpathrectangle{\pgfqpoint{0.407592in}{0.222904in}}{\pgfqpoint{1.411163in}{1.411163in}}%
\pgfusepath{clip}%
\pgfsetbuttcap%
\pgfsetroundjoin%
\definecolor{currentfill}{rgb}{0.556863,0.654902,0.847059}%
\pgfsetfillcolor{currentfill}%
\pgfsetfillopacity{0.130000}%
\pgfsetlinewidth{0.000000pt}%
\definecolor{currentstroke}{rgb}{0.000000,0.000000,0.000000}%
\pgfsetstrokecolor{currentstroke}%
\pgfsetdash{}{0pt}%
\pgfpathmoveto{\pgfqpoint{1.055511in}{0.804534in}}%
\pgfpathlineto{\pgfqpoint{1.072627in}{0.808373in}}%
\pgfpathlineto{\pgfqpoint{1.089802in}{0.812873in}}%
\pgfpathlineto{\pgfqpoint{1.098380in}{0.808029in}}%
\pgfpathlineto{\pgfqpoint{1.106920in}{0.803406in}}%
\pgfpathlineto{\pgfqpoint{1.089791in}{0.798439in}}%
\pgfpathlineto{\pgfqpoint{1.072722in}{0.794145in}}%
\pgfpathlineto{\pgfqpoint{1.064136in}{0.799238in}}%
\pgfpathlineto{\pgfqpoint{1.055511in}{0.804534in}}%
\pgfpathclose%
\pgfusepath{fill}%
\end{pgfscope}%
\begin{pgfscope}%
\pgfpathrectangle{\pgfqpoint{0.407592in}{0.222904in}}{\pgfqpoint{1.411163in}{1.411163in}}%
\pgfusepath{clip}%
\pgfsetbuttcap%
\pgfsetroundjoin%
\definecolor{currentfill}{rgb}{0.556863,0.654902,0.847059}%
\pgfsetfillcolor{currentfill}%
\pgfsetfillopacity{0.130000}%
\pgfsetlinewidth{0.000000pt}%
\definecolor{currentstroke}{rgb}{0.000000,0.000000,0.000000}%
\pgfsetstrokecolor{currentstroke}%
\pgfsetdash{}{0pt}%
\pgfpathmoveto{\pgfqpoint{1.586053in}{1.006306in}}%
\pgfpathlineto{\pgfqpoint{1.604333in}{1.019634in}}%
\pgfpathlineto{\pgfqpoint{1.622697in}{1.033061in}}%
\pgfpathlineto{\pgfqpoint{1.630384in}{1.038231in}}%
\pgfpathlineto{\pgfqpoint{1.638056in}{1.043703in}}%
\pgfpathlineto{\pgfqpoint{1.619717in}{1.030305in}}%
\pgfpathlineto{\pgfqpoint{1.601462in}{1.017006in}}%
\pgfpathlineto{\pgfqpoint{1.593765in}{1.011496in}}%
\pgfpathlineto{\pgfqpoint{1.586053in}{1.006306in}}%
\pgfpathclose%
\pgfusepath{fill}%
\end{pgfscope}%
\begin{pgfscope}%
\pgfpathrectangle{\pgfqpoint{0.407592in}{0.222904in}}{\pgfqpoint{1.411163in}{1.411163in}}%
\pgfusepath{clip}%
\pgfsetbuttcap%
\pgfsetroundjoin%
\definecolor{currentfill}{rgb}{0.556863,0.654902,0.847059}%
\pgfsetfillcolor{currentfill}%
\pgfsetfillopacity{0.130000}%
\pgfsetlinewidth{0.000000pt}%
\definecolor{currentstroke}{rgb}{0.000000,0.000000,0.000000}%
\pgfsetstrokecolor{currentstroke}%
\pgfsetdash{}{0pt}%
\pgfpathmoveto{\pgfqpoint{1.415156in}{0.900695in}}%
\pgfpathlineto{\pgfqpoint{1.433040in}{0.912935in}}%
\pgfpathlineto{\pgfqpoint{1.451006in}{0.925334in}}%
\pgfpathlineto{\pgfqpoint{1.458922in}{0.927665in}}%
\pgfpathlineto{\pgfqpoint{1.466817in}{0.930348in}}%
\pgfpathlineto{\pgfqpoint{1.448876in}{0.917782in}}%
\pgfpathlineto{\pgfqpoint{1.431018in}{0.905360in}}%
\pgfpathlineto{\pgfqpoint{1.423098in}{0.902840in}}%
\pgfpathlineto{\pgfqpoint{1.415156in}{0.900695in}}%
\pgfpathclose%
\pgfusepath{fill}%
\end{pgfscope}%
\begin{pgfscope}%
\pgfpathrectangle{\pgfqpoint{0.407592in}{0.222904in}}{\pgfqpoint{1.411163in}{1.411163in}}%
\pgfusepath{clip}%
\pgfsetbuttcap%
\pgfsetroundjoin%
\definecolor{currentfill}{rgb}{0.556863,0.654902,0.847059}%
\pgfsetfillcolor{currentfill}%
\pgfsetfillopacity{0.130000}%
\pgfsetlinewidth{0.000000pt}%
\definecolor{currentstroke}{rgb}{0.000000,0.000000,0.000000}%
\pgfsetstrokecolor{currentstroke}%
\pgfsetdash{}{0pt}%
\pgfpathmoveto{\pgfqpoint{1.638056in}{1.043703in}}%
\pgfpathlineto{\pgfqpoint{1.656479in}{1.057196in}}%
\pgfpathlineto{\pgfqpoint{1.664122in}{1.062941in}}%
\pgfpathlineto{\pgfqpoint{1.671749in}{1.068975in}}%
\pgfpathlineto{\pgfqpoint{1.653351in}{1.055544in}}%
\pgfpathlineto{\pgfqpoint{1.645711in}{1.049475in}}%
\pgfusepath{fill}%
\end{pgfscope}%
\begin{pgfscope}%
\pgfpathrectangle{\pgfqpoint{0.407592in}{0.222904in}}{\pgfqpoint{1.411163in}{1.411163in}}%
\pgfusepath{clip}%
\pgfsetbuttcap%
\pgfsetroundjoin%
\definecolor{currentfill}{rgb}{0.556863,0.654902,0.847059}%
\pgfsetfillcolor{currentfill}%
\pgfsetfillopacity{0.130000}%
\pgfsetlinewidth{0.000000pt}%
\definecolor{currentstroke}{rgb}{0.000000,0.000000,0.000000}%
\pgfsetstrokecolor{currentstroke}%
\pgfsetdash{}{0pt}%
\pgfpathmoveto{\pgfqpoint{1.003975in}{0.810916in}}%
\pgfpathlineto{\pgfqpoint{1.021031in}{0.812942in}}%
\pgfpathlineto{\pgfqpoint{1.038141in}{0.815682in}}%
\pgfpathlineto{\pgfqpoint{1.046846in}{0.810019in}}%
\pgfpathlineto{\pgfqpoint{1.055511in}{0.804534in}}%
\pgfpathlineto{\pgfqpoint{1.038451in}{0.801395in}}%
\pgfpathlineto{\pgfqpoint{1.021445in}{0.798995in}}%
\pgfpathlineto{\pgfqpoint{1.012730in}{0.804876in}}%
\pgfpathlineto{\pgfqpoint{1.003975in}{0.810916in}}%
\pgfpathclose%
\pgfusepath{fill}%
\end{pgfscope}%
\begin{pgfscope}%
\pgfpathrectangle{\pgfqpoint{0.407592in}{0.222904in}}{\pgfqpoint{1.411163in}{1.411163in}}%
\pgfusepath{clip}%
\pgfsetbuttcap%
\pgfsetroundjoin%
\definecolor{currentfill}{rgb}{0.556863,0.654902,0.847059}%
\pgfsetfillcolor{currentfill}%
\pgfsetfillopacity{0.130000}%
\pgfsetlinewidth{0.000000pt}%
\definecolor{currentstroke}{rgb}{0.000000,0.000000,0.000000}%
\pgfsetstrokecolor{currentstroke}%
\pgfsetdash{}{0pt}%
\pgfpathmoveto{\pgfqpoint{1.466817in}{0.930348in}}%
\pgfpathlineto{\pgfqpoint{1.484839in}{0.943052in}}%
\pgfpathlineto{\pgfqpoint{1.502944in}{0.955886in}}%
\pgfpathlineto{\pgfqpoint{1.510792in}{0.959046in}}%
\pgfpathlineto{\pgfqpoint{1.518620in}{0.962544in}}%
\pgfpathlineto{\pgfqpoint{1.500540in}{0.949606in}}%
\pgfpathlineto{\pgfqpoint{1.482543in}{0.936789in}}%
\pgfpathlineto{\pgfqpoint{1.474690in}{0.933389in}}%
\pgfpathlineto{\pgfqpoint{1.466817in}{0.930348in}}%
\pgfpathclose%
\pgfusepath{fill}%
\end{pgfscope}%
\begin{pgfscope}%
\pgfpathrectangle{\pgfqpoint{0.407592in}{0.222904in}}{\pgfqpoint{1.411163in}{1.411163in}}%
\pgfusepath{clip}%
\pgfsetbuttcap%
\pgfsetroundjoin%
\definecolor{currentfill}{rgb}{0.556863,0.654902,0.847059}%
\pgfsetfillcolor{currentfill}%
\pgfsetfillopacity{0.130000}%
\pgfsetlinewidth{0.000000pt}%
\definecolor{currentstroke}{rgb}{0.000000,0.000000,0.000000}%
\pgfsetstrokecolor{currentstroke}%
\pgfsetdash{}{0pt}%
\pgfpathmoveto{\pgfqpoint{0.952269in}{0.822230in}}%
\pgfpathlineto{\pgfqpoint{0.969281in}{0.822460in}}%
\pgfpathlineto{\pgfqpoint{0.986340in}{0.823431in}}%
\pgfpathlineto{\pgfqpoint{0.995178in}{0.817104in}}%
\pgfpathlineto{\pgfqpoint{1.003975in}{0.810916in}}%
\pgfpathlineto{\pgfqpoint{0.986968in}{0.809636in}}%
\pgfpathlineto{\pgfqpoint{0.970010in}{0.809134in}}%
\pgfpathlineto{\pgfqpoint{0.961161in}{0.815620in}}%
\pgfpathlineto{\pgfqpoint{0.952269in}{0.822230in}}%
\pgfpathclose%
\pgfusepath{fill}%
\end{pgfscope}%
\begin{pgfscope}%
\pgfpathrectangle{\pgfqpoint{0.407592in}{0.222904in}}{\pgfqpoint{1.411163in}{1.411163in}}%
\pgfusepath{clip}%
\pgfsetbuttcap%
\pgfsetroundjoin%
\definecolor{currentfill}{rgb}{0.556863,0.654902,0.847059}%
\pgfsetfillcolor{currentfill}%
\pgfsetfillopacity{0.130000}%
\pgfsetlinewidth{0.000000pt}%
\definecolor{currentstroke}{rgb}{0.000000,0.000000,0.000000}%
\pgfsetstrokecolor{currentstroke}%
\pgfsetdash{}{0pt}%
\pgfpathmoveto{\pgfqpoint{1.518620in}{0.962544in}}%
\pgfpathlineto{\pgfqpoint{1.536783in}{0.975599in}}%
\pgfpathlineto{\pgfqpoint{1.555029in}{0.988766in}}%
\pgfpathlineto{\pgfqpoint{1.562812in}{0.992667in}}%
\pgfpathlineto{\pgfqpoint{1.570577in}{0.996891in}}%
\pgfpathlineto{\pgfqpoint{1.552356in}{0.983674in}}%
\pgfpathlineto{\pgfqpoint{1.534218in}{0.970564in}}%
\pgfpathlineto{\pgfqpoint{1.526428in}{0.966383in}}%
\pgfpathlineto{\pgfqpoint{1.518620in}{0.962544in}}%
\pgfpathclose%
\pgfusepath{fill}%
\end{pgfscope}%
\begin{pgfscope}%
\pgfpathrectangle{\pgfqpoint{0.407592in}{0.222904in}}{\pgfqpoint{1.411163in}{1.411163in}}%
\pgfusepath{clip}%
\pgfsetbuttcap%
\pgfsetroundjoin%
\definecolor{currentfill}{rgb}{0.556863,0.654902,0.847059}%
\pgfsetfillcolor{currentfill}%
\pgfsetfillopacity{0.130000}%
\pgfsetlinewidth{0.000000pt}%
\definecolor{currentstroke}{rgb}{0.000000,0.000000,0.000000}%
\pgfsetstrokecolor{currentstroke}%
\pgfsetdash{}{0pt}%
\pgfpathmoveto{\pgfqpoint{1.192878in}{0.822594in}}%
\pgfpathlineto{\pgfqpoint{1.210299in}{0.830765in}}%
\pgfpathlineto{\pgfqpoint{1.227792in}{0.839360in}}%
\pgfpathlineto{\pgfqpoint{1.236103in}{0.836987in}}%
\pgfpathlineto{\pgfqpoint{1.244381in}{0.834920in}}%
\pgfpathlineto{\pgfqpoint{1.226924in}{0.825864in}}%
\pgfpathlineto{\pgfqpoint{1.209540in}{0.817214in}}%
\pgfpathlineto{\pgfqpoint{1.201226in}{0.819751in}}%
\pgfpathlineto{\pgfqpoint{1.192878in}{0.822594in}}%
\pgfpathclose%
\pgfusepath{fill}%
\end{pgfscope}%
\begin{pgfscope}%
\pgfpathrectangle{\pgfqpoint{0.407592in}{0.222904in}}{\pgfqpoint{1.411163in}{1.411163in}}%
\pgfusepath{clip}%
\pgfsetbuttcap%
\pgfsetroundjoin%
\definecolor{currentfill}{rgb}{0.556863,0.654902,0.847059}%
\pgfsetfillcolor{currentfill}%
\pgfsetfillopacity{0.130000}%
\pgfsetlinewidth{0.000000pt}%
\definecolor{currentstroke}{rgb}{0.000000,0.000000,0.000000}%
\pgfsetstrokecolor{currentstroke}%
\pgfsetdash{}{0pt}%
\pgfpathmoveto{\pgfqpoint{1.244381in}{0.834920in}}%
\pgfpathlineto{\pgfqpoint{1.261912in}{0.844349in}}%
\pgfpathlineto{\pgfqpoint{1.279519in}{0.854121in}}%
\pgfpathlineto{\pgfqpoint{1.287732in}{0.852803in}}%
\pgfpathlineto{\pgfqpoint{1.295914in}{0.851815in}}%
\pgfpathlineto{\pgfqpoint{1.278340in}{0.841630in}}%
\pgfpathlineto{\pgfqpoint{1.260842in}{0.831767in}}%
\pgfpathlineto{\pgfqpoint{1.252627in}{0.833175in}}%
\pgfpathlineto{\pgfqpoint{1.244381in}{0.834920in}}%
\pgfpathclose%
\pgfusepath{fill}%
\end{pgfscope}%
\begin{pgfscope}%
\pgfpathrectangle{\pgfqpoint{0.407592in}{0.222904in}}{\pgfqpoint{1.411163in}{1.411163in}}%
\pgfusepath{clip}%
\pgfsetbuttcap%
\pgfsetroundjoin%
\definecolor{currentfill}{rgb}{0.556863,0.654902,0.847059}%
\pgfsetfillcolor{currentfill}%
\pgfsetfillopacity{0.130000}%
\pgfsetlinewidth{0.000000pt}%
\definecolor{currentstroke}{rgb}{0.000000,0.000000,0.000000}%
\pgfsetstrokecolor{currentstroke}%
\pgfsetdash{}{0pt}%
\pgfpathmoveto{\pgfqpoint{1.141367in}{0.815180in}}%
\pgfpathlineto{\pgfqpoint{1.158688in}{0.821904in}}%
\pgfpathlineto{\pgfqpoint{1.176079in}{0.829133in}}%
\pgfpathlineto{\pgfqpoint{1.184496in}{0.825726in}}%
\pgfpathlineto{\pgfqpoint{1.192878in}{0.822594in}}%
\pgfpathlineto{\pgfqpoint{1.175528in}{0.814880in}}%
\pgfpathlineto{\pgfqpoint{1.158247in}{0.807663in}}%
\pgfpathlineto{\pgfqpoint{1.149825in}{0.811288in}}%
\pgfpathlineto{\pgfqpoint{1.141367in}{0.815180in}}%
\pgfpathclose%
\pgfusepath{fill}%
\end{pgfscope}%
\begin{pgfscope}%
\pgfpathrectangle{\pgfqpoint{0.407592in}{0.222904in}}{\pgfqpoint{1.411163in}{1.411163in}}%
\pgfusepath{clip}%
\pgfsetbuttcap%
\pgfsetroundjoin%
\definecolor{currentfill}{rgb}{0.556863,0.654902,0.847059}%
\pgfsetfillcolor{currentfill}%
\pgfsetfillopacity{0.130000}%
\pgfsetlinewidth{0.000000pt}%
\definecolor{currentstroke}{rgb}{0.000000,0.000000,0.000000}%
\pgfsetstrokecolor{currentstroke}%
\pgfsetdash{}{0pt}%
\pgfpathmoveto{\pgfqpoint{1.295914in}{0.851815in}}%
\pgfpathlineto{\pgfqpoint{1.313565in}{0.862297in}}%
\pgfpathlineto{\pgfqpoint{1.331294in}{0.873053in}}%
\pgfpathlineto{\pgfqpoint{1.339417in}{0.872774in}}%
\pgfpathlineto{\pgfqpoint{1.347511in}{0.872838in}}%
\pgfpathlineto{\pgfqpoint{1.329812in}{0.861734in}}%
\pgfpathlineto{\pgfqpoint{1.312192in}{0.850882in}}%
\pgfpathlineto{\pgfqpoint{1.304067in}{0.851171in}}%
\pgfpathlineto{\pgfqpoint{1.295914in}{0.851815in}}%
\pgfpathclose%
\pgfusepath{fill}%
\end{pgfscope}%
\begin{pgfscope}%
\pgfpathrectangle{\pgfqpoint{0.407592in}{0.222904in}}{\pgfqpoint{1.411163in}{1.411163in}}%
\pgfusepath{clip}%
\pgfsetbuttcap%
\pgfsetroundjoin%
\definecolor{currentfill}{rgb}{0.556863,0.654902,0.847059}%
\pgfsetfillcolor{currentfill}%
\pgfsetfillopacity{0.130000}%
\pgfsetlinewidth{0.000000pt}%
\definecolor{currentstroke}{rgb}{0.000000,0.000000,0.000000}%
\pgfsetstrokecolor{currentstroke}%
\pgfsetdash{}{0pt}%
\pgfpathmoveto{\pgfqpoint{1.347511in}{0.872838in}}%
\pgfpathlineto{\pgfqpoint{1.365290in}{0.884175in}}%
\pgfpathlineto{\pgfqpoint{1.383148in}{0.895730in}}%
\pgfpathlineto{\pgfqpoint{1.391187in}{0.896443in}}%
\pgfpathlineto{\pgfqpoint{1.399201in}{0.897504in}}%
\pgfpathlineto{\pgfqpoint{1.381371in}{0.885671in}}%
\pgfpathlineto{\pgfqpoint{1.363621in}{0.874036in}}%
\pgfpathlineto{\pgfqpoint{1.355579in}{0.873255in}}%
\pgfpathlineto{\pgfqpoint{1.347511in}{0.872838in}}%
\pgfpathclose%
\pgfusepath{fill}%
\end{pgfscope}%
\begin{pgfscope}%
\pgfpathrectangle{\pgfqpoint{0.407592in}{0.222904in}}{\pgfqpoint{1.411163in}{1.411163in}}%
\pgfusepath{clip}%
\pgfsetbuttcap%
\pgfsetroundjoin%
\definecolor{currentfill}{rgb}{0.556863,0.654902,0.847059}%
\pgfsetfillcolor{currentfill}%
\pgfsetfillopacity{0.130000}%
\pgfsetlinewidth{0.000000pt}%
\definecolor{currentstroke}{rgb}{0.000000,0.000000,0.000000}%
\pgfsetstrokecolor{currentstroke}%
\pgfsetdash{}{0pt}%
\pgfpathmoveto{\pgfqpoint{1.089802in}{0.812873in}}%
\pgfpathlineto{\pgfqpoint{1.107039in}{0.817994in}}%
\pgfpathlineto{\pgfqpoint{1.124339in}{0.823699in}}%
\pgfpathlineto{\pgfqpoint{1.132871in}{0.819322in}}%
\pgfpathlineto{\pgfqpoint{1.141367in}{0.815180in}}%
\pgfpathlineto{\pgfqpoint{1.124111in}{0.809000in}}%
\pgfpathlineto{\pgfqpoint{1.106920in}{0.803406in}}%
\pgfpathlineto{\pgfqpoint{1.098380in}{0.808029in}}%
\pgfpathlineto{\pgfqpoint{1.089802in}{0.812873in}}%
\pgfpathclose%
\pgfusepath{fill}%
\end{pgfscope}%
\begin{pgfscope}%
\pgfpathrectangle{\pgfqpoint{0.407592in}{0.222904in}}{\pgfqpoint{1.411163in}{1.411163in}}%
\pgfusepath{clip}%
\pgfsetbuttcap%
\pgfsetroundjoin%
\definecolor{currentfill}{rgb}{0.556863,0.654902,0.847059}%
\pgfsetfillcolor{currentfill}%
\pgfsetfillopacity{0.130000}%
\pgfsetlinewidth{0.000000pt}%
\definecolor{currentstroke}{rgb}{0.000000,0.000000,0.000000}%
\pgfsetstrokecolor{currentstroke}%
\pgfsetdash{}{0pt}%
\pgfpathmoveto{\pgfqpoint{1.570577in}{0.996891in}}%
\pgfpathlineto{\pgfqpoint{1.588882in}{1.010212in}}%
\pgfpathlineto{\pgfqpoint{1.607272in}{1.023632in}}%
\pgfpathlineto{\pgfqpoint{1.614993in}{1.028195in}}%
\pgfpathlineto{\pgfqpoint{1.622697in}{1.033061in}}%
\pgfpathlineto{\pgfqpoint{1.604333in}{1.019634in}}%
\pgfpathlineto{\pgfqpoint{1.586053in}{1.006306in}}%
\pgfpathlineto{\pgfqpoint{1.578324in}{1.001437in}}%
\pgfpathlineto{\pgfqpoint{1.570577in}{0.996891in}}%
\pgfpathclose%
\pgfusepath{fill}%
\end{pgfscope}%
\begin{pgfscope}%
\pgfpathrectangle{\pgfqpoint{0.407592in}{0.222904in}}{\pgfqpoint{1.411163in}{1.411163in}}%
\pgfusepath{clip}%
\pgfsetbuttcap%
\pgfsetroundjoin%
\definecolor{currentfill}{rgb}{0.556863,0.654902,0.847059}%
\pgfsetfillcolor{currentfill}%
\pgfsetfillopacity{0.130000}%
\pgfsetlinewidth{0.000000pt}%
\definecolor{currentstroke}{rgb}{0.000000,0.000000,0.000000}%
\pgfsetstrokecolor{currentstroke}%
\pgfsetdash{}{0pt}%
\pgfpathmoveto{\pgfqpoint{1.622697in}{1.033061in}}%
\pgfpathlineto{\pgfqpoint{1.641146in}{1.046584in}}%
\pgfpathlineto{\pgfqpoint{1.648821in}{1.051743in}}%
\pgfpathlineto{\pgfqpoint{1.656479in}{1.057196in}}%
\pgfpathlineto{\pgfqpoint{1.638056in}{1.043703in}}%
\pgfpathlineto{\pgfqpoint{1.630384in}{1.038231in}}%
\pgfusepath{fill}%
\end{pgfscope}%
\begin{pgfscope}%
\pgfpathrectangle{\pgfqpoint{0.407592in}{0.222904in}}{\pgfqpoint{1.411163in}{1.411163in}}%
\pgfusepath{clip}%
\pgfsetbuttcap%
\pgfsetroundjoin%
\definecolor{currentfill}{rgb}{0.556863,0.654902,0.847059}%
\pgfsetfillcolor{currentfill}%
\pgfsetfillopacity{0.130000}%
\pgfsetlinewidth{0.000000pt}%
\definecolor{currentstroke}{rgb}{0.000000,0.000000,0.000000}%
\pgfsetstrokecolor{currentstroke}%
\pgfsetdash{}{0pt}%
\pgfpathmoveto{\pgfqpoint{1.399201in}{0.897504in}}%
\pgfpathlineto{\pgfqpoint{1.417112in}{0.909521in}}%
\pgfpathlineto{\pgfqpoint{1.435105in}{0.921711in}}%
\pgfpathlineto{\pgfqpoint{1.443067in}{0.923352in}}%
\pgfpathlineto{\pgfqpoint{1.451006in}{0.925334in}}%
\pgfpathlineto{\pgfqpoint{1.433040in}{0.912935in}}%
\pgfpathlineto{\pgfqpoint{1.415156in}{0.900695in}}%
\pgfpathlineto{\pgfqpoint{1.407190in}{0.898918in}}%
\pgfpathlineto{\pgfqpoint{1.399201in}{0.897504in}}%
\pgfpathclose%
\pgfusepath{fill}%
\end{pgfscope}%
\begin{pgfscope}%
\pgfpathrectangle{\pgfqpoint{0.407592in}{0.222904in}}{\pgfqpoint{1.411163in}{1.411163in}}%
\pgfusepath{clip}%
\pgfsetbuttcap%
\pgfsetroundjoin%
\definecolor{currentfill}{rgb}{0.556863,0.654902,0.847059}%
\pgfsetfillcolor{currentfill}%
\pgfsetfillopacity{0.130000}%
\pgfsetlinewidth{0.000000pt}%
\definecolor{currentstroke}{rgb}{0.000000,0.000000,0.000000}%
\pgfsetstrokecolor{currentstroke}%
\pgfsetdash{}{0pt}%
\pgfpathmoveto{\pgfqpoint{1.038141in}{0.815682in}}%
\pgfpathlineto{\pgfqpoint{1.055307in}{0.819102in}}%
\pgfpathlineto{\pgfqpoint{1.072530in}{0.823169in}}%
\pgfpathlineto{\pgfqpoint{1.081186in}{0.817924in}}%
\pgfpathlineto{\pgfqpoint{1.089802in}{0.812873in}}%
\pgfpathlineto{\pgfqpoint{1.072627in}{0.808373in}}%
\pgfpathlineto{\pgfqpoint{1.055511in}{0.804534in}}%
\pgfpathlineto{\pgfqpoint{1.046846in}{0.810019in}}%
\pgfpathlineto{\pgfqpoint{1.038141in}{0.815682in}}%
\pgfpathclose%
\pgfusepath{fill}%
\end{pgfscope}%
\begin{pgfscope}%
\pgfpathrectangle{\pgfqpoint{0.407592in}{0.222904in}}{\pgfqpoint{1.411163in}{1.411163in}}%
\pgfusepath{clip}%
\pgfsetbuttcap%
\pgfsetroundjoin%
\definecolor{currentfill}{rgb}{0.556863,0.654902,0.847059}%
\pgfsetfillcolor{currentfill}%
\pgfsetfillopacity{0.130000}%
\pgfsetlinewidth{0.000000pt}%
\definecolor{currentstroke}{rgb}{0.000000,0.000000,0.000000}%
\pgfsetstrokecolor{currentstroke}%
\pgfsetdash{}{0pt}%
\pgfpathmoveto{\pgfqpoint{1.451006in}{0.925334in}}%
\pgfpathlineto{\pgfqpoint{1.469054in}{0.937883in}}%
\pgfpathlineto{\pgfqpoint{1.487185in}{0.950572in}}%
\pgfpathlineto{\pgfqpoint{1.495075in}{0.953063in}}%
\pgfpathlineto{\pgfqpoint{1.502944in}{0.955886in}}%
\pgfpathlineto{\pgfqpoint{1.484839in}{0.943052in}}%
\pgfpathlineto{\pgfqpoint{1.466817in}{0.930348in}}%
\pgfpathlineto{\pgfqpoint{1.458922in}{0.927665in}}%
\pgfpathlineto{\pgfqpoint{1.451006in}{0.925334in}}%
\pgfpathclose%
\pgfusepath{fill}%
\end{pgfscope}%
\begin{pgfscope}%
\pgfpathrectangle{\pgfqpoint{0.407592in}{0.222904in}}{\pgfqpoint{1.411163in}{1.411163in}}%
\pgfusepath{clip}%
\pgfsetbuttcap%
\pgfsetroundjoin%
\definecolor{currentfill}{rgb}{0.556863,0.654902,0.847059}%
\pgfsetfillcolor{currentfill}%
\pgfsetfillopacity{0.130000}%
\pgfsetlinewidth{0.000000pt}%
\definecolor{currentstroke}{rgb}{0.000000,0.000000,0.000000}%
\pgfsetstrokecolor{currentstroke}%
\pgfsetdash{}{0pt}%
\pgfpathmoveto{\pgfqpoint{0.986340in}{0.823431in}}%
\pgfpathlineto{\pgfqpoint{1.003449in}{0.825118in}}%
\pgfpathlineto{\pgfqpoint{1.020610in}{0.827494in}}%
\pgfpathlineto{\pgfqpoint{1.029396in}{0.821510in}}%
\pgfpathlineto{\pgfqpoint{1.038141in}{0.815682in}}%
\pgfpathlineto{\pgfqpoint{1.021031in}{0.812942in}}%
\pgfpathlineto{\pgfqpoint{1.003975in}{0.810916in}}%
\pgfpathlineto{\pgfqpoint{0.995178in}{0.817104in}}%
\pgfpathlineto{\pgfqpoint{0.986340in}{0.823431in}}%
\pgfpathclose%
\pgfusepath{fill}%
\end{pgfscope}%
\begin{pgfscope}%
\pgfpathrectangle{\pgfqpoint{0.407592in}{0.222904in}}{\pgfqpoint{1.411163in}{1.411163in}}%
\pgfusepath{clip}%
\pgfsetbuttcap%
\pgfsetroundjoin%
\definecolor{currentfill}{rgb}{0.556863,0.654902,0.847059}%
\pgfsetfillcolor{currentfill}%
\pgfsetfillopacity{0.130000}%
\pgfsetlinewidth{0.000000pt}%
\definecolor{currentstroke}{rgb}{0.000000,0.000000,0.000000}%
\pgfsetstrokecolor{currentstroke}%
\pgfsetdash{}{0pt}%
\pgfpathmoveto{\pgfqpoint{1.502944in}{0.955886in}}%
\pgfpathlineto{\pgfqpoint{1.521132in}{0.968845in}}%
\pgfpathlineto{\pgfqpoint{1.539404in}{0.981923in}}%
\pgfpathlineto{\pgfqpoint{1.547227in}{0.985185in}}%
\pgfpathlineto{\pgfqpoint{1.555029in}{0.988766in}}%
\pgfpathlineto{\pgfqpoint{1.536783in}{0.975599in}}%
\pgfpathlineto{\pgfqpoint{1.518620in}{0.962544in}}%
\pgfpathlineto{\pgfqpoint{1.510792in}{0.959046in}}%
\pgfpathlineto{\pgfqpoint{1.502944in}{0.955886in}}%
\pgfpathclose%
\pgfusepath{fill}%
\end{pgfscope}%
\begin{pgfscope}%
\pgfpathrectangle{\pgfqpoint{0.407592in}{0.222904in}}{\pgfqpoint{1.411163in}{1.411163in}}%
\pgfusepath{clip}%
\pgfsetbuttcap%
\pgfsetroundjoin%
\definecolor{currentfill}{rgb}{0.556863,0.654902,0.847059}%
\pgfsetfillcolor{currentfill}%
\pgfsetfillopacity{0.130000}%
\pgfsetlinewidth{0.000000pt}%
\definecolor{currentstroke}{rgb}{0.000000,0.000000,0.000000}%
\pgfsetstrokecolor{currentstroke}%
\pgfsetdash{}{0pt}%
\pgfpathmoveto{\pgfqpoint{1.227792in}{0.839360in}}%
\pgfpathlineto{\pgfqpoint{1.245359in}{0.848346in}}%
\pgfpathlineto{\pgfqpoint{1.263000in}{0.857695in}}%
\pgfpathlineto{\pgfqpoint{1.271275in}{0.855756in}}%
\pgfpathlineto{\pgfqpoint{1.279519in}{0.854121in}}%
\pgfpathlineto{\pgfqpoint{1.261912in}{0.844349in}}%
\pgfpathlineto{\pgfqpoint{1.244381in}{0.834920in}}%
\pgfpathlineto{\pgfqpoint{1.236103in}{0.836987in}}%
\pgfpathlineto{\pgfqpoint{1.227792in}{0.839360in}}%
\pgfpathclose%
\pgfusepath{fill}%
\end{pgfscope}%
\begin{pgfscope}%
\pgfpathrectangle{\pgfqpoint{0.407592in}{0.222904in}}{\pgfqpoint{1.411163in}{1.411163in}}%
\pgfusepath{clip}%
\pgfsetbuttcap%
\pgfsetroundjoin%
\definecolor{currentfill}{rgb}{0.556863,0.654902,0.847059}%
\pgfsetfillcolor{currentfill}%
\pgfsetfillopacity{0.130000}%
\pgfsetlinewidth{0.000000pt}%
\definecolor{currentstroke}{rgb}{0.000000,0.000000,0.000000}%
\pgfsetstrokecolor{currentstroke}%
\pgfsetdash{}{0pt}%
\pgfpathmoveto{\pgfqpoint{1.176079in}{0.829133in}}%
\pgfpathlineto{\pgfqpoint{1.193538in}{0.836833in}}%
\pgfpathlineto{\pgfqpoint{1.211070in}{0.844971in}}%
\pgfpathlineto{\pgfqpoint{1.219448in}{0.842026in}}%
\pgfpathlineto{\pgfqpoint{1.227792in}{0.839360in}}%
\pgfpathlineto{\pgfqpoint{1.210299in}{0.830765in}}%
\pgfpathlineto{\pgfqpoint{1.192878in}{0.822594in}}%
\pgfpathlineto{\pgfqpoint{1.184496in}{0.825726in}}%
\pgfpathlineto{\pgfqpoint{1.176079in}{0.829133in}}%
\pgfpathclose%
\pgfusepath{fill}%
\end{pgfscope}%
\begin{pgfscope}%
\pgfpathrectangle{\pgfqpoint{0.407592in}{0.222904in}}{\pgfqpoint{1.411163in}{1.411163in}}%
\pgfusepath{clip}%
\pgfsetbuttcap%
\pgfsetroundjoin%
\definecolor{currentfill}{rgb}{0.556863,0.654902,0.847059}%
\pgfsetfillcolor{currentfill}%
\pgfsetfillopacity{0.130000}%
\pgfsetlinewidth{0.000000pt}%
\definecolor{currentstroke}{rgb}{0.000000,0.000000,0.000000}%
\pgfsetstrokecolor{currentstroke}%
\pgfsetdash{}{0pt}%
\pgfpathmoveto{\pgfqpoint{1.279519in}{0.854121in}}%
\pgfpathlineto{\pgfqpoint{1.297202in}{0.864211in}}%
\pgfpathlineto{\pgfqpoint{1.314963in}{0.874596in}}%
\pgfpathlineto{\pgfqpoint{1.323143in}{0.873664in}}%
\pgfpathlineto{\pgfqpoint{1.331294in}{0.873053in}}%
\pgfpathlineto{\pgfqpoint{1.313565in}{0.862297in}}%
\pgfpathlineto{\pgfqpoint{1.295914in}{0.851815in}}%
\pgfpathlineto{\pgfqpoint{1.287732in}{0.852803in}}%
\pgfpathlineto{\pgfqpoint{1.279519in}{0.854121in}}%
\pgfpathclose%
\pgfusepath{fill}%
\end{pgfscope}%
\begin{pgfscope}%
\pgfpathrectangle{\pgfqpoint{0.407592in}{0.222904in}}{\pgfqpoint{1.411163in}{1.411163in}}%
\pgfusepath{clip}%
\pgfsetbuttcap%
\pgfsetroundjoin%
\definecolor{currentfill}{rgb}{0.556863,0.654902,0.847059}%
\pgfsetfillcolor{currentfill}%
\pgfsetfillopacity{0.130000}%
\pgfsetlinewidth{0.000000pt}%
\definecolor{currentstroke}{rgb}{0.000000,0.000000,0.000000}%
\pgfsetstrokecolor{currentstroke}%
\pgfsetdash{}{0pt}%
\pgfpathmoveto{\pgfqpoint{1.555029in}{0.988766in}}%
\pgfpathlineto{\pgfqpoint{1.573360in}{1.002040in}}%
\pgfpathlineto{\pgfqpoint{1.591775in}{1.015419in}}%
\pgfpathlineto{\pgfqpoint{1.599532in}{1.019374in}}%
\pgfpathlineto{\pgfqpoint{1.607272in}{1.023632in}}%
\pgfpathlineto{\pgfqpoint{1.588882in}{1.010212in}}%
\pgfpathlineto{\pgfqpoint{1.570577in}{0.996891in}}%
\pgfpathlineto{\pgfqpoint{1.562812in}{0.992667in}}%
\pgfpathlineto{\pgfqpoint{1.555029in}{0.988766in}}%
\pgfpathclose%
\pgfusepath{fill}%
\end{pgfscope}%
\begin{pgfscope}%
\pgfpathrectangle{\pgfqpoint{0.407592in}{0.222904in}}{\pgfqpoint{1.411163in}{1.411163in}}%
\pgfusepath{clip}%
\pgfsetbuttcap%
\pgfsetroundjoin%
\definecolor{currentfill}{rgb}{0.556863,0.654902,0.847059}%
\pgfsetfillcolor{currentfill}%
\pgfsetfillopacity{0.130000}%
\pgfsetlinewidth{0.000000pt}%
\definecolor{currentstroke}{rgb}{0.000000,0.000000,0.000000}%
\pgfsetstrokecolor{currentstroke}%
\pgfsetdash{}{0pt}%
\pgfpathmoveto{\pgfqpoint{1.607272in}{1.023632in}}%
\pgfpathlineto{\pgfqpoint{1.625746in}{1.037151in}}%
\pgfpathlineto{\pgfqpoint{1.633455in}{1.041720in}}%
\pgfpathlineto{\pgfqpoint{1.641146in}{1.046584in}}%
\pgfpathlineto{\pgfqpoint{1.622697in}{1.033061in}}%
\pgfpathlineto{\pgfqpoint{1.614993in}{1.028195in}}%
\pgfusepath{fill}%
\end{pgfscope}%
\begin{pgfscope}%
\pgfpathrectangle{\pgfqpoint{0.407592in}{0.222904in}}{\pgfqpoint{1.411163in}{1.411163in}}%
\pgfusepath{clip}%
\pgfsetbuttcap%
\pgfsetroundjoin%
\definecolor{currentfill}{rgb}{0.556863,0.654902,0.847059}%
\pgfsetfillcolor{currentfill}%
\pgfsetfillopacity{0.130000}%
\pgfsetlinewidth{0.000000pt}%
\definecolor{currentstroke}{rgb}{0.000000,0.000000,0.000000}%
\pgfsetstrokecolor{currentstroke}%
\pgfsetdash{}{0pt}%
\pgfpathmoveto{\pgfqpoint{1.124339in}{0.823699in}}%
\pgfpathlineto{\pgfqpoint{1.141704in}{0.829950in}}%
\pgfpathlineto{\pgfqpoint{1.159136in}{0.836713in}}%
\pgfpathlineto{\pgfqpoint{1.167625in}{0.832800in}}%
\pgfpathlineto{\pgfqpoint{1.176079in}{0.829133in}}%
\pgfpathlineto{\pgfqpoint{1.158688in}{0.821904in}}%
\pgfpathlineto{\pgfqpoint{1.141367in}{0.815180in}}%
\pgfpathlineto{\pgfqpoint{1.132871in}{0.819322in}}%
\pgfpathlineto{\pgfqpoint{1.124339in}{0.823699in}}%
\pgfpathclose%
\pgfusepath{fill}%
\end{pgfscope}%
\begin{pgfscope}%
\pgfpathrectangle{\pgfqpoint{0.407592in}{0.222904in}}{\pgfqpoint{1.411163in}{1.411163in}}%
\pgfusepath{clip}%
\pgfsetbuttcap%
\pgfsetroundjoin%
\definecolor{currentfill}{rgb}{0.556863,0.654902,0.847059}%
\pgfsetfillcolor{currentfill}%
\pgfsetfillopacity{0.130000}%
\pgfsetlinewidth{0.000000pt}%
\definecolor{currentstroke}{rgb}{0.000000,0.000000,0.000000}%
\pgfsetstrokecolor{currentstroke}%
\pgfsetdash{}{0pt}%
\pgfpathmoveto{\pgfqpoint{1.331294in}{0.873053in}}%
\pgfpathlineto{\pgfqpoint{1.349102in}{0.884062in}}%
\pgfpathlineto{\pgfqpoint{1.366990in}{0.895308in}}%
\pgfpathlineto{\pgfqpoint{1.375083in}{0.895354in}}%
\pgfpathlineto{\pgfqpoint{1.383148in}{0.895730in}}%
\pgfpathlineto{\pgfqpoint{1.365290in}{0.884175in}}%
\pgfpathlineto{\pgfqpoint{1.347511in}{0.872838in}}%
\pgfpathlineto{\pgfqpoint{1.339417in}{0.872774in}}%
\pgfpathlineto{\pgfqpoint{1.331294in}{0.873053in}}%
\pgfpathclose%
\pgfusepath{fill}%
\end{pgfscope}%
\begin{pgfscope}%
\pgfpathrectangle{\pgfqpoint{0.407592in}{0.222904in}}{\pgfqpoint{1.411163in}{1.411163in}}%
\pgfusepath{clip}%
\pgfsetbuttcap%
\pgfsetroundjoin%
\definecolor{currentfill}{rgb}{0.556863,0.654902,0.847059}%
\pgfsetfillcolor{currentfill}%
\pgfsetfillopacity{0.130000}%
\pgfsetlinewidth{0.000000pt}%
\definecolor{currentstroke}{rgb}{0.000000,0.000000,0.000000}%
\pgfsetstrokecolor{currentstroke}%
\pgfsetdash{}{0pt}%
\pgfpathmoveto{\pgfqpoint{1.383148in}{0.895730in}}%
\pgfpathlineto{\pgfqpoint{1.401087in}{0.907486in}}%
\pgfpathlineto{\pgfqpoint{1.419108in}{0.919433in}}%
\pgfpathlineto{\pgfqpoint{1.427119in}{0.920407in}}%
\pgfpathlineto{\pgfqpoint{1.435105in}{0.921711in}}%
\pgfpathlineto{\pgfqpoint{1.417112in}{0.909521in}}%
\pgfpathlineto{\pgfqpoint{1.399201in}{0.897504in}}%
\pgfpathlineto{\pgfqpoint{1.391187in}{0.896443in}}%
\pgfpathlineto{\pgfqpoint{1.383148in}{0.895730in}}%
\pgfpathclose%
\pgfusepath{fill}%
\end{pgfscope}%
\begin{pgfscope}%
\pgfpathrectangle{\pgfqpoint{0.407592in}{0.222904in}}{\pgfqpoint{1.411163in}{1.411163in}}%
\pgfusepath{clip}%
\pgfsetbuttcap%
\pgfsetroundjoin%
\definecolor{currentfill}{rgb}{0.556863,0.654902,0.847059}%
\pgfsetfillcolor{currentfill}%
\pgfsetfillopacity{0.130000}%
\pgfsetlinewidth{0.000000pt}%
\definecolor{currentstroke}{rgb}{0.000000,0.000000,0.000000}%
\pgfsetstrokecolor{currentstroke}%
\pgfsetdash{}{0pt}%
\pgfpathmoveto{\pgfqpoint{1.072530in}{0.823169in}}%
\pgfpathlineto{\pgfqpoint{1.089813in}{0.827849in}}%
\pgfpathlineto{\pgfqpoint{1.107159in}{0.833106in}}%
\pgfpathlineto{\pgfqpoint{1.115768in}{0.828298in}}%
\pgfpathlineto{\pgfqpoint{1.124339in}{0.823699in}}%
\pgfpathlineto{\pgfqpoint{1.107039in}{0.817994in}}%
\pgfpathlineto{\pgfqpoint{1.089802in}{0.812873in}}%
\pgfpathlineto{\pgfqpoint{1.081186in}{0.817924in}}%
\pgfpathlineto{\pgfqpoint{1.072530in}{0.823169in}}%
\pgfpathclose%
\pgfusepath{fill}%
\end{pgfscope}%
\begin{pgfscope}%
\pgfpathrectangle{\pgfqpoint{0.407592in}{0.222904in}}{\pgfqpoint{1.411163in}{1.411163in}}%
\pgfusepath{clip}%
\pgfsetbuttcap%
\pgfsetroundjoin%
\definecolor{currentfill}{rgb}{0.556863,0.654902,0.847059}%
\pgfsetfillcolor{currentfill}%
\pgfsetfillopacity{0.130000}%
\pgfsetlinewidth{0.000000pt}%
\definecolor{currentstroke}{rgb}{0.000000,0.000000,0.000000}%
\pgfsetstrokecolor{currentstroke}%
\pgfsetdash{}{0pt}%
\pgfpathmoveto{\pgfqpoint{1.435105in}{0.921711in}}%
\pgfpathlineto{\pgfqpoint{1.453180in}{0.934065in}}%
\pgfpathlineto{\pgfqpoint{1.471338in}{0.946572in}}%
\pgfpathlineto{\pgfqpoint{1.479273in}{0.948410in}}%
\pgfpathlineto{\pgfqpoint{1.487185in}{0.950572in}}%
\pgfpathlineto{\pgfqpoint{1.469054in}{0.937883in}}%
\pgfpathlineto{\pgfqpoint{1.451006in}{0.925334in}}%
\pgfpathlineto{\pgfqpoint{1.443067in}{0.923352in}}%
\pgfpathlineto{\pgfqpoint{1.435105in}{0.921711in}}%
\pgfpathclose%
\pgfusepath{fill}%
\end{pgfscope}%
\begin{pgfscope}%
\pgfpathrectangle{\pgfqpoint{0.407592in}{0.222904in}}{\pgfqpoint{1.411163in}{1.411163in}}%
\pgfusepath{clip}%
\pgfsetbuttcap%
\pgfsetroundjoin%
\definecolor{currentfill}{rgb}{0.556863,0.654902,0.847059}%
\pgfsetfillcolor{currentfill}%
\pgfsetfillopacity{0.130000}%
\pgfsetlinewidth{0.000000pt}%
\definecolor{currentstroke}{rgb}{0.000000,0.000000,0.000000}%
\pgfsetstrokecolor{currentstroke}%
\pgfsetdash{}{0pt}%
\pgfpathmoveto{\pgfqpoint{1.020610in}{0.827494in}}%
\pgfpathlineto{\pgfqpoint{1.037826in}{0.830530in}}%
\pgfpathlineto{\pgfqpoint{1.055099in}{0.834198in}}%
\pgfpathlineto{\pgfqpoint{1.063834in}{0.828598in}}%
\pgfpathlineto{\pgfqpoint{1.072530in}{0.823169in}}%
\pgfpathlineto{\pgfqpoint{1.055307in}{0.819102in}}%
\pgfpathlineto{\pgfqpoint{1.038141in}{0.815682in}}%
\pgfpathlineto{\pgfqpoint{1.029396in}{0.821510in}}%
\pgfpathlineto{\pgfqpoint{1.020610in}{0.827494in}}%
\pgfpathclose%
\pgfusepath{fill}%
\end{pgfscope}%
\begin{pgfscope}%
\pgfpathrectangle{\pgfqpoint{0.407592in}{0.222904in}}{\pgfqpoint{1.411163in}{1.411163in}}%
\pgfusepath{clip}%
\pgfsetbuttcap%
\pgfsetroundjoin%
\definecolor{currentfill}{rgb}{0.556863,0.654902,0.847059}%
\pgfsetfillcolor{currentfill}%
\pgfsetfillopacity{0.130000}%
\pgfsetlinewidth{0.000000pt}%
\definecolor{currentstroke}{rgb}{0.000000,0.000000,0.000000}%
\pgfsetstrokecolor{currentstroke}%
\pgfsetdash{}{0pt}%
\pgfpathmoveto{\pgfqpoint{1.487185in}{0.950572in}}%
\pgfpathlineto{\pgfqpoint{1.505400in}{0.963396in}}%
\pgfpathlineto{\pgfqpoint{1.523698in}{0.976349in}}%
\pgfpathlineto{\pgfqpoint{1.531562in}{0.978978in}}%
\pgfpathlineto{\pgfqpoint{1.539404in}{0.981923in}}%
\pgfpathlineto{\pgfqpoint{1.521132in}{0.968845in}}%
\pgfpathlineto{\pgfqpoint{1.502944in}{0.955886in}}%
\pgfpathlineto{\pgfqpoint{1.495075in}{0.953063in}}%
\pgfpathlineto{\pgfqpoint{1.487185in}{0.950572in}}%
\pgfpathclose%
\pgfusepath{fill}%
\end{pgfscope}%
\begin{pgfscope}%
\pgfpathrectangle{\pgfqpoint{0.407592in}{0.222904in}}{\pgfqpoint{1.411163in}{1.411163in}}%
\pgfusepath{clip}%
\pgfsetbuttcap%
\pgfsetroundjoin%
\definecolor{currentfill}{rgb}{0.556863,0.654902,0.847059}%
\pgfsetfillcolor{currentfill}%
\pgfsetfillopacity{0.130000}%
\pgfsetlinewidth{0.000000pt}%
\definecolor{currentstroke}{rgb}{0.000000,0.000000,0.000000}%
\pgfsetstrokecolor{currentstroke}%
\pgfsetdash{}{0pt}%
\pgfpathmoveto{\pgfqpoint{1.591775in}{1.015419in}}%
\pgfpathlineto{\pgfqpoint{1.610276in}{1.028900in}}%
\pgfpathlineto{\pgfqpoint{1.618020in}{1.032878in}}%
\pgfpathlineto{\pgfqpoint{1.625746in}{1.037151in}}%
\pgfpathlineto{\pgfqpoint{1.607272in}{1.023632in}}%
\pgfpathlineto{\pgfqpoint{1.599532in}{1.019374in}}%
\pgfusepath{fill}%
\end{pgfscope}%
\begin{pgfscope}%
\pgfpathrectangle{\pgfqpoint{0.407592in}{0.222904in}}{\pgfqpoint{1.411163in}{1.411163in}}%
\pgfusepath{clip}%
\pgfsetbuttcap%
\pgfsetroundjoin%
\definecolor{currentfill}{rgb}{0.556863,0.654902,0.847059}%
\pgfsetfillcolor{currentfill}%
\pgfsetfillopacity{0.130000}%
\pgfsetlinewidth{0.000000pt}%
\definecolor{currentstroke}{rgb}{0.000000,0.000000,0.000000}%
\pgfsetstrokecolor{currentstroke}%
\pgfsetdash{}{0pt}%
\pgfpathmoveto{\pgfqpoint{1.539404in}{0.981923in}}%
\pgfpathlineto{\pgfqpoint{1.557761in}{0.995115in}}%
\pgfpathlineto{\pgfqpoint{1.576203in}{1.008419in}}%
\pgfpathlineto{\pgfqpoint{1.583999in}{1.011768in}}%
\pgfpathlineto{\pgfqpoint{1.591775in}{1.015419in}}%
\pgfpathlineto{\pgfqpoint{1.573360in}{1.002040in}}%
\pgfpathlineto{\pgfqpoint{1.555029in}{0.988766in}}%
\pgfpathlineto{\pgfqpoint{1.547227in}{0.985185in}}%
\pgfpathlineto{\pgfqpoint{1.539404in}{0.981923in}}%
\pgfpathclose%
\pgfusepath{fill}%
\end{pgfscope}%
\begin{pgfscope}%
\pgfpathrectangle{\pgfqpoint{0.407592in}{0.222904in}}{\pgfqpoint{1.411163in}{1.411163in}}%
\pgfusepath{clip}%
\pgfsetbuttcap%
\pgfsetroundjoin%
\definecolor{currentfill}{rgb}{0.556863,0.654902,0.847059}%
\pgfsetfillcolor{currentfill}%
\pgfsetfillopacity{0.130000}%
\pgfsetlinewidth{0.000000pt}%
\definecolor{currentstroke}{rgb}{0.000000,0.000000,0.000000}%
\pgfsetstrokecolor{currentstroke}%
\pgfsetdash{}{0pt}%
\pgfpathmoveto{\pgfqpoint{1.211070in}{0.844971in}}%
\pgfpathlineto{\pgfqpoint{1.228674in}{0.853515in}}%
\pgfpathlineto{\pgfqpoint{1.246352in}{0.862437in}}%
\pgfpathlineto{\pgfqpoint{1.254693in}{0.859926in}}%
\pgfpathlineto{\pgfqpoint{1.263000in}{0.857695in}}%
\pgfpathlineto{\pgfqpoint{1.245359in}{0.848346in}}%
\pgfpathlineto{\pgfqpoint{1.227792in}{0.839360in}}%
\pgfpathlineto{\pgfqpoint{1.219448in}{0.842026in}}%
\pgfpathlineto{\pgfqpoint{1.211070in}{0.844971in}}%
\pgfpathclose%
\pgfusepath{fill}%
\end{pgfscope}%
\begin{pgfscope}%
\pgfpathrectangle{\pgfqpoint{0.407592in}{0.222904in}}{\pgfqpoint{1.411163in}{1.411163in}}%
\pgfusepath{clip}%
\pgfsetbuttcap%
\pgfsetroundjoin%
\definecolor{currentfill}{rgb}{0.556863,0.654902,0.847059}%
\pgfsetfillcolor{currentfill}%
\pgfsetfillopacity{0.130000}%
\pgfsetlinewidth{0.000000pt}%
\definecolor{currentstroke}{rgb}{0.000000,0.000000,0.000000}%
\pgfsetstrokecolor{currentstroke}%
\pgfsetdash{}{0pt}%
\pgfpathmoveto{\pgfqpoint{1.263000in}{0.857695in}}%
\pgfpathlineto{\pgfqpoint{1.280718in}{0.867381in}}%
\pgfpathlineto{\pgfqpoint{1.298512in}{0.877380in}}%
\pgfpathlineto{\pgfqpoint{1.306753in}{0.875838in}}%
\pgfpathlineto{\pgfqpoint{1.314963in}{0.874596in}}%
\pgfpathlineto{\pgfqpoint{1.297202in}{0.864211in}}%
\pgfpathlineto{\pgfqpoint{1.279519in}{0.854121in}}%
\pgfpathlineto{\pgfqpoint{1.271275in}{0.855756in}}%
\pgfpathlineto{\pgfqpoint{1.263000in}{0.857695in}}%
\pgfpathclose%
\pgfusepath{fill}%
\end{pgfscope}%
\begin{pgfscope}%
\pgfpathrectangle{\pgfqpoint{0.407592in}{0.222904in}}{\pgfqpoint{1.411163in}{1.411163in}}%
\pgfusepath{clip}%
\pgfsetbuttcap%
\pgfsetroundjoin%
\definecolor{currentfill}{rgb}{0.556863,0.654902,0.847059}%
\pgfsetfillcolor{currentfill}%
\pgfsetfillopacity{0.130000}%
\pgfsetlinewidth{0.000000pt}%
\definecolor{currentstroke}{rgb}{0.000000,0.000000,0.000000}%
\pgfsetstrokecolor{currentstroke}%
\pgfsetdash{}{0pt}%
\pgfpathmoveto{\pgfqpoint{1.159136in}{0.836713in}}%
\pgfpathlineto{\pgfqpoint{1.176637in}{0.843955in}}%
\pgfpathlineto{\pgfqpoint{1.194208in}{0.851644in}}%
\pgfpathlineto{\pgfqpoint{1.202657in}{0.848181in}}%
\pgfpathlineto{\pgfqpoint{1.211070in}{0.844971in}}%
\pgfpathlineto{\pgfqpoint{1.193538in}{0.836833in}}%
\pgfpathlineto{\pgfqpoint{1.176079in}{0.829133in}}%
\pgfpathlineto{\pgfqpoint{1.167625in}{0.832800in}}%
\pgfpathlineto{\pgfqpoint{1.159136in}{0.836713in}}%
\pgfpathclose%
\pgfusepath{fill}%
\end{pgfscope}%
\begin{pgfscope}%
\pgfpathrectangle{\pgfqpoint{0.407592in}{0.222904in}}{\pgfqpoint{1.411163in}{1.411163in}}%
\pgfusepath{clip}%
\pgfsetbuttcap%
\pgfsetroundjoin%
\definecolor{currentfill}{rgb}{0.556863,0.654902,0.847059}%
\pgfsetfillcolor{currentfill}%
\pgfsetfillopacity{0.130000}%
\pgfsetlinewidth{0.000000pt}%
\definecolor{currentstroke}{rgb}{0.000000,0.000000,0.000000}%
\pgfsetstrokecolor{currentstroke}%
\pgfsetdash{}{0pt}%
\pgfpathmoveto{\pgfqpoint{1.314963in}{0.874596in}}%
\pgfpathlineto{\pgfqpoint{1.332803in}{0.885254in}}%
\pgfpathlineto{\pgfqpoint{1.350722in}{0.896167in}}%
\pgfpathlineto{\pgfqpoint{1.358870in}{0.895582in}}%
\pgfpathlineto{\pgfqpoint{1.366990in}{0.895308in}}%
\pgfpathlineto{\pgfqpoint{1.349102in}{0.884062in}}%
\pgfpathlineto{\pgfqpoint{1.331294in}{0.873053in}}%
\pgfpathlineto{\pgfqpoint{1.323143in}{0.873664in}}%
\pgfpathlineto{\pgfqpoint{1.314963in}{0.874596in}}%
\pgfpathclose%
\pgfusepath{fill}%
\end{pgfscope}%
\begin{pgfscope}%
\pgfpathrectangle{\pgfqpoint{0.407592in}{0.222904in}}{\pgfqpoint{1.411163in}{1.411163in}}%
\pgfusepath{clip}%
\pgfsetbuttcap%
\pgfsetroundjoin%
\definecolor{currentfill}{rgb}{0.556863,0.654902,0.847059}%
\pgfsetfillcolor{currentfill}%
\pgfsetfillopacity{0.130000}%
\pgfsetlinewidth{0.000000pt}%
\definecolor{currentstroke}{rgb}{0.000000,0.000000,0.000000}%
\pgfsetstrokecolor{currentstroke}%
\pgfsetdash{}{0pt}%
\pgfpathmoveto{\pgfqpoint{1.107159in}{0.833106in}}%
\pgfpathlineto{\pgfqpoint{1.124569in}{0.838910in}}%
\pgfpathlineto{\pgfqpoint{1.142045in}{0.845227in}}%
\pgfpathlineto{\pgfqpoint{1.150609in}{0.840859in}}%
\pgfpathlineto{\pgfqpoint{1.159136in}{0.836713in}}%
\pgfpathlineto{\pgfqpoint{1.141704in}{0.829950in}}%
\pgfpathlineto{\pgfqpoint{1.124339in}{0.823699in}}%
\pgfpathlineto{\pgfqpoint{1.115768in}{0.828298in}}%
\pgfpathlineto{\pgfqpoint{1.107159in}{0.833106in}}%
\pgfpathclose%
\pgfusepath{fill}%
\end{pgfscope}%
\begin{pgfscope}%
\pgfpathrectangle{\pgfqpoint{0.407592in}{0.222904in}}{\pgfqpoint{1.411163in}{1.411163in}}%
\pgfusepath{clip}%
\pgfsetbuttcap%
\pgfsetroundjoin%
\definecolor{currentfill}{rgb}{0.556863,0.654902,0.847059}%
\pgfsetfillcolor{currentfill}%
\pgfsetfillopacity{0.130000}%
\pgfsetlinewidth{0.000000pt}%
\definecolor{currentstroke}{rgb}{0.000000,0.000000,0.000000}%
\pgfsetstrokecolor{currentstroke}%
\pgfsetdash{}{0pt}%
\pgfpathmoveto{\pgfqpoint{1.366990in}{0.895308in}}%
\pgfpathlineto{\pgfqpoint{1.384959in}{0.906774in}}%
\pgfpathlineto{\pgfqpoint{1.403009in}{0.918447in}}%
\pgfpathlineto{\pgfqpoint{1.411072in}{0.918782in}}%
\pgfpathlineto{\pgfqpoint{1.419108in}{0.919433in}}%
\pgfpathlineto{\pgfqpoint{1.401087in}{0.907486in}}%
\pgfpathlineto{\pgfqpoint{1.383148in}{0.895730in}}%
\pgfpathlineto{\pgfqpoint{1.375083in}{0.895354in}}%
\pgfpathlineto{\pgfqpoint{1.366990in}{0.895308in}}%
\pgfpathclose%
\pgfusepath{fill}%
\end{pgfscope}%
\begin{pgfscope}%
\pgfpathrectangle{\pgfqpoint{0.407592in}{0.222904in}}{\pgfqpoint{1.411163in}{1.411163in}}%
\pgfusepath{clip}%
\pgfsetbuttcap%
\pgfsetroundjoin%
\definecolor{currentfill}{rgb}{0.556863,0.654902,0.847059}%
\pgfsetfillcolor{currentfill}%
\pgfsetfillopacity{0.130000}%
\pgfsetlinewidth{0.000000pt}%
\definecolor{currentstroke}{rgb}{0.000000,0.000000,0.000000}%
\pgfsetstrokecolor{currentstroke}%
\pgfsetdash{}{0pt}%
\pgfpathmoveto{\pgfqpoint{1.419108in}{0.919433in}}%
\pgfpathlineto{\pgfqpoint{1.437211in}{0.931558in}}%
\pgfpathlineto{\pgfqpoint{1.455397in}{0.943851in}}%
\pgfpathlineto{\pgfqpoint{1.463380in}{0.945054in}}%
\pgfpathlineto{\pgfqpoint{1.471338in}{0.946572in}}%
\pgfpathlineto{\pgfqpoint{1.453180in}{0.934065in}}%
\pgfpathlineto{\pgfqpoint{1.435105in}{0.921711in}}%
\pgfpathlineto{\pgfqpoint{1.427119in}{0.920407in}}%
\pgfpathlineto{\pgfqpoint{1.419108in}{0.919433in}}%
\pgfpathclose%
\pgfusepath{fill}%
\end{pgfscope}%
\begin{pgfscope}%
\pgfpathrectangle{\pgfqpoint{0.407592in}{0.222904in}}{\pgfqpoint{1.411163in}{1.411163in}}%
\pgfusepath{clip}%
\pgfsetbuttcap%
\pgfsetroundjoin%
\definecolor{currentfill}{rgb}{0.556863,0.654902,0.847059}%
\pgfsetfillcolor{currentfill}%
\pgfsetfillopacity{0.130000}%
\pgfsetlinewidth{0.000000pt}%
\definecolor{currentstroke}{rgb}{0.000000,0.000000,0.000000}%
\pgfsetstrokecolor{currentstroke}%
\pgfsetdash{}{0pt}%
\pgfpathmoveto{\pgfqpoint{1.055099in}{0.834198in}}%
\pgfpathlineto{\pgfqpoint{1.072431in}{0.838466in}}%
\pgfpathlineto{\pgfqpoint{1.089824in}{0.843305in}}%
\pgfpathlineto{\pgfqpoint{1.098511in}{0.838112in}}%
\pgfpathlineto{\pgfqpoint{1.107159in}{0.833106in}}%
\pgfpathlineto{\pgfqpoint{1.089813in}{0.827849in}}%
\pgfpathlineto{\pgfqpoint{1.072530in}{0.823169in}}%
\pgfpathlineto{\pgfqpoint{1.063834in}{0.828598in}}%
\pgfpathlineto{\pgfqpoint{1.055099in}{0.834198in}}%
\pgfpathclose%
\pgfusepath{fill}%
\end{pgfscope}%
\begin{pgfscope}%
\pgfpathrectangle{\pgfqpoint{0.407592in}{0.222904in}}{\pgfqpoint{1.411163in}{1.411163in}}%
\pgfusepath{clip}%
\pgfsetbuttcap%
\pgfsetroundjoin%
\definecolor{currentfill}{rgb}{0.556863,0.654902,0.847059}%
\pgfsetfillcolor{currentfill}%
\pgfsetfillopacity{0.130000}%
\pgfsetlinewidth{0.000000pt}%
\definecolor{currentstroke}{rgb}{0.000000,0.000000,0.000000}%
\pgfsetstrokecolor{currentstroke}%
\pgfsetdash{}{0pt}%
\pgfpathmoveto{\pgfqpoint{1.471338in}{0.946572in}}%
\pgfpathlineto{\pgfqpoint{1.489580in}{0.959227in}}%
\pgfpathlineto{\pgfqpoint{1.507906in}{0.972020in}}%
\pgfpathlineto{\pgfqpoint{1.515813in}{0.974031in}}%
\pgfpathlineto{\pgfqpoint{1.523698in}{0.976349in}}%
\pgfpathlineto{\pgfqpoint{1.505400in}{0.963396in}}%
\pgfpathlineto{\pgfqpoint{1.487185in}{0.950572in}}%
\pgfpathlineto{\pgfqpoint{1.479273in}{0.948410in}}%
\pgfpathlineto{\pgfqpoint{1.471338in}{0.946572in}}%
\pgfpathclose%
\pgfusepath{fill}%
\end{pgfscope}%
\begin{pgfscope}%
\pgfpathrectangle{\pgfqpoint{0.407592in}{0.222904in}}{\pgfqpoint{1.411163in}{1.411163in}}%
\pgfusepath{clip}%
\pgfsetbuttcap%
\pgfsetroundjoin%
\definecolor{currentfill}{rgb}{0.556863,0.654902,0.847059}%
\pgfsetfillcolor{currentfill}%
\pgfsetfillopacity{0.130000}%
\pgfsetlinewidth{0.000000pt}%
\definecolor{currentstroke}{rgb}{0.000000,0.000000,0.000000}%
\pgfsetstrokecolor{currentstroke}%
\pgfsetdash{}{0pt}%
\pgfpathmoveto{\pgfqpoint{1.576203in}{1.008419in}}%
\pgfpathlineto{\pgfqpoint{1.594731in}{1.021829in}}%
\pgfpathlineto{\pgfqpoint{1.602513in}{1.025218in}}%
\pgfpathlineto{\pgfqpoint{1.610276in}{1.028900in}}%
\pgfpathlineto{\pgfqpoint{1.591775in}{1.015419in}}%
\pgfpathlineto{\pgfqpoint{1.583999in}{1.011768in}}%
\pgfusepath{fill}%
\end{pgfscope}%
\begin{pgfscope}%
\pgfpathrectangle{\pgfqpoint{0.407592in}{0.222904in}}{\pgfqpoint{1.411163in}{1.411163in}}%
\pgfusepath{clip}%
\pgfsetbuttcap%
\pgfsetroundjoin%
\definecolor{currentfill}{rgb}{0.556863,0.654902,0.847059}%
\pgfsetfillcolor{currentfill}%
\pgfsetfillopacity{0.130000}%
\pgfsetlinewidth{0.000000pt}%
\definecolor{currentstroke}{rgb}{0.000000,0.000000,0.000000}%
\pgfsetstrokecolor{currentstroke}%
\pgfsetdash{}{0pt}%
\pgfpathmoveto{\pgfqpoint{1.523698in}{0.976349in}}%
\pgfpathlineto{\pgfqpoint{1.542082in}{0.989425in}}%
\pgfpathlineto{\pgfqpoint{1.560551in}{1.002619in}}%
\pgfpathlineto{\pgfqpoint{1.568387in}{1.005370in}}%
\pgfpathlineto{\pgfqpoint{1.576203in}{1.008419in}}%
\pgfpathlineto{\pgfqpoint{1.557761in}{0.995115in}}%
\pgfpathlineto{\pgfqpoint{1.539404in}{0.981923in}}%
\pgfpathlineto{\pgfqpoint{1.531562in}{0.978978in}}%
\pgfpathlineto{\pgfqpoint{1.523698in}{0.976349in}}%
\pgfpathclose%
\pgfusepath{fill}%
\end{pgfscope}%
\begin{pgfscope}%
\pgfpathrectangle{\pgfqpoint{0.407592in}{0.222904in}}{\pgfqpoint{1.411163in}{1.411163in}}%
\pgfusepath{clip}%
\pgfsetbuttcap%
\pgfsetroundjoin%
\definecolor{currentfill}{rgb}{0.556863,0.654902,0.847059}%
\pgfsetfillcolor{currentfill}%
\pgfsetfillopacity{0.130000}%
\pgfsetlinewidth{0.000000pt}%
\definecolor{currentstroke}{rgb}{0.000000,0.000000,0.000000}%
\pgfsetstrokecolor{currentstroke}%
\pgfsetdash{}{0pt}%
\pgfpathmoveto{\pgfqpoint{1.246352in}{0.862437in}}%
\pgfpathlineto{\pgfqpoint{1.264105in}{0.871713in}}%
\pgfpathlineto{\pgfqpoint{1.281936in}{0.881319in}}%
\pgfpathlineto{\pgfqpoint{1.290240in}{0.879210in}}%
\pgfpathlineto{\pgfqpoint{1.298512in}{0.877380in}}%
\pgfpathlineto{\pgfqpoint{1.280718in}{0.867381in}}%
\pgfpathlineto{\pgfqpoint{1.263000in}{0.857695in}}%
\pgfpathlineto{\pgfqpoint{1.254693in}{0.859926in}}%
\pgfpathlineto{\pgfqpoint{1.246352in}{0.862437in}}%
\pgfpathclose%
\pgfusepath{fill}%
\end{pgfscope}%
\begin{pgfscope}%
\pgfpathrectangle{\pgfqpoint{0.407592in}{0.222904in}}{\pgfqpoint{1.411163in}{1.411163in}}%
\pgfusepath{clip}%
\pgfsetbuttcap%
\pgfsetroundjoin%
\definecolor{currentfill}{rgb}{0.556863,0.654902,0.847059}%
\pgfsetfillcolor{currentfill}%
\pgfsetfillopacity{0.130000}%
\pgfsetlinewidth{0.000000pt}%
\definecolor{currentstroke}{rgb}{0.000000,0.000000,0.000000}%
\pgfsetstrokecolor{currentstroke}%
\pgfsetdash{}{0pt}%
\pgfpathmoveto{\pgfqpoint{1.194208in}{0.851644in}}%
\pgfpathlineto{\pgfqpoint{1.211852in}{0.859752in}}%
\pgfpathlineto{\pgfqpoint{1.229568in}{0.868251in}}%
\pgfpathlineto{\pgfqpoint{1.237977in}{0.865216in}}%
\pgfpathlineto{\pgfqpoint{1.246352in}{0.862437in}}%
\pgfpathlineto{\pgfqpoint{1.228674in}{0.853515in}}%
\pgfpathlineto{\pgfqpoint{1.211070in}{0.844971in}}%
\pgfpathlineto{\pgfqpoint{1.202657in}{0.848181in}}%
\pgfpathlineto{\pgfqpoint{1.194208in}{0.851644in}}%
\pgfpathclose%
\pgfusepath{fill}%
\end{pgfscope}%
\begin{pgfscope}%
\pgfpathrectangle{\pgfqpoint{0.407592in}{0.222904in}}{\pgfqpoint{1.411163in}{1.411163in}}%
\pgfusepath{clip}%
\pgfsetbuttcap%
\pgfsetroundjoin%
\definecolor{currentfill}{rgb}{0.556863,0.654902,0.847059}%
\pgfsetfillcolor{currentfill}%
\pgfsetfillopacity{0.130000}%
\pgfsetlinewidth{0.000000pt}%
\definecolor{currentstroke}{rgb}{0.000000,0.000000,0.000000}%
\pgfsetstrokecolor{currentstroke}%
\pgfsetdash{}{0pt}%
\pgfpathmoveto{\pgfqpoint{1.298512in}{0.877380in}}%
\pgfpathlineto{\pgfqpoint{1.316385in}{0.887671in}}%
\pgfpathlineto{\pgfqpoint{1.334337in}{0.898235in}}%
\pgfpathlineto{\pgfqpoint{1.342544in}{0.897054in}}%
\pgfpathlineto{\pgfqpoint{1.350722in}{0.896167in}}%
\pgfpathlineto{\pgfqpoint{1.332803in}{0.885254in}}%
\pgfpathlineto{\pgfqpoint{1.314963in}{0.874596in}}%
\pgfpathlineto{\pgfqpoint{1.306753in}{0.875838in}}%
\pgfpathlineto{\pgfqpoint{1.298512in}{0.877380in}}%
\pgfpathclose%
\pgfusepath{fill}%
\end{pgfscope}%
\begin{pgfscope}%
\pgfpathrectangle{\pgfqpoint{0.407592in}{0.222904in}}{\pgfqpoint{1.411163in}{1.411163in}}%
\pgfusepath{clip}%
\pgfsetbuttcap%
\pgfsetroundjoin%
\definecolor{currentfill}{rgb}{0.556863,0.654902,0.847059}%
\pgfsetfillcolor{currentfill}%
\pgfsetfillopacity{0.130000}%
\pgfsetlinewidth{0.000000pt}%
\definecolor{currentstroke}{rgb}{0.000000,0.000000,0.000000}%
\pgfsetstrokecolor{currentstroke}%
\pgfsetdash{}{0pt}%
\pgfpathmoveto{\pgfqpoint{1.142045in}{0.845227in}}%
\pgfpathlineto{\pgfqpoint{1.159589in}{0.852027in}}%
\pgfpathlineto{\pgfqpoint{1.177203in}{0.859282in}}%
\pgfpathlineto{\pgfqpoint{1.185724in}{0.855348in}}%
\pgfpathlineto{\pgfqpoint{1.194208in}{0.851644in}}%
\pgfpathlineto{\pgfqpoint{1.176637in}{0.843955in}}%
\pgfpathlineto{\pgfqpoint{1.159136in}{0.836713in}}%
\pgfpathlineto{\pgfqpoint{1.150609in}{0.840859in}}%
\pgfpathlineto{\pgfqpoint{1.142045in}{0.845227in}}%
\pgfpathclose%
\pgfusepath{fill}%
\end{pgfscope}%
\begin{pgfscope}%
\pgfpathrectangle{\pgfqpoint{0.407592in}{0.222904in}}{\pgfqpoint{1.411163in}{1.411163in}}%
\pgfusepath{clip}%
\pgfsetbuttcap%
\pgfsetroundjoin%
\definecolor{currentfill}{rgb}{0.556863,0.654902,0.847059}%
\pgfsetfillcolor{currentfill}%
\pgfsetfillopacity{0.130000}%
\pgfsetlinewidth{0.000000pt}%
\definecolor{currentstroke}{rgb}{0.000000,0.000000,0.000000}%
\pgfsetstrokecolor{currentstroke}%
\pgfsetdash{}{0pt}%
\pgfpathmoveto{\pgfqpoint{1.350722in}{0.896167in}}%
\pgfpathlineto{\pgfqpoint{1.368722in}{0.907319in}}%
\pgfpathlineto{\pgfqpoint{1.386803in}{0.918694in}}%
\pgfpathlineto{\pgfqpoint{1.394920in}{0.918420in}}%
\pgfpathlineto{\pgfqpoint{1.403009in}{0.918447in}}%
\pgfpathlineto{\pgfqpoint{1.384959in}{0.906774in}}%
\pgfpathlineto{\pgfqpoint{1.366990in}{0.895308in}}%
\pgfpathlineto{\pgfqpoint{1.358870in}{0.895582in}}%
\pgfpathlineto{\pgfqpoint{1.350722in}{0.896167in}}%
\pgfpathclose%
\pgfusepath{fill}%
\end{pgfscope}%
\begin{pgfscope}%
\pgfpathrectangle{\pgfqpoint{0.407592in}{0.222904in}}{\pgfqpoint{1.411163in}{1.411163in}}%
\pgfusepath{clip}%
\pgfsetbuttcap%
\pgfsetroundjoin%
\definecolor{currentfill}{rgb}{0.556863,0.654902,0.847059}%
\pgfsetfillcolor{currentfill}%
\pgfsetfillopacity{0.130000}%
\pgfsetlinewidth{0.000000pt}%
\definecolor{currentstroke}{rgb}{0.000000,0.000000,0.000000}%
\pgfsetstrokecolor{currentstroke}%
\pgfsetdash{}{0pt}%
\pgfpathmoveto{\pgfqpoint{1.403009in}{0.918447in}}%
\pgfpathlineto{\pgfqpoint{1.421141in}{0.930314in}}%
\pgfpathlineto{\pgfqpoint{1.439357in}{0.942364in}}%
\pgfpathlineto{\pgfqpoint{1.447390in}{0.942956in}}%
\pgfpathlineto{\pgfqpoint{1.455397in}{0.943851in}}%
\pgfpathlineto{\pgfqpoint{1.437211in}{0.931558in}}%
\pgfpathlineto{\pgfqpoint{1.419108in}{0.919433in}}%
\pgfpathlineto{\pgfqpoint{1.411072in}{0.918782in}}%
\pgfpathlineto{\pgfqpoint{1.403009in}{0.918447in}}%
\pgfpathclose%
\pgfusepath{fill}%
\end{pgfscope}%
\begin{pgfscope}%
\pgfpathrectangle{\pgfqpoint{0.407592in}{0.222904in}}{\pgfqpoint{1.411163in}{1.411163in}}%
\pgfusepath{clip}%
\pgfsetbuttcap%
\pgfsetroundjoin%
\definecolor{currentfill}{rgb}{0.556863,0.654902,0.847059}%
\pgfsetfillcolor{currentfill}%
\pgfsetfillopacity{0.130000}%
\pgfsetlinewidth{0.000000pt}%
\definecolor{currentstroke}{rgb}{0.000000,0.000000,0.000000}%
\pgfsetstrokecolor{currentstroke}%
\pgfsetdash{}{0pt}%
\pgfpathmoveto{\pgfqpoint{1.089824in}{0.843305in}}%
\pgfpathlineto{\pgfqpoint{1.107280in}{0.848686in}}%
\pgfpathlineto{\pgfqpoint{1.124802in}{0.854579in}}%
\pgfpathlineto{\pgfqpoint{1.133443in}{0.849803in}}%
\pgfpathlineto{\pgfqpoint{1.142045in}{0.845227in}}%
\pgfpathlineto{\pgfqpoint{1.124569in}{0.838910in}}%
\pgfpathlineto{\pgfqpoint{1.107159in}{0.833106in}}%
\pgfpathlineto{\pgfqpoint{1.098511in}{0.838112in}}%
\pgfpathlineto{\pgfqpoint{1.089824in}{0.843305in}}%
\pgfpathclose%
\pgfusepath{fill}%
\end{pgfscope}%
\begin{pgfscope}%
\pgfpathrectangle{\pgfqpoint{0.407592in}{0.222904in}}{\pgfqpoint{1.411163in}{1.411163in}}%
\pgfusepath{clip}%
\pgfsetbuttcap%
\pgfsetroundjoin%
\definecolor{currentfill}{rgb}{0.556863,0.654902,0.847059}%
\pgfsetfillcolor{currentfill}%
\pgfsetfillopacity{0.130000}%
\pgfsetlinewidth{0.000000pt}%
\definecolor{currentstroke}{rgb}{0.000000,0.000000,0.000000}%
\pgfsetstrokecolor{currentstroke}%
\pgfsetdash{}{0pt}%
\pgfpathmoveto{\pgfqpoint{1.455397in}{0.943851in}}%
\pgfpathlineto{\pgfqpoint{1.473667in}{0.956304in}}%
\pgfpathlineto{\pgfqpoint{1.492021in}{0.968909in}}%
\pgfpathlineto{\pgfqpoint{1.499975in}{0.970315in}}%
\pgfpathlineto{\pgfqpoint{1.507906in}{0.972020in}}%
\pgfpathlineto{\pgfqpoint{1.489580in}{0.959227in}}%
\pgfpathlineto{\pgfqpoint{1.471338in}{0.946572in}}%
\pgfpathlineto{\pgfqpoint{1.463380in}{0.945054in}}%
\pgfpathlineto{\pgfqpoint{1.455397in}{0.943851in}}%
\pgfpathclose%
\pgfusepath{fill}%
\end{pgfscope}%
\begin{pgfscope}%
\pgfpathrectangle{\pgfqpoint{0.407592in}{0.222904in}}{\pgfqpoint{1.411163in}{1.411163in}}%
\pgfusepath{clip}%
\pgfsetbuttcap%
\pgfsetroundjoin%
\definecolor{currentfill}{rgb}{0.556863,0.654902,0.847059}%
\pgfsetfillcolor{currentfill}%
\pgfsetfillopacity{0.130000}%
\pgfsetlinewidth{0.000000pt}%
\definecolor{currentstroke}{rgb}{0.000000,0.000000,0.000000}%
\pgfsetstrokecolor{currentstroke}%
\pgfsetdash{}{0pt}%
\pgfpathmoveto{\pgfqpoint{1.560551in}{1.002619in}}%
\pgfpathlineto{\pgfqpoint{1.579106in}{1.015929in}}%
\pgfpathlineto{\pgfqpoint{1.586928in}{1.018733in}}%
\pgfpathlineto{\pgfqpoint{1.594731in}{1.021829in}}%
\pgfpathlineto{\pgfqpoint{1.576203in}{1.008419in}}%
\pgfpathlineto{\pgfqpoint{1.568387in}{1.005370in}}%
\pgfusepath{fill}%
\end{pgfscope}%
\begin{pgfscope}%
\pgfpathrectangle{\pgfqpoint{0.407592in}{0.222904in}}{\pgfqpoint{1.411163in}{1.411163in}}%
\pgfusepath{clip}%
\pgfsetbuttcap%
\pgfsetroundjoin%
\definecolor{currentfill}{rgb}{0.556863,0.654902,0.847059}%
\pgfsetfillcolor{currentfill}%
\pgfsetfillopacity{0.130000}%
\pgfsetlinewidth{0.000000pt}%
\definecolor{currentstroke}{rgb}{0.000000,0.000000,0.000000}%
\pgfsetstrokecolor{currentstroke}%
\pgfsetdash{}{0pt}%
\pgfpathmoveto{\pgfqpoint{1.507906in}{0.972020in}}%
\pgfpathlineto{\pgfqpoint{1.526317in}{0.984948in}}%
\pgfpathlineto{\pgfqpoint{1.544813in}{0.998003in}}%
\pgfpathlineto{\pgfqpoint{1.552693in}{1.000165in}}%
\pgfpathlineto{\pgfqpoint{1.560551in}{1.002619in}}%
\pgfpathlineto{\pgfqpoint{1.542082in}{0.989425in}}%
\pgfpathlineto{\pgfqpoint{1.523698in}{0.976349in}}%
\pgfpathlineto{\pgfqpoint{1.515813in}{0.974031in}}%
\pgfpathlineto{\pgfqpoint{1.507906in}{0.972020in}}%
\pgfpathclose%
\pgfusepath{fill}%
\end{pgfscope}%
\begin{pgfscope}%
\pgfpathrectangle{\pgfqpoint{0.407592in}{0.222904in}}{\pgfqpoint{1.411163in}{1.411163in}}%
\pgfusepath{clip}%
\pgfsetbuttcap%
\pgfsetroundjoin%
\definecolor{currentfill}{rgb}{0.556863,0.654902,0.847059}%
\pgfsetfillcolor{currentfill}%
\pgfsetfillopacity{0.130000}%
\pgfsetlinewidth{0.000000pt}%
\definecolor{currentstroke}{rgb}{0.000000,0.000000,0.000000}%
\pgfsetstrokecolor{currentstroke}%
\pgfsetdash{}{0pt}%
\pgfpathmoveto{\pgfqpoint{1.229568in}{0.868251in}}%
\pgfpathlineto{\pgfqpoint{1.247360in}{0.877118in}}%
\pgfpathlineto{\pgfqpoint{1.265228in}{0.886328in}}%
\pgfpathlineto{\pgfqpoint{1.273598in}{0.883695in}}%
\pgfpathlineto{\pgfqpoint{1.281936in}{0.881319in}}%
\pgfpathlineto{\pgfqpoint{1.264105in}{0.871713in}}%
\pgfpathlineto{\pgfqpoint{1.246352in}{0.862437in}}%
\pgfpathlineto{\pgfqpoint{1.237977in}{0.865216in}}%
\pgfpathlineto{\pgfqpoint{1.229568in}{0.868251in}}%
\pgfpathclose%
\pgfusepath{fill}%
\end{pgfscope}%
\begin{pgfscope}%
\pgfpathrectangle{\pgfqpoint{0.407592in}{0.222904in}}{\pgfqpoint{1.411163in}{1.411163in}}%
\pgfusepath{clip}%
\pgfsetbuttcap%
\pgfsetroundjoin%
\definecolor{currentfill}{rgb}{0.556863,0.654902,0.847059}%
\pgfsetfillcolor{currentfill}%
\pgfsetfillopacity{0.130000}%
\pgfsetlinewidth{0.000000pt}%
\definecolor{currentstroke}{rgb}{0.000000,0.000000,0.000000}%
\pgfsetstrokecolor{currentstroke}%
\pgfsetdash{}{0pt}%
\pgfpathmoveto{\pgfqpoint{1.281936in}{0.881319in}}%
\pgfpathlineto{\pgfqpoint{1.299843in}{0.891233in}}%
\pgfpathlineto{\pgfqpoint{1.317830in}{0.901437in}}%
\pgfpathlineto{\pgfqpoint{1.326099in}{0.899698in}}%
\pgfpathlineto{\pgfqpoint{1.334337in}{0.898235in}}%
\pgfpathlineto{\pgfqpoint{1.316385in}{0.887671in}}%
\pgfpathlineto{\pgfqpoint{1.298512in}{0.877380in}}%
\pgfpathlineto{\pgfqpoint{1.290240in}{0.879210in}}%
\pgfpathlineto{\pgfqpoint{1.281936in}{0.881319in}}%
\pgfpathclose%
\pgfusepath{fill}%
\end{pgfscope}%
\begin{pgfscope}%
\pgfpathrectangle{\pgfqpoint{0.407592in}{0.222904in}}{\pgfqpoint{1.411163in}{1.411163in}}%
\pgfusepath{clip}%
\pgfsetbuttcap%
\pgfsetroundjoin%
\definecolor{currentfill}{rgb}{0.556863,0.654902,0.847059}%
\pgfsetfillcolor{currentfill}%
\pgfsetfillopacity{0.130000}%
\pgfsetlinewidth{0.000000pt}%
\definecolor{currentstroke}{rgb}{0.000000,0.000000,0.000000}%
\pgfsetstrokecolor{currentstroke}%
\pgfsetdash{}{0pt}%
\pgfpathmoveto{\pgfqpoint{1.177203in}{0.859282in}}%
\pgfpathlineto{\pgfqpoint{1.194887in}{0.866964in}}%
\pgfpathlineto{\pgfqpoint{1.212645in}{0.875048in}}%
\pgfpathlineto{\pgfqpoint{1.221125in}{0.871532in}}%
\pgfpathlineto{\pgfqpoint{1.229568in}{0.868251in}}%
\pgfpathlineto{\pgfqpoint{1.211852in}{0.859752in}}%
\pgfpathlineto{\pgfqpoint{1.194208in}{0.851644in}}%
\pgfpathlineto{\pgfqpoint{1.185724in}{0.855348in}}%
\pgfpathlineto{\pgfqpoint{1.177203in}{0.859282in}}%
\pgfpathclose%
\pgfusepath{fill}%
\end{pgfscope}%
\begin{pgfscope}%
\pgfpathrectangle{\pgfqpoint{0.407592in}{0.222904in}}{\pgfqpoint{1.411163in}{1.411163in}}%
\pgfusepath{clip}%
\pgfsetbuttcap%
\pgfsetroundjoin%
\definecolor{currentfill}{rgb}{0.556863,0.654902,0.847059}%
\pgfsetfillcolor{currentfill}%
\pgfsetfillopacity{0.130000}%
\pgfsetlinewidth{0.000000pt}%
\definecolor{currentstroke}{rgb}{0.000000,0.000000,0.000000}%
\pgfsetstrokecolor{currentstroke}%
\pgfsetdash{}{0pt}%
\pgfpathmoveto{\pgfqpoint{1.334337in}{0.898235in}}%
\pgfpathlineto{\pgfqpoint{1.352369in}{0.909054in}}%
\pgfpathlineto{\pgfqpoint{1.370483in}{0.920114in}}%
\pgfpathlineto{\pgfqpoint{1.378657in}{0.919261in}}%
\pgfpathlineto{\pgfqpoint{1.386803in}{0.918694in}}%
\pgfpathlineto{\pgfqpoint{1.368722in}{0.907319in}}%
\pgfpathlineto{\pgfqpoint{1.350722in}{0.896167in}}%
\pgfpathlineto{\pgfqpoint{1.342544in}{0.897054in}}%
\pgfpathlineto{\pgfqpoint{1.334337in}{0.898235in}}%
\pgfpathclose%
\pgfusepath{fill}%
\end{pgfscope}%
\begin{pgfscope}%
\pgfpathrectangle{\pgfqpoint{0.407592in}{0.222904in}}{\pgfqpoint{1.411163in}{1.411163in}}%
\pgfusepath{clip}%
\pgfsetbuttcap%
\pgfsetroundjoin%
\definecolor{currentfill}{rgb}{0.556863,0.654902,0.847059}%
\pgfsetfillcolor{currentfill}%
\pgfsetfillopacity{0.130000}%
\pgfsetlinewidth{0.000000pt}%
\definecolor{currentstroke}{rgb}{0.000000,0.000000,0.000000}%
\pgfsetstrokecolor{currentstroke}%
\pgfsetdash{}{0pt}%
\pgfpathmoveto{\pgfqpoint{1.124802in}{0.854579in}}%
\pgfpathlineto{\pgfqpoint{1.142391in}{0.860958in}}%
\pgfpathlineto{\pgfqpoint{1.160049in}{0.867796in}}%
\pgfpathlineto{\pgfqpoint{1.168645in}{0.863434in}}%
\pgfpathlineto{\pgfqpoint{1.177203in}{0.859282in}}%
\pgfpathlineto{\pgfqpoint{1.159589in}{0.852027in}}%
\pgfpathlineto{\pgfqpoint{1.142045in}{0.845227in}}%
\pgfpathlineto{\pgfqpoint{1.133443in}{0.849803in}}%
\pgfpathlineto{\pgfqpoint{1.124802in}{0.854579in}}%
\pgfpathclose%
\pgfusepath{fill}%
\end{pgfscope}%
\begin{pgfscope}%
\pgfpathrectangle{\pgfqpoint{0.407592in}{0.222904in}}{\pgfqpoint{1.411163in}{1.411163in}}%
\pgfusepath{clip}%
\pgfsetbuttcap%
\pgfsetroundjoin%
\definecolor{currentfill}{rgb}{0.556863,0.654902,0.847059}%
\pgfsetfillcolor{currentfill}%
\pgfsetfillopacity{0.130000}%
\pgfsetlinewidth{0.000000pt}%
\definecolor{currentstroke}{rgb}{0.000000,0.000000,0.000000}%
\pgfsetstrokecolor{currentstroke}%
\pgfsetdash{}{0pt}%
\pgfpathmoveto{\pgfqpoint{1.386803in}{0.918694in}}%
\pgfpathlineto{\pgfqpoint{1.404966in}{0.930280in}}%
\pgfpathlineto{\pgfqpoint{1.423212in}{0.942065in}}%
\pgfpathlineto{\pgfqpoint{1.431298in}{0.942069in}}%
\pgfpathlineto{\pgfqpoint{1.439357in}{0.942364in}}%
\pgfpathlineto{\pgfqpoint{1.421141in}{0.930314in}}%
\pgfpathlineto{\pgfqpoint{1.403009in}{0.918447in}}%
\pgfpathlineto{\pgfqpoint{1.394920in}{0.918420in}}%
\pgfpathlineto{\pgfqpoint{1.386803in}{0.918694in}}%
\pgfpathclose%
\pgfusepath{fill}%
\end{pgfscope}%
\begin{pgfscope}%
\pgfpathrectangle{\pgfqpoint{0.407592in}{0.222904in}}{\pgfqpoint{1.411163in}{1.411163in}}%
\pgfusepath{clip}%
\pgfsetbuttcap%
\pgfsetroundjoin%
\definecolor{currentfill}{rgb}{0.556863,0.654902,0.847059}%
\pgfsetfillcolor{currentfill}%
\pgfsetfillopacity{0.130000}%
\pgfsetlinewidth{0.000000pt}%
\definecolor{currentstroke}{rgb}{0.000000,0.000000,0.000000}%
\pgfsetstrokecolor{currentstroke}%
\pgfsetdash{}{0pt}%
\pgfpathmoveto{\pgfqpoint{1.439357in}{0.942364in}}%
\pgfpathlineto{\pgfqpoint{1.457656in}{0.954589in}}%
\pgfpathlineto{\pgfqpoint{1.476040in}{0.966978in}}%
\pgfpathlineto{\pgfqpoint{1.484043in}{0.967798in}}%
\pgfpathlineto{\pgfqpoint{1.492021in}{0.968909in}}%
\pgfpathlineto{\pgfqpoint{1.473667in}{0.956304in}}%
\pgfpathlineto{\pgfqpoint{1.455397in}{0.943851in}}%
\pgfpathlineto{\pgfqpoint{1.447390in}{0.942956in}}%
\pgfpathlineto{\pgfqpoint{1.439357in}{0.942364in}}%
\pgfpathclose%
\pgfusepath{fill}%
\end{pgfscope}%
\begin{pgfscope}%
\pgfpathrectangle{\pgfqpoint{0.407592in}{0.222904in}}{\pgfqpoint{1.411163in}{1.411163in}}%
\pgfusepath{clip}%
\pgfsetbuttcap%
\pgfsetroundjoin%
\definecolor{currentfill}{rgb}{0.556863,0.654902,0.847059}%
\pgfsetfillcolor{currentfill}%
\pgfsetfillopacity{0.130000}%
\pgfsetlinewidth{0.000000pt}%
\definecolor{currentstroke}{rgb}{0.000000,0.000000,0.000000}%
\pgfsetstrokecolor{currentstroke}%
\pgfsetdash{}{0pt}%
\pgfpathmoveto{\pgfqpoint{1.544813in}{0.998003in}}%
\pgfpathlineto{\pgfqpoint{1.563397in}{1.011183in}}%
\pgfpathlineto{\pgfqpoint{1.571262in}{1.013413in}}%
\pgfpathlineto{\pgfqpoint{1.579106in}{1.015929in}}%
\pgfpathlineto{\pgfqpoint{1.560551in}{1.002619in}}%
\pgfpathlineto{\pgfqpoint{1.552693in}{1.000165in}}%
\pgfusepath{fill}%
\end{pgfscope}%
\begin{pgfscope}%
\pgfpathrectangle{\pgfqpoint{0.407592in}{0.222904in}}{\pgfqpoint{1.411163in}{1.411163in}}%
\pgfusepath{clip}%
\pgfsetbuttcap%
\pgfsetroundjoin%
\definecolor{currentfill}{rgb}{0.556863,0.654902,0.847059}%
\pgfsetfillcolor{currentfill}%
\pgfsetfillopacity{0.130000}%
\pgfsetlinewidth{0.000000pt}%
\definecolor{currentstroke}{rgb}{0.000000,0.000000,0.000000}%
\pgfsetstrokecolor{currentstroke}%
\pgfsetdash{}{0pt}%
\pgfpathmoveto{\pgfqpoint{1.492021in}{0.968909in}}%
\pgfpathlineto{\pgfqpoint{1.510461in}{0.981659in}}%
\pgfpathlineto{\pgfqpoint{1.528986in}{0.994547in}}%
\pgfpathlineto{\pgfqpoint{1.536912in}{0.996132in}}%
\pgfpathlineto{\pgfqpoint{1.544813in}{0.998003in}}%
\pgfpathlineto{\pgfqpoint{1.526317in}{0.984948in}}%
\pgfpathlineto{\pgfqpoint{1.507906in}{0.972020in}}%
\pgfpathlineto{\pgfqpoint{1.499975in}{0.970315in}}%
\pgfpathlineto{\pgfqpoint{1.492021in}{0.968909in}}%
\pgfpathclose%
\pgfusepath{fill}%
\end{pgfscope}%
\begin{pgfscope}%
\pgfpathrectangle{\pgfqpoint{0.407592in}{0.222904in}}{\pgfqpoint{1.411163in}{1.411163in}}%
\pgfusepath{clip}%
\pgfsetbuttcap%
\pgfsetroundjoin%
\definecolor{currentfill}{rgb}{0.556863,0.654902,0.847059}%
\pgfsetfillcolor{currentfill}%
\pgfsetfillopacity{0.130000}%
\pgfsetlinewidth{0.000000pt}%
\definecolor{currentstroke}{rgb}{0.000000,0.000000,0.000000}%
\pgfsetstrokecolor{currentstroke}%
\pgfsetdash{}{0pt}%
\pgfpathmoveto{\pgfqpoint{1.265228in}{0.886328in}}%
\pgfpathlineto{\pgfqpoint{1.283173in}{0.895863in}}%
\pgfpathlineto{\pgfqpoint{1.301196in}{0.905701in}}%
\pgfpathlineto{\pgfqpoint{1.309529in}{0.903440in}}%
\pgfpathlineto{\pgfqpoint{1.317830in}{0.901437in}}%
\pgfpathlineto{\pgfqpoint{1.299843in}{0.891233in}}%
\pgfpathlineto{\pgfqpoint{1.281936in}{0.881319in}}%
\pgfpathlineto{\pgfqpoint{1.273598in}{0.883695in}}%
\pgfpathlineto{\pgfqpoint{1.265228in}{0.886328in}}%
\pgfpathclose%
\pgfusepath{fill}%
\end{pgfscope}%
\begin{pgfscope}%
\pgfpathrectangle{\pgfqpoint{0.407592in}{0.222904in}}{\pgfqpoint{1.411163in}{1.411163in}}%
\pgfusepath{clip}%
\pgfsetbuttcap%
\pgfsetroundjoin%
\definecolor{currentfill}{rgb}{0.556863,0.654902,0.847059}%
\pgfsetfillcolor{currentfill}%
\pgfsetfillopacity{0.130000}%
\pgfsetlinewidth{0.000000pt}%
\definecolor{currentstroke}{rgb}{0.000000,0.000000,0.000000}%
\pgfsetstrokecolor{currentstroke}%
\pgfsetdash{}{0pt}%
\pgfpathmoveto{\pgfqpoint{1.212645in}{0.875048in}}%
\pgfpathlineto{\pgfqpoint{1.230476in}{0.883511in}}%
\pgfpathlineto{\pgfqpoint{1.248383in}{0.892330in}}%
\pgfpathlineto{\pgfqpoint{1.256823in}{0.889210in}}%
\pgfpathlineto{\pgfqpoint{1.265228in}{0.886328in}}%
\pgfpathlineto{\pgfqpoint{1.247360in}{0.877118in}}%
\pgfpathlineto{\pgfqpoint{1.229568in}{0.868251in}}%
\pgfpathlineto{\pgfqpoint{1.221125in}{0.871532in}}%
\pgfpathlineto{\pgfqpoint{1.212645in}{0.875048in}}%
\pgfpathclose%
\pgfusepath{fill}%
\end{pgfscope}%
\begin{pgfscope}%
\pgfpathrectangle{\pgfqpoint{0.407592in}{0.222904in}}{\pgfqpoint{1.411163in}{1.411163in}}%
\pgfusepath{clip}%
\pgfsetbuttcap%
\pgfsetroundjoin%
\definecolor{currentfill}{rgb}{0.556863,0.654902,0.847059}%
\pgfsetfillcolor{currentfill}%
\pgfsetfillopacity{0.130000}%
\pgfsetlinewidth{0.000000pt}%
\definecolor{currentstroke}{rgb}{0.000000,0.000000,0.000000}%
\pgfsetstrokecolor{currentstroke}%
\pgfsetdash{}{0pt}%
\pgfpathmoveto{\pgfqpoint{1.317830in}{0.901437in}}%
\pgfpathlineto{\pgfqpoint{1.335897in}{0.911912in}}%
\pgfpathlineto{\pgfqpoint{1.354044in}{0.922644in}}%
\pgfpathlineto{\pgfqpoint{1.362279in}{0.921244in}}%
\pgfpathlineto{\pgfqpoint{1.370483in}{0.920114in}}%
\pgfpathlineto{\pgfqpoint{1.352369in}{0.909054in}}%
\pgfpathlineto{\pgfqpoint{1.334337in}{0.898235in}}%
\pgfpathlineto{\pgfqpoint{1.326099in}{0.899698in}}%
\pgfpathlineto{\pgfqpoint{1.317830in}{0.901437in}}%
\pgfpathclose%
\pgfusepath{fill}%
\end{pgfscope}%
\begin{pgfscope}%
\pgfpathrectangle{\pgfqpoint{0.407592in}{0.222904in}}{\pgfqpoint{1.411163in}{1.411163in}}%
\pgfusepath{clip}%
\pgfsetbuttcap%
\pgfsetroundjoin%
\definecolor{currentfill}{rgb}{0.556863,0.654902,0.847059}%
\pgfsetfillcolor{currentfill}%
\pgfsetfillopacity{0.130000}%
\pgfsetlinewidth{0.000000pt}%
\definecolor{currentstroke}{rgb}{0.000000,0.000000,0.000000}%
\pgfsetstrokecolor{currentstroke}%
\pgfsetdash{}{0pt}%
\pgfpathmoveto{\pgfqpoint{1.160049in}{0.867796in}}%
\pgfpathlineto{\pgfqpoint{1.177777in}{0.875067in}}%
\pgfpathlineto{\pgfqpoint{1.195577in}{0.882748in}}%
\pgfpathlineto{\pgfqpoint{1.204129in}{0.878790in}}%
\pgfpathlineto{\pgfqpoint{1.212645in}{0.875048in}}%
\pgfpathlineto{\pgfqpoint{1.194887in}{0.866964in}}%
\pgfpathlineto{\pgfqpoint{1.177203in}{0.859282in}}%
\pgfpathlineto{\pgfqpoint{1.168645in}{0.863434in}}%
\pgfpathlineto{\pgfqpoint{1.160049in}{0.867796in}}%
\pgfpathclose%
\pgfusepath{fill}%
\end{pgfscope}%
\begin{pgfscope}%
\pgfpathrectangle{\pgfqpoint{0.407592in}{0.222904in}}{\pgfqpoint{1.411163in}{1.411163in}}%
\pgfusepath{clip}%
\pgfsetbuttcap%
\pgfsetroundjoin%
\definecolor{currentfill}{rgb}{0.556863,0.654902,0.847059}%
\pgfsetfillcolor{currentfill}%
\pgfsetfillopacity{0.130000}%
\pgfsetlinewidth{0.000000pt}%
\definecolor{currentstroke}{rgb}{0.000000,0.000000,0.000000}%
\pgfsetstrokecolor{currentstroke}%
\pgfsetdash{}{0pt}%
\pgfpathmoveto{\pgfqpoint{1.370483in}{0.920114in}}%
\pgfpathlineto{\pgfqpoint{1.388678in}{0.931400in}}%
\pgfpathlineto{\pgfqpoint{1.406957in}{0.942900in}}%
\pgfpathlineto{\pgfqpoint{1.415098in}{0.942344in}}%
\pgfpathlineto{\pgfqpoint{1.423212in}{0.942065in}}%
\pgfpathlineto{\pgfqpoint{1.404966in}{0.930280in}}%
\pgfpathlineto{\pgfqpoint{1.386803in}{0.918694in}}%
\pgfpathlineto{\pgfqpoint{1.378657in}{0.919261in}}%
\pgfpathlineto{\pgfqpoint{1.370483in}{0.920114in}}%
\pgfpathclose%
\pgfusepath{fill}%
\end{pgfscope}%
\begin{pgfscope}%
\pgfpathrectangle{\pgfqpoint{0.407592in}{0.222904in}}{\pgfqpoint{1.411163in}{1.411163in}}%
\pgfusepath{clip}%
\pgfsetbuttcap%
\pgfsetroundjoin%
\definecolor{currentfill}{rgb}{0.556863,0.654902,0.847059}%
\pgfsetfillcolor{currentfill}%
\pgfsetfillopacity{0.130000}%
\pgfsetlinewidth{0.000000pt}%
\definecolor{currentstroke}{rgb}{0.000000,0.000000,0.000000}%
\pgfsetstrokecolor{currentstroke}%
\pgfsetdash{}{0pt}%
\pgfpathmoveto{\pgfqpoint{1.423212in}{0.942065in}}%
\pgfpathlineto{\pgfqpoint{1.441542in}{0.954037in}}%
\pgfpathlineto{\pgfqpoint{1.459956in}{0.966189in}}%
\pgfpathlineto{\pgfqpoint{1.468011in}{0.966444in}}%
\pgfpathlineto{\pgfqpoint{1.476040in}{0.966978in}}%
\pgfpathlineto{\pgfqpoint{1.457656in}{0.954589in}}%
\pgfpathlineto{\pgfqpoint{1.439357in}{0.942364in}}%
\pgfpathlineto{\pgfqpoint{1.431298in}{0.942069in}}%
\pgfpathlineto{\pgfqpoint{1.423212in}{0.942065in}}%
\pgfpathclose%
\pgfusepath{fill}%
\end{pgfscope}%
\begin{pgfscope}%
\pgfpathrectangle{\pgfqpoint{0.407592in}{0.222904in}}{\pgfqpoint{1.411163in}{1.411163in}}%
\pgfusepath{clip}%
\pgfsetbuttcap%
\pgfsetroundjoin%
\definecolor{currentfill}{rgb}{0.556863,0.654902,0.847059}%
\pgfsetfillcolor{currentfill}%
\pgfsetfillopacity{0.130000}%
\pgfsetlinewidth{0.000000pt}%
\definecolor{currentstroke}{rgb}{0.000000,0.000000,0.000000}%
\pgfsetstrokecolor{currentstroke}%
\pgfsetdash{}{0pt}%
\pgfpathmoveto{\pgfqpoint{1.528986in}{0.994547in}}%
\pgfpathlineto{\pgfqpoint{1.547598in}{1.007570in}}%
\pgfpathlineto{\pgfqpoint{1.555509in}{1.009236in}}%
\pgfpathlineto{\pgfqpoint{1.563397in}{1.011183in}}%
\pgfpathlineto{\pgfqpoint{1.544813in}{0.998003in}}%
\pgfpathlineto{\pgfqpoint{1.536912in}{0.996132in}}%
\pgfusepath{fill}%
\end{pgfscope}%
\begin{pgfscope}%
\pgfpathrectangle{\pgfqpoint{0.407592in}{0.222904in}}{\pgfqpoint{1.411163in}{1.411163in}}%
\pgfusepath{clip}%
\pgfsetbuttcap%
\pgfsetroundjoin%
\definecolor{currentfill}{rgb}{0.556863,0.654902,0.847059}%
\pgfsetfillcolor{currentfill}%
\pgfsetfillopacity{0.130000}%
\pgfsetlinewidth{0.000000pt}%
\definecolor{currentstroke}{rgb}{0.000000,0.000000,0.000000}%
\pgfsetstrokecolor{currentstroke}%
\pgfsetdash{}{0pt}%
\pgfpathmoveto{\pgfqpoint{1.476040in}{0.966978in}}%
\pgfpathlineto{\pgfqpoint{1.494509in}{0.979525in}}%
\pgfpathlineto{\pgfqpoint{1.513064in}{0.992222in}}%
\pgfpathlineto{\pgfqpoint{1.521038in}{0.993246in}}%
\pgfpathlineto{\pgfqpoint{1.528986in}{0.994547in}}%
\pgfpathlineto{\pgfqpoint{1.510461in}{0.981659in}}%
\pgfpathlineto{\pgfqpoint{1.492021in}{0.968909in}}%
\pgfpathlineto{\pgfqpoint{1.484043in}{0.967798in}}%
\pgfpathlineto{\pgfqpoint{1.476040in}{0.966978in}}%
\pgfpathclose%
\pgfusepath{fill}%
\end{pgfscope}%
\begin{pgfscope}%
\pgfpathrectangle{\pgfqpoint{0.407592in}{0.222904in}}{\pgfqpoint{1.411163in}{1.411163in}}%
\pgfusepath{clip}%
\pgfsetbuttcap%
\pgfsetroundjoin%
\definecolor{currentfill}{rgb}{0.556863,0.654902,0.847059}%
\pgfsetfillcolor{currentfill}%
\pgfsetfillopacity{0.130000}%
\pgfsetlinewidth{0.000000pt}%
\definecolor{currentstroke}{rgb}{0.000000,0.000000,0.000000}%
\pgfsetstrokecolor{currentstroke}%
\pgfsetdash{}{0pt}%
\pgfpathmoveto{\pgfqpoint{1.248383in}{0.892330in}}%
\pgfpathlineto{\pgfqpoint{1.266367in}{0.901485in}}%
\pgfpathlineto{\pgfqpoint{1.284428in}{0.910956in}}%
\pgfpathlineto{\pgfqpoint{1.292829in}{0.908209in}}%
\pgfpathlineto{\pgfqpoint{1.301196in}{0.905701in}}%
\pgfpathlineto{\pgfqpoint{1.283173in}{0.895863in}}%
\pgfpathlineto{\pgfqpoint{1.265228in}{0.886328in}}%
\pgfpathlineto{\pgfqpoint{1.256823in}{0.889210in}}%
\pgfpathlineto{\pgfqpoint{1.248383in}{0.892330in}}%
\pgfpathclose%
\pgfusepath{fill}%
\end{pgfscope}%
\begin{pgfscope}%
\pgfpathrectangle{\pgfqpoint{0.407592in}{0.222904in}}{\pgfqpoint{1.411163in}{1.411163in}}%
\pgfusepath{clip}%
\pgfsetbuttcap%
\pgfsetroundjoin%
\definecolor{currentfill}{rgb}{0.556863,0.654902,0.847059}%
\pgfsetfillcolor{currentfill}%
\pgfsetfillopacity{0.130000}%
\pgfsetlinewidth{0.000000pt}%
\definecolor{currentstroke}{rgb}{0.000000,0.000000,0.000000}%
\pgfsetstrokecolor{currentstroke}%
\pgfsetdash{}{0pt}%
\pgfpathmoveto{\pgfqpoint{1.301196in}{0.905701in}}%
\pgfpathlineto{\pgfqpoint{1.319298in}{0.915825in}}%
\pgfpathlineto{\pgfqpoint{1.337482in}{0.926220in}}%
\pgfpathlineto{\pgfqpoint{1.345779in}{0.924305in}}%
\pgfpathlineto{\pgfqpoint{1.354044in}{0.922644in}}%
\pgfpathlineto{\pgfqpoint{1.335897in}{0.911912in}}%
\pgfpathlineto{\pgfqpoint{1.317830in}{0.901437in}}%
\pgfpathlineto{\pgfqpoint{1.309529in}{0.903440in}}%
\pgfpathlineto{\pgfqpoint{1.301196in}{0.905701in}}%
\pgfpathclose%
\pgfusepath{fill}%
\end{pgfscope}%
\begin{pgfscope}%
\pgfpathrectangle{\pgfqpoint{0.407592in}{0.222904in}}{\pgfqpoint{1.411163in}{1.411163in}}%
\pgfusepath{clip}%
\pgfsetbuttcap%
\pgfsetroundjoin%
\definecolor{currentfill}{rgb}{0.556863,0.654902,0.847059}%
\pgfsetfillcolor{currentfill}%
\pgfsetfillopacity{0.130000}%
\pgfsetlinewidth{0.000000pt}%
\definecolor{currentstroke}{rgb}{0.000000,0.000000,0.000000}%
\pgfsetstrokecolor{currentstroke}%
\pgfsetdash{}{0pt}%
\pgfpathmoveto{\pgfqpoint{1.195577in}{0.882748in}}%
\pgfpathlineto{\pgfqpoint{1.213450in}{0.890816in}}%
\pgfpathlineto{\pgfqpoint{1.231398in}{0.899249in}}%
\pgfpathlineto{\pgfqpoint{1.239909in}{0.895679in}}%
\pgfpathlineto{\pgfqpoint{1.248383in}{0.892330in}}%
\pgfpathlineto{\pgfqpoint{1.230476in}{0.883511in}}%
\pgfpathlineto{\pgfqpoint{1.212645in}{0.875048in}}%
\pgfpathlineto{\pgfqpoint{1.204129in}{0.878790in}}%
\pgfpathlineto{\pgfqpoint{1.195577in}{0.882748in}}%
\pgfpathclose%
\pgfusepath{fill}%
\end{pgfscope}%
\begin{pgfscope}%
\pgfpathrectangle{\pgfqpoint{0.407592in}{0.222904in}}{\pgfqpoint{1.411163in}{1.411163in}}%
\pgfusepath{clip}%
\pgfsetbuttcap%
\pgfsetroundjoin%
\definecolor{currentfill}{rgb}{0.556863,0.654902,0.847059}%
\pgfsetfillcolor{currentfill}%
\pgfsetfillopacity{0.130000}%
\pgfsetlinewidth{0.000000pt}%
\definecolor{currentstroke}{rgb}{0.000000,0.000000,0.000000}%
\pgfsetstrokecolor{currentstroke}%
\pgfsetdash{}{0pt}%
\pgfpathmoveto{\pgfqpoint{1.354044in}{0.922644in}}%
\pgfpathlineto{\pgfqpoint{1.372274in}{0.933616in}}%
\pgfpathlineto{\pgfqpoint{1.390585in}{0.944818in}}%
\pgfpathlineto{\pgfqpoint{1.398786in}{0.943727in}}%
\pgfpathlineto{\pgfqpoint{1.406957in}{0.942900in}}%
\pgfpathlineto{\pgfqpoint{1.388678in}{0.931400in}}%
\pgfpathlineto{\pgfqpoint{1.370483in}{0.920114in}}%
\pgfpathlineto{\pgfqpoint{1.362279in}{0.921244in}}%
\pgfpathlineto{\pgfqpoint{1.354044in}{0.922644in}}%
\pgfpathclose%
\pgfusepath{fill}%
\end{pgfscope}%
\begin{pgfscope}%
\pgfpathrectangle{\pgfqpoint{0.407592in}{0.222904in}}{\pgfqpoint{1.411163in}{1.411163in}}%
\pgfusepath{clip}%
\pgfsetbuttcap%
\pgfsetroundjoin%
\definecolor{currentfill}{rgb}{0.556863,0.654902,0.847059}%
\pgfsetfillcolor{currentfill}%
\pgfsetfillopacity{0.130000}%
\pgfsetlinewidth{0.000000pt}%
\definecolor{currentstroke}{rgb}{0.000000,0.000000,0.000000}%
\pgfsetstrokecolor{currentstroke}%
\pgfsetdash{}{0pt}%
\pgfpathmoveto{\pgfqpoint{1.406957in}{0.942900in}}%
\pgfpathlineto{\pgfqpoint{1.425319in}{0.954603in}}%
\pgfpathlineto{\pgfqpoint{1.443765in}{0.966498in}}%
\pgfpathlineto{\pgfqpoint{1.451875in}{0.966209in}}%
\pgfpathlineto{\pgfqpoint{1.459956in}{0.966189in}}%
\pgfpathlineto{\pgfqpoint{1.441542in}{0.954037in}}%
\pgfpathlineto{\pgfqpoint{1.423212in}{0.942065in}}%
\pgfpathlineto{\pgfqpoint{1.415098in}{0.942344in}}%
\pgfpathlineto{\pgfqpoint{1.406957in}{0.942900in}}%
\pgfpathclose%
\pgfusepath{fill}%
\end{pgfscope}%
\begin{pgfscope}%
\pgfpathrectangle{\pgfqpoint{0.407592in}{0.222904in}}{\pgfqpoint{1.411163in}{1.411163in}}%
\pgfusepath{clip}%
\pgfsetbuttcap%
\pgfsetroundjoin%
\definecolor{currentfill}{rgb}{0.556863,0.654902,0.847059}%
\pgfsetfillcolor{currentfill}%
\pgfsetfillopacity{0.130000}%
\pgfsetlinewidth{0.000000pt}%
\definecolor{currentstroke}{rgb}{0.000000,0.000000,0.000000}%
\pgfsetstrokecolor{currentstroke}%
\pgfsetdash{}{0pt}%
\pgfpathmoveto{\pgfqpoint{1.513064in}{0.992222in}}%
\pgfpathlineto{\pgfqpoint{1.531706in}{1.005064in}}%
\pgfpathlineto{\pgfqpoint{1.539664in}{1.006181in}}%
\pgfpathlineto{\pgfqpoint{1.547598in}{1.007570in}}%
\pgfpathlineto{\pgfqpoint{1.528986in}{0.994547in}}%
\pgfpathlineto{\pgfqpoint{1.521038in}{0.993246in}}%
\pgfusepath{fill}%
\end{pgfscope}%
\begin{pgfscope}%
\pgfpathrectangle{\pgfqpoint{0.407592in}{0.222904in}}{\pgfqpoint{1.411163in}{1.411163in}}%
\pgfusepath{clip}%
\pgfsetbuttcap%
\pgfsetroundjoin%
\definecolor{currentfill}{rgb}{0.556863,0.654902,0.847059}%
\pgfsetfillcolor{currentfill}%
\pgfsetfillopacity{0.130000}%
\pgfsetlinewidth{0.000000pt}%
\definecolor{currentstroke}{rgb}{0.000000,0.000000,0.000000}%
\pgfsetstrokecolor{currentstroke}%
\pgfsetdash{}{0pt}%
\pgfpathmoveto{\pgfqpoint{1.459956in}{0.966189in}}%
\pgfpathlineto{\pgfqpoint{1.478456in}{0.978511in}}%
\pgfpathlineto{\pgfqpoint{1.497042in}{0.990995in}}%
\pgfpathlineto{\pgfqpoint{1.505066in}{0.991474in}}%
\pgfpathlineto{\pgfqpoint{1.513064in}{0.992222in}}%
\pgfpathlineto{\pgfqpoint{1.494509in}{0.979525in}}%
\pgfpathlineto{\pgfqpoint{1.476040in}{0.966978in}}%
\pgfpathlineto{\pgfqpoint{1.468011in}{0.966444in}}%
\pgfpathlineto{\pgfqpoint{1.459956in}{0.966189in}}%
\pgfpathclose%
\pgfusepath{fill}%
\end{pgfscope}%
\begin{pgfscope}%
\pgfpathrectangle{\pgfqpoint{0.407592in}{0.222904in}}{\pgfqpoint{1.411163in}{1.411163in}}%
\pgfusepath{clip}%
\pgfsetbuttcap%
\pgfsetroundjoin%
\definecolor{currentfill}{rgb}{0.556863,0.654902,0.847059}%
\pgfsetfillcolor{currentfill}%
\pgfsetfillopacity{0.130000}%
\pgfsetlinewidth{0.000000pt}%
\definecolor{currentstroke}{rgb}{0.000000,0.000000,0.000000}%
\pgfsetstrokecolor{currentstroke}%
\pgfsetdash{}{0pt}%
\pgfpathmoveto{\pgfqpoint{1.284428in}{0.910956in}}%
\pgfpathlineto{\pgfqpoint{1.302569in}{0.920727in}}%
\pgfpathlineto{\pgfqpoint{1.320789in}{0.930781in}}%
\pgfpathlineto{\pgfqpoint{1.329152in}{0.928381in}}%
\pgfpathlineto{\pgfqpoint{1.337482in}{0.926220in}}%
\pgfpathlineto{\pgfqpoint{1.319298in}{0.915825in}}%
\pgfpathlineto{\pgfqpoint{1.301196in}{0.905701in}}%
\pgfpathlineto{\pgfqpoint{1.292829in}{0.908209in}}%
\pgfpathlineto{\pgfqpoint{1.284428in}{0.910956in}}%
\pgfpathclose%
\pgfusepath{fill}%
\end{pgfscope}%
\begin{pgfscope}%
\pgfpathrectangle{\pgfqpoint{0.407592in}{0.222904in}}{\pgfqpoint{1.411163in}{1.411163in}}%
\pgfusepath{clip}%
\pgfsetbuttcap%
\pgfsetroundjoin%
\definecolor{currentfill}{rgb}{0.556863,0.654902,0.847059}%
\pgfsetfillcolor{currentfill}%
\pgfsetfillopacity{0.130000}%
\pgfsetlinewidth{0.000000pt}%
\definecolor{currentstroke}{rgb}{0.000000,0.000000,0.000000}%
\pgfsetstrokecolor{currentstroke}%
\pgfsetdash{}{0pt}%
\pgfpathmoveto{\pgfqpoint{1.231398in}{0.899249in}}%
\pgfpathlineto{\pgfqpoint{1.249422in}{0.908029in}}%
\pgfpathlineto{\pgfqpoint{1.267523in}{0.917137in}}%
\pgfpathlineto{\pgfqpoint{1.275993in}{0.913935in}}%
\pgfpathlineto{\pgfqpoint{1.284428in}{0.910956in}}%
\pgfpathlineto{\pgfqpoint{1.266367in}{0.901485in}}%
\pgfpathlineto{\pgfqpoint{1.248383in}{0.892330in}}%
\pgfpathlineto{\pgfqpoint{1.239909in}{0.895679in}}%
\pgfpathlineto{\pgfqpoint{1.231398in}{0.899249in}}%
\pgfpathclose%
\pgfusepath{fill}%
\end{pgfscope}%
\begin{pgfscope}%
\pgfpathrectangle{\pgfqpoint{0.407592in}{0.222904in}}{\pgfqpoint{1.411163in}{1.411163in}}%
\pgfusepath{clip}%
\pgfsetbuttcap%
\pgfsetroundjoin%
\definecolor{currentfill}{rgb}{0.556863,0.654902,0.847059}%
\pgfsetfillcolor{currentfill}%
\pgfsetfillopacity{0.130000}%
\pgfsetlinewidth{0.000000pt}%
\definecolor{currentstroke}{rgb}{0.000000,0.000000,0.000000}%
\pgfsetstrokecolor{currentstroke}%
\pgfsetdash{}{0pt}%
\pgfpathmoveto{\pgfqpoint{1.337482in}{0.926220in}}%
\pgfpathlineto{\pgfqpoint{1.355746in}{0.936870in}}%
\pgfpathlineto{\pgfqpoint{1.374093in}{0.947763in}}%
\pgfpathlineto{\pgfqpoint{1.382355in}{0.946165in}}%
\pgfpathlineto{\pgfqpoint{1.390585in}{0.944818in}}%
\pgfpathlineto{\pgfqpoint{1.372274in}{0.933616in}}%
\pgfpathlineto{\pgfqpoint{1.354044in}{0.922644in}}%
\pgfpathlineto{\pgfqpoint{1.345779in}{0.924305in}}%
\pgfpathlineto{\pgfqpoint{1.337482in}{0.926220in}}%
\pgfpathclose%
\pgfusepath{fill}%
\end{pgfscope}%
\begin{pgfscope}%
\pgfpathrectangle{\pgfqpoint{0.407592in}{0.222904in}}{\pgfqpoint{1.411163in}{1.411163in}}%
\pgfusepath{clip}%
\pgfsetbuttcap%
\pgfsetroundjoin%
\definecolor{currentfill}{rgb}{0.556863,0.654902,0.847059}%
\pgfsetfillcolor{currentfill}%
\pgfsetfillopacity{0.130000}%
\pgfsetlinewidth{0.000000pt}%
\definecolor{currentstroke}{rgb}{0.000000,0.000000,0.000000}%
\pgfsetstrokecolor{currentstroke}%
\pgfsetdash{}{0pt}%
\pgfpathmoveto{\pgfqpoint{1.497042in}{0.990995in}}%
\pgfpathlineto{\pgfqpoint{1.515715in}{1.003635in}}%
\pgfpathlineto{\pgfqpoint{1.523724in}{1.004217in}}%
\pgfpathlineto{\pgfqpoint{1.531706in}{1.005064in}}%
\pgfpathlineto{\pgfqpoint{1.513064in}{0.992222in}}%
\pgfpathlineto{\pgfqpoint{1.505066in}{0.991474in}}%
\pgfusepath{fill}%
\end{pgfscope}%
\begin{pgfscope}%
\pgfpathrectangle{\pgfqpoint{0.407592in}{0.222904in}}{\pgfqpoint{1.411163in}{1.411163in}}%
\pgfusepath{clip}%
\pgfsetbuttcap%
\pgfsetroundjoin%
\definecolor{currentfill}{rgb}{0.556863,0.654902,0.847059}%
\pgfsetfillcolor{currentfill}%
\pgfsetfillopacity{0.130000}%
\pgfsetlinewidth{0.000000pt}%
\definecolor{currentstroke}{rgb}{0.000000,0.000000,0.000000}%
\pgfsetstrokecolor{currentstroke}%
\pgfsetdash{}{0pt}%
\pgfpathmoveto{\pgfqpoint{1.390585in}{0.944818in}}%
\pgfpathlineto{\pgfqpoint{1.408981in}{0.956235in}}%
\pgfpathlineto{\pgfqpoint{1.427461in}{0.967859in}}%
\pgfpathlineto{\pgfqpoint{1.435628in}{0.967050in}}%
\pgfpathlineto{\pgfqpoint{1.443765in}{0.966498in}}%
\pgfpathlineto{\pgfqpoint{1.425319in}{0.954603in}}%
\pgfpathlineto{\pgfqpoint{1.406957in}{0.942900in}}%
\pgfpathlineto{\pgfqpoint{1.398786in}{0.943727in}}%
\pgfpathlineto{\pgfqpoint{1.390585in}{0.944818in}}%
\pgfpathclose%
\pgfusepath{fill}%
\end{pgfscope}%
\begin{pgfscope}%
\pgfpathrectangle{\pgfqpoint{0.407592in}{0.222904in}}{\pgfqpoint{1.411163in}{1.411163in}}%
\pgfusepath{clip}%
\pgfsetbuttcap%
\pgfsetroundjoin%
\definecolor{currentfill}{rgb}{0.556863,0.654902,0.847059}%
\pgfsetfillcolor{currentfill}%
\pgfsetfillopacity{0.130000}%
\pgfsetlinewidth{0.000000pt}%
\definecolor{currentstroke}{rgb}{0.000000,0.000000,0.000000}%
\pgfsetstrokecolor{currentstroke}%
\pgfsetdash{}{0pt}%
\pgfpathmoveto{\pgfqpoint{1.443765in}{0.966498in}}%
\pgfpathlineto{\pgfqpoint{1.462297in}{0.978576in}}%
\pgfpathlineto{\pgfqpoint{1.480915in}{0.990829in}}%
\pgfpathlineto{\pgfqpoint{1.488992in}{0.990782in}}%
\pgfpathlineto{\pgfqpoint{1.497042in}{0.990995in}}%
\pgfpathlineto{\pgfqpoint{1.478456in}{0.978511in}}%
\pgfpathlineto{\pgfqpoint{1.459956in}{0.966189in}}%
\pgfpathlineto{\pgfqpoint{1.451875in}{0.966209in}}%
\pgfpathlineto{\pgfqpoint{1.443765in}{0.966498in}}%
\pgfpathclose%
\pgfusepath{fill}%
\end{pgfscope}%
\begin{pgfscope}%
\pgfpathrectangle{\pgfqpoint{0.407592in}{0.222904in}}{\pgfqpoint{1.411163in}{1.411163in}}%
\pgfusepath{clip}%
\pgfsetbuttcap%
\pgfsetroundjoin%
\definecolor{currentfill}{rgb}{0.556863,0.654902,0.847059}%
\pgfsetfillcolor{currentfill}%
\pgfsetfillopacity{0.130000}%
\pgfsetlinewidth{0.000000pt}%
\definecolor{currentstroke}{rgb}{0.000000,0.000000,0.000000}%
\pgfsetstrokecolor{currentstroke}%
\pgfsetdash{}{0pt}%
\pgfpathmoveto{\pgfqpoint{1.480915in}{0.990829in}}%
\pgfpathlineto{\pgfqpoint{1.499620in}{1.003250in}}%
\pgfpathlineto{\pgfqpoint{1.507681in}{1.003314in}}%
\pgfpathlineto{\pgfqpoint{1.515715in}{1.003635in}}%
\pgfpathlineto{\pgfqpoint{1.497042in}{0.990995in}}%
\pgfpathlineto{\pgfqpoint{1.488992in}{0.990782in}}%
\pgfusepath{fill}%
\end{pgfscope}%
\begin{pgfscope}%
\pgfpathrectangle{\pgfqpoint{0.407592in}{0.222904in}}{\pgfqpoint{1.411163in}{1.411163in}}%
\pgfusepath{clip}%
\pgfsetbuttcap%
\pgfsetroundjoin%
\definecolor{currentfill}{rgb}{0.556863,0.654902,0.847059}%
\pgfsetfillcolor{currentfill}%
\pgfsetfillopacity{0.130000}%
\pgfsetlinewidth{0.000000pt}%
\definecolor{currentstroke}{rgb}{0.000000,0.000000,0.000000}%
\pgfsetstrokecolor{currentstroke}%
\pgfsetdash{}{0pt}%
\pgfpathmoveto{\pgfqpoint{1.267523in}{0.917137in}}%
\pgfpathlineto{\pgfqpoint{1.285703in}{0.926555in}}%
\pgfpathlineto{\pgfqpoint{1.303963in}{0.936268in}}%
\pgfpathlineto{\pgfqpoint{1.312393in}{0.933413in}}%
\pgfpathlineto{\pgfqpoint{1.320789in}{0.930781in}}%
\pgfpathlineto{\pgfqpoint{1.302569in}{0.920727in}}%
\pgfpathlineto{\pgfqpoint{1.284428in}{0.910956in}}%
\pgfpathlineto{\pgfqpoint{1.275993in}{0.913935in}}%
\pgfpathlineto{\pgfqpoint{1.267523in}{0.917137in}}%
\pgfpathclose%
\pgfusepath{fill}%
\end{pgfscope}%
\begin{pgfscope}%
\pgfpathrectangle{\pgfqpoint{0.407592in}{0.222904in}}{\pgfqpoint{1.411163in}{1.411163in}}%
\pgfusepath{clip}%
\pgfsetbuttcap%
\pgfsetroundjoin%
\definecolor{currentfill}{rgb}{0.556863,0.654902,0.847059}%
\pgfsetfillcolor{currentfill}%
\pgfsetfillopacity{0.130000}%
\pgfsetlinewidth{0.000000pt}%
\definecolor{currentstroke}{rgb}{0.000000,0.000000,0.000000}%
\pgfsetstrokecolor{currentstroke}%
\pgfsetdash{}{0pt}%
\pgfpathmoveto{\pgfqpoint{1.320789in}{0.930781in}}%
\pgfpathlineto{\pgfqpoint{1.339091in}{0.941104in}}%
\pgfpathlineto{\pgfqpoint{1.357475in}{0.951681in}}%
\pgfpathlineto{\pgfqpoint{1.365800in}{0.949604in}}%
\pgfpathlineto{\pgfqpoint{1.374093in}{0.947763in}}%
\pgfpathlineto{\pgfqpoint{1.355746in}{0.936870in}}%
\pgfpathlineto{\pgfqpoint{1.337482in}{0.926220in}}%
\pgfpathlineto{\pgfqpoint{1.329152in}{0.928381in}}%
\pgfpathlineto{\pgfqpoint{1.320789in}{0.930781in}}%
\pgfpathclose%
\pgfusepath{fill}%
\end{pgfscope}%
\begin{pgfscope}%
\pgfpathrectangle{\pgfqpoint{0.407592in}{0.222904in}}{\pgfqpoint{1.411163in}{1.411163in}}%
\pgfusepath{clip}%
\pgfsetbuttcap%
\pgfsetroundjoin%
\definecolor{currentfill}{rgb}{0.556863,0.654902,0.847059}%
\pgfsetfillcolor{currentfill}%
\pgfsetfillopacity{0.130000}%
\pgfsetlinewidth{0.000000pt}%
\definecolor{currentstroke}{rgb}{0.000000,0.000000,0.000000}%
\pgfsetstrokecolor{currentstroke}%
\pgfsetdash{}{0pt}%
\pgfpathmoveto{\pgfqpoint{1.374093in}{0.947763in}}%
\pgfpathlineto{\pgfqpoint{1.392524in}{0.958885in}}%
\pgfpathlineto{\pgfqpoint{1.411039in}{0.970226in}}%
\pgfpathlineto{\pgfqpoint{1.419265in}{0.968920in}}%
\pgfpathlineto{\pgfqpoint{1.427461in}{0.967859in}}%
\pgfpathlineto{\pgfqpoint{1.408981in}{0.956235in}}%
\pgfpathlineto{\pgfqpoint{1.390585in}{0.944818in}}%
\pgfpathlineto{\pgfqpoint{1.382355in}{0.946165in}}%
\pgfpathlineto{\pgfqpoint{1.374093in}{0.947763in}}%
\pgfpathclose%
\pgfusepath{fill}%
\end{pgfscope}%
\begin{pgfscope}%
\pgfpathrectangle{\pgfqpoint{0.407592in}{0.222904in}}{\pgfqpoint{1.411163in}{1.411163in}}%
\pgfusepath{clip}%
\pgfsetbuttcap%
\pgfsetroundjoin%
\definecolor{currentfill}{rgb}{0.556863,0.654902,0.847059}%
\pgfsetfillcolor{currentfill}%
\pgfsetfillopacity{0.130000}%
\pgfsetlinewidth{0.000000pt}%
\definecolor{currentstroke}{rgb}{0.000000,0.000000,0.000000}%
\pgfsetstrokecolor{currentstroke}%
\pgfsetdash{}{0pt}%
\pgfpathmoveto{\pgfqpoint{1.427461in}{0.967859in}}%
\pgfpathlineto{\pgfqpoint{1.446026in}{0.979679in}}%
\pgfpathlineto{\pgfqpoint{1.464678in}{0.991685in}}%
\pgfpathlineto{\pgfqpoint{1.472811in}{0.991132in}}%
\pgfpathlineto{\pgfqpoint{1.480915in}{0.990829in}}%
\pgfpathlineto{\pgfqpoint{1.462297in}{0.978576in}}%
\pgfpathlineto{\pgfqpoint{1.443765in}{0.966498in}}%
\pgfpathlineto{\pgfqpoint{1.435628in}{0.967050in}}%
\pgfpathlineto{\pgfqpoint{1.427461in}{0.967859in}}%
\pgfpathclose%
\pgfusepath{fill}%
\end{pgfscope}%
\begin{pgfscope}%
\pgfpathrectangle{\pgfqpoint{0.407592in}{0.222904in}}{\pgfqpoint{1.411163in}{1.411163in}}%
\pgfusepath{clip}%
\pgfsetbuttcap%
\pgfsetroundjoin%
\definecolor{currentfill}{rgb}{0.556863,0.654902,0.847059}%
\pgfsetfillcolor{currentfill}%
\pgfsetfillopacity{0.130000}%
\pgfsetlinewidth{0.000000pt}%
\definecolor{currentstroke}{rgb}{0.000000,0.000000,0.000000}%
\pgfsetstrokecolor{currentstroke}%
\pgfsetdash{}{0pt}%
\pgfpathmoveto{\pgfqpoint{1.464678in}{0.991685in}}%
\pgfpathlineto{\pgfqpoint{1.483416in}{1.003871in}}%
\pgfpathlineto{\pgfqpoint{1.491532in}{1.003437in}}%
\pgfpathlineto{\pgfqpoint{1.499620in}{1.003250in}}%
\pgfpathlineto{\pgfqpoint{1.480915in}{0.990829in}}%
\pgfpathlineto{\pgfqpoint{1.472811in}{0.991132in}}%
\pgfusepath{fill}%
\end{pgfscope}%
\begin{pgfscope}%
\pgfpathrectangle{\pgfqpoint{0.407592in}{0.222904in}}{\pgfqpoint{1.411163in}{1.411163in}}%
\pgfusepath{clip}%
\pgfsetbuttcap%
\pgfsetroundjoin%
\definecolor{currentfill}{rgb}{0.556863,0.654902,0.847059}%
\pgfsetfillcolor{currentfill}%
\pgfsetfillopacity{0.130000}%
\pgfsetlinewidth{0.000000pt}%
\definecolor{currentstroke}{rgb}{0.000000,0.000000,0.000000}%
\pgfsetstrokecolor{currentstroke}%
\pgfsetdash{}{0pt}%
\pgfpathmoveto{\pgfqpoint{1.303963in}{0.936268in}}%
\pgfpathlineto{\pgfqpoint{1.322303in}{0.946261in}}%
\pgfpathlineto{\pgfqpoint{1.340726in}{0.956521in}}%
\pgfpathlineto{\pgfqpoint{1.349117in}{0.953989in}}%
\pgfpathlineto{\pgfqpoint{1.357475in}{0.951681in}}%
\pgfpathlineto{\pgfqpoint{1.339091in}{0.941104in}}%
\pgfpathlineto{\pgfqpoint{1.320789in}{0.930781in}}%
\pgfpathlineto{\pgfqpoint{1.312393in}{0.933413in}}%
\pgfpathlineto{\pgfqpoint{1.303963in}{0.936268in}}%
\pgfpathclose%
\pgfusepath{fill}%
\end{pgfscope}%
\begin{pgfscope}%
\pgfpathrectangle{\pgfqpoint{0.407592in}{0.222904in}}{\pgfqpoint{1.411163in}{1.411163in}}%
\pgfusepath{clip}%
\pgfsetbuttcap%
\pgfsetroundjoin%
\definecolor{currentfill}{rgb}{0.556863,0.654902,0.847059}%
\pgfsetfillcolor{currentfill}%
\pgfsetfillopacity{0.130000}%
\pgfsetlinewidth{0.000000pt}%
\definecolor{currentstroke}{rgb}{0.000000,0.000000,0.000000}%
\pgfsetstrokecolor{currentstroke}%
\pgfsetdash{}{0pt}%
\pgfpathmoveto{\pgfqpoint{1.357475in}{0.951681in}}%
\pgfpathlineto{\pgfqpoint{1.375942in}{0.962501in}}%
\pgfpathlineto{\pgfqpoint{1.394493in}{0.973552in}}%
\pgfpathlineto{\pgfqpoint{1.402782in}{0.971772in}}%
\pgfpathlineto{\pgfqpoint{1.411039in}{0.970226in}}%
\pgfpathlineto{\pgfqpoint{1.392524in}{0.958885in}}%
\pgfpathlineto{\pgfqpoint{1.374093in}{0.947763in}}%
\pgfpathlineto{\pgfqpoint{1.365800in}{0.949604in}}%
\pgfpathlineto{\pgfqpoint{1.357475in}{0.951681in}}%
\pgfpathclose%
\pgfusepath{fill}%
\end{pgfscope}%
\begin{pgfscope}%
\pgfpathrectangle{\pgfqpoint{0.407592in}{0.222904in}}{\pgfqpoint{1.411163in}{1.411163in}}%
\pgfusepath{clip}%
\pgfsetbuttcap%
\pgfsetroundjoin%
\definecolor{currentfill}{rgb}{0.556863,0.654902,0.847059}%
\pgfsetfillcolor{currentfill}%
\pgfsetfillopacity{0.130000}%
\pgfsetlinewidth{0.000000pt}%
\definecolor{currentstroke}{rgb}{0.000000,0.000000,0.000000}%
\pgfsetstrokecolor{currentstroke}%
\pgfsetdash{}{0pt}%
\pgfpathmoveto{\pgfqpoint{1.411039in}{0.970226in}}%
\pgfpathlineto{\pgfqpoint{1.429639in}{0.981775in}}%
\pgfpathlineto{\pgfqpoint{1.448325in}{0.993524in}}%
\pgfpathlineto{\pgfqpoint{1.456516in}{0.992484in}}%
\pgfpathlineto{\pgfqpoint{1.464678in}{0.991685in}}%
\pgfpathlineto{\pgfqpoint{1.446026in}{0.979679in}}%
\pgfpathlineto{\pgfqpoint{1.427461in}{0.967859in}}%
\pgfpathlineto{\pgfqpoint{1.419265in}{0.968920in}}%
\pgfpathlineto{\pgfqpoint{1.411039in}{0.970226in}}%
\pgfpathclose%
\pgfusepath{fill}%
\end{pgfscope}%
\begin{pgfscope}%
\pgfpathrectangle{\pgfqpoint{0.407592in}{0.222904in}}{\pgfqpoint{1.411163in}{1.411163in}}%
\pgfusepath{clip}%
\pgfsetbuttcap%
\pgfsetroundjoin%
\definecolor{currentfill}{rgb}{0.556863,0.654902,0.847059}%
\pgfsetfillcolor{currentfill}%
\pgfsetfillopacity{0.130000}%
\pgfsetlinewidth{0.000000pt}%
\definecolor{currentstroke}{rgb}{0.000000,0.000000,0.000000}%
\pgfsetstrokecolor{currentstroke}%
\pgfsetdash{}{0pt}%
\pgfpathmoveto{\pgfqpoint{1.448325in}{0.993524in}}%
\pgfpathlineto{\pgfqpoint{1.467098in}{1.005463in}}%
\pgfpathlineto{\pgfqpoint{1.475272in}{1.004548in}}%
\pgfpathlineto{\pgfqpoint{1.483416in}{1.003871in}}%
\pgfpathlineto{\pgfqpoint{1.464678in}{0.991685in}}%
\pgfpathlineto{\pgfqpoint{1.456516in}{0.992484in}}%
\pgfusepath{fill}%
\end{pgfscope}%
\begin{pgfscope}%
\pgfpathrectangle{\pgfqpoint{0.407592in}{0.222904in}}{\pgfqpoint{1.411163in}{1.411163in}}%
\pgfusepath{clip}%
\pgfsetbuttcap%
\pgfsetroundjoin%
\definecolor{currentfill}{rgb}{0.556863,0.654902,0.847059}%
\pgfsetfillcolor{currentfill}%
\pgfsetfillopacity{0.130000}%
\pgfsetlinewidth{0.000000pt}%
\definecolor{currentstroke}{rgb}{0.000000,0.000000,0.000000}%
\pgfsetstrokecolor{currentstroke}%
\pgfsetdash{}{0pt}%
\pgfpathmoveto{\pgfqpoint{1.340726in}{0.956521in}}%
\pgfpathlineto{\pgfqpoint{1.359231in}{0.967034in}}%
\pgfpathlineto{\pgfqpoint{1.377820in}{0.977789in}}%
\pgfpathlineto{\pgfqpoint{1.386173in}{0.975559in}}%
\pgfpathlineto{\pgfqpoint{1.394493in}{0.973552in}}%
\pgfpathlineto{\pgfqpoint{1.375942in}{0.962501in}}%
\pgfpathlineto{\pgfqpoint{1.357475in}{0.951681in}}%
\pgfpathlineto{\pgfqpoint{1.349117in}{0.953989in}}%
\pgfpathlineto{\pgfqpoint{1.340726in}{0.956521in}}%
\pgfpathclose%
\pgfusepath{fill}%
\end{pgfscope}%
\begin{pgfscope}%
\pgfpathrectangle{\pgfqpoint{0.407592in}{0.222904in}}{\pgfqpoint{1.411163in}{1.411163in}}%
\pgfusepath{clip}%
\pgfsetbuttcap%
\pgfsetroundjoin%
\definecolor{currentfill}{rgb}{0.556863,0.654902,0.847059}%
\pgfsetfillcolor{currentfill}%
\pgfsetfillopacity{0.130000}%
\pgfsetlinewidth{0.000000pt}%
\definecolor{currentstroke}{rgb}{0.000000,0.000000,0.000000}%
\pgfsetstrokecolor{currentstroke}%
\pgfsetdash{}{0pt}%
\pgfpathmoveto{\pgfqpoint{1.394493in}{0.973552in}}%
\pgfpathlineto{\pgfqpoint{1.413130in}{0.984822in}}%
\pgfpathlineto{\pgfqpoint{1.431852in}{0.996304in}}%
\pgfpathlineto{\pgfqpoint{1.440104in}{0.994799in}}%
\pgfpathlineto{\pgfqpoint{1.448325in}{0.993524in}}%
\pgfpathlineto{\pgfqpoint{1.429639in}{0.981775in}}%
\pgfpathlineto{\pgfqpoint{1.411039in}{0.970226in}}%
\pgfpathlineto{\pgfqpoint{1.402782in}{0.971772in}}%
\pgfpathlineto{\pgfqpoint{1.394493in}{0.973552in}}%
\pgfpathclose%
\pgfusepath{fill}%
\end{pgfscope}%
\begin{pgfscope}%
\pgfpathrectangle{\pgfqpoint{0.407592in}{0.222904in}}{\pgfqpoint{1.411163in}{1.411163in}}%
\pgfusepath{clip}%
\pgfsetbuttcap%
\pgfsetroundjoin%
\definecolor{currentfill}{rgb}{0.556863,0.654902,0.847059}%
\pgfsetfillcolor{currentfill}%
\pgfsetfillopacity{0.130000}%
\pgfsetlinewidth{0.000000pt}%
\definecolor{currentstroke}{rgb}{0.000000,0.000000,0.000000}%
\pgfsetstrokecolor{currentstroke}%
\pgfsetdash{}{0pt}%
\pgfpathmoveto{\pgfqpoint{1.431852in}{0.996304in}}%
\pgfpathlineto{\pgfqpoint{1.450661in}{1.007987in}}%
\pgfpathlineto{\pgfqpoint{1.458895in}{1.006611in}}%
\pgfpathlineto{\pgfqpoint{1.467098in}{1.005463in}}%
\pgfpathlineto{\pgfqpoint{1.448325in}{0.993524in}}%
\pgfpathlineto{\pgfqpoint{1.440104in}{0.994799in}}%
\pgfusepath{fill}%
\end{pgfscope}%
\begin{pgfscope}%
\pgfpathrectangle{\pgfqpoint{0.407592in}{0.222904in}}{\pgfqpoint{1.411163in}{1.411163in}}%
\pgfusepath{clip}%
\pgfsetbuttcap%
\pgfsetroundjoin%
\definecolor{currentfill}{rgb}{0.556863,0.654902,0.847059}%
\pgfsetfillcolor{currentfill}%
\pgfsetfillopacity{0.130000}%
\pgfsetlinewidth{0.000000pt}%
\definecolor{currentstroke}{rgb}{0.000000,0.000000,0.000000}%
\pgfsetstrokecolor{currentstroke}%
\pgfsetdash{}{0pt}%
\pgfpathmoveto{\pgfqpoint{1.377820in}{0.977789in}}%
\pgfpathlineto{\pgfqpoint{1.396494in}{0.988775in}}%
\pgfpathlineto{\pgfqpoint{1.415254in}{0.999983in}}%
\pgfpathlineto{\pgfqpoint{1.423569in}{0.998033in}}%
\pgfpathlineto{\pgfqpoint{1.431852in}{0.996304in}}%
\pgfpathlineto{\pgfqpoint{1.413130in}{0.984822in}}%
\pgfpathlineto{\pgfqpoint{1.394493in}{0.973552in}}%
\pgfpathlineto{\pgfqpoint{1.386173in}{0.975559in}}%
\pgfpathlineto{\pgfqpoint{1.377820in}{0.977789in}}%
\pgfpathclose%
\pgfusepath{fill}%
\end{pgfscope}%
\begin{pgfscope}%
\pgfpathrectangle{\pgfqpoint{0.407592in}{0.222904in}}{\pgfqpoint{1.411163in}{1.411163in}}%
\pgfusepath{clip}%
\pgfsetbuttcap%
\pgfsetroundjoin%
\definecolor{currentfill}{rgb}{0.556863,0.654902,0.847059}%
\pgfsetfillcolor{currentfill}%
\pgfsetfillopacity{0.130000}%
\pgfsetlinewidth{0.000000pt}%
\definecolor{currentstroke}{rgb}{0.000000,0.000000,0.000000}%
\pgfsetstrokecolor{currentstroke}%
\pgfsetdash{}{0pt}%
\pgfpathmoveto{\pgfqpoint{1.415254in}{0.999983in}}%
\pgfpathlineto{\pgfqpoint{1.434100in}{1.011403in}}%
\pgfpathlineto{\pgfqpoint{1.442396in}{1.009586in}}%
\pgfpathlineto{\pgfqpoint{1.450661in}{1.007987in}}%
\pgfpathlineto{\pgfqpoint{1.431852in}{0.996304in}}%
\pgfpathlineto{\pgfqpoint{1.423569in}{0.998033in}}%
\pgfusepath{fill}%
\end{pgfscope}%
\begin{pgfscope}%
\pgfpathrectangle{\pgfqpoint{0.407592in}{0.222904in}}{\pgfqpoint{1.411163in}{1.411163in}}%
\pgfusepath{clip}%
\pgfsetbuttcap%
\pgfsetroundjoin%
\definecolor{currentfill}{rgb}{0.909804,0.698039,0.423529}%
\pgfsetfillcolor{currentfill}%
\pgfsetfillopacity{0.120000}%
\pgfsetlinewidth{0.000000pt}%
\definecolor{currentstroke}{rgb}{0.000000,0.000000,0.000000}%
\pgfsetstrokecolor{currentstroke}%
\pgfsetdash{}{0pt}%
\pgfpathmoveto{\pgfqpoint{0.870800in}{1.442368in}}%
\pgfpathlineto{\pgfqpoint{0.886196in}{1.427129in}}%
\pgfpathlineto{\pgfqpoint{0.901605in}{1.412100in}}%
\pgfpathlineto{\pgfqpoint{0.909495in}{1.423095in}}%
\pgfpathlineto{\pgfqpoint{0.894100in}{1.437968in}}%
\pgfpathlineto{\pgfqpoint{0.878719in}{1.453053in}}%
\pgfusepath{fill}%
\end{pgfscope}%
\begin{pgfscope}%
\pgfpathrectangle{\pgfqpoint{0.407592in}{0.222904in}}{\pgfqpoint{1.411163in}{1.411163in}}%
\pgfusepath{clip}%
\pgfsetbuttcap%
\pgfsetroundjoin%
\definecolor{currentfill}{rgb}{0.909804,0.698039,0.423529}%
\pgfsetfillcolor{currentfill}%
\pgfsetfillopacity{0.120000}%
\pgfsetlinewidth{0.000000pt}%
\definecolor{currentstroke}{rgb}{0.000000,0.000000,0.000000}%
\pgfsetstrokecolor{currentstroke}%
\pgfsetdash{}{0pt}%
\pgfpathmoveto{\pgfqpoint{0.901605in}{1.412100in}}%
\pgfpathlineto{\pgfqpoint{0.917027in}{1.397294in}}%
\pgfpathlineto{\pgfqpoint{0.932464in}{1.382727in}}%
\pgfpathlineto{\pgfqpoint{0.940326in}{1.394042in}}%
\pgfpathlineto{\pgfqpoint{0.924904in}{1.408449in}}%
\pgfpathlineto{\pgfqpoint{0.909495in}{1.423095in}}%
\pgfusepath{fill}%
\end{pgfscope}%
\begin{pgfscope}%
\pgfpathrectangle{\pgfqpoint{0.407592in}{0.222904in}}{\pgfqpoint{1.411163in}{1.411163in}}%
\pgfusepath{clip}%
\pgfsetbuttcap%
\pgfsetroundjoin%
\definecolor{currentfill}{rgb}{0.909804,0.698039,0.423529}%
\pgfsetfillcolor{currentfill}%
\pgfsetfillopacity{0.120000}%
\pgfsetlinewidth{0.000000pt}%
\definecolor{currentstroke}{rgb}{0.000000,0.000000,0.000000}%
\pgfsetstrokecolor{currentstroke}%
\pgfsetdash{}{0pt}%
\pgfpathmoveto{\pgfqpoint{0.854899in}{1.421494in}}%
\pgfpathlineto{\pgfqpoint{0.870325in}{1.405949in}}%
\pgfpathlineto{\pgfqpoint{0.885764in}{1.390608in}}%
\pgfpathlineto{\pgfqpoint{0.893694in}{1.401268in}}%
\pgfpathlineto{\pgfqpoint{0.901605in}{1.412100in}}%
\pgfpathlineto{\pgfqpoint{0.886196in}{1.427129in}}%
\pgfpathlineto{\pgfqpoint{0.870800in}{1.442368in}}%
\pgfpathlineto{\pgfqpoint{0.862860in}{1.431846in}}%
\pgfpathlineto{\pgfqpoint{0.854899in}{1.421494in}}%
\pgfpathclose%
\pgfusepath{fill}%
\end{pgfscope}%
\begin{pgfscope}%
\pgfpathrectangle{\pgfqpoint{0.407592in}{0.222904in}}{\pgfqpoint{1.411163in}{1.411163in}}%
\pgfusepath{clip}%
\pgfsetbuttcap%
\pgfsetroundjoin%
\definecolor{currentfill}{rgb}{0.909804,0.698039,0.423529}%
\pgfsetfillcolor{currentfill}%
\pgfsetfillopacity{0.120000}%
\pgfsetlinewidth{0.000000pt}%
\definecolor{currentstroke}{rgb}{0.000000,0.000000,0.000000}%
\pgfsetstrokecolor{currentstroke}%
\pgfsetdash{}{0pt}%
\pgfpathmoveto{\pgfqpoint{0.932464in}{1.382727in}}%
\pgfpathlineto{\pgfqpoint{0.947915in}{1.368413in}}%
\pgfpathlineto{\pgfqpoint{0.963382in}{1.354369in}}%
\pgfpathlineto{\pgfqpoint{0.971217in}{1.366012in}}%
\pgfpathlineto{\pgfqpoint{0.955764in}{1.379892in}}%
\pgfpathlineto{\pgfqpoint{0.940326in}{1.394042in}}%
\pgfusepath{fill}%
\end{pgfscope}%
\begin{pgfscope}%
\pgfpathrectangle{\pgfqpoint{0.407592in}{0.222904in}}{\pgfqpoint{1.411163in}{1.411163in}}%
\pgfusepath{clip}%
\pgfsetbuttcap%
\pgfsetroundjoin%
\definecolor{currentfill}{rgb}{0.909804,0.698039,0.423529}%
\pgfsetfillcolor{currentfill}%
\pgfsetfillopacity{0.120000}%
\pgfsetlinewidth{0.000000pt}%
\definecolor{currentstroke}{rgb}{0.000000,0.000000,0.000000}%
\pgfsetstrokecolor{currentstroke}%
\pgfsetdash{}{0pt}%
\pgfpathmoveto{\pgfqpoint{0.885764in}{1.390608in}}%
\pgfpathlineto{\pgfqpoint{0.901216in}{1.375486in}}%
\pgfpathlineto{\pgfqpoint{0.916681in}{1.360597in}}%
\pgfpathlineto{\pgfqpoint{0.924582in}{1.371576in}}%
\pgfpathlineto{\pgfqpoint{0.932464in}{1.382727in}}%
\pgfpathlineto{\pgfqpoint{0.917027in}{1.397294in}}%
\pgfpathlineto{\pgfqpoint{0.901605in}{1.412100in}}%
\pgfpathlineto{\pgfqpoint{0.893694in}{1.401268in}}%
\pgfpathlineto{\pgfqpoint{0.885764in}{1.390608in}}%
\pgfpathclose%
\pgfusepath{fill}%
\end{pgfscope}%
\begin{pgfscope}%
\pgfpathrectangle{\pgfqpoint{0.407592in}{0.222904in}}{\pgfqpoint{1.411163in}{1.411163in}}%
\pgfusepath{clip}%
\pgfsetbuttcap%
\pgfsetroundjoin%
\definecolor{currentfill}{rgb}{0.909804,0.698039,0.423529}%
\pgfsetfillcolor{currentfill}%
\pgfsetfillopacity{0.120000}%
\pgfsetlinewidth{0.000000pt}%
\definecolor{currentstroke}{rgb}{0.000000,0.000000,0.000000}%
\pgfsetstrokecolor{currentstroke}%
\pgfsetdash{}{0pt}%
\pgfpathmoveto{\pgfqpoint{0.963382in}{1.354369in}}%
\pgfpathlineto{\pgfqpoint{0.978864in}{1.340612in}}%
\pgfpathlineto{\pgfqpoint{0.994363in}{1.327159in}}%
\pgfpathlineto{\pgfqpoint{1.002172in}{1.339134in}}%
\pgfpathlineto{\pgfqpoint{0.986686in}{1.352420in}}%
\pgfpathlineto{\pgfqpoint{0.971217in}{1.366012in}}%
\pgfusepath{fill}%
\end{pgfscope}%
\begin{pgfscope}%
\pgfpathrectangle{\pgfqpoint{0.407592in}{0.222904in}}{\pgfqpoint{1.411163in}{1.411163in}}%
\pgfusepath{clip}%
\pgfsetbuttcap%
\pgfsetroundjoin%
\definecolor{currentfill}{rgb}{0.909804,0.698039,0.423529}%
\pgfsetfillcolor{currentfill}%
\pgfsetfillopacity{0.120000}%
\pgfsetlinewidth{0.000000pt}%
\definecolor{currentstroke}{rgb}{0.000000,0.000000,0.000000}%
\pgfsetstrokecolor{currentstroke}%
\pgfsetdash{}{0pt}%
\pgfpathmoveto{\pgfqpoint{0.916681in}{1.360597in}}%
\pgfpathlineto{\pgfqpoint{0.932161in}{1.345958in}}%
\pgfpathlineto{\pgfqpoint{0.947655in}{1.331584in}}%
\pgfpathlineto{\pgfqpoint{0.955528in}{1.342890in}}%
\pgfpathlineto{\pgfqpoint{0.963382in}{1.354369in}}%
\pgfpathlineto{\pgfqpoint{0.947915in}{1.368413in}}%
\pgfpathlineto{\pgfqpoint{0.932464in}{1.382727in}}%
\pgfpathlineto{\pgfqpoint{0.924582in}{1.371576in}}%
\pgfpathlineto{\pgfqpoint{0.916681in}{1.360597in}}%
\pgfpathclose%
\pgfusepath{fill}%
\end{pgfscope}%
\begin{pgfscope}%
\pgfpathrectangle{\pgfqpoint{0.407592in}{0.222904in}}{\pgfqpoint{1.411163in}{1.411163in}}%
\pgfusepath{clip}%
\pgfsetbuttcap%
\pgfsetroundjoin%
\definecolor{currentfill}{rgb}{0.909804,0.698039,0.423529}%
\pgfsetfillcolor{currentfill}%
\pgfsetfillopacity{0.120000}%
\pgfsetlinewidth{0.000000pt}%
\definecolor{currentstroke}{rgb}{0.000000,0.000000,0.000000}%
\pgfsetstrokecolor{currentstroke}%
\pgfsetdash{}{0pt}%
\pgfpathmoveto{\pgfqpoint{0.838913in}{1.401324in}}%
\pgfpathlineto{\pgfqpoint{0.854371in}{1.385477in}}%
\pgfpathlineto{\pgfqpoint{0.869840in}{1.369829in}}%
\pgfpathlineto{\pgfqpoint{0.877813in}{1.380126in}}%
\pgfpathlineto{\pgfqpoint{0.885764in}{1.390608in}}%
\pgfpathlineto{\pgfqpoint{0.870325in}{1.405949in}}%
\pgfpathlineto{\pgfqpoint{0.854899in}{1.421494in}}%
\pgfpathlineto{\pgfqpoint{0.846917in}{1.411317in}}%
\pgfpathlineto{\pgfqpoint{0.838913in}{1.401324in}}%
\pgfpathclose%
\pgfusepath{fill}%
\end{pgfscope}%
\begin{pgfscope}%
\pgfpathrectangle{\pgfqpoint{0.407592in}{0.222904in}}{\pgfqpoint{1.411163in}{1.411163in}}%
\pgfusepath{clip}%
\pgfsetbuttcap%
\pgfsetroundjoin%
\definecolor{currentfill}{rgb}{0.909804,0.698039,0.423529}%
\pgfsetfillcolor{currentfill}%
\pgfsetfillopacity{0.120000}%
\pgfsetlinewidth{0.000000pt}%
\definecolor{currentstroke}{rgb}{0.000000,0.000000,0.000000}%
\pgfsetstrokecolor{currentstroke}%
\pgfsetdash{}{0pt}%
\pgfpathmoveto{\pgfqpoint{0.994363in}{1.327159in}}%
\pgfpathlineto{\pgfqpoint{1.009879in}{1.314030in}}%
\pgfpathlineto{\pgfqpoint{1.025412in}{1.301244in}}%
\pgfpathlineto{\pgfqpoint{1.033198in}{1.313552in}}%
\pgfpathlineto{\pgfqpoint{1.017676in}{1.326171in}}%
\pgfpathlineto{\pgfqpoint{1.002172in}{1.339134in}}%
\pgfusepath{fill}%
\end{pgfscope}%
\begin{pgfscope}%
\pgfpathrectangle{\pgfqpoint{0.407592in}{0.222904in}}{\pgfqpoint{1.411163in}{1.411163in}}%
\pgfusepath{clip}%
\pgfsetbuttcap%
\pgfsetroundjoin%
\definecolor{currentfill}{rgb}{0.909804,0.698039,0.423529}%
\pgfsetfillcolor{currentfill}%
\pgfsetfillopacity{0.120000}%
\pgfsetlinewidth{0.000000pt}%
\definecolor{currentstroke}{rgb}{0.000000,0.000000,0.000000}%
\pgfsetstrokecolor{currentstroke}%
\pgfsetdash{}{0pt}%
\pgfpathmoveto{\pgfqpoint{0.869840in}{1.369829in}}%
\pgfpathlineto{\pgfqpoint{0.885323in}{1.354394in}}%
\pgfpathlineto{\pgfqpoint{0.900818in}{1.339187in}}%
\pgfpathlineto{\pgfqpoint{0.908760in}{1.349798in}}%
\pgfpathlineto{\pgfqpoint{0.916681in}{1.360597in}}%
\pgfpathlineto{\pgfqpoint{0.901216in}{1.375486in}}%
\pgfpathlineto{\pgfqpoint{0.885764in}{1.390608in}}%
\pgfpathlineto{\pgfqpoint{0.877813in}{1.380126in}}%
\pgfpathlineto{\pgfqpoint{0.869840in}{1.369829in}}%
\pgfpathclose%
\pgfusepath{fill}%
\end{pgfscope}%
\begin{pgfscope}%
\pgfpathrectangle{\pgfqpoint{0.407592in}{0.222904in}}{\pgfqpoint{1.411163in}{1.411163in}}%
\pgfusepath{clip}%
\pgfsetbuttcap%
\pgfsetroundjoin%
\definecolor{currentfill}{rgb}{0.909804,0.698039,0.423529}%
\pgfsetfillcolor{currentfill}%
\pgfsetfillopacity{0.120000}%
\pgfsetlinewidth{0.000000pt}%
\definecolor{currentstroke}{rgb}{0.000000,0.000000,0.000000}%
\pgfsetstrokecolor{currentstroke}%
\pgfsetdash{}{0pt}%
\pgfpathmoveto{\pgfqpoint{0.947655in}{1.331584in}}%
\pgfpathlineto{\pgfqpoint{0.963164in}{1.317494in}}%
\pgfpathlineto{\pgfqpoint{0.978689in}{1.303706in}}%
\pgfpathlineto{\pgfqpoint{0.986535in}{1.315347in}}%
\pgfpathlineto{\pgfqpoint{0.994363in}{1.327159in}}%
\pgfpathlineto{\pgfqpoint{0.978864in}{1.340612in}}%
\pgfpathlineto{\pgfqpoint{0.963382in}{1.354369in}}%
\pgfpathlineto{\pgfqpoint{0.955528in}{1.342890in}}%
\pgfpathlineto{\pgfqpoint{0.947655in}{1.331584in}}%
\pgfpathclose%
\pgfusepath{fill}%
\end{pgfscope}%
\begin{pgfscope}%
\pgfpathrectangle{\pgfqpoint{0.407592in}{0.222904in}}{\pgfqpoint{1.411163in}{1.411163in}}%
\pgfusepath{clip}%
\pgfsetbuttcap%
\pgfsetroundjoin%
\definecolor{currentfill}{rgb}{0.909804,0.698039,0.423529}%
\pgfsetfillcolor{currentfill}%
\pgfsetfillopacity{0.120000}%
\pgfsetlinewidth{0.000000pt}%
\definecolor{currentstroke}{rgb}{0.000000,0.000000,0.000000}%
\pgfsetstrokecolor{currentstroke}%
\pgfsetdash{}{0pt}%
\pgfpathmoveto{\pgfqpoint{1.025412in}{1.301244in}}%
\pgfpathlineto{\pgfqpoint{1.040965in}{1.288822in}}%
\pgfpathlineto{\pgfqpoint{1.056537in}{1.276784in}}%
\pgfpathlineto{\pgfqpoint{1.064300in}{1.289422in}}%
\pgfpathlineto{\pgfqpoint{1.048739in}{1.301295in}}%
\pgfpathlineto{\pgfqpoint{1.033198in}{1.313552in}}%
\pgfusepath{fill}%
\end{pgfscope}%
\begin{pgfscope}%
\pgfpathrectangle{\pgfqpoint{0.407592in}{0.222904in}}{\pgfqpoint{1.411163in}{1.411163in}}%
\pgfusepath{clip}%
\pgfsetbuttcap%
\pgfsetroundjoin%
\definecolor{currentfill}{rgb}{0.909804,0.698039,0.423529}%
\pgfsetfillcolor{currentfill}%
\pgfsetfillopacity{0.120000}%
\pgfsetlinewidth{0.000000pt}%
\definecolor{currentstroke}{rgb}{0.000000,0.000000,0.000000}%
\pgfsetstrokecolor{currentstroke}%
\pgfsetdash{}{0pt}%
\pgfpathmoveto{\pgfqpoint{0.900818in}{1.339187in}}%
\pgfpathlineto{\pgfqpoint{0.916327in}{1.324223in}}%
\pgfpathlineto{\pgfqpoint{0.931850in}{1.309521in}}%
\pgfpathlineto{\pgfqpoint{0.939762in}{1.320458in}}%
\pgfpathlineto{\pgfqpoint{0.947655in}{1.331584in}}%
\pgfpathlineto{\pgfqpoint{0.932161in}{1.345958in}}%
\pgfpathlineto{\pgfqpoint{0.916681in}{1.360597in}}%
\pgfpathlineto{\pgfqpoint{0.908760in}{1.349798in}}%
\pgfpathlineto{\pgfqpoint{0.900818in}{1.339187in}}%
\pgfpathclose%
\pgfusepath{fill}%
\end{pgfscope}%
\begin{pgfscope}%
\pgfpathrectangle{\pgfqpoint{0.407592in}{0.222904in}}{\pgfqpoint{1.411163in}{1.411163in}}%
\pgfusepath{clip}%
\pgfsetbuttcap%
\pgfsetroundjoin%
\definecolor{currentfill}{rgb}{0.909804,0.698039,0.423529}%
\pgfsetfillcolor{currentfill}%
\pgfsetfillopacity{0.120000}%
\pgfsetlinewidth{0.000000pt}%
\definecolor{currentstroke}{rgb}{0.000000,0.000000,0.000000}%
\pgfsetstrokecolor{currentstroke}%
\pgfsetdash{}{0pt}%
\pgfpathmoveto{\pgfqpoint{0.822839in}{1.381909in}}%
\pgfpathlineto{\pgfqpoint{0.838329in}{1.365768in}}%
\pgfpathlineto{\pgfqpoint{0.853831in}{1.349820in}}%
\pgfpathlineto{\pgfqpoint{0.861847in}{1.359724in}}%
\pgfpathlineto{\pgfqpoint{0.869840in}{1.369829in}}%
\pgfpathlineto{\pgfqpoint{0.854371in}{1.385477in}}%
\pgfpathlineto{\pgfqpoint{0.838913in}{1.401324in}}%
\pgfpathlineto{\pgfqpoint{0.830887in}{1.391518in}}%
\pgfpathlineto{\pgfqpoint{0.822839in}{1.381909in}}%
\pgfpathclose%
\pgfusepath{fill}%
\end{pgfscope}%
\begin{pgfscope}%
\pgfpathrectangle{\pgfqpoint{0.407592in}{0.222904in}}{\pgfqpoint{1.411163in}{1.411163in}}%
\pgfusepath{clip}%
\pgfsetbuttcap%
\pgfsetroundjoin%
\definecolor{currentfill}{rgb}{0.909804,0.698039,0.423529}%
\pgfsetfillcolor{currentfill}%
\pgfsetfillopacity{0.120000}%
\pgfsetlinewidth{0.000000pt}%
\definecolor{currentstroke}{rgb}{0.000000,0.000000,0.000000}%
\pgfsetstrokecolor{currentstroke}%
\pgfsetdash{}{0pt}%
\pgfpathmoveto{\pgfqpoint{0.978689in}{1.303706in}}%
\pgfpathlineto{\pgfqpoint{0.994230in}{1.290241in}}%
\pgfpathlineto{\pgfqpoint{1.009789in}{1.277119in}}%
\pgfpathlineto{\pgfqpoint{1.017609in}{1.289096in}}%
\pgfpathlineto{\pgfqpoint{1.025412in}{1.301244in}}%
\pgfpathlineto{\pgfqpoint{1.009879in}{1.314030in}}%
\pgfpathlineto{\pgfqpoint{0.994363in}{1.327159in}}%
\pgfpathlineto{\pgfqpoint{0.986535in}{1.315347in}}%
\pgfpathlineto{\pgfqpoint{0.978689in}{1.303706in}}%
\pgfpathclose%
\pgfusepath{fill}%
\end{pgfscope}%
\begin{pgfscope}%
\pgfpathrectangle{\pgfqpoint{0.407592in}{0.222904in}}{\pgfqpoint{1.411163in}{1.411163in}}%
\pgfusepath{clip}%
\pgfsetbuttcap%
\pgfsetroundjoin%
\definecolor{currentfill}{rgb}{0.909804,0.698039,0.423529}%
\pgfsetfillcolor{currentfill}%
\pgfsetfillopacity{0.120000}%
\pgfsetlinewidth{0.000000pt}%
\definecolor{currentstroke}{rgb}{0.000000,0.000000,0.000000}%
\pgfsetstrokecolor{currentstroke}%
\pgfsetdash{}{0pt}%
\pgfpathmoveto{\pgfqpoint{0.931850in}{1.309521in}}%
\pgfpathlineto{\pgfqpoint{0.947387in}{1.295099in}}%
\pgfpathlineto{\pgfqpoint{0.962940in}{1.280976in}}%
\pgfpathlineto{\pgfqpoint{0.970824in}{1.292246in}}%
\pgfpathlineto{\pgfqpoint{0.978689in}{1.303706in}}%
\pgfpathlineto{\pgfqpoint{0.963164in}{1.317494in}}%
\pgfpathlineto{\pgfqpoint{0.947655in}{1.331584in}}%
\pgfpathlineto{\pgfqpoint{0.939762in}{1.320458in}}%
\pgfpathlineto{\pgfqpoint{0.931850in}{1.309521in}}%
\pgfpathclose%
\pgfusepath{fill}%
\end{pgfscope}%
\begin{pgfscope}%
\pgfpathrectangle{\pgfqpoint{0.407592in}{0.222904in}}{\pgfqpoint{1.411163in}{1.411163in}}%
\pgfusepath{clip}%
\pgfsetbuttcap%
\pgfsetroundjoin%
\definecolor{currentfill}{rgb}{0.909804,0.698039,0.423529}%
\pgfsetfillcolor{currentfill}%
\pgfsetfillopacity{0.120000}%
\pgfsetlinewidth{0.000000pt}%
\definecolor{currentstroke}{rgb}{0.000000,0.000000,0.000000}%
\pgfsetstrokecolor{currentstroke}%
\pgfsetdash{}{0pt}%
\pgfpathmoveto{\pgfqpoint{0.853831in}{1.349820in}}%
\pgfpathlineto{\pgfqpoint{0.869345in}{1.334077in}}%
\pgfpathlineto{\pgfqpoint{0.884871in}{1.318557in}}%
\pgfpathlineto{\pgfqpoint{0.892856in}{1.328770in}}%
\pgfpathlineto{\pgfqpoint{0.900818in}{1.339187in}}%
\pgfpathlineto{\pgfqpoint{0.885323in}{1.354394in}}%
\pgfpathlineto{\pgfqpoint{0.869840in}{1.369829in}}%
\pgfpathlineto{\pgfqpoint{0.861847in}{1.359724in}}%
\pgfpathlineto{\pgfqpoint{0.853831in}{1.349820in}}%
\pgfpathclose%
\pgfusepath{fill}%
\end{pgfscope}%
\begin{pgfscope}%
\pgfpathrectangle{\pgfqpoint{0.407592in}{0.222904in}}{\pgfqpoint{1.411163in}{1.411163in}}%
\pgfusepath{clip}%
\pgfsetbuttcap%
\pgfsetroundjoin%
\definecolor{currentfill}{rgb}{0.909804,0.698039,0.423529}%
\pgfsetfillcolor{currentfill}%
\pgfsetfillopacity{0.120000}%
\pgfsetlinewidth{0.000000pt}%
\definecolor{currentstroke}{rgb}{0.000000,0.000000,0.000000}%
\pgfsetstrokecolor{currentstroke}%
\pgfsetdash{}{0pt}%
\pgfpathmoveto{\pgfqpoint{1.056537in}{1.276784in}}%
\pgfpathlineto{\pgfqpoint{1.072130in}{1.265154in}}%
\pgfpathlineto{\pgfqpoint{1.087744in}{1.253952in}}%
\pgfpathlineto{\pgfqpoint{1.095486in}{1.266909in}}%
\pgfpathlineto{\pgfqpoint{1.079882in}{1.277953in}}%
\pgfpathlineto{\pgfqpoint{1.064300in}{1.289422in}}%
\pgfusepath{fill}%
\end{pgfscope}%
\begin{pgfscope}%
\pgfpathrectangle{\pgfqpoint{0.407592in}{0.222904in}}{\pgfqpoint{1.411163in}{1.411163in}}%
\pgfusepath{clip}%
\pgfsetbuttcap%
\pgfsetroundjoin%
\definecolor{currentfill}{rgb}{0.909804,0.698039,0.423529}%
\pgfsetfillcolor{currentfill}%
\pgfsetfillopacity{0.120000}%
\pgfsetlinewidth{0.000000pt}%
\definecolor{currentstroke}{rgb}{0.000000,0.000000,0.000000}%
\pgfsetstrokecolor{currentstroke}%
\pgfsetdash{}{0pt}%
\pgfpathmoveto{\pgfqpoint{1.009789in}{1.277119in}}%
\pgfpathlineto{\pgfqpoint{1.025365in}{1.264361in}}%
\pgfpathlineto{\pgfqpoint{1.040961in}{1.251991in}}%
\pgfpathlineto{\pgfqpoint{1.048757in}{1.264304in}}%
\pgfpathlineto{\pgfqpoint{1.056537in}{1.276784in}}%
\pgfpathlineto{\pgfqpoint{1.040965in}{1.288822in}}%
\pgfpathlineto{\pgfqpoint{1.025412in}{1.301244in}}%
\pgfpathlineto{\pgfqpoint{1.017609in}{1.289096in}}%
\pgfpathlineto{\pgfqpoint{1.009789in}{1.277119in}}%
\pgfpathclose%
\pgfusepath{fill}%
\end{pgfscope}%
\begin{pgfscope}%
\pgfpathrectangle{\pgfqpoint{0.407592in}{0.222904in}}{\pgfqpoint{1.411163in}{1.411163in}}%
\pgfusepath{clip}%
\pgfsetbuttcap%
\pgfsetroundjoin%
\definecolor{currentfill}{rgb}{0.909804,0.698039,0.423529}%
\pgfsetfillcolor{currentfill}%
\pgfsetfillopacity{0.120000}%
\pgfsetlinewidth{0.000000pt}%
\definecolor{currentstroke}{rgb}{0.000000,0.000000,0.000000}%
\pgfsetstrokecolor{currentstroke}%
\pgfsetdash{}{0pt}%
\pgfpathmoveto{\pgfqpoint{0.884871in}{1.318557in}}%
\pgfpathlineto{\pgfqpoint{0.900411in}{1.303275in}}%
\pgfpathlineto{\pgfqpoint{0.915964in}{1.288249in}}%
\pgfpathlineto{\pgfqpoint{0.923917in}{1.298782in}}%
\pgfpathlineto{\pgfqpoint{0.931850in}{1.309521in}}%
\pgfpathlineto{\pgfqpoint{0.916327in}{1.324223in}}%
\pgfpathlineto{\pgfqpoint{0.900818in}{1.339187in}}%
\pgfpathlineto{\pgfqpoint{0.892856in}{1.328770in}}%
\pgfpathlineto{\pgfqpoint{0.884871in}{1.318557in}}%
\pgfpathclose%
\pgfusepath{fill}%
\end{pgfscope}%
\begin{pgfscope}%
\pgfpathrectangle{\pgfqpoint{0.407592in}{0.222904in}}{\pgfqpoint{1.411163in}{1.411163in}}%
\pgfusepath{clip}%
\pgfsetbuttcap%
\pgfsetroundjoin%
\definecolor{currentfill}{rgb}{0.909804,0.698039,0.423529}%
\pgfsetfillcolor{currentfill}%
\pgfsetfillopacity{0.120000}%
\pgfsetlinewidth{0.000000pt}%
\definecolor{currentstroke}{rgb}{0.000000,0.000000,0.000000}%
\pgfsetstrokecolor{currentstroke}%
\pgfsetdash{}{0pt}%
\pgfpathmoveto{\pgfqpoint{0.962940in}{1.280976in}}%
\pgfpathlineto{\pgfqpoint{0.978508in}{1.267173in}}%
\pgfpathlineto{\pgfqpoint{0.994092in}{1.253712in}}%
\pgfpathlineto{\pgfqpoint{1.001950in}{1.265321in}}%
\pgfpathlineto{\pgfqpoint{1.009789in}{1.277119in}}%
\pgfpathlineto{\pgfqpoint{0.994230in}{1.290241in}}%
\pgfpathlineto{\pgfqpoint{0.978689in}{1.303706in}}%
\pgfpathlineto{\pgfqpoint{0.970824in}{1.292246in}}%
\pgfpathlineto{\pgfqpoint{0.962940in}{1.280976in}}%
\pgfpathclose%
\pgfusepath{fill}%
\end{pgfscope}%
\begin{pgfscope}%
\pgfpathrectangle{\pgfqpoint{0.407592in}{0.222904in}}{\pgfqpoint{1.411163in}{1.411163in}}%
\pgfusepath{clip}%
\pgfsetbuttcap%
\pgfsetroundjoin%
\definecolor{currentfill}{rgb}{0.909804,0.698039,0.423529}%
\pgfsetfillcolor{currentfill}%
\pgfsetfillopacity{0.120000}%
\pgfsetlinewidth{0.000000pt}%
\definecolor{currentstroke}{rgb}{0.000000,0.000000,0.000000}%
\pgfsetstrokecolor{currentstroke}%
\pgfsetdash{}{0pt}%
\pgfpathmoveto{\pgfqpoint{0.806672in}{1.363303in}}%
\pgfpathlineto{\pgfqpoint{0.822196in}{1.346878in}}%
\pgfpathlineto{\pgfqpoint{0.837731in}{1.330638in}}%
\pgfpathlineto{\pgfqpoint{0.845793in}{1.340122in}}%
\pgfpathlineto{\pgfqpoint{0.853831in}{1.349820in}}%
\pgfpathlineto{\pgfqpoint{0.838329in}{1.365768in}}%
\pgfpathlineto{\pgfqpoint{0.822839in}{1.381909in}}%
\pgfpathlineto{\pgfqpoint{0.814767in}{1.372501in}}%
\pgfpathlineto{\pgfqpoint{0.806672in}{1.363303in}}%
\pgfpathclose%
\pgfusepath{fill}%
\end{pgfscope}%
\begin{pgfscope}%
\pgfpathrectangle{\pgfqpoint{0.407592in}{0.222904in}}{\pgfqpoint{1.411163in}{1.411163in}}%
\pgfusepath{clip}%
\pgfsetbuttcap%
\pgfsetroundjoin%
\definecolor{currentfill}{rgb}{0.909804,0.698039,0.423529}%
\pgfsetfillcolor{currentfill}%
\pgfsetfillopacity{0.120000}%
\pgfsetlinewidth{0.000000pt}%
\definecolor{currentstroke}{rgb}{0.000000,0.000000,0.000000}%
\pgfsetstrokecolor{currentstroke}%
\pgfsetdash{}{0pt}%
\pgfpathmoveto{\pgfqpoint{1.087744in}{1.253952in}}%
\pgfpathlineto{\pgfqpoint{1.103382in}{1.243201in}}%
\pgfpathlineto{\pgfqpoint{1.119043in}{1.232925in}}%
\pgfpathlineto{\pgfqpoint{1.126765in}{1.246184in}}%
\pgfpathlineto{\pgfqpoint{1.111113in}{1.256312in}}%
\pgfpathlineto{\pgfqpoint{1.095486in}{1.266909in}}%
\pgfusepath{fill}%
\end{pgfscope}%
\begin{pgfscope}%
\pgfpathrectangle{\pgfqpoint{0.407592in}{0.222904in}}{\pgfqpoint{1.411163in}{1.411163in}}%
\pgfusepath{clip}%
\pgfsetbuttcap%
\pgfsetroundjoin%
\definecolor{currentfill}{rgb}{0.909804,0.698039,0.423529}%
\pgfsetfillcolor{currentfill}%
\pgfsetfillopacity{0.120000}%
\pgfsetlinewidth{0.000000pt}%
\definecolor{currentstroke}{rgb}{0.000000,0.000000,0.000000}%
\pgfsetstrokecolor{currentstroke}%
\pgfsetdash{}{0pt}%
\pgfpathmoveto{\pgfqpoint{1.040961in}{1.251991in}}%
\pgfpathlineto{\pgfqpoint{1.056576in}{1.240032in}}%
\pgfpathlineto{\pgfqpoint{1.072212in}{1.228508in}}%
\pgfpathlineto{\pgfqpoint{1.079986in}{1.241148in}}%
\pgfpathlineto{\pgfqpoint{1.087744in}{1.253952in}}%
\pgfpathlineto{\pgfqpoint{1.072130in}{1.265154in}}%
\pgfpathlineto{\pgfqpoint{1.056537in}{1.276784in}}%
\pgfpathlineto{\pgfqpoint{1.048757in}{1.264304in}}%
\pgfpathlineto{\pgfqpoint{1.040961in}{1.251991in}}%
\pgfpathclose%
\pgfusepath{fill}%
\end{pgfscope}%
\begin{pgfscope}%
\pgfpathrectangle{\pgfqpoint{0.407592in}{0.222904in}}{\pgfqpoint{1.411163in}{1.411163in}}%
\pgfusepath{clip}%
\pgfsetbuttcap%
\pgfsetroundjoin%
\definecolor{currentfill}{rgb}{0.909804,0.698039,0.423529}%
\pgfsetfillcolor{currentfill}%
\pgfsetfillopacity{0.120000}%
\pgfsetlinewidth{0.000000pt}%
\definecolor{currentstroke}{rgb}{0.000000,0.000000,0.000000}%
\pgfsetstrokecolor{currentstroke}%
\pgfsetdash{}{0pt}%
\pgfpathmoveto{\pgfqpoint{0.915964in}{1.288249in}}%
\pgfpathlineto{\pgfqpoint{0.931531in}{1.273497in}}%
\pgfpathlineto{\pgfqpoint{0.947112in}{1.259041in}}%
\pgfpathlineto{\pgfqpoint{0.955036in}{1.269904in}}%
\pgfpathlineto{\pgfqpoint{0.962940in}{1.280976in}}%
\pgfpathlineto{\pgfqpoint{0.947387in}{1.295099in}}%
\pgfpathlineto{\pgfqpoint{0.931850in}{1.309521in}}%
\pgfpathlineto{\pgfqpoint{0.923917in}{1.298782in}}%
\pgfpathlineto{\pgfqpoint{0.915964in}{1.288249in}}%
\pgfpathclose%
\pgfusepath{fill}%
\end{pgfscope}%
\begin{pgfscope}%
\pgfpathrectangle{\pgfqpoint{0.407592in}{0.222904in}}{\pgfqpoint{1.411163in}{1.411163in}}%
\pgfusepath{clip}%
\pgfsetbuttcap%
\pgfsetroundjoin%
\definecolor{currentfill}{rgb}{0.909804,0.698039,0.423529}%
\pgfsetfillcolor{currentfill}%
\pgfsetfillopacity{0.120000}%
\pgfsetlinewidth{0.000000pt}%
\definecolor{currentstroke}{rgb}{0.000000,0.000000,0.000000}%
\pgfsetstrokecolor{currentstroke}%
\pgfsetdash{}{0pt}%
\pgfpathmoveto{\pgfqpoint{0.837731in}{1.330638in}}%
\pgfpathlineto{\pgfqpoint{0.853278in}{1.314598in}}%
\pgfpathlineto{\pgfqpoint{0.868837in}{1.298773in}}%
\pgfpathlineto{\pgfqpoint{0.876866in}{1.308555in}}%
\pgfpathlineto{\pgfqpoint{0.884871in}{1.318557in}}%
\pgfpathlineto{\pgfqpoint{0.869345in}{1.334077in}}%
\pgfpathlineto{\pgfqpoint{0.853831in}{1.349820in}}%
\pgfpathlineto{\pgfqpoint{0.845793in}{1.340122in}}%
\pgfpathlineto{\pgfqpoint{0.837731in}{1.330638in}}%
\pgfpathclose%
\pgfusepath{fill}%
\end{pgfscope}%
\begin{pgfscope}%
\pgfpathrectangle{\pgfqpoint{0.407592in}{0.222904in}}{\pgfqpoint{1.411163in}{1.411163in}}%
\pgfusepath{clip}%
\pgfsetbuttcap%
\pgfsetroundjoin%
\definecolor{currentfill}{rgb}{0.909804,0.698039,0.423529}%
\pgfsetfillcolor{currentfill}%
\pgfsetfillopacity{0.120000}%
\pgfsetlinewidth{0.000000pt}%
\definecolor{currentstroke}{rgb}{0.000000,0.000000,0.000000}%
\pgfsetstrokecolor{currentstroke}%
\pgfsetdash{}{0pt}%
\pgfpathmoveto{\pgfqpoint{0.994092in}{1.253712in}}%
\pgfpathlineto{\pgfqpoint{1.009694in}{1.240615in}}%
\pgfpathlineto{\pgfqpoint{1.025314in}{1.227908in}}%
\pgfpathlineto{\pgfqpoint{1.033146in}{1.239856in}}%
\pgfpathlineto{\pgfqpoint{1.040961in}{1.251991in}}%
\pgfpathlineto{\pgfqpoint{1.025365in}{1.264361in}}%
\pgfpathlineto{\pgfqpoint{1.009789in}{1.277119in}}%
\pgfpathlineto{\pgfqpoint{1.001950in}{1.265321in}}%
\pgfpathlineto{\pgfqpoint{0.994092in}{1.253712in}}%
\pgfpathclose%
\pgfusepath{fill}%
\end{pgfscope}%
\begin{pgfscope}%
\pgfpathrectangle{\pgfqpoint{0.407592in}{0.222904in}}{\pgfqpoint{1.411163in}{1.411163in}}%
\pgfusepath{clip}%
\pgfsetbuttcap%
\pgfsetroundjoin%
\definecolor{currentfill}{rgb}{0.909804,0.698039,0.423529}%
\pgfsetfillcolor{currentfill}%
\pgfsetfillopacity{0.120000}%
\pgfsetlinewidth{0.000000pt}%
\definecolor{currentstroke}{rgb}{0.000000,0.000000,0.000000}%
\pgfsetstrokecolor{currentstroke}%
\pgfsetdash{}{0pt}%
\pgfpathmoveto{\pgfqpoint{1.072212in}{1.228508in}}%
\pgfpathlineto{\pgfqpoint{1.087870in}{1.217443in}}%
\pgfpathlineto{\pgfqpoint{1.103552in}{1.206863in}}%
\pgfpathlineto{\pgfqpoint{1.111306in}{1.219814in}}%
\pgfpathlineto{\pgfqpoint{1.119043in}{1.232925in}}%
\pgfpathlineto{\pgfqpoint{1.103382in}{1.243201in}}%
\pgfpathlineto{\pgfqpoint{1.087744in}{1.253952in}}%
\pgfpathlineto{\pgfqpoint{1.079986in}{1.241148in}}%
\pgfpathlineto{\pgfqpoint{1.072212in}{1.228508in}}%
\pgfpathclose%
\pgfusepath{fill}%
\end{pgfscope}%
\begin{pgfscope}%
\pgfpathrectangle{\pgfqpoint{0.407592in}{0.222904in}}{\pgfqpoint{1.411163in}{1.411163in}}%
\pgfusepath{clip}%
\pgfsetbuttcap%
\pgfsetroundjoin%
\definecolor{currentfill}{rgb}{0.909804,0.698039,0.423529}%
\pgfsetfillcolor{currentfill}%
\pgfsetfillopacity{0.120000}%
\pgfsetlinewidth{0.000000pt}%
\definecolor{currentstroke}{rgb}{0.000000,0.000000,0.000000}%
\pgfsetstrokecolor{currentstroke}%
\pgfsetdash{}{0pt}%
\pgfpathmoveto{\pgfqpoint{0.947112in}{1.259041in}}%
\pgfpathlineto{\pgfqpoint{0.962708in}{1.244902in}}%
\pgfpathlineto{\pgfqpoint{0.978320in}{1.231102in}}%
\pgfpathlineto{\pgfqpoint{0.986216in}{1.242302in}}%
\pgfpathlineto{\pgfqpoint{0.994092in}{1.253712in}}%
\pgfpathlineto{\pgfqpoint{0.978508in}{1.267173in}}%
\pgfpathlineto{\pgfqpoint{0.962940in}{1.280976in}}%
\pgfpathlineto{\pgfqpoint{0.955036in}{1.269904in}}%
\pgfpathlineto{\pgfqpoint{0.947112in}{1.259041in}}%
\pgfpathclose%
\pgfusepath{fill}%
\end{pgfscope}%
\begin{pgfscope}%
\pgfpathrectangle{\pgfqpoint{0.407592in}{0.222904in}}{\pgfqpoint{1.411163in}{1.411163in}}%
\pgfusepath{clip}%
\pgfsetbuttcap%
\pgfsetroundjoin%
\definecolor{currentfill}{rgb}{0.909804,0.698039,0.423529}%
\pgfsetfillcolor{currentfill}%
\pgfsetfillopacity{0.120000}%
\pgfsetlinewidth{0.000000pt}%
\definecolor{currentstroke}{rgb}{0.000000,0.000000,0.000000}%
\pgfsetstrokecolor{currentstroke}%
\pgfsetdash{}{0pt}%
\pgfpathmoveto{\pgfqpoint{0.868837in}{1.298773in}}%
\pgfpathlineto{\pgfqpoint{0.884409in}{1.283180in}}%
\pgfpathlineto{\pgfqpoint{0.899993in}{1.267837in}}%
\pgfpathlineto{\pgfqpoint{0.907989in}{1.277931in}}%
\pgfpathlineto{\pgfqpoint{0.915964in}{1.288249in}}%
\pgfpathlineto{\pgfqpoint{0.900411in}{1.303275in}}%
\pgfpathlineto{\pgfqpoint{0.884871in}{1.318557in}}%
\pgfpathlineto{\pgfqpoint{0.876866in}{1.308555in}}%
\pgfpathlineto{\pgfqpoint{0.868837in}{1.298773in}}%
\pgfpathclose%
\pgfusepath{fill}%
\end{pgfscope}%
\begin{pgfscope}%
\pgfpathrectangle{\pgfqpoint{0.407592in}{0.222904in}}{\pgfqpoint{1.411163in}{1.411163in}}%
\pgfusepath{clip}%
\pgfsetbuttcap%
\pgfsetroundjoin%
\definecolor{currentfill}{rgb}{0.909804,0.698039,0.423529}%
\pgfsetfillcolor{currentfill}%
\pgfsetfillopacity{0.120000}%
\pgfsetlinewidth{0.000000pt}%
\definecolor{currentstroke}{rgb}{0.000000,0.000000,0.000000}%
\pgfsetstrokecolor{currentstroke}%
\pgfsetdash{}{0pt}%
\pgfpathmoveto{\pgfqpoint{1.119043in}{1.232925in}}%
\pgfpathlineto{\pgfqpoint{1.134730in}{1.223145in}}%
\pgfpathlineto{\pgfqpoint{1.150444in}{1.213884in}}%
\pgfpathlineto{\pgfqpoint{1.158148in}{1.227420in}}%
\pgfpathlineto{\pgfqpoint{1.142443in}{1.236547in}}%
\pgfpathlineto{\pgfqpoint{1.126765in}{1.246184in}}%
\pgfusepath{fill}%
\end{pgfscope}%
\begin{pgfscope}%
\pgfpathrectangle{\pgfqpoint{0.407592in}{0.222904in}}{\pgfqpoint{1.411163in}{1.411163in}}%
\pgfusepath{clip}%
\pgfsetbuttcap%
\pgfsetroundjoin%
\definecolor{currentfill}{rgb}{0.909804,0.698039,0.423529}%
\pgfsetfillcolor{currentfill}%
\pgfsetfillopacity{0.120000}%
\pgfsetlinewidth{0.000000pt}%
\definecolor{currentstroke}{rgb}{0.000000,0.000000,0.000000}%
\pgfsetstrokecolor{currentstroke}%
\pgfsetdash{}{0pt}%
\pgfpathmoveto{\pgfqpoint{0.790410in}{1.345557in}}%
\pgfpathlineto{\pgfqpoint{0.805968in}{1.328861in}}%
\pgfpathlineto{\pgfqpoint{0.821538in}{1.312343in}}%
\pgfpathlineto{\pgfqpoint{0.829647in}{1.321376in}}%
\pgfpathlineto{\pgfqpoint{0.837731in}{1.330638in}}%
\pgfpathlineto{\pgfqpoint{0.822196in}{1.346878in}}%
\pgfpathlineto{\pgfqpoint{0.806672in}{1.363303in}}%
\pgfpathlineto{\pgfqpoint{0.798553in}{1.354319in}}%
\pgfpathlineto{\pgfqpoint{0.790410in}{1.345557in}}%
\pgfpathclose%
\pgfusepath{fill}%
\end{pgfscope}%
\begin{pgfscope}%
\pgfpathrectangle{\pgfqpoint{0.407592in}{0.222904in}}{\pgfqpoint{1.411163in}{1.411163in}}%
\pgfusepath{clip}%
\pgfsetbuttcap%
\pgfsetroundjoin%
\definecolor{currentfill}{rgb}{0.909804,0.698039,0.423529}%
\pgfsetfillcolor{currentfill}%
\pgfsetfillopacity{0.120000}%
\pgfsetlinewidth{0.000000pt}%
\definecolor{currentstroke}{rgb}{0.000000,0.000000,0.000000}%
\pgfsetstrokecolor{currentstroke}%
\pgfsetdash{}{0pt}%
\pgfpathmoveto{\pgfqpoint{1.025314in}{1.227908in}}%
\pgfpathlineto{\pgfqpoint{1.040953in}{1.215615in}}%
\pgfpathlineto{\pgfqpoint{1.056612in}{1.203762in}}%
\pgfpathlineto{\pgfqpoint{1.064421in}{1.216042in}}%
\pgfpathlineto{\pgfqpoint{1.072212in}{1.228508in}}%
\pgfpathlineto{\pgfqpoint{1.056576in}{1.240032in}}%
\pgfpathlineto{\pgfqpoint{1.040961in}{1.251991in}}%
\pgfpathlineto{\pgfqpoint{1.033146in}{1.239856in}}%
\pgfpathlineto{\pgfqpoint{1.025314in}{1.227908in}}%
\pgfpathclose%
\pgfusepath{fill}%
\end{pgfscope}%
\begin{pgfscope}%
\pgfpathrectangle{\pgfqpoint{0.407592in}{0.222904in}}{\pgfqpoint{1.411163in}{1.411163in}}%
\pgfusepath{clip}%
\pgfsetbuttcap%
\pgfsetroundjoin%
\definecolor{currentfill}{rgb}{0.909804,0.698039,0.423529}%
\pgfsetfillcolor{currentfill}%
\pgfsetfillopacity{0.120000}%
\pgfsetlinewidth{0.000000pt}%
\definecolor{currentstroke}{rgb}{0.000000,0.000000,0.000000}%
\pgfsetstrokecolor{currentstroke}%
\pgfsetdash{}{0pt}%
\pgfpathmoveto{\pgfqpoint{0.899993in}{1.267837in}}%
\pgfpathlineto{\pgfqpoint{0.915590in}{1.252764in}}%
\pgfpathlineto{\pgfqpoint{0.931202in}{1.237981in}}%
\pgfpathlineto{\pgfqpoint{0.939167in}{1.248397in}}%
\pgfpathlineto{\pgfqpoint{0.947112in}{1.259041in}}%
\pgfpathlineto{\pgfqpoint{0.931531in}{1.273497in}}%
\pgfpathlineto{\pgfqpoint{0.915964in}{1.288249in}}%
\pgfpathlineto{\pgfqpoint{0.907989in}{1.277931in}}%
\pgfpathlineto{\pgfqpoint{0.899993in}{1.267837in}}%
\pgfpathclose%
\pgfusepath{fill}%
\end{pgfscope}%
\begin{pgfscope}%
\pgfpathrectangle{\pgfqpoint{0.407592in}{0.222904in}}{\pgfqpoint{1.411163in}{1.411163in}}%
\pgfusepath{clip}%
\pgfsetbuttcap%
\pgfsetroundjoin%
\definecolor{currentfill}{rgb}{0.909804,0.698039,0.423529}%
\pgfsetfillcolor{currentfill}%
\pgfsetfillopacity{0.120000}%
\pgfsetlinewidth{0.000000pt}%
\definecolor{currentstroke}{rgb}{0.000000,0.000000,0.000000}%
\pgfsetstrokecolor{currentstroke}%
\pgfsetdash{}{0pt}%
\pgfpathmoveto{\pgfqpoint{0.978320in}{1.231102in}}%
\pgfpathlineto{\pgfqpoint{0.993948in}{1.217667in}}%
\pgfpathlineto{\pgfqpoint{1.009594in}{1.204621in}}%
\pgfpathlineto{\pgfqpoint{1.017463in}{1.216159in}}%
\pgfpathlineto{\pgfqpoint{1.025314in}{1.227908in}}%
\pgfpathlineto{\pgfqpoint{1.009694in}{1.240615in}}%
\pgfpathlineto{\pgfqpoint{0.994092in}{1.253712in}}%
\pgfpathlineto{\pgfqpoint{0.986216in}{1.242302in}}%
\pgfpathlineto{\pgfqpoint{0.978320in}{1.231102in}}%
\pgfpathclose%
\pgfusepath{fill}%
\end{pgfscope}%
\begin{pgfscope}%
\pgfpathrectangle{\pgfqpoint{0.407592in}{0.222904in}}{\pgfqpoint{1.411163in}{1.411163in}}%
\pgfusepath{clip}%
\pgfsetbuttcap%
\pgfsetroundjoin%
\definecolor{currentfill}{rgb}{0.909804,0.698039,0.423529}%
\pgfsetfillcolor{currentfill}%
\pgfsetfillopacity{0.120000}%
\pgfsetlinewidth{0.000000pt}%
\definecolor{currentstroke}{rgb}{0.000000,0.000000,0.000000}%
\pgfsetstrokecolor{currentstroke}%
\pgfsetdash{}{0pt}%
\pgfpathmoveto{\pgfqpoint{0.821538in}{1.312343in}}%
\pgfpathlineto{\pgfqpoint{0.837119in}{1.296017in}}%
\pgfpathlineto{\pgfqpoint{0.852711in}{1.279900in}}%
\pgfpathlineto{\pgfqpoint{0.860786in}{1.289218in}}%
\pgfpathlineto{\pgfqpoint{0.868837in}{1.298773in}}%
\pgfpathlineto{\pgfqpoint{0.853278in}{1.314598in}}%
\pgfpathlineto{\pgfqpoint{0.837731in}{1.330638in}}%
\pgfpathlineto{\pgfqpoint{0.829647in}{1.321376in}}%
\pgfpathlineto{\pgfqpoint{0.821538in}{1.312343in}}%
\pgfpathclose%
\pgfusepath{fill}%
\end{pgfscope}%
\begin{pgfscope}%
\pgfpathrectangle{\pgfqpoint{0.407592in}{0.222904in}}{\pgfqpoint{1.411163in}{1.411163in}}%
\pgfusepath{clip}%
\pgfsetbuttcap%
\pgfsetroundjoin%
\definecolor{currentfill}{rgb}{0.909804,0.698039,0.423529}%
\pgfsetfillcolor{currentfill}%
\pgfsetfillopacity{0.120000}%
\pgfsetlinewidth{0.000000pt}%
\definecolor{currentstroke}{rgb}{0.000000,0.000000,0.000000}%
\pgfsetstrokecolor{currentstroke}%
\pgfsetdash{}{0pt}%
\pgfpathmoveto{\pgfqpoint{1.103552in}{1.206863in}}%
\pgfpathlineto{\pgfqpoint{1.119258in}{1.196791in}}%
\pgfpathlineto{\pgfqpoint{1.134991in}{1.187253in}}%
\pgfpathlineto{\pgfqpoint{1.142725in}{1.200491in}}%
\pgfpathlineto{\pgfqpoint{1.150444in}{1.213884in}}%
\pgfpathlineto{\pgfqpoint{1.134730in}{1.223145in}}%
\pgfpathlineto{\pgfqpoint{1.119043in}{1.232925in}}%
\pgfpathlineto{\pgfqpoint{1.111306in}{1.219814in}}%
\pgfpathlineto{\pgfqpoint{1.103552in}{1.206863in}}%
\pgfpathclose%
\pgfusepath{fill}%
\end{pgfscope}%
\begin{pgfscope}%
\pgfpathrectangle{\pgfqpoint{0.407592in}{0.222904in}}{\pgfqpoint{1.411163in}{1.411163in}}%
\pgfusepath{clip}%
\pgfsetbuttcap%
\pgfsetroundjoin%
\definecolor{currentfill}{rgb}{0.909804,0.698039,0.423529}%
\pgfsetfillcolor{currentfill}%
\pgfsetfillopacity{0.120000}%
\pgfsetlinewidth{0.000000pt}%
\definecolor{currentstroke}{rgb}{0.000000,0.000000,0.000000}%
\pgfsetstrokecolor{currentstroke}%
\pgfsetdash{}{0pt}%
\pgfpathmoveto{\pgfqpoint{1.150444in}{1.213884in}}%
\pgfpathlineto{\pgfqpoint{1.166187in}{1.205163in}}%
\pgfpathlineto{\pgfqpoint{1.181959in}{1.197003in}}%
\pgfpathlineto{\pgfqpoint{1.189647in}{1.210783in}}%
\pgfpathlineto{\pgfqpoint{1.173882in}{1.218826in}}%
\pgfpathlineto{\pgfqpoint{1.158148in}{1.227420in}}%
\pgfusepath{fill}%
\end{pgfscope}%
\begin{pgfscope}%
\pgfpathrectangle{\pgfqpoint{0.407592in}{0.222904in}}{\pgfqpoint{1.411163in}{1.411163in}}%
\pgfusepath{clip}%
\pgfsetbuttcap%
\pgfsetroundjoin%
\definecolor{currentfill}{rgb}{0.909804,0.698039,0.423529}%
\pgfsetfillcolor{currentfill}%
\pgfsetfillopacity{0.120000}%
\pgfsetlinewidth{0.000000pt}%
\definecolor{currentstroke}{rgb}{0.000000,0.000000,0.000000}%
\pgfsetstrokecolor{currentstroke}%
\pgfsetdash{}{0pt}%
\pgfpathmoveto{\pgfqpoint{1.056612in}{1.203762in}}%
\pgfpathlineto{\pgfqpoint{1.072293in}{1.192376in}}%
\pgfpathlineto{\pgfqpoint{1.087995in}{1.181483in}}%
\pgfpathlineto{\pgfqpoint{1.095782in}{1.194082in}}%
\pgfpathlineto{\pgfqpoint{1.103552in}{1.206863in}}%
\pgfpathlineto{\pgfqpoint{1.087870in}{1.217443in}}%
\pgfpathlineto{\pgfqpoint{1.072212in}{1.228508in}}%
\pgfpathlineto{\pgfqpoint{1.064421in}{1.216042in}}%
\pgfpathlineto{\pgfqpoint{1.056612in}{1.203762in}}%
\pgfpathclose%
\pgfusepath{fill}%
\end{pgfscope}%
\begin{pgfscope}%
\pgfpathrectangle{\pgfqpoint{0.407592in}{0.222904in}}{\pgfqpoint{1.411163in}{1.411163in}}%
\pgfusepath{clip}%
\pgfsetbuttcap%
\pgfsetroundjoin%
\definecolor{currentfill}{rgb}{0.909804,0.698039,0.423529}%
\pgfsetfillcolor{currentfill}%
\pgfsetfillopacity{0.120000}%
\pgfsetlinewidth{0.000000pt}%
\definecolor{currentstroke}{rgb}{0.000000,0.000000,0.000000}%
\pgfsetstrokecolor{currentstroke}%
\pgfsetdash{}{0pt}%
\pgfpathmoveto{\pgfqpoint{0.931202in}{1.237981in}}%
\pgfpathlineto{\pgfqpoint{0.946827in}{1.223510in}}%
\pgfpathlineto{\pgfqpoint{0.962468in}{1.209377in}}%
\pgfpathlineto{\pgfqpoint{0.970404in}{1.220124in}}%
\pgfpathlineto{\pgfqpoint{0.978320in}{1.231102in}}%
\pgfpathlineto{\pgfqpoint{0.962708in}{1.244902in}}%
\pgfpathlineto{\pgfqpoint{0.947112in}{1.259041in}}%
\pgfpathlineto{\pgfqpoint{0.939167in}{1.248397in}}%
\pgfpathlineto{\pgfqpoint{0.931202in}{1.237981in}}%
\pgfpathclose%
\pgfusepath{fill}%
\end{pgfscope}%
\begin{pgfscope}%
\pgfpathrectangle{\pgfqpoint{0.407592in}{0.222904in}}{\pgfqpoint{1.411163in}{1.411163in}}%
\pgfusepath{clip}%
\pgfsetbuttcap%
\pgfsetroundjoin%
\definecolor{currentfill}{rgb}{0.909804,0.698039,0.423529}%
\pgfsetfillcolor{currentfill}%
\pgfsetfillopacity{0.120000}%
\pgfsetlinewidth{0.000000pt}%
\definecolor{currentstroke}{rgb}{0.000000,0.000000,0.000000}%
\pgfsetstrokecolor{currentstroke}%
\pgfsetdash{}{0pt}%
\pgfpathmoveto{\pgfqpoint{0.852711in}{1.279900in}}%
\pgfpathlineto{\pgfqpoint{0.868316in}{1.264008in}}%
\pgfpathlineto{\pgfqpoint{0.883933in}{1.248360in}}%
\pgfpathlineto{\pgfqpoint{0.891974in}{1.257977in}}%
\pgfpathlineto{\pgfqpoint{0.899993in}{1.267837in}}%
\pgfpathlineto{\pgfqpoint{0.884409in}{1.283180in}}%
\pgfpathlineto{\pgfqpoint{0.868837in}{1.298773in}}%
\pgfpathlineto{\pgfqpoint{0.860786in}{1.289218in}}%
\pgfpathlineto{\pgfqpoint{0.852711in}{1.279900in}}%
\pgfpathclose%
\pgfusepath{fill}%
\end{pgfscope}%
\begin{pgfscope}%
\pgfpathrectangle{\pgfqpoint{0.407592in}{0.222904in}}{\pgfqpoint{1.411163in}{1.411163in}}%
\pgfusepath{clip}%
\pgfsetbuttcap%
\pgfsetroundjoin%
\definecolor{currentfill}{rgb}{0.909804,0.698039,0.423529}%
\pgfsetfillcolor{currentfill}%
\pgfsetfillopacity{0.120000}%
\pgfsetlinewidth{0.000000pt}%
\definecolor{currentstroke}{rgb}{0.000000,0.000000,0.000000}%
\pgfsetstrokecolor{currentstroke}%
\pgfsetdash{}{0pt}%
\pgfpathmoveto{\pgfqpoint{1.009594in}{1.204621in}}%
\pgfpathlineto{\pgfqpoint{1.025258in}{1.191991in}}%
\pgfpathlineto{\pgfqpoint{1.040942in}{1.179806in}}%
\pgfpathlineto{\pgfqpoint{1.048786in}{1.191679in}}%
\pgfpathlineto{\pgfqpoint{1.056612in}{1.203762in}}%
\pgfpathlineto{\pgfqpoint{1.040953in}{1.215615in}}%
\pgfpathlineto{\pgfqpoint{1.025314in}{1.227908in}}%
\pgfpathlineto{\pgfqpoint{1.017463in}{1.216159in}}%
\pgfpathlineto{\pgfqpoint{1.009594in}{1.204621in}}%
\pgfpathclose%
\pgfusepath{fill}%
\end{pgfscope}%
\begin{pgfscope}%
\pgfpathrectangle{\pgfqpoint{0.407592in}{0.222904in}}{\pgfqpoint{1.411163in}{1.411163in}}%
\pgfusepath{clip}%
\pgfsetbuttcap%
\pgfsetroundjoin%
\definecolor{currentfill}{rgb}{0.909804,0.698039,0.423529}%
\pgfsetfillcolor{currentfill}%
\pgfsetfillopacity{0.120000}%
\pgfsetlinewidth{0.000000pt}%
\definecolor{currentstroke}{rgb}{0.000000,0.000000,0.000000}%
\pgfsetstrokecolor{currentstroke}%
\pgfsetdash{}{0pt}%
\pgfpathmoveto{\pgfqpoint{0.774047in}{1.328725in}}%
\pgfpathlineto{\pgfqpoint{0.789641in}{1.311773in}}%
\pgfpathlineto{\pgfqpoint{0.805246in}{1.294992in}}%
\pgfpathlineto{\pgfqpoint{0.813405in}{1.303546in}}%
\pgfpathlineto{\pgfqpoint{0.821538in}{1.312343in}}%
\pgfpathlineto{\pgfqpoint{0.805968in}{1.328861in}}%
\pgfpathlineto{\pgfqpoint{0.790410in}{1.345557in}}%
\pgfpathlineto{\pgfqpoint{0.782241in}{1.337024in}}%
\pgfpathlineto{\pgfqpoint{0.774047in}{1.328725in}}%
\pgfpathclose%
\pgfusepath{fill}%
\end{pgfscope}%
\begin{pgfscope}%
\pgfpathrectangle{\pgfqpoint{0.407592in}{0.222904in}}{\pgfqpoint{1.411163in}{1.411163in}}%
\pgfusepath{clip}%
\pgfsetbuttcap%
\pgfsetroundjoin%
\definecolor{currentfill}{rgb}{0.909804,0.698039,0.423529}%
\pgfsetfillcolor{currentfill}%
\pgfsetfillopacity{0.120000}%
\pgfsetlinewidth{0.000000pt}%
\definecolor{currentstroke}{rgb}{0.000000,0.000000,0.000000}%
\pgfsetstrokecolor{currentstroke}%
\pgfsetdash{}{0pt}%
\pgfpathmoveto{\pgfqpoint{1.134991in}{1.187253in}}%
\pgfpathlineto{\pgfqpoint{1.150751in}{1.178272in}}%
\pgfpathlineto{\pgfqpoint{1.166540in}{1.169872in}}%
\pgfpathlineto{\pgfqpoint{1.174257in}{1.183362in}}%
\pgfpathlineto{\pgfqpoint{1.181959in}{1.197003in}}%
\pgfpathlineto{\pgfqpoint{1.166187in}{1.205163in}}%
\pgfpathlineto{\pgfqpoint{1.150444in}{1.213884in}}%
\pgfpathlineto{\pgfqpoint{1.142725in}{1.200491in}}%
\pgfpathlineto{\pgfqpoint{1.134991in}{1.187253in}}%
\pgfpathclose%
\pgfusepath{fill}%
\end{pgfscope}%
\begin{pgfscope}%
\pgfpathrectangle{\pgfqpoint{0.407592in}{0.222904in}}{\pgfqpoint{1.411163in}{1.411163in}}%
\pgfusepath{clip}%
\pgfsetbuttcap%
\pgfsetroundjoin%
\definecolor{currentfill}{rgb}{0.909804,0.698039,0.423529}%
\pgfsetfillcolor{currentfill}%
\pgfsetfillopacity{0.120000}%
\pgfsetlinewidth{0.000000pt}%
\definecolor{currentstroke}{rgb}{0.000000,0.000000,0.000000}%
\pgfsetstrokecolor{currentstroke}%
\pgfsetdash{}{0pt}%
\pgfpathmoveto{\pgfqpoint{0.883933in}{1.248360in}}%
\pgfpathlineto{\pgfqpoint{0.899562in}{1.232975in}}%
\pgfpathlineto{\pgfqpoint{0.915205in}{1.217875in}}%
\pgfpathlineto{\pgfqpoint{0.923214in}{1.227803in}}%
\pgfpathlineto{\pgfqpoint{0.931202in}{1.237981in}}%
\pgfpathlineto{\pgfqpoint{0.915590in}{1.252764in}}%
\pgfpathlineto{\pgfqpoint{0.899993in}{1.267837in}}%
\pgfpathlineto{\pgfqpoint{0.891974in}{1.257977in}}%
\pgfpathlineto{\pgfqpoint{0.883933in}{1.248360in}}%
\pgfpathclose%
\pgfusepath{fill}%
\end{pgfscope}%
\begin{pgfscope}%
\pgfpathrectangle{\pgfqpoint{0.407592in}{0.222904in}}{\pgfqpoint{1.411163in}{1.411163in}}%
\pgfusepath{clip}%
\pgfsetbuttcap%
\pgfsetroundjoin%
\definecolor{currentfill}{rgb}{0.909804,0.698039,0.423529}%
\pgfsetfillcolor{currentfill}%
\pgfsetfillopacity{0.120000}%
\pgfsetlinewidth{0.000000pt}%
\definecolor{currentstroke}{rgb}{0.000000,0.000000,0.000000}%
\pgfsetstrokecolor{currentstroke}%
\pgfsetdash{}{0pt}%
\pgfpathmoveto{\pgfqpoint{0.962468in}{1.209377in}}%
\pgfpathlineto{\pgfqpoint{0.978125in}{1.195605in}}%
\pgfpathlineto{\pgfqpoint{0.993798in}{1.182222in}}%
\pgfpathlineto{\pgfqpoint{1.001706in}{1.193304in}}%
\pgfpathlineto{\pgfqpoint{1.009594in}{1.204621in}}%
\pgfpathlineto{\pgfqpoint{0.993948in}{1.217667in}}%
\pgfpathlineto{\pgfqpoint{0.978320in}{1.231102in}}%
\pgfpathlineto{\pgfqpoint{0.970404in}{1.220124in}}%
\pgfpathlineto{\pgfqpoint{0.962468in}{1.209377in}}%
\pgfpathclose%
\pgfusepath{fill}%
\end{pgfscope}%
\begin{pgfscope}%
\pgfpathrectangle{\pgfqpoint{0.407592in}{0.222904in}}{\pgfqpoint{1.411163in}{1.411163in}}%
\pgfusepath{clip}%
\pgfsetbuttcap%
\pgfsetroundjoin%
\definecolor{currentfill}{rgb}{0.909804,0.698039,0.423529}%
\pgfsetfillcolor{currentfill}%
\pgfsetfillopacity{0.120000}%
\pgfsetlinewidth{0.000000pt}%
\definecolor{currentstroke}{rgb}{0.000000,0.000000,0.000000}%
\pgfsetstrokecolor{currentstroke}%
\pgfsetdash{}{0pt}%
\pgfpathmoveto{\pgfqpoint{1.087995in}{1.181483in}}%
\pgfpathlineto{\pgfqpoint{1.103722in}{1.171111in}}%
\pgfpathlineto{\pgfqpoint{1.119474in}{1.161288in}}%
\pgfpathlineto{\pgfqpoint{1.127240in}{1.174181in}}%
\pgfpathlineto{\pgfqpoint{1.134991in}{1.187253in}}%
\pgfpathlineto{\pgfqpoint{1.119258in}{1.196791in}}%
\pgfpathlineto{\pgfqpoint{1.103552in}{1.206863in}}%
\pgfpathlineto{\pgfqpoint{1.095782in}{1.194082in}}%
\pgfpathlineto{\pgfqpoint{1.087995in}{1.181483in}}%
\pgfpathclose%
\pgfusepath{fill}%
\end{pgfscope}%
\begin{pgfscope}%
\pgfpathrectangle{\pgfqpoint{0.407592in}{0.222904in}}{\pgfqpoint{1.411163in}{1.411163in}}%
\pgfusepath{clip}%
\pgfsetbuttcap%
\pgfsetroundjoin%
\definecolor{currentfill}{rgb}{0.909804,0.698039,0.423529}%
\pgfsetfillcolor{currentfill}%
\pgfsetfillopacity{0.120000}%
\pgfsetlinewidth{0.000000pt}%
\definecolor{currentstroke}{rgb}{0.000000,0.000000,0.000000}%
\pgfsetstrokecolor{currentstroke}%
\pgfsetdash{}{0pt}%
\pgfpathmoveto{\pgfqpoint{1.181959in}{1.197003in}}%
\pgfpathlineto{\pgfqpoint{1.197763in}{1.189423in}}%
\pgfpathlineto{\pgfqpoint{1.213601in}{1.182442in}}%
\pgfpathlineto{\pgfqpoint{1.221273in}{1.196422in}}%
\pgfpathlineto{\pgfqpoint{1.205443in}{1.203309in}}%
\pgfpathlineto{\pgfqpoint{1.189647in}{1.210783in}}%
\pgfusepath{fill}%
\end{pgfscope}%
\begin{pgfscope}%
\pgfpathrectangle{\pgfqpoint{0.407592in}{0.222904in}}{\pgfqpoint{1.411163in}{1.411163in}}%
\pgfusepath{clip}%
\pgfsetbuttcap%
\pgfsetroundjoin%
\definecolor{currentfill}{rgb}{0.909804,0.698039,0.423529}%
\pgfsetfillcolor{currentfill}%
\pgfsetfillopacity{0.120000}%
\pgfsetlinewidth{0.000000pt}%
\definecolor{currentstroke}{rgb}{0.000000,0.000000,0.000000}%
\pgfsetstrokecolor{currentstroke}%
\pgfsetdash{}{0pt}%
\pgfpathmoveto{\pgfqpoint{0.805246in}{1.294992in}}%
\pgfpathlineto{\pgfqpoint{0.820862in}{1.278397in}}%
\pgfpathlineto{\pgfqpoint{0.836490in}{1.262004in}}%
\pgfpathlineto{\pgfqpoint{0.844613in}{1.270826in}}%
\pgfpathlineto{\pgfqpoint{0.852711in}{1.279900in}}%
\pgfpathlineto{\pgfqpoint{0.837119in}{1.296017in}}%
\pgfpathlineto{\pgfqpoint{0.821538in}{1.312343in}}%
\pgfpathlineto{\pgfqpoint{0.813405in}{1.303546in}}%
\pgfpathlineto{\pgfqpoint{0.805246in}{1.294992in}}%
\pgfpathclose%
\pgfusepath{fill}%
\end{pgfscope}%
\begin{pgfscope}%
\pgfpathrectangle{\pgfqpoint{0.407592in}{0.222904in}}{\pgfqpoint{1.411163in}{1.411163in}}%
\pgfusepath{clip}%
\pgfsetbuttcap%
\pgfsetroundjoin%
\definecolor{currentfill}{rgb}{0.909804,0.698039,0.423529}%
\pgfsetfillcolor{currentfill}%
\pgfsetfillopacity{0.120000}%
\pgfsetlinewidth{0.000000pt}%
\definecolor{currentstroke}{rgb}{0.000000,0.000000,0.000000}%
\pgfsetstrokecolor{currentstroke}%
\pgfsetdash{}{0pt}%
\pgfpathmoveto{\pgfqpoint{1.040942in}{1.179806in}}%
\pgfpathlineto{\pgfqpoint{1.056645in}{1.168093in}}%
\pgfpathlineto{\pgfqpoint{1.072371in}{1.156883in}}%
\pgfpathlineto{\pgfqpoint{1.080192in}{1.169079in}}%
\pgfpathlineto{\pgfqpoint{1.087995in}{1.181483in}}%
\pgfpathlineto{\pgfqpoint{1.072293in}{1.192376in}}%
\pgfpathlineto{\pgfqpoint{1.056612in}{1.203762in}}%
\pgfpathlineto{\pgfqpoint{1.048786in}{1.191679in}}%
\pgfpathlineto{\pgfqpoint{1.040942in}{1.179806in}}%
\pgfpathclose%
\pgfusepath{fill}%
\end{pgfscope}%
\begin{pgfscope}%
\pgfpathrectangle{\pgfqpoint{0.407592in}{0.222904in}}{\pgfqpoint{1.411163in}{1.411163in}}%
\pgfusepath{clip}%
\pgfsetbuttcap%
\pgfsetroundjoin%
\definecolor{currentfill}{rgb}{0.909804,0.698039,0.423529}%
\pgfsetfillcolor{currentfill}%
\pgfsetfillopacity{0.120000}%
\pgfsetlinewidth{0.000000pt}%
\definecolor{currentstroke}{rgb}{0.000000,0.000000,0.000000}%
\pgfsetstrokecolor{currentstroke}%
\pgfsetdash{}{0pt}%
\pgfpathmoveto{\pgfqpoint{0.915205in}{1.217875in}}%
\pgfpathlineto{\pgfqpoint{0.930862in}{1.203083in}}%
\pgfpathlineto{\pgfqpoint{0.946532in}{1.188624in}}%
\pgfpathlineto{\pgfqpoint{0.954511in}{1.198873in}}%
\pgfpathlineto{\pgfqpoint{0.962468in}{1.209377in}}%
\pgfpathlineto{\pgfqpoint{0.946827in}{1.223510in}}%
\pgfpathlineto{\pgfqpoint{0.931202in}{1.237981in}}%
\pgfpathlineto{\pgfqpoint{0.923214in}{1.227803in}}%
\pgfpathlineto{\pgfqpoint{0.915205in}{1.217875in}}%
\pgfpathclose%
\pgfusepath{fill}%
\end{pgfscope}%
\begin{pgfscope}%
\pgfpathrectangle{\pgfqpoint{0.407592in}{0.222904in}}{\pgfqpoint{1.411163in}{1.411163in}}%
\pgfusepath{clip}%
\pgfsetbuttcap%
\pgfsetroundjoin%
\definecolor{currentfill}{rgb}{0.909804,0.698039,0.423529}%
\pgfsetfillcolor{currentfill}%
\pgfsetfillopacity{0.120000}%
\pgfsetlinewidth{0.000000pt}%
\definecolor{currentstroke}{rgb}{0.000000,0.000000,0.000000}%
\pgfsetstrokecolor{currentstroke}%
\pgfsetdash{}{0pt}%
\pgfpathmoveto{\pgfqpoint{0.836490in}{1.262004in}}%
\pgfpathlineto{\pgfqpoint{0.852128in}{1.245829in}}%
\pgfpathlineto{\pgfqpoint{0.867779in}{1.229891in}}%
\pgfpathlineto{\pgfqpoint{0.875868in}{1.238995in}}%
\pgfpathlineto{\pgfqpoint{0.883933in}{1.248360in}}%
\pgfpathlineto{\pgfqpoint{0.868316in}{1.264008in}}%
\pgfpathlineto{\pgfqpoint{0.852711in}{1.279900in}}%
\pgfpathlineto{\pgfqpoint{0.844613in}{1.270826in}}%
\pgfpathlineto{\pgfqpoint{0.836490in}{1.262004in}}%
\pgfpathclose%
\pgfusepath{fill}%
\end{pgfscope}%
\begin{pgfscope}%
\pgfpathrectangle{\pgfqpoint{0.407592in}{0.222904in}}{\pgfqpoint{1.411163in}{1.411163in}}%
\pgfusepath{clip}%
\pgfsetbuttcap%
\pgfsetroundjoin%
\definecolor{currentfill}{rgb}{0.909804,0.698039,0.423529}%
\pgfsetfillcolor{currentfill}%
\pgfsetfillopacity{0.120000}%
\pgfsetlinewidth{0.000000pt}%
\definecolor{currentstroke}{rgb}{0.000000,0.000000,0.000000}%
\pgfsetstrokecolor{currentstroke}%
\pgfsetdash{}{0pt}%
\pgfpathmoveto{\pgfqpoint{0.993798in}{1.182222in}}%
\pgfpathlineto{\pgfqpoint{1.009488in}{1.169257in}}%
\pgfpathlineto{\pgfqpoint{1.025197in}{1.156739in}}%
\pgfpathlineto{\pgfqpoint{1.033079in}{1.168155in}}%
\pgfpathlineto{\pgfqpoint{1.040942in}{1.179806in}}%
\pgfpathlineto{\pgfqpoint{1.025258in}{1.191991in}}%
\pgfpathlineto{\pgfqpoint{1.009594in}{1.204621in}}%
\pgfpathlineto{\pgfqpoint{1.001706in}{1.193304in}}%
\pgfpathlineto{\pgfqpoint{0.993798in}{1.182222in}}%
\pgfpathclose%
\pgfusepath{fill}%
\end{pgfscope}%
\begin{pgfscope}%
\pgfpathrectangle{\pgfqpoint{0.407592in}{0.222904in}}{\pgfqpoint{1.411163in}{1.411163in}}%
\pgfusepath{clip}%
\pgfsetbuttcap%
\pgfsetroundjoin%
\definecolor{currentfill}{rgb}{0.909804,0.698039,0.423529}%
\pgfsetfillcolor{currentfill}%
\pgfsetfillopacity{0.120000}%
\pgfsetlinewidth{0.000000pt}%
\definecolor{currentstroke}{rgb}{0.000000,0.000000,0.000000}%
\pgfsetstrokecolor{currentstroke}%
\pgfsetdash{}{0pt}%
\pgfpathmoveto{\pgfqpoint{1.166540in}{1.169872in}}%
\pgfpathlineto{\pgfqpoint{1.182360in}{1.162074in}}%
\pgfpathlineto{\pgfqpoint{1.198213in}{1.154898in}}%
\pgfpathlineto{\pgfqpoint{1.205914in}{1.168597in}}%
\pgfpathlineto{\pgfqpoint{1.213601in}{1.182442in}}%
\pgfpathlineto{\pgfqpoint{1.197763in}{1.189423in}}%
\pgfpathlineto{\pgfqpoint{1.181959in}{1.197003in}}%
\pgfpathlineto{\pgfqpoint{1.174257in}{1.183362in}}%
\pgfpathlineto{\pgfqpoint{1.166540in}{1.169872in}}%
\pgfpathclose%
\pgfusepath{fill}%
\end{pgfscope}%
\begin{pgfscope}%
\pgfpathrectangle{\pgfqpoint{0.407592in}{0.222904in}}{\pgfqpoint{1.411163in}{1.411163in}}%
\pgfusepath{clip}%
\pgfsetbuttcap%
\pgfsetroundjoin%
\definecolor{currentfill}{rgb}{0.909804,0.698039,0.423529}%
\pgfsetfillcolor{currentfill}%
\pgfsetfillopacity{0.120000}%
\pgfsetlinewidth{0.000000pt}%
\definecolor{currentstroke}{rgb}{0.000000,0.000000,0.000000}%
\pgfsetstrokecolor{currentstroke}%
\pgfsetdash{}{0pt}%
\pgfpathmoveto{\pgfqpoint{1.119474in}{1.161288in}}%
\pgfpathlineto{\pgfqpoint{1.135253in}{1.152039in}}%
\pgfpathlineto{\pgfqpoint{1.151060in}{1.143390in}}%
\pgfpathlineto{\pgfqpoint{1.158808in}{1.156544in}}%
\pgfpathlineto{\pgfqpoint{1.166540in}{1.169872in}}%
\pgfpathlineto{\pgfqpoint{1.150751in}{1.178272in}}%
\pgfpathlineto{\pgfqpoint{1.134991in}{1.187253in}}%
\pgfpathlineto{\pgfqpoint{1.127240in}{1.174181in}}%
\pgfpathlineto{\pgfqpoint{1.119474in}{1.161288in}}%
\pgfpathclose%
\pgfusepath{fill}%
\end{pgfscope}%
\begin{pgfscope}%
\pgfpathrectangle{\pgfqpoint{0.407592in}{0.222904in}}{\pgfqpoint{1.411163in}{1.411163in}}%
\pgfusepath{clip}%
\pgfsetbuttcap%
\pgfsetroundjoin%
\definecolor{currentfill}{rgb}{0.909804,0.698039,0.423529}%
\pgfsetfillcolor{currentfill}%
\pgfsetfillopacity{0.120000}%
\pgfsetlinewidth{0.000000pt}%
\definecolor{currentstroke}{rgb}{0.000000,0.000000,0.000000}%
\pgfsetstrokecolor{currentstroke}%
\pgfsetdash{}{0pt}%
\pgfpathmoveto{\pgfqpoint{0.757579in}{1.312855in}}%
\pgfpathlineto{\pgfqpoint{0.773210in}{1.295666in}}%
\pgfpathlineto{\pgfqpoint{0.788852in}{1.278642in}}%
\pgfpathlineto{\pgfqpoint{0.797062in}{1.286689in}}%
\pgfpathlineto{\pgfqpoint{0.805246in}{1.294992in}}%
\pgfpathlineto{\pgfqpoint{0.789641in}{1.311773in}}%
\pgfpathlineto{\pgfqpoint{0.774047in}{1.328725in}}%
\pgfpathlineto{\pgfqpoint{0.765826in}{1.320667in}}%
\pgfpathlineto{\pgfqpoint{0.757579in}{1.312855in}}%
\pgfpathclose%
\pgfusepath{fill}%
\end{pgfscope}%
\begin{pgfscope}%
\pgfpathrectangle{\pgfqpoint{0.407592in}{0.222904in}}{\pgfqpoint{1.411163in}{1.411163in}}%
\pgfusepath{clip}%
\pgfsetbuttcap%
\pgfsetroundjoin%
\definecolor{currentfill}{rgb}{0.909804,0.698039,0.423529}%
\pgfsetfillcolor{currentfill}%
\pgfsetfillopacity{0.120000}%
\pgfsetlinewidth{0.000000pt}%
\definecolor{currentstroke}{rgb}{0.000000,0.000000,0.000000}%
\pgfsetstrokecolor{currentstroke}%
\pgfsetdash{}{0pt}%
\pgfpathmoveto{\pgfqpoint{1.072371in}{1.156883in}}%
\pgfpathlineto{\pgfqpoint{1.088119in}{1.146205in}}%
\pgfpathlineto{\pgfqpoint{1.103892in}{1.136089in}}%
\pgfpathlineto{\pgfqpoint{1.111691in}{1.148586in}}%
\pgfpathlineto{\pgfqpoint{1.119474in}{1.161288in}}%
\pgfpathlineto{\pgfqpoint{1.103722in}{1.171111in}}%
\pgfpathlineto{\pgfqpoint{1.087995in}{1.181483in}}%
\pgfpathlineto{\pgfqpoint{1.080192in}{1.169079in}}%
\pgfpathlineto{\pgfqpoint{1.072371in}{1.156883in}}%
\pgfpathclose%
\pgfusepath{fill}%
\end{pgfscope}%
\begin{pgfscope}%
\pgfpathrectangle{\pgfqpoint{0.407592in}{0.222904in}}{\pgfqpoint{1.411163in}{1.411163in}}%
\pgfusepath{clip}%
\pgfsetbuttcap%
\pgfsetroundjoin%
\definecolor{currentfill}{rgb}{0.909804,0.698039,0.423529}%
\pgfsetfillcolor{currentfill}%
\pgfsetfillopacity{0.120000}%
\pgfsetlinewidth{0.000000pt}%
\definecolor{currentstroke}{rgb}{0.000000,0.000000,0.000000}%
\pgfsetstrokecolor{currentstroke}%
\pgfsetdash{}{0pt}%
\pgfpathmoveto{\pgfqpoint{1.213601in}{1.182442in}}%
\pgfpathlineto{\pgfqpoint{1.229473in}{1.176074in}}%
\pgfpathlineto{\pgfqpoint{1.245382in}{1.170334in}}%
\pgfpathlineto{\pgfqpoint{1.253040in}{1.184463in}}%
\pgfpathlineto{\pgfqpoint{1.237138in}{1.190135in}}%
\pgfpathlineto{\pgfqpoint{1.221273in}{1.196422in}}%
\pgfusepath{fill}%
\end{pgfscope}%
\begin{pgfscope}%
\pgfpathrectangle{\pgfqpoint{0.407592in}{0.222904in}}{\pgfqpoint{1.411163in}{1.411163in}}%
\pgfusepath{clip}%
\pgfsetbuttcap%
\pgfsetroundjoin%
\definecolor{currentfill}{rgb}{0.909804,0.698039,0.423529}%
\pgfsetfillcolor{currentfill}%
\pgfsetfillopacity{0.120000}%
\pgfsetlinewidth{0.000000pt}%
\definecolor{currentstroke}{rgb}{0.000000,0.000000,0.000000}%
\pgfsetstrokecolor{currentstroke}%
\pgfsetdash{}{0pt}%
\pgfpathmoveto{\pgfqpoint{0.946532in}{1.188624in}}%
\pgfpathlineto{\pgfqpoint{0.962218in}{1.174523in}}%
\pgfpathlineto{\pgfqpoint{0.977920in}{1.160811in}}%
\pgfpathlineto{\pgfqpoint{0.985869in}{1.171387in}}%
\pgfpathlineto{\pgfqpoint{0.993798in}{1.182222in}}%
\pgfpathlineto{\pgfqpoint{0.978125in}{1.195605in}}%
\pgfpathlineto{\pgfqpoint{0.962468in}{1.209377in}}%
\pgfpathlineto{\pgfqpoint{0.954511in}{1.198873in}}%
\pgfpathlineto{\pgfqpoint{0.946532in}{1.188624in}}%
\pgfpathclose%
\pgfusepath{fill}%
\end{pgfscope}%
\begin{pgfscope}%
\pgfpathrectangle{\pgfqpoint{0.407592in}{0.222904in}}{\pgfqpoint{1.411163in}{1.411163in}}%
\pgfusepath{clip}%
\pgfsetbuttcap%
\pgfsetroundjoin%
\definecolor{currentfill}{rgb}{0.909804,0.698039,0.423529}%
\pgfsetfillcolor{currentfill}%
\pgfsetfillopacity{0.120000}%
\pgfsetlinewidth{0.000000pt}%
\definecolor{currentstroke}{rgb}{0.000000,0.000000,0.000000}%
\pgfsetstrokecolor{currentstroke}%
\pgfsetdash{}{0pt}%
\pgfpathmoveto{\pgfqpoint{0.867779in}{1.229891in}}%
\pgfpathlineto{\pgfqpoint{0.883442in}{1.214210in}}%
\pgfpathlineto{\pgfqpoint{0.899118in}{1.198808in}}%
\pgfpathlineto{\pgfqpoint{0.907173in}{1.208207in}}%
\pgfpathlineto{\pgfqpoint{0.915205in}{1.217875in}}%
\pgfpathlineto{\pgfqpoint{0.899562in}{1.232975in}}%
\pgfpathlineto{\pgfqpoint{0.883933in}{1.248360in}}%
\pgfpathlineto{\pgfqpoint{0.875868in}{1.238995in}}%
\pgfpathlineto{\pgfqpoint{0.867779in}{1.229891in}}%
\pgfpathclose%
\pgfusepath{fill}%
\end{pgfscope}%
\begin{pgfscope}%
\pgfpathrectangle{\pgfqpoint{0.407592in}{0.222904in}}{\pgfqpoint{1.411163in}{1.411163in}}%
\pgfusepath{clip}%
\pgfsetbuttcap%
\pgfsetroundjoin%
\definecolor{currentfill}{rgb}{0.909804,0.698039,0.423529}%
\pgfsetfillcolor{currentfill}%
\pgfsetfillopacity{0.120000}%
\pgfsetlinewidth{0.000000pt}%
\definecolor{currentstroke}{rgb}{0.000000,0.000000,0.000000}%
\pgfsetstrokecolor{currentstroke}%
\pgfsetdash{}{0pt}%
\pgfpathmoveto{\pgfqpoint{1.025197in}{1.156739in}}%
\pgfpathlineto{\pgfqpoint{1.040926in}{1.144699in}}%
\pgfpathlineto{\pgfqpoint{1.056675in}{1.133169in}}%
\pgfpathlineto{\pgfqpoint{1.064532in}{1.144908in}}%
\pgfpathlineto{\pgfqpoint{1.072371in}{1.156883in}}%
\pgfpathlineto{\pgfqpoint{1.056645in}{1.168093in}}%
\pgfpathlineto{\pgfqpoint{1.040942in}{1.179806in}}%
\pgfpathlineto{\pgfqpoint{1.033079in}{1.168155in}}%
\pgfpathlineto{\pgfqpoint{1.025197in}{1.156739in}}%
\pgfpathclose%
\pgfusepath{fill}%
\end{pgfscope}%
\begin{pgfscope}%
\pgfpathrectangle{\pgfqpoint{0.407592in}{0.222904in}}{\pgfqpoint{1.411163in}{1.411163in}}%
\pgfusepath{clip}%
\pgfsetbuttcap%
\pgfsetroundjoin%
\definecolor{currentfill}{rgb}{0.909804,0.698039,0.423529}%
\pgfsetfillcolor{currentfill}%
\pgfsetfillopacity{0.120000}%
\pgfsetlinewidth{0.000000pt}%
\definecolor{currentstroke}{rgb}{0.000000,0.000000,0.000000}%
\pgfsetstrokecolor{currentstroke}%
\pgfsetdash{}{0pt}%
\pgfpathmoveto{\pgfqpoint{0.788852in}{1.278642in}}%
\pgfpathlineto{\pgfqpoint{0.804504in}{1.261798in}}%
\pgfpathlineto{\pgfqpoint{0.820167in}{1.245148in}}%
\pgfpathlineto{\pgfqpoint{0.828341in}{1.253442in}}%
\pgfpathlineto{\pgfqpoint{0.836490in}{1.262004in}}%
\pgfpathlineto{\pgfqpoint{0.820862in}{1.278397in}}%
\pgfpathlineto{\pgfqpoint{0.805246in}{1.294992in}}%
\pgfpathlineto{\pgfqpoint{0.797062in}{1.286689in}}%
\pgfpathlineto{\pgfqpoint{0.788852in}{1.278642in}}%
\pgfpathclose%
\pgfusepath{fill}%
\end{pgfscope}%
\begin{pgfscope}%
\pgfpathrectangle{\pgfqpoint{0.407592in}{0.222904in}}{\pgfqpoint{1.411163in}{1.411163in}}%
\pgfusepath{clip}%
\pgfsetbuttcap%
\pgfsetroundjoin%
\definecolor{currentfill}{rgb}{0.909804,0.698039,0.423529}%
\pgfsetfillcolor{currentfill}%
\pgfsetfillopacity{0.120000}%
\pgfsetlinewidth{0.000000pt}%
\definecolor{currentstroke}{rgb}{0.000000,0.000000,0.000000}%
\pgfsetstrokecolor{currentstroke}%
\pgfsetdash{}{0pt}%
\pgfpathmoveto{\pgfqpoint{0.899118in}{1.198808in}}%
\pgfpathlineto{\pgfqpoint{0.914806in}{1.183708in}}%
\pgfpathlineto{\pgfqpoint{0.930509in}{1.168936in}}%
\pgfpathlineto{\pgfqpoint{0.938532in}{1.178641in}}%
\pgfpathlineto{\pgfqpoint{0.946532in}{1.188624in}}%
\pgfpathlineto{\pgfqpoint{0.930862in}{1.203083in}}%
\pgfpathlineto{\pgfqpoint{0.915205in}{1.217875in}}%
\pgfpathlineto{\pgfqpoint{0.907173in}{1.208207in}}%
\pgfpathlineto{\pgfqpoint{0.899118in}{1.198808in}}%
\pgfpathclose%
\pgfusepath{fill}%
\end{pgfscope}%
\begin{pgfscope}%
\pgfpathrectangle{\pgfqpoint{0.407592in}{0.222904in}}{\pgfqpoint{1.411163in}{1.411163in}}%
\pgfusepath{clip}%
\pgfsetbuttcap%
\pgfsetroundjoin%
\definecolor{currentfill}{rgb}{0.909804,0.698039,0.423529}%
\pgfsetfillcolor{currentfill}%
\pgfsetfillopacity{0.120000}%
\pgfsetlinewidth{0.000000pt}%
\definecolor{currentstroke}{rgb}{0.000000,0.000000,0.000000}%
\pgfsetstrokecolor{currentstroke}%
\pgfsetdash{}{0pt}%
\pgfpathmoveto{\pgfqpoint{0.977920in}{1.160811in}}%
\pgfpathlineto{\pgfqpoint{0.993639in}{1.147515in}}%
\pgfpathlineto{\pgfqpoint{1.009375in}{1.134669in}}%
\pgfpathlineto{\pgfqpoint{1.017296in}{1.145572in}}%
\pgfpathlineto{\pgfqpoint{1.025197in}{1.156739in}}%
\pgfpathlineto{\pgfqpoint{1.009488in}{1.169257in}}%
\pgfpathlineto{\pgfqpoint{0.993798in}{1.182222in}}%
\pgfpathlineto{\pgfqpoint{0.985869in}{1.171387in}}%
\pgfpathlineto{\pgfqpoint{0.977920in}{1.160811in}}%
\pgfpathclose%
\pgfusepath{fill}%
\end{pgfscope}%
\begin{pgfscope}%
\pgfpathrectangle{\pgfqpoint{0.407592in}{0.222904in}}{\pgfqpoint{1.411163in}{1.411163in}}%
\pgfusepath{clip}%
\pgfsetbuttcap%
\pgfsetroundjoin%
\definecolor{currentfill}{rgb}{0.909804,0.698039,0.423529}%
\pgfsetfillcolor{currentfill}%
\pgfsetfillopacity{0.120000}%
\pgfsetlinewidth{0.000000pt}%
\definecolor{currentstroke}{rgb}{0.000000,0.000000,0.000000}%
\pgfsetstrokecolor{currentstroke}%
\pgfsetdash{}{0pt}%
\pgfpathmoveto{\pgfqpoint{1.151060in}{1.143390in}}%
\pgfpathlineto{\pgfqpoint{1.166897in}{1.135365in}}%
\pgfpathlineto{\pgfqpoint{1.182767in}{1.127988in}}%
\pgfpathlineto{\pgfqpoint{1.190497in}{1.141358in}}%
\pgfpathlineto{\pgfqpoint{1.198213in}{1.154898in}}%
\pgfpathlineto{\pgfqpoint{1.182360in}{1.162074in}}%
\pgfpathlineto{\pgfqpoint{1.166540in}{1.169872in}}%
\pgfpathlineto{\pgfqpoint{1.158808in}{1.156544in}}%
\pgfpathlineto{\pgfqpoint{1.151060in}{1.143390in}}%
\pgfpathclose%
\pgfusepath{fill}%
\end{pgfscope}%
\begin{pgfscope}%
\pgfpathrectangle{\pgfqpoint{0.407592in}{0.222904in}}{\pgfqpoint{1.411163in}{1.411163in}}%
\pgfusepath{clip}%
\pgfsetbuttcap%
\pgfsetroundjoin%
\definecolor{currentfill}{rgb}{0.909804,0.698039,0.423529}%
\pgfsetfillcolor{currentfill}%
\pgfsetfillopacity{0.120000}%
\pgfsetlinewidth{0.000000pt}%
\definecolor{currentstroke}{rgb}{0.000000,0.000000,0.000000}%
\pgfsetstrokecolor{currentstroke}%
\pgfsetdash{}{0pt}%
\pgfpathmoveto{\pgfqpoint{1.198213in}{1.154898in}}%
\pgfpathlineto{\pgfqpoint{1.214101in}{1.148362in}}%
\pgfpathlineto{\pgfqpoint{1.230024in}{1.142482in}}%
\pgfpathlineto{\pgfqpoint{1.237711in}{1.156336in}}%
\pgfpathlineto{\pgfqpoint{1.245382in}{1.170334in}}%
\pgfpathlineto{\pgfqpoint{1.229473in}{1.176074in}}%
\pgfpathlineto{\pgfqpoint{1.213601in}{1.182442in}}%
\pgfpathlineto{\pgfqpoint{1.205914in}{1.168597in}}%
\pgfpathlineto{\pgfqpoint{1.198213in}{1.154898in}}%
\pgfpathclose%
\pgfusepath{fill}%
\end{pgfscope}%
\begin{pgfscope}%
\pgfpathrectangle{\pgfqpoint{0.407592in}{0.222904in}}{\pgfqpoint{1.411163in}{1.411163in}}%
\pgfusepath{clip}%
\pgfsetbuttcap%
\pgfsetroundjoin%
\definecolor{currentfill}{rgb}{0.909804,0.698039,0.423529}%
\pgfsetfillcolor{currentfill}%
\pgfsetfillopacity{0.120000}%
\pgfsetlinewidth{0.000000pt}%
\definecolor{currentstroke}{rgb}{0.000000,0.000000,0.000000}%
\pgfsetstrokecolor{currentstroke}%
\pgfsetdash{}{0pt}%
\pgfpathmoveto{\pgfqpoint{0.820167in}{1.245148in}}%
\pgfpathlineto{\pgfqpoint{0.835842in}{1.228710in}}%
\pgfpathlineto{\pgfqpoint{0.851527in}{1.212502in}}%
\pgfpathlineto{\pgfqpoint{0.859666in}{1.221057in}}%
\pgfpathlineto{\pgfqpoint{0.867779in}{1.229891in}}%
\pgfpathlineto{\pgfqpoint{0.852128in}{1.245829in}}%
\pgfpathlineto{\pgfqpoint{0.836490in}{1.262004in}}%
\pgfpathlineto{\pgfqpoint{0.828341in}{1.253442in}}%
\pgfpathlineto{\pgfqpoint{0.820167in}{1.245148in}}%
\pgfpathclose%
\pgfusepath{fill}%
\end{pgfscope}%
\begin{pgfscope}%
\pgfpathrectangle{\pgfqpoint{0.407592in}{0.222904in}}{\pgfqpoint{1.411163in}{1.411163in}}%
\pgfusepath{clip}%
\pgfsetbuttcap%
\pgfsetroundjoin%
\definecolor{currentfill}{rgb}{0.909804,0.698039,0.423529}%
\pgfsetfillcolor{currentfill}%
\pgfsetfillopacity{0.120000}%
\pgfsetlinewidth{0.000000pt}%
\definecolor{currentstroke}{rgb}{0.000000,0.000000,0.000000}%
\pgfsetstrokecolor{currentstroke}%
\pgfsetdash{}{0pt}%
\pgfpathmoveto{\pgfqpoint{1.103892in}{1.136089in}}%
\pgfpathlineto{\pgfqpoint{1.119690in}{1.126565in}}%
\pgfpathlineto{\pgfqpoint{1.135516in}{1.117660in}}%
\pgfpathlineto{\pgfqpoint{1.143296in}{1.130424in}}%
\pgfpathlineto{\pgfqpoint{1.151060in}{1.143390in}}%
\pgfpathlineto{\pgfqpoint{1.135253in}{1.152039in}}%
\pgfpathlineto{\pgfqpoint{1.119474in}{1.161288in}}%
\pgfpathlineto{\pgfqpoint{1.111691in}{1.148586in}}%
\pgfpathlineto{\pgfqpoint{1.103892in}{1.136089in}}%
\pgfpathclose%
\pgfusepath{fill}%
\end{pgfscope}%
\begin{pgfscope}%
\pgfpathrectangle{\pgfqpoint{0.407592in}{0.222904in}}{\pgfqpoint{1.411163in}{1.411163in}}%
\pgfusepath{clip}%
\pgfsetbuttcap%
\pgfsetroundjoin%
\definecolor{currentfill}{rgb}{0.909804,0.698039,0.423529}%
\pgfsetfillcolor{currentfill}%
\pgfsetfillopacity{0.120000}%
\pgfsetlinewidth{0.000000pt}%
\definecolor{currentstroke}{rgb}{0.000000,0.000000,0.000000}%
\pgfsetstrokecolor{currentstroke}%
\pgfsetdash{}{0pt}%
\pgfpathmoveto{\pgfqpoint{1.056675in}{1.133169in}}%
\pgfpathlineto{\pgfqpoint{1.072446in}{1.122182in}}%
\pgfpathlineto{\pgfqpoint{1.088241in}{1.111769in}}%
\pgfpathlineto{\pgfqpoint{1.096075in}{1.123812in}}%
\pgfpathlineto{\pgfqpoint{1.103892in}{1.136089in}}%
\pgfpathlineto{\pgfqpoint{1.088119in}{1.146205in}}%
\pgfpathlineto{\pgfqpoint{1.072371in}{1.156883in}}%
\pgfpathlineto{\pgfqpoint{1.064532in}{1.144908in}}%
\pgfpathlineto{\pgfqpoint{1.056675in}{1.133169in}}%
\pgfpathclose%
\pgfusepath{fill}%
\end{pgfscope}%
\begin{pgfscope}%
\pgfpathrectangle{\pgfqpoint{0.407592in}{0.222904in}}{\pgfqpoint{1.411163in}{1.411163in}}%
\pgfusepath{clip}%
\pgfsetbuttcap%
\pgfsetroundjoin%
\definecolor{currentfill}{rgb}{0.909804,0.698039,0.423529}%
\pgfsetfillcolor{currentfill}%
\pgfsetfillopacity{0.120000}%
\pgfsetlinewidth{0.000000pt}%
\definecolor{currentstroke}{rgb}{0.000000,0.000000,0.000000}%
\pgfsetstrokecolor{currentstroke}%
\pgfsetdash{}{0pt}%
\pgfpathmoveto{\pgfqpoint{1.245382in}{1.170334in}}%
\pgfpathlineto{\pgfqpoint{1.261331in}{1.165232in}}%
\pgfpathlineto{\pgfqpoint{1.277319in}{1.160778in}}%
\pgfpathlineto{\pgfqpoint{1.284964in}{1.175001in}}%
\pgfpathlineto{\pgfqpoint{1.268982in}{1.179416in}}%
\pgfpathlineto{\pgfqpoint{1.253040in}{1.184463in}}%
\pgfusepath{fill}%
\end{pgfscope}%
\begin{pgfscope}%
\pgfpathrectangle{\pgfqpoint{0.407592in}{0.222904in}}{\pgfqpoint{1.411163in}{1.411163in}}%
\pgfusepath{clip}%
\pgfsetbuttcap%
\pgfsetroundjoin%
\definecolor{currentfill}{rgb}{0.909804,0.698039,0.423529}%
\pgfsetfillcolor{currentfill}%
\pgfsetfillopacity{0.120000}%
\pgfsetlinewidth{0.000000pt}%
\definecolor{currentstroke}{rgb}{0.000000,0.000000,0.000000}%
\pgfsetstrokecolor{currentstroke}%
\pgfsetdash{}{0pt}%
\pgfpathmoveto{\pgfqpoint{0.741002in}{1.297996in}}%
\pgfpathlineto{\pgfqpoint{0.756671in}{1.280592in}}%
\pgfpathlineto{\pgfqpoint{0.772350in}{1.263348in}}%
\pgfpathlineto{\pgfqpoint{0.780615in}{1.270860in}}%
\pgfpathlineto{\pgfqpoint{0.788852in}{1.278642in}}%
\pgfpathlineto{\pgfqpoint{0.773210in}{1.295666in}}%
\pgfpathlineto{\pgfqpoint{0.757579in}{1.312855in}}%
\pgfpathlineto{\pgfqpoint{0.749305in}{1.305297in}}%
\pgfpathlineto{\pgfqpoint{0.741002in}{1.297996in}}%
\pgfpathclose%
\pgfusepath{fill}%
\end{pgfscope}%
\begin{pgfscope}%
\pgfpathrectangle{\pgfqpoint{0.407592in}{0.222904in}}{\pgfqpoint{1.411163in}{1.411163in}}%
\pgfusepath{clip}%
\pgfsetbuttcap%
\pgfsetroundjoin%
\definecolor{currentfill}{rgb}{0.909804,0.698039,0.423529}%
\pgfsetfillcolor{currentfill}%
\pgfsetfillopacity{0.120000}%
\pgfsetlinewidth{0.000000pt}%
\definecolor{currentstroke}{rgb}{0.000000,0.000000,0.000000}%
\pgfsetstrokecolor{currentstroke}%
\pgfsetdash{}{0pt}%
\pgfpathmoveto{\pgfqpoint{0.930509in}{1.168936in}}%
\pgfpathlineto{\pgfqpoint{0.946226in}{1.154520in}}%
\pgfpathlineto{\pgfqpoint{0.961958in}{1.140489in}}%
\pgfpathlineto{\pgfqpoint{0.969950in}{1.150507in}}%
\pgfpathlineto{\pgfqpoint{0.977920in}{1.160811in}}%
\pgfpathlineto{\pgfqpoint{0.962218in}{1.174523in}}%
\pgfpathlineto{\pgfqpoint{0.946532in}{1.188624in}}%
\pgfpathlineto{\pgfqpoint{0.938532in}{1.178641in}}%
\pgfpathlineto{\pgfqpoint{0.930509in}{1.168936in}}%
\pgfpathclose%
\pgfusepath{fill}%
\end{pgfscope}%
\begin{pgfscope}%
\pgfpathrectangle{\pgfqpoint{0.407592in}{0.222904in}}{\pgfqpoint{1.411163in}{1.411163in}}%
\pgfusepath{clip}%
\pgfsetbuttcap%
\pgfsetroundjoin%
\definecolor{currentfill}{rgb}{0.909804,0.698039,0.423529}%
\pgfsetfillcolor{currentfill}%
\pgfsetfillopacity{0.120000}%
\pgfsetlinewidth{0.000000pt}%
\definecolor{currentstroke}{rgb}{0.000000,0.000000,0.000000}%
\pgfsetstrokecolor{currentstroke}%
\pgfsetdash{}{0pt}%
\pgfpathmoveto{\pgfqpoint{0.851527in}{1.212502in}}%
\pgfpathlineto{\pgfqpoint{0.867225in}{1.196545in}}%
\pgfpathlineto{\pgfqpoint{0.882935in}{1.180860in}}%
\pgfpathlineto{\pgfqpoint{0.891038in}{1.189689in}}%
\pgfpathlineto{\pgfqpoint{0.899118in}{1.198808in}}%
\pgfpathlineto{\pgfqpoint{0.883442in}{1.214210in}}%
\pgfpathlineto{\pgfqpoint{0.867779in}{1.229891in}}%
\pgfpathlineto{\pgfqpoint{0.859666in}{1.221057in}}%
\pgfpathlineto{\pgfqpoint{0.851527in}{1.212502in}}%
\pgfpathclose%
\pgfusepath{fill}%
\end{pgfscope}%
\begin{pgfscope}%
\pgfpathrectangle{\pgfqpoint{0.407592in}{0.222904in}}{\pgfqpoint{1.411163in}{1.411163in}}%
\pgfusepath{clip}%
\pgfsetbuttcap%
\pgfsetroundjoin%
\definecolor{currentfill}{rgb}{0.909804,0.698039,0.423529}%
\pgfsetfillcolor{currentfill}%
\pgfsetfillopacity{0.120000}%
\pgfsetlinewidth{0.000000pt}%
\definecolor{currentstroke}{rgb}{0.000000,0.000000,0.000000}%
\pgfsetstrokecolor{currentstroke}%
\pgfsetdash{}{0pt}%
\pgfpathmoveto{\pgfqpoint{1.009375in}{1.134669in}}%
\pgfpathlineto{\pgfqpoint{1.025129in}{1.122305in}}%
\pgfpathlineto{\pgfqpoint{1.040904in}{1.110457in}}%
\pgfpathlineto{\pgfqpoint{1.048799in}{1.121680in}}%
\pgfpathlineto{\pgfqpoint{1.056675in}{1.133169in}}%
\pgfpathlineto{\pgfqpoint{1.040926in}{1.144699in}}%
\pgfpathlineto{\pgfqpoint{1.025197in}{1.156739in}}%
\pgfpathlineto{\pgfqpoint{1.017296in}{1.145572in}}%
\pgfpathlineto{\pgfqpoint{1.009375in}{1.134669in}}%
\pgfpathclose%
\pgfusepath{fill}%
\end{pgfscope}%
\begin{pgfscope}%
\pgfpathrectangle{\pgfqpoint{0.407592in}{0.222904in}}{\pgfqpoint{1.411163in}{1.411163in}}%
\pgfusepath{clip}%
\pgfsetbuttcap%
\pgfsetroundjoin%
\definecolor{currentfill}{rgb}{0.909804,0.698039,0.423529}%
\pgfsetfillcolor{currentfill}%
\pgfsetfillopacity{0.120000}%
\pgfsetlinewidth{0.000000pt}%
\definecolor{currentstroke}{rgb}{0.000000,0.000000,0.000000}%
\pgfsetstrokecolor{currentstroke}%
\pgfsetdash{}{0pt}%
\pgfpathmoveto{\pgfqpoint{0.772350in}{1.263348in}}%
\pgfpathlineto{\pgfqpoint{0.788040in}{1.246277in}}%
\pgfpathlineto{\pgfqpoint{0.803740in}{1.229394in}}%
\pgfpathlineto{\pgfqpoint{0.811967in}{1.237130in}}%
\pgfpathlineto{\pgfqpoint{0.820167in}{1.245148in}}%
\pgfpathlineto{\pgfqpoint{0.804504in}{1.261798in}}%
\pgfpathlineto{\pgfqpoint{0.788852in}{1.278642in}}%
\pgfpathlineto{\pgfqpoint{0.780615in}{1.270860in}}%
\pgfpathlineto{\pgfqpoint{0.772350in}{1.263348in}}%
\pgfpathclose%
\pgfusepath{fill}%
\end{pgfscope}%
\begin{pgfscope}%
\pgfpathrectangle{\pgfqpoint{0.407592in}{0.222904in}}{\pgfqpoint{1.411163in}{1.411163in}}%
\pgfusepath{clip}%
\pgfsetbuttcap%
\pgfsetroundjoin%
\definecolor{currentfill}{rgb}{0.909804,0.698039,0.423529}%
\pgfsetfillcolor{currentfill}%
\pgfsetfillopacity{0.120000}%
\pgfsetlinewidth{0.000000pt}%
\definecolor{currentstroke}{rgb}{0.000000,0.000000,0.000000}%
\pgfsetstrokecolor{currentstroke}%
\pgfsetdash{}{0pt}%
\pgfpathmoveto{\pgfqpoint{1.182767in}{1.127988in}}%
\pgfpathlineto{\pgfqpoint{1.198670in}{1.121278in}}%
\pgfpathlineto{\pgfqpoint{1.214609in}{1.115252in}}%
\pgfpathlineto{\pgfqpoint{1.222324in}{1.128783in}}%
\pgfpathlineto{\pgfqpoint{1.230024in}{1.142482in}}%
\pgfpathlineto{\pgfqpoint{1.214101in}{1.148362in}}%
\pgfpathlineto{\pgfqpoint{1.198213in}{1.154898in}}%
\pgfpathlineto{\pgfqpoint{1.190497in}{1.141358in}}%
\pgfpathlineto{\pgfqpoint{1.182767in}{1.127988in}}%
\pgfpathclose%
\pgfusepath{fill}%
\end{pgfscope}%
\begin{pgfscope}%
\pgfpathrectangle{\pgfqpoint{0.407592in}{0.222904in}}{\pgfqpoint{1.411163in}{1.411163in}}%
\pgfusepath{clip}%
\pgfsetbuttcap%
\pgfsetroundjoin%
\definecolor{currentfill}{rgb}{0.909804,0.698039,0.423529}%
\pgfsetfillcolor{currentfill}%
\pgfsetfillopacity{0.120000}%
\pgfsetlinewidth{0.000000pt}%
\definecolor{currentstroke}{rgb}{0.000000,0.000000,0.000000}%
\pgfsetstrokecolor{currentstroke}%
\pgfsetdash{}{0pt}%
\pgfpathmoveto{\pgfqpoint{1.135516in}{1.117660in}}%
\pgfpathlineto{\pgfqpoint{1.151372in}{1.109402in}}%
\pgfpathlineto{\pgfqpoint{1.167259in}{1.101817in}}%
\pgfpathlineto{\pgfqpoint{1.175021in}{1.114803in}}%
\pgfpathlineto{\pgfqpoint{1.182767in}{1.127988in}}%
\pgfpathlineto{\pgfqpoint{1.166897in}{1.135365in}}%
\pgfpathlineto{\pgfqpoint{1.151060in}{1.143390in}}%
\pgfpathlineto{\pgfqpoint{1.143296in}{1.130424in}}%
\pgfpathlineto{\pgfqpoint{1.135516in}{1.117660in}}%
\pgfpathclose%
\pgfusepath{fill}%
\end{pgfscope}%
\begin{pgfscope}%
\pgfpathrectangle{\pgfqpoint{0.407592in}{0.222904in}}{\pgfqpoint{1.411163in}{1.411163in}}%
\pgfusepath{clip}%
\pgfsetbuttcap%
\pgfsetroundjoin%
\definecolor{currentfill}{rgb}{0.909804,0.698039,0.423529}%
\pgfsetfillcolor{currentfill}%
\pgfsetfillopacity{0.120000}%
\pgfsetlinewidth{0.000000pt}%
\definecolor{currentstroke}{rgb}{0.000000,0.000000,0.000000}%
\pgfsetstrokecolor{currentstroke}%
\pgfsetdash{}{0pt}%
\pgfpathmoveto{\pgfqpoint{1.230024in}{1.142482in}}%
\pgfpathlineto{\pgfqpoint{1.245987in}{1.137269in}}%
\pgfpathlineto{\pgfqpoint{1.261989in}{1.132735in}}%
\pgfpathlineto{\pgfqpoint{1.269661in}{1.146685in}}%
\pgfpathlineto{\pgfqpoint{1.277319in}{1.160778in}}%
\pgfpathlineto{\pgfqpoint{1.261331in}{1.165232in}}%
\pgfpathlineto{\pgfqpoint{1.245382in}{1.170334in}}%
\pgfpathlineto{\pgfqpoint{1.237711in}{1.156336in}}%
\pgfpathlineto{\pgfqpoint{1.230024in}{1.142482in}}%
\pgfpathclose%
\pgfusepath{fill}%
\end{pgfscope}%
\begin{pgfscope}%
\pgfpathrectangle{\pgfqpoint{0.407592in}{0.222904in}}{\pgfqpoint{1.411163in}{1.411163in}}%
\pgfusepath{clip}%
\pgfsetbuttcap%
\pgfsetroundjoin%
\definecolor{currentfill}{rgb}{0.909804,0.698039,0.423529}%
\pgfsetfillcolor{currentfill}%
\pgfsetfillopacity{0.120000}%
\pgfsetlinewidth{0.000000pt}%
\definecolor{currentstroke}{rgb}{0.000000,0.000000,0.000000}%
\pgfsetstrokecolor{currentstroke}%
\pgfsetdash{}{0pt}%
\pgfpathmoveto{\pgfqpoint{0.882935in}{1.180860in}}%
\pgfpathlineto{\pgfqpoint{0.898657in}{1.165472in}}%
\pgfpathlineto{\pgfqpoint{0.914392in}{1.150408in}}%
\pgfpathlineto{\pgfqpoint{0.922462in}{1.159521in}}%
\pgfpathlineto{\pgfqpoint{0.930509in}{1.168936in}}%
\pgfpathlineto{\pgfqpoint{0.914806in}{1.183708in}}%
\pgfpathlineto{\pgfqpoint{0.899118in}{1.198808in}}%
\pgfpathlineto{\pgfqpoint{0.891038in}{1.189689in}}%
\pgfpathlineto{\pgfqpoint{0.882935in}{1.180860in}}%
\pgfpathclose%
\pgfusepath{fill}%
\end{pgfscope}%
\begin{pgfscope}%
\pgfpathrectangle{\pgfqpoint{0.407592in}{0.222904in}}{\pgfqpoint{1.411163in}{1.411163in}}%
\pgfusepath{clip}%
\pgfsetbuttcap%
\pgfsetroundjoin%
\definecolor{currentfill}{rgb}{0.909804,0.698039,0.423529}%
\pgfsetfillcolor{currentfill}%
\pgfsetfillopacity{0.120000}%
\pgfsetlinewidth{0.000000pt}%
\definecolor{currentstroke}{rgb}{0.000000,0.000000,0.000000}%
\pgfsetstrokecolor{currentstroke}%
\pgfsetdash{}{0pt}%
\pgfpathmoveto{\pgfqpoint{0.961958in}{1.140489in}}%
\pgfpathlineto{\pgfqpoint{0.977706in}{1.126874in}}%
\pgfpathlineto{\pgfqpoint{0.993470in}{1.113709in}}%
\pgfpathlineto{\pgfqpoint{1.001433in}{1.124043in}}%
\pgfpathlineto{\pgfqpoint{1.009375in}{1.134669in}}%
\pgfpathlineto{\pgfqpoint{0.993639in}{1.147515in}}%
\pgfpathlineto{\pgfqpoint{0.977920in}{1.160811in}}%
\pgfpathlineto{\pgfqpoint{0.969950in}{1.150507in}}%
\pgfpathlineto{\pgfqpoint{0.961958in}{1.140489in}}%
\pgfpathclose%
\pgfusepath{fill}%
\end{pgfscope}%
\begin{pgfscope}%
\pgfpathrectangle{\pgfqpoint{0.407592in}{0.222904in}}{\pgfqpoint{1.411163in}{1.411163in}}%
\pgfusepath{clip}%
\pgfsetbuttcap%
\pgfsetroundjoin%
\definecolor{currentfill}{rgb}{0.909804,0.698039,0.423529}%
\pgfsetfillcolor{currentfill}%
\pgfsetfillopacity{0.120000}%
\pgfsetlinewidth{0.000000pt}%
\definecolor{currentstroke}{rgb}{0.000000,0.000000,0.000000}%
\pgfsetstrokecolor{currentstroke}%
\pgfsetdash{}{0pt}%
\pgfpathmoveto{\pgfqpoint{1.088241in}{1.111769in}}%
\pgfpathlineto{\pgfqpoint{1.104061in}{1.101965in}}%
\pgfpathlineto{\pgfqpoint{1.119907in}{1.092800in}}%
\pgfpathlineto{\pgfqpoint{1.127720in}{1.105113in}}%
\pgfpathlineto{\pgfqpoint{1.135516in}{1.117660in}}%
\pgfpathlineto{\pgfqpoint{1.119690in}{1.126565in}}%
\pgfpathlineto{\pgfqpoint{1.103892in}{1.136089in}}%
\pgfpathlineto{\pgfqpoint{1.096075in}{1.123812in}}%
\pgfpathlineto{\pgfqpoint{1.088241in}{1.111769in}}%
\pgfpathclose%
\pgfusepath{fill}%
\end{pgfscope}%
\begin{pgfscope}%
\pgfpathrectangle{\pgfqpoint{0.407592in}{0.222904in}}{\pgfqpoint{1.411163in}{1.411163in}}%
\pgfusepath{clip}%
\pgfsetbuttcap%
\pgfsetroundjoin%
\definecolor{currentfill}{rgb}{0.909804,0.698039,0.423529}%
\pgfsetfillcolor{currentfill}%
\pgfsetfillopacity{0.120000}%
\pgfsetlinewidth{0.000000pt}%
\definecolor{currentstroke}{rgb}{0.000000,0.000000,0.000000}%
\pgfsetstrokecolor{currentstroke}%
\pgfsetdash{}{0pt}%
\pgfpathmoveto{\pgfqpoint{0.803740in}{1.229394in}}%
\pgfpathlineto{\pgfqpoint{0.819451in}{1.212717in}}%
\pgfpathlineto{\pgfqpoint{0.835173in}{1.196263in}}%
\pgfpathlineto{\pgfqpoint{0.843363in}{1.204235in}}%
\pgfpathlineto{\pgfqpoint{0.851527in}{1.212502in}}%
\pgfpathlineto{\pgfqpoint{0.835842in}{1.228710in}}%
\pgfpathlineto{\pgfqpoint{0.820167in}{1.245148in}}%
\pgfpathlineto{\pgfqpoint{0.811967in}{1.237130in}}%
\pgfpathlineto{\pgfqpoint{0.803740in}{1.229394in}}%
\pgfpathclose%
\pgfusepath{fill}%
\end{pgfscope}%
\begin{pgfscope}%
\pgfpathrectangle{\pgfqpoint{0.407592in}{0.222904in}}{\pgfqpoint{1.411163in}{1.411163in}}%
\pgfusepath{clip}%
\pgfsetbuttcap%
\pgfsetroundjoin%
\definecolor{currentfill}{rgb}{0.909804,0.698039,0.423529}%
\pgfsetfillcolor{currentfill}%
\pgfsetfillopacity{0.120000}%
\pgfsetlinewidth{0.000000pt}%
\definecolor{currentstroke}{rgb}{0.000000,0.000000,0.000000}%
\pgfsetstrokecolor{currentstroke}%
\pgfsetdash{}{0pt}%
\pgfpathmoveto{\pgfqpoint{1.040904in}{1.110457in}}%
\pgfpathlineto{\pgfqpoint{1.056700in}{1.099161in}}%
\pgfpathlineto{\pgfqpoint{1.072518in}{1.088453in}}%
\pgfpathlineto{\pgfqpoint{1.080389in}{1.099977in}}%
\pgfpathlineto{\pgfqpoint{1.088241in}{1.111769in}}%
\pgfpathlineto{\pgfqpoint{1.072446in}{1.122182in}}%
\pgfpathlineto{\pgfqpoint{1.056675in}{1.133169in}}%
\pgfpathlineto{\pgfqpoint{1.048799in}{1.121680in}}%
\pgfpathlineto{\pgfqpoint{1.040904in}{1.110457in}}%
\pgfpathclose%
\pgfusepath{fill}%
\end{pgfscope}%
\begin{pgfscope}%
\pgfpathrectangle{\pgfqpoint{0.407592in}{0.222904in}}{\pgfqpoint{1.411163in}{1.411163in}}%
\pgfusepath{clip}%
\pgfsetbuttcap%
\pgfsetroundjoin%
\definecolor{currentfill}{rgb}{0.909804,0.698039,0.423529}%
\pgfsetfillcolor{currentfill}%
\pgfsetfillopacity{0.120000}%
\pgfsetlinewidth{0.000000pt}%
\definecolor{currentstroke}{rgb}{0.000000,0.000000,0.000000}%
\pgfsetstrokecolor{currentstroke}%
\pgfsetdash{}{0pt}%
\pgfpathmoveto{\pgfqpoint{1.277319in}{1.160778in}}%
\pgfpathlineto{\pgfqpoint{1.293351in}{1.156977in}}%
\pgfpathlineto{\pgfqpoint{1.309427in}{1.153831in}}%
\pgfpathlineto{\pgfqpoint{1.317059in}{1.168089in}}%
\pgfpathlineto{\pgfqpoint{1.300989in}{1.171225in}}%
\pgfpathlineto{\pgfqpoint{1.284964in}{1.175001in}}%
\pgfusepath{fill}%
\end{pgfscope}%
\begin{pgfscope}%
\pgfpathrectangle{\pgfqpoint{0.407592in}{0.222904in}}{\pgfqpoint{1.411163in}{1.411163in}}%
\pgfusepath{clip}%
\pgfsetbuttcap%
\pgfsetroundjoin%
\definecolor{currentfill}{rgb}{0.909804,0.698039,0.423529}%
\pgfsetfillcolor{currentfill}%
\pgfsetfillopacity{0.120000}%
\pgfsetlinewidth{0.000000pt}%
\definecolor{currentstroke}{rgb}{0.000000,0.000000,0.000000}%
\pgfsetstrokecolor{currentstroke}%
\pgfsetdash{}{0pt}%
\pgfpathmoveto{\pgfqpoint{0.914392in}{1.150408in}}%
\pgfpathlineto{\pgfqpoint{0.930142in}{1.135694in}}%
\pgfpathlineto{\pgfqpoint{0.945905in}{1.121363in}}%
\pgfpathlineto{\pgfqpoint{0.953943in}{1.130770in}}%
\pgfpathlineto{\pgfqpoint{0.961958in}{1.140489in}}%
\pgfpathlineto{\pgfqpoint{0.946226in}{1.154520in}}%
\pgfpathlineto{\pgfqpoint{0.930509in}{1.168936in}}%
\pgfpathlineto{\pgfqpoint{0.922462in}{1.159521in}}%
\pgfpathlineto{\pgfqpoint{0.914392in}{1.150408in}}%
\pgfpathclose%
\pgfusepath{fill}%
\end{pgfscope}%
\begin{pgfscope}%
\pgfpathrectangle{\pgfqpoint{0.407592in}{0.222904in}}{\pgfqpoint{1.411163in}{1.411163in}}%
\pgfusepath{clip}%
\pgfsetbuttcap%
\pgfsetroundjoin%
\definecolor{currentfill}{rgb}{0.909804,0.698039,0.423529}%
\pgfsetfillcolor{currentfill}%
\pgfsetfillopacity{0.120000}%
\pgfsetlinewidth{0.000000pt}%
\definecolor{currentstroke}{rgb}{0.000000,0.000000,0.000000}%
\pgfsetstrokecolor{currentstroke}%
\pgfsetdash{}{0pt}%
\pgfpathmoveto{\pgfqpoint{0.724311in}{1.284190in}}%
\pgfpathlineto{\pgfqpoint{0.740019in}{1.266597in}}%
\pgfpathlineto{\pgfqpoint{0.755736in}{1.249158in}}%
\pgfpathlineto{\pgfqpoint{0.764058in}{1.256112in}}%
\pgfpathlineto{\pgfqpoint{0.772350in}{1.263348in}}%
\pgfpathlineto{\pgfqpoint{0.756671in}{1.280592in}}%
\pgfpathlineto{\pgfqpoint{0.741002in}{1.297996in}}%
\pgfpathlineto{\pgfqpoint{0.732671in}{1.290959in}}%
\pgfpathlineto{\pgfqpoint{0.724311in}{1.284190in}}%
\pgfpathclose%
\pgfusepath{fill}%
\end{pgfscope}%
\begin{pgfscope}%
\pgfpathrectangle{\pgfqpoint{0.407592in}{0.222904in}}{\pgfqpoint{1.411163in}{1.411163in}}%
\pgfusepath{clip}%
\pgfsetbuttcap%
\pgfsetroundjoin%
\definecolor{currentfill}{rgb}{0.909804,0.698039,0.423529}%
\pgfsetfillcolor{currentfill}%
\pgfsetfillopacity{0.120000}%
\pgfsetlinewidth{0.000000pt}%
\definecolor{currentstroke}{rgb}{0.000000,0.000000,0.000000}%
\pgfsetstrokecolor{currentstroke}%
\pgfsetdash{}{0pt}%
\pgfpathmoveto{\pgfqpoint{0.993470in}{1.113709in}}%
\pgfpathlineto{\pgfqpoint{1.009253in}{1.101030in}}%
\pgfpathlineto{\pgfqpoint{1.025054in}{1.088872in}}%
\pgfpathlineto{\pgfqpoint{1.032989in}{1.099515in}}%
\pgfpathlineto{\pgfqpoint{1.040904in}{1.110457in}}%
\pgfpathlineto{\pgfqpoint{1.025129in}{1.122305in}}%
\pgfpathlineto{\pgfqpoint{1.009375in}{1.134669in}}%
\pgfpathlineto{\pgfqpoint{1.001433in}{1.124043in}}%
\pgfpathlineto{\pgfqpoint{0.993470in}{1.113709in}}%
\pgfpathclose%
\pgfusepath{fill}%
\end{pgfscope}%
\begin{pgfscope}%
\pgfpathrectangle{\pgfqpoint{0.407592in}{0.222904in}}{\pgfqpoint{1.411163in}{1.411163in}}%
\pgfusepath{clip}%
\pgfsetbuttcap%
\pgfsetroundjoin%
\definecolor{currentfill}{rgb}{0.909804,0.698039,0.423529}%
\pgfsetfillcolor{currentfill}%
\pgfsetfillopacity{0.120000}%
\pgfsetlinewidth{0.000000pt}%
\definecolor{currentstroke}{rgb}{0.000000,0.000000,0.000000}%
\pgfsetstrokecolor{currentstroke}%
\pgfsetdash{}{0pt}%
\pgfpathmoveto{\pgfqpoint{0.835173in}{1.196263in}}%
\pgfpathlineto{\pgfqpoint{0.850906in}{1.180055in}}%
\pgfpathlineto{\pgfqpoint{0.866651in}{1.164113in}}%
\pgfpathlineto{\pgfqpoint{0.874806in}{1.172332in}}%
\pgfpathlineto{\pgfqpoint{0.882935in}{1.180860in}}%
\pgfpathlineto{\pgfqpoint{0.867225in}{1.196545in}}%
\pgfpathlineto{\pgfqpoint{0.851527in}{1.212502in}}%
\pgfpathlineto{\pgfqpoint{0.843363in}{1.204235in}}%
\pgfpathlineto{\pgfqpoint{0.835173in}{1.196263in}}%
\pgfpathclose%
\pgfusepath{fill}%
\end{pgfscope}%
\begin{pgfscope}%
\pgfpathrectangle{\pgfqpoint{0.407592in}{0.222904in}}{\pgfqpoint{1.411163in}{1.411163in}}%
\pgfusepath{clip}%
\pgfsetbuttcap%
\pgfsetroundjoin%
\definecolor{currentfill}{rgb}{0.909804,0.698039,0.423529}%
\pgfsetfillcolor{currentfill}%
\pgfsetfillopacity{0.120000}%
\pgfsetlinewidth{0.000000pt}%
\definecolor{currentstroke}{rgb}{0.000000,0.000000,0.000000}%
\pgfsetstrokecolor{currentstroke}%
\pgfsetdash{}{0pt}%
\pgfpathmoveto{\pgfqpoint{1.167259in}{1.101817in}}%
\pgfpathlineto{\pgfqpoint{1.183179in}{1.094928in}}%
\pgfpathlineto{\pgfqpoint{1.199134in}{1.088752in}}%
\pgfpathlineto{\pgfqpoint{1.206879in}{1.101904in}}%
\pgfpathlineto{\pgfqpoint{1.214609in}{1.115252in}}%
\pgfpathlineto{\pgfqpoint{1.198670in}{1.121278in}}%
\pgfpathlineto{\pgfqpoint{1.182767in}{1.127988in}}%
\pgfpathlineto{\pgfqpoint{1.175021in}{1.114803in}}%
\pgfpathlineto{\pgfqpoint{1.167259in}{1.101817in}}%
\pgfpathclose%
\pgfusepath{fill}%
\end{pgfscope}%
\begin{pgfscope}%
\pgfpathrectangle{\pgfqpoint{0.407592in}{0.222904in}}{\pgfqpoint{1.411163in}{1.411163in}}%
\pgfusepath{clip}%
\pgfsetbuttcap%
\pgfsetroundjoin%
\definecolor{currentfill}{rgb}{0.909804,0.698039,0.423529}%
\pgfsetfillcolor{currentfill}%
\pgfsetfillopacity{0.120000}%
\pgfsetlinewidth{0.000000pt}%
\definecolor{currentstroke}{rgb}{0.000000,0.000000,0.000000}%
\pgfsetstrokecolor{currentstroke}%
\pgfsetdash{}{0pt}%
\pgfpathmoveto{\pgfqpoint{1.119907in}{1.092800in}}%
\pgfpathlineto{\pgfqpoint{1.135782in}{1.084305in}}%
\pgfpathlineto{\pgfqpoint{1.151687in}{1.076507in}}%
\pgfpathlineto{\pgfqpoint{1.159481in}{1.089047in}}%
\pgfpathlineto{\pgfqpoint{1.167259in}{1.101817in}}%
\pgfpathlineto{\pgfqpoint{1.151372in}{1.109402in}}%
\pgfpathlineto{\pgfqpoint{1.135516in}{1.117660in}}%
\pgfpathlineto{\pgfqpoint{1.127720in}{1.105113in}}%
\pgfpathlineto{\pgfqpoint{1.119907in}{1.092800in}}%
\pgfpathclose%
\pgfusepath{fill}%
\end{pgfscope}%
\begin{pgfscope}%
\pgfpathrectangle{\pgfqpoint{0.407592in}{0.222904in}}{\pgfqpoint{1.411163in}{1.411163in}}%
\pgfusepath{clip}%
\pgfsetbuttcap%
\pgfsetroundjoin%
\definecolor{currentfill}{rgb}{0.909804,0.698039,0.423529}%
\pgfsetfillcolor{currentfill}%
\pgfsetfillopacity{0.120000}%
\pgfsetlinewidth{0.000000pt}%
\definecolor{currentstroke}{rgb}{0.000000,0.000000,0.000000}%
\pgfsetstrokecolor{currentstroke}%
\pgfsetdash{}{0pt}%
\pgfpathmoveto{\pgfqpoint{1.214609in}{1.115252in}}%
\pgfpathlineto{\pgfqpoint{1.230586in}{1.109925in}}%
\pgfpathlineto{\pgfqpoint{1.246603in}{1.105308in}}%
\pgfpathlineto{\pgfqpoint{1.254303in}{1.118938in}}%
\pgfpathlineto{\pgfqpoint{1.261989in}{1.132735in}}%
\pgfpathlineto{\pgfqpoint{1.245987in}{1.137269in}}%
\pgfpathlineto{\pgfqpoint{1.230024in}{1.142482in}}%
\pgfpathlineto{\pgfqpoint{1.222324in}{1.128783in}}%
\pgfpathlineto{\pgfqpoint{1.214609in}{1.115252in}}%
\pgfpathclose%
\pgfusepath{fill}%
\end{pgfscope}%
\begin{pgfscope}%
\pgfpathrectangle{\pgfqpoint{0.407592in}{0.222904in}}{\pgfqpoint{1.411163in}{1.411163in}}%
\pgfusepath{clip}%
\pgfsetbuttcap%
\pgfsetroundjoin%
\definecolor{currentfill}{rgb}{0.909804,0.698039,0.423529}%
\pgfsetfillcolor{currentfill}%
\pgfsetfillopacity{0.120000}%
\pgfsetlinewidth{0.000000pt}%
\definecolor{currentstroke}{rgb}{0.000000,0.000000,0.000000}%
\pgfsetstrokecolor{currentstroke}%
\pgfsetdash{}{0pt}%
\pgfpathmoveto{\pgfqpoint{0.755736in}{1.249158in}}%
\pgfpathlineto{\pgfqpoint{0.771464in}{1.231887in}}%
\pgfpathlineto{\pgfqpoint{0.787202in}{1.214799in}}%
\pgfpathlineto{\pgfqpoint{0.795485in}{1.221948in}}%
\pgfpathlineto{\pgfqpoint{0.803740in}{1.229394in}}%
\pgfpathlineto{\pgfqpoint{0.788040in}{1.246277in}}%
\pgfpathlineto{\pgfqpoint{0.772350in}{1.263348in}}%
\pgfpathlineto{\pgfqpoint{0.764058in}{1.256112in}}%
\pgfpathlineto{\pgfqpoint{0.755736in}{1.249158in}}%
\pgfpathclose%
\pgfusepath{fill}%
\end{pgfscope}%
\begin{pgfscope}%
\pgfpathrectangle{\pgfqpoint{0.407592in}{0.222904in}}{\pgfqpoint{1.411163in}{1.411163in}}%
\pgfusepath{clip}%
\pgfsetbuttcap%
\pgfsetroundjoin%
\definecolor{currentfill}{rgb}{0.909804,0.698039,0.423529}%
\pgfsetfillcolor{currentfill}%
\pgfsetfillopacity{0.120000}%
\pgfsetlinewidth{0.000000pt}%
\definecolor{currentstroke}{rgb}{0.000000,0.000000,0.000000}%
\pgfsetstrokecolor{currentstroke}%
\pgfsetdash{}{0pt}%
\pgfpathmoveto{\pgfqpoint{1.261989in}{1.132735in}}%
\pgfpathlineto{\pgfqpoint{1.278034in}{1.128884in}}%
\pgfpathlineto{\pgfqpoint{1.294124in}{1.125720in}}%
\pgfpathlineto{\pgfqpoint{1.301782in}{1.139705in}}%
\pgfpathlineto{\pgfqpoint{1.309427in}{1.153831in}}%
\pgfpathlineto{\pgfqpoint{1.293351in}{1.156977in}}%
\pgfpathlineto{\pgfqpoint{1.277319in}{1.160778in}}%
\pgfpathlineto{\pgfqpoint{1.269661in}{1.146685in}}%
\pgfpathlineto{\pgfqpoint{1.261989in}{1.132735in}}%
\pgfpathclose%
\pgfusepath{fill}%
\end{pgfscope}%
\begin{pgfscope}%
\pgfpathrectangle{\pgfqpoint{0.407592in}{0.222904in}}{\pgfqpoint{1.411163in}{1.411163in}}%
\pgfusepath{clip}%
\pgfsetbuttcap%
\pgfsetroundjoin%
\definecolor{currentfill}{rgb}{0.909804,0.698039,0.423529}%
\pgfsetfillcolor{currentfill}%
\pgfsetfillopacity{0.120000}%
\pgfsetlinewidth{0.000000pt}%
\definecolor{currentstroke}{rgb}{0.000000,0.000000,0.000000}%
\pgfsetstrokecolor{currentstroke}%
\pgfsetdash{}{0pt}%
\pgfpathmoveto{\pgfqpoint{1.072518in}{1.088453in}}%
\pgfpathlineto{\pgfqpoint{1.088360in}{1.078367in}}%
\pgfpathlineto{\pgfqpoint{1.104228in}{1.068940in}}%
\pgfpathlineto{\pgfqpoint{1.112077in}{1.080736in}}%
\pgfpathlineto{\pgfqpoint{1.119907in}{1.092800in}}%
\pgfpathlineto{\pgfqpoint{1.104061in}{1.101965in}}%
\pgfpathlineto{\pgfqpoint{1.088241in}{1.111769in}}%
\pgfpathlineto{\pgfqpoint{1.080389in}{1.099977in}}%
\pgfpathlineto{\pgfqpoint{1.072518in}{1.088453in}}%
\pgfpathclose%
\pgfusepath{fill}%
\end{pgfscope}%
\begin{pgfscope}%
\pgfpathrectangle{\pgfqpoint{0.407592in}{0.222904in}}{\pgfqpoint{1.411163in}{1.411163in}}%
\pgfusepath{clip}%
\pgfsetbuttcap%
\pgfsetroundjoin%
\definecolor{currentfill}{rgb}{0.909804,0.698039,0.423529}%
\pgfsetfillcolor{currentfill}%
\pgfsetfillopacity{0.120000}%
\pgfsetlinewidth{0.000000pt}%
\definecolor{currentstroke}{rgb}{0.000000,0.000000,0.000000}%
\pgfsetstrokecolor{currentstroke}%
\pgfsetdash{}{0pt}%
\pgfpathmoveto{\pgfqpoint{0.945905in}{1.121363in}}%
\pgfpathlineto{\pgfqpoint{0.961684in}{1.107445in}}%
\pgfpathlineto{\pgfqpoint{0.977479in}{1.093978in}}%
\pgfpathlineto{\pgfqpoint{0.985486in}{1.103683in}}%
\pgfpathlineto{\pgfqpoint{0.993470in}{1.113709in}}%
\pgfpathlineto{\pgfqpoint{0.977706in}{1.126874in}}%
\pgfpathlineto{\pgfqpoint{0.961958in}{1.140489in}}%
\pgfpathlineto{\pgfqpoint{0.953943in}{1.130770in}}%
\pgfpathlineto{\pgfqpoint{0.945905in}{1.121363in}}%
\pgfpathclose%
\pgfusepath{fill}%
\end{pgfscope}%
\begin{pgfscope}%
\pgfpathrectangle{\pgfqpoint{0.407592in}{0.222904in}}{\pgfqpoint{1.411163in}{1.411163in}}%
\pgfusepath{clip}%
\pgfsetbuttcap%
\pgfsetroundjoin%
\definecolor{currentfill}{rgb}{0.909804,0.698039,0.423529}%
\pgfsetfillcolor{currentfill}%
\pgfsetfillopacity{0.120000}%
\pgfsetlinewidth{0.000000pt}%
\definecolor{currentstroke}{rgb}{0.000000,0.000000,0.000000}%
\pgfsetstrokecolor{currentstroke}%
\pgfsetdash{}{0pt}%
\pgfpathmoveto{\pgfqpoint{0.866651in}{1.164113in}}%
\pgfpathlineto{\pgfqpoint{0.882408in}{1.148462in}}%
\pgfpathlineto{\pgfqpoint{0.898178in}{1.133129in}}%
\pgfpathlineto{\pgfqpoint{0.906298in}{1.141606in}}%
\pgfpathlineto{\pgfqpoint{0.914392in}{1.150408in}}%
\pgfpathlineto{\pgfqpoint{0.898657in}{1.165472in}}%
\pgfpathlineto{\pgfqpoint{0.882935in}{1.180860in}}%
\pgfpathlineto{\pgfqpoint{0.874806in}{1.172332in}}%
\pgfpathlineto{\pgfqpoint{0.866651in}{1.164113in}}%
\pgfpathclose%
\pgfusepath{fill}%
\end{pgfscope}%
\begin{pgfscope}%
\pgfpathrectangle{\pgfqpoint{0.407592in}{0.222904in}}{\pgfqpoint{1.411163in}{1.411163in}}%
\pgfusepath{clip}%
\pgfsetbuttcap%
\pgfsetroundjoin%
\definecolor{currentfill}{rgb}{0.909804,0.698039,0.423529}%
\pgfsetfillcolor{currentfill}%
\pgfsetfillopacity{0.120000}%
\pgfsetlinewidth{0.000000pt}%
\definecolor{currentstroke}{rgb}{0.000000,0.000000,0.000000}%
\pgfsetstrokecolor{currentstroke}%
\pgfsetdash{}{0pt}%
\pgfpathmoveto{\pgfqpoint{1.025054in}{1.088872in}}%
\pgfpathlineto{\pgfqpoint{1.040876in}{1.077273in}}%
\pgfpathlineto{\pgfqpoint{1.056719in}{1.066273in}}%
\pgfpathlineto{\pgfqpoint{1.064628in}{1.077212in}}%
\pgfpathlineto{\pgfqpoint{1.072518in}{1.088453in}}%
\pgfpathlineto{\pgfqpoint{1.056700in}{1.099161in}}%
\pgfpathlineto{\pgfqpoint{1.040904in}{1.110457in}}%
\pgfpathlineto{\pgfqpoint{1.032989in}{1.099515in}}%
\pgfpathlineto{\pgfqpoint{1.025054in}{1.088872in}}%
\pgfpathclose%
\pgfusepath{fill}%
\end{pgfscope}%
\begin{pgfscope}%
\pgfpathrectangle{\pgfqpoint{0.407592in}{0.222904in}}{\pgfqpoint{1.411163in}{1.411163in}}%
\pgfusepath{clip}%
\pgfsetbuttcap%
\pgfsetroundjoin%
\definecolor{currentfill}{rgb}{0.909804,0.698039,0.423529}%
\pgfsetfillcolor{currentfill}%
\pgfsetfillopacity{0.120000}%
\pgfsetlinewidth{0.000000pt}%
\definecolor{currentstroke}{rgb}{0.000000,0.000000,0.000000}%
\pgfsetstrokecolor{currentstroke}%
\pgfsetdash{}{0pt}%
\pgfpathmoveto{\pgfqpoint{0.787202in}{1.214799in}}%
\pgfpathlineto{\pgfqpoint{0.802950in}{1.197910in}}%
\pgfpathlineto{\pgfqpoint{0.818710in}{1.181241in}}%
\pgfpathlineto{\pgfqpoint{0.826955in}{1.188596in}}%
\pgfpathlineto{\pgfqpoint{0.835173in}{1.196263in}}%
\pgfpathlineto{\pgfqpoint{0.819451in}{1.212717in}}%
\pgfpathlineto{\pgfqpoint{0.803740in}{1.229394in}}%
\pgfpathlineto{\pgfqpoint{0.795485in}{1.221948in}}%
\pgfpathlineto{\pgfqpoint{0.787202in}{1.214799in}}%
\pgfpathclose%
\pgfusepath{fill}%
\end{pgfscope}%
\begin{pgfscope}%
\pgfpathrectangle{\pgfqpoint{0.407592in}{0.222904in}}{\pgfqpoint{1.411163in}{1.411163in}}%
\pgfusepath{clip}%
\pgfsetbuttcap%
\pgfsetroundjoin%
\definecolor{currentfill}{rgb}{0.909804,0.698039,0.423529}%
\pgfsetfillcolor{currentfill}%
\pgfsetfillopacity{0.120000}%
\pgfsetlinewidth{0.000000pt}%
\definecolor{currentstroke}{rgb}{0.000000,0.000000,0.000000}%
\pgfsetstrokecolor{currentstroke}%
\pgfsetdash{}{0pt}%
\pgfpathmoveto{\pgfqpoint{1.309427in}{1.153831in}}%
\pgfpathlineto{\pgfqpoint{1.325550in}{1.151341in}}%
\pgfpathlineto{\pgfqpoint{1.341720in}{1.149504in}}%
\pgfpathlineto{\pgfqpoint{1.349339in}{1.163737in}}%
\pgfpathlineto{\pgfqpoint{1.333175in}{1.165594in}}%
\pgfpathlineto{\pgfqpoint{1.317059in}{1.168089in}}%
\pgfusepath{fill}%
\end{pgfscope}%
\begin{pgfscope}%
\pgfpathrectangle{\pgfqpoint{0.407592in}{0.222904in}}{\pgfqpoint{1.411163in}{1.411163in}}%
\pgfusepath{clip}%
\pgfsetbuttcap%
\pgfsetroundjoin%
\definecolor{currentfill}{rgb}{0.909804,0.698039,0.423529}%
\pgfsetfillcolor{currentfill}%
\pgfsetfillopacity{0.120000}%
\pgfsetlinewidth{0.000000pt}%
\definecolor{currentstroke}{rgb}{0.000000,0.000000,0.000000}%
\pgfsetstrokecolor{currentstroke}%
\pgfsetdash{}{0pt}%
\pgfpathmoveto{\pgfqpoint{0.898178in}{1.133129in}}%
\pgfpathlineto{\pgfqpoint{0.913961in}{1.118143in}}%
\pgfpathlineto{\pgfqpoint{0.929758in}{1.103535in}}%
\pgfpathlineto{\pgfqpoint{0.937844in}{1.112280in}}%
\pgfpathlineto{\pgfqpoint{0.945905in}{1.121363in}}%
\pgfpathlineto{\pgfqpoint{0.930142in}{1.135694in}}%
\pgfpathlineto{\pgfqpoint{0.914392in}{1.150408in}}%
\pgfpathlineto{\pgfqpoint{0.906298in}{1.141606in}}%
\pgfpathlineto{\pgfqpoint{0.898178in}{1.133129in}}%
\pgfpathclose%
\pgfusepath{fill}%
\end{pgfscope}%
\begin{pgfscope}%
\pgfpathrectangle{\pgfqpoint{0.407592in}{0.222904in}}{\pgfqpoint{1.411163in}{1.411163in}}%
\pgfusepath{clip}%
\pgfsetbuttcap%
\pgfsetroundjoin%
\definecolor{currentfill}{rgb}{0.909804,0.698039,0.423529}%
\pgfsetfillcolor{currentfill}%
\pgfsetfillopacity{0.120000}%
\pgfsetlinewidth{0.000000pt}%
\definecolor{currentstroke}{rgb}{0.000000,0.000000,0.000000}%
\pgfsetstrokecolor{currentstroke}%
\pgfsetdash{}{0pt}%
\pgfpathmoveto{\pgfqpoint{0.977479in}{1.093978in}}%
\pgfpathlineto{\pgfqpoint{0.993291in}{1.080998in}}%
\pgfpathlineto{\pgfqpoint{1.009121in}{1.068543in}}%
\pgfpathlineto{\pgfqpoint{1.017098in}{1.078542in}}%
\pgfpathlineto{\pgfqpoint{1.025054in}{1.088872in}}%
\pgfpathlineto{\pgfqpoint{1.009253in}{1.101030in}}%
\pgfpathlineto{\pgfqpoint{0.993470in}{1.113709in}}%
\pgfpathlineto{\pgfqpoint{0.985486in}{1.103683in}}%
\pgfpathlineto{\pgfqpoint{0.977479in}{1.093978in}}%
\pgfpathclose%
\pgfusepath{fill}%
\end{pgfscope}%
\begin{pgfscope}%
\pgfpathrectangle{\pgfqpoint{0.407592in}{0.222904in}}{\pgfqpoint{1.411163in}{1.411163in}}%
\pgfusepath{clip}%
\pgfsetbuttcap%
\pgfsetroundjoin%
\definecolor{currentfill}{rgb}{0.909804,0.698039,0.423529}%
\pgfsetfillcolor{currentfill}%
\pgfsetfillopacity{0.120000}%
\pgfsetlinewidth{0.000000pt}%
\definecolor{currentstroke}{rgb}{0.000000,0.000000,0.000000}%
\pgfsetstrokecolor{currentstroke}%
\pgfsetdash{}{0pt}%
\pgfpathmoveto{\pgfqpoint{1.151687in}{1.076507in}}%
\pgfpathlineto{\pgfqpoint{1.167625in}{1.069434in}}%
\pgfpathlineto{\pgfqpoint{1.183598in}{1.063106in}}%
\pgfpathlineto{\pgfqpoint{1.191374in}{1.075814in}}%
\pgfpathlineto{\pgfqpoint{1.199134in}{1.088752in}}%
\pgfpathlineto{\pgfqpoint{1.183179in}{1.094928in}}%
\pgfpathlineto{\pgfqpoint{1.167259in}{1.101817in}}%
\pgfpathlineto{\pgfqpoint{1.159481in}{1.089047in}}%
\pgfpathlineto{\pgfqpoint{1.151687in}{1.076507in}}%
\pgfpathclose%
\pgfusepath{fill}%
\end{pgfscope}%
\begin{pgfscope}%
\pgfpathrectangle{\pgfqpoint{0.407592in}{0.222904in}}{\pgfqpoint{1.411163in}{1.411163in}}%
\pgfusepath{clip}%
\pgfsetbuttcap%
\pgfsetroundjoin%
\definecolor{currentfill}{rgb}{0.909804,0.698039,0.423529}%
\pgfsetfillcolor{currentfill}%
\pgfsetfillopacity{0.120000}%
\pgfsetlinewidth{0.000000pt}%
\definecolor{currentstroke}{rgb}{0.000000,0.000000,0.000000}%
\pgfsetstrokecolor{currentstroke}%
\pgfsetdash{}{0pt}%
\pgfpathmoveto{\pgfqpoint{1.199134in}{1.088752in}}%
\pgfpathlineto{\pgfqpoint{1.215127in}{1.083308in}}%
\pgfpathlineto{\pgfqpoint{1.231159in}{1.078608in}}%
\pgfpathlineto{\pgfqpoint{1.238888in}{1.091860in}}%
\pgfpathlineto{\pgfqpoint{1.246603in}{1.105308in}}%
\pgfpathlineto{\pgfqpoint{1.230586in}{1.109925in}}%
\pgfpathlineto{\pgfqpoint{1.214609in}{1.115252in}}%
\pgfpathlineto{\pgfqpoint{1.206879in}{1.101904in}}%
\pgfpathlineto{\pgfqpoint{1.199134in}{1.088752in}}%
\pgfpathclose%
\pgfusepath{fill}%
\end{pgfscope}%
\begin{pgfscope}%
\pgfpathrectangle{\pgfqpoint{0.407592in}{0.222904in}}{\pgfqpoint{1.411163in}{1.411163in}}%
\pgfusepath{clip}%
\pgfsetbuttcap%
\pgfsetroundjoin%
\definecolor{currentfill}{rgb}{0.909804,0.698039,0.423529}%
\pgfsetfillcolor{currentfill}%
\pgfsetfillopacity{0.120000}%
\pgfsetlinewidth{0.000000pt}%
\definecolor{currentstroke}{rgb}{0.000000,0.000000,0.000000}%
\pgfsetstrokecolor{currentstroke}%
\pgfsetdash{}{0pt}%
\pgfpathmoveto{\pgfqpoint{0.707502in}{1.271476in}}%
\pgfpathlineto{\pgfqpoint{0.723249in}{1.253721in}}%
\pgfpathlineto{\pgfqpoint{0.739005in}{1.236117in}}%
\pgfpathlineto{\pgfqpoint{0.747386in}{1.242491in}}%
\pgfpathlineto{\pgfqpoint{0.755736in}{1.249158in}}%
\pgfpathlineto{\pgfqpoint{0.740019in}{1.266597in}}%
\pgfpathlineto{\pgfqpoint{0.724311in}{1.284190in}}%
\pgfpathlineto{\pgfqpoint{0.715922in}{1.277694in}}%
\pgfpathlineto{\pgfqpoint{0.707502in}{1.271476in}}%
\pgfpathclose%
\pgfusepath{fill}%
\end{pgfscope}%
\begin{pgfscope}%
\pgfpathrectangle{\pgfqpoint{0.407592in}{0.222904in}}{\pgfqpoint{1.411163in}{1.411163in}}%
\pgfusepath{clip}%
\pgfsetbuttcap%
\pgfsetroundjoin%
\definecolor{currentfill}{rgb}{0.909804,0.698039,0.423529}%
\pgfsetfillcolor{currentfill}%
\pgfsetfillopacity{0.120000}%
\pgfsetlinewidth{0.000000pt}%
\definecolor{currentstroke}{rgb}{0.000000,0.000000,0.000000}%
\pgfsetstrokecolor{currentstroke}%
\pgfsetdash{}{0pt}%
\pgfpathmoveto{\pgfqpoint{1.104228in}{1.068940in}}%
\pgfpathlineto{\pgfqpoint{1.120124in}{1.060205in}}%
\pgfpathlineto{\pgfqpoint{1.136050in}{1.052193in}}%
\pgfpathlineto{\pgfqpoint{1.143877in}{1.064217in}}%
\pgfpathlineto{\pgfqpoint{1.151687in}{1.076507in}}%
\pgfpathlineto{\pgfqpoint{1.135782in}{1.084305in}}%
\pgfpathlineto{\pgfqpoint{1.119907in}{1.092800in}}%
\pgfpathlineto{\pgfqpoint{1.112077in}{1.080736in}}%
\pgfpathlineto{\pgfqpoint{1.104228in}{1.068940in}}%
\pgfpathclose%
\pgfusepath{fill}%
\end{pgfscope}%
\begin{pgfscope}%
\pgfpathrectangle{\pgfqpoint{0.407592in}{0.222904in}}{\pgfqpoint{1.411163in}{1.411163in}}%
\pgfusepath{clip}%
\pgfsetbuttcap%
\pgfsetroundjoin%
\definecolor{currentfill}{rgb}{0.909804,0.698039,0.423529}%
\pgfsetfillcolor{currentfill}%
\pgfsetfillopacity{0.120000}%
\pgfsetlinewidth{0.000000pt}%
\definecolor{currentstroke}{rgb}{0.000000,0.000000,0.000000}%
\pgfsetstrokecolor{currentstroke}%
\pgfsetdash{}{0pt}%
\pgfpathmoveto{\pgfqpoint{0.818710in}{1.181241in}}%
\pgfpathlineto{\pgfqpoint{0.834480in}{1.164810in}}%
\pgfpathlineto{\pgfqpoint{0.850261in}{1.148640in}}%
\pgfpathlineto{\pgfqpoint{0.858470in}{1.156213in}}%
\pgfpathlineto{\pgfqpoint{0.866651in}{1.164113in}}%
\pgfpathlineto{\pgfqpoint{0.850906in}{1.180055in}}%
\pgfpathlineto{\pgfqpoint{0.835173in}{1.196263in}}%
\pgfpathlineto{\pgfqpoint{0.826955in}{1.188596in}}%
\pgfpathlineto{\pgfqpoint{0.818710in}{1.181241in}}%
\pgfpathclose%
\pgfusepath{fill}%
\end{pgfscope}%
\begin{pgfscope}%
\pgfpathrectangle{\pgfqpoint{0.407592in}{0.222904in}}{\pgfqpoint{1.411163in}{1.411163in}}%
\pgfusepath{clip}%
\pgfsetbuttcap%
\pgfsetroundjoin%
\definecolor{currentfill}{rgb}{0.909804,0.698039,0.423529}%
\pgfsetfillcolor{currentfill}%
\pgfsetfillopacity{0.120000}%
\pgfsetlinewidth{0.000000pt}%
\definecolor{currentstroke}{rgb}{0.000000,0.000000,0.000000}%
\pgfsetstrokecolor{currentstroke}%
\pgfsetdash{}{0pt}%
\pgfpathmoveto{\pgfqpoint{1.246603in}{1.105308in}}%
\pgfpathlineto{\pgfqpoint{1.262662in}{1.101408in}}%
\pgfpathlineto{\pgfqpoint{1.278765in}{1.098228in}}%
\pgfpathlineto{\pgfqpoint{1.286451in}{1.111890in}}%
\pgfpathlineto{\pgfqpoint{1.294124in}{1.125720in}}%
\pgfpathlineto{\pgfqpoint{1.278034in}{1.128884in}}%
\pgfpathlineto{\pgfqpoint{1.261989in}{1.132735in}}%
\pgfpathlineto{\pgfqpoint{1.254303in}{1.118938in}}%
\pgfpathlineto{\pgfqpoint{1.246603in}{1.105308in}}%
\pgfpathclose%
\pgfusepath{fill}%
\end{pgfscope}%
\begin{pgfscope}%
\pgfpathrectangle{\pgfqpoint{0.407592in}{0.222904in}}{\pgfqpoint{1.411163in}{1.411163in}}%
\pgfusepath{clip}%
\pgfsetbuttcap%
\pgfsetroundjoin%
\definecolor{currentfill}{rgb}{0.909804,0.698039,0.423529}%
\pgfsetfillcolor{currentfill}%
\pgfsetfillopacity{0.120000}%
\pgfsetlinewidth{0.000000pt}%
\definecolor{currentstroke}{rgb}{0.000000,0.000000,0.000000}%
\pgfsetstrokecolor{currentstroke}%
\pgfsetdash{}{0pt}%
\pgfpathmoveto{\pgfqpoint{1.056719in}{1.066273in}}%
\pgfpathlineto{\pgfqpoint{1.072586in}{1.055911in}}%
\pgfpathlineto{\pgfqpoint{1.088477in}{1.046225in}}%
\pgfpathlineto{\pgfqpoint{1.096362in}{1.057430in}}%
\pgfpathlineto{\pgfqpoint{1.104228in}{1.068940in}}%
\pgfpathlineto{\pgfqpoint{1.088360in}{1.078367in}}%
\pgfpathlineto{\pgfqpoint{1.072518in}{1.088453in}}%
\pgfpathlineto{\pgfqpoint{1.064628in}{1.077212in}}%
\pgfpathlineto{\pgfqpoint{1.056719in}{1.066273in}}%
\pgfpathclose%
\pgfusepath{fill}%
\end{pgfscope}%
\begin{pgfscope}%
\pgfpathrectangle{\pgfqpoint{0.407592in}{0.222904in}}{\pgfqpoint{1.411163in}{1.411163in}}%
\pgfusepath{clip}%
\pgfsetbuttcap%
\pgfsetroundjoin%
\definecolor{currentfill}{rgb}{0.909804,0.698039,0.423529}%
\pgfsetfillcolor{currentfill}%
\pgfsetfillopacity{0.120000}%
\pgfsetlinewidth{0.000000pt}%
\definecolor{currentstroke}{rgb}{0.000000,0.000000,0.000000}%
\pgfsetstrokecolor{currentstroke}%
\pgfsetdash{}{0pt}%
\pgfpathmoveto{\pgfqpoint{1.294124in}{1.125720in}}%
\pgfpathlineto{\pgfqpoint{1.310259in}{1.123243in}}%
\pgfpathlineto{\pgfqpoint{1.326443in}{1.121449in}}%
\pgfpathlineto{\pgfqpoint{1.334089in}{1.135404in}}%
\pgfpathlineto{\pgfqpoint{1.341720in}{1.149504in}}%
\pgfpathlineto{\pgfqpoint{1.325550in}{1.151341in}}%
\pgfpathlineto{\pgfqpoint{1.309427in}{1.153831in}}%
\pgfpathlineto{\pgfqpoint{1.301782in}{1.139705in}}%
\pgfpathlineto{\pgfqpoint{1.294124in}{1.125720in}}%
\pgfpathclose%
\pgfusepath{fill}%
\end{pgfscope}%
\begin{pgfscope}%
\pgfpathrectangle{\pgfqpoint{0.407592in}{0.222904in}}{\pgfqpoint{1.411163in}{1.411163in}}%
\pgfusepath{clip}%
\pgfsetbuttcap%
\pgfsetroundjoin%
\definecolor{currentfill}{rgb}{0.909804,0.698039,0.423529}%
\pgfsetfillcolor{currentfill}%
\pgfsetfillopacity{0.120000}%
\pgfsetlinewidth{0.000000pt}%
\definecolor{currentstroke}{rgb}{0.000000,0.000000,0.000000}%
\pgfsetstrokecolor{currentstroke}%
\pgfsetdash{}{0pt}%
\pgfpathmoveto{\pgfqpoint{0.929758in}{1.103535in}}%
\pgfpathlineto{\pgfqpoint{0.945569in}{1.089340in}}%
\pgfpathlineto{\pgfqpoint{0.961396in}{1.075593in}}%
\pgfpathlineto{\pgfqpoint{0.969449in}{1.084610in}}%
\pgfpathlineto{\pgfqpoint{0.977479in}{1.093978in}}%
\pgfpathlineto{\pgfqpoint{0.961684in}{1.107445in}}%
\pgfpathlineto{\pgfqpoint{0.945905in}{1.121363in}}%
\pgfpathlineto{\pgfqpoint{0.937844in}{1.112280in}}%
\pgfpathlineto{\pgfqpoint{0.929758in}{1.103535in}}%
\pgfpathclose%
\pgfusepath{fill}%
\end{pgfscope}%
\begin{pgfscope}%
\pgfpathrectangle{\pgfqpoint{0.407592in}{0.222904in}}{\pgfqpoint{1.411163in}{1.411163in}}%
\pgfusepath{clip}%
\pgfsetbuttcap%
\pgfsetroundjoin%
\definecolor{currentfill}{rgb}{0.909804,0.698039,0.423529}%
\pgfsetfillcolor{currentfill}%
\pgfsetfillopacity{0.120000}%
\pgfsetlinewidth{0.000000pt}%
\definecolor{currentstroke}{rgb}{0.000000,0.000000,0.000000}%
\pgfsetstrokecolor{currentstroke}%
\pgfsetdash{}{0pt}%
\pgfpathmoveto{\pgfqpoint{0.739005in}{1.236117in}}%
\pgfpathlineto{\pgfqpoint{0.754772in}{1.218675in}}%
\pgfpathlineto{\pgfqpoint{0.770549in}{1.201413in}}%
\pgfpathlineto{\pgfqpoint{0.778890in}{1.207952in}}%
\pgfpathlineto{\pgfqpoint{0.787202in}{1.214799in}}%
\pgfpathlineto{\pgfqpoint{0.771464in}{1.231887in}}%
\pgfpathlineto{\pgfqpoint{0.755736in}{1.249158in}}%
\pgfpathlineto{\pgfqpoint{0.747386in}{1.242491in}}%
\pgfpathlineto{\pgfqpoint{0.739005in}{1.236117in}}%
\pgfpathclose%
\pgfusepath{fill}%
\end{pgfscope}%
\begin{pgfscope}%
\pgfpathrectangle{\pgfqpoint{0.407592in}{0.222904in}}{\pgfqpoint{1.411163in}{1.411163in}}%
\pgfusepath{clip}%
\pgfsetbuttcap%
\pgfsetroundjoin%
\definecolor{currentfill}{rgb}{0.909804,0.698039,0.423529}%
\pgfsetfillcolor{currentfill}%
\pgfsetfillopacity{0.120000}%
\pgfsetlinewidth{0.000000pt}%
\definecolor{currentstroke}{rgb}{0.000000,0.000000,0.000000}%
\pgfsetstrokecolor{currentstroke}%
\pgfsetdash{}{0pt}%
\pgfpathmoveto{\pgfqpoint{0.850261in}{1.148640in}}%
\pgfpathlineto{\pgfqpoint{0.866055in}{1.132757in}}%
\pgfpathlineto{\pgfqpoint{0.881860in}{1.117187in}}%
\pgfpathlineto{\pgfqpoint{0.890032in}{1.124986in}}%
\pgfpathlineto{\pgfqpoint{0.898178in}{1.133129in}}%
\pgfpathlineto{\pgfqpoint{0.882408in}{1.148462in}}%
\pgfpathlineto{\pgfqpoint{0.866651in}{1.164113in}}%
\pgfpathlineto{\pgfqpoint{0.858470in}{1.156213in}}%
\pgfpathlineto{\pgfqpoint{0.850261in}{1.148640in}}%
\pgfpathclose%
\pgfusepath{fill}%
\end{pgfscope}%
\begin{pgfscope}%
\pgfpathrectangle{\pgfqpoint{0.407592in}{0.222904in}}{\pgfqpoint{1.411163in}{1.411163in}}%
\pgfusepath{clip}%
\pgfsetbuttcap%
\pgfsetroundjoin%
\definecolor{currentfill}{rgb}{0.909804,0.698039,0.423529}%
\pgfsetfillcolor{currentfill}%
\pgfsetfillopacity{0.120000}%
\pgfsetlinewidth{0.000000pt}%
\definecolor{currentstroke}{rgb}{0.000000,0.000000,0.000000}%
\pgfsetstrokecolor{currentstroke}%
\pgfsetdash{}{0pt}%
\pgfpathmoveto{\pgfqpoint{1.009121in}{1.068543in}}%
\pgfpathlineto{\pgfqpoint{1.024970in}{1.056655in}}%
\pgfpathlineto{\pgfqpoint{1.040840in}{1.045374in}}%
\pgfpathlineto{\pgfqpoint{1.048790in}{1.055655in}}%
\pgfpathlineto{\pgfqpoint{1.056719in}{1.066273in}}%
\pgfpathlineto{\pgfqpoint{1.040876in}{1.077273in}}%
\pgfpathlineto{\pgfqpoint{1.025054in}{1.088872in}}%
\pgfpathlineto{\pgfqpoint{1.017098in}{1.078542in}}%
\pgfpathlineto{\pgfqpoint{1.009121in}{1.068543in}}%
\pgfpathclose%
\pgfusepath{fill}%
\end{pgfscope}%
\begin{pgfscope}%
\pgfpathrectangle{\pgfqpoint{0.407592in}{0.222904in}}{\pgfqpoint{1.411163in}{1.411163in}}%
\pgfusepath{clip}%
\pgfsetbuttcap%
\pgfsetroundjoin%
\definecolor{currentfill}{rgb}{0.909804,0.698039,0.423529}%
\pgfsetfillcolor{currentfill}%
\pgfsetfillopacity{0.120000}%
\pgfsetlinewidth{0.000000pt}%
\definecolor{currentstroke}{rgb}{0.000000,0.000000,0.000000}%
\pgfsetstrokecolor{currentstroke}%
\pgfsetdash{}{0pt}%
\pgfpathmoveto{\pgfqpoint{0.770549in}{1.201413in}}%
\pgfpathlineto{\pgfqpoint{0.786336in}{1.184345in}}%
\pgfpathlineto{\pgfqpoint{0.802133in}{1.167492in}}%
\pgfpathlineto{\pgfqpoint{0.810436in}{1.174204in}}%
\pgfpathlineto{\pgfqpoint{0.818710in}{1.181241in}}%
\pgfpathlineto{\pgfqpoint{0.802950in}{1.197910in}}%
\pgfpathlineto{\pgfqpoint{0.787202in}{1.214799in}}%
\pgfpathlineto{\pgfqpoint{0.778890in}{1.207952in}}%
\pgfpathlineto{\pgfqpoint{0.770549in}{1.201413in}}%
\pgfpathclose%
\pgfusepath{fill}%
\end{pgfscope}%
\begin{pgfscope}%
\pgfpathrectangle{\pgfqpoint{0.407592in}{0.222904in}}{\pgfqpoint{1.411163in}{1.411163in}}%
\pgfusepath{clip}%
\pgfsetbuttcap%
\pgfsetroundjoin%
\definecolor{currentfill}{rgb}{0.909804,0.698039,0.423529}%
\pgfsetfillcolor{currentfill}%
\pgfsetfillopacity{0.120000}%
\pgfsetlinewidth{0.000000pt}%
\definecolor{currentstroke}{rgb}{0.000000,0.000000,0.000000}%
\pgfsetstrokecolor{currentstroke}%
\pgfsetdash{}{0pt}%
\pgfpathmoveto{\pgfqpoint{1.341720in}{1.149504in}}%
\pgfpathlineto{\pgfqpoint{1.357942in}{1.148313in}}%
\pgfpathlineto{\pgfqpoint{1.374215in}{1.147760in}}%
\pgfpathlineto{\pgfqpoint{1.381820in}{1.161912in}}%
\pgfpathlineto{\pgfqpoint{1.365554in}{1.162512in}}%
\pgfpathlineto{\pgfqpoint{1.349339in}{1.163737in}}%
\pgfusepath{fill}%
\end{pgfscope}%
\begin{pgfscope}%
\pgfpathrectangle{\pgfqpoint{0.407592in}{0.222904in}}{\pgfqpoint{1.411163in}{1.411163in}}%
\pgfusepath{clip}%
\pgfsetbuttcap%
\pgfsetroundjoin%
\definecolor{currentfill}{rgb}{0.909804,0.698039,0.423529}%
\pgfsetfillcolor{currentfill}%
\pgfsetfillopacity{0.120000}%
\pgfsetlinewidth{0.000000pt}%
\definecolor{currentstroke}{rgb}{0.000000,0.000000,0.000000}%
\pgfsetstrokecolor{currentstroke}%
\pgfsetdash{}{0pt}%
\pgfpathmoveto{\pgfqpoint{1.183598in}{1.063106in}}%
\pgfpathlineto{\pgfqpoint{1.199607in}{1.057542in}}%
\pgfpathlineto{\pgfqpoint{1.215655in}{1.052757in}}%
\pgfpathlineto{\pgfqpoint{1.223415in}{1.065568in}}%
\pgfpathlineto{\pgfqpoint{1.231159in}{1.078608in}}%
\pgfpathlineto{\pgfqpoint{1.215127in}{1.083308in}}%
\pgfpathlineto{\pgfqpoint{1.199134in}{1.088752in}}%
\pgfpathlineto{\pgfqpoint{1.191374in}{1.075814in}}%
\pgfpathlineto{\pgfqpoint{1.183598in}{1.063106in}}%
\pgfpathclose%
\pgfusepath{fill}%
\end{pgfscope}%
\begin{pgfscope}%
\pgfpathrectangle{\pgfqpoint{0.407592in}{0.222904in}}{\pgfqpoint{1.411163in}{1.411163in}}%
\pgfusepath{clip}%
\pgfsetbuttcap%
\pgfsetroundjoin%
\definecolor{currentfill}{rgb}{0.909804,0.698039,0.423529}%
\pgfsetfillcolor{currentfill}%
\pgfsetfillopacity{0.120000}%
\pgfsetlinewidth{0.000000pt}%
\definecolor{currentstroke}{rgb}{0.000000,0.000000,0.000000}%
\pgfsetstrokecolor{currentstroke}%
\pgfsetdash{}{0pt}%
\pgfpathmoveto{\pgfqpoint{1.136050in}{1.052193in}}%
\pgfpathlineto{\pgfqpoint{1.152007in}{1.044934in}}%
\pgfpathlineto{\pgfqpoint{1.167997in}{1.038452in}}%
\pgfpathlineto{\pgfqpoint{1.175806in}{1.050646in}}%
\pgfpathlineto{\pgfqpoint{1.183598in}{1.063106in}}%
\pgfpathlineto{\pgfqpoint{1.167625in}{1.069434in}}%
\pgfpathlineto{\pgfqpoint{1.151687in}{1.076507in}}%
\pgfpathlineto{\pgfqpoint{1.143877in}{1.064217in}}%
\pgfpathlineto{\pgfqpoint{1.136050in}{1.052193in}}%
\pgfpathclose%
\pgfusepath{fill}%
\end{pgfscope}%
\begin{pgfscope}%
\pgfpathrectangle{\pgfqpoint{0.407592in}{0.222904in}}{\pgfqpoint{1.411163in}{1.411163in}}%
\pgfusepath{clip}%
\pgfsetbuttcap%
\pgfsetroundjoin%
\definecolor{currentfill}{rgb}{0.909804,0.698039,0.423529}%
\pgfsetfillcolor{currentfill}%
\pgfsetfillopacity{0.120000}%
\pgfsetlinewidth{0.000000pt}%
\definecolor{currentstroke}{rgb}{0.000000,0.000000,0.000000}%
\pgfsetstrokecolor{currentstroke}%
\pgfsetdash{}{0pt}%
\pgfpathmoveto{\pgfqpoint{0.961396in}{1.075593in}}%
\pgfpathlineto{\pgfqpoint{0.977238in}{1.062334in}}%
\pgfpathlineto{\pgfqpoint{0.993099in}{1.049603in}}%
\pgfpathlineto{\pgfqpoint{1.001121in}{1.058891in}}%
\pgfpathlineto{\pgfqpoint{1.009121in}{1.068543in}}%
\pgfpathlineto{\pgfqpoint{0.993291in}{1.080998in}}%
\pgfpathlineto{\pgfqpoint{0.977479in}{1.093978in}}%
\pgfpathlineto{\pgfqpoint{0.969449in}{1.084610in}}%
\pgfpathlineto{\pgfqpoint{0.961396in}{1.075593in}}%
\pgfpathclose%
\pgfusepath{fill}%
\end{pgfscope}%
\begin{pgfscope}%
\pgfpathrectangle{\pgfqpoint{0.407592in}{0.222904in}}{\pgfqpoint{1.411163in}{1.411163in}}%
\pgfusepath{clip}%
\pgfsetbuttcap%
\pgfsetroundjoin%
\definecolor{currentfill}{rgb}{0.909804,0.698039,0.423529}%
\pgfsetfillcolor{currentfill}%
\pgfsetfillopacity{0.120000}%
\pgfsetlinewidth{0.000000pt}%
\definecolor{currentstroke}{rgb}{0.000000,0.000000,0.000000}%
\pgfsetstrokecolor{currentstroke}%
\pgfsetdash{}{0pt}%
\pgfpathmoveto{\pgfqpoint{0.881860in}{1.117187in}}%
\pgfpathlineto{\pgfqpoint{0.897678in}{1.101959in}}%
\pgfpathlineto{\pgfqpoint{0.913509in}{1.087106in}}%
\pgfpathlineto{\pgfqpoint{0.921647in}{1.095140in}}%
\pgfpathlineto{\pgfqpoint{0.929758in}{1.103535in}}%
\pgfpathlineto{\pgfqpoint{0.913961in}{1.118143in}}%
\pgfpathlineto{\pgfqpoint{0.898178in}{1.133129in}}%
\pgfpathlineto{\pgfqpoint{0.890032in}{1.124986in}}%
\pgfpathlineto{\pgfqpoint{0.881860in}{1.117187in}}%
\pgfpathclose%
\pgfusepath{fill}%
\end{pgfscope}%
\begin{pgfscope}%
\pgfpathrectangle{\pgfqpoint{0.407592in}{0.222904in}}{\pgfqpoint{1.411163in}{1.411163in}}%
\pgfusepath{clip}%
\pgfsetbuttcap%
\pgfsetroundjoin%
\definecolor{currentfill}{rgb}{0.909804,0.698039,0.423529}%
\pgfsetfillcolor{currentfill}%
\pgfsetfillopacity{0.120000}%
\pgfsetlinewidth{0.000000pt}%
\definecolor{currentstroke}{rgb}{0.000000,0.000000,0.000000}%
\pgfsetstrokecolor{currentstroke}%
\pgfsetdash{}{0pt}%
\pgfpathmoveto{\pgfqpoint{1.231159in}{1.078608in}}%
\pgfpathlineto{\pgfqpoint{1.247233in}{1.074658in}}%
\pgfpathlineto{\pgfqpoint{1.263350in}{1.071463in}}%
\pgfpathlineto{\pgfqpoint{1.271065in}{1.084747in}}%
\pgfpathlineto{\pgfqpoint{1.278765in}{1.098228in}}%
\pgfpathlineto{\pgfqpoint{1.262662in}{1.101408in}}%
\pgfpathlineto{\pgfqpoint{1.246603in}{1.105308in}}%
\pgfpathlineto{\pgfqpoint{1.238888in}{1.091860in}}%
\pgfpathlineto{\pgfqpoint{1.231159in}{1.078608in}}%
\pgfpathclose%
\pgfusepath{fill}%
\end{pgfscope}%
\begin{pgfscope}%
\pgfpathrectangle{\pgfqpoint{0.407592in}{0.222904in}}{\pgfqpoint{1.411163in}{1.411163in}}%
\pgfusepath{clip}%
\pgfsetbuttcap%
\pgfsetroundjoin%
\definecolor{currentfill}{rgb}{0.909804,0.698039,0.423529}%
\pgfsetfillcolor{currentfill}%
\pgfsetfillopacity{0.120000}%
\pgfsetlinewidth{0.000000pt}%
\definecolor{currentstroke}{rgb}{0.000000,0.000000,0.000000}%
\pgfsetstrokecolor{currentstroke}%
\pgfsetdash{}{0pt}%
\pgfpathmoveto{\pgfqpoint{1.088477in}{1.046225in}}%
\pgfpathlineto{\pgfqpoint{1.104395in}{1.037252in}}%
\pgfpathlineto{\pgfqpoint{1.120342in}{1.029027in}}%
\pgfpathlineto{\pgfqpoint{1.128205in}{1.040457in}}%
\pgfpathlineto{\pgfqpoint{1.136050in}{1.052193in}}%
\pgfpathlineto{\pgfqpoint{1.120124in}{1.060205in}}%
\pgfpathlineto{\pgfqpoint{1.104228in}{1.068940in}}%
\pgfpathlineto{\pgfqpoint{1.096362in}{1.057430in}}%
\pgfpathlineto{\pgfqpoint{1.088477in}{1.046225in}}%
\pgfpathclose%
\pgfusepath{fill}%
\end{pgfscope}%
\begin{pgfscope}%
\pgfpathrectangle{\pgfqpoint{0.407592in}{0.222904in}}{\pgfqpoint{1.411163in}{1.411163in}}%
\pgfusepath{clip}%
\pgfsetbuttcap%
\pgfsetroundjoin%
\definecolor{currentfill}{rgb}{0.909804,0.698039,0.423529}%
\pgfsetfillcolor{currentfill}%
\pgfsetfillopacity{0.120000}%
\pgfsetlinewidth{0.000000pt}%
\definecolor{currentstroke}{rgb}{0.000000,0.000000,0.000000}%
\pgfsetstrokecolor{currentstroke}%
\pgfsetdash{}{0pt}%
\pgfpathmoveto{\pgfqpoint{1.278765in}{1.098228in}}%
\pgfpathlineto{\pgfqpoint{1.294915in}{1.095767in}}%
\pgfpathlineto{\pgfqpoint{1.311112in}{1.094022in}}%
\pgfpathlineto{\pgfqpoint{1.318785in}{1.107651in}}%
\pgfpathlineto{\pgfqpoint{1.326443in}{1.121449in}}%
\pgfpathlineto{\pgfqpoint{1.310259in}{1.123243in}}%
\pgfpathlineto{\pgfqpoint{1.294124in}{1.125720in}}%
\pgfpathlineto{\pgfqpoint{1.286451in}{1.111890in}}%
\pgfpathlineto{\pgfqpoint{1.278765in}{1.098228in}}%
\pgfpathclose%
\pgfusepath{fill}%
\end{pgfscope}%
\begin{pgfscope}%
\pgfpathrectangle{\pgfqpoint{0.407592in}{0.222904in}}{\pgfqpoint{1.411163in}{1.411163in}}%
\pgfusepath{clip}%
\pgfsetbuttcap%
\pgfsetroundjoin%
\definecolor{currentfill}{rgb}{0.909804,0.698039,0.423529}%
\pgfsetfillcolor{currentfill}%
\pgfsetfillopacity{0.120000}%
\pgfsetlinewidth{0.000000pt}%
\definecolor{currentstroke}{rgb}{0.000000,0.000000,0.000000}%
\pgfsetstrokecolor{currentstroke}%
\pgfsetdash{}{0pt}%
\pgfpathmoveto{\pgfqpoint{0.802133in}{1.167492in}}%
\pgfpathlineto{\pgfqpoint{0.817941in}{1.150873in}}%
\pgfpathlineto{\pgfqpoint{0.833760in}{1.134511in}}%
\pgfpathlineto{\pgfqpoint{0.842025in}{1.141404in}}%
\pgfpathlineto{\pgfqpoint{0.850261in}{1.148640in}}%
\pgfpathlineto{\pgfqpoint{0.834480in}{1.164810in}}%
\pgfpathlineto{\pgfqpoint{0.818710in}{1.181241in}}%
\pgfpathlineto{\pgfqpoint{0.810436in}{1.174204in}}%
\pgfpathlineto{\pgfqpoint{0.802133in}{1.167492in}}%
\pgfpathclose%
\pgfusepath{fill}%
\end{pgfscope}%
\begin{pgfscope}%
\pgfpathrectangle{\pgfqpoint{0.407592in}{0.222904in}}{\pgfqpoint{1.411163in}{1.411163in}}%
\pgfusepath{clip}%
\pgfsetbuttcap%
\pgfsetroundjoin%
\definecolor{currentfill}{rgb}{0.909804,0.698039,0.423529}%
\pgfsetfillcolor{currentfill}%
\pgfsetfillopacity{0.120000}%
\pgfsetlinewidth{0.000000pt}%
\definecolor{currentstroke}{rgb}{0.000000,0.000000,0.000000}%
\pgfsetstrokecolor{currentstroke}%
\pgfsetdash{}{0pt}%
\pgfpathmoveto{\pgfqpoint{1.040840in}{1.045374in}}%
\pgfpathlineto{\pgfqpoint{1.056732in}{1.034745in}}%
\pgfpathlineto{\pgfqpoint{1.072648in}{1.024807in}}%
\pgfpathlineto{\pgfqpoint{1.080573in}{1.035344in}}%
\pgfpathlineto{\pgfqpoint{1.088477in}{1.046225in}}%
\pgfpathlineto{\pgfqpoint{1.072586in}{1.055911in}}%
\pgfpathlineto{\pgfqpoint{1.056719in}{1.066273in}}%
\pgfpathlineto{\pgfqpoint{1.048790in}{1.055655in}}%
\pgfpathlineto{\pgfqpoint{1.040840in}{1.045374in}}%
\pgfpathclose%
\pgfusepath{fill}%
\end{pgfscope}%
\begin{pgfscope}%
\pgfpathrectangle{\pgfqpoint{0.407592in}{0.222904in}}{\pgfqpoint{1.411163in}{1.411163in}}%
\pgfusepath{clip}%
\pgfsetbuttcap%
\pgfsetroundjoin%
\definecolor{currentfill}{rgb}{0.909804,0.698039,0.423529}%
\pgfsetfillcolor{currentfill}%
\pgfsetfillopacity{0.120000}%
\pgfsetlinewidth{0.000000pt}%
\definecolor{currentstroke}{rgb}{0.000000,0.000000,0.000000}%
\pgfsetstrokecolor{currentstroke}%
\pgfsetdash{}{0pt}%
\pgfpathmoveto{\pgfqpoint{0.690568in}{1.259886in}}%
\pgfpathlineto{\pgfqpoint{0.706355in}{1.242001in}}%
\pgfpathlineto{\pgfqpoint{0.722152in}{1.224261in}}%
\pgfpathlineto{\pgfqpoint{0.730594in}{1.230039in}}%
\pgfpathlineto{\pgfqpoint{0.739005in}{1.236117in}}%
\pgfpathlineto{\pgfqpoint{0.723249in}{1.253721in}}%
\pgfpathlineto{\pgfqpoint{0.707502in}{1.271476in}}%
\pgfpathlineto{\pgfqpoint{0.699051in}{1.265539in}}%
\pgfpathlineto{\pgfqpoint{0.690568in}{1.259886in}}%
\pgfpathclose%
\pgfusepath{fill}%
\end{pgfscope}%
\begin{pgfscope}%
\pgfpathrectangle{\pgfqpoint{0.407592in}{0.222904in}}{\pgfqpoint{1.411163in}{1.411163in}}%
\pgfusepath{clip}%
\pgfsetbuttcap%
\pgfsetroundjoin%
\definecolor{currentfill}{rgb}{0.909804,0.698039,0.423529}%
\pgfsetfillcolor{currentfill}%
\pgfsetfillopacity{0.120000}%
\pgfsetlinewidth{0.000000pt}%
\definecolor{currentstroke}{rgb}{0.000000,0.000000,0.000000}%
\pgfsetstrokecolor{currentstroke}%
\pgfsetdash{}{0pt}%
\pgfpathmoveto{\pgfqpoint{1.326443in}{1.121449in}}%
\pgfpathlineto{\pgfqpoint{1.342677in}{1.120331in}}%
\pgfpathlineto{\pgfqpoint{1.358964in}{1.119878in}}%
\pgfpathlineto{\pgfqpoint{1.366596in}{1.133745in}}%
\pgfpathlineto{\pgfqpoint{1.374215in}{1.147760in}}%
\pgfpathlineto{\pgfqpoint{1.357942in}{1.148313in}}%
\pgfpathlineto{\pgfqpoint{1.341720in}{1.149504in}}%
\pgfpathlineto{\pgfqpoint{1.334089in}{1.135404in}}%
\pgfpathlineto{\pgfqpoint{1.326443in}{1.121449in}}%
\pgfpathclose%
\pgfusepath{fill}%
\end{pgfscope}%
\begin{pgfscope}%
\pgfpathrectangle{\pgfqpoint{0.407592in}{0.222904in}}{\pgfqpoint{1.411163in}{1.411163in}}%
\pgfusepath{clip}%
\pgfsetbuttcap%
\pgfsetroundjoin%
\definecolor{currentfill}{rgb}{0.909804,0.698039,0.423529}%
\pgfsetfillcolor{currentfill}%
\pgfsetfillopacity{0.120000}%
\pgfsetlinewidth{0.000000pt}%
\definecolor{currentstroke}{rgb}{0.000000,0.000000,0.000000}%
\pgfsetstrokecolor{currentstroke}%
\pgfsetdash{}{0pt}%
\pgfpathmoveto{\pgfqpoint{0.913509in}{1.087106in}}%
\pgfpathlineto{\pgfqpoint{0.929355in}{1.072663in}}%
\pgfpathlineto{\pgfqpoint{0.945215in}{1.058668in}}%
\pgfpathlineto{\pgfqpoint{0.953318in}{1.066941in}}%
\pgfpathlineto{\pgfqpoint{0.961396in}{1.075593in}}%
\pgfpathlineto{\pgfqpoint{0.945569in}{1.089340in}}%
\pgfpathlineto{\pgfqpoint{0.929758in}{1.103535in}}%
\pgfpathlineto{\pgfqpoint{0.921647in}{1.095140in}}%
\pgfpathlineto{\pgfqpoint{0.913509in}{1.087106in}}%
\pgfpathclose%
\pgfusepath{fill}%
\end{pgfscope}%
\begin{pgfscope}%
\pgfpathrectangle{\pgfqpoint{0.407592in}{0.222904in}}{\pgfqpoint{1.411163in}{1.411163in}}%
\pgfusepath{clip}%
\pgfsetbuttcap%
\pgfsetroundjoin%
\definecolor{currentfill}{rgb}{0.909804,0.698039,0.423529}%
\pgfsetfillcolor{currentfill}%
\pgfsetfillopacity{0.120000}%
\pgfsetlinewidth{0.000000pt}%
\definecolor{currentstroke}{rgb}{0.000000,0.000000,0.000000}%
\pgfsetstrokecolor{currentstroke}%
\pgfsetdash{}{0pt}%
\pgfpathmoveto{\pgfqpoint{0.993099in}{1.049603in}}%
\pgfpathlineto{\pgfqpoint{1.008977in}{1.037445in}}%
\pgfpathlineto{\pgfqpoint{1.024876in}{1.025903in}}%
\pgfpathlineto{\pgfqpoint{1.032869in}{1.035451in}}%
\pgfpathlineto{\pgfqpoint{1.040840in}{1.045374in}}%
\pgfpathlineto{\pgfqpoint{1.024970in}{1.056655in}}%
\pgfpathlineto{\pgfqpoint{1.009121in}{1.068543in}}%
\pgfpathlineto{\pgfqpoint{1.001121in}{1.058891in}}%
\pgfpathlineto{\pgfqpoint{0.993099in}{1.049603in}}%
\pgfpathclose%
\pgfusepath{fill}%
\end{pgfscope}%
\begin{pgfscope}%
\pgfpathrectangle{\pgfqpoint{0.407592in}{0.222904in}}{\pgfqpoint{1.411163in}{1.411163in}}%
\pgfusepath{clip}%
\pgfsetbuttcap%
\pgfsetroundjoin%
\definecolor{currentfill}{rgb}{0.909804,0.698039,0.423529}%
\pgfsetfillcolor{currentfill}%
\pgfsetfillopacity{0.120000}%
\pgfsetlinewidth{0.000000pt}%
\definecolor{currentstroke}{rgb}{0.000000,0.000000,0.000000}%
\pgfsetstrokecolor{currentstroke}%
\pgfsetdash{}{0pt}%
\pgfpathmoveto{\pgfqpoint{0.722152in}{1.224261in}}%
\pgfpathlineto{\pgfqpoint{0.737958in}{1.206682in}}%
\pgfpathlineto{\pgfqpoint{0.753774in}{1.189279in}}%
\pgfpathlineto{\pgfqpoint{0.762177in}{1.195187in}}%
\pgfpathlineto{\pgfqpoint{0.770549in}{1.201413in}}%
\pgfpathlineto{\pgfqpoint{0.754772in}{1.218675in}}%
\pgfpathlineto{\pgfqpoint{0.739005in}{1.236117in}}%
\pgfpathlineto{\pgfqpoint{0.730594in}{1.230039in}}%
\pgfpathlineto{\pgfqpoint{0.722152in}{1.224261in}}%
\pgfpathclose%
\pgfusepath{fill}%
\end{pgfscope}%
\begin{pgfscope}%
\pgfpathrectangle{\pgfqpoint{0.407592in}{0.222904in}}{\pgfqpoint{1.411163in}{1.411163in}}%
\pgfusepath{clip}%
\pgfsetbuttcap%
\pgfsetroundjoin%
\definecolor{currentfill}{rgb}{0.909804,0.698039,0.423529}%
\pgfsetfillcolor{currentfill}%
\pgfsetfillopacity{0.120000}%
\pgfsetlinewidth{0.000000pt}%
\definecolor{currentstroke}{rgb}{0.000000,0.000000,0.000000}%
\pgfsetstrokecolor{currentstroke}%
\pgfsetdash{}{0pt}%
\pgfpathmoveto{\pgfqpoint{0.833760in}{1.134511in}}%
\pgfpathlineto{\pgfqpoint{0.849591in}{1.118431in}}%
\pgfpathlineto{\pgfqpoint{0.865433in}{1.102660in}}%
\pgfpathlineto{\pgfqpoint{0.873660in}{1.109742in}}%
\pgfpathlineto{\pgfqpoint{0.881860in}{1.117187in}}%
\pgfpathlineto{\pgfqpoint{0.866055in}{1.132757in}}%
\pgfpathlineto{\pgfqpoint{0.850261in}{1.148640in}}%
\pgfpathlineto{\pgfqpoint{0.842025in}{1.141404in}}%
\pgfpathlineto{\pgfqpoint{0.833760in}{1.134511in}}%
\pgfpathclose%
\pgfusepath{fill}%
\end{pgfscope}%
\begin{pgfscope}%
\pgfpathrectangle{\pgfqpoint{0.407592in}{0.222904in}}{\pgfqpoint{1.411163in}{1.411163in}}%
\pgfusepath{clip}%
\pgfsetbuttcap%
\pgfsetroundjoin%
\definecolor{currentfill}{rgb}{0.909804,0.698039,0.423529}%
\pgfsetfillcolor{currentfill}%
\pgfsetfillopacity{0.120000}%
\pgfsetlinewidth{0.000000pt}%
\definecolor{currentstroke}{rgb}{0.000000,0.000000,0.000000}%
\pgfsetstrokecolor{currentstroke}%
\pgfsetdash{}{0pt}%
\pgfpathmoveto{\pgfqpoint{1.167997in}{1.038452in}}%
\pgfpathlineto{\pgfqpoint{1.184024in}{1.032768in}}%
\pgfpathlineto{\pgfqpoint{1.200090in}{1.027897in}}%
\pgfpathlineto{\pgfqpoint{1.207880in}{1.040194in}}%
\pgfpathlineto{\pgfqpoint{1.215655in}{1.052757in}}%
\pgfpathlineto{\pgfqpoint{1.199607in}{1.057542in}}%
\pgfpathlineto{\pgfqpoint{1.183598in}{1.063106in}}%
\pgfpathlineto{\pgfqpoint{1.175806in}{1.050646in}}%
\pgfpathlineto{\pgfqpoint{1.167997in}{1.038452in}}%
\pgfpathclose%
\pgfusepath{fill}%
\end{pgfscope}%
\begin{pgfscope}%
\pgfpathrectangle{\pgfqpoint{0.407592in}{0.222904in}}{\pgfqpoint{1.411163in}{1.411163in}}%
\pgfusepath{clip}%
\pgfsetbuttcap%
\pgfsetroundjoin%
\definecolor{currentfill}{rgb}{0.909804,0.698039,0.423529}%
\pgfsetfillcolor{currentfill}%
\pgfsetfillopacity{0.120000}%
\pgfsetlinewidth{0.000000pt}%
\definecolor{currentstroke}{rgb}{0.000000,0.000000,0.000000}%
\pgfsetstrokecolor{currentstroke}%
\pgfsetdash{}{0pt}%
\pgfpathmoveto{\pgfqpoint{1.120342in}{1.029027in}}%
\pgfpathlineto{\pgfqpoint{1.136319in}{1.021583in}}%
\pgfpathlineto{\pgfqpoint{1.152330in}{1.014947in}}%
\pgfpathlineto{\pgfqpoint{1.160172in}{1.026546in}}%
\pgfpathlineto{\pgfqpoint{1.167997in}{1.038452in}}%
\pgfpathlineto{\pgfqpoint{1.152007in}{1.044934in}}%
\pgfpathlineto{\pgfqpoint{1.136050in}{1.052193in}}%
\pgfpathlineto{\pgfqpoint{1.128205in}{1.040457in}}%
\pgfpathlineto{\pgfqpoint{1.120342in}{1.029027in}}%
\pgfpathclose%
\pgfusepath{fill}%
\end{pgfscope}%
\begin{pgfscope}%
\pgfpathrectangle{\pgfqpoint{0.407592in}{0.222904in}}{\pgfqpoint{1.411163in}{1.411163in}}%
\pgfusepath{clip}%
\pgfsetbuttcap%
\pgfsetroundjoin%
\definecolor{currentfill}{rgb}{0.909804,0.698039,0.423529}%
\pgfsetfillcolor{currentfill}%
\pgfsetfillopacity{0.120000}%
\pgfsetlinewidth{0.000000pt}%
\definecolor{currentstroke}{rgb}{0.000000,0.000000,0.000000}%
\pgfsetstrokecolor{currentstroke}%
\pgfsetdash{}{0pt}%
\pgfpathmoveto{\pgfqpoint{1.215655in}{1.052757in}}%
\pgfpathlineto{\pgfqpoint{1.231745in}{1.048758in}}%
\pgfpathlineto{\pgfqpoint{1.247878in}{1.045550in}}%
\pgfpathlineto{\pgfqpoint{1.255621in}{1.058392in}}%
\pgfpathlineto{\pgfqpoint{1.263350in}{1.071463in}}%
\pgfpathlineto{\pgfqpoint{1.247233in}{1.074658in}}%
\pgfpathlineto{\pgfqpoint{1.231159in}{1.078608in}}%
\pgfpathlineto{\pgfqpoint{1.223415in}{1.065568in}}%
\pgfpathlineto{\pgfqpoint{1.215655in}{1.052757in}}%
\pgfpathclose%
\pgfusepath{fill}%
\end{pgfscope}%
\begin{pgfscope}%
\pgfpathrectangle{\pgfqpoint{0.407592in}{0.222904in}}{\pgfqpoint{1.411163in}{1.411163in}}%
\pgfusepath{clip}%
\pgfsetbuttcap%
\pgfsetroundjoin%
\definecolor{currentfill}{rgb}{0.909804,0.698039,0.423529}%
\pgfsetfillcolor{currentfill}%
\pgfsetfillopacity{0.120000}%
\pgfsetlinewidth{0.000000pt}%
\definecolor{currentstroke}{rgb}{0.000000,0.000000,0.000000}%
\pgfsetstrokecolor{currentstroke}%
\pgfsetdash{}{0pt}%
\pgfpathmoveto{\pgfqpoint{1.072648in}{1.024807in}}%
\pgfpathlineto{\pgfqpoint{1.088590in}{1.015603in}}%
\pgfpathlineto{\pgfqpoint{1.104560in}{1.007171in}}%
\pgfpathlineto{\pgfqpoint{1.112460in}{1.017925in}}%
\pgfpathlineto{\pgfqpoint{1.120342in}{1.029027in}}%
\pgfpathlineto{\pgfqpoint{1.104395in}{1.037252in}}%
\pgfpathlineto{\pgfqpoint{1.088477in}{1.046225in}}%
\pgfpathlineto{\pgfqpoint{1.080573in}{1.035344in}}%
\pgfpathlineto{\pgfqpoint{1.072648in}{1.024807in}}%
\pgfpathclose%
\pgfusepath{fill}%
\end{pgfscope}%
\begin{pgfscope}%
\pgfpathrectangle{\pgfqpoint{0.407592in}{0.222904in}}{\pgfqpoint{1.411163in}{1.411163in}}%
\pgfusepath{clip}%
\pgfsetbuttcap%
\pgfsetroundjoin%
\definecolor{currentfill}{rgb}{0.909804,0.698039,0.423529}%
\pgfsetfillcolor{currentfill}%
\pgfsetfillopacity{0.120000}%
\pgfsetlinewidth{0.000000pt}%
\definecolor{currentstroke}{rgb}{0.000000,0.000000,0.000000}%
\pgfsetstrokecolor{currentstroke}%
\pgfsetdash{}{0pt}%
\pgfpathmoveto{\pgfqpoint{0.753774in}{1.189279in}}%
\pgfpathlineto{\pgfqpoint{0.769601in}{1.172068in}}%
\pgfpathlineto{\pgfqpoint{0.785437in}{1.155067in}}%
\pgfpathlineto{\pgfqpoint{0.793800in}{1.161112in}}%
\pgfpathlineto{\pgfqpoint{0.802133in}{1.167492in}}%
\pgfpathlineto{\pgfqpoint{0.786336in}{1.184345in}}%
\pgfpathlineto{\pgfqpoint{0.770549in}{1.201413in}}%
\pgfpathlineto{\pgfqpoint{0.762177in}{1.195187in}}%
\pgfpathlineto{\pgfqpoint{0.753774in}{1.189279in}}%
\pgfpathclose%
\pgfusepath{fill}%
\end{pgfscope}%
\begin{pgfscope}%
\pgfpathrectangle{\pgfqpoint{0.407592in}{0.222904in}}{\pgfqpoint{1.411163in}{1.411163in}}%
\pgfusepath{clip}%
\pgfsetbuttcap%
\pgfsetroundjoin%
\definecolor{currentfill}{rgb}{0.909804,0.698039,0.423529}%
\pgfsetfillcolor{currentfill}%
\pgfsetfillopacity{0.120000}%
\pgfsetlinewidth{0.000000pt}%
\definecolor{currentstroke}{rgb}{0.000000,0.000000,0.000000}%
\pgfsetstrokecolor{currentstroke}%
\pgfsetdash{}{0pt}%
\pgfpathmoveto{\pgfqpoint{1.374215in}{1.147760in}}%
\pgfpathlineto{\pgfqpoint{1.390541in}{1.147834in}}%
\pgfpathlineto{\pgfqpoint{1.406923in}{1.148521in}}%
\pgfpathlineto{\pgfqpoint{1.414516in}{1.162542in}}%
\pgfpathlineto{\pgfqpoint{1.398141in}{1.161926in}}%
\pgfpathlineto{\pgfqpoint{1.381820in}{1.161912in}}%
\pgfusepath{fill}%
\end{pgfscope}%
\begin{pgfscope}%
\pgfpathrectangle{\pgfqpoint{0.407592in}{0.222904in}}{\pgfqpoint{1.411163in}{1.411163in}}%
\pgfusepath{clip}%
\pgfsetbuttcap%
\pgfsetroundjoin%
\definecolor{currentfill}{rgb}{0.909804,0.698039,0.423529}%
\pgfsetfillcolor{currentfill}%
\pgfsetfillopacity{0.120000}%
\pgfsetlinewidth{0.000000pt}%
\definecolor{currentstroke}{rgb}{0.000000,0.000000,0.000000}%
\pgfsetstrokecolor{currentstroke}%
\pgfsetdash{}{0pt}%
\pgfpathmoveto{\pgfqpoint{0.945215in}{1.058668in}}%
\pgfpathlineto{\pgfqpoint{0.961090in}{1.045160in}}%
\pgfpathlineto{\pgfqpoint{0.976982in}{1.032183in}}%
\pgfpathlineto{\pgfqpoint{0.985052in}{1.040695in}}%
\pgfpathlineto{\pgfqpoint{0.993099in}{1.049603in}}%
\pgfpathlineto{\pgfqpoint{0.977238in}{1.062334in}}%
\pgfpathlineto{\pgfqpoint{0.961396in}{1.075593in}}%
\pgfpathlineto{\pgfqpoint{0.953318in}{1.066941in}}%
\pgfpathlineto{\pgfqpoint{0.945215in}{1.058668in}}%
\pgfpathclose%
\pgfusepath{fill}%
\end{pgfscope}%
\begin{pgfscope}%
\pgfpathrectangle{\pgfqpoint{0.407592in}{0.222904in}}{\pgfqpoint{1.411163in}{1.411163in}}%
\pgfusepath{clip}%
\pgfsetbuttcap%
\pgfsetroundjoin%
\definecolor{currentfill}{rgb}{0.909804,0.698039,0.423529}%
\pgfsetfillcolor{currentfill}%
\pgfsetfillopacity{0.120000}%
\pgfsetlinewidth{0.000000pt}%
\definecolor{currentstroke}{rgb}{0.000000,0.000000,0.000000}%
\pgfsetstrokecolor{currentstroke}%
\pgfsetdash{}{0pt}%
\pgfpathmoveto{\pgfqpoint{1.263350in}{1.071463in}}%
\pgfpathlineto{\pgfqpoint{1.279514in}{1.069022in}}%
\pgfpathlineto{\pgfqpoint{1.295726in}{1.067329in}}%
\pgfpathlineto{\pgfqpoint{1.303426in}{1.080576in}}%
\pgfpathlineto{\pgfqpoint{1.311112in}{1.094022in}}%
\pgfpathlineto{\pgfqpoint{1.294915in}{1.095767in}}%
\pgfpathlineto{\pgfqpoint{1.278765in}{1.098228in}}%
\pgfpathlineto{\pgfqpoint{1.271065in}{1.084747in}}%
\pgfpathlineto{\pgfqpoint{1.263350in}{1.071463in}}%
\pgfpathclose%
\pgfusepath{fill}%
\end{pgfscope}%
\begin{pgfscope}%
\pgfpathrectangle{\pgfqpoint{0.407592in}{0.222904in}}{\pgfqpoint{1.411163in}{1.411163in}}%
\pgfusepath{clip}%
\pgfsetbuttcap%
\pgfsetroundjoin%
\definecolor{currentfill}{rgb}{0.909804,0.698039,0.423529}%
\pgfsetfillcolor{currentfill}%
\pgfsetfillopacity{0.120000}%
\pgfsetlinewidth{0.000000pt}%
\definecolor{currentstroke}{rgb}{0.000000,0.000000,0.000000}%
\pgfsetstrokecolor{currentstroke}%
\pgfsetdash{}{0pt}%
\pgfpathmoveto{\pgfqpoint{0.865433in}{1.102660in}}%
\pgfpathlineto{\pgfqpoint{0.881287in}{1.087228in}}%
\pgfpathlineto{\pgfqpoint{0.897155in}{1.072168in}}%
\pgfpathlineto{\pgfqpoint{0.905346in}{1.079446in}}%
\pgfpathlineto{\pgfqpoint{0.913509in}{1.087106in}}%
\pgfpathlineto{\pgfqpoint{0.897678in}{1.101959in}}%
\pgfpathlineto{\pgfqpoint{0.881860in}{1.117187in}}%
\pgfpathlineto{\pgfqpoint{0.873660in}{1.109742in}}%
\pgfpathlineto{\pgfqpoint{0.865433in}{1.102660in}}%
\pgfpathclose%
\pgfusepath{fill}%
\end{pgfscope}%
\begin{pgfscope}%
\pgfpathrectangle{\pgfqpoint{0.407592in}{0.222904in}}{\pgfqpoint{1.411163in}{1.411163in}}%
\pgfusepath{clip}%
\pgfsetbuttcap%
\pgfsetroundjoin%
\definecolor{currentfill}{rgb}{0.909804,0.698039,0.423529}%
\pgfsetfillcolor{currentfill}%
\pgfsetfillopacity{0.120000}%
\pgfsetlinewidth{0.000000pt}%
\definecolor{currentstroke}{rgb}{0.000000,0.000000,0.000000}%
\pgfsetstrokecolor{currentstroke}%
\pgfsetdash{}{0pt}%
\pgfpathmoveto{\pgfqpoint{1.024876in}{1.025903in}}%
\pgfpathlineto{\pgfqpoint{1.040795in}{1.015022in}}%
\pgfpathlineto{\pgfqpoint{1.056738in}{1.004849in}}%
\pgfpathlineto{\pgfqpoint{1.064704in}{1.014636in}}%
\pgfpathlineto{\pgfqpoint{1.072648in}{1.024807in}}%
\pgfpathlineto{\pgfqpoint{1.056732in}{1.034745in}}%
\pgfpathlineto{\pgfqpoint{1.040840in}{1.045374in}}%
\pgfpathlineto{\pgfqpoint{1.032869in}{1.035451in}}%
\pgfpathlineto{\pgfqpoint{1.024876in}{1.025903in}}%
\pgfpathclose%
\pgfusepath{fill}%
\end{pgfscope}%
\begin{pgfscope}%
\pgfpathrectangle{\pgfqpoint{0.407592in}{0.222904in}}{\pgfqpoint{1.411163in}{1.411163in}}%
\pgfusepath{clip}%
\pgfsetbuttcap%
\pgfsetroundjoin%
\definecolor{currentfill}{rgb}{0.909804,0.698039,0.423529}%
\pgfsetfillcolor{currentfill}%
\pgfsetfillopacity{0.120000}%
\pgfsetlinewidth{0.000000pt}%
\definecolor{currentstroke}{rgb}{0.000000,0.000000,0.000000}%
\pgfsetstrokecolor{currentstroke}%
\pgfsetdash{}{0pt}%
\pgfpathmoveto{\pgfqpoint{1.311112in}{1.094022in}}%
\pgfpathlineto{\pgfqpoint{1.327360in}{1.092982in}}%
\pgfpathlineto{\pgfqpoint{1.343660in}{1.092638in}}%
\pgfpathlineto{\pgfqpoint{1.351319in}{1.106172in}}%
\pgfpathlineto{\pgfqpoint{1.358964in}{1.119878in}}%
\pgfpathlineto{\pgfqpoint{1.342677in}{1.120331in}}%
\pgfpathlineto{\pgfqpoint{1.326443in}{1.121449in}}%
\pgfpathlineto{\pgfqpoint{1.318785in}{1.107651in}}%
\pgfpathlineto{\pgfqpoint{1.311112in}{1.094022in}}%
\pgfpathclose%
\pgfusepath{fill}%
\end{pgfscope}%
\begin{pgfscope}%
\pgfpathrectangle{\pgfqpoint{0.407592in}{0.222904in}}{\pgfqpoint{1.411163in}{1.411163in}}%
\pgfusepath{clip}%
\pgfsetbuttcap%
\pgfsetroundjoin%
\definecolor{currentfill}{rgb}{0.909804,0.698039,0.423529}%
\pgfsetfillcolor{currentfill}%
\pgfsetfillopacity{0.120000}%
\pgfsetlinewidth{0.000000pt}%
\definecolor{currentstroke}{rgb}{0.000000,0.000000,0.000000}%
\pgfsetstrokecolor{currentstroke}%
\pgfsetdash{}{0pt}%
\pgfpathmoveto{\pgfqpoint{0.785437in}{1.155067in}}%
\pgfpathlineto{\pgfqpoint{0.801284in}{1.138298in}}%
\pgfpathlineto{\pgfqpoint{0.817142in}{1.121783in}}%
\pgfpathlineto{\pgfqpoint{0.825466in}{1.127969in}}%
\pgfpathlineto{\pgfqpoint{0.833760in}{1.134511in}}%
\pgfpathlineto{\pgfqpoint{0.817941in}{1.150873in}}%
\pgfpathlineto{\pgfqpoint{0.802133in}{1.167492in}}%
\pgfpathlineto{\pgfqpoint{0.793800in}{1.161112in}}%
\pgfpathlineto{\pgfqpoint{0.785437in}{1.155067in}}%
\pgfpathclose%
\pgfusepath{fill}%
\end{pgfscope}%
\begin{pgfscope}%
\pgfpathrectangle{\pgfqpoint{0.407592in}{0.222904in}}{\pgfqpoint{1.411163in}{1.411163in}}%
\pgfusepath{clip}%
\pgfsetbuttcap%
\pgfsetroundjoin%
\definecolor{currentfill}{rgb}{0.909804,0.698039,0.423529}%
\pgfsetfillcolor{currentfill}%
\pgfsetfillopacity{0.120000}%
\pgfsetlinewidth{0.000000pt}%
\definecolor{currentstroke}{rgb}{0.000000,0.000000,0.000000}%
\pgfsetstrokecolor{currentstroke}%
\pgfsetdash{}{0pt}%
\pgfpathmoveto{\pgfqpoint{0.673507in}{1.249448in}}%
\pgfpathlineto{\pgfqpoint{0.689334in}{1.231462in}}%
\pgfpathlineto{\pgfqpoint{0.705171in}{1.213622in}}%
\pgfpathlineto{\pgfqpoint{0.713678in}{1.218788in}}%
\pgfpathlineto{\pgfqpoint{0.722152in}{1.224261in}}%
\pgfpathlineto{\pgfqpoint{0.706355in}{1.242001in}}%
\pgfpathlineto{\pgfqpoint{0.690568in}{1.259886in}}%
\pgfpathlineto{\pgfqpoint{0.682054in}{1.254522in}}%
\pgfpathlineto{\pgfqpoint{0.673507in}{1.249448in}}%
\pgfpathclose%
\pgfusepath{fill}%
\end{pgfscope}%
\begin{pgfscope}%
\pgfpathrectangle{\pgfqpoint{0.407592in}{0.222904in}}{\pgfqpoint{1.411163in}{1.411163in}}%
\pgfusepath{clip}%
\pgfsetbuttcap%
\pgfsetroundjoin%
\definecolor{currentfill}{rgb}{0.909804,0.698039,0.423529}%
\pgfsetfillcolor{currentfill}%
\pgfsetfillopacity{0.120000}%
\pgfsetlinewidth{0.000000pt}%
\definecolor{currentstroke}{rgb}{0.000000,0.000000,0.000000}%
\pgfsetstrokecolor{currentstroke}%
\pgfsetdash{}{0pt}%
\pgfpathmoveto{\pgfqpoint{0.897155in}{1.072168in}}%
\pgfpathlineto{\pgfqpoint{0.913035in}{1.057515in}}%
\pgfpathlineto{\pgfqpoint{0.928930in}{1.043308in}}%
\pgfpathlineto{\pgfqpoint{0.937085in}{1.050786in}}%
\pgfpathlineto{\pgfqpoint{0.945215in}{1.058668in}}%
\pgfpathlineto{\pgfqpoint{0.929355in}{1.072663in}}%
\pgfpathlineto{\pgfqpoint{0.913509in}{1.087106in}}%
\pgfpathlineto{\pgfqpoint{0.905346in}{1.079446in}}%
\pgfpathlineto{\pgfqpoint{0.897155in}{1.072168in}}%
\pgfpathclose%
\pgfusepath{fill}%
\end{pgfscope}%
\begin{pgfscope}%
\pgfpathrectangle{\pgfqpoint{0.407592in}{0.222904in}}{\pgfqpoint{1.411163in}{1.411163in}}%
\pgfusepath{clip}%
\pgfsetbuttcap%
\pgfsetroundjoin%
\definecolor{currentfill}{rgb}{0.909804,0.698039,0.423529}%
\pgfsetfillcolor{currentfill}%
\pgfsetfillopacity{0.120000}%
\pgfsetlinewidth{0.000000pt}%
\definecolor{currentstroke}{rgb}{0.000000,0.000000,0.000000}%
\pgfsetstrokecolor{currentstroke}%
\pgfsetdash{}{0pt}%
\pgfpathmoveto{\pgfqpoint{1.358964in}{1.119878in}}%
\pgfpathlineto{\pgfqpoint{1.375304in}{1.120079in}}%
\pgfpathlineto{\pgfqpoint{1.391700in}{1.120917in}}%
\pgfpathlineto{\pgfqpoint{1.399318in}{1.134643in}}%
\pgfpathlineto{\pgfqpoint{1.406923in}{1.148521in}}%
\pgfpathlineto{\pgfqpoint{1.390541in}{1.147834in}}%
\pgfpathlineto{\pgfqpoint{1.374215in}{1.147760in}}%
\pgfpathlineto{\pgfqpoint{1.366596in}{1.133745in}}%
\pgfpathlineto{\pgfqpoint{1.358964in}{1.119878in}}%
\pgfpathclose%
\pgfusepath{fill}%
\end{pgfscope}%
\begin{pgfscope}%
\pgfpathrectangle{\pgfqpoint{0.407592in}{0.222904in}}{\pgfqpoint{1.411163in}{1.411163in}}%
\pgfusepath{clip}%
\pgfsetbuttcap%
\pgfsetroundjoin%
\definecolor{currentfill}{rgb}{0.909804,0.698039,0.423529}%
\pgfsetfillcolor{currentfill}%
\pgfsetfillopacity{0.120000}%
\pgfsetlinewidth{0.000000pt}%
\definecolor{currentstroke}{rgb}{0.000000,0.000000,0.000000}%
\pgfsetstrokecolor{currentstroke}%
\pgfsetdash{}{0pt}%
\pgfpathmoveto{\pgfqpoint{0.976982in}{1.032183in}}%
\pgfpathlineto{\pgfqpoint{0.992891in}{1.019782in}}%
\pgfpathlineto{\pgfqpoint{1.008820in}{1.008004in}}%
\pgfpathlineto{\pgfqpoint{1.016859in}{1.016748in}}%
\pgfpathlineto{\pgfqpoint{1.024876in}{1.025903in}}%
\pgfpathlineto{\pgfqpoint{1.008977in}{1.037445in}}%
\pgfpathlineto{\pgfqpoint{0.993099in}{1.049603in}}%
\pgfpathlineto{\pgfqpoint{0.985052in}{1.040695in}}%
\pgfpathlineto{\pgfqpoint{0.976982in}{1.032183in}}%
\pgfpathclose%
\pgfusepath{fill}%
\end{pgfscope}%
\begin{pgfscope}%
\pgfpathrectangle{\pgfqpoint{0.407592in}{0.222904in}}{\pgfqpoint{1.411163in}{1.411163in}}%
\pgfusepath{clip}%
\pgfsetbuttcap%
\pgfsetroundjoin%
\definecolor{currentfill}{rgb}{0.909804,0.698039,0.423529}%
\pgfsetfillcolor{currentfill}%
\pgfsetfillopacity{0.120000}%
\pgfsetlinewidth{0.000000pt}%
\definecolor{currentstroke}{rgb}{0.000000,0.000000,0.000000}%
\pgfsetstrokecolor{currentstroke}%
\pgfsetdash{}{0pt}%
\pgfpathmoveto{\pgfqpoint{1.152330in}{1.014947in}}%
\pgfpathlineto{\pgfqpoint{1.168376in}{1.009143in}}%
\pgfpathlineto{\pgfqpoint{1.184459in}{1.004188in}}%
\pgfpathlineto{\pgfqpoint{1.192283in}{1.015888in}}%
\pgfpathlineto{\pgfqpoint{1.200090in}{1.027897in}}%
\pgfpathlineto{\pgfqpoint{1.184024in}{1.032768in}}%
\pgfpathlineto{\pgfqpoint{1.167997in}{1.038452in}}%
\pgfpathlineto{\pgfqpoint{1.160172in}{1.026546in}}%
\pgfpathlineto{\pgfqpoint{1.152330in}{1.014947in}}%
\pgfpathclose%
\pgfusepath{fill}%
\end{pgfscope}%
\begin{pgfscope}%
\pgfpathrectangle{\pgfqpoint{0.407592in}{0.222904in}}{\pgfqpoint{1.411163in}{1.411163in}}%
\pgfusepath{clip}%
\pgfsetbuttcap%
\pgfsetroundjoin%
\definecolor{currentfill}{rgb}{0.909804,0.698039,0.423529}%
\pgfsetfillcolor{currentfill}%
\pgfsetfillopacity{0.120000}%
\pgfsetlinewidth{0.000000pt}%
\definecolor{currentstroke}{rgb}{0.000000,0.000000,0.000000}%
\pgfsetstrokecolor{currentstroke}%
\pgfsetdash{}{0pt}%
\pgfpathmoveto{\pgfqpoint{1.200090in}{1.027897in}}%
\pgfpathlineto{\pgfqpoint{1.216196in}{1.023850in}}%
\pgfpathlineto{\pgfqpoint{1.232345in}{1.020630in}}%
\pgfpathlineto{\pgfqpoint{1.240119in}{1.032956in}}%
\pgfpathlineto{\pgfqpoint{1.247878in}{1.045550in}}%
\pgfpathlineto{\pgfqpoint{1.231745in}{1.048758in}}%
\pgfpathlineto{\pgfqpoint{1.215655in}{1.052757in}}%
\pgfpathlineto{\pgfqpoint{1.207880in}{1.040194in}}%
\pgfpathlineto{\pgfqpoint{1.200090in}{1.027897in}}%
\pgfpathclose%
\pgfusepath{fill}%
\end{pgfscope}%
\begin{pgfscope}%
\pgfpathrectangle{\pgfqpoint{0.407592in}{0.222904in}}{\pgfqpoint{1.411163in}{1.411163in}}%
\pgfusepath{clip}%
\pgfsetbuttcap%
\pgfsetroundjoin%
\definecolor{currentfill}{rgb}{0.909804,0.698039,0.423529}%
\pgfsetfillcolor{currentfill}%
\pgfsetfillopacity{0.120000}%
\pgfsetlinewidth{0.000000pt}%
\definecolor{currentstroke}{rgb}{0.000000,0.000000,0.000000}%
\pgfsetstrokecolor{currentstroke}%
\pgfsetdash{}{0pt}%
\pgfpathmoveto{\pgfqpoint{1.104560in}{1.007171in}}%
\pgfpathlineto{\pgfqpoint{1.120559in}{0.999547in}}%
\pgfpathlineto{\pgfqpoint{1.136591in}{0.992761in}}%
\pgfpathlineto{\pgfqpoint{1.144469in}{1.003678in}}%
\pgfpathlineto{\pgfqpoint{1.152330in}{1.014947in}}%
\pgfpathlineto{\pgfqpoint{1.136319in}{1.021583in}}%
\pgfpathlineto{\pgfqpoint{1.120342in}{1.029027in}}%
\pgfpathlineto{\pgfqpoint{1.112460in}{1.017925in}}%
\pgfpathlineto{\pgfqpoint{1.104560in}{1.007171in}}%
\pgfpathclose%
\pgfusepath{fill}%
\end{pgfscope}%
\begin{pgfscope}%
\pgfpathrectangle{\pgfqpoint{0.407592in}{0.222904in}}{\pgfqpoint{1.411163in}{1.411163in}}%
\pgfusepath{clip}%
\pgfsetbuttcap%
\pgfsetroundjoin%
\definecolor{currentfill}{rgb}{0.909804,0.698039,0.423529}%
\pgfsetfillcolor{currentfill}%
\pgfsetfillopacity{0.120000}%
\pgfsetlinewidth{0.000000pt}%
\definecolor{currentstroke}{rgb}{0.000000,0.000000,0.000000}%
\pgfsetstrokecolor{currentstroke}%
\pgfsetdash{}{0pt}%
\pgfpathmoveto{\pgfqpoint{0.817142in}{1.121783in}}%
\pgfpathlineto{\pgfqpoint{0.833011in}{1.105546in}}%
\pgfpathlineto{\pgfqpoint{0.848891in}{1.089615in}}%
\pgfpathlineto{\pgfqpoint{0.857177in}{1.095949in}}%
\pgfpathlineto{\pgfqpoint{0.865433in}{1.102660in}}%
\pgfpathlineto{\pgfqpoint{0.849591in}{1.118431in}}%
\pgfpathlineto{\pgfqpoint{0.833760in}{1.134511in}}%
\pgfpathlineto{\pgfqpoint{0.825466in}{1.127969in}}%
\pgfpathlineto{\pgfqpoint{0.817142in}{1.121783in}}%
\pgfpathclose%
\pgfusepath{fill}%
\end{pgfscope}%
\begin{pgfscope}%
\pgfpathrectangle{\pgfqpoint{0.407592in}{0.222904in}}{\pgfqpoint{1.411163in}{1.411163in}}%
\pgfusepath{clip}%
\pgfsetbuttcap%
\pgfsetroundjoin%
\definecolor{currentfill}{rgb}{0.909804,0.698039,0.423529}%
\pgfsetfillcolor{currentfill}%
\pgfsetfillopacity{0.120000}%
\pgfsetlinewidth{0.000000pt}%
\definecolor{currentstroke}{rgb}{0.000000,0.000000,0.000000}%
\pgfsetstrokecolor{currentstroke}%
\pgfsetdash{}{0pt}%
\pgfpathmoveto{\pgfqpoint{0.705171in}{1.213622in}}%
\pgfpathlineto{\pgfqpoint{0.721018in}{1.195940in}}%
\pgfpathlineto{\pgfqpoint{0.736874in}{1.178432in}}%
\pgfpathlineto{\pgfqpoint{0.745340in}{1.183693in}}%
\pgfpathlineto{\pgfqpoint{0.753774in}{1.189279in}}%
\pgfpathlineto{\pgfqpoint{0.737958in}{1.206682in}}%
\pgfpathlineto{\pgfqpoint{0.722152in}{1.224261in}}%
\pgfpathlineto{\pgfqpoint{0.713678in}{1.218788in}}%
\pgfpathlineto{\pgfqpoint{0.705171in}{1.213622in}}%
\pgfpathclose%
\pgfusepath{fill}%
\end{pgfscope}%
\begin{pgfscope}%
\pgfpathrectangle{\pgfqpoint{0.407592in}{0.222904in}}{\pgfqpoint{1.411163in}{1.411163in}}%
\pgfusepath{clip}%
\pgfsetbuttcap%
\pgfsetroundjoin%
\definecolor{currentfill}{rgb}{0.909804,0.698039,0.423529}%
\pgfsetfillcolor{currentfill}%
\pgfsetfillopacity{0.120000}%
\pgfsetlinewidth{0.000000pt}%
\definecolor{currentstroke}{rgb}{0.000000,0.000000,0.000000}%
\pgfsetstrokecolor{currentstroke}%
\pgfsetdash{}{0pt}%
\pgfpathmoveto{\pgfqpoint{1.247878in}{1.045550in}}%
\pgfpathlineto{\pgfqpoint{1.264057in}{1.043131in}}%
\pgfpathlineto{\pgfqpoint{1.280284in}{1.041495in}}%
\pgfpathlineto{\pgfqpoint{1.288012in}{1.054297in}}%
\pgfpathlineto{\pgfqpoint{1.295726in}{1.067329in}}%
\pgfpathlineto{\pgfqpoint{1.279514in}{1.069022in}}%
\pgfpathlineto{\pgfqpoint{1.263350in}{1.071463in}}%
\pgfpathlineto{\pgfqpoint{1.255621in}{1.058392in}}%
\pgfpathlineto{\pgfqpoint{1.247878in}{1.045550in}}%
\pgfpathclose%
\pgfusepath{fill}%
\end{pgfscope}%
\begin{pgfscope}%
\pgfpathrectangle{\pgfqpoint{0.407592in}{0.222904in}}{\pgfqpoint{1.411163in}{1.411163in}}%
\pgfusepath{clip}%
\pgfsetbuttcap%
\pgfsetroundjoin%
\definecolor{currentfill}{rgb}{0.909804,0.698039,0.423529}%
\pgfsetfillcolor{currentfill}%
\pgfsetfillopacity{0.120000}%
\pgfsetlinewidth{0.000000pt}%
\definecolor{currentstroke}{rgb}{0.000000,0.000000,0.000000}%
\pgfsetstrokecolor{currentstroke}%
\pgfsetdash{}{0pt}%
\pgfpathmoveto{\pgfqpoint{1.056738in}{1.004849in}}%
\pgfpathlineto{\pgfqpoint{1.072705in}{0.995427in}}%
\pgfpathlineto{\pgfqpoint{1.088699in}{0.986799in}}%
\pgfpathlineto{\pgfqpoint{1.096640in}{0.996789in}}%
\pgfpathlineto{\pgfqpoint{1.104560in}{1.007171in}}%
\pgfpathlineto{\pgfqpoint{1.088590in}{1.015603in}}%
\pgfpathlineto{\pgfqpoint{1.072648in}{1.024807in}}%
\pgfpathlineto{\pgfqpoint{1.064704in}{1.014636in}}%
\pgfpathlineto{\pgfqpoint{1.056738in}{1.004849in}}%
\pgfpathclose%
\pgfusepath{fill}%
\end{pgfscope}%
\begin{pgfscope}%
\pgfpathrectangle{\pgfqpoint{0.407592in}{0.222904in}}{\pgfqpoint{1.411163in}{1.411163in}}%
\pgfusepath{clip}%
\pgfsetbuttcap%
\pgfsetroundjoin%
\definecolor{currentfill}{rgb}{0.909804,0.698039,0.423529}%
\pgfsetfillcolor{currentfill}%
\pgfsetfillopacity{0.120000}%
\pgfsetlinewidth{0.000000pt}%
\definecolor{currentstroke}{rgb}{0.000000,0.000000,0.000000}%
\pgfsetstrokecolor{currentstroke}%
\pgfsetdash{}{0pt}%
\pgfpathmoveto{\pgfqpoint{0.928930in}{1.043308in}}%
\pgfpathlineto{\pgfqpoint{0.944839in}{1.029589in}}%
\pgfpathlineto{\pgfqpoint{0.960764in}{1.016402in}}%
\pgfpathlineto{\pgfqpoint{0.968886in}{1.024081in}}%
\pgfpathlineto{\pgfqpoint{0.976982in}{1.032183in}}%
\pgfpathlineto{\pgfqpoint{0.961090in}{1.045160in}}%
\pgfpathlineto{\pgfqpoint{0.945215in}{1.058668in}}%
\pgfpathlineto{\pgfqpoint{0.937085in}{1.050786in}}%
\pgfpathlineto{\pgfqpoint{0.928930in}{1.043308in}}%
\pgfpathclose%
\pgfusepath{fill}%
\end{pgfscope}%
\begin{pgfscope}%
\pgfpathrectangle{\pgfqpoint{0.407592in}{0.222904in}}{\pgfqpoint{1.411163in}{1.411163in}}%
\pgfusepath{clip}%
\pgfsetbuttcap%
\pgfsetroundjoin%
\definecolor{currentfill}{rgb}{0.909804,0.698039,0.423529}%
\pgfsetfillcolor{currentfill}%
\pgfsetfillopacity{0.120000}%
\pgfsetlinewidth{0.000000pt}%
\definecolor{currentstroke}{rgb}{0.000000,0.000000,0.000000}%
\pgfsetstrokecolor{currentstroke}%
\pgfsetdash{}{0pt}%
\pgfpathmoveto{\pgfqpoint{1.295726in}{1.067329in}}%
\pgfpathlineto{\pgfqpoint{1.311989in}{1.066375in}}%
\pgfpathlineto{\pgfqpoint{1.328303in}{1.066145in}}%
\pgfpathlineto{\pgfqpoint{1.335989in}{1.079291in}}%
\pgfpathlineto{\pgfqpoint{1.343660in}{1.092638in}}%
\pgfpathlineto{\pgfqpoint{1.327360in}{1.092982in}}%
\pgfpathlineto{\pgfqpoint{1.311112in}{1.094022in}}%
\pgfpathlineto{\pgfqpoint{1.303426in}{1.080576in}}%
\pgfpathlineto{\pgfqpoint{1.295726in}{1.067329in}}%
\pgfpathclose%
\pgfusepath{fill}%
\end{pgfscope}%
\begin{pgfscope}%
\pgfpathrectangle{\pgfqpoint{0.407592in}{0.222904in}}{\pgfqpoint{1.411163in}{1.411163in}}%
\pgfusepath{clip}%
\pgfsetbuttcap%
\pgfsetroundjoin%
\definecolor{currentfill}{rgb}{0.909804,0.698039,0.423529}%
\pgfsetfillcolor{currentfill}%
\pgfsetfillopacity{0.120000}%
\pgfsetlinewidth{0.000000pt}%
\definecolor{currentstroke}{rgb}{0.000000,0.000000,0.000000}%
\pgfsetstrokecolor{currentstroke}%
\pgfsetdash{}{0pt}%
\pgfpathmoveto{\pgfqpoint{0.736874in}{1.178432in}}%
\pgfpathlineto{\pgfqpoint{0.752741in}{1.161114in}}%
\pgfpathlineto{\pgfqpoint{0.768617in}{1.144006in}}%
\pgfpathlineto{\pgfqpoint{0.777043in}{1.149364in}}%
\pgfpathlineto{\pgfqpoint{0.785437in}{1.155067in}}%
\pgfpathlineto{\pgfqpoint{0.769601in}{1.172068in}}%
\pgfpathlineto{\pgfqpoint{0.753774in}{1.189279in}}%
\pgfpathlineto{\pgfqpoint{0.745340in}{1.183693in}}%
\pgfpathlineto{\pgfqpoint{0.736874in}{1.178432in}}%
\pgfpathclose%
\pgfusepath{fill}%
\end{pgfscope}%
\begin{pgfscope}%
\pgfpathrectangle{\pgfqpoint{0.407592in}{0.222904in}}{\pgfqpoint{1.411163in}{1.411163in}}%
\pgfusepath{clip}%
\pgfsetbuttcap%
\pgfsetroundjoin%
\definecolor{currentfill}{rgb}{0.909804,0.698039,0.423529}%
\pgfsetfillcolor{currentfill}%
\pgfsetfillopacity{0.120000}%
\pgfsetlinewidth{0.000000pt}%
\definecolor{currentstroke}{rgb}{0.000000,0.000000,0.000000}%
\pgfsetstrokecolor{currentstroke}%
\pgfsetdash{}{0pt}%
\pgfpathmoveto{\pgfqpoint{1.406923in}{1.148521in}}%
\pgfpathlineto{\pgfqpoint{1.423363in}{1.149807in}}%
\pgfpathlineto{\pgfqpoint{1.439860in}{1.151675in}}%
\pgfpathlineto{\pgfqpoint{1.447439in}{1.165521in}}%
\pgfpathlineto{\pgfqpoint{1.430948in}{1.163745in}}%
\pgfpathlineto{\pgfqpoint{1.414516in}{1.162542in}}%
\pgfusepath{fill}%
\end{pgfscope}%
\begin{pgfscope}%
\pgfpathrectangle{\pgfqpoint{0.407592in}{0.222904in}}{\pgfqpoint{1.411163in}{1.411163in}}%
\pgfusepath{clip}%
\pgfsetbuttcap%
\pgfsetroundjoin%
\definecolor{currentfill}{rgb}{0.909804,0.698039,0.423529}%
\pgfsetfillcolor{currentfill}%
\pgfsetfillopacity{0.120000}%
\pgfsetlinewidth{0.000000pt}%
\definecolor{currentstroke}{rgb}{0.000000,0.000000,0.000000}%
\pgfsetstrokecolor{currentstroke}%
\pgfsetdash{}{0pt}%
\pgfpathmoveto{\pgfqpoint{0.848891in}{1.089615in}}%
\pgfpathlineto{\pgfqpoint{0.864783in}{1.074021in}}%
\pgfpathlineto{\pgfqpoint{0.880687in}{1.058797in}}%
\pgfpathlineto{\pgfqpoint{0.888935in}{1.065282in}}%
\pgfpathlineto{\pgfqpoint{0.897155in}{1.072168in}}%
\pgfpathlineto{\pgfqpoint{0.881287in}{1.087228in}}%
\pgfpathlineto{\pgfqpoint{0.865433in}{1.102660in}}%
\pgfpathlineto{\pgfqpoint{0.857177in}{1.095949in}}%
\pgfpathlineto{\pgfqpoint{0.848891in}{1.089615in}}%
\pgfpathclose%
\pgfusepath{fill}%
\end{pgfscope}%
\begin{pgfscope}%
\pgfpathrectangle{\pgfqpoint{0.407592in}{0.222904in}}{\pgfqpoint{1.411163in}{1.411163in}}%
\pgfusepath{clip}%
\pgfsetbuttcap%
\pgfsetroundjoin%
\definecolor{currentfill}{rgb}{0.909804,0.698039,0.423529}%
\pgfsetfillcolor{currentfill}%
\pgfsetfillopacity{0.120000}%
\pgfsetlinewidth{0.000000pt}%
\definecolor{currentstroke}{rgb}{0.000000,0.000000,0.000000}%
\pgfsetstrokecolor{currentstroke}%
\pgfsetdash{}{0pt}%
\pgfpathmoveto{\pgfqpoint{1.008820in}{1.008004in}}%
\pgfpathlineto{\pgfqpoint{1.024769in}{0.996898in}}%
\pgfpathlineto{\pgfqpoint{1.040740in}{0.986511in}}%
\pgfpathlineto{\pgfqpoint{1.048750in}{0.995467in}}%
\pgfpathlineto{\pgfqpoint{1.056738in}{1.004849in}}%
\pgfpathlineto{\pgfqpoint{1.040795in}{1.015022in}}%
\pgfpathlineto{\pgfqpoint{1.024876in}{1.025903in}}%
\pgfpathlineto{\pgfqpoint{1.016859in}{1.016748in}}%
\pgfpathlineto{\pgfqpoint{1.008820in}{1.008004in}}%
\pgfpathclose%
\pgfusepath{fill}%
\end{pgfscope}%
\begin{pgfscope}%
\pgfpathrectangle{\pgfqpoint{0.407592in}{0.222904in}}{\pgfqpoint{1.411163in}{1.411163in}}%
\pgfusepath{clip}%
\pgfsetbuttcap%
\pgfsetroundjoin%
\definecolor{currentfill}{rgb}{0.909804,0.698039,0.423529}%
\pgfsetfillcolor{currentfill}%
\pgfsetfillopacity{0.120000}%
\pgfsetlinewidth{0.000000pt}%
\definecolor{currentstroke}{rgb}{0.000000,0.000000,0.000000}%
\pgfsetstrokecolor{currentstroke}%
\pgfsetdash{}{0pt}%
\pgfpathmoveto{\pgfqpoint{1.343660in}{1.092638in}}%
\pgfpathlineto{\pgfqpoint{1.360015in}{1.092973in}}%
\pgfpathlineto{\pgfqpoint{1.376425in}{1.093971in}}%
\pgfpathlineto{\pgfqpoint{1.384069in}{1.107355in}}%
\pgfpathlineto{\pgfqpoint{1.391700in}{1.120917in}}%
\pgfpathlineto{\pgfqpoint{1.375304in}{1.120079in}}%
\pgfpathlineto{\pgfqpoint{1.358964in}{1.119878in}}%
\pgfpathlineto{\pgfqpoint{1.351319in}{1.106172in}}%
\pgfpathlineto{\pgfqpoint{1.343660in}{1.092638in}}%
\pgfpathclose%
\pgfusepath{fill}%
\end{pgfscope}%
\begin{pgfscope}%
\pgfpathrectangle{\pgfqpoint{0.407592in}{0.222904in}}{\pgfqpoint{1.411163in}{1.411163in}}%
\pgfusepath{clip}%
\pgfsetbuttcap%
\pgfsetroundjoin%
\definecolor{currentfill}{rgb}{0.909804,0.698039,0.423529}%
\pgfsetfillcolor{currentfill}%
\pgfsetfillopacity{0.120000}%
\pgfsetlinewidth{0.000000pt}%
\definecolor{currentstroke}{rgb}{0.000000,0.000000,0.000000}%
\pgfsetstrokecolor{currentstroke}%
\pgfsetdash{}{0pt}%
\pgfpathmoveto{\pgfqpoint{0.768617in}{1.144006in}}%
\pgfpathlineto{\pgfqpoint{0.784503in}{1.127127in}}%
\pgfpathlineto{\pgfqpoint{0.800401in}{1.110500in}}%
\pgfpathlineto{\pgfqpoint{0.808787in}{1.115958in}}%
\pgfpathlineto{\pgfqpoint{0.817142in}{1.121783in}}%
\pgfpathlineto{\pgfqpoint{0.801284in}{1.138298in}}%
\pgfpathlineto{\pgfqpoint{0.785437in}{1.155067in}}%
\pgfpathlineto{\pgfqpoint{0.777043in}{1.149364in}}%
\pgfpathlineto{\pgfqpoint{0.768617in}{1.144006in}}%
\pgfpathclose%
\pgfusepath{fill}%
\end{pgfscope}%
\begin{pgfscope}%
\pgfpathrectangle{\pgfqpoint{0.407592in}{0.222904in}}{\pgfqpoint{1.411163in}{1.411163in}}%
\pgfusepath{clip}%
\pgfsetbuttcap%
\pgfsetroundjoin%
\definecolor{currentfill}{rgb}{0.909804,0.698039,0.423529}%
\pgfsetfillcolor{currentfill}%
\pgfsetfillopacity{0.120000}%
\pgfsetlinewidth{0.000000pt}%
\definecolor{currentstroke}{rgb}{0.000000,0.000000,0.000000}%
\pgfsetstrokecolor{currentstroke}%
\pgfsetdash{}{0pt}%
\pgfpathmoveto{\pgfqpoint{1.136591in}{0.992761in}}%
\pgfpathlineto{\pgfqpoint{1.152657in}{0.986840in}}%
\pgfpathlineto{\pgfqpoint{1.168761in}{0.981802in}}%
\pgfpathlineto{\pgfqpoint{1.176619in}{0.992818in}}%
\pgfpathlineto{\pgfqpoint{1.184459in}{1.004188in}}%
\pgfpathlineto{\pgfqpoint{1.168376in}{1.009143in}}%
\pgfpathlineto{\pgfqpoint{1.152330in}{1.014947in}}%
\pgfpathlineto{\pgfqpoint{1.144469in}{1.003678in}}%
\pgfpathlineto{\pgfqpoint{1.136591in}{0.992761in}}%
\pgfpathclose%
\pgfusepath{fill}%
\end{pgfscope}%
\begin{pgfscope}%
\pgfpathrectangle{\pgfqpoint{0.407592in}{0.222904in}}{\pgfqpoint{1.411163in}{1.411163in}}%
\pgfusepath{clip}%
\pgfsetbuttcap%
\pgfsetroundjoin%
\definecolor{currentfill}{rgb}{0.909804,0.698039,0.423529}%
\pgfsetfillcolor{currentfill}%
\pgfsetfillopacity{0.120000}%
\pgfsetlinewidth{0.000000pt}%
\definecolor{currentstroke}{rgb}{0.000000,0.000000,0.000000}%
\pgfsetstrokecolor{currentstroke}%
\pgfsetdash{}{0pt}%
\pgfpathmoveto{\pgfqpoint{1.184459in}{1.004188in}}%
\pgfpathlineto{\pgfqpoint{1.200583in}{1.000093in}}%
\pgfpathlineto{\pgfqpoint{1.216750in}{0.996863in}}%
\pgfpathlineto{\pgfqpoint{1.224556in}{1.008592in}}%
\pgfpathlineto{\pgfqpoint{1.232345in}{1.020630in}}%
\pgfpathlineto{\pgfqpoint{1.216196in}{1.023850in}}%
\pgfpathlineto{\pgfqpoint{1.200090in}{1.027897in}}%
\pgfpathlineto{\pgfqpoint{1.192283in}{1.015888in}}%
\pgfpathlineto{\pgfqpoint{1.184459in}{1.004188in}}%
\pgfpathclose%
\pgfusepath{fill}%
\end{pgfscope}%
\begin{pgfscope}%
\pgfpathrectangle{\pgfqpoint{0.407592in}{0.222904in}}{\pgfqpoint{1.411163in}{1.411163in}}%
\pgfusepath{clip}%
\pgfsetbuttcap%
\pgfsetroundjoin%
\definecolor{currentfill}{rgb}{0.909804,0.698039,0.423529}%
\pgfsetfillcolor{currentfill}%
\pgfsetfillopacity{0.120000}%
\pgfsetlinewidth{0.000000pt}%
\definecolor{currentstroke}{rgb}{0.000000,0.000000,0.000000}%
\pgfsetstrokecolor{currentstroke}%
\pgfsetdash{}{0pt}%
\pgfpathmoveto{\pgfqpoint{0.656311in}{1.240179in}}%
\pgfpathlineto{\pgfqpoint{0.672180in}{1.222126in}}%
\pgfpathlineto{\pgfqpoint{0.688058in}{1.204219in}}%
\pgfpathlineto{\pgfqpoint{0.696632in}{1.208765in}}%
\pgfpathlineto{\pgfqpoint{0.705171in}{1.213622in}}%
\pgfpathlineto{\pgfqpoint{0.689334in}{1.231462in}}%
\pgfpathlineto{\pgfqpoint{0.673507in}{1.249448in}}%
\pgfpathlineto{\pgfqpoint{0.664926in}{1.244667in}}%
\pgfpathlineto{\pgfqpoint{0.656311in}{1.240179in}}%
\pgfpathclose%
\pgfusepath{fill}%
\end{pgfscope}%
\begin{pgfscope}%
\pgfpathrectangle{\pgfqpoint{0.407592in}{0.222904in}}{\pgfqpoint{1.411163in}{1.411163in}}%
\pgfusepath{clip}%
\pgfsetbuttcap%
\pgfsetroundjoin%
\definecolor{currentfill}{rgb}{0.909804,0.698039,0.423529}%
\pgfsetfillcolor{currentfill}%
\pgfsetfillopacity{0.120000}%
\pgfsetlinewidth{0.000000pt}%
\definecolor{currentstroke}{rgb}{0.000000,0.000000,0.000000}%
\pgfsetstrokecolor{currentstroke}%
\pgfsetdash{}{0pt}%
\pgfpathmoveto{\pgfqpoint{0.960764in}{1.016402in}}%
\pgfpathlineto{\pgfqpoint{0.976707in}{1.003795in}}%
\pgfpathlineto{\pgfqpoint{0.992667in}{0.991816in}}%
\pgfpathlineto{\pgfqpoint{1.000756in}{0.999688in}}%
\pgfpathlineto{\pgfqpoint{1.008820in}{1.008004in}}%
\pgfpathlineto{\pgfqpoint{0.992891in}{1.019782in}}%
\pgfpathlineto{\pgfqpoint{0.976982in}{1.032183in}}%
\pgfpathlineto{\pgfqpoint{0.968886in}{1.024081in}}%
\pgfpathlineto{\pgfqpoint{0.960764in}{1.016402in}}%
\pgfpathclose%
\pgfusepath{fill}%
\end{pgfscope}%
\begin{pgfscope}%
\pgfpathrectangle{\pgfqpoint{0.407592in}{0.222904in}}{\pgfqpoint{1.411163in}{1.411163in}}%
\pgfusepath{clip}%
\pgfsetbuttcap%
\pgfsetroundjoin%
\definecolor{currentfill}{rgb}{0.909804,0.698039,0.423529}%
\pgfsetfillcolor{currentfill}%
\pgfsetfillopacity{0.120000}%
\pgfsetlinewidth{0.000000pt}%
\definecolor{currentstroke}{rgb}{0.000000,0.000000,0.000000}%
\pgfsetstrokecolor{currentstroke}%
\pgfsetdash{}{0pt}%
\pgfpathmoveto{\pgfqpoint{0.880687in}{1.058797in}}%
\pgfpathlineto{\pgfqpoint{0.896604in}{1.043978in}}%
\pgfpathlineto{\pgfqpoint{0.912535in}{1.029605in}}%
\pgfpathlineto{\pgfqpoint{0.920746in}{1.036245in}}%
\pgfpathlineto{\pgfqpoint{0.928930in}{1.043308in}}%
\pgfpathlineto{\pgfqpoint{0.913035in}{1.057515in}}%
\pgfpathlineto{\pgfqpoint{0.897155in}{1.072168in}}%
\pgfpathlineto{\pgfqpoint{0.888935in}{1.065282in}}%
\pgfpathlineto{\pgfqpoint{0.880687in}{1.058797in}}%
\pgfpathclose%
\pgfusepath{fill}%
\end{pgfscope}%
\begin{pgfscope}%
\pgfpathrectangle{\pgfqpoint{0.407592in}{0.222904in}}{\pgfqpoint{1.411163in}{1.411163in}}%
\pgfusepath{clip}%
\pgfsetbuttcap%
\pgfsetroundjoin%
\definecolor{currentfill}{rgb}{0.909804,0.698039,0.423529}%
\pgfsetfillcolor{currentfill}%
\pgfsetfillopacity{0.120000}%
\pgfsetlinewidth{0.000000pt}%
\definecolor{currentstroke}{rgb}{0.000000,0.000000,0.000000}%
\pgfsetstrokecolor{currentstroke}%
\pgfsetdash{}{0pt}%
\pgfpathmoveto{\pgfqpoint{1.088699in}{0.986799in}}%
\pgfpathlineto{\pgfqpoint{1.104722in}{0.979004in}}%
\pgfpathlineto{\pgfqpoint{1.120777in}{0.972075in}}%
\pgfpathlineto{\pgfqpoint{1.128694in}{0.982219in}}%
\pgfpathlineto{\pgfqpoint{1.136591in}{0.992761in}}%
\pgfpathlineto{\pgfqpoint{1.120559in}{0.999547in}}%
\pgfpathlineto{\pgfqpoint{1.104560in}{1.007171in}}%
\pgfpathlineto{\pgfqpoint{1.096640in}{0.996789in}}%
\pgfpathlineto{\pgfqpoint{1.088699in}{0.986799in}}%
\pgfpathclose%
\pgfusepath{fill}%
\end{pgfscope}%
\begin{pgfscope}%
\pgfpathrectangle{\pgfqpoint{0.407592in}{0.222904in}}{\pgfqpoint{1.411163in}{1.411163in}}%
\pgfusepath{clip}%
\pgfsetbuttcap%
\pgfsetroundjoin%
\definecolor{currentfill}{rgb}{0.909804,0.698039,0.423529}%
\pgfsetfillcolor{currentfill}%
\pgfsetfillopacity{0.120000}%
\pgfsetlinewidth{0.000000pt}%
\definecolor{currentstroke}{rgb}{0.000000,0.000000,0.000000}%
\pgfsetstrokecolor{currentstroke}%
\pgfsetdash{}{0pt}%
\pgfpathmoveto{\pgfqpoint{1.232345in}{1.020630in}}%
\pgfpathlineto{\pgfqpoint{1.248540in}{1.018236in}}%
\pgfpathlineto{\pgfqpoint{1.264783in}{1.016660in}}%
\pgfpathlineto{\pgfqpoint{1.272541in}{1.028944in}}%
\pgfpathlineto{\pgfqpoint{1.280284in}{1.041495in}}%
\pgfpathlineto{\pgfqpoint{1.264057in}{1.043131in}}%
\pgfpathlineto{\pgfqpoint{1.247878in}{1.045550in}}%
\pgfpathlineto{\pgfqpoint{1.240119in}{1.032956in}}%
\pgfpathlineto{\pgfqpoint{1.232345in}{1.020630in}}%
\pgfpathclose%
\pgfusepath{fill}%
\end{pgfscope}%
\begin{pgfscope}%
\pgfpathrectangle{\pgfqpoint{0.407592in}{0.222904in}}{\pgfqpoint{1.411163in}{1.411163in}}%
\pgfusepath{clip}%
\pgfsetbuttcap%
\pgfsetroundjoin%
\definecolor{currentfill}{rgb}{0.909804,0.698039,0.423529}%
\pgfsetfillcolor{currentfill}%
\pgfsetfillopacity{0.120000}%
\pgfsetlinewidth{0.000000pt}%
\definecolor{currentstroke}{rgb}{0.000000,0.000000,0.000000}%
\pgfsetstrokecolor{currentstroke}%
\pgfsetdash{}{0pt}%
\pgfpathmoveto{\pgfqpoint{1.391700in}{1.120917in}}%
\pgfpathlineto{\pgfqpoint{1.408154in}{1.122375in}}%
\pgfpathlineto{\pgfqpoint{1.424666in}{1.124435in}}%
\pgfpathlineto{\pgfqpoint{1.432269in}{1.137976in}}%
\pgfpathlineto{\pgfqpoint{1.439860in}{1.151675in}}%
\pgfpathlineto{\pgfqpoint{1.423363in}{1.149807in}}%
\pgfpathlineto{\pgfqpoint{1.406923in}{1.148521in}}%
\pgfpathlineto{\pgfqpoint{1.399318in}{1.134643in}}%
\pgfpathlineto{\pgfqpoint{1.391700in}{1.120917in}}%
\pgfpathclose%
\pgfusepath{fill}%
\end{pgfscope}%
\begin{pgfscope}%
\pgfpathrectangle{\pgfqpoint{0.407592in}{0.222904in}}{\pgfqpoint{1.411163in}{1.411163in}}%
\pgfusepath{clip}%
\pgfsetbuttcap%
\pgfsetroundjoin%
\definecolor{currentfill}{rgb}{0.909804,0.698039,0.423529}%
\pgfsetfillcolor{currentfill}%
\pgfsetfillopacity{0.120000}%
\pgfsetlinewidth{0.000000pt}%
\definecolor{currentstroke}{rgb}{0.000000,0.000000,0.000000}%
\pgfsetstrokecolor{currentstroke}%
\pgfsetdash{}{0pt}%
\pgfpathmoveto{\pgfqpoint{0.800401in}{1.110500in}}%
\pgfpathlineto{\pgfqpoint{0.816308in}{1.094150in}}%
\pgfpathlineto{\pgfqpoint{0.832227in}{1.078106in}}%
\pgfpathlineto{\pgfqpoint{0.840575in}{1.083666in}}%
\pgfpathlineto{\pgfqpoint{0.848891in}{1.089615in}}%
\pgfpathlineto{\pgfqpoint{0.833011in}{1.105546in}}%
\pgfpathlineto{\pgfqpoint{0.817142in}{1.121783in}}%
\pgfpathlineto{\pgfqpoint{0.808787in}{1.115958in}}%
\pgfpathlineto{\pgfqpoint{0.800401in}{1.110500in}}%
\pgfpathclose%
\pgfusepath{fill}%
\end{pgfscope}%
\begin{pgfscope}%
\pgfpathrectangle{\pgfqpoint{0.407592in}{0.222904in}}{\pgfqpoint{1.411163in}{1.411163in}}%
\pgfusepath{clip}%
\pgfsetbuttcap%
\pgfsetroundjoin%
\definecolor{currentfill}{rgb}{0.909804,0.698039,0.423529}%
\pgfsetfillcolor{currentfill}%
\pgfsetfillopacity{0.120000}%
\pgfsetlinewidth{0.000000pt}%
\definecolor{currentstroke}{rgb}{0.000000,0.000000,0.000000}%
\pgfsetstrokecolor{currentstroke}%
\pgfsetdash{}{0pt}%
\pgfpathmoveto{\pgfqpoint{1.040740in}{0.986511in}}%
\pgfpathlineto{\pgfqpoint{1.056734in}{0.976891in}}%
\pgfpathlineto{\pgfqpoint{1.072755in}{0.968084in}}%
\pgfpathlineto{\pgfqpoint{1.080738in}{0.977224in}}%
\pgfpathlineto{\pgfqpoint{1.088699in}{0.986799in}}%
\pgfpathlineto{\pgfqpoint{1.072705in}{0.995427in}}%
\pgfpathlineto{\pgfqpoint{1.056738in}{1.004849in}}%
\pgfpathlineto{\pgfqpoint{1.048750in}{0.995467in}}%
\pgfpathlineto{\pgfqpoint{1.040740in}{0.986511in}}%
\pgfpathclose%
\pgfusepath{fill}%
\end{pgfscope}%
\begin{pgfscope}%
\pgfpathrectangle{\pgfqpoint{0.407592in}{0.222904in}}{\pgfqpoint{1.411163in}{1.411163in}}%
\pgfusepath{clip}%
\pgfsetbuttcap%
\pgfsetroundjoin%
\definecolor{currentfill}{rgb}{0.909804,0.698039,0.423529}%
\pgfsetfillcolor{currentfill}%
\pgfsetfillopacity{0.120000}%
\pgfsetlinewidth{0.000000pt}%
\definecolor{currentstroke}{rgb}{0.000000,0.000000,0.000000}%
\pgfsetstrokecolor{currentstroke}%
\pgfsetdash{}{0pt}%
\pgfpathmoveto{\pgfqpoint{1.280284in}{1.041495in}}%
\pgfpathlineto{\pgfqpoint{1.296561in}{1.040631in}}%
\pgfpathlineto{\pgfqpoint{1.312891in}{1.040522in}}%
\pgfpathlineto{\pgfqpoint{1.320604in}{1.053217in}}%
\pgfpathlineto{\pgfqpoint{1.328303in}{1.066145in}}%
\pgfpathlineto{\pgfqpoint{1.311989in}{1.066375in}}%
\pgfpathlineto{\pgfqpoint{1.295726in}{1.067329in}}%
\pgfpathlineto{\pgfqpoint{1.288012in}{1.054297in}}%
\pgfpathlineto{\pgfqpoint{1.280284in}{1.041495in}}%
\pgfpathclose%
\pgfusepath{fill}%
\end{pgfscope}%
\begin{pgfscope}%
\pgfpathrectangle{\pgfqpoint{0.407592in}{0.222904in}}{\pgfqpoint{1.411163in}{1.411163in}}%
\pgfusepath{clip}%
\pgfsetbuttcap%
\pgfsetroundjoin%
\definecolor{currentfill}{rgb}{0.909804,0.698039,0.423529}%
\pgfsetfillcolor{currentfill}%
\pgfsetfillopacity{0.120000}%
\pgfsetlinewidth{0.000000pt}%
\definecolor{currentstroke}{rgb}{0.000000,0.000000,0.000000}%
\pgfsetstrokecolor{currentstroke}%
\pgfsetdash{}{0pt}%
\pgfpathmoveto{\pgfqpoint{0.688058in}{1.204219in}}%
\pgfpathlineto{\pgfqpoint{0.703946in}{1.186469in}}%
\pgfpathlineto{\pgfqpoint{0.719843in}{1.168894in}}%
\pgfpathlineto{\pgfqpoint{0.728375in}{1.173498in}}%
\pgfpathlineto{\pgfqpoint{0.736874in}{1.178432in}}%
\pgfpathlineto{\pgfqpoint{0.721018in}{1.195940in}}%
\pgfpathlineto{\pgfqpoint{0.705171in}{1.213622in}}%
\pgfpathlineto{\pgfqpoint{0.696632in}{1.208765in}}%
\pgfpathlineto{\pgfqpoint{0.688058in}{1.204219in}}%
\pgfpathclose%
\pgfusepath{fill}%
\end{pgfscope}%
\begin{pgfscope}%
\pgfpathrectangle{\pgfqpoint{0.407592in}{0.222904in}}{\pgfqpoint{1.411163in}{1.411163in}}%
\pgfusepath{clip}%
\pgfsetbuttcap%
\pgfsetroundjoin%
\definecolor{currentfill}{rgb}{0.909804,0.698039,0.423529}%
\pgfsetfillcolor{currentfill}%
\pgfsetfillopacity{0.120000}%
\pgfsetlinewidth{0.000000pt}%
\definecolor{currentstroke}{rgb}{0.000000,0.000000,0.000000}%
\pgfsetstrokecolor{currentstroke}%
\pgfsetdash{}{0pt}%
\pgfpathmoveto{\pgfqpoint{0.912535in}{1.029605in}}%
\pgfpathlineto{\pgfqpoint{0.928480in}{1.015719in}}%
\pgfpathlineto{\pgfqpoint{0.944440in}{1.002367in}}%
\pgfpathlineto{\pgfqpoint{0.952616in}{1.009161in}}%
\pgfpathlineto{\pgfqpoint{0.960764in}{1.016402in}}%
\pgfpathlineto{\pgfqpoint{0.944839in}{1.029589in}}%
\pgfpathlineto{\pgfqpoint{0.928930in}{1.043308in}}%
\pgfpathlineto{\pgfqpoint{0.920746in}{1.036245in}}%
\pgfpathlineto{\pgfqpoint{0.912535in}{1.029605in}}%
\pgfpathclose%
\pgfusepath{fill}%
\end{pgfscope}%
\begin{pgfscope}%
\pgfpathrectangle{\pgfqpoint{0.407592in}{0.222904in}}{\pgfqpoint{1.411163in}{1.411163in}}%
\pgfusepath{clip}%
\pgfsetbuttcap%
\pgfsetroundjoin%
\definecolor{currentfill}{rgb}{0.909804,0.698039,0.423529}%
\pgfsetfillcolor{currentfill}%
\pgfsetfillopacity{0.120000}%
\pgfsetlinewidth{0.000000pt}%
\definecolor{currentstroke}{rgb}{0.000000,0.000000,0.000000}%
\pgfsetstrokecolor{currentstroke}%
\pgfsetdash{}{0pt}%
\pgfpathmoveto{\pgfqpoint{1.328303in}{1.066145in}}%
\pgfpathlineto{\pgfqpoint{1.344672in}{1.066624in}}%
\pgfpathlineto{\pgfqpoint{1.361097in}{1.067790in}}%
\pgfpathlineto{\pgfqpoint{1.368768in}{1.080778in}}%
\pgfpathlineto{\pgfqpoint{1.376425in}{1.093971in}}%
\pgfpathlineto{\pgfqpoint{1.360015in}{1.092973in}}%
\pgfpathlineto{\pgfqpoint{1.343660in}{1.092638in}}%
\pgfpathlineto{\pgfqpoint{1.335989in}{1.079291in}}%
\pgfpathlineto{\pgfqpoint{1.328303in}{1.066145in}}%
\pgfpathclose%
\pgfusepath{fill}%
\end{pgfscope}%
\begin{pgfscope}%
\pgfpathrectangle{\pgfqpoint{0.407592in}{0.222904in}}{\pgfqpoint{1.411163in}{1.411163in}}%
\pgfusepath{clip}%
\pgfsetbuttcap%
\pgfsetroundjoin%
\definecolor{currentfill}{rgb}{0.909804,0.698039,0.423529}%
\pgfsetfillcolor{currentfill}%
\pgfsetfillopacity{0.120000}%
\pgfsetlinewidth{0.000000pt}%
\definecolor{currentstroke}{rgb}{0.000000,0.000000,0.000000}%
\pgfsetstrokecolor{currentstroke}%
\pgfsetdash{}{0pt}%
\pgfpathmoveto{\pgfqpoint{0.992667in}{0.991816in}}%
\pgfpathlineto{\pgfqpoint{1.008647in}{0.980516in}}%
\pgfpathlineto{\pgfqpoint{1.024648in}{0.969945in}}%
\pgfpathlineto{\pgfqpoint{1.032706in}{0.977997in}}%
\pgfpathlineto{\pgfqpoint{1.040740in}{0.986511in}}%
\pgfpathlineto{\pgfqpoint{1.024769in}{0.996898in}}%
\pgfpathlineto{\pgfqpoint{1.008820in}{1.008004in}}%
\pgfpathlineto{\pgfqpoint{1.000756in}{0.999688in}}%
\pgfpathlineto{\pgfqpoint{0.992667in}{0.991816in}}%
\pgfpathclose%
\pgfusepath{fill}%
\end{pgfscope}%
\begin{pgfscope}%
\pgfpathrectangle{\pgfqpoint{0.407592in}{0.222904in}}{\pgfqpoint{1.411163in}{1.411163in}}%
\pgfusepath{clip}%
\pgfsetbuttcap%
\pgfsetroundjoin%
\definecolor{currentfill}{rgb}{0.909804,0.698039,0.423529}%
\pgfsetfillcolor{currentfill}%
\pgfsetfillopacity{0.120000}%
\pgfsetlinewidth{0.000000pt}%
\definecolor{currentstroke}{rgb}{0.000000,0.000000,0.000000}%
\pgfsetstrokecolor{currentstroke}%
\pgfsetdash{}{0pt}%
\pgfpathmoveto{\pgfqpoint{0.719843in}{1.168894in}}%
\pgfpathlineto{\pgfqpoint{0.735750in}{1.151509in}}%
\pgfpathlineto{\pgfqpoint{0.751666in}{1.134334in}}%
\pgfpathlineto{\pgfqpoint{0.760158in}{1.138995in}}%
\pgfpathlineto{\pgfqpoint{0.768617in}{1.144006in}}%
\pgfpathlineto{\pgfqpoint{0.752741in}{1.161114in}}%
\pgfpathlineto{\pgfqpoint{0.736874in}{1.178432in}}%
\pgfpathlineto{\pgfqpoint{0.728375in}{1.173498in}}%
\pgfpathlineto{\pgfqpoint{0.719843in}{1.168894in}}%
\pgfpathclose%
\pgfusepath{fill}%
\end{pgfscope}%
\begin{pgfscope}%
\pgfpathrectangle{\pgfqpoint{0.407592in}{0.222904in}}{\pgfqpoint{1.411163in}{1.411163in}}%
\pgfusepath{clip}%
\pgfsetbuttcap%
\pgfsetroundjoin%
\definecolor{currentfill}{rgb}{0.909804,0.698039,0.423529}%
\pgfsetfillcolor{currentfill}%
\pgfsetfillopacity{0.120000}%
\pgfsetlinewidth{0.000000pt}%
\definecolor{currentstroke}{rgb}{0.000000,0.000000,0.000000}%
\pgfsetstrokecolor{currentstroke}%
\pgfsetdash{}{0pt}%
\pgfpathmoveto{\pgfqpoint{0.832227in}{1.078106in}}%
\pgfpathlineto{\pgfqpoint{0.848158in}{1.062396in}}%
\pgfpathlineto{\pgfqpoint{0.864100in}{1.047056in}}%
\pgfpathlineto{\pgfqpoint{0.872409in}{1.052719in}}%
\pgfpathlineto{\pgfqpoint{0.880687in}{1.058797in}}%
\pgfpathlineto{\pgfqpoint{0.864783in}{1.074021in}}%
\pgfpathlineto{\pgfqpoint{0.848891in}{1.089615in}}%
\pgfpathlineto{\pgfqpoint{0.840575in}{1.083666in}}%
\pgfpathlineto{\pgfqpoint{0.832227in}{1.078106in}}%
\pgfpathclose%
\pgfusepath{fill}%
\end{pgfscope}%
\begin{pgfscope}%
\pgfpathrectangle{\pgfqpoint{0.407592in}{0.222904in}}{\pgfqpoint{1.411163in}{1.411163in}}%
\pgfusepath{clip}%
\pgfsetbuttcap%
\pgfsetroundjoin%
\definecolor{currentfill}{rgb}{0.909804,0.698039,0.423529}%
\pgfsetfillcolor{currentfill}%
\pgfsetfillopacity{0.120000}%
\pgfsetlinewidth{0.000000pt}%
\definecolor{currentstroke}{rgb}{0.000000,0.000000,0.000000}%
\pgfsetstrokecolor{currentstroke}%
\pgfsetdash{}{0pt}%
\pgfpathmoveto{\pgfqpoint{1.439860in}{1.151675in}}%
\pgfpathlineto{\pgfqpoint{1.456418in}{1.154106in}}%
\pgfpathlineto{\pgfqpoint{1.473037in}{1.157082in}}%
\pgfpathlineto{\pgfqpoint{1.480601in}{1.170718in}}%
\pgfpathlineto{\pgfqpoint{1.463989in}{1.167851in}}%
\pgfpathlineto{\pgfqpoint{1.447439in}{1.165521in}}%
\pgfusepath{fill}%
\end{pgfscope}%
\begin{pgfscope}%
\pgfpathrectangle{\pgfqpoint{0.407592in}{0.222904in}}{\pgfqpoint{1.411163in}{1.411163in}}%
\pgfusepath{clip}%
\pgfsetbuttcap%
\pgfsetroundjoin%
\definecolor{currentfill}{rgb}{0.909804,0.698039,0.423529}%
\pgfsetfillcolor{currentfill}%
\pgfsetfillopacity{0.120000}%
\pgfsetlinewidth{0.000000pt}%
\definecolor{currentstroke}{rgb}{0.000000,0.000000,0.000000}%
\pgfsetstrokecolor{currentstroke}%
\pgfsetdash{}{0pt}%
\pgfpathmoveto{\pgfqpoint{1.168761in}{0.981802in}}%
\pgfpathlineto{\pgfqpoint{1.184904in}{0.977661in}}%
\pgfpathlineto{\pgfqpoint{1.201089in}{0.974421in}}%
\pgfpathlineto{\pgfqpoint{1.208928in}{0.985465in}}%
\pgfpathlineto{\pgfqpoint{1.216750in}{0.996863in}}%
\pgfpathlineto{\pgfqpoint{1.200583in}{1.000093in}}%
\pgfpathlineto{\pgfqpoint{1.184459in}{1.004188in}}%
\pgfpathlineto{\pgfqpoint{1.176619in}{0.992818in}}%
\pgfpathlineto{\pgfqpoint{1.168761in}{0.981802in}}%
\pgfpathclose%
\pgfusepath{fill}%
\end{pgfscope}%
\begin{pgfscope}%
\pgfpathrectangle{\pgfqpoint{0.407592in}{0.222904in}}{\pgfqpoint{1.411163in}{1.411163in}}%
\pgfusepath{clip}%
\pgfsetbuttcap%
\pgfsetroundjoin%
\definecolor{currentfill}{rgb}{0.909804,0.698039,0.423529}%
\pgfsetfillcolor{currentfill}%
\pgfsetfillopacity{0.120000}%
\pgfsetlinewidth{0.000000pt}%
\definecolor{currentstroke}{rgb}{0.000000,0.000000,0.000000}%
\pgfsetstrokecolor{currentstroke}%
\pgfsetdash{}{0pt}%
\pgfpathmoveto{\pgfqpoint{1.120777in}{0.972075in}}%
\pgfpathlineto{\pgfqpoint{1.136865in}{0.966042in}}%
\pgfpathlineto{\pgfqpoint{1.152990in}{0.960926in}}%
\pgfpathlineto{\pgfqpoint{1.160885in}{0.971164in}}%
\pgfpathlineto{\pgfqpoint{1.168761in}{0.981802in}}%
\pgfpathlineto{\pgfqpoint{1.152657in}{0.986840in}}%
\pgfpathlineto{\pgfqpoint{1.136591in}{0.992761in}}%
\pgfpathlineto{\pgfqpoint{1.128694in}{0.982219in}}%
\pgfpathlineto{\pgfqpoint{1.120777in}{0.972075in}}%
\pgfpathclose%
\pgfusepath{fill}%
\end{pgfscope}%
\begin{pgfscope}%
\pgfpathrectangle{\pgfqpoint{0.407592in}{0.222904in}}{\pgfqpoint{1.411163in}{1.411163in}}%
\pgfusepath{clip}%
\pgfsetbuttcap%
\pgfsetroundjoin%
\definecolor{currentfill}{rgb}{0.909804,0.698039,0.423529}%
\pgfsetfillcolor{currentfill}%
\pgfsetfillopacity{0.120000}%
\pgfsetlinewidth{0.000000pt}%
\definecolor{currentstroke}{rgb}{0.000000,0.000000,0.000000}%
\pgfsetstrokecolor{currentstroke}%
\pgfsetdash{}{0pt}%
\pgfpathmoveto{\pgfqpoint{1.376425in}{1.093971in}}%
\pgfpathlineto{\pgfqpoint{1.392893in}{1.095612in}}%
\pgfpathlineto{\pgfqpoint{1.409420in}{1.097874in}}%
\pgfpathlineto{\pgfqpoint{1.417049in}{1.111064in}}%
\pgfpathlineto{\pgfqpoint{1.424666in}{1.124435in}}%
\pgfpathlineto{\pgfqpoint{1.408154in}{1.122375in}}%
\pgfpathlineto{\pgfqpoint{1.391700in}{1.120917in}}%
\pgfpathlineto{\pgfqpoint{1.384069in}{1.107355in}}%
\pgfpathlineto{\pgfqpoint{1.376425in}{1.093971in}}%
\pgfpathclose%
\pgfusepath{fill}%
\end{pgfscope}%
\begin{pgfscope}%
\pgfpathrectangle{\pgfqpoint{0.407592in}{0.222904in}}{\pgfqpoint{1.411163in}{1.411163in}}%
\pgfusepath{clip}%
\pgfsetbuttcap%
\pgfsetroundjoin%
\definecolor{currentfill}{rgb}{0.909804,0.698039,0.423529}%
\pgfsetfillcolor{currentfill}%
\pgfsetfillopacity{0.120000}%
\pgfsetlinewidth{0.000000pt}%
\definecolor{currentstroke}{rgb}{0.000000,0.000000,0.000000}%
\pgfsetstrokecolor{currentstroke}%
\pgfsetdash{}{0pt}%
\pgfpathmoveto{\pgfqpoint{1.216750in}{0.996863in}}%
\pgfpathlineto{\pgfqpoint{1.232962in}{0.994494in}}%
\pgfpathlineto{\pgfqpoint{1.249222in}{0.992980in}}%
\pgfpathlineto{\pgfqpoint{1.257010in}{1.004665in}}%
\pgfpathlineto{\pgfqpoint{1.264783in}{1.016660in}}%
\pgfpathlineto{\pgfqpoint{1.248540in}{1.018236in}}%
\pgfpathlineto{\pgfqpoint{1.232345in}{1.020630in}}%
\pgfpathlineto{\pgfqpoint{1.224556in}{1.008592in}}%
\pgfpathlineto{\pgfqpoint{1.216750in}{0.996863in}}%
\pgfpathclose%
\pgfusepath{fill}%
\end{pgfscope}%
\begin{pgfscope}%
\pgfpathrectangle{\pgfqpoint{0.407592in}{0.222904in}}{\pgfqpoint{1.411163in}{1.411163in}}%
\pgfusepath{clip}%
\pgfsetbuttcap%
\pgfsetroundjoin%
\definecolor{currentfill}{rgb}{0.909804,0.698039,0.423529}%
\pgfsetfillcolor{currentfill}%
\pgfsetfillopacity{0.120000}%
\pgfsetlinewidth{0.000000pt}%
\definecolor{currentstroke}{rgb}{0.000000,0.000000,0.000000}%
\pgfsetstrokecolor{currentstroke}%
\pgfsetdash{}{0pt}%
\pgfpathmoveto{\pgfqpoint{1.072755in}{0.968084in}}%
\pgfpathlineto{\pgfqpoint{1.088804in}{0.960132in}}%
\pgfpathlineto{\pgfqpoint{1.104883in}{0.953073in}}%
\pgfpathlineto{\pgfqpoint{1.112840in}{0.962352in}}%
\pgfpathlineto{\pgfqpoint{1.120777in}{0.972075in}}%
\pgfpathlineto{\pgfqpoint{1.104722in}{0.979004in}}%
\pgfpathlineto{\pgfqpoint{1.088699in}{0.986799in}}%
\pgfpathlineto{\pgfqpoint{1.080738in}{0.977224in}}%
\pgfpathlineto{\pgfqpoint{1.072755in}{0.968084in}}%
\pgfpathclose%
\pgfusepath{fill}%
\end{pgfscope}%
\begin{pgfscope}%
\pgfpathrectangle{\pgfqpoint{0.407592in}{0.222904in}}{\pgfqpoint{1.411163in}{1.411163in}}%
\pgfusepath{clip}%
\pgfsetbuttcap%
\pgfsetroundjoin%
\definecolor{currentfill}{rgb}{0.909804,0.698039,0.423529}%
\pgfsetfillcolor{currentfill}%
\pgfsetfillopacity{0.120000}%
\pgfsetlinewidth{0.000000pt}%
\definecolor{currentstroke}{rgb}{0.000000,0.000000,0.000000}%
\pgfsetstrokecolor{currentstroke}%
\pgfsetdash{}{0pt}%
\pgfpathmoveto{\pgfqpoint{0.751666in}{1.134334in}}%
\pgfpathlineto{\pgfqpoint{0.767593in}{1.117388in}}%
\pgfpathlineto{\pgfqpoint{0.783531in}{1.100694in}}%
\pgfpathlineto{\pgfqpoint{0.791982in}{1.105411in}}%
\pgfpathlineto{\pgfqpoint{0.800401in}{1.110500in}}%
\pgfpathlineto{\pgfqpoint{0.784503in}{1.127127in}}%
\pgfpathlineto{\pgfqpoint{0.768617in}{1.144006in}}%
\pgfpathlineto{\pgfqpoint{0.760158in}{1.138995in}}%
\pgfpathlineto{\pgfqpoint{0.751666in}{1.134334in}}%
\pgfpathclose%
\pgfusepath{fill}%
\end{pgfscope}%
\begin{pgfscope}%
\pgfpathrectangle{\pgfqpoint{0.407592in}{0.222904in}}{\pgfqpoint{1.411163in}{1.411163in}}%
\pgfusepath{clip}%
\pgfsetbuttcap%
\pgfsetroundjoin%
\definecolor{currentfill}{rgb}{0.909804,0.698039,0.423529}%
\pgfsetfillcolor{currentfill}%
\pgfsetfillopacity{0.120000}%
\pgfsetlinewidth{0.000000pt}%
\definecolor{currentstroke}{rgb}{0.000000,0.000000,0.000000}%
\pgfsetstrokecolor{currentstroke}%
\pgfsetdash{}{0pt}%
\pgfpathmoveto{\pgfqpoint{0.944440in}{1.002367in}}%
\pgfpathlineto{\pgfqpoint{0.960416in}{0.989596in}}%
\pgfpathlineto{\pgfqpoint{0.976410in}{0.977458in}}%
\pgfpathlineto{\pgfqpoint{0.984552in}{0.984402in}}%
\pgfpathlineto{\pgfqpoint{0.992667in}{0.991816in}}%
\pgfpathlineto{\pgfqpoint{0.976707in}{1.003795in}}%
\pgfpathlineto{\pgfqpoint{0.960764in}{1.016402in}}%
\pgfpathlineto{\pgfqpoint{0.952616in}{1.009161in}}%
\pgfpathlineto{\pgfqpoint{0.944440in}{1.002367in}}%
\pgfpathclose%
\pgfusepath{fill}%
\end{pgfscope}%
\begin{pgfscope}%
\pgfpathrectangle{\pgfqpoint{0.407592in}{0.222904in}}{\pgfqpoint{1.411163in}{1.411163in}}%
\pgfusepath{clip}%
\pgfsetbuttcap%
\pgfsetroundjoin%
\definecolor{currentfill}{rgb}{0.909804,0.698039,0.423529}%
\pgfsetfillcolor{currentfill}%
\pgfsetfillopacity{0.120000}%
\pgfsetlinewidth{0.000000pt}%
\definecolor{currentstroke}{rgb}{0.000000,0.000000,0.000000}%
\pgfsetstrokecolor{currentstroke}%
\pgfsetdash{}{0pt}%
\pgfpathmoveto{\pgfqpoint{1.264783in}{1.016660in}}%
\pgfpathlineto{\pgfqpoint{1.281076in}{1.015889in}}%
\pgfpathlineto{\pgfqpoint{1.297422in}{1.015905in}}%
\pgfpathlineto{\pgfqpoint{1.305163in}{1.028078in}}%
\pgfpathlineto{\pgfqpoint{1.312891in}{1.040522in}}%
\pgfpathlineto{\pgfqpoint{1.296561in}{1.040631in}}%
\pgfpathlineto{\pgfqpoint{1.280284in}{1.041495in}}%
\pgfpathlineto{\pgfqpoint{1.272541in}{1.028944in}}%
\pgfpathlineto{\pgfqpoint{1.264783in}{1.016660in}}%
\pgfpathclose%
\pgfusepath{fill}%
\end{pgfscope}%
\begin{pgfscope}%
\pgfpathrectangle{\pgfqpoint{0.407592in}{0.222904in}}{\pgfqpoint{1.411163in}{1.411163in}}%
\pgfusepath{clip}%
\pgfsetbuttcap%
\pgfsetroundjoin%
\definecolor{currentfill}{rgb}{0.909804,0.698039,0.423529}%
\pgfsetfillcolor{currentfill}%
\pgfsetfillopacity{0.120000}%
\pgfsetlinewidth{0.000000pt}%
\definecolor{currentstroke}{rgb}{0.000000,0.000000,0.000000}%
\pgfsetstrokecolor{currentstroke}%
\pgfsetdash{}{0pt}%
\pgfpathmoveto{\pgfqpoint{0.864100in}{1.047056in}}%
\pgfpathlineto{\pgfqpoint{0.880055in}{1.032120in}}%
\pgfpathlineto{\pgfqpoint{0.896023in}{1.017630in}}%
\pgfpathlineto{\pgfqpoint{0.904294in}{1.023398in}}%
\pgfpathlineto{\pgfqpoint{0.912535in}{1.029605in}}%
\pgfpathlineto{\pgfqpoint{0.896604in}{1.043978in}}%
\pgfpathlineto{\pgfqpoint{0.880687in}{1.058797in}}%
\pgfpathlineto{\pgfqpoint{0.872409in}{1.052719in}}%
\pgfpathlineto{\pgfqpoint{0.864100in}{1.047056in}}%
\pgfpathclose%
\pgfusepath{fill}%
\end{pgfscope}%
\begin{pgfscope}%
\pgfpathrectangle{\pgfqpoint{0.407592in}{0.222904in}}{\pgfqpoint{1.411163in}{1.411163in}}%
\pgfusepath{clip}%
\pgfsetbuttcap%
\pgfsetroundjoin%
\definecolor{currentfill}{rgb}{0.909804,0.698039,0.423529}%
\pgfsetfillcolor{currentfill}%
\pgfsetfillopacity{0.120000}%
\pgfsetlinewidth{0.000000pt}%
\definecolor{currentstroke}{rgb}{0.000000,0.000000,0.000000}%
\pgfsetstrokecolor{currentstroke}%
\pgfsetdash{}{0pt}%
\pgfpathmoveto{\pgfqpoint{0.638978in}{1.232091in}}%
\pgfpathlineto{\pgfqpoint{0.654888in}{1.214004in}}%
\pgfpathlineto{\pgfqpoint{0.670808in}{1.196064in}}%
\pgfpathlineto{\pgfqpoint{0.679450in}{1.199985in}}%
\pgfpathlineto{\pgfqpoint{0.688058in}{1.204219in}}%
\pgfpathlineto{\pgfqpoint{0.672180in}{1.222126in}}%
\pgfpathlineto{\pgfqpoint{0.656311in}{1.240179in}}%
\pgfpathlineto{\pgfqpoint{0.647662in}{1.235987in}}%
\pgfpathlineto{\pgfqpoint{0.638978in}{1.232091in}}%
\pgfpathclose%
\pgfusepath{fill}%
\end{pgfscope}%
\begin{pgfscope}%
\pgfpathrectangle{\pgfqpoint{0.407592in}{0.222904in}}{\pgfqpoint{1.411163in}{1.411163in}}%
\pgfusepath{clip}%
\pgfsetbuttcap%
\pgfsetroundjoin%
\definecolor{currentfill}{rgb}{0.909804,0.698039,0.423529}%
\pgfsetfillcolor{currentfill}%
\pgfsetfillopacity{0.120000}%
\pgfsetlinewidth{0.000000pt}%
\definecolor{currentstroke}{rgb}{0.000000,0.000000,0.000000}%
\pgfsetstrokecolor{currentstroke}%
\pgfsetdash{}{0pt}%
\pgfpathmoveto{\pgfqpoint{1.024648in}{0.969945in}}%
\pgfpathlineto{\pgfqpoint{1.040672in}{0.960155in}}%
\pgfpathlineto{\pgfqpoint{1.056720in}{0.951194in}}%
\pgfpathlineto{\pgfqpoint{1.064749in}{0.959401in}}%
\pgfpathlineto{\pgfqpoint{1.072755in}{0.968084in}}%
\pgfpathlineto{\pgfqpoint{1.056734in}{0.976891in}}%
\pgfpathlineto{\pgfqpoint{1.040740in}{0.986511in}}%
\pgfpathlineto{\pgfqpoint{1.032706in}{0.977997in}}%
\pgfpathlineto{\pgfqpoint{1.024648in}{0.969945in}}%
\pgfpathclose%
\pgfusepath{fill}%
\end{pgfscope}%
\begin{pgfscope}%
\pgfpathrectangle{\pgfqpoint{0.407592in}{0.222904in}}{\pgfqpoint{1.411163in}{1.411163in}}%
\pgfusepath{clip}%
\pgfsetbuttcap%
\pgfsetroundjoin%
\definecolor{currentfill}{rgb}{0.909804,0.698039,0.423529}%
\pgfsetfillcolor{currentfill}%
\pgfsetfillopacity{0.120000}%
\pgfsetlinewidth{0.000000pt}%
\definecolor{currentstroke}{rgb}{0.000000,0.000000,0.000000}%
\pgfsetstrokecolor{currentstroke}%
\pgfsetdash{}{0pt}%
\pgfpathmoveto{\pgfqpoint{1.424666in}{1.124435in}}%
\pgfpathlineto{\pgfqpoint{1.441239in}{1.127076in}}%
\pgfpathlineto{\pgfqpoint{1.457873in}{1.130277in}}%
\pgfpathlineto{\pgfqpoint{1.465462in}{1.143598in}}%
\pgfpathlineto{\pgfqpoint{1.473037in}{1.157082in}}%
\pgfpathlineto{\pgfqpoint{1.456418in}{1.154106in}}%
\pgfpathlineto{\pgfqpoint{1.439860in}{1.151675in}}%
\pgfpathlineto{\pgfqpoint{1.432269in}{1.137976in}}%
\pgfpathlineto{\pgfqpoint{1.424666in}{1.124435in}}%
\pgfpathclose%
\pgfusepath{fill}%
\end{pgfscope}%
\begin{pgfscope}%
\pgfpathrectangle{\pgfqpoint{0.407592in}{0.222904in}}{\pgfqpoint{1.411163in}{1.411163in}}%
\pgfusepath{clip}%
\pgfsetbuttcap%
\pgfsetroundjoin%
\definecolor{currentfill}{rgb}{0.909804,0.698039,0.423529}%
\pgfsetfillcolor{currentfill}%
\pgfsetfillopacity{0.120000}%
\pgfsetlinewidth{0.000000pt}%
\definecolor{currentstroke}{rgb}{0.000000,0.000000,0.000000}%
\pgfsetstrokecolor{currentstroke}%
\pgfsetdash{}{0pt}%
\pgfpathmoveto{\pgfqpoint{1.312891in}{1.040522in}}%
\pgfpathlineto{\pgfqpoint{1.329275in}{1.041149in}}%
\pgfpathlineto{\pgfqpoint{1.345716in}{1.042490in}}%
\pgfpathlineto{\pgfqpoint{1.353413in}{1.055022in}}%
\pgfpathlineto{\pgfqpoint{1.361097in}{1.067790in}}%
\pgfpathlineto{\pgfqpoint{1.344672in}{1.066624in}}%
\pgfpathlineto{\pgfqpoint{1.328303in}{1.066145in}}%
\pgfpathlineto{\pgfqpoint{1.320604in}{1.053217in}}%
\pgfpathlineto{\pgfqpoint{1.312891in}{1.040522in}}%
\pgfpathclose%
\pgfusepath{fill}%
\end{pgfscope}%
\begin{pgfscope}%
\pgfpathrectangle{\pgfqpoint{0.407592in}{0.222904in}}{\pgfqpoint{1.411163in}{1.411163in}}%
\pgfusepath{clip}%
\pgfsetbuttcap%
\pgfsetroundjoin%
\definecolor{currentfill}{rgb}{0.909804,0.698039,0.423529}%
\pgfsetfillcolor{currentfill}%
\pgfsetfillopacity{0.120000}%
\pgfsetlinewidth{0.000000pt}%
\definecolor{currentstroke}{rgb}{0.000000,0.000000,0.000000}%
\pgfsetstrokecolor{currentstroke}%
\pgfsetdash{}{0pt}%
\pgfpathmoveto{\pgfqpoint{0.783531in}{1.100694in}}%
\pgfpathlineto{\pgfqpoint{0.799478in}{1.084278in}}%
\pgfpathlineto{\pgfqpoint{0.815437in}{1.068166in}}%
\pgfpathlineto{\pgfqpoint{0.823848in}{1.072938in}}%
\pgfpathlineto{\pgfqpoint{0.832227in}{1.078106in}}%
\pgfpathlineto{\pgfqpoint{0.816308in}{1.094150in}}%
\pgfpathlineto{\pgfqpoint{0.800401in}{1.110500in}}%
\pgfpathlineto{\pgfqpoint{0.791982in}{1.105411in}}%
\pgfpathlineto{\pgfqpoint{0.783531in}{1.100694in}}%
\pgfpathclose%
\pgfusepath{fill}%
\end{pgfscope}%
\begin{pgfscope}%
\pgfpathrectangle{\pgfqpoint{0.407592in}{0.222904in}}{\pgfqpoint{1.411163in}{1.411163in}}%
\pgfusepath{clip}%
\pgfsetbuttcap%
\pgfsetroundjoin%
\definecolor{currentfill}{rgb}{0.909804,0.698039,0.423529}%
\pgfsetfillcolor{currentfill}%
\pgfsetfillopacity{0.120000}%
\pgfsetlinewidth{0.000000pt}%
\definecolor{currentstroke}{rgb}{0.000000,0.000000,0.000000}%
\pgfsetstrokecolor{currentstroke}%
\pgfsetdash{}{0pt}%
\pgfpathmoveto{\pgfqpoint{0.670808in}{1.196064in}}%
\pgfpathlineto{\pgfqpoint{0.686737in}{1.178284in}}%
\pgfpathlineto{\pgfqpoint{0.702675in}{1.160679in}}%
\pgfpathlineto{\pgfqpoint{0.711276in}{1.164621in}}%
\pgfpathlineto{\pgfqpoint{0.719843in}{1.168894in}}%
\pgfpathlineto{\pgfqpoint{0.703946in}{1.186469in}}%
\pgfpathlineto{\pgfqpoint{0.688058in}{1.204219in}}%
\pgfpathlineto{\pgfqpoint{0.679450in}{1.199985in}}%
\pgfpathlineto{\pgfqpoint{0.670808in}{1.196064in}}%
\pgfpathclose%
\pgfusepath{fill}%
\end{pgfscope}%
\begin{pgfscope}%
\pgfpathrectangle{\pgfqpoint{0.407592in}{0.222904in}}{\pgfqpoint{1.411163in}{1.411163in}}%
\pgfusepath{clip}%
\pgfsetbuttcap%
\pgfsetroundjoin%
\definecolor{currentfill}{rgb}{0.909804,0.698039,0.423529}%
\pgfsetfillcolor{currentfill}%
\pgfsetfillopacity{0.120000}%
\pgfsetlinewidth{0.000000pt}%
\definecolor{currentstroke}{rgb}{0.000000,0.000000,0.000000}%
\pgfsetstrokecolor{currentstroke}%
\pgfsetdash{}{0pt}%
\pgfpathmoveto{\pgfqpoint{0.896023in}{1.017630in}}%
\pgfpathlineto{\pgfqpoint{0.912005in}{1.003627in}}%
\pgfpathlineto{\pgfqpoint{0.928001in}{0.990159in}}%
\pgfpathlineto{\pgfqpoint{0.936235in}{0.996031in}}%
\pgfpathlineto{\pgfqpoint{0.944440in}{1.002367in}}%
\pgfpathlineto{\pgfqpoint{0.928480in}{1.015719in}}%
\pgfpathlineto{\pgfqpoint{0.912535in}{1.029605in}}%
\pgfpathlineto{\pgfqpoint{0.904294in}{1.023398in}}%
\pgfpathlineto{\pgfqpoint{0.896023in}{1.017630in}}%
\pgfpathclose%
\pgfusepath{fill}%
\end{pgfscope}%
\begin{pgfscope}%
\pgfpathrectangle{\pgfqpoint{0.407592in}{0.222904in}}{\pgfqpoint{1.411163in}{1.411163in}}%
\pgfusepath{clip}%
\pgfsetbuttcap%
\pgfsetroundjoin%
\definecolor{currentfill}{rgb}{0.909804,0.698039,0.423529}%
\pgfsetfillcolor{currentfill}%
\pgfsetfillopacity{0.120000}%
\pgfsetlinewidth{0.000000pt}%
\definecolor{currentstroke}{rgb}{0.000000,0.000000,0.000000}%
\pgfsetstrokecolor{currentstroke}%
\pgfsetdash{}{0pt}%
\pgfpathmoveto{\pgfqpoint{0.976410in}{0.977458in}}%
\pgfpathlineto{\pgfqpoint{0.992422in}{0.966004in}}%
\pgfpathlineto{\pgfqpoint{1.008455in}{0.955288in}}%
\pgfpathlineto{\pgfqpoint{1.016564in}{0.962371in}}%
\pgfpathlineto{\pgfqpoint{1.024648in}{0.969945in}}%
\pgfpathlineto{\pgfqpoint{1.008647in}{0.980516in}}%
\pgfpathlineto{\pgfqpoint{0.992667in}{0.991816in}}%
\pgfpathlineto{\pgfqpoint{0.984552in}{0.984402in}}%
\pgfpathlineto{\pgfqpoint{0.976410in}{0.977458in}}%
\pgfpathclose%
\pgfusepath{fill}%
\end{pgfscope}%
\begin{pgfscope}%
\pgfpathrectangle{\pgfqpoint{0.407592in}{0.222904in}}{\pgfqpoint{1.411163in}{1.411163in}}%
\pgfusepath{clip}%
\pgfsetbuttcap%
\pgfsetroundjoin%
\definecolor{currentfill}{rgb}{0.909804,0.698039,0.423529}%
\pgfsetfillcolor{currentfill}%
\pgfsetfillopacity{0.120000}%
\pgfsetlinewidth{0.000000pt}%
\definecolor{currentstroke}{rgb}{0.000000,0.000000,0.000000}%
\pgfsetstrokecolor{currentstroke}%
\pgfsetdash{}{0pt}%
\pgfpathmoveto{\pgfqpoint{1.361097in}{1.067790in}}%
\pgfpathlineto{\pgfqpoint{1.377580in}{1.069621in}}%
\pgfpathlineto{\pgfqpoint{1.394123in}{1.072094in}}%
\pgfpathlineto{\pgfqpoint{1.401778in}{1.084880in}}%
\pgfpathlineto{\pgfqpoint{1.409420in}{1.097874in}}%
\pgfpathlineto{\pgfqpoint{1.392893in}{1.095612in}}%
\pgfpathlineto{\pgfqpoint{1.376425in}{1.093971in}}%
\pgfpathlineto{\pgfqpoint{1.368768in}{1.080778in}}%
\pgfpathlineto{\pgfqpoint{1.361097in}{1.067790in}}%
\pgfpathclose%
\pgfusepath{fill}%
\end{pgfscope}%
\begin{pgfscope}%
\pgfpathrectangle{\pgfqpoint{0.407592in}{0.222904in}}{\pgfqpoint{1.411163in}{1.411163in}}%
\pgfusepath{clip}%
\pgfsetbuttcap%
\pgfsetroundjoin%
\definecolor{currentfill}{rgb}{0.909804,0.698039,0.423529}%
\pgfsetfillcolor{currentfill}%
\pgfsetfillopacity{0.120000}%
\pgfsetlinewidth{0.000000pt}%
\definecolor{currentstroke}{rgb}{0.000000,0.000000,0.000000}%
\pgfsetstrokecolor{currentstroke}%
\pgfsetdash{}{0pt}%
\pgfpathmoveto{\pgfqpoint{1.152990in}{0.960926in}}%
\pgfpathlineto{\pgfqpoint{1.169154in}{0.956740in}}%
\pgfpathlineto{\pgfqpoint{1.185359in}{0.953490in}}%
\pgfpathlineto{\pgfqpoint{1.193233in}{0.963755in}}%
\pgfpathlineto{\pgfqpoint{1.201089in}{0.974421in}}%
\pgfpathlineto{\pgfqpoint{1.184904in}{0.977661in}}%
\pgfpathlineto{\pgfqpoint{1.168761in}{0.981802in}}%
\pgfpathlineto{\pgfqpoint{1.160885in}{0.971164in}}%
\pgfpathlineto{\pgfqpoint{1.152990in}{0.960926in}}%
\pgfpathclose%
\pgfusepath{fill}%
\end{pgfscope}%
\begin{pgfscope}%
\pgfpathrectangle{\pgfqpoint{0.407592in}{0.222904in}}{\pgfqpoint{1.411163in}{1.411163in}}%
\pgfusepath{clip}%
\pgfsetbuttcap%
\pgfsetroundjoin%
\definecolor{currentfill}{rgb}{0.909804,0.698039,0.423529}%
\pgfsetfillcolor{currentfill}%
\pgfsetfillopacity{0.120000}%
\pgfsetlinewidth{0.000000pt}%
\definecolor{currentstroke}{rgb}{0.000000,0.000000,0.000000}%
\pgfsetstrokecolor{currentstroke}%
\pgfsetdash{}{0pt}%
\pgfpathmoveto{\pgfqpoint{0.815437in}{1.068166in}}%
\pgfpathlineto{\pgfqpoint{0.831407in}{1.052391in}}%
\pgfpathlineto{\pgfqpoint{0.847388in}{1.036985in}}%
\pgfpathlineto{\pgfqpoint{0.855760in}{1.041810in}}%
\pgfpathlineto{\pgfqpoint{0.864100in}{1.047056in}}%
\pgfpathlineto{\pgfqpoint{0.848158in}{1.062396in}}%
\pgfpathlineto{\pgfqpoint{0.832227in}{1.078106in}}%
\pgfpathlineto{\pgfqpoint{0.823848in}{1.072938in}}%
\pgfpathlineto{\pgfqpoint{0.815437in}{1.068166in}}%
\pgfpathclose%
\pgfusepath{fill}%
\end{pgfscope}%
\begin{pgfscope}%
\pgfpathrectangle{\pgfqpoint{0.407592in}{0.222904in}}{\pgfqpoint{1.411163in}{1.411163in}}%
\pgfusepath{clip}%
\pgfsetbuttcap%
\pgfsetroundjoin%
\definecolor{currentfill}{rgb}{0.909804,0.698039,0.423529}%
\pgfsetfillcolor{currentfill}%
\pgfsetfillopacity{0.120000}%
\pgfsetlinewidth{0.000000pt}%
\definecolor{currentstroke}{rgb}{0.000000,0.000000,0.000000}%
\pgfsetstrokecolor{currentstroke}%
\pgfsetdash{}{0pt}%
\pgfpathmoveto{\pgfqpoint{0.702675in}{1.160679in}}%
\pgfpathlineto{\pgfqpoint{0.718623in}{1.143267in}}%
\pgfpathlineto{\pgfqpoint{0.734581in}{1.126066in}}%
\pgfpathlineto{\pgfqpoint{0.743141in}{1.130024in}}%
\pgfpathlineto{\pgfqpoint{0.751666in}{1.134334in}}%
\pgfpathlineto{\pgfqpoint{0.735750in}{1.151509in}}%
\pgfpathlineto{\pgfqpoint{0.719843in}{1.168894in}}%
\pgfpathlineto{\pgfqpoint{0.711276in}{1.164621in}}%
\pgfpathlineto{\pgfqpoint{0.702675in}{1.160679in}}%
\pgfpathclose%
\pgfusepath{fill}%
\end{pgfscope}%
\begin{pgfscope}%
\pgfpathrectangle{\pgfqpoint{0.407592in}{0.222904in}}{\pgfqpoint{1.411163in}{1.411163in}}%
\pgfusepath{clip}%
\pgfsetbuttcap%
\pgfsetroundjoin%
\definecolor{currentfill}{rgb}{0.909804,0.698039,0.423529}%
\pgfsetfillcolor{currentfill}%
\pgfsetfillopacity{0.120000}%
\pgfsetlinewidth{0.000000pt}%
\definecolor{currentstroke}{rgb}{0.000000,0.000000,0.000000}%
\pgfsetstrokecolor{currentstroke}%
\pgfsetdash{}{0pt}%
\pgfpathmoveto{\pgfqpoint{1.201089in}{0.974421in}}%
\pgfpathlineto{\pgfqpoint{1.217319in}{0.972078in}}%
\pgfpathlineto{\pgfqpoint{1.233597in}{0.970624in}}%
\pgfpathlineto{\pgfqpoint{1.241417in}{0.981625in}}%
\pgfpathlineto{\pgfqpoint{1.249222in}{0.992980in}}%
\pgfpathlineto{\pgfqpoint{1.232962in}{0.994494in}}%
\pgfpathlineto{\pgfqpoint{1.216750in}{0.996863in}}%
\pgfpathlineto{\pgfqpoint{1.208928in}{0.985465in}}%
\pgfpathlineto{\pgfqpoint{1.201089in}{0.974421in}}%
\pgfpathclose%
\pgfusepath{fill}%
\end{pgfscope}%
\begin{pgfscope}%
\pgfpathrectangle{\pgfqpoint{0.407592in}{0.222904in}}{\pgfqpoint{1.411163in}{1.411163in}}%
\pgfusepath{clip}%
\pgfsetbuttcap%
\pgfsetroundjoin%
\definecolor{currentfill}{rgb}{0.909804,0.698039,0.423529}%
\pgfsetfillcolor{currentfill}%
\pgfsetfillopacity{0.120000}%
\pgfsetlinewidth{0.000000pt}%
\definecolor{currentstroke}{rgb}{0.000000,0.000000,0.000000}%
\pgfsetstrokecolor{currentstroke}%
\pgfsetdash{}{0pt}%
\pgfpathmoveto{\pgfqpoint{1.104883in}{0.953073in}}%
\pgfpathlineto{\pgfqpoint{1.120995in}{0.946937in}}%
\pgfpathlineto{\pgfqpoint{1.137142in}{0.941748in}}%
\pgfpathlineto{\pgfqpoint{1.145076in}{0.951113in}}%
\pgfpathlineto{\pgfqpoint{1.152990in}{0.960926in}}%
\pgfpathlineto{\pgfqpoint{1.136865in}{0.966042in}}%
\pgfpathlineto{\pgfqpoint{1.120777in}{0.972075in}}%
\pgfpathlineto{\pgfqpoint{1.112840in}{0.962352in}}%
\pgfpathlineto{\pgfqpoint{1.104883in}{0.953073in}}%
\pgfpathclose%
\pgfusepath{fill}%
\end{pgfscope}%
\begin{pgfscope}%
\pgfpathrectangle{\pgfqpoint{0.407592in}{0.222904in}}{\pgfqpoint{1.411163in}{1.411163in}}%
\pgfusepath{clip}%
\pgfsetbuttcap%
\pgfsetroundjoin%
\definecolor{currentfill}{rgb}{0.909804,0.698039,0.423529}%
\pgfsetfillcolor{currentfill}%
\pgfsetfillopacity{0.120000}%
\pgfsetlinewidth{0.000000pt}%
\definecolor{currentstroke}{rgb}{0.000000,0.000000,0.000000}%
\pgfsetstrokecolor{currentstroke}%
\pgfsetdash{}{0pt}%
\pgfpathmoveto{\pgfqpoint{1.473037in}{1.157082in}}%
\pgfpathlineto{\pgfqpoint{1.489719in}{1.160582in}}%
\pgfpathlineto{\pgfqpoint{1.506465in}{1.164588in}}%
\pgfpathlineto{\pgfqpoint{1.514013in}{1.177989in}}%
\pgfpathlineto{\pgfqpoint{1.497275in}{1.174104in}}%
\pgfpathlineto{\pgfqpoint{1.480601in}{1.170718in}}%
\pgfusepath{fill}%
\end{pgfscope}%
\begin{pgfscope}%
\pgfpathrectangle{\pgfqpoint{0.407592in}{0.222904in}}{\pgfqpoint{1.411163in}{1.411163in}}%
\pgfusepath{clip}%
\pgfsetbuttcap%
\pgfsetroundjoin%
\definecolor{currentfill}{rgb}{0.909804,0.698039,0.423529}%
\pgfsetfillcolor{currentfill}%
\pgfsetfillopacity{0.120000}%
\pgfsetlinewidth{0.000000pt}%
\definecolor{currentstroke}{rgb}{0.000000,0.000000,0.000000}%
\pgfsetstrokecolor{currentstroke}%
\pgfsetdash{}{0pt}%
\pgfpathmoveto{\pgfqpoint{1.249222in}{0.992980in}}%
\pgfpathlineto{\pgfqpoint{1.265531in}{0.992303in}}%
\pgfpathlineto{\pgfqpoint{1.281894in}{0.992444in}}%
\pgfpathlineto{\pgfqpoint{1.289665in}{1.004020in}}%
\pgfpathlineto{\pgfqpoint{1.297422in}{1.015905in}}%
\pgfpathlineto{\pgfqpoint{1.281076in}{1.015889in}}%
\pgfpathlineto{\pgfqpoint{1.264783in}{1.016660in}}%
\pgfpathlineto{\pgfqpoint{1.257010in}{1.004665in}}%
\pgfpathlineto{\pgfqpoint{1.249222in}{0.992980in}}%
\pgfpathclose%
\pgfusepath{fill}%
\end{pgfscope}%
\begin{pgfscope}%
\pgfpathrectangle{\pgfqpoint{0.407592in}{0.222904in}}{\pgfqpoint{1.411163in}{1.411163in}}%
\pgfusepath{clip}%
\pgfsetbuttcap%
\pgfsetroundjoin%
\definecolor{currentfill}{rgb}{0.909804,0.698039,0.423529}%
\pgfsetfillcolor{currentfill}%
\pgfsetfillopacity{0.120000}%
\pgfsetlinewidth{0.000000pt}%
\definecolor{currentstroke}{rgb}{0.000000,0.000000,0.000000}%
\pgfsetstrokecolor{currentstroke}%
\pgfsetdash{}{0pt}%
\pgfpathmoveto{\pgfqpoint{1.056720in}{0.951194in}}%
\pgfpathlineto{\pgfqpoint{1.072797in}{0.943107in}}%
\pgfpathlineto{\pgfqpoint{1.088902in}{0.935933in}}%
\pgfpathlineto{\pgfqpoint{1.096904in}{0.944260in}}%
\pgfpathlineto{\pgfqpoint{1.104883in}{0.953073in}}%
\pgfpathlineto{\pgfqpoint{1.088804in}{0.960132in}}%
\pgfpathlineto{\pgfqpoint{1.072755in}{0.968084in}}%
\pgfpathlineto{\pgfqpoint{1.064749in}{0.959401in}}%
\pgfpathlineto{\pgfqpoint{1.056720in}{0.951194in}}%
\pgfpathclose%
\pgfusepath{fill}%
\end{pgfscope}%
\begin{pgfscope}%
\pgfpathrectangle{\pgfqpoint{0.407592in}{0.222904in}}{\pgfqpoint{1.411163in}{1.411163in}}%
\pgfusepath{clip}%
\pgfsetbuttcap%
\pgfsetroundjoin%
\definecolor{currentfill}{rgb}{0.909804,0.698039,0.423529}%
\pgfsetfillcolor{currentfill}%
\pgfsetfillopacity{0.120000}%
\pgfsetlinewidth{0.000000pt}%
\definecolor{currentstroke}{rgb}{0.000000,0.000000,0.000000}%
\pgfsetstrokecolor{currentstroke}%
\pgfsetdash{}{0pt}%
\pgfpathmoveto{\pgfqpoint{1.409420in}{1.097874in}}%
\pgfpathlineto{\pgfqpoint{1.426008in}{1.100735in}}%
\pgfpathlineto{\pgfqpoint{1.442659in}{1.104169in}}%
\pgfpathlineto{\pgfqpoint{1.450272in}{1.117130in}}%
\pgfpathlineto{\pgfqpoint{1.457873in}{1.130277in}}%
\pgfpathlineto{\pgfqpoint{1.441239in}{1.127076in}}%
\pgfpathlineto{\pgfqpoint{1.424666in}{1.124435in}}%
\pgfpathlineto{\pgfqpoint{1.417049in}{1.111064in}}%
\pgfpathlineto{\pgfqpoint{1.409420in}{1.097874in}}%
\pgfpathclose%
\pgfusepath{fill}%
\end{pgfscope}%
\begin{pgfscope}%
\pgfpathrectangle{\pgfqpoint{0.407592in}{0.222904in}}{\pgfqpoint{1.411163in}{1.411163in}}%
\pgfusepath{clip}%
\pgfsetbuttcap%
\pgfsetroundjoin%
\definecolor{currentfill}{rgb}{0.909804,0.698039,0.423529}%
\pgfsetfillcolor{currentfill}%
\pgfsetfillopacity{0.120000}%
\pgfsetlinewidth{0.000000pt}%
\definecolor{currentstroke}{rgb}{0.000000,0.000000,0.000000}%
\pgfsetstrokecolor{currentstroke}%
\pgfsetdash{}{0pt}%
\pgfpathmoveto{\pgfqpoint{0.928001in}{0.990159in}}%
\pgfpathlineto{\pgfqpoint{0.944013in}{0.977275in}}%
\pgfpathlineto{\pgfqpoint{0.960042in}{0.965026in}}%
\pgfpathlineto{\pgfqpoint{0.968240in}{0.970996in}}%
\pgfpathlineto{\pgfqpoint{0.976410in}{0.977458in}}%
\pgfpathlineto{\pgfqpoint{0.960416in}{0.989596in}}%
\pgfpathlineto{\pgfqpoint{0.944440in}{1.002367in}}%
\pgfpathlineto{\pgfqpoint{0.936235in}{0.996031in}}%
\pgfpathlineto{\pgfqpoint{0.928001in}{0.990159in}}%
\pgfpathclose%
\pgfusepath{fill}%
\end{pgfscope}%
\begin{pgfscope}%
\pgfpathrectangle{\pgfqpoint{0.407592in}{0.222904in}}{\pgfqpoint{1.411163in}{1.411163in}}%
\pgfusepath{clip}%
\pgfsetbuttcap%
\pgfsetroundjoin%
\definecolor{currentfill}{rgb}{0.909804,0.698039,0.423529}%
\pgfsetfillcolor{currentfill}%
\pgfsetfillopacity{0.120000}%
\pgfsetlinewidth{0.000000pt}%
\definecolor{currentstroke}{rgb}{0.000000,0.000000,0.000000}%
\pgfsetstrokecolor{currentstroke}%
\pgfsetdash{}{0pt}%
\pgfpathmoveto{\pgfqpoint{0.734581in}{1.126066in}}%
\pgfpathlineto{\pgfqpoint{0.750548in}{1.109096in}}%
\pgfpathlineto{\pgfqpoint{0.766526in}{1.092380in}}%
\pgfpathlineto{\pgfqpoint{0.775045in}{1.096350in}}%
\pgfpathlineto{\pgfqpoint{0.783531in}{1.100694in}}%
\pgfpathlineto{\pgfqpoint{0.767593in}{1.117388in}}%
\pgfpathlineto{\pgfqpoint{0.751666in}{1.134334in}}%
\pgfpathlineto{\pgfqpoint{0.743141in}{1.130024in}}%
\pgfpathlineto{\pgfqpoint{0.734581in}{1.126066in}}%
\pgfpathclose%
\pgfusepath{fill}%
\end{pgfscope}%
\begin{pgfscope}%
\pgfpathrectangle{\pgfqpoint{0.407592in}{0.222904in}}{\pgfqpoint{1.411163in}{1.411163in}}%
\pgfusepath{clip}%
\pgfsetbuttcap%
\pgfsetroundjoin%
\definecolor{currentfill}{rgb}{0.909804,0.698039,0.423529}%
\pgfsetfillcolor{currentfill}%
\pgfsetfillopacity{0.120000}%
\pgfsetlinewidth{0.000000pt}%
\definecolor{currentstroke}{rgb}{0.000000,0.000000,0.000000}%
\pgfsetstrokecolor{currentstroke}%
\pgfsetdash{}{0pt}%
\pgfpathmoveto{\pgfqpoint{1.297422in}{1.015905in}}%
\pgfpathlineto{\pgfqpoint{1.313822in}{1.016685in}}%
\pgfpathlineto{\pgfqpoint{1.330278in}{1.018203in}}%
\pgfpathlineto{\pgfqpoint{1.338004in}{1.030211in}}%
\pgfpathlineto{\pgfqpoint{1.345716in}{1.042490in}}%
\pgfpathlineto{\pgfqpoint{1.329275in}{1.041149in}}%
\pgfpathlineto{\pgfqpoint{1.312891in}{1.040522in}}%
\pgfpathlineto{\pgfqpoint{1.305163in}{1.028078in}}%
\pgfpathlineto{\pgfqpoint{1.297422in}{1.015905in}}%
\pgfpathclose%
\pgfusepath{fill}%
\end{pgfscope}%
\begin{pgfscope}%
\pgfpathrectangle{\pgfqpoint{0.407592in}{0.222904in}}{\pgfqpoint{1.411163in}{1.411163in}}%
\pgfusepath{clip}%
\pgfsetbuttcap%
\pgfsetroundjoin%
\definecolor{currentfill}{rgb}{0.909804,0.698039,0.423529}%
\pgfsetfillcolor{currentfill}%
\pgfsetfillopacity{0.120000}%
\pgfsetlinewidth{0.000000pt}%
\definecolor{currentstroke}{rgb}{0.000000,0.000000,0.000000}%
\pgfsetstrokecolor{currentstroke}%
\pgfsetdash{}{0pt}%
\pgfpathmoveto{\pgfqpoint{0.847388in}{1.036985in}}%
\pgfpathlineto{\pgfqpoint{0.863382in}{1.021984in}}%
\pgfpathlineto{\pgfqpoint{0.879388in}{1.007429in}}%
\pgfpathlineto{\pgfqpoint{0.887721in}{1.012306in}}%
\pgfpathlineto{\pgfqpoint{0.896023in}{1.017630in}}%
\pgfpathlineto{\pgfqpoint{0.880055in}{1.032120in}}%
\pgfpathlineto{\pgfqpoint{0.864100in}{1.047056in}}%
\pgfpathlineto{\pgfqpoint{0.855760in}{1.041810in}}%
\pgfpathlineto{\pgfqpoint{0.847388in}{1.036985in}}%
\pgfpathclose%
\pgfusepath{fill}%
\end{pgfscope}%
\begin{pgfscope}%
\pgfpathrectangle{\pgfqpoint{0.407592in}{0.222904in}}{\pgfqpoint{1.411163in}{1.411163in}}%
\pgfusepath{clip}%
\pgfsetbuttcap%
\pgfsetroundjoin%
\definecolor{currentfill}{rgb}{0.909804,0.698039,0.423529}%
\pgfsetfillcolor{currentfill}%
\pgfsetfillopacity{0.120000}%
\pgfsetlinewidth{0.000000pt}%
\definecolor{currentstroke}{rgb}{0.000000,0.000000,0.000000}%
\pgfsetstrokecolor{currentstroke}%
\pgfsetdash{}{0pt}%
\pgfpathmoveto{\pgfqpoint{1.008455in}{0.955288in}}%
\pgfpathlineto{\pgfqpoint{1.024510in}{0.945362in}}%
\pgfpathlineto{\pgfqpoint{1.040590in}{0.936278in}}%
\pgfpathlineto{\pgfqpoint{1.048668in}{0.943481in}}%
\pgfpathlineto{\pgfqpoint{1.056720in}{0.951194in}}%
\pgfpathlineto{\pgfqpoint{1.040672in}{0.960155in}}%
\pgfpathlineto{\pgfqpoint{1.024648in}{0.969945in}}%
\pgfpathlineto{\pgfqpoint{1.016564in}{0.962371in}}%
\pgfpathlineto{\pgfqpoint{1.008455in}{0.955288in}}%
\pgfpathclose%
\pgfusepath{fill}%
\end{pgfscope}%
\begin{pgfscope}%
\pgfpathrectangle{\pgfqpoint{0.407592in}{0.222904in}}{\pgfqpoint{1.411163in}{1.411163in}}%
\pgfusepath{clip}%
\pgfsetbuttcap%
\pgfsetroundjoin%
\definecolor{currentfill}{rgb}{0.909804,0.698039,0.423529}%
\pgfsetfillcolor{currentfill}%
\pgfsetfillopacity{0.120000}%
\pgfsetlinewidth{0.000000pt}%
\definecolor{currentstroke}{rgb}{0.000000,0.000000,0.000000}%
\pgfsetstrokecolor{currentstroke}%
\pgfsetdash{}{0pt}%
\pgfpathmoveto{\pgfqpoint{0.621502in}{1.225186in}}%
\pgfpathlineto{\pgfqpoint{0.637454in}{1.207098in}}%
\pgfpathlineto{\pgfqpoint{0.653415in}{1.189160in}}%
\pgfpathlineto{\pgfqpoint{0.662129in}{1.192456in}}%
\pgfpathlineto{\pgfqpoint{0.670808in}{1.196064in}}%
\pgfpathlineto{\pgfqpoint{0.654888in}{1.214004in}}%
\pgfpathlineto{\pgfqpoint{0.638978in}{1.232091in}}%
\pgfpathlineto{\pgfqpoint{0.630258in}{1.228490in}}%
\pgfpathlineto{\pgfqpoint{0.621502in}{1.225186in}}%
\pgfpathclose%
\pgfusepath{fill}%
\end{pgfscope}%
\begin{pgfscope}%
\pgfpathrectangle{\pgfqpoint{0.407592in}{0.222904in}}{\pgfqpoint{1.411163in}{1.411163in}}%
\pgfusepath{clip}%
\pgfsetbuttcap%
\pgfsetroundjoin%
\definecolor{currentfill}{rgb}{0.909804,0.698039,0.423529}%
\pgfsetfillcolor{currentfill}%
\pgfsetfillopacity{0.120000}%
\pgfsetlinewidth{0.000000pt}%
\definecolor{currentstroke}{rgb}{0.000000,0.000000,0.000000}%
\pgfsetstrokecolor{currentstroke}%
\pgfsetdash{}{0pt}%
\pgfpathmoveto{\pgfqpoint{1.457873in}{1.130277in}}%
\pgfpathlineto{\pgfqpoint{1.474571in}{1.134015in}}%
\pgfpathlineto{\pgfqpoint{1.491333in}{1.138269in}}%
\pgfpathlineto{\pgfqpoint{1.498905in}{1.151344in}}%
\pgfpathlineto{\pgfqpoint{1.506465in}{1.164588in}}%
\pgfpathlineto{\pgfqpoint{1.489719in}{1.160582in}}%
\pgfpathlineto{\pgfqpoint{1.473037in}{1.157082in}}%
\pgfpathlineto{\pgfqpoint{1.465462in}{1.143598in}}%
\pgfpathlineto{\pgfqpoint{1.457873in}{1.130277in}}%
\pgfpathclose%
\pgfusepath{fill}%
\end{pgfscope}%
\begin{pgfscope}%
\pgfpathrectangle{\pgfqpoint{0.407592in}{0.222904in}}{\pgfqpoint{1.411163in}{1.411163in}}%
\pgfusepath{clip}%
\pgfsetbuttcap%
\pgfsetroundjoin%
\definecolor{currentfill}{rgb}{0.909804,0.698039,0.423529}%
\pgfsetfillcolor{currentfill}%
\pgfsetfillopacity{0.120000}%
\pgfsetlinewidth{0.000000pt}%
\definecolor{currentstroke}{rgb}{0.000000,0.000000,0.000000}%
\pgfsetstrokecolor{currentstroke}%
\pgfsetdash{}{0pt}%
\pgfpathmoveto{\pgfqpoint{1.345716in}{1.042490in}}%
\pgfpathlineto{\pgfqpoint{1.362214in}{1.044518in}}%
\pgfpathlineto{\pgfqpoint{1.378773in}{1.047207in}}%
\pgfpathlineto{\pgfqpoint{1.386455in}{1.059531in}}%
\pgfpathlineto{\pgfqpoint{1.394123in}{1.072094in}}%
\pgfpathlineto{\pgfqpoint{1.377580in}{1.069621in}}%
\pgfpathlineto{\pgfqpoint{1.361097in}{1.067790in}}%
\pgfpathlineto{\pgfqpoint{1.353413in}{1.055022in}}%
\pgfpathlineto{\pgfqpoint{1.345716in}{1.042490in}}%
\pgfpathclose%
\pgfusepath{fill}%
\end{pgfscope}%
\begin{pgfscope}%
\pgfpathrectangle{\pgfqpoint{0.407592in}{0.222904in}}{\pgfqpoint{1.411163in}{1.411163in}}%
\pgfusepath{clip}%
\pgfsetbuttcap%
\pgfsetroundjoin%
\definecolor{currentfill}{rgb}{0.909804,0.698039,0.423529}%
\pgfsetfillcolor{currentfill}%
\pgfsetfillopacity{0.120000}%
\pgfsetlinewidth{0.000000pt}%
\definecolor{currentstroke}{rgb}{0.000000,0.000000,0.000000}%
\pgfsetstrokecolor{currentstroke}%
\pgfsetdash{}{0pt}%
\pgfpathmoveto{\pgfqpoint{0.766526in}{1.092380in}}%
\pgfpathlineto{\pgfqpoint{0.782515in}{1.075944in}}%
\pgfpathlineto{\pgfqpoint{0.798513in}{1.059814in}}%
\pgfpathlineto{\pgfqpoint{0.806992in}{1.063792in}}%
\pgfpathlineto{\pgfqpoint{0.815437in}{1.068166in}}%
\pgfpathlineto{\pgfqpoint{0.799478in}{1.084278in}}%
\pgfpathlineto{\pgfqpoint{0.783531in}{1.100694in}}%
\pgfpathlineto{\pgfqpoint{0.775045in}{1.096350in}}%
\pgfpathlineto{\pgfqpoint{0.766526in}{1.092380in}}%
\pgfpathclose%
\pgfusepath{fill}%
\end{pgfscope}%
\begin{pgfscope}%
\pgfpathrectangle{\pgfqpoint{0.407592in}{0.222904in}}{\pgfqpoint{1.411163in}{1.411163in}}%
\pgfusepath{clip}%
\pgfsetbuttcap%
\pgfsetroundjoin%
\definecolor{currentfill}{rgb}{0.909804,0.698039,0.423529}%
\pgfsetfillcolor{currentfill}%
\pgfsetfillopacity{0.120000}%
\pgfsetlinewidth{0.000000pt}%
\definecolor{currentstroke}{rgb}{0.000000,0.000000,0.000000}%
\pgfsetstrokecolor{currentstroke}%
\pgfsetdash{}{0pt}%
\pgfpathmoveto{\pgfqpoint{0.653415in}{1.189160in}}%
\pgfpathlineto{\pgfqpoint{0.669386in}{1.171385in}}%
\pgfpathlineto{\pgfqpoint{0.685366in}{1.153789in}}%
\pgfpathlineto{\pgfqpoint{0.694038in}{1.157069in}}%
\pgfpathlineto{\pgfqpoint{0.702675in}{1.160679in}}%
\pgfpathlineto{\pgfqpoint{0.686737in}{1.178284in}}%
\pgfpathlineto{\pgfqpoint{0.670808in}{1.196064in}}%
\pgfpathlineto{\pgfqpoint{0.662129in}{1.192456in}}%
\pgfpathlineto{\pgfqpoint{0.653415in}{1.189160in}}%
\pgfpathclose%
\pgfusepath{fill}%
\end{pgfscope}%
\begin{pgfscope}%
\pgfpathrectangle{\pgfqpoint{0.407592in}{0.222904in}}{\pgfqpoint{1.411163in}{1.411163in}}%
\pgfusepath{clip}%
\pgfsetbuttcap%
\pgfsetroundjoin%
\definecolor{currentfill}{rgb}{0.909804,0.698039,0.423529}%
\pgfsetfillcolor{currentfill}%
\pgfsetfillopacity{0.120000}%
\pgfsetlinewidth{0.000000pt}%
\definecolor{currentstroke}{rgb}{0.000000,0.000000,0.000000}%
\pgfsetstrokecolor{currentstroke}%
\pgfsetdash{}{0pt}%
\pgfpathmoveto{\pgfqpoint{1.137142in}{0.941748in}}%
\pgfpathlineto{\pgfqpoint{1.153328in}{0.937520in}}%
\pgfpathlineto{\pgfqpoint{1.169555in}{0.934259in}}%
\pgfpathlineto{\pgfqpoint{1.177467in}{0.943650in}}%
\pgfpathlineto{\pgfqpoint{1.185359in}{0.953490in}}%
\pgfpathlineto{\pgfqpoint{1.169154in}{0.956740in}}%
\pgfpathlineto{\pgfqpoint{1.152990in}{0.960926in}}%
\pgfpathlineto{\pgfqpoint{1.145076in}{0.951113in}}%
\pgfpathlineto{\pgfqpoint{1.137142in}{0.941748in}}%
\pgfpathclose%
\pgfusepath{fill}%
\end{pgfscope}%
\begin{pgfscope}%
\pgfpathrectangle{\pgfqpoint{0.407592in}{0.222904in}}{\pgfqpoint{1.411163in}{1.411163in}}%
\pgfusepath{clip}%
\pgfsetbuttcap%
\pgfsetroundjoin%
\definecolor{currentfill}{rgb}{0.909804,0.698039,0.423529}%
\pgfsetfillcolor{currentfill}%
\pgfsetfillopacity{0.120000}%
\pgfsetlinewidth{0.000000pt}%
\definecolor{currentstroke}{rgb}{0.000000,0.000000,0.000000}%
\pgfsetstrokecolor{currentstroke}%
\pgfsetdash{}{0pt}%
\pgfpathmoveto{\pgfqpoint{1.185359in}{0.953490in}}%
\pgfpathlineto{\pgfqpoint{1.201609in}{0.951171in}}%
\pgfpathlineto{\pgfqpoint{1.217906in}{0.949774in}}%
\pgfpathlineto{\pgfqpoint{1.225760in}{0.959999in}}%
\pgfpathlineto{\pgfqpoint{1.233597in}{0.970624in}}%
\pgfpathlineto{\pgfqpoint{1.217319in}{0.972078in}}%
\pgfpathlineto{\pgfqpoint{1.201089in}{0.974421in}}%
\pgfpathlineto{\pgfqpoint{1.193233in}{0.963755in}}%
\pgfpathlineto{\pgfqpoint{1.185359in}{0.953490in}}%
\pgfpathclose%
\pgfusepath{fill}%
\end{pgfscope}%
\begin{pgfscope}%
\pgfpathrectangle{\pgfqpoint{0.407592in}{0.222904in}}{\pgfqpoint{1.411163in}{1.411163in}}%
\pgfusepath{clip}%
\pgfsetbuttcap%
\pgfsetroundjoin%
\definecolor{currentfill}{rgb}{0.909804,0.698039,0.423529}%
\pgfsetfillcolor{currentfill}%
\pgfsetfillopacity{0.120000}%
\pgfsetlinewidth{0.000000pt}%
\definecolor{currentstroke}{rgb}{0.000000,0.000000,0.000000}%
\pgfsetstrokecolor{currentstroke}%
\pgfsetdash{}{0pt}%
\pgfpathmoveto{\pgfqpoint{0.960042in}{0.965026in}}%
\pgfpathlineto{\pgfqpoint{0.976089in}{0.953465in}}%
\pgfpathlineto{\pgfqpoint{0.992156in}{0.942647in}}%
\pgfpathlineto{\pgfqpoint{1.000319in}{0.948710in}}%
\pgfpathlineto{\pgfqpoint{1.008455in}{0.955288in}}%
\pgfpathlineto{\pgfqpoint{0.992422in}{0.966004in}}%
\pgfpathlineto{\pgfqpoint{0.976410in}{0.977458in}}%
\pgfpathlineto{\pgfqpoint{0.968240in}{0.970996in}}%
\pgfpathlineto{\pgfqpoint{0.960042in}{0.965026in}}%
\pgfpathclose%
\pgfusepath{fill}%
\end{pgfscope}%
\begin{pgfscope}%
\pgfpathrectangle{\pgfqpoint{0.407592in}{0.222904in}}{\pgfqpoint{1.411163in}{1.411163in}}%
\pgfusepath{clip}%
\pgfsetbuttcap%
\pgfsetroundjoin%
\definecolor{currentfill}{rgb}{0.909804,0.698039,0.423529}%
\pgfsetfillcolor{currentfill}%
\pgfsetfillopacity{0.120000}%
\pgfsetlinewidth{0.000000pt}%
\definecolor{currentstroke}{rgb}{0.000000,0.000000,0.000000}%
\pgfsetstrokecolor{currentstroke}%
\pgfsetdash{}{0pt}%
\pgfpathmoveto{\pgfqpoint{0.879388in}{1.007429in}}%
\pgfpathlineto{\pgfqpoint{0.895408in}{0.993364in}}%
\pgfpathlineto{\pgfqpoint{0.911442in}{0.979834in}}%
\pgfpathlineto{\pgfqpoint{0.919737in}{0.984759in}}%
\pgfpathlineto{\pgfqpoint{0.928001in}{0.990159in}}%
\pgfpathlineto{\pgfqpoint{0.912005in}{1.003627in}}%
\pgfpathlineto{\pgfqpoint{0.896023in}{1.017630in}}%
\pgfpathlineto{\pgfqpoint{0.887721in}{1.012306in}}%
\pgfpathlineto{\pgfqpoint{0.879388in}{1.007429in}}%
\pgfpathclose%
\pgfusepath{fill}%
\end{pgfscope}%
\begin{pgfscope}%
\pgfpathrectangle{\pgfqpoint{0.407592in}{0.222904in}}{\pgfqpoint{1.411163in}{1.411163in}}%
\pgfusepath{clip}%
\pgfsetbuttcap%
\pgfsetroundjoin%
\definecolor{currentfill}{rgb}{0.909804,0.698039,0.423529}%
\pgfsetfillcolor{currentfill}%
\pgfsetfillopacity{0.120000}%
\pgfsetlinewidth{0.000000pt}%
\definecolor{currentstroke}{rgb}{0.000000,0.000000,0.000000}%
\pgfsetstrokecolor{currentstroke}%
\pgfsetdash{}{0pt}%
\pgfpathmoveto{\pgfqpoint{1.088902in}{0.935933in}}%
\pgfpathlineto{\pgfqpoint{1.105040in}{0.929706in}}%
\pgfpathlineto{\pgfqpoint{1.121212in}{0.924452in}}%
\pgfpathlineto{\pgfqpoint{1.129188in}{0.932854in}}%
\pgfpathlineto{\pgfqpoint{1.137142in}{0.941748in}}%
\pgfpathlineto{\pgfqpoint{1.120995in}{0.946937in}}%
\pgfpathlineto{\pgfqpoint{1.104883in}{0.953073in}}%
\pgfpathlineto{\pgfqpoint{1.096904in}{0.944260in}}%
\pgfpathlineto{\pgfqpoint{1.088902in}{0.935933in}}%
\pgfpathclose%
\pgfusepath{fill}%
\end{pgfscope}%
\begin{pgfscope}%
\pgfpathrectangle{\pgfqpoint{0.407592in}{0.222904in}}{\pgfqpoint{1.411163in}{1.411163in}}%
\pgfusepath{clip}%
\pgfsetbuttcap%
\pgfsetroundjoin%
\definecolor{currentfill}{rgb}{0.909804,0.698039,0.423529}%
\pgfsetfillcolor{currentfill}%
\pgfsetfillopacity{0.120000}%
\pgfsetlinewidth{0.000000pt}%
\definecolor{currentstroke}{rgb}{0.000000,0.000000,0.000000}%
\pgfsetstrokecolor{currentstroke}%
\pgfsetdash{}{0pt}%
\pgfpathmoveto{\pgfqpoint{1.233597in}{0.970624in}}%
\pgfpathlineto{\pgfqpoint{1.249924in}{0.970040in}}%
\pgfpathlineto{\pgfqpoint{1.266304in}{0.970304in}}%
\pgfpathlineto{\pgfqpoint{1.274107in}{0.981198in}}%
\pgfpathlineto{\pgfqpoint{1.281894in}{0.992444in}}%
\pgfpathlineto{\pgfqpoint{1.265531in}{0.992303in}}%
\pgfpathlineto{\pgfqpoint{1.249222in}{0.992980in}}%
\pgfpathlineto{\pgfqpoint{1.241417in}{0.981625in}}%
\pgfpathlineto{\pgfqpoint{1.233597in}{0.970624in}}%
\pgfpathclose%
\pgfusepath{fill}%
\end{pgfscope}%
\begin{pgfscope}%
\pgfpathrectangle{\pgfqpoint{0.407592in}{0.222904in}}{\pgfqpoint{1.411163in}{1.411163in}}%
\pgfusepath{clip}%
\pgfsetbuttcap%
\pgfsetroundjoin%
\definecolor{currentfill}{rgb}{0.909804,0.698039,0.423529}%
\pgfsetfillcolor{currentfill}%
\pgfsetfillopacity{0.120000}%
\pgfsetlinewidth{0.000000pt}%
\definecolor{currentstroke}{rgb}{0.000000,0.000000,0.000000}%
\pgfsetstrokecolor{currentstroke}%
\pgfsetdash{}{0pt}%
\pgfpathmoveto{\pgfqpoint{1.394123in}{1.072094in}}%
\pgfpathlineto{\pgfqpoint{1.410727in}{1.075181in}}%
\pgfpathlineto{\pgfqpoint{1.427394in}{1.078857in}}%
\pgfpathlineto{\pgfqpoint{1.435033in}{1.091407in}}%
\pgfpathlineto{\pgfqpoint{1.442659in}{1.104169in}}%
\pgfpathlineto{\pgfqpoint{1.426008in}{1.100735in}}%
\pgfpathlineto{\pgfqpoint{1.409420in}{1.097874in}}%
\pgfpathlineto{\pgfqpoint{1.401778in}{1.084880in}}%
\pgfpathlineto{\pgfqpoint{1.394123in}{1.072094in}}%
\pgfpathclose%
\pgfusepath{fill}%
\end{pgfscope}%
\begin{pgfscope}%
\pgfpathrectangle{\pgfqpoint{0.407592in}{0.222904in}}{\pgfqpoint{1.411163in}{1.411163in}}%
\pgfusepath{clip}%
\pgfsetbuttcap%
\pgfsetroundjoin%
\definecolor{currentfill}{rgb}{0.909804,0.698039,0.423529}%
\pgfsetfillcolor{currentfill}%
\pgfsetfillopacity{0.120000}%
\pgfsetlinewidth{0.000000pt}%
\definecolor{currentstroke}{rgb}{0.000000,0.000000,0.000000}%
\pgfsetstrokecolor{currentstroke}%
\pgfsetdash{}{0pt}%
\pgfpathmoveto{\pgfqpoint{0.798513in}{1.059814in}}%
\pgfpathlineto{\pgfqpoint{0.814523in}{1.044023in}}%
\pgfpathlineto{\pgfqpoint{0.830545in}{1.028602in}}%
\pgfpathlineto{\pgfqpoint{0.838983in}{1.032582in}}%
\pgfpathlineto{\pgfqpoint{0.847388in}{1.036985in}}%
\pgfpathlineto{\pgfqpoint{0.831407in}{1.052391in}}%
\pgfpathlineto{\pgfqpoint{0.815437in}{1.068166in}}%
\pgfpathlineto{\pgfqpoint{0.806992in}{1.063792in}}%
\pgfpathlineto{\pgfqpoint{0.798513in}{1.059814in}}%
\pgfpathclose%
\pgfusepath{fill}%
\end{pgfscope}%
\begin{pgfscope}%
\pgfpathrectangle{\pgfqpoint{0.407592in}{0.222904in}}{\pgfqpoint{1.411163in}{1.411163in}}%
\pgfusepath{clip}%
\pgfsetbuttcap%
\pgfsetroundjoin%
\definecolor{currentfill}{rgb}{0.909804,0.698039,0.423529}%
\pgfsetfillcolor{currentfill}%
\pgfsetfillopacity{0.120000}%
\pgfsetlinewidth{0.000000pt}%
\definecolor{currentstroke}{rgb}{0.000000,0.000000,0.000000}%
\pgfsetstrokecolor{currentstroke}%
\pgfsetdash{}{0pt}%
\pgfpathmoveto{\pgfqpoint{0.685366in}{1.153789in}}%
\pgfpathlineto{\pgfqpoint{0.701355in}{1.136388in}}%
\pgfpathlineto{\pgfqpoint{0.717354in}{1.119202in}}%
\pgfpathlineto{\pgfqpoint{0.725985in}{1.122459in}}%
\pgfpathlineto{\pgfqpoint{0.734581in}{1.126066in}}%
\pgfpathlineto{\pgfqpoint{0.718623in}{1.143267in}}%
\pgfpathlineto{\pgfqpoint{0.702675in}{1.160679in}}%
\pgfpathlineto{\pgfqpoint{0.694038in}{1.157069in}}%
\pgfpathlineto{\pgfqpoint{0.685366in}{1.153789in}}%
\pgfpathclose%
\pgfusepath{fill}%
\end{pgfscope}%
\begin{pgfscope}%
\pgfpathrectangle{\pgfqpoint{0.407592in}{0.222904in}}{\pgfqpoint{1.411163in}{1.411163in}}%
\pgfusepath{clip}%
\pgfsetbuttcap%
\pgfsetroundjoin%
\definecolor{currentfill}{rgb}{0.909804,0.698039,0.423529}%
\pgfsetfillcolor{currentfill}%
\pgfsetfillopacity{0.120000}%
\pgfsetlinewidth{0.000000pt}%
\definecolor{currentstroke}{rgb}{0.000000,0.000000,0.000000}%
\pgfsetstrokecolor{currentstroke}%
\pgfsetdash{}{0pt}%
\pgfpathmoveto{\pgfqpoint{1.040590in}{0.936278in}}%
\pgfpathlineto{\pgfqpoint{1.056695in}{0.928081in}}%
\pgfpathlineto{\pgfqpoint{1.072829in}{0.920814in}}%
\pgfpathlineto{\pgfqpoint{1.080878in}{0.928112in}}%
\pgfpathlineto{\pgfqpoint{1.088902in}{0.935933in}}%
\pgfpathlineto{\pgfqpoint{1.072797in}{0.943107in}}%
\pgfpathlineto{\pgfqpoint{1.056720in}{0.951194in}}%
\pgfpathlineto{\pgfqpoint{1.048668in}{0.943481in}}%
\pgfpathlineto{\pgfqpoint{1.040590in}{0.936278in}}%
\pgfpathclose%
\pgfusepath{fill}%
\end{pgfscope}%
\begin{pgfscope}%
\pgfpathrectangle{\pgfqpoint{0.407592in}{0.222904in}}{\pgfqpoint{1.411163in}{1.411163in}}%
\pgfusepath{clip}%
\pgfsetbuttcap%
\pgfsetroundjoin%
\definecolor{currentfill}{rgb}{0.909804,0.698039,0.423529}%
\pgfsetfillcolor{currentfill}%
\pgfsetfillopacity{0.120000}%
\pgfsetlinewidth{0.000000pt}%
\definecolor{currentstroke}{rgb}{0.000000,0.000000,0.000000}%
\pgfsetstrokecolor{currentstroke}%
\pgfsetdash{}{0pt}%
\pgfpathmoveto{\pgfqpoint{1.281894in}{0.992444in}}%
\pgfpathlineto{\pgfqpoint{1.298310in}{0.993377in}}%
\pgfpathlineto{\pgfqpoint{1.314784in}{0.995073in}}%
\pgfpathlineto{\pgfqpoint{1.322538in}{1.006484in}}%
\pgfpathlineto{\pgfqpoint{1.330278in}{1.018203in}}%
\pgfpathlineto{\pgfqpoint{1.313822in}{1.016685in}}%
\pgfpathlineto{\pgfqpoint{1.297422in}{1.015905in}}%
\pgfpathlineto{\pgfqpoint{1.289665in}{1.004020in}}%
\pgfpathlineto{\pgfqpoint{1.281894in}{0.992444in}}%
\pgfpathclose%
\pgfusepath{fill}%
\end{pgfscope}%
\begin{pgfscope}%
\pgfpathrectangle{\pgfqpoint{0.407592in}{0.222904in}}{\pgfqpoint{1.411163in}{1.411163in}}%
\pgfusepath{clip}%
\pgfsetbuttcap%
\pgfsetroundjoin%
\definecolor{currentfill}{rgb}{0.909804,0.698039,0.423529}%
\pgfsetfillcolor{currentfill}%
\pgfsetfillopacity{0.120000}%
\pgfsetlinewidth{0.000000pt}%
\definecolor{currentstroke}{rgb}{0.000000,0.000000,0.000000}%
\pgfsetstrokecolor{currentstroke}%
\pgfsetdash{}{0pt}%
\pgfpathmoveto{\pgfqpoint{1.506465in}{1.164588in}}%
\pgfpathlineto{\pgfqpoint{1.523277in}{1.169078in}}%
\pgfpathlineto{\pgfqpoint{1.540154in}{1.174033in}}%
\pgfpathlineto{\pgfqpoint{1.547685in}{1.187181in}}%
\pgfpathlineto{\pgfqpoint{1.530816in}{1.182354in}}%
\pgfpathlineto{\pgfqpoint{1.514013in}{1.177989in}}%
\pgfusepath{fill}%
\end{pgfscope}%
\begin{pgfscope}%
\pgfpathrectangle{\pgfqpoint{0.407592in}{0.222904in}}{\pgfqpoint{1.411163in}{1.411163in}}%
\pgfusepath{clip}%
\pgfsetbuttcap%
\pgfsetroundjoin%
\definecolor{currentfill}{rgb}{0.909804,0.698039,0.423529}%
\pgfsetfillcolor{currentfill}%
\pgfsetfillopacity{0.120000}%
\pgfsetlinewidth{0.000000pt}%
\definecolor{currentstroke}{rgb}{0.000000,0.000000,0.000000}%
\pgfsetstrokecolor{currentstroke}%
\pgfsetdash{}{0pt}%
\pgfpathmoveto{\pgfqpoint{0.911442in}{0.979834in}}%
\pgfpathlineto{\pgfqpoint{0.927491in}{0.966889in}}%
\pgfpathlineto{\pgfqpoint{0.943557in}{0.954581in}}%
\pgfpathlineto{\pgfqpoint{0.951814in}{0.959552in}}%
\pgfpathlineto{\pgfqpoint{0.960042in}{0.965026in}}%
\pgfpathlineto{\pgfqpoint{0.944013in}{0.977275in}}%
\pgfpathlineto{\pgfqpoint{0.928001in}{0.990159in}}%
\pgfpathlineto{\pgfqpoint{0.919737in}{0.984759in}}%
\pgfpathlineto{\pgfqpoint{0.911442in}{0.979834in}}%
\pgfpathclose%
\pgfusepath{fill}%
\end{pgfscope}%
\begin{pgfscope}%
\pgfpathrectangle{\pgfqpoint{0.407592in}{0.222904in}}{\pgfqpoint{1.411163in}{1.411163in}}%
\pgfusepath{clip}%
\pgfsetbuttcap%
\pgfsetroundjoin%
\definecolor{currentfill}{rgb}{0.909804,0.698039,0.423529}%
\pgfsetfillcolor{currentfill}%
\pgfsetfillopacity{0.120000}%
\pgfsetlinewidth{0.000000pt}%
\definecolor{currentstroke}{rgb}{0.000000,0.000000,0.000000}%
\pgfsetstrokecolor{currentstroke}%
\pgfsetdash{}{0pt}%
\pgfpathmoveto{\pgfqpoint{1.442659in}{1.104169in}}%
\pgfpathlineto{\pgfqpoint{1.459373in}{1.108154in}}%
\pgfpathlineto{\pgfqpoint{1.476152in}{1.112664in}}%
\pgfpathlineto{\pgfqpoint{1.483749in}{1.125371in}}%
\pgfpathlineto{\pgfqpoint{1.491333in}{1.138269in}}%
\pgfpathlineto{\pgfqpoint{1.474571in}{1.134015in}}%
\pgfpathlineto{\pgfqpoint{1.457873in}{1.130277in}}%
\pgfpathlineto{\pgfqpoint{1.450272in}{1.117130in}}%
\pgfpathlineto{\pgfqpoint{1.442659in}{1.104169in}}%
\pgfpathclose%
\pgfusepath{fill}%
\end{pgfscope}%
\begin{pgfscope}%
\pgfpathrectangle{\pgfqpoint{0.407592in}{0.222904in}}{\pgfqpoint{1.411163in}{1.411163in}}%
\pgfusepath{clip}%
\pgfsetbuttcap%
\pgfsetroundjoin%
\definecolor{currentfill}{rgb}{0.909804,0.698039,0.423529}%
\pgfsetfillcolor{currentfill}%
\pgfsetfillopacity{0.120000}%
\pgfsetlinewidth{0.000000pt}%
\definecolor{currentstroke}{rgb}{0.000000,0.000000,0.000000}%
\pgfsetstrokecolor{currentstroke}%
\pgfsetdash{}{0pt}%
\pgfpathmoveto{\pgfqpoint{1.330278in}{1.018203in}}%
\pgfpathlineto{\pgfqpoint{1.346793in}{1.020431in}}%
\pgfpathlineto{\pgfqpoint{1.363369in}{1.023338in}}%
\pgfpathlineto{\pgfqpoint{1.371078in}{1.035137in}}%
\pgfpathlineto{\pgfqpoint{1.378773in}{1.047207in}}%
\pgfpathlineto{\pgfqpoint{1.362214in}{1.044518in}}%
\pgfpathlineto{\pgfqpoint{1.345716in}{1.042490in}}%
\pgfpathlineto{\pgfqpoint{1.338004in}{1.030211in}}%
\pgfpathlineto{\pgfqpoint{1.330278in}{1.018203in}}%
\pgfpathclose%
\pgfusepath{fill}%
\end{pgfscope}%
\begin{pgfscope}%
\pgfpathrectangle{\pgfqpoint{0.407592in}{0.222904in}}{\pgfqpoint{1.411163in}{1.411163in}}%
\pgfusepath{clip}%
\pgfsetbuttcap%
\pgfsetroundjoin%
\definecolor{currentfill}{rgb}{0.909804,0.698039,0.423529}%
\pgfsetfillcolor{currentfill}%
\pgfsetfillopacity{0.120000}%
\pgfsetlinewidth{0.000000pt}%
\definecolor{currentstroke}{rgb}{0.000000,0.000000,0.000000}%
\pgfsetstrokecolor{currentstroke}%
\pgfsetdash{}{0pt}%
\pgfpathmoveto{\pgfqpoint{0.717354in}{1.119202in}}%
\pgfpathlineto{\pgfqpoint{0.733363in}{1.102251in}}%
\pgfpathlineto{\pgfqpoint{0.749382in}{1.085557in}}%
\pgfpathlineto{\pgfqpoint{0.757972in}{1.088783in}}%
\pgfpathlineto{\pgfqpoint{0.766526in}{1.092380in}}%
\pgfpathlineto{\pgfqpoint{0.750548in}{1.109096in}}%
\pgfpathlineto{\pgfqpoint{0.734581in}{1.126066in}}%
\pgfpathlineto{\pgfqpoint{0.725985in}{1.122459in}}%
\pgfpathlineto{\pgfqpoint{0.717354in}{1.119202in}}%
\pgfpathclose%
\pgfusepath{fill}%
\end{pgfscope}%
\begin{pgfscope}%
\pgfpathrectangle{\pgfqpoint{0.407592in}{0.222904in}}{\pgfqpoint{1.411163in}{1.411163in}}%
\pgfusepath{clip}%
\pgfsetbuttcap%
\pgfsetroundjoin%
\definecolor{currentfill}{rgb}{0.909804,0.698039,0.423529}%
\pgfsetfillcolor{currentfill}%
\pgfsetfillopacity{0.120000}%
\pgfsetlinewidth{0.000000pt}%
\definecolor{currentstroke}{rgb}{0.000000,0.000000,0.000000}%
\pgfsetstrokecolor{currentstroke}%
\pgfsetdash{}{0pt}%
\pgfpathmoveto{\pgfqpoint{0.992156in}{0.942647in}}%
\pgfpathlineto{\pgfqpoint{1.008244in}{0.932626in}}%
\pgfpathlineto{\pgfqpoint{1.024355in}{0.923453in}}%
\pgfpathlineto{\pgfqpoint{1.032486in}{0.929598in}}%
\pgfpathlineto{\pgfqpoint{1.040590in}{0.936278in}}%
\pgfpathlineto{\pgfqpoint{1.024510in}{0.945362in}}%
\pgfpathlineto{\pgfqpoint{1.008455in}{0.955288in}}%
\pgfpathlineto{\pgfqpoint{1.000319in}{0.948710in}}%
\pgfpathlineto{\pgfqpoint{0.992156in}{0.942647in}}%
\pgfpathclose%
\pgfusepath{fill}%
\end{pgfscope}%
\begin{pgfscope}%
\pgfpathrectangle{\pgfqpoint{0.407592in}{0.222904in}}{\pgfqpoint{1.411163in}{1.411163in}}%
\pgfusepath{clip}%
\pgfsetbuttcap%
\pgfsetroundjoin%
\definecolor{currentfill}{rgb}{0.909804,0.698039,0.423529}%
\pgfsetfillcolor{currentfill}%
\pgfsetfillopacity{0.120000}%
\pgfsetlinewidth{0.000000pt}%
\definecolor{currentstroke}{rgb}{0.000000,0.000000,0.000000}%
\pgfsetstrokecolor{currentstroke}%
\pgfsetdash{}{0pt}%
\pgfpathmoveto{\pgfqpoint{0.830545in}{1.028602in}}%
\pgfpathlineto{\pgfqpoint{0.846578in}{1.013590in}}%
\pgfpathlineto{\pgfqpoint{0.862624in}{0.999025in}}%
\pgfpathlineto{\pgfqpoint{0.871023in}{1.003002in}}%
\pgfpathlineto{\pgfqpoint{0.879388in}{1.007429in}}%
\pgfpathlineto{\pgfqpoint{0.863382in}{1.021984in}}%
\pgfpathlineto{\pgfqpoint{0.847388in}{1.036985in}}%
\pgfpathlineto{\pgfqpoint{0.838983in}{1.032582in}}%
\pgfpathlineto{\pgfqpoint{0.830545in}{1.028602in}}%
\pgfpathclose%
\pgfusepath{fill}%
\end{pgfscope}%
\begin{pgfscope}%
\pgfpathrectangle{\pgfqpoint{0.407592in}{0.222904in}}{\pgfqpoint{1.411163in}{1.411163in}}%
\pgfusepath{clip}%
\pgfsetbuttcap%
\pgfsetroundjoin%
\definecolor{currentfill}{rgb}{0.909804,0.698039,0.423529}%
\pgfsetfillcolor{currentfill}%
\pgfsetfillopacity{0.120000}%
\pgfsetlinewidth{0.000000pt}%
\definecolor{currentstroke}{rgb}{0.000000,0.000000,0.000000}%
\pgfsetstrokecolor{currentstroke}%
\pgfsetdash{}{0pt}%
\pgfpathmoveto{\pgfqpoint{0.603879in}{1.219459in}}%
\pgfpathlineto{\pgfqpoint{0.619873in}{1.201403in}}%
\pgfpathlineto{\pgfqpoint{0.635876in}{1.183500in}}%
\pgfpathlineto{\pgfqpoint{0.644664in}{1.186175in}}%
\pgfpathlineto{\pgfqpoint{0.653415in}{1.189160in}}%
\pgfpathlineto{\pgfqpoint{0.637454in}{1.207098in}}%
\pgfpathlineto{\pgfqpoint{0.621502in}{1.225186in}}%
\pgfpathlineto{\pgfqpoint{0.612709in}{1.222176in}}%
\pgfpathlineto{\pgfqpoint{0.603879in}{1.219459in}}%
\pgfpathclose%
\pgfusepath{fill}%
\end{pgfscope}%
\begin{pgfscope}%
\pgfpathrectangle{\pgfqpoint{0.407592in}{0.222904in}}{\pgfqpoint{1.411163in}{1.411163in}}%
\pgfusepath{clip}%
\pgfsetbuttcap%
\pgfsetroundjoin%
\definecolor{currentfill}{rgb}{0.909804,0.698039,0.423529}%
\pgfsetfillcolor{currentfill}%
\pgfsetfillopacity{0.120000}%
\pgfsetlinewidth{0.000000pt}%
\definecolor{currentstroke}{rgb}{0.000000,0.000000,0.000000}%
\pgfsetstrokecolor{currentstroke}%
\pgfsetdash{}{0pt}%
\pgfpathmoveto{\pgfqpoint{1.169555in}{0.934259in}}%
\pgfpathlineto{\pgfqpoint{1.185826in}{0.931960in}}%
\pgfpathlineto{\pgfqpoint{1.202143in}{0.930612in}}%
\pgfpathlineto{\pgfqpoint{1.210034in}{0.939970in}}%
\pgfpathlineto{\pgfqpoint{1.217906in}{0.949774in}}%
\pgfpathlineto{\pgfqpoint{1.201609in}{0.951171in}}%
\pgfpathlineto{\pgfqpoint{1.185359in}{0.953490in}}%
\pgfpathlineto{\pgfqpoint{1.177467in}{0.943650in}}%
\pgfpathlineto{\pgfqpoint{1.169555in}{0.934259in}}%
\pgfpathclose%
\pgfusepath{fill}%
\end{pgfscope}%
\begin{pgfscope}%
\pgfpathrectangle{\pgfqpoint{0.407592in}{0.222904in}}{\pgfqpoint{1.411163in}{1.411163in}}%
\pgfusepath{clip}%
\pgfsetbuttcap%
\pgfsetroundjoin%
\definecolor{currentfill}{rgb}{0.909804,0.698039,0.423529}%
\pgfsetfillcolor{currentfill}%
\pgfsetfillopacity{0.120000}%
\pgfsetlinewidth{0.000000pt}%
\definecolor{currentstroke}{rgb}{0.000000,0.000000,0.000000}%
\pgfsetstrokecolor{currentstroke}%
\pgfsetdash{}{0pt}%
\pgfpathmoveto{\pgfqpoint{1.121212in}{0.924452in}}%
\pgfpathlineto{\pgfqpoint{1.137422in}{0.920185in}}%
\pgfpathlineto{\pgfqpoint{1.153672in}{0.916911in}}%
\pgfpathlineto{\pgfqpoint{1.161624in}{0.925338in}}%
\pgfpathlineto{\pgfqpoint{1.169555in}{0.934259in}}%
\pgfpathlineto{\pgfqpoint{1.153328in}{0.937520in}}%
\pgfpathlineto{\pgfqpoint{1.137142in}{0.941748in}}%
\pgfpathlineto{\pgfqpoint{1.129188in}{0.932854in}}%
\pgfpathlineto{\pgfqpoint{1.121212in}{0.924452in}}%
\pgfpathclose%
\pgfusepath{fill}%
\end{pgfscope}%
\begin{pgfscope}%
\pgfpathrectangle{\pgfqpoint{0.407592in}{0.222904in}}{\pgfqpoint{1.411163in}{1.411163in}}%
\pgfusepath{clip}%
\pgfsetbuttcap%
\pgfsetroundjoin%
\definecolor{currentfill}{rgb}{0.909804,0.698039,0.423529}%
\pgfsetfillcolor{currentfill}%
\pgfsetfillopacity{0.120000}%
\pgfsetlinewidth{0.000000pt}%
\definecolor{currentstroke}{rgb}{0.000000,0.000000,0.000000}%
\pgfsetstrokecolor{currentstroke}%
\pgfsetdash{}{0pt}%
\pgfpathmoveto{\pgfqpoint{1.378773in}{1.047207in}}%
\pgfpathlineto{\pgfqpoint{1.395393in}{1.050527in}}%
\pgfpathlineto{\pgfqpoint{1.412077in}{1.054448in}}%
\pgfpathlineto{\pgfqpoint{1.419742in}{1.066533in}}%
\pgfpathlineto{\pgfqpoint{1.427394in}{1.078857in}}%
\pgfpathlineto{\pgfqpoint{1.410727in}{1.075181in}}%
\pgfpathlineto{\pgfqpoint{1.394123in}{1.072094in}}%
\pgfpathlineto{\pgfqpoint{1.386455in}{1.059531in}}%
\pgfpathlineto{\pgfqpoint{1.378773in}{1.047207in}}%
\pgfpathclose%
\pgfusepath{fill}%
\end{pgfscope}%
\begin{pgfscope}%
\pgfpathrectangle{\pgfqpoint{0.407592in}{0.222904in}}{\pgfqpoint{1.411163in}{1.411163in}}%
\pgfusepath{clip}%
\pgfsetbuttcap%
\pgfsetroundjoin%
\definecolor{currentfill}{rgb}{0.909804,0.698039,0.423529}%
\pgfsetfillcolor{currentfill}%
\pgfsetfillopacity{0.120000}%
\pgfsetlinewidth{0.000000pt}%
\definecolor{currentstroke}{rgb}{0.000000,0.000000,0.000000}%
\pgfsetstrokecolor{currentstroke}%
\pgfsetdash{}{0pt}%
\pgfpathmoveto{\pgfqpoint{1.491333in}{1.138269in}}%
\pgfpathlineto{\pgfqpoint{1.508161in}{1.143015in}}%
\pgfpathlineto{\pgfqpoint{1.525056in}{1.148233in}}%
\pgfpathlineto{\pgfqpoint{1.532611in}{1.161047in}}%
\pgfpathlineto{\pgfqpoint{1.540154in}{1.174033in}}%
\pgfpathlineto{\pgfqpoint{1.523277in}{1.169078in}}%
\pgfpathlineto{\pgfqpoint{1.506465in}{1.164588in}}%
\pgfpathlineto{\pgfqpoint{1.498905in}{1.151344in}}%
\pgfpathlineto{\pgfqpoint{1.491333in}{1.138269in}}%
\pgfpathclose%
\pgfusepath{fill}%
\end{pgfscope}%
\begin{pgfscope}%
\pgfpathrectangle{\pgfqpoint{0.407592in}{0.222904in}}{\pgfqpoint{1.411163in}{1.411163in}}%
\pgfusepath{clip}%
\pgfsetbuttcap%
\pgfsetroundjoin%
\definecolor{currentfill}{rgb}{0.909804,0.698039,0.423529}%
\pgfsetfillcolor{currentfill}%
\pgfsetfillopacity{0.120000}%
\pgfsetlinewidth{0.000000pt}%
\definecolor{currentstroke}{rgb}{0.000000,0.000000,0.000000}%
\pgfsetstrokecolor{currentstroke}%
\pgfsetdash{}{0pt}%
\pgfpathmoveto{\pgfqpoint{0.749382in}{1.085557in}}%
\pgfpathlineto{\pgfqpoint{0.765412in}{1.069147in}}%
\pgfpathlineto{\pgfqpoint{0.781452in}{1.053047in}}%
\pgfpathlineto{\pgfqpoint{0.790000in}{1.056233in}}%
\pgfpathlineto{\pgfqpoint{0.798513in}{1.059814in}}%
\pgfpathlineto{\pgfqpoint{0.782515in}{1.075944in}}%
\pgfpathlineto{\pgfqpoint{0.766526in}{1.092380in}}%
\pgfpathlineto{\pgfqpoint{0.757972in}{1.088783in}}%
\pgfpathlineto{\pgfqpoint{0.749382in}{1.085557in}}%
\pgfpathclose%
\pgfusepath{fill}%
\end{pgfscope}%
\begin{pgfscope}%
\pgfpathrectangle{\pgfqpoint{0.407592in}{0.222904in}}{\pgfqpoint{1.411163in}{1.411163in}}%
\pgfusepath{clip}%
\pgfsetbuttcap%
\pgfsetroundjoin%
\definecolor{currentfill}{rgb}{0.909804,0.698039,0.423529}%
\pgfsetfillcolor{currentfill}%
\pgfsetfillopacity{0.120000}%
\pgfsetlinewidth{0.000000pt}%
\definecolor{currentstroke}{rgb}{0.000000,0.000000,0.000000}%
\pgfsetstrokecolor{currentstroke}%
\pgfsetdash{}{0pt}%
\pgfpathmoveto{\pgfqpoint{1.217906in}{0.949774in}}%
\pgfpathlineto{\pgfqpoint{1.234252in}{0.949277in}}%
\pgfpathlineto{\pgfqpoint{1.250651in}{0.949655in}}%
\pgfpathlineto{\pgfqpoint{1.258486in}{0.959782in}}%
\pgfpathlineto{\pgfqpoint{1.266304in}{0.970304in}}%
\pgfpathlineto{\pgfqpoint{1.249924in}{0.970040in}}%
\pgfpathlineto{\pgfqpoint{1.233597in}{0.970624in}}%
\pgfpathlineto{\pgfqpoint{1.225760in}{0.959999in}}%
\pgfpathlineto{\pgfqpoint{1.217906in}{0.949774in}}%
\pgfpathclose%
\pgfusepath{fill}%
\end{pgfscope}%
\begin{pgfscope}%
\pgfpathrectangle{\pgfqpoint{0.407592in}{0.222904in}}{\pgfqpoint{1.411163in}{1.411163in}}%
\pgfusepath{clip}%
\pgfsetbuttcap%
\pgfsetroundjoin%
\definecolor{currentfill}{rgb}{0.909804,0.698039,0.423529}%
\pgfsetfillcolor{currentfill}%
\pgfsetfillopacity{0.120000}%
\pgfsetlinewidth{0.000000pt}%
\definecolor{currentstroke}{rgb}{0.000000,0.000000,0.000000}%
\pgfsetstrokecolor{currentstroke}%
\pgfsetdash{}{0pt}%
\pgfpathmoveto{\pgfqpoint{1.072829in}{0.920814in}}%
\pgfpathlineto{\pgfqpoint{1.088994in}{0.914513in}}%
\pgfpathlineto{\pgfqpoint{1.105193in}{0.909202in}}%
\pgfpathlineto{\pgfqpoint{1.113214in}{0.916562in}}%
\pgfpathlineto{\pgfqpoint{1.121212in}{0.924452in}}%
\pgfpathlineto{\pgfqpoint{1.105040in}{0.929706in}}%
\pgfpathlineto{\pgfqpoint{1.088902in}{0.935933in}}%
\pgfpathlineto{\pgfqpoint{1.080878in}{0.928112in}}%
\pgfpathlineto{\pgfqpoint{1.072829in}{0.920814in}}%
\pgfpathclose%
\pgfusepath{fill}%
\end{pgfscope}%
\begin{pgfscope}%
\pgfpathrectangle{\pgfqpoint{0.407592in}{0.222904in}}{\pgfqpoint{1.411163in}{1.411163in}}%
\pgfusepath{clip}%
\pgfsetbuttcap%
\pgfsetroundjoin%
\definecolor{currentfill}{rgb}{0.909804,0.698039,0.423529}%
\pgfsetfillcolor{currentfill}%
\pgfsetfillopacity{0.120000}%
\pgfsetlinewidth{0.000000pt}%
\definecolor{currentstroke}{rgb}{0.000000,0.000000,0.000000}%
\pgfsetstrokecolor{currentstroke}%
\pgfsetdash{}{0pt}%
\pgfpathmoveto{\pgfqpoint{0.943557in}{0.954581in}}%
\pgfpathlineto{\pgfqpoint{0.959640in}{0.942964in}}%
\pgfpathlineto{\pgfqpoint{0.975742in}{0.932092in}}%
\pgfpathlineto{\pgfqpoint{0.983963in}{0.937106in}}%
\pgfpathlineto{\pgfqpoint{0.992156in}{0.942647in}}%
\pgfpathlineto{\pgfqpoint{0.976089in}{0.953465in}}%
\pgfpathlineto{\pgfqpoint{0.960042in}{0.965026in}}%
\pgfpathlineto{\pgfqpoint{0.951814in}{0.959552in}}%
\pgfpathlineto{\pgfqpoint{0.943557in}{0.954581in}}%
\pgfpathclose%
\pgfusepath{fill}%
\end{pgfscope}%
\begin{pgfscope}%
\pgfpathrectangle{\pgfqpoint{0.407592in}{0.222904in}}{\pgfqpoint{1.411163in}{1.411163in}}%
\pgfusepath{clip}%
\pgfsetbuttcap%
\pgfsetroundjoin%
\definecolor{currentfill}{rgb}{0.909804,0.698039,0.423529}%
\pgfsetfillcolor{currentfill}%
\pgfsetfillopacity{0.120000}%
\pgfsetlinewidth{0.000000pt}%
\definecolor{currentstroke}{rgb}{0.000000,0.000000,0.000000}%
\pgfsetstrokecolor{currentstroke}%
\pgfsetdash{}{0pt}%
\pgfpathmoveto{\pgfqpoint{0.635876in}{1.183500in}}%
\pgfpathlineto{\pgfqpoint{0.651888in}{1.165765in}}%
\pgfpathlineto{\pgfqpoint{0.667910in}{1.148213in}}%
\pgfpathlineto{\pgfqpoint{0.676656in}{1.150838in}}%
\pgfpathlineto{\pgfqpoint{0.685366in}{1.153789in}}%
\pgfpathlineto{\pgfqpoint{0.669386in}{1.171385in}}%
\pgfpathlineto{\pgfqpoint{0.653415in}{1.189160in}}%
\pgfpathlineto{\pgfqpoint{0.644664in}{1.186175in}}%
\pgfpathlineto{\pgfqpoint{0.635876in}{1.183500in}}%
\pgfpathclose%
\pgfusepath{fill}%
\end{pgfscope}%
\begin{pgfscope}%
\pgfpathrectangle{\pgfqpoint{0.407592in}{0.222904in}}{\pgfqpoint{1.411163in}{1.411163in}}%
\pgfusepath{clip}%
\pgfsetbuttcap%
\pgfsetroundjoin%
\definecolor{currentfill}{rgb}{0.909804,0.698039,0.423529}%
\pgfsetfillcolor{currentfill}%
\pgfsetfillopacity{0.120000}%
\pgfsetlinewidth{0.000000pt}%
\definecolor{currentstroke}{rgb}{0.000000,0.000000,0.000000}%
\pgfsetstrokecolor{currentstroke}%
\pgfsetdash{}{0pt}%
\pgfpathmoveto{\pgfqpoint{0.862624in}{0.999025in}}%
\pgfpathlineto{\pgfqpoint{0.878683in}{0.984951in}}%
\pgfpathlineto{\pgfqpoint{0.894755in}{0.971414in}}%
\pgfpathlineto{\pgfqpoint{0.903115in}{0.975386in}}%
\pgfpathlineto{\pgfqpoint{0.911442in}{0.979834in}}%
\pgfpathlineto{\pgfqpoint{0.895408in}{0.993364in}}%
\pgfpathlineto{\pgfqpoint{0.879388in}{1.007429in}}%
\pgfpathlineto{\pgfqpoint{0.871023in}{1.003002in}}%
\pgfpathlineto{\pgfqpoint{0.862624in}{0.999025in}}%
\pgfpathclose%
\pgfusepath{fill}%
\end{pgfscope}%
\begin{pgfscope}%
\pgfpathrectangle{\pgfqpoint{0.407592in}{0.222904in}}{\pgfqpoint{1.411163in}{1.411163in}}%
\pgfusepath{clip}%
\pgfsetbuttcap%
\pgfsetroundjoin%
\definecolor{currentfill}{rgb}{0.909804,0.698039,0.423529}%
\pgfsetfillcolor{currentfill}%
\pgfsetfillopacity{0.120000}%
\pgfsetlinewidth{0.000000pt}%
\definecolor{currentstroke}{rgb}{0.000000,0.000000,0.000000}%
\pgfsetstrokecolor{currentstroke}%
\pgfsetdash{}{0pt}%
\pgfpathmoveto{\pgfqpoint{1.266304in}{0.970304in}}%
\pgfpathlineto{\pgfqpoint{1.282739in}{0.971386in}}%
\pgfpathlineto{\pgfqpoint{1.299230in}{0.973254in}}%
\pgfpathlineto{\pgfqpoint{1.307015in}{0.983990in}}%
\pgfpathlineto{\pgfqpoint{1.314784in}{0.995073in}}%
\pgfpathlineto{\pgfqpoint{1.298310in}{0.993377in}}%
\pgfpathlineto{\pgfqpoint{1.281894in}{0.992444in}}%
\pgfpathlineto{\pgfqpoint{1.274107in}{0.981198in}}%
\pgfpathlineto{\pgfqpoint{1.266304in}{0.970304in}}%
\pgfpathclose%
\pgfusepath{fill}%
\end{pgfscope}%
\begin{pgfscope}%
\pgfpathrectangle{\pgfqpoint{0.407592in}{0.222904in}}{\pgfqpoint{1.411163in}{1.411163in}}%
\pgfusepath{clip}%
\pgfsetbuttcap%
\pgfsetroundjoin%
\definecolor{currentfill}{rgb}{0.909804,0.698039,0.423529}%
\pgfsetfillcolor{currentfill}%
\pgfsetfillopacity{0.120000}%
\pgfsetlinewidth{0.000000pt}%
\definecolor{currentstroke}{rgb}{0.000000,0.000000,0.000000}%
\pgfsetstrokecolor{currentstroke}%
\pgfsetdash{}{0pt}%
\pgfpathmoveto{\pgfqpoint{1.024355in}{0.923453in}}%
\pgfpathlineto{\pgfqpoint{1.040492in}{0.915178in}}%
\pgfpathlineto{\pgfqpoint{1.056656in}{0.907844in}}%
\pgfpathlineto{\pgfqpoint{1.064756in}{0.914054in}}%
\pgfpathlineto{\pgfqpoint{1.072829in}{0.920814in}}%
\pgfpathlineto{\pgfqpoint{1.056695in}{0.928081in}}%
\pgfpathlineto{\pgfqpoint{1.040590in}{0.936278in}}%
\pgfpathlineto{\pgfqpoint{1.032486in}{0.929598in}}%
\pgfpathlineto{\pgfqpoint{1.024355in}{0.923453in}}%
\pgfpathclose%
\pgfusepath{fill}%
\end{pgfscope}%
\begin{pgfscope}%
\pgfpathrectangle{\pgfqpoint{0.407592in}{0.222904in}}{\pgfqpoint{1.411163in}{1.411163in}}%
\pgfusepath{clip}%
\pgfsetbuttcap%
\pgfsetroundjoin%
\definecolor{currentfill}{rgb}{0.909804,0.698039,0.423529}%
\pgfsetfillcolor{currentfill}%
\pgfsetfillopacity{0.120000}%
\pgfsetlinewidth{0.000000pt}%
\definecolor{currentstroke}{rgb}{0.000000,0.000000,0.000000}%
\pgfsetstrokecolor{currentstroke}%
\pgfsetdash{}{0pt}%
\pgfpathmoveto{\pgfqpoint{1.427394in}{1.078857in}}%
\pgfpathlineto{\pgfqpoint{1.444124in}{1.083094in}}%
\pgfpathlineto{\pgfqpoint{1.460920in}{1.087866in}}%
\pgfpathlineto{\pgfqpoint{1.468542in}{1.100158in}}%
\pgfpathlineto{\pgfqpoint{1.476152in}{1.112664in}}%
\pgfpathlineto{\pgfqpoint{1.459373in}{1.108154in}}%
\pgfpathlineto{\pgfqpoint{1.442659in}{1.104169in}}%
\pgfpathlineto{\pgfqpoint{1.435033in}{1.091407in}}%
\pgfpathlineto{\pgfqpoint{1.427394in}{1.078857in}}%
\pgfpathclose%
\pgfusepath{fill}%
\end{pgfscope}%
\begin{pgfscope}%
\pgfpathrectangle{\pgfqpoint{0.407592in}{0.222904in}}{\pgfqpoint{1.411163in}{1.411163in}}%
\pgfusepath{clip}%
\pgfsetbuttcap%
\pgfsetroundjoin%
\definecolor{currentfill}{rgb}{0.909804,0.698039,0.423529}%
\pgfsetfillcolor{currentfill}%
\pgfsetfillopacity{0.120000}%
\pgfsetlinewidth{0.000000pt}%
\definecolor{currentstroke}{rgb}{0.000000,0.000000,0.000000}%
\pgfsetstrokecolor{currentstroke}%
\pgfsetdash{}{0pt}%
\pgfpathmoveto{\pgfqpoint{0.781452in}{1.053047in}}%
\pgfpathlineto{\pgfqpoint{0.797502in}{1.037288in}}%
\pgfpathlineto{\pgfqpoint{0.813564in}{1.021904in}}%
\pgfpathlineto{\pgfqpoint{0.822072in}{1.025044in}}%
\pgfpathlineto{\pgfqpoint{0.830545in}{1.028602in}}%
\pgfpathlineto{\pgfqpoint{0.814523in}{1.044023in}}%
\pgfpathlineto{\pgfqpoint{0.798513in}{1.059814in}}%
\pgfpathlineto{\pgfqpoint{0.790000in}{1.056233in}}%
\pgfpathlineto{\pgfqpoint{0.781452in}{1.053047in}}%
\pgfpathclose%
\pgfusepath{fill}%
\end{pgfscope}%
\begin{pgfscope}%
\pgfpathrectangle{\pgfqpoint{0.407592in}{0.222904in}}{\pgfqpoint{1.411163in}{1.411163in}}%
\pgfusepath{clip}%
\pgfsetbuttcap%
\pgfsetroundjoin%
\definecolor{currentfill}{rgb}{0.909804,0.698039,0.423529}%
\pgfsetfillcolor{currentfill}%
\pgfsetfillopacity{0.120000}%
\pgfsetlinewidth{0.000000pt}%
\definecolor{currentstroke}{rgb}{0.000000,0.000000,0.000000}%
\pgfsetstrokecolor{currentstroke}%
\pgfsetdash{}{0pt}%
\pgfpathmoveto{\pgfqpoint{0.667910in}{1.148213in}}%
\pgfpathlineto{\pgfqpoint{0.683941in}{1.130862in}}%
\pgfpathlineto{\pgfqpoint{0.699982in}{1.113730in}}%
\pgfpathlineto{\pgfqpoint{0.708687in}{1.116293in}}%
\pgfpathlineto{\pgfqpoint{0.717354in}{1.119202in}}%
\pgfpathlineto{\pgfqpoint{0.701355in}{1.136388in}}%
\pgfpathlineto{\pgfqpoint{0.685366in}{1.153789in}}%
\pgfpathlineto{\pgfqpoint{0.676656in}{1.150838in}}%
\pgfpathlineto{\pgfqpoint{0.667910in}{1.148213in}}%
\pgfpathclose%
\pgfusepath{fill}%
\end{pgfscope}%
\begin{pgfscope}%
\pgfpathrectangle{\pgfqpoint{0.407592in}{0.222904in}}{\pgfqpoint{1.411163in}{1.411163in}}%
\pgfusepath{clip}%
\pgfsetbuttcap%
\pgfsetroundjoin%
\definecolor{currentfill}{rgb}{0.909804,0.698039,0.423529}%
\pgfsetfillcolor{currentfill}%
\pgfsetfillopacity{0.120000}%
\pgfsetlinewidth{0.000000pt}%
\definecolor{currentstroke}{rgb}{0.000000,0.000000,0.000000}%
\pgfsetstrokecolor{currentstroke}%
\pgfsetdash{}{0pt}%
\pgfpathmoveto{\pgfqpoint{1.314784in}{0.995073in}}%
\pgfpathlineto{\pgfqpoint{1.331316in}{0.997499in}}%
\pgfpathlineto{\pgfqpoint{1.347909in}{1.000622in}}%
\pgfpathlineto{\pgfqpoint{1.355646in}{1.011827in}}%
\pgfpathlineto{\pgfqpoint{1.363369in}{1.023338in}}%
\pgfpathlineto{\pgfqpoint{1.346793in}{1.020431in}}%
\pgfpathlineto{\pgfqpoint{1.330278in}{1.018203in}}%
\pgfpathlineto{\pgfqpoint{1.322538in}{1.006484in}}%
\pgfpathlineto{\pgfqpoint{1.314784in}{0.995073in}}%
\pgfpathclose%
\pgfusepath{fill}%
\end{pgfscope}%
\begin{pgfscope}%
\pgfpathrectangle{\pgfqpoint{0.407592in}{0.222904in}}{\pgfqpoint{1.411163in}{1.411163in}}%
\pgfusepath{clip}%
\pgfsetbuttcap%
\pgfsetroundjoin%
\definecolor{currentfill}{rgb}{0.909804,0.698039,0.423529}%
\pgfsetfillcolor{currentfill}%
\pgfsetfillopacity{0.120000}%
\pgfsetlinewidth{0.000000pt}%
\definecolor{currentstroke}{rgb}{0.000000,0.000000,0.000000}%
\pgfsetstrokecolor{currentstroke}%
\pgfsetdash{}{0pt}%
\pgfpathmoveto{\pgfqpoint{1.540154in}{1.174033in}}%
\pgfpathlineto{\pgfqpoint{1.557099in}{1.179432in}}%
\pgfpathlineto{\pgfqpoint{1.574112in}{1.185256in}}%
\pgfpathlineto{\pgfqpoint{1.581626in}{1.198139in}}%
\pgfpathlineto{\pgfqpoint{1.564622in}{1.192448in}}%
\pgfpathlineto{\pgfqpoint{1.547685in}{1.187181in}}%
\pgfusepath{fill}%
\end{pgfscope}%
\begin{pgfscope}%
\pgfpathrectangle{\pgfqpoint{0.407592in}{0.222904in}}{\pgfqpoint{1.411163in}{1.411163in}}%
\pgfusepath{clip}%
\pgfsetbuttcap%
\pgfsetroundjoin%
\definecolor{currentfill}{rgb}{0.909804,0.698039,0.423529}%
\pgfsetfillcolor{currentfill}%
\pgfsetfillopacity{0.120000}%
\pgfsetlinewidth{0.000000pt}%
\definecolor{currentstroke}{rgb}{0.000000,0.000000,0.000000}%
\pgfsetstrokecolor{currentstroke}%
\pgfsetdash{}{0pt}%
\pgfpathmoveto{\pgfqpoint{0.894755in}{0.971414in}}%
\pgfpathlineto{\pgfqpoint{0.910843in}{0.958463in}}%
\pgfpathlineto{\pgfqpoint{0.926946in}{0.946149in}}%
\pgfpathlineto{\pgfqpoint{0.935267in}{0.950114in}}%
\pgfpathlineto{\pgfqpoint{0.943557in}{0.954581in}}%
\pgfpathlineto{\pgfqpoint{0.927491in}{0.966889in}}%
\pgfpathlineto{\pgfqpoint{0.911442in}{0.979834in}}%
\pgfpathlineto{\pgfqpoint{0.903115in}{0.975386in}}%
\pgfpathlineto{\pgfqpoint{0.894755in}{0.971414in}}%
\pgfpathclose%
\pgfusepath{fill}%
\end{pgfscope}%
\begin{pgfscope}%
\pgfpathrectangle{\pgfqpoint{0.407592in}{0.222904in}}{\pgfqpoint{1.411163in}{1.411163in}}%
\pgfusepath{clip}%
\pgfsetbuttcap%
\pgfsetroundjoin%
\definecolor{currentfill}{rgb}{0.909804,0.698039,0.423529}%
\pgfsetfillcolor{currentfill}%
\pgfsetfillopacity{0.120000}%
\pgfsetlinewidth{0.000000pt}%
\definecolor{currentstroke}{rgb}{0.000000,0.000000,0.000000}%
\pgfsetstrokecolor{currentstroke}%
\pgfsetdash{}{0pt}%
\pgfpathmoveto{\pgfqpoint{0.975742in}{0.932092in}}%
\pgfpathlineto{\pgfqpoint{0.991864in}{0.922019in}}%
\pgfpathlineto{\pgfqpoint{1.008010in}{0.912799in}}%
\pgfpathlineto{\pgfqpoint{1.016197in}{0.917852in}}%
\pgfpathlineto{\pgfqpoint{1.024355in}{0.923453in}}%
\pgfpathlineto{\pgfqpoint{1.008244in}{0.932626in}}%
\pgfpathlineto{\pgfqpoint{0.992156in}{0.942647in}}%
\pgfpathlineto{\pgfqpoint{0.983963in}{0.937106in}}%
\pgfpathlineto{\pgfqpoint{0.975742in}{0.932092in}}%
\pgfpathclose%
\pgfusepath{fill}%
\end{pgfscope}%
\begin{pgfscope}%
\pgfpathrectangle{\pgfqpoint{0.407592in}{0.222904in}}{\pgfqpoint{1.411163in}{1.411163in}}%
\pgfusepath{clip}%
\pgfsetbuttcap%
\pgfsetroundjoin%
\definecolor{currentfill}{rgb}{0.909804,0.698039,0.423529}%
\pgfsetfillcolor{currentfill}%
\pgfsetfillopacity{0.120000}%
\pgfsetlinewidth{0.000000pt}%
\definecolor{currentstroke}{rgb}{0.000000,0.000000,0.000000}%
\pgfsetstrokecolor{currentstroke}%
\pgfsetdash{}{0pt}%
\pgfpathmoveto{\pgfqpoint{1.476152in}{1.112664in}}%
\pgfpathlineto{\pgfqpoint{1.492997in}{1.117674in}}%
\pgfpathlineto{\pgfqpoint{1.509909in}{1.123162in}}%
\pgfpathlineto{\pgfqpoint{1.517489in}{1.135601in}}%
\pgfpathlineto{\pgfqpoint{1.525056in}{1.148233in}}%
\pgfpathlineto{\pgfqpoint{1.508161in}{1.143015in}}%
\pgfpathlineto{\pgfqpoint{1.491333in}{1.138269in}}%
\pgfpathlineto{\pgfqpoint{1.483749in}{1.125371in}}%
\pgfpathlineto{\pgfqpoint{1.476152in}{1.112664in}}%
\pgfpathclose%
\pgfusepath{fill}%
\end{pgfscope}%
\begin{pgfscope}%
\pgfpathrectangle{\pgfqpoint{0.407592in}{0.222904in}}{\pgfqpoint{1.411163in}{1.411163in}}%
\pgfusepath{clip}%
\pgfsetbuttcap%
\pgfsetroundjoin%
\definecolor{currentfill}{rgb}{0.909804,0.698039,0.423529}%
\pgfsetfillcolor{currentfill}%
\pgfsetfillopacity{0.120000}%
\pgfsetlinewidth{0.000000pt}%
\definecolor{currentstroke}{rgb}{0.000000,0.000000,0.000000}%
\pgfsetstrokecolor{currentstroke}%
\pgfsetdash{}{0pt}%
\pgfpathmoveto{\pgfqpoint{1.363369in}{1.023338in}}%
\pgfpathlineto{\pgfqpoint{1.380006in}{1.026891in}}%
\pgfpathlineto{\pgfqpoint{1.396706in}{1.031059in}}%
\pgfpathlineto{\pgfqpoint{1.404398in}{1.042618in}}%
\pgfpathlineto{\pgfqpoint{1.412077in}{1.054448in}}%
\pgfpathlineto{\pgfqpoint{1.395393in}{1.050527in}}%
\pgfpathlineto{\pgfqpoint{1.378773in}{1.047207in}}%
\pgfpathlineto{\pgfqpoint{1.371078in}{1.035137in}}%
\pgfpathlineto{\pgfqpoint{1.363369in}{1.023338in}}%
\pgfpathclose%
\pgfusepath{fill}%
\end{pgfscope}%
\begin{pgfscope}%
\pgfpathrectangle{\pgfqpoint{0.407592in}{0.222904in}}{\pgfqpoint{1.411163in}{1.411163in}}%
\pgfusepath{clip}%
\pgfsetbuttcap%
\pgfsetroundjoin%
\definecolor{currentfill}{rgb}{0.909804,0.698039,0.423529}%
\pgfsetfillcolor{currentfill}%
\pgfsetfillopacity{0.120000}%
\pgfsetlinewidth{0.000000pt}%
\definecolor{currentstroke}{rgb}{0.000000,0.000000,0.000000}%
\pgfsetstrokecolor{currentstroke}%
\pgfsetdash{}{0pt}%
\pgfpathmoveto{\pgfqpoint{0.699982in}{1.113730in}}%
\pgfpathlineto{\pgfqpoint{0.716033in}{1.096838in}}%
\pgfpathlineto{\pgfqpoint{0.732094in}{1.080209in}}%
\pgfpathlineto{\pgfqpoint{0.740756in}{1.082700in}}%
\pgfpathlineto{\pgfqpoint{0.749382in}{1.085557in}}%
\pgfpathlineto{\pgfqpoint{0.733363in}{1.102251in}}%
\pgfpathlineto{\pgfqpoint{0.717354in}{1.119202in}}%
\pgfpathlineto{\pgfqpoint{0.708687in}{1.116293in}}%
\pgfpathlineto{\pgfqpoint{0.699982in}{1.113730in}}%
\pgfpathclose%
\pgfusepath{fill}%
\end{pgfscope}%
\begin{pgfscope}%
\pgfpathrectangle{\pgfqpoint{0.407592in}{0.222904in}}{\pgfqpoint{1.411163in}{1.411163in}}%
\pgfusepath{clip}%
\pgfsetbuttcap%
\pgfsetroundjoin%
\definecolor{currentfill}{rgb}{0.909804,0.698039,0.423529}%
\pgfsetfillcolor{currentfill}%
\pgfsetfillopacity{0.120000}%
\pgfsetlinewidth{0.000000pt}%
\definecolor{currentstroke}{rgb}{0.000000,0.000000,0.000000}%
\pgfsetstrokecolor{currentstroke}%
\pgfsetdash{}{0pt}%
\pgfpathmoveto{\pgfqpoint{0.813564in}{1.021904in}}%
\pgfpathlineto{\pgfqpoint{0.829638in}{1.006931in}}%
\pgfpathlineto{\pgfqpoint{0.845724in}{0.992408in}}%
\pgfpathlineto{\pgfqpoint{0.854191in}{0.995495in}}%
\pgfpathlineto{\pgfqpoint{0.862624in}{0.999025in}}%
\pgfpathlineto{\pgfqpoint{0.846578in}{1.013590in}}%
\pgfpathlineto{\pgfqpoint{0.830545in}{1.028602in}}%
\pgfpathlineto{\pgfqpoint{0.822072in}{1.025044in}}%
\pgfpathlineto{\pgfqpoint{0.813564in}{1.021904in}}%
\pgfpathclose%
\pgfusepath{fill}%
\end{pgfscope}%
\begin{pgfscope}%
\pgfpathrectangle{\pgfqpoint{0.407592in}{0.222904in}}{\pgfqpoint{1.411163in}{1.411163in}}%
\pgfusepath{clip}%
\pgfsetbuttcap%
\pgfsetroundjoin%
\definecolor{currentfill}{rgb}{0.909804,0.698039,0.423529}%
\pgfsetfillcolor{currentfill}%
\pgfsetfillopacity{0.120000}%
\pgfsetlinewidth{0.000000pt}%
\definecolor{currentstroke}{rgb}{0.000000,0.000000,0.000000}%
\pgfsetstrokecolor{currentstroke}%
\pgfsetdash{}{0pt}%
\pgfpathmoveto{\pgfqpoint{1.153672in}{0.916911in}}%
\pgfpathlineto{\pgfqpoint{1.169966in}{0.914626in}}%
\pgfpathlineto{\pgfqpoint{1.186306in}{0.913316in}}%
\pgfpathlineto{\pgfqpoint{1.194234in}{0.921720in}}%
\pgfpathlineto{\pgfqpoint{1.202143in}{0.930612in}}%
\pgfpathlineto{\pgfqpoint{1.185826in}{0.931960in}}%
\pgfpathlineto{\pgfqpoint{1.169555in}{0.934259in}}%
\pgfpathlineto{\pgfqpoint{1.161624in}{0.925338in}}%
\pgfpathlineto{\pgfqpoint{1.153672in}{0.916911in}}%
\pgfpathclose%
\pgfusepath{fill}%
\end{pgfscope}%
\begin{pgfscope}%
\pgfpathrectangle{\pgfqpoint{0.407592in}{0.222904in}}{\pgfqpoint{1.411163in}{1.411163in}}%
\pgfusepath{clip}%
\pgfsetbuttcap%
\pgfsetroundjoin%
\definecolor{currentfill}{rgb}{0.909804,0.698039,0.423529}%
\pgfsetfillcolor{currentfill}%
\pgfsetfillopacity{0.120000}%
\pgfsetlinewidth{0.000000pt}%
\definecolor{currentstroke}{rgb}{0.000000,0.000000,0.000000}%
\pgfsetstrokecolor{currentstroke}%
\pgfsetdash{}{0pt}%
\pgfpathmoveto{\pgfqpoint{1.105193in}{0.909202in}}%
\pgfpathlineto{\pgfqpoint{1.121429in}{0.904899in}}%
\pgfpathlineto{\pgfqpoint{1.137705in}{0.901610in}}%
\pgfpathlineto{\pgfqpoint{1.145700in}{0.908995in}}%
\pgfpathlineto{\pgfqpoint{1.153672in}{0.916911in}}%
\pgfpathlineto{\pgfqpoint{1.137422in}{0.920185in}}%
\pgfpathlineto{\pgfqpoint{1.121212in}{0.924452in}}%
\pgfpathlineto{\pgfqpoint{1.113214in}{0.916562in}}%
\pgfpathlineto{\pgfqpoint{1.105193in}{0.909202in}}%
\pgfpathclose%
\pgfusepath{fill}%
\end{pgfscope}%
\begin{pgfscope}%
\pgfpathrectangle{\pgfqpoint{0.407592in}{0.222904in}}{\pgfqpoint{1.411163in}{1.411163in}}%
\pgfusepath{clip}%
\pgfsetbuttcap%
\pgfsetroundjoin%
\definecolor{currentfill}{rgb}{0.909804,0.698039,0.423529}%
\pgfsetfillcolor{currentfill}%
\pgfsetfillopacity{0.120000}%
\pgfsetlinewidth{0.000000pt}%
\definecolor{currentstroke}{rgb}{0.000000,0.000000,0.000000}%
\pgfsetstrokecolor{currentstroke}%
\pgfsetdash{}{0pt}%
\pgfpathmoveto{\pgfqpoint{1.202143in}{0.930612in}}%
\pgfpathlineto{\pgfqpoint{1.218510in}{0.930192in}}%
\pgfpathlineto{\pgfqpoint{1.234929in}{0.930672in}}%
\pgfpathlineto{\pgfqpoint{1.242799in}{0.939945in}}%
\pgfpathlineto{\pgfqpoint{1.250651in}{0.949655in}}%
\pgfpathlineto{\pgfqpoint{1.234252in}{0.949277in}}%
\pgfpathlineto{\pgfqpoint{1.217906in}{0.949774in}}%
\pgfpathlineto{\pgfqpoint{1.210034in}{0.939970in}}%
\pgfpathlineto{\pgfqpoint{1.202143in}{0.930612in}}%
\pgfpathclose%
\pgfusepath{fill}%
\end{pgfscope}%
\begin{pgfscope}%
\pgfpathrectangle{\pgfqpoint{0.407592in}{0.222904in}}{\pgfqpoint{1.411163in}{1.411163in}}%
\pgfusepath{clip}%
\pgfsetbuttcap%
\pgfsetroundjoin%
\definecolor{currentfill}{rgb}{0.909804,0.698039,0.423529}%
\pgfsetfillcolor{currentfill}%
\pgfsetfillopacity{0.120000}%
\pgfsetlinewidth{0.000000pt}%
\definecolor{currentstroke}{rgb}{0.000000,0.000000,0.000000}%
\pgfsetstrokecolor{currentstroke}%
\pgfsetdash{}{0pt}%
\pgfpathmoveto{\pgfqpoint{0.586104in}{1.214899in}}%
\pgfpathlineto{\pgfqpoint{0.602140in}{1.196904in}}%
\pgfpathlineto{\pgfqpoint{0.618185in}{1.179069in}}%
\pgfpathlineto{\pgfqpoint{0.627050in}{1.181132in}}%
\pgfpathlineto{\pgfqpoint{0.635876in}{1.183500in}}%
\pgfpathlineto{\pgfqpoint{0.619873in}{1.201403in}}%
\pgfpathlineto{\pgfqpoint{0.603879in}{1.219459in}}%
\pgfpathlineto{\pgfqpoint{0.595011in}{1.217035in}}%
\pgfpathlineto{\pgfqpoint{0.586104in}{1.214899in}}%
\pgfpathclose%
\pgfusepath{fill}%
\end{pgfscope}%
\begin{pgfscope}%
\pgfpathrectangle{\pgfqpoint{0.407592in}{0.222904in}}{\pgfqpoint{1.411163in}{1.411163in}}%
\pgfusepath{clip}%
\pgfsetbuttcap%
\pgfsetroundjoin%
\definecolor{currentfill}{rgb}{0.909804,0.698039,0.423529}%
\pgfsetfillcolor{currentfill}%
\pgfsetfillopacity{0.120000}%
\pgfsetlinewidth{0.000000pt}%
\definecolor{currentstroke}{rgb}{0.000000,0.000000,0.000000}%
\pgfsetstrokecolor{currentstroke}%
\pgfsetdash{}{0pt}%
\pgfpathmoveto{\pgfqpoint{1.250651in}{0.949655in}}%
\pgfpathlineto{\pgfqpoint{1.267104in}{0.950877in}}%
\pgfpathlineto{\pgfqpoint{1.283614in}{0.952905in}}%
\pgfpathlineto{\pgfqpoint{1.291430in}{0.962886in}}%
\pgfpathlineto{\pgfqpoint{1.299230in}{0.973254in}}%
\pgfpathlineto{\pgfqpoint{1.282739in}{0.971386in}}%
\pgfpathlineto{\pgfqpoint{1.266304in}{0.970304in}}%
\pgfpathlineto{\pgfqpoint{1.258486in}{0.959782in}}%
\pgfpathlineto{\pgfqpoint{1.250651in}{0.949655in}}%
\pgfpathclose%
\pgfusepath{fill}%
\end{pgfscope}%
\begin{pgfscope}%
\pgfpathrectangle{\pgfqpoint{0.407592in}{0.222904in}}{\pgfqpoint{1.411163in}{1.411163in}}%
\pgfusepath{clip}%
\pgfsetbuttcap%
\pgfsetroundjoin%
\definecolor{currentfill}{rgb}{0.909804,0.698039,0.423529}%
\pgfsetfillcolor{currentfill}%
\pgfsetfillopacity{0.120000}%
\pgfsetlinewidth{0.000000pt}%
\definecolor{currentstroke}{rgb}{0.000000,0.000000,0.000000}%
\pgfsetstrokecolor{currentstroke}%
\pgfsetdash{}{0pt}%
\pgfpathmoveto{\pgfqpoint{1.056656in}{0.907844in}}%
\pgfpathlineto{\pgfqpoint{1.072851in}{0.901485in}}%
\pgfpathlineto{\pgfqpoint{1.089079in}{0.896130in}}%
\pgfpathlineto{\pgfqpoint{1.097149in}{0.902387in}}%
\pgfpathlineto{\pgfqpoint{1.105193in}{0.909202in}}%
\pgfpathlineto{\pgfqpoint{1.088994in}{0.914513in}}%
\pgfpathlineto{\pgfqpoint{1.072829in}{0.920814in}}%
\pgfpathlineto{\pgfqpoint{1.064756in}{0.914054in}}%
\pgfpathlineto{\pgfqpoint{1.056656in}{0.907844in}}%
\pgfpathclose%
\pgfusepath{fill}%
\end{pgfscope}%
\begin{pgfscope}%
\pgfpathrectangle{\pgfqpoint{0.407592in}{0.222904in}}{\pgfqpoint{1.411163in}{1.411163in}}%
\pgfusepath{clip}%
\pgfsetbuttcap%
\pgfsetroundjoin%
\definecolor{currentfill}{rgb}{0.909804,0.698039,0.423529}%
\pgfsetfillcolor{currentfill}%
\pgfsetfillopacity{0.120000}%
\pgfsetlinewidth{0.000000pt}%
\definecolor{currentstroke}{rgb}{0.000000,0.000000,0.000000}%
\pgfsetstrokecolor{currentstroke}%
\pgfsetdash{}{0pt}%
\pgfpathmoveto{\pgfqpoint{1.412077in}{1.054448in}}%
\pgfpathlineto{\pgfqpoint{1.428824in}{1.058942in}}%
\pgfpathlineto{\pgfqpoint{1.445638in}{1.063977in}}%
\pgfpathlineto{\pgfqpoint{1.453286in}{1.075801in}}%
\pgfpathlineto{\pgfqpoint{1.460920in}{1.087866in}}%
\pgfpathlineto{\pgfqpoint{1.444124in}{1.083094in}}%
\pgfpathlineto{\pgfqpoint{1.427394in}{1.078857in}}%
\pgfpathlineto{\pgfqpoint{1.419742in}{1.066533in}}%
\pgfpathlineto{\pgfqpoint{1.412077in}{1.054448in}}%
\pgfpathclose%
\pgfusepath{fill}%
\end{pgfscope}%
\begin{pgfscope}%
\pgfpathrectangle{\pgfqpoint{0.407592in}{0.222904in}}{\pgfqpoint{1.411163in}{1.411163in}}%
\pgfusepath{clip}%
\pgfsetbuttcap%
\pgfsetroundjoin%
\definecolor{currentfill}{rgb}{0.909804,0.698039,0.423529}%
\pgfsetfillcolor{currentfill}%
\pgfsetfillopacity{0.120000}%
\pgfsetlinewidth{0.000000pt}%
\definecolor{currentstroke}{rgb}{0.000000,0.000000,0.000000}%
\pgfsetstrokecolor{currentstroke}%
\pgfsetdash{}{0pt}%
\pgfpathmoveto{\pgfqpoint{0.732094in}{1.080209in}}%
\pgfpathlineto{\pgfqpoint{0.748165in}{1.063868in}}%
\pgfpathlineto{\pgfqpoint{0.764246in}{1.047843in}}%
\pgfpathlineto{\pgfqpoint{0.772867in}{1.050252in}}%
\pgfpathlineto{\pgfqpoint{0.781452in}{1.053047in}}%
\pgfpathlineto{\pgfqpoint{0.765412in}{1.069147in}}%
\pgfpathlineto{\pgfqpoint{0.749382in}{1.085557in}}%
\pgfpathlineto{\pgfqpoint{0.740756in}{1.082700in}}%
\pgfpathlineto{\pgfqpoint{0.732094in}{1.080209in}}%
\pgfpathclose%
\pgfusepath{fill}%
\end{pgfscope}%
\begin{pgfscope}%
\pgfpathrectangle{\pgfqpoint{0.407592in}{0.222904in}}{\pgfqpoint{1.411163in}{1.411163in}}%
\pgfusepath{clip}%
\pgfsetbuttcap%
\pgfsetroundjoin%
\definecolor{currentfill}{rgb}{0.909804,0.698039,0.423529}%
\pgfsetfillcolor{currentfill}%
\pgfsetfillopacity{0.120000}%
\pgfsetlinewidth{0.000000pt}%
\definecolor{currentstroke}{rgb}{0.000000,0.000000,0.000000}%
\pgfsetstrokecolor{currentstroke}%
\pgfsetdash{}{0pt}%
\pgfpathmoveto{\pgfqpoint{0.926946in}{0.946149in}}%
\pgfpathlineto{\pgfqpoint{0.943067in}{0.934527in}}%
\pgfpathlineto{\pgfqpoint{0.959206in}{0.923650in}}%
\pgfpathlineto{\pgfqpoint{0.967489in}{0.927607in}}%
\pgfpathlineto{\pgfqpoint{0.975742in}{0.932092in}}%
\pgfpathlineto{\pgfqpoint{0.959640in}{0.942964in}}%
\pgfpathlineto{\pgfqpoint{0.943557in}{0.954581in}}%
\pgfpathlineto{\pgfqpoint{0.935267in}{0.950114in}}%
\pgfpathlineto{\pgfqpoint{0.926946in}{0.946149in}}%
\pgfpathclose%
\pgfusepath{fill}%
\end{pgfscope}%
\begin{pgfscope}%
\pgfpathrectangle{\pgfqpoint{0.407592in}{0.222904in}}{\pgfqpoint{1.411163in}{1.411163in}}%
\pgfusepath{clip}%
\pgfsetbuttcap%
\pgfsetroundjoin%
\definecolor{currentfill}{rgb}{0.909804,0.698039,0.423529}%
\pgfsetfillcolor{currentfill}%
\pgfsetfillopacity{0.120000}%
\pgfsetlinewidth{0.000000pt}%
\definecolor{currentstroke}{rgb}{0.000000,0.000000,0.000000}%
\pgfsetstrokecolor{currentstroke}%
\pgfsetdash{}{0pt}%
\pgfpathmoveto{\pgfqpoint{1.525056in}{1.148233in}}%
\pgfpathlineto{\pgfqpoint{1.542019in}{1.153900in}}%
\pgfpathlineto{\pgfqpoint{1.559050in}{1.159996in}}%
\pgfpathlineto{\pgfqpoint{1.566587in}{1.172539in}}%
\pgfpathlineto{\pgfqpoint{1.574112in}{1.185256in}}%
\pgfpathlineto{\pgfqpoint{1.557099in}{1.179432in}}%
\pgfpathlineto{\pgfqpoint{1.540154in}{1.174033in}}%
\pgfpathlineto{\pgfqpoint{1.532611in}{1.161047in}}%
\pgfpathlineto{\pgfqpoint{1.525056in}{1.148233in}}%
\pgfpathclose%
\pgfusepath{fill}%
\end{pgfscope}%
\begin{pgfscope}%
\pgfpathrectangle{\pgfqpoint{0.407592in}{0.222904in}}{\pgfqpoint{1.411163in}{1.411163in}}%
\pgfusepath{clip}%
\pgfsetbuttcap%
\pgfsetroundjoin%
\definecolor{currentfill}{rgb}{0.909804,0.698039,0.423529}%
\pgfsetfillcolor{currentfill}%
\pgfsetfillopacity{0.120000}%
\pgfsetlinewidth{0.000000pt}%
\definecolor{currentstroke}{rgb}{0.000000,0.000000,0.000000}%
\pgfsetstrokecolor{currentstroke}%
\pgfsetdash{}{0pt}%
\pgfpathmoveto{\pgfqpoint{0.845724in}{0.992408in}}%
\pgfpathlineto{\pgfqpoint{0.861823in}{0.978379in}}%
\pgfpathlineto{\pgfqpoint{0.877935in}{0.964887in}}%
\pgfpathlineto{\pgfqpoint{0.886362in}{0.967916in}}%
\pgfpathlineto{\pgfqpoint{0.894755in}{0.971414in}}%
\pgfpathlineto{\pgfqpoint{0.878683in}{0.984951in}}%
\pgfpathlineto{\pgfqpoint{0.862624in}{0.999025in}}%
\pgfpathlineto{\pgfqpoint{0.854191in}{0.995495in}}%
\pgfpathlineto{\pgfqpoint{0.845724in}{0.992408in}}%
\pgfpathclose%
\pgfusepath{fill}%
\end{pgfscope}%
\begin{pgfscope}%
\pgfpathrectangle{\pgfqpoint{0.407592in}{0.222904in}}{\pgfqpoint{1.411163in}{1.411163in}}%
\pgfusepath{clip}%
\pgfsetbuttcap%
\pgfsetroundjoin%
\definecolor{currentfill}{rgb}{0.909804,0.698039,0.423529}%
\pgfsetfillcolor{currentfill}%
\pgfsetfillopacity{0.120000}%
\pgfsetlinewidth{0.000000pt}%
\definecolor{currentstroke}{rgb}{0.000000,0.000000,0.000000}%
\pgfsetstrokecolor{currentstroke}%
\pgfsetdash{}{0pt}%
\pgfpathmoveto{\pgfqpoint{0.618185in}{1.179069in}}%
\pgfpathlineto{\pgfqpoint{0.634240in}{1.161406in}}%
\pgfpathlineto{\pgfqpoint{0.650304in}{1.143933in}}%
\pgfpathlineto{\pgfqpoint{0.659126in}{1.145913in}}%
\pgfpathlineto{\pgfqpoint{0.667910in}{1.148213in}}%
\pgfpathlineto{\pgfqpoint{0.651888in}{1.165765in}}%
\pgfpathlineto{\pgfqpoint{0.635876in}{1.183500in}}%
\pgfpathlineto{\pgfqpoint{0.627050in}{1.181132in}}%
\pgfpathlineto{\pgfqpoint{0.618185in}{1.179069in}}%
\pgfpathclose%
\pgfusepath{fill}%
\end{pgfscope}%
\begin{pgfscope}%
\pgfpathrectangle{\pgfqpoint{0.407592in}{0.222904in}}{\pgfqpoint{1.411163in}{1.411163in}}%
\pgfusepath{clip}%
\pgfsetbuttcap%
\pgfsetroundjoin%
\definecolor{currentfill}{rgb}{0.909804,0.698039,0.423529}%
\pgfsetfillcolor{currentfill}%
\pgfsetfillopacity{0.120000}%
\pgfsetlinewidth{0.000000pt}%
\definecolor{currentstroke}{rgb}{0.000000,0.000000,0.000000}%
\pgfsetstrokecolor{currentstroke}%
\pgfsetdash{}{0pt}%
\pgfpathmoveto{\pgfqpoint{1.299230in}{0.973254in}}%
\pgfpathlineto{\pgfqpoint{1.315780in}{0.975872in}}%
\pgfpathlineto{\pgfqpoint{1.332390in}{0.979202in}}%
\pgfpathlineto{\pgfqpoint{1.340157in}{0.989741in}}%
\pgfpathlineto{\pgfqpoint{1.347909in}{1.000622in}}%
\pgfpathlineto{\pgfqpoint{1.331316in}{0.997499in}}%
\pgfpathlineto{\pgfqpoint{1.314784in}{0.995073in}}%
\pgfpathlineto{\pgfqpoint{1.307015in}{0.983990in}}%
\pgfpathlineto{\pgfqpoint{1.299230in}{0.973254in}}%
\pgfpathclose%
\pgfusepath{fill}%
\end{pgfscope}%
\begin{pgfscope}%
\pgfpathrectangle{\pgfqpoint{0.407592in}{0.222904in}}{\pgfqpoint{1.411163in}{1.411163in}}%
\pgfusepath{clip}%
\pgfsetbuttcap%
\pgfsetroundjoin%
\definecolor{currentfill}{rgb}{0.909804,0.698039,0.423529}%
\pgfsetfillcolor{currentfill}%
\pgfsetfillopacity{0.120000}%
\pgfsetlinewidth{0.000000pt}%
\definecolor{currentstroke}{rgb}{0.000000,0.000000,0.000000}%
\pgfsetstrokecolor{currentstroke}%
\pgfsetdash{}{0pt}%
\pgfpathmoveto{\pgfqpoint{1.008010in}{0.912799in}}%
\pgfpathlineto{\pgfqpoint{1.024180in}{0.904479in}}%
\pgfpathlineto{\pgfqpoint{1.040377in}{0.897104in}}%
\pgfpathlineto{\pgfqpoint{1.048530in}{0.902192in}}%
\pgfpathlineto{\pgfqpoint{1.056656in}{0.907844in}}%
\pgfpathlineto{\pgfqpoint{1.040492in}{0.915178in}}%
\pgfpathlineto{\pgfqpoint{1.024355in}{0.923453in}}%
\pgfpathlineto{\pgfqpoint{1.016197in}{0.917852in}}%
\pgfpathlineto{\pgfqpoint{1.008010in}{0.912799in}}%
\pgfpathclose%
\pgfusepath{fill}%
\end{pgfscope}%
\begin{pgfscope}%
\pgfpathrectangle{\pgfqpoint{0.407592in}{0.222904in}}{\pgfqpoint{1.411163in}{1.411163in}}%
\pgfusepath{clip}%
\pgfsetbuttcap%
\pgfsetroundjoin%
\definecolor{currentfill}{rgb}{0.909804,0.698039,0.423529}%
\pgfsetfillcolor{currentfill}%
\pgfsetfillopacity{0.120000}%
\pgfsetlinewidth{0.000000pt}%
\definecolor{currentstroke}{rgb}{0.000000,0.000000,0.000000}%
\pgfsetstrokecolor{currentstroke}%
\pgfsetdash{}{0pt}%
\pgfpathmoveto{\pgfqpoint{1.460920in}{1.087866in}}%
\pgfpathlineto{\pgfqpoint{1.477783in}{1.093145in}}%
\pgfpathlineto{\pgfqpoint{1.494713in}{1.098906in}}%
\pgfpathlineto{\pgfqpoint{1.502317in}{1.110926in}}%
\pgfpathlineto{\pgfqpoint{1.509909in}{1.123162in}}%
\pgfpathlineto{\pgfqpoint{1.492997in}{1.117674in}}%
\pgfpathlineto{\pgfqpoint{1.476152in}{1.112664in}}%
\pgfpathlineto{\pgfqpoint{1.468542in}{1.100158in}}%
\pgfpathlineto{\pgfqpoint{1.460920in}{1.087866in}}%
\pgfpathclose%
\pgfusepath{fill}%
\end{pgfscope}%
\begin{pgfscope}%
\pgfpathrectangle{\pgfqpoint{0.407592in}{0.222904in}}{\pgfqpoint{1.411163in}{1.411163in}}%
\pgfusepath{clip}%
\pgfsetbuttcap%
\pgfsetroundjoin%
\definecolor{currentfill}{rgb}{0.909804,0.698039,0.423529}%
\pgfsetfillcolor{currentfill}%
\pgfsetfillopacity{0.120000}%
\pgfsetlinewidth{0.000000pt}%
\definecolor{currentstroke}{rgb}{0.000000,0.000000,0.000000}%
\pgfsetstrokecolor{currentstroke}%
\pgfsetdash{}{0pt}%
\pgfpathmoveto{\pgfqpoint{0.764246in}{1.047843in}}%
\pgfpathlineto{\pgfqpoint{0.780338in}{1.032164in}}%
\pgfpathlineto{\pgfqpoint{0.796441in}{1.016864in}}%
\pgfpathlineto{\pgfqpoint{0.805021in}{1.019179in}}%
\pgfpathlineto{\pgfqpoint{0.813564in}{1.021904in}}%
\pgfpathlineto{\pgfqpoint{0.797502in}{1.037288in}}%
\pgfpathlineto{\pgfqpoint{0.781452in}{1.053047in}}%
\pgfpathlineto{\pgfqpoint{0.772867in}{1.050252in}}%
\pgfpathlineto{\pgfqpoint{0.764246in}{1.047843in}}%
\pgfpathclose%
\pgfusepath{fill}%
\end{pgfscope}%
\begin{pgfscope}%
\pgfpathrectangle{\pgfqpoint{0.407592in}{0.222904in}}{\pgfqpoint{1.411163in}{1.411163in}}%
\pgfusepath{clip}%
\pgfsetbuttcap%
\pgfsetroundjoin%
\definecolor{currentfill}{rgb}{0.909804,0.698039,0.423529}%
\pgfsetfillcolor{currentfill}%
\pgfsetfillopacity{0.120000}%
\pgfsetlinewidth{0.000000pt}%
\definecolor{currentstroke}{rgb}{0.000000,0.000000,0.000000}%
\pgfsetstrokecolor{currentstroke}%
\pgfsetdash{}{0pt}%
\pgfpathmoveto{\pgfqpoint{1.347909in}{1.000622in}}%
\pgfpathlineto{\pgfqpoint{1.364563in}{1.004406in}}%
\pgfpathlineto{\pgfqpoint{1.381282in}{1.008814in}}%
\pgfpathlineto{\pgfqpoint{1.389001in}{1.019785in}}%
\pgfpathlineto{\pgfqpoint{1.396706in}{1.031059in}}%
\pgfpathlineto{\pgfqpoint{1.380006in}{1.026891in}}%
\pgfpathlineto{\pgfqpoint{1.363369in}{1.023338in}}%
\pgfpathlineto{\pgfqpoint{1.355646in}{1.011827in}}%
\pgfpathlineto{\pgfqpoint{1.347909in}{1.000622in}}%
\pgfpathclose%
\pgfusepath{fill}%
\end{pgfscope}%
\begin{pgfscope}%
\pgfpathrectangle{\pgfqpoint{0.407592in}{0.222904in}}{\pgfqpoint{1.411163in}{1.411163in}}%
\pgfusepath{clip}%
\pgfsetbuttcap%
\pgfsetroundjoin%
\definecolor{currentfill}{rgb}{0.909804,0.698039,0.423529}%
\pgfsetfillcolor{currentfill}%
\pgfsetfillopacity{0.120000}%
\pgfsetlinewidth{0.000000pt}%
\definecolor{currentstroke}{rgb}{0.000000,0.000000,0.000000}%
\pgfsetstrokecolor{currentstroke}%
\pgfsetdash{}{0pt}%
\pgfpathmoveto{\pgfqpoint{0.650304in}{1.143933in}}%
\pgfpathlineto{\pgfqpoint{0.666377in}{1.126667in}}%
\pgfpathlineto{\pgfqpoint{0.682460in}{1.109626in}}%
\pgfpathlineto{\pgfqpoint{0.691240in}{1.111509in}}%
\pgfpathlineto{\pgfqpoint{0.699982in}{1.113730in}}%
\pgfpathlineto{\pgfqpoint{0.683941in}{1.130862in}}%
\pgfpathlineto{\pgfqpoint{0.667910in}{1.148213in}}%
\pgfpathlineto{\pgfqpoint{0.659126in}{1.145913in}}%
\pgfpathlineto{\pgfqpoint{0.650304in}{1.143933in}}%
\pgfpathclose%
\pgfusepath{fill}%
\end{pgfscope}%
\begin{pgfscope}%
\pgfpathrectangle{\pgfqpoint{0.407592in}{0.222904in}}{\pgfqpoint{1.411163in}{1.411163in}}%
\pgfusepath{clip}%
\pgfsetbuttcap%
\pgfsetroundjoin%
\definecolor{currentfill}{rgb}{0.909804,0.698039,0.423529}%
\pgfsetfillcolor{currentfill}%
\pgfsetfillopacity{0.120000}%
\pgfsetlinewidth{0.000000pt}%
\definecolor{currentstroke}{rgb}{0.000000,0.000000,0.000000}%
\pgfsetstrokecolor{currentstroke}%
\pgfsetdash{}{0pt}%
\pgfpathmoveto{\pgfqpoint{0.877935in}{0.964887in}}%
\pgfpathlineto{\pgfqpoint{0.894062in}{0.951983in}}%
\pgfpathlineto{\pgfqpoint{0.910205in}{0.939716in}}%
\pgfpathlineto{\pgfqpoint{0.918592in}{0.942685in}}%
\pgfpathlineto{\pgfqpoint{0.926946in}{0.946149in}}%
\pgfpathlineto{\pgfqpoint{0.910843in}{0.958463in}}%
\pgfpathlineto{\pgfqpoint{0.894755in}{0.971414in}}%
\pgfpathlineto{\pgfqpoint{0.886362in}{0.967916in}}%
\pgfpathlineto{\pgfqpoint{0.877935in}{0.964887in}}%
\pgfpathclose%
\pgfusepath{fill}%
\end{pgfscope}%
\begin{pgfscope}%
\pgfpathrectangle{\pgfqpoint{0.407592in}{0.222904in}}{\pgfqpoint{1.411163in}{1.411163in}}%
\pgfusepath{clip}%
\pgfsetbuttcap%
\pgfsetroundjoin%
\definecolor{currentfill}{rgb}{0.909804,0.698039,0.423529}%
\pgfsetfillcolor{currentfill}%
\pgfsetfillopacity{0.120000}%
\pgfsetlinewidth{0.000000pt}%
\definecolor{currentstroke}{rgb}{0.000000,0.000000,0.000000}%
\pgfsetstrokecolor{currentstroke}%
\pgfsetdash{}{0pt}%
\pgfpathmoveto{\pgfqpoint{0.959206in}{0.923650in}}%
\pgfpathlineto{\pgfqpoint{0.975365in}{0.913571in}}%
\pgfpathlineto{\pgfqpoint{0.991546in}{0.904343in}}%
\pgfpathlineto{\pgfqpoint{0.999793in}{0.908296in}}%
\pgfpathlineto{\pgfqpoint{1.008010in}{0.912799in}}%
\pgfpathlineto{\pgfqpoint{0.991864in}{0.922019in}}%
\pgfpathlineto{\pgfqpoint{0.975742in}{0.932092in}}%
\pgfpathlineto{\pgfqpoint{0.967489in}{0.927607in}}%
\pgfpathlineto{\pgfqpoint{0.959206in}{0.923650in}}%
\pgfpathclose%
\pgfusepath{fill}%
\end{pgfscope}%
\begin{pgfscope}%
\pgfpathrectangle{\pgfqpoint{0.407592in}{0.222904in}}{\pgfqpoint{1.411163in}{1.411163in}}%
\pgfusepath{clip}%
\pgfsetbuttcap%
\pgfsetroundjoin%
\definecolor{currentfill}{rgb}{0.909804,0.698039,0.423529}%
\pgfsetfillcolor{currentfill}%
\pgfsetfillopacity{0.120000}%
\pgfsetlinewidth{0.000000pt}%
\definecolor{currentstroke}{rgb}{0.000000,0.000000,0.000000}%
\pgfsetstrokecolor{currentstroke}%
\pgfsetdash{}{0pt}%
\pgfpathmoveto{\pgfqpoint{1.137705in}{0.901610in}}%
\pgfpathlineto{\pgfqpoint{1.154023in}{0.899330in}}%
\pgfpathlineto{\pgfqpoint{1.170387in}{0.898045in}}%
\pgfpathlineto{\pgfqpoint{1.178357in}{0.905419in}}%
\pgfpathlineto{\pgfqpoint{1.186306in}{0.913316in}}%
\pgfpathlineto{\pgfqpoint{1.169966in}{0.914626in}}%
\pgfpathlineto{\pgfqpoint{1.153672in}{0.916911in}}%
\pgfpathlineto{\pgfqpoint{1.145700in}{0.908995in}}%
\pgfpathlineto{\pgfqpoint{1.137705in}{0.901610in}}%
\pgfpathclose%
\pgfusepath{fill}%
\end{pgfscope}%
\begin{pgfscope}%
\pgfpathrectangle{\pgfqpoint{0.407592in}{0.222904in}}{\pgfqpoint{1.411163in}{1.411163in}}%
\pgfusepath{clip}%
\pgfsetbuttcap%
\pgfsetroundjoin%
\definecolor{currentfill}{rgb}{0.909804,0.698039,0.423529}%
\pgfsetfillcolor{currentfill}%
\pgfsetfillopacity{0.120000}%
\pgfsetlinewidth{0.000000pt}%
\definecolor{currentstroke}{rgb}{0.000000,0.000000,0.000000}%
\pgfsetstrokecolor{currentstroke}%
\pgfsetdash{}{0pt}%
\pgfpathmoveto{\pgfqpoint{1.186306in}{0.913316in}}%
\pgfpathlineto{\pgfqpoint{1.202694in}{0.912958in}}%
\pgfpathlineto{\pgfqpoint{1.219135in}{0.913520in}}%
\pgfpathlineto{\pgfqpoint{1.227041in}{0.921858in}}%
\pgfpathlineto{\pgfqpoint{1.234929in}{0.930672in}}%
\pgfpathlineto{\pgfqpoint{1.218510in}{0.930192in}}%
\pgfpathlineto{\pgfqpoint{1.202143in}{0.930612in}}%
\pgfpathlineto{\pgfqpoint{1.194234in}{0.921720in}}%
\pgfpathlineto{\pgfqpoint{1.186306in}{0.913316in}}%
\pgfpathclose%
\pgfusepath{fill}%
\end{pgfscope}%
\begin{pgfscope}%
\pgfpathrectangle{\pgfqpoint{0.407592in}{0.222904in}}{\pgfqpoint{1.411163in}{1.411163in}}%
\pgfusepath{clip}%
\pgfsetbuttcap%
\pgfsetroundjoin%
\definecolor{currentfill}{rgb}{0.909804,0.698039,0.423529}%
\pgfsetfillcolor{currentfill}%
\pgfsetfillopacity{0.120000}%
\pgfsetlinewidth{0.000000pt}%
\definecolor{currentstroke}{rgb}{0.000000,0.000000,0.000000}%
\pgfsetstrokecolor{currentstroke}%
\pgfsetdash{}{0pt}%
\pgfpathmoveto{\pgfqpoint{1.574112in}{1.185256in}}%
\pgfpathlineto{\pgfqpoint{1.591195in}{1.191486in}}%
\pgfpathlineto{\pgfqpoint{1.608348in}{1.198103in}}%
\pgfpathlineto{\pgfqpoint{1.615843in}{1.210716in}}%
\pgfpathlineto{\pgfqpoint{1.598699in}{1.204234in}}%
\pgfpathlineto{\pgfqpoint{1.581626in}{1.198139in}}%
\pgfusepath{fill}%
\end{pgfscope}%
\begin{pgfscope}%
\pgfpathrectangle{\pgfqpoint{0.407592in}{0.222904in}}{\pgfqpoint{1.411163in}{1.411163in}}%
\pgfusepath{clip}%
\pgfsetbuttcap%
\pgfsetroundjoin%
\definecolor{currentfill}{rgb}{0.909804,0.698039,0.423529}%
\pgfsetfillcolor{currentfill}%
\pgfsetfillopacity{0.120000}%
\pgfsetlinewidth{0.000000pt}%
\definecolor{currentstroke}{rgb}{0.000000,0.000000,0.000000}%
\pgfsetstrokecolor{currentstroke}%
\pgfsetdash{}{0pt}%
\pgfpathmoveto{\pgfqpoint{1.089079in}{0.896130in}}%
\pgfpathlineto{\pgfqpoint{1.105343in}{0.891796in}}%
\pgfpathlineto{\pgfqpoint{1.121647in}{0.888489in}}%
\pgfpathlineto{\pgfqpoint{1.129688in}{0.894771in}}%
\pgfpathlineto{\pgfqpoint{1.137705in}{0.901610in}}%
\pgfpathlineto{\pgfqpoint{1.121429in}{0.904899in}}%
\pgfpathlineto{\pgfqpoint{1.105193in}{0.909202in}}%
\pgfpathlineto{\pgfqpoint{1.097149in}{0.902387in}}%
\pgfpathlineto{\pgfqpoint{1.089079in}{0.896130in}}%
\pgfpathclose%
\pgfusepath{fill}%
\end{pgfscope}%
\begin{pgfscope}%
\pgfpathrectangle{\pgfqpoint{0.407592in}{0.222904in}}{\pgfqpoint{1.411163in}{1.411163in}}%
\pgfusepath{clip}%
\pgfsetbuttcap%
\pgfsetroundjoin%
\definecolor{currentfill}{rgb}{0.909804,0.698039,0.423529}%
\pgfsetfillcolor{currentfill}%
\pgfsetfillopacity{0.120000}%
\pgfsetlinewidth{0.000000pt}%
\definecolor{currentstroke}{rgb}{0.000000,0.000000,0.000000}%
\pgfsetstrokecolor{currentstroke}%
\pgfsetdash{}{0pt}%
\pgfpathmoveto{\pgfqpoint{1.396706in}{1.031059in}}%
\pgfpathlineto{\pgfqpoint{1.413472in}{1.035808in}}%
\pgfpathlineto{\pgfqpoint{1.430303in}{1.041105in}}%
\pgfpathlineto{\pgfqpoint{1.437977in}{1.052407in}}%
\pgfpathlineto{\pgfqpoint{1.445638in}{1.063977in}}%
\pgfpathlineto{\pgfqpoint{1.428824in}{1.058942in}}%
\pgfpathlineto{\pgfqpoint{1.412077in}{1.054448in}}%
\pgfpathlineto{\pgfqpoint{1.404398in}{1.042618in}}%
\pgfpathlineto{\pgfqpoint{1.396706in}{1.031059in}}%
\pgfpathclose%
\pgfusepath{fill}%
\end{pgfscope}%
\begin{pgfscope}%
\pgfpathrectangle{\pgfqpoint{0.407592in}{0.222904in}}{\pgfqpoint{1.411163in}{1.411163in}}%
\pgfusepath{clip}%
\pgfsetbuttcap%
\pgfsetroundjoin%
\definecolor{currentfill}{rgb}{0.909804,0.698039,0.423529}%
\pgfsetfillcolor{currentfill}%
\pgfsetfillopacity{0.120000}%
\pgfsetlinewidth{0.000000pt}%
\definecolor{currentstroke}{rgb}{0.000000,0.000000,0.000000}%
\pgfsetstrokecolor{currentstroke}%
\pgfsetdash{}{0pt}%
\pgfpathmoveto{\pgfqpoint{1.234929in}{0.930672in}}%
\pgfpathlineto{\pgfqpoint{1.251402in}{0.932017in}}%
\pgfpathlineto{\pgfqpoint{1.267932in}{0.934188in}}%
\pgfpathlineto{\pgfqpoint{1.275781in}{0.943333in}}%
\pgfpathlineto{\pgfqpoint{1.283614in}{0.952905in}}%
\pgfpathlineto{\pgfqpoint{1.267104in}{0.950877in}}%
\pgfpathlineto{\pgfqpoint{1.250651in}{0.949655in}}%
\pgfpathlineto{\pgfqpoint{1.242799in}{0.939945in}}%
\pgfpathlineto{\pgfqpoint{1.234929in}{0.930672in}}%
\pgfpathclose%
\pgfusepath{fill}%
\end{pgfscope}%
\begin{pgfscope}%
\pgfpathrectangle{\pgfqpoint{0.407592in}{0.222904in}}{\pgfqpoint{1.411163in}{1.411163in}}%
\pgfusepath{clip}%
\pgfsetbuttcap%
\pgfsetroundjoin%
\definecolor{currentfill}{rgb}{0.909804,0.698039,0.423529}%
\pgfsetfillcolor{currentfill}%
\pgfsetfillopacity{0.120000}%
\pgfsetlinewidth{0.000000pt}%
\definecolor{currentstroke}{rgb}{0.000000,0.000000,0.000000}%
\pgfsetstrokecolor{currentstroke}%
\pgfsetdash{}{0pt}%
\pgfpathmoveto{\pgfqpoint{1.509909in}{1.123162in}}%
\pgfpathlineto{\pgfqpoint{1.526890in}{1.129102in}}%
\pgfpathlineto{\pgfqpoint{1.543940in}{1.135473in}}%
\pgfpathlineto{\pgfqpoint{1.551501in}{1.147637in}}%
\pgfpathlineto{\pgfqpoint{1.559050in}{1.159996in}}%
\pgfpathlineto{\pgfqpoint{1.542019in}{1.153900in}}%
\pgfpathlineto{\pgfqpoint{1.525056in}{1.148233in}}%
\pgfpathlineto{\pgfqpoint{1.517489in}{1.135601in}}%
\pgfpathlineto{\pgfqpoint{1.509909in}{1.123162in}}%
\pgfpathclose%
\pgfusepath{fill}%
\end{pgfscope}%
\begin{pgfscope}%
\pgfpathrectangle{\pgfqpoint{0.407592in}{0.222904in}}{\pgfqpoint{1.411163in}{1.411163in}}%
\pgfusepath{clip}%
\pgfsetbuttcap%
\pgfsetroundjoin%
\definecolor{currentfill}{rgb}{0.909804,0.698039,0.423529}%
\pgfsetfillcolor{currentfill}%
\pgfsetfillopacity{0.120000}%
\pgfsetlinewidth{0.000000pt}%
\definecolor{currentstroke}{rgb}{0.000000,0.000000,0.000000}%
\pgfsetstrokecolor{currentstroke}%
\pgfsetdash{}{0pt}%
\pgfpathmoveto{\pgfqpoint{0.682460in}{1.109626in}}%
\pgfpathlineto{\pgfqpoint{0.698553in}{1.092831in}}%
\pgfpathlineto{\pgfqpoint{0.714655in}{1.076306in}}%
\pgfpathlineto{\pgfqpoint{0.723393in}{1.078079in}}%
\pgfpathlineto{\pgfqpoint{0.732094in}{1.080209in}}%
\pgfpathlineto{\pgfqpoint{0.716033in}{1.096838in}}%
\pgfpathlineto{\pgfqpoint{0.699982in}{1.113730in}}%
\pgfpathlineto{\pgfqpoint{0.691240in}{1.111509in}}%
\pgfpathlineto{\pgfqpoint{0.682460in}{1.109626in}}%
\pgfpathclose%
\pgfusepath{fill}%
\end{pgfscope}%
\begin{pgfscope}%
\pgfpathrectangle{\pgfqpoint{0.407592in}{0.222904in}}{\pgfqpoint{1.411163in}{1.411163in}}%
\pgfusepath{clip}%
\pgfsetbuttcap%
\pgfsetroundjoin%
\definecolor{currentfill}{rgb}{0.909804,0.698039,0.423529}%
\pgfsetfillcolor{currentfill}%
\pgfsetfillopacity{0.120000}%
\pgfsetlinewidth{0.000000pt}%
\definecolor{currentstroke}{rgb}{0.000000,0.000000,0.000000}%
\pgfsetstrokecolor{currentstroke}%
\pgfsetdash{}{0pt}%
\pgfpathmoveto{\pgfqpoint{0.796441in}{1.016864in}}%
\pgfpathlineto{\pgfqpoint{0.812556in}{1.001978in}}%
\pgfpathlineto{\pgfqpoint{0.828683in}{0.987546in}}%
\pgfpathlineto{\pgfqpoint{0.837221in}{0.989761in}}%
\pgfpathlineto{\pgfqpoint{0.845724in}{0.992408in}}%
\pgfpathlineto{\pgfqpoint{0.829638in}{1.006931in}}%
\pgfpathlineto{\pgfqpoint{0.813564in}{1.021904in}}%
\pgfpathlineto{\pgfqpoint{0.805021in}{1.019179in}}%
\pgfpathlineto{\pgfqpoint{0.796441in}{1.016864in}}%
\pgfpathclose%
\pgfusepath{fill}%
\end{pgfscope}%
\begin{pgfscope}%
\pgfpathrectangle{\pgfqpoint{0.407592in}{0.222904in}}{\pgfqpoint{1.411163in}{1.411163in}}%
\pgfusepath{clip}%
\pgfsetbuttcap%
\pgfsetroundjoin%
\definecolor{currentfill}{rgb}{0.909804,0.698039,0.423529}%
\pgfsetfillcolor{currentfill}%
\pgfsetfillopacity{0.120000}%
\pgfsetlinewidth{0.000000pt}%
\definecolor{currentstroke}{rgb}{0.000000,0.000000,0.000000}%
\pgfsetstrokecolor{currentstroke}%
\pgfsetdash{}{0pt}%
\pgfpathmoveto{\pgfqpoint{1.040377in}{0.897104in}}%
\pgfpathlineto{\pgfqpoint{1.056604in}{0.890710in}}%
\pgfpathlineto{\pgfqpoint{1.072863in}{0.885323in}}%
\pgfpathlineto{\pgfqpoint{1.080984in}{0.890440in}}%
\pgfpathlineto{\pgfqpoint{1.089079in}{0.896130in}}%
\pgfpathlineto{\pgfqpoint{1.072851in}{0.901485in}}%
\pgfpathlineto{\pgfqpoint{1.056656in}{0.907844in}}%
\pgfpathlineto{\pgfqpoint{1.048530in}{0.902192in}}%
\pgfpathlineto{\pgfqpoint{1.040377in}{0.897104in}}%
\pgfpathclose%
\pgfusepath{fill}%
\end{pgfscope}%
\begin{pgfscope}%
\pgfpathrectangle{\pgfqpoint{0.407592in}{0.222904in}}{\pgfqpoint{1.411163in}{1.411163in}}%
\pgfusepath{clip}%
\pgfsetbuttcap%
\pgfsetroundjoin%
\definecolor{currentfill}{rgb}{0.909804,0.698039,0.423529}%
\pgfsetfillcolor{currentfill}%
\pgfsetfillopacity{0.120000}%
\pgfsetlinewidth{0.000000pt}%
\definecolor{currentstroke}{rgb}{0.000000,0.000000,0.000000}%
\pgfsetstrokecolor{currentstroke}%
\pgfsetdash{}{0pt}%
\pgfpathmoveto{\pgfqpoint{0.568173in}{1.211487in}}%
\pgfpathlineto{\pgfqpoint{0.584252in}{1.193582in}}%
\pgfpathlineto{\pgfqpoint{0.600339in}{1.175842in}}%
\pgfpathlineto{\pgfqpoint{0.609282in}{1.177306in}}%
\pgfpathlineto{\pgfqpoint{0.618185in}{1.179069in}}%
\pgfpathlineto{\pgfqpoint{0.602140in}{1.196904in}}%
\pgfpathlineto{\pgfqpoint{0.586104in}{1.214899in}}%
\pgfpathlineto{\pgfqpoint{0.577158in}{1.213051in}}%
\pgfpathlineto{\pgfqpoint{0.568173in}{1.211487in}}%
\pgfpathclose%
\pgfusepath{fill}%
\end{pgfscope}%
\begin{pgfscope}%
\pgfpathrectangle{\pgfqpoint{0.407592in}{0.222904in}}{\pgfqpoint{1.411163in}{1.411163in}}%
\pgfusepath{clip}%
\pgfsetbuttcap%
\pgfsetroundjoin%
\definecolor{currentfill}{rgb}{0.909804,0.698039,0.423529}%
\pgfsetfillcolor{currentfill}%
\pgfsetfillopacity{0.120000}%
\pgfsetlinewidth{0.000000pt}%
\definecolor{currentstroke}{rgb}{0.000000,0.000000,0.000000}%
\pgfsetstrokecolor{currentstroke}%
\pgfsetdash{}{0pt}%
\pgfpathmoveto{\pgfqpoint{1.283614in}{0.952905in}}%
\pgfpathlineto{\pgfqpoint{1.300182in}{0.955702in}}%
\pgfpathlineto{\pgfqpoint{1.316812in}{0.959225in}}%
\pgfpathlineto{\pgfqpoint{1.324609in}{0.969024in}}%
\pgfpathlineto{\pgfqpoint{1.332390in}{0.979202in}}%
\pgfpathlineto{\pgfqpoint{1.315780in}{0.975872in}}%
\pgfpathlineto{\pgfqpoint{1.299230in}{0.973254in}}%
\pgfpathlineto{\pgfqpoint{1.291430in}{0.962886in}}%
\pgfpathlineto{\pgfqpoint{1.283614in}{0.952905in}}%
\pgfpathclose%
\pgfusepath{fill}%
\end{pgfscope}%
\begin{pgfscope}%
\pgfpathrectangle{\pgfqpoint{0.407592in}{0.222904in}}{\pgfqpoint{1.411163in}{1.411163in}}%
\pgfusepath{clip}%
\pgfsetbuttcap%
\pgfsetroundjoin%
\definecolor{currentfill}{rgb}{0.909804,0.698039,0.423529}%
\pgfsetfillcolor{currentfill}%
\pgfsetfillopacity{0.120000}%
\pgfsetlinewidth{0.000000pt}%
\definecolor{currentstroke}{rgb}{0.000000,0.000000,0.000000}%
\pgfsetstrokecolor{currentstroke}%
\pgfsetdash{}{0pt}%
\pgfpathmoveto{\pgfqpoint{0.910205in}{0.939716in}}%
\pgfpathlineto{\pgfqpoint{0.926364in}{0.928138in}}%
\pgfpathlineto{\pgfqpoint{0.942541in}{0.917303in}}%
\pgfpathlineto{\pgfqpoint{0.950890in}{0.920217in}}%
\pgfpathlineto{\pgfqpoint{0.959206in}{0.923650in}}%
\pgfpathlineto{\pgfqpoint{0.943067in}{0.934527in}}%
\pgfpathlineto{\pgfqpoint{0.926946in}{0.946149in}}%
\pgfpathlineto{\pgfqpoint{0.918592in}{0.942685in}}%
\pgfpathlineto{\pgfqpoint{0.910205in}{0.939716in}}%
\pgfpathclose%
\pgfusepath{fill}%
\end{pgfscope}%
\begin{pgfscope}%
\pgfpathrectangle{\pgfqpoint{0.407592in}{0.222904in}}{\pgfqpoint{1.411163in}{1.411163in}}%
\pgfusepath{clip}%
\pgfsetbuttcap%
\pgfsetroundjoin%
\definecolor{currentfill}{rgb}{0.909804,0.698039,0.423529}%
\pgfsetfillcolor{currentfill}%
\pgfsetfillopacity{0.120000}%
\pgfsetlinewidth{0.000000pt}%
\definecolor{currentstroke}{rgb}{0.000000,0.000000,0.000000}%
\pgfsetstrokecolor{currentstroke}%
\pgfsetdash{}{0pt}%
\pgfpathmoveto{\pgfqpoint{1.445638in}{1.063977in}}%
\pgfpathlineto{\pgfqpoint{1.462519in}{1.069526in}}%
\pgfpathlineto{\pgfqpoint{1.479467in}{1.075559in}}%
\pgfpathlineto{\pgfqpoint{1.487097in}{1.087113in}}%
\pgfpathlineto{\pgfqpoint{1.494713in}{1.098906in}}%
\pgfpathlineto{\pgfqpoint{1.477783in}{1.093145in}}%
\pgfpathlineto{\pgfqpoint{1.460920in}{1.087866in}}%
\pgfpathlineto{\pgfqpoint{1.453286in}{1.075801in}}%
\pgfpathlineto{\pgfqpoint{1.445638in}{1.063977in}}%
\pgfpathclose%
\pgfusepath{fill}%
\end{pgfscope}%
\begin{pgfscope}%
\pgfpathrectangle{\pgfqpoint{0.407592in}{0.222904in}}{\pgfqpoint{1.411163in}{1.411163in}}%
\pgfusepath{clip}%
\pgfsetbuttcap%
\pgfsetroundjoin%
\definecolor{currentfill}{rgb}{0.909804,0.698039,0.423529}%
\pgfsetfillcolor{currentfill}%
\pgfsetfillopacity{0.120000}%
\pgfsetlinewidth{0.000000pt}%
\definecolor{currentstroke}{rgb}{0.000000,0.000000,0.000000}%
\pgfsetstrokecolor{currentstroke}%
\pgfsetdash{}{0pt}%
\pgfpathmoveto{\pgfqpoint{0.714655in}{1.076306in}}%
\pgfpathlineto{\pgfqpoint{0.730768in}{1.060074in}}%
\pgfpathlineto{\pgfqpoint{0.746891in}{1.044165in}}%
\pgfpathlineto{\pgfqpoint{0.755587in}{1.045816in}}%
\pgfpathlineto{\pgfqpoint{0.764246in}{1.047843in}}%
\pgfpathlineto{\pgfqpoint{0.748165in}{1.063868in}}%
\pgfpathlineto{\pgfqpoint{0.732094in}{1.080209in}}%
\pgfpathlineto{\pgfqpoint{0.723393in}{1.078079in}}%
\pgfpathlineto{\pgfqpoint{0.714655in}{1.076306in}}%
\pgfpathclose%
\pgfusepath{fill}%
\end{pgfscope}%
\begin{pgfscope}%
\pgfpathrectangle{\pgfqpoint{0.407592in}{0.222904in}}{\pgfqpoint{1.411163in}{1.411163in}}%
\pgfusepath{clip}%
\pgfsetbuttcap%
\pgfsetroundjoin%
\definecolor{currentfill}{rgb}{0.909804,0.698039,0.423529}%
\pgfsetfillcolor{currentfill}%
\pgfsetfillopacity{0.120000}%
\pgfsetlinewidth{0.000000pt}%
\definecolor{currentstroke}{rgb}{0.000000,0.000000,0.000000}%
\pgfsetstrokecolor{currentstroke}%
\pgfsetdash{}{0pt}%
\pgfpathmoveto{\pgfqpoint{1.559050in}{1.159996in}}%
\pgfpathlineto{\pgfqpoint{1.576151in}{1.166498in}}%
\pgfpathlineto{\pgfqpoint{1.593323in}{1.173389in}}%
\pgfpathlineto{\pgfqpoint{1.600842in}{1.185658in}}%
\pgfpathlineto{\pgfqpoint{1.608348in}{1.198103in}}%
\pgfpathlineto{\pgfqpoint{1.591195in}{1.191486in}}%
\pgfpathlineto{\pgfqpoint{1.574112in}{1.185256in}}%
\pgfpathlineto{\pgfqpoint{1.566587in}{1.172539in}}%
\pgfpathlineto{\pgfqpoint{1.559050in}{1.159996in}}%
\pgfpathclose%
\pgfusepath{fill}%
\end{pgfscope}%
\begin{pgfscope}%
\pgfpathrectangle{\pgfqpoint{0.407592in}{0.222904in}}{\pgfqpoint{1.411163in}{1.411163in}}%
\pgfusepath{clip}%
\pgfsetbuttcap%
\pgfsetroundjoin%
\definecolor{currentfill}{rgb}{0.909804,0.698039,0.423529}%
\pgfsetfillcolor{currentfill}%
\pgfsetfillopacity{0.120000}%
\pgfsetlinewidth{0.000000pt}%
\definecolor{currentstroke}{rgb}{0.000000,0.000000,0.000000}%
\pgfsetstrokecolor{currentstroke}%
\pgfsetdash{}{0pt}%
\pgfpathmoveto{\pgfqpoint{0.828683in}{0.987546in}}%
\pgfpathlineto{\pgfqpoint{0.844822in}{0.973610in}}%
\pgfpathlineto{\pgfqpoint{0.860975in}{0.960214in}}%
\pgfpathlineto{\pgfqpoint{0.869473in}{0.962323in}}%
\pgfpathlineto{\pgfqpoint{0.877935in}{0.964887in}}%
\pgfpathlineto{\pgfqpoint{0.861823in}{0.978379in}}%
\pgfpathlineto{\pgfqpoint{0.845724in}{0.992408in}}%
\pgfpathlineto{\pgfqpoint{0.837221in}{0.989761in}}%
\pgfpathlineto{\pgfqpoint{0.828683in}{0.987546in}}%
\pgfpathclose%
\pgfusepath{fill}%
\end{pgfscope}%
\begin{pgfscope}%
\pgfpathrectangle{\pgfqpoint{0.407592in}{0.222904in}}{\pgfqpoint{1.411163in}{1.411163in}}%
\pgfusepath{clip}%
\pgfsetbuttcap%
\pgfsetroundjoin%
\definecolor{currentfill}{rgb}{0.909804,0.698039,0.423529}%
\pgfsetfillcolor{currentfill}%
\pgfsetfillopacity{0.120000}%
\pgfsetlinewidth{0.000000pt}%
\definecolor{currentstroke}{rgb}{0.000000,0.000000,0.000000}%
\pgfsetstrokecolor{currentstroke}%
\pgfsetdash{}{0pt}%
\pgfpathmoveto{\pgfqpoint{1.332390in}{0.979202in}}%
\pgfpathlineto{\pgfqpoint{1.349063in}{0.983206in}}%
\pgfpathlineto{\pgfqpoint{1.365800in}{0.987845in}}%
\pgfpathlineto{\pgfqpoint{1.373548in}{0.998162in}}%
\pgfpathlineto{\pgfqpoint{1.381282in}{1.008814in}}%
\pgfpathlineto{\pgfqpoint{1.364563in}{1.004406in}}%
\pgfpathlineto{\pgfqpoint{1.347909in}{1.000622in}}%
\pgfpathlineto{\pgfqpoint{1.340157in}{0.989741in}}%
\pgfpathlineto{\pgfqpoint{1.332390in}{0.979202in}}%
\pgfpathclose%
\pgfusepath{fill}%
\end{pgfscope}%
\begin{pgfscope}%
\pgfpathrectangle{\pgfqpoint{0.407592in}{0.222904in}}{\pgfqpoint{1.411163in}{1.411163in}}%
\pgfusepath{clip}%
\pgfsetbuttcap%
\pgfsetroundjoin%
\definecolor{currentfill}{rgb}{0.909804,0.698039,0.423529}%
\pgfsetfillcolor{currentfill}%
\pgfsetfillopacity{0.120000}%
\pgfsetlinewidth{0.000000pt}%
\definecolor{currentstroke}{rgb}{0.000000,0.000000,0.000000}%
\pgfsetstrokecolor{currentstroke}%
\pgfsetdash{}{0pt}%
\pgfpathmoveto{\pgfqpoint{0.991546in}{0.904343in}}%
\pgfpathlineto{\pgfqpoint{1.007751in}{0.896014in}}%
\pgfpathlineto{\pgfqpoint{1.023983in}{0.888627in}}%
\pgfpathlineto{\pgfqpoint{1.032195in}{0.892583in}}%
\pgfpathlineto{\pgfqpoint{1.040377in}{0.897104in}}%
\pgfpathlineto{\pgfqpoint{1.024180in}{0.904479in}}%
\pgfpathlineto{\pgfqpoint{1.008010in}{0.912799in}}%
\pgfpathlineto{\pgfqpoint{0.999793in}{0.908296in}}%
\pgfpathlineto{\pgfqpoint{0.991546in}{0.904343in}}%
\pgfpathclose%
\pgfusepath{fill}%
\end{pgfscope}%
\begin{pgfscope}%
\pgfpathrectangle{\pgfqpoint{0.407592in}{0.222904in}}{\pgfqpoint{1.411163in}{1.411163in}}%
\pgfusepath{clip}%
\pgfsetbuttcap%
\pgfsetroundjoin%
\definecolor{currentfill}{rgb}{0.909804,0.698039,0.423529}%
\pgfsetfillcolor{currentfill}%
\pgfsetfillopacity{0.120000}%
\pgfsetlinewidth{0.000000pt}%
\definecolor{currentstroke}{rgb}{0.000000,0.000000,0.000000}%
\pgfsetstrokecolor{currentstroke}%
\pgfsetdash{}{0pt}%
\pgfpathmoveto{\pgfqpoint{0.600339in}{1.175842in}}%
\pgfpathlineto{\pgfqpoint{0.616436in}{1.158283in}}%
\pgfpathlineto{\pgfqpoint{0.632542in}{1.140920in}}%
\pgfpathlineto{\pgfqpoint{0.641443in}{1.142270in}}%
\pgfpathlineto{\pgfqpoint{0.650304in}{1.143933in}}%
\pgfpathlineto{\pgfqpoint{0.634240in}{1.161406in}}%
\pgfpathlineto{\pgfqpoint{0.618185in}{1.179069in}}%
\pgfpathlineto{\pgfqpoint{0.609282in}{1.177306in}}%
\pgfpathlineto{\pgfqpoint{0.600339in}{1.175842in}}%
\pgfpathclose%
\pgfusepath{fill}%
\end{pgfscope}%
\begin{pgfscope}%
\pgfpathrectangle{\pgfqpoint{0.407592in}{0.222904in}}{\pgfqpoint{1.411163in}{1.411163in}}%
\pgfusepath{clip}%
\pgfsetbuttcap%
\pgfsetroundjoin%
\definecolor{currentfill}{rgb}{0.909804,0.698039,0.423529}%
\pgfsetfillcolor{currentfill}%
\pgfsetfillopacity{0.120000}%
\pgfsetlinewidth{0.000000pt}%
\definecolor{currentstroke}{rgb}{0.000000,0.000000,0.000000}%
\pgfsetstrokecolor{currentstroke}%
\pgfsetdash{}{0pt}%
\pgfpathmoveto{\pgfqpoint{0.746891in}{1.044165in}}%
\pgfpathlineto{\pgfqpoint{0.763025in}{1.028607in}}%
\pgfpathlineto{\pgfqpoint{0.779170in}{1.013433in}}%
\pgfpathlineto{\pgfqpoint{0.787824in}{1.014951in}}%
\pgfpathlineto{\pgfqpoint{0.796441in}{1.016864in}}%
\pgfpathlineto{\pgfqpoint{0.780338in}{1.032164in}}%
\pgfpathlineto{\pgfqpoint{0.764246in}{1.047843in}}%
\pgfpathlineto{\pgfqpoint{0.755587in}{1.045816in}}%
\pgfpathlineto{\pgfqpoint{0.746891in}{1.044165in}}%
\pgfpathclose%
\pgfusepath{fill}%
\end{pgfscope}%
\begin{pgfscope}%
\pgfpathrectangle{\pgfqpoint{0.407592in}{0.222904in}}{\pgfqpoint{1.411163in}{1.411163in}}%
\pgfusepath{clip}%
\pgfsetbuttcap%
\pgfsetroundjoin%
\definecolor{currentfill}{rgb}{0.909804,0.698039,0.423529}%
\pgfsetfillcolor{currentfill}%
\pgfsetfillopacity{0.120000}%
\pgfsetlinewidth{0.000000pt}%
\definecolor{currentstroke}{rgb}{0.000000,0.000000,0.000000}%
\pgfsetstrokecolor{currentstroke}%
\pgfsetdash{}{0pt}%
\pgfpathmoveto{\pgfqpoint{1.494713in}{1.098906in}}%
\pgfpathlineto{\pgfqpoint{1.511712in}{1.105122in}}%
\pgfpathlineto{\pgfqpoint{1.528781in}{1.111770in}}%
\pgfpathlineto{\pgfqpoint{1.536367in}{1.123513in}}%
\pgfpathlineto{\pgfqpoint{1.543940in}{1.135473in}}%
\pgfpathlineto{\pgfqpoint{1.526890in}{1.129102in}}%
\pgfpathlineto{\pgfqpoint{1.509909in}{1.123162in}}%
\pgfpathlineto{\pgfqpoint{1.502317in}{1.110926in}}%
\pgfpathlineto{\pgfqpoint{1.494713in}{1.098906in}}%
\pgfpathclose%
\pgfusepath{fill}%
\end{pgfscope}%
\begin{pgfscope}%
\pgfpathrectangle{\pgfqpoint{0.407592in}{0.222904in}}{\pgfqpoint{1.411163in}{1.411163in}}%
\pgfusepath{clip}%
\pgfsetbuttcap%
\pgfsetroundjoin%
\definecolor{currentfill}{rgb}{0.909804,0.698039,0.423529}%
\pgfsetfillcolor{currentfill}%
\pgfsetfillopacity{0.120000}%
\pgfsetlinewidth{0.000000pt}%
\definecolor{currentstroke}{rgb}{0.000000,0.000000,0.000000}%
\pgfsetstrokecolor{currentstroke}%
\pgfsetdash{}{0pt}%
\pgfpathmoveto{\pgfqpoint{1.381282in}{1.008814in}}%
\pgfpathlineto{\pgfqpoint{1.398065in}{1.013812in}}%
\pgfpathlineto{\pgfqpoint{1.414915in}{1.019364in}}%
\pgfpathlineto{\pgfqpoint{1.422616in}{1.030086in}}%
\pgfpathlineto{\pgfqpoint{1.430303in}{1.041105in}}%
\pgfpathlineto{\pgfqpoint{1.413472in}{1.035808in}}%
\pgfpathlineto{\pgfqpoint{1.396706in}{1.031059in}}%
\pgfpathlineto{\pgfqpoint{1.389001in}{1.019785in}}%
\pgfpathlineto{\pgfqpoint{1.381282in}{1.008814in}}%
\pgfpathclose%
\pgfusepath{fill}%
\end{pgfscope}%
\begin{pgfscope}%
\pgfpathrectangle{\pgfqpoint{0.407592in}{0.222904in}}{\pgfqpoint{1.411163in}{1.411163in}}%
\pgfusepath{clip}%
\pgfsetbuttcap%
\pgfsetroundjoin%
\definecolor{currentfill}{rgb}{0.909804,0.698039,0.423529}%
\pgfsetfillcolor{currentfill}%
\pgfsetfillopacity{0.120000}%
\pgfsetlinewidth{0.000000pt}%
\definecolor{currentstroke}{rgb}{0.000000,0.000000,0.000000}%
\pgfsetstrokecolor{currentstroke}%
\pgfsetdash{}{0pt}%
\pgfpathmoveto{\pgfqpoint{1.170387in}{0.898045in}}%
\pgfpathlineto{\pgfqpoint{1.186799in}{0.897728in}}%
\pgfpathlineto{\pgfqpoint{1.203263in}{0.898348in}}%
\pgfpathlineto{\pgfqpoint{1.211209in}{0.905679in}}%
\pgfpathlineto{\pgfqpoint{1.219135in}{0.913520in}}%
\pgfpathlineto{\pgfqpoint{1.202694in}{0.912958in}}%
\pgfpathlineto{\pgfqpoint{1.186306in}{0.913316in}}%
\pgfpathlineto{\pgfqpoint{1.178357in}{0.905419in}}%
\pgfpathlineto{\pgfqpoint{1.170387in}{0.898045in}}%
\pgfpathclose%
\pgfusepath{fill}%
\end{pgfscope}%
\begin{pgfscope}%
\pgfpathrectangle{\pgfqpoint{0.407592in}{0.222904in}}{\pgfqpoint{1.411163in}{1.411163in}}%
\pgfusepath{clip}%
\pgfsetbuttcap%
\pgfsetroundjoin%
\definecolor{currentfill}{rgb}{0.909804,0.698039,0.423529}%
\pgfsetfillcolor{currentfill}%
\pgfsetfillopacity{0.120000}%
\pgfsetlinewidth{0.000000pt}%
\definecolor{currentstroke}{rgb}{0.000000,0.000000,0.000000}%
\pgfsetstrokecolor{currentstroke}%
\pgfsetdash{}{0pt}%
\pgfpathmoveto{\pgfqpoint{1.121647in}{0.888489in}}%
\pgfpathlineto{\pgfqpoint{1.137991in}{0.886203in}}%
\pgfpathlineto{\pgfqpoint{1.154381in}{0.884923in}}%
\pgfpathlineto{\pgfqpoint{1.162396in}{0.891209in}}%
\pgfpathlineto{\pgfqpoint{1.170387in}{0.898045in}}%
\pgfpathlineto{\pgfqpoint{1.154023in}{0.899330in}}%
\pgfpathlineto{\pgfqpoint{1.137705in}{0.901610in}}%
\pgfpathlineto{\pgfqpoint{1.129688in}{0.894771in}}%
\pgfpathlineto{\pgfqpoint{1.121647in}{0.888489in}}%
\pgfpathclose%
\pgfusepath{fill}%
\end{pgfscope}%
\begin{pgfscope}%
\pgfpathrectangle{\pgfqpoint{0.407592in}{0.222904in}}{\pgfqpoint{1.411163in}{1.411163in}}%
\pgfusepath{clip}%
\pgfsetbuttcap%
\pgfsetroundjoin%
\definecolor{currentfill}{rgb}{0.909804,0.698039,0.423529}%
\pgfsetfillcolor{currentfill}%
\pgfsetfillopacity{0.120000}%
\pgfsetlinewidth{0.000000pt}%
\definecolor{currentstroke}{rgb}{0.000000,0.000000,0.000000}%
\pgfsetstrokecolor{currentstroke}%
\pgfsetdash{}{0pt}%
\pgfpathmoveto{\pgfqpoint{0.632542in}{1.140920in}}%
\pgfpathlineto{\pgfqpoint{0.648658in}{1.123770in}}%
\pgfpathlineto{\pgfqpoint{0.664783in}{1.106854in}}%
\pgfpathlineto{\pgfqpoint{0.673641in}{1.108076in}}%
\pgfpathlineto{\pgfqpoint{0.682460in}{1.109626in}}%
\pgfpathlineto{\pgfqpoint{0.666377in}{1.126667in}}%
\pgfpathlineto{\pgfqpoint{0.650304in}{1.143933in}}%
\pgfpathlineto{\pgfqpoint{0.641443in}{1.142270in}}%
\pgfpathlineto{\pgfqpoint{0.632542in}{1.140920in}}%
\pgfpathclose%
\pgfusepath{fill}%
\end{pgfscope}%
\begin{pgfscope}%
\pgfpathrectangle{\pgfqpoint{0.407592in}{0.222904in}}{\pgfqpoint{1.411163in}{1.411163in}}%
\pgfusepath{clip}%
\pgfsetbuttcap%
\pgfsetroundjoin%
\definecolor{currentfill}{rgb}{0.909804,0.698039,0.423529}%
\pgfsetfillcolor{currentfill}%
\pgfsetfillopacity{0.120000}%
\pgfsetlinewidth{0.000000pt}%
\definecolor{currentstroke}{rgb}{0.000000,0.000000,0.000000}%
\pgfsetstrokecolor{currentstroke}%
\pgfsetdash{}{0pt}%
\pgfpathmoveto{\pgfqpoint{1.219135in}{0.913520in}}%
\pgfpathlineto{\pgfqpoint{1.235629in}{0.914966in}}%
\pgfpathlineto{\pgfqpoint{1.252180in}{0.917253in}}%
\pgfpathlineto{\pgfqpoint{1.260065in}{0.925489in}}%
\pgfpathlineto{\pgfqpoint{1.267932in}{0.934188in}}%
\pgfpathlineto{\pgfqpoint{1.251402in}{0.932017in}}%
\pgfpathlineto{\pgfqpoint{1.234929in}{0.930672in}}%
\pgfpathlineto{\pgfqpoint{1.227041in}{0.921858in}}%
\pgfpathlineto{\pgfqpoint{1.219135in}{0.913520in}}%
\pgfpathclose%
\pgfusepath{fill}%
\end{pgfscope}%
\begin{pgfscope}%
\pgfpathrectangle{\pgfqpoint{0.407592in}{0.222904in}}{\pgfqpoint{1.411163in}{1.411163in}}%
\pgfusepath{clip}%
\pgfsetbuttcap%
\pgfsetroundjoin%
\definecolor{currentfill}{rgb}{0.909804,0.698039,0.423529}%
\pgfsetfillcolor{currentfill}%
\pgfsetfillopacity{0.120000}%
\pgfsetlinewidth{0.000000pt}%
\definecolor{currentstroke}{rgb}{0.000000,0.000000,0.000000}%
\pgfsetstrokecolor{currentstroke}%
\pgfsetdash{}{0pt}%
\pgfpathmoveto{\pgfqpoint{0.942541in}{0.917303in}}%
\pgfpathlineto{\pgfqpoint{0.958738in}{0.907261in}}%
\pgfpathlineto{\pgfqpoint{0.974957in}{0.898065in}}%
\pgfpathlineto{\pgfqpoint{0.983268in}{0.900935in}}%
\pgfpathlineto{\pgfqpoint{0.991546in}{0.904343in}}%
\pgfpathlineto{\pgfqpoint{0.975365in}{0.913571in}}%
\pgfpathlineto{\pgfqpoint{0.959206in}{0.923650in}}%
\pgfpathlineto{\pgfqpoint{0.950890in}{0.920217in}}%
\pgfpathlineto{\pgfqpoint{0.942541in}{0.917303in}}%
\pgfpathclose%
\pgfusepath{fill}%
\end{pgfscope}%
\begin{pgfscope}%
\pgfpathrectangle{\pgfqpoint{0.407592in}{0.222904in}}{\pgfqpoint{1.411163in}{1.411163in}}%
\pgfusepath{clip}%
\pgfsetbuttcap%
\pgfsetroundjoin%
\definecolor{currentfill}{rgb}{0.909804,0.698039,0.423529}%
\pgfsetfillcolor{currentfill}%
\pgfsetfillopacity{0.120000}%
\pgfsetlinewidth{0.000000pt}%
\definecolor{currentstroke}{rgb}{0.000000,0.000000,0.000000}%
\pgfsetstrokecolor{currentstroke}%
\pgfsetdash{}{0pt}%
\pgfpathmoveto{\pgfqpoint{0.860975in}{0.960214in}}%
\pgfpathlineto{\pgfqpoint{0.877143in}{0.947404in}}%
\pgfpathlineto{\pgfqpoint{0.893326in}{0.935230in}}%
\pgfpathlineto{\pgfqpoint{0.901783in}{0.937234in}}%
\pgfpathlineto{\pgfqpoint{0.910205in}{0.939716in}}%
\pgfpathlineto{\pgfqpoint{0.894062in}{0.951983in}}%
\pgfpathlineto{\pgfqpoint{0.877935in}{0.964887in}}%
\pgfpathlineto{\pgfqpoint{0.869473in}{0.962323in}}%
\pgfpathlineto{\pgfqpoint{0.860975in}{0.960214in}}%
\pgfpathclose%
\pgfusepath{fill}%
\end{pgfscope}%
\begin{pgfscope}%
\pgfpathrectangle{\pgfqpoint{0.407592in}{0.222904in}}{\pgfqpoint{1.411163in}{1.411163in}}%
\pgfusepath{clip}%
\pgfsetbuttcap%
\pgfsetroundjoin%
\definecolor{currentfill}{rgb}{0.909804,0.698039,0.423529}%
\pgfsetfillcolor{currentfill}%
\pgfsetfillopacity{0.120000}%
\pgfsetlinewidth{0.000000pt}%
\definecolor{currentstroke}{rgb}{0.000000,0.000000,0.000000}%
\pgfsetstrokecolor{currentstroke}%
\pgfsetdash{}{0pt}%
\pgfpathmoveto{\pgfqpoint{1.072863in}{0.885323in}}%
\pgfpathlineto{\pgfqpoint{1.089157in}{0.880962in}}%
\pgfpathlineto{\pgfqpoint{1.105490in}{0.877632in}}%
\pgfpathlineto{\pgfqpoint{1.113581in}{0.882774in}}%
\pgfpathlineto{\pgfqpoint{1.121647in}{0.888489in}}%
\pgfpathlineto{\pgfqpoint{1.105343in}{0.891796in}}%
\pgfpathlineto{\pgfqpoint{1.089079in}{0.896130in}}%
\pgfpathlineto{\pgfqpoint{1.080984in}{0.890440in}}%
\pgfpathlineto{\pgfqpoint{1.072863in}{0.885323in}}%
\pgfpathclose%
\pgfusepath{fill}%
\end{pgfscope}%
\begin{pgfscope}%
\pgfpathrectangle{\pgfqpoint{0.407592in}{0.222904in}}{\pgfqpoint{1.411163in}{1.411163in}}%
\pgfusepath{clip}%
\pgfsetbuttcap%
\pgfsetroundjoin%
\definecolor{currentfill}{rgb}{0.909804,0.698039,0.423529}%
\pgfsetfillcolor{currentfill}%
\pgfsetfillopacity{0.120000}%
\pgfsetlinewidth{0.000000pt}%
\definecolor{currentstroke}{rgb}{0.000000,0.000000,0.000000}%
\pgfsetstrokecolor{currentstroke}%
\pgfsetdash{}{0pt}%
\pgfpathmoveto{\pgfqpoint{1.267932in}{0.934188in}}%
\pgfpathlineto{\pgfqpoint{1.284520in}{0.937142in}}%
\pgfpathlineto{\pgfqpoint{1.301169in}{0.940834in}}%
\pgfpathlineto{\pgfqpoint{1.308999in}{0.949822in}}%
\pgfpathlineto{\pgfqpoint{1.316812in}{0.959225in}}%
\pgfpathlineto{\pgfqpoint{1.300182in}{0.955702in}}%
\pgfpathlineto{\pgfqpoint{1.283614in}{0.952905in}}%
\pgfpathlineto{\pgfqpoint{1.275781in}{0.943333in}}%
\pgfpathlineto{\pgfqpoint{1.267932in}{0.934188in}}%
\pgfpathclose%
\pgfusepath{fill}%
\end{pgfscope}%
\begin{pgfscope}%
\pgfpathrectangle{\pgfqpoint{0.407592in}{0.222904in}}{\pgfqpoint{1.411163in}{1.411163in}}%
\pgfusepath{clip}%
\pgfsetbuttcap%
\pgfsetroundjoin%
\definecolor{currentfill}{rgb}{0.909804,0.698039,0.423529}%
\pgfsetfillcolor{currentfill}%
\pgfsetfillopacity{0.120000}%
\pgfsetlinewidth{0.000000pt}%
\definecolor{currentstroke}{rgb}{0.000000,0.000000,0.000000}%
\pgfsetstrokecolor{currentstroke}%
\pgfsetdash{}{0pt}%
\pgfpathmoveto{\pgfqpoint{1.608348in}{1.198103in}}%
\pgfpathlineto{\pgfqpoint{1.625572in}{1.205090in}}%
\pgfpathlineto{\pgfqpoint{1.642868in}{1.212429in}}%
\pgfpathlineto{\pgfqpoint{1.650343in}{1.224772in}}%
\pgfpathlineto{\pgfqpoint{1.633057in}{1.217568in}}%
\pgfpathlineto{\pgfqpoint{1.615843in}{1.210716in}}%
\pgfusepath{fill}%
\end{pgfscope}%
\begin{pgfscope}%
\pgfpathrectangle{\pgfqpoint{0.407592in}{0.222904in}}{\pgfqpoint{1.411163in}{1.411163in}}%
\pgfusepath{clip}%
\pgfsetbuttcap%
\pgfsetroundjoin%
\definecolor{currentfill}{rgb}{0.909804,0.698039,0.423529}%
\pgfsetfillcolor{currentfill}%
\pgfsetfillopacity{0.120000}%
\pgfsetlinewidth{0.000000pt}%
\definecolor{currentstroke}{rgb}{0.000000,0.000000,0.000000}%
\pgfsetstrokecolor{currentstroke}%
\pgfsetdash{}{0pt}%
\pgfpathmoveto{\pgfqpoint{1.430303in}{1.041105in}}%
\pgfpathlineto{\pgfqpoint{1.447202in}{1.046920in}}%
\pgfpathlineto{\pgfqpoint{1.464170in}{1.053222in}}%
\pgfpathlineto{\pgfqpoint{1.471825in}{1.064258in}}%
\pgfpathlineto{\pgfqpoint{1.479467in}{1.075559in}}%
\pgfpathlineto{\pgfqpoint{1.462519in}{1.069526in}}%
\pgfpathlineto{\pgfqpoint{1.445638in}{1.063977in}}%
\pgfpathlineto{\pgfqpoint{1.437977in}{1.052407in}}%
\pgfpathlineto{\pgfqpoint{1.430303in}{1.041105in}}%
\pgfpathclose%
\pgfusepath{fill}%
\end{pgfscope}%
\begin{pgfscope}%
\pgfpathrectangle{\pgfqpoint{0.407592in}{0.222904in}}{\pgfqpoint{1.411163in}{1.411163in}}%
\pgfusepath{clip}%
\pgfsetbuttcap%
\pgfsetroundjoin%
\definecolor{currentfill}{rgb}{0.909804,0.698039,0.423529}%
\pgfsetfillcolor{currentfill}%
\pgfsetfillopacity{0.120000}%
\pgfsetlinewidth{0.000000pt}%
\definecolor{currentstroke}{rgb}{0.000000,0.000000,0.000000}%
\pgfsetstrokecolor{currentstroke}%
\pgfsetdash{}{0pt}%
\pgfpathmoveto{\pgfqpoint{0.779170in}{1.013433in}}%
\pgfpathlineto{\pgfqpoint{0.795327in}{0.998679in}}%
\pgfpathlineto{\pgfqpoint{0.811495in}{0.984382in}}%
\pgfpathlineto{\pgfqpoint{0.820108in}{0.985756in}}%
\pgfpathlineto{\pgfqpoint{0.828683in}{0.987546in}}%
\pgfpathlineto{\pgfqpoint{0.812556in}{1.001978in}}%
\pgfpathlineto{\pgfqpoint{0.796441in}{1.016864in}}%
\pgfpathlineto{\pgfqpoint{0.787824in}{1.014951in}}%
\pgfpathlineto{\pgfqpoint{0.779170in}{1.013433in}}%
\pgfpathclose%
\pgfusepath{fill}%
\end{pgfscope}%
\begin{pgfscope}%
\pgfpathrectangle{\pgfqpoint{0.407592in}{0.222904in}}{\pgfqpoint{1.411163in}{1.411163in}}%
\pgfusepath{clip}%
\pgfsetbuttcap%
\pgfsetroundjoin%
\definecolor{currentfill}{rgb}{0.909804,0.698039,0.423529}%
\pgfsetfillcolor{currentfill}%
\pgfsetfillopacity{0.120000}%
\pgfsetlinewidth{0.000000pt}%
\definecolor{currentstroke}{rgb}{0.000000,0.000000,0.000000}%
\pgfsetstrokecolor{currentstroke}%
\pgfsetdash{}{0pt}%
\pgfpathmoveto{\pgfqpoint{0.664783in}{1.106854in}}%
\pgfpathlineto{\pgfqpoint{0.680918in}{1.090190in}}%
\pgfpathlineto{\pgfqpoint{0.697063in}{1.073803in}}%
\pgfpathlineto{\pgfqpoint{0.705879in}{1.074882in}}%
\pgfpathlineto{\pgfqpoint{0.714655in}{1.076306in}}%
\pgfpathlineto{\pgfqpoint{0.698553in}{1.092831in}}%
\pgfpathlineto{\pgfqpoint{0.682460in}{1.109626in}}%
\pgfpathlineto{\pgfqpoint{0.673641in}{1.108076in}}%
\pgfpathlineto{\pgfqpoint{0.664783in}{1.106854in}}%
\pgfpathclose%
\pgfusepath{fill}%
\end{pgfscope}%
\begin{pgfscope}%
\pgfpathrectangle{\pgfqpoint{0.407592in}{0.222904in}}{\pgfqpoint{1.411163in}{1.411163in}}%
\pgfusepath{clip}%
\pgfsetbuttcap%
\pgfsetroundjoin%
\definecolor{currentfill}{rgb}{0.909804,0.698039,0.423529}%
\pgfsetfillcolor{currentfill}%
\pgfsetfillopacity{0.120000}%
\pgfsetlinewidth{0.000000pt}%
\definecolor{currentstroke}{rgb}{0.000000,0.000000,0.000000}%
\pgfsetstrokecolor{currentstroke}%
\pgfsetdash{}{0pt}%
\pgfpathmoveto{\pgfqpoint{1.543940in}{1.135473in}}%
\pgfpathlineto{\pgfqpoint{1.561060in}{1.142252in}}%
\pgfpathlineto{\pgfqpoint{1.578252in}{1.149419in}}%
\pgfpathlineto{\pgfqpoint{1.585794in}{1.161307in}}%
\pgfpathlineto{\pgfqpoint{1.593323in}{1.173389in}}%
\pgfpathlineto{\pgfqpoint{1.576151in}{1.166498in}}%
\pgfpathlineto{\pgfqpoint{1.559050in}{1.159996in}}%
\pgfpathlineto{\pgfqpoint{1.551501in}{1.147637in}}%
\pgfpathlineto{\pgfqpoint{1.543940in}{1.135473in}}%
\pgfpathclose%
\pgfusepath{fill}%
\end{pgfscope}%
\begin{pgfscope}%
\pgfpathrectangle{\pgfqpoint{0.407592in}{0.222904in}}{\pgfqpoint{1.411163in}{1.411163in}}%
\pgfusepath{clip}%
\pgfsetbuttcap%
\pgfsetroundjoin%
\definecolor{currentfill}{rgb}{0.909804,0.698039,0.423529}%
\pgfsetfillcolor{currentfill}%
\pgfsetfillopacity{0.120000}%
\pgfsetlinewidth{0.000000pt}%
\definecolor{currentstroke}{rgb}{0.000000,0.000000,0.000000}%
\pgfsetstrokecolor{currentstroke}%
\pgfsetdash{}{0pt}%
\pgfpathmoveto{\pgfqpoint{1.023983in}{0.888627in}}%
\pgfpathlineto{\pgfqpoint{1.040244in}{0.882217in}}%
\pgfpathlineto{\pgfqpoint{1.056536in}{0.876811in}}%
\pgfpathlineto{\pgfqpoint{1.064714in}{0.880781in}}%
\pgfpathlineto{\pgfqpoint{1.072863in}{0.885323in}}%
\pgfpathlineto{\pgfqpoint{1.056604in}{0.890710in}}%
\pgfpathlineto{\pgfqpoint{1.040377in}{0.897104in}}%
\pgfpathlineto{\pgfqpoint{1.032195in}{0.892583in}}%
\pgfpathlineto{\pgfqpoint{1.023983in}{0.888627in}}%
\pgfpathclose%
\pgfusepath{fill}%
\end{pgfscope}%
\begin{pgfscope}%
\pgfpathrectangle{\pgfqpoint{0.407592in}{0.222904in}}{\pgfqpoint{1.411163in}{1.411163in}}%
\pgfusepath{clip}%
\pgfsetbuttcap%
\pgfsetroundjoin%
\definecolor{currentfill}{rgb}{0.909804,0.698039,0.423529}%
\pgfsetfillcolor{currentfill}%
\pgfsetfillopacity{0.120000}%
\pgfsetlinewidth{0.000000pt}%
\definecolor{currentstroke}{rgb}{0.000000,0.000000,0.000000}%
\pgfsetstrokecolor{currentstroke}%
\pgfsetdash{}{0pt}%
\pgfpathmoveto{\pgfqpoint{1.316812in}{0.959225in}}%
\pgfpathlineto{\pgfqpoint{1.333504in}{0.963432in}}%
\pgfpathlineto{\pgfqpoint{1.350260in}{0.968282in}}%
\pgfpathlineto{\pgfqpoint{1.358038in}{0.977879in}}%
\pgfpathlineto{\pgfqpoint{1.365800in}{0.987845in}}%
\pgfpathlineto{\pgfqpoint{1.349063in}{0.983206in}}%
\pgfpathlineto{\pgfqpoint{1.332390in}{0.979202in}}%
\pgfpathlineto{\pgfqpoint{1.324609in}{0.969024in}}%
\pgfpathlineto{\pgfqpoint{1.316812in}{0.959225in}}%
\pgfpathclose%
\pgfusepath{fill}%
\end{pgfscope}%
\begin{pgfscope}%
\pgfpathrectangle{\pgfqpoint{0.407592in}{0.222904in}}{\pgfqpoint{1.411163in}{1.411163in}}%
\pgfusepath{clip}%
\pgfsetbuttcap%
\pgfsetroundjoin%
\definecolor{currentfill}{rgb}{0.909804,0.698039,0.423529}%
\pgfsetfillcolor{currentfill}%
\pgfsetfillopacity{0.120000}%
\pgfsetlinewidth{0.000000pt}%
\definecolor{currentstroke}{rgb}{0.000000,0.000000,0.000000}%
\pgfsetstrokecolor{currentstroke}%
\pgfsetdash{}{0pt}%
\pgfpathmoveto{\pgfqpoint{0.550083in}{1.209196in}}%
\pgfpathlineto{\pgfqpoint{0.566204in}{1.191407in}}%
\pgfpathlineto{\pgfqpoint{0.582334in}{1.173790in}}%
\pgfpathlineto{\pgfqpoint{0.591357in}{1.174672in}}%
\pgfpathlineto{\pgfqpoint{0.600339in}{1.175842in}}%
\pgfpathlineto{\pgfqpoint{0.584252in}{1.193582in}}%
\pgfpathlineto{\pgfqpoint{0.568173in}{1.211487in}}%
\pgfpathlineto{\pgfqpoint{0.559148in}{1.210203in}}%
\pgfpathlineto{\pgfqpoint{0.550083in}{1.209196in}}%
\pgfpathclose%
\pgfusepath{fill}%
\end{pgfscope}%
\begin{pgfscope}%
\pgfpathrectangle{\pgfqpoint{0.407592in}{0.222904in}}{\pgfqpoint{1.411163in}{1.411163in}}%
\pgfusepath{clip}%
\pgfsetbuttcap%
\pgfsetroundjoin%
\definecolor{currentfill}{rgb}{0.909804,0.698039,0.423529}%
\pgfsetfillcolor{currentfill}%
\pgfsetfillopacity{0.120000}%
\pgfsetlinewidth{0.000000pt}%
\definecolor{currentstroke}{rgb}{0.000000,0.000000,0.000000}%
\pgfsetstrokecolor{currentstroke}%
\pgfsetdash{}{0pt}%
\pgfpathmoveto{\pgfqpoint{0.893326in}{0.935230in}}%
\pgfpathlineto{\pgfqpoint{0.909525in}{0.923743in}}%
\pgfpathlineto{\pgfqpoint{0.925742in}{0.912991in}}%
\pgfpathlineto{\pgfqpoint{0.934159in}{0.914898in}}%
\pgfpathlineto{\pgfqpoint{0.942541in}{0.917303in}}%
\pgfpathlineto{\pgfqpoint{0.926364in}{0.928138in}}%
\pgfpathlineto{\pgfqpoint{0.910205in}{0.939716in}}%
\pgfpathlineto{\pgfqpoint{0.901783in}{0.937234in}}%
\pgfpathlineto{\pgfqpoint{0.893326in}{0.935230in}}%
\pgfpathclose%
\pgfusepath{fill}%
\end{pgfscope}%
\begin{pgfscope}%
\pgfpathrectangle{\pgfqpoint{0.407592in}{0.222904in}}{\pgfqpoint{1.411163in}{1.411163in}}%
\pgfusepath{clip}%
\pgfsetbuttcap%
\pgfsetroundjoin%
\definecolor{currentfill}{rgb}{0.909804,0.698039,0.423529}%
\pgfsetfillcolor{currentfill}%
\pgfsetfillopacity{0.120000}%
\pgfsetlinewidth{0.000000pt}%
\definecolor{currentstroke}{rgb}{0.000000,0.000000,0.000000}%
\pgfsetstrokecolor{currentstroke}%
\pgfsetdash{}{0pt}%
\pgfpathmoveto{\pgfqpoint{0.697063in}{1.073803in}}%
\pgfpathlineto{\pgfqpoint{0.713218in}{1.057717in}}%
\pgfpathlineto{\pgfqpoint{0.729383in}{1.041960in}}%
\pgfpathlineto{\pgfqpoint{0.738157in}{1.042882in}}%
\pgfpathlineto{\pgfqpoint{0.746891in}{1.044165in}}%
\pgfpathlineto{\pgfqpoint{0.730768in}{1.060074in}}%
\pgfpathlineto{\pgfqpoint{0.714655in}{1.076306in}}%
\pgfpathlineto{\pgfqpoint{0.705879in}{1.074882in}}%
\pgfpathlineto{\pgfqpoint{0.697063in}{1.073803in}}%
\pgfpathclose%
\pgfusepath{fill}%
\end{pgfscope}%
\begin{pgfscope}%
\pgfpathrectangle{\pgfqpoint{0.407592in}{0.222904in}}{\pgfqpoint{1.411163in}{1.411163in}}%
\pgfusepath{clip}%
\pgfsetbuttcap%
\pgfsetroundjoin%
\definecolor{currentfill}{rgb}{0.909804,0.698039,0.423529}%
\pgfsetfillcolor{currentfill}%
\pgfsetfillopacity{0.120000}%
\pgfsetlinewidth{0.000000pt}%
\definecolor{currentstroke}{rgb}{0.000000,0.000000,0.000000}%
\pgfsetstrokecolor{currentstroke}%
\pgfsetdash{}{0pt}%
\pgfpathmoveto{\pgfqpoint{1.479467in}{1.075559in}}%
\pgfpathlineto{\pgfqpoint{1.496485in}{1.082050in}}%
\pgfpathlineto{\pgfqpoint{1.513573in}{1.088973in}}%
\pgfpathlineto{\pgfqpoint{1.521184in}{1.100252in}}%
\pgfpathlineto{\pgfqpoint{1.528781in}{1.111770in}}%
\pgfpathlineto{\pgfqpoint{1.511712in}{1.105122in}}%
\pgfpathlineto{\pgfqpoint{1.494713in}{1.098906in}}%
\pgfpathlineto{\pgfqpoint{1.487097in}{1.087113in}}%
\pgfpathlineto{\pgfqpoint{1.479467in}{1.075559in}}%
\pgfpathclose%
\pgfusepath{fill}%
\end{pgfscope}%
\begin{pgfscope}%
\pgfpathrectangle{\pgfqpoint{0.407592in}{0.222904in}}{\pgfqpoint{1.411163in}{1.411163in}}%
\pgfusepath{clip}%
\pgfsetbuttcap%
\pgfsetroundjoin%
\definecolor{currentfill}{rgb}{0.909804,0.698039,0.423529}%
\pgfsetfillcolor{currentfill}%
\pgfsetfillopacity{0.120000}%
\pgfsetlinewidth{0.000000pt}%
\definecolor{currentstroke}{rgb}{0.000000,0.000000,0.000000}%
\pgfsetstrokecolor{currentstroke}%
\pgfsetdash{}{0pt}%
\pgfpathmoveto{\pgfqpoint{0.974957in}{0.898065in}}%
\pgfpathlineto{\pgfqpoint{0.991200in}{0.889760in}}%
\pgfpathlineto{\pgfqpoint{1.007468in}{0.882387in}}%
\pgfpathlineto{\pgfqpoint{1.015741in}{0.885231in}}%
\pgfpathlineto{\pgfqpoint{1.023983in}{0.888627in}}%
\pgfpathlineto{\pgfqpoint{1.007751in}{0.896014in}}%
\pgfpathlineto{\pgfqpoint{0.991546in}{0.904343in}}%
\pgfpathlineto{\pgfqpoint{0.983268in}{0.900935in}}%
\pgfpathlineto{\pgfqpoint{0.974957in}{0.898065in}}%
\pgfpathclose%
\pgfusepath{fill}%
\end{pgfscope}%
\begin{pgfscope}%
\pgfpathrectangle{\pgfqpoint{0.407592in}{0.222904in}}{\pgfqpoint{1.411163in}{1.411163in}}%
\pgfusepath{clip}%
\pgfsetbuttcap%
\pgfsetroundjoin%
\definecolor{currentfill}{rgb}{0.909804,0.698039,0.423529}%
\pgfsetfillcolor{currentfill}%
\pgfsetfillopacity{0.120000}%
\pgfsetlinewidth{0.000000pt}%
\definecolor{currentstroke}{rgb}{0.000000,0.000000,0.000000}%
\pgfsetstrokecolor{currentstroke}%
\pgfsetdash{}{0pt}%
\pgfpathmoveto{\pgfqpoint{1.365800in}{0.987845in}}%
\pgfpathlineto{\pgfqpoint{1.382603in}{0.993079in}}%
\pgfpathlineto{\pgfqpoint{1.399472in}{0.998872in}}%
\pgfpathlineto{\pgfqpoint{1.407201in}{1.008955in}}%
\pgfpathlineto{\pgfqpoint{1.414915in}{1.019364in}}%
\pgfpathlineto{\pgfqpoint{1.398065in}{1.013812in}}%
\pgfpathlineto{\pgfqpoint{1.381282in}{1.008814in}}%
\pgfpathlineto{\pgfqpoint{1.373548in}{0.998162in}}%
\pgfpathlineto{\pgfqpoint{1.365800in}{0.987845in}}%
\pgfpathclose%
\pgfusepath{fill}%
\end{pgfscope}%
\begin{pgfscope}%
\pgfpathrectangle{\pgfqpoint{0.407592in}{0.222904in}}{\pgfqpoint{1.411163in}{1.411163in}}%
\pgfusepath{clip}%
\pgfsetbuttcap%
\pgfsetroundjoin%
\definecolor{currentfill}{rgb}{0.909804,0.698039,0.423529}%
\pgfsetfillcolor{currentfill}%
\pgfsetfillopacity{0.120000}%
\pgfsetlinewidth{0.000000pt}%
\definecolor{currentstroke}{rgb}{0.000000,0.000000,0.000000}%
\pgfsetstrokecolor{currentstroke}%
\pgfsetdash{}{0pt}%
\pgfpathmoveto{\pgfqpoint{0.811495in}{0.984382in}}%
\pgfpathlineto{\pgfqpoint{0.827676in}{0.970582in}}%
\pgfpathlineto{\pgfqpoint{0.843871in}{0.957324in}}%
\pgfpathlineto{\pgfqpoint{0.852442in}{0.958551in}}%
\pgfpathlineto{\pgfqpoint{0.860975in}{0.960214in}}%
\pgfpathlineto{\pgfqpoint{0.844822in}{0.973610in}}%
\pgfpathlineto{\pgfqpoint{0.828683in}{0.987546in}}%
\pgfpathlineto{\pgfqpoint{0.820108in}{0.985756in}}%
\pgfpathlineto{\pgfqpoint{0.811495in}{0.984382in}}%
\pgfpathclose%
\pgfusepath{fill}%
\end{pgfscope}%
\begin{pgfscope}%
\pgfpathrectangle{\pgfqpoint{0.407592in}{0.222904in}}{\pgfqpoint{1.411163in}{1.411163in}}%
\pgfusepath{clip}%
\pgfsetbuttcap%
\pgfsetroundjoin%
\definecolor{currentfill}{rgb}{0.909804,0.698039,0.423529}%
\pgfsetfillcolor{currentfill}%
\pgfsetfillopacity{0.120000}%
\pgfsetlinewidth{0.000000pt}%
\definecolor{currentstroke}{rgb}{0.000000,0.000000,0.000000}%
\pgfsetstrokecolor{currentstroke}%
\pgfsetdash{}{0pt}%
\pgfpathmoveto{\pgfqpoint{0.582334in}{1.173790in}}%
\pgfpathlineto{\pgfqpoint{0.598473in}{1.156362in}}%
\pgfpathlineto{\pgfqpoint{0.614622in}{1.139137in}}%
\pgfpathlineto{\pgfqpoint{0.623602in}{1.139877in}}%
\pgfpathlineto{\pgfqpoint{0.632542in}{1.140920in}}%
\pgfpathlineto{\pgfqpoint{0.616436in}{1.158283in}}%
\pgfpathlineto{\pgfqpoint{0.600339in}{1.175842in}}%
\pgfpathlineto{\pgfqpoint{0.591357in}{1.174672in}}%
\pgfpathlineto{\pgfqpoint{0.582334in}{1.173790in}}%
\pgfpathclose%
\pgfusepath{fill}%
\end{pgfscope}%
\begin{pgfscope}%
\pgfpathrectangle{\pgfqpoint{0.407592in}{0.222904in}}{\pgfqpoint{1.411163in}{1.411163in}}%
\pgfusepath{clip}%
\pgfsetbuttcap%
\pgfsetroundjoin%
\definecolor{currentfill}{rgb}{0.909804,0.698039,0.423529}%
\pgfsetfillcolor{currentfill}%
\pgfsetfillopacity{0.120000}%
\pgfsetlinewidth{0.000000pt}%
\definecolor{currentstroke}{rgb}{0.000000,0.000000,0.000000}%
\pgfsetstrokecolor{currentstroke}%
\pgfsetdash{}{0pt}%
\pgfpathmoveto{\pgfqpoint{1.593323in}{1.173389in}}%
\pgfpathlineto{\pgfqpoint{1.610567in}{1.180649in}}%
\pgfpathlineto{\pgfqpoint{1.627883in}{1.188260in}}%
\pgfpathlineto{\pgfqpoint{1.635381in}{1.200256in}}%
\pgfpathlineto{\pgfqpoint{1.642868in}{1.212429in}}%
\pgfpathlineto{\pgfqpoint{1.625572in}{1.205090in}}%
\pgfpathlineto{\pgfqpoint{1.608348in}{1.198103in}}%
\pgfpathlineto{\pgfqpoint{1.600842in}{1.185658in}}%
\pgfpathlineto{\pgfqpoint{1.593323in}{1.173389in}}%
\pgfpathclose%
\pgfusepath{fill}%
\end{pgfscope}%
\begin{pgfscope}%
\pgfpathrectangle{\pgfqpoint{0.407592in}{0.222904in}}{\pgfqpoint{1.411163in}{1.411163in}}%
\pgfusepath{clip}%
\pgfsetbuttcap%
\pgfsetroundjoin%
\definecolor{currentfill}{rgb}{0.909804,0.698039,0.423529}%
\pgfsetfillcolor{currentfill}%
\pgfsetfillopacity{0.120000}%
\pgfsetlinewidth{0.000000pt}%
\definecolor{currentstroke}{rgb}{0.000000,0.000000,0.000000}%
\pgfsetstrokecolor{currentstroke}%
\pgfsetdash{}{0pt}%
\pgfpathmoveto{\pgfqpoint{1.154381in}{0.884923in}}%
\pgfpathlineto{\pgfqpoint{1.170819in}{0.884624in}}%
\pgfpathlineto{\pgfqpoint{1.187307in}{0.885272in}}%
\pgfpathlineto{\pgfqpoint{1.195296in}{0.891542in}}%
\pgfpathlineto{\pgfqpoint{1.203263in}{0.898348in}}%
\pgfpathlineto{\pgfqpoint{1.186799in}{0.897728in}}%
\pgfpathlineto{\pgfqpoint{1.170387in}{0.898045in}}%
\pgfpathlineto{\pgfqpoint{1.162396in}{0.891209in}}%
\pgfpathlineto{\pgfqpoint{1.154381in}{0.884923in}}%
\pgfpathclose%
\pgfusepath{fill}%
\end{pgfscope}%
\begin{pgfscope}%
\pgfpathrectangle{\pgfqpoint{0.407592in}{0.222904in}}{\pgfqpoint{1.411163in}{1.411163in}}%
\pgfusepath{clip}%
\pgfsetbuttcap%
\pgfsetroundjoin%
\definecolor{currentfill}{rgb}{0.909804,0.698039,0.423529}%
\pgfsetfillcolor{currentfill}%
\pgfsetfillopacity{0.120000}%
\pgfsetlinewidth{0.000000pt}%
\definecolor{currentstroke}{rgb}{0.000000,0.000000,0.000000}%
\pgfsetstrokecolor{currentstroke}%
\pgfsetdash{}{0pt}%
\pgfpathmoveto{\pgfqpoint{1.203263in}{0.898348in}}%
\pgfpathlineto{\pgfqpoint{1.219780in}{0.899865in}}%
\pgfpathlineto{\pgfqpoint{1.236353in}{0.902235in}}%
\pgfpathlineto{\pgfqpoint{1.244276in}{0.909497in}}%
\pgfpathlineto{\pgfqpoint{1.252180in}{0.917253in}}%
\pgfpathlineto{\pgfqpoint{1.235629in}{0.914966in}}%
\pgfpathlineto{\pgfqpoint{1.219135in}{0.913520in}}%
\pgfpathlineto{\pgfqpoint{1.211209in}{0.905679in}}%
\pgfpathlineto{\pgfqpoint{1.203263in}{0.898348in}}%
\pgfpathclose%
\pgfusepath{fill}%
\end{pgfscope}%
\begin{pgfscope}%
\pgfpathrectangle{\pgfqpoint{0.407592in}{0.222904in}}{\pgfqpoint{1.411163in}{1.411163in}}%
\pgfusepath{clip}%
\pgfsetbuttcap%
\pgfsetroundjoin%
\definecolor{currentfill}{rgb}{0.909804,0.698039,0.423529}%
\pgfsetfillcolor{currentfill}%
\pgfsetfillopacity{0.120000}%
\pgfsetlinewidth{0.000000pt}%
\definecolor{currentstroke}{rgb}{0.000000,0.000000,0.000000}%
\pgfsetstrokecolor{currentstroke}%
\pgfsetdash{}{0pt}%
\pgfpathmoveto{\pgfqpoint{1.105490in}{0.877632in}}%
\pgfpathlineto{\pgfqpoint{1.121863in}{0.875328in}}%
\pgfpathlineto{\pgfqpoint{1.138281in}{0.874035in}}%
\pgfpathlineto{\pgfqpoint{1.146344in}{0.879197in}}%
\pgfpathlineto{\pgfqpoint{1.154381in}{0.884923in}}%
\pgfpathlineto{\pgfqpoint{1.137991in}{0.886203in}}%
\pgfpathlineto{\pgfqpoint{1.121647in}{0.888489in}}%
\pgfpathlineto{\pgfqpoint{1.113581in}{0.882774in}}%
\pgfpathlineto{\pgfqpoint{1.105490in}{0.877632in}}%
\pgfpathclose%
\pgfusepath{fill}%
\end{pgfscope}%
\begin{pgfscope}%
\pgfpathrectangle{\pgfqpoint{0.407592in}{0.222904in}}{\pgfqpoint{1.411163in}{1.411163in}}%
\pgfusepath{clip}%
\pgfsetbuttcap%
\pgfsetroundjoin%
\definecolor{currentfill}{rgb}{0.909804,0.698039,0.423529}%
\pgfsetfillcolor{currentfill}%
\pgfsetfillopacity{0.120000}%
\pgfsetlinewidth{0.000000pt}%
\definecolor{currentstroke}{rgb}{0.000000,0.000000,0.000000}%
\pgfsetstrokecolor{currentstroke}%
\pgfsetdash{}{0pt}%
\pgfpathmoveto{\pgfqpoint{0.729383in}{1.041960in}}%
\pgfpathlineto{\pgfqpoint{0.745560in}{1.026560in}}%
\pgfpathlineto{\pgfqpoint{0.761747in}{1.011550in}}%
\pgfpathlineto{\pgfqpoint{0.770478in}{1.012303in}}%
\pgfpathlineto{\pgfqpoint{0.779170in}{1.013433in}}%
\pgfpathlineto{\pgfqpoint{0.763025in}{1.028607in}}%
\pgfpathlineto{\pgfqpoint{0.746891in}{1.044165in}}%
\pgfpathlineto{\pgfqpoint{0.738157in}{1.042882in}}%
\pgfpathlineto{\pgfqpoint{0.729383in}{1.041960in}}%
\pgfpathclose%
\pgfusepath{fill}%
\end{pgfscope}%
\begin{pgfscope}%
\pgfpathrectangle{\pgfqpoint{0.407592in}{0.222904in}}{\pgfqpoint{1.411163in}{1.411163in}}%
\pgfusepath{clip}%
\pgfsetbuttcap%
\pgfsetroundjoin%
\definecolor{currentfill}{rgb}{0.909804,0.698039,0.423529}%
\pgfsetfillcolor{currentfill}%
\pgfsetfillopacity{0.120000}%
\pgfsetlinewidth{0.000000pt}%
\definecolor{currentstroke}{rgb}{0.000000,0.000000,0.000000}%
\pgfsetstrokecolor{currentstroke}%
\pgfsetdash{}{0pt}%
\pgfpathmoveto{\pgfqpoint{1.252180in}{0.917253in}}%
\pgfpathlineto{\pgfqpoint{1.268789in}{0.920336in}}%
\pgfpathlineto{\pgfqpoint{1.285459in}{0.924167in}}%
\pgfpathlineto{\pgfqpoint{1.293323in}{0.932277in}}%
\pgfpathlineto{\pgfqpoint{1.301169in}{0.940834in}}%
\pgfpathlineto{\pgfqpoint{1.284520in}{0.937142in}}%
\pgfpathlineto{\pgfqpoint{1.267932in}{0.934188in}}%
\pgfpathlineto{\pgfqpoint{1.260065in}{0.925489in}}%
\pgfpathlineto{\pgfqpoint{1.252180in}{0.917253in}}%
\pgfpathclose%
\pgfusepath{fill}%
\end{pgfscope}%
\begin{pgfscope}%
\pgfpathrectangle{\pgfqpoint{0.407592in}{0.222904in}}{\pgfqpoint{1.411163in}{1.411163in}}%
\pgfusepath{clip}%
\pgfsetbuttcap%
\pgfsetroundjoin%
\definecolor{currentfill}{rgb}{0.909804,0.698039,0.423529}%
\pgfsetfillcolor{currentfill}%
\pgfsetfillopacity{0.120000}%
\pgfsetlinewidth{0.000000pt}%
\definecolor{currentstroke}{rgb}{0.000000,0.000000,0.000000}%
\pgfsetstrokecolor{currentstroke}%
\pgfsetdash{}{0pt}%
\pgfpathmoveto{\pgfqpoint{1.414915in}{1.019364in}}%
\pgfpathlineto{\pgfqpoint{1.431833in}{1.025437in}}%
\pgfpathlineto{\pgfqpoint{1.448820in}{1.031998in}}%
\pgfpathlineto{\pgfqpoint{1.456502in}{1.042465in}}%
\pgfpathlineto{\pgfqpoint{1.464170in}{1.053222in}}%
\pgfpathlineto{\pgfqpoint{1.447202in}{1.046920in}}%
\pgfpathlineto{\pgfqpoint{1.430303in}{1.041105in}}%
\pgfpathlineto{\pgfqpoint{1.422616in}{1.030086in}}%
\pgfpathlineto{\pgfqpoint{1.414915in}{1.019364in}}%
\pgfpathclose%
\pgfusepath{fill}%
\end{pgfscope}%
\begin{pgfscope}%
\pgfpathrectangle{\pgfqpoint{0.407592in}{0.222904in}}{\pgfqpoint{1.411163in}{1.411163in}}%
\pgfusepath{clip}%
\pgfsetbuttcap%
\pgfsetroundjoin%
\definecolor{currentfill}{rgb}{0.909804,0.698039,0.423529}%
\pgfsetfillcolor{currentfill}%
\pgfsetfillopacity{0.120000}%
\pgfsetlinewidth{0.000000pt}%
\definecolor{currentstroke}{rgb}{0.000000,0.000000,0.000000}%
\pgfsetstrokecolor{currentstroke}%
\pgfsetdash{}{0pt}%
\pgfpathmoveto{\pgfqpoint{1.528781in}{1.111770in}}%
\pgfpathlineto{\pgfqpoint{1.545921in}{1.118826in}}%
\pgfpathlineto{\pgfqpoint{1.563132in}{1.126268in}}%
\pgfpathlineto{\pgfqpoint{1.570698in}{1.137736in}}%
\pgfpathlineto{\pgfqpoint{1.578252in}{1.149419in}}%
\pgfpathlineto{\pgfqpoint{1.561060in}{1.142252in}}%
\pgfpathlineto{\pgfqpoint{1.543940in}{1.135473in}}%
\pgfpathlineto{\pgfqpoint{1.536367in}{1.123513in}}%
\pgfpathlineto{\pgfqpoint{1.528781in}{1.111770in}}%
\pgfpathclose%
\pgfusepath{fill}%
\end{pgfscope}%
\begin{pgfscope}%
\pgfpathrectangle{\pgfqpoint{0.407592in}{0.222904in}}{\pgfqpoint{1.411163in}{1.411163in}}%
\pgfusepath{clip}%
\pgfsetbuttcap%
\pgfsetroundjoin%
\definecolor{currentfill}{rgb}{0.909804,0.698039,0.423529}%
\pgfsetfillcolor{currentfill}%
\pgfsetfillopacity{0.120000}%
\pgfsetlinewidth{0.000000pt}%
\definecolor{currentstroke}{rgb}{0.000000,0.000000,0.000000}%
\pgfsetstrokecolor{currentstroke}%
\pgfsetdash{}{0pt}%
\pgfpathmoveto{\pgfqpoint{0.614622in}{1.139137in}}%
\pgfpathlineto{\pgfqpoint{0.630780in}{1.122133in}}%
\pgfpathlineto{\pgfqpoint{0.646948in}{1.105370in}}%
\pgfpathlineto{\pgfqpoint{0.655886in}{1.105954in}}%
\pgfpathlineto{\pgfqpoint{0.664783in}{1.106854in}}%
\pgfpathlineto{\pgfqpoint{0.648658in}{1.123770in}}%
\pgfpathlineto{\pgfqpoint{0.632542in}{1.140920in}}%
\pgfpathlineto{\pgfqpoint{0.623602in}{1.139877in}}%
\pgfpathlineto{\pgfqpoint{0.614622in}{1.139137in}}%
\pgfpathclose%
\pgfusepath{fill}%
\end{pgfscope}%
\begin{pgfscope}%
\pgfpathrectangle{\pgfqpoint{0.407592in}{0.222904in}}{\pgfqpoint{1.411163in}{1.411163in}}%
\pgfusepath{clip}%
\pgfsetbuttcap%
\pgfsetroundjoin%
\definecolor{currentfill}{rgb}{0.909804,0.698039,0.423529}%
\pgfsetfillcolor{currentfill}%
\pgfsetfillopacity{0.120000}%
\pgfsetlinewidth{0.000000pt}%
\definecolor{currentstroke}{rgb}{0.000000,0.000000,0.000000}%
\pgfsetstrokecolor{currentstroke}%
\pgfsetdash{}{0pt}%
\pgfpathmoveto{\pgfqpoint{1.056536in}{0.876811in}}%
\pgfpathlineto{\pgfqpoint{1.072863in}{0.872427in}}%
\pgfpathlineto{\pgfqpoint{1.089227in}{0.869071in}}%
\pgfpathlineto{\pgfqpoint{1.097372in}{0.873065in}}%
\pgfpathlineto{\pgfqpoint{1.105490in}{0.877632in}}%
\pgfpathlineto{\pgfqpoint{1.089157in}{0.880962in}}%
\pgfpathlineto{\pgfqpoint{1.072863in}{0.885323in}}%
\pgfpathlineto{\pgfqpoint{1.064714in}{0.880781in}}%
\pgfpathlineto{\pgfqpoint{1.056536in}{0.876811in}}%
\pgfpathclose%
\pgfusepath{fill}%
\end{pgfscope}%
\begin{pgfscope}%
\pgfpathrectangle{\pgfqpoint{0.407592in}{0.222904in}}{\pgfqpoint{1.411163in}{1.411163in}}%
\pgfusepath{clip}%
\pgfsetbuttcap%
\pgfsetroundjoin%
\definecolor{currentfill}{rgb}{0.909804,0.698039,0.423529}%
\pgfsetfillcolor{currentfill}%
\pgfsetfillopacity{0.120000}%
\pgfsetlinewidth{0.000000pt}%
\definecolor{currentstroke}{rgb}{0.000000,0.000000,0.000000}%
\pgfsetstrokecolor{currentstroke}%
\pgfsetdash{}{0pt}%
\pgfpathmoveto{\pgfqpoint{0.925742in}{0.912991in}}%
\pgfpathlineto{\pgfqpoint{0.941979in}{0.903026in}}%
\pgfpathlineto{\pgfqpoint{0.958237in}{0.893896in}}%
\pgfpathlineto{\pgfqpoint{0.966614in}{0.895723in}}%
\pgfpathlineto{\pgfqpoint{0.974957in}{0.898065in}}%
\pgfpathlineto{\pgfqpoint{0.958738in}{0.907261in}}%
\pgfpathlineto{\pgfqpoint{0.942541in}{0.917303in}}%
\pgfpathlineto{\pgfqpoint{0.934159in}{0.914898in}}%
\pgfpathlineto{\pgfqpoint{0.925742in}{0.912991in}}%
\pgfpathclose%
\pgfusepath{fill}%
\end{pgfscope}%
\begin{pgfscope}%
\pgfpathrectangle{\pgfqpoint{0.407592in}{0.222904in}}{\pgfqpoint{1.411163in}{1.411163in}}%
\pgfusepath{clip}%
\pgfsetbuttcap%
\pgfsetroundjoin%
\definecolor{currentfill}{rgb}{0.909804,0.698039,0.423529}%
\pgfsetfillcolor{currentfill}%
\pgfsetfillopacity{0.120000}%
\pgfsetlinewidth{0.000000pt}%
\definecolor{currentstroke}{rgb}{0.000000,0.000000,0.000000}%
\pgfsetstrokecolor{currentstroke}%
\pgfsetdash{}{0pt}%
\pgfpathmoveto{\pgfqpoint{0.843871in}{0.957324in}}%
\pgfpathlineto{\pgfqpoint{0.860079in}{0.944651in}}%
\pgfpathlineto{\pgfqpoint{0.876303in}{0.932611in}}%
\pgfpathlineto{\pgfqpoint{0.884833in}{0.933694in}}%
\pgfpathlineto{\pgfqpoint{0.893326in}{0.935230in}}%
\pgfpathlineto{\pgfqpoint{0.877143in}{0.947404in}}%
\pgfpathlineto{\pgfqpoint{0.860975in}{0.960214in}}%
\pgfpathlineto{\pgfqpoint{0.852442in}{0.958551in}}%
\pgfpathlineto{\pgfqpoint{0.843871in}{0.957324in}}%
\pgfpathclose%
\pgfusepath{fill}%
\end{pgfscope}%
\begin{pgfscope}%
\pgfpathrectangle{\pgfqpoint{0.407592in}{0.222904in}}{\pgfqpoint{1.411163in}{1.411163in}}%
\pgfusepath{clip}%
\pgfsetbuttcap%
\pgfsetroundjoin%
\definecolor{currentfill}{rgb}{0.909804,0.698039,0.423529}%
\pgfsetfillcolor{currentfill}%
\pgfsetfillopacity{0.120000}%
\pgfsetlinewidth{0.000000pt}%
\definecolor{currentstroke}{rgb}{0.000000,0.000000,0.000000}%
\pgfsetstrokecolor{currentstroke}%
\pgfsetdash{}{0pt}%
\pgfpathmoveto{\pgfqpoint{1.301169in}{0.940834in}}%
\pgfpathlineto{\pgfqpoint{1.317881in}{0.945221in}}%
\pgfpathlineto{\pgfqpoint{1.334657in}{0.950256in}}%
\pgfpathlineto{\pgfqpoint{1.342466in}{0.959069in}}%
\pgfpathlineto{\pgfqpoint{1.350260in}{0.968282in}}%
\pgfpathlineto{\pgfqpoint{1.333504in}{0.963432in}}%
\pgfpathlineto{\pgfqpoint{1.316812in}{0.959225in}}%
\pgfpathlineto{\pgfqpoint{1.308999in}{0.949822in}}%
\pgfpathlineto{\pgfqpoint{1.301169in}{0.940834in}}%
\pgfpathclose%
\pgfusepath{fill}%
\end{pgfscope}%
\begin{pgfscope}%
\pgfpathrectangle{\pgfqpoint{0.407592in}{0.222904in}}{\pgfqpoint{1.411163in}{1.411163in}}%
\pgfusepath{clip}%
\pgfsetbuttcap%
\pgfsetroundjoin%
\definecolor{currentfill}{rgb}{0.909804,0.698039,0.423529}%
\pgfsetfillcolor{currentfill}%
\pgfsetfillopacity{0.120000}%
\pgfsetlinewidth{0.000000pt}%
\definecolor{currentstroke}{rgb}{0.000000,0.000000,0.000000}%
\pgfsetstrokecolor{currentstroke}%
\pgfsetdash{}{0pt}%
\pgfpathmoveto{\pgfqpoint{1.007468in}{0.882387in}}%
\pgfpathlineto{\pgfqpoint{1.023765in}{0.875982in}}%
\pgfpathlineto{\pgfqpoint{1.040092in}{0.870570in}}%
\pgfpathlineto{\pgfqpoint{1.048329in}{0.873411in}}%
\pgfpathlineto{\pgfqpoint{1.056536in}{0.876811in}}%
\pgfpathlineto{\pgfqpoint{1.040244in}{0.882217in}}%
\pgfpathlineto{\pgfqpoint{1.023983in}{0.888627in}}%
\pgfpathlineto{\pgfqpoint{1.015741in}{0.885231in}}%
\pgfpathlineto{\pgfqpoint{1.007468in}{0.882387in}}%
\pgfpathclose%
\pgfusepath{fill}%
\end{pgfscope}%
\begin{pgfscope}%
\pgfpathrectangle{\pgfqpoint{0.407592in}{0.222904in}}{\pgfqpoint{1.411163in}{1.411163in}}%
\pgfusepath{clip}%
\pgfsetbuttcap%
\pgfsetroundjoin%
\definecolor{currentfill}{rgb}{0.909804,0.698039,0.423529}%
\pgfsetfillcolor{currentfill}%
\pgfsetfillopacity{0.120000}%
\pgfsetlinewidth{0.000000pt}%
\definecolor{currentstroke}{rgb}{0.000000,0.000000,0.000000}%
\pgfsetstrokecolor{currentstroke}%
\pgfsetdash{}{0pt}%
\pgfpathmoveto{\pgfqpoint{1.642868in}{1.212429in}}%
\pgfpathlineto{\pgfqpoint{1.660237in}{1.220105in}}%
\pgfpathlineto{\pgfqpoint{1.677679in}{1.228102in}}%
\pgfpathlineto{\pgfqpoint{1.685134in}{1.240178in}}%
\pgfpathlineto{\pgfqpoint{1.667702in}{1.232314in}}%
\pgfpathlineto{\pgfqpoint{1.650343in}{1.224772in}}%
\pgfusepath{fill}%
\end{pgfscope}%
\begin{pgfscope}%
\pgfpathrectangle{\pgfqpoint{0.407592in}{0.222904in}}{\pgfqpoint{1.411163in}{1.411163in}}%
\pgfusepath{clip}%
\pgfsetbuttcap%
\pgfsetroundjoin%
\definecolor{currentfill}{rgb}{0.909804,0.698039,0.423529}%
\pgfsetfillcolor{currentfill}%
\pgfsetfillopacity{0.120000}%
\pgfsetlinewidth{0.000000pt}%
\definecolor{currentstroke}{rgb}{0.000000,0.000000,0.000000}%
\pgfsetstrokecolor{currentstroke}%
\pgfsetdash{}{0pt}%
\pgfpathmoveto{\pgfqpoint{0.761747in}{1.011550in}}%
\pgfpathlineto{\pgfqpoint{0.777945in}{0.996964in}}%
\pgfpathlineto{\pgfqpoint{0.794156in}{0.982838in}}%
\pgfpathlineto{\pgfqpoint{0.802845in}{0.983413in}}%
\pgfpathlineto{\pgfqpoint{0.811495in}{0.984382in}}%
\pgfpathlineto{\pgfqpoint{0.795327in}{0.998679in}}%
\pgfpathlineto{\pgfqpoint{0.779170in}{1.013433in}}%
\pgfpathlineto{\pgfqpoint{0.770478in}{1.012303in}}%
\pgfpathlineto{\pgfqpoint{0.761747in}{1.011550in}}%
\pgfpathclose%
\pgfusepath{fill}%
\end{pgfscope}%
\begin{pgfscope}%
\pgfpathrectangle{\pgfqpoint{0.407592in}{0.222904in}}{\pgfqpoint{1.411163in}{1.411163in}}%
\pgfusepath{clip}%
\pgfsetbuttcap%
\pgfsetroundjoin%
\definecolor{currentfill}{rgb}{0.909804,0.698039,0.423529}%
\pgfsetfillcolor{currentfill}%
\pgfsetfillopacity{0.120000}%
\pgfsetlinewidth{0.000000pt}%
\definecolor{currentstroke}{rgb}{0.000000,0.000000,0.000000}%
\pgfsetstrokecolor{currentstroke}%
\pgfsetdash{}{0pt}%
\pgfpathmoveto{\pgfqpoint{0.646948in}{1.105370in}}%
\pgfpathlineto{\pgfqpoint{0.663125in}{1.088867in}}%
\pgfpathlineto{\pgfqpoint{0.679313in}{1.072649in}}%
\pgfpathlineto{\pgfqpoint{0.688208in}{1.073061in}}%
\pgfpathlineto{\pgfqpoint{0.697063in}{1.073803in}}%
\pgfpathlineto{\pgfqpoint{0.680918in}{1.090190in}}%
\pgfpathlineto{\pgfqpoint{0.664783in}{1.106854in}}%
\pgfpathlineto{\pgfqpoint{0.655886in}{1.105954in}}%
\pgfpathlineto{\pgfqpoint{0.646948in}{1.105370in}}%
\pgfpathclose%
\pgfusepath{fill}%
\end{pgfscope}%
\begin{pgfscope}%
\pgfpathrectangle{\pgfqpoint{0.407592in}{0.222904in}}{\pgfqpoint{1.411163in}{1.411163in}}%
\pgfusepath{clip}%
\pgfsetbuttcap%
\pgfsetroundjoin%
\definecolor{currentfill}{rgb}{0.909804,0.698039,0.423529}%
\pgfsetfillcolor{currentfill}%
\pgfsetfillopacity{0.120000}%
\pgfsetlinewidth{0.000000pt}%
\definecolor{currentstroke}{rgb}{0.000000,0.000000,0.000000}%
\pgfsetstrokecolor{currentstroke}%
\pgfsetdash{}{0pt}%
\pgfpathmoveto{\pgfqpoint{1.464170in}{1.053222in}}%
\pgfpathlineto{\pgfqpoint{1.481207in}{1.059983in}}%
\pgfpathlineto{\pgfqpoint{1.498315in}{1.067173in}}%
\pgfpathlineto{\pgfqpoint{1.505950in}{1.077943in}}%
\pgfpathlineto{\pgfqpoint{1.513573in}{1.088973in}}%
\pgfpathlineto{\pgfqpoint{1.496485in}{1.082050in}}%
\pgfpathlineto{\pgfqpoint{1.479467in}{1.075559in}}%
\pgfpathlineto{\pgfqpoint{1.471825in}{1.064258in}}%
\pgfpathlineto{\pgfqpoint{1.464170in}{1.053222in}}%
\pgfpathclose%
\pgfusepath{fill}%
\end{pgfscope}%
\begin{pgfscope}%
\pgfpathrectangle{\pgfqpoint{0.407592in}{0.222904in}}{\pgfqpoint{1.411163in}{1.411163in}}%
\pgfusepath{clip}%
\pgfsetbuttcap%
\pgfsetroundjoin%
\definecolor{currentfill}{rgb}{0.909804,0.698039,0.423529}%
\pgfsetfillcolor{currentfill}%
\pgfsetfillopacity{0.120000}%
\pgfsetlinewidth{0.000000pt}%
\definecolor{currentstroke}{rgb}{0.000000,0.000000,0.000000}%
\pgfsetstrokecolor{currentstroke}%
\pgfsetdash{}{0pt}%
\pgfpathmoveto{\pgfqpoint{1.578252in}{1.149419in}}%
\pgfpathlineto{\pgfqpoint{1.595515in}{1.156953in}}%
\pgfpathlineto{\pgfqpoint{1.612852in}{1.164837in}}%
\pgfpathlineto{\pgfqpoint{1.620373in}{1.176451in}}%
\pgfpathlineto{\pgfqpoint{1.627883in}{1.188260in}}%
\pgfpathlineto{\pgfqpoint{1.610567in}{1.180649in}}%
\pgfpathlineto{\pgfqpoint{1.593323in}{1.173389in}}%
\pgfpathlineto{\pgfqpoint{1.585794in}{1.161307in}}%
\pgfpathlineto{\pgfqpoint{1.578252in}{1.149419in}}%
\pgfpathclose%
\pgfusepath{fill}%
\end{pgfscope}%
\begin{pgfscope}%
\pgfpathrectangle{\pgfqpoint{0.407592in}{0.222904in}}{\pgfqpoint{1.411163in}{1.411163in}}%
\pgfusepath{clip}%
\pgfsetbuttcap%
\pgfsetroundjoin%
\definecolor{currentfill}{rgb}{0.909804,0.698039,0.423529}%
\pgfsetfillcolor{currentfill}%
\pgfsetfillopacity{0.120000}%
\pgfsetlinewidth{0.000000pt}%
\definecolor{currentstroke}{rgb}{0.000000,0.000000,0.000000}%
\pgfsetstrokecolor{currentstroke}%
\pgfsetdash{}{0pt}%
\pgfpathmoveto{\pgfqpoint{1.350260in}{0.968282in}}%
\pgfpathlineto{\pgfqpoint{1.367081in}{0.973734in}}%
\pgfpathlineto{\pgfqpoint{1.383970in}{0.979747in}}%
\pgfpathlineto{\pgfqpoint{1.391728in}{0.989132in}}%
\pgfpathlineto{\pgfqpoint{1.399472in}{0.998872in}}%
\pgfpathlineto{\pgfqpoint{1.382603in}{0.993079in}}%
\pgfpathlineto{\pgfqpoint{1.365800in}{0.987845in}}%
\pgfpathlineto{\pgfqpoint{1.358038in}{0.977879in}}%
\pgfpathlineto{\pgfqpoint{1.350260in}{0.968282in}}%
\pgfpathclose%
\pgfusepath{fill}%
\end{pgfscope}%
\begin{pgfscope}%
\pgfpathrectangle{\pgfqpoint{0.407592in}{0.222904in}}{\pgfqpoint{1.411163in}{1.411163in}}%
\pgfusepath{clip}%
\pgfsetbuttcap%
\pgfsetroundjoin%
\definecolor{currentfill}{rgb}{0.909804,0.698039,0.423529}%
\pgfsetfillcolor{currentfill}%
\pgfsetfillopacity{0.120000}%
\pgfsetlinewidth{0.000000pt}%
\definecolor{currentstroke}{rgb}{0.000000,0.000000,0.000000}%
\pgfsetstrokecolor{currentstroke}%
\pgfsetdash{}{0pt}%
\pgfpathmoveto{\pgfqpoint{0.531830in}{1.207998in}}%
\pgfpathlineto{\pgfqpoint{0.547993in}{1.190347in}}%
\pgfpathlineto{\pgfqpoint{0.564166in}{1.172877in}}%
\pgfpathlineto{\pgfqpoint{0.573270in}{1.173194in}}%
\pgfpathlineto{\pgfqpoint{0.582334in}{1.173790in}}%
\pgfpathlineto{\pgfqpoint{0.566204in}{1.191407in}}%
\pgfpathlineto{\pgfqpoint{0.550083in}{1.209196in}}%
\pgfpathlineto{\pgfqpoint{0.540977in}{1.208462in}}%
\pgfpathlineto{\pgfqpoint{0.531830in}{1.207998in}}%
\pgfpathclose%
\pgfusepath{fill}%
\end{pgfscope}%
\begin{pgfscope}%
\pgfpathrectangle{\pgfqpoint{0.407592in}{0.222904in}}{\pgfqpoint{1.411163in}{1.411163in}}%
\pgfusepath{clip}%
\pgfsetbuttcap%
\pgfsetroundjoin%
\definecolor{currentfill}{rgb}{0.909804,0.698039,0.423529}%
\pgfsetfillcolor{currentfill}%
\pgfsetfillopacity{0.120000}%
\pgfsetlinewidth{0.000000pt}%
\definecolor{currentstroke}{rgb}{0.000000,0.000000,0.000000}%
\pgfsetstrokecolor{currentstroke}%
\pgfsetdash{}{0pt}%
\pgfpathmoveto{\pgfqpoint{0.876303in}{0.932611in}}%
\pgfpathlineto{\pgfqpoint{0.892544in}{0.921252in}}%
\pgfpathlineto{\pgfqpoint{0.908802in}{0.910621in}}%
\pgfpathlineto{\pgfqpoint{0.917290in}{0.911570in}}%
\pgfpathlineto{\pgfqpoint{0.925742in}{0.912991in}}%
\pgfpathlineto{\pgfqpoint{0.909525in}{0.923743in}}%
\pgfpathlineto{\pgfqpoint{0.893326in}{0.935230in}}%
\pgfpathlineto{\pgfqpoint{0.884833in}{0.933694in}}%
\pgfpathlineto{\pgfqpoint{0.876303in}{0.932611in}}%
\pgfpathclose%
\pgfusepath{fill}%
\end{pgfscope}%
\begin{pgfscope}%
\pgfpathrectangle{\pgfqpoint{0.407592in}{0.222904in}}{\pgfqpoint{1.411163in}{1.411163in}}%
\pgfusepath{clip}%
\pgfsetbuttcap%
\pgfsetroundjoin%
\definecolor{currentfill}{rgb}{0.909804,0.698039,0.423529}%
\pgfsetfillcolor{currentfill}%
\pgfsetfillopacity{0.120000}%
\pgfsetlinewidth{0.000000pt}%
\definecolor{currentstroke}{rgb}{0.000000,0.000000,0.000000}%
\pgfsetstrokecolor{currentstroke}%
\pgfsetdash{}{0pt}%
\pgfpathmoveto{\pgfqpoint{0.958237in}{0.893896in}}%
\pgfpathlineto{\pgfqpoint{0.974518in}{0.885645in}}%
\pgfpathlineto{\pgfqpoint{0.990824in}{0.878311in}}%
\pgfpathlineto{\pgfqpoint{0.999163in}{0.880085in}}%
\pgfpathlineto{\pgfqpoint{1.007468in}{0.882387in}}%
\pgfpathlineto{\pgfqpoint{0.991200in}{0.889760in}}%
\pgfpathlineto{\pgfqpoint{0.974957in}{0.898065in}}%
\pgfpathlineto{\pgfqpoint{0.966614in}{0.895723in}}%
\pgfpathlineto{\pgfqpoint{0.958237in}{0.893896in}}%
\pgfpathclose%
\pgfusepath{fill}%
\end{pgfscope}%
\begin{pgfscope}%
\pgfpathrectangle{\pgfqpoint{0.407592in}{0.222904in}}{\pgfqpoint{1.411163in}{1.411163in}}%
\pgfusepath{clip}%
\pgfsetbuttcap%
\pgfsetroundjoin%
\definecolor{currentfill}{rgb}{0.909804,0.698039,0.423529}%
\pgfsetfillcolor{currentfill}%
\pgfsetfillopacity{0.120000}%
\pgfsetlinewidth{0.000000pt}%
\definecolor{currentstroke}{rgb}{0.000000,0.000000,0.000000}%
\pgfsetstrokecolor{currentstroke}%
\pgfsetdash{}{0pt}%
\pgfpathmoveto{\pgfqpoint{0.679313in}{1.072649in}}%
\pgfpathlineto{\pgfqpoint{0.695510in}{1.056739in}}%
\pgfpathlineto{\pgfqpoint{0.711718in}{1.041165in}}%
\pgfpathlineto{\pgfqpoint{0.720571in}{1.041390in}}%
\pgfpathlineto{\pgfqpoint{0.729383in}{1.041960in}}%
\pgfpathlineto{\pgfqpoint{0.713218in}{1.057717in}}%
\pgfpathlineto{\pgfqpoint{0.697063in}{1.073803in}}%
\pgfpathlineto{\pgfqpoint{0.688208in}{1.073061in}}%
\pgfpathlineto{\pgfqpoint{0.679313in}{1.072649in}}%
\pgfpathclose%
\pgfusepath{fill}%
\end{pgfscope}%
\begin{pgfscope}%
\pgfpathrectangle{\pgfqpoint{0.407592in}{0.222904in}}{\pgfqpoint{1.411163in}{1.411163in}}%
\pgfusepath{clip}%
\pgfsetbuttcap%
\pgfsetroundjoin%
\definecolor{currentfill}{rgb}{0.909804,0.698039,0.423529}%
\pgfsetfillcolor{currentfill}%
\pgfsetfillopacity{0.120000}%
\pgfsetlinewidth{0.000000pt}%
\definecolor{currentstroke}{rgb}{0.000000,0.000000,0.000000}%
\pgfsetstrokecolor{currentstroke}%
\pgfsetdash{}{0pt}%
\pgfpathmoveto{\pgfqpoint{1.138281in}{0.874035in}}%
\pgfpathlineto{\pgfqpoint{1.154747in}{0.873727in}}%
\pgfpathlineto{\pgfqpoint{1.171262in}{0.874370in}}%
\pgfpathlineto{\pgfqpoint{1.179296in}{0.879546in}}%
\pgfpathlineto{\pgfqpoint{1.187307in}{0.885272in}}%
\pgfpathlineto{\pgfqpoint{1.170819in}{0.884624in}}%
\pgfpathlineto{\pgfqpoint{1.154381in}{0.884923in}}%
\pgfpathlineto{\pgfqpoint{1.146344in}{0.879197in}}%
\pgfpathlineto{\pgfqpoint{1.138281in}{0.874035in}}%
\pgfpathclose%
\pgfusepath{fill}%
\end{pgfscope}%
\begin{pgfscope}%
\pgfpathrectangle{\pgfqpoint{0.407592in}{0.222904in}}{\pgfqpoint{1.411163in}{1.411163in}}%
\pgfusepath{clip}%
\pgfsetbuttcap%
\pgfsetroundjoin%
\definecolor{currentfill}{rgb}{0.909804,0.698039,0.423529}%
\pgfsetfillcolor{currentfill}%
\pgfsetfillopacity{0.120000}%
\pgfsetlinewidth{0.000000pt}%
\definecolor{currentstroke}{rgb}{0.000000,0.000000,0.000000}%
\pgfsetstrokecolor{currentstroke}%
\pgfsetdash{}{0pt}%
\pgfpathmoveto{\pgfqpoint{1.513573in}{1.088973in}}%
\pgfpathlineto{\pgfqpoint{1.530733in}{1.096303in}}%
\pgfpathlineto{\pgfqpoint{1.547964in}{1.104016in}}%
\pgfpathlineto{\pgfqpoint{1.555554in}{1.115024in}}%
\pgfpathlineto{\pgfqpoint{1.563132in}{1.126268in}}%
\pgfpathlineto{\pgfqpoint{1.545921in}{1.118826in}}%
\pgfpathlineto{\pgfqpoint{1.528781in}{1.111770in}}%
\pgfpathlineto{\pgfqpoint{1.521184in}{1.100252in}}%
\pgfpathlineto{\pgfqpoint{1.513573in}{1.088973in}}%
\pgfpathclose%
\pgfusepath{fill}%
\end{pgfscope}%
\begin{pgfscope}%
\pgfpathrectangle{\pgfqpoint{0.407592in}{0.222904in}}{\pgfqpoint{1.411163in}{1.411163in}}%
\pgfusepath{clip}%
\pgfsetbuttcap%
\pgfsetroundjoin%
\definecolor{currentfill}{rgb}{0.909804,0.698039,0.423529}%
\pgfsetfillcolor{currentfill}%
\pgfsetfillopacity{0.120000}%
\pgfsetlinewidth{0.000000pt}%
\definecolor{currentstroke}{rgb}{0.000000,0.000000,0.000000}%
\pgfsetstrokecolor{currentstroke}%
\pgfsetdash{}{0pt}%
\pgfpathmoveto{\pgfqpoint{1.187307in}{0.885272in}}%
\pgfpathlineto{\pgfqpoint{1.203849in}{0.886826in}}%
\pgfpathlineto{\pgfqpoint{1.220446in}{0.889240in}}%
\pgfpathlineto{\pgfqpoint{1.228410in}{0.895480in}}%
\pgfpathlineto{\pgfqpoint{1.236353in}{0.902235in}}%
\pgfpathlineto{\pgfqpoint{1.219780in}{0.899865in}}%
\pgfpathlineto{\pgfqpoint{1.203263in}{0.898348in}}%
\pgfpathlineto{\pgfqpoint{1.195296in}{0.891542in}}%
\pgfpathlineto{\pgfqpoint{1.187307in}{0.885272in}}%
\pgfpathclose%
\pgfusepath{fill}%
\end{pgfscope}%
\begin{pgfscope}%
\pgfpathrectangle{\pgfqpoint{0.407592in}{0.222904in}}{\pgfqpoint{1.411163in}{1.411163in}}%
\pgfusepath{clip}%
\pgfsetbuttcap%
\pgfsetroundjoin%
\definecolor{currentfill}{rgb}{0.909804,0.698039,0.423529}%
\pgfsetfillcolor{currentfill}%
\pgfsetfillopacity{0.120000}%
\pgfsetlinewidth{0.000000pt}%
\definecolor{currentstroke}{rgb}{0.000000,0.000000,0.000000}%
\pgfsetstrokecolor{currentstroke}%
\pgfsetdash{}{0pt}%
\pgfpathmoveto{\pgfqpoint{0.794156in}{0.982838in}}%
\pgfpathlineto{\pgfqpoint{0.810379in}{0.969212in}}%
\pgfpathlineto{\pgfqpoint{0.826616in}{0.956127in}}%
\pgfpathlineto{\pgfqpoint{0.835262in}{0.956520in}}%
\pgfpathlineto{\pgfqpoint{0.843871in}{0.957324in}}%
\pgfpathlineto{\pgfqpoint{0.827676in}{0.970582in}}%
\pgfpathlineto{\pgfqpoint{0.811495in}{0.984382in}}%
\pgfpathlineto{\pgfqpoint{0.802845in}{0.983413in}}%
\pgfpathlineto{\pgfqpoint{0.794156in}{0.982838in}}%
\pgfpathclose%
\pgfusepath{fill}%
\end{pgfscope}%
\begin{pgfscope}%
\pgfpathrectangle{\pgfqpoint{0.407592in}{0.222904in}}{\pgfqpoint{1.411163in}{1.411163in}}%
\pgfusepath{clip}%
\pgfsetbuttcap%
\pgfsetroundjoin%
\definecolor{currentfill}{rgb}{0.909804,0.698039,0.423529}%
\pgfsetfillcolor{currentfill}%
\pgfsetfillopacity{0.120000}%
\pgfsetlinewidth{0.000000pt}%
\definecolor{currentstroke}{rgb}{0.000000,0.000000,0.000000}%
\pgfsetstrokecolor{currentstroke}%
\pgfsetdash{}{0pt}%
\pgfpathmoveto{\pgfqpoint{1.399472in}{0.998872in}}%
\pgfpathlineto{\pgfqpoint{1.416409in}{1.005188in}}%
\pgfpathlineto{\pgfqpoint{1.433415in}{1.011992in}}%
\pgfpathlineto{\pgfqpoint{1.441124in}{1.021836in}}%
\pgfpathlineto{\pgfqpoint{1.448820in}{1.031998in}}%
\pgfpathlineto{\pgfqpoint{1.431833in}{1.025437in}}%
\pgfpathlineto{\pgfqpoint{1.414915in}{1.019364in}}%
\pgfpathlineto{\pgfqpoint{1.407201in}{1.008955in}}%
\pgfpathlineto{\pgfqpoint{1.399472in}{0.998872in}}%
\pgfpathclose%
\pgfusepath{fill}%
\end{pgfscope}%
\begin{pgfscope}%
\pgfpathrectangle{\pgfqpoint{0.407592in}{0.222904in}}{\pgfqpoint{1.411163in}{1.411163in}}%
\pgfusepath{clip}%
\pgfsetbuttcap%
\pgfsetroundjoin%
\definecolor{currentfill}{rgb}{0.909804,0.698039,0.423529}%
\pgfsetfillcolor{currentfill}%
\pgfsetfillopacity{0.120000}%
\pgfsetlinewidth{0.000000pt}%
\definecolor{currentstroke}{rgb}{0.000000,0.000000,0.000000}%
\pgfsetstrokecolor{currentstroke}%
\pgfsetdash{}{0pt}%
\pgfpathmoveto{\pgfqpoint{0.564166in}{1.172877in}}%
\pgfpathlineto{\pgfqpoint{0.580348in}{1.155603in}}%
\pgfpathlineto{\pgfqpoint{0.596539in}{1.138540in}}%
\pgfpathlineto{\pgfqpoint{0.605601in}{1.138693in}}%
\pgfpathlineto{\pgfqpoint{0.614622in}{1.139137in}}%
\pgfpathlineto{\pgfqpoint{0.598473in}{1.156362in}}%
\pgfpathlineto{\pgfqpoint{0.582334in}{1.173790in}}%
\pgfpathlineto{\pgfqpoint{0.573270in}{1.173194in}}%
\pgfpathlineto{\pgfqpoint{0.564166in}{1.172877in}}%
\pgfpathclose%
\pgfusepath{fill}%
\end{pgfscope}%
\begin{pgfscope}%
\pgfpathrectangle{\pgfqpoint{0.407592in}{0.222904in}}{\pgfqpoint{1.411163in}{1.411163in}}%
\pgfusepath{clip}%
\pgfsetbuttcap%
\pgfsetroundjoin%
\definecolor{currentfill}{rgb}{0.909804,0.698039,0.423529}%
\pgfsetfillcolor{currentfill}%
\pgfsetfillopacity{0.120000}%
\pgfsetlinewidth{0.000000pt}%
\definecolor{currentstroke}{rgb}{0.000000,0.000000,0.000000}%
\pgfsetstrokecolor{currentstroke}%
\pgfsetdash{}{0pt}%
\pgfpathmoveto{\pgfqpoint{1.089227in}{0.869071in}}%
\pgfpathlineto{\pgfqpoint{1.105632in}{0.866737in}}%
\pgfpathlineto{\pgfqpoint{1.122080in}{0.865411in}}%
\pgfpathlineto{\pgfqpoint{1.130194in}{0.869440in}}%
\pgfpathlineto{\pgfqpoint{1.138281in}{0.874035in}}%
\pgfpathlineto{\pgfqpoint{1.121863in}{0.875328in}}%
\pgfpathlineto{\pgfqpoint{1.105490in}{0.877632in}}%
\pgfpathlineto{\pgfqpoint{1.097372in}{0.873065in}}%
\pgfpathlineto{\pgfqpoint{1.089227in}{0.869071in}}%
\pgfpathclose%
\pgfusepath{fill}%
\end{pgfscope}%
\begin{pgfscope}%
\pgfpathrectangle{\pgfqpoint{0.407592in}{0.222904in}}{\pgfqpoint{1.411163in}{1.411163in}}%
\pgfusepath{clip}%
\pgfsetbuttcap%
\pgfsetroundjoin%
\definecolor{currentfill}{rgb}{0.909804,0.698039,0.423529}%
\pgfsetfillcolor{currentfill}%
\pgfsetfillopacity{0.120000}%
\pgfsetlinewidth{0.000000pt}%
\definecolor{currentstroke}{rgb}{0.000000,0.000000,0.000000}%
\pgfsetstrokecolor{currentstroke}%
\pgfsetdash{}{0pt}%
\pgfpathmoveto{\pgfqpoint{1.236353in}{0.902235in}}%
\pgfpathlineto{\pgfqpoint{1.252984in}{0.905410in}}%
\pgfpathlineto{\pgfqpoint{1.269677in}{0.909342in}}%
\pgfpathlineto{\pgfqpoint{1.277577in}{0.916518in}}%
\pgfpathlineto{\pgfqpoint{1.285459in}{0.924167in}}%
\pgfpathlineto{\pgfqpoint{1.268789in}{0.920336in}}%
\pgfpathlineto{\pgfqpoint{1.252180in}{0.917253in}}%
\pgfpathlineto{\pgfqpoint{1.244276in}{0.909497in}}%
\pgfpathlineto{\pgfqpoint{1.236353in}{0.902235in}}%
\pgfpathclose%
\pgfusepath{fill}%
\end{pgfscope}%
\begin{pgfscope}%
\pgfpathrectangle{\pgfqpoint{0.407592in}{0.222904in}}{\pgfqpoint{1.411163in}{1.411163in}}%
\pgfusepath{clip}%
\pgfsetbuttcap%
\pgfsetroundjoin%
\definecolor{currentfill}{rgb}{0.909804,0.698039,0.423529}%
\pgfsetfillcolor{currentfill}%
\pgfsetfillopacity{0.120000}%
\pgfsetlinewidth{0.000000pt}%
\definecolor{currentstroke}{rgb}{0.000000,0.000000,0.000000}%
\pgfsetstrokecolor{currentstroke}%
\pgfsetdash{}{0pt}%
\pgfpathmoveto{\pgfqpoint{1.627883in}{1.188260in}}%
\pgfpathlineto{\pgfqpoint{1.645272in}{1.196206in}}%
\pgfpathlineto{\pgfqpoint{1.662735in}{1.204469in}}%
\pgfpathlineto{\pgfqpoint{1.670213in}{1.216196in}}%
\pgfpathlineto{\pgfqpoint{1.677679in}{1.228102in}}%
\pgfpathlineto{\pgfqpoint{1.660237in}{1.220105in}}%
\pgfpathlineto{\pgfqpoint{1.642868in}{1.212429in}}%
\pgfpathlineto{\pgfqpoint{1.635381in}{1.200256in}}%
\pgfpathlineto{\pgfqpoint{1.627883in}{1.188260in}}%
\pgfpathclose%
\pgfusepath{fill}%
\end{pgfscope}%
\begin{pgfscope}%
\pgfpathrectangle{\pgfqpoint{0.407592in}{0.222904in}}{\pgfqpoint{1.411163in}{1.411163in}}%
\pgfusepath{clip}%
\pgfsetbuttcap%
\pgfsetroundjoin%
\definecolor{currentfill}{rgb}{0.909804,0.698039,0.423529}%
\pgfsetfillcolor{currentfill}%
\pgfsetfillopacity{0.120000}%
\pgfsetlinewidth{0.000000pt}%
\definecolor{currentstroke}{rgb}{0.000000,0.000000,0.000000}%
\pgfsetstrokecolor{currentstroke}%
\pgfsetdash{}{0pt}%
\pgfpathmoveto{\pgfqpoint{1.285459in}{0.924167in}}%
\pgfpathlineto{\pgfqpoint{1.302192in}{0.928700in}}%
\pgfpathlineto{\pgfqpoint{1.318989in}{0.933887in}}%
\pgfpathlineto{\pgfqpoint{1.326831in}{0.941857in}}%
\pgfpathlineto{\pgfqpoint{1.334657in}{0.950256in}}%
\pgfpathlineto{\pgfqpoint{1.317881in}{0.945221in}}%
\pgfpathlineto{\pgfqpoint{1.301169in}{0.940834in}}%
\pgfpathlineto{\pgfqpoint{1.293323in}{0.932277in}}%
\pgfpathlineto{\pgfqpoint{1.285459in}{0.924167in}}%
\pgfpathclose%
\pgfusepath{fill}%
\end{pgfscope}%
\begin{pgfscope}%
\pgfpathrectangle{\pgfqpoint{0.407592in}{0.222904in}}{\pgfqpoint{1.411163in}{1.411163in}}%
\pgfusepath{clip}%
\pgfsetbuttcap%
\pgfsetroundjoin%
\definecolor{currentfill}{rgb}{0.909804,0.698039,0.423529}%
\pgfsetfillcolor{currentfill}%
\pgfsetfillopacity{0.120000}%
\pgfsetlinewidth{0.000000pt}%
\definecolor{currentstroke}{rgb}{0.000000,0.000000,0.000000}%
\pgfsetstrokecolor{currentstroke}%
\pgfsetdash{}{0pt}%
\pgfpathmoveto{\pgfqpoint{1.040092in}{0.870570in}}%
\pgfpathlineto{\pgfqpoint{1.056454in}{0.866168in}}%
\pgfpathlineto{\pgfqpoint{1.072852in}{0.862781in}}%
\pgfpathlineto{\pgfqpoint{1.081054in}{0.865646in}}%
\pgfpathlineto{\pgfqpoint{1.089227in}{0.869071in}}%
\pgfpathlineto{\pgfqpoint{1.072863in}{0.872427in}}%
\pgfpathlineto{\pgfqpoint{1.056536in}{0.876811in}}%
\pgfpathlineto{\pgfqpoint{1.048329in}{0.873411in}}%
\pgfpathlineto{\pgfqpoint{1.040092in}{0.870570in}}%
\pgfpathclose%
\pgfusepath{fill}%
\end{pgfscope}%
\begin{pgfscope}%
\pgfpathrectangle{\pgfqpoint{0.407592in}{0.222904in}}{\pgfqpoint{1.411163in}{1.411163in}}%
\pgfusepath{clip}%
\pgfsetbuttcap%
\pgfsetroundjoin%
\definecolor{currentfill}{rgb}{0.909804,0.698039,0.423529}%
\pgfsetfillcolor{currentfill}%
\pgfsetfillopacity{0.120000}%
\pgfsetlinewidth{0.000000pt}%
\definecolor{currentstroke}{rgb}{0.000000,0.000000,0.000000}%
\pgfsetstrokecolor{currentstroke}%
\pgfsetdash{}{0pt}%
\pgfpathmoveto{\pgfqpoint{0.711718in}{1.041165in}}%
\pgfpathlineto{\pgfqpoint{0.727937in}{1.025954in}}%
\pgfpathlineto{\pgfqpoint{0.744167in}{1.011138in}}%
\pgfpathlineto{\pgfqpoint{0.752976in}{1.011165in}}%
\pgfpathlineto{\pgfqpoint{0.761747in}{1.011550in}}%
\pgfpathlineto{\pgfqpoint{0.745560in}{1.026560in}}%
\pgfpathlineto{\pgfqpoint{0.729383in}{1.041960in}}%
\pgfpathlineto{\pgfqpoint{0.720571in}{1.041390in}}%
\pgfpathlineto{\pgfqpoint{0.711718in}{1.041165in}}%
\pgfpathclose%
\pgfusepath{fill}%
\end{pgfscope}%
\begin{pgfscope}%
\pgfpathrectangle{\pgfqpoint{0.407592in}{0.222904in}}{\pgfqpoint{1.411163in}{1.411163in}}%
\pgfusepath{clip}%
\pgfsetbuttcap%
\pgfsetroundjoin%
\definecolor{currentfill}{rgb}{0.909804,0.698039,0.423529}%
\pgfsetfillcolor{currentfill}%
\pgfsetfillopacity{0.120000}%
\pgfsetlinewidth{0.000000pt}%
\definecolor{currentstroke}{rgb}{0.000000,0.000000,0.000000}%
\pgfsetstrokecolor{currentstroke}%
\pgfsetdash{}{0pt}%
\pgfpathmoveto{\pgfqpoint{1.448820in}{1.031998in}}%
\pgfpathlineto{\pgfqpoint{1.465877in}{1.039017in}}%
\pgfpathlineto{\pgfqpoint{1.483004in}{1.046464in}}%
\pgfpathlineto{\pgfqpoint{1.490666in}{1.056677in}}%
\pgfpathlineto{\pgfqpoint{1.498315in}{1.067173in}}%
\pgfpathlineto{\pgfqpoint{1.481207in}{1.059983in}}%
\pgfpathlineto{\pgfqpoint{1.464170in}{1.053222in}}%
\pgfpathlineto{\pgfqpoint{1.456502in}{1.042465in}}%
\pgfpathlineto{\pgfqpoint{1.448820in}{1.031998in}}%
\pgfpathclose%
\pgfusepath{fill}%
\end{pgfscope}%
\begin{pgfscope}%
\pgfpathrectangle{\pgfqpoint{0.407592in}{0.222904in}}{\pgfqpoint{1.411163in}{1.411163in}}%
\pgfusepath{clip}%
\pgfsetbuttcap%
\pgfsetroundjoin%
\definecolor{currentfill}{rgb}{0.909804,0.698039,0.423529}%
\pgfsetfillcolor{currentfill}%
\pgfsetfillopacity{0.120000}%
\pgfsetlinewidth{0.000000pt}%
\definecolor{currentstroke}{rgb}{0.000000,0.000000,0.000000}%
\pgfsetstrokecolor{currentstroke}%
\pgfsetdash{}{0pt}%
\pgfpathmoveto{\pgfqpoint{0.596539in}{1.138540in}}%
\pgfpathlineto{\pgfqpoint{0.612740in}{1.121707in}}%
\pgfpathlineto{\pgfqpoint{0.628950in}{1.105123in}}%
\pgfpathlineto{\pgfqpoint{0.637969in}{1.105095in}}%
\pgfpathlineto{\pgfqpoint{0.646948in}{1.105370in}}%
\pgfpathlineto{\pgfqpoint{0.630780in}{1.122133in}}%
\pgfpathlineto{\pgfqpoint{0.614622in}{1.139137in}}%
\pgfpathlineto{\pgfqpoint{0.605601in}{1.138693in}}%
\pgfpathlineto{\pgfqpoint{0.596539in}{1.138540in}}%
\pgfpathclose%
\pgfusepath{fill}%
\end{pgfscope}%
\begin{pgfscope}%
\pgfpathrectangle{\pgfqpoint{0.407592in}{0.222904in}}{\pgfqpoint{1.411163in}{1.411163in}}%
\pgfusepath{clip}%
\pgfsetbuttcap%
\pgfsetroundjoin%
\definecolor{currentfill}{rgb}{0.909804,0.698039,0.423529}%
\pgfsetfillcolor{currentfill}%
\pgfsetfillopacity{0.120000}%
\pgfsetlinewidth{0.000000pt}%
\definecolor{currentstroke}{rgb}{0.000000,0.000000,0.000000}%
\pgfsetstrokecolor{currentstroke}%
\pgfsetdash{}{0pt}%
\pgfpathmoveto{\pgfqpoint{1.563132in}{1.126268in}}%
\pgfpathlineto{\pgfqpoint{1.580416in}{1.134075in}}%
\pgfpathlineto{\pgfqpoint{1.597773in}{1.142228in}}%
\pgfpathlineto{\pgfqpoint{1.605318in}{1.153426in}}%
\pgfpathlineto{\pgfqpoint{1.612852in}{1.164837in}}%
\pgfpathlineto{\pgfqpoint{1.595515in}{1.156953in}}%
\pgfpathlineto{\pgfqpoint{1.578252in}{1.149419in}}%
\pgfpathlineto{\pgfqpoint{1.570698in}{1.137736in}}%
\pgfpathlineto{\pgfqpoint{1.563132in}{1.126268in}}%
\pgfpathclose%
\pgfusepath{fill}%
\end{pgfscope}%
\begin{pgfscope}%
\pgfpathrectangle{\pgfqpoint{0.407592in}{0.222904in}}{\pgfqpoint{1.411163in}{1.411163in}}%
\pgfusepath{clip}%
\pgfsetbuttcap%
\pgfsetroundjoin%
\definecolor{currentfill}{rgb}{0.909804,0.698039,0.423529}%
\pgfsetfillcolor{currentfill}%
\pgfsetfillopacity{0.120000}%
\pgfsetlinewidth{0.000000pt}%
\definecolor{currentstroke}{rgb}{0.000000,0.000000,0.000000}%
\pgfsetstrokecolor{currentstroke}%
\pgfsetdash{}{0pt}%
\pgfpathmoveto{\pgfqpoint{0.908802in}{0.910621in}}%
\pgfpathlineto{\pgfqpoint{0.925080in}{0.900766in}}%
\pgfpathlineto{\pgfqpoint{0.941378in}{0.891731in}}%
\pgfpathlineto{\pgfqpoint{0.949825in}{0.892571in}}%
\pgfpathlineto{\pgfqpoint{0.958237in}{0.893896in}}%
\pgfpathlineto{\pgfqpoint{0.941979in}{0.903026in}}%
\pgfpathlineto{\pgfqpoint{0.925742in}{0.912991in}}%
\pgfpathlineto{\pgfqpoint{0.917290in}{0.911570in}}%
\pgfpathlineto{\pgfqpoint{0.908802in}{0.910621in}}%
\pgfpathclose%
\pgfusepath{fill}%
\end{pgfscope}%
\begin{pgfscope}%
\pgfpathrectangle{\pgfqpoint{0.407592in}{0.222904in}}{\pgfqpoint{1.411163in}{1.411163in}}%
\pgfusepath{clip}%
\pgfsetbuttcap%
\pgfsetroundjoin%
\definecolor{currentfill}{rgb}{0.909804,0.698039,0.423529}%
\pgfsetfillcolor{currentfill}%
\pgfsetfillopacity{0.120000}%
\pgfsetlinewidth{0.000000pt}%
\definecolor{currentstroke}{rgb}{0.000000,0.000000,0.000000}%
\pgfsetstrokecolor{currentstroke}%
\pgfsetdash{}{0pt}%
\pgfpathmoveto{\pgfqpoint{0.826616in}{0.956127in}}%
\pgfpathlineto{\pgfqpoint{0.842867in}{0.943626in}}%
\pgfpathlineto{\pgfqpoint{0.859133in}{0.931753in}}%
\pgfpathlineto{\pgfqpoint{0.867737in}{0.931969in}}%
\pgfpathlineto{\pgfqpoint{0.876303in}{0.932611in}}%
\pgfpathlineto{\pgfqpoint{0.860079in}{0.944651in}}%
\pgfpathlineto{\pgfqpoint{0.843871in}{0.957324in}}%
\pgfpathlineto{\pgfqpoint{0.835262in}{0.956520in}}%
\pgfpathlineto{\pgfqpoint{0.826616in}{0.956127in}}%
\pgfpathclose%
\pgfusepath{fill}%
\end{pgfscope}%
\begin{pgfscope}%
\pgfpathrectangle{\pgfqpoint{0.407592in}{0.222904in}}{\pgfqpoint{1.411163in}{1.411163in}}%
\pgfusepath{clip}%
\pgfsetbuttcap%
\pgfsetroundjoin%
\definecolor{currentfill}{rgb}{0.909804,0.698039,0.423529}%
\pgfsetfillcolor{currentfill}%
\pgfsetfillopacity{0.120000}%
\pgfsetlinewidth{0.000000pt}%
\definecolor{currentstroke}{rgb}{0.000000,0.000000,0.000000}%
\pgfsetstrokecolor{currentstroke}%
\pgfsetdash{}{0pt}%
\pgfpathmoveto{\pgfqpoint{1.334657in}{0.950256in}}%
\pgfpathlineto{\pgfqpoint{1.351499in}{0.955898in}}%
\pgfpathlineto{\pgfqpoint{1.368408in}{0.962103in}}%
\pgfpathlineto{\pgfqpoint{1.376197in}{0.970733in}}%
\pgfpathlineto{\pgfqpoint{1.383970in}{0.979747in}}%
\pgfpathlineto{\pgfqpoint{1.367081in}{0.973734in}}%
\pgfpathlineto{\pgfqpoint{1.350260in}{0.968282in}}%
\pgfpathlineto{\pgfqpoint{1.342466in}{0.959069in}}%
\pgfpathlineto{\pgfqpoint{1.334657in}{0.950256in}}%
\pgfpathclose%
\pgfusepath{fill}%
\end{pgfscope}%
\begin{pgfscope}%
\pgfpathrectangle{\pgfqpoint{0.407592in}{0.222904in}}{\pgfqpoint{1.411163in}{1.411163in}}%
\pgfusepath{clip}%
\pgfsetbuttcap%
\pgfsetroundjoin%
\definecolor{currentfill}{rgb}{0.909804,0.698039,0.423529}%
\pgfsetfillcolor{currentfill}%
\pgfsetfillopacity{0.120000}%
\pgfsetlinewidth{0.000000pt}%
\definecolor{currentstroke}{rgb}{0.000000,0.000000,0.000000}%
\pgfsetstrokecolor{currentstroke}%
\pgfsetdash{}{0pt}%
\pgfpathmoveto{\pgfqpoint{0.990824in}{0.878311in}}%
\pgfpathlineto{\pgfqpoint{1.007159in}{0.871928in}}%
\pgfpathlineto{\pgfqpoint{1.023524in}{0.866521in}}%
\pgfpathlineto{\pgfqpoint{1.031824in}{0.868278in}}%
\pgfpathlineto{\pgfqpoint{1.040092in}{0.870570in}}%
\pgfpathlineto{\pgfqpoint{1.023765in}{0.875982in}}%
\pgfpathlineto{\pgfqpoint{1.007468in}{0.882387in}}%
\pgfpathlineto{\pgfqpoint{0.999163in}{0.880085in}}%
\pgfpathlineto{\pgfqpoint{0.990824in}{0.878311in}}%
\pgfpathclose%
\pgfusepath{fill}%
\end{pgfscope}%
\begin{pgfscope}%
\pgfpathrectangle{\pgfqpoint{0.407592in}{0.222904in}}{\pgfqpoint{1.411163in}{1.411163in}}%
\pgfusepath{clip}%
\pgfsetbuttcap%
\pgfsetroundjoin%
\definecolor{currentfill}{rgb}{0.909804,0.698039,0.423529}%
\pgfsetfillcolor{currentfill}%
\pgfsetfillopacity{0.120000}%
\pgfsetlinewidth{0.000000pt}%
\definecolor{currentstroke}{rgb}{0.000000,0.000000,0.000000}%
\pgfsetstrokecolor{currentstroke}%
\pgfsetdash{}{0pt}%
\pgfpathmoveto{\pgfqpoint{0.744167in}{1.011138in}}%
\pgfpathlineto{\pgfqpoint{0.760408in}{0.996751in}}%
\pgfpathlineto{\pgfqpoint{0.776661in}{0.982827in}}%
\pgfpathlineto{\pgfqpoint{0.785428in}{0.982647in}}%
\pgfpathlineto{\pgfqpoint{0.794156in}{0.982838in}}%
\pgfpathlineto{\pgfqpoint{0.777945in}{0.996964in}}%
\pgfpathlineto{\pgfqpoint{0.761747in}{1.011550in}}%
\pgfpathlineto{\pgfqpoint{0.752976in}{1.011165in}}%
\pgfpathlineto{\pgfqpoint{0.744167in}{1.011138in}}%
\pgfpathclose%
\pgfusepath{fill}%
\end{pgfscope}%
\begin{pgfscope}%
\pgfpathrectangle{\pgfqpoint{0.407592in}{0.222904in}}{\pgfqpoint{1.411163in}{1.411163in}}%
\pgfusepath{clip}%
\pgfsetbuttcap%
\pgfsetroundjoin%
\definecolor{currentfill}{rgb}{0.909804,0.698039,0.423529}%
\pgfsetfillcolor{currentfill}%
\pgfsetfillopacity{0.120000}%
\pgfsetlinewidth{0.000000pt}%
\definecolor{currentstroke}{rgb}{0.000000,0.000000,0.000000}%
\pgfsetstrokecolor{currentstroke}%
\pgfsetdash{}{0pt}%
\pgfpathmoveto{\pgfqpoint{0.628950in}{1.105123in}}%
\pgfpathlineto{\pgfqpoint{0.645170in}{1.088807in}}%
\pgfpathlineto{\pgfqpoint{0.661400in}{1.072783in}}%
\pgfpathlineto{\pgfqpoint{0.670377in}{1.072559in}}%
\pgfpathlineto{\pgfqpoint{0.679313in}{1.072649in}}%
\pgfpathlineto{\pgfqpoint{0.663125in}{1.088867in}}%
\pgfpathlineto{\pgfqpoint{0.646948in}{1.105370in}}%
\pgfpathlineto{\pgfqpoint{0.637969in}{1.105095in}}%
\pgfpathlineto{\pgfqpoint{0.628950in}{1.105123in}}%
\pgfpathclose%
\pgfusepath{fill}%
\end{pgfscope}%
\begin{pgfscope}%
\pgfpathrectangle{\pgfqpoint{0.407592in}{0.222904in}}{\pgfqpoint{1.411163in}{1.411163in}}%
\pgfusepath{clip}%
\pgfsetbuttcap%
\pgfsetroundjoin%
\definecolor{currentfill}{rgb}{0.909804,0.698039,0.423529}%
\pgfsetfillcolor{currentfill}%
\pgfsetfillopacity{0.120000}%
\pgfsetlinewidth{0.000000pt}%
\definecolor{currentstroke}{rgb}{0.000000,0.000000,0.000000}%
\pgfsetstrokecolor{currentstroke}%
\pgfsetdash{}{0pt}%
\pgfpathmoveto{\pgfqpoint{1.677679in}{1.228102in}}%
\pgfpathlineto{\pgfqpoint{1.695195in}{1.236405in}}%
\pgfpathlineto{\pgfqpoint{1.712787in}{1.245000in}}%
\pgfpathlineto{\pgfqpoint{1.720220in}{1.256815in}}%
\pgfpathlineto{\pgfqpoint{1.702639in}{1.248349in}}%
\pgfpathlineto{\pgfqpoint{1.685134in}{1.240178in}}%
\pgfusepath{fill}%
\end{pgfscope}%
\begin{pgfscope}%
\pgfpathrectangle{\pgfqpoint{0.407592in}{0.222904in}}{\pgfqpoint{1.411163in}{1.411163in}}%
\pgfusepath{clip}%
\pgfsetbuttcap%
\pgfsetroundjoin%
\definecolor{currentfill}{rgb}{0.909804,0.698039,0.423529}%
\pgfsetfillcolor{currentfill}%
\pgfsetfillopacity{0.120000}%
\pgfsetlinewidth{0.000000pt}%
\definecolor{currentstroke}{rgb}{0.000000,0.000000,0.000000}%
\pgfsetstrokecolor{currentstroke}%
\pgfsetdash{}{0pt}%
\pgfpathmoveto{\pgfqpoint{1.498315in}{1.067173in}}%
\pgfpathlineto{\pgfqpoint{1.515494in}{1.074769in}}%
\pgfpathlineto{\pgfqpoint{1.532746in}{1.082745in}}%
\pgfpathlineto{\pgfqpoint{1.540361in}{1.093253in}}%
\pgfpathlineto{\pgfqpoint{1.547964in}{1.104016in}}%
\pgfpathlineto{\pgfqpoint{1.530733in}{1.096303in}}%
\pgfpathlineto{\pgfqpoint{1.513573in}{1.088973in}}%
\pgfpathlineto{\pgfqpoint{1.505950in}{1.077943in}}%
\pgfpathlineto{\pgfqpoint{1.498315in}{1.067173in}}%
\pgfpathclose%
\pgfusepath{fill}%
\end{pgfscope}%
\begin{pgfscope}%
\pgfpathrectangle{\pgfqpoint{0.407592in}{0.222904in}}{\pgfqpoint{1.411163in}{1.411163in}}%
\pgfusepath{clip}%
\pgfsetbuttcap%
\pgfsetroundjoin%
\definecolor{currentfill}{rgb}{0.909804,0.698039,0.423529}%
\pgfsetfillcolor{currentfill}%
\pgfsetfillopacity{0.120000}%
\pgfsetlinewidth{0.000000pt}%
\definecolor{currentstroke}{rgb}{0.000000,0.000000,0.000000}%
\pgfsetstrokecolor{currentstroke}%
\pgfsetdash{}{0pt}%
\pgfpathmoveto{\pgfqpoint{1.383970in}{0.979747in}}%
\pgfpathlineto{\pgfqpoint{1.400927in}{0.986284in}}%
\pgfpathlineto{\pgfqpoint{1.417954in}{0.993309in}}%
\pgfpathlineto{\pgfqpoint{1.425692in}{1.002479in}}%
\pgfpathlineto{\pgfqpoint{1.433415in}{1.011992in}}%
\pgfpathlineto{\pgfqpoint{1.416409in}{1.005188in}}%
\pgfpathlineto{\pgfqpoint{1.399472in}{0.998872in}}%
\pgfpathlineto{\pgfqpoint{1.391728in}{0.989132in}}%
\pgfpathlineto{\pgfqpoint{1.383970in}{0.979747in}}%
\pgfpathclose%
\pgfusepath{fill}%
\end{pgfscope}%
\begin{pgfscope}%
\pgfpathrectangle{\pgfqpoint{0.407592in}{0.222904in}}{\pgfqpoint{1.411163in}{1.411163in}}%
\pgfusepath{clip}%
\pgfsetbuttcap%
\pgfsetroundjoin%
\definecolor{currentfill}{rgb}{0.909804,0.698039,0.423529}%
\pgfsetfillcolor{currentfill}%
\pgfsetfillopacity{0.120000}%
\pgfsetlinewidth{0.000000pt}%
\definecolor{currentstroke}{rgb}{0.000000,0.000000,0.000000}%
\pgfsetstrokecolor{currentstroke}%
\pgfsetdash{}{0pt}%
\pgfpathmoveto{\pgfqpoint{1.612852in}{1.164837in}}%
\pgfpathlineto{\pgfqpoint{1.630262in}{1.173051in}}%
\pgfpathlineto{\pgfqpoint{1.647746in}{1.181579in}}%
\pgfpathlineto{\pgfqpoint{1.655247in}{1.192927in}}%
\pgfpathlineto{\pgfqpoint{1.662735in}{1.204469in}}%
\pgfpathlineto{\pgfqpoint{1.645272in}{1.196206in}}%
\pgfpathlineto{\pgfqpoint{1.627883in}{1.188260in}}%
\pgfpathlineto{\pgfqpoint{1.620373in}{1.176451in}}%
\pgfpathlineto{\pgfqpoint{1.612852in}{1.164837in}}%
\pgfpathclose%
\pgfusepath{fill}%
\end{pgfscope}%
\begin{pgfscope}%
\pgfpathrectangle{\pgfqpoint{0.407592in}{0.222904in}}{\pgfqpoint{1.411163in}{1.411163in}}%
\pgfusepath{clip}%
\pgfsetbuttcap%
\pgfsetroundjoin%
\definecolor{currentfill}{rgb}{0.909804,0.698039,0.423529}%
\pgfsetfillcolor{currentfill}%
\pgfsetfillopacity{0.120000}%
\pgfsetlinewidth{0.000000pt}%
\definecolor{currentstroke}{rgb}{0.000000,0.000000,0.000000}%
\pgfsetstrokecolor{currentstroke}%
\pgfsetdash{}{0pt}%
\pgfpathmoveto{\pgfqpoint{0.513409in}{1.207856in}}%
\pgfpathlineto{\pgfqpoint{0.529615in}{1.190365in}}%
\pgfpathlineto{\pgfqpoint{0.545830in}{1.173062in}}%
\pgfpathlineto{\pgfqpoint{0.555019in}{1.172835in}}%
\pgfpathlineto{\pgfqpoint{0.564166in}{1.172877in}}%
\pgfpathlineto{\pgfqpoint{0.547993in}{1.190347in}}%
\pgfpathlineto{\pgfqpoint{0.531830in}{1.207998in}}%
\pgfpathlineto{\pgfqpoint{0.522640in}{1.207797in}}%
\pgfpathlineto{\pgfqpoint{0.513409in}{1.207856in}}%
\pgfpathclose%
\pgfusepath{fill}%
\end{pgfscope}%
\begin{pgfscope}%
\pgfpathrectangle{\pgfqpoint{0.407592in}{0.222904in}}{\pgfqpoint{1.411163in}{1.411163in}}%
\pgfusepath{clip}%
\pgfsetbuttcap%
\pgfsetroundjoin%
\definecolor{currentfill}{rgb}{0.909804,0.698039,0.423529}%
\pgfsetfillcolor{currentfill}%
\pgfsetfillopacity{0.120000}%
\pgfsetlinewidth{0.000000pt}%
\definecolor{currentstroke}{rgb}{0.000000,0.000000,0.000000}%
\pgfsetstrokecolor{currentstroke}%
\pgfsetdash{}{0pt}%
\pgfpathmoveto{\pgfqpoint{1.171262in}{0.874370in}}%
\pgfpathlineto{\pgfqpoint{1.187830in}{0.875922in}}%
\pgfpathlineto{\pgfqpoint{1.204453in}{0.878338in}}%
\pgfpathlineto{\pgfqpoint{1.212460in}{0.883525in}}%
\pgfpathlineto{\pgfqpoint{1.220446in}{0.889240in}}%
\pgfpathlineto{\pgfqpoint{1.203849in}{0.886826in}}%
\pgfpathlineto{\pgfqpoint{1.187307in}{0.885272in}}%
\pgfpathlineto{\pgfqpoint{1.179296in}{0.879546in}}%
\pgfpathlineto{\pgfqpoint{1.171262in}{0.874370in}}%
\pgfpathclose%
\pgfusepath{fill}%
\end{pgfscope}%
\begin{pgfscope}%
\pgfpathrectangle{\pgfqpoint{0.407592in}{0.222904in}}{\pgfqpoint{1.411163in}{1.411163in}}%
\pgfusepath{clip}%
\pgfsetbuttcap%
\pgfsetroundjoin%
\definecolor{currentfill}{rgb}{0.909804,0.698039,0.423529}%
\pgfsetfillcolor{currentfill}%
\pgfsetfillopacity{0.120000}%
\pgfsetlinewidth{0.000000pt}%
\definecolor{currentstroke}{rgb}{0.000000,0.000000,0.000000}%
\pgfsetstrokecolor{currentstroke}%
\pgfsetdash{}{0pt}%
\pgfpathmoveto{\pgfqpoint{1.122080in}{0.865411in}}%
\pgfpathlineto{\pgfqpoint{1.138575in}{0.865066in}}%
\pgfpathlineto{\pgfqpoint{1.155119in}{0.865670in}}%
\pgfpathlineto{\pgfqpoint{1.163203in}{0.869745in}}%
\pgfpathlineto{\pgfqpoint{1.171262in}{0.874370in}}%
\pgfpathlineto{\pgfqpoint{1.154747in}{0.873727in}}%
\pgfpathlineto{\pgfqpoint{1.138281in}{0.874035in}}%
\pgfpathlineto{\pgfqpoint{1.130194in}{0.869440in}}%
\pgfpathlineto{\pgfqpoint{1.122080in}{0.865411in}}%
\pgfpathclose%
\pgfusepath{fill}%
\end{pgfscope}%
\begin{pgfscope}%
\pgfpathrectangle{\pgfqpoint{0.407592in}{0.222904in}}{\pgfqpoint{1.411163in}{1.411163in}}%
\pgfusepath{clip}%
\pgfsetbuttcap%
\pgfsetroundjoin%
\definecolor{currentfill}{rgb}{0.909804,0.698039,0.423529}%
\pgfsetfillcolor{currentfill}%
\pgfsetfillopacity{0.120000}%
\pgfsetlinewidth{0.000000pt}%
\definecolor{currentstroke}{rgb}{0.000000,0.000000,0.000000}%
\pgfsetstrokecolor{currentstroke}%
\pgfsetdash{}{0pt}%
\pgfpathmoveto{\pgfqpoint{0.859133in}{0.931753in}}%
\pgfpathlineto{\pgfqpoint{0.875416in}{0.920553in}}%
\pgfpathlineto{\pgfqpoint{0.891716in}{0.910071in}}%
\pgfpathlineto{\pgfqpoint{0.900278in}{0.910127in}}%
\pgfpathlineto{\pgfqpoint{0.908802in}{0.910621in}}%
\pgfpathlineto{\pgfqpoint{0.892544in}{0.921252in}}%
\pgfpathlineto{\pgfqpoint{0.876303in}{0.932611in}}%
\pgfpathlineto{\pgfqpoint{0.867737in}{0.931969in}}%
\pgfpathlineto{\pgfqpoint{0.859133in}{0.931753in}}%
\pgfpathclose%
\pgfusepath{fill}%
\end{pgfscope}%
\begin{pgfscope}%
\pgfpathrectangle{\pgfqpoint{0.407592in}{0.222904in}}{\pgfqpoint{1.411163in}{1.411163in}}%
\pgfusepath{clip}%
\pgfsetbuttcap%
\pgfsetroundjoin%
\definecolor{currentfill}{rgb}{0.909804,0.698039,0.423529}%
\pgfsetfillcolor{currentfill}%
\pgfsetfillopacity{0.120000}%
\pgfsetlinewidth{0.000000pt}%
\definecolor{currentstroke}{rgb}{0.000000,0.000000,0.000000}%
\pgfsetstrokecolor{currentstroke}%
\pgfsetdash{}{0pt}%
\pgfpathmoveto{\pgfqpoint{1.220446in}{0.889240in}}%
\pgfpathlineto{\pgfqpoint{1.237101in}{0.892466in}}%
\pgfpathlineto{\pgfqpoint{1.253817in}{0.896454in}}%
\pgfpathlineto{\pgfqpoint{1.261757in}{0.902651in}}%
\pgfpathlineto{\pgfqpoint{1.269677in}{0.909342in}}%
\pgfpathlineto{\pgfqpoint{1.252984in}{0.905410in}}%
\pgfpathlineto{\pgfqpoint{1.236353in}{0.902235in}}%
\pgfpathlineto{\pgfqpoint{1.228410in}{0.895480in}}%
\pgfpathlineto{\pgfqpoint{1.220446in}{0.889240in}}%
\pgfpathclose%
\pgfusepath{fill}%
\end{pgfscope}%
\begin{pgfscope}%
\pgfpathrectangle{\pgfqpoint{0.407592in}{0.222904in}}{\pgfqpoint{1.411163in}{1.411163in}}%
\pgfusepath{clip}%
\pgfsetbuttcap%
\pgfsetroundjoin%
\definecolor{currentfill}{rgb}{0.909804,0.698039,0.423529}%
\pgfsetfillcolor{currentfill}%
\pgfsetfillopacity{0.120000}%
\pgfsetlinewidth{0.000000pt}%
\definecolor{currentstroke}{rgb}{0.000000,0.000000,0.000000}%
\pgfsetstrokecolor{currentstroke}%
\pgfsetdash{}{0pt}%
\pgfpathmoveto{\pgfqpoint{0.941378in}{0.891731in}}%
\pgfpathlineto{\pgfqpoint{0.957700in}{0.883558in}}%
\pgfpathlineto{\pgfqpoint{0.974046in}{0.876284in}}%
\pgfpathlineto{\pgfqpoint{0.982453in}{0.877050in}}%
\pgfpathlineto{\pgfqpoint{0.990824in}{0.878311in}}%
\pgfpathlineto{\pgfqpoint{0.974518in}{0.885645in}}%
\pgfpathlineto{\pgfqpoint{0.958237in}{0.893896in}}%
\pgfpathlineto{\pgfqpoint{0.949825in}{0.892571in}}%
\pgfpathlineto{\pgfqpoint{0.941378in}{0.891731in}}%
\pgfpathclose%
\pgfusepath{fill}%
\end{pgfscope}%
\begin{pgfscope}%
\pgfpathrectangle{\pgfqpoint{0.407592in}{0.222904in}}{\pgfqpoint{1.411163in}{1.411163in}}%
\pgfusepath{clip}%
\pgfsetbuttcap%
\pgfsetroundjoin%
\definecolor{currentfill}{rgb}{0.909804,0.698039,0.423529}%
\pgfsetfillcolor{currentfill}%
\pgfsetfillopacity{0.120000}%
\pgfsetlinewidth{0.000000pt}%
\definecolor{currentstroke}{rgb}{0.000000,0.000000,0.000000}%
\pgfsetstrokecolor{currentstroke}%
\pgfsetdash{}{0pt}%
\pgfpathmoveto{\pgfqpoint{0.661400in}{1.072783in}}%
\pgfpathlineto{\pgfqpoint{0.677641in}{1.057074in}}%
\pgfpathlineto{\pgfqpoint{0.693891in}{1.041708in}}%
\pgfpathlineto{\pgfqpoint{0.702825in}{1.041274in}}%
\pgfpathlineto{\pgfqpoint{0.711718in}{1.041165in}}%
\pgfpathlineto{\pgfqpoint{0.695510in}{1.056739in}}%
\pgfpathlineto{\pgfqpoint{0.679313in}{1.072649in}}%
\pgfpathlineto{\pgfqpoint{0.670377in}{1.072559in}}%
\pgfpathlineto{\pgfqpoint{0.661400in}{1.072783in}}%
\pgfpathclose%
\pgfusepath{fill}%
\end{pgfscope}%
\begin{pgfscope}%
\pgfpathrectangle{\pgfqpoint{0.407592in}{0.222904in}}{\pgfqpoint{1.411163in}{1.411163in}}%
\pgfusepath{clip}%
\pgfsetbuttcap%
\pgfsetroundjoin%
\definecolor{currentfill}{rgb}{0.909804,0.698039,0.423529}%
\pgfsetfillcolor{currentfill}%
\pgfsetfillopacity{0.120000}%
\pgfsetlinewidth{0.000000pt}%
\definecolor{currentstroke}{rgb}{0.000000,0.000000,0.000000}%
\pgfsetstrokecolor{currentstroke}%
\pgfsetdash{}{0pt}%
\pgfpathmoveto{\pgfqpoint{1.072852in}{0.862781in}}%
\pgfpathlineto{\pgfqpoint{1.089290in}{0.860406in}}%
\pgfpathlineto{\pgfqpoint{1.105771in}{0.859027in}}%
\pgfpathlineto{\pgfqpoint{1.113940in}{0.861942in}}%
\pgfpathlineto{\pgfqpoint{1.122080in}{0.865411in}}%
\pgfpathlineto{\pgfqpoint{1.105632in}{0.866737in}}%
\pgfpathlineto{\pgfqpoint{1.089227in}{0.869071in}}%
\pgfpathlineto{\pgfqpoint{1.081054in}{0.865646in}}%
\pgfpathlineto{\pgfqpoint{1.072852in}{0.862781in}}%
\pgfpathclose%
\pgfusepath{fill}%
\end{pgfscope}%
\begin{pgfscope}%
\pgfpathrectangle{\pgfqpoint{0.407592in}{0.222904in}}{\pgfqpoint{1.411163in}{1.411163in}}%
\pgfusepath{clip}%
\pgfsetbuttcap%
\pgfsetroundjoin%
\definecolor{currentfill}{rgb}{0.909804,0.698039,0.423529}%
\pgfsetfillcolor{currentfill}%
\pgfsetfillopacity{0.120000}%
\pgfsetlinewidth{0.000000pt}%
\definecolor{currentstroke}{rgb}{0.000000,0.000000,0.000000}%
\pgfsetstrokecolor{currentstroke}%
\pgfsetdash{}{0pt}%
\pgfpathmoveto{\pgfqpoint{0.776661in}{0.982827in}}%
\pgfpathlineto{\pgfqpoint{0.792928in}{0.969403in}}%
\pgfpathlineto{\pgfqpoint{0.809207in}{0.956520in}}%
\pgfpathlineto{\pgfqpoint{0.817931in}{0.956132in}}%
\pgfpathlineto{\pgfqpoint{0.826616in}{0.956127in}}%
\pgfpathlineto{\pgfqpoint{0.810379in}{0.969212in}}%
\pgfpathlineto{\pgfqpoint{0.794156in}{0.982838in}}%
\pgfpathlineto{\pgfqpoint{0.785428in}{0.982647in}}%
\pgfpathlineto{\pgfqpoint{0.776661in}{0.982827in}}%
\pgfpathclose%
\pgfusepath{fill}%
\end{pgfscope}%
\begin{pgfscope}%
\pgfpathrectangle{\pgfqpoint{0.407592in}{0.222904in}}{\pgfqpoint{1.411163in}{1.411163in}}%
\pgfusepath{clip}%
\pgfsetbuttcap%
\pgfsetroundjoin%
\definecolor{currentfill}{rgb}{0.909804,0.698039,0.423529}%
\pgfsetfillcolor{currentfill}%
\pgfsetfillopacity{0.120000}%
\pgfsetlinewidth{0.000000pt}%
\definecolor{currentstroke}{rgb}{0.000000,0.000000,0.000000}%
\pgfsetstrokecolor{currentstroke}%
\pgfsetdash{}{0pt}%
\pgfpathmoveto{\pgfqpoint{1.269677in}{0.909342in}}%
\pgfpathlineto{\pgfqpoint{1.286431in}{0.913981in}}%
\pgfpathlineto{\pgfqpoint{1.303250in}{0.919278in}}%
\pgfpathlineto{\pgfqpoint{1.311128in}{0.926356in}}%
\pgfpathlineto{\pgfqpoint{1.318989in}{0.933887in}}%
\pgfpathlineto{\pgfqpoint{1.302192in}{0.928700in}}%
\pgfpathlineto{\pgfqpoint{1.285459in}{0.924167in}}%
\pgfpathlineto{\pgfqpoint{1.277577in}{0.916518in}}%
\pgfpathlineto{\pgfqpoint{1.269677in}{0.909342in}}%
\pgfpathclose%
\pgfusepath{fill}%
\end{pgfscope}%
\begin{pgfscope}%
\pgfpathrectangle{\pgfqpoint{0.407592in}{0.222904in}}{\pgfqpoint{1.411163in}{1.411163in}}%
\pgfusepath{clip}%
\pgfsetbuttcap%
\pgfsetroundjoin%
\definecolor{currentfill}{rgb}{0.909804,0.698039,0.423529}%
\pgfsetfillcolor{currentfill}%
\pgfsetfillopacity{0.120000}%
\pgfsetlinewidth{0.000000pt}%
\definecolor{currentstroke}{rgb}{0.000000,0.000000,0.000000}%
\pgfsetstrokecolor{currentstroke}%
\pgfsetdash{}{0pt}%
\pgfpathmoveto{\pgfqpoint{1.433415in}{1.011992in}}%
\pgfpathlineto{\pgfqpoint{1.450492in}{1.019253in}}%
\pgfpathlineto{\pgfqpoint{1.467640in}{1.026940in}}%
\pgfpathlineto{\pgfqpoint{1.475329in}{1.036548in}}%
\pgfpathlineto{\pgfqpoint{1.483004in}{1.046464in}}%
\pgfpathlineto{\pgfqpoint{1.465877in}{1.039017in}}%
\pgfpathlineto{\pgfqpoint{1.448820in}{1.031998in}}%
\pgfpathlineto{\pgfqpoint{1.441124in}{1.021836in}}%
\pgfpathlineto{\pgfqpoint{1.433415in}{1.011992in}}%
\pgfpathclose%
\pgfusepath{fill}%
\end{pgfscope}%
\begin{pgfscope}%
\pgfpathrectangle{\pgfqpoint{0.407592in}{0.222904in}}{\pgfqpoint{1.411163in}{1.411163in}}%
\pgfusepath{clip}%
\pgfsetbuttcap%
\pgfsetroundjoin%
\definecolor{currentfill}{rgb}{0.909804,0.698039,0.423529}%
\pgfsetfillcolor{currentfill}%
\pgfsetfillopacity{0.120000}%
\pgfsetlinewidth{0.000000pt}%
\definecolor{currentstroke}{rgb}{0.000000,0.000000,0.000000}%
\pgfsetstrokecolor{currentstroke}%
\pgfsetdash{}{0pt}%
\pgfpathmoveto{\pgfqpoint{1.547964in}{1.104016in}}%
\pgfpathlineto{\pgfqpoint{1.565268in}{1.112091in}}%
\pgfpathlineto{\pgfqpoint{1.582646in}{1.120508in}}%
\pgfpathlineto{\pgfqpoint{1.590216in}{1.131252in}}%
\pgfpathlineto{\pgfqpoint{1.597773in}{1.142228in}}%
\pgfpathlineto{\pgfqpoint{1.580416in}{1.134075in}}%
\pgfpathlineto{\pgfqpoint{1.563132in}{1.126268in}}%
\pgfpathlineto{\pgfqpoint{1.555554in}{1.115024in}}%
\pgfpathlineto{\pgfqpoint{1.547964in}{1.104016in}}%
\pgfpathclose%
\pgfusepath{fill}%
\end{pgfscope}%
\begin{pgfscope}%
\pgfpathrectangle{\pgfqpoint{0.407592in}{0.222904in}}{\pgfqpoint{1.411163in}{1.411163in}}%
\pgfusepath{clip}%
\pgfsetbuttcap%
\pgfsetroundjoin%
\definecolor{currentfill}{rgb}{0.909804,0.698039,0.423529}%
\pgfsetfillcolor{currentfill}%
\pgfsetfillopacity{0.120000}%
\pgfsetlinewidth{0.000000pt}%
\definecolor{currentstroke}{rgb}{0.000000,0.000000,0.000000}%
\pgfsetstrokecolor{currentstroke}%
\pgfsetdash{}{0pt}%
\pgfpathmoveto{\pgfqpoint{0.545830in}{1.173062in}}%
\pgfpathlineto{\pgfqpoint{0.562055in}{1.155963in}}%
\pgfpathlineto{\pgfqpoint{0.578289in}{1.139084in}}%
\pgfpathlineto{\pgfqpoint{0.587435in}{1.138673in}}%
\pgfpathlineto{\pgfqpoint{0.596539in}{1.138540in}}%
\pgfpathlineto{\pgfqpoint{0.580348in}{1.155603in}}%
\pgfpathlineto{\pgfqpoint{0.564166in}{1.172877in}}%
\pgfpathlineto{\pgfqpoint{0.555019in}{1.172835in}}%
\pgfpathlineto{\pgfqpoint{0.545830in}{1.173062in}}%
\pgfpathclose%
\pgfusepath{fill}%
\end{pgfscope}%
\begin{pgfscope}%
\pgfpathrectangle{\pgfqpoint{0.407592in}{0.222904in}}{\pgfqpoint{1.411163in}{1.411163in}}%
\pgfusepath{clip}%
\pgfsetbuttcap%
\pgfsetroundjoin%
\definecolor{currentfill}{rgb}{0.909804,0.698039,0.423529}%
\pgfsetfillcolor{currentfill}%
\pgfsetfillopacity{0.120000}%
\pgfsetlinewidth{0.000000pt}%
\definecolor{currentstroke}{rgb}{0.000000,0.000000,0.000000}%
\pgfsetstrokecolor{currentstroke}%
\pgfsetdash{}{0pt}%
\pgfpathmoveto{\pgfqpoint{1.662735in}{1.204469in}}%
\pgfpathlineto{\pgfqpoint{1.680273in}{1.213035in}}%
\pgfpathlineto{\pgfqpoint{1.697887in}{1.221889in}}%
\pgfpathlineto{\pgfqpoint{1.705342in}{1.233356in}}%
\pgfpathlineto{\pgfqpoint{1.712787in}{1.245000in}}%
\pgfpathlineto{\pgfqpoint{1.695195in}{1.236405in}}%
\pgfpathlineto{\pgfqpoint{1.677679in}{1.228102in}}%
\pgfpathlineto{\pgfqpoint{1.670213in}{1.216196in}}%
\pgfpathlineto{\pgfqpoint{1.662735in}{1.204469in}}%
\pgfpathclose%
\pgfusepath{fill}%
\end{pgfscope}%
\begin{pgfscope}%
\pgfpathrectangle{\pgfqpoint{0.407592in}{0.222904in}}{\pgfqpoint{1.411163in}{1.411163in}}%
\pgfusepath{clip}%
\pgfsetbuttcap%
\pgfsetroundjoin%
\definecolor{currentfill}{rgb}{0.909804,0.698039,0.423529}%
\pgfsetfillcolor{currentfill}%
\pgfsetfillopacity{0.120000}%
\pgfsetlinewidth{0.000000pt}%
\definecolor{currentstroke}{rgb}{0.000000,0.000000,0.000000}%
\pgfsetstrokecolor{currentstroke}%
\pgfsetdash{}{0pt}%
\pgfpathmoveto{\pgfqpoint{1.023524in}{0.866521in}}%
\pgfpathlineto{\pgfqpoint{1.039922in}{0.862105in}}%
\pgfpathlineto{\pgfqpoint{1.056357in}{0.858685in}}%
\pgfpathlineto{\pgfqpoint{1.064620in}{0.860466in}}%
\pgfpathlineto{\pgfqpoint{1.072852in}{0.862781in}}%
\pgfpathlineto{\pgfqpoint{1.056454in}{0.866168in}}%
\pgfpathlineto{\pgfqpoint{1.040092in}{0.870570in}}%
\pgfpathlineto{\pgfqpoint{1.031824in}{0.868278in}}%
\pgfpathlineto{\pgfqpoint{1.023524in}{0.866521in}}%
\pgfpathclose%
\pgfusepath{fill}%
\end{pgfscope}%
\begin{pgfscope}%
\pgfpathrectangle{\pgfqpoint{0.407592in}{0.222904in}}{\pgfqpoint{1.411163in}{1.411163in}}%
\pgfusepath{clip}%
\pgfsetbuttcap%
\pgfsetroundjoin%
\definecolor{currentfill}{rgb}{0.909804,0.698039,0.423529}%
\pgfsetfillcolor{currentfill}%
\pgfsetfillopacity{0.120000}%
\pgfsetlinewidth{0.000000pt}%
\definecolor{currentstroke}{rgb}{0.000000,0.000000,0.000000}%
\pgfsetstrokecolor{currentstroke}%
\pgfsetdash{}{0pt}%
\pgfpathmoveto{\pgfqpoint{1.318989in}{0.933887in}}%
\pgfpathlineto{\pgfqpoint{1.335851in}{0.939682in}}%
\pgfpathlineto{\pgfqpoint{1.352782in}{0.946044in}}%
\pgfpathlineto{\pgfqpoint{1.360603in}{0.953869in}}%
\pgfpathlineto{\pgfqpoint{1.368408in}{0.962103in}}%
\pgfpathlineto{\pgfqpoint{1.351499in}{0.955898in}}%
\pgfpathlineto{\pgfqpoint{1.334657in}{0.950256in}}%
\pgfpathlineto{\pgfqpoint{1.326831in}{0.941857in}}%
\pgfpathlineto{\pgfqpoint{1.318989in}{0.933887in}}%
\pgfpathclose%
\pgfusepath{fill}%
\end{pgfscope}%
\begin{pgfscope}%
\pgfpathrectangle{\pgfqpoint{0.407592in}{0.222904in}}{\pgfqpoint{1.411163in}{1.411163in}}%
\pgfusepath{clip}%
\pgfsetbuttcap%
\pgfsetroundjoin%
\definecolor{currentfill}{rgb}{0.909804,0.698039,0.423529}%
\pgfsetfillcolor{currentfill}%
\pgfsetfillopacity{0.120000}%
\pgfsetlinewidth{0.000000pt}%
\definecolor{currentstroke}{rgb}{0.000000,0.000000,0.000000}%
\pgfsetstrokecolor{currentstroke}%
\pgfsetdash{}{0pt}%
\pgfpathmoveto{\pgfqpoint{0.693891in}{1.041708in}}%
\pgfpathlineto{\pgfqpoint{0.710153in}{1.026711in}}%
\pgfpathlineto{\pgfqpoint{0.726426in}{1.012114in}}%
\pgfpathlineto{\pgfqpoint{0.735317in}{1.011458in}}%
\pgfpathlineto{\pgfqpoint{0.744167in}{1.011138in}}%
\pgfpathlineto{\pgfqpoint{0.727937in}{1.025954in}}%
\pgfpathlineto{\pgfqpoint{0.711718in}{1.041165in}}%
\pgfpathlineto{\pgfqpoint{0.702825in}{1.041274in}}%
\pgfpathlineto{\pgfqpoint{0.693891in}{1.041708in}}%
\pgfpathclose%
\pgfusepath{fill}%
\end{pgfscope}%
\begin{pgfscope}%
\pgfpathrectangle{\pgfqpoint{0.407592in}{0.222904in}}{\pgfqpoint{1.411163in}{1.411163in}}%
\pgfusepath{clip}%
\pgfsetbuttcap%
\pgfsetroundjoin%
\definecolor{currentfill}{rgb}{0.909804,0.698039,0.423529}%
\pgfsetfillcolor{currentfill}%
\pgfsetfillopacity{0.120000}%
\pgfsetlinewidth{0.000000pt}%
\definecolor{currentstroke}{rgb}{0.000000,0.000000,0.000000}%
\pgfsetstrokecolor{currentstroke}%
\pgfsetdash{}{0pt}%
\pgfpathmoveto{\pgfqpoint{0.578289in}{1.139084in}}%
\pgfpathlineto{\pgfqpoint{0.594533in}{1.122442in}}%
\pgfpathlineto{\pgfqpoint{0.610786in}{1.106057in}}%
\pgfpathlineto{\pgfqpoint{0.619889in}{1.105446in}}%
\pgfpathlineto{\pgfqpoint{0.628950in}{1.105123in}}%
\pgfpathlineto{\pgfqpoint{0.612740in}{1.121707in}}%
\pgfpathlineto{\pgfqpoint{0.596539in}{1.138540in}}%
\pgfpathlineto{\pgfqpoint{0.587435in}{1.138673in}}%
\pgfpathlineto{\pgfqpoint{0.578289in}{1.139084in}}%
\pgfpathclose%
\pgfusepath{fill}%
\end{pgfscope}%
\begin{pgfscope}%
\pgfpathrectangle{\pgfqpoint{0.407592in}{0.222904in}}{\pgfqpoint{1.411163in}{1.411163in}}%
\pgfusepath{clip}%
\pgfsetbuttcap%
\pgfsetroundjoin%
\definecolor{currentfill}{rgb}{0.909804,0.698039,0.423529}%
\pgfsetfillcolor{currentfill}%
\pgfsetfillopacity{0.120000}%
\pgfsetlinewidth{0.000000pt}%
\definecolor{currentstroke}{rgb}{0.000000,0.000000,0.000000}%
\pgfsetstrokecolor{currentstroke}%
\pgfsetdash{}{0pt}%
\pgfpathmoveto{\pgfqpoint{1.483004in}{1.046464in}}%
\pgfpathlineto{\pgfqpoint{1.500204in}{1.054313in}}%
\pgfpathlineto{\pgfqpoint{1.517476in}{1.062539in}}%
\pgfpathlineto{\pgfqpoint{1.525117in}{1.072504in}}%
\pgfpathlineto{\pgfqpoint{1.532746in}{1.082745in}}%
\pgfpathlineto{\pgfqpoint{1.515494in}{1.074769in}}%
\pgfpathlineto{\pgfqpoint{1.498315in}{1.067173in}}%
\pgfpathlineto{\pgfqpoint{1.490666in}{1.056677in}}%
\pgfpathlineto{\pgfqpoint{1.483004in}{1.046464in}}%
\pgfpathclose%
\pgfusepath{fill}%
\end{pgfscope}%
\begin{pgfscope}%
\pgfpathrectangle{\pgfqpoint{0.407592in}{0.222904in}}{\pgfqpoint{1.411163in}{1.411163in}}%
\pgfusepath{clip}%
\pgfsetbuttcap%
\pgfsetroundjoin%
\definecolor{currentfill}{rgb}{0.909804,0.698039,0.423529}%
\pgfsetfillcolor{currentfill}%
\pgfsetfillopacity{0.120000}%
\pgfsetlinewidth{0.000000pt}%
\definecolor{currentstroke}{rgb}{0.000000,0.000000,0.000000}%
\pgfsetstrokecolor{currentstroke}%
\pgfsetdash{}{0pt}%
\pgfpathmoveto{\pgfqpoint{0.891716in}{0.910071in}}%
\pgfpathlineto{\pgfqpoint{0.908036in}{0.900351in}}%
\pgfpathlineto{\pgfqpoint{0.924376in}{0.891435in}}%
\pgfpathlineto{\pgfqpoint{0.932895in}{0.891358in}}%
\pgfpathlineto{\pgfqpoint{0.941378in}{0.891731in}}%
\pgfpathlineto{\pgfqpoint{0.925080in}{0.900766in}}%
\pgfpathlineto{\pgfqpoint{0.908802in}{0.910621in}}%
\pgfpathlineto{\pgfqpoint{0.900278in}{0.910127in}}%
\pgfpathlineto{\pgfqpoint{0.891716in}{0.910071in}}%
\pgfpathclose%
\pgfusepath{fill}%
\end{pgfscope}%
\begin{pgfscope}%
\pgfpathrectangle{\pgfqpoint{0.407592in}{0.222904in}}{\pgfqpoint{1.411163in}{1.411163in}}%
\pgfusepath{clip}%
\pgfsetbuttcap%
\pgfsetroundjoin%
\definecolor{currentfill}{rgb}{0.909804,0.698039,0.423529}%
\pgfsetfillcolor{currentfill}%
\pgfsetfillopacity{0.120000}%
\pgfsetlinewidth{0.000000pt}%
\definecolor{currentstroke}{rgb}{0.000000,0.000000,0.000000}%
\pgfsetstrokecolor{currentstroke}%
\pgfsetdash{}{0pt}%
\pgfpathmoveto{\pgfqpoint{1.597773in}{1.142228in}}%
\pgfpathlineto{\pgfqpoint{1.615204in}{1.150708in}}%
\pgfpathlineto{\pgfqpoint{1.632710in}{1.159498in}}%
\pgfpathlineto{\pgfqpoint{1.640234in}{1.170433in}}%
\pgfpathlineto{\pgfqpoint{1.647746in}{1.181579in}}%
\pgfpathlineto{\pgfqpoint{1.630262in}{1.173051in}}%
\pgfpathlineto{\pgfqpoint{1.612852in}{1.164837in}}%
\pgfpathlineto{\pgfqpoint{1.605318in}{1.153426in}}%
\pgfpathlineto{\pgfqpoint{1.597773in}{1.142228in}}%
\pgfpathclose%
\pgfusepath{fill}%
\end{pgfscope}%
\begin{pgfscope}%
\pgfpathrectangle{\pgfqpoint{0.407592in}{0.222904in}}{\pgfqpoint{1.411163in}{1.411163in}}%
\pgfusepath{clip}%
\pgfsetbuttcap%
\pgfsetroundjoin%
\definecolor{currentfill}{rgb}{0.909804,0.698039,0.423529}%
\pgfsetfillcolor{currentfill}%
\pgfsetfillopacity{0.120000}%
\pgfsetlinewidth{0.000000pt}%
\definecolor{currentstroke}{rgb}{0.000000,0.000000,0.000000}%
\pgfsetstrokecolor{currentstroke}%
\pgfsetdash{}{0pt}%
\pgfpathmoveto{\pgfqpoint{0.809207in}{0.956520in}}%
\pgfpathlineto{\pgfqpoint{0.825501in}{0.944217in}}%
\pgfpathlineto{\pgfqpoint{0.841810in}{0.932535in}}%
\pgfpathlineto{\pgfqpoint{0.850491in}{0.931947in}}%
\pgfpathlineto{\pgfqpoint{0.859133in}{0.931753in}}%
\pgfpathlineto{\pgfqpoint{0.842867in}{0.943626in}}%
\pgfpathlineto{\pgfqpoint{0.826616in}{0.956127in}}%
\pgfpathlineto{\pgfqpoint{0.817931in}{0.956132in}}%
\pgfpathlineto{\pgfqpoint{0.809207in}{0.956520in}}%
\pgfpathclose%
\pgfusepath{fill}%
\end{pgfscope}%
\begin{pgfscope}%
\pgfpathrectangle{\pgfqpoint{0.407592in}{0.222904in}}{\pgfqpoint{1.411163in}{1.411163in}}%
\pgfusepath{clip}%
\pgfsetbuttcap%
\pgfsetroundjoin%
\definecolor{currentfill}{rgb}{0.909804,0.698039,0.423529}%
\pgfsetfillcolor{currentfill}%
\pgfsetfillopacity{0.120000}%
\pgfsetlinewidth{0.000000pt}%
\definecolor{currentstroke}{rgb}{0.000000,0.000000,0.000000}%
\pgfsetstrokecolor{currentstroke}%
\pgfsetdash{}{0pt}%
\pgfpathmoveto{\pgfqpoint{1.368408in}{0.962103in}}%
\pgfpathlineto{\pgfqpoint{1.385386in}{0.968832in}}%
\pgfpathlineto{\pgfqpoint{1.402433in}{0.976048in}}%
\pgfpathlineto{\pgfqpoint{1.410201in}{0.984495in}}%
\pgfpathlineto{\pgfqpoint{1.417954in}{0.993309in}}%
\pgfpathlineto{\pgfqpoint{1.400927in}{0.986284in}}%
\pgfpathlineto{\pgfqpoint{1.383970in}{0.979747in}}%
\pgfpathlineto{\pgfqpoint{1.376197in}{0.970733in}}%
\pgfpathlineto{\pgfqpoint{1.368408in}{0.962103in}}%
\pgfpathclose%
\pgfusepath{fill}%
\end{pgfscope}%
\begin{pgfscope}%
\pgfpathrectangle{\pgfqpoint{0.407592in}{0.222904in}}{\pgfqpoint{1.411163in}{1.411163in}}%
\pgfusepath{clip}%
\pgfsetbuttcap%
\pgfsetroundjoin%
\definecolor{currentfill}{rgb}{0.909804,0.698039,0.423529}%
\pgfsetfillcolor{currentfill}%
\pgfsetfillopacity{0.120000}%
\pgfsetlinewidth{0.000000pt}%
\definecolor{currentstroke}{rgb}{0.000000,0.000000,0.000000}%
\pgfsetstrokecolor{currentstroke}%
\pgfsetdash{}{0pt}%
\pgfpathmoveto{\pgfqpoint{0.974046in}{0.876284in}}%
\pgfpathlineto{\pgfqpoint{0.990421in}{0.869938in}}%
\pgfpathlineto{\pgfqpoint{1.006825in}{0.864545in}}%
\pgfpathlineto{\pgfqpoint{1.015191in}{0.865283in}}%
\pgfpathlineto{\pgfqpoint{1.023524in}{0.866521in}}%
\pgfpathlineto{\pgfqpoint{1.007159in}{0.871928in}}%
\pgfpathlineto{\pgfqpoint{0.990824in}{0.878311in}}%
\pgfpathlineto{\pgfqpoint{0.982453in}{0.877050in}}%
\pgfpathlineto{\pgfqpoint{0.974046in}{0.876284in}}%
\pgfpathclose%
\pgfusepath{fill}%
\end{pgfscope}%
\begin{pgfscope}%
\pgfpathrectangle{\pgfqpoint{0.407592in}{0.222904in}}{\pgfqpoint{1.411163in}{1.411163in}}%
\pgfusepath{clip}%
\pgfsetbuttcap%
\pgfsetroundjoin%
\definecolor{currentfill}{rgb}{0.909804,0.698039,0.423529}%
\pgfsetfillcolor{currentfill}%
\pgfsetfillopacity{0.120000}%
\pgfsetlinewidth{0.000000pt}%
\definecolor{currentstroke}{rgb}{0.000000,0.000000,0.000000}%
\pgfsetstrokecolor{currentstroke}%
\pgfsetdash{}{0pt}%
\pgfpathmoveto{\pgfqpoint{0.610786in}{1.106057in}}%
\pgfpathlineto{\pgfqpoint{0.627049in}{1.089948in}}%
\pgfpathlineto{\pgfqpoint{0.643322in}{1.074138in}}%
\pgfpathlineto{\pgfqpoint{0.652382in}{1.073312in}}%
\pgfpathlineto{\pgfqpoint{0.661400in}{1.072783in}}%
\pgfpathlineto{\pgfqpoint{0.645170in}{1.088807in}}%
\pgfpathlineto{\pgfqpoint{0.628950in}{1.105123in}}%
\pgfpathlineto{\pgfqpoint{0.619889in}{1.105446in}}%
\pgfpathlineto{\pgfqpoint{0.610786in}{1.106057in}}%
\pgfpathclose%
\pgfusepath{fill}%
\end{pgfscope}%
\begin{pgfscope}%
\pgfpathrectangle{\pgfqpoint{0.407592in}{0.222904in}}{\pgfqpoint{1.411163in}{1.411163in}}%
\pgfusepath{clip}%
\pgfsetbuttcap%
\pgfsetroundjoin%
\definecolor{currentfill}{rgb}{0.909804,0.698039,0.423529}%
\pgfsetfillcolor{currentfill}%
\pgfsetfillopacity{0.120000}%
\pgfsetlinewidth{0.000000pt}%
\definecolor{currentstroke}{rgb}{0.000000,0.000000,0.000000}%
\pgfsetstrokecolor{currentstroke}%
\pgfsetdash{}{0pt}%
\pgfpathmoveto{\pgfqpoint{0.726426in}{1.012114in}}%
\pgfpathlineto{\pgfqpoint{0.742710in}{0.997948in}}%
\pgfpathlineto{\pgfqpoint{0.759007in}{0.984248in}}%
\pgfpathlineto{\pgfqpoint{0.767854in}{0.983365in}}%
\pgfpathlineto{\pgfqpoint{0.776661in}{0.982827in}}%
\pgfpathlineto{\pgfqpoint{0.760408in}{0.996751in}}%
\pgfpathlineto{\pgfqpoint{0.744167in}{1.011138in}}%
\pgfpathlineto{\pgfqpoint{0.735317in}{1.011458in}}%
\pgfpathlineto{\pgfqpoint{0.726426in}{1.012114in}}%
\pgfpathclose%
\pgfusepath{fill}%
\end{pgfscope}%
\begin{pgfscope}%
\pgfpathrectangle{\pgfqpoint{0.407592in}{0.222904in}}{\pgfqpoint{1.411163in}{1.411163in}}%
\pgfusepath{clip}%
\pgfsetbuttcap%
\pgfsetroundjoin%
\definecolor{currentfill}{rgb}{0.909804,0.698039,0.423529}%
\pgfsetfillcolor{currentfill}%
\pgfsetfillopacity{0.120000}%
\pgfsetlinewidth{0.000000pt}%
\definecolor{currentstroke}{rgb}{0.000000,0.000000,0.000000}%
\pgfsetstrokecolor{currentstroke}%
\pgfsetdash{}{0pt}%
\pgfpathmoveto{\pgfqpoint{1.155119in}{0.865670in}}%
\pgfpathlineto{\pgfqpoint{1.171716in}{0.867182in}}%
\pgfpathlineto{\pgfqpoint{1.188366in}{0.869556in}}%
\pgfpathlineto{\pgfqpoint{1.196422in}{0.873682in}}%
\pgfpathlineto{\pgfqpoint{1.204453in}{0.878338in}}%
\pgfpathlineto{\pgfqpoint{1.187830in}{0.875922in}}%
\pgfpathlineto{\pgfqpoint{1.171262in}{0.874370in}}%
\pgfpathlineto{\pgfqpoint{1.163203in}{0.869745in}}%
\pgfpathlineto{\pgfqpoint{1.155119in}{0.865670in}}%
\pgfpathclose%
\pgfusepath{fill}%
\end{pgfscope}%
\begin{pgfscope}%
\pgfpathrectangle{\pgfqpoint{0.407592in}{0.222904in}}{\pgfqpoint{1.411163in}{1.411163in}}%
\pgfusepath{clip}%
\pgfsetbuttcap%
\pgfsetroundjoin%
\definecolor{currentfill}{rgb}{0.909804,0.698039,0.423529}%
\pgfsetfillcolor{currentfill}%
\pgfsetfillopacity{0.120000}%
\pgfsetlinewidth{0.000000pt}%
\definecolor{currentstroke}{rgb}{0.000000,0.000000,0.000000}%
\pgfsetstrokecolor{currentstroke}%
\pgfsetdash{}{0pt}%
\pgfpathmoveto{\pgfqpoint{1.532746in}{1.082745in}}%
\pgfpathlineto{\pgfqpoint{1.550071in}{1.091079in}}%
\pgfpathlineto{\pgfqpoint{1.567470in}{1.099749in}}%
\pgfpathlineto{\pgfqpoint{1.575064in}{1.110003in}}%
\pgfpathlineto{\pgfqpoint{1.582646in}{1.120508in}}%
\pgfpathlineto{\pgfqpoint{1.565268in}{1.112091in}}%
\pgfpathlineto{\pgfqpoint{1.547964in}{1.104016in}}%
\pgfpathlineto{\pgfqpoint{1.540361in}{1.093253in}}%
\pgfpathlineto{\pgfqpoint{1.532746in}{1.082745in}}%
\pgfpathclose%
\pgfusepath{fill}%
\end{pgfscope}%
\begin{pgfscope}%
\pgfpathrectangle{\pgfqpoint{0.407592in}{0.222904in}}{\pgfqpoint{1.411163in}{1.411163in}}%
\pgfusepath{clip}%
\pgfsetbuttcap%
\pgfsetroundjoin%
\definecolor{currentfill}{rgb}{0.909804,0.698039,0.423529}%
\pgfsetfillcolor{currentfill}%
\pgfsetfillopacity{0.120000}%
\pgfsetlinewidth{0.000000pt}%
\definecolor{currentstroke}{rgb}{0.000000,0.000000,0.000000}%
\pgfsetstrokecolor{currentstroke}%
\pgfsetdash{}{0pt}%
\pgfpathmoveto{\pgfqpoint{1.712787in}{1.245000in}}%
\pgfpathlineto{\pgfqpoint{1.730454in}{1.253875in}}%
\pgfpathlineto{\pgfqpoint{1.748197in}{1.263018in}}%
\pgfpathlineto{\pgfqpoint{1.755608in}{1.274579in}}%
\pgfpathlineto{\pgfqpoint{1.737876in}{1.265562in}}%
\pgfpathlineto{\pgfqpoint{1.720220in}{1.256815in}}%
\pgfusepath{fill}%
\end{pgfscope}%
\begin{pgfscope}%
\pgfpathrectangle{\pgfqpoint{0.407592in}{0.222904in}}{\pgfqpoint{1.411163in}{1.411163in}}%
\pgfusepath{clip}%
\pgfsetbuttcap%
\pgfsetroundjoin%
\definecolor{currentfill}{rgb}{0.909804,0.698039,0.423529}%
\pgfsetfillcolor{currentfill}%
\pgfsetfillopacity{0.120000}%
\pgfsetlinewidth{0.000000pt}%
\definecolor{currentstroke}{rgb}{0.000000,0.000000,0.000000}%
\pgfsetstrokecolor{currentstroke}%
\pgfsetdash{}{0pt}%
\pgfpathmoveto{\pgfqpoint{1.204453in}{0.878338in}}%
\pgfpathlineto{\pgfqpoint{1.221133in}{0.881569in}}%
\pgfpathlineto{\pgfqpoint{1.237874in}{0.885564in}}%
\pgfpathlineto{\pgfqpoint{1.245856in}{0.890757in}}%
\pgfpathlineto{\pgfqpoint{1.253817in}{0.896454in}}%
\pgfpathlineto{\pgfqpoint{1.237101in}{0.892466in}}%
\pgfpathlineto{\pgfqpoint{1.220446in}{0.889240in}}%
\pgfpathlineto{\pgfqpoint{1.212460in}{0.883525in}}%
\pgfpathlineto{\pgfqpoint{1.204453in}{0.878338in}}%
\pgfpathclose%
\pgfusepath{fill}%
\end{pgfscope}%
\begin{pgfscope}%
\pgfpathrectangle{\pgfqpoint{0.407592in}{0.222904in}}{\pgfqpoint{1.411163in}{1.411163in}}%
\pgfusepath{clip}%
\pgfsetbuttcap%
\pgfsetroundjoin%
\definecolor{currentfill}{rgb}{0.909804,0.698039,0.423529}%
\pgfsetfillcolor{currentfill}%
\pgfsetfillopacity{0.120000}%
\pgfsetlinewidth{0.000000pt}%
\definecolor{currentstroke}{rgb}{0.000000,0.000000,0.000000}%
\pgfsetstrokecolor{currentstroke}%
\pgfsetdash{}{0pt}%
\pgfpathmoveto{\pgfqpoint{1.417954in}{0.993309in}}%
\pgfpathlineto{\pgfqpoint{1.435051in}{1.000789in}}%
\pgfpathlineto{\pgfqpoint{1.452220in}{1.008691in}}%
\pgfpathlineto{\pgfqpoint{1.459937in}{1.017650in}}%
\pgfpathlineto{\pgfqpoint{1.467640in}{1.026940in}}%
\pgfpathlineto{\pgfqpoint{1.450492in}{1.019253in}}%
\pgfpathlineto{\pgfqpoint{1.433415in}{1.011992in}}%
\pgfpathlineto{\pgfqpoint{1.425692in}{1.002479in}}%
\pgfpathlineto{\pgfqpoint{1.417954in}{0.993309in}}%
\pgfpathclose%
\pgfusepath{fill}%
\end{pgfscope}%
\begin{pgfscope}%
\pgfpathrectangle{\pgfqpoint{0.407592in}{0.222904in}}{\pgfqpoint{1.411163in}{1.411163in}}%
\pgfusepath{clip}%
\pgfsetbuttcap%
\pgfsetroundjoin%
\definecolor{currentfill}{rgb}{0.909804,0.698039,0.423529}%
\pgfsetfillcolor{currentfill}%
\pgfsetfillopacity{0.120000}%
\pgfsetlinewidth{0.000000pt}%
\definecolor{currentstroke}{rgb}{0.000000,0.000000,0.000000}%
\pgfsetstrokecolor{currentstroke}%
\pgfsetdash{}{0pt}%
\pgfpathmoveto{\pgfqpoint{1.105771in}{0.859027in}}%
\pgfpathlineto{\pgfqpoint{1.122298in}{0.858620in}}%
\pgfpathlineto{\pgfqpoint{1.138873in}{0.859153in}}%
\pgfpathlineto{\pgfqpoint{1.147010in}{0.862142in}}%
\pgfpathlineto{\pgfqpoint{1.155119in}{0.865670in}}%
\pgfpathlineto{\pgfqpoint{1.138575in}{0.865066in}}%
\pgfpathlineto{\pgfqpoint{1.122080in}{0.865411in}}%
\pgfpathlineto{\pgfqpoint{1.113940in}{0.861942in}}%
\pgfpathlineto{\pgfqpoint{1.105771in}{0.859027in}}%
\pgfpathclose%
\pgfusepath{fill}%
\end{pgfscope}%
\begin{pgfscope}%
\pgfpathrectangle{\pgfqpoint{0.407592in}{0.222904in}}{\pgfqpoint{1.411163in}{1.411163in}}%
\pgfusepath{clip}%
\pgfsetbuttcap%
\pgfsetroundjoin%
\definecolor{currentfill}{rgb}{0.909804,0.698039,0.423529}%
\pgfsetfillcolor{currentfill}%
\pgfsetfillopacity{0.120000}%
\pgfsetlinewidth{0.000000pt}%
\definecolor{currentstroke}{rgb}{0.000000,0.000000,0.000000}%
\pgfsetstrokecolor{currentstroke}%
\pgfsetdash{}{0pt}%
\pgfpathmoveto{\pgfqpoint{0.494818in}{1.208735in}}%
\pgfpathlineto{\pgfqpoint{0.511067in}{1.191420in}}%
\pgfpathlineto{\pgfqpoint{0.527325in}{1.174301in}}%
\pgfpathlineto{\pgfqpoint{0.536599in}{1.173553in}}%
\pgfpathlineto{\pgfqpoint{0.545830in}{1.173062in}}%
\pgfpathlineto{\pgfqpoint{0.529615in}{1.190365in}}%
\pgfpathlineto{\pgfqpoint{0.513409in}{1.207856in}}%
\pgfpathlineto{\pgfqpoint{0.504135in}{1.208171in}}%
\pgfpathlineto{\pgfqpoint{0.494818in}{1.208735in}}%
\pgfpathclose%
\pgfusepath{fill}%
\end{pgfscope}%
\begin{pgfscope}%
\pgfpathrectangle{\pgfqpoint{0.407592in}{0.222904in}}{\pgfqpoint{1.411163in}{1.411163in}}%
\pgfusepath{clip}%
\pgfsetbuttcap%
\pgfsetroundjoin%
\definecolor{currentfill}{rgb}{0.909804,0.698039,0.423529}%
\pgfsetfillcolor{currentfill}%
\pgfsetfillopacity{0.120000}%
\pgfsetlinewidth{0.000000pt}%
\definecolor{currentstroke}{rgb}{0.000000,0.000000,0.000000}%
\pgfsetstrokecolor{currentstroke}%
\pgfsetdash{}{0pt}%
\pgfpathmoveto{\pgfqpoint{1.647746in}{1.181579in}}%
\pgfpathlineto{\pgfqpoint{1.665305in}{1.190406in}}%
\pgfpathlineto{\pgfqpoint{1.682941in}{1.199517in}}%
\pgfpathlineto{\pgfqpoint{1.690419in}{1.210607in}}%
\pgfpathlineto{\pgfqpoint{1.697887in}{1.221889in}}%
\pgfpathlineto{\pgfqpoint{1.680273in}{1.213035in}}%
\pgfpathlineto{\pgfqpoint{1.662735in}{1.204469in}}%
\pgfpathlineto{\pgfqpoint{1.655247in}{1.192927in}}%
\pgfpathlineto{\pgfqpoint{1.647746in}{1.181579in}}%
\pgfpathclose%
\pgfusepath{fill}%
\end{pgfscope}%
\begin{pgfscope}%
\pgfpathrectangle{\pgfqpoint{0.407592in}{0.222904in}}{\pgfqpoint{1.411163in}{1.411163in}}%
\pgfusepath{clip}%
\pgfsetbuttcap%
\pgfsetroundjoin%
\definecolor{currentfill}{rgb}{0.909804,0.698039,0.423529}%
\pgfsetfillcolor{currentfill}%
\pgfsetfillopacity{0.120000}%
\pgfsetlinewidth{0.000000pt}%
\definecolor{currentstroke}{rgb}{0.000000,0.000000,0.000000}%
\pgfsetstrokecolor{currentstroke}%
\pgfsetdash{}{0pt}%
\pgfpathmoveto{\pgfqpoint{1.253817in}{0.896454in}}%
\pgfpathlineto{\pgfqpoint{1.270595in}{0.901152in}}%
\pgfpathlineto{\pgfqpoint{1.287437in}{0.906512in}}%
\pgfpathlineto{\pgfqpoint{1.295353in}{0.912660in}}%
\pgfpathlineto{\pgfqpoint{1.303250in}{0.919278in}}%
\pgfpathlineto{\pgfqpoint{1.286431in}{0.913981in}}%
\pgfpathlineto{\pgfqpoint{1.269677in}{0.909342in}}%
\pgfpathlineto{\pgfqpoint{1.261757in}{0.902651in}}%
\pgfpathlineto{\pgfqpoint{1.253817in}{0.896454in}}%
\pgfpathclose%
\pgfusepath{fill}%
\end{pgfscope}%
\begin{pgfscope}%
\pgfpathrectangle{\pgfqpoint{0.407592in}{0.222904in}}{\pgfqpoint{1.411163in}{1.411163in}}%
\pgfusepath{clip}%
\pgfsetbuttcap%
\pgfsetroundjoin%
\definecolor{currentfill}{rgb}{0.909804,0.698039,0.423529}%
\pgfsetfillcolor{currentfill}%
\pgfsetfillopacity{0.120000}%
\pgfsetlinewidth{0.000000pt}%
\definecolor{currentstroke}{rgb}{0.000000,0.000000,0.000000}%
\pgfsetstrokecolor{currentstroke}%
\pgfsetdash{}{0pt}%
\pgfpathmoveto{\pgfqpoint{0.841810in}{0.932535in}}%
\pgfpathlineto{\pgfqpoint{0.858136in}{0.921518in}}%
\pgfpathlineto{\pgfqpoint{0.874480in}{0.911207in}}%
\pgfpathlineto{\pgfqpoint{0.883117in}{0.910437in}}%
\pgfpathlineto{\pgfqpoint{0.891716in}{0.910071in}}%
\pgfpathlineto{\pgfqpoint{0.875416in}{0.920553in}}%
\pgfpathlineto{\pgfqpoint{0.859133in}{0.931753in}}%
\pgfpathlineto{\pgfqpoint{0.850491in}{0.931947in}}%
\pgfpathlineto{\pgfqpoint{0.841810in}{0.932535in}}%
\pgfpathclose%
\pgfusepath{fill}%
\end{pgfscope}%
\begin{pgfscope}%
\pgfpathrectangle{\pgfqpoint{0.407592in}{0.222904in}}{\pgfqpoint{1.411163in}{1.411163in}}%
\pgfusepath{clip}%
\pgfsetbuttcap%
\pgfsetroundjoin%
\definecolor{currentfill}{rgb}{0.909804,0.698039,0.423529}%
\pgfsetfillcolor{currentfill}%
\pgfsetfillopacity{0.120000}%
\pgfsetlinewidth{0.000000pt}%
\definecolor{currentstroke}{rgb}{0.000000,0.000000,0.000000}%
\pgfsetstrokecolor{currentstroke}%
\pgfsetdash{}{0pt}%
\pgfpathmoveto{\pgfqpoint{0.924376in}{0.891435in}}%
\pgfpathlineto{\pgfqpoint{0.940740in}{0.883360in}}%
\pgfpathlineto{\pgfqpoint{0.957128in}{0.876160in}}%
\pgfpathlineto{\pgfqpoint{0.965605in}{0.875994in}}%
\pgfpathlineto{\pgfqpoint{0.974046in}{0.876284in}}%
\pgfpathlineto{\pgfqpoint{0.957700in}{0.883558in}}%
\pgfpathlineto{\pgfqpoint{0.941378in}{0.891731in}}%
\pgfpathlineto{\pgfqpoint{0.932895in}{0.891358in}}%
\pgfpathlineto{\pgfqpoint{0.924376in}{0.891435in}}%
\pgfpathclose%
\pgfusepath{fill}%
\end{pgfscope}%
\begin{pgfscope}%
\pgfpathrectangle{\pgfqpoint{0.407592in}{0.222904in}}{\pgfqpoint{1.411163in}{1.411163in}}%
\pgfusepath{clip}%
\pgfsetbuttcap%
\pgfsetroundjoin%
\definecolor{currentfill}{rgb}{0.909804,0.698039,0.423529}%
\pgfsetfillcolor{currentfill}%
\pgfsetfillopacity{0.120000}%
\pgfsetlinewidth{0.000000pt}%
\definecolor{currentstroke}{rgb}{0.000000,0.000000,0.000000}%
\pgfsetstrokecolor{currentstroke}%
\pgfsetdash{}{0pt}%
\pgfpathmoveto{\pgfqpoint{1.056357in}{0.858685in}}%
\pgfpathlineto{\pgfqpoint{1.072831in}{0.856258in}}%
\pgfpathlineto{\pgfqpoint{1.089346in}{0.854809in}}%
\pgfpathlineto{\pgfqpoint{1.097574in}{0.856654in}}%
\pgfpathlineto{\pgfqpoint{1.105771in}{0.859027in}}%
\pgfpathlineto{\pgfqpoint{1.089290in}{0.860406in}}%
\pgfpathlineto{\pgfqpoint{1.072852in}{0.862781in}}%
\pgfpathlineto{\pgfqpoint{1.064620in}{0.860466in}}%
\pgfpathlineto{\pgfqpoint{1.056357in}{0.858685in}}%
\pgfpathclose%
\pgfusepath{fill}%
\end{pgfscope}%
\begin{pgfscope}%
\pgfpathrectangle{\pgfqpoint{0.407592in}{0.222904in}}{\pgfqpoint{1.411163in}{1.411163in}}%
\pgfusepath{clip}%
\pgfsetbuttcap%
\pgfsetroundjoin%
\definecolor{currentfill}{rgb}{0.909804,0.698039,0.423529}%
\pgfsetfillcolor{currentfill}%
\pgfsetfillopacity{0.120000}%
\pgfsetlinewidth{0.000000pt}%
\definecolor{currentstroke}{rgb}{0.000000,0.000000,0.000000}%
\pgfsetstrokecolor{currentstroke}%
\pgfsetdash{}{0pt}%
\pgfpathmoveto{\pgfqpoint{0.643322in}{1.074138in}}%
\pgfpathlineto{\pgfqpoint{0.659606in}{1.058651in}}%
\pgfpathlineto{\pgfqpoint{0.675900in}{1.043512in}}%
\pgfpathlineto{\pgfqpoint{0.684916in}{1.042457in}}%
\pgfpathlineto{\pgfqpoint{0.693891in}{1.041708in}}%
\pgfpathlineto{\pgfqpoint{0.677641in}{1.057074in}}%
\pgfpathlineto{\pgfqpoint{0.661400in}{1.072783in}}%
\pgfpathlineto{\pgfqpoint{0.652382in}{1.073312in}}%
\pgfpathlineto{\pgfqpoint{0.643322in}{1.074138in}}%
\pgfpathclose%
\pgfusepath{fill}%
\end{pgfscope}%
\begin{pgfscope}%
\pgfpathrectangle{\pgfqpoint{0.407592in}{0.222904in}}{\pgfqpoint{1.411163in}{1.411163in}}%
\pgfusepath{clip}%
\pgfsetbuttcap%
\pgfsetroundjoin%
\definecolor{currentfill}{rgb}{0.909804,0.698039,0.423529}%
\pgfsetfillcolor{currentfill}%
\pgfsetfillopacity{0.120000}%
\pgfsetlinewidth{0.000000pt}%
\definecolor{currentstroke}{rgb}{0.000000,0.000000,0.000000}%
\pgfsetstrokecolor{currentstroke}%
\pgfsetdash{}{0pt}%
\pgfpathmoveto{\pgfqpoint{1.303250in}{0.919278in}}%
\pgfpathlineto{\pgfqpoint{1.320135in}{0.925186in}}%
\pgfpathlineto{\pgfqpoint{1.337087in}{0.931661in}}%
\pgfpathlineto{\pgfqpoint{1.344943in}{0.938638in}}%
\pgfpathlineto{\pgfqpoint{1.352782in}{0.946044in}}%
\pgfpathlineto{\pgfqpoint{1.335851in}{0.939682in}}%
\pgfpathlineto{\pgfqpoint{1.318989in}{0.933887in}}%
\pgfpathlineto{\pgfqpoint{1.311128in}{0.926356in}}%
\pgfpathlineto{\pgfqpoint{1.303250in}{0.919278in}}%
\pgfpathclose%
\pgfusepath{fill}%
\end{pgfscope}%
\begin{pgfscope}%
\pgfpathrectangle{\pgfqpoint{0.407592in}{0.222904in}}{\pgfqpoint{1.411163in}{1.411163in}}%
\pgfusepath{clip}%
\pgfsetbuttcap%
\pgfsetroundjoin%
\definecolor{currentfill}{rgb}{0.909804,0.698039,0.423529}%
\pgfsetfillcolor{currentfill}%
\pgfsetfillopacity{0.120000}%
\pgfsetlinewidth{0.000000pt}%
\definecolor{currentstroke}{rgb}{0.000000,0.000000,0.000000}%
\pgfsetstrokecolor{currentstroke}%
\pgfsetdash{}{0pt}%
\pgfpathmoveto{\pgfqpoint{0.759007in}{0.984248in}}%
\pgfpathlineto{\pgfqpoint{0.775317in}{0.971049in}}%
\pgfpathlineto{\pgfqpoint{0.791640in}{0.958388in}}%
\pgfpathlineto{\pgfqpoint{0.800444in}{0.957276in}}%
\pgfpathlineto{\pgfqpoint{0.809207in}{0.956520in}}%
\pgfpathlineto{\pgfqpoint{0.792928in}{0.969403in}}%
\pgfpathlineto{\pgfqpoint{0.776661in}{0.982827in}}%
\pgfpathlineto{\pgfqpoint{0.767854in}{0.983365in}}%
\pgfpathlineto{\pgfqpoint{0.759007in}{0.984248in}}%
\pgfpathclose%
\pgfusepath{fill}%
\end{pgfscope}%
\begin{pgfscope}%
\pgfpathrectangle{\pgfqpoint{0.407592in}{0.222904in}}{\pgfqpoint{1.411163in}{1.411163in}}%
\pgfusepath{clip}%
\pgfsetbuttcap%
\pgfsetroundjoin%
\definecolor{currentfill}{rgb}{0.909804,0.698039,0.423529}%
\pgfsetfillcolor{currentfill}%
\pgfsetfillopacity{0.120000}%
\pgfsetlinewidth{0.000000pt}%
\definecolor{currentstroke}{rgb}{0.000000,0.000000,0.000000}%
\pgfsetstrokecolor{currentstroke}%
\pgfsetdash{}{0pt}%
\pgfpathmoveto{\pgfqpoint{1.467640in}{1.026940in}}%
\pgfpathlineto{\pgfqpoint{1.484860in}{1.035024in}}%
\pgfpathlineto{\pgfqpoint{1.502153in}{1.043481in}}%
\pgfpathlineto{\pgfqpoint{1.509822in}{1.052861in}}%
\pgfpathlineto{\pgfqpoint{1.517476in}{1.062539in}}%
\pgfpathlineto{\pgfqpoint{1.500204in}{1.054313in}}%
\pgfpathlineto{\pgfqpoint{1.483004in}{1.046464in}}%
\pgfpathlineto{\pgfqpoint{1.475329in}{1.036548in}}%
\pgfpathlineto{\pgfqpoint{1.467640in}{1.026940in}}%
\pgfpathclose%
\pgfusepath{fill}%
\end{pgfscope}%
\begin{pgfscope}%
\pgfpathrectangle{\pgfqpoint{0.407592in}{0.222904in}}{\pgfqpoint{1.411163in}{1.411163in}}%
\pgfusepath{clip}%
\pgfsetbuttcap%
\pgfsetroundjoin%
\definecolor{currentfill}{rgb}{0.909804,0.698039,0.423529}%
\pgfsetfillcolor{currentfill}%
\pgfsetfillopacity{0.120000}%
\pgfsetlinewidth{0.000000pt}%
\definecolor{currentstroke}{rgb}{0.000000,0.000000,0.000000}%
\pgfsetstrokecolor{currentstroke}%
\pgfsetdash{}{0pt}%
\pgfpathmoveto{\pgfqpoint{0.527325in}{1.174301in}}%
\pgfpathlineto{\pgfqpoint{0.543593in}{1.157394in}}%
\pgfpathlineto{\pgfqpoint{0.559870in}{1.140715in}}%
\pgfpathlineto{\pgfqpoint{0.569101in}{1.139767in}}%
\pgfpathlineto{\pgfqpoint{0.578289in}{1.139084in}}%
\pgfpathlineto{\pgfqpoint{0.562055in}{1.155963in}}%
\pgfpathlineto{\pgfqpoint{0.545830in}{1.173062in}}%
\pgfpathlineto{\pgfqpoint{0.536599in}{1.173553in}}%
\pgfpathlineto{\pgfqpoint{0.527325in}{1.174301in}}%
\pgfpathclose%
\pgfusepath{fill}%
\end{pgfscope}%
\begin{pgfscope}%
\pgfpathrectangle{\pgfqpoint{0.407592in}{0.222904in}}{\pgfqpoint{1.411163in}{1.411163in}}%
\pgfusepath{clip}%
\pgfsetbuttcap%
\pgfsetroundjoin%
\definecolor{currentfill}{rgb}{0.909804,0.698039,0.423529}%
\pgfsetfillcolor{currentfill}%
\pgfsetfillopacity{0.120000}%
\pgfsetlinewidth{0.000000pt}%
\definecolor{currentstroke}{rgb}{0.000000,0.000000,0.000000}%
\pgfsetstrokecolor{currentstroke}%
\pgfsetdash{}{0pt}%
\pgfpathmoveto{\pgfqpoint{1.582646in}{1.120508in}}%
\pgfpathlineto{\pgfqpoint{1.600098in}{1.129246in}}%
\pgfpathlineto{\pgfqpoint{1.617626in}{1.138290in}}%
\pgfpathlineto{\pgfqpoint{1.625174in}{1.148780in}}%
\pgfpathlineto{\pgfqpoint{1.632710in}{1.159498in}}%
\pgfpathlineto{\pgfqpoint{1.615204in}{1.150708in}}%
\pgfpathlineto{\pgfqpoint{1.597773in}{1.142228in}}%
\pgfpathlineto{\pgfqpoint{1.590216in}{1.131252in}}%
\pgfpathlineto{\pgfqpoint{1.582646in}{1.120508in}}%
\pgfpathclose%
\pgfusepath{fill}%
\end{pgfscope}%
\begin{pgfscope}%
\pgfpathrectangle{\pgfqpoint{0.407592in}{0.222904in}}{\pgfqpoint{1.411163in}{1.411163in}}%
\pgfusepath{clip}%
\pgfsetbuttcap%
\pgfsetroundjoin%
\definecolor{currentfill}{rgb}{0.909804,0.698039,0.423529}%
\pgfsetfillcolor{currentfill}%
\pgfsetfillopacity{0.120000}%
\pgfsetlinewidth{0.000000pt}%
\definecolor{currentstroke}{rgb}{0.000000,0.000000,0.000000}%
\pgfsetstrokecolor{currentstroke}%
\pgfsetdash{}{0pt}%
\pgfpathmoveto{\pgfqpoint{1.006825in}{0.864545in}}%
\pgfpathlineto{\pgfqpoint{1.023262in}{0.860118in}}%
\pgfpathlineto{\pgfqpoint{1.039735in}{0.856663in}}%
\pgfpathlineto{\pgfqpoint{1.048062in}{0.857423in}}%
\pgfpathlineto{\pgfqpoint{1.056357in}{0.858685in}}%
\pgfpathlineto{\pgfqpoint{1.039922in}{0.862105in}}%
\pgfpathlineto{\pgfqpoint{1.023524in}{0.866521in}}%
\pgfpathlineto{\pgfqpoint{1.015191in}{0.865283in}}%
\pgfpathlineto{\pgfqpoint{1.006825in}{0.864545in}}%
\pgfpathclose%
\pgfusepath{fill}%
\end{pgfscope}%
\begin{pgfscope}%
\pgfpathrectangle{\pgfqpoint{0.407592in}{0.222904in}}{\pgfqpoint{1.411163in}{1.411163in}}%
\pgfusepath{clip}%
\pgfsetbuttcap%
\pgfsetroundjoin%
\definecolor{currentfill}{rgb}{0.909804,0.698039,0.423529}%
\pgfsetfillcolor{currentfill}%
\pgfsetfillopacity{0.120000}%
\pgfsetlinewidth{0.000000pt}%
\definecolor{currentstroke}{rgb}{0.000000,0.000000,0.000000}%
\pgfsetstrokecolor{currentstroke}%
\pgfsetdash{}{0pt}%
\pgfpathmoveto{\pgfqpoint{1.697887in}{1.221889in}}%
\pgfpathlineto{\pgfqpoint{1.715576in}{1.231019in}}%
\pgfpathlineto{\pgfqpoint{1.733342in}{1.240411in}}%
\pgfpathlineto{\pgfqpoint{1.740775in}{1.251626in}}%
\pgfpathlineto{\pgfqpoint{1.748197in}{1.263018in}}%
\pgfpathlineto{\pgfqpoint{1.730454in}{1.253875in}}%
\pgfpathlineto{\pgfqpoint{1.712787in}{1.245000in}}%
\pgfpathlineto{\pgfqpoint{1.705342in}{1.233356in}}%
\pgfpathlineto{\pgfqpoint{1.697887in}{1.221889in}}%
\pgfpathclose%
\pgfusepath{fill}%
\end{pgfscope}%
\begin{pgfscope}%
\pgfpathrectangle{\pgfqpoint{0.407592in}{0.222904in}}{\pgfqpoint{1.411163in}{1.411163in}}%
\pgfusepath{clip}%
\pgfsetbuttcap%
\pgfsetroundjoin%
\definecolor{currentfill}{rgb}{0.909804,0.698039,0.423529}%
\pgfsetfillcolor{currentfill}%
\pgfsetfillopacity{0.120000}%
\pgfsetlinewidth{0.000000pt}%
\definecolor{currentstroke}{rgb}{0.000000,0.000000,0.000000}%
\pgfsetstrokecolor{currentstroke}%
\pgfsetdash{}{0pt}%
\pgfpathmoveto{\pgfqpoint{1.352782in}{0.946044in}}%
\pgfpathlineto{\pgfqpoint{1.369780in}{0.952930in}}%
\pgfpathlineto{\pgfqpoint{1.386849in}{0.960301in}}%
\pgfpathlineto{\pgfqpoint{1.394649in}{0.967980in}}%
\pgfpathlineto{\pgfqpoint{1.402433in}{0.976048in}}%
\pgfpathlineto{\pgfqpoint{1.385386in}{0.968832in}}%
\pgfpathlineto{\pgfqpoint{1.368408in}{0.962103in}}%
\pgfpathlineto{\pgfqpoint{1.360603in}{0.953869in}}%
\pgfpathlineto{\pgfqpoint{1.352782in}{0.946044in}}%
\pgfpathclose%
\pgfusepath{fill}%
\end{pgfscope}%
\begin{pgfscope}%
\pgfpathrectangle{\pgfqpoint{0.407592in}{0.222904in}}{\pgfqpoint{1.411163in}{1.411163in}}%
\pgfusepath{clip}%
\pgfsetbuttcap%
\pgfsetroundjoin%
\definecolor{currentfill}{rgb}{0.909804,0.698039,0.423529}%
\pgfsetfillcolor{currentfill}%
\pgfsetfillopacity{0.120000}%
\pgfsetlinewidth{0.000000pt}%
\definecolor{currentstroke}{rgb}{0.000000,0.000000,0.000000}%
\pgfsetstrokecolor{currentstroke}%
\pgfsetdash{}{0pt}%
\pgfpathmoveto{\pgfqpoint{0.675900in}{1.043512in}}%
\pgfpathlineto{\pgfqpoint{0.692205in}{1.028748in}}%
\pgfpathlineto{\pgfqpoint{0.708521in}{1.014388in}}%
\pgfpathlineto{\pgfqpoint{0.717494in}{1.013094in}}%
\pgfpathlineto{\pgfqpoint{0.726426in}{1.012114in}}%
\pgfpathlineto{\pgfqpoint{0.710153in}{1.026711in}}%
\pgfpathlineto{\pgfqpoint{0.693891in}{1.041708in}}%
\pgfpathlineto{\pgfqpoint{0.684916in}{1.042457in}}%
\pgfpathlineto{\pgfqpoint{0.675900in}{1.043512in}}%
\pgfpathclose%
\pgfusepath{fill}%
\end{pgfscope}%
\begin{pgfscope}%
\pgfpathrectangle{\pgfqpoint{0.407592in}{0.222904in}}{\pgfqpoint{1.411163in}{1.411163in}}%
\pgfusepath{clip}%
\pgfsetbuttcap%
\pgfsetroundjoin%
\definecolor{currentfill}{rgb}{0.909804,0.698039,0.423529}%
\pgfsetfillcolor{currentfill}%
\pgfsetfillopacity{0.120000}%
\pgfsetlinewidth{0.000000pt}%
\definecolor{currentstroke}{rgb}{0.000000,0.000000,0.000000}%
\pgfsetstrokecolor{currentstroke}%
\pgfsetdash{}{0pt}%
\pgfpathmoveto{\pgfqpoint{0.559870in}{1.140715in}}%
\pgfpathlineto{\pgfqpoint{0.576157in}{1.124282in}}%
\pgfpathlineto{\pgfqpoint{0.592453in}{1.108112in}}%
\pgfpathlineto{\pgfqpoint{0.601641in}{1.106948in}}%
\pgfpathlineto{\pgfqpoint{0.610786in}{1.106057in}}%
\pgfpathlineto{\pgfqpoint{0.594533in}{1.122442in}}%
\pgfpathlineto{\pgfqpoint{0.578289in}{1.139084in}}%
\pgfpathlineto{\pgfqpoint{0.569101in}{1.139767in}}%
\pgfpathlineto{\pgfqpoint{0.559870in}{1.140715in}}%
\pgfpathclose%
\pgfusepath{fill}%
\end{pgfscope}%
\begin{pgfscope}%
\pgfpathrectangle{\pgfqpoint{0.407592in}{0.222904in}}{\pgfqpoint{1.411163in}{1.411163in}}%
\pgfusepath{clip}%
\pgfsetbuttcap%
\pgfsetroundjoin%
\definecolor{currentfill}{rgb}{0.909804,0.698039,0.423529}%
\pgfsetfillcolor{currentfill}%
\pgfsetfillopacity{0.120000}%
\pgfsetlinewidth{0.000000pt}%
\definecolor{currentstroke}{rgb}{0.000000,0.000000,0.000000}%
\pgfsetstrokecolor{currentstroke}%
\pgfsetdash{}{0pt}%
\pgfpathmoveto{\pgfqpoint{0.874480in}{0.911207in}}%
\pgfpathlineto{\pgfqpoint{0.890842in}{0.901641in}}%
\pgfpathlineto{\pgfqpoint{0.907226in}{0.892859in}}%
\pgfpathlineto{\pgfqpoint{0.915820in}{0.891941in}}%
\pgfpathlineto{\pgfqpoint{0.924376in}{0.891435in}}%
\pgfpathlineto{\pgfqpoint{0.908036in}{0.900351in}}%
\pgfpathlineto{\pgfqpoint{0.891716in}{0.910071in}}%
\pgfpathlineto{\pgfqpoint{0.883117in}{0.910437in}}%
\pgfpathlineto{\pgfqpoint{0.874480in}{0.911207in}}%
\pgfpathclose%
\pgfusepath{fill}%
\end{pgfscope}%
\begin{pgfscope}%
\pgfpathrectangle{\pgfqpoint{0.407592in}{0.222904in}}{\pgfqpoint{1.411163in}{1.411163in}}%
\pgfusepath{clip}%
\pgfsetbuttcap%
\pgfsetroundjoin%
\definecolor{currentfill}{rgb}{0.909804,0.698039,0.423529}%
\pgfsetfillcolor{currentfill}%
\pgfsetfillopacity{0.120000}%
\pgfsetlinewidth{0.000000pt}%
\definecolor{currentstroke}{rgb}{0.000000,0.000000,0.000000}%
\pgfsetstrokecolor{currentstroke}%
\pgfsetdash{}{0pt}%
\pgfpathmoveto{\pgfqpoint{1.517476in}{1.062539in}}%
\pgfpathlineto{\pgfqpoint{1.534822in}{1.071118in}}%
\pgfpathlineto{\pgfqpoint{1.552242in}{1.080028in}}%
\pgfpathlineto{\pgfqpoint{1.559862in}{1.089754in}}%
\pgfpathlineto{\pgfqpoint{1.567470in}{1.099749in}}%
\pgfpathlineto{\pgfqpoint{1.550071in}{1.091079in}}%
\pgfpathlineto{\pgfqpoint{1.532746in}{1.082745in}}%
\pgfpathlineto{\pgfqpoint{1.525117in}{1.072504in}}%
\pgfpathlineto{\pgfqpoint{1.517476in}{1.062539in}}%
\pgfpathclose%
\pgfusepath{fill}%
\end{pgfscope}%
\begin{pgfscope}%
\pgfpathrectangle{\pgfqpoint{0.407592in}{0.222904in}}{\pgfqpoint{1.411163in}{1.411163in}}%
\pgfusepath{clip}%
\pgfsetbuttcap%
\pgfsetroundjoin%
\definecolor{currentfill}{rgb}{0.909804,0.698039,0.423529}%
\pgfsetfillcolor{currentfill}%
\pgfsetfillopacity{0.120000}%
\pgfsetlinewidth{0.000000pt}%
\definecolor{currentstroke}{rgb}{0.000000,0.000000,0.000000}%
\pgfsetstrokecolor{currentstroke}%
\pgfsetdash{}{0pt}%
\pgfpathmoveto{\pgfqpoint{1.632710in}{1.159498in}}%
\pgfpathlineto{\pgfqpoint{1.650291in}{1.168581in}}%
\pgfpathlineto{\pgfqpoint{1.667949in}{1.177943in}}%
\pgfpathlineto{\pgfqpoint{1.675450in}{1.188627in}}%
\pgfpathlineto{\pgfqpoint{1.682941in}{1.199517in}}%
\pgfpathlineto{\pgfqpoint{1.665305in}{1.190406in}}%
\pgfpathlineto{\pgfqpoint{1.647746in}{1.181579in}}%
\pgfpathlineto{\pgfqpoint{1.640234in}{1.170433in}}%
\pgfpathlineto{\pgfqpoint{1.632710in}{1.159498in}}%
\pgfpathclose%
\pgfusepath{fill}%
\end{pgfscope}%
\begin{pgfscope}%
\pgfpathrectangle{\pgfqpoint{0.407592in}{0.222904in}}{\pgfqpoint{1.411163in}{1.411163in}}%
\pgfusepath{clip}%
\pgfsetbuttcap%
\pgfsetroundjoin%
\definecolor{currentfill}{rgb}{0.909804,0.698039,0.423529}%
\pgfsetfillcolor{currentfill}%
\pgfsetfillopacity{0.120000}%
\pgfsetlinewidth{0.000000pt}%
\definecolor{currentstroke}{rgb}{0.000000,0.000000,0.000000}%
\pgfsetstrokecolor{currentstroke}%
\pgfsetdash{}{0pt}%
\pgfpathmoveto{\pgfqpoint{0.791640in}{0.958388in}}%
\pgfpathlineto{\pgfqpoint{0.807978in}{0.946302in}}%
\pgfpathlineto{\pgfqpoint{0.824331in}{0.934830in}}%
\pgfpathlineto{\pgfqpoint{0.833090in}{0.933502in}}%
\pgfpathlineto{\pgfqpoint{0.841810in}{0.932535in}}%
\pgfpathlineto{\pgfqpoint{0.825501in}{0.944217in}}%
\pgfpathlineto{\pgfqpoint{0.809207in}{0.956520in}}%
\pgfpathlineto{\pgfqpoint{0.800444in}{0.957276in}}%
\pgfpathlineto{\pgfqpoint{0.791640in}{0.958388in}}%
\pgfpathclose%
\pgfusepath{fill}%
\end{pgfscope}%
\begin{pgfscope}%
\pgfpathrectangle{\pgfqpoint{0.407592in}{0.222904in}}{\pgfqpoint{1.411163in}{1.411163in}}%
\pgfusepath{clip}%
\pgfsetbuttcap%
\pgfsetroundjoin%
\definecolor{currentfill}{rgb}{0.909804,0.698039,0.423529}%
\pgfsetfillcolor{currentfill}%
\pgfsetfillopacity{0.120000}%
\pgfsetlinewidth{0.000000pt}%
\definecolor{currentstroke}{rgb}{0.000000,0.000000,0.000000}%
\pgfsetstrokecolor{currentstroke}%
\pgfsetdash{}{0pt}%
\pgfpathmoveto{\pgfqpoint{1.402433in}{0.976048in}}%
\pgfpathlineto{\pgfqpoint{1.419551in}{0.983717in}}%
\pgfpathlineto{\pgfqpoint{1.436741in}{0.991806in}}%
\pgfpathlineto{\pgfqpoint{1.444488in}{1.000073in}}%
\pgfpathlineto{\pgfqpoint{1.452220in}{1.008691in}}%
\pgfpathlineto{\pgfqpoint{1.435051in}{1.000789in}}%
\pgfpathlineto{\pgfqpoint{1.417954in}{0.993309in}}%
\pgfpathlineto{\pgfqpoint{1.410201in}{0.984495in}}%
\pgfpathlineto{\pgfqpoint{1.402433in}{0.976048in}}%
\pgfpathclose%
\pgfusepath{fill}%
\end{pgfscope}%
\begin{pgfscope}%
\pgfpathrectangle{\pgfqpoint{0.407592in}{0.222904in}}{\pgfqpoint{1.411163in}{1.411163in}}%
\pgfusepath{clip}%
\pgfsetbuttcap%
\pgfsetroundjoin%
\definecolor{currentfill}{rgb}{0.909804,0.698039,0.423529}%
\pgfsetfillcolor{currentfill}%
\pgfsetfillopacity{0.120000}%
\pgfsetlinewidth{0.000000pt}%
\definecolor{currentstroke}{rgb}{0.000000,0.000000,0.000000}%
\pgfsetstrokecolor{currentstroke}%
\pgfsetdash{}{0pt}%
\pgfpathmoveto{\pgfqpoint{0.957128in}{0.876160in}}%
\pgfpathlineto{\pgfqpoint{0.973544in}{0.869864in}}%
\pgfpathlineto{\pgfqpoint{0.989989in}{0.864492in}}%
\pgfpathlineto{\pgfqpoint{0.998424in}{0.864288in}}%
\pgfpathlineto{\pgfqpoint{1.006825in}{0.864545in}}%
\pgfpathlineto{\pgfqpoint{0.990421in}{0.869938in}}%
\pgfpathlineto{\pgfqpoint{0.974046in}{0.876284in}}%
\pgfpathlineto{\pgfqpoint{0.965605in}{0.875994in}}%
\pgfpathlineto{\pgfqpoint{0.957128in}{0.876160in}}%
\pgfpathclose%
\pgfusepath{fill}%
\end{pgfscope}%
\begin{pgfscope}%
\pgfpathrectangle{\pgfqpoint{0.407592in}{0.222904in}}{\pgfqpoint{1.411163in}{1.411163in}}%
\pgfusepath{clip}%
\pgfsetbuttcap%
\pgfsetroundjoin%
\definecolor{currentfill}{rgb}{0.909804,0.698039,0.423529}%
\pgfsetfillcolor{currentfill}%
\pgfsetfillopacity{0.120000}%
\pgfsetlinewidth{0.000000pt}%
\definecolor{currentstroke}{rgb}{0.000000,0.000000,0.000000}%
\pgfsetstrokecolor{currentstroke}%
\pgfsetdash{}{0pt}%
\pgfpathmoveto{\pgfqpoint{1.188366in}{0.869556in}}%
\pgfpathlineto{\pgfqpoint{1.205074in}{0.872745in}}%
\pgfpathlineto{\pgfqpoint{1.221842in}{0.876697in}}%
\pgfpathlineto{\pgfqpoint{1.229870in}{0.880878in}}%
\pgfpathlineto{\pgfqpoint{1.237874in}{0.885564in}}%
\pgfpathlineto{\pgfqpoint{1.221133in}{0.881569in}}%
\pgfpathlineto{\pgfqpoint{1.204453in}{0.878338in}}%
\pgfpathlineto{\pgfqpoint{1.196422in}{0.873682in}}%
\pgfpathlineto{\pgfqpoint{1.188366in}{0.869556in}}%
\pgfpathclose%
\pgfusepath{fill}%
\end{pgfscope}%
\begin{pgfscope}%
\pgfpathrectangle{\pgfqpoint{0.407592in}{0.222904in}}{\pgfqpoint{1.411163in}{1.411163in}}%
\pgfusepath{clip}%
\pgfsetbuttcap%
\pgfsetroundjoin%
\definecolor{currentfill}{rgb}{0.909804,0.698039,0.423529}%
\pgfsetfillcolor{currentfill}%
\pgfsetfillopacity{0.120000}%
\pgfsetlinewidth{0.000000pt}%
\definecolor{currentstroke}{rgb}{0.000000,0.000000,0.000000}%
\pgfsetstrokecolor{currentstroke}%
\pgfsetdash{}{0pt}%
\pgfpathmoveto{\pgfqpoint{1.138873in}{0.859153in}}%
\pgfpathlineto{\pgfqpoint{1.155500in}{0.860586in}}%
\pgfpathlineto{\pgfqpoint{1.172180in}{0.862877in}}%
\pgfpathlineto{\pgfqpoint{1.180286in}{0.865957in}}%
\pgfpathlineto{\pgfqpoint{1.188366in}{0.869556in}}%
\pgfpathlineto{\pgfqpoint{1.171716in}{0.867182in}}%
\pgfpathlineto{\pgfqpoint{1.155119in}{0.865670in}}%
\pgfpathlineto{\pgfqpoint{1.147010in}{0.862142in}}%
\pgfpathlineto{\pgfqpoint{1.138873in}{0.859153in}}%
\pgfpathclose%
\pgfusepath{fill}%
\end{pgfscope}%
\begin{pgfscope}%
\pgfpathrectangle{\pgfqpoint{0.407592in}{0.222904in}}{\pgfqpoint{1.411163in}{1.411163in}}%
\pgfusepath{clip}%
\pgfsetbuttcap%
\pgfsetroundjoin%
\definecolor{currentfill}{rgb}{0.909804,0.698039,0.423529}%
\pgfsetfillcolor{currentfill}%
\pgfsetfillopacity{0.120000}%
\pgfsetlinewidth{0.000000pt}%
\definecolor{currentstroke}{rgb}{0.000000,0.000000,0.000000}%
\pgfsetstrokecolor{currentstroke}%
\pgfsetdash{}{0pt}%
\pgfpathmoveto{\pgfqpoint{1.237874in}{0.885564in}}%
\pgfpathlineto{\pgfqpoint{1.254677in}{0.890272in}}%
\pgfpathlineto{\pgfqpoint{1.271543in}{0.895643in}}%
\pgfpathlineto{\pgfqpoint{1.279501in}{0.900838in}}%
\pgfpathlineto{\pgfqpoint{1.287437in}{0.906512in}}%
\pgfpathlineto{\pgfqpoint{1.270595in}{0.901152in}}%
\pgfpathlineto{\pgfqpoint{1.253817in}{0.896454in}}%
\pgfpathlineto{\pgfqpoint{1.245856in}{0.890757in}}%
\pgfpathlineto{\pgfqpoint{1.237874in}{0.885564in}}%
\pgfpathclose%
\pgfusepath{fill}%
\end{pgfscope}%
\begin{pgfscope}%
\pgfpathrectangle{\pgfqpoint{0.407592in}{0.222904in}}{\pgfqpoint{1.411163in}{1.411163in}}%
\pgfusepath{clip}%
\pgfsetbuttcap%
\pgfsetroundjoin%
\definecolor{currentfill}{rgb}{0.909804,0.698039,0.423529}%
\pgfsetfillcolor{currentfill}%
\pgfsetfillopacity{0.120000}%
\pgfsetlinewidth{0.000000pt}%
\definecolor{currentstroke}{rgb}{0.000000,0.000000,0.000000}%
\pgfsetstrokecolor{currentstroke}%
\pgfsetdash{}{0pt}%
\pgfpathmoveto{\pgfqpoint{0.592453in}{1.108112in}}%
\pgfpathlineto{\pgfqpoint{0.608759in}{1.092227in}}%
\pgfpathlineto{\pgfqpoint{0.625076in}{1.076648in}}%
\pgfpathlineto{\pgfqpoint{0.634220in}{1.075253in}}%
\pgfpathlineto{\pgfqpoint{0.643322in}{1.074138in}}%
\pgfpathlineto{\pgfqpoint{0.627049in}{1.089948in}}%
\pgfpathlineto{\pgfqpoint{0.610786in}{1.106057in}}%
\pgfpathlineto{\pgfqpoint{0.601641in}{1.106948in}}%
\pgfpathlineto{\pgfqpoint{0.592453in}{1.108112in}}%
\pgfpathclose%
\pgfusepath{fill}%
\end{pgfscope}%
\begin{pgfscope}%
\pgfpathrectangle{\pgfqpoint{0.407592in}{0.222904in}}{\pgfqpoint{1.411163in}{1.411163in}}%
\pgfusepath{clip}%
\pgfsetbuttcap%
\pgfsetroundjoin%
\definecolor{currentfill}{rgb}{0.909804,0.698039,0.423529}%
\pgfsetfillcolor{currentfill}%
\pgfsetfillopacity{0.120000}%
\pgfsetlinewidth{0.000000pt}%
\definecolor{currentstroke}{rgb}{0.000000,0.000000,0.000000}%
\pgfsetstrokecolor{currentstroke}%
\pgfsetdash{}{0pt}%
\pgfpathmoveto{\pgfqpoint{0.708521in}{1.014388in}}%
\pgfpathlineto{\pgfqpoint{0.724850in}{1.000462in}}%
\pgfpathlineto{\pgfqpoint{0.741190in}{0.987002in}}%
\pgfpathlineto{\pgfqpoint{0.750119in}{0.985465in}}%
\pgfpathlineto{\pgfqpoint{0.759007in}{0.984248in}}%
\pgfpathlineto{\pgfqpoint{0.742710in}{0.997948in}}%
\pgfpathlineto{\pgfqpoint{0.726426in}{1.012114in}}%
\pgfpathlineto{\pgfqpoint{0.717494in}{1.013094in}}%
\pgfpathlineto{\pgfqpoint{0.708521in}{1.014388in}}%
\pgfpathclose%
\pgfusepath{fill}%
\end{pgfscope}%
\begin{pgfscope}%
\pgfpathrectangle{\pgfqpoint{0.407592in}{0.222904in}}{\pgfqpoint{1.411163in}{1.411163in}}%
\pgfusepath{clip}%
\pgfsetbuttcap%
\pgfsetroundjoin%
\definecolor{currentfill}{rgb}{0.909804,0.698039,0.423529}%
\pgfsetfillcolor{currentfill}%
\pgfsetfillopacity{0.120000}%
\pgfsetlinewidth{0.000000pt}%
\definecolor{currentstroke}{rgb}{0.000000,0.000000,0.000000}%
\pgfsetstrokecolor{currentstroke}%
\pgfsetdash{}{0pt}%
\pgfpathmoveto{\pgfqpoint{1.089346in}{0.854809in}}%
\pgfpathlineto{\pgfqpoint{1.105907in}{0.854317in}}%
\pgfpathlineto{\pgfqpoint{1.122515in}{0.854750in}}%
\pgfpathlineto{\pgfqpoint{1.130709in}{0.856693in}}%
\pgfpathlineto{\pgfqpoint{1.138873in}{0.859153in}}%
\pgfpathlineto{\pgfqpoint{1.122298in}{0.858620in}}%
\pgfpathlineto{\pgfqpoint{1.105771in}{0.859027in}}%
\pgfpathlineto{\pgfqpoint{1.097574in}{0.856654in}}%
\pgfpathlineto{\pgfqpoint{1.089346in}{0.854809in}}%
\pgfpathclose%
\pgfusepath{fill}%
\end{pgfscope}%
\begin{pgfscope}%
\pgfpathrectangle{\pgfqpoint{0.407592in}{0.222904in}}{\pgfqpoint{1.411163in}{1.411163in}}%
\pgfusepath{clip}%
\pgfsetbuttcap%
\pgfsetroundjoin%
\definecolor{currentfill}{rgb}{0.909804,0.698039,0.423529}%
\pgfsetfillcolor{currentfill}%
\pgfsetfillopacity{0.120000}%
\pgfsetlinewidth{0.000000pt}%
\definecolor{currentstroke}{rgb}{0.000000,0.000000,0.000000}%
\pgfsetstrokecolor{currentstroke}%
\pgfsetdash{}{0pt}%
\pgfpathmoveto{\pgfqpoint{1.567470in}{1.099749in}}%
\pgfpathlineto{\pgfqpoint{1.584943in}{1.108737in}}%
\pgfpathlineto{\pgfqpoint{1.602493in}{1.118023in}}%
\pgfpathlineto{\pgfqpoint{1.610065in}{1.128035in}}%
\pgfpathlineto{\pgfqpoint{1.617626in}{1.138290in}}%
\pgfpathlineto{\pgfqpoint{1.600098in}{1.129246in}}%
\pgfpathlineto{\pgfqpoint{1.582646in}{1.120508in}}%
\pgfpathlineto{\pgfqpoint{1.575064in}{1.110003in}}%
\pgfpathlineto{\pgfqpoint{1.567470in}{1.099749in}}%
\pgfpathclose%
\pgfusepath{fill}%
\end{pgfscope}%
\begin{pgfscope}%
\pgfpathrectangle{\pgfqpoint{0.407592in}{0.222904in}}{\pgfqpoint{1.411163in}{1.411163in}}%
\pgfusepath{clip}%
\pgfsetbuttcap%
\pgfsetroundjoin%
\definecolor{currentfill}{rgb}{0.909804,0.698039,0.423529}%
\pgfsetfillcolor{currentfill}%
\pgfsetfillopacity{0.120000}%
\pgfsetlinewidth{0.000000pt}%
\definecolor{currentstroke}{rgb}{0.000000,0.000000,0.000000}%
\pgfsetstrokecolor{currentstroke}%
\pgfsetdash{}{0pt}%
\pgfpathmoveto{\pgfqpoint{1.748197in}{1.263018in}}%
\pgfpathlineto{\pgfqpoint{1.766017in}{1.272415in}}%
\pgfpathlineto{\pgfqpoint{1.783915in}{1.282058in}}%
\pgfpathlineto{\pgfqpoint{1.791303in}{1.293374in}}%
\pgfpathlineto{\pgfqpoint{1.773417in}{1.283853in}}%
\pgfpathlineto{\pgfqpoint{1.755608in}{1.274579in}}%
\pgfusepath{fill}%
\end{pgfscope}%
\begin{pgfscope}%
\pgfpathrectangle{\pgfqpoint{0.407592in}{0.222904in}}{\pgfqpoint{1.411163in}{1.411163in}}%
\pgfusepath{clip}%
\pgfsetbuttcap%
\pgfsetroundjoin%
\definecolor{currentfill}{rgb}{0.909804,0.698039,0.423529}%
\pgfsetfillcolor{currentfill}%
\pgfsetfillopacity{0.120000}%
\pgfsetlinewidth{0.000000pt}%
\definecolor{currentstroke}{rgb}{0.000000,0.000000,0.000000}%
\pgfsetstrokecolor{currentstroke}%
\pgfsetdash{}{0pt}%
\pgfpathmoveto{\pgfqpoint{1.452220in}{1.008691in}}%
\pgfpathlineto{\pgfqpoint{1.469461in}{1.016988in}}%
\pgfpathlineto{\pgfqpoint{1.486776in}{1.025652in}}%
\pgfpathlineto{\pgfqpoint{1.494472in}{1.034408in}}%
\pgfpathlineto{\pgfqpoint{1.502153in}{1.043481in}}%
\pgfpathlineto{\pgfqpoint{1.484860in}{1.035024in}}%
\pgfpathlineto{\pgfqpoint{1.467640in}{1.026940in}}%
\pgfpathlineto{\pgfqpoint{1.459937in}{1.017650in}}%
\pgfpathlineto{\pgfqpoint{1.452220in}{1.008691in}}%
\pgfpathclose%
\pgfusepath{fill}%
\end{pgfscope}%
\begin{pgfscope}%
\pgfpathrectangle{\pgfqpoint{0.407592in}{0.222904in}}{\pgfqpoint{1.411163in}{1.411163in}}%
\pgfusepath{clip}%
\pgfsetbuttcap%
\pgfsetroundjoin%
\definecolor{currentfill}{rgb}{0.909804,0.698039,0.423529}%
\pgfsetfillcolor{currentfill}%
\pgfsetfillopacity{0.120000}%
\pgfsetlinewidth{0.000000pt}%
\definecolor{currentstroke}{rgb}{0.000000,0.000000,0.000000}%
\pgfsetstrokecolor{currentstroke}%
\pgfsetdash{}{0pt}%
\pgfpathmoveto{\pgfqpoint{1.287437in}{0.906512in}}%
\pgfpathlineto{\pgfqpoint{1.304345in}{0.912484in}}%
\pgfpathlineto{\pgfqpoint{1.321320in}{0.919025in}}%
\pgfpathlineto{\pgfqpoint{1.329213in}{0.925121in}}%
\pgfpathlineto{\pgfqpoint{1.337087in}{0.931661in}}%
\pgfpathlineto{\pgfqpoint{1.320135in}{0.925186in}}%
\pgfpathlineto{\pgfqpoint{1.303250in}{0.919278in}}%
\pgfpathlineto{\pgfqpoint{1.295353in}{0.912660in}}%
\pgfpathlineto{\pgfqpoint{1.287437in}{0.906512in}}%
\pgfpathclose%
\pgfusepath{fill}%
\end{pgfscope}%
\begin{pgfscope}%
\pgfpathrectangle{\pgfqpoint{0.407592in}{0.222904in}}{\pgfqpoint{1.411163in}{1.411163in}}%
\pgfusepath{clip}%
\pgfsetbuttcap%
\pgfsetroundjoin%
\definecolor{currentfill}{rgb}{0.909804,0.698039,0.423529}%
\pgfsetfillcolor{currentfill}%
\pgfsetfillopacity{0.120000}%
\pgfsetlinewidth{0.000000pt}%
\definecolor{currentstroke}{rgb}{0.000000,0.000000,0.000000}%
\pgfsetstrokecolor{currentstroke}%
\pgfsetdash{}{0pt}%
\pgfpathmoveto{\pgfqpoint{0.476054in}{1.210594in}}%
\pgfpathlineto{\pgfqpoint{0.492346in}{1.193469in}}%
\pgfpathlineto{\pgfqpoint{0.508647in}{1.176548in}}%
\pgfpathlineto{\pgfqpoint{0.518008in}{1.175302in}}%
\pgfpathlineto{\pgfqpoint{0.527325in}{1.174301in}}%
\pgfpathlineto{\pgfqpoint{0.511067in}{1.191420in}}%
\pgfpathlineto{\pgfqpoint{0.494818in}{1.208735in}}%
\pgfpathlineto{\pgfqpoint{0.485458in}{1.209544in}}%
\pgfpathlineto{\pgfqpoint{0.476054in}{1.210594in}}%
\pgfpathclose%
\pgfusepath{fill}%
\end{pgfscope}%
\begin{pgfscope}%
\pgfpathrectangle{\pgfqpoint{0.407592in}{0.222904in}}{\pgfqpoint{1.411163in}{1.411163in}}%
\pgfusepath{clip}%
\pgfsetbuttcap%
\pgfsetroundjoin%
\definecolor{currentfill}{rgb}{0.909804,0.698039,0.423529}%
\pgfsetfillcolor{currentfill}%
\pgfsetfillopacity{0.120000}%
\pgfsetlinewidth{0.000000pt}%
\definecolor{currentstroke}{rgb}{0.000000,0.000000,0.000000}%
\pgfsetstrokecolor{currentstroke}%
\pgfsetdash{}{0pt}%
\pgfpathmoveto{\pgfqpoint{1.682941in}{1.199517in}}%
\pgfpathlineto{\pgfqpoint{1.700652in}{1.208899in}}%
\pgfpathlineto{\pgfqpoint{1.718441in}{1.218537in}}%
\pgfpathlineto{\pgfqpoint{1.725897in}{1.229379in}}%
\pgfpathlineto{\pgfqpoint{1.733342in}{1.240411in}}%
\pgfpathlineto{\pgfqpoint{1.715576in}{1.231019in}}%
\pgfpathlineto{\pgfqpoint{1.697887in}{1.221889in}}%
\pgfpathlineto{\pgfqpoint{1.690419in}{1.210607in}}%
\pgfpathlineto{\pgfqpoint{1.682941in}{1.199517in}}%
\pgfpathclose%
\pgfusepath{fill}%
\end{pgfscope}%
\begin{pgfscope}%
\pgfpathrectangle{\pgfqpoint{0.407592in}{0.222904in}}{\pgfqpoint{1.411163in}{1.411163in}}%
\pgfusepath{clip}%
\pgfsetbuttcap%
\pgfsetroundjoin%
\definecolor{currentfill}{rgb}{0.909804,0.698039,0.423529}%
\pgfsetfillcolor{currentfill}%
\pgfsetfillopacity{0.120000}%
\pgfsetlinewidth{0.000000pt}%
\definecolor{currentstroke}{rgb}{0.000000,0.000000,0.000000}%
\pgfsetstrokecolor{currentstroke}%
\pgfsetdash{}{0pt}%
\pgfpathmoveto{\pgfqpoint{1.039735in}{0.856663in}}%
\pgfpathlineto{\pgfqpoint{1.056246in}{0.854176in}}%
\pgfpathlineto{\pgfqpoint{1.072799in}{0.852644in}}%
\pgfpathlineto{\pgfqpoint{1.081088in}{0.853479in}}%
\pgfpathlineto{\pgfqpoint{1.089346in}{0.854809in}}%
\pgfpathlineto{\pgfqpoint{1.072831in}{0.856258in}}%
\pgfpathlineto{\pgfqpoint{1.056357in}{0.858685in}}%
\pgfpathlineto{\pgfqpoint{1.048062in}{0.857423in}}%
\pgfpathlineto{\pgfqpoint{1.039735in}{0.856663in}}%
\pgfpathclose%
\pgfusepath{fill}%
\end{pgfscope}%
\begin{pgfscope}%
\pgfpathrectangle{\pgfqpoint{0.407592in}{0.222904in}}{\pgfqpoint{1.411163in}{1.411163in}}%
\pgfusepath{clip}%
\pgfsetbuttcap%
\pgfsetroundjoin%
\definecolor{currentfill}{rgb}{0.909804,0.698039,0.423529}%
\pgfsetfillcolor{currentfill}%
\pgfsetfillopacity{0.120000}%
\pgfsetlinewidth{0.000000pt}%
\definecolor{currentstroke}{rgb}{0.000000,0.000000,0.000000}%
\pgfsetstrokecolor{currentstroke}%
\pgfsetdash{}{0pt}%
\pgfpathmoveto{\pgfqpoint{0.824331in}{0.934830in}}%
\pgfpathlineto{\pgfqpoint{0.840701in}{0.924011in}}%
\pgfpathlineto{\pgfqpoint{0.857088in}{0.913884in}}%
\pgfpathlineto{\pgfqpoint{0.865804in}{0.912361in}}%
\pgfpathlineto{\pgfqpoint{0.874480in}{0.911207in}}%
\pgfpathlineto{\pgfqpoint{0.858136in}{0.921518in}}%
\pgfpathlineto{\pgfqpoint{0.841810in}{0.932535in}}%
\pgfpathlineto{\pgfqpoint{0.833090in}{0.933502in}}%
\pgfpathlineto{\pgfqpoint{0.824331in}{0.934830in}}%
\pgfpathclose%
\pgfusepath{fill}%
\end{pgfscope}%
\begin{pgfscope}%
\pgfpathrectangle{\pgfqpoint{0.407592in}{0.222904in}}{\pgfqpoint{1.411163in}{1.411163in}}%
\pgfusepath{clip}%
\pgfsetbuttcap%
\pgfsetroundjoin%
\definecolor{currentfill}{rgb}{0.909804,0.698039,0.423529}%
\pgfsetfillcolor{currentfill}%
\pgfsetfillopacity{0.120000}%
\pgfsetlinewidth{0.000000pt}%
\definecolor{currentstroke}{rgb}{0.000000,0.000000,0.000000}%
\pgfsetstrokecolor{currentstroke}%
\pgfsetdash{}{0pt}%
\pgfpathmoveto{\pgfqpoint{0.907226in}{0.892859in}}%
\pgfpathlineto{\pgfqpoint{0.923633in}{0.884895in}}%
\pgfpathlineto{\pgfqpoint{0.940064in}{0.877782in}}%
\pgfpathlineto{\pgfqpoint{0.948615in}{0.876763in}}%
\pgfpathlineto{\pgfqpoint{0.957128in}{0.876160in}}%
\pgfpathlineto{\pgfqpoint{0.940740in}{0.883360in}}%
\pgfpathlineto{\pgfqpoint{0.924376in}{0.891435in}}%
\pgfpathlineto{\pgfqpoint{0.915820in}{0.891941in}}%
\pgfpathlineto{\pgfqpoint{0.907226in}{0.892859in}}%
\pgfpathclose%
\pgfusepath{fill}%
\end{pgfscope}%
\begin{pgfscope}%
\pgfpathrectangle{\pgfqpoint{0.407592in}{0.222904in}}{\pgfqpoint{1.411163in}{1.411163in}}%
\pgfusepath{clip}%
\pgfsetbuttcap%
\pgfsetroundjoin%
\definecolor{currentfill}{rgb}{0.909804,0.698039,0.423529}%
\pgfsetfillcolor{currentfill}%
\pgfsetfillopacity{0.120000}%
\pgfsetlinewidth{0.000000pt}%
\definecolor{currentstroke}{rgb}{0.000000,0.000000,0.000000}%
\pgfsetstrokecolor{currentstroke}%
\pgfsetdash{}{0pt}%
\pgfpathmoveto{\pgfqpoint{1.337087in}{0.931661in}}%
\pgfpathlineto{\pgfqpoint{1.354108in}{0.938661in}}%
\pgfpathlineto{\pgfqpoint{1.371199in}{0.946146in}}%
\pgfpathlineto{\pgfqpoint{1.379032in}{0.953020in}}%
\pgfpathlineto{\pgfqpoint{1.386849in}{0.960301in}}%
\pgfpathlineto{\pgfqpoint{1.369780in}{0.952930in}}%
\pgfpathlineto{\pgfqpoint{1.352782in}{0.946044in}}%
\pgfpathlineto{\pgfqpoint{1.344943in}{0.938638in}}%
\pgfpathlineto{\pgfqpoint{1.337087in}{0.931661in}}%
\pgfpathclose%
\pgfusepath{fill}%
\end{pgfscope}%
\begin{pgfscope}%
\pgfpathrectangle{\pgfqpoint{0.407592in}{0.222904in}}{\pgfqpoint{1.411163in}{1.411163in}}%
\pgfusepath{clip}%
\pgfsetbuttcap%
\pgfsetroundjoin%
\definecolor{currentfill}{rgb}{0.909804,0.698039,0.423529}%
\pgfsetfillcolor{currentfill}%
\pgfsetfillopacity{0.120000}%
\pgfsetlinewidth{0.000000pt}%
\definecolor{currentstroke}{rgb}{0.000000,0.000000,0.000000}%
\pgfsetstrokecolor{currentstroke}%
\pgfsetdash{}{0pt}%
\pgfpathmoveto{\pgfqpoint{0.625076in}{1.076648in}}%
\pgfpathlineto{\pgfqpoint{0.641403in}{1.061397in}}%
\pgfpathlineto{\pgfqpoint{0.657740in}{1.046500in}}%
\pgfpathlineto{\pgfqpoint{0.666841in}{1.044863in}}%
\pgfpathlineto{\pgfqpoint{0.675900in}{1.043512in}}%
\pgfpathlineto{\pgfqpoint{0.659606in}{1.058651in}}%
\pgfpathlineto{\pgfqpoint{0.643322in}{1.074138in}}%
\pgfpathlineto{\pgfqpoint{0.634220in}{1.075253in}}%
\pgfpathlineto{\pgfqpoint{0.625076in}{1.076648in}}%
\pgfpathclose%
\pgfusepath{fill}%
\end{pgfscope}%
\begin{pgfscope}%
\pgfpathrectangle{\pgfqpoint{0.407592in}{0.222904in}}{\pgfqpoint{1.411163in}{1.411163in}}%
\pgfusepath{clip}%
\pgfsetbuttcap%
\pgfsetroundjoin%
\definecolor{currentfill}{rgb}{0.909804,0.698039,0.423529}%
\pgfsetfillcolor{currentfill}%
\pgfsetfillopacity{0.120000}%
\pgfsetlinewidth{0.000000pt}%
\definecolor{currentstroke}{rgb}{0.000000,0.000000,0.000000}%
\pgfsetstrokecolor{currentstroke}%
\pgfsetdash{}{0pt}%
\pgfpathmoveto{\pgfqpoint{0.741190in}{0.987002in}}%
\pgfpathlineto{\pgfqpoint{0.757544in}{0.974042in}}%
\pgfpathlineto{\pgfqpoint{0.773911in}{0.961616in}}%
\pgfpathlineto{\pgfqpoint{0.782796in}{0.959839in}}%
\pgfpathlineto{\pgfqpoint{0.791640in}{0.958388in}}%
\pgfpathlineto{\pgfqpoint{0.775317in}{0.971049in}}%
\pgfpathlineto{\pgfqpoint{0.759007in}{0.984248in}}%
\pgfpathlineto{\pgfqpoint{0.750119in}{0.985465in}}%
\pgfpathlineto{\pgfqpoint{0.741190in}{0.987002in}}%
\pgfpathclose%
\pgfusepath{fill}%
\end{pgfscope}%
\begin{pgfscope}%
\pgfpathrectangle{\pgfqpoint{0.407592in}{0.222904in}}{\pgfqpoint{1.411163in}{1.411163in}}%
\pgfusepath{clip}%
\pgfsetbuttcap%
\pgfsetroundjoin%
\definecolor{currentfill}{rgb}{0.909804,0.698039,0.423529}%
\pgfsetfillcolor{currentfill}%
\pgfsetfillopacity{0.120000}%
\pgfsetlinewidth{0.000000pt}%
\definecolor{currentstroke}{rgb}{0.000000,0.000000,0.000000}%
\pgfsetstrokecolor{currentstroke}%
\pgfsetdash{}{0pt}%
\pgfpathmoveto{\pgfqpoint{0.508647in}{1.176548in}}%
\pgfpathlineto{\pgfqpoint{0.524958in}{1.159848in}}%
\pgfpathlineto{\pgfqpoint{0.541278in}{1.143383in}}%
\pgfpathlineto{\pgfqpoint{0.550596in}{1.141923in}}%
\pgfpathlineto{\pgfqpoint{0.559870in}{1.140715in}}%
\pgfpathlineto{\pgfqpoint{0.543593in}{1.157394in}}%
\pgfpathlineto{\pgfqpoint{0.527325in}{1.174301in}}%
\pgfpathlineto{\pgfqpoint{0.518008in}{1.175302in}}%
\pgfpathlineto{\pgfqpoint{0.508647in}{1.176548in}}%
\pgfpathclose%
\pgfusepath{fill}%
\end{pgfscope}%
\begin{pgfscope}%
\pgfpathrectangle{\pgfqpoint{0.407592in}{0.222904in}}{\pgfqpoint{1.411163in}{1.411163in}}%
\pgfusepath{clip}%
\pgfsetbuttcap%
\pgfsetroundjoin%
\definecolor{currentfill}{rgb}{0.909804,0.698039,0.423529}%
\pgfsetfillcolor{currentfill}%
\pgfsetfillopacity{0.120000}%
\pgfsetlinewidth{0.000000pt}%
\definecolor{currentstroke}{rgb}{0.000000,0.000000,0.000000}%
\pgfsetstrokecolor{currentstroke}%
\pgfsetdash{}{0pt}%
\pgfpathmoveto{\pgfqpoint{1.502153in}{1.043481in}}%
\pgfpathlineto{\pgfqpoint{1.519521in}{1.052286in}}%
\pgfpathlineto{\pgfqpoint{1.536963in}{1.061417in}}%
\pgfpathlineto{\pgfqpoint{1.544609in}{1.070579in}}%
\pgfpathlineto{\pgfqpoint{1.552242in}{1.080028in}}%
\pgfpathlineto{\pgfqpoint{1.534822in}{1.071118in}}%
\pgfpathlineto{\pgfqpoint{1.517476in}{1.062539in}}%
\pgfpathlineto{\pgfqpoint{1.509822in}{1.052861in}}%
\pgfpathlineto{\pgfqpoint{1.502153in}{1.043481in}}%
\pgfpathclose%
\pgfusepath{fill}%
\end{pgfscope}%
\begin{pgfscope}%
\pgfpathrectangle{\pgfqpoint{0.407592in}{0.222904in}}{\pgfqpoint{1.411163in}{1.411163in}}%
\pgfusepath{clip}%
\pgfsetbuttcap%
\pgfsetroundjoin%
\definecolor{currentfill}{rgb}{0.909804,0.698039,0.423529}%
\pgfsetfillcolor{currentfill}%
\pgfsetfillopacity{0.120000}%
\pgfsetlinewidth{0.000000pt}%
\definecolor{currentstroke}{rgb}{0.000000,0.000000,0.000000}%
\pgfsetstrokecolor{currentstroke}%
\pgfsetdash{}{0pt}%
\pgfpathmoveto{\pgfqpoint{1.617626in}{1.138290in}}%
\pgfpathlineto{\pgfqpoint{1.635229in}{1.147622in}}%
\pgfpathlineto{\pgfqpoint{1.652909in}{1.157226in}}%
\pgfpathlineto{\pgfqpoint{1.660435in}{1.167474in}}%
\pgfpathlineto{\pgfqpoint{1.667949in}{1.177943in}}%
\pgfpathlineto{\pgfqpoint{1.650291in}{1.168581in}}%
\pgfpathlineto{\pgfqpoint{1.632710in}{1.159498in}}%
\pgfpathlineto{\pgfqpoint{1.625174in}{1.148780in}}%
\pgfpathlineto{\pgfqpoint{1.617626in}{1.138290in}}%
\pgfpathclose%
\pgfusepath{fill}%
\end{pgfscope}%
\begin{pgfscope}%
\pgfpathrectangle{\pgfqpoint{0.407592in}{0.222904in}}{\pgfqpoint{1.411163in}{1.411163in}}%
\pgfusepath{clip}%
\pgfsetbuttcap%
\pgfsetroundjoin%
\definecolor{currentfill}{rgb}{0.909804,0.698039,0.423529}%
\pgfsetfillcolor{currentfill}%
\pgfsetfillopacity{0.120000}%
\pgfsetlinewidth{0.000000pt}%
\definecolor{currentstroke}{rgb}{0.000000,0.000000,0.000000}%
\pgfsetstrokecolor{currentstroke}%
\pgfsetdash{}{0pt}%
\pgfpathmoveto{\pgfqpoint{0.989989in}{0.864492in}}%
\pgfpathlineto{\pgfqpoint{1.006467in}{0.860057in}}%
\pgfpathlineto{\pgfqpoint{1.022980in}{0.856565in}}%
\pgfpathlineto{\pgfqpoint{1.031374in}{0.856383in}}%
\pgfpathlineto{\pgfqpoint{1.039735in}{0.856663in}}%
\pgfpathlineto{\pgfqpoint{1.023262in}{0.860118in}}%
\pgfpathlineto{\pgfqpoint{1.006825in}{0.864545in}}%
\pgfpathlineto{\pgfqpoint{0.998424in}{0.864288in}}%
\pgfpathlineto{\pgfqpoint{0.989989in}{0.864492in}}%
\pgfpathclose%
\pgfusepath{fill}%
\end{pgfscope}%
\begin{pgfscope}%
\pgfpathrectangle{\pgfqpoint{0.407592in}{0.222904in}}{\pgfqpoint{1.411163in}{1.411163in}}%
\pgfusepath{clip}%
\pgfsetbuttcap%
\pgfsetroundjoin%
\definecolor{currentfill}{rgb}{0.909804,0.698039,0.423529}%
\pgfsetfillcolor{currentfill}%
\pgfsetfillopacity{0.120000}%
\pgfsetlinewidth{0.000000pt}%
\definecolor{currentstroke}{rgb}{0.000000,0.000000,0.000000}%
\pgfsetstrokecolor{currentstroke}%
\pgfsetdash{}{0pt}%
\pgfpathmoveto{\pgfqpoint{1.386849in}{0.960301in}}%
\pgfpathlineto{\pgfqpoint{1.403989in}{0.968123in}}%
\pgfpathlineto{\pgfqpoint{1.421200in}{0.976363in}}%
\pgfpathlineto{\pgfqpoint{1.428979in}{0.983900in}}%
\pgfpathlineto{\pgfqpoint{1.436741in}{0.991806in}}%
\pgfpathlineto{\pgfqpoint{1.419551in}{0.983717in}}%
\pgfpathlineto{\pgfqpoint{1.402433in}{0.976048in}}%
\pgfpathlineto{\pgfqpoint{1.394649in}{0.967980in}}%
\pgfpathlineto{\pgfqpoint{1.386849in}{0.960301in}}%
\pgfpathclose%
\pgfusepath{fill}%
\end{pgfscope}%
\begin{pgfscope}%
\pgfpathrectangle{\pgfqpoint{0.407592in}{0.222904in}}{\pgfqpoint{1.411163in}{1.411163in}}%
\pgfusepath{clip}%
\pgfsetbuttcap%
\pgfsetroundjoin%
\definecolor{currentfill}{rgb}{0.909804,0.698039,0.423529}%
\pgfsetfillcolor{currentfill}%
\pgfsetfillopacity{0.120000}%
\pgfsetlinewidth{0.000000pt}%
\definecolor{currentstroke}{rgb}{0.000000,0.000000,0.000000}%
\pgfsetstrokecolor{currentstroke}%
\pgfsetdash{}{0pt}%
\pgfpathmoveto{\pgfqpoint{1.733342in}{1.240411in}}%
\pgfpathlineto{\pgfqpoint{1.751185in}{1.250055in}}%
\pgfpathlineto{\pgfqpoint{1.769107in}{1.259938in}}%
\pgfpathlineto{\pgfqpoint{1.776517in}{1.270911in}}%
\pgfpathlineto{\pgfqpoint{1.783915in}{1.282058in}}%
\pgfpathlineto{\pgfqpoint{1.766017in}{1.272415in}}%
\pgfpathlineto{\pgfqpoint{1.748197in}{1.263018in}}%
\pgfpathlineto{\pgfqpoint{1.740775in}{1.251626in}}%
\pgfpathlineto{\pgfqpoint{1.733342in}{1.240411in}}%
\pgfpathclose%
\pgfusepath{fill}%
\end{pgfscope}%
\begin{pgfscope}%
\pgfpathrectangle{\pgfqpoint{0.407592in}{0.222904in}}{\pgfqpoint{1.411163in}{1.411163in}}%
\pgfusepath{clip}%
\pgfsetbuttcap%
\pgfsetroundjoin%
\definecolor{currentfill}{rgb}{0.909804,0.698039,0.423529}%
\pgfsetfillcolor{currentfill}%
\pgfsetfillopacity{0.120000}%
\pgfsetlinewidth{0.000000pt}%
\definecolor{currentstroke}{rgb}{0.000000,0.000000,0.000000}%
\pgfsetstrokecolor{currentstroke}%
\pgfsetdash{}{0pt}%
\pgfpathmoveto{\pgfqpoint{0.657740in}{1.046500in}}%
\pgfpathlineto{\pgfqpoint{0.674089in}{1.031982in}}%
\pgfpathlineto{\pgfqpoint{0.690450in}{1.017870in}}%
\pgfpathlineto{\pgfqpoint{0.699507in}{1.015984in}}%
\pgfpathlineto{\pgfqpoint{0.708521in}{1.014388in}}%
\pgfpathlineto{\pgfqpoint{0.692205in}{1.028748in}}%
\pgfpathlineto{\pgfqpoint{0.675900in}{1.043512in}}%
\pgfpathlineto{\pgfqpoint{0.666841in}{1.044863in}}%
\pgfpathlineto{\pgfqpoint{0.657740in}{1.046500in}}%
\pgfpathclose%
\pgfusepath{fill}%
\end{pgfscope}%
\begin{pgfscope}%
\pgfpathrectangle{\pgfqpoint{0.407592in}{0.222904in}}{\pgfqpoint{1.411163in}{1.411163in}}%
\pgfusepath{clip}%
\pgfsetbuttcap%
\pgfsetroundjoin%
\definecolor{currentfill}{rgb}{0.909804,0.698039,0.423529}%
\pgfsetfillcolor{currentfill}%
\pgfsetfillopacity{0.120000}%
\pgfsetlinewidth{0.000000pt}%
\definecolor{currentstroke}{rgb}{0.000000,0.000000,0.000000}%
\pgfsetstrokecolor{currentstroke}%
\pgfsetdash{}{0pt}%
\pgfpathmoveto{\pgfqpoint{0.541278in}{1.143383in}}%
\pgfpathlineto{\pgfqpoint{0.557608in}{1.127170in}}%
\pgfpathlineto{\pgfqpoint{0.573948in}{1.111229in}}%
\pgfpathlineto{\pgfqpoint{0.583222in}{1.109542in}}%
\pgfpathlineto{\pgfqpoint{0.592453in}{1.108112in}}%
\pgfpathlineto{\pgfqpoint{0.576157in}{1.124282in}}%
\pgfpathlineto{\pgfqpoint{0.559870in}{1.140715in}}%
\pgfpathlineto{\pgfqpoint{0.550596in}{1.141923in}}%
\pgfpathlineto{\pgfqpoint{0.541278in}{1.143383in}}%
\pgfpathclose%
\pgfusepath{fill}%
\end{pgfscope}%
\begin{pgfscope}%
\pgfpathrectangle{\pgfqpoint{0.407592in}{0.222904in}}{\pgfqpoint{1.411163in}{1.411163in}}%
\pgfusepath{clip}%
\pgfsetbuttcap%
\pgfsetroundjoin%
\definecolor{currentfill}{rgb}{0.909804,0.698039,0.423529}%
\pgfsetfillcolor{currentfill}%
\pgfsetfillopacity{0.120000}%
\pgfsetlinewidth{0.000000pt}%
\definecolor{currentstroke}{rgb}{0.000000,0.000000,0.000000}%
\pgfsetstrokecolor{currentstroke}%
\pgfsetdash{}{0pt}%
\pgfpathmoveto{\pgfqpoint{0.857088in}{0.913884in}}%
\pgfpathlineto{\pgfqpoint{0.873495in}{0.904484in}}%
\pgfpathlineto{\pgfqpoint{0.889924in}{0.895847in}}%
\pgfpathlineto{\pgfqpoint{0.898594in}{0.894167in}}%
\pgfpathlineto{\pgfqpoint{0.907226in}{0.892859in}}%
\pgfpathlineto{\pgfqpoint{0.890842in}{0.901641in}}%
\pgfpathlineto{\pgfqpoint{0.874480in}{0.911207in}}%
\pgfpathlineto{\pgfqpoint{0.865804in}{0.912361in}}%
\pgfpathlineto{\pgfqpoint{0.857088in}{0.913884in}}%
\pgfpathclose%
\pgfusepath{fill}%
\end{pgfscope}%
\begin{pgfscope}%
\pgfpathrectangle{\pgfqpoint{0.407592in}{0.222904in}}{\pgfqpoint{1.411163in}{1.411163in}}%
\pgfusepath{clip}%
\pgfsetbuttcap%
\pgfsetroundjoin%
\definecolor{currentfill}{rgb}{0.909804,0.698039,0.423529}%
\pgfsetfillcolor{currentfill}%
\pgfsetfillopacity{0.120000}%
\pgfsetlinewidth{0.000000pt}%
\definecolor{currentstroke}{rgb}{0.000000,0.000000,0.000000}%
\pgfsetstrokecolor{currentstroke}%
\pgfsetdash{}{0pt}%
\pgfpathmoveto{\pgfqpoint{1.552242in}{1.080028in}}%
\pgfpathlineto{\pgfqpoint{1.569738in}{1.089249in}}%
\pgfpathlineto{\pgfqpoint{1.587310in}{1.098764in}}%
\pgfpathlineto{\pgfqpoint{1.594908in}{1.108264in}}%
\pgfpathlineto{\pgfqpoint{1.602493in}{1.118023in}}%
\pgfpathlineto{\pgfqpoint{1.584943in}{1.108737in}}%
\pgfpathlineto{\pgfqpoint{1.567470in}{1.099749in}}%
\pgfpathlineto{\pgfqpoint{1.559862in}{1.089754in}}%
\pgfpathlineto{\pgfqpoint{1.552242in}{1.080028in}}%
\pgfpathclose%
\pgfusepath{fill}%
\end{pgfscope}%
\begin{pgfscope}%
\pgfpathrectangle{\pgfqpoint{0.407592in}{0.222904in}}{\pgfqpoint{1.411163in}{1.411163in}}%
\pgfusepath{clip}%
\pgfsetbuttcap%
\pgfsetroundjoin%
\definecolor{currentfill}{rgb}{0.909804,0.698039,0.423529}%
\pgfsetfillcolor{currentfill}%
\pgfsetfillopacity{0.120000}%
\pgfsetlinewidth{0.000000pt}%
\definecolor{currentstroke}{rgb}{0.000000,0.000000,0.000000}%
\pgfsetstrokecolor{currentstroke}%
\pgfsetdash{}{0pt}%
\pgfpathmoveto{\pgfqpoint{1.172180in}{0.862877in}}%
\pgfpathlineto{\pgfqpoint{1.188918in}{0.865978in}}%
\pgfpathlineto{\pgfqpoint{1.205714in}{0.869841in}}%
\pgfpathlineto{\pgfqpoint{1.213790in}{0.873020in}}%
\pgfpathlineto{\pgfqpoint{1.221842in}{0.876697in}}%
\pgfpathlineto{\pgfqpoint{1.205074in}{0.872745in}}%
\pgfpathlineto{\pgfqpoint{1.188366in}{0.869556in}}%
\pgfpathlineto{\pgfqpoint{1.180286in}{0.865957in}}%
\pgfpathlineto{\pgfqpoint{1.172180in}{0.862877in}}%
\pgfpathclose%
\pgfusepath{fill}%
\end{pgfscope}%
\begin{pgfscope}%
\pgfpathrectangle{\pgfqpoint{0.407592in}{0.222904in}}{\pgfqpoint{1.411163in}{1.411163in}}%
\pgfusepath{clip}%
\pgfsetbuttcap%
\pgfsetroundjoin%
\definecolor{currentfill}{rgb}{0.909804,0.698039,0.423529}%
\pgfsetfillcolor{currentfill}%
\pgfsetfillopacity{0.120000}%
\pgfsetlinewidth{0.000000pt}%
\definecolor{currentstroke}{rgb}{0.000000,0.000000,0.000000}%
\pgfsetstrokecolor{currentstroke}%
\pgfsetdash{}{0pt}%
\pgfpathmoveto{\pgfqpoint{0.773911in}{0.961616in}}%
\pgfpathlineto{\pgfqpoint{0.790293in}{0.949759in}}%
\pgfpathlineto{\pgfqpoint{0.806691in}{0.938508in}}%
\pgfpathlineto{\pgfqpoint{0.815531in}{0.936504in}}%
\pgfpathlineto{\pgfqpoint{0.824331in}{0.934830in}}%
\pgfpathlineto{\pgfqpoint{0.807978in}{0.946302in}}%
\pgfpathlineto{\pgfqpoint{0.791640in}{0.958388in}}%
\pgfpathlineto{\pgfqpoint{0.782796in}{0.959839in}}%
\pgfpathlineto{\pgfqpoint{0.773911in}{0.961616in}}%
\pgfpathclose%
\pgfusepath{fill}%
\end{pgfscope}%
\begin{pgfscope}%
\pgfpathrectangle{\pgfqpoint{0.407592in}{0.222904in}}{\pgfqpoint{1.411163in}{1.411163in}}%
\pgfusepath{clip}%
\pgfsetbuttcap%
\pgfsetroundjoin%
\definecolor{currentfill}{rgb}{0.909804,0.698039,0.423529}%
\pgfsetfillcolor{currentfill}%
\pgfsetfillopacity{0.120000}%
\pgfsetlinewidth{0.000000pt}%
\definecolor{currentstroke}{rgb}{0.000000,0.000000,0.000000}%
\pgfsetstrokecolor{currentstroke}%
\pgfsetdash{}{0pt}%
\pgfpathmoveto{\pgfqpoint{1.221842in}{0.876697in}}%
\pgfpathlineto{\pgfqpoint{1.238671in}{0.881363in}}%
\pgfpathlineto{\pgfqpoint{1.255564in}{0.886694in}}%
\pgfpathlineto{\pgfqpoint{1.263565in}{0.890928in}}%
\pgfpathlineto{\pgfqpoint{1.271543in}{0.895643in}}%
\pgfpathlineto{\pgfqpoint{1.254677in}{0.890272in}}%
\pgfpathlineto{\pgfqpoint{1.237874in}{0.885564in}}%
\pgfpathlineto{\pgfqpoint{1.229870in}{0.880878in}}%
\pgfpathlineto{\pgfqpoint{1.221842in}{0.876697in}}%
\pgfpathclose%
\pgfusepath{fill}%
\end{pgfscope}%
\begin{pgfscope}%
\pgfpathrectangle{\pgfqpoint{0.407592in}{0.222904in}}{\pgfqpoint{1.411163in}{1.411163in}}%
\pgfusepath{clip}%
\pgfsetbuttcap%
\pgfsetroundjoin%
\definecolor{currentfill}{rgb}{0.909804,0.698039,0.423529}%
\pgfsetfillcolor{currentfill}%
\pgfsetfillopacity{0.120000}%
\pgfsetlinewidth{0.000000pt}%
\definecolor{currentstroke}{rgb}{0.000000,0.000000,0.000000}%
\pgfsetstrokecolor{currentstroke}%
\pgfsetdash{}{0pt}%
\pgfpathmoveto{\pgfqpoint{1.122515in}{0.854750in}}%
\pgfpathlineto{\pgfqpoint{1.139175in}{0.856073in}}%
\pgfpathlineto{\pgfqpoint{1.155887in}{0.858243in}}%
\pgfpathlineto{\pgfqpoint{1.164048in}{0.860309in}}%
\pgfpathlineto{\pgfqpoint{1.172180in}{0.862877in}}%
\pgfpathlineto{\pgfqpoint{1.155500in}{0.860586in}}%
\pgfpathlineto{\pgfqpoint{1.138873in}{0.859153in}}%
\pgfpathlineto{\pgfqpoint{1.130709in}{0.856693in}}%
\pgfpathlineto{\pgfqpoint{1.122515in}{0.854750in}}%
\pgfpathclose%
\pgfusepath{fill}%
\end{pgfscope}%
\begin{pgfscope}%
\pgfpathrectangle{\pgfqpoint{0.407592in}{0.222904in}}{\pgfqpoint{1.411163in}{1.411163in}}%
\pgfusepath{clip}%
\pgfsetbuttcap%
\pgfsetroundjoin%
\definecolor{currentfill}{rgb}{0.909804,0.698039,0.423529}%
\pgfsetfillcolor{currentfill}%
\pgfsetfillopacity{0.120000}%
\pgfsetlinewidth{0.000000pt}%
\definecolor{currentstroke}{rgb}{0.000000,0.000000,0.000000}%
\pgfsetstrokecolor{currentstroke}%
\pgfsetdash{}{0pt}%
\pgfpathmoveto{\pgfqpoint{1.436741in}{0.991806in}}%
\pgfpathlineto{\pgfqpoint{1.454004in}{1.000286in}}%
\pgfpathlineto{\pgfqpoint{1.471340in}{1.009129in}}%
\pgfpathlineto{\pgfqpoint{1.479065in}{1.017223in}}%
\pgfpathlineto{\pgfqpoint{1.486776in}{1.025652in}}%
\pgfpathlineto{\pgfqpoint{1.469461in}{1.016988in}}%
\pgfpathlineto{\pgfqpoint{1.452220in}{1.008691in}}%
\pgfpathlineto{\pgfqpoint{1.444488in}{1.000073in}}%
\pgfpathlineto{\pgfqpoint{1.436741in}{0.991806in}}%
\pgfpathclose%
\pgfusepath{fill}%
\end{pgfscope}%
\begin{pgfscope}%
\pgfpathrectangle{\pgfqpoint{0.407592in}{0.222904in}}{\pgfqpoint{1.411163in}{1.411163in}}%
\pgfusepath{clip}%
\pgfsetbuttcap%
\pgfsetroundjoin%
\definecolor{currentfill}{rgb}{0.909804,0.698039,0.423529}%
\pgfsetfillcolor{currentfill}%
\pgfsetfillopacity{0.120000}%
\pgfsetlinewidth{0.000000pt}%
\definecolor{currentstroke}{rgb}{0.000000,0.000000,0.000000}%
\pgfsetstrokecolor{currentstroke}%
\pgfsetdash{}{0pt}%
\pgfpathmoveto{\pgfqpoint{1.667949in}{1.177943in}}%
\pgfpathlineto{\pgfqpoint{1.685683in}{1.187570in}}%
\pgfpathlineto{\pgfqpoint{1.703495in}{1.197448in}}%
\pgfpathlineto{\pgfqpoint{1.710974in}{1.207891in}}%
\pgfpathlineto{\pgfqpoint{1.718441in}{1.218537in}}%
\pgfpathlineto{\pgfqpoint{1.700652in}{1.208899in}}%
\pgfpathlineto{\pgfqpoint{1.682941in}{1.199517in}}%
\pgfpathlineto{\pgfqpoint{1.675450in}{1.188627in}}%
\pgfpathlineto{\pgfqpoint{1.667949in}{1.177943in}}%
\pgfpathclose%
\pgfusepath{fill}%
\end{pgfscope}%
\begin{pgfscope}%
\pgfpathrectangle{\pgfqpoint{0.407592in}{0.222904in}}{\pgfqpoint{1.411163in}{1.411163in}}%
\pgfusepath{clip}%
\pgfsetbuttcap%
\pgfsetroundjoin%
\definecolor{currentfill}{rgb}{0.909804,0.698039,0.423529}%
\pgfsetfillcolor{currentfill}%
\pgfsetfillopacity{0.120000}%
\pgfsetlinewidth{0.000000pt}%
\definecolor{currentstroke}{rgb}{0.000000,0.000000,0.000000}%
\pgfsetstrokecolor{currentstroke}%
\pgfsetdash{}{0pt}%
\pgfpathmoveto{\pgfqpoint{0.940064in}{0.877782in}}%
\pgfpathlineto{\pgfqpoint{0.956523in}{0.871542in}}%
\pgfpathlineto{\pgfqpoint{0.973011in}{0.866197in}}%
\pgfpathlineto{\pgfqpoint{0.981518in}{0.865135in}}%
\pgfpathlineto{\pgfqpoint{0.989989in}{0.864492in}}%
\pgfpathlineto{\pgfqpoint{0.973544in}{0.869864in}}%
\pgfpathlineto{\pgfqpoint{0.957128in}{0.876160in}}%
\pgfpathlineto{\pgfqpoint{0.948615in}{0.876763in}}%
\pgfpathlineto{\pgfqpoint{0.940064in}{0.877782in}}%
\pgfpathclose%
\pgfusepath{fill}%
\end{pgfscope}%
\begin{pgfscope}%
\pgfpathrectangle{\pgfqpoint{0.407592in}{0.222904in}}{\pgfqpoint{1.411163in}{1.411163in}}%
\pgfusepath{clip}%
\pgfsetbuttcap%
\pgfsetroundjoin%
\definecolor{currentfill}{rgb}{0.909804,0.698039,0.423529}%
\pgfsetfillcolor{currentfill}%
\pgfsetfillopacity{0.120000}%
\pgfsetlinewidth{0.000000pt}%
\definecolor{currentstroke}{rgb}{0.000000,0.000000,0.000000}%
\pgfsetstrokecolor{currentstroke}%
\pgfsetdash{}{0pt}%
\pgfpathmoveto{\pgfqpoint{1.271543in}{0.895643in}}%
\pgfpathlineto{\pgfqpoint{1.288476in}{0.901629in}}%
\pgfpathlineto{\pgfqpoint{1.305475in}{0.908184in}}%
\pgfpathlineto{\pgfqpoint{1.313407in}{0.913378in}}%
\pgfpathlineto{\pgfqpoint{1.321320in}{0.919025in}}%
\pgfpathlineto{\pgfqpoint{1.304345in}{0.912484in}}%
\pgfpathlineto{\pgfqpoint{1.287437in}{0.906512in}}%
\pgfpathlineto{\pgfqpoint{1.279501in}{0.900838in}}%
\pgfpathlineto{\pgfqpoint{1.271543in}{0.895643in}}%
\pgfpathclose%
\pgfusepath{fill}%
\end{pgfscope}%
\begin{pgfscope}%
\pgfpathrectangle{\pgfqpoint{0.407592in}{0.222904in}}{\pgfqpoint{1.411163in}{1.411163in}}%
\pgfusepath{clip}%
\pgfsetbuttcap%
\pgfsetroundjoin%
\definecolor{currentfill}{rgb}{0.909804,0.698039,0.423529}%
\pgfsetfillcolor{currentfill}%
\pgfsetfillopacity{0.120000}%
\pgfsetlinewidth{0.000000pt}%
\definecolor{currentstroke}{rgb}{0.000000,0.000000,0.000000}%
\pgfsetstrokecolor{currentstroke}%
\pgfsetdash{}{0pt}%
\pgfpathmoveto{\pgfqpoint{1.072799in}{0.852644in}}%
\pgfpathlineto{\pgfqpoint{1.089396in}{0.852048in}}%
\pgfpathlineto{\pgfqpoint{1.106040in}{0.852358in}}%
\pgfpathlineto{\pgfqpoint{1.114293in}{0.853311in}}%
\pgfpathlineto{\pgfqpoint{1.122515in}{0.854750in}}%
\pgfpathlineto{\pgfqpoint{1.105907in}{0.854317in}}%
\pgfpathlineto{\pgfqpoint{1.089346in}{0.854809in}}%
\pgfpathlineto{\pgfqpoint{1.081088in}{0.853479in}}%
\pgfpathlineto{\pgfqpoint{1.072799in}{0.852644in}}%
\pgfpathclose%
\pgfusepath{fill}%
\end{pgfscope}%
\begin{pgfscope}%
\pgfpathrectangle{\pgfqpoint{0.407592in}{0.222904in}}{\pgfqpoint{1.411163in}{1.411163in}}%
\pgfusepath{clip}%
\pgfsetbuttcap%
\pgfsetroundjoin%
\definecolor{currentfill}{rgb}{0.909804,0.698039,0.423529}%
\pgfsetfillcolor{currentfill}%
\pgfsetfillopacity{0.120000}%
\pgfsetlinewidth{0.000000pt}%
\definecolor{currentstroke}{rgb}{0.000000,0.000000,0.000000}%
\pgfsetstrokecolor{currentstroke}%
\pgfsetdash{}{0pt}%
\pgfpathmoveto{\pgfqpoint{0.573948in}{1.111229in}}%
\pgfpathlineto{\pgfqpoint{0.590297in}{1.095579in}}%
\pgfpathlineto{\pgfqpoint{0.606658in}{1.080241in}}%
\pgfpathlineto{\pgfqpoint{0.615888in}{1.078313in}}%
\pgfpathlineto{\pgfqpoint{0.625076in}{1.076648in}}%
\pgfpathlineto{\pgfqpoint{0.608759in}{1.092227in}}%
\pgfpathlineto{\pgfqpoint{0.592453in}{1.108112in}}%
\pgfpathlineto{\pgfqpoint{0.583222in}{1.109542in}}%
\pgfpathlineto{\pgfqpoint{0.573948in}{1.111229in}}%
\pgfpathclose%
\pgfusepath{fill}%
\end{pgfscope}%
\begin{pgfscope}%
\pgfpathrectangle{\pgfqpoint{0.407592in}{0.222904in}}{\pgfqpoint{1.411163in}{1.411163in}}%
\pgfusepath{clip}%
\pgfsetbuttcap%
\pgfsetroundjoin%
\definecolor{currentfill}{rgb}{0.909804,0.698039,0.423529}%
\pgfsetfillcolor{currentfill}%
\pgfsetfillopacity{0.120000}%
\pgfsetlinewidth{0.000000pt}%
\definecolor{currentstroke}{rgb}{0.000000,0.000000,0.000000}%
\pgfsetstrokecolor{currentstroke}%
\pgfsetdash{}{0pt}%
\pgfpathmoveto{\pgfqpoint{0.690450in}{1.017870in}}%
\pgfpathlineto{\pgfqpoint{0.706822in}{1.004195in}}%
\pgfpathlineto{\pgfqpoint{0.723207in}{0.990986in}}%
\pgfpathlineto{\pgfqpoint{0.732219in}{0.988846in}}%
\pgfpathlineto{\pgfqpoint{0.741190in}{0.987002in}}%
\pgfpathlineto{\pgfqpoint{0.724850in}{1.000462in}}%
\pgfpathlineto{\pgfqpoint{0.708521in}{1.014388in}}%
\pgfpathlineto{\pgfqpoint{0.699507in}{1.015984in}}%
\pgfpathlineto{\pgfqpoint{0.690450in}{1.017870in}}%
\pgfpathclose%
\pgfusepath{fill}%
\end{pgfscope}%
\begin{pgfscope}%
\pgfpathrectangle{\pgfqpoint{0.407592in}{0.222904in}}{\pgfqpoint{1.411163in}{1.411163in}}%
\pgfusepath{clip}%
\pgfsetbuttcap%
\pgfsetroundjoin%
\definecolor{currentfill}{rgb}{0.909804,0.698039,0.423529}%
\pgfsetfillcolor{currentfill}%
\pgfsetfillopacity{0.120000}%
\pgfsetlinewidth{0.000000pt}%
\definecolor{currentstroke}{rgb}{0.000000,0.000000,0.000000}%
\pgfsetstrokecolor{currentstroke}%
\pgfsetdash{}{0pt}%
\pgfpathmoveto{\pgfqpoint{1.321320in}{0.919025in}}%
\pgfpathlineto{\pgfqpoint{1.338364in}{0.926091in}}%
\pgfpathlineto{\pgfqpoint{1.355477in}{0.933643in}}%
\pgfpathlineto{\pgfqpoint{1.363347in}{0.939685in}}%
\pgfpathlineto{\pgfqpoint{1.371199in}{0.946146in}}%
\pgfpathlineto{\pgfqpoint{1.354108in}{0.938661in}}%
\pgfpathlineto{\pgfqpoint{1.337087in}{0.931661in}}%
\pgfpathlineto{\pgfqpoint{1.329213in}{0.925121in}}%
\pgfpathlineto{\pgfqpoint{1.321320in}{0.919025in}}%
\pgfpathclose%
\pgfusepath{fill}%
\end{pgfscope}%
\begin{pgfscope}%
\pgfpathrectangle{\pgfqpoint{0.407592in}{0.222904in}}{\pgfqpoint{1.411163in}{1.411163in}}%
\pgfusepath{clip}%
\pgfsetbuttcap%
\pgfsetroundjoin%
\definecolor{currentfill}{rgb}{0.909804,0.698039,0.423529}%
\pgfsetfillcolor{currentfill}%
\pgfsetfillopacity{0.120000}%
\pgfsetlinewidth{0.000000pt}%
\definecolor{currentstroke}{rgb}{0.000000,0.000000,0.000000}%
\pgfsetstrokecolor{currentstroke}%
\pgfsetdash{}{0pt}%
\pgfpathmoveto{\pgfqpoint{1.602493in}{1.118023in}}%
\pgfpathlineto{\pgfqpoint{1.620118in}{1.127592in}}%
\pgfpathlineto{\pgfqpoint{1.637821in}{1.137428in}}%
\pgfpathlineto{\pgfqpoint{1.645371in}{1.147209in}}%
\pgfpathlineto{\pgfqpoint{1.652909in}{1.157226in}}%
\pgfpathlineto{\pgfqpoint{1.635229in}{1.147622in}}%
\pgfpathlineto{\pgfqpoint{1.617626in}{1.138290in}}%
\pgfpathlineto{\pgfqpoint{1.610065in}{1.128035in}}%
\pgfpathlineto{\pgfqpoint{1.602493in}{1.118023in}}%
\pgfpathclose%
\pgfusepath{fill}%
\end{pgfscope}%
\begin{pgfscope}%
\pgfpathrectangle{\pgfqpoint{0.407592in}{0.222904in}}{\pgfqpoint{1.411163in}{1.411163in}}%
\pgfusepath{clip}%
\pgfsetbuttcap%
\pgfsetroundjoin%
\definecolor{currentfill}{rgb}{0.909804,0.698039,0.423529}%
\pgfsetfillcolor{currentfill}%
\pgfsetfillopacity{0.120000}%
\pgfsetlinewidth{0.000000pt}%
\definecolor{currentstroke}{rgb}{0.000000,0.000000,0.000000}%
\pgfsetstrokecolor{currentstroke}%
\pgfsetdash{}{0pt}%
\pgfpathmoveto{\pgfqpoint{1.486776in}{1.025652in}}%
\pgfpathlineto{\pgfqpoint{1.504165in}{1.034660in}}%
\pgfpathlineto{\pgfqpoint{1.521629in}{1.043989in}}%
\pgfpathlineto{\pgfqpoint{1.529303in}{1.052551in}}%
\pgfpathlineto{\pgfqpoint{1.536963in}{1.061417in}}%
\pgfpathlineto{\pgfqpoint{1.519521in}{1.052286in}}%
\pgfpathlineto{\pgfqpoint{1.502153in}{1.043481in}}%
\pgfpathlineto{\pgfqpoint{1.494472in}{1.034408in}}%
\pgfpathlineto{\pgfqpoint{1.486776in}{1.025652in}}%
\pgfpathclose%
\pgfusepath{fill}%
\end{pgfscope}%
\begin{pgfscope}%
\pgfpathrectangle{\pgfqpoint{0.407592in}{0.222904in}}{\pgfqpoint{1.411163in}{1.411163in}}%
\pgfusepath{clip}%
\pgfsetbuttcap%
\pgfsetroundjoin%
\definecolor{currentfill}{rgb}{0.909804,0.698039,0.423529}%
\pgfsetfillcolor{currentfill}%
\pgfsetfillopacity{0.120000}%
\pgfsetlinewidth{0.000000pt}%
\definecolor{currentstroke}{rgb}{0.000000,0.000000,0.000000}%
\pgfsetstrokecolor{currentstroke}%
\pgfsetdash{}{0pt}%
\pgfpathmoveto{\pgfqpoint{0.457113in}{1.213391in}}%
\pgfpathlineto{\pgfqpoint{0.473448in}{1.196468in}}%
\pgfpathlineto{\pgfqpoint{0.489793in}{1.179757in}}%
\pgfpathlineto{\pgfqpoint{0.499242in}{1.178035in}}%
\pgfpathlineto{\pgfqpoint{0.508647in}{1.176548in}}%
\pgfpathlineto{\pgfqpoint{0.492346in}{1.193469in}}%
\pgfpathlineto{\pgfqpoint{0.476054in}{1.210594in}}%
\pgfpathlineto{\pgfqpoint{0.466606in}{1.211878in}}%
\pgfpathlineto{\pgfqpoint{0.457113in}{1.213391in}}%
\pgfpathclose%
\pgfusepath{fill}%
\end{pgfscope}%
\begin{pgfscope}%
\pgfpathrectangle{\pgfqpoint{0.407592in}{0.222904in}}{\pgfqpoint{1.411163in}{1.411163in}}%
\pgfusepath{clip}%
\pgfsetbuttcap%
\pgfsetroundjoin%
\definecolor{currentfill}{rgb}{0.909804,0.698039,0.423529}%
\pgfsetfillcolor{currentfill}%
\pgfsetfillopacity{0.120000}%
\pgfsetlinewidth{0.000000pt}%
\definecolor{currentstroke}{rgb}{0.000000,0.000000,0.000000}%
\pgfsetstrokecolor{currentstroke}%
\pgfsetdash{}{0pt}%
\pgfpathmoveto{\pgfqpoint{1.783915in}{1.282058in}}%
\pgfpathlineto{\pgfqpoint{1.801892in}{1.291935in}}%
\pgfpathlineto{\pgfqpoint{1.819947in}{1.302038in}}%
\pgfpathlineto{\pgfqpoint{1.827310in}{1.313119in}}%
\pgfpathlineto{\pgfqpoint{1.809267in}{1.303133in}}%
\pgfpathlineto{\pgfqpoint{1.791303in}{1.293374in}}%
\pgfusepath{fill}%
\end{pgfscope}%
\begin{pgfscope}%
\pgfpathrectangle{\pgfqpoint{0.407592in}{0.222904in}}{\pgfqpoint{1.411163in}{1.411163in}}%
\pgfusepath{clip}%
\pgfsetbuttcap%
\pgfsetroundjoin%
\definecolor{currentfill}{rgb}{0.909804,0.698039,0.423529}%
\pgfsetfillcolor{currentfill}%
\pgfsetfillopacity{0.120000}%
\pgfsetlinewidth{0.000000pt}%
\definecolor{currentstroke}{rgb}{0.000000,0.000000,0.000000}%
\pgfsetstrokecolor{currentstroke}%
\pgfsetdash{}{0pt}%
\pgfpathmoveto{\pgfqpoint{1.022980in}{0.856565in}}%
\pgfpathlineto{\pgfqpoint{1.039531in}{0.854013in}}%
\pgfpathlineto{\pgfqpoint{1.056122in}{0.852389in}}%
\pgfpathlineto{\pgfqpoint{1.064477in}{0.852288in}}%
\pgfpathlineto{\pgfqpoint{1.072799in}{0.852644in}}%
\pgfpathlineto{\pgfqpoint{1.056246in}{0.854176in}}%
\pgfpathlineto{\pgfqpoint{1.039735in}{0.856663in}}%
\pgfpathlineto{\pgfqpoint{1.031374in}{0.856383in}}%
\pgfpathlineto{\pgfqpoint{1.022980in}{0.856565in}}%
\pgfpathclose%
\pgfusepath{fill}%
\end{pgfscope}%
\begin{pgfscope}%
\pgfpathrectangle{\pgfqpoint{0.407592in}{0.222904in}}{\pgfqpoint{1.411163in}{1.411163in}}%
\pgfusepath{clip}%
\pgfsetbuttcap%
\pgfsetroundjoin%
\definecolor{currentfill}{rgb}{0.909804,0.698039,0.423529}%
\pgfsetfillcolor{currentfill}%
\pgfsetfillopacity{0.120000}%
\pgfsetlinewidth{0.000000pt}%
\definecolor{currentstroke}{rgb}{0.000000,0.000000,0.000000}%
\pgfsetstrokecolor{currentstroke}%
\pgfsetdash{}{0pt}%
\pgfpathmoveto{\pgfqpoint{1.718441in}{1.218537in}}%
\pgfpathlineto{\pgfqpoint{1.736308in}{1.228421in}}%
\pgfpathlineto{\pgfqpoint{1.754253in}{1.238540in}}%
\pgfpathlineto{\pgfqpoint{1.761685in}{1.249145in}}%
\pgfpathlineto{\pgfqpoint{1.769107in}{1.259938in}}%
\pgfpathlineto{\pgfqpoint{1.751185in}{1.250055in}}%
\pgfpathlineto{\pgfqpoint{1.733342in}{1.240411in}}%
\pgfpathlineto{\pgfqpoint{1.725897in}{1.229379in}}%
\pgfpathlineto{\pgfqpoint{1.718441in}{1.218537in}}%
\pgfpathclose%
\pgfusepath{fill}%
\end{pgfscope}%
\begin{pgfscope}%
\pgfpathrectangle{\pgfqpoint{0.407592in}{0.222904in}}{\pgfqpoint{1.411163in}{1.411163in}}%
\pgfusepath{clip}%
\pgfsetbuttcap%
\pgfsetroundjoin%
\definecolor{currentfill}{rgb}{0.909804,0.698039,0.423529}%
\pgfsetfillcolor{currentfill}%
\pgfsetfillopacity{0.120000}%
\pgfsetlinewidth{0.000000pt}%
\definecolor{currentstroke}{rgb}{0.000000,0.000000,0.000000}%
\pgfsetstrokecolor{currentstroke}%
\pgfsetdash{}{0pt}%
\pgfpathmoveto{\pgfqpoint{0.806691in}{0.938508in}}%
\pgfpathlineto{\pgfqpoint{0.823106in}{0.927897in}}%
\pgfpathlineto{\pgfqpoint{0.839539in}{0.917962in}}%
\pgfpathlineto{\pgfqpoint{0.848333in}{0.915756in}}%
\pgfpathlineto{\pgfqpoint{0.857088in}{0.913884in}}%
\pgfpathlineto{\pgfqpoint{0.840701in}{0.924011in}}%
\pgfpathlineto{\pgfqpoint{0.824331in}{0.934830in}}%
\pgfpathlineto{\pgfqpoint{0.815531in}{0.936504in}}%
\pgfpathlineto{\pgfqpoint{0.806691in}{0.938508in}}%
\pgfpathclose%
\pgfusepath{fill}%
\end{pgfscope}%
\begin{pgfscope}%
\pgfpathrectangle{\pgfqpoint{0.407592in}{0.222904in}}{\pgfqpoint{1.411163in}{1.411163in}}%
\pgfusepath{clip}%
\pgfsetbuttcap%
\pgfsetroundjoin%
\definecolor{currentfill}{rgb}{0.909804,0.698039,0.423529}%
\pgfsetfillcolor{currentfill}%
\pgfsetfillopacity{0.120000}%
\pgfsetlinewidth{0.000000pt}%
\definecolor{currentstroke}{rgb}{0.000000,0.000000,0.000000}%
\pgfsetstrokecolor{currentstroke}%
\pgfsetdash{}{0pt}%
\pgfpathmoveto{\pgfqpoint{0.889924in}{0.895847in}}%
\pgfpathlineto{\pgfqpoint{0.906375in}{0.888005in}}%
\pgfpathlineto{\pgfqpoint{0.922850in}{0.880983in}}%
\pgfpathlineto{\pgfqpoint{0.931476in}{0.879195in}}%
\pgfpathlineto{\pgfqpoint{0.940064in}{0.877782in}}%
\pgfpathlineto{\pgfqpoint{0.923633in}{0.884895in}}%
\pgfpathlineto{\pgfqpoint{0.907226in}{0.892859in}}%
\pgfpathlineto{\pgfqpoint{0.898594in}{0.894167in}}%
\pgfpathlineto{\pgfqpoint{0.889924in}{0.895847in}}%
\pgfpathclose%
\pgfusepath{fill}%
\end{pgfscope}%
\begin{pgfscope}%
\pgfpathrectangle{\pgfqpoint{0.407592in}{0.222904in}}{\pgfqpoint{1.411163in}{1.411163in}}%
\pgfusepath{clip}%
\pgfsetbuttcap%
\pgfsetroundjoin%
\definecolor{currentfill}{rgb}{0.909804,0.698039,0.423529}%
\pgfsetfillcolor{currentfill}%
\pgfsetfillopacity{0.120000}%
\pgfsetlinewidth{0.000000pt}%
\definecolor{currentstroke}{rgb}{0.000000,0.000000,0.000000}%
\pgfsetstrokecolor{currentstroke}%
\pgfsetdash{}{0pt}%
\pgfpathmoveto{\pgfqpoint{1.371199in}{0.946146in}}%
\pgfpathlineto{\pgfqpoint{1.388361in}{0.954080in}}%
\pgfpathlineto{\pgfqpoint{1.405595in}{0.962430in}}%
\pgfpathlineto{\pgfqpoint{1.413406in}{0.969204in}}%
\pgfpathlineto{\pgfqpoint{1.421200in}{0.976363in}}%
\pgfpathlineto{\pgfqpoint{1.403989in}{0.968123in}}%
\pgfpathlineto{\pgfqpoint{1.386849in}{0.960301in}}%
\pgfpathlineto{\pgfqpoint{1.379032in}{0.953020in}}%
\pgfpathlineto{\pgfqpoint{1.371199in}{0.946146in}}%
\pgfpathclose%
\pgfusepath{fill}%
\end{pgfscope}%
\begin{pgfscope}%
\pgfpathrectangle{\pgfqpoint{0.407592in}{0.222904in}}{\pgfqpoint{1.411163in}{1.411163in}}%
\pgfusepath{clip}%
\pgfsetbuttcap%
\pgfsetroundjoin%
\definecolor{currentfill}{rgb}{0.909804,0.698039,0.423529}%
\pgfsetfillcolor{currentfill}%
\pgfsetfillopacity{0.120000}%
\pgfsetlinewidth{0.000000pt}%
\definecolor{currentstroke}{rgb}{0.000000,0.000000,0.000000}%
\pgfsetstrokecolor{currentstroke}%
\pgfsetdash{}{0pt}%
\pgfpathmoveto{\pgfqpoint{0.606658in}{1.080241in}}%
\pgfpathlineto{\pgfqpoint{0.623028in}{1.065238in}}%
\pgfpathlineto{\pgfqpoint{0.639410in}{1.050592in}}%
\pgfpathlineto{\pgfqpoint{0.648597in}{1.048412in}}%
\pgfpathlineto{\pgfqpoint{0.657740in}{1.046500in}}%
\pgfpathlineto{\pgfqpoint{0.641403in}{1.061397in}}%
\pgfpathlineto{\pgfqpoint{0.625076in}{1.076648in}}%
\pgfpathlineto{\pgfqpoint{0.615888in}{1.078313in}}%
\pgfpathlineto{\pgfqpoint{0.606658in}{1.080241in}}%
\pgfpathclose%
\pgfusepath{fill}%
\end{pgfscope}%
\begin{pgfscope}%
\pgfpathrectangle{\pgfqpoint{0.407592in}{0.222904in}}{\pgfqpoint{1.411163in}{1.411163in}}%
\pgfusepath{clip}%
\pgfsetbuttcap%
\pgfsetroundjoin%
\definecolor{currentfill}{rgb}{0.909804,0.698039,0.423529}%
\pgfsetfillcolor{currentfill}%
\pgfsetfillopacity{0.120000}%
\pgfsetlinewidth{0.000000pt}%
\definecolor{currentstroke}{rgb}{0.000000,0.000000,0.000000}%
\pgfsetstrokecolor{currentstroke}%
\pgfsetdash{}{0pt}%
\pgfpathmoveto{\pgfqpoint{0.489793in}{1.179757in}}%
\pgfpathlineto{\pgfqpoint{0.506147in}{1.163273in}}%
\pgfpathlineto{\pgfqpoint{0.522510in}{1.147032in}}%
\pgfpathlineto{\pgfqpoint{0.531916in}{1.145088in}}%
\pgfpathlineto{\pgfqpoint{0.541278in}{1.143383in}}%
\pgfpathlineto{\pgfqpoint{0.524958in}{1.159848in}}%
\pgfpathlineto{\pgfqpoint{0.508647in}{1.176548in}}%
\pgfpathlineto{\pgfqpoint{0.499242in}{1.178035in}}%
\pgfpathlineto{\pgfqpoint{0.489793in}{1.179757in}}%
\pgfpathclose%
\pgfusepath{fill}%
\end{pgfscope}%
\begin{pgfscope}%
\pgfpathrectangle{\pgfqpoint{0.407592in}{0.222904in}}{\pgfqpoint{1.411163in}{1.411163in}}%
\pgfusepath{clip}%
\pgfsetbuttcap%
\pgfsetroundjoin%
\definecolor{currentfill}{rgb}{0.909804,0.698039,0.423529}%
\pgfsetfillcolor{currentfill}%
\pgfsetfillopacity{0.120000}%
\pgfsetlinewidth{0.000000pt}%
\definecolor{currentstroke}{rgb}{0.000000,0.000000,0.000000}%
\pgfsetstrokecolor{currentstroke}%
\pgfsetdash{}{0pt}%
\pgfpathmoveto{\pgfqpoint{0.723207in}{0.990986in}}%
\pgfpathlineto{\pgfqpoint{0.739605in}{0.978274in}}%
\pgfpathlineto{\pgfqpoint{0.756017in}{0.966091in}}%
\pgfpathlineto{\pgfqpoint{0.764985in}{0.963704in}}%
\pgfpathlineto{\pgfqpoint{0.773911in}{0.961616in}}%
\pgfpathlineto{\pgfqpoint{0.757544in}{0.974042in}}%
\pgfpathlineto{\pgfqpoint{0.741190in}{0.987002in}}%
\pgfpathlineto{\pgfqpoint{0.732219in}{0.988846in}}%
\pgfpathlineto{\pgfqpoint{0.723207in}{0.990986in}}%
\pgfpathclose%
\pgfusepath{fill}%
\end{pgfscope}%
\begin{pgfscope}%
\pgfpathrectangle{\pgfqpoint{0.407592in}{0.222904in}}{\pgfqpoint{1.411163in}{1.411163in}}%
\pgfusepath{clip}%
\pgfsetbuttcap%
\pgfsetroundjoin%
\definecolor{currentfill}{rgb}{0.909804,0.698039,0.423529}%
\pgfsetfillcolor{currentfill}%
\pgfsetfillopacity{0.120000}%
\pgfsetlinewidth{0.000000pt}%
\definecolor{currentstroke}{rgb}{0.000000,0.000000,0.000000}%
\pgfsetstrokecolor{currentstroke}%
\pgfsetdash{}{0pt}%
\pgfpathmoveto{\pgfqpoint{1.536963in}{1.061417in}}%
\pgfpathlineto{\pgfqpoint{1.554480in}{1.070854in}}%
\pgfpathlineto{\pgfqpoint{1.572074in}{1.080578in}}%
\pgfpathlineto{\pgfqpoint{1.579698in}{1.089533in}}%
\pgfpathlineto{\pgfqpoint{1.587310in}{1.098764in}}%
\pgfpathlineto{\pgfqpoint{1.569738in}{1.089249in}}%
\pgfpathlineto{\pgfqpoint{1.552242in}{1.080028in}}%
\pgfpathlineto{\pgfqpoint{1.544609in}{1.070579in}}%
\pgfpathlineto{\pgfqpoint{1.536963in}{1.061417in}}%
\pgfpathclose%
\pgfusepath{fill}%
\end{pgfscope}%
\begin{pgfscope}%
\pgfpathrectangle{\pgfqpoint{0.407592in}{0.222904in}}{\pgfqpoint{1.411163in}{1.411163in}}%
\pgfusepath{clip}%
\pgfsetbuttcap%
\pgfsetroundjoin%
\definecolor{currentfill}{rgb}{0.909804,0.698039,0.423529}%
\pgfsetfillcolor{currentfill}%
\pgfsetfillopacity{0.120000}%
\pgfsetlinewidth{0.000000pt}%
\definecolor{currentstroke}{rgb}{0.000000,0.000000,0.000000}%
\pgfsetstrokecolor{currentstroke}%
\pgfsetdash{}{0pt}%
\pgfpathmoveto{\pgfqpoint{0.973011in}{0.866197in}}%
\pgfpathlineto{\pgfqpoint{0.989531in}{0.861757in}}%
\pgfpathlineto{\pgfqpoint{1.006086in}{0.858228in}}%
\pgfpathlineto{\pgfqpoint{1.014551in}{0.857187in}}%
\pgfpathlineto{\pgfqpoint{1.022980in}{0.856565in}}%
\pgfpathlineto{\pgfqpoint{1.006467in}{0.860057in}}%
\pgfpathlineto{\pgfqpoint{0.989989in}{0.864492in}}%
\pgfpathlineto{\pgfqpoint{0.981518in}{0.865135in}}%
\pgfpathlineto{\pgfqpoint{0.973011in}{0.866197in}}%
\pgfpathclose%
\pgfusepath{fill}%
\end{pgfscope}%
\begin{pgfscope}%
\pgfpathrectangle{\pgfqpoint{0.407592in}{0.222904in}}{\pgfqpoint{1.411163in}{1.411163in}}%
\pgfusepath{clip}%
\pgfsetbuttcap%
\pgfsetroundjoin%
\definecolor{currentfill}{rgb}{0.909804,0.698039,0.423529}%
\pgfsetfillcolor{currentfill}%
\pgfsetfillopacity{0.120000}%
\pgfsetlinewidth{0.000000pt}%
\definecolor{currentstroke}{rgb}{0.000000,0.000000,0.000000}%
\pgfsetstrokecolor{currentstroke}%
\pgfsetdash{}{0pt}%
\pgfpathmoveto{\pgfqpoint{1.652909in}{1.157226in}}%
\pgfpathlineto{\pgfqpoint{1.670666in}{1.167090in}}%
\pgfpathlineto{\pgfqpoint{1.688501in}{1.177199in}}%
\pgfpathlineto{\pgfqpoint{1.696004in}{1.187215in}}%
\pgfpathlineto{\pgfqpoint{1.703495in}{1.197448in}}%
\pgfpathlineto{\pgfqpoint{1.685683in}{1.187570in}}%
\pgfpathlineto{\pgfqpoint{1.667949in}{1.177943in}}%
\pgfpathlineto{\pgfqpoint{1.660435in}{1.167474in}}%
\pgfpathlineto{\pgfqpoint{1.652909in}{1.157226in}}%
\pgfpathclose%
\pgfusepath{fill}%
\end{pgfscope}%
\begin{pgfscope}%
\pgfpathrectangle{\pgfqpoint{0.407592in}{0.222904in}}{\pgfqpoint{1.411163in}{1.411163in}}%
\pgfusepath{clip}%
\pgfsetbuttcap%
\pgfsetroundjoin%
\definecolor{currentfill}{rgb}{0.909804,0.698039,0.423529}%
\pgfsetfillcolor{currentfill}%
\pgfsetfillopacity{0.120000}%
\pgfsetlinewidth{0.000000pt}%
\definecolor{currentstroke}{rgb}{0.000000,0.000000,0.000000}%
\pgfsetstrokecolor{currentstroke}%
\pgfsetdash{}{0pt}%
\pgfpathmoveto{\pgfqpoint{1.421200in}{0.976363in}}%
\pgfpathlineto{\pgfqpoint{1.438485in}{0.984991in}}%
\pgfpathlineto{\pgfqpoint{1.455844in}{0.993978in}}%
\pgfpathlineto{\pgfqpoint{1.463600in}{1.001378in}}%
\pgfpathlineto{\pgfqpoint{1.471340in}{1.009129in}}%
\pgfpathlineto{\pgfqpoint{1.454004in}{1.000286in}}%
\pgfpathlineto{\pgfqpoint{1.436741in}{0.991806in}}%
\pgfpathlineto{\pgfqpoint{1.428979in}{0.983900in}}%
\pgfpathlineto{\pgfqpoint{1.421200in}{0.976363in}}%
\pgfpathclose%
\pgfusepath{fill}%
\end{pgfscope}%
\begin{pgfscope}%
\pgfpathrectangle{\pgfqpoint{0.407592in}{0.222904in}}{\pgfqpoint{1.411163in}{1.411163in}}%
\pgfusepath{clip}%
\pgfsetbuttcap%
\pgfsetroundjoin%
\definecolor{currentfill}{rgb}{0.909804,0.698039,0.423529}%
\pgfsetfillcolor{currentfill}%
\pgfsetfillopacity{0.120000}%
\pgfsetlinewidth{0.000000pt}%
\definecolor{currentstroke}{rgb}{0.000000,0.000000,0.000000}%
\pgfsetstrokecolor{currentstroke}%
\pgfsetdash{}{0pt}%
\pgfpathmoveto{\pgfqpoint{1.769107in}{1.259938in}}%
\pgfpathlineto{\pgfqpoint{1.787107in}{1.270050in}}%
\pgfpathlineto{\pgfqpoint{1.805186in}{1.280383in}}%
\pgfpathlineto{\pgfqpoint{1.812572in}{1.291124in}}%
\pgfpathlineto{\pgfqpoint{1.819947in}{1.302038in}}%
\pgfpathlineto{\pgfqpoint{1.801892in}{1.291935in}}%
\pgfpathlineto{\pgfqpoint{1.783915in}{1.282058in}}%
\pgfpathlineto{\pgfqpoint{1.776517in}{1.270911in}}%
\pgfpathlineto{\pgfqpoint{1.769107in}{1.259938in}}%
\pgfpathclose%
\pgfusepath{fill}%
\end{pgfscope}%
\begin{pgfscope}%
\pgfpathrectangle{\pgfqpoint{0.407592in}{0.222904in}}{\pgfqpoint{1.411163in}{1.411163in}}%
\pgfusepath{clip}%
\pgfsetbuttcap%
\pgfsetroundjoin%
\definecolor{currentfill}{rgb}{0.909804,0.698039,0.423529}%
\pgfsetfillcolor{currentfill}%
\pgfsetfillopacity{0.120000}%
\pgfsetlinewidth{0.000000pt}%
\definecolor{currentstroke}{rgb}{0.000000,0.000000,0.000000}%
\pgfsetstrokecolor{currentstroke}%
\pgfsetdash{}{0pt}%
\pgfpathmoveto{\pgfqpoint{1.155887in}{0.858243in}}%
\pgfpathlineto{\pgfqpoint{1.172656in}{0.861216in}}%
\pgfpathlineto{\pgfqpoint{1.189483in}{0.864945in}}%
\pgfpathlineto{\pgfqpoint{1.197611in}{0.867152in}}%
\pgfpathlineto{\pgfqpoint{1.205714in}{0.869841in}}%
\pgfpathlineto{\pgfqpoint{1.188918in}{0.865978in}}%
\pgfpathlineto{\pgfqpoint{1.172180in}{0.862877in}}%
\pgfpathlineto{\pgfqpoint{1.164048in}{0.860309in}}%
\pgfpathlineto{\pgfqpoint{1.155887in}{0.858243in}}%
\pgfpathclose%
\pgfusepath{fill}%
\end{pgfscope}%
\begin{pgfscope}%
\pgfpathrectangle{\pgfqpoint{0.407592in}{0.222904in}}{\pgfqpoint{1.411163in}{1.411163in}}%
\pgfusepath{clip}%
\pgfsetbuttcap%
\pgfsetroundjoin%
\definecolor{currentfill}{rgb}{0.909804,0.698039,0.423529}%
\pgfsetfillcolor{currentfill}%
\pgfsetfillopacity{0.120000}%
\pgfsetlinewidth{0.000000pt}%
\definecolor{currentstroke}{rgb}{0.000000,0.000000,0.000000}%
\pgfsetstrokecolor{currentstroke}%
\pgfsetdash{}{0pt}%
\pgfpathmoveto{\pgfqpoint{1.205714in}{0.869841in}}%
\pgfpathlineto{\pgfqpoint{1.222571in}{0.874416in}}%
\pgfpathlineto{\pgfqpoint{1.239491in}{0.879655in}}%
\pgfpathlineto{\pgfqpoint{1.247539in}{0.882937in}}%
\pgfpathlineto{\pgfqpoint{1.255564in}{0.886694in}}%
\pgfpathlineto{\pgfqpoint{1.238671in}{0.881363in}}%
\pgfpathlineto{\pgfqpoint{1.221842in}{0.876697in}}%
\pgfpathlineto{\pgfqpoint{1.213790in}{0.873020in}}%
\pgfpathlineto{\pgfqpoint{1.205714in}{0.869841in}}%
\pgfpathclose%
\pgfusepath{fill}%
\end{pgfscope}%
\begin{pgfscope}%
\pgfpathrectangle{\pgfqpoint{0.407592in}{0.222904in}}{\pgfqpoint{1.411163in}{1.411163in}}%
\pgfusepath{clip}%
\pgfsetbuttcap%
\pgfsetroundjoin%
\definecolor{currentfill}{rgb}{0.909804,0.698039,0.423529}%
\pgfsetfillcolor{currentfill}%
\pgfsetfillopacity{0.120000}%
\pgfsetlinewidth{0.000000pt}%
\definecolor{currentstroke}{rgb}{0.000000,0.000000,0.000000}%
\pgfsetstrokecolor{currentstroke}%
\pgfsetdash{}{0pt}%
\pgfpathmoveto{\pgfqpoint{0.639410in}{1.050592in}}%
\pgfpathlineto{\pgfqpoint{0.655803in}{1.036328in}}%
\pgfpathlineto{\pgfqpoint{0.672208in}{1.022474in}}%
\pgfpathlineto{\pgfqpoint{0.681350in}{1.020037in}}%
\pgfpathlineto{\pgfqpoint{0.690450in}{1.017870in}}%
\pgfpathlineto{\pgfqpoint{0.674089in}{1.031982in}}%
\pgfpathlineto{\pgfqpoint{0.657740in}{1.046500in}}%
\pgfpathlineto{\pgfqpoint{0.648597in}{1.048412in}}%
\pgfpathlineto{\pgfqpoint{0.639410in}{1.050592in}}%
\pgfpathclose%
\pgfusepath{fill}%
\end{pgfscope}%
\begin{pgfscope}%
\pgfpathrectangle{\pgfqpoint{0.407592in}{0.222904in}}{\pgfqpoint{1.411163in}{1.411163in}}%
\pgfusepath{clip}%
\pgfsetbuttcap%
\pgfsetroundjoin%
\definecolor{currentfill}{rgb}{0.909804,0.698039,0.423529}%
\pgfsetfillcolor{currentfill}%
\pgfsetfillopacity{0.120000}%
\pgfsetlinewidth{0.000000pt}%
\definecolor{currentstroke}{rgb}{0.000000,0.000000,0.000000}%
\pgfsetstrokecolor{currentstroke}%
\pgfsetdash{}{0pt}%
\pgfpathmoveto{\pgfqpoint{0.522510in}{1.147032in}}%
\pgfpathlineto{\pgfqpoint{0.538884in}{1.131051in}}%
\pgfpathlineto{\pgfqpoint{0.555267in}{1.115347in}}%
\pgfpathlineto{\pgfqpoint{0.564629in}{1.113167in}}%
\pgfpathlineto{\pgfqpoint{0.573948in}{1.111229in}}%
\pgfpathlineto{\pgfqpoint{0.557608in}{1.127170in}}%
\pgfpathlineto{\pgfqpoint{0.541278in}{1.143383in}}%
\pgfpathlineto{\pgfqpoint{0.531916in}{1.145088in}}%
\pgfpathlineto{\pgfqpoint{0.522510in}{1.147032in}}%
\pgfpathclose%
\pgfusepath{fill}%
\end{pgfscope}%
\begin{pgfscope}%
\pgfpathrectangle{\pgfqpoint{0.407592in}{0.222904in}}{\pgfqpoint{1.411163in}{1.411163in}}%
\pgfusepath{clip}%
\pgfsetbuttcap%
\pgfsetroundjoin%
\definecolor{currentfill}{rgb}{0.909804,0.698039,0.423529}%
\pgfsetfillcolor{currentfill}%
\pgfsetfillopacity{0.120000}%
\pgfsetlinewidth{0.000000pt}%
\definecolor{currentstroke}{rgb}{0.000000,0.000000,0.000000}%
\pgfsetstrokecolor{currentstroke}%
\pgfsetdash{}{0pt}%
\pgfpathmoveto{\pgfqpoint{1.106040in}{0.852358in}}%
\pgfpathlineto{\pgfqpoint{1.122734in}{0.853542in}}%
\pgfpathlineto{\pgfqpoint{1.139481in}{0.855560in}}%
\pgfpathlineto{\pgfqpoint{1.147699in}{0.856664in}}%
\pgfpathlineto{\pgfqpoint{1.155887in}{0.858243in}}%
\pgfpathlineto{\pgfqpoint{1.139175in}{0.856073in}}%
\pgfpathlineto{\pgfqpoint{1.122515in}{0.854750in}}%
\pgfpathlineto{\pgfqpoint{1.114293in}{0.853311in}}%
\pgfpathlineto{\pgfqpoint{1.106040in}{0.852358in}}%
\pgfpathclose%
\pgfusepath{fill}%
\end{pgfscope}%
\begin{pgfscope}%
\pgfpathrectangle{\pgfqpoint{0.407592in}{0.222904in}}{\pgfqpoint{1.411163in}{1.411163in}}%
\pgfusepath{clip}%
\pgfsetbuttcap%
\pgfsetroundjoin%
\definecolor{currentfill}{rgb}{0.909804,0.698039,0.423529}%
\pgfsetfillcolor{currentfill}%
\pgfsetfillopacity{0.120000}%
\pgfsetlinewidth{0.000000pt}%
\definecolor{currentstroke}{rgb}{0.000000,0.000000,0.000000}%
\pgfsetstrokecolor{currentstroke}%
\pgfsetdash{}{0pt}%
\pgfpathmoveto{\pgfqpoint{0.839539in}{0.917962in}}%
\pgfpathlineto{\pgfqpoint{0.855991in}{0.908735in}}%
\pgfpathlineto{\pgfqpoint{0.872465in}{0.900249in}}%
\pgfpathlineto{\pgfqpoint{0.881214in}{0.897881in}}%
\pgfpathlineto{\pgfqpoint{0.889924in}{0.895847in}}%
\pgfpathlineto{\pgfqpoint{0.873495in}{0.904484in}}%
\pgfpathlineto{\pgfqpoint{0.857088in}{0.913884in}}%
\pgfpathlineto{\pgfqpoint{0.848333in}{0.915756in}}%
\pgfpathlineto{\pgfqpoint{0.839539in}{0.917962in}}%
\pgfpathclose%
\pgfusepath{fill}%
\end{pgfscope}%
\begin{pgfscope}%
\pgfpathrectangle{\pgfqpoint{0.407592in}{0.222904in}}{\pgfqpoint{1.411163in}{1.411163in}}%
\pgfusepath{clip}%
\pgfsetbuttcap%
\pgfsetroundjoin%
\definecolor{currentfill}{rgb}{0.909804,0.698039,0.423529}%
\pgfsetfillcolor{currentfill}%
\pgfsetfillopacity{0.120000}%
\pgfsetlinewidth{0.000000pt}%
\definecolor{currentstroke}{rgb}{0.000000,0.000000,0.000000}%
\pgfsetstrokecolor{currentstroke}%
\pgfsetdash{}{0pt}%
\pgfpathmoveto{\pgfqpoint{1.255564in}{0.886694in}}%
\pgfpathlineto{\pgfqpoint{1.272521in}{0.892640in}}%
\pgfpathlineto{\pgfqpoint{1.289546in}{0.899158in}}%
\pgfpathlineto{\pgfqpoint{1.297521in}{0.903444in}}%
\pgfpathlineto{\pgfqpoint{1.305475in}{0.908184in}}%
\pgfpathlineto{\pgfqpoint{1.288476in}{0.901629in}}%
\pgfpathlineto{\pgfqpoint{1.271543in}{0.895643in}}%
\pgfpathlineto{\pgfqpoint{1.263565in}{0.890928in}}%
\pgfpathlineto{\pgfqpoint{1.255564in}{0.886694in}}%
\pgfpathclose%
\pgfusepath{fill}%
\end{pgfscope}%
\begin{pgfscope}%
\pgfpathrectangle{\pgfqpoint{0.407592in}{0.222904in}}{\pgfqpoint{1.411163in}{1.411163in}}%
\pgfusepath{clip}%
\pgfsetbuttcap%
\pgfsetroundjoin%
\definecolor{currentfill}{rgb}{0.909804,0.698039,0.423529}%
\pgfsetfillcolor{currentfill}%
\pgfsetfillopacity{0.120000}%
\pgfsetlinewidth{0.000000pt}%
\definecolor{currentstroke}{rgb}{0.000000,0.000000,0.000000}%
\pgfsetstrokecolor{currentstroke}%
\pgfsetdash{}{0pt}%
\pgfpathmoveto{\pgfqpoint{1.587310in}{1.098764in}}%
\pgfpathlineto{\pgfqpoint{1.604958in}{1.108555in}}%
\pgfpathlineto{\pgfqpoint{1.622683in}{1.118607in}}%
\pgfpathlineto{\pgfqpoint{1.630258in}{1.127892in}}%
\pgfpathlineto{\pgfqpoint{1.637821in}{1.137428in}}%
\pgfpathlineto{\pgfqpoint{1.620118in}{1.127592in}}%
\pgfpathlineto{\pgfqpoint{1.602493in}{1.118023in}}%
\pgfpathlineto{\pgfqpoint{1.594908in}{1.108264in}}%
\pgfpathlineto{\pgfqpoint{1.587310in}{1.098764in}}%
\pgfpathclose%
\pgfusepath{fill}%
\end{pgfscope}%
\begin{pgfscope}%
\pgfpathrectangle{\pgfqpoint{0.407592in}{0.222904in}}{\pgfqpoint{1.411163in}{1.411163in}}%
\pgfusepath{clip}%
\pgfsetbuttcap%
\pgfsetroundjoin%
\definecolor{currentfill}{rgb}{0.909804,0.698039,0.423529}%
\pgfsetfillcolor{currentfill}%
\pgfsetfillopacity{0.120000}%
\pgfsetlinewidth{0.000000pt}%
\definecolor{currentstroke}{rgb}{0.000000,0.000000,0.000000}%
\pgfsetstrokecolor{currentstroke}%
\pgfsetdash{}{0pt}%
\pgfpathmoveto{\pgfqpoint{0.756017in}{0.966091in}}%
\pgfpathlineto{\pgfqpoint{0.772445in}{0.954470in}}%
\pgfpathlineto{\pgfqpoint{0.788888in}{0.943444in}}%
\pgfpathlineto{\pgfqpoint{0.797810in}{0.940826in}}%
\pgfpathlineto{\pgfqpoint{0.806691in}{0.938508in}}%
\pgfpathlineto{\pgfqpoint{0.790293in}{0.949759in}}%
\pgfpathlineto{\pgfqpoint{0.773911in}{0.961616in}}%
\pgfpathlineto{\pgfqpoint{0.764985in}{0.963704in}}%
\pgfpathlineto{\pgfqpoint{0.756017in}{0.966091in}}%
\pgfpathclose%
\pgfusepath{fill}%
\end{pgfscope}%
\begin{pgfscope}%
\pgfpathrectangle{\pgfqpoint{0.407592in}{0.222904in}}{\pgfqpoint{1.411163in}{1.411163in}}%
\pgfusepath{clip}%
\pgfsetbuttcap%
\pgfsetroundjoin%
\definecolor{currentfill}{rgb}{0.909804,0.698039,0.423529}%
\pgfsetfillcolor{currentfill}%
\pgfsetfillopacity{0.120000}%
\pgfsetlinewidth{0.000000pt}%
\definecolor{currentstroke}{rgb}{0.000000,0.000000,0.000000}%
\pgfsetstrokecolor{currentstroke}%
\pgfsetdash{}{0pt}%
\pgfpathmoveto{\pgfqpoint{0.922850in}{0.880983in}}%
\pgfpathlineto{\pgfqpoint{0.939353in}{0.874807in}}%
\pgfpathlineto{\pgfqpoint{0.955886in}{0.869491in}}%
\pgfpathlineto{\pgfqpoint{0.964467in}{0.867656in}}%
\pgfpathlineto{\pgfqpoint{0.973011in}{0.866197in}}%
\pgfpathlineto{\pgfqpoint{0.956523in}{0.871542in}}%
\pgfpathlineto{\pgfqpoint{0.940064in}{0.877782in}}%
\pgfpathlineto{\pgfqpoint{0.931476in}{0.879195in}}%
\pgfpathlineto{\pgfqpoint{0.922850in}{0.880983in}}%
\pgfpathclose%
\pgfusepath{fill}%
\end{pgfscope}%
\begin{pgfscope}%
\pgfpathrectangle{\pgfqpoint{0.407592in}{0.222904in}}{\pgfqpoint{1.411163in}{1.411163in}}%
\pgfusepath{clip}%
\pgfsetbuttcap%
\pgfsetroundjoin%
\definecolor{currentfill}{rgb}{0.909804,0.698039,0.423529}%
\pgfsetfillcolor{currentfill}%
\pgfsetfillopacity{0.120000}%
\pgfsetlinewidth{0.000000pt}%
\definecolor{currentstroke}{rgb}{0.000000,0.000000,0.000000}%
\pgfsetstrokecolor{currentstroke}%
\pgfsetdash{}{0pt}%
\pgfpathmoveto{\pgfqpoint{1.471340in}{1.009129in}}%
\pgfpathlineto{\pgfqpoint{1.488751in}{1.018310in}}%
\pgfpathlineto{\pgfqpoint{1.506237in}{1.027808in}}%
\pgfpathlineto{\pgfqpoint{1.513940in}{1.035738in}}%
\pgfpathlineto{\pgfqpoint{1.521629in}{1.043989in}}%
\pgfpathlineto{\pgfqpoint{1.504165in}{1.034660in}}%
\pgfpathlineto{\pgfqpoint{1.486776in}{1.025652in}}%
\pgfpathlineto{\pgfqpoint{1.479065in}{1.017223in}}%
\pgfpathlineto{\pgfqpoint{1.471340in}{1.009129in}}%
\pgfpathclose%
\pgfusepath{fill}%
\end{pgfscope}%
\begin{pgfscope}%
\pgfpathrectangle{\pgfqpoint{0.407592in}{0.222904in}}{\pgfqpoint{1.411163in}{1.411163in}}%
\pgfusepath{clip}%
\pgfsetbuttcap%
\pgfsetroundjoin%
\definecolor{currentfill}{rgb}{0.909804,0.698039,0.423529}%
\pgfsetfillcolor{currentfill}%
\pgfsetfillopacity{0.120000}%
\pgfsetlinewidth{0.000000pt}%
\definecolor{currentstroke}{rgb}{0.000000,0.000000,0.000000}%
\pgfsetstrokecolor{currentstroke}%
\pgfsetdash{}{0pt}%
\pgfpathmoveto{\pgfqpoint{1.703495in}{1.197448in}}%
\pgfpathlineto{\pgfqpoint{1.721385in}{1.207567in}}%
\pgfpathlineto{\pgfqpoint{1.739354in}{1.217914in}}%
\pgfpathlineto{\pgfqpoint{1.746809in}{1.228127in}}%
\pgfpathlineto{\pgfqpoint{1.754253in}{1.238540in}}%
\pgfpathlineto{\pgfqpoint{1.736308in}{1.228421in}}%
\pgfpathlineto{\pgfqpoint{1.718441in}{1.218537in}}%
\pgfpathlineto{\pgfqpoint{1.710974in}{1.207891in}}%
\pgfpathlineto{\pgfqpoint{1.703495in}{1.197448in}}%
\pgfpathclose%
\pgfusepath{fill}%
\end{pgfscope}%
\begin{pgfscope}%
\pgfpathrectangle{\pgfqpoint{0.407592in}{0.222904in}}{\pgfqpoint{1.411163in}{1.411163in}}%
\pgfusepath{clip}%
\pgfsetbuttcap%
\pgfsetroundjoin%
\definecolor{currentfill}{rgb}{0.909804,0.698039,0.423529}%
\pgfsetfillcolor{currentfill}%
\pgfsetfillopacity{0.120000}%
\pgfsetlinewidth{0.000000pt}%
\definecolor{currentstroke}{rgb}{0.000000,0.000000,0.000000}%
\pgfsetstrokecolor{currentstroke}%
\pgfsetdash{}{0pt}%
\pgfpathmoveto{\pgfqpoint{1.305475in}{0.908184in}}%
\pgfpathlineto{\pgfqpoint{1.322543in}{0.915266in}}%
\pgfpathlineto{\pgfqpoint{1.339680in}{0.922834in}}%
\pgfpathlineto{\pgfqpoint{1.347589in}{0.928025in}}%
\pgfpathlineto{\pgfqpoint{1.355477in}{0.933643in}}%
\pgfpathlineto{\pgfqpoint{1.338364in}{0.926091in}}%
\pgfpathlineto{\pgfqpoint{1.321320in}{0.919025in}}%
\pgfpathlineto{\pgfqpoint{1.313407in}{0.913378in}}%
\pgfpathlineto{\pgfqpoint{1.305475in}{0.908184in}}%
\pgfpathclose%
\pgfusepath{fill}%
\end{pgfscope}%
\begin{pgfscope}%
\pgfpathrectangle{\pgfqpoint{0.407592in}{0.222904in}}{\pgfqpoint{1.411163in}{1.411163in}}%
\pgfusepath{clip}%
\pgfsetbuttcap%
\pgfsetroundjoin%
\definecolor{currentfill}{rgb}{0.909804,0.698039,0.423529}%
\pgfsetfillcolor{currentfill}%
\pgfsetfillopacity{0.120000}%
\pgfsetlinewidth{0.000000pt}%
\definecolor{currentstroke}{rgb}{0.000000,0.000000,0.000000}%
\pgfsetstrokecolor{currentstroke}%
\pgfsetdash{}{0pt}%
\pgfpathmoveto{\pgfqpoint{1.819947in}{1.302038in}}%
\pgfpathlineto{\pgfqpoint{1.838081in}{1.312357in}}%
\pgfpathlineto{\pgfqpoint{1.845433in}{1.323323in}}%
\pgfpathlineto{\pgfqpoint{1.827310in}{1.313119in}}%
\pgfusepath{fill}%
\end{pgfscope}%
\begin{pgfscope}%
\pgfpathrectangle{\pgfqpoint{0.407592in}{0.222904in}}{\pgfqpoint{1.411163in}{1.411163in}}%
\pgfusepath{clip}%
\pgfsetbuttcap%
\pgfsetroundjoin%
\definecolor{currentfill}{rgb}{0.909804,0.698039,0.423529}%
\pgfsetfillcolor{currentfill}%
\pgfsetfillopacity{0.120000}%
\pgfsetlinewidth{0.000000pt}%
\definecolor{currentstroke}{rgb}{0.000000,0.000000,0.000000}%
\pgfsetstrokecolor{currentstroke}%
\pgfsetdash{}{0pt}%
\pgfpathmoveto{\pgfqpoint{1.056122in}{0.852389in}}%
\pgfpathlineto{\pgfqpoint{1.072758in}{0.851675in}}%
\pgfpathlineto{\pgfqpoint{1.089440in}{0.851845in}}%
\pgfpathlineto{\pgfqpoint{1.097756in}{0.851876in}}%
\pgfpathlineto{\pgfqpoint{1.106040in}{0.852358in}}%
\pgfpathlineto{\pgfqpoint{1.089396in}{0.852048in}}%
\pgfpathlineto{\pgfqpoint{1.072799in}{0.852644in}}%
\pgfpathlineto{\pgfqpoint{1.064477in}{0.852288in}}%
\pgfpathlineto{\pgfqpoint{1.056122in}{0.852389in}}%
\pgfpathclose%
\pgfusepath{fill}%
\end{pgfscope}%
\begin{pgfscope}%
\pgfpathrectangle{\pgfqpoint{0.407592in}{0.222904in}}{\pgfqpoint{1.411163in}{1.411163in}}%
\pgfusepath{clip}%
\pgfsetbuttcap%
\pgfsetroundjoin%
\definecolor{currentfill}{rgb}{0.909804,0.698039,0.423529}%
\pgfsetfillcolor{currentfill}%
\pgfsetfillopacity{0.120000}%
\pgfsetlinewidth{0.000000pt}%
\definecolor{currentstroke}{rgb}{0.000000,0.000000,0.000000}%
\pgfsetstrokecolor{currentstroke}%
\pgfsetdash{}{0pt}%
\pgfpathmoveto{\pgfqpoint{0.555267in}{1.115347in}}%
\pgfpathlineto{\pgfqpoint{0.571661in}{1.099941in}}%
\pgfpathlineto{\pgfqpoint{0.588065in}{1.084851in}}%
\pgfpathlineto{\pgfqpoint{0.597383in}{1.082424in}}%
\pgfpathlineto{\pgfqpoint{0.606658in}{1.080241in}}%
\pgfpathlineto{\pgfqpoint{0.590297in}{1.095579in}}%
\pgfpathlineto{\pgfqpoint{0.573948in}{1.111229in}}%
\pgfpathlineto{\pgfqpoint{0.564629in}{1.113167in}}%
\pgfpathlineto{\pgfqpoint{0.555267in}{1.115347in}}%
\pgfpathclose%
\pgfusepath{fill}%
\end{pgfscope}%
\begin{pgfscope}%
\pgfpathrectangle{\pgfqpoint{0.407592in}{0.222904in}}{\pgfqpoint{1.411163in}{1.411163in}}%
\pgfusepath{clip}%
\pgfsetbuttcap%
\pgfsetroundjoin%
\definecolor{currentfill}{rgb}{0.909804,0.698039,0.423529}%
\pgfsetfillcolor{currentfill}%
\pgfsetfillopacity{0.120000}%
\pgfsetlinewidth{0.000000pt}%
\definecolor{currentstroke}{rgb}{0.000000,0.000000,0.000000}%
\pgfsetstrokecolor{currentstroke}%
\pgfsetdash{}{0pt}%
\pgfpathmoveto{\pgfqpoint{0.672208in}{1.022474in}}%
\pgfpathlineto{\pgfqpoint{0.688625in}{1.009055in}}%
\pgfpathlineto{\pgfqpoint{0.705055in}{0.996102in}}%
\pgfpathlineto{\pgfqpoint{0.714152in}{0.993408in}}%
\pgfpathlineto{\pgfqpoint{0.723207in}{0.990986in}}%
\pgfpathlineto{\pgfqpoint{0.706822in}{1.004195in}}%
\pgfpathlineto{\pgfqpoint{0.690450in}{1.017870in}}%
\pgfpathlineto{\pgfqpoint{0.681350in}{1.020037in}}%
\pgfpathlineto{\pgfqpoint{0.672208in}{1.022474in}}%
\pgfpathclose%
\pgfusepath{fill}%
\end{pgfscope}%
\begin{pgfscope}%
\pgfpathrectangle{\pgfqpoint{0.407592in}{0.222904in}}{\pgfqpoint{1.411163in}{1.411163in}}%
\pgfusepath{clip}%
\pgfsetbuttcap%
\pgfsetroundjoin%
\definecolor{currentfill}{rgb}{0.909804,0.698039,0.423529}%
\pgfsetfillcolor{currentfill}%
\pgfsetfillopacity{0.120000}%
\pgfsetlinewidth{0.000000pt}%
\definecolor{currentstroke}{rgb}{0.000000,0.000000,0.000000}%
\pgfsetstrokecolor{currentstroke}%
\pgfsetdash{}{0pt}%
\pgfpathmoveto{\pgfqpoint{1.355477in}{0.933643in}}%
\pgfpathlineto{\pgfqpoint{1.372662in}{0.941644in}}%
\pgfpathlineto{\pgfqpoint{1.389919in}{0.950059in}}%
\pgfpathlineto{\pgfqpoint{1.397766in}{0.956047in}}%
\pgfpathlineto{\pgfqpoint{1.405595in}{0.962430in}}%
\pgfpathlineto{\pgfqpoint{1.388361in}{0.954080in}}%
\pgfpathlineto{\pgfqpoint{1.371199in}{0.946146in}}%
\pgfpathlineto{\pgfqpoint{1.363347in}{0.939685in}}%
\pgfpathlineto{\pgfqpoint{1.355477in}{0.933643in}}%
\pgfpathclose%
\pgfusepath{fill}%
\end{pgfscope}%
\begin{pgfscope}%
\pgfpathrectangle{\pgfqpoint{0.407592in}{0.222904in}}{\pgfqpoint{1.411163in}{1.411163in}}%
\pgfusepath{clip}%
\pgfsetbuttcap%
\pgfsetroundjoin%
\definecolor{currentfill}{rgb}{0.909804,0.698039,0.423529}%
\pgfsetfillcolor{currentfill}%
\pgfsetfillopacity{0.120000}%
\pgfsetlinewidth{0.000000pt}%
\definecolor{currentstroke}{rgb}{0.000000,0.000000,0.000000}%
\pgfsetstrokecolor{currentstroke}%
\pgfsetdash{}{0pt}%
\pgfpathmoveto{\pgfqpoint{0.437993in}{1.217086in}}%
\pgfpathlineto{\pgfqpoint{0.454371in}{1.200373in}}%
\pgfpathlineto{\pgfqpoint{0.470759in}{1.183880in}}%
\pgfpathlineto{\pgfqpoint{0.480298in}{1.181707in}}%
\pgfpathlineto{\pgfqpoint{0.489793in}{1.179757in}}%
\pgfpathlineto{\pgfqpoint{0.473448in}{1.196468in}}%
\pgfpathlineto{\pgfqpoint{0.457113in}{1.213391in}}%
\pgfpathlineto{\pgfqpoint{0.447576in}{1.215129in}}%
\pgfpathlineto{\pgfqpoint{0.437993in}{1.217086in}}%
\pgfpathclose%
\pgfusepath{fill}%
\end{pgfscope}%
\begin{pgfscope}%
\pgfpathrectangle{\pgfqpoint{0.407592in}{0.222904in}}{\pgfqpoint{1.411163in}{1.411163in}}%
\pgfusepath{clip}%
\pgfsetbuttcap%
\pgfsetroundjoin%
\definecolor{currentfill}{rgb}{0.909804,0.698039,0.423529}%
\pgfsetfillcolor{currentfill}%
\pgfsetfillopacity{0.120000}%
\pgfsetlinewidth{0.000000pt}%
\definecolor{currentstroke}{rgb}{0.000000,0.000000,0.000000}%
\pgfsetstrokecolor{currentstroke}%
\pgfsetdash{}{0pt}%
\pgfpathmoveto{\pgfqpoint{1.006086in}{0.858228in}}%
\pgfpathlineto{\pgfqpoint{1.022679in}{0.855607in}}%
\pgfpathlineto{\pgfqpoint{1.039311in}{0.853885in}}%
\pgfpathlineto{\pgfqpoint{1.047734in}{0.852928in}}%
\pgfpathlineto{\pgfqpoint{1.056122in}{0.852389in}}%
\pgfpathlineto{\pgfqpoint{1.039531in}{0.854013in}}%
\pgfpathlineto{\pgfqpoint{1.022980in}{0.856565in}}%
\pgfpathlineto{\pgfqpoint{1.014551in}{0.857187in}}%
\pgfpathlineto{\pgfqpoint{1.006086in}{0.858228in}}%
\pgfpathclose%
\pgfusepath{fill}%
\end{pgfscope}%
\begin{pgfscope}%
\pgfpathrectangle{\pgfqpoint{0.407592in}{0.222904in}}{\pgfqpoint{1.411163in}{1.411163in}}%
\pgfusepath{clip}%
\pgfsetbuttcap%
\pgfsetroundjoin%
\definecolor{currentfill}{rgb}{0.909804,0.698039,0.423529}%
\pgfsetfillcolor{currentfill}%
\pgfsetfillopacity{0.120000}%
\pgfsetlinewidth{0.000000pt}%
\definecolor{currentstroke}{rgb}{0.000000,0.000000,0.000000}%
\pgfsetstrokecolor{currentstroke}%
\pgfsetdash{}{0pt}%
\pgfpathmoveto{\pgfqpoint{1.637821in}{1.137428in}}%
\pgfpathlineto{\pgfqpoint{1.655601in}{1.147516in}}%
\pgfpathlineto{\pgfqpoint{1.673460in}{1.157845in}}%
\pgfpathlineto{\pgfqpoint{1.680987in}{1.167407in}}%
\pgfpathlineto{\pgfqpoint{1.688501in}{1.177199in}}%
\pgfpathlineto{\pgfqpoint{1.670666in}{1.167090in}}%
\pgfpathlineto{\pgfqpoint{1.652909in}{1.157226in}}%
\pgfpathlineto{\pgfqpoint{1.645371in}{1.147209in}}%
\pgfpathlineto{\pgfqpoint{1.637821in}{1.137428in}}%
\pgfpathclose%
\pgfusepath{fill}%
\end{pgfscope}%
\begin{pgfscope}%
\pgfpathrectangle{\pgfqpoint{0.407592in}{0.222904in}}{\pgfqpoint{1.411163in}{1.411163in}}%
\pgfusepath{clip}%
\pgfsetbuttcap%
\pgfsetroundjoin%
\definecolor{currentfill}{rgb}{0.909804,0.698039,0.423529}%
\pgfsetfillcolor{currentfill}%
\pgfsetfillopacity{0.120000}%
\pgfsetlinewidth{0.000000pt}%
\definecolor{currentstroke}{rgb}{0.000000,0.000000,0.000000}%
\pgfsetstrokecolor{currentstroke}%
\pgfsetdash{}{0pt}%
\pgfpathmoveto{\pgfqpoint{1.521629in}{1.043989in}}%
\pgfpathlineto{\pgfqpoint{1.539168in}{1.053617in}}%
\pgfpathlineto{\pgfqpoint{1.556785in}{1.063528in}}%
\pgfpathlineto{\pgfqpoint{1.564436in}{1.071907in}}%
\pgfpathlineto{\pgfqpoint{1.572074in}{1.080578in}}%
\pgfpathlineto{\pgfqpoint{1.554480in}{1.070854in}}%
\pgfpathlineto{\pgfqpoint{1.536963in}{1.061417in}}%
\pgfpathlineto{\pgfqpoint{1.529303in}{1.052551in}}%
\pgfpathlineto{\pgfqpoint{1.521629in}{1.043989in}}%
\pgfpathclose%
\pgfusepath{fill}%
\end{pgfscope}%
\begin{pgfscope}%
\pgfpathrectangle{\pgfqpoint{0.407592in}{0.222904in}}{\pgfqpoint{1.411163in}{1.411163in}}%
\pgfusepath{clip}%
\pgfsetbuttcap%
\pgfsetroundjoin%
\definecolor{currentfill}{rgb}{0.909804,0.698039,0.423529}%
\pgfsetfillcolor{currentfill}%
\pgfsetfillopacity{0.120000}%
\pgfsetlinewidth{0.000000pt}%
\definecolor{currentstroke}{rgb}{0.000000,0.000000,0.000000}%
\pgfsetstrokecolor{currentstroke}%
\pgfsetdash{}{0pt}%
\pgfpathmoveto{\pgfqpoint{1.754253in}{1.238540in}}%
\pgfpathlineto{\pgfqpoint{1.772277in}{1.248883in}}%
\pgfpathlineto{\pgfqpoint{1.790381in}{1.259440in}}%
\pgfpathlineto{\pgfqpoint{1.797789in}{1.269820in}}%
\pgfpathlineto{\pgfqpoint{1.805186in}{1.280383in}}%
\pgfpathlineto{\pgfqpoint{1.787107in}{1.270050in}}%
\pgfpathlineto{\pgfqpoint{1.769107in}{1.259938in}}%
\pgfpathlineto{\pgfqpoint{1.761685in}{1.249145in}}%
\pgfpathlineto{\pgfqpoint{1.754253in}{1.238540in}}%
\pgfpathclose%
\pgfusepath{fill}%
\end{pgfscope}%
\begin{pgfscope}%
\pgfpathrectangle{\pgfqpoint{0.407592in}{0.222904in}}{\pgfqpoint{1.411163in}{1.411163in}}%
\pgfusepath{clip}%
\pgfsetbuttcap%
\pgfsetroundjoin%
\definecolor{currentfill}{rgb}{0.909804,0.698039,0.423529}%
\pgfsetfillcolor{currentfill}%
\pgfsetfillopacity{0.120000}%
\pgfsetlinewidth{0.000000pt}%
\definecolor{currentstroke}{rgb}{0.000000,0.000000,0.000000}%
\pgfsetstrokecolor{currentstroke}%
\pgfsetdash{}{0pt}%
\pgfpathmoveto{\pgfqpoint{0.872465in}{0.900249in}}%
\pgfpathlineto{\pgfqpoint{0.888961in}{0.892531in}}%
\pgfpathlineto{\pgfqpoint{0.905483in}{0.885606in}}%
\pgfpathlineto{\pgfqpoint{0.914186in}{0.883127in}}%
\pgfpathlineto{\pgfqpoint{0.922850in}{0.880983in}}%
\pgfpathlineto{\pgfqpoint{0.906375in}{0.888005in}}%
\pgfpathlineto{\pgfqpoint{0.889924in}{0.895847in}}%
\pgfpathlineto{\pgfqpoint{0.881214in}{0.897881in}}%
\pgfpathlineto{\pgfqpoint{0.872465in}{0.900249in}}%
\pgfpathclose%
\pgfusepath{fill}%
\end{pgfscope}%
\begin{pgfscope}%
\pgfpathrectangle{\pgfqpoint{0.407592in}{0.222904in}}{\pgfqpoint{1.411163in}{1.411163in}}%
\pgfusepath{clip}%
\pgfsetbuttcap%
\pgfsetroundjoin%
\definecolor{currentfill}{rgb}{0.909804,0.698039,0.423529}%
\pgfsetfillcolor{currentfill}%
\pgfsetfillopacity{0.120000}%
\pgfsetlinewidth{0.000000pt}%
\definecolor{currentstroke}{rgb}{0.000000,0.000000,0.000000}%
\pgfsetstrokecolor{currentstroke}%
\pgfsetdash{}{0pt}%
\pgfpathmoveto{\pgfqpoint{0.788888in}{0.943444in}}%
\pgfpathlineto{\pgfqpoint{0.805348in}{0.933045in}}%
\pgfpathlineto{\pgfqpoint{0.821827in}{0.923305in}}%
\pgfpathlineto{\pgfqpoint{0.830703in}{0.920483in}}%
\pgfpathlineto{\pgfqpoint{0.839539in}{0.917962in}}%
\pgfpathlineto{\pgfqpoint{0.823106in}{0.927897in}}%
\pgfpathlineto{\pgfqpoint{0.806691in}{0.938508in}}%
\pgfpathlineto{\pgfqpoint{0.797810in}{0.940826in}}%
\pgfpathlineto{\pgfqpoint{0.788888in}{0.943444in}}%
\pgfpathclose%
\pgfusepath{fill}%
\end{pgfscope}%
\begin{pgfscope}%
\pgfpathrectangle{\pgfqpoint{0.407592in}{0.222904in}}{\pgfqpoint{1.411163in}{1.411163in}}%
\pgfusepath{clip}%
\pgfsetbuttcap%
\pgfsetroundjoin%
\definecolor{currentfill}{rgb}{0.909804,0.698039,0.423529}%
\pgfsetfillcolor{currentfill}%
\pgfsetfillopacity{0.120000}%
\pgfsetlinewidth{0.000000pt}%
\definecolor{currentstroke}{rgb}{0.000000,0.000000,0.000000}%
\pgfsetstrokecolor{currentstroke}%
\pgfsetdash{}{0pt}%
\pgfpathmoveto{\pgfqpoint{1.405595in}{0.962430in}}%
\pgfpathlineto{\pgfqpoint{1.422902in}{0.971166in}}%
\pgfpathlineto{\pgfqpoint{1.440283in}{0.980259in}}%
\pgfpathlineto{\pgfqpoint{1.448072in}{0.986936in}}%
\pgfpathlineto{\pgfqpoint{1.455844in}{0.993978in}}%
\pgfpathlineto{\pgfqpoint{1.438485in}{0.984991in}}%
\pgfpathlineto{\pgfqpoint{1.421200in}{0.976363in}}%
\pgfpathlineto{\pgfqpoint{1.413406in}{0.969204in}}%
\pgfpathlineto{\pgfqpoint{1.405595in}{0.962430in}}%
\pgfpathclose%
\pgfusepath{fill}%
\end{pgfscope}%
\begin{pgfscope}%
\pgfpathrectangle{\pgfqpoint{0.407592in}{0.222904in}}{\pgfqpoint{1.411163in}{1.411163in}}%
\pgfusepath{clip}%
\pgfsetbuttcap%
\pgfsetroundjoin%
\definecolor{currentfill}{rgb}{0.909804,0.698039,0.423529}%
\pgfsetfillcolor{currentfill}%
\pgfsetfillopacity{0.120000}%
\pgfsetlinewidth{0.000000pt}%
\definecolor{currentstroke}{rgb}{0.000000,0.000000,0.000000}%
\pgfsetstrokecolor{currentstroke}%
\pgfsetdash{}{0pt}%
\pgfpathmoveto{\pgfqpoint{0.588065in}{1.084851in}}%
\pgfpathlineto{\pgfqpoint{0.604480in}{1.070101in}}%
\pgfpathlineto{\pgfqpoint{0.620906in}{1.055712in}}%
\pgfpathlineto{\pgfqpoint{0.630180in}{1.053028in}}%
\pgfpathlineto{\pgfqpoint{0.639410in}{1.050592in}}%
\pgfpathlineto{\pgfqpoint{0.623028in}{1.065238in}}%
\pgfpathlineto{\pgfqpoint{0.606658in}{1.080241in}}%
\pgfpathlineto{\pgfqpoint{0.597383in}{1.082424in}}%
\pgfpathlineto{\pgfqpoint{0.588065in}{1.084851in}}%
\pgfpathclose%
\pgfusepath{fill}%
\end{pgfscope}%
\begin{pgfscope}%
\pgfpathrectangle{\pgfqpoint{0.407592in}{0.222904in}}{\pgfqpoint{1.411163in}{1.411163in}}%
\pgfusepath{clip}%
\pgfsetbuttcap%
\pgfsetroundjoin%
\definecolor{currentfill}{rgb}{0.909804,0.698039,0.423529}%
\pgfsetfillcolor{currentfill}%
\pgfsetfillopacity{0.120000}%
\pgfsetlinewidth{0.000000pt}%
\definecolor{currentstroke}{rgb}{0.000000,0.000000,0.000000}%
\pgfsetstrokecolor{currentstroke}%
\pgfsetdash{}{0pt}%
\pgfpathmoveto{\pgfqpoint{1.805186in}{1.280383in}}%
\pgfpathlineto{\pgfqpoint{1.823345in}{1.290927in}}%
\pgfpathlineto{\pgfqpoint{1.830719in}{1.301556in}}%
\pgfpathlineto{\pgfqpoint{1.838081in}{1.312357in}}%
\pgfpathlineto{\pgfqpoint{1.819947in}{1.302038in}}%
\pgfpathlineto{\pgfqpoint{1.812572in}{1.291124in}}%
\pgfusepath{fill}%
\end{pgfscope}%
\begin{pgfscope}%
\pgfpathrectangle{\pgfqpoint{0.407592in}{0.222904in}}{\pgfqpoint{1.411163in}{1.411163in}}%
\pgfusepath{clip}%
\pgfsetbuttcap%
\pgfsetroundjoin%
\definecolor{currentfill}{rgb}{0.909804,0.698039,0.423529}%
\pgfsetfillcolor{currentfill}%
\pgfsetfillopacity{0.120000}%
\pgfsetlinewidth{0.000000pt}%
\definecolor{currentstroke}{rgb}{0.000000,0.000000,0.000000}%
\pgfsetstrokecolor{currentstroke}%
\pgfsetdash{}{0pt}%
\pgfpathmoveto{\pgfqpoint{0.470759in}{1.183880in}}%
\pgfpathlineto{\pgfqpoint{0.487157in}{1.167621in}}%
\pgfpathlineto{\pgfqpoint{0.503564in}{1.151611in}}%
\pgfpathlineto{\pgfqpoint{0.513059in}{1.149208in}}%
\pgfpathlineto{\pgfqpoint{0.522510in}{1.147032in}}%
\pgfpathlineto{\pgfqpoint{0.506147in}{1.163273in}}%
\pgfpathlineto{\pgfqpoint{0.489793in}{1.179757in}}%
\pgfpathlineto{\pgfqpoint{0.480298in}{1.181707in}}%
\pgfpathlineto{\pgfqpoint{0.470759in}{1.183880in}}%
\pgfpathclose%
\pgfusepath{fill}%
\end{pgfscope}%
\begin{pgfscope}%
\pgfpathrectangle{\pgfqpoint{0.407592in}{0.222904in}}{\pgfqpoint{1.411163in}{1.411163in}}%
\pgfusepath{clip}%
\pgfsetbuttcap%
\pgfsetroundjoin%
\definecolor{currentfill}{rgb}{0.909804,0.698039,0.423529}%
\pgfsetfillcolor{currentfill}%
\pgfsetfillopacity{0.120000}%
\pgfsetlinewidth{0.000000pt}%
\definecolor{currentstroke}{rgb}{0.000000,0.000000,0.000000}%
\pgfsetstrokecolor{currentstroke}%
\pgfsetdash{}{0pt}%
\pgfpathmoveto{\pgfqpoint{0.705055in}{0.996102in}}%
\pgfpathlineto{\pgfqpoint{0.721498in}{0.983641in}}%
\pgfpathlineto{\pgfqpoint{0.737956in}{0.971704in}}%
\pgfpathlineto{\pgfqpoint{0.747008in}{0.968762in}}%
\pgfpathlineto{\pgfqpoint{0.756017in}{0.966091in}}%
\pgfpathlineto{\pgfqpoint{0.739605in}{0.978274in}}%
\pgfpathlineto{\pgfqpoint{0.723207in}{0.990986in}}%
\pgfpathlineto{\pgfqpoint{0.714152in}{0.993408in}}%
\pgfpathlineto{\pgfqpoint{0.705055in}{0.996102in}}%
\pgfpathclose%
\pgfusepath{fill}%
\end{pgfscope}%
\begin{pgfscope}%
\pgfpathrectangle{\pgfqpoint{0.407592in}{0.222904in}}{\pgfqpoint{1.411163in}{1.411163in}}%
\pgfusepath{clip}%
\pgfsetbuttcap%
\pgfsetroundjoin%
\definecolor{currentfill}{rgb}{0.909804,0.698039,0.423529}%
\pgfsetfillcolor{currentfill}%
\pgfsetfillopacity{0.120000}%
\pgfsetlinewidth{0.000000pt}%
\definecolor{currentstroke}{rgb}{0.000000,0.000000,0.000000}%
\pgfsetstrokecolor{currentstroke}%
\pgfsetdash{}{0pt}%
\pgfpathmoveto{\pgfqpoint{0.955886in}{0.869491in}}%
\pgfpathlineto{\pgfqpoint{0.972450in}{0.865049in}}%
\pgfpathlineto{\pgfqpoint{0.989049in}{0.861483in}}%
\pgfpathlineto{\pgfqpoint{0.997586in}{0.859667in}}%
\pgfpathlineto{\pgfqpoint{1.006086in}{0.858228in}}%
\pgfpathlineto{\pgfqpoint{0.989531in}{0.861757in}}%
\pgfpathlineto{\pgfqpoint{0.973011in}{0.866197in}}%
\pgfpathlineto{\pgfqpoint{0.964467in}{0.867656in}}%
\pgfpathlineto{\pgfqpoint{0.955886in}{0.869491in}}%
\pgfpathclose%
\pgfusepath{fill}%
\end{pgfscope}%
\begin{pgfscope}%
\pgfpathrectangle{\pgfqpoint{0.407592in}{0.222904in}}{\pgfqpoint{1.411163in}{1.411163in}}%
\pgfusepath{clip}%
\pgfsetbuttcap%
\pgfsetroundjoin%
\definecolor{currentfill}{rgb}{0.909804,0.698039,0.423529}%
\pgfsetfillcolor{currentfill}%
\pgfsetfillopacity{0.120000}%
\pgfsetlinewidth{0.000000pt}%
\definecolor{currentstroke}{rgb}{0.000000,0.000000,0.000000}%
\pgfsetstrokecolor{currentstroke}%
\pgfsetdash{}{0pt}%
\pgfpathmoveto{\pgfqpoint{1.572074in}{1.080578in}}%
\pgfpathlineto{\pgfqpoint{1.589745in}{1.090573in}}%
\pgfpathlineto{\pgfqpoint{1.607493in}{1.100822in}}%
\pgfpathlineto{\pgfqpoint{1.615095in}{1.109582in}}%
\pgfpathlineto{\pgfqpoint{1.622683in}{1.118607in}}%
\pgfpathlineto{\pgfqpoint{1.604958in}{1.108555in}}%
\pgfpathlineto{\pgfqpoint{1.587310in}{1.098764in}}%
\pgfpathlineto{\pgfqpoint{1.579698in}{1.089533in}}%
\pgfpathlineto{\pgfqpoint{1.572074in}{1.080578in}}%
\pgfpathclose%
\pgfusepath{fill}%
\end{pgfscope}%
\begin{pgfscope}%
\pgfpathrectangle{\pgfqpoint{0.407592in}{0.222904in}}{\pgfqpoint{1.411163in}{1.411163in}}%
\pgfusepath{clip}%
\pgfsetbuttcap%
\pgfsetroundjoin%
\definecolor{currentfill}{rgb}{0.909804,0.698039,0.423529}%
\pgfsetfillcolor{currentfill}%
\pgfsetfillopacity{0.120000}%
\pgfsetlinewidth{0.000000pt}%
\definecolor{currentstroke}{rgb}{0.000000,0.000000,0.000000}%
\pgfsetstrokecolor{currentstroke}%
\pgfsetdash{}{0pt}%
\pgfpathmoveto{\pgfqpoint{1.688501in}{1.177199in}}%
\pgfpathlineto{\pgfqpoint{1.706415in}{1.187543in}}%
\pgfpathlineto{\pgfqpoint{1.724408in}{1.198109in}}%
\pgfpathlineto{\pgfqpoint{1.731887in}{1.207906in}}%
\pgfpathlineto{\pgfqpoint{1.739354in}{1.217914in}}%
\pgfpathlineto{\pgfqpoint{1.721385in}{1.207567in}}%
\pgfpathlineto{\pgfqpoint{1.703495in}{1.197448in}}%
\pgfpathlineto{\pgfqpoint{1.696004in}{1.187215in}}%
\pgfpathlineto{\pgfqpoint{1.688501in}{1.177199in}}%
\pgfpathclose%
\pgfusepath{fill}%
\end{pgfscope}%
\begin{pgfscope}%
\pgfpathrectangle{\pgfqpoint{0.407592in}{0.222904in}}{\pgfqpoint{1.411163in}{1.411163in}}%
\pgfusepath{clip}%
\pgfsetbuttcap%
\pgfsetroundjoin%
\definecolor{currentfill}{rgb}{0.909804,0.698039,0.423529}%
\pgfsetfillcolor{currentfill}%
\pgfsetfillopacity{0.120000}%
\pgfsetlinewidth{0.000000pt}%
\definecolor{currentstroke}{rgb}{0.000000,0.000000,0.000000}%
\pgfsetstrokecolor{currentstroke}%
\pgfsetdash{}{0pt}%
\pgfpathmoveto{\pgfqpoint{1.189483in}{0.864945in}}%
\pgfpathlineto{\pgfqpoint{1.206370in}{0.869384in}}%
\pgfpathlineto{\pgfqpoint{1.223319in}{0.874486in}}%
\pgfpathlineto{\pgfqpoint{1.231418in}{0.876840in}}%
\pgfpathlineto{\pgfqpoint{1.239491in}{0.879655in}}%
\pgfpathlineto{\pgfqpoint{1.222571in}{0.874416in}}%
\pgfpathlineto{\pgfqpoint{1.205714in}{0.869841in}}%
\pgfpathlineto{\pgfqpoint{1.197611in}{0.867152in}}%
\pgfpathlineto{\pgfqpoint{1.189483in}{0.864945in}}%
\pgfpathclose%
\pgfusepath{fill}%
\end{pgfscope}%
\begin{pgfscope}%
\pgfpathrectangle{\pgfqpoint{0.407592in}{0.222904in}}{\pgfqpoint{1.411163in}{1.411163in}}%
\pgfusepath{clip}%
\pgfsetbuttcap%
\pgfsetroundjoin%
\definecolor{currentfill}{rgb}{0.909804,0.698039,0.423529}%
\pgfsetfillcolor{currentfill}%
\pgfsetfillopacity{0.120000}%
\pgfsetlinewidth{0.000000pt}%
\definecolor{currentstroke}{rgb}{0.000000,0.000000,0.000000}%
\pgfsetstrokecolor{currentstroke}%
\pgfsetdash{}{0pt}%
\pgfpathmoveto{\pgfqpoint{1.455844in}{0.993978in}}%
\pgfpathlineto{\pgfqpoint{1.473277in}{1.003300in}}%
\pgfpathlineto{\pgfqpoint{1.490786in}{1.012935in}}%
\pgfpathlineto{\pgfqpoint{1.498519in}{1.020205in}}%
\pgfpathlineto{\pgfqpoint{1.506237in}{1.027808in}}%
\pgfpathlineto{\pgfqpoint{1.488751in}{1.018310in}}%
\pgfpathlineto{\pgfqpoint{1.471340in}{1.009129in}}%
\pgfpathlineto{\pgfqpoint{1.463600in}{1.001378in}}%
\pgfpathlineto{\pgfqpoint{1.455844in}{0.993978in}}%
\pgfpathclose%
\pgfusepath{fill}%
\end{pgfscope}%
\begin{pgfscope}%
\pgfpathrectangle{\pgfqpoint{0.407592in}{0.222904in}}{\pgfqpoint{1.411163in}{1.411163in}}%
\pgfusepath{clip}%
\pgfsetbuttcap%
\pgfsetroundjoin%
\definecolor{currentfill}{rgb}{0.909804,0.698039,0.423529}%
\pgfsetfillcolor{currentfill}%
\pgfsetfillopacity{0.120000}%
\pgfsetlinewidth{0.000000pt}%
\definecolor{currentstroke}{rgb}{0.000000,0.000000,0.000000}%
\pgfsetstrokecolor{currentstroke}%
\pgfsetdash{}{0pt}%
\pgfpathmoveto{\pgfqpoint{1.139481in}{0.855560in}}%
\pgfpathlineto{\pgfqpoint{1.156282in}{0.858370in}}%
\pgfpathlineto{\pgfqpoint{1.173142in}{0.861930in}}%
\pgfpathlineto{\pgfqpoint{1.181326in}{0.863208in}}%
\pgfpathlineto{\pgfqpoint{1.189483in}{0.864945in}}%
\pgfpathlineto{\pgfqpoint{1.172656in}{0.861216in}}%
\pgfpathlineto{\pgfqpoint{1.155887in}{0.858243in}}%
\pgfpathlineto{\pgfqpoint{1.147699in}{0.856664in}}%
\pgfpathlineto{\pgfqpoint{1.139481in}{0.855560in}}%
\pgfpathclose%
\pgfusepath{fill}%
\end{pgfscope}%
\begin{pgfscope}%
\pgfpathrectangle{\pgfqpoint{0.407592in}{0.222904in}}{\pgfqpoint{1.411163in}{1.411163in}}%
\pgfusepath{clip}%
\pgfsetbuttcap%
\pgfsetroundjoin%
\definecolor{currentfill}{rgb}{0.909804,0.698039,0.423529}%
\pgfsetfillcolor{currentfill}%
\pgfsetfillopacity{0.120000}%
\pgfsetlinewidth{0.000000pt}%
\definecolor{currentstroke}{rgb}{0.000000,0.000000,0.000000}%
\pgfsetstrokecolor{currentstroke}%
\pgfsetdash{}{0pt}%
\pgfpathmoveto{\pgfqpoint{1.239491in}{0.879655in}}%
\pgfpathlineto{\pgfqpoint{1.256476in}{0.885511in}}%
\pgfpathlineto{\pgfqpoint{1.273528in}{0.891941in}}%
\pgfpathlineto{\pgfqpoint{1.281548in}{0.895325in}}%
\pgfpathlineto{\pgfqpoint{1.289546in}{0.899158in}}%
\pgfpathlineto{\pgfqpoint{1.272521in}{0.892640in}}%
\pgfpathlineto{\pgfqpoint{1.255564in}{0.886694in}}%
\pgfpathlineto{\pgfqpoint{1.247539in}{0.882937in}}%
\pgfpathlineto{\pgfqpoint{1.239491in}{0.879655in}}%
\pgfpathclose%
\pgfusepath{fill}%
\end{pgfscope}%
\begin{pgfscope}%
\pgfpathrectangle{\pgfqpoint{0.407592in}{0.222904in}}{\pgfqpoint{1.411163in}{1.411163in}}%
\pgfusepath{clip}%
\pgfsetbuttcap%
\pgfsetroundjoin%
\definecolor{currentfill}{rgb}{0.909804,0.698039,0.423529}%
\pgfsetfillcolor{currentfill}%
\pgfsetfillopacity{0.120000}%
\pgfsetlinewidth{0.000000pt}%
\definecolor{currentstroke}{rgb}{0.000000,0.000000,0.000000}%
\pgfsetstrokecolor{currentstroke}%
\pgfsetdash{}{0pt}%
\pgfpathmoveto{\pgfqpoint{1.089440in}{0.851845in}}%
\pgfpathlineto{\pgfqpoint{1.106171in}{0.852869in}}%
\pgfpathlineto{\pgfqpoint{1.122953in}{0.854710in}}%
\pgfpathlineto{\pgfqpoint{1.131232in}{0.854914in}}%
\pgfpathlineto{\pgfqpoint{1.139481in}{0.855560in}}%
\pgfpathlineto{\pgfqpoint{1.122734in}{0.853542in}}%
\pgfpathlineto{\pgfqpoint{1.106040in}{0.852358in}}%
\pgfpathlineto{\pgfqpoint{1.097756in}{0.851876in}}%
\pgfpathlineto{\pgfqpoint{1.089440in}{0.851845in}}%
\pgfpathclose%
\pgfusepath{fill}%
\end{pgfscope}%
\begin{pgfscope}%
\pgfpathrectangle{\pgfqpoint{0.407592in}{0.222904in}}{\pgfqpoint{1.411163in}{1.411163in}}%
\pgfusepath{clip}%
\pgfsetbuttcap%
\pgfsetroundjoin%
\definecolor{currentfill}{rgb}{0.909804,0.698039,0.423529}%
\pgfsetfillcolor{currentfill}%
\pgfsetfillopacity{0.120000}%
\pgfsetlinewidth{0.000000pt}%
\definecolor{currentstroke}{rgb}{0.000000,0.000000,0.000000}%
\pgfsetstrokecolor{currentstroke}%
\pgfsetdash{}{0pt}%
\pgfpathmoveto{\pgfqpoint{0.620906in}{1.055712in}}%
\pgfpathlineto{\pgfqpoint{0.637343in}{1.041707in}}%
\pgfpathlineto{\pgfqpoint{0.653793in}{1.028113in}}%
\pgfpathlineto{\pgfqpoint{0.663022in}{1.025169in}}%
\pgfpathlineto{\pgfqpoint{0.672208in}{1.022474in}}%
\pgfpathlineto{\pgfqpoint{0.655803in}{1.036328in}}%
\pgfpathlineto{\pgfqpoint{0.639410in}{1.050592in}}%
\pgfpathlineto{\pgfqpoint{0.630180in}{1.053028in}}%
\pgfpathlineto{\pgfqpoint{0.620906in}{1.055712in}}%
\pgfpathclose%
\pgfusepath{fill}%
\end{pgfscope}%
\begin{pgfscope}%
\pgfpathrectangle{\pgfqpoint{0.407592in}{0.222904in}}{\pgfqpoint{1.411163in}{1.411163in}}%
\pgfusepath{clip}%
\pgfsetbuttcap%
\pgfsetroundjoin%
\definecolor{currentfill}{rgb}{0.909804,0.698039,0.423529}%
\pgfsetfillcolor{currentfill}%
\pgfsetfillopacity{0.120000}%
\pgfsetlinewidth{0.000000pt}%
\definecolor{currentstroke}{rgb}{0.000000,0.000000,0.000000}%
\pgfsetstrokecolor{currentstroke}%
\pgfsetdash{}{0pt}%
\pgfpathmoveto{\pgfqpoint{1.289546in}{0.899158in}}%
\pgfpathlineto{\pgfqpoint{1.306639in}{0.906204in}}%
\pgfpathlineto{\pgfqpoint{1.323802in}{0.913739in}}%
\pgfpathlineto{\pgfqpoint{1.331751in}{0.918072in}}%
\pgfpathlineto{\pgfqpoint{1.339680in}{0.922834in}}%
\pgfpathlineto{\pgfqpoint{1.322543in}{0.915266in}}%
\pgfpathlineto{\pgfqpoint{1.305475in}{0.908184in}}%
\pgfpathlineto{\pgfqpoint{1.297521in}{0.903444in}}%
\pgfpathlineto{\pgfqpoint{1.289546in}{0.899158in}}%
\pgfpathclose%
\pgfusepath{fill}%
\end{pgfscope}%
\begin{pgfscope}%
\pgfpathrectangle{\pgfqpoint{0.407592in}{0.222904in}}{\pgfqpoint{1.411163in}{1.411163in}}%
\pgfusepath{clip}%
\pgfsetbuttcap%
\pgfsetroundjoin%
\definecolor{currentfill}{rgb}{0.909804,0.698039,0.423529}%
\pgfsetfillcolor{currentfill}%
\pgfsetfillopacity{0.120000}%
\pgfsetlinewidth{0.000000pt}%
\definecolor{currentstroke}{rgb}{0.000000,0.000000,0.000000}%
\pgfsetstrokecolor{currentstroke}%
\pgfsetdash{}{0pt}%
\pgfpathmoveto{\pgfqpoint{0.503564in}{1.151611in}}%
\pgfpathlineto{\pgfqpoint{0.519981in}{1.135867in}}%
\pgfpathlineto{\pgfqpoint{0.536408in}{1.120407in}}%
\pgfpathlineto{\pgfqpoint{0.545860in}{1.117763in}}%
\pgfpathlineto{\pgfqpoint{0.555267in}{1.115347in}}%
\pgfpathlineto{\pgfqpoint{0.538884in}{1.131051in}}%
\pgfpathlineto{\pgfqpoint{0.522510in}{1.147032in}}%
\pgfpathlineto{\pgfqpoint{0.513059in}{1.149208in}}%
\pgfpathlineto{\pgfqpoint{0.503564in}{1.151611in}}%
\pgfpathclose%
\pgfusepath{fill}%
\end{pgfscope}%
\begin{pgfscope}%
\pgfpathrectangle{\pgfqpoint{0.407592in}{0.222904in}}{\pgfqpoint{1.411163in}{1.411163in}}%
\pgfusepath{clip}%
\pgfsetbuttcap%
\pgfsetroundjoin%
\definecolor{currentfill}{rgb}{0.909804,0.698039,0.423529}%
\pgfsetfillcolor{currentfill}%
\pgfsetfillopacity{0.120000}%
\pgfsetlinewidth{0.000000pt}%
\definecolor{currentstroke}{rgb}{0.000000,0.000000,0.000000}%
\pgfsetstrokecolor{currentstroke}%
\pgfsetdash{}{0pt}%
\pgfpathmoveto{\pgfqpoint{0.821827in}{0.923305in}}%
\pgfpathlineto{\pgfqpoint{0.838326in}{0.914253in}}%
\pgfpathlineto{\pgfqpoint{0.854846in}{0.905919in}}%
\pgfpathlineto{\pgfqpoint{0.863675in}{0.902934in}}%
\pgfpathlineto{\pgfqpoint{0.872465in}{0.900249in}}%
\pgfpathlineto{\pgfqpoint{0.855991in}{0.908735in}}%
\pgfpathlineto{\pgfqpoint{0.839539in}{0.917962in}}%
\pgfpathlineto{\pgfqpoint{0.830703in}{0.920483in}}%
\pgfpathlineto{\pgfqpoint{0.821827in}{0.923305in}}%
\pgfpathclose%
\pgfusepath{fill}%
\end{pgfscope}%
\begin{pgfscope}%
\pgfpathrectangle{\pgfqpoint{0.407592in}{0.222904in}}{\pgfqpoint{1.411163in}{1.411163in}}%
\pgfusepath{clip}%
\pgfsetbuttcap%
\pgfsetroundjoin%
\definecolor{currentfill}{rgb}{0.909804,0.698039,0.423529}%
\pgfsetfillcolor{currentfill}%
\pgfsetfillopacity{0.120000}%
\pgfsetlinewidth{0.000000pt}%
\definecolor{currentstroke}{rgb}{0.000000,0.000000,0.000000}%
\pgfsetstrokecolor{currentstroke}%
\pgfsetdash{}{0pt}%
\pgfpathmoveto{\pgfqpoint{1.622683in}{1.118607in}}%
\pgfpathlineto{\pgfqpoint{1.640486in}{1.128906in}}%
\pgfpathlineto{\pgfqpoint{1.658368in}{1.139438in}}%
\pgfpathlineto{\pgfqpoint{1.665920in}{1.148520in}}%
\pgfpathlineto{\pgfqpoint{1.673460in}{1.157845in}}%
\pgfpathlineto{\pgfqpoint{1.655601in}{1.147516in}}%
\pgfpathlineto{\pgfqpoint{1.637821in}{1.137428in}}%
\pgfpathlineto{\pgfqpoint{1.630258in}{1.127892in}}%
\pgfpathlineto{\pgfqpoint{1.622683in}{1.118607in}}%
\pgfpathclose%
\pgfusepath{fill}%
\end{pgfscope}%
\begin{pgfscope}%
\pgfpathrectangle{\pgfqpoint{0.407592in}{0.222904in}}{\pgfqpoint{1.411163in}{1.411163in}}%
\pgfusepath{clip}%
\pgfsetbuttcap%
\pgfsetroundjoin%
\definecolor{currentfill}{rgb}{0.909804,0.698039,0.423529}%
\pgfsetfillcolor{currentfill}%
\pgfsetfillopacity{0.120000}%
\pgfsetlinewidth{0.000000pt}%
\definecolor{currentstroke}{rgb}{0.000000,0.000000,0.000000}%
\pgfsetstrokecolor{currentstroke}%
\pgfsetdash{}{0pt}%
\pgfpathmoveto{\pgfqpoint{0.905483in}{0.885606in}}%
\pgfpathlineto{\pgfqpoint{0.922031in}{0.879494in}}%
\pgfpathlineto{\pgfqpoint{0.938609in}{0.874211in}}%
\pgfpathlineto{\pgfqpoint{0.947267in}{0.871683in}}%
\pgfpathlineto{\pgfqpoint{0.955886in}{0.869491in}}%
\pgfpathlineto{\pgfqpoint{0.939353in}{0.874807in}}%
\pgfpathlineto{\pgfqpoint{0.922850in}{0.880983in}}%
\pgfpathlineto{\pgfqpoint{0.914186in}{0.883127in}}%
\pgfpathlineto{\pgfqpoint{0.905483in}{0.885606in}}%
\pgfpathclose%
\pgfusepath{fill}%
\end{pgfscope}%
\begin{pgfscope}%
\pgfpathrectangle{\pgfqpoint{0.407592in}{0.222904in}}{\pgfqpoint{1.411163in}{1.411163in}}%
\pgfusepath{clip}%
\pgfsetbuttcap%
\pgfsetroundjoin%
\definecolor{currentfill}{rgb}{0.909804,0.698039,0.423529}%
\pgfsetfillcolor{currentfill}%
\pgfsetfillopacity{0.120000}%
\pgfsetlinewidth{0.000000pt}%
\definecolor{currentstroke}{rgb}{0.000000,0.000000,0.000000}%
\pgfsetstrokecolor{currentstroke}%
\pgfsetdash{}{0pt}%
\pgfpathmoveto{\pgfqpoint{0.737956in}{0.971704in}}%
\pgfpathlineto{\pgfqpoint{0.754429in}{0.960321in}}%
\pgfpathlineto{\pgfqpoint{0.770919in}{0.949521in}}%
\pgfpathlineto{\pgfqpoint{0.779924in}{0.946347in}}%
\pgfpathlineto{\pgfqpoint{0.788888in}{0.943444in}}%
\pgfpathlineto{\pgfqpoint{0.772445in}{0.954470in}}%
\pgfpathlineto{\pgfqpoint{0.756017in}{0.966091in}}%
\pgfpathlineto{\pgfqpoint{0.747008in}{0.968762in}}%
\pgfpathlineto{\pgfqpoint{0.737956in}{0.971704in}}%
\pgfpathclose%
\pgfusepath{fill}%
\end{pgfscope}%
\begin{pgfscope}%
\pgfpathrectangle{\pgfqpoint{0.407592in}{0.222904in}}{\pgfqpoint{1.411163in}{1.411163in}}%
\pgfusepath{clip}%
\pgfsetbuttcap%
\pgfsetroundjoin%
\definecolor{currentfill}{rgb}{0.909804,0.698039,0.423529}%
\pgfsetfillcolor{currentfill}%
\pgfsetfillopacity{0.120000}%
\pgfsetlinewidth{0.000000pt}%
\definecolor{currentstroke}{rgb}{0.000000,0.000000,0.000000}%
\pgfsetstrokecolor{currentstroke}%
\pgfsetdash{}{0pt}%
\pgfpathmoveto{\pgfqpoint{1.506237in}{1.027808in}}%
\pgfpathlineto{\pgfqpoint{1.523800in}{1.037601in}}%
\pgfpathlineto{\pgfqpoint{1.541439in}{1.047671in}}%
\pgfpathlineto{\pgfqpoint{1.549119in}{1.055447in}}%
\pgfpathlineto{\pgfqpoint{1.556785in}{1.063528in}}%
\pgfpathlineto{\pgfqpoint{1.539168in}{1.053617in}}%
\pgfpathlineto{\pgfqpoint{1.521629in}{1.043989in}}%
\pgfpathlineto{\pgfqpoint{1.513940in}{1.035738in}}%
\pgfpathlineto{\pgfqpoint{1.506237in}{1.027808in}}%
\pgfpathclose%
\pgfusepath{fill}%
\end{pgfscope}%
\begin{pgfscope}%
\pgfpathrectangle{\pgfqpoint{0.407592in}{0.222904in}}{\pgfqpoint{1.411163in}{1.411163in}}%
\pgfusepath{clip}%
\pgfsetbuttcap%
\pgfsetroundjoin%
\definecolor{currentfill}{rgb}{0.909804,0.698039,0.423529}%
\pgfsetfillcolor{currentfill}%
\pgfsetfillopacity{0.120000}%
\pgfsetlinewidth{0.000000pt}%
\definecolor{currentstroke}{rgb}{0.000000,0.000000,0.000000}%
\pgfsetstrokecolor{currentstroke}%
\pgfsetdash{}{0pt}%
\pgfpathmoveto{\pgfqpoint{1.039311in}{0.853885in}}%
\pgfpathlineto{\pgfqpoint{1.055987in}{0.853044in}}%
\pgfpathlineto{\pgfqpoint{1.072708in}{0.853062in}}%
\pgfpathlineto{\pgfqpoint{1.081091in}{0.852247in}}%
\pgfpathlineto{\pgfqpoint{1.089440in}{0.851845in}}%
\pgfpathlineto{\pgfqpoint{1.072758in}{0.851675in}}%
\pgfpathlineto{\pgfqpoint{1.056122in}{0.852389in}}%
\pgfpathlineto{\pgfqpoint{1.047734in}{0.852928in}}%
\pgfpathlineto{\pgfqpoint{1.039311in}{0.853885in}}%
\pgfpathclose%
\pgfusepath{fill}%
\end{pgfscope}%
\begin{pgfscope}%
\pgfpathrectangle{\pgfqpoint{0.407592in}{0.222904in}}{\pgfqpoint{1.411163in}{1.411163in}}%
\pgfusepath{clip}%
\pgfsetbuttcap%
\pgfsetroundjoin%
\definecolor{currentfill}{rgb}{0.909804,0.698039,0.423529}%
\pgfsetfillcolor{currentfill}%
\pgfsetfillopacity{0.120000}%
\pgfsetlinewidth{0.000000pt}%
\definecolor{currentstroke}{rgb}{0.000000,0.000000,0.000000}%
\pgfsetstrokecolor{currentstroke}%
\pgfsetdash{}{0pt}%
\pgfpathmoveto{\pgfqpoint{1.339680in}{0.922834in}}%
\pgfpathlineto{\pgfqpoint{1.356889in}{0.930853in}}%
\pgfpathlineto{\pgfqpoint{1.374170in}{0.939287in}}%
\pgfpathlineto{\pgfqpoint{1.382054in}{0.944472in}}%
\pgfpathlineto{\pgfqpoint{1.389919in}{0.950059in}}%
\pgfpathlineto{\pgfqpoint{1.372662in}{0.941644in}}%
\pgfpathlineto{\pgfqpoint{1.355477in}{0.933643in}}%
\pgfpathlineto{\pgfqpoint{1.347589in}{0.928025in}}%
\pgfpathlineto{\pgfqpoint{1.339680in}{0.922834in}}%
\pgfpathclose%
\pgfusepath{fill}%
\end{pgfscope}%
\begin{pgfscope}%
\pgfpathrectangle{\pgfqpoint{0.407592in}{0.222904in}}{\pgfqpoint{1.411163in}{1.411163in}}%
\pgfusepath{clip}%
\pgfsetbuttcap%
\pgfsetroundjoin%
\definecolor{currentfill}{rgb}{0.909804,0.698039,0.423529}%
\pgfsetfillcolor{currentfill}%
\pgfsetfillopacity{0.120000}%
\pgfsetlinewidth{0.000000pt}%
\definecolor{currentstroke}{rgb}{0.000000,0.000000,0.000000}%
\pgfsetstrokecolor{currentstroke}%
\pgfsetdash{}{0pt}%
\pgfpathmoveto{\pgfqpoint{1.739354in}{1.217914in}}%
\pgfpathlineto{\pgfqpoint{1.757402in}{1.228479in}}%
\pgfpathlineto{\pgfqpoint{1.775531in}{1.239254in}}%
\pgfpathlineto{\pgfqpoint{1.782961in}{1.249249in}}%
\pgfpathlineto{\pgfqpoint{1.790381in}{1.259440in}}%
\pgfpathlineto{\pgfqpoint{1.772277in}{1.248883in}}%
\pgfpathlineto{\pgfqpoint{1.754253in}{1.238540in}}%
\pgfpathlineto{\pgfqpoint{1.746809in}{1.228127in}}%
\pgfpathlineto{\pgfqpoint{1.739354in}{1.217914in}}%
\pgfpathclose%
\pgfusepath{fill}%
\end{pgfscope}%
\begin{pgfscope}%
\pgfpathrectangle{\pgfqpoint{0.407592in}{0.222904in}}{\pgfqpoint{1.411163in}{1.411163in}}%
\pgfusepath{clip}%
\pgfsetbuttcap%
\pgfsetroundjoin%
\definecolor{currentfill}{rgb}{0.909804,0.698039,0.423529}%
\pgfsetfillcolor{currentfill}%
\pgfsetfillopacity{0.120000}%
\pgfsetlinewidth{0.000000pt}%
\definecolor{currentstroke}{rgb}{0.000000,0.000000,0.000000}%
\pgfsetstrokecolor{currentstroke}%
\pgfsetdash{}{0pt}%
\pgfpathmoveto{\pgfqpoint{1.790381in}{1.259440in}}%
\pgfpathlineto{\pgfqpoint{1.808565in}{1.270203in}}%
\pgfpathlineto{\pgfqpoint{1.815961in}{1.280474in}}%
\pgfpathlineto{\pgfqpoint{1.823345in}{1.290927in}}%
\pgfpathlineto{\pgfqpoint{1.805186in}{1.280383in}}%
\pgfpathlineto{\pgfqpoint{1.797789in}{1.269820in}}%
\pgfusepath{fill}%
\end{pgfscope}%
\begin{pgfscope}%
\pgfpathrectangle{\pgfqpoint{0.407592in}{0.222904in}}{\pgfqpoint{1.411163in}{1.411163in}}%
\pgfusepath{clip}%
\pgfsetbuttcap%
\pgfsetroundjoin%
\definecolor{currentfill}{rgb}{0.909804,0.698039,0.423529}%
\pgfsetfillcolor{currentfill}%
\pgfsetfillopacity{0.120000}%
\pgfsetlinewidth{0.000000pt}%
\definecolor{currentstroke}{rgb}{0.000000,0.000000,0.000000}%
\pgfsetstrokecolor{currentstroke}%
\pgfsetdash{}{0pt}%
\pgfpathmoveto{\pgfqpoint{0.536408in}{1.120407in}}%
\pgfpathlineto{\pgfqpoint{0.552846in}{1.105248in}}%
\pgfpathlineto{\pgfqpoint{0.569295in}{1.090412in}}%
\pgfpathlineto{\pgfqpoint{0.578702in}{1.087517in}}%
\pgfpathlineto{\pgfqpoint{0.588065in}{1.084851in}}%
\pgfpathlineto{\pgfqpoint{0.571661in}{1.099941in}}%
\pgfpathlineto{\pgfqpoint{0.555267in}{1.115347in}}%
\pgfpathlineto{\pgfqpoint{0.545860in}{1.117763in}}%
\pgfpathlineto{\pgfqpoint{0.536408in}{1.120407in}}%
\pgfpathclose%
\pgfusepath{fill}%
\end{pgfscope}%
\begin{pgfscope}%
\pgfpathrectangle{\pgfqpoint{0.407592in}{0.222904in}}{\pgfqpoint{1.411163in}{1.411163in}}%
\pgfusepath{clip}%
\pgfsetbuttcap%
\pgfsetroundjoin%
\definecolor{currentfill}{rgb}{0.909804,0.698039,0.423529}%
\pgfsetfillcolor{currentfill}%
\pgfsetfillopacity{0.120000}%
\pgfsetlinewidth{0.000000pt}%
\definecolor{currentstroke}{rgb}{0.000000,0.000000,0.000000}%
\pgfsetstrokecolor{currentstroke}%
\pgfsetdash{}{0pt}%
\pgfpathmoveto{\pgfqpoint{0.653793in}{1.028113in}}%
\pgfpathlineto{\pgfqpoint{0.670255in}{1.014954in}}%
\pgfpathlineto{\pgfqpoint{0.686731in}{1.002256in}}%
\pgfpathlineto{\pgfqpoint{0.695914in}{0.999055in}}%
\pgfpathlineto{\pgfqpoint{0.705055in}{0.996102in}}%
\pgfpathlineto{\pgfqpoint{0.688625in}{1.009055in}}%
\pgfpathlineto{\pgfqpoint{0.672208in}{1.022474in}}%
\pgfpathlineto{\pgfqpoint{0.663022in}{1.025169in}}%
\pgfpathlineto{\pgfqpoint{0.653793in}{1.028113in}}%
\pgfpathclose%
\pgfusepath{fill}%
\end{pgfscope}%
\begin{pgfscope}%
\pgfpathrectangle{\pgfqpoint{0.407592in}{0.222904in}}{\pgfqpoint{1.411163in}{1.411163in}}%
\pgfusepath{clip}%
\pgfsetbuttcap%
\pgfsetroundjoin%
\definecolor{currentfill}{rgb}{0.909804,0.698039,0.423529}%
\pgfsetfillcolor{currentfill}%
\pgfsetfillopacity{0.120000}%
\pgfsetlinewidth{0.000000pt}%
\definecolor{currentstroke}{rgb}{0.000000,0.000000,0.000000}%
\pgfsetstrokecolor{currentstroke}%
\pgfsetdash{}{0pt}%
\pgfpathmoveto{\pgfqpoint{1.389919in}{0.950059in}}%
\pgfpathlineto{\pgfqpoint{1.407250in}{0.958860in}}%
\pgfpathlineto{\pgfqpoint{1.424654in}{0.968017in}}%
\pgfpathlineto{\pgfqpoint{1.432477in}{0.973951in}}%
\pgfpathlineto{\pgfqpoint{1.440283in}{0.980259in}}%
\pgfpathlineto{\pgfqpoint{1.422902in}{0.971166in}}%
\pgfpathlineto{\pgfqpoint{1.405595in}{0.962430in}}%
\pgfpathlineto{\pgfqpoint{1.397766in}{0.956047in}}%
\pgfpathlineto{\pgfqpoint{1.389919in}{0.950059in}}%
\pgfpathclose%
\pgfusepath{fill}%
\end{pgfscope}%
\begin{pgfscope}%
\pgfpathrectangle{\pgfqpoint{0.407592in}{0.222904in}}{\pgfqpoint{1.411163in}{1.411163in}}%
\pgfusepath{clip}%
\pgfsetbuttcap%
\pgfsetroundjoin%
\definecolor{currentfill}{rgb}{0.909804,0.698039,0.423529}%
\pgfsetfillcolor{currentfill}%
\pgfsetfillopacity{0.120000}%
\pgfsetlinewidth{0.000000pt}%
\definecolor{currentstroke}{rgb}{0.000000,0.000000,0.000000}%
\pgfsetstrokecolor{currentstroke}%
\pgfsetdash{}{0pt}%
\pgfpathmoveto{\pgfqpoint{0.418690in}{1.221637in}}%
\pgfpathlineto{\pgfqpoint{0.435112in}{1.205141in}}%
\pgfpathlineto{\pgfqpoint{0.451544in}{1.188871in}}%
\pgfpathlineto{\pgfqpoint{0.461174in}{1.186270in}}%
\pgfpathlineto{\pgfqpoint{0.470759in}{1.183880in}}%
\pgfpathlineto{\pgfqpoint{0.454371in}{1.200373in}}%
\pgfpathlineto{\pgfqpoint{0.437993in}{1.217086in}}%
\pgfpathlineto{\pgfqpoint{0.428365in}{1.219257in}}%
\pgfpathlineto{\pgfqpoint{0.418690in}{1.221637in}}%
\pgfpathclose%
\pgfusepath{fill}%
\end{pgfscope}%
\begin{pgfscope}%
\pgfpathrectangle{\pgfqpoint{0.407592in}{0.222904in}}{\pgfqpoint{1.411163in}{1.411163in}}%
\pgfusepath{clip}%
\pgfsetbuttcap%
\pgfsetroundjoin%
\definecolor{currentfill}{rgb}{0.909804,0.698039,0.423529}%
\pgfsetfillcolor{currentfill}%
\pgfsetfillopacity{0.120000}%
\pgfsetlinewidth{0.000000pt}%
\definecolor{currentstroke}{rgb}{0.000000,0.000000,0.000000}%
\pgfsetstrokecolor{currentstroke}%
\pgfsetdash{}{0pt}%
\pgfpathmoveto{\pgfqpoint{0.989049in}{0.861483in}}%
\pgfpathlineto{\pgfqpoint{1.005684in}{0.858792in}}%
\pgfpathlineto{\pgfqpoint{1.022360in}{0.856968in}}%
\pgfpathlineto{\pgfqpoint{1.030853in}{0.855238in}}%
\pgfpathlineto{\pgfqpoint{1.039311in}{0.853885in}}%
\pgfpathlineto{\pgfqpoint{1.022679in}{0.855607in}}%
\pgfpathlineto{\pgfqpoint{1.006086in}{0.858228in}}%
\pgfpathlineto{\pgfqpoint{0.997586in}{0.859667in}}%
\pgfpathlineto{\pgfqpoint{0.989049in}{0.861483in}}%
\pgfpathclose%
\pgfusepath{fill}%
\end{pgfscope}%
\begin{pgfscope}%
\pgfpathrectangle{\pgfqpoint{0.407592in}{0.222904in}}{\pgfqpoint{1.411163in}{1.411163in}}%
\pgfusepath{clip}%
\pgfsetbuttcap%
\pgfsetroundjoin%
\definecolor{currentfill}{rgb}{0.909804,0.698039,0.423529}%
\pgfsetfillcolor{currentfill}%
\pgfsetfillopacity{0.120000}%
\pgfsetlinewidth{0.000000pt}%
\definecolor{currentstroke}{rgb}{0.000000,0.000000,0.000000}%
\pgfsetstrokecolor{currentstroke}%
\pgfsetdash{}{0pt}%
\pgfpathmoveto{\pgfqpoint{1.556785in}{1.063528in}}%
\pgfpathlineto{\pgfqpoint{1.574478in}{1.073703in}}%
\pgfpathlineto{\pgfqpoint{1.592250in}{1.084128in}}%
\pgfpathlineto{\pgfqpoint{1.599879in}{1.092336in}}%
\pgfpathlineto{\pgfqpoint{1.607493in}{1.100822in}}%
\pgfpathlineto{\pgfqpoint{1.589745in}{1.090573in}}%
\pgfpathlineto{\pgfqpoint{1.572074in}{1.080578in}}%
\pgfpathlineto{\pgfqpoint{1.564436in}{1.071907in}}%
\pgfpathlineto{\pgfqpoint{1.556785in}{1.063528in}}%
\pgfpathclose%
\pgfusepath{fill}%
\end{pgfscope}%
\begin{pgfscope}%
\pgfpathrectangle{\pgfqpoint{0.407592in}{0.222904in}}{\pgfqpoint{1.411163in}{1.411163in}}%
\pgfusepath{clip}%
\pgfsetbuttcap%
\pgfsetroundjoin%
\definecolor{currentfill}{rgb}{0.909804,0.698039,0.423529}%
\pgfsetfillcolor{currentfill}%
\pgfsetfillopacity{0.120000}%
\pgfsetlinewidth{0.000000pt}%
\definecolor{currentstroke}{rgb}{0.000000,0.000000,0.000000}%
\pgfsetstrokecolor{currentstroke}%
\pgfsetdash{}{0pt}%
\pgfpathmoveto{\pgfqpoint{1.673460in}{1.157845in}}%
\pgfpathlineto{\pgfqpoint{1.691397in}{1.168401in}}%
\pgfpathlineto{\pgfqpoint{1.709414in}{1.179174in}}%
\pgfpathlineto{\pgfqpoint{1.716917in}{1.188529in}}%
\pgfpathlineto{\pgfqpoint{1.724408in}{1.198109in}}%
\pgfpathlineto{\pgfqpoint{1.706415in}{1.187543in}}%
\pgfpathlineto{\pgfqpoint{1.688501in}{1.177199in}}%
\pgfpathlineto{\pgfqpoint{1.680987in}{1.167407in}}%
\pgfpathlineto{\pgfqpoint{1.673460in}{1.157845in}}%
\pgfpathclose%
\pgfusepath{fill}%
\end{pgfscope}%
\begin{pgfscope}%
\pgfpathrectangle{\pgfqpoint{0.407592in}{0.222904in}}{\pgfqpoint{1.411163in}{1.411163in}}%
\pgfusepath{clip}%
\pgfsetbuttcap%
\pgfsetroundjoin%
\definecolor{currentfill}{rgb}{0.909804,0.698039,0.423529}%
\pgfsetfillcolor{currentfill}%
\pgfsetfillopacity{0.120000}%
\pgfsetlinewidth{0.000000pt}%
\definecolor{currentstroke}{rgb}{0.000000,0.000000,0.000000}%
\pgfsetstrokecolor{currentstroke}%
\pgfsetdash{}{0pt}%
\pgfpathmoveto{\pgfqpoint{0.854846in}{0.905919in}}%
\pgfpathlineto{\pgfqpoint{0.871389in}{0.898327in}}%
\pgfpathlineto{\pgfqpoint{0.887957in}{0.891498in}}%
\pgfpathlineto{\pgfqpoint{0.896740in}{0.888402in}}%
\pgfpathlineto{\pgfqpoint{0.905483in}{0.885606in}}%
\pgfpathlineto{\pgfqpoint{0.888961in}{0.892531in}}%
\pgfpathlineto{\pgfqpoint{0.872465in}{0.900249in}}%
\pgfpathlineto{\pgfqpoint{0.863675in}{0.902934in}}%
\pgfpathlineto{\pgfqpoint{0.854846in}{0.905919in}}%
\pgfpathclose%
\pgfusepath{fill}%
\end{pgfscope}%
\begin{pgfscope}%
\pgfpathrectangle{\pgfqpoint{0.407592in}{0.222904in}}{\pgfqpoint{1.411163in}{1.411163in}}%
\pgfusepath{clip}%
\pgfsetbuttcap%
\pgfsetroundjoin%
\definecolor{currentfill}{rgb}{0.909804,0.698039,0.423529}%
\pgfsetfillcolor{currentfill}%
\pgfsetfillopacity{0.120000}%
\pgfsetlinewidth{0.000000pt}%
\definecolor{currentstroke}{rgb}{0.000000,0.000000,0.000000}%
\pgfsetstrokecolor{currentstroke}%
\pgfsetdash{}{0pt}%
\pgfpathmoveto{\pgfqpoint{1.440283in}{0.980259in}}%
\pgfpathlineto{\pgfqpoint{1.457739in}{0.989684in}}%
\pgfpathlineto{\pgfqpoint{1.475271in}{0.999420in}}%
\pgfpathlineto{\pgfqpoint{1.483037in}{1.006005in}}%
\pgfpathlineto{\pgfqpoint{1.490786in}{1.012935in}}%
\pgfpathlineto{\pgfqpoint{1.473277in}{1.003300in}}%
\pgfpathlineto{\pgfqpoint{1.455844in}{0.993978in}}%
\pgfpathlineto{\pgfqpoint{1.448072in}{0.986936in}}%
\pgfpathlineto{\pgfqpoint{1.440283in}{0.980259in}}%
\pgfpathclose%
\pgfusepath{fill}%
\end{pgfscope}%
\begin{pgfscope}%
\pgfpathrectangle{\pgfqpoint{0.407592in}{0.222904in}}{\pgfqpoint{1.411163in}{1.411163in}}%
\pgfusepath{clip}%
\pgfsetbuttcap%
\pgfsetroundjoin%
\definecolor{currentfill}{rgb}{0.909804,0.698039,0.423529}%
\pgfsetfillcolor{currentfill}%
\pgfsetfillopacity{0.120000}%
\pgfsetlinewidth{0.000000pt}%
\definecolor{currentstroke}{rgb}{0.000000,0.000000,0.000000}%
\pgfsetstrokecolor{currentstroke}%
\pgfsetdash{}{0pt}%
\pgfpathmoveto{\pgfqpoint{0.770919in}{0.949521in}}%
\pgfpathlineto{\pgfqpoint{0.787425in}{0.939333in}}%
\pgfpathlineto{\pgfqpoint{0.803951in}{0.929787in}}%
\pgfpathlineto{\pgfqpoint{0.812910in}{0.926411in}}%
\pgfpathlineto{\pgfqpoint{0.821827in}{0.923305in}}%
\pgfpathlineto{\pgfqpoint{0.805348in}{0.933045in}}%
\pgfpathlineto{\pgfqpoint{0.788888in}{0.943444in}}%
\pgfpathlineto{\pgfqpoint{0.779924in}{0.946347in}}%
\pgfpathlineto{\pgfqpoint{0.770919in}{0.949521in}}%
\pgfpathclose%
\pgfusepath{fill}%
\end{pgfscope}%
\begin{pgfscope}%
\pgfpathrectangle{\pgfqpoint{0.407592in}{0.222904in}}{\pgfqpoint{1.411163in}{1.411163in}}%
\pgfusepath{clip}%
\pgfsetbuttcap%
\pgfsetroundjoin%
\definecolor{currentfill}{rgb}{0.909804,0.698039,0.423529}%
\pgfsetfillcolor{currentfill}%
\pgfsetfillopacity{0.120000}%
\pgfsetlinewidth{0.000000pt}%
\definecolor{currentstroke}{rgb}{0.000000,0.000000,0.000000}%
\pgfsetstrokecolor{currentstroke}%
\pgfsetdash{}{0pt}%
\pgfpathmoveto{\pgfqpoint{0.569295in}{1.090412in}}%
\pgfpathlineto{\pgfqpoint{0.585754in}{1.075918in}}%
\pgfpathlineto{\pgfqpoint{0.602225in}{1.061787in}}%
\pgfpathlineto{\pgfqpoint{0.611588in}{1.058634in}}%
\pgfpathlineto{\pgfqpoint{0.620906in}{1.055712in}}%
\pgfpathlineto{\pgfqpoint{0.604480in}{1.070101in}}%
\pgfpathlineto{\pgfqpoint{0.588065in}{1.084851in}}%
\pgfpathlineto{\pgfqpoint{0.578702in}{1.087517in}}%
\pgfpathlineto{\pgfqpoint{0.569295in}{1.090412in}}%
\pgfpathclose%
\pgfusepath{fill}%
\end{pgfscope}%
\begin{pgfscope}%
\pgfpathrectangle{\pgfqpoint{0.407592in}{0.222904in}}{\pgfqpoint{1.411163in}{1.411163in}}%
\pgfusepath{clip}%
\pgfsetbuttcap%
\pgfsetroundjoin%
\definecolor{currentfill}{rgb}{0.909804,0.698039,0.423529}%
\pgfsetfillcolor{currentfill}%
\pgfsetfillopacity{0.120000}%
\pgfsetlinewidth{0.000000pt}%
\definecolor{currentstroke}{rgb}{0.000000,0.000000,0.000000}%
\pgfsetstrokecolor{currentstroke}%
\pgfsetdash{}{0pt}%
\pgfpathmoveto{\pgfqpoint{0.451544in}{1.188871in}}%
\pgfpathlineto{\pgfqpoint{0.467985in}{1.172841in}}%
\pgfpathlineto{\pgfqpoint{0.484436in}{1.157067in}}%
\pgfpathlineto{\pgfqpoint{0.494023in}{1.154232in}}%
\pgfpathlineto{\pgfqpoint{0.503564in}{1.151611in}}%
\pgfpathlineto{\pgfqpoint{0.487157in}{1.167621in}}%
\pgfpathlineto{\pgfqpoint{0.470759in}{1.183880in}}%
\pgfpathlineto{\pgfqpoint{0.461174in}{1.186270in}}%
\pgfpathlineto{\pgfqpoint{0.451544in}{1.188871in}}%
\pgfpathclose%
\pgfusepath{fill}%
\end{pgfscope}%
\begin{pgfscope}%
\pgfpathrectangle{\pgfqpoint{0.407592in}{0.222904in}}{\pgfqpoint{1.411163in}{1.411163in}}%
\pgfusepath{clip}%
\pgfsetbuttcap%
\pgfsetroundjoin%
\definecolor{currentfill}{rgb}{0.909804,0.698039,0.423529}%
\pgfsetfillcolor{currentfill}%
\pgfsetfillopacity{0.120000}%
\pgfsetlinewidth{0.000000pt}%
\definecolor{currentstroke}{rgb}{0.000000,0.000000,0.000000}%
\pgfsetstrokecolor{currentstroke}%
\pgfsetdash{}{0pt}%
\pgfpathmoveto{\pgfqpoint{1.173142in}{0.861930in}}%
\pgfpathlineto{\pgfqpoint{1.190061in}{0.866194in}}%
\pgfpathlineto{\pgfqpoint{1.207042in}{0.871119in}}%
\pgfpathlineto{\pgfqpoint{1.215194in}{0.872582in}}%
\pgfpathlineto{\pgfqpoint{1.223319in}{0.874486in}}%
\pgfpathlineto{\pgfqpoint{1.206370in}{0.869384in}}%
\pgfpathlineto{\pgfqpoint{1.189483in}{0.864945in}}%
\pgfpathlineto{\pgfqpoint{1.181326in}{0.863208in}}%
\pgfpathlineto{\pgfqpoint{1.173142in}{0.861930in}}%
\pgfpathclose%
\pgfusepath{fill}%
\end{pgfscope}%
\begin{pgfscope}%
\pgfpathrectangle{\pgfqpoint{0.407592in}{0.222904in}}{\pgfqpoint{1.411163in}{1.411163in}}%
\pgfusepath{clip}%
\pgfsetbuttcap%
\pgfsetroundjoin%
\definecolor{currentfill}{rgb}{0.909804,0.698039,0.423529}%
\pgfsetfillcolor{currentfill}%
\pgfsetfillopacity{0.120000}%
\pgfsetlinewidth{0.000000pt}%
\definecolor{currentstroke}{rgb}{0.000000,0.000000,0.000000}%
\pgfsetstrokecolor{currentstroke}%
\pgfsetdash{}{0pt}%
\pgfpathmoveto{\pgfqpoint{1.223319in}{0.874486in}}%
\pgfpathlineto{\pgfqpoint{1.240333in}{0.880205in}}%
\pgfpathlineto{\pgfqpoint{1.257414in}{0.886499in}}%
\pgfpathlineto{\pgfqpoint{1.265483in}{0.889001in}}%
\pgfpathlineto{\pgfqpoint{1.273528in}{0.891941in}}%
\pgfpathlineto{\pgfqpoint{1.256476in}{0.885511in}}%
\pgfpathlineto{\pgfqpoint{1.239491in}{0.879655in}}%
\pgfpathlineto{\pgfqpoint{1.231418in}{0.876840in}}%
\pgfpathlineto{\pgfqpoint{1.223319in}{0.874486in}}%
\pgfpathclose%
\pgfusepath{fill}%
\end{pgfscope}%
\begin{pgfscope}%
\pgfpathrectangle{\pgfqpoint{0.407592in}{0.222904in}}{\pgfqpoint{1.411163in}{1.411163in}}%
\pgfusepath{clip}%
\pgfsetbuttcap%
\pgfsetroundjoin%
\definecolor{currentfill}{rgb}{0.909804,0.698039,0.423529}%
\pgfsetfillcolor{currentfill}%
\pgfsetfillopacity{0.120000}%
\pgfsetlinewidth{0.000000pt}%
\definecolor{currentstroke}{rgb}{0.000000,0.000000,0.000000}%
\pgfsetstrokecolor{currentstroke}%
\pgfsetdash{}{0pt}%
\pgfpathmoveto{\pgfqpoint{0.686731in}{1.002256in}}%
\pgfpathlineto{\pgfqpoint{0.703220in}{0.990048in}}%
\pgfpathlineto{\pgfqpoint{0.719724in}{0.978355in}}%
\pgfpathlineto{\pgfqpoint{0.728861in}{0.974906in}}%
\pgfpathlineto{\pgfqpoint{0.737956in}{0.971704in}}%
\pgfpathlineto{\pgfqpoint{0.721498in}{0.983641in}}%
\pgfpathlineto{\pgfqpoint{0.705055in}{0.996102in}}%
\pgfpathlineto{\pgfqpoint{0.695914in}{0.999055in}}%
\pgfpathlineto{\pgfqpoint{0.686731in}{1.002256in}}%
\pgfpathclose%
\pgfusepath{fill}%
\end{pgfscope}%
\begin{pgfscope}%
\pgfpathrectangle{\pgfqpoint{0.407592in}{0.222904in}}{\pgfqpoint{1.411163in}{1.411163in}}%
\pgfusepath{clip}%
\pgfsetbuttcap%
\pgfsetroundjoin%
\definecolor{currentfill}{rgb}{0.909804,0.698039,0.423529}%
\pgfsetfillcolor{currentfill}%
\pgfsetfillopacity{0.120000}%
\pgfsetlinewidth{0.000000pt}%
\definecolor{currentstroke}{rgb}{0.000000,0.000000,0.000000}%
\pgfsetstrokecolor{currentstroke}%
\pgfsetdash{}{0pt}%
\pgfpathmoveto{\pgfqpoint{0.938609in}{0.874211in}}%
\pgfpathlineto{\pgfqpoint{0.955219in}{0.869767in}}%
\pgfpathlineto{\pgfqpoint{0.971863in}{0.866165in}}%
\pgfpathlineto{\pgfqpoint{0.980475in}{0.863655in}}%
\pgfpathlineto{\pgfqpoint{0.989049in}{0.861483in}}%
\pgfpathlineto{\pgfqpoint{0.972450in}{0.865049in}}%
\pgfpathlineto{\pgfqpoint{0.955886in}{0.869491in}}%
\pgfpathlineto{\pgfqpoint{0.947267in}{0.871683in}}%
\pgfpathlineto{\pgfqpoint{0.938609in}{0.874211in}}%
\pgfpathclose%
\pgfusepath{fill}%
\end{pgfscope}%
\begin{pgfscope}%
\pgfpathrectangle{\pgfqpoint{0.407592in}{0.222904in}}{\pgfqpoint{1.411163in}{1.411163in}}%
\pgfusepath{clip}%
\pgfsetbuttcap%
\pgfsetroundjoin%
\definecolor{currentfill}{rgb}{0.909804,0.698039,0.423529}%
\pgfsetfillcolor{currentfill}%
\pgfsetfillopacity{0.120000}%
\pgfsetlinewidth{0.000000pt}%
\definecolor{currentstroke}{rgb}{0.000000,0.000000,0.000000}%
\pgfsetstrokecolor{currentstroke}%
\pgfsetdash{}{0pt}%
\pgfpathmoveto{\pgfqpoint{1.122953in}{0.854710in}}%
\pgfpathlineto{\pgfqpoint{1.139791in}{0.857332in}}%
\pgfpathlineto{\pgfqpoint{1.156685in}{0.860693in}}%
\pgfpathlineto{\pgfqpoint{1.164928in}{0.861096in}}%
\pgfpathlineto{\pgfqpoint{1.173142in}{0.861930in}}%
\pgfpathlineto{\pgfqpoint{1.156282in}{0.858370in}}%
\pgfpathlineto{\pgfqpoint{1.139481in}{0.855560in}}%
\pgfpathlineto{\pgfqpoint{1.131232in}{0.854914in}}%
\pgfpathlineto{\pgfqpoint{1.122953in}{0.854710in}}%
\pgfpathclose%
\pgfusepath{fill}%
\end{pgfscope}%
\begin{pgfscope}%
\pgfpathrectangle{\pgfqpoint{0.407592in}{0.222904in}}{\pgfqpoint{1.411163in}{1.411163in}}%
\pgfusepath{clip}%
\pgfsetbuttcap%
\pgfsetroundjoin%
\definecolor{currentfill}{rgb}{0.909804,0.698039,0.423529}%
\pgfsetfillcolor{currentfill}%
\pgfsetfillopacity{0.120000}%
\pgfsetlinewidth{0.000000pt}%
\definecolor{currentstroke}{rgb}{0.000000,0.000000,0.000000}%
\pgfsetstrokecolor{currentstroke}%
\pgfsetdash{}{0pt}%
\pgfpathmoveto{\pgfqpoint{1.607493in}{1.100822in}}%
\pgfpathlineto{\pgfqpoint{1.625320in}{1.111312in}}%
\pgfpathlineto{\pgfqpoint{1.643226in}{1.122030in}}%
\pgfpathlineto{\pgfqpoint{1.650804in}{1.130606in}}%
\pgfpathlineto{\pgfqpoint{1.658368in}{1.139438in}}%
\pgfpathlineto{\pgfqpoint{1.640486in}{1.128906in}}%
\pgfpathlineto{\pgfqpoint{1.622683in}{1.118607in}}%
\pgfpathlineto{\pgfqpoint{1.615095in}{1.109582in}}%
\pgfpathlineto{\pgfqpoint{1.607493in}{1.100822in}}%
\pgfpathclose%
\pgfusepath{fill}%
\end{pgfscope}%
\begin{pgfscope}%
\pgfpathrectangle{\pgfqpoint{0.407592in}{0.222904in}}{\pgfqpoint{1.411163in}{1.411163in}}%
\pgfusepath{clip}%
\pgfsetbuttcap%
\pgfsetroundjoin%
\definecolor{currentfill}{rgb}{0.909804,0.698039,0.423529}%
\pgfsetfillcolor{currentfill}%
\pgfsetfillopacity{0.120000}%
\pgfsetlinewidth{0.000000pt}%
\definecolor{currentstroke}{rgb}{0.000000,0.000000,0.000000}%
\pgfsetstrokecolor{currentstroke}%
\pgfsetdash{}{0pt}%
\pgfpathmoveto{\pgfqpoint{1.273528in}{0.891941in}}%
\pgfpathlineto{\pgfqpoint{1.290647in}{0.898901in}}%
\pgfpathlineto{\pgfqpoint{1.307836in}{0.906352in}}%
\pgfpathlineto{\pgfqpoint{1.315830in}{0.909833in}}%
\pgfpathlineto{\pgfqpoint{1.323802in}{0.913739in}}%
\pgfpathlineto{\pgfqpoint{1.306639in}{0.906204in}}%
\pgfpathlineto{\pgfqpoint{1.289546in}{0.899158in}}%
\pgfpathlineto{\pgfqpoint{1.281548in}{0.895325in}}%
\pgfpathlineto{\pgfqpoint{1.273528in}{0.891941in}}%
\pgfpathclose%
\pgfusepath{fill}%
\end{pgfscope}%
\begin{pgfscope}%
\pgfpathrectangle{\pgfqpoint{0.407592in}{0.222904in}}{\pgfqpoint{1.411163in}{1.411163in}}%
\pgfusepath{clip}%
\pgfsetbuttcap%
\pgfsetroundjoin%
\definecolor{currentfill}{rgb}{0.909804,0.698039,0.423529}%
\pgfsetfillcolor{currentfill}%
\pgfsetfillopacity{0.120000}%
\pgfsetlinewidth{0.000000pt}%
\definecolor{currentstroke}{rgb}{0.000000,0.000000,0.000000}%
\pgfsetstrokecolor{currentstroke}%
\pgfsetdash{}{0pt}%
\pgfpathmoveto{\pgfqpoint{1.724408in}{1.198109in}}%
\pgfpathlineto{\pgfqpoint{1.742481in}{1.208887in}}%
\pgfpathlineto{\pgfqpoint{1.760634in}{1.219869in}}%
\pgfpathlineto{\pgfqpoint{1.768088in}{1.229458in}}%
\pgfpathlineto{\pgfqpoint{1.775531in}{1.239254in}}%
\pgfpathlineto{\pgfqpoint{1.757402in}{1.228479in}}%
\pgfpathlineto{\pgfqpoint{1.739354in}{1.217914in}}%
\pgfpathlineto{\pgfqpoint{1.731887in}{1.207906in}}%
\pgfpathlineto{\pgfqpoint{1.724408in}{1.198109in}}%
\pgfpathclose%
\pgfusepath{fill}%
\end{pgfscope}%
\begin{pgfscope}%
\pgfpathrectangle{\pgfqpoint{0.407592in}{0.222904in}}{\pgfqpoint{1.411163in}{1.411163in}}%
\pgfusepath{clip}%
\pgfsetbuttcap%
\pgfsetroundjoin%
\definecolor{currentfill}{rgb}{0.909804,0.698039,0.423529}%
\pgfsetfillcolor{currentfill}%
\pgfsetfillopacity{0.120000}%
\pgfsetlinewidth{0.000000pt}%
\definecolor{currentstroke}{rgb}{0.000000,0.000000,0.000000}%
\pgfsetstrokecolor{currentstroke}%
\pgfsetdash{}{0pt}%
\pgfpathmoveto{\pgfqpoint{1.490786in}{1.012935in}}%
\pgfpathlineto{\pgfqpoint{1.508371in}{1.022861in}}%
\pgfpathlineto{\pgfqpoint{1.526034in}{1.033060in}}%
\pgfpathlineto{\pgfqpoint{1.533744in}{1.040207in}}%
\pgfpathlineto{\pgfqpoint{1.541439in}{1.047671in}}%
\pgfpathlineto{\pgfqpoint{1.523800in}{1.037601in}}%
\pgfpathlineto{\pgfqpoint{1.506237in}{1.027808in}}%
\pgfpathlineto{\pgfqpoint{1.498519in}{1.020205in}}%
\pgfpathlineto{\pgfqpoint{1.490786in}{1.012935in}}%
\pgfpathclose%
\pgfusepath{fill}%
\end{pgfscope}%
\begin{pgfscope}%
\pgfpathrectangle{\pgfqpoint{0.407592in}{0.222904in}}{\pgfqpoint{1.411163in}{1.411163in}}%
\pgfusepath{clip}%
\pgfsetbuttcap%
\pgfsetroundjoin%
\definecolor{currentfill}{rgb}{0.909804,0.698039,0.423529}%
\pgfsetfillcolor{currentfill}%
\pgfsetfillopacity{0.120000}%
\pgfsetlinewidth{0.000000pt}%
\definecolor{currentstroke}{rgb}{0.000000,0.000000,0.000000}%
\pgfsetstrokecolor{currentstroke}%
\pgfsetdash{}{0pt}%
\pgfpathmoveto{\pgfqpoint{1.072708in}{0.853062in}}%
\pgfpathlineto{\pgfqpoint{1.089478in}{0.853911in}}%
\pgfpathlineto{\pgfqpoint{1.106299in}{0.855560in}}%
\pgfpathlineto{\pgfqpoint{1.114643in}{0.854932in}}%
\pgfpathlineto{\pgfqpoint{1.122953in}{0.854710in}}%
\pgfpathlineto{\pgfqpoint{1.106171in}{0.852869in}}%
\pgfpathlineto{\pgfqpoint{1.089440in}{0.851845in}}%
\pgfpathlineto{\pgfqpoint{1.081091in}{0.852247in}}%
\pgfpathlineto{\pgfqpoint{1.072708in}{0.853062in}}%
\pgfpathclose%
\pgfusepath{fill}%
\end{pgfscope}%
\begin{pgfscope}%
\pgfpathrectangle{\pgfqpoint{0.407592in}{0.222904in}}{\pgfqpoint{1.411163in}{1.411163in}}%
\pgfusepath{clip}%
\pgfsetbuttcap%
\pgfsetroundjoin%
\definecolor{currentfill}{rgb}{0.909804,0.698039,0.423529}%
\pgfsetfillcolor{currentfill}%
\pgfsetfillopacity{0.120000}%
\pgfsetlinewidth{0.000000pt}%
\definecolor{currentstroke}{rgb}{0.000000,0.000000,0.000000}%
\pgfsetstrokecolor{currentstroke}%
\pgfsetdash{}{0pt}%
\pgfpathmoveto{\pgfqpoint{1.775531in}{1.239254in}}%
\pgfpathlineto{\pgfqpoint{1.793740in}{1.250228in}}%
\pgfpathlineto{\pgfqpoint{1.801158in}{1.260119in}}%
\pgfpathlineto{\pgfqpoint{1.808565in}{1.270203in}}%
\pgfpathlineto{\pgfqpoint{1.790381in}{1.259440in}}%
\pgfpathlineto{\pgfqpoint{1.782961in}{1.249249in}}%
\pgfusepath{fill}%
\end{pgfscope}%
\begin{pgfscope}%
\pgfpathrectangle{\pgfqpoint{0.407592in}{0.222904in}}{\pgfqpoint{1.411163in}{1.411163in}}%
\pgfusepath{clip}%
\pgfsetbuttcap%
\pgfsetroundjoin%
\definecolor{currentfill}{rgb}{0.909804,0.698039,0.423529}%
\pgfsetfillcolor{currentfill}%
\pgfsetfillopacity{0.120000}%
\pgfsetlinewidth{0.000000pt}%
\definecolor{currentstroke}{rgb}{0.000000,0.000000,0.000000}%
\pgfsetstrokecolor{currentstroke}%
\pgfsetdash{}{0pt}%
\pgfpathmoveto{\pgfqpoint{1.323802in}{0.913739in}}%
\pgfpathlineto{\pgfqpoint{1.341036in}{0.921725in}}%
\pgfpathlineto{\pgfqpoint{1.358341in}{0.930130in}}%
\pgfpathlineto{\pgfqpoint{1.366266in}{0.934506in}}%
\pgfpathlineto{\pgfqpoint{1.374170in}{0.939287in}}%
\pgfpathlineto{\pgfqpoint{1.356889in}{0.930853in}}%
\pgfpathlineto{\pgfqpoint{1.339680in}{0.922834in}}%
\pgfpathlineto{\pgfqpoint{1.331751in}{0.918072in}}%
\pgfpathlineto{\pgfqpoint{1.323802in}{0.913739in}}%
\pgfpathclose%
\pgfusepath{fill}%
\end{pgfscope}%
\begin{pgfscope}%
\pgfpathrectangle{\pgfqpoint{0.407592in}{0.222904in}}{\pgfqpoint{1.411163in}{1.411163in}}%
\pgfusepath{clip}%
\pgfsetbuttcap%
\pgfsetroundjoin%
\definecolor{currentfill}{rgb}{0.909804,0.698039,0.423529}%
\pgfsetfillcolor{currentfill}%
\pgfsetfillopacity{0.120000}%
\pgfsetlinewidth{0.000000pt}%
\definecolor{currentstroke}{rgb}{0.000000,0.000000,0.000000}%
\pgfsetstrokecolor{currentstroke}%
\pgfsetdash{}{0pt}%
\pgfpathmoveto{\pgfqpoint{0.602225in}{1.061787in}}%
\pgfpathlineto{\pgfqpoint{0.618708in}{1.048043in}}%
\pgfpathlineto{\pgfqpoint{0.635203in}{1.034708in}}%
\pgfpathlineto{\pgfqpoint{0.644520in}{1.031296in}}%
\pgfpathlineto{\pgfqpoint{0.653793in}{1.028113in}}%
\pgfpathlineto{\pgfqpoint{0.637343in}{1.041707in}}%
\pgfpathlineto{\pgfqpoint{0.620906in}{1.055712in}}%
\pgfpathlineto{\pgfqpoint{0.611588in}{1.058634in}}%
\pgfpathlineto{\pgfqpoint{0.602225in}{1.061787in}}%
\pgfpathclose%
\pgfusepath{fill}%
\end{pgfscope}%
\begin{pgfscope}%
\pgfpathrectangle{\pgfqpoint{0.407592in}{0.222904in}}{\pgfqpoint{1.411163in}{1.411163in}}%
\pgfusepath{clip}%
\pgfsetbuttcap%
\pgfsetroundjoin%
\definecolor{currentfill}{rgb}{0.909804,0.698039,0.423529}%
\pgfsetfillcolor{currentfill}%
\pgfsetfillopacity{0.120000}%
\pgfsetlinewidth{0.000000pt}%
\definecolor{currentstroke}{rgb}{0.000000,0.000000,0.000000}%
\pgfsetstrokecolor{currentstroke}%
\pgfsetdash{}{0pt}%
\pgfpathmoveto{\pgfqpoint{0.484436in}{1.157067in}}%
\pgfpathlineto{\pgfqpoint{0.500898in}{1.141564in}}%
\pgfpathlineto{\pgfqpoint{0.517369in}{1.126350in}}%
\pgfpathlineto{\pgfqpoint{0.526912in}{1.123271in}}%
\pgfpathlineto{\pgfqpoint{0.536408in}{1.120407in}}%
\pgfpathlineto{\pgfqpoint{0.519981in}{1.135867in}}%
\pgfpathlineto{\pgfqpoint{0.503564in}{1.151611in}}%
\pgfpathlineto{\pgfqpoint{0.494023in}{1.154232in}}%
\pgfpathlineto{\pgfqpoint{0.484436in}{1.157067in}}%
\pgfpathclose%
\pgfusepath{fill}%
\end{pgfscope}%
\begin{pgfscope}%
\pgfpathrectangle{\pgfqpoint{0.407592in}{0.222904in}}{\pgfqpoint{1.411163in}{1.411163in}}%
\pgfusepath{clip}%
\pgfsetbuttcap%
\pgfsetroundjoin%
\definecolor{currentfill}{rgb}{0.909804,0.698039,0.423529}%
\pgfsetfillcolor{currentfill}%
\pgfsetfillopacity{0.120000}%
\pgfsetlinewidth{0.000000pt}%
\definecolor{currentstroke}{rgb}{0.000000,0.000000,0.000000}%
\pgfsetstrokecolor{currentstroke}%
\pgfsetdash{}{0pt}%
\pgfpathmoveto{\pgfqpoint{0.803951in}{0.929787in}}%
\pgfpathlineto{\pgfqpoint{0.820497in}{0.920909in}}%
\pgfpathlineto{\pgfqpoint{0.837064in}{0.912725in}}%
\pgfpathlineto{\pgfqpoint{0.845975in}{0.909188in}}%
\pgfpathlineto{\pgfqpoint{0.854846in}{0.905919in}}%
\pgfpathlineto{\pgfqpoint{0.838326in}{0.914253in}}%
\pgfpathlineto{\pgfqpoint{0.821827in}{0.923305in}}%
\pgfpathlineto{\pgfqpoint{0.812910in}{0.926411in}}%
\pgfpathlineto{\pgfqpoint{0.803951in}{0.929787in}}%
\pgfpathclose%
\pgfusepath{fill}%
\end{pgfscope}%
\begin{pgfscope}%
\pgfpathrectangle{\pgfqpoint{0.407592in}{0.222904in}}{\pgfqpoint{1.411163in}{1.411163in}}%
\pgfusepath{clip}%
\pgfsetbuttcap%
\pgfsetroundjoin%
\definecolor{currentfill}{rgb}{0.909804,0.698039,0.423529}%
\pgfsetfillcolor{currentfill}%
\pgfsetfillopacity{0.120000}%
\pgfsetlinewidth{0.000000pt}%
\definecolor{currentstroke}{rgb}{0.000000,0.000000,0.000000}%
\pgfsetstrokecolor{currentstroke}%
\pgfsetdash{}{0pt}%
\pgfpathmoveto{\pgfqpoint{1.022360in}{0.856968in}}%
\pgfpathlineto{\pgfqpoint{1.039078in}{0.855995in}}%
\pgfpathlineto{\pgfqpoint{1.055841in}{0.855855in}}%
\pgfpathlineto{\pgfqpoint{1.064292in}{0.854271in}}%
\pgfpathlineto{\pgfqpoint{1.072708in}{0.853062in}}%
\pgfpathlineto{\pgfqpoint{1.055987in}{0.853044in}}%
\pgfpathlineto{\pgfqpoint{1.039311in}{0.853885in}}%
\pgfpathlineto{\pgfqpoint{1.030853in}{0.855238in}}%
\pgfpathlineto{\pgfqpoint{1.022360in}{0.856968in}}%
\pgfpathclose%
\pgfusepath{fill}%
\end{pgfscope}%
\begin{pgfscope}%
\pgfpathrectangle{\pgfqpoint{0.407592in}{0.222904in}}{\pgfqpoint{1.411163in}{1.411163in}}%
\pgfusepath{clip}%
\pgfsetbuttcap%
\pgfsetroundjoin%
\definecolor{currentfill}{rgb}{0.909804,0.698039,0.423529}%
\pgfsetfillcolor{currentfill}%
\pgfsetfillopacity{0.120000}%
\pgfsetlinewidth{0.000000pt}%
\definecolor{currentstroke}{rgb}{0.000000,0.000000,0.000000}%
\pgfsetstrokecolor{currentstroke}%
\pgfsetdash{}{0pt}%
\pgfpathmoveto{\pgfqpoint{1.658368in}{1.139438in}}%
\pgfpathlineto{\pgfqpoint{1.676330in}{1.150192in}}%
\pgfpathlineto{\pgfqpoint{1.694371in}{1.161157in}}%
\pgfpathlineto{\pgfqpoint{1.701899in}{1.170048in}}%
\pgfpathlineto{\pgfqpoint{1.709414in}{1.179174in}}%
\pgfpathlineto{\pgfqpoint{1.691397in}{1.168401in}}%
\pgfpathlineto{\pgfqpoint{1.673460in}{1.157845in}}%
\pgfpathlineto{\pgfqpoint{1.665920in}{1.148520in}}%
\pgfpathlineto{\pgfqpoint{1.658368in}{1.139438in}}%
\pgfpathclose%
\pgfusepath{fill}%
\end{pgfscope}%
\begin{pgfscope}%
\pgfpathrectangle{\pgfqpoint{0.407592in}{0.222904in}}{\pgfqpoint{1.411163in}{1.411163in}}%
\pgfusepath{clip}%
\pgfsetbuttcap%
\pgfsetroundjoin%
\definecolor{currentfill}{rgb}{0.909804,0.698039,0.423529}%
\pgfsetfillcolor{currentfill}%
\pgfsetfillopacity{0.120000}%
\pgfsetlinewidth{0.000000pt}%
\definecolor{currentstroke}{rgb}{0.000000,0.000000,0.000000}%
\pgfsetstrokecolor{currentstroke}%
\pgfsetdash{}{0pt}%
\pgfpathmoveto{\pgfqpoint{1.374170in}{0.939287in}}%
\pgfpathlineto{\pgfqpoint{1.391524in}{0.948107in}}%
\pgfpathlineto{\pgfqpoint{1.408953in}{0.957283in}}%
\pgfpathlineto{\pgfqpoint{1.416813in}{0.962460in}}%
\pgfpathlineto{\pgfqpoint{1.424654in}{0.968017in}}%
\pgfpathlineto{\pgfqpoint{1.407250in}{0.958860in}}%
\pgfpathlineto{\pgfqpoint{1.389919in}{0.950059in}}%
\pgfpathlineto{\pgfqpoint{1.382054in}{0.944472in}}%
\pgfpathlineto{\pgfqpoint{1.374170in}{0.939287in}}%
\pgfpathclose%
\pgfusepath{fill}%
\end{pgfscope}%
\begin{pgfscope}%
\pgfpathrectangle{\pgfqpoint{0.407592in}{0.222904in}}{\pgfqpoint{1.411163in}{1.411163in}}%
\pgfusepath{clip}%
\pgfsetbuttcap%
\pgfsetroundjoin%
\definecolor{currentfill}{rgb}{0.909804,0.698039,0.423529}%
\pgfsetfillcolor{currentfill}%
\pgfsetfillopacity{0.120000}%
\pgfsetlinewidth{0.000000pt}%
\definecolor{currentstroke}{rgb}{0.000000,0.000000,0.000000}%
\pgfsetstrokecolor{currentstroke}%
\pgfsetdash{}{0pt}%
\pgfpathmoveto{\pgfqpoint{0.887957in}{0.891498in}}%
\pgfpathlineto{\pgfqpoint{0.904553in}{0.885453in}}%
\pgfpathlineto{\pgfqpoint{0.921178in}{0.880203in}}%
\pgfpathlineto{\pgfqpoint{0.929913in}{0.877057in}}%
\pgfpathlineto{\pgfqpoint{0.938609in}{0.874211in}}%
\pgfpathlineto{\pgfqpoint{0.922031in}{0.879494in}}%
\pgfpathlineto{\pgfqpoint{0.905483in}{0.885606in}}%
\pgfpathlineto{\pgfqpoint{0.896740in}{0.888402in}}%
\pgfpathlineto{\pgfqpoint{0.887957in}{0.891498in}}%
\pgfpathclose%
\pgfusepath{fill}%
\end{pgfscope}%
\begin{pgfscope}%
\pgfpathrectangle{\pgfqpoint{0.407592in}{0.222904in}}{\pgfqpoint{1.411163in}{1.411163in}}%
\pgfusepath{clip}%
\pgfsetbuttcap%
\pgfsetroundjoin%
\definecolor{currentfill}{rgb}{0.909804,0.698039,0.423529}%
\pgfsetfillcolor{currentfill}%
\pgfsetfillopacity{0.120000}%
\pgfsetlinewidth{0.000000pt}%
\definecolor{currentstroke}{rgb}{0.000000,0.000000,0.000000}%
\pgfsetstrokecolor{currentstroke}%
\pgfsetdash{}{0pt}%
\pgfpathmoveto{\pgfqpoint{1.541439in}{1.047671in}}%
\pgfpathlineto{\pgfqpoint{1.559156in}{1.058001in}}%
\pgfpathlineto{\pgfqpoint{1.576951in}{1.068576in}}%
\pgfpathlineto{\pgfqpoint{1.584608in}{1.076206in}}%
\pgfpathlineto{\pgfqpoint{1.592250in}{1.084128in}}%
\pgfpathlineto{\pgfqpoint{1.574478in}{1.073703in}}%
\pgfpathlineto{\pgfqpoint{1.556785in}{1.063528in}}%
\pgfpathlineto{\pgfqpoint{1.549119in}{1.055447in}}%
\pgfpathlineto{\pgfqpoint{1.541439in}{1.047671in}}%
\pgfpathclose%
\pgfusepath{fill}%
\end{pgfscope}%
\begin{pgfscope}%
\pgfpathrectangle{\pgfqpoint{0.407592in}{0.222904in}}{\pgfqpoint{1.411163in}{1.411163in}}%
\pgfusepath{clip}%
\pgfsetbuttcap%
\pgfsetroundjoin%
\definecolor{currentfill}{rgb}{0.909804,0.698039,0.423529}%
\pgfsetfillcolor{currentfill}%
\pgfsetfillopacity{0.120000}%
\pgfsetlinewidth{0.000000pt}%
\definecolor{currentstroke}{rgb}{0.000000,0.000000,0.000000}%
\pgfsetstrokecolor{currentstroke}%
\pgfsetdash{}{0pt}%
\pgfpathmoveto{\pgfqpoint{0.719724in}{0.978355in}}%
\pgfpathlineto{\pgfqpoint{0.736244in}{0.967207in}}%
\pgfpathlineto{\pgfqpoint{0.752780in}{0.956630in}}%
\pgfpathlineto{\pgfqpoint{0.761871in}{0.952952in}}%
\pgfpathlineto{\pgfqpoint{0.770919in}{0.949521in}}%
\pgfpathlineto{\pgfqpoint{0.754429in}{0.960321in}}%
\pgfpathlineto{\pgfqpoint{0.737956in}{0.971704in}}%
\pgfpathlineto{\pgfqpoint{0.728861in}{0.974906in}}%
\pgfpathlineto{\pgfqpoint{0.719724in}{0.978355in}}%
\pgfpathclose%
\pgfusepath{fill}%
\end{pgfscope}%
\begin{pgfscope}%
\pgfpathrectangle{\pgfqpoint{0.407592in}{0.222904in}}{\pgfqpoint{1.411163in}{1.411163in}}%
\pgfusepath{clip}%
\pgfsetbuttcap%
\pgfsetroundjoin%
\definecolor{currentfill}{rgb}{0.909804,0.698039,0.423529}%
\pgfsetfillcolor{currentfill}%
\pgfsetfillopacity{0.120000}%
\pgfsetlinewidth{0.000000pt}%
\definecolor{currentstroke}{rgb}{0.000000,0.000000,0.000000}%
\pgfsetstrokecolor{currentstroke}%
\pgfsetdash{}{0pt}%
\pgfpathmoveto{\pgfqpoint{0.517369in}{1.126350in}}%
\pgfpathlineto{\pgfqpoint{0.533852in}{1.111442in}}%
\pgfpathlineto{\pgfqpoint{0.550345in}{1.096860in}}%
\pgfpathlineto{\pgfqpoint{0.559842in}{1.093529in}}%
\pgfpathlineto{\pgfqpoint{0.569295in}{1.090412in}}%
\pgfpathlineto{\pgfqpoint{0.552846in}{1.105248in}}%
\pgfpathlineto{\pgfqpoint{0.536408in}{1.120407in}}%
\pgfpathlineto{\pgfqpoint{0.526912in}{1.123271in}}%
\pgfpathlineto{\pgfqpoint{0.517369in}{1.126350in}}%
\pgfpathclose%
\pgfusepath{fill}%
\end{pgfscope}%
\begin{pgfscope}%
\pgfpathrectangle{\pgfqpoint{0.407592in}{0.222904in}}{\pgfqpoint{1.411163in}{1.411163in}}%
\pgfusepath{clip}%
\pgfsetbuttcap%
\pgfsetroundjoin%
\definecolor{currentfill}{rgb}{0.909804,0.698039,0.423529}%
\pgfsetfillcolor{currentfill}%
\pgfsetfillopacity{0.120000}%
\pgfsetlinewidth{0.000000pt}%
\definecolor{currentstroke}{rgb}{0.000000,0.000000,0.000000}%
\pgfsetstrokecolor{currentstroke}%
\pgfsetdash{}{0pt}%
\pgfpathmoveto{\pgfqpoint{1.424654in}{0.968017in}}%
\pgfpathlineto{\pgfqpoint{1.442134in}{0.977505in}}%
\pgfpathlineto{\pgfqpoint{1.459689in}{0.987302in}}%
\pgfpathlineto{\pgfqpoint{1.467489in}{0.993184in}}%
\pgfpathlineto{\pgfqpoint{1.475271in}{0.999420in}}%
\pgfpathlineto{\pgfqpoint{1.457739in}{0.989684in}}%
\pgfpathlineto{\pgfqpoint{1.440283in}{0.980259in}}%
\pgfpathlineto{\pgfqpoint{1.432477in}{0.973951in}}%
\pgfpathlineto{\pgfqpoint{1.424654in}{0.968017in}}%
\pgfpathclose%
\pgfusepath{fill}%
\end{pgfscope}%
\begin{pgfscope}%
\pgfpathrectangle{\pgfqpoint{0.407592in}{0.222904in}}{\pgfqpoint{1.411163in}{1.411163in}}%
\pgfusepath{clip}%
\pgfsetbuttcap%
\pgfsetroundjoin%
\definecolor{currentfill}{rgb}{0.909804,0.698039,0.423529}%
\pgfsetfillcolor{currentfill}%
\pgfsetfillopacity{0.120000}%
\pgfsetlinewidth{0.000000pt}%
\definecolor{currentstroke}{rgb}{0.000000,0.000000,0.000000}%
\pgfsetstrokecolor{currentstroke}%
\pgfsetdash{}{0pt}%
\pgfpathmoveto{\pgfqpoint{0.635203in}{1.034708in}}%
\pgfpathlineto{\pgfqpoint{0.651711in}{1.021807in}}%
\pgfpathlineto{\pgfqpoint{0.668232in}{1.009363in}}%
\pgfpathlineto{\pgfqpoint{0.677503in}{1.005696in}}%
\pgfpathlineto{\pgfqpoint{0.686731in}{1.002256in}}%
\pgfpathlineto{\pgfqpoint{0.670255in}{1.014954in}}%
\pgfpathlineto{\pgfqpoint{0.653793in}{1.028113in}}%
\pgfpathlineto{\pgfqpoint{0.644520in}{1.031296in}}%
\pgfpathlineto{\pgfqpoint{0.635203in}{1.034708in}}%
\pgfpathclose%
\pgfusepath{fill}%
\end{pgfscope}%
\begin{pgfscope}%
\pgfpathrectangle{\pgfqpoint{0.407592in}{0.222904in}}{\pgfqpoint{1.411163in}{1.411163in}}%
\pgfusepath{clip}%
\pgfsetbuttcap%
\pgfsetroundjoin%
\definecolor{currentfill}{rgb}{0.909804,0.698039,0.423529}%
\pgfsetfillcolor{currentfill}%
\pgfsetfillopacity{0.120000}%
\pgfsetlinewidth{0.000000pt}%
\definecolor{currentstroke}{rgb}{0.000000,0.000000,0.000000}%
\pgfsetstrokecolor{currentstroke}%
\pgfsetdash{}{0pt}%
\pgfpathmoveto{\pgfqpoint{0.971863in}{0.866165in}}%
\pgfpathlineto{\pgfqpoint{0.988543in}{0.863405in}}%
\pgfpathlineto{\pgfqpoint{1.005263in}{0.861478in}}%
\pgfpathlineto{\pgfqpoint{1.013830in}{0.859054in}}%
\pgfpathlineto{\pgfqpoint{1.022360in}{0.856968in}}%
\pgfpathlineto{\pgfqpoint{1.005684in}{0.858792in}}%
\pgfpathlineto{\pgfqpoint{0.989049in}{0.861483in}}%
\pgfpathlineto{\pgfqpoint{0.980475in}{0.863655in}}%
\pgfpathlineto{\pgfqpoint{0.971863in}{0.866165in}}%
\pgfpathclose%
\pgfusepath{fill}%
\end{pgfscope}%
\begin{pgfscope}%
\pgfpathrectangle{\pgfqpoint{0.407592in}{0.222904in}}{\pgfqpoint{1.411163in}{1.411163in}}%
\pgfusepath{clip}%
\pgfsetbuttcap%
\pgfsetroundjoin%
\definecolor{currentfill}{rgb}{0.909804,0.698039,0.423529}%
\pgfsetfillcolor{currentfill}%
\pgfsetfillopacity{0.120000}%
\pgfsetlinewidth{0.000000pt}%
\definecolor{currentstroke}{rgb}{0.000000,0.000000,0.000000}%
\pgfsetstrokecolor{currentstroke}%
\pgfsetdash{}{0pt}%
\pgfpathmoveto{\pgfqpoint{1.592250in}{1.084128in}}%
\pgfpathlineto{\pgfqpoint{1.610100in}{1.094788in}}%
\pgfpathlineto{\pgfqpoint{1.628030in}{1.105671in}}%
\pgfpathlineto{\pgfqpoint{1.635635in}{1.113717in}}%
\pgfpathlineto{\pgfqpoint{1.643226in}{1.122030in}}%
\pgfpathlineto{\pgfqpoint{1.625320in}{1.111312in}}%
\pgfpathlineto{\pgfqpoint{1.607493in}{1.100822in}}%
\pgfpathlineto{\pgfqpoint{1.599879in}{1.092336in}}%
\pgfpathlineto{\pgfqpoint{1.592250in}{1.084128in}}%
\pgfpathclose%
\pgfusepath{fill}%
\end{pgfscope}%
\begin{pgfscope}%
\pgfpathrectangle{\pgfqpoint{0.407592in}{0.222904in}}{\pgfqpoint{1.411163in}{1.411163in}}%
\pgfusepath{clip}%
\pgfsetbuttcap%
\pgfsetroundjoin%
\definecolor{currentfill}{rgb}{0.909804,0.698039,0.423529}%
\pgfsetfillcolor{currentfill}%
\pgfsetfillopacity{0.120000}%
\pgfsetlinewidth{0.000000pt}%
\definecolor{currentstroke}{rgb}{0.000000,0.000000,0.000000}%
\pgfsetstrokecolor{currentstroke}%
\pgfsetdash{}{0pt}%
\pgfpathmoveto{\pgfqpoint{1.709414in}{1.179174in}}%
\pgfpathlineto{\pgfqpoint{1.727511in}{1.190153in}}%
\pgfpathlineto{\pgfqpoint{1.745689in}{1.201330in}}%
\pgfpathlineto{\pgfqpoint{1.753168in}{1.210491in}}%
\pgfpathlineto{\pgfqpoint{1.760634in}{1.219869in}}%
\pgfpathlineto{\pgfqpoint{1.742481in}{1.208887in}}%
\pgfpathlineto{\pgfqpoint{1.724408in}{1.198109in}}%
\pgfpathlineto{\pgfqpoint{1.716917in}{1.188529in}}%
\pgfpathlineto{\pgfqpoint{1.709414in}{1.179174in}}%
\pgfpathclose%
\pgfusepath{fill}%
\end{pgfscope}%
\begin{pgfscope}%
\pgfpathrectangle{\pgfqpoint{0.407592in}{0.222904in}}{\pgfqpoint{1.411163in}{1.411163in}}%
\pgfusepath{clip}%
\pgfsetbuttcap%
\pgfsetroundjoin%
\definecolor{currentfill}{rgb}{0.909804,0.698039,0.423529}%
\pgfsetfillcolor{currentfill}%
\pgfsetfillopacity{0.120000}%
\pgfsetlinewidth{0.000000pt}%
\definecolor{currentstroke}{rgb}{0.000000,0.000000,0.000000}%
\pgfsetstrokecolor{currentstroke}%
\pgfsetdash{}{0pt}%
\pgfpathmoveto{\pgfqpoint{1.760634in}{1.219869in}}%
\pgfpathlineto{\pgfqpoint{1.778868in}{1.231045in}}%
\pgfpathlineto{\pgfqpoint{1.786310in}{1.240535in}}%
\pgfpathlineto{\pgfqpoint{1.793740in}{1.250228in}}%
\pgfpathlineto{\pgfqpoint{1.775531in}{1.239254in}}%
\pgfpathlineto{\pgfqpoint{1.768088in}{1.229458in}}%
\pgfusepath{fill}%
\end{pgfscope}%
\begin{pgfscope}%
\pgfpathrectangle{\pgfqpoint{0.407592in}{0.222904in}}{\pgfqpoint{1.411163in}{1.411163in}}%
\pgfusepath{clip}%
\pgfsetbuttcap%
\pgfsetroundjoin%
\definecolor{currentfill}{rgb}{0.909804,0.698039,0.423529}%
\pgfsetfillcolor{currentfill}%
\pgfsetfillopacity{0.120000}%
\pgfsetlinewidth{0.000000pt}%
\definecolor{currentstroke}{rgb}{0.000000,0.000000,0.000000}%
\pgfsetstrokecolor{currentstroke}%
\pgfsetdash{}{0pt}%
\pgfpathmoveto{\pgfqpoint{1.475271in}{0.999420in}}%
\pgfpathlineto{\pgfqpoint{1.492880in}{1.009444in}}%
\pgfpathlineto{\pgfqpoint{1.510566in}{1.019738in}}%
\pgfpathlineto{\pgfqpoint{1.518308in}{1.026235in}}%
\pgfpathlineto{\pgfqpoint{1.526034in}{1.033060in}}%
\pgfpathlineto{\pgfqpoint{1.508371in}{1.022861in}}%
\pgfpathlineto{\pgfqpoint{1.490786in}{1.012935in}}%
\pgfpathlineto{\pgfqpoint{1.483037in}{1.006005in}}%
\pgfpathlineto{\pgfqpoint{1.475271in}{0.999420in}}%
\pgfpathclose%
\pgfusepath{fill}%
\end{pgfscope}%
\begin{pgfscope}%
\pgfpathrectangle{\pgfqpoint{0.407592in}{0.222904in}}{\pgfqpoint{1.411163in}{1.411163in}}%
\pgfusepath{clip}%
\pgfsetbuttcap%
\pgfsetroundjoin%
\definecolor{currentfill}{rgb}{0.909804,0.698039,0.423529}%
\pgfsetfillcolor{currentfill}%
\pgfsetfillopacity{0.120000}%
\pgfsetlinewidth{0.000000pt}%
\definecolor{currentstroke}{rgb}{0.000000,0.000000,0.000000}%
\pgfsetstrokecolor{currentstroke}%
\pgfsetdash{}{0pt}%
\pgfpathmoveto{\pgfqpoint{0.837064in}{0.912725in}}%
\pgfpathlineto{\pgfqpoint{0.853655in}{0.905257in}}%
\pgfpathlineto{\pgfqpoint{0.870271in}{0.898525in}}%
\pgfpathlineto{\pgfqpoint{0.879135in}{0.894878in}}%
\pgfpathlineto{\pgfqpoint{0.887957in}{0.891498in}}%
\pgfpathlineto{\pgfqpoint{0.871389in}{0.898327in}}%
\pgfpathlineto{\pgfqpoint{0.854846in}{0.905919in}}%
\pgfpathlineto{\pgfqpoint{0.845975in}{0.909188in}}%
\pgfpathlineto{\pgfqpoint{0.837064in}{0.912725in}}%
\pgfpathclose%
\pgfusepath{fill}%
\end{pgfscope}%
\begin{pgfscope}%
\pgfpathrectangle{\pgfqpoint{0.407592in}{0.222904in}}{\pgfqpoint{1.411163in}{1.411163in}}%
\pgfusepath{clip}%
\pgfsetbuttcap%
\pgfsetroundjoin%
\definecolor{currentfill}{rgb}{0.909804,0.698039,0.423529}%
\pgfsetfillcolor{currentfill}%
\pgfsetfillopacity{0.120000}%
\pgfsetlinewidth{0.000000pt}%
\definecolor{currentstroke}{rgb}{0.000000,0.000000,0.000000}%
\pgfsetstrokecolor{currentstroke}%
\pgfsetdash{}{0pt}%
\pgfpathmoveto{\pgfqpoint{1.207042in}{0.871119in}}%
\pgfpathlineto{\pgfqpoint{1.224087in}{0.876661in}}%
\pgfpathlineto{\pgfqpoint{1.241197in}{0.882778in}}%
\pgfpathlineto{\pgfqpoint{1.249318in}{0.884428in}}%
\pgfpathlineto{\pgfqpoint{1.257414in}{0.886499in}}%
\pgfpathlineto{\pgfqpoint{1.240333in}{0.880205in}}%
\pgfpathlineto{\pgfqpoint{1.223319in}{0.874486in}}%
\pgfpathlineto{\pgfqpoint{1.215194in}{0.872582in}}%
\pgfpathlineto{\pgfqpoint{1.207042in}{0.871119in}}%
\pgfpathclose%
\pgfusepath{fill}%
\end{pgfscope}%
\begin{pgfscope}%
\pgfpathrectangle{\pgfqpoint{0.407592in}{0.222904in}}{\pgfqpoint{1.411163in}{1.411163in}}%
\pgfusepath{clip}%
\pgfsetbuttcap%
\pgfsetroundjoin%
\definecolor{currentfill}{rgb}{0.909804,0.698039,0.423529}%
\pgfsetfillcolor{currentfill}%
\pgfsetfillopacity{0.120000}%
\pgfsetlinewidth{0.000000pt}%
\definecolor{currentstroke}{rgb}{0.000000,0.000000,0.000000}%
\pgfsetstrokecolor{currentstroke}%
\pgfsetdash{}{0pt}%
\pgfpathmoveto{\pgfqpoint{0.752780in}{0.956630in}}%
\pgfpathlineto{\pgfqpoint{0.769334in}{0.946650in}}%
\pgfpathlineto{\pgfqpoint{0.785907in}{0.937294in}}%
\pgfpathlineto{\pgfqpoint{0.794950in}{0.933419in}}%
\pgfpathlineto{\pgfqpoint{0.803951in}{0.929787in}}%
\pgfpathlineto{\pgfqpoint{0.787425in}{0.939333in}}%
\pgfpathlineto{\pgfqpoint{0.770919in}{0.949521in}}%
\pgfpathlineto{\pgfqpoint{0.761871in}{0.952952in}}%
\pgfpathlineto{\pgfqpoint{0.752780in}{0.956630in}}%
\pgfpathclose%
\pgfusepath{fill}%
\end{pgfscope}%
\begin{pgfscope}%
\pgfpathrectangle{\pgfqpoint{0.407592in}{0.222904in}}{\pgfqpoint{1.411163in}{1.411163in}}%
\pgfusepath{clip}%
\pgfsetbuttcap%
\pgfsetroundjoin%
\definecolor{currentfill}{rgb}{0.909804,0.698039,0.423529}%
\pgfsetfillcolor{currentfill}%
\pgfsetfillopacity{0.120000}%
\pgfsetlinewidth{0.000000pt}%
\definecolor{currentstroke}{rgb}{0.000000,0.000000,0.000000}%
\pgfsetstrokecolor{currentstroke}%
\pgfsetdash{}{0pt}%
\pgfpathmoveto{\pgfqpoint{1.156685in}{0.860693in}}%
\pgfpathlineto{\pgfqpoint{1.173638in}{0.864751in}}%
\pgfpathlineto{\pgfqpoint{1.190652in}{0.869466in}}%
\pgfpathlineto{\pgfqpoint{1.198861in}{0.870084in}}%
\pgfpathlineto{\pgfqpoint{1.207042in}{0.871119in}}%
\pgfpathlineto{\pgfqpoint{1.190061in}{0.866194in}}%
\pgfpathlineto{\pgfqpoint{1.173142in}{0.861930in}}%
\pgfpathlineto{\pgfqpoint{1.164928in}{0.861096in}}%
\pgfpathlineto{\pgfqpoint{1.156685in}{0.860693in}}%
\pgfpathclose%
\pgfusepath{fill}%
\end{pgfscope}%
\begin{pgfscope}%
\pgfpathrectangle{\pgfqpoint{0.407592in}{0.222904in}}{\pgfqpoint{1.411163in}{1.411163in}}%
\pgfusepath{clip}%
\pgfsetbuttcap%
\pgfsetroundjoin%
\definecolor{currentfill}{rgb}{0.909804,0.698039,0.423529}%
\pgfsetfillcolor{currentfill}%
\pgfsetfillopacity{0.120000}%
\pgfsetlinewidth{0.000000pt}%
\definecolor{currentstroke}{rgb}{0.000000,0.000000,0.000000}%
\pgfsetstrokecolor{currentstroke}%
\pgfsetdash{}{0pt}%
\pgfpathmoveto{\pgfqpoint{0.550345in}{1.096860in}}%
\pgfpathlineto{\pgfqpoint{0.566849in}{1.082622in}}%
\pgfpathlineto{\pgfqpoint{0.583365in}{1.068749in}}%
\pgfpathlineto{\pgfqpoint{0.592818in}{1.065161in}}%
\pgfpathlineto{\pgfqpoint{0.602225in}{1.061787in}}%
\pgfpathlineto{\pgfqpoint{0.585754in}{1.075918in}}%
\pgfpathlineto{\pgfqpoint{0.569295in}{1.090412in}}%
\pgfpathlineto{\pgfqpoint{0.559842in}{1.093529in}}%
\pgfpathlineto{\pgfqpoint{0.550345in}{1.096860in}}%
\pgfpathclose%
\pgfusepath{fill}%
\end{pgfscope}%
\begin{pgfscope}%
\pgfpathrectangle{\pgfqpoint{0.407592in}{0.222904in}}{\pgfqpoint{1.411163in}{1.411163in}}%
\pgfusepath{clip}%
\pgfsetbuttcap%
\pgfsetroundjoin%
\definecolor{currentfill}{rgb}{0.909804,0.698039,0.423529}%
\pgfsetfillcolor{currentfill}%
\pgfsetfillopacity{0.120000}%
\pgfsetlinewidth{0.000000pt}%
\definecolor{currentstroke}{rgb}{0.000000,0.000000,0.000000}%
\pgfsetstrokecolor{currentstroke}%
\pgfsetdash{}{0pt}%
\pgfpathmoveto{\pgfqpoint{1.257414in}{0.886499in}}%
\pgfpathlineto{\pgfqpoint{1.274561in}{0.893327in}}%
\pgfpathlineto{\pgfqpoint{1.291778in}{0.900649in}}%
\pgfpathlineto{\pgfqpoint{1.299819in}{0.903293in}}%
\pgfpathlineto{\pgfqpoint{1.307836in}{0.906352in}}%
\pgfpathlineto{\pgfqpoint{1.290647in}{0.898901in}}%
\pgfpathlineto{\pgfqpoint{1.273528in}{0.891941in}}%
\pgfpathlineto{\pgfqpoint{1.265483in}{0.889001in}}%
\pgfpathlineto{\pgfqpoint{1.257414in}{0.886499in}}%
\pgfpathclose%
\pgfusepath{fill}%
\end{pgfscope}%
\begin{pgfscope}%
\pgfpathrectangle{\pgfqpoint{0.407592in}{0.222904in}}{\pgfqpoint{1.411163in}{1.411163in}}%
\pgfusepath{clip}%
\pgfsetbuttcap%
\pgfsetroundjoin%
\definecolor{currentfill}{rgb}{0.909804,0.698039,0.423529}%
\pgfsetfillcolor{currentfill}%
\pgfsetfillopacity{0.120000}%
\pgfsetlinewidth{0.000000pt}%
\definecolor{currentstroke}{rgb}{0.000000,0.000000,0.000000}%
\pgfsetstrokecolor{currentstroke}%
\pgfsetdash{}{0pt}%
\pgfpathmoveto{\pgfqpoint{1.106299in}{0.855560in}}%
\pgfpathlineto{\pgfqpoint{1.123174in}{0.857973in}}%
\pgfpathlineto{\pgfqpoint{1.140105in}{0.861114in}}%
\pgfpathlineto{\pgfqpoint{1.148411in}{0.860704in}}%
\pgfpathlineto{\pgfqpoint{1.156685in}{0.860693in}}%
\pgfpathlineto{\pgfqpoint{1.139791in}{0.857332in}}%
\pgfpathlineto{\pgfqpoint{1.122953in}{0.854710in}}%
\pgfpathlineto{\pgfqpoint{1.114643in}{0.854932in}}%
\pgfpathlineto{\pgfqpoint{1.106299in}{0.855560in}}%
\pgfpathclose%
\pgfusepath{fill}%
\end{pgfscope}%
\begin{pgfscope}%
\pgfpathrectangle{\pgfqpoint{0.407592in}{0.222904in}}{\pgfqpoint{1.411163in}{1.411163in}}%
\pgfusepath{clip}%
\pgfsetbuttcap%
\pgfsetroundjoin%
\definecolor{currentfill}{rgb}{0.909804,0.698039,0.423529}%
\pgfsetfillcolor{currentfill}%
\pgfsetfillopacity{0.120000}%
\pgfsetlinewidth{0.000000pt}%
\definecolor{currentstroke}{rgb}{0.000000,0.000000,0.000000}%
\pgfsetstrokecolor{currentstroke}%
\pgfsetdash{}{0pt}%
\pgfpathmoveto{\pgfqpoint{0.921178in}{0.880203in}}%
\pgfpathlineto{\pgfqpoint{0.937834in}{0.875758in}}%
\pgfpathlineto{\pgfqpoint{0.954524in}{0.872122in}}%
\pgfpathlineto{\pgfqpoint{0.963213in}{0.868993in}}%
\pgfpathlineto{\pgfqpoint{0.971863in}{0.866165in}}%
\pgfpathlineto{\pgfqpoint{0.955219in}{0.869767in}}%
\pgfpathlineto{\pgfqpoint{0.938609in}{0.874211in}}%
\pgfpathlineto{\pgfqpoint{0.929913in}{0.877057in}}%
\pgfpathlineto{\pgfqpoint{0.921178in}{0.880203in}}%
\pgfpathclose%
\pgfusepath{fill}%
\end{pgfscope}%
\begin{pgfscope}%
\pgfpathrectangle{\pgfqpoint{0.407592in}{0.222904in}}{\pgfqpoint{1.411163in}{1.411163in}}%
\pgfusepath{clip}%
\pgfsetbuttcap%
\pgfsetroundjoin%
\definecolor{currentfill}{rgb}{0.909804,0.698039,0.423529}%
\pgfsetfillcolor{currentfill}%
\pgfsetfillopacity{0.120000}%
\pgfsetlinewidth{0.000000pt}%
\definecolor{currentstroke}{rgb}{0.000000,0.000000,0.000000}%
\pgfsetstrokecolor{currentstroke}%
\pgfsetdash{}{0pt}%
\pgfpathmoveto{\pgfqpoint{1.307836in}{0.906352in}}%
\pgfpathlineto{\pgfqpoint{1.325096in}{0.914259in}}%
\pgfpathlineto{\pgfqpoint{1.342428in}{0.922588in}}%
\pgfpathlineto{\pgfqpoint{1.350396in}{0.926158in}}%
\pgfpathlineto{\pgfqpoint{1.358341in}{0.930130in}}%
\pgfpathlineto{\pgfqpoint{1.341036in}{0.921725in}}%
\pgfpathlineto{\pgfqpoint{1.323802in}{0.913739in}}%
\pgfpathlineto{\pgfqpoint{1.315830in}{0.909833in}}%
\pgfpathlineto{\pgfqpoint{1.307836in}{0.906352in}}%
\pgfpathclose%
\pgfusepath{fill}%
\end{pgfscope}%
\begin{pgfscope}%
\pgfpathrectangle{\pgfqpoint{0.407592in}{0.222904in}}{\pgfqpoint{1.411163in}{1.411163in}}%
\pgfusepath{clip}%
\pgfsetbuttcap%
\pgfsetroundjoin%
\definecolor{currentfill}{rgb}{0.909804,0.698039,0.423529}%
\pgfsetfillcolor{currentfill}%
\pgfsetfillopacity{0.120000}%
\pgfsetlinewidth{0.000000pt}%
\definecolor{currentstroke}{rgb}{0.000000,0.000000,0.000000}%
\pgfsetstrokecolor{currentstroke}%
\pgfsetdash{}{0pt}%
\pgfpathmoveto{\pgfqpoint{0.668232in}{1.009363in}}%
\pgfpathlineto{\pgfqpoint{0.684768in}{0.997403in}}%
\pgfpathlineto{\pgfqpoint{0.701319in}{0.985951in}}%
\pgfpathlineto{\pgfqpoint{0.710543in}{0.982041in}}%
\pgfpathlineto{\pgfqpoint{0.719724in}{0.978355in}}%
\pgfpathlineto{\pgfqpoint{0.703220in}{0.990048in}}%
\pgfpathlineto{\pgfqpoint{0.686731in}{1.002256in}}%
\pgfpathlineto{\pgfqpoint{0.677503in}{1.005696in}}%
\pgfpathlineto{\pgfqpoint{0.668232in}{1.009363in}}%
\pgfpathclose%
\pgfusepath{fill}%
\end{pgfscope}%
\begin{pgfscope}%
\pgfpathrectangle{\pgfqpoint{0.407592in}{0.222904in}}{\pgfqpoint{1.411163in}{1.411163in}}%
\pgfusepath{clip}%
\pgfsetbuttcap%
\pgfsetroundjoin%
\definecolor{currentfill}{rgb}{0.909804,0.698039,0.423529}%
\pgfsetfillcolor{currentfill}%
\pgfsetfillopacity{0.120000}%
\pgfsetlinewidth{0.000000pt}%
\definecolor{currentstroke}{rgb}{0.000000,0.000000,0.000000}%
\pgfsetstrokecolor{currentstroke}%
\pgfsetdash{}{0pt}%
\pgfpathmoveto{\pgfqpoint{1.643226in}{1.122030in}}%
\pgfpathlineto{\pgfqpoint{1.661211in}{1.132964in}}%
\pgfpathlineto{\pgfqpoint{1.679277in}{1.144104in}}%
\pgfpathlineto{\pgfqpoint{1.686830in}{1.152507in}}%
\pgfpathlineto{\pgfqpoint{1.694371in}{1.161157in}}%
\pgfpathlineto{\pgfqpoint{1.676330in}{1.150192in}}%
\pgfpathlineto{\pgfqpoint{1.658368in}{1.139438in}}%
\pgfpathlineto{\pgfqpoint{1.650804in}{1.130606in}}%
\pgfpathlineto{\pgfqpoint{1.643226in}{1.122030in}}%
\pgfpathclose%
\pgfusepath{fill}%
\end{pgfscope}%
\begin{pgfscope}%
\pgfpathrectangle{\pgfqpoint{0.407592in}{0.222904in}}{\pgfqpoint{1.411163in}{1.411163in}}%
\pgfusepath{clip}%
\pgfsetbuttcap%
\pgfsetroundjoin%
\definecolor{currentfill}{rgb}{0.909804,0.698039,0.423529}%
\pgfsetfillcolor{currentfill}%
\pgfsetfillopacity{0.120000}%
\pgfsetlinewidth{0.000000pt}%
\definecolor{currentstroke}{rgb}{0.000000,0.000000,0.000000}%
\pgfsetstrokecolor{currentstroke}%
\pgfsetdash{}{0pt}%
\pgfpathmoveto{\pgfqpoint{1.055841in}{0.855855in}}%
\pgfpathlineto{\pgfqpoint{1.072652in}{0.856521in}}%
\pgfpathlineto{\pgfqpoint{1.089513in}{0.857967in}}%
\pgfpathlineto{\pgfqpoint{1.097923in}{0.856578in}}%
\pgfpathlineto{\pgfqpoint{1.106299in}{0.855560in}}%
\pgfpathlineto{\pgfqpoint{1.089478in}{0.853911in}}%
\pgfpathlineto{\pgfqpoint{1.072708in}{0.853062in}}%
\pgfpathlineto{\pgfqpoint{1.064292in}{0.854271in}}%
\pgfpathlineto{\pgfqpoint{1.055841in}{0.855855in}}%
\pgfpathclose%
\pgfusepath{fill}%
\end{pgfscope}%
\begin{pgfscope}%
\pgfpathrectangle{\pgfqpoint{0.407592in}{0.222904in}}{\pgfqpoint{1.411163in}{1.411163in}}%
\pgfusepath{clip}%
\pgfsetbuttcap%
\pgfsetroundjoin%
\definecolor{currentfill}{rgb}{0.909804,0.698039,0.423529}%
\pgfsetfillcolor{currentfill}%
\pgfsetfillopacity{0.120000}%
\pgfsetlinewidth{0.000000pt}%
\definecolor{currentstroke}{rgb}{0.000000,0.000000,0.000000}%
\pgfsetstrokecolor{currentstroke}%
\pgfsetdash{}{0pt}%
\pgfpathmoveto{\pgfqpoint{1.526034in}{1.033060in}}%
\pgfpathlineto{\pgfqpoint{1.543774in}{1.043515in}}%
\pgfpathlineto{\pgfqpoint{1.561593in}{1.054211in}}%
\pgfpathlineto{\pgfqpoint{1.569279in}{1.061242in}}%
\pgfpathlineto{\pgfqpoint{1.576951in}{1.068576in}}%
\pgfpathlineto{\pgfqpoint{1.559156in}{1.058001in}}%
\pgfpathlineto{\pgfqpoint{1.541439in}{1.047671in}}%
\pgfpathlineto{\pgfqpoint{1.533744in}{1.040207in}}%
\pgfpathlineto{\pgfqpoint{1.526034in}{1.033060in}}%
\pgfpathclose%
\pgfusepath{fill}%
\end{pgfscope}%
\begin{pgfscope}%
\pgfpathrectangle{\pgfqpoint{0.407592in}{0.222904in}}{\pgfqpoint{1.411163in}{1.411163in}}%
\pgfusepath{clip}%
\pgfsetbuttcap%
\pgfsetroundjoin%
\definecolor{currentfill}{rgb}{0.909804,0.698039,0.423529}%
\pgfsetfillcolor{currentfill}%
\pgfsetfillopacity{0.120000}%
\pgfsetlinewidth{0.000000pt}%
\definecolor{currentstroke}{rgb}{0.000000,0.000000,0.000000}%
\pgfsetstrokecolor{currentstroke}%
\pgfsetdash{}{0pt}%
\pgfpathmoveto{\pgfqpoint{1.358341in}{0.930130in}}%
\pgfpathlineto{\pgfqpoint{1.375721in}{0.938923in}}%
\pgfpathlineto{\pgfqpoint{1.393174in}{0.948074in}}%
\pgfpathlineto{\pgfqpoint{1.401073in}{0.952488in}}%
\pgfpathlineto{\pgfqpoint{1.408953in}{0.957283in}}%
\pgfpathlineto{\pgfqpoint{1.391524in}{0.948107in}}%
\pgfpathlineto{\pgfqpoint{1.374170in}{0.939287in}}%
\pgfpathlineto{\pgfqpoint{1.366266in}{0.934506in}}%
\pgfpathlineto{\pgfqpoint{1.358341in}{0.930130in}}%
\pgfpathclose%
\pgfusepath{fill}%
\end{pgfscope}%
\begin{pgfscope}%
\pgfpathrectangle{\pgfqpoint{0.407592in}{0.222904in}}{\pgfqpoint{1.411163in}{1.411163in}}%
\pgfusepath{clip}%
\pgfsetbuttcap%
\pgfsetroundjoin%
\definecolor{currentfill}{rgb}{0.909804,0.698039,0.423529}%
\pgfsetfillcolor{currentfill}%
\pgfsetfillopacity{0.120000}%
\pgfsetlinewidth{0.000000pt}%
\definecolor{currentstroke}{rgb}{0.000000,0.000000,0.000000}%
\pgfsetstrokecolor{currentstroke}%
\pgfsetdash{}{0pt}%
\pgfpathmoveto{\pgfqpoint{0.583365in}{1.068749in}}%
\pgfpathlineto{\pgfqpoint{0.599894in}{1.055263in}}%
\pgfpathlineto{\pgfqpoint{0.616435in}{1.042185in}}%
\pgfpathlineto{\pgfqpoint{0.625841in}{1.038340in}}%
\pgfpathlineto{\pgfqpoint{0.635203in}{1.034708in}}%
\pgfpathlineto{\pgfqpoint{0.618708in}{1.048043in}}%
\pgfpathlineto{\pgfqpoint{0.602225in}{1.061787in}}%
\pgfpathlineto{\pgfqpoint{0.592818in}{1.065161in}}%
\pgfpathlineto{\pgfqpoint{0.583365in}{1.068749in}}%
\pgfpathclose%
\pgfusepath{fill}%
\end{pgfscope}%
\begin{pgfscope}%
\pgfpathrectangle{\pgfqpoint{0.407592in}{0.222904in}}{\pgfqpoint{1.411163in}{1.411163in}}%
\pgfusepath{clip}%
\pgfsetbuttcap%
\pgfsetroundjoin%
\definecolor{currentfill}{rgb}{0.909804,0.698039,0.423529}%
\pgfsetfillcolor{currentfill}%
\pgfsetfillopacity{0.120000}%
\pgfsetlinewidth{0.000000pt}%
\definecolor{currentstroke}{rgb}{0.000000,0.000000,0.000000}%
\pgfsetstrokecolor{currentstroke}%
\pgfsetdash{}{0pt}%
\pgfpathmoveto{\pgfqpoint{1.005263in}{0.861478in}}%
\pgfpathlineto{\pgfqpoint{1.022025in}{0.860373in}}%
\pgfpathlineto{\pgfqpoint{1.038832in}{0.860072in}}%
\pgfpathlineto{\pgfqpoint{1.047354in}{0.857795in}}%
\pgfpathlineto{\pgfqpoint{1.055841in}{0.855855in}}%
\pgfpathlineto{\pgfqpoint{1.039078in}{0.855995in}}%
\pgfpathlineto{\pgfqpoint{1.022360in}{0.856968in}}%
\pgfpathlineto{\pgfqpoint{1.013830in}{0.859054in}}%
\pgfpathlineto{\pgfqpoint{1.005263in}{0.861478in}}%
\pgfpathclose%
\pgfusepath{fill}%
\end{pgfscope}%
\begin{pgfscope}%
\pgfpathrectangle{\pgfqpoint{0.407592in}{0.222904in}}{\pgfqpoint{1.411163in}{1.411163in}}%
\pgfusepath{clip}%
\pgfsetbuttcap%
\pgfsetroundjoin%
\definecolor{currentfill}{rgb}{0.909804,0.698039,0.423529}%
\pgfsetfillcolor{currentfill}%
\pgfsetfillopacity{0.120000}%
\pgfsetlinewidth{0.000000pt}%
\definecolor{currentstroke}{rgb}{0.000000,0.000000,0.000000}%
\pgfsetstrokecolor{currentstroke}%
\pgfsetdash{}{0pt}%
\pgfpathmoveto{\pgfqpoint{0.785907in}{0.937294in}}%
\pgfpathlineto{\pgfqpoint{0.802501in}{0.928586in}}%
\pgfpathlineto{\pgfqpoint{0.819116in}{0.920549in}}%
\pgfpathlineto{\pgfqpoint{0.828111in}{0.916517in}}%
\pgfpathlineto{\pgfqpoint{0.837064in}{0.912725in}}%
\pgfpathlineto{\pgfqpoint{0.820497in}{0.920909in}}%
\pgfpathlineto{\pgfqpoint{0.803951in}{0.929787in}}%
\pgfpathlineto{\pgfqpoint{0.794950in}{0.933419in}}%
\pgfpathlineto{\pgfqpoint{0.785907in}{0.937294in}}%
\pgfpathclose%
\pgfusepath{fill}%
\end{pgfscope}%
\begin{pgfscope}%
\pgfpathrectangle{\pgfqpoint{0.407592in}{0.222904in}}{\pgfqpoint{1.411163in}{1.411163in}}%
\pgfusepath{clip}%
\pgfsetbuttcap%
\pgfsetroundjoin%
\definecolor{currentfill}{rgb}{0.909804,0.698039,0.423529}%
\pgfsetfillcolor{currentfill}%
\pgfsetfillopacity{0.120000}%
\pgfsetlinewidth{0.000000pt}%
\definecolor{currentstroke}{rgb}{0.000000,0.000000,0.000000}%
\pgfsetstrokecolor{currentstroke}%
\pgfsetdash{}{0pt}%
\pgfpathmoveto{\pgfqpoint{1.408953in}{0.957283in}}%
\pgfpathlineto{\pgfqpoint{1.426457in}{0.966792in}}%
\pgfpathlineto{\pgfqpoint{1.444036in}{0.976610in}}%
\pgfpathlineto{\pgfqpoint{1.451872in}{0.981776in}}%
\pgfpathlineto{\pgfqpoint{1.459689in}{0.987302in}}%
\pgfpathlineto{\pgfqpoint{1.442134in}{0.977505in}}%
\pgfpathlineto{\pgfqpoint{1.424654in}{0.968017in}}%
\pgfpathlineto{\pgfqpoint{1.416813in}{0.962460in}}%
\pgfpathlineto{\pgfqpoint{1.408953in}{0.957283in}}%
\pgfpathclose%
\pgfusepath{fill}%
\end{pgfscope}%
\begin{pgfscope}%
\pgfpathrectangle{\pgfqpoint{0.407592in}{0.222904in}}{\pgfqpoint{1.411163in}{1.411163in}}%
\pgfusepath{clip}%
\pgfsetbuttcap%
\pgfsetroundjoin%
\definecolor{currentfill}{rgb}{0.909804,0.698039,0.423529}%
\pgfsetfillcolor{currentfill}%
\pgfsetfillopacity{0.120000}%
\pgfsetlinewidth{0.000000pt}%
\definecolor{currentstroke}{rgb}{0.000000,0.000000,0.000000}%
\pgfsetstrokecolor{currentstroke}%
\pgfsetdash{}{0pt}%
\pgfpathmoveto{\pgfqpoint{1.694371in}{1.161157in}}%
\pgfpathlineto{\pgfqpoint{1.712493in}{1.172323in}}%
\pgfpathlineto{\pgfqpoint{1.730696in}{1.183680in}}%
\pgfpathlineto{\pgfqpoint{1.738199in}{1.192391in}}%
\pgfpathlineto{\pgfqpoint{1.745689in}{1.201330in}}%
\pgfpathlineto{\pgfqpoint{1.727511in}{1.190153in}}%
\pgfpathlineto{\pgfqpoint{1.709414in}{1.179174in}}%
\pgfpathlineto{\pgfqpoint{1.701899in}{1.170048in}}%
\pgfpathlineto{\pgfqpoint{1.694371in}{1.161157in}}%
\pgfpathclose%
\pgfusepath{fill}%
\end{pgfscope}%
\begin{pgfscope}%
\pgfpathrectangle{\pgfqpoint{0.407592in}{0.222904in}}{\pgfqpoint{1.411163in}{1.411163in}}%
\pgfusepath{clip}%
\pgfsetbuttcap%
\pgfsetroundjoin%
\definecolor{currentfill}{rgb}{0.909804,0.698039,0.423529}%
\pgfsetfillcolor{currentfill}%
\pgfsetfillopacity{0.120000}%
\pgfsetlinewidth{0.000000pt}%
\definecolor{currentstroke}{rgb}{0.000000,0.000000,0.000000}%
\pgfsetstrokecolor{currentstroke}%
\pgfsetdash{}{0pt}%
\pgfpathmoveto{\pgfqpoint{0.870271in}{0.898525in}}%
\pgfpathlineto{\pgfqpoint{0.886914in}{0.892544in}}%
\pgfpathlineto{\pgfqpoint{0.903587in}{0.887328in}}%
\pgfpathlineto{\pgfqpoint{0.912402in}{0.883632in}}%
\pgfpathlineto{\pgfqpoint{0.921178in}{0.880203in}}%
\pgfpathlineto{\pgfqpoint{0.904553in}{0.885453in}}%
\pgfpathlineto{\pgfqpoint{0.887957in}{0.891498in}}%
\pgfpathlineto{\pgfqpoint{0.879135in}{0.894878in}}%
\pgfpathlineto{\pgfqpoint{0.870271in}{0.898525in}}%
\pgfpathclose%
\pgfusepath{fill}%
\end{pgfscope}%
\begin{pgfscope}%
\pgfpathrectangle{\pgfqpoint{0.407592in}{0.222904in}}{\pgfqpoint{1.411163in}{1.411163in}}%
\pgfusepath{clip}%
\pgfsetbuttcap%
\pgfsetroundjoin%
\definecolor{currentfill}{rgb}{0.909804,0.698039,0.423529}%
\pgfsetfillcolor{currentfill}%
\pgfsetfillopacity{0.120000}%
\pgfsetlinewidth{0.000000pt}%
\definecolor{currentstroke}{rgb}{0.000000,0.000000,0.000000}%
\pgfsetstrokecolor{currentstroke}%
\pgfsetdash{}{0pt}%
\pgfpathmoveto{\pgfqpoint{1.576951in}{1.068576in}}%
\pgfpathlineto{\pgfqpoint{1.594825in}{1.079381in}}%
\pgfpathlineto{\pgfqpoint{1.612779in}{1.090404in}}%
\pgfpathlineto{\pgfqpoint{1.620411in}{1.097898in}}%
\pgfpathlineto{\pgfqpoint{1.628030in}{1.105671in}}%
\pgfpathlineto{\pgfqpoint{1.610100in}{1.094788in}}%
\pgfpathlineto{\pgfqpoint{1.592250in}{1.084128in}}%
\pgfpathlineto{\pgfqpoint{1.584608in}{1.076206in}}%
\pgfpathlineto{\pgfqpoint{1.576951in}{1.068576in}}%
\pgfpathclose%
\pgfusepath{fill}%
\end{pgfscope}%
\begin{pgfscope}%
\pgfpathrectangle{\pgfqpoint{0.407592in}{0.222904in}}{\pgfqpoint{1.411163in}{1.411163in}}%
\pgfusepath{clip}%
\pgfsetbuttcap%
\pgfsetroundjoin%
\definecolor{currentfill}{rgb}{0.909804,0.698039,0.423529}%
\pgfsetfillcolor{currentfill}%
\pgfsetfillopacity{0.120000}%
\pgfsetlinewidth{0.000000pt}%
\definecolor{currentstroke}{rgb}{0.000000,0.000000,0.000000}%
\pgfsetstrokecolor{currentstroke}%
\pgfsetdash{}{0pt}%
\pgfpathmoveto{\pgfqpoint{0.701319in}{0.985951in}}%
\pgfpathlineto{\pgfqpoint{0.717886in}{0.975033in}}%
\pgfpathlineto{\pgfqpoint{0.734469in}{0.964673in}}%
\pgfpathlineto{\pgfqpoint{0.743646in}{0.960540in}}%
\pgfpathlineto{\pgfqpoint{0.752780in}{0.956630in}}%
\pgfpathlineto{\pgfqpoint{0.736244in}{0.967207in}}%
\pgfpathlineto{\pgfqpoint{0.719724in}{0.978355in}}%
\pgfpathlineto{\pgfqpoint{0.710543in}{0.982041in}}%
\pgfpathlineto{\pgfqpoint{0.701319in}{0.985951in}}%
\pgfpathclose%
\pgfusepath{fill}%
\end{pgfscope}%
\begin{pgfscope}%
\pgfpathrectangle{\pgfqpoint{0.407592in}{0.222904in}}{\pgfqpoint{1.411163in}{1.411163in}}%
\pgfusepath{clip}%
\pgfsetbuttcap%
\pgfsetroundjoin%
\definecolor{currentfill}{rgb}{0.909804,0.698039,0.423529}%
\pgfsetfillcolor{currentfill}%
\pgfsetfillopacity{0.120000}%
\pgfsetlinewidth{0.000000pt}%
\definecolor{currentstroke}{rgb}{0.000000,0.000000,0.000000}%
\pgfsetstrokecolor{currentstroke}%
\pgfsetdash{}{0pt}%
\pgfpathmoveto{\pgfqpoint{1.745689in}{1.201330in}}%
\pgfpathlineto{\pgfqpoint{1.763949in}{1.212696in}}%
\pgfpathlineto{\pgfqpoint{1.771414in}{1.221764in}}%
\pgfpathlineto{\pgfqpoint{1.778868in}{1.231045in}}%
\pgfpathlineto{\pgfqpoint{1.760634in}{1.219869in}}%
\pgfpathlineto{\pgfqpoint{1.753168in}{1.210491in}}%
\pgfusepath{fill}%
\end{pgfscope}%
\begin{pgfscope}%
\pgfpathrectangle{\pgfqpoint{0.407592in}{0.222904in}}{\pgfqpoint{1.411163in}{1.411163in}}%
\pgfusepath{clip}%
\pgfsetbuttcap%
\pgfsetroundjoin%
\definecolor{currentfill}{rgb}{0.909804,0.698039,0.423529}%
\pgfsetfillcolor{currentfill}%
\pgfsetfillopacity{0.120000}%
\pgfsetlinewidth{0.000000pt}%
\definecolor{currentstroke}{rgb}{0.000000,0.000000,0.000000}%
\pgfsetstrokecolor{currentstroke}%
\pgfsetdash{}{0pt}%
\pgfpathmoveto{\pgfqpoint{1.459689in}{0.987302in}}%
\pgfpathlineto{\pgfqpoint{1.477322in}{0.997386in}}%
\pgfpathlineto{\pgfqpoint{1.495032in}{1.007740in}}%
\pgfpathlineto{\pgfqpoint{1.502807in}{1.013572in}}%
\pgfpathlineto{\pgfqpoint{1.510566in}{1.019738in}}%
\pgfpathlineto{\pgfqpoint{1.492880in}{1.009444in}}%
\pgfpathlineto{\pgfqpoint{1.475271in}{0.999420in}}%
\pgfpathlineto{\pgfqpoint{1.467489in}{0.993184in}}%
\pgfpathlineto{\pgfqpoint{1.459689in}{0.987302in}}%
\pgfpathclose%
\pgfusepath{fill}%
\end{pgfscope}%
\begin{pgfscope}%
\pgfpathrectangle{\pgfqpoint{0.407592in}{0.222904in}}{\pgfqpoint{1.411163in}{1.411163in}}%
\pgfusepath{clip}%
\pgfsetbuttcap%
\pgfsetroundjoin%
\definecolor{currentfill}{rgb}{0.909804,0.698039,0.423529}%
\pgfsetfillcolor{currentfill}%
\pgfsetfillopacity{0.120000}%
\pgfsetlinewidth{0.000000pt}%
\definecolor{currentstroke}{rgb}{0.000000,0.000000,0.000000}%
\pgfsetstrokecolor{currentstroke}%
\pgfsetdash{}{0pt}%
\pgfpathmoveto{\pgfqpoint{0.616435in}{1.042185in}}%
\pgfpathlineto{\pgfqpoint{0.632989in}{1.029537in}}%
\pgfpathlineto{\pgfqpoint{0.649557in}{1.017343in}}%
\pgfpathlineto{\pgfqpoint{0.658917in}{1.013249in}}%
\pgfpathlineto{\pgfqpoint{0.668232in}{1.009363in}}%
\pgfpathlineto{\pgfqpoint{0.651711in}{1.021807in}}%
\pgfpathlineto{\pgfqpoint{0.635203in}{1.034708in}}%
\pgfpathlineto{\pgfqpoint{0.625841in}{1.038340in}}%
\pgfpathlineto{\pgfqpoint{0.616435in}{1.042185in}}%
\pgfpathclose%
\pgfusepath{fill}%
\end{pgfscope}%
\begin{pgfscope}%
\pgfpathrectangle{\pgfqpoint{0.407592in}{0.222904in}}{\pgfqpoint{1.411163in}{1.411163in}}%
\pgfusepath{clip}%
\pgfsetbuttcap%
\pgfsetroundjoin%
\definecolor{currentfill}{rgb}{0.909804,0.698039,0.423529}%
\pgfsetfillcolor{currentfill}%
\pgfsetfillopacity{0.120000}%
\pgfsetlinewidth{0.000000pt}%
\definecolor{currentstroke}{rgb}{0.000000,0.000000,0.000000}%
\pgfsetstrokecolor{currentstroke}%
\pgfsetdash{}{0pt}%
\pgfpathmoveto{\pgfqpoint{0.954524in}{0.872122in}}%
\pgfpathlineto{\pgfqpoint{0.971251in}{0.869294in}}%
\pgfpathlineto{\pgfqpoint{0.988017in}{0.867267in}}%
\pgfpathlineto{\pgfqpoint{0.996659in}{0.864222in}}%
\pgfpathlineto{\pgfqpoint{1.005263in}{0.861478in}}%
\pgfpathlineto{\pgfqpoint{0.988543in}{0.863405in}}%
\pgfpathlineto{\pgfqpoint{0.971863in}{0.866165in}}%
\pgfpathlineto{\pgfqpoint{0.963213in}{0.868993in}}%
\pgfpathlineto{\pgfqpoint{0.954524in}{0.872122in}}%
\pgfpathclose%
\pgfusepath{fill}%
\end{pgfscope}%
\begin{pgfscope}%
\pgfpathrectangle{\pgfqpoint{0.407592in}{0.222904in}}{\pgfqpoint{1.411163in}{1.411163in}}%
\pgfusepath{clip}%
\pgfsetbuttcap%
\pgfsetroundjoin%
\definecolor{currentfill}{rgb}{0.909804,0.698039,0.423529}%
\pgfsetfillcolor{currentfill}%
\pgfsetfillopacity{0.120000}%
\pgfsetlinewidth{0.000000pt}%
\definecolor{currentstroke}{rgb}{0.000000,0.000000,0.000000}%
\pgfsetstrokecolor{currentstroke}%
\pgfsetdash{}{0pt}%
\pgfpathmoveto{\pgfqpoint{1.628030in}{1.105671in}}%
\pgfpathlineto{\pgfqpoint{1.646040in}{1.116764in}}%
\pgfpathlineto{\pgfqpoint{1.664130in}{1.128058in}}%
\pgfpathlineto{\pgfqpoint{1.671710in}{1.135952in}}%
\pgfpathlineto{\pgfqpoint{1.679277in}{1.144104in}}%
\pgfpathlineto{\pgfqpoint{1.661211in}{1.132964in}}%
\pgfpathlineto{\pgfqpoint{1.643226in}{1.122030in}}%
\pgfpathlineto{\pgfqpoint{1.635635in}{1.113717in}}%
\pgfpathlineto{\pgfqpoint{1.628030in}{1.105671in}}%
\pgfpathclose%
\pgfusepath{fill}%
\end{pgfscope}%
\begin{pgfscope}%
\pgfpathrectangle{\pgfqpoint{0.407592in}{0.222904in}}{\pgfqpoint{1.411163in}{1.411163in}}%
\pgfusepath{clip}%
\pgfsetbuttcap%
\pgfsetroundjoin%
\definecolor{currentfill}{rgb}{0.909804,0.698039,0.423529}%
\pgfsetfillcolor{currentfill}%
\pgfsetfillopacity{0.120000}%
\pgfsetlinewidth{0.000000pt}%
\definecolor{currentstroke}{rgb}{0.000000,0.000000,0.000000}%
\pgfsetstrokecolor{currentstroke}%
\pgfsetdash{}{0pt}%
\pgfpathmoveto{\pgfqpoint{1.190652in}{0.869466in}}%
\pgfpathlineto{\pgfqpoint{1.207730in}{0.874796in}}%
\pgfpathlineto{\pgfqpoint{1.224872in}{0.880702in}}%
\pgfpathlineto{\pgfqpoint{1.233048in}{0.881540in}}%
\pgfpathlineto{\pgfqpoint{1.241197in}{0.882778in}}%
\pgfpathlineto{\pgfqpoint{1.224087in}{0.876661in}}%
\pgfpathlineto{\pgfqpoint{1.207042in}{0.871119in}}%
\pgfpathlineto{\pgfqpoint{1.198861in}{0.870084in}}%
\pgfpathlineto{\pgfqpoint{1.190652in}{0.869466in}}%
\pgfpathclose%
\pgfusepath{fill}%
\end{pgfscope}%
\begin{pgfscope}%
\pgfpathrectangle{\pgfqpoint{0.407592in}{0.222904in}}{\pgfqpoint{1.411163in}{1.411163in}}%
\pgfusepath{clip}%
\pgfsetbuttcap%
\pgfsetroundjoin%
\definecolor{currentfill}{rgb}{0.909804,0.698039,0.423529}%
\pgfsetfillcolor{currentfill}%
\pgfsetfillopacity{0.120000}%
\pgfsetlinewidth{0.000000pt}%
\definecolor{currentstroke}{rgb}{0.000000,0.000000,0.000000}%
\pgfsetstrokecolor{currentstroke}%
\pgfsetdash{}{0pt}%
\pgfpathmoveto{\pgfqpoint{1.241197in}{0.882778in}}%
\pgfpathlineto{\pgfqpoint{1.258375in}{0.889431in}}%
\pgfpathlineto{\pgfqpoint{1.275621in}{0.896583in}}%
\pgfpathlineto{\pgfqpoint{1.283712in}{0.898415in}}%
\pgfpathlineto{\pgfqpoint{1.291778in}{0.900649in}}%
\pgfpathlineto{\pgfqpoint{1.274561in}{0.893327in}}%
\pgfpathlineto{\pgfqpoint{1.257414in}{0.886499in}}%
\pgfpathlineto{\pgfqpoint{1.249318in}{0.884428in}}%
\pgfpathlineto{\pgfqpoint{1.241197in}{0.882778in}}%
\pgfpathclose%
\pgfusepath{fill}%
\end{pgfscope}%
\begin{pgfscope}%
\pgfpathrectangle{\pgfqpoint{0.407592in}{0.222904in}}{\pgfqpoint{1.411163in}{1.411163in}}%
\pgfusepath{clip}%
\pgfsetbuttcap%
\pgfsetroundjoin%
\definecolor{currentfill}{rgb}{0.909804,0.698039,0.423529}%
\pgfsetfillcolor{currentfill}%
\pgfsetfillopacity{0.120000}%
\pgfsetlinewidth{0.000000pt}%
\definecolor{currentstroke}{rgb}{0.000000,0.000000,0.000000}%
\pgfsetstrokecolor{currentstroke}%
\pgfsetdash{}{0pt}%
\pgfpathmoveto{\pgfqpoint{1.140105in}{0.861114in}}%
\pgfpathlineto{\pgfqpoint{1.157094in}{0.864943in}}%
\pgfpathlineto{\pgfqpoint{1.174144in}{0.869424in}}%
\pgfpathlineto{\pgfqpoint{1.182413in}{0.869250in}}%
\pgfpathlineto{\pgfqpoint{1.190652in}{0.869466in}}%
\pgfpathlineto{\pgfqpoint{1.173638in}{0.864751in}}%
\pgfpathlineto{\pgfqpoint{1.156685in}{0.860693in}}%
\pgfpathlineto{\pgfqpoint{1.148411in}{0.860704in}}%
\pgfpathlineto{\pgfqpoint{1.140105in}{0.861114in}}%
\pgfpathclose%
\pgfusepath{fill}%
\end{pgfscope}%
\begin{pgfscope}%
\pgfpathrectangle{\pgfqpoint{0.407592in}{0.222904in}}{\pgfqpoint{1.411163in}{1.411163in}}%
\pgfusepath{clip}%
\pgfsetbuttcap%
\pgfsetroundjoin%
\definecolor{currentfill}{rgb}{0.909804,0.698039,0.423529}%
\pgfsetfillcolor{currentfill}%
\pgfsetfillopacity{0.120000}%
\pgfsetlinewidth{0.000000pt}%
\definecolor{currentstroke}{rgb}{0.000000,0.000000,0.000000}%
\pgfsetstrokecolor{currentstroke}%
\pgfsetdash{}{0pt}%
\pgfpathmoveto{\pgfqpoint{1.291778in}{0.900649in}}%
\pgfpathlineto{\pgfqpoint{1.309066in}{0.908431in}}%
\pgfpathlineto{\pgfqpoint{1.326425in}{0.916638in}}%
\pgfpathlineto{\pgfqpoint{1.334438in}{0.919416in}}%
\pgfpathlineto{\pgfqpoint{1.342428in}{0.922588in}}%
\pgfpathlineto{\pgfqpoint{1.325096in}{0.914259in}}%
\pgfpathlineto{\pgfqpoint{1.307836in}{0.906352in}}%
\pgfpathlineto{\pgfqpoint{1.299819in}{0.903293in}}%
\pgfpathlineto{\pgfqpoint{1.291778in}{0.900649in}}%
\pgfpathclose%
\pgfusepath{fill}%
\end{pgfscope}%
\begin{pgfscope}%
\pgfpathrectangle{\pgfqpoint{0.407592in}{0.222904in}}{\pgfqpoint{1.411163in}{1.411163in}}%
\pgfusepath{clip}%
\pgfsetbuttcap%
\pgfsetroundjoin%
\definecolor{currentfill}{rgb}{0.909804,0.698039,0.423529}%
\pgfsetfillcolor{currentfill}%
\pgfsetfillopacity{0.120000}%
\pgfsetlinewidth{0.000000pt}%
\definecolor{currentstroke}{rgb}{0.000000,0.000000,0.000000}%
\pgfsetstrokecolor{currentstroke}%
\pgfsetdash{}{0pt}%
\pgfpathmoveto{\pgfqpoint{1.510566in}{1.019738in}}%
\pgfpathlineto{\pgfqpoint{1.528330in}{1.030286in}}%
\pgfpathlineto{\pgfqpoint{1.546173in}{1.041071in}}%
\pgfpathlineto{\pgfqpoint{1.553891in}{1.047486in}}%
\pgfpathlineto{\pgfqpoint{1.561593in}{1.054211in}}%
\pgfpathlineto{\pgfqpoint{1.543774in}{1.043515in}}%
\pgfpathlineto{\pgfqpoint{1.526034in}{1.033060in}}%
\pgfpathlineto{\pgfqpoint{1.518308in}{1.026235in}}%
\pgfpathlineto{\pgfqpoint{1.510566in}{1.019738in}}%
\pgfpathclose%
\pgfusepath{fill}%
\end{pgfscope}%
\begin{pgfscope}%
\pgfpathrectangle{\pgfqpoint{0.407592in}{0.222904in}}{\pgfqpoint{1.411163in}{1.411163in}}%
\pgfusepath{clip}%
\pgfsetbuttcap%
\pgfsetroundjoin%
\definecolor{currentfill}{rgb}{0.909804,0.698039,0.423529}%
\pgfsetfillcolor{currentfill}%
\pgfsetfillopacity{0.120000}%
\pgfsetlinewidth{0.000000pt}%
\definecolor{currentstroke}{rgb}{0.000000,0.000000,0.000000}%
\pgfsetstrokecolor{currentstroke}%
\pgfsetdash{}{0pt}%
\pgfpathmoveto{\pgfqpoint{0.819116in}{0.920549in}}%
\pgfpathlineto{\pgfqpoint{0.835756in}{0.913201in}}%
\pgfpathlineto{\pgfqpoint{0.852421in}{0.906562in}}%
\pgfpathlineto{\pgfqpoint{0.861367in}{0.902424in}}%
\pgfpathlineto{\pgfqpoint{0.870271in}{0.898525in}}%
\pgfpathlineto{\pgfqpoint{0.853655in}{0.905257in}}%
\pgfpathlineto{\pgfqpoint{0.837064in}{0.912725in}}%
\pgfpathlineto{\pgfqpoint{0.828111in}{0.916517in}}%
\pgfpathlineto{\pgfqpoint{0.819116in}{0.920549in}}%
\pgfpathclose%
\pgfusepath{fill}%
\end{pgfscope}%
\begin{pgfscope}%
\pgfpathrectangle{\pgfqpoint{0.407592in}{0.222904in}}{\pgfqpoint{1.411163in}{1.411163in}}%
\pgfusepath{clip}%
\pgfsetbuttcap%
\pgfsetroundjoin%
\definecolor{currentfill}{rgb}{0.909804,0.698039,0.423529}%
\pgfsetfillcolor{currentfill}%
\pgfsetfillopacity{0.120000}%
\pgfsetlinewidth{0.000000pt}%
\definecolor{currentstroke}{rgb}{0.000000,0.000000,0.000000}%
\pgfsetstrokecolor{currentstroke}%
\pgfsetdash{}{0pt}%
\pgfpathmoveto{\pgfqpoint{0.734469in}{0.964673in}}%
\pgfpathlineto{\pgfqpoint{0.751071in}{0.954896in}}%
\pgfpathlineto{\pgfqpoint{0.767693in}{0.945724in}}%
\pgfpathlineto{\pgfqpoint{0.776821in}{0.941400in}}%
\pgfpathlineto{\pgfqpoint{0.785907in}{0.937294in}}%
\pgfpathlineto{\pgfqpoint{0.769334in}{0.946650in}}%
\pgfpathlineto{\pgfqpoint{0.752780in}{0.956630in}}%
\pgfpathlineto{\pgfqpoint{0.743646in}{0.960540in}}%
\pgfpathlineto{\pgfqpoint{0.734469in}{0.964673in}}%
\pgfpathclose%
\pgfusepath{fill}%
\end{pgfscope}%
\begin{pgfscope}%
\pgfpathrectangle{\pgfqpoint{0.407592in}{0.222904in}}{\pgfqpoint{1.411163in}{1.411163in}}%
\pgfusepath{clip}%
\pgfsetbuttcap%
\pgfsetroundjoin%
\definecolor{currentfill}{rgb}{0.909804,0.698039,0.423529}%
\pgfsetfillcolor{currentfill}%
\pgfsetfillopacity{0.120000}%
\pgfsetlinewidth{0.000000pt}%
\definecolor{currentstroke}{rgb}{0.000000,0.000000,0.000000}%
\pgfsetstrokecolor{currentstroke}%
\pgfsetdash{}{0pt}%
\pgfpathmoveto{\pgfqpoint{1.089513in}{0.857967in}}%
\pgfpathlineto{\pgfqpoint{1.106427in}{0.860160in}}%
\pgfpathlineto{\pgfqpoint{1.123397in}{0.863066in}}%
\pgfpathlineto{\pgfqpoint{1.131767in}{0.861907in}}%
\pgfpathlineto{\pgfqpoint{1.140105in}{0.861114in}}%
\pgfpathlineto{\pgfqpoint{1.123174in}{0.857973in}}%
\pgfpathlineto{\pgfqpoint{1.106299in}{0.855560in}}%
\pgfpathlineto{\pgfqpoint{1.097923in}{0.856578in}}%
\pgfpathlineto{\pgfqpoint{1.089513in}{0.857967in}}%
\pgfpathclose%
\pgfusepath{fill}%
\end{pgfscope}%
\begin{pgfscope}%
\pgfpathrectangle{\pgfqpoint{0.407592in}{0.222904in}}{\pgfqpoint{1.411163in}{1.411163in}}%
\pgfusepath{clip}%
\pgfsetbuttcap%
\pgfsetroundjoin%
\definecolor{currentfill}{rgb}{0.909804,0.698039,0.423529}%
\pgfsetfillcolor{currentfill}%
\pgfsetfillopacity{0.120000}%
\pgfsetlinewidth{0.000000pt}%
\definecolor{currentstroke}{rgb}{0.000000,0.000000,0.000000}%
\pgfsetstrokecolor{currentstroke}%
\pgfsetdash{}{0pt}%
\pgfpathmoveto{\pgfqpoint{1.342428in}{0.922588in}}%
\pgfpathlineto{\pgfqpoint{1.359833in}{0.931307in}}%
\pgfpathlineto{\pgfqpoint{1.377313in}{0.940390in}}%
\pgfpathlineto{\pgfqpoint{1.385254in}{0.944042in}}%
\pgfpathlineto{\pgfqpoint{1.393174in}{0.948074in}}%
\pgfpathlineto{\pgfqpoint{1.375721in}{0.938923in}}%
\pgfpathlineto{\pgfqpoint{1.358341in}{0.930130in}}%
\pgfpathlineto{\pgfqpoint{1.350396in}{0.926158in}}%
\pgfpathlineto{\pgfqpoint{1.342428in}{0.922588in}}%
\pgfpathclose%
\pgfusepath{fill}%
\end{pgfscope}%
\begin{pgfscope}%
\pgfpathrectangle{\pgfqpoint{0.407592in}{0.222904in}}{\pgfqpoint{1.411163in}{1.411163in}}%
\pgfusepath{clip}%
\pgfsetbuttcap%
\pgfsetroundjoin%
\definecolor{currentfill}{rgb}{0.909804,0.698039,0.423529}%
\pgfsetfillcolor{currentfill}%
\pgfsetfillopacity{0.120000}%
\pgfsetlinewidth{0.000000pt}%
\definecolor{currentstroke}{rgb}{0.000000,0.000000,0.000000}%
\pgfsetstrokecolor{currentstroke}%
\pgfsetdash{}{0pt}%
\pgfpathmoveto{\pgfqpoint{0.903587in}{0.887328in}}%
\pgfpathlineto{\pgfqpoint{0.920291in}{0.882884in}}%
\pgfpathlineto{\pgfqpoint{0.937030in}{0.879215in}}%
\pgfpathlineto{\pgfqpoint{0.945797in}{0.875535in}}%
\pgfpathlineto{\pgfqpoint{0.954524in}{0.872122in}}%
\pgfpathlineto{\pgfqpoint{0.937834in}{0.875758in}}%
\pgfpathlineto{\pgfqpoint{0.921178in}{0.880203in}}%
\pgfpathlineto{\pgfqpoint{0.912402in}{0.883632in}}%
\pgfpathlineto{\pgfqpoint{0.903587in}{0.887328in}}%
\pgfpathclose%
\pgfusepath{fill}%
\end{pgfscope}%
\begin{pgfscope}%
\pgfpathrectangle{\pgfqpoint{0.407592in}{0.222904in}}{\pgfqpoint{1.411163in}{1.411163in}}%
\pgfusepath{clip}%
\pgfsetbuttcap%
\pgfsetroundjoin%
\definecolor{currentfill}{rgb}{0.909804,0.698039,0.423529}%
\pgfsetfillcolor{currentfill}%
\pgfsetfillopacity{0.120000}%
\pgfsetlinewidth{0.000000pt}%
\definecolor{currentstroke}{rgb}{0.000000,0.000000,0.000000}%
\pgfsetstrokecolor{currentstroke}%
\pgfsetdash{}{0pt}%
\pgfpathmoveto{\pgfqpoint{0.649557in}{1.017343in}}%
\pgfpathlineto{\pgfqpoint{0.666140in}{1.005626in}}%
\pgfpathlineto{\pgfqpoint{0.682738in}{0.994408in}}%
\pgfpathlineto{\pgfqpoint{0.692050in}{0.990077in}}%
\pgfpathlineto{\pgfqpoint{0.701319in}{0.985951in}}%
\pgfpathlineto{\pgfqpoint{0.684768in}{0.997403in}}%
\pgfpathlineto{\pgfqpoint{0.668232in}{1.009363in}}%
\pgfpathlineto{\pgfqpoint{0.658917in}{1.013249in}}%
\pgfpathlineto{\pgfqpoint{0.649557in}{1.017343in}}%
\pgfpathclose%
\pgfusepath{fill}%
\end{pgfscope}%
\begin{pgfscope}%
\pgfpathrectangle{\pgfqpoint{0.407592in}{0.222904in}}{\pgfqpoint{1.411163in}{1.411163in}}%
\pgfusepath{clip}%
\pgfsetbuttcap%
\pgfsetroundjoin%
\definecolor{currentfill}{rgb}{0.909804,0.698039,0.423529}%
\pgfsetfillcolor{currentfill}%
\pgfsetfillopacity{0.120000}%
\pgfsetlinewidth{0.000000pt}%
\definecolor{currentstroke}{rgb}{0.000000,0.000000,0.000000}%
\pgfsetstrokecolor{currentstroke}%
\pgfsetdash{}{0pt}%
\pgfpathmoveto{\pgfqpoint{1.679277in}{1.144104in}}%
\pgfpathlineto{\pgfqpoint{1.697423in}{1.155439in}}%
\pgfpathlineto{\pgfqpoint{1.715651in}{1.166960in}}%
\pgfpathlineto{\pgfqpoint{1.723180in}{1.175201in}}%
\pgfpathlineto{\pgfqpoint{1.730696in}{1.183680in}}%
\pgfpathlineto{\pgfqpoint{1.712493in}{1.172323in}}%
\pgfpathlineto{\pgfqpoint{1.694371in}{1.161157in}}%
\pgfpathlineto{\pgfqpoint{1.686830in}{1.152507in}}%
\pgfpathlineto{\pgfqpoint{1.679277in}{1.144104in}}%
\pgfpathclose%
\pgfusepath{fill}%
\end{pgfscope}%
\begin{pgfscope}%
\pgfpathrectangle{\pgfqpoint{0.407592in}{0.222904in}}{\pgfqpoint{1.411163in}{1.411163in}}%
\pgfusepath{clip}%
\pgfsetbuttcap%
\pgfsetroundjoin%
\definecolor{currentfill}{rgb}{0.909804,0.698039,0.423529}%
\pgfsetfillcolor{currentfill}%
\pgfsetfillopacity{0.120000}%
\pgfsetlinewidth{0.000000pt}%
\definecolor{currentstroke}{rgb}{0.000000,0.000000,0.000000}%
\pgfsetstrokecolor{currentstroke}%
\pgfsetdash{}{0pt}%
\pgfpathmoveto{\pgfqpoint{1.038832in}{0.860072in}}%
\pgfpathlineto{\pgfqpoint{1.055685in}{0.860553in}}%
\pgfpathlineto{\pgfqpoint{1.072588in}{0.861791in}}%
\pgfpathlineto{\pgfqpoint{1.081068in}{0.859710in}}%
\pgfpathlineto{\pgfqpoint{1.089513in}{0.857967in}}%
\pgfpathlineto{\pgfqpoint{1.072652in}{0.856521in}}%
\pgfpathlineto{\pgfqpoint{1.055841in}{0.855855in}}%
\pgfpathlineto{\pgfqpoint{1.047354in}{0.857795in}}%
\pgfpathlineto{\pgfqpoint{1.038832in}{0.860072in}}%
\pgfpathclose%
\pgfusepath{fill}%
\end{pgfscope}%
\begin{pgfscope}%
\pgfpathrectangle{\pgfqpoint{0.407592in}{0.222904in}}{\pgfqpoint{1.411163in}{1.411163in}}%
\pgfusepath{clip}%
\pgfsetbuttcap%
\pgfsetroundjoin%
\definecolor{currentfill}{rgb}{0.909804,0.698039,0.423529}%
\pgfsetfillcolor{currentfill}%
\pgfsetfillopacity{0.120000}%
\pgfsetlinewidth{0.000000pt}%
\definecolor{currentstroke}{rgb}{0.000000,0.000000,0.000000}%
\pgfsetstrokecolor{currentstroke}%
\pgfsetdash{}{0pt}%
\pgfpathmoveto{\pgfqpoint{1.561593in}{1.054211in}}%
\pgfpathlineto{\pgfqpoint{1.579491in}{1.065133in}}%
\pgfpathlineto{\pgfqpoint{1.597469in}{1.076270in}}%
\pgfpathlineto{\pgfqpoint{1.605131in}{1.083193in}}%
\pgfpathlineto{\pgfqpoint{1.612779in}{1.090404in}}%
\pgfpathlineto{\pgfqpoint{1.594825in}{1.079381in}}%
\pgfpathlineto{\pgfqpoint{1.576951in}{1.068576in}}%
\pgfpathlineto{\pgfqpoint{1.569279in}{1.061242in}}%
\pgfpathlineto{\pgfqpoint{1.561593in}{1.054211in}}%
\pgfpathclose%
\pgfusepath{fill}%
\end{pgfscope}%
\begin{pgfscope}%
\pgfpathrectangle{\pgfqpoint{0.407592in}{0.222904in}}{\pgfqpoint{1.411163in}{1.411163in}}%
\pgfusepath{clip}%
\pgfsetbuttcap%
\pgfsetroundjoin%
\definecolor{currentfill}{rgb}{0.909804,0.698039,0.423529}%
\pgfsetfillcolor{currentfill}%
\pgfsetfillopacity{0.120000}%
\pgfsetlinewidth{0.000000pt}%
\definecolor{currentstroke}{rgb}{0.000000,0.000000,0.000000}%
\pgfsetstrokecolor{currentstroke}%
\pgfsetdash{}{0pt}%
\pgfpathmoveto{\pgfqpoint{1.393174in}{0.948074in}}%
\pgfpathlineto{\pgfqpoint{1.410703in}{0.957560in}}%
\pgfpathlineto{\pgfqpoint{1.428308in}{0.967357in}}%
\pgfpathlineto{\pgfqpoint{1.436182in}{0.971803in}}%
\pgfpathlineto{\pgfqpoint{1.444036in}{0.976610in}}%
\pgfpathlineto{\pgfqpoint{1.426457in}{0.966792in}}%
\pgfpathlineto{\pgfqpoint{1.408953in}{0.957283in}}%
\pgfpathlineto{\pgfqpoint{1.401073in}{0.952488in}}%
\pgfpathlineto{\pgfqpoint{1.393174in}{0.948074in}}%
\pgfpathclose%
\pgfusepath{fill}%
\end{pgfscope}%
\begin{pgfscope}%
\pgfpathrectangle{\pgfqpoint{0.407592in}{0.222904in}}{\pgfqpoint{1.411163in}{1.411163in}}%
\pgfusepath{clip}%
\pgfsetbuttcap%
\pgfsetroundjoin%
\definecolor{currentfill}{rgb}{0.909804,0.698039,0.423529}%
\pgfsetfillcolor{currentfill}%
\pgfsetfillopacity{0.120000}%
\pgfsetlinewidth{0.000000pt}%
\definecolor{currentstroke}{rgb}{0.000000,0.000000,0.000000}%
\pgfsetstrokecolor{currentstroke}%
\pgfsetdash{}{0pt}%
\pgfpathmoveto{\pgfqpoint{1.730696in}{1.183680in}}%
\pgfpathlineto{\pgfqpoint{1.748980in}{1.195221in}}%
\pgfpathlineto{\pgfqpoint{1.756471in}{1.203847in}}%
\pgfpathlineto{\pgfqpoint{1.763949in}{1.212696in}}%
\pgfpathlineto{\pgfqpoint{1.745689in}{1.201330in}}%
\pgfpathlineto{\pgfqpoint{1.738199in}{1.192391in}}%
\pgfusepath{fill}%
\end{pgfscope}%
\begin{pgfscope}%
\pgfpathrectangle{\pgfqpoint{0.407592in}{0.222904in}}{\pgfqpoint{1.411163in}{1.411163in}}%
\pgfusepath{clip}%
\pgfsetbuttcap%
\pgfsetroundjoin%
\definecolor{currentfill}{rgb}{0.909804,0.698039,0.423529}%
\pgfsetfillcolor{currentfill}%
\pgfsetfillopacity{0.120000}%
\pgfsetlinewidth{0.000000pt}%
\definecolor{currentstroke}{rgb}{0.000000,0.000000,0.000000}%
\pgfsetstrokecolor{currentstroke}%
\pgfsetdash{}{0pt}%
\pgfpathmoveto{\pgfqpoint{0.988017in}{0.867267in}}%
\pgfpathlineto{\pgfqpoint{1.004824in}{0.866031in}}%
\pgfpathlineto{\pgfqpoint{1.021676in}{0.865571in}}%
\pgfpathlineto{\pgfqpoint{1.030272in}{0.862670in}}%
\pgfpathlineto{\pgfqpoint{1.038832in}{0.860072in}}%
\pgfpathlineto{\pgfqpoint{1.022025in}{0.860373in}}%
\pgfpathlineto{\pgfqpoint{1.005263in}{0.861478in}}%
\pgfpathlineto{\pgfqpoint{0.996659in}{0.864222in}}%
\pgfpathlineto{\pgfqpoint{0.988017in}{0.867267in}}%
\pgfpathclose%
\pgfusepath{fill}%
\end{pgfscope}%
\begin{pgfscope}%
\pgfpathrectangle{\pgfqpoint{0.407592in}{0.222904in}}{\pgfqpoint{1.411163in}{1.411163in}}%
\pgfusepath{clip}%
\pgfsetbuttcap%
\pgfsetroundjoin%
\definecolor{currentfill}{rgb}{0.909804,0.698039,0.423529}%
\pgfsetfillcolor{currentfill}%
\pgfsetfillopacity{0.120000}%
\pgfsetlinewidth{0.000000pt}%
\definecolor{currentstroke}{rgb}{0.000000,0.000000,0.000000}%
\pgfsetstrokecolor{currentstroke}%
\pgfsetdash{}{0pt}%
\pgfpathmoveto{\pgfqpoint{1.444036in}{0.976610in}}%
\pgfpathlineto{\pgfqpoint{1.461693in}{0.986715in}}%
\pgfpathlineto{\pgfqpoint{1.479428in}{0.997090in}}%
\pgfpathlineto{\pgfqpoint{1.487239in}{1.002245in}}%
\pgfpathlineto{\pgfqpoint{1.495032in}{1.007740in}}%
\pgfpathlineto{\pgfqpoint{1.477322in}{0.997386in}}%
\pgfpathlineto{\pgfqpoint{1.459689in}{0.987302in}}%
\pgfpathlineto{\pgfqpoint{1.451872in}{0.981776in}}%
\pgfpathlineto{\pgfqpoint{1.444036in}{0.976610in}}%
\pgfpathclose%
\pgfusepath{fill}%
\end{pgfscope}%
\begin{pgfscope}%
\pgfpathrectangle{\pgfqpoint{0.407592in}{0.222904in}}{\pgfqpoint{1.411163in}{1.411163in}}%
\pgfusepath{clip}%
\pgfsetbuttcap%
\pgfsetroundjoin%
\definecolor{currentfill}{rgb}{0.909804,0.698039,0.423529}%
\pgfsetfillcolor{currentfill}%
\pgfsetfillopacity{0.120000}%
\pgfsetlinewidth{0.000000pt}%
\definecolor{currentstroke}{rgb}{0.000000,0.000000,0.000000}%
\pgfsetstrokecolor{currentstroke}%
\pgfsetdash{}{0pt}%
\pgfpathmoveto{\pgfqpoint{0.767693in}{0.945724in}}%
\pgfpathlineto{\pgfqpoint{0.784335in}{0.937180in}}%
\pgfpathlineto{\pgfqpoint{0.801000in}{0.929284in}}%
\pgfpathlineto{\pgfqpoint{0.810079in}{0.924808in}}%
\pgfpathlineto{\pgfqpoint{0.819116in}{0.920549in}}%
\pgfpathlineto{\pgfqpoint{0.802501in}{0.928586in}}%
\pgfpathlineto{\pgfqpoint{0.785907in}{0.937294in}}%
\pgfpathlineto{\pgfqpoint{0.776821in}{0.941400in}}%
\pgfpathlineto{\pgfqpoint{0.767693in}{0.945724in}}%
\pgfpathclose%
\pgfusepath{fill}%
\end{pgfscope}%
\begin{pgfscope}%
\pgfpathrectangle{\pgfqpoint{0.407592in}{0.222904in}}{\pgfqpoint{1.411163in}{1.411163in}}%
\pgfusepath{clip}%
\pgfsetbuttcap%
\pgfsetroundjoin%
\definecolor{currentfill}{rgb}{0.909804,0.698039,0.423529}%
\pgfsetfillcolor{currentfill}%
\pgfsetfillopacity{0.120000}%
\pgfsetlinewidth{0.000000pt}%
\definecolor{currentstroke}{rgb}{0.000000,0.000000,0.000000}%
\pgfsetstrokecolor{currentstroke}%
\pgfsetdash{}{0pt}%
\pgfpathmoveto{\pgfqpoint{0.852421in}{0.906562in}}%
\pgfpathlineto{\pgfqpoint{0.869113in}{0.900646in}}%
\pgfpathlineto{\pgfqpoint{0.885835in}{0.895462in}}%
\pgfpathlineto{\pgfqpoint{0.894731in}{0.891276in}}%
\pgfpathlineto{\pgfqpoint{0.903587in}{0.887328in}}%
\pgfpathlineto{\pgfqpoint{0.886914in}{0.892544in}}%
\pgfpathlineto{\pgfqpoint{0.870271in}{0.898525in}}%
\pgfpathlineto{\pgfqpoint{0.861367in}{0.902424in}}%
\pgfpathlineto{\pgfqpoint{0.852421in}{0.906562in}}%
\pgfpathclose%
\pgfusepath{fill}%
\end{pgfscope}%
\begin{pgfscope}%
\pgfpathrectangle{\pgfqpoint{0.407592in}{0.222904in}}{\pgfqpoint{1.411163in}{1.411163in}}%
\pgfusepath{clip}%
\pgfsetbuttcap%
\pgfsetroundjoin%
\definecolor{currentfill}{rgb}{0.909804,0.698039,0.423529}%
\pgfsetfillcolor{currentfill}%
\pgfsetfillopacity{0.120000}%
\pgfsetlinewidth{0.000000pt}%
\definecolor{currentstroke}{rgb}{0.000000,0.000000,0.000000}%
\pgfsetstrokecolor{currentstroke}%
\pgfsetdash{}{0pt}%
\pgfpathmoveto{\pgfqpoint{1.612779in}{1.090404in}}%
\pgfpathlineto{\pgfqpoint{1.630813in}{1.101634in}}%
\pgfpathlineto{\pgfqpoint{1.648927in}{1.113059in}}%
\pgfpathlineto{\pgfqpoint{1.656536in}{1.120425in}}%
\pgfpathlineto{\pgfqpoint{1.664130in}{1.128058in}}%
\pgfpathlineto{\pgfqpoint{1.646040in}{1.116764in}}%
\pgfpathlineto{\pgfqpoint{1.628030in}{1.105671in}}%
\pgfpathlineto{\pgfqpoint{1.620411in}{1.097898in}}%
\pgfpathlineto{\pgfqpoint{1.612779in}{1.090404in}}%
\pgfpathclose%
\pgfusepath{fill}%
\end{pgfscope}%
\begin{pgfscope}%
\pgfpathrectangle{\pgfqpoint{0.407592in}{0.222904in}}{\pgfqpoint{1.411163in}{1.411163in}}%
\pgfusepath{clip}%
\pgfsetbuttcap%
\pgfsetroundjoin%
\definecolor{currentfill}{rgb}{0.909804,0.698039,0.423529}%
\pgfsetfillcolor{currentfill}%
\pgfsetfillopacity{0.120000}%
\pgfsetlinewidth{0.000000pt}%
\definecolor{currentstroke}{rgb}{0.000000,0.000000,0.000000}%
\pgfsetstrokecolor{currentstroke}%
\pgfsetdash{}{0pt}%
\pgfpathmoveto{\pgfqpoint{0.682738in}{0.994408in}}%
\pgfpathlineto{\pgfqpoint{0.699353in}{0.983713in}}%
\pgfpathlineto{\pgfqpoint{0.715985in}{0.973563in}}%
\pgfpathlineto{\pgfqpoint{0.725249in}{0.969017in}}%
\pgfpathlineto{\pgfqpoint{0.734469in}{0.964673in}}%
\pgfpathlineto{\pgfqpoint{0.717886in}{0.975033in}}%
\pgfpathlineto{\pgfqpoint{0.701319in}{0.985951in}}%
\pgfpathlineto{\pgfqpoint{0.692050in}{0.990077in}}%
\pgfpathlineto{\pgfqpoint{0.682738in}{0.994408in}}%
\pgfpathclose%
\pgfusepath{fill}%
\end{pgfscope}%
\begin{pgfscope}%
\pgfpathrectangle{\pgfqpoint{0.407592in}{0.222904in}}{\pgfqpoint{1.411163in}{1.411163in}}%
\pgfusepath{clip}%
\pgfsetbuttcap%
\pgfsetroundjoin%
\definecolor{currentfill}{rgb}{0.909804,0.698039,0.423529}%
\pgfsetfillcolor{currentfill}%
\pgfsetfillopacity{0.120000}%
\pgfsetlinewidth{0.000000pt}%
\definecolor{currentstroke}{rgb}{0.000000,0.000000,0.000000}%
\pgfsetstrokecolor{currentstroke}%
\pgfsetdash{}{0pt}%
\pgfpathmoveto{\pgfqpoint{1.495032in}{1.007740in}}%
\pgfpathlineto{\pgfqpoint{1.512820in}{1.018345in}}%
\pgfpathlineto{\pgfqpoint{1.530687in}{1.029188in}}%
\pgfpathlineto{\pgfqpoint{1.538438in}{1.034971in}}%
\pgfpathlineto{\pgfqpoint{1.546173in}{1.041071in}}%
\pgfpathlineto{\pgfqpoint{1.528330in}{1.030286in}}%
\pgfpathlineto{\pgfqpoint{1.510566in}{1.019738in}}%
\pgfpathlineto{\pgfqpoint{1.502807in}{1.013572in}}%
\pgfpathlineto{\pgfqpoint{1.495032in}{1.007740in}}%
\pgfpathclose%
\pgfusepath{fill}%
\end{pgfscope}%
\begin{pgfscope}%
\pgfpathrectangle{\pgfqpoint{0.407592in}{0.222904in}}{\pgfqpoint{1.411163in}{1.411163in}}%
\pgfusepath{clip}%
\pgfsetbuttcap%
\pgfsetroundjoin%
\definecolor{currentfill}{rgb}{0.909804,0.698039,0.423529}%
\pgfsetfillcolor{currentfill}%
\pgfsetfillopacity{0.120000}%
\pgfsetlinewidth{0.000000pt}%
\definecolor{currentstroke}{rgb}{0.000000,0.000000,0.000000}%
\pgfsetstrokecolor{currentstroke}%
\pgfsetdash{}{0pt}%
\pgfpathmoveto{\pgfqpoint{1.224872in}{0.880702in}}%
\pgfpathlineto{\pgfqpoint{1.242081in}{0.887146in}}%
\pgfpathlineto{\pgfqpoint{1.259358in}{0.894093in}}%
\pgfpathlineto{\pgfqpoint{1.267503in}{0.895145in}}%
\pgfpathlineto{\pgfqpoint{1.275621in}{0.896583in}}%
\pgfpathlineto{\pgfqpoint{1.258375in}{0.889431in}}%
\pgfpathlineto{\pgfqpoint{1.241197in}{0.882778in}}%
\pgfpathlineto{\pgfqpoint{1.233048in}{0.881540in}}%
\pgfpathlineto{\pgfqpoint{1.224872in}{0.880702in}}%
\pgfpathclose%
\pgfusepath{fill}%
\end{pgfscope}%
\begin{pgfscope}%
\pgfpathrectangle{\pgfqpoint{0.407592in}{0.222904in}}{\pgfqpoint{1.411163in}{1.411163in}}%
\pgfusepath{clip}%
\pgfsetbuttcap%
\pgfsetroundjoin%
\definecolor{currentfill}{rgb}{0.909804,0.698039,0.423529}%
\pgfsetfillcolor{currentfill}%
\pgfsetfillopacity{0.120000}%
\pgfsetlinewidth{0.000000pt}%
\definecolor{currentstroke}{rgb}{0.000000,0.000000,0.000000}%
\pgfsetstrokecolor{currentstroke}%
\pgfsetdash{}{0pt}%
\pgfpathmoveto{\pgfqpoint{1.174144in}{0.869424in}}%
\pgfpathlineto{\pgfqpoint{1.191256in}{0.874516in}}%
\pgfpathlineto{\pgfqpoint{1.208432in}{0.880184in}}%
\pgfpathlineto{\pgfqpoint{1.216667in}{0.880254in}}%
\pgfpathlineto{\pgfqpoint{1.224872in}{0.880702in}}%
\pgfpathlineto{\pgfqpoint{1.207730in}{0.874796in}}%
\pgfpathlineto{\pgfqpoint{1.190652in}{0.869466in}}%
\pgfpathlineto{\pgfqpoint{1.182413in}{0.869250in}}%
\pgfpathlineto{\pgfqpoint{1.174144in}{0.869424in}}%
\pgfpathclose%
\pgfusepath{fill}%
\end{pgfscope}%
\begin{pgfscope}%
\pgfpathrectangle{\pgfqpoint{0.407592in}{0.222904in}}{\pgfqpoint{1.411163in}{1.411163in}}%
\pgfusepath{clip}%
\pgfsetbuttcap%
\pgfsetroundjoin%
\definecolor{currentfill}{rgb}{0.909804,0.698039,0.423529}%
\pgfsetfillcolor{currentfill}%
\pgfsetfillopacity{0.120000}%
\pgfsetlinewidth{0.000000pt}%
\definecolor{currentstroke}{rgb}{0.000000,0.000000,0.000000}%
\pgfsetstrokecolor{currentstroke}%
\pgfsetdash{}{0pt}%
\pgfpathmoveto{\pgfqpoint{0.937030in}{0.879215in}}%
\pgfpathlineto{\pgfqpoint{0.953804in}{0.876321in}}%
\pgfpathlineto{\pgfqpoint{0.970617in}{0.874197in}}%
\pgfpathlineto{\pgfqpoint{0.979337in}{0.870598in}}%
\pgfpathlineto{\pgfqpoint{0.988017in}{0.867267in}}%
\pgfpathlineto{\pgfqpoint{0.971251in}{0.869294in}}%
\pgfpathlineto{\pgfqpoint{0.954524in}{0.872122in}}%
\pgfpathlineto{\pgfqpoint{0.945797in}{0.875535in}}%
\pgfpathlineto{\pgfqpoint{0.937030in}{0.879215in}}%
\pgfpathclose%
\pgfusepath{fill}%
\end{pgfscope}%
\begin{pgfscope}%
\pgfpathrectangle{\pgfqpoint{0.407592in}{0.222904in}}{\pgfqpoint{1.411163in}{1.411163in}}%
\pgfusepath{clip}%
\pgfsetbuttcap%
\pgfsetroundjoin%
\definecolor{currentfill}{rgb}{0.909804,0.698039,0.423529}%
\pgfsetfillcolor{currentfill}%
\pgfsetfillopacity{0.120000}%
\pgfsetlinewidth{0.000000pt}%
\definecolor{currentstroke}{rgb}{0.000000,0.000000,0.000000}%
\pgfsetstrokecolor{currentstroke}%
\pgfsetdash{}{0pt}%
\pgfpathmoveto{\pgfqpoint{1.275621in}{0.896583in}}%
\pgfpathlineto{\pgfqpoint{1.292937in}{0.904199in}}%
\pgfpathlineto{\pgfqpoint{1.310325in}{0.912246in}}%
\pgfpathlineto{\pgfqpoint{1.318388in}{0.914250in}}%
\pgfpathlineto{\pgfqpoint{1.326425in}{0.916638in}}%
\pgfpathlineto{\pgfqpoint{1.309066in}{0.908431in}}%
\pgfpathlineto{\pgfqpoint{1.291778in}{0.900649in}}%
\pgfpathlineto{\pgfqpoint{1.283712in}{0.898415in}}%
\pgfpathlineto{\pgfqpoint{1.275621in}{0.896583in}}%
\pgfpathclose%
\pgfusepath{fill}%
\end{pgfscope}%
\begin{pgfscope}%
\pgfpathrectangle{\pgfqpoint{0.407592in}{0.222904in}}{\pgfqpoint{1.411163in}{1.411163in}}%
\pgfusepath{clip}%
\pgfsetbuttcap%
\pgfsetroundjoin%
\definecolor{currentfill}{rgb}{0.909804,0.698039,0.423529}%
\pgfsetfillcolor{currentfill}%
\pgfsetfillopacity{0.120000}%
\pgfsetlinewidth{0.000000pt}%
\definecolor{currentstroke}{rgb}{0.000000,0.000000,0.000000}%
\pgfsetstrokecolor{currentstroke}%
\pgfsetdash{}{0pt}%
\pgfpathmoveto{\pgfqpoint{1.123397in}{0.863066in}}%
\pgfpathlineto{\pgfqpoint{1.140424in}{0.866651in}}%
\pgfpathlineto{\pgfqpoint{1.157510in}{0.870880in}}%
\pgfpathlineto{\pgfqpoint{1.165843in}{0.869971in}}%
\pgfpathlineto{\pgfqpoint{1.174144in}{0.869424in}}%
\pgfpathlineto{\pgfqpoint{1.157094in}{0.864943in}}%
\pgfpathlineto{\pgfqpoint{1.140105in}{0.861114in}}%
\pgfpathlineto{\pgfqpoint{1.131767in}{0.861907in}}%
\pgfpathlineto{\pgfqpoint{1.123397in}{0.863066in}}%
\pgfpathclose%
\pgfusepath{fill}%
\end{pgfscope}%
\begin{pgfscope}%
\pgfpathrectangle{\pgfqpoint{0.407592in}{0.222904in}}{\pgfqpoint{1.411163in}{1.411163in}}%
\pgfusepath{clip}%
\pgfsetbuttcap%
\pgfsetroundjoin%
\definecolor{currentfill}{rgb}{0.909804,0.698039,0.423529}%
\pgfsetfillcolor{currentfill}%
\pgfsetfillopacity{0.120000}%
\pgfsetlinewidth{0.000000pt}%
\definecolor{currentstroke}{rgb}{0.000000,0.000000,0.000000}%
\pgfsetstrokecolor{currentstroke}%
\pgfsetdash{}{0pt}%
\pgfpathmoveto{\pgfqpoint{1.664130in}{1.128058in}}%
\pgfpathlineto{\pgfqpoint{1.682301in}{1.139542in}}%
\pgfpathlineto{\pgfqpoint{1.700554in}{1.151209in}}%
\pgfpathlineto{\pgfqpoint{1.708110in}{1.158961in}}%
\pgfpathlineto{\pgfqpoint{1.715651in}{1.166960in}}%
\pgfpathlineto{\pgfqpoint{1.697423in}{1.155439in}}%
\pgfpathlineto{\pgfqpoint{1.679277in}{1.144104in}}%
\pgfpathlineto{\pgfqpoint{1.671710in}{1.135952in}}%
\pgfpathlineto{\pgfqpoint{1.664130in}{1.128058in}}%
\pgfpathclose%
\pgfusepath{fill}%
\end{pgfscope}%
\begin{pgfscope}%
\pgfpathrectangle{\pgfqpoint{0.407592in}{0.222904in}}{\pgfqpoint{1.411163in}{1.411163in}}%
\pgfusepath{clip}%
\pgfsetbuttcap%
\pgfsetroundjoin%
\definecolor{currentfill}{rgb}{0.909804,0.698039,0.423529}%
\pgfsetfillcolor{currentfill}%
\pgfsetfillopacity{0.120000}%
\pgfsetlinewidth{0.000000pt}%
\definecolor{currentstroke}{rgb}{0.000000,0.000000,0.000000}%
\pgfsetstrokecolor{currentstroke}%
\pgfsetdash{}{0pt}%
\pgfpathmoveto{\pgfqpoint{1.326425in}{0.916638in}}%
\pgfpathlineto{\pgfqpoint{1.343857in}{0.925242in}}%
\pgfpathlineto{\pgfqpoint{1.361364in}{0.934214in}}%
\pgfpathlineto{\pgfqpoint{1.369350in}{0.937115in}}%
\pgfpathlineto{\pgfqpoint{1.377313in}{0.940390in}}%
\pgfpathlineto{\pgfqpoint{1.359833in}{0.931307in}}%
\pgfpathlineto{\pgfqpoint{1.342428in}{0.922588in}}%
\pgfpathlineto{\pgfqpoint{1.334438in}{0.919416in}}%
\pgfpathlineto{\pgfqpoint{1.326425in}{0.916638in}}%
\pgfpathclose%
\pgfusepath{fill}%
\end{pgfscope}%
\begin{pgfscope}%
\pgfpathrectangle{\pgfqpoint{0.407592in}{0.222904in}}{\pgfqpoint{1.411163in}{1.411163in}}%
\pgfusepath{clip}%
\pgfsetbuttcap%
\pgfsetroundjoin%
\definecolor{currentfill}{rgb}{0.909804,0.698039,0.423529}%
\pgfsetfillcolor{currentfill}%
\pgfsetfillopacity{0.120000}%
\pgfsetlinewidth{0.000000pt}%
\definecolor{currentstroke}{rgb}{0.000000,0.000000,0.000000}%
\pgfsetstrokecolor{currentstroke}%
\pgfsetdash{}{0pt}%
\pgfpathmoveto{\pgfqpoint{1.715651in}{1.166960in}}%
\pgfpathlineto{\pgfqpoint{1.733962in}{1.178660in}}%
\pgfpathlineto{\pgfqpoint{1.741478in}{1.186824in}}%
\pgfpathlineto{\pgfqpoint{1.748980in}{1.195221in}}%
\pgfpathlineto{\pgfqpoint{1.730696in}{1.183680in}}%
\pgfpathlineto{\pgfqpoint{1.723180in}{1.175201in}}%
\pgfusepath{fill}%
\end{pgfscope}%
\begin{pgfscope}%
\pgfpathrectangle{\pgfqpoint{0.407592in}{0.222904in}}{\pgfqpoint{1.411163in}{1.411163in}}%
\pgfusepath{clip}%
\pgfsetbuttcap%
\pgfsetroundjoin%
\definecolor{currentfill}{rgb}{0.909804,0.698039,0.423529}%
\pgfsetfillcolor{currentfill}%
\pgfsetfillopacity{0.120000}%
\pgfsetlinewidth{0.000000pt}%
\definecolor{currentstroke}{rgb}{0.000000,0.000000,0.000000}%
\pgfsetstrokecolor{currentstroke}%
\pgfsetdash{}{0pt}%
\pgfpathmoveto{\pgfqpoint{1.072588in}{0.861791in}}%
\pgfpathlineto{\pgfqpoint{1.089544in}{0.863757in}}%
\pgfpathlineto{\pgfqpoint{1.106554in}{0.866422in}}%
\pgfpathlineto{\pgfqpoint{1.114992in}{0.864577in}}%
\pgfpathlineto{\pgfqpoint{1.123397in}{0.863066in}}%
\pgfpathlineto{\pgfqpoint{1.106427in}{0.860160in}}%
\pgfpathlineto{\pgfqpoint{1.089513in}{0.857967in}}%
\pgfpathlineto{\pgfqpoint{1.081068in}{0.859710in}}%
\pgfpathlineto{\pgfqpoint{1.072588in}{0.861791in}}%
\pgfpathclose%
\pgfusepath{fill}%
\end{pgfscope}%
\begin{pgfscope}%
\pgfpathrectangle{\pgfqpoint{0.407592in}{0.222904in}}{\pgfqpoint{1.411163in}{1.411163in}}%
\pgfusepath{clip}%
\pgfsetbuttcap%
\pgfsetroundjoin%
\definecolor{currentfill}{rgb}{0.909804,0.698039,0.423529}%
\pgfsetfillcolor{currentfill}%
\pgfsetfillopacity{0.120000}%
\pgfsetlinewidth{0.000000pt}%
\definecolor{currentstroke}{rgb}{0.000000,0.000000,0.000000}%
\pgfsetstrokecolor{currentstroke}%
\pgfsetdash{}{0pt}%
\pgfpathmoveto{\pgfqpoint{1.546173in}{1.041071in}}%
\pgfpathlineto{\pgfqpoint{1.564095in}{1.052081in}}%
\pgfpathlineto{\pgfqpoint{1.582098in}{1.063303in}}%
\pgfpathlineto{\pgfqpoint{1.589791in}{1.069639in}}%
\pgfpathlineto{\pgfqpoint{1.597469in}{1.076270in}}%
\pgfpathlineto{\pgfqpoint{1.579491in}{1.065133in}}%
\pgfpathlineto{\pgfqpoint{1.561593in}{1.054211in}}%
\pgfpathlineto{\pgfqpoint{1.553891in}{1.047486in}}%
\pgfpathlineto{\pgfqpoint{1.546173in}{1.041071in}}%
\pgfpathclose%
\pgfusepath{fill}%
\end{pgfscope}%
\begin{pgfscope}%
\pgfpathrectangle{\pgfqpoint{0.407592in}{0.222904in}}{\pgfqpoint{1.411163in}{1.411163in}}%
\pgfusepath{clip}%
\pgfsetbuttcap%
\pgfsetroundjoin%
\definecolor{currentfill}{rgb}{0.909804,0.698039,0.423529}%
\pgfsetfillcolor{currentfill}%
\pgfsetfillopacity{0.120000}%
\pgfsetlinewidth{0.000000pt}%
\definecolor{currentstroke}{rgb}{0.000000,0.000000,0.000000}%
\pgfsetstrokecolor{currentstroke}%
\pgfsetdash{}{0pt}%
\pgfpathmoveto{\pgfqpoint{0.801000in}{0.929284in}}%
\pgfpathlineto{\pgfqpoint{0.817689in}{0.922053in}}%
\pgfpathlineto{\pgfqpoint{0.834403in}{0.915504in}}%
\pgfpathlineto{\pgfqpoint{0.843433in}{0.910926in}}%
\pgfpathlineto{\pgfqpoint{0.852421in}{0.906562in}}%
\pgfpathlineto{\pgfqpoint{0.835756in}{0.913201in}}%
\pgfpathlineto{\pgfqpoint{0.819116in}{0.920549in}}%
\pgfpathlineto{\pgfqpoint{0.810079in}{0.924808in}}%
\pgfpathlineto{\pgfqpoint{0.801000in}{0.929284in}}%
\pgfpathclose%
\pgfusepath{fill}%
\end{pgfscope}%
\begin{pgfscope}%
\pgfpathrectangle{\pgfqpoint{0.407592in}{0.222904in}}{\pgfqpoint{1.411163in}{1.411163in}}%
\pgfusepath{clip}%
\pgfsetbuttcap%
\pgfsetroundjoin%
\definecolor{currentfill}{rgb}{0.909804,0.698039,0.423529}%
\pgfsetfillcolor{currentfill}%
\pgfsetfillopacity{0.120000}%
\pgfsetlinewidth{0.000000pt}%
\definecolor{currentstroke}{rgb}{0.000000,0.000000,0.000000}%
\pgfsetstrokecolor{currentstroke}%
\pgfsetdash{}{0pt}%
\pgfpathmoveto{\pgfqpoint{0.715985in}{0.973563in}}%
\pgfpathlineto{\pgfqpoint{0.732635in}{0.963981in}}%
\pgfpathlineto{\pgfqpoint{0.749306in}{0.954986in}}%
\pgfpathlineto{\pgfqpoint{0.758521in}{0.950256in}}%
\pgfpathlineto{\pgfqpoint{0.767693in}{0.945724in}}%
\pgfpathlineto{\pgfqpoint{0.751071in}{0.954896in}}%
\pgfpathlineto{\pgfqpoint{0.734469in}{0.964673in}}%
\pgfpathlineto{\pgfqpoint{0.725249in}{0.969017in}}%
\pgfpathlineto{\pgfqpoint{0.715985in}{0.973563in}}%
\pgfpathclose%
\pgfusepath{fill}%
\end{pgfscope}%
\begin{pgfscope}%
\pgfpathrectangle{\pgfqpoint{0.407592in}{0.222904in}}{\pgfqpoint{1.411163in}{1.411163in}}%
\pgfusepath{clip}%
\pgfsetbuttcap%
\pgfsetroundjoin%
\definecolor{currentfill}{rgb}{0.909804,0.698039,0.423529}%
\pgfsetfillcolor{currentfill}%
\pgfsetfillopacity{0.120000}%
\pgfsetlinewidth{0.000000pt}%
\definecolor{currentstroke}{rgb}{0.000000,0.000000,0.000000}%
\pgfsetstrokecolor{currentstroke}%
\pgfsetdash{}{0pt}%
\pgfpathmoveto{\pgfqpoint{1.377313in}{0.940390in}}%
\pgfpathlineto{\pgfqpoint{1.394867in}{0.949810in}}%
\pgfpathlineto{\pgfqpoint{1.412498in}{0.959546in}}%
\pgfpathlineto{\pgfqpoint{1.420413in}{0.963272in}}%
\pgfpathlineto{\pgfqpoint{1.428308in}{0.967357in}}%
\pgfpathlineto{\pgfqpoint{1.410703in}{0.957560in}}%
\pgfpathlineto{\pgfqpoint{1.393174in}{0.948074in}}%
\pgfpathlineto{\pgfqpoint{1.385254in}{0.944042in}}%
\pgfpathlineto{\pgfqpoint{1.377313in}{0.940390in}}%
\pgfpathclose%
\pgfusepath{fill}%
\end{pgfscope}%
\begin{pgfscope}%
\pgfpathrectangle{\pgfqpoint{0.407592in}{0.222904in}}{\pgfqpoint{1.411163in}{1.411163in}}%
\pgfusepath{clip}%
\pgfsetbuttcap%
\pgfsetroundjoin%
\definecolor{currentfill}{rgb}{0.909804,0.698039,0.423529}%
\pgfsetfillcolor{currentfill}%
\pgfsetfillopacity{0.120000}%
\pgfsetlinewidth{0.000000pt}%
\definecolor{currentstroke}{rgb}{0.000000,0.000000,0.000000}%
\pgfsetstrokecolor{currentstroke}%
\pgfsetdash{}{0pt}%
\pgfpathmoveto{\pgfqpoint{0.885835in}{0.895462in}}%
\pgfpathlineto{\pgfqpoint{0.902588in}{0.891019in}}%
\pgfpathlineto{\pgfqpoint{0.919375in}{0.887319in}}%
\pgfpathlineto{\pgfqpoint{0.928222in}{0.883148in}}%
\pgfpathlineto{\pgfqpoint{0.937030in}{0.879215in}}%
\pgfpathlineto{\pgfqpoint{0.920291in}{0.882884in}}%
\pgfpathlineto{\pgfqpoint{0.903587in}{0.887328in}}%
\pgfpathlineto{\pgfqpoint{0.894731in}{0.891276in}}%
\pgfpathlineto{\pgfqpoint{0.885835in}{0.895462in}}%
\pgfpathclose%
\pgfusepath{fill}%
\end{pgfscope}%
\begin{pgfscope}%
\pgfpathrectangle{\pgfqpoint{0.407592in}{0.222904in}}{\pgfqpoint{1.411163in}{1.411163in}}%
\pgfusepath{clip}%
\pgfsetbuttcap%
\pgfsetroundjoin%
\definecolor{currentfill}{rgb}{0.909804,0.698039,0.423529}%
\pgfsetfillcolor{currentfill}%
\pgfsetfillopacity{0.120000}%
\pgfsetlinewidth{0.000000pt}%
\definecolor{currentstroke}{rgb}{0.000000,0.000000,0.000000}%
\pgfsetstrokecolor{currentstroke}%
\pgfsetdash{}{0pt}%
\pgfpathmoveto{\pgfqpoint{1.021676in}{0.865571in}}%
\pgfpathlineto{\pgfqpoint{1.038574in}{0.865867in}}%
\pgfpathlineto{\pgfqpoint{1.055521in}{0.866897in}}%
\pgfpathlineto{\pgfqpoint{1.064073in}{0.864192in}}%
\pgfpathlineto{\pgfqpoint{1.072588in}{0.861791in}}%
\pgfpathlineto{\pgfqpoint{1.055685in}{0.860553in}}%
\pgfpathlineto{\pgfqpoint{1.038832in}{0.860072in}}%
\pgfpathlineto{\pgfqpoint{1.030272in}{0.862670in}}%
\pgfpathlineto{\pgfqpoint{1.021676in}{0.865571in}}%
\pgfpathclose%
\pgfusepath{fill}%
\end{pgfscope}%
\begin{pgfscope}%
\pgfpathrectangle{\pgfqpoint{0.407592in}{0.222904in}}{\pgfqpoint{1.411163in}{1.411163in}}%
\pgfusepath{clip}%
\pgfsetbuttcap%
\pgfsetroundjoin%
\definecolor{currentfill}{rgb}{0.909804,0.698039,0.423529}%
\pgfsetfillcolor{currentfill}%
\pgfsetfillopacity{0.120000}%
\pgfsetlinewidth{0.000000pt}%
\definecolor{currentstroke}{rgb}{0.000000,0.000000,0.000000}%
\pgfsetstrokecolor{currentstroke}%
\pgfsetdash{}{0pt}%
\pgfpathmoveto{\pgfqpoint{1.428308in}{0.967357in}}%
\pgfpathlineto{\pgfqpoint{1.445989in}{0.977444in}}%
\pgfpathlineto{\pgfqpoint{1.463749in}{0.987802in}}%
\pgfpathlineto{\pgfqpoint{1.471598in}{0.992276in}}%
\pgfpathlineto{\pgfqpoint{1.479428in}{0.997090in}}%
\pgfpathlineto{\pgfqpoint{1.461693in}{0.986715in}}%
\pgfpathlineto{\pgfqpoint{1.444036in}{0.976610in}}%
\pgfpathlineto{\pgfqpoint{1.436182in}{0.971803in}}%
\pgfpathlineto{\pgfqpoint{1.428308in}{0.967357in}}%
\pgfpathclose%
\pgfusepath{fill}%
\end{pgfscope}%
\begin{pgfscope}%
\pgfpathrectangle{\pgfqpoint{0.407592in}{0.222904in}}{\pgfqpoint{1.411163in}{1.411163in}}%
\pgfusepath{clip}%
\pgfsetbuttcap%
\pgfsetroundjoin%
\definecolor{currentfill}{rgb}{0.909804,0.698039,0.423529}%
\pgfsetfillcolor{currentfill}%
\pgfsetfillopacity{0.120000}%
\pgfsetlinewidth{0.000000pt}%
\definecolor{currentstroke}{rgb}{0.000000,0.000000,0.000000}%
\pgfsetstrokecolor{currentstroke}%
\pgfsetdash{}{0pt}%
\pgfpathmoveto{\pgfqpoint{1.597469in}{1.076270in}}%
\pgfpathlineto{\pgfqpoint{1.615527in}{1.087610in}}%
\pgfpathlineto{\pgfqpoint{1.633667in}{1.099142in}}%
\pgfpathlineto{\pgfqpoint{1.641305in}{1.105963in}}%
\pgfpathlineto{\pgfqpoint{1.648927in}{1.113059in}}%
\pgfpathlineto{\pgfqpoint{1.630813in}{1.101634in}}%
\pgfpathlineto{\pgfqpoint{1.612779in}{1.090404in}}%
\pgfpathlineto{\pgfqpoint{1.605131in}{1.083193in}}%
\pgfpathlineto{\pgfqpoint{1.597469in}{1.076270in}}%
\pgfpathclose%
\pgfusepath{fill}%
\end{pgfscope}%
\begin{pgfscope}%
\pgfpathrectangle{\pgfqpoint{0.407592in}{0.222904in}}{\pgfqpoint{1.411163in}{1.411163in}}%
\pgfusepath{clip}%
\pgfsetbuttcap%
\pgfsetroundjoin%
\definecolor{currentfill}{rgb}{0.909804,0.698039,0.423529}%
\pgfsetfillcolor{currentfill}%
\pgfsetfillopacity{0.120000}%
\pgfsetlinewidth{0.000000pt}%
\definecolor{currentstroke}{rgb}{0.000000,0.000000,0.000000}%
\pgfsetstrokecolor{currentstroke}%
\pgfsetdash{}{0pt}%
\pgfpathmoveto{\pgfqpoint{0.970617in}{0.874197in}}%
\pgfpathlineto{\pgfqpoint{0.987471in}{0.872834in}}%
\pgfpathlineto{\pgfqpoint{1.004369in}{0.872218in}}%
\pgfpathlineto{\pgfqpoint{1.013042in}{0.868758in}}%
\pgfpathlineto{\pgfqpoint{1.021676in}{0.865571in}}%
\pgfpathlineto{\pgfqpoint{1.004824in}{0.866031in}}%
\pgfpathlineto{\pgfqpoint{0.988017in}{0.867267in}}%
\pgfpathlineto{\pgfqpoint{0.979337in}{0.870598in}}%
\pgfpathlineto{\pgfqpoint{0.970617in}{0.874197in}}%
\pgfpathclose%
\pgfusepath{fill}%
\end{pgfscope}%
\begin{pgfscope}%
\pgfpathrectangle{\pgfqpoint{0.407592in}{0.222904in}}{\pgfqpoint{1.411163in}{1.411163in}}%
\pgfusepath{clip}%
\pgfsetbuttcap%
\pgfsetroundjoin%
\definecolor{currentfill}{rgb}{0.909804,0.698039,0.423529}%
\pgfsetfillcolor{currentfill}%
\pgfsetfillopacity{0.120000}%
\pgfsetlinewidth{0.000000pt}%
\definecolor{currentstroke}{rgb}{0.000000,0.000000,0.000000}%
\pgfsetstrokecolor{currentstroke}%
\pgfsetdash{}{0pt}%
\pgfpathmoveto{\pgfqpoint{1.479428in}{0.997090in}}%
\pgfpathlineto{\pgfqpoint{1.497241in}{1.007717in}}%
\pgfpathlineto{\pgfqpoint{1.515133in}{1.018582in}}%
\pgfpathlineto{\pgfqpoint{1.522919in}{1.023724in}}%
\pgfpathlineto{\pgfqpoint{1.530687in}{1.029188in}}%
\pgfpathlineto{\pgfqpoint{1.512820in}{1.018345in}}%
\pgfpathlineto{\pgfqpoint{1.495032in}{1.007740in}}%
\pgfpathlineto{\pgfqpoint{1.487239in}{1.002245in}}%
\pgfpathlineto{\pgfqpoint{1.479428in}{0.997090in}}%
\pgfpathclose%
\pgfusepath{fill}%
\end{pgfscope}%
\begin{pgfscope}%
\pgfpathrectangle{\pgfqpoint{0.407592in}{0.222904in}}{\pgfqpoint{1.411163in}{1.411163in}}%
\pgfusepath{clip}%
\pgfsetbuttcap%
\pgfsetroundjoin%
\definecolor{currentfill}{rgb}{0.909804,0.698039,0.423529}%
\pgfsetfillcolor{currentfill}%
\pgfsetfillopacity{0.120000}%
\pgfsetlinewidth{0.000000pt}%
\definecolor{currentstroke}{rgb}{0.000000,0.000000,0.000000}%
\pgfsetstrokecolor{currentstroke}%
\pgfsetdash{}{0pt}%
\pgfpathmoveto{\pgfqpoint{0.749306in}{0.954986in}}%
\pgfpathlineto{\pgfqpoint{0.765998in}{0.946600in}}%
\pgfpathlineto{\pgfqpoint{0.782712in}{0.938839in}}%
\pgfpathlineto{\pgfqpoint{0.791878in}{0.933964in}}%
\pgfpathlineto{\pgfqpoint{0.801000in}{0.929284in}}%
\pgfpathlineto{\pgfqpoint{0.784335in}{0.937180in}}%
\pgfpathlineto{\pgfqpoint{0.767693in}{0.945724in}}%
\pgfpathlineto{\pgfqpoint{0.758521in}{0.950256in}}%
\pgfpathlineto{\pgfqpoint{0.749306in}{0.954986in}}%
\pgfpathclose%
\pgfusepath{fill}%
\end{pgfscope}%
\begin{pgfscope}%
\pgfpathrectangle{\pgfqpoint{0.407592in}{0.222904in}}{\pgfqpoint{1.411163in}{1.411163in}}%
\pgfusepath{clip}%
\pgfsetbuttcap%
\pgfsetroundjoin%
\definecolor{currentfill}{rgb}{0.909804,0.698039,0.423529}%
\pgfsetfillcolor{currentfill}%
\pgfsetfillopacity{0.120000}%
\pgfsetlinewidth{0.000000pt}%
\definecolor{currentstroke}{rgb}{0.000000,0.000000,0.000000}%
\pgfsetstrokecolor{currentstroke}%
\pgfsetdash{}{0pt}%
\pgfpathmoveto{\pgfqpoint{1.648927in}{1.113059in}}%
\pgfpathlineto{\pgfqpoint{1.667124in}{1.124671in}}%
\pgfpathlineto{\pgfqpoint{1.685403in}{1.136461in}}%
\pgfpathlineto{\pgfqpoint{1.692986in}{1.143707in}}%
\pgfpathlineto{\pgfqpoint{1.700554in}{1.151209in}}%
\pgfpathlineto{\pgfqpoint{1.682301in}{1.139542in}}%
\pgfpathlineto{\pgfqpoint{1.664130in}{1.128058in}}%
\pgfpathlineto{\pgfqpoint{1.656536in}{1.120425in}}%
\pgfpathlineto{\pgfqpoint{1.648927in}{1.113059in}}%
\pgfpathclose%
\pgfusepath{fill}%
\end{pgfscope}%
\begin{pgfscope}%
\pgfpathrectangle{\pgfqpoint{0.407592in}{0.222904in}}{\pgfqpoint{1.411163in}{1.411163in}}%
\pgfusepath{clip}%
\pgfsetbuttcap%
\pgfsetroundjoin%
\definecolor{currentfill}{rgb}{0.909804,0.698039,0.423529}%
\pgfsetfillcolor{currentfill}%
\pgfsetfillopacity{0.120000}%
\pgfsetlinewidth{0.000000pt}%
\definecolor{currentstroke}{rgb}{0.000000,0.000000,0.000000}%
\pgfsetstrokecolor{currentstroke}%
\pgfsetdash{}{0pt}%
\pgfpathmoveto{\pgfqpoint{0.834403in}{0.915504in}}%
\pgfpathlineto{\pgfqpoint{0.851145in}{0.909648in}}%
\pgfpathlineto{\pgfqpoint{0.867917in}{0.904496in}}%
\pgfpathlineto{\pgfqpoint{0.876896in}{0.899873in}}%
\pgfpathlineto{\pgfqpoint{0.885835in}{0.895462in}}%
\pgfpathlineto{\pgfqpoint{0.869113in}{0.900646in}}%
\pgfpathlineto{\pgfqpoint{0.852421in}{0.906562in}}%
\pgfpathlineto{\pgfqpoint{0.843433in}{0.910926in}}%
\pgfpathlineto{\pgfqpoint{0.834403in}{0.915504in}}%
\pgfpathclose%
\pgfusepath{fill}%
\end{pgfscope}%
\begin{pgfscope}%
\pgfpathrectangle{\pgfqpoint{0.407592in}{0.222904in}}{\pgfqpoint{1.411163in}{1.411163in}}%
\pgfusepath{clip}%
\pgfsetbuttcap%
\pgfsetroundjoin%
\definecolor{currentfill}{rgb}{0.909804,0.698039,0.423529}%
\pgfsetfillcolor{currentfill}%
\pgfsetfillopacity{0.120000}%
\pgfsetlinewidth{0.000000pt}%
\definecolor{currentstroke}{rgb}{0.000000,0.000000,0.000000}%
\pgfsetstrokecolor{currentstroke}%
\pgfsetdash{}{0pt}%
\pgfpathmoveto{\pgfqpoint{1.700554in}{1.151209in}}%
\pgfpathlineto{\pgfqpoint{1.718890in}{1.163049in}}%
\pgfpathlineto{\pgfqpoint{1.726433in}{1.170734in}}%
\pgfpathlineto{\pgfqpoint{1.733962in}{1.178660in}}%
\pgfpathlineto{\pgfqpoint{1.715651in}{1.166960in}}%
\pgfpathlineto{\pgfqpoint{1.708110in}{1.158961in}}%
\pgfusepath{fill}%
\end{pgfscope}%
\begin{pgfscope}%
\pgfpathrectangle{\pgfqpoint{0.407592in}{0.222904in}}{\pgfqpoint{1.411163in}{1.411163in}}%
\pgfusepath{clip}%
\pgfsetbuttcap%
\pgfsetroundjoin%
\definecolor{currentfill}{rgb}{0.909804,0.698039,0.423529}%
\pgfsetfillcolor{currentfill}%
\pgfsetfillopacity{0.120000}%
\pgfsetlinewidth{0.000000pt}%
\definecolor{currentstroke}{rgb}{0.000000,0.000000,0.000000}%
\pgfsetstrokecolor{currentstroke}%
\pgfsetdash{}{0pt}%
\pgfpathmoveto{\pgfqpoint{1.208432in}{0.880184in}}%
\pgfpathlineto{\pgfqpoint{1.225674in}{0.886391in}}%
\pgfpathlineto{\pgfqpoint{1.242984in}{0.893102in}}%
\pgfpathlineto{\pgfqpoint{1.251186in}{0.893415in}}%
\pgfpathlineto{\pgfqpoint{1.259358in}{0.894093in}}%
\pgfpathlineto{\pgfqpoint{1.242081in}{0.887146in}}%
\pgfpathlineto{\pgfqpoint{1.224872in}{0.880702in}}%
\pgfpathlineto{\pgfqpoint{1.216667in}{0.880254in}}%
\pgfpathlineto{\pgfqpoint{1.208432in}{0.880184in}}%
\pgfpathclose%
\pgfusepath{fill}%
\end{pgfscope}%
\begin{pgfscope}%
\pgfpathrectangle{\pgfqpoint{0.407592in}{0.222904in}}{\pgfqpoint{1.411163in}{1.411163in}}%
\pgfusepath{clip}%
\pgfsetbuttcap%
\pgfsetroundjoin%
\definecolor{currentfill}{rgb}{0.909804,0.698039,0.423529}%
\pgfsetfillcolor{currentfill}%
\pgfsetfillopacity{0.120000}%
\pgfsetlinewidth{0.000000pt}%
\definecolor{currentstroke}{rgb}{0.000000,0.000000,0.000000}%
\pgfsetstrokecolor{currentstroke}%
\pgfsetdash{}{0pt}%
\pgfpathmoveto{\pgfqpoint{1.259358in}{0.894093in}}%
\pgfpathlineto{\pgfqpoint{1.276706in}{0.901507in}}%
\pgfpathlineto{\pgfqpoint{1.294124in}{0.909358in}}%
\pgfpathlineto{\pgfqpoint{1.302238in}{0.910618in}}%
\pgfpathlineto{\pgfqpoint{1.310325in}{0.912246in}}%
\pgfpathlineto{\pgfqpoint{1.292937in}{0.904199in}}%
\pgfpathlineto{\pgfqpoint{1.275621in}{0.896583in}}%
\pgfpathlineto{\pgfqpoint{1.267503in}{0.895145in}}%
\pgfpathlineto{\pgfqpoint{1.259358in}{0.894093in}}%
\pgfpathclose%
\pgfusepath{fill}%
\end{pgfscope}%
\begin{pgfscope}%
\pgfpathrectangle{\pgfqpoint{0.407592in}{0.222904in}}{\pgfqpoint{1.411163in}{1.411163in}}%
\pgfusepath{clip}%
\pgfsetbuttcap%
\pgfsetroundjoin%
\definecolor{currentfill}{rgb}{0.909804,0.698039,0.423529}%
\pgfsetfillcolor{currentfill}%
\pgfsetfillopacity{0.120000}%
\pgfsetlinewidth{0.000000pt}%
\definecolor{currentstroke}{rgb}{0.000000,0.000000,0.000000}%
\pgfsetstrokecolor{currentstroke}%
\pgfsetdash{}{0pt}%
\pgfpathmoveto{\pgfqpoint{1.157510in}{0.870880in}}%
\pgfpathlineto{\pgfqpoint{1.174659in}{0.875716in}}%
\pgfpathlineto{\pgfqpoint{1.191871in}{0.881125in}}%
\pgfpathlineto{\pgfqpoint{1.200167in}{0.880478in}}%
\pgfpathlineto{\pgfqpoint{1.208432in}{0.880184in}}%
\pgfpathlineto{\pgfqpoint{1.191256in}{0.874516in}}%
\pgfpathlineto{\pgfqpoint{1.174144in}{0.869424in}}%
\pgfpathlineto{\pgfqpoint{1.165843in}{0.869971in}}%
\pgfpathlineto{\pgfqpoint{1.157510in}{0.870880in}}%
\pgfpathclose%
\pgfusepath{fill}%
\end{pgfscope}%
\begin{pgfscope}%
\pgfpathrectangle{\pgfqpoint{0.407592in}{0.222904in}}{\pgfqpoint{1.411163in}{1.411163in}}%
\pgfusepath{clip}%
\pgfsetbuttcap%
\pgfsetroundjoin%
\definecolor{currentfill}{rgb}{0.909804,0.698039,0.423529}%
\pgfsetfillcolor{currentfill}%
\pgfsetfillopacity{0.120000}%
\pgfsetlinewidth{0.000000pt}%
\definecolor{currentstroke}{rgb}{0.000000,0.000000,0.000000}%
\pgfsetstrokecolor{currentstroke}%
\pgfsetdash{}{0pt}%
\pgfpathmoveto{\pgfqpoint{1.530687in}{1.029188in}}%
\pgfpathlineto{\pgfqpoint{1.548635in}{1.040253in}}%
\pgfpathlineto{\pgfqpoint{1.566662in}{1.051529in}}%
\pgfpathlineto{\pgfqpoint{1.574388in}{1.057265in}}%
\pgfpathlineto{\pgfqpoint{1.582098in}{1.063303in}}%
\pgfpathlineto{\pgfqpoint{1.564095in}{1.052081in}}%
\pgfpathlineto{\pgfqpoint{1.546173in}{1.041071in}}%
\pgfpathlineto{\pgfqpoint{1.538438in}{1.034971in}}%
\pgfpathlineto{\pgfqpoint{1.530687in}{1.029188in}}%
\pgfpathclose%
\pgfusepath{fill}%
\end{pgfscope}%
\begin{pgfscope}%
\pgfpathrectangle{\pgfqpoint{0.407592in}{0.222904in}}{\pgfqpoint{1.411163in}{1.411163in}}%
\pgfusepath{clip}%
\pgfsetbuttcap%
\pgfsetroundjoin%
\definecolor{currentfill}{rgb}{0.909804,0.698039,0.423529}%
\pgfsetfillcolor{currentfill}%
\pgfsetfillopacity{0.120000}%
\pgfsetlinewidth{0.000000pt}%
\definecolor{currentstroke}{rgb}{0.000000,0.000000,0.000000}%
\pgfsetstrokecolor{currentstroke}%
\pgfsetdash{}{0pt}%
\pgfpathmoveto{\pgfqpoint{1.310325in}{0.912246in}}%
\pgfpathlineto{\pgfqpoint{1.327786in}{0.920695in}}%
\pgfpathlineto{\pgfqpoint{1.345321in}{0.929517in}}%
\pgfpathlineto{\pgfqpoint{1.353354in}{0.931683in}}%
\pgfpathlineto{\pgfqpoint{1.361364in}{0.934214in}}%
\pgfpathlineto{\pgfqpoint{1.343857in}{0.925242in}}%
\pgfpathlineto{\pgfqpoint{1.326425in}{0.916638in}}%
\pgfpathlineto{\pgfqpoint{1.318388in}{0.914250in}}%
\pgfpathlineto{\pgfqpoint{1.310325in}{0.912246in}}%
\pgfpathclose%
\pgfusepath{fill}%
\end{pgfscope}%
\begin{pgfscope}%
\pgfpathrectangle{\pgfqpoint{0.407592in}{0.222904in}}{\pgfqpoint{1.411163in}{1.411163in}}%
\pgfusepath{clip}%
\pgfsetbuttcap%
\pgfsetroundjoin%
\definecolor{currentfill}{rgb}{0.909804,0.698039,0.423529}%
\pgfsetfillcolor{currentfill}%
\pgfsetfillopacity{0.120000}%
\pgfsetlinewidth{0.000000pt}%
\definecolor{currentstroke}{rgb}{0.000000,0.000000,0.000000}%
\pgfsetstrokecolor{currentstroke}%
\pgfsetdash{}{0pt}%
\pgfpathmoveto{\pgfqpoint{1.106554in}{0.866422in}}%
\pgfpathlineto{\pgfqpoint{1.123621in}{0.869754in}}%
\pgfpathlineto{\pgfqpoint{1.140747in}{0.873720in}}%
\pgfpathlineto{\pgfqpoint{1.149145in}{0.872134in}}%
\pgfpathlineto{\pgfqpoint{1.157510in}{0.870880in}}%
\pgfpathlineto{\pgfqpoint{1.140424in}{0.866651in}}%
\pgfpathlineto{\pgfqpoint{1.123397in}{0.863066in}}%
\pgfpathlineto{\pgfqpoint{1.114992in}{0.864577in}}%
\pgfpathlineto{\pgfqpoint{1.106554in}{0.866422in}}%
\pgfpathclose%
\pgfusepath{fill}%
\end{pgfscope}%
\begin{pgfscope}%
\pgfpathrectangle{\pgfqpoint{0.407592in}{0.222904in}}{\pgfqpoint{1.411163in}{1.411163in}}%
\pgfusepath{clip}%
\pgfsetbuttcap%
\pgfsetroundjoin%
\definecolor{currentfill}{rgb}{0.909804,0.698039,0.423529}%
\pgfsetfillcolor{currentfill}%
\pgfsetfillopacity{0.120000}%
\pgfsetlinewidth{0.000000pt}%
\definecolor{currentstroke}{rgb}{0.000000,0.000000,0.000000}%
\pgfsetstrokecolor{currentstroke}%
\pgfsetdash{}{0pt}%
\pgfpathmoveto{\pgfqpoint{0.919375in}{0.887319in}}%
\pgfpathlineto{\pgfqpoint{0.936198in}{0.884363in}}%
\pgfpathlineto{\pgfqpoint{0.953060in}{0.882146in}}%
\pgfpathlineto{\pgfqpoint{0.961858in}{0.878052in}}%
\pgfpathlineto{\pgfqpoint{0.970617in}{0.874197in}}%
\pgfpathlineto{\pgfqpoint{0.953804in}{0.876321in}}%
\pgfpathlineto{\pgfqpoint{0.937030in}{0.879215in}}%
\pgfpathlineto{\pgfqpoint{0.928222in}{0.883148in}}%
\pgfpathlineto{\pgfqpoint{0.919375in}{0.887319in}}%
\pgfpathclose%
\pgfusepath{fill}%
\end{pgfscope}%
\begin{pgfscope}%
\pgfpathrectangle{\pgfqpoint{0.407592in}{0.222904in}}{\pgfqpoint{1.411163in}{1.411163in}}%
\pgfusepath{clip}%
\pgfsetbuttcap%
\pgfsetroundjoin%
\definecolor{currentfill}{rgb}{0.909804,0.698039,0.423529}%
\pgfsetfillcolor{currentfill}%
\pgfsetfillopacity{0.120000}%
\pgfsetlinewidth{0.000000pt}%
\definecolor{currentstroke}{rgb}{0.000000,0.000000,0.000000}%
\pgfsetstrokecolor{currentstroke}%
\pgfsetdash{}{0pt}%
\pgfpathmoveto{\pgfqpoint{1.361364in}{0.934214in}}%
\pgfpathlineto{\pgfqpoint{1.378945in}{0.943529in}}%
\pgfpathlineto{\pgfqpoint{1.396603in}{0.953164in}}%
\pgfpathlineto{\pgfqpoint{1.404561in}{0.956177in}}%
\pgfpathlineto{\pgfqpoint{1.412498in}{0.959546in}}%
\pgfpathlineto{\pgfqpoint{1.394867in}{0.949810in}}%
\pgfpathlineto{\pgfqpoint{1.377313in}{0.940390in}}%
\pgfpathlineto{\pgfqpoint{1.369350in}{0.937115in}}%
\pgfpathlineto{\pgfqpoint{1.361364in}{0.934214in}}%
\pgfpathclose%
\pgfusepath{fill}%
\end{pgfscope}%
\begin{pgfscope}%
\pgfpathrectangle{\pgfqpoint{0.407592in}{0.222904in}}{\pgfqpoint{1.411163in}{1.411163in}}%
\pgfusepath{clip}%
\pgfsetbuttcap%
\pgfsetroundjoin%
\definecolor{currentfill}{rgb}{0.909804,0.698039,0.423529}%
\pgfsetfillcolor{currentfill}%
\pgfsetfillopacity{0.120000}%
\pgfsetlinewidth{0.000000pt}%
\definecolor{currentstroke}{rgb}{0.000000,0.000000,0.000000}%
\pgfsetstrokecolor{currentstroke}%
\pgfsetdash{}{0pt}%
\pgfpathmoveto{\pgfqpoint{1.055521in}{0.866897in}}%
\pgfpathlineto{\pgfqpoint{1.072519in}{0.868636in}}%
\pgfpathlineto{\pgfqpoint{1.089572in}{0.871057in}}%
\pgfpathlineto{\pgfqpoint{1.098081in}{0.868587in}}%
\pgfpathlineto{\pgfqpoint{1.106554in}{0.866422in}}%
\pgfpathlineto{\pgfqpoint{1.089544in}{0.863757in}}%
\pgfpathlineto{\pgfqpoint{1.072588in}{0.861791in}}%
\pgfpathlineto{\pgfqpoint{1.064073in}{0.864192in}}%
\pgfpathlineto{\pgfqpoint{1.055521in}{0.866897in}}%
\pgfpathclose%
\pgfusepath{fill}%
\end{pgfscope}%
\begin{pgfscope}%
\pgfpathrectangle{\pgfqpoint{0.407592in}{0.222904in}}{\pgfqpoint{1.411163in}{1.411163in}}%
\pgfusepath{clip}%
\pgfsetbuttcap%
\pgfsetroundjoin%
\definecolor{currentfill}{rgb}{0.909804,0.698039,0.423529}%
\pgfsetfillcolor{currentfill}%
\pgfsetfillopacity{0.120000}%
\pgfsetlinewidth{0.000000pt}%
\definecolor{currentstroke}{rgb}{0.000000,0.000000,0.000000}%
\pgfsetstrokecolor{currentstroke}%
\pgfsetdash{}{0pt}%
\pgfpathmoveto{\pgfqpoint{1.582098in}{1.063303in}}%
\pgfpathlineto{\pgfqpoint{1.600181in}{1.074725in}}%
\pgfpathlineto{\pgfqpoint{1.618346in}{1.086338in}}%
\pgfpathlineto{\pgfqpoint{1.626015in}{1.092599in}}%
\pgfpathlineto{\pgfqpoint{1.633667in}{1.099142in}}%
\pgfpathlineto{\pgfqpoint{1.615527in}{1.087610in}}%
\pgfpathlineto{\pgfqpoint{1.597469in}{1.076270in}}%
\pgfpathlineto{\pgfqpoint{1.589791in}{1.069639in}}%
\pgfpathlineto{\pgfqpoint{1.582098in}{1.063303in}}%
\pgfpathclose%
\pgfusepath{fill}%
\end{pgfscope}%
\begin{pgfscope}%
\pgfpathrectangle{\pgfqpoint{0.407592in}{0.222904in}}{\pgfqpoint{1.411163in}{1.411163in}}%
\pgfusepath{clip}%
\pgfsetbuttcap%
\pgfsetroundjoin%
\definecolor{currentfill}{rgb}{0.909804,0.698039,0.423529}%
\pgfsetfillcolor{currentfill}%
\pgfsetfillopacity{0.120000}%
\pgfsetlinewidth{0.000000pt}%
\definecolor{currentstroke}{rgb}{0.000000,0.000000,0.000000}%
\pgfsetstrokecolor{currentstroke}%
\pgfsetdash{}{0pt}%
\pgfpathmoveto{\pgfqpoint{1.412498in}{0.959546in}}%
\pgfpathlineto{\pgfqpoint{1.430206in}{0.969575in}}%
\pgfpathlineto{\pgfqpoint{1.447992in}{0.979879in}}%
\pgfpathlineto{\pgfqpoint{1.455881in}{0.983670in}}%
\pgfpathlineto{\pgfqpoint{1.463749in}{0.987802in}}%
\pgfpathlineto{\pgfqpoint{1.445989in}{0.977444in}}%
\pgfpathlineto{\pgfqpoint{1.428308in}{0.967357in}}%
\pgfpathlineto{\pgfqpoint{1.420413in}{0.963272in}}%
\pgfpathlineto{\pgfqpoint{1.412498in}{0.959546in}}%
\pgfpathclose%
\pgfusepath{fill}%
\end{pgfscope}%
\begin{pgfscope}%
\pgfpathrectangle{\pgfqpoint{0.407592in}{0.222904in}}{\pgfqpoint{1.411163in}{1.411163in}}%
\pgfusepath{clip}%
\pgfsetbuttcap%
\pgfsetroundjoin%
\definecolor{currentfill}{rgb}{0.909804,0.698039,0.423529}%
\pgfsetfillcolor{currentfill}%
\pgfsetfillopacity{0.120000}%
\pgfsetlinewidth{0.000000pt}%
\definecolor{currentstroke}{rgb}{0.000000,0.000000,0.000000}%
\pgfsetstrokecolor{currentstroke}%
\pgfsetdash{}{0pt}%
\pgfpathmoveto{\pgfqpoint{0.782712in}{0.938839in}}%
\pgfpathlineto{\pgfqpoint{0.799451in}{0.931719in}}%
\pgfpathlineto{\pgfqpoint{0.816216in}{0.925255in}}%
\pgfpathlineto{\pgfqpoint{0.825331in}{0.920284in}}%
\pgfpathlineto{\pgfqpoint{0.834403in}{0.915504in}}%
\pgfpathlineto{\pgfqpoint{0.817689in}{0.922053in}}%
\pgfpathlineto{\pgfqpoint{0.801000in}{0.929284in}}%
\pgfpathlineto{\pgfqpoint{0.791878in}{0.933964in}}%
\pgfpathlineto{\pgfqpoint{0.782712in}{0.938839in}}%
\pgfpathclose%
\pgfusepath{fill}%
\end{pgfscope}%
\begin{pgfscope}%
\pgfpathrectangle{\pgfqpoint{0.407592in}{0.222904in}}{\pgfqpoint{1.411163in}{1.411163in}}%
\pgfusepath{clip}%
\pgfsetbuttcap%
\pgfsetroundjoin%
\definecolor{currentfill}{rgb}{0.909804,0.698039,0.423529}%
\pgfsetfillcolor{currentfill}%
\pgfsetfillopacity{0.120000}%
\pgfsetlinewidth{0.000000pt}%
\definecolor{currentstroke}{rgb}{0.000000,0.000000,0.000000}%
\pgfsetstrokecolor{currentstroke}%
\pgfsetdash{}{0pt}%
\pgfpathmoveto{\pgfqpoint{1.004369in}{0.872218in}}%
\pgfpathlineto{\pgfqpoint{1.021313in}{0.872333in}}%
\pgfpathlineto{\pgfqpoint{1.038306in}{0.873158in}}%
\pgfpathlineto{\pgfqpoint{1.046932in}{0.869890in}}%
\pgfpathlineto{\pgfqpoint{1.055521in}{0.866897in}}%
\pgfpathlineto{\pgfqpoint{1.038574in}{0.865867in}}%
\pgfpathlineto{\pgfqpoint{1.021676in}{0.865571in}}%
\pgfpathlineto{\pgfqpoint{1.013042in}{0.868758in}}%
\pgfpathlineto{\pgfqpoint{1.004369in}{0.872218in}}%
\pgfpathclose%
\pgfusepath{fill}%
\end{pgfscope}%
\begin{pgfscope}%
\pgfpathrectangle{\pgfqpoint{0.407592in}{0.222904in}}{\pgfqpoint{1.411163in}{1.411163in}}%
\pgfusepath{clip}%
\pgfsetbuttcap%
\pgfsetroundjoin%
\definecolor{currentfill}{rgb}{0.909804,0.698039,0.423529}%
\pgfsetfillcolor{currentfill}%
\pgfsetfillopacity{0.120000}%
\pgfsetlinewidth{0.000000pt}%
\definecolor{currentstroke}{rgb}{0.000000,0.000000,0.000000}%
\pgfsetstrokecolor{currentstroke}%
\pgfsetdash{}{0pt}%
\pgfpathmoveto{\pgfqpoint{0.867917in}{0.904496in}}%
\pgfpathlineto{\pgfqpoint{0.884720in}{0.900055in}}%
\pgfpathlineto{\pgfqpoint{0.901557in}{0.896327in}}%
\pgfpathlineto{\pgfqpoint{0.910487in}{0.891716in}}%
\pgfpathlineto{\pgfqpoint{0.919375in}{0.887319in}}%
\pgfpathlineto{\pgfqpoint{0.902588in}{0.891019in}}%
\pgfpathlineto{\pgfqpoint{0.885835in}{0.895462in}}%
\pgfpathlineto{\pgfqpoint{0.876896in}{0.899873in}}%
\pgfpathlineto{\pgfqpoint{0.867917in}{0.904496in}}%
\pgfpathclose%
\pgfusepath{fill}%
\end{pgfscope}%
\begin{pgfscope}%
\pgfpathrectangle{\pgfqpoint{0.407592in}{0.222904in}}{\pgfqpoint{1.411163in}{1.411163in}}%
\pgfusepath{clip}%
\pgfsetbuttcap%
\pgfsetroundjoin%
\definecolor{currentfill}{rgb}{0.909804,0.698039,0.423529}%
\pgfsetfillcolor{currentfill}%
\pgfsetfillopacity{0.120000}%
\pgfsetlinewidth{0.000000pt}%
\definecolor{currentstroke}{rgb}{0.000000,0.000000,0.000000}%
\pgfsetstrokecolor{currentstroke}%
\pgfsetdash{}{0pt}%
\pgfpathmoveto{\pgfqpoint{1.463749in}{0.987802in}}%
\pgfpathlineto{\pgfqpoint{1.481588in}{0.998415in}}%
\pgfpathlineto{\pgfqpoint{1.499505in}{1.009267in}}%
\pgfpathlineto{\pgfqpoint{1.507329in}{1.013763in}}%
\pgfpathlineto{\pgfqpoint{1.515133in}{1.018582in}}%
\pgfpathlineto{\pgfqpoint{1.497241in}{1.007717in}}%
\pgfpathlineto{\pgfqpoint{1.479428in}{0.997090in}}%
\pgfpathlineto{\pgfqpoint{1.471598in}{0.992276in}}%
\pgfpathlineto{\pgfqpoint{1.463749in}{0.987802in}}%
\pgfpathclose%
\pgfusepath{fill}%
\end{pgfscope}%
\begin{pgfscope}%
\pgfpathrectangle{\pgfqpoint{0.407592in}{0.222904in}}{\pgfqpoint{1.411163in}{1.411163in}}%
\pgfusepath{clip}%
\pgfsetbuttcap%
\pgfsetroundjoin%
\definecolor{currentfill}{rgb}{0.909804,0.698039,0.423529}%
\pgfsetfillcolor{currentfill}%
\pgfsetfillopacity{0.120000}%
\pgfsetlinewidth{0.000000pt}%
\definecolor{currentstroke}{rgb}{0.000000,0.000000,0.000000}%
\pgfsetstrokecolor{currentstroke}%
\pgfsetdash{}{0pt}%
\pgfpathmoveto{\pgfqpoint{1.633667in}{1.099142in}}%
\pgfpathlineto{\pgfqpoint{1.651889in}{1.110858in}}%
\pgfpathlineto{\pgfqpoint{1.670193in}{1.122749in}}%
\pgfpathlineto{\pgfqpoint{1.677805in}{1.129474in}}%
\pgfpathlineto{\pgfqpoint{1.685403in}{1.136461in}}%
\pgfpathlineto{\pgfqpoint{1.667124in}{1.124671in}}%
\pgfpathlineto{\pgfqpoint{1.648927in}{1.113059in}}%
\pgfpathlineto{\pgfqpoint{1.641305in}{1.105963in}}%
\pgfpathlineto{\pgfqpoint{1.633667in}{1.099142in}}%
\pgfpathclose%
\pgfusepath{fill}%
\end{pgfscope}%
\begin{pgfscope}%
\pgfpathrectangle{\pgfqpoint{0.407592in}{0.222904in}}{\pgfqpoint{1.411163in}{1.411163in}}%
\pgfusepath{clip}%
\pgfsetbuttcap%
\pgfsetroundjoin%
\definecolor{currentfill}{rgb}{0.909804,0.698039,0.423529}%
\pgfsetfillcolor{currentfill}%
\pgfsetfillopacity{0.120000}%
\pgfsetlinewidth{0.000000pt}%
\definecolor{currentstroke}{rgb}{0.000000,0.000000,0.000000}%
\pgfsetstrokecolor{currentstroke}%
\pgfsetdash{}{0pt}%
\pgfpathmoveto{\pgfqpoint{1.685403in}{1.136461in}}%
\pgfpathlineto{\pgfqpoint{1.703764in}{1.148422in}}%
\pgfpathlineto{\pgfqpoint{1.711334in}{1.155610in}}%
\pgfpathlineto{\pgfqpoint{1.718890in}{1.163049in}}%
\pgfpathlineto{\pgfqpoint{1.700554in}{1.151209in}}%
\pgfpathlineto{\pgfqpoint{1.692986in}{1.143707in}}%
\pgfusepath{fill}%
\end{pgfscope}%
\begin{pgfscope}%
\pgfpathrectangle{\pgfqpoint{0.407592in}{0.222904in}}{\pgfqpoint{1.411163in}{1.411163in}}%
\pgfusepath{clip}%
\pgfsetbuttcap%
\pgfsetroundjoin%
\definecolor{currentfill}{rgb}{0.909804,0.698039,0.423529}%
\pgfsetfillcolor{currentfill}%
\pgfsetfillopacity{0.120000}%
\pgfsetlinewidth{0.000000pt}%
\definecolor{currentstroke}{rgb}{0.000000,0.000000,0.000000}%
\pgfsetstrokecolor{currentstroke}%
\pgfsetdash{}{0pt}%
\pgfpathmoveto{\pgfqpoint{0.953060in}{0.882146in}}%
\pgfpathlineto{\pgfqpoint{0.969962in}{0.880661in}}%
\pgfpathlineto{\pgfqpoint{0.986908in}{0.879895in}}%
\pgfpathlineto{\pgfqpoint{0.995658in}{0.875934in}}%
\pgfpathlineto{\pgfqpoint{1.004369in}{0.872218in}}%
\pgfpathlineto{\pgfqpoint{0.987471in}{0.872834in}}%
\pgfpathlineto{\pgfqpoint{0.970617in}{0.874197in}}%
\pgfpathlineto{\pgfqpoint{0.961858in}{0.878052in}}%
\pgfpathlineto{\pgfqpoint{0.953060in}{0.882146in}}%
\pgfpathclose%
\pgfusepath{fill}%
\end{pgfscope}%
\begin{pgfscope}%
\pgfpathrectangle{\pgfqpoint{0.407592in}{0.222904in}}{\pgfqpoint{1.411163in}{1.411163in}}%
\pgfusepath{clip}%
\pgfsetbuttcap%
\pgfsetroundjoin%
\definecolor{currentfill}{rgb}{0.909804,0.698039,0.423529}%
\pgfsetfillcolor{currentfill}%
\pgfsetfillopacity{0.120000}%
\pgfsetlinewidth{0.000000pt}%
\definecolor{currentstroke}{rgb}{0.000000,0.000000,0.000000}%
\pgfsetstrokecolor{currentstroke}%
\pgfsetdash{}{0pt}%
\pgfpathmoveto{\pgfqpoint{1.515133in}{1.018582in}}%
\pgfpathlineto{\pgfqpoint{1.533105in}{1.029670in}}%
\pgfpathlineto{\pgfqpoint{1.551158in}{1.040969in}}%
\pgfpathlineto{\pgfqpoint{1.558919in}{1.046096in}}%
\pgfpathlineto{\pgfqpoint{1.566662in}{1.051529in}}%
\pgfpathlineto{\pgfqpoint{1.548635in}{1.040253in}}%
\pgfpathlineto{\pgfqpoint{1.530687in}{1.029188in}}%
\pgfpathlineto{\pgfqpoint{1.522919in}{1.023724in}}%
\pgfpathlineto{\pgfqpoint{1.515133in}{1.018582in}}%
\pgfpathclose%
\pgfusepath{fill}%
\end{pgfscope}%
\begin{pgfscope}%
\pgfpathrectangle{\pgfqpoint{0.407592in}{0.222904in}}{\pgfqpoint{1.411163in}{1.411163in}}%
\pgfusepath{clip}%
\pgfsetbuttcap%
\pgfsetroundjoin%
\definecolor{currentfill}{rgb}{0.909804,0.698039,0.423529}%
\pgfsetfillcolor{currentfill}%
\pgfsetfillopacity{0.120000}%
\pgfsetlinewidth{0.000000pt}%
\definecolor{currentstroke}{rgb}{0.000000,0.000000,0.000000}%
\pgfsetstrokecolor{currentstroke}%
\pgfsetdash{}{0pt}%
\pgfpathmoveto{\pgfqpoint{1.242984in}{0.893102in}}%
\pgfpathlineto{\pgfqpoint{1.260364in}{0.900287in}}%
\pgfpathlineto{\pgfqpoint{1.277814in}{0.907913in}}%
\pgfpathlineto{\pgfqpoint{1.285983in}{0.908460in}}%
\pgfpathlineto{\pgfqpoint{1.294124in}{0.909358in}}%
\pgfpathlineto{\pgfqpoint{1.276706in}{0.901507in}}%
\pgfpathlineto{\pgfqpoint{1.259358in}{0.894093in}}%
\pgfpathlineto{\pgfqpoint{1.251186in}{0.893415in}}%
\pgfpathlineto{\pgfqpoint{1.242984in}{0.893102in}}%
\pgfpathclose%
\pgfusepath{fill}%
\end{pgfscope}%
\begin{pgfscope}%
\pgfpathrectangle{\pgfqpoint{0.407592in}{0.222904in}}{\pgfqpoint{1.411163in}{1.411163in}}%
\pgfusepath{clip}%
\pgfsetbuttcap%
\pgfsetroundjoin%
\definecolor{currentfill}{rgb}{0.909804,0.698039,0.423529}%
\pgfsetfillcolor{currentfill}%
\pgfsetfillopacity{0.120000}%
\pgfsetlinewidth{0.000000pt}%
\definecolor{currentstroke}{rgb}{0.000000,0.000000,0.000000}%
\pgfsetstrokecolor{currentstroke}%
\pgfsetdash{}{0pt}%
\pgfpathmoveto{\pgfqpoint{1.191871in}{0.881125in}}%
\pgfpathlineto{\pgfqpoint{1.209148in}{0.887074in}}%
\pgfpathlineto{\pgfqpoint{1.226493in}{0.893529in}}%
\pgfpathlineto{\pgfqpoint{1.234754in}{0.893144in}}%
\pgfpathlineto{\pgfqpoint{1.242984in}{0.893102in}}%
\pgfpathlineto{\pgfqpoint{1.225674in}{0.886391in}}%
\pgfpathlineto{\pgfqpoint{1.208432in}{0.880184in}}%
\pgfpathlineto{\pgfqpoint{1.200167in}{0.880478in}}%
\pgfpathlineto{\pgfqpoint{1.191871in}{0.881125in}}%
\pgfpathclose%
\pgfusepath{fill}%
\end{pgfscope}%
\begin{pgfscope}%
\pgfpathrectangle{\pgfqpoint{0.407592in}{0.222904in}}{\pgfqpoint{1.411163in}{1.411163in}}%
\pgfusepath{clip}%
\pgfsetbuttcap%
\pgfsetroundjoin%
\definecolor{currentfill}{rgb}{0.909804,0.698039,0.423529}%
\pgfsetfillcolor{currentfill}%
\pgfsetfillopacity{0.120000}%
\pgfsetlinewidth{0.000000pt}%
\definecolor{currentstroke}{rgb}{0.000000,0.000000,0.000000}%
\pgfsetstrokecolor{currentstroke}%
\pgfsetdash{}{0pt}%
\pgfpathmoveto{\pgfqpoint{0.816216in}{0.925255in}}%
\pgfpathlineto{\pgfqpoint{0.833008in}{0.919459in}}%
\pgfpathlineto{\pgfqpoint{0.849831in}{0.914337in}}%
\pgfpathlineto{\pgfqpoint{0.858895in}{0.909321in}}%
\pgfpathlineto{\pgfqpoint{0.867917in}{0.904496in}}%
\pgfpathlineto{\pgfqpoint{0.851145in}{0.909648in}}%
\pgfpathlineto{\pgfqpoint{0.834403in}{0.915504in}}%
\pgfpathlineto{\pgfqpoint{0.825331in}{0.920284in}}%
\pgfpathlineto{\pgfqpoint{0.816216in}{0.925255in}}%
\pgfpathclose%
\pgfusepath{fill}%
\end{pgfscope}%
\begin{pgfscope}%
\pgfpathrectangle{\pgfqpoint{0.407592in}{0.222904in}}{\pgfqpoint{1.411163in}{1.411163in}}%
\pgfusepath{clip}%
\pgfsetbuttcap%
\pgfsetroundjoin%
\definecolor{currentfill}{rgb}{0.909804,0.698039,0.423529}%
\pgfsetfillcolor{currentfill}%
\pgfsetfillopacity{0.120000}%
\pgfsetlinewidth{0.000000pt}%
\definecolor{currentstroke}{rgb}{0.000000,0.000000,0.000000}%
\pgfsetstrokecolor{currentstroke}%
\pgfsetdash{}{0pt}%
\pgfpathmoveto{\pgfqpoint{1.294124in}{0.909358in}}%
\pgfpathlineto{\pgfqpoint{1.311614in}{0.917618in}}%
\pgfpathlineto{\pgfqpoint{1.329179in}{0.926257in}}%
\pgfpathlineto{\pgfqpoint{1.337263in}{0.927711in}}%
\pgfpathlineto{\pgfqpoint{1.345321in}{0.929517in}}%
\pgfpathlineto{\pgfqpoint{1.327786in}{0.920695in}}%
\pgfpathlineto{\pgfqpoint{1.310325in}{0.912246in}}%
\pgfpathlineto{\pgfqpoint{1.302238in}{0.910618in}}%
\pgfpathlineto{\pgfqpoint{1.294124in}{0.909358in}}%
\pgfpathclose%
\pgfusepath{fill}%
\end{pgfscope}%
\begin{pgfscope}%
\pgfpathrectangle{\pgfqpoint{0.407592in}{0.222904in}}{\pgfqpoint{1.411163in}{1.411163in}}%
\pgfusepath{clip}%
\pgfsetbuttcap%
\pgfsetroundjoin%
\definecolor{currentfill}{rgb}{0.909804,0.698039,0.423529}%
\pgfsetfillcolor{currentfill}%
\pgfsetfillopacity{0.120000}%
\pgfsetlinewidth{0.000000pt}%
\definecolor{currentstroke}{rgb}{0.000000,0.000000,0.000000}%
\pgfsetstrokecolor{currentstroke}%
\pgfsetdash{}{0pt}%
\pgfpathmoveto{\pgfqpoint{1.140747in}{0.873720in}}%
\pgfpathlineto{\pgfqpoint{1.157933in}{0.878288in}}%
\pgfpathlineto{\pgfqpoint{1.175183in}{0.883426in}}%
\pgfpathlineto{\pgfqpoint{1.183543in}{0.882112in}}%
\pgfpathlineto{\pgfqpoint{1.191871in}{0.881125in}}%
\pgfpathlineto{\pgfqpoint{1.174659in}{0.875716in}}%
\pgfpathlineto{\pgfqpoint{1.157510in}{0.870880in}}%
\pgfpathlineto{\pgfqpoint{1.149145in}{0.872134in}}%
\pgfpathlineto{\pgfqpoint{1.140747in}{0.873720in}}%
\pgfpathclose%
\pgfusepath{fill}%
\end{pgfscope}%
\begin{pgfscope}%
\pgfpathrectangle{\pgfqpoint{0.407592in}{0.222904in}}{\pgfqpoint{1.411163in}{1.411163in}}%
\pgfusepath{clip}%
\pgfsetbuttcap%
\pgfsetroundjoin%
\definecolor{currentfill}{rgb}{0.909804,0.698039,0.423529}%
\pgfsetfillcolor{currentfill}%
\pgfsetfillopacity{0.120000}%
\pgfsetlinewidth{0.000000pt}%
\definecolor{currentstroke}{rgb}{0.000000,0.000000,0.000000}%
\pgfsetstrokecolor{currentstroke}%
\pgfsetdash{}{0pt}%
\pgfpathmoveto{\pgfqpoint{1.345321in}{0.929517in}}%
\pgfpathlineto{\pgfqpoint{1.362930in}{0.938689in}}%
\pgfpathlineto{\pgfqpoint{1.380616in}{0.948187in}}%
\pgfpathlineto{\pgfqpoint{1.388621in}{0.950502in}}%
\pgfpathlineto{\pgfqpoint{1.396603in}{0.953164in}}%
\pgfpathlineto{\pgfqpoint{1.378945in}{0.943529in}}%
\pgfpathlineto{\pgfqpoint{1.361364in}{0.934214in}}%
\pgfpathlineto{\pgfqpoint{1.353354in}{0.931683in}}%
\pgfpathlineto{\pgfqpoint{1.345321in}{0.929517in}}%
\pgfpathclose%
\pgfusepath{fill}%
\end{pgfscope}%
\begin{pgfscope}%
\pgfpathrectangle{\pgfqpoint{0.407592in}{0.222904in}}{\pgfqpoint{1.411163in}{1.411163in}}%
\pgfusepath{clip}%
\pgfsetbuttcap%
\pgfsetroundjoin%
\definecolor{currentfill}{rgb}{0.909804,0.698039,0.423529}%
\pgfsetfillcolor{currentfill}%
\pgfsetfillopacity{0.120000}%
\pgfsetlinewidth{0.000000pt}%
\definecolor{currentstroke}{rgb}{0.000000,0.000000,0.000000}%
\pgfsetstrokecolor{currentstroke}%
\pgfsetdash{}{0pt}%
\pgfpathmoveto{\pgfqpoint{1.566662in}{1.051529in}}%
\pgfpathlineto{\pgfqpoint{1.584771in}{1.063005in}}%
\pgfpathlineto{\pgfqpoint{1.602961in}{1.074670in}}%
\pgfpathlineto{\pgfqpoint{1.610662in}{1.080360in}}%
\pgfpathlineto{\pgfqpoint{1.618346in}{1.086338in}}%
\pgfpathlineto{\pgfqpoint{1.600181in}{1.074725in}}%
\pgfpathlineto{\pgfqpoint{1.582098in}{1.063303in}}%
\pgfpathlineto{\pgfqpoint{1.574388in}{1.057265in}}%
\pgfpathlineto{\pgfqpoint{1.566662in}{1.051529in}}%
\pgfpathclose%
\pgfusepath{fill}%
\end{pgfscope}%
\begin{pgfscope}%
\pgfpathrectangle{\pgfqpoint{0.407592in}{0.222904in}}{\pgfqpoint{1.411163in}{1.411163in}}%
\pgfusepath{clip}%
\pgfsetbuttcap%
\pgfsetroundjoin%
\definecolor{currentfill}{rgb}{0.909804,0.698039,0.423529}%
\pgfsetfillcolor{currentfill}%
\pgfsetfillopacity{0.120000}%
\pgfsetlinewidth{0.000000pt}%
\definecolor{currentstroke}{rgb}{0.000000,0.000000,0.000000}%
\pgfsetstrokecolor{currentstroke}%
\pgfsetdash{}{0pt}%
\pgfpathmoveto{\pgfqpoint{1.089572in}{0.871057in}}%
\pgfpathlineto{\pgfqpoint{1.106680in}{0.874132in}}%
\pgfpathlineto{\pgfqpoint{1.123847in}{0.877831in}}%
\pgfpathlineto{\pgfqpoint{1.132314in}{0.875624in}}%
\pgfpathlineto{\pgfqpoint{1.140747in}{0.873720in}}%
\pgfpathlineto{\pgfqpoint{1.123621in}{0.869754in}}%
\pgfpathlineto{\pgfqpoint{1.106554in}{0.866422in}}%
\pgfpathlineto{\pgfqpoint{1.098081in}{0.868587in}}%
\pgfpathlineto{\pgfqpoint{1.089572in}{0.871057in}}%
\pgfpathclose%
\pgfusepath{fill}%
\end{pgfscope}%
\begin{pgfscope}%
\pgfpathrectangle{\pgfqpoint{0.407592in}{0.222904in}}{\pgfqpoint{1.411163in}{1.411163in}}%
\pgfusepath{clip}%
\pgfsetbuttcap%
\pgfsetroundjoin%
\definecolor{currentfill}{rgb}{0.909804,0.698039,0.423529}%
\pgfsetfillcolor{currentfill}%
\pgfsetfillopacity{0.120000}%
\pgfsetlinewidth{0.000000pt}%
\definecolor{currentstroke}{rgb}{0.000000,0.000000,0.000000}%
\pgfsetstrokecolor{currentstroke}%
\pgfsetdash{}{0pt}%
\pgfpathmoveto{\pgfqpoint{0.901557in}{0.896327in}}%
\pgfpathlineto{\pgfqpoint{0.918430in}{0.893312in}}%
\pgfpathlineto{\pgfqpoint{0.935341in}{0.891006in}}%
\pgfpathlineto{\pgfqpoint{0.944221in}{0.886469in}}%
\pgfpathlineto{\pgfqpoint{0.953060in}{0.882146in}}%
\pgfpathlineto{\pgfqpoint{0.936198in}{0.884363in}}%
\pgfpathlineto{\pgfqpoint{0.919375in}{0.887319in}}%
\pgfpathlineto{\pgfqpoint{0.910487in}{0.891716in}}%
\pgfpathlineto{\pgfqpoint{0.901557in}{0.896327in}}%
\pgfpathclose%
\pgfusepath{fill}%
\end{pgfscope}%
\begin{pgfscope}%
\pgfpathrectangle{\pgfqpoint{0.407592in}{0.222904in}}{\pgfqpoint{1.411163in}{1.411163in}}%
\pgfusepath{clip}%
\pgfsetbuttcap%
\pgfsetroundjoin%
\definecolor{currentfill}{rgb}{0.909804,0.698039,0.423529}%
\pgfsetfillcolor{currentfill}%
\pgfsetfillopacity{0.120000}%
\pgfsetlinewidth{0.000000pt}%
\definecolor{currentstroke}{rgb}{0.000000,0.000000,0.000000}%
\pgfsetstrokecolor{currentstroke}%
\pgfsetdash{}{0pt}%
\pgfpathmoveto{\pgfqpoint{1.396603in}{0.953164in}}%
\pgfpathlineto{\pgfqpoint{1.414337in}{0.963098in}}%
\pgfpathlineto{\pgfqpoint{1.432150in}{0.973311in}}%
\pgfpathlineto{\pgfqpoint{1.440082in}{0.976427in}}%
\pgfpathlineto{\pgfqpoint{1.447992in}{0.979879in}}%
\pgfpathlineto{\pgfqpoint{1.430206in}{0.969575in}}%
\pgfpathlineto{\pgfqpoint{1.412498in}{0.959546in}}%
\pgfpathlineto{\pgfqpoint{1.404561in}{0.956177in}}%
\pgfpathlineto{\pgfqpoint{1.396603in}{0.953164in}}%
\pgfpathclose%
\pgfusepath{fill}%
\end{pgfscope}%
\begin{pgfscope}%
\pgfpathrectangle{\pgfqpoint{0.407592in}{0.222904in}}{\pgfqpoint{1.411163in}{1.411163in}}%
\pgfusepath{clip}%
\pgfsetbuttcap%
\pgfsetroundjoin%
\definecolor{currentfill}{rgb}{0.909804,0.698039,0.423529}%
\pgfsetfillcolor{currentfill}%
\pgfsetfillopacity{0.120000}%
\pgfsetlinewidth{0.000000pt}%
\definecolor{currentstroke}{rgb}{0.000000,0.000000,0.000000}%
\pgfsetstrokecolor{currentstroke}%
\pgfsetdash{}{0pt}%
\pgfpathmoveto{\pgfqpoint{1.038306in}{0.873158in}}%
\pgfpathlineto{\pgfqpoint{1.055349in}{0.874673in}}%
\pgfpathlineto{\pgfqpoint{1.072445in}{0.876853in}}%
\pgfpathlineto{\pgfqpoint{1.081027in}{0.873817in}}%
\pgfpathlineto{\pgfqpoint{1.089572in}{0.871057in}}%
\pgfpathlineto{\pgfqpoint{1.072519in}{0.868636in}}%
\pgfpathlineto{\pgfqpoint{1.055521in}{0.866897in}}%
\pgfpathlineto{\pgfqpoint{1.046932in}{0.869890in}}%
\pgfpathlineto{\pgfqpoint{1.038306in}{0.873158in}}%
\pgfpathclose%
\pgfusepath{fill}%
\end{pgfscope}%
\begin{pgfscope}%
\pgfpathrectangle{\pgfqpoint{0.407592in}{0.222904in}}{\pgfqpoint{1.411163in}{1.411163in}}%
\pgfusepath{clip}%
\pgfsetbuttcap%
\pgfsetroundjoin%
\definecolor{currentfill}{rgb}{0.909804,0.698039,0.423529}%
\pgfsetfillcolor{currentfill}%
\pgfsetfillopacity{0.120000}%
\pgfsetlinewidth{0.000000pt}%
\definecolor{currentstroke}{rgb}{0.000000,0.000000,0.000000}%
\pgfsetstrokecolor{currentstroke}%
\pgfsetdash{}{0pt}%
\pgfpathmoveto{\pgfqpoint{1.618346in}{1.086338in}}%
\pgfpathlineto{\pgfqpoint{1.636594in}{1.098132in}}%
\pgfpathlineto{\pgfqpoint{1.654924in}{1.110099in}}%
\pgfpathlineto{\pgfqpoint{1.662566in}{1.116290in}}%
\pgfpathlineto{\pgfqpoint{1.670193in}{1.122749in}}%
\pgfpathlineto{\pgfqpoint{1.651889in}{1.110858in}}%
\pgfpathlineto{\pgfqpoint{1.633667in}{1.099142in}}%
\pgfpathlineto{\pgfqpoint{1.626015in}{1.092599in}}%
\pgfpathlineto{\pgfqpoint{1.618346in}{1.086338in}}%
\pgfpathclose%
\pgfusepath{fill}%
\end{pgfscope}%
\begin{pgfscope}%
\pgfpathrectangle{\pgfqpoint{0.407592in}{0.222904in}}{\pgfqpoint{1.411163in}{1.411163in}}%
\pgfusepath{clip}%
\pgfsetbuttcap%
\pgfsetroundjoin%
\definecolor{currentfill}{rgb}{0.909804,0.698039,0.423529}%
\pgfsetfillcolor{currentfill}%
\pgfsetfillopacity{0.120000}%
\pgfsetlinewidth{0.000000pt}%
\definecolor{currentstroke}{rgb}{0.000000,0.000000,0.000000}%
\pgfsetstrokecolor{currentstroke}%
\pgfsetdash{}{0pt}%
\pgfpathmoveto{\pgfqpoint{1.447992in}{0.979879in}}%
\pgfpathlineto{\pgfqpoint{1.465856in}{0.990441in}}%
\pgfpathlineto{\pgfqpoint{1.483800in}{1.001245in}}%
\pgfpathlineto{\pgfqpoint{1.491663in}{1.005095in}}%
\pgfpathlineto{\pgfqpoint{1.499505in}{1.009267in}}%
\pgfpathlineto{\pgfqpoint{1.481588in}{0.998415in}}%
\pgfpathlineto{\pgfqpoint{1.463749in}{0.987802in}}%
\pgfpathlineto{\pgfqpoint{1.455881in}{0.983670in}}%
\pgfpathlineto{\pgfqpoint{1.447992in}{0.979879in}}%
\pgfpathclose%
\pgfusepath{fill}%
\end{pgfscope}%
\begin{pgfscope}%
\pgfpathrectangle{\pgfqpoint{0.407592in}{0.222904in}}{\pgfqpoint{1.411163in}{1.411163in}}%
\pgfusepath{clip}%
\pgfsetbuttcap%
\pgfsetroundjoin%
\definecolor{currentfill}{rgb}{0.909804,0.698039,0.423529}%
\pgfsetfillcolor{currentfill}%
\pgfsetfillopacity{0.120000}%
\pgfsetlinewidth{0.000000pt}%
\definecolor{currentstroke}{rgb}{0.000000,0.000000,0.000000}%
\pgfsetstrokecolor{currentstroke}%
\pgfsetdash{}{0pt}%
\pgfpathmoveto{\pgfqpoint{1.670193in}{1.122749in}}%
\pgfpathlineto{\pgfqpoint{1.688581in}{1.134807in}}%
\pgfpathlineto{\pgfqpoint{1.696180in}{1.141486in}}%
\pgfpathlineto{\pgfqpoint{1.703764in}{1.148422in}}%
\pgfpathlineto{\pgfqpoint{1.685403in}{1.136461in}}%
\pgfpathlineto{\pgfqpoint{1.677805in}{1.129474in}}%
\pgfusepath{fill}%
\end{pgfscope}%
\begin{pgfscope}%
\pgfpathrectangle{\pgfqpoint{0.407592in}{0.222904in}}{\pgfqpoint{1.411163in}{1.411163in}}%
\pgfusepath{clip}%
\pgfsetbuttcap%
\pgfsetroundjoin%
\definecolor{currentfill}{rgb}{0.909804,0.698039,0.423529}%
\pgfsetfillcolor{currentfill}%
\pgfsetfillopacity{0.120000}%
\pgfsetlinewidth{0.000000pt}%
\definecolor{currentstroke}{rgb}{0.000000,0.000000,0.000000}%
\pgfsetstrokecolor{currentstroke}%
\pgfsetdash{}{0pt}%
\pgfpathmoveto{\pgfqpoint{0.986908in}{0.879895in}}%
\pgfpathlineto{\pgfqpoint{1.003899in}{0.879834in}}%
\pgfpathlineto{\pgfqpoint{1.020938in}{0.880462in}}%
\pgfpathlineto{\pgfqpoint{1.029641in}{0.876687in}}%
\pgfpathlineto{\pgfqpoint{1.038306in}{0.873158in}}%
\pgfpathlineto{\pgfqpoint{1.021313in}{0.872333in}}%
\pgfpathlineto{\pgfqpoint{1.004369in}{0.872218in}}%
\pgfpathlineto{\pgfqpoint{0.995658in}{0.875934in}}%
\pgfpathlineto{\pgfqpoint{0.986908in}{0.879895in}}%
\pgfpathclose%
\pgfusepath{fill}%
\end{pgfscope}%
\begin{pgfscope}%
\pgfpathrectangle{\pgfqpoint{0.407592in}{0.222904in}}{\pgfqpoint{1.411163in}{1.411163in}}%
\pgfusepath{clip}%
\pgfsetbuttcap%
\pgfsetroundjoin%
\definecolor{currentfill}{rgb}{0.909804,0.698039,0.423529}%
\pgfsetfillcolor{currentfill}%
\pgfsetfillopacity{0.120000}%
\pgfsetlinewidth{0.000000pt}%
\definecolor{currentstroke}{rgb}{0.000000,0.000000,0.000000}%
\pgfsetstrokecolor{currentstroke}%
\pgfsetdash{}{0pt}%
\pgfpathmoveto{\pgfqpoint{0.849831in}{0.914337in}}%
\pgfpathlineto{\pgfqpoint{0.866685in}{0.909898in}}%
\pgfpathlineto{\pgfqpoint{0.883573in}{0.906143in}}%
\pgfpathlineto{\pgfqpoint{0.892586in}{0.901139in}}%
\pgfpathlineto{\pgfqpoint{0.901557in}{0.896327in}}%
\pgfpathlineto{\pgfqpoint{0.884720in}{0.900055in}}%
\pgfpathlineto{\pgfqpoint{0.867917in}{0.904496in}}%
\pgfpathlineto{\pgfqpoint{0.858895in}{0.909321in}}%
\pgfpathlineto{\pgfqpoint{0.849831in}{0.914337in}}%
\pgfpathclose%
\pgfusepath{fill}%
\end{pgfscope}%
\begin{pgfscope}%
\pgfpathrectangle{\pgfqpoint{0.407592in}{0.222904in}}{\pgfqpoint{1.411163in}{1.411163in}}%
\pgfusepath{clip}%
\pgfsetbuttcap%
\pgfsetroundjoin%
\definecolor{currentfill}{rgb}{0.909804,0.698039,0.423529}%
\pgfsetfillcolor{currentfill}%
\pgfsetfillopacity{0.120000}%
\pgfsetlinewidth{0.000000pt}%
\definecolor{currentstroke}{rgb}{0.000000,0.000000,0.000000}%
\pgfsetstrokecolor{currentstroke}%
\pgfsetdash{}{0pt}%
\pgfpathmoveto{\pgfqpoint{1.499505in}{1.009267in}}%
\pgfpathlineto{\pgfqpoint{1.517503in}{1.020344in}}%
\pgfpathlineto{\pgfqpoint{1.535582in}{1.031634in}}%
\pgfpathlineto{\pgfqpoint{1.543380in}{1.036148in}}%
\pgfpathlineto{\pgfqpoint{1.551158in}{1.040969in}}%
\pgfpathlineto{\pgfqpoint{1.533105in}{1.029670in}}%
\pgfpathlineto{\pgfqpoint{1.515133in}{1.018582in}}%
\pgfpathlineto{\pgfqpoint{1.507329in}{1.013763in}}%
\pgfpathlineto{\pgfqpoint{1.499505in}{1.009267in}}%
\pgfpathclose%
\pgfusepath{fill}%
\end{pgfscope}%
\begin{pgfscope}%
\pgfpathrectangle{\pgfqpoint{0.407592in}{0.222904in}}{\pgfqpoint{1.411163in}{1.411163in}}%
\pgfusepath{clip}%
\pgfsetbuttcap%
\pgfsetroundjoin%
\definecolor{currentfill}{rgb}{0.909804,0.698039,0.423529}%
\pgfsetfillcolor{currentfill}%
\pgfsetfillopacity{0.120000}%
\pgfsetlinewidth{0.000000pt}%
\definecolor{currentstroke}{rgb}{0.000000,0.000000,0.000000}%
\pgfsetstrokecolor{currentstroke}%
\pgfsetdash{}{0pt}%
\pgfpathmoveto{\pgfqpoint{1.226493in}{0.893529in}}%
\pgfpathlineto{\pgfqpoint{1.243907in}{0.900461in}}%
\pgfpathlineto{\pgfqpoint{1.261390in}{0.907840in}}%
\pgfpathlineto{\pgfqpoint{1.269617in}{0.907710in}}%
\pgfpathlineto{\pgfqpoint{1.277814in}{0.907913in}}%
\pgfpathlineto{\pgfqpoint{1.260364in}{0.900287in}}%
\pgfpathlineto{\pgfqpoint{1.242984in}{0.893102in}}%
\pgfpathlineto{\pgfqpoint{1.234754in}{0.893144in}}%
\pgfpathlineto{\pgfqpoint{1.226493in}{0.893529in}}%
\pgfpathclose%
\pgfusepath{fill}%
\end{pgfscope}%
\begin{pgfscope}%
\pgfpathrectangle{\pgfqpoint{0.407592in}{0.222904in}}{\pgfqpoint{1.411163in}{1.411163in}}%
\pgfusepath{clip}%
\pgfsetbuttcap%
\pgfsetroundjoin%
\definecolor{currentfill}{rgb}{0.909804,0.698039,0.423529}%
\pgfsetfillcolor{currentfill}%
\pgfsetfillopacity{0.120000}%
\pgfsetlinewidth{0.000000pt}%
\definecolor{currentstroke}{rgb}{0.000000,0.000000,0.000000}%
\pgfsetstrokecolor{currentstroke}%
\pgfsetdash{}{0pt}%
\pgfpathmoveto{\pgfqpoint{0.935341in}{0.891006in}}%
\pgfpathlineto{\pgfqpoint{0.952293in}{0.889404in}}%
\pgfpathlineto{\pgfqpoint{0.969288in}{0.888495in}}%
\pgfpathlineto{\pgfqpoint{0.978118in}{0.884085in}}%
\pgfpathlineto{\pgfqpoint{0.986908in}{0.879895in}}%
\pgfpathlineto{\pgfqpoint{0.969962in}{0.880661in}}%
\pgfpathlineto{\pgfqpoint{0.953060in}{0.882146in}}%
\pgfpathlineto{\pgfqpoint{0.944221in}{0.886469in}}%
\pgfpathlineto{\pgfqpoint{0.935341in}{0.891006in}}%
\pgfpathclose%
\pgfusepath{fill}%
\end{pgfscope}%
\begin{pgfscope}%
\pgfpathrectangle{\pgfqpoint{0.407592in}{0.222904in}}{\pgfqpoint{1.411163in}{1.411163in}}%
\pgfusepath{clip}%
\pgfsetbuttcap%
\pgfsetroundjoin%
\definecolor{currentfill}{rgb}{0.909804,0.698039,0.423529}%
\pgfsetfillcolor{currentfill}%
\pgfsetfillopacity{0.120000}%
\pgfsetlinewidth{0.000000pt}%
\definecolor{currentstroke}{rgb}{0.000000,0.000000,0.000000}%
\pgfsetstrokecolor{currentstroke}%
\pgfsetdash{}{0pt}%
\pgfpathmoveto{\pgfqpoint{1.175183in}{0.883426in}}%
\pgfpathlineto{\pgfqpoint{1.192498in}{0.889102in}}%
\pgfpathlineto{\pgfqpoint{1.209879in}{0.895285in}}%
\pgfpathlineto{\pgfqpoint{1.218202in}{0.894247in}}%
\pgfpathlineto{\pgfqpoint{1.226493in}{0.893529in}}%
\pgfpathlineto{\pgfqpoint{1.209148in}{0.887074in}}%
\pgfpathlineto{\pgfqpoint{1.191871in}{0.881125in}}%
\pgfpathlineto{\pgfqpoint{1.183543in}{0.882112in}}%
\pgfpathlineto{\pgfqpoint{1.175183in}{0.883426in}}%
\pgfpathclose%
\pgfusepath{fill}%
\end{pgfscope}%
\begin{pgfscope}%
\pgfpathrectangle{\pgfqpoint{0.407592in}{0.222904in}}{\pgfqpoint{1.411163in}{1.411163in}}%
\pgfusepath{clip}%
\pgfsetbuttcap%
\pgfsetroundjoin%
\definecolor{currentfill}{rgb}{0.909804,0.698039,0.423529}%
\pgfsetfillcolor{currentfill}%
\pgfsetfillopacity{0.120000}%
\pgfsetlinewidth{0.000000pt}%
\definecolor{currentstroke}{rgb}{0.000000,0.000000,0.000000}%
\pgfsetstrokecolor{currentstroke}%
\pgfsetdash{}{0pt}%
\pgfpathmoveto{\pgfqpoint{1.277814in}{0.907913in}}%
\pgfpathlineto{\pgfqpoint{1.295336in}{0.915954in}}%
\pgfpathlineto{\pgfqpoint{1.312932in}{0.924381in}}%
\pgfpathlineto{\pgfqpoint{1.321069in}{0.925150in}}%
\pgfpathlineto{\pgfqpoint{1.329179in}{0.926257in}}%
\pgfpathlineto{\pgfqpoint{1.311614in}{0.917618in}}%
\pgfpathlineto{\pgfqpoint{1.294124in}{0.909358in}}%
\pgfpathlineto{\pgfqpoint{1.285983in}{0.908460in}}%
\pgfpathlineto{\pgfqpoint{1.277814in}{0.907913in}}%
\pgfpathclose%
\pgfusepath{fill}%
\end{pgfscope}%
\begin{pgfscope}%
\pgfpathrectangle{\pgfqpoint{0.407592in}{0.222904in}}{\pgfqpoint{1.411163in}{1.411163in}}%
\pgfusepath{clip}%
\pgfsetbuttcap%
\pgfsetroundjoin%
\definecolor{currentfill}{rgb}{0.909804,0.698039,0.423529}%
\pgfsetfillcolor{currentfill}%
\pgfsetfillopacity{0.120000}%
\pgfsetlinewidth{0.000000pt}%
\definecolor{currentstroke}{rgb}{0.000000,0.000000,0.000000}%
\pgfsetstrokecolor{currentstroke}%
\pgfsetdash{}{0pt}%
\pgfpathmoveto{\pgfqpoint{1.551158in}{1.040969in}}%
\pgfpathlineto{\pgfqpoint{1.569293in}{1.052467in}}%
\pgfpathlineto{\pgfqpoint{1.587509in}{1.064156in}}%
\pgfpathlineto{\pgfqpoint{1.595244in}{1.069268in}}%
\pgfpathlineto{\pgfqpoint{1.602961in}{1.074670in}}%
\pgfpathlineto{\pgfqpoint{1.584771in}{1.063005in}}%
\pgfpathlineto{\pgfqpoint{1.566662in}{1.051529in}}%
\pgfpathlineto{\pgfqpoint{1.558919in}{1.046096in}}%
\pgfpathlineto{\pgfqpoint{1.551158in}{1.040969in}}%
\pgfpathclose%
\pgfusepath{fill}%
\end{pgfscope}%
\begin{pgfscope}%
\pgfpathrectangle{\pgfqpoint{0.407592in}{0.222904in}}{\pgfqpoint{1.411163in}{1.411163in}}%
\pgfusepath{clip}%
\pgfsetbuttcap%
\pgfsetroundjoin%
\definecolor{currentfill}{rgb}{0.909804,0.698039,0.423529}%
\pgfsetfillcolor{currentfill}%
\pgfsetfillopacity{0.120000}%
\pgfsetlinewidth{0.000000pt}%
\definecolor{currentstroke}{rgb}{0.000000,0.000000,0.000000}%
\pgfsetstrokecolor{currentstroke}%
\pgfsetdash{}{0pt}%
\pgfpathmoveto{\pgfqpoint{1.329179in}{0.926257in}}%
\pgfpathlineto{\pgfqpoint{1.346818in}{0.935253in}}%
\pgfpathlineto{\pgfqpoint{1.364533in}{0.944582in}}%
\pgfpathlineto{\pgfqpoint{1.372587in}{0.946216in}}%
\pgfpathlineto{\pgfqpoint{1.380616in}{0.948187in}}%
\pgfpathlineto{\pgfqpoint{1.362930in}{0.938689in}}%
\pgfpathlineto{\pgfqpoint{1.345321in}{0.929517in}}%
\pgfpathlineto{\pgfqpoint{1.337263in}{0.927711in}}%
\pgfpathlineto{\pgfqpoint{1.329179in}{0.926257in}}%
\pgfpathclose%
\pgfusepath{fill}%
\end{pgfscope}%
\begin{pgfscope}%
\pgfpathrectangle{\pgfqpoint{0.407592in}{0.222904in}}{\pgfqpoint{1.411163in}{1.411163in}}%
\pgfusepath{clip}%
\pgfsetbuttcap%
\pgfsetroundjoin%
\definecolor{currentfill}{rgb}{0.909804,0.698039,0.423529}%
\pgfsetfillcolor{currentfill}%
\pgfsetfillopacity{0.120000}%
\pgfsetlinewidth{0.000000pt}%
\definecolor{currentstroke}{rgb}{0.000000,0.000000,0.000000}%
\pgfsetstrokecolor{currentstroke}%
\pgfsetdash{}{0pt}%
\pgfpathmoveto{\pgfqpoint{1.123847in}{0.877831in}}%
\pgfpathlineto{\pgfqpoint{1.141074in}{0.882125in}}%
\pgfpathlineto{\pgfqpoint{1.158363in}{0.886984in}}%
\pgfpathlineto{\pgfqpoint{1.166790in}{0.885054in}}%
\pgfpathlineto{\pgfqpoint{1.175183in}{0.883426in}}%
\pgfpathlineto{\pgfqpoint{1.157933in}{0.878288in}}%
\pgfpathlineto{\pgfqpoint{1.140747in}{0.873720in}}%
\pgfpathlineto{\pgfqpoint{1.132314in}{0.875624in}}%
\pgfpathlineto{\pgfqpoint{1.123847in}{0.877831in}}%
\pgfpathclose%
\pgfusepath{fill}%
\end{pgfscope}%
\begin{pgfscope}%
\pgfpathrectangle{\pgfqpoint{0.407592in}{0.222904in}}{\pgfqpoint{1.411163in}{1.411163in}}%
\pgfusepath{clip}%
\pgfsetbuttcap%
\pgfsetroundjoin%
\definecolor{currentfill}{rgb}{0.909804,0.698039,0.423529}%
\pgfsetfillcolor{currentfill}%
\pgfsetfillopacity{0.120000}%
\pgfsetlinewidth{0.000000pt}%
\definecolor{currentstroke}{rgb}{0.000000,0.000000,0.000000}%
\pgfsetstrokecolor{currentstroke}%
\pgfsetdash{}{0pt}%
\pgfpathmoveto{\pgfqpoint{1.380616in}{0.948187in}}%
\pgfpathlineto{\pgfqpoint{1.398379in}{0.957991in}}%
\pgfpathlineto{\pgfqpoint{1.416219in}{0.968080in}}%
\pgfpathlineto{\pgfqpoint{1.424196in}{0.970530in}}%
\pgfpathlineto{\pgfqpoint{1.432150in}{0.973311in}}%
\pgfpathlineto{\pgfqpoint{1.414337in}{0.963098in}}%
\pgfpathlineto{\pgfqpoint{1.396603in}{0.953164in}}%
\pgfpathlineto{\pgfqpoint{1.388621in}{0.950502in}}%
\pgfpathlineto{\pgfqpoint{1.380616in}{0.948187in}}%
\pgfpathclose%
\pgfusepath{fill}%
\end{pgfscope}%
\begin{pgfscope}%
\pgfpathrectangle{\pgfqpoint{0.407592in}{0.222904in}}{\pgfqpoint{1.411163in}{1.411163in}}%
\pgfusepath{clip}%
\pgfsetbuttcap%
\pgfsetroundjoin%
\definecolor{currentfill}{rgb}{0.909804,0.698039,0.423529}%
\pgfsetfillcolor{currentfill}%
\pgfsetfillopacity{0.120000}%
\pgfsetlinewidth{0.000000pt}%
\definecolor{currentstroke}{rgb}{0.000000,0.000000,0.000000}%
\pgfsetstrokecolor{currentstroke}%
\pgfsetdash{}{0pt}%
\pgfpathmoveto{\pgfqpoint{1.072445in}{0.876853in}}%
\pgfpathlineto{\pgfqpoint{1.089597in}{0.879673in}}%
\pgfpathlineto{\pgfqpoint{1.106806in}{0.883105in}}%
\pgfpathlineto{\pgfqpoint{1.115345in}{0.880329in}}%
\pgfpathlineto{\pgfqpoint{1.123847in}{0.877831in}}%
\pgfpathlineto{\pgfqpoint{1.106680in}{0.874132in}}%
\pgfpathlineto{\pgfqpoint{1.089572in}{0.871057in}}%
\pgfpathlineto{\pgfqpoint{1.081027in}{0.873817in}}%
\pgfpathlineto{\pgfqpoint{1.072445in}{0.876853in}}%
\pgfpathclose%
\pgfusepath{fill}%
\end{pgfscope}%
\begin{pgfscope}%
\pgfpathrectangle{\pgfqpoint{0.407592in}{0.222904in}}{\pgfqpoint{1.411163in}{1.411163in}}%
\pgfusepath{clip}%
\pgfsetbuttcap%
\pgfsetroundjoin%
\definecolor{currentfill}{rgb}{0.909804,0.698039,0.423529}%
\pgfsetfillcolor{currentfill}%
\pgfsetfillopacity{0.120000}%
\pgfsetlinewidth{0.000000pt}%
\definecolor{currentstroke}{rgb}{0.000000,0.000000,0.000000}%
\pgfsetstrokecolor{currentstroke}%
\pgfsetdash{}{0pt}%
\pgfpathmoveto{\pgfqpoint{1.602961in}{1.074670in}}%
\pgfpathlineto{\pgfqpoint{1.621235in}{1.086515in}}%
\pgfpathlineto{\pgfqpoint{1.639591in}{1.098532in}}%
\pgfpathlineto{\pgfqpoint{1.647265in}{1.104179in}}%
\pgfpathlineto{\pgfqpoint{1.654924in}{1.110099in}}%
\pgfpathlineto{\pgfqpoint{1.636594in}{1.098132in}}%
\pgfpathlineto{\pgfqpoint{1.618346in}{1.086338in}}%
\pgfpathlineto{\pgfqpoint{1.610662in}{1.080360in}}%
\pgfpathlineto{\pgfqpoint{1.602961in}{1.074670in}}%
\pgfpathclose%
\pgfusepath{fill}%
\end{pgfscope}%
\begin{pgfscope}%
\pgfpathrectangle{\pgfqpoint{0.407592in}{0.222904in}}{\pgfqpoint{1.411163in}{1.411163in}}%
\pgfusepath{clip}%
\pgfsetbuttcap%
\pgfsetroundjoin%
\definecolor{currentfill}{rgb}{0.909804,0.698039,0.423529}%
\pgfsetfillcolor{currentfill}%
\pgfsetfillopacity{0.120000}%
\pgfsetlinewidth{0.000000pt}%
\definecolor{currentstroke}{rgb}{0.000000,0.000000,0.000000}%
\pgfsetstrokecolor{currentstroke}%
\pgfsetdash{}{0pt}%
\pgfpathmoveto{\pgfqpoint{1.654924in}{1.110099in}}%
\pgfpathlineto{\pgfqpoint{1.673337in}{1.122231in}}%
\pgfpathlineto{\pgfqpoint{1.680967in}{1.128388in}}%
\pgfpathlineto{\pgfqpoint{1.688581in}{1.134807in}}%
\pgfpathlineto{\pgfqpoint{1.670193in}{1.122749in}}%
\pgfpathlineto{\pgfqpoint{1.662566in}{1.116290in}}%
\pgfusepath{fill}%
\end{pgfscope}%
\begin{pgfscope}%
\pgfpathrectangle{\pgfqpoint{0.407592in}{0.222904in}}{\pgfqpoint{1.411163in}{1.411163in}}%
\pgfusepath{clip}%
\pgfsetbuttcap%
\pgfsetroundjoin%
\definecolor{currentfill}{rgb}{0.909804,0.698039,0.423529}%
\pgfsetfillcolor{currentfill}%
\pgfsetfillopacity{0.120000}%
\pgfsetlinewidth{0.000000pt}%
\definecolor{currentstroke}{rgb}{0.000000,0.000000,0.000000}%
\pgfsetstrokecolor{currentstroke}%
\pgfsetdash{}{0pt}%
\pgfpathmoveto{\pgfqpoint{0.883573in}{0.906143in}}%
\pgfpathlineto{\pgfqpoint{0.900497in}{0.903072in}}%
\pgfpathlineto{\pgfqpoint{0.917459in}{0.900684in}}%
\pgfpathlineto{\pgfqpoint{0.926421in}{0.895748in}}%
\pgfpathlineto{\pgfqpoint{0.935341in}{0.891006in}}%
\pgfpathlineto{\pgfqpoint{0.918430in}{0.893312in}}%
\pgfpathlineto{\pgfqpoint{0.901557in}{0.896327in}}%
\pgfpathlineto{\pgfqpoint{0.892586in}{0.901139in}}%
\pgfpathlineto{\pgfqpoint{0.883573in}{0.906143in}}%
\pgfpathclose%
\pgfusepath{fill}%
\end{pgfscope}%
\begin{pgfscope}%
\pgfpathrectangle{\pgfqpoint{0.407592in}{0.222904in}}{\pgfqpoint{1.411163in}{1.411163in}}%
\pgfusepath{clip}%
\pgfsetbuttcap%
\pgfsetroundjoin%
\definecolor{currentfill}{rgb}{0.909804,0.698039,0.423529}%
\pgfsetfillcolor{currentfill}%
\pgfsetfillopacity{0.120000}%
\pgfsetlinewidth{0.000000pt}%
\definecolor{currentstroke}{rgb}{0.000000,0.000000,0.000000}%
\pgfsetstrokecolor{currentstroke}%
\pgfsetdash{}{0pt}%
\pgfpathmoveto{\pgfqpoint{1.432150in}{0.973311in}}%
\pgfpathlineto{\pgfqpoint{1.450041in}{0.983788in}}%
\pgfpathlineto{\pgfqpoint{1.468012in}{0.994511in}}%
\pgfpathlineto{\pgfqpoint{1.475917in}{0.997718in}}%
\pgfpathlineto{\pgfqpoint{1.483800in}{1.001245in}}%
\pgfpathlineto{\pgfqpoint{1.465856in}{0.990441in}}%
\pgfpathlineto{\pgfqpoint{1.447992in}{0.979879in}}%
\pgfpathlineto{\pgfqpoint{1.440082in}{0.976427in}}%
\pgfpathlineto{\pgfqpoint{1.432150in}{0.973311in}}%
\pgfpathclose%
\pgfusepath{fill}%
\end{pgfscope}%
\begin{pgfscope}%
\pgfpathrectangle{\pgfqpoint{0.407592in}{0.222904in}}{\pgfqpoint{1.411163in}{1.411163in}}%
\pgfusepath{clip}%
\pgfsetbuttcap%
\pgfsetroundjoin%
\definecolor{currentfill}{rgb}{0.909804,0.698039,0.423529}%
\pgfsetfillcolor{currentfill}%
\pgfsetfillopacity{0.120000}%
\pgfsetlinewidth{0.000000pt}%
\definecolor{currentstroke}{rgb}{0.000000,0.000000,0.000000}%
\pgfsetstrokecolor{currentstroke}%
\pgfsetdash{}{0pt}%
\pgfpathmoveto{\pgfqpoint{1.020938in}{0.880462in}}%
\pgfpathlineto{\pgfqpoint{1.038028in}{0.881758in}}%
\pgfpathlineto{\pgfqpoint{1.055170in}{0.883702in}}%
\pgfpathlineto{\pgfqpoint{1.063826in}{0.880152in}}%
\pgfpathlineto{\pgfqpoint{1.072445in}{0.876853in}}%
\pgfpathlineto{\pgfqpoint{1.055349in}{0.874673in}}%
\pgfpathlineto{\pgfqpoint{1.038306in}{0.873158in}}%
\pgfpathlineto{\pgfqpoint{1.029641in}{0.876687in}}%
\pgfpathlineto{\pgfqpoint{1.020938in}{0.880462in}}%
\pgfpathclose%
\pgfusepath{fill}%
\end{pgfscope}%
\begin{pgfscope}%
\pgfpathrectangle{\pgfqpoint{0.407592in}{0.222904in}}{\pgfqpoint{1.411163in}{1.411163in}}%
\pgfusepath{clip}%
\pgfsetbuttcap%
\pgfsetroundjoin%
\definecolor{currentfill}{rgb}{0.909804,0.698039,0.423529}%
\pgfsetfillcolor{currentfill}%
\pgfsetfillopacity{0.120000}%
\pgfsetlinewidth{0.000000pt}%
\definecolor{currentstroke}{rgb}{0.000000,0.000000,0.000000}%
\pgfsetstrokecolor{currentstroke}%
\pgfsetdash{}{0pt}%
\pgfpathmoveto{\pgfqpoint{1.483800in}{1.001245in}}%
\pgfpathlineto{\pgfqpoint{1.501825in}{1.012278in}}%
\pgfpathlineto{\pgfqpoint{1.519930in}{1.023527in}}%
\pgfpathlineto{\pgfqpoint{1.527766in}{1.027427in}}%
\pgfpathlineto{\pgfqpoint{1.535582in}{1.031634in}}%
\pgfpathlineto{\pgfqpoint{1.517503in}{1.020344in}}%
\pgfpathlineto{\pgfqpoint{1.499505in}{1.009267in}}%
\pgfpathlineto{\pgfqpoint{1.491663in}{1.005095in}}%
\pgfpathlineto{\pgfqpoint{1.483800in}{1.001245in}}%
\pgfpathclose%
\pgfusepath{fill}%
\end{pgfscope}%
\begin{pgfscope}%
\pgfpathrectangle{\pgfqpoint{0.407592in}{0.222904in}}{\pgfqpoint{1.411163in}{1.411163in}}%
\pgfusepath{clip}%
\pgfsetbuttcap%
\pgfsetroundjoin%
\definecolor{currentfill}{rgb}{0.909804,0.698039,0.423529}%
\pgfsetfillcolor{currentfill}%
\pgfsetfillopacity{0.120000}%
\pgfsetlinewidth{0.000000pt}%
\definecolor{currentstroke}{rgb}{0.000000,0.000000,0.000000}%
\pgfsetstrokecolor{currentstroke}%
\pgfsetdash{}{0pt}%
\pgfpathmoveto{\pgfqpoint{0.969288in}{0.888495in}}%
\pgfpathlineto{\pgfqpoint{0.986328in}{0.888266in}}%
\pgfpathlineto{\pgfqpoint{1.003415in}{0.888704in}}%
\pgfpathlineto{\pgfqpoint{1.012197in}{0.884471in}}%
\pgfpathlineto{\pgfqpoint{1.020938in}{0.880462in}}%
\pgfpathlineto{\pgfqpoint{1.003899in}{0.879834in}}%
\pgfpathlineto{\pgfqpoint{0.986908in}{0.879895in}}%
\pgfpathlineto{\pgfqpoint{0.978118in}{0.884085in}}%
\pgfpathlineto{\pgfqpoint{0.969288in}{0.888495in}}%
\pgfpathclose%
\pgfusepath{fill}%
\end{pgfscope}%
\begin{pgfscope}%
\pgfpathrectangle{\pgfqpoint{0.407592in}{0.222904in}}{\pgfqpoint{1.411163in}{1.411163in}}%
\pgfusepath{clip}%
\pgfsetbuttcap%
\pgfsetroundjoin%
\definecolor{currentfill}{rgb}{0.909804,0.698039,0.423529}%
\pgfsetfillcolor{currentfill}%
\pgfsetfillopacity{0.120000}%
\pgfsetlinewidth{0.000000pt}%
\definecolor{currentstroke}{rgb}{0.000000,0.000000,0.000000}%
\pgfsetstrokecolor{currentstroke}%
\pgfsetdash{}{0pt}%
\pgfpathmoveto{\pgfqpoint{1.535582in}{1.031634in}}%
\pgfpathlineto{\pgfqpoint{1.553743in}{1.043125in}}%
\pgfpathlineto{\pgfqpoint{1.571986in}{1.054807in}}%
\pgfpathlineto{\pgfqpoint{1.579757in}{1.059335in}}%
\pgfpathlineto{\pgfqpoint{1.587509in}{1.064156in}}%
\pgfpathlineto{\pgfqpoint{1.569293in}{1.052467in}}%
\pgfpathlineto{\pgfqpoint{1.551158in}{1.040969in}}%
\pgfpathlineto{\pgfqpoint{1.543380in}{1.036148in}}%
\pgfpathlineto{\pgfqpoint{1.535582in}{1.031634in}}%
\pgfpathclose%
\pgfusepath{fill}%
\end{pgfscope}%
\begin{pgfscope}%
\pgfpathrectangle{\pgfqpoint{0.407592in}{0.222904in}}{\pgfqpoint{1.411163in}{1.411163in}}%
\pgfusepath{clip}%
\pgfsetbuttcap%
\pgfsetroundjoin%
\definecolor{currentfill}{rgb}{0.909804,0.698039,0.423529}%
\pgfsetfillcolor{currentfill}%
\pgfsetfillopacity{0.120000}%
\pgfsetlinewidth{0.000000pt}%
\definecolor{currentstroke}{rgb}{0.000000,0.000000,0.000000}%
\pgfsetstrokecolor{currentstroke}%
\pgfsetdash{}{0pt}%
\pgfpathmoveto{\pgfqpoint{1.209879in}{0.895285in}}%
\pgfpathlineto{\pgfqpoint{1.227328in}{0.901948in}}%
\pgfpathlineto{\pgfqpoint{1.244847in}{0.909061in}}%
\pgfpathlineto{\pgfqpoint{1.253134in}{0.908293in}}%
\pgfpathlineto{\pgfqpoint{1.261390in}{0.907840in}}%
\pgfpathlineto{\pgfqpoint{1.243907in}{0.900461in}}%
\pgfpathlineto{\pgfqpoint{1.226493in}{0.893529in}}%
\pgfpathlineto{\pgfqpoint{1.218202in}{0.894247in}}%
\pgfpathlineto{\pgfqpoint{1.209879in}{0.895285in}}%
\pgfpathclose%
\pgfusepath{fill}%
\end{pgfscope}%
\begin{pgfscope}%
\pgfpathrectangle{\pgfqpoint{0.407592in}{0.222904in}}{\pgfqpoint{1.411163in}{1.411163in}}%
\pgfusepath{clip}%
\pgfsetbuttcap%
\pgfsetroundjoin%
\definecolor{currentfill}{rgb}{0.909804,0.698039,0.423529}%
\pgfsetfillcolor{currentfill}%
\pgfsetfillopacity{0.120000}%
\pgfsetlinewidth{0.000000pt}%
\definecolor{currentstroke}{rgb}{0.000000,0.000000,0.000000}%
\pgfsetstrokecolor{currentstroke}%
\pgfsetdash{}{0pt}%
\pgfpathmoveto{\pgfqpoint{1.261390in}{0.907840in}}%
\pgfpathlineto{\pgfqpoint{1.278945in}{0.915638in}}%
\pgfpathlineto{\pgfqpoint{1.296574in}{0.923830in}}%
\pgfpathlineto{\pgfqpoint{1.304767in}{0.923944in}}%
\pgfpathlineto{\pgfqpoint{1.312932in}{0.924381in}}%
\pgfpathlineto{\pgfqpoint{1.295336in}{0.915954in}}%
\pgfpathlineto{\pgfqpoint{1.277814in}{0.907913in}}%
\pgfpathlineto{\pgfqpoint{1.269617in}{0.907710in}}%
\pgfpathlineto{\pgfqpoint{1.261390in}{0.907840in}}%
\pgfpathclose%
\pgfusepath{fill}%
\end{pgfscope}%
\begin{pgfscope}%
\pgfpathrectangle{\pgfqpoint{0.407592in}{0.222904in}}{\pgfqpoint{1.411163in}{1.411163in}}%
\pgfusepath{clip}%
\pgfsetbuttcap%
\pgfsetroundjoin%
\definecolor{currentfill}{rgb}{0.909804,0.698039,0.423529}%
\pgfsetfillcolor{currentfill}%
\pgfsetfillopacity{0.120000}%
\pgfsetlinewidth{0.000000pt}%
\definecolor{currentstroke}{rgb}{0.000000,0.000000,0.000000}%
\pgfsetstrokecolor{currentstroke}%
\pgfsetdash{}{0pt}%
\pgfpathmoveto{\pgfqpoint{1.158363in}{0.886984in}}%
\pgfpathlineto{\pgfqpoint{1.175716in}{0.892379in}}%
\pgfpathlineto{\pgfqpoint{1.193136in}{0.898281in}}%
\pgfpathlineto{\pgfqpoint{1.201524in}{0.896634in}}%
\pgfpathlineto{\pgfqpoint{1.209879in}{0.895285in}}%
\pgfpathlineto{\pgfqpoint{1.192498in}{0.889102in}}%
\pgfpathlineto{\pgfqpoint{1.175183in}{0.883426in}}%
\pgfpathlineto{\pgfqpoint{1.166790in}{0.885054in}}%
\pgfpathlineto{\pgfqpoint{1.158363in}{0.886984in}}%
\pgfpathclose%
\pgfusepath{fill}%
\end{pgfscope}%
\begin{pgfscope}%
\pgfpathrectangle{\pgfqpoint{0.407592in}{0.222904in}}{\pgfqpoint{1.411163in}{1.411163in}}%
\pgfusepath{clip}%
\pgfsetbuttcap%
\pgfsetroundjoin%
\definecolor{currentfill}{rgb}{0.909804,0.698039,0.423529}%
\pgfsetfillcolor{currentfill}%
\pgfsetfillopacity{0.120000}%
\pgfsetlinewidth{0.000000pt}%
\definecolor{currentstroke}{rgb}{0.000000,0.000000,0.000000}%
\pgfsetstrokecolor{currentstroke}%
\pgfsetdash{}{0pt}%
\pgfpathmoveto{\pgfqpoint{1.312932in}{0.924381in}}%
\pgfpathlineto{\pgfqpoint{1.330601in}{0.933172in}}%
\pgfpathlineto{\pgfqpoint{1.348347in}{0.942304in}}%
\pgfpathlineto{\pgfqpoint{1.356453in}{0.943280in}}%
\pgfpathlineto{\pgfqpoint{1.364533in}{0.944582in}}%
\pgfpathlineto{\pgfqpoint{1.346818in}{0.935253in}}%
\pgfpathlineto{\pgfqpoint{1.329179in}{0.926257in}}%
\pgfpathlineto{\pgfqpoint{1.321069in}{0.925150in}}%
\pgfpathlineto{\pgfqpoint{1.312932in}{0.924381in}}%
\pgfpathclose%
\pgfusepath{fill}%
\end{pgfscope}%
\begin{pgfscope}%
\pgfpathrectangle{\pgfqpoint{0.407592in}{0.222904in}}{\pgfqpoint{1.411163in}{1.411163in}}%
\pgfusepath{clip}%
\pgfsetbuttcap%
\pgfsetroundjoin%
\definecolor{currentfill}{rgb}{0.909804,0.698039,0.423529}%
\pgfsetfillcolor{currentfill}%
\pgfsetfillopacity{0.120000}%
\pgfsetlinewidth{0.000000pt}%
\definecolor{currentstroke}{rgb}{0.000000,0.000000,0.000000}%
\pgfsetstrokecolor{currentstroke}%
\pgfsetdash{}{0pt}%
\pgfpathmoveto{\pgfqpoint{0.917459in}{0.900684in}}%
\pgfpathlineto{\pgfqpoint{0.934461in}{0.898971in}}%
\pgfpathlineto{\pgfqpoint{0.951506in}{0.897926in}}%
\pgfpathlineto{\pgfqpoint{0.960417in}{0.893112in}}%
\pgfpathlineto{\pgfqpoint{0.969288in}{0.888495in}}%
\pgfpathlineto{\pgfqpoint{0.952293in}{0.889404in}}%
\pgfpathlineto{\pgfqpoint{0.935341in}{0.891006in}}%
\pgfpathlineto{\pgfqpoint{0.926421in}{0.895748in}}%
\pgfpathlineto{\pgfqpoint{0.917459in}{0.900684in}}%
\pgfpathclose%
\pgfusepath{fill}%
\end{pgfscope}%
\begin{pgfscope}%
\pgfpathrectangle{\pgfqpoint{0.407592in}{0.222904in}}{\pgfqpoint{1.411163in}{1.411163in}}%
\pgfusepath{clip}%
\pgfsetbuttcap%
\pgfsetroundjoin%
\definecolor{currentfill}{rgb}{0.909804,0.698039,0.423529}%
\pgfsetfillcolor{currentfill}%
\pgfsetfillopacity{0.120000}%
\pgfsetlinewidth{0.000000pt}%
\definecolor{currentstroke}{rgb}{0.000000,0.000000,0.000000}%
\pgfsetstrokecolor{currentstroke}%
\pgfsetdash{}{0pt}%
\pgfpathmoveto{\pgfqpoint{1.587509in}{1.064156in}}%
\pgfpathlineto{\pgfqpoint{1.605808in}{1.076025in}}%
\pgfpathlineto{\pgfqpoint{1.624191in}{1.088066in}}%
\pgfpathlineto{\pgfqpoint{1.631900in}{1.093161in}}%
\pgfpathlineto{\pgfqpoint{1.639591in}{1.098532in}}%
\pgfpathlineto{\pgfqpoint{1.621235in}{1.086515in}}%
\pgfpathlineto{\pgfqpoint{1.602961in}{1.074670in}}%
\pgfpathlineto{\pgfqpoint{1.595244in}{1.069268in}}%
\pgfpathlineto{\pgfqpoint{1.587509in}{1.064156in}}%
\pgfpathclose%
\pgfusepath{fill}%
\end{pgfscope}%
\begin{pgfscope}%
\pgfpathrectangle{\pgfqpoint{0.407592in}{0.222904in}}{\pgfqpoint{1.411163in}{1.411163in}}%
\pgfusepath{clip}%
\pgfsetbuttcap%
\pgfsetroundjoin%
\definecolor{currentfill}{rgb}{0.909804,0.698039,0.423529}%
\pgfsetfillcolor{currentfill}%
\pgfsetfillopacity{0.120000}%
\pgfsetlinewidth{0.000000pt}%
\definecolor{currentstroke}{rgb}{0.000000,0.000000,0.000000}%
\pgfsetstrokecolor{currentstroke}%
\pgfsetdash{}{0pt}%
\pgfpathmoveto{\pgfqpoint{1.639591in}{1.098532in}}%
\pgfpathlineto{\pgfqpoint{1.658031in}{1.110714in}}%
\pgfpathlineto{\pgfqpoint{1.665692in}{1.116339in}}%
\pgfpathlineto{\pgfqpoint{1.673337in}{1.122231in}}%
\pgfpathlineto{\pgfqpoint{1.654924in}{1.110099in}}%
\pgfpathlineto{\pgfqpoint{1.647265in}{1.104179in}}%
\pgfusepath{fill}%
\end{pgfscope}%
\begin{pgfscope}%
\pgfpathrectangle{\pgfqpoint{0.407592in}{0.222904in}}{\pgfqpoint{1.411163in}{1.411163in}}%
\pgfusepath{clip}%
\pgfsetbuttcap%
\pgfsetroundjoin%
\definecolor{currentfill}{rgb}{0.909804,0.698039,0.423529}%
\pgfsetfillcolor{currentfill}%
\pgfsetfillopacity{0.120000}%
\pgfsetlinewidth{0.000000pt}%
\definecolor{currentstroke}{rgb}{0.000000,0.000000,0.000000}%
\pgfsetstrokecolor{currentstroke}%
\pgfsetdash{}{0pt}%
\pgfpathmoveto{\pgfqpoint{1.364533in}{0.944582in}}%
\pgfpathlineto{\pgfqpoint{1.382324in}{0.954223in}}%
\pgfpathlineto{\pgfqpoint{1.400194in}{0.964157in}}%
\pgfpathlineto{\pgfqpoint{1.408219in}{0.965957in}}%
\pgfpathlineto{\pgfqpoint{1.416219in}{0.968080in}}%
\pgfpathlineto{\pgfqpoint{1.398379in}{0.957991in}}%
\pgfpathlineto{\pgfqpoint{1.380616in}{0.948187in}}%
\pgfpathlineto{\pgfqpoint{1.372587in}{0.946216in}}%
\pgfpathlineto{\pgfqpoint{1.364533in}{0.944582in}}%
\pgfpathclose%
\pgfusepath{fill}%
\end{pgfscope}%
\begin{pgfscope}%
\pgfpathrectangle{\pgfqpoint{0.407592in}{0.222904in}}{\pgfqpoint{1.411163in}{1.411163in}}%
\pgfusepath{clip}%
\pgfsetbuttcap%
\pgfsetroundjoin%
\definecolor{currentfill}{rgb}{0.909804,0.698039,0.423529}%
\pgfsetfillcolor{currentfill}%
\pgfsetfillopacity{0.120000}%
\pgfsetlinewidth{0.000000pt}%
\definecolor{currentstroke}{rgb}{0.000000,0.000000,0.000000}%
\pgfsetstrokecolor{currentstroke}%
\pgfsetdash{}{0pt}%
\pgfpathmoveto{\pgfqpoint{1.106806in}{0.883105in}}%
\pgfpathlineto{\pgfqpoint{1.124075in}{0.887124in}}%
\pgfpathlineto{\pgfqpoint{1.141406in}{0.891701in}}%
\pgfpathlineto{\pgfqpoint{1.149902in}{0.889204in}}%
\pgfpathlineto{\pgfqpoint{1.158363in}{0.886984in}}%
\pgfpathlineto{\pgfqpoint{1.141074in}{0.882125in}}%
\pgfpathlineto{\pgfqpoint{1.123847in}{0.877831in}}%
\pgfpathlineto{\pgfqpoint{1.115345in}{0.880329in}}%
\pgfpathlineto{\pgfqpoint{1.106806in}{0.883105in}}%
\pgfpathclose%
\pgfusepath{fill}%
\end{pgfscope}%
\begin{pgfscope}%
\pgfpathrectangle{\pgfqpoint{0.407592in}{0.222904in}}{\pgfqpoint{1.411163in}{1.411163in}}%
\pgfusepath{clip}%
\pgfsetbuttcap%
\pgfsetroundjoin%
\definecolor{currentfill}{rgb}{0.909804,0.698039,0.423529}%
\pgfsetfillcolor{currentfill}%
\pgfsetfillopacity{0.120000}%
\pgfsetlinewidth{0.000000pt}%
\definecolor{currentstroke}{rgb}{0.000000,0.000000,0.000000}%
\pgfsetstrokecolor{currentstroke}%
\pgfsetdash{}{0pt}%
\pgfpathmoveto{\pgfqpoint{1.416219in}{0.968080in}}%
\pgfpathlineto{\pgfqpoint{1.434139in}{0.978439in}}%
\pgfpathlineto{\pgfqpoint{1.452138in}{0.989050in}}%
\pgfpathlineto{\pgfqpoint{1.460086in}{0.991623in}}%
\pgfpathlineto{\pgfqpoint{1.468012in}{0.994511in}}%
\pgfpathlineto{\pgfqpoint{1.450041in}{0.983788in}}%
\pgfpathlineto{\pgfqpoint{1.432150in}{0.973311in}}%
\pgfpathlineto{\pgfqpoint{1.424196in}{0.970530in}}%
\pgfpathlineto{\pgfqpoint{1.416219in}{0.968080in}}%
\pgfpathclose%
\pgfusepath{fill}%
\end{pgfscope}%
\begin{pgfscope}%
\pgfpathrectangle{\pgfqpoint{0.407592in}{0.222904in}}{\pgfqpoint{1.411163in}{1.411163in}}%
\pgfusepath{clip}%
\pgfsetbuttcap%
\pgfsetroundjoin%
\definecolor{currentfill}{rgb}{0.909804,0.698039,0.423529}%
\pgfsetfillcolor{currentfill}%
\pgfsetfillopacity{0.120000}%
\pgfsetlinewidth{0.000000pt}%
\definecolor{currentstroke}{rgb}{0.000000,0.000000,0.000000}%
\pgfsetstrokecolor{currentstroke}%
\pgfsetdash{}{0pt}%
\pgfpathmoveto{\pgfqpoint{1.055170in}{0.883702in}}%
\pgfpathlineto{\pgfqpoint{1.072367in}{0.886271in}}%
\pgfpathlineto{\pgfqpoint{1.089620in}{0.889440in}}%
\pgfpathlineto{\pgfqpoint{1.098232in}{0.886146in}}%
\pgfpathlineto{\pgfqpoint{1.106806in}{0.883105in}}%
\pgfpathlineto{\pgfqpoint{1.089597in}{0.879673in}}%
\pgfpathlineto{\pgfqpoint{1.072445in}{0.876853in}}%
\pgfpathlineto{\pgfqpoint{1.063826in}{0.880152in}}%
\pgfpathlineto{\pgfqpoint{1.055170in}{0.883702in}}%
\pgfpathclose%
\pgfusepath{fill}%
\end{pgfscope}%
\begin{pgfscope}%
\pgfpathrectangle{\pgfqpoint{0.407592in}{0.222904in}}{\pgfqpoint{1.411163in}{1.411163in}}%
\pgfusepath{clip}%
\pgfsetbuttcap%
\pgfsetroundjoin%
\definecolor{currentfill}{rgb}{0.909804,0.698039,0.423529}%
\pgfsetfillcolor{currentfill}%
\pgfsetfillopacity{0.120000}%
\pgfsetlinewidth{0.000000pt}%
\definecolor{currentstroke}{rgb}{0.000000,0.000000,0.000000}%
\pgfsetstrokecolor{currentstroke}%
\pgfsetdash{}{0pt}%
\pgfpathmoveto{\pgfqpoint{1.468012in}{0.994511in}}%
\pgfpathlineto{\pgfqpoint{1.486064in}{1.005468in}}%
\pgfpathlineto{\pgfqpoint{1.504197in}{1.016645in}}%
\pgfpathlineto{\pgfqpoint{1.512074in}{1.019934in}}%
\pgfpathlineto{\pgfqpoint{1.519930in}{1.023527in}}%
\pgfpathlineto{\pgfqpoint{1.501825in}{1.012278in}}%
\pgfpathlineto{\pgfqpoint{1.483800in}{1.001245in}}%
\pgfpathlineto{\pgfqpoint{1.475917in}{0.997718in}}%
\pgfpathlineto{\pgfqpoint{1.468012in}{0.994511in}}%
\pgfpathclose%
\pgfusepath{fill}%
\end{pgfscope}%
\begin{pgfscope}%
\pgfpathrectangle{\pgfqpoint{0.407592in}{0.222904in}}{\pgfqpoint{1.411163in}{1.411163in}}%
\pgfusepath{clip}%
\pgfsetbuttcap%
\pgfsetroundjoin%
\definecolor{currentfill}{rgb}{0.909804,0.698039,0.423529}%
\pgfsetfillcolor{currentfill}%
\pgfsetfillopacity{0.120000}%
\pgfsetlinewidth{0.000000pt}%
\definecolor{currentstroke}{rgb}{0.000000,0.000000,0.000000}%
\pgfsetstrokecolor{currentstroke}%
\pgfsetdash{}{0pt}%
\pgfpathmoveto{\pgfqpoint{1.003415in}{0.888704in}}%
\pgfpathlineto{\pgfqpoint{1.020552in}{0.889789in}}%
\pgfpathlineto{\pgfqpoint{1.037741in}{0.891505in}}%
\pgfpathlineto{\pgfqpoint{1.046475in}{0.887490in}}%
\pgfpathlineto{\pgfqpoint{1.055170in}{0.883702in}}%
\pgfpathlineto{\pgfqpoint{1.038028in}{0.881758in}}%
\pgfpathlineto{\pgfqpoint{1.020938in}{0.880462in}}%
\pgfpathlineto{\pgfqpoint{1.012197in}{0.884471in}}%
\pgfpathlineto{\pgfqpoint{1.003415in}{0.888704in}}%
\pgfpathclose%
\pgfusepath{fill}%
\end{pgfscope}%
\begin{pgfscope}%
\pgfpathrectangle{\pgfqpoint{0.407592in}{0.222904in}}{\pgfqpoint{1.411163in}{1.411163in}}%
\pgfusepath{clip}%
\pgfsetbuttcap%
\pgfsetroundjoin%
\definecolor{currentfill}{rgb}{0.909804,0.698039,0.423529}%
\pgfsetfillcolor{currentfill}%
\pgfsetfillopacity{0.120000}%
\pgfsetlinewidth{0.000000pt}%
\definecolor{currentstroke}{rgb}{0.000000,0.000000,0.000000}%
\pgfsetstrokecolor{currentstroke}%
\pgfsetdash{}{0pt}%
\pgfpathmoveto{\pgfqpoint{1.519930in}{1.023527in}}%
\pgfpathlineto{\pgfqpoint{1.538117in}{1.034980in}}%
\pgfpathlineto{\pgfqpoint{1.556387in}{1.046628in}}%
\pgfpathlineto{\pgfqpoint{1.564196in}{1.050571in}}%
\pgfpathlineto{\pgfqpoint{1.571986in}{1.054807in}}%
\pgfpathlineto{\pgfqpoint{1.553743in}{1.043125in}}%
\pgfpathlineto{\pgfqpoint{1.535582in}{1.031634in}}%
\pgfpathlineto{\pgfqpoint{1.527766in}{1.027427in}}%
\pgfpathlineto{\pgfqpoint{1.519930in}{1.023527in}}%
\pgfpathclose%
\pgfusepath{fill}%
\end{pgfscope}%
\begin{pgfscope}%
\pgfpathrectangle{\pgfqpoint{0.407592in}{0.222904in}}{\pgfqpoint{1.411163in}{1.411163in}}%
\pgfusepath{clip}%
\pgfsetbuttcap%
\pgfsetroundjoin%
\definecolor{currentfill}{rgb}{0.909804,0.698039,0.423529}%
\pgfsetfillcolor{currentfill}%
\pgfsetfillopacity{0.120000}%
\pgfsetlinewidth{0.000000pt}%
\definecolor{currentstroke}{rgb}{0.000000,0.000000,0.000000}%
\pgfsetstrokecolor{currentstroke}%
\pgfsetdash{}{0pt}%
\pgfpathmoveto{\pgfqpoint{0.951506in}{0.897926in}}%
\pgfpathlineto{\pgfqpoint{0.968596in}{0.897538in}}%
\pgfpathlineto{\pgfqpoint{0.985732in}{0.897793in}}%
\pgfpathlineto{\pgfqpoint{0.994594in}{0.893148in}}%
\pgfpathlineto{\pgfqpoint{1.003415in}{0.888704in}}%
\pgfpathlineto{\pgfqpoint{0.986328in}{0.888266in}}%
\pgfpathlineto{\pgfqpoint{0.969288in}{0.888495in}}%
\pgfpathlineto{\pgfqpoint{0.960417in}{0.893112in}}%
\pgfpathlineto{\pgfqpoint{0.951506in}{0.897926in}}%
\pgfpathclose%
\pgfusepath{fill}%
\end{pgfscope}%
\begin{pgfscope}%
\pgfpathrectangle{\pgfqpoint{0.407592in}{0.222904in}}{\pgfqpoint{1.411163in}{1.411163in}}%
\pgfusepath{clip}%
\pgfsetbuttcap%
\pgfsetroundjoin%
\definecolor{currentfill}{rgb}{0.909804,0.698039,0.423529}%
\pgfsetfillcolor{currentfill}%
\pgfsetfillopacity{0.120000}%
\pgfsetlinewidth{0.000000pt}%
\definecolor{currentstroke}{rgb}{0.000000,0.000000,0.000000}%
\pgfsetstrokecolor{currentstroke}%
\pgfsetdash{}{0pt}%
\pgfpathmoveto{\pgfqpoint{1.624191in}{1.088066in}}%
\pgfpathlineto{\pgfqpoint{1.642658in}{1.100272in}}%
\pgfpathlineto{\pgfqpoint{1.650353in}{1.105358in}}%
\pgfpathlineto{\pgfqpoint{1.658031in}{1.110714in}}%
\pgfpathlineto{\pgfqpoint{1.639591in}{1.098532in}}%
\pgfpathlineto{\pgfqpoint{1.631900in}{1.093161in}}%
\pgfusepath{fill}%
\end{pgfscope}%
\begin{pgfscope}%
\pgfpathrectangle{\pgfqpoint{0.407592in}{0.222904in}}{\pgfqpoint{1.411163in}{1.411163in}}%
\pgfusepath{clip}%
\pgfsetbuttcap%
\pgfsetroundjoin%
\definecolor{currentfill}{rgb}{0.909804,0.698039,0.423529}%
\pgfsetfillcolor{currentfill}%
\pgfsetfillopacity{0.120000}%
\pgfsetlinewidth{0.000000pt}%
\definecolor{currentstroke}{rgb}{0.000000,0.000000,0.000000}%
\pgfsetstrokecolor{currentstroke}%
\pgfsetdash{}{0pt}%
\pgfpathmoveto{\pgfqpoint{1.571986in}{1.054807in}}%
\pgfpathlineto{\pgfqpoint{1.590312in}{1.066672in}}%
\pgfpathlineto{\pgfqpoint{1.608721in}{1.078710in}}%
\pgfpathlineto{\pgfqpoint{1.616465in}{1.083249in}}%
\pgfpathlineto{\pgfqpoint{1.624191in}{1.088066in}}%
\pgfpathlineto{\pgfqpoint{1.605808in}{1.076025in}}%
\pgfpathlineto{\pgfqpoint{1.587509in}{1.064156in}}%
\pgfpathlineto{\pgfqpoint{1.579757in}{1.059335in}}%
\pgfpathlineto{\pgfqpoint{1.571986in}{1.054807in}}%
\pgfpathclose%
\pgfusepath{fill}%
\end{pgfscope}%
\begin{pgfscope}%
\pgfpathrectangle{\pgfqpoint{0.407592in}{0.222904in}}{\pgfqpoint{1.411163in}{1.411163in}}%
\pgfusepath{clip}%
\pgfsetbuttcap%
\pgfsetroundjoin%
\definecolor{currentfill}{rgb}{0.909804,0.698039,0.423529}%
\pgfsetfillcolor{currentfill}%
\pgfsetfillopacity{0.120000}%
\pgfsetlinewidth{0.000000pt}%
\definecolor{currentstroke}{rgb}{0.000000,0.000000,0.000000}%
\pgfsetstrokecolor{currentstroke}%
\pgfsetdash{}{0pt}%
\pgfpathmoveto{\pgfqpoint{1.244847in}{0.909061in}}%
\pgfpathlineto{\pgfqpoint{1.262437in}{0.916599in}}%
\pgfpathlineto{\pgfqpoint{1.280099in}{0.924536in}}%
\pgfpathlineto{\pgfqpoint{1.288351in}{0.924030in}}%
\pgfpathlineto{\pgfqpoint{1.296574in}{0.923830in}}%
\pgfpathlineto{\pgfqpoint{1.278945in}{0.915638in}}%
\pgfpathlineto{\pgfqpoint{1.261390in}{0.907840in}}%
\pgfpathlineto{\pgfqpoint{1.253134in}{0.908293in}}%
\pgfpathlineto{\pgfqpoint{1.244847in}{0.909061in}}%
\pgfpathclose%
\pgfusepath{fill}%
\end{pgfscope}%
\begin{pgfscope}%
\pgfpathrectangle{\pgfqpoint{0.407592in}{0.222904in}}{\pgfqpoint{1.411163in}{1.411163in}}%
\pgfusepath{clip}%
\pgfsetbuttcap%
\pgfsetroundjoin%
\definecolor{currentfill}{rgb}{0.909804,0.698039,0.423529}%
\pgfsetfillcolor{currentfill}%
\pgfsetfillopacity{0.120000}%
\pgfsetlinewidth{0.000000pt}%
\definecolor{currentstroke}{rgb}{0.000000,0.000000,0.000000}%
\pgfsetstrokecolor{currentstroke}%
\pgfsetdash{}{0pt}%
\pgfpathmoveto{\pgfqpoint{1.193136in}{0.898281in}}%
\pgfpathlineto{\pgfqpoint{1.210622in}{0.904663in}}%
\pgfpathlineto{\pgfqpoint{1.228178in}{0.911499in}}%
\pgfpathlineto{\pgfqpoint{1.236528in}{0.910133in}}%
\pgfpathlineto{\pgfqpoint{1.244847in}{0.909061in}}%
\pgfpathlineto{\pgfqpoint{1.227328in}{0.901948in}}%
\pgfpathlineto{\pgfqpoint{1.209879in}{0.895285in}}%
\pgfpathlineto{\pgfqpoint{1.201524in}{0.896634in}}%
\pgfpathlineto{\pgfqpoint{1.193136in}{0.898281in}}%
\pgfpathclose%
\pgfusepath{fill}%
\end{pgfscope}%
\begin{pgfscope}%
\pgfpathrectangle{\pgfqpoint{0.407592in}{0.222904in}}{\pgfqpoint{1.411163in}{1.411163in}}%
\pgfusepath{clip}%
\pgfsetbuttcap%
\pgfsetroundjoin%
\definecolor{currentfill}{rgb}{0.909804,0.698039,0.423529}%
\pgfsetfillcolor{currentfill}%
\pgfsetfillopacity{0.120000}%
\pgfsetlinewidth{0.000000pt}%
\definecolor{currentstroke}{rgb}{0.000000,0.000000,0.000000}%
\pgfsetstrokecolor{currentstroke}%
\pgfsetdash{}{0pt}%
\pgfpathmoveto{\pgfqpoint{1.296574in}{0.923830in}}%
\pgfpathlineto{\pgfqpoint{1.314276in}{0.932391in}}%
\pgfpathlineto{\pgfqpoint{1.332053in}{0.941301in}}%
\pgfpathlineto{\pgfqpoint{1.340214in}{0.941646in}}%
\pgfpathlineto{\pgfqpoint{1.348347in}{0.942304in}}%
\pgfpathlineto{\pgfqpoint{1.330601in}{0.933172in}}%
\pgfpathlineto{\pgfqpoint{1.312932in}{0.924381in}}%
\pgfpathlineto{\pgfqpoint{1.304767in}{0.923944in}}%
\pgfpathlineto{\pgfqpoint{1.296574in}{0.923830in}}%
\pgfpathclose%
\pgfusepath{fill}%
\end{pgfscope}%
\begin{pgfscope}%
\pgfpathrectangle{\pgfqpoint{0.407592in}{0.222904in}}{\pgfqpoint{1.411163in}{1.411163in}}%
\pgfusepath{clip}%
\pgfsetbuttcap%
\pgfsetroundjoin%
\definecolor{currentfill}{rgb}{0.909804,0.698039,0.423529}%
\pgfsetfillcolor{currentfill}%
\pgfsetfillopacity{0.120000}%
\pgfsetlinewidth{0.000000pt}%
\definecolor{currentstroke}{rgb}{0.000000,0.000000,0.000000}%
\pgfsetstrokecolor{currentstroke}%
\pgfsetdash{}{0pt}%
\pgfpathmoveto{\pgfqpoint{1.141406in}{0.891701in}}%
\pgfpathlineto{\pgfqpoint{1.158799in}{0.896811in}}%
\pgfpathlineto{\pgfqpoint{1.176259in}{0.902427in}}%
\pgfpathlineto{\pgfqpoint{1.184714in}{0.900215in}}%
\pgfpathlineto{\pgfqpoint{1.193136in}{0.898281in}}%
\pgfpathlineto{\pgfqpoint{1.175716in}{0.892379in}}%
\pgfpathlineto{\pgfqpoint{1.158363in}{0.886984in}}%
\pgfpathlineto{\pgfqpoint{1.149902in}{0.889204in}}%
\pgfpathlineto{\pgfqpoint{1.141406in}{0.891701in}}%
\pgfpathclose%
\pgfusepath{fill}%
\end{pgfscope}%
\begin{pgfscope}%
\pgfpathrectangle{\pgfqpoint{0.407592in}{0.222904in}}{\pgfqpoint{1.411163in}{1.411163in}}%
\pgfusepath{clip}%
\pgfsetbuttcap%
\pgfsetroundjoin%
\definecolor{currentfill}{rgb}{0.909804,0.698039,0.423529}%
\pgfsetfillcolor{currentfill}%
\pgfsetfillopacity{0.120000}%
\pgfsetlinewidth{0.000000pt}%
\definecolor{currentstroke}{rgb}{0.000000,0.000000,0.000000}%
\pgfsetstrokecolor{currentstroke}%
\pgfsetdash{}{0pt}%
\pgfpathmoveto{\pgfqpoint{1.348347in}{0.942304in}}%
\pgfpathlineto{\pgfqpoint{1.366169in}{0.951755in}}%
\pgfpathlineto{\pgfqpoint{1.384069in}{0.961507in}}%
\pgfpathlineto{\pgfqpoint{1.392144in}{0.962676in}}%
\pgfpathlineto{\pgfqpoint{1.400194in}{0.964157in}}%
\pgfpathlineto{\pgfqpoint{1.382324in}{0.954223in}}%
\pgfpathlineto{\pgfqpoint{1.364533in}{0.944582in}}%
\pgfpathlineto{\pgfqpoint{1.356453in}{0.943280in}}%
\pgfpathlineto{\pgfqpoint{1.348347in}{0.942304in}}%
\pgfpathclose%
\pgfusepath{fill}%
\end{pgfscope}%
\begin{pgfscope}%
\pgfpathrectangle{\pgfqpoint{0.407592in}{0.222904in}}{\pgfqpoint{1.411163in}{1.411163in}}%
\pgfusepath{clip}%
\pgfsetbuttcap%
\pgfsetroundjoin%
\definecolor{currentfill}{rgb}{0.909804,0.698039,0.423529}%
\pgfsetfillcolor{currentfill}%
\pgfsetfillopacity{0.120000}%
\pgfsetlinewidth{0.000000pt}%
\definecolor{currentstroke}{rgb}{0.000000,0.000000,0.000000}%
\pgfsetstrokecolor{currentstroke}%
\pgfsetdash{}{0pt}%
\pgfpathmoveto{\pgfqpoint{1.400194in}{0.964157in}}%
\pgfpathlineto{\pgfqpoint{1.418142in}{0.974368in}}%
\pgfpathlineto{\pgfqpoint{1.436170in}{0.984838in}}%
\pgfpathlineto{\pgfqpoint{1.444166in}{0.986789in}}%
\pgfpathlineto{\pgfqpoint{1.452138in}{0.989050in}}%
\pgfpathlineto{\pgfqpoint{1.434139in}{0.978439in}}%
\pgfpathlineto{\pgfqpoint{1.416219in}{0.968080in}}%
\pgfpathlineto{\pgfqpoint{1.408219in}{0.965957in}}%
\pgfpathlineto{\pgfqpoint{1.400194in}{0.964157in}}%
\pgfpathclose%
\pgfusepath{fill}%
\end{pgfscope}%
\begin{pgfscope}%
\pgfpathrectangle{\pgfqpoint{0.407592in}{0.222904in}}{\pgfqpoint{1.411163in}{1.411163in}}%
\pgfusepath{clip}%
\pgfsetbuttcap%
\pgfsetroundjoin%
\definecolor{currentfill}{rgb}{0.909804,0.698039,0.423529}%
\pgfsetfillcolor{currentfill}%
\pgfsetfillopacity{0.120000}%
\pgfsetlinewidth{0.000000pt}%
\definecolor{currentstroke}{rgb}{0.000000,0.000000,0.000000}%
\pgfsetstrokecolor{currentstroke}%
\pgfsetdash{}{0pt}%
\pgfpathmoveto{\pgfqpoint{1.089620in}{0.889440in}}%
\pgfpathlineto{\pgfqpoint{1.106933in}{0.893187in}}%
\pgfpathlineto{\pgfqpoint{1.124306in}{0.897485in}}%
\pgfpathlineto{\pgfqpoint{1.132874in}{0.894466in}}%
\pgfpathlineto{\pgfqpoint{1.141406in}{0.891701in}}%
\pgfpathlineto{\pgfqpoint{1.124075in}{0.887124in}}%
\pgfpathlineto{\pgfqpoint{1.106806in}{0.883105in}}%
\pgfpathlineto{\pgfqpoint{1.098232in}{0.886146in}}%
\pgfpathlineto{\pgfqpoint{1.089620in}{0.889440in}}%
\pgfpathclose%
\pgfusepath{fill}%
\end{pgfscope}%
\begin{pgfscope}%
\pgfpathrectangle{\pgfqpoint{0.407592in}{0.222904in}}{\pgfqpoint{1.411163in}{1.411163in}}%
\pgfusepath{clip}%
\pgfsetbuttcap%
\pgfsetroundjoin%
\definecolor{currentfill}{rgb}{0.909804,0.698039,0.423529}%
\pgfsetfillcolor{currentfill}%
\pgfsetfillopacity{0.120000}%
\pgfsetlinewidth{0.000000pt}%
\definecolor{currentstroke}{rgb}{0.000000,0.000000,0.000000}%
\pgfsetstrokecolor{currentstroke}%
\pgfsetdash{}{0pt}%
\pgfpathmoveto{\pgfqpoint{1.452138in}{0.989050in}}%
\pgfpathlineto{\pgfqpoint{1.470217in}{0.999900in}}%
\pgfpathlineto{\pgfqpoint{1.488378in}{1.010977in}}%
\pgfpathlineto{\pgfqpoint{1.496298in}{1.013660in}}%
\pgfpathlineto{\pgfqpoint{1.504197in}{1.016645in}}%
\pgfpathlineto{\pgfqpoint{1.486064in}{1.005468in}}%
\pgfpathlineto{\pgfqpoint{1.468012in}{0.994511in}}%
\pgfpathlineto{\pgfqpoint{1.460086in}{0.991623in}}%
\pgfpathlineto{\pgfqpoint{1.452138in}{0.989050in}}%
\pgfpathclose%
\pgfusepath{fill}%
\end{pgfscope}%
\begin{pgfscope}%
\pgfpathrectangle{\pgfqpoint{0.407592in}{0.222904in}}{\pgfqpoint{1.411163in}{1.411163in}}%
\pgfusepath{clip}%
\pgfsetbuttcap%
\pgfsetroundjoin%
\definecolor{currentfill}{rgb}{0.909804,0.698039,0.423529}%
\pgfsetfillcolor{currentfill}%
\pgfsetfillopacity{0.120000}%
\pgfsetlinewidth{0.000000pt}%
\definecolor{currentstroke}{rgb}{0.000000,0.000000,0.000000}%
\pgfsetstrokecolor{currentstroke}%
\pgfsetdash{}{0pt}%
\pgfpathmoveto{\pgfqpoint{1.037741in}{0.891505in}}%
\pgfpathlineto{\pgfqpoint{1.054984in}{0.893830in}}%
\pgfpathlineto{\pgfqpoint{1.072284in}{0.896744in}}%
\pgfpathlineto{\pgfqpoint{1.080971in}{0.892976in}}%
\pgfpathlineto{\pgfqpoint{1.089620in}{0.889440in}}%
\pgfpathlineto{\pgfqpoint{1.072367in}{0.886271in}}%
\pgfpathlineto{\pgfqpoint{1.055170in}{0.883702in}}%
\pgfpathlineto{\pgfqpoint{1.046475in}{0.887490in}}%
\pgfpathlineto{\pgfqpoint{1.037741in}{0.891505in}}%
\pgfpathclose%
\pgfusepath{fill}%
\end{pgfscope}%
\begin{pgfscope}%
\pgfpathrectangle{\pgfqpoint{0.407592in}{0.222904in}}{\pgfqpoint{1.411163in}{1.411163in}}%
\pgfusepath{clip}%
\pgfsetbuttcap%
\pgfsetroundjoin%
\definecolor{currentfill}{rgb}{0.909804,0.698039,0.423529}%
\pgfsetfillcolor{currentfill}%
\pgfsetfillopacity{0.120000}%
\pgfsetlinewidth{0.000000pt}%
\definecolor{currentstroke}{rgb}{0.000000,0.000000,0.000000}%
\pgfsetstrokecolor{currentstroke}%
\pgfsetdash{}{0pt}%
\pgfpathmoveto{\pgfqpoint{1.504197in}{1.016645in}}%
\pgfpathlineto{\pgfqpoint{1.522412in}{1.028031in}}%
\pgfpathlineto{\pgfqpoint{1.540709in}{1.039616in}}%
\pgfpathlineto{\pgfqpoint{1.548558in}{1.042976in}}%
\pgfpathlineto{\pgfqpoint{1.556387in}{1.046628in}}%
\pgfpathlineto{\pgfqpoint{1.538117in}{1.034980in}}%
\pgfpathlineto{\pgfqpoint{1.519930in}{1.023527in}}%
\pgfpathlineto{\pgfqpoint{1.512074in}{1.019934in}}%
\pgfpathlineto{\pgfqpoint{1.504197in}{1.016645in}}%
\pgfpathclose%
\pgfusepath{fill}%
\end{pgfscope}%
\begin{pgfscope}%
\pgfpathrectangle{\pgfqpoint{0.407592in}{0.222904in}}{\pgfqpoint{1.411163in}{1.411163in}}%
\pgfusepath{clip}%
\pgfsetbuttcap%
\pgfsetroundjoin%
\definecolor{currentfill}{rgb}{0.909804,0.698039,0.423529}%
\pgfsetfillcolor{currentfill}%
\pgfsetfillopacity{0.120000}%
\pgfsetlinewidth{0.000000pt}%
\definecolor{currentstroke}{rgb}{0.000000,0.000000,0.000000}%
\pgfsetstrokecolor{currentstroke}%
\pgfsetdash{}{0pt}%
\pgfpathmoveto{\pgfqpoint{0.985732in}{0.897793in}}%
\pgfpathlineto{\pgfqpoint{1.002918in}{0.898678in}}%
\pgfpathlineto{\pgfqpoint{1.020156in}{0.900174in}}%
\pgfpathlineto{\pgfqpoint{1.028968in}{0.895736in}}%
\pgfpathlineto{\pgfqpoint{1.037741in}{0.891505in}}%
\pgfpathlineto{\pgfqpoint{1.020552in}{0.889789in}}%
\pgfpathlineto{\pgfqpoint{1.003415in}{0.888704in}}%
\pgfpathlineto{\pgfqpoint{0.994594in}{0.893148in}}%
\pgfpathlineto{\pgfqpoint{0.985732in}{0.897793in}}%
\pgfpathclose%
\pgfusepath{fill}%
\end{pgfscope}%
\begin{pgfscope}%
\pgfpathrectangle{\pgfqpoint{0.407592in}{0.222904in}}{\pgfqpoint{1.411163in}{1.411163in}}%
\pgfusepath{clip}%
\pgfsetbuttcap%
\pgfsetroundjoin%
\definecolor{currentfill}{rgb}{0.909804,0.698039,0.423529}%
\pgfsetfillcolor{currentfill}%
\pgfsetfillopacity{0.120000}%
\pgfsetlinewidth{0.000000pt}%
\definecolor{currentstroke}{rgb}{0.000000,0.000000,0.000000}%
\pgfsetstrokecolor{currentstroke}%
\pgfsetdash{}{0pt}%
\pgfpathmoveto{\pgfqpoint{1.608721in}{1.078710in}}%
\pgfpathlineto{\pgfqpoint{1.627215in}{1.090915in}}%
\pgfpathlineto{\pgfqpoint{1.634946in}{1.095457in}}%
\pgfpathlineto{\pgfqpoint{1.642658in}{1.100272in}}%
\pgfpathlineto{\pgfqpoint{1.624191in}{1.088066in}}%
\pgfpathlineto{\pgfqpoint{1.616465in}{1.083249in}}%
\pgfusepath{fill}%
\end{pgfscope}%
\begin{pgfscope}%
\pgfpathrectangle{\pgfqpoint{0.407592in}{0.222904in}}{\pgfqpoint{1.411163in}{1.411163in}}%
\pgfusepath{clip}%
\pgfsetbuttcap%
\pgfsetroundjoin%
\definecolor{currentfill}{rgb}{0.909804,0.698039,0.423529}%
\pgfsetfillcolor{currentfill}%
\pgfsetfillopacity{0.120000}%
\pgfsetlinewidth{0.000000pt}%
\definecolor{currentstroke}{rgb}{0.000000,0.000000,0.000000}%
\pgfsetstrokecolor{currentstroke}%
\pgfsetdash{}{0pt}%
\pgfpathmoveto{\pgfqpoint{1.556387in}{1.046628in}}%
\pgfpathlineto{\pgfqpoint{1.574740in}{1.058460in}}%
\pgfpathlineto{\pgfqpoint{1.593177in}{1.070469in}}%
\pgfpathlineto{\pgfqpoint{1.600959in}{1.074450in}}%
\pgfpathlineto{\pgfqpoint{1.608721in}{1.078710in}}%
\pgfpathlineto{\pgfqpoint{1.590312in}{1.066672in}}%
\pgfpathlineto{\pgfqpoint{1.571986in}{1.054807in}}%
\pgfpathlineto{\pgfqpoint{1.564196in}{1.050571in}}%
\pgfpathlineto{\pgfqpoint{1.556387in}{1.046628in}}%
\pgfpathclose%
\pgfusepath{fill}%
\end{pgfscope}%
\begin{pgfscope}%
\pgfpathrectangle{\pgfqpoint{0.407592in}{0.222904in}}{\pgfqpoint{1.411163in}{1.411163in}}%
\pgfusepath{clip}%
\pgfsetbuttcap%
\pgfsetroundjoin%
\definecolor{currentfill}{rgb}{0.909804,0.698039,0.423529}%
\pgfsetfillcolor{currentfill}%
\pgfsetfillopacity{0.120000}%
\pgfsetlinewidth{0.000000pt}%
\definecolor{currentstroke}{rgb}{0.000000,0.000000,0.000000}%
\pgfsetstrokecolor{currentstroke}%
\pgfsetdash{}{0pt}%
\pgfpathmoveto{\pgfqpoint{1.228178in}{0.911499in}}%
\pgfpathlineto{\pgfqpoint{1.245805in}{0.918763in}}%
\pgfpathlineto{\pgfqpoint{1.263503in}{0.926433in}}%
\pgfpathlineto{\pgfqpoint{1.271817in}{0.925340in}}%
\pgfpathlineto{\pgfqpoint{1.280099in}{0.924536in}}%
\pgfpathlineto{\pgfqpoint{1.262437in}{0.916599in}}%
\pgfpathlineto{\pgfqpoint{1.244847in}{0.909061in}}%
\pgfpathlineto{\pgfqpoint{1.236528in}{0.910133in}}%
\pgfpathlineto{\pgfqpoint{1.228178in}{0.911499in}}%
\pgfpathclose%
\pgfusepath{fill}%
\end{pgfscope}%
\begin{pgfscope}%
\pgfpathrectangle{\pgfqpoint{0.407592in}{0.222904in}}{\pgfqpoint{1.411163in}{1.411163in}}%
\pgfusepath{clip}%
\pgfsetbuttcap%
\pgfsetroundjoin%
\definecolor{currentfill}{rgb}{0.909804,0.698039,0.423529}%
\pgfsetfillcolor{currentfill}%
\pgfsetfillopacity{0.120000}%
\pgfsetlinewidth{0.000000pt}%
\definecolor{currentstroke}{rgb}{0.000000,0.000000,0.000000}%
\pgfsetstrokecolor{currentstroke}%
\pgfsetdash{}{0pt}%
\pgfpathmoveto{\pgfqpoint{1.280099in}{0.924536in}}%
\pgfpathlineto{\pgfqpoint{1.297835in}{0.932850in}}%
\pgfpathlineto{\pgfqpoint{1.315646in}{0.941519in}}%
\pgfpathlineto{\pgfqpoint{1.323864in}{0.941262in}}%
\pgfpathlineto{\pgfqpoint{1.332053in}{0.941301in}}%
\pgfpathlineto{\pgfqpoint{1.314276in}{0.932391in}}%
\pgfpathlineto{\pgfqpoint{1.296574in}{0.923830in}}%
\pgfpathlineto{\pgfqpoint{1.288351in}{0.924030in}}%
\pgfpathlineto{\pgfqpoint{1.280099in}{0.924536in}}%
\pgfpathclose%
\pgfusepath{fill}%
\end{pgfscope}%
\begin{pgfscope}%
\pgfpathrectangle{\pgfqpoint{0.407592in}{0.222904in}}{\pgfqpoint{1.411163in}{1.411163in}}%
\pgfusepath{clip}%
\pgfsetbuttcap%
\pgfsetroundjoin%
\definecolor{currentfill}{rgb}{0.909804,0.698039,0.423529}%
\pgfsetfillcolor{currentfill}%
\pgfsetfillopacity{0.120000}%
\pgfsetlinewidth{0.000000pt}%
\definecolor{currentstroke}{rgb}{0.000000,0.000000,0.000000}%
\pgfsetstrokecolor{currentstroke}%
\pgfsetdash{}{0pt}%
\pgfpathmoveto{\pgfqpoint{1.176259in}{0.902427in}}%
\pgfpathlineto{\pgfqpoint{1.193784in}{0.908523in}}%
\pgfpathlineto{\pgfqpoint{1.211379in}{0.915074in}}%
\pgfpathlineto{\pgfqpoint{1.219795in}{0.913149in}}%
\pgfpathlineto{\pgfqpoint{1.228178in}{0.911499in}}%
\pgfpathlineto{\pgfqpoint{1.210622in}{0.904663in}}%
\pgfpathlineto{\pgfqpoint{1.193136in}{0.898281in}}%
\pgfpathlineto{\pgfqpoint{1.184714in}{0.900215in}}%
\pgfpathlineto{\pgfqpoint{1.176259in}{0.902427in}}%
\pgfpathclose%
\pgfusepath{fill}%
\end{pgfscope}%
\begin{pgfscope}%
\pgfpathrectangle{\pgfqpoint{0.407592in}{0.222904in}}{\pgfqpoint{1.411163in}{1.411163in}}%
\pgfusepath{clip}%
\pgfsetbuttcap%
\pgfsetroundjoin%
\definecolor{currentfill}{rgb}{0.909804,0.698039,0.423529}%
\pgfsetfillcolor{currentfill}%
\pgfsetfillopacity{0.120000}%
\pgfsetlinewidth{0.000000pt}%
\definecolor{currentstroke}{rgb}{0.000000,0.000000,0.000000}%
\pgfsetstrokecolor{currentstroke}%
\pgfsetdash{}{0pt}%
\pgfpathmoveto{\pgfqpoint{1.332053in}{0.941301in}}%
\pgfpathlineto{\pgfqpoint{1.349907in}{0.950539in}}%
\pgfpathlineto{\pgfqpoint{1.367839in}{0.960085in}}%
\pgfpathlineto{\pgfqpoint{1.375967in}{0.960645in}}%
\pgfpathlineto{\pgfqpoint{1.384069in}{0.961507in}}%
\pgfpathlineto{\pgfqpoint{1.366169in}{0.951755in}}%
\pgfpathlineto{\pgfqpoint{1.348347in}{0.942304in}}%
\pgfpathlineto{\pgfqpoint{1.340214in}{0.941646in}}%
\pgfpathlineto{\pgfqpoint{1.332053in}{0.941301in}}%
\pgfpathclose%
\pgfusepath{fill}%
\end{pgfscope}%
\begin{pgfscope}%
\pgfpathrectangle{\pgfqpoint{0.407592in}{0.222904in}}{\pgfqpoint{1.411163in}{1.411163in}}%
\pgfusepath{clip}%
\pgfsetbuttcap%
\pgfsetroundjoin%
\definecolor{currentfill}{rgb}{0.909804,0.698039,0.423529}%
\pgfsetfillcolor{currentfill}%
\pgfsetfillopacity{0.120000}%
\pgfsetlinewidth{0.000000pt}%
\definecolor{currentstroke}{rgb}{0.000000,0.000000,0.000000}%
\pgfsetstrokecolor{currentstroke}%
\pgfsetdash{}{0pt}%
\pgfpathmoveto{\pgfqpoint{1.124306in}{0.897485in}}%
\pgfpathlineto{\pgfqpoint{1.141742in}{0.902311in}}%
\pgfpathlineto{\pgfqpoint{1.159242in}{0.907639in}}%
\pgfpathlineto{\pgfqpoint{1.167768in}{0.904905in}}%
\pgfpathlineto{\pgfqpoint{1.176259in}{0.902427in}}%
\pgfpathlineto{\pgfqpoint{1.158799in}{0.896811in}}%
\pgfpathlineto{\pgfqpoint{1.141406in}{0.891701in}}%
\pgfpathlineto{\pgfqpoint{1.132874in}{0.894466in}}%
\pgfpathlineto{\pgfqpoint{1.124306in}{0.897485in}}%
\pgfpathclose%
\pgfusepath{fill}%
\end{pgfscope}%
\begin{pgfscope}%
\pgfpathrectangle{\pgfqpoint{0.407592in}{0.222904in}}{\pgfqpoint{1.411163in}{1.411163in}}%
\pgfusepath{clip}%
\pgfsetbuttcap%
\pgfsetroundjoin%
\definecolor{currentfill}{rgb}{0.909804,0.698039,0.423529}%
\pgfsetfillcolor{currentfill}%
\pgfsetfillopacity{0.120000}%
\pgfsetlinewidth{0.000000pt}%
\definecolor{currentstroke}{rgb}{0.000000,0.000000,0.000000}%
\pgfsetstrokecolor{currentstroke}%
\pgfsetdash{}{0pt}%
\pgfpathmoveto{\pgfqpoint{1.384069in}{0.961507in}}%
\pgfpathlineto{\pgfqpoint{1.402048in}{0.971542in}}%
\pgfpathlineto{\pgfqpoint{1.420106in}{0.981845in}}%
\pgfpathlineto{\pgfqpoint{1.428151in}{0.983191in}}%
\pgfpathlineto{\pgfqpoint{1.436170in}{0.984838in}}%
\pgfpathlineto{\pgfqpoint{1.418142in}{0.974368in}}%
\pgfpathlineto{\pgfqpoint{1.400194in}{0.964157in}}%
\pgfpathlineto{\pgfqpoint{1.392144in}{0.962676in}}%
\pgfpathlineto{\pgfqpoint{1.384069in}{0.961507in}}%
\pgfpathclose%
\pgfusepath{fill}%
\end{pgfscope}%
\begin{pgfscope}%
\pgfpathrectangle{\pgfqpoint{0.407592in}{0.222904in}}{\pgfqpoint{1.411163in}{1.411163in}}%
\pgfusepath{clip}%
\pgfsetbuttcap%
\pgfsetroundjoin%
\definecolor{currentfill}{rgb}{0.909804,0.698039,0.423529}%
\pgfsetfillcolor{currentfill}%
\pgfsetfillopacity{0.120000}%
\pgfsetlinewidth{0.000000pt}%
\definecolor{currentstroke}{rgb}{0.000000,0.000000,0.000000}%
\pgfsetstrokecolor{currentstroke}%
\pgfsetdash{}{0pt}%
\pgfpathmoveto{\pgfqpoint{1.436170in}{0.984838in}}%
\pgfpathlineto{\pgfqpoint{1.454279in}{0.995554in}}%
\pgfpathlineto{\pgfqpoint{1.472469in}{1.006502in}}%
\pgfpathlineto{\pgfqpoint{1.480435in}{1.008591in}}%
\pgfpathlineto{\pgfqpoint{1.488378in}{1.010977in}}%
\pgfpathlineto{\pgfqpoint{1.470217in}{0.999900in}}%
\pgfpathlineto{\pgfqpoint{1.452138in}{0.989050in}}%
\pgfpathlineto{\pgfqpoint{1.444166in}{0.986789in}}%
\pgfpathlineto{\pgfqpoint{1.436170in}{0.984838in}}%
\pgfpathclose%
\pgfusepath{fill}%
\end{pgfscope}%
\begin{pgfscope}%
\pgfpathrectangle{\pgfqpoint{0.407592in}{0.222904in}}{\pgfqpoint{1.411163in}{1.411163in}}%
\pgfusepath{clip}%
\pgfsetbuttcap%
\pgfsetroundjoin%
\definecolor{currentfill}{rgb}{0.909804,0.698039,0.423529}%
\pgfsetfillcolor{currentfill}%
\pgfsetfillopacity{0.120000}%
\pgfsetlinewidth{0.000000pt}%
\definecolor{currentstroke}{rgb}{0.000000,0.000000,0.000000}%
\pgfsetstrokecolor{currentstroke}%
\pgfsetdash{}{0pt}%
\pgfpathmoveto{\pgfqpoint{1.072284in}{0.896744in}}%
\pgfpathlineto{\pgfqpoint{1.089641in}{0.900223in}}%
\pgfpathlineto{\pgfqpoint{1.107059in}{0.904247in}}%
\pgfpathlineto{\pgfqpoint{1.115701in}{0.900749in}}%
\pgfpathlineto{\pgfqpoint{1.124306in}{0.897485in}}%
\pgfpathlineto{\pgfqpoint{1.106933in}{0.893187in}}%
\pgfpathlineto{\pgfqpoint{1.089620in}{0.889440in}}%
\pgfpathlineto{\pgfqpoint{1.080971in}{0.892976in}}%
\pgfpathlineto{\pgfqpoint{1.072284in}{0.896744in}}%
\pgfpathclose%
\pgfusepath{fill}%
\end{pgfscope}%
\begin{pgfscope}%
\pgfpathrectangle{\pgfqpoint{0.407592in}{0.222904in}}{\pgfqpoint{1.411163in}{1.411163in}}%
\pgfusepath{clip}%
\pgfsetbuttcap%
\pgfsetroundjoin%
\definecolor{currentfill}{rgb}{0.909804,0.698039,0.423529}%
\pgfsetfillcolor{currentfill}%
\pgfsetfillopacity{0.120000}%
\pgfsetlinewidth{0.000000pt}%
\definecolor{currentstroke}{rgb}{0.000000,0.000000,0.000000}%
\pgfsetstrokecolor{currentstroke}%
\pgfsetdash{}{0pt}%
\pgfpathmoveto{\pgfqpoint{1.488378in}{1.010977in}}%
\pgfpathlineto{\pgfqpoint{1.506621in}{1.022267in}}%
\pgfpathlineto{\pgfqpoint{1.524947in}{1.033761in}}%
\pgfpathlineto{\pgfqpoint{1.532839in}{1.036544in}}%
\pgfpathlineto{\pgfqpoint{1.540709in}{1.039616in}}%
\pgfpathlineto{\pgfqpoint{1.522412in}{1.028031in}}%
\pgfpathlineto{\pgfqpoint{1.504197in}{1.016645in}}%
\pgfpathlineto{\pgfqpoint{1.496298in}{1.013660in}}%
\pgfpathlineto{\pgfqpoint{1.488378in}{1.010977in}}%
\pgfpathclose%
\pgfusepath{fill}%
\end{pgfscope}%
\begin{pgfscope}%
\pgfpathrectangle{\pgfqpoint{0.407592in}{0.222904in}}{\pgfqpoint{1.411163in}{1.411163in}}%
\pgfusepath{clip}%
\pgfsetbuttcap%
\pgfsetroundjoin%
\definecolor{currentfill}{rgb}{0.909804,0.698039,0.423529}%
\pgfsetfillcolor{currentfill}%
\pgfsetfillopacity{0.120000}%
\pgfsetlinewidth{0.000000pt}%
\definecolor{currentstroke}{rgb}{0.000000,0.000000,0.000000}%
\pgfsetstrokecolor{currentstroke}%
\pgfsetdash{}{0pt}%
\pgfpathmoveto{\pgfqpoint{1.020156in}{0.900174in}}%
\pgfpathlineto{\pgfqpoint{1.037447in}{0.902265in}}%
\pgfpathlineto{\pgfqpoint{1.054793in}{0.904931in}}%
\pgfpathlineto{\pgfqpoint{1.063558in}{0.900732in}}%
\pgfpathlineto{\pgfqpoint{1.072284in}{0.896744in}}%
\pgfpathlineto{\pgfqpoint{1.054984in}{0.893830in}}%
\pgfpathlineto{\pgfqpoint{1.037741in}{0.891505in}}%
\pgfpathlineto{\pgfqpoint{1.028968in}{0.895736in}}%
\pgfpathlineto{\pgfqpoint{1.020156in}{0.900174in}}%
\pgfpathclose%
\pgfusepath{fill}%
\end{pgfscope}%
\begin{pgfscope}%
\pgfpathrectangle{\pgfqpoint{0.407592in}{0.222904in}}{\pgfqpoint{1.411163in}{1.411163in}}%
\pgfusepath{clip}%
\pgfsetbuttcap%
\pgfsetroundjoin%
\definecolor{currentfill}{rgb}{0.909804,0.698039,0.423529}%
\pgfsetfillcolor{currentfill}%
\pgfsetfillopacity{0.120000}%
\pgfsetlinewidth{0.000000pt}%
\definecolor{currentstroke}{rgb}{0.000000,0.000000,0.000000}%
\pgfsetstrokecolor{currentstroke}%
\pgfsetdash{}{0pt}%
\pgfpathmoveto{\pgfqpoint{1.593177in}{1.070469in}}%
\pgfpathlineto{\pgfqpoint{1.611699in}{1.082647in}}%
\pgfpathlineto{\pgfqpoint{1.619467in}{1.086644in}}%
\pgfpathlineto{\pgfqpoint{1.627215in}{1.090915in}}%
\pgfpathlineto{\pgfqpoint{1.608721in}{1.078710in}}%
\pgfpathlineto{\pgfqpoint{1.600959in}{1.074450in}}%
\pgfusepath{fill}%
\end{pgfscope}%
\begin{pgfscope}%
\pgfpathrectangle{\pgfqpoint{0.407592in}{0.222904in}}{\pgfqpoint{1.411163in}{1.411163in}}%
\pgfusepath{clip}%
\pgfsetbuttcap%
\pgfsetroundjoin%
\definecolor{currentfill}{rgb}{0.909804,0.698039,0.423529}%
\pgfsetfillcolor{currentfill}%
\pgfsetfillopacity{0.120000}%
\pgfsetlinewidth{0.000000pt}%
\definecolor{currentstroke}{rgb}{0.000000,0.000000,0.000000}%
\pgfsetstrokecolor{currentstroke}%
\pgfsetdash{}{0pt}%
\pgfpathmoveto{\pgfqpoint{1.540709in}{1.039616in}}%
\pgfpathlineto{\pgfqpoint{1.559090in}{1.051389in}}%
\pgfpathlineto{\pgfqpoint{1.577555in}{1.063342in}}%
\pgfpathlineto{\pgfqpoint{1.585376in}{1.066767in}}%
\pgfpathlineto{\pgfqpoint{1.593177in}{1.070469in}}%
\pgfpathlineto{\pgfqpoint{1.574740in}{1.058460in}}%
\pgfpathlineto{\pgfqpoint{1.556387in}{1.046628in}}%
\pgfpathlineto{\pgfqpoint{1.548558in}{1.042976in}}%
\pgfpathlineto{\pgfqpoint{1.540709in}{1.039616in}}%
\pgfpathclose%
\pgfusepath{fill}%
\end{pgfscope}%
\begin{pgfscope}%
\pgfpathrectangle{\pgfqpoint{0.407592in}{0.222904in}}{\pgfqpoint{1.411163in}{1.411163in}}%
\pgfusepath{clip}%
\pgfsetbuttcap%
\pgfsetroundjoin%
\definecolor{currentfill}{rgb}{0.909804,0.698039,0.423529}%
\pgfsetfillcolor{currentfill}%
\pgfsetfillopacity{0.120000}%
\pgfsetlinewidth{0.000000pt}%
\definecolor{currentstroke}{rgb}{0.000000,0.000000,0.000000}%
\pgfsetstrokecolor{currentstroke}%
\pgfsetdash{}{0pt}%
\pgfpathmoveto{\pgfqpoint{1.263503in}{0.926433in}}%
\pgfpathlineto{\pgfqpoint{1.281275in}{0.934485in}}%
\pgfpathlineto{\pgfqpoint{1.299121in}{0.942899in}}%
\pgfpathlineto{\pgfqpoint{1.307399in}{0.942068in}}%
\pgfpathlineto{\pgfqpoint{1.315646in}{0.941519in}}%
\pgfpathlineto{\pgfqpoint{1.297835in}{0.932850in}}%
\pgfpathlineto{\pgfqpoint{1.280099in}{0.924536in}}%
\pgfpathlineto{\pgfqpoint{1.271817in}{0.925340in}}%
\pgfpathlineto{\pgfqpoint{1.263503in}{0.926433in}}%
\pgfpathclose%
\pgfusepath{fill}%
\end{pgfscope}%
\begin{pgfscope}%
\pgfpathrectangle{\pgfqpoint{0.407592in}{0.222904in}}{\pgfqpoint{1.411163in}{1.411163in}}%
\pgfusepath{clip}%
\pgfsetbuttcap%
\pgfsetroundjoin%
\definecolor{currentfill}{rgb}{0.909804,0.698039,0.423529}%
\pgfsetfillcolor{currentfill}%
\pgfsetfillopacity{0.120000}%
\pgfsetlinewidth{0.000000pt}%
\definecolor{currentstroke}{rgb}{0.000000,0.000000,0.000000}%
\pgfsetstrokecolor{currentstroke}%
\pgfsetdash{}{0pt}%
\pgfpathmoveto{\pgfqpoint{1.211379in}{0.915074in}}%
\pgfpathlineto{\pgfqpoint{1.229043in}{0.922057in}}%
\pgfpathlineto{\pgfqpoint{1.246779in}{0.929450in}}%
\pgfpathlineto{\pgfqpoint{1.255157in}{0.927805in}}%
\pgfpathlineto{\pgfqpoint{1.263503in}{0.926433in}}%
\pgfpathlineto{\pgfqpoint{1.245805in}{0.918763in}}%
\pgfpathlineto{\pgfqpoint{1.228178in}{0.911499in}}%
\pgfpathlineto{\pgfqpoint{1.219795in}{0.913149in}}%
\pgfpathlineto{\pgfqpoint{1.211379in}{0.915074in}}%
\pgfpathclose%
\pgfusepath{fill}%
\end{pgfscope}%
\begin{pgfscope}%
\pgfpathrectangle{\pgfqpoint{0.407592in}{0.222904in}}{\pgfqpoint{1.411163in}{1.411163in}}%
\pgfusepath{clip}%
\pgfsetbuttcap%
\pgfsetroundjoin%
\definecolor{currentfill}{rgb}{0.909804,0.698039,0.423529}%
\pgfsetfillcolor{currentfill}%
\pgfsetfillopacity{0.120000}%
\pgfsetlinewidth{0.000000pt}%
\definecolor{currentstroke}{rgb}{0.000000,0.000000,0.000000}%
\pgfsetstrokecolor{currentstroke}%
\pgfsetdash{}{0pt}%
\pgfpathmoveto{\pgfqpoint{1.315646in}{0.941519in}}%
\pgfpathlineto{\pgfqpoint{1.333534in}{0.950524in}}%
\pgfpathlineto{\pgfqpoint{1.351498in}{0.959845in}}%
\pgfpathlineto{\pgfqpoint{1.359683in}{0.959820in}}%
\pgfpathlineto{\pgfqpoint{1.367839in}{0.960085in}}%
\pgfpathlineto{\pgfqpoint{1.349907in}{0.950539in}}%
\pgfpathlineto{\pgfqpoint{1.332053in}{0.941301in}}%
\pgfpathlineto{\pgfqpoint{1.323864in}{0.941262in}}%
\pgfpathlineto{\pgfqpoint{1.315646in}{0.941519in}}%
\pgfpathclose%
\pgfusepath{fill}%
\end{pgfscope}%
\begin{pgfscope}%
\pgfpathrectangle{\pgfqpoint{0.407592in}{0.222904in}}{\pgfqpoint{1.411163in}{1.411163in}}%
\pgfusepath{clip}%
\pgfsetbuttcap%
\pgfsetroundjoin%
\definecolor{currentfill}{rgb}{0.909804,0.698039,0.423529}%
\pgfsetfillcolor{currentfill}%
\pgfsetfillopacity{0.120000}%
\pgfsetlinewidth{0.000000pt}%
\definecolor{currentstroke}{rgb}{0.000000,0.000000,0.000000}%
\pgfsetstrokecolor{currentstroke}%
\pgfsetdash{}{0pt}%
\pgfpathmoveto{\pgfqpoint{1.159242in}{0.907639in}}%
\pgfpathlineto{\pgfqpoint{1.176809in}{0.913447in}}%
\pgfpathlineto{\pgfqpoint{1.194444in}{0.919711in}}%
\pgfpathlineto{\pgfqpoint{1.202929in}{0.917265in}}%
\pgfpathlineto{\pgfqpoint{1.211379in}{0.915074in}}%
\pgfpathlineto{\pgfqpoint{1.193784in}{0.908523in}}%
\pgfpathlineto{\pgfqpoint{1.176259in}{0.902427in}}%
\pgfpathlineto{\pgfqpoint{1.167768in}{0.904905in}}%
\pgfpathlineto{\pgfqpoint{1.159242in}{0.907639in}}%
\pgfpathclose%
\pgfusepath{fill}%
\end{pgfscope}%
\begin{pgfscope}%
\pgfpathrectangle{\pgfqpoint{0.407592in}{0.222904in}}{\pgfqpoint{1.411163in}{1.411163in}}%
\pgfusepath{clip}%
\pgfsetbuttcap%
\pgfsetroundjoin%
\definecolor{currentfill}{rgb}{0.909804,0.698039,0.423529}%
\pgfsetfillcolor{currentfill}%
\pgfsetfillopacity{0.120000}%
\pgfsetlinewidth{0.000000pt}%
\definecolor{currentstroke}{rgb}{0.000000,0.000000,0.000000}%
\pgfsetstrokecolor{currentstroke}%
\pgfsetdash{}{0pt}%
\pgfpathmoveto{\pgfqpoint{1.367839in}{0.960085in}}%
\pgfpathlineto{\pgfqpoint{1.385849in}{0.969922in}}%
\pgfpathlineto{\pgfqpoint{1.403939in}{0.980035in}}%
\pgfpathlineto{\pgfqpoint{1.412035in}{0.980795in}}%
\pgfpathlineto{\pgfqpoint{1.420106in}{0.981845in}}%
\pgfpathlineto{\pgfqpoint{1.402048in}{0.971542in}}%
\pgfpathlineto{\pgfqpoint{1.384069in}{0.961507in}}%
\pgfpathlineto{\pgfqpoint{1.375967in}{0.960645in}}%
\pgfpathlineto{\pgfqpoint{1.367839in}{0.960085in}}%
\pgfpathclose%
\pgfusepath{fill}%
\end{pgfscope}%
\begin{pgfscope}%
\pgfpathrectangle{\pgfqpoint{0.407592in}{0.222904in}}{\pgfqpoint{1.411163in}{1.411163in}}%
\pgfusepath{clip}%
\pgfsetbuttcap%
\pgfsetroundjoin%
\definecolor{currentfill}{rgb}{0.909804,0.698039,0.423529}%
\pgfsetfillcolor{currentfill}%
\pgfsetfillopacity{0.120000}%
\pgfsetlinewidth{0.000000pt}%
\definecolor{currentstroke}{rgb}{0.000000,0.000000,0.000000}%
\pgfsetstrokecolor{currentstroke}%
\pgfsetdash{}{0pt}%
\pgfpathmoveto{\pgfqpoint{1.107059in}{0.904247in}}%
\pgfpathlineto{\pgfqpoint{1.124539in}{0.908792in}}%
\pgfpathlineto{\pgfqpoint{1.142083in}{0.913837in}}%
\pgfpathlineto{\pgfqpoint{1.150681in}{0.910620in}}%
\pgfpathlineto{\pgfqpoint{1.159242in}{0.907639in}}%
\pgfpathlineto{\pgfqpoint{1.141742in}{0.902311in}}%
\pgfpathlineto{\pgfqpoint{1.124306in}{0.897485in}}%
\pgfpathlineto{\pgfqpoint{1.115701in}{0.900749in}}%
\pgfpathlineto{\pgfqpoint{1.107059in}{0.904247in}}%
\pgfpathclose%
\pgfusepath{fill}%
\end{pgfscope}%
\begin{pgfscope}%
\pgfpathrectangle{\pgfqpoint{0.407592in}{0.222904in}}{\pgfqpoint{1.411163in}{1.411163in}}%
\pgfusepath{clip}%
\pgfsetbuttcap%
\pgfsetroundjoin%
\definecolor{currentfill}{rgb}{0.909804,0.698039,0.423529}%
\pgfsetfillcolor{currentfill}%
\pgfsetfillopacity{0.120000}%
\pgfsetlinewidth{0.000000pt}%
\definecolor{currentstroke}{rgb}{0.000000,0.000000,0.000000}%
\pgfsetstrokecolor{currentstroke}%
\pgfsetdash{}{0pt}%
\pgfpathmoveto{\pgfqpoint{1.420106in}{0.981845in}}%
\pgfpathlineto{\pgfqpoint{1.438244in}{0.992401in}}%
\pgfpathlineto{\pgfqpoint{1.456465in}{1.003196in}}%
\pgfpathlineto{\pgfqpoint{1.464479in}{1.004705in}}%
\pgfpathlineto{\pgfqpoint{1.472469in}{1.006502in}}%
\pgfpathlineto{\pgfqpoint{1.454279in}{0.995554in}}%
\pgfpathlineto{\pgfqpoint{1.436170in}{0.984838in}}%
\pgfpathlineto{\pgfqpoint{1.428151in}{0.983191in}}%
\pgfpathlineto{\pgfqpoint{1.420106in}{0.981845in}}%
\pgfpathclose%
\pgfusepath{fill}%
\end{pgfscope}%
\begin{pgfscope}%
\pgfpathrectangle{\pgfqpoint{0.407592in}{0.222904in}}{\pgfqpoint{1.411163in}{1.411163in}}%
\pgfusepath{clip}%
\pgfsetbuttcap%
\pgfsetroundjoin%
\definecolor{currentfill}{rgb}{0.909804,0.698039,0.423529}%
\pgfsetfillcolor{currentfill}%
\pgfsetfillopacity{0.120000}%
\pgfsetlinewidth{0.000000pt}%
\definecolor{currentstroke}{rgb}{0.000000,0.000000,0.000000}%
\pgfsetstrokecolor{currentstroke}%
\pgfsetdash{}{0pt}%
\pgfpathmoveto{\pgfqpoint{1.472469in}{1.006502in}}%
\pgfpathlineto{\pgfqpoint{1.490741in}{1.017671in}}%
\pgfpathlineto{\pgfqpoint{1.509096in}{1.029049in}}%
\pgfpathlineto{\pgfqpoint{1.517033in}{1.031263in}}%
\pgfpathlineto{\pgfqpoint{1.524947in}{1.033761in}}%
\pgfpathlineto{\pgfqpoint{1.506621in}{1.022267in}}%
\pgfpathlineto{\pgfqpoint{1.488378in}{1.010977in}}%
\pgfpathlineto{\pgfqpoint{1.480435in}{1.008591in}}%
\pgfpathlineto{\pgfqpoint{1.472469in}{1.006502in}}%
\pgfpathclose%
\pgfusepath{fill}%
\end{pgfscope}%
\begin{pgfscope}%
\pgfpathrectangle{\pgfqpoint{0.407592in}{0.222904in}}{\pgfqpoint{1.411163in}{1.411163in}}%
\pgfusepath{clip}%
\pgfsetbuttcap%
\pgfsetroundjoin%
\definecolor{currentfill}{rgb}{0.909804,0.698039,0.423529}%
\pgfsetfillcolor{currentfill}%
\pgfsetfillopacity{0.120000}%
\pgfsetlinewidth{0.000000pt}%
\definecolor{currentstroke}{rgb}{0.000000,0.000000,0.000000}%
\pgfsetstrokecolor{currentstroke}%
\pgfsetdash{}{0pt}%
\pgfpathmoveto{\pgfqpoint{1.054793in}{0.904931in}}%
\pgfpathlineto{\pgfqpoint{1.072197in}{0.908152in}}%
\pgfpathlineto{\pgfqpoint{1.089661in}{0.911908in}}%
\pgfpathlineto{\pgfqpoint{1.098379in}{0.907970in}}%
\pgfpathlineto{\pgfqpoint{1.107059in}{0.904247in}}%
\pgfpathlineto{\pgfqpoint{1.089641in}{0.900223in}}%
\pgfpathlineto{\pgfqpoint{1.072284in}{0.896744in}}%
\pgfpathlineto{\pgfqpoint{1.063558in}{0.900732in}}%
\pgfpathlineto{\pgfqpoint{1.054793in}{0.904931in}}%
\pgfpathclose%
\pgfusepath{fill}%
\end{pgfscope}%
\begin{pgfscope}%
\pgfpathrectangle{\pgfqpoint{0.407592in}{0.222904in}}{\pgfqpoint{1.411163in}{1.411163in}}%
\pgfusepath{clip}%
\pgfsetbuttcap%
\pgfsetroundjoin%
\definecolor{currentfill}{rgb}{0.909804,0.698039,0.423529}%
\pgfsetfillcolor{currentfill}%
\pgfsetfillopacity{0.120000}%
\pgfsetlinewidth{0.000000pt}%
\definecolor{currentstroke}{rgb}{0.000000,0.000000,0.000000}%
\pgfsetstrokecolor{currentstroke}%
\pgfsetdash{}{0pt}%
\pgfpathmoveto{\pgfqpoint{1.577555in}{1.063342in}}%
\pgfpathlineto{\pgfqpoint{1.596105in}{1.075467in}}%
\pgfpathlineto{\pgfqpoint{1.603912in}{1.078921in}}%
\pgfpathlineto{\pgfqpoint{1.611699in}{1.082647in}}%
\pgfpathlineto{\pgfqpoint{1.593177in}{1.070469in}}%
\pgfpathlineto{\pgfqpoint{1.585376in}{1.066767in}}%
\pgfusepath{fill}%
\end{pgfscope}%
\begin{pgfscope}%
\pgfpathrectangle{\pgfqpoint{0.407592in}{0.222904in}}{\pgfqpoint{1.411163in}{1.411163in}}%
\pgfusepath{clip}%
\pgfsetbuttcap%
\pgfsetroundjoin%
\definecolor{currentfill}{rgb}{0.909804,0.698039,0.423529}%
\pgfsetfillcolor{currentfill}%
\pgfsetfillopacity{0.120000}%
\pgfsetlinewidth{0.000000pt}%
\definecolor{currentstroke}{rgb}{0.000000,0.000000,0.000000}%
\pgfsetstrokecolor{currentstroke}%
\pgfsetdash{}{0pt}%
\pgfpathmoveto{\pgfqpoint{1.524947in}{1.033761in}}%
\pgfpathlineto{\pgfqpoint{1.543356in}{1.045449in}}%
\pgfpathlineto{\pgfqpoint{1.561850in}{1.057321in}}%
\pgfpathlineto{\pgfqpoint{1.569713in}{1.060194in}}%
\pgfpathlineto{\pgfqpoint{1.577555in}{1.063342in}}%
\pgfpathlineto{\pgfqpoint{1.559090in}{1.051389in}}%
\pgfpathlineto{\pgfqpoint{1.540709in}{1.039616in}}%
\pgfpathlineto{\pgfqpoint{1.532839in}{1.036544in}}%
\pgfpathlineto{\pgfqpoint{1.524947in}{1.033761in}}%
\pgfpathclose%
\pgfusepath{fill}%
\end{pgfscope}%
\begin{pgfscope}%
\pgfpathrectangle{\pgfqpoint{0.407592in}{0.222904in}}{\pgfqpoint{1.411163in}{1.411163in}}%
\pgfusepath{clip}%
\pgfsetbuttcap%
\pgfsetroundjoin%
\definecolor{currentfill}{rgb}{0.909804,0.698039,0.423529}%
\pgfsetfillcolor{currentfill}%
\pgfsetfillopacity{0.120000}%
\pgfsetlinewidth{0.000000pt}%
\definecolor{currentstroke}{rgb}{0.000000,0.000000,0.000000}%
\pgfsetstrokecolor{currentstroke}%
\pgfsetdash{}{0pt}%
\pgfpathmoveto{\pgfqpoint{1.246779in}{0.929450in}}%
\pgfpathlineto{\pgfqpoint{1.264588in}{0.937230in}}%
\pgfpathlineto{\pgfqpoint{1.282471in}{0.945378in}}%
\pgfpathlineto{\pgfqpoint{1.290812in}{0.944004in}}%
\pgfpathlineto{\pgfqpoint{1.299121in}{0.942899in}}%
\pgfpathlineto{\pgfqpoint{1.281275in}{0.934485in}}%
\pgfpathlineto{\pgfqpoint{1.263503in}{0.926433in}}%
\pgfpathlineto{\pgfqpoint{1.255157in}{0.927805in}}%
\pgfpathlineto{\pgfqpoint{1.246779in}{0.929450in}}%
\pgfpathclose%
\pgfusepath{fill}%
\end{pgfscope}%
\begin{pgfscope}%
\pgfpathrectangle{\pgfqpoint{0.407592in}{0.222904in}}{\pgfqpoint{1.411163in}{1.411163in}}%
\pgfusepath{clip}%
\pgfsetbuttcap%
\pgfsetroundjoin%
\definecolor{currentfill}{rgb}{0.909804,0.698039,0.423529}%
\pgfsetfillcolor{currentfill}%
\pgfsetfillopacity{0.120000}%
\pgfsetlinewidth{0.000000pt}%
\definecolor{currentstroke}{rgb}{0.000000,0.000000,0.000000}%
\pgfsetstrokecolor{currentstroke}%
\pgfsetdash{}{0pt}%
\pgfpathmoveto{\pgfqpoint{1.299121in}{0.942899in}}%
\pgfpathlineto{\pgfqpoint{1.317043in}{0.951655in}}%
\pgfpathlineto{\pgfqpoint{1.335042in}{0.960734in}}%
\pgfpathlineto{\pgfqpoint{1.343285in}{0.960152in}}%
\pgfpathlineto{\pgfqpoint{1.351498in}{0.959845in}}%
\pgfpathlineto{\pgfqpoint{1.333534in}{0.950524in}}%
\pgfpathlineto{\pgfqpoint{1.315646in}{0.941519in}}%
\pgfpathlineto{\pgfqpoint{1.307399in}{0.942068in}}%
\pgfpathlineto{\pgfqpoint{1.299121in}{0.942899in}}%
\pgfpathclose%
\pgfusepath{fill}%
\end{pgfscope}%
\begin{pgfscope}%
\pgfpathrectangle{\pgfqpoint{0.407592in}{0.222904in}}{\pgfqpoint{1.411163in}{1.411163in}}%
\pgfusepath{clip}%
\pgfsetbuttcap%
\pgfsetroundjoin%
\definecolor{currentfill}{rgb}{0.909804,0.698039,0.423529}%
\pgfsetfillcolor{currentfill}%
\pgfsetfillopacity{0.120000}%
\pgfsetlinewidth{0.000000pt}%
\definecolor{currentstroke}{rgb}{0.000000,0.000000,0.000000}%
\pgfsetstrokecolor{currentstroke}%
\pgfsetdash{}{0pt}%
\pgfpathmoveto{\pgfqpoint{1.194444in}{0.919711in}}%
\pgfpathlineto{\pgfqpoint{1.212148in}{0.926409in}}%
\pgfpathlineto{\pgfqpoint{1.229923in}{0.933519in}}%
\pgfpathlineto{\pgfqpoint{1.238368in}{0.931357in}}%
\pgfpathlineto{\pgfqpoint{1.246779in}{0.929450in}}%
\pgfpathlineto{\pgfqpoint{1.229043in}{0.922057in}}%
\pgfpathlineto{\pgfqpoint{1.211379in}{0.915074in}}%
\pgfpathlineto{\pgfqpoint{1.202929in}{0.917265in}}%
\pgfpathlineto{\pgfqpoint{1.194444in}{0.919711in}}%
\pgfpathclose%
\pgfusepath{fill}%
\end{pgfscope}%
\begin{pgfscope}%
\pgfpathrectangle{\pgfqpoint{0.407592in}{0.222904in}}{\pgfqpoint{1.411163in}{1.411163in}}%
\pgfusepath{clip}%
\pgfsetbuttcap%
\pgfsetroundjoin%
\definecolor{currentfill}{rgb}{0.909804,0.698039,0.423529}%
\pgfsetfillcolor{currentfill}%
\pgfsetfillopacity{0.120000}%
\pgfsetlinewidth{0.000000pt}%
\definecolor{currentstroke}{rgb}{0.000000,0.000000,0.000000}%
\pgfsetstrokecolor{currentstroke}%
\pgfsetdash{}{0pt}%
\pgfpathmoveto{\pgfqpoint{1.351498in}{0.959845in}}%
\pgfpathlineto{\pgfqpoint{1.369541in}{0.969465in}}%
\pgfpathlineto{\pgfqpoint{1.387663in}{0.979368in}}%
\pgfpathlineto{\pgfqpoint{1.395815in}{0.979561in}}%
\pgfpathlineto{\pgfqpoint{1.403939in}{0.980035in}}%
\pgfpathlineto{\pgfqpoint{1.385849in}{0.969922in}}%
\pgfpathlineto{\pgfqpoint{1.367839in}{0.960085in}}%
\pgfpathlineto{\pgfqpoint{1.359683in}{0.959820in}}%
\pgfpathlineto{\pgfqpoint{1.351498in}{0.959845in}}%
\pgfpathclose%
\pgfusepath{fill}%
\end{pgfscope}%
\begin{pgfscope}%
\pgfpathrectangle{\pgfqpoint{0.407592in}{0.222904in}}{\pgfqpoint{1.411163in}{1.411163in}}%
\pgfusepath{clip}%
\pgfsetbuttcap%
\pgfsetroundjoin%
\definecolor{currentfill}{rgb}{0.909804,0.698039,0.423529}%
\pgfsetfillcolor{currentfill}%
\pgfsetfillopacity{0.120000}%
\pgfsetlinewidth{0.000000pt}%
\definecolor{currentstroke}{rgb}{0.000000,0.000000,0.000000}%
\pgfsetstrokecolor{currentstroke}%
\pgfsetdash{}{0pt}%
\pgfpathmoveto{\pgfqpoint{1.142083in}{0.913837in}}%
\pgfpathlineto{\pgfqpoint{1.159692in}{0.919358in}}%
\pgfpathlineto{\pgfqpoint{1.177369in}{0.925336in}}%
\pgfpathlineto{\pgfqpoint{1.185924in}{0.922404in}}%
\pgfpathlineto{\pgfqpoint{1.194444in}{0.919711in}}%
\pgfpathlineto{\pgfqpoint{1.176809in}{0.913447in}}%
\pgfpathlineto{\pgfqpoint{1.159242in}{0.907639in}}%
\pgfpathlineto{\pgfqpoint{1.150681in}{0.910620in}}%
\pgfpathlineto{\pgfqpoint{1.142083in}{0.913837in}}%
\pgfpathclose%
\pgfusepath{fill}%
\end{pgfscope}%
\begin{pgfscope}%
\pgfpathrectangle{\pgfqpoint{0.407592in}{0.222904in}}{\pgfqpoint{1.411163in}{1.411163in}}%
\pgfusepath{clip}%
\pgfsetbuttcap%
\pgfsetroundjoin%
\definecolor{currentfill}{rgb}{0.909804,0.698039,0.423529}%
\pgfsetfillcolor{currentfill}%
\pgfsetfillopacity{0.120000}%
\pgfsetlinewidth{0.000000pt}%
\definecolor{currentstroke}{rgb}{0.000000,0.000000,0.000000}%
\pgfsetstrokecolor{currentstroke}%
\pgfsetdash{}{0pt}%
\pgfpathmoveto{\pgfqpoint{1.403939in}{0.980035in}}%
\pgfpathlineto{\pgfqpoint{1.422108in}{0.990409in}}%
\pgfpathlineto{\pgfqpoint{1.440360in}{1.001030in}}%
\pgfpathlineto{\pgfqpoint{1.448425in}{1.001973in}}%
\pgfpathlineto{\pgfqpoint{1.456465in}{1.003196in}}%
\pgfpathlineto{\pgfqpoint{1.438244in}{0.992401in}}%
\pgfpathlineto{\pgfqpoint{1.420106in}{0.981845in}}%
\pgfpathlineto{\pgfqpoint{1.412035in}{0.980795in}}%
\pgfpathlineto{\pgfqpoint{1.403939in}{0.980035in}}%
\pgfpathclose%
\pgfusepath{fill}%
\end{pgfscope}%
\begin{pgfscope}%
\pgfpathrectangle{\pgfqpoint{0.407592in}{0.222904in}}{\pgfqpoint{1.411163in}{1.411163in}}%
\pgfusepath{clip}%
\pgfsetbuttcap%
\pgfsetroundjoin%
\definecolor{currentfill}{rgb}{0.909804,0.698039,0.423529}%
\pgfsetfillcolor{currentfill}%
\pgfsetfillopacity{0.120000}%
\pgfsetlinewidth{0.000000pt}%
\definecolor{currentstroke}{rgb}{0.000000,0.000000,0.000000}%
\pgfsetstrokecolor{currentstroke}%
\pgfsetdash{}{0pt}%
\pgfpathmoveto{\pgfqpoint{1.089661in}{0.911908in}}%
\pgfpathlineto{\pgfqpoint{1.107186in}{0.916179in}}%
\pgfpathlineto{\pgfqpoint{1.124775in}{0.920945in}}%
\pgfpathlineto{\pgfqpoint{1.133447in}{0.917282in}}%
\pgfpathlineto{\pgfqpoint{1.142083in}{0.913837in}}%
\pgfpathlineto{\pgfqpoint{1.124539in}{0.908792in}}%
\pgfpathlineto{\pgfqpoint{1.107059in}{0.904247in}}%
\pgfpathlineto{\pgfqpoint{1.098379in}{0.907970in}}%
\pgfpathlineto{\pgfqpoint{1.089661in}{0.911908in}}%
\pgfpathclose%
\pgfusepath{fill}%
\end{pgfscope}%
\begin{pgfscope}%
\pgfpathrectangle{\pgfqpoint{0.407592in}{0.222904in}}{\pgfqpoint{1.411163in}{1.411163in}}%
\pgfusepath{clip}%
\pgfsetbuttcap%
\pgfsetroundjoin%
\definecolor{currentfill}{rgb}{0.909804,0.698039,0.423529}%
\pgfsetfillcolor{currentfill}%
\pgfsetfillopacity{0.120000}%
\pgfsetlinewidth{0.000000pt}%
\definecolor{currentstroke}{rgb}{0.000000,0.000000,0.000000}%
\pgfsetstrokecolor{currentstroke}%
\pgfsetdash{}{0pt}%
\pgfpathmoveto{\pgfqpoint{1.456465in}{1.003196in}}%
\pgfpathlineto{\pgfqpoint{1.474767in}{1.014219in}}%
\pgfpathlineto{\pgfqpoint{1.493152in}{1.025459in}}%
\pgfpathlineto{\pgfqpoint{1.501136in}{1.027115in}}%
\pgfpathlineto{\pgfqpoint{1.509096in}{1.029049in}}%
\pgfpathlineto{\pgfqpoint{1.490741in}{1.017671in}}%
\pgfpathlineto{\pgfqpoint{1.472469in}{1.006502in}}%
\pgfpathlineto{\pgfqpoint{1.464479in}{1.004705in}}%
\pgfpathlineto{\pgfqpoint{1.456465in}{1.003196in}}%
\pgfpathclose%
\pgfusepath{fill}%
\end{pgfscope}%
\begin{pgfscope}%
\pgfpathrectangle{\pgfqpoint{0.407592in}{0.222904in}}{\pgfqpoint{1.411163in}{1.411163in}}%
\pgfusepath{clip}%
\pgfsetbuttcap%
\pgfsetroundjoin%
\definecolor{currentfill}{rgb}{0.909804,0.698039,0.423529}%
\pgfsetfillcolor{currentfill}%
\pgfsetfillopacity{0.120000}%
\pgfsetlinewidth{0.000000pt}%
\definecolor{currentstroke}{rgb}{0.000000,0.000000,0.000000}%
\pgfsetstrokecolor{currentstroke}%
\pgfsetdash{}{0pt}%
\pgfpathmoveto{\pgfqpoint{1.561850in}{1.057321in}}%
\pgfpathlineto{\pgfqpoint{1.580430in}{1.069371in}}%
\pgfpathlineto{\pgfqpoint{1.588278in}{1.072284in}}%
\pgfpathlineto{\pgfqpoint{1.596105in}{1.075467in}}%
\pgfpathlineto{\pgfqpoint{1.577555in}{1.063342in}}%
\pgfpathlineto{\pgfqpoint{1.569713in}{1.060194in}}%
\pgfusepath{fill}%
\end{pgfscope}%
\begin{pgfscope}%
\pgfpathrectangle{\pgfqpoint{0.407592in}{0.222904in}}{\pgfqpoint{1.411163in}{1.411163in}}%
\pgfusepath{clip}%
\pgfsetbuttcap%
\pgfsetroundjoin%
\definecolor{currentfill}{rgb}{0.909804,0.698039,0.423529}%
\pgfsetfillcolor{currentfill}%
\pgfsetfillopacity{0.120000}%
\pgfsetlinewidth{0.000000pt}%
\definecolor{currentstroke}{rgb}{0.000000,0.000000,0.000000}%
\pgfsetstrokecolor{currentstroke}%
\pgfsetdash{}{0pt}%
\pgfpathmoveto{\pgfqpoint{1.509096in}{1.029049in}}%
\pgfpathlineto{\pgfqpoint{1.527535in}{1.040627in}}%
\pgfpathlineto{\pgfqpoint{1.546059in}{1.052395in}}%
\pgfpathlineto{\pgfqpoint{1.553966in}{1.054722in}}%
\pgfpathlineto{\pgfqpoint{1.561850in}{1.057321in}}%
\pgfpathlineto{\pgfqpoint{1.543356in}{1.045449in}}%
\pgfpathlineto{\pgfqpoint{1.524947in}{1.033761in}}%
\pgfpathlineto{\pgfqpoint{1.517033in}{1.031263in}}%
\pgfpathlineto{\pgfqpoint{1.509096in}{1.029049in}}%
\pgfpathclose%
\pgfusepath{fill}%
\end{pgfscope}%
\begin{pgfscope}%
\pgfpathrectangle{\pgfqpoint{0.407592in}{0.222904in}}{\pgfqpoint{1.411163in}{1.411163in}}%
\pgfusepath{clip}%
\pgfsetbuttcap%
\pgfsetroundjoin%
\definecolor{currentfill}{rgb}{0.909804,0.698039,0.423529}%
\pgfsetfillcolor{currentfill}%
\pgfsetfillopacity{0.120000}%
\pgfsetlinewidth{0.000000pt}%
\definecolor{currentstroke}{rgb}{0.000000,0.000000,0.000000}%
\pgfsetstrokecolor{currentstroke}%
\pgfsetdash{}{0pt}%
\pgfpathmoveto{\pgfqpoint{1.282471in}{0.945378in}}%
\pgfpathlineto{\pgfqpoint{1.300430in}{0.953874in}}%
\pgfpathlineto{\pgfqpoint{1.318465in}{0.962701in}}%
\pgfpathlineto{\pgfqpoint{1.326769in}{0.961587in}}%
\pgfpathlineto{\pgfqpoint{1.335042in}{0.960734in}}%
\pgfpathlineto{\pgfqpoint{1.317043in}{0.951655in}}%
\pgfpathlineto{\pgfqpoint{1.299121in}{0.942899in}}%
\pgfpathlineto{\pgfqpoint{1.290812in}{0.944004in}}%
\pgfpathlineto{\pgfqpoint{1.282471in}{0.945378in}}%
\pgfpathclose%
\pgfusepath{fill}%
\end{pgfscope}%
\begin{pgfscope}%
\pgfpathrectangle{\pgfqpoint{0.407592in}{0.222904in}}{\pgfqpoint{1.411163in}{1.411163in}}%
\pgfusepath{clip}%
\pgfsetbuttcap%
\pgfsetroundjoin%
\definecolor{currentfill}{rgb}{0.909804,0.698039,0.423529}%
\pgfsetfillcolor{currentfill}%
\pgfsetfillopacity{0.120000}%
\pgfsetlinewidth{0.000000pt}%
\definecolor{currentstroke}{rgb}{0.000000,0.000000,0.000000}%
\pgfsetstrokecolor{currentstroke}%
\pgfsetdash{}{0pt}%
\pgfpathmoveto{\pgfqpoint{1.229923in}{0.933519in}}%
\pgfpathlineto{\pgfqpoint{1.247771in}{0.941021in}}%
\pgfpathlineto{\pgfqpoint{1.265692in}{0.948896in}}%
\pgfpathlineto{\pgfqpoint{1.274098in}{0.947011in}}%
\pgfpathlineto{\pgfqpoint{1.282471in}{0.945378in}}%
\pgfpathlineto{\pgfqpoint{1.264588in}{0.937230in}}%
\pgfpathlineto{\pgfqpoint{1.246779in}{0.929450in}}%
\pgfpathlineto{\pgfqpoint{1.238368in}{0.931357in}}%
\pgfpathlineto{\pgfqpoint{1.229923in}{0.933519in}}%
\pgfpathclose%
\pgfusepath{fill}%
\end{pgfscope}%
\begin{pgfscope}%
\pgfpathrectangle{\pgfqpoint{0.407592in}{0.222904in}}{\pgfqpoint{1.411163in}{1.411163in}}%
\pgfusepath{clip}%
\pgfsetbuttcap%
\pgfsetroundjoin%
\definecolor{currentfill}{rgb}{0.909804,0.698039,0.423529}%
\pgfsetfillcolor{currentfill}%
\pgfsetfillopacity{0.120000}%
\pgfsetlinewidth{0.000000pt}%
\definecolor{currentstroke}{rgb}{0.000000,0.000000,0.000000}%
\pgfsetstrokecolor{currentstroke}%
\pgfsetdash{}{0pt}%
\pgfpathmoveto{\pgfqpoint{1.335042in}{0.960734in}}%
\pgfpathlineto{\pgfqpoint{1.353119in}{0.970121in}}%
\pgfpathlineto{\pgfqpoint{1.371275in}{0.979799in}}%
\pgfpathlineto{\pgfqpoint{1.379484in}{0.979449in}}%
\pgfpathlineto{\pgfqpoint{1.387663in}{0.979368in}}%
\pgfpathlineto{\pgfqpoint{1.369541in}{0.969465in}}%
\pgfpathlineto{\pgfqpoint{1.351498in}{0.959845in}}%
\pgfpathlineto{\pgfqpoint{1.343285in}{0.960152in}}%
\pgfpathlineto{\pgfqpoint{1.335042in}{0.960734in}}%
\pgfpathclose%
\pgfusepath{fill}%
\end{pgfscope}%
\begin{pgfscope}%
\pgfpathrectangle{\pgfqpoint{0.407592in}{0.222904in}}{\pgfqpoint{1.411163in}{1.411163in}}%
\pgfusepath{clip}%
\pgfsetbuttcap%
\pgfsetroundjoin%
\definecolor{currentfill}{rgb}{0.909804,0.698039,0.423529}%
\pgfsetfillcolor{currentfill}%
\pgfsetfillopacity{0.120000}%
\pgfsetlinewidth{0.000000pt}%
\definecolor{currentstroke}{rgb}{0.000000,0.000000,0.000000}%
\pgfsetstrokecolor{currentstroke}%
\pgfsetdash{}{0pt}%
\pgfpathmoveto{\pgfqpoint{1.177369in}{0.925336in}}%
\pgfpathlineto{\pgfqpoint{1.195114in}{0.931747in}}%
\pgfpathlineto{\pgfqpoint{1.212930in}{0.938574in}}%
\pgfpathlineto{\pgfqpoint{1.221444in}{0.935927in}}%
\pgfpathlineto{\pgfqpoint{1.229923in}{0.933519in}}%
\pgfpathlineto{\pgfqpoint{1.212148in}{0.926409in}}%
\pgfpathlineto{\pgfqpoint{1.194444in}{0.919711in}}%
\pgfpathlineto{\pgfqpoint{1.185924in}{0.922404in}}%
\pgfpathlineto{\pgfqpoint{1.177369in}{0.925336in}}%
\pgfpathclose%
\pgfusepath{fill}%
\end{pgfscope}%
\begin{pgfscope}%
\pgfpathrectangle{\pgfqpoint{0.407592in}{0.222904in}}{\pgfqpoint{1.411163in}{1.411163in}}%
\pgfusepath{clip}%
\pgfsetbuttcap%
\pgfsetroundjoin%
\definecolor{currentfill}{rgb}{0.909804,0.698039,0.423529}%
\pgfsetfillcolor{currentfill}%
\pgfsetfillopacity{0.120000}%
\pgfsetlinewidth{0.000000pt}%
\definecolor{currentstroke}{rgb}{0.000000,0.000000,0.000000}%
\pgfsetstrokecolor{currentstroke}%
\pgfsetdash{}{0pt}%
\pgfpathmoveto{\pgfqpoint{1.387663in}{0.979368in}}%
\pgfpathlineto{\pgfqpoint{1.405866in}{0.989540in}}%
\pgfpathlineto{\pgfqpoint{1.424150in}{0.999967in}}%
\pgfpathlineto{\pgfqpoint{1.432268in}{1.000363in}}%
\pgfpathlineto{\pgfqpoint{1.440360in}{1.001030in}}%
\pgfpathlineto{\pgfqpoint{1.422108in}{0.990409in}}%
\pgfpathlineto{\pgfqpoint{1.403939in}{0.980035in}}%
\pgfpathlineto{\pgfqpoint{1.395815in}{0.979561in}}%
\pgfpathlineto{\pgfqpoint{1.387663in}{0.979368in}}%
\pgfpathclose%
\pgfusepath{fill}%
\end{pgfscope}%
\begin{pgfscope}%
\pgfpathrectangle{\pgfqpoint{0.407592in}{0.222904in}}{\pgfqpoint{1.411163in}{1.411163in}}%
\pgfusepath{clip}%
\pgfsetbuttcap%
\pgfsetroundjoin%
\definecolor{currentfill}{rgb}{0.909804,0.698039,0.423529}%
\pgfsetfillcolor{currentfill}%
\pgfsetfillopacity{0.120000}%
\pgfsetlinewidth{0.000000pt}%
\definecolor{currentstroke}{rgb}{0.000000,0.000000,0.000000}%
\pgfsetstrokecolor{currentstroke}%
\pgfsetdash{}{0pt}%
\pgfpathmoveto{\pgfqpoint{1.124775in}{0.920945in}}%
\pgfpathlineto{\pgfqpoint{1.142428in}{0.926184in}}%
\pgfpathlineto{\pgfqpoint{1.160148in}{0.931878in}}%
\pgfpathlineto{\pgfqpoint{1.168777in}{0.928496in}}%
\pgfpathlineto{\pgfqpoint{1.177369in}{0.925336in}}%
\pgfpathlineto{\pgfqpoint{1.159692in}{0.919358in}}%
\pgfpathlineto{\pgfqpoint{1.142083in}{0.913837in}}%
\pgfpathlineto{\pgfqpoint{1.133447in}{0.917282in}}%
\pgfpathlineto{\pgfqpoint{1.124775in}{0.920945in}}%
\pgfpathclose%
\pgfusepath{fill}%
\end{pgfscope}%
\begin{pgfscope}%
\pgfpathrectangle{\pgfqpoint{0.407592in}{0.222904in}}{\pgfqpoint{1.411163in}{1.411163in}}%
\pgfusepath{clip}%
\pgfsetbuttcap%
\pgfsetroundjoin%
\definecolor{currentfill}{rgb}{0.909804,0.698039,0.423529}%
\pgfsetfillcolor{currentfill}%
\pgfsetfillopacity{0.120000}%
\pgfsetlinewidth{0.000000pt}%
\definecolor{currentstroke}{rgb}{0.000000,0.000000,0.000000}%
\pgfsetstrokecolor{currentstroke}%
\pgfsetdash{}{0pt}%
\pgfpathmoveto{\pgfqpoint{1.440360in}{1.001030in}}%
\pgfpathlineto{\pgfqpoint{1.458693in}{1.011885in}}%
\pgfpathlineto{\pgfqpoint{1.477110in}{1.022965in}}%
\pgfpathlineto{\pgfqpoint{1.485144in}{1.024076in}}%
\pgfpathlineto{\pgfqpoint{1.493152in}{1.025459in}}%
\pgfpathlineto{\pgfqpoint{1.474767in}{1.014219in}}%
\pgfpathlineto{\pgfqpoint{1.456465in}{1.003196in}}%
\pgfpathlineto{\pgfqpoint{1.448425in}{1.001973in}}%
\pgfpathlineto{\pgfqpoint{1.440360in}{1.001030in}}%
\pgfpathclose%
\pgfusepath{fill}%
\end{pgfscope}%
\begin{pgfscope}%
\pgfpathrectangle{\pgfqpoint{0.407592in}{0.222904in}}{\pgfqpoint{1.411163in}{1.411163in}}%
\pgfusepath{clip}%
\pgfsetbuttcap%
\pgfsetroundjoin%
\definecolor{currentfill}{rgb}{0.909804,0.698039,0.423529}%
\pgfsetfillcolor{currentfill}%
\pgfsetfillopacity{0.120000}%
\pgfsetlinewidth{0.000000pt}%
\definecolor{currentstroke}{rgb}{0.000000,0.000000,0.000000}%
\pgfsetstrokecolor{currentstroke}%
\pgfsetdash{}{0pt}%
\pgfpathmoveto{\pgfqpoint{1.546059in}{1.052395in}}%
\pgfpathlineto{\pgfqpoint{1.564668in}{1.064346in}}%
\pgfpathlineto{\pgfqpoint{1.572560in}{1.066725in}}%
\pgfpathlineto{\pgfqpoint{1.580430in}{1.069371in}}%
\pgfpathlineto{\pgfqpoint{1.561850in}{1.057321in}}%
\pgfpathlineto{\pgfqpoint{1.553966in}{1.054722in}}%
\pgfusepath{fill}%
\end{pgfscope}%
\begin{pgfscope}%
\pgfpathrectangle{\pgfqpoint{0.407592in}{0.222904in}}{\pgfqpoint{1.411163in}{1.411163in}}%
\pgfusepath{clip}%
\pgfsetbuttcap%
\pgfsetroundjoin%
\definecolor{currentfill}{rgb}{0.909804,0.698039,0.423529}%
\pgfsetfillcolor{currentfill}%
\pgfsetfillopacity{0.120000}%
\pgfsetlinewidth{0.000000pt}%
\definecolor{currentstroke}{rgb}{0.000000,0.000000,0.000000}%
\pgfsetstrokecolor{currentstroke}%
\pgfsetdash{}{0pt}%
\pgfpathmoveto{\pgfqpoint{1.493152in}{1.025459in}}%
\pgfpathlineto{\pgfqpoint{1.511622in}{1.036904in}}%
\pgfpathlineto{\pgfqpoint{1.530176in}{1.048547in}}%
\pgfpathlineto{\pgfqpoint{1.538129in}{1.050337in}}%
\pgfpathlineto{\pgfqpoint{1.546059in}{1.052395in}}%
\pgfpathlineto{\pgfqpoint{1.527535in}{1.040627in}}%
\pgfpathlineto{\pgfqpoint{1.509096in}{1.029049in}}%
\pgfpathlineto{\pgfqpoint{1.501136in}{1.027115in}}%
\pgfpathlineto{\pgfqpoint{1.493152in}{1.025459in}}%
\pgfpathclose%
\pgfusepath{fill}%
\end{pgfscope}%
\begin{pgfscope}%
\pgfpathrectangle{\pgfqpoint{0.407592in}{0.222904in}}{\pgfqpoint{1.411163in}{1.411163in}}%
\pgfusepath{clip}%
\pgfsetbuttcap%
\pgfsetroundjoin%
\definecolor{currentfill}{rgb}{0.909804,0.698039,0.423529}%
\pgfsetfillcolor{currentfill}%
\pgfsetfillopacity{0.120000}%
\pgfsetlinewidth{0.000000pt}%
\definecolor{currentstroke}{rgb}{0.000000,0.000000,0.000000}%
\pgfsetstrokecolor{currentstroke}%
\pgfsetdash{}{0pt}%
\pgfpathmoveto{\pgfqpoint{1.265692in}{0.948896in}}%
\pgfpathlineto{\pgfqpoint{1.283689in}{0.957125in}}%
\pgfpathlineto{\pgfqpoint{1.301761in}{0.965691in}}%
\pgfpathlineto{\pgfqpoint{1.310129in}{0.964072in}}%
\pgfpathlineto{\pgfqpoint{1.318465in}{0.962701in}}%
\pgfpathlineto{\pgfqpoint{1.300430in}{0.953874in}}%
\pgfpathlineto{\pgfqpoint{1.282471in}{0.945378in}}%
\pgfpathlineto{\pgfqpoint{1.274098in}{0.947011in}}%
\pgfpathlineto{\pgfqpoint{1.265692in}{0.948896in}}%
\pgfpathclose%
\pgfusepath{fill}%
\end{pgfscope}%
\begin{pgfscope}%
\pgfpathrectangle{\pgfqpoint{0.407592in}{0.222904in}}{\pgfqpoint{1.411163in}{1.411163in}}%
\pgfusepath{clip}%
\pgfsetbuttcap%
\pgfsetroundjoin%
\definecolor{currentfill}{rgb}{0.909804,0.698039,0.423529}%
\pgfsetfillcolor{currentfill}%
\pgfsetfillopacity{0.120000}%
\pgfsetlinewidth{0.000000pt}%
\definecolor{currentstroke}{rgb}{0.000000,0.000000,0.000000}%
\pgfsetstrokecolor{currentstroke}%
\pgfsetdash{}{0pt}%
\pgfpathmoveto{\pgfqpoint{1.318465in}{0.962701in}}%
\pgfpathlineto{\pgfqpoint{1.336577in}{0.971843in}}%
\pgfpathlineto{\pgfqpoint{1.354769in}{0.981283in}}%
\pgfpathlineto{\pgfqpoint{1.363037in}{0.980412in}}%
\pgfpathlineto{\pgfqpoint{1.371275in}{0.979799in}}%
\pgfpathlineto{\pgfqpoint{1.353119in}{0.970121in}}%
\pgfpathlineto{\pgfqpoint{1.335042in}{0.960734in}}%
\pgfpathlineto{\pgfqpoint{1.326769in}{0.961587in}}%
\pgfpathlineto{\pgfqpoint{1.318465in}{0.962701in}}%
\pgfpathclose%
\pgfusepath{fill}%
\end{pgfscope}%
\begin{pgfscope}%
\pgfpathrectangle{\pgfqpoint{0.407592in}{0.222904in}}{\pgfqpoint{1.411163in}{1.411163in}}%
\pgfusepath{clip}%
\pgfsetbuttcap%
\pgfsetroundjoin%
\definecolor{currentfill}{rgb}{0.909804,0.698039,0.423529}%
\pgfsetfillcolor{currentfill}%
\pgfsetfillopacity{0.120000}%
\pgfsetlinewidth{0.000000pt}%
\definecolor{currentstroke}{rgb}{0.000000,0.000000,0.000000}%
\pgfsetstrokecolor{currentstroke}%
\pgfsetdash{}{0pt}%
\pgfpathmoveto{\pgfqpoint{1.212930in}{0.938574in}}%
\pgfpathlineto{\pgfqpoint{1.230818in}{0.945795in}}%
\pgfpathlineto{\pgfqpoint{1.248779in}{0.953394in}}%
\pgfpathlineto{\pgfqpoint{1.257253in}{0.951026in}}%
\pgfpathlineto{\pgfqpoint{1.265692in}{0.948896in}}%
\pgfpathlineto{\pgfqpoint{1.247771in}{0.941021in}}%
\pgfpathlineto{\pgfqpoint{1.229923in}{0.933519in}}%
\pgfpathlineto{\pgfqpoint{1.221444in}{0.935927in}}%
\pgfpathlineto{\pgfqpoint{1.212930in}{0.938574in}}%
\pgfpathclose%
\pgfusepath{fill}%
\end{pgfscope}%
\begin{pgfscope}%
\pgfpathrectangle{\pgfqpoint{0.407592in}{0.222904in}}{\pgfqpoint{1.411163in}{1.411163in}}%
\pgfusepath{clip}%
\pgfsetbuttcap%
\pgfsetroundjoin%
\definecolor{currentfill}{rgb}{0.909804,0.698039,0.423529}%
\pgfsetfillcolor{currentfill}%
\pgfsetfillopacity{0.120000}%
\pgfsetlinewidth{0.000000pt}%
\definecolor{currentstroke}{rgb}{0.000000,0.000000,0.000000}%
\pgfsetstrokecolor{currentstroke}%
\pgfsetdash{}{0pt}%
\pgfpathmoveto{\pgfqpoint{1.371275in}{0.979799in}}%
\pgfpathlineto{\pgfqpoint{1.389512in}{0.989754in}}%
\pgfpathlineto{\pgfqpoint{1.407829in}{0.999971in}}%
\pgfpathlineto{\pgfqpoint{1.416004in}{0.999838in}}%
\pgfpathlineto{\pgfqpoint{1.424150in}{0.999967in}}%
\pgfpathlineto{\pgfqpoint{1.405866in}{0.989540in}}%
\pgfpathlineto{\pgfqpoint{1.387663in}{0.979368in}}%
\pgfpathlineto{\pgfqpoint{1.379484in}{0.979449in}}%
\pgfpathlineto{\pgfqpoint{1.371275in}{0.979799in}}%
\pgfpathclose%
\pgfusepath{fill}%
\end{pgfscope}%
\begin{pgfscope}%
\pgfpathrectangle{\pgfqpoint{0.407592in}{0.222904in}}{\pgfqpoint{1.411163in}{1.411163in}}%
\pgfusepath{clip}%
\pgfsetbuttcap%
\pgfsetroundjoin%
\definecolor{currentfill}{rgb}{0.909804,0.698039,0.423529}%
\pgfsetfillcolor{currentfill}%
\pgfsetfillopacity{0.120000}%
\pgfsetlinewidth{0.000000pt}%
\definecolor{currentstroke}{rgb}{0.000000,0.000000,0.000000}%
\pgfsetstrokecolor{currentstroke}%
\pgfsetdash{}{0pt}%
\pgfpathmoveto{\pgfqpoint{1.160148in}{0.931878in}}%
\pgfpathlineto{\pgfqpoint{1.177937in}{0.938006in}}%
\pgfpathlineto{\pgfqpoint{1.195795in}{0.944550in}}%
\pgfpathlineto{\pgfqpoint{1.204381in}{0.941450in}}%
\pgfpathlineto{\pgfqpoint{1.212930in}{0.938574in}}%
\pgfpathlineto{\pgfqpoint{1.195114in}{0.931747in}}%
\pgfpathlineto{\pgfqpoint{1.177369in}{0.925336in}}%
\pgfpathlineto{\pgfqpoint{1.168777in}{0.928496in}}%
\pgfpathlineto{\pgfqpoint{1.160148in}{0.931878in}}%
\pgfpathclose%
\pgfusepath{fill}%
\end{pgfscope}%
\begin{pgfscope}%
\pgfpathrectangle{\pgfqpoint{0.407592in}{0.222904in}}{\pgfqpoint{1.411163in}{1.411163in}}%
\pgfusepath{clip}%
\pgfsetbuttcap%
\pgfsetroundjoin%
\definecolor{currentfill}{rgb}{0.909804,0.698039,0.423529}%
\pgfsetfillcolor{currentfill}%
\pgfsetfillopacity{0.120000}%
\pgfsetlinewidth{0.000000pt}%
\definecolor{currentstroke}{rgb}{0.000000,0.000000,0.000000}%
\pgfsetstrokecolor{currentstroke}%
\pgfsetdash{}{0pt}%
\pgfpathmoveto{\pgfqpoint{1.424150in}{0.999967in}}%
\pgfpathlineto{\pgfqpoint{1.442516in}{1.010637in}}%
\pgfpathlineto{\pgfqpoint{1.460965in}{1.021538in}}%
\pgfpathlineto{\pgfqpoint{1.469051in}{1.022120in}}%
\pgfpathlineto{\pgfqpoint{1.477110in}{1.022965in}}%
\pgfpathlineto{\pgfqpoint{1.458693in}{1.011885in}}%
\pgfpathlineto{\pgfqpoint{1.440360in}{1.001030in}}%
\pgfpathlineto{\pgfqpoint{1.432268in}{1.000363in}}%
\pgfpathlineto{\pgfqpoint{1.424150in}{0.999967in}}%
\pgfpathclose%
\pgfusepath{fill}%
\end{pgfscope}%
\begin{pgfscope}%
\pgfpathrectangle{\pgfqpoint{0.407592in}{0.222904in}}{\pgfqpoint{1.411163in}{1.411163in}}%
\pgfusepath{clip}%
\pgfsetbuttcap%
\pgfsetroundjoin%
\definecolor{currentfill}{rgb}{0.909804,0.698039,0.423529}%
\pgfsetfillcolor{currentfill}%
\pgfsetfillopacity{0.120000}%
\pgfsetlinewidth{0.000000pt}%
\definecolor{currentstroke}{rgb}{0.000000,0.000000,0.000000}%
\pgfsetstrokecolor{currentstroke}%
\pgfsetdash{}{0pt}%
\pgfpathmoveto{\pgfqpoint{1.530176in}{1.048547in}}%
\pgfpathlineto{\pgfqpoint{1.548815in}{1.060377in}}%
\pgfpathlineto{\pgfqpoint{1.556753in}{1.062231in}}%
\pgfpathlineto{\pgfqpoint{1.564668in}{1.064346in}}%
\pgfpathlineto{\pgfqpoint{1.546059in}{1.052395in}}%
\pgfpathlineto{\pgfqpoint{1.538129in}{1.050337in}}%
\pgfusepath{fill}%
\end{pgfscope}%
\begin{pgfscope}%
\pgfpathrectangle{\pgfqpoint{0.407592in}{0.222904in}}{\pgfqpoint{1.411163in}{1.411163in}}%
\pgfusepath{clip}%
\pgfsetbuttcap%
\pgfsetroundjoin%
\definecolor{currentfill}{rgb}{0.909804,0.698039,0.423529}%
\pgfsetfillcolor{currentfill}%
\pgfsetfillopacity{0.120000}%
\pgfsetlinewidth{0.000000pt}%
\definecolor{currentstroke}{rgb}{0.000000,0.000000,0.000000}%
\pgfsetstrokecolor{currentstroke}%
\pgfsetdash{}{0pt}%
\pgfpathmoveto{\pgfqpoint{1.477110in}{1.022965in}}%
\pgfpathlineto{\pgfqpoint{1.495611in}{1.034258in}}%
\pgfpathlineto{\pgfqpoint{1.514196in}{1.045754in}}%
\pgfpathlineto{\pgfqpoint{1.522198in}{1.047020in}}%
\pgfpathlineto{\pgfqpoint{1.530176in}{1.048547in}}%
\pgfpathlineto{\pgfqpoint{1.511622in}{1.036904in}}%
\pgfpathlineto{\pgfqpoint{1.493152in}{1.025459in}}%
\pgfpathlineto{\pgfqpoint{1.485144in}{1.024076in}}%
\pgfpathlineto{\pgfqpoint{1.477110in}{1.022965in}}%
\pgfpathclose%
\pgfusepath{fill}%
\end{pgfscope}%
\begin{pgfscope}%
\pgfpathrectangle{\pgfqpoint{0.407592in}{0.222904in}}{\pgfqpoint{1.411163in}{1.411163in}}%
\pgfusepath{clip}%
\pgfsetbuttcap%
\pgfsetroundjoin%
\definecolor{currentfill}{rgb}{0.909804,0.698039,0.423529}%
\pgfsetfillcolor{currentfill}%
\pgfsetfillopacity{0.120000}%
\pgfsetlinewidth{0.000000pt}%
\definecolor{currentstroke}{rgb}{0.000000,0.000000,0.000000}%
\pgfsetstrokecolor{currentstroke}%
\pgfsetdash{}{0pt}%
\pgfpathmoveto{\pgfqpoint{1.301761in}{0.965691in}}%
\pgfpathlineto{\pgfqpoint{1.319911in}{0.974578in}}%
\pgfpathlineto{\pgfqpoint{1.338140in}{0.983771in}}%
\pgfpathlineto{\pgfqpoint{1.346470in}{0.982404in}}%
\pgfpathlineto{\pgfqpoint{1.354769in}{0.981283in}}%
\pgfpathlineto{\pgfqpoint{1.336577in}{0.971843in}}%
\pgfpathlineto{\pgfqpoint{1.318465in}{0.962701in}}%
\pgfpathlineto{\pgfqpoint{1.310129in}{0.964072in}}%
\pgfpathlineto{\pgfqpoint{1.301761in}{0.965691in}}%
\pgfpathclose%
\pgfusepath{fill}%
\end{pgfscope}%
\begin{pgfscope}%
\pgfpathrectangle{\pgfqpoint{0.407592in}{0.222904in}}{\pgfqpoint{1.411163in}{1.411163in}}%
\pgfusepath{clip}%
\pgfsetbuttcap%
\pgfsetroundjoin%
\definecolor{currentfill}{rgb}{0.909804,0.698039,0.423529}%
\pgfsetfillcolor{currentfill}%
\pgfsetfillopacity{0.120000}%
\pgfsetlinewidth{0.000000pt}%
\definecolor{currentstroke}{rgb}{0.000000,0.000000,0.000000}%
\pgfsetstrokecolor{currentstroke}%
\pgfsetdash{}{0pt}%
\pgfpathmoveto{\pgfqpoint{1.248779in}{0.953394in}}%
\pgfpathlineto{\pgfqpoint{1.266815in}{0.961351in}}%
\pgfpathlineto{\pgfqpoint{1.284926in}{0.969651in}}%
\pgfpathlineto{\pgfqpoint{1.293360in}{0.967553in}}%
\pgfpathlineto{\pgfqpoint{1.301761in}{0.965691in}}%
\pgfpathlineto{\pgfqpoint{1.283689in}{0.957125in}}%
\pgfpathlineto{\pgfqpoint{1.265692in}{0.948896in}}%
\pgfpathlineto{\pgfqpoint{1.257253in}{0.951026in}}%
\pgfpathlineto{\pgfqpoint{1.248779in}{0.953394in}}%
\pgfpathclose%
\pgfusepath{fill}%
\end{pgfscope}%
\begin{pgfscope}%
\pgfpathrectangle{\pgfqpoint{0.407592in}{0.222904in}}{\pgfqpoint{1.411163in}{1.411163in}}%
\pgfusepath{clip}%
\pgfsetbuttcap%
\pgfsetroundjoin%
\definecolor{currentfill}{rgb}{0.909804,0.698039,0.423529}%
\pgfsetfillcolor{currentfill}%
\pgfsetfillopacity{0.120000}%
\pgfsetlinewidth{0.000000pt}%
\definecolor{currentstroke}{rgb}{0.000000,0.000000,0.000000}%
\pgfsetstrokecolor{currentstroke}%
\pgfsetdash{}{0pt}%
\pgfpathmoveto{\pgfqpoint{1.354769in}{0.981283in}}%
\pgfpathlineto{\pgfqpoint{1.373041in}{0.991007in}}%
\pgfpathlineto{\pgfqpoint{1.391393in}{1.001002in}}%
\pgfpathlineto{\pgfqpoint{1.399626in}{1.000361in}}%
\pgfpathlineto{\pgfqpoint{1.407829in}{0.999971in}}%
\pgfpathlineto{\pgfqpoint{1.389512in}{0.989754in}}%
\pgfpathlineto{\pgfqpoint{1.371275in}{0.979799in}}%
\pgfpathlineto{\pgfqpoint{1.363037in}{0.980412in}}%
\pgfpathlineto{\pgfqpoint{1.354769in}{0.981283in}}%
\pgfpathclose%
\pgfusepath{fill}%
\end{pgfscope}%
\begin{pgfscope}%
\pgfpathrectangle{\pgfqpoint{0.407592in}{0.222904in}}{\pgfqpoint{1.411163in}{1.411163in}}%
\pgfusepath{clip}%
\pgfsetbuttcap%
\pgfsetroundjoin%
\definecolor{currentfill}{rgb}{0.909804,0.698039,0.423529}%
\pgfsetfillcolor{currentfill}%
\pgfsetfillopacity{0.120000}%
\pgfsetlinewidth{0.000000pt}%
\definecolor{currentstroke}{rgb}{0.000000,0.000000,0.000000}%
\pgfsetstrokecolor{currentstroke}%
\pgfsetdash{}{0pt}%
\pgfpathmoveto{\pgfqpoint{1.195795in}{0.944550in}}%
\pgfpathlineto{\pgfqpoint{1.213725in}{0.951491in}}%
\pgfpathlineto{\pgfqpoint{1.231727in}{0.958812in}}%
\pgfpathlineto{\pgfqpoint{1.240271in}{0.955991in}}%
\pgfpathlineto{\pgfqpoint{1.248779in}{0.953394in}}%
\pgfpathlineto{\pgfqpoint{1.230818in}{0.945795in}}%
\pgfpathlineto{\pgfqpoint{1.212930in}{0.938574in}}%
\pgfpathlineto{\pgfqpoint{1.204381in}{0.941450in}}%
\pgfpathlineto{\pgfqpoint{1.195795in}{0.944550in}}%
\pgfpathclose%
\pgfusepath{fill}%
\end{pgfscope}%
\begin{pgfscope}%
\pgfpathrectangle{\pgfqpoint{0.407592in}{0.222904in}}{\pgfqpoint{1.411163in}{1.411163in}}%
\pgfusepath{clip}%
\pgfsetbuttcap%
\pgfsetroundjoin%
\definecolor{currentfill}{rgb}{0.909804,0.698039,0.423529}%
\pgfsetfillcolor{currentfill}%
\pgfsetfillopacity{0.120000}%
\pgfsetlinewidth{0.000000pt}%
\definecolor{currentstroke}{rgb}{0.000000,0.000000,0.000000}%
\pgfsetstrokecolor{currentstroke}%
\pgfsetdash{}{0pt}%
\pgfpathmoveto{\pgfqpoint{1.514196in}{1.045754in}}%
\pgfpathlineto{\pgfqpoint{1.532867in}{1.057446in}}%
\pgfpathlineto{\pgfqpoint{1.540853in}{1.058783in}}%
\pgfpathlineto{\pgfqpoint{1.548815in}{1.060377in}}%
\pgfpathlineto{\pgfqpoint{1.530176in}{1.048547in}}%
\pgfpathlineto{\pgfqpoint{1.522198in}{1.047020in}}%
\pgfusepath{fill}%
\end{pgfscope}%
\begin{pgfscope}%
\pgfpathrectangle{\pgfqpoint{0.407592in}{0.222904in}}{\pgfqpoint{1.411163in}{1.411163in}}%
\pgfusepath{clip}%
\pgfsetbuttcap%
\pgfsetroundjoin%
\definecolor{currentfill}{rgb}{0.909804,0.698039,0.423529}%
\pgfsetfillcolor{currentfill}%
\pgfsetfillopacity{0.120000}%
\pgfsetlinewidth{0.000000pt}%
\definecolor{currentstroke}{rgb}{0.000000,0.000000,0.000000}%
\pgfsetstrokecolor{currentstroke}%
\pgfsetdash{}{0pt}%
\pgfpathmoveto{\pgfqpoint{1.407829in}{0.999971in}}%
\pgfpathlineto{\pgfqpoint{1.426229in}{1.010439in}}%
\pgfpathlineto{\pgfqpoint{1.444712in}{1.021146in}}%
\pgfpathlineto{\pgfqpoint{1.452852in}{1.021215in}}%
\pgfpathlineto{\pgfqpoint{1.460965in}{1.021538in}}%
\pgfpathlineto{\pgfqpoint{1.442516in}{1.010637in}}%
\pgfpathlineto{\pgfqpoint{1.424150in}{0.999967in}}%
\pgfpathlineto{\pgfqpoint{1.416004in}{0.999838in}}%
\pgfpathlineto{\pgfqpoint{1.407829in}{0.999971in}}%
\pgfpathclose%
\pgfusepath{fill}%
\end{pgfscope}%
\begin{pgfscope}%
\pgfpathrectangle{\pgfqpoint{0.407592in}{0.222904in}}{\pgfqpoint{1.411163in}{1.411163in}}%
\pgfusepath{clip}%
\pgfsetbuttcap%
\pgfsetroundjoin%
\definecolor{currentfill}{rgb}{0.909804,0.698039,0.423529}%
\pgfsetfillcolor{currentfill}%
\pgfsetfillopacity{0.120000}%
\pgfsetlinewidth{0.000000pt}%
\definecolor{currentstroke}{rgb}{0.000000,0.000000,0.000000}%
\pgfsetstrokecolor{currentstroke}%
\pgfsetdash{}{0pt}%
\pgfpathmoveto{\pgfqpoint{1.460965in}{1.021538in}}%
\pgfpathlineto{\pgfqpoint{1.479498in}{1.032660in}}%
\pgfpathlineto{\pgfqpoint{1.498116in}{1.043993in}}%
\pgfpathlineto{\pgfqpoint{1.506169in}{1.044746in}}%
\pgfpathlineto{\pgfqpoint{1.514196in}{1.045754in}}%
\pgfpathlineto{\pgfqpoint{1.495611in}{1.034258in}}%
\pgfpathlineto{\pgfqpoint{1.477110in}{1.022965in}}%
\pgfpathlineto{\pgfqpoint{1.469051in}{1.022120in}}%
\pgfpathlineto{\pgfqpoint{1.460965in}{1.021538in}}%
\pgfpathclose%
\pgfusepath{fill}%
\end{pgfscope}%
\begin{pgfscope}%
\pgfpathrectangle{\pgfqpoint{0.407592in}{0.222904in}}{\pgfqpoint{1.411163in}{1.411163in}}%
\pgfusepath{clip}%
\pgfsetbuttcap%
\pgfsetroundjoin%
\definecolor{currentfill}{rgb}{0.909804,0.698039,0.423529}%
\pgfsetfillcolor{currentfill}%
\pgfsetfillopacity{0.120000}%
\pgfsetlinewidth{0.000000pt}%
\definecolor{currentstroke}{rgb}{0.000000,0.000000,0.000000}%
\pgfsetstrokecolor{currentstroke}%
\pgfsetdash{}{0pt}%
\pgfpathmoveto{\pgfqpoint{1.284926in}{0.969651in}}%
\pgfpathlineto{\pgfqpoint{1.303115in}{0.978277in}}%
\pgfpathlineto{\pgfqpoint{1.321382in}{0.987216in}}%
\pgfpathlineto{\pgfqpoint{1.329777in}{0.985377in}}%
\pgfpathlineto{\pgfqpoint{1.338140in}{0.983771in}}%
\pgfpathlineto{\pgfqpoint{1.319911in}{0.974578in}}%
\pgfpathlineto{\pgfqpoint{1.301761in}{0.965691in}}%
\pgfpathlineto{\pgfqpoint{1.293360in}{0.967553in}}%
\pgfpathlineto{\pgfqpoint{1.284926in}{0.969651in}}%
\pgfpathclose%
\pgfusepath{fill}%
\end{pgfscope}%
\begin{pgfscope}%
\pgfpathrectangle{\pgfqpoint{0.407592in}{0.222904in}}{\pgfqpoint{1.411163in}{1.411163in}}%
\pgfusepath{clip}%
\pgfsetbuttcap%
\pgfsetroundjoin%
\definecolor{currentfill}{rgb}{0.909804,0.698039,0.423529}%
\pgfsetfillcolor{currentfill}%
\pgfsetfillopacity{0.120000}%
\pgfsetlinewidth{0.000000pt}%
\definecolor{currentstroke}{rgb}{0.000000,0.000000,0.000000}%
\pgfsetstrokecolor{currentstroke}%
\pgfsetdash{}{0pt}%
\pgfpathmoveto{\pgfqpoint{1.498116in}{1.043993in}}%
\pgfpathlineto{\pgfqpoint{1.516819in}{1.055528in}}%
\pgfpathlineto{\pgfqpoint{1.524856in}{1.056362in}}%
\pgfpathlineto{\pgfqpoint{1.532867in}{1.057446in}}%
\pgfpathlineto{\pgfqpoint{1.514196in}{1.045754in}}%
\pgfpathlineto{\pgfqpoint{1.506169in}{1.044746in}}%
\pgfusepath{fill}%
\end{pgfscope}%
\begin{pgfscope}%
\pgfpathrectangle{\pgfqpoint{0.407592in}{0.222904in}}{\pgfqpoint{1.411163in}{1.411163in}}%
\pgfusepath{clip}%
\pgfsetbuttcap%
\pgfsetroundjoin%
\definecolor{currentfill}{rgb}{0.909804,0.698039,0.423529}%
\pgfsetfillcolor{currentfill}%
\pgfsetfillopacity{0.120000}%
\pgfsetlinewidth{0.000000pt}%
\definecolor{currentstroke}{rgb}{0.000000,0.000000,0.000000}%
\pgfsetstrokecolor{currentstroke}%
\pgfsetdash{}{0pt}%
\pgfpathmoveto{\pgfqpoint{1.338140in}{0.983771in}}%
\pgfpathlineto{\pgfqpoint{1.356448in}{0.993255in}}%
\pgfpathlineto{\pgfqpoint{1.374837in}{1.003016in}}%
\pgfpathlineto{\pgfqpoint{1.383130in}{1.001888in}}%
\pgfpathlineto{\pgfqpoint{1.391393in}{1.001002in}}%
\pgfpathlineto{\pgfqpoint{1.373041in}{0.991007in}}%
\pgfpathlineto{\pgfqpoint{1.354769in}{0.981283in}}%
\pgfpathlineto{\pgfqpoint{1.346470in}{0.982404in}}%
\pgfpathlineto{\pgfqpoint{1.338140in}{0.983771in}}%
\pgfpathclose%
\pgfusepath{fill}%
\end{pgfscope}%
\begin{pgfscope}%
\pgfpathrectangle{\pgfqpoint{0.407592in}{0.222904in}}{\pgfqpoint{1.411163in}{1.411163in}}%
\pgfusepath{clip}%
\pgfsetbuttcap%
\pgfsetroundjoin%
\definecolor{currentfill}{rgb}{0.909804,0.698039,0.423529}%
\pgfsetfillcolor{currentfill}%
\pgfsetfillopacity{0.120000}%
\pgfsetlinewidth{0.000000pt}%
\definecolor{currentstroke}{rgb}{0.000000,0.000000,0.000000}%
\pgfsetstrokecolor{currentstroke}%
\pgfsetdash{}{0pt}%
\pgfpathmoveto{\pgfqpoint{1.231727in}{0.958812in}}%
\pgfpathlineto{\pgfqpoint{1.249804in}{0.966496in}}%
\pgfpathlineto{\pgfqpoint{1.267956in}{0.974527in}}%
\pgfpathlineto{\pgfqpoint{1.276459in}{0.971978in}}%
\pgfpathlineto{\pgfqpoint{1.284926in}{0.969651in}}%
\pgfpathlineto{\pgfqpoint{1.266815in}{0.961351in}}%
\pgfpathlineto{\pgfqpoint{1.248779in}{0.953394in}}%
\pgfpathlineto{\pgfqpoint{1.240271in}{0.955991in}}%
\pgfpathlineto{\pgfqpoint{1.231727in}{0.958812in}}%
\pgfpathclose%
\pgfusepath{fill}%
\end{pgfscope}%
\begin{pgfscope}%
\pgfpathrectangle{\pgfqpoint{0.407592in}{0.222904in}}{\pgfqpoint{1.411163in}{1.411163in}}%
\pgfusepath{clip}%
\pgfsetbuttcap%
\pgfsetroundjoin%
\definecolor{currentfill}{rgb}{0.909804,0.698039,0.423529}%
\pgfsetfillcolor{currentfill}%
\pgfsetfillopacity{0.120000}%
\pgfsetlinewidth{0.000000pt}%
\definecolor{currentstroke}{rgb}{0.000000,0.000000,0.000000}%
\pgfsetstrokecolor{currentstroke}%
\pgfsetdash{}{0pt}%
\pgfpathmoveto{\pgfqpoint{1.391393in}{1.001002in}}%
\pgfpathlineto{\pgfqpoint{1.409828in}{1.011254in}}%
\pgfpathlineto{\pgfqpoint{1.428345in}{1.021754in}}%
\pgfpathlineto{\pgfqpoint{1.436543in}{1.021327in}}%
\pgfpathlineto{\pgfqpoint{1.444712in}{1.021146in}}%
\pgfpathlineto{\pgfqpoint{1.426229in}{1.010439in}}%
\pgfpathlineto{\pgfqpoint{1.407829in}{0.999971in}}%
\pgfpathlineto{\pgfqpoint{1.399626in}{1.000361in}}%
\pgfpathlineto{\pgfqpoint{1.391393in}{1.001002in}}%
\pgfpathclose%
\pgfusepath{fill}%
\end{pgfscope}%
\begin{pgfscope}%
\pgfpathrectangle{\pgfqpoint{0.407592in}{0.222904in}}{\pgfqpoint{1.411163in}{1.411163in}}%
\pgfusepath{clip}%
\pgfsetbuttcap%
\pgfsetroundjoin%
\definecolor{currentfill}{rgb}{0.909804,0.698039,0.423529}%
\pgfsetfillcolor{currentfill}%
\pgfsetfillopacity{0.120000}%
\pgfsetlinewidth{0.000000pt}%
\definecolor{currentstroke}{rgb}{0.000000,0.000000,0.000000}%
\pgfsetstrokecolor{currentstroke}%
\pgfsetdash{}{0pt}%
\pgfpathmoveto{\pgfqpoint{1.444712in}{1.021146in}}%
\pgfpathlineto{\pgfqpoint{1.463278in}{1.032081in}}%
\pgfpathlineto{\pgfqpoint{1.481930in}{1.043235in}}%
\pgfpathlineto{\pgfqpoint{1.490036in}{1.043490in}}%
\pgfpathlineto{\pgfqpoint{1.498116in}{1.043993in}}%
\pgfpathlineto{\pgfqpoint{1.479498in}{1.032660in}}%
\pgfpathlineto{\pgfqpoint{1.460965in}{1.021538in}}%
\pgfpathlineto{\pgfqpoint{1.452852in}{1.021215in}}%
\pgfpathlineto{\pgfqpoint{1.444712in}{1.021146in}}%
\pgfpathclose%
\pgfusepath{fill}%
\end{pgfscope}%
\begin{pgfscope}%
\pgfpathrectangle{\pgfqpoint{0.407592in}{0.222904in}}{\pgfqpoint{1.411163in}{1.411163in}}%
\pgfusepath{clip}%
\pgfsetbuttcap%
\pgfsetroundjoin%
\definecolor{currentfill}{rgb}{0.909804,0.698039,0.423529}%
\pgfsetfillcolor{currentfill}%
\pgfsetfillopacity{0.120000}%
\pgfsetlinewidth{0.000000pt}%
\definecolor{currentstroke}{rgb}{0.000000,0.000000,0.000000}%
\pgfsetstrokecolor{currentstroke}%
\pgfsetdash{}{0pt}%
\pgfpathmoveto{\pgfqpoint{1.481930in}{1.043235in}}%
\pgfpathlineto{\pgfqpoint{1.500667in}{1.054598in}}%
\pgfpathlineto{\pgfqpoint{1.508756in}{1.054941in}}%
\pgfpathlineto{\pgfqpoint{1.516819in}{1.055528in}}%
\pgfpathlineto{\pgfqpoint{1.498116in}{1.043993in}}%
\pgfpathlineto{\pgfqpoint{1.490036in}{1.043490in}}%
\pgfusepath{fill}%
\end{pgfscope}%
\begin{pgfscope}%
\pgfpathrectangle{\pgfqpoint{0.407592in}{0.222904in}}{\pgfqpoint{1.411163in}{1.411163in}}%
\pgfusepath{clip}%
\pgfsetbuttcap%
\pgfsetroundjoin%
\definecolor{currentfill}{rgb}{0.909804,0.698039,0.423529}%
\pgfsetfillcolor{currentfill}%
\pgfsetfillopacity{0.120000}%
\pgfsetlinewidth{0.000000pt}%
\definecolor{currentstroke}{rgb}{0.000000,0.000000,0.000000}%
\pgfsetstrokecolor{currentstroke}%
\pgfsetdash{}{0pt}%
\pgfpathmoveto{\pgfqpoint{1.321382in}{0.987216in}}%
\pgfpathlineto{\pgfqpoint{1.339728in}{0.996452in}}%
\pgfpathlineto{\pgfqpoint{1.358155in}{1.005973in}}%
\pgfpathlineto{\pgfqpoint{1.366512in}{1.004380in}}%
\pgfpathlineto{\pgfqpoint{1.374837in}{1.003016in}}%
\pgfpathlineto{\pgfqpoint{1.356448in}{0.993255in}}%
\pgfpathlineto{\pgfqpoint{1.338140in}{0.983771in}}%
\pgfpathlineto{\pgfqpoint{1.329777in}{0.985377in}}%
\pgfpathlineto{\pgfqpoint{1.321382in}{0.987216in}}%
\pgfpathclose%
\pgfusepath{fill}%
\end{pgfscope}%
\begin{pgfscope}%
\pgfpathrectangle{\pgfqpoint{0.407592in}{0.222904in}}{\pgfqpoint{1.411163in}{1.411163in}}%
\pgfusepath{clip}%
\pgfsetbuttcap%
\pgfsetroundjoin%
\definecolor{currentfill}{rgb}{0.909804,0.698039,0.423529}%
\pgfsetfillcolor{currentfill}%
\pgfsetfillopacity{0.120000}%
\pgfsetlinewidth{0.000000pt}%
\definecolor{currentstroke}{rgb}{0.000000,0.000000,0.000000}%
\pgfsetstrokecolor{currentstroke}%
\pgfsetdash{}{0pt}%
\pgfpathmoveto{\pgfqpoint{1.267956in}{0.974527in}}%
\pgfpathlineto{\pgfqpoint{1.286185in}{0.982890in}}%
\pgfpathlineto{\pgfqpoint{1.304492in}{0.991571in}}%
\pgfpathlineto{\pgfqpoint{1.312954in}{0.989282in}}%
\pgfpathlineto{\pgfqpoint{1.321382in}{0.987216in}}%
\pgfpathlineto{\pgfqpoint{1.303115in}{0.978277in}}%
\pgfpathlineto{\pgfqpoint{1.284926in}{0.969651in}}%
\pgfpathlineto{\pgfqpoint{1.276459in}{0.971978in}}%
\pgfpathlineto{\pgfqpoint{1.267956in}{0.974527in}}%
\pgfpathclose%
\pgfusepath{fill}%
\end{pgfscope}%
\begin{pgfscope}%
\pgfpathrectangle{\pgfqpoint{0.407592in}{0.222904in}}{\pgfqpoint{1.411163in}{1.411163in}}%
\pgfusepath{clip}%
\pgfsetbuttcap%
\pgfsetroundjoin%
\definecolor{currentfill}{rgb}{0.909804,0.698039,0.423529}%
\pgfsetfillcolor{currentfill}%
\pgfsetfillopacity{0.120000}%
\pgfsetlinewidth{0.000000pt}%
\definecolor{currentstroke}{rgb}{0.000000,0.000000,0.000000}%
\pgfsetstrokecolor{currentstroke}%
\pgfsetdash{}{0pt}%
\pgfpathmoveto{\pgfqpoint{1.374837in}{1.003016in}}%
\pgfpathlineto{\pgfqpoint{1.393308in}{1.013044in}}%
\pgfpathlineto{\pgfqpoint{1.411862in}{1.023325in}}%
\pgfpathlineto{\pgfqpoint{1.420118in}{1.022421in}}%
\pgfpathlineto{\pgfqpoint{1.428345in}{1.021754in}}%
\pgfpathlineto{\pgfqpoint{1.409828in}{1.011254in}}%
\pgfpathlineto{\pgfqpoint{1.391393in}{1.001002in}}%
\pgfpathlineto{\pgfqpoint{1.383130in}{1.001888in}}%
\pgfpathlineto{\pgfqpoint{1.374837in}{1.003016in}}%
\pgfpathclose%
\pgfusepath{fill}%
\end{pgfscope}%
\begin{pgfscope}%
\pgfpathrectangle{\pgfqpoint{0.407592in}{0.222904in}}{\pgfqpoint{1.411163in}{1.411163in}}%
\pgfusepath{clip}%
\pgfsetbuttcap%
\pgfsetroundjoin%
\definecolor{currentfill}{rgb}{0.909804,0.698039,0.423529}%
\pgfsetfillcolor{currentfill}%
\pgfsetfillopacity{0.120000}%
\pgfsetlinewidth{0.000000pt}%
\definecolor{currentstroke}{rgb}{0.000000,0.000000,0.000000}%
\pgfsetstrokecolor{currentstroke}%
\pgfsetdash{}{0pt}%
\pgfpathmoveto{\pgfqpoint{1.428345in}{1.021754in}}%
\pgfpathlineto{\pgfqpoint{1.446947in}{1.032489in}}%
\pgfpathlineto{\pgfqpoint{1.465633in}{1.043450in}}%
\pgfpathlineto{\pgfqpoint{1.473795in}{1.043223in}}%
\pgfpathlineto{\pgfqpoint{1.481930in}{1.043235in}}%
\pgfpathlineto{\pgfqpoint{1.463278in}{1.032081in}}%
\pgfpathlineto{\pgfqpoint{1.444712in}{1.021146in}}%
\pgfpathlineto{\pgfqpoint{1.436543in}{1.021327in}}%
\pgfpathlineto{\pgfqpoint{1.428345in}{1.021754in}}%
\pgfpathclose%
\pgfusepath{fill}%
\end{pgfscope}%
\begin{pgfscope}%
\pgfpathrectangle{\pgfqpoint{0.407592in}{0.222904in}}{\pgfqpoint{1.411163in}{1.411163in}}%
\pgfusepath{clip}%
\pgfsetbuttcap%
\pgfsetroundjoin%
\definecolor{currentfill}{rgb}{0.909804,0.698039,0.423529}%
\pgfsetfillcolor{currentfill}%
\pgfsetfillopacity{0.120000}%
\pgfsetlinewidth{0.000000pt}%
\definecolor{currentstroke}{rgb}{0.000000,0.000000,0.000000}%
\pgfsetstrokecolor{currentstroke}%
\pgfsetdash{}{0pt}%
\pgfpathmoveto{\pgfqpoint{1.465633in}{1.043450in}}%
\pgfpathlineto{\pgfqpoint{1.484405in}{1.054628in}}%
\pgfpathlineto{\pgfqpoint{1.492550in}{1.054495in}}%
\pgfpathlineto{\pgfqpoint{1.500667in}{1.054598in}}%
\pgfpathlineto{\pgfqpoint{1.481930in}{1.043235in}}%
\pgfpathlineto{\pgfqpoint{1.473795in}{1.043223in}}%
\pgfusepath{fill}%
\end{pgfscope}%
\begin{pgfscope}%
\pgfpathrectangle{\pgfqpoint{0.407592in}{0.222904in}}{\pgfqpoint{1.411163in}{1.411163in}}%
\pgfusepath{clip}%
\pgfsetbuttcap%
\pgfsetroundjoin%
\definecolor{currentfill}{rgb}{0.909804,0.698039,0.423529}%
\pgfsetfillcolor{currentfill}%
\pgfsetfillopacity{0.120000}%
\pgfsetlinewidth{0.000000pt}%
\definecolor{currentstroke}{rgb}{0.000000,0.000000,0.000000}%
\pgfsetstrokecolor{currentstroke}%
\pgfsetdash{}{0pt}%
\pgfpathmoveto{\pgfqpoint{1.304492in}{0.991571in}}%
\pgfpathlineto{\pgfqpoint{1.322878in}{1.000555in}}%
\pgfpathlineto{\pgfqpoint{1.341344in}{1.009830in}}%
\pgfpathlineto{\pgfqpoint{1.349766in}{1.007792in}}%
\pgfpathlineto{\pgfqpoint{1.358155in}{1.005973in}}%
\pgfpathlineto{\pgfqpoint{1.339728in}{0.996452in}}%
\pgfpathlineto{\pgfqpoint{1.321382in}{0.987216in}}%
\pgfpathlineto{\pgfqpoint{1.312954in}{0.989282in}}%
\pgfpathlineto{\pgfqpoint{1.304492in}{0.991571in}}%
\pgfpathclose%
\pgfusepath{fill}%
\end{pgfscope}%
\begin{pgfscope}%
\pgfpathrectangle{\pgfqpoint{0.407592in}{0.222904in}}{\pgfqpoint{1.411163in}{1.411163in}}%
\pgfusepath{clip}%
\pgfsetbuttcap%
\pgfsetroundjoin%
\definecolor{currentfill}{rgb}{0.909804,0.698039,0.423529}%
\pgfsetfillcolor{currentfill}%
\pgfsetfillopacity{0.120000}%
\pgfsetlinewidth{0.000000pt}%
\definecolor{currentstroke}{rgb}{0.000000,0.000000,0.000000}%
\pgfsetstrokecolor{currentstroke}%
\pgfsetdash{}{0pt}%
\pgfpathmoveto{\pgfqpoint{1.358155in}{1.005973in}}%
\pgfpathlineto{\pgfqpoint{1.376664in}{1.015767in}}%
\pgfpathlineto{\pgfqpoint{1.395255in}{1.025821in}}%
\pgfpathlineto{\pgfqpoint{1.403574in}{1.024460in}}%
\pgfpathlineto{\pgfqpoint{1.411862in}{1.023325in}}%
\pgfpathlineto{\pgfqpoint{1.393308in}{1.013044in}}%
\pgfpathlineto{\pgfqpoint{1.374837in}{1.003016in}}%
\pgfpathlineto{\pgfqpoint{1.366512in}{1.004380in}}%
\pgfpathlineto{\pgfqpoint{1.358155in}{1.005973in}}%
\pgfpathclose%
\pgfusepath{fill}%
\end{pgfscope}%
\begin{pgfscope}%
\pgfpathrectangle{\pgfqpoint{0.407592in}{0.222904in}}{\pgfqpoint{1.411163in}{1.411163in}}%
\pgfusepath{clip}%
\pgfsetbuttcap%
\pgfsetroundjoin%
\definecolor{currentfill}{rgb}{0.909804,0.698039,0.423529}%
\pgfsetfillcolor{currentfill}%
\pgfsetfillopacity{0.120000}%
\pgfsetlinewidth{0.000000pt}%
\definecolor{currentstroke}{rgb}{0.000000,0.000000,0.000000}%
\pgfsetstrokecolor{currentstroke}%
\pgfsetdash{}{0pt}%
\pgfpathmoveto{\pgfqpoint{1.411862in}{1.023325in}}%
\pgfpathlineto{\pgfqpoint{1.430499in}{1.033849in}}%
\pgfpathlineto{\pgfqpoint{1.449221in}{1.044607in}}%
\pgfpathlineto{\pgfqpoint{1.457442in}{1.043913in}}%
\pgfpathlineto{\pgfqpoint{1.465633in}{1.043450in}}%
\pgfpathlineto{\pgfqpoint{1.446947in}{1.032489in}}%
\pgfpathlineto{\pgfqpoint{1.428345in}{1.021754in}}%
\pgfpathlineto{\pgfqpoint{1.420118in}{1.022421in}}%
\pgfpathlineto{\pgfqpoint{1.411862in}{1.023325in}}%
\pgfpathclose%
\pgfusepath{fill}%
\end{pgfscope}%
\begin{pgfscope}%
\pgfpathrectangle{\pgfqpoint{0.407592in}{0.222904in}}{\pgfqpoint{1.411163in}{1.411163in}}%
\pgfusepath{clip}%
\pgfsetbuttcap%
\pgfsetroundjoin%
\definecolor{currentfill}{rgb}{0.909804,0.698039,0.423529}%
\pgfsetfillcolor{currentfill}%
\pgfsetfillopacity{0.120000}%
\pgfsetlinewidth{0.000000pt}%
\definecolor{currentstroke}{rgb}{0.000000,0.000000,0.000000}%
\pgfsetstrokecolor{currentstroke}%
\pgfsetdash{}{0pt}%
\pgfpathmoveto{\pgfqpoint{1.449221in}{1.044607in}}%
\pgfpathlineto{\pgfqpoint{1.468029in}{1.055588in}}%
\pgfpathlineto{\pgfqpoint{1.476232in}{1.054994in}}%
\pgfpathlineto{\pgfqpoint{1.484405in}{1.054628in}}%
\pgfpathlineto{\pgfqpoint{1.465633in}{1.043450in}}%
\pgfpathlineto{\pgfqpoint{1.457442in}{1.043913in}}%
\pgfusepath{fill}%
\end{pgfscope}%
\begin{pgfscope}%
\pgfpathrectangle{\pgfqpoint{0.407592in}{0.222904in}}{\pgfqpoint{1.411163in}{1.411163in}}%
\pgfusepath{clip}%
\pgfsetbuttcap%
\pgfsetroundjoin%
\definecolor{currentfill}{rgb}{0.909804,0.698039,0.423529}%
\pgfsetfillcolor{currentfill}%
\pgfsetfillopacity{0.120000}%
\pgfsetlinewidth{0.000000pt}%
\definecolor{currentstroke}{rgb}{0.000000,0.000000,0.000000}%
\pgfsetstrokecolor{currentstroke}%
\pgfsetdash{}{0pt}%
\pgfpathmoveto{\pgfqpoint{1.341344in}{1.009830in}}%
\pgfpathlineto{\pgfqpoint{1.359891in}{1.019384in}}%
\pgfpathlineto{\pgfqpoint{1.378521in}{1.029206in}}%
\pgfpathlineto{\pgfqpoint{1.386904in}{1.027405in}}%
\pgfpathlineto{\pgfqpoint{1.395255in}{1.025821in}}%
\pgfpathlineto{\pgfqpoint{1.376664in}{1.015767in}}%
\pgfpathlineto{\pgfqpoint{1.358155in}{1.005973in}}%
\pgfpathlineto{\pgfqpoint{1.349766in}{1.007792in}}%
\pgfpathlineto{\pgfqpoint{1.341344in}{1.009830in}}%
\pgfpathclose%
\pgfusepath{fill}%
\end{pgfscope}%
\begin{pgfscope}%
\pgfpathrectangle{\pgfqpoint{0.407592in}{0.222904in}}{\pgfqpoint{1.411163in}{1.411163in}}%
\pgfusepath{clip}%
\pgfsetbuttcap%
\pgfsetroundjoin%
\definecolor{currentfill}{rgb}{0.909804,0.698039,0.423529}%
\pgfsetfillcolor{currentfill}%
\pgfsetfillopacity{0.120000}%
\pgfsetlinewidth{0.000000pt}%
\definecolor{currentstroke}{rgb}{0.000000,0.000000,0.000000}%
\pgfsetstrokecolor{currentstroke}%
\pgfsetdash{}{0pt}%
\pgfpathmoveto{\pgfqpoint{1.395255in}{1.025821in}}%
\pgfpathlineto{\pgfqpoint{1.413930in}{1.036126in}}%
\pgfpathlineto{\pgfqpoint{1.432689in}{1.046671in}}%
\pgfpathlineto{\pgfqpoint{1.440971in}{1.045528in}}%
\pgfpathlineto{\pgfqpoint{1.449221in}{1.044607in}}%
\pgfpathlineto{\pgfqpoint{1.430499in}{1.033849in}}%
\pgfpathlineto{\pgfqpoint{1.411862in}{1.023325in}}%
\pgfpathlineto{\pgfqpoint{1.403574in}{1.024460in}}%
\pgfpathlineto{\pgfqpoint{1.395255in}{1.025821in}}%
\pgfpathclose%
\pgfusepath{fill}%
\end{pgfscope}%
\begin{pgfscope}%
\pgfpathrectangle{\pgfqpoint{0.407592in}{0.222904in}}{\pgfqpoint{1.411163in}{1.411163in}}%
\pgfusepath{clip}%
\pgfsetbuttcap%
\pgfsetroundjoin%
\definecolor{currentfill}{rgb}{0.909804,0.698039,0.423529}%
\pgfsetfillcolor{currentfill}%
\pgfsetfillopacity{0.120000}%
\pgfsetlinewidth{0.000000pt}%
\definecolor{currentstroke}{rgb}{0.000000,0.000000,0.000000}%
\pgfsetstrokecolor{currentstroke}%
\pgfsetdash{}{0pt}%
\pgfpathmoveto{\pgfqpoint{1.432689in}{1.046671in}}%
\pgfpathlineto{\pgfqpoint{1.451535in}{1.057447in}}%
\pgfpathlineto{\pgfqpoint{1.459797in}{1.056407in}}%
\pgfpathlineto{\pgfqpoint{1.468029in}{1.055588in}}%
\pgfpathlineto{\pgfqpoint{1.449221in}{1.044607in}}%
\pgfpathlineto{\pgfqpoint{1.440971in}{1.045528in}}%
\pgfusepath{fill}%
\end{pgfscope}%
\begin{pgfscope}%
\pgfpathrectangle{\pgfqpoint{0.407592in}{0.222904in}}{\pgfqpoint{1.411163in}{1.411163in}}%
\pgfusepath{clip}%
\pgfsetbuttcap%
\pgfsetroundjoin%
\definecolor{currentfill}{rgb}{0.909804,0.698039,0.423529}%
\pgfsetfillcolor{currentfill}%
\pgfsetfillopacity{0.120000}%
\pgfsetlinewidth{0.000000pt}%
\definecolor{currentstroke}{rgb}{0.000000,0.000000,0.000000}%
\pgfsetstrokecolor{currentstroke}%
\pgfsetdash{}{0pt}%
\pgfpathmoveto{\pgfqpoint{1.378521in}{1.029206in}}%
\pgfpathlineto{\pgfqpoint{1.397235in}{1.039285in}}%
\pgfpathlineto{\pgfqpoint{1.416033in}{1.049611in}}%
\pgfpathlineto{\pgfqpoint{1.424377in}{1.048034in}}%
\pgfpathlineto{\pgfqpoint{1.432689in}{1.046671in}}%
\pgfpathlineto{\pgfqpoint{1.413930in}{1.036126in}}%
\pgfpathlineto{\pgfqpoint{1.395255in}{1.025821in}}%
\pgfpathlineto{\pgfqpoint{1.386904in}{1.027405in}}%
\pgfpathlineto{\pgfqpoint{1.378521in}{1.029206in}}%
\pgfpathclose%
\pgfusepath{fill}%
\end{pgfscope}%
\begin{pgfscope}%
\pgfpathrectangle{\pgfqpoint{0.407592in}{0.222904in}}{\pgfqpoint{1.411163in}{1.411163in}}%
\pgfusepath{clip}%
\pgfsetbuttcap%
\pgfsetroundjoin%
\definecolor{currentfill}{rgb}{0.909804,0.698039,0.423529}%
\pgfsetfillcolor{currentfill}%
\pgfsetfillopacity{0.120000}%
\pgfsetlinewidth{0.000000pt}%
\definecolor{currentstroke}{rgb}{0.000000,0.000000,0.000000}%
\pgfsetstrokecolor{currentstroke}%
\pgfsetdash{}{0pt}%
\pgfpathmoveto{\pgfqpoint{1.416033in}{1.049611in}}%
\pgfpathlineto{\pgfqpoint{1.434916in}{1.060174in}}%
\pgfpathlineto{\pgfqpoint{1.443241in}{1.058704in}}%
\pgfpathlineto{\pgfqpoint{1.451535in}{1.057447in}}%
\pgfpathlineto{\pgfqpoint{1.432689in}{1.046671in}}%
\pgfpathlineto{\pgfqpoint{1.424377in}{1.048034in}}%
\pgfusepath{fill}%
\end{pgfscope}%
\begin{pgfscope}%
\pgfpathrectangle{\pgfqpoint{0.407592in}{0.222904in}}{\pgfqpoint{1.411163in}{1.411163in}}%
\pgfusepath{clip}%
\pgfsetbuttcap%
\pgfsetroundjoin%
\definecolor{currentfill}{rgb}{0.749020,0.772549,0.792157}%
\pgfsetfillcolor{currentfill}%
\pgfsetfillopacity{0.160000}%
\pgfsetlinewidth{0.000000pt}%
\definecolor{currentstroke}{rgb}{0.000000,0.000000,0.000000}%
\pgfsetstrokecolor{currentstroke}%
\pgfsetdash{}{0pt}%
\pgfpathmoveto{\pgfqpoint{1.207276in}{1.516005in}}%
\pgfpathlineto{\pgfqpoint{1.209257in}{1.531330in}}%
\pgfpathlineto{\pgfqpoint{1.162616in}{1.532749in}}%
\pgfpathlineto{\pgfqpoint{1.115711in}{1.533155in}}%
\pgfpathlineto{\pgfqpoint{1.115287in}{1.517809in}}%
\pgfpathlineto{\pgfqpoint{1.161411in}{1.517408in}}%
\pgfusepath{fill}%
\end{pgfscope}%
\begin{pgfscope}%
\pgfpathrectangle{\pgfqpoint{0.407592in}{0.222904in}}{\pgfqpoint{1.411163in}{1.411163in}}%
\pgfusepath{clip}%
\pgfsetbuttcap%
\pgfsetroundjoin%
\definecolor{currentfill}{rgb}{0.749020,0.772549,0.792157}%
\pgfsetfillcolor{currentfill}%
\pgfsetfillopacity{0.160000}%
\pgfsetlinewidth{0.000000pt}%
\definecolor{currentstroke}{rgb}{0.000000,0.000000,0.000000}%
\pgfsetstrokecolor{currentstroke}%
\pgfsetdash{}{0pt}%
\pgfpathmoveto{\pgfqpoint{1.115287in}{1.517809in}}%
\pgfpathlineto{\pgfqpoint{1.115711in}{1.533155in}}%
\pgfpathlineto{\pgfqpoint{1.068695in}{1.532546in}}%
\pgfpathlineto{\pgfqpoint{1.021719in}{1.530924in}}%
\pgfpathlineto{\pgfqpoint{1.022860in}{1.515603in}}%
\pgfpathlineto{\pgfqpoint{1.069053in}{1.517207in}}%
\pgfusepath{fill}%
\end{pgfscope}%
\begin{pgfscope}%
\pgfpathrectangle{\pgfqpoint{0.407592in}{0.222904in}}{\pgfqpoint{1.411163in}{1.411163in}}%
\pgfusepath{clip}%
\pgfsetbuttcap%
\pgfsetroundjoin%
\definecolor{currentfill}{rgb}{0.749020,0.772549,0.792157}%
\pgfsetfillcolor{currentfill}%
\pgfsetfillopacity{0.160000}%
\pgfsetlinewidth{0.000000pt}%
\definecolor{currentstroke}{rgb}{0.000000,0.000000,0.000000}%
\pgfsetstrokecolor{currentstroke}%
\pgfsetdash{}{0pt}%
\pgfpathmoveto{\pgfqpoint{1.203303in}{1.485270in}}%
\pgfpathlineto{\pgfqpoint{1.205291in}{1.500652in}}%
\pgfpathlineto{\pgfqpoint{1.207276in}{1.516005in}}%
\pgfpathlineto{\pgfqpoint{1.161411in}{1.517408in}}%
\pgfpathlineto{\pgfqpoint{1.115287in}{1.517809in}}%
\pgfpathlineto{\pgfqpoint{1.114863in}{1.502435in}}%
\pgfpathlineto{\pgfqpoint{1.114437in}{1.487033in}}%
\pgfpathlineto{\pgfqpoint{1.158996in}{1.486641in}}%
\pgfpathlineto{\pgfqpoint{1.203303in}{1.485270in}}%
\pgfpathclose%
\pgfusepath{fill}%
\end{pgfscope}%
\begin{pgfscope}%
\pgfpathrectangle{\pgfqpoint{0.407592in}{0.222904in}}{\pgfqpoint{1.411163in}{1.411163in}}%
\pgfusepath{clip}%
\pgfsetbuttcap%
\pgfsetroundjoin%
\definecolor{currentfill}{rgb}{0.749020,0.772549,0.792157}%
\pgfsetfillcolor{currentfill}%
\pgfsetfillopacity{0.160000}%
\pgfsetlinewidth{0.000000pt}%
\definecolor{currentstroke}{rgb}{0.000000,0.000000,0.000000}%
\pgfsetstrokecolor{currentstroke}%
\pgfsetdash{}{0pt}%
\pgfpathmoveto{\pgfqpoint{1.297627in}{1.510206in}}%
\pgfpathlineto{\pgfqpoint{1.301145in}{1.525467in}}%
\pgfpathlineto{\pgfqpoint{1.255484in}{1.528901in}}%
\pgfpathlineto{\pgfqpoint{1.209257in}{1.531330in}}%
\pgfpathlineto{\pgfqpoint{1.207276in}{1.516005in}}%
\pgfpathlineto{\pgfqpoint{1.252731in}{1.513603in}}%
\pgfusepath{fill}%
\end{pgfscope}%
\begin{pgfscope}%
\pgfpathrectangle{\pgfqpoint{0.407592in}{0.222904in}}{\pgfqpoint{1.411163in}{1.411163in}}%
\pgfusepath{clip}%
\pgfsetbuttcap%
\pgfsetroundjoin%
\definecolor{currentfill}{rgb}{0.749020,0.772549,0.792157}%
\pgfsetfillcolor{currentfill}%
\pgfsetfillopacity{0.160000}%
\pgfsetlinewidth{0.000000pt}%
\definecolor{currentstroke}{rgb}{0.000000,0.000000,0.000000}%
\pgfsetstrokecolor{currentstroke}%
\pgfsetdash{}{0pt}%
\pgfpathmoveto{\pgfqpoint{1.114437in}{1.487033in}}%
\pgfpathlineto{\pgfqpoint{1.114863in}{1.502435in}}%
\pgfpathlineto{\pgfqpoint{1.115287in}{1.517809in}}%
\pgfpathlineto{\pgfqpoint{1.069053in}{1.517207in}}%
\pgfpathlineto{\pgfqpoint{1.022860in}{1.515603in}}%
\pgfpathlineto{\pgfqpoint{1.024002in}{1.500254in}}%
\pgfpathlineto{\pgfqpoint{1.025147in}{1.484877in}}%
\pgfpathlineto{\pgfqpoint{1.069772in}{1.486445in}}%
\pgfpathlineto{\pgfqpoint{1.114437in}{1.487033in}}%
\pgfpathclose%
\pgfusepath{fill}%
\end{pgfscope}%
\begin{pgfscope}%
\pgfpathrectangle{\pgfqpoint{0.407592in}{0.222904in}}{\pgfqpoint{1.411163in}{1.411163in}}%
\pgfusepath{clip}%
\pgfsetbuttcap%
\pgfsetroundjoin%
\definecolor{currentfill}{rgb}{0.749020,0.772549,0.792157}%
\pgfsetfillcolor{currentfill}%
\pgfsetfillopacity{0.160000}%
\pgfsetlinewidth{0.000000pt}%
\definecolor{currentstroke}{rgb}{0.000000,0.000000,0.000000}%
\pgfsetstrokecolor{currentstroke}%
\pgfsetdash{}{0pt}%
\pgfpathmoveto{\pgfqpoint{1.022860in}{1.515603in}}%
\pgfpathlineto{\pgfqpoint{1.021719in}{1.530924in}}%
\pgfpathlineto{\pgfqpoint{0.974935in}{1.528293in}}%
\pgfpathlineto{\pgfqpoint{0.928496in}{1.524657in}}%
\pgfpathlineto{\pgfqpoint{0.931195in}{1.509405in}}%
\pgfpathlineto{\pgfqpoint{0.976856in}{1.513001in}}%
\pgfusepath{fill}%
\end{pgfscope}%
\begin{pgfscope}%
\pgfpathrectangle{\pgfqpoint{0.407592in}{0.222904in}}{\pgfqpoint{1.411163in}{1.411163in}}%
\pgfusepath{clip}%
\pgfsetbuttcap%
\pgfsetroundjoin%
\definecolor{currentfill}{rgb}{0.749020,0.772549,0.792157}%
\pgfsetfillcolor{currentfill}%
\pgfsetfillopacity{0.160000}%
\pgfsetlinewidth{0.000000pt}%
\definecolor{currentstroke}{rgb}{0.000000,0.000000,0.000000}%
\pgfsetstrokecolor{currentstroke}%
\pgfsetdash{}{0pt}%
\pgfpathmoveto{\pgfqpoint{1.290572in}{1.479604in}}%
\pgfpathlineto{\pgfqpoint{1.294103in}{1.494918in}}%
\pgfpathlineto{\pgfqpoint{1.297627in}{1.510206in}}%
\pgfpathlineto{\pgfqpoint{1.252731in}{1.513603in}}%
\pgfpathlineto{\pgfqpoint{1.207276in}{1.516005in}}%
\pgfpathlineto{\pgfqpoint{1.205291in}{1.500652in}}%
\pgfpathlineto{\pgfqpoint{1.203303in}{1.485270in}}%
\pgfpathlineto{\pgfqpoint{1.247210in}{1.482922in}}%
\pgfpathlineto{\pgfqpoint{1.290572in}{1.479604in}}%
\pgfpathclose%
\pgfusepath{fill}%
\end{pgfscope}%
\begin{pgfscope}%
\pgfpathrectangle{\pgfqpoint{0.407592in}{0.222904in}}{\pgfqpoint{1.411163in}{1.411163in}}%
\pgfusepath{clip}%
\pgfsetbuttcap%
\pgfsetroundjoin%
\definecolor{currentfill}{rgb}{0.749020,0.772549,0.792157}%
\pgfsetfillcolor{currentfill}%
\pgfsetfillopacity{0.160000}%
\pgfsetlinewidth{0.000000pt}%
\definecolor{currentstroke}{rgb}{0.000000,0.000000,0.000000}%
\pgfsetstrokecolor{currentstroke}%
\pgfsetdash{}{0pt}%
\pgfpathmoveto{\pgfqpoint{1.025147in}{1.484877in}}%
\pgfpathlineto{\pgfqpoint{1.024002in}{1.500254in}}%
\pgfpathlineto{\pgfqpoint{1.022860in}{1.515603in}}%
\pgfpathlineto{\pgfqpoint{0.976856in}{1.513001in}}%
\pgfpathlineto{\pgfqpoint{0.931195in}{1.509405in}}%
\pgfpathlineto{\pgfqpoint{0.933898in}{1.494127in}}%
\pgfpathlineto{\pgfqpoint{0.936606in}{1.478821in}}%
\pgfpathlineto{\pgfqpoint{0.980709in}{1.482334in}}%
\pgfpathlineto{\pgfqpoint{1.025147in}{1.484877in}}%
\pgfpathclose%
\pgfusepath{fill}%
\end{pgfscope}%
\begin{pgfscope}%
\pgfpathrectangle{\pgfqpoint{0.407592in}{0.222904in}}{\pgfqpoint{1.411163in}{1.411163in}}%
\pgfusepath{clip}%
\pgfsetbuttcap%
\pgfsetroundjoin%
\definecolor{currentfill}{rgb}{0.749020,0.772549,0.792157}%
\pgfsetfillcolor{currentfill}%
\pgfsetfillopacity{0.160000}%
\pgfsetlinewidth{0.000000pt}%
\definecolor{currentstroke}{rgb}{0.000000,0.000000,0.000000}%
\pgfsetstrokecolor{currentstroke}%
\pgfsetdash{}{0pt}%
\pgfpathmoveto{\pgfqpoint{1.199315in}{1.454422in}}%
\pgfpathlineto{\pgfqpoint{1.201311in}{1.469860in}}%
\pgfpathlineto{\pgfqpoint{1.203303in}{1.485270in}}%
\pgfpathlineto{\pgfqpoint{1.158996in}{1.486641in}}%
\pgfpathlineto{\pgfqpoint{1.114437in}{1.487033in}}%
\pgfpathlineto{\pgfqpoint{1.114011in}{1.471602in}}%
\pgfpathlineto{\pgfqpoint{1.113584in}{1.456142in}}%
\pgfpathlineto{\pgfqpoint{1.156573in}{1.455760in}}%
\pgfpathlineto{\pgfqpoint{1.199315in}{1.454422in}}%
\pgfpathclose%
\pgfusepath{fill}%
\end{pgfscope}%
\begin{pgfscope}%
\pgfpathrectangle{\pgfqpoint{0.407592in}{0.222904in}}{\pgfqpoint{1.411163in}{1.411163in}}%
\pgfusepath{clip}%
\pgfsetbuttcap%
\pgfsetroundjoin%
\definecolor{currentfill}{rgb}{0.749020,0.772549,0.792157}%
\pgfsetfillcolor{currentfill}%
\pgfsetfillopacity{0.160000}%
\pgfsetlinewidth{0.000000pt}%
\definecolor{currentstroke}{rgb}{0.000000,0.000000,0.000000}%
\pgfsetstrokecolor{currentstroke}%
\pgfsetdash{}{0pt}%
\pgfpathmoveto{\pgfqpoint{1.113584in}{1.456142in}}%
\pgfpathlineto{\pgfqpoint{1.114011in}{1.471602in}}%
\pgfpathlineto{\pgfqpoint{1.114437in}{1.487033in}}%
\pgfpathlineto{\pgfqpoint{1.069772in}{1.486445in}}%
\pgfpathlineto{\pgfqpoint{1.025147in}{1.484877in}}%
\pgfpathlineto{\pgfqpoint{1.026293in}{1.469472in}}%
\pgfpathlineto{\pgfqpoint{1.027442in}{1.454039in}}%
\pgfpathlineto{\pgfqpoint{1.070493in}{1.455568in}}%
\pgfpathlineto{\pgfqpoint{1.113584in}{1.456142in}}%
\pgfpathclose%
\pgfusepath{fill}%
\end{pgfscope}%
\begin{pgfscope}%
\pgfpathrectangle{\pgfqpoint{0.407592in}{0.222904in}}{\pgfqpoint{1.411163in}{1.411163in}}%
\pgfusepath{clip}%
\pgfsetbuttcap%
\pgfsetroundjoin%
\definecolor{currentfill}{rgb}{0.749020,0.772549,0.792157}%
\pgfsetfillcolor{currentfill}%
\pgfsetfillopacity{0.160000}%
\pgfsetlinewidth{0.000000pt}%
\definecolor{currentstroke}{rgb}{0.000000,0.000000,0.000000}%
\pgfsetstrokecolor{currentstroke}%
\pgfsetdash{}{0pt}%
\pgfpathmoveto{\pgfqpoint{1.385145in}{1.500462in}}%
\pgfpathlineto{\pgfqpoint{1.390166in}{1.515613in}}%
\pgfpathlineto{\pgfqpoint{1.346090in}{1.521034in}}%
\pgfpathlineto{\pgfqpoint{1.301145in}{1.525467in}}%
\pgfpathlineto{\pgfqpoint{1.297627in}{1.510206in}}%
\pgfpathlineto{\pgfqpoint{1.341815in}{1.505822in}}%
\pgfusepath{fill}%
\end{pgfscope}%
\begin{pgfscope}%
\pgfpathrectangle{\pgfqpoint{0.407592in}{0.222904in}}{\pgfqpoint{1.411163in}{1.411163in}}%
\pgfusepath{clip}%
\pgfsetbuttcap%
\pgfsetroundjoin%
\definecolor{currentfill}{rgb}{0.749020,0.772549,0.792157}%
\pgfsetfillcolor{currentfill}%
\pgfsetfillopacity{0.160000}%
\pgfsetlinewidth{0.000000pt}%
\definecolor{currentstroke}{rgb}{0.000000,0.000000,0.000000}%
\pgfsetstrokecolor{currentstroke}%
\pgfsetdash{}{0pt}%
\pgfpathmoveto{\pgfqpoint{1.283493in}{1.448893in}}%
\pgfpathlineto{\pgfqpoint{1.287036in}{1.464262in}}%
\pgfpathlineto{\pgfqpoint{1.290572in}{1.479604in}}%
\pgfpathlineto{\pgfqpoint{1.247210in}{1.482922in}}%
\pgfpathlineto{\pgfqpoint{1.203303in}{1.485270in}}%
\pgfpathlineto{\pgfqpoint{1.201311in}{1.469860in}}%
\pgfpathlineto{\pgfqpoint{1.199315in}{1.454422in}}%
\pgfpathlineto{\pgfqpoint{1.241669in}{1.452131in}}%
\pgfpathlineto{\pgfqpoint{1.283493in}{1.448893in}}%
\pgfpathclose%
\pgfusepath{fill}%
\end{pgfscope}%
\begin{pgfscope}%
\pgfpathrectangle{\pgfqpoint{0.407592in}{0.222904in}}{\pgfqpoint{1.411163in}{1.411163in}}%
\pgfusepath{clip}%
\pgfsetbuttcap%
\pgfsetroundjoin%
\definecolor{currentfill}{rgb}{0.749020,0.772549,0.792157}%
\pgfsetfillcolor{currentfill}%
\pgfsetfillopacity{0.160000}%
\pgfsetlinewidth{0.000000pt}%
\definecolor{currentstroke}{rgb}{0.000000,0.000000,0.000000}%
\pgfsetstrokecolor{currentstroke}%
\pgfsetdash{}{0pt}%
\pgfpathmoveto{\pgfqpoint{0.931195in}{1.509405in}}%
\pgfpathlineto{\pgfqpoint{0.928496in}{1.524657in}}%
\pgfpathlineto{\pgfqpoint{0.882555in}{1.520026in}}%
\pgfpathlineto{\pgfqpoint{0.837265in}{1.514408in}}%
\pgfpathlineto{\pgfqpoint{0.841501in}{1.499270in}}%
\pgfpathlineto{\pgfqpoint{0.886026in}{1.504825in}}%
\pgfusepath{fill}%
\end{pgfscope}%
\begin{pgfscope}%
\pgfpathrectangle{\pgfqpoint{0.407592in}{0.222904in}}{\pgfqpoint{1.411163in}{1.411163in}}%
\pgfusepath{clip}%
\pgfsetbuttcap%
\pgfsetroundjoin%
\definecolor{currentfill}{rgb}{0.749020,0.772549,0.792157}%
\pgfsetfillcolor{currentfill}%
\pgfsetfillopacity{0.160000}%
\pgfsetlinewidth{0.000000pt}%
\definecolor{currentstroke}{rgb}{0.000000,0.000000,0.000000}%
\pgfsetstrokecolor{currentstroke}%
\pgfsetdash{}{0pt}%
\pgfpathmoveto{\pgfqpoint{1.027442in}{1.454039in}}%
\pgfpathlineto{\pgfqpoint{1.026293in}{1.469472in}}%
\pgfpathlineto{\pgfqpoint{1.025147in}{1.484877in}}%
\pgfpathlineto{\pgfqpoint{0.980709in}{1.482334in}}%
\pgfpathlineto{\pgfqpoint{0.936606in}{1.478821in}}%
\pgfpathlineto{\pgfqpoint{0.939319in}{1.463489in}}%
\pgfpathlineto{\pgfqpoint{0.942036in}{1.448130in}}%
\pgfpathlineto{\pgfqpoint{0.984575in}{1.451557in}}%
\pgfpathlineto{\pgfqpoint{1.027442in}{1.454039in}}%
\pgfpathclose%
\pgfusepath{fill}%
\end{pgfscope}%
\begin{pgfscope}%
\pgfpathrectangle{\pgfqpoint{0.407592in}{0.222904in}}{\pgfqpoint{1.411163in}{1.411163in}}%
\pgfusepath{clip}%
\pgfsetbuttcap%
\pgfsetroundjoin%
\definecolor{currentfill}{rgb}{0.749020,0.772549,0.792157}%
\pgfsetfillcolor{currentfill}%
\pgfsetfillopacity{0.160000}%
\pgfsetlinewidth{0.000000pt}%
\definecolor{currentstroke}{rgb}{0.000000,0.000000,0.000000}%
\pgfsetstrokecolor{currentstroke}%
\pgfsetdash{}{0pt}%
\pgfpathmoveto{\pgfqpoint{1.195313in}{1.423460in}}%
\pgfpathlineto{\pgfqpoint{1.197316in}{1.438955in}}%
\pgfpathlineto{\pgfqpoint{1.199315in}{1.454422in}}%
\pgfpathlineto{\pgfqpoint{1.156573in}{1.455760in}}%
\pgfpathlineto{\pgfqpoint{1.113584in}{1.456142in}}%
\pgfpathlineto{\pgfqpoint{1.113156in}{1.440654in}}%
\pgfpathlineto{\pgfqpoint{1.112728in}{1.425136in}}%
\pgfpathlineto{\pgfqpoint{1.154140in}{1.424763in}}%
\pgfpathlineto{\pgfqpoint{1.195313in}{1.423460in}}%
\pgfpathclose%
\pgfusepath{fill}%
\end{pgfscope}%
\begin{pgfscope}%
\pgfpathrectangle{\pgfqpoint{0.407592in}{0.222904in}}{\pgfqpoint{1.411163in}{1.411163in}}%
\pgfusepath{clip}%
\pgfsetbuttcap%
\pgfsetroundjoin%
\definecolor{currentfill}{rgb}{0.749020,0.772549,0.792157}%
\pgfsetfillcolor{currentfill}%
\pgfsetfillopacity{0.160000}%
\pgfsetlinewidth{0.000000pt}%
\definecolor{currentstroke}{rgb}{0.000000,0.000000,0.000000}%
\pgfsetstrokecolor{currentstroke}%
\pgfsetdash{}{0pt}%
\pgfpathmoveto{\pgfqpoint{1.112728in}{1.425136in}}%
\pgfpathlineto{\pgfqpoint{1.113156in}{1.440654in}}%
\pgfpathlineto{\pgfqpoint{1.113584in}{1.456142in}}%
\pgfpathlineto{\pgfqpoint{1.070493in}{1.455568in}}%
\pgfpathlineto{\pgfqpoint{1.027442in}{1.454039in}}%
\pgfpathlineto{\pgfqpoint{1.028593in}{1.438577in}}%
\pgfpathlineto{\pgfqpoint{1.029746in}{1.423087in}}%
\pgfpathlineto{\pgfqpoint{1.071217in}{1.424577in}}%
\pgfpathlineto{\pgfqpoint{1.112728in}{1.425136in}}%
\pgfpathclose%
\pgfusepath{fill}%
\end{pgfscope}%
\begin{pgfscope}%
\pgfpathrectangle{\pgfqpoint{0.407592in}{0.222904in}}{\pgfqpoint{1.411163in}{1.411163in}}%
\pgfusepath{clip}%
\pgfsetbuttcap%
\pgfsetroundjoin%
\definecolor{currentfill}{rgb}{0.749020,0.772549,0.792157}%
\pgfsetfillcolor{currentfill}%
\pgfsetfillopacity{0.160000}%
\pgfsetlinewidth{0.000000pt}%
\definecolor{currentstroke}{rgb}{0.000000,0.000000,0.000000}%
\pgfsetstrokecolor{currentstroke}%
\pgfsetdash{}{0pt}%
\pgfpathmoveto{\pgfqpoint{1.375079in}{1.470084in}}%
\pgfpathlineto{\pgfqpoint{1.380116in}{1.485285in}}%
\pgfpathlineto{\pgfqpoint{1.385145in}{1.500462in}}%
\pgfpathlineto{\pgfqpoint{1.341815in}{1.505822in}}%
\pgfpathlineto{\pgfqpoint{1.297627in}{1.510206in}}%
\pgfpathlineto{\pgfqpoint{1.294103in}{1.494918in}}%
\pgfpathlineto{\pgfqpoint{1.290572in}{1.479604in}}%
\pgfpathlineto{\pgfqpoint{1.333244in}{1.475321in}}%
\pgfpathlineto{\pgfqpoint{1.375079in}{1.470084in}}%
\pgfpathclose%
\pgfusepath{fill}%
\end{pgfscope}%
\begin{pgfscope}%
\pgfpathrectangle{\pgfqpoint{0.407592in}{0.222904in}}{\pgfqpoint{1.411163in}{1.411163in}}%
\pgfusepath{clip}%
\pgfsetbuttcap%
\pgfsetroundjoin%
\definecolor{currentfill}{rgb}{0.749020,0.772549,0.792157}%
\pgfsetfillcolor{currentfill}%
\pgfsetfillopacity{0.160000}%
\pgfsetlinewidth{0.000000pt}%
\definecolor{currentstroke}{rgb}{0.000000,0.000000,0.000000}%
\pgfsetstrokecolor{currentstroke}%
\pgfsetdash{}{0pt}%
\pgfpathmoveto{\pgfqpoint{0.936606in}{1.478821in}}%
\pgfpathlineto{\pgfqpoint{0.933898in}{1.494127in}}%
\pgfpathlineto{\pgfqpoint{0.931195in}{1.509405in}}%
\pgfpathlineto{\pgfqpoint{0.886026in}{1.504825in}}%
\pgfpathlineto{\pgfqpoint{0.841501in}{1.499270in}}%
\pgfpathlineto{\pgfqpoint{0.845745in}{1.484107in}}%
\pgfpathlineto{\pgfqpoint{0.849995in}{1.468920in}}%
\pgfpathlineto{\pgfqpoint{0.892985in}{1.474347in}}%
\pgfpathlineto{\pgfqpoint{0.936606in}{1.478821in}}%
\pgfpathclose%
\pgfusepath{fill}%
\end{pgfscope}%
\begin{pgfscope}%
\pgfpathrectangle{\pgfqpoint{0.407592in}{0.222904in}}{\pgfqpoint{1.411163in}{1.411163in}}%
\pgfusepath{clip}%
\pgfsetbuttcap%
\pgfsetroundjoin%
\definecolor{currentfill}{rgb}{0.749020,0.772549,0.792157}%
\pgfsetfillcolor{currentfill}%
\pgfsetfillopacity{0.160000}%
\pgfsetlinewidth{0.000000pt}%
\definecolor{currentstroke}{rgb}{0.000000,0.000000,0.000000}%
\pgfsetstrokecolor{currentstroke}%
\pgfsetdash{}{0pt}%
\pgfpathmoveto{\pgfqpoint{1.276388in}{1.418074in}}%
\pgfpathlineto{\pgfqpoint{1.279944in}{1.433497in}}%
\pgfpathlineto{\pgfqpoint{1.283493in}{1.448893in}}%
\pgfpathlineto{\pgfqpoint{1.241669in}{1.452131in}}%
\pgfpathlineto{\pgfqpoint{1.199315in}{1.454422in}}%
\pgfpathlineto{\pgfqpoint{1.197316in}{1.438955in}}%
\pgfpathlineto{\pgfqpoint{1.195313in}{1.423460in}}%
\pgfpathlineto{\pgfqpoint{1.236109in}{1.421228in}}%
\pgfpathlineto{\pgfqpoint{1.276388in}{1.418074in}}%
\pgfpathclose%
\pgfusepath{fill}%
\end{pgfscope}%
\begin{pgfscope}%
\pgfpathrectangle{\pgfqpoint{0.407592in}{0.222904in}}{\pgfqpoint{1.411163in}{1.411163in}}%
\pgfusepath{clip}%
\pgfsetbuttcap%
\pgfsetroundjoin%
\definecolor{currentfill}{rgb}{0.749020,0.772549,0.792157}%
\pgfsetfillcolor{currentfill}%
\pgfsetfillopacity{0.160000}%
\pgfsetlinewidth{0.000000pt}%
\definecolor{currentstroke}{rgb}{0.000000,0.000000,0.000000}%
\pgfsetstrokecolor{currentstroke}%
\pgfsetdash{}{0pt}%
\pgfpathmoveto{\pgfqpoint{1.029746in}{1.423087in}}%
\pgfpathlineto{\pgfqpoint{1.028593in}{1.438577in}}%
\pgfpathlineto{\pgfqpoint{1.027442in}{1.454039in}}%
\pgfpathlineto{\pgfqpoint{0.984575in}{1.451557in}}%
\pgfpathlineto{\pgfqpoint{0.942036in}{1.448130in}}%
\pgfpathlineto{\pgfqpoint{0.944758in}{1.432744in}}%
\pgfpathlineto{\pgfqpoint{0.947485in}{1.417331in}}%
\pgfpathlineto{\pgfqpoint{0.988456in}{1.420669in}}%
\pgfpathlineto{\pgfqpoint{1.029746in}{1.423087in}}%
\pgfpathclose%
\pgfusepath{fill}%
\end{pgfscope}%
\begin{pgfscope}%
\pgfpathrectangle{\pgfqpoint{0.407592in}{0.222904in}}{\pgfqpoint{1.411163in}{1.411163in}}%
\pgfusepath{clip}%
\pgfsetbuttcap%
\pgfsetroundjoin%
\definecolor{currentfill}{rgb}{0.749020,0.772549,0.792157}%
\pgfsetfillcolor{currentfill}%
\pgfsetfillopacity{0.160000}%
\pgfsetlinewidth{0.000000pt}%
\definecolor{currentstroke}{rgb}{0.000000,0.000000,0.000000}%
\pgfsetstrokecolor{currentstroke}%
\pgfsetdash{}{0pt}%
\pgfpathmoveto{\pgfqpoint{1.191296in}{1.392384in}}%
\pgfpathlineto{\pgfqpoint{1.193306in}{1.407936in}}%
\pgfpathlineto{\pgfqpoint{1.195313in}{1.423460in}}%
\pgfpathlineto{\pgfqpoint{1.154140in}{1.424763in}}%
\pgfpathlineto{\pgfqpoint{1.112728in}{1.425136in}}%
\pgfpathlineto{\pgfqpoint{1.112298in}{1.409589in}}%
\pgfpathlineto{\pgfqpoint{1.111868in}{1.394014in}}%
\pgfpathlineto{\pgfqpoint{1.151698in}{1.393652in}}%
\pgfpathlineto{\pgfqpoint{1.191296in}{1.392384in}}%
\pgfpathclose%
\pgfusepath{fill}%
\end{pgfscope}%
\begin{pgfscope}%
\pgfpathrectangle{\pgfqpoint{0.407592in}{0.222904in}}{\pgfqpoint{1.411163in}{1.411163in}}%
\pgfusepath{clip}%
\pgfsetbuttcap%
\pgfsetroundjoin%
\definecolor{currentfill}{rgb}{0.749020,0.772549,0.792157}%
\pgfsetfillcolor{currentfill}%
\pgfsetfillopacity{0.160000}%
\pgfsetlinewidth{0.000000pt}%
\definecolor{currentstroke}{rgb}{0.000000,0.000000,0.000000}%
\pgfsetstrokecolor{currentstroke}%
\pgfsetdash{}{0pt}%
\pgfpathmoveto{\pgfqpoint{1.111868in}{1.394014in}}%
\pgfpathlineto{\pgfqpoint{1.112298in}{1.409589in}}%
\pgfpathlineto{\pgfqpoint{1.112728in}{1.425136in}}%
\pgfpathlineto{\pgfqpoint{1.071217in}{1.424577in}}%
\pgfpathlineto{\pgfqpoint{1.029746in}{1.423087in}}%
\pgfpathlineto{\pgfqpoint{1.030901in}{1.407568in}}%
\pgfpathlineto{\pgfqpoint{1.032059in}{1.392021in}}%
\pgfpathlineto{\pgfqpoint{1.071943in}{1.393470in}}%
\pgfpathlineto{\pgfqpoint{1.111868in}{1.394014in}}%
\pgfpathclose%
\pgfusepath{fill}%
\end{pgfscope}%
\begin{pgfscope}%
\pgfpathrectangle{\pgfqpoint{0.407592in}{0.222904in}}{\pgfqpoint{1.411163in}{1.411163in}}%
\pgfusepath{clip}%
\pgfsetbuttcap%
\pgfsetroundjoin%
\definecolor{currentfill}{rgb}{0.749020,0.772549,0.792157}%
\pgfsetfillcolor{currentfill}%
\pgfsetfillopacity{0.160000}%
\pgfsetlinewidth{0.000000pt}%
\definecolor{currentstroke}{rgb}{0.000000,0.000000,0.000000}%
\pgfsetstrokecolor{currentstroke}%
\pgfsetdash{}{0pt}%
\pgfpathmoveto{\pgfqpoint{1.364979in}{1.439607in}}%
\pgfpathlineto{\pgfqpoint{1.370033in}{1.454858in}}%
\pgfpathlineto{\pgfqpoint{1.375079in}{1.470084in}}%
\pgfpathlineto{\pgfqpoint{1.333244in}{1.475321in}}%
\pgfpathlineto{\pgfqpoint{1.290572in}{1.479604in}}%
\pgfpathlineto{\pgfqpoint{1.287036in}{1.464262in}}%
\pgfpathlineto{\pgfqpoint{1.283493in}{1.448893in}}%
\pgfpathlineto{\pgfqpoint{1.324643in}{1.444715in}}%
\pgfpathlineto{\pgfqpoint{1.364979in}{1.439607in}}%
\pgfpathclose%
\pgfusepath{fill}%
\end{pgfscope}%
\begin{pgfscope}%
\pgfpathrectangle{\pgfqpoint{0.407592in}{0.222904in}}{\pgfqpoint{1.411163in}{1.411163in}}%
\pgfusepath{clip}%
\pgfsetbuttcap%
\pgfsetroundjoin%
\definecolor{currentfill}{rgb}{0.749020,0.772549,0.792157}%
\pgfsetfillcolor{currentfill}%
\pgfsetfillopacity{0.160000}%
\pgfsetlinewidth{0.000000pt}%
\definecolor{currentstroke}{rgb}{0.000000,0.000000,0.000000}%
\pgfsetstrokecolor{currentstroke}%
\pgfsetdash{}{0pt}%
\pgfpathmoveto{\pgfqpoint{0.942036in}{1.448130in}}%
\pgfpathlineto{\pgfqpoint{0.939319in}{1.463489in}}%
\pgfpathlineto{\pgfqpoint{0.936606in}{1.478821in}}%
\pgfpathlineto{\pgfqpoint{0.892985in}{1.474347in}}%
\pgfpathlineto{\pgfqpoint{0.849995in}{1.468920in}}%
\pgfpathlineto{\pgfqpoint{0.854253in}{1.453708in}}%
\pgfpathlineto{\pgfqpoint{0.858517in}{1.438471in}}%
\pgfpathlineto{\pgfqpoint{0.899968in}{1.443765in}}%
\pgfpathlineto{\pgfqpoint{0.942036in}{1.448130in}}%
\pgfpathclose%
\pgfusepath{fill}%
\end{pgfscope}%
\begin{pgfscope}%
\pgfpathrectangle{\pgfqpoint{0.407592in}{0.222904in}}{\pgfqpoint{1.411163in}{1.411163in}}%
\pgfusepath{clip}%
\pgfsetbuttcap%
\pgfsetroundjoin%
\definecolor{currentfill}{rgb}{0.749020,0.772549,0.792157}%
\pgfsetfillcolor{currentfill}%
\pgfsetfillopacity{0.160000}%
\pgfsetlinewidth{0.000000pt}%
\definecolor{currentstroke}{rgb}{0.000000,0.000000,0.000000}%
\pgfsetstrokecolor{currentstroke}%
\pgfsetdash{}{0pt}%
\pgfpathmoveto{\pgfqpoint{1.269258in}{1.387146in}}%
\pgfpathlineto{\pgfqpoint{1.272827in}{1.402624in}}%
\pgfpathlineto{\pgfqpoint{1.276388in}{1.418074in}}%
\pgfpathlineto{\pgfqpoint{1.236109in}{1.421228in}}%
\pgfpathlineto{\pgfqpoint{1.195313in}{1.423460in}}%
\pgfpathlineto{\pgfqpoint{1.193306in}{1.407936in}}%
\pgfpathlineto{\pgfqpoint{1.191296in}{1.392384in}}%
\pgfpathlineto{\pgfqpoint{1.230528in}{1.390213in}}%
\pgfpathlineto{\pgfqpoint{1.269258in}{1.387146in}}%
\pgfpathclose%
\pgfusepath{fill}%
\end{pgfscope}%
\begin{pgfscope}%
\pgfpathrectangle{\pgfqpoint{0.407592in}{0.222904in}}{\pgfqpoint{1.411163in}{1.411163in}}%
\pgfusepath{clip}%
\pgfsetbuttcap%
\pgfsetroundjoin%
\definecolor{currentfill}{rgb}{0.749020,0.772549,0.792157}%
\pgfsetfillcolor{currentfill}%
\pgfsetfillopacity{0.160000}%
\pgfsetlinewidth{0.000000pt}%
\definecolor{currentstroke}{rgb}{0.000000,0.000000,0.000000}%
\pgfsetstrokecolor{currentstroke}%
\pgfsetdash{}{0pt}%
\pgfpathmoveto{\pgfqpoint{1.032059in}{1.392021in}}%
\pgfpathlineto{\pgfqpoint{1.030901in}{1.407568in}}%
\pgfpathlineto{\pgfqpoint{1.029746in}{1.423087in}}%
\pgfpathlineto{\pgfqpoint{0.988456in}{1.420669in}}%
\pgfpathlineto{\pgfqpoint{0.947485in}{1.417331in}}%
\pgfpathlineto{\pgfqpoint{0.950217in}{1.401891in}}%
\pgfpathlineto{\pgfqpoint{0.952954in}{1.386423in}}%
\pgfpathlineto{\pgfqpoint{0.992350in}{1.389670in}}%
\pgfpathlineto{\pgfqpoint{1.032059in}{1.392021in}}%
\pgfpathclose%
\pgfusepath{fill}%
\end{pgfscope}%
\begin{pgfscope}%
\pgfpathrectangle{\pgfqpoint{0.407592in}{0.222904in}}{\pgfqpoint{1.411163in}{1.411163in}}%
\pgfusepath{clip}%
\pgfsetbuttcap%
\pgfsetroundjoin%
\definecolor{currentfill}{rgb}{0.749020,0.772549,0.792157}%
\pgfsetfillcolor{currentfill}%
\pgfsetfillopacity{0.160000}%
\pgfsetlinewidth{0.000000pt}%
\definecolor{currentstroke}{rgb}{0.000000,0.000000,0.000000}%
\pgfsetstrokecolor{currentstroke}%
\pgfsetdash{}{0pt}%
\pgfpathmoveto{\pgfqpoint{1.468638in}{1.486859in}}%
\pgfpathlineto{\pgfqpoint{1.475114in}{1.501856in}}%
\pgfpathlineto{\pgfqpoint{1.433225in}{1.509216in}}%
\pgfpathlineto{\pgfqpoint{1.390166in}{1.515613in}}%
\pgfpathlineto{\pgfqpoint{1.385145in}{1.500462in}}%
\pgfpathlineto{\pgfqpoint{1.427469in}{1.494136in}}%
\pgfusepath{fill}%
\end{pgfscope}%
\begin{pgfscope}%
\pgfpathrectangle{\pgfqpoint{0.407592in}{0.222904in}}{\pgfqpoint{1.411163in}{1.411163in}}%
\pgfusepath{clip}%
\pgfsetbuttcap%
\pgfsetroundjoin%
\definecolor{currentfill}{rgb}{0.749020,0.772549,0.792157}%
\pgfsetfillcolor{currentfill}%
\pgfsetfillopacity{0.160000}%
\pgfsetlinewidth{0.000000pt}%
\definecolor{currentstroke}{rgb}{0.000000,0.000000,0.000000}%
\pgfsetstrokecolor{currentstroke}%
\pgfsetdash{}{0pt}%
\pgfpathmoveto{\pgfqpoint{1.187264in}{1.361193in}}%
\pgfpathlineto{\pgfqpoint{1.189282in}{1.376803in}}%
\pgfpathlineto{\pgfqpoint{1.191296in}{1.392384in}}%
\pgfpathlineto{\pgfqpoint{1.151698in}{1.393652in}}%
\pgfpathlineto{\pgfqpoint{1.111868in}{1.394014in}}%
\pgfpathlineto{\pgfqpoint{1.111437in}{1.378409in}}%
\pgfpathlineto{\pgfqpoint{1.111005in}{1.362775in}}%
\pgfpathlineto{\pgfqpoint{1.149246in}{1.362423in}}%
\pgfpathlineto{\pgfqpoint{1.187264in}{1.361193in}}%
\pgfpathclose%
\pgfusepath{fill}%
\end{pgfscope}%
\begin{pgfscope}%
\pgfpathrectangle{\pgfqpoint{0.407592in}{0.222904in}}{\pgfqpoint{1.411163in}{1.411163in}}%
\pgfusepath{clip}%
\pgfsetbuttcap%
\pgfsetroundjoin%
\definecolor{currentfill}{rgb}{0.749020,0.772549,0.792157}%
\pgfsetfillcolor{currentfill}%
\pgfsetfillopacity{0.160000}%
\pgfsetlinewidth{0.000000pt}%
\definecolor{currentstroke}{rgb}{0.000000,0.000000,0.000000}%
\pgfsetstrokecolor{currentstroke}%
\pgfsetdash{}{0pt}%
\pgfpathmoveto{\pgfqpoint{1.111005in}{1.362775in}}%
\pgfpathlineto{\pgfqpoint{1.111437in}{1.378409in}}%
\pgfpathlineto{\pgfqpoint{1.111868in}{1.394014in}}%
\pgfpathlineto{\pgfqpoint{1.071943in}{1.393470in}}%
\pgfpathlineto{\pgfqpoint{1.032059in}{1.392021in}}%
\pgfpathlineto{\pgfqpoint{1.033218in}{1.376445in}}%
\pgfpathlineto{\pgfqpoint{1.034380in}{1.360840in}}%
\pgfpathlineto{\pgfqpoint{1.072673in}{1.362247in}}%
\pgfpathlineto{\pgfqpoint{1.111005in}{1.362775in}}%
\pgfpathclose%
\pgfusepath{fill}%
\end{pgfscope}%
\begin{pgfscope}%
\pgfpathrectangle{\pgfqpoint{0.407592in}{0.222904in}}{\pgfqpoint{1.411163in}{1.411163in}}%
\pgfusepath{clip}%
\pgfsetbuttcap%
\pgfsetroundjoin%
\definecolor{currentfill}{rgb}{0.749020,0.772549,0.792157}%
\pgfsetfillcolor{currentfill}%
\pgfsetfillopacity{0.160000}%
\pgfsetlinewidth{0.000000pt}%
\definecolor{currentstroke}{rgb}{0.000000,0.000000,0.000000}%
\pgfsetstrokecolor{currentstroke}%
\pgfsetdash{}{0pt}%
\pgfpathmoveto{\pgfqpoint{1.354846in}{1.409030in}}%
\pgfpathlineto{\pgfqpoint{1.359916in}{1.424331in}}%
\pgfpathlineto{\pgfqpoint{1.364979in}{1.439607in}}%
\pgfpathlineto{\pgfqpoint{1.324643in}{1.444715in}}%
\pgfpathlineto{\pgfqpoint{1.283493in}{1.448893in}}%
\pgfpathlineto{\pgfqpoint{1.279944in}{1.433497in}}%
\pgfpathlineto{\pgfqpoint{1.276388in}{1.418074in}}%
\pgfpathlineto{\pgfqpoint{1.316013in}{1.414005in}}%
\pgfpathlineto{\pgfqpoint{1.354846in}{1.409030in}}%
\pgfpathclose%
\pgfusepath{fill}%
\end{pgfscope}%
\begin{pgfscope}%
\pgfpathrectangle{\pgfqpoint{0.407592in}{0.222904in}}{\pgfqpoint{1.411163in}{1.411163in}}%
\pgfusepath{clip}%
\pgfsetbuttcap%
\pgfsetroundjoin%
\definecolor{currentfill}{rgb}{0.749020,0.772549,0.792157}%
\pgfsetfillcolor{currentfill}%
\pgfsetfillopacity{0.160000}%
\pgfsetlinewidth{0.000000pt}%
\definecolor{currentstroke}{rgb}{0.000000,0.000000,0.000000}%
\pgfsetstrokecolor{currentstroke}%
\pgfsetdash{}{0pt}%
\pgfpathmoveto{\pgfqpoint{0.841501in}{1.499270in}}%
\pgfpathlineto{\pgfqpoint{0.837265in}{1.514408in}}%
\pgfpathlineto{\pgfqpoint{0.792780in}{1.507816in}}%
\pgfpathlineto{\pgfqpoint{0.749254in}{1.500265in}}%
\pgfpathlineto{\pgfqpoint{0.754995in}{1.485286in}}%
\pgfpathlineto{\pgfqpoint{0.797774in}{1.492752in}}%
\pgfusepath{fill}%
\end{pgfscope}%
\begin{pgfscope}%
\pgfpathrectangle{\pgfqpoint{0.407592in}{0.222904in}}{\pgfqpoint{1.411163in}{1.411163in}}%
\pgfusepath{clip}%
\pgfsetbuttcap%
\pgfsetroundjoin%
\definecolor{currentfill}{rgb}{0.749020,0.772549,0.792157}%
\pgfsetfillcolor{currentfill}%
\pgfsetfillopacity{0.160000}%
\pgfsetlinewidth{0.000000pt}%
\definecolor{currentstroke}{rgb}{0.000000,0.000000,0.000000}%
\pgfsetstrokecolor{currentstroke}%
\pgfsetdash{}{0pt}%
\pgfpathmoveto{\pgfqpoint{0.947485in}{1.417331in}}%
\pgfpathlineto{\pgfqpoint{0.944758in}{1.432744in}}%
\pgfpathlineto{\pgfqpoint{0.942036in}{1.448130in}}%
\pgfpathlineto{\pgfqpoint{0.899968in}{1.443765in}}%
\pgfpathlineto{\pgfqpoint{0.858517in}{1.438471in}}%
\pgfpathlineto{\pgfqpoint{0.862788in}{1.423210in}}%
\pgfpathlineto{\pgfqpoint{0.867066in}{1.407924in}}%
\pgfpathlineto{\pgfqpoint{0.906975in}{1.413079in}}%
\pgfpathlineto{\pgfqpoint{0.947485in}{1.417331in}}%
\pgfpathclose%
\pgfusepath{fill}%
\end{pgfscope}%
\begin{pgfscope}%
\pgfpathrectangle{\pgfqpoint{0.407592in}{0.222904in}}{\pgfqpoint{1.411163in}{1.411163in}}%
\pgfusepath{clip}%
\pgfsetbuttcap%
\pgfsetroundjoin%
\definecolor{currentfill}{rgb}{0.749020,0.772549,0.792157}%
\pgfsetfillcolor{currentfill}%
\pgfsetfillopacity{0.160000}%
\pgfsetlinewidth{0.000000pt}%
\definecolor{currentstroke}{rgb}{0.000000,0.000000,0.000000}%
\pgfsetstrokecolor{currentstroke}%
\pgfsetdash{}{0pt}%
\pgfpathmoveto{\pgfqpoint{1.455658in}{1.456799in}}%
\pgfpathlineto{\pgfqpoint{1.462153in}{1.471840in}}%
\pgfpathlineto{\pgfqpoint{1.468638in}{1.486859in}}%
\pgfpathlineto{\pgfqpoint{1.427469in}{1.494136in}}%
\pgfpathlineto{\pgfqpoint{1.385145in}{1.500462in}}%
\pgfpathlineto{\pgfqpoint{1.380116in}{1.485285in}}%
\pgfpathlineto{\pgfqpoint{1.375079in}{1.470084in}}%
\pgfpathlineto{\pgfqpoint{1.415932in}{1.463906in}}%
\pgfpathlineto{\pgfqpoint{1.455658in}{1.456799in}}%
\pgfpathclose%
\pgfusepath{fill}%
\end{pgfscope}%
\begin{pgfscope}%
\pgfpathrectangle{\pgfqpoint{0.407592in}{0.222904in}}{\pgfqpoint{1.411163in}{1.411163in}}%
\pgfusepath{clip}%
\pgfsetbuttcap%
\pgfsetroundjoin%
\definecolor{currentfill}{rgb}{0.749020,0.772549,0.792157}%
\pgfsetfillcolor{currentfill}%
\pgfsetfillopacity{0.160000}%
\pgfsetlinewidth{0.000000pt}%
\definecolor{currentstroke}{rgb}{0.000000,0.000000,0.000000}%
\pgfsetstrokecolor{currentstroke}%
\pgfsetdash{}{0pt}%
\pgfpathmoveto{\pgfqpoint{1.262103in}{1.356108in}}%
\pgfpathlineto{\pgfqpoint{1.265684in}{1.371641in}}%
\pgfpathlineto{\pgfqpoint{1.269258in}{1.387146in}}%
\pgfpathlineto{\pgfqpoint{1.230528in}{1.390213in}}%
\pgfpathlineto{\pgfqpoint{1.191296in}{1.392384in}}%
\pgfpathlineto{\pgfqpoint{1.189282in}{1.376803in}}%
\pgfpathlineto{\pgfqpoint{1.187264in}{1.361193in}}%
\pgfpathlineto{\pgfqpoint{1.224927in}{1.359086in}}%
\pgfpathlineto{\pgfqpoint{1.262103in}{1.356108in}}%
\pgfpathclose%
\pgfusepath{fill}%
\end{pgfscope}%
\begin{pgfscope}%
\pgfpathrectangle{\pgfqpoint{0.407592in}{0.222904in}}{\pgfqpoint{1.411163in}{1.411163in}}%
\pgfusepath{clip}%
\pgfsetbuttcap%
\pgfsetroundjoin%
\definecolor{currentfill}{rgb}{0.749020,0.772549,0.792157}%
\pgfsetfillcolor{currentfill}%
\pgfsetfillopacity{0.160000}%
\pgfsetlinewidth{0.000000pt}%
\definecolor{currentstroke}{rgb}{0.000000,0.000000,0.000000}%
\pgfsetstrokecolor{currentstroke}%
\pgfsetdash{}{0pt}%
\pgfpathmoveto{\pgfqpoint{1.034380in}{1.360840in}}%
\pgfpathlineto{\pgfqpoint{1.033218in}{1.376445in}}%
\pgfpathlineto{\pgfqpoint{1.032059in}{1.392021in}}%
\pgfpathlineto{\pgfqpoint{0.992350in}{1.389670in}}%
\pgfpathlineto{\pgfqpoint{0.952954in}{1.386423in}}%
\pgfpathlineto{\pgfqpoint{0.955695in}{1.370928in}}%
\pgfpathlineto{\pgfqpoint{0.958441in}{1.355406in}}%
\pgfpathlineto{\pgfqpoint{0.996258in}{1.358558in}}%
\pgfpathlineto{\pgfqpoint{1.034380in}{1.360840in}}%
\pgfpathclose%
\pgfusepath{fill}%
\end{pgfscope}%
\begin{pgfscope}%
\pgfpathrectangle{\pgfqpoint{0.407592in}{0.222904in}}{\pgfqpoint{1.411163in}{1.411163in}}%
\pgfusepath{clip}%
\pgfsetbuttcap%
\pgfsetroundjoin%
\definecolor{currentfill}{rgb}{0.749020,0.772549,0.792157}%
\pgfsetfillcolor{currentfill}%
\pgfsetfillopacity{0.160000}%
\pgfsetlinewidth{0.000000pt}%
\definecolor{currentstroke}{rgb}{0.000000,0.000000,0.000000}%
\pgfsetstrokecolor{currentstroke}%
\pgfsetdash{}{0pt}%
\pgfpathmoveto{\pgfqpoint{0.849995in}{1.468920in}}%
\pgfpathlineto{\pgfqpoint{0.845745in}{1.484107in}}%
\pgfpathlineto{\pgfqpoint{0.841501in}{1.499270in}}%
\pgfpathlineto{\pgfqpoint{0.797774in}{1.492752in}}%
\pgfpathlineto{\pgfqpoint{0.754995in}{1.485286in}}%
\pgfpathlineto{\pgfqpoint{0.760743in}{1.470286in}}%
\pgfpathlineto{\pgfqpoint{0.766500in}{1.455264in}}%
\pgfpathlineto{\pgfqpoint{0.807784in}{1.462554in}}%
\pgfpathlineto{\pgfqpoint{0.849995in}{1.468920in}}%
\pgfpathclose%
\pgfusepath{fill}%
\end{pgfscope}%
\begin{pgfscope}%
\pgfpathrectangle{\pgfqpoint{0.407592in}{0.222904in}}{\pgfqpoint{1.411163in}{1.411163in}}%
\pgfusepath{clip}%
\pgfsetbuttcap%
\pgfsetroundjoin%
\definecolor{currentfill}{rgb}{0.749020,0.772549,0.792157}%
\pgfsetfillcolor{currentfill}%
\pgfsetfillopacity{0.160000}%
\pgfsetlinewidth{0.000000pt}%
\definecolor{currentstroke}{rgb}{0.000000,0.000000,0.000000}%
\pgfsetstrokecolor{currentstroke}%
\pgfsetdash{}{0pt}%
\pgfpathmoveto{\pgfqpoint{1.183217in}{1.329886in}}%
\pgfpathlineto{\pgfqpoint{1.185242in}{1.345554in}}%
\pgfpathlineto{\pgfqpoint{1.187264in}{1.361193in}}%
\pgfpathlineto{\pgfqpoint{1.149246in}{1.362423in}}%
\pgfpathlineto{\pgfqpoint{1.111005in}{1.362775in}}%
\pgfpathlineto{\pgfqpoint{1.110573in}{1.347112in}}%
\pgfpathlineto{\pgfqpoint{1.110139in}{1.331419in}}%
\pgfpathlineto{\pgfqpoint{1.146786in}{1.331078in}}%
\pgfpathlineto{\pgfqpoint{1.183217in}{1.329886in}}%
\pgfpathclose%
\pgfusepath{fill}%
\end{pgfscope}%
\begin{pgfscope}%
\pgfpathrectangle{\pgfqpoint{0.407592in}{0.222904in}}{\pgfqpoint{1.411163in}{1.411163in}}%
\pgfusepath{clip}%
\pgfsetbuttcap%
\pgfsetroundjoin%
\definecolor{currentfill}{rgb}{0.749020,0.772549,0.792157}%
\pgfsetfillcolor{currentfill}%
\pgfsetfillopacity{0.160000}%
\pgfsetlinewidth{0.000000pt}%
\definecolor{currentstroke}{rgb}{0.000000,0.000000,0.000000}%
\pgfsetstrokecolor{currentstroke}%
\pgfsetdash{}{0pt}%
\pgfpathmoveto{\pgfqpoint{1.110139in}{1.331419in}}%
\pgfpathlineto{\pgfqpoint{1.110573in}{1.347112in}}%
\pgfpathlineto{\pgfqpoint{1.111005in}{1.362775in}}%
\pgfpathlineto{\pgfqpoint{1.072673in}{1.362247in}}%
\pgfpathlineto{\pgfqpoint{1.034380in}{1.360840in}}%
\pgfpathlineto{\pgfqpoint{1.035543in}{1.345206in}}%
\pgfpathlineto{\pgfqpoint{1.036709in}{1.329544in}}%
\pgfpathlineto{\pgfqpoint{1.073405in}{1.330907in}}%
\pgfpathlineto{\pgfqpoint{1.110139in}{1.331419in}}%
\pgfpathclose%
\pgfusepath{fill}%
\end{pgfscope}%
\begin{pgfscope}%
\pgfpathrectangle{\pgfqpoint{0.407592in}{0.222904in}}{\pgfqpoint{1.411163in}{1.411163in}}%
\pgfusepath{clip}%
\pgfsetbuttcap%
\pgfsetroundjoin%
\definecolor{currentfill}{rgb}{0.749020,0.772549,0.792157}%
\pgfsetfillcolor{currentfill}%
\pgfsetfillopacity{0.160000}%
\pgfsetlinewidth{0.000000pt}%
\definecolor{currentstroke}{rgb}{0.000000,0.000000,0.000000}%
\pgfsetstrokecolor{currentstroke}%
\pgfsetdash{}{0pt}%
\pgfpathmoveto{\pgfqpoint{1.344679in}{1.378352in}}%
\pgfpathlineto{\pgfqpoint{1.349767in}{1.393704in}}%
\pgfpathlineto{\pgfqpoint{1.354846in}{1.409030in}}%
\pgfpathlineto{\pgfqpoint{1.316013in}{1.414005in}}%
\pgfpathlineto{\pgfqpoint{1.276388in}{1.418074in}}%
\pgfpathlineto{\pgfqpoint{1.272827in}{1.402624in}}%
\pgfpathlineto{\pgfqpoint{1.269258in}{1.387146in}}%
\pgfpathlineto{\pgfqpoint{1.307354in}{1.383189in}}%
\pgfpathlineto{\pgfqpoint{1.344679in}{1.378352in}}%
\pgfpathclose%
\pgfusepath{fill}%
\end{pgfscope}%
\begin{pgfscope}%
\pgfpathrectangle{\pgfqpoint{0.407592in}{0.222904in}}{\pgfqpoint{1.411163in}{1.411163in}}%
\pgfusepath{clip}%
\pgfsetbuttcap%
\pgfsetroundjoin%
\definecolor{currentfill}{rgb}{0.749020,0.772549,0.792157}%
\pgfsetfillcolor{currentfill}%
\pgfsetfillopacity{0.160000}%
\pgfsetlinewidth{0.000000pt}%
\definecolor{currentstroke}{rgb}{0.000000,0.000000,0.000000}%
\pgfsetstrokecolor{currentstroke}%
\pgfsetdash{}{0pt}%
\pgfpathmoveto{\pgfqpoint{0.952954in}{1.386423in}}%
\pgfpathlineto{\pgfqpoint{0.950217in}{1.401891in}}%
\pgfpathlineto{\pgfqpoint{0.947485in}{1.417331in}}%
\pgfpathlineto{\pgfqpoint{0.906975in}{1.413079in}}%
\pgfpathlineto{\pgfqpoint{0.867066in}{1.407924in}}%
\pgfpathlineto{\pgfqpoint{0.871351in}{1.392613in}}%
\pgfpathlineto{\pgfqpoint{0.875643in}{1.377277in}}%
\pgfpathlineto{\pgfqpoint{0.914006in}{1.382289in}}%
\pgfpathlineto{\pgfqpoint{0.952954in}{1.386423in}}%
\pgfpathclose%
\pgfusepath{fill}%
\end{pgfscope}%
\begin{pgfscope}%
\pgfpathrectangle{\pgfqpoint{0.407592in}{0.222904in}}{\pgfqpoint{1.411163in}{1.411163in}}%
\pgfusepath{clip}%
\pgfsetbuttcap%
\pgfsetroundjoin%
\definecolor{currentfill}{rgb}{0.749020,0.772549,0.792157}%
\pgfsetfillcolor{currentfill}%
\pgfsetfillopacity{0.160000}%
\pgfsetlinewidth{0.000000pt}%
\definecolor{currentstroke}{rgb}{0.000000,0.000000,0.000000}%
\pgfsetstrokecolor{currentstroke}%
\pgfsetdash{}{0pt}%
\pgfpathmoveto{\pgfqpoint{1.442640in}{1.426652in}}%
\pgfpathlineto{\pgfqpoint{1.449154in}{1.441737in}}%
\pgfpathlineto{\pgfqpoint{1.455658in}{1.456799in}}%
\pgfpathlineto{\pgfqpoint{1.415932in}{1.463906in}}%
\pgfpathlineto{\pgfqpoint{1.375079in}{1.470084in}}%
\pgfpathlineto{\pgfqpoint{1.370033in}{1.454858in}}%
\pgfpathlineto{\pgfqpoint{1.364979in}{1.439607in}}%
\pgfpathlineto{\pgfqpoint{1.404358in}{1.433581in}}%
\pgfpathlineto{\pgfqpoint{1.442640in}{1.426652in}}%
\pgfpathclose%
\pgfusepath{fill}%
\end{pgfscope}%
\begin{pgfscope}%
\pgfpathrectangle{\pgfqpoint{0.407592in}{0.222904in}}{\pgfqpoint{1.411163in}{1.411163in}}%
\pgfusepath{clip}%
\pgfsetbuttcap%
\pgfsetroundjoin%
\definecolor{currentfill}{rgb}{0.749020,0.772549,0.792157}%
\pgfsetfillcolor{currentfill}%
\pgfsetfillopacity{0.160000}%
\pgfsetlinewidth{0.000000pt}%
\definecolor{currentstroke}{rgb}{0.000000,0.000000,0.000000}%
\pgfsetstrokecolor{currentstroke}%
\pgfsetdash{}{0pt}%
\pgfpathmoveto{\pgfqpoint{1.254923in}{1.324960in}}%
\pgfpathlineto{\pgfqpoint{1.258516in}{1.340548in}}%
\pgfpathlineto{\pgfqpoint{1.262103in}{1.356108in}}%
\pgfpathlineto{\pgfqpoint{1.224927in}{1.359086in}}%
\pgfpathlineto{\pgfqpoint{1.187264in}{1.361193in}}%
\pgfpathlineto{\pgfqpoint{1.185242in}{1.345554in}}%
\pgfpathlineto{\pgfqpoint{1.183217in}{1.329886in}}%
\pgfpathlineto{\pgfqpoint{1.219305in}{1.327844in}}%
\pgfpathlineto{\pgfqpoint{1.254923in}{1.324960in}}%
\pgfpathclose%
\pgfusepath{fill}%
\end{pgfscope}%
\begin{pgfscope}%
\pgfpathrectangle{\pgfqpoint{0.407592in}{0.222904in}}{\pgfqpoint{1.411163in}{1.411163in}}%
\pgfusepath{clip}%
\pgfsetbuttcap%
\pgfsetroundjoin%
\definecolor{currentfill}{rgb}{0.749020,0.772549,0.792157}%
\pgfsetfillcolor{currentfill}%
\pgfsetfillopacity{0.160000}%
\pgfsetlinewidth{0.000000pt}%
\definecolor{currentstroke}{rgb}{0.000000,0.000000,0.000000}%
\pgfsetstrokecolor{currentstroke}%
\pgfsetdash{}{0pt}%
\pgfpathmoveto{\pgfqpoint{0.858517in}{1.438471in}}%
\pgfpathlineto{\pgfqpoint{0.854253in}{1.453708in}}%
\pgfpathlineto{\pgfqpoint{0.849995in}{1.468920in}}%
\pgfpathlineto{\pgfqpoint{0.807784in}{1.462554in}}%
\pgfpathlineto{\pgfqpoint{0.766500in}{1.455264in}}%
\pgfpathlineto{\pgfqpoint{0.772266in}{1.440220in}}%
\pgfpathlineto{\pgfqpoint{0.778039in}{1.425154in}}%
\pgfpathlineto{\pgfqpoint{0.817826in}{1.432263in}}%
\pgfpathlineto{\pgfqpoint{0.858517in}{1.438471in}}%
\pgfpathclose%
\pgfusepath{fill}%
\end{pgfscope}%
\begin{pgfscope}%
\pgfpathrectangle{\pgfqpoint{0.407592in}{0.222904in}}{\pgfqpoint{1.411163in}{1.411163in}}%
\pgfusepath{clip}%
\pgfsetbuttcap%
\pgfsetroundjoin%
\definecolor{currentfill}{rgb}{0.749020,0.772549,0.792157}%
\pgfsetfillcolor{currentfill}%
\pgfsetfillopacity{0.160000}%
\pgfsetlinewidth{0.000000pt}%
\definecolor{currentstroke}{rgb}{0.000000,0.000000,0.000000}%
\pgfsetstrokecolor{currentstroke}%
\pgfsetdash{}{0pt}%
\pgfpathmoveto{\pgfqpoint{1.036709in}{1.329544in}}%
\pgfpathlineto{\pgfqpoint{1.035543in}{1.345206in}}%
\pgfpathlineto{\pgfqpoint{1.034380in}{1.360840in}}%
\pgfpathlineto{\pgfqpoint{0.996258in}{1.358558in}}%
\pgfpathlineto{\pgfqpoint{0.958441in}{1.355406in}}%
\pgfpathlineto{\pgfqpoint{0.961192in}{1.339857in}}%
\pgfpathlineto{\pgfqpoint{0.963948in}{1.324280in}}%
\pgfpathlineto{\pgfqpoint{1.000181in}{1.327333in}}%
\pgfpathlineto{\pgfqpoint{1.036709in}{1.329544in}}%
\pgfpathclose%
\pgfusepath{fill}%
\end{pgfscope}%
\begin{pgfscope}%
\pgfpathrectangle{\pgfqpoint{0.407592in}{0.222904in}}{\pgfqpoint{1.411163in}{1.411163in}}%
\pgfusepath{clip}%
\pgfsetbuttcap%
\pgfsetroundjoin%
\definecolor{currentfill}{rgb}{0.749020,0.772549,0.792157}%
\pgfsetfillcolor{currentfill}%
\pgfsetfillopacity{0.160000}%
\pgfsetlinewidth{0.000000pt}%
\definecolor{currentstroke}{rgb}{0.000000,0.000000,0.000000}%
\pgfsetstrokecolor{currentstroke}%
\pgfsetdash{}{0pt}%
\pgfpathmoveto{\pgfqpoint{1.334479in}{1.347573in}}%
\pgfpathlineto{\pgfqpoint{1.339583in}{1.362975in}}%
\pgfpathlineto{\pgfqpoint{1.344679in}{1.378352in}}%
\pgfpathlineto{\pgfqpoint{1.307354in}{1.383189in}}%
\pgfpathlineto{\pgfqpoint{1.269258in}{1.387146in}}%
\pgfpathlineto{\pgfqpoint{1.265684in}{1.371641in}}%
\pgfpathlineto{\pgfqpoint{1.262103in}{1.356108in}}%
\pgfpathlineto{\pgfqpoint{1.298664in}{1.352267in}}%
\pgfpathlineto{\pgfqpoint{1.334479in}{1.347573in}}%
\pgfpathclose%
\pgfusepath{fill}%
\end{pgfscope}%
\begin{pgfscope}%
\pgfpathrectangle{\pgfqpoint{0.407592in}{0.222904in}}{\pgfqpoint{1.411163in}{1.411163in}}%
\pgfusepath{clip}%
\pgfsetbuttcap%
\pgfsetroundjoin%
\definecolor{currentfill}{rgb}{0.749020,0.772549,0.792157}%
\pgfsetfillcolor{currentfill}%
\pgfsetfillopacity{0.160000}%
\pgfsetlinewidth{0.000000pt}%
\definecolor{currentstroke}{rgb}{0.000000,0.000000,0.000000}%
\pgfsetstrokecolor{currentstroke}%
\pgfsetdash{}{0pt}%
\pgfpathmoveto{\pgfqpoint{1.179155in}{1.298463in}}%
\pgfpathlineto{\pgfqpoint{1.181188in}{1.314189in}}%
\pgfpathlineto{\pgfqpoint{1.183217in}{1.329886in}}%
\pgfpathlineto{\pgfqpoint{1.146786in}{1.331078in}}%
\pgfpathlineto{\pgfqpoint{1.110139in}{1.331419in}}%
\pgfpathlineto{\pgfqpoint{1.109705in}{1.315697in}}%
\pgfpathlineto{\pgfqpoint{1.109270in}{1.299945in}}%
\pgfpathlineto{\pgfqpoint{1.144317in}{1.299616in}}%
\pgfpathlineto{\pgfqpoint{1.179155in}{1.298463in}}%
\pgfpathclose%
\pgfusepath{fill}%
\end{pgfscope}%
\begin{pgfscope}%
\pgfpathrectangle{\pgfqpoint{0.407592in}{0.222904in}}{\pgfqpoint{1.411163in}{1.411163in}}%
\pgfusepath{clip}%
\pgfsetbuttcap%
\pgfsetroundjoin%
\definecolor{currentfill}{rgb}{0.749020,0.772549,0.792157}%
\pgfsetfillcolor{currentfill}%
\pgfsetfillopacity{0.160000}%
\pgfsetlinewidth{0.000000pt}%
\definecolor{currentstroke}{rgb}{0.000000,0.000000,0.000000}%
\pgfsetstrokecolor{currentstroke}%
\pgfsetdash{}{0pt}%
\pgfpathmoveto{\pgfqpoint{1.109270in}{1.299945in}}%
\pgfpathlineto{\pgfqpoint{1.109705in}{1.315697in}}%
\pgfpathlineto{\pgfqpoint{1.110139in}{1.331419in}}%
\pgfpathlineto{\pgfqpoint{1.073405in}{1.330907in}}%
\pgfpathlineto{\pgfqpoint{1.036709in}{1.329544in}}%
\pgfpathlineto{\pgfqpoint{1.037877in}{1.313853in}}%
\pgfpathlineto{\pgfqpoint{1.039047in}{1.298132in}}%
\pgfpathlineto{\pgfqpoint{1.074139in}{1.299451in}}%
\pgfpathlineto{\pgfqpoint{1.109270in}{1.299945in}}%
\pgfpathclose%
\pgfusepath{fill}%
\end{pgfscope}%
\begin{pgfscope}%
\pgfpathrectangle{\pgfqpoint{0.407592in}{0.222904in}}{\pgfqpoint{1.411163in}{1.411163in}}%
\pgfusepath{clip}%
\pgfsetbuttcap%
\pgfsetroundjoin%
\definecolor{currentfill}{rgb}{0.749020,0.772549,0.792157}%
\pgfsetfillcolor{currentfill}%
\pgfsetfillopacity{0.160000}%
\pgfsetlinewidth{0.000000pt}%
\definecolor{currentstroke}{rgb}{0.000000,0.000000,0.000000}%
\pgfsetstrokecolor{currentstroke}%
\pgfsetdash{}{0pt}%
\pgfpathmoveto{\pgfqpoint{0.958441in}{1.355406in}}%
\pgfpathlineto{\pgfqpoint{0.955695in}{1.370928in}}%
\pgfpathlineto{\pgfqpoint{0.952954in}{1.386423in}}%
\pgfpathlineto{\pgfqpoint{0.914006in}{1.382289in}}%
\pgfpathlineto{\pgfqpoint{0.875643in}{1.377277in}}%
\pgfpathlineto{\pgfqpoint{0.879942in}{1.361916in}}%
\pgfpathlineto{\pgfqpoint{0.884248in}{1.346529in}}%
\pgfpathlineto{\pgfqpoint{0.921061in}{1.351393in}}%
\pgfpathlineto{\pgfqpoint{0.958441in}{1.355406in}}%
\pgfpathclose%
\pgfusepath{fill}%
\end{pgfscope}%
\begin{pgfscope}%
\pgfpathrectangle{\pgfqpoint{0.407592in}{0.222904in}}{\pgfqpoint{1.411163in}{1.411163in}}%
\pgfusepath{clip}%
\pgfsetbuttcap%
\pgfsetroundjoin%
\definecolor{currentfill}{rgb}{0.749020,0.772549,0.792157}%
\pgfsetfillcolor{currentfill}%
\pgfsetfillopacity{0.160000}%
\pgfsetlinewidth{0.000000pt}%
\definecolor{currentstroke}{rgb}{0.000000,0.000000,0.000000}%
\pgfsetstrokecolor{currentstroke}%
\pgfsetdash{}{0pt}%
\pgfpathmoveto{\pgfqpoint{1.429584in}{1.396416in}}%
\pgfpathlineto{\pgfqpoint{1.436117in}{1.411545in}}%
\pgfpathlineto{\pgfqpoint{1.442640in}{1.426652in}}%
\pgfpathlineto{\pgfqpoint{1.404358in}{1.433581in}}%
\pgfpathlineto{\pgfqpoint{1.364979in}{1.439607in}}%
\pgfpathlineto{\pgfqpoint{1.359916in}{1.424331in}}%
\pgfpathlineto{\pgfqpoint{1.354846in}{1.409030in}}%
\pgfpathlineto{\pgfqpoint{1.392748in}{1.403162in}}%
\pgfpathlineto{\pgfqpoint{1.429584in}{1.396416in}}%
\pgfpathclose%
\pgfusepath{fill}%
\end{pgfscope}%
\begin{pgfscope}%
\pgfpathrectangle{\pgfqpoint{0.407592in}{0.222904in}}{\pgfqpoint{1.411163in}{1.411163in}}%
\pgfusepath{clip}%
\pgfsetbuttcap%
\pgfsetroundjoin%
\definecolor{currentfill}{rgb}{0.749020,0.772549,0.792157}%
\pgfsetfillcolor{currentfill}%
\pgfsetfillopacity{0.160000}%
\pgfsetlinewidth{0.000000pt}%
\definecolor{currentstroke}{rgb}{0.000000,0.000000,0.000000}%
\pgfsetstrokecolor{currentstroke}%
\pgfsetdash{}{0pt}%
\pgfpathmoveto{\pgfqpoint{0.867066in}{1.407924in}}%
\pgfpathlineto{\pgfqpoint{0.862788in}{1.423210in}}%
\pgfpathlineto{\pgfqpoint{0.858517in}{1.438471in}}%
\pgfpathlineto{\pgfqpoint{0.817826in}{1.432263in}}%
\pgfpathlineto{\pgfqpoint{0.778039in}{1.425154in}}%
\pgfpathlineto{\pgfqpoint{0.783821in}{1.410067in}}%
\pgfpathlineto{\pgfqpoint{0.789612in}{1.394958in}}%
\pgfpathlineto{\pgfqpoint{0.827898in}{1.401879in}}%
\pgfpathlineto{\pgfqpoint{0.867066in}{1.407924in}}%
\pgfpathclose%
\pgfusepath{fill}%
\end{pgfscope}%
\begin{pgfscope}%
\pgfpathrectangle{\pgfqpoint{0.407592in}{0.222904in}}{\pgfqpoint{1.411163in}{1.411163in}}%
\pgfusepath{clip}%
\pgfsetbuttcap%
\pgfsetroundjoin%
\definecolor{currentfill}{rgb}{0.749020,0.772549,0.792157}%
\pgfsetfillcolor{currentfill}%
\pgfsetfillopacity{0.160000}%
\pgfsetlinewidth{0.000000pt}%
\definecolor{currentstroke}{rgb}{0.000000,0.000000,0.000000}%
\pgfsetstrokecolor{currentstroke}%
\pgfsetdash{}{0pt}%
\pgfpathmoveto{\pgfqpoint{1.247717in}{1.293700in}}%
\pgfpathlineto{\pgfqpoint{1.251323in}{1.309344in}}%
\pgfpathlineto{\pgfqpoint{1.254923in}{1.324960in}}%
\pgfpathlineto{\pgfqpoint{1.219305in}{1.327844in}}%
\pgfpathlineto{\pgfqpoint{1.183217in}{1.329886in}}%
\pgfpathlineto{\pgfqpoint{1.181188in}{1.314189in}}%
\pgfpathlineto{\pgfqpoint{1.179155in}{1.298463in}}%
\pgfpathlineto{\pgfqpoint{1.213663in}{1.296489in}}%
\pgfpathlineto{\pgfqpoint{1.247717in}{1.293700in}}%
\pgfpathclose%
\pgfusepath{fill}%
\end{pgfscope}%
\begin{pgfscope}%
\pgfpathrectangle{\pgfqpoint{0.407592in}{0.222904in}}{\pgfqpoint{1.411163in}{1.411163in}}%
\pgfusepath{clip}%
\pgfsetbuttcap%
\pgfsetroundjoin%
\definecolor{currentfill}{rgb}{0.749020,0.772549,0.792157}%
\pgfsetfillcolor{currentfill}%
\pgfsetfillopacity{0.160000}%
\pgfsetlinewidth{0.000000pt}%
\definecolor{currentstroke}{rgb}{0.000000,0.000000,0.000000}%
\pgfsetstrokecolor{currentstroke}%
\pgfsetdash{}{0pt}%
\pgfpathmoveto{\pgfqpoint{1.039047in}{1.298132in}}%
\pgfpathlineto{\pgfqpoint{1.037877in}{1.313853in}}%
\pgfpathlineto{\pgfqpoint{1.036709in}{1.329544in}}%
\pgfpathlineto{\pgfqpoint{1.000181in}{1.327333in}}%
\pgfpathlineto{\pgfqpoint{0.963948in}{1.324280in}}%
\pgfpathlineto{\pgfqpoint{0.966709in}{1.308675in}}%
\pgfpathlineto{\pgfqpoint{0.969475in}{1.293043in}}%
\pgfpathlineto{\pgfqpoint{1.004118in}{1.295995in}}%
\pgfpathlineto{\pgfqpoint{1.039047in}{1.298132in}}%
\pgfpathclose%
\pgfusepath{fill}%
\end{pgfscope}%
\begin{pgfscope}%
\pgfpathrectangle{\pgfqpoint{0.407592in}{0.222904in}}{\pgfqpoint{1.411163in}{1.411163in}}%
\pgfusepath{clip}%
\pgfsetbuttcap%
\pgfsetroundjoin%
\definecolor{currentfill}{rgb}{0.749020,0.772549,0.792157}%
\pgfsetfillcolor{currentfill}%
\pgfsetfillopacity{0.160000}%
\pgfsetlinewidth{0.000000pt}%
\definecolor{currentstroke}{rgb}{0.000000,0.000000,0.000000}%
\pgfsetstrokecolor{currentstroke}%
\pgfsetdash{}{0pt}%
\pgfpathmoveto{\pgfqpoint{1.324246in}{1.316692in}}%
\pgfpathlineto{\pgfqpoint{1.329367in}{1.332145in}}%
\pgfpathlineto{\pgfqpoint{1.334479in}{1.347573in}}%
\pgfpathlineto{\pgfqpoint{1.298664in}{1.352267in}}%
\pgfpathlineto{\pgfqpoint{1.262103in}{1.356108in}}%
\pgfpathlineto{\pgfqpoint{1.258516in}{1.340548in}}%
\pgfpathlineto{\pgfqpoint{1.254923in}{1.324960in}}%
\pgfpathlineto{\pgfqpoint{1.289945in}{1.321239in}}%
\pgfpathlineto{\pgfqpoint{1.324246in}{1.316692in}}%
\pgfpathclose%
\pgfusepath{fill}%
\end{pgfscope}%
\begin{pgfscope}%
\pgfpathrectangle{\pgfqpoint{0.407592in}{0.222904in}}{\pgfqpoint{1.411163in}{1.411163in}}%
\pgfusepath{clip}%
\pgfsetbuttcap%
\pgfsetroundjoin%
\definecolor{currentfill}{rgb}{0.749020,0.772549,0.792157}%
\pgfsetfillcolor{currentfill}%
\pgfsetfillopacity{0.160000}%
\pgfsetlinewidth{0.000000pt}%
\definecolor{currentstroke}{rgb}{0.000000,0.000000,0.000000}%
\pgfsetstrokecolor{currentstroke}%
\pgfsetdash{}{0pt}%
\pgfpathmoveto{\pgfqpoint{1.175078in}{1.266923in}}%
\pgfpathlineto{\pgfqpoint{1.177118in}{1.282707in}}%
\pgfpathlineto{\pgfqpoint{1.179155in}{1.298463in}}%
\pgfpathlineto{\pgfqpoint{1.144317in}{1.299616in}}%
\pgfpathlineto{\pgfqpoint{1.109270in}{1.299945in}}%
\pgfpathlineto{\pgfqpoint{1.108834in}{1.284164in}}%
\pgfpathlineto{\pgfqpoint{1.108398in}{1.268353in}}%
\pgfpathlineto{\pgfqpoint{1.141838in}{1.268035in}}%
\pgfpathlineto{\pgfqpoint{1.175078in}{1.266923in}}%
\pgfpathclose%
\pgfusepath{fill}%
\end{pgfscope}%
\begin{pgfscope}%
\pgfpathrectangle{\pgfqpoint{0.407592in}{0.222904in}}{\pgfqpoint{1.411163in}{1.411163in}}%
\pgfusepath{clip}%
\pgfsetbuttcap%
\pgfsetroundjoin%
\definecolor{currentfill}{rgb}{0.749020,0.772549,0.792157}%
\pgfsetfillcolor{currentfill}%
\pgfsetfillopacity{0.160000}%
\pgfsetlinewidth{0.000000pt}%
\definecolor{currentstroke}{rgb}{0.000000,0.000000,0.000000}%
\pgfsetstrokecolor{currentstroke}%
\pgfsetdash{}{0pt}%
\pgfpathmoveto{\pgfqpoint{1.108398in}{1.268353in}}%
\pgfpathlineto{\pgfqpoint{1.108834in}{1.284164in}}%
\pgfpathlineto{\pgfqpoint{1.109270in}{1.299945in}}%
\pgfpathlineto{\pgfqpoint{1.074139in}{1.299451in}}%
\pgfpathlineto{\pgfqpoint{1.039047in}{1.298132in}}%
\pgfpathlineto{\pgfqpoint{1.040220in}{1.282383in}}%
\pgfpathlineto{\pgfqpoint{1.041394in}{1.266604in}}%
\pgfpathlineto{\pgfqpoint{1.074877in}{1.267876in}}%
\pgfpathlineto{\pgfqpoint{1.108398in}{1.268353in}}%
\pgfpathclose%
\pgfusepath{fill}%
\end{pgfscope}%
\begin{pgfscope}%
\pgfpathrectangle{\pgfqpoint{0.407592in}{0.222904in}}{\pgfqpoint{1.411163in}{1.411163in}}%
\pgfusepath{clip}%
\pgfsetbuttcap%
\pgfsetroundjoin%
\definecolor{currentfill}{rgb}{0.749020,0.772549,0.792157}%
\pgfsetfillcolor{currentfill}%
\pgfsetfillopacity{0.160000}%
\pgfsetlinewidth{0.000000pt}%
\definecolor{currentstroke}{rgb}{0.000000,0.000000,0.000000}%
\pgfsetstrokecolor{currentstroke}%
\pgfsetdash{}{0pt}%
\pgfpathmoveto{\pgfqpoint{0.963948in}{1.324280in}}%
\pgfpathlineto{\pgfqpoint{0.961192in}{1.339857in}}%
\pgfpathlineto{\pgfqpoint{0.958441in}{1.355406in}}%
\pgfpathlineto{\pgfqpoint{0.921061in}{1.351393in}}%
\pgfpathlineto{\pgfqpoint{0.884248in}{1.346529in}}%
\pgfpathlineto{\pgfqpoint{0.888561in}{1.331118in}}%
\pgfpathlineto{\pgfqpoint{0.892882in}{1.315681in}}%
\pgfpathlineto{\pgfqpoint{0.928139in}{1.320393in}}%
\pgfpathlineto{\pgfqpoint{0.963948in}{1.324280in}}%
\pgfpathclose%
\pgfusepath{fill}%
\end{pgfscope}%
\begin{pgfscope}%
\pgfpathrectangle{\pgfqpoint{0.407592in}{0.222904in}}{\pgfqpoint{1.411163in}{1.411163in}}%
\pgfusepath{clip}%
\pgfsetbuttcap%
\pgfsetroundjoin%
\definecolor{currentfill}{rgb}{0.749020,0.772549,0.792157}%
\pgfsetfillcolor{currentfill}%
\pgfsetfillopacity{0.160000}%
\pgfsetlinewidth{0.000000pt}%
\definecolor{currentstroke}{rgb}{0.000000,0.000000,0.000000}%
\pgfsetstrokecolor{currentstroke}%
\pgfsetdash{}{0pt}%
\pgfpathmoveto{\pgfqpoint{1.546915in}{1.469523in}}%
\pgfpathlineto{\pgfqpoint{1.554781in}{1.484319in}}%
\pgfpathlineto{\pgfqpoint{1.515683in}{1.493550in}}%
\pgfpathlineto{\pgfqpoint{1.475114in}{1.501856in}}%
\pgfpathlineto{\pgfqpoint{1.468638in}{1.486859in}}%
\pgfpathlineto{\pgfqpoint{1.508503in}{1.478648in}}%
\pgfusepath{fill}%
\end{pgfscope}%
\begin{pgfscope}%
\pgfpathrectangle{\pgfqpoint{0.407592in}{0.222904in}}{\pgfqpoint{1.411163in}{1.411163in}}%
\pgfusepath{clip}%
\pgfsetbuttcap%
\pgfsetroundjoin%
\definecolor{currentfill}{rgb}{0.749020,0.772549,0.792157}%
\pgfsetfillcolor{currentfill}%
\pgfsetfillopacity{0.160000}%
\pgfsetlinewidth{0.000000pt}%
\definecolor{currentstroke}{rgb}{0.000000,0.000000,0.000000}%
\pgfsetstrokecolor{currentstroke}%
\pgfsetdash{}{0pt}%
\pgfpathmoveto{\pgfqpoint{1.416489in}{1.366090in}}%
\pgfpathlineto{\pgfqpoint{1.423041in}{1.381264in}}%
\pgfpathlineto{\pgfqpoint{1.429584in}{1.396416in}}%
\pgfpathlineto{\pgfqpoint{1.392748in}{1.403162in}}%
\pgfpathlineto{\pgfqpoint{1.354846in}{1.409030in}}%
\pgfpathlineto{\pgfqpoint{1.349767in}{1.393704in}}%
\pgfpathlineto{\pgfqpoint{1.344679in}{1.378352in}}%
\pgfpathlineto{\pgfqpoint{1.381102in}{1.372648in}}%
\pgfpathlineto{\pgfqpoint{1.416489in}{1.366090in}}%
\pgfpathclose%
\pgfusepath{fill}%
\end{pgfscope}%
\begin{pgfscope}%
\pgfpathrectangle{\pgfqpoint{0.407592in}{0.222904in}}{\pgfqpoint{1.411163in}{1.411163in}}%
\pgfusepath{clip}%
\pgfsetbuttcap%
\pgfsetroundjoin%
\definecolor{currentfill}{rgb}{0.749020,0.772549,0.792157}%
\pgfsetfillcolor{currentfill}%
\pgfsetfillopacity{0.160000}%
\pgfsetlinewidth{0.000000pt}%
\definecolor{currentstroke}{rgb}{0.000000,0.000000,0.000000}%
\pgfsetstrokecolor{currentstroke}%
\pgfsetdash{}{0pt}%
\pgfpathmoveto{\pgfqpoint{0.875643in}{1.377277in}}%
\pgfpathlineto{\pgfqpoint{0.871351in}{1.392613in}}%
\pgfpathlineto{\pgfqpoint{0.867066in}{1.407924in}}%
\pgfpathlineto{\pgfqpoint{0.827898in}{1.401879in}}%
\pgfpathlineto{\pgfqpoint{0.789612in}{1.394958in}}%
\pgfpathlineto{\pgfqpoint{0.795410in}{1.379827in}}%
\pgfpathlineto{\pgfqpoint{0.801217in}{1.364674in}}%
\pgfpathlineto{\pgfqpoint{0.838002in}{1.371400in}}%
\pgfpathlineto{\pgfqpoint{0.875643in}{1.377277in}}%
\pgfpathclose%
\pgfusepath{fill}%
\end{pgfscope}%
\begin{pgfscope}%
\pgfpathrectangle{\pgfqpoint{0.407592in}{0.222904in}}{\pgfqpoint{1.411163in}{1.411163in}}%
\pgfusepath{clip}%
\pgfsetbuttcap%
\pgfsetroundjoin%
\definecolor{currentfill}{rgb}{0.749020,0.772549,0.792157}%
\pgfsetfillcolor{currentfill}%
\pgfsetfillopacity{0.160000}%
\pgfsetlinewidth{0.000000pt}%
\definecolor{currentstroke}{rgb}{0.000000,0.000000,0.000000}%
\pgfsetstrokecolor{currentstroke}%
\pgfsetdash{}{0pt}%
\pgfpathmoveto{\pgfqpoint{0.754995in}{1.485286in}}%
\pgfpathlineto{\pgfqpoint{0.749254in}{1.500265in}}%
\pgfpathlineto{\pgfqpoint{0.706844in}{1.491772in}}%
\pgfpathlineto{\pgfqpoint{0.665703in}{1.482357in}}%
\pgfpathlineto{\pgfqpoint{0.672897in}{1.467584in}}%
\pgfpathlineto{\pgfqpoint{0.713318in}{1.476890in}}%
\pgfusepath{fill}%
\end{pgfscope}%
\begin{pgfscope}%
\pgfpathrectangle{\pgfqpoint{0.407592in}{0.222904in}}{\pgfqpoint{1.411163in}{1.411163in}}%
\pgfusepath{clip}%
\pgfsetbuttcap%
\pgfsetroundjoin%
\definecolor{currentfill}{rgb}{0.749020,0.772549,0.792157}%
\pgfsetfillcolor{currentfill}%
\pgfsetfillopacity{0.160000}%
\pgfsetlinewidth{0.000000pt}%
\definecolor{currentstroke}{rgb}{0.000000,0.000000,0.000000}%
\pgfsetstrokecolor{currentstroke}%
\pgfsetdash{}{0pt}%
\pgfpathmoveto{\pgfqpoint{1.240485in}{1.262329in}}%
\pgfpathlineto{\pgfqpoint{1.244104in}{1.278029in}}%
\pgfpathlineto{\pgfqpoint{1.247717in}{1.293700in}}%
\pgfpathlineto{\pgfqpoint{1.213663in}{1.296489in}}%
\pgfpathlineto{\pgfqpoint{1.179155in}{1.298463in}}%
\pgfpathlineto{\pgfqpoint{1.177118in}{1.282707in}}%
\pgfpathlineto{\pgfqpoint{1.175078in}{1.266923in}}%
\pgfpathlineto{\pgfqpoint{1.208000in}{1.265019in}}%
\pgfpathlineto{\pgfqpoint{1.240485in}{1.262329in}}%
\pgfpathclose%
\pgfusepath{fill}%
\end{pgfscope}%
\begin{pgfscope}%
\pgfpathrectangle{\pgfqpoint{0.407592in}{0.222904in}}{\pgfqpoint{1.411163in}{1.411163in}}%
\pgfusepath{clip}%
\pgfsetbuttcap%
\pgfsetroundjoin%
\definecolor{currentfill}{rgb}{0.749020,0.772549,0.792157}%
\pgfsetfillcolor{currentfill}%
\pgfsetfillopacity{0.160000}%
\pgfsetlinewidth{0.000000pt}%
\definecolor{currentstroke}{rgb}{0.000000,0.000000,0.000000}%
\pgfsetstrokecolor{currentstroke}%
\pgfsetdash{}{0pt}%
\pgfpathmoveto{\pgfqpoint{1.531153in}{1.439875in}}%
\pgfpathlineto{\pgfqpoint{1.539039in}{1.454708in}}%
\pgfpathlineto{\pgfqpoint{1.546915in}{1.469523in}}%
\pgfpathlineto{\pgfqpoint{1.508503in}{1.478648in}}%
\pgfpathlineto{\pgfqpoint{1.468638in}{1.486859in}}%
\pgfpathlineto{\pgfqpoint{1.462153in}{1.471840in}}%
\pgfpathlineto{\pgfqpoint{1.455658in}{1.456799in}}%
\pgfpathlineto{\pgfqpoint{1.494113in}{1.448783in}}%
\pgfpathlineto{\pgfqpoint{1.531153in}{1.439875in}}%
\pgfpathclose%
\pgfusepath{fill}%
\end{pgfscope}%
\begin{pgfscope}%
\pgfpathrectangle{\pgfqpoint{0.407592in}{0.222904in}}{\pgfqpoint{1.411163in}{1.411163in}}%
\pgfusepath{clip}%
\pgfsetbuttcap%
\pgfsetroundjoin%
\definecolor{currentfill}{rgb}{0.749020,0.772549,0.792157}%
\pgfsetfillcolor{currentfill}%
\pgfsetfillopacity{0.160000}%
\pgfsetlinewidth{0.000000pt}%
\definecolor{currentstroke}{rgb}{0.000000,0.000000,0.000000}%
\pgfsetstrokecolor{currentstroke}%
\pgfsetdash{}{0pt}%
\pgfpathmoveto{\pgfqpoint{1.041394in}{1.266604in}}%
\pgfpathlineto{\pgfqpoint{1.040220in}{1.282383in}}%
\pgfpathlineto{\pgfqpoint{1.039047in}{1.298132in}}%
\pgfpathlineto{\pgfqpoint{1.004118in}{1.295995in}}%
\pgfpathlineto{\pgfqpoint{0.969475in}{1.293043in}}%
\pgfpathlineto{\pgfqpoint{0.972246in}{1.277383in}}%
\pgfpathlineto{\pgfqpoint{0.975021in}{1.261696in}}%
\pgfpathlineto{\pgfqpoint{1.008069in}{1.264542in}}%
\pgfpathlineto{\pgfqpoint{1.041394in}{1.266604in}}%
\pgfpathclose%
\pgfusepath{fill}%
\end{pgfscope}%
\begin{pgfscope}%
\pgfpathrectangle{\pgfqpoint{0.407592in}{0.222904in}}{\pgfqpoint{1.411163in}{1.411163in}}%
\pgfusepath{clip}%
\pgfsetbuttcap%
\pgfsetroundjoin%
\definecolor{currentfill}{rgb}{0.749020,0.772549,0.792157}%
\pgfsetfillcolor{currentfill}%
\pgfsetfillopacity{0.160000}%
\pgfsetlinewidth{0.000000pt}%
\definecolor{currentstroke}{rgb}{0.000000,0.000000,0.000000}%
\pgfsetstrokecolor{currentstroke}%
\pgfsetdash{}{0pt}%
\pgfpathmoveto{\pgfqpoint{1.313978in}{1.285709in}}%
\pgfpathlineto{\pgfqpoint{1.319116in}{1.301213in}}%
\pgfpathlineto{\pgfqpoint{1.324246in}{1.316692in}}%
\pgfpathlineto{\pgfqpoint{1.289945in}{1.321239in}}%
\pgfpathlineto{\pgfqpoint{1.254923in}{1.324960in}}%
\pgfpathlineto{\pgfqpoint{1.251323in}{1.309344in}}%
\pgfpathlineto{\pgfqpoint{1.247717in}{1.293700in}}%
\pgfpathlineto{\pgfqpoint{1.281196in}{1.290104in}}%
\pgfpathlineto{\pgfqpoint{1.313978in}{1.285709in}}%
\pgfpathclose%
\pgfusepath{fill}%
\end{pgfscope}%
\begin{pgfscope}%
\pgfpathrectangle{\pgfqpoint{0.407592in}{0.222904in}}{\pgfqpoint{1.411163in}{1.411163in}}%
\pgfusepath{clip}%
\pgfsetbuttcap%
\pgfsetroundjoin%
\definecolor{currentfill}{rgb}{0.749020,0.772549,0.792157}%
\pgfsetfillcolor{currentfill}%
\pgfsetfillopacity{0.160000}%
\pgfsetlinewidth{0.000000pt}%
\definecolor{currentstroke}{rgb}{0.000000,0.000000,0.000000}%
\pgfsetstrokecolor{currentstroke}%
\pgfsetdash{}{0pt}%
\pgfpathmoveto{\pgfqpoint{1.403355in}{1.335676in}}%
\pgfpathlineto{\pgfqpoint{1.409927in}{1.350894in}}%
\pgfpathlineto{\pgfqpoint{1.416489in}{1.366090in}}%
\pgfpathlineto{\pgfqpoint{1.381102in}{1.372648in}}%
\pgfpathlineto{\pgfqpoint{1.344679in}{1.378352in}}%
\pgfpathlineto{\pgfqpoint{1.339583in}{1.362975in}}%
\pgfpathlineto{\pgfqpoint{1.334479in}{1.347573in}}%
\pgfpathlineto{\pgfqpoint{1.369419in}{1.342038in}}%
\pgfpathlineto{\pgfqpoint{1.403355in}{1.335676in}}%
\pgfpathclose%
\pgfusepath{fill}%
\end{pgfscope}%
\begin{pgfscope}%
\pgfpathrectangle{\pgfqpoint{0.407592in}{0.222904in}}{\pgfqpoint{1.411163in}{1.411163in}}%
\pgfusepath{clip}%
\pgfsetbuttcap%
\pgfsetroundjoin%
\definecolor{currentfill}{rgb}{0.749020,0.772549,0.792157}%
\pgfsetfillcolor{currentfill}%
\pgfsetfillopacity{0.160000}%
\pgfsetlinewidth{0.000000pt}%
\definecolor{currentstroke}{rgb}{0.000000,0.000000,0.000000}%
\pgfsetstrokecolor{currentstroke}%
\pgfsetdash{}{0pt}%
\pgfpathmoveto{\pgfqpoint{0.969475in}{1.293043in}}%
\pgfpathlineto{\pgfqpoint{0.966709in}{1.308675in}}%
\pgfpathlineto{\pgfqpoint{0.963948in}{1.324280in}}%
\pgfpathlineto{\pgfqpoint{0.928139in}{1.320393in}}%
\pgfpathlineto{\pgfqpoint{0.892882in}{1.315681in}}%
\pgfpathlineto{\pgfqpoint{0.897209in}{1.300220in}}%
\pgfpathlineto{\pgfqpoint{0.901543in}{1.284732in}}%
\pgfpathlineto{\pgfqpoint{0.935242in}{1.289286in}}%
\pgfpathlineto{\pgfqpoint{0.969475in}{1.293043in}}%
\pgfpathclose%
\pgfusepath{fill}%
\end{pgfscope}%
\begin{pgfscope}%
\pgfpathrectangle{\pgfqpoint{0.407592in}{0.222904in}}{\pgfqpoint{1.411163in}{1.411163in}}%
\pgfusepath{clip}%
\pgfsetbuttcap%
\pgfsetroundjoin%
\definecolor{currentfill}{rgb}{0.749020,0.772549,0.792157}%
\pgfsetfillcolor{currentfill}%
\pgfsetfillopacity{0.160000}%
\pgfsetlinewidth{0.000000pt}%
\definecolor{currentstroke}{rgb}{0.000000,0.000000,0.000000}%
\pgfsetstrokecolor{currentstroke}%
\pgfsetdash{}{0pt}%
\pgfpathmoveto{\pgfqpoint{0.766500in}{1.455264in}}%
\pgfpathlineto{\pgfqpoint{0.760743in}{1.470286in}}%
\pgfpathlineto{\pgfqpoint{0.754995in}{1.485286in}}%
\pgfpathlineto{\pgfqpoint{0.713318in}{1.476890in}}%
\pgfpathlineto{\pgfqpoint{0.672897in}{1.467584in}}%
\pgfpathlineto{\pgfqpoint{0.680099in}{1.452792in}}%
\pgfpathlineto{\pgfqpoint{0.687311in}{1.437983in}}%
\pgfpathlineto{\pgfqpoint{0.726293in}{1.447066in}}%
\pgfpathlineto{\pgfqpoint{0.766500in}{1.455264in}}%
\pgfpathclose%
\pgfusepath{fill}%
\end{pgfscope}%
\begin{pgfscope}%
\pgfpathrectangle{\pgfqpoint{0.407592in}{0.222904in}}{\pgfqpoint{1.411163in}{1.411163in}}%
\pgfusepath{clip}%
\pgfsetbuttcap%
\pgfsetroundjoin%
\definecolor{currentfill}{rgb}{0.749020,0.772549,0.792157}%
\pgfsetfillcolor{currentfill}%
\pgfsetfillopacity{0.160000}%
\pgfsetlinewidth{0.000000pt}%
\definecolor{currentstroke}{rgb}{0.000000,0.000000,0.000000}%
\pgfsetstrokecolor{currentstroke}%
\pgfsetdash{}{0pt}%
\pgfpathmoveto{\pgfqpoint{1.170986in}{1.235265in}}%
\pgfpathlineto{\pgfqpoint{1.173034in}{1.251109in}}%
\pgfpathlineto{\pgfqpoint{1.175078in}{1.266923in}}%
\pgfpathlineto{\pgfqpoint{1.141838in}{1.268035in}}%
\pgfpathlineto{\pgfqpoint{1.108398in}{1.268353in}}%
\pgfpathlineto{\pgfqpoint{1.107960in}{1.252512in}}%
\pgfpathlineto{\pgfqpoint{1.107522in}{1.236641in}}%
\pgfpathlineto{\pgfqpoint{1.139350in}{1.236336in}}%
\pgfpathlineto{\pgfqpoint{1.170986in}{1.235265in}}%
\pgfpathclose%
\pgfusepath{fill}%
\end{pgfscope}%
\begin{pgfscope}%
\pgfpathrectangle{\pgfqpoint{0.407592in}{0.222904in}}{\pgfqpoint{1.411163in}{1.411163in}}%
\pgfusepath{clip}%
\pgfsetbuttcap%
\pgfsetroundjoin%
\definecolor{currentfill}{rgb}{0.749020,0.772549,0.792157}%
\pgfsetfillcolor{currentfill}%
\pgfsetfillopacity{0.160000}%
\pgfsetlinewidth{0.000000pt}%
\definecolor{currentstroke}{rgb}{0.000000,0.000000,0.000000}%
\pgfsetstrokecolor{currentstroke}%
\pgfsetdash{}{0pt}%
\pgfpathmoveto{\pgfqpoint{1.107522in}{1.236641in}}%
\pgfpathlineto{\pgfqpoint{1.107960in}{1.252512in}}%
\pgfpathlineto{\pgfqpoint{1.108398in}{1.268353in}}%
\pgfpathlineto{\pgfqpoint{1.074877in}{1.267876in}}%
\pgfpathlineto{\pgfqpoint{1.041394in}{1.266604in}}%
\pgfpathlineto{\pgfqpoint{1.042571in}{1.250796in}}%
\pgfpathlineto{\pgfqpoint{1.043750in}{1.234959in}}%
\pgfpathlineto{\pgfqpoint{1.075617in}{1.236182in}}%
\pgfpathlineto{\pgfqpoint{1.107522in}{1.236641in}}%
\pgfpathclose%
\pgfusepath{fill}%
\end{pgfscope}%
\begin{pgfscope}%
\pgfpathrectangle{\pgfqpoint{0.407592in}{0.222904in}}{\pgfqpoint{1.411163in}{1.411163in}}%
\pgfusepath{clip}%
\pgfsetbuttcap%
\pgfsetroundjoin%
\definecolor{currentfill}{rgb}{0.749020,0.772549,0.792157}%
\pgfsetfillcolor{currentfill}%
\pgfsetfillopacity{0.160000}%
\pgfsetlinewidth{0.000000pt}%
\definecolor{currentstroke}{rgb}{0.000000,0.000000,0.000000}%
\pgfsetstrokecolor{currentstroke}%
\pgfsetdash{}{0pt}%
\pgfpathmoveto{\pgfqpoint{0.884248in}{1.346529in}}%
\pgfpathlineto{\pgfqpoint{0.879942in}{1.361916in}}%
\pgfpathlineto{\pgfqpoint{0.875643in}{1.377277in}}%
\pgfpathlineto{\pgfqpoint{0.838002in}{1.371400in}}%
\pgfpathlineto{\pgfqpoint{0.801217in}{1.364674in}}%
\pgfpathlineto{\pgfqpoint{0.807033in}{1.349499in}}%
\pgfpathlineto{\pgfqpoint{0.812857in}{1.334301in}}%
\pgfpathlineto{\pgfqpoint{0.848137in}{1.340827in}}%
\pgfpathlineto{\pgfqpoint{0.884248in}{1.346529in}}%
\pgfpathclose%
\pgfusepath{fill}%
\end{pgfscope}%
\begin{pgfscope}%
\pgfpathrectangle{\pgfqpoint{0.407592in}{0.222904in}}{\pgfqpoint{1.411163in}{1.411163in}}%
\pgfusepath{clip}%
\pgfsetbuttcap%
\pgfsetroundjoin%
\definecolor{currentfill}{rgb}{0.749020,0.772549,0.792157}%
\pgfsetfillcolor{currentfill}%
\pgfsetfillopacity{0.160000}%
\pgfsetlinewidth{0.000000pt}%
\definecolor{currentstroke}{rgb}{0.000000,0.000000,0.000000}%
\pgfsetstrokecolor{currentstroke}%
\pgfsetdash{}{0pt}%
\pgfpathmoveto{\pgfqpoint{1.515351in}{1.410154in}}%
\pgfpathlineto{\pgfqpoint{1.523257in}{1.425024in}}%
\pgfpathlineto{\pgfqpoint{1.531153in}{1.439875in}}%
\pgfpathlineto{\pgfqpoint{1.494113in}{1.448783in}}%
\pgfpathlineto{\pgfqpoint{1.455658in}{1.456799in}}%
\pgfpathlineto{\pgfqpoint{1.449154in}{1.441737in}}%
\pgfpathlineto{\pgfqpoint{1.442640in}{1.426652in}}%
\pgfpathlineto{\pgfqpoint{1.479684in}{1.418836in}}%
\pgfpathlineto{\pgfqpoint{1.515351in}{1.410154in}}%
\pgfpathclose%
\pgfusepath{fill}%
\end{pgfscope}%
\begin{pgfscope}%
\pgfpathrectangle{\pgfqpoint{0.407592in}{0.222904in}}{\pgfqpoint{1.411163in}{1.411163in}}%
\pgfusepath{clip}%
\pgfsetbuttcap%
\pgfsetroundjoin%
\definecolor{currentfill}{rgb}{0.749020,0.772549,0.792157}%
\pgfsetfillcolor{currentfill}%
\pgfsetfillopacity{0.160000}%
\pgfsetlinewidth{0.000000pt}%
\definecolor{currentstroke}{rgb}{0.000000,0.000000,0.000000}%
\pgfsetstrokecolor{currentstroke}%
\pgfsetdash{}{0pt}%
\pgfpathmoveto{\pgfqpoint{1.233227in}{1.230846in}}%
\pgfpathlineto{\pgfqpoint{1.236860in}{1.246602in}}%
\pgfpathlineto{\pgfqpoint{1.240485in}{1.262329in}}%
\pgfpathlineto{\pgfqpoint{1.208000in}{1.265019in}}%
\pgfpathlineto{\pgfqpoint{1.175078in}{1.266923in}}%
\pgfpathlineto{\pgfqpoint{1.173034in}{1.251109in}}%
\pgfpathlineto{\pgfqpoint{1.170986in}{1.235265in}}%
\pgfpathlineto{\pgfqpoint{1.202316in}{1.233434in}}%
\pgfpathlineto{\pgfqpoint{1.233227in}{1.230846in}}%
\pgfpathclose%
\pgfusepath{fill}%
\end{pgfscope}%
\begin{pgfscope}%
\pgfpathrectangle{\pgfqpoint{0.407592in}{0.222904in}}{\pgfqpoint{1.411163in}{1.411163in}}%
\pgfusepath{clip}%
\pgfsetbuttcap%
\pgfsetroundjoin%
\definecolor{currentfill}{rgb}{0.749020,0.772549,0.792157}%
\pgfsetfillcolor{currentfill}%
\pgfsetfillopacity{0.160000}%
\pgfsetlinewidth{0.000000pt}%
\definecolor{currentstroke}{rgb}{0.000000,0.000000,0.000000}%
\pgfsetstrokecolor{currentstroke}%
\pgfsetdash{}{0pt}%
\pgfpathmoveto{\pgfqpoint{1.043750in}{1.234959in}}%
\pgfpathlineto{\pgfqpoint{1.042571in}{1.250796in}}%
\pgfpathlineto{\pgfqpoint{1.041394in}{1.266604in}}%
\pgfpathlineto{\pgfqpoint{1.008069in}{1.264542in}}%
\pgfpathlineto{\pgfqpoint{0.975021in}{1.261696in}}%
\pgfpathlineto{\pgfqpoint{0.977802in}{1.245980in}}%
\pgfpathlineto{\pgfqpoint{0.980587in}{1.230237in}}%
\pgfpathlineto{\pgfqpoint{1.012035in}{1.232975in}}%
\pgfpathlineto{\pgfqpoint{1.043750in}{1.234959in}}%
\pgfpathclose%
\pgfusepath{fill}%
\end{pgfscope}%
\begin{pgfscope}%
\pgfpathrectangle{\pgfqpoint{0.407592in}{0.222904in}}{\pgfqpoint{1.411163in}{1.411163in}}%
\pgfusepath{clip}%
\pgfsetbuttcap%
\pgfsetroundjoin%
\definecolor{currentfill}{rgb}{0.749020,0.772549,0.792157}%
\pgfsetfillcolor{currentfill}%
\pgfsetfillopacity{0.160000}%
\pgfsetlinewidth{0.000000pt}%
\definecolor{currentstroke}{rgb}{0.000000,0.000000,0.000000}%
\pgfsetstrokecolor{currentstroke}%
\pgfsetdash{}{0pt}%
\pgfpathmoveto{\pgfqpoint{0.778039in}{1.425154in}}%
\pgfpathlineto{\pgfqpoint{0.772266in}{1.440220in}}%
\pgfpathlineto{\pgfqpoint{0.766500in}{1.455264in}}%
\pgfpathlineto{\pgfqpoint{0.726293in}{1.447066in}}%
\pgfpathlineto{\pgfqpoint{0.687311in}{1.437983in}}%
\pgfpathlineto{\pgfqpoint{0.694531in}{1.423155in}}%
\pgfpathlineto{\pgfqpoint{0.701760in}{1.408309in}}%
\pgfpathlineto{\pgfqpoint{0.739303in}{1.417163in}}%
\pgfpathlineto{\pgfqpoint{0.778039in}{1.425154in}}%
\pgfpathclose%
\pgfusepath{fill}%
\end{pgfscope}%
\begin{pgfscope}%
\pgfpathrectangle{\pgfqpoint{0.407592in}{0.222904in}}{\pgfqpoint{1.411163in}{1.411163in}}%
\pgfusepath{clip}%
\pgfsetbuttcap%
\pgfsetroundjoin%
\definecolor{currentfill}{rgb}{0.749020,0.772549,0.792157}%
\pgfsetfillcolor{currentfill}%
\pgfsetfillopacity{0.160000}%
\pgfsetlinewidth{0.000000pt}%
\definecolor{currentstroke}{rgb}{0.000000,0.000000,0.000000}%
\pgfsetstrokecolor{currentstroke}%
\pgfsetdash{}{0pt}%
\pgfpathmoveto{\pgfqpoint{1.303676in}{1.254623in}}%
\pgfpathlineto{\pgfqpoint{1.308831in}{1.270179in}}%
\pgfpathlineto{\pgfqpoint{1.313978in}{1.285709in}}%
\pgfpathlineto{\pgfqpoint{1.281196in}{1.290104in}}%
\pgfpathlineto{\pgfqpoint{1.247717in}{1.293700in}}%
\pgfpathlineto{\pgfqpoint{1.244104in}{1.278029in}}%
\pgfpathlineto{\pgfqpoint{1.240485in}{1.262329in}}%
\pgfpathlineto{\pgfqpoint{1.272416in}{1.258861in}}%
\pgfpathlineto{\pgfqpoint{1.303676in}{1.254623in}}%
\pgfpathclose%
\pgfusepath{fill}%
\end{pgfscope}%
\begin{pgfscope}%
\pgfpathrectangle{\pgfqpoint{0.407592in}{0.222904in}}{\pgfqpoint{1.411163in}{1.411163in}}%
\pgfusepath{clip}%
\pgfsetbuttcap%
\pgfsetroundjoin%
\definecolor{currentfill}{rgb}{0.749020,0.772549,0.792157}%
\pgfsetfillcolor{currentfill}%
\pgfsetfillopacity{0.160000}%
\pgfsetlinewidth{0.000000pt}%
\definecolor{currentstroke}{rgb}{0.000000,0.000000,0.000000}%
\pgfsetstrokecolor{currentstroke}%
\pgfsetdash{}{0pt}%
\pgfpathmoveto{\pgfqpoint{1.390183in}{1.305172in}}%
\pgfpathlineto{\pgfqpoint{1.396774in}{1.320435in}}%
\pgfpathlineto{\pgfqpoint{1.403355in}{1.335676in}}%
\pgfpathlineto{\pgfqpoint{1.369419in}{1.342038in}}%
\pgfpathlineto{\pgfqpoint{1.334479in}{1.347573in}}%
\pgfpathlineto{\pgfqpoint{1.329367in}{1.332145in}}%
\pgfpathlineto{\pgfqpoint{1.324246in}{1.316692in}}%
\pgfpathlineto{\pgfqpoint{1.357700in}{1.311332in}}%
\pgfpathlineto{\pgfqpoint{1.390183in}{1.305172in}}%
\pgfpathclose%
\pgfusepath{fill}%
\end{pgfscope}%
\begin{pgfscope}%
\pgfpathrectangle{\pgfqpoint{0.407592in}{0.222904in}}{\pgfqpoint{1.411163in}{1.411163in}}%
\pgfusepath{clip}%
\pgfsetbuttcap%
\pgfsetroundjoin%
\definecolor{currentfill}{rgb}{0.749020,0.772549,0.792157}%
\pgfsetfillcolor{currentfill}%
\pgfsetfillopacity{0.160000}%
\pgfsetlinewidth{0.000000pt}%
\definecolor{currentstroke}{rgb}{0.000000,0.000000,0.000000}%
\pgfsetstrokecolor{currentstroke}%
\pgfsetdash{}{0pt}%
\pgfpathmoveto{\pgfqpoint{0.975021in}{1.261696in}}%
\pgfpathlineto{\pgfqpoint{0.972246in}{1.277383in}}%
\pgfpathlineto{\pgfqpoint{0.969475in}{1.293043in}}%
\pgfpathlineto{\pgfqpoint{0.935242in}{1.289286in}}%
\pgfpathlineto{\pgfqpoint{0.901543in}{1.284732in}}%
\pgfpathlineto{\pgfqpoint{0.905885in}{1.269220in}}%
\pgfpathlineto{\pgfqpoint{0.910233in}{1.253682in}}%
\pgfpathlineto{\pgfqpoint{0.942370in}{1.258072in}}%
\pgfpathlineto{\pgfqpoint{0.975021in}{1.261696in}}%
\pgfpathclose%
\pgfusepath{fill}%
\end{pgfscope}%
\begin{pgfscope}%
\pgfpathrectangle{\pgfqpoint{0.407592in}{0.222904in}}{\pgfqpoint{1.411163in}{1.411163in}}%
\pgfusepath{clip}%
\pgfsetbuttcap%
\pgfsetroundjoin%
\definecolor{currentfill}{rgb}{0.749020,0.772549,0.792157}%
\pgfsetfillcolor{currentfill}%
\pgfsetfillopacity{0.160000}%
\pgfsetlinewidth{0.000000pt}%
\definecolor{currentstroke}{rgb}{0.000000,0.000000,0.000000}%
\pgfsetstrokecolor{currentstroke}%
\pgfsetdash{}{0pt}%
\pgfpathmoveto{\pgfqpoint{0.892882in}{1.315681in}}%
\pgfpathlineto{\pgfqpoint{0.888561in}{1.331118in}}%
\pgfpathlineto{\pgfqpoint{0.884248in}{1.346529in}}%
\pgfpathlineto{\pgfqpoint{0.848137in}{1.340827in}}%
\pgfpathlineto{\pgfqpoint{0.812857in}{1.334301in}}%
\pgfpathlineto{\pgfqpoint{0.818690in}{1.319082in}}%
\pgfpathlineto{\pgfqpoint{0.824531in}{1.303841in}}%
\pgfpathlineto{\pgfqpoint{0.858303in}{1.310159in}}%
\pgfpathlineto{\pgfqpoint{0.892882in}{1.315681in}}%
\pgfpathclose%
\pgfusepath{fill}%
\end{pgfscope}%
\begin{pgfscope}%
\pgfpathrectangle{\pgfqpoint{0.407592in}{0.222904in}}{\pgfqpoint{1.411163in}{1.411163in}}%
\pgfusepath{clip}%
\pgfsetbuttcap%
\pgfsetroundjoin%
\definecolor{currentfill}{rgb}{0.749020,0.772549,0.792157}%
\pgfsetfillcolor{currentfill}%
\pgfsetfillopacity{0.160000}%
\pgfsetlinewidth{0.000000pt}%
\definecolor{currentstroke}{rgb}{0.000000,0.000000,0.000000}%
\pgfsetstrokecolor{currentstroke}%
\pgfsetdash{}{0pt}%
\pgfpathmoveto{\pgfqpoint{1.166878in}{1.203490in}}%
\pgfpathlineto{\pgfqpoint{1.168934in}{1.219392in}}%
\pgfpathlineto{\pgfqpoint{1.170986in}{1.235265in}}%
\pgfpathlineto{\pgfqpoint{1.139350in}{1.236336in}}%
\pgfpathlineto{\pgfqpoint{1.107522in}{1.236641in}}%
\pgfpathlineto{\pgfqpoint{1.107083in}{1.220741in}}%
\pgfpathlineto{\pgfqpoint{1.106643in}{1.204810in}}%
\pgfpathlineto{\pgfqpoint{1.136852in}{1.204516in}}%
\pgfpathlineto{\pgfqpoint{1.166878in}{1.203490in}}%
\pgfpathclose%
\pgfusepath{fill}%
\end{pgfscope}%
\begin{pgfscope}%
\pgfpathrectangle{\pgfqpoint{0.407592in}{0.222904in}}{\pgfqpoint{1.411163in}{1.411163in}}%
\pgfusepath{clip}%
\pgfsetbuttcap%
\pgfsetroundjoin%
\definecolor{currentfill}{rgb}{0.749020,0.772549,0.792157}%
\pgfsetfillcolor{currentfill}%
\pgfsetfillopacity{0.160000}%
\pgfsetlinewidth{0.000000pt}%
\definecolor{currentstroke}{rgb}{0.000000,0.000000,0.000000}%
\pgfsetstrokecolor{currentstroke}%
\pgfsetdash{}{0pt}%
\pgfpathmoveto{\pgfqpoint{1.106643in}{1.204810in}}%
\pgfpathlineto{\pgfqpoint{1.107083in}{1.220741in}}%
\pgfpathlineto{\pgfqpoint{1.107522in}{1.236641in}}%
\pgfpathlineto{\pgfqpoint{1.075617in}{1.236182in}}%
\pgfpathlineto{\pgfqpoint{1.043750in}{1.234959in}}%
\pgfpathlineto{\pgfqpoint{1.044931in}{1.219092in}}%
\pgfpathlineto{\pgfqpoint{1.046114in}{1.203196in}}%
\pgfpathlineto{\pgfqpoint{1.076360in}{1.204369in}}%
\pgfpathlineto{\pgfqpoint{1.106643in}{1.204810in}}%
\pgfpathclose%
\pgfusepath{fill}%
\end{pgfscope}%
\begin{pgfscope}%
\pgfpathrectangle{\pgfqpoint{0.407592in}{0.222904in}}{\pgfqpoint{1.411163in}{1.411163in}}%
\pgfusepath{clip}%
\pgfsetbuttcap%
\pgfsetroundjoin%
\definecolor{currentfill}{rgb}{0.749020,0.772549,0.792157}%
\pgfsetfillcolor{currentfill}%
\pgfsetfillopacity{0.160000}%
\pgfsetlinewidth{0.000000pt}%
\definecolor{currentstroke}{rgb}{0.000000,0.000000,0.000000}%
\pgfsetstrokecolor{currentstroke}%
\pgfsetdash{}{0pt}%
\pgfpathmoveto{\pgfqpoint{1.499510in}{1.380358in}}%
\pgfpathlineto{\pgfqpoint{1.507436in}{1.395265in}}%
\pgfpathlineto{\pgfqpoint{1.515351in}{1.410154in}}%
\pgfpathlineto{\pgfqpoint{1.479684in}{1.418836in}}%
\pgfpathlineto{\pgfqpoint{1.442640in}{1.426652in}}%
\pgfpathlineto{\pgfqpoint{1.436117in}{1.411545in}}%
\pgfpathlineto{\pgfqpoint{1.429584in}{1.396416in}}%
\pgfpathlineto{\pgfqpoint{1.465216in}{1.388808in}}%
\pgfpathlineto{\pgfqpoint{1.499510in}{1.380358in}}%
\pgfpathclose%
\pgfusepath{fill}%
\end{pgfscope}%
\begin{pgfscope}%
\pgfpathrectangle{\pgfqpoint{0.407592in}{0.222904in}}{\pgfqpoint{1.411163in}{1.411163in}}%
\pgfusepath{clip}%
\pgfsetbuttcap%
\pgfsetroundjoin%
\definecolor{currentfill}{rgb}{0.749020,0.772549,0.792157}%
\pgfsetfillcolor{currentfill}%
\pgfsetfillopacity{0.160000}%
\pgfsetlinewidth{0.000000pt}%
\definecolor{currentstroke}{rgb}{0.000000,0.000000,0.000000}%
\pgfsetstrokecolor{currentstroke}%
\pgfsetdash{}{0pt}%
\pgfpathmoveto{\pgfqpoint{0.789612in}{1.394958in}}%
\pgfpathlineto{\pgfqpoint{0.783821in}{1.410067in}}%
\pgfpathlineto{\pgfqpoint{0.778039in}{1.425154in}}%
\pgfpathlineto{\pgfqpoint{0.739303in}{1.417163in}}%
\pgfpathlineto{\pgfqpoint{0.701760in}{1.408309in}}%
\pgfpathlineto{\pgfqpoint{0.708998in}{1.393446in}}%
\pgfpathlineto{\pgfqpoint{0.716244in}{1.378564in}}%
\pgfpathlineto{\pgfqpoint{0.752347in}{1.387179in}}%
\pgfpathlineto{\pgfqpoint{0.789612in}{1.394958in}}%
\pgfpathclose%
\pgfusepath{fill}%
\end{pgfscope}%
\begin{pgfscope}%
\pgfpathrectangle{\pgfqpoint{0.407592in}{0.222904in}}{\pgfqpoint{1.411163in}{1.411163in}}%
\pgfusepath{clip}%
\pgfsetbuttcap%
\pgfsetroundjoin%
\definecolor{currentfill}{rgb}{0.749020,0.772549,0.792157}%
\pgfsetfillcolor{currentfill}%
\pgfsetfillopacity{0.160000}%
\pgfsetlinewidth{0.000000pt}%
\definecolor{currentstroke}{rgb}{0.000000,0.000000,0.000000}%
\pgfsetstrokecolor{currentstroke}%
\pgfsetdash{}{0pt}%
\pgfpathmoveto{\pgfqpoint{1.225944in}{1.199251in}}%
\pgfpathlineto{\pgfqpoint{1.229589in}{1.215063in}}%
\pgfpathlineto{\pgfqpoint{1.233227in}{1.230846in}}%
\pgfpathlineto{\pgfqpoint{1.202316in}{1.233434in}}%
\pgfpathlineto{\pgfqpoint{1.170986in}{1.235265in}}%
\pgfpathlineto{\pgfqpoint{1.168934in}{1.219392in}}%
\pgfpathlineto{\pgfqpoint{1.166878in}{1.203490in}}%
\pgfpathlineto{\pgfqpoint{1.196612in}{1.201733in}}%
\pgfpathlineto{\pgfqpoint{1.225944in}{1.199251in}}%
\pgfpathclose%
\pgfusepath{fill}%
\end{pgfscope}%
\begin{pgfscope}%
\pgfpathrectangle{\pgfqpoint{0.407592in}{0.222904in}}{\pgfqpoint{1.411163in}{1.411163in}}%
\pgfusepath{clip}%
\pgfsetbuttcap%
\pgfsetroundjoin%
\definecolor{currentfill}{rgb}{0.749020,0.772549,0.792157}%
\pgfsetfillcolor{currentfill}%
\pgfsetfillopacity{0.160000}%
\pgfsetlinewidth{0.000000pt}%
\definecolor{currentstroke}{rgb}{0.000000,0.000000,0.000000}%
\pgfsetstrokecolor{currentstroke}%
\pgfsetdash{}{0pt}%
\pgfpathmoveto{\pgfqpoint{1.046114in}{1.203196in}}%
\pgfpathlineto{\pgfqpoint{1.044931in}{1.219092in}}%
\pgfpathlineto{\pgfqpoint{1.043750in}{1.234959in}}%
\pgfpathlineto{\pgfqpoint{1.012035in}{1.232975in}}%
\pgfpathlineto{\pgfqpoint{0.980587in}{1.230237in}}%
\pgfpathlineto{\pgfqpoint{0.983378in}{1.214465in}}%
\pgfpathlineto{\pgfqpoint{0.986173in}{1.198666in}}%
\pgfpathlineto{\pgfqpoint{1.016015in}{1.201293in}}%
\pgfpathlineto{\pgfqpoint{1.046114in}{1.203196in}}%
\pgfpathclose%
\pgfusepath{fill}%
\end{pgfscope}%
\begin{pgfscope}%
\pgfpathrectangle{\pgfqpoint{0.407592in}{0.222904in}}{\pgfqpoint{1.411163in}{1.411163in}}%
\pgfusepath{clip}%
\pgfsetbuttcap%
\pgfsetroundjoin%
\definecolor{currentfill}{rgb}{0.749020,0.772549,0.792157}%
\pgfsetfillcolor{currentfill}%
\pgfsetfillopacity{0.160000}%
\pgfsetlinewidth{0.000000pt}%
\definecolor{currentstroke}{rgb}{0.000000,0.000000,0.000000}%
\pgfsetstrokecolor{currentstroke}%
\pgfsetdash{}{0pt}%
\pgfpathmoveto{\pgfqpoint{1.376972in}{1.274577in}}%
\pgfpathlineto{\pgfqpoint{1.383583in}{1.289886in}}%
\pgfpathlineto{\pgfqpoint{1.390183in}{1.305172in}}%
\pgfpathlineto{\pgfqpoint{1.357700in}{1.311332in}}%
\pgfpathlineto{\pgfqpoint{1.324246in}{1.316692in}}%
\pgfpathlineto{\pgfqpoint{1.319116in}{1.301213in}}%
\pgfpathlineto{\pgfqpoint{1.313978in}{1.285709in}}%
\pgfpathlineto{\pgfqpoint{1.345943in}{1.280529in}}%
\pgfpathlineto{\pgfqpoint{1.376972in}{1.274577in}}%
\pgfpathclose%
\pgfusepath{fill}%
\end{pgfscope}%
\begin{pgfscope}%
\pgfpathrectangle{\pgfqpoint{0.407592in}{0.222904in}}{\pgfqpoint{1.411163in}{1.411163in}}%
\pgfusepath{clip}%
\pgfsetbuttcap%
\pgfsetroundjoin%
\definecolor{currentfill}{rgb}{0.749020,0.772549,0.792157}%
\pgfsetfillcolor{currentfill}%
\pgfsetfillopacity{0.160000}%
\pgfsetlinewidth{0.000000pt}%
\definecolor{currentstroke}{rgb}{0.000000,0.000000,0.000000}%
\pgfsetstrokecolor{currentstroke}%
\pgfsetdash{}{0pt}%
\pgfpathmoveto{\pgfqpoint{1.293341in}{1.223435in}}%
\pgfpathlineto{\pgfqpoint{1.298513in}{1.239042in}}%
\pgfpathlineto{\pgfqpoint{1.303676in}{1.254623in}}%
\pgfpathlineto{\pgfqpoint{1.272416in}{1.258861in}}%
\pgfpathlineto{\pgfqpoint{1.240485in}{1.262329in}}%
\pgfpathlineto{\pgfqpoint{1.236860in}{1.246602in}}%
\pgfpathlineto{\pgfqpoint{1.233227in}{1.230846in}}%
\pgfpathlineto{\pgfqpoint{1.263606in}{1.227510in}}%
\pgfpathlineto{\pgfqpoint{1.293341in}{1.223435in}}%
\pgfpathclose%
\pgfusepath{fill}%
\end{pgfscope}%
\begin{pgfscope}%
\pgfpathrectangle{\pgfqpoint{0.407592in}{0.222904in}}{\pgfqpoint{1.411163in}{1.411163in}}%
\pgfusepath{clip}%
\pgfsetbuttcap%
\pgfsetroundjoin%
\definecolor{currentfill}{rgb}{0.749020,0.772549,0.792157}%
\pgfsetfillcolor{currentfill}%
\pgfsetfillopacity{0.160000}%
\pgfsetlinewidth{0.000000pt}%
\definecolor{currentstroke}{rgb}{0.000000,0.000000,0.000000}%
\pgfsetstrokecolor{currentstroke}%
\pgfsetdash{}{0pt}%
\pgfpathmoveto{\pgfqpoint{0.980587in}{1.230237in}}%
\pgfpathlineto{\pgfqpoint{0.977802in}{1.245980in}}%
\pgfpathlineto{\pgfqpoint{0.975021in}{1.261696in}}%
\pgfpathlineto{\pgfqpoint{0.942370in}{1.258072in}}%
\pgfpathlineto{\pgfqpoint{0.910233in}{1.253682in}}%
\pgfpathlineto{\pgfqpoint{0.914589in}{1.238118in}}%
\pgfpathlineto{\pgfqpoint{0.918952in}{1.222529in}}%
\pgfpathlineto{\pgfqpoint{0.949521in}{1.226751in}}%
\pgfpathlineto{\pgfqpoint{0.980587in}{1.230237in}}%
\pgfpathclose%
\pgfusepath{fill}%
\end{pgfscope}%
\begin{pgfscope}%
\pgfpathrectangle{\pgfqpoint{0.407592in}{0.222904in}}{\pgfqpoint{1.411163in}{1.411163in}}%
\pgfusepath{clip}%
\pgfsetbuttcap%
\pgfsetroundjoin%
\definecolor{currentfill}{rgb}{0.749020,0.772549,0.792157}%
\pgfsetfillcolor{currentfill}%
\pgfsetfillopacity{0.160000}%
\pgfsetlinewidth{0.000000pt}%
\definecolor{currentstroke}{rgb}{0.000000,0.000000,0.000000}%
\pgfsetstrokecolor{currentstroke}%
\pgfsetdash{}{0pt}%
\pgfpathmoveto{\pgfqpoint{1.483630in}{1.350488in}}%
\pgfpathlineto{\pgfqpoint{1.491575in}{1.365433in}}%
\pgfpathlineto{\pgfqpoint{1.499510in}{1.380358in}}%
\pgfpathlineto{\pgfqpoint{1.465216in}{1.388808in}}%
\pgfpathlineto{\pgfqpoint{1.429584in}{1.396416in}}%
\pgfpathlineto{\pgfqpoint{1.423041in}{1.381264in}}%
\pgfpathlineto{\pgfqpoint{1.416489in}{1.366090in}}%
\pgfpathlineto{\pgfqpoint{1.450708in}{1.358697in}}%
\pgfpathlineto{\pgfqpoint{1.483630in}{1.350488in}}%
\pgfpathclose%
\pgfusepath{fill}%
\end{pgfscope}%
\begin{pgfscope}%
\pgfpathrectangle{\pgfqpoint{0.407592in}{0.222904in}}{\pgfqpoint{1.411163in}{1.411163in}}%
\pgfusepath{clip}%
\pgfsetbuttcap%
\pgfsetroundjoin%
\definecolor{currentfill}{rgb}{0.749020,0.772549,0.792157}%
\pgfsetfillcolor{currentfill}%
\pgfsetfillopacity{0.160000}%
\pgfsetlinewidth{0.000000pt}%
\definecolor{currentstroke}{rgb}{0.000000,0.000000,0.000000}%
\pgfsetstrokecolor{currentstroke}%
\pgfsetdash{}{0pt}%
\pgfpathmoveto{\pgfqpoint{0.901543in}{1.284732in}}%
\pgfpathlineto{\pgfqpoint{0.897209in}{1.300220in}}%
\pgfpathlineto{\pgfqpoint{0.892882in}{1.315681in}}%
\pgfpathlineto{\pgfqpoint{0.858303in}{1.310159in}}%
\pgfpathlineto{\pgfqpoint{0.824531in}{1.303841in}}%
\pgfpathlineto{\pgfqpoint{0.830380in}{1.288577in}}%
\pgfpathlineto{\pgfqpoint{0.836238in}{1.273292in}}%
\pgfpathlineto{\pgfqpoint{0.868501in}{1.279396in}}%
\pgfpathlineto{\pgfqpoint{0.901543in}{1.284732in}}%
\pgfpathclose%
\pgfusepath{fill}%
\end{pgfscope}%
\begin{pgfscope}%
\pgfpathrectangle{\pgfqpoint{0.407592in}{0.222904in}}{\pgfqpoint{1.411163in}{1.411163in}}%
\pgfusepath{clip}%
\pgfsetbuttcap%
\pgfsetroundjoin%
\definecolor{currentfill}{rgb}{0.749020,0.772549,0.792157}%
\pgfsetfillcolor{currentfill}%
\pgfsetfillopacity{0.160000}%
\pgfsetlinewidth{0.000000pt}%
\definecolor{currentstroke}{rgb}{0.000000,0.000000,0.000000}%
\pgfsetstrokecolor{currentstroke}%
\pgfsetdash{}{0pt}%
\pgfpathmoveto{\pgfqpoint{1.162756in}{1.171595in}}%
\pgfpathlineto{\pgfqpoint{1.164819in}{1.187557in}}%
\pgfpathlineto{\pgfqpoint{1.166878in}{1.203490in}}%
\pgfpathlineto{\pgfqpoint{1.136852in}{1.204516in}}%
\pgfpathlineto{\pgfqpoint{1.106643in}{1.204810in}}%
\pgfpathlineto{\pgfqpoint{1.106202in}{1.188849in}}%
\pgfpathlineto{\pgfqpoint{1.105760in}{1.172858in}}%
\pgfpathlineto{\pgfqpoint{1.134345in}{1.172577in}}%
\pgfpathlineto{\pgfqpoint{1.162756in}{1.171595in}}%
\pgfpathclose%
\pgfusepath{fill}%
\end{pgfscope}%
\begin{pgfscope}%
\pgfpathrectangle{\pgfqpoint{0.407592in}{0.222904in}}{\pgfqpoint{1.411163in}{1.411163in}}%
\pgfusepath{clip}%
\pgfsetbuttcap%
\pgfsetroundjoin%
\definecolor{currentfill}{rgb}{0.749020,0.772549,0.792157}%
\pgfsetfillcolor{currentfill}%
\pgfsetfillopacity{0.160000}%
\pgfsetlinewidth{0.000000pt}%
\definecolor{currentstroke}{rgb}{0.000000,0.000000,0.000000}%
\pgfsetstrokecolor{currentstroke}%
\pgfsetdash{}{0pt}%
\pgfpathmoveto{\pgfqpoint{1.105760in}{1.172858in}}%
\pgfpathlineto{\pgfqpoint{1.106202in}{1.188849in}}%
\pgfpathlineto{\pgfqpoint{1.106643in}{1.204810in}}%
\pgfpathlineto{\pgfqpoint{1.076360in}{1.204369in}}%
\pgfpathlineto{\pgfqpoint{1.046114in}{1.203196in}}%
\pgfpathlineto{\pgfqpoint{1.047299in}{1.187270in}}%
\pgfpathlineto{\pgfqpoint{1.048487in}{1.171314in}}%
\pgfpathlineto{\pgfqpoint{1.077106in}{1.172436in}}%
\pgfpathlineto{\pgfqpoint{1.105760in}{1.172858in}}%
\pgfpathclose%
\pgfusepath{fill}%
\end{pgfscope}%
\begin{pgfscope}%
\pgfpathrectangle{\pgfqpoint{0.407592in}{0.222904in}}{\pgfqpoint{1.411163in}{1.411163in}}%
\pgfusepath{clip}%
\pgfsetbuttcap%
\pgfsetroundjoin%
\definecolor{currentfill}{rgb}{0.749020,0.772549,0.792157}%
\pgfsetfillcolor{currentfill}%
\pgfsetfillopacity{0.160000}%
\pgfsetlinewidth{0.000000pt}%
\definecolor{currentstroke}{rgb}{0.000000,0.000000,0.000000}%
\pgfsetstrokecolor{currentstroke}%
\pgfsetdash{}{0pt}%
\pgfpathmoveto{\pgfqpoint{0.801217in}{1.364674in}}%
\pgfpathlineto{\pgfqpoint{0.795410in}{1.379827in}}%
\pgfpathlineto{\pgfqpoint{0.789612in}{1.394958in}}%
\pgfpathlineto{\pgfqpoint{0.752347in}{1.387179in}}%
\pgfpathlineto{\pgfqpoint{0.716244in}{1.378564in}}%
\pgfpathlineto{\pgfqpoint{0.723500in}{1.363663in}}%
\pgfpathlineto{\pgfqpoint{0.730764in}{1.348745in}}%
\pgfpathlineto{\pgfqpoint{0.765427in}{1.357115in}}%
\pgfpathlineto{\pgfqpoint{0.801217in}{1.364674in}}%
\pgfpathclose%
\pgfusepath{fill}%
\end{pgfscope}%
\begin{pgfscope}%
\pgfpathrectangle{\pgfqpoint{0.407592in}{0.222904in}}{\pgfqpoint{1.411163in}{1.411163in}}%
\pgfusepath{clip}%
\pgfsetbuttcap%
\pgfsetroundjoin%
\definecolor{currentfill}{rgb}{0.749020,0.772549,0.792157}%
\pgfsetfillcolor{currentfill}%
\pgfsetfillopacity{0.160000}%
\pgfsetlinewidth{0.000000pt}%
\definecolor{currentstroke}{rgb}{0.000000,0.000000,0.000000}%
\pgfsetstrokecolor{currentstroke}%
\pgfsetdash{}{0pt}%
\pgfpathmoveto{\pgfqpoint{1.363722in}{1.243892in}}%
\pgfpathlineto{\pgfqpoint{1.370352in}{1.259246in}}%
\pgfpathlineto{\pgfqpoint{1.376972in}{1.274577in}}%
\pgfpathlineto{\pgfqpoint{1.345943in}{1.280529in}}%
\pgfpathlineto{\pgfqpoint{1.313978in}{1.285709in}}%
\pgfpathlineto{\pgfqpoint{1.308831in}{1.270179in}}%
\pgfpathlineto{\pgfqpoint{1.303676in}{1.254623in}}%
\pgfpathlineto{\pgfqpoint{1.334150in}{1.249629in}}%
\pgfpathlineto{\pgfqpoint{1.363722in}{1.243892in}}%
\pgfpathclose%
\pgfusepath{fill}%
\end{pgfscope}%
\begin{pgfscope}%
\pgfpathrectangle{\pgfqpoint{0.407592in}{0.222904in}}{\pgfqpoint{1.411163in}{1.411163in}}%
\pgfusepath{clip}%
\pgfsetbuttcap%
\pgfsetroundjoin%
\definecolor{currentfill}{rgb}{0.749020,0.772549,0.792157}%
\pgfsetfillcolor{currentfill}%
\pgfsetfillopacity{0.160000}%
\pgfsetlinewidth{0.000000pt}%
\definecolor{currentstroke}{rgb}{0.000000,0.000000,0.000000}%
\pgfsetstrokecolor{currentstroke}%
\pgfsetdash{}{0pt}%
\pgfpathmoveto{\pgfqpoint{1.218634in}{1.167541in}}%
\pgfpathlineto{\pgfqpoint{1.222292in}{1.183410in}}%
\pgfpathlineto{\pgfqpoint{1.225944in}{1.199251in}}%
\pgfpathlineto{\pgfqpoint{1.196612in}{1.201733in}}%
\pgfpathlineto{\pgfqpoint{1.166878in}{1.203490in}}%
\pgfpathlineto{\pgfqpoint{1.164819in}{1.187557in}}%
\pgfpathlineto{\pgfqpoint{1.162756in}{1.171595in}}%
\pgfpathlineto{\pgfqpoint{1.190887in}{1.169915in}}%
\pgfpathlineto{\pgfqpoint{1.218634in}{1.167541in}}%
\pgfpathclose%
\pgfusepath{fill}%
\end{pgfscope}%
\begin{pgfscope}%
\pgfpathrectangle{\pgfqpoint{0.407592in}{0.222904in}}{\pgfqpoint{1.411163in}{1.411163in}}%
\pgfusepath{clip}%
\pgfsetbuttcap%
\pgfsetroundjoin%
\definecolor{currentfill}{rgb}{0.749020,0.772549,0.792157}%
\pgfsetfillcolor{currentfill}%
\pgfsetfillopacity{0.160000}%
\pgfsetlinewidth{0.000000pt}%
\definecolor{currentstroke}{rgb}{0.000000,0.000000,0.000000}%
\pgfsetstrokecolor{currentstroke}%
\pgfsetdash{}{0pt}%
\pgfpathmoveto{\pgfqpoint{1.467710in}{1.320544in}}%
\pgfpathlineto{\pgfqpoint{1.475675in}{1.335525in}}%
\pgfpathlineto{\pgfqpoint{1.483630in}{1.350488in}}%
\pgfpathlineto{\pgfqpoint{1.450708in}{1.358697in}}%
\pgfpathlineto{\pgfqpoint{1.416489in}{1.366090in}}%
\pgfpathlineto{\pgfqpoint{1.409927in}{1.350894in}}%
\pgfpathlineto{\pgfqpoint{1.403355in}{1.335676in}}%
\pgfpathlineto{\pgfqpoint{1.436161in}{1.328505in}}%
\pgfpathlineto{\pgfqpoint{1.467710in}{1.320544in}}%
\pgfpathclose%
\pgfusepath{fill}%
\end{pgfscope}%
\begin{pgfscope}%
\pgfpathrectangle{\pgfqpoint{0.407592in}{0.222904in}}{\pgfqpoint{1.411163in}{1.411163in}}%
\pgfusepath{clip}%
\pgfsetbuttcap%
\pgfsetroundjoin%
\definecolor{currentfill}{rgb}{0.749020,0.772549,0.792157}%
\pgfsetfillcolor{currentfill}%
\pgfsetfillopacity{0.160000}%
\pgfsetlinewidth{0.000000pt}%
\definecolor{currentstroke}{rgb}{0.000000,0.000000,0.000000}%
\pgfsetstrokecolor{currentstroke}%
\pgfsetdash{}{0pt}%
\pgfpathmoveto{\pgfqpoint{1.282971in}{1.192142in}}%
\pgfpathlineto{\pgfqpoint{1.288160in}{1.207801in}}%
\pgfpathlineto{\pgfqpoint{1.293341in}{1.223435in}}%
\pgfpathlineto{\pgfqpoint{1.263606in}{1.227510in}}%
\pgfpathlineto{\pgfqpoint{1.233227in}{1.230846in}}%
\pgfpathlineto{\pgfqpoint{1.229589in}{1.215063in}}%
\pgfpathlineto{\pgfqpoint{1.225944in}{1.199251in}}%
\pgfpathlineto{\pgfqpoint{1.254766in}{1.196051in}}%
\pgfpathlineto{\pgfqpoint{1.282971in}{1.192142in}}%
\pgfpathclose%
\pgfusepath{fill}%
\end{pgfscope}%
\begin{pgfscope}%
\pgfpathrectangle{\pgfqpoint{0.407592in}{0.222904in}}{\pgfqpoint{1.411163in}{1.411163in}}%
\pgfusepath{clip}%
\pgfsetbuttcap%
\pgfsetroundjoin%
\definecolor{currentfill}{rgb}{0.749020,0.772549,0.792157}%
\pgfsetfillcolor{currentfill}%
\pgfsetfillopacity{0.160000}%
\pgfsetlinewidth{0.000000pt}%
\definecolor{currentstroke}{rgb}{0.000000,0.000000,0.000000}%
\pgfsetstrokecolor{currentstroke}%
\pgfsetdash{}{0pt}%
\pgfpathmoveto{\pgfqpoint{1.048487in}{1.171314in}}%
\pgfpathlineto{\pgfqpoint{1.047299in}{1.187270in}}%
\pgfpathlineto{\pgfqpoint{1.046114in}{1.203196in}}%
\pgfpathlineto{\pgfqpoint{1.016015in}{1.201293in}}%
\pgfpathlineto{\pgfqpoint{0.986173in}{1.198666in}}%
\pgfpathlineto{\pgfqpoint{0.988973in}{1.182838in}}%
\pgfpathlineto{\pgfqpoint{0.991779in}{1.166982in}}%
\pgfpathlineto{\pgfqpoint{1.020010in}{1.169494in}}%
\pgfpathlineto{\pgfqpoint{1.048487in}{1.171314in}}%
\pgfpathclose%
\pgfusepath{fill}%
\end{pgfscope}%
\begin{pgfscope}%
\pgfpathrectangle{\pgfqpoint{0.407592in}{0.222904in}}{\pgfqpoint{1.411163in}{1.411163in}}%
\pgfusepath{clip}%
\pgfsetbuttcap%
\pgfsetroundjoin%
\definecolor{currentfill}{rgb}{0.749020,0.772549,0.792157}%
\pgfsetfillcolor{currentfill}%
\pgfsetfillopacity{0.160000}%
\pgfsetlinewidth{0.000000pt}%
\definecolor{currentstroke}{rgb}{0.000000,0.000000,0.000000}%
\pgfsetstrokecolor{currentstroke}%
\pgfsetdash{}{0pt}%
\pgfpathmoveto{\pgfqpoint{0.910233in}{1.253682in}}%
\pgfpathlineto{\pgfqpoint{0.905885in}{1.269220in}}%
\pgfpathlineto{\pgfqpoint{0.901543in}{1.284732in}}%
\pgfpathlineto{\pgfqpoint{0.868501in}{1.279396in}}%
\pgfpathlineto{\pgfqpoint{0.836238in}{1.273292in}}%
\pgfpathlineto{\pgfqpoint{0.842105in}{1.257983in}}%
\pgfpathlineto{\pgfqpoint{0.847980in}{1.242653in}}%
\pgfpathlineto{\pgfqpoint{0.878731in}{1.248537in}}%
\pgfpathlineto{\pgfqpoint{0.910233in}{1.253682in}}%
\pgfpathclose%
\pgfusepath{fill}%
\end{pgfscope}%
\begin{pgfscope}%
\pgfpathrectangle{\pgfqpoint{0.407592in}{0.222904in}}{\pgfqpoint{1.411163in}{1.411163in}}%
\pgfusepath{clip}%
\pgfsetbuttcap%
\pgfsetroundjoin%
\definecolor{currentfill}{rgb}{0.749020,0.772549,0.792157}%
\pgfsetfillcolor{currentfill}%
\pgfsetfillopacity{0.160000}%
\pgfsetlinewidth{0.000000pt}%
\definecolor{currentstroke}{rgb}{0.000000,0.000000,0.000000}%
\pgfsetstrokecolor{currentstroke}%
\pgfsetdash{}{0pt}%
\pgfpathmoveto{\pgfqpoint{0.986173in}{1.198666in}}%
\pgfpathlineto{\pgfqpoint{0.983378in}{1.214465in}}%
\pgfpathlineto{\pgfqpoint{0.980587in}{1.230237in}}%
\pgfpathlineto{\pgfqpoint{0.949521in}{1.226751in}}%
\pgfpathlineto{\pgfqpoint{0.918952in}{1.222529in}}%
\pgfpathlineto{\pgfqpoint{0.923322in}{1.206914in}}%
\pgfpathlineto{\pgfqpoint{0.927699in}{1.191274in}}%
\pgfpathlineto{\pgfqpoint{0.956698in}{1.195323in}}%
\pgfpathlineto{\pgfqpoint{0.986173in}{1.198666in}}%
\pgfpathclose%
\pgfusepath{fill}%
\end{pgfscope}%
\begin{pgfscope}%
\pgfpathrectangle{\pgfqpoint{0.407592in}{0.222904in}}{\pgfqpoint{1.411163in}{1.411163in}}%
\pgfusepath{clip}%
\pgfsetbuttcap%
\pgfsetroundjoin%
\definecolor{currentfill}{rgb}{0.749020,0.772549,0.792157}%
\pgfsetfillcolor{currentfill}%
\pgfsetfillopacity{0.160000}%
\pgfsetlinewidth{0.000000pt}%
\definecolor{currentstroke}{rgb}{0.000000,0.000000,0.000000}%
\pgfsetstrokecolor{currentstroke}%
\pgfsetdash{}{0pt}%
\pgfpathmoveto{\pgfqpoint{0.812857in}{1.334301in}}%
\pgfpathlineto{\pgfqpoint{0.807033in}{1.349499in}}%
\pgfpathlineto{\pgfqpoint{0.801217in}{1.364674in}}%
\pgfpathlineto{\pgfqpoint{0.765427in}{1.357115in}}%
\pgfpathlineto{\pgfqpoint{0.730764in}{1.348745in}}%
\pgfpathlineto{\pgfqpoint{0.738037in}{1.333808in}}%
\pgfpathlineto{\pgfqpoint{0.745320in}{1.318853in}}%
\pgfpathlineto{\pgfqpoint{0.778541in}{1.326970in}}%
\pgfpathlineto{\pgfqpoint{0.812857in}{1.334301in}}%
\pgfpathclose%
\pgfusepath{fill}%
\end{pgfscope}%
\begin{pgfscope}%
\pgfpathrectangle{\pgfqpoint{0.407592in}{0.222904in}}{\pgfqpoint{1.411163in}{1.411163in}}%
\pgfusepath{clip}%
\pgfsetbuttcap%
\pgfsetroundjoin%
\definecolor{currentfill}{rgb}{0.749020,0.772549,0.792157}%
\pgfsetfillcolor{currentfill}%
\pgfsetfillopacity{0.160000}%
\pgfsetlinewidth{0.000000pt}%
\definecolor{currentstroke}{rgb}{0.000000,0.000000,0.000000}%
\pgfsetstrokecolor{currentstroke}%
\pgfsetdash{}{0pt}%
\pgfpathmoveto{\pgfqpoint{1.618793in}{1.448621in}}%
\pgfpathlineto{\pgfqpoint{1.627969in}{1.463170in}}%
\pgfpathlineto{\pgfqpoint{1.592260in}{1.474184in}}%
\pgfpathlineto{\pgfqpoint{1.554781in}{1.484319in}}%
\pgfpathlineto{\pgfqpoint{1.546915in}{1.469523in}}%
\pgfpathlineto{\pgfqpoint{1.583727in}{1.459505in}}%
\pgfusepath{fill}%
\end{pgfscope}%
\begin{pgfscope}%
\pgfpathrectangle{\pgfqpoint{0.407592in}{0.222904in}}{\pgfqpoint{1.411163in}{1.411163in}}%
\pgfusepath{clip}%
\pgfsetbuttcap%
\pgfsetroundjoin%
\definecolor{currentfill}{rgb}{0.749020,0.772549,0.792157}%
\pgfsetfillcolor{currentfill}%
\pgfsetfillopacity{0.160000}%
\pgfsetlinewidth{0.000000pt}%
\definecolor{currentstroke}{rgb}{0.000000,0.000000,0.000000}%
\pgfsetstrokecolor{currentstroke}%
\pgfsetdash{}{0pt}%
\pgfpathmoveto{\pgfqpoint{1.158617in}{1.139581in}}%
\pgfpathlineto{\pgfqpoint{1.160688in}{1.155603in}}%
\pgfpathlineto{\pgfqpoint{1.162756in}{1.171595in}}%
\pgfpathlineto{\pgfqpoint{1.134345in}{1.172577in}}%
\pgfpathlineto{\pgfqpoint{1.105760in}{1.172858in}}%
\pgfpathlineto{\pgfqpoint{1.105318in}{1.156836in}}%
\pgfpathlineto{\pgfqpoint{1.104874in}{1.140784in}}%
\pgfpathlineto{\pgfqpoint{1.131829in}{1.140517in}}%
\pgfpathlineto{\pgfqpoint{1.158617in}{1.139581in}}%
\pgfpathclose%
\pgfusepath{fill}%
\end{pgfscope}%
\begin{pgfscope}%
\pgfpathrectangle{\pgfqpoint{0.407592in}{0.222904in}}{\pgfqpoint{1.411163in}{1.411163in}}%
\pgfusepath{clip}%
\pgfsetbuttcap%
\pgfsetroundjoin%
\definecolor{currentfill}{rgb}{0.749020,0.772549,0.792157}%
\pgfsetfillcolor{currentfill}%
\pgfsetfillopacity{0.160000}%
\pgfsetlinewidth{0.000000pt}%
\definecolor{currentstroke}{rgb}{0.000000,0.000000,0.000000}%
\pgfsetstrokecolor{currentstroke}%
\pgfsetdash{}{0pt}%
\pgfpathmoveto{\pgfqpoint{1.104874in}{1.140784in}}%
\pgfpathlineto{\pgfqpoint{1.105318in}{1.156836in}}%
\pgfpathlineto{\pgfqpoint{1.105760in}{1.172858in}}%
\pgfpathlineto{\pgfqpoint{1.077106in}{1.172436in}}%
\pgfpathlineto{\pgfqpoint{1.048487in}{1.171314in}}%
\pgfpathlineto{\pgfqpoint{1.049677in}{1.155328in}}%
\pgfpathlineto{\pgfqpoint{1.050869in}{1.139313in}}%
\pgfpathlineto{\pgfqpoint{1.077854in}{1.140383in}}%
\pgfpathlineto{\pgfqpoint{1.104874in}{1.140784in}}%
\pgfpathclose%
\pgfusepath{fill}%
\end{pgfscope}%
\begin{pgfscope}%
\pgfpathrectangle{\pgfqpoint{0.407592in}{0.222904in}}{\pgfqpoint{1.411163in}{1.411163in}}%
\pgfusepath{clip}%
\pgfsetbuttcap%
\pgfsetroundjoin%
\definecolor{currentfill}{rgb}{0.749020,0.772549,0.792157}%
\pgfsetfillcolor{currentfill}%
\pgfsetfillopacity{0.160000}%
\pgfsetlinewidth{0.000000pt}%
\definecolor{currentstroke}{rgb}{0.000000,0.000000,0.000000}%
\pgfsetstrokecolor{currentstroke}%
\pgfsetdash{}{0pt}%
\pgfpathmoveto{\pgfqpoint{1.451750in}{1.290524in}}%
\pgfpathlineto{\pgfqpoint{1.459735in}{1.305543in}}%
\pgfpathlineto{\pgfqpoint{1.467710in}{1.320544in}}%
\pgfpathlineto{\pgfqpoint{1.436161in}{1.328505in}}%
\pgfpathlineto{\pgfqpoint{1.403355in}{1.335676in}}%
\pgfpathlineto{\pgfqpoint{1.396774in}{1.320435in}}%
\pgfpathlineto{\pgfqpoint{1.390183in}{1.305172in}}%
\pgfpathlineto{\pgfqpoint{1.421574in}{1.298229in}}%
\pgfpathlineto{\pgfqpoint{1.451750in}{1.290524in}}%
\pgfpathclose%
\pgfusepath{fill}%
\end{pgfscope}%
\begin{pgfscope}%
\pgfpathrectangle{\pgfqpoint{0.407592in}{0.222904in}}{\pgfqpoint{1.411163in}{1.411163in}}%
\pgfusepath{clip}%
\pgfsetbuttcap%
\pgfsetroundjoin%
\definecolor{currentfill}{rgb}{0.749020,0.772549,0.792157}%
\pgfsetfillcolor{currentfill}%
\pgfsetfillopacity{0.160000}%
\pgfsetlinewidth{0.000000pt}%
\definecolor{currentstroke}{rgb}{0.000000,0.000000,0.000000}%
\pgfsetstrokecolor{currentstroke}%
\pgfsetdash{}{0pt}%
\pgfpathmoveto{\pgfqpoint{0.672897in}{1.467584in}}%
\pgfpathlineto{\pgfqpoint{0.665703in}{1.482357in}}%
\pgfpathlineto{\pgfqpoint{0.625987in}{1.472043in}}%
\pgfpathlineto{\pgfqpoint{0.587852in}{1.460857in}}%
\pgfpathlineto{\pgfqpoint{0.596432in}{1.446335in}}%
\pgfpathlineto{\pgfqpoint{0.633884in}{1.457390in}}%
\pgfusepath{fill}%
\end{pgfscope}%
\begin{pgfscope}%
\pgfpathrectangle{\pgfqpoint{0.407592in}{0.222904in}}{\pgfqpoint{1.411163in}{1.411163in}}%
\pgfusepath{clip}%
\pgfsetbuttcap%
\pgfsetroundjoin%
\definecolor{currentfill}{rgb}{0.749020,0.772549,0.792157}%
\pgfsetfillcolor{currentfill}%
\pgfsetfillopacity{0.160000}%
\pgfsetlinewidth{0.000000pt}%
\definecolor{currentstroke}{rgb}{0.000000,0.000000,0.000000}%
\pgfsetstrokecolor{currentstroke}%
\pgfsetdash{}{0pt}%
\pgfpathmoveto{\pgfqpoint{1.350432in}{1.213116in}}%
\pgfpathlineto{\pgfqpoint{1.357082in}{1.228515in}}%
\pgfpathlineto{\pgfqpoint{1.363722in}{1.243892in}}%
\pgfpathlineto{\pgfqpoint{1.334150in}{1.249629in}}%
\pgfpathlineto{\pgfqpoint{1.303676in}{1.254623in}}%
\pgfpathlineto{\pgfqpoint{1.298513in}{1.239042in}}%
\pgfpathlineto{\pgfqpoint{1.293341in}{1.223435in}}%
\pgfpathlineto{\pgfqpoint{1.322319in}{1.218632in}}%
\pgfpathlineto{\pgfqpoint{1.350432in}{1.213116in}}%
\pgfpathclose%
\pgfusepath{fill}%
\end{pgfscope}%
\begin{pgfscope}%
\pgfpathrectangle{\pgfqpoint{0.407592in}{0.222904in}}{\pgfqpoint{1.411163in}{1.411163in}}%
\pgfusepath{clip}%
\pgfsetbuttcap%
\pgfsetroundjoin%
\definecolor{currentfill}{rgb}{0.749020,0.772549,0.792157}%
\pgfsetfillcolor{currentfill}%
\pgfsetfillopacity{0.160000}%
\pgfsetlinewidth{0.000000pt}%
\definecolor{currentstroke}{rgb}{0.000000,0.000000,0.000000}%
\pgfsetstrokecolor{currentstroke}%
\pgfsetdash{}{0pt}%
\pgfpathmoveto{\pgfqpoint{1.600415in}{1.419480in}}%
\pgfpathlineto{\pgfqpoint{1.609608in}{1.434057in}}%
\pgfpathlineto{\pgfqpoint{1.618793in}{1.448621in}}%
\pgfpathlineto{\pgfqpoint{1.583727in}{1.459505in}}%
\pgfpathlineto{\pgfqpoint{1.546915in}{1.469523in}}%
\pgfpathlineto{\pgfqpoint{1.539039in}{1.454708in}}%
\pgfpathlineto{\pgfqpoint{1.531153in}{1.439875in}}%
\pgfpathlineto{\pgfqpoint{1.566634in}{1.430099in}}%
\pgfpathlineto{\pgfqpoint{1.600415in}{1.419480in}}%
\pgfpathclose%
\pgfusepath{fill}%
\end{pgfscope}%
\begin{pgfscope}%
\pgfpathrectangle{\pgfqpoint{0.407592in}{0.222904in}}{\pgfqpoint{1.411163in}{1.411163in}}%
\pgfusepath{clip}%
\pgfsetbuttcap%
\pgfsetroundjoin%
\definecolor{currentfill}{rgb}{0.749020,0.772549,0.792157}%
\pgfsetfillcolor{currentfill}%
\pgfsetfillopacity{0.160000}%
\pgfsetlinewidth{0.000000pt}%
\definecolor{currentstroke}{rgb}{0.000000,0.000000,0.000000}%
\pgfsetstrokecolor{currentstroke}%
\pgfsetdash{}{0pt}%
\pgfpathmoveto{\pgfqpoint{0.918952in}{1.222529in}}%
\pgfpathlineto{\pgfqpoint{0.914589in}{1.238118in}}%
\pgfpathlineto{\pgfqpoint{0.910233in}{1.253682in}}%
\pgfpathlineto{\pgfqpoint{0.878731in}{1.248537in}}%
\pgfpathlineto{\pgfqpoint{0.847980in}{1.242653in}}%
\pgfpathlineto{\pgfqpoint{0.853864in}{1.227300in}}%
\pgfpathlineto{\pgfqpoint{0.859756in}{1.211925in}}%
\pgfpathlineto{\pgfqpoint{0.888992in}{1.217582in}}%
\pgfpathlineto{\pgfqpoint{0.918952in}{1.222529in}}%
\pgfpathclose%
\pgfusepath{fill}%
\end{pgfscope}%
\begin{pgfscope}%
\pgfpathrectangle{\pgfqpoint{0.407592in}{0.222904in}}{\pgfqpoint{1.411163in}{1.411163in}}%
\pgfusepath{clip}%
\pgfsetbuttcap%
\pgfsetroundjoin%
\definecolor{currentfill}{rgb}{0.749020,0.772549,0.792157}%
\pgfsetfillcolor{currentfill}%
\pgfsetfillopacity{0.160000}%
\pgfsetlinewidth{0.000000pt}%
\definecolor{currentstroke}{rgb}{0.000000,0.000000,0.000000}%
\pgfsetstrokecolor{currentstroke}%
\pgfsetdash{}{0pt}%
\pgfpathmoveto{\pgfqpoint{1.211298in}{1.135718in}}%
\pgfpathlineto{\pgfqpoint{1.214969in}{1.151644in}}%
\pgfpathlineto{\pgfqpoint{1.218634in}{1.167541in}}%
\pgfpathlineto{\pgfqpoint{1.190887in}{1.169915in}}%
\pgfpathlineto{\pgfqpoint{1.162756in}{1.171595in}}%
\pgfpathlineto{\pgfqpoint{1.160688in}{1.155603in}}%
\pgfpathlineto{\pgfqpoint{1.158617in}{1.139581in}}%
\pgfpathlineto{\pgfqpoint{1.185140in}{1.137980in}}%
\pgfpathlineto{\pgfqpoint{1.211298in}{1.135718in}}%
\pgfpathclose%
\pgfusepath{fill}%
\end{pgfscope}%
\begin{pgfscope}%
\pgfpathrectangle{\pgfqpoint{0.407592in}{0.222904in}}{\pgfqpoint{1.411163in}{1.411163in}}%
\pgfusepath{clip}%
\pgfsetbuttcap%
\pgfsetroundjoin%
\definecolor{currentfill}{rgb}{0.749020,0.772549,0.792157}%
\pgfsetfillcolor{currentfill}%
\pgfsetfillopacity{0.160000}%
\pgfsetlinewidth{0.000000pt}%
\definecolor{currentstroke}{rgb}{0.000000,0.000000,0.000000}%
\pgfsetstrokecolor{currentstroke}%
\pgfsetdash{}{0pt}%
\pgfpathmoveto{\pgfqpoint{1.272566in}{1.160745in}}%
\pgfpathlineto{\pgfqpoint{1.277773in}{1.176457in}}%
\pgfpathlineto{\pgfqpoint{1.282971in}{1.192142in}}%
\pgfpathlineto{\pgfqpoint{1.254766in}{1.196051in}}%
\pgfpathlineto{\pgfqpoint{1.225944in}{1.199251in}}%
\pgfpathlineto{\pgfqpoint{1.222292in}{1.183410in}}%
\pgfpathlineto{\pgfqpoint{1.218634in}{1.167541in}}%
\pgfpathlineto{\pgfqpoint{1.245895in}{1.164482in}}%
\pgfpathlineto{\pgfqpoint{1.272566in}{1.160745in}}%
\pgfpathclose%
\pgfusepath{fill}%
\end{pgfscope}%
\begin{pgfscope}%
\pgfpathrectangle{\pgfqpoint{0.407592in}{0.222904in}}{\pgfqpoint{1.411163in}{1.411163in}}%
\pgfusepath{clip}%
\pgfsetbuttcap%
\pgfsetroundjoin%
\definecolor{currentfill}{rgb}{0.749020,0.772549,0.792157}%
\pgfsetfillcolor{currentfill}%
\pgfsetfillopacity{0.160000}%
\pgfsetlinewidth{0.000000pt}%
\definecolor{currentstroke}{rgb}{0.000000,0.000000,0.000000}%
\pgfsetstrokecolor{currentstroke}%
\pgfsetdash{}{0pt}%
\pgfpathmoveto{\pgfqpoint{0.824531in}{1.303841in}}%
\pgfpathlineto{\pgfqpoint{0.818690in}{1.319082in}}%
\pgfpathlineto{\pgfqpoint{0.812857in}{1.334301in}}%
\pgfpathlineto{\pgfqpoint{0.778541in}{1.326970in}}%
\pgfpathlineto{\pgfqpoint{0.745320in}{1.318853in}}%
\pgfpathlineto{\pgfqpoint{0.752611in}{1.303880in}}%
\pgfpathlineto{\pgfqpoint{0.759911in}{1.288888in}}%
\pgfpathlineto{\pgfqpoint{0.791691in}{1.296744in}}%
\pgfpathlineto{\pgfqpoint{0.824531in}{1.303841in}}%
\pgfpathclose%
\pgfusepath{fill}%
\end{pgfscope}%
\begin{pgfscope}%
\pgfpathrectangle{\pgfqpoint{0.407592in}{0.222904in}}{\pgfqpoint{1.411163in}{1.411163in}}%
\pgfusepath{clip}%
\pgfsetbuttcap%
\pgfsetroundjoin%
\definecolor{currentfill}{rgb}{0.749020,0.772549,0.792157}%
\pgfsetfillcolor{currentfill}%
\pgfsetfillopacity{0.160000}%
\pgfsetlinewidth{0.000000pt}%
\definecolor{currentstroke}{rgb}{0.000000,0.000000,0.000000}%
\pgfsetstrokecolor{currentstroke}%
\pgfsetdash{}{0pt}%
\pgfpathmoveto{\pgfqpoint{1.050869in}{1.139313in}}%
\pgfpathlineto{\pgfqpoint{1.049677in}{1.155328in}}%
\pgfpathlineto{\pgfqpoint{1.048487in}{1.171314in}}%
\pgfpathlineto{\pgfqpoint{1.020010in}{1.169494in}}%
\pgfpathlineto{\pgfqpoint{0.991779in}{1.166982in}}%
\pgfpathlineto{\pgfqpoint{0.994589in}{1.151098in}}%
\pgfpathlineto{\pgfqpoint{0.997404in}{1.135186in}}%
\pgfpathlineto{\pgfqpoint{1.024019in}{1.137579in}}%
\pgfpathlineto{\pgfqpoint{1.050869in}{1.139313in}}%
\pgfpathclose%
\pgfusepath{fill}%
\end{pgfscope}%
\begin{pgfscope}%
\pgfpathrectangle{\pgfqpoint{0.407592in}{0.222904in}}{\pgfqpoint{1.411163in}{1.411163in}}%
\pgfusepath{clip}%
\pgfsetbuttcap%
\pgfsetroundjoin%
\definecolor{currentfill}{rgb}{0.749020,0.772549,0.792157}%
\pgfsetfillcolor{currentfill}%
\pgfsetfillopacity{0.160000}%
\pgfsetlinewidth{0.000000pt}%
\definecolor{currentstroke}{rgb}{0.000000,0.000000,0.000000}%
\pgfsetstrokecolor{currentstroke}%
\pgfsetdash{}{0pt}%
\pgfpathmoveto{\pgfqpoint{0.991779in}{1.166982in}}%
\pgfpathlineto{\pgfqpoint{0.988973in}{1.182838in}}%
\pgfpathlineto{\pgfqpoint{0.986173in}{1.198666in}}%
\pgfpathlineto{\pgfqpoint{0.956698in}{1.195323in}}%
\pgfpathlineto{\pgfqpoint{0.927699in}{1.191274in}}%
\pgfpathlineto{\pgfqpoint{0.932084in}{1.175607in}}%
\pgfpathlineto{\pgfqpoint{0.936476in}{1.159915in}}%
\pgfpathlineto{\pgfqpoint{0.963899in}{1.163786in}}%
\pgfpathlineto{\pgfqpoint{0.991779in}{1.166982in}}%
\pgfpathclose%
\pgfusepath{fill}%
\end{pgfscope}%
\begin{pgfscope}%
\pgfpathrectangle{\pgfqpoint{0.407592in}{0.222904in}}{\pgfqpoint{1.411163in}{1.411163in}}%
\pgfusepath{clip}%
\pgfsetbuttcap%
\pgfsetroundjoin%
\definecolor{currentfill}{rgb}{0.749020,0.772549,0.792157}%
\pgfsetfillcolor{currentfill}%
\pgfsetfillopacity{0.160000}%
\pgfsetlinewidth{0.000000pt}%
\definecolor{currentstroke}{rgb}{0.000000,0.000000,0.000000}%
\pgfsetstrokecolor{currentstroke}%
\pgfsetdash{}{0pt}%
\pgfpathmoveto{\pgfqpoint{0.687311in}{1.437983in}}%
\pgfpathlineto{\pgfqpoint{0.680099in}{1.452792in}}%
\pgfpathlineto{\pgfqpoint{0.672897in}{1.467584in}}%
\pgfpathlineto{\pgfqpoint{0.633884in}{1.457390in}}%
\pgfpathlineto{\pgfqpoint{0.596432in}{1.446335in}}%
\pgfpathlineto{\pgfqpoint{0.605020in}{1.431799in}}%
\pgfpathlineto{\pgfqpoint{0.613616in}{1.417250in}}%
\pgfpathlineto{\pgfqpoint{0.649702in}{1.428035in}}%
\pgfpathlineto{\pgfqpoint{0.687311in}{1.437983in}}%
\pgfpathclose%
\pgfusepath{fill}%
\end{pgfscope}%
\begin{pgfscope}%
\pgfpathrectangle{\pgfqpoint{0.407592in}{0.222904in}}{\pgfqpoint{1.411163in}{1.411163in}}%
\pgfusepath{clip}%
\pgfsetbuttcap%
\pgfsetroundjoin%
\definecolor{currentfill}{rgb}{0.749020,0.772549,0.792157}%
\pgfsetfillcolor{currentfill}%
\pgfsetfillopacity{0.160000}%
\pgfsetlinewidth{0.000000pt}%
\definecolor{currentstroke}{rgb}{0.000000,0.000000,0.000000}%
\pgfsetstrokecolor{currentstroke}%
\pgfsetdash{}{0pt}%
\pgfpathmoveto{\pgfqpoint{1.582000in}{1.390282in}}%
\pgfpathlineto{\pgfqpoint{1.591212in}{1.404888in}}%
\pgfpathlineto{\pgfqpoint{1.600415in}{1.419480in}}%
\pgfpathlineto{\pgfqpoint{1.566634in}{1.430099in}}%
\pgfpathlineto{\pgfqpoint{1.531153in}{1.439875in}}%
\pgfpathlineto{\pgfqpoint{1.523257in}{1.425024in}}%
\pgfpathlineto{\pgfqpoint{1.515351in}{1.410154in}}%
\pgfpathlineto{\pgfqpoint{1.549502in}{1.400627in}}%
\pgfpathlineto{\pgfqpoint{1.582000in}{1.390282in}}%
\pgfpathclose%
\pgfusepath{fill}%
\end{pgfscope}%
\begin{pgfscope}%
\pgfpathrectangle{\pgfqpoint{0.407592in}{0.222904in}}{\pgfqpoint{1.411163in}{1.411163in}}%
\pgfusepath{clip}%
\pgfsetbuttcap%
\pgfsetroundjoin%
\definecolor{currentfill}{rgb}{0.749020,0.772549,0.792157}%
\pgfsetfillcolor{currentfill}%
\pgfsetfillopacity{0.160000}%
\pgfsetlinewidth{0.000000pt}%
\definecolor{currentstroke}{rgb}{0.000000,0.000000,0.000000}%
\pgfsetstrokecolor{currentstroke}%
\pgfsetdash{}{0pt}%
\pgfpathmoveto{\pgfqpoint{1.435749in}{1.260429in}}%
\pgfpathlineto{\pgfqpoint{1.443755in}{1.275486in}}%
\pgfpathlineto{\pgfqpoint{1.451750in}{1.290524in}}%
\pgfpathlineto{\pgfqpoint{1.421574in}{1.298229in}}%
\pgfpathlineto{\pgfqpoint{1.390183in}{1.305172in}}%
\pgfpathlineto{\pgfqpoint{1.383583in}{1.289886in}}%
\pgfpathlineto{\pgfqpoint{1.376972in}{1.274577in}}%
\pgfpathlineto{\pgfqpoint{1.406946in}{1.267871in}}%
\pgfpathlineto{\pgfqpoint{1.435749in}{1.260429in}}%
\pgfpathclose%
\pgfusepath{fill}%
\end{pgfscope}%
\begin{pgfscope}%
\pgfpathrectangle{\pgfqpoint{0.407592in}{0.222904in}}{\pgfqpoint{1.411163in}{1.411163in}}%
\pgfusepath{clip}%
\pgfsetbuttcap%
\pgfsetroundjoin%
\definecolor{currentfill}{rgb}{0.749020,0.772549,0.792157}%
\pgfsetfillcolor{currentfill}%
\pgfsetfillopacity{0.160000}%
\pgfsetlinewidth{0.000000pt}%
\definecolor{currentstroke}{rgb}{0.000000,0.000000,0.000000}%
\pgfsetstrokecolor{currentstroke}%
\pgfsetdash{}{0pt}%
\pgfpathmoveto{\pgfqpoint{1.154463in}{1.107447in}}%
\pgfpathlineto{\pgfqpoint{1.156542in}{1.123529in}}%
\pgfpathlineto{\pgfqpoint{1.158617in}{1.139581in}}%
\pgfpathlineto{\pgfqpoint{1.131829in}{1.140517in}}%
\pgfpathlineto{\pgfqpoint{1.104874in}{1.140784in}}%
\pgfpathlineto{\pgfqpoint{1.104430in}{1.124702in}}%
\pgfpathlineto{\pgfqpoint{1.103985in}{1.108589in}}%
\pgfpathlineto{\pgfqpoint{1.129303in}{1.108335in}}%
\pgfpathlineto{\pgfqpoint{1.154463in}{1.107447in}}%
\pgfpathclose%
\pgfusepath{fill}%
\end{pgfscope}%
\begin{pgfscope}%
\pgfpathrectangle{\pgfqpoint{0.407592in}{0.222904in}}{\pgfqpoint{1.411163in}{1.411163in}}%
\pgfusepath{clip}%
\pgfsetbuttcap%
\pgfsetroundjoin%
\definecolor{currentfill}{rgb}{0.749020,0.772549,0.792157}%
\pgfsetfillcolor{currentfill}%
\pgfsetfillopacity{0.160000}%
\pgfsetlinewidth{0.000000pt}%
\definecolor{currentstroke}{rgb}{0.000000,0.000000,0.000000}%
\pgfsetstrokecolor{currentstroke}%
\pgfsetdash{}{0pt}%
\pgfpathmoveto{\pgfqpoint{1.103985in}{1.108589in}}%
\pgfpathlineto{\pgfqpoint{1.104430in}{1.124702in}}%
\pgfpathlineto{\pgfqpoint{1.104874in}{1.140784in}}%
\pgfpathlineto{\pgfqpoint{1.077854in}{1.140383in}}%
\pgfpathlineto{\pgfqpoint{1.050869in}{1.139313in}}%
\pgfpathlineto{\pgfqpoint{1.052063in}{1.123268in}}%
\pgfpathlineto{\pgfqpoint{1.053260in}{1.107192in}}%
\pgfpathlineto{\pgfqpoint{1.078606in}{1.108208in}}%
\pgfpathlineto{\pgfqpoint{1.103985in}{1.108589in}}%
\pgfpathclose%
\pgfusepath{fill}%
\end{pgfscope}%
\begin{pgfscope}%
\pgfpathrectangle{\pgfqpoint{0.407592in}{0.222904in}}{\pgfqpoint{1.411163in}{1.411163in}}%
\pgfusepath{clip}%
\pgfsetbuttcap%
\pgfsetroundjoin%
\definecolor{currentfill}{rgb}{0.749020,0.772549,0.792157}%
\pgfsetfillcolor{currentfill}%
\pgfsetfillopacity{0.160000}%
\pgfsetlinewidth{0.000000pt}%
\definecolor{currentstroke}{rgb}{0.000000,0.000000,0.000000}%
\pgfsetstrokecolor{currentstroke}%
\pgfsetdash{}{0pt}%
\pgfpathmoveto{\pgfqpoint{1.337103in}{1.182248in}}%
\pgfpathlineto{\pgfqpoint{1.343773in}{1.197694in}}%
\pgfpathlineto{\pgfqpoint{1.350432in}{1.213116in}}%
\pgfpathlineto{\pgfqpoint{1.322319in}{1.218632in}}%
\pgfpathlineto{\pgfqpoint{1.293341in}{1.223435in}}%
\pgfpathlineto{\pgfqpoint{1.288160in}{1.207801in}}%
\pgfpathlineto{\pgfqpoint{1.282971in}{1.192142in}}%
\pgfpathlineto{\pgfqpoint{1.310451in}{1.187537in}}%
\pgfpathlineto{\pgfqpoint{1.337103in}{1.182248in}}%
\pgfpathclose%
\pgfusepath{fill}%
\end{pgfscope}%
\begin{pgfscope}%
\pgfpathrectangle{\pgfqpoint{0.407592in}{0.222904in}}{\pgfqpoint{1.411163in}{1.411163in}}%
\pgfusepath{clip}%
\pgfsetbuttcap%
\pgfsetroundjoin%
\definecolor{currentfill}{rgb}{0.749020,0.772549,0.792157}%
\pgfsetfillcolor{currentfill}%
\pgfsetfillopacity{0.160000}%
\pgfsetlinewidth{0.000000pt}%
\definecolor{currentstroke}{rgb}{0.000000,0.000000,0.000000}%
\pgfsetstrokecolor{currentstroke}%
\pgfsetdash{}{0pt}%
\pgfpathmoveto{\pgfqpoint{0.836238in}{1.273292in}}%
\pgfpathlineto{\pgfqpoint{0.830380in}{1.288577in}}%
\pgfpathlineto{\pgfqpoint{0.824531in}{1.303841in}}%
\pgfpathlineto{\pgfqpoint{0.791691in}{1.296744in}}%
\pgfpathlineto{\pgfqpoint{0.759911in}{1.288888in}}%
\pgfpathlineto{\pgfqpoint{0.767220in}{1.273878in}}%
\pgfpathlineto{\pgfqpoint{0.774538in}{1.258849in}}%
\pgfpathlineto{\pgfqpoint{0.804877in}{1.266436in}}%
\pgfpathlineto{\pgfqpoint{0.836238in}{1.273292in}}%
\pgfpathclose%
\pgfusepath{fill}%
\end{pgfscope}%
\begin{pgfscope}%
\pgfpathrectangle{\pgfqpoint{0.407592in}{0.222904in}}{\pgfqpoint{1.411163in}{1.411163in}}%
\pgfusepath{clip}%
\pgfsetbuttcap%
\pgfsetroundjoin%
\definecolor{currentfill}{rgb}{0.749020,0.772549,0.792157}%
\pgfsetfillcolor{currentfill}%
\pgfsetfillopacity{0.160000}%
\pgfsetlinewidth{0.000000pt}%
\definecolor{currentstroke}{rgb}{0.000000,0.000000,0.000000}%
\pgfsetstrokecolor{currentstroke}%
\pgfsetdash{}{0pt}%
\pgfpathmoveto{\pgfqpoint{0.701760in}{1.408309in}}%
\pgfpathlineto{\pgfqpoint{0.694531in}{1.423155in}}%
\pgfpathlineto{\pgfqpoint{0.687311in}{1.437983in}}%
\pgfpathlineto{\pgfqpoint{0.649702in}{1.428035in}}%
\pgfpathlineto{\pgfqpoint{0.613616in}{1.417250in}}%
\pgfpathlineto{\pgfqpoint{0.622220in}{1.402687in}}%
\pgfpathlineto{\pgfqpoint{0.630832in}{1.388110in}}%
\pgfpathlineto{\pgfqpoint{0.665555in}{1.398616in}}%
\pgfpathlineto{\pgfqpoint{0.701760in}{1.408309in}}%
\pgfpathclose%
\pgfusepath{fill}%
\end{pgfscope}%
\begin{pgfscope}%
\pgfpathrectangle{\pgfqpoint{0.407592in}{0.222904in}}{\pgfqpoint{1.411163in}{1.411163in}}%
\pgfusepath{clip}%
\pgfsetbuttcap%
\pgfsetroundjoin%
\definecolor{currentfill}{rgb}{0.749020,0.772549,0.792157}%
\pgfsetfillcolor{currentfill}%
\pgfsetfillopacity{0.160000}%
\pgfsetlinewidth{0.000000pt}%
\definecolor{currentstroke}{rgb}{0.000000,0.000000,0.000000}%
\pgfsetstrokecolor{currentstroke}%
\pgfsetdash{}{0pt}%
\pgfpathmoveto{\pgfqpoint{0.927699in}{1.191274in}}%
\pgfpathlineto{\pgfqpoint{0.923322in}{1.206914in}}%
\pgfpathlineto{\pgfqpoint{0.918952in}{1.222529in}}%
\pgfpathlineto{\pgfqpoint{0.888992in}{1.217582in}}%
\pgfpathlineto{\pgfqpoint{0.859756in}{1.211925in}}%
\pgfpathlineto{\pgfqpoint{0.865657in}{1.196527in}}%
\pgfpathlineto{\pgfqpoint{0.871567in}{1.181107in}}%
\pgfpathlineto{\pgfqpoint{0.899286in}{1.186530in}}%
\pgfpathlineto{\pgfqpoint{0.927699in}{1.191274in}}%
\pgfpathclose%
\pgfusepath{fill}%
\end{pgfscope}%
\begin{pgfscope}%
\pgfpathrectangle{\pgfqpoint{0.407592in}{0.222904in}}{\pgfqpoint{1.411163in}{1.411163in}}%
\pgfusepath{clip}%
\pgfsetbuttcap%
\pgfsetroundjoin%
\definecolor{currentfill}{rgb}{0.749020,0.772549,0.792157}%
\pgfsetfillcolor{currentfill}%
\pgfsetfillopacity{0.160000}%
\pgfsetlinewidth{0.000000pt}%
\definecolor{currentstroke}{rgb}{0.000000,0.000000,0.000000}%
\pgfsetstrokecolor{currentstroke}%
\pgfsetdash{}{0pt}%
\pgfpathmoveto{\pgfqpoint{1.262126in}{1.129244in}}%
\pgfpathlineto{\pgfqpoint{1.267351in}{1.145008in}}%
\pgfpathlineto{\pgfqpoint{1.272566in}{1.160745in}}%
\pgfpathlineto{\pgfqpoint{1.245895in}{1.164482in}}%
\pgfpathlineto{\pgfqpoint{1.218634in}{1.167541in}}%
\pgfpathlineto{\pgfqpoint{1.214969in}{1.151644in}}%
\pgfpathlineto{\pgfqpoint{1.211298in}{1.135718in}}%
\pgfpathlineto{\pgfqpoint{1.236993in}{1.132803in}}%
\pgfpathlineto{\pgfqpoint{1.262126in}{1.129244in}}%
\pgfpathclose%
\pgfusepath{fill}%
\end{pgfscope}%
\begin{pgfscope}%
\pgfpathrectangle{\pgfqpoint{0.407592in}{0.222904in}}{\pgfqpoint{1.411163in}{1.411163in}}%
\pgfusepath{clip}%
\pgfsetbuttcap%
\pgfsetroundjoin%
\definecolor{currentfill}{rgb}{0.749020,0.772549,0.792157}%
\pgfsetfillcolor{currentfill}%
\pgfsetfillopacity{0.160000}%
\pgfsetlinewidth{0.000000pt}%
\definecolor{currentstroke}{rgb}{0.000000,0.000000,0.000000}%
\pgfsetstrokecolor{currentstroke}%
\pgfsetdash{}{0pt}%
\pgfpathmoveto{\pgfqpoint{1.203936in}{1.103781in}}%
\pgfpathlineto{\pgfqpoint{1.207620in}{1.119764in}}%
\pgfpathlineto{\pgfqpoint{1.211298in}{1.135718in}}%
\pgfpathlineto{\pgfqpoint{1.185140in}{1.137980in}}%
\pgfpathlineto{\pgfqpoint{1.158617in}{1.139581in}}%
\pgfpathlineto{\pgfqpoint{1.156542in}{1.123529in}}%
\pgfpathlineto{\pgfqpoint{1.154463in}{1.107447in}}%
\pgfpathlineto{\pgfqpoint{1.179372in}{1.105927in}}%
\pgfpathlineto{\pgfqpoint{1.203936in}{1.103781in}}%
\pgfpathclose%
\pgfusepath{fill}%
\end{pgfscope}%
\begin{pgfscope}%
\pgfpathrectangle{\pgfqpoint{0.407592in}{0.222904in}}{\pgfqpoint{1.411163in}{1.411163in}}%
\pgfusepath{clip}%
\pgfsetbuttcap%
\pgfsetroundjoin%
\definecolor{currentfill}{rgb}{0.749020,0.772549,0.792157}%
\pgfsetfillcolor{currentfill}%
\pgfsetfillopacity{0.160000}%
\pgfsetlinewidth{0.000000pt}%
\definecolor{currentstroke}{rgb}{0.000000,0.000000,0.000000}%
\pgfsetstrokecolor{currentstroke}%
\pgfsetdash{}{0pt}%
\pgfpathmoveto{\pgfqpoint{1.563549in}{1.361026in}}%
\pgfpathlineto{\pgfqpoint{1.572779in}{1.375661in}}%
\pgfpathlineto{\pgfqpoint{1.582000in}{1.390282in}}%
\pgfpathlineto{\pgfqpoint{1.549502in}{1.400627in}}%
\pgfpathlineto{\pgfqpoint{1.515351in}{1.410154in}}%
\pgfpathlineto{\pgfqpoint{1.507436in}{1.395265in}}%
\pgfpathlineto{\pgfqpoint{1.499510in}{1.380358in}}%
\pgfpathlineto{\pgfqpoint{1.532332in}{1.371089in}}%
\pgfpathlineto{\pgfqpoint{1.563549in}{1.361026in}}%
\pgfpathclose%
\pgfusepath{fill}%
\end{pgfscope}%
\begin{pgfscope}%
\pgfpathrectangle{\pgfqpoint{0.407592in}{0.222904in}}{\pgfqpoint{1.411163in}{1.411163in}}%
\pgfusepath{clip}%
\pgfsetbuttcap%
\pgfsetroundjoin%
\definecolor{currentfill}{rgb}{0.749020,0.772549,0.792157}%
\pgfsetfillcolor{currentfill}%
\pgfsetfillopacity{0.160000}%
\pgfsetlinewidth{0.000000pt}%
\definecolor{currentstroke}{rgb}{0.000000,0.000000,0.000000}%
\pgfsetstrokecolor{currentstroke}%
\pgfsetdash{}{0pt}%
\pgfpathmoveto{\pgfqpoint{1.053260in}{1.107192in}}%
\pgfpathlineto{\pgfqpoint{1.052063in}{1.123268in}}%
\pgfpathlineto{\pgfqpoint{1.050869in}{1.139313in}}%
\pgfpathlineto{\pgfqpoint{1.024019in}{1.137579in}}%
\pgfpathlineto{\pgfqpoint{0.997404in}{1.135186in}}%
\pgfpathlineto{\pgfqpoint{1.000225in}{1.119245in}}%
\pgfpathlineto{\pgfqpoint{1.003050in}{1.103275in}}%
\pgfpathlineto{\pgfqpoint{1.028043in}{1.105547in}}%
\pgfpathlineto{\pgfqpoint{1.053260in}{1.107192in}}%
\pgfpathclose%
\pgfusepath{fill}%
\end{pgfscope}%
\begin{pgfscope}%
\pgfpathrectangle{\pgfqpoint{0.407592in}{0.222904in}}{\pgfqpoint{1.411163in}{1.411163in}}%
\pgfusepath{clip}%
\pgfsetbuttcap%
\pgfsetroundjoin%
\definecolor{currentfill}{rgb}{0.749020,0.772549,0.792157}%
\pgfsetfillcolor{currentfill}%
\pgfsetfillopacity{0.160000}%
\pgfsetlinewidth{0.000000pt}%
\definecolor{currentstroke}{rgb}{0.000000,0.000000,0.000000}%
\pgfsetstrokecolor{currentstroke}%
\pgfsetdash{}{0pt}%
\pgfpathmoveto{\pgfqpoint{0.997404in}{1.135186in}}%
\pgfpathlineto{\pgfqpoint{0.994589in}{1.151098in}}%
\pgfpathlineto{\pgfqpoint{0.991779in}{1.166982in}}%
\pgfpathlineto{\pgfqpoint{0.963899in}{1.163786in}}%
\pgfpathlineto{\pgfqpoint{0.936476in}{1.159915in}}%
\pgfpathlineto{\pgfqpoint{0.940875in}{1.144197in}}%
\pgfpathlineto{\pgfqpoint{0.945281in}{1.128453in}}%
\pgfpathlineto{\pgfqpoint{0.971125in}{1.132140in}}%
\pgfpathlineto{\pgfqpoint{0.997404in}{1.135186in}}%
\pgfpathclose%
\pgfusepath{fill}%
\end{pgfscope}%
\begin{pgfscope}%
\pgfpathrectangle{\pgfqpoint{0.407592in}{0.222904in}}{\pgfqpoint{1.411163in}{1.411163in}}%
\pgfusepath{clip}%
\pgfsetbuttcap%
\pgfsetroundjoin%
\definecolor{currentfill}{rgb}{0.749020,0.772549,0.792157}%
\pgfsetfillcolor{currentfill}%
\pgfsetfillopacity{0.160000}%
\pgfsetlinewidth{0.000000pt}%
\definecolor{currentstroke}{rgb}{0.000000,0.000000,0.000000}%
\pgfsetstrokecolor{currentstroke}%
\pgfsetdash{}{0pt}%
\pgfpathmoveto{\pgfqpoint{1.419709in}{1.230258in}}%
\pgfpathlineto{\pgfqpoint{1.427734in}{1.245353in}}%
\pgfpathlineto{\pgfqpoint{1.435749in}{1.260429in}}%
\pgfpathlineto{\pgfqpoint{1.406946in}{1.267871in}}%
\pgfpathlineto{\pgfqpoint{1.376972in}{1.274577in}}%
\pgfpathlineto{\pgfqpoint{1.370352in}{1.259246in}}%
\pgfpathlineto{\pgfqpoint{1.363722in}{1.243892in}}%
\pgfpathlineto{\pgfqpoint{1.392279in}{1.237429in}}%
\pgfpathlineto{\pgfqpoint{1.419709in}{1.230258in}}%
\pgfpathclose%
\pgfusepath{fill}%
\end{pgfscope}%
\begin{pgfscope}%
\pgfpathrectangle{\pgfqpoint{0.407592in}{0.222904in}}{\pgfqpoint{1.411163in}{1.411163in}}%
\pgfusepath{clip}%
\pgfsetbuttcap%
\pgfsetroundjoin%
\definecolor{currentfill}{rgb}{0.749020,0.772549,0.792157}%
\pgfsetfillcolor{currentfill}%
\pgfsetfillopacity{0.160000}%
\pgfsetlinewidth{0.000000pt}%
\definecolor{currentstroke}{rgb}{0.000000,0.000000,0.000000}%
\pgfsetstrokecolor{currentstroke}%
\pgfsetdash{}{0pt}%
\pgfpathmoveto{\pgfqpoint{0.716244in}{1.378564in}}%
\pgfpathlineto{\pgfqpoint{0.708998in}{1.393446in}}%
\pgfpathlineto{\pgfqpoint{0.701760in}{1.408309in}}%
\pgfpathlineto{\pgfqpoint{0.665555in}{1.398616in}}%
\pgfpathlineto{\pgfqpoint{0.630832in}{1.388110in}}%
\pgfpathlineto{\pgfqpoint{0.639452in}{1.373519in}}%
\pgfpathlineto{\pgfqpoint{0.648081in}{1.358914in}}%
\pgfpathlineto{\pgfqpoint{0.681443in}{1.369133in}}%
\pgfpathlineto{\pgfqpoint{0.716244in}{1.378564in}}%
\pgfpathclose%
\pgfusepath{fill}%
\end{pgfscope}%
\begin{pgfscope}%
\pgfpathrectangle{\pgfqpoint{0.407592in}{0.222904in}}{\pgfqpoint{1.411163in}{1.411163in}}%
\pgfusepath{clip}%
\pgfsetbuttcap%
\pgfsetroundjoin%
\definecolor{currentfill}{rgb}{0.749020,0.772549,0.792157}%
\pgfsetfillcolor{currentfill}%
\pgfsetfillopacity{0.160000}%
\pgfsetlinewidth{0.000000pt}%
\definecolor{currentstroke}{rgb}{0.000000,0.000000,0.000000}%
\pgfsetstrokecolor{currentstroke}%
\pgfsetdash{}{0pt}%
\pgfpathmoveto{\pgfqpoint{1.323735in}{1.151289in}}%
\pgfpathlineto{\pgfqpoint{1.330424in}{1.166780in}}%
\pgfpathlineto{\pgfqpoint{1.337103in}{1.182248in}}%
\pgfpathlineto{\pgfqpoint{1.310451in}{1.187537in}}%
\pgfpathlineto{\pgfqpoint{1.282971in}{1.192142in}}%
\pgfpathlineto{\pgfqpoint{1.277773in}{1.176457in}}%
\pgfpathlineto{\pgfqpoint{1.272566in}{1.160745in}}%
\pgfpathlineto{\pgfqpoint{1.298546in}{1.156343in}}%
\pgfpathlineto{\pgfqpoint{1.323735in}{1.151289in}}%
\pgfpathclose%
\pgfusepath{fill}%
\end{pgfscope}%
\begin{pgfscope}%
\pgfpathrectangle{\pgfqpoint{0.407592in}{0.222904in}}{\pgfqpoint{1.411163in}{1.411163in}}%
\pgfusepath{clip}%
\pgfsetbuttcap%
\pgfsetroundjoin%
\definecolor{currentfill}{rgb}{0.749020,0.772549,0.792157}%
\pgfsetfillcolor{currentfill}%
\pgfsetfillopacity{0.160000}%
\pgfsetlinewidth{0.000000pt}%
\definecolor{currentstroke}{rgb}{0.000000,0.000000,0.000000}%
\pgfsetstrokecolor{currentstroke}%
\pgfsetdash{}{0pt}%
\pgfpathmoveto{\pgfqpoint{0.847980in}{1.242653in}}%
\pgfpathlineto{\pgfqpoint{0.842105in}{1.257983in}}%
\pgfpathlineto{\pgfqpoint{0.836238in}{1.273292in}}%
\pgfpathlineto{\pgfqpoint{0.804877in}{1.266436in}}%
\pgfpathlineto{\pgfqpoint{0.774538in}{1.258849in}}%
\pgfpathlineto{\pgfqpoint{0.781865in}{1.243802in}}%
\pgfpathlineto{\pgfqpoint{0.789201in}{1.228737in}}%
\pgfpathlineto{\pgfqpoint{0.818098in}{1.236046in}}%
\pgfpathlineto{\pgfqpoint{0.847980in}{1.242653in}}%
\pgfpathclose%
\pgfusepath{fill}%
\end{pgfscope}%
\begin{pgfscope}%
\pgfpathrectangle{\pgfqpoint{0.407592in}{0.222904in}}{\pgfqpoint{1.411163in}{1.411163in}}%
\pgfusepath{clip}%
\pgfsetbuttcap%
\pgfsetroundjoin%
\definecolor{currentfill}{rgb}{0.749020,0.772549,0.792157}%
\pgfsetfillcolor{currentfill}%
\pgfsetfillopacity{0.160000}%
\pgfsetlinewidth{0.000000pt}%
\definecolor{currentstroke}{rgb}{0.000000,0.000000,0.000000}%
\pgfsetstrokecolor{currentstroke}%
\pgfsetdash{}{0pt}%
\pgfpathmoveto{\pgfqpoint{1.150294in}{1.075192in}}%
\pgfpathlineto{\pgfqpoint{1.152381in}{1.091334in}}%
\pgfpathlineto{\pgfqpoint{1.154463in}{1.107447in}}%
\pgfpathlineto{\pgfqpoint{1.129303in}{1.108335in}}%
\pgfpathlineto{\pgfqpoint{1.103985in}{1.108589in}}%
\pgfpathlineto{\pgfqpoint{1.103539in}{1.092445in}}%
\pgfpathlineto{\pgfqpoint{1.103093in}{1.076271in}}%
\pgfpathlineto{\pgfqpoint{1.126767in}{1.076031in}}%
\pgfpathlineto{\pgfqpoint{1.150294in}{1.075192in}}%
\pgfpathclose%
\pgfusepath{fill}%
\end{pgfscope}%
\begin{pgfscope}%
\pgfpathrectangle{\pgfqpoint{0.407592in}{0.222904in}}{\pgfqpoint{1.411163in}{1.411163in}}%
\pgfusepath{clip}%
\pgfsetbuttcap%
\pgfsetroundjoin%
\definecolor{currentfill}{rgb}{0.749020,0.772549,0.792157}%
\pgfsetfillcolor{currentfill}%
\pgfsetfillopacity{0.160000}%
\pgfsetlinewidth{0.000000pt}%
\definecolor{currentstroke}{rgb}{0.000000,0.000000,0.000000}%
\pgfsetstrokecolor{currentstroke}%
\pgfsetdash{}{0pt}%
\pgfpathmoveto{\pgfqpoint{1.103093in}{1.076271in}}%
\pgfpathlineto{\pgfqpoint{1.103539in}{1.092445in}}%
\pgfpathlineto{\pgfqpoint{1.103985in}{1.108589in}}%
\pgfpathlineto{\pgfqpoint{1.078606in}{1.108208in}}%
\pgfpathlineto{\pgfqpoint{1.053260in}{1.107192in}}%
\pgfpathlineto{\pgfqpoint{1.054459in}{1.091087in}}%
\pgfpathlineto{\pgfqpoint{1.055660in}{1.074951in}}%
\pgfpathlineto{\pgfqpoint{1.079360in}{1.075910in}}%
\pgfpathlineto{\pgfqpoint{1.103093in}{1.076271in}}%
\pgfpathclose%
\pgfusepath{fill}%
\end{pgfscope}%
\begin{pgfscope}%
\pgfpathrectangle{\pgfqpoint{0.407592in}{0.222904in}}{\pgfqpoint{1.411163in}{1.411163in}}%
\pgfusepath{clip}%
\pgfsetbuttcap%
\pgfsetroundjoin%
\definecolor{currentfill}{rgb}{0.749020,0.772549,0.792157}%
\pgfsetfillcolor{currentfill}%
\pgfsetfillopacity{0.160000}%
\pgfsetlinewidth{0.000000pt}%
\definecolor{currentstroke}{rgb}{0.000000,0.000000,0.000000}%
\pgfsetstrokecolor{currentstroke}%
\pgfsetdash{}{0pt}%
\pgfpathmoveto{\pgfqpoint{1.545062in}{1.331713in}}%
\pgfpathlineto{\pgfqpoint{1.554310in}{1.346377in}}%
\pgfpathlineto{\pgfqpoint{1.563549in}{1.361026in}}%
\pgfpathlineto{\pgfqpoint{1.532332in}{1.371089in}}%
\pgfpathlineto{\pgfqpoint{1.499510in}{1.380358in}}%
\pgfpathlineto{\pgfqpoint{1.491575in}{1.365433in}}%
\pgfpathlineto{\pgfqpoint{1.483630in}{1.350488in}}%
\pgfpathlineto{\pgfqpoint{1.515124in}{1.341485in}}%
\pgfpathlineto{\pgfqpoint{1.545062in}{1.331713in}}%
\pgfpathclose%
\pgfusepath{fill}%
\end{pgfscope}%
\begin{pgfscope}%
\pgfpathrectangle{\pgfqpoint{0.407592in}{0.222904in}}{\pgfqpoint{1.411163in}{1.411163in}}%
\pgfusepath{clip}%
\pgfsetbuttcap%
\pgfsetroundjoin%
\definecolor{currentfill}{rgb}{0.749020,0.772549,0.792157}%
\pgfsetfillcolor{currentfill}%
\pgfsetfillopacity{0.160000}%
\pgfsetlinewidth{0.000000pt}%
\definecolor{currentstroke}{rgb}{0.000000,0.000000,0.000000}%
\pgfsetstrokecolor{currentstroke}%
\pgfsetdash{}{0pt}%
\pgfpathmoveto{\pgfqpoint{0.936476in}{1.159915in}}%
\pgfpathlineto{\pgfqpoint{0.932084in}{1.175607in}}%
\pgfpathlineto{\pgfqpoint{0.927699in}{1.191274in}}%
\pgfpathlineto{\pgfqpoint{0.899286in}{1.186530in}}%
\pgfpathlineto{\pgfqpoint{0.871567in}{1.181107in}}%
\pgfpathlineto{\pgfqpoint{0.877485in}{1.165664in}}%
\pgfpathlineto{\pgfqpoint{0.883412in}{1.150198in}}%
\pgfpathlineto{\pgfqpoint{0.909612in}{1.155381in}}%
\pgfpathlineto{\pgfqpoint{0.936476in}{1.159915in}}%
\pgfpathclose%
\pgfusepath{fill}%
\end{pgfscope}%
\begin{pgfscope}%
\pgfpathrectangle{\pgfqpoint{0.407592in}{0.222904in}}{\pgfqpoint{1.411163in}{1.411163in}}%
\pgfusepath{clip}%
\pgfsetbuttcap%
\pgfsetroundjoin%
\definecolor{currentfill}{rgb}{0.749020,0.772549,0.792157}%
\pgfsetfillcolor{currentfill}%
\pgfsetfillopacity{0.160000}%
\pgfsetlinewidth{0.000000pt}%
\definecolor{currentstroke}{rgb}{0.000000,0.000000,0.000000}%
\pgfsetstrokecolor{currentstroke}%
\pgfsetdash{}{0pt}%
\pgfpathmoveto{\pgfqpoint{1.251652in}{1.097637in}}%
\pgfpathlineto{\pgfqpoint{1.256894in}{1.113454in}}%
\pgfpathlineto{\pgfqpoint{1.262126in}{1.129244in}}%
\pgfpathlineto{\pgfqpoint{1.236993in}{1.132803in}}%
\pgfpathlineto{\pgfqpoint{1.211298in}{1.135718in}}%
\pgfpathlineto{\pgfqpoint{1.207620in}{1.119764in}}%
\pgfpathlineto{\pgfqpoint{1.203936in}{1.103781in}}%
\pgfpathlineto{\pgfqpoint{1.228060in}{1.101014in}}%
\pgfpathlineto{\pgfqpoint{1.251652in}{1.097637in}}%
\pgfpathclose%
\pgfusepath{fill}%
\end{pgfscope}%
\begin{pgfscope}%
\pgfpathrectangle{\pgfqpoint{0.407592in}{0.222904in}}{\pgfqpoint{1.411163in}{1.411163in}}%
\pgfusepath{clip}%
\pgfsetbuttcap%
\pgfsetroundjoin%
\definecolor{currentfill}{rgb}{0.749020,0.772549,0.792157}%
\pgfsetfillcolor{currentfill}%
\pgfsetfillopacity{0.160000}%
\pgfsetlinewidth{0.000000pt}%
\definecolor{currentstroke}{rgb}{0.000000,0.000000,0.000000}%
\pgfsetstrokecolor{currentstroke}%
\pgfsetdash{}{0pt}%
\pgfpathmoveto{\pgfqpoint{1.403629in}{1.200012in}}%
\pgfpathlineto{\pgfqpoint{1.411674in}{1.215145in}}%
\pgfpathlineto{\pgfqpoint{1.419709in}{1.230258in}}%
\pgfpathlineto{\pgfqpoint{1.392279in}{1.237429in}}%
\pgfpathlineto{\pgfqpoint{1.363722in}{1.243892in}}%
\pgfpathlineto{\pgfqpoint{1.357082in}{1.228515in}}%
\pgfpathlineto{\pgfqpoint{1.350432in}{1.213116in}}%
\pgfpathlineto{\pgfqpoint{1.377571in}{1.206903in}}%
\pgfpathlineto{\pgfqpoint{1.403629in}{1.200012in}}%
\pgfpathclose%
\pgfusepath{fill}%
\end{pgfscope}%
\begin{pgfscope}%
\pgfpathrectangle{\pgfqpoint{0.407592in}{0.222904in}}{\pgfqpoint{1.411163in}{1.411163in}}%
\pgfusepath{clip}%
\pgfsetbuttcap%
\pgfsetroundjoin%
\definecolor{currentfill}{rgb}{0.749020,0.772549,0.792157}%
\pgfsetfillcolor{currentfill}%
\pgfsetfillopacity{0.160000}%
\pgfsetlinewidth{0.000000pt}%
\definecolor{currentstroke}{rgb}{0.000000,0.000000,0.000000}%
\pgfsetstrokecolor{currentstroke}%
\pgfsetdash{}{0pt}%
\pgfpathmoveto{\pgfqpoint{1.196547in}{1.071728in}}%
\pgfpathlineto{\pgfqpoint{1.200244in}{1.087769in}}%
\pgfpathlineto{\pgfqpoint{1.203936in}{1.103781in}}%
\pgfpathlineto{\pgfqpoint{1.179372in}{1.105927in}}%
\pgfpathlineto{\pgfqpoint{1.154463in}{1.107447in}}%
\pgfpathlineto{\pgfqpoint{1.152381in}{1.091334in}}%
\pgfpathlineto{\pgfqpoint{1.150294in}{1.075192in}}%
\pgfpathlineto{\pgfqpoint{1.173583in}{1.073756in}}%
\pgfpathlineto{\pgfqpoint{1.196547in}{1.071728in}}%
\pgfpathclose%
\pgfusepath{fill}%
\end{pgfscope}%
\begin{pgfscope}%
\pgfpathrectangle{\pgfqpoint{0.407592in}{0.222904in}}{\pgfqpoint{1.411163in}{1.411163in}}%
\pgfusepath{clip}%
\pgfsetbuttcap%
\pgfsetroundjoin%
\definecolor{currentfill}{rgb}{0.749020,0.772549,0.792157}%
\pgfsetfillcolor{currentfill}%
\pgfsetfillopacity{0.160000}%
\pgfsetlinewidth{0.000000pt}%
\definecolor{currentstroke}{rgb}{0.000000,0.000000,0.000000}%
\pgfsetstrokecolor{currentstroke}%
\pgfsetdash{}{0pt}%
\pgfpathmoveto{\pgfqpoint{1.003050in}{1.103275in}}%
\pgfpathlineto{\pgfqpoint{1.000225in}{1.119245in}}%
\pgfpathlineto{\pgfqpoint{0.997404in}{1.135186in}}%
\pgfpathlineto{\pgfqpoint{0.971125in}{1.132140in}}%
\pgfpathlineto{\pgfqpoint{0.945281in}{1.128453in}}%
\pgfpathlineto{\pgfqpoint{0.949694in}{1.112683in}}%
\pgfpathlineto{\pgfqpoint{0.954115in}{1.096887in}}%
\pgfpathlineto{\pgfqpoint{0.978376in}{1.100385in}}%
\pgfpathlineto{\pgfqpoint{1.003050in}{1.103275in}}%
\pgfpathclose%
\pgfusepath{fill}%
\end{pgfscope}%
\begin{pgfscope}%
\pgfpathrectangle{\pgfqpoint{0.407592in}{0.222904in}}{\pgfqpoint{1.411163in}{1.411163in}}%
\pgfusepath{clip}%
\pgfsetbuttcap%
\pgfsetroundjoin%
\definecolor{currentfill}{rgb}{0.749020,0.772549,0.792157}%
\pgfsetfillcolor{currentfill}%
\pgfsetfillopacity{0.160000}%
\pgfsetlinewidth{0.000000pt}%
\definecolor{currentstroke}{rgb}{0.000000,0.000000,0.000000}%
\pgfsetstrokecolor{currentstroke}%
\pgfsetdash{}{0pt}%
\pgfpathmoveto{\pgfqpoint{0.730764in}{1.348745in}}%
\pgfpathlineto{\pgfqpoint{0.723500in}{1.363663in}}%
\pgfpathlineto{\pgfqpoint{0.716244in}{1.378564in}}%
\pgfpathlineto{\pgfqpoint{0.681443in}{1.369133in}}%
\pgfpathlineto{\pgfqpoint{0.648081in}{1.358914in}}%
\pgfpathlineto{\pgfqpoint{0.656718in}{1.344295in}}%
\pgfpathlineto{\pgfqpoint{0.665363in}{1.329663in}}%
\pgfpathlineto{\pgfqpoint{0.697365in}{1.339586in}}%
\pgfpathlineto{\pgfqpoint{0.730764in}{1.348745in}}%
\pgfpathclose%
\pgfusepath{fill}%
\end{pgfscope}%
\begin{pgfscope}%
\pgfpathrectangle{\pgfqpoint{0.407592in}{0.222904in}}{\pgfqpoint{1.411163in}{1.411163in}}%
\pgfusepath{clip}%
\pgfsetbuttcap%
\pgfsetroundjoin%
\definecolor{currentfill}{rgb}{0.749020,0.772549,0.792157}%
\pgfsetfillcolor{currentfill}%
\pgfsetfillopacity{0.160000}%
\pgfsetlinewidth{0.000000pt}%
\definecolor{currentstroke}{rgb}{0.000000,0.000000,0.000000}%
\pgfsetstrokecolor{currentstroke}%
\pgfsetdash{}{0pt}%
\pgfpathmoveto{\pgfqpoint{1.055660in}{1.074951in}}%
\pgfpathlineto{\pgfqpoint{1.054459in}{1.091087in}}%
\pgfpathlineto{\pgfqpoint{1.053260in}{1.107192in}}%
\pgfpathlineto{\pgfqpoint{1.028043in}{1.105547in}}%
\pgfpathlineto{\pgfqpoint{1.003050in}{1.103275in}}%
\pgfpathlineto{\pgfqpoint{1.005881in}{1.087277in}}%
\pgfpathlineto{\pgfqpoint{1.008716in}{1.071250in}}%
\pgfpathlineto{\pgfqpoint{1.032082in}{1.073396in}}%
\pgfpathlineto{\pgfqpoint{1.055660in}{1.074951in}}%
\pgfpathclose%
\pgfusepath{fill}%
\end{pgfscope}%
\begin{pgfscope}%
\pgfpathrectangle{\pgfqpoint{0.407592in}{0.222904in}}{\pgfqpoint{1.411163in}{1.411163in}}%
\pgfusepath{clip}%
\pgfsetbuttcap%
\pgfsetroundjoin%
\definecolor{currentfill}{rgb}{0.749020,0.772549,0.792157}%
\pgfsetfillcolor{currentfill}%
\pgfsetfillopacity{0.160000}%
\pgfsetlinewidth{0.000000pt}%
\definecolor{currentstroke}{rgb}{0.000000,0.000000,0.000000}%
\pgfsetstrokecolor{currentstroke}%
\pgfsetdash{}{0pt}%
\pgfpathmoveto{\pgfqpoint{0.859756in}{1.211925in}}%
\pgfpathlineto{\pgfqpoint{0.853864in}{1.227300in}}%
\pgfpathlineto{\pgfqpoint{0.847980in}{1.242653in}}%
\pgfpathlineto{\pgfqpoint{0.818098in}{1.236046in}}%
\pgfpathlineto{\pgfqpoint{0.789201in}{1.228737in}}%
\pgfpathlineto{\pgfqpoint{0.796546in}{1.213653in}}%
\pgfpathlineto{\pgfqpoint{0.803900in}{1.198550in}}%
\pgfpathlineto{\pgfqpoint{0.831355in}{1.205574in}}%
\pgfpathlineto{\pgfqpoint{0.859756in}{1.211925in}}%
\pgfpathclose%
\pgfusepath{fill}%
\end{pgfscope}%
\begin{pgfscope}%
\pgfpathrectangle{\pgfqpoint{0.407592in}{0.222904in}}{\pgfqpoint{1.411163in}{1.411163in}}%
\pgfusepath{clip}%
\pgfsetbuttcap%
\pgfsetroundjoin%
\definecolor{currentfill}{rgb}{0.749020,0.772549,0.792157}%
\pgfsetfillcolor{currentfill}%
\pgfsetfillopacity{0.160000}%
\pgfsetlinewidth{0.000000pt}%
\definecolor{currentstroke}{rgb}{0.000000,0.000000,0.000000}%
\pgfsetstrokecolor{currentstroke}%
\pgfsetdash{}{0pt}%
\pgfpathmoveto{\pgfqpoint{1.526539in}{1.302343in}}%
\pgfpathlineto{\pgfqpoint{1.535805in}{1.317036in}}%
\pgfpathlineto{\pgfqpoint{1.545062in}{1.331713in}}%
\pgfpathlineto{\pgfqpoint{1.515124in}{1.341485in}}%
\pgfpathlineto{\pgfqpoint{1.483630in}{1.350488in}}%
\pgfpathlineto{\pgfqpoint{1.475675in}{1.335525in}}%
\pgfpathlineto{\pgfqpoint{1.467710in}{1.320544in}}%
\pgfpathlineto{\pgfqpoint{1.497877in}{1.311815in}}%
\pgfpathlineto{\pgfqpoint{1.526539in}{1.302343in}}%
\pgfpathclose%
\pgfusepath{fill}%
\end{pgfscope}%
\begin{pgfscope}%
\pgfpathrectangle{\pgfqpoint{0.407592in}{0.222904in}}{\pgfqpoint{1.411163in}{1.411163in}}%
\pgfusepath{clip}%
\pgfsetbuttcap%
\pgfsetroundjoin%
\definecolor{currentfill}{rgb}{0.749020,0.772549,0.792157}%
\pgfsetfillcolor{currentfill}%
\pgfsetfillopacity{0.160000}%
\pgfsetlinewidth{0.000000pt}%
\definecolor{currentstroke}{rgb}{0.000000,0.000000,0.000000}%
\pgfsetstrokecolor{currentstroke}%
\pgfsetdash{}{0pt}%
\pgfpathmoveto{\pgfqpoint{1.310327in}{1.120238in}}%
\pgfpathlineto{\pgfqpoint{1.317036in}{1.135775in}}%
\pgfpathlineto{\pgfqpoint{1.323735in}{1.151289in}}%
\pgfpathlineto{\pgfqpoint{1.298546in}{1.156343in}}%
\pgfpathlineto{\pgfqpoint{1.272566in}{1.160745in}}%
\pgfpathlineto{\pgfqpoint{1.267351in}{1.145008in}}%
\pgfpathlineto{\pgfqpoint{1.262126in}{1.129244in}}%
\pgfpathlineto{\pgfqpoint{1.286603in}{1.125051in}}%
\pgfpathlineto{\pgfqpoint{1.310327in}{1.120238in}}%
\pgfpathclose%
\pgfusepath{fill}%
\end{pgfscope}%
\begin{pgfscope}%
\pgfpathrectangle{\pgfqpoint{0.407592in}{0.222904in}}{\pgfqpoint{1.411163in}{1.411163in}}%
\pgfusepath{clip}%
\pgfsetbuttcap%
\pgfsetroundjoin%
\definecolor{currentfill}{rgb}{0.749020,0.772549,0.792157}%
\pgfsetfillcolor{currentfill}%
\pgfsetfillopacity{0.160000}%
\pgfsetlinewidth{0.000000pt}%
\definecolor{currentstroke}{rgb}{0.000000,0.000000,0.000000}%
\pgfsetstrokecolor{currentstroke}%
\pgfsetdash{}{0pt}%
\pgfpathmoveto{\pgfqpoint{0.945281in}{1.128453in}}%
\pgfpathlineto{\pgfqpoint{0.940875in}{1.144197in}}%
\pgfpathlineto{\pgfqpoint{0.936476in}{1.159915in}}%
\pgfpathlineto{\pgfqpoint{0.909612in}{1.155381in}}%
\pgfpathlineto{\pgfqpoint{0.883412in}{1.150198in}}%
\pgfpathlineto{\pgfqpoint{0.889347in}{1.134710in}}%
\pgfpathlineto{\pgfqpoint{0.895292in}{1.119199in}}%
\pgfpathlineto{\pgfqpoint{0.919970in}{1.124135in}}%
\pgfpathlineto{\pgfqpoint{0.945281in}{1.128453in}}%
\pgfpathclose%
\pgfusepath{fill}%
\end{pgfscope}%
\begin{pgfscope}%
\pgfpathrectangle{\pgfqpoint{0.407592in}{0.222904in}}{\pgfqpoint{1.411163in}{1.411163in}}%
\pgfusepath{clip}%
\pgfsetbuttcap%
\pgfsetroundjoin%
\definecolor{currentfill}{rgb}{0.749020,0.772549,0.792157}%
\pgfsetfillcolor{currentfill}%
\pgfsetfillopacity{0.160000}%
\pgfsetlinewidth{0.000000pt}%
\definecolor{currentstroke}{rgb}{0.000000,0.000000,0.000000}%
\pgfsetstrokecolor{currentstroke}%
\pgfsetdash{}{0pt}%
\pgfpathmoveto{\pgfqpoint{0.745320in}{1.318853in}}%
\pgfpathlineto{\pgfqpoint{0.738037in}{1.333808in}}%
\pgfpathlineto{\pgfqpoint{0.730764in}{1.348745in}}%
\pgfpathlineto{\pgfqpoint{0.697365in}{1.339586in}}%
\pgfpathlineto{\pgfqpoint{0.665363in}{1.329663in}}%
\pgfpathlineto{\pgfqpoint{0.674016in}{1.315016in}}%
\pgfpathlineto{\pgfqpoint{0.682677in}{1.300356in}}%
\pgfpathlineto{\pgfqpoint{0.713322in}{1.309973in}}%
\pgfpathlineto{\pgfqpoint{0.745320in}{1.318853in}}%
\pgfpathclose%
\pgfusepath{fill}%
\end{pgfscope}%
\begin{pgfscope}%
\pgfpathrectangle{\pgfqpoint{0.407592in}{0.222904in}}{\pgfqpoint{1.411163in}{1.411163in}}%
\pgfusepath{clip}%
\pgfsetbuttcap%
\pgfsetroundjoin%
\definecolor{currentfill}{rgb}{0.749020,0.772549,0.792157}%
\pgfsetfillcolor{currentfill}%
\pgfsetfillopacity{0.160000}%
\pgfsetlinewidth{0.000000pt}%
\definecolor{currentstroke}{rgb}{0.000000,0.000000,0.000000}%
\pgfsetstrokecolor{currentstroke}%
\pgfsetdash{}{0pt}%
\pgfpathmoveto{\pgfqpoint{1.146109in}{1.042815in}}%
\pgfpathlineto{\pgfqpoint{1.148203in}{1.059018in}}%
\pgfpathlineto{\pgfqpoint{1.150294in}{1.075192in}}%
\pgfpathlineto{\pgfqpoint{1.126767in}{1.076031in}}%
\pgfpathlineto{\pgfqpoint{1.103093in}{1.076271in}}%
\pgfpathlineto{\pgfqpoint{1.102645in}{1.060065in}}%
\pgfpathlineto{\pgfqpoint{1.102197in}{1.043829in}}%
\pgfpathlineto{\pgfqpoint{1.124222in}{1.043604in}}%
\pgfpathlineto{\pgfqpoint{1.146109in}{1.042815in}}%
\pgfpathclose%
\pgfusepath{fill}%
\end{pgfscope}%
\begin{pgfscope}%
\pgfpathrectangle{\pgfqpoint{0.407592in}{0.222904in}}{\pgfqpoint{1.411163in}{1.411163in}}%
\pgfusepath{clip}%
\pgfsetbuttcap%
\pgfsetroundjoin%
\definecolor{currentfill}{rgb}{0.749020,0.772549,0.792157}%
\pgfsetfillcolor{currentfill}%
\pgfsetfillopacity{0.160000}%
\pgfsetlinewidth{0.000000pt}%
\definecolor{currentstroke}{rgb}{0.000000,0.000000,0.000000}%
\pgfsetstrokecolor{currentstroke}%
\pgfsetdash{}{0pt}%
\pgfpathmoveto{\pgfqpoint{1.102197in}{1.043829in}}%
\pgfpathlineto{\pgfqpoint{1.102645in}{1.060065in}}%
\pgfpathlineto{\pgfqpoint{1.103093in}{1.076271in}}%
\pgfpathlineto{\pgfqpoint{1.079360in}{1.075910in}}%
\pgfpathlineto{\pgfqpoint{1.055660in}{1.074951in}}%
\pgfpathlineto{\pgfqpoint{1.056863in}{1.058785in}}%
\pgfpathlineto{\pgfqpoint{1.058069in}{1.042589in}}%
\pgfpathlineto{\pgfqpoint{1.080117in}{1.043491in}}%
\pgfpathlineto{\pgfqpoint{1.102197in}{1.043829in}}%
\pgfpathclose%
\pgfusepath{fill}%
\end{pgfscope}%
\begin{pgfscope}%
\pgfpathrectangle{\pgfqpoint{0.407592in}{0.222904in}}{\pgfqpoint{1.411163in}{1.411163in}}%
\pgfusepath{clip}%
\pgfsetbuttcap%
\pgfsetroundjoin%
\definecolor{currentfill}{rgb}{0.749020,0.772549,0.792157}%
\pgfsetfillcolor{currentfill}%
\pgfsetfillopacity{0.160000}%
\pgfsetlinewidth{0.000000pt}%
\definecolor{currentstroke}{rgb}{0.000000,0.000000,0.000000}%
\pgfsetstrokecolor{currentstroke}%
\pgfsetdash{}{0pt}%
\pgfpathmoveto{\pgfqpoint{1.387508in}{1.169690in}}%
\pgfpathlineto{\pgfqpoint{1.395573in}{1.184860in}}%
\pgfpathlineto{\pgfqpoint{1.403629in}{1.200012in}}%
\pgfpathlineto{\pgfqpoint{1.377571in}{1.206903in}}%
\pgfpathlineto{\pgfqpoint{1.350432in}{1.213116in}}%
\pgfpathlineto{\pgfqpoint{1.343773in}{1.197694in}}%
\pgfpathlineto{\pgfqpoint{1.337103in}{1.182248in}}%
\pgfpathlineto{\pgfqpoint{1.362823in}{1.176293in}}%
\pgfpathlineto{\pgfqpoint{1.387508in}{1.169690in}}%
\pgfpathclose%
\pgfusepath{fill}%
\end{pgfscope}%
\begin{pgfscope}%
\pgfpathrectangle{\pgfqpoint{0.407592in}{0.222904in}}{\pgfqpoint{1.411163in}{1.411163in}}%
\pgfusepath{clip}%
\pgfsetbuttcap%
\pgfsetroundjoin%
\definecolor{currentfill}{rgb}{0.749020,0.772549,0.792157}%
\pgfsetfillcolor{currentfill}%
\pgfsetfillopacity{0.160000}%
\pgfsetlinewidth{0.000000pt}%
\definecolor{currentstroke}{rgb}{0.000000,0.000000,0.000000}%
\pgfsetstrokecolor{currentstroke}%
\pgfsetdash{}{0pt}%
\pgfpathmoveto{\pgfqpoint{1.241143in}{1.065924in}}%
\pgfpathlineto{\pgfqpoint{1.246402in}{1.081794in}}%
\pgfpathlineto{\pgfqpoint{1.251652in}{1.097637in}}%
\pgfpathlineto{\pgfqpoint{1.228060in}{1.101014in}}%
\pgfpathlineto{\pgfqpoint{1.203936in}{1.103781in}}%
\pgfpathlineto{\pgfqpoint{1.200244in}{1.087769in}}%
\pgfpathlineto{\pgfqpoint{1.196547in}{1.071728in}}%
\pgfpathlineto{\pgfqpoint{1.219096in}{1.069114in}}%
\pgfpathlineto{\pgfqpoint{1.241143in}{1.065924in}}%
\pgfpathclose%
\pgfusepath{fill}%
\end{pgfscope}%
\begin{pgfscope}%
\pgfpathrectangle{\pgfqpoint{0.407592in}{0.222904in}}{\pgfqpoint{1.411163in}{1.411163in}}%
\pgfusepath{clip}%
\pgfsetbuttcap%
\pgfsetroundjoin%
\definecolor{currentfill}{rgb}{0.749020,0.772549,0.792157}%
\pgfsetfillcolor{currentfill}%
\pgfsetfillopacity{0.160000}%
\pgfsetlinewidth{0.000000pt}%
\definecolor{currentstroke}{rgb}{0.000000,0.000000,0.000000}%
\pgfsetstrokecolor{currentstroke}%
\pgfsetdash{}{0pt}%
\pgfpathmoveto{\pgfqpoint{1.507980in}{1.272915in}}%
\pgfpathlineto{\pgfqpoint{1.517264in}{1.287636in}}%
\pgfpathlineto{\pgfqpoint{1.526539in}{1.302343in}}%
\pgfpathlineto{\pgfqpoint{1.497877in}{1.311815in}}%
\pgfpathlineto{\pgfqpoint{1.467710in}{1.320544in}}%
\pgfpathlineto{\pgfqpoint{1.459735in}{1.305543in}}%
\pgfpathlineto{\pgfqpoint{1.451750in}{1.290524in}}%
\pgfpathlineto{\pgfqpoint{1.480591in}{1.282078in}}%
\pgfpathlineto{\pgfqpoint{1.507980in}{1.272915in}}%
\pgfpathclose%
\pgfusepath{fill}%
\end{pgfscope}%
\begin{pgfscope}%
\pgfpathrectangle{\pgfqpoint{0.407592in}{0.222904in}}{\pgfqpoint{1.411163in}{1.411163in}}%
\pgfusepath{clip}%
\pgfsetbuttcap%
\pgfsetroundjoin%
\definecolor{currentfill}{rgb}{0.749020,0.772549,0.792157}%
\pgfsetfillcolor{currentfill}%
\pgfsetfillopacity{0.160000}%
\pgfsetlinewidth{0.000000pt}%
\definecolor{currentstroke}{rgb}{0.000000,0.000000,0.000000}%
\pgfsetstrokecolor{currentstroke}%
\pgfsetdash{}{0pt}%
\pgfpathmoveto{\pgfqpoint{1.008716in}{1.071250in}}%
\pgfpathlineto{\pgfqpoint{1.005881in}{1.087277in}}%
\pgfpathlineto{\pgfqpoint{1.003050in}{1.103275in}}%
\pgfpathlineto{\pgfqpoint{0.978376in}{1.100385in}}%
\pgfpathlineto{\pgfqpoint{0.954115in}{1.096887in}}%
\pgfpathlineto{\pgfqpoint{0.958543in}{1.081064in}}%
\pgfpathlineto{\pgfqpoint{0.962979in}{1.065216in}}%
\pgfpathlineto{\pgfqpoint{0.985652in}{1.068520in}}%
\pgfpathlineto{\pgfqpoint{1.008716in}{1.071250in}}%
\pgfpathclose%
\pgfusepath{fill}%
\end{pgfscope}%
\begin{pgfscope}%
\pgfpathrectangle{\pgfqpoint{0.407592in}{0.222904in}}{\pgfqpoint{1.411163in}{1.411163in}}%
\pgfusepath{clip}%
\pgfsetbuttcap%
\pgfsetroundjoin%
\definecolor{currentfill}{rgb}{0.749020,0.772549,0.792157}%
\pgfsetfillcolor{currentfill}%
\pgfsetfillopacity{0.160000}%
\pgfsetlinewidth{0.000000pt}%
\definecolor{currentstroke}{rgb}{0.000000,0.000000,0.000000}%
\pgfsetstrokecolor{currentstroke}%
\pgfsetdash{}{0pt}%
\pgfpathmoveto{\pgfqpoint{0.871567in}{1.181107in}}%
\pgfpathlineto{\pgfqpoint{0.865657in}{1.196527in}}%
\pgfpathlineto{\pgfqpoint{0.859756in}{1.211925in}}%
\pgfpathlineto{\pgfqpoint{0.831355in}{1.205574in}}%
\pgfpathlineto{\pgfqpoint{0.803900in}{1.198550in}}%
\pgfpathlineto{\pgfqpoint{0.811263in}{1.183428in}}%
\pgfpathlineto{\pgfqpoint{0.818636in}{1.168289in}}%
\pgfpathlineto{\pgfqpoint{0.844648in}{1.175020in}}%
\pgfpathlineto{\pgfqpoint{0.871567in}{1.181107in}}%
\pgfpathclose%
\pgfusepath{fill}%
\end{pgfscope}%
\begin{pgfscope}%
\pgfpathrectangle{\pgfqpoint{0.407592in}{0.222904in}}{\pgfqpoint{1.411163in}{1.411163in}}%
\pgfusepath{clip}%
\pgfsetbuttcap%
\pgfsetroundjoin%
\definecolor{currentfill}{rgb}{0.749020,0.772549,0.792157}%
\pgfsetfillcolor{currentfill}%
\pgfsetfillopacity{0.160000}%
\pgfsetlinewidth{0.000000pt}%
\definecolor{currentstroke}{rgb}{0.000000,0.000000,0.000000}%
\pgfsetstrokecolor{currentstroke}%
\pgfsetdash{}{0pt}%
\pgfpathmoveto{\pgfqpoint{1.189131in}{1.039559in}}%
\pgfpathlineto{\pgfqpoint{1.192842in}{1.055658in}}%
\pgfpathlineto{\pgfqpoint{1.196547in}{1.071728in}}%
\pgfpathlineto{\pgfqpoint{1.173583in}{1.073756in}}%
\pgfpathlineto{\pgfqpoint{1.150294in}{1.075192in}}%
\pgfpathlineto{\pgfqpoint{1.148203in}{1.059018in}}%
\pgfpathlineto{\pgfqpoint{1.146109in}{1.042815in}}%
\pgfpathlineto{\pgfqpoint{1.167773in}{1.041465in}}%
\pgfpathlineto{\pgfqpoint{1.189131in}{1.039559in}}%
\pgfpathclose%
\pgfusepath{fill}%
\end{pgfscope}%
\begin{pgfscope}%
\pgfpathrectangle{\pgfqpoint{0.407592in}{0.222904in}}{\pgfqpoint{1.411163in}{1.411163in}}%
\pgfusepath{clip}%
\pgfsetbuttcap%
\pgfsetroundjoin%
\definecolor{currentfill}{rgb}{0.749020,0.772549,0.792157}%
\pgfsetfillcolor{currentfill}%
\pgfsetfillopacity{0.160000}%
\pgfsetlinewidth{0.000000pt}%
\definecolor{currentstroke}{rgb}{0.000000,0.000000,0.000000}%
\pgfsetstrokecolor{currentstroke}%
\pgfsetdash{}{0pt}%
\pgfpathmoveto{\pgfqpoint{1.058069in}{1.042589in}}%
\pgfpathlineto{\pgfqpoint{1.056863in}{1.058785in}}%
\pgfpathlineto{\pgfqpoint{1.055660in}{1.074951in}}%
\pgfpathlineto{\pgfqpoint{1.032082in}{1.073396in}}%
\pgfpathlineto{\pgfqpoint{1.008716in}{1.071250in}}%
\pgfpathlineto{\pgfqpoint{1.011557in}{1.055195in}}%
\pgfpathlineto{\pgfqpoint{1.014403in}{1.039110in}}%
\pgfpathlineto{\pgfqpoint{1.036136in}{1.041127in}}%
\pgfpathlineto{\pgfqpoint{1.058069in}{1.042589in}}%
\pgfpathclose%
\pgfusepath{fill}%
\end{pgfscope}%
\begin{pgfscope}%
\pgfpathrectangle{\pgfqpoint{0.407592in}{0.222904in}}{\pgfqpoint{1.411163in}{1.411163in}}%
\pgfusepath{clip}%
\pgfsetbuttcap%
\pgfsetroundjoin%
\definecolor{currentfill}{rgb}{0.749020,0.772549,0.792157}%
\pgfsetfillcolor{currentfill}%
\pgfsetfillopacity{0.160000}%
\pgfsetlinewidth{0.000000pt}%
\definecolor{currentstroke}{rgb}{0.000000,0.000000,0.000000}%
\pgfsetstrokecolor{currentstroke}%
\pgfsetdash{}{0pt}%
\pgfpathmoveto{\pgfqpoint{0.759911in}{1.288888in}}%
\pgfpathlineto{\pgfqpoint{0.752611in}{1.303880in}}%
\pgfpathlineto{\pgfqpoint{0.745320in}{1.318853in}}%
\pgfpathlineto{\pgfqpoint{0.713322in}{1.309973in}}%
\pgfpathlineto{\pgfqpoint{0.682677in}{1.300356in}}%
\pgfpathlineto{\pgfqpoint{0.691347in}{1.285681in}}%
\pgfpathlineto{\pgfqpoint{0.700025in}{1.270993in}}%
\pgfpathlineto{\pgfqpoint{0.729314in}{1.280296in}}%
\pgfpathlineto{\pgfqpoint{0.759911in}{1.288888in}}%
\pgfpathclose%
\pgfusepath{fill}%
\end{pgfscope}%
\begin{pgfscope}%
\pgfpathrectangle{\pgfqpoint{0.407592in}{0.222904in}}{\pgfqpoint{1.411163in}{1.411163in}}%
\pgfusepath{clip}%
\pgfsetbuttcap%
\pgfsetroundjoin%
\definecolor{currentfill}{rgb}{0.749020,0.772549,0.792157}%
\pgfsetfillcolor{currentfill}%
\pgfsetfillopacity{0.160000}%
\pgfsetlinewidth{0.000000pt}%
\definecolor{currentstroke}{rgb}{0.000000,0.000000,0.000000}%
\pgfsetstrokecolor{currentstroke}%
\pgfsetdash{}{0pt}%
\pgfpathmoveto{\pgfqpoint{1.296878in}{1.089094in}}%
\pgfpathlineto{\pgfqpoint{1.303607in}{1.104677in}}%
\pgfpathlineto{\pgfqpoint{1.310327in}{1.120238in}}%
\pgfpathlineto{\pgfqpoint{1.286603in}{1.125051in}}%
\pgfpathlineto{\pgfqpoint{1.262126in}{1.129244in}}%
\pgfpathlineto{\pgfqpoint{1.256894in}{1.113454in}}%
\pgfpathlineto{\pgfqpoint{1.251652in}{1.097637in}}%
\pgfpathlineto{\pgfqpoint{1.274622in}{1.093659in}}%
\pgfpathlineto{\pgfqpoint{1.296878in}{1.089094in}}%
\pgfpathclose%
\pgfusepath{fill}%
\end{pgfscope}%
\begin{pgfscope}%
\pgfpathrectangle{\pgfqpoint{0.407592in}{0.222904in}}{\pgfqpoint{1.411163in}{1.411163in}}%
\pgfusepath{clip}%
\pgfsetbuttcap%
\pgfsetroundjoin%
\definecolor{currentfill}{rgb}{0.749020,0.772549,0.792157}%
\pgfsetfillcolor{currentfill}%
\pgfsetfillopacity{0.160000}%
\pgfsetlinewidth{0.000000pt}%
\definecolor{currentstroke}{rgb}{0.000000,0.000000,0.000000}%
\pgfsetstrokecolor{currentstroke}%
\pgfsetdash{}{0pt}%
\pgfpathmoveto{\pgfqpoint{0.954115in}{1.096887in}}%
\pgfpathlineto{\pgfqpoint{0.949694in}{1.112683in}}%
\pgfpathlineto{\pgfqpoint{0.945281in}{1.128453in}}%
\pgfpathlineto{\pgfqpoint{0.919970in}{1.124135in}}%
\pgfpathlineto{\pgfqpoint{0.895292in}{1.119199in}}%
\pgfpathlineto{\pgfqpoint{0.901245in}{1.103665in}}%
\pgfpathlineto{\pgfqpoint{0.907207in}{1.088109in}}%
\pgfpathlineto{\pgfqpoint{0.930361in}{1.092790in}}%
\pgfpathlineto{\pgfqpoint{0.954115in}{1.096887in}}%
\pgfpathclose%
\pgfusepath{fill}%
\end{pgfscope}%
\begin{pgfscope}%
\pgfpathrectangle{\pgfqpoint{0.407592in}{0.222904in}}{\pgfqpoint{1.411163in}{1.411163in}}%
\pgfusepath{clip}%
\pgfsetbuttcap%
\pgfsetroundjoin%
\definecolor{currentfill}{rgb}{0.749020,0.772549,0.792157}%
\pgfsetfillcolor{currentfill}%
\pgfsetfillopacity{0.160000}%
\pgfsetlinewidth{0.000000pt}%
\definecolor{currentstroke}{rgb}{0.000000,0.000000,0.000000}%
\pgfsetstrokecolor{currentstroke}%
\pgfsetdash{}{0pt}%
\pgfpathmoveto{\pgfqpoint{1.489383in}{1.243429in}}%
\pgfpathlineto{\pgfqpoint{1.498686in}{1.258179in}}%
\pgfpathlineto{\pgfqpoint{1.507980in}{1.272915in}}%
\pgfpathlineto{\pgfqpoint{1.480591in}{1.282078in}}%
\pgfpathlineto{\pgfqpoint{1.451750in}{1.290524in}}%
\pgfpathlineto{\pgfqpoint{1.443755in}{1.275486in}}%
\pgfpathlineto{\pgfqpoint{1.435749in}{1.260429in}}%
\pgfpathlineto{\pgfqpoint{1.463266in}{1.252274in}}%
\pgfpathlineto{\pgfqpoint{1.489383in}{1.243429in}}%
\pgfpathclose%
\pgfusepath{fill}%
\end{pgfscope}%
\begin{pgfscope}%
\pgfpathrectangle{\pgfqpoint{0.407592in}{0.222904in}}{\pgfqpoint{1.411163in}{1.411163in}}%
\pgfusepath{clip}%
\pgfsetbuttcap%
\pgfsetroundjoin%
\definecolor{currentfill}{rgb}{0.749020,0.772549,0.792157}%
\pgfsetfillcolor{currentfill}%
\pgfsetfillopacity{0.160000}%
\pgfsetlinewidth{0.000000pt}%
\definecolor{currentstroke}{rgb}{0.000000,0.000000,0.000000}%
\pgfsetstrokecolor{currentstroke}%
\pgfsetdash{}{0pt}%
\pgfpathmoveto{\pgfqpoint{1.371346in}{1.139291in}}%
\pgfpathlineto{\pgfqpoint{1.379432in}{1.154500in}}%
\pgfpathlineto{\pgfqpoint{1.387508in}{1.169690in}}%
\pgfpathlineto{\pgfqpoint{1.362823in}{1.176293in}}%
\pgfpathlineto{\pgfqpoint{1.337103in}{1.182248in}}%
\pgfpathlineto{\pgfqpoint{1.330424in}{1.166780in}}%
\pgfpathlineto{\pgfqpoint{1.323735in}{1.151289in}}%
\pgfpathlineto{\pgfqpoint{1.348034in}{1.145599in}}%
\pgfpathlineto{\pgfqpoint{1.371346in}{1.139291in}}%
\pgfpathclose%
\pgfusepath{fill}%
\end{pgfscope}%
\begin{pgfscope}%
\pgfpathrectangle{\pgfqpoint{0.407592in}{0.222904in}}{\pgfqpoint{1.411163in}{1.411163in}}%
\pgfusepath{clip}%
\pgfsetbuttcap%
\pgfsetroundjoin%
\definecolor{currentfill}{rgb}{0.749020,0.772549,0.792157}%
\pgfsetfillcolor{currentfill}%
\pgfsetfillopacity{0.160000}%
\pgfsetlinewidth{0.000000pt}%
\definecolor{currentstroke}{rgb}{0.000000,0.000000,0.000000}%
\pgfsetstrokecolor{currentstroke}%
\pgfsetdash{}{0pt}%
\pgfpathmoveto{\pgfqpoint{1.141908in}{1.010316in}}%
\pgfpathlineto{\pgfqpoint{1.144010in}{1.026581in}}%
\pgfpathlineto{\pgfqpoint{1.146109in}{1.042815in}}%
\pgfpathlineto{\pgfqpoint{1.124222in}{1.043604in}}%
\pgfpathlineto{\pgfqpoint{1.102197in}{1.043829in}}%
\pgfpathlineto{\pgfqpoint{1.101747in}{1.027562in}}%
\pgfpathlineto{\pgfqpoint{1.101297in}{1.011264in}}%
\pgfpathlineto{\pgfqpoint{1.121667in}{1.011053in}}%
\pgfpathlineto{\pgfqpoint{1.141908in}{1.010316in}}%
\pgfpathclose%
\pgfusepath{fill}%
\end{pgfscope}%
\begin{pgfscope}%
\pgfpathrectangle{\pgfqpoint{0.407592in}{0.222904in}}{\pgfqpoint{1.411163in}{1.411163in}}%
\pgfusepath{clip}%
\pgfsetbuttcap%
\pgfsetroundjoin%
\definecolor{currentfill}{rgb}{0.749020,0.772549,0.792157}%
\pgfsetfillcolor{currentfill}%
\pgfsetfillopacity{0.160000}%
\pgfsetlinewidth{0.000000pt}%
\definecolor{currentstroke}{rgb}{0.000000,0.000000,0.000000}%
\pgfsetstrokecolor{currentstroke}%
\pgfsetdash{}{0pt}%
\pgfpathmoveto{\pgfqpoint{1.101297in}{1.011264in}}%
\pgfpathlineto{\pgfqpoint{1.101747in}{1.027562in}}%
\pgfpathlineto{\pgfqpoint{1.102197in}{1.043829in}}%
\pgfpathlineto{\pgfqpoint{1.080117in}{1.043491in}}%
\pgfpathlineto{\pgfqpoint{1.058069in}{1.042589in}}%
\pgfpathlineto{\pgfqpoint{1.059277in}{1.026362in}}%
\pgfpathlineto{\pgfqpoint{1.060487in}{1.010105in}}%
\pgfpathlineto{\pgfqpoint{1.080877in}{1.010947in}}%
\pgfpathlineto{\pgfqpoint{1.101297in}{1.011264in}}%
\pgfpathclose%
\pgfusepath{fill}%
\end{pgfscope}%
\begin{pgfscope}%
\pgfpathrectangle{\pgfqpoint{0.407592in}{0.222904in}}{\pgfqpoint{1.411163in}{1.411163in}}%
\pgfusepath{clip}%
\pgfsetbuttcap%
\pgfsetroundjoin%
\definecolor{currentfill}{rgb}{0.749020,0.772549,0.792157}%
\pgfsetfillcolor{currentfill}%
\pgfsetfillopacity{0.160000}%
\pgfsetlinewidth{0.000000pt}%
\definecolor{currentstroke}{rgb}{0.000000,0.000000,0.000000}%
\pgfsetstrokecolor{currentstroke}%
\pgfsetdash{}{0pt}%
\pgfpathmoveto{\pgfqpoint{0.883412in}{1.150198in}}%
\pgfpathlineto{\pgfqpoint{0.877485in}{1.165664in}}%
\pgfpathlineto{\pgfqpoint{0.871567in}{1.181107in}}%
\pgfpathlineto{\pgfqpoint{0.844648in}{1.175020in}}%
\pgfpathlineto{\pgfqpoint{0.818636in}{1.168289in}}%
\pgfpathlineto{\pgfqpoint{0.826017in}{1.153130in}}%
\pgfpathlineto{\pgfqpoint{0.833408in}{1.137953in}}%
\pgfpathlineto{\pgfqpoint{0.857977in}{1.144383in}}%
\pgfpathlineto{\pgfqpoint{0.883412in}{1.150198in}}%
\pgfpathclose%
\pgfusepath{fill}%
\end{pgfscope}%
\begin{pgfscope}%
\pgfpathrectangle{\pgfqpoint{0.407592in}{0.222904in}}{\pgfqpoint{1.411163in}{1.411163in}}%
\pgfusepath{clip}%
\pgfsetbuttcap%
\pgfsetroundjoin%
\definecolor{currentfill}{rgb}{0.749020,0.772549,0.792157}%
\pgfsetfillcolor{currentfill}%
\pgfsetfillopacity{0.160000}%
\pgfsetlinewidth{0.000000pt}%
\definecolor{currentstroke}{rgb}{0.000000,0.000000,0.000000}%
\pgfsetstrokecolor{currentstroke}%
\pgfsetdash{}{0pt}%
\pgfpathmoveto{\pgfqpoint{1.230598in}{1.034105in}}%
\pgfpathlineto{\pgfqpoint{1.235875in}{1.050028in}}%
\pgfpathlineto{\pgfqpoint{1.241143in}{1.065924in}}%
\pgfpathlineto{\pgfqpoint{1.219096in}{1.069114in}}%
\pgfpathlineto{\pgfqpoint{1.196547in}{1.071728in}}%
\pgfpathlineto{\pgfqpoint{1.192842in}{1.055658in}}%
\pgfpathlineto{\pgfqpoint{1.189131in}{1.039559in}}%
\pgfpathlineto{\pgfqpoint{1.210100in}{1.037103in}}%
\pgfpathlineto{\pgfqpoint{1.230598in}{1.034105in}}%
\pgfpathclose%
\pgfusepath{fill}%
\end{pgfscope}%
\begin{pgfscope}%
\pgfpathrectangle{\pgfqpoint{0.407592in}{0.222904in}}{\pgfqpoint{1.411163in}{1.411163in}}%
\pgfusepath{clip}%
\pgfsetbuttcap%
\pgfsetroundjoin%
\definecolor{currentfill}{rgb}{0.749020,0.772549,0.792157}%
\pgfsetfillcolor{currentfill}%
\pgfsetfillopacity{0.160000}%
\pgfsetlinewidth{0.000000pt}%
\definecolor{currentstroke}{rgb}{0.000000,0.000000,0.000000}%
\pgfsetstrokecolor{currentstroke}%
\pgfsetdash{}{0pt}%
\pgfpathmoveto{\pgfqpoint{0.774538in}{1.258849in}}%
\pgfpathlineto{\pgfqpoint{0.767220in}{1.273878in}}%
\pgfpathlineto{\pgfqpoint{0.759911in}{1.288888in}}%
\pgfpathlineto{\pgfqpoint{0.729314in}{1.280296in}}%
\pgfpathlineto{\pgfqpoint{0.700025in}{1.270993in}}%
\pgfpathlineto{\pgfqpoint{0.708711in}{1.256291in}}%
\pgfpathlineto{\pgfqpoint{0.717406in}{1.241574in}}%
\pgfpathlineto{\pgfqpoint{0.745342in}{1.250554in}}%
\pgfpathlineto{\pgfqpoint{0.774538in}{1.258849in}}%
\pgfpathclose%
\pgfusepath{fill}%
\end{pgfscope}%
\begin{pgfscope}%
\pgfpathrectangle{\pgfqpoint{0.407592in}{0.222904in}}{\pgfqpoint{1.411163in}{1.411163in}}%
\pgfusepath{clip}%
\pgfsetbuttcap%
\pgfsetroundjoin%
\definecolor{currentfill}{rgb}{0.749020,0.772549,0.792157}%
\pgfsetfillcolor{currentfill}%
\pgfsetfillopacity{0.160000}%
\pgfsetlinewidth{0.000000pt}%
\definecolor{currentstroke}{rgb}{0.000000,0.000000,0.000000}%
\pgfsetstrokecolor{currentstroke}%
\pgfsetdash{}{0pt}%
\pgfpathmoveto{\pgfqpoint{1.014403in}{1.039110in}}%
\pgfpathlineto{\pgfqpoint{1.011557in}{1.055195in}}%
\pgfpathlineto{\pgfqpoint{1.008716in}{1.071250in}}%
\pgfpathlineto{\pgfqpoint{0.985652in}{1.068520in}}%
\pgfpathlineto{\pgfqpoint{0.962979in}{1.065216in}}%
\pgfpathlineto{\pgfqpoint{0.967422in}{1.049341in}}%
\pgfpathlineto{\pgfqpoint{0.971872in}{1.033439in}}%
\pgfpathlineto{\pgfqpoint{0.992953in}{1.036545in}}%
\pgfpathlineto{\pgfqpoint{1.014403in}{1.039110in}}%
\pgfpathclose%
\pgfusepath{fill}%
\end{pgfscope}%
\begin{pgfscope}%
\pgfpathrectangle{\pgfqpoint{0.407592in}{0.222904in}}{\pgfqpoint{1.411163in}{1.411163in}}%
\pgfusepath{clip}%
\pgfsetbuttcap%
\pgfsetroundjoin%
\definecolor{currentfill}{rgb}{0.749020,0.772549,0.792157}%
\pgfsetfillcolor{currentfill}%
\pgfsetfillopacity{0.160000}%
\pgfsetlinewidth{0.000000pt}%
\definecolor{currentstroke}{rgb}{0.000000,0.000000,0.000000}%
\pgfsetstrokecolor{currentstroke}%
\pgfsetdash{}{0pt}%
\pgfpathmoveto{\pgfqpoint{1.683104in}{1.424367in}}%
\pgfpathlineto{\pgfqpoint{1.693489in}{1.438622in}}%
\pgfpathlineto{\pgfqpoint{1.661761in}{1.451305in}}%
\pgfpathlineto{\pgfqpoint{1.627969in}{1.463170in}}%
\pgfpathlineto{\pgfqpoint{1.618793in}{1.448621in}}%
\pgfpathlineto{\pgfqpoint{1.651967in}{1.436897in}}%
\pgfusepath{fill}%
\end{pgfscope}%
\begin{pgfscope}%
\pgfpathrectangle{\pgfqpoint{0.407592in}{0.222904in}}{\pgfqpoint{1.411163in}{1.411163in}}%
\pgfusepath{clip}%
\pgfsetbuttcap%
\pgfsetroundjoin%
\definecolor{currentfill}{rgb}{0.749020,0.772549,0.792157}%
\pgfsetfillcolor{currentfill}%
\pgfsetfillopacity{0.160000}%
\pgfsetlinewidth{0.000000pt}%
\definecolor{currentstroke}{rgb}{0.000000,0.000000,0.000000}%
\pgfsetstrokecolor{currentstroke}%
\pgfsetdash{}{0pt}%
\pgfpathmoveto{\pgfqpoint{1.181688in}{1.007274in}}%
\pgfpathlineto{\pgfqpoint{1.185413in}{1.023431in}}%
\pgfpathlineto{\pgfqpoint{1.189131in}{1.039559in}}%
\pgfpathlineto{\pgfqpoint{1.167773in}{1.041465in}}%
\pgfpathlineto{\pgfqpoint{1.146109in}{1.042815in}}%
\pgfpathlineto{\pgfqpoint{1.144010in}{1.026581in}}%
\pgfpathlineto{\pgfqpoint{1.141908in}{1.010316in}}%
\pgfpathlineto{\pgfqpoint{1.161941in}{1.009055in}}%
\pgfpathlineto{\pgfqpoint{1.181688in}{1.007274in}}%
\pgfpathclose%
\pgfusepath{fill}%
\end{pgfscope}%
\begin{pgfscope}%
\pgfpathrectangle{\pgfqpoint{0.407592in}{0.222904in}}{\pgfqpoint{1.411163in}{1.411163in}}%
\pgfusepath{clip}%
\pgfsetbuttcap%
\pgfsetroundjoin%
\definecolor{currentfill}{rgb}{0.749020,0.772549,0.792157}%
\pgfsetfillcolor{currentfill}%
\pgfsetfillopacity{0.160000}%
\pgfsetlinewidth{0.000000pt}%
\definecolor{currentstroke}{rgb}{0.000000,0.000000,0.000000}%
\pgfsetstrokecolor{currentstroke}%
\pgfsetdash{}{0pt}%
\pgfpathmoveto{\pgfqpoint{1.060487in}{1.010105in}}%
\pgfpathlineto{\pgfqpoint{1.059277in}{1.026362in}}%
\pgfpathlineto{\pgfqpoint{1.058069in}{1.042589in}}%
\pgfpathlineto{\pgfqpoint{1.036136in}{1.041127in}}%
\pgfpathlineto{\pgfqpoint{1.014403in}{1.039110in}}%
\pgfpathlineto{\pgfqpoint{1.017254in}{1.022997in}}%
\pgfpathlineto{\pgfqpoint{1.020110in}{1.006855in}}%
\pgfpathlineto{\pgfqpoint{1.040205in}{1.008739in}}%
\pgfpathlineto{\pgfqpoint{1.060487in}{1.010105in}}%
\pgfpathclose%
\pgfusepath{fill}%
\end{pgfscope}%
\begin{pgfscope}%
\pgfpathrectangle{\pgfqpoint{0.407592in}{0.222904in}}{\pgfqpoint{1.411163in}{1.411163in}}%
\pgfusepath{clip}%
\pgfsetbuttcap%
\pgfsetroundjoin%
\definecolor{currentfill}{rgb}{0.749020,0.772549,0.792157}%
\pgfsetfillcolor{currentfill}%
\pgfsetfillopacity{0.160000}%
\pgfsetlinewidth{0.000000pt}%
\definecolor{currentstroke}{rgb}{0.000000,0.000000,0.000000}%
\pgfsetstrokecolor{currentstroke}%
\pgfsetdash{}{0pt}%
\pgfpathmoveto{\pgfqpoint{1.470751in}{1.213885in}}%
\pgfpathlineto{\pgfqpoint{1.480072in}{1.228664in}}%
\pgfpathlineto{\pgfqpoint{1.489383in}{1.243429in}}%
\pgfpathlineto{\pgfqpoint{1.463266in}{1.252274in}}%
\pgfpathlineto{\pgfqpoint{1.435749in}{1.260429in}}%
\pgfpathlineto{\pgfqpoint{1.427734in}{1.245353in}}%
\pgfpathlineto{\pgfqpoint{1.419709in}{1.230258in}}%
\pgfpathlineto{\pgfqpoint{1.445902in}{1.222403in}}%
\pgfpathlineto{\pgfqpoint{1.470751in}{1.213885in}}%
\pgfpathclose%
\pgfusepath{fill}%
\end{pgfscope}%
\begin{pgfscope}%
\pgfpathrectangle{\pgfqpoint{0.407592in}{0.222904in}}{\pgfqpoint{1.411163in}{1.411163in}}%
\pgfusepath{clip}%
\pgfsetbuttcap%
\pgfsetroundjoin%
\definecolor{currentfill}{rgb}{0.749020,0.772549,0.792157}%
\pgfsetfillcolor{currentfill}%
\pgfsetfillopacity{0.160000}%
\pgfsetlinewidth{0.000000pt}%
\definecolor{currentstroke}{rgb}{0.000000,0.000000,0.000000}%
\pgfsetstrokecolor{currentstroke}%
\pgfsetdash{}{0pt}%
\pgfpathmoveto{\pgfqpoint{1.283390in}{1.057857in}}%
\pgfpathlineto{\pgfqpoint{1.290139in}{1.073487in}}%
\pgfpathlineto{\pgfqpoint{1.296878in}{1.089094in}}%
\pgfpathlineto{\pgfqpoint{1.274622in}{1.093659in}}%
\pgfpathlineto{\pgfqpoint{1.251652in}{1.097637in}}%
\pgfpathlineto{\pgfqpoint{1.246402in}{1.081794in}}%
\pgfpathlineto{\pgfqpoint{1.241143in}{1.065924in}}%
\pgfpathlineto{\pgfqpoint{1.262602in}{1.062168in}}%
\pgfpathlineto{\pgfqpoint{1.283390in}{1.057857in}}%
\pgfpathclose%
\pgfusepath{fill}%
\end{pgfscope}%
\begin{pgfscope}%
\pgfpathrectangle{\pgfqpoint{0.407592in}{0.222904in}}{\pgfqpoint{1.411163in}{1.411163in}}%
\pgfusepath{clip}%
\pgfsetbuttcap%
\pgfsetroundjoin%
\definecolor{currentfill}{rgb}{0.749020,0.772549,0.792157}%
\pgfsetfillcolor{currentfill}%
\pgfsetfillopacity{0.160000}%
\pgfsetlinewidth{0.000000pt}%
\definecolor{currentstroke}{rgb}{0.000000,0.000000,0.000000}%
\pgfsetstrokecolor{currentstroke}%
\pgfsetdash{}{0pt}%
\pgfpathmoveto{\pgfqpoint{0.962979in}{1.065216in}}%
\pgfpathlineto{\pgfqpoint{0.958543in}{1.081064in}}%
\pgfpathlineto{\pgfqpoint{0.954115in}{1.096887in}}%
\pgfpathlineto{\pgfqpoint{0.930361in}{1.092790in}}%
\pgfpathlineto{\pgfqpoint{0.907207in}{1.088109in}}%
\pgfpathlineto{\pgfqpoint{0.913177in}{1.072529in}}%
\pgfpathlineto{\pgfqpoint{0.919157in}{1.056927in}}%
\pgfpathlineto{\pgfqpoint{0.940784in}{1.061347in}}%
\pgfpathlineto{\pgfqpoint{0.962979in}{1.065216in}}%
\pgfpathclose%
\pgfusepath{fill}%
\end{pgfscope}%
\begin{pgfscope}%
\pgfpathrectangle{\pgfqpoint{0.407592in}{0.222904in}}{\pgfqpoint{1.411163in}{1.411163in}}%
\pgfusepath{clip}%
\pgfsetbuttcap%
\pgfsetroundjoin%
\definecolor{currentfill}{rgb}{0.749020,0.772549,0.792157}%
\pgfsetfillcolor{currentfill}%
\pgfsetfillopacity{0.160000}%
\pgfsetlinewidth{0.000000pt}%
\definecolor{currentstroke}{rgb}{0.000000,0.000000,0.000000}%
\pgfsetstrokecolor{currentstroke}%
\pgfsetdash{}{0pt}%
\pgfpathmoveto{\pgfqpoint{1.355143in}{1.108815in}}%
\pgfpathlineto{\pgfqpoint{1.363250in}{1.124063in}}%
\pgfpathlineto{\pgfqpoint{1.371346in}{1.139291in}}%
\pgfpathlineto{\pgfqpoint{1.348034in}{1.145599in}}%
\pgfpathlineto{\pgfqpoint{1.323735in}{1.151289in}}%
\pgfpathlineto{\pgfqpoint{1.317036in}{1.135775in}}%
\pgfpathlineto{\pgfqpoint{1.310327in}{1.120238in}}%
\pgfpathlineto{\pgfqpoint{1.333204in}{1.114820in}}%
\pgfpathlineto{\pgfqpoint{1.355143in}{1.108815in}}%
\pgfpathclose%
\pgfusepath{fill}%
\end{pgfscope}%
\begin{pgfscope}%
\pgfpathrectangle{\pgfqpoint{0.407592in}{0.222904in}}{\pgfqpoint{1.411163in}{1.411163in}}%
\pgfusepath{clip}%
\pgfsetbuttcap%
\pgfsetroundjoin%
\definecolor{currentfill}{rgb}{0.749020,0.772549,0.792157}%
\pgfsetfillcolor{currentfill}%
\pgfsetfillopacity{0.160000}%
\pgfsetlinewidth{0.000000pt}%
\definecolor{currentstroke}{rgb}{0.000000,0.000000,0.000000}%
\pgfsetstrokecolor{currentstroke}%
\pgfsetdash{}{0pt}%
\pgfpathmoveto{\pgfqpoint{1.662313in}{1.395827in}}%
\pgfpathlineto{\pgfqpoint{1.672712in}{1.410102in}}%
\pgfpathlineto{\pgfqpoint{1.683104in}{1.424367in}}%
\pgfpathlineto{\pgfqpoint{1.651967in}{1.436897in}}%
\pgfpathlineto{\pgfqpoint{1.618793in}{1.448621in}}%
\pgfpathlineto{\pgfqpoint{1.609608in}{1.434057in}}%
\pgfpathlineto{\pgfqpoint{1.600415in}{1.419480in}}%
\pgfpathlineto{\pgfqpoint{1.632354in}{1.408045in}}%
\pgfpathlineto{\pgfqpoint{1.662313in}{1.395827in}}%
\pgfpathclose%
\pgfusepath{fill}%
\end{pgfscope}%
\begin{pgfscope}%
\pgfpathrectangle{\pgfqpoint{0.407592in}{0.222904in}}{\pgfqpoint{1.411163in}{1.411163in}}%
\pgfusepath{clip}%
\pgfsetbuttcap%
\pgfsetroundjoin%
\definecolor{currentfill}{rgb}{0.749020,0.772549,0.792157}%
\pgfsetfillcolor{currentfill}%
\pgfsetfillopacity{0.160000}%
\pgfsetlinewidth{0.000000pt}%
\definecolor{currentstroke}{rgb}{0.000000,0.000000,0.000000}%
\pgfsetstrokecolor{currentstroke}%
\pgfsetdash{}{0pt}%
\pgfpathmoveto{\pgfqpoint{0.596432in}{1.446335in}}%
\pgfpathlineto{\pgfqpoint{0.587852in}{1.460857in}}%
\pgfpathlineto{\pgfqpoint{0.551453in}{1.448826in}}%
\pgfpathlineto{\pgfqpoint{0.516943in}{1.435982in}}%
\pgfpathlineto{\pgfqpoint{0.526822in}{1.421759in}}%
\pgfpathlineto{\pgfqpoint{0.560694in}{1.434447in}}%
\pgfusepath{fill}%
\end{pgfscope}%
\begin{pgfscope}%
\pgfpathrectangle{\pgfqpoint{0.407592in}{0.222904in}}{\pgfqpoint{1.411163in}{1.411163in}}%
\pgfusepath{clip}%
\pgfsetbuttcap%
\pgfsetroundjoin%
\definecolor{currentfill}{rgb}{0.749020,0.772549,0.792157}%
\pgfsetfillcolor{currentfill}%
\pgfsetfillopacity{0.160000}%
\pgfsetlinewidth{0.000000pt}%
\definecolor{currentstroke}{rgb}{0.000000,0.000000,0.000000}%
\pgfsetstrokecolor{currentstroke}%
\pgfsetdash{}{0pt}%
\pgfpathmoveto{\pgfqpoint{0.789201in}{1.228737in}}%
\pgfpathlineto{\pgfqpoint{0.781865in}{1.243802in}}%
\pgfpathlineto{\pgfqpoint{0.774538in}{1.258849in}}%
\pgfpathlineto{\pgfqpoint{0.745342in}{1.250554in}}%
\pgfpathlineto{\pgfqpoint{0.717406in}{1.241574in}}%
\pgfpathlineto{\pgfqpoint{0.726109in}{1.226844in}}%
\pgfpathlineto{\pgfqpoint{0.734820in}{1.212099in}}%
\pgfpathlineto{\pgfqpoint{0.761404in}{1.220746in}}%
\pgfpathlineto{\pgfqpoint{0.789201in}{1.228737in}}%
\pgfpathclose%
\pgfusepath{fill}%
\end{pgfscope}%
\begin{pgfscope}%
\pgfpathrectangle{\pgfqpoint{0.407592in}{0.222904in}}{\pgfqpoint{1.411163in}{1.411163in}}%
\pgfusepath{clip}%
\pgfsetbuttcap%
\pgfsetroundjoin%
\definecolor{currentfill}{rgb}{0.749020,0.772549,0.792157}%
\pgfsetfillcolor{currentfill}%
\pgfsetfillopacity{0.160000}%
\pgfsetlinewidth{0.000000pt}%
\definecolor{currentstroke}{rgb}{0.000000,0.000000,0.000000}%
\pgfsetstrokecolor{currentstroke}%
\pgfsetdash{}{0pt}%
\pgfpathmoveto{\pgfqpoint{0.895292in}{1.119199in}}%
\pgfpathlineto{\pgfqpoint{0.889347in}{1.134710in}}%
\pgfpathlineto{\pgfqpoint{0.883412in}{1.150198in}}%
\pgfpathlineto{\pgfqpoint{0.857977in}{1.144383in}}%
\pgfpathlineto{\pgfqpoint{0.833408in}{1.137953in}}%
\pgfpathlineto{\pgfqpoint{0.840807in}{1.122756in}}%
\pgfpathlineto{\pgfqpoint{0.848216in}{1.107541in}}%
\pgfpathlineto{\pgfqpoint{0.871342in}{1.113662in}}%
\pgfpathlineto{\pgfqpoint{0.895292in}{1.119199in}}%
\pgfpathclose%
\pgfusepath{fill}%
\end{pgfscope}%
\begin{pgfscope}%
\pgfpathrectangle{\pgfqpoint{0.407592in}{0.222904in}}{\pgfqpoint{1.411163in}{1.411163in}}%
\pgfusepath{clip}%
\pgfsetbuttcap%
\pgfsetroundjoin%
\definecolor{currentfill}{rgb}{0.749020,0.772549,0.792157}%
\pgfsetfillcolor{currentfill}%
\pgfsetfillopacity{0.160000}%
\pgfsetlinewidth{0.000000pt}%
\definecolor{currentstroke}{rgb}{0.000000,0.000000,0.000000}%
\pgfsetstrokecolor{currentstroke}%
\pgfsetdash{}{0pt}%
\pgfpathmoveto{\pgfqpoint{1.137691in}{0.977694in}}%
\pgfpathlineto{\pgfqpoint{1.139801in}{0.994020in}}%
\pgfpathlineto{\pgfqpoint{1.141908in}{1.010316in}}%
\pgfpathlineto{\pgfqpoint{1.121667in}{1.011053in}}%
\pgfpathlineto{\pgfqpoint{1.101297in}{1.011264in}}%
\pgfpathlineto{\pgfqpoint{1.100846in}{0.994934in}}%
\pgfpathlineto{\pgfqpoint{1.100394in}{0.978573in}}%
\pgfpathlineto{\pgfqpoint{1.119102in}{0.978378in}}%
\pgfpathlineto{\pgfqpoint{1.137691in}{0.977694in}}%
\pgfpathclose%
\pgfusepath{fill}%
\end{pgfscope}%
\begin{pgfscope}%
\pgfpathrectangle{\pgfqpoint{0.407592in}{0.222904in}}{\pgfqpoint{1.411163in}{1.411163in}}%
\pgfusepath{clip}%
\pgfsetbuttcap%
\pgfsetroundjoin%
\definecolor{currentfill}{rgb}{0.749020,0.772549,0.792157}%
\pgfsetfillcolor{currentfill}%
\pgfsetfillopacity{0.160000}%
\pgfsetlinewidth{0.000000pt}%
\definecolor{currentstroke}{rgb}{0.000000,0.000000,0.000000}%
\pgfsetstrokecolor{currentstroke}%
\pgfsetdash{}{0pt}%
\pgfpathmoveto{\pgfqpoint{1.641494in}{1.367249in}}%
\pgfpathlineto{\pgfqpoint{1.651907in}{1.381543in}}%
\pgfpathlineto{\pgfqpoint{1.662313in}{1.395827in}}%
\pgfpathlineto{\pgfqpoint{1.632354in}{1.408045in}}%
\pgfpathlineto{\pgfqpoint{1.600415in}{1.419480in}}%
\pgfpathlineto{\pgfqpoint{1.591212in}{1.404888in}}%
\pgfpathlineto{\pgfqpoint{1.582000in}{1.390282in}}%
\pgfpathlineto{\pgfqpoint{1.612709in}{1.379145in}}%
\pgfpathlineto{\pgfqpoint{1.641494in}{1.367249in}}%
\pgfpathclose%
\pgfusepath{fill}%
\end{pgfscope}%
\begin{pgfscope}%
\pgfpathrectangle{\pgfqpoint{0.407592in}{0.222904in}}{\pgfqpoint{1.411163in}{1.411163in}}%
\pgfusepath{clip}%
\pgfsetbuttcap%
\pgfsetroundjoin%
\definecolor{currentfill}{rgb}{0.749020,0.772549,0.792157}%
\pgfsetfillcolor{currentfill}%
\pgfsetfillopacity{0.160000}%
\pgfsetlinewidth{0.000000pt}%
\definecolor{currentstroke}{rgb}{0.000000,0.000000,0.000000}%
\pgfsetstrokecolor{currentstroke}%
\pgfsetdash{}{0pt}%
\pgfpathmoveto{\pgfqpoint{1.452081in}{1.184283in}}%
\pgfpathlineto{\pgfqpoint{1.461420in}{1.199091in}}%
\pgfpathlineto{\pgfqpoint{1.470751in}{1.213885in}}%
\pgfpathlineto{\pgfqpoint{1.445902in}{1.222403in}}%
\pgfpathlineto{\pgfqpoint{1.419709in}{1.230258in}}%
\pgfpathlineto{\pgfqpoint{1.411674in}{1.215145in}}%
\pgfpathlineto{\pgfqpoint{1.403629in}{1.200012in}}%
\pgfpathlineto{\pgfqpoint{1.428500in}{1.192464in}}%
\pgfpathlineto{\pgfqpoint{1.452081in}{1.184283in}}%
\pgfpathclose%
\pgfusepath{fill}%
\end{pgfscope}%
\begin{pgfscope}%
\pgfpathrectangle{\pgfqpoint{0.407592in}{0.222904in}}{\pgfqpoint{1.411163in}{1.411163in}}%
\pgfusepath{clip}%
\pgfsetbuttcap%
\pgfsetroundjoin%
\definecolor{currentfill}{rgb}{0.749020,0.772549,0.792157}%
\pgfsetfillcolor{currentfill}%
\pgfsetfillopacity{0.160000}%
\pgfsetlinewidth{0.000000pt}%
\definecolor{currentstroke}{rgb}{0.000000,0.000000,0.000000}%
\pgfsetstrokecolor{currentstroke}%
\pgfsetdash{}{0pt}%
\pgfpathmoveto{\pgfqpoint{1.100394in}{0.978573in}}%
\pgfpathlineto{\pgfqpoint{1.100846in}{0.994934in}}%
\pgfpathlineto{\pgfqpoint{1.101297in}{1.011264in}}%
\pgfpathlineto{\pgfqpoint{1.080877in}{1.010947in}}%
\pgfpathlineto{\pgfqpoint{1.060487in}{1.010105in}}%
\pgfpathlineto{\pgfqpoint{1.061699in}{0.993817in}}%
\pgfpathlineto{\pgfqpoint{1.062914in}{0.977498in}}%
\pgfpathlineto{\pgfqpoint{1.081640in}{0.978280in}}%
\pgfpathlineto{\pgfqpoint{1.100394in}{0.978573in}}%
\pgfpathclose%
\pgfusepath{fill}%
\end{pgfscope}%
\begin{pgfscope}%
\pgfpathrectangle{\pgfqpoint{0.407592in}{0.222904in}}{\pgfqpoint{1.411163in}{1.411163in}}%
\pgfusepath{clip}%
\pgfsetbuttcap%
\pgfsetroundjoin%
\definecolor{currentfill}{rgb}{0.749020,0.772549,0.792157}%
\pgfsetfillcolor{currentfill}%
\pgfsetfillopacity{0.160000}%
\pgfsetlinewidth{0.000000pt}%
\definecolor{currentstroke}{rgb}{0.000000,0.000000,0.000000}%
\pgfsetstrokecolor{currentstroke}%
\pgfsetdash{}{0pt}%
\pgfpathmoveto{\pgfqpoint{0.613616in}{1.417250in}}%
\pgfpathlineto{\pgfqpoint{0.605020in}{1.431799in}}%
\pgfpathlineto{\pgfqpoint{0.596432in}{1.446335in}}%
\pgfpathlineto{\pgfqpoint{0.560694in}{1.434447in}}%
\pgfpathlineto{\pgfqpoint{0.526822in}{1.421759in}}%
\pgfpathlineto{\pgfqpoint{0.536706in}{1.407526in}}%
\pgfpathlineto{\pgfqpoint{0.546597in}{1.393285in}}%
\pgfpathlineto{\pgfqpoint{0.579199in}{1.405656in}}%
\pgfpathlineto{\pgfqpoint{0.613616in}{1.417250in}}%
\pgfpathclose%
\pgfusepath{fill}%
\end{pgfscope}%
\begin{pgfscope}%
\pgfpathrectangle{\pgfqpoint{0.407592in}{0.222904in}}{\pgfqpoint{1.411163in}{1.411163in}}%
\pgfusepath{clip}%
\pgfsetbuttcap%
\pgfsetroundjoin%
\definecolor{currentfill}{rgb}{0.749020,0.772549,0.792157}%
\pgfsetfillcolor{currentfill}%
\pgfsetfillopacity{0.160000}%
\pgfsetlinewidth{0.000000pt}%
\definecolor{currentstroke}{rgb}{0.000000,0.000000,0.000000}%
\pgfsetstrokecolor{currentstroke}%
\pgfsetdash{}{0pt}%
\pgfpathmoveto{\pgfqpoint{1.220018in}{1.002179in}}%
\pgfpathlineto{\pgfqpoint{1.225313in}{1.018156in}}%
\pgfpathlineto{\pgfqpoint{1.230598in}{1.034105in}}%
\pgfpathlineto{\pgfqpoint{1.210100in}{1.037103in}}%
\pgfpathlineto{\pgfqpoint{1.189131in}{1.039559in}}%
\pgfpathlineto{\pgfqpoint{1.185413in}{1.023431in}}%
\pgfpathlineto{\pgfqpoint{1.181688in}{1.007274in}}%
\pgfpathlineto{\pgfqpoint{1.201073in}{1.004980in}}%
\pgfpathlineto{\pgfqpoint{1.220018in}{1.002179in}}%
\pgfpathclose%
\pgfusepath{fill}%
\end{pgfscope}%
\begin{pgfscope}%
\pgfpathrectangle{\pgfqpoint{0.407592in}{0.222904in}}{\pgfqpoint{1.411163in}{1.411163in}}%
\pgfusepath{clip}%
\pgfsetbuttcap%
\pgfsetroundjoin%
\definecolor{currentfill}{rgb}{0.749020,0.772549,0.792157}%
\pgfsetfillcolor{currentfill}%
\pgfsetfillopacity{0.160000}%
\pgfsetlinewidth{0.000000pt}%
\definecolor{currentstroke}{rgb}{0.000000,0.000000,0.000000}%
\pgfsetstrokecolor{currentstroke}%
\pgfsetdash{}{0pt}%
\pgfpathmoveto{\pgfqpoint{1.020110in}{1.006855in}}%
\pgfpathlineto{\pgfqpoint{1.017254in}{1.022997in}}%
\pgfpathlineto{\pgfqpoint{1.014403in}{1.039110in}}%
\pgfpathlineto{\pgfqpoint{0.992953in}{1.036545in}}%
\pgfpathlineto{\pgfqpoint{0.971872in}{1.033439in}}%
\pgfpathlineto{\pgfqpoint{0.976330in}{1.017512in}}%
\pgfpathlineto{\pgfqpoint{0.980795in}{1.001558in}}%
\pgfpathlineto{\pgfqpoint{1.000280in}{1.004458in}}%
\pgfpathlineto{\pgfqpoint{1.020110in}{1.006855in}}%
\pgfpathclose%
\pgfusepath{fill}%
\end{pgfscope}%
\begin{pgfscope}%
\pgfpathrectangle{\pgfqpoint{0.407592in}{0.222904in}}{\pgfqpoint{1.411163in}{1.411163in}}%
\pgfusepath{clip}%
\pgfsetbuttcap%
\pgfsetroundjoin%
\definecolor{currentfill}{rgb}{0.749020,0.772549,0.792157}%
\pgfsetfillcolor{currentfill}%
\pgfsetfillopacity{0.160000}%
\pgfsetlinewidth{0.000000pt}%
\definecolor{currentstroke}{rgb}{0.000000,0.000000,0.000000}%
\pgfsetstrokecolor{currentstroke}%
\pgfsetdash{}{0pt}%
\pgfpathmoveto{\pgfqpoint{1.269861in}{1.026526in}}%
\pgfpathlineto{\pgfqpoint{1.276630in}{1.042203in}}%
\pgfpathlineto{\pgfqpoint{1.283390in}{1.057857in}}%
\pgfpathlineto{\pgfqpoint{1.262602in}{1.062168in}}%
\pgfpathlineto{\pgfqpoint{1.241143in}{1.065924in}}%
\pgfpathlineto{\pgfqpoint{1.235875in}{1.050028in}}%
\pgfpathlineto{\pgfqpoint{1.230598in}{1.034105in}}%
\pgfpathlineto{\pgfqpoint{1.250545in}{1.030576in}}%
\pgfpathlineto{\pgfqpoint{1.269861in}{1.026526in}}%
\pgfpathclose%
\pgfusepath{fill}%
\end{pgfscope}%
\begin{pgfscope}%
\pgfpathrectangle{\pgfqpoint{0.407592in}{0.222904in}}{\pgfqpoint{1.411163in}{1.411163in}}%
\pgfusepath{clip}%
\pgfsetbuttcap%
\pgfsetroundjoin%
\definecolor{currentfill}{rgb}{0.749020,0.772549,0.792157}%
\pgfsetfillcolor{currentfill}%
\pgfsetfillopacity{0.160000}%
\pgfsetlinewidth{0.000000pt}%
\definecolor{currentstroke}{rgb}{0.000000,0.000000,0.000000}%
\pgfsetstrokecolor{currentstroke}%
\pgfsetdash{}{0pt}%
\pgfpathmoveto{\pgfqpoint{1.338900in}{1.078263in}}%
\pgfpathlineto{\pgfqpoint{1.347027in}{1.093549in}}%
\pgfpathlineto{\pgfqpoint{1.355143in}{1.108815in}}%
\pgfpathlineto{\pgfqpoint{1.333204in}{1.114820in}}%
\pgfpathlineto{\pgfqpoint{1.310327in}{1.120238in}}%
\pgfpathlineto{\pgfqpoint{1.303607in}{1.104677in}}%
\pgfpathlineto{\pgfqpoint{1.296878in}{1.089094in}}%
\pgfpathlineto{\pgfqpoint{1.318333in}{1.083956in}}%
\pgfpathlineto{\pgfqpoint{1.338900in}{1.078263in}}%
\pgfpathclose%
\pgfusepath{fill}%
\end{pgfscope}%
\begin{pgfscope}%
\pgfpathrectangle{\pgfqpoint{0.407592in}{0.222904in}}{\pgfqpoint{1.411163in}{1.411163in}}%
\pgfusepath{clip}%
\pgfsetbuttcap%
\pgfsetroundjoin%
\definecolor{currentfill}{rgb}{0.749020,0.772549,0.792157}%
\pgfsetfillcolor{currentfill}%
\pgfsetfillopacity{0.160000}%
\pgfsetlinewidth{0.000000pt}%
\definecolor{currentstroke}{rgb}{0.000000,0.000000,0.000000}%
\pgfsetstrokecolor{currentstroke}%
\pgfsetdash{}{0pt}%
\pgfpathmoveto{\pgfqpoint{0.803900in}{1.198550in}}%
\pgfpathlineto{\pgfqpoint{0.796546in}{1.213653in}}%
\pgfpathlineto{\pgfqpoint{0.789201in}{1.228737in}}%
\pgfpathlineto{\pgfqpoint{0.761404in}{1.220746in}}%
\pgfpathlineto{\pgfqpoint{0.734820in}{1.212099in}}%
\pgfpathlineto{\pgfqpoint{0.743540in}{1.197340in}}%
\pgfpathlineto{\pgfqpoint{0.752268in}{1.182567in}}%
\pgfpathlineto{\pgfqpoint{0.777502in}{1.190873in}}%
\pgfpathlineto{\pgfqpoint{0.803900in}{1.198550in}}%
\pgfpathclose%
\pgfusepath{fill}%
\end{pgfscope}%
\begin{pgfscope}%
\pgfpathrectangle{\pgfqpoint{0.407592in}{0.222904in}}{\pgfqpoint{1.411163in}{1.411163in}}%
\pgfusepath{clip}%
\pgfsetbuttcap%
\pgfsetroundjoin%
\definecolor{currentfill}{rgb}{0.749020,0.772549,0.792157}%
\pgfsetfillcolor{currentfill}%
\pgfsetfillopacity{0.160000}%
\pgfsetlinewidth{0.000000pt}%
\definecolor{currentstroke}{rgb}{0.000000,0.000000,0.000000}%
\pgfsetstrokecolor{currentstroke}%
\pgfsetdash{}{0pt}%
\pgfpathmoveto{\pgfqpoint{1.174219in}{0.974872in}}%
\pgfpathlineto{\pgfqpoint{1.177957in}{0.991087in}}%
\pgfpathlineto{\pgfqpoint{1.181688in}{1.007274in}}%
\pgfpathlineto{\pgfqpoint{1.161941in}{1.009055in}}%
\pgfpathlineto{\pgfqpoint{1.141908in}{1.010316in}}%
\pgfpathlineto{\pgfqpoint{1.139801in}{0.994020in}}%
\pgfpathlineto{\pgfqpoint{1.137691in}{0.977694in}}%
\pgfpathlineto{\pgfqpoint{1.156087in}{0.976524in}}%
\pgfpathlineto{\pgfqpoint{1.174219in}{0.974872in}}%
\pgfpathclose%
\pgfusepath{fill}%
\end{pgfscope}%
\begin{pgfscope}%
\pgfpathrectangle{\pgfqpoint{0.407592in}{0.222904in}}{\pgfqpoint{1.411163in}{1.411163in}}%
\pgfusepath{clip}%
\pgfsetbuttcap%
\pgfsetroundjoin%
\definecolor{currentfill}{rgb}{0.749020,0.772549,0.792157}%
\pgfsetfillcolor{currentfill}%
\pgfsetfillopacity{0.160000}%
\pgfsetlinewidth{0.000000pt}%
\definecolor{currentstroke}{rgb}{0.000000,0.000000,0.000000}%
\pgfsetstrokecolor{currentstroke}%
\pgfsetdash{}{0pt}%
\pgfpathmoveto{\pgfqpoint{1.062914in}{0.977498in}}%
\pgfpathlineto{\pgfqpoint{1.061699in}{0.993817in}}%
\pgfpathlineto{\pgfqpoint{1.060487in}{1.010105in}}%
\pgfpathlineto{\pgfqpoint{1.040205in}{1.008739in}}%
\pgfpathlineto{\pgfqpoint{1.020110in}{1.006855in}}%
\pgfpathlineto{\pgfqpoint{1.022971in}{0.990683in}}%
\pgfpathlineto{\pgfqpoint{1.025837in}{0.974483in}}%
\pgfpathlineto{\pgfqpoint{1.044288in}{0.976231in}}%
\pgfpathlineto{\pgfqpoint{1.062914in}{0.977498in}}%
\pgfpathclose%
\pgfusepath{fill}%
\end{pgfscope}%
\begin{pgfscope}%
\pgfpathrectangle{\pgfqpoint{0.407592in}{0.222904in}}{\pgfqpoint{1.411163in}{1.411163in}}%
\pgfusepath{clip}%
\pgfsetbuttcap%
\pgfsetroundjoin%
\definecolor{currentfill}{rgb}{0.749020,0.772549,0.792157}%
\pgfsetfillcolor{currentfill}%
\pgfsetfillopacity{0.160000}%
\pgfsetlinewidth{0.000000pt}%
\definecolor{currentstroke}{rgb}{0.000000,0.000000,0.000000}%
\pgfsetstrokecolor{currentstroke}%
\pgfsetdash{}{0pt}%
\pgfpathmoveto{\pgfqpoint{0.971872in}{1.033439in}}%
\pgfpathlineto{\pgfqpoint{0.967422in}{1.049341in}}%
\pgfpathlineto{\pgfqpoint{0.962979in}{1.065216in}}%
\pgfpathlineto{\pgfqpoint{0.940784in}{1.061347in}}%
\pgfpathlineto{\pgfqpoint{0.919157in}{1.056927in}}%
\pgfpathlineto{\pgfqpoint{0.925145in}{1.041301in}}%
\pgfpathlineto{\pgfqpoint{0.931142in}{1.025653in}}%
\pgfpathlineto{\pgfqpoint{0.951241in}{1.029805in}}%
\pgfpathlineto{\pgfqpoint{0.971872in}{1.033439in}}%
\pgfpathclose%
\pgfusepath{fill}%
\end{pgfscope}%
\begin{pgfscope}%
\pgfpathrectangle{\pgfqpoint{0.407592in}{0.222904in}}{\pgfqpoint{1.411163in}{1.411163in}}%
\pgfusepath{clip}%
\pgfsetbuttcap%
\pgfsetroundjoin%
\definecolor{currentfill}{rgb}{0.749020,0.772549,0.792157}%
\pgfsetfillcolor{currentfill}%
\pgfsetfillopacity{0.160000}%
\pgfsetlinewidth{0.000000pt}%
\definecolor{currentstroke}{rgb}{0.000000,0.000000,0.000000}%
\pgfsetstrokecolor{currentstroke}%
\pgfsetdash{}{0pt}%
\pgfpathmoveto{\pgfqpoint{1.620647in}{1.338633in}}%
\pgfpathlineto{\pgfqpoint{1.631074in}{1.352946in}}%
\pgfpathlineto{\pgfqpoint{1.641494in}{1.367249in}}%
\pgfpathlineto{\pgfqpoint{1.612709in}{1.379145in}}%
\pgfpathlineto{\pgfqpoint{1.582000in}{1.390282in}}%
\pgfpathlineto{\pgfqpoint{1.572779in}{1.375661in}}%
\pgfpathlineto{\pgfqpoint{1.563549in}{1.361026in}}%
\pgfpathlineto{\pgfqpoint{1.593031in}{1.350197in}}%
\pgfpathlineto{\pgfqpoint{1.620647in}{1.338633in}}%
\pgfpathclose%
\pgfusepath{fill}%
\end{pgfscope}%
\begin{pgfscope}%
\pgfpathrectangle{\pgfqpoint{0.407592in}{0.222904in}}{\pgfqpoint{1.411163in}{1.411163in}}%
\pgfusepath{clip}%
\pgfsetbuttcap%
\pgfsetroundjoin%
\definecolor{currentfill}{rgb}{0.749020,0.772549,0.792157}%
\pgfsetfillcolor{currentfill}%
\pgfsetfillopacity{0.160000}%
\pgfsetlinewidth{0.000000pt}%
\definecolor{currentstroke}{rgb}{0.000000,0.000000,0.000000}%
\pgfsetstrokecolor{currentstroke}%
\pgfsetdash{}{0pt}%
\pgfpathmoveto{\pgfqpoint{0.630832in}{1.388110in}}%
\pgfpathlineto{\pgfqpoint{0.622220in}{1.402687in}}%
\pgfpathlineto{\pgfqpoint{0.613616in}{1.417250in}}%
\pgfpathlineto{\pgfqpoint{0.579199in}{1.405656in}}%
\pgfpathlineto{\pgfqpoint{0.546597in}{1.393285in}}%
\pgfpathlineto{\pgfqpoint{0.556495in}{1.379034in}}%
\pgfpathlineto{\pgfqpoint{0.566398in}{1.364774in}}%
\pgfpathlineto{\pgfqpoint{0.597732in}{1.376818in}}%
\pgfpathlineto{\pgfqpoint{0.630832in}{1.388110in}}%
\pgfpathclose%
\pgfusepath{fill}%
\end{pgfscope}%
\begin{pgfscope}%
\pgfpathrectangle{\pgfqpoint{0.407592in}{0.222904in}}{\pgfqpoint{1.411163in}{1.411163in}}%
\pgfusepath{clip}%
\pgfsetbuttcap%
\pgfsetroundjoin%
\definecolor{currentfill}{rgb}{0.749020,0.772549,0.792157}%
\pgfsetfillcolor{currentfill}%
\pgfsetfillopacity{0.160000}%
\pgfsetlinewidth{0.000000pt}%
\definecolor{currentstroke}{rgb}{0.000000,0.000000,0.000000}%
\pgfsetstrokecolor{currentstroke}%
\pgfsetdash{}{0pt}%
\pgfpathmoveto{\pgfqpoint{0.907207in}{1.088109in}}%
\pgfpathlineto{\pgfqpoint{0.901245in}{1.103665in}}%
\pgfpathlineto{\pgfqpoint{0.895292in}{1.119199in}}%
\pgfpathlineto{\pgfqpoint{0.871342in}{1.113662in}}%
\pgfpathlineto{\pgfqpoint{0.848216in}{1.107541in}}%
\pgfpathlineto{\pgfqpoint{0.855634in}{1.092308in}}%
\pgfpathlineto{\pgfqpoint{0.863061in}{1.077055in}}%
\pgfpathlineto{\pgfqpoint{0.884743in}{1.082858in}}%
\pgfpathlineto{\pgfqpoint{0.907207in}{1.088109in}}%
\pgfpathclose%
\pgfusepath{fill}%
\end{pgfscope}%
\begin{pgfscope}%
\pgfpathrectangle{\pgfqpoint{0.407592in}{0.222904in}}{\pgfqpoint{1.411163in}{1.411163in}}%
\pgfusepath{clip}%
\pgfsetbuttcap%
\pgfsetroundjoin%
\definecolor{currentfill}{rgb}{0.749020,0.772549,0.792157}%
\pgfsetfillcolor{currentfill}%
\pgfsetfillopacity{0.160000}%
\pgfsetlinewidth{0.000000pt}%
\definecolor{currentstroke}{rgb}{0.000000,0.000000,0.000000}%
\pgfsetstrokecolor{currentstroke}%
\pgfsetdash{}{0pt}%
\pgfpathmoveto{\pgfqpoint{1.433375in}{1.154622in}}%
\pgfpathlineto{\pgfqpoint{1.442732in}{1.169460in}}%
\pgfpathlineto{\pgfqpoint{1.452081in}{1.184283in}}%
\pgfpathlineto{\pgfqpoint{1.428500in}{1.192464in}}%
\pgfpathlineto{\pgfqpoint{1.403629in}{1.200012in}}%
\pgfpathlineto{\pgfqpoint{1.395573in}{1.184860in}}%
\pgfpathlineto{\pgfqpoint{1.387508in}{1.169690in}}%
\pgfpathlineto{\pgfqpoint{1.411057in}{1.162458in}}%
\pgfpathlineto{\pgfqpoint{1.433375in}{1.154622in}}%
\pgfpathclose%
\pgfusepath{fill}%
\end{pgfscope}%
\begin{pgfscope}%
\pgfpathrectangle{\pgfqpoint{0.407592in}{0.222904in}}{\pgfqpoint{1.411163in}{1.411163in}}%
\pgfusepath{clip}%
\pgfsetbuttcap%
\pgfsetroundjoin%
\definecolor{currentfill}{rgb}{0.749020,0.772549,0.792157}%
\pgfsetfillcolor{currentfill}%
\pgfsetfillopacity{0.160000}%
\pgfsetlinewidth{0.000000pt}%
\definecolor{currentstroke}{rgb}{0.000000,0.000000,0.000000}%
\pgfsetstrokecolor{currentstroke}%
\pgfsetdash{}{0pt}%
\pgfpathmoveto{\pgfqpoint{0.648081in}{1.358914in}}%
\pgfpathlineto{\pgfqpoint{0.639452in}{1.373519in}}%
\pgfpathlineto{\pgfqpoint{0.630832in}{1.388110in}}%
\pgfpathlineto{\pgfqpoint{0.597732in}{1.376818in}}%
\pgfpathlineto{\pgfqpoint{0.566398in}{1.364774in}}%
\pgfpathlineto{\pgfqpoint{0.576308in}{1.350505in}}%
\pgfpathlineto{\pgfqpoint{0.586224in}{1.336227in}}%
\pgfpathlineto{\pgfqpoint{0.616296in}{1.347935in}}%
\pgfpathlineto{\pgfqpoint{0.648081in}{1.358914in}}%
\pgfpathclose%
\pgfusepath{fill}%
\end{pgfscope}%
\begin{pgfscope}%
\pgfpathrectangle{\pgfqpoint{0.407592in}{0.222904in}}{\pgfqpoint{1.411163in}{1.411163in}}%
\pgfusepath{clip}%
\pgfsetbuttcap%
\pgfsetroundjoin%
\definecolor{currentfill}{rgb}{0.749020,0.772549,0.792157}%
\pgfsetfillcolor{currentfill}%
\pgfsetfillopacity{0.160000}%
\pgfsetlinewidth{0.000000pt}%
\definecolor{currentstroke}{rgb}{0.000000,0.000000,0.000000}%
\pgfsetstrokecolor{currentstroke}%
\pgfsetdash{}{0pt}%
\pgfpathmoveto{\pgfqpoint{0.818636in}{1.168289in}}%
\pgfpathlineto{\pgfqpoint{0.811263in}{1.183428in}}%
\pgfpathlineto{\pgfqpoint{0.803900in}{1.198550in}}%
\pgfpathlineto{\pgfqpoint{0.777502in}{1.190873in}}%
\pgfpathlineto{\pgfqpoint{0.752268in}{1.182567in}}%
\pgfpathlineto{\pgfqpoint{0.761004in}{1.167780in}}%
\pgfpathlineto{\pgfqpoint{0.769749in}{1.152979in}}%
\pgfpathlineto{\pgfqpoint{0.793635in}{1.160934in}}%
\pgfpathlineto{\pgfqpoint{0.818636in}{1.168289in}}%
\pgfpathclose%
\pgfusepath{fill}%
\end{pgfscope}%
\begin{pgfscope}%
\pgfpathrectangle{\pgfqpoint{0.407592in}{0.222904in}}{\pgfqpoint{1.411163in}{1.411163in}}%
\pgfusepath{clip}%
\pgfsetbuttcap%
\pgfsetroundjoin%
\definecolor{currentfill}{rgb}{0.749020,0.772549,0.792157}%
\pgfsetfillcolor{currentfill}%
\pgfsetfillopacity{0.160000}%
\pgfsetlinewidth{0.000000pt}%
\definecolor{currentstroke}{rgb}{0.000000,0.000000,0.000000}%
\pgfsetstrokecolor{currentstroke}%
\pgfsetdash{}{0pt}%
\pgfpathmoveto{\pgfqpoint{1.599772in}{1.309978in}}%
\pgfpathlineto{\pgfqpoint{1.610213in}{1.324310in}}%
\pgfpathlineto{\pgfqpoint{1.620647in}{1.338633in}}%
\pgfpathlineto{\pgfqpoint{1.593031in}{1.350197in}}%
\pgfpathlineto{\pgfqpoint{1.563549in}{1.361026in}}%
\pgfpathlineto{\pgfqpoint{1.554310in}{1.346377in}}%
\pgfpathlineto{\pgfqpoint{1.545062in}{1.331713in}}%
\pgfpathlineto{\pgfqpoint{1.573320in}{1.321201in}}%
\pgfpathlineto{\pgfqpoint{1.599772in}{1.309978in}}%
\pgfpathclose%
\pgfusepath{fill}%
\end{pgfscope}%
\begin{pgfscope}%
\pgfpathrectangle{\pgfqpoint{0.407592in}{0.222904in}}{\pgfqpoint{1.411163in}{1.411163in}}%
\pgfusepath{clip}%
\pgfsetbuttcap%
\pgfsetroundjoin%
\definecolor{currentfill}{rgb}{0.749020,0.772549,0.792157}%
\pgfsetfillcolor{currentfill}%
\pgfsetfillopacity{0.160000}%
\pgfsetlinewidth{0.000000pt}%
\definecolor{currentstroke}{rgb}{0.000000,0.000000,0.000000}%
\pgfsetstrokecolor{currentstroke}%
\pgfsetdash{}{0pt}%
\pgfpathmoveto{\pgfqpoint{1.209403in}{0.970146in}}%
\pgfpathlineto{\pgfqpoint{1.214715in}{0.986176in}}%
\pgfpathlineto{\pgfqpoint{1.220018in}{1.002179in}}%
\pgfpathlineto{\pgfqpoint{1.201073in}{1.004980in}}%
\pgfpathlineto{\pgfqpoint{1.181688in}{1.007274in}}%
\pgfpathlineto{\pgfqpoint{1.177957in}{0.991087in}}%
\pgfpathlineto{\pgfqpoint{1.174219in}{0.974872in}}%
\pgfpathlineto{\pgfqpoint{1.192014in}{0.972743in}}%
\pgfpathlineto{\pgfqpoint{1.209403in}{0.970146in}}%
\pgfpathclose%
\pgfusepath{fill}%
\end{pgfscope}%
\begin{pgfscope}%
\pgfpathrectangle{\pgfqpoint{0.407592in}{0.222904in}}{\pgfqpoint{1.411163in}{1.411163in}}%
\pgfusepath{clip}%
\pgfsetbuttcap%
\pgfsetroundjoin%
\definecolor{currentfill}{rgb}{0.749020,0.772549,0.792157}%
\pgfsetfillcolor{currentfill}%
\pgfsetfillopacity{0.160000}%
\pgfsetlinewidth{0.000000pt}%
\definecolor{currentstroke}{rgb}{0.000000,0.000000,0.000000}%
\pgfsetstrokecolor{currentstroke}%
\pgfsetdash{}{0pt}%
\pgfpathmoveto{\pgfqpoint{1.133458in}{0.944948in}}%
\pgfpathlineto{\pgfqpoint{1.135576in}{0.961336in}}%
\pgfpathlineto{\pgfqpoint{1.137691in}{0.977694in}}%
\pgfpathlineto{\pgfqpoint{1.119102in}{0.978378in}}%
\pgfpathlineto{\pgfqpoint{1.100394in}{0.978573in}}%
\pgfpathlineto{\pgfqpoint{1.099942in}{0.962181in}}%
\pgfpathlineto{\pgfqpoint{1.099488in}{0.945757in}}%
\pgfpathlineto{\pgfqpoint{1.116528in}{0.945577in}}%
\pgfpathlineto{\pgfqpoint{1.133458in}{0.944948in}}%
\pgfpathclose%
\pgfusepath{fill}%
\end{pgfscope}%
\begin{pgfscope}%
\pgfpathrectangle{\pgfqpoint{0.407592in}{0.222904in}}{\pgfqpoint{1.411163in}{1.411163in}}%
\pgfusepath{clip}%
\pgfsetbuttcap%
\pgfsetroundjoin%
\definecolor{currentfill}{rgb}{0.749020,0.772549,0.792157}%
\pgfsetfillcolor{currentfill}%
\pgfsetfillopacity{0.160000}%
\pgfsetlinewidth{0.000000pt}%
\definecolor{currentstroke}{rgb}{0.000000,0.000000,0.000000}%
\pgfsetstrokecolor{currentstroke}%
\pgfsetdash{}{0pt}%
\pgfpathmoveto{\pgfqpoint{1.322616in}{1.047633in}}%
\pgfpathlineto{\pgfqpoint{1.330763in}{1.062957in}}%
\pgfpathlineto{\pgfqpoint{1.338900in}{1.078263in}}%
\pgfpathlineto{\pgfqpoint{1.318333in}{1.083956in}}%
\pgfpathlineto{\pgfqpoint{1.296878in}{1.089094in}}%
\pgfpathlineto{\pgfqpoint{1.290139in}{1.073487in}}%
\pgfpathlineto{\pgfqpoint{1.283390in}{1.057857in}}%
\pgfpathlineto{\pgfqpoint{1.303421in}{1.053007in}}%
\pgfpathlineto{\pgfqpoint{1.322616in}{1.047633in}}%
\pgfpathclose%
\pgfusepath{fill}%
\end{pgfscope}%
\begin{pgfscope}%
\pgfpathrectangle{\pgfqpoint{0.407592in}{0.222904in}}{\pgfqpoint{1.411163in}{1.411163in}}%
\pgfusepath{clip}%
\pgfsetbuttcap%
\pgfsetroundjoin%
\definecolor{currentfill}{rgb}{0.749020,0.772549,0.792157}%
\pgfsetfillcolor{currentfill}%
\pgfsetfillopacity{0.160000}%
\pgfsetlinewidth{0.000000pt}%
\definecolor{currentstroke}{rgb}{0.000000,0.000000,0.000000}%
\pgfsetstrokecolor{currentstroke}%
\pgfsetdash{}{0pt}%
\pgfpathmoveto{\pgfqpoint{1.099488in}{0.945757in}}%
\pgfpathlineto{\pgfqpoint{1.099942in}{0.962181in}}%
\pgfpathlineto{\pgfqpoint{1.100394in}{0.978573in}}%
\pgfpathlineto{\pgfqpoint{1.081640in}{0.978280in}}%
\pgfpathlineto{\pgfqpoint{1.062914in}{0.977498in}}%
\pgfpathlineto{\pgfqpoint{1.064131in}{0.961148in}}%
\pgfpathlineto{\pgfqpoint{1.065350in}{0.944768in}}%
\pgfpathlineto{\pgfqpoint{1.082406in}{0.945487in}}%
\pgfpathlineto{\pgfqpoint{1.099488in}{0.945757in}}%
\pgfpathclose%
\pgfusepath{fill}%
\end{pgfscope}%
\begin{pgfscope}%
\pgfpathrectangle{\pgfqpoint{0.407592in}{0.222904in}}{\pgfqpoint{1.411163in}{1.411163in}}%
\pgfusepath{clip}%
\pgfsetbuttcap%
\pgfsetroundjoin%
\definecolor{currentfill}{rgb}{0.749020,0.772549,0.792157}%
\pgfsetfillcolor{currentfill}%
\pgfsetfillopacity{0.160000}%
\pgfsetlinewidth{0.000000pt}%
\definecolor{currentstroke}{rgb}{0.000000,0.000000,0.000000}%
\pgfsetstrokecolor{currentstroke}%
\pgfsetdash{}{0pt}%
\pgfpathmoveto{\pgfqpoint{1.025837in}{0.974483in}}%
\pgfpathlineto{\pgfqpoint{1.022971in}{0.990683in}}%
\pgfpathlineto{\pgfqpoint{1.020110in}{1.006855in}}%
\pgfpathlineto{\pgfqpoint{1.000280in}{1.004458in}}%
\pgfpathlineto{\pgfqpoint{0.980795in}{1.001558in}}%
\pgfpathlineto{\pgfqpoint{0.985267in}{0.985577in}}%
\pgfpathlineto{\pgfqpoint{0.989747in}{0.969569in}}%
\pgfpathlineto{\pgfqpoint{1.007632in}{0.972260in}}%
\pgfpathlineto{\pgfqpoint{1.025837in}{0.974483in}}%
\pgfpathclose%
\pgfusepath{fill}%
\end{pgfscope}%
\begin{pgfscope}%
\pgfpathrectangle{\pgfqpoint{0.407592in}{0.222904in}}{\pgfqpoint{1.411163in}{1.411163in}}%
\pgfusepath{clip}%
\pgfsetbuttcap%
\pgfsetroundjoin%
\definecolor{currentfill}{rgb}{0.749020,0.772549,0.792157}%
\pgfsetfillcolor{currentfill}%
\pgfsetfillopacity{0.160000}%
\pgfsetlinewidth{0.000000pt}%
\definecolor{currentstroke}{rgb}{0.000000,0.000000,0.000000}%
\pgfsetstrokecolor{currentstroke}%
\pgfsetdash{}{0pt}%
\pgfpathmoveto{\pgfqpoint{1.256291in}{0.995102in}}%
\pgfpathlineto{\pgfqpoint{1.263081in}{1.010826in}}%
\pgfpathlineto{\pgfqpoint{1.269861in}{1.026526in}}%
\pgfpathlineto{\pgfqpoint{1.250545in}{1.030576in}}%
\pgfpathlineto{\pgfqpoint{1.230598in}{1.034105in}}%
\pgfpathlineto{\pgfqpoint{1.225313in}{1.018156in}}%
\pgfpathlineto{\pgfqpoint{1.220018in}{1.002179in}}%
\pgfpathlineto{\pgfqpoint{1.238449in}{0.998883in}}%
\pgfpathlineto{\pgfqpoint{1.256291in}{0.995102in}}%
\pgfpathclose%
\pgfusepath{fill}%
\end{pgfscope}%
\begin{pgfscope}%
\pgfpathrectangle{\pgfqpoint{0.407592in}{0.222904in}}{\pgfqpoint{1.411163in}{1.411163in}}%
\pgfusepath{clip}%
\pgfsetbuttcap%
\pgfsetroundjoin%
\definecolor{currentfill}{rgb}{0.749020,0.772549,0.792157}%
\pgfsetfillcolor{currentfill}%
\pgfsetfillopacity{0.160000}%
\pgfsetlinewidth{0.000000pt}%
\definecolor{currentstroke}{rgb}{0.000000,0.000000,0.000000}%
\pgfsetstrokecolor{currentstroke}%
\pgfsetdash{}{0pt}%
\pgfpathmoveto{\pgfqpoint{1.414631in}{1.124902in}}%
\pgfpathlineto{\pgfqpoint{1.424008in}{1.139770in}}%
\pgfpathlineto{\pgfqpoint{1.433375in}{1.154622in}}%
\pgfpathlineto{\pgfqpoint{1.411057in}{1.162458in}}%
\pgfpathlineto{\pgfqpoint{1.387508in}{1.169690in}}%
\pgfpathlineto{\pgfqpoint{1.379432in}{1.154500in}}%
\pgfpathlineto{\pgfqpoint{1.371346in}{1.139291in}}%
\pgfpathlineto{\pgfqpoint{1.393576in}{1.132384in}}%
\pgfpathlineto{\pgfqpoint{1.414631in}{1.124902in}}%
\pgfpathclose%
\pgfusepath{fill}%
\end{pgfscope}%
\begin{pgfscope}%
\pgfpathrectangle{\pgfqpoint{0.407592in}{0.222904in}}{\pgfqpoint{1.411163in}{1.411163in}}%
\pgfusepath{clip}%
\pgfsetbuttcap%
\pgfsetroundjoin%
\definecolor{currentfill}{rgb}{0.749020,0.772549,0.792157}%
\pgfsetfillcolor{currentfill}%
\pgfsetfillopacity{0.160000}%
\pgfsetlinewidth{0.000000pt}%
\definecolor{currentstroke}{rgb}{0.000000,0.000000,0.000000}%
\pgfsetstrokecolor{currentstroke}%
\pgfsetdash{}{0pt}%
\pgfpathmoveto{\pgfqpoint{0.919157in}{1.056927in}}%
\pgfpathlineto{\pgfqpoint{0.913177in}{1.072529in}}%
\pgfpathlineto{\pgfqpoint{0.907207in}{1.088109in}}%
\pgfpathlineto{\pgfqpoint{0.884743in}{1.082858in}}%
\pgfpathlineto{\pgfqpoint{0.863061in}{1.077055in}}%
\pgfpathlineto{\pgfqpoint{0.870497in}{1.061784in}}%
\pgfpathlineto{\pgfqpoint{0.877943in}{1.046493in}}%
\pgfpathlineto{\pgfqpoint{0.898181in}{1.051970in}}%
\pgfpathlineto{\pgfqpoint{0.919157in}{1.056927in}}%
\pgfpathclose%
\pgfusepath{fill}%
\end{pgfscope}%
\begin{pgfscope}%
\pgfpathrectangle{\pgfqpoint{0.407592in}{0.222904in}}{\pgfqpoint{1.411163in}{1.411163in}}%
\pgfusepath{clip}%
\pgfsetbuttcap%
\pgfsetroundjoin%
\definecolor{currentfill}{rgb}{0.749020,0.772549,0.792157}%
\pgfsetfillcolor{currentfill}%
\pgfsetfillopacity{0.160000}%
\pgfsetlinewidth{0.000000pt}%
\definecolor{currentstroke}{rgb}{0.000000,0.000000,0.000000}%
\pgfsetstrokecolor{currentstroke}%
\pgfsetdash{}{0pt}%
\pgfpathmoveto{\pgfqpoint{0.980795in}{1.001558in}}%
\pgfpathlineto{\pgfqpoint{0.976330in}{1.017512in}}%
\pgfpathlineto{\pgfqpoint{0.971872in}{1.033439in}}%
\pgfpathlineto{\pgfqpoint{0.951241in}{1.029805in}}%
\pgfpathlineto{\pgfqpoint{0.931142in}{1.025653in}}%
\pgfpathlineto{\pgfqpoint{0.937148in}{1.009981in}}%
\pgfpathlineto{\pgfqpoint{0.943162in}{0.994286in}}%
\pgfpathlineto{\pgfqpoint{0.961730in}{0.998163in}}%
\pgfpathlineto{\pgfqpoint{0.980795in}{1.001558in}}%
\pgfpathclose%
\pgfusepath{fill}%
\end{pgfscope}%
\begin{pgfscope}%
\pgfpathrectangle{\pgfqpoint{0.407592in}{0.222904in}}{\pgfqpoint{1.411163in}{1.411163in}}%
\pgfusepath{clip}%
\pgfsetbuttcap%
\pgfsetroundjoin%
\definecolor{currentfill}{rgb}{0.749020,0.772549,0.792157}%
\pgfsetfillcolor{currentfill}%
\pgfsetfillopacity{0.160000}%
\pgfsetlinewidth{0.000000pt}%
\definecolor{currentstroke}{rgb}{0.000000,0.000000,0.000000}%
\pgfsetstrokecolor{currentstroke}%
\pgfsetdash{}{0pt}%
\pgfpathmoveto{\pgfqpoint{0.665363in}{1.329663in}}%
\pgfpathlineto{\pgfqpoint{0.656718in}{1.344295in}}%
\pgfpathlineto{\pgfqpoint{0.648081in}{1.358914in}}%
\pgfpathlineto{\pgfqpoint{0.616296in}{1.347935in}}%
\pgfpathlineto{\pgfqpoint{0.586224in}{1.336227in}}%
\pgfpathlineto{\pgfqpoint{0.596147in}{1.321940in}}%
\pgfpathlineto{\pgfqpoint{0.606076in}{1.307644in}}%
\pgfpathlineto{\pgfqpoint{0.634889in}{1.319005in}}%
\pgfpathlineto{\pgfqpoint{0.665363in}{1.329663in}}%
\pgfpathclose%
\pgfusepath{fill}%
\end{pgfscope}%
\begin{pgfscope}%
\pgfpathrectangle{\pgfqpoint{0.407592in}{0.222904in}}{\pgfqpoint{1.411163in}{1.411163in}}%
\pgfusepath{clip}%
\pgfsetbuttcap%
\pgfsetroundjoin%
\definecolor{currentfill}{rgb}{0.749020,0.772549,0.792157}%
\pgfsetfillcolor{currentfill}%
\pgfsetfillopacity{0.160000}%
\pgfsetlinewidth{0.000000pt}%
\definecolor{currentstroke}{rgb}{0.000000,0.000000,0.000000}%
\pgfsetstrokecolor{currentstroke}%
\pgfsetdash{}{0pt}%
\pgfpathmoveto{\pgfqpoint{1.166722in}{0.942352in}}%
\pgfpathlineto{\pgfqpoint{1.170474in}{0.958626in}}%
\pgfpathlineto{\pgfqpoint{1.174219in}{0.974872in}}%
\pgfpathlineto{\pgfqpoint{1.156087in}{0.976524in}}%
\pgfpathlineto{\pgfqpoint{1.137691in}{0.977694in}}%
\pgfpathlineto{\pgfqpoint{1.135576in}{0.961336in}}%
\pgfpathlineto{\pgfqpoint{1.133458in}{0.944948in}}%
\pgfpathlineto{\pgfqpoint{1.150212in}{0.943872in}}%
\pgfpathlineto{\pgfqpoint{1.166722in}{0.942352in}}%
\pgfpathclose%
\pgfusepath{fill}%
\end{pgfscope}%
\begin{pgfscope}%
\pgfpathrectangle{\pgfqpoint{0.407592in}{0.222904in}}{\pgfqpoint{1.411163in}{1.411163in}}%
\pgfusepath{clip}%
\pgfsetbuttcap%
\pgfsetroundjoin%
\definecolor{currentfill}{rgb}{0.749020,0.772549,0.792157}%
\pgfsetfillcolor{currentfill}%
\pgfsetfillopacity{0.160000}%
\pgfsetlinewidth{0.000000pt}%
\definecolor{currentstroke}{rgb}{0.000000,0.000000,0.000000}%
\pgfsetstrokecolor{currentstroke}%
\pgfsetdash{}{0pt}%
\pgfpathmoveto{\pgfqpoint{1.578870in}{1.281285in}}%
\pgfpathlineto{\pgfqpoint{1.589325in}{1.295637in}}%
\pgfpathlineto{\pgfqpoint{1.599772in}{1.309978in}}%
\pgfpathlineto{\pgfqpoint{1.573320in}{1.321201in}}%
\pgfpathlineto{\pgfqpoint{1.545062in}{1.331713in}}%
\pgfpathlineto{\pgfqpoint{1.535805in}{1.317036in}}%
\pgfpathlineto{\pgfqpoint{1.526539in}{1.302343in}}%
\pgfpathlineto{\pgfqpoint{1.553576in}{1.292156in}}%
\pgfpathlineto{\pgfqpoint{1.578870in}{1.281285in}}%
\pgfpathclose%
\pgfusepath{fill}%
\end{pgfscope}%
\begin{pgfscope}%
\pgfpathrectangle{\pgfqpoint{0.407592in}{0.222904in}}{\pgfqpoint{1.411163in}{1.411163in}}%
\pgfusepath{clip}%
\pgfsetbuttcap%
\pgfsetroundjoin%
\definecolor{currentfill}{rgb}{0.749020,0.772549,0.792157}%
\pgfsetfillcolor{currentfill}%
\pgfsetfillopacity{0.160000}%
\pgfsetlinewidth{0.000000pt}%
\definecolor{currentstroke}{rgb}{0.000000,0.000000,0.000000}%
\pgfsetstrokecolor{currentstroke}%
\pgfsetdash{}{0pt}%
\pgfpathmoveto{\pgfqpoint{1.065350in}{0.944768in}}%
\pgfpathlineto{\pgfqpoint{1.064131in}{0.961148in}}%
\pgfpathlineto{\pgfqpoint{1.062914in}{0.977498in}}%
\pgfpathlineto{\pgfqpoint{1.044288in}{0.976231in}}%
\pgfpathlineto{\pgfqpoint{1.025837in}{0.974483in}}%
\pgfpathlineto{\pgfqpoint{1.028708in}{0.958253in}}%
\pgfpathlineto{\pgfqpoint{1.031585in}{0.941994in}}%
\pgfpathlineto{\pgfqpoint{1.048387in}{0.943602in}}%
\pgfpathlineto{\pgfqpoint{1.065350in}{0.944768in}}%
\pgfpathclose%
\pgfusepath{fill}%
\end{pgfscope}%
\begin{pgfscope}%
\pgfpathrectangle{\pgfqpoint{0.407592in}{0.222904in}}{\pgfqpoint{1.411163in}{1.411163in}}%
\pgfusepath{clip}%
\pgfsetbuttcap%
\pgfsetroundjoin%
\definecolor{currentfill}{rgb}{0.749020,0.772549,0.792157}%
\pgfsetfillcolor{currentfill}%
\pgfsetfillopacity{0.160000}%
\pgfsetlinewidth{0.000000pt}%
\definecolor{currentstroke}{rgb}{0.000000,0.000000,0.000000}%
\pgfsetstrokecolor{currentstroke}%
\pgfsetdash{}{0pt}%
\pgfpathmoveto{\pgfqpoint{0.833408in}{1.137953in}}%
\pgfpathlineto{\pgfqpoint{0.826017in}{1.153130in}}%
\pgfpathlineto{\pgfqpoint{0.818636in}{1.168289in}}%
\pgfpathlineto{\pgfqpoint{0.793635in}{1.160934in}}%
\pgfpathlineto{\pgfqpoint{0.769749in}{1.152979in}}%
\pgfpathlineto{\pgfqpoint{0.778502in}{1.138164in}}%
\pgfpathlineto{\pgfqpoint{0.787263in}{1.123334in}}%
\pgfpathlineto{\pgfqpoint{0.809804in}{1.130929in}}%
\pgfpathlineto{\pgfqpoint{0.833408in}{1.137953in}}%
\pgfpathclose%
\pgfusepath{fill}%
\end{pgfscope}%
\begin{pgfscope}%
\pgfpathrectangle{\pgfqpoint{0.407592in}{0.222904in}}{\pgfqpoint{1.411163in}{1.411163in}}%
\pgfusepath{clip}%
\pgfsetbuttcap%
\pgfsetroundjoin%
\definecolor{currentfill}{rgb}{0.749020,0.772549,0.792157}%
\pgfsetfillcolor{currentfill}%
\pgfsetfillopacity{0.160000}%
\pgfsetlinewidth{0.000000pt}%
\definecolor{currentstroke}{rgb}{0.000000,0.000000,0.000000}%
\pgfsetstrokecolor{currentstroke}%
\pgfsetdash{}{0pt}%
\pgfpathmoveto{\pgfqpoint{0.682677in}{1.300356in}}%
\pgfpathlineto{\pgfqpoint{0.674016in}{1.315016in}}%
\pgfpathlineto{\pgfqpoint{0.665363in}{1.329663in}}%
\pgfpathlineto{\pgfqpoint{0.634889in}{1.319005in}}%
\pgfpathlineto{\pgfqpoint{0.606076in}{1.307644in}}%
\pgfpathlineto{\pgfqpoint{0.616011in}{1.293339in}}%
\pgfpathlineto{\pgfqpoint{0.625953in}{1.279025in}}%
\pgfpathlineto{\pgfqpoint{0.653512in}{1.290029in}}%
\pgfpathlineto{\pgfqpoint{0.682677in}{1.300356in}}%
\pgfpathclose%
\pgfusepath{fill}%
\end{pgfscope}%
\begin{pgfscope}%
\pgfpathrectangle{\pgfqpoint{0.407592in}{0.222904in}}{\pgfqpoint{1.411163in}{1.411163in}}%
\pgfusepath{clip}%
\pgfsetbuttcap%
\pgfsetroundjoin%
\definecolor{currentfill}{rgb}{0.749020,0.772549,0.792157}%
\pgfsetfillcolor{currentfill}%
\pgfsetfillopacity{0.160000}%
\pgfsetlinewidth{0.000000pt}%
\definecolor{currentstroke}{rgb}{0.000000,0.000000,0.000000}%
\pgfsetstrokecolor{currentstroke}%
\pgfsetdash{}{0pt}%
\pgfpathmoveto{\pgfqpoint{1.306290in}{1.016925in}}%
\pgfpathlineto{\pgfqpoint{1.314458in}{1.032289in}}%
\pgfpathlineto{\pgfqpoint{1.322616in}{1.047633in}}%
\pgfpathlineto{\pgfqpoint{1.303421in}{1.053007in}}%
\pgfpathlineto{\pgfqpoint{1.283390in}{1.057857in}}%
\pgfpathlineto{\pgfqpoint{1.276630in}{1.042203in}}%
\pgfpathlineto{\pgfqpoint{1.269861in}{1.026526in}}%
\pgfpathlineto{\pgfqpoint{1.288468in}{1.021971in}}%
\pgfpathlineto{\pgfqpoint{1.306290in}{1.016925in}}%
\pgfpathclose%
\pgfusepath{fill}%
\end{pgfscope}%
\begin{pgfscope}%
\pgfpathrectangle{\pgfqpoint{0.407592in}{0.222904in}}{\pgfqpoint{1.411163in}{1.411163in}}%
\pgfusepath{clip}%
\pgfsetbuttcap%
\pgfsetroundjoin%
\definecolor{currentfill}{rgb}{0.749020,0.772549,0.792157}%
\pgfsetfillcolor{currentfill}%
\pgfsetfillopacity{0.160000}%
\pgfsetlinewidth{0.000000pt}%
\definecolor{currentstroke}{rgb}{0.000000,0.000000,0.000000}%
\pgfsetstrokecolor{currentstroke}%
\pgfsetdash{}{0pt}%
\pgfpathmoveto{\pgfqpoint{1.557939in}{1.252553in}}%
\pgfpathlineto{\pgfqpoint{1.568408in}{1.266924in}}%
\pgfpathlineto{\pgfqpoint{1.578870in}{1.281285in}}%
\pgfpathlineto{\pgfqpoint{1.553576in}{1.292156in}}%
\pgfpathlineto{\pgfqpoint{1.526539in}{1.302343in}}%
\pgfpathlineto{\pgfqpoint{1.517264in}{1.287636in}}%
\pgfpathlineto{\pgfqpoint{1.507980in}{1.272915in}}%
\pgfpathlineto{\pgfqpoint{1.533800in}{1.263064in}}%
\pgfpathlineto{\pgfqpoint{1.557939in}{1.252553in}}%
\pgfpathclose%
\pgfusepath{fill}%
\end{pgfscope}%
\begin{pgfscope}%
\pgfpathrectangle{\pgfqpoint{0.407592in}{0.222904in}}{\pgfqpoint{1.411163in}{1.411163in}}%
\pgfusepath{clip}%
\pgfsetbuttcap%
\pgfsetroundjoin%
\definecolor{currentfill}{rgb}{0.749020,0.772549,0.792157}%
\pgfsetfillcolor{currentfill}%
\pgfsetfillopacity{0.160000}%
\pgfsetlinewidth{0.000000pt}%
\definecolor{currentstroke}{rgb}{0.000000,0.000000,0.000000}%
\pgfsetstrokecolor{currentstroke}%
\pgfsetdash{}{0pt}%
\pgfpathmoveto{\pgfqpoint{1.395851in}{1.095124in}}%
\pgfpathlineto{\pgfqpoint{1.405246in}{1.110021in}}%
\pgfpathlineto{\pgfqpoint{1.414631in}{1.124902in}}%
\pgfpathlineto{\pgfqpoint{1.393576in}{1.132384in}}%
\pgfpathlineto{\pgfqpoint{1.371346in}{1.139291in}}%
\pgfpathlineto{\pgfqpoint{1.363250in}{1.124063in}}%
\pgfpathlineto{\pgfqpoint{1.355143in}{1.108815in}}%
\pgfpathlineto{\pgfqpoint{1.376055in}{1.102243in}}%
\pgfpathlineto{\pgfqpoint{1.395851in}{1.095124in}}%
\pgfpathclose%
\pgfusepath{fill}%
\end{pgfscope}%
\begin{pgfscope}%
\pgfpathrectangle{\pgfqpoint{0.407592in}{0.222904in}}{\pgfqpoint{1.411163in}{1.411163in}}%
\pgfusepath{clip}%
\pgfsetbuttcap%
\pgfsetroundjoin%
\definecolor{currentfill}{rgb}{0.749020,0.772549,0.792157}%
\pgfsetfillcolor{currentfill}%
\pgfsetfillopacity{0.160000}%
\pgfsetlinewidth{0.000000pt}%
\definecolor{currentstroke}{rgb}{0.000000,0.000000,0.000000}%
\pgfsetstrokecolor{currentstroke}%
\pgfsetdash{}{0pt}%
\pgfpathmoveto{\pgfqpoint{1.198751in}{0.938005in}}%
\pgfpathlineto{\pgfqpoint{1.204082in}{0.954089in}}%
\pgfpathlineto{\pgfqpoint{1.209403in}{0.970146in}}%
\pgfpathlineto{\pgfqpoint{1.192014in}{0.972743in}}%
\pgfpathlineto{\pgfqpoint{1.174219in}{0.974872in}}%
\pgfpathlineto{\pgfqpoint{1.170474in}{0.958626in}}%
\pgfpathlineto{\pgfqpoint{1.166722in}{0.942352in}}%
\pgfpathlineto{\pgfqpoint{1.182924in}{0.940394in}}%
\pgfpathlineto{\pgfqpoint{1.198751in}{0.938005in}}%
\pgfpathclose%
\pgfusepath{fill}%
\end{pgfscope}%
\begin{pgfscope}%
\pgfpathrectangle{\pgfqpoint{0.407592in}{0.222904in}}{\pgfqpoint{1.411163in}{1.411163in}}%
\pgfusepath{clip}%
\pgfsetbuttcap%
\pgfsetroundjoin%
\definecolor{currentfill}{rgb}{0.749020,0.772549,0.792157}%
\pgfsetfillcolor{currentfill}%
\pgfsetfillopacity{0.160000}%
\pgfsetlinewidth{0.000000pt}%
\definecolor{currentstroke}{rgb}{0.000000,0.000000,0.000000}%
\pgfsetstrokecolor{currentstroke}%
\pgfsetdash{}{0pt}%
\pgfpathmoveto{\pgfqpoint{0.931142in}{1.025653in}}%
\pgfpathlineto{\pgfqpoint{0.925145in}{1.041301in}}%
\pgfpathlineto{\pgfqpoint{0.919157in}{1.056927in}}%
\pgfpathlineto{\pgfqpoint{0.898181in}{1.051970in}}%
\pgfpathlineto{\pgfqpoint{0.877943in}{1.046493in}}%
\pgfpathlineto{\pgfqpoint{0.885397in}{1.031184in}}%
\pgfpathlineto{\pgfqpoint{0.892861in}{1.015856in}}%
\pgfpathlineto{\pgfqpoint{0.911656in}{1.020998in}}%
\pgfpathlineto{\pgfqpoint{0.931142in}{1.025653in}}%
\pgfpathclose%
\pgfusepath{fill}%
\end{pgfscope}%
\begin{pgfscope}%
\pgfpathrectangle{\pgfqpoint{0.407592in}{0.222904in}}{\pgfqpoint{1.411163in}{1.411163in}}%
\pgfusepath{clip}%
\pgfsetbuttcap%
\pgfsetroundjoin%
\definecolor{currentfill}{rgb}{0.749020,0.772549,0.792157}%
\pgfsetfillcolor{currentfill}%
\pgfsetfillopacity{0.160000}%
\pgfsetlinewidth{0.000000pt}%
\definecolor{currentstroke}{rgb}{0.000000,0.000000,0.000000}%
\pgfsetstrokecolor{currentstroke}%
\pgfsetdash{}{0pt}%
\pgfpathmoveto{\pgfqpoint{1.242681in}{0.963583in}}%
\pgfpathlineto{\pgfqpoint{1.249491in}{0.979355in}}%
\pgfpathlineto{\pgfqpoint{1.256291in}{0.995102in}}%
\pgfpathlineto{\pgfqpoint{1.238449in}{0.998883in}}%
\pgfpathlineto{\pgfqpoint{1.220018in}{1.002179in}}%
\pgfpathlineto{\pgfqpoint{1.214715in}{0.986176in}}%
\pgfpathlineto{\pgfqpoint{1.209403in}{0.970146in}}%
\pgfpathlineto{\pgfqpoint{1.226314in}{0.967089in}}%
\pgfpathlineto{\pgfqpoint{1.242681in}{0.963583in}}%
\pgfpathclose%
\pgfusepath{fill}%
\end{pgfscope}%
\begin{pgfscope}%
\pgfpathrectangle{\pgfqpoint{0.407592in}{0.222904in}}{\pgfqpoint{1.411163in}{1.411163in}}%
\pgfusepath{clip}%
\pgfsetbuttcap%
\pgfsetroundjoin%
\definecolor{currentfill}{rgb}{0.749020,0.772549,0.792157}%
\pgfsetfillcolor{currentfill}%
\pgfsetfillopacity{0.160000}%
\pgfsetlinewidth{0.000000pt}%
\definecolor{currentstroke}{rgb}{0.000000,0.000000,0.000000}%
\pgfsetstrokecolor{currentstroke}%
\pgfsetdash{}{0pt}%
\pgfpathmoveto{\pgfqpoint{1.031585in}{0.941994in}}%
\pgfpathlineto{\pgfqpoint{1.028708in}{0.958253in}}%
\pgfpathlineto{\pgfqpoint{1.025837in}{0.974483in}}%
\pgfpathlineto{\pgfqpoint{1.007632in}{0.972260in}}%
\pgfpathlineto{\pgfqpoint{0.989747in}{0.969569in}}%
\pgfpathlineto{\pgfqpoint{0.994234in}{0.953535in}}%
\pgfpathlineto{\pgfqpoint{0.998729in}{0.937475in}}%
\pgfpathlineto{\pgfqpoint{1.015010in}{0.939949in}}%
\pgfpathlineto{\pgfqpoint{1.031585in}{0.941994in}}%
\pgfpathclose%
\pgfusepath{fill}%
\end{pgfscope}%
\begin{pgfscope}%
\pgfpathrectangle{\pgfqpoint{0.407592in}{0.222904in}}{\pgfqpoint{1.411163in}{1.411163in}}%
\pgfusepath{clip}%
\pgfsetbuttcap%
\pgfsetroundjoin%
\definecolor{currentfill}{rgb}{0.749020,0.772549,0.792157}%
\pgfsetfillcolor{currentfill}%
\pgfsetfillopacity{0.160000}%
\pgfsetlinewidth{0.000000pt}%
\definecolor{currentstroke}{rgb}{0.000000,0.000000,0.000000}%
\pgfsetstrokecolor{currentstroke}%
\pgfsetdash{}{0pt}%
\pgfpathmoveto{\pgfqpoint{1.129209in}{0.912078in}}%
\pgfpathlineto{\pgfqpoint{1.131335in}{0.928529in}}%
\pgfpathlineto{\pgfqpoint{1.133458in}{0.944948in}}%
\pgfpathlineto{\pgfqpoint{1.116528in}{0.945577in}}%
\pgfpathlineto{\pgfqpoint{1.099488in}{0.945757in}}%
\pgfpathlineto{\pgfqpoint{1.099034in}{0.929302in}}%
\pgfpathlineto{\pgfqpoint{1.098578in}{0.912815in}}%
\pgfpathlineto{\pgfqpoint{1.113943in}{0.912651in}}%
\pgfpathlineto{\pgfqpoint{1.129209in}{0.912078in}}%
\pgfpathclose%
\pgfusepath{fill}%
\end{pgfscope}%
\begin{pgfscope}%
\pgfpathrectangle{\pgfqpoint{0.407592in}{0.222904in}}{\pgfqpoint{1.411163in}{1.411163in}}%
\pgfusepath{clip}%
\pgfsetbuttcap%
\pgfsetroundjoin%
\definecolor{currentfill}{rgb}{0.749020,0.772549,0.792157}%
\pgfsetfillcolor{currentfill}%
\pgfsetfillopacity{0.160000}%
\pgfsetlinewidth{0.000000pt}%
\definecolor{currentstroke}{rgb}{0.000000,0.000000,0.000000}%
\pgfsetstrokecolor{currentstroke}%
\pgfsetdash{}{0pt}%
\pgfpathmoveto{\pgfqpoint{1.098578in}{0.912815in}}%
\pgfpathlineto{\pgfqpoint{1.099034in}{0.929302in}}%
\pgfpathlineto{\pgfqpoint{1.099488in}{0.945757in}}%
\pgfpathlineto{\pgfqpoint{1.082406in}{0.945487in}}%
\pgfpathlineto{\pgfqpoint{1.065350in}{0.944768in}}%
\pgfpathlineto{\pgfqpoint{1.066572in}{0.928356in}}%
\pgfpathlineto{\pgfqpoint{1.067796in}{0.911914in}}%
\pgfpathlineto{\pgfqpoint{1.083175in}{0.912569in}}%
\pgfpathlineto{\pgfqpoint{1.098578in}{0.912815in}}%
\pgfpathclose%
\pgfusepath{fill}%
\end{pgfscope}%
\begin{pgfscope}%
\pgfpathrectangle{\pgfqpoint{0.407592in}{0.222904in}}{\pgfqpoint{1.411163in}{1.411163in}}%
\pgfusepath{clip}%
\pgfsetbuttcap%
\pgfsetroundjoin%
\definecolor{currentfill}{rgb}{0.749020,0.772549,0.792157}%
\pgfsetfillcolor{currentfill}%
\pgfsetfillopacity{0.160000}%
\pgfsetlinewidth{0.000000pt}%
\definecolor{currentstroke}{rgb}{0.000000,0.000000,0.000000}%
\pgfsetstrokecolor{currentstroke}%
\pgfsetdash{}{0pt}%
\pgfpathmoveto{\pgfqpoint{0.848216in}{1.107541in}}%
\pgfpathlineto{\pgfqpoint{0.840807in}{1.122756in}}%
\pgfpathlineto{\pgfqpoint{0.833408in}{1.137953in}}%
\pgfpathlineto{\pgfqpoint{0.809804in}{1.130929in}}%
\pgfpathlineto{\pgfqpoint{0.787263in}{1.123334in}}%
\pgfpathlineto{\pgfqpoint{0.796033in}{1.108491in}}%
\pgfpathlineto{\pgfqpoint{0.804811in}{1.093633in}}%
\pgfpathlineto{\pgfqpoint{0.826008in}{1.100858in}}%
\pgfpathlineto{\pgfqpoint{0.848216in}{1.107541in}}%
\pgfpathclose%
\pgfusepath{fill}%
\end{pgfscope}%
\begin{pgfscope}%
\pgfpathrectangle{\pgfqpoint{0.407592in}{0.222904in}}{\pgfqpoint{1.411163in}{1.411163in}}%
\pgfusepath{clip}%
\pgfsetbuttcap%
\pgfsetroundjoin%
\definecolor{currentfill}{rgb}{0.749020,0.772549,0.792157}%
\pgfsetfillcolor{currentfill}%
\pgfsetfillopacity{0.160000}%
\pgfsetlinewidth{0.000000pt}%
\definecolor{currentstroke}{rgb}{0.000000,0.000000,0.000000}%
\pgfsetstrokecolor{currentstroke}%
\pgfsetdash{}{0pt}%
\pgfpathmoveto{\pgfqpoint{0.989747in}{0.969569in}}%
\pgfpathlineto{\pgfqpoint{0.985267in}{0.985577in}}%
\pgfpathlineto{\pgfqpoint{0.980795in}{1.001558in}}%
\pgfpathlineto{\pgfqpoint{0.961730in}{0.998163in}}%
\pgfpathlineto{\pgfqpoint{0.943162in}{0.994286in}}%
\pgfpathlineto{\pgfqpoint{0.949186in}{0.978568in}}%
\pgfpathlineto{\pgfqpoint{0.955219in}{0.962827in}}%
\pgfpathlineto{\pgfqpoint{0.972252in}{0.966422in}}%
\pgfpathlineto{\pgfqpoint{0.989747in}{0.969569in}}%
\pgfpathclose%
\pgfusepath{fill}%
\end{pgfscope}%
\begin{pgfscope}%
\pgfpathrectangle{\pgfqpoint{0.407592in}{0.222904in}}{\pgfqpoint{1.411163in}{1.411163in}}%
\pgfusepath{clip}%
\pgfsetbuttcap%
\pgfsetroundjoin%
\definecolor{currentfill}{rgb}{0.749020,0.772549,0.792157}%
\pgfsetfillcolor{currentfill}%
\pgfsetfillopacity{0.160000}%
\pgfsetlinewidth{0.000000pt}%
\definecolor{currentstroke}{rgb}{0.000000,0.000000,0.000000}%
\pgfsetstrokecolor{currentstroke}%
\pgfsetdash{}{0pt}%
\pgfpathmoveto{\pgfqpoint{0.700025in}{1.270993in}}%
\pgfpathlineto{\pgfqpoint{0.691347in}{1.285681in}}%
\pgfpathlineto{\pgfqpoint{0.682677in}{1.300356in}}%
\pgfpathlineto{\pgfqpoint{0.653512in}{1.290029in}}%
\pgfpathlineto{\pgfqpoint{0.625953in}{1.279025in}}%
\pgfpathlineto{\pgfqpoint{0.635901in}{1.264701in}}%
\pgfpathlineto{\pgfqpoint{0.645855in}{1.250368in}}%
\pgfpathlineto{\pgfqpoint{0.672165in}{1.261007in}}%
\pgfpathlineto{\pgfqpoint{0.700025in}{1.270993in}}%
\pgfpathclose%
\pgfusepath{fill}%
\end{pgfscope}%
\begin{pgfscope}%
\pgfpathrectangle{\pgfqpoint{0.407592in}{0.222904in}}{\pgfqpoint{1.411163in}{1.411163in}}%
\pgfusepath{clip}%
\pgfsetbuttcap%
\pgfsetroundjoin%
\definecolor{currentfill}{rgb}{0.749020,0.772549,0.792157}%
\pgfsetfillcolor{currentfill}%
\pgfsetfillopacity{0.160000}%
\pgfsetlinewidth{0.000000pt}%
\definecolor{currentstroke}{rgb}{0.000000,0.000000,0.000000}%
\pgfsetstrokecolor{currentstroke}%
\pgfsetdash{}{0pt}%
\pgfpathmoveto{\pgfqpoint{1.536980in}{1.223783in}}%
\pgfpathlineto{\pgfqpoint{1.547463in}{1.238173in}}%
\pgfpathlineto{\pgfqpoint{1.557939in}{1.252553in}}%
\pgfpathlineto{\pgfqpoint{1.533800in}{1.263064in}}%
\pgfpathlineto{\pgfqpoint{1.507980in}{1.272915in}}%
\pgfpathlineto{\pgfqpoint{1.498686in}{1.258179in}}%
\pgfpathlineto{\pgfqpoint{1.489383in}{1.243429in}}%
\pgfpathlineto{\pgfqpoint{1.513991in}{1.233922in}}%
\pgfpathlineto{\pgfqpoint{1.536980in}{1.223783in}}%
\pgfpathclose%
\pgfusepath{fill}%
\end{pgfscope}%
\begin{pgfscope}%
\pgfpathrectangle{\pgfqpoint{0.407592in}{0.222904in}}{\pgfqpoint{1.411163in}{1.411163in}}%
\pgfusepath{clip}%
\pgfsetbuttcap%
\pgfsetroundjoin%
\definecolor{currentfill}{rgb}{0.749020,0.772549,0.792157}%
\pgfsetfillcolor{currentfill}%
\pgfsetfillopacity{0.160000}%
\pgfsetlinewidth{0.000000pt}%
\definecolor{currentstroke}{rgb}{0.000000,0.000000,0.000000}%
\pgfsetstrokecolor{currentstroke}%
\pgfsetdash{}{0pt}%
\pgfpathmoveto{\pgfqpoint{1.159198in}{0.909713in}}%
\pgfpathlineto{\pgfqpoint{1.162964in}{0.926047in}}%
\pgfpathlineto{\pgfqpoint{1.166722in}{0.942352in}}%
\pgfpathlineto{\pgfqpoint{1.150212in}{0.943872in}}%
\pgfpathlineto{\pgfqpoint{1.133458in}{0.944948in}}%
\pgfpathlineto{\pgfqpoint{1.131335in}{0.928529in}}%
\pgfpathlineto{\pgfqpoint{1.129209in}{0.912078in}}%
\pgfpathlineto{\pgfqpoint{1.144314in}{0.911098in}}%
\pgfpathlineto{\pgfqpoint{1.159198in}{0.909713in}}%
\pgfpathclose%
\pgfusepath{fill}%
\end{pgfscope}%
\begin{pgfscope}%
\pgfpathrectangle{\pgfqpoint{0.407592in}{0.222904in}}{\pgfqpoint{1.411163in}{1.411163in}}%
\pgfusepath{clip}%
\pgfsetbuttcap%
\pgfsetroundjoin%
\definecolor{currentfill}{rgb}{0.749020,0.772549,0.792157}%
\pgfsetfillcolor{currentfill}%
\pgfsetfillopacity{0.160000}%
\pgfsetlinewidth{0.000000pt}%
\definecolor{currentstroke}{rgb}{0.000000,0.000000,0.000000}%
\pgfsetstrokecolor{currentstroke}%
\pgfsetdash{}{0pt}%
\pgfpathmoveto{\pgfqpoint{1.067796in}{0.911914in}}%
\pgfpathlineto{\pgfqpoint{1.066572in}{0.928356in}}%
\pgfpathlineto{\pgfqpoint{1.065350in}{0.944768in}}%
\pgfpathlineto{\pgfqpoint{1.048387in}{0.943602in}}%
\pgfpathlineto{\pgfqpoint{1.031585in}{0.941994in}}%
\pgfpathlineto{\pgfqpoint{1.034467in}{0.925705in}}%
\pgfpathlineto{\pgfqpoint{1.037354in}{0.909387in}}%
\pgfpathlineto{\pgfqpoint{1.052502in}{0.910852in}}%
\pgfpathlineto{\pgfqpoint{1.067796in}{0.911914in}}%
\pgfpathclose%
\pgfusepath{fill}%
\end{pgfscope}%
\begin{pgfscope}%
\pgfpathrectangle{\pgfqpoint{0.407592in}{0.222904in}}{\pgfqpoint{1.411163in}{1.411163in}}%
\pgfusepath{clip}%
\pgfsetbuttcap%
\pgfsetroundjoin%
\definecolor{currentfill}{rgb}{0.749020,0.772549,0.792157}%
\pgfsetfillcolor{currentfill}%
\pgfsetfillopacity{0.160000}%
\pgfsetlinewidth{0.000000pt}%
\definecolor{currentstroke}{rgb}{0.000000,0.000000,0.000000}%
\pgfsetstrokecolor{currentstroke}%
\pgfsetdash{}{0pt}%
\pgfpathmoveto{\pgfqpoint{1.377033in}{1.065287in}}%
\pgfpathlineto{\pgfqpoint{1.386447in}{1.080213in}}%
\pgfpathlineto{\pgfqpoint{1.395851in}{1.095124in}}%
\pgfpathlineto{\pgfqpoint{1.376055in}{1.102243in}}%
\pgfpathlineto{\pgfqpoint{1.355143in}{1.108815in}}%
\pgfpathlineto{\pgfqpoint{1.347027in}{1.093549in}}%
\pgfpathlineto{\pgfqpoint{1.338900in}{1.078263in}}%
\pgfpathlineto{\pgfqpoint{1.358494in}{1.072033in}}%
\pgfpathlineto{\pgfqpoint{1.377033in}{1.065287in}}%
\pgfpathclose%
\pgfusepath{fill}%
\end{pgfscope}%
\begin{pgfscope}%
\pgfpathrectangle{\pgfqpoint{0.407592in}{0.222904in}}{\pgfqpoint{1.411163in}{1.411163in}}%
\pgfusepath{clip}%
\pgfsetbuttcap%
\pgfsetroundjoin%
\definecolor{currentfill}{rgb}{0.749020,0.772549,0.792157}%
\pgfsetfillcolor{currentfill}%
\pgfsetfillopacity{0.160000}%
\pgfsetlinewidth{0.000000pt}%
\definecolor{currentstroke}{rgb}{0.000000,0.000000,0.000000}%
\pgfsetstrokecolor{currentstroke}%
\pgfsetdash{}{0pt}%
\pgfpathmoveto{\pgfqpoint{1.289923in}{0.986140in}}%
\pgfpathlineto{\pgfqpoint{1.298111in}{1.001543in}}%
\pgfpathlineto{\pgfqpoint{1.306290in}{1.016925in}}%
\pgfpathlineto{\pgfqpoint{1.288468in}{1.021971in}}%
\pgfpathlineto{\pgfqpoint{1.269861in}{1.026526in}}%
\pgfpathlineto{\pgfqpoint{1.263081in}{1.010826in}}%
\pgfpathlineto{\pgfqpoint{1.256291in}{0.995102in}}%
\pgfpathlineto{\pgfqpoint{1.273473in}{0.990849in}}%
\pgfpathlineto{\pgfqpoint{1.289923in}{0.986140in}}%
\pgfpathclose%
\pgfusepath{fill}%
\end{pgfscope}%
\begin{pgfscope}%
\pgfpathrectangle{\pgfqpoint{0.407592in}{0.222904in}}{\pgfqpoint{1.411163in}{1.411163in}}%
\pgfusepath{clip}%
\pgfsetbuttcap%
\pgfsetroundjoin%
\definecolor{currentfill}{rgb}{0.749020,0.772549,0.792157}%
\pgfsetfillcolor{currentfill}%
\pgfsetfillopacity{0.160000}%
\pgfsetlinewidth{0.000000pt}%
\definecolor{currentstroke}{rgb}{0.000000,0.000000,0.000000}%
\pgfsetstrokecolor{currentstroke}%
\pgfsetdash{}{0pt}%
\pgfpathmoveto{\pgfqpoint{0.717406in}{1.241574in}}%
\pgfpathlineto{\pgfqpoint{0.708711in}{1.256291in}}%
\pgfpathlineto{\pgfqpoint{0.700025in}{1.270993in}}%
\pgfpathlineto{\pgfqpoint{0.672165in}{1.261007in}}%
\pgfpathlineto{\pgfqpoint{0.645855in}{1.250368in}}%
\pgfpathlineto{\pgfqpoint{0.655815in}{1.236026in}}%
\pgfpathlineto{\pgfqpoint{0.665782in}{1.221675in}}%
\pgfpathlineto{\pgfqpoint{0.690848in}{1.231938in}}%
\pgfpathlineto{\pgfqpoint{0.717406in}{1.241574in}}%
\pgfpathclose%
\pgfusepath{fill}%
\end{pgfscope}%
\begin{pgfscope}%
\pgfpathrectangle{\pgfqpoint{0.407592in}{0.222904in}}{\pgfqpoint{1.411163in}{1.411163in}}%
\pgfusepath{clip}%
\pgfsetbuttcap%
\pgfsetroundjoin%
\definecolor{currentfill}{rgb}{0.749020,0.772549,0.792157}%
\pgfsetfillcolor{currentfill}%
\pgfsetfillopacity{0.160000}%
\pgfsetlinewidth{0.000000pt}%
\definecolor{currentstroke}{rgb}{0.000000,0.000000,0.000000}%
\pgfsetstrokecolor{currentstroke}%
\pgfsetdash{}{0pt}%
\pgfpathmoveto{\pgfqpoint{0.863061in}{1.077055in}}%
\pgfpathlineto{\pgfqpoint{0.855634in}{1.092308in}}%
\pgfpathlineto{\pgfqpoint{0.848216in}{1.107541in}}%
\pgfpathlineto{\pgfqpoint{0.826008in}{1.100858in}}%
\pgfpathlineto{\pgfqpoint{0.804811in}{1.093633in}}%
\pgfpathlineto{\pgfqpoint{0.813598in}{1.078760in}}%
\pgfpathlineto{\pgfqpoint{0.822393in}{1.063874in}}%
\pgfpathlineto{\pgfqpoint{0.842249in}{1.070720in}}%
\pgfpathlineto{\pgfqpoint{0.863061in}{1.077055in}}%
\pgfpathclose%
\pgfusepath{fill}%
\end{pgfscope}%
\begin{pgfscope}%
\pgfpathrectangle{\pgfqpoint{0.407592in}{0.222904in}}{\pgfqpoint{1.411163in}{1.411163in}}%
\pgfusepath{clip}%
\pgfsetbuttcap%
\pgfsetroundjoin%
\definecolor{currentfill}{rgb}{0.749020,0.772549,0.792157}%
\pgfsetfillcolor{currentfill}%
\pgfsetfillopacity{0.160000}%
\pgfsetlinewidth{0.000000pt}%
\definecolor{currentstroke}{rgb}{0.000000,0.000000,0.000000}%
\pgfsetstrokecolor{currentstroke}%
\pgfsetdash{}{0pt}%
\pgfpathmoveto{\pgfqpoint{1.515992in}{1.194974in}}%
\pgfpathlineto{\pgfqpoint{1.526489in}{1.209383in}}%
\pgfpathlineto{\pgfqpoint{1.536980in}{1.223783in}}%
\pgfpathlineto{\pgfqpoint{1.513991in}{1.233922in}}%
\pgfpathlineto{\pgfqpoint{1.489383in}{1.243429in}}%
\pgfpathlineto{\pgfqpoint{1.480072in}{1.228664in}}%
\pgfpathlineto{\pgfqpoint{1.470751in}{1.213885in}}%
\pgfpathlineto{\pgfqpoint{1.494148in}{1.204732in}}%
\pgfpathlineto{\pgfqpoint{1.515992in}{1.194974in}}%
\pgfpathclose%
\pgfusepath{fill}%
\end{pgfscope}%
\begin{pgfscope}%
\pgfpathrectangle{\pgfqpoint{0.407592in}{0.222904in}}{\pgfqpoint{1.411163in}{1.411163in}}%
\pgfusepath{clip}%
\pgfsetbuttcap%
\pgfsetroundjoin%
\definecolor{currentfill}{rgb}{0.749020,0.772549,0.792157}%
\pgfsetfillcolor{currentfill}%
\pgfsetfillopacity{0.160000}%
\pgfsetlinewidth{0.000000pt}%
\definecolor{currentstroke}{rgb}{0.000000,0.000000,0.000000}%
\pgfsetstrokecolor{currentstroke}%
\pgfsetdash{}{0pt}%
\pgfpathmoveto{\pgfqpoint{0.943162in}{0.994286in}}%
\pgfpathlineto{\pgfqpoint{0.937148in}{1.009981in}}%
\pgfpathlineto{\pgfqpoint{0.931142in}{1.025653in}}%
\pgfpathlineto{\pgfqpoint{0.911656in}{1.020998in}}%
\pgfpathlineto{\pgfqpoint{0.892861in}{1.015856in}}%
\pgfpathlineto{\pgfqpoint{0.900335in}{1.000508in}}%
\pgfpathlineto{\pgfqpoint{0.907817in}{0.985142in}}%
\pgfpathlineto{\pgfqpoint{0.925167in}{0.989941in}}%
\pgfpathlineto{\pgfqpoint{0.943162in}{0.994286in}}%
\pgfpathclose%
\pgfusepath{fill}%
\end{pgfscope}%
\begin{pgfscope}%
\pgfpathrectangle{\pgfqpoint{0.407592in}{0.222904in}}{\pgfqpoint{1.411163in}{1.411163in}}%
\pgfusepath{clip}%
\pgfsetbuttcap%
\pgfsetroundjoin%
\definecolor{currentfill}{rgb}{0.749020,0.772549,0.792157}%
\pgfsetfillcolor{currentfill}%
\pgfsetfillopacity{0.160000}%
\pgfsetlinewidth{0.000000pt}%
\definecolor{currentstroke}{rgb}{0.000000,0.000000,0.000000}%
\pgfsetstrokecolor{currentstroke}%
\pgfsetdash{}{0pt}%
\pgfpathmoveto{\pgfqpoint{1.229030in}{0.931970in}}%
\pgfpathlineto{\pgfqpoint{1.235861in}{0.947789in}}%
\pgfpathlineto{\pgfqpoint{1.242681in}{0.963583in}}%
\pgfpathlineto{\pgfqpoint{1.226314in}{0.967089in}}%
\pgfpathlineto{\pgfqpoint{1.209403in}{0.970146in}}%
\pgfpathlineto{\pgfqpoint{1.204082in}{0.954089in}}%
\pgfpathlineto{\pgfqpoint{1.198751in}{0.938005in}}%
\pgfpathlineto{\pgfqpoint{1.214141in}{0.935194in}}%
\pgfpathlineto{\pgfqpoint{1.229030in}{0.931970in}}%
\pgfpathclose%
\pgfusepath{fill}%
\end{pgfscope}%
\begin{pgfscope}%
\pgfpathrectangle{\pgfqpoint{0.407592in}{0.222904in}}{\pgfqpoint{1.411163in}{1.411163in}}%
\pgfusepath{clip}%
\pgfsetbuttcap%
\pgfsetroundjoin%
\definecolor{currentfill}{rgb}{0.749020,0.772549,0.792157}%
\pgfsetfillcolor{currentfill}%
\pgfsetfillopacity{0.160000}%
\pgfsetlinewidth{0.000000pt}%
\definecolor{currentstroke}{rgb}{0.000000,0.000000,0.000000}%
\pgfsetstrokecolor{currentstroke}%
\pgfsetdash{}{0pt}%
\pgfpathmoveto{\pgfqpoint{1.188064in}{0.905755in}}%
\pgfpathlineto{\pgfqpoint{1.193412in}{0.921894in}}%
\pgfpathlineto{\pgfqpoint{1.198751in}{0.938005in}}%
\pgfpathlineto{\pgfqpoint{1.182924in}{0.940394in}}%
\pgfpathlineto{\pgfqpoint{1.166722in}{0.942352in}}%
\pgfpathlineto{\pgfqpoint{1.162964in}{0.926047in}}%
\pgfpathlineto{\pgfqpoint{1.159198in}{0.909713in}}%
\pgfpathlineto{\pgfqpoint{1.173801in}{0.907930in}}%
\pgfpathlineto{\pgfqpoint{1.188064in}{0.905755in}}%
\pgfpathclose%
\pgfusepath{fill}%
\end{pgfscope}%
\begin{pgfscope}%
\pgfpathrectangle{\pgfqpoint{0.407592in}{0.222904in}}{\pgfqpoint{1.411163in}{1.411163in}}%
\pgfusepath{clip}%
\pgfsetbuttcap%
\pgfsetroundjoin%
\definecolor{currentfill}{rgb}{0.749020,0.772549,0.792157}%
\pgfsetfillcolor{currentfill}%
\pgfsetfillopacity{0.160000}%
\pgfsetlinewidth{0.000000pt}%
\definecolor{currentstroke}{rgb}{0.000000,0.000000,0.000000}%
\pgfsetstrokecolor{currentstroke}%
\pgfsetdash{}{0pt}%
\pgfpathmoveto{\pgfqpoint{0.998729in}{0.937475in}}%
\pgfpathlineto{\pgfqpoint{0.994234in}{0.953535in}}%
\pgfpathlineto{\pgfqpoint{0.989747in}{0.969569in}}%
\pgfpathlineto{\pgfqpoint{0.972252in}{0.966422in}}%
\pgfpathlineto{\pgfqpoint{0.955219in}{0.962827in}}%
\pgfpathlineto{\pgfqpoint{0.961260in}{0.947063in}}%
\pgfpathlineto{\pgfqpoint{0.967310in}{0.931275in}}%
\pgfpathlineto{\pgfqpoint{0.982808in}{0.934580in}}%
\pgfpathlineto{\pgfqpoint{0.998729in}{0.937475in}}%
\pgfpathclose%
\pgfusepath{fill}%
\end{pgfscope}%
\begin{pgfscope}%
\pgfpathrectangle{\pgfqpoint{0.407592in}{0.222904in}}{\pgfqpoint{1.411163in}{1.411163in}}%
\pgfusepath{clip}%
\pgfsetbuttcap%
\pgfsetroundjoin%
\definecolor{currentfill}{rgb}{0.749020,0.772549,0.792157}%
\pgfsetfillcolor{currentfill}%
\pgfsetfillopacity{0.160000}%
\pgfsetlinewidth{0.000000pt}%
\definecolor{currentstroke}{rgb}{0.000000,0.000000,0.000000}%
\pgfsetstrokecolor{currentstroke}%
\pgfsetdash{}{0pt}%
\pgfpathmoveto{\pgfqpoint{1.037354in}{0.909387in}}%
\pgfpathlineto{\pgfqpoint{1.034467in}{0.925705in}}%
\pgfpathlineto{\pgfqpoint{1.031585in}{0.941994in}}%
\pgfpathlineto{\pgfqpoint{1.015010in}{0.939949in}}%
\pgfpathlineto{\pgfqpoint{0.998729in}{0.937475in}}%
\pgfpathlineto{\pgfqpoint{1.003232in}{0.921387in}}%
\pgfpathlineto{\pgfqpoint{1.007742in}{0.905272in}}%
\pgfpathlineto{\pgfqpoint{1.022414in}{0.907525in}}%
\pgfpathlineto{\pgfqpoint{1.037354in}{0.909387in}}%
\pgfpathclose%
\pgfusepath{fill}%
\end{pgfscope}%
\begin{pgfscope}%
\pgfpathrectangle{\pgfqpoint{0.407592in}{0.222904in}}{\pgfqpoint{1.411163in}{1.411163in}}%
\pgfusepath{clip}%
\pgfsetbuttcap%
\pgfsetroundjoin%
\definecolor{currentfill}{rgb}{0.749020,0.772549,0.792157}%
\pgfsetfillcolor{currentfill}%
\pgfsetfillopacity{0.160000}%
\pgfsetlinewidth{0.000000pt}%
\definecolor{currentstroke}{rgb}{0.000000,0.000000,0.000000}%
\pgfsetstrokecolor{currentstroke}%
\pgfsetdash{}{0pt}%
\pgfpathmoveto{\pgfqpoint{0.734820in}{1.212099in}}%
\pgfpathlineto{\pgfqpoint{0.726109in}{1.226844in}}%
\pgfpathlineto{\pgfqpoint{0.717406in}{1.241574in}}%
\pgfpathlineto{\pgfqpoint{0.690848in}{1.231938in}}%
\pgfpathlineto{\pgfqpoint{0.665782in}{1.221675in}}%
\pgfpathlineto{\pgfqpoint{0.675756in}{1.207315in}}%
\pgfpathlineto{\pgfqpoint{0.685736in}{1.192946in}}%
\pgfpathlineto{\pgfqpoint{0.709561in}{1.202822in}}%
\pgfpathlineto{\pgfqpoint{0.734820in}{1.212099in}}%
\pgfpathclose%
\pgfusepath{fill}%
\end{pgfscope}%
\begin{pgfscope}%
\pgfpathrectangle{\pgfqpoint{0.407592in}{0.222904in}}{\pgfqpoint{1.411163in}{1.411163in}}%
\pgfusepath{clip}%
\pgfsetbuttcap%
\pgfsetroundjoin%
\definecolor{currentfill}{rgb}{0.749020,0.772549,0.792157}%
\pgfsetfillcolor{currentfill}%
\pgfsetfillopacity{0.160000}%
\pgfsetlinewidth{0.000000pt}%
\definecolor{currentstroke}{rgb}{0.000000,0.000000,0.000000}%
\pgfsetstrokecolor{currentstroke}%
\pgfsetdash{}{0pt}%
\pgfpathmoveto{\pgfqpoint{1.358178in}{1.035391in}}%
\pgfpathlineto{\pgfqpoint{1.367610in}{1.050346in}}%
\pgfpathlineto{\pgfqpoint{1.377033in}{1.065287in}}%
\pgfpathlineto{\pgfqpoint{1.358494in}{1.072033in}}%
\pgfpathlineto{\pgfqpoint{1.338900in}{1.078263in}}%
\pgfpathlineto{\pgfqpoint{1.330763in}{1.062957in}}%
\pgfpathlineto{\pgfqpoint{1.322616in}{1.047633in}}%
\pgfpathlineto{\pgfqpoint{1.340894in}{1.041754in}}%
\pgfpathlineto{\pgfqpoint{1.358178in}{1.035391in}}%
\pgfpathclose%
\pgfusepath{fill}%
\end{pgfscope}%
\begin{pgfscope}%
\pgfpathrectangle{\pgfqpoint{0.407592in}{0.222904in}}{\pgfqpoint{1.411163in}{1.411163in}}%
\pgfusepath{clip}%
\pgfsetbuttcap%
\pgfsetroundjoin%
\definecolor{currentfill}{rgb}{0.749020,0.772549,0.792157}%
\pgfsetfillcolor{currentfill}%
\pgfsetfillopacity{0.160000}%
\pgfsetlinewidth{0.000000pt}%
\definecolor{currentstroke}{rgb}{0.000000,0.000000,0.000000}%
\pgfsetstrokecolor{currentstroke}%
\pgfsetdash{}{0pt}%
\pgfpathmoveto{\pgfqpoint{1.124944in}{0.879083in}}%
\pgfpathlineto{\pgfqpoint{1.127078in}{0.895596in}}%
\pgfpathlineto{\pgfqpoint{1.129209in}{0.912078in}}%
\pgfpathlineto{\pgfqpoint{1.113943in}{0.912651in}}%
\pgfpathlineto{\pgfqpoint{1.098578in}{0.912815in}}%
\pgfpathlineto{\pgfqpoint{1.098122in}{0.896297in}}%
\pgfpathlineto{\pgfqpoint{1.097665in}{0.879746in}}%
\pgfpathlineto{\pgfqpoint{1.111349in}{0.879599in}}%
\pgfpathlineto{\pgfqpoint{1.124944in}{0.879083in}}%
\pgfpathclose%
\pgfusepath{fill}%
\end{pgfscope}%
\begin{pgfscope}%
\pgfpathrectangle{\pgfqpoint{0.407592in}{0.222904in}}{\pgfqpoint{1.411163in}{1.411163in}}%
\pgfusepath{clip}%
\pgfsetbuttcap%
\pgfsetroundjoin%
\definecolor{currentfill}{rgb}{0.749020,0.772549,0.792157}%
\pgfsetfillcolor{currentfill}%
\pgfsetfillopacity{0.160000}%
\pgfsetlinewidth{0.000000pt}%
\definecolor{currentstroke}{rgb}{0.000000,0.000000,0.000000}%
\pgfsetstrokecolor{currentstroke}%
\pgfsetdash{}{0pt}%
\pgfpathmoveto{\pgfqpoint{1.097665in}{0.879746in}}%
\pgfpathlineto{\pgfqpoint{1.098122in}{0.896297in}}%
\pgfpathlineto{\pgfqpoint{1.098578in}{0.912815in}}%
\pgfpathlineto{\pgfqpoint{1.083175in}{0.912569in}}%
\pgfpathlineto{\pgfqpoint{1.067796in}{0.911914in}}%
\pgfpathlineto{\pgfqpoint{1.069022in}{0.895440in}}%
\pgfpathlineto{\pgfqpoint{1.070251in}{0.878935in}}%
\pgfpathlineto{\pgfqpoint{1.083947in}{0.879525in}}%
\pgfpathlineto{\pgfqpoint{1.097665in}{0.879746in}}%
\pgfpathclose%
\pgfusepath{fill}%
\end{pgfscope}%
\begin{pgfscope}%
\pgfpathrectangle{\pgfqpoint{0.407592in}{0.222904in}}{\pgfqpoint{1.411163in}{1.411163in}}%
\pgfusepath{clip}%
\pgfsetbuttcap%
\pgfsetroundjoin%
\definecolor{currentfill}{rgb}{0.749020,0.772549,0.792157}%
\pgfsetfillcolor{currentfill}%
\pgfsetfillopacity{0.160000}%
\pgfsetlinewidth{0.000000pt}%
\definecolor{currentstroke}{rgb}{0.000000,0.000000,0.000000}%
\pgfsetstrokecolor{currentstroke}%
\pgfsetdash{}{0pt}%
\pgfpathmoveto{\pgfqpoint{1.494977in}{1.166126in}}%
\pgfpathlineto{\pgfqpoint{1.505488in}{1.180555in}}%
\pgfpathlineto{\pgfqpoint{1.515992in}{1.194974in}}%
\pgfpathlineto{\pgfqpoint{1.494148in}{1.204732in}}%
\pgfpathlineto{\pgfqpoint{1.470751in}{1.213885in}}%
\pgfpathlineto{\pgfqpoint{1.461420in}{1.199091in}}%
\pgfpathlineto{\pgfqpoint{1.452081in}{1.184283in}}%
\pgfpathlineto{\pgfqpoint{1.474273in}{1.175494in}}%
\pgfpathlineto{\pgfqpoint{1.494977in}{1.166126in}}%
\pgfpathclose%
\pgfusepath{fill}%
\end{pgfscope}%
\begin{pgfscope}%
\pgfpathrectangle{\pgfqpoint{0.407592in}{0.222904in}}{\pgfqpoint{1.411163in}{1.411163in}}%
\pgfusepath{clip}%
\pgfsetbuttcap%
\pgfsetroundjoin%
\definecolor{currentfill}{rgb}{0.749020,0.772549,0.792157}%
\pgfsetfillcolor{currentfill}%
\pgfsetfillopacity{0.160000}%
\pgfsetlinewidth{0.000000pt}%
\definecolor{currentstroke}{rgb}{0.000000,0.000000,0.000000}%
\pgfsetstrokecolor{currentstroke}%
\pgfsetdash{}{0pt}%
\pgfpathmoveto{\pgfqpoint{0.877943in}{1.046493in}}%
\pgfpathlineto{\pgfqpoint{0.870497in}{1.061784in}}%
\pgfpathlineto{\pgfqpoint{0.863061in}{1.077055in}}%
\pgfpathlineto{\pgfqpoint{0.842249in}{1.070720in}}%
\pgfpathlineto{\pgfqpoint{0.822393in}{1.063874in}}%
\pgfpathlineto{\pgfqpoint{0.831196in}{1.048973in}}%
\pgfpathlineto{\pgfqpoint{0.840008in}{1.034058in}}%
\pgfpathlineto{\pgfqpoint{0.858525in}{1.040516in}}%
\pgfpathlineto{\pgfqpoint{0.877943in}{1.046493in}}%
\pgfpathclose%
\pgfusepath{fill}%
\end{pgfscope}%
\begin{pgfscope}%
\pgfpathrectangle{\pgfqpoint{0.407592in}{0.222904in}}{\pgfqpoint{1.411163in}{1.411163in}}%
\pgfusepath{clip}%
\pgfsetbuttcap%
\pgfsetroundjoin%
\definecolor{currentfill}{rgb}{0.749020,0.772549,0.792157}%
\pgfsetfillcolor{currentfill}%
\pgfsetfillopacity{0.160000}%
\pgfsetlinewidth{0.000000pt}%
\definecolor{currentstroke}{rgb}{0.000000,0.000000,0.000000}%
\pgfsetstrokecolor{currentstroke}%
\pgfsetdash{}{0pt}%
\pgfpathmoveto{\pgfqpoint{1.273514in}{0.955277in}}%
\pgfpathlineto{\pgfqpoint{1.281724in}{0.970718in}}%
\pgfpathlineto{\pgfqpoint{1.289923in}{0.986140in}}%
\pgfpathlineto{\pgfqpoint{1.273473in}{0.990849in}}%
\pgfpathlineto{\pgfqpoint{1.256291in}{0.995102in}}%
\pgfpathlineto{\pgfqpoint{1.249491in}{0.979355in}}%
\pgfpathlineto{\pgfqpoint{1.242681in}{0.963583in}}%
\pgfpathlineto{\pgfqpoint{1.258436in}{0.959641in}}%
\pgfpathlineto{\pgfqpoint{1.273514in}{0.955277in}}%
\pgfpathclose%
\pgfusepath{fill}%
\end{pgfscope}%
\begin{pgfscope}%
\pgfpathrectangle{\pgfqpoint{0.407592in}{0.222904in}}{\pgfqpoint{1.411163in}{1.411163in}}%
\pgfusepath{clip}%
\pgfsetbuttcap%
\pgfsetroundjoin%
\definecolor{currentfill}{rgb}{0.749020,0.772549,0.792157}%
\pgfsetfillcolor{currentfill}%
\pgfsetfillopacity{0.160000}%
\pgfsetlinewidth{0.000000pt}%
\definecolor{currentstroke}{rgb}{0.000000,0.000000,0.000000}%
\pgfsetstrokecolor{currentstroke}%
\pgfsetdash{}{0pt}%
\pgfpathmoveto{\pgfqpoint{1.151647in}{0.876956in}}%
\pgfpathlineto{\pgfqpoint{1.155426in}{0.893350in}}%
\pgfpathlineto{\pgfqpoint{1.159198in}{0.909713in}}%
\pgfpathlineto{\pgfqpoint{1.144314in}{0.911098in}}%
\pgfpathlineto{\pgfqpoint{1.129209in}{0.912078in}}%
\pgfpathlineto{\pgfqpoint{1.127078in}{0.895596in}}%
\pgfpathlineto{\pgfqpoint{1.124944in}{0.879083in}}%
\pgfpathlineto{\pgfqpoint{1.138395in}{0.878201in}}%
\pgfpathlineto{\pgfqpoint{1.151647in}{0.876956in}}%
\pgfpathclose%
\pgfusepath{fill}%
\end{pgfscope}%
\begin{pgfscope}%
\pgfpathrectangle{\pgfqpoint{0.407592in}{0.222904in}}{\pgfqpoint{1.411163in}{1.411163in}}%
\pgfusepath{clip}%
\pgfsetbuttcap%
\pgfsetroundjoin%
\definecolor{currentfill}{rgb}{0.749020,0.772549,0.792157}%
\pgfsetfillcolor{currentfill}%
\pgfsetfillopacity{0.160000}%
\pgfsetlinewidth{0.000000pt}%
\definecolor{currentstroke}{rgb}{0.000000,0.000000,0.000000}%
\pgfsetstrokecolor{currentstroke}%
\pgfsetdash{}{0pt}%
\pgfpathmoveto{\pgfqpoint{1.070251in}{0.878935in}}%
\pgfpathlineto{\pgfqpoint{1.069022in}{0.895440in}}%
\pgfpathlineto{\pgfqpoint{1.067796in}{0.911914in}}%
\pgfpathlineto{\pgfqpoint{1.052502in}{0.910852in}}%
\pgfpathlineto{\pgfqpoint{1.037354in}{0.909387in}}%
\pgfpathlineto{\pgfqpoint{1.040246in}{0.893040in}}%
\pgfpathlineto{\pgfqpoint{1.043144in}{0.876663in}}%
\pgfpathlineto{\pgfqpoint{1.056631in}{0.877980in}}%
\pgfpathlineto{\pgfqpoint{1.070251in}{0.878935in}}%
\pgfpathclose%
\pgfusepath{fill}%
\end{pgfscope}%
\begin{pgfscope}%
\pgfpathrectangle{\pgfqpoint{0.407592in}{0.222904in}}{\pgfqpoint{1.411163in}{1.411163in}}%
\pgfusepath{clip}%
\pgfsetbuttcap%
\pgfsetroundjoin%
\definecolor{currentfill}{rgb}{0.749020,0.772549,0.792157}%
\pgfsetfillcolor{currentfill}%
\pgfsetfillopacity{0.160000}%
\pgfsetlinewidth{0.000000pt}%
\definecolor{currentstroke}{rgb}{0.000000,0.000000,0.000000}%
\pgfsetstrokecolor{currentstroke}%
\pgfsetdash{}{0pt}%
\pgfpathmoveto{\pgfqpoint{0.955219in}{0.962827in}}%
\pgfpathlineto{\pgfqpoint{0.949186in}{0.978568in}}%
\pgfpathlineto{\pgfqpoint{0.943162in}{0.994286in}}%
\pgfpathlineto{\pgfqpoint{0.925167in}{0.989941in}}%
\pgfpathlineto{\pgfqpoint{0.907817in}{0.985142in}}%
\pgfpathlineto{\pgfqpoint{0.915309in}{0.969756in}}%
\pgfpathlineto{\pgfqpoint{0.922810in}{0.954352in}}%
\pgfpathlineto{\pgfqpoint{0.938715in}{0.958799in}}%
\pgfpathlineto{\pgfqpoint{0.955219in}{0.962827in}}%
\pgfpathclose%
\pgfusepath{fill}%
\end{pgfscope}%
\begin{pgfscope}%
\pgfpathrectangle{\pgfqpoint{0.407592in}{0.222904in}}{\pgfqpoint{1.411163in}{1.411163in}}%
\pgfusepath{clip}%
\pgfsetbuttcap%
\pgfsetroundjoin%
\definecolor{currentfill}{rgb}{0.749020,0.772549,0.792157}%
\pgfsetfillcolor{currentfill}%
\pgfsetfillopacity{0.160000}%
\pgfsetlinewidth{0.000000pt}%
\definecolor{currentstroke}{rgb}{0.000000,0.000000,0.000000}%
\pgfsetstrokecolor{currentstroke}%
\pgfsetdash{}{0pt}%
\pgfpathmoveto{\pgfqpoint{0.752268in}{1.182567in}}%
\pgfpathlineto{\pgfqpoint{0.743540in}{1.197340in}}%
\pgfpathlineto{\pgfqpoint{0.734820in}{1.212099in}}%
\pgfpathlineto{\pgfqpoint{0.709561in}{1.202822in}}%
\pgfpathlineto{\pgfqpoint{0.685736in}{1.192946in}}%
\pgfpathlineto{\pgfqpoint{0.695722in}{1.178567in}}%
\pgfpathlineto{\pgfqpoint{0.705714in}{1.164179in}}%
\pgfpathlineto{\pgfqpoint{0.728304in}{1.173660in}}%
\pgfpathlineto{\pgfqpoint{0.752268in}{1.182567in}}%
\pgfpathclose%
\pgfusepath{fill}%
\end{pgfscope}%
\begin{pgfscope}%
\pgfpathrectangle{\pgfqpoint{0.407592in}{0.222904in}}{\pgfqpoint{1.411163in}{1.411163in}}%
\pgfusepath{clip}%
\pgfsetbuttcap%
\pgfsetroundjoin%
\definecolor{currentfill}{rgb}{0.749020,0.772549,0.792157}%
\pgfsetfillcolor{currentfill}%
\pgfsetfillopacity{0.160000}%
\pgfsetlinewidth{0.000000pt}%
\definecolor{currentstroke}{rgb}{0.000000,0.000000,0.000000}%
\pgfsetstrokecolor{currentstroke}%
\pgfsetdash{}{0pt}%
\pgfpathmoveto{\pgfqpoint{1.473933in}{1.137239in}}%
\pgfpathlineto{\pgfqpoint{1.484458in}{1.151687in}}%
\pgfpathlineto{\pgfqpoint{1.494977in}{1.166126in}}%
\pgfpathlineto{\pgfqpoint{1.474273in}{1.175494in}}%
\pgfpathlineto{\pgfqpoint{1.452081in}{1.184283in}}%
\pgfpathlineto{\pgfqpoint{1.442732in}{1.169460in}}%
\pgfpathlineto{\pgfqpoint{1.433375in}{1.154622in}}%
\pgfpathlineto{\pgfqpoint{1.454364in}{1.146206in}}%
\pgfpathlineto{\pgfqpoint{1.473933in}{1.137239in}}%
\pgfpathclose%
\pgfusepath{fill}%
\end{pgfscope}%
\begin{pgfscope}%
\pgfpathrectangle{\pgfqpoint{0.407592in}{0.222904in}}{\pgfqpoint{1.411163in}{1.411163in}}%
\pgfusepath{clip}%
\pgfsetbuttcap%
\pgfsetroundjoin%
\definecolor{currentfill}{rgb}{0.749020,0.772549,0.792157}%
\pgfsetfillcolor{currentfill}%
\pgfsetfillopacity{0.160000}%
\pgfsetlinewidth{0.000000pt}%
\definecolor{currentstroke}{rgb}{0.000000,0.000000,0.000000}%
\pgfsetstrokecolor{currentstroke}%
\pgfsetdash{}{0pt}%
\pgfpathmoveto{\pgfqpoint{1.339286in}{1.005435in}}%
\pgfpathlineto{\pgfqpoint{1.348737in}{1.020420in}}%
\pgfpathlineto{\pgfqpoint{1.358178in}{1.035391in}}%
\pgfpathlineto{\pgfqpoint{1.340894in}{1.041754in}}%
\pgfpathlineto{\pgfqpoint{1.322616in}{1.047633in}}%
\pgfpathlineto{\pgfqpoint{1.314458in}{1.032289in}}%
\pgfpathlineto{\pgfqpoint{1.306290in}{1.016925in}}%
\pgfpathlineto{\pgfqpoint{1.323253in}{1.011407in}}%
\pgfpathlineto{\pgfqpoint{1.339286in}{1.005435in}}%
\pgfpathclose%
\pgfusepath{fill}%
\end{pgfscope}%
\begin{pgfscope}%
\pgfpathrectangle{\pgfqpoint{0.407592in}{0.222904in}}{\pgfqpoint{1.411163in}{1.411163in}}%
\pgfusepath{clip}%
\pgfsetbuttcap%
\pgfsetroundjoin%
\definecolor{currentfill}{rgb}{0.749020,0.772549,0.792157}%
\pgfsetfillcolor{currentfill}%
\pgfsetfillopacity{0.160000}%
\pgfsetlinewidth{0.000000pt}%
\definecolor{currentstroke}{rgb}{0.000000,0.000000,0.000000}%
\pgfsetstrokecolor{currentstroke}%
\pgfsetdash{}{0pt}%
\pgfpathmoveto{\pgfqpoint{1.215338in}{0.900262in}}%
\pgfpathlineto{\pgfqpoint{1.222189in}{0.916128in}}%
\pgfpathlineto{\pgfqpoint{1.229030in}{0.931970in}}%
\pgfpathlineto{\pgfqpoint{1.214141in}{0.935194in}}%
\pgfpathlineto{\pgfqpoint{1.198751in}{0.938005in}}%
\pgfpathlineto{\pgfqpoint{1.193412in}{0.921894in}}%
\pgfpathlineto{\pgfqpoint{1.188064in}{0.905755in}}%
\pgfpathlineto{\pgfqpoint{1.201929in}{0.903196in}}%
\pgfpathlineto{\pgfqpoint{1.215338in}{0.900262in}}%
\pgfpathclose%
\pgfusepath{fill}%
\end{pgfscope}%
\begin{pgfscope}%
\pgfpathrectangle{\pgfqpoint{0.407592in}{0.222904in}}{\pgfqpoint{1.411163in}{1.411163in}}%
\pgfusepath{clip}%
\pgfsetbuttcap%
\pgfsetroundjoin%
\definecolor{currentfill}{rgb}{0.749020,0.772549,0.792157}%
\pgfsetfillcolor{currentfill}%
\pgfsetfillopacity{0.160000}%
\pgfsetlinewidth{0.000000pt}%
\definecolor{currentstroke}{rgb}{0.000000,0.000000,0.000000}%
\pgfsetstrokecolor{currentstroke}%
\pgfsetdash{}{0pt}%
\pgfpathmoveto{\pgfqpoint{1.007742in}{0.905272in}}%
\pgfpathlineto{\pgfqpoint{1.003232in}{0.921387in}}%
\pgfpathlineto{\pgfqpoint{0.998729in}{0.937475in}}%
\pgfpathlineto{\pgfqpoint{0.982808in}{0.934580in}}%
\pgfpathlineto{\pgfqpoint{0.967310in}{0.931275in}}%
\pgfpathlineto{\pgfqpoint{0.973370in}{0.915464in}}%
\pgfpathlineto{\pgfqpoint{0.979438in}{0.899629in}}%
\pgfpathlineto{\pgfqpoint{0.993397in}{0.902637in}}%
\pgfpathlineto{\pgfqpoint{1.007742in}{0.905272in}}%
\pgfpathclose%
\pgfusepath{fill}%
\end{pgfscope}%
\begin{pgfscope}%
\pgfpathrectangle{\pgfqpoint{0.407592in}{0.222904in}}{\pgfqpoint{1.411163in}{1.411163in}}%
\pgfusepath{clip}%
\pgfsetbuttcap%
\pgfsetroundjoin%
\definecolor{currentfill}{rgb}{0.749020,0.772549,0.792157}%
\pgfsetfillcolor{currentfill}%
\pgfsetfillopacity{0.160000}%
\pgfsetlinewidth{0.000000pt}%
\definecolor{currentstroke}{rgb}{0.000000,0.000000,0.000000}%
\pgfsetstrokecolor{currentstroke}%
\pgfsetdash{}{0pt}%
\pgfpathmoveto{\pgfqpoint{1.177341in}{0.873396in}}%
\pgfpathlineto{\pgfqpoint{1.182707in}{0.889589in}}%
\pgfpathlineto{\pgfqpoint{1.188064in}{0.905755in}}%
\pgfpathlineto{\pgfqpoint{1.173801in}{0.907930in}}%
\pgfpathlineto{\pgfqpoint{1.159198in}{0.909713in}}%
\pgfpathlineto{\pgfqpoint{1.155426in}{0.893350in}}%
\pgfpathlineto{\pgfqpoint{1.151647in}{0.876956in}}%
\pgfpathlineto{\pgfqpoint{1.164646in}{0.875352in}}%
\pgfpathlineto{\pgfqpoint{1.177341in}{0.873396in}}%
\pgfpathclose%
\pgfusepath{fill}%
\end{pgfscope}%
\begin{pgfscope}%
\pgfpathrectangle{\pgfqpoint{0.407592in}{0.222904in}}{\pgfqpoint{1.411163in}{1.411163in}}%
\pgfusepath{clip}%
\pgfsetbuttcap%
\pgfsetroundjoin%
\definecolor{currentfill}{rgb}{0.749020,0.772549,0.792157}%
\pgfsetfillcolor{currentfill}%
\pgfsetfillopacity{0.160000}%
\pgfsetlinewidth{0.000000pt}%
\definecolor{currentstroke}{rgb}{0.000000,0.000000,0.000000}%
\pgfsetstrokecolor{currentstroke}%
\pgfsetdash{}{0pt}%
\pgfpathmoveto{\pgfqpoint{0.892861in}{1.015856in}}%
\pgfpathlineto{\pgfqpoint{0.885397in}{1.031184in}}%
\pgfpathlineto{\pgfqpoint{0.877943in}{1.046493in}}%
\pgfpathlineto{\pgfqpoint{0.858525in}{1.040516in}}%
\pgfpathlineto{\pgfqpoint{0.840008in}{1.034058in}}%
\pgfpathlineto{\pgfqpoint{0.848829in}{1.019128in}}%
\pgfpathlineto{\pgfqpoint{0.857658in}{1.004184in}}%
\pgfpathlineto{\pgfqpoint{0.874837in}{1.010245in}}%
\pgfpathlineto{\pgfqpoint{0.892861in}{1.015856in}}%
\pgfpathclose%
\pgfusepath{fill}%
\end{pgfscope}%
\begin{pgfscope}%
\pgfpathrectangle{\pgfqpoint{0.407592in}{0.222904in}}{\pgfqpoint{1.411163in}{1.411163in}}%
\pgfusepath{clip}%
\pgfsetbuttcap%
\pgfsetroundjoin%
\definecolor{currentfill}{rgb}{0.749020,0.772549,0.792157}%
\pgfsetfillcolor{currentfill}%
\pgfsetfillopacity{0.160000}%
\pgfsetlinewidth{0.000000pt}%
\definecolor{currentstroke}{rgb}{0.000000,0.000000,0.000000}%
\pgfsetstrokecolor{currentstroke}%
\pgfsetdash{}{0pt}%
\pgfpathmoveto{\pgfqpoint{1.043144in}{0.876663in}}%
\pgfpathlineto{\pgfqpoint{1.040246in}{0.893040in}}%
\pgfpathlineto{\pgfqpoint{1.037354in}{0.909387in}}%
\pgfpathlineto{\pgfqpoint{1.022414in}{0.907525in}}%
\pgfpathlineto{\pgfqpoint{1.007742in}{0.905272in}}%
\pgfpathlineto{\pgfqpoint{1.012259in}{0.889131in}}%
\pgfpathlineto{\pgfqpoint{1.016784in}{0.872962in}}%
\pgfpathlineto{\pgfqpoint{1.029844in}{0.874988in}}%
\pgfpathlineto{\pgfqpoint{1.043144in}{0.876663in}}%
\pgfpathclose%
\pgfusepath{fill}%
\end{pgfscope}%
\begin{pgfscope}%
\pgfpathrectangle{\pgfqpoint{0.407592in}{0.222904in}}{\pgfqpoint{1.411163in}{1.411163in}}%
\pgfusepath{clip}%
\pgfsetbuttcap%
\pgfsetroundjoin%
\definecolor{currentfill}{rgb}{0.749020,0.772549,0.792157}%
\pgfsetfillcolor{currentfill}%
\pgfsetfillopacity{0.160000}%
\pgfsetlinewidth{0.000000pt}%
\definecolor{currentstroke}{rgb}{0.000000,0.000000,0.000000}%
\pgfsetstrokecolor{currentstroke}%
\pgfsetdash{}{0pt}%
\pgfpathmoveto{\pgfqpoint{0.769749in}{1.152979in}}%
\pgfpathlineto{\pgfqpoint{0.761004in}{1.167780in}}%
\pgfpathlineto{\pgfqpoint{0.752268in}{1.182567in}}%
\pgfpathlineto{\pgfqpoint{0.728304in}{1.173660in}}%
\pgfpathlineto{\pgfqpoint{0.705714in}{1.164179in}}%
\pgfpathlineto{\pgfqpoint{0.715713in}{1.149782in}}%
\pgfpathlineto{\pgfqpoint{0.725719in}{1.135376in}}%
\pgfpathlineto{\pgfqpoint{0.747077in}{1.144450in}}%
\pgfpathlineto{\pgfqpoint{0.769749in}{1.152979in}}%
\pgfpathclose%
\pgfusepath{fill}%
\end{pgfscope}%
\begin{pgfscope}%
\pgfpathrectangle{\pgfqpoint{0.407592in}{0.222904in}}{\pgfqpoint{1.411163in}{1.411163in}}%
\pgfusepath{clip}%
\pgfsetbuttcap%
\pgfsetroundjoin%
\definecolor{currentfill}{rgb}{0.749020,0.772549,0.792157}%
\pgfsetfillcolor{currentfill}%
\pgfsetfillopacity{0.160000}%
\pgfsetlinewidth{0.000000pt}%
\definecolor{currentstroke}{rgb}{0.000000,0.000000,0.000000}%
\pgfsetstrokecolor{currentstroke}%
\pgfsetdash{}{0pt}%
\pgfpathmoveto{\pgfqpoint{1.257064in}{0.924335in}}%
\pgfpathlineto{\pgfqpoint{1.265294in}{0.939816in}}%
\pgfpathlineto{\pgfqpoint{1.273514in}{0.955277in}}%
\pgfpathlineto{\pgfqpoint{1.258436in}{0.959641in}}%
\pgfpathlineto{\pgfqpoint{1.242681in}{0.963583in}}%
\pgfpathlineto{\pgfqpoint{1.235861in}{0.947789in}}%
\pgfpathlineto{\pgfqpoint{1.229030in}{0.931970in}}%
\pgfpathlineto{\pgfqpoint{1.243358in}{0.928346in}}%
\pgfpathlineto{\pgfqpoint{1.257064in}{0.924335in}}%
\pgfpathclose%
\pgfusepath{fill}%
\end{pgfscope}%
\begin{pgfscope}%
\pgfpathrectangle{\pgfqpoint{0.407592in}{0.222904in}}{\pgfqpoint{1.411163in}{1.411163in}}%
\pgfusepath{clip}%
\pgfsetbuttcap%
\pgfsetroundjoin%
\definecolor{currentfill}{rgb}{0.749020,0.772549,0.792157}%
\pgfsetfillcolor{currentfill}%
\pgfsetfillopacity{0.160000}%
\pgfsetlinewidth{0.000000pt}%
\definecolor{currentstroke}{rgb}{0.000000,0.000000,0.000000}%
\pgfsetstrokecolor{currentstroke}%
\pgfsetdash{}{0pt}%
\pgfpathmoveto{\pgfqpoint{1.452860in}{1.108313in}}%
\pgfpathlineto{\pgfqpoint{1.463400in}{1.122781in}}%
\pgfpathlineto{\pgfqpoint{1.473933in}{1.137239in}}%
\pgfpathlineto{\pgfqpoint{1.454364in}{1.146206in}}%
\pgfpathlineto{\pgfqpoint{1.433375in}{1.154622in}}%
\pgfpathlineto{\pgfqpoint{1.424008in}{1.139770in}}%
\pgfpathlineto{\pgfqpoint{1.414631in}{1.124902in}}%
\pgfpathlineto{\pgfqpoint{1.434422in}{1.116870in}}%
\pgfpathlineto{\pgfqpoint{1.452860in}{1.108313in}}%
\pgfpathclose%
\pgfusepath{fill}%
\end{pgfscope}%
\begin{pgfscope}%
\pgfpathrectangle{\pgfqpoint{0.407592in}{0.222904in}}{\pgfqpoint{1.411163in}{1.411163in}}%
\pgfusepath{clip}%
\pgfsetbuttcap%
\pgfsetroundjoin%
\definecolor{currentfill}{rgb}{0.749020,0.772549,0.792157}%
\pgfsetfillcolor{currentfill}%
\pgfsetfillopacity{0.160000}%
\pgfsetlinewidth{0.000000pt}%
\definecolor{currentstroke}{rgb}{0.000000,0.000000,0.000000}%
\pgfsetstrokecolor{currentstroke}%
\pgfsetdash{}{0pt}%
\pgfpathmoveto{\pgfqpoint{1.120662in}{0.845962in}}%
\pgfpathlineto{\pgfqpoint{1.122805in}{0.862538in}}%
\pgfpathlineto{\pgfqpoint{1.124944in}{0.879083in}}%
\pgfpathlineto{\pgfqpoint{1.111349in}{0.879599in}}%
\pgfpathlineto{\pgfqpoint{1.097665in}{0.879746in}}%
\pgfpathlineto{\pgfqpoint{1.097207in}{0.863164in}}%
\pgfpathlineto{\pgfqpoint{1.096748in}{0.846549in}}%
\pgfpathlineto{\pgfqpoint{1.108745in}{0.846419in}}%
\pgfpathlineto{\pgfqpoint{1.120662in}{0.845962in}}%
\pgfpathclose%
\pgfusepath{fill}%
\end{pgfscope}%
\begin{pgfscope}%
\pgfpathrectangle{\pgfqpoint{0.407592in}{0.222904in}}{\pgfqpoint{1.411163in}{1.411163in}}%
\pgfusepath{clip}%
\pgfsetbuttcap%
\pgfsetroundjoin%
\definecolor{currentfill}{rgb}{0.749020,0.772549,0.792157}%
\pgfsetfillcolor{currentfill}%
\pgfsetfillopacity{0.160000}%
\pgfsetlinewidth{0.000000pt}%
\definecolor{currentstroke}{rgb}{0.000000,0.000000,0.000000}%
\pgfsetstrokecolor{currentstroke}%
\pgfsetdash{}{0pt}%
\pgfpathmoveto{\pgfqpoint{1.096748in}{0.846549in}}%
\pgfpathlineto{\pgfqpoint{1.097207in}{0.863164in}}%
\pgfpathlineto{\pgfqpoint{1.097665in}{0.879746in}}%
\pgfpathlineto{\pgfqpoint{1.083947in}{0.879525in}}%
\pgfpathlineto{\pgfqpoint{1.070251in}{0.878935in}}%
\pgfpathlineto{\pgfqpoint{1.071482in}{0.862399in}}%
\pgfpathlineto{\pgfqpoint{1.072715in}{0.845831in}}%
\pgfpathlineto{\pgfqpoint{1.084722in}{0.846353in}}%
\pgfpathlineto{\pgfqpoint{1.096748in}{0.846549in}}%
\pgfpathclose%
\pgfusepath{fill}%
\end{pgfscope}%
\begin{pgfscope}%
\pgfpathrectangle{\pgfqpoint{0.407592in}{0.222904in}}{\pgfqpoint{1.411163in}{1.411163in}}%
\pgfusepath{clip}%
\pgfsetbuttcap%
\pgfsetroundjoin%
\definecolor{currentfill}{rgb}{0.749020,0.772549,0.792157}%
\pgfsetfillcolor{currentfill}%
\pgfsetfillopacity{0.160000}%
\pgfsetlinewidth{0.000000pt}%
\definecolor{currentstroke}{rgb}{0.000000,0.000000,0.000000}%
\pgfsetstrokecolor{currentstroke}%
\pgfsetdash{}{0pt}%
\pgfpathmoveto{\pgfqpoint{0.967310in}{0.931275in}}%
\pgfpathlineto{\pgfqpoint{0.961260in}{0.947063in}}%
\pgfpathlineto{\pgfqpoint{0.955219in}{0.962827in}}%
\pgfpathlineto{\pgfqpoint{0.938715in}{0.958799in}}%
\pgfpathlineto{\pgfqpoint{0.922810in}{0.954352in}}%
\pgfpathlineto{\pgfqpoint{0.930321in}{0.938928in}}%
\pgfpathlineto{\pgfqpoint{0.937841in}{0.923485in}}%
\pgfpathlineto{\pgfqpoint{0.952301in}{0.927572in}}%
\pgfpathlineto{\pgfqpoint{0.967310in}{0.931275in}}%
\pgfpathclose%
\pgfusepath{fill}%
\end{pgfscope}%
\begin{pgfscope}%
\pgfpathrectangle{\pgfqpoint{0.407592in}{0.222904in}}{\pgfqpoint{1.411163in}{1.411163in}}%
\pgfusepath{clip}%
\pgfsetbuttcap%
\pgfsetroundjoin%
\definecolor{currentfill}{rgb}{0.749020,0.772549,0.792157}%
\pgfsetfillcolor{currentfill}%
\pgfsetfillopacity{0.160000}%
\pgfsetlinewidth{0.000000pt}%
\definecolor{currentstroke}{rgb}{0.000000,0.000000,0.000000}%
\pgfsetstrokecolor{currentstroke}%
\pgfsetdash{}{0pt}%
\pgfpathmoveto{\pgfqpoint{1.320356in}{0.975420in}}%
\pgfpathlineto{\pgfqpoint{1.329826in}{0.990435in}}%
\pgfpathlineto{\pgfqpoint{1.339286in}{1.005435in}}%
\pgfpathlineto{\pgfqpoint{1.323253in}{1.011407in}}%
\pgfpathlineto{\pgfqpoint{1.306290in}{1.016925in}}%
\pgfpathlineto{\pgfqpoint{1.298111in}{1.001543in}}%
\pgfpathlineto{\pgfqpoint{1.289923in}{0.986140in}}%
\pgfpathlineto{\pgfqpoint{1.305573in}{0.980991in}}%
\pgfpathlineto{\pgfqpoint{1.320356in}{0.975420in}}%
\pgfpathclose%
\pgfusepath{fill}%
\end{pgfscope}%
\begin{pgfscope}%
\pgfpathrectangle{\pgfqpoint{0.407592in}{0.222904in}}{\pgfqpoint{1.411163in}{1.411163in}}%
\pgfusepath{clip}%
\pgfsetbuttcap%
\pgfsetroundjoin%
\definecolor{currentfill}{rgb}{0.749020,0.772549,0.792157}%
\pgfsetfillcolor{currentfill}%
\pgfsetfillopacity{0.160000}%
\pgfsetlinewidth{0.000000pt}%
\definecolor{currentstroke}{rgb}{0.000000,0.000000,0.000000}%
\pgfsetstrokecolor{currentstroke}%
\pgfsetdash{}{0pt}%
\pgfpathmoveto{\pgfqpoint{1.738704in}{1.397027in}}%
\pgfpathlineto{\pgfqpoint{1.750178in}{1.410942in}}%
\pgfpathlineto{\pgfqpoint{1.723009in}{1.425155in}}%
\pgfpathlineto{\pgfqpoint{1.693489in}{1.438622in}}%
\pgfpathlineto{\pgfqpoint{1.683104in}{1.424367in}}%
\pgfpathlineto{\pgfqpoint{1.712063in}{1.411064in}}%
\pgfusepath{fill}%
\end{pgfscope}%
\begin{pgfscope}%
\pgfpathrectangle{\pgfqpoint{0.407592in}{0.222904in}}{\pgfqpoint{1.411163in}{1.411163in}}%
\pgfusepath{clip}%
\pgfsetbuttcap%
\pgfsetroundjoin%
\definecolor{currentfill}{rgb}{0.749020,0.772549,0.792157}%
\pgfsetfillcolor{currentfill}%
\pgfsetfillopacity{0.160000}%
\pgfsetlinewidth{0.000000pt}%
\definecolor{currentstroke}{rgb}{0.000000,0.000000,0.000000}%
\pgfsetstrokecolor{currentstroke}%
\pgfsetdash{}{0pt}%
\pgfpathmoveto{\pgfqpoint{0.787263in}{1.123334in}}%
\pgfpathlineto{\pgfqpoint{0.778502in}{1.138164in}}%
\pgfpathlineto{\pgfqpoint{0.769749in}{1.152979in}}%
\pgfpathlineto{\pgfqpoint{0.747077in}{1.144450in}}%
\pgfpathlineto{\pgfqpoint{0.725719in}{1.135376in}}%
\pgfpathlineto{\pgfqpoint{0.735730in}{1.120961in}}%
\pgfpathlineto{\pgfqpoint{0.745749in}{1.106536in}}%
\pgfpathlineto{\pgfqpoint{0.765880in}{1.115194in}}%
\pgfpathlineto{\pgfqpoint{0.787263in}{1.123334in}}%
\pgfpathclose%
\pgfusepath{fill}%
\end{pgfscope}%
\begin{pgfscope}%
\pgfpathrectangle{\pgfqpoint{0.407592in}{0.222904in}}{\pgfqpoint{1.411163in}{1.411163in}}%
\pgfusepath{clip}%
\pgfsetbuttcap%
\pgfsetroundjoin%
\definecolor{currentfill}{rgb}{0.749020,0.772549,0.792157}%
\pgfsetfillcolor{currentfill}%
\pgfsetfillopacity{0.160000}%
\pgfsetlinewidth{0.000000pt}%
\definecolor{currentstroke}{rgb}{0.000000,0.000000,0.000000}%
\pgfsetstrokecolor{currentstroke}%
\pgfsetdash{}{0pt}%
\pgfpathmoveto{\pgfqpoint{1.144068in}{0.844079in}}%
\pgfpathlineto{\pgfqpoint{1.147861in}{0.860532in}}%
\pgfpathlineto{\pgfqpoint{1.151647in}{0.876956in}}%
\pgfpathlineto{\pgfqpoint{1.138395in}{0.878201in}}%
\pgfpathlineto{\pgfqpoint{1.124944in}{0.879083in}}%
\pgfpathlineto{\pgfqpoint{1.122805in}{0.862538in}}%
\pgfpathlineto{\pgfqpoint{1.120662in}{0.845962in}}%
\pgfpathlineto{\pgfqpoint{1.132453in}{0.845181in}}%
\pgfpathlineto{\pgfqpoint{1.144068in}{0.844079in}}%
\pgfpathclose%
\pgfusepath{fill}%
\end{pgfscope}%
\begin{pgfscope}%
\pgfpathrectangle{\pgfqpoint{0.407592in}{0.222904in}}{\pgfqpoint{1.411163in}{1.411163in}}%
\pgfusepath{clip}%
\pgfsetbuttcap%
\pgfsetroundjoin%
\definecolor{currentfill}{rgb}{0.749020,0.772549,0.792157}%
\pgfsetfillcolor{currentfill}%
\pgfsetfillopacity{0.160000}%
\pgfsetlinewidth{0.000000pt}%
\definecolor{currentstroke}{rgb}{0.000000,0.000000,0.000000}%
\pgfsetstrokecolor{currentstroke}%
\pgfsetdash{}{0pt}%
\pgfpathmoveto{\pgfqpoint{1.072715in}{0.845831in}}%
\pgfpathlineto{\pgfqpoint{1.071482in}{0.862399in}}%
\pgfpathlineto{\pgfqpoint{1.070251in}{0.878935in}}%
\pgfpathlineto{\pgfqpoint{1.056631in}{0.877980in}}%
\pgfpathlineto{\pgfqpoint{1.043144in}{0.876663in}}%
\pgfpathlineto{\pgfqpoint{1.046047in}{0.860256in}}%
\pgfpathlineto{\pgfqpoint{1.048955in}{0.843819in}}%
\pgfpathlineto{\pgfqpoint{1.060776in}{0.844986in}}%
\pgfpathlineto{\pgfqpoint{1.072715in}{0.845831in}}%
\pgfpathclose%
\pgfusepath{fill}%
\end{pgfscope}%
\begin{pgfscope}%
\pgfpathrectangle{\pgfqpoint{0.407592in}{0.222904in}}{\pgfqpoint{1.411163in}{1.411163in}}%
\pgfusepath{clip}%
\pgfsetbuttcap%
\pgfsetroundjoin%
\definecolor{currentfill}{rgb}{0.749020,0.772549,0.792157}%
\pgfsetfillcolor{currentfill}%
\pgfsetfillopacity{0.160000}%
\pgfsetlinewidth{0.000000pt}%
\definecolor{currentstroke}{rgb}{0.000000,0.000000,0.000000}%
\pgfsetstrokecolor{currentstroke}%
\pgfsetdash{}{0pt}%
\pgfpathmoveto{\pgfqpoint{0.907817in}{0.985142in}}%
\pgfpathlineto{\pgfqpoint{0.900335in}{1.000508in}}%
\pgfpathlineto{\pgfqpoint{0.892861in}{1.015856in}}%
\pgfpathlineto{\pgfqpoint{0.874837in}{1.010245in}}%
\pgfpathlineto{\pgfqpoint{0.857658in}{1.004184in}}%
\pgfpathlineto{\pgfqpoint{0.866495in}{0.989226in}}%
\pgfpathlineto{\pgfqpoint{0.875341in}{0.974253in}}%
\pgfpathlineto{\pgfqpoint{0.891185in}{0.979906in}}%
\pgfpathlineto{\pgfqpoint{0.907817in}{0.985142in}}%
\pgfpathclose%
\pgfusepath{fill}%
\end{pgfscope}%
\begin{pgfscope}%
\pgfpathrectangle{\pgfqpoint{0.407592in}{0.222904in}}{\pgfqpoint{1.411163in}{1.411163in}}%
\pgfusepath{clip}%
\pgfsetbuttcap%
\pgfsetroundjoin%
\definecolor{currentfill}{rgb}{0.749020,0.772549,0.792157}%
\pgfsetfillcolor{currentfill}%
\pgfsetfillopacity{0.160000}%
\pgfsetlinewidth{0.000000pt}%
\definecolor{currentstroke}{rgb}{0.000000,0.000000,0.000000}%
\pgfsetstrokecolor{currentstroke}%
\pgfsetdash{}{0pt}%
\pgfpathmoveto{\pgfqpoint{1.431759in}{1.079348in}}%
\pgfpathlineto{\pgfqpoint{1.442313in}{1.093835in}}%
\pgfpathlineto{\pgfqpoint{1.452860in}{1.108313in}}%
\pgfpathlineto{\pgfqpoint{1.434422in}{1.116870in}}%
\pgfpathlineto{\pgfqpoint{1.414631in}{1.124902in}}%
\pgfpathlineto{\pgfqpoint{1.405246in}{1.110021in}}%
\pgfpathlineto{\pgfqpoint{1.395851in}{1.095124in}}%
\pgfpathlineto{\pgfqpoint{1.414446in}{1.087484in}}%
\pgfpathlineto{\pgfqpoint{1.431759in}{1.079348in}}%
\pgfpathclose%
\pgfusepath{fill}%
\end{pgfscope}%
\begin{pgfscope}%
\pgfpathrectangle{\pgfqpoint{0.407592in}{0.222904in}}{\pgfqpoint{1.411163in}{1.411163in}}%
\pgfusepath{clip}%
\pgfsetbuttcap%
\pgfsetroundjoin%
\definecolor{currentfill}{rgb}{0.749020,0.772549,0.792157}%
\pgfsetfillcolor{currentfill}%
\pgfsetfillopacity{0.160000}%
\pgfsetlinewidth{0.000000pt}%
\definecolor{currentstroke}{rgb}{0.000000,0.000000,0.000000}%
\pgfsetstrokecolor{currentstroke}%
\pgfsetdash{}{0pt}%
\pgfpathmoveto{\pgfqpoint{1.715746in}{1.369183in}}%
\pgfpathlineto{\pgfqpoint{1.727227in}{1.383107in}}%
\pgfpathlineto{\pgfqpoint{1.738704in}{1.397027in}}%
\pgfpathlineto{\pgfqpoint{1.712063in}{1.411064in}}%
\pgfpathlineto{\pgfqpoint{1.683104in}{1.424367in}}%
\pgfpathlineto{\pgfqpoint{1.672712in}{1.410102in}}%
\pgfpathlineto{\pgfqpoint{1.662313in}{1.395827in}}%
\pgfpathlineto{\pgfqpoint{1.690155in}{1.382860in}}%
\pgfpathlineto{\pgfqpoint{1.715746in}{1.369183in}}%
\pgfpathclose%
\pgfusepath{fill}%
\end{pgfscope}%
\begin{pgfscope}%
\pgfpathrectangle{\pgfqpoint{0.407592in}{0.222904in}}{\pgfqpoint{1.411163in}{1.411163in}}%
\pgfusepath{clip}%
\pgfsetbuttcap%
\pgfsetroundjoin%
\definecolor{currentfill}{rgb}{0.749020,0.772549,0.792157}%
\pgfsetfillcolor{currentfill}%
\pgfsetfillopacity{0.160000}%
\pgfsetlinewidth{0.000000pt}%
\definecolor{currentstroke}{rgb}{0.000000,0.000000,0.000000}%
\pgfsetstrokecolor{currentstroke}%
\pgfsetdash{}{0pt}%
\pgfpathmoveto{\pgfqpoint{1.201605in}{0.868458in}}%
\pgfpathlineto{\pgfqpoint{1.208477in}{0.884372in}}%
\pgfpathlineto{\pgfqpoint{1.215338in}{0.900262in}}%
\pgfpathlineto{\pgfqpoint{1.201929in}{0.903196in}}%
\pgfpathlineto{\pgfqpoint{1.188064in}{0.905755in}}%
\pgfpathlineto{\pgfqpoint{1.182707in}{0.889589in}}%
\pgfpathlineto{\pgfqpoint{1.177341in}{0.873396in}}%
\pgfpathlineto{\pgfqpoint{1.189677in}{0.871095in}}%
\pgfpathlineto{\pgfqpoint{1.201605in}{0.868458in}}%
\pgfpathclose%
\pgfusepath{fill}%
\end{pgfscope}%
\begin{pgfscope}%
\pgfpathrectangle{\pgfqpoint{0.407592in}{0.222904in}}{\pgfqpoint{1.411163in}{1.411163in}}%
\pgfusepath{clip}%
\pgfsetbuttcap%
\pgfsetroundjoin%
\definecolor{currentfill}{rgb}{0.749020,0.772549,0.792157}%
\pgfsetfillcolor{currentfill}%
\pgfsetfillopacity{0.160000}%
\pgfsetlinewidth{0.000000pt}%
\definecolor{currentstroke}{rgb}{0.000000,0.000000,0.000000}%
\pgfsetstrokecolor{currentstroke}%
\pgfsetdash{}{0pt}%
\pgfpathmoveto{\pgfqpoint{1.016784in}{0.872962in}}%
\pgfpathlineto{\pgfqpoint{1.012259in}{0.889131in}}%
\pgfpathlineto{\pgfqpoint{1.007742in}{0.905272in}}%
\pgfpathlineto{\pgfqpoint{0.993397in}{0.902637in}}%
\pgfpathlineto{\pgfqpoint{0.979438in}{0.899629in}}%
\pgfpathlineto{\pgfqpoint{0.985516in}{0.883771in}}%
\pgfpathlineto{\pgfqpoint{0.991602in}{0.867889in}}%
\pgfpathlineto{\pgfqpoint{1.004020in}{0.870593in}}%
\pgfpathlineto{\pgfqpoint{1.016784in}{0.872962in}}%
\pgfpathclose%
\pgfusepath{fill}%
\end{pgfscope}%
\begin{pgfscope}%
\pgfpathrectangle{\pgfqpoint{0.407592in}{0.222904in}}{\pgfqpoint{1.411163in}{1.411163in}}%
\pgfusepath{clip}%
\pgfsetbuttcap%
\pgfsetroundjoin%
\definecolor{currentfill}{rgb}{0.749020,0.772549,0.792157}%
\pgfsetfillcolor{currentfill}%
\pgfsetfillopacity{0.160000}%
\pgfsetlinewidth{0.000000pt}%
\definecolor{currentstroke}{rgb}{0.000000,0.000000,0.000000}%
\pgfsetstrokecolor{currentstroke}%
\pgfsetdash{}{0pt}%
\pgfpathmoveto{\pgfqpoint{0.804811in}{1.093633in}}%
\pgfpathlineto{\pgfqpoint{0.796033in}{1.108491in}}%
\pgfpathlineto{\pgfqpoint{0.787263in}{1.123334in}}%
\pgfpathlineto{\pgfqpoint{0.765880in}{1.115194in}}%
\pgfpathlineto{\pgfqpoint{0.745749in}{1.106536in}}%
\pgfpathlineto{\pgfqpoint{0.755773in}{1.092102in}}%
\pgfpathlineto{\pgfqpoint{0.765804in}{1.077659in}}%
\pgfpathlineto{\pgfqpoint{0.784714in}{1.085891in}}%
\pgfpathlineto{\pgfqpoint{0.804811in}{1.093633in}}%
\pgfpathclose%
\pgfusepath{fill}%
\end{pgfscope}%
\begin{pgfscope}%
\pgfpathrectangle{\pgfqpoint{0.407592in}{0.222904in}}{\pgfqpoint{1.411163in}{1.411163in}}%
\pgfusepath{clip}%
\pgfsetbuttcap%
\pgfsetroundjoin%
\definecolor{currentfill}{rgb}{0.749020,0.772549,0.792157}%
\pgfsetfillcolor{currentfill}%
\pgfsetfillopacity{0.160000}%
\pgfsetlinewidth{0.000000pt}%
\definecolor{currentstroke}{rgb}{0.000000,0.000000,0.000000}%
\pgfsetstrokecolor{currentstroke}%
\pgfsetdash{}{0pt}%
\pgfpathmoveto{\pgfqpoint{1.166581in}{0.840928in}}%
\pgfpathlineto{\pgfqpoint{1.171965in}{0.857176in}}%
\pgfpathlineto{\pgfqpoint{1.177341in}{0.873396in}}%
\pgfpathlineto{\pgfqpoint{1.164646in}{0.875352in}}%
\pgfpathlineto{\pgfqpoint{1.151647in}{0.876956in}}%
\pgfpathlineto{\pgfqpoint{1.147861in}{0.860532in}}%
\pgfpathlineto{\pgfqpoint{1.144068in}{0.844079in}}%
\pgfpathlineto{\pgfqpoint{1.155459in}{0.842659in}}%
\pgfpathlineto{\pgfqpoint{1.166581in}{0.840928in}}%
\pgfpathclose%
\pgfusepath{fill}%
\end{pgfscope}%
\begin{pgfscope}%
\pgfpathrectangle{\pgfqpoint{0.407592in}{0.222904in}}{\pgfqpoint{1.411163in}{1.411163in}}%
\pgfusepath{clip}%
\pgfsetbuttcap%
\pgfsetroundjoin%
\definecolor{currentfill}{rgb}{0.749020,0.772549,0.792157}%
\pgfsetfillcolor{currentfill}%
\pgfsetfillopacity{0.160000}%
\pgfsetlinewidth{0.000000pt}%
\definecolor{currentstroke}{rgb}{0.000000,0.000000,0.000000}%
\pgfsetstrokecolor{currentstroke}%
\pgfsetdash{}{0pt}%
\pgfpathmoveto{\pgfqpoint{1.240571in}{0.893314in}}%
\pgfpathlineto{\pgfqpoint{1.248823in}{0.908834in}}%
\pgfpathlineto{\pgfqpoint{1.257064in}{0.924335in}}%
\pgfpathlineto{\pgfqpoint{1.243358in}{0.928346in}}%
\pgfpathlineto{\pgfqpoint{1.229030in}{0.931970in}}%
\pgfpathlineto{\pgfqpoint{1.222189in}{0.916128in}}%
\pgfpathlineto{\pgfqpoint{1.215338in}{0.900262in}}%
\pgfpathlineto{\pgfqpoint{1.228237in}{0.896964in}}%
\pgfpathlineto{\pgfqpoint{1.240571in}{0.893314in}}%
\pgfpathclose%
\pgfusepath{fill}%
\end{pgfscope}%
\begin{pgfscope}%
\pgfpathrectangle{\pgfqpoint{0.407592in}{0.222904in}}{\pgfqpoint{1.411163in}{1.411163in}}%
\pgfusepath{clip}%
\pgfsetbuttcap%
\pgfsetroundjoin%
\definecolor{currentfill}{rgb}{0.749020,0.772549,0.792157}%
\pgfsetfillcolor{currentfill}%
\pgfsetfillopacity{0.160000}%
\pgfsetlinewidth{0.000000pt}%
\definecolor{currentstroke}{rgb}{0.000000,0.000000,0.000000}%
\pgfsetstrokecolor{currentstroke}%
\pgfsetdash{}{0pt}%
\pgfpathmoveto{\pgfqpoint{1.692774in}{1.341320in}}%
\pgfpathlineto{\pgfqpoint{1.704262in}{1.355254in}}%
\pgfpathlineto{\pgfqpoint{1.715746in}{1.369183in}}%
\pgfpathlineto{\pgfqpoint{1.690155in}{1.382860in}}%
\pgfpathlineto{\pgfqpoint{1.662313in}{1.395827in}}%
\pgfpathlineto{\pgfqpoint{1.651907in}{1.381543in}}%
\pgfpathlineto{\pgfqpoint{1.641494in}{1.367249in}}%
\pgfpathlineto{\pgfqpoint{1.668225in}{1.354628in}}%
\pgfpathlineto{\pgfqpoint{1.692774in}{1.341320in}}%
\pgfpathclose%
\pgfusepath{fill}%
\end{pgfscope}%
\begin{pgfscope}%
\pgfpathrectangle{\pgfqpoint{0.407592in}{0.222904in}}{\pgfqpoint{1.411163in}{1.411163in}}%
\pgfusepath{clip}%
\pgfsetbuttcap%
\pgfsetroundjoin%
\definecolor{currentfill}{rgb}{0.749020,0.772549,0.792157}%
\pgfsetfillcolor{currentfill}%
\pgfsetfillopacity{0.160000}%
\pgfsetlinewidth{0.000000pt}%
\definecolor{currentstroke}{rgb}{0.000000,0.000000,0.000000}%
\pgfsetstrokecolor{currentstroke}%
\pgfsetdash{}{0pt}%
\pgfpathmoveto{\pgfqpoint{1.048955in}{0.843819in}}%
\pgfpathlineto{\pgfqpoint{1.046047in}{0.860256in}}%
\pgfpathlineto{\pgfqpoint{1.043144in}{0.876663in}}%
\pgfpathlineto{\pgfqpoint{1.029844in}{0.874988in}}%
\pgfpathlineto{\pgfqpoint{1.016784in}{0.872962in}}%
\pgfpathlineto{\pgfqpoint{1.021317in}{0.856767in}}%
\pgfpathlineto{\pgfqpoint{1.025857in}{0.840544in}}%
\pgfpathlineto{\pgfqpoint{1.037299in}{0.842337in}}%
\pgfpathlineto{\pgfqpoint{1.048955in}{0.843819in}}%
\pgfpathclose%
\pgfusepath{fill}%
\end{pgfscope}%
\begin{pgfscope}%
\pgfpathrectangle{\pgfqpoint{0.407592in}{0.222904in}}{\pgfqpoint{1.411163in}{1.411163in}}%
\pgfusepath{clip}%
\pgfsetbuttcap%
\pgfsetroundjoin%
\definecolor{currentfill}{rgb}{0.749020,0.772549,0.792157}%
\pgfsetfillcolor{currentfill}%
\pgfsetfillopacity{0.160000}%
\pgfsetlinewidth{0.000000pt}%
\definecolor{currentstroke}{rgb}{0.000000,0.000000,0.000000}%
\pgfsetstrokecolor{currentstroke}%
\pgfsetdash{}{0pt}%
\pgfpathmoveto{\pgfqpoint{0.526822in}{1.421759in}}%
\pgfpathlineto{\pgfqpoint{0.516943in}{1.435982in}}%
\pgfpathlineto{\pgfqpoint{0.484477in}{1.422362in}}%
\pgfpathlineto{\pgfqpoint{0.454204in}{1.408005in}}%
\pgfpathlineto{\pgfqpoint{0.465273in}{1.394126in}}%
\pgfpathlineto{\pgfqpoint{0.494965in}{1.408305in}}%
\pgfusepath{fill}%
\end{pgfscope}%
\begin{pgfscope}%
\pgfpathrectangle{\pgfqpoint{0.407592in}{0.222904in}}{\pgfqpoint{1.411163in}{1.411163in}}%
\pgfusepath{clip}%
\pgfsetbuttcap%
\pgfsetroundjoin%
\definecolor{currentfill}{rgb}{0.749020,0.772549,0.792157}%
\pgfsetfillcolor{currentfill}%
\pgfsetfillopacity{0.160000}%
\pgfsetlinewidth{0.000000pt}%
\definecolor{currentstroke}{rgb}{0.000000,0.000000,0.000000}%
\pgfsetstrokecolor{currentstroke}%
\pgfsetdash{}{0pt}%
\pgfpathmoveto{\pgfqpoint{1.301388in}{0.945345in}}%
\pgfpathlineto{\pgfqpoint{1.310877in}{0.960390in}}%
\pgfpathlineto{\pgfqpoint{1.320356in}{0.975420in}}%
\pgfpathlineto{\pgfqpoint{1.305573in}{0.980991in}}%
\pgfpathlineto{\pgfqpoint{1.289923in}{0.986140in}}%
\pgfpathlineto{\pgfqpoint{1.281724in}{0.970718in}}%
\pgfpathlineto{\pgfqpoint{1.273514in}{0.955277in}}%
\pgfpathlineto{\pgfqpoint{1.287852in}{0.950506in}}%
\pgfpathlineto{\pgfqpoint{1.301388in}{0.945345in}}%
\pgfpathclose%
\pgfusepath{fill}%
\end{pgfscope}%
\begin{pgfscope}%
\pgfpathrectangle{\pgfqpoint{0.407592in}{0.222904in}}{\pgfqpoint{1.411163in}{1.411163in}}%
\pgfusepath{clip}%
\pgfsetbuttcap%
\pgfsetroundjoin%
\definecolor{currentfill}{rgb}{0.749020,0.772549,0.792157}%
\pgfsetfillcolor{currentfill}%
\pgfsetfillopacity{0.160000}%
\pgfsetlinewidth{0.000000pt}%
\definecolor{currentstroke}{rgb}{0.000000,0.000000,0.000000}%
\pgfsetstrokecolor{currentstroke}%
\pgfsetdash{}{0pt}%
\pgfpathmoveto{\pgfqpoint{0.979438in}{0.899629in}}%
\pgfpathlineto{\pgfqpoint{0.973370in}{0.915464in}}%
\pgfpathlineto{\pgfqpoint{0.967310in}{0.931275in}}%
\pgfpathlineto{\pgfqpoint{0.952301in}{0.927572in}}%
\pgfpathlineto{\pgfqpoint{0.937841in}{0.923485in}}%
\pgfpathlineto{\pgfqpoint{0.945370in}{0.908022in}}%
\pgfpathlineto{\pgfqpoint{0.952909in}{0.892541in}}%
\pgfpathlineto{\pgfqpoint{0.965923in}{0.896259in}}%
\pgfpathlineto{\pgfqpoint{0.979438in}{0.899629in}}%
\pgfpathclose%
\pgfusepath{fill}%
\end{pgfscope}%
\begin{pgfscope}%
\pgfpathrectangle{\pgfqpoint{0.407592in}{0.222904in}}{\pgfqpoint{1.411163in}{1.411163in}}%
\pgfusepath{clip}%
\pgfsetbuttcap%
\pgfsetroundjoin%
\definecolor{currentfill}{rgb}{0.749020,0.772549,0.792157}%
\pgfsetfillcolor{currentfill}%
\pgfsetfillopacity{0.160000}%
\pgfsetlinewidth{0.000000pt}%
\definecolor{currentstroke}{rgb}{0.000000,0.000000,0.000000}%
\pgfsetstrokecolor{currentstroke}%
\pgfsetdash{}{0pt}%
\pgfpathmoveto{\pgfqpoint{1.410630in}{1.050344in}}%
\pgfpathlineto{\pgfqpoint{1.421198in}{1.064851in}}%
\pgfpathlineto{\pgfqpoint{1.431759in}{1.079348in}}%
\pgfpathlineto{\pgfqpoint{1.414446in}{1.087484in}}%
\pgfpathlineto{\pgfqpoint{1.395851in}{1.095124in}}%
\pgfpathlineto{\pgfqpoint{1.386447in}{1.080213in}}%
\pgfpathlineto{\pgfqpoint{1.377033in}{1.065287in}}%
\pgfpathlineto{\pgfqpoint{1.394437in}{1.058049in}}%
\pgfpathlineto{\pgfqpoint{1.410630in}{1.050344in}}%
\pgfpathclose%
\pgfusepath{fill}%
\end{pgfscope}%
\begin{pgfscope}%
\pgfpathrectangle{\pgfqpoint{0.407592in}{0.222904in}}{\pgfqpoint{1.411163in}{1.411163in}}%
\pgfusepath{clip}%
\pgfsetbuttcap%
\pgfsetroundjoin%
\definecolor{currentfill}{rgb}{0.749020,0.772549,0.792157}%
\pgfsetfillcolor{currentfill}%
\pgfsetfillopacity{0.160000}%
\pgfsetlinewidth{0.000000pt}%
\definecolor{currentstroke}{rgb}{0.000000,0.000000,0.000000}%
\pgfsetstrokecolor{currentstroke}%
\pgfsetdash{}{0pt}%
\pgfpathmoveto{\pgfqpoint{0.546597in}{1.393285in}}%
\pgfpathlineto{\pgfqpoint{0.536706in}{1.407526in}}%
\pgfpathlineto{\pgfqpoint{0.526822in}{1.421759in}}%
\pgfpathlineto{\pgfqpoint{0.494965in}{1.408305in}}%
\pgfpathlineto{\pgfqpoint{0.465273in}{1.394126in}}%
\pgfpathlineto{\pgfqpoint{0.476345in}{1.380243in}}%
\pgfpathlineto{\pgfqpoint{0.487421in}{1.366357in}}%
\pgfpathlineto{\pgfqpoint{0.515957in}{1.380172in}}%
\pgfpathlineto{\pgfqpoint{0.546597in}{1.393285in}}%
\pgfpathclose%
\pgfusepath{fill}%
\end{pgfscope}%
\begin{pgfscope}%
\pgfpathrectangle{\pgfqpoint{0.407592in}{0.222904in}}{\pgfqpoint{1.411163in}{1.411163in}}%
\pgfusepath{clip}%
\pgfsetbuttcap%
\pgfsetroundjoin%
\definecolor{currentfill}{rgb}{0.749020,0.772549,0.792157}%
\pgfsetfillcolor{currentfill}%
\pgfsetfillopacity{0.160000}%
\pgfsetlinewidth{0.000000pt}%
\definecolor{currentstroke}{rgb}{0.000000,0.000000,0.000000}%
\pgfsetstrokecolor{currentstroke}%
\pgfsetdash{}{0pt}%
\pgfpathmoveto{\pgfqpoint{0.922810in}{0.954352in}}%
\pgfpathlineto{\pgfqpoint{0.915309in}{0.969756in}}%
\pgfpathlineto{\pgfqpoint{0.907817in}{0.985142in}}%
\pgfpathlineto{\pgfqpoint{0.891185in}{0.979906in}}%
\pgfpathlineto{\pgfqpoint{0.875341in}{0.974253in}}%
\pgfpathlineto{\pgfqpoint{0.884195in}{0.959266in}}%
\pgfpathlineto{\pgfqpoint{0.893058in}{0.944265in}}%
\pgfpathlineto{\pgfqpoint{0.907570in}{0.949501in}}%
\pgfpathlineto{\pgfqpoint{0.922810in}{0.954352in}}%
\pgfpathclose%
\pgfusepath{fill}%
\end{pgfscope}%
\begin{pgfscope}%
\pgfpathrectangle{\pgfqpoint{0.407592in}{0.222904in}}{\pgfqpoint{1.411163in}{1.411163in}}%
\pgfusepath{clip}%
\pgfsetbuttcap%
\pgfsetroundjoin%
\definecolor{currentfill}{rgb}{0.749020,0.772549,0.792157}%
\pgfsetfillcolor{currentfill}%
\pgfsetfillopacity{0.160000}%
\pgfsetlinewidth{0.000000pt}%
\definecolor{currentstroke}{rgb}{0.000000,0.000000,0.000000}%
\pgfsetstrokecolor{currentstroke}%
\pgfsetdash{}{0pt}%
\pgfpathmoveto{\pgfqpoint{1.669786in}{1.313440in}}%
\pgfpathlineto{\pgfqpoint{1.681282in}{1.327382in}}%
\pgfpathlineto{\pgfqpoint{1.692774in}{1.341320in}}%
\pgfpathlineto{\pgfqpoint{1.668225in}{1.354628in}}%
\pgfpathlineto{\pgfqpoint{1.641494in}{1.367249in}}%
\pgfpathlineto{\pgfqpoint{1.631074in}{1.352946in}}%
\pgfpathlineto{\pgfqpoint{1.620647in}{1.338633in}}%
\pgfpathlineto{\pgfqpoint{1.646273in}{1.326368in}}%
\pgfpathlineto{\pgfqpoint{1.669786in}{1.313440in}}%
\pgfpathclose%
\pgfusepath{fill}%
\end{pgfscope}%
\begin{pgfscope}%
\pgfpathrectangle{\pgfqpoint{0.407592in}{0.222904in}}{\pgfqpoint{1.411163in}{1.411163in}}%
\pgfusepath{clip}%
\pgfsetbuttcap%
\pgfsetroundjoin%
\definecolor{currentfill}{rgb}{0.749020,0.772549,0.792157}%
\pgfsetfillcolor{currentfill}%
\pgfsetfillopacity{0.160000}%
\pgfsetlinewidth{0.000000pt}%
\definecolor{currentstroke}{rgb}{0.000000,0.000000,0.000000}%
\pgfsetstrokecolor{currentstroke}%
\pgfsetdash{}{0pt}%
\pgfpathmoveto{\pgfqpoint{1.116365in}{0.812715in}}%
\pgfpathlineto{\pgfqpoint{1.118516in}{0.829354in}}%
\pgfpathlineto{\pgfqpoint{1.120662in}{0.845962in}}%
\pgfpathlineto{\pgfqpoint{1.108745in}{0.846419in}}%
\pgfpathlineto{\pgfqpoint{1.096748in}{0.846549in}}%
\pgfpathlineto{\pgfqpoint{1.096288in}{0.829903in}}%
\pgfpathlineto{\pgfqpoint{1.095828in}{0.813224in}}%
\pgfpathlineto{\pgfqpoint{1.106130in}{0.813111in}}%
\pgfpathlineto{\pgfqpoint{1.116365in}{0.812715in}}%
\pgfpathclose%
\pgfusepath{fill}%
\end{pgfscope}%
\begin{pgfscope}%
\pgfpathrectangle{\pgfqpoint{0.407592in}{0.222904in}}{\pgfqpoint{1.411163in}{1.411163in}}%
\pgfusepath{clip}%
\pgfsetbuttcap%
\pgfsetroundjoin%
\definecolor{currentfill}{rgb}{0.749020,0.772549,0.792157}%
\pgfsetfillcolor{currentfill}%
\pgfsetfillopacity{0.160000}%
\pgfsetlinewidth{0.000000pt}%
\definecolor{currentstroke}{rgb}{0.000000,0.000000,0.000000}%
\pgfsetstrokecolor{currentstroke}%
\pgfsetdash{}{0pt}%
\pgfpathmoveto{\pgfqpoint{1.095828in}{0.813224in}}%
\pgfpathlineto{\pgfqpoint{1.096288in}{0.829903in}}%
\pgfpathlineto{\pgfqpoint{1.096748in}{0.846549in}}%
\pgfpathlineto{\pgfqpoint{1.084722in}{0.846353in}}%
\pgfpathlineto{\pgfqpoint{1.072715in}{0.845831in}}%
\pgfpathlineto{\pgfqpoint{1.073950in}{0.829232in}}%
\pgfpathlineto{\pgfqpoint{1.075188in}{0.812601in}}%
\pgfpathlineto{\pgfqpoint{1.085499in}{0.813054in}}%
\pgfpathlineto{\pgfqpoint{1.095828in}{0.813224in}}%
\pgfpathclose%
\pgfusepath{fill}%
\end{pgfscope}%
\begin{pgfscope}%
\pgfpathrectangle{\pgfqpoint{0.407592in}{0.222904in}}{\pgfqpoint{1.411163in}{1.411163in}}%
\pgfusepath{clip}%
\pgfsetbuttcap%
\pgfsetroundjoin%
\definecolor{currentfill}{rgb}{0.749020,0.772549,0.792157}%
\pgfsetfillcolor{currentfill}%
\pgfsetfillopacity{0.160000}%
\pgfsetlinewidth{0.000000pt}%
\definecolor{currentstroke}{rgb}{0.000000,0.000000,0.000000}%
\pgfsetstrokecolor{currentstroke}%
\pgfsetdash{}{0pt}%
\pgfpathmoveto{\pgfqpoint{0.822393in}{1.063874in}}%
\pgfpathlineto{\pgfqpoint{0.813598in}{1.078760in}}%
\pgfpathlineto{\pgfqpoint{0.804811in}{1.093633in}}%
\pgfpathlineto{\pgfqpoint{0.784714in}{1.085891in}}%
\pgfpathlineto{\pgfqpoint{0.765804in}{1.077659in}}%
\pgfpathlineto{\pgfqpoint{0.775842in}{1.063206in}}%
\pgfpathlineto{\pgfqpoint{0.785886in}{1.048744in}}%
\pgfpathlineto{\pgfqpoint{0.803578in}{1.056540in}}%
\pgfpathlineto{\pgfqpoint{0.822393in}{1.063874in}}%
\pgfpathclose%
\pgfusepath{fill}%
\end{pgfscope}%
\begin{pgfscope}%
\pgfpathrectangle{\pgfqpoint{0.407592in}{0.222904in}}{\pgfqpoint{1.411163in}{1.411163in}}%
\pgfusepath{clip}%
\pgfsetbuttcap%
\pgfsetroundjoin%
\definecolor{currentfill}{rgb}{0.749020,0.772549,0.792157}%
\pgfsetfillcolor{currentfill}%
\pgfsetfillopacity{0.160000}%
\pgfsetlinewidth{0.000000pt}%
\definecolor{currentstroke}{rgb}{0.000000,0.000000,0.000000}%
\pgfsetstrokecolor{currentstroke}%
\pgfsetdash{}{0pt}%
\pgfpathmoveto{\pgfqpoint{0.566398in}{1.364774in}}%
\pgfpathlineto{\pgfqpoint{0.556495in}{1.379034in}}%
\pgfpathlineto{\pgfqpoint{0.546597in}{1.393285in}}%
\pgfpathlineto{\pgfqpoint{0.515957in}{1.380172in}}%
\pgfpathlineto{\pgfqpoint{0.487421in}{1.366357in}}%
\pgfpathlineto{\pgfqpoint{0.498499in}{1.352466in}}%
\pgfpathlineto{\pgfqpoint{0.509581in}{1.338571in}}%
\pgfpathlineto{\pgfqpoint{0.536969in}{1.352012in}}%
\pgfpathlineto{\pgfqpoint{0.566398in}{1.364774in}}%
\pgfpathclose%
\pgfusepath{fill}%
\end{pgfscope}%
\begin{pgfscope}%
\pgfpathrectangle{\pgfqpoint{0.407592in}{0.222904in}}{\pgfqpoint{1.411163in}{1.411163in}}%
\pgfusepath{clip}%
\pgfsetbuttcap%
\pgfsetroundjoin%
\definecolor{currentfill}{rgb}{0.749020,0.772549,0.792157}%
\pgfsetfillcolor{currentfill}%
\pgfsetfillopacity{0.160000}%
\pgfsetlinewidth{0.000000pt}%
\definecolor{currentstroke}{rgb}{0.000000,0.000000,0.000000}%
\pgfsetstrokecolor{currentstroke}%
\pgfsetdash{}{0pt}%
\pgfpathmoveto{\pgfqpoint{1.646783in}{1.285542in}}%
\pgfpathlineto{\pgfqpoint{1.658286in}{1.299493in}}%
\pgfpathlineto{\pgfqpoint{1.669786in}{1.313440in}}%
\pgfpathlineto{\pgfqpoint{1.646273in}{1.326368in}}%
\pgfpathlineto{\pgfqpoint{1.620647in}{1.338633in}}%
\pgfpathlineto{\pgfqpoint{1.610213in}{1.324310in}}%
\pgfpathlineto{\pgfqpoint{1.599772in}{1.309978in}}%
\pgfpathlineto{\pgfqpoint{1.624299in}{1.298080in}}%
\pgfpathlineto{\pgfqpoint{1.646783in}{1.285542in}}%
\pgfpathclose%
\pgfusepath{fill}%
\end{pgfscope}%
\begin{pgfscope}%
\pgfpathrectangle{\pgfqpoint{0.407592in}{0.222904in}}{\pgfqpoint{1.411163in}{1.411163in}}%
\pgfusepath{clip}%
\pgfsetbuttcap%
\pgfsetroundjoin%
\definecolor{currentfill}{rgb}{0.749020,0.772549,0.792157}%
\pgfsetfillcolor{currentfill}%
\pgfsetfillopacity{0.160000}%
\pgfsetlinewidth{0.000000pt}%
\definecolor{currentstroke}{rgb}{0.000000,0.000000,0.000000}%
\pgfsetstrokecolor{currentstroke}%
\pgfsetdash{}{0pt}%
\pgfpathmoveto{\pgfqpoint{1.136461in}{0.811081in}}%
\pgfpathlineto{\pgfqpoint{1.140268in}{0.827595in}}%
\pgfpathlineto{\pgfqpoint{1.144068in}{0.844079in}}%
\pgfpathlineto{\pgfqpoint{1.132453in}{0.845181in}}%
\pgfpathlineto{\pgfqpoint{1.120662in}{0.845962in}}%
\pgfpathlineto{\pgfqpoint{1.118516in}{0.829354in}}%
\pgfpathlineto{\pgfqpoint{1.116365in}{0.812715in}}%
\pgfpathlineto{\pgfqpoint{1.126489in}{0.812037in}}%
\pgfpathlineto{\pgfqpoint{1.136461in}{0.811081in}}%
\pgfpathclose%
\pgfusepath{fill}%
\end{pgfscope}%
\begin{pgfscope}%
\pgfpathrectangle{\pgfqpoint{0.407592in}{0.222904in}}{\pgfqpoint{1.411163in}{1.411163in}}%
\pgfusepath{clip}%
\pgfsetbuttcap%
\pgfsetroundjoin%
\definecolor{currentfill}{rgb}{0.749020,0.772549,0.792157}%
\pgfsetfillcolor{currentfill}%
\pgfsetfillopacity{0.160000}%
\pgfsetlinewidth{0.000000pt}%
\definecolor{currentstroke}{rgb}{0.000000,0.000000,0.000000}%
\pgfsetstrokecolor{currentstroke}%
\pgfsetdash{}{0pt}%
\pgfpathmoveto{\pgfqpoint{1.075188in}{0.812601in}}%
\pgfpathlineto{\pgfqpoint{1.073950in}{0.829232in}}%
\pgfpathlineto{\pgfqpoint{1.072715in}{0.845831in}}%
\pgfpathlineto{\pgfqpoint{1.060776in}{0.844986in}}%
\pgfpathlineto{\pgfqpoint{1.048955in}{0.843819in}}%
\pgfpathlineto{\pgfqpoint{1.051868in}{0.827353in}}%
\pgfpathlineto{\pgfqpoint{1.054787in}{0.810856in}}%
\pgfpathlineto{\pgfqpoint{1.064937in}{0.811868in}}%
\pgfpathlineto{\pgfqpoint{1.075188in}{0.812601in}}%
\pgfpathclose%
\pgfusepath{fill}%
\end{pgfscope}%
\begin{pgfscope}%
\pgfpathrectangle{\pgfqpoint{0.407592in}{0.222904in}}{\pgfqpoint{1.411163in}{1.411163in}}%
\pgfusepath{clip}%
\pgfsetbuttcap%
\pgfsetroundjoin%
\definecolor{currentfill}{rgb}{0.749020,0.772549,0.792157}%
\pgfsetfillcolor{currentfill}%
\pgfsetfillopacity{0.160000}%
\pgfsetlinewidth{0.000000pt}%
\definecolor{currentstroke}{rgb}{0.000000,0.000000,0.000000}%
\pgfsetstrokecolor{currentstroke}%
\pgfsetdash{}{0pt}%
\pgfpathmoveto{\pgfqpoint{1.389472in}{1.021301in}}%
\pgfpathlineto{\pgfqpoint{1.400055in}{1.035827in}}%
\pgfpathlineto{\pgfqpoint{1.410630in}{1.050344in}}%
\pgfpathlineto{\pgfqpoint{1.394437in}{1.058049in}}%
\pgfpathlineto{\pgfqpoint{1.377033in}{1.065287in}}%
\pgfpathlineto{\pgfqpoint{1.367610in}{1.050346in}}%
\pgfpathlineto{\pgfqpoint{1.358178in}{1.035391in}}%
\pgfpathlineto{\pgfqpoint{1.374395in}{1.028565in}}%
\pgfpathlineto{\pgfqpoint{1.389472in}{1.021301in}}%
\pgfpathclose%
\pgfusepath{fill}%
\end{pgfscope}%
\begin{pgfscope}%
\pgfpathrectangle{\pgfqpoint{0.407592in}{0.222904in}}{\pgfqpoint{1.411163in}{1.411163in}}%
\pgfusepath{clip}%
\pgfsetbuttcap%
\pgfsetroundjoin%
\definecolor{currentfill}{rgb}{0.749020,0.772549,0.792157}%
\pgfsetfillcolor{currentfill}%
\pgfsetfillopacity{0.160000}%
\pgfsetlinewidth{0.000000pt}%
\definecolor{currentstroke}{rgb}{0.000000,0.000000,0.000000}%
\pgfsetstrokecolor{currentstroke}%
\pgfsetdash{}{0pt}%
\pgfpathmoveto{\pgfqpoint{1.187830in}{0.836558in}}%
\pgfpathlineto{\pgfqpoint{1.194722in}{0.852520in}}%
\pgfpathlineto{\pgfqpoint{1.201605in}{0.868458in}}%
\pgfpathlineto{\pgfqpoint{1.189677in}{0.871095in}}%
\pgfpathlineto{\pgfqpoint{1.177341in}{0.873396in}}%
\pgfpathlineto{\pgfqpoint{1.171965in}{0.857176in}}%
\pgfpathlineto{\pgfqpoint{1.166581in}{0.840928in}}%
\pgfpathlineto{\pgfqpoint{1.177386in}{0.838892in}}%
\pgfpathlineto{\pgfqpoint{1.187830in}{0.836558in}}%
\pgfpathclose%
\pgfusepath{fill}%
\end{pgfscope}%
\begin{pgfscope}%
\pgfpathrectangle{\pgfqpoint{0.407592in}{0.222904in}}{\pgfqpoint{1.411163in}{1.411163in}}%
\pgfusepath{clip}%
\pgfsetbuttcap%
\pgfsetroundjoin%
\definecolor{currentfill}{rgb}{0.749020,0.772549,0.792157}%
\pgfsetfillcolor{currentfill}%
\pgfsetfillopacity{0.160000}%
\pgfsetlinewidth{0.000000pt}%
\definecolor{currentstroke}{rgb}{0.000000,0.000000,0.000000}%
\pgfsetstrokecolor{currentstroke}%
\pgfsetdash{}{0pt}%
\pgfpathmoveto{\pgfqpoint{1.282383in}{0.915211in}}%
\pgfpathlineto{\pgfqpoint{1.291891in}{0.930285in}}%
\pgfpathlineto{\pgfqpoint{1.301388in}{0.945345in}}%
\pgfpathlineto{\pgfqpoint{1.287852in}{0.950506in}}%
\pgfpathlineto{\pgfqpoint{1.273514in}{0.955277in}}%
\pgfpathlineto{\pgfqpoint{1.265294in}{0.939816in}}%
\pgfpathlineto{\pgfqpoint{1.257064in}{0.924335in}}%
\pgfpathlineto{\pgfqpoint{1.270091in}{0.919951in}}%
\pgfpathlineto{\pgfqpoint{1.282383in}{0.915211in}}%
\pgfpathclose%
\pgfusepath{fill}%
\end{pgfscope}%
\begin{pgfscope}%
\pgfpathrectangle{\pgfqpoint{0.407592in}{0.222904in}}{\pgfqpoint{1.411163in}{1.411163in}}%
\pgfusepath{clip}%
\pgfsetbuttcap%
\pgfsetroundjoin%
\definecolor{currentfill}{rgb}{0.749020,0.772549,0.792157}%
\pgfsetfillcolor{currentfill}%
\pgfsetfillopacity{0.160000}%
\pgfsetlinewidth{0.000000pt}%
\definecolor{currentstroke}{rgb}{0.000000,0.000000,0.000000}%
\pgfsetstrokecolor{currentstroke}%
\pgfsetdash{}{0pt}%
\pgfpathmoveto{\pgfqpoint{1.025857in}{0.840544in}}%
\pgfpathlineto{\pgfqpoint{1.021317in}{0.856767in}}%
\pgfpathlineto{\pgfqpoint{1.016784in}{0.872962in}}%
\pgfpathlineto{\pgfqpoint{1.004020in}{0.870593in}}%
\pgfpathlineto{\pgfqpoint{0.991602in}{0.867889in}}%
\pgfpathlineto{\pgfqpoint{0.997698in}{0.851984in}}%
\pgfpathlineto{\pgfqpoint{1.003802in}{0.836054in}}%
\pgfpathlineto{\pgfqpoint{1.014676in}{0.838447in}}%
\pgfpathlineto{\pgfqpoint{1.025857in}{0.840544in}}%
\pgfpathclose%
\pgfusepath{fill}%
\end{pgfscope}%
\begin{pgfscope}%
\pgfpathrectangle{\pgfqpoint{0.407592in}{0.222904in}}{\pgfqpoint{1.411163in}{1.411163in}}%
\pgfusepath{clip}%
\pgfsetbuttcap%
\pgfsetroundjoin%
\definecolor{currentfill}{rgb}{0.749020,0.772549,0.792157}%
\pgfsetfillcolor{currentfill}%
\pgfsetfillopacity{0.160000}%
\pgfsetlinewidth{0.000000pt}%
\definecolor{currentstroke}{rgb}{0.000000,0.000000,0.000000}%
\pgfsetstrokecolor{currentstroke}%
\pgfsetdash{}{0pt}%
\pgfpathmoveto{\pgfqpoint{1.224037in}{0.862214in}}%
\pgfpathlineto{\pgfqpoint{1.232309in}{0.877774in}}%
\pgfpathlineto{\pgfqpoint{1.240571in}{0.893314in}}%
\pgfpathlineto{\pgfqpoint{1.228237in}{0.896964in}}%
\pgfpathlineto{\pgfqpoint{1.215338in}{0.900262in}}%
\pgfpathlineto{\pgfqpoint{1.208477in}{0.884372in}}%
\pgfpathlineto{\pgfqpoint{1.201605in}{0.868458in}}%
\pgfpathlineto{\pgfqpoint{1.213074in}{0.865494in}}%
\pgfpathlineto{\pgfqpoint{1.224037in}{0.862214in}}%
\pgfpathclose%
\pgfusepath{fill}%
\end{pgfscope}%
\begin{pgfscope}%
\pgfpathrectangle{\pgfqpoint{0.407592in}{0.222904in}}{\pgfqpoint{1.411163in}{1.411163in}}%
\pgfusepath{clip}%
\pgfsetbuttcap%
\pgfsetroundjoin%
\definecolor{currentfill}{rgb}{0.749020,0.772549,0.792157}%
\pgfsetfillcolor{currentfill}%
\pgfsetfillopacity{0.160000}%
\pgfsetlinewidth{0.000000pt}%
\definecolor{currentstroke}{rgb}{0.000000,0.000000,0.000000}%
\pgfsetstrokecolor{currentstroke}%
\pgfsetdash{}{0pt}%
\pgfpathmoveto{\pgfqpoint{0.586224in}{1.336227in}}%
\pgfpathlineto{\pgfqpoint{0.576308in}{1.350505in}}%
\pgfpathlineto{\pgfqpoint{0.566398in}{1.364774in}}%
\pgfpathlineto{\pgfqpoint{0.536969in}{1.352012in}}%
\pgfpathlineto{\pgfqpoint{0.509581in}{1.338571in}}%
\pgfpathlineto{\pgfqpoint{0.520666in}{1.324673in}}%
\pgfpathlineto{\pgfqpoint{0.531754in}{1.310770in}}%
\pgfpathlineto{\pgfqpoint{0.558000in}{1.323826in}}%
\pgfpathlineto{\pgfqpoint{0.586224in}{1.336227in}}%
\pgfpathclose%
\pgfusepath{fill}%
\end{pgfscope}%
\begin{pgfscope}%
\pgfpathrectangle{\pgfqpoint{0.407592in}{0.222904in}}{\pgfqpoint{1.411163in}{1.411163in}}%
\pgfusepath{clip}%
\pgfsetbuttcap%
\pgfsetroundjoin%
\definecolor{currentfill}{rgb}{0.749020,0.772549,0.792157}%
\pgfsetfillcolor{currentfill}%
\pgfsetfillopacity{0.160000}%
\pgfsetlinewidth{0.000000pt}%
\definecolor{currentstroke}{rgb}{0.000000,0.000000,0.000000}%
\pgfsetstrokecolor{currentstroke}%
\pgfsetdash{}{0pt}%
\pgfpathmoveto{\pgfqpoint{1.623766in}{1.257625in}}%
\pgfpathlineto{\pgfqpoint{1.635277in}{1.271586in}}%
\pgfpathlineto{\pgfqpoint{1.646783in}{1.285542in}}%
\pgfpathlineto{\pgfqpoint{1.624299in}{1.298080in}}%
\pgfpathlineto{\pgfqpoint{1.599772in}{1.309978in}}%
\pgfpathlineto{\pgfqpoint{1.589325in}{1.295637in}}%
\pgfpathlineto{\pgfqpoint{1.578870in}{1.281285in}}%
\pgfpathlineto{\pgfqpoint{1.602304in}{1.269763in}}%
\pgfpathlineto{\pgfqpoint{1.623766in}{1.257625in}}%
\pgfpathclose%
\pgfusepath{fill}%
\end{pgfscope}%
\begin{pgfscope}%
\pgfpathrectangle{\pgfqpoint{0.407592in}{0.222904in}}{\pgfqpoint{1.411163in}{1.411163in}}%
\pgfusepath{clip}%
\pgfsetbuttcap%
\pgfsetroundjoin%
\definecolor{currentfill}{rgb}{0.749020,0.772549,0.792157}%
\pgfsetfillcolor{currentfill}%
\pgfsetfillopacity{0.160000}%
\pgfsetlinewidth{0.000000pt}%
\definecolor{currentstroke}{rgb}{0.000000,0.000000,0.000000}%
\pgfsetstrokecolor{currentstroke}%
\pgfsetdash{}{0pt}%
\pgfpathmoveto{\pgfqpoint{0.840008in}{1.034058in}}%
\pgfpathlineto{\pgfqpoint{0.831196in}{1.048973in}}%
\pgfpathlineto{\pgfqpoint{0.822393in}{1.063874in}}%
\pgfpathlineto{\pgfqpoint{0.803578in}{1.056540in}}%
\pgfpathlineto{\pgfqpoint{0.785886in}{1.048744in}}%
\pgfpathlineto{\pgfqpoint{0.795936in}{1.034273in}}%
\pgfpathlineto{\pgfqpoint{0.805993in}{1.019793in}}%
\pgfpathlineto{\pgfqpoint{0.822472in}{1.027142in}}%
\pgfpathlineto{\pgfqpoint{0.840008in}{1.034058in}}%
\pgfpathclose%
\pgfusepath{fill}%
\end{pgfscope}%
\begin{pgfscope}%
\pgfpathrectangle{\pgfqpoint{0.407592in}{0.222904in}}{\pgfqpoint{1.411163in}{1.411163in}}%
\pgfusepath{clip}%
\pgfsetbuttcap%
\pgfsetroundjoin%
\definecolor{currentfill}{rgb}{0.749020,0.772549,0.792157}%
\pgfsetfillcolor{currentfill}%
\pgfsetfillopacity{0.160000}%
\pgfsetlinewidth{0.000000pt}%
\definecolor{currentstroke}{rgb}{0.000000,0.000000,0.000000}%
\pgfsetstrokecolor{currentstroke}%
\pgfsetdash{}{0pt}%
\pgfpathmoveto{\pgfqpoint{0.991602in}{0.867889in}}%
\pgfpathlineto{\pgfqpoint{0.985516in}{0.883771in}}%
\pgfpathlineto{\pgfqpoint{0.979438in}{0.899629in}}%
\pgfpathlineto{\pgfqpoint{0.965923in}{0.896259in}}%
\pgfpathlineto{\pgfqpoint{0.952909in}{0.892541in}}%
\pgfpathlineto{\pgfqpoint{0.960457in}{0.877040in}}%
\pgfpathlineto{\pgfqpoint{0.968014in}{0.861519in}}%
\pgfpathlineto{\pgfqpoint{0.979583in}{0.864861in}}%
\pgfpathlineto{\pgfqpoint{0.991602in}{0.867889in}}%
\pgfpathclose%
\pgfusepath{fill}%
\end{pgfscope}%
\begin{pgfscope}%
\pgfpathrectangle{\pgfqpoint{0.407592in}{0.222904in}}{\pgfqpoint{1.411163in}{1.411163in}}%
\pgfusepath{clip}%
\pgfsetbuttcap%
\pgfsetroundjoin%
\definecolor{currentfill}{rgb}{0.749020,0.772549,0.792157}%
\pgfsetfillcolor{currentfill}%
\pgfsetfillopacity{0.160000}%
\pgfsetlinewidth{0.000000pt}%
\definecolor{currentstroke}{rgb}{0.000000,0.000000,0.000000}%
\pgfsetstrokecolor{currentstroke}%
\pgfsetdash{}{0pt}%
\pgfpathmoveto{\pgfqpoint{0.937841in}{0.923485in}}%
\pgfpathlineto{\pgfqpoint{0.930321in}{0.938928in}}%
\pgfpathlineto{\pgfqpoint{0.922810in}{0.954352in}}%
\pgfpathlineto{\pgfqpoint{0.907570in}{0.949501in}}%
\pgfpathlineto{\pgfqpoint{0.893058in}{0.944265in}}%
\pgfpathlineto{\pgfqpoint{0.901930in}{0.929249in}}%
\pgfpathlineto{\pgfqpoint{0.910810in}{0.914218in}}%
\pgfpathlineto{\pgfqpoint{0.923991in}{0.919028in}}%
\pgfpathlineto{\pgfqpoint{0.937841in}{0.923485in}}%
\pgfpathclose%
\pgfusepath{fill}%
\end{pgfscope}%
\begin{pgfscope}%
\pgfpathrectangle{\pgfqpoint{0.407592in}{0.222904in}}{\pgfqpoint{1.411163in}{1.411163in}}%
\pgfusepath{clip}%
\pgfsetbuttcap%
\pgfsetroundjoin%
\definecolor{currentfill}{rgb}{0.749020,0.772549,0.792157}%
\pgfsetfillcolor{currentfill}%
\pgfsetfillopacity{0.160000}%
\pgfsetlinewidth{0.000000pt}%
\definecolor{currentstroke}{rgb}{0.000000,0.000000,0.000000}%
\pgfsetstrokecolor{currentstroke}%
\pgfsetdash{}{0pt}%
\pgfpathmoveto{\pgfqpoint{1.155785in}{0.808349in}}%
\pgfpathlineto{\pgfqpoint{1.161187in}{0.824653in}}%
\pgfpathlineto{\pgfqpoint{1.166581in}{0.840928in}}%
\pgfpathlineto{\pgfqpoint{1.155459in}{0.842659in}}%
\pgfpathlineto{\pgfqpoint{1.144068in}{0.844079in}}%
\pgfpathlineto{\pgfqpoint{1.140268in}{0.827595in}}%
\pgfpathlineto{\pgfqpoint{1.136461in}{0.811081in}}%
\pgfpathlineto{\pgfqpoint{1.146240in}{0.809851in}}%
\pgfpathlineto{\pgfqpoint{1.155785in}{0.808349in}}%
\pgfpathclose%
\pgfusepath{fill}%
\end{pgfscope}%
\begin{pgfscope}%
\pgfpathrectangle{\pgfqpoint{0.407592in}{0.222904in}}{\pgfqpoint{1.411163in}{1.411163in}}%
\pgfusepath{clip}%
\pgfsetbuttcap%
\pgfsetroundjoin%
\definecolor{currentfill}{rgb}{0.749020,0.772549,0.792157}%
\pgfsetfillcolor{currentfill}%
\pgfsetfillopacity{0.160000}%
\pgfsetlinewidth{0.000000pt}%
\definecolor{currentstroke}{rgb}{0.000000,0.000000,0.000000}%
\pgfsetstrokecolor{currentstroke}%
\pgfsetdash{}{0pt}%
\pgfpathmoveto{\pgfqpoint{1.054787in}{0.810856in}}%
\pgfpathlineto{\pgfqpoint{1.051868in}{0.827353in}}%
\pgfpathlineto{\pgfqpoint{1.048955in}{0.843819in}}%
\pgfpathlineto{\pgfqpoint{1.037299in}{0.842337in}}%
\pgfpathlineto{\pgfqpoint{1.025857in}{0.840544in}}%
\pgfpathlineto{\pgfqpoint{1.030405in}{0.824294in}}%
\pgfpathlineto{\pgfqpoint{1.034960in}{0.808016in}}%
\pgfpathlineto{\pgfqpoint{1.044781in}{0.809571in}}%
\pgfpathlineto{\pgfqpoint{1.054787in}{0.810856in}}%
\pgfpathclose%
\pgfusepath{fill}%
\end{pgfscope}%
\begin{pgfscope}%
\pgfpathrectangle{\pgfqpoint{0.407592in}{0.222904in}}{\pgfqpoint{1.411163in}{1.411163in}}%
\pgfusepath{clip}%
\pgfsetbuttcap%
\pgfsetroundjoin%
\definecolor{currentfill}{rgb}{0.749020,0.772549,0.792157}%
\pgfsetfillcolor{currentfill}%
\pgfsetfillopacity{0.160000}%
\pgfsetlinewidth{0.000000pt}%
\definecolor{currentstroke}{rgb}{0.000000,0.000000,0.000000}%
\pgfsetstrokecolor{currentstroke}%
\pgfsetdash{}{0pt}%
\pgfpathmoveto{\pgfqpoint{0.606076in}{1.307644in}}%
\pgfpathlineto{\pgfqpoint{0.596147in}{1.321940in}}%
\pgfpathlineto{\pgfqpoint{0.586224in}{1.336227in}}%
\pgfpathlineto{\pgfqpoint{0.558000in}{1.323826in}}%
\pgfpathlineto{\pgfqpoint{0.531754in}{1.310770in}}%
\pgfpathlineto{\pgfqpoint{0.542845in}{1.296863in}}%
\pgfpathlineto{\pgfqpoint{0.553940in}{1.282953in}}%
\pgfpathlineto{\pgfqpoint{0.579051in}{1.295614in}}%
\pgfpathlineto{\pgfqpoint{0.606076in}{1.307644in}}%
\pgfpathclose%
\pgfusepath{fill}%
\end{pgfscope}%
\begin{pgfscope}%
\pgfpathrectangle{\pgfqpoint{0.407592in}{0.222904in}}{\pgfqpoint{1.411163in}{1.411163in}}%
\pgfusepath{clip}%
\pgfsetbuttcap%
\pgfsetroundjoin%
\definecolor{currentfill}{rgb}{0.749020,0.772549,0.792157}%
\pgfsetfillcolor{currentfill}%
\pgfsetfillopacity{0.160000}%
\pgfsetlinewidth{0.000000pt}%
\definecolor{currentstroke}{rgb}{0.000000,0.000000,0.000000}%
\pgfsetstrokecolor{currentstroke}%
\pgfsetdash{}{0pt}%
\pgfpathmoveto{\pgfqpoint{1.368286in}{0.992218in}}%
\pgfpathlineto{\pgfqpoint{1.378882in}{1.006764in}}%
\pgfpathlineto{\pgfqpoint{1.389472in}{1.021301in}}%
\pgfpathlineto{\pgfqpoint{1.374395in}{1.028565in}}%
\pgfpathlineto{\pgfqpoint{1.358178in}{1.035391in}}%
\pgfpathlineto{\pgfqpoint{1.348737in}{1.020420in}}%
\pgfpathlineto{\pgfqpoint{1.339286in}{1.005435in}}%
\pgfpathlineto{\pgfqpoint{1.354319in}{0.999031in}}%
\pgfpathlineto{\pgfqpoint{1.368286in}{0.992218in}}%
\pgfpathclose%
\pgfusepath{fill}%
\end{pgfscope}%
\begin{pgfscope}%
\pgfpathrectangle{\pgfqpoint{0.407592in}{0.222904in}}{\pgfqpoint{1.411163in}{1.411163in}}%
\pgfusepath{clip}%
\pgfsetbuttcap%
\pgfsetroundjoin%
\definecolor{currentfill}{rgb}{0.749020,0.772549,0.792157}%
\pgfsetfillcolor{currentfill}%
\pgfsetfillopacity{0.160000}%
\pgfsetlinewidth{0.000000pt}%
\definecolor{currentstroke}{rgb}{0.000000,0.000000,0.000000}%
\pgfsetstrokecolor{currentstroke}%
\pgfsetdash{}{0pt}%
\pgfpathmoveto{\pgfqpoint{1.600734in}{1.229691in}}%
\pgfpathlineto{\pgfqpoint{1.612252in}{1.243660in}}%
\pgfpathlineto{\pgfqpoint{1.623766in}{1.257625in}}%
\pgfpathlineto{\pgfqpoint{1.602304in}{1.269763in}}%
\pgfpathlineto{\pgfqpoint{1.578870in}{1.281285in}}%
\pgfpathlineto{\pgfqpoint{1.568408in}{1.266924in}}%
\pgfpathlineto{\pgfqpoint{1.557939in}{1.252553in}}%
\pgfpathlineto{\pgfqpoint{1.580285in}{1.241417in}}%
\pgfpathlineto{\pgfqpoint{1.600734in}{1.229691in}}%
\pgfpathclose%
\pgfusepath{fill}%
\end{pgfscope}%
\begin{pgfscope}%
\pgfpathrectangle{\pgfqpoint{0.407592in}{0.222904in}}{\pgfqpoint{1.411163in}{1.411163in}}%
\pgfusepath{clip}%
\pgfsetbuttcap%
\pgfsetroundjoin%
\definecolor{currentfill}{rgb}{0.749020,0.772549,0.792157}%
\pgfsetfillcolor{currentfill}%
\pgfsetfillopacity{0.160000}%
\pgfsetlinewidth{0.000000pt}%
\definecolor{currentstroke}{rgb}{0.000000,0.000000,0.000000}%
\pgfsetstrokecolor{currentstroke}%
\pgfsetdash{}{0pt}%
\pgfpathmoveto{\pgfqpoint{0.625953in}{1.279025in}}%
\pgfpathlineto{\pgfqpoint{0.616011in}{1.293339in}}%
\pgfpathlineto{\pgfqpoint{0.606076in}{1.307644in}}%
\pgfpathlineto{\pgfqpoint{0.579051in}{1.295614in}}%
\pgfpathlineto{\pgfqpoint{0.553940in}{1.282953in}}%
\pgfpathlineto{\pgfqpoint{0.565037in}{1.269038in}}%
\pgfpathlineto{\pgfqpoint{0.576138in}{1.255120in}}%
\pgfpathlineto{\pgfqpoint{0.600121in}{1.267376in}}%
\pgfpathlineto{\pgfqpoint{0.625953in}{1.279025in}}%
\pgfpathclose%
\pgfusepath{fill}%
\end{pgfscope}%
\begin{pgfscope}%
\pgfpathrectangle{\pgfqpoint{0.407592in}{0.222904in}}{\pgfqpoint{1.411163in}{1.411163in}}%
\pgfusepath{clip}%
\pgfsetbuttcap%
\pgfsetroundjoin%
\definecolor{currentfill}{rgb}{0.749020,0.772549,0.792157}%
\pgfsetfillcolor{currentfill}%
\pgfsetfillopacity{0.160000}%
\pgfsetlinewidth{0.000000pt}%
\definecolor{currentstroke}{rgb}{0.000000,0.000000,0.000000}%
\pgfsetstrokecolor{currentstroke}%
\pgfsetdash{}{0pt}%
\pgfpathmoveto{\pgfqpoint{0.857658in}{1.004184in}}%
\pgfpathlineto{\pgfqpoint{0.848829in}{1.019128in}}%
\pgfpathlineto{\pgfqpoint{0.840008in}{1.034058in}}%
\pgfpathlineto{\pgfqpoint{0.822472in}{1.027142in}}%
\pgfpathlineto{\pgfqpoint{0.805993in}{1.019793in}}%
\pgfpathlineto{\pgfqpoint{0.816056in}{1.005303in}}%
\pgfpathlineto{\pgfqpoint{0.826126in}{0.990804in}}%
\pgfpathlineto{\pgfqpoint{0.841397in}{0.997696in}}%
\pgfpathlineto{\pgfqpoint{0.857658in}{1.004184in}}%
\pgfpathclose%
\pgfusepath{fill}%
\end{pgfscope}%
\begin{pgfscope}%
\pgfpathrectangle{\pgfqpoint{0.407592in}{0.222904in}}{\pgfqpoint{1.411163in}{1.411163in}}%
\pgfusepath{clip}%
\pgfsetbuttcap%
\pgfsetroundjoin%
\definecolor{currentfill}{rgb}{0.749020,0.772549,0.792157}%
\pgfsetfillcolor{currentfill}%
\pgfsetfillopacity{0.160000}%
\pgfsetlinewidth{0.000000pt}%
\definecolor{currentstroke}{rgb}{0.000000,0.000000,0.000000}%
\pgfsetstrokecolor{currentstroke}%
\pgfsetdash{}{0pt}%
\pgfpathmoveto{\pgfqpoint{1.263340in}{0.885016in}}%
\pgfpathlineto{\pgfqpoint{1.272866in}{0.900121in}}%
\pgfpathlineto{\pgfqpoint{1.282383in}{0.915211in}}%
\pgfpathlineto{\pgfqpoint{1.270091in}{0.919951in}}%
\pgfpathlineto{\pgfqpoint{1.257064in}{0.924335in}}%
\pgfpathlineto{\pgfqpoint{1.248823in}{0.908834in}}%
\pgfpathlineto{\pgfqpoint{1.240571in}{0.893314in}}%
\pgfpathlineto{\pgfqpoint{1.252289in}{0.889326in}}%
\pgfpathlineto{\pgfqpoint{1.263340in}{0.885016in}}%
\pgfpathclose%
\pgfusepath{fill}%
\end{pgfscope}%
\begin{pgfscope}%
\pgfpathrectangle{\pgfqpoint{0.407592in}{0.222904in}}{\pgfqpoint{1.411163in}{1.411163in}}%
\pgfusepath{clip}%
\pgfsetbuttcap%
\pgfsetroundjoin%
\definecolor{currentfill}{rgb}{0.749020,0.772549,0.792157}%
\pgfsetfillcolor{currentfill}%
\pgfsetfillopacity{0.160000}%
\pgfsetlinewidth{0.000000pt}%
\definecolor{currentstroke}{rgb}{0.000000,0.000000,0.000000}%
\pgfsetstrokecolor{currentstroke}%
\pgfsetdash{}{0pt}%
\pgfpathmoveto{\pgfqpoint{1.577687in}{1.201739in}}%
\pgfpathlineto{\pgfqpoint{1.589212in}{1.215717in}}%
\pgfpathlineto{\pgfqpoint{1.600734in}{1.229691in}}%
\pgfpathlineto{\pgfqpoint{1.580285in}{1.241417in}}%
\pgfpathlineto{\pgfqpoint{1.557939in}{1.252553in}}%
\pgfpathlineto{\pgfqpoint{1.547463in}{1.238173in}}%
\pgfpathlineto{\pgfqpoint{1.536980in}{1.223783in}}%
\pgfpathlineto{\pgfqpoint{1.558245in}{1.213043in}}%
\pgfpathlineto{\pgfqpoint{1.577687in}{1.201739in}}%
\pgfpathclose%
\pgfusepath{fill}%
\end{pgfscope}%
\begin{pgfscope}%
\pgfpathrectangle{\pgfqpoint{0.407592in}{0.222904in}}{\pgfqpoint{1.411163in}{1.411163in}}%
\pgfusepath{clip}%
\pgfsetbuttcap%
\pgfsetroundjoin%
\definecolor{currentfill}{rgb}{0.749020,0.772549,0.792157}%
\pgfsetfillcolor{currentfill}%
\pgfsetfillopacity{0.160000}%
\pgfsetlinewidth{0.000000pt}%
\definecolor{currentstroke}{rgb}{0.000000,0.000000,0.000000}%
\pgfsetstrokecolor{currentstroke}%
\pgfsetdash{}{0pt}%
\pgfpathmoveto{\pgfqpoint{1.112050in}{0.779340in}}%
\pgfpathlineto{\pgfqpoint{1.114210in}{0.796043in}}%
\pgfpathlineto{\pgfqpoint{1.116365in}{0.812715in}}%
\pgfpathlineto{\pgfqpoint{1.106130in}{0.813111in}}%
\pgfpathlineto{\pgfqpoint{1.095828in}{0.813224in}}%
\pgfpathlineto{\pgfqpoint{1.095366in}{0.796513in}}%
\pgfpathlineto{\pgfqpoint{1.094904in}{0.779770in}}%
\pgfpathlineto{\pgfqpoint{1.103506in}{0.779674in}}%
\pgfpathlineto{\pgfqpoint{1.112050in}{0.779340in}}%
\pgfpathclose%
\pgfusepath{fill}%
\end{pgfscope}%
\begin{pgfscope}%
\pgfpathrectangle{\pgfqpoint{0.407592in}{0.222904in}}{\pgfqpoint{1.411163in}{1.411163in}}%
\pgfusepath{clip}%
\pgfsetbuttcap%
\pgfsetroundjoin%
\definecolor{currentfill}{rgb}{0.749020,0.772549,0.792157}%
\pgfsetfillcolor{currentfill}%
\pgfsetfillopacity{0.160000}%
\pgfsetlinewidth{0.000000pt}%
\definecolor{currentstroke}{rgb}{0.000000,0.000000,0.000000}%
\pgfsetstrokecolor{currentstroke}%
\pgfsetdash{}{0pt}%
\pgfpathmoveto{\pgfqpoint{1.094904in}{0.779770in}}%
\pgfpathlineto{\pgfqpoint{1.095366in}{0.796513in}}%
\pgfpathlineto{\pgfqpoint{1.095828in}{0.813224in}}%
\pgfpathlineto{\pgfqpoint{1.085499in}{0.813054in}}%
\pgfpathlineto{\pgfqpoint{1.075188in}{0.812601in}}%
\pgfpathlineto{\pgfqpoint{1.076429in}{0.795939in}}%
\pgfpathlineto{\pgfqpoint{1.077671in}{0.779245in}}%
\pgfpathlineto{\pgfqpoint{1.086280in}{0.779626in}}%
\pgfpathlineto{\pgfqpoint{1.094904in}{0.779770in}}%
\pgfpathclose%
\pgfusepath{fill}%
\end{pgfscope}%
\begin{pgfscope}%
\pgfpathrectangle{\pgfqpoint{0.407592in}{0.222904in}}{\pgfqpoint{1.411163in}{1.411163in}}%
\pgfusepath{clip}%
\pgfsetbuttcap%
\pgfsetroundjoin%
\definecolor{currentfill}{rgb}{0.749020,0.772549,0.792157}%
\pgfsetfillcolor{currentfill}%
\pgfsetfillopacity{0.160000}%
\pgfsetlinewidth{0.000000pt}%
\definecolor{currentstroke}{rgb}{0.000000,0.000000,0.000000}%
\pgfsetstrokecolor{currentstroke}%
\pgfsetdash{}{0pt}%
\pgfpathmoveto{\pgfqpoint{0.952909in}{0.892541in}}%
\pgfpathlineto{\pgfqpoint{0.945370in}{0.908022in}}%
\pgfpathlineto{\pgfqpoint{0.937841in}{0.923485in}}%
\pgfpathlineto{\pgfqpoint{0.923991in}{0.919028in}}%
\pgfpathlineto{\pgfqpoint{0.910810in}{0.914218in}}%
\pgfpathlineto{\pgfqpoint{0.919699in}{0.899173in}}%
\pgfpathlineto{\pgfqpoint{0.928596in}{0.884114in}}%
\pgfpathlineto{\pgfqpoint{0.940448in}{0.888487in}}%
\pgfpathlineto{\pgfqpoint{0.952909in}{0.892541in}}%
\pgfpathclose%
\pgfusepath{fill}%
\end{pgfscope}%
\begin{pgfscope}%
\pgfpathrectangle{\pgfqpoint{0.407592in}{0.222904in}}{\pgfqpoint{1.411163in}{1.411163in}}%
\pgfusepath{clip}%
\pgfsetbuttcap%
\pgfsetroundjoin%
\definecolor{currentfill}{rgb}{0.749020,0.772549,0.792157}%
\pgfsetfillcolor{currentfill}%
\pgfsetfillopacity{0.160000}%
\pgfsetlinewidth{0.000000pt}%
\definecolor{currentstroke}{rgb}{0.000000,0.000000,0.000000}%
\pgfsetstrokecolor{currentstroke}%
\pgfsetdash{}{0pt}%
\pgfpathmoveto{\pgfqpoint{1.207460in}{0.831035in}}%
\pgfpathlineto{\pgfqpoint{1.215754in}{0.846635in}}%
\pgfpathlineto{\pgfqpoint{1.224037in}{0.862214in}}%
\pgfpathlineto{\pgfqpoint{1.213074in}{0.865494in}}%
\pgfpathlineto{\pgfqpoint{1.201605in}{0.868458in}}%
\pgfpathlineto{\pgfqpoint{1.194722in}{0.852520in}}%
\pgfpathlineto{\pgfqpoint{1.187830in}{0.836558in}}%
\pgfpathlineto{\pgfqpoint{1.197869in}{0.833936in}}%
\pgfpathlineto{\pgfqpoint{1.207460in}{0.831035in}}%
\pgfpathclose%
\pgfusepath{fill}%
\end{pgfscope}%
\begin{pgfscope}%
\pgfpathrectangle{\pgfqpoint{0.407592in}{0.222904in}}{\pgfqpoint{1.411163in}{1.411163in}}%
\pgfusepath{clip}%
\pgfsetbuttcap%
\pgfsetroundjoin%
\definecolor{currentfill}{rgb}{0.749020,0.772549,0.792157}%
\pgfsetfillcolor{currentfill}%
\pgfsetfillopacity{0.160000}%
\pgfsetlinewidth{0.000000pt}%
\definecolor{currentstroke}{rgb}{0.000000,0.000000,0.000000}%
\pgfsetstrokecolor{currentstroke}%
\pgfsetdash{}{0pt}%
\pgfpathmoveto{\pgfqpoint{1.347070in}{0.963096in}}%
\pgfpathlineto{\pgfqpoint{1.357682in}{0.977662in}}%
\pgfpathlineto{\pgfqpoint{1.368286in}{0.992218in}}%
\pgfpathlineto{\pgfqpoint{1.354319in}{0.999031in}}%
\pgfpathlineto{\pgfqpoint{1.339286in}{1.005435in}}%
\pgfpathlineto{\pgfqpoint{1.329826in}{0.990435in}}%
\pgfpathlineto{\pgfqpoint{1.320356in}{0.975420in}}%
\pgfpathlineto{\pgfqpoint{1.334209in}{0.969448in}}%
\pgfpathlineto{\pgfqpoint{1.347070in}{0.963096in}}%
\pgfpathclose%
\pgfusepath{fill}%
\end{pgfscope}%
\begin{pgfscope}%
\pgfpathrectangle{\pgfqpoint{0.407592in}{0.222904in}}{\pgfqpoint{1.411163in}{1.411163in}}%
\pgfusepath{clip}%
\pgfsetbuttcap%
\pgfsetroundjoin%
\definecolor{currentfill}{rgb}{0.749020,0.772549,0.792157}%
\pgfsetfillcolor{currentfill}%
\pgfsetfillopacity{0.160000}%
\pgfsetlinewidth{0.000000pt}%
\definecolor{currentstroke}{rgb}{0.000000,0.000000,0.000000}%
\pgfsetstrokecolor{currentstroke}%
\pgfsetdash{}{0pt}%
\pgfpathmoveto{\pgfqpoint{0.645855in}{1.250368in}}%
\pgfpathlineto{\pgfqpoint{0.635901in}{1.264701in}}%
\pgfpathlineto{\pgfqpoint{0.625953in}{1.279025in}}%
\pgfpathlineto{\pgfqpoint{0.600121in}{1.267376in}}%
\pgfpathlineto{\pgfqpoint{0.576138in}{1.255120in}}%
\pgfpathlineto{\pgfqpoint{0.587242in}{1.241197in}}%
\pgfpathlineto{\pgfqpoint{0.598349in}{1.227271in}}%
\pgfpathlineto{\pgfqpoint{0.621211in}{1.239111in}}%
\pgfpathlineto{\pgfqpoint{0.645855in}{1.250368in}}%
\pgfpathclose%
\pgfusepath{fill}%
\end{pgfscope}%
\begin{pgfscope}%
\pgfpathrectangle{\pgfqpoint{0.407592in}{0.222904in}}{\pgfqpoint{1.411163in}{1.411163in}}%
\pgfusepath{clip}%
\pgfsetbuttcap%
\pgfsetroundjoin%
\definecolor{currentfill}{rgb}{0.749020,0.772549,0.792157}%
\pgfsetfillcolor{currentfill}%
\pgfsetfillopacity{0.160000}%
\pgfsetlinewidth{0.000000pt}%
\definecolor{currentstroke}{rgb}{0.000000,0.000000,0.000000}%
\pgfsetstrokecolor{currentstroke}%
\pgfsetdash{}{0pt}%
\pgfpathmoveto{\pgfqpoint{1.174013in}{0.804561in}}%
\pgfpathlineto{\pgfqpoint{1.180927in}{0.820572in}}%
\pgfpathlineto{\pgfqpoint{1.187830in}{0.836558in}}%
\pgfpathlineto{\pgfqpoint{1.177386in}{0.838892in}}%
\pgfpathlineto{\pgfqpoint{1.166581in}{0.840928in}}%
\pgfpathlineto{\pgfqpoint{1.161187in}{0.824653in}}%
\pgfpathlineto{\pgfqpoint{1.155785in}{0.808349in}}%
\pgfpathlineto{\pgfqpoint{1.165055in}{0.806584in}}%
\pgfpathlineto{\pgfqpoint{1.174013in}{0.804561in}}%
\pgfpathclose%
\pgfusepath{fill}%
\end{pgfscope}%
\begin{pgfscope}%
\pgfpathrectangle{\pgfqpoint{0.407592in}{0.222904in}}{\pgfqpoint{1.411163in}{1.411163in}}%
\pgfusepath{clip}%
\pgfsetbuttcap%
\pgfsetroundjoin%
\definecolor{currentfill}{rgb}{0.749020,0.772549,0.792157}%
\pgfsetfillcolor{currentfill}%
\pgfsetfillopacity{0.160000}%
\pgfsetlinewidth{0.000000pt}%
\definecolor{currentstroke}{rgb}{0.000000,0.000000,0.000000}%
\pgfsetstrokecolor{currentstroke}%
\pgfsetdash{}{0pt}%
\pgfpathmoveto{\pgfqpoint{1.034960in}{0.808016in}}%
\pgfpathlineto{\pgfqpoint{1.030405in}{0.824294in}}%
\pgfpathlineto{\pgfqpoint{1.025857in}{0.840544in}}%
\pgfpathlineto{\pgfqpoint{1.014676in}{0.838447in}}%
\pgfpathlineto{\pgfqpoint{1.003802in}{0.836054in}}%
\pgfpathlineto{\pgfqpoint{1.009916in}{0.820102in}}%
\pgfpathlineto{\pgfqpoint{1.016039in}{0.804125in}}%
\pgfpathlineto{\pgfqpoint{1.025366in}{0.806199in}}%
\pgfpathlineto{\pgfqpoint{1.034960in}{0.808016in}}%
\pgfpathclose%
\pgfusepath{fill}%
\end{pgfscope}%
\begin{pgfscope}%
\pgfpathrectangle{\pgfqpoint{0.407592in}{0.222904in}}{\pgfqpoint{1.411163in}{1.411163in}}%
\pgfusepath{clip}%
\pgfsetbuttcap%
\pgfsetroundjoin%
\definecolor{currentfill}{rgb}{0.749020,0.772549,0.792157}%
\pgfsetfillcolor{currentfill}%
\pgfsetfillopacity{0.160000}%
\pgfsetlinewidth{0.000000pt}%
\definecolor{currentstroke}{rgb}{0.000000,0.000000,0.000000}%
\pgfsetstrokecolor{currentstroke}%
\pgfsetdash{}{0pt}%
\pgfpathmoveto{\pgfqpoint{1.128826in}{0.777963in}}%
\pgfpathlineto{\pgfqpoint{1.132647in}{0.794537in}}%
\pgfpathlineto{\pgfqpoint{1.136461in}{0.811081in}}%
\pgfpathlineto{\pgfqpoint{1.126489in}{0.812037in}}%
\pgfpathlineto{\pgfqpoint{1.116365in}{0.812715in}}%
\pgfpathlineto{\pgfqpoint{1.114210in}{0.796043in}}%
\pgfpathlineto{\pgfqpoint{1.112050in}{0.779340in}}%
\pgfpathlineto{\pgfqpoint{1.120503in}{0.778769in}}%
\pgfpathlineto{\pgfqpoint{1.128826in}{0.777963in}}%
\pgfpathclose%
\pgfusepath{fill}%
\end{pgfscope}%
\begin{pgfscope}%
\pgfpathrectangle{\pgfqpoint{0.407592in}{0.222904in}}{\pgfqpoint{1.411163in}{1.411163in}}%
\pgfusepath{clip}%
\pgfsetbuttcap%
\pgfsetroundjoin%
\definecolor{currentfill}{rgb}{0.749020,0.772549,0.792157}%
\pgfsetfillcolor{currentfill}%
\pgfsetfillopacity{0.160000}%
\pgfsetlinewidth{0.000000pt}%
\definecolor{currentstroke}{rgb}{0.000000,0.000000,0.000000}%
\pgfsetstrokecolor{currentstroke}%
\pgfsetdash{}{0pt}%
\pgfpathmoveto{\pgfqpoint{1.003802in}{0.836054in}}%
\pgfpathlineto{\pgfqpoint{0.997698in}{0.851984in}}%
\pgfpathlineto{\pgfqpoint{0.991602in}{0.867889in}}%
\pgfpathlineto{\pgfqpoint{0.979583in}{0.864861in}}%
\pgfpathlineto{\pgfqpoint{0.968014in}{0.861519in}}%
\pgfpathlineto{\pgfqpoint{0.975581in}{0.845980in}}%
\pgfpathlineto{\pgfqpoint{0.983157in}{0.830421in}}%
\pgfpathlineto{\pgfqpoint{0.993281in}{0.833376in}}%
\pgfpathlineto{\pgfqpoint{1.003802in}{0.836054in}}%
\pgfpathclose%
\pgfusepath{fill}%
\end{pgfscope}%
\begin{pgfscope}%
\pgfpathrectangle{\pgfqpoint{0.407592in}{0.222904in}}{\pgfqpoint{1.411163in}{1.411163in}}%
\pgfusepath{clip}%
\pgfsetbuttcap%
\pgfsetroundjoin%
\definecolor{currentfill}{rgb}{0.749020,0.772549,0.792157}%
\pgfsetfillcolor{currentfill}%
\pgfsetfillopacity{0.160000}%
\pgfsetlinewidth{0.000000pt}%
\definecolor{currentstroke}{rgb}{0.000000,0.000000,0.000000}%
\pgfsetstrokecolor{currentstroke}%
\pgfsetdash{}{0pt}%
\pgfpathmoveto{\pgfqpoint{1.077671in}{0.779245in}}%
\pgfpathlineto{\pgfqpoint{1.076429in}{0.795939in}}%
\pgfpathlineto{\pgfqpoint{1.075188in}{0.812601in}}%
\pgfpathlineto{\pgfqpoint{1.064937in}{0.811868in}}%
\pgfpathlineto{\pgfqpoint{1.054787in}{0.810856in}}%
\pgfpathlineto{\pgfqpoint{1.057711in}{0.794330in}}%
\pgfpathlineto{\pgfqpoint{1.060640in}{0.777773in}}%
\pgfpathlineto{\pgfqpoint{1.069113in}{0.778626in}}%
\pgfpathlineto{\pgfqpoint{1.077671in}{0.779245in}}%
\pgfpathclose%
\pgfusepath{fill}%
\end{pgfscope}%
\begin{pgfscope}%
\pgfpathrectangle{\pgfqpoint{0.407592in}{0.222904in}}{\pgfqpoint{1.411163in}{1.411163in}}%
\pgfusepath{clip}%
\pgfsetbuttcap%
\pgfsetroundjoin%
\definecolor{currentfill}{rgb}{0.749020,0.772549,0.792157}%
\pgfsetfillcolor{currentfill}%
\pgfsetfillopacity{0.160000}%
\pgfsetlinewidth{0.000000pt}%
\definecolor{currentstroke}{rgb}{0.000000,0.000000,0.000000}%
\pgfsetstrokecolor{currentstroke}%
\pgfsetdash{}{0pt}%
\pgfpathmoveto{\pgfqpoint{1.554624in}{1.173768in}}%
\pgfpathlineto{\pgfqpoint{1.566157in}{1.187756in}}%
\pgfpathlineto{\pgfqpoint{1.577687in}{1.201739in}}%
\pgfpathlineto{\pgfqpoint{1.558245in}{1.213043in}}%
\pgfpathlineto{\pgfqpoint{1.536980in}{1.223783in}}%
\pgfpathlineto{\pgfqpoint{1.526489in}{1.209383in}}%
\pgfpathlineto{\pgfqpoint{1.515992in}{1.194974in}}%
\pgfpathlineto{\pgfqpoint{1.536183in}{1.184641in}}%
\pgfpathlineto{\pgfqpoint{1.554624in}{1.173768in}}%
\pgfpathclose%
\pgfusepath{fill}%
\end{pgfscope}%
\begin{pgfscope}%
\pgfpathrectangle{\pgfqpoint{0.407592in}{0.222904in}}{\pgfqpoint{1.411163in}{1.411163in}}%
\pgfusepath{clip}%
\pgfsetbuttcap%
\pgfsetroundjoin%
\definecolor{currentfill}{rgb}{0.749020,0.772549,0.792157}%
\pgfsetfillcolor{currentfill}%
\pgfsetfillopacity{0.160000}%
\pgfsetlinewidth{0.000000pt}%
\definecolor{currentstroke}{rgb}{0.000000,0.000000,0.000000}%
\pgfsetstrokecolor{currentstroke}%
\pgfsetdash{}{0pt}%
\pgfpathmoveto{\pgfqpoint{0.875341in}{0.974253in}}%
\pgfpathlineto{\pgfqpoint{0.866495in}{0.989226in}}%
\pgfpathlineto{\pgfqpoint{0.857658in}{1.004184in}}%
\pgfpathlineto{\pgfqpoint{0.841397in}{0.997696in}}%
\pgfpathlineto{\pgfqpoint{0.826126in}{0.990804in}}%
\pgfpathlineto{\pgfqpoint{0.836202in}{0.976296in}}%
\pgfpathlineto{\pgfqpoint{0.846285in}{0.961778in}}%
\pgfpathlineto{\pgfqpoint{0.860352in}{0.968203in}}%
\pgfpathlineto{\pgfqpoint{0.875341in}{0.974253in}}%
\pgfpathclose%
\pgfusepath{fill}%
\end{pgfscope}%
\begin{pgfscope}%
\pgfpathrectangle{\pgfqpoint{0.407592in}{0.222904in}}{\pgfqpoint{1.411163in}{1.411163in}}%
\pgfusepath{clip}%
\pgfsetbuttcap%
\pgfsetroundjoin%
\definecolor{currentfill}{rgb}{0.749020,0.772549,0.792157}%
\pgfsetfillcolor{currentfill}%
\pgfsetfillopacity{0.160000}%
\pgfsetlinewidth{0.000000pt}%
\definecolor{currentstroke}{rgb}{0.000000,0.000000,0.000000}%
\pgfsetstrokecolor{currentstroke}%
\pgfsetdash{}{0pt}%
\pgfpathmoveto{\pgfqpoint{0.665782in}{1.221675in}}%
\pgfpathlineto{\pgfqpoint{0.655815in}{1.236026in}}%
\pgfpathlineto{\pgfqpoint{0.645855in}{1.250368in}}%
\pgfpathlineto{\pgfqpoint{0.621211in}{1.239111in}}%
\pgfpathlineto{\pgfqpoint{0.598349in}{1.227271in}}%
\pgfpathlineto{\pgfqpoint{0.609459in}{1.213340in}}%
\pgfpathlineto{\pgfqpoint{0.620573in}{1.199406in}}%
\pgfpathlineto{\pgfqpoint{0.642321in}{1.210819in}}%
\pgfpathlineto{\pgfqpoint{0.665782in}{1.221675in}}%
\pgfpathclose%
\pgfusepath{fill}%
\end{pgfscope}%
\begin{pgfscope}%
\pgfpathrectangle{\pgfqpoint{0.407592in}{0.222904in}}{\pgfqpoint{1.411163in}{1.411163in}}%
\pgfusepath{clip}%
\pgfsetbuttcap%
\pgfsetroundjoin%
\definecolor{currentfill}{rgb}{0.749020,0.772549,0.792157}%
\pgfsetfillcolor{currentfill}%
\pgfsetfillopacity{0.160000}%
\pgfsetlinewidth{0.000000pt}%
\definecolor{currentstroke}{rgb}{0.000000,0.000000,0.000000}%
\pgfsetstrokecolor{currentstroke}%
\pgfsetdash{}{0pt}%
\pgfpathmoveto{\pgfqpoint{1.531547in}{1.145779in}}%
\pgfpathlineto{\pgfqpoint{1.543088in}{1.159776in}}%
\pgfpathlineto{\pgfqpoint{1.554624in}{1.173768in}}%
\pgfpathlineto{\pgfqpoint{1.536183in}{1.184641in}}%
\pgfpathlineto{\pgfqpoint{1.515992in}{1.194974in}}%
\pgfpathlineto{\pgfqpoint{1.505488in}{1.180555in}}%
\pgfpathlineto{\pgfqpoint{1.494977in}{1.166126in}}%
\pgfpathlineto{\pgfqpoint{1.514098in}{1.156210in}}%
\pgfpathlineto{\pgfqpoint{1.531547in}{1.145779in}}%
\pgfpathclose%
\pgfusepath{fill}%
\end{pgfscope}%
\begin{pgfscope}%
\pgfpathrectangle{\pgfqpoint{0.407592in}{0.222904in}}{\pgfqpoint{1.411163in}{1.411163in}}%
\pgfusepath{clip}%
\pgfsetbuttcap%
\pgfsetroundjoin%
\definecolor{currentfill}{rgb}{0.749020,0.772549,0.792157}%
\pgfsetfillcolor{currentfill}%
\pgfsetfillopacity{0.160000}%
\pgfsetlinewidth{0.000000pt}%
\definecolor{currentstroke}{rgb}{0.000000,0.000000,0.000000}%
\pgfsetstrokecolor{currentstroke}%
\pgfsetdash{}{0pt}%
\pgfpathmoveto{\pgfqpoint{1.144952in}{0.775660in}}%
\pgfpathlineto{\pgfqpoint{1.150373in}{0.792019in}}%
\pgfpathlineto{\pgfqpoint{1.155785in}{0.808349in}}%
\pgfpathlineto{\pgfqpoint{1.146240in}{0.809851in}}%
\pgfpathlineto{\pgfqpoint{1.136461in}{0.811081in}}%
\pgfpathlineto{\pgfqpoint{1.132647in}{0.794537in}}%
\pgfpathlineto{\pgfqpoint{1.128826in}{0.777963in}}%
\pgfpathlineto{\pgfqpoint{1.136988in}{0.776925in}}%
\pgfpathlineto{\pgfqpoint{1.144952in}{0.775660in}}%
\pgfpathclose%
\pgfusepath{fill}%
\end{pgfscope}%
\begin{pgfscope}%
\pgfpathrectangle{\pgfqpoint{0.407592in}{0.222904in}}{\pgfqpoint{1.411163in}{1.411163in}}%
\pgfusepath{clip}%
\pgfsetbuttcap%
\pgfsetroundjoin%
\definecolor{currentfill}{rgb}{0.749020,0.772549,0.792157}%
\pgfsetfillcolor{currentfill}%
\pgfsetfillopacity{0.160000}%
\pgfsetlinewidth{0.000000pt}%
\definecolor{currentstroke}{rgb}{0.000000,0.000000,0.000000}%
\pgfsetstrokecolor{currentstroke}%
\pgfsetdash{}{0pt}%
\pgfpathmoveto{\pgfqpoint{1.244259in}{0.854761in}}%
\pgfpathlineto{\pgfqpoint{1.253804in}{0.869896in}}%
\pgfpathlineto{\pgfqpoint{1.263340in}{0.885016in}}%
\pgfpathlineto{\pgfqpoint{1.252289in}{0.889326in}}%
\pgfpathlineto{\pgfqpoint{1.240571in}{0.893314in}}%
\pgfpathlineto{\pgfqpoint{1.232309in}{0.877774in}}%
\pgfpathlineto{\pgfqpoint{1.224037in}{0.862214in}}%
\pgfpathlineto{\pgfqpoint{1.234447in}{0.858632in}}%
\pgfpathlineto{\pgfqpoint{1.244259in}{0.854761in}}%
\pgfpathclose%
\pgfusepath{fill}%
\end{pgfscope}%
\begin{pgfscope}%
\pgfpathrectangle{\pgfqpoint{0.407592in}{0.222904in}}{\pgfqpoint{1.411163in}{1.411163in}}%
\pgfusepath{clip}%
\pgfsetbuttcap%
\pgfsetroundjoin%
\definecolor{currentfill}{rgb}{0.749020,0.772549,0.792157}%
\pgfsetfillcolor{currentfill}%
\pgfsetfillopacity{0.160000}%
\pgfsetlinewidth{0.000000pt}%
\definecolor{currentstroke}{rgb}{0.000000,0.000000,0.000000}%
\pgfsetstrokecolor{currentstroke}%
\pgfsetdash{}{0pt}%
\pgfpathmoveto{\pgfqpoint{1.060640in}{0.777773in}}%
\pgfpathlineto{\pgfqpoint{1.057711in}{0.794330in}}%
\pgfpathlineto{\pgfqpoint{1.054787in}{0.810856in}}%
\pgfpathlineto{\pgfqpoint{1.044781in}{0.809571in}}%
\pgfpathlineto{\pgfqpoint{1.034960in}{0.808016in}}%
\pgfpathlineto{\pgfqpoint{1.039524in}{0.791712in}}%
\pgfpathlineto{\pgfqpoint{1.044094in}{0.775379in}}%
\pgfpathlineto{\pgfqpoint{1.052289in}{0.776690in}}%
\pgfpathlineto{\pgfqpoint{1.060640in}{0.777773in}}%
\pgfpathclose%
\pgfusepath{fill}%
\end{pgfscope}%
\begin{pgfscope}%
\pgfpathrectangle{\pgfqpoint{0.407592in}{0.222904in}}{\pgfqpoint{1.411163in}{1.411163in}}%
\pgfusepath{clip}%
\pgfsetbuttcap%
\pgfsetroundjoin%
\definecolor{currentfill}{rgb}{0.749020,0.772549,0.792157}%
\pgfsetfillcolor{currentfill}%
\pgfsetfillopacity{0.160000}%
\pgfsetlinewidth{0.000000pt}%
\definecolor{currentstroke}{rgb}{0.000000,0.000000,0.000000}%
\pgfsetstrokecolor{currentstroke}%
\pgfsetdash{}{0pt}%
\pgfpathmoveto{\pgfqpoint{1.325827in}{0.933935in}}%
\pgfpathlineto{\pgfqpoint{1.336452in}{0.948521in}}%
\pgfpathlineto{\pgfqpoint{1.347070in}{0.963096in}}%
\pgfpathlineto{\pgfqpoint{1.334209in}{0.969448in}}%
\pgfpathlineto{\pgfqpoint{1.320356in}{0.975420in}}%
\pgfpathlineto{\pgfqpoint{1.310877in}{0.960390in}}%
\pgfpathlineto{\pgfqpoint{1.301388in}{0.945345in}}%
\pgfpathlineto{\pgfqpoint{1.314065in}{0.939815in}}%
\pgfpathlineto{\pgfqpoint{1.325827in}{0.933935in}}%
\pgfpathclose%
\pgfusepath{fill}%
\end{pgfscope}%
\begin{pgfscope}%
\pgfpathrectangle{\pgfqpoint{0.407592in}{0.222904in}}{\pgfqpoint{1.411163in}{1.411163in}}%
\pgfusepath{clip}%
\pgfsetbuttcap%
\pgfsetroundjoin%
\definecolor{currentfill}{rgb}{0.749020,0.772549,0.792157}%
\pgfsetfillcolor{currentfill}%
\pgfsetfillopacity{0.160000}%
\pgfsetlinewidth{0.000000pt}%
\definecolor{currentstroke}{rgb}{0.000000,0.000000,0.000000}%
\pgfsetstrokecolor{currentstroke}%
\pgfsetdash{}{0pt}%
\pgfpathmoveto{\pgfqpoint{0.685736in}{1.192946in}}%
\pgfpathlineto{\pgfqpoint{0.675756in}{1.207315in}}%
\pgfpathlineto{\pgfqpoint{0.665782in}{1.221675in}}%
\pgfpathlineto{\pgfqpoint{0.642321in}{1.210819in}}%
\pgfpathlineto{\pgfqpoint{0.620573in}{1.199406in}}%
\pgfpathlineto{\pgfqpoint{0.631689in}{1.185467in}}%
\pgfpathlineto{\pgfqpoint{0.642809in}{1.171525in}}%
\pgfpathlineto{\pgfqpoint{0.663451in}{1.182502in}}%
\pgfpathlineto{\pgfqpoint{0.685736in}{1.192946in}}%
\pgfpathclose%
\pgfusepath{fill}%
\end{pgfscope}%
\begin{pgfscope}%
\pgfpathrectangle{\pgfqpoint{0.407592in}{0.222904in}}{\pgfqpoint{1.411163in}{1.411163in}}%
\pgfusepath{clip}%
\pgfsetbuttcap%
\pgfsetroundjoin%
\definecolor{currentfill}{rgb}{0.749020,0.772549,0.792157}%
\pgfsetfillcolor{currentfill}%
\pgfsetfillopacity{0.160000}%
\pgfsetlinewidth{0.000000pt}%
\definecolor{currentstroke}{rgb}{0.000000,0.000000,0.000000}%
\pgfsetstrokecolor{currentstroke}%
\pgfsetdash{}{0pt}%
\pgfpathmoveto{\pgfqpoint{0.968014in}{0.861519in}}%
\pgfpathlineto{\pgfqpoint{0.960457in}{0.877040in}}%
\pgfpathlineto{\pgfqpoint{0.952909in}{0.892541in}}%
\pgfpathlineto{\pgfqpoint{0.940448in}{0.888487in}}%
\pgfpathlineto{\pgfqpoint{0.928596in}{0.884114in}}%
\pgfpathlineto{\pgfqpoint{0.937502in}{0.869040in}}%
\pgfpathlineto{\pgfqpoint{0.946416in}{0.853951in}}%
\pgfpathlineto{\pgfqpoint{0.956943in}{0.857878in}}%
\pgfpathlineto{\pgfqpoint{0.968014in}{0.861519in}}%
\pgfpathclose%
\pgfusepath{fill}%
\end{pgfscope}%
\begin{pgfscope}%
\pgfpathrectangle{\pgfqpoint{0.407592in}{0.222904in}}{\pgfqpoint{1.411163in}{1.411163in}}%
\pgfusepath{clip}%
\pgfsetbuttcap%
\pgfsetroundjoin%
\definecolor{currentfill}{rgb}{0.749020,0.772549,0.792157}%
\pgfsetfillcolor{currentfill}%
\pgfsetfillopacity{0.160000}%
\pgfsetlinewidth{0.000000pt}%
\definecolor{currentstroke}{rgb}{0.000000,0.000000,0.000000}%
\pgfsetstrokecolor{currentstroke}%
\pgfsetdash{}{0pt}%
\pgfpathmoveto{\pgfqpoint{0.893058in}{0.944265in}}%
\pgfpathlineto{\pgfqpoint{0.884195in}{0.959266in}}%
\pgfpathlineto{\pgfqpoint{0.875341in}{0.974253in}}%
\pgfpathlineto{\pgfqpoint{0.860352in}{0.968203in}}%
\pgfpathlineto{\pgfqpoint{0.846285in}{0.961778in}}%
\pgfpathlineto{\pgfqpoint{0.856374in}{0.947251in}}%
\pgfpathlineto{\pgfqpoint{0.866469in}{0.932715in}}%
\pgfpathlineto{\pgfqpoint{0.879338in}{0.938663in}}%
\pgfpathlineto{\pgfqpoint{0.893058in}{0.944265in}}%
\pgfpathclose%
\pgfusepath{fill}%
\end{pgfscope}%
\begin{pgfscope}%
\pgfpathrectangle{\pgfqpoint{0.407592in}{0.222904in}}{\pgfqpoint{1.411163in}{1.411163in}}%
\pgfusepath{clip}%
\pgfsetbuttcap%
\pgfsetroundjoin%
\definecolor{currentfill}{rgb}{0.749020,0.772549,0.792157}%
\pgfsetfillcolor{currentfill}%
\pgfsetfillopacity{0.160000}%
\pgfsetlinewidth{0.000000pt}%
\definecolor{currentstroke}{rgb}{0.000000,0.000000,0.000000}%
\pgfsetstrokecolor{currentstroke}%
\pgfsetdash{}{0pt}%
\pgfpathmoveto{\pgfqpoint{1.508456in}{1.117773in}}%
\pgfpathlineto{\pgfqpoint{1.520003in}{1.131778in}}%
\pgfpathlineto{\pgfqpoint{1.531547in}{1.145779in}}%
\pgfpathlineto{\pgfqpoint{1.514098in}{1.156210in}}%
\pgfpathlineto{\pgfqpoint{1.494977in}{1.166126in}}%
\pgfpathlineto{\pgfqpoint{1.484458in}{1.151687in}}%
\pgfpathlineto{\pgfqpoint{1.473933in}{1.137239in}}%
\pgfpathlineto{\pgfqpoint{1.491992in}{1.127750in}}%
\pgfpathlineto{\pgfqpoint{1.508456in}{1.117773in}}%
\pgfpathclose%
\pgfusepath{fill}%
\end{pgfscope}%
\begin{pgfscope}%
\pgfpathrectangle{\pgfqpoint{0.407592in}{0.222904in}}{\pgfqpoint{1.411163in}{1.411163in}}%
\pgfusepath{clip}%
\pgfsetbuttcap%
\pgfsetroundjoin%
\definecolor{currentfill}{rgb}{0.749020,0.772549,0.792157}%
\pgfsetfillcolor{currentfill}%
\pgfsetfillopacity{0.160000}%
\pgfsetlinewidth{0.000000pt}%
\definecolor{currentstroke}{rgb}{0.000000,0.000000,0.000000}%
\pgfsetstrokecolor{currentstroke}%
\pgfsetdash{}{0pt}%
\pgfpathmoveto{\pgfqpoint{1.190842in}{0.799776in}}%
\pgfpathlineto{\pgfqpoint{1.199156in}{0.815416in}}%
\pgfpathlineto{\pgfqpoint{1.207460in}{0.831035in}}%
\pgfpathlineto{\pgfqpoint{1.197869in}{0.833936in}}%
\pgfpathlineto{\pgfqpoint{1.187830in}{0.836558in}}%
\pgfpathlineto{\pgfqpoint{1.180927in}{0.820572in}}%
\pgfpathlineto{\pgfqpoint{1.174013in}{0.804561in}}%
\pgfpathlineto{\pgfqpoint{1.182621in}{0.802289in}}%
\pgfpathlineto{\pgfqpoint{1.190842in}{0.799776in}}%
\pgfpathclose%
\pgfusepath{fill}%
\end{pgfscope}%
\begin{pgfscope}%
\pgfpathrectangle{\pgfqpoint{0.407592in}{0.222904in}}{\pgfqpoint{1.411163in}{1.411163in}}%
\pgfusepath{clip}%
\pgfsetbuttcap%
\pgfsetroundjoin%
\definecolor{currentfill}{rgb}{0.749020,0.772549,0.792157}%
\pgfsetfillcolor{currentfill}%
\pgfsetfillopacity{0.160000}%
\pgfsetlinewidth{0.000000pt}%
\definecolor{currentstroke}{rgb}{0.000000,0.000000,0.000000}%
\pgfsetstrokecolor{currentstroke}%
\pgfsetdash{}{0pt}%
\pgfpathmoveto{\pgfqpoint{0.705714in}{1.164179in}}%
\pgfpathlineto{\pgfqpoint{0.695722in}{1.178567in}}%
\pgfpathlineto{\pgfqpoint{0.685736in}{1.192946in}}%
\pgfpathlineto{\pgfqpoint{0.663451in}{1.182502in}}%
\pgfpathlineto{\pgfqpoint{0.642809in}{1.171525in}}%
\pgfpathlineto{\pgfqpoint{0.653932in}{1.157578in}}%
\pgfpathlineto{\pgfqpoint{0.665058in}{1.143628in}}%
\pgfpathlineto{\pgfqpoint{0.684600in}{1.154157in}}%
\pgfpathlineto{\pgfqpoint{0.705714in}{1.164179in}}%
\pgfpathclose%
\pgfusepath{fill}%
\end{pgfscope}%
\begin{pgfscope}%
\pgfpathrectangle{\pgfqpoint{0.407592in}{0.222904in}}{\pgfqpoint{1.411163in}{1.411163in}}%
\pgfusepath{clip}%
\pgfsetbuttcap%
\pgfsetroundjoin%
\definecolor{currentfill}{rgb}{0.749020,0.772549,0.792157}%
\pgfsetfillcolor{currentfill}%
\pgfsetfillopacity{0.160000}%
\pgfsetlinewidth{0.000000pt}%
\definecolor{currentstroke}{rgb}{0.000000,0.000000,0.000000}%
\pgfsetstrokecolor{currentstroke}%
\pgfsetdash{}{0pt}%
\pgfpathmoveto{\pgfqpoint{1.016039in}{0.804125in}}%
\pgfpathlineto{\pgfqpoint{1.009916in}{0.820102in}}%
\pgfpathlineto{\pgfqpoint{1.003802in}{0.836054in}}%
\pgfpathlineto{\pgfqpoint{0.993281in}{0.833376in}}%
\pgfpathlineto{\pgfqpoint{0.983157in}{0.830421in}}%
\pgfpathlineto{\pgfqpoint{0.990743in}{0.814842in}}%
\pgfpathlineto{\pgfqpoint{0.998338in}{0.799244in}}%
\pgfpathlineto{\pgfqpoint{1.007017in}{0.801804in}}%
\pgfpathlineto{\pgfqpoint{1.016039in}{0.804125in}}%
\pgfpathclose%
\pgfusepath{fill}%
\end{pgfscope}%
\begin{pgfscope}%
\pgfpathrectangle{\pgfqpoint{0.407592in}{0.222904in}}{\pgfqpoint{1.411163in}{1.411163in}}%
\pgfusepath{clip}%
\pgfsetbuttcap%
\pgfsetroundjoin%
\definecolor{currentfill}{rgb}{0.749020,0.772549,0.792157}%
\pgfsetfillcolor{currentfill}%
\pgfsetfillopacity{0.160000}%
\pgfsetlinewidth{0.000000pt}%
\definecolor{currentstroke}{rgb}{0.000000,0.000000,0.000000}%
\pgfsetstrokecolor{currentstroke}%
\pgfsetdash{}{0pt}%
\pgfpathmoveto{\pgfqpoint{1.160155in}{0.772468in}}%
\pgfpathlineto{\pgfqpoint{1.167090in}{0.788527in}}%
\pgfpathlineto{\pgfqpoint{1.174013in}{0.804561in}}%
\pgfpathlineto{\pgfqpoint{1.165055in}{0.806584in}}%
\pgfpathlineto{\pgfqpoint{1.155785in}{0.808349in}}%
\pgfpathlineto{\pgfqpoint{1.150373in}{0.792019in}}%
\pgfpathlineto{\pgfqpoint{1.144952in}{0.775660in}}%
\pgfpathlineto{\pgfqpoint{1.152685in}{0.774172in}}%
\pgfpathlineto{\pgfqpoint{1.160155in}{0.772468in}}%
\pgfpathclose%
\pgfusepath{fill}%
\end{pgfscope}%
\begin{pgfscope}%
\pgfpathrectangle{\pgfqpoint{0.407592in}{0.222904in}}{\pgfqpoint{1.411163in}{1.411163in}}%
\pgfusepath{clip}%
\pgfsetbuttcap%
\pgfsetroundjoin%
\definecolor{currentfill}{rgb}{0.749020,0.772549,0.792157}%
\pgfsetfillcolor{currentfill}%
\pgfsetfillopacity{0.160000}%
\pgfsetlinewidth{0.000000pt}%
\definecolor{currentstroke}{rgb}{0.000000,0.000000,0.000000}%
\pgfsetstrokecolor{currentstroke}%
\pgfsetdash{}{0pt}%
\pgfpathmoveto{\pgfqpoint{1.304554in}{0.904734in}}%
\pgfpathlineto{\pgfqpoint{1.315194in}{0.919339in}}%
\pgfpathlineto{\pgfqpoint{1.325827in}{0.933935in}}%
\pgfpathlineto{\pgfqpoint{1.314065in}{0.939815in}}%
\pgfpathlineto{\pgfqpoint{1.301388in}{0.945345in}}%
\pgfpathlineto{\pgfqpoint{1.291891in}{0.930285in}}%
\pgfpathlineto{\pgfqpoint{1.282383in}{0.915211in}}%
\pgfpathlineto{\pgfqpoint{1.293888in}{0.910132in}}%
\pgfpathlineto{\pgfqpoint{1.304554in}{0.904734in}}%
\pgfpathclose%
\pgfusepath{fill}%
\end{pgfscope}%
\begin{pgfscope}%
\pgfpathrectangle{\pgfqpoint{0.407592in}{0.222904in}}{\pgfqpoint{1.411163in}{1.411163in}}%
\pgfusepath{clip}%
\pgfsetbuttcap%
\pgfsetroundjoin%
\definecolor{currentfill}{rgb}{0.749020,0.772549,0.792157}%
\pgfsetfillcolor{currentfill}%
\pgfsetfillopacity{0.160000}%
\pgfsetlinewidth{0.000000pt}%
\definecolor{currentstroke}{rgb}{0.000000,0.000000,0.000000}%
\pgfsetstrokecolor{currentstroke}%
\pgfsetdash{}{0pt}%
\pgfpathmoveto{\pgfqpoint{1.485349in}{1.089748in}}%
\pgfpathlineto{\pgfqpoint{1.496904in}{1.103763in}}%
\pgfpathlineto{\pgfqpoint{1.508456in}{1.117773in}}%
\pgfpathlineto{\pgfqpoint{1.491992in}{1.127750in}}%
\pgfpathlineto{\pgfqpoint{1.473933in}{1.137239in}}%
\pgfpathlineto{\pgfqpoint{1.463400in}{1.122781in}}%
\pgfpathlineto{\pgfqpoint{1.452860in}{1.108313in}}%
\pgfpathlineto{\pgfqpoint{1.469863in}{1.099262in}}%
\pgfpathlineto{\pgfqpoint{1.485349in}{1.089748in}}%
\pgfpathclose%
\pgfusepath{fill}%
\end{pgfscope}%
\begin{pgfscope}%
\pgfpathrectangle{\pgfqpoint{0.407592in}{0.222904in}}{\pgfqpoint{1.411163in}{1.411163in}}%
\pgfusepath{clip}%
\pgfsetbuttcap%
\pgfsetroundjoin%
\definecolor{currentfill}{rgb}{0.749020,0.772549,0.792157}%
\pgfsetfillcolor{currentfill}%
\pgfsetfillopacity{0.160000}%
\pgfsetlinewidth{0.000000pt}%
\definecolor{currentstroke}{rgb}{0.000000,0.000000,0.000000}%
\pgfsetstrokecolor{currentstroke}%
\pgfsetdash{}{0pt}%
\pgfpathmoveto{\pgfqpoint{1.044094in}{0.775379in}}%
\pgfpathlineto{\pgfqpoint{1.039524in}{0.791712in}}%
\pgfpathlineto{\pgfqpoint{1.034960in}{0.808016in}}%
\pgfpathlineto{\pgfqpoint{1.025366in}{0.806199in}}%
\pgfpathlineto{\pgfqpoint{1.016039in}{0.804125in}}%
\pgfpathlineto{\pgfqpoint{1.022170in}{0.788125in}}%
\pgfpathlineto{\pgfqpoint{1.028311in}{0.772101in}}%
\pgfpathlineto{\pgfqpoint{1.036091in}{0.773848in}}%
\pgfpathlineto{\pgfqpoint{1.044094in}{0.775379in}}%
\pgfpathclose%
\pgfusepath{fill}%
\end{pgfscope}%
\begin{pgfscope}%
\pgfpathrectangle{\pgfqpoint{0.407592in}{0.222904in}}{\pgfqpoint{1.411163in}{1.411163in}}%
\pgfusepath{clip}%
\pgfsetbuttcap%
\pgfsetroundjoin%
\definecolor{currentfill}{rgb}{0.749020,0.772549,0.792157}%
\pgfsetfillcolor{currentfill}%
\pgfsetfillopacity{0.160000}%
\pgfsetlinewidth{0.000000pt}%
\definecolor{currentstroke}{rgb}{0.000000,0.000000,0.000000}%
\pgfsetstrokecolor{currentstroke}%
\pgfsetdash{}{0pt}%
\pgfpathmoveto{\pgfqpoint{1.107720in}{0.745838in}}%
\pgfpathlineto{\pgfqpoint{1.109887in}{0.762605in}}%
\pgfpathlineto{\pgfqpoint{1.112050in}{0.779340in}}%
\pgfpathlineto{\pgfqpoint{1.103506in}{0.779674in}}%
\pgfpathlineto{\pgfqpoint{1.094904in}{0.779770in}}%
\pgfpathlineto{\pgfqpoint{1.094440in}{0.762994in}}%
\pgfpathlineto{\pgfqpoint{1.093976in}{0.746185in}}%
\pgfpathlineto{\pgfqpoint{1.100871in}{0.746108in}}%
\pgfpathlineto{\pgfqpoint{1.107720in}{0.745838in}}%
\pgfpathclose%
\pgfusepath{fill}%
\end{pgfscope}%
\begin{pgfscope}%
\pgfpathrectangle{\pgfqpoint{0.407592in}{0.222904in}}{\pgfqpoint{1.411163in}{1.411163in}}%
\pgfusepath{clip}%
\pgfsetbuttcap%
\pgfsetroundjoin%
\definecolor{currentfill}{rgb}{0.749020,0.772549,0.792157}%
\pgfsetfillcolor{currentfill}%
\pgfsetfillopacity{0.160000}%
\pgfsetlinewidth{0.000000pt}%
\definecolor{currentstroke}{rgb}{0.000000,0.000000,0.000000}%
\pgfsetstrokecolor{currentstroke}%
\pgfsetdash{}{0pt}%
\pgfpathmoveto{\pgfqpoint{1.093976in}{0.746185in}}%
\pgfpathlineto{\pgfqpoint{1.094440in}{0.762994in}}%
\pgfpathlineto{\pgfqpoint{1.094904in}{0.779770in}}%
\pgfpathlineto{\pgfqpoint{1.086280in}{0.779626in}}%
\pgfpathlineto{\pgfqpoint{1.077671in}{0.779245in}}%
\pgfpathlineto{\pgfqpoint{1.078916in}{0.762518in}}%
\pgfpathlineto{\pgfqpoint{1.080164in}{0.745760in}}%
\pgfpathlineto{\pgfqpoint{1.087064in}{0.746069in}}%
\pgfpathlineto{\pgfqpoint{1.093976in}{0.746185in}}%
\pgfpathclose%
\pgfusepath{fill}%
\end{pgfscope}%
\begin{pgfscope}%
\pgfpathrectangle{\pgfqpoint{0.407592in}{0.222904in}}{\pgfqpoint{1.411163in}{1.411163in}}%
\pgfusepath{clip}%
\pgfsetbuttcap%
\pgfsetroundjoin%
\definecolor{currentfill}{rgb}{0.749020,0.772549,0.792157}%
\pgfsetfillcolor{currentfill}%
\pgfsetfillopacity{0.160000}%
\pgfsetlinewidth{0.000000pt}%
\definecolor{currentstroke}{rgb}{0.000000,0.000000,0.000000}%
\pgfsetstrokecolor{currentstroke}%
\pgfsetdash{}{0pt}%
\pgfpathmoveto{\pgfqpoint{0.725719in}{1.135376in}}%
\pgfpathlineto{\pgfqpoint{0.715713in}{1.149782in}}%
\pgfpathlineto{\pgfqpoint{0.705714in}{1.164179in}}%
\pgfpathlineto{\pgfqpoint{0.684600in}{1.154157in}}%
\pgfpathlineto{\pgfqpoint{0.665058in}{1.143628in}}%
\pgfpathlineto{\pgfqpoint{0.676188in}{1.129674in}}%
\pgfpathlineto{\pgfqpoint{0.687320in}{1.115715in}}%
\pgfpathlineto{\pgfqpoint{0.705769in}{1.125787in}}%
\pgfpathlineto{\pgfqpoint{0.725719in}{1.135376in}}%
\pgfpathclose%
\pgfusepath{fill}%
\end{pgfscope}%
\begin{pgfscope}%
\pgfpathrectangle{\pgfqpoint{0.407592in}{0.222904in}}{\pgfqpoint{1.411163in}{1.411163in}}%
\pgfusepath{clip}%
\pgfsetbuttcap%
\pgfsetroundjoin%
\definecolor{currentfill}{rgb}{0.749020,0.772549,0.792157}%
\pgfsetfillcolor{currentfill}%
\pgfsetfillopacity{0.160000}%
\pgfsetlinewidth{0.000000pt}%
\definecolor{currentstroke}{rgb}{0.000000,0.000000,0.000000}%
\pgfsetstrokecolor{currentstroke}%
\pgfsetdash{}{0pt}%
\pgfpathmoveto{\pgfqpoint{1.225140in}{0.824446in}}%
\pgfpathlineto{\pgfqpoint{1.234704in}{0.839611in}}%
\pgfpathlineto{\pgfqpoint{1.244259in}{0.854761in}}%
\pgfpathlineto{\pgfqpoint{1.234447in}{0.858632in}}%
\pgfpathlineto{\pgfqpoint{1.224037in}{0.862214in}}%
\pgfpathlineto{\pgfqpoint{1.215754in}{0.846635in}}%
\pgfpathlineto{\pgfqpoint{1.207460in}{0.831035in}}%
\pgfpathlineto{\pgfqpoint{1.216564in}{0.827868in}}%
\pgfpathlineto{\pgfqpoint{1.225140in}{0.824446in}}%
\pgfpathclose%
\pgfusepath{fill}%
\end{pgfscope}%
\begin{pgfscope}%
\pgfpathrectangle{\pgfqpoint{0.407592in}{0.222904in}}{\pgfqpoint{1.411163in}{1.411163in}}%
\pgfusepath{clip}%
\pgfsetbuttcap%
\pgfsetroundjoin%
\definecolor{currentfill}{rgb}{0.749020,0.772549,0.792157}%
\pgfsetfillcolor{currentfill}%
\pgfsetfillopacity{0.160000}%
\pgfsetlinewidth{0.000000pt}%
\definecolor{currentstroke}{rgb}{0.000000,0.000000,0.000000}%
\pgfsetstrokecolor{currentstroke}%
\pgfsetdash{}{0pt}%
\pgfpathmoveto{\pgfqpoint{0.910810in}{0.914218in}}%
\pgfpathlineto{\pgfqpoint{0.901930in}{0.929249in}}%
\pgfpathlineto{\pgfqpoint{0.893058in}{0.944265in}}%
\pgfpathlineto{\pgfqpoint{0.879338in}{0.938663in}}%
\pgfpathlineto{\pgfqpoint{0.866469in}{0.932715in}}%
\pgfpathlineto{\pgfqpoint{0.876572in}{0.918169in}}%
\pgfpathlineto{\pgfqpoint{0.886680in}{0.903614in}}%
\pgfpathlineto{\pgfqpoint{0.898355in}{0.909074in}}%
\pgfpathlineto{\pgfqpoint{0.910810in}{0.914218in}}%
\pgfpathclose%
\pgfusepath{fill}%
\end{pgfscope}%
\begin{pgfscope}%
\pgfpathrectangle{\pgfqpoint{0.407592in}{0.222904in}}{\pgfqpoint{1.411163in}{1.411163in}}%
\pgfusepath{clip}%
\pgfsetbuttcap%
\pgfsetroundjoin%
\definecolor{currentfill}{rgb}{0.749020,0.772549,0.792157}%
\pgfsetfillcolor{currentfill}%
\pgfsetfillopacity{0.160000}%
\pgfsetlinewidth{0.000000pt}%
\definecolor{currentstroke}{rgb}{0.000000,0.000000,0.000000}%
\pgfsetstrokecolor{currentstroke}%
\pgfsetdash{}{0pt}%
\pgfpathmoveto{\pgfqpoint{1.121164in}{0.744723in}}%
\pgfpathlineto{\pgfqpoint{1.124999in}{0.761358in}}%
\pgfpathlineto{\pgfqpoint{1.128826in}{0.777963in}}%
\pgfpathlineto{\pgfqpoint{1.120503in}{0.778769in}}%
\pgfpathlineto{\pgfqpoint{1.112050in}{0.779340in}}%
\pgfpathlineto{\pgfqpoint{1.109887in}{0.762605in}}%
\pgfpathlineto{\pgfqpoint{1.107720in}{0.745838in}}%
\pgfpathlineto{\pgfqpoint{1.114494in}{0.745375in}}%
\pgfpathlineto{\pgfqpoint{1.121164in}{0.744723in}}%
\pgfpathclose%
\pgfusepath{fill}%
\end{pgfscope}%
\begin{pgfscope}%
\pgfpathrectangle{\pgfqpoint{0.407592in}{0.222904in}}{\pgfqpoint{1.411163in}{1.411163in}}%
\pgfusepath{clip}%
\pgfsetbuttcap%
\pgfsetroundjoin%
\definecolor{currentfill}{rgb}{0.749020,0.772549,0.792157}%
\pgfsetfillcolor{currentfill}%
\pgfsetfillopacity{0.160000}%
\pgfsetlinewidth{0.000000pt}%
\definecolor{currentstroke}{rgb}{0.000000,0.000000,0.000000}%
\pgfsetstrokecolor{currentstroke}%
\pgfsetdash{}{0pt}%
\pgfpathmoveto{\pgfqpoint{1.080164in}{0.745760in}}%
\pgfpathlineto{\pgfqpoint{1.078916in}{0.762518in}}%
\pgfpathlineto{\pgfqpoint{1.077671in}{0.779245in}}%
\pgfpathlineto{\pgfqpoint{1.069113in}{0.778626in}}%
\pgfpathlineto{\pgfqpoint{1.060640in}{0.777773in}}%
\pgfpathlineto{\pgfqpoint{1.063575in}{0.761187in}}%
\pgfpathlineto{\pgfqpoint{1.066515in}{0.744570in}}%
\pgfpathlineto{\pgfqpoint{1.073304in}{0.745260in}}%
\pgfpathlineto{\pgfqpoint{1.080164in}{0.745760in}}%
\pgfpathclose%
\pgfusepath{fill}%
\end{pgfscope}%
\begin{pgfscope}%
\pgfpathrectangle{\pgfqpoint{0.407592in}{0.222904in}}{\pgfqpoint{1.411163in}{1.411163in}}%
\pgfusepath{clip}%
\pgfsetbuttcap%
\pgfsetroundjoin%
\definecolor{currentfill}{rgb}{0.749020,0.772549,0.792157}%
\pgfsetfillcolor{currentfill}%
\pgfsetfillopacity{0.160000}%
\pgfsetlinewidth{0.000000pt}%
\definecolor{currentstroke}{rgb}{0.000000,0.000000,0.000000}%
\pgfsetstrokecolor{currentstroke}%
\pgfsetdash{}{0pt}%
\pgfpathmoveto{\pgfqpoint{0.983157in}{0.830421in}}%
\pgfpathlineto{\pgfqpoint{0.975581in}{0.845980in}}%
\pgfpathlineto{\pgfqpoint{0.968014in}{0.861519in}}%
\pgfpathlineto{\pgfqpoint{0.956943in}{0.857878in}}%
\pgfpathlineto{\pgfqpoint{0.946416in}{0.853951in}}%
\pgfpathlineto{\pgfqpoint{0.955340in}{0.838848in}}%
\pgfpathlineto{\pgfqpoint{0.964271in}{0.823730in}}%
\pgfpathlineto{\pgfqpoint{0.973474in}{0.827201in}}%
\pgfpathlineto{\pgfqpoint{0.983157in}{0.830421in}}%
\pgfpathclose%
\pgfusepath{fill}%
\end{pgfscope}%
\begin{pgfscope}%
\pgfpathrectangle{\pgfqpoint{0.407592in}{0.222904in}}{\pgfqpoint{1.411163in}{1.411163in}}%
\pgfusepath{clip}%
\pgfsetbuttcap%
\pgfsetroundjoin%
\definecolor{currentfill}{rgb}{0.749020,0.772549,0.792157}%
\pgfsetfillcolor{currentfill}%
\pgfsetfillopacity{0.160000}%
\pgfsetlinewidth{0.000000pt}%
\definecolor{currentstroke}{rgb}{0.000000,0.000000,0.000000}%
\pgfsetstrokecolor{currentstroke}%
\pgfsetdash{}{0pt}%
\pgfpathmoveto{\pgfqpoint{1.462227in}{1.061705in}}%
\pgfpathlineto{\pgfqpoint{1.473790in}{1.075729in}}%
\pgfpathlineto{\pgfqpoint{1.485349in}{1.089748in}}%
\pgfpathlineto{\pgfqpoint{1.469863in}{1.099262in}}%
\pgfpathlineto{\pgfqpoint{1.452860in}{1.108313in}}%
\pgfpathlineto{\pgfqpoint{1.442313in}{1.093835in}}%
\pgfpathlineto{\pgfqpoint{1.431759in}{1.079348in}}%
\pgfpathlineto{\pgfqpoint{1.447711in}{1.070745in}}%
\pgfpathlineto{\pgfqpoint{1.462227in}{1.061705in}}%
\pgfpathclose%
\pgfusepath{fill}%
\end{pgfscope}%
\begin{pgfscope}%
\pgfpathrectangle{\pgfqpoint{0.407592in}{0.222904in}}{\pgfqpoint{1.411163in}{1.411163in}}%
\pgfusepath{clip}%
\pgfsetbuttcap%
\pgfsetroundjoin%
\definecolor{currentfill}{rgb}{0.749020,0.772549,0.792157}%
\pgfsetfillcolor{currentfill}%
\pgfsetfillopacity{0.160000}%
\pgfsetlinewidth{0.000000pt}%
\definecolor{currentstroke}{rgb}{0.000000,0.000000,0.000000}%
\pgfsetstrokecolor{currentstroke}%
\pgfsetdash{}{0pt}%
\pgfpathmoveto{\pgfqpoint{0.745749in}{1.106536in}}%
\pgfpathlineto{\pgfqpoint{0.735730in}{1.120961in}}%
\pgfpathlineto{\pgfqpoint{0.725719in}{1.135376in}}%
\pgfpathlineto{\pgfqpoint{0.705769in}{1.125787in}}%
\pgfpathlineto{\pgfqpoint{0.687320in}{1.115715in}}%
\pgfpathlineto{\pgfqpoint{0.698456in}{1.101753in}}%
\pgfpathlineto{\pgfqpoint{0.709595in}{1.087786in}}%
\pgfpathlineto{\pgfqpoint{0.726958in}{1.097389in}}%
\pgfpathlineto{\pgfqpoint{0.745749in}{1.106536in}}%
\pgfpathclose%
\pgfusepath{fill}%
\end{pgfscope}%
\begin{pgfscope}%
\pgfpathrectangle{\pgfqpoint{0.407592in}{0.222904in}}{\pgfqpoint{1.411163in}{1.411163in}}%
\pgfusepath{clip}%
\pgfsetbuttcap%
\pgfsetroundjoin%
\definecolor{currentfill}{rgb}{0.749020,0.772549,0.792157}%
\pgfsetfillcolor{currentfill}%
\pgfsetfillopacity{0.160000}%
\pgfsetlinewidth{0.000000pt}%
\definecolor{currentstroke}{rgb}{0.000000,0.000000,0.000000}%
\pgfsetstrokecolor{currentstroke}%
\pgfsetdash{}{0pt}%
\pgfpathmoveto{\pgfqpoint{1.134082in}{0.742860in}}%
\pgfpathlineto{\pgfqpoint{1.139521in}{0.759274in}}%
\pgfpathlineto{\pgfqpoint{1.144952in}{0.775660in}}%
\pgfpathlineto{\pgfqpoint{1.136988in}{0.776925in}}%
\pgfpathlineto{\pgfqpoint{1.128826in}{0.777963in}}%
\pgfpathlineto{\pgfqpoint{1.124999in}{0.761358in}}%
\pgfpathlineto{\pgfqpoint{1.121164in}{0.744723in}}%
\pgfpathlineto{\pgfqpoint{1.127702in}{0.743883in}}%
\pgfpathlineto{\pgfqpoint{1.134082in}{0.742860in}}%
\pgfpathclose%
\pgfusepath{fill}%
\end{pgfscope}%
\begin{pgfscope}%
\pgfpathrectangle{\pgfqpoint{0.407592in}{0.222904in}}{\pgfqpoint{1.411163in}{1.411163in}}%
\pgfusepath{clip}%
\pgfsetbuttcap%
\pgfsetroundjoin%
\definecolor{currentfill}{rgb}{0.749020,0.772549,0.792157}%
\pgfsetfillcolor{currentfill}%
\pgfsetfillopacity{0.160000}%
\pgfsetlinewidth{0.000000pt}%
\definecolor{currentstroke}{rgb}{0.000000,0.000000,0.000000}%
\pgfsetstrokecolor{currentstroke}%
\pgfsetdash{}{0pt}%
\pgfpathmoveto{\pgfqpoint{1.283252in}{0.875494in}}%
\pgfpathlineto{\pgfqpoint{1.293907in}{0.890119in}}%
\pgfpathlineto{\pgfqpoint{1.304554in}{0.904734in}}%
\pgfpathlineto{\pgfqpoint{1.293888in}{0.910132in}}%
\pgfpathlineto{\pgfqpoint{1.282383in}{0.915211in}}%
\pgfpathlineto{\pgfqpoint{1.272866in}{0.900121in}}%
\pgfpathlineto{\pgfqpoint{1.263340in}{0.885016in}}%
\pgfpathlineto{\pgfqpoint{1.273676in}{0.880399in}}%
\pgfpathlineto{\pgfqpoint{1.283252in}{0.875494in}}%
\pgfpathclose%
\pgfusepath{fill}%
\end{pgfscope}%
\begin{pgfscope}%
\pgfpathrectangle{\pgfqpoint{0.407592in}{0.222904in}}{\pgfqpoint{1.411163in}{1.411163in}}%
\pgfusepath{clip}%
\pgfsetbuttcap%
\pgfsetroundjoin%
\definecolor{currentfill}{rgb}{0.749020,0.772549,0.792157}%
\pgfsetfillcolor{currentfill}%
\pgfsetfillopacity{0.160000}%
\pgfsetlinewidth{0.000000pt}%
\definecolor{currentstroke}{rgb}{0.000000,0.000000,0.000000}%
\pgfsetstrokecolor{currentstroke}%
\pgfsetdash{}{0pt}%
\pgfpathmoveto{\pgfqpoint{1.066515in}{0.744570in}}%
\pgfpathlineto{\pgfqpoint{1.063575in}{0.761187in}}%
\pgfpathlineto{\pgfqpoint{1.060640in}{0.777773in}}%
\pgfpathlineto{\pgfqpoint{1.052289in}{0.776690in}}%
\pgfpathlineto{\pgfqpoint{1.044094in}{0.775379in}}%
\pgfpathlineto{\pgfqpoint{1.048673in}{0.759020in}}%
\pgfpathlineto{\pgfqpoint{1.053259in}{0.742632in}}%
\pgfpathlineto{\pgfqpoint{1.059824in}{0.743693in}}%
\pgfpathlineto{\pgfqpoint{1.066515in}{0.744570in}}%
\pgfpathclose%
\pgfusepath{fill}%
\end{pgfscope}%
\begin{pgfscope}%
\pgfpathrectangle{\pgfqpoint{0.407592in}{0.222904in}}{\pgfqpoint{1.411163in}{1.411163in}}%
\pgfusepath{clip}%
\pgfsetbuttcap%
\pgfsetroundjoin%
\definecolor{currentfill}{rgb}{0.749020,0.772549,0.792157}%
\pgfsetfillcolor{currentfill}%
\pgfsetfillopacity{0.160000}%
\pgfsetlinewidth{0.000000pt}%
\definecolor{currentstroke}{rgb}{0.000000,0.000000,0.000000}%
\pgfsetstrokecolor{currentstroke}%
\pgfsetdash{}{0pt}%
\pgfpathmoveto{\pgfqpoint{1.439090in}{1.033644in}}%
\pgfpathlineto{\pgfqpoint{1.450660in}{1.047676in}}%
\pgfpathlineto{\pgfqpoint{1.462227in}{1.061705in}}%
\pgfpathlineto{\pgfqpoint{1.447711in}{1.070745in}}%
\pgfpathlineto{\pgfqpoint{1.431759in}{1.079348in}}%
\pgfpathlineto{\pgfqpoint{1.421198in}{1.064851in}}%
\pgfpathlineto{\pgfqpoint{1.410630in}{1.050344in}}%
\pgfpathlineto{\pgfqpoint{1.425538in}{1.042199in}}%
\pgfpathlineto{\pgfqpoint{1.439090in}{1.033644in}}%
\pgfpathclose%
\pgfusepath{fill}%
\end{pgfscope}%
\begin{pgfscope}%
\pgfpathrectangle{\pgfqpoint{0.407592in}{0.222904in}}{\pgfqpoint{1.411163in}{1.411163in}}%
\pgfusepath{clip}%
\pgfsetbuttcap%
\pgfsetroundjoin%
\definecolor{currentfill}{rgb}{0.749020,0.772549,0.792157}%
\pgfsetfillcolor{currentfill}%
\pgfsetfillopacity{0.160000}%
\pgfsetlinewidth{0.000000pt}%
\definecolor{currentstroke}{rgb}{0.000000,0.000000,0.000000}%
\pgfsetstrokecolor{currentstroke}%
\pgfsetdash{}{0pt}%
\pgfpathmoveto{\pgfqpoint{1.174180in}{0.768437in}}%
\pgfpathlineto{\pgfqpoint{1.182516in}{0.784117in}}%
\pgfpathlineto{\pgfqpoint{1.190842in}{0.799776in}}%
\pgfpathlineto{\pgfqpoint{1.182621in}{0.802289in}}%
\pgfpathlineto{\pgfqpoint{1.174013in}{0.804561in}}%
\pgfpathlineto{\pgfqpoint{1.167090in}{0.788527in}}%
\pgfpathlineto{\pgfqpoint{1.160155in}{0.772468in}}%
\pgfpathlineto{\pgfqpoint{1.167330in}{0.770554in}}%
\pgfpathlineto{\pgfqpoint{1.174180in}{0.768437in}}%
\pgfpathclose%
\pgfusepath{fill}%
\end{pgfscope}%
\begin{pgfscope}%
\pgfpathrectangle{\pgfqpoint{0.407592in}{0.222904in}}{\pgfqpoint{1.411163in}{1.411163in}}%
\pgfusepath{clip}%
\pgfsetbuttcap%
\pgfsetroundjoin%
\definecolor{currentfill}{rgb}{0.749020,0.772549,0.792157}%
\pgfsetfillcolor{currentfill}%
\pgfsetfillopacity{0.160000}%
\pgfsetlinewidth{0.000000pt}%
\definecolor{currentstroke}{rgb}{0.000000,0.000000,0.000000}%
\pgfsetstrokecolor{currentstroke}%
\pgfsetdash{}{0pt}%
\pgfpathmoveto{\pgfqpoint{0.765804in}{1.077659in}}%
\pgfpathlineto{\pgfqpoint{0.755773in}{1.092102in}}%
\pgfpathlineto{\pgfqpoint{0.745749in}{1.106536in}}%
\pgfpathlineto{\pgfqpoint{0.726958in}{1.097389in}}%
\pgfpathlineto{\pgfqpoint{0.709595in}{1.087786in}}%
\pgfpathlineto{\pgfqpoint{0.720737in}{1.073816in}}%
\pgfpathlineto{\pgfqpoint{0.731883in}{1.059841in}}%
\pgfpathlineto{\pgfqpoint{0.748167in}{1.068965in}}%
\pgfpathlineto{\pgfqpoint{0.765804in}{1.077659in}}%
\pgfpathclose%
\pgfusepath{fill}%
\end{pgfscope}%
\begin{pgfscope}%
\pgfpathrectangle{\pgfqpoint{0.407592in}{0.222904in}}{\pgfqpoint{1.411163in}{1.411163in}}%
\pgfusepath{clip}%
\pgfsetbuttcap%
\pgfsetroundjoin%
\definecolor{currentfill}{rgb}{0.749020,0.772549,0.792157}%
\pgfsetfillcolor{currentfill}%
\pgfsetfillopacity{0.160000}%
\pgfsetlinewidth{0.000000pt}%
\definecolor{currentstroke}{rgb}{0.000000,0.000000,0.000000}%
\pgfsetstrokecolor{currentstroke}%
\pgfsetdash{}{0pt}%
\pgfpathmoveto{\pgfqpoint{0.928596in}{0.884114in}}%
\pgfpathlineto{\pgfqpoint{0.919699in}{0.899173in}}%
\pgfpathlineto{\pgfqpoint{0.910810in}{0.914218in}}%
\pgfpathlineto{\pgfqpoint{0.898355in}{0.909074in}}%
\pgfpathlineto{\pgfqpoint{0.886680in}{0.903614in}}%
\pgfpathlineto{\pgfqpoint{0.896795in}{0.889050in}}%
\pgfpathlineto{\pgfqpoint{0.906917in}{0.874476in}}%
\pgfpathlineto{\pgfqpoint{0.917403in}{0.879437in}}%
\pgfpathlineto{\pgfqpoint{0.928596in}{0.884114in}}%
\pgfpathclose%
\pgfusepath{fill}%
\end{pgfscope}%
\begin{pgfscope}%
\pgfpathrectangle{\pgfqpoint{0.407592in}{0.222904in}}{\pgfqpoint{1.411163in}{1.411163in}}%
\pgfusepath{clip}%
\pgfsetbuttcap%
\pgfsetroundjoin%
\definecolor{currentfill}{rgb}{0.749020,0.772549,0.792157}%
\pgfsetfillcolor{currentfill}%
\pgfsetfillopacity{0.160000}%
\pgfsetlinewidth{0.000000pt}%
\definecolor{currentstroke}{rgb}{0.000000,0.000000,0.000000}%
\pgfsetstrokecolor{currentstroke}%
\pgfsetdash{}{0pt}%
\pgfpathmoveto{\pgfqpoint{1.028311in}{0.772101in}}%
\pgfpathlineto{\pgfqpoint{1.022170in}{0.788125in}}%
\pgfpathlineto{\pgfqpoint{1.016039in}{0.804125in}}%
\pgfpathlineto{\pgfqpoint{1.007017in}{0.801804in}}%
\pgfpathlineto{\pgfqpoint{0.998338in}{0.799244in}}%
\pgfpathlineto{\pgfqpoint{1.005943in}{0.783626in}}%
\pgfpathlineto{\pgfqpoint{1.013558in}{0.767989in}}%
\pgfpathlineto{\pgfqpoint{1.020790in}{0.770145in}}%
\pgfpathlineto{\pgfqpoint{1.028311in}{0.772101in}}%
\pgfpathclose%
\pgfusepath{fill}%
\end{pgfscope}%
\begin{pgfscope}%
\pgfpathrectangle{\pgfqpoint{0.407592in}{0.222904in}}{\pgfqpoint{1.411163in}{1.411163in}}%
\pgfusepath{clip}%
\pgfsetbuttcap%
\pgfsetroundjoin%
\definecolor{currentfill}{rgb}{0.749020,0.772549,0.792157}%
\pgfsetfillcolor{currentfill}%
\pgfsetfillopacity{0.160000}%
\pgfsetlinewidth{0.000000pt}%
\definecolor{currentstroke}{rgb}{0.000000,0.000000,0.000000}%
\pgfsetstrokecolor{currentstroke}%
\pgfsetdash{}{0pt}%
\pgfpathmoveto{\pgfqpoint{1.205982in}{0.794070in}}%
\pgfpathlineto{\pgfqpoint{1.215566in}{0.809266in}}%
\pgfpathlineto{\pgfqpoint{1.225140in}{0.824446in}}%
\pgfpathlineto{\pgfqpoint{1.216564in}{0.827868in}}%
\pgfpathlineto{\pgfqpoint{1.207460in}{0.831035in}}%
\pgfpathlineto{\pgfqpoint{1.199156in}{0.815416in}}%
\pgfpathlineto{\pgfqpoint{1.190842in}{0.799776in}}%
\pgfpathlineto{\pgfqpoint{1.198640in}{0.797033in}}%
\pgfpathlineto{\pgfqpoint{1.205982in}{0.794070in}}%
\pgfpathclose%
\pgfusepath{fill}%
\end{pgfscope}%
\begin{pgfscope}%
\pgfpathrectangle{\pgfqpoint{0.407592in}{0.222904in}}{\pgfqpoint{1.411163in}{1.411163in}}%
\pgfusepath{clip}%
\pgfsetbuttcap%
\pgfsetroundjoin%
\definecolor{currentfill}{rgb}{0.749020,0.772549,0.792157}%
\pgfsetfillcolor{currentfill}%
\pgfsetfillopacity{0.160000}%
\pgfsetlinewidth{0.000000pt}%
\definecolor{currentstroke}{rgb}{0.000000,0.000000,0.000000}%
\pgfsetstrokecolor{currentstroke}%
\pgfsetdash{}{0pt}%
\pgfpathmoveto{\pgfqpoint{1.415938in}{1.005564in}}%
\pgfpathlineto{\pgfqpoint{1.427516in}{1.019606in}}%
\pgfpathlineto{\pgfqpoint{1.439090in}{1.033644in}}%
\pgfpathlineto{\pgfqpoint{1.425538in}{1.042199in}}%
\pgfpathlineto{\pgfqpoint{1.410630in}{1.050344in}}%
\pgfpathlineto{\pgfqpoint{1.400055in}{1.035827in}}%
\pgfpathlineto{\pgfqpoint{1.389472in}{1.021301in}}%
\pgfpathlineto{\pgfqpoint{1.403342in}{1.013625in}}%
\pgfpathlineto{\pgfqpoint{1.415938in}{1.005564in}}%
\pgfpathclose%
\pgfusepath{fill}%
\end{pgfscope}%
\begin{pgfscope}%
\pgfpathrectangle{\pgfqpoint{0.407592in}{0.222904in}}{\pgfqpoint{1.411163in}{1.411163in}}%
\pgfusepath{clip}%
\pgfsetbuttcap%
\pgfsetroundjoin%
\definecolor{currentfill}{rgb}{0.749020,0.772549,0.792157}%
\pgfsetfillcolor{currentfill}%
\pgfsetfillopacity{0.160000}%
\pgfsetlinewidth{0.000000pt}%
\definecolor{currentstroke}{rgb}{0.000000,0.000000,0.000000}%
\pgfsetstrokecolor{currentstroke}%
\pgfsetdash{}{0pt}%
\pgfpathmoveto{\pgfqpoint{1.146255in}{0.740277in}}%
\pgfpathlineto{\pgfqpoint{1.153210in}{0.756385in}}%
\pgfpathlineto{\pgfqpoint{1.160155in}{0.772468in}}%
\pgfpathlineto{\pgfqpoint{1.152685in}{0.774172in}}%
\pgfpathlineto{\pgfqpoint{1.144952in}{0.775660in}}%
\pgfpathlineto{\pgfqpoint{1.139521in}{0.759274in}}%
\pgfpathlineto{\pgfqpoint{1.134082in}{0.742860in}}%
\pgfpathlineto{\pgfqpoint{1.140275in}{0.741656in}}%
\pgfpathlineto{\pgfqpoint{1.146255in}{0.740277in}}%
\pgfpathclose%
\pgfusepath{fill}%
\end{pgfscope}%
\begin{pgfscope}%
\pgfpathrectangle{\pgfqpoint{0.407592in}{0.222904in}}{\pgfqpoint{1.411163in}{1.411163in}}%
\pgfusepath{clip}%
\pgfsetbuttcap%
\pgfsetroundjoin%
\definecolor{currentfill}{rgb}{0.749020,0.772549,0.792157}%
\pgfsetfillcolor{currentfill}%
\pgfsetfillopacity{0.160000}%
\pgfsetlinewidth{0.000000pt}%
\definecolor{currentstroke}{rgb}{0.000000,0.000000,0.000000}%
\pgfsetstrokecolor{currentstroke}%
\pgfsetdash{}{0pt}%
\pgfpathmoveto{\pgfqpoint{0.998338in}{0.799244in}}%
\pgfpathlineto{\pgfqpoint{0.990743in}{0.814842in}}%
\pgfpathlineto{\pgfqpoint{0.983157in}{0.830421in}}%
\pgfpathlineto{\pgfqpoint{0.973474in}{0.827201in}}%
\pgfpathlineto{\pgfqpoint{0.964271in}{0.823730in}}%
\pgfpathlineto{\pgfqpoint{0.973212in}{0.808597in}}%
\pgfpathlineto{\pgfqpoint{0.982161in}{0.793450in}}%
\pgfpathlineto{\pgfqpoint{0.990041in}{0.796456in}}%
\pgfpathlineto{\pgfqpoint{0.998338in}{0.799244in}}%
\pgfpathclose%
\pgfusepath{fill}%
\end{pgfscope}%
\begin{pgfscope}%
\pgfpathrectangle{\pgfqpoint{0.407592in}{0.222904in}}{\pgfqpoint{1.411163in}{1.411163in}}%
\pgfusepath{clip}%
\pgfsetbuttcap%
\pgfsetroundjoin%
\definecolor{currentfill}{rgb}{0.749020,0.772549,0.792157}%
\pgfsetfillcolor{currentfill}%
\pgfsetfillopacity{0.160000}%
\pgfsetlinewidth{0.000000pt}%
\definecolor{currentstroke}{rgb}{0.000000,0.000000,0.000000}%
\pgfsetstrokecolor{currentstroke}%
\pgfsetdash{}{0pt}%
\pgfpathmoveto{\pgfqpoint{0.785886in}{1.048744in}}%
\pgfpathlineto{\pgfqpoint{0.775842in}{1.063206in}}%
\pgfpathlineto{\pgfqpoint{0.765804in}{1.077659in}}%
\pgfpathlineto{\pgfqpoint{0.748167in}{1.068965in}}%
\pgfpathlineto{\pgfqpoint{0.731883in}{1.059841in}}%
\pgfpathlineto{\pgfqpoint{0.743031in}{1.045863in}}%
\pgfpathlineto{\pgfqpoint{0.754183in}{1.031880in}}%
\pgfpathlineto{\pgfqpoint{0.769395in}{1.040515in}}%
\pgfpathlineto{\pgfqpoint{0.785886in}{1.048744in}}%
\pgfpathclose%
\pgfusepath{fill}%
\end{pgfscope}%
\begin{pgfscope}%
\pgfpathrectangle{\pgfqpoint{0.407592in}{0.222904in}}{\pgfqpoint{1.411163in}{1.411163in}}%
\pgfusepath{clip}%
\pgfsetbuttcap%
\pgfsetroundjoin%
\definecolor{currentfill}{rgb}{0.749020,0.772549,0.792157}%
\pgfsetfillcolor{currentfill}%
\pgfsetfillopacity{0.160000}%
\pgfsetlinewidth{0.000000pt}%
\definecolor{currentstroke}{rgb}{0.000000,0.000000,0.000000}%
\pgfsetstrokecolor{currentstroke}%
\pgfsetdash{}{0pt}%
\pgfpathmoveto{\pgfqpoint{1.053259in}{0.742632in}}%
\pgfpathlineto{\pgfqpoint{1.048673in}{0.759020in}}%
\pgfpathlineto{\pgfqpoint{1.044094in}{0.775379in}}%
\pgfpathlineto{\pgfqpoint{1.036091in}{0.773848in}}%
\pgfpathlineto{\pgfqpoint{1.028311in}{0.772101in}}%
\pgfpathlineto{\pgfqpoint{1.034462in}{0.756052in}}%
\pgfpathlineto{\pgfqpoint{1.040621in}{0.739980in}}%
\pgfpathlineto{\pgfqpoint{1.046849in}{0.741393in}}%
\pgfpathlineto{\pgfqpoint{1.053259in}{0.742632in}}%
\pgfpathclose%
\pgfusepath{fill}%
\end{pgfscope}%
\begin{pgfscope}%
\pgfpathrectangle{\pgfqpoint{0.407592in}{0.222904in}}{\pgfqpoint{1.411163in}{1.411163in}}%
\pgfusepath{clip}%
\pgfsetbuttcap%
\pgfsetroundjoin%
\definecolor{currentfill}{rgb}{0.749020,0.772549,0.792157}%
\pgfsetfillcolor{currentfill}%
\pgfsetfillopacity{0.160000}%
\pgfsetlinewidth{0.000000pt}%
\definecolor{currentstroke}{rgb}{0.000000,0.000000,0.000000}%
\pgfsetstrokecolor{currentstroke}%
\pgfsetdash{}{0pt}%
\pgfpathmoveto{\pgfqpoint{1.261922in}{0.846213in}}%
\pgfpathlineto{\pgfqpoint{1.272591in}{0.860858in}}%
\pgfpathlineto{\pgfqpoint{1.283252in}{0.875494in}}%
\pgfpathlineto{\pgfqpoint{1.273676in}{0.880399in}}%
\pgfpathlineto{\pgfqpoint{1.263340in}{0.885016in}}%
\pgfpathlineto{\pgfqpoint{1.253804in}{0.869896in}}%
\pgfpathlineto{\pgfqpoint{1.244259in}{0.854761in}}%
\pgfpathlineto{\pgfqpoint{1.253431in}{0.850616in}}%
\pgfpathlineto{\pgfqpoint{1.261922in}{0.846213in}}%
\pgfpathclose%
\pgfusepath{fill}%
\end{pgfscope}%
\begin{pgfscope}%
\pgfpathrectangle{\pgfqpoint{0.407592in}{0.222904in}}{\pgfqpoint{1.411163in}{1.411163in}}%
\pgfusepath{clip}%
\pgfsetbuttcap%
\pgfsetroundjoin%
\definecolor{currentfill}{rgb}{0.749020,0.772549,0.792157}%
\pgfsetfillcolor{currentfill}%
\pgfsetfillopacity{0.160000}%
\pgfsetlinewidth{0.000000pt}%
\definecolor{currentstroke}{rgb}{0.000000,0.000000,0.000000}%
\pgfsetstrokecolor{currentstroke}%
\pgfsetdash{}{0pt}%
\pgfpathmoveto{\pgfqpoint{0.946416in}{0.853951in}}%
\pgfpathlineto{\pgfqpoint{0.937502in}{0.869040in}}%
\pgfpathlineto{\pgfqpoint{0.928596in}{0.884114in}}%
\pgfpathlineto{\pgfqpoint{0.917403in}{0.879437in}}%
\pgfpathlineto{\pgfqpoint{0.906917in}{0.874476in}}%
\pgfpathlineto{\pgfqpoint{0.917045in}{0.859893in}}%
\pgfpathlineto{\pgfqpoint{0.927180in}{0.845300in}}%
\pgfpathlineto{\pgfqpoint{0.936481in}{0.849753in}}%
\pgfpathlineto{\pgfqpoint{0.946416in}{0.853951in}}%
\pgfpathclose%
\pgfusepath{fill}%
\end{pgfscope}%
\begin{pgfscope}%
\pgfpathrectangle{\pgfqpoint{0.407592in}{0.222904in}}{\pgfqpoint{1.411163in}{1.411163in}}%
\pgfusepath{clip}%
\pgfsetbuttcap%
\pgfsetroundjoin%
\definecolor{currentfill}{rgb}{0.749020,0.772549,0.792157}%
\pgfsetfillcolor{currentfill}%
\pgfsetfillopacity{0.160000}%
\pgfsetlinewidth{0.000000pt}%
\definecolor{currentstroke}{rgb}{0.000000,0.000000,0.000000}%
\pgfsetstrokecolor{currentstroke}%
\pgfsetdash{}{0pt}%
\pgfpathmoveto{\pgfqpoint{1.392771in}{0.977467in}}%
\pgfpathlineto{\pgfqpoint{1.404357in}{0.991518in}}%
\pgfpathlineto{\pgfqpoint{1.415938in}{1.005564in}}%
\pgfpathlineto{\pgfqpoint{1.403342in}{1.013625in}}%
\pgfpathlineto{\pgfqpoint{1.389472in}{1.021301in}}%
\pgfpathlineto{\pgfqpoint{1.378882in}{1.006764in}}%
\pgfpathlineto{\pgfqpoint{1.368286in}{0.992218in}}%
\pgfpathlineto{\pgfqpoint{1.381123in}{0.985021in}}%
\pgfpathlineto{\pgfqpoint{1.392771in}{0.977467in}}%
\pgfpathclose%
\pgfusepath{fill}%
\end{pgfscope}%
\begin{pgfscope}%
\pgfpathrectangle{\pgfqpoint{0.407592in}{0.222904in}}{\pgfqpoint{1.411163in}{1.411163in}}%
\pgfusepath{clip}%
\pgfsetbuttcap%
\pgfsetroundjoin%
\definecolor{currentfill}{rgb}{0.749020,0.772549,0.792157}%
\pgfsetfillcolor{currentfill}%
\pgfsetfillopacity{0.160000}%
\pgfsetlinewidth{0.000000pt}%
\definecolor{currentstroke}{rgb}{0.000000,0.000000,0.000000}%
\pgfsetstrokecolor{currentstroke}%
\pgfsetdash{}{0pt}%
\pgfpathmoveto{\pgfqpoint{1.103372in}{0.712206in}}%
\pgfpathlineto{\pgfqpoint{1.105548in}{0.729038in}}%
\pgfpathlineto{\pgfqpoint{1.107720in}{0.745838in}}%
\pgfpathlineto{\pgfqpoint{1.100871in}{0.746108in}}%
\pgfpathlineto{\pgfqpoint{1.093976in}{0.746185in}}%
\pgfpathlineto{\pgfqpoint{1.093511in}{0.729344in}}%
\pgfpathlineto{\pgfqpoint{1.093045in}{0.712470in}}%
\pgfpathlineto{\pgfqpoint{1.098226in}{0.712411in}}%
\pgfpathlineto{\pgfqpoint{1.103372in}{0.712206in}}%
\pgfpathclose%
\pgfusepath{fill}%
\end{pgfscope}%
\begin{pgfscope}%
\pgfpathrectangle{\pgfqpoint{0.407592in}{0.222904in}}{\pgfqpoint{1.411163in}{1.411163in}}%
\pgfusepath{clip}%
\pgfsetbuttcap%
\pgfsetroundjoin%
\definecolor{currentfill}{rgb}{0.749020,0.772549,0.792157}%
\pgfsetfillcolor{currentfill}%
\pgfsetfillopacity{0.160000}%
\pgfsetlinewidth{0.000000pt}%
\definecolor{currentstroke}{rgb}{0.000000,0.000000,0.000000}%
\pgfsetstrokecolor{currentstroke}%
\pgfsetdash{}{0pt}%
\pgfpathmoveto{\pgfqpoint{1.093045in}{0.712470in}}%
\pgfpathlineto{\pgfqpoint{1.093511in}{0.729344in}}%
\pgfpathlineto{\pgfqpoint{1.093976in}{0.746185in}}%
\pgfpathlineto{\pgfqpoint{1.087064in}{0.746069in}}%
\pgfpathlineto{\pgfqpoint{1.080164in}{0.745760in}}%
\pgfpathlineto{\pgfqpoint{1.081414in}{0.728970in}}%
\pgfpathlineto{\pgfqpoint{1.082666in}{0.712148in}}%
\pgfpathlineto{\pgfqpoint{1.087851in}{0.712382in}}%
\pgfpathlineto{\pgfqpoint{1.093045in}{0.712470in}}%
\pgfpathclose%
\pgfusepath{fill}%
\end{pgfscope}%
\begin{pgfscope}%
\pgfpathrectangle{\pgfqpoint{0.407592in}{0.222904in}}{\pgfqpoint{1.411163in}{1.411163in}}%
\pgfusepath{clip}%
\pgfsetbuttcap%
\pgfsetroundjoin%
\definecolor{currentfill}{rgb}{0.749020,0.772549,0.792157}%
\pgfsetfillcolor{currentfill}%
\pgfsetfillopacity{0.160000}%
\pgfsetlinewidth{0.000000pt}%
\definecolor{currentstroke}{rgb}{0.000000,0.000000,0.000000}%
\pgfsetstrokecolor{currentstroke}%
\pgfsetdash{}{0pt}%
\pgfpathmoveto{\pgfqpoint{0.805993in}{1.019793in}}%
\pgfpathlineto{\pgfqpoint{0.795936in}{1.034273in}}%
\pgfpathlineto{\pgfqpoint{0.785886in}{1.048744in}}%
\pgfpathlineto{\pgfqpoint{0.769395in}{1.040515in}}%
\pgfpathlineto{\pgfqpoint{0.754183in}{1.031880in}}%
\pgfpathlineto{\pgfqpoint{0.765338in}{1.017894in}}%
\pgfpathlineto{\pgfqpoint{0.776496in}{1.003903in}}%
\pgfpathlineto{\pgfqpoint{0.790644in}{1.012037in}}%
\pgfpathlineto{\pgfqpoint{0.805993in}{1.019793in}}%
\pgfpathclose%
\pgfusepath{fill}%
\end{pgfscope}%
\begin{pgfscope}%
\pgfpathrectangle{\pgfqpoint{0.407592in}{0.222904in}}{\pgfqpoint{1.411163in}{1.411163in}}%
\pgfusepath{clip}%
\pgfsetbuttcap%
\pgfsetroundjoin%
\definecolor{currentfill}{rgb}{0.749020,0.772549,0.792157}%
\pgfsetfillcolor{currentfill}%
\pgfsetfillopacity{0.160000}%
\pgfsetlinewidth{0.000000pt}%
\definecolor{currentstroke}{rgb}{0.000000,0.000000,0.000000}%
\pgfsetstrokecolor{currentstroke}%
\pgfsetdash{}{0pt}%
\pgfpathmoveto{\pgfqpoint{1.113473in}{0.711361in}}%
\pgfpathlineto{\pgfqpoint{1.117322in}{0.728057in}}%
\pgfpathlineto{\pgfqpoint{1.121164in}{0.744723in}}%
\pgfpathlineto{\pgfqpoint{1.114494in}{0.745375in}}%
\pgfpathlineto{\pgfqpoint{1.107720in}{0.745838in}}%
\pgfpathlineto{\pgfqpoint{1.105548in}{0.729038in}}%
\pgfpathlineto{\pgfqpoint{1.103372in}{0.712206in}}%
\pgfpathlineto{\pgfqpoint{1.108462in}{0.711856in}}%
\pgfpathlineto{\pgfqpoint{1.113473in}{0.711361in}}%
\pgfpathclose%
\pgfusepath{fill}%
\end{pgfscope}%
\begin{pgfscope}%
\pgfpathrectangle{\pgfqpoint{0.407592in}{0.222904in}}{\pgfqpoint{1.411163in}{1.411163in}}%
\pgfusepath{clip}%
\pgfsetbuttcap%
\pgfsetroundjoin%
\definecolor{currentfill}{rgb}{0.749020,0.772549,0.792157}%
\pgfsetfillcolor{currentfill}%
\pgfsetfillopacity{0.160000}%
\pgfsetlinewidth{0.000000pt}%
\definecolor{currentstroke}{rgb}{0.000000,0.000000,0.000000}%
\pgfsetstrokecolor{currentstroke}%
\pgfsetdash{}{0pt}%
\pgfpathmoveto{\pgfqpoint{1.082666in}{0.712148in}}%
\pgfpathlineto{\pgfqpoint{1.081414in}{0.728970in}}%
\pgfpathlineto{\pgfqpoint{1.080164in}{0.745760in}}%
\pgfpathlineto{\pgfqpoint{1.073304in}{0.745260in}}%
\pgfpathlineto{\pgfqpoint{1.066515in}{0.744570in}}%
\pgfpathlineto{\pgfqpoint{1.069460in}{0.727922in}}%
\pgfpathlineto{\pgfqpoint{1.072411in}{0.711244in}}%
\pgfpathlineto{\pgfqpoint{1.077512in}{0.711768in}}%
\pgfpathlineto{\pgfqpoint{1.082666in}{0.712148in}}%
\pgfpathclose%
\pgfusepath{fill}%
\end{pgfscope}%
\begin{pgfscope}%
\pgfpathrectangle{\pgfqpoint{0.407592in}{0.222904in}}{\pgfqpoint{1.411163in}{1.411163in}}%
\pgfusepath{clip}%
\pgfsetbuttcap%
\pgfsetroundjoin%
\definecolor{currentfill}{rgb}{0.749020,0.772549,0.792157}%
\pgfsetfillcolor{currentfill}%
\pgfsetfillopacity{0.160000}%
\pgfsetlinewidth{0.000000pt}%
\definecolor{currentstroke}{rgb}{0.000000,0.000000,0.000000}%
\pgfsetstrokecolor{currentstroke}%
\pgfsetdash{}{0pt}%
\pgfpathmoveto{\pgfqpoint{1.240562in}{0.816893in}}%
\pgfpathlineto{\pgfqpoint{1.251246in}{0.831558in}}%
\pgfpathlineto{\pgfqpoint{1.261922in}{0.846213in}}%
\pgfpathlineto{\pgfqpoint{1.253431in}{0.850616in}}%
\pgfpathlineto{\pgfqpoint{1.244259in}{0.854761in}}%
\pgfpathlineto{\pgfqpoint{1.234704in}{0.839611in}}%
\pgfpathlineto{\pgfqpoint{1.225140in}{0.824446in}}%
\pgfpathlineto{\pgfqpoint{1.233151in}{0.820783in}}%
\pgfpathlineto{\pgfqpoint{1.240562in}{0.816893in}}%
\pgfpathclose%
\pgfusepath{fill}%
\end{pgfscope}%
\begin{pgfscope}%
\pgfpathrectangle{\pgfqpoint{0.407592in}{0.222904in}}{\pgfqpoint{1.411163in}{1.411163in}}%
\pgfusepath{clip}%
\pgfsetbuttcap%
\pgfsetroundjoin%
\definecolor{currentfill}{rgb}{0.749020,0.772549,0.792157}%
\pgfsetfillcolor{currentfill}%
\pgfsetfillopacity{0.160000}%
\pgfsetlinewidth{0.000000pt}%
\definecolor{currentstroke}{rgb}{0.000000,0.000000,0.000000}%
\pgfsetstrokecolor{currentstroke}%
\pgfsetdash{}{0pt}%
\pgfpathmoveto{\pgfqpoint{1.157476in}{0.737018in}}%
\pgfpathlineto{\pgfqpoint{1.165833in}{0.752738in}}%
\pgfpathlineto{\pgfqpoint{1.174180in}{0.768437in}}%
\pgfpathlineto{\pgfqpoint{1.167330in}{0.770554in}}%
\pgfpathlineto{\pgfqpoint{1.160155in}{0.772468in}}%
\pgfpathlineto{\pgfqpoint{1.153210in}{0.756385in}}%
\pgfpathlineto{\pgfqpoint{1.146255in}{0.740277in}}%
\pgfpathlineto{\pgfqpoint{1.151997in}{0.738729in}}%
\pgfpathlineto{\pgfqpoint{1.157476in}{0.737018in}}%
\pgfpathclose%
\pgfusepath{fill}%
\end{pgfscope}%
\begin{pgfscope}%
\pgfpathrectangle{\pgfqpoint{0.407592in}{0.222904in}}{\pgfqpoint{1.411163in}{1.411163in}}%
\pgfusepath{clip}%
\pgfsetbuttcap%
\pgfsetroundjoin%
\definecolor{currentfill}{rgb}{0.749020,0.772549,0.792157}%
\pgfsetfillcolor{currentfill}%
\pgfsetfillopacity{0.160000}%
\pgfsetlinewidth{0.000000pt}%
\definecolor{currentstroke}{rgb}{0.000000,0.000000,0.000000}%
\pgfsetstrokecolor{currentstroke}%
\pgfsetdash{}{0pt}%
\pgfpathmoveto{\pgfqpoint{0.826126in}{0.990804in}}%
\pgfpathlineto{\pgfqpoint{0.816056in}{1.005303in}}%
\pgfpathlineto{\pgfqpoint{0.805993in}{1.019793in}}%
\pgfpathlineto{\pgfqpoint{0.790644in}{1.012037in}}%
\pgfpathlineto{\pgfqpoint{0.776496in}{1.003903in}}%
\pgfpathlineto{\pgfqpoint{0.787657in}{0.989909in}}%
\pgfpathlineto{\pgfqpoint{0.798822in}{0.975910in}}%
\pgfpathlineto{\pgfqpoint{0.811912in}{0.983533in}}%
\pgfpathlineto{\pgfqpoint{0.826126in}{0.990804in}}%
\pgfpathclose%
\pgfusepath{fill}%
\end{pgfscope}%
\begin{pgfscope}%
\pgfpathrectangle{\pgfqpoint{0.407592in}{0.222904in}}{\pgfqpoint{1.411163in}{1.411163in}}%
\pgfusepath{clip}%
\pgfsetbuttcap%
\pgfsetroundjoin%
\definecolor{currentfill}{rgb}{0.749020,0.772549,0.792157}%
\pgfsetfillcolor{currentfill}%
\pgfsetfillopacity{0.160000}%
\pgfsetlinewidth{0.000000pt}%
\definecolor{currentstroke}{rgb}{0.000000,0.000000,0.000000}%
\pgfsetstrokecolor{currentstroke}%
\pgfsetdash{}{0pt}%
\pgfpathmoveto{\pgfqpoint{1.186787in}{0.763633in}}%
\pgfpathlineto{\pgfqpoint{1.196389in}{0.778859in}}%
\pgfpathlineto{\pgfqpoint{1.205982in}{0.794070in}}%
\pgfpathlineto{\pgfqpoint{1.198640in}{0.797033in}}%
\pgfpathlineto{\pgfqpoint{1.190842in}{0.799776in}}%
\pgfpathlineto{\pgfqpoint{1.182516in}{0.784117in}}%
\pgfpathlineto{\pgfqpoint{1.174180in}{0.768437in}}%
\pgfpathlineto{\pgfqpoint{1.180675in}{0.766128in}}%
\pgfpathlineto{\pgfqpoint{1.186787in}{0.763633in}}%
\pgfpathclose%
\pgfusepath{fill}%
\end{pgfscope}%
\begin{pgfscope}%
\pgfpathrectangle{\pgfqpoint{0.407592in}{0.222904in}}{\pgfqpoint{1.411163in}{1.411163in}}%
\pgfusepath{clip}%
\pgfsetbuttcap%
\pgfsetroundjoin%
\definecolor{currentfill}{rgb}{0.749020,0.772549,0.792157}%
\pgfsetfillcolor{currentfill}%
\pgfsetfillopacity{0.160000}%
\pgfsetlinewidth{0.000000pt}%
\definecolor{currentstroke}{rgb}{0.000000,0.000000,0.000000}%
\pgfsetstrokecolor{currentstroke}%
\pgfsetdash{}{0pt}%
\pgfpathmoveto{\pgfqpoint{1.369589in}{0.949351in}}%
\pgfpathlineto{\pgfqpoint{1.381182in}{0.963411in}}%
\pgfpathlineto{\pgfqpoint{1.392771in}{0.977467in}}%
\pgfpathlineto{\pgfqpoint{1.381123in}{0.985021in}}%
\pgfpathlineto{\pgfqpoint{1.368286in}{0.992218in}}%
\pgfpathlineto{\pgfqpoint{1.357682in}{0.977662in}}%
\pgfpathlineto{\pgfqpoint{1.347070in}{0.963096in}}%
\pgfpathlineto{\pgfqpoint{1.358882in}{0.956389in}}%
\pgfpathlineto{\pgfqpoint{1.369589in}{0.949351in}}%
\pgfpathclose%
\pgfusepath{fill}%
\end{pgfscope}%
\begin{pgfscope}%
\pgfpathrectangle{\pgfqpoint{0.407592in}{0.222904in}}{\pgfqpoint{1.411163in}{1.411163in}}%
\pgfusepath{clip}%
\pgfsetbuttcap%
\pgfsetroundjoin%
\definecolor{currentfill}{rgb}{0.749020,0.772549,0.792157}%
\pgfsetfillcolor{currentfill}%
\pgfsetfillopacity{0.160000}%
\pgfsetlinewidth{0.000000pt}%
\definecolor{currentstroke}{rgb}{0.000000,0.000000,0.000000}%
\pgfsetstrokecolor{currentstroke}%
\pgfsetdash{}{0pt}%
\pgfpathmoveto{\pgfqpoint{1.123175in}{0.709947in}}%
\pgfpathlineto{\pgfqpoint{1.128633in}{0.726417in}}%
\pgfpathlineto{\pgfqpoint{1.134082in}{0.742860in}}%
\pgfpathlineto{\pgfqpoint{1.127702in}{0.743883in}}%
\pgfpathlineto{\pgfqpoint{1.121164in}{0.744723in}}%
\pgfpathlineto{\pgfqpoint{1.117322in}{0.728057in}}%
\pgfpathlineto{\pgfqpoint{1.113473in}{0.711361in}}%
\pgfpathlineto{\pgfqpoint{1.118384in}{0.710724in}}%
\pgfpathlineto{\pgfqpoint{1.123175in}{0.709947in}}%
\pgfpathclose%
\pgfusepath{fill}%
\end{pgfscope}%
\begin{pgfscope}%
\pgfpathrectangle{\pgfqpoint{0.407592in}{0.222904in}}{\pgfqpoint{1.411163in}{1.411163in}}%
\pgfusepath{clip}%
\pgfsetbuttcap%
\pgfsetroundjoin%
\definecolor{currentfill}{rgb}{0.749020,0.772549,0.792157}%
\pgfsetfillcolor{currentfill}%
\pgfsetfillopacity{0.160000}%
\pgfsetlinewidth{0.000000pt}%
\definecolor{currentstroke}{rgb}{0.000000,0.000000,0.000000}%
\pgfsetstrokecolor{currentstroke}%
\pgfsetdash{}{0pt}%
\pgfpathmoveto{\pgfqpoint{1.040621in}{0.739980in}}%
\pgfpathlineto{\pgfqpoint{1.034462in}{0.756052in}}%
\pgfpathlineto{\pgfqpoint{1.028311in}{0.772101in}}%
\pgfpathlineto{\pgfqpoint{1.020790in}{0.770145in}}%
\pgfpathlineto{\pgfqpoint{1.013558in}{0.767989in}}%
\pgfpathlineto{\pgfqpoint{1.021182in}{0.752333in}}%
\pgfpathlineto{\pgfqpoint{1.028815in}{0.736656in}}%
\pgfpathlineto{\pgfqpoint{1.034601in}{0.738399in}}%
\pgfpathlineto{\pgfqpoint{1.040621in}{0.739980in}}%
\pgfpathclose%
\pgfusepath{fill}%
\end{pgfscope}%
\begin{pgfscope}%
\pgfpathrectangle{\pgfqpoint{0.407592in}{0.222904in}}{\pgfqpoint{1.411163in}{1.411163in}}%
\pgfusepath{clip}%
\pgfsetbuttcap%
\pgfsetroundjoin%
\definecolor{currentfill}{rgb}{0.749020,0.772549,0.792157}%
\pgfsetfillcolor{currentfill}%
\pgfsetfillopacity{0.160000}%
\pgfsetlinewidth{0.000000pt}%
\definecolor{currentstroke}{rgb}{0.000000,0.000000,0.000000}%
\pgfsetstrokecolor{currentstroke}%
\pgfsetdash{}{0pt}%
\pgfpathmoveto{\pgfqpoint{1.072411in}{0.711244in}}%
\pgfpathlineto{\pgfqpoint{1.069460in}{0.727922in}}%
\pgfpathlineto{\pgfqpoint{1.066515in}{0.744570in}}%
\pgfpathlineto{\pgfqpoint{1.059824in}{0.743693in}}%
\pgfpathlineto{\pgfqpoint{1.053259in}{0.742632in}}%
\pgfpathlineto{\pgfqpoint{1.057853in}{0.726217in}}%
\pgfpathlineto{\pgfqpoint{1.062455in}{0.709775in}}%
\pgfpathlineto{\pgfqpoint{1.067385in}{0.710579in}}%
\pgfpathlineto{\pgfqpoint{1.072411in}{0.711244in}}%
\pgfpathclose%
\pgfusepath{fill}%
\end{pgfscope}%
\begin{pgfscope}%
\pgfpathrectangle{\pgfqpoint{0.407592in}{0.222904in}}{\pgfqpoint{1.411163in}{1.411163in}}%
\pgfusepath{clip}%
\pgfsetbuttcap%
\pgfsetroundjoin%
\definecolor{currentfill}{rgb}{0.749020,0.772549,0.792157}%
\pgfsetfillcolor{currentfill}%
\pgfsetfillopacity{0.160000}%
\pgfsetlinewidth{0.000000pt}%
\definecolor{currentstroke}{rgb}{0.000000,0.000000,0.000000}%
\pgfsetstrokecolor{currentstroke}%
\pgfsetdash{}{0pt}%
\pgfpathmoveto{\pgfqpoint{1.013558in}{0.767989in}}%
\pgfpathlineto{\pgfqpoint{1.005943in}{0.783626in}}%
\pgfpathlineto{\pgfqpoint{0.998338in}{0.799244in}}%
\pgfpathlineto{\pgfqpoint{0.990041in}{0.796456in}}%
\pgfpathlineto{\pgfqpoint{0.982161in}{0.793450in}}%
\pgfpathlineto{\pgfqpoint{0.991119in}{0.778288in}}%
\pgfpathlineto{\pgfqpoint{1.000085in}{0.763112in}}%
\pgfpathlineto{\pgfqpoint{1.006646in}{0.765642in}}%
\pgfpathlineto{\pgfqpoint{1.013558in}{0.767989in}}%
\pgfpathclose%
\pgfusepath{fill}%
\end{pgfscope}%
\begin{pgfscope}%
\pgfpathrectangle{\pgfqpoint{0.407592in}{0.222904in}}{\pgfqpoint{1.411163in}{1.411163in}}%
\pgfusepath{clip}%
\pgfsetbuttcap%
\pgfsetroundjoin%
\definecolor{currentfill}{rgb}{0.749020,0.772549,0.792157}%
\pgfsetfillcolor{currentfill}%
\pgfsetfillopacity{0.160000}%
\pgfsetlinewidth{0.000000pt}%
\definecolor{currentstroke}{rgb}{0.000000,0.000000,0.000000}%
\pgfsetstrokecolor{currentstroke}%
\pgfsetdash{}{0pt}%
\pgfpathmoveto{\pgfqpoint{0.964271in}{0.823730in}}%
\pgfpathlineto{\pgfqpoint{0.955340in}{0.838848in}}%
\pgfpathlineto{\pgfqpoint{0.946416in}{0.853951in}}%
\pgfpathlineto{\pgfqpoint{0.936481in}{0.849753in}}%
\pgfpathlineto{\pgfqpoint{0.927180in}{0.845300in}}%
\pgfpathlineto{\pgfqpoint{0.937321in}{0.830698in}}%
\pgfpathlineto{\pgfqpoint{0.947469in}{0.816087in}}%
\pgfpathlineto{\pgfqpoint{0.955591in}{0.820020in}}%
\pgfpathlineto{\pgfqpoint{0.964271in}{0.823730in}}%
\pgfpathclose%
\pgfusepath{fill}%
\end{pgfscope}%
\begin{pgfscope}%
\pgfpathrectangle{\pgfqpoint{0.407592in}{0.222904in}}{\pgfqpoint{1.411163in}{1.411163in}}%
\pgfusepath{clip}%
\pgfsetbuttcap%
\pgfsetroundjoin%
\definecolor{currentfill}{rgb}{0.749020,0.772549,0.792157}%
\pgfsetfillcolor{currentfill}%
\pgfsetfillopacity{0.160000}%
\pgfsetlinewidth{0.000000pt}%
\definecolor{currentstroke}{rgb}{0.000000,0.000000,0.000000}%
\pgfsetstrokecolor{currentstroke}%
\pgfsetdash{}{0pt}%
\pgfpathmoveto{\pgfqpoint{0.846285in}{0.961778in}}%
\pgfpathlineto{\pgfqpoint{0.836202in}{0.976296in}}%
\pgfpathlineto{\pgfqpoint{0.826126in}{0.990804in}}%
\pgfpathlineto{\pgfqpoint{0.811912in}{0.983533in}}%
\pgfpathlineto{\pgfqpoint{0.798822in}{0.975910in}}%
\pgfpathlineto{\pgfqpoint{0.809990in}{0.961908in}}%
\pgfpathlineto{\pgfqpoint{0.821161in}{0.947901in}}%
\pgfpathlineto{\pgfqpoint{0.833201in}{0.955003in}}%
\pgfpathlineto{\pgfqpoint{0.846285in}{0.961778in}}%
\pgfpathclose%
\pgfusepath{fill}%
\end{pgfscope}%
\begin{pgfscope}%
\pgfpathrectangle{\pgfqpoint{0.407592in}{0.222904in}}{\pgfqpoint{1.411163in}{1.411163in}}%
\pgfusepath{clip}%
\pgfsetbuttcap%
\pgfsetroundjoin%
\definecolor{currentfill}{rgb}{0.749020,0.772549,0.792157}%
\pgfsetfillcolor{currentfill}%
\pgfsetfillopacity{0.160000}%
\pgfsetlinewidth{0.000000pt}%
\definecolor{currentstroke}{rgb}{0.000000,0.000000,0.000000}%
\pgfsetstrokecolor{currentstroke}%
\pgfsetdash{}{0pt}%
\pgfpathmoveto{\pgfqpoint{1.346392in}{0.921217in}}%
\pgfpathlineto{\pgfqpoint{1.357993in}{0.935286in}}%
\pgfpathlineto{\pgfqpoint{1.369589in}{0.949351in}}%
\pgfpathlineto{\pgfqpoint{1.358882in}{0.956389in}}%
\pgfpathlineto{\pgfqpoint{1.347070in}{0.963096in}}%
\pgfpathlineto{\pgfqpoint{1.336452in}{0.948521in}}%
\pgfpathlineto{\pgfqpoint{1.325827in}{0.933935in}}%
\pgfpathlineto{\pgfqpoint{1.336619in}{0.927728in}}%
\pgfpathlineto{\pgfqpoint{1.346392in}{0.921217in}}%
\pgfpathclose%
\pgfusepath{fill}%
\end{pgfscope}%
\begin{pgfscope}%
\pgfpathrectangle{\pgfqpoint{0.407592in}{0.222904in}}{\pgfqpoint{1.411163in}{1.411163in}}%
\pgfusepath{clip}%
\pgfsetbuttcap%
\pgfsetroundjoin%
\definecolor{currentfill}{rgb}{0.749020,0.772549,0.792157}%
\pgfsetfillcolor{currentfill}%
\pgfsetfillopacity{0.160000}%
\pgfsetlinewidth{0.000000pt}%
\definecolor{currentstroke}{rgb}{0.000000,0.000000,0.000000}%
\pgfsetstrokecolor{currentstroke}%
\pgfsetdash{}{0pt}%
\pgfpathmoveto{\pgfqpoint{1.132312in}{0.707989in}}%
\pgfpathlineto{\pgfqpoint{1.139289in}{0.724145in}}%
\pgfpathlineto{\pgfqpoint{1.146255in}{0.740277in}}%
\pgfpathlineto{\pgfqpoint{1.140275in}{0.741656in}}%
\pgfpathlineto{\pgfqpoint{1.134082in}{0.742860in}}%
\pgfpathlineto{\pgfqpoint{1.128633in}{0.726417in}}%
\pgfpathlineto{\pgfqpoint{1.123175in}{0.709947in}}%
\pgfpathlineto{\pgfqpoint{1.127824in}{0.709034in}}%
\pgfpathlineto{\pgfqpoint{1.132312in}{0.707989in}}%
\pgfpathclose%
\pgfusepath{fill}%
\end{pgfscope}%
\begin{pgfscope}%
\pgfpathrectangle{\pgfqpoint{0.407592in}{0.222904in}}{\pgfqpoint{1.411163in}{1.411163in}}%
\pgfusepath{clip}%
\pgfsetbuttcap%
\pgfsetroundjoin%
\definecolor{currentfill}{rgb}{0.749020,0.772549,0.792157}%
\pgfsetfillcolor{currentfill}%
\pgfsetfillopacity{0.160000}%
\pgfsetlinewidth{0.000000pt}%
\definecolor{currentstroke}{rgb}{0.000000,0.000000,0.000000}%
\pgfsetstrokecolor{currentstroke}%
\pgfsetdash{}{0pt}%
\pgfpathmoveto{\pgfqpoint{1.062455in}{0.709775in}}%
\pgfpathlineto{\pgfqpoint{1.057853in}{0.726217in}}%
\pgfpathlineto{\pgfqpoint{1.053259in}{0.742632in}}%
\pgfpathlineto{\pgfqpoint{1.046849in}{0.741393in}}%
\pgfpathlineto{\pgfqpoint{1.040621in}{0.739980in}}%
\pgfpathlineto{\pgfqpoint{1.046790in}{0.723884in}}%
\pgfpathlineto{\pgfqpoint{1.052967in}{0.707764in}}%
\pgfpathlineto{\pgfqpoint{1.057642in}{0.708835in}}%
\pgfpathlineto{\pgfqpoint{1.062455in}{0.709775in}}%
\pgfpathclose%
\pgfusepath{fill}%
\end{pgfscope}%
\begin{pgfscope}%
\pgfpathrectangle{\pgfqpoint{0.407592in}{0.222904in}}{\pgfqpoint{1.411163in}{1.411163in}}%
\pgfusepath{clip}%
\pgfsetbuttcap%
\pgfsetroundjoin%
\definecolor{currentfill}{rgb}{0.749020,0.772549,0.792157}%
\pgfsetfillcolor{currentfill}%
\pgfsetfillopacity{0.160000}%
\pgfsetlinewidth{0.000000pt}%
\definecolor{currentstroke}{rgb}{0.000000,0.000000,0.000000}%
\pgfsetstrokecolor{currentstroke}%
\pgfsetdash{}{0pt}%
\pgfpathmoveto{\pgfqpoint{1.219174in}{0.787533in}}%
\pgfpathlineto{\pgfqpoint{1.229872in}{0.802218in}}%
\pgfpathlineto{\pgfqpoint{1.240562in}{0.816893in}}%
\pgfpathlineto{\pgfqpoint{1.233151in}{0.820783in}}%
\pgfpathlineto{\pgfqpoint{1.225140in}{0.824446in}}%
\pgfpathlineto{\pgfqpoint{1.215566in}{0.809266in}}%
\pgfpathlineto{\pgfqpoint{1.205982in}{0.794070in}}%
\pgfpathlineto{\pgfqpoint{1.212837in}{0.790899in}}%
\pgfpathlineto{\pgfqpoint{1.219174in}{0.787533in}}%
\pgfpathclose%
\pgfusepath{fill}%
\end{pgfscope}%
\begin{pgfscope}%
\pgfpathrectangle{\pgfqpoint{0.407592in}{0.222904in}}{\pgfqpoint{1.411163in}{1.411163in}}%
\pgfusepath{clip}%
\pgfsetbuttcap%
\pgfsetroundjoin%
\definecolor{currentfill}{rgb}{0.749020,0.772549,0.792157}%
\pgfsetfillcolor{currentfill}%
\pgfsetfillopacity{0.160000}%
\pgfsetlinewidth{0.000000pt}%
\definecolor{currentstroke}{rgb}{0.000000,0.000000,0.000000}%
\pgfsetstrokecolor{currentstroke}%
\pgfsetdash{}{0pt}%
\pgfpathmoveto{\pgfqpoint{0.866469in}{0.932715in}}%
\pgfpathlineto{\pgfqpoint{0.856374in}{0.947251in}}%
\pgfpathlineto{\pgfqpoint{0.846285in}{0.961778in}}%
\pgfpathlineto{\pgfqpoint{0.833201in}{0.955003in}}%
\pgfpathlineto{\pgfqpoint{0.821161in}{0.947901in}}%
\pgfpathlineto{\pgfqpoint{0.832335in}{0.933891in}}%
\pgfpathlineto{\pgfqpoint{0.843512in}{0.919876in}}%
\pgfpathlineto{\pgfqpoint{0.854509in}{0.926445in}}%
\pgfpathlineto{\pgfqpoint{0.866469in}{0.932715in}}%
\pgfpathclose%
\pgfusepath{fill}%
\end{pgfscope}%
\begin{pgfscope}%
\pgfpathrectangle{\pgfqpoint{0.407592in}{0.222904in}}{\pgfqpoint{1.411163in}{1.411163in}}%
\pgfusepath{clip}%
\pgfsetbuttcap%
\pgfsetroundjoin%
\definecolor{currentfill}{rgb}{0.749020,0.772549,0.792157}%
\pgfsetfillcolor{currentfill}%
\pgfsetfillopacity{0.160000}%
\pgfsetlinewidth{0.000000pt}%
\definecolor{currentstroke}{rgb}{0.000000,0.000000,0.000000}%
\pgfsetstrokecolor{currentstroke}%
\pgfsetdash{}{0pt}%
\pgfpathmoveto{\pgfqpoint{1.323181in}{0.893064in}}%
\pgfpathlineto{\pgfqpoint{1.334788in}{0.907143in}}%
\pgfpathlineto{\pgfqpoint{1.346392in}{0.921217in}}%
\pgfpathlineto{\pgfqpoint{1.336619in}{0.927728in}}%
\pgfpathlineto{\pgfqpoint{1.325827in}{0.933935in}}%
\pgfpathlineto{\pgfqpoint{1.315194in}{0.919339in}}%
\pgfpathlineto{\pgfqpoint{1.304554in}{0.904734in}}%
\pgfpathlineto{\pgfqpoint{1.314333in}{0.899038in}}%
\pgfpathlineto{\pgfqpoint{1.323181in}{0.893064in}}%
\pgfpathclose%
\pgfusepath{fill}%
\end{pgfscope}%
\begin{pgfscope}%
\pgfpathrectangle{\pgfqpoint{0.407592in}{0.222904in}}{\pgfqpoint{1.411163in}{1.411163in}}%
\pgfusepath{clip}%
\pgfsetbuttcap%
\pgfsetroundjoin%
\definecolor{currentfill}{rgb}{0.749020,0.772549,0.792157}%
\pgfsetfillcolor{currentfill}%
\pgfsetfillopacity{0.160000}%
\pgfsetlinewidth{0.000000pt}%
\definecolor{currentstroke}{rgb}{0.000000,0.000000,0.000000}%
\pgfsetstrokecolor{currentstroke}%
\pgfsetdash{}{0pt}%
\pgfpathmoveto{\pgfqpoint{0.982161in}{0.793450in}}%
\pgfpathlineto{\pgfqpoint{0.973212in}{0.808597in}}%
\pgfpathlineto{\pgfqpoint{0.964271in}{0.823730in}}%
\pgfpathlineto{\pgfqpoint{0.955591in}{0.820020in}}%
\pgfpathlineto{\pgfqpoint{0.947469in}{0.816087in}}%
\pgfpathlineto{\pgfqpoint{0.957624in}{0.801466in}}%
\pgfpathlineto{\pgfqpoint{0.967785in}{0.786836in}}%
\pgfpathlineto{\pgfqpoint{0.974731in}{0.790239in}}%
\pgfpathlineto{\pgfqpoint{0.982161in}{0.793450in}}%
\pgfpathclose%
\pgfusepath{fill}%
\end{pgfscope}%
\begin{pgfscope}%
\pgfpathrectangle{\pgfqpoint{0.407592in}{0.222904in}}{\pgfqpoint{1.411163in}{1.411163in}}%
\pgfusepath{clip}%
\pgfsetbuttcap%
\pgfsetroundjoin%
\definecolor{currentfill}{rgb}{0.749020,0.772549,0.792157}%
\pgfsetfillcolor{currentfill}%
\pgfsetfillopacity{0.160000}%
\pgfsetlinewidth{0.000000pt}%
\definecolor{currentstroke}{rgb}{0.000000,0.000000,0.000000}%
\pgfsetstrokecolor{currentstroke}%
\pgfsetdash{}{0pt}%
\pgfpathmoveto{\pgfqpoint{1.167553in}{0.733136in}}%
\pgfpathlineto{\pgfqpoint{1.177174in}{0.748392in}}%
\pgfpathlineto{\pgfqpoint{1.186787in}{0.763633in}}%
\pgfpathlineto{\pgfqpoint{1.180675in}{0.766128in}}%
\pgfpathlineto{\pgfqpoint{1.174180in}{0.768437in}}%
\pgfpathlineto{\pgfqpoint{1.165833in}{0.752738in}}%
\pgfpathlineto{\pgfqpoint{1.157476in}{0.737018in}}%
\pgfpathlineto{\pgfqpoint{1.162669in}{0.735151in}}%
\pgfpathlineto{\pgfqpoint{1.167553in}{0.733136in}}%
\pgfpathclose%
\pgfusepath{fill}%
\end{pgfscope}%
\begin{pgfscope}%
\pgfpathrectangle{\pgfqpoint{0.407592in}{0.222904in}}{\pgfqpoint{1.411163in}{1.411163in}}%
\pgfusepath{clip}%
\pgfsetbuttcap%
\pgfsetroundjoin%
\definecolor{currentfill}{rgb}{0.749020,0.772549,0.792157}%
\pgfsetfillcolor{currentfill}%
\pgfsetfillopacity{0.160000}%
\pgfsetlinewidth{0.000000pt}%
\definecolor{currentstroke}{rgb}{0.000000,0.000000,0.000000}%
\pgfsetstrokecolor{currentstroke}%
\pgfsetdash{}{0pt}%
\pgfpathmoveto{\pgfqpoint{1.028815in}{0.736656in}}%
\pgfpathlineto{\pgfqpoint{1.021182in}{0.752333in}}%
\pgfpathlineto{\pgfqpoint{1.013558in}{0.767989in}}%
\pgfpathlineto{\pgfqpoint{1.006646in}{0.765642in}}%
\pgfpathlineto{\pgfqpoint{1.000085in}{0.763112in}}%
\pgfpathlineto{\pgfqpoint{1.009060in}{0.747920in}}%
\pgfpathlineto{\pgfqpoint{1.018044in}{0.732714in}}%
\pgfpathlineto{\pgfqpoint{1.023288in}{0.734759in}}%
\pgfpathlineto{\pgfqpoint{1.028815in}{0.736656in}}%
\pgfpathclose%
\pgfusepath{fill}%
\end{pgfscope}%
\begin{pgfscope}%
\pgfpathrectangle{\pgfqpoint{0.407592in}{0.222904in}}{\pgfqpoint{1.411163in}{1.411163in}}%
\pgfusepath{clip}%
\pgfsetbuttcap%
\pgfsetroundjoin%
\definecolor{currentfill}{rgb}{0.749020,0.772549,0.792157}%
\pgfsetfillcolor{currentfill}%
\pgfsetfillopacity{0.160000}%
\pgfsetlinewidth{0.000000pt}%
\definecolor{currentstroke}{rgb}{0.000000,0.000000,0.000000}%
\pgfsetstrokecolor{currentstroke}%
\pgfsetdash{}{0pt}%
\pgfpathmoveto{\pgfqpoint{1.140729in}{0.705519in}}%
\pgfpathlineto{\pgfqpoint{1.149108in}{0.721279in}}%
\pgfpathlineto{\pgfqpoint{1.157476in}{0.737018in}}%
\pgfpathlineto{\pgfqpoint{1.151997in}{0.738729in}}%
\pgfpathlineto{\pgfqpoint{1.146255in}{0.740277in}}%
\pgfpathlineto{\pgfqpoint{1.139289in}{0.724145in}}%
\pgfpathlineto{\pgfqpoint{1.132312in}{0.707989in}}%
\pgfpathlineto{\pgfqpoint{1.136620in}{0.706815in}}%
\pgfpathlineto{\pgfqpoint{1.140729in}{0.705519in}}%
\pgfpathclose%
\pgfusepath{fill}%
\end{pgfscope}%
\begin{pgfscope}%
\pgfpathrectangle{\pgfqpoint{0.407592in}{0.222904in}}{\pgfqpoint{1.411163in}{1.411163in}}%
\pgfusepath{clip}%
\pgfsetbuttcap%
\pgfsetroundjoin%
\definecolor{currentfill}{rgb}{0.749020,0.772549,0.792157}%
\pgfsetfillcolor{currentfill}%
\pgfsetfillopacity{0.160000}%
\pgfsetlinewidth{0.000000pt}%
\definecolor{currentstroke}{rgb}{0.000000,0.000000,0.000000}%
\pgfsetstrokecolor{currentstroke}%
\pgfsetdash{}{0pt}%
\pgfpathmoveto{\pgfqpoint{0.886680in}{0.903614in}}%
\pgfpathlineto{\pgfqpoint{0.876572in}{0.918169in}}%
\pgfpathlineto{\pgfqpoint{0.866469in}{0.932715in}}%
\pgfpathlineto{\pgfqpoint{0.854509in}{0.926445in}}%
\pgfpathlineto{\pgfqpoint{0.843512in}{0.919876in}}%
\pgfpathlineto{\pgfqpoint{0.854693in}{0.905857in}}%
\pgfpathlineto{\pgfqpoint{0.865877in}{0.891835in}}%
\pgfpathlineto{\pgfqpoint{0.875838in}{0.897861in}}%
\pgfpathlineto{\pgfqpoint{0.886680in}{0.903614in}}%
\pgfpathclose%
\pgfusepath{fill}%
\end{pgfscope}%
\begin{pgfscope}%
\pgfpathrectangle{\pgfqpoint{0.407592in}{0.222904in}}{\pgfqpoint{1.411163in}{1.411163in}}%
\pgfusepath{clip}%
\pgfsetbuttcap%
\pgfsetroundjoin%
\definecolor{currentfill}{rgb}{0.749020,0.772549,0.792157}%
\pgfsetfillcolor{currentfill}%
\pgfsetfillopacity{0.160000}%
\pgfsetlinewidth{0.000000pt}%
\definecolor{currentstroke}{rgb}{0.000000,0.000000,0.000000}%
\pgfsetstrokecolor{currentstroke}%
\pgfsetdash{}{0pt}%
\pgfpathmoveto{\pgfqpoint{1.099008in}{0.678446in}}%
\pgfpathlineto{\pgfqpoint{1.101192in}{0.695342in}}%
\pgfpathlineto{\pgfqpoint{1.103372in}{0.712206in}}%
\pgfpathlineto{\pgfqpoint{1.098226in}{0.712411in}}%
\pgfpathlineto{\pgfqpoint{1.093045in}{0.712470in}}%
\pgfpathlineto{\pgfqpoint{1.092578in}{0.695563in}}%
\pgfpathlineto{\pgfqpoint{1.092110in}{0.678624in}}%
\pgfpathlineto{\pgfqpoint{1.095571in}{0.678584in}}%
\pgfpathlineto{\pgfqpoint{1.099008in}{0.678446in}}%
\pgfpathclose%
\pgfusepath{fill}%
\end{pgfscope}%
\begin{pgfscope}%
\pgfpathrectangle{\pgfqpoint{0.407592in}{0.222904in}}{\pgfqpoint{1.411163in}{1.411163in}}%
\pgfusepath{clip}%
\pgfsetbuttcap%
\pgfsetroundjoin%
\definecolor{currentfill}{rgb}{0.749020,0.772549,0.792157}%
\pgfsetfillcolor{currentfill}%
\pgfsetfillopacity{0.160000}%
\pgfsetlinewidth{0.000000pt}%
\definecolor{currentstroke}{rgb}{0.000000,0.000000,0.000000}%
\pgfsetstrokecolor{currentstroke}%
\pgfsetdash{}{0pt}%
\pgfpathmoveto{\pgfqpoint{1.092110in}{0.678624in}}%
\pgfpathlineto{\pgfqpoint{1.092578in}{0.695563in}}%
\pgfpathlineto{\pgfqpoint{1.093045in}{0.712470in}}%
\pgfpathlineto{\pgfqpoint{1.087851in}{0.712382in}}%
\pgfpathlineto{\pgfqpoint{1.082666in}{0.712148in}}%
\pgfpathlineto{\pgfqpoint{1.083920in}{0.695293in}}%
\pgfpathlineto{\pgfqpoint{1.085177in}{0.678406in}}%
\pgfpathlineto{\pgfqpoint{1.088641in}{0.678564in}}%
\pgfpathlineto{\pgfqpoint{1.092110in}{0.678624in}}%
\pgfpathclose%
\pgfusepath{fill}%
\end{pgfscope}%
\begin{pgfscope}%
\pgfpathrectangle{\pgfqpoint{0.407592in}{0.222904in}}{\pgfqpoint{1.411163in}{1.411163in}}%
\pgfusepath{clip}%
\pgfsetbuttcap%
\pgfsetroundjoin%
\definecolor{currentfill}{rgb}{0.749020,0.772549,0.792157}%
\pgfsetfillcolor{currentfill}%
\pgfsetfillopacity{0.160000}%
\pgfsetlinewidth{0.000000pt}%
\definecolor{currentstroke}{rgb}{0.000000,0.000000,0.000000}%
\pgfsetstrokecolor{currentstroke}%
\pgfsetdash{}{0pt}%
\pgfpathmoveto{\pgfqpoint{1.052967in}{0.707764in}}%
\pgfpathlineto{\pgfqpoint{1.046790in}{0.723884in}}%
\pgfpathlineto{\pgfqpoint{1.040621in}{0.739980in}}%
\pgfpathlineto{\pgfqpoint{1.034601in}{0.738399in}}%
\pgfpathlineto{\pgfqpoint{1.028815in}{0.736656in}}%
\pgfpathlineto{\pgfqpoint{1.036458in}{0.720960in}}%
\pgfpathlineto{\pgfqpoint{1.044111in}{0.705244in}}%
\pgfpathlineto{\pgfqpoint{1.048450in}{0.706565in}}%
\pgfpathlineto{\pgfqpoint{1.052967in}{0.707764in}}%
\pgfpathclose%
\pgfusepath{fill}%
\end{pgfscope}%
\begin{pgfscope}%
\pgfpathrectangle{\pgfqpoint{0.407592in}{0.222904in}}{\pgfqpoint{1.411163in}{1.411163in}}%
\pgfusepath{clip}%
\pgfsetbuttcap%
\pgfsetroundjoin%
\definecolor{currentfill}{rgb}{0.749020,0.772549,0.792157}%
\pgfsetfillcolor{currentfill}%
\pgfsetfillopacity{0.160000}%
\pgfsetlinewidth{0.000000pt}%
\definecolor{currentstroke}{rgb}{0.000000,0.000000,0.000000}%
\pgfsetstrokecolor{currentstroke}%
\pgfsetdash{}{0pt}%
\pgfpathmoveto{\pgfqpoint{1.299953in}{0.864894in}}%
\pgfpathlineto{\pgfqpoint{1.311569in}{0.878981in}}%
\pgfpathlineto{\pgfqpoint{1.323181in}{0.893064in}}%
\pgfpathlineto{\pgfqpoint{1.314333in}{0.899038in}}%
\pgfpathlineto{\pgfqpoint{1.304554in}{0.904734in}}%
\pgfpathlineto{\pgfqpoint{1.293907in}{0.890119in}}%
\pgfpathlineto{\pgfqpoint{1.283252in}{0.875494in}}%
\pgfpathlineto{\pgfqpoint{1.292025in}{0.870318in}}%
\pgfpathlineto{\pgfqpoint{1.299953in}{0.864894in}}%
\pgfpathclose%
\pgfusepath{fill}%
\end{pgfscope}%
\begin{pgfscope}%
\pgfpathrectangle{\pgfqpoint{0.407592in}{0.222904in}}{\pgfqpoint{1.411163in}{1.411163in}}%
\pgfusepath{clip}%
\pgfsetbuttcap%
\pgfsetroundjoin%
\definecolor{currentfill}{rgb}{0.749020,0.772549,0.792157}%
\pgfsetfillcolor{currentfill}%
\pgfsetfillopacity{0.160000}%
\pgfsetlinewidth{0.000000pt}%
\definecolor{currentstroke}{rgb}{0.000000,0.000000,0.000000}%
\pgfsetstrokecolor{currentstroke}%
\pgfsetdash{}{0pt}%
\pgfpathmoveto{\pgfqpoint{1.105754in}{0.677875in}}%
\pgfpathlineto{\pgfqpoint{1.109617in}{0.694633in}}%
\pgfpathlineto{\pgfqpoint{1.113473in}{0.711361in}}%
\pgfpathlineto{\pgfqpoint{1.108462in}{0.711856in}}%
\pgfpathlineto{\pgfqpoint{1.103372in}{0.712206in}}%
\pgfpathlineto{\pgfqpoint{1.101192in}{0.695342in}}%
\pgfpathlineto{\pgfqpoint{1.099008in}{0.678446in}}%
\pgfpathlineto{\pgfqpoint{1.102407in}{0.678209in}}%
\pgfpathlineto{\pgfqpoint{1.105754in}{0.677875in}}%
\pgfpathclose%
\pgfusepath{fill}%
\end{pgfscope}%
\begin{pgfscope}%
\pgfpathrectangle{\pgfqpoint{0.407592in}{0.222904in}}{\pgfqpoint{1.411163in}{1.411163in}}%
\pgfusepath{clip}%
\pgfsetbuttcap%
\pgfsetroundjoin%
\definecolor{currentfill}{rgb}{0.749020,0.772549,0.792157}%
\pgfsetfillcolor{currentfill}%
\pgfsetfillopacity{0.160000}%
\pgfsetlinewidth{0.000000pt}%
\definecolor{currentstroke}{rgb}{0.000000,0.000000,0.000000}%
\pgfsetstrokecolor{currentstroke}%
\pgfsetdash{}{0pt}%
\pgfpathmoveto{\pgfqpoint{1.085177in}{0.678406in}}%
\pgfpathlineto{\pgfqpoint{1.083920in}{0.695293in}}%
\pgfpathlineto{\pgfqpoint{1.082666in}{0.712148in}}%
\pgfpathlineto{\pgfqpoint{1.077512in}{0.711768in}}%
\pgfpathlineto{\pgfqpoint{1.072411in}{0.711244in}}%
\pgfpathlineto{\pgfqpoint{1.075367in}{0.694536in}}%
\pgfpathlineto{\pgfqpoint{1.078329in}{0.677797in}}%
\pgfpathlineto{\pgfqpoint{1.081735in}{0.678150in}}%
\pgfpathlineto{\pgfqpoint{1.085177in}{0.678406in}}%
\pgfpathclose%
\pgfusepath{fill}%
\end{pgfscope}%
\begin{pgfscope}%
\pgfpathrectangle{\pgfqpoint{0.407592in}{0.222904in}}{\pgfqpoint{1.411163in}{1.411163in}}%
\pgfusepath{clip}%
\pgfsetbuttcap%
\pgfsetroundjoin%
\definecolor{currentfill}{rgb}{0.749020,0.772549,0.792157}%
\pgfsetfillcolor{currentfill}%
\pgfsetfillopacity{0.160000}%
\pgfsetlinewidth{0.000000pt}%
\definecolor{currentstroke}{rgb}{0.000000,0.000000,0.000000}%
\pgfsetstrokecolor{currentstroke}%
\pgfsetdash{}{0pt}%
\pgfpathmoveto{\pgfqpoint{1.197756in}{0.758134in}}%
\pgfpathlineto{\pgfqpoint{1.208469in}{0.772839in}}%
\pgfpathlineto{\pgfqpoint{1.219174in}{0.787533in}}%
\pgfpathlineto{\pgfqpoint{1.212837in}{0.790899in}}%
\pgfpathlineto{\pgfqpoint{1.205982in}{0.794070in}}%
\pgfpathlineto{\pgfqpoint{1.196389in}{0.778859in}}%
\pgfpathlineto{\pgfqpoint{1.186787in}{0.763633in}}%
\pgfpathlineto{\pgfqpoint{1.192489in}{0.760965in}}%
\pgfpathlineto{\pgfqpoint{1.197756in}{0.758134in}}%
\pgfpathclose%
\pgfusepath{fill}%
\end{pgfscope}%
\begin{pgfscope}%
\pgfpathrectangle{\pgfqpoint{0.407592in}{0.222904in}}{\pgfqpoint{1.411163in}{1.411163in}}%
\pgfusepath{clip}%
\pgfsetbuttcap%
\pgfsetroundjoin%
\definecolor{currentfill}{rgb}{0.749020,0.772549,0.792157}%
\pgfsetfillcolor{currentfill}%
\pgfsetfillopacity{0.160000}%
\pgfsetlinewidth{0.000000pt}%
\definecolor{currentstroke}{rgb}{0.000000,0.000000,0.000000}%
\pgfsetstrokecolor{currentstroke}%
\pgfsetdash{}{0pt}%
\pgfpathmoveto{\pgfqpoint{1.000085in}{0.763112in}}%
\pgfpathlineto{\pgfqpoint{0.991119in}{0.778288in}}%
\pgfpathlineto{\pgfqpoint{0.982161in}{0.793450in}}%
\pgfpathlineto{\pgfqpoint{0.974731in}{0.790239in}}%
\pgfpathlineto{\pgfqpoint{0.967785in}{0.786836in}}%
\pgfpathlineto{\pgfqpoint{0.977952in}{0.772196in}}%
\pgfpathlineto{\pgfqpoint{0.988126in}{0.757547in}}%
\pgfpathlineto{\pgfqpoint{0.993903in}{0.760410in}}%
\pgfpathlineto{\pgfqpoint{1.000085in}{0.763112in}}%
\pgfpathclose%
\pgfusepath{fill}%
\end{pgfscope}%
\begin{pgfscope}%
\pgfpathrectangle{\pgfqpoint{0.407592in}{0.222904in}}{\pgfqpoint{1.411163in}{1.411163in}}%
\pgfusepath{clip}%
\pgfsetbuttcap%
\pgfsetroundjoin%
\definecolor{currentfill}{rgb}{0.749020,0.772549,0.792157}%
\pgfsetfillcolor{currentfill}%
\pgfsetfillopacity{0.160000}%
\pgfsetlinewidth{0.000000pt}%
\definecolor{currentstroke}{rgb}{0.000000,0.000000,0.000000}%
\pgfsetstrokecolor{currentstroke}%
\pgfsetdash{}{0pt}%
\pgfpathmoveto{\pgfqpoint{1.112231in}{0.676922in}}%
\pgfpathlineto{\pgfqpoint{1.117707in}{0.693449in}}%
\pgfpathlineto{\pgfqpoint{1.123175in}{0.709947in}}%
\pgfpathlineto{\pgfqpoint{1.118384in}{0.710724in}}%
\pgfpathlineto{\pgfqpoint{1.113473in}{0.711361in}}%
\pgfpathlineto{\pgfqpoint{1.109617in}{0.694633in}}%
\pgfpathlineto{\pgfqpoint{1.105754in}{0.677875in}}%
\pgfpathlineto{\pgfqpoint{1.109033in}{0.677446in}}%
\pgfpathlineto{\pgfqpoint{1.112231in}{0.676922in}}%
\pgfpathclose%
\pgfusepath{fill}%
\end{pgfscope}%
\begin{pgfscope}%
\pgfpathrectangle{\pgfqpoint{0.407592in}{0.222904in}}{\pgfqpoint{1.411163in}{1.411163in}}%
\pgfusepath{clip}%
\pgfsetbuttcap%
\pgfsetroundjoin%
\definecolor{currentfill}{rgb}{0.749020,0.772549,0.792157}%
\pgfsetfillcolor{currentfill}%
\pgfsetfillopacity{0.160000}%
\pgfsetlinewidth{0.000000pt}%
\definecolor{currentstroke}{rgb}{0.000000,0.000000,0.000000}%
\pgfsetstrokecolor{currentstroke}%
\pgfsetdash{}{0pt}%
\pgfpathmoveto{\pgfqpoint{0.906917in}{0.874476in}}%
\pgfpathlineto{\pgfqpoint{0.896795in}{0.889050in}}%
\pgfpathlineto{\pgfqpoint{0.886680in}{0.903614in}}%
\pgfpathlineto{\pgfqpoint{0.875838in}{0.897861in}}%
\pgfpathlineto{\pgfqpoint{0.865877in}{0.891835in}}%
\pgfpathlineto{\pgfqpoint{0.877064in}{0.877808in}}%
\pgfpathlineto{\pgfqpoint{0.888254in}{0.863777in}}%
\pgfpathlineto{\pgfqpoint{0.897186in}{0.869249in}}%
\pgfpathlineto{\pgfqpoint{0.906917in}{0.874476in}}%
\pgfpathclose%
\pgfusepath{fill}%
\end{pgfscope}%
\begin{pgfscope}%
\pgfpathrectangle{\pgfqpoint{0.407592in}{0.222904in}}{\pgfqpoint{1.411163in}{1.411163in}}%
\pgfusepath{clip}%
\pgfsetbuttcap%
\pgfsetroundjoin%
\definecolor{currentfill}{rgb}{0.749020,0.772549,0.792157}%
\pgfsetfillcolor{currentfill}%
\pgfsetfillopacity{0.160000}%
\pgfsetlinewidth{0.000000pt}%
\definecolor{currentstroke}{rgb}{0.000000,0.000000,0.000000}%
\pgfsetstrokecolor{currentstroke}%
\pgfsetdash{}{0pt}%
\pgfpathmoveto{\pgfqpoint{1.078329in}{0.677797in}}%
\pgfpathlineto{\pgfqpoint{1.075367in}{0.694536in}}%
\pgfpathlineto{\pgfqpoint{1.072411in}{0.711244in}}%
\pgfpathlineto{\pgfqpoint{1.067385in}{0.710579in}}%
\pgfpathlineto{\pgfqpoint{1.062455in}{0.709775in}}%
\pgfpathlineto{\pgfqpoint{1.067065in}{0.693304in}}%
\pgfpathlineto{\pgfqpoint{1.071682in}{0.676806in}}%
\pgfpathlineto{\pgfqpoint{1.074973in}{0.677348in}}%
\pgfpathlineto{\pgfqpoint{1.078329in}{0.677797in}}%
\pgfpathclose%
\pgfusepath{fill}%
\end{pgfscope}%
\begin{pgfscope}%
\pgfpathrectangle{\pgfqpoint{0.407592in}{0.222904in}}{\pgfqpoint{1.411163in}{1.411163in}}%
\pgfusepath{clip}%
\pgfsetbuttcap%
\pgfsetroundjoin%
\definecolor{currentfill}{rgb}{0.749020,0.772549,0.792157}%
\pgfsetfillcolor{currentfill}%
\pgfsetfillopacity{0.160000}%
\pgfsetlinewidth{0.000000pt}%
\definecolor{currentstroke}{rgb}{0.000000,0.000000,0.000000}%
\pgfsetstrokecolor{currentstroke}%
\pgfsetdash{}{0pt}%
\pgfpathmoveto{\pgfqpoint{1.276711in}{0.836705in}}%
\pgfpathlineto{\pgfqpoint{1.288334in}{0.850802in}}%
\pgfpathlineto{\pgfqpoint{1.299953in}{0.864894in}}%
\pgfpathlineto{\pgfqpoint{1.292025in}{0.870318in}}%
\pgfpathlineto{\pgfqpoint{1.283252in}{0.875494in}}%
\pgfpathlineto{\pgfqpoint{1.272591in}{0.860858in}}%
\pgfpathlineto{\pgfqpoint{1.261922in}{0.846213in}}%
\pgfpathlineto{\pgfqpoint{1.269694in}{0.841570in}}%
\pgfpathlineto{\pgfqpoint{1.276711in}{0.836705in}}%
\pgfpathclose%
\pgfusepath{fill}%
\end{pgfscope}%
\begin{pgfscope}%
\pgfpathrectangle{\pgfqpoint{0.407592in}{0.222904in}}{\pgfqpoint{1.411163in}{1.411163in}}%
\pgfusepath{clip}%
\pgfsetbuttcap%
\pgfsetroundjoin%
\definecolor{currentfill}{rgb}{0.749020,0.772549,0.792157}%
\pgfsetfillcolor{currentfill}%
\pgfsetfillopacity{0.160000}%
\pgfsetlinewidth{0.000000pt}%
\definecolor{currentstroke}{rgb}{0.000000,0.000000,0.000000}%
\pgfsetstrokecolor{currentstroke}%
\pgfsetdash{}{0pt}%
\pgfpathmoveto{\pgfqpoint{1.118327in}{0.675602in}}%
\pgfpathlineto{\pgfqpoint{1.125325in}{0.691808in}}%
\pgfpathlineto{\pgfqpoint{1.132312in}{0.707989in}}%
\pgfpathlineto{\pgfqpoint{1.127824in}{0.709034in}}%
\pgfpathlineto{\pgfqpoint{1.123175in}{0.709947in}}%
\pgfpathlineto{\pgfqpoint{1.117707in}{0.693449in}}%
\pgfpathlineto{\pgfqpoint{1.112231in}{0.676922in}}%
\pgfpathlineto{\pgfqpoint{1.115333in}{0.676307in}}%
\pgfpathlineto{\pgfqpoint{1.118327in}{0.675602in}}%
\pgfpathclose%
\pgfusepath{fill}%
\end{pgfscope}%
\begin{pgfscope}%
\pgfpathrectangle{\pgfqpoint{0.407592in}{0.222904in}}{\pgfqpoint{1.411163in}{1.411163in}}%
\pgfusepath{clip}%
\pgfsetbuttcap%
\pgfsetroundjoin%
\definecolor{currentfill}{rgb}{0.749020,0.772549,0.792157}%
\pgfsetfillcolor{currentfill}%
\pgfsetfillopacity{0.160000}%
\pgfsetlinewidth{0.000000pt}%
\definecolor{currentstroke}{rgb}{0.000000,0.000000,0.000000}%
\pgfsetstrokecolor{currentstroke}%
\pgfsetdash{}{0pt}%
\pgfpathmoveto{\pgfqpoint{1.148280in}{0.702577in}}%
\pgfpathlineto{\pgfqpoint{1.157921in}{0.717864in}}%
\pgfpathlineto{\pgfqpoint{1.167553in}{0.733136in}}%
\pgfpathlineto{\pgfqpoint{1.162669in}{0.735151in}}%
\pgfpathlineto{\pgfqpoint{1.157476in}{0.737018in}}%
\pgfpathlineto{\pgfqpoint{1.149108in}{0.721279in}}%
\pgfpathlineto{\pgfqpoint{1.140729in}{0.705519in}}%
\pgfpathlineto{\pgfqpoint{1.144621in}{0.704104in}}%
\pgfpathlineto{\pgfqpoint{1.148280in}{0.702577in}}%
\pgfpathclose%
\pgfusepath{fill}%
\end{pgfscope}%
\begin{pgfscope}%
\pgfpathrectangle{\pgfqpoint{0.407592in}{0.222904in}}{\pgfqpoint{1.411163in}{1.411163in}}%
\pgfusepath{clip}%
\pgfsetbuttcap%
\pgfsetroundjoin%
\definecolor{currentfill}{rgb}{0.749020,0.772549,0.792157}%
\pgfsetfillcolor{currentfill}%
\pgfsetfillopacity{0.160000}%
\pgfsetlinewidth{0.000000pt}%
\definecolor{currentstroke}{rgb}{0.000000,0.000000,0.000000}%
\pgfsetstrokecolor{currentstroke}%
\pgfsetdash{}{0pt}%
\pgfpathmoveto{\pgfqpoint{1.071682in}{0.676806in}}%
\pgfpathlineto{\pgfqpoint{1.067065in}{0.693304in}}%
\pgfpathlineto{\pgfqpoint{1.062455in}{0.709775in}}%
\pgfpathlineto{\pgfqpoint{1.057642in}{0.708835in}}%
\pgfpathlineto{\pgfqpoint{1.052967in}{0.707764in}}%
\pgfpathlineto{\pgfqpoint{1.059155in}{0.691619in}}%
\pgfpathlineto{\pgfqpoint{1.065351in}{0.675451in}}%
\pgfpathlineto{\pgfqpoint{1.068470in}{0.676173in}}%
\pgfpathlineto{\pgfqpoint{1.071682in}{0.676806in}}%
\pgfpathclose%
\pgfusepath{fill}%
\end{pgfscope}%
\begin{pgfscope}%
\pgfpathrectangle{\pgfqpoint{0.407592in}{0.222904in}}{\pgfqpoint{1.411163in}{1.411163in}}%
\pgfusepath{clip}%
\pgfsetbuttcap%
\pgfsetroundjoin%
\definecolor{currentfill}{rgb}{0.749020,0.772549,0.792157}%
\pgfsetfillcolor{currentfill}%
\pgfsetfillopacity{0.160000}%
\pgfsetlinewidth{0.000000pt}%
\definecolor{currentstroke}{rgb}{0.000000,0.000000,0.000000}%
\pgfsetstrokecolor{currentstroke}%
\pgfsetdash{}{0pt}%
\pgfpathmoveto{\pgfqpoint{0.927180in}{0.845300in}}%
\pgfpathlineto{\pgfqpoint{0.917045in}{0.859893in}}%
\pgfpathlineto{\pgfqpoint{0.906917in}{0.874476in}}%
\pgfpathlineto{\pgfqpoint{0.897186in}{0.869249in}}%
\pgfpathlineto{\pgfqpoint{0.888254in}{0.863777in}}%
\pgfpathlineto{\pgfqpoint{0.899447in}{0.849743in}}%
\pgfpathlineto{\pgfqpoint{0.910644in}{0.835704in}}%
\pgfpathlineto{\pgfqpoint{0.918555in}{0.840611in}}%
\pgfpathlineto{\pgfqpoint{0.927180in}{0.845300in}}%
\pgfpathclose%
\pgfusepath{fill}%
\end{pgfscope}%
\begin{pgfscope}%
\pgfpathrectangle{\pgfqpoint{0.407592in}{0.222904in}}{\pgfqpoint{1.411163in}{1.411163in}}%
\pgfusepath{clip}%
\pgfsetbuttcap%
\pgfsetroundjoin%
\definecolor{currentfill}{rgb}{0.749020,0.772549,0.792157}%
\pgfsetfillcolor{currentfill}%
\pgfsetfillopacity{0.160000}%
\pgfsetlinewidth{0.000000pt}%
\definecolor{currentstroke}{rgb}{0.000000,0.000000,0.000000}%
\pgfsetstrokecolor{currentstroke}%
\pgfsetdash{}{0pt}%
\pgfpathmoveto{\pgfqpoint{1.044111in}{0.705244in}}%
\pgfpathlineto{\pgfqpoint{1.036458in}{0.720960in}}%
\pgfpathlineto{\pgfqpoint{1.028815in}{0.736656in}}%
\pgfpathlineto{\pgfqpoint{1.023288in}{0.734759in}}%
\pgfpathlineto{\pgfqpoint{1.018044in}{0.732714in}}%
\pgfpathlineto{\pgfqpoint{1.027036in}{0.717494in}}%
\pgfpathlineto{\pgfqpoint{1.036038in}{0.702258in}}%
\pgfpathlineto{\pgfqpoint{1.039967in}{0.703806in}}%
\pgfpathlineto{\pgfqpoint{1.044111in}{0.705244in}}%
\pgfpathclose%
\pgfusepath{fill}%
\end{pgfscope}%
\begin{pgfscope}%
\pgfpathrectangle{\pgfqpoint{0.407592in}{0.222904in}}{\pgfqpoint{1.411163in}{1.411163in}}%
\pgfusepath{clip}%
\pgfsetbuttcap%
\pgfsetroundjoin%
\definecolor{currentfill}{rgb}{0.749020,0.772549,0.792157}%
\pgfsetfillcolor{currentfill}%
\pgfsetfillopacity{0.160000}%
\pgfsetlinewidth{0.000000pt}%
\definecolor{currentstroke}{rgb}{0.000000,0.000000,0.000000}%
\pgfsetstrokecolor{currentstroke}%
\pgfsetdash{}{0pt}%
\pgfpathmoveto{\pgfqpoint{1.176309in}{0.728694in}}%
\pgfpathlineto{\pgfqpoint{1.187036in}{0.743419in}}%
\pgfpathlineto{\pgfqpoint{1.197756in}{0.758134in}}%
\pgfpathlineto{\pgfqpoint{1.192489in}{0.760965in}}%
\pgfpathlineto{\pgfqpoint{1.186787in}{0.763633in}}%
\pgfpathlineto{\pgfqpoint{1.177174in}{0.748392in}}%
\pgfpathlineto{\pgfqpoint{1.167553in}{0.733136in}}%
\pgfpathlineto{\pgfqpoint{1.172106in}{0.730980in}}%
\pgfpathlineto{\pgfqpoint{1.176309in}{0.728694in}}%
\pgfpathclose%
\pgfusepath{fill}%
\end{pgfscope}%
\begin{pgfscope}%
\pgfpathrectangle{\pgfqpoint{0.407592in}{0.222904in}}{\pgfqpoint{1.411163in}{1.411163in}}%
\pgfusepath{clip}%
\pgfsetbuttcap%
\pgfsetroundjoin%
\definecolor{currentfill}{rgb}{0.749020,0.772549,0.792157}%
\pgfsetfillcolor{currentfill}%
\pgfsetfillopacity{0.160000}%
\pgfsetlinewidth{0.000000pt}%
\definecolor{currentstroke}{rgb}{0.000000,0.000000,0.000000}%
\pgfsetstrokecolor{currentstroke}%
\pgfsetdash{}{0pt}%
\pgfpathmoveto{\pgfqpoint{1.253454in}{0.808498in}}%
\pgfpathlineto{\pgfqpoint{1.265085in}{0.822604in}}%
\pgfpathlineto{\pgfqpoint{1.276711in}{0.836705in}}%
\pgfpathlineto{\pgfqpoint{1.269694in}{0.841570in}}%
\pgfpathlineto{\pgfqpoint{1.261922in}{0.846213in}}%
\pgfpathlineto{\pgfqpoint{1.251246in}{0.831558in}}%
\pgfpathlineto{\pgfqpoint{1.240562in}{0.816893in}}%
\pgfpathlineto{\pgfqpoint{1.247340in}{0.812793in}}%
\pgfpathlineto{\pgfqpoint{1.253454in}{0.808498in}}%
\pgfpathclose%
\pgfusepath{fill}%
\end{pgfscope}%
\begin{pgfscope}%
\pgfpathrectangle{\pgfqpoint{0.407592in}{0.222904in}}{\pgfqpoint{1.411163in}{1.411163in}}%
\pgfusepath{clip}%
\pgfsetbuttcap%
\pgfsetroundjoin%
\definecolor{currentfill}{rgb}{0.749020,0.772549,0.792157}%
\pgfsetfillcolor{currentfill}%
\pgfsetfillopacity{0.160000}%
\pgfsetlinewidth{0.000000pt}%
\definecolor{currentstroke}{rgb}{0.000000,0.000000,0.000000}%
\pgfsetstrokecolor{currentstroke}%
\pgfsetdash{}{0pt}%
\pgfpathmoveto{\pgfqpoint{1.784494in}{1.366922in}}%
\pgfpathlineto{\pgfqpoint{1.796912in}{1.380452in}}%
\pgfpathlineto{\pgfqpoint{1.774857in}{1.396026in}}%
\pgfpathlineto{\pgfqpoint{1.750178in}{1.410942in}}%
\pgfpathlineto{\pgfqpoint{1.738704in}{1.397027in}}%
\pgfpathlineto{\pgfqpoint{1.762892in}{1.382298in}}%
\pgfusepath{fill}%
\end{pgfscope}%
\begin{pgfscope}%
\pgfpathrectangle{\pgfqpoint{0.407592in}{0.222904in}}{\pgfqpoint{1.411163in}{1.411163in}}%
\pgfusepath{clip}%
\pgfsetbuttcap%
\pgfsetroundjoin%
\definecolor{currentfill}{rgb}{0.749020,0.772549,0.792157}%
\pgfsetfillcolor{currentfill}%
\pgfsetfillopacity{0.160000}%
\pgfsetlinewidth{0.000000pt}%
\definecolor{currentstroke}{rgb}{0.000000,0.000000,0.000000}%
\pgfsetstrokecolor{currentstroke}%
\pgfsetdash{}{0pt}%
\pgfpathmoveto{\pgfqpoint{1.759659in}{1.339865in}}%
\pgfpathlineto{\pgfqpoint{1.772076in}{1.353393in}}%
\pgfpathlineto{\pgfqpoint{1.784494in}{1.366922in}}%
\pgfpathlineto{\pgfqpoint{1.762892in}{1.382298in}}%
\pgfpathlineto{\pgfqpoint{1.738704in}{1.397027in}}%
\pgfpathlineto{\pgfqpoint{1.727227in}{1.383107in}}%
\pgfpathlineto{\pgfqpoint{1.715746in}{1.369183in}}%
\pgfpathlineto{\pgfqpoint{1.738956in}{1.354836in}}%
\pgfpathlineto{\pgfqpoint{1.759659in}{1.339865in}}%
\pgfpathclose%
\pgfusepath{fill}%
\end{pgfscope}%
\begin{pgfscope}%
\pgfpathrectangle{\pgfqpoint{0.407592in}{0.222904in}}{\pgfqpoint{1.411163in}{1.411163in}}%
\pgfusepath{clip}%
\pgfsetbuttcap%
\pgfsetroundjoin%
\definecolor{currentfill}{rgb}{0.749020,0.772549,0.792157}%
\pgfsetfillcolor{currentfill}%
\pgfsetfillopacity{0.160000}%
\pgfsetlinewidth{0.000000pt}%
\definecolor{currentstroke}{rgb}{0.000000,0.000000,0.000000}%
\pgfsetstrokecolor{currentstroke}%
\pgfsetdash{}{0pt}%
\pgfpathmoveto{\pgfqpoint{1.018044in}{0.732714in}}%
\pgfpathlineto{\pgfqpoint{1.009060in}{0.747920in}}%
\pgfpathlineto{\pgfqpoint{1.000085in}{0.763112in}}%
\pgfpathlineto{\pgfqpoint{0.993903in}{0.760410in}}%
\pgfpathlineto{\pgfqpoint{0.988126in}{0.757547in}}%
\pgfpathlineto{\pgfqpoint{0.998307in}{0.742888in}}%
\pgfpathlineto{\pgfqpoint{1.008494in}{0.728220in}}%
\pgfpathlineto{\pgfqpoint{1.013105in}{0.730532in}}%
\pgfpathlineto{\pgfqpoint{1.018044in}{0.732714in}}%
\pgfpathclose%
\pgfusepath{fill}%
\end{pgfscope}%
\begin{pgfscope}%
\pgfpathrectangle{\pgfqpoint{0.407592in}{0.222904in}}{\pgfqpoint{1.411163in}{1.411163in}}%
\pgfusepath{clip}%
\pgfsetbuttcap%
\pgfsetroundjoin%
\definecolor{currentfill}{rgb}{0.749020,0.772549,0.792157}%
\pgfsetfillcolor{currentfill}%
\pgfsetfillopacity{0.160000}%
\pgfsetlinewidth{0.000000pt}%
\definecolor{currentstroke}{rgb}{0.000000,0.000000,0.000000}%
\pgfsetstrokecolor{currentstroke}%
\pgfsetdash{}{0pt}%
\pgfpathmoveto{\pgfqpoint{1.734827in}{1.312811in}}%
\pgfpathlineto{\pgfqpoint{1.747242in}{1.326337in}}%
\pgfpathlineto{\pgfqpoint{1.759659in}{1.339865in}}%
\pgfpathlineto{\pgfqpoint{1.738956in}{1.354836in}}%
\pgfpathlineto{\pgfqpoint{1.715746in}{1.369183in}}%
\pgfpathlineto{\pgfqpoint{1.704262in}{1.355254in}}%
\pgfpathlineto{\pgfqpoint{1.692774in}{1.341320in}}%
\pgfpathlineto{\pgfqpoint{1.715014in}{1.327366in}}%
\pgfpathlineto{\pgfqpoint{1.734827in}{1.312811in}}%
\pgfpathclose%
\pgfusepath{fill}%
\end{pgfscope}%
\begin{pgfscope}%
\pgfpathrectangle{\pgfqpoint{0.407592in}{0.222904in}}{\pgfqpoint{1.411163in}{1.411163in}}%
\pgfusepath{clip}%
\pgfsetbuttcap%
\pgfsetroundjoin%
\definecolor{currentfill}{rgb}{0.749020,0.772549,0.792157}%
\pgfsetfillcolor{currentfill}%
\pgfsetfillopacity{0.160000}%
\pgfsetlinewidth{0.000000pt}%
\definecolor{currentstroke}{rgb}{0.000000,0.000000,0.000000}%
\pgfsetstrokecolor{currentstroke}%
\pgfsetdash{}{0pt}%
\pgfpathmoveto{\pgfqpoint{1.123939in}{0.673938in}}%
\pgfpathlineto{\pgfqpoint{1.132340in}{0.689739in}}%
\pgfpathlineto{\pgfqpoint{1.140729in}{0.705519in}}%
\pgfpathlineto{\pgfqpoint{1.136620in}{0.706815in}}%
\pgfpathlineto{\pgfqpoint{1.132312in}{0.707989in}}%
\pgfpathlineto{\pgfqpoint{1.125325in}{0.691808in}}%
\pgfpathlineto{\pgfqpoint{1.118327in}{0.675602in}}%
\pgfpathlineto{\pgfqpoint{1.121200in}{0.674812in}}%
\pgfpathlineto{\pgfqpoint{1.123939in}{0.673938in}}%
\pgfpathclose%
\pgfusepath{fill}%
\end{pgfscope}%
\begin{pgfscope}%
\pgfpathrectangle{\pgfqpoint{0.407592in}{0.222904in}}{\pgfqpoint{1.411163in}{1.411163in}}%
\pgfusepath{clip}%
\pgfsetbuttcap%
\pgfsetroundjoin%
\definecolor{currentfill}{rgb}{0.749020,0.772549,0.792157}%
\pgfsetfillcolor{currentfill}%
\pgfsetfillopacity{0.160000}%
\pgfsetlinewidth{0.000000pt}%
\definecolor{currentstroke}{rgb}{0.000000,0.000000,0.000000}%
\pgfsetstrokecolor{currentstroke}%
\pgfsetdash{}{0pt}%
\pgfpathmoveto{\pgfqpoint{0.947469in}{0.816087in}}%
\pgfpathlineto{\pgfqpoint{0.937321in}{0.830698in}}%
\pgfpathlineto{\pgfqpoint{0.927180in}{0.845300in}}%
\pgfpathlineto{\pgfqpoint{0.918555in}{0.840611in}}%
\pgfpathlineto{\pgfqpoint{0.910644in}{0.835704in}}%
\pgfpathlineto{\pgfqpoint{0.921844in}{0.821661in}}%
\pgfpathlineto{\pgfqpoint{0.933047in}{0.807614in}}%
\pgfpathlineto{\pgfqpoint{0.939943in}{0.811946in}}%
\pgfpathlineto{\pgfqpoint{0.947469in}{0.816087in}}%
\pgfpathclose%
\pgfusepath{fill}%
\end{pgfscope}%
\begin{pgfscope}%
\pgfpathrectangle{\pgfqpoint{0.407592in}{0.222904in}}{\pgfqpoint{1.411163in}{1.411163in}}%
\pgfusepath{clip}%
\pgfsetbuttcap%
\pgfsetroundjoin%
\definecolor{currentfill}{rgb}{0.749020,0.772549,0.792157}%
\pgfsetfillcolor{currentfill}%
\pgfsetfillopacity{0.160000}%
\pgfsetlinewidth{0.000000pt}%
\definecolor{currentstroke}{rgb}{0.000000,0.000000,0.000000}%
\pgfsetstrokecolor{currentstroke}%
\pgfsetdash{}{0pt}%
\pgfpathmoveto{\pgfqpoint{1.065351in}{0.675451in}}%
\pgfpathlineto{\pgfqpoint{1.059155in}{0.691619in}}%
\pgfpathlineto{\pgfqpoint{1.052967in}{0.707764in}}%
\pgfpathlineto{\pgfqpoint{1.048450in}{0.706565in}}%
\pgfpathlineto{\pgfqpoint{1.044111in}{0.705244in}}%
\pgfpathlineto{\pgfqpoint{1.051773in}{0.689509in}}%
\pgfpathlineto{\pgfqpoint{1.059445in}{0.673753in}}%
\pgfpathlineto{\pgfqpoint{1.062338in}{0.674643in}}%
\pgfpathlineto{\pgfqpoint{1.065351in}{0.675451in}}%
\pgfpathclose%
\pgfusepath{fill}%
\end{pgfscope}%
\begin{pgfscope}%
\pgfpathrectangle{\pgfqpoint{0.407592in}{0.222904in}}{\pgfqpoint{1.411163in}{1.411163in}}%
\pgfusepath{clip}%
\pgfsetbuttcap%
\pgfsetroundjoin%
\definecolor{currentfill}{rgb}{0.749020,0.772549,0.792157}%
\pgfsetfillcolor{currentfill}%
\pgfsetfillopacity{0.160000}%
\pgfsetlinewidth{0.000000pt}%
\definecolor{currentstroke}{rgb}{0.000000,0.000000,0.000000}%
\pgfsetstrokecolor{currentstroke}%
\pgfsetdash{}{0pt}%
\pgfpathmoveto{\pgfqpoint{1.709998in}{1.285760in}}%
\pgfpathlineto{\pgfqpoint{1.722412in}{1.299285in}}%
\pgfpathlineto{\pgfqpoint{1.734827in}{1.312811in}}%
\pgfpathlineto{\pgfqpoint{1.715014in}{1.327366in}}%
\pgfpathlineto{\pgfqpoint{1.692774in}{1.341320in}}%
\pgfpathlineto{\pgfqpoint{1.681282in}{1.327382in}}%
\pgfpathlineto{\pgfqpoint{1.669786in}{1.313440in}}%
\pgfpathlineto{\pgfqpoint{1.691066in}{1.299889in}}%
\pgfpathlineto{\pgfqpoint{1.709998in}{1.285760in}}%
\pgfpathclose%
\pgfusepath{fill}%
\end{pgfscope}%
\begin{pgfscope}%
\pgfpathrectangle{\pgfqpoint{0.407592in}{0.222904in}}{\pgfqpoint{1.411163in}{1.411163in}}%
\pgfusepath{clip}%
\pgfsetbuttcap%
\pgfsetroundjoin%
\definecolor{currentfill}{rgb}{0.749020,0.772549,0.792157}%
\pgfsetfillcolor{currentfill}%
\pgfsetfillopacity{0.160000}%
\pgfsetlinewidth{0.000000pt}%
\definecolor{currentstroke}{rgb}{0.000000,0.000000,0.000000}%
\pgfsetstrokecolor{currentstroke}%
\pgfsetdash{}{0pt}%
\pgfpathmoveto{\pgfqpoint{1.230182in}{0.780272in}}%
\pgfpathlineto{\pgfqpoint{1.241820in}{0.794387in}}%
\pgfpathlineto{\pgfqpoint{1.253454in}{0.808498in}}%
\pgfpathlineto{\pgfqpoint{1.247340in}{0.812793in}}%
\pgfpathlineto{\pgfqpoint{1.240562in}{0.816893in}}%
\pgfpathlineto{\pgfqpoint{1.229872in}{0.802218in}}%
\pgfpathlineto{\pgfqpoint{1.219174in}{0.787533in}}%
\pgfpathlineto{\pgfqpoint{1.224964in}{0.783986in}}%
\pgfpathlineto{\pgfqpoint{1.230182in}{0.780272in}}%
\pgfpathclose%
\pgfusepath{fill}%
\end{pgfscope}%
\begin{pgfscope}%
\pgfpathrectangle{\pgfqpoint{0.407592in}{0.222904in}}{\pgfqpoint{1.411163in}{1.411163in}}%
\pgfusepath{clip}%
\pgfsetbuttcap%
\pgfsetroundjoin%
\definecolor{currentfill}{rgb}{0.749020,0.772549,0.792157}%
\pgfsetfillcolor{currentfill}%
\pgfsetfillopacity{0.160000}%
\pgfsetlinewidth{0.000000pt}%
\definecolor{currentstroke}{rgb}{0.000000,0.000000,0.000000}%
\pgfsetstrokecolor{currentstroke}%
\pgfsetdash{}{0pt}%
\pgfpathmoveto{\pgfqpoint{1.685171in}{1.258712in}}%
\pgfpathlineto{\pgfqpoint{1.697584in}{1.272236in}}%
\pgfpathlineto{\pgfqpoint{1.709998in}{1.285760in}}%
\pgfpathlineto{\pgfqpoint{1.691066in}{1.299889in}}%
\pgfpathlineto{\pgfqpoint{1.669786in}{1.313440in}}%
\pgfpathlineto{\pgfqpoint{1.658286in}{1.299493in}}%
\pgfpathlineto{\pgfqpoint{1.646783in}{1.285542in}}%
\pgfpathlineto{\pgfqpoint{1.667110in}{1.272405in}}%
\pgfpathlineto{\pgfqpoint{1.685171in}{1.258712in}}%
\pgfpathclose%
\pgfusepath{fill}%
\end{pgfscope}%
\begin{pgfscope}%
\pgfpathrectangle{\pgfqpoint{0.407592in}{0.222904in}}{\pgfqpoint{1.411163in}{1.411163in}}%
\pgfusepath{clip}%
\pgfsetbuttcap%
\pgfsetroundjoin%
\definecolor{currentfill}{rgb}{0.749020,0.772549,0.792157}%
\pgfsetfillcolor{currentfill}%
\pgfsetfillopacity{0.160000}%
\pgfsetlinewidth{0.000000pt}%
\definecolor{currentstroke}{rgb}{0.000000,0.000000,0.000000}%
\pgfsetstrokecolor{currentstroke}%
\pgfsetdash{}{0pt}%
\pgfpathmoveto{\pgfqpoint{1.660348in}{1.231668in}}%
\pgfpathlineto{\pgfqpoint{1.672759in}{1.245190in}}%
\pgfpathlineto{\pgfqpoint{1.685171in}{1.258712in}}%
\pgfpathlineto{\pgfqpoint{1.667110in}{1.272405in}}%
\pgfpathlineto{\pgfqpoint{1.646783in}{1.285542in}}%
\pgfpathlineto{\pgfqpoint{1.635277in}{1.271586in}}%
\pgfpathlineto{\pgfqpoint{1.623766in}{1.257625in}}%
\pgfpathlineto{\pgfqpoint{1.643149in}{1.244913in}}%
\pgfpathlineto{\pgfqpoint{1.660348in}{1.231668in}}%
\pgfpathclose%
\pgfusepath{fill}%
\end{pgfscope}%
\begin{pgfscope}%
\pgfpathrectangle{\pgfqpoint{0.407592in}{0.222904in}}{\pgfqpoint{1.411163in}{1.411163in}}%
\pgfusepath{clip}%
\pgfsetbuttcap%
\pgfsetroundjoin%
\definecolor{currentfill}{rgb}{0.749020,0.772549,0.792157}%
\pgfsetfillcolor{currentfill}%
\pgfsetfillopacity{0.160000}%
\pgfsetlinewidth{0.000000pt}%
\definecolor{currentstroke}{rgb}{0.000000,0.000000,0.000000}%
\pgfsetstrokecolor{currentstroke}%
\pgfsetdash{}{0pt}%
\pgfpathmoveto{\pgfqpoint{0.967785in}{0.786836in}}%
\pgfpathlineto{\pgfqpoint{0.957624in}{0.801466in}}%
\pgfpathlineto{\pgfqpoint{0.947469in}{0.816087in}}%
\pgfpathlineto{\pgfqpoint{0.939943in}{0.811946in}}%
\pgfpathlineto{\pgfqpoint{0.933047in}{0.807614in}}%
\pgfpathlineto{\pgfqpoint{0.944253in}{0.793563in}}%
\pgfpathlineto{\pgfqpoint{0.955463in}{0.779508in}}%
\pgfpathlineto{\pgfqpoint{0.961352in}{0.783254in}}%
\pgfpathlineto{\pgfqpoint{0.967785in}{0.786836in}}%
\pgfpathclose%
\pgfusepath{fill}%
\end{pgfscope}%
\begin{pgfscope}%
\pgfpathrectangle{\pgfqpoint{0.407592in}{0.222904in}}{\pgfqpoint{1.411163in}{1.411163in}}%
\pgfusepath{clip}%
\pgfsetbuttcap%
\pgfsetroundjoin%
\definecolor{currentfill}{rgb}{0.749020,0.772549,0.792157}%
\pgfsetfillcolor{currentfill}%
\pgfsetfillopacity{0.160000}%
\pgfsetlinewidth{0.000000pt}%
\definecolor{currentstroke}{rgb}{0.000000,0.000000,0.000000}%
\pgfsetstrokecolor{currentstroke}%
\pgfsetdash{}{0pt}%
\pgfpathmoveto{\pgfqpoint{1.635528in}{1.204627in}}%
\pgfpathlineto{\pgfqpoint{1.647938in}{1.218147in}}%
\pgfpathlineto{\pgfqpoint{1.660348in}{1.231668in}}%
\pgfpathlineto{\pgfqpoint{1.643149in}{1.244913in}}%
\pgfpathlineto{\pgfqpoint{1.623766in}{1.257625in}}%
\pgfpathlineto{\pgfqpoint{1.612252in}{1.243660in}}%
\pgfpathlineto{\pgfqpoint{1.600734in}{1.229691in}}%
\pgfpathlineto{\pgfqpoint{1.619181in}{1.217414in}}%
\pgfpathlineto{\pgfqpoint{1.635528in}{1.204627in}}%
\pgfpathclose%
\pgfusepath{fill}%
\end{pgfscope}%
\begin{pgfscope}%
\pgfpathrectangle{\pgfqpoint{0.407592in}{0.222904in}}{\pgfqpoint{1.411163in}{1.411163in}}%
\pgfusepath{clip}%
\pgfsetbuttcap%
\pgfsetroundjoin%
\definecolor{currentfill}{rgb}{0.749020,0.772549,0.792157}%
\pgfsetfillcolor{currentfill}%
\pgfsetfillopacity{0.160000}%
\pgfsetlinewidth{0.000000pt}%
\definecolor{currentstroke}{rgb}{0.000000,0.000000,0.000000}%
\pgfsetstrokecolor{currentstroke}%
\pgfsetdash{}{0pt}%
\pgfpathmoveto{\pgfqpoint{1.154833in}{0.699214in}}%
\pgfpathlineto{\pgfqpoint{1.165575in}{0.713959in}}%
\pgfpathlineto{\pgfqpoint{1.176309in}{0.728694in}}%
\pgfpathlineto{\pgfqpoint{1.172106in}{0.730980in}}%
\pgfpathlineto{\pgfqpoint{1.167553in}{0.733136in}}%
\pgfpathlineto{\pgfqpoint{1.157921in}{0.717864in}}%
\pgfpathlineto{\pgfqpoint{1.148280in}{0.702577in}}%
\pgfpathlineto{\pgfqpoint{1.151689in}{0.700945in}}%
\pgfpathlineto{\pgfqpoint{1.154833in}{0.699214in}}%
\pgfpathclose%
\pgfusepath{fill}%
\end{pgfscope}%
\begin{pgfscope}%
\pgfpathrectangle{\pgfqpoint{0.407592in}{0.222904in}}{\pgfqpoint{1.411163in}{1.411163in}}%
\pgfusepath{clip}%
\pgfsetbuttcap%
\pgfsetroundjoin%
\definecolor{currentfill}{rgb}{0.749020,0.772549,0.792157}%
\pgfsetfillcolor{currentfill}%
\pgfsetfillopacity{0.160000}%
\pgfsetlinewidth{0.000000pt}%
\definecolor{currentstroke}{rgb}{0.000000,0.000000,0.000000}%
\pgfsetstrokecolor{currentstroke}%
\pgfsetdash{}{0pt}%
\pgfpathmoveto{\pgfqpoint{1.094627in}{0.644555in}}%
\pgfpathlineto{\pgfqpoint{1.096820in}{0.661517in}}%
\pgfpathlineto{\pgfqpoint{1.099008in}{0.678446in}}%
\pgfpathlineto{\pgfqpoint{1.095571in}{0.678584in}}%
\pgfpathlineto{\pgfqpoint{1.092110in}{0.678624in}}%
\pgfpathlineto{\pgfqpoint{1.091641in}{0.661651in}}%
\pgfpathlineto{\pgfqpoint{1.091172in}{0.644645in}}%
\pgfpathlineto{\pgfqpoint{1.092905in}{0.644625in}}%
\pgfpathlineto{\pgfqpoint{1.094627in}{0.644555in}}%
\pgfpathclose%
\pgfusepath{fill}%
\end{pgfscope}%
\begin{pgfscope}%
\pgfpathrectangle{\pgfqpoint{0.407592in}{0.222904in}}{\pgfqpoint{1.411163in}{1.411163in}}%
\pgfusepath{clip}%
\pgfsetbuttcap%
\pgfsetroundjoin%
\definecolor{currentfill}{rgb}{0.749020,0.772549,0.792157}%
\pgfsetfillcolor{currentfill}%
\pgfsetfillopacity{0.160000}%
\pgfsetlinewidth{0.000000pt}%
\definecolor{currentstroke}{rgb}{0.000000,0.000000,0.000000}%
\pgfsetstrokecolor{currentstroke}%
\pgfsetdash{}{0pt}%
\pgfpathmoveto{\pgfqpoint{1.091172in}{0.644645in}}%
\pgfpathlineto{\pgfqpoint{1.091641in}{0.661651in}}%
\pgfpathlineto{\pgfqpoint{1.092110in}{0.678624in}}%
\pgfpathlineto{\pgfqpoint{1.088641in}{0.678564in}}%
\pgfpathlineto{\pgfqpoint{1.085177in}{0.678406in}}%
\pgfpathlineto{\pgfqpoint{1.086437in}{0.661487in}}%
\pgfpathlineto{\pgfqpoint{1.087699in}{0.644535in}}%
\pgfpathlineto{\pgfqpoint{1.089434in}{0.644615in}}%
\pgfpathlineto{\pgfqpoint{1.091172in}{0.644645in}}%
\pgfpathclose%
\pgfusepath{fill}%
\end{pgfscope}%
\begin{pgfscope}%
\pgfpathrectangle{\pgfqpoint{0.407592in}{0.222904in}}{\pgfqpoint{1.411163in}{1.411163in}}%
\pgfusepath{clip}%
\pgfsetbuttcap%
\pgfsetroundjoin%
\definecolor{currentfill}{rgb}{0.749020,0.772549,0.792157}%
\pgfsetfillcolor{currentfill}%
\pgfsetfillopacity{0.160000}%
\pgfsetlinewidth{0.000000pt}%
\definecolor{currentstroke}{rgb}{0.000000,0.000000,0.000000}%
\pgfsetstrokecolor{currentstroke}%
\pgfsetdash{}{0pt}%
\pgfpathmoveto{\pgfqpoint{1.128969in}{0.671957in}}%
\pgfpathlineto{\pgfqpoint{1.138629in}{0.687275in}}%
\pgfpathlineto{\pgfqpoint{1.148280in}{0.702577in}}%
\pgfpathlineto{\pgfqpoint{1.144621in}{0.704104in}}%
\pgfpathlineto{\pgfqpoint{1.140729in}{0.705519in}}%
\pgfpathlineto{\pgfqpoint{1.132340in}{0.689739in}}%
\pgfpathlineto{\pgfqpoint{1.123939in}{0.673938in}}%
\pgfpathlineto{\pgfqpoint{1.126532in}{0.672985in}}%
\pgfpathlineto{\pgfqpoint{1.128969in}{0.671957in}}%
\pgfpathclose%
\pgfusepath{fill}%
\end{pgfscope}%
\begin{pgfscope}%
\pgfpathrectangle{\pgfqpoint{0.407592in}{0.222904in}}{\pgfqpoint{1.411163in}{1.411163in}}%
\pgfusepath{clip}%
\pgfsetbuttcap%
\pgfsetroundjoin%
\definecolor{currentfill}{rgb}{0.749020,0.772549,0.792157}%
\pgfsetfillcolor{currentfill}%
\pgfsetfillopacity{0.160000}%
\pgfsetlinewidth{0.000000pt}%
\definecolor{currentstroke}{rgb}{0.000000,0.000000,0.000000}%
\pgfsetstrokecolor{currentstroke}%
\pgfsetdash{}{0pt}%
\pgfpathmoveto{\pgfqpoint{1.610711in}{1.177589in}}%
\pgfpathlineto{\pgfqpoint{1.623119in}{1.191107in}}%
\pgfpathlineto{\pgfqpoint{1.635528in}{1.204627in}}%
\pgfpathlineto{\pgfqpoint{1.619181in}{1.217414in}}%
\pgfpathlineto{\pgfqpoint{1.600734in}{1.229691in}}%
\pgfpathlineto{\pgfqpoint{1.589212in}{1.215717in}}%
\pgfpathlineto{\pgfqpoint{1.577687in}{1.201739in}}%
\pgfpathlineto{\pgfqpoint{1.595206in}{1.189907in}}%
\pgfpathlineto{\pgfqpoint{1.610711in}{1.177589in}}%
\pgfpathclose%
\pgfusepath{fill}%
\end{pgfscope}%
\begin{pgfscope}%
\pgfpathrectangle{\pgfqpoint{0.407592in}{0.222904in}}{\pgfqpoint{1.411163in}{1.411163in}}%
\pgfusepath{clip}%
\pgfsetbuttcap%
\pgfsetroundjoin%
\definecolor{currentfill}{rgb}{0.749020,0.772549,0.792157}%
\pgfsetfillcolor{currentfill}%
\pgfsetfillopacity{0.160000}%
\pgfsetlinewidth{0.000000pt}%
\definecolor{currentstroke}{rgb}{0.000000,0.000000,0.000000}%
\pgfsetstrokecolor{currentstroke}%
\pgfsetdash{}{0pt}%
\pgfpathmoveto{\pgfqpoint{1.036038in}{0.702258in}}%
\pgfpathlineto{\pgfqpoint{1.027036in}{0.717494in}}%
\pgfpathlineto{\pgfqpoint{1.018044in}{0.732714in}}%
\pgfpathlineto{\pgfqpoint{1.013105in}{0.730532in}}%
\pgfpathlineto{\pgfqpoint{1.008494in}{0.728220in}}%
\pgfpathlineto{\pgfqpoint{1.018688in}{0.713543in}}%
\pgfpathlineto{\pgfqpoint{1.028888in}{0.698855in}}%
\pgfpathlineto{\pgfqpoint{1.032339in}{0.700605in}}%
\pgfpathlineto{\pgfqpoint{1.036038in}{0.702258in}}%
\pgfpathclose%
\pgfusepath{fill}%
\end{pgfscope}%
\begin{pgfscope}%
\pgfpathrectangle{\pgfqpoint{0.407592in}{0.222904in}}{\pgfqpoint{1.411163in}{1.411163in}}%
\pgfusepath{clip}%
\pgfsetbuttcap%
\pgfsetroundjoin%
\definecolor{currentfill}{rgb}{0.749020,0.772549,0.792157}%
\pgfsetfillcolor{currentfill}%
\pgfsetfillopacity{0.160000}%
\pgfsetlinewidth{0.000000pt}%
\definecolor{currentstroke}{rgb}{0.000000,0.000000,0.000000}%
\pgfsetstrokecolor{currentstroke}%
\pgfsetdash{}{0pt}%
\pgfpathmoveto{\pgfqpoint{1.098006in}{0.644266in}}%
\pgfpathlineto{\pgfqpoint{1.101883in}{0.661086in}}%
\pgfpathlineto{\pgfqpoint{1.105754in}{0.677875in}}%
\pgfpathlineto{\pgfqpoint{1.102407in}{0.678209in}}%
\pgfpathlineto{\pgfqpoint{1.099008in}{0.678446in}}%
\pgfpathlineto{\pgfqpoint{1.096820in}{0.661517in}}%
\pgfpathlineto{\pgfqpoint{1.094627in}{0.644555in}}%
\pgfpathlineto{\pgfqpoint{1.096330in}{0.644435in}}%
\pgfpathlineto{\pgfqpoint{1.098006in}{0.644266in}}%
\pgfpathclose%
\pgfusepath{fill}%
\end{pgfscope}%
\begin{pgfscope}%
\pgfpathrectangle{\pgfqpoint{0.407592in}{0.222904in}}{\pgfqpoint{1.411163in}{1.411163in}}%
\pgfusepath{clip}%
\pgfsetbuttcap%
\pgfsetroundjoin%
\definecolor{currentfill}{rgb}{0.749020,0.772549,0.792157}%
\pgfsetfillcolor{currentfill}%
\pgfsetfillopacity{0.160000}%
\pgfsetlinewidth{0.000000pt}%
\definecolor{currentstroke}{rgb}{0.000000,0.000000,0.000000}%
\pgfsetstrokecolor{currentstroke}%
\pgfsetdash{}{0pt}%
\pgfpathmoveto{\pgfqpoint{1.087699in}{0.644535in}}%
\pgfpathlineto{\pgfqpoint{1.086437in}{0.661487in}}%
\pgfpathlineto{\pgfqpoint{1.085177in}{0.678406in}}%
\pgfpathlineto{\pgfqpoint{1.081735in}{0.678150in}}%
\pgfpathlineto{\pgfqpoint{1.078329in}{0.677797in}}%
\pgfpathlineto{\pgfqpoint{1.081296in}{0.661027in}}%
\pgfpathlineto{\pgfqpoint{1.084268in}{0.644227in}}%
\pgfpathlineto{\pgfqpoint{1.085974in}{0.644405in}}%
\pgfpathlineto{\pgfqpoint{1.087699in}{0.644535in}}%
\pgfpathclose%
\pgfusepath{fill}%
\end{pgfscope}%
\begin{pgfscope}%
\pgfpathrectangle{\pgfqpoint{0.407592in}{0.222904in}}{\pgfqpoint{1.411163in}{1.411163in}}%
\pgfusepath{clip}%
\pgfsetbuttcap%
\pgfsetroundjoin%
\definecolor{currentfill}{rgb}{0.749020,0.772549,0.792157}%
\pgfsetfillcolor{currentfill}%
\pgfsetfillopacity{0.160000}%
\pgfsetlinewidth{0.000000pt}%
\definecolor{currentstroke}{rgb}{0.000000,0.000000,0.000000}%
\pgfsetstrokecolor{currentstroke}%
\pgfsetdash{}{0pt}%
\pgfpathmoveto{\pgfqpoint{1.059445in}{0.673753in}}%
\pgfpathlineto{\pgfqpoint{1.051773in}{0.689509in}}%
\pgfpathlineto{\pgfqpoint{1.044111in}{0.705244in}}%
\pgfpathlineto{\pgfqpoint{1.039967in}{0.703806in}}%
\pgfpathlineto{\pgfqpoint{1.036038in}{0.702258in}}%
\pgfpathlineto{\pgfqpoint{1.045048in}{0.687008in}}%
\pgfpathlineto{\pgfqpoint{1.054067in}{0.671742in}}%
\pgfpathlineto{\pgfqpoint{1.056684in}{0.672785in}}%
\pgfpathlineto{\pgfqpoint{1.059445in}{0.673753in}}%
\pgfpathclose%
\pgfusepath{fill}%
\end{pgfscope}%
\begin{pgfscope}%
\pgfpathrectangle{\pgfqpoint{0.407592in}{0.222904in}}{\pgfqpoint{1.411163in}{1.411163in}}%
\pgfusepath{clip}%
\pgfsetbuttcap%
\pgfsetroundjoin%
\definecolor{currentfill}{rgb}{0.749020,0.772549,0.792157}%
\pgfsetfillcolor{currentfill}%
\pgfsetfillopacity{0.160000}%
\pgfsetlinewidth{0.000000pt}%
\definecolor{currentstroke}{rgb}{0.000000,0.000000,0.000000}%
\pgfsetstrokecolor{currentstroke}%
\pgfsetdash{}{0pt}%
\pgfpathmoveto{\pgfqpoint{1.206894in}{0.752028in}}%
\pgfpathlineto{\pgfqpoint{1.218540in}{0.766152in}}%
\pgfpathlineto{\pgfqpoint{1.230182in}{0.780272in}}%
\pgfpathlineto{\pgfqpoint{1.224964in}{0.783986in}}%
\pgfpathlineto{\pgfqpoint{1.219174in}{0.787533in}}%
\pgfpathlineto{\pgfqpoint{1.208469in}{0.772839in}}%
\pgfpathlineto{\pgfqpoint{1.197756in}{0.758134in}}%
\pgfpathlineto{\pgfqpoint{1.202565in}{0.755151in}}%
\pgfpathlineto{\pgfqpoint{1.206894in}{0.752028in}}%
\pgfpathclose%
\pgfusepath{fill}%
\end{pgfscope}%
\begin{pgfscope}%
\pgfpathrectangle{\pgfqpoint{0.407592in}{0.222904in}}{\pgfqpoint{1.411163in}{1.411163in}}%
\pgfusepath{clip}%
\pgfsetbuttcap%
\pgfsetroundjoin%
\definecolor{currentfill}{rgb}{0.749020,0.772549,0.792157}%
\pgfsetfillcolor{currentfill}%
\pgfsetfillopacity{0.160000}%
\pgfsetlinewidth{0.000000pt}%
\definecolor{currentstroke}{rgb}{0.000000,0.000000,0.000000}%
\pgfsetstrokecolor{currentstroke}%
\pgfsetdash{}{0pt}%
\pgfpathmoveto{\pgfqpoint{1.585897in}{1.150554in}}%
\pgfpathlineto{\pgfqpoint{1.598303in}{1.164071in}}%
\pgfpathlineto{\pgfqpoint{1.610711in}{1.177589in}}%
\pgfpathlineto{\pgfqpoint{1.595206in}{1.189907in}}%
\pgfpathlineto{\pgfqpoint{1.577687in}{1.201739in}}%
\pgfpathlineto{\pgfqpoint{1.566157in}{1.187756in}}%
\pgfpathlineto{\pgfqpoint{1.554624in}{1.173768in}}%
\pgfpathlineto{\pgfqpoint{1.571225in}{1.162392in}}%
\pgfpathlineto{\pgfqpoint{1.585897in}{1.150554in}}%
\pgfpathclose%
\pgfusepath{fill}%
\end{pgfscope}%
\begin{pgfscope}%
\pgfpathrectangle{\pgfqpoint{0.407592in}{0.222904in}}{\pgfqpoint{1.411163in}{1.411163in}}%
\pgfusepath{clip}%
\pgfsetbuttcap%
\pgfsetroundjoin%
\definecolor{currentfill}{rgb}{0.749020,0.772549,0.792157}%
\pgfsetfillcolor{currentfill}%
\pgfsetfillopacity{0.160000}%
\pgfsetlinewidth{0.000000pt}%
\definecolor{currentstroke}{rgb}{0.000000,0.000000,0.000000}%
\pgfsetstrokecolor{currentstroke}%
\pgfsetdash{}{0pt}%
\pgfpathmoveto{\pgfqpoint{0.465273in}{1.394126in}}%
\pgfpathlineto{\pgfqpoint{0.454204in}{1.408005in}}%
\pgfpathlineto{\pgfqpoint{0.426275in}{1.392953in}}%
\pgfpathlineto{\pgfqpoint{0.400835in}{1.377253in}}%
\pgfpathlineto{\pgfqpoint{0.412963in}{1.363764in}}%
\pgfpathlineto{\pgfqpoint{0.437892in}{1.379263in}}%
\pgfusepath{fill}%
\end{pgfscope}%
\begin{pgfscope}%
\pgfpathrectangle{\pgfqpoint{0.407592in}{0.222904in}}{\pgfqpoint{1.411163in}{1.411163in}}%
\pgfusepath{clip}%
\pgfsetbuttcap%
\pgfsetroundjoin%
\definecolor{currentfill}{rgb}{0.749020,0.772549,0.792157}%
\pgfsetfillcolor{currentfill}%
\pgfsetfillopacity{0.160000}%
\pgfsetlinewidth{0.000000pt}%
\definecolor{currentstroke}{rgb}{0.000000,0.000000,0.000000}%
\pgfsetstrokecolor{currentstroke}%
\pgfsetdash{}{0pt}%
\pgfpathmoveto{\pgfqpoint{1.101249in}{0.643784in}}%
\pgfpathlineto{\pgfqpoint{1.106744in}{0.660368in}}%
\pgfpathlineto{\pgfqpoint{1.112231in}{0.676922in}}%
\pgfpathlineto{\pgfqpoint{1.109033in}{0.677446in}}%
\pgfpathlineto{\pgfqpoint{1.105754in}{0.677875in}}%
\pgfpathlineto{\pgfqpoint{1.101883in}{0.661086in}}%
\pgfpathlineto{\pgfqpoint{1.098006in}{0.644266in}}%
\pgfpathlineto{\pgfqpoint{1.099648in}{0.644049in}}%
\pgfpathlineto{\pgfqpoint{1.101249in}{0.643784in}}%
\pgfpathclose%
\pgfusepath{fill}%
\end{pgfscope}%
\begin{pgfscope}%
\pgfpathrectangle{\pgfqpoint{0.407592in}{0.222904in}}{\pgfqpoint{1.411163in}{1.411163in}}%
\pgfusepath{clip}%
\pgfsetbuttcap%
\pgfsetroundjoin%
\definecolor{currentfill}{rgb}{0.749020,0.772549,0.792157}%
\pgfsetfillcolor{currentfill}%
\pgfsetfillopacity{0.160000}%
\pgfsetlinewidth{0.000000pt}%
\definecolor{currentstroke}{rgb}{0.000000,0.000000,0.000000}%
\pgfsetstrokecolor{currentstroke}%
\pgfsetdash{}{0pt}%
\pgfpathmoveto{\pgfqpoint{1.084268in}{0.644227in}}%
\pgfpathlineto{\pgfqpoint{1.081296in}{0.661027in}}%
\pgfpathlineto{\pgfqpoint{1.078329in}{0.677797in}}%
\pgfpathlineto{\pgfqpoint{1.074973in}{0.677348in}}%
\pgfpathlineto{\pgfqpoint{1.071682in}{0.676806in}}%
\pgfpathlineto{\pgfqpoint{1.076307in}{0.660280in}}%
\pgfpathlineto{\pgfqpoint{1.080940in}{0.643726in}}%
\pgfpathlineto{\pgfqpoint{1.082588in}{0.644000in}}%
\pgfpathlineto{\pgfqpoint{1.084268in}{0.644227in}}%
\pgfpathclose%
\pgfusepath{fill}%
\end{pgfscope}%
\begin{pgfscope}%
\pgfpathrectangle{\pgfqpoint{0.407592in}{0.222904in}}{\pgfqpoint{1.411163in}{1.411163in}}%
\pgfusepath{clip}%
\pgfsetbuttcap%
\pgfsetroundjoin%
\definecolor{currentfill}{rgb}{0.749020,0.772549,0.792157}%
\pgfsetfillcolor{currentfill}%
\pgfsetfillopacity{0.160000}%
\pgfsetlinewidth{0.000000pt}%
\definecolor{currentstroke}{rgb}{0.000000,0.000000,0.000000}%
\pgfsetstrokecolor{currentstroke}%
\pgfsetdash{}{0pt}%
\pgfpathmoveto{\pgfqpoint{1.561085in}{1.123522in}}%
\pgfpathlineto{\pgfqpoint{1.573491in}{1.137038in}}%
\pgfpathlineto{\pgfqpoint{1.585897in}{1.150554in}}%
\pgfpathlineto{\pgfqpoint{1.571225in}{1.162392in}}%
\pgfpathlineto{\pgfqpoint{1.554624in}{1.173768in}}%
\pgfpathlineto{\pgfqpoint{1.543088in}{1.159776in}}%
\pgfpathlineto{\pgfqpoint{1.531547in}{1.145779in}}%
\pgfpathlineto{\pgfqpoint{1.547237in}{1.134871in}}%
\pgfpathlineto{\pgfqpoint{1.561085in}{1.123522in}}%
\pgfpathclose%
\pgfusepath{fill}%
\end{pgfscope}%
\begin{pgfscope}%
\pgfpathrectangle{\pgfqpoint{0.407592in}{0.222904in}}{\pgfqpoint{1.411163in}{1.411163in}}%
\pgfusepath{clip}%
\pgfsetbuttcap%
\pgfsetroundjoin%
\definecolor{currentfill}{rgb}{0.749020,0.772549,0.792157}%
\pgfsetfillcolor{currentfill}%
\pgfsetfillopacity{0.160000}%
\pgfsetlinewidth{0.000000pt}%
\definecolor{currentstroke}{rgb}{0.000000,0.000000,0.000000}%
\pgfsetstrokecolor{currentstroke}%
\pgfsetdash{}{0pt}%
\pgfpathmoveto{\pgfqpoint{0.487421in}{1.366357in}}%
\pgfpathlineto{\pgfqpoint{0.476345in}{1.380243in}}%
\pgfpathlineto{\pgfqpoint{0.465273in}{1.394126in}}%
\pgfpathlineto{\pgfqpoint{0.437892in}{1.379263in}}%
\pgfpathlineto{\pgfqpoint{0.412963in}{1.363764in}}%
\pgfpathlineto{\pgfqpoint{0.425090in}{1.350276in}}%
\pgfpathlineto{\pgfqpoint{0.437216in}{1.336790in}}%
\pgfpathlineto{\pgfqpoint{0.461128in}{1.351881in}}%
\pgfpathlineto{\pgfqpoint{0.487421in}{1.366357in}}%
\pgfpathclose%
\pgfusepath{fill}%
\end{pgfscope}%
\begin{pgfscope}%
\pgfpathrectangle{\pgfqpoint{0.407592in}{0.222904in}}{\pgfqpoint{1.411163in}{1.411163in}}%
\pgfusepath{clip}%
\pgfsetbuttcap%
\pgfsetroundjoin%
\definecolor{currentfill}{rgb}{0.749020,0.772549,0.792157}%
\pgfsetfillcolor{currentfill}%
\pgfsetfillopacity{0.160000}%
\pgfsetlinewidth{0.000000pt}%
\definecolor{currentstroke}{rgb}{0.000000,0.000000,0.000000}%
\pgfsetstrokecolor{currentstroke}%
\pgfsetdash{}{0pt}%
\pgfpathmoveto{\pgfqpoint{0.988126in}{0.757547in}}%
\pgfpathlineto{\pgfqpoint{0.977952in}{0.772196in}}%
\pgfpathlineto{\pgfqpoint{0.967785in}{0.786836in}}%
\pgfpathlineto{\pgfqpoint{0.961352in}{0.783254in}}%
\pgfpathlineto{\pgfqpoint{0.955463in}{0.779508in}}%
\pgfpathlineto{\pgfqpoint{0.966676in}{0.765449in}}%
\pgfpathlineto{\pgfqpoint{0.977892in}{0.751386in}}%
\pgfpathlineto{\pgfqpoint{0.982781in}{0.754535in}}%
\pgfpathlineto{\pgfqpoint{0.988126in}{0.757547in}}%
\pgfpathclose%
\pgfusepath{fill}%
\end{pgfscope}%
\begin{pgfscope}%
\pgfpathrectangle{\pgfqpoint{0.407592in}{0.222904in}}{\pgfqpoint{1.411163in}{1.411163in}}%
\pgfusepath{clip}%
\pgfsetbuttcap%
\pgfsetroundjoin%
\definecolor{currentfill}{rgb}{0.749020,0.772549,0.792157}%
\pgfsetfillcolor{currentfill}%
\pgfsetfillopacity{0.160000}%
\pgfsetlinewidth{0.000000pt}%
\definecolor{currentstroke}{rgb}{0.000000,0.000000,0.000000}%
\pgfsetstrokecolor{currentstroke}%
\pgfsetdash{}{0pt}%
\pgfpathmoveto{\pgfqpoint{0.509581in}{1.338571in}}%
\pgfpathlineto{\pgfqpoint{0.498499in}{1.352466in}}%
\pgfpathlineto{\pgfqpoint{0.487421in}{1.366357in}}%
\pgfpathlineto{\pgfqpoint{0.461128in}{1.351881in}}%
\pgfpathlineto{\pgfqpoint{0.437216in}{1.336790in}}%
\pgfpathlineto{\pgfqpoint{0.449341in}{1.323306in}}%
\pgfpathlineto{\pgfqpoint{0.461465in}{1.309822in}}%
\pgfpathlineto{\pgfqpoint{0.484369in}{1.324493in}}%
\pgfpathlineto{\pgfqpoint{0.509581in}{1.338571in}}%
\pgfpathclose%
\pgfusepath{fill}%
\end{pgfscope}%
\begin{pgfscope}%
\pgfpathrectangle{\pgfqpoint{0.407592in}{0.222904in}}{\pgfqpoint{1.411163in}{1.411163in}}%
\pgfusepath{clip}%
\pgfsetbuttcap%
\pgfsetroundjoin%
\definecolor{currentfill}{rgb}{0.749020,0.772549,0.792157}%
\pgfsetfillcolor{currentfill}%
\pgfsetfillopacity{0.160000}%
\pgfsetlinewidth{0.000000pt}%
\definecolor{currentstroke}{rgb}{0.000000,0.000000,0.000000}%
\pgfsetstrokecolor{currentstroke}%
\pgfsetdash{}{0pt}%
\pgfpathmoveto{\pgfqpoint{1.536277in}{1.096494in}}%
\pgfpathlineto{\pgfqpoint{1.548681in}{1.110008in}}%
\pgfpathlineto{\pgfqpoint{1.561085in}{1.123522in}}%
\pgfpathlineto{\pgfqpoint{1.547237in}{1.134871in}}%
\pgfpathlineto{\pgfqpoint{1.531547in}{1.145779in}}%
\pgfpathlineto{\pgfqpoint{1.520003in}{1.131778in}}%
\pgfpathlineto{\pgfqpoint{1.508456in}{1.117773in}}%
\pgfpathlineto{\pgfqpoint{1.523243in}{1.107341in}}%
\pgfpathlineto{\pgfqpoint{1.536277in}{1.096494in}}%
\pgfpathclose%
\pgfusepath{fill}%
\end{pgfscope}%
\begin{pgfscope}%
\pgfpathrectangle{\pgfqpoint{0.407592in}{0.222904in}}{\pgfqpoint{1.411163in}{1.411163in}}%
\pgfusepath{clip}%
\pgfsetbuttcap%
\pgfsetroundjoin%
\definecolor{currentfill}{rgb}{0.749020,0.772549,0.792157}%
\pgfsetfillcolor{currentfill}%
\pgfsetfillopacity{0.160000}%
\pgfsetlinewidth{0.000000pt}%
\definecolor{currentstroke}{rgb}{0.000000,0.000000,0.000000}%
\pgfsetstrokecolor{currentstroke}%
\pgfsetdash{}{0pt}%
\pgfpathmoveto{\pgfqpoint{1.104300in}{0.643117in}}%
\pgfpathlineto{\pgfqpoint{1.111319in}{0.659372in}}%
\pgfpathlineto{\pgfqpoint{1.118327in}{0.675602in}}%
\pgfpathlineto{\pgfqpoint{1.115333in}{0.676307in}}%
\pgfpathlineto{\pgfqpoint{1.112231in}{0.676922in}}%
\pgfpathlineto{\pgfqpoint{1.106744in}{0.660368in}}%
\pgfpathlineto{\pgfqpoint{1.101249in}{0.643784in}}%
\pgfpathlineto{\pgfqpoint{1.102802in}{0.643473in}}%
\pgfpathlineto{\pgfqpoint{1.104300in}{0.643117in}}%
\pgfpathclose%
\pgfusepath{fill}%
\end{pgfscope}%
\begin{pgfscope}%
\pgfpathrectangle{\pgfqpoint{0.407592in}{0.222904in}}{\pgfqpoint{1.411163in}{1.411163in}}%
\pgfusepath{clip}%
\pgfsetbuttcap%
\pgfsetroundjoin%
\definecolor{currentfill}{rgb}{0.749020,0.772549,0.792157}%
\pgfsetfillcolor{currentfill}%
\pgfsetfillopacity{0.160000}%
\pgfsetlinewidth{0.000000pt}%
\definecolor{currentstroke}{rgb}{0.000000,0.000000,0.000000}%
\pgfsetstrokecolor{currentstroke}%
\pgfsetdash{}{0pt}%
\pgfpathmoveto{\pgfqpoint{0.531754in}{1.310770in}}%
\pgfpathlineto{\pgfqpoint{0.520666in}{1.324673in}}%
\pgfpathlineto{\pgfqpoint{0.509581in}{1.338571in}}%
\pgfpathlineto{\pgfqpoint{0.484369in}{1.324493in}}%
\pgfpathlineto{\pgfqpoint{0.461465in}{1.309822in}}%
\pgfpathlineto{\pgfqpoint{0.473587in}{1.296340in}}%
\pgfpathlineto{\pgfqpoint{0.485708in}{1.282860in}}%
\pgfpathlineto{\pgfqpoint{0.507615in}{1.297099in}}%
\pgfpathlineto{\pgfqpoint{0.531754in}{1.310770in}}%
\pgfpathclose%
\pgfusepath{fill}%
\end{pgfscope}%
\begin{pgfscope}%
\pgfpathrectangle{\pgfqpoint{0.407592in}{0.222904in}}{\pgfqpoint{1.411163in}{1.411163in}}%
\pgfusepath{clip}%
\pgfsetbuttcap%
\pgfsetroundjoin%
\definecolor{currentfill}{rgb}{0.749020,0.772549,0.792157}%
\pgfsetfillcolor{currentfill}%
\pgfsetfillopacity{0.160000}%
\pgfsetlinewidth{0.000000pt}%
\definecolor{currentstroke}{rgb}{0.000000,0.000000,0.000000}%
\pgfsetstrokecolor{currentstroke}%
\pgfsetdash{}{0pt}%
\pgfpathmoveto{\pgfqpoint{1.080940in}{0.643726in}}%
\pgfpathlineto{\pgfqpoint{1.076307in}{0.660280in}}%
\pgfpathlineto{\pgfqpoint{1.071682in}{0.676806in}}%
\pgfpathlineto{\pgfqpoint{1.068470in}{0.676173in}}%
\pgfpathlineto{\pgfqpoint{1.065351in}{0.675451in}}%
\pgfpathlineto{\pgfqpoint{1.071557in}{0.659258in}}%
\pgfpathlineto{\pgfqpoint{1.077772in}{0.643040in}}%
\pgfpathlineto{\pgfqpoint{1.079332in}{0.643405in}}%
\pgfpathlineto{\pgfqpoint{1.080940in}{0.643726in}}%
\pgfpathclose%
\pgfusepath{fill}%
\end{pgfscope}%
\begin{pgfscope}%
\pgfpathrectangle{\pgfqpoint{0.407592in}{0.222904in}}{\pgfqpoint{1.411163in}{1.411163in}}%
\pgfusepath{clip}%
\pgfsetbuttcap%
\pgfsetroundjoin%
\definecolor{currentfill}{rgb}{0.749020,0.772549,0.792157}%
\pgfsetfillcolor{currentfill}%
\pgfsetfillopacity{0.160000}%
\pgfsetlinewidth{0.000000pt}%
\definecolor{currentstroke}{rgb}{0.000000,0.000000,0.000000}%
\pgfsetstrokecolor{currentstroke}%
\pgfsetdash{}{0pt}%
\pgfpathmoveto{\pgfqpoint{1.511471in}{1.069469in}}%
\pgfpathlineto{\pgfqpoint{1.523874in}{1.082981in}}%
\pgfpathlineto{\pgfqpoint{1.536277in}{1.096494in}}%
\pgfpathlineto{\pgfqpoint{1.523243in}{1.107341in}}%
\pgfpathlineto{\pgfqpoint{1.508456in}{1.117773in}}%
\pgfpathlineto{\pgfqpoint{1.496904in}{1.103763in}}%
\pgfpathlineto{\pgfqpoint{1.485349in}{1.089748in}}%
\pgfpathlineto{\pgfqpoint{1.499242in}{1.079805in}}%
\pgfpathlineto{\pgfqpoint{1.511471in}{1.069469in}}%
\pgfpathclose%
\pgfusepath{fill}%
\end{pgfscope}%
\begin{pgfscope}%
\pgfpathrectangle{\pgfqpoint{0.407592in}{0.222904in}}{\pgfqpoint{1.411163in}{1.411163in}}%
\pgfusepath{clip}%
\pgfsetbuttcap%
\pgfsetroundjoin%
\definecolor{currentfill}{rgb}{0.749020,0.772549,0.792157}%
\pgfsetfillcolor{currentfill}%
\pgfsetfillopacity{0.160000}%
\pgfsetlinewidth{0.000000pt}%
\definecolor{currentstroke}{rgb}{0.000000,0.000000,0.000000}%
\pgfsetstrokecolor{currentstroke}%
\pgfsetdash{}{0pt}%
\pgfpathmoveto{\pgfqpoint{0.553940in}{1.282953in}}%
\pgfpathlineto{\pgfqpoint{0.542845in}{1.296863in}}%
\pgfpathlineto{\pgfqpoint{0.531754in}{1.310770in}}%
\pgfpathlineto{\pgfqpoint{0.507615in}{1.297099in}}%
\pgfpathlineto{\pgfqpoint{0.485708in}{1.282860in}}%
\pgfpathlineto{\pgfqpoint{0.497828in}{1.269380in}}%
\pgfpathlineto{\pgfqpoint{0.509947in}{1.255902in}}%
\pgfpathlineto{\pgfqpoint{0.530865in}{1.269701in}}%
\pgfpathlineto{\pgfqpoint{0.553940in}{1.282953in}}%
\pgfpathclose%
\pgfusepath{fill}%
\end{pgfscope}%
\begin{pgfscope}%
\pgfpathrectangle{\pgfqpoint{0.407592in}{0.222904in}}{\pgfqpoint{1.411163in}{1.411163in}}%
\pgfusepath{clip}%
\pgfsetbuttcap%
\pgfsetroundjoin%
\definecolor{currentfill}{rgb}{0.749020,0.772549,0.792157}%
\pgfsetfillcolor{currentfill}%
\pgfsetfillopacity{0.160000}%
\pgfsetlinewidth{0.000000pt}%
\definecolor{currentstroke}{rgb}{0.000000,0.000000,0.000000}%
\pgfsetstrokecolor{currentstroke}%
\pgfsetdash{}{0pt}%
\pgfpathmoveto{\pgfqpoint{1.183592in}{0.723766in}}%
\pgfpathlineto{\pgfqpoint{1.195245in}{0.737899in}}%
\pgfpathlineto{\pgfqpoint{1.206894in}{0.752028in}}%
\pgfpathlineto{\pgfqpoint{1.202565in}{0.755151in}}%
\pgfpathlineto{\pgfqpoint{1.197756in}{0.758134in}}%
\pgfpathlineto{\pgfqpoint{1.187036in}{0.743419in}}%
\pgfpathlineto{\pgfqpoint{1.176309in}{0.728694in}}%
\pgfpathlineto{\pgfqpoint{1.180144in}{0.726286in}}%
\pgfpathlineto{\pgfqpoint{1.183592in}{0.723766in}}%
\pgfpathclose%
\pgfusepath{fill}%
\end{pgfscope}%
\begin{pgfscope}%
\pgfpathrectangle{\pgfqpoint{0.407592in}{0.222904in}}{\pgfqpoint{1.411163in}{1.411163in}}%
\pgfusepath{clip}%
\pgfsetbuttcap%
\pgfsetroundjoin%
\definecolor{currentfill}{rgb}{0.749020,0.772549,0.792157}%
\pgfsetfillcolor{currentfill}%
\pgfsetfillopacity{0.160000}%
\pgfsetlinewidth{0.000000pt}%
\definecolor{currentstroke}{rgb}{0.000000,0.000000,0.000000}%
\pgfsetstrokecolor{currentstroke}%
\pgfsetdash{}{0pt}%
\pgfpathmoveto{\pgfqpoint{0.576138in}{1.255120in}}%
\pgfpathlineto{\pgfqpoint{0.565037in}{1.269038in}}%
\pgfpathlineto{\pgfqpoint{0.553940in}{1.282953in}}%
\pgfpathlineto{\pgfqpoint{0.530865in}{1.269701in}}%
\pgfpathlineto{\pgfqpoint{0.509947in}{1.255902in}}%
\pgfpathlineto{\pgfqpoint{0.522064in}{1.242426in}}%
\pgfpathlineto{\pgfqpoint{0.534180in}{1.228950in}}%
\pgfpathlineto{\pgfqpoint{0.554120in}{1.242297in}}%
\pgfpathlineto{\pgfqpoint{0.576138in}{1.255120in}}%
\pgfpathclose%
\pgfusepath{fill}%
\end{pgfscope}%
\begin{pgfscope}%
\pgfpathrectangle{\pgfqpoint{0.407592in}{0.222904in}}{\pgfqpoint{1.411163in}{1.411163in}}%
\pgfusepath{clip}%
\pgfsetbuttcap%
\pgfsetroundjoin%
\definecolor{currentfill}{rgb}{0.749020,0.772549,0.792157}%
\pgfsetfillcolor{currentfill}%
\pgfsetfillopacity{0.160000}%
\pgfsetlinewidth{0.000000pt}%
\definecolor{currentstroke}{rgb}{0.000000,0.000000,0.000000}%
\pgfsetstrokecolor{currentstroke}%
\pgfsetdash{}{0pt}%
\pgfpathmoveto{\pgfqpoint{1.486669in}{1.042447in}}%
\pgfpathlineto{\pgfqpoint{1.499070in}{1.055958in}}%
\pgfpathlineto{\pgfqpoint{1.511471in}{1.069469in}}%
\pgfpathlineto{\pgfqpoint{1.499242in}{1.079805in}}%
\pgfpathlineto{\pgfqpoint{1.485349in}{1.089748in}}%
\pgfpathlineto{\pgfqpoint{1.473790in}{1.075729in}}%
\pgfpathlineto{\pgfqpoint{1.462227in}{1.061705in}}%
\pgfpathlineto{\pgfqpoint{1.475235in}{1.052261in}}%
\pgfpathlineto{\pgfqpoint{1.486669in}{1.042447in}}%
\pgfpathclose%
\pgfusepath{fill}%
\end{pgfscope}%
\begin{pgfscope}%
\pgfpathrectangle{\pgfqpoint{0.407592in}{0.222904in}}{\pgfqpoint{1.411163in}{1.411163in}}%
\pgfusepath{clip}%
\pgfsetbuttcap%
\pgfsetroundjoin%
\definecolor{currentfill}{rgb}{0.749020,0.772549,0.792157}%
\pgfsetfillcolor{currentfill}%
\pgfsetfillopacity{0.160000}%
\pgfsetlinewidth{0.000000pt}%
\definecolor{currentstroke}{rgb}{0.000000,0.000000,0.000000}%
\pgfsetstrokecolor{currentstroke}%
\pgfsetdash{}{0pt}%
\pgfpathmoveto{\pgfqpoint{0.598349in}{1.227271in}}%
\pgfpathlineto{\pgfqpoint{0.587242in}{1.241197in}}%
\pgfpathlineto{\pgfqpoint{0.576138in}{1.255120in}}%
\pgfpathlineto{\pgfqpoint{0.554120in}{1.242297in}}%
\pgfpathlineto{\pgfqpoint{0.534180in}{1.228950in}}%
\pgfpathlineto{\pgfqpoint{0.546296in}{1.215477in}}%
\pgfpathlineto{\pgfqpoint{0.558409in}{1.202004in}}%
\pgfpathlineto{\pgfqpoint{0.577379in}{1.214888in}}%
\pgfpathlineto{\pgfqpoint{0.598349in}{1.227271in}}%
\pgfpathclose%
\pgfusepath{fill}%
\end{pgfscope}%
\begin{pgfscope}%
\pgfpathrectangle{\pgfqpoint{0.407592in}{0.222904in}}{\pgfqpoint{1.411163in}{1.411163in}}%
\pgfusepath{clip}%
\pgfsetbuttcap%
\pgfsetroundjoin%
\definecolor{currentfill}{rgb}{0.749020,0.772549,0.792157}%
\pgfsetfillcolor{currentfill}%
\pgfsetfillopacity{0.160000}%
\pgfsetlinewidth{0.000000pt}%
\definecolor{currentstroke}{rgb}{0.000000,0.000000,0.000000}%
\pgfsetstrokecolor{currentstroke}%
\pgfsetdash{}{0pt}%
\pgfpathmoveto{\pgfqpoint{1.133328in}{0.669694in}}%
\pgfpathlineto{\pgfqpoint{1.144084in}{0.684459in}}%
\pgfpathlineto{\pgfqpoint{1.154833in}{0.699214in}}%
\pgfpathlineto{\pgfqpoint{1.151689in}{0.700945in}}%
\pgfpathlineto{\pgfqpoint{1.148280in}{0.702577in}}%
\pgfpathlineto{\pgfqpoint{1.138629in}{0.687275in}}%
\pgfpathlineto{\pgfqpoint{1.128969in}{0.671957in}}%
\pgfpathlineto{\pgfqpoint{1.131237in}{0.670859in}}%
\pgfpathlineto{\pgfqpoint{1.133328in}{0.669694in}}%
\pgfpathclose%
\pgfusepath{fill}%
\end{pgfscope}%
\begin{pgfscope}%
\pgfpathrectangle{\pgfqpoint{0.407592in}{0.222904in}}{\pgfqpoint{1.411163in}{1.411163in}}%
\pgfusepath{clip}%
\pgfsetbuttcap%
\pgfsetroundjoin%
\definecolor{currentfill}{rgb}{0.749020,0.772549,0.792157}%
\pgfsetfillcolor{currentfill}%
\pgfsetfillopacity{0.160000}%
\pgfsetlinewidth{0.000000pt}%
\definecolor{currentstroke}{rgb}{0.000000,0.000000,0.000000}%
\pgfsetstrokecolor{currentstroke}%
\pgfsetdash{}{0pt}%
\pgfpathmoveto{\pgfqpoint{1.107106in}{0.642276in}}%
\pgfpathlineto{\pgfqpoint{1.115528in}{0.658117in}}%
\pgfpathlineto{\pgfqpoint{1.123939in}{0.673938in}}%
\pgfpathlineto{\pgfqpoint{1.121200in}{0.674812in}}%
\pgfpathlineto{\pgfqpoint{1.118327in}{0.675602in}}%
\pgfpathlineto{\pgfqpoint{1.111319in}{0.659372in}}%
\pgfpathlineto{\pgfqpoint{1.104300in}{0.643117in}}%
\pgfpathlineto{\pgfqpoint{1.105737in}{0.642718in}}%
\pgfpathlineto{\pgfqpoint{1.107106in}{0.642276in}}%
\pgfpathclose%
\pgfusepath{fill}%
\end{pgfscope}%
\begin{pgfscope}%
\pgfpathrectangle{\pgfqpoint{0.407592in}{0.222904in}}{\pgfqpoint{1.411163in}{1.411163in}}%
\pgfusepath{clip}%
\pgfsetbuttcap%
\pgfsetroundjoin%
\definecolor{currentfill}{rgb}{0.749020,0.772549,0.792157}%
\pgfsetfillcolor{currentfill}%
\pgfsetfillopacity{0.160000}%
\pgfsetlinewidth{0.000000pt}%
\definecolor{currentstroke}{rgb}{0.000000,0.000000,0.000000}%
\pgfsetstrokecolor{currentstroke}%
\pgfsetdash{}{0pt}%
\pgfpathmoveto{\pgfqpoint{1.008494in}{0.728220in}}%
\pgfpathlineto{\pgfqpoint{0.998307in}{0.742888in}}%
\pgfpathlineto{\pgfqpoint{0.988126in}{0.757547in}}%
\pgfpathlineto{\pgfqpoint{0.982781in}{0.754535in}}%
\pgfpathlineto{\pgfqpoint{0.977892in}{0.751386in}}%
\pgfpathlineto{\pgfqpoint{0.989111in}{0.737319in}}%
\pgfpathlineto{\pgfqpoint{1.000333in}{0.723248in}}%
\pgfpathlineto{\pgfqpoint{1.004230in}{0.725789in}}%
\pgfpathlineto{\pgfqpoint{1.008494in}{0.728220in}}%
\pgfpathclose%
\pgfusepath{fill}%
\end{pgfscope}%
\begin{pgfscope}%
\pgfpathrectangle{\pgfqpoint{0.407592in}{0.222904in}}{\pgfqpoint{1.411163in}{1.411163in}}%
\pgfusepath{clip}%
\pgfsetbuttcap%
\pgfsetroundjoin%
\definecolor{currentfill}{rgb}{0.749020,0.772549,0.792157}%
\pgfsetfillcolor{currentfill}%
\pgfsetfillopacity{0.160000}%
\pgfsetlinewidth{0.000000pt}%
\definecolor{currentstroke}{rgb}{0.000000,0.000000,0.000000}%
\pgfsetstrokecolor{currentstroke}%
\pgfsetdash{}{0pt}%
\pgfpathmoveto{\pgfqpoint{1.077772in}{0.643040in}}%
\pgfpathlineto{\pgfqpoint{1.071557in}{0.659258in}}%
\pgfpathlineto{\pgfqpoint{1.065351in}{0.675451in}}%
\pgfpathlineto{\pgfqpoint{1.062338in}{0.674643in}}%
\pgfpathlineto{\pgfqpoint{1.059445in}{0.673753in}}%
\pgfpathlineto{\pgfqpoint{1.067127in}{0.657978in}}%
\pgfpathlineto{\pgfqpoint{1.074818in}{0.642183in}}%
\pgfpathlineto{\pgfqpoint{1.076265in}{0.642632in}}%
\pgfpathlineto{\pgfqpoint{1.077772in}{0.643040in}}%
\pgfpathclose%
\pgfusepath{fill}%
\end{pgfscope}%
\begin{pgfscope}%
\pgfpathrectangle{\pgfqpoint{0.407592in}{0.222904in}}{\pgfqpoint{1.411163in}{1.411163in}}%
\pgfusepath{clip}%
\pgfsetbuttcap%
\pgfsetroundjoin%
\definecolor{currentfill}{rgb}{0.749020,0.772549,0.792157}%
\pgfsetfillcolor{currentfill}%
\pgfsetfillopacity{0.160000}%
\pgfsetlinewidth{0.000000pt}%
\definecolor{currentstroke}{rgb}{0.000000,0.000000,0.000000}%
\pgfsetstrokecolor{currentstroke}%
\pgfsetdash{}{0pt}%
\pgfpathmoveto{\pgfqpoint{1.461869in}{1.015428in}}%
\pgfpathlineto{\pgfqpoint{1.474269in}{1.028937in}}%
\pgfpathlineto{\pgfqpoint{1.486669in}{1.042447in}}%
\pgfpathlineto{\pgfqpoint{1.475235in}{1.052261in}}%
\pgfpathlineto{\pgfqpoint{1.462227in}{1.061705in}}%
\pgfpathlineto{\pgfqpoint{1.450660in}{1.047676in}}%
\pgfpathlineto{\pgfqpoint{1.439090in}{1.033644in}}%
\pgfpathlineto{\pgfqpoint{1.451221in}{1.024709in}}%
\pgfpathlineto{\pgfqpoint{1.461869in}{1.015428in}}%
\pgfpathclose%
\pgfusepath{fill}%
\end{pgfscope}%
\begin{pgfscope}%
\pgfpathrectangle{\pgfqpoint{0.407592in}{0.222904in}}{\pgfqpoint{1.411163in}{1.411163in}}%
\pgfusepath{clip}%
\pgfsetbuttcap%
\pgfsetroundjoin%
\definecolor{currentfill}{rgb}{0.749020,0.772549,0.792157}%
\pgfsetfillcolor{currentfill}%
\pgfsetfillopacity{0.160000}%
\pgfsetlinewidth{0.000000pt}%
\definecolor{currentstroke}{rgb}{0.000000,0.000000,0.000000}%
\pgfsetstrokecolor{currentstroke}%
\pgfsetdash{}{0pt}%
\pgfpathmoveto{\pgfqpoint{0.620573in}{1.199406in}}%
\pgfpathlineto{\pgfqpoint{0.609459in}{1.213340in}}%
\pgfpathlineto{\pgfqpoint{0.598349in}{1.227271in}}%
\pgfpathlineto{\pgfqpoint{0.577379in}{1.214888in}}%
\pgfpathlineto{\pgfqpoint{0.558409in}{1.202004in}}%
\pgfpathlineto{\pgfqpoint{0.570522in}{1.188533in}}%
\pgfpathlineto{\pgfqpoint{0.582634in}{1.175063in}}%
\pgfpathlineto{\pgfqpoint{0.600643in}{1.187473in}}%
\pgfpathlineto{\pgfqpoint{0.620573in}{1.199406in}}%
\pgfpathclose%
\pgfusepath{fill}%
\end{pgfscope}%
\begin{pgfscope}%
\pgfpathrectangle{\pgfqpoint{0.407592in}{0.222904in}}{\pgfqpoint{1.411163in}{1.411163in}}%
\pgfusepath{clip}%
\pgfsetbuttcap%
\pgfsetroundjoin%
\definecolor{currentfill}{rgb}{0.749020,0.772549,0.792157}%
\pgfsetfillcolor{currentfill}%
\pgfsetfillopacity{0.160000}%
\pgfsetlinewidth{0.000000pt}%
\definecolor{currentstroke}{rgb}{0.000000,0.000000,0.000000}%
\pgfsetstrokecolor{currentstroke}%
\pgfsetdash{}{0pt}%
\pgfpathmoveto{\pgfqpoint{1.054067in}{0.671742in}}%
\pgfpathlineto{\pgfqpoint{1.045048in}{0.687008in}}%
\pgfpathlineto{\pgfqpoint{1.036038in}{0.702258in}}%
\pgfpathlineto{\pgfqpoint{1.032339in}{0.700605in}}%
\pgfpathlineto{\pgfqpoint{1.028888in}{0.698855in}}%
\pgfpathlineto{\pgfqpoint{1.039095in}{0.684159in}}%
\pgfpathlineto{\pgfqpoint{1.049309in}{0.669453in}}%
\pgfpathlineto{\pgfqpoint{1.051605in}{0.670630in}}%
\pgfpathlineto{\pgfqpoint{1.054067in}{0.671742in}}%
\pgfpathclose%
\pgfusepath{fill}%
\end{pgfscope}%
\begin{pgfscope}%
\pgfpathrectangle{\pgfqpoint{0.407592in}{0.222904in}}{\pgfqpoint{1.411163in}{1.411163in}}%
\pgfusepath{clip}%
\pgfsetbuttcap%
\pgfsetroundjoin%
\definecolor{currentfill}{rgb}{0.749020,0.772549,0.792157}%
\pgfsetfillcolor{currentfill}%
\pgfsetfillopacity{0.160000}%
\pgfsetlinewidth{0.000000pt}%
\definecolor{currentstroke}{rgb}{0.000000,0.000000,0.000000}%
\pgfsetstrokecolor{currentstroke}%
\pgfsetdash{}{0pt}%
\pgfpathmoveto{\pgfqpoint{0.642809in}{1.171525in}}%
\pgfpathlineto{\pgfqpoint{0.631689in}{1.185467in}}%
\pgfpathlineto{\pgfqpoint{0.620573in}{1.199406in}}%
\pgfpathlineto{\pgfqpoint{0.600643in}{1.187473in}}%
\pgfpathlineto{\pgfqpoint{0.582634in}{1.175063in}}%
\pgfpathlineto{\pgfqpoint{0.594744in}{1.161594in}}%
\pgfpathlineto{\pgfqpoint{0.606853in}{1.148127in}}%
\pgfpathlineto{\pgfqpoint{0.623911in}{1.160053in}}%
\pgfpathlineto{\pgfqpoint{0.642809in}{1.171525in}}%
\pgfpathclose%
\pgfusepath{fill}%
\end{pgfscope}%
\begin{pgfscope}%
\pgfpathrectangle{\pgfqpoint{0.407592in}{0.222904in}}{\pgfqpoint{1.411163in}{1.411163in}}%
\pgfusepath{clip}%
\pgfsetbuttcap%
\pgfsetroundjoin%
\definecolor{currentfill}{rgb}{0.749020,0.772549,0.792157}%
\pgfsetfillcolor{currentfill}%
\pgfsetfillopacity{0.160000}%
\pgfsetlinewidth{0.000000pt}%
\definecolor{currentstroke}{rgb}{0.000000,0.000000,0.000000}%
\pgfsetstrokecolor{currentstroke}%
\pgfsetdash{}{0pt}%
\pgfpathmoveto{\pgfqpoint{1.437073in}{0.988413in}}%
\pgfpathlineto{\pgfqpoint{1.449471in}{1.001920in}}%
\pgfpathlineto{\pgfqpoint{1.461869in}{1.015428in}}%
\pgfpathlineto{\pgfqpoint{1.451221in}{1.024709in}}%
\pgfpathlineto{\pgfqpoint{1.439090in}{1.033644in}}%
\pgfpathlineto{\pgfqpoint{1.427516in}{1.019606in}}%
\pgfpathlineto{\pgfqpoint{1.415938in}{1.005564in}}%
\pgfpathlineto{\pgfqpoint{1.427201in}{0.997150in}}%
\pgfpathlineto{\pgfqpoint{1.437073in}{0.988413in}}%
\pgfpathclose%
\pgfusepath{fill}%
\end{pgfscope}%
\begin{pgfscope}%
\pgfpathrectangle{\pgfqpoint{0.407592in}{0.222904in}}{\pgfqpoint{1.411163in}{1.411163in}}%
\pgfusepath{clip}%
\pgfsetbuttcap%
\pgfsetroundjoin%
\definecolor{currentfill}{rgb}{0.749020,0.772549,0.792157}%
\pgfsetfillcolor{currentfill}%
\pgfsetfillopacity{0.160000}%
\pgfsetlinewidth{0.000000pt}%
\definecolor{currentstroke}{rgb}{0.000000,0.000000,0.000000}%
\pgfsetstrokecolor{currentstroke}%
\pgfsetdash{}{0pt}%
\pgfpathmoveto{\pgfqpoint{0.665058in}{1.143628in}}%
\pgfpathlineto{\pgfqpoint{0.653932in}{1.157578in}}%
\pgfpathlineto{\pgfqpoint{0.642809in}{1.171525in}}%
\pgfpathlineto{\pgfqpoint{0.623911in}{1.160053in}}%
\pgfpathlineto{\pgfqpoint{0.606853in}{1.148127in}}%
\pgfpathlineto{\pgfqpoint{0.618961in}{1.134661in}}%
\pgfpathlineto{\pgfqpoint{0.631067in}{1.121196in}}%
\pgfpathlineto{\pgfqpoint{0.647184in}{1.132628in}}%
\pgfpathlineto{\pgfqpoint{0.665058in}{1.143628in}}%
\pgfpathclose%
\pgfusepath{fill}%
\end{pgfscope}%
\begin{pgfscope}%
\pgfpathrectangle{\pgfqpoint{0.407592in}{0.222904in}}{\pgfqpoint{1.411163in}{1.411163in}}%
\pgfusepath{clip}%
\pgfsetbuttcap%
\pgfsetroundjoin%
\definecolor{currentfill}{rgb}{0.749020,0.772549,0.792157}%
\pgfsetfillcolor{currentfill}%
\pgfsetfillopacity{0.160000}%
\pgfsetlinewidth{0.000000pt}%
\definecolor{currentstroke}{rgb}{0.000000,0.000000,0.000000}%
\pgfsetstrokecolor{currentstroke}%
\pgfsetdash{}{0pt}%
\pgfpathmoveto{\pgfqpoint{1.412279in}{0.961401in}}%
\pgfpathlineto{\pgfqpoint{1.424676in}{0.974906in}}%
\pgfpathlineto{\pgfqpoint{1.437073in}{0.988413in}}%
\pgfpathlineto{\pgfqpoint{1.427201in}{0.997150in}}%
\pgfpathlineto{\pgfqpoint{1.415938in}{1.005564in}}%
\pgfpathlineto{\pgfqpoint{1.404357in}{0.991518in}}%
\pgfpathlineto{\pgfqpoint{1.392771in}{0.977467in}}%
\pgfpathlineto{\pgfqpoint{1.403174in}{0.969583in}}%
\pgfpathlineto{\pgfqpoint{1.412279in}{0.961401in}}%
\pgfpathclose%
\pgfusepath{fill}%
\end{pgfscope}%
\begin{pgfscope}%
\pgfpathrectangle{\pgfqpoint{0.407592in}{0.222904in}}{\pgfqpoint{1.411163in}{1.411163in}}%
\pgfusepath{clip}%
\pgfsetbuttcap%
\pgfsetroundjoin%
\definecolor{currentfill}{rgb}{0.749020,0.772549,0.792157}%
\pgfsetfillcolor{currentfill}%
\pgfsetfillopacity{0.160000}%
\pgfsetlinewidth{0.000000pt}%
\definecolor{currentstroke}{rgb}{0.000000,0.000000,0.000000}%
\pgfsetstrokecolor{currentstroke}%
\pgfsetdash{}{0pt}%
\pgfpathmoveto{\pgfqpoint{1.160274in}{0.695485in}}%
\pgfpathlineto{\pgfqpoint{1.171935in}{0.709628in}}%
\pgfpathlineto{\pgfqpoint{1.183592in}{0.723766in}}%
\pgfpathlineto{\pgfqpoint{1.180144in}{0.726286in}}%
\pgfpathlineto{\pgfqpoint{1.176309in}{0.728694in}}%
\pgfpathlineto{\pgfqpoint{1.165575in}{0.713959in}}%
\pgfpathlineto{\pgfqpoint{1.154833in}{0.699214in}}%
\pgfpathlineto{\pgfqpoint{1.157699in}{0.697392in}}%
\pgfpathlineto{\pgfqpoint{1.160274in}{0.695485in}}%
\pgfpathclose%
\pgfusepath{fill}%
\end{pgfscope}%
\begin{pgfscope}%
\pgfpathrectangle{\pgfqpoint{0.407592in}{0.222904in}}{\pgfqpoint{1.411163in}{1.411163in}}%
\pgfusepath{clip}%
\pgfsetbuttcap%
\pgfsetroundjoin%
\definecolor{currentfill}{rgb}{0.749020,0.772549,0.792157}%
\pgfsetfillcolor{currentfill}%
\pgfsetfillopacity{0.160000}%
\pgfsetlinewidth{0.000000pt}%
\definecolor{currentstroke}{rgb}{0.000000,0.000000,0.000000}%
\pgfsetstrokecolor{currentstroke}%
\pgfsetdash{}{0pt}%
\pgfpathmoveto{\pgfqpoint{0.687320in}{1.115715in}}%
\pgfpathlineto{\pgfqpoint{0.676188in}{1.129674in}}%
\pgfpathlineto{\pgfqpoint{0.665058in}{1.143628in}}%
\pgfpathlineto{\pgfqpoint{0.647184in}{1.132628in}}%
\pgfpathlineto{\pgfqpoint{0.631067in}{1.121196in}}%
\pgfpathlineto{\pgfqpoint{0.643173in}{1.107733in}}%
\pgfpathlineto{\pgfqpoint{0.655277in}{1.094271in}}%
\pgfpathlineto{\pgfqpoint{0.670462in}{1.105197in}}%
\pgfpathlineto{\pgfqpoint{0.687320in}{1.115715in}}%
\pgfpathclose%
\pgfusepath{fill}%
\end{pgfscope}%
\begin{pgfscope}%
\pgfpathrectangle{\pgfqpoint{0.407592in}{0.222904in}}{\pgfqpoint{1.411163in}{1.411163in}}%
\pgfusepath{clip}%
\pgfsetbuttcap%
\pgfsetroundjoin%
\definecolor{currentfill}{rgb}{0.749020,0.772549,0.792157}%
\pgfsetfillcolor{currentfill}%
\pgfsetfillopacity{0.160000}%
\pgfsetlinewidth{0.000000pt}%
\definecolor{currentstroke}{rgb}{0.000000,0.000000,0.000000}%
\pgfsetstrokecolor{currentstroke}%
\pgfsetdash{}{0pt}%
\pgfpathmoveto{\pgfqpoint{1.109619in}{0.641276in}}%
\pgfpathlineto{\pgfqpoint{1.119298in}{0.656624in}}%
\pgfpathlineto{\pgfqpoint{1.128969in}{0.671957in}}%
\pgfpathlineto{\pgfqpoint{1.126532in}{0.672985in}}%
\pgfpathlineto{\pgfqpoint{1.123939in}{0.673938in}}%
\pgfpathlineto{\pgfqpoint{1.115528in}{0.658117in}}%
\pgfpathlineto{\pgfqpoint{1.107106in}{0.642276in}}%
\pgfpathlineto{\pgfqpoint{1.108402in}{0.641795in}}%
\pgfpathlineto{\pgfqpoint{1.109619in}{0.641276in}}%
\pgfpathclose%
\pgfusepath{fill}%
\end{pgfscope}%
\begin{pgfscope}%
\pgfpathrectangle{\pgfqpoint{0.407592in}{0.222904in}}{\pgfqpoint{1.411163in}{1.411163in}}%
\pgfusepath{clip}%
\pgfsetbuttcap%
\pgfsetroundjoin%
\definecolor{currentfill}{rgb}{0.749020,0.772549,0.792157}%
\pgfsetfillcolor{currentfill}%
\pgfsetfillopacity{0.160000}%
\pgfsetlinewidth{0.000000pt}%
\definecolor{currentstroke}{rgb}{0.000000,0.000000,0.000000}%
\pgfsetstrokecolor{currentstroke}%
\pgfsetdash{}{0pt}%
\pgfpathmoveto{\pgfqpoint{1.074818in}{0.642183in}}%
\pgfpathlineto{\pgfqpoint{1.067127in}{0.657978in}}%
\pgfpathlineto{\pgfqpoint{1.059445in}{0.673753in}}%
\pgfpathlineto{\pgfqpoint{1.056684in}{0.672785in}}%
\pgfpathlineto{\pgfqpoint{1.054067in}{0.671742in}}%
\pgfpathlineto{\pgfqpoint{1.063094in}{0.656462in}}%
\pgfpathlineto{\pgfqpoint{1.072130in}{0.641167in}}%
\pgfpathlineto{\pgfqpoint{1.073438in}{0.641694in}}%
\pgfpathlineto{\pgfqpoint{1.074818in}{0.642183in}}%
\pgfpathclose%
\pgfusepath{fill}%
\end{pgfscope}%
\begin{pgfscope}%
\pgfpathrectangle{\pgfqpoint{0.407592in}{0.222904in}}{\pgfqpoint{1.411163in}{1.411163in}}%
\pgfusepath{clip}%
\pgfsetbuttcap%
\pgfsetroundjoin%
\definecolor{currentfill}{rgb}{0.749020,0.772549,0.792157}%
\pgfsetfillcolor{currentfill}%
\pgfsetfillopacity{0.160000}%
\pgfsetlinewidth{0.000000pt}%
\definecolor{currentstroke}{rgb}{0.000000,0.000000,0.000000}%
\pgfsetstrokecolor{currentstroke}%
\pgfsetdash{}{0pt}%
\pgfpathmoveto{\pgfqpoint{1.387488in}{0.934391in}}%
\pgfpathlineto{\pgfqpoint{1.399883in}{0.947896in}}%
\pgfpathlineto{\pgfqpoint{1.412279in}{0.961401in}}%
\pgfpathlineto{\pgfqpoint{1.403174in}{0.969583in}}%
\pgfpathlineto{\pgfqpoint{1.392771in}{0.977467in}}%
\pgfpathlineto{\pgfqpoint{1.381182in}{0.963411in}}%
\pgfpathlineto{\pgfqpoint{1.369589in}{0.949351in}}%
\pgfpathlineto{\pgfqpoint{1.379141in}{0.942009in}}%
\pgfpathlineto{\pgfqpoint{1.387488in}{0.934391in}}%
\pgfpathclose%
\pgfusepath{fill}%
\end{pgfscope}%
\begin{pgfscope}%
\pgfpathrectangle{\pgfqpoint{0.407592in}{0.222904in}}{\pgfqpoint{1.411163in}{1.411163in}}%
\pgfusepath{clip}%
\pgfsetbuttcap%
\pgfsetroundjoin%
\definecolor{currentfill}{rgb}{0.749020,0.772549,0.792157}%
\pgfsetfillcolor{currentfill}%
\pgfsetfillopacity{0.160000}%
\pgfsetlinewidth{0.000000pt}%
\definecolor{currentstroke}{rgb}{0.000000,0.000000,0.000000}%
\pgfsetstrokecolor{currentstroke}%
\pgfsetdash{}{0pt}%
\pgfpathmoveto{\pgfqpoint{0.709595in}{1.087786in}}%
\pgfpathlineto{\pgfqpoint{0.698456in}{1.101753in}}%
\pgfpathlineto{\pgfqpoint{0.687320in}{1.115715in}}%
\pgfpathlineto{\pgfqpoint{0.670462in}{1.105197in}}%
\pgfpathlineto{\pgfqpoint{0.655277in}{1.094271in}}%
\pgfpathlineto{\pgfqpoint{0.667380in}{1.080811in}}%
\pgfpathlineto{\pgfqpoint{0.679482in}{1.067351in}}%
\pgfpathlineto{\pgfqpoint{0.693743in}{1.077761in}}%
\pgfpathlineto{\pgfqpoint{0.709595in}{1.087786in}}%
\pgfpathclose%
\pgfusepath{fill}%
\end{pgfscope}%
\begin{pgfscope}%
\pgfpathrectangle{\pgfqpoint{0.407592in}{0.222904in}}{\pgfqpoint{1.411163in}{1.411163in}}%
\pgfusepath{clip}%
\pgfsetbuttcap%
\pgfsetroundjoin%
\definecolor{currentfill}{rgb}{0.749020,0.772549,0.792157}%
\pgfsetfillcolor{currentfill}%
\pgfsetfillopacity{0.160000}%
\pgfsetlinewidth{0.000000pt}%
\definecolor{currentstroke}{rgb}{0.000000,0.000000,0.000000}%
\pgfsetstrokecolor{currentstroke}%
\pgfsetdash{}{0pt}%
\pgfpathmoveto{\pgfqpoint{1.028888in}{0.698855in}}%
\pgfpathlineto{\pgfqpoint{1.018688in}{0.713543in}}%
\pgfpathlineto{\pgfqpoint{1.008494in}{0.728220in}}%
\pgfpathlineto{\pgfqpoint{1.004230in}{0.725789in}}%
\pgfpathlineto{\pgfqpoint{1.000333in}{0.723248in}}%
\pgfpathlineto{\pgfqpoint{1.011559in}{0.709173in}}%
\pgfpathlineto{\pgfqpoint{1.022788in}{0.695094in}}%
\pgfpathlineto{\pgfqpoint{1.025700in}{0.697016in}}%
\pgfpathlineto{\pgfqpoint{1.028888in}{0.698855in}}%
\pgfpathclose%
\pgfusepath{fill}%
\end{pgfscope}%
\begin{pgfscope}%
\pgfpathrectangle{\pgfqpoint{0.407592in}{0.222904in}}{\pgfqpoint{1.411163in}{1.411163in}}%
\pgfusepath{clip}%
\pgfsetbuttcap%
\pgfsetroundjoin%
\definecolor{currentfill}{rgb}{0.749020,0.772549,0.792157}%
\pgfsetfillcolor{currentfill}%
\pgfsetfillopacity{0.160000}%
\pgfsetlinewidth{0.000000pt}%
\definecolor{currentstroke}{rgb}{0.000000,0.000000,0.000000}%
\pgfsetstrokecolor{currentstroke}%
\pgfsetdash{}{0pt}%
\pgfpathmoveto{\pgfqpoint{0.731883in}{1.059841in}}%
\pgfpathlineto{\pgfqpoint{0.720737in}{1.073816in}}%
\pgfpathlineto{\pgfqpoint{0.709595in}{1.087786in}}%
\pgfpathlineto{\pgfqpoint{0.693743in}{1.077761in}}%
\pgfpathlineto{\pgfqpoint{0.679482in}{1.067351in}}%
\pgfpathlineto{\pgfqpoint{0.691583in}{1.053893in}}%
\pgfpathlineto{\pgfqpoint{0.703682in}{1.040437in}}%
\pgfpathlineto{\pgfqpoint{0.717030in}{1.050320in}}%
\pgfpathlineto{\pgfqpoint{0.731883in}{1.059841in}}%
\pgfpathclose%
\pgfusepath{fill}%
\end{pgfscope}%
\begin{pgfscope}%
\pgfpathrectangle{\pgfqpoint{0.407592in}{0.222904in}}{\pgfqpoint{1.411163in}{1.411163in}}%
\pgfusepath{clip}%
\pgfsetbuttcap%
\pgfsetroundjoin%
\definecolor{currentfill}{rgb}{0.749020,0.772549,0.792157}%
\pgfsetfillcolor{currentfill}%
\pgfsetfillopacity{0.160000}%
\pgfsetlinewidth{0.000000pt}%
\definecolor{currentstroke}{rgb}{0.000000,0.000000,0.000000}%
\pgfsetstrokecolor{currentstroke}%
\pgfsetdash{}{0pt}%
\pgfpathmoveto{\pgfqpoint{1.362701in}{0.907386in}}%
\pgfpathlineto{\pgfqpoint{1.375094in}{0.920888in}}%
\pgfpathlineto{\pgfqpoint{1.387488in}{0.934391in}}%
\pgfpathlineto{\pgfqpoint{1.379141in}{0.942009in}}%
\pgfpathlineto{\pgfqpoint{1.369589in}{0.949351in}}%
\pgfpathlineto{\pgfqpoint{1.357993in}{0.935286in}}%
\pgfpathlineto{\pgfqpoint{1.346392in}{0.921217in}}%
\pgfpathlineto{\pgfqpoint{1.355101in}{0.914427in}}%
\pgfpathlineto{\pgfqpoint{1.362701in}{0.907386in}}%
\pgfpathclose%
\pgfusepath{fill}%
\end{pgfscope}%
\begin{pgfscope}%
\pgfpathrectangle{\pgfqpoint{0.407592in}{0.222904in}}{\pgfqpoint{1.411163in}{1.411163in}}%
\pgfusepath{clip}%
\pgfsetbuttcap%
\pgfsetroundjoin%
\definecolor{currentfill}{rgb}{0.749020,0.772549,0.792157}%
\pgfsetfillcolor{currentfill}%
\pgfsetfillopacity{0.160000}%
\pgfsetlinewidth{0.000000pt}%
\definecolor{currentstroke}{rgb}{0.000000,0.000000,0.000000}%
\pgfsetstrokecolor{currentstroke}%
\pgfsetdash{}{0pt}%
\pgfpathmoveto{\pgfqpoint{0.754183in}{1.031880in}}%
\pgfpathlineto{\pgfqpoint{0.743031in}{1.045863in}}%
\pgfpathlineto{\pgfqpoint{0.731883in}{1.059841in}}%
\pgfpathlineto{\pgfqpoint{0.717030in}{1.050320in}}%
\pgfpathlineto{\pgfqpoint{0.703682in}{1.040437in}}%
\pgfpathlineto{\pgfqpoint{0.715780in}{1.026982in}}%
\pgfpathlineto{\pgfqpoint{0.727877in}{1.013528in}}%
\pgfpathlineto{\pgfqpoint{0.740321in}{1.022874in}}%
\pgfpathlineto{\pgfqpoint{0.754183in}{1.031880in}}%
\pgfpathclose%
\pgfusepath{fill}%
\end{pgfscope}%
\begin{pgfscope}%
\pgfpathrectangle{\pgfqpoint{0.407592in}{0.222904in}}{\pgfqpoint{1.411163in}{1.411163in}}%
\pgfusepath{clip}%
\pgfsetbuttcap%
\pgfsetroundjoin%
\definecolor{currentfill}{rgb}{0.749020,0.772549,0.792157}%
\pgfsetfillcolor{currentfill}%
\pgfsetfillopacity{0.160000}%
\pgfsetlinewidth{0.000000pt}%
\definecolor{currentstroke}{rgb}{0.000000,0.000000,0.000000}%
\pgfsetstrokecolor{currentstroke}%
\pgfsetdash{}{0pt}%
\pgfpathmoveto{\pgfqpoint{1.337916in}{0.880383in}}%
\pgfpathlineto{\pgfqpoint{1.350308in}{0.893884in}}%
\pgfpathlineto{\pgfqpoint{1.362701in}{0.907386in}}%
\pgfpathlineto{\pgfqpoint{1.355101in}{0.914427in}}%
\pgfpathlineto{\pgfqpoint{1.346392in}{0.921217in}}%
\pgfpathlineto{\pgfqpoint{1.334788in}{0.907143in}}%
\pgfpathlineto{\pgfqpoint{1.323181in}{0.893064in}}%
\pgfpathlineto{\pgfqpoint{1.331054in}{0.886838in}}%
\pgfpathlineto{\pgfqpoint{1.337916in}{0.880383in}}%
\pgfpathclose%
\pgfusepath{fill}%
\end{pgfscope}%
\begin{pgfscope}%
\pgfpathrectangle{\pgfqpoint{0.407592in}{0.222904in}}{\pgfqpoint{1.411163in}{1.411163in}}%
\pgfusepath{clip}%
\pgfsetbuttcap%
\pgfsetroundjoin%
\definecolor{currentfill}{rgb}{0.749020,0.772549,0.792157}%
\pgfsetfillcolor{currentfill}%
\pgfsetfillopacity{0.160000}%
\pgfsetlinewidth{0.000000pt}%
\definecolor{currentstroke}{rgb}{0.000000,0.000000,0.000000}%
\pgfsetstrokecolor{currentstroke}%
\pgfsetdash{}{0pt}%
\pgfpathmoveto{\pgfqpoint{0.776496in}{1.003903in}}%
\pgfpathlineto{\pgfqpoint{0.765338in}{1.017894in}}%
\pgfpathlineto{\pgfqpoint{0.754183in}{1.031880in}}%
\pgfpathlineto{\pgfqpoint{0.740321in}{1.022874in}}%
\pgfpathlineto{\pgfqpoint{0.727877in}{1.013528in}}%
\pgfpathlineto{\pgfqpoint{0.739973in}{1.000075in}}%
\pgfpathlineto{\pgfqpoint{0.752068in}{0.986624in}}%
\pgfpathlineto{\pgfqpoint{0.763616in}{0.995422in}}%
\pgfpathlineto{\pgfqpoint{0.776496in}{1.003903in}}%
\pgfpathclose%
\pgfusepath{fill}%
\end{pgfscope}%
\begin{pgfscope}%
\pgfpathrectangle{\pgfqpoint{0.407592in}{0.222904in}}{\pgfqpoint{1.411163in}{1.411163in}}%
\pgfusepath{clip}%
\pgfsetbuttcap%
\pgfsetroundjoin%
\definecolor{currentfill}{rgb}{0.749020,0.772549,0.792157}%
\pgfsetfillcolor{currentfill}%
\pgfsetfillopacity{0.160000}%
\pgfsetlinewidth{0.000000pt}%
\definecolor{currentstroke}{rgb}{0.000000,0.000000,0.000000}%
\pgfsetstrokecolor{currentstroke}%
\pgfsetdash{}{0pt}%
\pgfpathmoveto{\pgfqpoint{0.798822in}{0.975910in}}%
\pgfpathlineto{\pgfqpoint{0.787657in}{0.989909in}}%
\pgfpathlineto{\pgfqpoint{0.776496in}{1.003903in}}%
\pgfpathlineto{\pgfqpoint{0.763616in}{0.995422in}}%
\pgfpathlineto{\pgfqpoint{0.752068in}{0.986624in}}%
\pgfpathlineto{\pgfqpoint{0.764161in}{0.973174in}}%
\pgfpathlineto{\pgfqpoint{0.776253in}{0.959726in}}%
\pgfpathlineto{\pgfqpoint{0.786916in}{0.967964in}}%
\pgfpathlineto{\pgfqpoint{0.798822in}{0.975910in}}%
\pgfpathclose%
\pgfusepath{fill}%
\end{pgfscope}%
\begin{pgfscope}%
\pgfpathrectangle{\pgfqpoint{0.407592in}{0.222904in}}{\pgfqpoint{1.411163in}{1.411163in}}%
\pgfusepath{clip}%
\pgfsetbuttcap%
\pgfsetroundjoin%
\definecolor{currentfill}{rgb}{0.749020,0.772549,0.792157}%
\pgfsetfillcolor{currentfill}%
\pgfsetfillopacity{0.160000}%
\pgfsetlinewidth{0.000000pt}%
\definecolor{currentstroke}{rgb}{0.000000,0.000000,0.000000}%
\pgfsetstrokecolor{currentstroke}%
\pgfsetdash{}{0pt}%
\pgfpathmoveto{\pgfqpoint{1.313134in}{0.853384in}}%
\pgfpathlineto{\pgfqpoint{1.325524in}{0.866883in}}%
\pgfpathlineto{\pgfqpoint{1.337916in}{0.880383in}}%
\pgfpathlineto{\pgfqpoint{1.331054in}{0.886838in}}%
\pgfpathlineto{\pgfqpoint{1.323181in}{0.893064in}}%
\pgfpathlineto{\pgfqpoint{1.311569in}{0.878981in}}%
\pgfpathlineto{\pgfqpoint{1.299953in}{0.864894in}}%
\pgfpathlineto{\pgfqpoint{1.307001in}{0.859241in}}%
\pgfpathlineto{\pgfqpoint{1.313134in}{0.853384in}}%
\pgfpathclose%
\pgfusepath{fill}%
\end{pgfscope}%
\begin{pgfscope}%
\pgfpathrectangle{\pgfqpoint{0.407592in}{0.222904in}}{\pgfqpoint{1.411163in}{1.411163in}}%
\pgfusepath{clip}%
\pgfsetbuttcap%
\pgfsetroundjoin%
\definecolor{currentfill}{rgb}{0.749020,0.772549,0.792157}%
\pgfsetfillcolor{currentfill}%
\pgfsetfillopacity{0.160000}%
\pgfsetlinewidth{0.000000pt}%
\definecolor{currentstroke}{rgb}{0.000000,0.000000,0.000000}%
\pgfsetstrokecolor{currentstroke}%
\pgfsetdash{}{0pt}%
\pgfpathmoveto{\pgfqpoint{1.111794in}{0.640134in}}%
\pgfpathlineto{\pgfqpoint{1.122564in}{0.654919in}}%
\pgfpathlineto{\pgfqpoint{1.133328in}{0.669694in}}%
\pgfpathlineto{\pgfqpoint{1.131237in}{0.670859in}}%
\pgfpathlineto{\pgfqpoint{1.128969in}{0.671957in}}%
\pgfpathlineto{\pgfqpoint{1.119298in}{0.656624in}}%
\pgfpathlineto{\pgfqpoint{1.109619in}{0.641276in}}%
\pgfpathlineto{\pgfqpoint{1.110751in}{0.640721in}}%
\pgfpathlineto{\pgfqpoint{1.111794in}{0.640134in}}%
\pgfpathclose%
\pgfusepath{fill}%
\end{pgfscope}%
\begin{pgfscope}%
\pgfpathrectangle{\pgfqpoint{0.407592in}{0.222904in}}{\pgfqpoint{1.411163in}{1.411163in}}%
\pgfusepath{clip}%
\pgfsetbuttcap%
\pgfsetroundjoin%
\definecolor{currentfill}{rgb}{0.749020,0.772549,0.792157}%
\pgfsetfillcolor{currentfill}%
\pgfsetfillopacity{0.160000}%
\pgfsetlinewidth{0.000000pt}%
\definecolor{currentstroke}{rgb}{0.000000,0.000000,0.000000}%
\pgfsetstrokecolor{currentstroke}%
\pgfsetdash{}{0pt}%
\pgfpathmoveto{\pgfqpoint{1.136941in}{0.667186in}}%
\pgfpathlineto{\pgfqpoint{1.148609in}{0.681338in}}%
\pgfpathlineto{\pgfqpoint{1.160274in}{0.695485in}}%
\pgfpathlineto{\pgfqpoint{1.157699in}{0.697392in}}%
\pgfpathlineto{\pgfqpoint{1.154833in}{0.699214in}}%
\pgfpathlineto{\pgfqpoint{1.144084in}{0.684459in}}%
\pgfpathlineto{\pgfqpoint{1.133328in}{0.669694in}}%
\pgfpathlineto{\pgfqpoint{1.135232in}{0.668468in}}%
\pgfpathlineto{\pgfqpoint{1.136941in}{0.667186in}}%
\pgfpathclose%
\pgfusepath{fill}%
\end{pgfscope}%
\begin{pgfscope}%
\pgfpathrectangle{\pgfqpoint{0.407592in}{0.222904in}}{\pgfqpoint{1.411163in}{1.411163in}}%
\pgfusepath{clip}%
\pgfsetbuttcap%
\pgfsetroundjoin%
\definecolor{currentfill}{rgb}{0.749020,0.772549,0.792157}%
\pgfsetfillcolor{currentfill}%
\pgfsetfillopacity{0.160000}%
\pgfsetlinewidth{0.000000pt}%
\definecolor{currentstroke}{rgb}{0.000000,0.000000,0.000000}%
\pgfsetstrokecolor{currentstroke}%
\pgfsetdash{}{0pt}%
\pgfpathmoveto{\pgfqpoint{1.072130in}{0.641167in}}%
\pgfpathlineto{\pgfqpoint{1.063094in}{0.656462in}}%
\pgfpathlineto{\pgfqpoint{1.054067in}{0.671742in}}%
\pgfpathlineto{\pgfqpoint{1.051605in}{0.670630in}}%
\pgfpathlineto{\pgfqpoint{1.049309in}{0.669453in}}%
\pgfpathlineto{\pgfqpoint{1.059529in}{0.654737in}}%
\pgfpathlineto{\pgfqpoint{1.069756in}{0.640012in}}%
\pgfpathlineto{\pgfqpoint{1.070901in}{0.640606in}}%
\pgfpathlineto{\pgfqpoint{1.072130in}{0.641167in}}%
\pgfpathclose%
\pgfusepath{fill}%
\end{pgfscope}%
\begin{pgfscope}%
\pgfpathrectangle{\pgfqpoint{0.407592in}{0.222904in}}{\pgfqpoint{1.411163in}{1.411163in}}%
\pgfusepath{clip}%
\pgfsetbuttcap%
\pgfsetroundjoin%
\definecolor{currentfill}{rgb}{0.749020,0.772549,0.792157}%
\pgfsetfillcolor{currentfill}%
\pgfsetfillopacity{0.160000}%
\pgfsetlinewidth{0.000000pt}%
\definecolor{currentstroke}{rgb}{0.000000,0.000000,0.000000}%
\pgfsetstrokecolor{currentstroke}%
\pgfsetdash{}{0pt}%
\pgfpathmoveto{\pgfqpoint{0.821161in}{0.947901in}}%
\pgfpathlineto{\pgfqpoint{0.809990in}{0.961908in}}%
\pgfpathlineto{\pgfqpoint{0.798822in}{0.975910in}}%
\pgfpathlineto{\pgfqpoint{0.786916in}{0.967964in}}%
\pgfpathlineto{\pgfqpoint{0.776253in}{0.959726in}}%
\pgfpathlineto{\pgfqpoint{0.788344in}{0.946279in}}%
\pgfpathlineto{\pgfqpoint{0.800434in}{0.932833in}}%
\pgfpathlineto{\pgfqpoint{0.810221in}{0.940502in}}%
\pgfpathlineto{\pgfqpoint{0.821161in}{0.947901in}}%
\pgfpathclose%
\pgfusepath{fill}%
\end{pgfscope}%
\begin{pgfscope}%
\pgfpathrectangle{\pgfqpoint{0.407592in}{0.222904in}}{\pgfqpoint{1.411163in}{1.411163in}}%
\pgfusepath{clip}%
\pgfsetbuttcap%
\pgfsetroundjoin%
\definecolor{currentfill}{rgb}{0.749020,0.772549,0.792157}%
\pgfsetfillcolor{currentfill}%
\pgfsetfillopacity{0.160000}%
\pgfsetlinewidth{0.000000pt}%
\definecolor{currentstroke}{rgb}{0.000000,0.000000,0.000000}%
\pgfsetstrokecolor{currentstroke}%
\pgfsetdash{}{0pt}%
\pgfpathmoveto{\pgfqpoint{1.049309in}{0.669453in}}%
\pgfpathlineto{\pgfqpoint{1.039095in}{0.684159in}}%
\pgfpathlineto{\pgfqpoint{1.028888in}{0.698855in}}%
\pgfpathlineto{\pgfqpoint{1.025700in}{0.697016in}}%
\pgfpathlineto{\pgfqpoint{1.022788in}{0.695094in}}%
\pgfpathlineto{\pgfqpoint{1.034020in}{0.681010in}}%
\pgfpathlineto{\pgfqpoint{1.045256in}{0.666923in}}%
\pgfpathlineto{\pgfqpoint{1.047189in}{0.668215in}}%
\pgfpathlineto{\pgfqpoint{1.049309in}{0.669453in}}%
\pgfpathclose%
\pgfusepath{fill}%
\end{pgfscope}%
\begin{pgfscope}%
\pgfpathrectangle{\pgfqpoint{0.407592in}{0.222904in}}{\pgfqpoint{1.411163in}{1.411163in}}%
\pgfusepath{clip}%
\pgfsetbuttcap%
\pgfsetroundjoin%
\definecolor{currentfill}{rgb}{0.749020,0.772549,0.792157}%
\pgfsetfillcolor{currentfill}%
\pgfsetfillopacity{0.160000}%
\pgfsetlinewidth{0.000000pt}%
\definecolor{currentstroke}{rgb}{0.000000,0.000000,0.000000}%
\pgfsetstrokecolor{currentstroke}%
\pgfsetdash{}{0pt}%
\pgfpathmoveto{\pgfqpoint{1.288355in}{0.826387in}}%
\pgfpathlineto{\pgfqpoint{1.300744in}{0.839885in}}%
\pgfpathlineto{\pgfqpoint{1.313134in}{0.853384in}}%
\pgfpathlineto{\pgfqpoint{1.307001in}{0.859241in}}%
\pgfpathlineto{\pgfqpoint{1.299953in}{0.864894in}}%
\pgfpathlineto{\pgfqpoint{1.288334in}{0.850802in}}%
\pgfpathlineto{\pgfqpoint{1.276711in}{0.836705in}}%
\pgfpathlineto{\pgfqpoint{1.282942in}{0.831637in}}%
\pgfpathlineto{\pgfqpoint{1.288355in}{0.826387in}}%
\pgfpathclose%
\pgfusepath{fill}%
\end{pgfscope}%
\begin{pgfscope}%
\pgfpathrectangle{\pgfqpoint{0.407592in}{0.222904in}}{\pgfqpoint{1.411163in}{1.411163in}}%
\pgfusepath{clip}%
\pgfsetbuttcap%
\pgfsetroundjoin%
\definecolor{currentfill}{rgb}{0.749020,0.772549,0.792157}%
\pgfsetfillcolor{currentfill}%
\pgfsetfillopacity{0.160000}%
\pgfsetlinewidth{0.000000pt}%
\definecolor{currentstroke}{rgb}{0.000000,0.000000,0.000000}%
\pgfsetstrokecolor{currentstroke}%
\pgfsetdash{}{0pt}%
\pgfpathmoveto{\pgfqpoint{0.843512in}{0.919876in}}%
\pgfpathlineto{\pgfqpoint{0.832335in}{0.933891in}}%
\pgfpathlineto{\pgfqpoint{0.821161in}{0.947901in}}%
\pgfpathlineto{\pgfqpoint{0.810221in}{0.940502in}}%
\pgfpathlineto{\pgfqpoint{0.800434in}{0.932833in}}%
\pgfpathlineto{\pgfqpoint{0.812523in}{0.919388in}}%
\pgfpathlineto{\pgfqpoint{0.824610in}{0.905945in}}%
\pgfpathlineto{\pgfqpoint{0.833530in}{0.913034in}}%
\pgfpathlineto{\pgfqpoint{0.843512in}{0.919876in}}%
\pgfpathclose%
\pgfusepath{fill}%
\end{pgfscope}%
\begin{pgfscope}%
\pgfpathrectangle{\pgfqpoint{0.407592in}{0.222904in}}{\pgfqpoint{1.411163in}{1.411163in}}%
\pgfusepath{clip}%
\pgfsetbuttcap%
\pgfsetroundjoin%
\definecolor{currentfill}{rgb}{0.749020,0.772549,0.792157}%
\pgfsetfillcolor{currentfill}%
\pgfsetfillopacity{0.160000}%
\pgfsetlinewidth{0.000000pt}%
\definecolor{currentstroke}{rgb}{0.000000,0.000000,0.000000}%
\pgfsetstrokecolor{currentstroke}%
\pgfsetdash{}{0pt}%
\pgfpathmoveto{\pgfqpoint{1.263579in}{0.799394in}}%
\pgfpathlineto{\pgfqpoint{1.275967in}{0.812890in}}%
\pgfpathlineto{\pgfqpoint{1.288355in}{0.826387in}}%
\pgfpathlineto{\pgfqpoint{1.282942in}{0.831637in}}%
\pgfpathlineto{\pgfqpoint{1.276711in}{0.836705in}}%
\pgfpathlineto{\pgfqpoint{1.265085in}{0.822604in}}%
\pgfpathlineto{\pgfqpoint{1.253454in}{0.808498in}}%
\pgfpathlineto{\pgfqpoint{1.258876in}{0.804025in}}%
\pgfpathlineto{\pgfqpoint{1.263579in}{0.799394in}}%
\pgfpathclose%
\pgfusepath{fill}%
\end{pgfscope}%
\begin{pgfscope}%
\pgfpathrectangle{\pgfqpoint{0.407592in}{0.222904in}}{\pgfqpoint{1.411163in}{1.411163in}}%
\pgfusepath{clip}%
\pgfsetbuttcap%
\pgfsetroundjoin%
\definecolor{currentfill}{rgb}{0.749020,0.772549,0.792157}%
\pgfsetfillcolor{currentfill}%
\pgfsetfillopacity{0.160000}%
\pgfsetlinewidth{0.000000pt}%
\definecolor{currentstroke}{rgb}{0.000000,0.000000,0.000000}%
\pgfsetstrokecolor{currentstroke}%
\pgfsetdash{}{0pt}%
\pgfpathmoveto{\pgfqpoint{0.865877in}{0.891835in}}%
\pgfpathlineto{\pgfqpoint{0.854693in}{0.905857in}}%
\pgfpathlineto{\pgfqpoint{0.843512in}{0.919876in}}%
\pgfpathlineto{\pgfqpoint{0.833530in}{0.913034in}}%
\pgfpathlineto{\pgfqpoint{0.824610in}{0.905945in}}%
\pgfpathlineto{\pgfqpoint{0.836696in}{0.892503in}}%
\pgfpathlineto{\pgfqpoint{0.848781in}{0.879063in}}%
\pgfpathlineto{\pgfqpoint{0.856844in}{0.885560in}}%
\pgfpathlineto{\pgfqpoint{0.865877in}{0.891835in}}%
\pgfpathclose%
\pgfusepath{fill}%
\end{pgfscope}%
\begin{pgfscope}%
\pgfpathrectangle{\pgfqpoint{0.407592in}{0.222904in}}{\pgfqpoint{1.411163in}{1.411163in}}%
\pgfusepath{clip}%
\pgfsetbuttcap%
\pgfsetroundjoin%
\definecolor{currentfill}{rgb}{0.749020,0.772549,0.792157}%
\pgfsetfillcolor{currentfill}%
\pgfsetfillopacity{0.160000}%
\pgfsetlinewidth{0.000000pt}%
\definecolor{currentstroke}{rgb}{0.000000,0.000000,0.000000}%
\pgfsetstrokecolor{currentstroke}%
\pgfsetdash{}{0pt}%
\pgfpathmoveto{\pgfqpoint{1.090230in}{0.610533in}}%
\pgfpathlineto{\pgfqpoint{1.092431in}{0.627560in}}%
\pgfpathlineto{\pgfqpoint{1.094627in}{0.644555in}}%
\pgfpathlineto{\pgfqpoint{1.092905in}{0.644625in}}%
\pgfpathlineto{\pgfqpoint{1.091172in}{0.644645in}}%
\pgfpathlineto{\pgfqpoint{1.090701in}{0.627606in}}%
\pgfpathlineto{\pgfqpoint{1.090230in}{0.610533in}}%
\pgfpathlineto{\pgfqpoint{1.090230in}{0.610533in}}%
\pgfpathlineto{\pgfqpoint{1.090230in}{0.610533in}}%
\pgfpathclose%
\pgfusepath{fill}%
\end{pgfscope}%
\begin{pgfscope}%
\pgfpathrectangle{\pgfqpoint{0.407592in}{0.222904in}}{\pgfqpoint{1.411163in}{1.411163in}}%
\pgfusepath{clip}%
\pgfsetbuttcap%
\pgfsetroundjoin%
\definecolor{currentfill}{rgb}{0.749020,0.772549,0.792157}%
\pgfsetfillcolor{currentfill}%
\pgfsetfillopacity{0.160000}%
\pgfsetlinewidth{0.000000pt}%
\definecolor{currentstroke}{rgb}{0.000000,0.000000,0.000000}%
\pgfsetstrokecolor{currentstroke}%
\pgfsetdash{}{0pt}%
\pgfpathmoveto{\pgfqpoint{1.090230in}{0.610533in}}%
\pgfpathlineto{\pgfqpoint{1.090701in}{0.627606in}}%
\pgfpathlineto{\pgfqpoint{1.091172in}{0.644645in}}%
\pgfpathlineto{\pgfqpoint{1.089434in}{0.644615in}}%
\pgfpathlineto{\pgfqpoint{1.087699in}{0.644535in}}%
\pgfpathlineto{\pgfqpoint{1.088963in}{0.627550in}}%
\pgfpathlineto{\pgfqpoint{1.090230in}{0.610533in}}%
\pgfpathlineto{\pgfqpoint{1.090230in}{0.610533in}}%
\pgfpathlineto{\pgfqpoint{1.090230in}{0.610533in}}%
\pgfpathclose%
\pgfusepath{fill}%
\end{pgfscope}%
\begin{pgfscope}%
\pgfpathrectangle{\pgfqpoint{0.407592in}{0.222904in}}{\pgfqpoint{1.411163in}{1.411163in}}%
\pgfusepath{clip}%
\pgfsetbuttcap%
\pgfsetroundjoin%
\definecolor{currentfill}{rgb}{0.749020,0.772549,0.792157}%
\pgfsetfillcolor{currentfill}%
\pgfsetfillopacity{0.160000}%
\pgfsetlinewidth{0.000000pt}%
\definecolor{currentstroke}{rgb}{0.000000,0.000000,0.000000}%
\pgfsetstrokecolor{currentstroke}%
\pgfsetdash{}{0pt}%
\pgfpathmoveto{\pgfqpoint{1.090230in}{0.610533in}}%
\pgfpathlineto{\pgfqpoint{1.094121in}{0.627415in}}%
\pgfpathlineto{\pgfqpoint{1.098006in}{0.644266in}}%
\pgfpathlineto{\pgfqpoint{1.096330in}{0.644435in}}%
\pgfpathlineto{\pgfqpoint{1.094627in}{0.644555in}}%
\pgfpathlineto{\pgfqpoint{1.092431in}{0.627560in}}%
\pgfpathlineto{\pgfqpoint{1.090230in}{0.610533in}}%
\pgfpathlineto{\pgfqpoint{1.090230in}{0.610533in}}%
\pgfpathlineto{\pgfqpoint{1.090230in}{0.610533in}}%
\pgfpathclose%
\pgfusepath{fill}%
\end{pgfscope}%
\begin{pgfscope}%
\pgfpathrectangle{\pgfqpoint{0.407592in}{0.222904in}}{\pgfqpoint{1.411163in}{1.411163in}}%
\pgfusepath{clip}%
\pgfsetbuttcap%
\pgfsetroundjoin%
\definecolor{currentfill}{rgb}{0.749020,0.772549,0.792157}%
\pgfsetfillcolor{currentfill}%
\pgfsetfillopacity{0.160000}%
\pgfsetlinewidth{0.000000pt}%
\definecolor{currentstroke}{rgb}{0.000000,0.000000,0.000000}%
\pgfsetstrokecolor{currentstroke}%
\pgfsetdash{}{0pt}%
\pgfpathmoveto{\pgfqpoint{1.090230in}{0.610533in}}%
\pgfpathlineto{\pgfqpoint{1.088963in}{0.627550in}}%
\pgfpathlineto{\pgfqpoint{1.087699in}{0.644535in}}%
\pgfpathlineto{\pgfqpoint{1.085974in}{0.644405in}}%
\pgfpathlineto{\pgfqpoint{1.084268in}{0.644227in}}%
\pgfpathlineto{\pgfqpoint{1.087246in}{0.627395in}}%
\pgfpathlineto{\pgfqpoint{1.090230in}{0.610533in}}%
\pgfpathlineto{\pgfqpoint{1.090230in}{0.610533in}}%
\pgfpathlineto{\pgfqpoint{1.090230in}{0.610533in}}%
\pgfpathclose%
\pgfusepath{fill}%
\end{pgfscope}%
\begin{pgfscope}%
\pgfpathrectangle{\pgfqpoint{0.407592in}{0.222904in}}{\pgfqpoint{1.411163in}{1.411163in}}%
\pgfusepath{clip}%
\pgfsetbuttcap%
\pgfsetroundjoin%
\definecolor{currentfill}{rgb}{0.749020,0.772549,0.792157}%
\pgfsetfillcolor{currentfill}%
\pgfsetfillopacity{0.160000}%
\pgfsetlinewidth{0.000000pt}%
\definecolor{currentstroke}{rgb}{0.000000,0.000000,0.000000}%
\pgfsetstrokecolor{currentstroke}%
\pgfsetdash{}{0pt}%
\pgfpathmoveto{\pgfqpoint{0.888254in}{0.863777in}}%
\pgfpathlineto{\pgfqpoint{0.877064in}{0.877808in}}%
\pgfpathlineto{\pgfqpoint{0.865877in}{0.891835in}}%
\pgfpathlineto{\pgfqpoint{0.856844in}{0.885560in}}%
\pgfpathlineto{\pgfqpoint{0.848781in}{0.879063in}}%
\pgfpathlineto{\pgfqpoint{0.860865in}{0.865624in}}%
\pgfpathlineto{\pgfqpoint{0.872948in}{0.852186in}}%
\pgfpathlineto{\pgfqpoint{0.880162in}{0.858082in}}%
\pgfpathlineto{\pgfqpoint{0.888254in}{0.863777in}}%
\pgfpathclose%
\pgfusepath{fill}%
\end{pgfscope}%
\begin{pgfscope}%
\pgfpathrectangle{\pgfqpoint{0.407592in}{0.222904in}}{\pgfqpoint{1.411163in}{1.411163in}}%
\pgfusepath{clip}%
\pgfsetbuttcap%
\pgfsetroundjoin%
\definecolor{currentfill}{rgb}{0.749020,0.772549,0.792157}%
\pgfsetfillcolor{currentfill}%
\pgfsetfillopacity{0.160000}%
\pgfsetlinewidth{0.000000pt}%
\definecolor{currentstroke}{rgb}{0.000000,0.000000,0.000000}%
\pgfsetstrokecolor{currentstroke}%
\pgfsetdash{}{0pt}%
\pgfpathmoveto{\pgfqpoint{1.090230in}{0.610533in}}%
\pgfpathlineto{\pgfqpoint{1.095744in}{0.627173in}}%
\pgfpathlineto{\pgfqpoint{1.101249in}{0.643784in}}%
\pgfpathlineto{\pgfqpoint{1.099648in}{0.644049in}}%
\pgfpathlineto{\pgfqpoint{1.098006in}{0.644266in}}%
\pgfpathlineto{\pgfqpoint{1.094121in}{0.627415in}}%
\pgfpathlineto{\pgfqpoint{1.090230in}{0.610533in}}%
\pgfpathlineto{\pgfqpoint{1.090230in}{0.610533in}}%
\pgfpathlineto{\pgfqpoint{1.090230in}{0.610533in}}%
\pgfpathclose%
\pgfusepath{fill}%
\end{pgfscope}%
\begin{pgfscope}%
\pgfpathrectangle{\pgfqpoint{0.407592in}{0.222904in}}{\pgfqpoint{1.411163in}{1.411163in}}%
\pgfusepath{clip}%
\pgfsetbuttcap%
\pgfsetroundjoin%
\definecolor{currentfill}{rgb}{0.749020,0.772549,0.792157}%
\pgfsetfillcolor{currentfill}%
\pgfsetfillopacity{0.160000}%
\pgfsetlinewidth{0.000000pt}%
\definecolor{currentstroke}{rgb}{0.000000,0.000000,0.000000}%
\pgfsetstrokecolor{currentstroke}%
\pgfsetdash{}{0pt}%
\pgfpathmoveto{\pgfqpoint{1.090230in}{0.610533in}}%
\pgfpathlineto{\pgfqpoint{1.087246in}{0.627395in}}%
\pgfpathlineto{\pgfqpoint{1.084268in}{0.644227in}}%
\pgfpathlineto{\pgfqpoint{1.082588in}{0.644000in}}%
\pgfpathlineto{\pgfqpoint{1.080940in}{0.643726in}}%
\pgfpathlineto{\pgfqpoint{1.085581in}{0.627143in}}%
\pgfpathlineto{\pgfqpoint{1.090230in}{0.610533in}}%
\pgfpathlineto{\pgfqpoint{1.090230in}{0.610533in}}%
\pgfpathlineto{\pgfqpoint{1.090230in}{0.610533in}}%
\pgfpathclose%
\pgfusepath{fill}%
\end{pgfscope}%
\begin{pgfscope}%
\pgfpathrectangle{\pgfqpoint{0.407592in}{0.222904in}}{\pgfqpoint{1.411163in}{1.411163in}}%
\pgfusepath{clip}%
\pgfsetbuttcap%
\pgfsetroundjoin%
\definecolor{currentfill}{rgb}{0.749020,0.772549,0.792157}%
\pgfsetfillcolor{currentfill}%
\pgfsetfillopacity{0.160000}%
\pgfsetlinewidth{0.000000pt}%
\definecolor{currentstroke}{rgb}{0.000000,0.000000,0.000000}%
\pgfsetstrokecolor{currentstroke}%
\pgfsetdash{}{0pt}%
\pgfpathmoveto{\pgfqpoint{1.238806in}{0.772404in}}%
\pgfpathlineto{\pgfqpoint{1.251192in}{0.785899in}}%
\pgfpathlineto{\pgfqpoint{1.263579in}{0.799394in}}%
\pgfpathlineto{\pgfqpoint{1.258876in}{0.804025in}}%
\pgfpathlineto{\pgfqpoint{1.253454in}{0.808498in}}%
\pgfpathlineto{\pgfqpoint{1.241820in}{0.794387in}}%
\pgfpathlineto{\pgfqpoint{1.230182in}{0.780272in}}%
\pgfpathlineto{\pgfqpoint{1.234803in}{0.776406in}}%
\pgfpathlineto{\pgfqpoint{1.238806in}{0.772404in}}%
\pgfpathclose%
\pgfusepath{fill}%
\end{pgfscope}%
\begin{pgfscope}%
\pgfpathrectangle{\pgfqpoint{0.407592in}{0.222904in}}{\pgfqpoint{1.411163in}{1.411163in}}%
\pgfusepath{clip}%
\pgfsetbuttcap%
\pgfsetroundjoin%
\definecolor{currentfill}{rgb}{0.749020,0.772549,0.792157}%
\pgfsetfillcolor{currentfill}%
\pgfsetfillopacity{0.160000}%
\pgfsetlinewidth{0.000000pt}%
\definecolor{currentstroke}{rgb}{0.000000,0.000000,0.000000}%
\pgfsetstrokecolor{currentstroke}%
\pgfsetdash{}{0pt}%
\pgfpathmoveto{\pgfqpoint{1.090230in}{0.610533in}}%
\pgfpathlineto{\pgfqpoint{1.097270in}{0.626837in}}%
\pgfpathlineto{\pgfqpoint{1.104300in}{0.643117in}}%
\pgfpathlineto{\pgfqpoint{1.102802in}{0.643473in}}%
\pgfpathlineto{\pgfqpoint{1.101249in}{0.643784in}}%
\pgfpathlineto{\pgfqpoint{1.095744in}{0.627173in}}%
\pgfpathlineto{\pgfqpoint{1.090230in}{0.610533in}}%
\pgfpathlineto{\pgfqpoint{1.090230in}{0.610533in}}%
\pgfpathlineto{\pgfqpoint{1.090230in}{0.610533in}}%
\pgfpathclose%
\pgfusepath{fill}%
\end{pgfscope}%
\begin{pgfscope}%
\pgfpathrectangle{\pgfqpoint{0.407592in}{0.222904in}}{\pgfqpoint{1.411163in}{1.411163in}}%
\pgfusepath{clip}%
\pgfsetbuttcap%
\pgfsetroundjoin%
\definecolor{currentfill}{rgb}{0.749020,0.772549,0.792157}%
\pgfsetfillcolor{currentfill}%
\pgfsetfillopacity{0.160000}%
\pgfsetlinewidth{0.000000pt}%
\definecolor{currentstroke}{rgb}{0.000000,0.000000,0.000000}%
\pgfsetstrokecolor{currentstroke}%
\pgfsetdash{}{0pt}%
\pgfpathmoveto{\pgfqpoint{0.910644in}{0.835704in}}%
\pgfpathlineto{\pgfqpoint{0.899447in}{0.849743in}}%
\pgfpathlineto{\pgfqpoint{0.888254in}{0.863777in}}%
\pgfpathlineto{\pgfqpoint{0.880162in}{0.858082in}}%
\pgfpathlineto{\pgfqpoint{0.872948in}{0.852186in}}%
\pgfpathlineto{\pgfqpoint{0.885029in}{0.838749in}}%
\pgfpathlineto{\pgfqpoint{0.897109in}{0.825314in}}%
\pgfpathlineto{\pgfqpoint{0.903484in}{0.830598in}}%
\pgfpathlineto{\pgfqpoint{0.910644in}{0.835704in}}%
\pgfpathclose%
\pgfusepath{fill}%
\end{pgfscope}%
\begin{pgfscope}%
\pgfpathrectangle{\pgfqpoint{0.407592in}{0.222904in}}{\pgfqpoint{1.411163in}{1.411163in}}%
\pgfusepath{clip}%
\pgfsetbuttcap%
\pgfsetroundjoin%
\definecolor{currentfill}{rgb}{0.749020,0.772549,0.792157}%
\pgfsetfillcolor{currentfill}%
\pgfsetfillopacity{0.160000}%
\pgfsetlinewidth{0.000000pt}%
\definecolor{currentstroke}{rgb}{0.000000,0.000000,0.000000}%
\pgfsetstrokecolor{currentstroke}%
\pgfsetdash{}{0pt}%
\pgfpathmoveto{\pgfqpoint{1.090230in}{0.610533in}}%
\pgfpathlineto{\pgfqpoint{1.085581in}{0.627143in}}%
\pgfpathlineto{\pgfqpoint{1.080940in}{0.643726in}}%
\pgfpathlineto{\pgfqpoint{1.079332in}{0.643405in}}%
\pgfpathlineto{\pgfqpoint{1.077772in}{0.643040in}}%
\pgfpathlineto{\pgfqpoint{1.083996in}{0.626799in}}%
\pgfpathlineto{\pgfqpoint{1.090230in}{0.610533in}}%
\pgfpathlineto{\pgfqpoint{1.090230in}{0.610533in}}%
\pgfpathlineto{\pgfqpoint{1.090230in}{0.610533in}}%
\pgfpathclose%
\pgfusepath{fill}%
\end{pgfscope}%
\begin{pgfscope}%
\pgfpathrectangle{\pgfqpoint{0.407592in}{0.222904in}}{\pgfqpoint{1.411163in}{1.411163in}}%
\pgfusepath{clip}%
\pgfsetbuttcap%
\pgfsetroundjoin%
\definecolor{currentfill}{rgb}{0.749020,0.772549,0.792157}%
\pgfsetfillcolor{currentfill}%
\pgfsetfillopacity{0.160000}%
\pgfsetlinewidth{0.000000pt}%
\definecolor{currentstroke}{rgb}{0.000000,0.000000,0.000000}%
\pgfsetstrokecolor{currentstroke}%
\pgfsetdash{}{0pt}%
\pgfpathmoveto{\pgfqpoint{1.113593in}{0.638869in}}%
\pgfpathlineto{\pgfqpoint{1.125269in}{0.653030in}}%
\pgfpathlineto{\pgfqpoint{1.136941in}{0.667186in}}%
\pgfpathlineto{\pgfqpoint{1.135232in}{0.668468in}}%
\pgfpathlineto{\pgfqpoint{1.133328in}{0.669694in}}%
\pgfpathlineto{\pgfqpoint{1.122564in}{0.654919in}}%
\pgfpathlineto{\pgfqpoint{1.111794in}{0.640134in}}%
\pgfpathlineto{\pgfqpoint{1.112742in}{0.639515in}}%
\pgfpathlineto{\pgfqpoint{1.113593in}{0.638869in}}%
\pgfpathclose%
\pgfusepath{fill}%
\end{pgfscope}%
\begin{pgfscope}%
\pgfpathrectangle{\pgfqpoint{0.407592in}{0.222904in}}{\pgfqpoint{1.411163in}{1.411163in}}%
\pgfusepath{clip}%
\pgfsetbuttcap%
\pgfsetroundjoin%
\definecolor{currentfill}{rgb}{0.749020,0.772549,0.792157}%
\pgfsetfillcolor{currentfill}%
\pgfsetfillopacity{0.160000}%
\pgfsetlinewidth{0.000000pt}%
\definecolor{currentstroke}{rgb}{0.000000,0.000000,0.000000}%
\pgfsetstrokecolor{currentstroke}%
\pgfsetdash{}{0pt}%
\pgfpathmoveto{\pgfqpoint{1.214036in}{0.745418in}}%
\pgfpathlineto{\pgfqpoint{1.226421in}{0.758911in}}%
\pgfpathlineto{\pgfqpoint{1.238806in}{0.772404in}}%
\pgfpathlineto{\pgfqpoint{1.234803in}{0.776406in}}%
\pgfpathlineto{\pgfqpoint{1.230182in}{0.780272in}}%
\pgfpathlineto{\pgfqpoint{1.218540in}{0.766152in}}%
\pgfpathlineto{\pgfqpoint{1.206894in}{0.752028in}}%
\pgfpathlineto{\pgfqpoint{1.210724in}{0.748779in}}%
\pgfpathlineto{\pgfqpoint{1.214036in}{0.745418in}}%
\pgfpathclose%
\pgfusepath{fill}%
\end{pgfscope}%
\begin{pgfscope}%
\pgfpathrectangle{\pgfqpoint{0.407592in}{0.222904in}}{\pgfqpoint{1.411163in}{1.411163in}}%
\pgfusepath{clip}%
\pgfsetbuttcap%
\pgfsetroundjoin%
\definecolor{currentfill}{rgb}{0.749020,0.772549,0.792157}%
\pgfsetfillcolor{currentfill}%
\pgfsetfillopacity{0.160000}%
\pgfsetlinewidth{0.000000pt}%
\definecolor{currentstroke}{rgb}{0.000000,0.000000,0.000000}%
\pgfsetstrokecolor{currentstroke}%
\pgfsetdash{}{0pt}%
\pgfpathmoveto{\pgfqpoint{0.933047in}{0.807614in}}%
\pgfpathlineto{\pgfqpoint{0.921844in}{0.821661in}}%
\pgfpathlineto{\pgfqpoint{0.910644in}{0.835704in}}%
\pgfpathlineto{\pgfqpoint{0.903484in}{0.830598in}}%
\pgfpathlineto{\pgfqpoint{0.897109in}{0.825314in}}%
\pgfpathlineto{\pgfqpoint{0.909188in}{0.811880in}}%
\pgfpathlineto{\pgfqpoint{0.921266in}{0.798448in}}%
\pgfpathlineto{\pgfqpoint{0.926812in}{0.803108in}}%
\pgfpathlineto{\pgfqpoint{0.933047in}{0.807614in}}%
\pgfpathclose%
\pgfusepath{fill}%
\end{pgfscope}%
\begin{pgfscope}%
\pgfpathrectangle{\pgfqpoint{0.407592in}{0.222904in}}{\pgfqpoint{1.411163in}{1.411163in}}%
\pgfusepath{clip}%
\pgfsetbuttcap%
\pgfsetroundjoin%
\definecolor{currentfill}{rgb}{0.749020,0.772549,0.792157}%
\pgfsetfillcolor{currentfill}%
\pgfsetfillopacity{0.160000}%
\pgfsetlinewidth{0.000000pt}%
\definecolor{currentstroke}{rgb}{0.000000,0.000000,0.000000}%
\pgfsetstrokecolor{currentstroke}%
\pgfsetdash{}{0pt}%
\pgfpathmoveto{\pgfqpoint{1.090230in}{0.610533in}}%
\pgfpathlineto{\pgfqpoint{1.098673in}{0.626415in}}%
\pgfpathlineto{\pgfqpoint{1.107106in}{0.642276in}}%
\pgfpathlineto{\pgfqpoint{1.105737in}{0.642718in}}%
\pgfpathlineto{\pgfqpoint{1.104300in}{0.643117in}}%
\pgfpathlineto{\pgfqpoint{1.097270in}{0.626837in}}%
\pgfpathlineto{\pgfqpoint{1.090230in}{0.610533in}}%
\pgfpathlineto{\pgfqpoint{1.090230in}{0.610533in}}%
\pgfpathlineto{\pgfqpoint{1.090230in}{0.610533in}}%
\pgfpathclose%
\pgfusepath{fill}%
\end{pgfscope}%
\begin{pgfscope}%
\pgfpathrectangle{\pgfqpoint{0.407592in}{0.222904in}}{\pgfqpoint{1.411163in}{1.411163in}}%
\pgfusepath{clip}%
\pgfsetbuttcap%
\pgfsetroundjoin%
\definecolor{currentfill}{rgb}{0.749020,0.772549,0.792157}%
\pgfsetfillcolor{currentfill}%
\pgfsetfillopacity{0.160000}%
\pgfsetlinewidth{0.000000pt}%
\definecolor{currentstroke}{rgb}{0.000000,0.000000,0.000000}%
\pgfsetstrokecolor{currentstroke}%
\pgfsetdash{}{0pt}%
\pgfpathmoveto{\pgfqpoint{1.069756in}{0.640012in}}%
\pgfpathlineto{\pgfqpoint{1.059529in}{0.654737in}}%
\pgfpathlineto{\pgfqpoint{1.049309in}{0.669453in}}%
\pgfpathlineto{\pgfqpoint{1.047189in}{0.668215in}}%
\pgfpathlineto{\pgfqpoint{1.045256in}{0.666923in}}%
\pgfpathlineto{\pgfqpoint{1.056494in}{0.652831in}}%
\pgfpathlineto{\pgfqpoint{1.067736in}{0.638736in}}%
\pgfpathlineto{\pgfqpoint{1.068699in}{0.639388in}}%
\pgfpathlineto{\pgfqpoint{1.069756in}{0.640012in}}%
\pgfpathclose%
\pgfusepath{fill}%
\end{pgfscope}%
\begin{pgfscope}%
\pgfpathrectangle{\pgfqpoint{0.407592in}{0.222904in}}{\pgfqpoint{1.411163in}{1.411163in}}%
\pgfusepath{clip}%
\pgfsetbuttcap%
\pgfsetroundjoin%
\definecolor{currentfill}{rgb}{0.749020,0.772549,0.792157}%
\pgfsetfillcolor{currentfill}%
\pgfsetfillopacity{0.160000}%
\pgfsetlinewidth{0.000000pt}%
\definecolor{currentstroke}{rgb}{0.000000,0.000000,0.000000}%
\pgfsetstrokecolor{currentstroke}%
\pgfsetdash{}{0pt}%
\pgfpathmoveto{\pgfqpoint{1.090230in}{0.610533in}}%
\pgfpathlineto{\pgfqpoint{1.083996in}{0.626799in}}%
\pgfpathlineto{\pgfqpoint{1.077772in}{0.643040in}}%
\pgfpathlineto{\pgfqpoint{1.076265in}{0.642632in}}%
\pgfpathlineto{\pgfqpoint{1.074818in}{0.642183in}}%
\pgfpathlineto{\pgfqpoint{1.082519in}{0.626368in}}%
\pgfpathlineto{\pgfqpoint{1.090230in}{0.610533in}}%
\pgfpathlineto{\pgfqpoint{1.090230in}{0.610533in}}%
\pgfpathlineto{\pgfqpoint{1.090230in}{0.610533in}}%
\pgfpathclose%
\pgfusepath{fill}%
\end{pgfscope}%
\begin{pgfscope}%
\pgfpathrectangle{\pgfqpoint{0.407592in}{0.222904in}}{\pgfqpoint{1.411163in}{1.411163in}}%
\pgfusepath{clip}%
\pgfsetbuttcap%
\pgfsetroundjoin%
\definecolor{currentfill}{rgb}{0.749020,0.772549,0.792157}%
\pgfsetfillcolor{currentfill}%
\pgfsetfillopacity{0.160000}%
\pgfsetlinewidth{0.000000pt}%
\definecolor{currentstroke}{rgb}{0.000000,0.000000,0.000000}%
\pgfsetstrokecolor{currentstroke}%
\pgfsetdash{}{0pt}%
\pgfpathmoveto{\pgfqpoint{0.955463in}{0.779508in}}%
\pgfpathlineto{\pgfqpoint{0.944253in}{0.793563in}}%
\pgfpathlineto{\pgfqpoint{0.933047in}{0.807614in}}%
\pgfpathlineto{\pgfqpoint{0.926812in}{0.803108in}}%
\pgfpathlineto{\pgfqpoint{0.921266in}{0.798448in}}%
\pgfpathlineto{\pgfqpoint{0.933343in}{0.785017in}}%
\pgfpathlineto{\pgfqpoint{0.945418in}{0.771587in}}%
\pgfpathlineto{\pgfqpoint{0.950143in}{0.775614in}}%
\pgfpathlineto{\pgfqpoint{0.955463in}{0.779508in}}%
\pgfpathclose%
\pgfusepath{fill}%
\end{pgfscope}%
\begin{pgfscope}%
\pgfpathrectangle{\pgfqpoint{0.407592in}{0.222904in}}{\pgfqpoint{1.411163in}{1.411163in}}%
\pgfusepath{clip}%
\pgfsetbuttcap%
\pgfsetroundjoin%
\definecolor{currentfill}{rgb}{0.749020,0.772549,0.792157}%
\pgfsetfillcolor{currentfill}%
\pgfsetfillopacity{0.160000}%
\pgfsetlinewidth{0.000000pt}%
\definecolor{currentstroke}{rgb}{0.000000,0.000000,0.000000}%
\pgfsetstrokecolor{currentstroke}%
\pgfsetdash{}{0pt}%
\pgfpathmoveto{\pgfqpoint{1.189269in}{0.718434in}}%
\pgfpathlineto{\pgfqpoint{1.201652in}{0.731926in}}%
\pgfpathlineto{\pgfqpoint{1.214036in}{0.745418in}}%
\pgfpathlineto{\pgfqpoint{1.210724in}{0.748779in}}%
\pgfpathlineto{\pgfqpoint{1.206894in}{0.752028in}}%
\pgfpathlineto{\pgfqpoint{1.195245in}{0.737899in}}%
\pgfpathlineto{\pgfqpoint{1.183592in}{0.723766in}}%
\pgfpathlineto{\pgfqpoint{1.186638in}{0.721145in}}%
\pgfpathlineto{\pgfqpoint{1.189269in}{0.718434in}}%
\pgfpathclose%
\pgfusepath{fill}%
\end{pgfscope}%
\begin{pgfscope}%
\pgfpathrectangle{\pgfqpoint{0.407592in}{0.222904in}}{\pgfqpoint{1.411163in}{1.411163in}}%
\pgfusepath{clip}%
\pgfsetbuttcap%
\pgfsetroundjoin%
\definecolor{currentfill}{rgb}{0.749020,0.772549,0.792157}%
\pgfsetfillcolor{currentfill}%
\pgfsetfillopacity{0.160000}%
\pgfsetlinewidth{0.000000pt}%
\definecolor{currentstroke}{rgb}{0.000000,0.000000,0.000000}%
\pgfsetstrokecolor{currentstroke}%
\pgfsetdash{}{0pt}%
\pgfpathmoveto{\pgfqpoint{1.090230in}{0.610533in}}%
\pgfpathlineto{\pgfqpoint{1.099929in}{0.625912in}}%
\pgfpathlineto{\pgfqpoint{1.109619in}{0.641276in}}%
\pgfpathlineto{\pgfqpoint{1.108402in}{0.641795in}}%
\pgfpathlineto{\pgfqpoint{1.107106in}{0.642276in}}%
\pgfpathlineto{\pgfqpoint{1.098673in}{0.626415in}}%
\pgfpathlineto{\pgfqpoint{1.090230in}{0.610533in}}%
\pgfpathlineto{\pgfqpoint{1.090230in}{0.610533in}}%
\pgfpathlineto{\pgfqpoint{1.090230in}{0.610533in}}%
\pgfpathclose%
\pgfusepath{fill}%
\end{pgfscope}%
\begin{pgfscope}%
\pgfpathrectangle{\pgfqpoint{0.407592in}{0.222904in}}{\pgfqpoint{1.411163in}{1.411163in}}%
\pgfusepath{clip}%
\pgfsetbuttcap%
\pgfsetroundjoin%
\definecolor{currentfill}{rgb}{0.749020,0.772549,0.792157}%
\pgfsetfillcolor{currentfill}%
\pgfsetfillopacity{0.160000}%
\pgfsetlinewidth{0.000000pt}%
\definecolor{currentstroke}{rgb}{0.000000,0.000000,0.000000}%
\pgfsetstrokecolor{currentstroke}%
\pgfsetdash{}{0pt}%
\pgfpathmoveto{\pgfqpoint{1.090230in}{0.610533in}}%
\pgfpathlineto{\pgfqpoint{1.082519in}{0.626368in}}%
\pgfpathlineto{\pgfqpoint{1.074818in}{0.642183in}}%
\pgfpathlineto{\pgfqpoint{1.073438in}{0.641694in}}%
\pgfpathlineto{\pgfqpoint{1.072130in}{0.641167in}}%
\pgfpathlineto{\pgfqpoint{1.081176in}{0.625858in}}%
\pgfpathlineto{\pgfqpoint{1.090230in}{0.610533in}}%
\pgfpathlineto{\pgfqpoint{1.090230in}{0.610533in}}%
\pgfpathlineto{\pgfqpoint{1.090230in}{0.610533in}}%
\pgfpathclose%
\pgfusepath{fill}%
\end{pgfscope}%
\begin{pgfscope}%
\pgfpathrectangle{\pgfqpoint{0.407592in}{0.222904in}}{\pgfqpoint{1.411163in}{1.411163in}}%
\pgfusepath{clip}%
\pgfsetbuttcap%
\pgfsetroundjoin%
\definecolor{currentfill}{rgb}{0.749020,0.772549,0.792157}%
\pgfsetfillcolor{currentfill}%
\pgfsetfillopacity{0.160000}%
\pgfsetlinewidth{0.000000pt}%
\definecolor{currentstroke}{rgb}{0.000000,0.000000,0.000000}%
\pgfsetstrokecolor{currentstroke}%
\pgfsetdash{}{0pt}%
\pgfpathmoveto{\pgfqpoint{0.977892in}{0.751386in}}%
\pgfpathlineto{\pgfqpoint{0.966676in}{0.765449in}}%
\pgfpathlineto{\pgfqpoint{0.955463in}{0.779508in}}%
\pgfpathlineto{\pgfqpoint{0.950143in}{0.775614in}}%
\pgfpathlineto{\pgfqpoint{0.945418in}{0.771587in}}%
\pgfpathlineto{\pgfqpoint{0.957492in}{0.758158in}}%
\pgfpathlineto{\pgfqpoint{0.969565in}{0.744731in}}%
\pgfpathlineto{\pgfqpoint{0.973480in}{0.748114in}}%
\pgfpathlineto{\pgfqpoint{0.977892in}{0.751386in}}%
\pgfpathclose%
\pgfusepath{fill}%
\end{pgfscope}%
\begin{pgfscope}%
\pgfpathrectangle{\pgfqpoint{0.407592in}{0.222904in}}{\pgfqpoint{1.411163in}{1.411163in}}%
\pgfusepath{clip}%
\pgfsetbuttcap%
\pgfsetroundjoin%
\definecolor{currentfill}{rgb}{0.749020,0.772549,0.792157}%
\pgfsetfillcolor{currentfill}%
\pgfsetfillopacity{0.160000}%
\pgfsetlinewidth{0.000000pt}%
\definecolor{currentstroke}{rgb}{0.000000,0.000000,0.000000}%
\pgfsetstrokecolor{currentstroke}%
\pgfsetdash{}{0pt}%
\pgfpathmoveto{\pgfqpoint{1.164505in}{0.691454in}}%
\pgfpathlineto{\pgfqpoint{1.176886in}{0.704944in}}%
\pgfpathlineto{\pgfqpoint{1.189269in}{0.718434in}}%
\pgfpathlineto{\pgfqpoint{1.186638in}{0.721145in}}%
\pgfpathlineto{\pgfqpoint{1.183592in}{0.723766in}}%
\pgfpathlineto{\pgfqpoint{1.171935in}{0.709628in}}%
\pgfpathlineto{\pgfqpoint{1.160274in}{0.695485in}}%
\pgfpathlineto{\pgfqpoint{1.162546in}{0.693503in}}%
\pgfpathlineto{\pgfqpoint{1.164505in}{0.691454in}}%
\pgfpathclose%
\pgfusepath{fill}%
\end{pgfscope}%
\begin{pgfscope}%
\pgfpathrectangle{\pgfqpoint{0.407592in}{0.222904in}}{\pgfqpoint{1.411163in}{1.411163in}}%
\pgfusepath{clip}%
\pgfsetbuttcap%
\pgfsetroundjoin%
\definecolor{currentfill}{rgb}{0.749020,0.772549,0.792157}%
\pgfsetfillcolor{currentfill}%
\pgfsetfillopacity{0.160000}%
\pgfsetlinewidth{0.000000pt}%
\definecolor{currentstroke}{rgb}{0.000000,0.000000,0.000000}%
\pgfsetstrokecolor{currentstroke}%
\pgfsetdash{}{0pt}%
\pgfpathmoveto{\pgfqpoint{1.000333in}{0.723248in}}%
\pgfpathlineto{\pgfqpoint{0.989111in}{0.737319in}}%
\pgfpathlineto{\pgfqpoint{0.977892in}{0.751386in}}%
\pgfpathlineto{\pgfqpoint{0.973480in}{0.748114in}}%
\pgfpathlineto{\pgfqpoint{0.969565in}{0.744731in}}%
\pgfpathlineto{\pgfqpoint{0.981637in}{0.731305in}}%
\pgfpathlineto{\pgfqpoint{0.993708in}{0.717881in}}%
\pgfpathlineto{\pgfqpoint{0.996821in}{0.720608in}}%
\pgfpathlineto{\pgfqpoint{1.000333in}{0.723248in}}%
\pgfpathclose%
\pgfusepath{fill}%
\end{pgfscope}%
\begin{pgfscope}%
\pgfpathrectangle{\pgfqpoint{0.407592in}{0.222904in}}{\pgfqpoint{1.411163in}{1.411163in}}%
\pgfusepath{clip}%
\pgfsetbuttcap%
\pgfsetroundjoin%
\definecolor{currentfill}{rgb}{0.749020,0.772549,0.792157}%
\pgfsetfillcolor{currentfill}%
\pgfsetfillopacity{0.160000}%
\pgfsetlinewidth{0.000000pt}%
\definecolor{currentstroke}{rgb}{0.000000,0.000000,0.000000}%
\pgfsetstrokecolor{currentstroke}%
\pgfsetdash{}{0pt}%
\pgfpathmoveto{\pgfqpoint{1.090230in}{0.610533in}}%
\pgfpathlineto{\pgfqpoint{1.101015in}{0.625338in}}%
\pgfpathlineto{\pgfqpoint{1.111794in}{0.640134in}}%
\pgfpathlineto{\pgfqpoint{1.110751in}{0.640721in}}%
\pgfpathlineto{\pgfqpoint{1.109619in}{0.641276in}}%
\pgfpathlineto{\pgfqpoint{1.099929in}{0.625912in}}%
\pgfpathlineto{\pgfqpoint{1.090230in}{0.610533in}}%
\pgfpathlineto{\pgfqpoint{1.090230in}{0.610533in}}%
\pgfpathlineto{\pgfqpoint{1.090230in}{0.610533in}}%
\pgfpathclose%
\pgfusepath{fill}%
\end{pgfscope}%
\begin{pgfscope}%
\pgfpathrectangle{\pgfqpoint{0.407592in}{0.222904in}}{\pgfqpoint{1.411163in}{1.411163in}}%
\pgfusepath{clip}%
\pgfsetbuttcap%
\pgfsetroundjoin%
\definecolor{currentfill}{rgb}{0.749020,0.772549,0.792157}%
\pgfsetfillcolor{currentfill}%
\pgfsetfillopacity{0.160000}%
\pgfsetlinewidth{0.000000pt}%
\definecolor{currentstroke}{rgb}{0.000000,0.000000,0.000000}%
\pgfsetstrokecolor{currentstroke}%
\pgfsetdash{}{0pt}%
\pgfpathmoveto{\pgfqpoint{1.090230in}{0.610533in}}%
\pgfpathlineto{\pgfqpoint{1.081176in}{0.625858in}}%
\pgfpathlineto{\pgfqpoint{1.072130in}{0.641167in}}%
\pgfpathlineto{\pgfqpoint{1.070901in}{0.640606in}}%
\pgfpathlineto{\pgfqpoint{1.069756in}{0.640012in}}%
\pgfpathlineto{\pgfqpoint{1.079989in}{0.625277in}}%
\pgfpathlineto{\pgfqpoint{1.090230in}{0.610533in}}%
\pgfpathlineto{\pgfqpoint{1.090230in}{0.610533in}}%
\pgfpathlineto{\pgfqpoint{1.090230in}{0.610533in}}%
\pgfpathclose%
\pgfusepath{fill}%
\end{pgfscope}%
\begin{pgfscope}%
\pgfpathrectangle{\pgfqpoint{0.407592in}{0.222904in}}{\pgfqpoint{1.411163in}{1.411163in}}%
\pgfusepath{clip}%
\pgfsetbuttcap%
\pgfsetroundjoin%
\definecolor{currentfill}{rgb}{0.749020,0.772549,0.792157}%
\pgfsetfillcolor{currentfill}%
\pgfsetfillopacity{0.160000}%
\pgfsetlinewidth{0.000000pt}%
\definecolor{currentstroke}{rgb}{0.000000,0.000000,0.000000}%
\pgfsetstrokecolor{currentstroke}%
\pgfsetdash{}{0pt}%
\pgfpathmoveto{\pgfqpoint{1.139743in}{0.664477in}}%
\pgfpathlineto{\pgfqpoint{1.152124in}{0.677965in}}%
\pgfpathlineto{\pgfqpoint{1.164505in}{0.691454in}}%
\pgfpathlineto{\pgfqpoint{1.162546in}{0.693503in}}%
\pgfpathlineto{\pgfqpoint{1.160274in}{0.695485in}}%
\pgfpathlineto{\pgfqpoint{1.148609in}{0.681338in}}%
\pgfpathlineto{\pgfqpoint{1.136941in}{0.667186in}}%
\pgfpathlineto{\pgfqpoint{1.138447in}{0.665854in}}%
\pgfpathlineto{\pgfqpoint{1.139743in}{0.664477in}}%
\pgfpathclose%
\pgfusepath{fill}%
\end{pgfscope}%
\begin{pgfscope}%
\pgfpathrectangle{\pgfqpoint{0.407592in}{0.222904in}}{\pgfqpoint{1.411163in}{1.411163in}}%
\pgfusepath{clip}%
\pgfsetbuttcap%
\pgfsetroundjoin%
\definecolor{currentfill}{rgb}{0.749020,0.772549,0.792157}%
\pgfsetfillcolor{currentfill}%
\pgfsetfillopacity{0.160000}%
\pgfsetlinewidth{0.000000pt}%
\definecolor{currentstroke}{rgb}{0.000000,0.000000,0.000000}%
\pgfsetstrokecolor{currentstroke}%
\pgfsetdash{}{0pt}%
\pgfpathmoveto{\pgfqpoint{1.022788in}{0.695094in}}%
\pgfpathlineto{\pgfqpoint{1.011559in}{0.709173in}}%
\pgfpathlineto{\pgfqpoint{1.000333in}{0.723248in}}%
\pgfpathlineto{\pgfqpoint{0.996821in}{0.720608in}}%
\pgfpathlineto{\pgfqpoint{0.993708in}{0.717881in}}%
\pgfpathlineto{\pgfqpoint{1.005777in}{0.704458in}}%
\pgfpathlineto{\pgfqpoint{1.017845in}{0.691036in}}%
\pgfpathlineto{\pgfqpoint{1.020166in}{0.693097in}}%
\pgfpathlineto{\pgfqpoint{1.022788in}{0.695094in}}%
\pgfpathclose%
\pgfusepath{fill}%
\end{pgfscope}%
\begin{pgfscope}%
\pgfpathrectangle{\pgfqpoint{0.407592in}{0.222904in}}{\pgfqpoint{1.411163in}{1.411163in}}%
\pgfusepath{clip}%
\pgfsetbuttcap%
\pgfsetroundjoin%
\definecolor{currentfill}{rgb}{0.749020,0.772549,0.792157}%
\pgfsetfillcolor{currentfill}%
\pgfsetfillopacity{0.160000}%
\pgfsetlinewidth{0.000000pt}%
\definecolor{currentstroke}{rgb}{0.000000,0.000000,0.000000}%
\pgfsetstrokecolor{currentstroke}%
\pgfsetdash{}{0pt}%
\pgfpathmoveto{\pgfqpoint{1.090230in}{0.610533in}}%
\pgfpathlineto{\pgfqpoint{1.101913in}{0.624703in}}%
\pgfpathlineto{\pgfqpoint{1.113593in}{0.638869in}}%
\pgfpathlineto{\pgfqpoint{1.112742in}{0.639515in}}%
\pgfpathlineto{\pgfqpoint{1.111794in}{0.640134in}}%
\pgfpathlineto{\pgfqpoint{1.101015in}{0.625338in}}%
\pgfpathlineto{\pgfqpoint{1.090230in}{0.610533in}}%
\pgfpathlineto{\pgfqpoint{1.090230in}{0.610533in}}%
\pgfpathlineto{\pgfqpoint{1.090230in}{0.610533in}}%
\pgfpathclose%
\pgfusepath{fill}%
\end{pgfscope}%
\begin{pgfscope}%
\pgfpathrectangle{\pgfqpoint{0.407592in}{0.222904in}}{\pgfqpoint{1.411163in}{1.411163in}}%
\pgfusepath{clip}%
\pgfsetbuttcap%
\pgfsetroundjoin%
\definecolor{currentfill}{rgb}{0.749020,0.772549,0.792157}%
\pgfsetfillcolor{currentfill}%
\pgfsetfillopacity{0.160000}%
\pgfsetlinewidth{0.000000pt}%
\definecolor{currentstroke}{rgb}{0.000000,0.000000,0.000000}%
\pgfsetstrokecolor{currentstroke}%
\pgfsetdash{}{0pt}%
\pgfpathmoveto{\pgfqpoint{1.045256in}{0.666923in}}%
\pgfpathlineto{\pgfqpoint{1.034020in}{0.681010in}}%
\pgfpathlineto{\pgfqpoint{1.022788in}{0.695094in}}%
\pgfpathlineto{\pgfqpoint{1.020166in}{0.693097in}}%
\pgfpathlineto{\pgfqpoint{1.017845in}{0.691036in}}%
\pgfpathlineto{\pgfqpoint{1.029912in}{0.677615in}}%
\pgfpathlineto{\pgfqpoint{1.041978in}{0.664196in}}%
\pgfpathlineto{\pgfqpoint{1.043516in}{0.665581in}}%
\pgfpathlineto{\pgfqpoint{1.045256in}{0.666923in}}%
\pgfpathclose%
\pgfusepath{fill}%
\end{pgfscope}%
\begin{pgfscope}%
\pgfpathrectangle{\pgfqpoint{0.407592in}{0.222904in}}{\pgfqpoint{1.411163in}{1.411163in}}%
\pgfusepath{clip}%
\pgfsetbuttcap%
\pgfsetroundjoin%
\definecolor{currentfill}{rgb}{0.749020,0.772549,0.792157}%
\pgfsetfillcolor{currentfill}%
\pgfsetfillopacity{0.160000}%
\pgfsetlinewidth{0.000000pt}%
\definecolor{currentstroke}{rgb}{0.000000,0.000000,0.000000}%
\pgfsetstrokecolor{currentstroke}%
\pgfsetdash{}{0pt}%
\pgfpathmoveto{\pgfqpoint{1.090230in}{0.610533in}}%
\pgfpathlineto{\pgfqpoint{1.079989in}{0.625277in}}%
\pgfpathlineto{\pgfqpoint{1.069756in}{0.640012in}}%
\pgfpathlineto{\pgfqpoint{1.068699in}{0.639388in}}%
\pgfpathlineto{\pgfqpoint{1.067736in}{0.638736in}}%
\pgfpathlineto{\pgfqpoint{1.078981in}{0.624636in}}%
\pgfpathlineto{\pgfqpoint{1.090230in}{0.610533in}}%
\pgfpathlineto{\pgfqpoint{1.090230in}{0.610533in}}%
\pgfpathlineto{\pgfqpoint{1.090230in}{0.610533in}}%
\pgfpathclose%
\pgfusepath{fill}%
\end{pgfscope}%
\begin{pgfscope}%
\pgfpathrectangle{\pgfqpoint{0.407592in}{0.222904in}}{\pgfqpoint{1.411163in}{1.411163in}}%
\pgfusepath{clip}%
\pgfsetbuttcap%
\pgfsetroundjoin%
\definecolor{currentfill}{rgb}{0.749020,0.772549,0.792157}%
\pgfsetfillcolor{currentfill}%
\pgfsetfillopacity{0.160000}%
\pgfsetlinewidth{0.000000pt}%
\definecolor{currentstroke}{rgb}{0.000000,0.000000,0.000000}%
\pgfsetstrokecolor{currentstroke}%
\pgfsetdash{}{0pt}%
\pgfpathmoveto{\pgfqpoint{1.114985in}{0.637503in}}%
\pgfpathlineto{\pgfqpoint{1.127364in}{0.650990in}}%
\pgfpathlineto{\pgfqpoint{1.139743in}{0.664477in}}%
\pgfpathlineto{\pgfqpoint{1.138447in}{0.665854in}}%
\pgfpathlineto{\pgfqpoint{1.136941in}{0.667186in}}%
\pgfpathlineto{\pgfqpoint{1.125269in}{0.653030in}}%
\pgfpathlineto{\pgfqpoint{1.113593in}{0.638869in}}%
\pgfpathlineto{\pgfqpoint{1.114342in}{0.638197in}}%
\pgfpathlineto{\pgfqpoint{1.114985in}{0.637503in}}%
\pgfpathclose%
\pgfusepath{fill}%
\end{pgfscope}%
\begin{pgfscope}%
\pgfpathrectangle{\pgfqpoint{0.407592in}{0.222904in}}{\pgfqpoint{1.411163in}{1.411163in}}%
\pgfusepath{clip}%
\pgfsetbuttcap%
\pgfsetroundjoin%
\definecolor{currentfill}{rgb}{0.749020,0.772549,0.792157}%
\pgfsetfillcolor{currentfill}%
\pgfsetfillopacity{0.160000}%
\pgfsetlinewidth{0.000000pt}%
\definecolor{currentstroke}{rgb}{0.000000,0.000000,0.000000}%
\pgfsetstrokecolor{currentstroke}%
\pgfsetdash{}{0pt}%
\pgfpathmoveto{\pgfqpoint{1.067736in}{0.638736in}}%
\pgfpathlineto{\pgfqpoint{1.056494in}{0.652831in}}%
\pgfpathlineto{\pgfqpoint{1.045256in}{0.666923in}}%
\pgfpathlineto{\pgfqpoint{1.043516in}{0.665581in}}%
\pgfpathlineto{\pgfqpoint{1.041978in}{0.664196in}}%
\pgfpathlineto{\pgfqpoint{1.054043in}{0.650778in}}%
\pgfpathlineto{\pgfqpoint{1.066106in}{0.637362in}}%
\pgfpathlineto{\pgfqpoint{1.066870in}{0.638060in}}%
\pgfpathlineto{\pgfqpoint{1.067736in}{0.638736in}}%
\pgfpathclose%
\pgfusepath{fill}%
\end{pgfscope}%
\begin{pgfscope}%
\pgfpathrectangle{\pgfqpoint{0.407592in}{0.222904in}}{\pgfqpoint{1.411163in}{1.411163in}}%
\pgfusepath{clip}%
\pgfsetbuttcap%
\pgfsetroundjoin%
\definecolor{currentfill}{rgb}{0.749020,0.772549,0.792157}%
\pgfsetfillcolor{currentfill}%
\pgfsetfillopacity{0.160000}%
\pgfsetlinewidth{0.000000pt}%
\definecolor{currentstroke}{rgb}{0.000000,0.000000,0.000000}%
\pgfsetstrokecolor{currentstroke}%
\pgfsetdash{}{0pt}%
\pgfpathmoveto{\pgfqpoint{1.090230in}{0.610533in}}%
\pgfpathlineto{\pgfqpoint{1.102607in}{0.624018in}}%
\pgfpathlineto{\pgfqpoint{1.114985in}{0.637503in}}%
\pgfpathlineto{\pgfqpoint{1.114342in}{0.638197in}}%
\pgfpathlineto{\pgfqpoint{1.113593in}{0.638869in}}%
\pgfpathlineto{\pgfqpoint{1.101913in}{0.624703in}}%
\pgfpathlineto{\pgfqpoint{1.090230in}{0.610533in}}%
\pgfpathlineto{\pgfqpoint{1.090230in}{0.610533in}}%
\pgfpathlineto{\pgfqpoint{1.090230in}{0.610533in}}%
\pgfpathclose%
\pgfusepath{fill}%
\end{pgfscope}%
\begin{pgfscope}%
\pgfpathrectangle{\pgfqpoint{0.407592in}{0.222904in}}{\pgfqpoint{1.411163in}{1.411163in}}%
\pgfusepath{clip}%
\pgfsetbuttcap%
\pgfsetroundjoin%
\definecolor{currentfill}{rgb}{0.749020,0.772549,0.792157}%
\pgfsetfillcolor{currentfill}%
\pgfsetfillopacity{0.160000}%
\pgfsetlinewidth{0.000000pt}%
\definecolor{currentstroke}{rgb}{0.000000,0.000000,0.000000}%
\pgfsetstrokecolor{currentstroke}%
\pgfsetdash{}{0pt}%
\pgfpathmoveto{\pgfqpoint{1.090230in}{0.610533in}}%
\pgfpathlineto{\pgfqpoint{1.078981in}{0.624636in}}%
\pgfpathlineto{\pgfqpoint{1.067736in}{0.638736in}}%
\pgfpathlineto{\pgfqpoint{1.066870in}{0.638060in}}%
\pgfpathlineto{\pgfqpoint{1.066106in}{0.637362in}}%
\pgfpathlineto{\pgfqpoint{1.078169in}{0.623947in}}%
\pgfpathlineto{\pgfqpoint{1.090230in}{0.610533in}}%
\pgfpathlineto{\pgfqpoint{1.090230in}{0.610533in}}%
\pgfpathlineto{\pgfqpoint{1.090230in}{0.610533in}}%
\pgfpathclose%
\pgfusepath{fill}%
\end{pgfscope}%
\begin{pgfscope}%
\pgfpathrectangle{\pgfqpoint{0.407592in}{0.222904in}}{\pgfqpoint{1.411163in}{1.411163in}}%
\pgfusepath{clip}%
\pgfsetbuttcap%
\pgfsetroundjoin%
\definecolor{currentfill}{rgb}{0.749020,0.772549,0.792157}%
\pgfsetfillcolor{currentfill}%
\pgfsetfillopacity{0.160000}%
\pgfsetlinewidth{0.000000pt}%
\definecolor{currentstroke}{rgb}{0.000000,0.000000,0.000000}%
\pgfsetstrokecolor{currentstroke}%
\pgfsetdash{}{0pt}%
\pgfpathmoveto{\pgfqpoint{1.090230in}{0.610533in}}%
\pgfpathlineto{\pgfqpoint{1.103084in}{0.623294in}}%
\pgfpathlineto{\pgfqpoint{1.115945in}{0.636061in}}%
\pgfpathlineto{\pgfqpoint{1.115520in}{0.636790in}}%
\pgfpathlineto{\pgfqpoint{1.114985in}{0.637503in}}%
\pgfpathlineto{\pgfqpoint{1.102607in}{0.624018in}}%
\pgfpathlineto{\pgfqpoint{1.090230in}{0.610533in}}%
\pgfpathlineto{\pgfqpoint{1.090230in}{0.610533in}}%
\pgfpathlineto{\pgfqpoint{1.090230in}{0.610533in}}%
\pgfpathclose%
\pgfusepath{fill}%
\end{pgfscope}%
\begin{pgfscope}%
\pgfpathrectangle{\pgfqpoint{0.407592in}{0.222904in}}{\pgfqpoint{1.411163in}{1.411163in}}%
\pgfusepath{clip}%
\pgfsetbuttcap%
\pgfsetroundjoin%
\definecolor{currentfill}{rgb}{0.749020,0.772549,0.792157}%
\pgfsetfillcolor{currentfill}%
\pgfsetfillopacity{0.160000}%
\pgfsetlinewidth{0.000000pt}%
\definecolor{currentstroke}{rgb}{0.000000,0.000000,0.000000}%
\pgfsetstrokecolor{currentstroke}%
\pgfsetdash{}{0pt}%
\pgfpathmoveto{\pgfqpoint{1.090230in}{0.610533in}}%
\pgfpathlineto{\pgfqpoint{1.078169in}{0.623947in}}%
\pgfpathlineto{\pgfqpoint{1.066106in}{0.637362in}}%
\pgfpathlineto{\pgfqpoint{1.065447in}{0.636645in}}%
\pgfpathlineto{\pgfqpoint{1.064895in}{0.635913in}}%
\pgfpathlineto{\pgfqpoint{1.077566in}{0.623220in}}%
\pgfpathlineto{\pgfqpoint{1.090230in}{0.610533in}}%
\pgfpathlineto{\pgfqpoint{1.090230in}{0.610533in}}%
\pgfpathlineto{\pgfqpoint{1.090230in}{0.610533in}}%
\pgfpathclose%
\pgfusepath{fill}%
\end{pgfscope}%
\begin{pgfscope}%
\pgfpathrectangle{\pgfqpoint{0.407592in}{0.222904in}}{\pgfqpoint{1.411163in}{1.411163in}}%
\pgfusepath{clip}%
\pgfsetbuttcap%
\pgfsetroundjoin%
\definecolor{currentfill}{rgb}{0.749020,0.772549,0.792157}%
\pgfsetfillcolor{currentfill}%
\pgfsetfillopacity{0.160000}%
\pgfsetlinewidth{0.000000pt}%
\definecolor{currentstroke}{rgb}{0.000000,0.000000,0.000000}%
\pgfsetstrokecolor{currentstroke}%
\pgfsetdash{}{0pt}%
\pgfpathmoveto{\pgfqpoint{1.090230in}{0.610533in}}%
\pgfpathlineto{\pgfqpoint{1.103337in}{0.622545in}}%
\pgfpathlineto{\pgfqpoint{1.116456in}{0.634567in}}%
\pgfpathlineto{\pgfqpoint{1.116257in}{0.635319in}}%
\pgfpathlineto{\pgfqpoint{1.115945in}{0.636061in}}%
\pgfpathlineto{\pgfqpoint{1.103084in}{0.623294in}}%
\pgfpathlineto{\pgfqpoint{1.090230in}{0.610533in}}%
\pgfpathlineto{\pgfqpoint{1.090230in}{0.610533in}}%
\pgfpathlineto{\pgfqpoint{1.090230in}{0.610533in}}%
\pgfpathclose%
\pgfusepath{fill}%
\end{pgfscope}%
\begin{pgfscope}%
\pgfpathrectangle{\pgfqpoint{0.407592in}{0.222904in}}{\pgfqpoint{1.411163in}{1.411163in}}%
\pgfusepath{clip}%
\pgfsetbuttcap%
\pgfsetroundjoin%
\definecolor{currentfill}{rgb}{0.749020,0.772549,0.792157}%
\pgfsetfillcolor{currentfill}%
\pgfsetfillopacity{0.160000}%
\pgfsetlinewidth{0.000000pt}%
\definecolor{currentstroke}{rgb}{0.000000,0.000000,0.000000}%
\pgfsetstrokecolor{currentstroke}%
\pgfsetdash{}{0pt}%
\pgfpathmoveto{\pgfqpoint{1.090230in}{0.610533in}}%
\pgfpathlineto{\pgfqpoint{1.077566in}{0.623220in}}%
\pgfpathlineto{\pgfqpoint{1.064895in}{0.635913in}}%
\pgfpathlineto{\pgfqpoint{1.064454in}{0.635169in}}%
\pgfpathlineto{\pgfqpoint{1.064125in}{0.634416in}}%
\pgfpathlineto{\pgfqpoint{1.077183in}{0.622469in}}%
\pgfpathlineto{\pgfqpoint{1.090230in}{0.610533in}}%
\pgfpathlineto{\pgfqpoint{1.090230in}{0.610533in}}%
\pgfpathlineto{\pgfqpoint{1.090230in}{0.610533in}}%
\pgfpathclose%
\pgfusepath{fill}%
\end{pgfscope}%
\begin{pgfscope}%
\pgfpathrectangle{\pgfqpoint{0.407592in}{0.222904in}}{\pgfqpoint{1.411163in}{1.411163in}}%
\pgfusepath{clip}%
\pgfsetbuttcap%
\pgfsetroundjoin%
\definecolor{currentfill}{rgb}{0.749020,0.772549,0.792157}%
\pgfsetfillcolor{currentfill}%
\pgfsetfillopacity{0.160000}%
\pgfsetlinewidth{0.000000pt}%
\definecolor{currentstroke}{rgb}{0.000000,0.000000,0.000000}%
\pgfsetstrokecolor{currentstroke}%
\pgfsetdash{}{0pt}%
\pgfpathmoveto{\pgfqpoint{1.115945in}{0.636061in}}%
\pgfpathlineto{\pgfqpoint{1.128812in}{0.648834in}}%
\pgfpathlineto{\pgfqpoint{1.141684in}{0.661613in}}%
\pgfpathlineto{\pgfqpoint{1.140824in}{0.663062in}}%
\pgfpathlineto{\pgfqpoint{1.139743in}{0.664477in}}%
\pgfpathlineto{\pgfqpoint{1.127364in}{0.650990in}}%
\pgfpathlineto{\pgfqpoint{1.114985in}{0.637503in}}%
\pgfpathlineto{\pgfqpoint{1.115520in}{0.636790in}}%
\pgfpathlineto{\pgfqpoint{1.115945in}{0.636061in}}%
\pgfpathclose%
\pgfusepath{fill}%
\end{pgfscope}%
\begin{pgfscope}%
\pgfpathrectangle{\pgfqpoint{0.407592in}{0.222904in}}{\pgfqpoint{1.411163in}{1.411163in}}%
\pgfusepath{clip}%
\pgfsetbuttcap%
\pgfsetroundjoin%
\definecolor{currentfill}{rgb}{0.749020,0.772549,0.792157}%
\pgfsetfillcolor{currentfill}%
\pgfsetfillopacity{0.160000}%
\pgfsetlinewidth{0.000000pt}%
\definecolor{currentstroke}{rgb}{0.000000,0.000000,0.000000}%
\pgfsetstrokecolor{currentstroke}%
\pgfsetdash{}{0pt}%
\pgfpathmoveto{\pgfqpoint{1.066106in}{0.637362in}}%
\pgfpathlineto{\pgfqpoint{1.054043in}{0.650778in}}%
\pgfpathlineto{\pgfqpoint{1.041978in}{0.664196in}}%
\pgfpathlineto{\pgfqpoint{1.040649in}{0.662774in}}%
\pgfpathlineto{\pgfqpoint{1.039536in}{0.661320in}}%
\pgfpathlineto{\pgfqpoint{1.052219in}{0.648613in}}%
\pgfpathlineto{\pgfqpoint{1.064895in}{0.635913in}}%
\pgfpathlineto{\pgfqpoint{1.065447in}{0.636645in}}%
\pgfpathlineto{\pgfqpoint{1.066106in}{0.637362in}}%
\pgfpathclose%
\pgfusepath{fill}%
\end{pgfscope}%
\begin{pgfscope}%
\pgfpathrectangle{\pgfqpoint{0.407592in}{0.222904in}}{\pgfqpoint{1.411163in}{1.411163in}}%
\pgfusepath{clip}%
\pgfsetbuttcap%
\pgfsetroundjoin%
\definecolor{currentfill}{rgb}{0.749020,0.772549,0.792157}%
\pgfsetfillcolor{currentfill}%
\pgfsetfillopacity{0.160000}%
\pgfsetlinewidth{0.000000pt}%
\definecolor{currentstroke}{rgb}{0.000000,0.000000,0.000000}%
\pgfsetstrokecolor{currentstroke}%
\pgfsetdash{}{0pt}%
\pgfpathmoveto{\pgfqpoint{1.090230in}{0.610533in}}%
\pgfpathlineto{\pgfqpoint{1.103360in}{0.621783in}}%
\pgfpathlineto{\pgfqpoint{1.116508in}{0.633047in}}%
\pgfpathlineto{\pgfqpoint{1.116539in}{0.633809in}}%
\pgfpathlineto{\pgfqpoint{1.116456in}{0.634567in}}%
\pgfpathlineto{\pgfqpoint{1.103337in}{0.622545in}}%
\pgfpathlineto{\pgfqpoint{1.090230in}{0.610533in}}%
\pgfpathlineto{\pgfqpoint{1.090230in}{0.610533in}}%
\pgfpathlineto{\pgfqpoint{1.090230in}{0.610533in}}%
\pgfpathclose%
\pgfusepath{fill}%
\end{pgfscope}%
\begin{pgfscope}%
\pgfpathrectangle{\pgfqpoint{0.407592in}{0.222904in}}{\pgfqpoint{1.411163in}{1.411163in}}%
\pgfusepath{clip}%
\pgfsetbuttcap%
\pgfsetroundjoin%
\definecolor{currentfill}{rgb}{0.749020,0.772549,0.792157}%
\pgfsetfillcolor{currentfill}%
\pgfsetfillopacity{0.160000}%
\pgfsetlinewidth{0.000000pt}%
\definecolor{currentstroke}{rgb}{0.000000,0.000000,0.000000}%
\pgfsetstrokecolor{currentstroke}%
\pgfsetdash{}{0pt}%
\pgfpathmoveto{\pgfqpoint{1.090230in}{0.610533in}}%
\pgfpathlineto{\pgfqpoint{1.077183in}{0.622469in}}%
\pgfpathlineto{\pgfqpoint{1.064125in}{0.634416in}}%
\pgfpathlineto{\pgfqpoint{1.063910in}{0.633656in}}%
\pgfpathlineto{\pgfqpoint{1.063809in}{0.632895in}}%
\pgfpathlineto{\pgfqpoint{1.077028in}{0.621706in}}%
\pgfpathlineto{\pgfqpoint{1.090230in}{0.610533in}}%
\pgfpathlineto{\pgfqpoint{1.090230in}{0.610533in}}%
\pgfpathlineto{\pgfqpoint{1.090230in}{0.610533in}}%
\pgfpathclose%
\pgfusepath{fill}%
\end{pgfscope}%
\begin{pgfscope}%
\pgfpathrectangle{\pgfqpoint{0.407592in}{0.222904in}}{\pgfqpoint{1.411163in}{1.411163in}}%
\pgfusepath{clip}%
\pgfsetbuttcap%
\pgfsetroundjoin%
\definecolor{currentfill}{rgb}{0.749020,0.772549,0.792157}%
\pgfsetfillcolor{currentfill}%
\pgfsetfillopacity{0.160000}%
\pgfsetlinewidth{0.000000pt}%
\definecolor{currentstroke}{rgb}{0.000000,0.000000,0.000000}%
\pgfsetstrokecolor{currentstroke}%
\pgfsetdash{}{0pt}%
\pgfpathmoveto{\pgfqpoint{1.141684in}{0.661613in}}%
\pgfpathlineto{\pgfqpoint{1.154563in}{0.674399in}}%
\pgfpathlineto{\pgfqpoint{1.167448in}{0.687189in}}%
\pgfpathlineto{\pgfqpoint{1.166141in}{0.689347in}}%
\pgfpathlineto{\pgfqpoint{1.164505in}{0.691454in}}%
\pgfpathlineto{\pgfqpoint{1.152124in}{0.677965in}}%
\pgfpathlineto{\pgfqpoint{1.139743in}{0.664477in}}%
\pgfpathlineto{\pgfqpoint{1.140824in}{0.663062in}}%
\pgfpathlineto{\pgfqpoint{1.141684in}{0.661613in}}%
\pgfpathclose%
\pgfusepath{fill}%
\end{pgfscope}%
\begin{pgfscope}%
\pgfpathrectangle{\pgfqpoint{0.407592in}{0.222904in}}{\pgfqpoint{1.411163in}{1.411163in}}%
\pgfusepath{clip}%
\pgfsetbuttcap%
\pgfsetroundjoin%
\definecolor{currentfill}{rgb}{0.749020,0.772549,0.792157}%
\pgfsetfillcolor{currentfill}%
\pgfsetfillopacity{0.160000}%
\pgfsetlinewidth{0.000000pt}%
\definecolor{currentstroke}{rgb}{0.000000,0.000000,0.000000}%
\pgfsetstrokecolor{currentstroke}%
\pgfsetdash{}{0pt}%
\pgfpathmoveto{\pgfqpoint{1.041978in}{0.664196in}}%
\pgfpathlineto{\pgfqpoint{1.029912in}{0.677615in}}%
\pgfpathlineto{\pgfqpoint{1.017845in}{0.691036in}}%
\pgfpathlineto{\pgfqpoint{1.015837in}{0.688918in}}%
\pgfpathlineto{\pgfqpoint{1.014150in}{0.686752in}}%
\pgfpathlineto{\pgfqpoint{1.026846in}{0.674032in}}%
\pgfpathlineto{\pgfqpoint{1.039536in}{0.661320in}}%
\pgfpathlineto{\pgfqpoint{1.040649in}{0.662774in}}%
\pgfpathlineto{\pgfqpoint{1.041978in}{0.664196in}}%
\pgfpathclose%
\pgfusepath{fill}%
\end{pgfscope}%
\begin{pgfscope}%
\pgfpathrectangle{\pgfqpoint{0.407592in}{0.222904in}}{\pgfqpoint{1.411163in}{1.411163in}}%
\pgfusepath{clip}%
\pgfsetbuttcap%
\pgfsetroundjoin%
\definecolor{currentfill}{rgb}{0.749020,0.772549,0.792157}%
\pgfsetfillcolor{currentfill}%
\pgfsetfillopacity{0.160000}%
\pgfsetlinewidth{0.000000pt}%
\definecolor{currentstroke}{rgb}{0.000000,0.000000,0.000000}%
\pgfsetstrokecolor{currentstroke}%
\pgfsetdash{}{0pt}%
\pgfpathmoveto{\pgfqpoint{1.090230in}{0.610533in}}%
\pgfpathlineto{\pgfqpoint{1.103153in}{0.621021in}}%
\pgfpathlineto{\pgfqpoint{1.116100in}{0.631528in}}%
\pgfpathlineto{\pgfqpoint{1.116361in}{0.632286in}}%
\pgfpathlineto{\pgfqpoint{1.116508in}{0.633047in}}%
\pgfpathlineto{\pgfqpoint{1.103360in}{0.621783in}}%
\pgfpathlineto{\pgfqpoint{1.090230in}{0.610533in}}%
\pgfpathlineto{\pgfqpoint{1.090230in}{0.610533in}}%
\pgfpathlineto{\pgfqpoint{1.090230in}{0.610533in}}%
\pgfpathclose%
\pgfusepath{fill}%
\end{pgfscope}%
\begin{pgfscope}%
\pgfpathrectangle{\pgfqpoint{0.407592in}{0.222904in}}{\pgfqpoint{1.411163in}{1.411163in}}%
\pgfusepath{clip}%
\pgfsetbuttcap%
\pgfsetroundjoin%
\definecolor{currentfill}{rgb}{0.749020,0.772549,0.792157}%
\pgfsetfillcolor{currentfill}%
\pgfsetfillopacity{0.160000}%
\pgfsetlinewidth{0.000000pt}%
\definecolor{currentstroke}{rgb}{0.000000,0.000000,0.000000}%
\pgfsetstrokecolor{currentstroke}%
\pgfsetdash{}{0pt}%
\pgfpathmoveto{\pgfqpoint{1.090230in}{0.610533in}}%
\pgfpathlineto{\pgfqpoint{1.077028in}{0.621706in}}%
\pgfpathlineto{\pgfqpoint{1.063809in}{0.632895in}}%
\pgfpathlineto{\pgfqpoint{1.063824in}{0.632134in}}%
\pgfpathlineto{\pgfqpoint{1.063954in}{0.631377in}}%
\pgfpathlineto{\pgfqpoint{1.077104in}{0.620946in}}%
\pgfpathlineto{\pgfqpoint{1.090230in}{0.610533in}}%
\pgfpathlineto{\pgfqpoint{1.090230in}{0.610533in}}%
\pgfpathlineto{\pgfqpoint{1.090230in}{0.610533in}}%
\pgfpathclose%
\pgfusepath{fill}%
\end{pgfscope}%
\begin{pgfscope}%
\pgfpathrectangle{\pgfqpoint{0.407592in}{0.222904in}}{\pgfqpoint{1.411163in}{1.411163in}}%
\pgfusepath{clip}%
\pgfsetbuttcap%
\pgfsetroundjoin%
\definecolor{currentfill}{rgb}{0.749020,0.772549,0.792157}%
\pgfsetfillcolor{currentfill}%
\pgfsetfillopacity{0.160000}%
\pgfsetlinewidth{0.000000pt}%
\definecolor{currentstroke}{rgb}{0.000000,0.000000,0.000000}%
\pgfsetstrokecolor{currentstroke}%
\pgfsetdash{}{0pt}%
\pgfpathmoveto{\pgfqpoint{1.167448in}{0.687189in}}%
\pgfpathlineto{\pgfqpoint{1.180339in}{0.699986in}}%
\pgfpathlineto{\pgfqpoint{1.193235in}{0.712789in}}%
\pgfpathlineto{\pgfqpoint{1.191471in}{0.715645in}}%
\pgfpathlineto{\pgfqpoint{1.189269in}{0.718434in}}%
\pgfpathlineto{\pgfqpoint{1.176886in}{0.704944in}}%
\pgfpathlineto{\pgfqpoint{1.164505in}{0.691454in}}%
\pgfpathlineto{\pgfqpoint{1.166141in}{0.689347in}}%
\pgfpathlineto{\pgfqpoint{1.167448in}{0.687189in}}%
\pgfpathclose%
\pgfusepath{fill}%
\end{pgfscope}%
\begin{pgfscope}%
\pgfpathrectangle{\pgfqpoint{0.407592in}{0.222904in}}{\pgfqpoint{1.411163in}{1.411163in}}%
\pgfusepath{clip}%
\pgfsetbuttcap%
\pgfsetroundjoin%
\definecolor{currentfill}{rgb}{0.749020,0.772549,0.792157}%
\pgfsetfillcolor{currentfill}%
\pgfsetfillopacity{0.160000}%
\pgfsetlinewidth{0.000000pt}%
\definecolor{currentstroke}{rgb}{0.000000,0.000000,0.000000}%
\pgfsetstrokecolor{currentstroke}%
\pgfsetdash{}{0pt}%
\pgfpathmoveto{\pgfqpoint{1.090230in}{0.610533in}}%
\pgfpathlineto{\pgfqpoint{1.102720in}{0.620274in}}%
\pgfpathlineto{\pgfqpoint{1.115238in}{0.630036in}}%
\pgfpathlineto{\pgfqpoint{1.115725in}{0.630777in}}%
\pgfpathlineto{\pgfqpoint{1.116100in}{0.631528in}}%
\pgfpathlineto{\pgfqpoint{1.103153in}{0.621021in}}%
\pgfpathlineto{\pgfqpoint{1.090230in}{0.610533in}}%
\pgfpathlineto{\pgfqpoint{1.090230in}{0.610533in}}%
\pgfpathlineto{\pgfqpoint{1.090230in}{0.610533in}}%
\pgfpathclose%
\pgfusepath{fill}%
\end{pgfscope}%
\begin{pgfscope}%
\pgfpathrectangle{\pgfqpoint{0.407592in}{0.222904in}}{\pgfqpoint{1.411163in}{1.411163in}}%
\pgfusepath{clip}%
\pgfsetbuttcap%
\pgfsetroundjoin%
\definecolor{currentfill}{rgb}{0.749020,0.772549,0.792157}%
\pgfsetfillcolor{currentfill}%
\pgfsetfillopacity{0.160000}%
\pgfsetlinewidth{0.000000pt}%
\definecolor{currentstroke}{rgb}{0.000000,0.000000,0.000000}%
\pgfsetstrokecolor{currentstroke}%
\pgfsetdash{}{0pt}%
\pgfpathmoveto{\pgfqpoint{1.090230in}{0.610533in}}%
\pgfpathlineto{\pgfqpoint{1.077104in}{0.620946in}}%
\pgfpathlineto{\pgfqpoint{1.063954in}{0.631377in}}%
\pgfpathlineto{\pgfqpoint{1.064199in}{0.630628in}}%
\pgfpathlineto{\pgfqpoint{1.064558in}{0.629889in}}%
\pgfpathlineto{\pgfqpoint{1.077408in}{0.620200in}}%
\pgfpathlineto{\pgfqpoint{1.090230in}{0.610533in}}%
\pgfpathlineto{\pgfqpoint{1.090230in}{0.610533in}}%
\pgfpathlineto{\pgfqpoint{1.090230in}{0.610533in}}%
\pgfpathclose%
\pgfusepath{fill}%
\end{pgfscope}%
\begin{pgfscope}%
\pgfpathrectangle{\pgfqpoint{0.407592in}{0.222904in}}{\pgfqpoint{1.411163in}{1.411163in}}%
\pgfusepath{clip}%
\pgfsetbuttcap%
\pgfsetroundjoin%
\definecolor{currentfill}{rgb}{0.749020,0.772549,0.792157}%
\pgfsetfillcolor{currentfill}%
\pgfsetfillopacity{0.160000}%
\pgfsetlinewidth{0.000000pt}%
\definecolor{currentstroke}{rgb}{0.000000,0.000000,0.000000}%
\pgfsetstrokecolor{currentstroke}%
\pgfsetdash{}{0pt}%
\pgfpathmoveto{\pgfqpoint{1.116456in}{0.634567in}}%
\pgfpathlineto{\pgfqpoint{1.129586in}{0.646600in}}%
\pgfpathlineto{\pgfqpoint{1.142729in}{0.658644in}}%
\pgfpathlineto{\pgfqpoint{1.142320in}{0.660139in}}%
\pgfpathlineto{\pgfqpoint{1.141684in}{0.661613in}}%
\pgfpathlineto{\pgfqpoint{1.128812in}{0.648834in}}%
\pgfpathlineto{\pgfqpoint{1.115945in}{0.636061in}}%
\pgfpathlineto{\pgfqpoint{1.116257in}{0.635319in}}%
\pgfpathlineto{\pgfqpoint{1.116456in}{0.634567in}}%
\pgfpathclose%
\pgfusepath{fill}%
\end{pgfscope}%
\begin{pgfscope}%
\pgfpathrectangle{\pgfqpoint{0.407592in}{0.222904in}}{\pgfqpoint{1.411163in}{1.411163in}}%
\pgfusepath{clip}%
\pgfsetbuttcap%
\pgfsetroundjoin%
\definecolor{currentfill}{rgb}{0.749020,0.772549,0.792157}%
\pgfsetfillcolor{currentfill}%
\pgfsetfillopacity{0.160000}%
\pgfsetlinewidth{0.000000pt}%
\definecolor{currentstroke}{rgb}{0.000000,0.000000,0.000000}%
\pgfsetstrokecolor{currentstroke}%
\pgfsetdash{}{0pt}%
\pgfpathmoveto{\pgfqpoint{1.017845in}{0.691036in}}%
\pgfpathlineto{\pgfqpoint{1.005777in}{0.704458in}}%
\pgfpathlineto{\pgfqpoint{0.993708in}{0.717881in}}%
\pgfpathlineto{\pgfqpoint{0.991010in}{0.715077in}}%
\pgfpathlineto{\pgfqpoint{0.988739in}{0.712209in}}%
\pgfpathlineto{\pgfqpoint{1.001448in}{0.699477in}}%
\pgfpathlineto{\pgfqpoint{1.014150in}{0.686752in}}%
\pgfpathlineto{\pgfqpoint{1.015837in}{0.688918in}}%
\pgfpathlineto{\pgfqpoint{1.017845in}{0.691036in}}%
\pgfpathclose%
\pgfusepath{fill}%
\end{pgfscope}%
\begin{pgfscope}%
\pgfpathrectangle{\pgfqpoint{0.407592in}{0.222904in}}{\pgfqpoint{1.411163in}{1.411163in}}%
\pgfusepath{clip}%
\pgfsetbuttcap%
\pgfsetroundjoin%
\definecolor{currentfill}{rgb}{0.749020,0.772549,0.792157}%
\pgfsetfillcolor{currentfill}%
\pgfsetfillopacity{0.160000}%
\pgfsetlinewidth{0.000000pt}%
\definecolor{currentstroke}{rgb}{0.000000,0.000000,0.000000}%
\pgfsetstrokecolor{currentstroke}%
\pgfsetdash{}{0pt}%
\pgfpathmoveto{\pgfqpoint{1.064895in}{0.635913in}}%
\pgfpathlineto{\pgfqpoint{1.052219in}{0.648613in}}%
\pgfpathlineto{\pgfqpoint{1.039536in}{0.661320in}}%
\pgfpathlineto{\pgfqpoint{1.038642in}{0.659840in}}%
\pgfpathlineto{\pgfqpoint{1.037972in}{0.658343in}}%
\pgfpathlineto{\pgfqpoint{1.051055in}{0.646373in}}%
\pgfpathlineto{\pgfqpoint{1.064125in}{0.634416in}}%
\pgfpathlineto{\pgfqpoint{1.064454in}{0.635169in}}%
\pgfpathlineto{\pgfqpoint{1.064895in}{0.635913in}}%
\pgfpathclose%
\pgfusepath{fill}%
\end{pgfscope}%
\begin{pgfscope}%
\pgfpathrectangle{\pgfqpoint{0.407592in}{0.222904in}}{\pgfqpoint{1.411163in}{1.411163in}}%
\pgfusepath{clip}%
\pgfsetbuttcap%
\pgfsetroundjoin%
\definecolor{currentfill}{rgb}{0.749020,0.772549,0.792157}%
\pgfsetfillcolor{currentfill}%
\pgfsetfillopacity{0.160000}%
\pgfsetlinewidth{0.000000pt}%
\definecolor{currentstroke}{rgb}{0.000000,0.000000,0.000000}%
\pgfsetstrokecolor{currentstroke}%
\pgfsetdash{}{0pt}%
\pgfpathmoveto{\pgfqpoint{1.193235in}{0.712789in}}%
\pgfpathlineto{\pgfqpoint{1.206138in}{0.725598in}}%
\pgfpathlineto{\pgfqpoint{1.219047in}{0.738413in}}%
\pgfpathlineto{\pgfqpoint{1.216815in}{0.741958in}}%
\pgfpathlineto{\pgfqpoint{1.214036in}{0.745418in}}%
\pgfpathlineto{\pgfqpoint{1.201652in}{0.731926in}}%
\pgfpathlineto{\pgfqpoint{1.189269in}{0.718434in}}%
\pgfpathlineto{\pgfqpoint{1.191471in}{0.715645in}}%
\pgfpathlineto{\pgfqpoint{1.193235in}{0.712789in}}%
\pgfpathclose%
\pgfusepath{fill}%
\end{pgfscope}%
\begin{pgfscope}%
\pgfpathrectangle{\pgfqpoint{0.407592in}{0.222904in}}{\pgfqpoint{1.411163in}{1.411163in}}%
\pgfusepath{clip}%
\pgfsetbuttcap%
\pgfsetroundjoin%
\definecolor{currentfill}{rgb}{0.749020,0.772549,0.792157}%
\pgfsetfillcolor{currentfill}%
\pgfsetfillopacity{0.160000}%
\pgfsetlinewidth{0.000000pt}%
\definecolor{currentstroke}{rgb}{0.000000,0.000000,0.000000}%
\pgfsetstrokecolor{currentstroke}%
\pgfsetdash{}{0pt}%
\pgfpathmoveto{\pgfqpoint{1.090230in}{0.610533in}}%
\pgfpathlineto{\pgfqpoint{1.102068in}{0.619554in}}%
\pgfpathlineto{\pgfqpoint{1.113937in}{0.628598in}}%
\pgfpathlineto{\pgfqpoint{1.114641in}{0.629309in}}%
\pgfpathlineto{\pgfqpoint{1.115238in}{0.630036in}}%
\pgfpathlineto{\pgfqpoint{1.102720in}{0.620274in}}%
\pgfpathlineto{\pgfqpoint{1.090230in}{0.610533in}}%
\pgfpathlineto{\pgfqpoint{1.090230in}{0.610533in}}%
\pgfpathlineto{\pgfqpoint{1.090230in}{0.610533in}}%
\pgfpathclose%
\pgfusepath{fill}%
\end{pgfscope}%
\begin{pgfscope}%
\pgfpathrectangle{\pgfqpoint{0.407592in}{0.222904in}}{\pgfqpoint{1.411163in}{1.411163in}}%
\pgfusepath{clip}%
\pgfsetbuttcap%
\pgfsetroundjoin%
\definecolor{currentfill}{rgb}{0.749020,0.772549,0.792157}%
\pgfsetfillcolor{currentfill}%
\pgfsetfillopacity{0.160000}%
\pgfsetlinewidth{0.000000pt}%
\definecolor{currentstroke}{rgb}{0.000000,0.000000,0.000000}%
\pgfsetstrokecolor{currentstroke}%
\pgfsetdash{}{0pt}%
\pgfpathmoveto{\pgfqpoint{1.090230in}{0.610533in}}%
\pgfpathlineto{\pgfqpoint{1.077408in}{0.620200in}}%
\pgfpathlineto{\pgfqpoint{1.064558in}{0.629889in}}%
\pgfpathlineto{\pgfqpoint{1.065029in}{0.629165in}}%
\pgfpathlineto{\pgfqpoint{1.065610in}{0.628457in}}%
\pgfpathlineto{\pgfqpoint{1.077936in}{0.619484in}}%
\pgfpathlineto{\pgfqpoint{1.090230in}{0.610533in}}%
\pgfpathlineto{\pgfqpoint{1.090230in}{0.610533in}}%
\pgfpathlineto{\pgfqpoint{1.090230in}{0.610533in}}%
\pgfpathclose%
\pgfusepath{fill}%
\end{pgfscope}%
\begin{pgfscope}%
\pgfpathrectangle{\pgfqpoint{0.407592in}{0.222904in}}{\pgfqpoint{1.411163in}{1.411163in}}%
\pgfusepath{clip}%
\pgfsetbuttcap%
\pgfsetroundjoin%
\definecolor{currentfill}{rgb}{0.749020,0.772549,0.792157}%
\pgfsetfillcolor{currentfill}%
\pgfsetfillopacity{0.160000}%
\pgfsetlinewidth{0.000000pt}%
\definecolor{currentstroke}{rgb}{0.000000,0.000000,0.000000}%
\pgfsetstrokecolor{currentstroke}%
\pgfsetdash{}{0pt}%
\pgfpathmoveto{\pgfqpoint{0.993708in}{0.717881in}}%
\pgfpathlineto{\pgfqpoint{0.981637in}{0.731305in}}%
\pgfpathlineto{\pgfqpoint{0.969565in}{0.744731in}}%
\pgfpathlineto{\pgfqpoint{0.966167in}{0.741253in}}%
\pgfpathlineto{\pgfqpoint{0.963301in}{0.737693in}}%
\pgfpathlineto{\pgfqpoint{0.976023in}{0.724948in}}%
\pgfpathlineto{\pgfqpoint{0.988739in}{0.712209in}}%
\pgfpathlineto{\pgfqpoint{0.991010in}{0.715077in}}%
\pgfpathlineto{\pgfqpoint{0.993708in}{0.717881in}}%
\pgfpathclose%
\pgfusepath{fill}%
\end{pgfscope}%
\begin{pgfscope}%
\pgfpathrectangle{\pgfqpoint{0.407592in}{0.222904in}}{\pgfqpoint{1.411163in}{1.411163in}}%
\pgfusepath{clip}%
\pgfsetbuttcap%
\pgfsetroundjoin%
\definecolor{currentfill}{rgb}{0.749020,0.772549,0.792157}%
\pgfsetfillcolor{currentfill}%
\pgfsetfillopacity{0.160000}%
\pgfsetlinewidth{0.000000pt}%
\definecolor{currentstroke}{rgb}{0.000000,0.000000,0.000000}%
\pgfsetstrokecolor{currentstroke}%
\pgfsetdash{}{0pt}%
\pgfpathmoveto{\pgfqpoint{1.090230in}{0.610533in}}%
\pgfpathlineto{\pgfqpoint{1.101208in}{0.618873in}}%
\pgfpathlineto{\pgfqpoint{1.112219in}{0.627238in}}%
\pgfpathlineto{\pgfqpoint{1.113128in}{0.627907in}}%
\pgfpathlineto{\pgfqpoint{1.113937in}{0.628598in}}%
\pgfpathlineto{\pgfqpoint{1.102068in}{0.619554in}}%
\pgfpathlineto{\pgfqpoint{1.090230in}{0.610533in}}%
\pgfpathlineto{\pgfqpoint{1.090230in}{0.610533in}}%
\pgfpathlineto{\pgfqpoint{1.090230in}{0.610533in}}%
\pgfpathclose%
\pgfusepath{fill}%
\end{pgfscope}%
\begin{pgfscope}%
\pgfpathrectangle{\pgfqpoint{0.407592in}{0.222904in}}{\pgfqpoint{1.411163in}{1.411163in}}%
\pgfusepath{clip}%
\pgfsetbuttcap%
\pgfsetroundjoin%
\definecolor{currentfill}{rgb}{0.749020,0.772549,0.792157}%
\pgfsetfillcolor{currentfill}%
\pgfsetfillopacity{0.160000}%
\pgfsetlinewidth{0.000000pt}%
\definecolor{currentstroke}{rgb}{0.000000,0.000000,0.000000}%
\pgfsetstrokecolor{currentstroke}%
\pgfsetdash{}{0pt}%
\pgfpathmoveto{\pgfqpoint{1.090230in}{0.610533in}}%
\pgfpathlineto{\pgfqpoint{1.077936in}{0.619484in}}%
\pgfpathlineto{\pgfqpoint{1.065610in}{0.628457in}}%
\pgfpathlineto{\pgfqpoint{1.066299in}{0.627771in}}%
\pgfpathlineto{\pgfqpoint{1.067093in}{0.627107in}}%
\pgfpathlineto{\pgfqpoint{1.078678in}{0.618808in}}%
\pgfpathlineto{\pgfqpoint{1.090230in}{0.610533in}}%
\pgfpathlineto{\pgfqpoint{1.090230in}{0.610533in}}%
\pgfpathlineto{\pgfqpoint{1.090230in}{0.610533in}}%
\pgfpathclose%
\pgfusepath{fill}%
\end{pgfscope}%
\begin{pgfscope}%
\pgfpathrectangle{\pgfqpoint{0.407592in}{0.222904in}}{\pgfqpoint{1.411163in}{1.411163in}}%
\pgfusepath{clip}%
\pgfsetbuttcap%
\pgfsetroundjoin%
\definecolor{currentfill}{rgb}{0.749020,0.772549,0.792157}%
\pgfsetfillcolor{currentfill}%
\pgfsetfillopacity{0.160000}%
\pgfsetlinewidth{0.000000pt}%
\definecolor{currentstroke}{rgb}{0.000000,0.000000,0.000000}%
\pgfsetstrokecolor{currentstroke}%
\pgfsetdash{}{0pt}%
\pgfpathmoveto{\pgfqpoint{1.219047in}{0.738413in}}%
\pgfpathlineto{\pgfqpoint{1.231962in}{0.751234in}}%
\pgfpathlineto{\pgfqpoint{1.244882in}{0.764061in}}%
\pgfpathlineto{\pgfqpoint{1.242171in}{0.768284in}}%
\pgfpathlineto{\pgfqpoint{1.238806in}{0.772404in}}%
\pgfpathlineto{\pgfqpoint{1.226421in}{0.758911in}}%
\pgfpathlineto{\pgfqpoint{1.214036in}{0.745418in}}%
\pgfpathlineto{\pgfqpoint{1.216815in}{0.741958in}}%
\pgfpathlineto{\pgfqpoint{1.219047in}{0.738413in}}%
\pgfpathclose%
\pgfusepath{fill}%
\end{pgfscope}%
\begin{pgfscope}%
\pgfpathrectangle{\pgfqpoint{0.407592in}{0.222904in}}{\pgfqpoint{1.411163in}{1.411163in}}%
\pgfusepath{clip}%
\pgfsetbuttcap%
\pgfsetroundjoin%
\definecolor{currentfill}{rgb}{0.749020,0.772549,0.792157}%
\pgfsetfillcolor{currentfill}%
\pgfsetfillopacity{0.160000}%
\pgfsetlinewidth{0.000000pt}%
\definecolor{currentstroke}{rgb}{0.000000,0.000000,0.000000}%
\pgfsetstrokecolor{currentstroke}%
\pgfsetdash{}{0pt}%
\pgfpathmoveto{\pgfqpoint{1.090230in}{0.610533in}}%
\pgfpathlineto{\pgfqpoint{1.100705in}{0.618552in}}%
\pgfpathlineto{\pgfqpoint{1.111212in}{0.626595in}}%
\pgfpathlineto{\pgfqpoint{1.112219in}{0.627238in}}%
\pgfpathlineto{\pgfqpoint{1.101208in}{0.618873in}}%
\pgfpathlineto{\pgfqpoint{1.090230in}{0.610533in}}%
\pgfusepath{fill}%
\end{pgfscope}%
\begin{pgfscope}%
\pgfpathrectangle{\pgfqpoint{0.407592in}{0.222904in}}{\pgfqpoint{1.411163in}{1.411163in}}%
\pgfusepath{clip}%
\pgfsetbuttcap%
\pgfsetroundjoin%
\definecolor{currentfill}{rgb}{0.749020,0.772549,0.792157}%
\pgfsetfillcolor{currentfill}%
\pgfsetfillopacity{0.160000}%
\pgfsetlinewidth{0.000000pt}%
\definecolor{currentstroke}{rgb}{0.000000,0.000000,0.000000}%
\pgfsetstrokecolor{currentstroke}%
\pgfsetdash{}{0pt}%
\pgfpathmoveto{\pgfqpoint{1.090230in}{0.610533in}}%
\pgfpathlineto{\pgfqpoint{1.078678in}{0.618808in}}%
\pgfpathlineto{\pgfqpoint{1.067093in}{0.627107in}}%
\pgfpathlineto{\pgfqpoint{1.067988in}{0.626470in}}%
\pgfpathlineto{\pgfqpoint{1.068981in}{0.625862in}}%
\pgfpathlineto{\pgfqpoint{1.079623in}{0.618184in}}%
\pgfpathlineto{\pgfqpoint{1.090230in}{0.610533in}}%
\pgfpathlineto{\pgfqpoint{1.090230in}{0.610533in}}%
\pgfpathlineto{\pgfqpoint{1.090230in}{0.610533in}}%
\pgfpathclose%
\pgfusepath{fill}%
\end{pgfscope}%
\begin{pgfscope}%
\pgfpathrectangle{\pgfqpoint{0.407592in}{0.222904in}}{\pgfqpoint{1.411163in}{1.411163in}}%
\pgfusepath{clip}%
\pgfsetbuttcap%
\pgfsetroundjoin%
\definecolor{currentfill}{rgb}{0.749020,0.772549,0.792157}%
\pgfsetfillcolor{currentfill}%
\pgfsetfillopacity{0.160000}%
\pgfsetlinewidth{0.000000pt}%
\definecolor{currentstroke}{rgb}{0.000000,0.000000,0.000000}%
\pgfsetstrokecolor{currentstroke}%
\pgfsetdash{}{0pt}%
\pgfpathmoveto{\pgfqpoint{1.090230in}{0.610533in}}%
\pgfpathlineto{\pgfqpoint{1.099562in}{0.617953in}}%
\pgfpathlineto{\pgfqpoint{1.108927in}{0.625398in}}%
\pgfpathlineto{\pgfqpoint{1.110114in}{0.625981in}}%
\pgfpathlineto{\pgfqpoint{1.111212in}{0.626595in}}%
\pgfpathlineto{\pgfqpoint{1.100705in}{0.618552in}}%
\pgfpathlineto{\pgfqpoint{1.090230in}{0.610533in}}%
\pgfpathlineto{\pgfqpoint{1.090230in}{0.610533in}}%
\pgfpathlineto{\pgfqpoint{1.090230in}{0.610533in}}%
\pgfpathclose%
\pgfusepath{fill}%
\end{pgfscope}%
\begin{pgfscope}%
\pgfpathrectangle{\pgfqpoint{0.407592in}{0.222904in}}{\pgfqpoint{1.411163in}{1.411163in}}%
\pgfusepath{clip}%
\pgfsetbuttcap%
\pgfsetroundjoin%
\definecolor{currentfill}{rgb}{0.749020,0.772549,0.792157}%
\pgfsetfillcolor{currentfill}%
\pgfsetfillopacity{0.160000}%
\pgfsetlinewidth{0.000000pt}%
\definecolor{currentstroke}{rgb}{0.000000,0.000000,0.000000}%
\pgfsetstrokecolor{currentstroke}%
\pgfsetdash{}{0pt}%
\pgfpathmoveto{\pgfqpoint{1.116508in}{0.633047in}}%
\pgfpathlineto{\pgfqpoint{1.129673in}{0.644327in}}%
\pgfpathlineto{\pgfqpoint{1.142855in}{0.655621in}}%
\pgfpathlineto{\pgfqpoint{1.142907in}{0.657136in}}%
\pgfpathlineto{\pgfqpoint{1.142729in}{0.658644in}}%
\pgfpathlineto{\pgfqpoint{1.129586in}{0.646600in}}%
\pgfpathlineto{\pgfqpoint{1.116456in}{0.634567in}}%
\pgfpathlineto{\pgfqpoint{1.116539in}{0.633809in}}%
\pgfpathlineto{\pgfqpoint{1.116508in}{0.633047in}}%
\pgfpathclose%
\pgfusepath{fill}%
\end{pgfscope}%
\begin{pgfscope}%
\pgfpathrectangle{\pgfqpoint{0.407592in}{0.222904in}}{\pgfqpoint{1.411163in}{1.411163in}}%
\pgfusepath{clip}%
\pgfsetbuttcap%
\pgfsetroundjoin%
\definecolor{currentfill}{rgb}{0.749020,0.772549,0.792157}%
\pgfsetfillcolor{currentfill}%
\pgfsetfillopacity{0.160000}%
\pgfsetlinewidth{0.000000pt}%
\definecolor{currentstroke}{rgb}{0.000000,0.000000,0.000000}%
\pgfsetstrokecolor{currentstroke}%
\pgfsetdash{}{0pt}%
\pgfpathmoveto{\pgfqpoint{0.969565in}{0.744731in}}%
\pgfpathlineto{\pgfqpoint{0.957492in}{0.758158in}}%
\pgfpathlineto{\pgfqpoint{0.945418in}{0.771587in}}%
\pgfpathlineto{\pgfqpoint{0.941310in}{0.767444in}}%
\pgfpathlineto{\pgfqpoint{0.937838in}{0.763203in}}%
\pgfpathlineto{\pgfqpoint{0.950573in}{0.750445in}}%
\pgfpathlineto{\pgfqpoint{0.963301in}{0.737693in}}%
\pgfpathlineto{\pgfqpoint{0.966167in}{0.741253in}}%
\pgfpathlineto{\pgfqpoint{0.969565in}{0.744731in}}%
\pgfpathclose%
\pgfusepath{fill}%
\end{pgfscope}%
\begin{pgfscope}%
\pgfpathrectangle{\pgfqpoint{0.407592in}{0.222904in}}{\pgfqpoint{1.411163in}{1.411163in}}%
\pgfusepath{clip}%
\pgfsetbuttcap%
\pgfsetroundjoin%
\definecolor{currentfill}{rgb}{0.749020,0.772549,0.792157}%
\pgfsetfillcolor{currentfill}%
\pgfsetfillopacity{0.160000}%
\pgfsetlinewidth{0.000000pt}%
\definecolor{currentstroke}{rgb}{0.000000,0.000000,0.000000}%
\pgfsetstrokecolor{currentstroke}%
\pgfsetdash{}{0pt}%
\pgfpathmoveto{\pgfqpoint{1.142729in}{0.658644in}}%
\pgfpathlineto{\pgfqpoint{1.155882in}{0.670699in}}%
\pgfpathlineto{\pgfqpoint{1.169048in}{0.682764in}}%
\pgfpathlineto{\pgfqpoint{1.168418in}{0.684992in}}%
\pgfpathlineto{\pgfqpoint{1.167448in}{0.687189in}}%
\pgfpathlineto{\pgfqpoint{1.154563in}{0.674399in}}%
\pgfpathlineto{\pgfqpoint{1.141684in}{0.661613in}}%
\pgfpathlineto{\pgfqpoint{1.142320in}{0.660139in}}%
\pgfpathlineto{\pgfqpoint{1.142729in}{0.658644in}}%
\pgfpathclose%
\pgfusepath{fill}%
\end{pgfscope}%
\begin{pgfscope}%
\pgfpathrectangle{\pgfqpoint{0.407592in}{0.222904in}}{\pgfqpoint{1.411163in}{1.411163in}}%
\pgfusepath{clip}%
\pgfsetbuttcap%
\pgfsetroundjoin%
\definecolor{currentfill}{rgb}{0.749020,0.772549,0.792157}%
\pgfsetfillcolor{currentfill}%
\pgfsetfillopacity{0.160000}%
\pgfsetlinewidth{0.000000pt}%
\definecolor{currentstroke}{rgb}{0.000000,0.000000,0.000000}%
\pgfsetstrokecolor{currentstroke}%
\pgfsetdash{}{0pt}%
\pgfpathmoveto{\pgfqpoint{1.244882in}{0.764061in}}%
\pgfpathlineto{\pgfqpoint{1.257809in}{0.776894in}}%
\pgfpathlineto{\pgfqpoint{1.270742in}{0.789733in}}%
\pgfpathlineto{\pgfqpoint{1.267541in}{0.794623in}}%
\pgfpathlineto{\pgfqpoint{1.263579in}{0.799394in}}%
\pgfpathlineto{\pgfqpoint{1.251192in}{0.785899in}}%
\pgfpathlineto{\pgfqpoint{1.238806in}{0.772404in}}%
\pgfpathlineto{\pgfqpoint{1.242171in}{0.768284in}}%
\pgfpathlineto{\pgfqpoint{1.244882in}{0.764061in}}%
\pgfpathclose%
\pgfusepath{fill}%
\end{pgfscope}%
\begin{pgfscope}%
\pgfpathrectangle{\pgfqpoint{0.407592in}{0.222904in}}{\pgfqpoint{1.411163in}{1.411163in}}%
\pgfusepath{clip}%
\pgfsetbuttcap%
\pgfsetroundjoin%
\definecolor{currentfill}{rgb}{0.749020,0.772549,0.792157}%
\pgfsetfillcolor{currentfill}%
\pgfsetfillopacity{0.160000}%
\pgfsetlinewidth{0.000000pt}%
\definecolor{currentstroke}{rgb}{0.000000,0.000000,0.000000}%
\pgfsetstrokecolor{currentstroke}%
\pgfsetdash{}{0pt}%
\pgfpathmoveto{\pgfqpoint{1.064125in}{0.634416in}}%
\pgfpathlineto{\pgfqpoint{1.051055in}{0.646373in}}%
\pgfpathlineto{\pgfqpoint{1.037972in}{0.658343in}}%
\pgfpathlineto{\pgfqpoint{1.037530in}{0.656833in}}%
\pgfpathlineto{\pgfqpoint{1.037317in}{0.655317in}}%
\pgfpathlineto{\pgfqpoint{1.050572in}{0.644098in}}%
\pgfpathlineto{\pgfqpoint{1.063809in}{0.632895in}}%
\pgfpathlineto{\pgfqpoint{1.063910in}{0.633656in}}%
\pgfpathlineto{\pgfqpoint{1.064125in}{0.634416in}}%
\pgfpathclose%
\pgfusepath{fill}%
\end{pgfscope}%
\begin{pgfscope}%
\pgfpathrectangle{\pgfqpoint{0.407592in}{0.222904in}}{\pgfqpoint{1.411163in}{1.411163in}}%
\pgfusepath{clip}%
\pgfsetbuttcap%
\pgfsetroundjoin%
\definecolor{currentfill}{rgb}{0.749020,0.772549,0.792157}%
\pgfsetfillcolor{currentfill}%
\pgfsetfillopacity{0.160000}%
\pgfsetlinewidth{0.000000pt}%
\definecolor{currentstroke}{rgb}{0.000000,0.000000,0.000000}%
\pgfsetstrokecolor{currentstroke}%
\pgfsetdash{}{0pt}%
\pgfpathmoveto{\pgfqpoint{1.090230in}{0.610533in}}%
\pgfpathlineto{\pgfqpoint{1.079623in}{0.618184in}}%
\pgfpathlineto{\pgfqpoint{1.068981in}{0.625862in}}%
\pgfpathlineto{\pgfqpoint{1.070067in}{0.625285in}}%
\pgfpathlineto{\pgfqpoint{1.071242in}{0.624743in}}%
\pgfpathlineto{\pgfqpoint{1.080753in}{0.617625in}}%
\pgfpathlineto{\pgfqpoint{1.090230in}{0.610533in}}%
\pgfpathlineto{\pgfqpoint{1.090230in}{0.610533in}}%
\pgfpathlineto{\pgfqpoint{1.090230in}{0.610533in}}%
\pgfpathclose%
\pgfusepath{fill}%
\end{pgfscope}%
\begin{pgfscope}%
\pgfpathrectangle{\pgfqpoint{0.407592in}{0.222904in}}{\pgfqpoint{1.411163in}{1.411163in}}%
\pgfusepath{clip}%
\pgfsetbuttcap%
\pgfsetroundjoin%
\definecolor{currentfill}{rgb}{0.749020,0.772549,0.792157}%
\pgfsetfillcolor{currentfill}%
\pgfsetfillopacity{0.160000}%
\pgfsetlinewidth{0.000000pt}%
\definecolor{currentstroke}{rgb}{0.000000,0.000000,0.000000}%
\pgfsetstrokecolor{currentstroke}%
\pgfsetdash{}{0pt}%
\pgfpathmoveto{\pgfqpoint{1.039536in}{0.661320in}}%
\pgfpathlineto{\pgfqpoint{1.026846in}{0.674032in}}%
\pgfpathlineto{\pgfqpoint{1.014150in}{0.686752in}}%
\pgfpathlineto{\pgfqpoint{1.012792in}{0.684547in}}%
\pgfpathlineto{\pgfqpoint{1.011770in}{0.682314in}}%
\pgfpathlineto{\pgfqpoint{1.024877in}{0.670323in}}%
\pgfpathlineto{\pgfqpoint{1.037972in}{0.658343in}}%
\pgfpathlineto{\pgfqpoint{1.038642in}{0.659840in}}%
\pgfpathlineto{\pgfqpoint{1.039536in}{0.661320in}}%
\pgfpathclose%
\pgfusepath{fill}%
\end{pgfscope}%
\begin{pgfscope}%
\pgfpathrectangle{\pgfqpoint{0.407592in}{0.222904in}}{\pgfqpoint{1.411163in}{1.411163in}}%
\pgfusepath{clip}%
\pgfsetbuttcap%
\pgfsetroundjoin%
\definecolor{currentfill}{rgb}{0.749020,0.772549,0.792157}%
\pgfsetfillcolor{currentfill}%
\pgfsetfillopacity{0.160000}%
\pgfsetlinewidth{0.000000pt}%
\definecolor{currentstroke}{rgb}{0.000000,0.000000,0.000000}%
\pgfsetstrokecolor{currentstroke}%
\pgfsetdash{}{0pt}%
\pgfpathmoveto{\pgfqpoint{1.090230in}{0.610533in}}%
\pgfpathlineto{\pgfqpoint{1.098256in}{0.617422in}}%
\pgfpathlineto{\pgfqpoint{1.106312in}{0.624336in}}%
\pgfpathlineto{\pgfqpoint{1.107658in}{0.624849in}}%
\pgfpathlineto{\pgfqpoint{1.108927in}{0.625398in}}%
\pgfpathlineto{\pgfqpoint{1.099562in}{0.617953in}}%
\pgfpathlineto{\pgfqpoint{1.090230in}{0.610533in}}%
\pgfpathlineto{\pgfqpoint{1.090230in}{0.610533in}}%
\pgfpathlineto{\pgfqpoint{1.090230in}{0.610533in}}%
\pgfpathclose%
\pgfusepath{fill}%
\end{pgfscope}%
\begin{pgfscope}%
\pgfpathrectangle{\pgfqpoint{0.407592in}{0.222904in}}{\pgfqpoint{1.411163in}{1.411163in}}%
\pgfusepath{clip}%
\pgfsetbuttcap%
\pgfsetroundjoin%
\definecolor{currentfill}{rgb}{0.749020,0.772549,0.792157}%
\pgfsetfillcolor{currentfill}%
\pgfsetfillopacity{0.160000}%
\pgfsetlinewidth{0.000000pt}%
\definecolor{currentstroke}{rgb}{0.000000,0.000000,0.000000}%
\pgfsetstrokecolor{currentstroke}%
\pgfsetdash{}{0pt}%
\pgfpathmoveto{\pgfqpoint{1.090230in}{0.610533in}}%
\pgfpathlineto{\pgfqpoint{1.080753in}{0.617625in}}%
\pgfpathlineto{\pgfqpoint{1.071242in}{0.624743in}}%
\pgfpathlineto{\pgfqpoint{1.072499in}{0.624238in}}%
\pgfpathlineto{\pgfqpoint{1.073835in}{0.623771in}}%
\pgfpathlineto{\pgfqpoint{1.082048in}{0.617139in}}%
\pgfpathlineto{\pgfqpoint{1.090230in}{0.610533in}}%
\pgfpathlineto{\pgfqpoint{1.090230in}{0.610533in}}%
\pgfpathlineto{\pgfqpoint{1.090230in}{0.610533in}}%
\pgfpathclose%
\pgfusepath{fill}%
\end{pgfscope}%
\begin{pgfscope}%
\pgfpathrectangle{\pgfqpoint{0.407592in}{0.222904in}}{\pgfqpoint{1.411163in}{1.411163in}}%
\pgfusepath{clip}%
\pgfsetbuttcap%
\pgfsetroundjoin%
\definecolor{currentfill}{rgb}{0.749020,0.772549,0.792157}%
\pgfsetfillcolor{currentfill}%
\pgfsetfillopacity{0.160000}%
\pgfsetlinewidth{0.000000pt}%
\definecolor{currentstroke}{rgb}{0.000000,0.000000,0.000000}%
\pgfsetstrokecolor{currentstroke}%
\pgfsetdash{}{0pt}%
\pgfpathmoveto{\pgfqpoint{1.090230in}{0.610533in}}%
\pgfpathlineto{\pgfqpoint{1.096808in}{0.616968in}}%
\pgfpathlineto{\pgfqpoint{1.103412in}{0.623428in}}%
\pgfpathlineto{\pgfqpoint{1.104895in}{0.623862in}}%
\pgfpathlineto{\pgfqpoint{1.106312in}{0.624336in}}%
\pgfpathlineto{\pgfqpoint{1.098256in}{0.617422in}}%
\pgfpathlineto{\pgfqpoint{1.090230in}{0.610533in}}%
\pgfpathlineto{\pgfqpoint{1.090230in}{0.610533in}}%
\pgfpathlineto{\pgfqpoint{1.090230in}{0.610533in}}%
\pgfpathclose%
\pgfusepath{fill}%
\end{pgfscope}%
\begin{pgfscope}%
\pgfpathrectangle{\pgfqpoint{0.407592in}{0.222904in}}{\pgfqpoint{1.411163in}{1.411163in}}%
\pgfusepath{clip}%
\pgfsetbuttcap%
\pgfsetroundjoin%
\definecolor{currentfill}{rgb}{0.749020,0.772549,0.792157}%
\pgfsetfillcolor{currentfill}%
\pgfsetfillopacity{0.160000}%
\pgfsetlinewidth{0.000000pt}%
\definecolor{currentstroke}{rgb}{0.000000,0.000000,0.000000}%
\pgfsetstrokecolor{currentstroke}%
\pgfsetdash{}{0pt}%
\pgfpathmoveto{\pgfqpoint{1.270742in}{0.789733in}}%
\pgfpathlineto{\pgfqpoint{1.283681in}{0.802577in}}%
\pgfpathlineto{\pgfqpoint{1.296626in}{0.815428in}}%
\pgfpathlineto{\pgfqpoint{1.292924in}{0.820977in}}%
\pgfpathlineto{\pgfqpoint{1.288355in}{0.826387in}}%
\pgfpathlineto{\pgfqpoint{1.275967in}{0.812890in}}%
\pgfpathlineto{\pgfqpoint{1.263579in}{0.799394in}}%
\pgfpathlineto{\pgfqpoint{1.267541in}{0.794623in}}%
\pgfpathlineto{\pgfqpoint{1.270742in}{0.789733in}}%
\pgfpathclose%
\pgfusepath{fill}%
\end{pgfscope}%
\begin{pgfscope}%
\pgfpathrectangle{\pgfqpoint{0.407592in}{0.222904in}}{\pgfqpoint{1.411163in}{1.411163in}}%
\pgfusepath{clip}%
\pgfsetbuttcap%
\pgfsetroundjoin%
\definecolor{currentfill}{rgb}{0.749020,0.772549,0.792157}%
\pgfsetfillcolor{currentfill}%
\pgfsetfillopacity{0.160000}%
\pgfsetlinewidth{0.000000pt}%
\definecolor{currentstroke}{rgb}{0.000000,0.000000,0.000000}%
\pgfsetstrokecolor{currentstroke}%
\pgfsetdash{}{0pt}%
\pgfpathmoveto{\pgfqpoint{0.945418in}{0.771587in}}%
\pgfpathlineto{\pgfqpoint{0.933343in}{0.785017in}}%
\pgfpathlineto{\pgfqpoint{0.921266in}{0.798448in}}%
\pgfpathlineto{\pgfqpoint{0.916437in}{0.793651in}}%
\pgfpathlineto{\pgfqpoint{0.912349in}{0.788739in}}%
\pgfpathlineto{\pgfqpoint{0.925097in}{0.775968in}}%
\pgfpathlineto{\pgfqpoint{0.937838in}{0.763203in}}%
\pgfpathlineto{\pgfqpoint{0.941310in}{0.767444in}}%
\pgfpathlineto{\pgfqpoint{0.945418in}{0.771587in}}%
\pgfpathclose%
\pgfusepath{fill}%
\end{pgfscope}%
\begin{pgfscope}%
\pgfpathrectangle{\pgfqpoint{0.407592in}{0.222904in}}{\pgfqpoint{1.411163in}{1.411163in}}%
\pgfusepath{clip}%
\pgfsetbuttcap%
\pgfsetroundjoin%
\definecolor{currentfill}{rgb}{0.749020,0.772549,0.792157}%
\pgfsetfillcolor{currentfill}%
\pgfsetfillopacity{0.160000}%
\pgfsetlinewidth{0.000000pt}%
\definecolor{currentstroke}{rgb}{0.000000,0.000000,0.000000}%
\pgfsetstrokecolor{currentstroke}%
\pgfsetdash{}{0pt}%
\pgfpathmoveto{\pgfqpoint{1.090230in}{0.610533in}}%
\pgfpathlineto{\pgfqpoint{1.082048in}{0.617139in}}%
\pgfpathlineto{\pgfqpoint{1.073835in}{0.623771in}}%
\pgfpathlineto{\pgfqpoint{1.075243in}{0.623346in}}%
\pgfpathlineto{\pgfqpoint{1.076716in}{0.622964in}}%
\pgfpathlineto{\pgfqpoint{1.083487in}{0.616735in}}%
\pgfpathlineto{\pgfqpoint{1.090230in}{0.610533in}}%
\pgfpathlineto{\pgfqpoint{1.090230in}{0.610533in}}%
\pgfpathlineto{\pgfqpoint{1.090230in}{0.610533in}}%
\pgfpathclose%
\pgfusepath{fill}%
\end{pgfscope}%
\begin{pgfscope}%
\pgfpathrectangle{\pgfqpoint{0.407592in}{0.222904in}}{\pgfqpoint{1.411163in}{1.411163in}}%
\pgfusepath{clip}%
\pgfsetbuttcap%
\pgfsetroundjoin%
\definecolor{currentfill}{rgb}{0.749020,0.772549,0.792157}%
\pgfsetfillcolor{currentfill}%
\pgfsetfillopacity{0.160000}%
\pgfsetlinewidth{0.000000pt}%
\definecolor{currentstroke}{rgb}{0.000000,0.000000,0.000000}%
\pgfsetstrokecolor{currentstroke}%
\pgfsetdash{}{0pt}%
\pgfpathmoveto{\pgfqpoint{1.090230in}{0.610533in}}%
\pgfpathlineto{\pgfqpoint{1.095244in}{0.616599in}}%
\pgfpathlineto{\pgfqpoint{1.100280in}{0.622690in}}%
\pgfpathlineto{\pgfqpoint{1.101872in}{0.623037in}}%
\pgfpathlineto{\pgfqpoint{1.103412in}{0.623428in}}%
\pgfpathlineto{\pgfqpoint{1.096808in}{0.616968in}}%
\pgfpathlineto{\pgfqpoint{1.090230in}{0.610533in}}%
\pgfpathlineto{\pgfqpoint{1.090230in}{0.610533in}}%
\pgfpathlineto{\pgfqpoint{1.090230in}{0.610533in}}%
\pgfpathclose%
\pgfusepath{fill}%
\end{pgfscope}%
\begin{pgfscope}%
\pgfpathrectangle{\pgfqpoint{0.407592in}{0.222904in}}{\pgfqpoint{1.411163in}{1.411163in}}%
\pgfusepath{clip}%
\pgfsetbuttcap%
\pgfsetroundjoin%
\definecolor{currentfill}{rgb}{0.749020,0.772549,0.792157}%
\pgfsetfillcolor{currentfill}%
\pgfsetfillopacity{0.160000}%
\pgfsetlinewidth{0.000000pt}%
\definecolor{currentstroke}{rgb}{0.000000,0.000000,0.000000}%
\pgfsetstrokecolor{currentstroke}%
\pgfsetdash{}{0pt}%
\pgfpathmoveto{\pgfqpoint{1.090230in}{0.610533in}}%
\pgfpathlineto{\pgfqpoint{1.083487in}{0.616735in}}%
\pgfpathlineto{\pgfqpoint{1.076716in}{0.622964in}}%
\pgfpathlineto{\pgfqpoint{1.078249in}{0.622626in}}%
\pgfpathlineto{\pgfqpoint{1.079835in}{0.622335in}}%
\pgfpathlineto{\pgfqpoint{1.085043in}{0.616421in}}%
\pgfpathlineto{\pgfqpoint{1.090230in}{0.610533in}}%
\pgfpathlineto{\pgfqpoint{1.090230in}{0.610533in}}%
\pgfpathlineto{\pgfqpoint{1.090230in}{0.610533in}}%
\pgfpathclose%
\pgfusepath{fill}%
\end{pgfscope}%
\begin{pgfscope}%
\pgfpathrectangle{\pgfqpoint{0.407592in}{0.222904in}}{\pgfqpoint{1.411163in}{1.411163in}}%
\pgfusepath{clip}%
\pgfsetbuttcap%
\pgfsetroundjoin%
\definecolor{currentfill}{rgb}{0.749020,0.772549,0.792157}%
\pgfsetfillcolor{currentfill}%
\pgfsetfillopacity{0.160000}%
\pgfsetlinewidth{0.000000pt}%
\definecolor{currentstroke}{rgb}{0.000000,0.000000,0.000000}%
\pgfsetstrokecolor{currentstroke}%
\pgfsetdash{}{0pt}%
\pgfpathmoveto{\pgfqpoint{1.090230in}{0.610533in}}%
\pgfpathlineto{\pgfqpoint{1.093592in}{0.616322in}}%
\pgfpathlineto{\pgfqpoint{1.096969in}{0.622136in}}%
\pgfpathlineto{\pgfqpoint{1.098643in}{0.622389in}}%
\pgfpathlineto{\pgfqpoint{1.100280in}{0.622690in}}%
\pgfpathlineto{\pgfqpoint{1.095244in}{0.616599in}}%
\pgfpathlineto{\pgfqpoint{1.090230in}{0.610533in}}%
\pgfpathlineto{\pgfqpoint{1.090230in}{0.610533in}}%
\pgfpathlineto{\pgfqpoint{1.090230in}{0.610533in}}%
\pgfpathclose%
\pgfusepath{fill}%
\end{pgfscope}%
\begin{pgfscope}%
\pgfpathrectangle{\pgfqpoint{0.407592in}{0.222904in}}{\pgfqpoint{1.411163in}{1.411163in}}%
\pgfusepath{clip}%
\pgfsetbuttcap%
\pgfsetroundjoin%
\definecolor{currentfill}{rgb}{0.749020,0.772549,0.792157}%
\pgfsetfillcolor{currentfill}%
\pgfsetfillopacity{0.160000}%
\pgfsetlinewidth{0.000000pt}%
\definecolor{currentstroke}{rgb}{0.000000,0.000000,0.000000}%
\pgfsetstrokecolor{currentstroke}%
\pgfsetdash{}{0pt}%
\pgfpathmoveto{\pgfqpoint{1.090230in}{0.610533in}}%
\pgfpathlineto{\pgfqpoint{1.085043in}{0.616421in}}%
\pgfpathlineto{\pgfqpoint{1.079835in}{0.622335in}}%
\pgfpathlineto{\pgfqpoint{1.081466in}{0.622091in}}%
\pgfpathlineto{\pgfqpoint{1.083135in}{0.621895in}}%
\pgfpathlineto{\pgfqpoint{1.086690in}{0.616201in}}%
\pgfpathlineto{\pgfqpoint{1.090230in}{0.610533in}}%
\pgfpathlineto{\pgfqpoint{1.090230in}{0.610533in}}%
\pgfpathlineto{\pgfqpoint{1.090230in}{0.610533in}}%
\pgfpathclose%
\pgfusepath{fill}%
\end{pgfscope}%
\begin{pgfscope}%
\pgfpathrectangle{\pgfqpoint{0.407592in}{0.222904in}}{\pgfqpoint{1.411163in}{1.411163in}}%
\pgfusepath{clip}%
\pgfsetbuttcap%
\pgfsetroundjoin%
\definecolor{currentfill}{rgb}{0.749020,0.772549,0.792157}%
\pgfsetfillcolor{currentfill}%
\pgfsetfillopacity{0.160000}%
\pgfsetlinewidth{0.000000pt}%
\definecolor{currentstroke}{rgb}{0.000000,0.000000,0.000000}%
\pgfsetstrokecolor{currentstroke}%
\pgfsetdash{}{0pt}%
\pgfpathmoveto{\pgfqpoint{1.090230in}{0.610533in}}%
\pgfpathlineto{\pgfqpoint{1.091880in}{0.616141in}}%
\pgfpathlineto{\pgfqpoint{1.093538in}{0.621774in}}%
\pgfpathlineto{\pgfqpoint{1.095265in}{0.621930in}}%
\pgfpathlineto{\pgfqpoint{1.096969in}{0.622136in}}%
\pgfpathlineto{\pgfqpoint{1.093592in}{0.616322in}}%
\pgfpathlineto{\pgfqpoint{1.090230in}{0.610533in}}%
\pgfpathlineto{\pgfqpoint{1.090230in}{0.610533in}}%
\pgfpathlineto{\pgfqpoint{1.090230in}{0.610533in}}%
\pgfpathclose%
\pgfusepath{fill}%
\end{pgfscope}%
\begin{pgfscope}%
\pgfpathrectangle{\pgfqpoint{0.407592in}{0.222904in}}{\pgfqpoint{1.411163in}{1.411163in}}%
\pgfusepath{clip}%
\pgfsetbuttcap%
\pgfsetroundjoin%
\definecolor{currentfill}{rgb}{0.749020,0.772549,0.792157}%
\pgfsetfillcolor{currentfill}%
\pgfsetfillopacity{0.160000}%
\pgfsetlinewidth{0.000000pt}%
\definecolor{currentstroke}{rgb}{0.000000,0.000000,0.000000}%
\pgfsetstrokecolor{currentstroke}%
\pgfsetdash{}{0pt}%
\pgfpathmoveto{\pgfqpoint{1.296626in}{0.815428in}}%
\pgfpathlineto{\pgfqpoint{1.309577in}{0.828285in}}%
\pgfpathlineto{\pgfqpoint{1.322534in}{0.841148in}}%
\pgfpathlineto{\pgfqpoint{1.318321in}{0.847344in}}%
\pgfpathlineto{\pgfqpoint{1.313134in}{0.853384in}}%
\pgfpathlineto{\pgfqpoint{1.300744in}{0.839885in}}%
\pgfpathlineto{\pgfqpoint{1.288355in}{0.826387in}}%
\pgfpathlineto{\pgfqpoint{1.292924in}{0.820977in}}%
\pgfpathlineto{\pgfqpoint{1.296626in}{0.815428in}}%
\pgfpathclose%
\pgfusepath{fill}%
\end{pgfscope}%
\begin{pgfscope}%
\pgfpathrectangle{\pgfqpoint{0.407592in}{0.222904in}}{\pgfqpoint{1.411163in}{1.411163in}}%
\pgfusepath{clip}%
\pgfsetbuttcap%
\pgfsetroundjoin%
\definecolor{currentfill}{rgb}{0.749020,0.772549,0.792157}%
\pgfsetfillcolor{currentfill}%
\pgfsetfillopacity{0.160000}%
\pgfsetlinewidth{0.000000pt}%
\definecolor{currentstroke}{rgb}{0.000000,0.000000,0.000000}%
\pgfsetstrokecolor{currentstroke}%
\pgfsetdash{}{0pt}%
\pgfpathmoveto{\pgfqpoint{1.090230in}{0.610533in}}%
\pgfpathlineto{\pgfqpoint{1.086690in}{0.616201in}}%
\pgfpathlineto{\pgfqpoint{1.083135in}{0.621895in}}%
\pgfpathlineto{\pgfqpoint{1.084836in}{0.621749in}}%
\pgfpathlineto{\pgfqpoint{1.086560in}{0.621653in}}%
\pgfpathlineto{\pgfqpoint{1.088399in}{0.616080in}}%
\pgfpathlineto{\pgfqpoint{1.090230in}{0.610533in}}%
\pgfpathlineto{\pgfqpoint{1.090230in}{0.610533in}}%
\pgfpathlineto{\pgfqpoint{1.090230in}{0.610533in}}%
\pgfpathclose%
\pgfusepath{fill}%
\end{pgfscope}%
\begin{pgfscope}%
\pgfpathrectangle{\pgfqpoint{0.407592in}{0.222904in}}{\pgfqpoint{1.411163in}{1.411163in}}%
\pgfusepath{clip}%
\pgfsetbuttcap%
\pgfsetroundjoin%
\definecolor{currentfill}{rgb}{0.749020,0.772549,0.792157}%
\pgfsetfillcolor{currentfill}%
\pgfsetfillopacity{0.160000}%
\pgfsetlinewidth{0.000000pt}%
\definecolor{currentstroke}{rgb}{0.000000,0.000000,0.000000}%
\pgfsetstrokecolor{currentstroke}%
\pgfsetdash{}{0pt}%
\pgfpathmoveto{\pgfqpoint{1.090230in}{0.610533in}}%
\pgfpathlineto{\pgfqpoint{1.090139in}{0.616060in}}%
\pgfpathlineto{\pgfqpoint{1.090048in}{0.621612in}}%
\pgfpathlineto{\pgfqpoint{1.091796in}{0.621668in}}%
\pgfpathlineto{\pgfqpoint{1.093538in}{0.621774in}}%
\pgfpathlineto{\pgfqpoint{1.091880in}{0.616141in}}%
\pgfpathlineto{\pgfqpoint{1.090230in}{0.610533in}}%
\pgfpathlineto{\pgfqpoint{1.090230in}{0.610533in}}%
\pgfpathlineto{\pgfqpoint{1.090230in}{0.610533in}}%
\pgfpathclose%
\pgfusepath{fill}%
\end{pgfscope}%
\begin{pgfscope}%
\pgfpathrectangle{\pgfqpoint{0.407592in}{0.222904in}}{\pgfqpoint{1.411163in}{1.411163in}}%
\pgfusepath{clip}%
\pgfsetbuttcap%
\pgfsetroundjoin%
\definecolor{currentfill}{rgb}{0.749020,0.772549,0.792157}%
\pgfsetfillcolor{currentfill}%
\pgfsetfillopacity{0.160000}%
\pgfsetlinewidth{0.000000pt}%
\definecolor{currentstroke}{rgb}{0.000000,0.000000,0.000000}%
\pgfsetstrokecolor{currentstroke}%
\pgfsetdash{}{0pt}%
\pgfpathmoveto{\pgfqpoint{1.090230in}{0.610533in}}%
\pgfpathlineto{\pgfqpoint{1.088399in}{0.616080in}}%
\pgfpathlineto{\pgfqpoint{1.086560in}{0.621653in}}%
\pgfpathlineto{\pgfqpoint{1.088300in}{0.621607in}}%
\pgfpathlineto{\pgfqpoint{1.090048in}{0.621612in}}%
\pgfpathlineto{\pgfqpoint{1.090139in}{0.616060in}}%
\pgfpathlineto{\pgfqpoint{1.090230in}{0.610533in}}%
\pgfpathlineto{\pgfqpoint{1.090230in}{0.610533in}}%
\pgfpathlineto{\pgfqpoint{1.090230in}{0.610533in}}%
\pgfpathclose%
\pgfusepath{fill}%
\end{pgfscope}%
\begin{pgfscope}%
\pgfpathrectangle{\pgfqpoint{0.407592in}{0.222904in}}{\pgfqpoint{1.411163in}{1.411163in}}%
\pgfusepath{clip}%
\pgfsetbuttcap%
\pgfsetroundjoin%
\definecolor{currentfill}{rgb}{0.749020,0.772549,0.792157}%
\pgfsetfillcolor{currentfill}%
\pgfsetfillopacity{0.160000}%
\pgfsetlinewidth{0.000000pt}%
\definecolor{currentstroke}{rgb}{0.000000,0.000000,0.000000}%
\pgfsetstrokecolor{currentstroke}%
\pgfsetdash{}{0pt}%
\pgfpathmoveto{\pgfqpoint{0.921266in}{0.798448in}}%
\pgfpathlineto{\pgfqpoint{0.909188in}{0.811880in}}%
\pgfpathlineto{\pgfqpoint{0.897109in}{0.825314in}}%
\pgfpathlineto{\pgfqpoint{0.891550in}{0.819874in}}%
\pgfpathlineto{\pgfqpoint{0.886834in}{0.814300in}}%
\pgfpathlineto{\pgfqpoint{0.899595in}{0.801516in}}%
\pgfpathlineto{\pgfqpoint{0.912349in}{0.788739in}}%
\pgfpathlineto{\pgfqpoint{0.916437in}{0.793651in}}%
\pgfpathlineto{\pgfqpoint{0.921266in}{0.798448in}}%
\pgfpathclose%
\pgfusepath{fill}%
\end{pgfscope}%
\begin{pgfscope}%
\pgfpathrectangle{\pgfqpoint{0.407592in}{0.222904in}}{\pgfqpoint{1.411163in}{1.411163in}}%
\pgfusepath{clip}%
\pgfsetbuttcap%
\pgfsetroundjoin%
\definecolor{currentfill}{rgb}{0.749020,0.772549,0.792157}%
\pgfsetfillcolor{currentfill}%
\pgfsetfillopacity{0.160000}%
\pgfsetlinewidth{0.000000pt}%
\definecolor{currentstroke}{rgb}{0.000000,0.000000,0.000000}%
\pgfsetstrokecolor{currentstroke}%
\pgfsetdash{}{0pt}%
\pgfpathmoveto{\pgfqpoint{1.116100in}{0.631528in}}%
\pgfpathlineto{\pgfqpoint{1.129068in}{0.642053in}}%
\pgfpathlineto{\pgfqpoint{1.142060in}{0.652597in}}%
\pgfpathlineto{\pgfqpoint{1.142573in}{0.654106in}}%
\pgfpathlineto{\pgfqpoint{1.142855in}{0.655621in}}%
\pgfpathlineto{\pgfqpoint{1.129673in}{0.644327in}}%
\pgfpathlineto{\pgfqpoint{1.116508in}{0.633047in}}%
\pgfpathlineto{\pgfqpoint{1.116361in}{0.632286in}}%
\pgfpathlineto{\pgfqpoint{1.116100in}{0.631528in}}%
\pgfpathclose%
\pgfusepath{fill}%
\end{pgfscope}%
\begin{pgfscope}%
\pgfpathrectangle{\pgfqpoint{0.407592in}{0.222904in}}{\pgfqpoint{1.411163in}{1.411163in}}%
\pgfusepath{clip}%
\pgfsetbuttcap%
\pgfsetroundjoin%
\definecolor{currentfill}{rgb}{0.749020,0.772549,0.792157}%
\pgfsetfillcolor{currentfill}%
\pgfsetfillopacity{0.160000}%
\pgfsetlinewidth{0.000000pt}%
\definecolor{currentstroke}{rgb}{0.000000,0.000000,0.000000}%
\pgfsetstrokecolor{currentstroke}%
\pgfsetdash{}{0pt}%
\pgfpathmoveto{\pgfqpoint{1.063809in}{0.632895in}}%
\pgfpathlineto{\pgfqpoint{1.050572in}{0.644098in}}%
\pgfpathlineto{\pgfqpoint{1.037317in}{0.655317in}}%
\pgfpathlineto{\pgfqpoint{1.037335in}{0.653802in}}%
\pgfpathlineto{\pgfqpoint{1.037585in}{0.652295in}}%
\pgfpathlineto{\pgfqpoint{1.050781in}{0.641827in}}%
\pgfpathlineto{\pgfqpoint{1.063954in}{0.631377in}}%
\pgfpathlineto{\pgfqpoint{1.063824in}{0.632134in}}%
\pgfpathlineto{\pgfqpoint{1.063809in}{0.632895in}}%
\pgfpathclose%
\pgfusepath{fill}%
\end{pgfscope}%
\begin{pgfscope}%
\pgfpathrectangle{\pgfqpoint{0.407592in}{0.222904in}}{\pgfqpoint{1.411163in}{1.411163in}}%
\pgfusepath{clip}%
\pgfsetbuttcap%
\pgfsetroundjoin%
\definecolor{currentfill}{rgb}{0.749020,0.772549,0.792157}%
\pgfsetfillcolor{currentfill}%
\pgfsetfillopacity{0.160000}%
\pgfsetlinewidth{0.000000pt}%
\definecolor{currentstroke}{rgb}{0.000000,0.000000,0.000000}%
\pgfsetstrokecolor{currentstroke}%
\pgfsetdash{}{0pt}%
\pgfpathmoveto{\pgfqpoint{1.169048in}{0.682764in}}%
\pgfpathlineto{\pgfqpoint{1.182225in}{0.694840in}}%
\pgfpathlineto{\pgfqpoint{1.195414in}{0.706927in}}%
\pgfpathlineto{\pgfqpoint{1.194552in}{0.709879in}}%
\pgfpathlineto{\pgfqpoint{1.193235in}{0.712789in}}%
\pgfpathlineto{\pgfqpoint{1.180339in}{0.699986in}}%
\pgfpathlineto{\pgfqpoint{1.167448in}{0.687189in}}%
\pgfpathlineto{\pgfqpoint{1.168418in}{0.684992in}}%
\pgfpathlineto{\pgfqpoint{1.169048in}{0.682764in}}%
\pgfpathclose%
\pgfusepath{fill}%
\end{pgfscope}%
\begin{pgfscope}%
\pgfpathrectangle{\pgfqpoint{0.407592in}{0.222904in}}{\pgfqpoint{1.411163in}{1.411163in}}%
\pgfusepath{clip}%
\pgfsetbuttcap%
\pgfsetroundjoin%
\definecolor{currentfill}{rgb}{0.749020,0.772549,0.792157}%
\pgfsetfillcolor{currentfill}%
\pgfsetfillopacity{0.160000}%
\pgfsetlinewidth{0.000000pt}%
\definecolor{currentstroke}{rgb}{0.000000,0.000000,0.000000}%
\pgfsetstrokecolor{currentstroke}%
\pgfsetdash{}{0pt}%
\pgfpathmoveto{\pgfqpoint{1.322534in}{0.841148in}}%
\pgfpathlineto{\pgfqpoint{1.335497in}{0.854017in}}%
\pgfpathlineto{\pgfqpoint{1.348467in}{0.866892in}}%
\pgfpathlineto{\pgfqpoint{1.343730in}{0.873725in}}%
\pgfpathlineto{\pgfqpoint{1.337916in}{0.880383in}}%
\pgfpathlineto{\pgfqpoint{1.325524in}{0.866883in}}%
\pgfpathlineto{\pgfqpoint{1.313134in}{0.853384in}}%
\pgfpathlineto{\pgfqpoint{1.318321in}{0.847344in}}%
\pgfpathlineto{\pgfqpoint{1.322534in}{0.841148in}}%
\pgfpathclose%
\pgfusepath{fill}%
\end{pgfscope}%
\begin{pgfscope}%
\pgfpathrectangle{\pgfqpoint{0.407592in}{0.222904in}}{\pgfqpoint{1.411163in}{1.411163in}}%
\pgfusepath{clip}%
\pgfsetbuttcap%
\pgfsetroundjoin%
\definecolor{currentfill}{rgb}{0.749020,0.772549,0.792157}%
\pgfsetfillcolor{currentfill}%
\pgfsetfillopacity{0.160000}%
\pgfsetlinewidth{0.000000pt}%
\definecolor{currentstroke}{rgb}{0.000000,0.000000,0.000000}%
\pgfsetstrokecolor{currentstroke}%
\pgfsetdash{}{0pt}%
\pgfpathmoveto{\pgfqpoint{1.014150in}{0.686752in}}%
\pgfpathlineto{\pgfqpoint{1.001448in}{0.699477in}}%
\pgfpathlineto{\pgfqpoint{0.988739in}{0.712209in}}%
\pgfpathlineto{\pgfqpoint{0.986906in}{0.709290in}}%
\pgfpathlineto{\pgfqpoint{0.985520in}{0.706331in}}%
\pgfpathlineto{\pgfqpoint{0.998651in}{0.694317in}}%
\pgfpathlineto{\pgfqpoint{1.011770in}{0.682314in}}%
\pgfpathlineto{\pgfqpoint{1.012792in}{0.684547in}}%
\pgfpathlineto{\pgfqpoint{1.014150in}{0.686752in}}%
\pgfpathclose%
\pgfusepath{fill}%
\end{pgfscope}%
\begin{pgfscope}%
\pgfpathrectangle{\pgfqpoint{0.407592in}{0.222904in}}{\pgfqpoint{1.411163in}{1.411163in}}%
\pgfusepath{clip}%
\pgfsetbuttcap%
\pgfsetroundjoin%
\definecolor{currentfill}{rgb}{0.749020,0.772549,0.792157}%
\pgfsetfillcolor{currentfill}%
\pgfsetfillopacity{0.160000}%
\pgfsetlinewidth{0.000000pt}%
\definecolor{currentstroke}{rgb}{0.000000,0.000000,0.000000}%
\pgfsetstrokecolor{currentstroke}%
\pgfsetdash{}{0pt}%
\pgfpathmoveto{\pgfqpoint{0.897109in}{0.825314in}}%
\pgfpathlineto{\pgfqpoint{0.885029in}{0.838749in}}%
\pgfpathlineto{\pgfqpoint{0.872948in}{0.852186in}}%
\pgfpathlineto{\pgfqpoint{0.866648in}{0.846113in}}%
\pgfpathlineto{\pgfqpoint{0.861293in}{0.839888in}}%
\pgfpathlineto{\pgfqpoint{0.874067in}{0.827091in}}%
\pgfpathlineto{\pgfqpoint{0.886834in}{0.814300in}}%
\pgfpathlineto{\pgfqpoint{0.891550in}{0.819874in}}%
\pgfpathlineto{\pgfqpoint{0.897109in}{0.825314in}}%
\pgfpathclose%
\pgfusepath{fill}%
\end{pgfscope}%
\begin{pgfscope}%
\pgfpathrectangle{\pgfqpoint{0.407592in}{0.222904in}}{\pgfqpoint{1.411163in}{1.411163in}}%
\pgfusepath{clip}%
\pgfsetbuttcap%
\pgfsetroundjoin%
\definecolor{currentfill}{rgb}{0.749020,0.772549,0.792157}%
\pgfsetfillcolor{currentfill}%
\pgfsetfillopacity{0.160000}%
\pgfsetlinewidth{0.000000pt}%
\definecolor{currentstroke}{rgb}{0.000000,0.000000,0.000000}%
\pgfsetstrokecolor{currentstroke}%
\pgfsetdash{}{0pt}%
\pgfpathmoveto{\pgfqpoint{1.142855in}{0.655621in}}%
\pgfpathlineto{\pgfqpoint{1.156055in}{0.666930in}}%
\pgfpathlineto{\pgfqpoint{1.169272in}{0.678255in}}%
\pgfpathlineto{\pgfqpoint{1.169333in}{0.680515in}}%
\pgfpathlineto{\pgfqpoint{1.169048in}{0.682764in}}%
\pgfpathlineto{\pgfqpoint{1.155882in}{0.670699in}}%
\pgfpathlineto{\pgfqpoint{1.142729in}{0.658644in}}%
\pgfpathlineto{\pgfqpoint{1.142907in}{0.657136in}}%
\pgfpathlineto{\pgfqpoint{1.142855in}{0.655621in}}%
\pgfpathclose%
\pgfusepath{fill}%
\end{pgfscope}%
\begin{pgfscope}%
\pgfpathrectangle{\pgfqpoint{0.407592in}{0.222904in}}{\pgfqpoint{1.411163in}{1.411163in}}%
\pgfusepath{clip}%
\pgfsetbuttcap%
\pgfsetroundjoin%
\definecolor{currentfill}{rgb}{0.749020,0.772549,0.792157}%
\pgfsetfillcolor{currentfill}%
\pgfsetfillopacity{0.160000}%
\pgfsetlinewidth{0.000000pt}%
\definecolor{currentstroke}{rgb}{0.000000,0.000000,0.000000}%
\pgfsetstrokecolor{currentstroke}%
\pgfsetdash{}{0pt}%
\pgfpathmoveto{\pgfqpoint{1.348467in}{0.866892in}}%
\pgfpathlineto{\pgfqpoint{1.361442in}{0.879773in}}%
\pgfpathlineto{\pgfqpoint{1.374423in}{0.892660in}}%
\pgfpathlineto{\pgfqpoint{1.369153in}{0.900120in}}%
\pgfpathlineto{\pgfqpoint{1.362701in}{0.907386in}}%
\pgfpathlineto{\pgfqpoint{1.350308in}{0.893884in}}%
\pgfpathlineto{\pgfqpoint{1.337916in}{0.880383in}}%
\pgfpathlineto{\pgfqpoint{1.343730in}{0.873725in}}%
\pgfpathlineto{\pgfqpoint{1.348467in}{0.866892in}}%
\pgfpathclose%
\pgfusepath{fill}%
\end{pgfscope}%
\begin{pgfscope}%
\pgfpathrectangle{\pgfqpoint{0.407592in}{0.222904in}}{\pgfqpoint{1.411163in}{1.411163in}}%
\pgfusepath{clip}%
\pgfsetbuttcap%
\pgfsetroundjoin%
\definecolor{currentfill}{rgb}{0.749020,0.772549,0.792157}%
\pgfsetfillcolor{currentfill}%
\pgfsetfillopacity{0.160000}%
\pgfsetlinewidth{0.000000pt}%
\definecolor{currentstroke}{rgb}{0.000000,0.000000,0.000000}%
\pgfsetstrokecolor{currentstroke}%
\pgfsetdash{}{0pt}%
\pgfpathmoveto{\pgfqpoint{1.037972in}{0.658343in}}%
\pgfpathlineto{\pgfqpoint{1.024877in}{0.670323in}}%
\pgfpathlineto{\pgfqpoint{1.011770in}{0.682314in}}%
\pgfpathlineto{\pgfqpoint{1.011089in}{0.680062in}}%
\pgfpathlineto{\pgfqpoint{1.010753in}{0.677801in}}%
\pgfpathlineto{\pgfqpoint{1.024044in}{0.666551in}}%
\pgfpathlineto{\pgfqpoint{1.037317in}{0.655317in}}%
\pgfpathlineto{\pgfqpoint{1.037530in}{0.656833in}}%
\pgfpathlineto{\pgfqpoint{1.037972in}{0.658343in}}%
\pgfpathclose%
\pgfusepath{fill}%
\end{pgfscope}%
\begin{pgfscope}%
\pgfpathrectangle{\pgfqpoint{0.407592in}{0.222904in}}{\pgfqpoint{1.411163in}{1.411163in}}%
\pgfusepath{clip}%
\pgfsetbuttcap%
\pgfsetroundjoin%
\definecolor{currentfill}{rgb}{0.749020,0.772549,0.792157}%
\pgfsetfillcolor{currentfill}%
\pgfsetfillopacity{0.160000}%
\pgfsetlinewidth{0.000000pt}%
\definecolor{currentstroke}{rgb}{0.000000,0.000000,0.000000}%
\pgfsetstrokecolor{currentstroke}%
\pgfsetdash{}{0pt}%
\pgfpathmoveto{\pgfqpoint{0.872948in}{0.852186in}}%
\pgfpathlineto{\pgfqpoint{0.860865in}{0.865624in}}%
\pgfpathlineto{\pgfqpoint{0.848781in}{0.879063in}}%
\pgfpathlineto{\pgfqpoint{0.841730in}{0.872367in}}%
\pgfpathlineto{\pgfqpoint{0.835726in}{0.865502in}}%
\pgfpathlineto{\pgfqpoint{0.848513in}{0.852692in}}%
\pgfpathlineto{\pgfqpoint{0.861293in}{0.839888in}}%
\pgfpathlineto{\pgfqpoint{0.866648in}{0.846113in}}%
\pgfpathlineto{\pgfqpoint{0.872948in}{0.852186in}}%
\pgfpathclose%
\pgfusepath{fill}%
\end{pgfscope}%
\begin{pgfscope}%
\pgfpathrectangle{\pgfqpoint{0.407592in}{0.222904in}}{\pgfqpoint{1.411163in}{1.411163in}}%
\pgfusepath{clip}%
\pgfsetbuttcap%
\pgfsetroundjoin%
\definecolor{currentfill}{rgb}{0.749020,0.772549,0.792157}%
\pgfsetfillcolor{currentfill}%
\pgfsetfillopacity{0.160000}%
\pgfsetlinewidth{0.000000pt}%
\definecolor{currentstroke}{rgb}{0.000000,0.000000,0.000000}%
\pgfsetstrokecolor{currentstroke}%
\pgfsetdash{}{0pt}%
\pgfpathmoveto{\pgfqpoint{1.115238in}{0.630036in}}%
\pgfpathlineto{\pgfqpoint{1.127782in}{0.639820in}}%
\pgfpathlineto{\pgfqpoint{1.140354in}{0.649624in}}%
\pgfpathlineto{\pgfqpoint{1.141320in}{0.651101in}}%
\pgfpathlineto{\pgfqpoint{1.142060in}{0.652597in}}%
\pgfpathlineto{\pgfqpoint{1.129068in}{0.642053in}}%
\pgfpathlineto{\pgfqpoint{1.116100in}{0.631528in}}%
\pgfpathlineto{\pgfqpoint{1.115725in}{0.630777in}}%
\pgfpathlineto{\pgfqpoint{1.115238in}{0.630036in}}%
\pgfpathclose%
\pgfusepath{fill}%
\end{pgfscope}%
\begin{pgfscope}%
\pgfpathrectangle{\pgfqpoint{0.407592in}{0.222904in}}{\pgfqpoint{1.411163in}{1.411163in}}%
\pgfusepath{clip}%
\pgfsetbuttcap%
\pgfsetroundjoin%
\definecolor{currentfill}{rgb}{0.749020,0.772549,0.792157}%
\pgfsetfillcolor{currentfill}%
\pgfsetfillopacity{0.160000}%
\pgfsetlinewidth{0.000000pt}%
\definecolor{currentstroke}{rgb}{0.000000,0.000000,0.000000}%
\pgfsetstrokecolor{currentstroke}%
\pgfsetdash{}{0pt}%
\pgfpathmoveto{\pgfqpoint{1.063954in}{0.631377in}}%
\pgfpathlineto{\pgfqpoint{1.050781in}{0.641827in}}%
\pgfpathlineto{\pgfqpoint{1.037585in}{0.652295in}}%
\pgfpathlineto{\pgfqpoint{1.038064in}{0.650803in}}%
\pgfpathlineto{\pgfqpoint{1.038772in}{0.649331in}}%
\pgfpathlineto{\pgfqpoint{1.051679in}{0.639599in}}%
\pgfpathlineto{\pgfqpoint{1.064558in}{0.629889in}}%
\pgfpathlineto{\pgfqpoint{1.064199in}{0.630628in}}%
\pgfpathlineto{\pgfqpoint{1.063954in}{0.631377in}}%
\pgfpathclose%
\pgfusepath{fill}%
\end{pgfscope}%
\begin{pgfscope}%
\pgfpathrectangle{\pgfqpoint{0.407592in}{0.222904in}}{\pgfqpoint{1.411163in}{1.411163in}}%
\pgfusepath{clip}%
\pgfsetbuttcap%
\pgfsetroundjoin%
\definecolor{currentfill}{rgb}{0.749020,0.772549,0.792157}%
\pgfsetfillcolor{currentfill}%
\pgfsetfillopacity{0.160000}%
\pgfsetlinewidth{0.000000pt}%
\definecolor{currentstroke}{rgb}{0.000000,0.000000,0.000000}%
\pgfsetstrokecolor{currentstroke}%
\pgfsetdash{}{0pt}%
\pgfpathmoveto{\pgfqpoint{1.374423in}{0.892660in}}%
\pgfpathlineto{\pgfqpoint{1.387411in}{0.905553in}}%
\pgfpathlineto{\pgfqpoint{1.400404in}{0.918452in}}%
\pgfpathlineto{\pgfqpoint{1.394589in}{0.926529in}}%
\pgfpathlineto{\pgfqpoint{1.387488in}{0.934391in}}%
\pgfpathlineto{\pgfqpoint{1.375094in}{0.920888in}}%
\pgfpathlineto{\pgfqpoint{1.362701in}{0.907386in}}%
\pgfpathlineto{\pgfqpoint{1.369153in}{0.900120in}}%
\pgfpathlineto{\pgfqpoint{1.374423in}{0.892660in}}%
\pgfpathclose%
\pgfusepath{fill}%
\end{pgfscope}%
\begin{pgfscope}%
\pgfpathrectangle{\pgfqpoint{0.407592in}{0.222904in}}{\pgfqpoint{1.411163in}{1.411163in}}%
\pgfusepath{clip}%
\pgfsetbuttcap%
\pgfsetroundjoin%
\definecolor{currentfill}{rgb}{0.749020,0.772549,0.792157}%
\pgfsetfillcolor{currentfill}%
\pgfsetfillopacity{0.160000}%
\pgfsetlinewidth{0.000000pt}%
\definecolor{currentstroke}{rgb}{0.000000,0.000000,0.000000}%
\pgfsetstrokecolor{currentstroke}%
\pgfsetdash{}{0pt}%
\pgfpathmoveto{\pgfqpoint{1.195414in}{0.706927in}}%
\pgfpathlineto{\pgfqpoint{1.208615in}{0.719024in}}%
\pgfpathlineto{\pgfqpoint{1.221828in}{0.731133in}}%
\pgfpathlineto{\pgfqpoint{1.220721in}{0.734800in}}%
\pgfpathlineto{\pgfqpoint{1.219047in}{0.738413in}}%
\pgfpathlineto{\pgfqpoint{1.206138in}{0.725598in}}%
\pgfpathlineto{\pgfqpoint{1.193235in}{0.712789in}}%
\pgfpathlineto{\pgfqpoint{1.194552in}{0.709879in}}%
\pgfpathlineto{\pgfqpoint{1.195414in}{0.706927in}}%
\pgfpathclose%
\pgfusepath{fill}%
\end{pgfscope}%
\begin{pgfscope}%
\pgfpathrectangle{\pgfqpoint{0.407592in}{0.222904in}}{\pgfqpoint{1.411163in}{1.411163in}}%
\pgfusepath{clip}%
\pgfsetbuttcap%
\pgfsetroundjoin%
\definecolor{currentfill}{rgb}{0.749020,0.772549,0.792157}%
\pgfsetfillcolor{currentfill}%
\pgfsetfillopacity{0.160000}%
\pgfsetlinewidth{0.000000pt}%
\definecolor{currentstroke}{rgb}{0.000000,0.000000,0.000000}%
\pgfsetstrokecolor{currentstroke}%
\pgfsetdash{}{0pt}%
\pgfpathmoveto{\pgfqpoint{0.848781in}{0.879063in}}%
\pgfpathlineto{\pgfqpoint{0.836696in}{0.892503in}}%
\pgfpathlineto{\pgfqpoint{0.824610in}{0.905945in}}%
\pgfpathlineto{\pgfqpoint{0.816798in}{0.898638in}}%
\pgfpathlineto{\pgfqpoint{0.810133in}{0.891142in}}%
\pgfpathlineto{\pgfqpoint{0.822932in}{0.878319in}}%
\pgfpathlineto{\pgfqpoint{0.835726in}{0.865502in}}%
\pgfpathlineto{\pgfqpoint{0.841730in}{0.872367in}}%
\pgfpathlineto{\pgfqpoint{0.848781in}{0.879063in}}%
\pgfpathclose%
\pgfusepath{fill}%
\end{pgfscope}%
\begin{pgfscope}%
\pgfpathrectangle{\pgfqpoint{0.407592in}{0.222904in}}{\pgfqpoint{1.411163in}{1.411163in}}%
\pgfusepath{clip}%
\pgfsetbuttcap%
\pgfsetroundjoin%
\definecolor{currentfill}{rgb}{0.749020,0.772549,0.792157}%
\pgfsetfillcolor{currentfill}%
\pgfsetfillopacity{0.160000}%
\pgfsetlinewidth{0.000000pt}%
\definecolor{currentstroke}{rgb}{0.000000,0.000000,0.000000}%
\pgfsetstrokecolor{currentstroke}%
\pgfsetdash{}{0pt}%
\pgfpathmoveto{\pgfqpoint{0.988739in}{0.712209in}}%
\pgfpathlineto{\pgfqpoint{0.976023in}{0.724948in}}%
\pgfpathlineto{\pgfqpoint{0.963301in}{0.737693in}}%
\pgfpathlineto{\pgfqpoint{0.960982in}{0.734068in}}%
\pgfpathlineto{\pgfqpoint{0.959220in}{0.730392in}}%
\pgfpathlineto{\pgfqpoint{0.972376in}{0.718356in}}%
\pgfpathlineto{\pgfqpoint{0.985520in}{0.706331in}}%
\pgfpathlineto{\pgfqpoint{0.986906in}{0.709290in}}%
\pgfpathlineto{\pgfqpoint{0.988739in}{0.712209in}}%
\pgfpathclose%
\pgfusepath{fill}%
\end{pgfscope}%
\begin{pgfscope}%
\pgfpathrectangle{\pgfqpoint{0.407592in}{0.222904in}}{\pgfqpoint{1.411163in}{1.411163in}}%
\pgfusepath{clip}%
\pgfsetbuttcap%
\pgfsetroundjoin%
\definecolor{currentfill}{rgb}{0.749020,0.772549,0.792157}%
\pgfsetfillcolor{currentfill}%
\pgfsetfillopacity{0.160000}%
\pgfsetlinewidth{0.000000pt}%
\definecolor{currentstroke}{rgb}{0.000000,0.000000,0.000000}%
\pgfsetstrokecolor{currentstroke}%
\pgfsetdash{}{0pt}%
\pgfpathmoveto{\pgfqpoint{1.400404in}{0.918452in}}%
\pgfpathlineto{\pgfqpoint{1.413404in}{0.931357in}}%
\pgfpathlineto{\pgfqpoint{1.426410in}{0.944268in}}%
\pgfpathlineto{\pgfqpoint{1.420039in}{0.952951in}}%
\pgfpathlineto{\pgfqpoint{1.412279in}{0.961401in}}%
\pgfpathlineto{\pgfqpoint{1.399883in}{0.947896in}}%
\pgfpathlineto{\pgfqpoint{1.387488in}{0.934391in}}%
\pgfpathlineto{\pgfqpoint{1.394589in}{0.926529in}}%
\pgfpathlineto{\pgfqpoint{1.400404in}{0.918452in}}%
\pgfpathclose%
\pgfusepath{fill}%
\end{pgfscope}%
\begin{pgfscope}%
\pgfpathrectangle{\pgfqpoint{0.407592in}{0.222904in}}{\pgfqpoint{1.411163in}{1.411163in}}%
\pgfusepath{clip}%
\pgfsetbuttcap%
\pgfsetroundjoin%
\definecolor{currentfill}{rgb}{0.749020,0.772549,0.792157}%
\pgfsetfillcolor{currentfill}%
\pgfsetfillopacity{0.160000}%
\pgfsetlinewidth{0.000000pt}%
\definecolor{currentstroke}{rgb}{0.000000,0.000000,0.000000}%
\pgfsetstrokecolor{currentstroke}%
\pgfsetdash{}{0pt}%
\pgfpathmoveto{\pgfqpoint{0.824610in}{0.905945in}}%
\pgfpathlineto{\pgfqpoint{0.812523in}{0.919388in}}%
\pgfpathlineto{\pgfqpoint{0.800434in}{0.932833in}}%
\pgfpathlineto{\pgfqpoint{0.791850in}{0.924924in}}%
\pgfpathlineto{\pgfqpoint{0.784513in}{0.916808in}}%
\pgfpathlineto{\pgfqpoint{0.797326in}{0.903972in}}%
\pgfpathlineto{\pgfqpoint{0.810133in}{0.891142in}}%
\pgfpathlineto{\pgfqpoint{0.816798in}{0.898638in}}%
\pgfpathlineto{\pgfqpoint{0.824610in}{0.905945in}}%
\pgfpathclose%
\pgfusepath{fill}%
\end{pgfscope}%
\begin{pgfscope}%
\pgfpathrectangle{\pgfqpoint{0.407592in}{0.222904in}}{\pgfqpoint{1.411163in}{1.411163in}}%
\pgfusepath{clip}%
\pgfsetbuttcap%
\pgfsetroundjoin%
\definecolor{currentfill}{rgb}{0.749020,0.772549,0.792157}%
\pgfsetfillcolor{currentfill}%
\pgfsetfillopacity{0.160000}%
\pgfsetlinewidth{0.000000pt}%
\definecolor{currentstroke}{rgb}{0.000000,0.000000,0.000000}%
\pgfsetstrokecolor{currentstroke}%
\pgfsetdash{}{0pt}%
\pgfpathmoveto{\pgfqpoint{1.426410in}{0.944268in}}%
\pgfpathlineto{\pgfqpoint{1.439422in}{0.957185in}}%
\pgfpathlineto{\pgfqpoint{1.452439in}{0.970108in}}%
\pgfpathlineto{\pgfqpoint{1.445501in}{0.979387in}}%
\pgfpathlineto{\pgfqpoint{1.437073in}{0.988413in}}%
\pgfpathlineto{\pgfqpoint{1.424676in}{0.974906in}}%
\pgfpathlineto{\pgfqpoint{1.412279in}{0.961401in}}%
\pgfpathlineto{\pgfqpoint{1.420039in}{0.952951in}}%
\pgfpathlineto{\pgfqpoint{1.426410in}{0.944268in}}%
\pgfpathclose%
\pgfusepath{fill}%
\end{pgfscope}%
\begin{pgfscope}%
\pgfpathrectangle{\pgfqpoint{0.407592in}{0.222904in}}{\pgfqpoint{1.411163in}{1.411163in}}%
\pgfusepath{clip}%
\pgfsetbuttcap%
\pgfsetroundjoin%
\definecolor{currentfill}{rgb}{0.749020,0.772549,0.792157}%
\pgfsetfillcolor{currentfill}%
\pgfsetfillopacity{0.160000}%
\pgfsetlinewidth{0.000000pt}%
\definecolor{currentstroke}{rgb}{0.000000,0.000000,0.000000}%
\pgfsetstrokecolor{currentstroke}%
\pgfsetdash{}{0pt}%
\pgfpathmoveto{\pgfqpoint{1.113937in}{0.628598in}}%
\pgfpathlineto{\pgfqpoint{1.125836in}{0.637665in}}%
\pgfpathlineto{\pgfqpoint{1.137766in}{0.646756in}}%
\pgfpathlineto{\pgfqpoint{1.139168in}{0.648174in}}%
\pgfpathlineto{\pgfqpoint{1.140354in}{0.649624in}}%
\pgfpathlineto{\pgfqpoint{1.127782in}{0.639820in}}%
\pgfpathlineto{\pgfqpoint{1.115238in}{0.630036in}}%
\pgfpathlineto{\pgfqpoint{1.114641in}{0.629309in}}%
\pgfpathlineto{\pgfqpoint{1.113937in}{0.628598in}}%
\pgfpathclose%
\pgfusepath{fill}%
\end{pgfscope}%
\begin{pgfscope}%
\pgfpathrectangle{\pgfqpoint{0.407592in}{0.222904in}}{\pgfqpoint{1.411163in}{1.411163in}}%
\pgfusepath{clip}%
\pgfsetbuttcap%
\pgfsetroundjoin%
\definecolor{currentfill}{rgb}{0.749020,0.772549,0.792157}%
\pgfsetfillcolor{currentfill}%
\pgfsetfillopacity{0.160000}%
\pgfsetlinewidth{0.000000pt}%
\definecolor{currentstroke}{rgb}{0.000000,0.000000,0.000000}%
\pgfsetstrokecolor{currentstroke}%
\pgfsetdash{}{0pt}%
\pgfpathmoveto{\pgfqpoint{1.064558in}{0.629889in}}%
\pgfpathlineto{\pgfqpoint{1.051679in}{0.639599in}}%
\pgfpathlineto{\pgfqpoint{1.038772in}{0.649331in}}%
\pgfpathlineto{\pgfqpoint{1.039706in}{0.647887in}}%
\pgfpathlineto{\pgfqpoint{1.040861in}{0.646476in}}%
\pgfpathlineto{\pgfqpoint{1.053252in}{0.637455in}}%
\pgfpathlineto{\pgfqpoint{1.065610in}{0.628457in}}%
\pgfpathlineto{\pgfqpoint{1.065029in}{0.629165in}}%
\pgfpathlineto{\pgfqpoint{1.064558in}{0.629889in}}%
\pgfpathclose%
\pgfusepath{fill}%
\end{pgfscope}%
\begin{pgfscope}%
\pgfpathrectangle{\pgfqpoint{0.407592in}{0.222904in}}{\pgfqpoint{1.411163in}{1.411163in}}%
\pgfusepath{clip}%
\pgfsetbuttcap%
\pgfsetroundjoin%
\definecolor{currentfill}{rgb}{0.749020,0.772549,0.792157}%
\pgfsetfillcolor{currentfill}%
\pgfsetfillopacity{0.160000}%
\pgfsetlinewidth{0.000000pt}%
\definecolor{currentstroke}{rgb}{0.000000,0.000000,0.000000}%
\pgfsetstrokecolor{currentstroke}%
\pgfsetdash{}{0pt}%
\pgfpathmoveto{\pgfqpoint{1.221828in}{0.731133in}}%
\pgfpathlineto{\pgfqpoint{1.235052in}{0.743252in}}%
\pgfpathlineto{\pgfqpoint{1.248288in}{0.755381in}}%
\pgfpathlineto{\pgfqpoint{1.246925in}{0.759754in}}%
\pgfpathlineto{\pgfqpoint{1.244882in}{0.764061in}}%
\pgfpathlineto{\pgfqpoint{1.231962in}{0.751234in}}%
\pgfpathlineto{\pgfqpoint{1.219047in}{0.738413in}}%
\pgfpathlineto{\pgfqpoint{1.220721in}{0.734800in}}%
\pgfpathlineto{\pgfqpoint{1.221828in}{0.731133in}}%
\pgfpathclose%
\pgfusepath{fill}%
\end{pgfscope}%
\begin{pgfscope}%
\pgfpathrectangle{\pgfqpoint{0.407592in}{0.222904in}}{\pgfqpoint{1.411163in}{1.411163in}}%
\pgfusepath{clip}%
\pgfsetbuttcap%
\pgfsetroundjoin%
\definecolor{currentfill}{rgb}{0.749020,0.772549,0.792157}%
\pgfsetfillcolor{currentfill}%
\pgfsetfillopacity{0.160000}%
\pgfsetlinewidth{0.000000pt}%
\definecolor{currentstroke}{rgb}{0.000000,0.000000,0.000000}%
\pgfsetstrokecolor{currentstroke}%
\pgfsetdash{}{0pt}%
\pgfpathmoveto{\pgfqpoint{0.800434in}{0.932833in}}%
\pgfpathlineto{\pgfqpoint{0.788344in}{0.946279in}}%
\pgfpathlineto{\pgfqpoint{0.776253in}{0.959726in}}%
\pgfpathlineto{\pgfqpoint{0.766887in}{0.951227in}}%
\pgfpathlineto{\pgfqpoint{0.758868in}{0.942501in}}%
\pgfpathlineto{\pgfqpoint{0.771694in}{0.929651in}}%
\pgfpathlineto{\pgfqpoint{0.784513in}{0.916808in}}%
\pgfpathlineto{\pgfqpoint{0.791850in}{0.924924in}}%
\pgfpathlineto{\pgfqpoint{0.800434in}{0.932833in}}%
\pgfpathclose%
\pgfusepath{fill}%
\end{pgfscope}%
\begin{pgfscope}%
\pgfpathrectangle{\pgfqpoint{0.407592in}{0.222904in}}{\pgfqpoint{1.411163in}{1.411163in}}%
\pgfusepath{clip}%
\pgfsetbuttcap%
\pgfsetroundjoin%
\definecolor{currentfill}{rgb}{0.749020,0.772549,0.792157}%
\pgfsetfillcolor{currentfill}%
\pgfsetfillopacity{0.160000}%
\pgfsetlinewidth{0.000000pt}%
\definecolor{currentstroke}{rgb}{0.000000,0.000000,0.000000}%
\pgfsetstrokecolor{currentstroke}%
\pgfsetdash{}{0pt}%
\pgfpathmoveto{\pgfqpoint{1.452439in}{0.970108in}}%
\pgfpathlineto{\pgfqpoint{1.465463in}{0.983037in}}%
\pgfpathlineto{\pgfqpoint{1.478494in}{0.995973in}}%
\pgfpathlineto{\pgfqpoint{1.470977in}{1.005837in}}%
\pgfpathlineto{\pgfqpoint{1.461869in}{1.015428in}}%
\pgfpathlineto{\pgfqpoint{1.449471in}{1.001920in}}%
\pgfpathlineto{\pgfqpoint{1.437073in}{0.988413in}}%
\pgfpathlineto{\pgfqpoint{1.445501in}{0.979387in}}%
\pgfpathlineto{\pgfqpoint{1.452439in}{0.970108in}}%
\pgfpathclose%
\pgfusepath{fill}%
\end{pgfscope}%
\begin{pgfscope}%
\pgfpathrectangle{\pgfqpoint{0.407592in}{0.222904in}}{\pgfqpoint{1.411163in}{1.411163in}}%
\pgfusepath{clip}%
\pgfsetbuttcap%
\pgfsetroundjoin%
\definecolor{currentfill}{rgb}{0.749020,0.772549,0.792157}%
\pgfsetfillcolor{currentfill}%
\pgfsetfillopacity{0.160000}%
\pgfsetlinewidth{0.000000pt}%
\definecolor{currentstroke}{rgb}{0.000000,0.000000,0.000000}%
\pgfsetstrokecolor{currentstroke}%
\pgfsetdash{}{0pt}%
\pgfpathmoveto{\pgfqpoint{1.142060in}{0.652597in}}%
\pgfpathlineto{\pgfqpoint{1.155074in}{0.663159in}}%
\pgfpathlineto{\pgfqpoint{1.168111in}{0.673739in}}%
\pgfpathlineto{\pgfqpoint{1.168865in}{0.675993in}}%
\pgfpathlineto{\pgfqpoint{1.169272in}{0.678255in}}%
\pgfpathlineto{\pgfqpoint{1.156055in}{0.666930in}}%
\pgfpathlineto{\pgfqpoint{1.142855in}{0.655621in}}%
\pgfpathlineto{\pgfqpoint{1.142573in}{0.654106in}}%
\pgfpathlineto{\pgfqpoint{1.142060in}{0.652597in}}%
\pgfpathclose%
\pgfusepath{fill}%
\end{pgfscope}%
\begin{pgfscope}%
\pgfpathrectangle{\pgfqpoint{0.407592in}{0.222904in}}{\pgfqpoint{1.411163in}{1.411163in}}%
\pgfusepath{clip}%
\pgfsetbuttcap%
\pgfsetroundjoin%
\definecolor{currentfill}{rgb}{0.749020,0.772549,0.792157}%
\pgfsetfillcolor{currentfill}%
\pgfsetfillopacity{0.160000}%
\pgfsetlinewidth{0.000000pt}%
\definecolor{currentstroke}{rgb}{0.000000,0.000000,0.000000}%
\pgfsetstrokecolor{currentstroke}%
\pgfsetdash{}{0pt}%
\pgfpathmoveto{\pgfqpoint{1.169272in}{0.678255in}}%
\pgfpathlineto{\pgfqpoint{1.182507in}{0.689594in}}%
\pgfpathlineto{\pgfqpoint{1.195759in}{0.700948in}}%
\pgfpathlineto{\pgfqpoint{1.195818in}{0.703945in}}%
\pgfpathlineto{\pgfqpoint{1.195414in}{0.706927in}}%
\pgfpathlineto{\pgfqpoint{1.182225in}{0.694840in}}%
\pgfpathlineto{\pgfqpoint{1.169048in}{0.682764in}}%
\pgfpathlineto{\pgfqpoint{1.169333in}{0.680515in}}%
\pgfpathlineto{\pgfqpoint{1.169272in}{0.678255in}}%
\pgfpathclose%
\pgfusepath{fill}%
\end{pgfscope}%
\begin{pgfscope}%
\pgfpathrectangle{\pgfqpoint{0.407592in}{0.222904in}}{\pgfqpoint{1.411163in}{1.411163in}}%
\pgfusepath{clip}%
\pgfsetbuttcap%
\pgfsetroundjoin%
\definecolor{currentfill}{rgb}{0.749020,0.772549,0.792157}%
\pgfsetfillcolor{currentfill}%
\pgfsetfillopacity{0.160000}%
\pgfsetlinewidth{0.000000pt}%
\definecolor{currentstroke}{rgb}{0.000000,0.000000,0.000000}%
\pgfsetstrokecolor{currentstroke}%
\pgfsetdash{}{0pt}%
\pgfpathmoveto{\pgfqpoint{1.037317in}{0.655317in}}%
\pgfpathlineto{\pgfqpoint{1.024044in}{0.666551in}}%
\pgfpathlineto{\pgfqpoint{1.010753in}{0.677801in}}%
\pgfpathlineto{\pgfqpoint{1.010763in}{0.675540in}}%
\pgfpathlineto{\pgfqpoint{1.011121in}{0.673289in}}%
\pgfpathlineto{\pgfqpoint{1.024364in}{0.662783in}}%
\pgfpathlineto{\pgfqpoint{1.037585in}{0.652295in}}%
\pgfpathlineto{\pgfqpoint{1.037335in}{0.653802in}}%
\pgfpathlineto{\pgfqpoint{1.037317in}{0.655317in}}%
\pgfpathclose%
\pgfusepath{fill}%
\end{pgfscope}%
\begin{pgfscope}%
\pgfpathrectangle{\pgfqpoint{0.407592in}{0.222904in}}{\pgfqpoint{1.411163in}{1.411163in}}%
\pgfusepath{clip}%
\pgfsetbuttcap%
\pgfsetroundjoin%
\definecolor{currentfill}{rgb}{0.749020,0.772549,0.792157}%
\pgfsetfillcolor{currentfill}%
\pgfsetfillopacity{0.160000}%
\pgfsetlinewidth{0.000000pt}%
\definecolor{currentstroke}{rgb}{0.000000,0.000000,0.000000}%
\pgfsetstrokecolor{currentstroke}%
\pgfsetdash{}{0pt}%
\pgfpathmoveto{\pgfqpoint{0.963301in}{0.737693in}}%
\pgfpathlineto{\pgfqpoint{0.950573in}{0.750445in}}%
\pgfpathlineto{\pgfqpoint{0.937838in}{0.763203in}}%
\pgfpathlineto{\pgfqpoint{0.935020in}{0.758881in}}%
\pgfpathlineto{\pgfqpoint{0.932871in}{0.754498in}}%
\pgfpathlineto{\pgfqpoint{0.946052in}{0.742439in}}%
\pgfpathlineto{\pgfqpoint{0.959220in}{0.730392in}}%
\pgfpathlineto{\pgfqpoint{0.960982in}{0.734068in}}%
\pgfpathlineto{\pgfqpoint{0.963301in}{0.737693in}}%
\pgfpathclose%
\pgfusepath{fill}%
\end{pgfscope}%
\begin{pgfscope}%
\pgfpathrectangle{\pgfqpoint{0.407592in}{0.222904in}}{\pgfqpoint{1.411163in}{1.411163in}}%
\pgfusepath{clip}%
\pgfsetbuttcap%
\pgfsetroundjoin%
\definecolor{currentfill}{rgb}{0.749020,0.772549,0.792157}%
\pgfsetfillcolor{currentfill}%
\pgfsetfillopacity{0.160000}%
\pgfsetlinewidth{0.000000pt}%
\definecolor{currentstroke}{rgb}{0.000000,0.000000,0.000000}%
\pgfsetstrokecolor{currentstroke}%
\pgfsetdash{}{0pt}%
\pgfpathmoveto{\pgfqpoint{1.011770in}{0.682314in}}%
\pgfpathlineto{\pgfqpoint{0.998651in}{0.694317in}}%
\pgfpathlineto{\pgfqpoint{0.985520in}{0.706331in}}%
\pgfpathlineto{\pgfqpoint{0.984588in}{0.703345in}}%
\pgfpathlineto{\pgfqpoint{0.984116in}{0.700345in}}%
\pgfpathlineto{\pgfqpoint{0.997444in}{0.689065in}}%
\pgfpathlineto{\pgfqpoint{1.010753in}{0.677801in}}%
\pgfpathlineto{\pgfqpoint{1.011089in}{0.680062in}}%
\pgfpathlineto{\pgfqpoint{1.011770in}{0.682314in}}%
\pgfpathclose%
\pgfusepath{fill}%
\end{pgfscope}%
\begin{pgfscope}%
\pgfpathrectangle{\pgfqpoint{0.407592in}{0.222904in}}{\pgfqpoint{1.411163in}{1.411163in}}%
\pgfusepath{clip}%
\pgfsetbuttcap%
\pgfsetroundjoin%
\definecolor{currentfill}{rgb}{0.749020,0.772549,0.792157}%
\pgfsetfillcolor{currentfill}%
\pgfsetfillopacity{0.160000}%
\pgfsetlinewidth{0.000000pt}%
\definecolor{currentstroke}{rgb}{0.000000,0.000000,0.000000}%
\pgfsetstrokecolor{currentstroke}%
\pgfsetdash{}{0pt}%
\pgfpathmoveto{\pgfqpoint{1.478494in}{0.995973in}}%
\pgfpathlineto{\pgfqpoint{1.491530in}{1.008914in}}%
\pgfpathlineto{\pgfqpoint{1.504572in}{1.021862in}}%
\pgfpathlineto{\pgfqpoint{1.496467in}{1.032301in}}%
\pgfpathlineto{\pgfqpoint{1.486669in}{1.042447in}}%
\pgfpathlineto{\pgfqpoint{1.474269in}{1.028937in}}%
\pgfpathlineto{\pgfqpoint{1.461869in}{1.015428in}}%
\pgfpathlineto{\pgfqpoint{1.470977in}{1.005837in}}%
\pgfpathlineto{\pgfqpoint{1.478494in}{0.995973in}}%
\pgfpathclose%
\pgfusepath{fill}%
\end{pgfscope}%
\begin{pgfscope}%
\pgfpathrectangle{\pgfqpoint{0.407592in}{0.222904in}}{\pgfqpoint{1.411163in}{1.411163in}}%
\pgfusepath{clip}%
\pgfsetbuttcap%
\pgfsetroundjoin%
\definecolor{currentfill}{rgb}{0.749020,0.772549,0.792157}%
\pgfsetfillcolor{currentfill}%
\pgfsetfillopacity{0.160000}%
\pgfsetlinewidth{0.000000pt}%
\definecolor{currentstroke}{rgb}{0.000000,0.000000,0.000000}%
\pgfsetstrokecolor{currentstroke}%
\pgfsetdash{}{0pt}%
\pgfpathmoveto{\pgfqpoint{0.776253in}{0.959726in}}%
\pgfpathlineto{\pgfqpoint{0.764161in}{0.973174in}}%
\pgfpathlineto{\pgfqpoint{0.752068in}{0.986624in}}%
\pgfpathlineto{\pgfqpoint{0.741910in}{0.977545in}}%
\pgfpathlineto{\pgfqpoint{0.733196in}{0.968219in}}%
\pgfpathlineto{\pgfqpoint{0.746035in}{0.955357in}}%
\pgfpathlineto{\pgfqpoint{0.758868in}{0.942501in}}%
\pgfpathlineto{\pgfqpoint{0.766887in}{0.951227in}}%
\pgfpathlineto{\pgfqpoint{0.776253in}{0.959726in}}%
\pgfpathclose%
\pgfusepath{fill}%
\end{pgfscope}%
\begin{pgfscope}%
\pgfpathrectangle{\pgfqpoint{0.407592in}{0.222904in}}{\pgfqpoint{1.411163in}{1.411163in}}%
\pgfusepath{clip}%
\pgfsetbuttcap%
\pgfsetroundjoin%
\definecolor{currentfill}{rgb}{0.749020,0.772549,0.792157}%
\pgfsetfillcolor{currentfill}%
\pgfsetfillopacity{0.160000}%
\pgfsetlinewidth{0.000000pt}%
\definecolor{currentstroke}{rgb}{0.000000,0.000000,0.000000}%
\pgfsetstrokecolor{currentstroke}%
\pgfsetdash{}{0pt}%
\pgfpathmoveto{\pgfqpoint{1.112219in}{0.627238in}}%
\pgfpathlineto{\pgfqpoint{1.123262in}{0.635628in}}%
\pgfpathlineto{\pgfqpoint{1.134338in}{0.644042in}}%
\pgfpathlineto{\pgfqpoint{1.136154in}{0.645377in}}%
\pgfpathlineto{\pgfqpoint{1.137766in}{0.646756in}}%
\pgfpathlineto{\pgfqpoint{1.125836in}{0.637665in}}%
\pgfpathlineto{\pgfqpoint{1.113937in}{0.628598in}}%
\pgfpathlineto{\pgfqpoint{1.113128in}{0.627907in}}%
\pgfpathlineto{\pgfqpoint{1.112219in}{0.627238in}}%
\pgfpathclose%
\pgfusepath{fill}%
\end{pgfscope}%
\begin{pgfscope}%
\pgfpathrectangle{\pgfqpoint{0.407592in}{0.222904in}}{\pgfqpoint{1.411163in}{1.411163in}}%
\pgfusepath{clip}%
\pgfsetbuttcap%
\pgfsetroundjoin%
\definecolor{currentfill}{rgb}{0.749020,0.772549,0.792157}%
\pgfsetfillcolor{currentfill}%
\pgfsetfillopacity{0.160000}%
\pgfsetlinewidth{0.000000pt}%
\definecolor{currentstroke}{rgb}{0.000000,0.000000,0.000000}%
\pgfsetstrokecolor{currentstroke}%
\pgfsetdash{}{0pt}%
\pgfpathmoveto{\pgfqpoint{1.065610in}{0.628457in}}%
\pgfpathlineto{\pgfqpoint{1.053252in}{0.637455in}}%
\pgfpathlineto{\pgfqpoint{1.040861in}{0.646476in}}%
\pgfpathlineto{\pgfqpoint{1.042234in}{0.645105in}}%
\pgfpathlineto{\pgfqpoint{1.043817in}{0.643781in}}%
\pgfpathlineto{\pgfqpoint{1.055472in}{0.635431in}}%
\pgfpathlineto{\pgfqpoint{1.067093in}{0.627107in}}%
\pgfpathlineto{\pgfqpoint{1.066299in}{0.627771in}}%
\pgfpathlineto{\pgfqpoint{1.065610in}{0.628457in}}%
\pgfpathclose%
\pgfusepath{fill}%
\end{pgfscope}%
\begin{pgfscope}%
\pgfpathrectangle{\pgfqpoint{0.407592in}{0.222904in}}{\pgfqpoint{1.411163in}{1.411163in}}%
\pgfusepath{clip}%
\pgfsetbuttcap%
\pgfsetroundjoin%
\definecolor{currentfill}{rgb}{0.749020,0.772549,0.792157}%
\pgfsetfillcolor{currentfill}%
\pgfsetfillopacity{0.160000}%
\pgfsetlinewidth{0.000000pt}%
\definecolor{currentstroke}{rgb}{0.000000,0.000000,0.000000}%
\pgfsetstrokecolor{currentstroke}%
\pgfsetdash{}{0pt}%
\pgfpathmoveto{\pgfqpoint{1.504572in}{1.021862in}}%
\pgfpathlineto{\pgfqpoint{1.517621in}{1.034815in}}%
\pgfpathlineto{\pgfqpoint{1.530675in}{1.047775in}}%
\pgfpathlineto{\pgfqpoint{1.521969in}{1.058779in}}%
\pgfpathlineto{\pgfqpoint{1.511471in}{1.069469in}}%
\pgfpathlineto{\pgfqpoint{1.499070in}{1.055958in}}%
\pgfpathlineto{\pgfqpoint{1.486669in}{1.042447in}}%
\pgfpathlineto{\pgfqpoint{1.496467in}{1.032301in}}%
\pgfpathlineto{\pgfqpoint{1.504572in}{1.021862in}}%
\pgfpathclose%
\pgfusepath{fill}%
\end{pgfscope}%
\begin{pgfscope}%
\pgfpathrectangle{\pgfqpoint{0.407592in}{0.222904in}}{\pgfqpoint{1.411163in}{1.411163in}}%
\pgfusepath{clip}%
\pgfsetbuttcap%
\pgfsetroundjoin%
\definecolor{currentfill}{rgb}{0.749020,0.772549,0.792157}%
\pgfsetfillcolor{currentfill}%
\pgfsetfillopacity{0.160000}%
\pgfsetlinewidth{0.000000pt}%
\definecolor{currentstroke}{rgb}{0.000000,0.000000,0.000000}%
\pgfsetstrokecolor{currentstroke}%
\pgfsetdash{}{0pt}%
\pgfpathmoveto{\pgfqpoint{1.248288in}{0.755381in}}%
\pgfpathlineto{\pgfqpoint{1.261536in}{0.767522in}}%
\pgfpathlineto{\pgfqpoint{1.274795in}{0.779674in}}%
\pgfpathlineto{\pgfqpoint{1.273165in}{0.784742in}}%
\pgfpathlineto{\pgfqpoint{1.270742in}{0.789733in}}%
\pgfpathlineto{\pgfqpoint{1.257809in}{0.776894in}}%
\pgfpathlineto{\pgfqpoint{1.244882in}{0.764061in}}%
\pgfpathlineto{\pgfqpoint{1.246925in}{0.759754in}}%
\pgfpathlineto{\pgfqpoint{1.248288in}{0.755381in}}%
\pgfpathclose%
\pgfusepath{fill}%
\end{pgfscope}%
\begin{pgfscope}%
\pgfpathrectangle{\pgfqpoint{0.407592in}{0.222904in}}{\pgfqpoint{1.411163in}{1.411163in}}%
\pgfusepath{clip}%
\pgfsetbuttcap%
\pgfsetroundjoin%
\definecolor{currentfill}{rgb}{0.749020,0.772549,0.792157}%
\pgfsetfillcolor{currentfill}%
\pgfsetfillopacity{0.160000}%
\pgfsetlinewidth{0.000000pt}%
\definecolor{currentstroke}{rgb}{0.000000,0.000000,0.000000}%
\pgfsetstrokecolor{currentstroke}%
\pgfsetdash{}{0pt}%
\pgfpathmoveto{\pgfqpoint{0.752068in}{0.986624in}}%
\pgfpathlineto{\pgfqpoint{0.739973in}{1.000075in}}%
\pgfpathlineto{\pgfqpoint{0.727877in}{1.013528in}}%
\pgfpathlineto{\pgfqpoint{0.716917in}{1.003879in}}%
\pgfpathlineto{\pgfqpoint{0.707498in}{0.993964in}}%
\pgfpathlineto{\pgfqpoint{0.720350in}{0.981088in}}%
\pgfpathlineto{\pgfqpoint{0.733196in}{0.968219in}}%
\pgfpathlineto{\pgfqpoint{0.741910in}{0.977545in}}%
\pgfpathlineto{\pgfqpoint{0.752068in}{0.986624in}}%
\pgfpathclose%
\pgfusepath{fill}%
\end{pgfscope}%
\begin{pgfscope}%
\pgfpathrectangle{\pgfqpoint{0.407592in}{0.222904in}}{\pgfqpoint{1.411163in}{1.411163in}}%
\pgfusepath{clip}%
\pgfsetbuttcap%
\pgfsetroundjoin%
\definecolor{currentfill}{rgb}{0.749020,0.772549,0.792157}%
\pgfsetfillcolor{currentfill}%
\pgfsetfillopacity{0.160000}%
\pgfsetlinewidth{0.000000pt}%
\definecolor{currentstroke}{rgb}{0.000000,0.000000,0.000000}%
\pgfsetstrokecolor{currentstroke}%
\pgfsetdash{}{0pt}%
\pgfpathmoveto{\pgfqpoint{1.111212in}{0.626595in}}%
\pgfpathlineto{\pgfqpoint{1.121753in}{0.634664in}}%
\pgfpathlineto{\pgfqpoint{1.132327in}{0.642759in}}%
\pgfpathlineto{\pgfqpoint{1.134338in}{0.644042in}}%
\pgfpathlineto{\pgfqpoint{1.123262in}{0.635628in}}%
\pgfpathlineto{\pgfqpoint{1.112219in}{0.627238in}}%
\pgfusepath{fill}%
\end{pgfscope}%
\begin{pgfscope}%
\pgfpathrectangle{\pgfqpoint{0.407592in}{0.222904in}}{\pgfqpoint{1.411163in}{1.411163in}}%
\pgfusepath{clip}%
\pgfsetbuttcap%
\pgfsetroundjoin%
\definecolor{currentfill}{rgb}{0.749020,0.772549,0.792157}%
\pgfsetfillcolor{currentfill}%
\pgfsetfillopacity{0.160000}%
\pgfsetlinewidth{0.000000pt}%
\definecolor{currentstroke}{rgb}{0.000000,0.000000,0.000000}%
\pgfsetstrokecolor{currentstroke}%
\pgfsetdash{}{0pt}%
\pgfpathmoveto{\pgfqpoint{1.530675in}{1.047775in}}%
\pgfpathlineto{\pgfqpoint{1.543736in}{1.060741in}}%
\pgfpathlineto{\pgfqpoint{1.556803in}{1.073712in}}%
\pgfpathlineto{\pgfqpoint{1.547486in}{1.085270in}}%
\pgfpathlineto{\pgfqpoint{1.536277in}{1.096494in}}%
\pgfpathlineto{\pgfqpoint{1.523874in}{1.082981in}}%
\pgfpathlineto{\pgfqpoint{1.511471in}{1.069469in}}%
\pgfpathlineto{\pgfqpoint{1.521969in}{1.058779in}}%
\pgfpathlineto{\pgfqpoint{1.530675in}{1.047775in}}%
\pgfpathclose%
\pgfusepath{fill}%
\end{pgfscope}%
\begin{pgfscope}%
\pgfpathrectangle{\pgfqpoint{0.407592in}{0.222904in}}{\pgfqpoint{1.411163in}{1.411163in}}%
\pgfusepath{clip}%
\pgfsetbuttcap%
\pgfsetroundjoin%
\definecolor{currentfill}{rgb}{0.749020,0.772549,0.792157}%
\pgfsetfillcolor{currentfill}%
\pgfsetfillopacity{0.160000}%
\pgfsetlinewidth{0.000000pt}%
\definecolor{currentstroke}{rgb}{0.000000,0.000000,0.000000}%
\pgfsetstrokecolor{currentstroke}%
\pgfsetdash{}{0pt}%
\pgfpathmoveto{\pgfqpoint{0.937838in}{0.763203in}}%
\pgfpathlineto{\pgfqpoint{0.925097in}{0.775968in}}%
\pgfpathlineto{\pgfqpoint{0.912349in}{0.788739in}}%
\pgfpathlineto{\pgfqpoint{0.909022in}{0.783731in}}%
\pgfpathlineto{\pgfqpoint{0.906473in}{0.778649in}}%
\pgfpathlineto{\pgfqpoint{0.919678in}{0.766568in}}%
\pgfpathlineto{\pgfqpoint{0.932871in}{0.754498in}}%
\pgfpathlineto{\pgfqpoint{0.935020in}{0.758881in}}%
\pgfpathlineto{\pgfqpoint{0.937838in}{0.763203in}}%
\pgfpathclose%
\pgfusepath{fill}%
\end{pgfscope}%
\begin{pgfscope}%
\pgfpathrectangle{\pgfqpoint{0.407592in}{0.222904in}}{\pgfqpoint{1.411163in}{1.411163in}}%
\pgfusepath{clip}%
\pgfsetbuttcap%
\pgfsetroundjoin%
\definecolor{currentfill}{rgb}{0.749020,0.772549,0.792157}%
\pgfsetfillcolor{currentfill}%
\pgfsetfillopacity{0.160000}%
\pgfsetlinewidth{0.000000pt}%
\definecolor{currentstroke}{rgb}{0.000000,0.000000,0.000000}%
\pgfsetstrokecolor{currentstroke}%
\pgfsetdash{}{0pt}%
\pgfpathmoveto{\pgfqpoint{0.727877in}{1.013528in}}%
\pgfpathlineto{\pgfqpoint{0.715780in}{1.026982in}}%
\pgfpathlineto{\pgfqpoint{0.703682in}{1.040437in}}%
\pgfpathlineto{\pgfqpoint{0.691909in}{1.030229in}}%
\pgfpathlineto{\pgfqpoint{0.681774in}{1.019736in}}%
\pgfpathlineto{\pgfqpoint{0.694639in}{1.006847in}}%
\pgfpathlineto{\pgfqpoint{0.707498in}{0.993964in}}%
\pgfpathlineto{\pgfqpoint{0.716917in}{1.003879in}}%
\pgfpathlineto{\pgfqpoint{0.727877in}{1.013528in}}%
\pgfpathclose%
\pgfusepath{fill}%
\end{pgfscope}%
\begin{pgfscope}%
\pgfpathrectangle{\pgfqpoint{0.407592in}{0.222904in}}{\pgfqpoint{1.411163in}{1.411163in}}%
\pgfusepath{clip}%
\pgfsetbuttcap%
\pgfsetroundjoin%
\definecolor{currentfill}{rgb}{0.749020,0.772549,0.792157}%
\pgfsetfillcolor{currentfill}%
\pgfsetfillopacity{0.160000}%
\pgfsetlinewidth{0.000000pt}%
\definecolor{currentstroke}{rgb}{0.000000,0.000000,0.000000}%
\pgfsetstrokecolor{currentstroke}%
\pgfsetdash{}{0pt}%
\pgfpathmoveto{\pgfqpoint{1.067093in}{0.627107in}}%
\pgfpathlineto{\pgfqpoint{1.055472in}{0.635431in}}%
\pgfpathlineto{\pgfqpoint{1.043817in}{0.643781in}}%
\pgfpathlineto{\pgfqpoint{1.045605in}{0.642508in}}%
\pgfpathlineto{\pgfqpoint{1.047590in}{0.641293in}}%
\pgfpathlineto{\pgfqpoint{1.058303in}{0.633564in}}%
\pgfpathlineto{\pgfqpoint{1.068981in}{0.625862in}}%
\pgfpathlineto{\pgfqpoint{1.067988in}{0.626470in}}%
\pgfpathlineto{\pgfqpoint{1.067093in}{0.627107in}}%
\pgfpathclose%
\pgfusepath{fill}%
\end{pgfscope}%
\begin{pgfscope}%
\pgfpathrectangle{\pgfqpoint{0.407592in}{0.222904in}}{\pgfqpoint{1.411163in}{1.411163in}}%
\pgfusepath{clip}%
\pgfsetbuttcap%
\pgfsetroundjoin%
\definecolor{currentfill}{rgb}{0.749020,0.772549,0.792157}%
\pgfsetfillcolor{currentfill}%
\pgfsetfillopacity{0.160000}%
\pgfsetlinewidth{0.000000pt}%
\definecolor{currentstroke}{rgb}{0.000000,0.000000,0.000000}%
\pgfsetstrokecolor{currentstroke}%
\pgfsetdash{}{0pt}%
\pgfpathmoveto{\pgfqpoint{1.556803in}{1.073712in}}%
\pgfpathlineto{\pgfqpoint{1.569876in}{1.086690in}}%
\pgfpathlineto{\pgfqpoint{1.582955in}{1.099674in}}%
\pgfpathlineto{\pgfqpoint{1.573015in}{1.111776in}}%
\pgfpathlineto{\pgfqpoint{1.561085in}{1.123522in}}%
\pgfpathlineto{\pgfqpoint{1.548681in}{1.110008in}}%
\pgfpathlineto{\pgfqpoint{1.536277in}{1.096494in}}%
\pgfpathlineto{\pgfqpoint{1.547486in}{1.085270in}}%
\pgfpathlineto{\pgfqpoint{1.556803in}{1.073712in}}%
\pgfpathclose%
\pgfusepath{fill}%
\end{pgfscope}%
\begin{pgfscope}%
\pgfpathrectangle{\pgfqpoint{0.407592in}{0.222904in}}{\pgfqpoint{1.411163in}{1.411163in}}%
\pgfusepath{clip}%
\pgfsetbuttcap%
\pgfsetroundjoin%
\definecolor{currentfill}{rgb}{0.749020,0.772549,0.792157}%
\pgfsetfillcolor{currentfill}%
\pgfsetfillopacity{0.160000}%
\pgfsetlinewidth{0.000000pt}%
\definecolor{currentstroke}{rgb}{0.000000,0.000000,0.000000}%
\pgfsetstrokecolor{currentstroke}%
\pgfsetdash{}{0pt}%
\pgfpathmoveto{\pgfqpoint{1.140354in}{0.649624in}}%
\pgfpathlineto{\pgfqpoint{1.152954in}{0.659450in}}%
\pgfpathlineto{\pgfqpoint{1.165580in}{0.669297in}}%
\pgfpathlineto{\pgfqpoint{1.167015in}{0.671504in}}%
\pgfpathlineto{\pgfqpoint{1.168111in}{0.673739in}}%
\pgfpathlineto{\pgfqpoint{1.155074in}{0.663159in}}%
\pgfpathlineto{\pgfqpoint{1.142060in}{0.652597in}}%
\pgfpathlineto{\pgfqpoint{1.141320in}{0.651101in}}%
\pgfpathlineto{\pgfqpoint{1.140354in}{0.649624in}}%
\pgfpathclose%
\pgfusepath{fill}%
\end{pgfscope}%
\begin{pgfscope}%
\pgfpathrectangle{\pgfqpoint{0.407592in}{0.222904in}}{\pgfqpoint{1.411163in}{1.411163in}}%
\pgfusepath{clip}%
\pgfsetbuttcap%
\pgfsetroundjoin%
\definecolor{currentfill}{rgb}{0.749020,0.772549,0.792157}%
\pgfsetfillcolor{currentfill}%
\pgfsetfillopacity{0.160000}%
\pgfsetlinewidth{0.000000pt}%
\definecolor{currentstroke}{rgb}{0.000000,0.000000,0.000000}%
\pgfsetstrokecolor{currentstroke}%
\pgfsetdash{}{0pt}%
\pgfpathmoveto{\pgfqpoint{1.108927in}{0.625398in}}%
\pgfpathlineto{\pgfqpoint{1.118325in}{0.632869in}}%
\pgfpathlineto{\pgfqpoint{1.127755in}{0.640367in}}%
\pgfpathlineto{\pgfqpoint{1.130130in}{0.641532in}}%
\pgfpathlineto{\pgfqpoint{1.132327in}{0.642759in}}%
\pgfpathlineto{\pgfqpoint{1.121753in}{0.634664in}}%
\pgfpathlineto{\pgfqpoint{1.111212in}{0.626595in}}%
\pgfpathlineto{\pgfqpoint{1.110114in}{0.625981in}}%
\pgfpathlineto{\pgfqpoint{1.108927in}{0.625398in}}%
\pgfpathclose%
\pgfusepath{fill}%
\end{pgfscope}%
\begin{pgfscope}%
\pgfpathrectangle{\pgfqpoint{0.407592in}{0.222904in}}{\pgfqpoint{1.411163in}{1.411163in}}%
\pgfusepath{clip}%
\pgfsetbuttcap%
\pgfsetroundjoin%
\definecolor{currentfill}{rgb}{0.749020,0.772549,0.792157}%
\pgfsetfillcolor{currentfill}%
\pgfsetfillopacity{0.160000}%
\pgfsetlinewidth{0.000000pt}%
\definecolor{currentstroke}{rgb}{0.000000,0.000000,0.000000}%
\pgfsetstrokecolor{currentstroke}%
\pgfsetdash{}{0pt}%
\pgfpathmoveto{\pgfqpoint{1.195759in}{0.700948in}}%
\pgfpathlineto{\pgfqpoint{1.209029in}{0.712317in}}%
\pgfpathlineto{\pgfqpoint{1.222316in}{0.723701in}}%
\pgfpathlineto{\pgfqpoint{1.222361in}{0.727428in}}%
\pgfpathlineto{\pgfqpoint{1.221828in}{0.731133in}}%
\pgfpathlineto{\pgfqpoint{1.208615in}{0.719024in}}%
\pgfpathlineto{\pgfqpoint{1.195414in}{0.706927in}}%
\pgfpathlineto{\pgfqpoint{1.195818in}{0.703945in}}%
\pgfpathlineto{\pgfqpoint{1.195759in}{0.700948in}}%
\pgfpathclose%
\pgfusepath{fill}%
\end{pgfscope}%
\begin{pgfscope}%
\pgfpathrectangle{\pgfqpoint{0.407592in}{0.222904in}}{\pgfqpoint{1.411163in}{1.411163in}}%
\pgfusepath{clip}%
\pgfsetbuttcap%
\pgfsetroundjoin%
\definecolor{currentfill}{rgb}{0.749020,0.772549,0.792157}%
\pgfsetfillcolor{currentfill}%
\pgfsetfillopacity{0.160000}%
\pgfsetlinewidth{0.000000pt}%
\definecolor{currentstroke}{rgb}{0.000000,0.000000,0.000000}%
\pgfsetstrokecolor{currentstroke}%
\pgfsetdash{}{0pt}%
\pgfpathmoveto{\pgfqpoint{1.037585in}{0.652295in}}%
\pgfpathlineto{\pgfqpoint{1.024364in}{0.662783in}}%
\pgfpathlineto{\pgfqpoint{1.011121in}{0.673289in}}%
\pgfpathlineto{\pgfqpoint{1.011825in}{0.671059in}}%
\pgfpathlineto{\pgfqpoint{1.012872in}{0.668859in}}%
\pgfpathlineto{\pgfqpoint{1.025837in}{0.659084in}}%
\pgfpathlineto{\pgfqpoint{1.038772in}{0.649331in}}%
\pgfpathlineto{\pgfqpoint{1.038064in}{0.650803in}}%
\pgfpathlineto{\pgfqpoint{1.037585in}{0.652295in}}%
\pgfpathclose%
\pgfusepath{fill}%
\end{pgfscope}%
\begin{pgfscope}%
\pgfpathrectangle{\pgfqpoint{0.407592in}{0.222904in}}{\pgfqpoint{1.411163in}{1.411163in}}%
\pgfusepath{clip}%
\pgfsetbuttcap%
\pgfsetroundjoin%
\definecolor{currentfill}{rgb}{0.749020,0.772549,0.792157}%
\pgfsetfillcolor{currentfill}%
\pgfsetfillopacity{0.160000}%
\pgfsetlinewidth{0.000000pt}%
\definecolor{currentstroke}{rgb}{0.000000,0.000000,0.000000}%
\pgfsetstrokecolor{currentstroke}%
\pgfsetdash{}{0pt}%
\pgfpathmoveto{\pgfqpoint{0.703682in}{1.040437in}}%
\pgfpathlineto{\pgfqpoint{0.691583in}{1.053893in}}%
\pgfpathlineto{\pgfqpoint{0.679482in}{1.067351in}}%
\pgfpathlineto{\pgfqpoint{0.666885in}{1.056595in}}%
\pgfpathlineto{\pgfqpoint{0.656023in}{1.045533in}}%
\pgfpathlineto{\pgfqpoint{0.668902in}{1.032631in}}%
\pgfpathlineto{\pgfqpoint{0.681774in}{1.019736in}}%
\pgfpathlineto{\pgfqpoint{0.691909in}{1.030229in}}%
\pgfpathlineto{\pgfqpoint{0.703682in}{1.040437in}}%
\pgfpathclose%
\pgfusepath{fill}%
\end{pgfscope}%
\begin{pgfscope}%
\pgfpathrectangle{\pgfqpoint{0.407592in}{0.222904in}}{\pgfqpoint{1.411163in}{1.411163in}}%
\pgfusepath{clip}%
\pgfsetbuttcap%
\pgfsetroundjoin%
\definecolor{currentfill}{rgb}{0.749020,0.772549,0.792157}%
\pgfsetfillcolor{currentfill}%
\pgfsetfillopacity{0.160000}%
\pgfsetlinewidth{0.000000pt}%
\definecolor{currentstroke}{rgb}{0.000000,0.000000,0.000000}%
\pgfsetstrokecolor{currentstroke}%
\pgfsetdash{}{0pt}%
\pgfpathmoveto{\pgfqpoint{1.274795in}{0.779674in}}%
\pgfpathlineto{\pgfqpoint{1.288067in}{0.791836in}}%
\pgfpathlineto{\pgfqpoint{1.301350in}{0.804009in}}%
\pgfpathlineto{\pgfqpoint{1.299440in}{0.809764in}}%
\pgfpathlineto{\pgfqpoint{1.296626in}{0.815428in}}%
\pgfpathlineto{\pgfqpoint{1.283681in}{0.802577in}}%
\pgfpathlineto{\pgfqpoint{1.270742in}{0.789733in}}%
\pgfpathlineto{\pgfqpoint{1.273165in}{0.784742in}}%
\pgfpathlineto{\pgfqpoint{1.274795in}{0.779674in}}%
\pgfpathclose%
\pgfusepath{fill}%
\end{pgfscope}%
\begin{pgfscope}%
\pgfpathrectangle{\pgfqpoint{0.407592in}{0.222904in}}{\pgfqpoint{1.411163in}{1.411163in}}%
\pgfusepath{clip}%
\pgfsetbuttcap%
\pgfsetroundjoin%
\definecolor{currentfill}{rgb}{0.749020,0.772549,0.792157}%
\pgfsetfillcolor{currentfill}%
\pgfsetfillopacity{0.160000}%
\pgfsetlinewidth{0.000000pt}%
\definecolor{currentstroke}{rgb}{0.000000,0.000000,0.000000}%
\pgfsetstrokecolor{currentstroke}%
\pgfsetdash{}{0pt}%
\pgfpathmoveto{\pgfqpoint{1.582955in}{1.099674in}}%
\pgfpathlineto{\pgfqpoint{1.596040in}{1.112664in}}%
\pgfpathlineto{\pgfqpoint{1.609132in}{1.125661in}}%
\pgfpathlineto{\pgfqpoint{1.598558in}{1.138295in}}%
\pgfpathlineto{\pgfqpoint{1.585897in}{1.150554in}}%
\pgfpathlineto{\pgfqpoint{1.573491in}{1.137038in}}%
\pgfpathlineto{\pgfqpoint{1.561085in}{1.123522in}}%
\pgfpathlineto{\pgfqpoint{1.573015in}{1.111776in}}%
\pgfpathlineto{\pgfqpoint{1.582955in}{1.099674in}}%
\pgfpathclose%
\pgfusepath{fill}%
\end{pgfscope}%
\begin{pgfscope}%
\pgfpathrectangle{\pgfqpoint{0.407592in}{0.222904in}}{\pgfqpoint{1.411163in}{1.411163in}}%
\pgfusepath{clip}%
\pgfsetbuttcap%
\pgfsetroundjoin%
\definecolor{currentfill}{rgb}{0.749020,0.772549,0.792157}%
\pgfsetfillcolor{currentfill}%
\pgfsetfillopacity{0.160000}%
\pgfsetlinewidth{0.000000pt}%
\definecolor{currentstroke}{rgb}{0.000000,0.000000,0.000000}%
\pgfsetstrokecolor{currentstroke}%
\pgfsetdash{}{0pt}%
\pgfpathmoveto{\pgfqpoint{0.985520in}{0.706331in}}%
\pgfpathlineto{\pgfqpoint{0.972376in}{0.718356in}}%
\pgfpathlineto{\pgfqpoint{0.959220in}{0.730392in}}%
\pgfpathlineto{\pgfqpoint{0.958026in}{0.726681in}}%
\pgfpathlineto{\pgfqpoint{0.957406in}{0.722952in}}%
\pgfpathlineto{\pgfqpoint{0.970770in}{0.711641in}}%
\pgfpathlineto{\pgfqpoint{0.984116in}{0.700345in}}%
\pgfpathlineto{\pgfqpoint{0.984588in}{0.703345in}}%
\pgfpathlineto{\pgfqpoint{0.985520in}{0.706331in}}%
\pgfpathclose%
\pgfusepath{fill}%
\end{pgfscope}%
\begin{pgfscope}%
\pgfpathrectangle{\pgfqpoint{0.407592in}{0.222904in}}{\pgfqpoint{1.411163in}{1.411163in}}%
\pgfusepath{clip}%
\pgfsetbuttcap%
\pgfsetroundjoin%
\definecolor{currentfill}{rgb}{0.749020,0.772549,0.792157}%
\pgfsetfillcolor{currentfill}%
\pgfsetfillopacity{0.160000}%
\pgfsetlinewidth{0.000000pt}%
\definecolor{currentstroke}{rgb}{0.000000,0.000000,0.000000}%
\pgfsetstrokecolor{currentstroke}%
\pgfsetdash{}{0pt}%
\pgfpathmoveto{\pgfqpoint{1.068981in}{0.625862in}}%
\pgfpathlineto{\pgfqpoint{1.058303in}{0.633564in}}%
\pgfpathlineto{\pgfqpoint{1.047590in}{0.641293in}}%
\pgfpathlineto{\pgfqpoint{1.049763in}{0.640141in}}%
\pgfpathlineto{\pgfqpoint{1.052114in}{0.639058in}}%
\pgfpathlineto{\pgfqpoint{1.061695in}{0.631887in}}%
\pgfpathlineto{\pgfqpoint{1.071242in}{0.624743in}}%
\pgfpathlineto{\pgfqpoint{1.070067in}{0.625285in}}%
\pgfpathlineto{\pgfqpoint{1.068981in}{0.625862in}}%
\pgfpathclose%
\pgfusepath{fill}%
\end{pgfscope}%
\begin{pgfscope}%
\pgfpathrectangle{\pgfqpoint{0.407592in}{0.222904in}}{\pgfqpoint{1.411163in}{1.411163in}}%
\pgfusepath{clip}%
\pgfsetbuttcap%
\pgfsetroundjoin%
\definecolor{currentfill}{rgb}{0.749020,0.772549,0.792157}%
\pgfsetfillcolor{currentfill}%
\pgfsetfillopacity{0.160000}%
\pgfsetlinewidth{0.000000pt}%
\definecolor{currentstroke}{rgb}{0.000000,0.000000,0.000000}%
\pgfsetstrokecolor{currentstroke}%
\pgfsetdash{}{0pt}%
\pgfpathmoveto{\pgfqpoint{0.679482in}{1.067351in}}%
\pgfpathlineto{\pgfqpoint{0.667380in}{1.080811in}}%
\pgfpathlineto{\pgfqpoint{0.655277in}{1.094271in}}%
\pgfpathlineto{\pgfqpoint{0.641847in}{1.082977in}}%
\pgfpathlineto{\pgfqpoint{0.630246in}{1.071357in}}%
\pgfpathlineto{\pgfqpoint{0.643138in}{1.058442in}}%
\pgfpathlineto{\pgfqpoint{0.656023in}{1.045533in}}%
\pgfpathlineto{\pgfqpoint{0.666885in}{1.056595in}}%
\pgfpathlineto{\pgfqpoint{0.679482in}{1.067351in}}%
\pgfpathclose%
\pgfusepath{fill}%
\end{pgfscope}%
\begin{pgfscope}%
\pgfpathrectangle{\pgfqpoint{0.407592in}{0.222904in}}{\pgfqpoint{1.411163in}{1.411163in}}%
\pgfusepath{clip}%
\pgfsetbuttcap%
\pgfsetroundjoin%
\definecolor{currentfill}{rgb}{0.749020,0.772549,0.792157}%
\pgfsetfillcolor{currentfill}%
\pgfsetfillopacity{0.160000}%
\pgfsetlinewidth{0.000000pt}%
\definecolor{currentstroke}{rgb}{0.000000,0.000000,0.000000}%
\pgfsetstrokecolor{currentstroke}%
\pgfsetdash{}{0pt}%
\pgfpathmoveto{\pgfqpoint{1.168111in}{0.673739in}}%
\pgfpathlineto{\pgfqpoint{1.181171in}{0.684338in}}%
\pgfpathlineto{\pgfqpoint{1.194254in}{0.694956in}}%
\pgfpathlineto{\pgfqpoint{1.195237in}{0.697947in}}%
\pgfpathlineto{\pgfqpoint{1.195759in}{0.700948in}}%
\pgfpathlineto{\pgfqpoint{1.182507in}{0.689594in}}%
\pgfpathlineto{\pgfqpoint{1.169272in}{0.678255in}}%
\pgfpathlineto{\pgfqpoint{1.168865in}{0.675993in}}%
\pgfpathlineto{\pgfqpoint{1.168111in}{0.673739in}}%
\pgfpathclose%
\pgfusepath{fill}%
\end{pgfscope}%
\begin{pgfscope}%
\pgfpathrectangle{\pgfqpoint{0.407592in}{0.222904in}}{\pgfqpoint{1.411163in}{1.411163in}}%
\pgfusepath{clip}%
\pgfsetbuttcap%
\pgfsetroundjoin%
\definecolor{currentfill}{rgb}{0.749020,0.772549,0.792157}%
\pgfsetfillcolor{currentfill}%
\pgfsetfillopacity{0.160000}%
\pgfsetlinewidth{0.000000pt}%
\definecolor{currentstroke}{rgb}{0.000000,0.000000,0.000000}%
\pgfsetstrokecolor{currentstroke}%
\pgfsetdash{}{0pt}%
\pgfpathmoveto{\pgfqpoint{1.609132in}{1.125661in}}%
\pgfpathlineto{\pgfqpoint{1.622229in}{1.138663in}}%
\pgfpathlineto{\pgfqpoint{1.635333in}{1.151671in}}%
\pgfpathlineto{\pgfqpoint{1.624114in}{1.164828in}}%
\pgfpathlineto{\pgfqpoint{1.610711in}{1.177589in}}%
\pgfpathlineto{\pgfqpoint{1.598303in}{1.164071in}}%
\pgfpathlineto{\pgfqpoint{1.585897in}{1.150554in}}%
\pgfpathlineto{\pgfqpoint{1.598558in}{1.138295in}}%
\pgfpathlineto{\pgfqpoint{1.609132in}{1.125661in}}%
\pgfpathclose%
\pgfusepath{fill}%
\end{pgfscope}%
\begin{pgfscope}%
\pgfpathrectangle{\pgfqpoint{0.407592in}{0.222904in}}{\pgfqpoint{1.411163in}{1.411163in}}%
\pgfusepath{clip}%
\pgfsetbuttcap%
\pgfsetroundjoin%
\definecolor{currentfill}{rgb}{0.749020,0.772549,0.792157}%
\pgfsetfillcolor{currentfill}%
\pgfsetfillopacity{0.160000}%
\pgfsetlinewidth{0.000000pt}%
\definecolor{currentstroke}{rgb}{0.000000,0.000000,0.000000}%
\pgfsetstrokecolor{currentstroke}%
\pgfsetdash{}{0pt}%
\pgfpathmoveto{\pgfqpoint{0.912349in}{0.788739in}}%
\pgfpathlineto{\pgfqpoint{0.899595in}{0.801516in}}%
\pgfpathlineto{\pgfqpoint{0.886834in}{0.814300in}}%
\pgfpathlineto{\pgfqpoint{0.882986in}{0.808616in}}%
\pgfpathlineto{\pgfqpoint{0.880026in}{0.802846in}}%
\pgfpathlineto{\pgfqpoint{0.893256in}{0.790742in}}%
\pgfpathlineto{\pgfqpoint{0.906473in}{0.778649in}}%
\pgfpathlineto{\pgfqpoint{0.909022in}{0.783731in}}%
\pgfpathlineto{\pgfqpoint{0.912349in}{0.788739in}}%
\pgfpathclose%
\pgfusepath{fill}%
\end{pgfscope}%
\begin{pgfscope}%
\pgfpathrectangle{\pgfqpoint{0.407592in}{0.222904in}}{\pgfqpoint{1.411163in}{1.411163in}}%
\pgfusepath{clip}%
\pgfsetbuttcap%
\pgfsetroundjoin%
\definecolor{currentfill}{rgb}{0.749020,0.772549,0.792157}%
\pgfsetfillcolor{currentfill}%
\pgfsetfillopacity{0.160000}%
\pgfsetlinewidth{0.000000pt}%
\definecolor{currentstroke}{rgb}{0.000000,0.000000,0.000000}%
\pgfsetstrokecolor{currentstroke}%
\pgfsetdash{}{0pt}%
\pgfpathmoveto{\pgfqpoint{1.106312in}{0.624336in}}%
\pgfpathlineto{\pgfqpoint{1.114399in}{0.631277in}}%
\pgfpathlineto{\pgfqpoint{1.122516in}{0.638244in}}%
\pgfpathlineto{\pgfqpoint{1.125213in}{0.639269in}}%
\pgfpathlineto{\pgfqpoint{1.127755in}{0.640367in}}%
\pgfpathlineto{\pgfqpoint{1.118325in}{0.632869in}}%
\pgfpathlineto{\pgfqpoint{1.108927in}{0.625398in}}%
\pgfpathlineto{\pgfqpoint{1.107658in}{0.624849in}}%
\pgfpathlineto{\pgfqpoint{1.106312in}{0.624336in}}%
\pgfpathclose%
\pgfusepath{fill}%
\end{pgfscope}%
\begin{pgfscope}%
\pgfpathrectangle{\pgfqpoint{0.407592in}{0.222904in}}{\pgfqpoint{1.411163in}{1.411163in}}%
\pgfusepath{clip}%
\pgfsetbuttcap%
\pgfsetroundjoin%
\definecolor{currentfill}{rgb}{0.749020,0.772549,0.792157}%
\pgfsetfillcolor{currentfill}%
\pgfsetfillopacity{0.160000}%
\pgfsetlinewidth{0.000000pt}%
\definecolor{currentstroke}{rgb}{0.000000,0.000000,0.000000}%
\pgfsetstrokecolor{currentstroke}%
\pgfsetdash{}{0pt}%
\pgfpathmoveto{\pgfqpoint{1.010753in}{0.677801in}}%
\pgfpathlineto{\pgfqpoint{0.997444in}{0.689065in}}%
\pgfpathlineto{\pgfqpoint{0.984116in}{0.700345in}}%
\pgfpathlineto{\pgfqpoint{0.984107in}{0.697345in}}%
\pgfpathlineto{\pgfqpoint{0.984562in}{0.694358in}}%
\pgfpathlineto{\pgfqpoint{0.997853in}{0.683814in}}%
\pgfpathlineto{\pgfqpoint{1.011121in}{0.673289in}}%
\pgfpathlineto{\pgfqpoint{1.010763in}{0.675540in}}%
\pgfpathlineto{\pgfqpoint{1.010753in}{0.677801in}}%
\pgfpathclose%
\pgfusepath{fill}%
\end{pgfscope}%
\begin{pgfscope}%
\pgfpathrectangle{\pgfqpoint{0.407592in}{0.222904in}}{\pgfqpoint{1.411163in}{1.411163in}}%
\pgfusepath{clip}%
\pgfsetbuttcap%
\pgfsetroundjoin%
\definecolor{currentfill}{rgb}{0.749020,0.772549,0.792157}%
\pgfsetfillcolor{currentfill}%
\pgfsetfillopacity{0.160000}%
\pgfsetlinewidth{0.000000pt}%
\definecolor{currentstroke}{rgb}{0.000000,0.000000,0.000000}%
\pgfsetstrokecolor{currentstroke}%
\pgfsetdash{}{0pt}%
\pgfpathmoveto{\pgfqpoint{0.655277in}{1.094271in}}%
\pgfpathlineto{\pgfqpoint{0.643173in}{1.107733in}}%
\pgfpathlineto{\pgfqpoint{0.631067in}{1.121196in}}%
\pgfpathlineto{\pgfqpoint{0.616794in}{1.109375in}}%
\pgfpathlineto{\pgfqpoint{0.604443in}{1.097208in}}%
\pgfpathlineto{\pgfqpoint{0.617348in}{1.084279in}}%
\pgfpathlineto{\pgfqpoint{0.630246in}{1.071357in}}%
\pgfpathlineto{\pgfqpoint{0.641847in}{1.082977in}}%
\pgfpathlineto{\pgfqpoint{0.655277in}{1.094271in}}%
\pgfpathclose%
\pgfusepath{fill}%
\end{pgfscope}%
\begin{pgfscope}%
\pgfpathrectangle{\pgfqpoint{0.407592in}{0.222904in}}{\pgfqpoint{1.411163in}{1.411163in}}%
\pgfusepath{clip}%
\pgfsetbuttcap%
\pgfsetroundjoin%
\definecolor{currentfill}{rgb}{0.749020,0.772549,0.792157}%
\pgfsetfillcolor{currentfill}%
\pgfsetfillopacity{0.160000}%
\pgfsetlinewidth{0.000000pt}%
\definecolor{currentstroke}{rgb}{0.000000,0.000000,0.000000}%
\pgfsetstrokecolor{currentstroke}%
\pgfsetdash{}{0pt}%
\pgfpathmoveto{\pgfqpoint{1.635333in}{1.151671in}}%
\pgfpathlineto{\pgfqpoint{1.648443in}{1.164686in}}%
\pgfpathlineto{\pgfqpoint{1.661559in}{1.177707in}}%
\pgfpathlineto{\pgfqpoint{1.649683in}{1.191375in}}%
\pgfpathlineto{\pgfqpoint{1.635528in}{1.204627in}}%
\pgfpathlineto{\pgfqpoint{1.623119in}{1.191107in}}%
\pgfpathlineto{\pgfqpoint{1.610711in}{1.177589in}}%
\pgfpathlineto{\pgfqpoint{1.624114in}{1.164828in}}%
\pgfpathlineto{\pgfqpoint{1.635333in}{1.151671in}}%
\pgfpathclose%
\pgfusepath{fill}%
\end{pgfscope}%
\begin{pgfscope}%
\pgfpathrectangle{\pgfqpoint{0.407592in}{0.222904in}}{\pgfqpoint{1.411163in}{1.411163in}}%
\pgfusepath{clip}%
\pgfsetbuttcap%
\pgfsetroundjoin%
\definecolor{currentfill}{rgb}{0.749020,0.772549,0.792157}%
\pgfsetfillcolor{currentfill}%
\pgfsetfillopacity{0.160000}%
\pgfsetlinewidth{0.000000pt}%
\definecolor{currentstroke}{rgb}{0.000000,0.000000,0.000000}%
\pgfsetstrokecolor{currentstroke}%
\pgfsetdash{}{0pt}%
\pgfpathmoveto{\pgfqpoint{1.071242in}{0.624743in}}%
\pgfpathlineto{\pgfqpoint{1.061695in}{0.631887in}}%
\pgfpathlineto{\pgfqpoint{1.052114in}{0.639058in}}%
\pgfpathlineto{\pgfqpoint{1.054634in}{0.638047in}}%
\pgfpathlineto{\pgfqpoint{1.057311in}{0.637114in}}%
\pgfpathlineto{\pgfqpoint{1.065590in}{0.630430in}}%
\pgfpathlineto{\pgfqpoint{1.073835in}{0.623771in}}%
\pgfpathlineto{\pgfqpoint{1.072499in}{0.624238in}}%
\pgfpathlineto{\pgfqpoint{1.071242in}{0.624743in}}%
\pgfpathclose%
\pgfusepath{fill}%
\end{pgfscope}%
\begin{pgfscope}%
\pgfpathrectangle{\pgfqpoint{0.407592in}{0.222904in}}{\pgfqpoint{1.411163in}{1.411163in}}%
\pgfusepath{clip}%
\pgfsetbuttcap%
\pgfsetroundjoin%
\definecolor{currentfill}{rgb}{0.749020,0.772549,0.792157}%
\pgfsetfillcolor{currentfill}%
\pgfsetfillopacity{0.160000}%
\pgfsetlinewidth{0.000000pt}%
\definecolor{currentstroke}{rgb}{0.000000,0.000000,0.000000}%
\pgfsetstrokecolor{currentstroke}%
\pgfsetdash{}{0pt}%
\pgfpathmoveto{\pgfqpoint{1.301350in}{0.804009in}}%
\pgfpathlineto{\pgfqpoint{1.314645in}{0.816193in}}%
\pgfpathlineto{\pgfqpoint{1.327952in}{0.828388in}}%
\pgfpathlineto{\pgfqpoint{1.325751in}{0.834820in}}%
\pgfpathlineto{\pgfqpoint{1.322534in}{0.841148in}}%
\pgfpathlineto{\pgfqpoint{1.309577in}{0.828285in}}%
\pgfpathlineto{\pgfqpoint{1.296626in}{0.815428in}}%
\pgfpathlineto{\pgfqpoint{1.299440in}{0.809764in}}%
\pgfpathlineto{\pgfqpoint{1.301350in}{0.804009in}}%
\pgfpathclose%
\pgfusepath{fill}%
\end{pgfscope}%
\begin{pgfscope}%
\pgfpathrectangle{\pgfqpoint{0.407592in}{0.222904in}}{\pgfqpoint{1.411163in}{1.411163in}}%
\pgfusepath{clip}%
\pgfsetbuttcap%
\pgfsetroundjoin%
\definecolor{currentfill}{rgb}{0.749020,0.772549,0.792157}%
\pgfsetfillcolor{currentfill}%
\pgfsetfillopacity{0.160000}%
\pgfsetlinewidth{0.000000pt}%
\definecolor{currentstroke}{rgb}{0.000000,0.000000,0.000000}%
\pgfsetstrokecolor{currentstroke}%
\pgfsetdash{}{0pt}%
\pgfpathmoveto{\pgfqpoint{1.103412in}{0.623428in}}%
\pgfpathlineto{\pgfqpoint{1.110044in}{0.629915in}}%
\pgfpathlineto{\pgfqpoint{1.116702in}{0.636428in}}%
\pgfpathlineto{\pgfqpoint{1.119675in}{0.637295in}}%
\pgfpathlineto{\pgfqpoint{1.122516in}{0.638244in}}%
\pgfpathlineto{\pgfqpoint{1.114399in}{0.631277in}}%
\pgfpathlineto{\pgfqpoint{1.106312in}{0.624336in}}%
\pgfpathlineto{\pgfqpoint{1.104895in}{0.623862in}}%
\pgfpathlineto{\pgfqpoint{1.103412in}{0.623428in}}%
\pgfpathclose%
\pgfusepath{fill}%
\end{pgfscope}%
\begin{pgfscope}%
\pgfpathrectangle{\pgfqpoint{0.407592in}{0.222904in}}{\pgfqpoint{1.411163in}{1.411163in}}%
\pgfusepath{clip}%
\pgfsetbuttcap%
\pgfsetroundjoin%
\definecolor{currentfill}{rgb}{0.749020,0.772549,0.792157}%
\pgfsetfillcolor{currentfill}%
\pgfsetfillopacity{0.160000}%
\pgfsetlinewidth{0.000000pt}%
\definecolor{currentstroke}{rgb}{0.000000,0.000000,0.000000}%
\pgfsetstrokecolor{currentstroke}%
\pgfsetdash{}{0pt}%
\pgfpathmoveto{\pgfqpoint{1.137766in}{0.646756in}}%
\pgfpathlineto{\pgfqpoint{1.149727in}{0.655870in}}%
\pgfpathlineto{\pgfqpoint{1.161718in}{0.665008in}}%
\pgfpathlineto{\pgfqpoint{1.163812in}{0.667129in}}%
\pgfpathlineto{\pgfqpoint{1.165580in}{0.669297in}}%
\pgfpathlineto{\pgfqpoint{1.152954in}{0.659450in}}%
\pgfpathlineto{\pgfqpoint{1.140354in}{0.649624in}}%
\pgfpathlineto{\pgfqpoint{1.139168in}{0.648174in}}%
\pgfpathlineto{\pgfqpoint{1.137766in}{0.646756in}}%
\pgfpathclose%
\pgfusepath{fill}%
\end{pgfscope}%
\begin{pgfscope}%
\pgfpathrectangle{\pgfqpoint{0.407592in}{0.222904in}}{\pgfqpoint{1.411163in}{1.411163in}}%
\pgfusepath{clip}%
\pgfsetbuttcap%
\pgfsetroundjoin%
\definecolor{currentfill}{rgb}{0.749020,0.772549,0.792157}%
\pgfsetfillcolor{currentfill}%
\pgfsetfillopacity{0.160000}%
\pgfsetlinewidth{0.000000pt}%
\definecolor{currentstroke}{rgb}{0.000000,0.000000,0.000000}%
\pgfsetstrokecolor{currentstroke}%
\pgfsetdash{}{0pt}%
\pgfpathmoveto{\pgfqpoint{1.661559in}{1.177707in}}%
\pgfpathlineto{\pgfqpoint{1.674681in}{1.190733in}}%
\pgfpathlineto{\pgfqpoint{1.687810in}{1.203766in}}%
\pgfpathlineto{\pgfqpoint{1.675266in}{1.217936in}}%
\pgfpathlineto{\pgfqpoint{1.660348in}{1.231668in}}%
\pgfpathlineto{\pgfqpoint{1.647938in}{1.218147in}}%
\pgfpathlineto{\pgfqpoint{1.635528in}{1.204627in}}%
\pgfpathlineto{\pgfqpoint{1.649683in}{1.191375in}}%
\pgfpathlineto{\pgfqpoint{1.661559in}{1.177707in}}%
\pgfpathclose%
\pgfusepath{fill}%
\end{pgfscope}%
\begin{pgfscope}%
\pgfpathrectangle{\pgfqpoint{0.407592in}{0.222904in}}{\pgfqpoint{1.411163in}{1.411163in}}%
\pgfusepath{clip}%
\pgfsetbuttcap%
\pgfsetroundjoin%
\definecolor{currentfill}{rgb}{0.749020,0.772549,0.792157}%
\pgfsetfillcolor{currentfill}%
\pgfsetfillopacity{0.160000}%
\pgfsetlinewidth{0.000000pt}%
\definecolor{currentstroke}{rgb}{0.000000,0.000000,0.000000}%
\pgfsetstrokecolor{currentstroke}%
\pgfsetdash{}{0pt}%
\pgfpathmoveto{\pgfqpoint{0.631067in}{1.121196in}}%
\pgfpathlineto{\pgfqpoint{0.618961in}{1.134661in}}%
\pgfpathlineto{\pgfqpoint{0.606853in}{1.148127in}}%
\pgfpathlineto{\pgfqpoint{0.591725in}{1.135789in}}%
\pgfpathlineto{\pgfqpoint{0.578613in}{1.123085in}}%
\pgfpathlineto{\pgfqpoint{0.591531in}{1.110143in}}%
\pgfpathlineto{\pgfqpoint{0.604443in}{1.097208in}}%
\pgfpathlineto{\pgfqpoint{0.616794in}{1.109375in}}%
\pgfpathlineto{\pgfqpoint{0.631067in}{1.121196in}}%
\pgfpathclose%
\pgfusepath{fill}%
\end{pgfscope}%
\begin{pgfscope}%
\pgfpathrectangle{\pgfqpoint{0.407592in}{0.222904in}}{\pgfqpoint{1.411163in}{1.411163in}}%
\pgfusepath{clip}%
\pgfsetbuttcap%
\pgfsetroundjoin%
\definecolor{currentfill}{rgb}{0.749020,0.772549,0.792157}%
\pgfsetfillcolor{currentfill}%
\pgfsetfillopacity{0.160000}%
\pgfsetlinewidth{0.000000pt}%
\definecolor{currentstroke}{rgb}{0.000000,0.000000,0.000000}%
\pgfsetstrokecolor{currentstroke}%
\pgfsetdash{}{0pt}%
\pgfpathmoveto{\pgfqpoint{1.038772in}{0.649331in}}%
\pgfpathlineto{\pgfqpoint{1.025837in}{0.659084in}}%
\pgfpathlineto{\pgfqpoint{1.012872in}{0.668859in}}%
\pgfpathlineto{\pgfqpoint{1.014260in}{0.666699in}}%
\pgfpathlineto{\pgfqpoint{1.015983in}{0.664589in}}%
\pgfpathlineto{\pgfqpoint{1.028438in}{0.655520in}}%
\pgfpathlineto{\pgfqpoint{1.040861in}{0.646476in}}%
\pgfpathlineto{\pgfqpoint{1.039706in}{0.647887in}}%
\pgfpathlineto{\pgfqpoint{1.038772in}{0.649331in}}%
\pgfpathclose%
\pgfusepath{fill}%
\end{pgfscope}%
\begin{pgfscope}%
\pgfpathrectangle{\pgfqpoint{0.407592in}{0.222904in}}{\pgfqpoint{1.411163in}{1.411163in}}%
\pgfusepath{clip}%
\pgfsetbuttcap%
\pgfsetroundjoin%
\definecolor{currentfill}{rgb}{0.749020,0.772549,0.792157}%
\pgfsetfillcolor{currentfill}%
\pgfsetfillopacity{0.160000}%
\pgfsetlinewidth{0.000000pt}%
\definecolor{currentstroke}{rgb}{0.000000,0.000000,0.000000}%
\pgfsetstrokecolor{currentstroke}%
\pgfsetdash{}{0pt}%
\pgfpathmoveto{\pgfqpoint{1.073835in}{0.623771in}}%
\pgfpathlineto{\pgfqpoint{1.065590in}{0.630430in}}%
\pgfpathlineto{\pgfqpoint{1.057311in}{0.637114in}}%
\pgfpathlineto{\pgfqpoint{1.060134in}{0.636264in}}%
\pgfpathlineto{\pgfqpoint{1.063090in}{0.635499in}}%
\pgfpathlineto{\pgfqpoint{1.069918in}{0.629218in}}%
\pgfpathlineto{\pgfqpoint{1.076716in}{0.622964in}}%
\pgfpathlineto{\pgfqpoint{1.075243in}{0.623346in}}%
\pgfpathlineto{\pgfqpoint{1.073835in}{0.623771in}}%
\pgfpathclose%
\pgfusepath{fill}%
\end{pgfscope}%
\begin{pgfscope}%
\pgfpathrectangle{\pgfqpoint{0.407592in}{0.222904in}}{\pgfqpoint{1.411163in}{1.411163in}}%
\pgfusepath{clip}%
\pgfsetbuttcap%
\pgfsetroundjoin%
\definecolor{currentfill}{rgb}{0.749020,0.772549,0.792157}%
\pgfsetfillcolor{currentfill}%
\pgfsetfillopacity{0.160000}%
\pgfsetlinewidth{0.000000pt}%
\definecolor{currentstroke}{rgb}{0.000000,0.000000,0.000000}%
\pgfsetstrokecolor{currentstroke}%
\pgfsetdash{}{0pt}%
\pgfpathmoveto{\pgfqpoint{1.222316in}{0.723701in}}%
\pgfpathlineto{\pgfqpoint{1.235621in}{0.735100in}}%
\pgfpathlineto{\pgfqpoint{1.248943in}{0.746514in}}%
\pgfpathlineto{\pgfqpoint{1.248962in}{0.750962in}}%
\pgfpathlineto{\pgfqpoint{1.248288in}{0.755381in}}%
\pgfpathlineto{\pgfqpoint{1.235052in}{0.743252in}}%
\pgfpathlineto{\pgfqpoint{1.221828in}{0.731133in}}%
\pgfpathlineto{\pgfqpoint{1.222361in}{0.727428in}}%
\pgfpathlineto{\pgfqpoint{1.222316in}{0.723701in}}%
\pgfpathclose%
\pgfusepath{fill}%
\end{pgfscope}%
\begin{pgfscope}%
\pgfpathrectangle{\pgfqpoint{0.407592in}{0.222904in}}{\pgfqpoint{1.411163in}{1.411163in}}%
\pgfusepath{clip}%
\pgfsetbuttcap%
\pgfsetroundjoin%
\definecolor{currentfill}{rgb}{0.749020,0.772549,0.792157}%
\pgfsetfillcolor{currentfill}%
\pgfsetfillopacity{0.160000}%
\pgfsetlinewidth{0.000000pt}%
\definecolor{currentstroke}{rgb}{0.000000,0.000000,0.000000}%
\pgfsetstrokecolor{currentstroke}%
\pgfsetdash{}{0pt}%
\pgfpathmoveto{\pgfqpoint{0.886834in}{0.814300in}}%
\pgfpathlineto{\pgfqpoint{0.874067in}{0.827091in}}%
\pgfpathlineto{\pgfqpoint{0.861293in}{0.839888in}}%
\pgfpathlineto{\pgfqpoint{0.856912in}{0.833537in}}%
\pgfpathlineto{\pgfqpoint{0.853529in}{0.827087in}}%
\pgfpathlineto{\pgfqpoint{0.866784in}{0.814961in}}%
\pgfpathlineto{\pgfqpoint{0.880026in}{0.802846in}}%
\pgfpathlineto{\pgfqpoint{0.882986in}{0.808616in}}%
\pgfpathlineto{\pgfqpoint{0.886834in}{0.814300in}}%
\pgfpathclose%
\pgfusepath{fill}%
\end{pgfscope}%
\begin{pgfscope}%
\pgfpathrectangle{\pgfqpoint{0.407592in}{0.222904in}}{\pgfqpoint{1.411163in}{1.411163in}}%
\pgfusepath{clip}%
\pgfsetbuttcap%
\pgfsetroundjoin%
\definecolor{currentfill}{rgb}{0.749020,0.772549,0.792157}%
\pgfsetfillcolor{currentfill}%
\pgfsetfillopacity{0.160000}%
\pgfsetlinewidth{0.000000pt}%
\definecolor{currentstroke}{rgb}{0.000000,0.000000,0.000000}%
\pgfsetstrokecolor{currentstroke}%
\pgfsetdash{}{0pt}%
\pgfpathmoveto{\pgfqpoint{1.687810in}{1.203766in}}%
\pgfpathlineto{\pgfqpoint{1.700944in}{1.216805in}}%
\pgfpathlineto{\pgfqpoint{1.714085in}{1.229851in}}%
\pgfpathlineto{\pgfqpoint{1.700862in}{1.244511in}}%
\pgfpathlineto{\pgfqpoint{1.685171in}{1.258712in}}%
\pgfpathlineto{\pgfqpoint{1.672759in}{1.245190in}}%
\pgfpathlineto{\pgfqpoint{1.660348in}{1.231668in}}%
\pgfpathlineto{\pgfqpoint{1.675266in}{1.217936in}}%
\pgfpathlineto{\pgfqpoint{1.687810in}{1.203766in}}%
\pgfpathclose%
\pgfusepath{fill}%
\end{pgfscope}%
\begin{pgfscope}%
\pgfpathrectangle{\pgfqpoint{0.407592in}{0.222904in}}{\pgfqpoint{1.411163in}{1.411163in}}%
\pgfusepath{clip}%
\pgfsetbuttcap%
\pgfsetroundjoin%
\definecolor{currentfill}{rgb}{0.749020,0.772549,0.792157}%
\pgfsetfillcolor{currentfill}%
\pgfsetfillopacity{0.160000}%
\pgfsetlinewidth{0.000000pt}%
\definecolor{currentstroke}{rgb}{0.000000,0.000000,0.000000}%
\pgfsetstrokecolor{currentstroke}%
\pgfsetdash{}{0pt}%
\pgfpathmoveto{\pgfqpoint{1.100280in}{0.622690in}}%
\pgfpathlineto{\pgfqpoint{1.105337in}{0.628808in}}%
\pgfpathlineto{\pgfqpoint{1.110416in}{0.634951in}}%
\pgfpathlineto{\pgfqpoint{1.113611in}{0.635645in}}%
\pgfpathlineto{\pgfqpoint{1.116702in}{0.636428in}}%
\pgfpathlineto{\pgfqpoint{1.110044in}{0.629915in}}%
\pgfpathlineto{\pgfqpoint{1.103412in}{0.623428in}}%
\pgfpathlineto{\pgfqpoint{1.101872in}{0.623037in}}%
\pgfpathlineto{\pgfqpoint{1.100280in}{0.622690in}}%
\pgfpathclose%
\pgfusepath{fill}%
\end{pgfscope}%
\begin{pgfscope}%
\pgfpathrectangle{\pgfqpoint{0.407592in}{0.222904in}}{\pgfqpoint{1.411163in}{1.411163in}}%
\pgfusepath{clip}%
\pgfsetbuttcap%
\pgfsetroundjoin%
\definecolor{currentfill}{rgb}{0.749020,0.772549,0.792157}%
\pgfsetfillcolor{currentfill}%
\pgfsetfillopacity{0.160000}%
\pgfsetlinewidth{0.000000pt}%
\definecolor{currentstroke}{rgb}{0.000000,0.000000,0.000000}%
\pgfsetstrokecolor{currentstroke}%
\pgfsetdash{}{0pt}%
\pgfpathmoveto{\pgfqpoint{0.606853in}{1.148127in}}%
\pgfpathlineto{\pgfqpoint{0.594744in}{1.161594in}}%
\pgfpathlineto{\pgfqpoint{0.582634in}{1.175063in}}%
\pgfpathlineto{\pgfqpoint{0.566641in}{1.162219in}}%
\pgfpathlineto{\pgfqpoint{0.552757in}{1.148988in}}%
\pgfpathlineto{\pgfqpoint{0.565688in}{1.136033in}}%
\pgfpathlineto{\pgfqpoint{0.578613in}{1.123085in}}%
\pgfpathlineto{\pgfqpoint{0.591725in}{1.135789in}}%
\pgfpathlineto{\pgfqpoint{0.606853in}{1.148127in}}%
\pgfpathclose%
\pgfusepath{fill}%
\end{pgfscope}%
\begin{pgfscope}%
\pgfpathrectangle{\pgfqpoint{0.407592in}{0.222904in}}{\pgfqpoint{1.411163in}{1.411163in}}%
\pgfusepath{clip}%
\pgfsetbuttcap%
\pgfsetroundjoin%
\definecolor{currentfill}{rgb}{0.749020,0.772549,0.792157}%
\pgfsetfillcolor{currentfill}%
\pgfsetfillopacity{0.160000}%
\pgfsetlinewidth{0.000000pt}%
\definecolor{currentstroke}{rgb}{0.000000,0.000000,0.000000}%
\pgfsetstrokecolor{currentstroke}%
\pgfsetdash{}{0pt}%
\pgfpathmoveto{\pgfqpoint{0.959220in}{0.730392in}}%
\pgfpathlineto{\pgfqpoint{0.946052in}{0.742439in}}%
\pgfpathlineto{\pgfqpoint{0.932871in}{0.754498in}}%
\pgfpathlineto{\pgfqpoint{0.931403in}{0.750071in}}%
\pgfpathlineto{\pgfqpoint{0.930624in}{0.745620in}}%
\pgfpathlineto{\pgfqpoint{0.944024in}{0.734278in}}%
\pgfpathlineto{\pgfqpoint{0.957406in}{0.722952in}}%
\pgfpathlineto{\pgfqpoint{0.958026in}{0.726681in}}%
\pgfpathlineto{\pgfqpoint{0.959220in}{0.730392in}}%
\pgfpathclose%
\pgfusepath{fill}%
\end{pgfscope}%
\begin{pgfscope}%
\pgfpathrectangle{\pgfqpoint{0.407592in}{0.222904in}}{\pgfqpoint{1.411163in}{1.411163in}}%
\pgfusepath{clip}%
\pgfsetbuttcap%
\pgfsetroundjoin%
\definecolor{currentfill}{rgb}{0.749020,0.772549,0.792157}%
\pgfsetfillcolor{currentfill}%
\pgfsetfillopacity{0.160000}%
\pgfsetlinewidth{0.000000pt}%
\definecolor{currentstroke}{rgb}{0.000000,0.000000,0.000000}%
\pgfsetstrokecolor{currentstroke}%
\pgfsetdash{}{0pt}%
\pgfpathmoveto{\pgfqpoint{1.076716in}{0.622964in}}%
\pgfpathlineto{\pgfqpoint{1.069918in}{0.629218in}}%
\pgfpathlineto{\pgfqpoint{1.063090in}{0.635499in}}%
\pgfpathlineto{\pgfqpoint{1.066167in}{0.634823in}}%
\pgfpathlineto{\pgfqpoint{1.069349in}{0.634239in}}%
\pgfpathlineto{\pgfqpoint{1.074603in}{0.628274in}}%
\pgfpathlineto{\pgfqpoint{1.079835in}{0.622335in}}%
\pgfpathlineto{\pgfqpoint{1.078249in}{0.622626in}}%
\pgfpathlineto{\pgfqpoint{1.076716in}{0.622964in}}%
\pgfpathclose%
\pgfusepath{fill}%
\end{pgfscope}%
\begin{pgfscope}%
\pgfpathrectangle{\pgfqpoint{0.407592in}{0.222904in}}{\pgfqpoint{1.411163in}{1.411163in}}%
\pgfusepath{clip}%
\pgfsetbuttcap%
\pgfsetroundjoin%
\definecolor{currentfill}{rgb}{0.749020,0.772549,0.792157}%
\pgfsetfillcolor{currentfill}%
\pgfsetfillopacity{0.160000}%
\pgfsetlinewidth{0.000000pt}%
\definecolor{currentstroke}{rgb}{0.000000,0.000000,0.000000}%
\pgfsetstrokecolor{currentstroke}%
\pgfsetdash{}{0pt}%
\pgfpathmoveto{\pgfqpoint{1.327952in}{0.828388in}}%
\pgfpathlineto{\pgfqpoint{1.341271in}{0.840594in}}%
\pgfpathlineto{\pgfqpoint{1.354602in}{0.852811in}}%
\pgfpathlineto{\pgfqpoint{1.352098in}{0.859910in}}%
\pgfpathlineto{\pgfqpoint{1.348467in}{0.866892in}}%
\pgfpathlineto{\pgfqpoint{1.335497in}{0.854017in}}%
\pgfpathlineto{\pgfqpoint{1.322534in}{0.841148in}}%
\pgfpathlineto{\pgfqpoint{1.325751in}{0.834820in}}%
\pgfpathlineto{\pgfqpoint{1.327952in}{0.828388in}}%
\pgfpathclose%
\pgfusepath{fill}%
\end{pgfscope}%
\begin{pgfscope}%
\pgfpathrectangle{\pgfqpoint{0.407592in}{0.222904in}}{\pgfqpoint{1.411163in}{1.411163in}}%
\pgfusepath{clip}%
\pgfsetbuttcap%
\pgfsetroundjoin%
\definecolor{currentfill}{rgb}{0.749020,0.772549,0.792157}%
\pgfsetfillcolor{currentfill}%
\pgfsetfillopacity{0.160000}%
\pgfsetlinewidth{0.000000pt}%
\definecolor{currentstroke}{rgb}{0.000000,0.000000,0.000000}%
\pgfsetstrokecolor{currentstroke}%
\pgfsetdash{}{0pt}%
\pgfpathmoveto{\pgfqpoint{1.714085in}{1.229851in}}%
\pgfpathlineto{\pgfqpoint{1.727232in}{1.242902in}}%
\pgfpathlineto{\pgfqpoint{1.740385in}{1.255959in}}%
\pgfpathlineto{\pgfqpoint{1.726472in}{1.271100in}}%
\pgfpathlineto{\pgfqpoint{1.709998in}{1.285760in}}%
\pgfpathlineto{\pgfqpoint{1.697584in}{1.272236in}}%
\pgfpathlineto{\pgfqpoint{1.685171in}{1.258712in}}%
\pgfpathlineto{\pgfqpoint{1.700862in}{1.244511in}}%
\pgfpathlineto{\pgfqpoint{1.714085in}{1.229851in}}%
\pgfpathclose%
\pgfusepath{fill}%
\end{pgfscope}%
\begin{pgfscope}%
\pgfpathrectangle{\pgfqpoint{0.407592in}{0.222904in}}{\pgfqpoint{1.411163in}{1.411163in}}%
\pgfusepath{clip}%
\pgfsetbuttcap%
\pgfsetroundjoin%
\definecolor{currentfill}{rgb}{0.749020,0.772549,0.792157}%
\pgfsetfillcolor{currentfill}%
\pgfsetfillopacity{0.160000}%
\pgfsetlinewidth{0.000000pt}%
\definecolor{currentstroke}{rgb}{0.000000,0.000000,0.000000}%
\pgfsetstrokecolor{currentstroke}%
\pgfsetdash{}{0pt}%
\pgfpathmoveto{\pgfqpoint{1.096969in}{0.622136in}}%
\pgfpathlineto{\pgfqpoint{1.100361in}{0.627975in}}%
\pgfpathlineto{\pgfqpoint{1.103768in}{0.633841in}}%
\pgfpathlineto{\pgfqpoint{1.107130in}{0.634349in}}%
\pgfpathlineto{\pgfqpoint{1.110416in}{0.634951in}}%
\pgfpathlineto{\pgfqpoint{1.105337in}{0.628808in}}%
\pgfpathlineto{\pgfqpoint{1.100280in}{0.622690in}}%
\pgfpathlineto{\pgfqpoint{1.098643in}{0.622389in}}%
\pgfpathlineto{\pgfqpoint{1.096969in}{0.622136in}}%
\pgfpathclose%
\pgfusepath{fill}%
\end{pgfscope}%
\begin{pgfscope}%
\pgfpathrectangle{\pgfqpoint{0.407592in}{0.222904in}}{\pgfqpoint{1.411163in}{1.411163in}}%
\pgfusepath{clip}%
\pgfsetbuttcap%
\pgfsetroundjoin%
\definecolor{currentfill}{rgb}{0.749020,0.772549,0.792157}%
\pgfsetfillcolor{currentfill}%
\pgfsetfillopacity{0.160000}%
\pgfsetlinewidth{0.000000pt}%
\definecolor{currentstroke}{rgb}{0.000000,0.000000,0.000000}%
\pgfsetstrokecolor{currentstroke}%
\pgfsetdash{}{0pt}%
\pgfpathmoveto{\pgfqpoint{0.582634in}{1.175063in}}%
\pgfpathlineto{\pgfqpoint{0.570522in}{1.188533in}}%
\pgfpathlineto{\pgfqpoint{0.558409in}{1.202004in}}%
\pgfpathlineto{\pgfqpoint{0.541542in}{1.188665in}}%
\pgfpathlineto{\pgfqpoint{0.526874in}{1.174919in}}%
\pgfpathlineto{\pgfqpoint{0.539819in}{1.161950in}}%
\pgfpathlineto{\pgfqpoint{0.552757in}{1.148988in}}%
\pgfpathlineto{\pgfqpoint{0.566641in}{1.162219in}}%
\pgfpathlineto{\pgfqpoint{0.582634in}{1.175063in}}%
\pgfpathclose%
\pgfusepath{fill}%
\end{pgfscope}%
\begin{pgfscope}%
\pgfpathrectangle{\pgfqpoint{0.407592in}{0.222904in}}{\pgfqpoint{1.411163in}{1.411163in}}%
\pgfusepath{clip}%
\pgfsetbuttcap%
\pgfsetroundjoin%
\definecolor{currentfill}{rgb}{0.749020,0.772549,0.792157}%
\pgfsetfillcolor{currentfill}%
\pgfsetfillopacity{0.160000}%
\pgfsetlinewidth{0.000000pt}%
\definecolor{currentstroke}{rgb}{0.000000,0.000000,0.000000}%
\pgfsetstrokecolor{currentstroke}%
\pgfsetdash{}{0pt}%
\pgfpathmoveto{\pgfqpoint{1.079835in}{0.622335in}}%
\pgfpathlineto{\pgfqpoint{1.074603in}{0.628274in}}%
\pgfpathlineto{\pgfqpoint{1.069349in}{0.634239in}}%
\pgfpathlineto{\pgfqpoint{1.072624in}{0.633751in}}%
\pgfpathlineto{\pgfqpoint{1.075977in}{0.633359in}}%
\pgfpathlineto{\pgfqpoint{1.079564in}{0.627614in}}%
\pgfpathlineto{\pgfqpoint{1.083135in}{0.621895in}}%
\pgfpathlineto{\pgfqpoint{1.081466in}{0.622091in}}%
\pgfpathlineto{\pgfqpoint{1.079835in}{0.622335in}}%
\pgfpathclose%
\pgfusepath{fill}%
\end{pgfscope}%
\begin{pgfscope}%
\pgfpathrectangle{\pgfqpoint{0.407592in}{0.222904in}}{\pgfqpoint{1.411163in}{1.411163in}}%
\pgfusepath{clip}%
\pgfsetbuttcap%
\pgfsetroundjoin%
\definecolor{currentfill}{rgb}{0.749020,0.772549,0.792157}%
\pgfsetfillcolor{currentfill}%
\pgfsetfillopacity{0.160000}%
\pgfsetlinewidth{0.000000pt}%
\definecolor{currentstroke}{rgb}{0.000000,0.000000,0.000000}%
\pgfsetstrokecolor{currentstroke}%
\pgfsetdash{}{0pt}%
\pgfpathmoveto{\pgfqpoint{1.093538in}{0.621774in}}%
\pgfpathlineto{\pgfqpoint{1.095203in}{0.627433in}}%
\pgfpathlineto{\pgfqpoint{1.096876in}{0.633117in}}%
\pgfpathlineto{\pgfqpoint{1.100345in}{0.633430in}}%
\pgfpathlineto{\pgfqpoint{1.103768in}{0.633841in}}%
\pgfpathlineto{\pgfqpoint{1.100361in}{0.627975in}}%
\pgfpathlineto{\pgfqpoint{1.096969in}{0.622136in}}%
\pgfpathlineto{\pgfqpoint{1.095265in}{0.621930in}}%
\pgfpathlineto{\pgfqpoint{1.093538in}{0.621774in}}%
\pgfpathclose%
\pgfusepath{fill}%
\end{pgfscope}%
\begin{pgfscope}%
\pgfpathrectangle{\pgfqpoint{0.407592in}{0.222904in}}{\pgfqpoint{1.411163in}{1.411163in}}%
\pgfusepath{clip}%
\pgfsetbuttcap%
\pgfsetroundjoin%
\definecolor{currentfill}{rgb}{0.749020,0.772549,0.792157}%
\pgfsetfillcolor{currentfill}%
\pgfsetfillopacity{0.160000}%
\pgfsetlinewidth{0.000000pt}%
\definecolor{currentstroke}{rgb}{0.000000,0.000000,0.000000}%
\pgfsetstrokecolor{currentstroke}%
\pgfsetdash{}{0pt}%
\pgfpathmoveto{\pgfqpoint{1.740385in}{1.255959in}}%
\pgfpathlineto{\pgfqpoint{1.753545in}{1.269023in}}%
\pgfpathlineto{\pgfqpoint{1.766710in}{1.282093in}}%
\pgfpathlineto{\pgfqpoint{1.752095in}{1.297702in}}%
\pgfpathlineto{\pgfqpoint{1.734827in}{1.312811in}}%
\pgfpathlineto{\pgfqpoint{1.722412in}{1.299285in}}%
\pgfpathlineto{\pgfqpoint{1.709998in}{1.285760in}}%
\pgfpathlineto{\pgfqpoint{1.726472in}{1.271100in}}%
\pgfpathlineto{\pgfqpoint{1.740385in}{1.255959in}}%
\pgfpathclose%
\pgfusepath{fill}%
\end{pgfscope}%
\begin{pgfscope}%
\pgfpathrectangle{\pgfqpoint{0.407592in}{0.222904in}}{\pgfqpoint{1.411163in}{1.411163in}}%
\pgfusepath{clip}%
\pgfsetbuttcap%
\pgfsetroundjoin%
\definecolor{currentfill}{rgb}{0.749020,0.772549,0.792157}%
\pgfsetfillcolor{currentfill}%
\pgfsetfillopacity{0.160000}%
\pgfsetlinewidth{0.000000pt}%
\definecolor{currentstroke}{rgb}{0.000000,0.000000,0.000000}%
\pgfsetstrokecolor{currentstroke}%
\pgfsetdash{}{0pt}%
\pgfpathmoveto{\pgfqpoint{1.083135in}{0.621895in}}%
\pgfpathlineto{\pgfqpoint{1.079564in}{0.627614in}}%
\pgfpathlineto{\pgfqpoint{1.075977in}{0.633359in}}%
\pgfpathlineto{\pgfqpoint{1.079393in}{0.633066in}}%
\pgfpathlineto{\pgfqpoint{1.082856in}{0.632874in}}%
\pgfpathlineto{\pgfqpoint{1.084712in}{0.627251in}}%
\pgfpathlineto{\pgfqpoint{1.086560in}{0.621653in}}%
\pgfpathlineto{\pgfqpoint{1.084836in}{0.621749in}}%
\pgfpathlineto{\pgfqpoint{1.083135in}{0.621895in}}%
\pgfpathclose%
\pgfusepath{fill}%
\end{pgfscope}%
\begin{pgfscope}%
\pgfpathrectangle{\pgfqpoint{0.407592in}{0.222904in}}{\pgfqpoint{1.411163in}{1.411163in}}%
\pgfusepath{clip}%
\pgfsetbuttcap%
\pgfsetroundjoin%
\definecolor{currentfill}{rgb}{0.749020,0.772549,0.792157}%
\pgfsetfillcolor{currentfill}%
\pgfsetfillopacity{0.160000}%
\pgfsetlinewidth{0.000000pt}%
\definecolor{currentstroke}{rgb}{0.000000,0.000000,0.000000}%
\pgfsetstrokecolor{currentstroke}%
\pgfsetdash{}{0pt}%
\pgfpathmoveto{\pgfqpoint{1.090048in}{0.621612in}}%
\pgfpathlineto{\pgfqpoint{1.089956in}{0.627190in}}%
\pgfpathlineto{\pgfqpoint{1.089864in}{0.632793in}}%
\pgfpathlineto{\pgfqpoint{1.093378in}{0.632904in}}%
\pgfpathlineto{\pgfqpoint{1.096876in}{0.633117in}}%
\pgfpathlineto{\pgfqpoint{1.095203in}{0.627433in}}%
\pgfpathlineto{\pgfqpoint{1.093538in}{0.621774in}}%
\pgfpathlineto{\pgfqpoint{1.091796in}{0.621668in}}%
\pgfpathlineto{\pgfqpoint{1.090048in}{0.621612in}}%
\pgfpathclose%
\pgfusepath{fill}%
\end{pgfscope}%
\begin{pgfscope}%
\pgfpathrectangle{\pgfqpoint{0.407592in}{0.222904in}}{\pgfqpoint{1.411163in}{1.411163in}}%
\pgfusepath{clip}%
\pgfsetbuttcap%
\pgfsetroundjoin%
\definecolor{currentfill}{rgb}{0.749020,0.772549,0.792157}%
\pgfsetfillcolor{currentfill}%
\pgfsetfillopacity{0.160000}%
\pgfsetlinewidth{0.000000pt}%
\definecolor{currentstroke}{rgb}{0.000000,0.000000,0.000000}%
\pgfsetstrokecolor{currentstroke}%
\pgfsetdash{}{0pt}%
\pgfpathmoveto{\pgfqpoint{1.086560in}{0.621653in}}%
\pgfpathlineto{\pgfqpoint{1.084712in}{0.627251in}}%
\pgfpathlineto{\pgfqpoint{1.082856in}{0.632874in}}%
\pgfpathlineto{\pgfqpoint{1.086352in}{0.632783in}}%
\pgfpathlineto{\pgfqpoint{1.089864in}{0.632793in}}%
\pgfpathlineto{\pgfqpoint{1.089956in}{0.627190in}}%
\pgfpathlineto{\pgfqpoint{1.090048in}{0.621612in}}%
\pgfpathlineto{\pgfqpoint{1.088300in}{0.621607in}}%
\pgfpathlineto{\pgfqpoint{1.086560in}{0.621653in}}%
\pgfpathclose%
\pgfusepath{fill}%
\end{pgfscope}%
\begin{pgfscope}%
\pgfpathrectangle{\pgfqpoint{0.407592in}{0.222904in}}{\pgfqpoint{1.411163in}{1.411163in}}%
\pgfusepath{clip}%
\pgfsetbuttcap%
\pgfsetroundjoin%
\definecolor{currentfill}{rgb}{0.749020,0.772549,0.792157}%
\pgfsetfillcolor{currentfill}%
\pgfsetfillopacity{0.160000}%
\pgfsetlinewidth{0.000000pt}%
\definecolor{currentstroke}{rgb}{0.000000,0.000000,0.000000}%
\pgfsetstrokecolor{currentstroke}%
\pgfsetdash{}{0pt}%
\pgfpathmoveto{\pgfqpoint{0.861293in}{0.839888in}}%
\pgfpathlineto{\pgfqpoint{0.848513in}{0.852692in}}%
\pgfpathlineto{\pgfqpoint{0.835726in}{0.865502in}}%
\pgfpathlineto{\pgfqpoint{0.830800in}{0.858495in}}%
\pgfpathlineto{\pgfqpoint{0.826982in}{0.851374in}}%
\pgfpathlineto{\pgfqpoint{0.840262in}{0.839225in}}%
\pgfpathlineto{\pgfqpoint{0.853529in}{0.827087in}}%
\pgfpathlineto{\pgfqpoint{0.856912in}{0.833537in}}%
\pgfpathlineto{\pgfqpoint{0.861293in}{0.839888in}}%
\pgfpathclose%
\pgfusepath{fill}%
\end{pgfscope}%
\begin{pgfscope}%
\pgfpathrectangle{\pgfqpoint{0.407592in}{0.222904in}}{\pgfqpoint{1.411163in}{1.411163in}}%
\pgfusepath{clip}%
\pgfsetbuttcap%
\pgfsetroundjoin%
\definecolor{currentfill}{rgb}{0.749020,0.772549,0.792157}%
\pgfsetfillcolor{currentfill}%
\pgfsetfillopacity{0.160000}%
\pgfsetlinewidth{0.000000pt}%
\definecolor{currentstroke}{rgb}{0.000000,0.000000,0.000000}%
\pgfsetstrokecolor{currentstroke}%
\pgfsetdash{}{0pt}%
\pgfpathmoveto{\pgfqpoint{1.165580in}{0.669297in}}%
\pgfpathlineto{\pgfqpoint{1.178234in}{0.679166in}}%
\pgfpathlineto{\pgfqpoint{1.190916in}{0.689056in}}%
\pgfpathlineto{\pgfqpoint{1.192812in}{0.691988in}}%
\pgfpathlineto{\pgfqpoint{1.194254in}{0.694956in}}%
\pgfpathlineto{\pgfqpoint{1.181171in}{0.684338in}}%
\pgfpathlineto{\pgfqpoint{1.168111in}{0.673739in}}%
\pgfpathlineto{\pgfqpoint{1.167015in}{0.671504in}}%
\pgfpathlineto{\pgfqpoint{1.165580in}{0.669297in}}%
\pgfpathclose%
\pgfusepath{fill}%
\end{pgfscope}%
\begin{pgfscope}%
\pgfpathrectangle{\pgfqpoint{0.407592in}{0.222904in}}{\pgfqpoint{1.411163in}{1.411163in}}%
\pgfusepath{clip}%
\pgfsetbuttcap%
\pgfsetroundjoin%
\definecolor{currentfill}{rgb}{0.749020,0.772549,0.792157}%
\pgfsetfillcolor{currentfill}%
\pgfsetfillopacity{0.160000}%
\pgfsetlinewidth{0.000000pt}%
\definecolor{currentstroke}{rgb}{0.000000,0.000000,0.000000}%
\pgfsetstrokecolor{currentstroke}%
\pgfsetdash{}{0pt}%
\pgfpathmoveto{\pgfqpoint{0.558409in}{1.202004in}}%
\pgfpathlineto{\pgfqpoint{0.546296in}{1.215477in}}%
\pgfpathlineto{\pgfqpoint{0.534180in}{1.228950in}}%
\pgfpathlineto{\pgfqpoint{0.516428in}{1.215127in}}%
\pgfpathlineto{\pgfqpoint{0.500965in}{1.200875in}}%
\pgfpathlineto{\pgfqpoint{0.513923in}{1.187894in}}%
\pgfpathlineto{\pgfqpoint{0.526874in}{1.174919in}}%
\pgfpathlineto{\pgfqpoint{0.541542in}{1.188665in}}%
\pgfpathlineto{\pgfqpoint{0.558409in}{1.202004in}}%
\pgfpathclose%
\pgfusepath{fill}%
\end{pgfscope}%
\begin{pgfscope}%
\pgfpathrectangle{\pgfqpoint{0.407592in}{0.222904in}}{\pgfqpoint{1.411163in}{1.411163in}}%
\pgfusepath{clip}%
\pgfsetbuttcap%
\pgfsetroundjoin%
\definecolor{currentfill}{rgb}{0.749020,0.772549,0.792157}%
\pgfsetfillcolor{currentfill}%
\pgfsetfillopacity{0.160000}%
\pgfsetlinewidth{0.000000pt}%
\definecolor{currentstroke}{rgb}{0.000000,0.000000,0.000000}%
\pgfsetstrokecolor{currentstroke}%
\pgfsetdash{}{0pt}%
\pgfpathmoveto{\pgfqpoint{1.194254in}{0.694956in}}%
\pgfpathlineto{\pgfqpoint{1.207360in}{0.705592in}}%
\pgfpathlineto{\pgfqpoint{1.220488in}{0.716247in}}%
\pgfpathlineto{\pgfqpoint{1.221691in}{0.719968in}}%
\pgfpathlineto{\pgfqpoint{1.222316in}{0.723701in}}%
\pgfpathlineto{\pgfqpoint{1.209029in}{0.712317in}}%
\pgfpathlineto{\pgfqpoint{1.195759in}{0.700948in}}%
\pgfpathlineto{\pgfqpoint{1.195237in}{0.697947in}}%
\pgfpathlineto{\pgfqpoint{1.194254in}{0.694956in}}%
\pgfpathclose%
\pgfusepath{fill}%
\end{pgfscope}%
\begin{pgfscope}%
\pgfpathrectangle{\pgfqpoint{0.407592in}{0.222904in}}{\pgfqpoint{1.411163in}{1.411163in}}%
\pgfusepath{clip}%
\pgfsetbuttcap%
\pgfsetroundjoin%
\definecolor{currentfill}{rgb}{0.749020,0.772549,0.792157}%
\pgfsetfillcolor{currentfill}%
\pgfsetfillopacity{0.160000}%
\pgfsetlinewidth{0.000000pt}%
\definecolor{currentstroke}{rgb}{0.000000,0.000000,0.000000}%
\pgfsetstrokecolor{currentstroke}%
\pgfsetdash{}{0pt}%
\pgfpathmoveto{\pgfqpoint{1.134338in}{0.644042in}}%
\pgfpathlineto{\pgfqpoint{1.145447in}{0.652482in}}%
\pgfpathlineto{\pgfqpoint{1.156589in}{0.660946in}}%
\pgfpathlineto{\pgfqpoint{1.159307in}{0.662944in}}%
\pgfpathlineto{\pgfqpoint{1.161718in}{0.665008in}}%
\pgfpathlineto{\pgfqpoint{1.149727in}{0.655870in}}%
\pgfpathlineto{\pgfqpoint{1.137766in}{0.646756in}}%
\pgfpathlineto{\pgfqpoint{1.136154in}{0.645377in}}%
\pgfpathlineto{\pgfqpoint{1.134338in}{0.644042in}}%
\pgfpathclose%
\pgfusepath{fill}%
\end{pgfscope}%
\begin{pgfscope}%
\pgfpathrectangle{\pgfqpoint{0.407592in}{0.222904in}}{\pgfqpoint{1.411163in}{1.411163in}}%
\pgfusepath{clip}%
\pgfsetbuttcap%
\pgfsetroundjoin%
\definecolor{currentfill}{rgb}{0.749020,0.772549,0.792157}%
\pgfsetfillcolor{currentfill}%
\pgfsetfillopacity{0.160000}%
\pgfsetlinewidth{0.000000pt}%
\definecolor{currentstroke}{rgb}{0.000000,0.000000,0.000000}%
\pgfsetstrokecolor{currentstroke}%
\pgfsetdash{}{0pt}%
\pgfpathmoveto{\pgfqpoint{1.766710in}{1.282093in}}%
\pgfpathlineto{\pgfqpoint{1.779882in}{1.295169in}}%
\pgfpathlineto{\pgfqpoint{1.793060in}{1.308251in}}%
\pgfpathlineto{\pgfqpoint{1.777732in}{1.324319in}}%
\pgfpathlineto{\pgfqpoint{1.759659in}{1.339865in}}%
\pgfpathlineto{\pgfqpoint{1.747242in}{1.326337in}}%
\pgfpathlineto{\pgfqpoint{1.734827in}{1.312811in}}%
\pgfpathlineto{\pgfqpoint{1.752095in}{1.297702in}}%
\pgfpathlineto{\pgfqpoint{1.766710in}{1.282093in}}%
\pgfpathclose%
\pgfusepath{fill}%
\end{pgfscope}%
\begin{pgfscope}%
\pgfpathrectangle{\pgfqpoint{0.407592in}{0.222904in}}{\pgfqpoint{1.411163in}{1.411163in}}%
\pgfusepath{clip}%
\pgfsetbuttcap%
\pgfsetroundjoin%
\definecolor{currentfill}{rgb}{0.749020,0.772549,0.792157}%
\pgfsetfillcolor{currentfill}%
\pgfsetfillopacity{0.160000}%
\pgfsetlinewidth{0.000000pt}%
\definecolor{currentstroke}{rgb}{0.000000,0.000000,0.000000}%
\pgfsetstrokecolor{currentstroke}%
\pgfsetdash{}{0pt}%
\pgfpathmoveto{\pgfqpoint{1.011121in}{0.673289in}}%
\pgfpathlineto{\pgfqpoint{0.997853in}{0.683814in}}%
\pgfpathlineto{\pgfqpoint{0.984562in}{0.694358in}}%
\pgfpathlineto{\pgfqpoint{0.985479in}{0.691396in}}%
\pgfpathlineto{\pgfqpoint{0.986858in}{0.688474in}}%
\pgfpathlineto{\pgfqpoint{0.999880in}{0.678655in}}%
\pgfpathlineto{\pgfqpoint{1.012872in}{0.668859in}}%
\pgfpathlineto{\pgfqpoint{1.011825in}{0.671059in}}%
\pgfpathlineto{\pgfqpoint{1.011121in}{0.673289in}}%
\pgfpathclose%
\pgfusepath{fill}%
\end{pgfscope}%
\begin{pgfscope}%
\pgfpathrectangle{\pgfqpoint{0.407592in}{0.222904in}}{\pgfqpoint{1.411163in}{1.411163in}}%
\pgfusepath{clip}%
\pgfsetbuttcap%
\pgfsetroundjoin%
\definecolor{currentfill}{rgb}{0.749020,0.772549,0.792157}%
\pgfsetfillcolor{currentfill}%
\pgfsetfillopacity{0.160000}%
\pgfsetlinewidth{0.000000pt}%
\definecolor{currentstroke}{rgb}{0.000000,0.000000,0.000000}%
\pgfsetstrokecolor{currentstroke}%
\pgfsetdash{}{0pt}%
\pgfpathmoveto{\pgfqpoint{1.040861in}{0.646476in}}%
\pgfpathlineto{\pgfqpoint{1.028438in}{0.655520in}}%
\pgfpathlineto{\pgfqpoint{1.015983in}{0.664589in}}%
\pgfpathlineto{\pgfqpoint{1.018033in}{0.662537in}}%
\pgfpathlineto{\pgfqpoint{1.020402in}{0.660554in}}%
\pgfpathlineto{\pgfqpoint{1.032127in}{0.652155in}}%
\pgfpathlineto{\pgfqpoint{1.043817in}{0.643781in}}%
\pgfpathlineto{\pgfqpoint{1.042234in}{0.645105in}}%
\pgfpathlineto{\pgfqpoint{1.040861in}{0.646476in}}%
\pgfpathclose%
\pgfusepath{fill}%
\end{pgfscope}%
\begin{pgfscope}%
\pgfpathrectangle{\pgfqpoint{0.407592in}{0.222904in}}{\pgfqpoint{1.411163in}{1.411163in}}%
\pgfusepath{clip}%
\pgfsetbuttcap%
\pgfsetroundjoin%
\definecolor{currentfill}{rgb}{0.749020,0.772549,0.792157}%
\pgfsetfillcolor{currentfill}%
\pgfsetfillopacity{0.160000}%
\pgfsetlinewidth{0.000000pt}%
\definecolor{currentstroke}{rgb}{0.000000,0.000000,0.000000}%
\pgfsetstrokecolor{currentstroke}%
\pgfsetdash{}{0pt}%
\pgfpathmoveto{\pgfqpoint{0.984116in}{0.700345in}}%
\pgfpathlineto{\pgfqpoint{0.970770in}{0.711641in}}%
\pgfpathlineto{\pgfqpoint{0.957406in}{0.722952in}}%
\pgfpathlineto{\pgfqpoint{0.957366in}{0.719220in}}%
\pgfpathlineto{\pgfqpoint{0.957907in}{0.715503in}}%
\pgfpathlineto{\pgfqpoint{0.971246in}{0.704921in}}%
\pgfpathlineto{\pgfqpoint{0.984562in}{0.694358in}}%
\pgfpathlineto{\pgfqpoint{0.984107in}{0.697345in}}%
\pgfpathlineto{\pgfqpoint{0.984116in}{0.700345in}}%
\pgfpathclose%
\pgfusepath{fill}%
\end{pgfscope}%
\begin{pgfscope}%
\pgfpathrectangle{\pgfqpoint{0.407592in}{0.222904in}}{\pgfqpoint{1.411163in}{1.411163in}}%
\pgfusepath{clip}%
\pgfsetbuttcap%
\pgfsetroundjoin%
\definecolor{currentfill}{rgb}{0.749020,0.772549,0.792157}%
\pgfsetfillcolor{currentfill}%
\pgfsetfillopacity{0.160000}%
\pgfsetlinewidth{0.000000pt}%
\definecolor{currentstroke}{rgb}{0.000000,0.000000,0.000000}%
\pgfsetstrokecolor{currentstroke}%
\pgfsetdash{}{0pt}%
\pgfpathmoveto{\pgfqpoint{1.354602in}{0.852811in}}%
\pgfpathlineto{\pgfqpoint{1.367945in}{0.865038in}}%
\pgfpathlineto{\pgfqpoint{1.381300in}{0.877277in}}%
\pgfpathlineto{\pgfqpoint{1.378481in}{0.885035in}}%
\pgfpathlineto{\pgfqpoint{1.374423in}{0.892660in}}%
\pgfpathlineto{\pgfqpoint{1.361442in}{0.879773in}}%
\pgfpathlineto{\pgfqpoint{1.348467in}{0.866892in}}%
\pgfpathlineto{\pgfqpoint{1.352098in}{0.859910in}}%
\pgfpathlineto{\pgfqpoint{1.354602in}{0.852811in}}%
\pgfpathclose%
\pgfusepath{fill}%
\end{pgfscope}%
\begin{pgfscope}%
\pgfpathrectangle{\pgfqpoint{0.407592in}{0.222904in}}{\pgfqpoint{1.411163in}{1.411163in}}%
\pgfusepath{clip}%
\pgfsetbuttcap%
\pgfsetroundjoin%
\definecolor{currentfill}{rgb}{0.749020,0.772549,0.792157}%
\pgfsetfillcolor{currentfill}%
\pgfsetfillopacity{0.160000}%
\pgfsetlinewidth{0.000000pt}%
\definecolor{currentstroke}{rgb}{0.000000,0.000000,0.000000}%
\pgfsetstrokecolor{currentstroke}%
\pgfsetdash{}{0pt}%
\pgfpathmoveto{\pgfqpoint{0.534180in}{1.228950in}}%
\pgfpathlineto{\pgfqpoint{0.522064in}{1.242426in}}%
\pgfpathlineto{\pgfqpoint{0.509947in}{1.255902in}}%
\pgfpathlineto{\pgfqpoint{0.491298in}{1.241605in}}%
\pgfpathlineto{\pgfqpoint{0.475029in}{1.226859in}}%
\pgfpathlineto{\pgfqpoint{0.488000in}{1.213864in}}%
\pgfpathlineto{\pgfqpoint{0.500965in}{1.200875in}}%
\pgfpathlineto{\pgfqpoint{0.516428in}{1.215127in}}%
\pgfpathlineto{\pgfqpoint{0.534180in}{1.228950in}}%
\pgfpathclose%
\pgfusepath{fill}%
\end{pgfscope}%
\begin{pgfscope}%
\pgfpathrectangle{\pgfqpoint{0.407592in}{0.222904in}}{\pgfqpoint{1.411163in}{1.411163in}}%
\pgfusepath{clip}%
\pgfsetbuttcap%
\pgfsetroundjoin%
\definecolor{currentfill}{rgb}{0.749020,0.772549,0.792157}%
\pgfsetfillcolor{currentfill}%
\pgfsetfillopacity{0.160000}%
\pgfsetlinewidth{0.000000pt}%
\definecolor{currentstroke}{rgb}{0.000000,0.000000,0.000000}%
\pgfsetstrokecolor{currentstroke}%
\pgfsetdash{}{0pt}%
\pgfpathmoveto{\pgfqpoint{1.248943in}{0.746514in}}%
\pgfpathlineto{\pgfqpoint{1.262283in}{0.757944in}}%
\pgfpathlineto{\pgfqpoint{1.275641in}{0.769388in}}%
\pgfpathlineto{\pgfqpoint{1.275623in}{0.774548in}}%
\pgfpathlineto{\pgfqpoint{1.274795in}{0.779674in}}%
\pgfpathlineto{\pgfqpoint{1.261536in}{0.767522in}}%
\pgfpathlineto{\pgfqpoint{1.248288in}{0.755381in}}%
\pgfpathlineto{\pgfqpoint{1.248962in}{0.750962in}}%
\pgfpathlineto{\pgfqpoint{1.248943in}{0.746514in}}%
\pgfpathclose%
\pgfusepath{fill}%
\end{pgfscope}%
\begin{pgfscope}%
\pgfpathrectangle{\pgfqpoint{0.407592in}{0.222904in}}{\pgfqpoint{1.411163in}{1.411163in}}%
\pgfusepath{clip}%
\pgfsetbuttcap%
\pgfsetroundjoin%
\definecolor{currentfill}{rgb}{0.749020,0.772549,0.792157}%
\pgfsetfillcolor{currentfill}%
\pgfsetfillopacity{0.160000}%
\pgfsetlinewidth{0.000000pt}%
\definecolor{currentstroke}{rgb}{0.000000,0.000000,0.000000}%
\pgfsetstrokecolor{currentstroke}%
\pgfsetdash{}{0pt}%
\pgfpathmoveto{\pgfqpoint{1.793060in}{1.308251in}}%
\pgfpathlineto{\pgfqpoint{1.806244in}{1.321339in}}%
\pgfpathlineto{\pgfqpoint{1.819435in}{1.334434in}}%
\pgfpathlineto{\pgfqpoint{1.803382in}{1.350950in}}%
\pgfpathlineto{\pgfqpoint{1.784494in}{1.366922in}}%
\pgfpathlineto{\pgfqpoint{1.772076in}{1.353393in}}%
\pgfpathlineto{\pgfqpoint{1.759659in}{1.339865in}}%
\pgfpathlineto{\pgfqpoint{1.777732in}{1.324319in}}%
\pgfpathlineto{\pgfqpoint{1.793060in}{1.308251in}}%
\pgfpathclose%
\pgfusepath{fill}%
\end{pgfscope}%
\begin{pgfscope}%
\pgfpathrectangle{\pgfqpoint{0.407592in}{0.222904in}}{\pgfqpoint{1.411163in}{1.411163in}}%
\pgfusepath{clip}%
\pgfsetbuttcap%
\pgfsetroundjoin%
\definecolor{currentfill}{rgb}{0.749020,0.772549,0.792157}%
\pgfsetfillcolor{currentfill}%
\pgfsetfillopacity{0.160000}%
\pgfsetlinewidth{0.000000pt}%
\definecolor{currentstroke}{rgb}{0.000000,0.000000,0.000000}%
\pgfsetstrokecolor{currentstroke}%
\pgfsetdash{}{0pt}%
\pgfpathmoveto{\pgfqpoint{1.819435in}{1.334434in}}%
\pgfpathlineto{\pgfqpoint{1.832631in}{1.347535in}}%
\pgfpathlineto{\pgfqpoint{1.816212in}{1.364270in}}%
\pgfpathlineto{\pgfqpoint{1.796912in}{1.380452in}}%
\pgfpathlineto{\pgfqpoint{1.784494in}{1.366922in}}%
\pgfpathlineto{\pgfqpoint{1.803382in}{1.350950in}}%
\pgfusepath{fill}%
\end{pgfscope}%
\begin{pgfscope}%
\pgfpathrectangle{\pgfqpoint{0.407592in}{0.222904in}}{\pgfqpoint{1.411163in}{1.411163in}}%
\pgfusepath{clip}%
\pgfsetbuttcap%
\pgfsetroundjoin%
\definecolor{currentfill}{rgb}{0.749020,0.772549,0.792157}%
\pgfsetfillcolor{currentfill}%
\pgfsetfillopacity{0.160000}%
\pgfsetlinewidth{0.000000pt}%
\definecolor{currentstroke}{rgb}{0.000000,0.000000,0.000000}%
\pgfsetstrokecolor{currentstroke}%
\pgfsetdash{}{0pt}%
\pgfpathmoveto{\pgfqpoint{0.932871in}{0.754498in}}%
\pgfpathlineto{\pgfqpoint{0.919678in}{0.766568in}}%
\pgfpathlineto{\pgfqpoint{0.906473in}{0.778649in}}%
\pgfpathlineto{\pgfqpoint{0.904718in}{0.773515in}}%
\pgfpathlineto{\pgfqpoint{0.903768in}{0.768351in}}%
\pgfpathlineto{\pgfqpoint{0.917205in}{0.756978in}}%
\pgfpathlineto{\pgfqpoint{0.930624in}{0.745620in}}%
\pgfpathlineto{\pgfqpoint{0.931403in}{0.750071in}}%
\pgfpathlineto{\pgfqpoint{0.932871in}{0.754498in}}%
\pgfpathclose%
\pgfusepath{fill}%
\end{pgfscope}%
\begin{pgfscope}%
\pgfpathrectangle{\pgfqpoint{0.407592in}{0.222904in}}{\pgfqpoint{1.411163in}{1.411163in}}%
\pgfusepath{clip}%
\pgfsetbuttcap%
\pgfsetroundjoin%
\definecolor{currentfill}{rgb}{0.749020,0.772549,0.792157}%
\pgfsetfillcolor{currentfill}%
\pgfsetfillopacity{0.160000}%
\pgfsetlinewidth{0.000000pt}%
\definecolor{currentstroke}{rgb}{0.000000,0.000000,0.000000}%
\pgfsetstrokecolor{currentstroke}%
\pgfsetdash{}{0pt}%
\pgfpathmoveto{\pgfqpoint{0.509947in}{1.255902in}}%
\pgfpathlineto{\pgfqpoint{0.497828in}{1.269380in}}%
\pgfpathlineto{\pgfqpoint{0.485708in}{1.282860in}}%
\pgfpathlineto{\pgfqpoint{0.466153in}{1.268099in}}%
\pgfpathlineto{\pgfqpoint{0.449066in}{1.252869in}}%
\pgfpathlineto{\pgfqpoint{0.462051in}{1.239861in}}%
\pgfpathlineto{\pgfqpoint{0.475029in}{1.226859in}}%
\pgfpathlineto{\pgfqpoint{0.491298in}{1.241605in}}%
\pgfpathlineto{\pgfqpoint{0.509947in}{1.255902in}}%
\pgfpathclose%
\pgfusepath{fill}%
\end{pgfscope}%
\begin{pgfscope}%
\pgfpathrectangle{\pgfqpoint{0.407592in}{0.222904in}}{\pgfqpoint{1.411163in}{1.411163in}}%
\pgfusepath{clip}%
\pgfsetbuttcap%
\pgfsetroundjoin%
\definecolor{currentfill}{rgb}{0.749020,0.772549,0.792157}%
\pgfsetfillcolor{currentfill}%
\pgfsetfillopacity{0.160000}%
\pgfsetlinewidth{0.000000pt}%
\definecolor{currentstroke}{rgb}{0.000000,0.000000,0.000000}%
\pgfsetstrokecolor{currentstroke}%
\pgfsetdash{}{0pt}%
\pgfpathmoveto{\pgfqpoint{0.835726in}{0.865502in}}%
\pgfpathlineto{\pgfqpoint{0.822932in}{0.878319in}}%
\pgfpathlineto{\pgfqpoint{0.810133in}{0.891142in}}%
\pgfpathlineto{\pgfqpoint{0.804651in}{0.883488in}}%
\pgfpathlineto{\pgfqpoint{0.800386in}{0.875707in}}%
\pgfpathlineto{\pgfqpoint{0.813690in}{0.863535in}}%
\pgfpathlineto{\pgfqpoint{0.826982in}{0.851374in}}%
\pgfpathlineto{\pgfqpoint{0.830800in}{0.858495in}}%
\pgfpathlineto{\pgfqpoint{0.835726in}{0.865502in}}%
\pgfpathclose%
\pgfusepath{fill}%
\end{pgfscope}%
\begin{pgfscope}%
\pgfpathrectangle{\pgfqpoint{0.407592in}{0.222904in}}{\pgfqpoint{1.411163in}{1.411163in}}%
\pgfusepath{clip}%
\pgfsetbuttcap%
\pgfsetroundjoin%
\definecolor{currentfill}{rgb}{0.749020,0.772549,0.792157}%
\pgfsetfillcolor{currentfill}%
\pgfsetfillopacity{0.160000}%
\pgfsetlinewidth{0.000000pt}%
\definecolor{currentstroke}{rgb}{0.000000,0.000000,0.000000}%
\pgfsetstrokecolor{currentstroke}%
\pgfsetdash{}{0pt}%
\pgfpathmoveto{\pgfqpoint{1.132327in}{0.642759in}}%
\pgfpathlineto{\pgfqpoint{1.142935in}{0.650879in}}%
\pgfpathlineto{\pgfqpoint{1.153575in}{0.659024in}}%
\pgfpathlineto{\pgfqpoint{1.156589in}{0.660946in}}%
\pgfpathlineto{\pgfqpoint{1.145447in}{0.652482in}}%
\pgfpathlineto{\pgfqpoint{1.134338in}{0.644042in}}%
\pgfusepath{fill}%
\end{pgfscope}%
\begin{pgfscope}%
\pgfpathrectangle{\pgfqpoint{0.407592in}{0.222904in}}{\pgfqpoint{1.411163in}{1.411163in}}%
\pgfusepath{clip}%
\pgfsetbuttcap%
\pgfsetroundjoin%
\definecolor{currentfill}{rgb}{0.749020,0.772549,0.792157}%
\pgfsetfillcolor{currentfill}%
\pgfsetfillopacity{0.160000}%
\pgfsetlinewidth{0.000000pt}%
\definecolor{currentstroke}{rgb}{0.000000,0.000000,0.000000}%
\pgfsetstrokecolor{currentstroke}%
\pgfsetdash{}{0pt}%
\pgfpathmoveto{\pgfqpoint{0.485708in}{1.282860in}}%
\pgfpathlineto{\pgfqpoint{0.473587in}{1.296340in}}%
\pgfpathlineto{\pgfqpoint{0.461465in}{1.309822in}}%
\pgfpathlineto{\pgfqpoint{0.440993in}{1.294609in}}%
\pgfpathlineto{\pgfqpoint{0.423077in}{1.278906in}}%
\pgfpathlineto{\pgfqpoint{0.436075in}{1.265884in}}%
\pgfpathlineto{\pgfqpoint{0.449066in}{1.252869in}}%
\pgfpathlineto{\pgfqpoint{0.466153in}{1.268099in}}%
\pgfpathlineto{\pgfqpoint{0.485708in}{1.282860in}}%
\pgfpathclose%
\pgfusepath{fill}%
\end{pgfscope}%
\begin{pgfscope}%
\pgfpathrectangle{\pgfqpoint{0.407592in}{0.222904in}}{\pgfqpoint{1.411163in}{1.411163in}}%
\pgfusepath{clip}%
\pgfsetbuttcap%
\pgfsetroundjoin%
\definecolor{currentfill}{rgb}{0.749020,0.772549,0.792157}%
\pgfsetfillcolor{currentfill}%
\pgfsetfillopacity{0.160000}%
\pgfsetlinewidth{0.000000pt}%
\definecolor{currentstroke}{rgb}{0.000000,0.000000,0.000000}%
\pgfsetstrokecolor{currentstroke}%
\pgfsetdash{}{0pt}%
\pgfpathmoveto{\pgfqpoint{1.381300in}{0.877277in}}%
\pgfpathlineto{\pgfqpoint{1.394666in}{0.889526in}}%
\pgfpathlineto{\pgfqpoint{1.408045in}{0.901787in}}%
\pgfpathlineto{\pgfqpoint{1.404899in}{0.910193in}}%
\pgfpathlineto{\pgfqpoint{1.400404in}{0.918452in}}%
\pgfpathlineto{\pgfqpoint{1.387411in}{0.905553in}}%
\pgfpathlineto{\pgfqpoint{1.374423in}{0.892660in}}%
\pgfpathlineto{\pgfqpoint{1.378481in}{0.885035in}}%
\pgfpathlineto{\pgfqpoint{1.381300in}{0.877277in}}%
\pgfpathclose%
\pgfusepath{fill}%
\end{pgfscope}%
\begin{pgfscope}%
\pgfpathrectangle{\pgfqpoint{0.407592in}{0.222904in}}{\pgfqpoint{1.411163in}{1.411163in}}%
\pgfusepath{clip}%
\pgfsetbuttcap%
\pgfsetroundjoin%
\definecolor{currentfill}{rgb}{0.749020,0.772549,0.792157}%
\pgfsetfillcolor{currentfill}%
\pgfsetfillopacity{0.160000}%
\pgfsetlinewidth{0.000000pt}%
\definecolor{currentstroke}{rgb}{0.000000,0.000000,0.000000}%
\pgfsetstrokecolor{currentstroke}%
\pgfsetdash{}{0pt}%
\pgfpathmoveto{\pgfqpoint{1.043817in}{0.643781in}}%
\pgfpathlineto{\pgfqpoint{1.032127in}{0.652155in}}%
\pgfpathlineto{\pgfqpoint{1.020402in}{0.660554in}}%
\pgfpathlineto{\pgfqpoint{1.023079in}{0.658648in}}%
\pgfpathlineto{\pgfqpoint{1.026055in}{0.656828in}}%
\pgfpathlineto{\pgfqpoint{1.036840in}{0.649047in}}%
\pgfpathlineto{\pgfqpoint{1.047590in}{0.641293in}}%
\pgfpathlineto{\pgfqpoint{1.045605in}{0.642508in}}%
\pgfpathlineto{\pgfqpoint{1.043817in}{0.643781in}}%
\pgfpathclose%
\pgfusepath{fill}%
\end{pgfscope}%
\begin{pgfscope}%
\pgfpathrectangle{\pgfqpoint{0.407592in}{0.222904in}}{\pgfqpoint{1.411163in}{1.411163in}}%
\pgfusepath{clip}%
\pgfsetbuttcap%
\pgfsetroundjoin%
\definecolor{currentfill}{rgb}{0.749020,0.772549,0.792157}%
\pgfsetfillcolor{currentfill}%
\pgfsetfillopacity{0.160000}%
\pgfsetlinewidth{0.000000pt}%
\definecolor{currentstroke}{rgb}{0.000000,0.000000,0.000000}%
\pgfsetstrokecolor{currentstroke}%
\pgfsetdash{}{0pt}%
\pgfpathmoveto{\pgfqpoint{0.461465in}{1.309822in}}%
\pgfpathlineto{\pgfqpoint{0.449341in}{1.323306in}}%
\pgfpathlineto{\pgfqpoint{0.437216in}{1.336790in}}%
\pgfpathlineto{\pgfqpoint{0.415818in}{1.321136in}}%
\pgfpathlineto{\pgfqpoint{0.397060in}{1.304970in}}%
\pgfpathlineto{\pgfqpoint{0.410072in}{1.291934in}}%
\pgfpathlineto{\pgfqpoint{0.423077in}{1.278906in}}%
\pgfpathlineto{\pgfqpoint{0.440993in}{1.294609in}}%
\pgfpathlineto{\pgfqpoint{0.461465in}{1.309822in}}%
\pgfpathclose%
\pgfusepath{fill}%
\end{pgfscope}%
\begin{pgfscope}%
\pgfpathrectangle{\pgfqpoint{0.407592in}{0.222904in}}{\pgfqpoint{1.411163in}{1.411163in}}%
\pgfusepath{clip}%
\pgfsetbuttcap%
\pgfsetroundjoin%
\definecolor{currentfill}{rgb}{0.749020,0.772549,0.792157}%
\pgfsetfillcolor{currentfill}%
\pgfsetfillopacity{0.160000}%
\pgfsetlinewidth{0.000000pt}%
\definecolor{currentstroke}{rgb}{0.000000,0.000000,0.000000}%
\pgfsetstrokecolor{currentstroke}%
\pgfsetdash{}{0pt}%
\pgfpathmoveto{\pgfqpoint{1.127755in}{0.640367in}}%
\pgfpathlineto{\pgfqpoint{1.137218in}{0.647890in}}%
\pgfpathlineto{\pgfqpoint{1.146714in}{0.655440in}}%
\pgfpathlineto{\pgfqpoint{1.150279in}{0.657186in}}%
\pgfpathlineto{\pgfqpoint{1.153575in}{0.659024in}}%
\pgfpathlineto{\pgfqpoint{1.142935in}{0.650879in}}%
\pgfpathlineto{\pgfqpoint{1.132327in}{0.642759in}}%
\pgfpathlineto{\pgfqpoint{1.130130in}{0.641532in}}%
\pgfpathlineto{\pgfqpoint{1.127755in}{0.640367in}}%
\pgfpathclose%
\pgfusepath{fill}%
\end{pgfscope}%
\begin{pgfscope}%
\pgfpathrectangle{\pgfqpoint{0.407592in}{0.222904in}}{\pgfqpoint{1.411163in}{1.411163in}}%
\pgfusepath{clip}%
\pgfsetbuttcap%
\pgfsetroundjoin%
\definecolor{currentfill}{rgb}{0.749020,0.772549,0.792157}%
\pgfsetfillcolor{currentfill}%
\pgfsetfillopacity{0.160000}%
\pgfsetlinewidth{0.000000pt}%
\definecolor{currentstroke}{rgb}{0.000000,0.000000,0.000000}%
\pgfsetstrokecolor{currentstroke}%
\pgfsetdash{}{0pt}%
\pgfpathmoveto{\pgfqpoint{1.275641in}{0.769388in}}%
\pgfpathlineto{\pgfqpoint{1.289016in}{0.780848in}}%
\pgfpathlineto{\pgfqpoint{1.302409in}{0.792323in}}%
\pgfpathlineto{\pgfqpoint{1.302343in}{0.798187in}}%
\pgfpathlineto{\pgfqpoint{1.301350in}{0.804009in}}%
\pgfpathlineto{\pgfqpoint{1.288067in}{0.791836in}}%
\pgfpathlineto{\pgfqpoint{1.274795in}{0.779674in}}%
\pgfpathlineto{\pgfqpoint{1.275623in}{0.774548in}}%
\pgfpathlineto{\pgfqpoint{1.275641in}{0.769388in}}%
\pgfpathclose%
\pgfusepath{fill}%
\end{pgfscope}%
\begin{pgfscope}%
\pgfpathrectangle{\pgfqpoint{0.407592in}{0.222904in}}{\pgfqpoint{1.411163in}{1.411163in}}%
\pgfusepath{clip}%
\pgfsetbuttcap%
\pgfsetroundjoin%
\definecolor{currentfill}{rgb}{0.749020,0.772549,0.792157}%
\pgfsetfillcolor{currentfill}%
\pgfsetfillopacity{0.160000}%
\pgfsetlinewidth{0.000000pt}%
\definecolor{currentstroke}{rgb}{0.000000,0.000000,0.000000}%
\pgfsetstrokecolor{currentstroke}%
\pgfsetdash{}{0pt}%
\pgfpathmoveto{\pgfqpoint{0.810133in}{0.891142in}}%
\pgfpathlineto{\pgfqpoint{0.797326in}{0.903972in}}%
\pgfpathlineto{\pgfqpoint{0.784513in}{0.916808in}}%
\pgfpathlineto{\pgfqpoint{0.778464in}{0.908517in}}%
\pgfpathlineto{\pgfqpoint{0.773739in}{0.900086in}}%
\pgfpathlineto{\pgfqpoint{0.787069in}{0.887891in}}%
\pgfpathlineto{\pgfqpoint{0.800386in}{0.875707in}}%
\pgfpathlineto{\pgfqpoint{0.804651in}{0.883488in}}%
\pgfpathlineto{\pgfqpoint{0.810133in}{0.891142in}}%
\pgfpathclose%
\pgfusepath{fill}%
\end{pgfscope}%
\begin{pgfscope}%
\pgfpathrectangle{\pgfqpoint{0.407592in}{0.222904in}}{\pgfqpoint{1.411163in}{1.411163in}}%
\pgfusepath{clip}%
\pgfsetbuttcap%
\pgfsetroundjoin%
\definecolor{currentfill}{rgb}{0.749020,0.772549,0.792157}%
\pgfsetfillcolor{currentfill}%
\pgfsetfillopacity{0.160000}%
\pgfsetlinewidth{0.000000pt}%
\definecolor{currentstroke}{rgb}{0.000000,0.000000,0.000000}%
\pgfsetstrokecolor{currentstroke}%
\pgfsetdash{}{0pt}%
\pgfpathmoveto{\pgfqpoint{1.161718in}{0.665008in}}%
\pgfpathlineto{\pgfqpoint{1.173741in}{0.674169in}}%
\pgfpathlineto{\pgfqpoint{1.185795in}{0.683354in}}%
\pgfpathlineto{\pgfqpoint{1.188574in}{0.686174in}}%
\pgfpathlineto{\pgfqpoint{1.190916in}{0.689056in}}%
\pgfpathlineto{\pgfqpoint{1.178234in}{0.679166in}}%
\pgfpathlineto{\pgfqpoint{1.165580in}{0.669297in}}%
\pgfpathlineto{\pgfqpoint{1.163812in}{0.667129in}}%
\pgfpathlineto{\pgfqpoint{1.161718in}{0.665008in}}%
\pgfpathclose%
\pgfusepath{fill}%
\end{pgfscope}%
\begin{pgfscope}%
\pgfpathrectangle{\pgfqpoint{0.407592in}{0.222904in}}{\pgfqpoint{1.411163in}{1.411163in}}%
\pgfusepath{clip}%
\pgfsetbuttcap%
\pgfsetroundjoin%
\definecolor{currentfill}{rgb}{0.749020,0.772549,0.792157}%
\pgfsetfillcolor{currentfill}%
\pgfsetfillopacity{0.160000}%
\pgfsetlinewidth{0.000000pt}%
\definecolor{currentstroke}{rgb}{0.000000,0.000000,0.000000}%
\pgfsetstrokecolor{currentstroke}%
\pgfsetdash{}{0pt}%
\pgfpathmoveto{\pgfqpoint{0.437216in}{1.336790in}}%
\pgfpathlineto{\pgfqpoint{0.425090in}{1.350276in}}%
\pgfpathlineto{\pgfqpoint{0.412963in}{1.363764in}}%
\pgfpathlineto{\pgfqpoint{0.390627in}{1.347678in}}%
\pgfpathlineto{\pgfqpoint{0.371018in}{1.331060in}}%
\pgfpathlineto{\pgfqpoint{0.384042in}{1.318011in}}%
\pgfpathlineto{\pgfqpoint{0.397060in}{1.304970in}}%
\pgfpathlineto{\pgfqpoint{0.415818in}{1.321136in}}%
\pgfpathlineto{\pgfqpoint{0.437216in}{1.336790in}}%
\pgfpathclose%
\pgfusepath{fill}%
\end{pgfscope}%
\begin{pgfscope}%
\pgfpathrectangle{\pgfqpoint{0.407592in}{0.222904in}}{\pgfqpoint{1.411163in}{1.411163in}}%
\pgfusepath{clip}%
\pgfsetbuttcap%
\pgfsetroundjoin%
\definecolor{currentfill}{rgb}{0.749020,0.772549,0.792157}%
\pgfsetfillcolor{currentfill}%
\pgfsetfillopacity{0.160000}%
\pgfsetlinewidth{0.000000pt}%
\definecolor{currentstroke}{rgb}{0.000000,0.000000,0.000000}%
\pgfsetstrokecolor{currentstroke}%
\pgfsetdash{}{0pt}%
\pgfpathmoveto{\pgfqpoint{1.220488in}{0.716247in}}%
\pgfpathlineto{\pgfqpoint{1.233640in}{0.726920in}}%
\pgfpathlineto{\pgfqpoint{1.246815in}{0.737613in}}%
\pgfpathlineto{\pgfqpoint{1.248227in}{0.742058in}}%
\pgfpathlineto{\pgfqpoint{1.248943in}{0.746514in}}%
\pgfpathlineto{\pgfqpoint{1.235621in}{0.735100in}}%
\pgfpathlineto{\pgfqpoint{1.222316in}{0.723701in}}%
\pgfpathlineto{\pgfqpoint{1.221691in}{0.719968in}}%
\pgfpathlineto{\pgfqpoint{1.220488in}{0.716247in}}%
\pgfpathclose%
\pgfusepath{fill}%
\end{pgfscope}%
\begin{pgfscope}%
\pgfpathrectangle{\pgfqpoint{0.407592in}{0.222904in}}{\pgfqpoint{1.411163in}{1.411163in}}%
\pgfusepath{clip}%
\pgfsetbuttcap%
\pgfsetroundjoin%
\definecolor{currentfill}{rgb}{0.749020,0.772549,0.792157}%
\pgfsetfillcolor{currentfill}%
\pgfsetfillopacity{0.160000}%
\pgfsetlinewidth{0.000000pt}%
\definecolor{currentstroke}{rgb}{0.000000,0.000000,0.000000}%
\pgfsetstrokecolor{currentstroke}%
\pgfsetdash{}{0pt}%
\pgfpathmoveto{\pgfqpoint{1.012872in}{0.668859in}}%
\pgfpathlineto{\pgfqpoint{0.999880in}{0.678655in}}%
\pgfpathlineto{\pgfqpoint{0.986858in}{0.688474in}}%
\pgfpathlineto{\pgfqpoint{0.988691in}{0.685603in}}%
\pgfpathlineto{\pgfqpoint{0.990974in}{0.682797in}}%
\pgfpathlineto{\pgfqpoint{1.003495in}{0.673681in}}%
\pgfpathlineto{\pgfqpoint{1.015983in}{0.664589in}}%
\pgfpathlineto{\pgfqpoint{1.014260in}{0.666699in}}%
\pgfpathlineto{\pgfqpoint{1.012872in}{0.668859in}}%
\pgfpathclose%
\pgfusepath{fill}%
\end{pgfscope}%
\begin{pgfscope}%
\pgfpathrectangle{\pgfqpoint{0.407592in}{0.222904in}}{\pgfqpoint{1.411163in}{1.411163in}}%
\pgfusepath{clip}%
\pgfsetbuttcap%
\pgfsetroundjoin%
\definecolor{currentfill}{rgb}{0.749020,0.772549,0.792157}%
\pgfsetfillcolor{currentfill}%
\pgfsetfillopacity{0.160000}%
\pgfsetlinewidth{0.000000pt}%
\definecolor{currentstroke}{rgb}{0.000000,0.000000,0.000000}%
\pgfsetstrokecolor{currentstroke}%
\pgfsetdash{}{0pt}%
\pgfpathmoveto{\pgfqpoint{1.408045in}{0.901787in}}%
\pgfpathlineto{\pgfqpoint{1.421436in}{0.914058in}}%
\pgfpathlineto{\pgfqpoint{1.434838in}{0.926341in}}%
\pgfpathlineto{\pgfqpoint{1.431354in}{0.935386in}}%
\pgfpathlineto{\pgfqpoint{1.426410in}{0.944268in}}%
\pgfpathlineto{\pgfqpoint{1.413404in}{0.931357in}}%
\pgfpathlineto{\pgfqpoint{1.400404in}{0.918452in}}%
\pgfpathlineto{\pgfqpoint{1.404899in}{0.910193in}}%
\pgfpathlineto{\pgfqpoint{1.408045in}{0.901787in}}%
\pgfpathclose%
\pgfusepath{fill}%
\end{pgfscope}%
\begin{pgfscope}%
\pgfpathrectangle{\pgfqpoint{0.407592in}{0.222904in}}{\pgfqpoint{1.411163in}{1.411163in}}%
\pgfusepath{clip}%
\pgfsetbuttcap%
\pgfsetroundjoin%
\definecolor{currentfill}{rgb}{0.749020,0.772549,0.792157}%
\pgfsetfillcolor{currentfill}%
\pgfsetfillopacity{0.160000}%
\pgfsetlinewidth{0.000000pt}%
\definecolor{currentstroke}{rgb}{0.000000,0.000000,0.000000}%
\pgfsetstrokecolor{currentstroke}%
\pgfsetdash{}{0pt}%
\pgfpathmoveto{\pgfqpoint{0.412963in}{1.363764in}}%
\pgfpathlineto{\pgfqpoint{0.400835in}{1.377253in}}%
\pgfpathlineto{\pgfqpoint{0.378026in}{1.360955in}}%
\pgfpathlineto{\pgfqpoint{0.357986in}{1.344115in}}%
\pgfpathlineto{\pgfqpoint{0.371018in}{1.331060in}}%
\pgfpathlineto{\pgfqpoint{0.390627in}{1.347678in}}%
\pgfusepath{fill}%
\end{pgfscope}%
\begin{pgfscope}%
\pgfpathrectangle{\pgfqpoint{0.407592in}{0.222904in}}{\pgfqpoint{1.411163in}{1.411163in}}%
\pgfusepath{clip}%
\pgfsetbuttcap%
\pgfsetroundjoin%
\definecolor{currentfill}{rgb}{0.749020,0.772549,0.792157}%
\pgfsetfillcolor{currentfill}%
\pgfsetfillopacity{0.160000}%
\pgfsetlinewidth{0.000000pt}%
\definecolor{currentstroke}{rgb}{0.000000,0.000000,0.000000}%
\pgfsetstrokecolor{currentstroke}%
\pgfsetdash{}{0pt}%
\pgfpathmoveto{\pgfqpoint{0.906473in}{0.778649in}}%
\pgfpathlineto{\pgfqpoint{0.893256in}{0.790742in}}%
\pgfpathlineto{\pgfqpoint{0.880026in}{0.802846in}}%
\pgfpathlineto{\pgfqpoint{0.877972in}{0.797013in}}%
\pgfpathlineto{\pgfqpoint{0.876838in}{0.791143in}}%
\pgfpathlineto{\pgfqpoint{0.890312in}{0.779739in}}%
\pgfpathlineto{\pgfqpoint{0.903768in}{0.768351in}}%
\pgfpathlineto{\pgfqpoint{0.904718in}{0.773515in}}%
\pgfpathlineto{\pgfqpoint{0.906473in}{0.778649in}}%
\pgfpathclose%
\pgfusepath{fill}%
\end{pgfscope}%
\begin{pgfscope}%
\pgfpathrectangle{\pgfqpoint{0.407592in}{0.222904in}}{\pgfqpoint{1.411163in}{1.411163in}}%
\pgfusepath{clip}%
\pgfsetbuttcap%
\pgfsetroundjoin%
\definecolor{currentfill}{rgb}{0.749020,0.772549,0.792157}%
\pgfsetfillcolor{currentfill}%
\pgfsetfillopacity{0.160000}%
\pgfsetlinewidth{0.000000pt}%
\definecolor{currentstroke}{rgb}{0.000000,0.000000,0.000000}%
\pgfsetstrokecolor{currentstroke}%
\pgfsetdash{}{0pt}%
\pgfpathmoveto{\pgfqpoint{0.957406in}{0.722952in}}%
\pgfpathlineto{\pgfqpoint{0.944024in}{0.734278in}}%
\pgfpathlineto{\pgfqpoint{0.930624in}{0.745620in}}%
\pgfpathlineto{\pgfqpoint{0.930541in}{0.741165in}}%
\pgfpathlineto{\pgfqpoint{0.931156in}{0.736724in}}%
\pgfpathlineto{\pgfqpoint{0.944544in}{0.726104in}}%
\pgfpathlineto{\pgfqpoint{0.957907in}{0.715503in}}%
\pgfpathlineto{\pgfqpoint{0.957366in}{0.719220in}}%
\pgfpathlineto{\pgfqpoint{0.957406in}{0.722952in}}%
\pgfpathclose%
\pgfusepath{fill}%
\end{pgfscope}%
\begin{pgfscope}%
\pgfpathrectangle{\pgfqpoint{0.407592in}{0.222904in}}{\pgfqpoint{1.411163in}{1.411163in}}%
\pgfusepath{clip}%
\pgfsetbuttcap%
\pgfsetroundjoin%
\definecolor{currentfill}{rgb}{0.749020,0.772549,0.792157}%
\pgfsetfillcolor{currentfill}%
\pgfsetfillopacity{0.160000}%
\pgfsetlinewidth{0.000000pt}%
\definecolor{currentstroke}{rgb}{0.000000,0.000000,0.000000}%
\pgfsetstrokecolor{currentstroke}%
\pgfsetdash{}{0pt}%
\pgfpathmoveto{\pgfqpoint{1.047590in}{0.641293in}}%
\pgfpathlineto{\pgfqpoint{1.036840in}{0.649047in}}%
\pgfpathlineto{\pgfqpoint{1.026055in}{0.656828in}}%
\pgfpathlineto{\pgfqpoint{1.029315in}{0.655102in}}%
\pgfpathlineto{\pgfqpoint{1.032846in}{0.653477in}}%
\pgfpathlineto{\pgfqpoint{1.042498in}{0.646254in}}%
\pgfpathlineto{\pgfqpoint{1.052114in}{0.639058in}}%
\pgfpathlineto{\pgfqpoint{1.049763in}{0.640141in}}%
\pgfpathlineto{\pgfqpoint{1.047590in}{0.641293in}}%
\pgfpathclose%
\pgfusepath{fill}%
\end{pgfscope}%
\begin{pgfscope}%
\pgfpathrectangle{\pgfqpoint{0.407592in}{0.222904in}}{\pgfqpoint{1.411163in}{1.411163in}}%
\pgfusepath{clip}%
\pgfsetbuttcap%
\pgfsetroundjoin%
\definecolor{currentfill}{rgb}{0.749020,0.772549,0.792157}%
\pgfsetfillcolor{currentfill}%
\pgfsetfillopacity{0.160000}%
\pgfsetlinewidth{0.000000pt}%
\definecolor{currentstroke}{rgb}{0.000000,0.000000,0.000000}%
\pgfsetstrokecolor{currentstroke}%
\pgfsetdash{}{0pt}%
\pgfpathmoveto{\pgfqpoint{1.190916in}{0.689056in}}%
\pgfpathlineto{\pgfqpoint{1.203625in}{0.698968in}}%
\pgfpathlineto{\pgfqpoint{1.216362in}{0.708901in}}%
\pgfpathlineto{\pgfqpoint{1.218710in}{0.712552in}}%
\pgfpathlineto{\pgfqpoint{1.220488in}{0.716247in}}%
\pgfpathlineto{\pgfqpoint{1.207360in}{0.705592in}}%
\pgfpathlineto{\pgfqpoint{1.194254in}{0.694956in}}%
\pgfpathlineto{\pgfqpoint{1.192812in}{0.691988in}}%
\pgfpathlineto{\pgfqpoint{1.190916in}{0.689056in}}%
\pgfpathclose%
\pgfusepath{fill}%
\end{pgfscope}%
\begin{pgfscope}%
\pgfpathrectangle{\pgfqpoint{0.407592in}{0.222904in}}{\pgfqpoint{1.411163in}{1.411163in}}%
\pgfusepath{clip}%
\pgfsetbuttcap%
\pgfsetroundjoin%
\definecolor{currentfill}{rgb}{0.749020,0.772549,0.792157}%
\pgfsetfillcolor{currentfill}%
\pgfsetfillopacity{0.160000}%
\pgfsetlinewidth{0.000000pt}%
\definecolor{currentstroke}{rgb}{0.000000,0.000000,0.000000}%
\pgfsetstrokecolor{currentstroke}%
\pgfsetdash{}{0pt}%
\pgfpathmoveto{\pgfqpoint{0.984562in}{0.694358in}}%
\pgfpathlineto{\pgfqpoint{0.971246in}{0.704921in}}%
\pgfpathlineto{\pgfqpoint{0.957907in}{0.715503in}}%
\pgfpathlineto{\pgfqpoint{0.959028in}{0.711816in}}%
\pgfpathlineto{\pgfqpoint{0.960727in}{0.708176in}}%
\pgfpathlineto{\pgfqpoint{0.973807in}{0.698314in}}%
\pgfpathlineto{\pgfqpoint{0.986858in}{0.688474in}}%
\pgfpathlineto{\pgfqpoint{0.985479in}{0.691396in}}%
\pgfpathlineto{\pgfqpoint{0.984562in}{0.694358in}}%
\pgfpathclose%
\pgfusepath{fill}%
\end{pgfscope}%
\begin{pgfscope}%
\pgfpathrectangle{\pgfqpoint{0.407592in}{0.222904in}}{\pgfqpoint{1.411163in}{1.411163in}}%
\pgfusepath{clip}%
\pgfsetbuttcap%
\pgfsetroundjoin%
\definecolor{currentfill}{rgb}{0.749020,0.772549,0.792157}%
\pgfsetfillcolor{currentfill}%
\pgfsetfillopacity{0.160000}%
\pgfsetlinewidth{0.000000pt}%
\definecolor{currentstroke}{rgb}{0.000000,0.000000,0.000000}%
\pgfsetstrokecolor{currentstroke}%
\pgfsetdash{}{0pt}%
\pgfpathmoveto{\pgfqpoint{0.784513in}{0.916808in}}%
\pgfpathlineto{\pgfqpoint{0.771694in}{0.929651in}}%
\pgfpathlineto{\pgfqpoint{0.758868in}{0.942501in}}%
\pgfpathlineto{\pgfqpoint{0.752240in}{0.933583in}}%
\pgfpathlineto{\pgfqpoint{0.747043in}{0.924510in}}%
\pgfpathlineto{\pgfqpoint{0.760397in}{0.912292in}}%
\pgfpathlineto{\pgfqpoint{0.773739in}{0.900086in}}%
\pgfpathlineto{\pgfqpoint{0.778464in}{0.908517in}}%
\pgfpathlineto{\pgfqpoint{0.784513in}{0.916808in}}%
\pgfpathclose%
\pgfusepath{fill}%
\end{pgfscope}%
\begin{pgfscope}%
\pgfpathrectangle{\pgfqpoint{0.407592in}{0.222904in}}{\pgfqpoint{1.411163in}{1.411163in}}%
\pgfusepath{clip}%
\pgfsetbuttcap%
\pgfsetroundjoin%
\definecolor{currentfill}{rgb}{0.749020,0.772549,0.792157}%
\pgfsetfillcolor{currentfill}%
\pgfsetfillopacity{0.160000}%
\pgfsetlinewidth{0.000000pt}%
\definecolor{currentstroke}{rgb}{0.000000,0.000000,0.000000}%
\pgfsetstrokecolor{currentstroke}%
\pgfsetdash{}{0pt}%
\pgfpathmoveto{\pgfqpoint{1.122516in}{0.638244in}}%
\pgfpathlineto{\pgfqpoint{1.130664in}{0.645237in}}%
\pgfpathlineto{\pgfqpoint{1.138843in}{0.652257in}}%
\pgfpathlineto{\pgfqpoint{1.142897in}{0.653795in}}%
\pgfpathlineto{\pgfqpoint{1.146714in}{0.655440in}}%
\pgfpathlineto{\pgfqpoint{1.137218in}{0.647890in}}%
\pgfpathlineto{\pgfqpoint{1.127755in}{0.640367in}}%
\pgfpathlineto{\pgfqpoint{1.125213in}{0.639269in}}%
\pgfpathlineto{\pgfqpoint{1.122516in}{0.638244in}}%
\pgfpathclose%
\pgfusepath{fill}%
\end{pgfscope}%
\begin{pgfscope}%
\pgfpathrectangle{\pgfqpoint{0.407592in}{0.222904in}}{\pgfqpoint{1.411163in}{1.411163in}}%
\pgfusepath{clip}%
\pgfsetbuttcap%
\pgfsetroundjoin%
\definecolor{currentfill}{rgb}{0.749020,0.772549,0.792157}%
\pgfsetfillcolor{currentfill}%
\pgfsetfillopacity{0.160000}%
\pgfsetlinewidth{0.000000pt}%
\definecolor{currentstroke}{rgb}{0.000000,0.000000,0.000000}%
\pgfsetstrokecolor{currentstroke}%
\pgfsetdash{}{0pt}%
\pgfpathmoveto{\pgfqpoint{1.434838in}{0.926341in}}%
\pgfpathlineto{\pgfqpoint{1.448253in}{0.938635in}}%
\pgfpathlineto{\pgfqpoint{1.461680in}{0.950939in}}%
\pgfpathlineto{\pgfqpoint{1.457844in}{0.960613in}}%
\pgfpathlineto{\pgfqpoint{1.452439in}{0.970108in}}%
\pgfpathlineto{\pgfqpoint{1.439422in}{0.957185in}}%
\pgfpathlineto{\pgfqpoint{1.426410in}{0.944268in}}%
\pgfpathlineto{\pgfqpoint{1.431354in}{0.935386in}}%
\pgfpathlineto{\pgfqpoint{1.434838in}{0.926341in}}%
\pgfpathclose%
\pgfusepath{fill}%
\end{pgfscope}%
\begin{pgfscope}%
\pgfpathrectangle{\pgfqpoint{0.407592in}{0.222904in}}{\pgfqpoint{1.411163in}{1.411163in}}%
\pgfusepath{clip}%
\pgfsetbuttcap%
\pgfsetroundjoin%
\definecolor{currentfill}{rgb}{0.749020,0.772549,0.792157}%
\pgfsetfillcolor{currentfill}%
\pgfsetfillopacity{0.160000}%
\pgfsetlinewidth{0.000000pt}%
\definecolor{currentstroke}{rgb}{0.000000,0.000000,0.000000}%
\pgfsetstrokecolor{currentstroke}%
\pgfsetdash{}{0pt}%
\pgfpathmoveto{\pgfqpoint{1.302409in}{0.792323in}}%
\pgfpathlineto{\pgfqpoint{1.315820in}{0.803813in}}%
\pgfpathlineto{\pgfqpoint{1.329249in}{0.815318in}}%
\pgfpathlineto{\pgfqpoint{1.329122in}{0.821878in}}%
\pgfpathlineto{\pgfqpoint{1.327952in}{0.828388in}}%
\pgfpathlineto{\pgfqpoint{1.314645in}{0.816193in}}%
\pgfpathlineto{\pgfqpoint{1.301350in}{0.804009in}}%
\pgfpathlineto{\pgfqpoint{1.302343in}{0.798187in}}%
\pgfpathlineto{\pgfqpoint{1.302409in}{0.792323in}}%
\pgfpathclose%
\pgfusepath{fill}%
\end{pgfscope}%
\begin{pgfscope}%
\pgfpathrectangle{\pgfqpoint{0.407592in}{0.222904in}}{\pgfqpoint{1.411163in}{1.411163in}}%
\pgfusepath{clip}%
\pgfsetbuttcap%
\pgfsetroundjoin%
\definecolor{currentfill}{rgb}{0.749020,0.772549,0.792157}%
\pgfsetfillcolor{currentfill}%
\pgfsetfillopacity{0.160000}%
\pgfsetlinewidth{0.000000pt}%
\definecolor{currentstroke}{rgb}{0.000000,0.000000,0.000000}%
\pgfsetstrokecolor{currentstroke}%
\pgfsetdash{}{0pt}%
\pgfpathmoveto{\pgfqpoint{1.052114in}{0.639058in}}%
\pgfpathlineto{\pgfqpoint{1.042498in}{0.646254in}}%
\pgfpathlineto{\pgfqpoint{1.032846in}{0.653477in}}%
\pgfpathlineto{\pgfqpoint{1.036632in}{0.651962in}}%
\pgfpathlineto{\pgfqpoint{1.040657in}{0.650563in}}%
\pgfpathlineto{\pgfqpoint{1.049000in}{0.643826in}}%
\pgfpathlineto{\pgfqpoint{1.057311in}{0.637114in}}%
\pgfpathlineto{\pgfqpoint{1.054634in}{0.638047in}}%
\pgfpathlineto{\pgfqpoint{1.052114in}{0.639058in}}%
\pgfpathclose%
\pgfusepath{fill}%
\end{pgfscope}%
\begin{pgfscope}%
\pgfpathrectangle{\pgfqpoint{0.407592in}{0.222904in}}{\pgfqpoint{1.411163in}{1.411163in}}%
\pgfusepath{clip}%
\pgfsetbuttcap%
\pgfsetroundjoin%
\definecolor{currentfill}{rgb}{0.749020,0.772549,0.792157}%
\pgfsetfillcolor{currentfill}%
\pgfsetfillopacity{0.160000}%
\pgfsetlinewidth{0.000000pt}%
\definecolor{currentstroke}{rgb}{0.000000,0.000000,0.000000}%
\pgfsetstrokecolor{currentstroke}%
\pgfsetdash{}{0pt}%
\pgfpathmoveto{\pgfqpoint{1.156589in}{0.660946in}}%
\pgfpathlineto{\pgfqpoint{1.167764in}{0.669436in}}%
\pgfpathlineto{\pgfqpoint{1.178973in}{0.677951in}}%
\pgfpathlineto{\pgfqpoint{1.182590in}{0.680609in}}%
\pgfpathlineto{\pgfqpoint{1.185795in}{0.683354in}}%
\pgfpathlineto{\pgfqpoint{1.173741in}{0.674169in}}%
\pgfpathlineto{\pgfqpoint{1.161718in}{0.665008in}}%
\pgfpathlineto{\pgfqpoint{1.159307in}{0.662944in}}%
\pgfpathlineto{\pgfqpoint{1.156589in}{0.660946in}}%
\pgfpathclose%
\pgfusepath{fill}%
\end{pgfscope}%
\begin{pgfscope}%
\pgfpathrectangle{\pgfqpoint{0.407592in}{0.222904in}}{\pgfqpoint{1.411163in}{1.411163in}}%
\pgfusepath{clip}%
\pgfsetbuttcap%
\pgfsetroundjoin%
\definecolor{currentfill}{rgb}{0.749020,0.772549,0.792157}%
\pgfsetfillcolor{currentfill}%
\pgfsetfillopacity{0.160000}%
\pgfsetlinewidth{0.000000pt}%
\definecolor{currentstroke}{rgb}{0.000000,0.000000,0.000000}%
\pgfsetstrokecolor{currentstroke}%
\pgfsetdash{}{0pt}%
\pgfpathmoveto{\pgfqpoint{0.758868in}{0.942501in}}%
\pgfpathlineto{\pgfqpoint{0.746035in}{0.955357in}}%
\pgfpathlineto{\pgfqpoint{0.733196in}{0.968219in}}%
\pgfpathlineto{\pgfqpoint{0.725977in}{0.958685in}}%
\pgfpathlineto{\pgfqpoint{0.720296in}{0.948980in}}%
\pgfpathlineto{\pgfqpoint{0.733675in}{0.936740in}}%
\pgfpathlineto{\pgfqpoint{0.747043in}{0.924510in}}%
\pgfpathlineto{\pgfqpoint{0.752240in}{0.933583in}}%
\pgfpathlineto{\pgfqpoint{0.758868in}{0.942501in}}%
\pgfpathclose%
\pgfusepath{fill}%
\end{pgfscope}%
\begin{pgfscope}%
\pgfpathrectangle{\pgfqpoint{0.407592in}{0.222904in}}{\pgfqpoint{1.411163in}{1.411163in}}%
\pgfusepath{clip}%
\pgfsetbuttcap%
\pgfsetroundjoin%
\definecolor{currentfill}{rgb}{0.749020,0.772549,0.792157}%
\pgfsetfillcolor{currentfill}%
\pgfsetfillopacity{0.160000}%
\pgfsetlinewidth{0.000000pt}%
\definecolor{currentstroke}{rgb}{0.000000,0.000000,0.000000}%
\pgfsetstrokecolor{currentstroke}%
\pgfsetdash{}{0pt}%
\pgfpathmoveto{\pgfqpoint{1.116702in}{0.636428in}}%
\pgfpathlineto{\pgfqpoint{1.123387in}{0.642967in}}%
\pgfpathlineto{\pgfqpoint{1.130099in}{0.649533in}}%
\pgfpathlineto{\pgfqpoint{1.134571in}{0.650834in}}%
\pgfpathlineto{\pgfqpoint{1.138843in}{0.652257in}}%
\pgfpathlineto{\pgfqpoint{1.130664in}{0.645237in}}%
\pgfpathlineto{\pgfqpoint{1.122516in}{0.638244in}}%
\pgfpathlineto{\pgfqpoint{1.119675in}{0.637295in}}%
\pgfpathlineto{\pgfqpoint{1.116702in}{0.636428in}}%
\pgfpathclose%
\pgfusepath{fill}%
\end{pgfscope}%
\begin{pgfscope}%
\pgfpathrectangle{\pgfqpoint{0.407592in}{0.222904in}}{\pgfqpoint{1.411163in}{1.411163in}}%
\pgfusepath{clip}%
\pgfsetbuttcap%
\pgfsetroundjoin%
\definecolor{currentfill}{rgb}{0.749020,0.772549,0.792157}%
\pgfsetfillcolor{currentfill}%
\pgfsetfillopacity{0.160000}%
\pgfsetlinewidth{0.000000pt}%
\definecolor{currentstroke}{rgb}{0.000000,0.000000,0.000000}%
\pgfsetstrokecolor{currentstroke}%
\pgfsetdash{}{0pt}%
\pgfpathmoveto{\pgfqpoint{0.880026in}{0.802846in}}%
\pgfpathlineto{\pgfqpoint{0.866784in}{0.814961in}}%
\pgfpathlineto{\pgfqpoint{0.853529in}{0.827087in}}%
\pgfpathlineto{\pgfqpoint{0.851164in}{0.820565in}}%
\pgfpathlineto{\pgfqpoint{0.849835in}{0.813999in}}%
\pgfpathlineto{\pgfqpoint{0.863346in}{0.802563in}}%
\pgfpathlineto{\pgfqpoint{0.876838in}{0.791143in}}%
\pgfpathlineto{\pgfqpoint{0.877972in}{0.797013in}}%
\pgfpathlineto{\pgfqpoint{0.880026in}{0.802846in}}%
\pgfpathclose%
\pgfusepath{fill}%
\end{pgfscope}%
\begin{pgfscope}%
\pgfpathrectangle{\pgfqpoint{0.407592in}{0.222904in}}{\pgfqpoint{1.411163in}{1.411163in}}%
\pgfusepath{clip}%
\pgfsetbuttcap%
\pgfsetroundjoin%
\definecolor{currentfill}{rgb}{0.749020,0.772549,0.792157}%
\pgfsetfillcolor{currentfill}%
\pgfsetfillopacity{0.160000}%
\pgfsetlinewidth{0.000000pt}%
\definecolor{currentstroke}{rgb}{0.000000,0.000000,0.000000}%
\pgfsetstrokecolor{currentstroke}%
\pgfsetdash{}{0pt}%
\pgfpathmoveto{\pgfqpoint{1.015983in}{0.664589in}}%
\pgfpathlineto{\pgfqpoint{1.003495in}{0.673681in}}%
\pgfpathlineto{\pgfqpoint{0.990974in}{0.682797in}}%
\pgfpathlineto{\pgfqpoint{0.993695in}{0.680068in}}%
\pgfpathlineto{\pgfqpoint{0.996845in}{0.677429in}}%
\pgfpathlineto{\pgfqpoint{1.008641in}{0.668979in}}%
\pgfpathlineto{\pgfqpoint{1.020402in}{0.660554in}}%
\pgfpathlineto{\pgfqpoint{1.018033in}{0.662537in}}%
\pgfpathlineto{\pgfqpoint{1.015983in}{0.664589in}}%
\pgfpathclose%
\pgfusepath{fill}%
\end{pgfscope}%
\begin{pgfscope}%
\pgfpathrectangle{\pgfqpoint{0.407592in}{0.222904in}}{\pgfqpoint{1.411163in}{1.411163in}}%
\pgfusepath{clip}%
\pgfsetbuttcap%
\pgfsetroundjoin%
\definecolor{currentfill}{rgb}{0.749020,0.772549,0.792157}%
\pgfsetfillcolor{currentfill}%
\pgfsetfillopacity{0.160000}%
\pgfsetlinewidth{0.000000pt}%
\definecolor{currentstroke}{rgb}{0.000000,0.000000,0.000000}%
\pgfsetstrokecolor{currentstroke}%
\pgfsetdash{}{0pt}%
\pgfpathmoveto{\pgfqpoint{1.461680in}{0.950939in}}%
\pgfpathlineto{\pgfqpoint{1.475119in}{0.963255in}}%
\pgfpathlineto{\pgfqpoint{1.488569in}{0.975581in}}%
\pgfpathlineto{\pgfqpoint{1.484371in}{0.985874in}}%
\pgfpathlineto{\pgfqpoint{1.478494in}{0.995973in}}%
\pgfpathlineto{\pgfqpoint{1.465463in}{0.983037in}}%
\pgfpathlineto{\pgfqpoint{1.452439in}{0.970108in}}%
\pgfpathlineto{\pgfqpoint{1.457844in}{0.960613in}}%
\pgfpathlineto{\pgfqpoint{1.461680in}{0.950939in}}%
\pgfpathclose%
\pgfusepath{fill}%
\end{pgfscope}%
\begin{pgfscope}%
\pgfpathrectangle{\pgfqpoint{0.407592in}{0.222904in}}{\pgfqpoint{1.411163in}{1.411163in}}%
\pgfusepath{clip}%
\pgfsetbuttcap%
\pgfsetroundjoin%
\definecolor{currentfill}{rgb}{0.749020,0.772549,0.792157}%
\pgfsetfillcolor{currentfill}%
\pgfsetfillopacity{0.160000}%
\pgfsetlinewidth{0.000000pt}%
\definecolor{currentstroke}{rgb}{0.000000,0.000000,0.000000}%
\pgfsetstrokecolor{currentstroke}%
\pgfsetdash{}{0pt}%
\pgfpathmoveto{\pgfqpoint{1.246815in}{0.737613in}}%
\pgfpathlineto{\pgfqpoint{1.260013in}{0.748324in}}%
\pgfpathlineto{\pgfqpoint{1.273234in}{0.759054in}}%
\pgfpathlineto{\pgfqpoint{1.274845in}{0.764216in}}%
\pgfpathlineto{\pgfqpoint{1.275641in}{0.769388in}}%
\pgfpathlineto{\pgfqpoint{1.262283in}{0.757944in}}%
\pgfpathlineto{\pgfqpoint{1.248943in}{0.746514in}}%
\pgfpathlineto{\pgfqpoint{1.248227in}{0.742058in}}%
\pgfpathlineto{\pgfqpoint{1.246815in}{0.737613in}}%
\pgfpathclose%
\pgfusepath{fill}%
\end{pgfscope}%
\begin{pgfscope}%
\pgfpathrectangle{\pgfqpoint{0.407592in}{0.222904in}}{\pgfqpoint{1.411163in}{1.411163in}}%
\pgfusepath{clip}%
\pgfsetbuttcap%
\pgfsetroundjoin%
\definecolor{currentfill}{rgb}{0.749020,0.772549,0.792157}%
\pgfsetfillcolor{currentfill}%
\pgfsetfillopacity{0.160000}%
\pgfsetlinewidth{0.000000pt}%
\definecolor{currentstroke}{rgb}{0.000000,0.000000,0.000000}%
\pgfsetstrokecolor{currentstroke}%
\pgfsetdash{}{0pt}%
\pgfpathmoveto{\pgfqpoint{1.057311in}{0.637114in}}%
\pgfpathlineto{\pgfqpoint{1.049000in}{0.643826in}}%
\pgfpathlineto{\pgfqpoint{1.040657in}{0.650563in}}%
\pgfpathlineto{\pgfqpoint{1.044902in}{0.649287in}}%
\pgfpathlineto{\pgfqpoint{1.049350in}{0.648138in}}%
\pgfpathlineto{\pgfqpoint{1.056235in}{0.641805in}}%
\pgfpathlineto{\pgfqpoint{1.063090in}{0.635499in}}%
\pgfpathlineto{\pgfqpoint{1.060134in}{0.636264in}}%
\pgfpathlineto{\pgfqpoint{1.057311in}{0.637114in}}%
\pgfpathclose%
\pgfusepath{fill}%
\end{pgfscope}%
\begin{pgfscope}%
\pgfpathrectangle{\pgfqpoint{0.407592in}{0.222904in}}{\pgfqpoint{1.411163in}{1.411163in}}%
\pgfusepath{clip}%
\pgfsetbuttcap%
\pgfsetroundjoin%
\definecolor{currentfill}{rgb}{0.749020,0.772549,0.792157}%
\pgfsetfillcolor{currentfill}%
\pgfsetfillopacity{0.160000}%
\pgfsetlinewidth{0.000000pt}%
\definecolor{currentstroke}{rgb}{0.000000,0.000000,0.000000}%
\pgfsetstrokecolor{currentstroke}%
\pgfsetdash{}{0pt}%
\pgfpathmoveto{\pgfqpoint{0.930624in}{0.745620in}}%
\pgfpathlineto{\pgfqpoint{0.917205in}{0.756978in}}%
\pgfpathlineto{\pgfqpoint{0.903768in}{0.768351in}}%
\pgfpathlineto{\pgfqpoint{0.903630in}{0.763179in}}%
\pgfpathlineto{\pgfqpoint{0.904309in}{0.758021in}}%
\pgfpathlineto{\pgfqpoint{0.917745in}{0.747363in}}%
\pgfpathlineto{\pgfqpoint{0.931156in}{0.736724in}}%
\pgfpathlineto{\pgfqpoint{0.930541in}{0.741165in}}%
\pgfpathlineto{\pgfqpoint{0.930624in}{0.745620in}}%
\pgfpathclose%
\pgfusepath{fill}%
\end{pgfscope}%
\begin{pgfscope}%
\pgfpathrectangle{\pgfqpoint{0.407592in}{0.222904in}}{\pgfqpoint{1.411163in}{1.411163in}}%
\pgfusepath{clip}%
\pgfsetbuttcap%
\pgfsetroundjoin%
\definecolor{currentfill}{rgb}{0.749020,0.772549,0.792157}%
\pgfsetfillcolor{currentfill}%
\pgfsetfillopacity{0.160000}%
\pgfsetlinewidth{0.000000pt}%
\definecolor{currentstroke}{rgb}{0.000000,0.000000,0.000000}%
\pgfsetstrokecolor{currentstroke}%
\pgfsetdash{}{0pt}%
\pgfpathmoveto{\pgfqpoint{1.110416in}{0.634951in}}%
\pgfpathlineto{\pgfqpoint{1.115516in}{0.641121in}}%
\pgfpathlineto{\pgfqpoint{1.120638in}{0.647317in}}%
\pgfpathlineto{\pgfqpoint{1.125448in}{0.648358in}}%
\pgfpathlineto{\pgfqpoint{1.130099in}{0.649533in}}%
\pgfpathlineto{\pgfqpoint{1.123387in}{0.642967in}}%
\pgfpathlineto{\pgfqpoint{1.116702in}{0.636428in}}%
\pgfpathlineto{\pgfqpoint{1.113611in}{0.635645in}}%
\pgfpathlineto{\pgfqpoint{1.110416in}{0.634951in}}%
\pgfpathclose%
\pgfusepath{fill}%
\end{pgfscope}%
\begin{pgfscope}%
\pgfpathrectangle{\pgfqpoint{0.407592in}{0.222904in}}{\pgfqpoint{1.411163in}{1.411163in}}%
\pgfusepath{clip}%
\pgfsetbuttcap%
\pgfsetroundjoin%
\definecolor{currentfill}{rgb}{0.749020,0.772549,0.792157}%
\pgfsetfillcolor{currentfill}%
\pgfsetfillopacity{0.160000}%
\pgfsetlinewidth{0.000000pt}%
\definecolor{currentstroke}{rgb}{0.000000,0.000000,0.000000}%
\pgfsetstrokecolor{currentstroke}%
\pgfsetdash{}{0pt}%
\pgfpathmoveto{\pgfqpoint{0.733196in}{0.968219in}}%
\pgfpathlineto{\pgfqpoint{0.720350in}{0.981088in}}%
\pgfpathlineto{\pgfqpoint{0.707498in}{0.993964in}}%
\pgfpathlineto{\pgfqpoint{0.699676in}{0.983823in}}%
\pgfpathlineto{\pgfqpoint{0.693499in}{0.973497in}}%
\pgfpathlineto{\pgfqpoint{0.706903in}{0.961233in}}%
\pgfpathlineto{\pgfqpoint{0.720296in}{0.948980in}}%
\pgfpathlineto{\pgfqpoint{0.725977in}{0.958685in}}%
\pgfpathlineto{\pgfqpoint{0.733196in}{0.968219in}}%
\pgfpathclose%
\pgfusepath{fill}%
\end{pgfscope}%
\begin{pgfscope}%
\pgfpathrectangle{\pgfqpoint{0.407592in}{0.222904in}}{\pgfqpoint{1.411163in}{1.411163in}}%
\pgfusepath{clip}%
\pgfsetbuttcap%
\pgfsetroundjoin%
\definecolor{currentfill}{rgb}{0.749020,0.772549,0.792157}%
\pgfsetfillcolor{currentfill}%
\pgfsetfillopacity{0.160000}%
\pgfsetlinewidth{0.000000pt}%
\definecolor{currentstroke}{rgb}{0.000000,0.000000,0.000000}%
\pgfsetstrokecolor{currentstroke}%
\pgfsetdash{}{0pt}%
\pgfpathmoveto{\pgfqpoint{1.329249in}{0.815318in}}%
\pgfpathlineto{\pgfqpoint{1.342696in}{0.826839in}}%
\pgfpathlineto{\pgfqpoint{1.356160in}{0.838375in}}%
\pgfpathlineto{\pgfqpoint{1.355961in}{0.845622in}}%
\pgfpathlineto{\pgfqpoint{1.354602in}{0.852811in}}%
\pgfpathlineto{\pgfqpoint{1.341271in}{0.840594in}}%
\pgfpathlineto{\pgfqpoint{1.327952in}{0.828388in}}%
\pgfpathlineto{\pgfqpoint{1.329122in}{0.821878in}}%
\pgfpathlineto{\pgfqpoint{1.329249in}{0.815318in}}%
\pgfpathclose%
\pgfusepath{fill}%
\end{pgfscope}%
\begin{pgfscope}%
\pgfpathrectangle{\pgfqpoint{0.407592in}{0.222904in}}{\pgfqpoint{1.411163in}{1.411163in}}%
\pgfusepath{clip}%
\pgfsetbuttcap%
\pgfsetroundjoin%
\definecolor{currentfill}{rgb}{0.749020,0.772549,0.792157}%
\pgfsetfillcolor{currentfill}%
\pgfsetfillopacity{0.160000}%
\pgfsetlinewidth{0.000000pt}%
\definecolor{currentstroke}{rgb}{0.000000,0.000000,0.000000}%
\pgfsetstrokecolor{currentstroke}%
\pgfsetdash{}{0pt}%
\pgfpathmoveto{\pgfqpoint{1.488569in}{0.975581in}}%
\pgfpathlineto{\pgfqpoint{1.502032in}{0.987919in}}%
\pgfpathlineto{\pgfqpoint{1.515507in}{1.000268in}}%
\pgfpathlineto{\pgfqpoint{1.510934in}{1.011170in}}%
\pgfpathlineto{\pgfqpoint{1.504572in}{1.021862in}}%
\pgfpathlineto{\pgfqpoint{1.491530in}{1.008914in}}%
\pgfpathlineto{\pgfqpoint{1.478494in}{0.995973in}}%
\pgfpathlineto{\pgfqpoint{1.484371in}{0.985874in}}%
\pgfpathlineto{\pgfqpoint{1.488569in}{0.975581in}}%
\pgfpathclose%
\pgfusepath{fill}%
\end{pgfscope}%
\begin{pgfscope}%
\pgfpathrectangle{\pgfqpoint{0.407592in}{0.222904in}}{\pgfqpoint{1.411163in}{1.411163in}}%
\pgfusepath{clip}%
\pgfsetbuttcap%
\pgfsetroundjoin%
\definecolor{currentfill}{rgb}{0.749020,0.772549,0.792157}%
\pgfsetfillcolor{currentfill}%
\pgfsetfillopacity{0.160000}%
\pgfsetlinewidth{0.000000pt}%
\definecolor{currentstroke}{rgb}{0.000000,0.000000,0.000000}%
\pgfsetstrokecolor{currentstroke}%
\pgfsetdash{}{0pt}%
\pgfpathmoveto{\pgfqpoint{1.063090in}{0.635499in}}%
\pgfpathlineto{\pgfqpoint{1.056235in}{0.641805in}}%
\pgfpathlineto{\pgfqpoint{1.049350in}{0.648138in}}%
\pgfpathlineto{\pgfqpoint{1.053980in}{0.647124in}}%
\pgfpathlineto{\pgfqpoint{1.058772in}{0.646248in}}%
\pgfpathlineto{\pgfqpoint{1.064072in}{0.640230in}}%
\pgfpathlineto{\pgfqpoint{1.069349in}{0.634239in}}%
\pgfpathlineto{\pgfqpoint{1.066167in}{0.634823in}}%
\pgfpathlineto{\pgfqpoint{1.063090in}{0.635499in}}%
\pgfpathclose%
\pgfusepath{fill}%
\end{pgfscope}%
\begin{pgfscope}%
\pgfpathrectangle{\pgfqpoint{0.407592in}{0.222904in}}{\pgfqpoint{1.411163in}{1.411163in}}%
\pgfusepath{clip}%
\pgfsetbuttcap%
\pgfsetroundjoin%
\definecolor{currentfill}{rgb}{0.749020,0.772549,0.792157}%
\pgfsetfillcolor{currentfill}%
\pgfsetfillopacity{0.160000}%
\pgfsetlinewidth{0.000000pt}%
\definecolor{currentstroke}{rgb}{0.000000,0.000000,0.000000}%
\pgfsetstrokecolor{currentstroke}%
\pgfsetdash{}{0pt}%
\pgfpathmoveto{\pgfqpoint{1.153575in}{0.659024in}}%
\pgfpathlineto{\pgfqpoint{1.164250in}{0.667195in}}%
\pgfpathlineto{\pgfqpoint{1.174958in}{0.675392in}}%
\pgfpathlineto{\pgfqpoint{1.178973in}{0.677951in}}%
\pgfpathlineto{\pgfqpoint{1.167764in}{0.669436in}}%
\pgfpathlineto{\pgfqpoint{1.156589in}{0.660946in}}%
\pgfusepath{fill}%
\end{pgfscope}%
\begin{pgfscope}%
\pgfpathrectangle{\pgfqpoint{0.407592in}{0.222904in}}{\pgfqpoint{1.411163in}{1.411163in}}%
\pgfusepath{clip}%
\pgfsetbuttcap%
\pgfsetroundjoin%
\definecolor{currentfill}{rgb}{0.749020,0.772549,0.792157}%
\pgfsetfillcolor{currentfill}%
\pgfsetfillopacity{0.160000}%
\pgfsetlinewidth{0.000000pt}%
\definecolor{currentstroke}{rgb}{0.000000,0.000000,0.000000}%
\pgfsetstrokecolor{currentstroke}%
\pgfsetdash{}{0pt}%
\pgfpathmoveto{\pgfqpoint{1.185795in}{0.683354in}}%
\pgfpathlineto{\pgfqpoint{1.197880in}{0.692563in}}%
\pgfpathlineto{\pgfqpoint{1.209996in}{0.701796in}}%
\pgfpathlineto{\pgfqpoint{1.213454in}{0.705310in}}%
\pgfpathlineto{\pgfqpoint{1.216362in}{0.708901in}}%
\pgfpathlineto{\pgfqpoint{1.203625in}{0.698968in}}%
\pgfpathlineto{\pgfqpoint{1.190916in}{0.689056in}}%
\pgfpathlineto{\pgfqpoint{1.188574in}{0.686174in}}%
\pgfpathlineto{\pgfqpoint{1.185795in}{0.683354in}}%
\pgfpathclose%
\pgfusepath{fill}%
\end{pgfscope}%
\begin{pgfscope}%
\pgfpathrectangle{\pgfqpoint{0.407592in}{0.222904in}}{\pgfqpoint{1.411163in}{1.411163in}}%
\pgfusepath{clip}%
\pgfsetbuttcap%
\pgfsetroundjoin%
\definecolor{currentfill}{rgb}{0.749020,0.772549,0.792157}%
\pgfsetfillcolor{currentfill}%
\pgfsetfillopacity{0.160000}%
\pgfsetlinewidth{0.000000pt}%
\definecolor{currentstroke}{rgb}{0.000000,0.000000,0.000000}%
\pgfsetstrokecolor{currentstroke}%
\pgfsetdash{}{0pt}%
\pgfpathmoveto{\pgfqpoint{1.216362in}{0.708901in}}%
\pgfpathlineto{\pgfqpoint{1.229127in}{0.718857in}}%
\pgfpathlineto{\pgfqpoint{1.241920in}{0.728833in}}%
\pgfpathlineto{\pgfqpoint{1.244710in}{0.733198in}}%
\pgfpathlineto{\pgfqpoint{1.246815in}{0.737613in}}%
\pgfpathlineto{\pgfqpoint{1.233640in}{0.726920in}}%
\pgfpathlineto{\pgfqpoint{1.220488in}{0.716247in}}%
\pgfpathlineto{\pgfqpoint{1.218710in}{0.712552in}}%
\pgfpathlineto{\pgfqpoint{1.216362in}{0.708901in}}%
\pgfpathclose%
\pgfusepath{fill}%
\end{pgfscope}%
\begin{pgfscope}%
\pgfpathrectangle{\pgfqpoint{0.407592in}{0.222904in}}{\pgfqpoint{1.411163in}{1.411163in}}%
\pgfusepath{clip}%
\pgfsetbuttcap%
\pgfsetroundjoin%
\definecolor{currentfill}{rgb}{0.749020,0.772549,0.792157}%
\pgfsetfillcolor{currentfill}%
\pgfsetfillopacity{0.160000}%
\pgfsetlinewidth{0.000000pt}%
\definecolor{currentstroke}{rgb}{0.000000,0.000000,0.000000}%
\pgfsetstrokecolor{currentstroke}%
\pgfsetdash{}{0pt}%
\pgfpathmoveto{\pgfqpoint{1.103768in}{0.633841in}}%
\pgfpathlineto{\pgfqpoint{1.107190in}{0.639732in}}%
\pgfpathlineto{\pgfqpoint{1.110627in}{0.645650in}}%
\pgfpathlineto{\pgfqpoint{1.115690in}{0.646412in}}%
\pgfpathlineto{\pgfqpoint{1.120638in}{0.647317in}}%
\pgfpathlineto{\pgfqpoint{1.115516in}{0.641121in}}%
\pgfpathlineto{\pgfqpoint{1.110416in}{0.634951in}}%
\pgfpathlineto{\pgfqpoint{1.107130in}{0.634349in}}%
\pgfpathlineto{\pgfqpoint{1.103768in}{0.633841in}}%
\pgfpathclose%
\pgfusepath{fill}%
\end{pgfscope}%
\begin{pgfscope}%
\pgfpathrectangle{\pgfqpoint{0.407592in}{0.222904in}}{\pgfqpoint{1.411163in}{1.411163in}}%
\pgfusepath{clip}%
\pgfsetbuttcap%
\pgfsetroundjoin%
\definecolor{currentfill}{rgb}{0.749020,0.772549,0.792157}%
\pgfsetfillcolor{currentfill}%
\pgfsetfillopacity{0.160000}%
\pgfsetlinewidth{0.000000pt}%
\definecolor{currentstroke}{rgb}{0.000000,0.000000,0.000000}%
\pgfsetstrokecolor{currentstroke}%
\pgfsetdash{}{0pt}%
\pgfpathmoveto{\pgfqpoint{0.986858in}{0.688474in}}%
\pgfpathlineto{\pgfqpoint{0.973807in}{0.698314in}}%
\pgfpathlineto{\pgfqpoint{0.960727in}{0.708176in}}%
\pgfpathlineto{\pgfqpoint{0.962998in}{0.704599in}}%
\pgfpathlineto{\pgfqpoint{0.965832in}{0.701101in}}%
\pgfpathlineto{\pgfqpoint{0.978420in}{0.691937in}}%
\pgfpathlineto{\pgfqpoint{0.990974in}{0.682797in}}%
\pgfpathlineto{\pgfqpoint{0.988691in}{0.685603in}}%
\pgfpathlineto{\pgfqpoint{0.986858in}{0.688474in}}%
\pgfpathclose%
\pgfusepath{fill}%
\end{pgfscope}%
\begin{pgfscope}%
\pgfpathrectangle{\pgfqpoint{0.407592in}{0.222904in}}{\pgfqpoint{1.411163in}{1.411163in}}%
\pgfusepath{clip}%
\pgfsetbuttcap%
\pgfsetroundjoin%
\definecolor{currentfill}{rgb}{0.749020,0.772549,0.792157}%
\pgfsetfillcolor{currentfill}%
\pgfsetfillopacity{0.160000}%
\pgfsetlinewidth{0.000000pt}%
\definecolor{currentstroke}{rgb}{0.000000,0.000000,0.000000}%
\pgfsetstrokecolor{currentstroke}%
\pgfsetdash{}{0pt}%
\pgfpathmoveto{\pgfqpoint{0.853529in}{0.827087in}}%
\pgfpathlineto{\pgfqpoint{0.840262in}{0.839225in}}%
\pgfpathlineto{\pgfqpoint{0.826982in}{0.851374in}}%
\pgfpathlineto{\pgfqpoint{0.824295in}{0.844172in}}%
\pgfpathlineto{\pgfqpoint{0.822757in}{0.836917in}}%
\pgfpathlineto{\pgfqpoint{0.836305in}{0.825450in}}%
\pgfpathlineto{\pgfqpoint{0.849835in}{0.813999in}}%
\pgfpathlineto{\pgfqpoint{0.851164in}{0.820565in}}%
\pgfpathlineto{\pgfqpoint{0.853529in}{0.827087in}}%
\pgfpathclose%
\pgfusepath{fill}%
\end{pgfscope}%
\begin{pgfscope}%
\pgfpathrectangle{\pgfqpoint{0.407592in}{0.222904in}}{\pgfqpoint{1.411163in}{1.411163in}}%
\pgfusepath{clip}%
\pgfsetbuttcap%
\pgfsetroundjoin%
\definecolor{currentfill}{rgb}{0.749020,0.772549,0.792157}%
\pgfsetfillcolor{currentfill}%
\pgfsetfillopacity{0.160000}%
\pgfsetlinewidth{0.000000pt}%
\definecolor{currentstroke}{rgb}{0.000000,0.000000,0.000000}%
\pgfsetstrokecolor{currentstroke}%
\pgfsetdash{}{0pt}%
\pgfpathmoveto{\pgfqpoint{1.069349in}{0.634239in}}%
\pgfpathlineto{\pgfqpoint{1.064072in}{0.640230in}}%
\pgfpathlineto{\pgfqpoint{1.058772in}{0.646248in}}%
\pgfpathlineto{\pgfqpoint{1.063704in}{0.645514in}}%
\pgfpathlineto{\pgfqpoint{1.068755in}{0.644926in}}%
\pgfpathlineto{\pgfqpoint{1.072374in}{0.639130in}}%
\pgfpathlineto{\pgfqpoint{1.075977in}{0.633359in}}%
\pgfpathlineto{\pgfqpoint{1.072624in}{0.633751in}}%
\pgfpathlineto{\pgfqpoint{1.069349in}{0.634239in}}%
\pgfpathclose%
\pgfusepath{fill}%
\end{pgfscope}%
\begin{pgfscope}%
\pgfpathrectangle{\pgfqpoint{0.407592in}{0.222904in}}{\pgfqpoint{1.411163in}{1.411163in}}%
\pgfusepath{clip}%
\pgfsetbuttcap%
\pgfsetroundjoin%
\definecolor{currentfill}{rgb}{0.749020,0.772549,0.792157}%
\pgfsetfillcolor{currentfill}%
\pgfsetfillopacity{0.160000}%
\pgfsetlinewidth{0.000000pt}%
\definecolor{currentstroke}{rgb}{0.000000,0.000000,0.000000}%
\pgfsetstrokecolor{currentstroke}%
\pgfsetdash{}{0pt}%
\pgfpathmoveto{\pgfqpoint{0.957907in}{0.715503in}}%
\pgfpathlineto{\pgfqpoint{0.944544in}{0.726104in}}%
\pgfpathlineto{\pgfqpoint{0.931156in}{0.736724in}}%
\pgfpathlineto{\pgfqpoint{0.932470in}{0.732318in}}%
\pgfpathlineto{\pgfqpoint{0.934480in}{0.727966in}}%
\pgfpathlineto{\pgfqpoint{0.947618in}{0.718060in}}%
\pgfpathlineto{\pgfqpoint{0.960727in}{0.708176in}}%
\pgfpathlineto{\pgfqpoint{0.959028in}{0.711816in}}%
\pgfpathlineto{\pgfqpoint{0.957907in}{0.715503in}}%
\pgfpathclose%
\pgfusepath{fill}%
\end{pgfscope}%
\begin{pgfscope}%
\pgfpathrectangle{\pgfqpoint{0.407592in}{0.222904in}}{\pgfqpoint{1.411163in}{1.411163in}}%
\pgfusepath{clip}%
\pgfsetbuttcap%
\pgfsetroundjoin%
\definecolor{currentfill}{rgb}{0.749020,0.772549,0.792157}%
\pgfsetfillcolor{currentfill}%
\pgfsetfillopacity{0.160000}%
\pgfsetlinewidth{0.000000pt}%
\definecolor{currentstroke}{rgb}{0.000000,0.000000,0.000000}%
\pgfsetstrokecolor{currentstroke}%
\pgfsetdash{}{0pt}%
\pgfpathmoveto{\pgfqpoint{1.020402in}{0.660554in}}%
\pgfpathlineto{\pgfqpoint{1.008641in}{0.668979in}}%
\pgfpathlineto{\pgfqpoint{0.996845in}{0.677429in}}%
\pgfpathlineto{\pgfqpoint{1.000410in}{0.674892in}}%
\pgfpathlineto{\pgfqpoint{1.004374in}{0.672468in}}%
\pgfpathlineto{\pgfqpoint{1.015233in}{0.664635in}}%
\pgfpathlineto{\pgfqpoint{1.026055in}{0.656828in}}%
\pgfpathlineto{\pgfqpoint{1.023079in}{0.658648in}}%
\pgfpathlineto{\pgfqpoint{1.020402in}{0.660554in}}%
\pgfpathclose%
\pgfusepath{fill}%
\end{pgfscope}%
\begin{pgfscope}%
\pgfpathrectangle{\pgfqpoint{0.407592in}{0.222904in}}{\pgfqpoint{1.411163in}{1.411163in}}%
\pgfusepath{clip}%
\pgfsetbuttcap%
\pgfsetroundjoin%
\definecolor{currentfill}{rgb}{0.749020,0.772549,0.792157}%
\pgfsetfillcolor{currentfill}%
\pgfsetfillopacity{0.160000}%
\pgfsetlinewidth{0.000000pt}%
\definecolor{currentstroke}{rgb}{0.000000,0.000000,0.000000}%
\pgfsetstrokecolor{currentstroke}%
\pgfsetdash{}{0pt}%
\pgfpathmoveto{\pgfqpoint{1.096876in}{0.633117in}}%
\pgfpathlineto{\pgfqpoint{1.098557in}{0.638827in}}%
\pgfpathlineto{\pgfqpoint{1.100245in}{0.644563in}}%
\pgfpathlineto{\pgfqpoint{1.105471in}{0.645032in}}%
\pgfpathlineto{\pgfqpoint{1.110627in}{0.645650in}}%
\pgfpathlineto{\pgfqpoint{1.107190in}{0.639732in}}%
\pgfpathlineto{\pgfqpoint{1.103768in}{0.633841in}}%
\pgfpathlineto{\pgfqpoint{1.100345in}{0.633430in}}%
\pgfpathlineto{\pgfqpoint{1.096876in}{0.633117in}}%
\pgfpathclose%
\pgfusepath{fill}%
\end{pgfscope}%
\begin{pgfscope}%
\pgfpathrectangle{\pgfqpoint{0.407592in}{0.222904in}}{\pgfqpoint{1.411163in}{1.411163in}}%
\pgfusepath{clip}%
\pgfsetbuttcap%
\pgfsetroundjoin%
\definecolor{currentfill}{rgb}{0.749020,0.772549,0.792157}%
\pgfsetfillcolor{currentfill}%
\pgfsetfillopacity{0.160000}%
\pgfsetlinewidth{0.000000pt}%
\definecolor{currentstroke}{rgb}{0.000000,0.000000,0.000000}%
\pgfsetstrokecolor{currentstroke}%
\pgfsetdash{}{0pt}%
\pgfpathmoveto{\pgfqpoint{0.707498in}{0.993964in}}%
\pgfpathlineto{\pgfqpoint{0.694639in}{1.006847in}}%
\pgfpathlineto{\pgfqpoint{0.681774in}{1.019736in}}%
\pgfpathlineto{\pgfqpoint{0.673337in}{1.008998in}}%
\pgfpathlineto{\pgfqpoint{0.666651in}{0.998060in}}%
\pgfpathlineto{\pgfqpoint{0.680081in}{0.985772in}}%
\pgfpathlineto{\pgfqpoint{0.693499in}{0.973497in}}%
\pgfpathlineto{\pgfqpoint{0.699676in}{0.983823in}}%
\pgfpathlineto{\pgfqpoint{0.707498in}{0.993964in}}%
\pgfpathclose%
\pgfusepath{fill}%
\end{pgfscope}%
\begin{pgfscope}%
\pgfpathrectangle{\pgfqpoint{0.407592in}{0.222904in}}{\pgfqpoint{1.411163in}{1.411163in}}%
\pgfusepath{clip}%
\pgfsetbuttcap%
\pgfsetroundjoin%
\definecolor{currentfill}{rgb}{0.749020,0.772549,0.792157}%
\pgfsetfillcolor{currentfill}%
\pgfsetfillopacity{0.160000}%
\pgfsetlinewidth{0.000000pt}%
\definecolor{currentstroke}{rgb}{0.000000,0.000000,0.000000}%
\pgfsetstrokecolor{currentstroke}%
\pgfsetdash{}{0pt}%
\pgfpathmoveto{\pgfqpoint{1.515507in}{1.000268in}}%
\pgfpathlineto{\pgfqpoint{1.528995in}{1.012628in}}%
\pgfpathlineto{\pgfqpoint{1.542494in}{1.024999in}}%
\pgfpathlineto{\pgfqpoint{1.537533in}{1.036500in}}%
\pgfpathlineto{\pgfqpoint{1.530675in}{1.047775in}}%
\pgfpathlineto{\pgfqpoint{1.517621in}{1.034815in}}%
\pgfpathlineto{\pgfqpoint{1.504572in}{1.021862in}}%
\pgfpathlineto{\pgfqpoint{1.510934in}{1.011170in}}%
\pgfpathlineto{\pgfqpoint{1.515507in}{1.000268in}}%
\pgfpathclose%
\pgfusepath{fill}%
\end{pgfscope}%
\begin{pgfscope}%
\pgfpathrectangle{\pgfqpoint{0.407592in}{0.222904in}}{\pgfqpoint{1.411163in}{1.411163in}}%
\pgfusepath{clip}%
\pgfsetbuttcap%
\pgfsetroundjoin%
\definecolor{currentfill}{rgb}{0.749020,0.772549,0.792157}%
\pgfsetfillcolor{currentfill}%
\pgfsetfillopacity{0.160000}%
\pgfsetlinewidth{0.000000pt}%
\definecolor{currentstroke}{rgb}{0.000000,0.000000,0.000000}%
\pgfsetstrokecolor{currentstroke}%
\pgfsetdash{}{0pt}%
\pgfpathmoveto{\pgfqpoint{1.075977in}{0.633359in}}%
\pgfpathlineto{\pgfqpoint{1.072374in}{0.639130in}}%
\pgfpathlineto{\pgfqpoint{1.068755in}{0.644926in}}%
\pgfpathlineto{\pgfqpoint{1.073901in}{0.644486in}}%
\pgfpathlineto{\pgfqpoint{1.079119in}{0.644197in}}%
\pgfpathlineto{\pgfqpoint{1.080992in}{0.638523in}}%
\pgfpathlineto{\pgfqpoint{1.082856in}{0.632874in}}%
\pgfpathlineto{\pgfqpoint{1.079393in}{0.633066in}}%
\pgfpathlineto{\pgfqpoint{1.075977in}{0.633359in}}%
\pgfpathclose%
\pgfusepath{fill}%
\end{pgfscope}%
\begin{pgfscope}%
\pgfpathrectangle{\pgfqpoint{0.407592in}{0.222904in}}{\pgfqpoint{1.411163in}{1.411163in}}%
\pgfusepath{clip}%
\pgfsetbuttcap%
\pgfsetroundjoin%
\definecolor{currentfill}{rgb}{0.749020,0.772549,0.792157}%
\pgfsetfillcolor{currentfill}%
\pgfsetfillopacity{0.160000}%
\pgfsetlinewidth{0.000000pt}%
\definecolor{currentstroke}{rgb}{0.000000,0.000000,0.000000}%
\pgfsetstrokecolor{currentstroke}%
\pgfsetdash{}{0pt}%
\pgfpathmoveto{\pgfqpoint{1.089864in}{0.632793in}}%
\pgfpathlineto{\pgfqpoint{1.089772in}{0.638421in}}%
\pgfpathlineto{\pgfqpoint{1.089679in}{0.644076in}}%
\pgfpathlineto{\pgfqpoint{1.094973in}{0.644243in}}%
\pgfpathlineto{\pgfqpoint{1.100245in}{0.644563in}}%
\pgfpathlineto{\pgfqpoint{1.098557in}{0.638827in}}%
\pgfpathlineto{\pgfqpoint{1.096876in}{0.633117in}}%
\pgfpathlineto{\pgfqpoint{1.093378in}{0.632904in}}%
\pgfpathlineto{\pgfqpoint{1.089864in}{0.632793in}}%
\pgfpathclose%
\pgfusepath{fill}%
\end{pgfscope}%
\begin{pgfscope}%
\pgfpathrectangle{\pgfqpoint{0.407592in}{0.222904in}}{\pgfqpoint{1.411163in}{1.411163in}}%
\pgfusepath{clip}%
\pgfsetbuttcap%
\pgfsetroundjoin%
\definecolor{currentfill}{rgb}{0.749020,0.772549,0.792157}%
\pgfsetfillcolor{currentfill}%
\pgfsetfillopacity{0.160000}%
\pgfsetlinewidth{0.000000pt}%
\definecolor{currentstroke}{rgb}{0.000000,0.000000,0.000000}%
\pgfsetstrokecolor{currentstroke}%
\pgfsetdash{}{0pt}%
\pgfpathmoveto{\pgfqpoint{1.082856in}{0.632874in}}%
\pgfpathlineto{\pgfqpoint{1.080992in}{0.638523in}}%
\pgfpathlineto{\pgfqpoint{1.079119in}{0.644197in}}%
\pgfpathlineto{\pgfqpoint{1.084387in}{0.644060in}}%
\pgfpathlineto{\pgfqpoint{1.089679in}{0.644076in}}%
\pgfpathlineto{\pgfqpoint{1.089772in}{0.638421in}}%
\pgfpathlineto{\pgfqpoint{1.089864in}{0.632793in}}%
\pgfpathlineto{\pgfqpoint{1.086352in}{0.632783in}}%
\pgfpathlineto{\pgfqpoint{1.082856in}{0.632874in}}%
\pgfpathclose%
\pgfusepath{fill}%
\end{pgfscope}%
\begin{pgfscope}%
\pgfpathrectangle{\pgfqpoint{0.407592in}{0.222904in}}{\pgfqpoint{1.411163in}{1.411163in}}%
\pgfusepath{clip}%
\pgfsetbuttcap%
\pgfsetroundjoin%
\definecolor{currentfill}{rgb}{0.749020,0.772549,0.792157}%
\pgfsetfillcolor{currentfill}%
\pgfsetfillopacity{0.160000}%
\pgfsetlinewidth{0.000000pt}%
\definecolor{currentstroke}{rgb}{0.000000,0.000000,0.000000}%
\pgfsetstrokecolor{currentstroke}%
\pgfsetdash{}{0pt}%
\pgfpathmoveto{\pgfqpoint{1.273234in}{0.759054in}}%
\pgfpathlineto{\pgfqpoint{1.286479in}{0.769803in}}%
\pgfpathlineto{\pgfqpoint{1.299747in}{0.780571in}}%
\pgfpathlineto{\pgfqpoint{1.301544in}{0.786442in}}%
\pgfpathlineto{\pgfqpoint{1.302409in}{0.792323in}}%
\pgfpathlineto{\pgfqpoint{1.289016in}{0.780848in}}%
\pgfpathlineto{\pgfqpoint{1.275641in}{0.769388in}}%
\pgfpathlineto{\pgfqpoint{1.274845in}{0.764216in}}%
\pgfpathlineto{\pgfqpoint{1.273234in}{0.759054in}}%
\pgfpathclose%
\pgfusepath{fill}%
\end{pgfscope}%
\begin{pgfscope}%
\pgfpathrectangle{\pgfqpoint{0.407592in}{0.222904in}}{\pgfqpoint{1.411163in}{1.411163in}}%
\pgfusepath{clip}%
\pgfsetbuttcap%
\pgfsetroundjoin%
\definecolor{currentfill}{rgb}{0.749020,0.772549,0.792157}%
\pgfsetfillcolor{currentfill}%
\pgfsetfillopacity{0.160000}%
\pgfsetlinewidth{0.000000pt}%
\definecolor{currentstroke}{rgb}{0.000000,0.000000,0.000000}%
\pgfsetstrokecolor{currentstroke}%
\pgfsetdash{}{0pt}%
\pgfpathmoveto{\pgfqpoint{1.146714in}{0.655440in}}%
\pgfpathlineto{\pgfqpoint{1.156244in}{0.663017in}}%
\pgfpathlineto{\pgfqpoint{1.165807in}{0.670619in}}%
\pgfpathlineto{\pgfqpoint{1.170563in}{0.672945in}}%
\pgfpathlineto{\pgfqpoint{1.174958in}{0.675392in}}%
\pgfpathlineto{\pgfqpoint{1.164250in}{0.667195in}}%
\pgfpathlineto{\pgfqpoint{1.153575in}{0.659024in}}%
\pgfpathlineto{\pgfqpoint{1.150279in}{0.657186in}}%
\pgfpathlineto{\pgfqpoint{1.146714in}{0.655440in}}%
\pgfpathclose%
\pgfusepath{fill}%
\end{pgfscope}%
\begin{pgfscope}%
\pgfpathrectangle{\pgfqpoint{0.407592in}{0.222904in}}{\pgfqpoint{1.411163in}{1.411163in}}%
\pgfusepath{clip}%
\pgfsetbuttcap%
\pgfsetroundjoin%
\definecolor{currentfill}{rgb}{0.749020,0.772549,0.792157}%
\pgfsetfillcolor{currentfill}%
\pgfsetfillopacity{0.160000}%
\pgfsetlinewidth{0.000000pt}%
\definecolor{currentstroke}{rgb}{0.000000,0.000000,0.000000}%
\pgfsetstrokecolor{currentstroke}%
\pgfsetdash{}{0pt}%
\pgfpathmoveto{\pgfqpoint{1.356160in}{0.838375in}}%
\pgfpathlineto{\pgfqpoint{1.369643in}{0.849927in}}%
\pgfpathlineto{\pgfqpoint{1.383143in}{0.861493in}}%
\pgfpathlineto{\pgfqpoint{1.382859in}{0.869419in}}%
\pgfpathlineto{\pgfqpoint{1.381300in}{0.877277in}}%
\pgfpathlineto{\pgfqpoint{1.367945in}{0.865038in}}%
\pgfpathlineto{\pgfqpoint{1.354602in}{0.852811in}}%
\pgfpathlineto{\pgfqpoint{1.355961in}{0.845622in}}%
\pgfpathlineto{\pgfqpoint{1.356160in}{0.838375in}}%
\pgfpathclose%
\pgfusepath{fill}%
\end{pgfscope}%
\begin{pgfscope}%
\pgfpathrectangle{\pgfqpoint{0.407592in}{0.222904in}}{\pgfqpoint{1.411163in}{1.411163in}}%
\pgfusepath{clip}%
\pgfsetbuttcap%
\pgfsetroundjoin%
\definecolor{currentfill}{rgb}{0.749020,0.772549,0.792157}%
\pgfsetfillcolor{currentfill}%
\pgfsetfillopacity{0.160000}%
\pgfsetlinewidth{0.000000pt}%
\definecolor{currentstroke}{rgb}{0.000000,0.000000,0.000000}%
\pgfsetstrokecolor{currentstroke}%
\pgfsetdash{}{0pt}%
\pgfpathmoveto{\pgfqpoint{0.903768in}{0.768351in}}%
\pgfpathlineto{\pgfqpoint{0.890312in}{0.779739in}}%
\pgfpathlineto{\pgfqpoint{0.876838in}{0.791143in}}%
\pgfpathlineto{\pgfqpoint{0.876634in}{0.785263in}}%
\pgfpathlineto{\pgfqpoint{0.877364in}{0.779396in}}%
\pgfpathlineto{\pgfqpoint{0.890849in}{0.768699in}}%
\pgfpathlineto{\pgfqpoint{0.904309in}{0.758021in}}%
\pgfpathlineto{\pgfqpoint{0.903630in}{0.763179in}}%
\pgfpathlineto{\pgfqpoint{0.903768in}{0.768351in}}%
\pgfpathclose%
\pgfusepath{fill}%
\end{pgfscope}%
\begin{pgfscope}%
\pgfpathrectangle{\pgfqpoint{0.407592in}{0.222904in}}{\pgfqpoint{1.411163in}{1.411163in}}%
\pgfusepath{clip}%
\pgfsetbuttcap%
\pgfsetroundjoin%
\definecolor{currentfill}{rgb}{0.749020,0.772549,0.792157}%
\pgfsetfillcolor{currentfill}%
\pgfsetfillopacity{0.160000}%
\pgfsetlinewidth{0.000000pt}%
\definecolor{currentstroke}{rgb}{0.000000,0.000000,0.000000}%
\pgfsetstrokecolor{currentstroke}%
\pgfsetdash{}{0pt}%
\pgfpathmoveto{\pgfqpoint{1.542494in}{1.024999in}}%
\pgfpathlineto{\pgfqpoint{1.556005in}{1.037381in}}%
\pgfpathlineto{\pgfqpoint{1.569529in}{1.049775in}}%
\pgfpathlineto{\pgfqpoint{1.564168in}{1.061865in}}%
\pgfpathlineto{\pgfqpoint{1.556803in}{1.073712in}}%
\pgfpathlineto{\pgfqpoint{1.543736in}{1.060741in}}%
\pgfpathlineto{\pgfqpoint{1.530675in}{1.047775in}}%
\pgfpathlineto{\pgfqpoint{1.537533in}{1.036500in}}%
\pgfpathlineto{\pgfqpoint{1.542494in}{1.024999in}}%
\pgfpathclose%
\pgfusepath{fill}%
\end{pgfscope}%
\begin{pgfscope}%
\pgfpathrectangle{\pgfqpoint{0.407592in}{0.222904in}}{\pgfqpoint{1.411163in}{1.411163in}}%
\pgfusepath{clip}%
\pgfsetbuttcap%
\pgfsetroundjoin%
\definecolor{currentfill}{rgb}{0.749020,0.772549,0.792157}%
\pgfsetfillcolor{currentfill}%
\pgfsetfillopacity{0.160000}%
\pgfsetlinewidth{0.000000pt}%
\definecolor{currentstroke}{rgb}{0.000000,0.000000,0.000000}%
\pgfsetstrokecolor{currentstroke}%
\pgfsetdash{}{0pt}%
\pgfpathmoveto{\pgfqpoint{0.681774in}{1.019736in}}%
\pgfpathlineto{\pgfqpoint{0.668902in}{1.032631in}}%
\pgfpathlineto{\pgfqpoint{0.656023in}{1.045533in}}%
\pgfpathlineto{\pgfqpoint{0.646960in}{1.034209in}}%
\pgfpathlineto{\pgfqpoint{0.639752in}{1.022669in}}%
\pgfpathlineto{\pgfqpoint{0.653208in}{1.010358in}}%
\pgfpathlineto{\pgfqpoint{0.666651in}{0.998060in}}%
\pgfpathlineto{\pgfqpoint{0.673337in}{1.008998in}}%
\pgfpathlineto{\pgfqpoint{0.681774in}{1.019736in}}%
\pgfpathclose%
\pgfusepath{fill}%
\end{pgfscope}%
\begin{pgfscope}%
\pgfpathrectangle{\pgfqpoint{0.407592in}{0.222904in}}{\pgfqpoint{1.411163in}{1.411163in}}%
\pgfusepath{clip}%
\pgfsetbuttcap%
\pgfsetroundjoin%
\definecolor{currentfill}{rgb}{0.749020,0.772549,0.792157}%
\pgfsetfillcolor{currentfill}%
\pgfsetfillopacity{0.160000}%
\pgfsetlinewidth{0.000000pt}%
\definecolor{currentstroke}{rgb}{0.000000,0.000000,0.000000}%
\pgfsetstrokecolor{currentstroke}%
\pgfsetdash{}{0pt}%
\pgfpathmoveto{\pgfqpoint{0.826982in}{0.851374in}}%
\pgfpathlineto{\pgfqpoint{0.813690in}{0.863535in}}%
\pgfpathlineto{\pgfqpoint{0.800386in}{0.875707in}}%
\pgfpathlineto{\pgfqpoint{0.797363in}{0.867833in}}%
\pgfpathlineto{\pgfqpoint{0.795604in}{0.859898in}}%
\pgfpathlineto{\pgfqpoint{0.809190in}{0.848400in}}%
\pgfpathlineto{\pgfqpoint{0.822757in}{0.836917in}}%
\pgfpathlineto{\pgfqpoint{0.824295in}{0.844172in}}%
\pgfpathlineto{\pgfqpoint{0.826982in}{0.851374in}}%
\pgfpathclose%
\pgfusepath{fill}%
\end{pgfscope}%
\begin{pgfscope}%
\pgfpathrectangle{\pgfqpoint{0.407592in}{0.222904in}}{\pgfqpoint{1.411163in}{1.411163in}}%
\pgfusepath{clip}%
\pgfsetbuttcap%
\pgfsetroundjoin%
\definecolor{currentfill}{rgb}{0.749020,0.772549,0.792157}%
\pgfsetfillcolor{currentfill}%
\pgfsetfillopacity{0.160000}%
\pgfsetlinewidth{0.000000pt}%
\definecolor{currentstroke}{rgb}{0.000000,0.000000,0.000000}%
\pgfsetstrokecolor{currentstroke}%
\pgfsetdash{}{0pt}%
\pgfpathmoveto{\pgfqpoint{1.026055in}{0.656828in}}%
\pgfpathlineto{\pgfqpoint{1.015233in}{0.664635in}}%
\pgfpathlineto{\pgfqpoint{1.004374in}{0.672468in}}%
\pgfpathlineto{\pgfqpoint{1.008723in}{0.670169in}}%
\pgfpathlineto{\pgfqpoint{1.013435in}{0.668004in}}%
\pgfpathlineto{\pgfqpoint{1.023158in}{0.660727in}}%
\pgfpathlineto{\pgfqpoint{1.032846in}{0.653477in}}%
\pgfpathlineto{\pgfqpoint{1.029315in}{0.655102in}}%
\pgfpathlineto{\pgfqpoint{1.026055in}{0.656828in}}%
\pgfpathclose%
\pgfusepath{fill}%
\end{pgfscope}%
\begin{pgfscope}%
\pgfpathrectangle{\pgfqpoint{0.407592in}{0.222904in}}{\pgfqpoint{1.411163in}{1.411163in}}%
\pgfusepath{clip}%
\pgfsetbuttcap%
\pgfsetroundjoin%
\definecolor{currentfill}{rgb}{0.749020,0.772549,0.792157}%
\pgfsetfillcolor{currentfill}%
\pgfsetfillopacity{0.160000}%
\pgfsetlinewidth{0.000000pt}%
\definecolor{currentstroke}{rgb}{0.000000,0.000000,0.000000}%
\pgfsetstrokecolor{currentstroke}%
\pgfsetdash{}{0pt}%
\pgfpathmoveto{\pgfqpoint{1.178973in}{0.677951in}}%
\pgfpathlineto{\pgfqpoint{1.190214in}{0.686491in}}%
\pgfpathlineto{\pgfqpoint{1.201490in}{0.695057in}}%
\pgfpathlineto{\pgfqpoint{1.206003in}{0.698373in}}%
\pgfpathlineto{\pgfqpoint{1.209996in}{0.701796in}}%
\pgfpathlineto{\pgfqpoint{1.197880in}{0.692563in}}%
\pgfpathlineto{\pgfqpoint{1.185795in}{0.683354in}}%
\pgfpathlineto{\pgfqpoint{1.182590in}{0.680609in}}%
\pgfpathlineto{\pgfqpoint{1.178973in}{0.677951in}}%
\pgfpathclose%
\pgfusepath{fill}%
\end{pgfscope}%
\begin{pgfscope}%
\pgfpathrectangle{\pgfqpoint{0.407592in}{0.222904in}}{\pgfqpoint{1.411163in}{1.411163in}}%
\pgfusepath{clip}%
\pgfsetbuttcap%
\pgfsetroundjoin%
\definecolor{currentfill}{rgb}{0.749020,0.772549,0.792157}%
\pgfsetfillcolor{currentfill}%
\pgfsetfillopacity{0.160000}%
\pgfsetlinewidth{0.000000pt}%
\definecolor{currentstroke}{rgb}{0.000000,0.000000,0.000000}%
\pgfsetstrokecolor{currentstroke}%
\pgfsetdash{}{0pt}%
\pgfpathmoveto{\pgfqpoint{1.569529in}{1.049775in}}%
\pgfpathlineto{\pgfqpoint{1.583065in}{1.062179in}}%
\pgfpathlineto{\pgfqpoint{1.596613in}{1.074595in}}%
\pgfpathlineto{\pgfqpoint{1.590840in}{1.087265in}}%
\pgfpathlineto{\pgfqpoint{1.582955in}{1.099674in}}%
\pgfpathlineto{\pgfqpoint{1.569876in}{1.086690in}}%
\pgfpathlineto{\pgfqpoint{1.556803in}{1.073712in}}%
\pgfpathlineto{\pgfqpoint{1.564168in}{1.061865in}}%
\pgfpathlineto{\pgfqpoint{1.569529in}{1.049775in}}%
\pgfpathclose%
\pgfusepath{fill}%
\end{pgfscope}%
\begin{pgfscope}%
\pgfpathrectangle{\pgfqpoint{0.407592in}{0.222904in}}{\pgfqpoint{1.411163in}{1.411163in}}%
\pgfusepath{clip}%
\pgfsetbuttcap%
\pgfsetroundjoin%
\definecolor{currentfill}{rgb}{0.749020,0.772549,0.792157}%
\pgfsetfillcolor{currentfill}%
\pgfsetfillopacity{0.160000}%
\pgfsetlinewidth{0.000000pt}%
\definecolor{currentstroke}{rgb}{0.000000,0.000000,0.000000}%
\pgfsetstrokecolor{currentstroke}%
\pgfsetdash{}{0pt}%
\pgfpathmoveto{\pgfqpoint{0.990974in}{0.682797in}}%
\pgfpathlineto{\pgfqpoint{0.978420in}{0.691937in}}%
\pgfpathlineto{\pgfqpoint{0.965832in}{0.701101in}}%
\pgfpathlineto{\pgfqpoint{0.969219in}{0.697698in}}%
\pgfpathlineto{\pgfqpoint{0.973145in}{0.694406in}}%
\pgfpathlineto{\pgfqpoint{0.985013in}{0.685905in}}%
\pgfpathlineto{\pgfqpoint{0.996845in}{0.677429in}}%
\pgfpathlineto{\pgfqpoint{0.993695in}{0.680068in}}%
\pgfpathlineto{\pgfqpoint{0.990974in}{0.682797in}}%
\pgfpathclose%
\pgfusepath{fill}%
\end{pgfscope}%
\begin{pgfscope}%
\pgfpathrectangle{\pgfqpoint{0.407592in}{0.222904in}}{\pgfqpoint{1.411163in}{1.411163in}}%
\pgfusepath{clip}%
\pgfsetbuttcap%
\pgfsetroundjoin%
\definecolor{currentfill}{rgb}{0.749020,0.772549,0.792157}%
\pgfsetfillcolor{currentfill}%
\pgfsetfillopacity{0.160000}%
\pgfsetlinewidth{0.000000pt}%
\definecolor{currentstroke}{rgb}{0.000000,0.000000,0.000000}%
\pgfsetstrokecolor{currentstroke}%
\pgfsetdash{}{0pt}%
\pgfpathmoveto{\pgfqpoint{0.656023in}{1.045533in}}%
\pgfpathlineto{\pgfqpoint{0.643138in}{1.058442in}}%
\pgfpathlineto{\pgfqpoint{0.630246in}{1.071357in}}%
\pgfpathlineto{\pgfqpoint{0.620544in}{1.059457in}}%
\pgfpathlineto{\pgfqpoint{0.612803in}{1.047324in}}%
\pgfpathlineto{\pgfqpoint{0.626284in}{1.034991in}}%
\pgfpathlineto{\pgfqpoint{0.639752in}{1.022669in}}%
\pgfpathlineto{\pgfqpoint{0.646960in}{1.034209in}}%
\pgfpathlineto{\pgfqpoint{0.656023in}{1.045533in}}%
\pgfpathclose%
\pgfusepath{fill}%
\end{pgfscope}%
\begin{pgfscope}%
\pgfpathrectangle{\pgfqpoint{0.407592in}{0.222904in}}{\pgfqpoint{1.411163in}{1.411163in}}%
\pgfusepath{clip}%
\pgfsetbuttcap%
\pgfsetroundjoin%
\definecolor{currentfill}{rgb}{0.749020,0.772549,0.792157}%
\pgfsetfillcolor{currentfill}%
\pgfsetfillopacity{0.160000}%
\pgfsetlinewidth{0.000000pt}%
\definecolor{currentstroke}{rgb}{0.000000,0.000000,0.000000}%
\pgfsetstrokecolor{currentstroke}%
\pgfsetdash{}{0pt}%
\pgfpathmoveto{\pgfqpoint{1.138843in}{0.652257in}}%
\pgfpathlineto{\pgfqpoint{1.147053in}{0.659304in}}%
\pgfpathlineto{\pgfqpoint{1.155295in}{0.666377in}}%
\pgfpathlineto{\pgfqpoint{1.160710in}{0.668427in}}%
\pgfpathlineto{\pgfqpoint{1.165807in}{0.670619in}}%
\pgfpathlineto{\pgfqpoint{1.156244in}{0.663017in}}%
\pgfpathlineto{\pgfqpoint{1.146714in}{0.655440in}}%
\pgfpathlineto{\pgfqpoint{1.142897in}{0.653795in}}%
\pgfpathlineto{\pgfqpoint{1.138843in}{0.652257in}}%
\pgfpathclose%
\pgfusepath{fill}%
\end{pgfscope}%
\begin{pgfscope}%
\pgfpathrectangle{\pgfqpoint{0.407592in}{0.222904in}}{\pgfqpoint{1.411163in}{1.411163in}}%
\pgfusepath{clip}%
\pgfsetbuttcap%
\pgfsetroundjoin%
\definecolor{currentfill}{rgb}{0.749020,0.772549,0.792157}%
\pgfsetfillcolor{currentfill}%
\pgfsetfillopacity{0.160000}%
\pgfsetlinewidth{0.000000pt}%
\definecolor{currentstroke}{rgb}{0.000000,0.000000,0.000000}%
\pgfsetstrokecolor{currentstroke}%
\pgfsetdash{}{0pt}%
\pgfpathmoveto{\pgfqpoint{1.383143in}{0.861493in}}%
\pgfpathlineto{\pgfqpoint{1.396662in}{0.873076in}}%
\pgfpathlineto{\pgfqpoint{1.410198in}{0.884673in}}%
\pgfpathlineto{\pgfqpoint{1.409818in}{0.893268in}}%
\pgfpathlineto{\pgfqpoint{1.408045in}{0.901787in}}%
\pgfpathlineto{\pgfqpoint{1.394666in}{0.889526in}}%
\pgfpathlineto{\pgfqpoint{1.381300in}{0.877277in}}%
\pgfpathlineto{\pgfqpoint{1.382859in}{0.869419in}}%
\pgfpathlineto{\pgfqpoint{1.383143in}{0.861493in}}%
\pgfpathclose%
\pgfusepath{fill}%
\end{pgfscope}%
\begin{pgfscope}%
\pgfpathrectangle{\pgfqpoint{0.407592in}{0.222904in}}{\pgfqpoint{1.411163in}{1.411163in}}%
\pgfusepath{clip}%
\pgfsetbuttcap%
\pgfsetroundjoin%
\definecolor{currentfill}{rgb}{0.749020,0.772549,0.792157}%
\pgfsetfillcolor{currentfill}%
\pgfsetfillopacity{0.160000}%
\pgfsetlinewidth{0.000000pt}%
\definecolor{currentstroke}{rgb}{0.000000,0.000000,0.000000}%
\pgfsetstrokecolor{currentstroke}%
\pgfsetdash{}{0pt}%
\pgfpathmoveto{\pgfqpoint{1.241920in}{0.728833in}}%
\pgfpathlineto{\pgfqpoint{1.254741in}{0.738832in}}%
\pgfpathlineto{\pgfqpoint{1.267590in}{0.748853in}}%
\pgfpathlineto{\pgfqpoint{1.270814in}{0.753925in}}%
\pgfpathlineto{\pgfqpoint{1.273234in}{0.759054in}}%
\pgfpathlineto{\pgfqpoint{1.260013in}{0.748324in}}%
\pgfpathlineto{\pgfqpoint{1.246815in}{0.737613in}}%
\pgfpathlineto{\pgfqpoint{1.244710in}{0.733198in}}%
\pgfpathlineto{\pgfqpoint{1.241920in}{0.728833in}}%
\pgfpathclose%
\pgfusepath{fill}%
\end{pgfscope}%
\begin{pgfscope}%
\pgfpathrectangle{\pgfqpoint{0.407592in}{0.222904in}}{\pgfqpoint{1.411163in}{1.411163in}}%
\pgfusepath{clip}%
\pgfsetbuttcap%
\pgfsetroundjoin%
\definecolor{currentfill}{rgb}{0.749020,0.772549,0.792157}%
\pgfsetfillcolor{currentfill}%
\pgfsetfillopacity{0.160000}%
\pgfsetlinewidth{0.000000pt}%
\definecolor{currentstroke}{rgb}{0.000000,0.000000,0.000000}%
\pgfsetstrokecolor{currentstroke}%
\pgfsetdash{}{0pt}%
\pgfpathmoveto{\pgfqpoint{0.931156in}{0.736724in}}%
\pgfpathlineto{\pgfqpoint{0.917745in}{0.747363in}}%
\pgfpathlineto{\pgfqpoint{0.904309in}{0.758021in}}%
\pgfpathlineto{\pgfqpoint{0.905805in}{0.752902in}}%
\pgfpathlineto{\pgfqpoint{0.908116in}{0.747844in}}%
\pgfpathlineto{\pgfqpoint{0.921313in}{0.737894in}}%
\pgfpathlineto{\pgfqpoint{0.934480in}{0.727966in}}%
\pgfpathlineto{\pgfqpoint{0.932470in}{0.732318in}}%
\pgfpathlineto{\pgfqpoint{0.931156in}{0.736724in}}%
\pgfpathclose%
\pgfusepath{fill}%
\end{pgfscope}%
\begin{pgfscope}%
\pgfpathrectangle{\pgfqpoint{0.407592in}{0.222904in}}{\pgfqpoint{1.411163in}{1.411163in}}%
\pgfusepath{clip}%
\pgfsetbuttcap%
\pgfsetroundjoin%
\definecolor{currentfill}{rgb}{0.749020,0.772549,0.792157}%
\pgfsetfillcolor{currentfill}%
\pgfsetfillopacity{0.160000}%
\pgfsetlinewidth{0.000000pt}%
\definecolor{currentstroke}{rgb}{0.000000,0.000000,0.000000}%
\pgfsetstrokecolor{currentstroke}%
\pgfsetdash{}{0pt}%
\pgfpathmoveto{\pgfqpoint{1.299747in}{0.780571in}}%
\pgfpathlineto{\pgfqpoint{1.313039in}{0.791358in}}%
\pgfpathlineto{\pgfqpoint{1.326354in}{0.802164in}}%
\pgfpathlineto{\pgfqpoint{1.328327in}{0.808737in}}%
\pgfpathlineto{\pgfqpoint{1.329249in}{0.815318in}}%
\pgfpathlineto{\pgfqpoint{1.315820in}{0.803813in}}%
\pgfpathlineto{\pgfqpoint{1.302409in}{0.792323in}}%
\pgfpathlineto{\pgfqpoint{1.301544in}{0.786442in}}%
\pgfpathlineto{\pgfqpoint{1.299747in}{0.780571in}}%
\pgfpathclose%
\pgfusepath{fill}%
\end{pgfscope}%
\begin{pgfscope}%
\pgfpathrectangle{\pgfqpoint{0.407592in}{0.222904in}}{\pgfqpoint{1.411163in}{1.411163in}}%
\pgfusepath{clip}%
\pgfsetbuttcap%
\pgfsetroundjoin%
\definecolor{currentfill}{rgb}{0.749020,0.772549,0.792157}%
\pgfsetfillcolor{currentfill}%
\pgfsetfillopacity{0.160000}%
\pgfsetlinewidth{0.000000pt}%
\definecolor{currentstroke}{rgb}{0.000000,0.000000,0.000000}%
\pgfsetstrokecolor{currentstroke}%
\pgfsetdash{}{0pt}%
\pgfpathmoveto{\pgfqpoint{1.209996in}{0.701796in}}%
\pgfpathlineto{\pgfqpoint{1.222144in}{0.711052in}}%
\pgfpathlineto{\pgfqpoint{1.234324in}{0.720333in}}%
\pgfpathlineto{\pgfqpoint{1.238453in}{0.724539in}}%
\pgfpathlineto{\pgfqpoint{1.241920in}{0.728833in}}%
\pgfpathlineto{\pgfqpoint{1.229127in}{0.718857in}}%
\pgfpathlineto{\pgfqpoint{1.216362in}{0.708901in}}%
\pgfpathlineto{\pgfqpoint{1.213454in}{0.705310in}}%
\pgfpathlineto{\pgfqpoint{1.209996in}{0.701796in}}%
\pgfpathclose%
\pgfusepath{fill}%
\end{pgfscope}%
\begin{pgfscope}%
\pgfpathrectangle{\pgfqpoint{0.407592in}{0.222904in}}{\pgfqpoint{1.411163in}{1.411163in}}%
\pgfusepath{clip}%
\pgfsetbuttcap%
\pgfsetroundjoin%
\definecolor{currentfill}{rgb}{0.749020,0.772549,0.792157}%
\pgfsetfillcolor{currentfill}%
\pgfsetfillopacity{0.160000}%
\pgfsetlinewidth{0.000000pt}%
\definecolor{currentstroke}{rgb}{0.000000,0.000000,0.000000}%
\pgfsetstrokecolor{currentstroke}%
\pgfsetdash{}{0pt}%
\pgfpathmoveto{\pgfqpoint{0.800386in}{0.875707in}}%
\pgfpathlineto{\pgfqpoint{0.787069in}{0.887891in}}%
\pgfpathlineto{\pgfqpoint{0.773739in}{0.900086in}}%
\pgfpathlineto{\pgfqpoint{0.770368in}{0.891549in}}%
\pgfpathlineto{\pgfqpoint{0.768377in}{0.882943in}}%
\pgfpathlineto{\pgfqpoint{0.782000in}{0.871413in}}%
\pgfpathlineto{\pgfqpoint{0.795604in}{0.859898in}}%
\pgfpathlineto{\pgfqpoint{0.797363in}{0.867833in}}%
\pgfpathlineto{\pgfqpoint{0.800386in}{0.875707in}}%
\pgfpathclose%
\pgfusepath{fill}%
\end{pgfscope}%
\begin{pgfscope}%
\pgfpathrectangle{\pgfqpoint{0.407592in}{0.222904in}}{\pgfqpoint{1.411163in}{1.411163in}}%
\pgfusepath{clip}%
\pgfsetbuttcap%
\pgfsetroundjoin%
\definecolor{currentfill}{rgb}{0.749020,0.772549,0.792157}%
\pgfsetfillcolor{currentfill}%
\pgfsetfillopacity{0.160000}%
\pgfsetlinewidth{0.000000pt}%
\definecolor{currentstroke}{rgb}{0.000000,0.000000,0.000000}%
\pgfsetstrokecolor{currentstroke}%
\pgfsetdash{}{0pt}%
\pgfpathmoveto{\pgfqpoint{1.596613in}{1.074595in}}%
\pgfpathlineto{\pgfqpoint{1.610173in}{1.087022in}}%
\pgfpathlineto{\pgfqpoint{1.623745in}{1.099460in}}%
\pgfpathlineto{\pgfqpoint{1.617549in}{1.112699in}}%
\pgfpathlineto{\pgfqpoint{1.609132in}{1.125661in}}%
\pgfpathlineto{\pgfqpoint{1.596040in}{1.112664in}}%
\pgfpathlineto{\pgfqpoint{1.582955in}{1.099674in}}%
\pgfpathlineto{\pgfqpoint{1.590840in}{1.087265in}}%
\pgfpathlineto{\pgfqpoint{1.596613in}{1.074595in}}%
\pgfpathclose%
\pgfusepath{fill}%
\end{pgfscope}%
\begin{pgfscope}%
\pgfpathrectangle{\pgfqpoint{0.407592in}{0.222904in}}{\pgfqpoint{1.411163in}{1.411163in}}%
\pgfusepath{clip}%
\pgfsetbuttcap%
\pgfsetroundjoin%
\definecolor{currentfill}{rgb}{0.749020,0.772549,0.792157}%
\pgfsetfillcolor{currentfill}%
\pgfsetfillopacity{0.160000}%
\pgfsetlinewidth{0.000000pt}%
\definecolor{currentstroke}{rgb}{0.000000,0.000000,0.000000}%
\pgfsetstrokecolor{currentstroke}%
\pgfsetdash{}{0pt}%
\pgfpathmoveto{\pgfqpoint{0.960727in}{0.708176in}}%
\pgfpathlineto{\pgfqpoint{0.947618in}{0.718060in}}%
\pgfpathlineto{\pgfqpoint{0.934480in}{0.727966in}}%
\pgfpathlineto{\pgfqpoint{0.937179in}{0.723687in}}%
\pgfpathlineto{\pgfqpoint{0.940558in}{0.719502in}}%
\pgfpathlineto{\pgfqpoint{0.953212in}{0.710289in}}%
\pgfpathlineto{\pgfqpoint{0.965832in}{0.701101in}}%
\pgfpathlineto{\pgfqpoint{0.962998in}{0.704599in}}%
\pgfpathlineto{\pgfqpoint{0.960727in}{0.708176in}}%
\pgfpathclose%
\pgfusepath{fill}%
\end{pgfscope}%
\begin{pgfscope}%
\pgfpathrectangle{\pgfqpoint{0.407592in}{0.222904in}}{\pgfqpoint{1.411163in}{1.411163in}}%
\pgfusepath{clip}%
\pgfsetbuttcap%
\pgfsetroundjoin%
\definecolor{currentfill}{rgb}{0.749020,0.772549,0.792157}%
\pgfsetfillcolor{currentfill}%
\pgfsetfillopacity{0.160000}%
\pgfsetlinewidth{0.000000pt}%
\definecolor{currentstroke}{rgb}{0.000000,0.000000,0.000000}%
\pgfsetstrokecolor{currentstroke}%
\pgfsetdash{}{0pt}%
\pgfpathmoveto{\pgfqpoint{1.032846in}{0.653477in}}%
\pgfpathlineto{\pgfqpoint{1.023158in}{0.660727in}}%
\pgfpathlineto{\pgfqpoint{1.013435in}{0.668004in}}%
\pgfpathlineto{\pgfqpoint{1.018491in}{0.665984in}}%
\pgfpathlineto{\pgfqpoint{1.023870in}{0.664119in}}%
\pgfpathlineto{\pgfqpoint{1.032280in}{0.657328in}}%
\pgfpathlineto{\pgfqpoint{1.040657in}{0.650563in}}%
\pgfpathlineto{\pgfqpoint{1.036632in}{0.651962in}}%
\pgfpathlineto{\pgfqpoint{1.032846in}{0.653477in}}%
\pgfpathclose%
\pgfusepath{fill}%
\end{pgfscope}%
\begin{pgfscope}%
\pgfpathrectangle{\pgfqpoint{0.407592in}{0.222904in}}{\pgfqpoint{1.411163in}{1.411163in}}%
\pgfusepath{clip}%
\pgfsetbuttcap%
\pgfsetroundjoin%
\definecolor{currentfill}{rgb}{0.749020,0.772549,0.792157}%
\pgfsetfillcolor{currentfill}%
\pgfsetfillopacity{0.160000}%
\pgfsetlinewidth{0.000000pt}%
\definecolor{currentstroke}{rgb}{0.000000,0.000000,0.000000}%
\pgfsetstrokecolor{currentstroke}%
\pgfsetdash{}{0pt}%
\pgfpathmoveto{\pgfqpoint{0.630246in}{1.071357in}}%
\pgfpathlineto{\pgfqpoint{0.617348in}{1.084279in}}%
\pgfpathlineto{\pgfqpoint{0.604443in}{1.097208in}}%
\pgfpathlineto{\pgfqpoint{0.594090in}{1.084742in}}%
\pgfpathlineto{\pgfqpoint{0.585803in}{1.072026in}}%
\pgfpathlineto{\pgfqpoint{0.599309in}{1.059669in}}%
\pgfpathlineto{\pgfqpoint{0.612803in}{1.047324in}}%
\pgfpathlineto{\pgfqpoint{0.620544in}{1.059457in}}%
\pgfpathlineto{\pgfqpoint{0.630246in}{1.071357in}}%
\pgfpathclose%
\pgfusepath{fill}%
\end{pgfscope}%
\begin{pgfscope}%
\pgfpathrectangle{\pgfqpoint{0.407592in}{0.222904in}}{\pgfqpoint{1.411163in}{1.411163in}}%
\pgfusepath{clip}%
\pgfsetbuttcap%
\pgfsetroundjoin%
\definecolor{currentfill}{rgb}{0.749020,0.772549,0.792157}%
\pgfsetfillcolor{currentfill}%
\pgfsetfillopacity{0.160000}%
\pgfsetlinewidth{0.000000pt}%
\definecolor{currentstroke}{rgb}{0.000000,0.000000,0.000000}%
\pgfsetstrokecolor{currentstroke}%
\pgfsetdash{}{0pt}%
\pgfpathmoveto{\pgfqpoint{0.876838in}{0.791143in}}%
\pgfpathlineto{\pgfqpoint{0.863346in}{0.802563in}}%
\pgfpathlineto{\pgfqpoint{0.849835in}{0.813999in}}%
\pgfpathlineto{\pgfqpoint{0.849551in}{0.807417in}}%
\pgfpathlineto{\pgfqpoint{0.850322in}{0.800848in}}%
\pgfpathlineto{\pgfqpoint{0.863855in}{0.790113in}}%
\pgfpathlineto{\pgfqpoint{0.877364in}{0.779396in}}%
\pgfpathlineto{\pgfqpoint{0.876634in}{0.785263in}}%
\pgfpathlineto{\pgfqpoint{0.876838in}{0.791143in}}%
\pgfpathclose%
\pgfusepath{fill}%
\end{pgfscope}%
\begin{pgfscope}%
\pgfpathrectangle{\pgfqpoint{0.407592in}{0.222904in}}{\pgfqpoint{1.411163in}{1.411163in}}%
\pgfusepath{clip}%
\pgfsetbuttcap%
\pgfsetroundjoin%
\definecolor{currentfill}{rgb}{0.749020,0.772549,0.792157}%
\pgfsetfillcolor{currentfill}%
\pgfsetfillopacity{0.160000}%
\pgfsetlinewidth{0.000000pt}%
\definecolor{currentstroke}{rgb}{0.000000,0.000000,0.000000}%
\pgfsetstrokecolor{currentstroke}%
\pgfsetdash{}{0pt}%
\pgfpathmoveto{\pgfqpoint{1.130099in}{0.649533in}}%
\pgfpathlineto{\pgfqpoint{1.136839in}{0.656125in}}%
\pgfpathlineto{\pgfqpoint{1.143606in}{0.662744in}}%
\pgfpathlineto{\pgfqpoint{1.149585in}{0.664480in}}%
\pgfpathlineto{\pgfqpoint{1.155295in}{0.666377in}}%
\pgfpathlineto{\pgfqpoint{1.147053in}{0.659304in}}%
\pgfpathlineto{\pgfqpoint{1.138843in}{0.652257in}}%
\pgfpathlineto{\pgfqpoint{1.134571in}{0.650834in}}%
\pgfpathlineto{\pgfqpoint{1.130099in}{0.649533in}}%
\pgfpathclose%
\pgfusepath{fill}%
\end{pgfscope}%
\begin{pgfscope}%
\pgfpathrectangle{\pgfqpoint{0.407592in}{0.222904in}}{\pgfqpoint{1.411163in}{1.411163in}}%
\pgfusepath{clip}%
\pgfsetbuttcap%
\pgfsetroundjoin%
\definecolor{currentfill}{rgb}{0.749020,0.772549,0.792157}%
\pgfsetfillcolor{currentfill}%
\pgfsetfillopacity{0.160000}%
\pgfsetlinewidth{0.000000pt}%
\definecolor{currentstroke}{rgb}{0.000000,0.000000,0.000000}%
\pgfsetstrokecolor{currentstroke}%
\pgfsetdash{}{0pt}%
\pgfpathmoveto{\pgfqpoint{1.174958in}{0.675392in}}%
\pgfpathlineto{\pgfqpoint{1.185700in}{0.683615in}}%
\pgfpathlineto{\pgfqpoint{1.196476in}{0.691865in}}%
\pgfpathlineto{\pgfqpoint{1.201490in}{0.695057in}}%
\pgfpathlineto{\pgfqpoint{1.190214in}{0.686491in}}%
\pgfpathlineto{\pgfqpoint{1.178973in}{0.677951in}}%
\pgfusepath{fill}%
\end{pgfscope}%
\begin{pgfscope}%
\pgfpathrectangle{\pgfqpoint{0.407592in}{0.222904in}}{\pgfqpoint{1.411163in}{1.411163in}}%
\pgfusepath{clip}%
\pgfsetbuttcap%
\pgfsetroundjoin%
\definecolor{currentfill}{rgb}{0.749020,0.772549,0.792157}%
\pgfsetfillcolor{currentfill}%
\pgfsetfillopacity{0.160000}%
\pgfsetlinewidth{0.000000pt}%
\definecolor{currentstroke}{rgb}{0.000000,0.000000,0.000000}%
\pgfsetstrokecolor{currentstroke}%
\pgfsetdash{}{0pt}%
\pgfpathmoveto{\pgfqpoint{1.410198in}{0.884673in}}%
\pgfpathlineto{\pgfqpoint{1.423753in}{0.896287in}}%
\pgfpathlineto{\pgfqpoint{1.437326in}{0.907915in}}%
\pgfpathlineto{\pgfqpoint{1.436836in}{0.917171in}}%
\pgfpathlineto{\pgfqpoint{1.434838in}{0.926341in}}%
\pgfpathlineto{\pgfqpoint{1.421436in}{0.914058in}}%
\pgfpathlineto{\pgfqpoint{1.408045in}{0.901787in}}%
\pgfpathlineto{\pgfqpoint{1.409818in}{0.893268in}}%
\pgfpathlineto{\pgfqpoint{1.410198in}{0.884673in}}%
\pgfpathclose%
\pgfusepath{fill}%
\end{pgfscope}%
\begin{pgfscope}%
\pgfpathrectangle{\pgfqpoint{0.407592in}{0.222904in}}{\pgfqpoint{1.411163in}{1.411163in}}%
\pgfusepath{clip}%
\pgfsetbuttcap%
\pgfsetroundjoin%
\definecolor{currentfill}{rgb}{0.749020,0.772549,0.792157}%
\pgfsetfillcolor{currentfill}%
\pgfsetfillopacity{0.160000}%
\pgfsetlinewidth{0.000000pt}%
\definecolor{currentstroke}{rgb}{0.000000,0.000000,0.000000}%
\pgfsetstrokecolor{currentstroke}%
\pgfsetdash{}{0pt}%
\pgfpathmoveto{\pgfqpoint{1.623745in}{1.099460in}}%
\pgfpathlineto{\pgfqpoint{1.637330in}{1.111909in}}%
\pgfpathlineto{\pgfqpoint{1.650927in}{1.124370in}}%
\pgfpathlineto{\pgfqpoint{1.644293in}{1.138168in}}%
\pgfpathlineto{\pgfqpoint{1.635333in}{1.151671in}}%
\pgfpathlineto{\pgfqpoint{1.622229in}{1.138663in}}%
\pgfpathlineto{\pgfqpoint{1.609132in}{1.125661in}}%
\pgfpathlineto{\pgfqpoint{1.617549in}{1.112699in}}%
\pgfpathlineto{\pgfqpoint{1.623745in}{1.099460in}}%
\pgfpathclose%
\pgfusepath{fill}%
\end{pgfscope}%
\begin{pgfscope}%
\pgfpathrectangle{\pgfqpoint{0.407592in}{0.222904in}}{\pgfqpoint{1.411163in}{1.411163in}}%
\pgfusepath{clip}%
\pgfsetbuttcap%
\pgfsetroundjoin%
\definecolor{currentfill}{rgb}{0.749020,0.772549,0.792157}%
\pgfsetfillcolor{currentfill}%
\pgfsetfillopacity{0.160000}%
\pgfsetlinewidth{0.000000pt}%
\definecolor{currentstroke}{rgb}{0.000000,0.000000,0.000000}%
\pgfsetstrokecolor{currentstroke}%
\pgfsetdash{}{0pt}%
\pgfpathmoveto{\pgfqpoint{0.604443in}{1.097208in}}%
\pgfpathlineto{\pgfqpoint{0.591531in}{1.110143in}}%
\pgfpathlineto{\pgfqpoint{0.578613in}{1.123085in}}%
\pgfpathlineto{\pgfqpoint{0.567598in}{1.110063in}}%
\pgfpathlineto{\pgfqpoint{0.558751in}{1.096775in}}%
\pgfpathlineto{\pgfqpoint{0.572283in}{1.084395in}}%
\pgfpathlineto{\pgfqpoint{0.585803in}{1.072026in}}%
\pgfpathlineto{\pgfqpoint{0.594090in}{1.084742in}}%
\pgfpathlineto{\pgfqpoint{0.604443in}{1.097208in}}%
\pgfpathclose%
\pgfusepath{fill}%
\end{pgfscope}%
\begin{pgfscope}%
\pgfpathrectangle{\pgfqpoint{0.407592in}{0.222904in}}{\pgfqpoint{1.411163in}{1.411163in}}%
\pgfusepath{clip}%
\pgfsetbuttcap%
\pgfsetroundjoin%
\definecolor{currentfill}{rgb}{0.749020,0.772549,0.792157}%
\pgfsetfillcolor{currentfill}%
\pgfsetfillopacity{0.160000}%
\pgfsetlinewidth{0.000000pt}%
\definecolor{currentstroke}{rgb}{0.000000,0.000000,0.000000}%
\pgfsetstrokecolor{currentstroke}%
\pgfsetdash{}{0pt}%
\pgfpathmoveto{\pgfqpoint{1.040657in}{0.650563in}}%
\pgfpathlineto{\pgfqpoint{1.032280in}{0.657328in}}%
\pgfpathlineto{\pgfqpoint{1.023870in}{0.664119in}}%
\pgfpathlineto{\pgfqpoint{1.029545in}{0.662416in}}%
\pgfpathlineto{\pgfqpoint{1.035494in}{0.660885in}}%
\pgfpathlineto{\pgfqpoint{1.042437in}{0.654498in}}%
\pgfpathlineto{\pgfqpoint{1.049350in}{0.648138in}}%
\pgfpathlineto{\pgfqpoint{1.044902in}{0.649287in}}%
\pgfpathlineto{\pgfqpoint{1.040657in}{0.650563in}}%
\pgfpathclose%
\pgfusepath{fill}%
\end{pgfscope}%
\begin{pgfscope}%
\pgfpathrectangle{\pgfqpoint{0.407592in}{0.222904in}}{\pgfqpoint{1.411163in}{1.411163in}}%
\pgfusepath{clip}%
\pgfsetbuttcap%
\pgfsetroundjoin%
\definecolor{currentfill}{rgb}{0.749020,0.772549,0.792157}%
\pgfsetfillcolor{currentfill}%
\pgfsetfillopacity{0.160000}%
\pgfsetlinewidth{0.000000pt}%
\definecolor{currentstroke}{rgb}{0.000000,0.000000,0.000000}%
\pgfsetstrokecolor{currentstroke}%
\pgfsetdash{}{0pt}%
\pgfpathmoveto{\pgfqpoint{0.996845in}{0.677429in}}%
\pgfpathlineto{\pgfqpoint{0.985013in}{0.685905in}}%
\pgfpathlineto{\pgfqpoint{0.973145in}{0.694406in}}%
\pgfpathlineto{\pgfqpoint{0.977594in}{0.691240in}}%
\pgfpathlineto{\pgfqpoint{0.982547in}{0.688214in}}%
\pgfpathlineto{\pgfqpoint{0.993479in}{0.680328in}}%
\pgfpathlineto{\pgfqpoint{1.004374in}{0.672468in}}%
\pgfpathlineto{\pgfqpoint{1.000410in}{0.674892in}}%
\pgfpathlineto{\pgfqpoint{0.996845in}{0.677429in}}%
\pgfpathclose%
\pgfusepath{fill}%
\end{pgfscope}%
\begin{pgfscope}%
\pgfpathrectangle{\pgfqpoint{0.407592in}{0.222904in}}{\pgfqpoint{1.411163in}{1.411163in}}%
\pgfusepath{clip}%
\pgfsetbuttcap%
\pgfsetroundjoin%
\definecolor{currentfill}{rgb}{0.749020,0.772549,0.792157}%
\pgfsetfillcolor{currentfill}%
\pgfsetfillopacity{0.160000}%
\pgfsetlinewidth{0.000000pt}%
\definecolor{currentstroke}{rgb}{0.000000,0.000000,0.000000}%
\pgfsetstrokecolor{currentstroke}%
\pgfsetdash{}{0pt}%
\pgfpathmoveto{\pgfqpoint{0.773739in}{0.900086in}}%
\pgfpathlineto{\pgfqpoint{0.760397in}{0.912292in}}%
\pgfpathlineto{\pgfqpoint{0.747043in}{0.924510in}}%
\pgfpathlineto{\pgfqpoint{0.743311in}{0.915320in}}%
\pgfpathlineto{\pgfqpoint{0.741074in}{0.906052in}}%
\pgfpathlineto{\pgfqpoint{0.754735in}{0.894489in}}%
\pgfpathlineto{\pgfqpoint{0.768377in}{0.882943in}}%
\pgfpathlineto{\pgfqpoint{0.770368in}{0.891549in}}%
\pgfpathlineto{\pgfqpoint{0.773739in}{0.900086in}}%
\pgfpathclose%
\pgfusepath{fill}%
\end{pgfscope}%
\begin{pgfscope}%
\pgfpathrectangle{\pgfqpoint{0.407592in}{0.222904in}}{\pgfqpoint{1.411163in}{1.411163in}}%
\pgfusepath{clip}%
\pgfsetbuttcap%
\pgfsetroundjoin%
\definecolor{currentfill}{rgb}{0.749020,0.772549,0.792157}%
\pgfsetfillcolor{currentfill}%
\pgfsetfillopacity{0.160000}%
\pgfsetlinewidth{0.000000pt}%
\definecolor{currentstroke}{rgb}{0.000000,0.000000,0.000000}%
\pgfsetstrokecolor{currentstroke}%
\pgfsetdash{}{0pt}%
\pgfpathmoveto{\pgfqpoint{1.120638in}{0.647317in}}%
\pgfpathlineto{\pgfqpoint{1.125782in}{0.653539in}}%
\pgfpathlineto{\pgfqpoint{1.130948in}{0.659788in}}%
\pgfpathlineto{\pgfqpoint{1.137384in}{0.661178in}}%
\pgfpathlineto{\pgfqpoint{1.143606in}{0.662744in}}%
\pgfpathlineto{\pgfqpoint{1.136839in}{0.656125in}}%
\pgfpathlineto{\pgfqpoint{1.130099in}{0.649533in}}%
\pgfpathlineto{\pgfqpoint{1.125448in}{0.648358in}}%
\pgfpathlineto{\pgfqpoint{1.120638in}{0.647317in}}%
\pgfpathclose%
\pgfusepath{fill}%
\end{pgfscope}%
\begin{pgfscope}%
\pgfpathrectangle{\pgfqpoint{0.407592in}{0.222904in}}{\pgfqpoint{1.411163in}{1.411163in}}%
\pgfusepath{clip}%
\pgfsetbuttcap%
\pgfsetroundjoin%
\definecolor{currentfill}{rgb}{0.749020,0.772549,0.792157}%
\pgfsetfillcolor{currentfill}%
\pgfsetfillopacity{0.160000}%
\pgfsetlinewidth{0.000000pt}%
\definecolor{currentstroke}{rgb}{0.000000,0.000000,0.000000}%
\pgfsetstrokecolor{currentstroke}%
\pgfsetdash{}{0pt}%
\pgfpathmoveto{\pgfqpoint{1.326354in}{0.802164in}}%
\pgfpathlineto{\pgfqpoint{1.339692in}{0.812989in}}%
\pgfpathlineto{\pgfqpoint{1.353055in}{0.823833in}}%
\pgfpathlineto{\pgfqpoint{1.355193in}{0.831102in}}%
\pgfpathlineto{\pgfqpoint{1.356160in}{0.838375in}}%
\pgfpathlineto{\pgfqpoint{1.342696in}{0.826839in}}%
\pgfpathlineto{\pgfqpoint{1.329249in}{0.815318in}}%
\pgfpathlineto{\pgfqpoint{1.328327in}{0.808737in}}%
\pgfpathlineto{\pgfqpoint{1.326354in}{0.802164in}}%
\pgfpathclose%
\pgfusepath{fill}%
\end{pgfscope}%
\begin{pgfscope}%
\pgfpathrectangle{\pgfqpoint{0.407592in}{0.222904in}}{\pgfqpoint{1.411163in}{1.411163in}}%
\pgfusepath{clip}%
\pgfsetbuttcap%
\pgfsetroundjoin%
\definecolor{currentfill}{rgb}{0.749020,0.772549,0.792157}%
\pgfsetfillcolor{currentfill}%
\pgfsetfillopacity{0.160000}%
\pgfsetlinewidth{0.000000pt}%
\definecolor{currentstroke}{rgb}{0.000000,0.000000,0.000000}%
\pgfsetstrokecolor{currentstroke}%
\pgfsetdash{}{0pt}%
\pgfpathmoveto{\pgfqpoint{1.650927in}{1.124370in}}%
\pgfpathlineto{\pgfqpoint{1.664536in}{1.136842in}}%
\pgfpathlineto{\pgfqpoint{1.678158in}{1.149325in}}%
\pgfpathlineto{\pgfqpoint{1.671075in}{1.163672in}}%
\pgfpathlineto{\pgfqpoint{1.661559in}{1.177707in}}%
\pgfpathlineto{\pgfqpoint{1.648443in}{1.164686in}}%
\pgfpathlineto{\pgfqpoint{1.635333in}{1.151671in}}%
\pgfpathlineto{\pgfqpoint{1.644293in}{1.138168in}}%
\pgfpathlineto{\pgfqpoint{1.650927in}{1.124370in}}%
\pgfpathclose%
\pgfusepath{fill}%
\end{pgfscope}%
\begin{pgfscope}%
\pgfpathrectangle{\pgfqpoint{0.407592in}{0.222904in}}{\pgfqpoint{1.411163in}{1.411163in}}%
\pgfusepath{clip}%
\pgfsetbuttcap%
\pgfsetroundjoin%
\definecolor{currentfill}{rgb}{0.749020,0.772549,0.792157}%
\pgfsetfillcolor{currentfill}%
\pgfsetfillopacity{0.160000}%
\pgfsetlinewidth{0.000000pt}%
\definecolor{currentstroke}{rgb}{0.000000,0.000000,0.000000}%
\pgfsetstrokecolor{currentstroke}%
\pgfsetdash{}{0pt}%
\pgfpathmoveto{\pgfqpoint{1.267590in}{0.748853in}}%
\pgfpathlineto{\pgfqpoint{1.280467in}{0.758895in}}%
\pgfpathlineto{\pgfqpoint{1.293372in}{0.768960in}}%
\pgfpathlineto{\pgfqpoint{1.297021in}{0.774735in}}%
\pgfpathlineto{\pgfqpoint{1.299747in}{0.780571in}}%
\pgfpathlineto{\pgfqpoint{1.286479in}{0.769803in}}%
\pgfpathlineto{\pgfqpoint{1.273234in}{0.759054in}}%
\pgfpathlineto{\pgfqpoint{1.270814in}{0.753925in}}%
\pgfpathlineto{\pgfqpoint{1.267590in}{0.748853in}}%
\pgfpathclose%
\pgfusepath{fill}%
\end{pgfscope}%
\begin{pgfscope}%
\pgfpathrectangle{\pgfqpoint{0.407592in}{0.222904in}}{\pgfqpoint{1.411163in}{1.411163in}}%
\pgfusepath{clip}%
\pgfsetbuttcap%
\pgfsetroundjoin%
\definecolor{currentfill}{rgb}{0.749020,0.772549,0.792157}%
\pgfsetfillcolor{currentfill}%
\pgfsetfillopacity{0.160000}%
\pgfsetlinewidth{0.000000pt}%
\definecolor{currentstroke}{rgb}{0.000000,0.000000,0.000000}%
\pgfsetstrokecolor{currentstroke}%
\pgfsetdash{}{0pt}%
\pgfpathmoveto{\pgfqpoint{0.578613in}{1.123085in}}%
\pgfpathlineto{\pgfqpoint{0.565688in}{1.136033in}}%
\pgfpathlineto{\pgfqpoint{0.552757in}{1.148988in}}%
\pgfpathlineto{\pgfqpoint{0.541067in}{1.135422in}}%
\pgfpathlineto{\pgfqpoint{0.531649in}{1.121571in}}%
\pgfpathlineto{\pgfqpoint{0.545207in}{1.109167in}}%
\pgfpathlineto{\pgfqpoint{0.558751in}{1.096775in}}%
\pgfpathlineto{\pgfqpoint{0.567598in}{1.110063in}}%
\pgfpathlineto{\pgfqpoint{0.578613in}{1.123085in}}%
\pgfpathclose%
\pgfusepath{fill}%
\end{pgfscope}%
\begin{pgfscope}%
\pgfpathrectangle{\pgfqpoint{0.407592in}{0.222904in}}{\pgfqpoint{1.411163in}{1.411163in}}%
\pgfusepath{clip}%
\pgfsetbuttcap%
\pgfsetroundjoin%
\definecolor{currentfill}{rgb}{0.749020,0.772549,0.792157}%
\pgfsetfillcolor{currentfill}%
\pgfsetfillopacity{0.160000}%
\pgfsetlinewidth{0.000000pt}%
\definecolor{currentstroke}{rgb}{0.000000,0.000000,0.000000}%
\pgfsetstrokecolor{currentstroke}%
\pgfsetdash{}{0pt}%
\pgfpathmoveto{\pgfqpoint{1.049350in}{0.648138in}}%
\pgfpathlineto{\pgfqpoint{1.042437in}{0.654498in}}%
\pgfpathlineto{\pgfqpoint{1.035494in}{0.660885in}}%
\pgfpathlineto{\pgfqpoint{1.041689in}{0.659531in}}%
\pgfpathlineto{\pgfqpoint{1.048102in}{0.658362in}}%
\pgfpathlineto{\pgfqpoint{1.053449in}{0.652292in}}%
\pgfpathlineto{\pgfqpoint{1.058772in}{0.646248in}}%
\pgfpathlineto{\pgfqpoint{1.053980in}{0.647124in}}%
\pgfpathlineto{\pgfqpoint{1.049350in}{0.648138in}}%
\pgfpathclose%
\pgfusepath{fill}%
\end{pgfscope}%
\begin{pgfscope}%
\pgfpathrectangle{\pgfqpoint{0.407592in}{0.222904in}}{\pgfqpoint{1.411163in}{1.411163in}}%
\pgfusepath{clip}%
\pgfsetbuttcap%
\pgfsetroundjoin%
\definecolor{currentfill}{rgb}{0.749020,0.772549,0.792157}%
\pgfsetfillcolor{currentfill}%
\pgfsetfillopacity{0.160000}%
\pgfsetlinewidth{0.000000pt}%
\definecolor{currentstroke}{rgb}{0.000000,0.000000,0.000000}%
\pgfsetstrokecolor{currentstroke}%
\pgfsetdash{}{0pt}%
\pgfpathmoveto{\pgfqpoint{1.165807in}{0.670619in}}%
\pgfpathlineto{\pgfqpoint{1.175404in}{0.678249in}}%
\pgfpathlineto{\pgfqpoint{1.185034in}{0.685906in}}%
\pgfpathlineto{\pgfqpoint{1.190983in}{0.688809in}}%
\pgfpathlineto{\pgfqpoint{1.196476in}{0.691865in}}%
\pgfpathlineto{\pgfqpoint{1.185700in}{0.683615in}}%
\pgfpathlineto{\pgfqpoint{1.174958in}{0.675392in}}%
\pgfpathlineto{\pgfqpoint{1.170563in}{0.672945in}}%
\pgfpathlineto{\pgfqpoint{1.165807in}{0.670619in}}%
\pgfpathclose%
\pgfusepath{fill}%
\end{pgfscope}%
\begin{pgfscope}%
\pgfpathrectangle{\pgfqpoint{0.407592in}{0.222904in}}{\pgfqpoint{1.411163in}{1.411163in}}%
\pgfusepath{clip}%
\pgfsetbuttcap%
\pgfsetroundjoin%
\definecolor{currentfill}{rgb}{0.749020,0.772549,0.792157}%
\pgfsetfillcolor{currentfill}%
\pgfsetfillopacity{0.160000}%
\pgfsetlinewidth{0.000000pt}%
\definecolor{currentstroke}{rgb}{0.000000,0.000000,0.000000}%
\pgfsetstrokecolor{currentstroke}%
\pgfsetdash{}{0pt}%
\pgfpathmoveto{\pgfqpoint{1.437326in}{0.907915in}}%
\pgfpathlineto{\pgfqpoint{1.450917in}{0.919560in}}%
\pgfpathlineto{\pgfqpoint{1.464526in}{0.931220in}}%
\pgfpathlineto{\pgfqpoint{1.463915in}{0.941128in}}%
\pgfpathlineto{\pgfqpoint{1.461680in}{0.950939in}}%
\pgfpathlineto{\pgfqpoint{1.448253in}{0.938635in}}%
\pgfpathlineto{\pgfqpoint{1.434838in}{0.926341in}}%
\pgfpathlineto{\pgfqpoint{1.436836in}{0.917171in}}%
\pgfpathlineto{\pgfqpoint{1.437326in}{0.907915in}}%
\pgfpathclose%
\pgfusepath{fill}%
\end{pgfscope}%
\begin{pgfscope}%
\pgfpathrectangle{\pgfqpoint{0.407592in}{0.222904in}}{\pgfqpoint{1.411163in}{1.411163in}}%
\pgfusepath{clip}%
\pgfsetbuttcap%
\pgfsetroundjoin%
\definecolor{currentfill}{rgb}{0.749020,0.772549,0.792157}%
\pgfsetfillcolor{currentfill}%
\pgfsetfillopacity{0.160000}%
\pgfsetlinewidth{0.000000pt}%
\definecolor{currentstroke}{rgb}{0.000000,0.000000,0.000000}%
\pgfsetstrokecolor{currentstroke}%
\pgfsetdash{}{0pt}%
\pgfpathmoveto{\pgfqpoint{0.849835in}{0.813999in}}%
\pgfpathlineto{\pgfqpoint{0.836305in}{0.825450in}}%
\pgfpathlineto{\pgfqpoint{0.822757in}{0.836917in}}%
\pgfpathlineto{\pgfqpoint{0.822383in}{0.829642in}}%
\pgfpathlineto{\pgfqpoint{0.823182in}{0.822378in}}%
\pgfpathlineto{\pgfqpoint{0.836764in}{0.811604in}}%
\pgfpathlineto{\pgfqpoint{0.850322in}{0.800848in}}%
\pgfpathlineto{\pgfqpoint{0.849551in}{0.807417in}}%
\pgfpathlineto{\pgfqpoint{0.849835in}{0.813999in}}%
\pgfpathclose%
\pgfusepath{fill}%
\end{pgfscope}%
\begin{pgfscope}%
\pgfpathrectangle{\pgfqpoint{0.407592in}{0.222904in}}{\pgfqpoint{1.411163in}{1.411163in}}%
\pgfusepath{clip}%
\pgfsetbuttcap%
\pgfsetroundjoin%
\definecolor{currentfill}{rgb}{0.749020,0.772549,0.792157}%
\pgfsetfillcolor{currentfill}%
\pgfsetfillopacity{0.160000}%
\pgfsetlinewidth{0.000000pt}%
\definecolor{currentstroke}{rgb}{0.000000,0.000000,0.000000}%
\pgfsetstrokecolor{currentstroke}%
\pgfsetdash{}{0pt}%
\pgfpathmoveto{\pgfqpoint{0.904309in}{0.758021in}}%
\pgfpathlineto{\pgfqpoint{0.890849in}{0.768699in}}%
\pgfpathlineto{\pgfqpoint{0.877364in}{0.779396in}}%
\pgfpathlineto{\pgfqpoint{0.879032in}{0.773570in}}%
\pgfpathlineto{\pgfqpoint{0.881633in}{0.767811in}}%
\pgfpathlineto{\pgfqpoint{0.894889in}{0.757816in}}%
\pgfpathlineto{\pgfqpoint{0.908116in}{0.747844in}}%
\pgfpathlineto{\pgfqpoint{0.905805in}{0.752902in}}%
\pgfpathlineto{\pgfqpoint{0.904309in}{0.758021in}}%
\pgfpathclose%
\pgfusepath{fill}%
\end{pgfscope}%
\begin{pgfscope}%
\pgfpathrectangle{\pgfqpoint{0.407592in}{0.222904in}}{\pgfqpoint{1.411163in}{1.411163in}}%
\pgfusepath{clip}%
\pgfsetbuttcap%
\pgfsetroundjoin%
\definecolor{currentfill}{rgb}{0.749020,0.772549,0.792157}%
\pgfsetfillcolor{currentfill}%
\pgfsetfillopacity{0.160000}%
\pgfsetlinewidth{0.000000pt}%
\definecolor{currentstroke}{rgb}{0.000000,0.000000,0.000000}%
\pgfsetstrokecolor{currentstroke}%
\pgfsetdash{}{0pt}%
\pgfpathmoveto{\pgfqpoint{1.110627in}{0.645650in}}%
\pgfpathlineto{\pgfqpoint{1.114079in}{0.651593in}}%
\pgfpathlineto{\pgfqpoint{1.117547in}{0.657564in}}%
\pgfpathlineto{\pgfqpoint{1.124325in}{0.658581in}}%
\pgfpathlineto{\pgfqpoint{1.130948in}{0.659788in}}%
\pgfpathlineto{\pgfqpoint{1.125782in}{0.653539in}}%
\pgfpathlineto{\pgfqpoint{1.120638in}{0.647317in}}%
\pgfpathlineto{\pgfqpoint{1.115690in}{0.646412in}}%
\pgfpathlineto{\pgfqpoint{1.110627in}{0.645650in}}%
\pgfpathclose%
\pgfusepath{fill}%
\end{pgfscope}%
\begin{pgfscope}%
\pgfpathrectangle{\pgfqpoint{0.407592in}{0.222904in}}{\pgfqpoint{1.411163in}{1.411163in}}%
\pgfusepath{clip}%
\pgfsetbuttcap%
\pgfsetroundjoin%
\definecolor{currentfill}{rgb}{0.749020,0.772549,0.792157}%
\pgfsetfillcolor{currentfill}%
\pgfsetfillopacity{0.160000}%
\pgfsetlinewidth{0.000000pt}%
\definecolor{currentstroke}{rgb}{0.000000,0.000000,0.000000}%
\pgfsetstrokecolor{currentstroke}%
\pgfsetdash{}{0pt}%
\pgfpathmoveto{\pgfqpoint{1.201490in}{0.695057in}}%
\pgfpathlineto{\pgfqpoint{1.212799in}{0.703649in}}%
\pgfpathlineto{\pgfqpoint{1.224142in}{0.712266in}}%
\pgfpathlineto{\pgfqpoint{1.229547in}{0.716236in}}%
\pgfpathlineto{\pgfqpoint{1.234324in}{0.720333in}}%
\pgfpathlineto{\pgfqpoint{1.222144in}{0.711052in}}%
\pgfpathlineto{\pgfqpoint{1.209996in}{0.701796in}}%
\pgfpathlineto{\pgfqpoint{1.206003in}{0.698373in}}%
\pgfpathlineto{\pgfqpoint{1.201490in}{0.695057in}}%
\pgfpathclose%
\pgfusepath{fill}%
\end{pgfscope}%
\begin{pgfscope}%
\pgfpathrectangle{\pgfqpoint{0.407592in}{0.222904in}}{\pgfqpoint{1.411163in}{1.411163in}}%
\pgfusepath{clip}%
\pgfsetbuttcap%
\pgfsetroundjoin%
\definecolor{currentfill}{rgb}{0.749020,0.772549,0.792157}%
\pgfsetfillcolor{currentfill}%
\pgfsetfillopacity{0.160000}%
\pgfsetlinewidth{0.000000pt}%
\definecolor{currentstroke}{rgb}{0.000000,0.000000,0.000000}%
\pgfsetstrokecolor{currentstroke}%
\pgfsetdash{}{0pt}%
\pgfpathmoveto{\pgfqpoint{1.678158in}{1.149325in}}%
\pgfpathlineto{\pgfqpoint{1.691792in}{1.161819in}}%
\pgfpathlineto{\pgfqpoint{1.705438in}{1.174325in}}%
\pgfpathlineto{\pgfqpoint{1.697893in}{1.189211in}}%
\pgfpathlineto{\pgfqpoint{1.687810in}{1.203766in}}%
\pgfpathlineto{\pgfqpoint{1.674681in}{1.190733in}}%
\pgfpathlineto{\pgfqpoint{1.661559in}{1.177707in}}%
\pgfpathlineto{\pgfqpoint{1.671075in}{1.163672in}}%
\pgfpathlineto{\pgfqpoint{1.678158in}{1.149325in}}%
\pgfpathclose%
\pgfusepath{fill}%
\end{pgfscope}%
\begin{pgfscope}%
\pgfpathrectangle{\pgfqpoint{0.407592in}{0.222904in}}{\pgfqpoint{1.411163in}{1.411163in}}%
\pgfusepath{clip}%
\pgfsetbuttcap%
\pgfsetroundjoin%
\definecolor{currentfill}{rgb}{0.749020,0.772549,0.792157}%
\pgfsetfillcolor{currentfill}%
\pgfsetfillopacity{0.160000}%
\pgfsetlinewidth{0.000000pt}%
\definecolor{currentstroke}{rgb}{0.000000,0.000000,0.000000}%
\pgfsetstrokecolor{currentstroke}%
\pgfsetdash{}{0pt}%
\pgfpathmoveto{\pgfqpoint{1.058772in}{0.646248in}}%
\pgfpathlineto{\pgfqpoint{1.053449in}{0.652292in}}%
\pgfpathlineto{\pgfqpoint{1.048102in}{0.658362in}}%
\pgfpathlineto{\pgfqpoint{1.054705in}{0.657383in}}%
\pgfpathlineto{\pgfqpoint{1.061467in}{0.656598in}}%
\pgfpathlineto{\pgfqpoint{1.065119in}{0.650749in}}%
\pgfpathlineto{\pgfqpoint{1.068755in}{0.644926in}}%
\pgfpathlineto{\pgfqpoint{1.063704in}{0.645514in}}%
\pgfpathlineto{\pgfqpoint{1.058772in}{0.646248in}}%
\pgfpathclose%
\pgfusepath{fill}%
\end{pgfscope}%
\begin{pgfscope}%
\pgfpathrectangle{\pgfqpoint{0.407592in}{0.222904in}}{\pgfqpoint{1.411163in}{1.411163in}}%
\pgfusepath{clip}%
\pgfsetbuttcap%
\pgfsetroundjoin%
\definecolor{currentfill}{rgb}{0.749020,0.772549,0.792157}%
\pgfsetfillcolor{currentfill}%
\pgfsetfillopacity{0.160000}%
\pgfsetlinewidth{0.000000pt}%
\definecolor{currentstroke}{rgb}{0.000000,0.000000,0.000000}%
\pgfsetstrokecolor{currentstroke}%
\pgfsetdash{}{0pt}%
\pgfpathmoveto{\pgfqpoint{1.234324in}{0.720333in}}%
\pgfpathlineto{\pgfqpoint{1.246535in}{0.729638in}}%
\pgfpathlineto{\pgfqpoint{1.258778in}{0.738968in}}%
\pgfpathlineto{\pgfqpoint{1.263573in}{0.743859in}}%
\pgfpathlineto{\pgfqpoint{1.267590in}{0.748853in}}%
\pgfpathlineto{\pgfqpoint{1.254741in}{0.738832in}}%
\pgfpathlineto{\pgfqpoint{1.241920in}{0.728833in}}%
\pgfpathlineto{\pgfqpoint{1.238453in}{0.724539in}}%
\pgfpathlineto{\pgfqpoint{1.234324in}{0.720333in}}%
\pgfpathclose%
\pgfusepath{fill}%
\end{pgfscope}%
\begin{pgfscope}%
\pgfpathrectangle{\pgfqpoint{0.407592in}{0.222904in}}{\pgfqpoint{1.411163in}{1.411163in}}%
\pgfusepath{clip}%
\pgfsetbuttcap%
\pgfsetroundjoin%
\definecolor{currentfill}{rgb}{0.749020,0.772549,0.792157}%
\pgfsetfillcolor{currentfill}%
\pgfsetfillopacity{0.160000}%
\pgfsetlinewidth{0.000000pt}%
\definecolor{currentstroke}{rgb}{0.000000,0.000000,0.000000}%
\pgfsetstrokecolor{currentstroke}%
\pgfsetdash{}{0pt}%
\pgfpathmoveto{\pgfqpoint{0.965832in}{0.701101in}}%
\pgfpathlineto{\pgfqpoint{0.953212in}{0.710289in}}%
\pgfpathlineto{\pgfqpoint{0.940558in}{0.719502in}}%
\pgfpathlineto{\pgfqpoint{0.944605in}{0.715429in}}%
\pgfpathlineto{\pgfqpoint{0.949302in}{0.711487in}}%
\pgfpathlineto{\pgfqpoint{0.961242in}{0.702934in}}%
\pgfpathlineto{\pgfqpoint{0.973145in}{0.694406in}}%
\pgfpathlineto{\pgfqpoint{0.969219in}{0.697698in}}%
\pgfpathlineto{\pgfqpoint{0.965832in}{0.701101in}}%
\pgfpathclose%
\pgfusepath{fill}%
\end{pgfscope}%
\begin{pgfscope}%
\pgfpathrectangle{\pgfqpoint{0.407592in}{0.222904in}}{\pgfqpoint{1.411163in}{1.411163in}}%
\pgfusepath{clip}%
\pgfsetbuttcap%
\pgfsetroundjoin%
\definecolor{currentfill}{rgb}{0.749020,0.772549,0.792157}%
\pgfsetfillcolor{currentfill}%
\pgfsetfillopacity{0.160000}%
\pgfsetlinewidth{0.000000pt}%
\definecolor{currentstroke}{rgb}{0.000000,0.000000,0.000000}%
\pgfsetstrokecolor{currentstroke}%
\pgfsetdash{}{0pt}%
\pgfpathmoveto{\pgfqpoint{0.747043in}{0.924510in}}%
\pgfpathlineto{\pgfqpoint{0.733675in}{0.936740in}}%
\pgfpathlineto{\pgfqpoint{0.720296in}{0.948980in}}%
\pgfpathlineto{\pgfqpoint{0.716191in}{0.939146in}}%
\pgfpathlineto{\pgfqpoint{0.713696in}{0.929224in}}%
\pgfpathlineto{\pgfqpoint{0.727395in}{0.917630in}}%
\pgfpathlineto{\pgfqpoint{0.741074in}{0.906052in}}%
\pgfpathlineto{\pgfqpoint{0.743311in}{0.915320in}}%
\pgfpathlineto{\pgfqpoint{0.747043in}{0.924510in}}%
\pgfpathclose%
\pgfusepath{fill}%
\end{pgfscope}%
\begin{pgfscope}%
\pgfpathrectangle{\pgfqpoint{0.407592in}{0.222904in}}{\pgfqpoint{1.411163in}{1.411163in}}%
\pgfusepath{clip}%
\pgfsetbuttcap%
\pgfsetroundjoin%
\definecolor{currentfill}{rgb}{0.749020,0.772549,0.792157}%
\pgfsetfillcolor{currentfill}%
\pgfsetfillopacity{0.160000}%
\pgfsetlinewidth{0.000000pt}%
\definecolor{currentstroke}{rgb}{0.000000,0.000000,0.000000}%
\pgfsetstrokecolor{currentstroke}%
\pgfsetdash{}{0pt}%
\pgfpathmoveto{\pgfqpoint{0.552757in}{1.148988in}}%
\pgfpathlineto{\pgfqpoint{0.539819in}{1.161950in}}%
\pgfpathlineto{\pgfqpoint{0.526874in}{1.174919in}}%
\pgfpathlineto{\pgfqpoint{0.514497in}{1.160817in}}%
\pgfpathlineto{\pgfqpoint{0.504495in}{1.146414in}}%
\pgfpathlineto{\pgfqpoint{0.518078in}{1.133987in}}%
\pgfpathlineto{\pgfqpoint{0.531649in}{1.121571in}}%
\pgfpathlineto{\pgfqpoint{0.541067in}{1.135422in}}%
\pgfpathlineto{\pgfqpoint{0.552757in}{1.148988in}}%
\pgfpathclose%
\pgfusepath{fill}%
\end{pgfscope}%
\begin{pgfscope}%
\pgfpathrectangle{\pgfqpoint{0.407592in}{0.222904in}}{\pgfqpoint{1.411163in}{1.411163in}}%
\pgfusepath{clip}%
\pgfsetbuttcap%
\pgfsetroundjoin%
\definecolor{currentfill}{rgb}{0.749020,0.772549,0.792157}%
\pgfsetfillcolor{currentfill}%
\pgfsetfillopacity{0.160000}%
\pgfsetlinewidth{0.000000pt}%
\definecolor{currentstroke}{rgb}{0.000000,0.000000,0.000000}%
\pgfsetstrokecolor{currentstroke}%
\pgfsetdash{}{0pt}%
\pgfpathmoveto{\pgfqpoint{1.100245in}{0.644563in}}%
\pgfpathlineto{\pgfqpoint{1.101941in}{0.650324in}}%
\pgfpathlineto{\pgfqpoint{1.103644in}{0.656113in}}%
\pgfpathlineto{\pgfqpoint{1.110643in}{0.656739in}}%
\pgfpathlineto{\pgfqpoint{1.117547in}{0.657564in}}%
\pgfpathlineto{\pgfqpoint{1.114079in}{0.651593in}}%
\pgfpathlineto{\pgfqpoint{1.110627in}{0.645650in}}%
\pgfpathlineto{\pgfqpoint{1.105471in}{0.645032in}}%
\pgfpathlineto{\pgfqpoint{1.100245in}{0.644563in}}%
\pgfpathclose%
\pgfusepath{fill}%
\end{pgfscope}%
\begin{pgfscope}%
\pgfpathrectangle{\pgfqpoint{0.407592in}{0.222904in}}{\pgfqpoint{1.411163in}{1.411163in}}%
\pgfusepath{clip}%
\pgfsetbuttcap%
\pgfsetroundjoin%
\definecolor{currentfill}{rgb}{0.749020,0.772549,0.792157}%
\pgfsetfillcolor{currentfill}%
\pgfsetfillopacity{0.160000}%
\pgfsetlinewidth{0.000000pt}%
\definecolor{currentstroke}{rgb}{0.000000,0.000000,0.000000}%
\pgfsetstrokecolor{currentstroke}%
\pgfsetdash{}{0pt}%
\pgfpathmoveto{\pgfqpoint{0.934480in}{0.727966in}}%
\pgfpathlineto{\pgfqpoint{0.921313in}{0.737894in}}%
\pgfpathlineto{\pgfqpoint{0.908116in}{0.747844in}}%
\pgfpathlineto{\pgfqpoint{0.911234in}{0.742869in}}%
\pgfpathlineto{\pgfqpoint{0.915150in}{0.738000in}}%
\pgfpathlineto{\pgfqpoint{0.927871in}{0.728739in}}%
\pgfpathlineto{\pgfqpoint{0.940558in}{0.719502in}}%
\pgfpathlineto{\pgfqpoint{0.937179in}{0.723687in}}%
\pgfpathlineto{\pgfqpoint{0.934480in}{0.727966in}}%
\pgfpathclose%
\pgfusepath{fill}%
\end{pgfscope}%
\begin{pgfscope}%
\pgfpathrectangle{\pgfqpoint{0.407592in}{0.222904in}}{\pgfqpoint{1.411163in}{1.411163in}}%
\pgfusepath{clip}%
\pgfsetbuttcap%
\pgfsetroundjoin%
\definecolor{currentfill}{rgb}{0.749020,0.772549,0.792157}%
\pgfsetfillcolor{currentfill}%
\pgfsetfillopacity{0.160000}%
\pgfsetlinewidth{0.000000pt}%
\definecolor{currentstroke}{rgb}{0.000000,0.000000,0.000000}%
\pgfsetstrokecolor{currentstroke}%
\pgfsetdash{}{0pt}%
\pgfpathmoveto{\pgfqpoint{1.068755in}{0.644926in}}%
\pgfpathlineto{\pgfqpoint{1.065119in}{0.650749in}}%
\pgfpathlineto{\pgfqpoint{1.061467in}{0.656598in}}%
\pgfpathlineto{\pgfqpoint{1.068358in}{0.656011in}}%
\pgfpathlineto{\pgfqpoint{1.075348in}{0.655625in}}%
\pgfpathlineto{\pgfqpoint{1.077238in}{0.649898in}}%
\pgfpathlineto{\pgfqpoint{1.079119in}{0.644197in}}%
\pgfpathlineto{\pgfqpoint{1.073901in}{0.644486in}}%
\pgfpathlineto{\pgfqpoint{1.068755in}{0.644926in}}%
\pgfpathclose%
\pgfusepath{fill}%
\end{pgfscope}%
\begin{pgfscope}%
\pgfpathrectangle{\pgfqpoint{0.407592in}{0.222904in}}{\pgfqpoint{1.411163in}{1.411163in}}%
\pgfusepath{clip}%
\pgfsetbuttcap%
\pgfsetroundjoin%
\definecolor{currentfill}{rgb}{0.749020,0.772549,0.792157}%
\pgfsetfillcolor{currentfill}%
\pgfsetfillopacity{0.160000}%
\pgfsetlinewidth{0.000000pt}%
\definecolor{currentstroke}{rgb}{0.000000,0.000000,0.000000}%
\pgfsetstrokecolor{currentstroke}%
\pgfsetdash{}{0pt}%
\pgfpathmoveto{\pgfqpoint{1.004374in}{0.672468in}}%
\pgfpathlineto{\pgfqpoint{0.993479in}{0.680328in}}%
\pgfpathlineto{\pgfqpoint{0.982547in}{0.688214in}}%
\pgfpathlineto{\pgfqpoint{0.987984in}{0.685342in}}%
\pgfpathlineto{\pgfqpoint{0.993880in}{0.682638in}}%
\pgfpathlineto{\pgfqpoint{1.003676in}{0.675308in}}%
\pgfpathlineto{\pgfqpoint{1.013435in}{0.668004in}}%
\pgfpathlineto{\pgfqpoint{1.008723in}{0.670169in}}%
\pgfpathlineto{\pgfqpoint{1.004374in}{0.672468in}}%
\pgfpathclose%
\pgfusepath{fill}%
\end{pgfscope}%
\begin{pgfscope}%
\pgfpathrectangle{\pgfqpoint{0.407592in}{0.222904in}}{\pgfqpoint{1.411163in}{1.411163in}}%
\pgfusepath{clip}%
\pgfsetbuttcap%
\pgfsetroundjoin%
\definecolor{currentfill}{rgb}{0.749020,0.772549,0.792157}%
\pgfsetfillcolor{currentfill}%
\pgfsetfillopacity{0.160000}%
\pgfsetlinewidth{0.000000pt}%
\definecolor{currentstroke}{rgb}{0.000000,0.000000,0.000000}%
\pgfsetstrokecolor{currentstroke}%
\pgfsetdash{}{0pt}%
\pgfpathmoveto{\pgfqpoint{1.089679in}{0.644076in}}%
\pgfpathlineto{\pgfqpoint{1.089586in}{0.649756in}}%
\pgfpathlineto{\pgfqpoint{1.089492in}{0.655462in}}%
\pgfpathlineto{\pgfqpoint{1.096584in}{0.655686in}}%
\pgfpathlineto{\pgfqpoint{1.103644in}{0.656113in}}%
\pgfpathlineto{\pgfqpoint{1.101941in}{0.650324in}}%
\pgfpathlineto{\pgfqpoint{1.100245in}{0.644563in}}%
\pgfpathlineto{\pgfqpoint{1.094973in}{0.644243in}}%
\pgfpathlineto{\pgfqpoint{1.089679in}{0.644076in}}%
\pgfpathclose%
\pgfusepath{fill}%
\end{pgfscope}%
\begin{pgfscope}%
\pgfpathrectangle{\pgfqpoint{0.407592in}{0.222904in}}{\pgfqpoint{1.411163in}{1.411163in}}%
\pgfusepath{clip}%
\pgfsetbuttcap%
\pgfsetroundjoin%
\definecolor{currentfill}{rgb}{0.749020,0.772549,0.792157}%
\pgfsetfillcolor{currentfill}%
\pgfsetfillopacity{0.160000}%
\pgfsetlinewidth{0.000000pt}%
\definecolor{currentstroke}{rgb}{0.000000,0.000000,0.000000}%
\pgfsetstrokecolor{currentstroke}%
\pgfsetdash{}{0pt}%
\pgfpathmoveto{\pgfqpoint{1.079119in}{0.644197in}}%
\pgfpathlineto{\pgfqpoint{1.077238in}{0.649898in}}%
\pgfpathlineto{\pgfqpoint{1.075348in}{0.655625in}}%
\pgfpathlineto{\pgfqpoint{1.082403in}{0.655442in}}%
\pgfpathlineto{\pgfqpoint{1.089492in}{0.655462in}}%
\pgfpathlineto{\pgfqpoint{1.089586in}{0.649756in}}%
\pgfpathlineto{\pgfqpoint{1.089679in}{0.644076in}}%
\pgfpathlineto{\pgfqpoint{1.084387in}{0.644060in}}%
\pgfpathlineto{\pgfqpoint{1.079119in}{0.644197in}}%
\pgfpathclose%
\pgfusepath{fill}%
\end{pgfscope}%
\begin{pgfscope}%
\pgfpathrectangle{\pgfqpoint{0.407592in}{0.222904in}}{\pgfqpoint{1.411163in}{1.411163in}}%
\pgfusepath{clip}%
\pgfsetbuttcap%
\pgfsetroundjoin%
\definecolor{currentfill}{rgb}{0.749020,0.772549,0.792157}%
\pgfsetfillcolor{currentfill}%
\pgfsetfillopacity{0.160000}%
\pgfsetlinewidth{0.000000pt}%
\definecolor{currentstroke}{rgb}{0.000000,0.000000,0.000000}%
\pgfsetstrokecolor{currentstroke}%
\pgfsetdash{}{0pt}%
\pgfpathmoveto{\pgfqpoint{1.705438in}{1.174325in}}%
\pgfpathlineto{\pgfqpoint{1.719097in}{1.186842in}}%
\pgfpathlineto{\pgfqpoint{1.732768in}{1.199371in}}%
\pgfpathlineto{\pgfqpoint{1.724748in}{1.214785in}}%
\pgfpathlineto{\pgfqpoint{1.714085in}{1.229851in}}%
\pgfpathlineto{\pgfqpoint{1.700944in}{1.216805in}}%
\pgfpathlineto{\pgfqpoint{1.687810in}{1.203766in}}%
\pgfpathlineto{\pgfqpoint{1.697893in}{1.189211in}}%
\pgfpathlineto{\pgfqpoint{1.705438in}{1.174325in}}%
\pgfpathclose%
\pgfusepath{fill}%
\end{pgfscope}%
\begin{pgfscope}%
\pgfpathrectangle{\pgfqpoint{0.407592in}{0.222904in}}{\pgfqpoint{1.411163in}{1.411163in}}%
\pgfusepath{clip}%
\pgfsetbuttcap%
\pgfsetroundjoin%
\definecolor{currentfill}{rgb}{0.749020,0.772549,0.792157}%
\pgfsetfillcolor{currentfill}%
\pgfsetfillopacity{0.160000}%
\pgfsetlinewidth{0.000000pt}%
\definecolor{currentstroke}{rgb}{0.000000,0.000000,0.000000}%
\pgfsetstrokecolor{currentstroke}%
\pgfsetdash{}{0pt}%
\pgfpathmoveto{\pgfqpoint{1.464526in}{0.931220in}}%
\pgfpathlineto{\pgfqpoint{1.478153in}{0.942895in}}%
\pgfpathlineto{\pgfqpoint{1.491799in}{0.954586in}}%
\pgfpathlineto{\pgfqpoint{1.491055in}{0.965138in}}%
\pgfpathlineto{\pgfqpoint{1.488569in}{0.975581in}}%
\pgfpathlineto{\pgfqpoint{1.475119in}{0.963255in}}%
\pgfpathlineto{\pgfqpoint{1.461680in}{0.950939in}}%
\pgfpathlineto{\pgfqpoint{1.463915in}{0.941128in}}%
\pgfpathlineto{\pgfqpoint{1.464526in}{0.931220in}}%
\pgfpathclose%
\pgfusepath{fill}%
\end{pgfscope}%
\begin{pgfscope}%
\pgfpathrectangle{\pgfqpoint{0.407592in}{0.222904in}}{\pgfqpoint{1.411163in}{1.411163in}}%
\pgfusepath{clip}%
\pgfsetbuttcap%
\pgfsetroundjoin%
\definecolor{currentfill}{rgb}{0.749020,0.772549,0.792157}%
\pgfsetfillcolor{currentfill}%
\pgfsetfillopacity{0.160000}%
\pgfsetlinewidth{0.000000pt}%
\definecolor{currentstroke}{rgb}{0.000000,0.000000,0.000000}%
\pgfsetstrokecolor{currentstroke}%
\pgfsetdash{}{0pt}%
\pgfpathmoveto{\pgfqpoint{1.353055in}{0.823833in}}%
\pgfpathlineto{\pgfqpoint{1.366440in}{0.834697in}}%
\pgfpathlineto{\pgfqpoint{1.379850in}{0.845580in}}%
\pgfpathlineto{\pgfqpoint{1.382142in}{0.853535in}}%
\pgfpathlineto{\pgfqpoint{1.383143in}{0.861493in}}%
\pgfpathlineto{\pgfqpoint{1.369643in}{0.849927in}}%
\pgfpathlineto{\pgfqpoint{1.356160in}{0.838375in}}%
\pgfpathlineto{\pgfqpoint{1.355193in}{0.831102in}}%
\pgfpathlineto{\pgfqpoint{1.353055in}{0.823833in}}%
\pgfpathclose%
\pgfusepath{fill}%
\end{pgfscope}%
\begin{pgfscope}%
\pgfpathrectangle{\pgfqpoint{0.407592in}{0.222904in}}{\pgfqpoint{1.411163in}{1.411163in}}%
\pgfusepath{clip}%
\pgfsetbuttcap%
\pgfsetroundjoin%
\definecolor{currentfill}{rgb}{0.749020,0.772549,0.792157}%
\pgfsetfillcolor{currentfill}%
\pgfsetfillopacity{0.160000}%
\pgfsetlinewidth{0.000000pt}%
\definecolor{currentstroke}{rgb}{0.000000,0.000000,0.000000}%
\pgfsetstrokecolor{currentstroke}%
\pgfsetdash{}{0pt}%
\pgfpathmoveto{\pgfqpoint{0.526874in}{1.174919in}}%
\pgfpathlineto{\pgfqpoint{0.513923in}{1.187894in}}%
\pgfpathlineto{\pgfqpoint{0.500965in}{1.200875in}}%
\pgfpathlineto{\pgfqpoint{0.487888in}{1.186249in}}%
\pgfpathlineto{\pgfqpoint{0.477289in}{1.171304in}}%
\pgfpathlineto{\pgfqpoint{0.490898in}{1.158853in}}%
\pgfpathlineto{\pgfqpoint{0.504495in}{1.146414in}}%
\pgfpathlineto{\pgfqpoint{0.514497in}{1.160817in}}%
\pgfpathlineto{\pgfqpoint{0.526874in}{1.174919in}}%
\pgfpathclose%
\pgfusepath{fill}%
\end{pgfscope}%
\begin{pgfscope}%
\pgfpathrectangle{\pgfqpoint{0.407592in}{0.222904in}}{\pgfqpoint{1.411163in}{1.411163in}}%
\pgfusepath{clip}%
\pgfsetbuttcap%
\pgfsetroundjoin%
\definecolor{currentfill}{rgb}{0.749020,0.772549,0.792157}%
\pgfsetfillcolor{currentfill}%
\pgfsetfillopacity{0.160000}%
\pgfsetlinewidth{0.000000pt}%
\definecolor{currentstroke}{rgb}{0.000000,0.000000,0.000000}%
\pgfsetstrokecolor{currentstroke}%
\pgfsetdash{}{0pt}%
\pgfpathmoveto{\pgfqpoint{1.155295in}{0.666377in}}%
\pgfpathlineto{\pgfqpoint{1.163568in}{0.673478in}}%
\pgfpathlineto{\pgfqpoint{1.171873in}{0.680606in}}%
\pgfpathlineto{\pgfqpoint{1.178654in}{0.683167in}}%
\pgfpathlineto{\pgfqpoint{1.185034in}{0.685906in}}%
\pgfpathlineto{\pgfqpoint{1.175404in}{0.678249in}}%
\pgfpathlineto{\pgfqpoint{1.165807in}{0.670619in}}%
\pgfpathlineto{\pgfqpoint{1.160710in}{0.668427in}}%
\pgfpathlineto{\pgfqpoint{1.155295in}{0.666377in}}%
\pgfpathclose%
\pgfusepath{fill}%
\end{pgfscope}%
\begin{pgfscope}%
\pgfpathrectangle{\pgfqpoint{0.407592in}{0.222904in}}{\pgfqpoint{1.411163in}{1.411163in}}%
\pgfusepath{clip}%
\pgfsetbuttcap%
\pgfsetroundjoin%
\definecolor{currentfill}{rgb}{0.749020,0.772549,0.792157}%
\pgfsetfillcolor{currentfill}%
\pgfsetfillopacity{0.160000}%
\pgfsetlinewidth{0.000000pt}%
\definecolor{currentstroke}{rgb}{0.000000,0.000000,0.000000}%
\pgfsetstrokecolor{currentstroke}%
\pgfsetdash{}{0pt}%
\pgfpathmoveto{\pgfqpoint{0.822757in}{0.836917in}}%
\pgfpathlineto{\pgfqpoint{0.809190in}{0.848400in}}%
\pgfpathlineto{\pgfqpoint{0.795604in}{0.859898in}}%
\pgfpathlineto{\pgfqpoint{0.795127in}{0.851938in}}%
\pgfpathlineto{\pgfqpoint{0.795943in}{0.843987in}}%
\pgfpathlineto{\pgfqpoint{0.809575in}{0.833173in}}%
\pgfpathlineto{\pgfqpoint{0.823182in}{0.822378in}}%
\pgfpathlineto{\pgfqpoint{0.822383in}{0.829642in}}%
\pgfpathlineto{\pgfqpoint{0.822757in}{0.836917in}}%
\pgfpathclose%
\pgfusepath{fill}%
\end{pgfscope}%
\begin{pgfscope}%
\pgfpathrectangle{\pgfqpoint{0.407592in}{0.222904in}}{\pgfqpoint{1.411163in}{1.411163in}}%
\pgfusepath{clip}%
\pgfsetbuttcap%
\pgfsetroundjoin%
\definecolor{currentfill}{rgb}{0.749020,0.772549,0.792157}%
\pgfsetfillcolor{currentfill}%
\pgfsetfillopacity{0.160000}%
\pgfsetlinewidth{0.000000pt}%
\definecolor{currentstroke}{rgb}{0.000000,0.000000,0.000000}%
\pgfsetstrokecolor{currentstroke}%
\pgfsetdash{}{0pt}%
\pgfpathmoveto{\pgfqpoint{1.732768in}{1.199371in}}%
\pgfpathlineto{\pgfqpoint{1.746451in}{1.211910in}}%
\pgfpathlineto{\pgfqpoint{1.760147in}{1.224461in}}%
\pgfpathlineto{\pgfqpoint{1.751640in}{1.240394in}}%
\pgfpathlineto{\pgfqpoint{1.740385in}{1.255959in}}%
\pgfpathlineto{\pgfqpoint{1.727232in}{1.242902in}}%
\pgfpathlineto{\pgfqpoint{1.714085in}{1.229851in}}%
\pgfpathlineto{\pgfqpoint{1.724748in}{1.214785in}}%
\pgfpathlineto{\pgfqpoint{1.732768in}{1.199371in}}%
\pgfpathclose%
\pgfusepath{fill}%
\end{pgfscope}%
\begin{pgfscope}%
\pgfpathrectangle{\pgfqpoint{0.407592in}{0.222904in}}{\pgfqpoint{1.411163in}{1.411163in}}%
\pgfusepath{clip}%
\pgfsetbuttcap%
\pgfsetroundjoin%
\definecolor{currentfill}{rgb}{0.749020,0.772549,0.792157}%
\pgfsetfillcolor{currentfill}%
\pgfsetfillopacity{0.160000}%
\pgfsetlinewidth{0.000000pt}%
\definecolor{currentstroke}{rgb}{0.000000,0.000000,0.000000}%
\pgfsetstrokecolor{currentstroke}%
\pgfsetdash{}{0pt}%
\pgfpathmoveto{\pgfqpoint{0.720296in}{0.948980in}}%
\pgfpathlineto{\pgfqpoint{0.706903in}{0.961233in}}%
\pgfpathlineto{\pgfqpoint{0.693499in}{0.973497in}}%
\pgfpathlineto{\pgfqpoint{0.689009in}{0.963028in}}%
\pgfpathlineto{\pgfqpoint{0.686242in}{0.952460in}}%
\pgfpathlineto{\pgfqpoint{0.699979in}{0.940834in}}%
\pgfpathlineto{\pgfqpoint{0.713696in}{0.929224in}}%
\pgfpathlineto{\pgfqpoint{0.716191in}{0.939146in}}%
\pgfpathlineto{\pgfqpoint{0.720296in}{0.948980in}}%
\pgfpathclose%
\pgfusepath{fill}%
\end{pgfscope}%
\begin{pgfscope}%
\pgfpathrectangle{\pgfqpoint{0.407592in}{0.222904in}}{\pgfqpoint{1.411163in}{1.411163in}}%
\pgfusepath{clip}%
\pgfsetbuttcap%
\pgfsetroundjoin%
\definecolor{currentfill}{rgb}{0.749020,0.772549,0.792157}%
\pgfsetfillcolor{currentfill}%
\pgfsetfillopacity{0.160000}%
\pgfsetlinewidth{0.000000pt}%
\definecolor{currentstroke}{rgb}{0.000000,0.000000,0.000000}%
\pgfsetstrokecolor{currentstroke}%
\pgfsetdash{}{0pt}%
\pgfpathmoveto{\pgfqpoint{1.293372in}{0.768960in}}%
\pgfpathlineto{\pgfqpoint{1.306306in}{0.779047in}}%
\pgfpathlineto{\pgfqpoint{1.319268in}{0.789156in}}%
\pgfpathlineto{\pgfqpoint{1.323332in}{0.795627in}}%
\pgfpathlineto{\pgfqpoint{1.326354in}{0.802164in}}%
\pgfpathlineto{\pgfqpoint{1.313039in}{0.791358in}}%
\pgfpathlineto{\pgfqpoint{1.299747in}{0.780571in}}%
\pgfpathlineto{\pgfqpoint{1.297021in}{0.774735in}}%
\pgfpathlineto{\pgfqpoint{1.293372in}{0.768960in}}%
\pgfpathclose%
\pgfusepath{fill}%
\end{pgfscope}%
\begin{pgfscope}%
\pgfpathrectangle{\pgfqpoint{0.407592in}{0.222904in}}{\pgfqpoint{1.411163in}{1.411163in}}%
\pgfusepath{clip}%
\pgfsetbuttcap%
\pgfsetroundjoin%
\definecolor{currentfill}{rgb}{0.749020,0.772549,0.792157}%
\pgfsetfillcolor{currentfill}%
\pgfsetfillopacity{0.160000}%
\pgfsetlinewidth{0.000000pt}%
\definecolor{currentstroke}{rgb}{0.000000,0.000000,0.000000}%
\pgfsetstrokecolor{currentstroke}%
\pgfsetdash{}{0pt}%
\pgfpathmoveto{\pgfqpoint{1.196476in}{0.691865in}}%
\pgfpathlineto{\pgfqpoint{1.207287in}{0.700140in}}%
\pgfpathlineto{\pgfqpoint{1.218131in}{0.708442in}}%
\pgfpathlineto{\pgfqpoint{1.224142in}{0.712266in}}%
\pgfpathlineto{\pgfqpoint{1.212799in}{0.703649in}}%
\pgfpathlineto{\pgfqpoint{1.201490in}{0.695057in}}%
\pgfusepath{fill}%
\end{pgfscope}%
\begin{pgfscope}%
\pgfpathrectangle{\pgfqpoint{0.407592in}{0.222904in}}{\pgfqpoint{1.411163in}{1.411163in}}%
\pgfusepath{clip}%
\pgfsetbuttcap%
\pgfsetroundjoin%
\definecolor{currentfill}{rgb}{0.749020,0.772549,0.792157}%
\pgfsetfillcolor{currentfill}%
\pgfsetfillopacity{0.160000}%
\pgfsetlinewidth{0.000000pt}%
\definecolor{currentstroke}{rgb}{0.000000,0.000000,0.000000}%
\pgfsetstrokecolor{currentstroke}%
\pgfsetdash{}{0pt}%
\pgfpathmoveto{\pgfqpoint{0.500965in}{1.200875in}}%
\pgfpathlineto{\pgfqpoint{0.488000in}{1.213864in}}%
\pgfpathlineto{\pgfqpoint{0.475029in}{1.226859in}}%
\pgfpathlineto{\pgfqpoint{0.461241in}{1.211719in}}%
\pgfpathlineto{\pgfqpoint{0.450032in}{1.196241in}}%
\pgfpathlineto{\pgfqpoint{0.463667in}{1.183767in}}%
\pgfpathlineto{\pgfqpoint{0.477289in}{1.171304in}}%
\pgfpathlineto{\pgfqpoint{0.487888in}{1.186249in}}%
\pgfpathlineto{\pgfqpoint{0.500965in}{1.200875in}}%
\pgfpathclose%
\pgfusepath{fill}%
\end{pgfscope}%
\begin{pgfscope}%
\pgfpathrectangle{\pgfqpoint{0.407592in}{0.222904in}}{\pgfqpoint{1.411163in}{1.411163in}}%
\pgfusepath{clip}%
\pgfsetbuttcap%
\pgfsetroundjoin%
\definecolor{currentfill}{rgb}{0.749020,0.772549,0.792157}%
\pgfsetfillcolor{currentfill}%
\pgfsetfillopacity{0.160000}%
\pgfsetlinewidth{0.000000pt}%
\definecolor{currentstroke}{rgb}{0.000000,0.000000,0.000000}%
\pgfsetstrokecolor{currentstroke}%
\pgfsetdash{}{0pt}%
\pgfpathmoveto{\pgfqpoint{0.877364in}{0.779396in}}%
\pgfpathlineto{\pgfqpoint{0.863855in}{0.790113in}}%
\pgfpathlineto{\pgfqpoint{0.850322in}{0.800848in}}%
\pgfpathlineto{\pgfqpoint{0.852150in}{0.794322in}}%
\pgfpathlineto{\pgfqpoint{0.855032in}{0.787868in}}%
\pgfpathlineto{\pgfqpoint{0.868348in}{0.777829in}}%
\pgfpathlineto{\pgfqpoint{0.881633in}{0.767811in}}%
\pgfpathlineto{\pgfqpoint{0.879032in}{0.773570in}}%
\pgfpathlineto{\pgfqpoint{0.877364in}{0.779396in}}%
\pgfpathclose%
\pgfusepath{fill}%
\end{pgfscope}%
\begin{pgfscope}%
\pgfpathrectangle{\pgfqpoint{0.407592in}{0.222904in}}{\pgfqpoint{1.411163in}{1.411163in}}%
\pgfusepath{clip}%
\pgfsetbuttcap%
\pgfsetroundjoin%
\definecolor{currentfill}{rgb}{0.749020,0.772549,0.792157}%
\pgfsetfillcolor{currentfill}%
\pgfsetfillopacity{0.160000}%
\pgfsetlinewidth{0.000000pt}%
\definecolor{currentstroke}{rgb}{0.000000,0.000000,0.000000}%
\pgfsetstrokecolor{currentstroke}%
\pgfsetdash{}{0pt}%
\pgfpathmoveto{\pgfqpoint{1.013435in}{0.668004in}}%
\pgfpathlineto{\pgfqpoint{1.003676in}{0.675308in}}%
\pgfpathlineto{\pgfqpoint{0.993880in}{0.682638in}}%
\pgfpathlineto{\pgfqpoint{1.000211in}{0.680115in}}%
\pgfpathlineto{\pgfqpoint{1.006948in}{0.677783in}}%
\pgfpathlineto{\pgfqpoint{1.015426in}{0.670937in}}%
\pgfpathlineto{\pgfqpoint{1.023870in}{0.664119in}}%
\pgfpathlineto{\pgfqpoint{1.018491in}{0.665984in}}%
\pgfpathlineto{\pgfqpoint{1.013435in}{0.668004in}}%
\pgfpathclose%
\pgfusepath{fill}%
\end{pgfscope}%
\begin{pgfscope}%
\pgfpathrectangle{\pgfqpoint{0.407592in}{0.222904in}}{\pgfqpoint{1.411163in}{1.411163in}}%
\pgfusepath{clip}%
\pgfsetbuttcap%
\pgfsetroundjoin%
\definecolor{currentfill}{rgb}{0.749020,0.772549,0.792157}%
\pgfsetfillcolor{currentfill}%
\pgfsetfillopacity{0.160000}%
\pgfsetlinewidth{0.000000pt}%
\definecolor{currentstroke}{rgb}{0.000000,0.000000,0.000000}%
\pgfsetstrokecolor{currentstroke}%
\pgfsetdash{}{0pt}%
\pgfpathmoveto{\pgfqpoint{1.491799in}{0.954586in}}%
\pgfpathlineto{\pgfqpoint{1.505463in}{0.966293in}}%
\pgfpathlineto{\pgfqpoint{1.519145in}{0.978016in}}%
\pgfpathlineto{\pgfqpoint{1.518255in}{0.989201in}}%
\pgfpathlineto{\pgfqpoint{1.515507in}{1.000268in}}%
\pgfpathlineto{\pgfqpoint{1.502032in}{0.987919in}}%
\pgfpathlineto{\pgfqpoint{1.488569in}{0.975581in}}%
\pgfpathlineto{\pgfqpoint{1.491055in}{0.965138in}}%
\pgfpathlineto{\pgfqpoint{1.491799in}{0.954586in}}%
\pgfpathclose%
\pgfusepath{fill}%
\end{pgfscope}%
\begin{pgfscope}%
\pgfpathrectangle{\pgfqpoint{0.407592in}{0.222904in}}{\pgfqpoint{1.411163in}{1.411163in}}%
\pgfusepath{clip}%
\pgfsetbuttcap%
\pgfsetroundjoin%
\definecolor{currentfill}{rgb}{0.749020,0.772549,0.792157}%
\pgfsetfillcolor{currentfill}%
\pgfsetfillopacity{0.160000}%
\pgfsetlinewidth{0.000000pt}%
\definecolor{currentstroke}{rgb}{0.000000,0.000000,0.000000}%
\pgfsetstrokecolor{currentstroke}%
\pgfsetdash{}{0pt}%
\pgfpathmoveto{\pgfqpoint{1.760147in}{1.224461in}}%
\pgfpathlineto{\pgfqpoint{1.773855in}{1.237024in}}%
\pgfpathlineto{\pgfqpoint{1.787575in}{1.249598in}}%
\pgfpathlineto{\pgfqpoint{1.778569in}{1.266038in}}%
\pgfpathlineto{\pgfqpoint{1.766710in}{1.282093in}}%
\pgfpathlineto{\pgfqpoint{1.753545in}{1.269023in}}%
\pgfpathlineto{\pgfqpoint{1.740385in}{1.255959in}}%
\pgfpathlineto{\pgfqpoint{1.751640in}{1.240394in}}%
\pgfpathlineto{\pgfqpoint{1.760147in}{1.224461in}}%
\pgfpathclose%
\pgfusepath{fill}%
\end{pgfscope}%
\begin{pgfscope}%
\pgfpathrectangle{\pgfqpoint{0.407592in}{0.222904in}}{\pgfqpoint{1.411163in}{1.411163in}}%
\pgfusepath{clip}%
\pgfsetbuttcap%
\pgfsetroundjoin%
\definecolor{currentfill}{rgb}{0.749020,0.772549,0.792157}%
\pgfsetfillcolor{currentfill}%
\pgfsetfillopacity{0.160000}%
\pgfsetlinewidth{0.000000pt}%
\definecolor{currentstroke}{rgb}{0.000000,0.000000,0.000000}%
\pgfsetstrokecolor{currentstroke}%
\pgfsetdash{}{0pt}%
\pgfpathmoveto{\pgfqpoint{0.973145in}{0.694406in}}%
\pgfpathlineto{\pgfqpoint{0.961242in}{0.702934in}}%
\pgfpathlineto{\pgfqpoint{0.949302in}{0.711487in}}%
\pgfpathlineto{\pgfqpoint{0.954632in}{0.707694in}}%
\pgfpathlineto{\pgfqpoint{0.960572in}{0.704067in}}%
\pgfpathlineto{\pgfqpoint{0.971578in}{0.696127in}}%
\pgfpathlineto{\pgfqpoint{0.982547in}{0.688214in}}%
\pgfpathlineto{\pgfqpoint{0.977594in}{0.691240in}}%
\pgfpathlineto{\pgfqpoint{0.973145in}{0.694406in}}%
\pgfpathclose%
\pgfusepath{fill}%
\end{pgfscope}%
\begin{pgfscope}%
\pgfpathrectangle{\pgfqpoint{0.407592in}{0.222904in}}{\pgfqpoint{1.411163in}{1.411163in}}%
\pgfusepath{clip}%
\pgfsetbuttcap%
\pgfsetroundjoin%
\definecolor{currentfill}{rgb}{0.749020,0.772549,0.792157}%
\pgfsetfillcolor{currentfill}%
\pgfsetfillopacity{0.160000}%
\pgfsetlinewidth{0.000000pt}%
\definecolor{currentstroke}{rgb}{0.000000,0.000000,0.000000}%
\pgfsetstrokecolor{currentstroke}%
\pgfsetdash{}{0pt}%
\pgfpathmoveto{\pgfqpoint{1.143606in}{0.662744in}}%
\pgfpathlineto{\pgfqpoint{1.150401in}{0.669391in}}%
\pgfpathlineto{\pgfqpoint{1.157223in}{0.676065in}}%
\pgfpathlineto{\pgfqpoint{1.164718in}{0.678235in}}%
\pgfpathlineto{\pgfqpoint{1.171873in}{0.680606in}}%
\pgfpathlineto{\pgfqpoint{1.163568in}{0.673478in}}%
\pgfpathlineto{\pgfqpoint{1.155295in}{0.666377in}}%
\pgfpathlineto{\pgfqpoint{1.149585in}{0.664480in}}%
\pgfpathlineto{\pgfqpoint{1.143606in}{0.662744in}}%
\pgfpathclose%
\pgfusepath{fill}%
\end{pgfscope}%
\begin{pgfscope}%
\pgfpathrectangle{\pgfqpoint{0.407592in}{0.222904in}}{\pgfqpoint{1.411163in}{1.411163in}}%
\pgfusepath{clip}%
\pgfsetbuttcap%
\pgfsetroundjoin%
\definecolor{currentfill}{rgb}{0.749020,0.772549,0.792157}%
\pgfsetfillcolor{currentfill}%
\pgfsetfillopacity{0.160000}%
\pgfsetlinewidth{0.000000pt}%
\definecolor{currentstroke}{rgb}{0.000000,0.000000,0.000000}%
\pgfsetstrokecolor{currentstroke}%
\pgfsetdash{}{0pt}%
\pgfpathmoveto{\pgfqpoint{1.379850in}{0.845580in}}%
\pgfpathlineto{\pgfqpoint{1.393284in}{0.856482in}}%
\pgfpathlineto{\pgfqpoint{1.406741in}{0.867403in}}%
\pgfpathlineto{\pgfqpoint{1.409175in}{0.876039in}}%
\pgfpathlineto{\pgfqpoint{1.410198in}{0.884673in}}%
\pgfpathlineto{\pgfqpoint{1.396662in}{0.873076in}}%
\pgfpathlineto{\pgfqpoint{1.383143in}{0.861493in}}%
\pgfpathlineto{\pgfqpoint{1.382142in}{0.853535in}}%
\pgfpathlineto{\pgfqpoint{1.379850in}{0.845580in}}%
\pgfpathclose%
\pgfusepath{fill}%
\end{pgfscope}%
\begin{pgfscope}%
\pgfpathrectangle{\pgfqpoint{0.407592in}{0.222904in}}{\pgfqpoint{1.411163in}{1.411163in}}%
\pgfusepath{clip}%
\pgfsetbuttcap%
\pgfsetroundjoin%
\definecolor{currentfill}{rgb}{0.749020,0.772549,0.792157}%
\pgfsetfillcolor{currentfill}%
\pgfsetfillopacity{0.160000}%
\pgfsetlinewidth{0.000000pt}%
\definecolor{currentstroke}{rgb}{0.000000,0.000000,0.000000}%
\pgfsetstrokecolor{currentstroke}%
\pgfsetdash{}{0pt}%
\pgfpathmoveto{\pgfqpoint{0.475029in}{1.226859in}}%
\pgfpathlineto{\pgfqpoint{0.462051in}{1.239861in}}%
\pgfpathlineto{\pgfqpoint{0.449066in}{1.252869in}}%
\pgfpathlineto{\pgfqpoint{0.434555in}{1.237225in}}%
\pgfpathlineto{\pgfqpoint{0.422723in}{1.221226in}}%
\pgfpathlineto{\pgfqpoint{0.436384in}{1.208728in}}%
\pgfpathlineto{\pgfqpoint{0.450032in}{1.196241in}}%
\pgfpathlineto{\pgfqpoint{0.461241in}{1.211719in}}%
\pgfpathlineto{\pgfqpoint{0.475029in}{1.226859in}}%
\pgfpathclose%
\pgfusepath{fill}%
\end{pgfscope}%
\begin{pgfscope}%
\pgfpathrectangle{\pgfqpoint{0.407592in}{0.222904in}}{\pgfqpoint{1.411163in}{1.411163in}}%
\pgfusepath{clip}%
\pgfsetbuttcap%
\pgfsetroundjoin%
\definecolor{currentfill}{rgb}{0.749020,0.772549,0.792157}%
\pgfsetfillcolor{currentfill}%
\pgfsetfillopacity{0.160000}%
\pgfsetlinewidth{0.000000pt}%
\definecolor{currentstroke}{rgb}{0.000000,0.000000,0.000000}%
\pgfsetstrokecolor{currentstroke}%
\pgfsetdash{}{0pt}%
\pgfpathmoveto{\pgfqpoint{1.258778in}{0.738968in}}%
\pgfpathlineto{\pgfqpoint{1.271053in}{0.748321in}}%
\pgfpathlineto{\pgfqpoint{1.283360in}{0.757700in}}%
\pgfpathlineto{\pgfqpoint{1.288814in}{0.763273in}}%
\pgfpathlineto{\pgfqpoint{1.293372in}{0.768960in}}%
\pgfpathlineto{\pgfqpoint{1.280467in}{0.758895in}}%
\pgfpathlineto{\pgfqpoint{1.267590in}{0.748853in}}%
\pgfpathlineto{\pgfqpoint{1.263573in}{0.743859in}}%
\pgfpathlineto{\pgfqpoint{1.258778in}{0.738968in}}%
\pgfpathclose%
\pgfusepath{fill}%
\end{pgfscope}%
\begin{pgfscope}%
\pgfpathrectangle{\pgfqpoint{0.407592in}{0.222904in}}{\pgfqpoint{1.411163in}{1.411163in}}%
\pgfusepath{clip}%
\pgfsetbuttcap%
\pgfsetroundjoin%
\definecolor{currentfill}{rgb}{0.749020,0.772549,0.792157}%
\pgfsetfillcolor{currentfill}%
\pgfsetfillopacity{0.160000}%
\pgfsetlinewidth{0.000000pt}%
\definecolor{currentstroke}{rgb}{0.000000,0.000000,0.000000}%
\pgfsetstrokecolor{currentstroke}%
\pgfsetdash{}{0pt}%
\pgfpathmoveto{\pgfqpoint{0.693499in}{0.973497in}}%
\pgfpathlineto{\pgfqpoint{0.680081in}{0.985772in}}%
\pgfpathlineto{\pgfqpoint{0.666651in}{0.998060in}}%
\pgfpathlineto{\pgfqpoint{0.661762in}{0.986965in}}%
\pgfpathlineto{\pgfqpoint{0.658712in}{0.975761in}}%
\pgfpathlineto{\pgfqpoint{0.672487in}{0.964103in}}%
\pgfpathlineto{\pgfqpoint{0.686242in}{0.952460in}}%
\pgfpathlineto{\pgfqpoint{0.689009in}{0.963028in}}%
\pgfpathlineto{\pgfqpoint{0.693499in}{0.973497in}}%
\pgfpathclose%
\pgfusepath{fill}%
\end{pgfscope}%
\begin{pgfscope}%
\pgfpathrectangle{\pgfqpoint{0.407592in}{0.222904in}}{\pgfqpoint{1.411163in}{1.411163in}}%
\pgfusepath{clip}%
\pgfsetbuttcap%
\pgfsetroundjoin%
\definecolor{currentfill}{rgb}{0.749020,0.772549,0.792157}%
\pgfsetfillcolor{currentfill}%
\pgfsetfillopacity{0.160000}%
\pgfsetlinewidth{0.000000pt}%
\definecolor{currentstroke}{rgb}{0.000000,0.000000,0.000000}%
\pgfsetstrokecolor{currentstroke}%
\pgfsetdash{}{0pt}%
\pgfpathmoveto{\pgfqpoint{1.224142in}{0.712266in}}%
\pgfpathlineto{\pgfqpoint{1.235519in}{0.720909in}}%
\pgfpathlineto{\pgfqpoint{1.246930in}{0.729579in}}%
\pgfpathlineto{\pgfqpoint{1.253223in}{0.734200in}}%
\pgfpathlineto{\pgfqpoint{1.258778in}{0.738968in}}%
\pgfpathlineto{\pgfqpoint{1.246535in}{0.729638in}}%
\pgfpathlineto{\pgfqpoint{1.234324in}{0.720333in}}%
\pgfpathlineto{\pgfqpoint{1.229547in}{0.716236in}}%
\pgfpathlineto{\pgfqpoint{1.224142in}{0.712266in}}%
\pgfpathclose%
\pgfusepath{fill}%
\end{pgfscope}%
\begin{pgfscope}%
\pgfpathrectangle{\pgfqpoint{0.407592in}{0.222904in}}{\pgfqpoint{1.411163in}{1.411163in}}%
\pgfusepath{clip}%
\pgfsetbuttcap%
\pgfsetroundjoin%
\definecolor{currentfill}{rgb}{0.749020,0.772549,0.792157}%
\pgfsetfillcolor{currentfill}%
\pgfsetfillopacity{0.160000}%
\pgfsetlinewidth{0.000000pt}%
\definecolor{currentstroke}{rgb}{0.000000,0.000000,0.000000}%
\pgfsetstrokecolor{currentstroke}%
\pgfsetdash{}{0pt}%
\pgfpathmoveto{\pgfqpoint{1.787575in}{1.249598in}}%
\pgfpathlineto{\pgfqpoint{1.801308in}{1.262183in}}%
\pgfpathlineto{\pgfqpoint{1.815054in}{1.274780in}}%
\pgfpathlineto{\pgfqpoint{1.805534in}{1.291718in}}%
\pgfpathlineto{\pgfqpoint{1.793060in}{1.308251in}}%
\pgfpathlineto{\pgfqpoint{1.779882in}{1.295169in}}%
\pgfpathlineto{\pgfqpoint{1.766710in}{1.282093in}}%
\pgfpathlineto{\pgfqpoint{1.778569in}{1.266038in}}%
\pgfpathlineto{\pgfqpoint{1.787575in}{1.249598in}}%
\pgfpathclose%
\pgfusepath{fill}%
\end{pgfscope}%
\begin{pgfscope}%
\pgfpathrectangle{\pgfqpoint{0.407592in}{0.222904in}}{\pgfqpoint{1.411163in}{1.411163in}}%
\pgfusepath{clip}%
\pgfsetbuttcap%
\pgfsetroundjoin%
\definecolor{currentfill}{rgb}{0.749020,0.772549,0.792157}%
\pgfsetfillcolor{currentfill}%
\pgfsetfillopacity{0.160000}%
\pgfsetlinewidth{0.000000pt}%
\definecolor{currentstroke}{rgb}{0.000000,0.000000,0.000000}%
\pgfsetstrokecolor{currentstroke}%
\pgfsetdash{}{0pt}%
\pgfpathmoveto{\pgfqpoint{0.908116in}{0.747844in}}%
\pgfpathlineto{\pgfqpoint{0.894889in}{0.757816in}}%
\pgfpathlineto{\pgfqpoint{0.881633in}{0.767811in}}%
\pgfpathlineto{\pgfqpoint{0.885162in}{0.762145in}}%
\pgfpathlineto{\pgfqpoint{0.889607in}{0.756597in}}%
\pgfpathlineto{\pgfqpoint{0.902396in}{0.747287in}}%
\pgfpathlineto{\pgfqpoint{0.915150in}{0.738000in}}%
\pgfpathlineto{\pgfqpoint{0.911234in}{0.742869in}}%
\pgfpathlineto{\pgfqpoint{0.908116in}{0.747844in}}%
\pgfpathclose%
\pgfusepath{fill}%
\end{pgfscope}%
\begin{pgfscope}%
\pgfpathrectangle{\pgfqpoint{0.407592in}{0.222904in}}{\pgfqpoint{1.411163in}{1.411163in}}%
\pgfusepath{clip}%
\pgfsetbuttcap%
\pgfsetroundjoin%
\definecolor{currentfill}{rgb}{0.749020,0.772549,0.792157}%
\pgfsetfillcolor{currentfill}%
\pgfsetfillopacity{0.160000}%
\pgfsetlinewidth{0.000000pt}%
\definecolor{currentstroke}{rgb}{0.000000,0.000000,0.000000}%
\pgfsetstrokecolor{currentstroke}%
\pgfsetdash{}{0pt}%
\pgfpathmoveto{\pgfqpoint{0.795604in}{0.859898in}}%
\pgfpathlineto{\pgfqpoint{0.782000in}{0.871413in}}%
\pgfpathlineto{\pgfqpoint{0.768377in}{0.882943in}}%
\pgfpathlineto{\pgfqpoint{0.767784in}{0.874305in}}%
\pgfpathlineto{\pgfqpoint{0.768605in}{0.865674in}}%
\pgfpathlineto{\pgfqpoint{0.782286in}{0.854821in}}%
\pgfpathlineto{\pgfqpoint{0.795943in}{0.843987in}}%
\pgfpathlineto{\pgfqpoint{0.795127in}{0.851938in}}%
\pgfpathlineto{\pgfqpoint{0.795604in}{0.859898in}}%
\pgfpathclose%
\pgfusepath{fill}%
\end{pgfscope}%
\begin{pgfscope}%
\pgfpathrectangle{\pgfqpoint{0.407592in}{0.222904in}}{\pgfqpoint{1.411163in}{1.411163in}}%
\pgfusepath{clip}%
\pgfsetbuttcap%
\pgfsetroundjoin%
\definecolor{currentfill}{rgb}{0.749020,0.772549,0.792157}%
\pgfsetfillcolor{currentfill}%
\pgfsetfillopacity{0.160000}%
\pgfsetlinewidth{0.000000pt}%
\definecolor{currentstroke}{rgb}{0.000000,0.000000,0.000000}%
\pgfsetstrokecolor{currentstroke}%
\pgfsetdash{}{0pt}%
\pgfpathmoveto{\pgfqpoint{1.185034in}{0.685906in}}%
\pgfpathlineto{\pgfqpoint{1.194698in}{0.693589in}}%
\pgfpathlineto{\pgfqpoint{1.204397in}{0.701300in}}%
\pgfpathlineto{\pgfqpoint{1.211540in}{0.704780in}}%
\pgfpathlineto{\pgfqpoint{1.218131in}{0.708442in}}%
\pgfpathlineto{\pgfqpoint{1.207287in}{0.700140in}}%
\pgfpathlineto{\pgfqpoint{1.196476in}{0.691865in}}%
\pgfpathlineto{\pgfqpoint{1.190983in}{0.688809in}}%
\pgfpathlineto{\pgfqpoint{1.185034in}{0.685906in}}%
\pgfpathclose%
\pgfusepath{fill}%
\end{pgfscope}%
\begin{pgfscope}%
\pgfpathrectangle{\pgfqpoint{0.407592in}{0.222904in}}{\pgfqpoint{1.411163in}{1.411163in}}%
\pgfusepath{clip}%
\pgfsetbuttcap%
\pgfsetroundjoin%
\definecolor{currentfill}{rgb}{0.749020,0.772549,0.792157}%
\pgfsetfillcolor{currentfill}%
\pgfsetfillopacity{0.160000}%
\pgfsetlinewidth{0.000000pt}%
\definecolor{currentstroke}{rgb}{0.000000,0.000000,0.000000}%
\pgfsetstrokecolor{currentstroke}%
\pgfsetdash{}{0pt}%
\pgfpathmoveto{\pgfqpoint{1.023870in}{0.664119in}}%
\pgfpathlineto{\pgfqpoint{1.015426in}{0.670937in}}%
\pgfpathlineto{\pgfqpoint{1.006948in}{0.677783in}}%
\pgfpathlineto{\pgfqpoint{1.014062in}{0.675654in}}%
\pgfpathlineto{\pgfqpoint{1.021521in}{0.673739in}}%
\pgfpathlineto{\pgfqpoint{1.028522in}{0.667298in}}%
\pgfpathlineto{\pgfqpoint{1.035494in}{0.660885in}}%
\pgfpathlineto{\pgfqpoint{1.029545in}{0.662416in}}%
\pgfpathlineto{\pgfqpoint{1.023870in}{0.664119in}}%
\pgfpathclose%
\pgfusepath{fill}%
\end{pgfscope}%
\begin{pgfscope}%
\pgfpathrectangle{\pgfqpoint{0.407592in}{0.222904in}}{\pgfqpoint{1.411163in}{1.411163in}}%
\pgfusepath{clip}%
\pgfsetbuttcap%
\pgfsetroundjoin%
\definecolor{currentfill}{rgb}{0.749020,0.772549,0.792157}%
\pgfsetfillcolor{currentfill}%
\pgfsetfillopacity{0.160000}%
\pgfsetlinewidth{0.000000pt}%
\definecolor{currentstroke}{rgb}{0.000000,0.000000,0.000000}%
\pgfsetstrokecolor{currentstroke}%
\pgfsetdash{}{0pt}%
\pgfpathmoveto{\pgfqpoint{0.940558in}{0.719502in}}%
\pgfpathlineto{\pgfqpoint{0.927871in}{0.728739in}}%
\pgfpathlineto{\pgfqpoint{0.915150in}{0.738000in}}%
\pgfpathlineto{\pgfqpoint{0.919849in}{0.733260in}}%
\pgfpathlineto{\pgfqpoint{0.925313in}{0.728671in}}%
\pgfpathlineto{\pgfqpoint{0.937326in}{0.720066in}}%
\pgfpathlineto{\pgfqpoint{0.949302in}{0.711487in}}%
\pgfpathlineto{\pgfqpoint{0.944605in}{0.715429in}}%
\pgfpathlineto{\pgfqpoint{0.940558in}{0.719502in}}%
\pgfpathclose%
\pgfusepath{fill}%
\end{pgfscope}%
\begin{pgfscope}%
\pgfpathrectangle{\pgfqpoint{0.407592in}{0.222904in}}{\pgfqpoint{1.411163in}{1.411163in}}%
\pgfusepath{clip}%
\pgfsetbuttcap%
\pgfsetroundjoin%
\definecolor{currentfill}{rgb}{0.749020,0.772549,0.792157}%
\pgfsetfillcolor{currentfill}%
\pgfsetfillopacity{0.160000}%
\pgfsetlinewidth{0.000000pt}%
\definecolor{currentstroke}{rgb}{0.000000,0.000000,0.000000}%
\pgfsetstrokecolor{currentstroke}%
\pgfsetdash{}{0pt}%
\pgfpathmoveto{\pgfqpoint{1.519145in}{0.978016in}}%
\pgfpathlineto{\pgfqpoint{1.532845in}{0.989754in}}%
\pgfpathlineto{\pgfqpoint{1.546564in}{1.001508in}}%
\pgfpathlineto{\pgfqpoint{1.545516in}{1.013319in}}%
\pgfpathlineto{\pgfqpoint{1.542494in}{1.024999in}}%
\pgfpathlineto{\pgfqpoint{1.528995in}{1.012628in}}%
\pgfpathlineto{\pgfqpoint{1.515507in}{1.000268in}}%
\pgfpathlineto{\pgfqpoint{1.518255in}{0.989201in}}%
\pgfpathlineto{\pgfqpoint{1.519145in}{0.978016in}}%
\pgfpathclose%
\pgfusepath{fill}%
\end{pgfscope}%
\begin{pgfscope}%
\pgfpathrectangle{\pgfqpoint{0.407592in}{0.222904in}}{\pgfqpoint{1.411163in}{1.411163in}}%
\pgfusepath{clip}%
\pgfsetbuttcap%
\pgfsetroundjoin%
\definecolor{currentfill}{rgb}{0.749020,0.772549,0.792157}%
\pgfsetfillcolor{currentfill}%
\pgfsetfillopacity{0.160000}%
\pgfsetlinewidth{0.000000pt}%
\definecolor{currentstroke}{rgb}{0.000000,0.000000,0.000000}%
\pgfsetstrokecolor{currentstroke}%
\pgfsetdash{}{0pt}%
\pgfpathmoveto{\pgfqpoint{0.449066in}{1.252869in}}%
\pgfpathlineto{\pgfqpoint{0.436075in}{1.265884in}}%
\pgfpathlineto{\pgfqpoint{0.423077in}{1.278906in}}%
\pgfpathlineto{\pgfqpoint{0.407830in}{1.262769in}}%
\pgfpathlineto{\pgfqpoint{0.395362in}{1.246258in}}%
\pgfpathlineto{\pgfqpoint{0.409049in}{1.233736in}}%
\pgfpathlineto{\pgfqpoint{0.422723in}{1.221226in}}%
\pgfpathlineto{\pgfqpoint{0.434555in}{1.237225in}}%
\pgfpathlineto{\pgfqpoint{0.449066in}{1.252869in}}%
\pgfpathclose%
\pgfusepath{fill}%
\end{pgfscope}%
\begin{pgfscope}%
\pgfpathrectangle{\pgfqpoint{0.407592in}{0.222904in}}{\pgfqpoint{1.411163in}{1.411163in}}%
\pgfusepath{clip}%
\pgfsetbuttcap%
\pgfsetroundjoin%
\definecolor{currentfill}{rgb}{0.749020,0.772549,0.792157}%
\pgfsetfillcolor{currentfill}%
\pgfsetfillopacity{0.160000}%
\pgfsetlinewidth{0.000000pt}%
\definecolor{currentstroke}{rgb}{0.000000,0.000000,0.000000}%
\pgfsetstrokecolor{currentstroke}%
\pgfsetdash{}{0pt}%
\pgfpathmoveto{\pgfqpoint{1.130948in}{0.659788in}}%
\pgfpathlineto{\pgfqpoint{1.136136in}{0.666064in}}%
\pgfpathlineto{\pgfqpoint{1.141347in}{0.672367in}}%
\pgfpathlineto{\pgfqpoint{1.149421in}{0.674106in}}%
\pgfpathlineto{\pgfqpoint{1.157223in}{0.676065in}}%
\pgfpathlineto{\pgfqpoint{1.150401in}{0.669391in}}%
\pgfpathlineto{\pgfqpoint{1.143606in}{0.662744in}}%
\pgfpathlineto{\pgfqpoint{1.137384in}{0.661178in}}%
\pgfpathlineto{\pgfqpoint{1.130948in}{0.659788in}}%
\pgfpathclose%
\pgfusepath{fill}%
\end{pgfscope}%
\begin{pgfscope}%
\pgfpathrectangle{\pgfqpoint{0.407592in}{0.222904in}}{\pgfqpoint{1.411163in}{1.411163in}}%
\pgfusepath{clip}%
\pgfsetbuttcap%
\pgfsetroundjoin%
\definecolor{currentfill}{rgb}{0.749020,0.772549,0.792157}%
\pgfsetfillcolor{currentfill}%
\pgfsetfillopacity{0.160000}%
\pgfsetlinewidth{0.000000pt}%
\definecolor{currentstroke}{rgb}{0.000000,0.000000,0.000000}%
\pgfsetstrokecolor{currentstroke}%
\pgfsetdash{}{0pt}%
\pgfpathmoveto{\pgfqpoint{1.319268in}{0.789156in}}%
\pgfpathlineto{\pgfqpoint{1.332259in}{0.799287in}}%
\pgfpathlineto{\pgfqpoint{1.345279in}{0.809441in}}%
\pgfpathlineto{\pgfqpoint{1.349747in}{0.816602in}}%
\pgfpathlineto{\pgfqpoint{1.353055in}{0.823833in}}%
\pgfpathlineto{\pgfqpoint{1.339692in}{0.812989in}}%
\pgfpathlineto{\pgfqpoint{1.326354in}{0.802164in}}%
\pgfpathlineto{\pgfqpoint{1.323332in}{0.795627in}}%
\pgfpathlineto{\pgfqpoint{1.319268in}{0.789156in}}%
\pgfpathclose%
\pgfusepath{fill}%
\end{pgfscope}%
\begin{pgfscope}%
\pgfpathrectangle{\pgfqpoint{0.407592in}{0.222904in}}{\pgfqpoint{1.411163in}{1.411163in}}%
\pgfusepath{clip}%
\pgfsetbuttcap%
\pgfsetroundjoin%
\definecolor{currentfill}{rgb}{0.749020,0.772549,0.792157}%
\pgfsetfillcolor{currentfill}%
\pgfsetfillopacity{0.160000}%
\pgfsetlinewidth{0.000000pt}%
\definecolor{currentstroke}{rgb}{0.000000,0.000000,0.000000}%
\pgfsetstrokecolor{currentstroke}%
\pgfsetdash{}{0pt}%
\pgfpathmoveto{\pgfqpoint{1.815054in}{1.274780in}}%
\pgfpathlineto{\pgfqpoint{1.828812in}{1.287388in}}%
\pgfpathlineto{\pgfqpoint{1.842582in}{1.300007in}}%
\pgfpathlineto{\pgfqpoint{1.832537in}{1.317433in}}%
\pgfpathlineto{\pgfqpoint{1.819435in}{1.334434in}}%
\pgfpathlineto{\pgfqpoint{1.806244in}{1.321339in}}%
\pgfpathlineto{\pgfqpoint{1.793060in}{1.308251in}}%
\pgfpathlineto{\pgfqpoint{1.805534in}{1.291718in}}%
\pgfpathlineto{\pgfqpoint{1.815054in}{1.274780in}}%
\pgfpathclose%
\pgfusepath{fill}%
\end{pgfscope}%
\begin{pgfscope}%
\pgfpathrectangle{\pgfqpoint{0.407592in}{0.222904in}}{\pgfqpoint{1.411163in}{1.411163in}}%
\pgfusepath{clip}%
\pgfsetbuttcap%
\pgfsetroundjoin%
\definecolor{currentfill}{rgb}{0.749020,0.772549,0.792157}%
\pgfsetfillcolor{currentfill}%
\pgfsetfillopacity{0.160000}%
\pgfsetlinewidth{0.000000pt}%
\definecolor{currentstroke}{rgb}{0.000000,0.000000,0.000000}%
\pgfsetstrokecolor{currentstroke}%
\pgfsetdash{}{0pt}%
\pgfpathmoveto{\pgfqpoint{0.666651in}{0.998060in}}%
\pgfpathlineto{\pgfqpoint{0.653208in}{1.010358in}}%
\pgfpathlineto{\pgfqpoint{0.639752in}{1.022669in}}%
\pgfpathlineto{\pgfqpoint{0.634453in}{1.010958in}}%
\pgfpathlineto{\pgfqpoint{0.631105in}{0.999127in}}%
\pgfpathlineto{\pgfqpoint{0.644918in}{0.987436in}}%
\pgfpathlineto{\pgfqpoint{0.658712in}{0.975761in}}%
\pgfpathlineto{\pgfqpoint{0.661762in}{0.986965in}}%
\pgfpathlineto{\pgfqpoint{0.666651in}{0.998060in}}%
\pgfpathclose%
\pgfusepath{fill}%
\end{pgfscope}%
\begin{pgfscope}%
\pgfpathrectangle{\pgfqpoint{0.407592in}{0.222904in}}{\pgfqpoint{1.411163in}{1.411163in}}%
\pgfusepath{clip}%
\pgfsetbuttcap%
\pgfsetroundjoin%
\definecolor{currentfill}{rgb}{0.749020,0.772549,0.792157}%
\pgfsetfillcolor{currentfill}%
\pgfsetfillopacity{0.160000}%
\pgfsetlinewidth{0.000000pt}%
\definecolor{currentstroke}{rgb}{0.000000,0.000000,0.000000}%
\pgfsetstrokecolor{currentstroke}%
\pgfsetdash{}{0pt}%
\pgfpathmoveto{\pgfqpoint{0.850322in}{0.800848in}}%
\pgfpathlineto{\pgfqpoint{0.836764in}{0.811604in}}%
\pgfpathlineto{\pgfqpoint{0.823182in}{0.822378in}}%
\pgfpathlineto{\pgfqpoint{0.825158in}{0.815159in}}%
\pgfpathlineto{\pgfqpoint{0.828311in}{0.808016in}}%
\pgfpathlineto{\pgfqpoint{0.841686in}{0.797931in}}%
\pgfpathlineto{\pgfqpoint{0.855032in}{0.787868in}}%
\pgfpathlineto{\pgfqpoint{0.852150in}{0.794322in}}%
\pgfpathlineto{\pgfqpoint{0.850322in}{0.800848in}}%
\pgfpathclose%
\pgfusepath{fill}%
\end{pgfscope}%
\begin{pgfscope}%
\pgfpathrectangle{\pgfqpoint{0.407592in}{0.222904in}}{\pgfqpoint{1.411163in}{1.411163in}}%
\pgfusepath{clip}%
\pgfsetbuttcap%
\pgfsetroundjoin%
\definecolor{currentfill}{rgb}{0.749020,0.772549,0.792157}%
\pgfsetfillcolor{currentfill}%
\pgfsetfillopacity{0.160000}%
\pgfsetlinewidth{0.000000pt}%
\definecolor{currentstroke}{rgb}{0.000000,0.000000,0.000000}%
\pgfsetstrokecolor{currentstroke}%
\pgfsetdash{}{0pt}%
\pgfpathmoveto{\pgfqpoint{1.035494in}{0.660885in}}%
\pgfpathlineto{\pgfqpoint{1.028522in}{0.667298in}}%
\pgfpathlineto{\pgfqpoint{1.021521in}{0.673739in}}%
\pgfpathlineto{\pgfqpoint{1.029291in}{0.672046in}}%
\pgfpathlineto{\pgfqpoint{1.037338in}{0.670583in}}%
\pgfpathlineto{\pgfqpoint{1.042732in}{0.664459in}}%
\pgfpathlineto{\pgfqpoint{1.048102in}{0.658362in}}%
\pgfpathlineto{\pgfqpoint{1.041689in}{0.659531in}}%
\pgfpathlineto{\pgfqpoint{1.035494in}{0.660885in}}%
\pgfpathclose%
\pgfusepath{fill}%
\end{pgfscope}%
\begin{pgfscope}%
\pgfpathrectangle{\pgfqpoint{0.407592in}{0.222904in}}{\pgfqpoint{1.411163in}{1.411163in}}%
\pgfusepath{clip}%
\pgfsetbuttcap%
\pgfsetroundjoin%
\definecolor{currentfill}{rgb}{0.749020,0.772549,0.792157}%
\pgfsetfillcolor{currentfill}%
\pgfsetfillopacity{0.160000}%
\pgfsetlinewidth{0.000000pt}%
\definecolor{currentstroke}{rgb}{0.000000,0.000000,0.000000}%
\pgfsetstrokecolor{currentstroke}%
\pgfsetdash{}{0pt}%
\pgfpathmoveto{\pgfqpoint{0.423077in}{1.278906in}}%
\pgfpathlineto{\pgfqpoint{0.410072in}{1.291934in}}%
\pgfpathlineto{\pgfqpoint{0.397060in}{1.304970in}}%
\pgfpathlineto{\pgfqpoint{0.381066in}{1.288350in}}%
\pgfpathlineto{\pgfqpoint{0.367948in}{1.271338in}}%
\pgfpathlineto{\pgfqpoint{0.381662in}{1.258792in}}%
\pgfpathlineto{\pgfqpoint{0.395362in}{1.246258in}}%
\pgfpathlineto{\pgfqpoint{0.407830in}{1.262769in}}%
\pgfpathlineto{\pgfqpoint{0.423077in}{1.278906in}}%
\pgfpathclose%
\pgfusepath{fill}%
\end{pgfscope}%
\begin{pgfscope}%
\pgfpathrectangle{\pgfqpoint{0.407592in}{0.222904in}}{\pgfqpoint{1.411163in}{1.411163in}}%
\pgfusepath{clip}%
\pgfsetbuttcap%
\pgfsetroundjoin%
\definecolor{currentfill}{rgb}{0.749020,0.772549,0.792157}%
\pgfsetfillcolor{currentfill}%
\pgfsetfillopacity{0.160000}%
\pgfsetlinewidth{0.000000pt}%
\definecolor{currentstroke}{rgb}{0.000000,0.000000,0.000000}%
\pgfsetstrokecolor{currentstroke}%
\pgfsetdash{}{0pt}%
\pgfpathmoveto{\pgfqpoint{1.406741in}{0.867403in}}%
\pgfpathlineto{\pgfqpoint{1.420222in}{0.878344in}}%
\pgfpathlineto{\pgfqpoint{1.433727in}{0.889305in}}%
\pgfpathlineto{\pgfqpoint{1.436292in}{0.898613in}}%
\pgfpathlineto{\pgfqpoint{1.437326in}{0.907915in}}%
\pgfpathlineto{\pgfqpoint{1.423753in}{0.896287in}}%
\pgfpathlineto{\pgfqpoint{1.410198in}{0.884673in}}%
\pgfpathlineto{\pgfqpoint{1.409175in}{0.876039in}}%
\pgfpathlineto{\pgfqpoint{1.406741in}{0.867403in}}%
\pgfpathclose%
\pgfusepath{fill}%
\end{pgfscope}%
\begin{pgfscope}%
\pgfpathrectangle{\pgfqpoint{0.407592in}{0.222904in}}{\pgfqpoint{1.411163in}{1.411163in}}%
\pgfusepath{clip}%
\pgfsetbuttcap%
\pgfsetroundjoin%
\definecolor{currentfill}{rgb}{0.749020,0.772549,0.792157}%
\pgfsetfillcolor{currentfill}%
\pgfsetfillopacity{0.160000}%
\pgfsetlinewidth{0.000000pt}%
\definecolor{currentstroke}{rgb}{0.000000,0.000000,0.000000}%
\pgfsetstrokecolor{currentstroke}%
\pgfsetdash{}{0pt}%
\pgfpathmoveto{\pgfqpoint{1.842582in}{1.300007in}}%
\pgfpathlineto{\pgfqpoint{1.856365in}{1.312638in}}%
\pgfpathlineto{\pgfqpoint{1.846053in}{1.330303in}}%
\pgfpathlineto{\pgfqpoint{1.832631in}{1.347535in}}%
\pgfpathlineto{\pgfqpoint{1.819435in}{1.334434in}}%
\pgfpathlineto{\pgfqpoint{1.832537in}{1.317433in}}%
\pgfusepath{fill}%
\end{pgfscope}%
\begin{pgfscope}%
\pgfpathrectangle{\pgfqpoint{0.407592in}{0.222904in}}{\pgfqpoint{1.411163in}{1.411163in}}%
\pgfusepath{clip}%
\pgfsetbuttcap%
\pgfsetroundjoin%
\definecolor{currentfill}{rgb}{0.749020,0.772549,0.792157}%
\pgfsetfillcolor{currentfill}%
\pgfsetfillopacity{0.160000}%
\pgfsetlinewidth{0.000000pt}%
\definecolor{currentstroke}{rgb}{0.000000,0.000000,0.000000}%
\pgfsetstrokecolor{currentstroke}%
\pgfsetdash{}{0pt}%
\pgfpathmoveto{\pgfqpoint{0.982547in}{0.688214in}}%
\pgfpathlineto{\pgfqpoint{0.971578in}{0.696127in}}%
\pgfpathlineto{\pgfqpoint{0.960572in}{0.704067in}}%
\pgfpathlineto{\pgfqpoint{0.967097in}{0.700624in}}%
\pgfpathlineto{\pgfqpoint{0.974180in}{0.697382in}}%
\pgfpathlineto{\pgfqpoint{0.984048in}{0.689996in}}%
\pgfpathlineto{\pgfqpoint{0.993880in}{0.682638in}}%
\pgfpathlineto{\pgfqpoint{0.987984in}{0.685342in}}%
\pgfpathlineto{\pgfqpoint{0.982547in}{0.688214in}}%
\pgfpathclose%
\pgfusepath{fill}%
\end{pgfscope}%
\begin{pgfscope}%
\pgfpathrectangle{\pgfqpoint{0.407592in}{0.222904in}}{\pgfqpoint{1.411163in}{1.411163in}}%
\pgfusepath{clip}%
\pgfsetbuttcap%
\pgfsetroundjoin%
\definecolor{currentfill}{rgb}{0.749020,0.772549,0.792157}%
\pgfsetfillcolor{currentfill}%
\pgfsetfillopacity{0.160000}%
\pgfsetlinewidth{0.000000pt}%
\definecolor{currentstroke}{rgb}{0.000000,0.000000,0.000000}%
\pgfsetstrokecolor{currentstroke}%
\pgfsetdash{}{0pt}%
\pgfpathmoveto{\pgfqpoint{1.546564in}{1.001508in}}%
\pgfpathlineto{\pgfqpoint{1.560302in}{1.013278in}}%
\pgfpathlineto{\pgfqpoint{1.574058in}{1.025063in}}%
\pgfpathlineto{\pgfqpoint{1.572838in}{1.037491in}}%
\pgfpathlineto{\pgfqpoint{1.569529in}{1.049775in}}%
\pgfpathlineto{\pgfqpoint{1.556005in}{1.037381in}}%
\pgfpathlineto{\pgfqpoint{1.542494in}{1.024999in}}%
\pgfpathlineto{\pgfqpoint{1.545516in}{1.013319in}}%
\pgfpathlineto{\pgfqpoint{1.546564in}{1.001508in}}%
\pgfpathclose%
\pgfusepath{fill}%
\end{pgfscope}%
\begin{pgfscope}%
\pgfpathrectangle{\pgfqpoint{0.407592in}{0.222904in}}{\pgfqpoint{1.411163in}{1.411163in}}%
\pgfusepath{clip}%
\pgfsetbuttcap%
\pgfsetroundjoin%
\definecolor{currentfill}{rgb}{0.749020,0.772549,0.792157}%
\pgfsetfillcolor{currentfill}%
\pgfsetfillopacity{0.160000}%
\pgfsetlinewidth{0.000000pt}%
\definecolor{currentstroke}{rgb}{0.000000,0.000000,0.000000}%
\pgfsetstrokecolor{currentstroke}%
\pgfsetdash{}{0pt}%
\pgfpathmoveto{\pgfqpoint{1.117547in}{0.657564in}}%
\pgfpathlineto{\pgfqpoint{1.121030in}{0.663560in}}%
\pgfpathlineto{\pgfqpoint{1.124529in}{0.669584in}}%
\pgfpathlineto{\pgfqpoint{1.133037in}{0.670858in}}%
\pgfpathlineto{\pgfqpoint{1.141347in}{0.672367in}}%
\pgfpathlineto{\pgfqpoint{1.136136in}{0.666064in}}%
\pgfpathlineto{\pgfqpoint{1.130948in}{0.659788in}}%
\pgfpathlineto{\pgfqpoint{1.124325in}{0.658581in}}%
\pgfpathlineto{\pgfqpoint{1.117547in}{0.657564in}}%
\pgfpathclose%
\pgfusepath{fill}%
\end{pgfscope}%
\begin{pgfscope}%
\pgfpathrectangle{\pgfqpoint{0.407592in}{0.222904in}}{\pgfqpoint{1.411163in}{1.411163in}}%
\pgfusepath{clip}%
\pgfsetbuttcap%
\pgfsetroundjoin%
\definecolor{currentfill}{rgb}{0.749020,0.772549,0.792157}%
\pgfsetfillcolor{currentfill}%
\pgfsetfillopacity{0.160000}%
\pgfsetlinewidth{0.000000pt}%
\definecolor{currentstroke}{rgb}{0.000000,0.000000,0.000000}%
\pgfsetstrokecolor{currentstroke}%
\pgfsetdash{}{0pt}%
\pgfpathmoveto{\pgfqpoint{0.768377in}{0.882943in}}%
\pgfpathlineto{\pgfqpoint{0.754735in}{0.894489in}}%
\pgfpathlineto{\pgfqpoint{0.741074in}{0.906052in}}%
\pgfpathlineto{\pgfqpoint{0.740354in}{0.896745in}}%
\pgfpathlineto{\pgfqpoint{0.741166in}{0.887440in}}%
\pgfpathlineto{\pgfqpoint{0.754898in}{0.876547in}}%
\pgfpathlineto{\pgfqpoint{0.768605in}{0.865674in}}%
\pgfpathlineto{\pgfqpoint{0.767784in}{0.874305in}}%
\pgfpathlineto{\pgfqpoint{0.768377in}{0.882943in}}%
\pgfpathclose%
\pgfusepath{fill}%
\end{pgfscope}%
\begin{pgfscope}%
\pgfpathrectangle{\pgfqpoint{0.407592in}{0.222904in}}{\pgfqpoint{1.411163in}{1.411163in}}%
\pgfusepath{clip}%
\pgfsetbuttcap%
\pgfsetroundjoin%
\definecolor{currentfill}{rgb}{0.749020,0.772549,0.792157}%
\pgfsetfillcolor{currentfill}%
\pgfsetfillopacity{0.160000}%
\pgfsetlinewidth{0.000000pt}%
\definecolor{currentstroke}{rgb}{0.000000,0.000000,0.000000}%
\pgfsetstrokecolor{currentstroke}%
\pgfsetdash{}{0pt}%
\pgfpathmoveto{\pgfqpoint{1.048102in}{0.658362in}}%
\pgfpathlineto{\pgfqpoint{1.042732in}{0.664459in}}%
\pgfpathlineto{\pgfqpoint{1.037338in}{0.670583in}}%
\pgfpathlineto{\pgfqpoint{1.045624in}{0.669358in}}%
\pgfpathlineto{\pgfqpoint{1.054113in}{0.668375in}}%
\pgfpathlineto{\pgfqpoint{1.057798in}{0.662473in}}%
\pgfpathlineto{\pgfqpoint{1.061467in}{0.656598in}}%
\pgfpathlineto{\pgfqpoint{1.054705in}{0.657383in}}%
\pgfpathlineto{\pgfqpoint{1.048102in}{0.658362in}}%
\pgfpathclose%
\pgfusepath{fill}%
\end{pgfscope}%
\begin{pgfscope}%
\pgfpathrectangle{\pgfqpoint{0.407592in}{0.222904in}}{\pgfqpoint{1.411163in}{1.411163in}}%
\pgfusepath{clip}%
\pgfsetbuttcap%
\pgfsetroundjoin%
\definecolor{currentfill}{rgb}{0.749020,0.772549,0.792157}%
\pgfsetfillcolor{currentfill}%
\pgfsetfillopacity{0.160000}%
\pgfsetlinewidth{0.000000pt}%
\definecolor{currentstroke}{rgb}{0.000000,0.000000,0.000000}%
\pgfsetstrokecolor{currentstroke}%
\pgfsetdash{}{0pt}%
\pgfpathmoveto{\pgfqpoint{1.218131in}{0.708442in}}%
\pgfpathlineto{\pgfqpoint{1.229011in}{0.716770in}}%
\pgfpathlineto{\pgfqpoint{1.239925in}{0.725125in}}%
\pgfpathlineto{\pgfqpoint{1.246930in}{0.729579in}}%
\pgfpathlineto{\pgfqpoint{1.235519in}{0.720909in}}%
\pgfpathlineto{\pgfqpoint{1.224142in}{0.712266in}}%
\pgfusepath{fill}%
\end{pgfscope}%
\begin{pgfscope}%
\pgfpathrectangle{\pgfqpoint{0.407592in}{0.222904in}}{\pgfqpoint{1.411163in}{1.411163in}}%
\pgfusepath{clip}%
\pgfsetbuttcap%
\pgfsetroundjoin%
\definecolor{currentfill}{rgb}{0.749020,0.772549,0.792157}%
\pgfsetfillcolor{currentfill}%
\pgfsetfillopacity{0.160000}%
\pgfsetlinewidth{0.000000pt}%
\definecolor{currentstroke}{rgb}{0.000000,0.000000,0.000000}%
\pgfsetstrokecolor{currentstroke}%
\pgfsetdash{}{0pt}%
\pgfpathmoveto{\pgfqpoint{0.397060in}{1.304970in}}%
\pgfpathlineto{\pgfqpoint{0.384042in}{1.318011in}}%
\pgfpathlineto{\pgfqpoint{0.371018in}{1.331060in}}%
\pgfpathlineto{\pgfqpoint{0.354262in}{1.313969in}}%
\pgfpathlineto{\pgfqpoint{0.340483in}{1.296466in}}%
\pgfpathlineto{\pgfqpoint{0.354222in}{1.283896in}}%
\pgfpathlineto{\pgfqpoint{0.367948in}{1.271338in}}%
\pgfpathlineto{\pgfqpoint{0.381066in}{1.288350in}}%
\pgfpathlineto{\pgfqpoint{0.397060in}{1.304970in}}%
\pgfpathclose%
\pgfusepath{fill}%
\end{pgfscope}%
\begin{pgfscope}%
\pgfpathrectangle{\pgfqpoint{0.407592in}{0.222904in}}{\pgfqpoint{1.411163in}{1.411163in}}%
\pgfusepath{clip}%
\pgfsetbuttcap%
\pgfsetroundjoin%
\definecolor{currentfill}{rgb}{0.749020,0.772549,0.792157}%
\pgfsetfillcolor{currentfill}%
\pgfsetfillopacity{0.160000}%
\pgfsetlinewidth{0.000000pt}%
\definecolor{currentstroke}{rgb}{0.000000,0.000000,0.000000}%
\pgfsetstrokecolor{currentstroke}%
\pgfsetdash{}{0pt}%
\pgfpathmoveto{\pgfqpoint{1.171873in}{0.680606in}}%
\pgfpathlineto{\pgfqpoint{1.180209in}{0.687761in}}%
\pgfpathlineto{\pgfqpoint{1.188578in}{0.694944in}}%
\pgfpathlineto{\pgfqpoint{1.196731in}{0.698015in}}%
\pgfpathlineto{\pgfqpoint{1.204397in}{0.701300in}}%
\pgfpathlineto{\pgfqpoint{1.194698in}{0.693589in}}%
\pgfpathlineto{\pgfqpoint{1.185034in}{0.685906in}}%
\pgfpathlineto{\pgfqpoint{1.178654in}{0.683167in}}%
\pgfpathlineto{\pgfqpoint{1.171873in}{0.680606in}}%
\pgfpathclose%
\pgfusepath{fill}%
\end{pgfscope}%
\begin{pgfscope}%
\pgfpathrectangle{\pgfqpoint{0.407592in}{0.222904in}}{\pgfqpoint{1.411163in}{1.411163in}}%
\pgfusepath{clip}%
\pgfsetbuttcap%
\pgfsetroundjoin%
\definecolor{currentfill}{rgb}{0.749020,0.772549,0.792157}%
\pgfsetfillcolor{currentfill}%
\pgfsetfillopacity{0.160000}%
\pgfsetlinewidth{0.000000pt}%
\definecolor{currentstroke}{rgb}{0.000000,0.000000,0.000000}%
\pgfsetstrokecolor{currentstroke}%
\pgfsetdash{}{0pt}%
\pgfpathmoveto{\pgfqpoint{1.103644in}{0.656113in}}%
\pgfpathlineto{\pgfqpoint{1.105356in}{0.661927in}}%
\pgfpathlineto{\pgfqpoint{1.107075in}{0.667768in}}%
\pgfpathlineto{\pgfqpoint{1.115861in}{0.668553in}}%
\pgfpathlineto{\pgfqpoint{1.124529in}{0.669584in}}%
\pgfpathlineto{\pgfqpoint{1.121030in}{0.663560in}}%
\pgfpathlineto{\pgfqpoint{1.117547in}{0.657564in}}%
\pgfpathlineto{\pgfqpoint{1.110643in}{0.656739in}}%
\pgfpathlineto{\pgfqpoint{1.103644in}{0.656113in}}%
\pgfpathclose%
\pgfusepath{fill}%
\end{pgfscope}%
\begin{pgfscope}%
\pgfpathrectangle{\pgfqpoint{0.407592in}{0.222904in}}{\pgfqpoint{1.411163in}{1.411163in}}%
\pgfusepath{clip}%
\pgfsetbuttcap%
\pgfsetroundjoin%
\definecolor{currentfill}{rgb}{0.749020,0.772549,0.792157}%
\pgfsetfillcolor{currentfill}%
\pgfsetfillopacity{0.160000}%
\pgfsetlinewidth{0.000000pt}%
\definecolor{currentstroke}{rgb}{0.000000,0.000000,0.000000}%
\pgfsetstrokecolor{currentstroke}%
\pgfsetdash{}{0pt}%
\pgfpathmoveto{\pgfqpoint{0.639752in}{1.022669in}}%
\pgfpathlineto{\pgfqpoint{0.626284in}{1.034991in}}%
\pgfpathlineto{\pgfqpoint{0.612803in}{1.047324in}}%
\pgfpathlineto{\pgfqpoint{0.607079in}{1.035007in}}%
\pgfpathlineto{\pgfqpoint{0.603422in}{1.022558in}}%
\pgfpathlineto{\pgfqpoint{0.617273in}{1.010834in}}%
\pgfpathlineto{\pgfqpoint{0.631105in}{0.999127in}}%
\pgfpathlineto{\pgfqpoint{0.634453in}{1.010958in}}%
\pgfpathlineto{\pgfqpoint{0.639752in}{1.022669in}}%
\pgfpathclose%
\pgfusepath{fill}%
\end{pgfscope}%
\begin{pgfscope}%
\pgfpathrectangle{\pgfqpoint{0.407592in}{0.222904in}}{\pgfqpoint{1.411163in}{1.411163in}}%
\pgfusepath{clip}%
\pgfsetbuttcap%
\pgfsetroundjoin%
\definecolor{currentfill}{rgb}{0.749020,0.772549,0.792157}%
\pgfsetfillcolor{currentfill}%
\pgfsetfillopacity{0.160000}%
\pgfsetlinewidth{0.000000pt}%
\definecolor{currentstroke}{rgb}{0.000000,0.000000,0.000000}%
\pgfsetstrokecolor{currentstroke}%
\pgfsetdash{}{0pt}%
\pgfpathmoveto{\pgfqpoint{1.283360in}{0.757700in}}%
\pgfpathlineto{\pgfqpoint{1.295700in}{0.767102in}}%
\pgfpathlineto{\pgfqpoint{1.308072in}{0.776530in}}%
\pgfpathlineto{\pgfqpoint{1.314176in}{0.782781in}}%
\pgfpathlineto{\pgfqpoint{1.319268in}{0.789156in}}%
\pgfpathlineto{\pgfqpoint{1.306306in}{0.779047in}}%
\pgfpathlineto{\pgfqpoint{1.293372in}{0.768960in}}%
\pgfpathlineto{\pgfqpoint{1.288814in}{0.763273in}}%
\pgfpathlineto{\pgfqpoint{1.283360in}{0.757700in}}%
\pgfpathclose%
\pgfusepath{fill}%
\end{pgfscope}%
\begin{pgfscope}%
\pgfpathrectangle{\pgfqpoint{0.407592in}{0.222904in}}{\pgfqpoint{1.411163in}{1.411163in}}%
\pgfusepath{clip}%
\pgfsetbuttcap%
\pgfsetroundjoin%
\definecolor{currentfill}{rgb}{0.749020,0.772549,0.792157}%
\pgfsetfillcolor{currentfill}%
\pgfsetfillopacity{0.160000}%
\pgfsetlinewidth{0.000000pt}%
\definecolor{currentstroke}{rgb}{0.000000,0.000000,0.000000}%
\pgfsetstrokecolor{currentstroke}%
\pgfsetdash{}{0pt}%
\pgfpathmoveto{\pgfqpoint{1.061467in}{0.656598in}}%
\pgfpathlineto{\pgfqpoint{1.057798in}{0.662473in}}%
\pgfpathlineto{\pgfqpoint{1.054113in}{0.668375in}}%
\pgfpathlineto{\pgfqpoint{1.062765in}{0.667641in}}%
\pgfpathlineto{\pgfqpoint{1.071541in}{0.667158in}}%
\pgfpathlineto{\pgfqpoint{1.073449in}{0.661378in}}%
\pgfpathlineto{\pgfqpoint{1.075348in}{0.655625in}}%
\pgfpathlineto{\pgfqpoint{1.068358in}{0.656011in}}%
\pgfpathlineto{\pgfqpoint{1.061467in}{0.656598in}}%
\pgfpathclose%
\pgfusepath{fill}%
\end{pgfscope}%
\begin{pgfscope}%
\pgfpathrectangle{\pgfqpoint{0.407592in}{0.222904in}}{\pgfqpoint{1.411163in}{1.411163in}}%
\pgfusepath{clip}%
\pgfsetbuttcap%
\pgfsetroundjoin%
\definecolor{currentfill}{rgb}{0.749020,0.772549,0.792157}%
\pgfsetfillcolor{currentfill}%
\pgfsetfillopacity{0.160000}%
\pgfsetlinewidth{0.000000pt}%
\definecolor{currentstroke}{rgb}{0.000000,0.000000,0.000000}%
\pgfsetstrokecolor{currentstroke}%
\pgfsetdash{}{0pt}%
\pgfpathmoveto{\pgfqpoint{1.089492in}{0.655462in}}%
\pgfpathlineto{\pgfqpoint{1.089398in}{0.661195in}}%
\pgfpathlineto{\pgfqpoint{1.089304in}{0.666954in}}%
\pgfpathlineto{\pgfqpoint{1.098209in}{0.667234in}}%
\pgfpathlineto{\pgfqpoint{1.107075in}{0.667768in}}%
\pgfpathlineto{\pgfqpoint{1.105356in}{0.661927in}}%
\pgfpathlineto{\pgfqpoint{1.103644in}{0.656113in}}%
\pgfpathlineto{\pgfqpoint{1.096584in}{0.655686in}}%
\pgfpathlineto{\pgfqpoint{1.089492in}{0.655462in}}%
\pgfpathclose%
\pgfusepath{fill}%
\end{pgfscope}%
\begin{pgfscope}%
\pgfpathrectangle{\pgfqpoint{0.407592in}{0.222904in}}{\pgfqpoint{1.411163in}{1.411163in}}%
\pgfusepath{clip}%
\pgfsetbuttcap%
\pgfsetroundjoin%
\definecolor{currentfill}{rgb}{0.749020,0.772549,0.792157}%
\pgfsetfillcolor{currentfill}%
\pgfsetfillopacity{0.160000}%
\pgfsetlinewidth{0.000000pt}%
\definecolor{currentstroke}{rgb}{0.000000,0.000000,0.000000}%
\pgfsetstrokecolor{currentstroke}%
\pgfsetdash{}{0pt}%
\pgfpathmoveto{\pgfqpoint{0.371018in}{1.331060in}}%
\pgfpathlineto{\pgfqpoint{0.357986in}{1.344115in}}%
\pgfpathlineto{\pgfqpoint{0.340846in}{1.326792in}}%
\pgfpathlineto{\pgfqpoint{0.326731in}{1.309048in}}%
\pgfpathlineto{\pgfqpoint{0.340483in}{1.296466in}}%
\pgfpathlineto{\pgfqpoint{0.354262in}{1.313969in}}%
\pgfusepath{fill}%
\end{pgfscope}%
\begin{pgfscope}%
\pgfpathrectangle{\pgfqpoint{0.407592in}{0.222904in}}{\pgfqpoint{1.411163in}{1.411163in}}%
\pgfusepath{clip}%
\pgfsetbuttcap%
\pgfsetroundjoin%
\definecolor{currentfill}{rgb}{0.749020,0.772549,0.792157}%
\pgfsetfillcolor{currentfill}%
\pgfsetfillopacity{0.160000}%
\pgfsetlinewidth{0.000000pt}%
\definecolor{currentstroke}{rgb}{0.000000,0.000000,0.000000}%
\pgfsetstrokecolor{currentstroke}%
\pgfsetdash{}{0pt}%
\pgfpathmoveto{\pgfqpoint{1.075348in}{0.655625in}}%
\pgfpathlineto{\pgfqpoint{1.073449in}{0.661378in}}%
\pgfpathlineto{\pgfqpoint{1.071541in}{0.667158in}}%
\pgfpathlineto{\pgfqpoint{1.080401in}{0.666929in}}%
\pgfpathlineto{\pgfqpoint{1.089304in}{0.666954in}}%
\pgfpathlineto{\pgfqpoint{1.089398in}{0.661195in}}%
\pgfpathlineto{\pgfqpoint{1.089492in}{0.655462in}}%
\pgfpathlineto{\pgfqpoint{1.082403in}{0.655442in}}%
\pgfpathlineto{\pgfqpoint{1.075348in}{0.655625in}}%
\pgfpathclose%
\pgfusepath{fill}%
\end{pgfscope}%
\begin{pgfscope}%
\pgfpathrectangle{\pgfqpoint{0.407592in}{0.222904in}}{\pgfqpoint{1.411163in}{1.411163in}}%
\pgfusepath{clip}%
\pgfsetbuttcap%
\pgfsetroundjoin%
\definecolor{currentfill}{rgb}{0.749020,0.772549,0.792157}%
\pgfsetfillcolor{currentfill}%
\pgfsetfillopacity{0.160000}%
\pgfsetlinewidth{0.000000pt}%
\definecolor{currentstroke}{rgb}{0.000000,0.000000,0.000000}%
\pgfsetstrokecolor{currentstroke}%
\pgfsetdash{}{0pt}%
\pgfpathmoveto{\pgfqpoint{1.574058in}{1.025063in}}%
\pgfpathlineto{\pgfqpoint{1.587832in}{1.036865in}}%
\pgfpathlineto{\pgfqpoint{1.601625in}{1.048682in}}%
\pgfpathlineto{\pgfqpoint{1.600222in}{1.061717in}}%
\pgfpathlineto{\pgfqpoint{1.596613in}{1.074595in}}%
\pgfpathlineto{\pgfqpoint{1.583065in}{1.062179in}}%
\pgfpathlineto{\pgfqpoint{1.569529in}{1.049775in}}%
\pgfpathlineto{\pgfqpoint{1.572838in}{1.037491in}}%
\pgfpathlineto{\pgfqpoint{1.574058in}{1.025063in}}%
\pgfpathclose%
\pgfusepath{fill}%
\end{pgfscope}%
\begin{pgfscope}%
\pgfpathrectangle{\pgfqpoint{0.407592in}{0.222904in}}{\pgfqpoint{1.411163in}{1.411163in}}%
\pgfusepath{clip}%
\pgfsetbuttcap%
\pgfsetroundjoin%
\definecolor{currentfill}{rgb}{0.749020,0.772549,0.792157}%
\pgfsetfillcolor{currentfill}%
\pgfsetfillopacity{0.160000}%
\pgfsetlinewidth{0.000000pt}%
\definecolor{currentstroke}{rgb}{0.000000,0.000000,0.000000}%
\pgfsetstrokecolor{currentstroke}%
\pgfsetdash{}{0pt}%
\pgfpathmoveto{\pgfqpoint{0.881633in}{0.767811in}}%
\pgfpathlineto{\pgfqpoint{0.868348in}{0.777829in}}%
\pgfpathlineto{\pgfqpoint{0.855032in}{0.787868in}}%
\pgfpathlineto{\pgfqpoint{0.858962in}{0.781516in}}%
\pgfpathlineto{\pgfqpoint{0.863928in}{0.775293in}}%
\pgfpathlineto{\pgfqpoint{0.876785in}{0.765933in}}%
\pgfpathlineto{\pgfqpoint{0.889607in}{0.756597in}}%
\pgfpathlineto{\pgfqpoint{0.885162in}{0.762145in}}%
\pgfpathlineto{\pgfqpoint{0.881633in}{0.767811in}}%
\pgfpathclose%
\pgfusepath{fill}%
\end{pgfscope}%
\begin{pgfscope}%
\pgfpathrectangle{\pgfqpoint{0.407592in}{0.222904in}}{\pgfqpoint{1.411163in}{1.411163in}}%
\pgfusepath{clip}%
\pgfsetbuttcap%
\pgfsetroundjoin%
\definecolor{currentfill}{rgb}{0.749020,0.772549,0.792157}%
\pgfsetfillcolor{currentfill}%
\pgfsetfillopacity{0.160000}%
\pgfsetlinewidth{0.000000pt}%
\definecolor{currentstroke}{rgb}{0.000000,0.000000,0.000000}%
\pgfsetstrokecolor{currentstroke}%
\pgfsetdash{}{0pt}%
\pgfpathmoveto{\pgfqpoint{1.345279in}{0.809441in}}%
\pgfpathlineto{\pgfqpoint{1.358327in}{0.819617in}}%
\pgfpathlineto{\pgfqpoint{1.371404in}{0.829816in}}%
\pgfpathlineto{\pgfqpoint{1.376269in}{0.837661in}}%
\pgfpathlineto{\pgfqpoint{1.379850in}{0.845580in}}%
\pgfpathlineto{\pgfqpoint{1.366440in}{0.834697in}}%
\pgfpathlineto{\pgfqpoint{1.353055in}{0.823833in}}%
\pgfpathlineto{\pgfqpoint{1.349747in}{0.816602in}}%
\pgfpathlineto{\pgfqpoint{1.345279in}{0.809441in}}%
\pgfpathclose%
\pgfusepath{fill}%
\end{pgfscope}%
\begin{pgfscope}%
\pgfpathrectangle{\pgfqpoint{0.407592in}{0.222904in}}{\pgfqpoint{1.411163in}{1.411163in}}%
\pgfusepath{clip}%
\pgfsetbuttcap%
\pgfsetroundjoin%
\definecolor{currentfill}{rgb}{0.749020,0.772549,0.792157}%
\pgfsetfillcolor{currentfill}%
\pgfsetfillopacity{0.160000}%
\pgfsetlinewidth{0.000000pt}%
\definecolor{currentstroke}{rgb}{0.000000,0.000000,0.000000}%
\pgfsetstrokecolor{currentstroke}%
\pgfsetdash{}{0pt}%
\pgfpathmoveto{\pgfqpoint{0.949302in}{0.711487in}}%
\pgfpathlineto{\pgfqpoint{0.937326in}{0.720066in}}%
\pgfpathlineto{\pgfqpoint{0.925313in}{0.728671in}}%
\pgfpathlineto{\pgfqpoint{0.931521in}{0.724253in}}%
\pgfpathlineto{\pgfqpoint{0.938446in}{0.720028in}}%
\pgfpathlineto{\pgfqpoint{0.949528in}{0.712034in}}%
\pgfpathlineto{\pgfqpoint{0.960572in}{0.704067in}}%
\pgfpathlineto{\pgfqpoint{0.954632in}{0.707694in}}%
\pgfpathlineto{\pgfqpoint{0.949302in}{0.711487in}}%
\pgfpathclose%
\pgfusepath{fill}%
\end{pgfscope}%
\begin{pgfscope}%
\pgfpathrectangle{\pgfqpoint{0.407592in}{0.222904in}}{\pgfqpoint{1.411163in}{1.411163in}}%
\pgfusepath{clip}%
\pgfsetbuttcap%
\pgfsetroundjoin%
\definecolor{currentfill}{rgb}{0.749020,0.772549,0.792157}%
\pgfsetfillcolor{currentfill}%
\pgfsetfillopacity{0.160000}%
\pgfsetlinewidth{0.000000pt}%
\definecolor{currentstroke}{rgb}{0.000000,0.000000,0.000000}%
\pgfsetstrokecolor{currentstroke}%
\pgfsetdash{}{0pt}%
\pgfpathmoveto{\pgfqpoint{1.433727in}{0.889305in}}%
\pgfpathlineto{\pgfqpoint{1.447257in}{0.900285in}}%
\pgfpathlineto{\pgfqpoint{1.460810in}{0.911284in}}%
\pgfpathlineto{\pgfqpoint{1.463494in}{0.921257in}}%
\pgfpathlineto{\pgfqpoint{1.464526in}{0.931220in}}%
\pgfpathlineto{\pgfqpoint{1.450917in}{0.919560in}}%
\pgfpathlineto{\pgfqpoint{1.437326in}{0.907915in}}%
\pgfpathlineto{\pgfqpoint{1.436292in}{0.898613in}}%
\pgfpathlineto{\pgfqpoint{1.433727in}{0.889305in}}%
\pgfpathclose%
\pgfusepath{fill}%
\end{pgfscope}%
\begin{pgfscope}%
\pgfpathrectangle{\pgfqpoint{0.407592in}{0.222904in}}{\pgfqpoint{1.411163in}{1.411163in}}%
\pgfusepath{clip}%
\pgfsetbuttcap%
\pgfsetroundjoin%
\definecolor{currentfill}{rgb}{0.749020,0.772549,0.792157}%
\pgfsetfillcolor{currentfill}%
\pgfsetfillopacity{0.160000}%
\pgfsetlinewidth{0.000000pt}%
\definecolor{currentstroke}{rgb}{0.000000,0.000000,0.000000}%
\pgfsetstrokecolor{currentstroke}%
\pgfsetdash{}{0pt}%
\pgfpathmoveto{\pgfqpoint{1.246930in}{0.729579in}}%
\pgfpathlineto{\pgfqpoint{1.258376in}{0.738274in}}%
\pgfpathlineto{\pgfqpoint{1.269856in}{0.746995in}}%
\pgfpathlineto{\pgfqpoint{1.277033in}{0.752265in}}%
\pgfpathlineto{\pgfqpoint{1.283360in}{0.757700in}}%
\pgfpathlineto{\pgfqpoint{1.271053in}{0.748321in}}%
\pgfpathlineto{\pgfqpoint{1.258778in}{0.738968in}}%
\pgfpathlineto{\pgfqpoint{1.253223in}{0.734200in}}%
\pgfpathlineto{\pgfqpoint{1.246930in}{0.729579in}}%
\pgfpathclose%
\pgfusepath{fill}%
\end{pgfscope}%
\begin{pgfscope}%
\pgfpathrectangle{\pgfqpoint{0.407592in}{0.222904in}}{\pgfqpoint{1.411163in}{1.411163in}}%
\pgfusepath{clip}%
\pgfsetbuttcap%
\pgfsetroundjoin%
\definecolor{currentfill}{rgb}{0.749020,0.772549,0.792157}%
\pgfsetfillcolor{currentfill}%
\pgfsetfillopacity{0.160000}%
\pgfsetlinewidth{0.000000pt}%
\definecolor{currentstroke}{rgb}{0.000000,0.000000,0.000000}%
\pgfsetstrokecolor{currentstroke}%
\pgfsetdash{}{0pt}%
\pgfpathmoveto{\pgfqpoint{0.993880in}{0.682638in}}%
\pgfpathlineto{\pgfqpoint{0.984048in}{0.689996in}}%
\pgfpathlineto{\pgfqpoint{0.974180in}{0.697382in}}%
\pgfpathlineto{\pgfqpoint{0.981789in}{0.694354in}}%
\pgfpathlineto{\pgfqpoint{0.989891in}{0.691556in}}%
\pgfpathlineto{\pgfqpoint{0.998437in}{0.684656in}}%
\pgfpathlineto{\pgfqpoint{1.006948in}{0.677783in}}%
\pgfpathlineto{\pgfqpoint{1.000211in}{0.680115in}}%
\pgfpathlineto{\pgfqpoint{0.993880in}{0.682638in}}%
\pgfpathclose%
\pgfusepath{fill}%
\end{pgfscope}%
\begin{pgfscope}%
\pgfpathrectangle{\pgfqpoint{0.407592in}{0.222904in}}{\pgfqpoint{1.411163in}{1.411163in}}%
\pgfusepath{clip}%
\pgfsetbuttcap%
\pgfsetroundjoin%
\definecolor{currentfill}{rgb}{0.749020,0.772549,0.792157}%
\pgfsetfillcolor{currentfill}%
\pgfsetfillopacity{0.160000}%
\pgfsetlinewidth{0.000000pt}%
\definecolor{currentstroke}{rgb}{0.000000,0.000000,0.000000}%
\pgfsetstrokecolor{currentstroke}%
\pgfsetdash{}{0pt}%
\pgfpathmoveto{\pgfqpoint{0.823182in}{0.822378in}}%
\pgfpathlineto{\pgfqpoint{0.809575in}{0.833173in}}%
\pgfpathlineto{\pgfqpoint{0.795943in}{0.843987in}}%
\pgfpathlineto{\pgfqpoint{0.798057in}{0.836080in}}%
\pgfpathlineto{\pgfqpoint{0.801469in}{0.828254in}}%
\pgfpathlineto{\pgfqpoint{0.814905in}{0.818123in}}%
\pgfpathlineto{\pgfqpoint{0.828311in}{0.808016in}}%
\pgfpathlineto{\pgfqpoint{0.825158in}{0.815159in}}%
\pgfpathlineto{\pgfqpoint{0.823182in}{0.822378in}}%
\pgfpathclose%
\pgfusepath{fill}%
\end{pgfscope}%
\begin{pgfscope}%
\pgfpathrectangle{\pgfqpoint{0.407592in}{0.222904in}}{\pgfqpoint{1.411163in}{1.411163in}}%
\pgfusepath{clip}%
\pgfsetbuttcap%
\pgfsetroundjoin%
\definecolor{currentfill}{rgb}{0.749020,0.772549,0.792157}%
\pgfsetfillcolor{currentfill}%
\pgfsetfillopacity{0.160000}%
\pgfsetlinewidth{0.000000pt}%
\definecolor{currentstroke}{rgb}{0.000000,0.000000,0.000000}%
\pgfsetstrokecolor{currentstroke}%
\pgfsetdash{}{0pt}%
\pgfpathmoveto{\pgfqpoint{0.915150in}{0.738000in}}%
\pgfpathlineto{\pgfqpoint{0.902396in}{0.747287in}}%
\pgfpathlineto{\pgfqpoint{0.889607in}{0.756597in}}%
\pgfpathlineto{\pgfqpoint{0.894953in}{0.751194in}}%
\pgfpathlineto{\pgfqpoint{0.901178in}{0.745960in}}%
\pgfpathlineto{\pgfqpoint{0.913264in}{0.737302in}}%
\pgfpathlineto{\pgfqpoint{0.925313in}{0.728671in}}%
\pgfpathlineto{\pgfqpoint{0.919849in}{0.733260in}}%
\pgfpathlineto{\pgfqpoint{0.915150in}{0.738000in}}%
\pgfpathclose%
\pgfusepath{fill}%
\end{pgfscope}%
\begin{pgfscope}%
\pgfpathrectangle{\pgfqpoint{0.407592in}{0.222904in}}{\pgfqpoint{1.411163in}{1.411163in}}%
\pgfusepath{clip}%
\pgfsetbuttcap%
\pgfsetroundjoin%
\definecolor{currentfill}{rgb}{0.749020,0.772549,0.792157}%
\pgfsetfillcolor{currentfill}%
\pgfsetfillopacity{0.160000}%
\pgfsetlinewidth{0.000000pt}%
\definecolor{currentstroke}{rgb}{0.000000,0.000000,0.000000}%
\pgfsetstrokecolor{currentstroke}%
\pgfsetdash{}{0pt}%
\pgfpathmoveto{\pgfqpoint{0.612803in}{1.047324in}}%
\pgfpathlineto{\pgfqpoint{0.599309in}{1.059669in}}%
\pgfpathlineto{\pgfqpoint{0.585803in}{1.072026in}}%
\pgfpathlineto{\pgfqpoint{0.579642in}{1.059112in}}%
\pgfpathlineto{\pgfqpoint{0.575662in}{1.046054in}}%
\pgfpathlineto{\pgfqpoint{0.589552in}{1.034297in}}%
\pgfpathlineto{\pgfqpoint{0.603422in}{1.022558in}}%
\pgfpathlineto{\pgfqpoint{0.607079in}{1.035007in}}%
\pgfpathlineto{\pgfqpoint{0.612803in}{1.047324in}}%
\pgfpathclose%
\pgfusepath{fill}%
\end{pgfscope}%
\begin{pgfscope}%
\pgfpathrectangle{\pgfqpoint{0.407592in}{0.222904in}}{\pgfqpoint{1.411163in}{1.411163in}}%
\pgfusepath{clip}%
\pgfsetbuttcap%
\pgfsetroundjoin%
\definecolor{currentfill}{rgb}{0.749020,0.772549,0.792157}%
\pgfsetfillcolor{currentfill}%
\pgfsetfillopacity{0.160000}%
\pgfsetlinewidth{0.000000pt}%
\definecolor{currentstroke}{rgb}{0.000000,0.000000,0.000000}%
\pgfsetstrokecolor{currentstroke}%
\pgfsetdash{}{0pt}%
\pgfpathmoveto{\pgfqpoint{0.741074in}{0.906052in}}%
\pgfpathlineto{\pgfqpoint{0.727395in}{0.917630in}}%
\pgfpathlineto{\pgfqpoint{0.713696in}{0.929224in}}%
\pgfpathlineto{\pgfqpoint{0.712835in}{0.919256in}}%
\pgfpathlineto{\pgfqpoint{0.713628in}{0.909286in}}%
\pgfpathlineto{\pgfqpoint{0.727410in}{0.898353in}}%
\pgfpathlineto{\pgfqpoint{0.741166in}{0.887440in}}%
\pgfpathlineto{\pgfqpoint{0.740354in}{0.896745in}}%
\pgfpathlineto{\pgfqpoint{0.741074in}{0.906052in}}%
\pgfpathclose%
\pgfusepath{fill}%
\end{pgfscope}%
\begin{pgfscope}%
\pgfpathrectangle{\pgfqpoint{0.407592in}{0.222904in}}{\pgfqpoint{1.411163in}{1.411163in}}%
\pgfusepath{clip}%
\pgfsetbuttcap%
\pgfsetroundjoin%
\definecolor{currentfill}{rgb}{0.749020,0.772549,0.792157}%
\pgfsetfillcolor{currentfill}%
\pgfsetfillopacity{0.160000}%
\pgfsetlinewidth{0.000000pt}%
\definecolor{currentstroke}{rgb}{0.000000,0.000000,0.000000}%
\pgfsetstrokecolor{currentstroke}%
\pgfsetdash{}{0pt}%
\pgfpathmoveto{\pgfqpoint{1.157223in}{0.676065in}}%
\pgfpathlineto{\pgfqpoint{1.164074in}{0.682765in}}%
\pgfpathlineto{\pgfqpoint{1.170952in}{0.689494in}}%
\pgfpathlineto{\pgfqpoint{1.179972in}{0.692099in}}%
\pgfpathlineto{\pgfqpoint{1.188578in}{0.694944in}}%
\pgfpathlineto{\pgfqpoint{1.180209in}{0.687761in}}%
\pgfpathlineto{\pgfqpoint{1.171873in}{0.680606in}}%
\pgfpathlineto{\pgfqpoint{1.164718in}{0.678235in}}%
\pgfpathlineto{\pgfqpoint{1.157223in}{0.676065in}}%
\pgfpathclose%
\pgfusepath{fill}%
\end{pgfscope}%
\begin{pgfscope}%
\pgfpathrectangle{\pgfqpoint{0.407592in}{0.222904in}}{\pgfqpoint{1.411163in}{1.411163in}}%
\pgfusepath{clip}%
\pgfsetbuttcap%
\pgfsetroundjoin%
\definecolor{currentfill}{rgb}{0.749020,0.772549,0.792157}%
\pgfsetfillcolor{currentfill}%
\pgfsetfillopacity{0.160000}%
\pgfsetlinewidth{0.000000pt}%
\definecolor{currentstroke}{rgb}{0.000000,0.000000,0.000000}%
\pgfsetstrokecolor{currentstroke}%
\pgfsetdash{}{0pt}%
\pgfpathmoveto{\pgfqpoint{1.204397in}{0.701300in}}%
\pgfpathlineto{\pgfqpoint{1.214129in}{0.709037in}}%
\pgfpathlineto{\pgfqpoint{1.223897in}{0.716803in}}%
\pgfpathlineto{\pgfqpoint{1.232236in}{0.720859in}}%
\pgfpathlineto{\pgfqpoint{1.239925in}{0.725125in}}%
\pgfpathlineto{\pgfqpoint{1.229011in}{0.716770in}}%
\pgfpathlineto{\pgfqpoint{1.218131in}{0.708442in}}%
\pgfpathlineto{\pgfqpoint{1.211540in}{0.704780in}}%
\pgfpathlineto{\pgfqpoint{1.204397in}{0.701300in}}%
\pgfpathclose%
\pgfusepath{fill}%
\end{pgfscope}%
\begin{pgfscope}%
\pgfpathrectangle{\pgfqpoint{0.407592in}{0.222904in}}{\pgfqpoint{1.411163in}{1.411163in}}%
\pgfusepath{clip}%
\pgfsetbuttcap%
\pgfsetroundjoin%
\definecolor{currentfill}{rgb}{0.749020,0.772549,0.792157}%
\pgfsetfillcolor{currentfill}%
\pgfsetfillopacity{0.160000}%
\pgfsetlinewidth{0.000000pt}%
\definecolor{currentstroke}{rgb}{0.000000,0.000000,0.000000}%
\pgfsetstrokecolor{currentstroke}%
\pgfsetdash{}{0pt}%
\pgfpathmoveto{\pgfqpoint{1.601625in}{1.048682in}}%
\pgfpathlineto{\pgfqpoint{1.615437in}{1.060516in}}%
\pgfpathlineto{\pgfqpoint{1.629267in}{1.072365in}}%
\pgfpathlineto{\pgfqpoint{1.627667in}{1.085997in}}%
\pgfpathlineto{\pgfqpoint{1.623745in}{1.099460in}}%
\pgfpathlineto{\pgfqpoint{1.610173in}{1.087022in}}%
\pgfpathlineto{\pgfqpoint{1.596613in}{1.074595in}}%
\pgfpathlineto{\pgfqpoint{1.600222in}{1.061717in}}%
\pgfpathlineto{\pgfqpoint{1.601625in}{1.048682in}}%
\pgfpathclose%
\pgfusepath{fill}%
\end{pgfscope}%
\begin{pgfscope}%
\pgfpathrectangle{\pgfqpoint{0.407592in}{0.222904in}}{\pgfqpoint{1.411163in}{1.411163in}}%
\pgfusepath{clip}%
\pgfsetbuttcap%
\pgfsetroundjoin%
\definecolor{currentfill}{rgb}{0.749020,0.772549,0.792157}%
\pgfsetfillcolor{currentfill}%
\pgfsetfillopacity{0.160000}%
\pgfsetlinewidth{0.000000pt}%
\definecolor{currentstroke}{rgb}{0.000000,0.000000,0.000000}%
\pgfsetstrokecolor{currentstroke}%
\pgfsetdash{}{0pt}%
\pgfpathmoveto{\pgfqpoint{1.460810in}{0.911284in}}%
\pgfpathlineto{\pgfqpoint{1.474388in}{0.922304in}}%
\pgfpathlineto{\pgfqpoint{1.487990in}{0.933343in}}%
\pgfpathlineto{\pgfqpoint{1.490781in}{0.943972in}}%
\pgfpathlineto{\pgfqpoint{1.491799in}{0.954586in}}%
\pgfpathlineto{\pgfqpoint{1.478153in}{0.942895in}}%
\pgfpathlineto{\pgfqpoint{1.464526in}{0.931220in}}%
\pgfpathlineto{\pgfqpoint{1.463494in}{0.921257in}}%
\pgfpathlineto{\pgfqpoint{1.460810in}{0.911284in}}%
\pgfpathclose%
\pgfusepath{fill}%
\end{pgfscope}%
\begin{pgfscope}%
\pgfpathrectangle{\pgfqpoint{0.407592in}{0.222904in}}{\pgfqpoint{1.411163in}{1.411163in}}%
\pgfusepath{clip}%
\pgfsetbuttcap%
\pgfsetroundjoin%
\definecolor{currentfill}{rgb}{0.749020,0.772549,0.792157}%
\pgfsetfillcolor{currentfill}%
\pgfsetfillopacity{0.160000}%
\pgfsetlinewidth{0.000000pt}%
\definecolor{currentstroke}{rgb}{0.000000,0.000000,0.000000}%
\pgfsetstrokecolor{currentstroke}%
\pgfsetdash{}{0pt}%
\pgfpathmoveto{\pgfqpoint{1.006948in}{0.677783in}}%
\pgfpathlineto{\pgfqpoint{0.998437in}{0.684656in}}%
\pgfpathlineto{\pgfqpoint{0.989891in}{0.691556in}}%
\pgfpathlineto{\pgfqpoint{0.998451in}{0.689002in}}%
\pgfpathlineto{\pgfqpoint{1.007429in}{0.686702in}}%
\pgfpathlineto{\pgfqpoint{1.014490in}{0.680207in}}%
\pgfpathlineto{\pgfqpoint{1.021521in}{0.673739in}}%
\pgfpathlineto{\pgfqpoint{1.014062in}{0.675654in}}%
\pgfpathlineto{\pgfqpoint{1.006948in}{0.677783in}}%
\pgfpathclose%
\pgfusepath{fill}%
\end{pgfscope}%
\begin{pgfscope}%
\pgfpathrectangle{\pgfqpoint{0.407592in}{0.222904in}}{\pgfqpoint{1.411163in}{1.411163in}}%
\pgfusepath{clip}%
\pgfsetbuttcap%
\pgfsetroundjoin%
\definecolor{currentfill}{rgb}{0.749020,0.772549,0.792157}%
\pgfsetfillcolor{currentfill}%
\pgfsetfillopacity{0.160000}%
\pgfsetlinewidth{0.000000pt}%
\definecolor{currentstroke}{rgb}{0.000000,0.000000,0.000000}%
\pgfsetstrokecolor{currentstroke}%
\pgfsetdash{}{0pt}%
\pgfpathmoveto{\pgfqpoint{0.585803in}{1.072026in}}%
\pgfpathlineto{\pgfqpoint{0.572283in}{1.084395in}}%
\pgfpathlineto{\pgfqpoint{0.558751in}{1.096775in}}%
\pgfpathlineto{\pgfqpoint{0.552141in}{1.083274in}}%
\pgfpathlineto{\pgfqpoint{0.547824in}{1.069615in}}%
\pgfpathlineto{\pgfqpoint{0.561752in}{1.057826in}}%
\pgfpathlineto{\pgfqpoint{0.575662in}{1.046054in}}%
\pgfpathlineto{\pgfqpoint{0.579642in}{1.059112in}}%
\pgfpathlineto{\pgfqpoint{0.585803in}{1.072026in}}%
\pgfpathclose%
\pgfusepath{fill}%
\end{pgfscope}%
\begin{pgfscope}%
\pgfpathrectangle{\pgfqpoint{0.407592in}{0.222904in}}{\pgfqpoint{1.411163in}{1.411163in}}%
\pgfusepath{clip}%
\pgfsetbuttcap%
\pgfsetroundjoin%
\definecolor{currentfill}{rgb}{0.749020,0.772549,0.792157}%
\pgfsetfillcolor{currentfill}%
\pgfsetfillopacity{0.160000}%
\pgfsetlinewidth{0.000000pt}%
\definecolor{currentstroke}{rgb}{0.000000,0.000000,0.000000}%
\pgfsetstrokecolor{currentstroke}%
\pgfsetdash{}{0pt}%
\pgfpathmoveto{\pgfqpoint{1.371404in}{0.829816in}}%
\pgfpathlineto{\pgfqpoint{1.384511in}{0.840037in}}%
\pgfpathlineto{\pgfqpoint{1.397646in}{0.850281in}}%
\pgfpathlineto{\pgfqpoint{1.402896in}{0.858804in}}%
\pgfpathlineto{\pgfqpoint{1.406741in}{0.867403in}}%
\pgfpathlineto{\pgfqpoint{1.393284in}{0.856482in}}%
\pgfpathlineto{\pgfqpoint{1.379850in}{0.845580in}}%
\pgfpathlineto{\pgfqpoint{1.376269in}{0.837661in}}%
\pgfpathlineto{\pgfqpoint{1.371404in}{0.829816in}}%
\pgfpathclose%
\pgfusepath{fill}%
\end{pgfscope}%
\begin{pgfscope}%
\pgfpathrectangle{\pgfqpoint{0.407592in}{0.222904in}}{\pgfqpoint{1.411163in}{1.411163in}}%
\pgfusepath{clip}%
\pgfsetbuttcap%
\pgfsetroundjoin%
\definecolor{currentfill}{rgb}{0.749020,0.772549,0.792157}%
\pgfsetfillcolor{currentfill}%
\pgfsetfillopacity{0.160000}%
\pgfsetlinewidth{0.000000pt}%
\definecolor{currentstroke}{rgb}{0.000000,0.000000,0.000000}%
\pgfsetstrokecolor{currentstroke}%
\pgfsetdash{}{0pt}%
\pgfpathmoveto{\pgfqpoint{1.308072in}{0.776530in}}%
\pgfpathlineto{\pgfqpoint{1.320476in}{0.785982in}}%
\pgfpathlineto{\pgfqpoint{1.332913in}{0.795459in}}%
\pgfpathlineto{\pgfqpoint{1.339661in}{0.802383in}}%
\pgfpathlineto{\pgfqpoint{1.345279in}{0.809441in}}%
\pgfpathlineto{\pgfqpoint{1.332259in}{0.799287in}}%
\pgfpathlineto{\pgfqpoint{1.319268in}{0.789156in}}%
\pgfpathlineto{\pgfqpoint{1.314176in}{0.782781in}}%
\pgfpathlineto{\pgfqpoint{1.308072in}{0.776530in}}%
\pgfpathclose%
\pgfusepath{fill}%
\end{pgfscope}%
\begin{pgfscope}%
\pgfpathrectangle{\pgfqpoint{0.407592in}{0.222904in}}{\pgfqpoint{1.411163in}{1.411163in}}%
\pgfusepath{clip}%
\pgfsetbuttcap%
\pgfsetroundjoin%
\definecolor{currentfill}{rgb}{0.749020,0.772549,0.792157}%
\pgfsetfillcolor{currentfill}%
\pgfsetfillopacity{0.160000}%
\pgfsetlinewidth{0.000000pt}%
\definecolor{currentstroke}{rgb}{0.000000,0.000000,0.000000}%
\pgfsetstrokecolor{currentstroke}%
\pgfsetdash{}{0pt}%
\pgfpathmoveto{\pgfqpoint{1.629267in}{1.072365in}}%
\pgfpathlineto{\pgfqpoint{1.643116in}{1.084231in}}%
\pgfpathlineto{\pgfqpoint{1.656983in}{1.096112in}}%
\pgfpathlineto{\pgfqpoint{1.655174in}{1.110332in}}%
\pgfpathlineto{\pgfqpoint{1.650927in}{1.124370in}}%
\pgfpathlineto{\pgfqpoint{1.637330in}{1.111909in}}%
\pgfpathlineto{\pgfqpoint{1.623745in}{1.099460in}}%
\pgfpathlineto{\pgfqpoint{1.627667in}{1.085997in}}%
\pgfpathlineto{\pgfqpoint{1.629267in}{1.072365in}}%
\pgfpathclose%
\pgfusepath{fill}%
\end{pgfscope}%
\begin{pgfscope}%
\pgfpathrectangle{\pgfqpoint{0.407592in}{0.222904in}}{\pgfqpoint{1.411163in}{1.411163in}}%
\pgfusepath{clip}%
\pgfsetbuttcap%
\pgfsetroundjoin%
\definecolor{currentfill}{rgb}{0.749020,0.772549,0.792157}%
\pgfsetfillcolor{currentfill}%
\pgfsetfillopacity{0.160000}%
\pgfsetlinewidth{0.000000pt}%
\definecolor{currentstroke}{rgb}{0.000000,0.000000,0.000000}%
\pgfsetstrokecolor{currentstroke}%
\pgfsetdash{}{0pt}%
\pgfpathmoveto{\pgfqpoint{1.141347in}{0.672367in}}%
\pgfpathlineto{\pgfqpoint{1.146580in}{0.678697in}}%
\pgfpathlineto{\pgfqpoint{1.151835in}{0.685055in}}%
\pgfpathlineto{\pgfqpoint{1.161559in}{0.687142in}}%
\pgfpathlineto{\pgfqpoint{1.170952in}{0.689494in}}%
\pgfpathlineto{\pgfqpoint{1.164074in}{0.682765in}}%
\pgfpathlineto{\pgfqpoint{1.157223in}{0.676065in}}%
\pgfpathlineto{\pgfqpoint{1.149421in}{0.674106in}}%
\pgfpathlineto{\pgfqpoint{1.141347in}{0.672367in}}%
\pgfpathclose%
\pgfusepath{fill}%
\end{pgfscope}%
\begin{pgfscope}%
\pgfpathrectangle{\pgfqpoint{0.407592in}{0.222904in}}{\pgfqpoint{1.411163in}{1.411163in}}%
\pgfusepath{clip}%
\pgfsetbuttcap%
\pgfsetroundjoin%
\definecolor{currentfill}{rgb}{0.749020,0.772549,0.792157}%
\pgfsetfillcolor{currentfill}%
\pgfsetfillopacity{0.160000}%
\pgfsetlinewidth{0.000000pt}%
\definecolor{currentstroke}{rgb}{0.000000,0.000000,0.000000}%
\pgfsetstrokecolor{currentstroke}%
\pgfsetdash{}{0pt}%
\pgfpathmoveto{\pgfqpoint{0.960572in}{0.704067in}}%
\pgfpathlineto{\pgfqpoint{0.949528in}{0.712034in}}%
\pgfpathlineto{\pgfqpoint{0.938446in}{0.720028in}}%
\pgfpathlineto{\pgfqpoint{0.946061in}{0.716016in}}%
\pgfpathlineto{\pgfqpoint{0.954332in}{0.712235in}}%
\pgfpathlineto{\pgfqpoint{0.964274in}{0.704795in}}%
\pgfpathlineto{\pgfqpoint{0.974180in}{0.697382in}}%
\pgfpathlineto{\pgfqpoint{0.967097in}{0.700624in}}%
\pgfpathlineto{\pgfqpoint{0.960572in}{0.704067in}}%
\pgfpathclose%
\pgfusepath{fill}%
\end{pgfscope}%
\begin{pgfscope}%
\pgfpathrectangle{\pgfqpoint{0.407592in}{0.222904in}}{\pgfqpoint{1.411163in}{1.411163in}}%
\pgfusepath{clip}%
\pgfsetbuttcap%
\pgfsetroundjoin%
\definecolor{currentfill}{rgb}{0.749020,0.772549,0.792157}%
\pgfsetfillcolor{currentfill}%
\pgfsetfillopacity{0.160000}%
\pgfsetlinewidth{0.000000pt}%
\definecolor{currentstroke}{rgb}{0.000000,0.000000,0.000000}%
\pgfsetstrokecolor{currentstroke}%
\pgfsetdash{}{0pt}%
\pgfpathmoveto{\pgfqpoint{1.239925in}{0.725125in}}%
\pgfpathlineto{\pgfqpoint{1.250874in}{0.733506in}}%
\pgfpathlineto{\pgfqpoint{1.261858in}{0.741915in}}%
\pgfpathlineto{\pgfqpoint{1.269856in}{0.746995in}}%
\pgfpathlineto{\pgfqpoint{1.258376in}{0.738274in}}%
\pgfpathlineto{\pgfqpoint{1.246930in}{0.729579in}}%
\pgfusepath{fill}%
\end{pgfscope}%
\begin{pgfscope}%
\pgfpathrectangle{\pgfqpoint{0.407592in}{0.222904in}}{\pgfqpoint{1.411163in}{1.411163in}}%
\pgfusepath{clip}%
\pgfsetbuttcap%
\pgfsetroundjoin%
\definecolor{currentfill}{rgb}{0.749020,0.772549,0.792157}%
\pgfsetfillcolor{currentfill}%
\pgfsetfillopacity{0.160000}%
\pgfsetlinewidth{0.000000pt}%
\definecolor{currentstroke}{rgb}{0.000000,0.000000,0.000000}%
\pgfsetstrokecolor{currentstroke}%
\pgfsetdash{}{0pt}%
\pgfpathmoveto{\pgfqpoint{0.713696in}{0.929224in}}%
\pgfpathlineto{\pgfqpoint{0.699979in}{0.940834in}}%
\pgfpathlineto{\pgfqpoint{0.686242in}{0.952460in}}%
\pgfpathlineto{\pgfqpoint{0.685228in}{0.941840in}}%
\pgfpathlineto{\pgfqpoint{0.685989in}{0.931212in}}%
\pgfpathlineto{\pgfqpoint{0.699821in}{0.920239in}}%
\pgfpathlineto{\pgfqpoint{0.713628in}{0.909286in}}%
\pgfpathlineto{\pgfqpoint{0.712835in}{0.919256in}}%
\pgfpathlineto{\pgfqpoint{0.713696in}{0.929224in}}%
\pgfpathclose%
\pgfusepath{fill}%
\end{pgfscope}%
\begin{pgfscope}%
\pgfpathrectangle{\pgfqpoint{0.407592in}{0.222904in}}{\pgfqpoint{1.411163in}{1.411163in}}%
\pgfusepath{clip}%
\pgfsetbuttcap%
\pgfsetroundjoin%
\definecolor{currentfill}{rgb}{0.749020,0.772549,0.792157}%
\pgfsetfillcolor{currentfill}%
\pgfsetfillopacity{0.160000}%
\pgfsetlinewidth{0.000000pt}%
\definecolor{currentstroke}{rgb}{0.000000,0.000000,0.000000}%
\pgfsetstrokecolor{currentstroke}%
\pgfsetdash{}{0pt}%
\pgfpathmoveto{\pgfqpoint{0.855032in}{0.787868in}}%
\pgfpathlineto{\pgfqpoint{0.841686in}{0.797931in}}%
\pgfpathlineto{\pgfqpoint{0.828311in}{0.808016in}}%
\pgfpathlineto{\pgfqpoint{0.832633in}{0.800982in}}%
\pgfpathlineto{\pgfqpoint{0.838112in}{0.794089in}}%
\pgfpathlineto{\pgfqpoint{0.851037in}{0.784679in}}%
\pgfpathlineto{\pgfqpoint{0.863928in}{0.775293in}}%
\pgfpathlineto{\pgfqpoint{0.858962in}{0.781516in}}%
\pgfpathlineto{\pgfqpoint{0.855032in}{0.787868in}}%
\pgfpathclose%
\pgfusepath{fill}%
\end{pgfscope}%
\begin{pgfscope}%
\pgfpathrectangle{\pgfqpoint{0.407592in}{0.222904in}}{\pgfqpoint{1.411163in}{1.411163in}}%
\pgfusepath{clip}%
\pgfsetbuttcap%
\pgfsetroundjoin%
\definecolor{currentfill}{rgb}{0.749020,0.772549,0.792157}%
\pgfsetfillcolor{currentfill}%
\pgfsetfillopacity{0.160000}%
\pgfsetlinewidth{0.000000pt}%
\definecolor{currentstroke}{rgb}{0.000000,0.000000,0.000000}%
\pgfsetstrokecolor{currentstroke}%
\pgfsetdash{}{0pt}%
\pgfpathmoveto{\pgfqpoint{0.795943in}{0.843987in}}%
\pgfpathlineto{\pgfqpoint{0.782286in}{0.854821in}}%
\pgfpathlineto{\pgfqpoint{0.768605in}{0.865674in}}%
\pgfpathlineto{\pgfqpoint{0.770845in}{0.857087in}}%
\pgfpathlineto{\pgfqpoint{0.774506in}{0.848584in}}%
\pgfpathlineto{\pgfqpoint{0.788003in}{0.838407in}}%
\pgfpathlineto{\pgfqpoint{0.801469in}{0.828254in}}%
\pgfpathlineto{\pgfqpoint{0.798057in}{0.836080in}}%
\pgfpathlineto{\pgfqpoint{0.795943in}{0.843987in}}%
\pgfpathclose%
\pgfusepath{fill}%
\end{pgfscope}%
\begin{pgfscope}%
\pgfpathrectangle{\pgfqpoint{0.407592in}{0.222904in}}{\pgfqpoint{1.411163in}{1.411163in}}%
\pgfusepath{clip}%
\pgfsetbuttcap%
\pgfsetroundjoin%
\definecolor{currentfill}{rgb}{0.749020,0.772549,0.792157}%
\pgfsetfillcolor{currentfill}%
\pgfsetfillopacity{0.160000}%
\pgfsetlinewidth{0.000000pt}%
\definecolor{currentstroke}{rgb}{0.000000,0.000000,0.000000}%
\pgfsetstrokecolor{currentstroke}%
\pgfsetdash{}{0pt}%
\pgfpathmoveto{\pgfqpoint{1.021521in}{0.673739in}}%
\pgfpathlineto{\pgfqpoint{1.014490in}{0.680207in}}%
\pgfpathlineto{\pgfqpoint{1.007429in}{0.686702in}}%
\pgfpathlineto{\pgfqpoint{1.016786in}{0.684669in}}%
\pgfpathlineto{\pgfqpoint{1.026478in}{0.682912in}}%
\pgfpathlineto{\pgfqpoint{1.031920in}{0.676734in}}%
\pgfpathlineto{\pgfqpoint{1.037338in}{0.670583in}}%
\pgfpathlineto{\pgfqpoint{1.029291in}{0.672046in}}%
\pgfpathlineto{\pgfqpoint{1.021521in}{0.673739in}}%
\pgfpathclose%
\pgfusepath{fill}%
\end{pgfscope}%
\begin{pgfscope}%
\pgfpathrectangle{\pgfqpoint{0.407592in}{0.222904in}}{\pgfqpoint{1.411163in}{1.411163in}}%
\pgfusepath{clip}%
\pgfsetbuttcap%
\pgfsetroundjoin%
\definecolor{currentfill}{rgb}{0.749020,0.772549,0.792157}%
\pgfsetfillcolor{currentfill}%
\pgfsetfillopacity{0.160000}%
\pgfsetlinewidth{0.000000pt}%
\definecolor{currentstroke}{rgb}{0.000000,0.000000,0.000000}%
\pgfsetstrokecolor{currentstroke}%
\pgfsetdash{}{0pt}%
\pgfpathmoveto{\pgfqpoint{1.269856in}{0.746995in}}%
\pgfpathlineto{\pgfqpoint{1.281371in}{0.755743in}}%
\pgfpathlineto{\pgfqpoint{1.292921in}{0.764517in}}%
\pgfpathlineto{\pgfqpoint{1.300978in}{0.770432in}}%
\pgfpathlineto{\pgfqpoint{1.308072in}{0.776530in}}%
\pgfpathlineto{\pgfqpoint{1.295700in}{0.767102in}}%
\pgfpathlineto{\pgfqpoint{1.283360in}{0.757700in}}%
\pgfpathlineto{\pgfqpoint{1.277033in}{0.752265in}}%
\pgfpathlineto{\pgfqpoint{1.269856in}{0.746995in}}%
\pgfpathclose%
\pgfusepath{fill}%
\end{pgfscope}%
\begin{pgfscope}%
\pgfpathrectangle{\pgfqpoint{0.407592in}{0.222904in}}{\pgfqpoint{1.411163in}{1.411163in}}%
\pgfusepath{clip}%
\pgfsetbuttcap%
\pgfsetroundjoin%
\definecolor{currentfill}{rgb}{0.749020,0.772549,0.792157}%
\pgfsetfillcolor{currentfill}%
\pgfsetfillopacity{0.160000}%
\pgfsetlinewidth{0.000000pt}%
\definecolor{currentstroke}{rgb}{0.000000,0.000000,0.000000}%
\pgfsetstrokecolor{currentstroke}%
\pgfsetdash{}{0pt}%
\pgfpathmoveto{\pgfqpoint{0.558751in}{1.096775in}}%
\pgfpathlineto{\pgfqpoint{0.545207in}{1.109167in}}%
\pgfpathlineto{\pgfqpoint{0.531649in}{1.121571in}}%
\pgfpathlineto{\pgfqpoint{0.524575in}{1.107492in}}%
\pgfpathlineto{\pgfqpoint{0.519908in}{1.093242in}}%
\pgfpathlineto{\pgfqpoint{0.533876in}{1.081420in}}%
\pgfpathlineto{\pgfqpoint{0.547824in}{1.069615in}}%
\pgfpathlineto{\pgfqpoint{0.552141in}{1.083274in}}%
\pgfpathlineto{\pgfqpoint{0.558751in}{1.096775in}}%
\pgfpathclose%
\pgfusepath{fill}%
\end{pgfscope}%
\begin{pgfscope}%
\pgfpathrectangle{\pgfqpoint{0.407592in}{0.222904in}}{\pgfqpoint{1.411163in}{1.411163in}}%
\pgfusepath{clip}%
\pgfsetbuttcap%
\pgfsetroundjoin%
\definecolor{currentfill}{rgb}{0.749020,0.772549,0.792157}%
\pgfsetfillcolor{currentfill}%
\pgfsetfillopacity{0.160000}%
\pgfsetlinewidth{0.000000pt}%
\definecolor{currentstroke}{rgb}{0.000000,0.000000,0.000000}%
\pgfsetstrokecolor{currentstroke}%
\pgfsetdash{}{0pt}%
\pgfpathmoveto{\pgfqpoint{1.188578in}{0.694944in}}%
\pgfpathlineto{\pgfqpoint{1.196979in}{0.702154in}}%
\pgfpathlineto{\pgfqpoint{1.205412in}{0.709392in}}%
\pgfpathlineto{\pgfqpoint{1.214942in}{0.712974in}}%
\pgfpathlineto{\pgfqpoint{1.223897in}{0.716803in}}%
\pgfpathlineto{\pgfqpoint{1.214129in}{0.709037in}}%
\pgfpathlineto{\pgfqpoint{1.204397in}{0.701300in}}%
\pgfpathlineto{\pgfqpoint{1.196731in}{0.698015in}}%
\pgfpathlineto{\pgfqpoint{1.188578in}{0.694944in}}%
\pgfpathclose%
\pgfusepath{fill}%
\end{pgfscope}%
\begin{pgfscope}%
\pgfpathrectangle{\pgfqpoint{0.407592in}{0.222904in}}{\pgfqpoint{1.411163in}{1.411163in}}%
\pgfusepath{clip}%
\pgfsetbuttcap%
\pgfsetroundjoin%
\definecolor{currentfill}{rgb}{0.749020,0.772549,0.792157}%
\pgfsetfillcolor{currentfill}%
\pgfsetfillopacity{0.160000}%
\pgfsetlinewidth{0.000000pt}%
\definecolor{currentstroke}{rgb}{0.000000,0.000000,0.000000}%
\pgfsetstrokecolor{currentstroke}%
\pgfsetdash{}{0pt}%
\pgfpathmoveto{\pgfqpoint{1.124529in}{0.669584in}}%
\pgfpathlineto{\pgfqpoint{1.128043in}{0.675635in}}%
\pgfpathlineto{\pgfqpoint{1.131573in}{0.681712in}}%
\pgfpathlineto{\pgfqpoint{1.141825in}{0.683242in}}%
\pgfpathlineto{\pgfqpoint{1.151835in}{0.685055in}}%
\pgfpathlineto{\pgfqpoint{1.146580in}{0.678697in}}%
\pgfpathlineto{\pgfqpoint{1.141347in}{0.672367in}}%
\pgfpathlineto{\pgfqpoint{1.133037in}{0.670858in}}%
\pgfpathlineto{\pgfqpoint{1.124529in}{0.669584in}}%
\pgfpathclose%
\pgfusepath{fill}%
\end{pgfscope}%
\begin{pgfscope}%
\pgfpathrectangle{\pgfqpoint{0.407592in}{0.222904in}}{\pgfqpoint{1.411163in}{1.411163in}}%
\pgfusepath{clip}%
\pgfsetbuttcap%
\pgfsetroundjoin%
\definecolor{currentfill}{rgb}{0.749020,0.772549,0.792157}%
\pgfsetfillcolor{currentfill}%
\pgfsetfillopacity{0.160000}%
\pgfsetlinewidth{0.000000pt}%
\definecolor{currentstroke}{rgb}{0.000000,0.000000,0.000000}%
\pgfsetstrokecolor{currentstroke}%
\pgfsetdash{}{0pt}%
\pgfpathmoveto{\pgfqpoint{1.656983in}{1.096112in}}%
\pgfpathlineto{\pgfqpoint{1.670870in}{1.108009in}}%
\pgfpathlineto{\pgfqpoint{1.684775in}{1.119923in}}%
\pgfpathlineto{\pgfqpoint{1.682743in}{1.134722in}}%
\pgfpathlineto{\pgfqpoint{1.678158in}{1.149325in}}%
\pgfpathlineto{\pgfqpoint{1.664536in}{1.136842in}}%
\pgfpathlineto{\pgfqpoint{1.650927in}{1.124370in}}%
\pgfpathlineto{\pgfqpoint{1.655174in}{1.110332in}}%
\pgfpathlineto{\pgfqpoint{1.656983in}{1.096112in}}%
\pgfpathclose%
\pgfusepath{fill}%
\end{pgfscope}%
\begin{pgfscope}%
\pgfpathrectangle{\pgfqpoint{0.407592in}{0.222904in}}{\pgfqpoint{1.411163in}{1.411163in}}%
\pgfusepath{clip}%
\pgfsetbuttcap%
\pgfsetroundjoin%
\definecolor{currentfill}{rgb}{0.749020,0.772549,0.792157}%
\pgfsetfillcolor{currentfill}%
\pgfsetfillopacity{0.160000}%
\pgfsetlinewidth{0.000000pt}%
\definecolor{currentstroke}{rgb}{0.000000,0.000000,0.000000}%
\pgfsetstrokecolor{currentstroke}%
\pgfsetdash{}{0pt}%
\pgfpathmoveto{\pgfqpoint{0.925313in}{0.728671in}}%
\pgfpathlineto{\pgfqpoint{0.913264in}{0.737302in}}%
\pgfpathlineto{\pgfqpoint{0.901178in}{0.745960in}}%
\pgfpathlineto{\pgfqpoint{0.908260in}{0.740920in}}%
\pgfpathlineto{\pgfqpoint{0.916170in}{0.736098in}}%
\pgfpathlineto{\pgfqpoint{0.927327in}{0.728050in}}%
\pgfpathlineto{\pgfqpoint{0.938446in}{0.720028in}}%
\pgfpathlineto{\pgfqpoint{0.931521in}{0.724253in}}%
\pgfpathlineto{\pgfqpoint{0.925313in}{0.728671in}}%
\pgfpathclose%
\pgfusepath{fill}%
\end{pgfscope}%
\begin{pgfscope}%
\pgfpathrectangle{\pgfqpoint{0.407592in}{0.222904in}}{\pgfqpoint{1.411163in}{1.411163in}}%
\pgfusepath{clip}%
\pgfsetbuttcap%
\pgfsetroundjoin%
\definecolor{currentfill}{rgb}{0.749020,0.772549,0.792157}%
\pgfsetfillcolor{currentfill}%
\pgfsetfillopacity{0.160000}%
\pgfsetlinewidth{0.000000pt}%
\definecolor{currentstroke}{rgb}{0.000000,0.000000,0.000000}%
\pgfsetstrokecolor{currentstroke}%
\pgfsetdash{}{0pt}%
\pgfpathmoveto{\pgfqpoint{1.487990in}{0.933343in}}%
\pgfpathlineto{\pgfqpoint{1.501616in}{0.944401in}}%
\pgfpathlineto{\pgfqpoint{1.515267in}{0.955480in}}%
\pgfpathlineto{\pgfqpoint{1.518154in}{0.966759in}}%
\pgfpathlineto{\pgfqpoint{1.519145in}{0.978016in}}%
\pgfpathlineto{\pgfqpoint{1.505463in}{0.966293in}}%
\pgfpathlineto{\pgfqpoint{1.491799in}{0.954586in}}%
\pgfpathlineto{\pgfqpoint{1.490781in}{0.943972in}}%
\pgfpathlineto{\pgfqpoint{1.487990in}{0.933343in}}%
\pgfpathclose%
\pgfusepath{fill}%
\end{pgfscope}%
\begin{pgfscope}%
\pgfpathrectangle{\pgfqpoint{0.407592in}{0.222904in}}{\pgfqpoint{1.411163in}{1.411163in}}%
\pgfusepath{clip}%
\pgfsetbuttcap%
\pgfsetroundjoin%
\definecolor{currentfill}{rgb}{0.749020,0.772549,0.792157}%
\pgfsetfillcolor{currentfill}%
\pgfsetfillopacity{0.160000}%
\pgfsetlinewidth{0.000000pt}%
\definecolor{currentstroke}{rgb}{0.000000,0.000000,0.000000}%
\pgfsetstrokecolor{currentstroke}%
\pgfsetdash{}{0pt}%
\pgfpathmoveto{\pgfqpoint{0.889607in}{0.756597in}}%
\pgfpathlineto{\pgfqpoint{0.876785in}{0.765933in}}%
\pgfpathlineto{\pgfqpoint{0.863928in}{0.775293in}}%
\pgfpathlineto{\pgfqpoint{0.869913in}{0.769230in}}%
\pgfpathlineto{\pgfqpoint{0.876896in}{0.763355in}}%
\pgfpathlineto{\pgfqpoint{0.889056in}{0.754644in}}%
\pgfpathlineto{\pgfqpoint{0.901178in}{0.745960in}}%
\pgfpathlineto{\pgfqpoint{0.894953in}{0.751194in}}%
\pgfpathlineto{\pgfqpoint{0.889607in}{0.756597in}}%
\pgfpathclose%
\pgfusepath{fill}%
\end{pgfscope}%
\begin{pgfscope}%
\pgfpathrectangle{\pgfqpoint{0.407592in}{0.222904in}}{\pgfqpoint{1.411163in}{1.411163in}}%
\pgfusepath{clip}%
\pgfsetbuttcap%
\pgfsetroundjoin%
\definecolor{currentfill}{rgb}{0.749020,0.772549,0.792157}%
\pgfsetfillcolor{currentfill}%
\pgfsetfillopacity{0.160000}%
\pgfsetlinewidth{0.000000pt}%
\definecolor{currentstroke}{rgb}{0.000000,0.000000,0.000000}%
\pgfsetstrokecolor{currentstroke}%
\pgfsetdash{}{0pt}%
\pgfpathmoveto{\pgfqpoint{1.037338in}{0.670583in}}%
\pgfpathlineto{\pgfqpoint{1.031920in}{0.676734in}}%
\pgfpathlineto{\pgfqpoint{1.026478in}{0.682912in}}%
\pgfpathlineto{\pgfqpoint{1.036462in}{0.681440in}}%
\pgfpathlineto{\pgfqpoint{1.046692in}{0.680260in}}%
\pgfpathlineto{\pgfqpoint{1.050411in}{0.674304in}}%
\pgfpathlineto{\pgfqpoint{1.054113in}{0.668375in}}%
\pgfpathlineto{\pgfqpoint{1.045624in}{0.669358in}}%
\pgfpathlineto{\pgfqpoint{1.037338in}{0.670583in}}%
\pgfpathclose%
\pgfusepath{fill}%
\end{pgfscope}%
\begin{pgfscope}%
\pgfpathrectangle{\pgfqpoint{0.407592in}{0.222904in}}{\pgfqpoint{1.411163in}{1.411163in}}%
\pgfusepath{clip}%
\pgfsetbuttcap%
\pgfsetroundjoin%
\definecolor{currentfill}{rgb}{0.749020,0.772549,0.792157}%
\pgfsetfillcolor{currentfill}%
\pgfsetfillopacity{0.160000}%
\pgfsetlinewidth{0.000000pt}%
\definecolor{currentstroke}{rgb}{0.000000,0.000000,0.000000}%
\pgfsetstrokecolor{currentstroke}%
\pgfsetdash{}{0pt}%
\pgfpathmoveto{\pgfqpoint{1.107075in}{0.667768in}}%
\pgfpathlineto{\pgfqpoint{1.108802in}{0.673636in}}%
\pgfpathlineto{\pgfqpoint{1.110537in}{0.679531in}}%
\pgfpathlineto{\pgfqpoint{1.121128in}{0.680473in}}%
\pgfpathlineto{\pgfqpoint{1.131573in}{0.681712in}}%
\pgfpathlineto{\pgfqpoint{1.128043in}{0.675635in}}%
\pgfpathlineto{\pgfqpoint{1.124529in}{0.669584in}}%
\pgfpathlineto{\pgfqpoint{1.115861in}{0.668553in}}%
\pgfpathlineto{\pgfqpoint{1.107075in}{0.667768in}}%
\pgfpathclose%
\pgfusepath{fill}%
\end{pgfscope}%
\begin{pgfscope}%
\pgfpathrectangle{\pgfqpoint{0.407592in}{0.222904in}}{\pgfqpoint{1.411163in}{1.411163in}}%
\pgfusepath{clip}%
\pgfsetbuttcap%
\pgfsetroundjoin%
\definecolor{currentfill}{rgb}{0.749020,0.772549,0.792157}%
\pgfsetfillcolor{currentfill}%
\pgfsetfillopacity{0.160000}%
\pgfsetlinewidth{0.000000pt}%
\definecolor{currentstroke}{rgb}{0.000000,0.000000,0.000000}%
\pgfsetstrokecolor{currentstroke}%
\pgfsetdash{}{0pt}%
\pgfpathmoveto{\pgfqpoint{0.686242in}{0.952460in}}%
\pgfpathlineto{\pgfqpoint{0.672487in}{0.964103in}}%
\pgfpathlineto{\pgfqpoint{0.658712in}{0.975761in}}%
\pgfpathlineto{\pgfqpoint{0.657532in}{0.964496in}}%
\pgfpathlineto{\pgfqpoint{0.658248in}{0.953219in}}%
\pgfpathlineto{\pgfqpoint{0.672131in}{0.942205in}}%
\pgfpathlineto{\pgfqpoint{0.685989in}{0.931212in}}%
\pgfpathlineto{\pgfqpoint{0.685228in}{0.941840in}}%
\pgfpathlineto{\pgfqpoint{0.686242in}{0.952460in}}%
\pgfpathclose%
\pgfusepath{fill}%
\end{pgfscope}%
\begin{pgfscope}%
\pgfpathrectangle{\pgfqpoint{0.407592in}{0.222904in}}{\pgfqpoint{1.411163in}{1.411163in}}%
\pgfusepath{clip}%
\pgfsetbuttcap%
\pgfsetroundjoin%
\definecolor{currentfill}{rgb}{0.749020,0.772549,0.792157}%
\pgfsetfillcolor{currentfill}%
\pgfsetfillopacity{0.160000}%
\pgfsetlinewidth{0.000000pt}%
\definecolor{currentstroke}{rgb}{0.000000,0.000000,0.000000}%
\pgfsetstrokecolor{currentstroke}%
\pgfsetdash{}{0pt}%
\pgfpathmoveto{\pgfqpoint{1.397646in}{0.850281in}}%
\pgfpathlineto{\pgfqpoint{1.410810in}{0.860548in}}%
\pgfpathlineto{\pgfqpoint{1.424004in}{0.870837in}}%
\pgfpathlineto{\pgfqpoint{1.429629in}{0.880032in}}%
\pgfpathlineto{\pgfqpoint{1.433727in}{0.889305in}}%
\pgfpathlineto{\pgfqpoint{1.420222in}{0.878344in}}%
\pgfpathlineto{\pgfqpoint{1.406741in}{0.867403in}}%
\pgfpathlineto{\pgfqpoint{1.402896in}{0.858804in}}%
\pgfpathlineto{\pgfqpoint{1.397646in}{0.850281in}}%
\pgfpathclose%
\pgfusepath{fill}%
\end{pgfscope}%
\begin{pgfscope}%
\pgfpathrectangle{\pgfqpoint{0.407592in}{0.222904in}}{\pgfqpoint{1.411163in}{1.411163in}}%
\pgfusepath{clip}%
\pgfsetbuttcap%
\pgfsetroundjoin%
\definecolor{currentfill}{rgb}{0.749020,0.772549,0.792157}%
\pgfsetfillcolor{currentfill}%
\pgfsetfillopacity{0.160000}%
\pgfsetlinewidth{0.000000pt}%
\definecolor{currentstroke}{rgb}{0.000000,0.000000,0.000000}%
\pgfsetstrokecolor{currentstroke}%
\pgfsetdash{}{0pt}%
\pgfpathmoveto{\pgfqpoint{0.974180in}{0.697382in}}%
\pgfpathlineto{\pgfqpoint{0.964274in}{0.704795in}}%
\pgfpathlineto{\pgfqpoint{0.954332in}{0.712235in}}%
\pgfpathlineto{\pgfqpoint{0.963224in}{0.708704in}}%
\pgfpathlineto{\pgfqpoint{0.972697in}{0.705441in}}%
\pgfpathlineto{\pgfqpoint{0.981312in}{0.698485in}}%
\pgfpathlineto{\pgfqpoint{0.989891in}{0.691556in}}%
\pgfpathlineto{\pgfqpoint{0.981789in}{0.694354in}}%
\pgfpathlineto{\pgfqpoint{0.974180in}{0.697382in}}%
\pgfpathclose%
\pgfusepath{fill}%
\end{pgfscope}%
\begin{pgfscope}%
\pgfpathrectangle{\pgfqpoint{0.407592in}{0.222904in}}{\pgfqpoint{1.411163in}{1.411163in}}%
\pgfusepath{clip}%
\pgfsetbuttcap%
\pgfsetroundjoin%
\definecolor{currentfill}{rgb}{0.749020,0.772549,0.792157}%
\pgfsetfillcolor{currentfill}%
\pgfsetfillopacity{0.160000}%
\pgfsetlinewidth{0.000000pt}%
\definecolor{currentstroke}{rgb}{0.000000,0.000000,0.000000}%
\pgfsetstrokecolor{currentstroke}%
\pgfsetdash{}{0pt}%
\pgfpathmoveto{\pgfqpoint{0.531649in}{1.121571in}}%
\pgfpathlineto{\pgfqpoint{0.518078in}{1.133987in}}%
\pgfpathlineto{\pgfqpoint{0.504495in}{1.146414in}}%
\pgfpathlineto{\pgfqpoint{0.496945in}{1.131767in}}%
\pgfpathlineto{\pgfqpoint{0.491914in}{1.116936in}}%
\pgfpathlineto{\pgfqpoint{0.505921in}{1.105081in}}%
\pgfpathlineto{\pgfqpoint{0.519908in}{1.093242in}}%
\pgfpathlineto{\pgfqpoint{0.524575in}{1.107492in}}%
\pgfpathlineto{\pgfqpoint{0.531649in}{1.121571in}}%
\pgfpathclose%
\pgfusepath{fill}%
\end{pgfscope}%
\begin{pgfscope}%
\pgfpathrectangle{\pgfqpoint{0.407592in}{0.222904in}}{\pgfqpoint{1.411163in}{1.411163in}}%
\pgfusepath{clip}%
\pgfsetbuttcap%
\pgfsetroundjoin%
\definecolor{currentfill}{rgb}{0.749020,0.772549,0.792157}%
\pgfsetfillcolor{currentfill}%
\pgfsetfillopacity{0.160000}%
\pgfsetlinewidth{0.000000pt}%
\definecolor{currentstroke}{rgb}{0.000000,0.000000,0.000000}%
\pgfsetstrokecolor{currentstroke}%
\pgfsetdash{}{0pt}%
\pgfpathmoveto{\pgfqpoint{1.054113in}{0.668375in}}%
\pgfpathlineto{\pgfqpoint{1.050411in}{0.674304in}}%
\pgfpathlineto{\pgfqpoint{1.046692in}{0.680260in}}%
\pgfpathlineto{\pgfqpoint{1.057121in}{0.679378in}}%
\pgfpathlineto{\pgfqpoint{1.067700in}{0.678798in}}%
\pgfpathlineto{\pgfqpoint{1.069625in}{0.672964in}}%
\pgfpathlineto{\pgfqpoint{1.071541in}{0.667158in}}%
\pgfpathlineto{\pgfqpoint{1.062765in}{0.667641in}}%
\pgfpathlineto{\pgfqpoint{1.054113in}{0.668375in}}%
\pgfpathclose%
\pgfusepath{fill}%
\end{pgfscope}%
\begin{pgfscope}%
\pgfpathrectangle{\pgfqpoint{0.407592in}{0.222904in}}{\pgfqpoint{1.411163in}{1.411163in}}%
\pgfusepath{clip}%
\pgfsetbuttcap%
\pgfsetroundjoin%
\definecolor{currentfill}{rgb}{0.749020,0.772549,0.792157}%
\pgfsetfillcolor{currentfill}%
\pgfsetfillopacity{0.160000}%
\pgfsetlinewidth{0.000000pt}%
\definecolor{currentstroke}{rgb}{0.000000,0.000000,0.000000}%
\pgfsetstrokecolor{currentstroke}%
\pgfsetdash{}{0pt}%
\pgfpathmoveto{\pgfqpoint{1.223897in}{0.716803in}}%
\pgfpathlineto{\pgfqpoint{1.233699in}{0.724596in}}%
\pgfpathlineto{\pgfqpoint{1.243536in}{0.732416in}}%
\pgfpathlineto{\pgfqpoint{1.253072in}{0.737047in}}%
\pgfpathlineto{\pgfqpoint{1.261858in}{0.741915in}}%
\pgfpathlineto{\pgfqpoint{1.250874in}{0.733506in}}%
\pgfpathlineto{\pgfqpoint{1.239925in}{0.725125in}}%
\pgfpathlineto{\pgfqpoint{1.232236in}{0.720859in}}%
\pgfpathlineto{\pgfqpoint{1.223897in}{0.716803in}}%
\pgfpathclose%
\pgfusepath{fill}%
\end{pgfscope}%
\begin{pgfscope}%
\pgfpathrectangle{\pgfqpoint{0.407592in}{0.222904in}}{\pgfqpoint{1.411163in}{1.411163in}}%
\pgfusepath{clip}%
\pgfsetbuttcap%
\pgfsetroundjoin%
\definecolor{currentfill}{rgb}{0.749020,0.772549,0.792157}%
\pgfsetfillcolor{currentfill}%
\pgfsetfillopacity{0.160000}%
\pgfsetlinewidth{0.000000pt}%
\definecolor{currentstroke}{rgb}{0.000000,0.000000,0.000000}%
\pgfsetstrokecolor{currentstroke}%
\pgfsetdash{}{0pt}%
\pgfpathmoveto{\pgfqpoint{1.089304in}{0.666954in}}%
\pgfpathlineto{\pgfqpoint{1.089209in}{0.672740in}}%
\pgfpathlineto{\pgfqpoint{1.089114in}{0.678553in}}%
\pgfpathlineto{\pgfqpoint{1.099849in}{0.678890in}}%
\pgfpathlineto{\pgfqpoint{1.110537in}{0.679531in}}%
\pgfpathlineto{\pgfqpoint{1.108802in}{0.673636in}}%
\pgfpathlineto{\pgfqpoint{1.107075in}{0.667768in}}%
\pgfpathlineto{\pgfqpoint{1.098209in}{0.667234in}}%
\pgfpathlineto{\pgfqpoint{1.089304in}{0.666954in}}%
\pgfpathclose%
\pgfusepath{fill}%
\end{pgfscope}%
\begin{pgfscope}%
\pgfpathrectangle{\pgfqpoint{0.407592in}{0.222904in}}{\pgfqpoint{1.411163in}{1.411163in}}%
\pgfusepath{clip}%
\pgfsetbuttcap%
\pgfsetroundjoin%
\definecolor{currentfill}{rgb}{0.749020,0.772549,0.792157}%
\pgfsetfillcolor{currentfill}%
\pgfsetfillopacity{0.160000}%
\pgfsetlinewidth{0.000000pt}%
\definecolor{currentstroke}{rgb}{0.000000,0.000000,0.000000}%
\pgfsetstrokecolor{currentstroke}%
\pgfsetdash{}{0pt}%
\pgfpathmoveto{\pgfqpoint{1.684775in}{1.119923in}}%
\pgfpathlineto{\pgfqpoint{1.698699in}{1.131853in}}%
\pgfpathlineto{\pgfqpoint{1.712642in}{1.143798in}}%
\pgfpathlineto{\pgfqpoint{1.710375in}{1.159167in}}%
\pgfpathlineto{\pgfqpoint{1.705438in}{1.174325in}}%
\pgfpathlineto{\pgfqpoint{1.691792in}{1.161819in}}%
\pgfpathlineto{\pgfqpoint{1.678158in}{1.149325in}}%
\pgfpathlineto{\pgfqpoint{1.682743in}{1.134722in}}%
\pgfpathlineto{\pgfqpoint{1.684775in}{1.119923in}}%
\pgfpathclose%
\pgfusepath{fill}%
\end{pgfscope}%
\begin{pgfscope}%
\pgfpathrectangle{\pgfqpoint{0.407592in}{0.222904in}}{\pgfqpoint{1.411163in}{1.411163in}}%
\pgfusepath{clip}%
\pgfsetbuttcap%
\pgfsetroundjoin%
\definecolor{currentfill}{rgb}{0.749020,0.772549,0.792157}%
\pgfsetfillcolor{currentfill}%
\pgfsetfillopacity{0.160000}%
\pgfsetlinewidth{0.000000pt}%
\definecolor{currentstroke}{rgb}{0.000000,0.000000,0.000000}%
\pgfsetstrokecolor{currentstroke}%
\pgfsetdash{}{0pt}%
\pgfpathmoveto{\pgfqpoint{1.332913in}{0.795459in}}%
\pgfpathlineto{\pgfqpoint{1.345383in}{0.804961in}}%
\pgfpathlineto{\pgfqpoint{1.357886in}{0.814488in}}%
\pgfpathlineto{\pgfqpoint{1.365270in}{0.822080in}}%
\pgfpathlineto{\pgfqpoint{1.371404in}{0.829816in}}%
\pgfpathlineto{\pgfqpoint{1.358327in}{0.819617in}}%
\pgfpathlineto{\pgfqpoint{1.345279in}{0.809441in}}%
\pgfpathlineto{\pgfqpoint{1.339661in}{0.802383in}}%
\pgfpathlineto{\pgfqpoint{1.332913in}{0.795459in}}%
\pgfpathclose%
\pgfusepath{fill}%
\end{pgfscope}%
\begin{pgfscope}%
\pgfpathrectangle{\pgfqpoint{0.407592in}{0.222904in}}{\pgfqpoint{1.411163in}{1.411163in}}%
\pgfusepath{clip}%
\pgfsetbuttcap%
\pgfsetroundjoin%
\definecolor{currentfill}{rgb}{0.749020,0.772549,0.792157}%
\pgfsetfillcolor{currentfill}%
\pgfsetfillopacity{0.160000}%
\pgfsetlinewidth{0.000000pt}%
\definecolor{currentstroke}{rgb}{0.000000,0.000000,0.000000}%
\pgfsetstrokecolor{currentstroke}%
\pgfsetdash{}{0pt}%
\pgfpathmoveto{\pgfqpoint{1.071541in}{0.667158in}}%
\pgfpathlineto{\pgfqpoint{1.069625in}{0.672964in}}%
\pgfpathlineto{\pgfqpoint{1.067700in}{0.678798in}}%
\pgfpathlineto{\pgfqpoint{1.078381in}{0.678522in}}%
\pgfpathlineto{\pgfqpoint{1.089114in}{0.678553in}}%
\pgfpathlineto{\pgfqpoint{1.089209in}{0.672740in}}%
\pgfpathlineto{\pgfqpoint{1.089304in}{0.666954in}}%
\pgfpathlineto{\pgfqpoint{1.080401in}{0.666929in}}%
\pgfpathlineto{\pgfqpoint{1.071541in}{0.667158in}}%
\pgfpathclose%
\pgfusepath{fill}%
\end{pgfscope}%
\begin{pgfscope}%
\pgfpathrectangle{\pgfqpoint{0.407592in}{0.222904in}}{\pgfqpoint{1.411163in}{1.411163in}}%
\pgfusepath{clip}%
\pgfsetbuttcap%
\pgfsetroundjoin%
\definecolor{currentfill}{rgb}{0.749020,0.772549,0.792157}%
\pgfsetfillcolor{currentfill}%
\pgfsetfillopacity{0.160000}%
\pgfsetlinewidth{0.000000pt}%
\definecolor{currentstroke}{rgb}{0.000000,0.000000,0.000000}%
\pgfsetstrokecolor{currentstroke}%
\pgfsetdash{}{0pt}%
\pgfpathmoveto{\pgfqpoint{0.768605in}{0.865674in}}%
\pgfpathlineto{\pgfqpoint{0.754898in}{0.876547in}}%
\pgfpathlineto{\pgfqpoint{0.741166in}{0.887440in}}%
\pgfpathlineto{\pgfqpoint{0.743521in}{0.878180in}}%
\pgfpathlineto{\pgfqpoint{0.747421in}{0.869006in}}%
\pgfpathlineto{\pgfqpoint{0.760979in}{0.858783in}}%
\pgfpathlineto{\pgfqpoint{0.774506in}{0.848584in}}%
\pgfpathlineto{\pgfqpoint{0.770845in}{0.857087in}}%
\pgfpathlineto{\pgfqpoint{0.768605in}{0.865674in}}%
\pgfpathclose%
\pgfusepath{fill}%
\end{pgfscope}%
\begin{pgfscope}%
\pgfpathrectangle{\pgfqpoint{0.407592in}{0.222904in}}{\pgfqpoint{1.411163in}{1.411163in}}%
\pgfusepath{clip}%
\pgfsetbuttcap%
\pgfsetroundjoin%
\definecolor{currentfill}{rgb}{0.749020,0.772549,0.792157}%
\pgfsetfillcolor{currentfill}%
\pgfsetfillopacity{0.160000}%
\pgfsetlinewidth{0.000000pt}%
\definecolor{currentstroke}{rgb}{0.000000,0.000000,0.000000}%
\pgfsetstrokecolor{currentstroke}%
\pgfsetdash{}{0pt}%
\pgfpathmoveto{\pgfqpoint{1.515267in}{0.955480in}}%
\pgfpathlineto{\pgfqpoint{1.528942in}{0.966578in}}%
\pgfpathlineto{\pgfqpoint{1.542642in}{0.977696in}}%
\pgfpathlineto{\pgfqpoint{1.545612in}{0.989616in}}%
\pgfpathlineto{\pgfqpoint{1.546564in}{1.001508in}}%
\pgfpathlineto{\pgfqpoint{1.532845in}{0.989754in}}%
\pgfpathlineto{\pgfqpoint{1.519145in}{0.978016in}}%
\pgfpathlineto{\pgfqpoint{1.518154in}{0.966759in}}%
\pgfpathlineto{\pgfqpoint{1.515267in}{0.955480in}}%
\pgfpathclose%
\pgfusepath{fill}%
\end{pgfscope}%
\begin{pgfscope}%
\pgfpathrectangle{\pgfqpoint{0.407592in}{0.222904in}}{\pgfqpoint{1.411163in}{1.411163in}}%
\pgfusepath{clip}%
\pgfsetbuttcap%
\pgfsetroundjoin%
\definecolor{currentfill}{rgb}{0.749020,0.772549,0.792157}%
\pgfsetfillcolor{currentfill}%
\pgfsetfillopacity{0.160000}%
\pgfsetlinewidth{0.000000pt}%
\definecolor{currentstroke}{rgb}{0.000000,0.000000,0.000000}%
\pgfsetstrokecolor{currentstroke}%
\pgfsetdash{}{0pt}%
\pgfpathmoveto{\pgfqpoint{0.828311in}{0.808016in}}%
\pgfpathlineto{\pgfqpoint{0.814905in}{0.818123in}}%
\pgfpathlineto{\pgfqpoint{0.801469in}{0.828254in}}%
\pgfpathlineto{\pgfqpoint{0.806173in}{0.820544in}}%
\pgfpathlineto{\pgfqpoint{0.812157in}{0.812986in}}%
\pgfpathlineto{\pgfqpoint{0.825152in}{0.803525in}}%
\pgfpathlineto{\pgfqpoint{0.838112in}{0.794089in}}%
\pgfpathlineto{\pgfqpoint{0.832633in}{0.800982in}}%
\pgfpathlineto{\pgfqpoint{0.828311in}{0.808016in}}%
\pgfpathclose%
\pgfusepath{fill}%
\end{pgfscope}%
\begin{pgfscope}%
\pgfpathrectangle{\pgfqpoint{0.407592in}{0.222904in}}{\pgfqpoint{1.411163in}{1.411163in}}%
\pgfusepath{clip}%
\pgfsetbuttcap%
\pgfsetroundjoin%
\definecolor{currentfill}{rgb}{0.749020,0.772549,0.792157}%
\pgfsetfillcolor{currentfill}%
\pgfsetfillopacity{0.160000}%
\pgfsetlinewidth{0.000000pt}%
\definecolor{currentstroke}{rgb}{0.000000,0.000000,0.000000}%
\pgfsetstrokecolor{currentstroke}%
\pgfsetdash{}{0pt}%
\pgfpathmoveto{\pgfqpoint{1.170952in}{0.689494in}}%
\pgfpathlineto{\pgfqpoint{1.177859in}{0.696250in}}%
\pgfpathlineto{\pgfqpoint{1.184795in}{0.703034in}}%
\pgfpathlineto{\pgfqpoint{1.195348in}{0.706073in}}%
\pgfpathlineto{\pgfqpoint{1.205412in}{0.709392in}}%
\pgfpathlineto{\pgfqpoint{1.196979in}{0.702154in}}%
\pgfpathlineto{\pgfqpoint{1.188578in}{0.694944in}}%
\pgfpathlineto{\pgfqpoint{1.179972in}{0.692099in}}%
\pgfpathlineto{\pgfqpoint{1.170952in}{0.689494in}}%
\pgfpathclose%
\pgfusepath{fill}%
\end{pgfscope}%
\begin{pgfscope}%
\pgfpathrectangle{\pgfqpoint{0.407592in}{0.222904in}}{\pgfqpoint{1.411163in}{1.411163in}}%
\pgfusepath{clip}%
\pgfsetbuttcap%
\pgfsetroundjoin%
\definecolor{currentfill}{rgb}{0.749020,0.772549,0.792157}%
\pgfsetfillcolor{currentfill}%
\pgfsetfillopacity{0.160000}%
\pgfsetlinewidth{0.000000pt}%
\definecolor{currentstroke}{rgb}{0.000000,0.000000,0.000000}%
\pgfsetstrokecolor{currentstroke}%
\pgfsetdash{}{0pt}%
\pgfpathmoveto{\pgfqpoint{0.504495in}{1.146414in}}%
\pgfpathlineto{\pgfqpoint{0.490898in}{1.158853in}}%
\pgfpathlineto{\pgfqpoint{0.477289in}{1.171304in}}%
\pgfpathlineto{\pgfqpoint{0.469249in}{1.156099in}}%
\pgfpathlineto{\pgfqpoint{0.463842in}{1.140695in}}%
\pgfpathlineto{\pgfqpoint{0.477888in}{1.128807in}}%
\pgfpathlineto{\pgfqpoint{0.491914in}{1.116936in}}%
\pgfpathlineto{\pgfqpoint{0.496945in}{1.131767in}}%
\pgfpathlineto{\pgfqpoint{0.504495in}{1.146414in}}%
\pgfpathclose%
\pgfusepath{fill}%
\end{pgfscope}%
\begin{pgfscope}%
\pgfpathrectangle{\pgfqpoint{0.407592in}{0.222904in}}{\pgfqpoint{1.411163in}{1.411163in}}%
\pgfusepath{clip}%
\pgfsetbuttcap%
\pgfsetroundjoin%
\definecolor{currentfill}{rgb}{0.749020,0.772549,0.792157}%
\pgfsetfillcolor{currentfill}%
\pgfsetfillopacity{0.160000}%
\pgfsetlinewidth{0.000000pt}%
\definecolor{currentstroke}{rgb}{0.000000,0.000000,0.000000}%
\pgfsetstrokecolor{currentstroke}%
\pgfsetdash{}{0pt}%
\pgfpathmoveto{\pgfqpoint{1.261858in}{0.741915in}}%
\pgfpathlineto{\pgfqpoint{1.272878in}{0.750350in}}%
\pgfpathlineto{\pgfqpoint{1.283933in}{0.758812in}}%
\pgfpathlineto{\pgfqpoint{1.292921in}{0.764517in}}%
\pgfpathlineto{\pgfqpoint{1.281371in}{0.755743in}}%
\pgfpathlineto{\pgfqpoint{1.269856in}{0.746995in}}%
\pgfusepath{fill}%
\end{pgfscope}%
\begin{pgfscope}%
\pgfpathrectangle{\pgfqpoint{0.407592in}{0.222904in}}{\pgfqpoint{1.411163in}{1.411163in}}%
\pgfusepath{clip}%
\pgfsetbuttcap%
\pgfsetroundjoin%
\definecolor{currentfill}{rgb}{0.749020,0.772549,0.792157}%
\pgfsetfillcolor{currentfill}%
\pgfsetfillopacity{0.160000}%
\pgfsetlinewidth{0.000000pt}%
\definecolor{currentstroke}{rgb}{0.000000,0.000000,0.000000}%
\pgfsetstrokecolor{currentstroke}%
\pgfsetdash{}{0pt}%
\pgfpathmoveto{\pgfqpoint{1.712642in}{1.143798in}}%
\pgfpathlineto{\pgfqpoint{1.726603in}{1.155760in}}%
\pgfpathlineto{\pgfqpoint{1.740584in}{1.167739in}}%
\pgfpathlineto{\pgfqpoint{1.738068in}{1.183667in}}%
\pgfpathlineto{\pgfqpoint{1.732768in}{1.199371in}}%
\pgfpathlineto{\pgfqpoint{1.719097in}{1.186842in}}%
\pgfpathlineto{\pgfqpoint{1.705438in}{1.174325in}}%
\pgfpathlineto{\pgfqpoint{1.710375in}{1.159167in}}%
\pgfpathlineto{\pgfqpoint{1.712642in}{1.143798in}}%
\pgfpathclose%
\pgfusepath{fill}%
\end{pgfscope}%
\begin{pgfscope}%
\pgfpathrectangle{\pgfqpoint{0.407592in}{0.222904in}}{\pgfqpoint{1.411163in}{1.411163in}}%
\pgfusepath{clip}%
\pgfsetbuttcap%
\pgfsetroundjoin%
\definecolor{currentfill}{rgb}{0.749020,0.772549,0.792157}%
\pgfsetfillcolor{currentfill}%
\pgfsetfillopacity{0.160000}%
\pgfsetlinewidth{0.000000pt}%
\definecolor{currentstroke}{rgb}{0.000000,0.000000,0.000000}%
\pgfsetstrokecolor{currentstroke}%
\pgfsetdash{}{0pt}%
\pgfpathmoveto{\pgfqpoint{0.658712in}{0.975761in}}%
\pgfpathlineto{\pgfqpoint{0.644918in}{0.987436in}}%
\pgfpathlineto{\pgfqpoint{0.631105in}{0.999127in}}%
\pgfpathlineto{\pgfqpoint{0.629747in}{0.987226in}}%
\pgfpathlineto{\pgfqpoint{0.630405in}{0.975306in}}%
\pgfpathlineto{\pgfqpoint{0.644339in}{0.964252in}}%
\pgfpathlineto{\pgfqpoint{0.658248in}{0.953219in}}%
\pgfpathlineto{\pgfqpoint{0.657532in}{0.964496in}}%
\pgfpathlineto{\pgfqpoint{0.658712in}{0.975761in}}%
\pgfpathclose%
\pgfusepath{fill}%
\end{pgfscope}%
\begin{pgfscope}%
\pgfpathrectangle{\pgfqpoint{0.407592in}{0.222904in}}{\pgfqpoint{1.411163in}{1.411163in}}%
\pgfusepath{clip}%
\pgfsetbuttcap%
\pgfsetroundjoin%
\definecolor{currentfill}{rgb}{0.749020,0.772549,0.792157}%
\pgfsetfillcolor{currentfill}%
\pgfsetfillopacity{0.160000}%
\pgfsetlinewidth{0.000000pt}%
\definecolor{currentstroke}{rgb}{0.000000,0.000000,0.000000}%
\pgfsetstrokecolor{currentstroke}%
\pgfsetdash{}{0pt}%
\pgfpathmoveto{\pgfqpoint{0.938446in}{0.720028in}}%
\pgfpathlineto{\pgfqpoint{0.927327in}{0.728050in}}%
\pgfpathlineto{\pgfqpoint{0.916170in}{0.736098in}}%
\pgfpathlineto{\pgfqpoint{0.924874in}{0.731518in}}%
\pgfpathlineto{\pgfqpoint{0.934335in}{0.727200in}}%
\pgfpathlineto{\pgfqpoint{0.944352in}{0.719704in}}%
\pgfpathlineto{\pgfqpoint{0.954332in}{0.712235in}}%
\pgfpathlineto{\pgfqpoint{0.946061in}{0.716016in}}%
\pgfpathlineto{\pgfqpoint{0.938446in}{0.720028in}}%
\pgfpathclose%
\pgfusepath{fill}%
\end{pgfscope}%
\begin{pgfscope}%
\pgfpathrectangle{\pgfqpoint{0.407592in}{0.222904in}}{\pgfqpoint{1.411163in}{1.411163in}}%
\pgfusepath{clip}%
\pgfsetbuttcap%
\pgfsetroundjoin%
\definecolor{currentfill}{rgb}{0.749020,0.772549,0.792157}%
\pgfsetfillcolor{currentfill}%
\pgfsetfillopacity{0.160000}%
\pgfsetlinewidth{0.000000pt}%
\definecolor{currentstroke}{rgb}{0.000000,0.000000,0.000000}%
\pgfsetstrokecolor{currentstroke}%
\pgfsetdash{}{0pt}%
\pgfpathmoveto{\pgfqpoint{1.292921in}{0.764517in}}%
\pgfpathlineto{\pgfqpoint{1.304505in}{0.773318in}}%
\pgfpathlineto{\pgfqpoint{1.316125in}{0.782146in}}%
\pgfpathlineto{\pgfqpoint{1.325058in}{0.788703in}}%
\pgfpathlineto{\pgfqpoint{1.332913in}{0.795459in}}%
\pgfpathlineto{\pgfqpoint{1.320476in}{0.785982in}}%
\pgfpathlineto{\pgfqpoint{1.308072in}{0.776530in}}%
\pgfpathlineto{\pgfqpoint{1.300978in}{0.770432in}}%
\pgfpathlineto{\pgfqpoint{1.292921in}{0.764517in}}%
\pgfpathclose%
\pgfusepath{fill}%
\end{pgfscope}%
\begin{pgfscope}%
\pgfpathrectangle{\pgfqpoint{0.407592in}{0.222904in}}{\pgfqpoint{1.411163in}{1.411163in}}%
\pgfusepath{clip}%
\pgfsetbuttcap%
\pgfsetroundjoin%
\definecolor{currentfill}{rgb}{0.749020,0.772549,0.792157}%
\pgfsetfillcolor{currentfill}%
\pgfsetfillopacity{0.160000}%
\pgfsetlinewidth{0.000000pt}%
\definecolor{currentstroke}{rgb}{0.000000,0.000000,0.000000}%
\pgfsetstrokecolor{currentstroke}%
\pgfsetdash{}{0pt}%
\pgfpathmoveto{\pgfqpoint{0.989891in}{0.691556in}}%
\pgfpathlineto{\pgfqpoint{0.981312in}{0.698485in}}%
\pgfpathlineto{\pgfqpoint{0.972697in}{0.705441in}}%
\pgfpathlineto{\pgfqpoint{0.982710in}{0.702460in}}%
\pgfpathlineto{\pgfqpoint{0.993217in}{0.699776in}}%
\pgfpathlineto{\pgfqpoint{1.000338in}{0.693225in}}%
\pgfpathlineto{\pgfqpoint{1.007429in}{0.686702in}}%
\pgfpathlineto{\pgfqpoint{0.998451in}{0.689002in}}%
\pgfpathlineto{\pgfqpoint{0.989891in}{0.691556in}}%
\pgfpathclose%
\pgfusepath{fill}%
\end{pgfscope}%
\begin{pgfscope}%
\pgfpathrectangle{\pgfqpoint{0.407592in}{0.222904in}}{\pgfqpoint{1.411163in}{1.411163in}}%
\pgfusepath{clip}%
\pgfsetbuttcap%
\pgfsetroundjoin%
\definecolor{currentfill}{rgb}{0.749020,0.772549,0.792157}%
\pgfsetfillcolor{currentfill}%
\pgfsetfillopacity{0.160000}%
\pgfsetlinewidth{0.000000pt}%
\definecolor{currentstroke}{rgb}{0.000000,0.000000,0.000000}%
\pgfsetstrokecolor{currentstroke}%
\pgfsetdash{}{0pt}%
\pgfpathmoveto{\pgfqpoint{1.424004in}{0.870837in}}%
\pgfpathlineto{\pgfqpoint{1.437227in}{0.881150in}}%
\pgfpathlineto{\pgfqpoint{1.450479in}{0.891485in}}%
\pgfpathlineto{\pgfqpoint{1.456470in}{0.901345in}}%
\pgfpathlineto{\pgfqpoint{1.460810in}{0.911284in}}%
\pgfpathlineto{\pgfqpoint{1.447257in}{0.900285in}}%
\pgfpathlineto{\pgfqpoint{1.433727in}{0.889305in}}%
\pgfpathlineto{\pgfqpoint{1.429629in}{0.880032in}}%
\pgfpathlineto{\pgfqpoint{1.424004in}{0.870837in}}%
\pgfpathclose%
\pgfusepath{fill}%
\end{pgfscope}%
\begin{pgfscope}%
\pgfpathrectangle{\pgfqpoint{0.407592in}{0.222904in}}{\pgfqpoint{1.411163in}{1.411163in}}%
\pgfusepath{clip}%
\pgfsetbuttcap%
\pgfsetroundjoin%
\definecolor{currentfill}{rgb}{0.749020,0.772549,0.792157}%
\pgfsetfillcolor{currentfill}%
\pgfsetfillopacity{0.160000}%
\pgfsetlinewidth{0.000000pt}%
\definecolor{currentstroke}{rgb}{0.000000,0.000000,0.000000}%
\pgfsetstrokecolor{currentstroke}%
\pgfsetdash{}{0pt}%
\pgfpathmoveto{\pgfqpoint{0.863928in}{0.775293in}}%
\pgfpathlineto{\pgfqpoint{0.851037in}{0.784679in}}%
\pgfpathlineto{\pgfqpoint{0.838112in}{0.794089in}}%
\pgfpathlineto{\pgfqpoint{0.844730in}{0.787370in}}%
\pgfpathlineto{\pgfqpoint{0.852464in}{0.780857in}}%
\pgfpathlineto{\pgfqpoint{0.864698in}{0.772093in}}%
\pgfpathlineto{\pgfqpoint{0.876896in}{0.763355in}}%
\pgfpathlineto{\pgfqpoint{0.869913in}{0.769230in}}%
\pgfpathlineto{\pgfqpoint{0.863928in}{0.775293in}}%
\pgfpathclose%
\pgfusepath{fill}%
\end{pgfscope}%
\begin{pgfscope}%
\pgfpathrectangle{\pgfqpoint{0.407592in}{0.222904in}}{\pgfqpoint{1.411163in}{1.411163in}}%
\pgfusepath{clip}%
\pgfsetbuttcap%
\pgfsetroundjoin%
\definecolor{currentfill}{rgb}{0.749020,0.772549,0.792157}%
\pgfsetfillcolor{currentfill}%
\pgfsetfillopacity{0.160000}%
\pgfsetlinewidth{0.000000pt}%
\definecolor{currentstroke}{rgb}{0.000000,0.000000,0.000000}%
\pgfsetstrokecolor{currentstroke}%
\pgfsetdash{}{0pt}%
\pgfpathmoveto{\pgfqpoint{1.542642in}{0.977696in}}%
\pgfpathlineto{\pgfqpoint{1.556366in}{0.988834in}}%
\pgfpathlineto{\pgfqpoint{1.570115in}{0.999993in}}%
\pgfpathlineto{\pgfqpoint{1.573157in}{1.012546in}}%
\pgfpathlineto{\pgfqpoint{1.574058in}{1.025063in}}%
\pgfpathlineto{\pgfqpoint{1.560302in}{1.013278in}}%
\pgfpathlineto{\pgfqpoint{1.546564in}{1.001508in}}%
\pgfpathlineto{\pgfqpoint{1.545612in}{0.989616in}}%
\pgfpathlineto{\pgfqpoint{1.542642in}{0.977696in}}%
\pgfpathclose%
\pgfusepath{fill}%
\end{pgfscope}%
\begin{pgfscope}%
\pgfpathrectangle{\pgfqpoint{0.407592in}{0.222904in}}{\pgfqpoint{1.411163in}{1.411163in}}%
\pgfusepath{clip}%
\pgfsetbuttcap%
\pgfsetroundjoin%
\definecolor{currentfill}{rgb}{0.749020,0.772549,0.792157}%
\pgfsetfillcolor{currentfill}%
\pgfsetfillopacity{0.160000}%
\pgfsetlinewidth{0.000000pt}%
\definecolor{currentstroke}{rgb}{0.000000,0.000000,0.000000}%
\pgfsetstrokecolor{currentstroke}%
\pgfsetdash{}{0pt}%
\pgfpathmoveto{\pgfqpoint{1.151835in}{0.685055in}}%
\pgfpathlineto{\pgfqpoint{1.157114in}{0.691440in}}%
\pgfpathlineto{\pgfqpoint{1.162415in}{0.697853in}}%
\pgfpathlineto{\pgfqpoint{1.173801in}{0.700290in}}%
\pgfpathlineto{\pgfqpoint{1.184795in}{0.703034in}}%
\pgfpathlineto{\pgfqpoint{1.177859in}{0.696250in}}%
\pgfpathlineto{\pgfqpoint{1.170952in}{0.689494in}}%
\pgfpathlineto{\pgfqpoint{1.161559in}{0.687142in}}%
\pgfpathlineto{\pgfqpoint{1.151835in}{0.685055in}}%
\pgfpathclose%
\pgfusepath{fill}%
\end{pgfscope}%
\begin{pgfscope}%
\pgfpathrectangle{\pgfqpoint{0.407592in}{0.222904in}}{\pgfqpoint{1.411163in}{1.411163in}}%
\pgfusepath{clip}%
\pgfsetbuttcap%
\pgfsetroundjoin%
\definecolor{currentfill}{rgb}{0.749020,0.772549,0.792157}%
\pgfsetfillcolor{currentfill}%
\pgfsetfillopacity{0.160000}%
\pgfsetlinewidth{0.000000pt}%
\definecolor{currentstroke}{rgb}{0.000000,0.000000,0.000000}%
\pgfsetstrokecolor{currentstroke}%
\pgfsetdash{}{0pt}%
\pgfpathmoveto{\pgfqpoint{0.901178in}{0.745960in}}%
\pgfpathlineto{\pgfqpoint{0.889056in}{0.754644in}}%
\pgfpathlineto{\pgfqpoint{0.876896in}{0.763355in}}%
\pgfpathlineto{\pgfqpoint{0.884849in}{0.757696in}}%
\pgfpathlineto{\pgfqpoint{0.893740in}{0.752279in}}%
\pgfpathlineto{\pgfqpoint{0.904974in}{0.744175in}}%
\pgfpathlineto{\pgfqpoint{0.916170in}{0.736098in}}%
\pgfpathlineto{\pgfqpoint{0.908260in}{0.740920in}}%
\pgfpathlineto{\pgfqpoint{0.901178in}{0.745960in}}%
\pgfpathclose%
\pgfusepath{fill}%
\end{pgfscope}%
\begin{pgfscope}%
\pgfpathrectangle{\pgfqpoint{0.407592in}{0.222904in}}{\pgfqpoint{1.411163in}{1.411163in}}%
\pgfusepath{clip}%
\pgfsetbuttcap%
\pgfsetroundjoin%
\definecolor{currentfill}{rgb}{0.749020,0.772549,0.792157}%
\pgfsetfillcolor{currentfill}%
\pgfsetfillopacity{0.160000}%
\pgfsetlinewidth{0.000000pt}%
\definecolor{currentstroke}{rgb}{0.000000,0.000000,0.000000}%
\pgfsetstrokecolor{currentstroke}%
\pgfsetdash{}{0pt}%
\pgfpathmoveto{\pgfqpoint{0.477289in}{1.171304in}}%
\pgfpathlineto{\pgfqpoint{0.463667in}{1.183767in}}%
\pgfpathlineto{\pgfqpoint{0.450032in}{1.196241in}}%
\pgfpathlineto{\pgfqpoint{0.441489in}{1.180488in}}%
\pgfpathlineto{\pgfqpoint{0.435692in}{1.164522in}}%
\pgfpathlineto{\pgfqpoint{0.449777in}{1.152600in}}%
\pgfpathlineto{\pgfqpoint{0.463842in}{1.140695in}}%
\pgfpathlineto{\pgfqpoint{0.469249in}{1.156099in}}%
\pgfpathlineto{\pgfqpoint{0.477289in}{1.171304in}}%
\pgfpathclose%
\pgfusepath{fill}%
\end{pgfscope}%
\begin{pgfscope}%
\pgfpathrectangle{\pgfqpoint{0.407592in}{0.222904in}}{\pgfqpoint{1.411163in}{1.411163in}}%
\pgfusepath{clip}%
\pgfsetbuttcap%
\pgfsetroundjoin%
\definecolor{currentfill}{rgb}{0.749020,0.772549,0.792157}%
\pgfsetfillcolor{currentfill}%
\pgfsetfillopacity{0.160000}%
\pgfsetlinewidth{0.000000pt}%
\definecolor{currentstroke}{rgb}{0.000000,0.000000,0.000000}%
\pgfsetstrokecolor{currentstroke}%
\pgfsetdash{}{0pt}%
\pgfpathmoveto{\pgfqpoint{0.741166in}{0.887440in}}%
\pgfpathlineto{\pgfqpoint{0.727410in}{0.898353in}}%
\pgfpathlineto{\pgfqpoint{0.713628in}{0.909286in}}%
\pgfpathlineto{\pgfqpoint{0.716086in}{0.899359in}}%
\pgfpathlineto{\pgfqpoint{0.720213in}{0.889520in}}%
\pgfpathlineto{\pgfqpoint{0.733832in}{0.879251in}}%
\pgfpathlineto{\pgfqpoint{0.747421in}{0.869006in}}%
\pgfpathlineto{\pgfqpoint{0.743521in}{0.878180in}}%
\pgfpathlineto{\pgfqpoint{0.741166in}{0.887440in}}%
\pgfpathclose%
\pgfusepath{fill}%
\end{pgfscope}%
\begin{pgfscope}%
\pgfpathrectangle{\pgfqpoint{0.407592in}{0.222904in}}{\pgfqpoint{1.411163in}{1.411163in}}%
\pgfusepath{clip}%
\pgfsetbuttcap%
\pgfsetroundjoin%
\definecolor{currentfill}{rgb}{0.749020,0.772549,0.792157}%
\pgfsetfillcolor{currentfill}%
\pgfsetfillopacity{0.160000}%
\pgfsetlinewidth{0.000000pt}%
\definecolor{currentstroke}{rgb}{0.000000,0.000000,0.000000}%
\pgfsetstrokecolor{currentstroke}%
\pgfsetdash{}{0pt}%
\pgfpathmoveto{\pgfqpoint{1.740584in}{1.167739in}}%
\pgfpathlineto{\pgfqpoint{1.754584in}{1.179733in}}%
\pgfpathlineto{\pgfqpoint{1.768602in}{1.191744in}}%
\pgfpathlineto{\pgfqpoint{1.765824in}{1.208223in}}%
\pgfpathlineto{\pgfqpoint{1.760147in}{1.224461in}}%
\pgfpathlineto{\pgfqpoint{1.746451in}{1.211910in}}%
\pgfpathlineto{\pgfqpoint{1.732768in}{1.199371in}}%
\pgfpathlineto{\pgfqpoint{1.738068in}{1.183667in}}%
\pgfpathlineto{\pgfqpoint{1.740584in}{1.167739in}}%
\pgfpathclose%
\pgfusepath{fill}%
\end{pgfscope}%
\begin{pgfscope}%
\pgfpathrectangle{\pgfqpoint{0.407592in}{0.222904in}}{\pgfqpoint{1.411163in}{1.411163in}}%
\pgfusepath{clip}%
\pgfsetbuttcap%
\pgfsetroundjoin%
\definecolor{currentfill}{rgb}{0.749020,0.772549,0.792157}%
\pgfsetfillcolor{currentfill}%
\pgfsetfillopacity{0.160000}%
\pgfsetlinewidth{0.000000pt}%
\definecolor{currentstroke}{rgb}{0.000000,0.000000,0.000000}%
\pgfsetstrokecolor{currentstroke}%
\pgfsetdash{}{0pt}%
\pgfpathmoveto{\pgfqpoint{1.205412in}{0.709392in}}%
\pgfpathlineto{\pgfqpoint{1.213878in}{0.716658in}}%
\pgfpathlineto{\pgfqpoint{1.222376in}{0.723952in}}%
\pgfpathlineto{\pgfqpoint{1.233289in}{0.728044in}}%
\pgfpathlineto{\pgfqpoint{1.243536in}{0.732416in}}%
\pgfpathlineto{\pgfqpoint{1.233699in}{0.724596in}}%
\pgfpathlineto{\pgfqpoint{1.223897in}{0.716803in}}%
\pgfpathlineto{\pgfqpoint{1.214942in}{0.712974in}}%
\pgfpathlineto{\pgfqpoint{1.205412in}{0.709392in}}%
\pgfpathclose%
\pgfusepath{fill}%
\end{pgfscope}%
\begin{pgfscope}%
\pgfpathrectangle{\pgfqpoint{0.407592in}{0.222904in}}{\pgfqpoint{1.411163in}{1.411163in}}%
\pgfusepath{clip}%
\pgfsetbuttcap%
\pgfsetroundjoin%
\definecolor{currentfill}{rgb}{0.749020,0.772549,0.792157}%
\pgfsetfillcolor{currentfill}%
\pgfsetfillopacity{0.160000}%
\pgfsetlinewidth{0.000000pt}%
\definecolor{currentstroke}{rgb}{0.000000,0.000000,0.000000}%
\pgfsetstrokecolor{currentstroke}%
\pgfsetdash{}{0pt}%
\pgfpathmoveto{\pgfqpoint{1.357886in}{0.814488in}}%
\pgfpathlineto{\pgfqpoint{1.370421in}{0.824040in}}%
\pgfpathlineto{\pgfqpoint{1.382990in}{0.833618in}}%
\pgfpathlineto{\pgfqpoint{1.391004in}{0.841873in}}%
\pgfpathlineto{\pgfqpoint{1.397646in}{0.850281in}}%
\pgfpathlineto{\pgfqpoint{1.384511in}{0.840037in}}%
\pgfpathlineto{\pgfqpoint{1.371404in}{0.829816in}}%
\pgfpathlineto{\pgfqpoint{1.365270in}{0.822080in}}%
\pgfpathlineto{\pgfqpoint{1.357886in}{0.814488in}}%
\pgfpathclose%
\pgfusepath{fill}%
\end{pgfscope}%
\begin{pgfscope}%
\pgfpathrectangle{\pgfqpoint{0.407592in}{0.222904in}}{\pgfqpoint{1.411163in}{1.411163in}}%
\pgfusepath{clip}%
\pgfsetbuttcap%
\pgfsetroundjoin%
\definecolor{currentfill}{rgb}{0.749020,0.772549,0.792157}%
\pgfsetfillcolor{currentfill}%
\pgfsetfillopacity{0.160000}%
\pgfsetlinewidth{0.000000pt}%
\definecolor{currentstroke}{rgb}{0.000000,0.000000,0.000000}%
\pgfsetstrokecolor{currentstroke}%
\pgfsetdash{}{0pt}%
\pgfpathmoveto{\pgfqpoint{0.631105in}{0.999127in}}%
\pgfpathlineto{\pgfqpoint{0.617273in}{1.010834in}}%
\pgfpathlineto{\pgfqpoint{0.603422in}{1.022558in}}%
\pgfpathlineto{\pgfqpoint{0.601871in}{1.010029in}}%
\pgfpathlineto{\pgfqpoint{0.602459in}{0.997475in}}%
\pgfpathlineto{\pgfqpoint{0.616445in}{0.986380in}}%
\pgfpathlineto{\pgfqpoint{0.630405in}{0.975306in}}%
\pgfpathlineto{\pgfqpoint{0.629747in}{0.987226in}}%
\pgfpathlineto{\pgfqpoint{0.631105in}{0.999127in}}%
\pgfpathclose%
\pgfusepath{fill}%
\end{pgfscope}%
\begin{pgfscope}%
\pgfpathrectangle{\pgfqpoint{0.407592in}{0.222904in}}{\pgfqpoint{1.411163in}{1.411163in}}%
\pgfusepath{clip}%
\pgfsetbuttcap%
\pgfsetroundjoin%
\definecolor{currentfill}{rgb}{0.749020,0.772549,0.792157}%
\pgfsetfillcolor{currentfill}%
\pgfsetfillopacity{0.160000}%
\pgfsetlinewidth{0.000000pt}%
\definecolor{currentstroke}{rgb}{0.000000,0.000000,0.000000}%
\pgfsetstrokecolor{currentstroke}%
\pgfsetdash{}{0pt}%
\pgfpathmoveto{\pgfqpoint{1.007429in}{0.686702in}}%
\pgfpathlineto{\pgfqpoint{1.000338in}{0.693225in}}%
\pgfpathlineto{\pgfqpoint{0.993217in}{0.699776in}}%
\pgfpathlineto{\pgfqpoint{1.004171in}{0.697403in}}%
\pgfpathlineto{\pgfqpoint{1.015522in}{0.695351in}}%
\pgfpathlineto{\pgfqpoint{1.021012in}{0.689118in}}%
\pgfpathlineto{\pgfqpoint{1.026478in}{0.682912in}}%
\pgfpathlineto{\pgfqpoint{1.016786in}{0.684669in}}%
\pgfpathlineto{\pgfqpoint{1.007429in}{0.686702in}}%
\pgfpathclose%
\pgfusepath{fill}%
\end{pgfscope}%
\begin{pgfscope}%
\pgfpathrectangle{\pgfqpoint{0.407592in}{0.222904in}}{\pgfqpoint{1.411163in}{1.411163in}}%
\pgfusepath{clip}%
\pgfsetbuttcap%
\pgfsetroundjoin%
\definecolor{currentfill}{rgb}{0.749020,0.772549,0.792157}%
\pgfsetfillcolor{currentfill}%
\pgfsetfillopacity{0.160000}%
\pgfsetlinewidth{0.000000pt}%
\definecolor{currentstroke}{rgb}{0.000000,0.000000,0.000000}%
\pgfsetstrokecolor{currentstroke}%
\pgfsetdash{}{0pt}%
\pgfpathmoveto{\pgfqpoint{0.801469in}{0.828254in}}%
\pgfpathlineto{\pgfqpoint{0.788003in}{0.838407in}}%
\pgfpathlineto{\pgfqpoint{0.774506in}{0.848584in}}%
\pgfpathlineto{\pgfqpoint{0.779583in}{0.840203in}}%
\pgfpathlineto{\pgfqpoint{0.786063in}{0.831984in}}%
\pgfpathlineto{\pgfqpoint{0.799127in}{0.822472in}}%
\pgfpathlineto{\pgfqpoint{0.812157in}{0.812986in}}%
\pgfpathlineto{\pgfqpoint{0.806173in}{0.820544in}}%
\pgfpathlineto{\pgfqpoint{0.801469in}{0.828254in}}%
\pgfpathclose%
\pgfusepath{fill}%
\end{pgfscope}%
\begin{pgfscope}%
\pgfpathrectangle{\pgfqpoint{0.407592in}{0.222904in}}{\pgfqpoint{1.411163in}{1.411163in}}%
\pgfusepath{clip}%
\pgfsetbuttcap%
\pgfsetroundjoin%
\definecolor{currentfill}{rgb}{0.749020,0.772549,0.792157}%
\pgfsetfillcolor{currentfill}%
\pgfsetfillopacity{0.160000}%
\pgfsetlinewidth{0.000000pt}%
\definecolor{currentstroke}{rgb}{0.000000,0.000000,0.000000}%
\pgfsetstrokecolor{currentstroke}%
\pgfsetdash{}{0pt}%
\pgfpathmoveto{\pgfqpoint{1.131573in}{0.681712in}}%
\pgfpathlineto{\pgfqpoint{1.135119in}{0.687818in}}%
\pgfpathlineto{\pgfqpoint{1.138681in}{0.693950in}}%
\pgfpathlineto{\pgfqpoint{1.150691in}{0.695737in}}%
\pgfpathlineto{\pgfqpoint{1.162415in}{0.697853in}}%
\pgfpathlineto{\pgfqpoint{1.157114in}{0.691440in}}%
\pgfpathlineto{\pgfqpoint{1.151835in}{0.685055in}}%
\pgfpathlineto{\pgfqpoint{1.141825in}{0.683242in}}%
\pgfpathlineto{\pgfqpoint{1.131573in}{0.681712in}}%
\pgfpathclose%
\pgfusepath{fill}%
\end{pgfscope}%
\begin{pgfscope}%
\pgfpathrectangle{\pgfqpoint{0.407592in}{0.222904in}}{\pgfqpoint{1.411163in}{1.411163in}}%
\pgfusepath{clip}%
\pgfsetbuttcap%
\pgfsetroundjoin%
\definecolor{currentfill}{rgb}{0.749020,0.772549,0.792157}%
\pgfsetfillcolor{currentfill}%
\pgfsetfillopacity{0.160000}%
\pgfsetlinewidth{0.000000pt}%
\definecolor{currentstroke}{rgb}{0.000000,0.000000,0.000000}%
\pgfsetstrokecolor{currentstroke}%
\pgfsetdash{}{0pt}%
\pgfpathmoveto{\pgfqpoint{1.243536in}{0.732416in}}%
\pgfpathlineto{\pgfqpoint{1.253408in}{0.740265in}}%
\pgfpathlineto{\pgfqpoint{1.263315in}{0.748141in}}%
\pgfpathlineto{\pgfqpoint{1.274050in}{0.753345in}}%
\pgfpathlineto{\pgfqpoint{1.283933in}{0.758812in}}%
\pgfpathlineto{\pgfqpoint{1.272878in}{0.750350in}}%
\pgfpathlineto{\pgfqpoint{1.261858in}{0.741915in}}%
\pgfpathlineto{\pgfqpoint{1.253072in}{0.737047in}}%
\pgfpathlineto{\pgfqpoint{1.243536in}{0.732416in}}%
\pgfpathclose%
\pgfusepath{fill}%
\end{pgfscope}%
\begin{pgfscope}%
\pgfpathrectangle{\pgfqpoint{0.407592in}{0.222904in}}{\pgfqpoint{1.411163in}{1.411163in}}%
\pgfusepath{clip}%
\pgfsetbuttcap%
\pgfsetroundjoin%
\definecolor{currentfill}{rgb}{0.749020,0.772549,0.792157}%
\pgfsetfillcolor{currentfill}%
\pgfsetfillopacity{0.160000}%
\pgfsetlinewidth{0.000000pt}%
\definecolor{currentstroke}{rgb}{0.000000,0.000000,0.000000}%
\pgfsetstrokecolor{currentstroke}%
\pgfsetdash{}{0pt}%
\pgfpathmoveto{\pgfqpoint{0.450032in}{1.196241in}}%
\pgfpathlineto{\pgfqpoint{0.436384in}{1.208728in}}%
\pgfpathlineto{\pgfqpoint{0.422723in}{1.221226in}}%
\pgfpathlineto{\pgfqpoint{0.413663in}{1.204934in}}%
\pgfpathlineto{\pgfqpoint{0.407462in}{1.188415in}}%
\pgfpathlineto{\pgfqpoint{0.421586in}{1.176460in}}%
\pgfpathlineto{\pgfqpoint{0.435692in}{1.164522in}}%
\pgfpathlineto{\pgfqpoint{0.441489in}{1.180488in}}%
\pgfpathlineto{\pgfqpoint{0.450032in}{1.196241in}}%
\pgfpathclose%
\pgfusepath{fill}%
\end{pgfscope}%
\begin{pgfscope}%
\pgfpathrectangle{\pgfqpoint{0.407592in}{0.222904in}}{\pgfqpoint{1.411163in}{1.411163in}}%
\pgfusepath{clip}%
\pgfsetbuttcap%
\pgfsetroundjoin%
\definecolor{currentfill}{rgb}{0.749020,0.772549,0.792157}%
\pgfsetfillcolor{currentfill}%
\pgfsetfillopacity{0.160000}%
\pgfsetlinewidth{0.000000pt}%
\definecolor{currentstroke}{rgb}{0.000000,0.000000,0.000000}%
\pgfsetstrokecolor{currentstroke}%
\pgfsetdash{}{0pt}%
\pgfpathmoveto{\pgfqpoint{1.768602in}{1.191744in}}%
\pgfpathlineto{\pgfqpoint{1.782640in}{1.203771in}}%
\pgfpathlineto{\pgfqpoint{1.796697in}{1.215815in}}%
\pgfpathlineto{\pgfqpoint{1.793644in}{1.232834in}}%
\pgfpathlineto{\pgfqpoint{1.787575in}{1.249598in}}%
\pgfpathlineto{\pgfqpoint{1.773855in}{1.237024in}}%
\pgfpathlineto{\pgfqpoint{1.760147in}{1.224461in}}%
\pgfpathlineto{\pgfqpoint{1.765824in}{1.208223in}}%
\pgfpathlineto{\pgfqpoint{1.768602in}{1.191744in}}%
\pgfpathclose%
\pgfusepath{fill}%
\end{pgfscope}%
\begin{pgfscope}%
\pgfpathrectangle{\pgfqpoint{0.407592in}{0.222904in}}{\pgfqpoint{1.411163in}{1.411163in}}%
\pgfusepath{clip}%
\pgfsetbuttcap%
\pgfsetroundjoin%
\definecolor{currentfill}{rgb}{0.749020,0.772549,0.792157}%
\pgfsetfillcolor{currentfill}%
\pgfsetfillopacity{0.160000}%
\pgfsetlinewidth{0.000000pt}%
\definecolor{currentstroke}{rgb}{0.000000,0.000000,0.000000}%
\pgfsetstrokecolor{currentstroke}%
\pgfsetdash{}{0pt}%
\pgfpathmoveto{\pgfqpoint{1.570115in}{0.999993in}}%
\pgfpathlineto{\pgfqpoint{1.583889in}{1.011171in}}%
\pgfpathlineto{\pgfqpoint{1.597687in}{1.022369in}}%
\pgfpathlineto{\pgfqpoint{1.600788in}{1.035548in}}%
\pgfpathlineto{\pgfqpoint{1.601625in}{1.048682in}}%
\pgfpathlineto{\pgfqpoint{1.587832in}{1.036865in}}%
\pgfpathlineto{\pgfqpoint{1.574058in}{1.025063in}}%
\pgfpathlineto{\pgfqpoint{1.573157in}{1.012546in}}%
\pgfpathlineto{\pgfqpoint{1.570115in}{0.999993in}}%
\pgfpathclose%
\pgfusepath{fill}%
\end{pgfscope}%
\begin{pgfscope}%
\pgfpathrectangle{\pgfqpoint{0.407592in}{0.222904in}}{\pgfqpoint{1.411163in}{1.411163in}}%
\pgfusepath{clip}%
\pgfsetbuttcap%
\pgfsetroundjoin%
\definecolor{currentfill}{rgb}{0.749020,0.772549,0.792157}%
\pgfsetfillcolor{currentfill}%
\pgfsetfillopacity{0.160000}%
\pgfsetlinewidth{0.000000pt}%
\definecolor{currentstroke}{rgb}{0.000000,0.000000,0.000000}%
\pgfsetstrokecolor{currentstroke}%
\pgfsetdash{}{0pt}%
\pgfpathmoveto{\pgfqpoint{1.450479in}{0.891485in}}%
\pgfpathlineto{\pgfqpoint{1.463761in}{0.901843in}}%
\pgfpathlineto{\pgfqpoint{1.477073in}{0.912225in}}%
\pgfpathlineto{\pgfqpoint{1.483419in}{0.922744in}}%
\pgfpathlineto{\pgfqpoint{1.487990in}{0.933343in}}%
\pgfpathlineto{\pgfqpoint{1.474388in}{0.922304in}}%
\pgfpathlineto{\pgfqpoint{1.460810in}{0.911284in}}%
\pgfpathlineto{\pgfqpoint{1.456470in}{0.901345in}}%
\pgfpathlineto{\pgfqpoint{1.450479in}{0.891485in}}%
\pgfpathclose%
\pgfusepath{fill}%
\end{pgfscope}%
\begin{pgfscope}%
\pgfpathrectangle{\pgfqpoint{0.407592in}{0.222904in}}{\pgfqpoint{1.411163in}{1.411163in}}%
\pgfusepath{clip}%
\pgfsetbuttcap%
\pgfsetroundjoin%
\definecolor{currentfill}{rgb}{0.749020,0.772549,0.792157}%
\pgfsetfillcolor{currentfill}%
\pgfsetfillopacity{0.160000}%
\pgfsetlinewidth{0.000000pt}%
\definecolor{currentstroke}{rgb}{0.000000,0.000000,0.000000}%
\pgfsetstrokecolor{currentstroke}%
\pgfsetdash{}{0pt}%
\pgfpathmoveto{\pgfqpoint{0.954332in}{0.712235in}}%
\pgfpathlineto{\pgfqpoint{0.944352in}{0.719704in}}%
\pgfpathlineto{\pgfqpoint{0.934335in}{0.727200in}}%
\pgfpathlineto{\pgfqpoint{0.944514in}{0.723167in}}%
\pgfpathlineto{\pgfqpoint{0.955364in}{0.719437in}}%
\pgfpathlineto{\pgfqpoint{0.964048in}{0.712425in}}%
\pgfpathlineto{\pgfqpoint{0.972697in}{0.705441in}}%
\pgfpathlineto{\pgfqpoint{0.963224in}{0.708704in}}%
\pgfpathlineto{\pgfqpoint{0.954332in}{0.712235in}}%
\pgfpathclose%
\pgfusepath{fill}%
\end{pgfscope}%
\begin{pgfscope}%
\pgfpathrectangle{\pgfqpoint{0.407592in}{0.222904in}}{\pgfqpoint{1.411163in}{1.411163in}}%
\pgfusepath{clip}%
\pgfsetbuttcap%
\pgfsetroundjoin%
\definecolor{currentfill}{rgb}{0.749020,0.772549,0.792157}%
\pgfsetfillcolor{currentfill}%
\pgfsetfillopacity{0.160000}%
\pgfsetlinewidth{0.000000pt}%
\definecolor{currentstroke}{rgb}{0.000000,0.000000,0.000000}%
\pgfsetstrokecolor{currentstroke}%
\pgfsetdash{}{0pt}%
\pgfpathmoveto{\pgfqpoint{1.026478in}{0.682912in}}%
\pgfpathlineto{\pgfqpoint{1.021012in}{0.689118in}}%
\pgfpathlineto{\pgfqpoint{1.015522in}{0.695351in}}%
\pgfpathlineto{\pgfqpoint{1.027217in}{0.693632in}}%
\pgfpathlineto{\pgfqpoint{1.039203in}{0.692254in}}%
\pgfpathlineto{\pgfqpoint{1.042956in}{0.686244in}}%
\pgfpathlineto{\pgfqpoint{1.046692in}{0.680260in}}%
\pgfpathlineto{\pgfqpoint{1.036462in}{0.681440in}}%
\pgfpathlineto{\pgfqpoint{1.026478in}{0.682912in}}%
\pgfpathclose%
\pgfusepath{fill}%
\end{pgfscope}%
\begin{pgfscope}%
\pgfpathrectangle{\pgfqpoint{0.407592in}{0.222904in}}{\pgfqpoint{1.411163in}{1.411163in}}%
\pgfusepath{clip}%
\pgfsetbuttcap%
\pgfsetroundjoin%
\definecolor{currentfill}{rgb}{0.749020,0.772549,0.792157}%
\pgfsetfillcolor{currentfill}%
\pgfsetfillopacity{0.160000}%
\pgfsetlinewidth{0.000000pt}%
\definecolor{currentstroke}{rgb}{0.000000,0.000000,0.000000}%
\pgfsetstrokecolor{currentstroke}%
\pgfsetdash{}{0pt}%
\pgfpathmoveto{\pgfqpoint{1.110537in}{0.679531in}}%
\pgfpathlineto{\pgfqpoint{1.112280in}{0.685453in}}%
\pgfpathlineto{\pgfqpoint{1.114031in}{0.691402in}}%
\pgfpathlineto{\pgfqpoint{1.126442in}{0.692503in}}%
\pgfpathlineto{\pgfqpoint{1.138681in}{0.693950in}}%
\pgfpathlineto{\pgfqpoint{1.135119in}{0.687818in}}%
\pgfpathlineto{\pgfqpoint{1.131573in}{0.681712in}}%
\pgfpathlineto{\pgfqpoint{1.121128in}{0.680473in}}%
\pgfpathlineto{\pgfqpoint{1.110537in}{0.679531in}}%
\pgfpathclose%
\pgfusepath{fill}%
\end{pgfscope}%
\begin{pgfscope}%
\pgfpathrectangle{\pgfqpoint{0.407592in}{0.222904in}}{\pgfqpoint{1.411163in}{1.411163in}}%
\pgfusepath{clip}%
\pgfsetbuttcap%
\pgfsetroundjoin%
\definecolor{currentfill}{rgb}{0.749020,0.772549,0.792157}%
\pgfsetfillcolor{currentfill}%
\pgfsetfillopacity{0.160000}%
\pgfsetlinewidth{0.000000pt}%
\definecolor{currentstroke}{rgb}{0.000000,0.000000,0.000000}%
\pgfsetstrokecolor{currentstroke}%
\pgfsetdash{}{0pt}%
\pgfpathmoveto{\pgfqpoint{1.283933in}{0.758812in}}%
\pgfpathlineto{\pgfqpoint{1.295023in}{0.767302in}}%
\pgfpathlineto{\pgfqpoint{1.306149in}{0.775819in}}%
\pgfpathlineto{\pgfqpoint{1.316125in}{0.782146in}}%
\pgfpathlineto{\pgfqpoint{1.304505in}{0.773318in}}%
\pgfpathlineto{\pgfqpoint{1.292921in}{0.764517in}}%
\pgfusepath{fill}%
\end{pgfscope}%
\begin{pgfscope}%
\pgfpathrectangle{\pgfqpoint{0.407592in}{0.222904in}}{\pgfqpoint{1.411163in}{1.411163in}}%
\pgfusepath{clip}%
\pgfsetbuttcap%
\pgfsetroundjoin%
\definecolor{currentfill}{rgb}{0.749020,0.772549,0.792157}%
\pgfsetfillcolor{currentfill}%
\pgfsetfillopacity{0.160000}%
\pgfsetlinewidth{0.000000pt}%
\definecolor{currentstroke}{rgb}{0.000000,0.000000,0.000000}%
\pgfsetstrokecolor{currentstroke}%
\pgfsetdash{}{0pt}%
\pgfpathmoveto{\pgfqpoint{0.713628in}{0.909286in}}%
\pgfpathlineto{\pgfqpoint{0.699821in}{0.920239in}}%
\pgfpathlineto{\pgfqpoint{0.685989in}{0.931212in}}%
\pgfpathlineto{\pgfqpoint{0.688537in}{0.920626in}}%
\pgfpathlineto{\pgfqpoint{0.692881in}{0.910128in}}%
\pgfpathlineto{\pgfqpoint{0.706562in}{0.899812in}}%
\pgfpathlineto{\pgfqpoint{0.720213in}{0.889520in}}%
\pgfpathlineto{\pgfqpoint{0.716086in}{0.899359in}}%
\pgfpathlineto{\pgfqpoint{0.713628in}{0.909286in}}%
\pgfpathclose%
\pgfusepath{fill}%
\end{pgfscope}%
\begin{pgfscope}%
\pgfpathrectangle{\pgfqpoint{0.407592in}{0.222904in}}{\pgfqpoint{1.411163in}{1.411163in}}%
\pgfusepath{clip}%
\pgfsetbuttcap%
\pgfsetroundjoin%
\definecolor{currentfill}{rgb}{0.749020,0.772549,0.792157}%
\pgfsetfillcolor{currentfill}%
\pgfsetfillopacity{0.160000}%
\pgfsetlinewidth{0.000000pt}%
\definecolor{currentstroke}{rgb}{0.000000,0.000000,0.000000}%
\pgfsetstrokecolor{currentstroke}%
\pgfsetdash{}{0pt}%
\pgfpathmoveto{\pgfqpoint{1.316125in}{0.782146in}}%
\pgfpathlineto{\pgfqpoint{1.327780in}{0.791000in}}%
\pgfpathlineto{\pgfqpoint{1.339471in}{0.799881in}}%
\pgfpathlineto{\pgfqpoint{1.349275in}{0.807077in}}%
\pgfpathlineto{\pgfqpoint{1.357886in}{0.814488in}}%
\pgfpathlineto{\pgfqpoint{1.345383in}{0.804961in}}%
\pgfpathlineto{\pgfqpoint{1.332913in}{0.795459in}}%
\pgfpathlineto{\pgfqpoint{1.325058in}{0.788703in}}%
\pgfpathlineto{\pgfqpoint{1.316125in}{0.782146in}}%
\pgfpathclose%
\pgfusepath{fill}%
\end{pgfscope}%
\begin{pgfscope}%
\pgfpathrectangle{\pgfqpoint{0.407592in}{0.222904in}}{\pgfqpoint{1.411163in}{1.411163in}}%
\pgfusepath{clip}%
\pgfsetbuttcap%
\pgfsetroundjoin%
\definecolor{currentfill}{rgb}{0.749020,0.772549,0.792157}%
\pgfsetfillcolor{currentfill}%
\pgfsetfillopacity{0.160000}%
\pgfsetlinewidth{0.000000pt}%
\definecolor{currentstroke}{rgb}{0.000000,0.000000,0.000000}%
\pgfsetstrokecolor{currentstroke}%
\pgfsetdash{}{0pt}%
\pgfpathmoveto{\pgfqpoint{0.603422in}{1.022558in}}%
\pgfpathlineto{\pgfqpoint{0.589552in}{1.034297in}}%
\pgfpathlineto{\pgfqpoint{0.575662in}{1.046054in}}%
\pgfpathlineto{\pgfqpoint{0.573906in}{1.032905in}}%
\pgfpathlineto{\pgfqpoint{0.574411in}{1.019725in}}%
\pgfpathlineto{\pgfqpoint{0.588448in}{1.008590in}}%
\pgfpathlineto{\pgfqpoint{0.602459in}{0.997475in}}%
\pgfpathlineto{\pgfqpoint{0.601871in}{1.010029in}}%
\pgfpathlineto{\pgfqpoint{0.603422in}{1.022558in}}%
\pgfpathclose%
\pgfusepath{fill}%
\end{pgfscope}%
\begin{pgfscope}%
\pgfpathrectangle{\pgfqpoint{0.407592in}{0.222904in}}{\pgfqpoint{1.411163in}{1.411163in}}%
\pgfusepath{clip}%
\pgfsetbuttcap%
\pgfsetroundjoin%
\definecolor{currentfill}{rgb}{0.749020,0.772549,0.792157}%
\pgfsetfillcolor{currentfill}%
\pgfsetfillopacity{0.160000}%
\pgfsetlinewidth{0.000000pt}%
\definecolor{currentstroke}{rgb}{0.000000,0.000000,0.000000}%
\pgfsetstrokecolor{currentstroke}%
\pgfsetdash{}{0pt}%
\pgfpathmoveto{\pgfqpoint{1.046692in}{0.680260in}}%
\pgfpathlineto{\pgfqpoint{1.042956in}{0.686244in}}%
\pgfpathlineto{\pgfqpoint{1.039203in}{0.692254in}}%
\pgfpathlineto{\pgfqpoint{1.051424in}{0.691224in}}%
\pgfpathlineto{\pgfqpoint{1.063823in}{0.690546in}}%
\pgfpathlineto{\pgfqpoint{1.065766in}{0.684658in}}%
\pgfpathlineto{\pgfqpoint{1.067700in}{0.678798in}}%
\pgfpathlineto{\pgfqpoint{1.057121in}{0.679378in}}%
\pgfpathlineto{\pgfqpoint{1.046692in}{0.680260in}}%
\pgfpathclose%
\pgfusepath{fill}%
\end{pgfscope}%
\begin{pgfscope}%
\pgfpathrectangle{\pgfqpoint{0.407592in}{0.222904in}}{\pgfqpoint{1.411163in}{1.411163in}}%
\pgfusepath{clip}%
\pgfsetbuttcap%
\pgfsetroundjoin%
\definecolor{currentfill}{rgb}{0.749020,0.772549,0.792157}%
\pgfsetfillcolor{currentfill}%
\pgfsetfillopacity{0.160000}%
\pgfsetlinewidth{0.000000pt}%
\definecolor{currentstroke}{rgb}{0.000000,0.000000,0.000000}%
\pgfsetstrokecolor{currentstroke}%
\pgfsetdash{}{0pt}%
\pgfpathmoveto{\pgfqpoint{1.184795in}{0.703034in}}%
\pgfpathlineto{\pgfqpoint{1.191759in}{0.709846in}}%
\pgfpathlineto{\pgfqpoint{1.198752in}{0.716687in}}%
\pgfpathlineto{\pgfqpoint{1.210847in}{0.720160in}}%
\pgfpathlineto{\pgfqpoint{1.222376in}{0.723952in}}%
\pgfpathlineto{\pgfqpoint{1.213878in}{0.716658in}}%
\pgfpathlineto{\pgfqpoint{1.205412in}{0.709392in}}%
\pgfpathlineto{\pgfqpoint{1.195348in}{0.706073in}}%
\pgfpathlineto{\pgfqpoint{1.184795in}{0.703034in}}%
\pgfpathclose%
\pgfusepath{fill}%
\end{pgfscope}%
\begin{pgfscope}%
\pgfpathrectangle{\pgfqpoint{0.407592in}{0.222904in}}{\pgfqpoint{1.411163in}{1.411163in}}%
\pgfusepath{clip}%
\pgfsetbuttcap%
\pgfsetroundjoin%
\definecolor{currentfill}{rgb}{0.749020,0.772549,0.792157}%
\pgfsetfillcolor{currentfill}%
\pgfsetfillopacity{0.160000}%
\pgfsetlinewidth{0.000000pt}%
\definecolor{currentstroke}{rgb}{0.000000,0.000000,0.000000}%
\pgfsetstrokecolor{currentstroke}%
\pgfsetdash{}{0pt}%
\pgfpathmoveto{\pgfqpoint{0.422723in}{1.221226in}}%
\pgfpathlineto{\pgfqpoint{0.409049in}{1.233736in}}%
\pgfpathlineto{\pgfqpoint{0.395362in}{1.246258in}}%
\pgfpathlineto{\pgfqpoint{0.385772in}{1.229438in}}%
\pgfpathlineto{\pgfqpoint{0.379152in}{1.212375in}}%
\pgfpathlineto{\pgfqpoint{0.393317in}{1.200387in}}%
\pgfpathlineto{\pgfqpoint{0.407462in}{1.188415in}}%
\pgfpathlineto{\pgfqpoint{0.413663in}{1.204934in}}%
\pgfpathlineto{\pgfqpoint{0.422723in}{1.221226in}}%
\pgfpathclose%
\pgfusepath{fill}%
\end{pgfscope}%
\begin{pgfscope}%
\pgfpathrectangle{\pgfqpoint{0.407592in}{0.222904in}}{\pgfqpoint{1.411163in}{1.411163in}}%
\pgfusepath{clip}%
\pgfsetbuttcap%
\pgfsetroundjoin%
\definecolor{currentfill}{rgb}{0.749020,0.772549,0.792157}%
\pgfsetfillcolor{currentfill}%
\pgfsetfillopacity{0.160000}%
\pgfsetlinewidth{0.000000pt}%
\definecolor{currentstroke}{rgb}{0.000000,0.000000,0.000000}%
\pgfsetstrokecolor{currentstroke}%
\pgfsetdash{}{0pt}%
\pgfpathmoveto{\pgfqpoint{1.796697in}{1.215815in}}%
\pgfpathlineto{\pgfqpoint{1.810773in}{1.227874in}}%
\pgfpathlineto{\pgfqpoint{1.824868in}{1.239951in}}%
\pgfpathlineto{\pgfqpoint{1.821526in}{1.257501in}}%
\pgfpathlineto{\pgfqpoint{1.815054in}{1.274780in}}%
\pgfpathlineto{\pgfqpoint{1.801308in}{1.262183in}}%
\pgfpathlineto{\pgfqpoint{1.787575in}{1.249598in}}%
\pgfpathlineto{\pgfqpoint{1.793644in}{1.232834in}}%
\pgfpathlineto{\pgfqpoint{1.796697in}{1.215815in}}%
\pgfpathclose%
\pgfusepath{fill}%
\end{pgfscope}%
\begin{pgfscope}%
\pgfpathrectangle{\pgfqpoint{0.407592in}{0.222904in}}{\pgfqpoint{1.411163in}{1.411163in}}%
\pgfusepath{clip}%
\pgfsetbuttcap%
\pgfsetroundjoin%
\definecolor{currentfill}{rgb}{0.749020,0.772549,0.792157}%
\pgfsetfillcolor{currentfill}%
\pgfsetfillopacity{0.160000}%
\pgfsetlinewidth{0.000000pt}%
\definecolor{currentstroke}{rgb}{0.000000,0.000000,0.000000}%
\pgfsetstrokecolor{currentstroke}%
\pgfsetdash{}{0pt}%
\pgfpathmoveto{\pgfqpoint{1.089114in}{0.678553in}}%
\pgfpathlineto{\pgfqpoint{1.089018in}{0.684393in}}%
\pgfpathlineto{\pgfqpoint{1.088921in}{0.690260in}}%
\pgfpathlineto{\pgfqpoint{1.101504in}{0.690653in}}%
\pgfpathlineto{\pgfqpoint{1.114031in}{0.691402in}}%
\pgfpathlineto{\pgfqpoint{1.112280in}{0.685453in}}%
\pgfpathlineto{\pgfqpoint{1.110537in}{0.679531in}}%
\pgfpathlineto{\pgfqpoint{1.099849in}{0.678890in}}%
\pgfpathlineto{\pgfqpoint{1.089114in}{0.678553in}}%
\pgfpathclose%
\pgfusepath{fill}%
\end{pgfscope}%
\begin{pgfscope}%
\pgfpathrectangle{\pgfqpoint{0.407592in}{0.222904in}}{\pgfqpoint{1.411163in}{1.411163in}}%
\pgfusepath{clip}%
\pgfsetbuttcap%
\pgfsetroundjoin%
\definecolor{currentfill}{rgb}{0.749020,0.772549,0.792157}%
\pgfsetfillcolor{currentfill}%
\pgfsetfillopacity{0.160000}%
\pgfsetlinewidth{0.000000pt}%
\definecolor{currentstroke}{rgb}{0.000000,0.000000,0.000000}%
\pgfsetstrokecolor{currentstroke}%
\pgfsetdash{}{0pt}%
\pgfpathmoveto{\pgfqpoint{0.838112in}{0.794089in}}%
\pgfpathlineto{\pgfqpoint{0.825152in}{0.803525in}}%
\pgfpathlineto{\pgfqpoint{0.812157in}{0.812986in}}%
\pgfpathlineto{\pgfqpoint{0.819401in}{0.805615in}}%
\pgfpathlineto{\pgfqpoint{0.827881in}{0.798467in}}%
\pgfpathlineto{\pgfqpoint{0.840191in}{0.789648in}}%
\pgfpathlineto{\pgfqpoint{0.852464in}{0.780857in}}%
\pgfpathlineto{\pgfqpoint{0.844730in}{0.787370in}}%
\pgfpathlineto{\pgfqpoint{0.838112in}{0.794089in}}%
\pgfpathclose%
\pgfusepath{fill}%
\end{pgfscope}%
\begin{pgfscope}%
\pgfpathrectangle{\pgfqpoint{0.407592in}{0.222904in}}{\pgfqpoint{1.411163in}{1.411163in}}%
\pgfusepath{clip}%
\pgfsetbuttcap%
\pgfsetroundjoin%
\definecolor{currentfill}{rgb}{0.749020,0.772549,0.792157}%
\pgfsetfillcolor{currentfill}%
\pgfsetfillopacity{0.160000}%
\pgfsetlinewidth{0.000000pt}%
\definecolor{currentstroke}{rgb}{0.000000,0.000000,0.000000}%
\pgfsetstrokecolor{currentstroke}%
\pgfsetdash{}{0pt}%
\pgfpathmoveto{\pgfqpoint{1.067700in}{0.678798in}}%
\pgfpathlineto{\pgfqpoint{1.065766in}{0.684658in}}%
\pgfpathlineto{\pgfqpoint{1.063823in}{0.690546in}}%
\pgfpathlineto{\pgfqpoint{1.076341in}{0.690224in}}%
\pgfpathlineto{\pgfqpoint{1.088921in}{0.690260in}}%
\pgfpathlineto{\pgfqpoint{1.089018in}{0.684393in}}%
\pgfpathlineto{\pgfqpoint{1.089114in}{0.678553in}}%
\pgfpathlineto{\pgfqpoint{1.078381in}{0.678522in}}%
\pgfpathlineto{\pgfqpoint{1.067700in}{0.678798in}}%
\pgfpathclose%
\pgfusepath{fill}%
\end{pgfscope}%
\begin{pgfscope}%
\pgfpathrectangle{\pgfqpoint{0.407592in}{0.222904in}}{\pgfqpoint{1.411163in}{1.411163in}}%
\pgfusepath{clip}%
\pgfsetbuttcap%
\pgfsetroundjoin%
\definecolor{currentfill}{rgb}{0.749020,0.772549,0.792157}%
\pgfsetfillcolor{currentfill}%
\pgfsetfillopacity{0.160000}%
\pgfsetlinewidth{0.000000pt}%
\definecolor{currentstroke}{rgb}{0.000000,0.000000,0.000000}%
\pgfsetstrokecolor{currentstroke}%
\pgfsetdash{}{0pt}%
\pgfpathmoveto{\pgfqpoint{1.382990in}{0.833618in}}%
\pgfpathlineto{\pgfqpoint{1.395592in}{0.843221in}}%
\pgfpathlineto{\pgfqpoint{1.408228in}{0.852849in}}%
\pgfpathlineto{\pgfqpoint{1.416863in}{0.861762in}}%
\pgfpathlineto{\pgfqpoint{1.424004in}{0.870837in}}%
\pgfpathlineto{\pgfqpoint{1.410810in}{0.860548in}}%
\pgfpathlineto{\pgfqpoint{1.397646in}{0.850281in}}%
\pgfpathlineto{\pgfqpoint{1.391004in}{0.841873in}}%
\pgfpathlineto{\pgfqpoint{1.382990in}{0.833618in}}%
\pgfpathclose%
\pgfusepath{fill}%
\end{pgfscope}%
\begin{pgfscope}%
\pgfpathrectangle{\pgfqpoint{0.407592in}{0.222904in}}{\pgfqpoint{1.411163in}{1.411163in}}%
\pgfusepath{clip}%
\pgfsetbuttcap%
\pgfsetroundjoin%
\definecolor{currentfill}{rgb}{0.749020,0.772549,0.792157}%
\pgfsetfillcolor{currentfill}%
\pgfsetfillopacity{0.160000}%
\pgfsetlinewidth{0.000000pt}%
\definecolor{currentstroke}{rgb}{0.000000,0.000000,0.000000}%
\pgfsetstrokecolor{currentstroke}%
\pgfsetdash{}{0pt}%
\pgfpathmoveto{\pgfqpoint{0.916170in}{0.736098in}}%
\pgfpathlineto{\pgfqpoint{0.904974in}{0.744175in}}%
\pgfpathlineto{\pgfqpoint{0.893740in}{0.752279in}}%
\pgfpathlineto{\pgfqpoint{0.903534in}{0.747131in}}%
\pgfpathlineto{\pgfqpoint{0.914188in}{0.742278in}}%
\pgfpathlineto{\pgfqpoint{0.924281in}{0.734725in}}%
\pgfpathlineto{\pgfqpoint{0.934335in}{0.727200in}}%
\pgfpathlineto{\pgfqpoint{0.924874in}{0.731518in}}%
\pgfpathlineto{\pgfqpoint{0.916170in}{0.736098in}}%
\pgfpathclose%
\pgfusepath{fill}%
\end{pgfscope}%
\begin{pgfscope}%
\pgfpathrectangle{\pgfqpoint{0.407592in}{0.222904in}}{\pgfqpoint{1.411163in}{1.411163in}}%
\pgfusepath{clip}%
\pgfsetbuttcap%
\pgfsetroundjoin%
\definecolor{currentfill}{rgb}{0.749020,0.772549,0.792157}%
\pgfsetfillcolor{currentfill}%
\pgfsetfillopacity{0.160000}%
\pgfsetlinewidth{0.000000pt}%
\definecolor{currentstroke}{rgb}{0.000000,0.000000,0.000000}%
\pgfsetstrokecolor{currentstroke}%
\pgfsetdash{}{0pt}%
\pgfpathmoveto{\pgfqpoint{1.597687in}{1.022369in}}%
\pgfpathlineto{\pgfqpoint{1.611511in}{1.033588in}}%
\pgfpathlineto{\pgfqpoint{1.625359in}{1.044827in}}%
\pgfpathlineto{\pgfqpoint{1.628507in}{1.058622in}}%
\pgfpathlineto{\pgfqpoint{1.629267in}{1.072365in}}%
\pgfpathlineto{\pgfqpoint{1.615437in}{1.060516in}}%
\pgfpathlineto{\pgfqpoint{1.601625in}{1.048682in}}%
\pgfpathlineto{\pgfqpoint{1.600788in}{1.035548in}}%
\pgfpathlineto{\pgfqpoint{1.597687in}{1.022369in}}%
\pgfpathclose%
\pgfusepath{fill}%
\end{pgfscope}%
\begin{pgfscope}%
\pgfpathrectangle{\pgfqpoint{0.407592in}{0.222904in}}{\pgfqpoint{1.411163in}{1.411163in}}%
\pgfusepath{clip}%
\pgfsetbuttcap%
\pgfsetroundjoin%
\definecolor{currentfill}{rgb}{0.749020,0.772549,0.792157}%
\pgfsetfillcolor{currentfill}%
\pgfsetfillopacity{0.160000}%
\pgfsetlinewidth{0.000000pt}%
\definecolor{currentstroke}{rgb}{0.000000,0.000000,0.000000}%
\pgfsetstrokecolor{currentstroke}%
\pgfsetdash{}{0pt}%
\pgfpathmoveto{\pgfqpoint{0.876896in}{0.763355in}}%
\pgfpathlineto{\pgfqpoint{0.864698in}{0.772093in}}%
\pgfpathlineto{\pgfqpoint{0.852464in}{0.780857in}}%
\pgfpathlineto{\pgfqpoint{0.861284in}{0.774581in}}%
\pgfpathlineto{\pgfqpoint{0.871156in}{0.768571in}}%
\pgfpathlineto{\pgfqpoint{0.882467in}{0.760411in}}%
\pgfpathlineto{\pgfqpoint{0.893740in}{0.752279in}}%
\pgfpathlineto{\pgfqpoint{0.884849in}{0.757696in}}%
\pgfpathlineto{\pgfqpoint{0.876896in}{0.763355in}}%
\pgfpathclose%
\pgfusepath{fill}%
\end{pgfscope}%
\begin{pgfscope}%
\pgfpathrectangle{\pgfqpoint{0.407592in}{0.222904in}}{\pgfqpoint{1.411163in}{1.411163in}}%
\pgfusepath{clip}%
\pgfsetbuttcap%
\pgfsetroundjoin%
\definecolor{currentfill}{rgb}{0.749020,0.772549,0.792157}%
\pgfsetfillcolor{currentfill}%
\pgfsetfillopacity{0.160000}%
\pgfsetlinewidth{0.000000pt}%
\definecolor{currentstroke}{rgb}{0.000000,0.000000,0.000000}%
\pgfsetstrokecolor{currentstroke}%
\pgfsetdash{}{0pt}%
\pgfpathmoveto{\pgfqpoint{0.774506in}{0.848584in}}%
\pgfpathlineto{\pgfqpoint{0.760979in}{0.858783in}}%
\pgfpathlineto{\pgfqpoint{0.747421in}{0.869006in}}%
\pgfpathlineto{\pgfqpoint{0.752860in}{0.859959in}}%
\pgfpathlineto{\pgfqpoint{0.759828in}{0.851084in}}%
\pgfpathlineto{\pgfqpoint{0.772963in}{0.841521in}}%
\pgfpathlineto{\pgfqpoint{0.786063in}{0.831984in}}%
\pgfpathlineto{\pgfqpoint{0.779583in}{0.840203in}}%
\pgfpathlineto{\pgfqpoint{0.774506in}{0.848584in}}%
\pgfpathclose%
\pgfusepath{fill}%
\end{pgfscope}%
\begin{pgfscope}%
\pgfpathrectangle{\pgfqpoint{0.407592in}{0.222904in}}{\pgfqpoint{1.411163in}{1.411163in}}%
\pgfusepath{clip}%
\pgfsetbuttcap%
\pgfsetroundjoin%
\definecolor{currentfill}{rgb}{0.749020,0.772549,0.792157}%
\pgfsetfillcolor{currentfill}%
\pgfsetfillopacity{0.160000}%
\pgfsetlinewidth{0.000000pt}%
\definecolor{currentstroke}{rgb}{0.000000,0.000000,0.000000}%
\pgfsetstrokecolor{currentstroke}%
\pgfsetdash{}{0pt}%
\pgfpathmoveto{\pgfqpoint{0.972697in}{0.705441in}}%
\pgfpathlineto{\pgfqpoint{0.964048in}{0.712425in}}%
\pgfpathlineto{\pgfqpoint{0.955364in}{0.719437in}}%
\pgfpathlineto{\pgfqpoint{0.966838in}{0.716030in}}%
\pgfpathlineto{\pgfqpoint{0.978884in}{0.712962in}}%
\pgfpathlineto{\pgfqpoint{0.986066in}{0.706355in}}%
\pgfpathlineto{\pgfqpoint{0.993217in}{0.699776in}}%
\pgfpathlineto{\pgfqpoint{0.982710in}{0.702460in}}%
\pgfpathlineto{\pgfqpoint{0.972697in}{0.705441in}}%
\pgfpathclose%
\pgfusepath{fill}%
\end{pgfscope}%
\begin{pgfscope}%
\pgfpathrectangle{\pgfqpoint{0.407592in}{0.222904in}}{\pgfqpoint{1.411163in}{1.411163in}}%
\pgfusepath{clip}%
\pgfsetbuttcap%
\pgfsetroundjoin%
\definecolor{currentfill}{rgb}{0.749020,0.772549,0.792157}%
\pgfsetfillcolor{currentfill}%
\pgfsetfillopacity{0.160000}%
\pgfsetlinewidth{0.000000pt}%
\definecolor{currentstroke}{rgb}{0.000000,0.000000,0.000000}%
\pgfsetstrokecolor{currentstroke}%
\pgfsetdash{}{0pt}%
\pgfpathmoveto{\pgfqpoint{1.477073in}{0.912225in}}%
\pgfpathlineto{\pgfqpoint{1.490414in}{0.922629in}}%
\pgfpathlineto{\pgfqpoint{1.503785in}{0.933057in}}%
\pgfpathlineto{\pgfqpoint{1.510477in}{0.944229in}}%
\pgfpathlineto{\pgfqpoint{1.515267in}{0.955480in}}%
\pgfpathlineto{\pgfqpoint{1.501616in}{0.944401in}}%
\pgfpathlineto{\pgfqpoint{1.487990in}{0.933343in}}%
\pgfpathlineto{\pgfqpoint{1.483419in}{0.922744in}}%
\pgfpathlineto{\pgfqpoint{1.477073in}{0.912225in}}%
\pgfpathclose%
\pgfusepath{fill}%
\end{pgfscope}%
\begin{pgfscope}%
\pgfpathrectangle{\pgfqpoint{0.407592in}{0.222904in}}{\pgfqpoint{1.411163in}{1.411163in}}%
\pgfusepath{clip}%
\pgfsetbuttcap%
\pgfsetroundjoin%
\definecolor{currentfill}{rgb}{0.749020,0.772549,0.792157}%
\pgfsetfillcolor{currentfill}%
\pgfsetfillopacity{0.160000}%
\pgfsetlinewidth{0.000000pt}%
\definecolor{currentstroke}{rgb}{0.000000,0.000000,0.000000}%
\pgfsetstrokecolor{currentstroke}%
\pgfsetdash{}{0pt}%
\pgfpathmoveto{\pgfqpoint{1.824868in}{1.239951in}}%
\pgfpathlineto{\pgfqpoint{1.838982in}{1.252044in}}%
\pgfpathlineto{\pgfqpoint{1.853116in}{1.264153in}}%
\pgfpathlineto{\pgfqpoint{1.849471in}{1.282224in}}%
\pgfpathlineto{\pgfqpoint{1.842582in}{1.300007in}}%
\pgfpathlineto{\pgfqpoint{1.828812in}{1.287388in}}%
\pgfpathlineto{\pgfqpoint{1.815054in}{1.274780in}}%
\pgfpathlineto{\pgfqpoint{1.821526in}{1.257501in}}%
\pgfpathlineto{\pgfqpoint{1.824868in}{1.239951in}}%
\pgfpathclose%
\pgfusepath{fill}%
\end{pgfscope}%
\begin{pgfscope}%
\pgfpathrectangle{\pgfqpoint{0.407592in}{0.222904in}}{\pgfqpoint{1.411163in}{1.411163in}}%
\pgfusepath{clip}%
\pgfsetbuttcap%
\pgfsetroundjoin%
\definecolor{currentfill}{rgb}{0.749020,0.772549,0.792157}%
\pgfsetfillcolor{currentfill}%
\pgfsetfillopacity{0.160000}%
\pgfsetlinewidth{0.000000pt}%
\definecolor{currentstroke}{rgb}{0.000000,0.000000,0.000000}%
\pgfsetstrokecolor{currentstroke}%
\pgfsetdash{}{0pt}%
\pgfpathmoveto{\pgfqpoint{0.395362in}{1.246258in}}%
\pgfpathlineto{\pgfqpoint{0.381662in}{1.258792in}}%
\pgfpathlineto{\pgfqpoint{0.367948in}{1.271338in}}%
\pgfpathlineto{\pgfqpoint{0.357815in}{1.254000in}}%
\pgfpathlineto{\pgfqpoint{0.350763in}{1.236403in}}%
\pgfpathlineto{\pgfqpoint{0.364968in}{1.224381in}}%
\pgfpathlineto{\pgfqpoint{0.379152in}{1.212375in}}%
\pgfpathlineto{\pgfqpoint{0.385772in}{1.229438in}}%
\pgfpathlineto{\pgfqpoint{0.395362in}{1.246258in}}%
\pgfpathclose%
\pgfusepath{fill}%
\end{pgfscope}%
\begin{pgfscope}%
\pgfpathrectangle{\pgfqpoint{0.407592in}{0.222904in}}{\pgfqpoint{1.411163in}{1.411163in}}%
\pgfusepath{clip}%
\pgfsetbuttcap%
\pgfsetroundjoin%
\definecolor{currentfill}{rgb}{0.749020,0.772549,0.792157}%
\pgfsetfillcolor{currentfill}%
\pgfsetfillopacity{0.160000}%
\pgfsetlinewidth{0.000000pt}%
\definecolor{currentstroke}{rgb}{0.000000,0.000000,0.000000}%
\pgfsetstrokecolor{currentstroke}%
\pgfsetdash{}{0pt}%
\pgfpathmoveto{\pgfqpoint{0.575662in}{1.046054in}}%
\pgfpathlineto{\pgfqpoint{0.561752in}{1.057826in}}%
\pgfpathlineto{\pgfqpoint{0.547824in}{1.069615in}}%
\pgfpathlineto{\pgfqpoint{0.545850in}{1.055857in}}%
\pgfpathlineto{\pgfqpoint{0.546258in}{1.042058in}}%
\pgfpathlineto{\pgfqpoint{0.560348in}{1.030882in}}%
\pgfpathlineto{\pgfqpoint{0.574411in}{1.019725in}}%
\pgfpathlineto{\pgfqpoint{0.573906in}{1.032905in}}%
\pgfpathlineto{\pgfqpoint{0.575662in}{1.046054in}}%
\pgfpathclose%
\pgfusepath{fill}%
\end{pgfscope}%
\begin{pgfscope}%
\pgfpathrectangle{\pgfqpoint{0.407592in}{0.222904in}}{\pgfqpoint{1.411163in}{1.411163in}}%
\pgfusepath{clip}%
\pgfsetbuttcap%
\pgfsetroundjoin%
\definecolor{currentfill}{rgb}{0.749020,0.772549,0.792157}%
\pgfsetfillcolor{currentfill}%
\pgfsetfillopacity{0.160000}%
\pgfsetlinewidth{0.000000pt}%
\definecolor{currentstroke}{rgb}{0.000000,0.000000,0.000000}%
\pgfsetstrokecolor{currentstroke}%
\pgfsetdash{}{0pt}%
\pgfpathmoveto{\pgfqpoint{1.222376in}{0.723952in}}%
\pgfpathlineto{\pgfqpoint{1.230908in}{0.731275in}}%
\pgfpathlineto{\pgfqpoint{1.239473in}{0.738626in}}%
\pgfpathlineto{\pgfqpoint{1.251772in}{0.743227in}}%
\pgfpathlineto{\pgfqpoint{1.263315in}{0.748141in}}%
\pgfpathlineto{\pgfqpoint{1.253408in}{0.740265in}}%
\pgfpathlineto{\pgfqpoint{1.243536in}{0.732416in}}%
\pgfpathlineto{\pgfqpoint{1.233289in}{0.728044in}}%
\pgfpathlineto{\pgfqpoint{1.222376in}{0.723952in}}%
\pgfpathclose%
\pgfusepath{fill}%
\end{pgfscope}%
\begin{pgfscope}%
\pgfpathrectangle{\pgfqpoint{0.407592in}{0.222904in}}{\pgfqpoint{1.411163in}{1.411163in}}%
\pgfusepath{clip}%
\pgfsetbuttcap%
\pgfsetroundjoin%
\definecolor{currentfill}{rgb}{0.749020,0.772549,0.792157}%
\pgfsetfillcolor{currentfill}%
\pgfsetfillopacity{0.160000}%
\pgfsetlinewidth{0.000000pt}%
\definecolor{currentstroke}{rgb}{0.000000,0.000000,0.000000}%
\pgfsetstrokecolor{currentstroke}%
\pgfsetdash{}{0pt}%
\pgfpathmoveto{\pgfqpoint{1.162415in}{0.697853in}}%
\pgfpathlineto{\pgfqpoint{1.167740in}{0.704294in}}%
\pgfpathlineto{\pgfqpoint{1.173088in}{0.710763in}}%
\pgfpathlineto{\pgfqpoint{1.186146in}{0.713549in}}%
\pgfpathlineto{\pgfqpoint{1.198752in}{0.716687in}}%
\pgfpathlineto{\pgfqpoint{1.191759in}{0.709846in}}%
\pgfpathlineto{\pgfqpoint{1.184795in}{0.703034in}}%
\pgfpathlineto{\pgfqpoint{1.173801in}{0.700290in}}%
\pgfpathlineto{\pgfqpoint{1.162415in}{0.697853in}}%
\pgfpathclose%
\pgfusepath{fill}%
\end{pgfscope}%
\begin{pgfscope}%
\pgfpathrectangle{\pgfqpoint{0.407592in}{0.222904in}}{\pgfqpoint{1.411163in}{1.411163in}}%
\pgfusepath{clip}%
\pgfsetbuttcap%
\pgfsetroundjoin%
\definecolor{currentfill}{rgb}{0.749020,0.772549,0.792157}%
\pgfsetfillcolor{currentfill}%
\pgfsetfillopacity{0.160000}%
\pgfsetlinewidth{0.000000pt}%
\definecolor{currentstroke}{rgb}{0.000000,0.000000,0.000000}%
\pgfsetstrokecolor{currentstroke}%
\pgfsetdash{}{0pt}%
\pgfpathmoveto{\pgfqpoint{0.685989in}{0.931212in}}%
\pgfpathlineto{\pgfqpoint{0.672131in}{0.942205in}}%
\pgfpathlineto{\pgfqpoint{0.658248in}{0.953219in}}%
\pgfpathlineto{\pgfqpoint{0.660876in}{0.941980in}}%
\pgfpathlineto{\pgfqpoint{0.665424in}{0.930830in}}%
\pgfpathlineto{\pgfqpoint{0.679168in}{0.920467in}}%
\pgfpathlineto{\pgfqpoint{0.692881in}{0.910128in}}%
\pgfpathlineto{\pgfqpoint{0.688537in}{0.920626in}}%
\pgfpathlineto{\pgfqpoint{0.685989in}{0.931212in}}%
\pgfpathclose%
\pgfusepath{fill}%
\end{pgfscope}%
\begin{pgfscope}%
\pgfpathrectangle{\pgfqpoint{0.407592in}{0.222904in}}{\pgfqpoint{1.411163in}{1.411163in}}%
\pgfusepath{clip}%
\pgfsetbuttcap%
\pgfsetroundjoin%
\definecolor{currentfill}{rgb}{0.749020,0.772549,0.792157}%
\pgfsetfillcolor{currentfill}%
\pgfsetfillopacity{0.160000}%
\pgfsetlinewidth{0.000000pt}%
\definecolor{currentstroke}{rgb}{0.000000,0.000000,0.000000}%
\pgfsetstrokecolor{currentstroke}%
\pgfsetdash{}{0pt}%
\pgfpathmoveto{\pgfqpoint{1.853116in}{1.264153in}}%
\pgfpathlineto{\pgfqpoint{1.867269in}{1.276279in}}%
\pgfpathlineto{\pgfqpoint{1.863468in}{1.294607in}}%
\pgfpathlineto{\pgfqpoint{1.856365in}{1.312638in}}%
\pgfpathlineto{\pgfqpoint{1.842582in}{1.300007in}}%
\pgfpathlineto{\pgfqpoint{1.849471in}{1.282224in}}%
\pgfusepath{fill}%
\end{pgfscope}%
\begin{pgfscope}%
\pgfpathrectangle{\pgfqpoint{0.407592in}{0.222904in}}{\pgfqpoint{1.411163in}{1.411163in}}%
\pgfusepath{clip}%
\pgfsetbuttcap%
\pgfsetroundjoin%
\definecolor{currentfill}{rgb}{0.749020,0.772549,0.792157}%
\pgfsetfillcolor{currentfill}%
\pgfsetfillopacity{0.160000}%
\pgfsetlinewidth{0.000000pt}%
\definecolor{currentstroke}{rgb}{0.000000,0.000000,0.000000}%
\pgfsetstrokecolor{currentstroke}%
\pgfsetdash{}{0pt}%
\pgfpathmoveto{\pgfqpoint{1.263315in}{0.748141in}}%
\pgfpathlineto{\pgfqpoint{1.273257in}{0.756046in}}%
\pgfpathlineto{\pgfqpoint{1.283236in}{0.763979in}}%
\pgfpathlineto{\pgfqpoint{1.295171in}{0.769754in}}%
\pgfpathlineto{\pgfqpoint{1.306149in}{0.775819in}}%
\pgfpathlineto{\pgfqpoint{1.295023in}{0.767302in}}%
\pgfpathlineto{\pgfqpoint{1.283933in}{0.758812in}}%
\pgfpathlineto{\pgfqpoint{1.274050in}{0.753345in}}%
\pgfpathlineto{\pgfqpoint{1.263315in}{0.748141in}}%
\pgfpathclose%
\pgfusepath{fill}%
\end{pgfscope}%
\begin{pgfscope}%
\pgfpathrectangle{\pgfqpoint{0.407592in}{0.222904in}}{\pgfqpoint{1.411163in}{1.411163in}}%
\pgfusepath{clip}%
\pgfsetbuttcap%
\pgfsetroundjoin%
\definecolor{currentfill}{rgb}{0.749020,0.772549,0.792157}%
\pgfsetfillcolor{currentfill}%
\pgfsetfillopacity{0.160000}%
\pgfsetlinewidth{0.000000pt}%
\definecolor{currentstroke}{rgb}{0.000000,0.000000,0.000000}%
\pgfsetstrokecolor{currentstroke}%
\pgfsetdash{}{0pt}%
\pgfpathmoveto{\pgfqpoint{1.625359in}{1.044827in}}%
\pgfpathlineto{\pgfqpoint{1.639232in}{1.056086in}}%
\pgfpathlineto{\pgfqpoint{1.653131in}{1.067365in}}%
\pgfpathlineto{\pgfqpoint{1.656313in}{1.081769in}}%
\pgfpathlineto{\pgfqpoint{1.656983in}{1.096112in}}%
\pgfpathlineto{\pgfqpoint{1.643116in}{1.084231in}}%
\pgfpathlineto{\pgfqpoint{1.629267in}{1.072365in}}%
\pgfpathlineto{\pgfqpoint{1.628507in}{1.058622in}}%
\pgfpathlineto{\pgfqpoint{1.625359in}{1.044827in}}%
\pgfpathclose%
\pgfusepath{fill}%
\end{pgfscope}%
\begin{pgfscope}%
\pgfpathrectangle{\pgfqpoint{0.407592in}{0.222904in}}{\pgfqpoint{1.411163in}{1.411163in}}%
\pgfusepath{clip}%
\pgfsetbuttcap%
\pgfsetroundjoin%
\definecolor{currentfill}{rgb}{0.749020,0.772549,0.792157}%
\pgfsetfillcolor{currentfill}%
\pgfsetfillopacity{0.160000}%
\pgfsetlinewidth{0.000000pt}%
\definecolor{currentstroke}{rgb}{0.000000,0.000000,0.000000}%
\pgfsetstrokecolor{currentstroke}%
\pgfsetdash{}{0pt}%
\pgfpathmoveto{\pgfqpoint{0.367948in}{1.271338in}}%
\pgfpathlineto{\pgfqpoint{0.354222in}{1.283896in}}%
\pgfpathlineto{\pgfqpoint{0.340483in}{1.296466in}}%
\pgfpathlineto{\pgfqpoint{0.329792in}{1.278619in}}%
\pgfpathlineto{\pgfqpoint{0.322294in}{1.260499in}}%
\pgfpathlineto{\pgfqpoint{0.336539in}{1.248443in}}%
\pgfpathlineto{\pgfqpoint{0.350763in}{1.236403in}}%
\pgfpathlineto{\pgfqpoint{0.357815in}{1.254000in}}%
\pgfpathlineto{\pgfqpoint{0.367948in}{1.271338in}}%
\pgfpathclose%
\pgfusepath{fill}%
\end{pgfscope}%
\begin{pgfscope}%
\pgfpathrectangle{\pgfqpoint{0.407592in}{0.222904in}}{\pgfqpoint{1.411163in}{1.411163in}}%
\pgfusepath{clip}%
\pgfsetbuttcap%
\pgfsetroundjoin%
\definecolor{currentfill}{rgb}{0.749020,0.772549,0.792157}%
\pgfsetfillcolor{currentfill}%
\pgfsetfillopacity{0.160000}%
\pgfsetlinewidth{0.000000pt}%
\definecolor{currentstroke}{rgb}{0.000000,0.000000,0.000000}%
\pgfsetstrokecolor{currentstroke}%
\pgfsetdash{}{0pt}%
\pgfpathmoveto{\pgfqpoint{1.339471in}{0.799881in}}%
\pgfpathlineto{\pgfqpoint{1.351197in}{0.808790in}}%
\pgfpathlineto{\pgfqpoint{1.362959in}{0.817725in}}%
\pgfpathlineto{\pgfqpoint{1.373631in}{0.825556in}}%
\pgfpathlineto{\pgfqpoint{1.382990in}{0.833618in}}%
\pgfpathlineto{\pgfqpoint{1.370421in}{0.824040in}}%
\pgfpathlineto{\pgfqpoint{1.357886in}{0.814488in}}%
\pgfpathlineto{\pgfqpoint{1.349275in}{0.807077in}}%
\pgfpathlineto{\pgfqpoint{1.339471in}{0.799881in}}%
\pgfpathclose%
\pgfusepath{fill}%
\end{pgfscope}%
\begin{pgfscope}%
\pgfpathrectangle{\pgfqpoint{0.407592in}{0.222904in}}{\pgfqpoint{1.411163in}{1.411163in}}%
\pgfusepath{clip}%
\pgfsetbuttcap%
\pgfsetroundjoin%
\definecolor{currentfill}{rgb}{0.749020,0.772549,0.792157}%
\pgfsetfillcolor{currentfill}%
\pgfsetfillopacity{0.160000}%
\pgfsetlinewidth{0.000000pt}%
\definecolor{currentstroke}{rgb}{0.000000,0.000000,0.000000}%
\pgfsetstrokecolor{currentstroke}%
\pgfsetdash{}{0pt}%
\pgfpathmoveto{\pgfqpoint{0.993217in}{0.699776in}}%
\pgfpathlineto{\pgfqpoint{0.986066in}{0.706355in}}%
\pgfpathlineto{\pgfqpoint{0.978884in}{0.712962in}}%
\pgfpathlineto{\pgfqpoint{0.991446in}{0.710248in}}%
\pgfpathlineto{\pgfqpoint{1.004468in}{0.707902in}}%
\pgfpathlineto{\pgfqpoint{1.010007in}{0.701613in}}%
\pgfpathlineto{\pgfqpoint{1.015522in}{0.695351in}}%
\pgfpathlineto{\pgfqpoint{1.004171in}{0.697403in}}%
\pgfpathlineto{\pgfqpoint{0.993217in}{0.699776in}}%
\pgfpathclose%
\pgfusepath{fill}%
\end{pgfscope}%
\begin{pgfscope}%
\pgfpathrectangle{\pgfqpoint{0.407592in}{0.222904in}}{\pgfqpoint{1.411163in}{1.411163in}}%
\pgfusepath{clip}%
\pgfsetbuttcap%
\pgfsetroundjoin%
\definecolor{currentfill}{rgb}{0.749020,0.772549,0.792157}%
\pgfsetfillcolor{currentfill}%
\pgfsetfillopacity{0.160000}%
\pgfsetlinewidth{0.000000pt}%
\definecolor{currentstroke}{rgb}{0.000000,0.000000,0.000000}%
\pgfsetstrokecolor{currentstroke}%
\pgfsetdash{}{0pt}%
\pgfpathmoveto{\pgfqpoint{1.306149in}{0.775819in}}%
\pgfpathlineto{\pgfqpoint{1.317311in}{0.784364in}}%
\pgfpathlineto{\pgfqpoint{1.328510in}{0.792936in}}%
\pgfpathlineto{\pgfqpoint{1.339471in}{0.799881in}}%
\pgfpathlineto{\pgfqpoint{1.327780in}{0.791000in}}%
\pgfpathlineto{\pgfqpoint{1.316125in}{0.782146in}}%
\pgfusepath{fill}%
\end{pgfscope}%
\begin{pgfscope}%
\pgfpathrectangle{\pgfqpoint{0.407592in}{0.222904in}}{\pgfqpoint{1.411163in}{1.411163in}}%
\pgfusepath{clip}%
\pgfsetbuttcap%
\pgfsetroundjoin%
\definecolor{currentfill}{rgb}{0.749020,0.772549,0.792157}%
\pgfsetfillcolor{currentfill}%
\pgfsetfillopacity{0.160000}%
\pgfsetlinewidth{0.000000pt}%
\definecolor{currentstroke}{rgb}{0.000000,0.000000,0.000000}%
\pgfsetstrokecolor{currentstroke}%
\pgfsetdash{}{0pt}%
\pgfpathmoveto{\pgfqpoint{1.408228in}{0.852849in}}%
\pgfpathlineto{\pgfqpoint{1.420897in}{0.862503in}}%
\pgfpathlineto{\pgfqpoint{1.433600in}{0.872182in}}%
\pgfpathlineto{\pgfqpoint{1.442849in}{0.881749in}}%
\pgfpathlineto{\pgfqpoint{1.450479in}{0.891485in}}%
\pgfpathlineto{\pgfqpoint{1.437227in}{0.881150in}}%
\pgfpathlineto{\pgfqpoint{1.424004in}{0.870837in}}%
\pgfpathlineto{\pgfqpoint{1.416863in}{0.861762in}}%
\pgfpathlineto{\pgfqpoint{1.408228in}{0.852849in}}%
\pgfpathclose%
\pgfusepath{fill}%
\end{pgfscope}%
\begin{pgfscope}%
\pgfpathrectangle{\pgfqpoint{0.407592in}{0.222904in}}{\pgfqpoint{1.411163in}{1.411163in}}%
\pgfusepath{clip}%
\pgfsetbuttcap%
\pgfsetroundjoin%
\definecolor{currentfill}{rgb}{0.749020,0.772549,0.792157}%
\pgfsetfillcolor{currentfill}%
\pgfsetfillopacity{0.160000}%
\pgfsetlinewidth{0.000000pt}%
\definecolor{currentstroke}{rgb}{0.000000,0.000000,0.000000}%
\pgfsetstrokecolor{currentstroke}%
\pgfsetdash{}{0pt}%
\pgfpathmoveto{\pgfqpoint{0.812157in}{0.812986in}}%
\pgfpathlineto{\pgfqpoint{0.799127in}{0.822472in}}%
\pgfpathlineto{\pgfqpoint{0.786063in}{0.831984in}}%
\pgfpathlineto{\pgfqpoint{0.793926in}{0.823965in}}%
\pgfpathlineto{\pgfqpoint{0.803146in}{0.816186in}}%
\pgfpathlineto{\pgfqpoint{0.815533in}{0.807313in}}%
\pgfpathlineto{\pgfqpoint{0.827881in}{0.798467in}}%
\pgfpathlineto{\pgfqpoint{0.819401in}{0.805615in}}%
\pgfpathlineto{\pgfqpoint{0.812157in}{0.812986in}}%
\pgfpathclose%
\pgfusepath{fill}%
\end{pgfscope}%
\begin{pgfscope}%
\pgfpathrectangle{\pgfqpoint{0.407592in}{0.222904in}}{\pgfqpoint{1.411163in}{1.411163in}}%
\pgfusepath{clip}%
\pgfsetbuttcap%
\pgfsetroundjoin%
\definecolor{currentfill}{rgb}{0.749020,0.772549,0.792157}%
\pgfsetfillcolor{currentfill}%
\pgfsetfillopacity{0.160000}%
\pgfsetlinewidth{0.000000pt}%
\definecolor{currentstroke}{rgb}{0.000000,0.000000,0.000000}%
\pgfsetstrokecolor{currentstroke}%
\pgfsetdash{}{0pt}%
\pgfpathmoveto{\pgfqpoint{0.934335in}{0.727200in}}%
\pgfpathlineto{\pgfqpoint{0.924281in}{0.734725in}}%
\pgfpathlineto{\pgfqpoint{0.914188in}{0.742278in}}%
\pgfpathlineto{\pgfqpoint{0.925657in}{0.737742in}}%
\pgfpathlineto{\pgfqpoint{0.937890in}{0.733547in}}%
\pgfpathlineto{\pgfqpoint{0.946645in}{0.726478in}}%
\pgfpathlineto{\pgfqpoint{0.955364in}{0.719437in}}%
\pgfpathlineto{\pgfqpoint{0.944514in}{0.723167in}}%
\pgfpathlineto{\pgfqpoint{0.934335in}{0.727200in}}%
\pgfpathclose%
\pgfusepath{fill}%
\end{pgfscope}%
\begin{pgfscope}%
\pgfpathrectangle{\pgfqpoint{0.407592in}{0.222904in}}{\pgfqpoint{1.411163in}{1.411163in}}%
\pgfusepath{clip}%
\pgfsetbuttcap%
\pgfsetroundjoin%
\definecolor{currentfill}{rgb}{0.749020,0.772549,0.792157}%
\pgfsetfillcolor{currentfill}%
\pgfsetfillopacity{0.160000}%
\pgfsetlinewidth{0.000000pt}%
\definecolor{currentstroke}{rgb}{0.000000,0.000000,0.000000}%
\pgfsetstrokecolor{currentstroke}%
\pgfsetdash{}{0pt}%
\pgfpathmoveto{\pgfqpoint{1.503785in}{0.933057in}}%
\pgfpathlineto{\pgfqpoint{1.517186in}{0.943508in}}%
\pgfpathlineto{\pgfqpoint{1.530618in}{0.953983in}}%
\pgfpathlineto{\pgfqpoint{1.537644in}{0.965800in}}%
\pgfpathlineto{\pgfqpoint{1.542642in}{0.977696in}}%
\pgfpathlineto{\pgfqpoint{1.528942in}{0.966578in}}%
\pgfpathlineto{\pgfqpoint{1.515267in}{0.955480in}}%
\pgfpathlineto{\pgfqpoint{1.510477in}{0.944229in}}%
\pgfpathlineto{\pgfqpoint{1.503785in}{0.933057in}}%
\pgfpathclose%
\pgfusepath{fill}%
\end{pgfscope}%
\begin{pgfscope}%
\pgfpathrectangle{\pgfqpoint{0.407592in}{0.222904in}}{\pgfqpoint{1.411163in}{1.411163in}}%
\pgfusepath{clip}%
\pgfsetbuttcap%
\pgfsetroundjoin%
\definecolor{currentfill}{rgb}{0.749020,0.772549,0.792157}%
\pgfsetfillcolor{currentfill}%
\pgfsetfillopacity{0.160000}%
\pgfsetlinewidth{0.000000pt}%
\definecolor{currentstroke}{rgb}{0.000000,0.000000,0.000000}%
\pgfsetstrokecolor{currentstroke}%
\pgfsetdash{}{0pt}%
\pgfpathmoveto{\pgfqpoint{0.547824in}{1.069615in}}%
\pgfpathlineto{\pgfqpoint{0.533876in}{1.081420in}}%
\pgfpathlineto{\pgfqpoint{0.519908in}{1.093242in}}%
\pgfpathlineto{\pgfqpoint{0.517702in}{1.078882in}}%
\pgfpathlineto{\pgfqpoint{0.518002in}{1.064474in}}%
\pgfpathlineto{\pgfqpoint{0.532143in}{1.053256in}}%
\pgfpathlineto{\pgfqpoint{0.546258in}{1.042058in}}%
\pgfpathlineto{\pgfqpoint{0.545850in}{1.055857in}}%
\pgfpathlineto{\pgfqpoint{0.547824in}{1.069615in}}%
\pgfpathclose%
\pgfusepath{fill}%
\end{pgfscope}%
\begin{pgfscope}%
\pgfpathrectangle{\pgfqpoint{0.407592in}{0.222904in}}{\pgfqpoint{1.411163in}{1.411163in}}%
\pgfusepath{clip}%
\pgfsetbuttcap%
\pgfsetroundjoin%
\definecolor{currentfill}{rgb}{0.749020,0.772549,0.792157}%
\pgfsetfillcolor{currentfill}%
\pgfsetfillopacity{0.160000}%
\pgfsetlinewidth{0.000000pt}%
\definecolor{currentstroke}{rgb}{0.000000,0.000000,0.000000}%
\pgfsetstrokecolor{currentstroke}%
\pgfsetdash{}{0pt}%
\pgfpathmoveto{\pgfqpoint{1.138681in}{0.693950in}}%
\pgfpathlineto{\pgfqpoint{1.142260in}{0.700111in}}%
\pgfpathlineto{\pgfqpoint{1.145854in}{0.706299in}}%
\pgfpathlineto{\pgfqpoint{1.159636in}{0.708342in}}%
\pgfpathlineto{\pgfqpoint{1.173088in}{0.710763in}}%
\pgfpathlineto{\pgfqpoint{1.167740in}{0.704294in}}%
\pgfpathlineto{\pgfqpoint{1.162415in}{0.697853in}}%
\pgfpathlineto{\pgfqpoint{1.150691in}{0.695737in}}%
\pgfpathlineto{\pgfqpoint{1.138681in}{0.693950in}}%
\pgfpathclose%
\pgfusepath{fill}%
\end{pgfscope}%
\begin{pgfscope}%
\pgfpathrectangle{\pgfqpoint{0.407592in}{0.222904in}}{\pgfqpoint{1.411163in}{1.411163in}}%
\pgfusepath{clip}%
\pgfsetbuttcap%
\pgfsetroundjoin%
\definecolor{currentfill}{rgb}{0.749020,0.772549,0.792157}%
\pgfsetfillcolor{currentfill}%
\pgfsetfillopacity{0.160000}%
\pgfsetlinewidth{0.000000pt}%
\definecolor{currentstroke}{rgb}{0.000000,0.000000,0.000000}%
\pgfsetstrokecolor{currentstroke}%
\pgfsetdash{}{0pt}%
\pgfpathmoveto{\pgfqpoint{0.340483in}{1.296466in}}%
\pgfpathlineto{\pgfqpoint{0.326731in}{1.309048in}}%
\pgfpathlineto{\pgfqpoint{0.315756in}{1.290951in}}%
\pgfpathlineto{\pgfqpoint{0.308030in}{1.272572in}}%
\pgfpathlineto{\pgfqpoint{0.322294in}{1.260499in}}%
\pgfpathlineto{\pgfqpoint{0.329792in}{1.278619in}}%
\pgfusepath{fill}%
\end{pgfscope}%
\begin{pgfscope}%
\pgfpathrectangle{\pgfqpoint{0.407592in}{0.222904in}}{\pgfqpoint{1.411163in}{1.411163in}}%
\pgfusepath{clip}%
\pgfsetbuttcap%
\pgfsetroundjoin%
\definecolor{currentfill}{rgb}{0.749020,0.772549,0.792157}%
\pgfsetfillcolor{currentfill}%
\pgfsetfillopacity{0.160000}%
\pgfsetlinewidth{0.000000pt}%
\definecolor{currentstroke}{rgb}{0.000000,0.000000,0.000000}%
\pgfsetstrokecolor{currentstroke}%
\pgfsetdash{}{0pt}%
\pgfpathmoveto{\pgfqpoint{0.747421in}{0.869006in}}%
\pgfpathlineto{\pgfqpoint{0.733832in}{0.879251in}}%
\pgfpathlineto{\pgfqpoint{0.720213in}{0.889520in}}%
\pgfpathlineto{\pgfqpoint{0.726005in}{0.879814in}}%
\pgfpathlineto{\pgfqpoint{0.733451in}{0.870288in}}%
\pgfpathlineto{\pgfqpoint{0.746657in}{0.860673in}}%
\pgfpathlineto{\pgfqpoint{0.759828in}{0.851084in}}%
\pgfpathlineto{\pgfqpoint{0.752860in}{0.859959in}}%
\pgfpathlineto{\pgfqpoint{0.747421in}{0.869006in}}%
\pgfpathclose%
\pgfusepath{fill}%
\end{pgfscope}%
\begin{pgfscope}%
\pgfpathrectangle{\pgfqpoint{0.407592in}{0.222904in}}{\pgfqpoint{1.411163in}{1.411163in}}%
\pgfusepath{clip}%
\pgfsetbuttcap%
\pgfsetroundjoin%
\definecolor{currentfill}{rgb}{0.749020,0.772549,0.792157}%
\pgfsetfillcolor{currentfill}%
\pgfsetfillopacity{0.160000}%
\pgfsetlinewidth{0.000000pt}%
\definecolor{currentstroke}{rgb}{0.000000,0.000000,0.000000}%
\pgfsetstrokecolor{currentstroke}%
\pgfsetdash{}{0pt}%
\pgfpathmoveto{\pgfqpoint{1.015522in}{0.695351in}}%
\pgfpathlineto{\pgfqpoint{1.010007in}{0.701613in}}%
\pgfpathlineto{\pgfqpoint{1.004468in}{0.707902in}}%
\pgfpathlineto{\pgfqpoint{1.017888in}{0.705935in}}%
\pgfpathlineto{\pgfqpoint{1.031645in}{0.704359in}}%
\pgfpathlineto{\pgfqpoint{1.035433in}{0.698292in}}%
\pgfpathlineto{\pgfqpoint{1.039203in}{0.692254in}}%
\pgfpathlineto{\pgfqpoint{1.027217in}{0.693632in}}%
\pgfpathlineto{\pgfqpoint{1.015522in}{0.695351in}}%
\pgfpathclose%
\pgfusepath{fill}%
\end{pgfscope}%
\begin{pgfscope}%
\pgfpathrectangle{\pgfqpoint{0.407592in}{0.222904in}}{\pgfqpoint{1.411163in}{1.411163in}}%
\pgfusepath{clip}%
\pgfsetbuttcap%
\pgfsetroundjoin%
\definecolor{currentfill}{rgb}{0.749020,0.772549,0.792157}%
\pgfsetfillcolor{currentfill}%
\pgfsetfillopacity{0.160000}%
\pgfsetlinewidth{0.000000pt}%
\definecolor{currentstroke}{rgb}{0.000000,0.000000,0.000000}%
\pgfsetstrokecolor{currentstroke}%
\pgfsetdash{}{0pt}%
\pgfpathmoveto{\pgfqpoint{1.653131in}{1.067365in}}%
\pgfpathlineto{\pgfqpoint{1.667054in}{1.078665in}}%
\pgfpathlineto{\pgfqpoint{1.681003in}{1.089985in}}%
\pgfpathlineto{\pgfqpoint{1.684207in}{1.104989in}}%
\pgfpathlineto{\pgfqpoint{1.684775in}{1.119923in}}%
\pgfpathlineto{\pgfqpoint{1.670870in}{1.108009in}}%
\pgfpathlineto{\pgfqpoint{1.656983in}{1.096112in}}%
\pgfpathlineto{\pgfqpoint{1.656313in}{1.081769in}}%
\pgfpathlineto{\pgfqpoint{1.653131in}{1.067365in}}%
\pgfpathclose%
\pgfusepath{fill}%
\end{pgfscope}%
\begin{pgfscope}%
\pgfpathrectangle{\pgfqpoint{0.407592in}{0.222904in}}{\pgfqpoint{1.411163in}{1.411163in}}%
\pgfusepath{clip}%
\pgfsetbuttcap%
\pgfsetroundjoin%
\definecolor{currentfill}{rgb}{0.749020,0.772549,0.792157}%
\pgfsetfillcolor{currentfill}%
\pgfsetfillopacity{0.160000}%
\pgfsetlinewidth{0.000000pt}%
\definecolor{currentstroke}{rgb}{0.000000,0.000000,0.000000}%
\pgfsetstrokecolor{currentstroke}%
\pgfsetdash{}{0pt}%
\pgfpathmoveto{\pgfqpoint{0.658248in}{0.953219in}}%
\pgfpathlineto{\pgfqpoint{0.644339in}{0.964252in}}%
\pgfpathlineto{\pgfqpoint{0.630405in}{0.975306in}}%
\pgfpathlineto{\pgfqpoint{0.633100in}{0.963422in}}%
\pgfpathlineto{\pgfqpoint{0.637842in}{0.951627in}}%
\pgfpathlineto{\pgfqpoint{0.651648in}{0.941217in}}%
\pgfpathlineto{\pgfqpoint{0.665424in}{0.930830in}}%
\pgfpathlineto{\pgfqpoint{0.660876in}{0.941980in}}%
\pgfpathlineto{\pgfqpoint{0.658248in}{0.953219in}}%
\pgfpathclose%
\pgfusepath{fill}%
\end{pgfscope}%
\begin{pgfscope}%
\pgfpathrectangle{\pgfqpoint{0.407592in}{0.222904in}}{\pgfqpoint{1.411163in}{1.411163in}}%
\pgfusepath{clip}%
\pgfsetbuttcap%
\pgfsetroundjoin%
\definecolor{currentfill}{rgb}{0.749020,0.772549,0.792157}%
\pgfsetfillcolor{currentfill}%
\pgfsetfillopacity{0.160000}%
\pgfsetlinewidth{0.000000pt}%
\definecolor{currentstroke}{rgb}{0.000000,0.000000,0.000000}%
\pgfsetstrokecolor{currentstroke}%
\pgfsetdash{}{0pt}%
\pgfpathmoveto{\pgfqpoint{1.198752in}{0.716687in}}%
\pgfpathlineto{\pgfqpoint{1.205774in}{0.723556in}}%
\pgfpathlineto{\pgfqpoint{1.212825in}{0.730453in}}%
\pgfpathlineto{\pgfqpoint{1.226471in}{0.734361in}}%
\pgfpathlineto{\pgfqpoint{1.239473in}{0.738626in}}%
\pgfpathlineto{\pgfqpoint{1.230908in}{0.731275in}}%
\pgfpathlineto{\pgfqpoint{1.222376in}{0.723952in}}%
\pgfpathlineto{\pgfqpoint{1.210847in}{0.720160in}}%
\pgfpathlineto{\pgfqpoint{1.198752in}{0.716687in}}%
\pgfpathclose%
\pgfusepath{fill}%
\end{pgfscope}%
\begin{pgfscope}%
\pgfpathrectangle{\pgfqpoint{0.407592in}{0.222904in}}{\pgfqpoint{1.411163in}{1.411163in}}%
\pgfusepath{clip}%
\pgfsetbuttcap%
\pgfsetroundjoin%
\definecolor{currentfill}{rgb}{0.749020,0.772549,0.792157}%
\pgfsetfillcolor{currentfill}%
\pgfsetfillopacity{0.160000}%
\pgfsetlinewidth{0.000000pt}%
\definecolor{currentstroke}{rgb}{0.000000,0.000000,0.000000}%
\pgfsetstrokecolor{currentstroke}%
\pgfsetdash{}{0pt}%
\pgfpathmoveto{\pgfqpoint{0.893740in}{0.752279in}}%
\pgfpathlineto{\pgfqpoint{0.882467in}{0.760411in}}%
\pgfpathlineto{\pgfqpoint{0.871156in}{0.768571in}}%
\pgfpathlineto{\pgfqpoint{0.882039in}{0.762858in}}%
\pgfpathlineto{\pgfqpoint{0.893888in}{0.757470in}}%
\pgfpathlineto{\pgfqpoint{0.904057in}{0.749860in}}%
\pgfpathlineto{\pgfqpoint{0.914188in}{0.742278in}}%
\pgfpathlineto{\pgfqpoint{0.903534in}{0.747131in}}%
\pgfpathlineto{\pgfqpoint{0.893740in}{0.752279in}}%
\pgfpathclose%
\pgfusepath{fill}%
\end{pgfscope}%
\begin{pgfscope}%
\pgfpathrectangle{\pgfqpoint{0.407592in}{0.222904in}}{\pgfqpoint{1.411163in}{1.411163in}}%
\pgfusepath{clip}%
\pgfsetbuttcap%
\pgfsetroundjoin%
\definecolor{currentfill}{rgb}{0.749020,0.772549,0.792157}%
\pgfsetfillcolor{currentfill}%
\pgfsetfillopacity{0.160000}%
\pgfsetlinewidth{0.000000pt}%
\definecolor{currentstroke}{rgb}{0.000000,0.000000,0.000000}%
\pgfsetstrokecolor{currentstroke}%
\pgfsetdash{}{0pt}%
\pgfpathmoveto{\pgfqpoint{0.852464in}{0.780857in}}%
\pgfpathlineto{\pgfqpoint{0.840191in}{0.789648in}}%
\pgfpathlineto{\pgfqpoint{0.827881in}{0.798467in}}%
\pgfpathlineto{\pgfqpoint{0.837565in}{0.791576in}}%
\pgfpathlineto{\pgfqpoint{0.848416in}{0.784976in}}%
\pgfpathlineto{\pgfqpoint{0.859805in}{0.776759in}}%
\pgfpathlineto{\pgfqpoint{0.871156in}{0.768571in}}%
\pgfpathlineto{\pgfqpoint{0.861284in}{0.774581in}}%
\pgfpathlineto{\pgfqpoint{0.852464in}{0.780857in}}%
\pgfpathclose%
\pgfusepath{fill}%
\end{pgfscope}%
\begin{pgfscope}%
\pgfpathrectangle{\pgfqpoint{0.407592in}{0.222904in}}{\pgfqpoint{1.411163in}{1.411163in}}%
\pgfusepath{clip}%
\pgfsetbuttcap%
\pgfsetroundjoin%
\definecolor{currentfill}{rgb}{0.749020,0.772549,0.792157}%
\pgfsetfillcolor{currentfill}%
\pgfsetfillopacity{0.160000}%
\pgfsetlinewidth{0.000000pt}%
\definecolor{currentstroke}{rgb}{0.000000,0.000000,0.000000}%
\pgfsetstrokecolor{currentstroke}%
\pgfsetdash{}{0pt}%
\pgfpathmoveto{\pgfqpoint{1.114031in}{0.691402in}}%
\pgfpathlineto{\pgfqpoint{1.115790in}{0.697379in}}%
\pgfpathlineto{\pgfqpoint{1.117557in}{0.703384in}}%
\pgfpathlineto{\pgfqpoint{1.131806in}{0.704643in}}%
\pgfpathlineto{\pgfqpoint{1.145854in}{0.706299in}}%
\pgfpathlineto{\pgfqpoint{1.142260in}{0.700111in}}%
\pgfpathlineto{\pgfqpoint{1.138681in}{0.693950in}}%
\pgfpathlineto{\pgfqpoint{1.126442in}{0.692503in}}%
\pgfpathlineto{\pgfqpoint{1.114031in}{0.691402in}}%
\pgfpathclose%
\pgfusepath{fill}%
\end{pgfscope}%
\begin{pgfscope}%
\pgfpathrectangle{\pgfqpoint{0.407592in}{0.222904in}}{\pgfqpoint{1.411163in}{1.411163in}}%
\pgfusepath{clip}%
\pgfsetbuttcap%
\pgfsetroundjoin%
\definecolor{currentfill}{rgb}{0.749020,0.772549,0.792157}%
\pgfsetfillcolor{currentfill}%
\pgfsetfillopacity{0.160000}%
\pgfsetlinewidth{0.000000pt}%
\definecolor{currentstroke}{rgb}{0.000000,0.000000,0.000000}%
\pgfsetstrokecolor{currentstroke}%
\pgfsetdash{}{0pt}%
\pgfpathmoveto{\pgfqpoint{1.039203in}{0.692254in}}%
\pgfpathlineto{\pgfqpoint{1.035433in}{0.698292in}}%
\pgfpathlineto{\pgfqpoint{1.031645in}{0.704359in}}%
\pgfpathlineto{\pgfqpoint{1.045674in}{0.703179in}}%
\pgfpathlineto{\pgfqpoint{1.059909in}{0.702404in}}%
\pgfpathlineto{\pgfqpoint{1.061870in}{0.696461in}}%
\pgfpathlineto{\pgfqpoint{1.063823in}{0.690546in}}%
\pgfpathlineto{\pgfqpoint{1.051424in}{0.691224in}}%
\pgfpathlineto{\pgfqpoint{1.039203in}{0.692254in}}%
\pgfpathclose%
\pgfusepath{fill}%
\end{pgfscope}%
\begin{pgfscope}%
\pgfpathrectangle{\pgfqpoint{0.407592in}{0.222904in}}{\pgfqpoint{1.411163in}{1.411163in}}%
\pgfusepath{clip}%
\pgfsetbuttcap%
\pgfsetroundjoin%
\definecolor{currentfill}{rgb}{0.749020,0.772549,0.792157}%
\pgfsetfillcolor{currentfill}%
\pgfsetfillopacity{0.160000}%
\pgfsetlinewidth{0.000000pt}%
\definecolor{currentstroke}{rgb}{0.000000,0.000000,0.000000}%
\pgfsetstrokecolor{currentstroke}%
\pgfsetdash{}{0pt}%
\pgfpathmoveto{\pgfqpoint{1.088921in}{0.690260in}}%
\pgfpathlineto{\pgfqpoint{1.088825in}{0.696154in}}%
\pgfpathlineto{\pgfqpoint{1.088728in}{0.702077in}}%
\pgfpathlineto{\pgfqpoint{1.103175in}{0.702527in}}%
\pgfpathlineto{\pgfqpoint{1.117557in}{0.703384in}}%
\pgfpathlineto{\pgfqpoint{1.115790in}{0.697379in}}%
\pgfpathlineto{\pgfqpoint{1.114031in}{0.691402in}}%
\pgfpathlineto{\pgfqpoint{1.101504in}{0.690653in}}%
\pgfpathlineto{\pgfqpoint{1.088921in}{0.690260in}}%
\pgfpathclose%
\pgfusepath{fill}%
\end{pgfscope}%
\begin{pgfscope}%
\pgfpathrectangle{\pgfqpoint{0.407592in}{0.222904in}}{\pgfqpoint{1.411163in}{1.411163in}}%
\pgfusepath{clip}%
\pgfsetbuttcap%
\pgfsetroundjoin%
\definecolor{currentfill}{rgb}{0.749020,0.772549,0.792157}%
\pgfsetfillcolor{currentfill}%
\pgfsetfillopacity{0.160000}%
\pgfsetlinewidth{0.000000pt}%
\definecolor{currentstroke}{rgb}{0.000000,0.000000,0.000000}%
\pgfsetstrokecolor{currentstroke}%
\pgfsetdash{}{0pt}%
\pgfpathmoveto{\pgfqpoint{0.519908in}{1.093242in}}%
\pgfpathlineto{\pgfqpoint{0.505921in}{1.105081in}}%
\pgfpathlineto{\pgfqpoint{0.491914in}{1.116936in}}%
\pgfpathlineto{\pgfqpoint{0.489463in}{1.101983in}}%
\pgfpathlineto{\pgfqpoint{0.489640in}{1.086973in}}%
\pgfpathlineto{\pgfqpoint{0.503834in}{1.075713in}}%
\pgfpathlineto{\pgfqpoint{0.518002in}{1.064474in}}%
\pgfpathlineto{\pgfqpoint{0.517702in}{1.078882in}}%
\pgfpathlineto{\pgfqpoint{0.519908in}{1.093242in}}%
\pgfpathclose%
\pgfusepath{fill}%
\end{pgfscope}%
\begin{pgfscope}%
\pgfpathrectangle{\pgfqpoint{0.407592in}{0.222904in}}{\pgfqpoint{1.411163in}{1.411163in}}%
\pgfusepath{clip}%
\pgfsetbuttcap%
\pgfsetroundjoin%
\definecolor{currentfill}{rgb}{0.749020,0.772549,0.792157}%
\pgfsetfillcolor{currentfill}%
\pgfsetfillopacity{0.160000}%
\pgfsetlinewidth{0.000000pt}%
\definecolor{currentstroke}{rgb}{0.000000,0.000000,0.000000}%
\pgfsetstrokecolor{currentstroke}%
\pgfsetdash{}{0pt}%
\pgfpathmoveto{\pgfqpoint{1.063823in}{0.690546in}}%
\pgfpathlineto{\pgfqpoint{1.061870in}{0.696461in}}%
\pgfpathlineto{\pgfqpoint{1.059909in}{0.702404in}}%
\pgfpathlineto{\pgfqpoint{1.074283in}{0.702035in}}%
\pgfpathlineto{\pgfqpoint{1.088728in}{0.702077in}}%
\pgfpathlineto{\pgfqpoint{1.088825in}{0.696154in}}%
\pgfpathlineto{\pgfqpoint{1.088921in}{0.690260in}}%
\pgfpathlineto{\pgfqpoint{1.076341in}{0.690224in}}%
\pgfpathlineto{\pgfqpoint{1.063823in}{0.690546in}}%
\pgfpathclose%
\pgfusepath{fill}%
\end{pgfscope}%
\begin{pgfscope}%
\pgfpathrectangle{\pgfqpoint{0.407592in}{0.222904in}}{\pgfqpoint{1.411163in}{1.411163in}}%
\pgfusepath{clip}%
\pgfsetbuttcap%
\pgfsetroundjoin%
\definecolor{currentfill}{rgb}{0.749020,0.772549,0.792157}%
\pgfsetfillcolor{currentfill}%
\pgfsetfillopacity{0.160000}%
\pgfsetlinewidth{0.000000pt}%
\definecolor{currentstroke}{rgb}{0.000000,0.000000,0.000000}%
\pgfsetstrokecolor{currentstroke}%
\pgfsetdash{}{0pt}%
\pgfpathmoveto{\pgfqpoint{1.530618in}{0.953983in}}%
\pgfpathlineto{\pgfqpoint{1.544079in}{0.964481in}}%
\pgfpathlineto{\pgfqpoint{1.557570in}{0.975003in}}%
\pgfpathlineto{\pgfqpoint{1.564920in}{0.987459in}}%
\pgfpathlineto{\pgfqpoint{1.570115in}{0.999993in}}%
\pgfpathlineto{\pgfqpoint{1.556366in}{0.988834in}}%
\pgfpathlineto{\pgfqpoint{1.542642in}{0.977696in}}%
\pgfpathlineto{\pgfqpoint{1.537644in}{0.965800in}}%
\pgfpathlineto{\pgfqpoint{1.530618in}{0.953983in}}%
\pgfpathclose%
\pgfusepath{fill}%
\end{pgfscope}%
\begin{pgfscope}%
\pgfpathrectangle{\pgfqpoint{0.407592in}{0.222904in}}{\pgfqpoint{1.411163in}{1.411163in}}%
\pgfusepath{clip}%
\pgfsetbuttcap%
\pgfsetroundjoin%
\definecolor{currentfill}{rgb}{0.749020,0.772549,0.792157}%
\pgfsetfillcolor{currentfill}%
\pgfsetfillopacity{0.160000}%
\pgfsetlinewidth{0.000000pt}%
\definecolor{currentstroke}{rgb}{0.000000,0.000000,0.000000}%
\pgfsetstrokecolor{currentstroke}%
\pgfsetdash{}{0pt}%
\pgfpathmoveto{\pgfqpoint{0.955364in}{0.719437in}}%
\pgfpathlineto{\pgfqpoint{0.946645in}{0.726478in}}%
\pgfpathlineto{\pgfqpoint{0.937890in}{0.733547in}}%
\pgfpathlineto{\pgfqpoint{0.950833in}{0.729714in}}%
\pgfpathlineto{\pgfqpoint{0.964427in}{0.726261in}}%
\pgfpathlineto{\pgfqpoint{0.971671in}{0.719597in}}%
\pgfpathlineto{\pgfqpoint{0.978884in}{0.712962in}}%
\pgfpathlineto{\pgfqpoint{0.966838in}{0.716030in}}%
\pgfpathlineto{\pgfqpoint{0.955364in}{0.719437in}}%
\pgfpathclose%
\pgfusepath{fill}%
\end{pgfscope}%
\begin{pgfscope}%
\pgfpathrectangle{\pgfqpoint{0.407592in}{0.222904in}}{\pgfqpoint{1.411163in}{1.411163in}}%
\pgfusepath{clip}%
\pgfsetbuttcap%
\pgfsetroundjoin%
\definecolor{currentfill}{rgb}{0.749020,0.772549,0.792157}%
\pgfsetfillcolor{currentfill}%
\pgfsetfillopacity{0.160000}%
\pgfsetlinewidth{0.000000pt}%
\definecolor{currentstroke}{rgb}{0.000000,0.000000,0.000000}%
\pgfsetstrokecolor{currentstroke}%
\pgfsetdash{}{0pt}%
\pgfpathmoveto{\pgfqpoint{1.433600in}{0.872182in}}%
\pgfpathlineto{\pgfqpoint{1.446336in}{0.881888in}}%
\pgfpathlineto{\pgfqpoint{1.459107in}{0.891619in}}%
\pgfpathlineto{\pgfqpoint{1.468962in}{0.901833in}}%
\pgfpathlineto{\pgfqpoint{1.477073in}{0.912225in}}%
\pgfpathlineto{\pgfqpoint{1.463761in}{0.901843in}}%
\pgfpathlineto{\pgfqpoint{1.450479in}{0.891485in}}%
\pgfpathlineto{\pgfqpoint{1.442849in}{0.881749in}}%
\pgfpathlineto{\pgfqpoint{1.433600in}{0.872182in}}%
\pgfpathclose%
\pgfusepath{fill}%
\end{pgfscope}%
\begin{pgfscope}%
\pgfpathrectangle{\pgfqpoint{0.407592in}{0.222904in}}{\pgfqpoint{1.411163in}{1.411163in}}%
\pgfusepath{clip}%
\pgfsetbuttcap%
\pgfsetroundjoin%
\definecolor{currentfill}{rgb}{0.749020,0.772549,0.792157}%
\pgfsetfillcolor{currentfill}%
\pgfsetfillopacity{0.160000}%
\pgfsetlinewidth{0.000000pt}%
\definecolor{currentstroke}{rgb}{0.000000,0.000000,0.000000}%
\pgfsetstrokecolor{currentstroke}%
\pgfsetdash{}{0pt}%
\pgfpathmoveto{\pgfqpoint{1.239473in}{0.738626in}}%
\pgfpathlineto{\pgfqpoint{1.248071in}{0.746005in}}%
\pgfpathlineto{\pgfqpoint{1.256703in}{0.753414in}}%
\pgfpathlineto{\pgfqpoint{1.270394in}{0.758523in}}%
\pgfpathlineto{\pgfqpoint{1.283236in}{0.763979in}}%
\pgfpathlineto{\pgfqpoint{1.273257in}{0.756046in}}%
\pgfpathlineto{\pgfqpoint{1.263315in}{0.748141in}}%
\pgfpathlineto{\pgfqpoint{1.251772in}{0.743227in}}%
\pgfpathlineto{\pgfqpoint{1.239473in}{0.738626in}}%
\pgfpathclose%
\pgfusepath{fill}%
\end{pgfscope}%
\begin{pgfscope}%
\pgfpathrectangle{\pgfqpoint{0.407592in}{0.222904in}}{\pgfqpoint{1.411163in}{1.411163in}}%
\pgfusepath{clip}%
\pgfsetbuttcap%
\pgfsetroundjoin%
\definecolor{currentfill}{rgb}{0.749020,0.772549,0.792157}%
\pgfsetfillcolor{currentfill}%
\pgfsetfillopacity{0.160000}%
\pgfsetlinewidth{0.000000pt}%
\definecolor{currentstroke}{rgb}{0.000000,0.000000,0.000000}%
\pgfsetstrokecolor{currentstroke}%
\pgfsetdash{}{0pt}%
\pgfpathmoveto{\pgfqpoint{1.681003in}{1.089985in}}%
\pgfpathlineto{\pgfqpoint{1.694977in}{1.101326in}}%
\pgfpathlineto{\pgfqpoint{1.708976in}{1.112688in}}%
\pgfpathlineto{\pgfqpoint{1.712189in}{1.128283in}}%
\pgfpathlineto{\pgfqpoint{1.712642in}{1.143798in}}%
\pgfpathlineto{\pgfqpoint{1.698699in}{1.131853in}}%
\pgfpathlineto{\pgfqpoint{1.684775in}{1.119923in}}%
\pgfpathlineto{\pgfqpoint{1.684207in}{1.104989in}}%
\pgfpathlineto{\pgfqpoint{1.681003in}{1.089985in}}%
\pgfpathclose%
\pgfusepath{fill}%
\end{pgfscope}%
\begin{pgfscope}%
\pgfpathrectangle{\pgfqpoint{0.407592in}{0.222904in}}{\pgfqpoint{1.411163in}{1.411163in}}%
\pgfusepath{clip}%
\pgfsetbuttcap%
\pgfsetroundjoin%
\definecolor{currentfill}{rgb}{0.749020,0.772549,0.792157}%
\pgfsetfillcolor{currentfill}%
\pgfsetfillopacity{0.160000}%
\pgfsetlinewidth{0.000000pt}%
\definecolor{currentstroke}{rgb}{0.000000,0.000000,0.000000}%
\pgfsetstrokecolor{currentstroke}%
\pgfsetdash{}{0pt}%
\pgfpathmoveto{\pgfqpoint{1.362959in}{0.817725in}}%
\pgfpathlineto{\pgfqpoint{1.374757in}{0.826688in}}%
\pgfpathlineto{\pgfqpoint{1.386591in}{0.835679in}}%
\pgfpathlineto{\pgfqpoint{1.398125in}{0.844141in}}%
\pgfpathlineto{\pgfqpoint{1.408228in}{0.852849in}}%
\pgfpathlineto{\pgfqpoint{1.395592in}{0.843221in}}%
\pgfpathlineto{\pgfqpoint{1.382990in}{0.833618in}}%
\pgfpathlineto{\pgfqpoint{1.373631in}{0.825556in}}%
\pgfpathlineto{\pgfqpoint{1.362959in}{0.817725in}}%
\pgfpathclose%
\pgfusepath{fill}%
\end{pgfscope}%
\begin{pgfscope}%
\pgfpathrectangle{\pgfqpoint{0.407592in}{0.222904in}}{\pgfqpoint{1.411163in}{1.411163in}}%
\pgfusepath{clip}%
\pgfsetbuttcap%
\pgfsetroundjoin%
\definecolor{currentfill}{rgb}{0.749020,0.772549,0.792157}%
\pgfsetfillcolor{currentfill}%
\pgfsetfillopacity{0.160000}%
\pgfsetlinewidth{0.000000pt}%
\definecolor{currentstroke}{rgb}{0.000000,0.000000,0.000000}%
\pgfsetstrokecolor{currentstroke}%
\pgfsetdash{}{0pt}%
\pgfpathmoveto{\pgfqpoint{1.283236in}{0.763979in}}%
\pgfpathlineto{\pgfqpoint{1.293250in}{0.771941in}}%
\pgfpathlineto{\pgfqpoint{1.303300in}{0.779931in}}%
\pgfpathlineto{\pgfqpoint{1.316436in}{0.786275in}}%
\pgfpathlineto{\pgfqpoint{1.328510in}{0.792936in}}%
\pgfpathlineto{\pgfqpoint{1.317311in}{0.784364in}}%
\pgfpathlineto{\pgfqpoint{1.306149in}{0.775819in}}%
\pgfpathlineto{\pgfqpoint{1.295171in}{0.769754in}}%
\pgfpathlineto{\pgfqpoint{1.283236in}{0.763979in}}%
\pgfpathclose%
\pgfusepath{fill}%
\end{pgfscope}%
\begin{pgfscope}%
\pgfpathrectangle{\pgfqpoint{0.407592in}{0.222904in}}{\pgfqpoint{1.411163in}{1.411163in}}%
\pgfusepath{clip}%
\pgfsetbuttcap%
\pgfsetroundjoin%
\definecolor{currentfill}{rgb}{0.749020,0.772549,0.792157}%
\pgfsetfillcolor{currentfill}%
\pgfsetfillopacity{0.160000}%
\pgfsetlinewidth{0.000000pt}%
\definecolor{currentstroke}{rgb}{0.000000,0.000000,0.000000}%
\pgfsetstrokecolor{currentstroke}%
\pgfsetdash{}{0pt}%
\pgfpathmoveto{\pgfqpoint{0.720213in}{0.889520in}}%
\pgfpathlineto{\pgfqpoint{0.706562in}{0.899812in}}%
\pgfpathlineto{\pgfqpoint{0.692881in}{0.910128in}}%
\pgfpathlineto{\pgfqpoint{0.699016in}{0.899768in}}%
\pgfpathlineto{\pgfqpoint{0.706932in}{0.889595in}}%
\pgfpathlineto{\pgfqpoint{0.720209in}{0.879928in}}%
\pgfpathlineto{\pgfqpoint{0.733451in}{0.870288in}}%
\pgfpathlineto{\pgfqpoint{0.726005in}{0.879814in}}%
\pgfpathlineto{\pgfqpoint{0.720213in}{0.889520in}}%
\pgfpathclose%
\pgfusepath{fill}%
\end{pgfscope}%
\begin{pgfscope}%
\pgfpathrectangle{\pgfqpoint{0.407592in}{0.222904in}}{\pgfqpoint{1.411163in}{1.411163in}}%
\pgfusepath{clip}%
\pgfsetbuttcap%
\pgfsetroundjoin%
\definecolor{currentfill}{rgb}{0.749020,0.772549,0.792157}%
\pgfsetfillcolor{currentfill}%
\pgfsetfillopacity{0.160000}%
\pgfsetlinewidth{0.000000pt}%
\definecolor{currentstroke}{rgb}{0.000000,0.000000,0.000000}%
\pgfsetstrokecolor{currentstroke}%
\pgfsetdash{}{0pt}%
\pgfpathmoveto{\pgfqpoint{0.630405in}{0.975306in}}%
\pgfpathlineto{\pgfqpoint{0.616445in}{0.986380in}}%
\pgfpathlineto{\pgfqpoint{0.602459in}{0.997475in}}%
\pgfpathlineto{\pgfqpoint{0.605209in}{0.984952in}}%
\pgfpathlineto{\pgfqpoint{0.610133in}{0.972519in}}%
\pgfpathlineto{\pgfqpoint{0.624003in}{0.962061in}}%
\pgfpathlineto{\pgfqpoint{0.637842in}{0.951627in}}%
\pgfpathlineto{\pgfqpoint{0.633100in}{0.963422in}}%
\pgfpathlineto{\pgfqpoint{0.630405in}{0.975306in}}%
\pgfpathclose%
\pgfusepath{fill}%
\end{pgfscope}%
\begin{pgfscope}%
\pgfpathrectangle{\pgfqpoint{0.407592in}{0.222904in}}{\pgfqpoint{1.411163in}{1.411163in}}%
\pgfusepath{clip}%
\pgfsetbuttcap%
\pgfsetroundjoin%
\definecolor{currentfill}{rgb}{0.749020,0.772549,0.792157}%
\pgfsetfillcolor{currentfill}%
\pgfsetfillopacity{0.160000}%
\pgfsetlinewidth{0.000000pt}%
\definecolor{currentstroke}{rgb}{0.000000,0.000000,0.000000}%
\pgfsetstrokecolor{currentstroke}%
\pgfsetdash{}{0pt}%
\pgfpathmoveto{\pgfqpoint{1.328510in}{0.792936in}}%
\pgfpathlineto{\pgfqpoint{1.339744in}{0.801536in}}%
\pgfpathlineto{\pgfqpoint{1.351015in}{0.810164in}}%
\pgfpathlineto{\pgfqpoint{1.362959in}{0.817725in}}%
\pgfpathlineto{\pgfqpoint{1.351197in}{0.808790in}}%
\pgfpathlineto{\pgfqpoint{1.339471in}{0.799881in}}%
\pgfusepath{fill}%
\end{pgfscope}%
\begin{pgfscope}%
\pgfpathrectangle{\pgfqpoint{0.407592in}{0.222904in}}{\pgfqpoint{1.411163in}{1.411163in}}%
\pgfusepath{clip}%
\pgfsetbuttcap%
\pgfsetroundjoin%
\definecolor{currentfill}{rgb}{0.749020,0.772549,0.792157}%
\pgfsetfillcolor{currentfill}%
\pgfsetfillopacity{0.160000}%
\pgfsetlinewidth{0.000000pt}%
\definecolor{currentstroke}{rgb}{0.000000,0.000000,0.000000}%
\pgfsetstrokecolor{currentstroke}%
\pgfsetdash{}{0pt}%
\pgfpathmoveto{\pgfqpoint{0.786063in}{0.831984in}}%
\pgfpathlineto{\pgfqpoint{0.772963in}{0.841521in}}%
\pgfpathlineto{\pgfqpoint{0.759828in}{0.851084in}}%
\pgfpathlineto{\pgfqpoint{0.768303in}{0.842422in}}%
\pgfpathlineto{\pgfqpoint{0.778258in}{0.834014in}}%
\pgfpathlineto{\pgfqpoint{0.790721in}{0.825086in}}%
\pgfpathlineto{\pgfqpoint{0.803146in}{0.816186in}}%
\pgfpathlineto{\pgfqpoint{0.793926in}{0.823965in}}%
\pgfpathlineto{\pgfqpoint{0.786063in}{0.831984in}}%
\pgfpathclose%
\pgfusepath{fill}%
\end{pgfscope}%
\begin{pgfscope}%
\pgfpathrectangle{\pgfqpoint{0.407592in}{0.222904in}}{\pgfqpoint{1.411163in}{1.411163in}}%
\pgfusepath{clip}%
\pgfsetbuttcap%
\pgfsetroundjoin%
\definecolor{currentfill}{rgb}{0.749020,0.772549,0.792157}%
\pgfsetfillcolor{currentfill}%
\pgfsetfillopacity{0.160000}%
\pgfsetlinewidth{0.000000pt}%
\definecolor{currentstroke}{rgb}{0.000000,0.000000,0.000000}%
\pgfsetstrokecolor{currentstroke}%
\pgfsetdash{}{0pt}%
\pgfpathmoveto{\pgfqpoint{1.173088in}{0.710763in}}%
\pgfpathlineto{\pgfqpoint{1.178459in}{0.717260in}}%
\pgfpathlineto{\pgfqpoint{1.183854in}{0.723786in}}%
\pgfpathlineto{\pgfqpoint{1.198598in}{0.726922in}}%
\pgfpathlineto{\pgfqpoint{1.212825in}{0.730453in}}%
\pgfpathlineto{\pgfqpoint{1.205774in}{0.723556in}}%
\pgfpathlineto{\pgfqpoint{1.198752in}{0.716687in}}%
\pgfpathlineto{\pgfqpoint{1.186146in}{0.713549in}}%
\pgfpathlineto{\pgfqpoint{1.173088in}{0.710763in}}%
\pgfpathclose%
\pgfusepath{fill}%
\end{pgfscope}%
\begin{pgfscope}%
\pgfpathrectangle{\pgfqpoint{0.407592in}{0.222904in}}{\pgfqpoint{1.411163in}{1.411163in}}%
\pgfusepath{clip}%
\pgfsetbuttcap%
\pgfsetroundjoin%
\definecolor{currentfill}{rgb}{0.749020,0.772549,0.792157}%
\pgfsetfillcolor{currentfill}%
\pgfsetfillopacity{0.160000}%
\pgfsetlinewidth{0.000000pt}%
\definecolor{currentstroke}{rgb}{0.000000,0.000000,0.000000}%
\pgfsetstrokecolor{currentstroke}%
\pgfsetdash{}{0pt}%
\pgfpathmoveto{\pgfqpoint{0.491914in}{1.116936in}}%
\pgfpathlineto{\pgfqpoint{0.477888in}{1.128807in}}%
\pgfpathlineto{\pgfqpoint{0.463842in}{1.140695in}}%
\pgfpathlineto{\pgfqpoint{0.461132in}{1.125158in}}%
\pgfpathlineto{\pgfqpoint{0.461173in}{1.109555in}}%
\pgfpathlineto{\pgfqpoint{0.475420in}{1.098254in}}%
\pgfpathlineto{\pgfqpoint{0.489640in}{1.086973in}}%
\pgfpathlineto{\pgfqpoint{0.489463in}{1.101983in}}%
\pgfpathlineto{\pgfqpoint{0.491914in}{1.116936in}}%
\pgfpathclose%
\pgfusepath{fill}%
\end{pgfscope}%
\begin{pgfscope}%
\pgfpathrectangle{\pgfqpoint{0.407592in}{0.222904in}}{\pgfqpoint{1.411163in}{1.411163in}}%
\pgfusepath{clip}%
\pgfsetbuttcap%
\pgfsetroundjoin%
\definecolor{currentfill}{rgb}{0.749020,0.772549,0.792157}%
\pgfsetfillcolor{currentfill}%
\pgfsetfillopacity{0.160000}%
\pgfsetlinewidth{0.000000pt}%
\definecolor{currentstroke}{rgb}{0.000000,0.000000,0.000000}%
\pgfsetstrokecolor{currentstroke}%
\pgfsetdash{}{0pt}%
\pgfpathmoveto{\pgfqpoint{1.708976in}{1.112688in}}%
\pgfpathlineto{\pgfqpoint{1.723001in}{1.124070in}}%
\pgfpathlineto{\pgfqpoint{1.737051in}{1.135472in}}%
\pgfpathlineto{\pgfqpoint{1.740260in}{1.151651in}}%
\pgfpathlineto{\pgfqpoint{1.740584in}{1.167739in}}%
\pgfpathlineto{\pgfqpoint{1.726603in}{1.155760in}}%
\pgfpathlineto{\pgfqpoint{1.712642in}{1.143798in}}%
\pgfpathlineto{\pgfqpoint{1.712189in}{1.128283in}}%
\pgfpathlineto{\pgfqpoint{1.708976in}{1.112688in}}%
\pgfpathclose%
\pgfusepath{fill}%
\end{pgfscope}%
\begin{pgfscope}%
\pgfpathrectangle{\pgfqpoint{0.407592in}{0.222904in}}{\pgfqpoint{1.411163in}{1.411163in}}%
\pgfusepath{clip}%
\pgfsetbuttcap%
\pgfsetroundjoin%
\definecolor{currentfill}{rgb}{0.749020,0.772549,0.792157}%
\pgfsetfillcolor{currentfill}%
\pgfsetfillopacity{0.160000}%
\pgfsetlinewidth{0.000000pt}%
\definecolor{currentstroke}{rgb}{0.000000,0.000000,0.000000}%
\pgfsetstrokecolor{currentstroke}%
\pgfsetdash{}{0pt}%
\pgfpathmoveto{\pgfqpoint{0.914188in}{0.742278in}}%
\pgfpathlineto{\pgfqpoint{0.904057in}{0.749860in}}%
\pgfpathlineto{\pgfqpoint{0.893888in}{0.757470in}}%
\pgfpathlineto{\pgfqpoint{0.906652in}{0.752433in}}%
\pgfpathlineto{\pgfqpoint{0.920274in}{0.747773in}}%
\pgfpathlineto{\pgfqpoint{0.929100in}{0.740646in}}%
\pgfpathlineto{\pgfqpoint{0.937890in}{0.733547in}}%
\pgfpathlineto{\pgfqpoint{0.925657in}{0.737742in}}%
\pgfpathlineto{\pgfqpoint{0.914188in}{0.742278in}}%
\pgfpathclose%
\pgfusepath{fill}%
\end{pgfscope}%
\begin{pgfscope}%
\pgfpathrectangle{\pgfqpoint{0.407592in}{0.222904in}}{\pgfqpoint{1.411163in}{1.411163in}}%
\pgfusepath{clip}%
\pgfsetbuttcap%
\pgfsetroundjoin%
\definecolor{currentfill}{rgb}{0.749020,0.772549,0.792157}%
\pgfsetfillcolor{currentfill}%
\pgfsetfillopacity{0.160000}%
\pgfsetlinewidth{0.000000pt}%
\definecolor{currentstroke}{rgb}{0.000000,0.000000,0.000000}%
\pgfsetstrokecolor{currentstroke}%
\pgfsetdash{}{0pt}%
\pgfpathmoveto{\pgfqpoint{0.978884in}{0.712962in}}%
\pgfpathlineto{\pgfqpoint{0.971671in}{0.719597in}}%
\pgfpathlineto{\pgfqpoint{0.964427in}{0.726261in}}%
\pgfpathlineto{\pgfqpoint{0.978609in}{0.723206in}}%
\pgfpathlineto{\pgfqpoint{0.993314in}{0.720565in}}%
\pgfpathlineto{\pgfqpoint{0.998904in}{0.714219in}}%
\pgfpathlineto{\pgfqpoint{1.004468in}{0.707902in}}%
\pgfpathlineto{\pgfqpoint{0.991446in}{0.710248in}}%
\pgfpathlineto{\pgfqpoint{0.978884in}{0.712962in}}%
\pgfpathclose%
\pgfusepath{fill}%
\end{pgfscope}%
\begin{pgfscope}%
\pgfpathrectangle{\pgfqpoint{0.407592in}{0.222904in}}{\pgfqpoint{1.411163in}{1.411163in}}%
\pgfusepath{clip}%
\pgfsetbuttcap%
\pgfsetroundjoin%
\definecolor{currentfill}{rgb}{0.749020,0.772549,0.792157}%
\pgfsetfillcolor{currentfill}%
\pgfsetfillopacity{0.160000}%
\pgfsetlinewidth{0.000000pt}%
\definecolor{currentstroke}{rgb}{0.000000,0.000000,0.000000}%
\pgfsetstrokecolor{currentstroke}%
\pgfsetdash{}{0pt}%
\pgfpathmoveto{\pgfqpoint{1.557570in}{0.975003in}}%
\pgfpathlineto{\pgfqpoint{1.571092in}{0.985549in}}%
\pgfpathlineto{\pgfqpoint{1.584645in}{0.996118in}}%
\pgfpathlineto{\pgfqpoint{1.592308in}{1.009206in}}%
\pgfpathlineto{\pgfqpoint{1.597687in}{1.022369in}}%
\pgfpathlineto{\pgfqpoint{1.583889in}{1.011171in}}%
\pgfpathlineto{\pgfqpoint{1.570115in}{0.999993in}}%
\pgfpathlineto{\pgfqpoint{1.564920in}{0.987459in}}%
\pgfpathlineto{\pgfqpoint{1.557570in}{0.975003in}}%
\pgfpathclose%
\pgfusepath{fill}%
\end{pgfscope}%
\begin{pgfscope}%
\pgfpathrectangle{\pgfqpoint{0.407592in}{0.222904in}}{\pgfqpoint{1.411163in}{1.411163in}}%
\pgfusepath{clip}%
\pgfsetbuttcap%
\pgfsetroundjoin%
\definecolor{currentfill}{rgb}{0.749020,0.772549,0.792157}%
\pgfsetfillcolor{currentfill}%
\pgfsetfillopacity{0.160000}%
\pgfsetlinewidth{0.000000pt}%
\definecolor{currentstroke}{rgb}{0.000000,0.000000,0.000000}%
\pgfsetstrokecolor{currentstroke}%
\pgfsetdash{}{0pt}%
\pgfpathmoveto{\pgfqpoint{0.827881in}{0.798467in}}%
\pgfpathlineto{\pgfqpoint{0.815533in}{0.807313in}}%
\pgfpathlineto{\pgfqpoint{0.803146in}{0.816186in}}%
\pgfpathlineto{\pgfqpoint{0.813690in}{0.808683in}}%
\pgfpathlineto{\pgfqpoint{0.825518in}{0.801494in}}%
\pgfpathlineto{\pgfqpoint{0.836986in}{0.793221in}}%
\pgfpathlineto{\pgfqpoint{0.848416in}{0.784976in}}%
\pgfpathlineto{\pgfqpoint{0.837565in}{0.791576in}}%
\pgfpathlineto{\pgfqpoint{0.827881in}{0.798467in}}%
\pgfpathclose%
\pgfusepath{fill}%
\end{pgfscope}%
\begin{pgfscope}%
\pgfpathrectangle{\pgfqpoint{0.407592in}{0.222904in}}{\pgfqpoint{1.411163in}{1.411163in}}%
\pgfusepath{clip}%
\pgfsetbuttcap%
\pgfsetroundjoin%
\definecolor{currentfill}{rgb}{0.749020,0.772549,0.792157}%
\pgfsetfillcolor{currentfill}%
\pgfsetfillopacity{0.160000}%
\pgfsetlinewidth{0.000000pt}%
\definecolor{currentstroke}{rgb}{0.000000,0.000000,0.000000}%
\pgfsetstrokecolor{currentstroke}%
\pgfsetdash{}{0pt}%
\pgfpathmoveto{\pgfqpoint{1.145854in}{0.706299in}}%
\pgfpathlineto{\pgfqpoint{1.149465in}{0.712515in}}%
\pgfpathlineto{\pgfqpoint{1.153092in}{0.718760in}}%
\pgfpathlineto{\pgfqpoint{1.168661in}{0.721061in}}%
\pgfpathlineto{\pgfqpoint{1.183854in}{0.723786in}}%
\pgfpathlineto{\pgfqpoint{1.178459in}{0.717260in}}%
\pgfpathlineto{\pgfqpoint{1.173088in}{0.710763in}}%
\pgfpathlineto{\pgfqpoint{1.159636in}{0.708342in}}%
\pgfpathlineto{\pgfqpoint{1.145854in}{0.706299in}}%
\pgfpathclose%
\pgfusepath{fill}%
\end{pgfscope}%
\begin{pgfscope}%
\pgfpathrectangle{\pgfqpoint{0.407592in}{0.222904in}}{\pgfqpoint{1.411163in}{1.411163in}}%
\pgfusepath{clip}%
\pgfsetbuttcap%
\pgfsetroundjoin%
\definecolor{currentfill}{rgb}{0.749020,0.772549,0.792157}%
\pgfsetfillcolor{currentfill}%
\pgfsetfillopacity{0.160000}%
\pgfsetlinewidth{0.000000pt}%
\definecolor{currentstroke}{rgb}{0.000000,0.000000,0.000000}%
\pgfsetstrokecolor{currentstroke}%
\pgfsetdash{}{0pt}%
\pgfpathmoveto{\pgfqpoint{1.459107in}{0.891619in}}%
\pgfpathlineto{\pgfqpoint{1.471911in}{0.901376in}}%
\pgfpathlineto{\pgfqpoint{1.484750in}{0.911159in}}%
\pgfpathlineto{\pgfqpoint{1.495203in}{0.922017in}}%
\pgfpathlineto{\pgfqpoint{1.503785in}{0.933057in}}%
\pgfpathlineto{\pgfqpoint{1.490414in}{0.922629in}}%
\pgfpathlineto{\pgfqpoint{1.477073in}{0.912225in}}%
\pgfpathlineto{\pgfqpoint{1.468962in}{0.901833in}}%
\pgfpathlineto{\pgfqpoint{1.459107in}{0.891619in}}%
\pgfpathclose%
\pgfusepath{fill}%
\end{pgfscope}%
\begin{pgfscope}%
\pgfpathrectangle{\pgfqpoint{0.407592in}{0.222904in}}{\pgfqpoint{1.411163in}{1.411163in}}%
\pgfusepath{clip}%
\pgfsetbuttcap%
\pgfsetroundjoin%
\definecolor{currentfill}{rgb}{0.749020,0.772549,0.792157}%
\pgfsetfillcolor{currentfill}%
\pgfsetfillopacity{0.160000}%
\pgfsetlinewidth{0.000000pt}%
\definecolor{currentstroke}{rgb}{0.000000,0.000000,0.000000}%
\pgfsetstrokecolor{currentstroke}%
\pgfsetdash{}{0pt}%
\pgfpathmoveto{\pgfqpoint{0.871156in}{0.768571in}}%
\pgfpathlineto{\pgfqpoint{0.859805in}{0.776759in}}%
\pgfpathlineto{\pgfqpoint{0.848416in}{0.784976in}}%
\pgfpathlineto{\pgfqpoint{0.860389in}{0.778699in}}%
\pgfpathlineto{\pgfqpoint{0.873435in}{0.772777in}}%
\pgfpathlineto{\pgfqpoint{0.883681in}{0.765109in}}%
\pgfpathlineto{\pgfqpoint{0.893888in}{0.757470in}}%
\pgfpathlineto{\pgfqpoint{0.882039in}{0.762858in}}%
\pgfpathlineto{\pgfqpoint{0.871156in}{0.768571in}}%
\pgfpathclose%
\pgfusepath{fill}%
\end{pgfscope}%
\begin{pgfscope}%
\pgfpathrectangle{\pgfqpoint{0.407592in}{0.222904in}}{\pgfqpoint{1.411163in}{1.411163in}}%
\pgfusepath{clip}%
\pgfsetbuttcap%
\pgfsetroundjoin%
\definecolor{currentfill}{rgb}{0.749020,0.772549,0.792157}%
\pgfsetfillcolor{currentfill}%
\pgfsetfillopacity{0.160000}%
\pgfsetlinewidth{0.000000pt}%
\definecolor{currentstroke}{rgb}{0.000000,0.000000,0.000000}%
\pgfsetstrokecolor{currentstroke}%
\pgfsetdash{}{0pt}%
\pgfpathmoveto{\pgfqpoint{0.602459in}{0.997475in}}%
\pgfpathlineto{\pgfqpoint{0.588448in}{1.008590in}}%
\pgfpathlineto{\pgfqpoint{0.574411in}{1.019725in}}%
\pgfpathlineto{\pgfqpoint{0.577202in}{1.006572in}}%
\pgfpathlineto{\pgfqpoint{0.582297in}{0.993507in}}%
\pgfpathlineto{\pgfqpoint{0.596231in}{0.983001in}}%
\pgfpathlineto{\pgfqpoint{0.610133in}{0.972519in}}%
\pgfpathlineto{\pgfqpoint{0.605209in}{0.984952in}}%
\pgfpathlineto{\pgfqpoint{0.602459in}{0.997475in}}%
\pgfpathclose%
\pgfusepath{fill}%
\end{pgfscope}%
\begin{pgfscope}%
\pgfpathrectangle{\pgfqpoint{0.407592in}{0.222904in}}{\pgfqpoint{1.411163in}{1.411163in}}%
\pgfusepath{clip}%
\pgfsetbuttcap%
\pgfsetroundjoin%
\definecolor{currentfill}{rgb}{0.749020,0.772549,0.792157}%
\pgfsetfillcolor{currentfill}%
\pgfsetfillopacity{0.160000}%
\pgfsetlinewidth{0.000000pt}%
\definecolor{currentstroke}{rgb}{0.000000,0.000000,0.000000}%
\pgfsetstrokecolor{currentstroke}%
\pgfsetdash{}{0pt}%
\pgfpathmoveto{\pgfqpoint{1.212825in}{0.730453in}}%
\pgfpathlineto{\pgfqpoint{1.219906in}{0.737379in}}%
\pgfpathlineto{\pgfqpoint{1.227016in}{0.744334in}}%
\pgfpathlineto{\pgfqpoint{1.242222in}{0.748676in}}%
\pgfpathlineto{\pgfqpoint{1.256703in}{0.753414in}}%
\pgfpathlineto{\pgfqpoint{1.248071in}{0.746005in}}%
\pgfpathlineto{\pgfqpoint{1.239473in}{0.738626in}}%
\pgfpathlineto{\pgfqpoint{1.226471in}{0.734361in}}%
\pgfpathlineto{\pgfqpoint{1.212825in}{0.730453in}}%
\pgfpathclose%
\pgfusepath{fill}%
\end{pgfscope}%
\begin{pgfscope}%
\pgfpathrectangle{\pgfqpoint{0.407592in}{0.222904in}}{\pgfqpoint{1.411163in}{1.411163in}}%
\pgfusepath{clip}%
\pgfsetbuttcap%
\pgfsetroundjoin%
\definecolor{currentfill}{rgb}{0.749020,0.772549,0.792157}%
\pgfsetfillcolor{currentfill}%
\pgfsetfillopacity{0.160000}%
\pgfsetlinewidth{0.000000pt}%
\definecolor{currentstroke}{rgb}{0.000000,0.000000,0.000000}%
\pgfsetstrokecolor{currentstroke}%
\pgfsetdash{}{0pt}%
\pgfpathmoveto{\pgfqpoint{0.463842in}{1.140695in}}%
\pgfpathlineto{\pgfqpoint{0.449777in}{1.152600in}}%
\pgfpathlineto{\pgfqpoint{0.435692in}{1.164522in}}%
\pgfpathlineto{\pgfqpoint{0.432709in}{1.148410in}}%
\pgfpathlineto{\pgfqpoint{0.432600in}{1.132222in}}%
\pgfpathlineto{\pgfqpoint{0.446900in}{1.120878in}}%
\pgfpathlineto{\pgfqpoint{0.461173in}{1.109555in}}%
\pgfpathlineto{\pgfqpoint{0.461132in}{1.125158in}}%
\pgfpathlineto{\pgfqpoint{0.463842in}{1.140695in}}%
\pgfpathclose%
\pgfusepath{fill}%
\end{pgfscope}%
\begin{pgfscope}%
\pgfpathrectangle{\pgfqpoint{0.407592in}{0.222904in}}{\pgfqpoint{1.411163in}{1.411163in}}%
\pgfusepath{clip}%
\pgfsetbuttcap%
\pgfsetroundjoin%
\definecolor{currentfill}{rgb}{0.749020,0.772549,0.792157}%
\pgfsetfillcolor{currentfill}%
\pgfsetfillopacity{0.160000}%
\pgfsetlinewidth{0.000000pt}%
\definecolor{currentstroke}{rgb}{0.000000,0.000000,0.000000}%
\pgfsetstrokecolor{currentstroke}%
\pgfsetdash{}{0pt}%
\pgfpathmoveto{\pgfqpoint{1.386591in}{0.835679in}}%
\pgfpathlineto{\pgfqpoint{1.398461in}{0.844696in}}%
\pgfpathlineto{\pgfqpoint{1.410368in}{0.853742in}}%
\pgfpathlineto{\pgfqpoint{1.422760in}{0.862832in}}%
\pgfpathlineto{\pgfqpoint{1.433600in}{0.872182in}}%
\pgfpathlineto{\pgfqpoint{1.420897in}{0.862503in}}%
\pgfpathlineto{\pgfqpoint{1.408228in}{0.852849in}}%
\pgfpathlineto{\pgfqpoint{1.398125in}{0.844141in}}%
\pgfpathlineto{\pgfqpoint{1.386591in}{0.835679in}}%
\pgfpathclose%
\pgfusepath{fill}%
\end{pgfscope}%
\begin{pgfscope}%
\pgfpathrectangle{\pgfqpoint{0.407592in}{0.222904in}}{\pgfqpoint{1.411163in}{1.411163in}}%
\pgfusepath{clip}%
\pgfsetbuttcap%
\pgfsetroundjoin%
\definecolor{currentfill}{rgb}{0.749020,0.772549,0.792157}%
\pgfsetfillcolor{currentfill}%
\pgfsetfillopacity{0.160000}%
\pgfsetlinewidth{0.000000pt}%
\definecolor{currentstroke}{rgb}{0.000000,0.000000,0.000000}%
\pgfsetstrokecolor{currentstroke}%
\pgfsetdash{}{0pt}%
\pgfpathmoveto{\pgfqpoint{1.004468in}{0.707902in}}%
\pgfpathlineto{\pgfqpoint{0.998904in}{0.714219in}}%
\pgfpathlineto{\pgfqpoint{0.993314in}{0.720565in}}%
\pgfpathlineto{\pgfqpoint{1.008474in}{0.718351in}}%
\pgfpathlineto{\pgfqpoint{1.024018in}{0.716575in}}%
\pgfpathlineto{\pgfqpoint{1.027840in}{0.710453in}}%
\pgfpathlineto{\pgfqpoint{1.031645in}{0.704359in}}%
\pgfpathlineto{\pgfqpoint{1.017888in}{0.705935in}}%
\pgfpathlineto{\pgfqpoint{1.004468in}{0.707902in}}%
\pgfpathclose%
\pgfusepath{fill}%
\end{pgfscope}%
\begin{pgfscope}%
\pgfpathrectangle{\pgfqpoint{0.407592in}{0.222904in}}{\pgfqpoint{1.411163in}{1.411163in}}%
\pgfusepath{clip}%
\pgfsetbuttcap%
\pgfsetroundjoin%
\definecolor{currentfill}{rgb}{0.749020,0.772549,0.792157}%
\pgfsetfillcolor{currentfill}%
\pgfsetfillopacity{0.160000}%
\pgfsetlinewidth{0.000000pt}%
\definecolor{currentstroke}{rgb}{0.000000,0.000000,0.000000}%
\pgfsetstrokecolor{currentstroke}%
\pgfsetdash{}{0pt}%
\pgfpathmoveto{\pgfqpoint{0.692881in}{0.910128in}}%
\pgfpathlineto{\pgfqpoint{0.679168in}{0.920467in}}%
\pgfpathlineto{\pgfqpoint{0.665424in}{0.930830in}}%
\pgfpathlineto{\pgfqpoint{0.671892in}{0.919822in}}%
\pgfpathlineto{\pgfqpoint{0.680268in}{0.909008in}}%
\pgfpathlineto{\pgfqpoint{0.693618in}{0.899288in}}%
\pgfpathlineto{\pgfqpoint{0.706932in}{0.889595in}}%
\pgfpathlineto{\pgfqpoint{0.699016in}{0.899768in}}%
\pgfpathlineto{\pgfqpoint{0.692881in}{0.910128in}}%
\pgfpathclose%
\pgfusepath{fill}%
\end{pgfscope}%
\begin{pgfscope}%
\pgfpathrectangle{\pgfqpoint{0.407592in}{0.222904in}}{\pgfqpoint{1.411163in}{1.411163in}}%
\pgfusepath{clip}%
\pgfsetbuttcap%
\pgfsetroundjoin%
\definecolor{currentfill}{rgb}{0.749020,0.772549,0.792157}%
\pgfsetfillcolor{currentfill}%
\pgfsetfillopacity{0.160000}%
\pgfsetlinewidth{0.000000pt}%
\definecolor{currentstroke}{rgb}{0.000000,0.000000,0.000000}%
\pgfsetstrokecolor{currentstroke}%
\pgfsetdash{}{0pt}%
\pgfpathmoveto{\pgfqpoint{1.737051in}{1.135472in}}%
\pgfpathlineto{\pgfqpoint{1.751127in}{1.146896in}}%
\pgfpathlineto{\pgfqpoint{1.765228in}{1.158340in}}%
\pgfpathlineto{\pgfqpoint{1.768420in}{1.175092in}}%
\pgfpathlineto{\pgfqpoint{1.768602in}{1.191744in}}%
\pgfpathlineto{\pgfqpoint{1.754584in}{1.179733in}}%
\pgfpathlineto{\pgfqpoint{1.740584in}{1.167739in}}%
\pgfpathlineto{\pgfqpoint{1.740260in}{1.151651in}}%
\pgfpathlineto{\pgfqpoint{1.737051in}{1.135472in}}%
\pgfpathclose%
\pgfusepath{fill}%
\end{pgfscope}%
\begin{pgfscope}%
\pgfpathrectangle{\pgfqpoint{0.407592in}{0.222904in}}{\pgfqpoint{1.411163in}{1.411163in}}%
\pgfusepath{clip}%
\pgfsetbuttcap%
\pgfsetroundjoin%
\definecolor{currentfill}{rgb}{0.749020,0.772549,0.792157}%
\pgfsetfillcolor{currentfill}%
\pgfsetfillopacity{0.160000}%
\pgfsetlinewidth{0.000000pt}%
\definecolor{currentstroke}{rgb}{0.000000,0.000000,0.000000}%
\pgfsetstrokecolor{currentstroke}%
\pgfsetdash{}{0pt}%
\pgfpathmoveto{\pgfqpoint{1.117557in}{0.703384in}}%
\pgfpathlineto{\pgfqpoint{1.119332in}{0.709416in}}%
\pgfpathlineto{\pgfqpoint{1.121116in}{0.715477in}}%
\pgfpathlineto{\pgfqpoint{1.137218in}{0.716895in}}%
\pgfpathlineto{\pgfqpoint{1.153092in}{0.718760in}}%
\pgfpathlineto{\pgfqpoint{1.149465in}{0.712515in}}%
\pgfpathlineto{\pgfqpoint{1.145854in}{0.706299in}}%
\pgfpathlineto{\pgfqpoint{1.131806in}{0.704643in}}%
\pgfpathlineto{\pgfqpoint{1.117557in}{0.703384in}}%
\pgfpathclose%
\pgfusepath{fill}%
\end{pgfscope}%
\begin{pgfscope}%
\pgfpathrectangle{\pgfqpoint{0.407592in}{0.222904in}}{\pgfqpoint{1.411163in}{1.411163in}}%
\pgfusepath{clip}%
\pgfsetbuttcap%
\pgfsetroundjoin%
\definecolor{currentfill}{rgb}{0.749020,0.772549,0.792157}%
\pgfsetfillcolor{currentfill}%
\pgfsetfillopacity{0.160000}%
\pgfsetlinewidth{0.000000pt}%
\definecolor{currentstroke}{rgb}{0.000000,0.000000,0.000000}%
\pgfsetstrokecolor{currentstroke}%
\pgfsetdash{}{0pt}%
\pgfpathmoveto{\pgfqpoint{0.759828in}{0.851084in}}%
\pgfpathlineto{\pgfqpoint{0.746657in}{0.860673in}}%
\pgfpathlineto{\pgfqpoint{0.733451in}{0.870288in}}%
\pgfpathlineto{\pgfqpoint{0.742531in}{0.860986in}}%
\pgfpathlineto{\pgfqpoint{0.753215in}{0.851954in}}%
\pgfpathlineto{\pgfqpoint{0.765756in}{0.842970in}}%
\pgfpathlineto{\pgfqpoint{0.778258in}{0.834014in}}%
\pgfpathlineto{\pgfqpoint{0.768303in}{0.842422in}}%
\pgfpathlineto{\pgfqpoint{0.759828in}{0.851084in}}%
\pgfpathclose%
\pgfusepath{fill}%
\end{pgfscope}%
\begin{pgfscope}%
\pgfpathrectangle{\pgfqpoint{0.407592in}{0.222904in}}{\pgfqpoint{1.411163in}{1.411163in}}%
\pgfusepath{clip}%
\pgfsetbuttcap%
\pgfsetroundjoin%
\definecolor{currentfill}{rgb}{0.749020,0.772549,0.792157}%
\pgfsetfillcolor{currentfill}%
\pgfsetfillopacity{0.160000}%
\pgfsetlinewidth{0.000000pt}%
\definecolor{currentstroke}{rgb}{0.000000,0.000000,0.000000}%
\pgfsetstrokecolor{currentstroke}%
\pgfsetdash{}{0pt}%
\pgfpathmoveto{\pgfqpoint{1.303300in}{0.779931in}}%
\pgfpathlineto{\pgfqpoint{1.313386in}{0.787950in}}%
\pgfpathlineto{\pgfqpoint{1.323509in}{0.795998in}}%
\pgfpathlineto{\pgfqpoint{1.337847in}{0.802910in}}%
\pgfpathlineto{\pgfqpoint{1.351015in}{0.810164in}}%
\pgfpathlineto{\pgfqpoint{1.339744in}{0.801536in}}%
\pgfpathlineto{\pgfqpoint{1.328510in}{0.792936in}}%
\pgfpathlineto{\pgfqpoint{1.316436in}{0.786275in}}%
\pgfpathlineto{\pgfqpoint{1.303300in}{0.779931in}}%
\pgfpathclose%
\pgfusepath{fill}%
\end{pgfscope}%
\begin{pgfscope}%
\pgfpathrectangle{\pgfqpoint{0.407592in}{0.222904in}}{\pgfqpoint{1.411163in}{1.411163in}}%
\pgfusepath{clip}%
\pgfsetbuttcap%
\pgfsetroundjoin%
\definecolor{currentfill}{rgb}{0.749020,0.772549,0.792157}%
\pgfsetfillcolor{currentfill}%
\pgfsetfillopacity{0.160000}%
\pgfsetlinewidth{0.000000pt}%
\definecolor{currentstroke}{rgb}{0.000000,0.000000,0.000000}%
\pgfsetstrokecolor{currentstroke}%
\pgfsetdash{}{0pt}%
\pgfpathmoveto{\pgfqpoint{1.256703in}{0.753414in}}%
\pgfpathlineto{\pgfqpoint{1.265368in}{0.760851in}}%
\pgfpathlineto{\pgfqpoint{1.274068in}{0.768318in}}%
\pgfpathlineto{\pgfqpoint{1.289156in}{0.773935in}}%
\pgfpathlineto{\pgfqpoint{1.303300in}{0.779931in}}%
\pgfpathlineto{\pgfqpoint{1.293250in}{0.771941in}}%
\pgfpathlineto{\pgfqpoint{1.283236in}{0.763979in}}%
\pgfpathlineto{\pgfqpoint{1.270394in}{0.758523in}}%
\pgfpathlineto{\pgfqpoint{1.256703in}{0.753414in}}%
\pgfpathclose%
\pgfusepath{fill}%
\end{pgfscope}%
\begin{pgfscope}%
\pgfpathrectangle{\pgfqpoint{0.407592in}{0.222904in}}{\pgfqpoint{1.411163in}{1.411163in}}%
\pgfusepath{clip}%
\pgfsetbuttcap%
\pgfsetroundjoin%
\definecolor{currentfill}{rgb}{0.749020,0.772549,0.792157}%
\pgfsetfillcolor{currentfill}%
\pgfsetfillopacity{0.160000}%
\pgfsetlinewidth{0.000000pt}%
\definecolor{currentstroke}{rgb}{0.000000,0.000000,0.000000}%
\pgfsetstrokecolor{currentstroke}%
\pgfsetdash{}{0pt}%
\pgfpathmoveto{\pgfqpoint{1.351015in}{0.810164in}}%
\pgfpathlineto{\pgfqpoint{1.362323in}{0.818820in}}%
\pgfpathlineto{\pgfqpoint{1.373668in}{0.827504in}}%
\pgfpathlineto{\pgfqpoint{1.386591in}{0.835679in}}%
\pgfpathlineto{\pgfqpoint{1.374757in}{0.826688in}}%
\pgfpathlineto{\pgfqpoint{1.362959in}{0.817725in}}%
\pgfusepath{fill}%
\end{pgfscope}%
\begin{pgfscope}%
\pgfpathrectangle{\pgfqpoint{0.407592in}{0.222904in}}{\pgfqpoint{1.411163in}{1.411163in}}%
\pgfusepath{clip}%
\pgfsetbuttcap%
\pgfsetroundjoin%
\definecolor{currentfill}{rgb}{0.749020,0.772549,0.792157}%
\pgfsetfillcolor{currentfill}%
\pgfsetfillopacity{0.160000}%
\pgfsetlinewidth{0.000000pt}%
\definecolor{currentstroke}{rgb}{0.000000,0.000000,0.000000}%
\pgfsetstrokecolor{currentstroke}%
\pgfsetdash{}{0pt}%
\pgfpathmoveto{\pgfqpoint{1.031645in}{0.704359in}}%
\pgfpathlineto{\pgfqpoint{1.027840in}{0.710453in}}%
\pgfpathlineto{\pgfqpoint{1.024018in}{0.716575in}}%
\pgfpathlineto{\pgfqpoint{1.039871in}{0.715247in}}%
\pgfpathlineto{\pgfqpoint{1.055959in}{0.714373in}}%
\pgfpathlineto{\pgfqpoint{1.057939in}{0.708374in}}%
\pgfpathlineto{\pgfqpoint{1.059909in}{0.702404in}}%
\pgfpathlineto{\pgfqpoint{1.045674in}{0.703179in}}%
\pgfpathlineto{\pgfqpoint{1.031645in}{0.704359in}}%
\pgfpathclose%
\pgfusepath{fill}%
\end{pgfscope}%
\begin{pgfscope}%
\pgfpathrectangle{\pgfqpoint{0.407592in}{0.222904in}}{\pgfqpoint{1.411163in}{1.411163in}}%
\pgfusepath{clip}%
\pgfsetbuttcap%
\pgfsetroundjoin%
\definecolor{currentfill}{rgb}{0.749020,0.772549,0.792157}%
\pgfsetfillcolor{currentfill}%
\pgfsetfillopacity{0.160000}%
\pgfsetlinewidth{0.000000pt}%
\definecolor{currentstroke}{rgb}{0.000000,0.000000,0.000000}%
\pgfsetstrokecolor{currentstroke}%
\pgfsetdash{}{0pt}%
\pgfpathmoveto{\pgfqpoint{1.584645in}{0.996118in}}%
\pgfpathlineto{\pgfqpoint{1.598228in}{1.006711in}}%
\pgfpathlineto{\pgfqpoint{1.611841in}{1.017328in}}%
\pgfpathlineto{\pgfqpoint{1.619806in}{1.031041in}}%
\pgfpathlineto{\pgfqpoint{1.625359in}{1.044827in}}%
\pgfpathlineto{\pgfqpoint{1.611511in}{1.033588in}}%
\pgfpathlineto{\pgfqpoint{1.597687in}{1.022369in}}%
\pgfpathlineto{\pgfqpoint{1.592308in}{1.009206in}}%
\pgfpathlineto{\pgfqpoint{1.584645in}{0.996118in}}%
\pgfpathclose%
\pgfusepath{fill}%
\end{pgfscope}%
\begin{pgfscope}%
\pgfpathrectangle{\pgfqpoint{0.407592in}{0.222904in}}{\pgfqpoint{1.411163in}{1.411163in}}%
\pgfusepath{clip}%
\pgfsetbuttcap%
\pgfsetroundjoin%
\definecolor{currentfill}{rgb}{0.749020,0.772549,0.792157}%
\pgfsetfillcolor{currentfill}%
\pgfsetfillopacity{0.160000}%
\pgfsetlinewidth{0.000000pt}%
\definecolor{currentstroke}{rgb}{0.000000,0.000000,0.000000}%
\pgfsetstrokecolor{currentstroke}%
\pgfsetdash{}{0pt}%
\pgfpathmoveto{\pgfqpoint{1.088728in}{0.702077in}}%
\pgfpathlineto{\pgfqpoint{1.088630in}{0.708026in}}%
\pgfpathlineto{\pgfqpoint{1.088532in}{0.714004in}}%
\pgfpathlineto{\pgfqpoint{1.104862in}{0.714512in}}%
\pgfpathlineto{\pgfqpoint{1.121116in}{0.715477in}}%
\pgfpathlineto{\pgfqpoint{1.119332in}{0.709416in}}%
\pgfpathlineto{\pgfqpoint{1.117557in}{0.703384in}}%
\pgfpathlineto{\pgfqpoint{1.103175in}{0.702527in}}%
\pgfpathlineto{\pgfqpoint{1.088728in}{0.702077in}}%
\pgfpathclose%
\pgfusepath{fill}%
\end{pgfscope}%
\begin{pgfscope}%
\pgfpathrectangle{\pgfqpoint{0.407592in}{0.222904in}}{\pgfqpoint{1.411163in}{1.411163in}}%
\pgfusepath{clip}%
\pgfsetbuttcap%
\pgfsetroundjoin%
\definecolor{currentfill}{rgb}{0.749020,0.772549,0.792157}%
\pgfsetfillcolor{currentfill}%
\pgfsetfillopacity{0.160000}%
\pgfsetlinewidth{0.000000pt}%
\definecolor{currentstroke}{rgb}{0.000000,0.000000,0.000000}%
\pgfsetstrokecolor{currentstroke}%
\pgfsetdash{}{0pt}%
\pgfpathmoveto{\pgfqpoint{0.937890in}{0.733547in}}%
\pgfpathlineto{\pgfqpoint{0.929100in}{0.740646in}}%
\pgfpathlineto{\pgfqpoint{0.920274in}{0.747773in}}%
\pgfpathlineto{\pgfqpoint{0.934693in}{0.743513in}}%
\pgfpathlineto{\pgfqpoint{0.949845in}{0.739675in}}%
\pgfpathlineto{\pgfqpoint{0.957151in}{0.732953in}}%
\pgfpathlineto{\pgfqpoint{0.964427in}{0.726261in}}%
\pgfpathlineto{\pgfqpoint{0.950833in}{0.729714in}}%
\pgfpathlineto{\pgfqpoint{0.937890in}{0.733547in}}%
\pgfpathclose%
\pgfusepath{fill}%
\end{pgfscope}%
\begin{pgfscope}%
\pgfpathrectangle{\pgfqpoint{0.407592in}{0.222904in}}{\pgfqpoint{1.411163in}{1.411163in}}%
\pgfusepath{clip}%
\pgfsetbuttcap%
\pgfsetroundjoin%
\definecolor{currentfill}{rgb}{0.749020,0.772549,0.792157}%
\pgfsetfillcolor{currentfill}%
\pgfsetfillopacity{0.160000}%
\pgfsetlinewidth{0.000000pt}%
\definecolor{currentstroke}{rgb}{0.000000,0.000000,0.000000}%
\pgfsetstrokecolor{currentstroke}%
\pgfsetdash{}{0pt}%
\pgfpathmoveto{\pgfqpoint{1.059909in}{0.702404in}}%
\pgfpathlineto{\pgfqpoint{1.057939in}{0.708374in}}%
\pgfpathlineto{\pgfqpoint{1.055959in}{0.714373in}}%
\pgfpathlineto{\pgfqpoint{1.072205in}{0.713958in}}%
\pgfpathlineto{\pgfqpoint{1.088532in}{0.714004in}}%
\pgfpathlineto{\pgfqpoint{1.088630in}{0.708026in}}%
\pgfpathlineto{\pgfqpoint{1.088728in}{0.702077in}}%
\pgfpathlineto{\pgfqpoint{1.074283in}{0.702035in}}%
\pgfpathlineto{\pgfqpoint{1.059909in}{0.702404in}}%
\pgfpathclose%
\pgfusepath{fill}%
\end{pgfscope}%
\begin{pgfscope}%
\pgfpathrectangle{\pgfqpoint{0.407592in}{0.222904in}}{\pgfqpoint{1.411163in}{1.411163in}}%
\pgfusepath{clip}%
\pgfsetbuttcap%
\pgfsetroundjoin%
\definecolor{currentfill}{rgb}{0.749020,0.772549,0.792157}%
\pgfsetfillcolor{currentfill}%
\pgfsetfillopacity{0.160000}%
\pgfsetlinewidth{0.000000pt}%
\definecolor{currentstroke}{rgb}{0.000000,0.000000,0.000000}%
\pgfsetstrokecolor{currentstroke}%
\pgfsetdash{}{0pt}%
\pgfpathmoveto{\pgfqpoint{0.435692in}{1.164522in}}%
\pgfpathlineto{\pgfqpoint{0.421586in}{1.176460in}}%
\pgfpathlineto{\pgfqpoint{0.407462in}{1.188415in}}%
\pgfpathlineto{\pgfqpoint{0.404193in}{1.171737in}}%
\pgfpathlineto{\pgfqpoint{0.403920in}{1.154973in}}%
\pgfpathlineto{\pgfqpoint{0.418274in}{1.143587in}}%
\pgfpathlineto{\pgfqpoint{0.432600in}{1.132222in}}%
\pgfpathlineto{\pgfqpoint{0.432709in}{1.148410in}}%
\pgfpathlineto{\pgfqpoint{0.435692in}{1.164522in}}%
\pgfpathclose%
\pgfusepath{fill}%
\end{pgfscope}%
\begin{pgfscope}%
\pgfpathrectangle{\pgfqpoint{0.407592in}{0.222904in}}{\pgfqpoint{1.411163in}{1.411163in}}%
\pgfusepath{clip}%
\pgfsetbuttcap%
\pgfsetroundjoin%
\definecolor{currentfill}{rgb}{0.749020,0.772549,0.792157}%
\pgfsetfillcolor{currentfill}%
\pgfsetfillopacity{0.160000}%
\pgfsetlinewidth{0.000000pt}%
\definecolor{currentstroke}{rgb}{0.000000,0.000000,0.000000}%
\pgfsetstrokecolor{currentstroke}%
\pgfsetdash{}{0pt}%
\pgfpathmoveto{\pgfqpoint{0.574411in}{1.019725in}}%
\pgfpathlineto{\pgfqpoint{0.560348in}{1.030882in}}%
\pgfpathlineto{\pgfqpoint{0.546258in}{1.042058in}}%
\pgfpathlineto{\pgfqpoint{0.549080in}{1.028282in}}%
\pgfpathlineto{\pgfqpoint{0.554333in}{1.014591in}}%
\pgfpathlineto{\pgfqpoint{0.568331in}{1.004037in}}%
\pgfpathlineto{\pgfqpoint{0.582297in}{0.993507in}}%
\pgfpathlineto{\pgfqpoint{0.577202in}{1.006572in}}%
\pgfpathlineto{\pgfqpoint{0.574411in}{1.019725in}}%
\pgfpathclose%
\pgfusepath{fill}%
\end{pgfscope}%
\begin{pgfscope}%
\pgfpathrectangle{\pgfqpoint{0.407592in}{0.222904in}}{\pgfqpoint{1.411163in}{1.411163in}}%
\pgfusepath{clip}%
\pgfsetbuttcap%
\pgfsetroundjoin%
\definecolor{currentfill}{rgb}{0.749020,0.772549,0.792157}%
\pgfsetfillcolor{currentfill}%
\pgfsetfillopacity{0.160000}%
\pgfsetlinewidth{0.000000pt}%
\definecolor{currentstroke}{rgb}{0.000000,0.000000,0.000000}%
\pgfsetstrokecolor{currentstroke}%
\pgfsetdash{}{0pt}%
\pgfpathmoveto{\pgfqpoint{1.765228in}{1.158340in}}%
\pgfpathlineto{\pgfqpoint{1.779355in}{1.169805in}}%
\pgfpathlineto{\pgfqpoint{1.793508in}{1.181291in}}%
\pgfpathlineto{\pgfqpoint{1.796670in}{1.198609in}}%
\pgfpathlineto{\pgfqpoint{1.796697in}{1.215815in}}%
\pgfpathlineto{\pgfqpoint{1.782640in}{1.203771in}}%
\pgfpathlineto{\pgfqpoint{1.768602in}{1.191744in}}%
\pgfpathlineto{\pgfqpoint{1.768420in}{1.175092in}}%
\pgfpathlineto{\pgfqpoint{1.765228in}{1.158340in}}%
\pgfpathclose%
\pgfusepath{fill}%
\end{pgfscope}%
\begin{pgfscope}%
\pgfpathrectangle{\pgfqpoint{0.407592in}{0.222904in}}{\pgfqpoint{1.411163in}{1.411163in}}%
\pgfusepath{clip}%
\pgfsetbuttcap%
\pgfsetroundjoin%
\definecolor{currentfill}{rgb}{0.749020,0.772549,0.792157}%
\pgfsetfillcolor{currentfill}%
\pgfsetfillopacity{0.160000}%
\pgfsetlinewidth{0.000000pt}%
\definecolor{currentstroke}{rgb}{0.000000,0.000000,0.000000}%
\pgfsetstrokecolor{currentstroke}%
\pgfsetdash{}{0pt}%
\pgfpathmoveto{\pgfqpoint{1.484750in}{0.911159in}}%
\pgfpathlineto{\pgfqpoint{1.497623in}{0.920969in}}%
\pgfpathlineto{\pgfqpoint{1.510531in}{0.930804in}}%
\pgfpathlineto{\pgfqpoint{1.521573in}{0.942299in}}%
\pgfpathlineto{\pgfqpoint{1.530618in}{0.953983in}}%
\pgfpathlineto{\pgfqpoint{1.517186in}{0.943508in}}%
\pgfpathlineto{\pgfqpoint{1.503785in}{0.933057in}}%
\pgfpathlineto{\pgfqpoint{1.495203in}{0.922017in}}%
\pgfpathlineto{\pgfqpoint{1.484750in}{0.911159in}}%
\pgfpathclose%
\pgfusepath{fill}%
\end{pgfscope}%
\begin{pgfscope}%
\pgfpathrectangle{\pgfqpoint{0.407592in}{0.222904in}}{\pgfqpoint{1.411163in}{1.411163in}}%
\pgfusepath{clip}%
\pgfsetbuttcap%
\pgfsetroundjoin%
\definecolor{currentfill}{rgb}{0.749020,0.772549,0.792157}%
\pgfsetfillcolor{currentfill}%
\pgfsetfillopacity{0.160000}%
\pgfsetlinewidth{0.000000pt}%
\definecolor{currentstroke}{rgb}{0.000000,0.000000,0.000000}%
\pgfsetstrokecolor{currentstroke}%
\pgfsetdash{}{0pt}%
\pgfpathmoveto{\pgfqpoint{1.183854in}{0.723786in}}%
\pgfpathlineto{\pgfqpoint{1.189272in}{0.730341in}}%
\pgfpathlineto{\pgfqpoint{1.194714in}{0.736924in}}%
\pgfpathlineto{\pgfqpoint{1.211156in}{0.740410in}}%
\pgfpathlineto{\pgfqpoint{1.227016in}{0.744334in}}%
\pgfpathlineto{\pgfqpoint{1.219906in}{0.737379in}}%
\pgfpathlineto{\pgfqpoint{1.212825in}{0.730453in}}%
\pgfpathlineto{\pgfqpoint{1.198598in}{0.726922in}}%
\pgfpathlineto{\pgfqpoint{1.183854in}{0.723786in}}%
\pgfpathclose%
\pgfusepath{fill}%
\end{pgfscope}%
\begin{pgfscope}%
\pgfpathrectangle{\pgfqpoint{0.407592in}{0.222904in}}{\pgfqpoint{1.411163in}{1.411163in}}%
\pgfusepath{clip}%
\pgfsetbuttcap%
\pgfsetroundjoin%
\definecolor{currentfill}{rgb}{0.749020,0.772549,0.792157}%
\pgfsetfillcolor{currentfill}%
\pgfsetfillopacity{0.160000}%
\pgfsetlinewidth{0.000000pt}%
\definecolor{currentstroke}{rgb}{0.000000,0.000000,0.000000}%
\pgfsetstrokecolor{currentstroke}%
\pgfsetdash{}{0pt}%
\pgfpathmoveto{\pgfqpoint{0.803146in}{0.816186in}}%
\pgfpathlineto{\pgfqpoint{0.790721in}{0.825086in}}%
\pgfpathlineto{\pgfqpoint{0.778258in}{0.834014in}}%
\pgfpathlineto{\pgfqpoint{0.789658in}{0.825903in}}%
\pgfpathlineto{\pgfqpoint{0.802460in}{0.818128in}}%
\pgfpathlineto{\pgfqpoint{0.814009in}{0.809797in}}%
\pgfpathlineto{\pgfqpoint{0.825518in}{0.801494in}}%
\pgfpathlineto{\pgfqpoint{0.813690in}{0.808683in}}%
\pgfpathlineto{\pgfqpoint{0.803146in}{0.816186in}}%
\pgfpathclose%
\pgfusepath{fill}%
\end{pgfscope}%
\begin{pgfscope}%
\pgfpathrectangle{\pgfqpoint{0.407592in}{0.222904in}}{\pgfqpoint{1.411163in}{1.411163in}}%
\pgfusepath{clip}%
\pgfsetbuttcap%
\pgfsetroundjoin%
\definecolor{currentfill}{rgb}{0.749020,0.772549,0.792157}%
\pgfsetfillcolor{currentfill}%
\pgfsetfillopacity{0.160000}%
\pgfsetlinewidth{0.000000pt}%
\definecolor{currentstroke}{rgb}{0.000000,0.000000,0.000000}%
\pgfsetstrokecolor{currentstroke}%
\pgfsetdash{}{0pt}%
\pgfpathmoveto{\pgfqpoint{0.665424in}{0.930830in}}%
\pgfpathlineto{\pgfqpoint{0.651648in}{0.941217in}}%
\pgfpathlineto{\pgfqpoint{0.637842in}{0.951627in}}%
\pgfpathlineto{\pgfqpoint{0.644631in}{0.939976in}}%
\pgfpathlineto{\pgfqpoint{0.653460in}{0.928526in}}%
\pgfpathlineto{\pgfqpoint{0.666882in}{0.918754in}}%
\pgfpathlineto{\pgfqpoint{0.680268in}{0.909008in}}%
\pgfpathlineto{\pgfqpoint{0.671892in}{0.919822in}}%
\pgfpathlineto{\pgfqpoint{0.665424in}{0.930830in}}%
\pgfpathclose%
\pgfusepath{fill}%
\end{pgfscope}%
\begin{pgfscope}%
\pgfpathrectangle{\pgfqpoint{0.407592in}{0.222904in}}{\pgfqpoint{1.411163in}{1.411163in}}%
\pgfusepath{clip}%
\pgfsetbuttcap%
\pgfsetroundjoin%
\definecolor{currentfill}{rgb}{0.749020,0.772549,0.792157}%
\pgfsetfillcolor{currentfill}%
\pgfsetfillopacity{0.160000}%
\pgfsetlinewidth{0.000000pt}%
\definecolor{currentstroke}{rgb}{0.000000,0.000000,0.000000}%
\pgfsetstrokecolor{currentstroke}%
\pgfsetdash{}{0pt}%
\pgfpathmoveto{\pgfqpoint{0.893888in}{0.757470in}}%
\pgfpathlineto{\pgfqpoint{0.883681in}{0.765109in}}%
\pgfpathlineto{\pgfqpoint{0.873435in}{0.772777in}}%
\pgfpathlineto{\pgfqpoint{0.887497in}{0.767239in}}%
\pgfpathlineto{\pgfqpoint{0.902514in}{0.762114in}}%
\pgfpathlineto{\pgfqpoint{0.911412in}{0.754929in}}%
\pgfpathlineto{\pgfqpoint{0.920274in}{0.747773in}}%
\pgfpathlineto{\pgfqpoint{0.906652in}{0.752433in}}%
\pgfpathlineto{\pgfqpoint{0.893888in}{0.757470in}}%
\pgfpathclose%
\pgfusepath{fill}%
\end{pgfscope}%
\begin{pgfscope}%
\pgfpathrectangle{\pgfqpoint{0.407592in}{0.222904in}}{\pgfqpoint{1.411163in}{1.411163in}}%
\pgfusepath{clip}%
\pgfsetbuttcap%
\pgfsetroundjoin%
\definecolor{currentfill}{rgb}{0.749020,0.772549,0.792157}%
\pgfsetfillcolor{currentfill}%
\pgfsetfillopacity{0.160000}%
\pgfsetlinewidth{0.000000pt}%
\definecolor{currentstroke}{rgb}{0.000000,0.000000,0.000000}%
\pgfsetstrokecolor{currentstroke}%
\pgfsetdash{}{0pt}%
\pgfpathmoveto{\pgfqpoint{1.410368in}{0.853742in}}%
\pgfpathlineto{\pgfqpoint{1.422311in}{0.862815in}}%
\pgfpathlineto{\pgfqpoint{1.434292in}{0.871917in}}%
\pgfpathlineto{\pgfqpoint{1.447537in}{0.881631in}}%
\pgfpathlineto{\pgfqpoint{1.459107in}{0.891619in}}%
\pgfpathlineto{\pgfqpoint{1.446336in}{0.881888in}}%
\pgfpathlineto{\pgfqpoint{1.433600in}{0.872182in}}%
\pgfpathlineto{\pgfqpoint{1.422760in}{0.862832in}}%
\pgfpathlineto{\pgfqpoint{1.410368in}{0.853742in}}%
\pgfpathclose%
\pgfusepath{fill}%
\end{pgfscope}%
\begin{pgfscope}%
\pgfpathrectangle{\pgfqpoint{0.407592in}{0.222904in}}{\pgfqpoint{1.411163in}{1.411163in}}%
\pgfusepath{clip}%
\pgfsetbuttcap%
\pgfsetroundjoin%
\definecolor{currentfill}{rgb}{0.749020,0.772549,0.792157}%
\pgfsetfillcolor{currentfill}%
\pgfsetfillopacity{0.160000}%
\pgfsetlinewidth{0.000000pt}%
\definecolor{currentstroke}{rgb}{0.000000,0.000000,0.000000}%
\pgfsetstrokecolor{currentstroke}%
\pgfsetdash{}{0pt}%
\pgfpathmoveto{\pgfqpoint{1.611841in}{1.017328in}}%
\pgfpathlineto{\pgfqpoint{1.625485in}{1.027969in}}%
\pgfpathlineto{\pgfqpoint{1.639161in}{1.038634in}}%
\pgfpathlineto{\pgfqpoint{1.647416in}{1.052965in}}%
\pgfpathlineto{\pgfqpoint{1.653131in}{1.067365in}}%
\pgfpathlineto{\pgfqpoint{1.639232in}{1.056086in}}%
\pgfpathlineto{\pgfqpoint{1.625359in}{1.044827in}}%
\pgfpathlineto{\pgfqpoint{1.619806in}{1.031041in}}%
\pgfpathlineto{\pgfqpoint{1.611841in}{1.017328in}}%
\pgfpathclose%
\pgfusepath{fill}%
\end{pgfscope}%
\begin{pgfscope}%
\pgfpathrectangle{\pgfqpoint{0.407592in}{0.222904in}}{\pgfqpoint{1.411163in}{1.411163in}}%
\pgfusepath{clip}%
\pgfsetbuttcap%
\pgfsetroundjoin%
\definecolor{currentfill}{rgb}{0.749020,0.772549,0.792157}%
\pgfsetfillcolor{currentfill}%
\pgfsetfillopacity{0.160000}%
\pgfsetlinewidth{0.000000pt}%
\definecolor{currentstroke}{rgb}{0.000000,0.000000,0.000000}%
\pgfsetstrokecolor{currentstroke}%
\pgfsetdash{}{0pt}%
\pgfpathmoveto{\pgfqpoint{0.407462in}{1.188415in}}%
\pgfpathlineto{\pgfqpoint{0.393317in}{1.200387in}}%
\pgfpathlineto{\pgfqpoint{0.379152in}{1.212375in}}%
\pgfpathlineto{\pgfqpoint{0.375583in}{1.195141in}}%
\pgfpathlineto{\pgfqpoint{0.375133in}{1.177810in}}%
\pgfpathlineto{\pgfqpoint{0.389540in}{1.166381in}}%
\pgfpathlineto{\pgfqpoint{0.403920in}{1.154973in}}%
\pgfpathlineto{\pgfqpoint{0.404193in}{1.171737in}}%
\pgfpathlineto{\pgfqpoint{0.407462in}{1.188415in}}%
\pgfpathclose%
\pgfusepath{fill}%
\end{pgfscope}%
\begin{pgfscope}%
\pgfpathrectangle{\pgfqpoint{0.407592in}{0.222904in}}{\pgfqpoint{1.411163in}{1.411163in}}%
\pgfusepath{clip}%
\pgfsetbuttcap%
\pgfsetroundjoin%
\definecolor{currentfill}{rgb}{0.749020,0.772549,0.792157}%
\pgfsetfillcolor{currentfill}%
\pgfsetfillopacity{0.160000}%
\pgfsetlinewidth{0.000000pt}%
\definecolor{currentstroke}{rgb}{0.000000,0.000000,0.000000}%
\pgfsetstrokecolor{currentstroke}%
\pgfsetdash{}{0pt}%
\pgfpathmoveto{\pgfqpoint{0.848416in}{0.784976in}}%
\pgfpathlineto{\pgfqpoint{0.836986in}{0.793221in}}%
\pgfpathlineto{\pgfqpoint{0.825518in}{0.801494in}}%
\pgfpathlineto{\pgfqpoint{0.838581in}{0.794655in}}%
\pgfpathlineto{\pgfqpoint{0.852825in}{0.788200in}}%
\pgfpathlineto{\pgfqpoint{0.863150in}{0.780474in}}%
\pgfpathlineto{\pgfqpoint{0.873435in}{0.772777in}}%
\pgfpathlineto{\pgfqpoint{0.860389in}{0.778699in}}%
\pgfpathlineto{\pgfqpoint{0.848416in}{0.784976in}}%
\pgfpathclose%
\pgfusepath{fill}%
\end{pgfscope}%
\begin{pgfscope}%
\pgfpathrectangle{\pgfqpoint{0.407592in}{0.222904in}}{\pgfqpoint{1.411163in}{1.411163in}}%
\pgfusepath{clip}%
\pgfsetbuttcap%
\pgfsetroundjoin%
\definecolor{currentfill}{rgb}{0.749020,0.772549,0.792157}%
\pgfsetfillcolor{currentfill}%
\pgfsetfillopacity{0.160000}%
\pgfsetlinewidth{0.000000pt}%
\definecolor{currentstroke}{rgb}{0.000000,0.000000,0.000000}%
\pgfsetstrokecolor{currentstroke}%
\pgfsetdash{}{0pt}%
\pgfpathmoveto{\pgfqpoint{0.964427in}{0.726261in}}%
\pgfpathlineto{\pgfqpoint{0.957151in}{0.732953in}}%
\pgfpathlineto{\pgfqpoint{0.949845in}{0.739675in}}%
\pgfpathlineto{\pgfqpoint{0.965658in}{0.736279in}}%
\pgfpathlineto{\pgfqpoint{0.982060in}{0.733342in}}%
\pgfpathlineto{\pgfqpoint{0.987700in}{0.726939in}}%
\pgfpathlineto{\pgfqpoint{0.993314in}{0.720565in}}%
\pgfpathlineto{\pgfqpoint{0.978609in}{0.723206in}}%
\pgfpathlineto{\pgfqpoint{0.964427in}{0.726261in}}%
\pgfpathclose%
\pgfusepath{fill}%
\end{pgfscope}%
\begin{pgfscope}%
\pgfpathrectangle{\pgfqpoint{0.407592in}{0.222904in}}{\pgfqpoint{1.411163in}{1.411163in}}%
\pgfusepath{clip}%
\pgfsetbuttcap%
\pgfsetroundjoin%
\definecolor{currentfill}{rgb}{0.749020,0.772549,0.792157}%
\pgfsetfillcolor{currentfill}%
\pgfsetfillopacity{0.160000}%
\pgfsetlinewidth{0.000000pt}%
\definecolor{currentstroke}{rgb}{0.000000,0.000000,0.000000}%
\pgfsetstrokecolor{currentstroke}%
\pgfsetdash{}{0pt}%
\pgfpathmoveto{\pgfqpoint{0.733451in}{0.870288in}}%
\pgfpathlineto{\pgfqpoint{0.720209in}{0.879928in}}%
\pgfpathlineto{\pgfqpoint{0.706932in}{0.889595in}}%
\pgfpathlineto{\pgfqpoint{0.716608in}{0.879658in}}%
\pgfpathlineto{\pgfqpoint{0.728016in}{0.870005in}}%
\pgfpathlineto{\pgfqpoint{0.740635in}{0.860966in}}%
\pgfpathlineto{\pgfqpoint{0.753215in}{0.851954in}}%
\pgfpathlineto{\pgfqpoint{0.742531in}{0.860986in}}%
\pgfpathlineto{\pgfqpoint{0.733451in}{0.870288in}}%
\pgfpathclose%
\pgfusepath{fill}%
\end{pgfscope}%
\begin{pgfscope}%
\pgfpathrectangle{\pgfqpoint{0.407592in}{0.222904in}}{\pgfqpoint{1.411163in}{1.411163in}}%
\pgfusepath{clip}%
\pgfsetbuttcap%
\pgfsetroundjoin%
\definecolor{currentfill}{rgb}{0.749020,0.772549,0.792157}%
\pgfsetfillcolor{currentfill}%
\pgfsetfillopacity{0.160000}%
\pgfsetlinewidth{0.000000pt}%
\definecolor{currentstroke}{rgb}{0.000000,0.000000,0.000000}%
\pgfsetstrokecolor{currentstroke}%
\pgfsetdash{}{0pt}%
\pgfpathmoveto{\pgfqpoint{1.793508in}{1.181291in}}%
\pgfpathlineto{\pgfqpoint{1.807686in}{1.192798in}}%
\pgfpathlineto{\pgfqpoint{1.821891in}{1.204326in}}%
\pgfpathlineto{\pgfqpoint{1.825009in}{1.222200in}}%
\pgfpathlineto{\pgfqpoint{1.824868in}{1.239951in}}%
\pgfpathlineto{\pgfqpoint{1.810773in}{1.227874in}}%
\pgfpathlineto{\pgfqpoint{1.796697in}{1.215815in}}%
\pgfpathlineto{\pgfqpoint{1.796670in}{1.198609in}}%
\pgfpathlineto{\pgfqpoint{1.793508in}{1.181291in}}%
\pgfpathclose%
\pgfusepath{fill}%
\end{pgfscope}%
\begin{pgfscope}%
\pgfpathrectangle{\pgfqpoint{0.407592in}{0.222904in}}{\pgfqpoint{1.411163in}{1.411163in}}%
\pgfusepath{clip}%
\pgfsetbuttcap%
\pgfsetroundjoin%
\definecolor{currentfill}{rgb}{0.749020,0.772549,0.792157}%
\pgfsetfillcolor{currentfill}%
\pgfsetfillopacity{0.160000}%
\pgfsetlinewidth{0.000000pt}%
\definecolor{currentstroke}{rgb}{0.000000,0.000000,0.000000}%
\pgfsetstrokecolor{currentstroke}%
\pgfsetdash{}{0pt}%
\pgfpathmoveto{\pgfqpoint{1.227016in}{0.744334in}}%
\pgfpathlineto{\pgfqpoint{1.234156in}{0.751318in}}%
\pgfpathlineto{\pgfqpoint{1.241326in}{0.758332in}}%
\pgfpathlineto{\pgfqpoint{1.258100in}{0.763108in}}%
\pgfpathlineto{\pgfqpoint{1.274068in}{0.768318in}}%
\pgfpathlineto{\pgfqpoint{1.265368in}{0.760851in}}%
\pgfpathlineto{\pgfqpoint{1.256703in}{0.753414in}}%
\pgfpathlineto{\pgfqpoint{1.242222in}{0.748676in}}%
\pgfpathlineto{\pgfqpoint{1.227016in}{0.744334in}}%
\pgfpathclose%
\pgfusepath{fill}%
\end{pgfscope}%
\begin{pgfscope}%
\pgfpathrectangle{\pgfqpoint{0.407592in}{0.222904in}}{\pgfqpoint{1.411163in}{1.411163in}}%
\pgfusepath{clip}%
\pgfsetbuttcap%
\pgfsetroundjoin%
\definecolor{currentfill}{rgb}{0.749020,0.772549,0.792157}%
\pgfsetfillcolor{currentfill}%
\pgfsetfillopacity{0.160000}%
\pgfsetlinewidth{0.000000pt}%
\definecolor{currentstroke}{rgb}{0.000000,0.000000,0.000000}%
\pgfsetstrokecolor{currentstroke}%
\pgfsetdash{}{0pt}%
\pgfpathmoveto{\pgfqpoint{1.153092in}{0.718760in}}%
\pgfpathlineto{\pgfqpoint{1.156735in}{0.725033in}}%
\pgfpathlineto{\pgfqpoint{1.160396in}{0.731335in}}%
\pgfpathlineto{\pgfqpoint{1.177768in}{0.733894in}}%
\pgfpathlineto{\pgfqpoint{1.194714in}{0.736924in}}%
\pgfpathlineto{\pgfqpoint{1.189272in}{0.730341in}}%
\pgfpathlineto{\pgfqpoint{1.183854in}{0.723786in}}%
\pgfpathlineto{\pgfqpoint{1.168661in}{0.721061in}}%
\pgfpathlineto{\pgfqpoint{1.153092in}{0.718760in}}%
\pgfpathclose%
\pgfusepath{fill}%
\end{pgfscope}%
\begin{pgfscope}%
\pgfpathrectangle{\pgfqpoint{0.407592in}{0.222904in}}{\pgfqpoint{1.411163in}{1.411163in}}%
\pgfusepath{clip}%
\pgfsetbuttcap%
\pgfsetroundjoin%
\definecolor{currentfill}{rgb}{0.749020,0.772549,0.792157}%
\pgfsetfillcolor{currentfill}%
\pgfsetfillopacity{0.160000}%
\pgfsetlinewidth{0.000000pt}%
\definecolor{currentstroke}{rgb}{0.000000,0.000000,0.000000}%
\pgfsetstrokecolor{currentstroke}%
\pgfsetdash{}{0pt}%
\pgfpathmoveto{\pgfqpoint{0.546258in}{1.042058in}}%
\pgfpathlineto{\pgfqpoint{0.532143in}{1.053256in}}%
\pgfpathlineto{\pgfqpoint{0.518002in}{1.064474in}}%
\pgfpathlineto{\pgfqpoint{0.520840in}{1.050082in}}%
\pgfpathlineto{\pgfqpoint{0.526239in}{1.035773in}}%
\pgfpathlineto{\pgfqpoint{0.540302in}{1.025170in}}%
\pgfpathlineto{\pgfqpoint{0.554333in}{1.014591in}}%
\pgfpathlineto{\pgfqpoint{0.549080in}{1.028282in}}%
\pgfpathlineto{\pgfqpoint{0.546258in}{1.042058in}}%
\pgfpathclose%
\pgfusepath{fill}%
\end{pgfscope}%
\begin{pgfscope}%
\pgfpathrectangle{\pgfqpoint{0.407592in}{0.222904in}}{\pgfqpoint{1.411163in}{1.411163in}}%
\pgfusepath{clip}%
\pgfsetbuttcap%
\pgfsetroundjoin%
\definecolor{currentfill}{rgb}{0.749020,0.772549,0.792157}%
\pgfsetfillcolor{currentfill}%
\pgfsetfillopacity{0.160000}%
\pgfsetlinewidth{0.000000pt}%
\definecolor{currentstroke}{rgb}{0.000000,0.000000,0.000000}%
\pgfsetstrokecolor{currentstroke}%
\pgfsetdash{}{0pt}%
\pgfpathmoveto{\pgfqpoint{1.373668in}{0.827504in}}%
\pgfpathlineto{\pgfqpoint{1.385049in}{0.836217in}}%
\pgfpathlineto{\pgfqpoint{1.396468in}{0.844958in}}%
\pgfpathlineto{\pgfqpoint{1.410368in}{0.853742in}}%
\pgfpathlineto{\pgfqpoint{1.398461in}{0.844696in}}%
\pgfpathlineto{\pgfqpoint{1.386591in}{0.835679in}}%
\pgfusepath{fill}%
\end{pgfscope}%
\begin{pgfscope}%
\pgfpathrectangle{\pgfqpoint{0.407592in}{0.222904in}}{\pgfqpoint{1.411163in}{1.411163in}}%
\pgfusepath{clip}%
\pgfsetbuttcap%
\pgfsetroundjoin%
\definecolor{currentfill}{rgb}{0.749020,0.772549,0.792157}%
\pgfsetfillcolor{currentfill}%
\pgfsetfillopacity{0.160000}%
\pgfsetlinewidth{0.000000pt}%
\definecolor{currentstroke}{rgb}{0.000000,0.000000,0.000000}%
\pgfsetstrokecolor{currentstroke}%
\pgfsetdash{}{0pt}%
\pgfpathmoveto{\pgfqpoint{1.323509in}{0.795998in}}%
\pgfpathlineto{\pgfqpoint{1.333668in}{0.804075in}}%
\pgfpathlineto{\pgfqpoint{1.343864in}{0.812181in}}%
\pgfpathlineto{\pgfqpoint{1.359406in}{0.819659in}}%
\pgfpathlineto{\pgfqpoint{1.373668in}{0.827504in}}%
\pgfpathlineto{\pgfqpoint{1.362323in}{0.818820in}}%
\pgfpathlineto{\pgfqpoint{1.351015in}{0.810164in}}%
\pgfpathlineto{\pgfqpoint{1.337847in}{0.802910in}}%
\pgfpathlineto{\pgfqpoint{1.323509in}{0.795998in}}%
\pgfpathclose%
\pgfusepath{fill}%
\end{pgfscope}%
\begin{pgfscope}%
\pgfpathrectangle{\pgfqpoint{0.407592in}{0.222904in}}{\pgfqpoint{1.411163in}{1.411163in}}%
\pgfusepath{clip}%
\pgfsetbuttcap%
\pgfsetroundjoin%
\definecolor{currentfill}{rgb}{0.749020,0.772549,0.792157}%
\pgfsetfillcolor{currentfill}%
\pgfsetfillopacity{0.160000}%
\pgfsetlinewidth{0.000000pt}%
\definecolor{currentstroke}{rgb}{0.000000,0.000000,0.000000}%
\pgfsetstrokecolor{currentstroke}%
\pgfsetdash{}{0pt}%
\pgfpathmoveto{\pgfqpoint{1.510531in}{0.930804in}}%
\pgfpathlineto{\pgfqpoint{1.523473in}{0.940667in}}%
\pgfpathlineto{\pgfqpoint{1.536450in}{0.950555in}}%
\pgfpathlineto{\pgfqpoint{1.548074in}{0.962682in}}%
\pgfpathlineto{\pgfqpoint{1.557570in}{0.975003in}}%
\pgfpathlineto{\pgfqpoint{1.544079in}{0.964481in}}%
\pgfpathlineto{\pgfqpoint{1.530618in}{0.953983in}}%
\pgfpathlineto{\pgfqpoint{1.521573in}{0.942299in}}%
\pgfpathlineto{\pgfqpoint{1.510531in}{0.930804in}}%
\pgfpathclose%
\pgfusepath{fill}%
\end{pgfscope}%
\begin{pgfscope}%
\pgfpathrectangle{\pgfqpoint{0.407592in}{0.222904in}}{\pgfqpoint{1.411163in}{1.411163in}}%
\pgfusepath{clip}%
\pgfsetbuttcap%
\pgfsetroundjoin%
\definecolor{currentfill}{rgb}{0.749020,0.772549,0.792157}%
\pgfsetfillcolor{currentfill}%
\pgfsetfillopacity{0.160000}%
\pgfsetlinewidth{0.000000pt}%
\definecolor{currentstroke}{rgb}{0.000000,0.000000,0.000000}%
\pgfsetstrokecolor{currentstroke}%
\pgfsetdash{}{0pt}%
\pgfpathmoveto{\pgfqpoint{1.274068in}{0.768318in}}%
\pgfpathlineto{\pgfqpoint{1.282801in}{0.775814in}}%
\pgfpathlineto{\pgfqpoint{1.291569in}{0.783339in}}%
\pgfpathlineto{\pgfqpoint{1.308060in}{0.789463in}}%
\pgfpathlineto{\pgfqpoint{1.323509in}{0.795998in}}%
\pgfpathlineto{\pgfqpoint{1.313386in}{0.787950in}}%
\pgfpathlineto{\pgfqpoint{1.303300in}{0.779931in}}%
\pgfpathlineto{\pgfqpoint{1.289156in}{0.773935in}}%
\pgfpathlineto{\pgfqpoint{1.274068in}{0.768318in}}%
\pgfpathclose%
\pgfusepath{fill}%
\end{pgfscope}%
\begin{pgfscope}%
\pgfpathrectangle{\pgfqpoint{0.407592in}{0.222904in}}{\pgfqpoint{1.411163in}{1.411163in}}%
\pgfusepath{clip}%
\pgfsetbuttcap%
\pgfsetroundjoin%
\definecolor{currentfill}{rgb}{0.749020,0.772549,0.792157}%
\pgfsetfillcolor{currentfill}%
\pgfsetfillopacity{0.160000}%
\pgfsetlinewidth{0.000000pt}%
\definecolor{currentstroke}{rgb}{0.000000,0.000000,0.000000}%
\pgfsetstrokecolor{currentstroke}%
\pgfsetdash{}{0pt}%
\pgfpathmoveto{\pgfqpoint{0.993314in}{0.720565in}}%
\pgfpathlineto{\pgfqpoint{0.987700in}{0.726939in}}%
\pgfpathlineto{\pgfqpoint{0.982060in}{0.733342in}}%
\pgfpathlineto{\pgfqpoint{0.998974in}{0.730880in}}%
\pgfpathlineto{\pgfqpoint{1.016319in}{0.728905in}}%
\pgfpathlineto{\pgfqpoint{1.020177in}{0.722725in}}%
\pgfpathlineto{\pgfqpoint{1.024018in}{0.716575in}}%
\pgfpathlineto{\pgfqpoint{1.008474in}{0.718351in}}%
\pgfpathlineto{\pgfqpoint{0.993314in}{0.720565in}}%
\pgfpathclose%
\pgfusepath{fill}%
\end{pgfscope}%
\begin{pgfscope}%
\pgfpathrectangle{\pgfqpoint{0.407592in}{0.222904in}}{\pgfqpoint{1.411163in}{1.411163in}}%
\pgfusepath{clip}%
\pgfsetbuttcap%
\pgfsetroundjoin%
\definecolor{currentfill}{rgb}{0.749020,0.772549,0.792157}%
\pgfsetfillcolor{currentfill}%
\pgfsetfillopacity{0.160000}%
\pgfsetlinewidth{0.000000pt}%
\definecolor{currentstroke}{rgb}{0.000000,0.000000,0.000000}%
\pgfsetstrokecolor{currentstroke}%
\pgfsetdash{}{0pt}%
\pgfpathmoveto{\pgfqpoint{0.379152in}{1.212375in}}%
\pgfpathlineto{\pgfqpoint{0.364968in}{1.224381in}}%
\pgfpathlineto{\pgfqpoint{0.350763in}{1.236403in}}%
\pgfpathlineto{\pgfqpoint{0.346879in}{1.218622in}}%
\pgfpathlineto{\pgfqpoint{0.346239in}{1.200731in}}%
\pgfpathlineto{\pgfqpoint{0.360700in}{1.189260in}}%
\pgfpathlineto{\pgfqpoint{0.375133in}{1.177810in}}%
\pgfpathlineto{\pgfqpoint{0.375583in}{1.195141in}}%
\pgfpathlineto{\pgfqpoint{0.379152in}{1.212375in}}%
\pgfpathclose%
\pgfusepath{fill}%
\end{pgfscope}%
\begin{pgfscope}%
\pgfpathrectangle{\pgfqpoint{0.407592in}{0.222904in}}{\pgfqpoint{1.411163in}{1.411163in}}%
\pgfusepath{clip}%
\pgfsetbuttcap%
\pgfsetroundjoin%
\definecolor{currentfill}{rgb}{0.749020,0.772549,0.792157}%
\pgfsetfillcolor{currentfill}%
\pgfsetfillopacity{0.160000}%
\pgfsetlinewidth{0.000000pt}%
\definecolor{currentstroke}{rgb}{0.000000,0.000000,0.000000}%
\pgfsetstrokecolor{currentstroke}%
\pgfsetdash{}{0pt}%
\pgfpathmoveto{\pgfqpoint{1.639161in}{1.038634in}}%
\pgfpathlineto{\pgfqpoint{1.652867in}{1.049323in}}%
\pgfpathlineto{\pgfqpoint{1.666604in}{1.060037in}}%
\pgfpathlineto{\pgfqpoint{1.675139in}{1.074978in}}%
\pgfpathlineto{\pgfqpoint{1.681003in}{1.089985in}}%
\pgfpathlineto{\pgfqpoint{1.667054in}{1.078665in}}%
\pgfpathlineto{\pgfqpoint{1.653131in}{1.067365in}}%
\pgfpathlineto{\pgfqpoint{1.647416in}{1.052965in}}%
\pgfpathlineto{\pgfqpoint{1.639161in}{1.038634in}}%
\pgfpathclose%
\pgfusepath{fill}%
\end{pgfscope}%
\begin{pgfscope}%
\pgfpathrectangle{\pgfqpoint{0.407592in}{0.222904in}}{\pgfqpoint{1.411163in}{1.411163in}}%
\pgfusepath{clip}%
\pgfsetbuttcap%
\pgfsetroundjoin%
\definecolor{currentfill}{rgb}{0.749020,0.772549,0.792157}%
\pgfsetfillcolor{currentfill}%
\pgfsetfillopacity{0.160000}%
\pgfsetlinewidth{0.000000pt}%
\definecolor{currentstroke}{rgb}{0.000000,0.000000,0.000000}%
\pgfsetstrokecolor{currentstroke}%
\pgfsetdash{}{0pt}%
\pgfpathmoveto{\pgfqpoint{1.821891in}{1.204326in}}%
\pgfpathlineto{\pgfqpoint{1.836121in}{1.215875in}}%
\pgfpathlineto{\pgfqpoint{1.850378in}{1.227445in}}%
\pgfpathlineto{\pgfqpoint{1.853439in}{1.245867in}}%
\pgfpathlineto{\pgfqpoint{1.853116in}{1.264153in}}%
\pgfpathlineto{\pgfqpoint{1.838982in}{1.252044in}}%
\pgfpathlineto{\pgfqpoint{1.824868in}{1.239951in}}%
\pgfpathlineto{\pgfqpoint{1.825009in}{1.222200in}}%
\pgfpathlineto{\pgfqpoint{1.821891in}{1.204326in}}%
\pgfpathclose%
\pgfusepath{fill}%
\end{pgfscope}%
\begin{pgfscope}%
\pgfpathrectangle{\pgfqpoint{0.407592in}{0.222904in}}{\pgfqpoint{1.411163in}{1.411163in}}%
\pgfusepath{clip}%
\pgfsetbuttcap%
\pgfsetroundjoin%
\definecolor{currentfill}{rgb}{0.749020,0.772549,0.792157}%
\pgfsetfillcolor{currentfill}%
\pgfsetfillopacity{0.160000}%
\pgfsetlinewidth{0.000000pt}%
\definecolor{currentstroke}{rgb}{0.000000,0.000000,0.000000}%
\pgfsetstrokecolor{currentstroke}%
\pgfsetdash{}{0pt}%
\pgfpathmoveto{\pgfqpoint{1.121116in}{0.715477in}}%
\pgfpathlineto{\pgfqpoint{1.122908in}{0.721566in}}%
\pgfpathlineto{\pgfqpoint{1.124709in}{0.727683in}}%
\pgfpathlineto{\pgfqpoint{1.142681in}{0.729261in}}%
\pgfpathlineto{\pgfqpoint{1.160396in}{0.731335in}}%
\pgfpathlineto{\pgfqpoint{1.156735in}{0.725033in}}%
\pgfpathlineto{\pgfqpoint{1.153092in}{0.718760in}}%
\pgfpathlineto{\pgfqpoint{1.137218in}{0.716895in}}%
\pgfpathlineto{\pgfqpoint{1.121116in}{0.715477in}}%
\pgfpathclose%
\pgfusepath{fill}%
\end{pgfscope}%
\begin{pgfscope}%
\pgfpathrectangle{\pgfqpoint{0.407592in}{0.222904in}}{\pgfqpoint{1.411163in}{1.411163in}}%
\pgfusepath{clip}%
\pgfsetbuttcap%
\pgfsetroundjoin%
\definecolor{currentfill}{rgb}{0.749020,0.772549,0.792157}%
\pgfsetfillcolor{currentfill}%
\pgfsetfillopacity{0.160000}%
\pgfsetlinewidth{0.000000pt}%
\definecolor{currentstroke}{rgb}{0.000000,0.000000,0.000000}%
\pgfsetstrokecolor{currentstroke}%
\pgfsetdash{}{0pt}%
\pgfpathmoveto{\pgfqpoint{0.637842in}{0.951627in}}%
\pgfpathlineto{\pgfqpoint{0.624003in}{0.962061in}}%
\pgfpathlineto{\pgfqpoint{0.610133in}{0.972519in}}%
\pgfpathlineto{\pgfqpoint{0.617234in}{0.960232in}}%
\pgfpathlineto{\pgfqpoint{0.626505in}{0.948151in}}%
\pgfpathlineto{\pgfqpoint{0.640001in}{0.938325in}}%
\pgfpathlineto{\pgfqpoint{0.653460in}{0.928526in}}%
\pgfpathlineto{\pgfqpoint{0.644631in}{0.939976in}}%
\pgfpathlineto{\pgfqpoint{0.637842in}{0.951627in}}%
\pgfpathclose%
\pgfusepath{fill}%
\end{pgfscope}%
\begin{pgfscope}%
\pgfpathrectangle{\pgfqpoint{0.407592in}{0.222904in}}{\pgfqpoint{1.411163in}{1.411163in}}%
\pgfusepath{clip}%
\pgfsetbuttcap%
\pgfsetroundjoin%
\definecolor{currentfill}{rgb}{0.749020,0.772549,0.792157}%
\pgfsetfillcolor{currentfill}%
\pgfsetfillopacity{0.160000}%
\pgfsetlinewidth{0.000000pt}%
\definecolor{currentstroke}{rgb}{0.000000,0.000000,0.000000}%
\pgfsetstrokecolor{currentstroke}%
\pgfsetdash{}{0pt}%
\pgfpathmoveto{\pgfqpoint{0.920274in}{0.747773in}}%
\pgfpathlineto{\pgfqpoint{0.911412in}{0.754929in}}%
\pgfpathlineto{\pgfqpoint{0.902514in}{0.762114in}}%
\pgfpathlineto{\pgfqpoint{0.918418in}{0.757428in}}%
\pgfpathlineto{\pgfqpoint{0.935136in}{0.753205in}}%
\pgfpathlineto{\pgfqpoint{0.942506in}{0.746425in}}%
\pgfpathlineto{\pgfqpoint{0.949845in}{0.739675in}}%
\pgfpathlineto{\pgfqpoint{0.934693in}{0.743513in}}%
\pgfpathlineto{\pgfqpoint{0.920274in}{0.747773in}}%
\pgfpathclose%
\pgfusepath{fill}%
\end{pgfscope}%
\begin{pgfscope}%
\pgfpathrectangle{\pgfqpoint{0.407592in}{0.222904in}}{\pgfqpoint{1.411163in}{1.411163in}}%
\pgfusepath{clip}%
\pgfsetbuttcap%
\pgfsetroundjoin%
\definecolor{currentfill}{rgb}{0.749020,0.772549,0.792157}%
\pgfsetfillcolor{currentfill}%
\pgfsetfillopacity{0.160000}%
\pgfsetlinewidth{0.000000pt}%
\definecolor{currentstroke}{rgb}{0.000000,0.000000,0.000000}%
\pgfsetstrokecolor{currentstroke}%
\pgfsetdash{}{0pt}%
\pgfpathmoveto{\pgfqpoint{1.024018in}{0.716575in}}%
\pgfpathlineto{\pgfqpoint{1.020177in}{0.722725in}}%
\pgfpathlineto{\pgfqpoint{1.016319in}{0.728905in}}%
\pgfpathlineto{\pgfqpoint{1.034013in}{0.727427in}}%
\pgfpathlineto{\pgfqpoint{1.051971in}{0.726455in}}%
\pgfpathlineto{\pgfqpoint{1.053970in}{0.720400in}}%
\pgfpathlineto{\pgfqpoint{1.055959in}{0.714373in}}%
\pgfpathlineto{\pgfqpoint{1.039871in}{0.715247in}}%
\pgfpathlineto{\pgfqpoint{1.024018in}{0.716575in}}%
\pgfpathclose%
\pgfusepath{fill}%
\end{pgfscope}%
\begin{pgfscope}%
\pgfpathrectangle{\pgfqpoint{0.407592in}{0.222904in}}{\pgfqpoint{1.411163in}{1.411163in}}%
\pgfusepath{clip}%
\pgfsetbuttcap%
\pgfsetroundjoin%
\definecolor{currentfill}{rgb}{0.749020,0.772549,0.792157}%
\pgfsetfillcolor{currentfill}%
\pgfsetfillopacity{0.160000}%
\pgfsetlinewidth{0.000000pt}%
\definecolor{currentstroke}{rgb}{0.000000,0.000000,0.000000}%
\pgfsetstrokecolor{currentstroke}%
\pgfsetdash{}{0pt}%
\pgfpathmoveto{\pgfqpoint{0.778258in}{0.834014in}}%
\pgfpathlineto{\pgfqpoint{0.765756in}{0.842970in}}%
\pgfpathlineto{\pgfqpoint{0.753215in}{0.851954in}}%
\pgfpathlineto{\pgfqpoint{0.765467in}{0.843236in}}%
\pgfpathlineto{\pgfqpoint{0.779242in}{0.834877in}}%
\pgfpathlineto{\pgfqpoint{0.790871in}{0.826488in}}%
\pgfpathlineto{\pgfqpoint{0.802460in}{0.818128in}}%
\pgfpathlineto{\pgfqpoint{0.789658in}{0.825903in}}%
\pgfpathlineto{\pgfqpoint{0.778258in}{0.834014in}}%
\pgfpathclose%
\pgfusepath{fill}%
\end{pgfscope}%
\begin{pgfscope}%
\pgfpathrectangle{\pgfqpoint{0.407592in}{0.222904in}}{\pgfqpoint{1.411163in}{1.411163in}}%
\pgfusepath{clip}%
\pgfsetbuttcap%
\pgfsetroundjoin%
\definecolor{currentfill}{rgb}{0.749020,0.772549,0.792157}%
\pgfsetfillcolor{currentfill}%
\pgfsetfillopacity{0.160000}%
\pgfsetlinewidth{0.000000pt}%
\definecolor{currentstroke}{rgb}{0.000000,0.000000,0.000000}%
\pgfsetstrokecolor{currentstroke}%
\pgfsetdash{}{0pt}%
\pgfpathmoveto{\pgfqpoint{1.434292in}{0.871917in}}%
\pgfpathlineto{\pgfqpoint{1.446309in}{0.881046in}}%
\pgfpathlineto{\pgfqpoint{1.458363in}{0.890204in}}%
\pgfpathlineto{\pgfqpoint{1.472457in}{0.900538in}}%
\pgfpathlineto{\pgfqpoint{1.484750in}{0.911159in}}%
\pgfpathlineto{\pgfqpoint{1.471911in}{0.901376in}}%
\pgfpathlineto{\pgfqpoint{1.459107in}{0.891619in}}%
\pgfpathlineto{\pgfqpoint{1.447537in}{0.881631in}}%
\pgfpathlineto{\pgfqpoint{1.434292in}{0.871917in}}%
\pgfpathclose%
\pgfusepath{fill}%
\end{pgfscope}%
\begin{pgfscope}%
\pgfpathrectangle{\pgfqpoint{0.407592in}{0.222904in}}{\pgfqpoint{1.411163in}{1.411163in}}%
\pgfusepath{clip}%
\pgfsetbuttcap%
\pgfsetroundjoin%
\definecolor{currentfill}{rgb}{0.749020,0.772549,0.792157}%
\pgfsetfillcolor{currentfill}%
\pgfsetfillopacity{0.160000}%
\pgfsetlinewidth{0.000000pt}%
\definecolor{currentstroke}{rgb}{0.000000,0.000000,0.000000}%
\pgfsetstrokecolor{currentstroke}%
\pgfsetdash{}{0pt}%
\pgfpathmoveto{\pgfqpoint{1.088532in}{0.714004in}}%
\pgfpathlineto{\pgfqpoint{1.088433in}{0.720011in}}%
\pgfpathlineto{\pgfqpoint{1.088334in}{0.726045in}}%
\pgfpathlineto{\pgfqpoint{1.106564in}{0.726610in}}%
\pgfpathlineto{\pgfqpoint{1.124709in}{0.727683in}}%
\pgfpathlineto{\pgfqpoint{1.122908in}{0.721566in}}%
\pgfpathlineto{\pgfqpoint{1.121116in}{0.715477in}}%
\pgfpathlineto{\pgfqpoint{1.104862in}{0.714512in}}%
\pgfpathlineto{\pgfqpoint{1.088532in}{0.714004in}}%
\pgfpathclose%
\pgfusepath{fill}%
\end{pgfscope}%
\begin{pgfscope}%
\pgfpathrectangle{\pgfqpoint{0.407592in}{0.222904in}}{\pgfqpoint{1.411163in}{1.411163in}}%
\pgfusepath{clip}%
\pgfsetbuttcap%
\pgfsetroundjoin%
\definecolor{currentfill}{rgb}{0.749020,0.772549,0.792157}%
\pgfsetfillcolor{currentfill}%
\pgfsetfillopacity{0.160000}%
\pgfsetlinewidth{0.000000pt}%
\definecolor{currentstroke}{rgb}{0.000000,0.000000,0.000000}%
\pgfsetstrokecolor{currentstroke}%
\pgfsetdash{}{0pt}%
\pgfpathmoveto{\pgfqpoint{0.518002in}{1.064474in}}%
\pgfpathlineto{\pgfqpoint{0.503834in}{1.075713in}}%
\pgfpathlineto{\pgfqpoint{0.489640in}{1.086973in}}%
\pgfpathlineto{\pgfqpoint{0.492483in}{1.071973in}}%
\pgfpathlineto{\pgfqpoint{0.498016in}{1.057053in}}%
\pgfpathlineto{\pgfqpoint{0.512144in}{1.046401in}}%
\pgfpathlineto{\pgfqpoint{0.526239in}{1.035773in}}%
\pgfpathlineto{\pgfqpoint{0.520840in}{1.050082in}}%
\pgfpathlineto{\pgfqpoint{0.518002in}{1.064474in}}%
\pgfpathclose%
\pgfusepath{fill}%
\end{pgfscope}%
\begin{pgfscope}%
\pgfpathrectangle{\pgfqpoint{0.407592in}{0.222904in}}{\pgfqpoint{1.411163in}{1.411163in}}%
\pgfusepath{clip}%
\pgfsetbuttcap%
\pgfsetroundjoin%
\definecolor{currentfill}{rgb}{0.749020,0.772549,0.792157}%
\pgfsetfillcolor{currentfill}%
\pgfsetfillopacity{0.160000}%
\pgfsetlinewidth{0.000000pt}%
\definecolor{currentstroke}{rgb}{0.000000,0.000000,0.000000}%
\pgfsetstrokecolor{currentstroke}%
\pgfsetdash{}{0pt}%
\pgfpathmoveto{\pgfqpoint{1.055959in}{0.714373in}}%
\pgfpathlineto{\pgfqpoint{1.053970in}{0.720400in}}%
\pgfpathlineto{\pgfqpoint{1.051971in}{0.726455in}}%
\pgfpathlineto{\pgfqpoint{1.070107in}{0.725994in}}%
\pgfpathlineto{\pgfqpoint{1.088334in}{0.726045in}}%
\pgfpathlineto{\pgfqpoint{1.088433in}{0.720011in}}%
\pgfpathlineto{\pgfqpoint{1.088532in}{0.714004in}}%
\pgfpathlineto{\pgfqpoint{1.072205in}{0.713958in}}%
\pgfpathlineto{\pgfqpoint{1.055959in}{0.714373in}}%
\pgfpathclose%
\pgfusepath{fill}%
\end{pgfscope}%
\begin{pgfscope}%
\pgfpathrectangle{\pgfqpoint{0.407592in}{0.222904in}}{\pgfqpoint{1.411163in}{1.411163in}}%
\pgfusepath{clip}%
\pgfsetbuttcap%
\pgfsetroundjoin%
\definecolor{currentfill}{rgb}{0.749020,0.772549,0.792157}%
\pgfsetfillcolor{currentfill}%
\pgfsetfillopacity{0.160000}%
\pgfsetlinewidth{0.000000pt}%
\definecolor{currentstroke}{rgb}{0.000000,0.000000,0.000000}%
\pgfsetstrokecolor{currentstroke}%
\pgfsetdash{}{0pt}%
\pgfpathmoveto{\pgfqpoint{1.850378in}{1.227445in}}%
\pgfpathlineto{\pgfqpoint{1.864660in}{1.239036in}}%
\pgfpathlineto{\pgfqpoint{1.867689in}{1.257729in}}%
\pgfpathlineto{\pgfqpoint{1.867269in}{1.276279in}}%
\pgfpathlineto{\pgfqpoint{1.853116in}{1.264153in}}%
\pgfpathlineto{\pgfqpoint{1.853439in}{1.245867in}}%
\pgfusepath{fill}%
\end{pgfscope}%
\begin{pgfscope}%
\pgfpathrectangle{\pgfqpoint{0.407592in}{0.222904in}}{\pgfqpoint{1.411163in}{1.411163in}}%
\pgfusepath{clip}%
\pgfsetbuttcap%
\pgfsetroundjoin%
\definecolor{currentfill}{rgb}{0.749020,0.772549,0.792157}%
\pgfsetfillcolor{currentfill}%
\pgfsetfillopacity{0.160000}%
\pgfsetlinewidth{0.000000pt}%
\definecolor{currentstroke}{rgb}{0.000000,0.000000,0.000000}%
\pgfsetstrokecolor{currentstroke}%
\pgfsetdash{}{0pt}%
\pgfpathmoveto{\pgfqpoint{0.350763in}{1.236403in}}%
\pgfpathlineto{\pgfqpoint{0.336539in}{1.248443in}}%
\pgfpathlineto{\pgfqpoint{0.322294in}{1.260499in}}%
\pgfpathlineto{\pgfqpoint{0.318082in}{1.242180in}}%
\pgfpathlineto{\pgfqpoint{0.317235in}{1.223739in}}%
\pgfpathlineto{\pgfqpoint{0.331751in}{1.212225in}}%
\pgfpathlineto{\pgfqpoint{0.346239in}{1.200731in}}%
\pgfpathlineto{\pgfqpoint{0.346879in}{1.218622in}}%
\pgfpathlineto{\pgfqpoint{0.350763in}{1.236403in}}%
\pgfpathclose%
\pgfusepath{fill}%
\end{pgfscope}%
\begin{pgfscope}%
\pgfpathrectangle{\pgfqpoint{0.407592in}{0.222904in}}{\pgfqpoint{1.411163in}{1.411163in}}%
\pgfusepath{clip}%
\pgfsetbuttcap%
\pgfsetroundjoin%
\definecolor{currentfill}{rgb}{0.749020,0.772549,0.792157}%
\pgfsetfillcolor{currentfill}%
\pgfsetfillopacity{0.160000}%
\pgfsetlinewidth{0.000000pt}%
\definecolor{currentstroke}{rgb}{0.000000,0.000000,0.000000}%
\pgfsetstrokecolor{currentstroke}%
\pgfsetdash{}{0pt}%
\pgfpathmoveto{\pgfqpoint{1.194714in}{0.736924in}}%
\pgfpathlineto{\pgfqpoint{1.200181in}{0.743536in}}%
\pgfpathlineto{\pgfqpoint{1.205671in}{0.750178in}}%
\pgfpathlineto{\pgfqpoint{1.223823in}{0.754014in}}%
\pgfpathlineto{\pgfqpoint{1.241326in}{0.758332in}}%
\pgfpathlineto{\pgfqpoint{1.234156in}{0.751318in}}%
\pgfpathlineto{\pgfqpoint{1.227016in}{0.744334in}}%
\pgfpathlineto{\pgfqpoint{1.211156in}{0.740410in}}%
\pgfpathlineto{\pgfqpoint{1.194714in}{0.736924in}}%
\pgfpathclose%
\pgfusepath{fill}%
\end{pgfscope}%
\begin{pgfscope}%
\pgfpathrectangle{\pgfqpoint{0.407592in}{0.222904in}}{\pgfqpoint{1.411163in}{1.411163in}}%
\pgfusepath{clip}%
\pgfsetbuttcap%
\pgfsetroundjoin%
\definecolor{currentfill}{rgb}{0.749020,0.772549,0.792157}%
\pgfsetfillcolor{currentfill}%
\pgfsetfillopacity{0.160000}%
\pgfsetlinewidth{0.000000pt}%
\definecolor{currentstroke}{rgb}{0.000000,0.000000,0.000000}%
\pgfsetstrokecolor{currentstroke}%
\pgfsetdash{}{0pt}%
\pgfpathmoveto{\pgfqpoint{0.706932in}{0.889595in}}%
\pgfpathlineto{\pgfqpoint{0.693618in}{0.899288in}}%
\pgfpathlineto{\pgfqpoint{0.680268in}{0.909008in}}%
\pgfpathlineto{\pgfqpoint{0.690534in}{0.898440in}}%
\pgfpathlineto{\pgfqpoint{0.702659in}{0.888170in}}%
\pgfpathlineto{\pgfqpoint{0.715357in}{0.879073in}}%
\pgfpathlineto{\pgfqpoint{0.728016in}{0.870005in}}%
\pgfpathlineto{\pgfqpoint{0.716608in}{0.879658in}}%
\pgfpathlineto{\pgfqpoint{0.706932in}{0.889595in}}%
\pgfpathclose%
\pgfusepath{fill}%
\end{pgfscope}%
\begin{pgfscope}%
\pgfpathrectangle{\pgfqpoint{0.407592in}{0.222904in}}{\pgfqpoint{1.411163in}{1.411163in}}%
\pgfusepath{clip}%
\pgfsetbuttcap%
\pgfsetroundjoin%
\definecolor{currentfill}{rgb}{0.749020,0.772549,0.792157}%
\pgfsetfillcolor{currentfill}%
\pgfsetfillopacity{0.160000}%
\pgfsetlinewidth{0.000000pt}%
\definecolor{currentstroke}{rgb}{0.000000,0.000000,0.000000}%
\pgfsetstrokecolor{currentstroke}%
\pgfsetdash{}{0pt}%
\pgfpathmoveto{\pgfqpoint{0.873435in}{0.772777in}}%
\pgfpathlineto{\pgfqpoint{0.863150in}{0.780474in}}%
\pgfpathlineto{\pgfqpoint{0.852825in}{0.788200in}}%
\pgfpathlineto{\pgfqpoint{0.868190in}{0.782163in}}%
\pgfpathlineto{\pgfqpoint{0.884607in}{0.776574in}}%
\pgfpathlineto{\pgfqpoint{0.893579in}{0.769329in}}%
\pgfpathlineto{\pgfqpoint{0.902514in}{0.762114in}}%
\pgfpathlineto{\pgfqpoint{0.887497in}{0.767239in}}%
\pgfpathlineto{\pgfqpoint{0.873435in}{0.772777in}}%
\pgfpathclose%
\pgfusepath{fill}%
\end{pgfscope}%
\begin{pgfscope}%
\pgfpathrectangle{\pgfqpoint{0.407592in}{0.222904in}}{\pgfqpoint{1.411163in}{1.411163in}}%
\pgfusepath{clip}%
\pgfsetbuttcap%
\pgfsetroundjoin%
\definecolor{currentfill}{rgb}{0.749020,0.772549,0.792157}%
\pgfsetfillcolor{currentfill}%
\pgfsetfillopacity{0.160000}%
\pgfsetlinewidth{0.000000pt}%
\definecolor{currentstroke}{rgb}{0.000000,0.000000,0.000000}%
\pgfsetstrokecolor{currentstroke}%
\pgfsetdash{}{0pt}%
\pgfpathmoveto{\pgfqpoint{0.825518in}{0.801494in}}%
\pgfpathlineto{\pgfqpoint{0.814009in}{0.809797in}}%
\pgfpathlineto{\pgfqpoint{0.802460in}{0.818128in}}%
\pgfpathlineto{\pgfqpoint{0.816613in}{0.810728in}}%
\pgfpathlineto{\pgfqpoint{0.832058in}{0.803742in}}%
\pgfpathlineto{\pgfqpoint{0.842461in}{0.795956in}}%
\pgfpathlineto{\pgfqpoint{0.852825in}{0.788200in}}%
\pgfpathlineto{\pgfqpoint{0.838581in}{0.794655in}}%
\pgfpathlineto{\pgfqpoint{0.825518in}{0.801494in}}%
\pgfpathclose%
\pgfusepath{fill}%
\end{pgfscope}%
\begin{pgfscope}%
\pgfpathrectangle{\pgfqpoint{0.407592in}{0.222904in}}{\pgfqpoint{1.411163in}{1.411163in}}%
\pgfusepath{clip}%
\pgfsetbuttcap%
\pgfsetroundjoin%
\definecolor{currentfill}{rgb}{0.749020,0.772549,0.792157}%
\pgfsetfillcolor{currentfill}%
\pgfsetfillopacity{0.160000}%
\pgfsetlinewidth{0.000000pt}%
\definecolor{currentstroke}{rgb}{0.000000,0.000000,0.000000}%
\pgfsetstrokecolor{currentstroke}%
\pgfsetdash{}{0pt}%
\pgfpathmoveto{\pgfqpoint{1.536450in}{0.950555in}}%
\pgfpathlineto{\pgfqpoint{1.549462in}{0.960470in}}%
\pgfpathlineto{\pgfqpoint{1.562509in}{0.970412in}}%
\pgfpathlineto{\pgfqpoint{1.574706in}{0.983166in}}%
\pgfpathlineto{\pgfqpoint{1.584645in}{0.996118in}}%
\pgfpathlineto{\pgfqpoint{1.571092in}{0.985549in}}%
\pgfpathlineto{\pgfqpoint{1.557570in}{0.975003in}}%
\pgfpathlineto{\pgfqpoint{1.548074in}{0.962682in}}%
\pgfpathlineto{\pgfqpoint{1.536450in}{0.950555in}}%
\pgfpathclose%
\pgfusepath{fill}%
\end{pgfscope}%
\begin{pgfscope}%
\pgfpathrectangle{\pgfqpoint{0.407592in}{0.222904in}}{\pgfqpoint{1.411163in}{1.411163in}}%
\pgfusepath{clip}%
\pgfsetbuttcap%
\pgfsetroundjoin%
\definecolor{currentfill}{rgb}{0.749020,0.772549,0.792157}%
\pgfsetfillcolor{currentfill}%
\pgfsetfillopacity{0.160000}%
\pgfsetlinewidth{0.000000pt}%
\definecolor{currentstroke}{rgb}{0.000000,0.000000,0.000000}%
\pgfsetstrokecolor{currentstroke}%
\pgfsetdash{}{0pt}%
\pgfpathmoveto{\pgfqpoint{1.666604in}{1.060037in}}%
\pgfpathlineto{\pgfqpoint{1.680372in}{1.070774in}}%
\pgfpathlineto{\pgfqpoint{1.694172in}{1.081536in}}%
\pgfpathlineto{\pgfqpoint{1.702976in}{1.097082in}}%
\pgfpathlineto{\pgfqpoint{1.708976in}{1.112688in}}%
\pgfpathlineto{\pgfqpoint{1.694977in}{1.101326in}}%
\pgfpathlineto{\pgfqpoint{1.681003in}{1.089985in}}%
\pgfpathlineto{\pgfqpoint{1.675139in}{1.074978in}}%
\pgfpathlineto{\pgfqpoint{1.666604in}{1.060037in}}%
\pgfpathclose%
\pgfusepath{fill}%
\end{pgfscope}%
\begin{pgfscope}%
\pgfpathrectangle{\pgfqpoint{0.407592in}{0.222904in}}{\pgfqpoint{1.411163in}{1.411163in}}%
\pgfusepath{clip}%
\pgfsetbuttcap%
\pgfsetroundjoin%
\definecolor{currentfill}{rgb}{0.749020,0.772549,0.792157}%
\pgfsetfillcolor{currentfill}%
\pgfsetfillopacity{0.160000}%
\pgfsetlinewidth{0.000000pt}%
\definecolor{currentstroke}{rgb}{0.000000,0.000000,0.000000}%
\pgfsetstrokecolor{currentstroke}%
\pgfsetdash{}{0pt}%
\pgfpathmoveto{\pgfqpoint{1.396468in}{0.844958in}}%
\pgfpathlineto{\pgfqpoint{1.407924in}{0.853728in}}%
\pgfpathlineto{\pgfqpoint{1.419417in}{0.862526in}}%
\pgfpathlineto{\pgfqpoint{1.434292in}{0.871917in}}%
\pgfpathlineto{\pgfqpoint{1.422311in}{0.862815in}}%
\pgfpathlineto{\pgfqpoint{1.410368in}{0.853742in}}%
\pgfusepath{fill}%
\end{pgfscope}%
\begin{pgfscope}%
\pgfpathrectangle{\pgfqpoint{0.407592in}{0.222904in}}{\pgfqpoint{1.411163in}{1.411163in}}%
\pgfusepath{clip}%
\pgfsetbuttcap%
\pgfsetroundjoin%
\definecolor{currentfill}{rgb}{0.749020,0.772549,0.792157}%
\pgfsetfillcolor{currentfill}%
\pgfsetfillopacity{0.160000}%
\pgfsetlinewidth{0.000000pt}%
\definecolor{currentstroke}{rgb}{0.000000,0.000000,0.000000}%
\pgfsetstrokecolor{currentstroke}%
\pgfsetdash{}{0pt}%
\pgfpathmoveto{\pgfqpoint{0.322294in}{1.260499in}}%
\pgfpathlineto{\pgfqpoint{0.308030in}{1.272572in}}%
\pgfpathlineto{\pgfqpoint{0.303647in}{1.253987in}}%
\pgfpathlineto{\pgfqpoint{0.302693in}{1.235276in}}%
\pgfpathlineto{\pgfqpoint{0.317235in}{1.223739in}}%
\pgfpathlineto{\pgfqpoint{0.318082in}{1.242180in}}%
\pgfusepath{fill}%
\end{pgfscope}%
\begin{pgfscope}%
\pgfpathrectangle{\pgfqpoint{0.407592in}{0.222904in}}{\pgfqpoint{1.411163in}{1.411163in}}%
\pgfusepath{clip}%
\pgfsetbuttcap%
\pgfsetroundjoin%
\definecolor{currentfill}{rgb}{0.749020,0.772549,0.792157}%
\pgfsetfillcolor{currentfill}%
\pgfsetfillopacity{0.160000}%
\pgfsetlinewidth{0.000000pt}%
\definecolor{currentstroke}{rgb}{0.000000,0.000000,0.000000}%
\pgfsetstrokecolor{currentstroke}%
\pgfsetdash{}{0pt}%
\pgfpathmoveto{\pgfqpoint{0.949845in}{0.739675in}}%
\pgfpathlineto{\pgfqpoint{0.942506in}{0.746425in}}%
\pgfpathlineto{\pgfqpoint{0.935136in}{0.753205in}}%
\pgfpathlineto{\pgfqpoint{0.952592in}{0.749468in}}%
\pgfpathlineto{\pgfqpoint{0.970704in}{0.746235in}}%
\pgfpathlineto{\pgfqpoint{0.976395in}{0.739774in}}%
\pgfpathlineto{\pgfqpoint{0.982060in}{0.733342in}}%
\pgfpathlineto{\pgfqpoint{0.965658in}{0.736279in}}%
\pgfpathlineto{\pgfqpoint{0.949845in}{0.739675in}}%
\pgfpathclose%
\pgfusepath{fill}%
\end{pgfscope}%
\begin{pgfscope}%
\pgfpathrectangle{\pgfqpoint{0.407592in}{0.222904in}}{\pgfqpoint{1.411163in}{1.411163in}}%
\pgfusepath{clip}%
\pgfsetbuttcap%
\pgfsetroundjoin%
\definecolor{currentfill}{rgb}{0.749020,0.772549,0.792157}%
\pgfsetfillcolor{currentfill}%
\pgfsetfillopacity{0.160000}%
\pgfsetlinewidth{0.000000pt}%
\definecolor{currentstroke}{rgb}{0.000000,0.000000,0.000000}%
\pgfsetstrokecolor{currentstroke}%
\pgfsetdash{}{0pt}%
\pgfpathmoveto{\pgfqpoint{0.610133in}{0.972519in}}%
\pgfpathlineto{\pgfqpoint{0.596231in}{0.983001in}}%
\pgfpathlineto{\pgfqpoint{0.582297in}{0.993507in}}%
\pgfpathlineto{\pgfqpoint{0.589699in}{0.980589in}}%
\pgfpathlineto{\pgfqpoint{0.599402in}{0.967883in}}%
\pgfpathlineto{\pgfqpoint{0.612972in}{0.958003in}}%
\pgfpathlineto{\pgfqpoint{0.626505in}{0.948151in}}%
\pgfpathlineto{\pgfqpoint{0.617234in}{0.960232in}}%
\pgfpathlineto{\pgfqpoint{0.610133in}{0.972519in}}%
\pgfpathclose%
\pgfusepath{fill}%
\end{pgfscope}%
\begin{pgfscope}%
\pgfpathrectangle{\pgfqpoint{0.407592in}{0.222904in}}{\pgfqpoint{1.411163in}{1.411163in}}%
\pgfusepath{clip}%
\pgfsetbuttcap%
\pgfsetroundjoin%
\definecolor{currentfill}{rgb}{0.749020,0.772549,0.792157}%
\pgfsetfillcolor{currentfill}%
\pgfsetfillopacity{0.160000}%
\pgfsetlinewidth{0.000000pt}%
\definecolor{currentstroke}{rgb}{0.000000,0.000000,0.000000}%
\pgfsetstrokecolor{currentstroke}%
\pgfsetdash{}{0pt}%
\pgfpathmoveto{\pgfqpoint{1.343864in}{0.812181in}}%
\pgfpathlineto{\pgfqpoint{1.354097in}{0.820316in}}%
\pgfpathlineto{\pgfqpoint{1.364368in}{0.828482in}}%
\pgfpathlineto{\pgfqpoint{1.381114in}{0.836524in}}%
\pgfpathlineto{\pgfqpoint{1.396468in}{0.844958in}}%
\pgfpathlineto{\pgfqpoint{1.385049in}{0.836217in}}%
\pgfpathlineto{\pgfqpoint{1.373668in}{0.827504in}}%
\pgfpathlineto{\pgfqpoint{1.359406in}{0.819659in}}%
\pgfpathlineto{\pgfqpoint{1.343864in}{0.812181in}}%
\pgfpathclose%
\pgfusepath{fill}%
\end{pgfscope}%
\begin{pgfscope}%
\pgfpathrectangle{\pgfqpoint{0.407592in}{0.222904in}}{\pgfqpoint{1.411163in}{1.411163in}}%
\pgfusepath{clip}%
\pgfsetbuttcap%
\pgfsetroundjoin%
\definecolor{currentfill}{rgb}{0.749020,0.772549,0.792157}%
\pgfsetfillcolor{currentfill}%
\pgfsetfillopacity{0.160000}%
\pgfsetlinewidth{0.000000pt}%
\definecolor{currentstroke}{rgb}{0.000000,0.000000,0.000000}%
\pgfsetstrokecolor{currentstroke}%
\pgfsetdash{}{0pt}%
\pgfpathmoveto{\pgfqpoint{1.241326in}{0.758332in}}%
\pgfpathlineto{\pgfqpoint{1.248526in}{0.765375in}}%
\pgfpathlineto{\pgfqpoint{1.255757in}{0.772447in}}%
\pgfpathlineto{\pgfqpoint{1.274108in}{0.777657in}}%
\pgfpathlineto{\pgfqpoint{1.291569in}{0.783339in}}%
\pgfpathlineto{\pgfqpoint{1.282801in}{0.775814in}}%
\pgfpathlineto{\pgfqpoint{1.274068in}{0.768318in}}%
\pgfpathlineto{\pgfqpoint{1.258100in}{0.763108in}}%
\pgfpathlineto{\pgfqpoint{1.241326in}{0.758332in}}%
\pgfpathclose%
\pgfusepath{fill}%
\end{pgfscope}%
\begin{pgfscope}%
\pgfpathrectangle{\pgfqpoint{0.407592in}{0.222904in}}{\pgfqpoint{1.411163in}{1.411163in}}%
\pgfusepath{clip}%
\pgfsetbuttcap%
\pgfsetroundjoin%
\definecolor{currentfill}{rgb}{0.749020,0.772549,0.792157}%
\pgfsetfillcolor{currentfill}%
\pgfsetfillopacity{0.160000}%
\pgfsetlinewidth{0.000000pt}%
\definecolor{currentstroke}{rgb}{0.000000,0.000000,0.000000}%
\pgfsetstrokecolor{currentstroke}%
\pgfsetdash{}{0pt}%
\pgfpathmoveto{\pgfqpoint{0.489640in}{1.086973in}}%
\pgfpathlineto{\pgfqpoint{0.475420in}{1.098254in}}%
\pgfpathlineto{\pgfqpoint{0.461173in}{1.109555in}}%
\pgfpathlineto{\pgfqpoint{0.464007in}{1.093956in}}%
\pgfpathlineto{\pgfqpoint{0.469662in}{1.078431in}}%
\pgfpathlineto{\pgfqpoint{0.483856in}{1.067730in}}%
\pgfpathlineto{\pgfqpoint{0.498016in}{1.057053in}}%
\pgfpathlineto{\pgfqpoint{0.492483in}{1.071973in}}%
\pgfpathlineto{\pgfqpoint{0.489640in}{1.086973in}}%
\pgfpathclose%
\pgfusepath{fill}%
\end{pgfscope}%
\begin{pgfscope}%
\pgfpathrectangle{\pgfqpoint{0.407592in}{0.222904in}}{\pgfqpoint{1.411163in}{1.411163in}}%
\pgfusepath{clip}%
\pgfsetbuttcap%
\pgfsetroundjoin%
\definecolor{currentfill}{rgb}{0.749020,0.772549,0.792157}%
\pgfsetfillcolor{currentfill}%
\pgfsetfillopacity{0.160000}%
\pgfsetlinewidth{0.000000pt}%
\definecolor{currentstroke}{rgb}{0.000000,0.000000,0.000000}%
\pgfsetstrokecolor{currentstroke}%
\pgfsetdash{}{0pt}%
\pgfpathmoveto{\pgfqpoint{1.160396in}{0.731335in}}%
\pgfpathlineto{\pgfqpoint{1.164073in}{0.737666in}}%
\pgfpathlineto{\pgfqpoint{1.167767in}{0.744025in}}%
\pgfpathlineto{\pgfqpoint{1.186956in}{0.746842in}}%
\pgfpathlineto{\pgfqpoint{1.205671in}{0.750178in}}%
\pgfpathlineto{\pgfqpoint{1.200181in}{0.743536in}}%
\pgfpathlineto{\pgfqpoint{1.194714in}{0.736924in}}%
\pgfpathlineto{\pgfqpoint{1.177768in}{0.733894in}}%
\pgfpathlineto{\pgfqpoint{1.160396in}{0.731335in}}%
\pgfpathclose%
\pgfusepath{fill}%
\end{pgfscope}%
\begin{pgfscope}%
\pgfpathrectangle{\pgfqpoint{0.407592in}{0.222904in}}{\pgfqpoint{1.411163in}{1.411163in}}%
\pgfusepath{clip}%
\pgfsetbuttcap%
\pgfsetroundjoin%
\definecolor{currentfill}{rgb}{0.749020,0.772549,0.792157}%
\pgfsetfillcolor{currentfill}%
\pgfsetfillopacity{0.160000}%
\pgfsetlinewidth{0.000000pt}%
\definecolor{currentstroke}{rgb}{0.000000,0.000000,0.000000}%
\pgfsetstrokecolor{currentstroke}%
\pgfsetdash{}{0pt}%
\pgfpathmoveto{\pgfqpoint{1.291569in}{0.783339in}}%
\pgfpathlineto{\pgfqpoint{1.300372in}{0.790894in}}%
\pgfpathlineto{\pgfqpoint{1.309209in}{0.798479in}}%
\pgfpathlineto{\pgfqpoint{1.327108in}{0.805109in}}%
\pgfpathlineto{\pgfqpoint{1.343864in}{0.812181in}}%
\pgfpathlineto{\pgfqpoint{1.333668in}{0.804075in}}%
\pgfpathlineto{\pgfqpoint{1.323509in}{0.795998in}}%
\pgfpathlineto{\pgfqpoint{1.308060in}{0.789463in}}%
\pgfpathlineto{\pgfqpoint{1.291569in}{0.783339in}}%
\pgfpathclose%
\pgfusepath{fill}%
\end{pgfscope}%
\begin{pgfscope}%
\pgfpathrectangle{\pgfqpoint{0.407592in}{0.222904in}}{\pgfqpoint{1.411163in}{1.411163in}}%
\pgfusepath{clip}%
\pgfsetbuttcap%
\pgfsetroundjoin%
\definecolor{currentfill}{rgb}{0.749020,0.772549,0.792157}%
\pgfsetfillcolor{currentfill}%
\pgfsetfillopacity{0.160000}%
\pgfsetlinewidth{0.000000pt}%
\definecolor{currentstroke}{rgb}{0.000000,0.000000,0.000000}%
\pgfsetstrokecolor{currentstroke}%
\pgfsetdash{}{0pt}%
\pgfpathmoveto{\pgfqpoint{1.458363in}{0.890204in}}%
\pgfpathlineto{\pgfqpoint{1.470455in}{0.899390in}}%
\pgfpathlineto{\pgfqpoint{1.482584in}{0.908604in}}%
\pgfpathlineto{\pgfqpoint{1.497521in}{0.919554in}}%
\pgfpathlineto{\pgfqpoint{1.510531in}{0.930804in}}%
\pgfpathlineto{\pgfqpoint{1.497623in}{0.920969in}}%
\pgfpathlineto{\pgfqpoint{1.484750in}{0.911159in}}%
\pgfpathlineto{\pgfqpoint{1.472457in}{0.900538in}}%
\pgfpathlineto{\pgfqpoint{1.458363in}{0.890204in}}%
\pgfpathclose%
\pgfusepath{fill}%
\end{pgfscope}%
\begin{pgfscope}%
\pgfpathrectangle{\pgfqpoint{0.407592in}{0.222904in}}{\pgfqpoint{1.411163in}{1.411163in}}%
\pgfusepath{clip}%
\pgfsetbuttcap%
\pgfsetroundjoin%
\definecolor{currentfill}{rgb}{0.749020,0.772549,0.792157}%
\pgfsetfillcolor{currentfill}%
\pgfsetfillopacity{0.160000}%
\pgfsetlinewidth{0.000000pt}%
\definecolor{currentstroke}{rgb}{0.000000,0.000000,0.000000}%
\pgfsetstrokecolor{currentstroke}%
\pgfsetdash{}{0pt}%
\pgfpathmoveto{\pgfqpoint{0.753215in}{0.851954in}}%
\pgfpathlineto{\pgfqpoint{0.740635in}{0.860966in}}%
\pgfpathlineto{\pgfqpoint{0.728016in}{0.870005in}}%
\pgfpathlineto{\pgfqpoint{0.741116in}{0.860685in}}%
\pgfpathlineto{\pgfqpoint{0.755860in}{0.851745in}}%
\pgfpathlineto{\pgfqpoint{0.767572in}{0.843296in}}%
\pgfpathlineto{\pgfqpoint{0.779242in}{0.834877in}}%
\pgfpathlineto{\pgfqpoint{0.765467in}{0.843236in}}%
\pgfpathlineto{\pgfqpoint{0.753215in}{0.851954in}}%
\pgfpathclose%
\pgfusepath{fill}%
\end{pgfscope}%
\begin{pgfscope}%
\pgfpathrectangle{\pgfqpoint{0.407592in}{0.222904in}}{\pgfqpoint{1.411163in}{1.411163in}}%
\pgfusepath{clip}%
\pgfsetbuttcap%
\pgfsetroundjoin%
\definecolor{currentfill}{rgb}{0.749020,0.772549,0.792157}%
\pgfsetfillcolor{currentfill}%
\pgfsetfillopacity{0.160000}%
\pgfsetlinewidth{0.000000pt}%
\definecolor{currentstroke}{rgb}{0.000000,0.000000,0.000000}%
\pgfsetstrokecolor{currentstroke}%
\pgfsetdash{}{0pt}%
\pgfpathmoveto{\pgfqpoint{1.694172in}{1.081536in}}%
\pgfpathlineto{\pgfqpoint{1.708003in}{1.092323in}}%
\pgfpathlineto{\pgfqpoint{1.721866in}{1.103134in}}%
\pgfpathlineto{\pgfqpoint{1.730926in}{1.119275in}}%
\pgfpathlineto{\pgfqpoint{1.737051in}{1.135472in}}%
\pgfpathlineto{\pgfqpoint{1.723001in}{1.124070in}}%
\pgfpathlineto{\pgfqpoint{1.708976in}{1.112688in}}%
\pgfpathlineto{\pgfqpoint{1.702976in}{1.097082in}}%
\pgfpathlineto{\pgfqpoint{1.694172in}{1.081536in}}%
\pgfpathclose%
\pgfusepath{fill}%
\end{pgfscope}%
\begin{pgfscope}%
\pgfpathrectangle{\pgfqpoint{0.407592in}{0.222904in}}{\pgfqpoint{1.411163in}{1.411163in}}%
\pgfusepath{clip}%
\pgfsetbuttcap%
\pgfsetroundjoin%
\definecolor{currentfill}{rgb}{0.749020,0.772549,0.792157}%
\pgfsetfillcolor{currentfill}%
\pgfsetfillopacity{0.160000}%
\pgfsetlinewidth{0.000000pt}%
\definecolor{currentstroke}{rgb}{0.000000,0.000000,0.000000}%
\pgfsetstrokecolor{currentstroke}%
\pgfsetdash{}{0pt}%
\pgfpathmoveto{\pgfqpoint{0.982060in}{0.733342in}}%
\pgfpathlineto{\pgfqpoint{0.976395in}{0.739774in}}%
\pgfpathlineto{\pgfqpoint{0.970704in}{0.746235in}}%
\pgfpathlineto{\pgfqpoint{0.989386in}{0.743524in}}%
\pgfpathlineto{\pgfqpoint{1.008549in}{0.741349in}}%
\pgfpathlineto{\pgfqpoint{1.012443in}{0.735113in}}%
\pgfpathlineto{\pgfqpoint{1.016319in}{0.728905in}}%
\pgfpathlineto{\pgfqpoint{0.998974in}{0.730880in}}%
\pgfpathlineto{\pgfqpoint{0.982060in}{0.733342in}}%
\pgfpathclose%
\pgfusepath{fill}%
\end{pgfscope}%
\begin{pgfscope}%
\pgfpathrectangle{\pgfqpoint{0.407592in}{0.222904in}}{\pgfqpoint{1.411163in}{1.411163in}}%
\pgfusepath{clip}%
\pgfsetbuttcap%
\pgfsetroundjoin%
\definecolor{currentfill}{rgb}{0.749020,0.772549,0.792157}%
\pgfsetfillcolor{currentfill}%
\pgfsetfillopacity{0.160000}%
\pgfsetlinewidth{0.000000pt}%
\definecolor{currentstroke}{rgb}{0.000000,0.000000,0.000000}%
\pgfsetstrokecolor{currentstroke}%
\pgfsetdash{}{0pt}%
\pgfpathmoveto{\pgfqpoint{1.562509in}{0.970412in}}%
\pgfpathlineto{\pgfqpoint{1.575592in}{0.980381in}}%
\pgfpathlineto{\pgfqpoint{1.588709in}{0.990377in}}%
\pgfpathlineto{\pgfqpoint{1.601470in}{1.003751in}}%
\pgfpathlineto{\pgfqpoint{1.611841in}{1.017328in}}%
\pgfpathlineto{\pgfqpoint{1.598228in}{1.006711in}}%
\pgfpathlineto{\pgfqpoint{1.584645in}{0.996118in}}%
\pgfpathlineto{\pgfqpoint{1.574706in}{0.983166in}}%
\pgfpathlineto{\pgfqpoint{1.562509in}{0.970412in}}%
\pgfpathclose%
\pgfusepath{fill}%
\end{pgfscope}%
\begin{pgfscope}%
\pgfpathrectangle{\pgfqpoint{0.407592in}{0.222904in}}{\pgfqpoint{1.411163in}{1.411163in}}%
\pgfusepath{clip}%
\pgfsetbuttcap%
\pgfsetroundjoin%
\definecolor{currentfill}{rgb}{0.749020,0.772549,0.792157}%
\pgfsetfillcolor{currentfill}%
\pgfsetfillopacity{0.160000}%
\pgfsetlinewidth{0.000000pt}%
\definecolor{currentstroke}{rgb}{0.000000,0.000000,0.000000}%
\pgfsetstrokecolor{currentstroke}%
\pgfsetdash{}{0pt}%
\pgfpathmoveto{\pgfqpoint{0.680268in}{0.909008in}}%
\pgfpathlineto{\pgfqpoint{0.666882in}{0.918754in}}%
\pgfpathlineto{\pgfqpoint{0.653460in}{0.928526in}}%
\pgfpathlineto{\pgfqpoint{0.664307in}{0.917332in}}%
\pgfpathlineto{\pgfqpoint{0.677143in}{0.906448in}}%
\pgfpathlineto{\pgfqpoint{0.689921in}{0.897295in}}%
\pgfpathlineto{\pgfqpoint{0.702659in}{0.888170in}}%
\pgfpathlineto{\pgfqpoint{0.690534in}{0.898440in}}%
\pgfpathlineto{\pgfqpoint{0.680268in}{0.909008in}}%
\pgfpathclose%
\pgfusepath{fill}%
\end{pgfscope}%
\begin{pgfscope}%
\pgfpathrectangle{\pgfqpoint{0.407592in}{0.222904in}}{\pgfqpoint{1.411163in}{1.411163in}}%
\pgfusepath{clip}%
\pgfsetbuttcap%
\pgfsetroundjoin%
\definecolor{currentfill}{rgb}{0.749020,0.772549,0.792157}%
\pgfsetfillcolor{currentfill}%
\pgfsetfillopacity{0.160000}%
\pgfsetlinewidth{0.000000pt}%
\definecolor{currentstroke}{rgb}{0.000000,0.000000,0.000000}%
\pgfsetstrokecolor{currentstroke}%
\pgfsetdash{}{0pt}%
\pgfpathmoveto{\pgfqpoint{0.902514in}{0.762114in}}%
\pgfpathlineto{\pgfqpoint{0.893579in}{0.769329in}}%
\pgfpathlineto{\pgfqpoint{0.884607in}{0.776574in}}%
\pgfpathlineto{\pgfqpoint{0.902004in}{0.771462in}}%
\pgfpathlineto{\pgfqpoint{0.920300in}{0.766854in}}%
\pgfpathlineto{\pgfqpoint{0.927734in}{0.760015in}}%
\pgfpathlineto{\pgfqpoint{0.935136in}{0.753205in}}%
\pgfpathlineto{\pgfqpoint{0.918418in}{0.757428in}}%
\pgfpathlineto{\pgfqpoint{0.902514in}{0.762114in}}%
\pgfpathclose%
\pgfusepath{fill}%
\end{pgfscope}%
\begin{pgfscope}%
\pgfpathrectangle{\pgfqpoint{0.407592in}{0.222904in}}{\pgfqpoint{1.411163in}{1.411163in}}%
\pgfusepath{clip}%
\pgfsetbuttcap%
\pgfsetroundjoin%
\definecolor{currentfill}{rgb}{0.749020,0.772549,0.792157}%
\pgfsetfillcolor{currentfill}%
\pgfsetfillopacity{0.160000}%
\pgfsetlinewidth{0.000000pt}%
\definecolor{currentstroke}{rgb}{0.000000,0.000000,0.000000}%
\pgfsetstrokecolor{currentstroke}%
\pgfsetdash{}{0pt}%
\pgfpathmoveto{\pgfqpoint{1.124709in}{0.727683in}}%
\pgfpathlineto{\pgfqpoint{1.126518in}{0.733830in}}%
\pgfpathlineto{\pgfqpoint{1.128335in}{0.740005in}}%
\pgfpathlineto{\pgfqpoint{1.148195in}{0.741742in}}%
\pgfpathlineto{\pgfqpoint{1.167767in}{0.744025in}}%
\pgfpathlineto{\pgfqpoint{1.164073in}{0.737666in}}%
\pgfpathlineto{\pgfqpoint{1.160396in}{0.731335in}}%
\pgfpathlineto{\pgfqpoint{1.142681in}{0.729261in}}%
\pgfpathlineto{\pgfqpoint{1.124709in}{0.727683in}}%
\pgfpathclose%
\pgfusepath{fill}%
\end{pgfscope}%
\begin{pgfscope}%
\pgfpathrectangle{\pgfqpoint{0.407592in}{0.222904in}}{\pgfqpoint{1.411163in}{1.411163in}}%
\pgfusepath{clip}%
\pgfsetbuttcap%
\pgfsetroundjoin%
\definecolor{currentfill}{rgb}{0.749020,0.772549,0.792157}%
\pgfsetfillcolor{currentfill}%
\pgfsetfillopacity{0.160000}%
\pgfsetlinewidth{0.000000pt}%
\definecolor{currentstroke}{rgb}{0.000000,0.000000,0.000000}%
\pgfsetstrokecolor{currentstroke}%
\pgfsetdash{}{0pt}%
\pgfpathmoveto{\pgfqpoint{0.461173in}{1.109555in}}%
\pgfpathlineto{\pgfqpoint{0.446900in}{1.120878in}}%
\pgfpathlineto{\pgfqpoint{0.432600in}{1.132222in}}%
\pgfpathlineto{\pgfqpoint{0.435411in}{1.116030in}}%
\pgfpathlineto{\pgfqpoint{0.441176in}{1.099909in}}%
\pgfpathlineto{\pgfqpoint{0.455436in}{1.089158in}}%
\pgfpathlineto{\pgfqpoint{0.469662in}{1.078431in}}%
\pgfpathlineto{\pgfqpoint{0.464007in}{1.093956in}}%
\pgfpathlineto{\pgfqpoint{0.461173in}{1.109555in}}%
\pgfpathclose%
\pgfusepath{fill}%
\end{pgfscope}%
\begin{pgfscope}%
\pgfpathrectangle{\pgfqpoint{0.407592in}{0.222904in}}{\pgfqpoint{1.411163in}{1.411163in}}%
\pgfusepath{clip}%
\pgfsetbuttcap%
\pgfsetroundjoin%
\definecolor{currentfill}{rgb}{0.749020,0.772549,0.792157}%
\pgfsetfillcolor{currentfill}%
\pgfsetfillopacity{0.160000}%
\pgfsetlinewidth{0.000000pt}%
\definecolor{currentstroke}{rgb}{0.000000,0.000000,0.000000}%
\pgfsetstrokecolor{currentstroke}%
\pgfsetdash{}{0pt}%
\pgfpathmoveto{\pgfqpoint{0.582297in}{0.993507in}}%
\pgfpathlineto{\pgfqpoint{0.568331in}{1.004037in}}%
\pgfpathlineto{\pgfqpoint{0.554333in}{1.014591in}}%
\pgfpathlineto{\pgfqpoint{0.562025in}{1.001050in}}%
\pgfpathlineto{\pgfqpoint{0.572151in}{0.987724in}}%
\pgfpathlineto{\pgfqpoint{0.585795in}{0.977790in}}%
\pgfpathlineto{\pgfqpoint{0.599402in}{0.967883in}}%
\pgfpathlineto{\pgfqpoint{0.589699in}{0.980589in}}%
\pgfpathlineto{\pgfqpoint{0.582297in}{0.993507in}}%
\pgfpathclose%
\pgfusepath{fill}%
\end{pgfscope}%
\begin{pgfscope}%
\pgfpathrectangle{\pgfqpoint{0.407592in}{0.222904in}}{\pgfqpoint{1.411163in}{1.411163in}}%
\pgfusepath{clip}%
\pgfsetbuttcap%
\pgfsetroundjoin%
\definecolor{currentfill}{rgb}{0.749020,0.772549,0.792157}%
\pgfsetfillcolor{currentfill}%
\pgfsetfillopacity{0.160000}%
\pgfsetlinewidth{0.000000pt}%
\definecolor{currentstroke}{rgb}{0.000000,0.000000,0.000000}%
\pgfsetstrokecolor{currentstroke}%
\pgfsetdash{}{0pt}%
\pgfpathmoveto{\pgfqpoint{1.016319in}{0.728905in}}%
\pgfpathlineto{\pgfqpoint{1.012443in}{0.735113in}}%
\pgfpathlineto{\pgfqpoint{1.008549in}{0.741349in}}%
\pgfpathlineto{\pgfqpoint{1.028100in}{0.739723in}}%
\pgfpathlineto{\pgfqpoint{1.047946in}{0.738652in}}%
\pgfpathlineto{\pgfqpoint{1.049963in}{0.732539in}}%
\pgfpathlineto{\pgfqpoint{1.051971in}{0.726455in}}%
\pgfpathlineto{\pgfqpoint{1.034013in}{0.727427in}}%
\pgfpathlineto{\pgfqpoint{1.016319in}{0.728905in}}%
\pgfpathclose%
\pgfusepath{fill}%
\end{pgfscope}%
\begin{pgfscope}%
\pgfpathrectangle{\pgfqpoint{0.407592in}{0.222904in}}{\pgfqpoint{1.411163in}{1.411163in}}%
\pgfusepath{clip}%
\pgfsetbuttcap%
\pgfsetroundjoin%
\definecolor{currentfill}{rgb}{0.749020,0.772549,0.792157}%
\pgfsetfillcolor{currentfill}%
\pgfsetfillopacity{0.160000}%
\pgfsetlinewidth{0.000000pt}%
\definecolor{currentstroke}{rgb}{0.000000,0.000000,0.000000}%
\pgfsetstrokecolor{currentstroke}%
\pgfsetdash{}{0pt}%
\pgfpathmoveto{\pgfqpoint{1.419417in}{0.862526in}}%
\pgfpathlineto{\pgfqpoint{1.430949in}{0.871353in}}%
\pgfpathlineto{\pgfqpoint{1.442518in}{0.880209in}}%
\pgfpathlineto{\pgfqpoint{1.458363in}{0.890204in}}%
\pgfpathlineto{\pgfqpoint{1.446309in}{0.881046in}}%
\pgfpathlineto{\pgfqpoint{1.434292in}{0.871917in}}%
\pgfusepath{fill}%
\end{pgfscope}%
\begin{pgfscope}%
\pgfpathrectangle{\pgfqpoint{0.407592in}{0.222904in}}{\pgfqpoint{1.411163in}{1.411163in}}%
\pgfusepath{clip}%
\pgfsetbuttcap%
\pgfsetroundjoin%
\definecolor{currentfill}{rgb}{0.749020,0.772549,0.792157}%
\pgfsetfillcolor{currentfill}%
\pgfsetfillopacity{0.160000}%
\pgfsetlinewidth{0.000000pt}%
\definecolor{currentstroke}{rgb}{0.000000,0.000000,0.000000}%
\pgfsetstrokecolor{currentstroke}%
\pgfsetdash{}{0pt}%
\pgfpathmoveto{\pgfqpoint{0.802460in}{0.818128in}}%
\pgfpathlineto{\pgfqpoint{0.790871in}{0.826488in}}%
\pgfpathlineto{\pgfqpoint{0.779242in}{0.834877in}}%
\pgfpathlineto{\pgfqpoint{0.794484in}{0.826919in}}%
\pgfpathlineto{\pgfqpoint{0.811131in}{0.819403in}}%
\pgfpathlineto{\pgfqpoint{0.821615in}{0.811558in}}%
\pgfpathlineto{\pgfqpoint{0.832058in}{0.803742in}}%
\pgfpathlineto{\pgfqpoint{0.816613in}{0.810728in}}%
\pgfpathlineto{\pgfqpoint{0.802460in}{0.818128in}}%
\pgfpathclose%
\pgfusepath{fill}%
\end{pgfscope}%
\begin{pgfscope}%
\pgfpathrectangle{\pgfqpoint{0.407592in}{0.222904in}}{\pgfqpoint{1.411163in}{1.411163in}}%
\pgfusepath{clip}%
\pgfsetbuttcap%
\pgfsetroundjoin%
\definecolor{currentfill}{rgb}{0.749020,0.772549,0.792157}%
\pgfsetfillcolor{currentfill}%
\pgfsetfillopacity{0.160000}%
\pgfsetlinewidth{0.000000pt}%
\definecolor{currentstroke}{rgb}{0.000000,0.000000,0.000000}%
\pgfsetstrokecolor{currentstroke}%
\pgfsetdash{}{0pt}%
\pgfpathmoveto{\pgfqpoint{1.088334in}{0.726045in}}%
\pgfpathlineto{\pgfqpoint{1.088235in}{0.732109in}}%
\pgfpathlineto{\pgfqpoint{1.088135in}{0.738201in}}%
\pgfpathlineto{\pgfqpoint{1.108283in}{0.738822in}}%
\pgfpathlineto{\pgfqpoint{1.128335in}{0.740005in}}%
\pgfpathlineto{\pgfqpoint{1.126518in}{0.733830in}}%
\pgfpathlineto{\pgfqpoint{1.124709in}{0.727683in}}%
\pgfpathlineto{\pgfqpoint{1.106564in}{0.726610in}}%
\pgfpathlineto{\pgfqpoint{1.088334in}{0.726045in}}%
\pgfpathclose%
\pgfusepath{fill}%
\end{pgfscope}%
\begin{pgfscope}%
\pgfpathrectangle{\pgfqpoint{0.407592in}{0.222904in}}{\pgfqpoint{1.411163in}{1.411163in}}%
\pgfusepath{clip}%
\pgfsetbuttcap%
\pgfsetroundjoin%
\definecolor{currentfill}{rgb}{0.749020,0.772549,0.792157}%
\pgfsetfillcolor{currentfill}%
\pgfsetfillopacity{0.160000}%
\pgfsetlinewidth{0.000000pt}%
\definecolor{currentstroke}{rgb}{0.000000,0.000000,0.000000}%
\pgfsetstrokecolor{currentstroke}%
\pgfsetdash{}{0pt}%
\pgfpathmoveto{\pgfqpoint{0.852825in}{0.788200in}}%
\pgfpathlineto{\pgfqpoint{0.842461in}{0.795956in}}%
\pgfpathlineto{\pgfqpoint{0.832058in}{0.803742in}}%
\pgfpathlineto{\pgfqpoint{0.848729in}{0.797205in}}%
\pgfpathlineto{\pgfqpoint{0.866553in}{0.791152in}}%
\pgfpathlineto{\pgfqpoint{0.875599in}{0.783848in}}%
\pgfpathlineto{\pgfqpoint{0.884607in}{0.776574in}}%
\pgfpathlineto{\pgfqpoint{0.868190in}{0.782163in}}%
\pgfpathlineto{\pgfqpoint{0.852825in}{0.788200in}}%
\pgfpathclose%
\pgfusepath{fill}%
\end{pgfscope}%
\begin{pgfscope}%
\pgfpathrectangle{\pgfqpoint{0.407592in}{0.222904in}}{\pgfqpoint{1.411163in}{1.411163in}}%
\pgfusepath{clip}%
\pgfsetbuttcap%
\pgfsetroundjoin%
\definecolor{currentfill}{rgb}{0.749020,0.772549,0.792157}%
\pgfsetfillcolor{currentfill}%
\pgfsetfillopacity{0.160000}%
\pgfsetlinewidth{0.000000pt}%
\definecolor{currentstroke}{rgb}{0.000000,0.000000,0.000000}%
\pgfsetstrokecolor{currentstroke}%
\pgfsetdash{}{0pt}%
\pgfpathmoveto{\pgfqpoint{1.205671in}{0.750178in}}%
\pgfpathlineto{\pgfqpoint{1.211186in}{0.756849in}}%
\pgfpathlineto{\pgfqpoint{1.216725in}{0.763549in}}%
\pgfpathlineto{\pgfqpoint{1.236599in}{0.767736in}}%
\pgfpathlineto{\pgfqpoint{1.255757in}{0.772447in}}%
\pgfpathlineto{\pgfqpoint{1.248526in}{0.765375in}}%
\pgfpathlineto{\pgfqpoint{1.241326in}{0.758332in}}%
\pgfpathlineto{\pgfqpoint{1.223823in}{0.754014in}}%
\pgfpathlineto{\pgfqpoint{1.205671in}{0.750178in}}%
\pgfpathclose%
\pgfusepath{fill}%
\end{pgfscope}%
\begin{pgfscope}%
\pgfpathrectangle{\pgfqpoint{0.407592in}{0.222904in}}{\pgfqpoint{1.411163in}{1.411163in}}%
\pgfusepath{clip}%
\pgfsetbuttcap%
\pgfsetroundjoin%
\definecolor{currentfill}{rgb}{0.749020,0.772549,0.792157}%
\pgfsetfillcolor{currentfill}%
\pgfsetfillopacity{0.160000}%
\pgfsetlinewidth{0.000000pt}%
\definecolor{currentstroke}{rgb}{0.000000,0.000000,0.000000}%
\pgfsetstrokecolor{currentstroke}%
\pgfsetdash{}{0pt}%
\pgfpathmoveto{\pgfqpoint{1.051971in}{0.726455in}}%
\pgfpathlineto{\pgfqpoint{1.049963in}{0.732539in}}%
\pgfpathlineto{\pgfqpoint{1.047946in}{0.738652in}}%
\pgfpathlineto{\pgfqpoint{1.067990in}{0.738144in}}%
\pgfpathlineto{\pgfqpoint{1.088135in}{0.738201in}}%
\pgfpathlineto{\pgfqpoint{1.088235in}{0.732109in}}%
\pgfpathlineto{\pgfqpoint{1.088334in}{0.726045in}}%
\pgfpathlineto{\pgfqpoint{1.070107in}{0.725994in}}%
\pgfpathlineto{\pgfqpoint{1.051971in}{0.726455in}}%
\pgfpathclose%
\pgfusepath{fill}%
\end{pgfscope}%
\begin{pgfscope}%
\pgfpathrectangle{\pgfqpoint{0.407592in}{0.222904in}}{\pgfqpoint{1.411163in}{1.411163in}}%
\pgfusepath{clip}%
\pgfsetbuttcap%
\pgfsetroundjoin%
\definecolor{currentfill}{rgb}{0.749020,0.772549,0.792157}%
\pgfsetfillcolor{currentfill}%
\pgfsetfillopacity{0.160000}%
\pgfsetlinewidth{0.000000pt}%
\definecolor{currentstroke}{rgb}{0.000000,0.000000,0.000000}%
\pgfsetstrokecolor{currentstroke}%
\pgfsetdash{}{0pt}%
\pgfpathmoveto{\pgfqpoint{1.364368in}{0.828482in}}%
\pgfpathlineto{\pgfqpoint{1.374675in}{0.836677in}}%
\pgfpathlineto{\pgfqpoint{1.385021in}{0.844902in}}%
\pgfpathlineto{\pgfqpoint{1.402972in}{0.853506in}}%
\pgfpathlineto{\pgfqpoint{1.419417in}{0.862526in}}%
\pgfpathlineto{\pgfqpoint{1.407924in}{0.853728in}}%
\pgfpathlineto{\pgfqpoint{1.396468in}{0.844958in}}%
\pgfpathlineto{\pgfqpoint{1.381114in}{0.836524in}}%
\pgfpathlineto{\pgfqpoint{1.364368in}{0.828482in}}%
\pgfpathclose%
\pgfusepath{fill}%
\end{pgfscope}%
\begin{pgfscope}%
\pgfpathrectangle{\pgfqpoint{0.407592in}{0.222904in}}{\pgfqpoint{1.411163in}{1.411163in}}%
\pgfusepath{clip}%
\pgfsetbuttcap%
\pgfsetroundjoin%
\definecolor{currentfill}{rgb}{0.749020,0.772549,0.792157}%
\pgfsetfillcolor{currentfill}%
\pgfsetfillopacity{0.160000}%
\pgfsetlinewidth{0.000000pt}%
\definecolor{currentstroke}{rgb}{0.000000,0.000000,0.000000}%
\pgfsetstrokecolor{currentstroke}%
\pgfsetdash{}{0pt}%
\pgfpathmoveto{\pgfqpoint{1.721866in}{1.103134in}}%
\pgfpathlineto{\pgfqpoint{1.735760in}{1.113970in}}%
\pgfpathlineto{\pgfqpoint{1.749686in}{1.124831in}}%
\pgfpathlineto{\pgfqpoint{1.758991in}{1.141560in}}%
\pgfpathlineto{\pgfqpoint{1.765228in}{1.158340in}}%
\pgfpathlineto{\pgfqpoint{1.751127in}{1.146896in}}%
\pgfpathlineto{\pgfqpoint{1.737051in}{1.135472in}}%
\pgfpathlineto{\pgfqpoint{1.730926in}{1.119275in}}%
\pgfpathlineto{\pgfqpoint{1.721866in}{1.103134in}}%
\pgfpathclose%
\pgfusepath{fill}%
\end{pgfscope}%
\begin{pgfscope}%
\pgfpathrectangle{\pgfqpoint{0.407592in}{0.222904in}}{\pgfqpoint{1.411163in}{1.411163in}}%
\pgfusepath{clip}%
\pgfsetbuttcap%
\pgfsetroundjoin%
\definecolor{currentfill}{rgb}{0.749020,0.772549,0.792157}%
\pgfsetfillcolor{currentfill}%
\pgfsetfillopacity{0.160000}%
\pgfsetlinewidth{0.000000pt}%
\definecolor{currentstroke}{rgb}{0.000000,0.000000,0.000000}%
\pgfsetstrokecolor{currentstroke}%
\pgfsetdash{}{0pt}%
\pgfpathmoveto{\pgfqpoint{1.482584in}{0.908604in}}%
\pgfpathlineto{\pgfqpoint{1.494750in}{0.917847in}}%
\pgfpathlineto{\pgfqpoint{1.506955in}{0.927119in}}%
\pgfpathlineto{\pgfqpoint{1.522730in}{0.938681in}}%
\pgfpathlineto{\pgfqpoint{1.536450in}{0.950555in}}%
\pgfpathlineto{\pgfqpoint{1.523473in}{0.940667in}}%
\pgfpathlineto{\pgfqpoint{1.510531in}{0.930804in}}%
\pgfpathlineto{\pgfqpoint{1.497521in}{0.919554in}}%
\pgfpathlineto{\pgfqpoint{1.482584in}{0.908604in}}%
\pgfpathclose%
\pgfusepath{fill}%
\end{pgfscope}%
\begin{pgfscope}%
\pgfpathrectangle{\pgfqpoint{0.407592in}{0.222904in}}{\pgfqpoint{1.411163in}{1.411163in}}%
\pgfusepath{clip}%
\pgfsetbuttcap%
\pgfsetroundjoin%
\definecolor{currentfill}{rgb}{0.749020,0.772549,0.792157}%
\pgfsetfillcolor{currentfill}%
\pgfsetfillopacity{0.160000}%
\pgfsetlinewidth{0.000000pt}%
\definecolor{currentstroke}{rgb}{0.000000,0.000000,0.000000}%
\pgfsetstrokecolor{currentstroke}%
\pgfsetdash{}{0pt}%
\pgfpathmoveto{\pgfqpoint{1.588709in}{0.990377in}}%
\pgfpathlineto{\pgfqpoint{1.601862in}{1.000399in}}%
\pgfpathlineto{\pgfqpoint{1.615051in}{1.010449in}}%
\pgfpathlineto{\pgfqpoint{1.628367in}{1.024439in}}%
\pgfpathlineto{\pgfqpoint{1.639161in}{1.038634in}}%
\pgfpathlineto{\pgfqpoint{1.625485in}{1.027969in}}%
\pgfpathlineto{\pgfqpoint{1.611841in}{1.017328in}}%
\pgfpathlineto{\pgfqpoint{1.601470in}{1.003751in}}%
\pgfpathlineto{\pgfqpoint{1.588709in}{0.990377in}}%
\pgfpathclose%
\pgfusepath{fill}%
\end{pgfscope}%
\begin{pgfscope}%
\pgfpathrectangle{\pgfqpoint{0.407592in}{0.222904in}}{\pgfqpoint{1.411163in}{1.411163in}}%
\pgfusepath{clip}%
\pgfsetbuttcap%
\pgfsetroundjoin%
\definecolor{currentfill}{rgb}{0.749020,0.772549,0.792157}%
\pgfsetfillcolor{currentfill}%
\pgfsetfillopacity{0.160000}%
\pgfsetlinewidth{0.000000pt}%
\definecolor{currentstroke}{rgb}{0.000000,0.000000,0.000000}%
\pgfsetstrokecolor{currentstroke}%
\pgfsetdash{}{0pt}%
\pgfpathmoveto{\pgfqpoint{1.255757in}{0.772447in}}%
\pgfpathlineto{\pgfqpoint{1.263018in}{0.779550in}}%
\pgfpathlineto{\pgfqpoint{1.270310in}{0.786683in}}%
\pgfpathlineto{\pgfqpoint{1.290248in}{0.792326in}}%
\pgfpathlineto{\pgfqpoint{1.309209in}{0.798479in}}%
\pgfpathlineto{\pgfqpoint{1.300372in}{0.790894in}}%
\pgfpathlineto{\pgfqpoint{1.291569in}{0.783339in}}%
\pgfpathlineto{\pgfqpoint{1.274108in}{0.777657in}}%
\pgfpathlineto{\pgfqpoint{1.255757in}{0.772447in}}%
\pgfpathclose%
\pgfusepath{fill}%
\end{pgfscope}%
\begin{pgfscope}%
\pgfpathrectangle{\pgfqpoint{0.407592in}{0.222904in}}{\pgfqpoint{1.411163in}{1.411163in}}%
\pgfusepath{clip}%
\pgfsetbuttcap%
\pgfsetroundjoin%
\definecolor{currentfill}{rgb}{0.749020,0.772549,0.792157}%
\pgfsetfillcolor{currentfill}%
\pgfsetfillopacity{0.160000}%
\pgfsetlinewidth{0.000000pt}%
\definecolor{currentstroke}{rgb}{0.000000,0.000000,0.000000}%
\pgfsetstrokecolor{currentstroke}%
\pgfsetdash{}{0pt}%
\pgfpathmoveto{\pgfqpoint{0.935136in}{0.753205in}}%
\pgfpathlineto{\pgfqpoint{0.927734in}{0.760015in}}%
\pgfpathlineto{\pgfqpoint{0.920300in}{0.766854in}}%
\pgfpathlineto{\pgfqpoint{0.939410in}{0.762774in}}%
\pgfpathlineto{\pgfqpoint{0.959245in}{0.759245in}}%
\pgfpathlineto{\pgfqpoint{0.964987in}{0.752725in}}%
\pgfpathlineto{\pgfqpoint{0.970704in}{0.746235in}}%
\pgfpathlineto{\pgfqpoint{0.952592in}{0.749468in}}%
\pgfpathlineto{\pgfqpoint{0.935136in}{0.753205in}}%
\pgfpathclose%
\pgfusepath{fill}%
\end{pgfscope}%
\begin{pgfscope}%
\pgfpathrectangle{\pgfqpoint{0.407592in}{0.222904in}}{\pgfqpoint{1.411163in}{1.411163in}}%
\pgfusepath{clip}%
\pgfsetbuttcap%
\pgfsetroundjoin%
\definecolor{currentfill}{rgb}{0.749020,0.772549,0.792157}%
\pgfsetfillcolor{currentfill}%
\pgfsetfillopacity{0.160000}%
\pgfsetlinewidth{0.000000pt}%
\definecolor{currentstroke}{rgb}{0.000000,0.000000,0.000000}%
\pgfsetstrokecolor{currentstroke}%
\pgfsetdash{}{0pt}%
\pgfpathmoveto{\pgfqpoint{1.309209in}{0.798479in}}%
\pgfpathlineto{\pgfqpoint{1.318081in}{0.806094in}}%
\pgfpathlineto{\pgfqpoint{1.326989in}{0.813739in}}%
\pgfpathlineto{\pgfqpoint{1.346300in}{0.820874in}}%
\pgfpathlineto{\pgfqpoint{1.364368in}{0.828482in}}%
\pgfpathlineto{\pgfqpoint{1.354097in}{0.820316in}}%
\pgfpathlineto{\pgfqpoint{1.343864in}{0.812181in}}%
\pgfpathlineto{\pgfqpoint{1.327108in}{0.805109in}}%
\pgfpathlineto{\pgfqpoint{1.309209in}{0.798479in}}%
\pgfpathclose%
\pgfusepath{fill}%
\end{pgfscope}%
\begin{pgfscope}%
\pgfpathrectangle{\pgfqpoint{0.407592in}{0.222904in}}{\pgfqpoint{1.411163in}{1.411163in}}%
\pgfusepath{clip}%
\pgfsetbuttcap%
\pgfsetroundjoin%
\definecolor{currentfill}{rgb}{0.749020,0.772549,0.792157}%
\pgfsetfillcolor{currentfill}%
\pgfsetfillopacity{0.160000}%
\pgfsetlinewidth{0.000000pt}%
\definecolor{currentstroke}{rgb}{0.000000,0.000000,0.000000}%
\pgfsetstrokecolor{currentstroke}%
\pgfsetdash{}{0pt}%
\pgfpathmoveto{\pgfqpoint{0.432600in}{1.132222in}}%
\pgfpathlineto{\pgfqpoint{0.418274in}{1.143587in}}%
\pgfpathlineto{\pgfqpoint{0.403920in}{1.154973in}}%
\pgfpathlineto{\pgfqpoint{0.406696in}{1.138197in}}%
\pgfpathlineto{\pgfqpoint{0.412558in}{1.121487in}}%
\pgfpathlineto{\pgfqpoint{0.426884in}{1.110686in}}%
\pgfpathlineto{\pgfqpoint{0.441176in}{1.099909in}}%
\pgfpathlineto{\pgfqpoint{0.435411in}{1.116030in}}%
\pgfpathlineto{\pgfqpoint{0.432600in}{1.132222in}}%
\pgfpathclose%
\pgfusepath{fill}%
\end{pgfscope}%
\begin{pgfscope}%
\pgfpathrectangle{\pgfqpoint{0.407592in}{0.222904in}}{\pgfqpoint{1.411163in}{1.411163in}}%
\pgfusepath{clip}%
\pgfsetbuttcap%
\pgfsetroundjoin%
\definecolor{currentfill}{rgb}{0.749020,0.772549,0.792157}%
\pgfsetfillcolor{currentfill}%
\pgfsetfillopacity{0.160000}%
\pgfsetlinewidth{0.000000pt}%
\definecolor{currentstroke}{rgb}{0.000000,0.000000,0.000000}%
\pgfsetstrokecolor{currentstroke}%
\pgfsetdash{}{0pt}%
\pgfpathmoveto{\pgfqpoint{0.653460in}{0.928526in}}%
\pgfpathlineto{\pgfqpoint{0.640001in}{0.938325in}}%
\pgfpathlineto{\pgfqpoint{0.626505in}{0.948151in}}%
\pgfpathlineto{\pgfqpoint{0.637925in}{0.936335in}}%
\pgfpathlineto{\pgfqpoint{0.651466in}{0.924842in}}%
\pgfpathlineto{\pgfqpoint{0.664324in}{0.915631in}}%
\pgfpathlineto{\pgfqpoint{0.677143in}{0.906448in}}%
\pgfpathlineto{\pgfqpoint{0.664307in}{0.917332in}}%
\pgfpathlineto{\pgfqpoint{0.653460in}{0.928526in}}%
\pgfpathclose%
\pgfusepath{fill}%
\end{pgfscope}%
\begin{pgfscope}%
\pgfpathrectangle{\pgfqpoint{0.407592in}{0.222904in}}{\pgfqpoint{1.411163in}{1.411163in}}%
\pgfusepath{clip}%
\pgfsetbuttcap%
\pgfsetroundjoin%
\definecolor{currentfill}{rgb}{0.749020,0.772549,0.792157}%
\pgfsetfillcolor{currentfill}%
\pgfsetfillopacity{0.160000}%
\pgfsetlinewidth{0.000000pt}%
\definecolor{currentstroke}{rgb}{0.000000,0.000000,0.000000}%
\pgfsetstrokecolor{currentstroke}%
\pgfsetdash{}{0pt}%
\pgfpathmoveto{\pgfqpoint{0.728016in}{0.870005in}}%
\pgfpathlineto{\pgfqpoint{0.715357in}{0.879073in}}%
\pgfpathlineto{\pgfqpoint{0.702659in}{0.888170in}}%
\pgfpathlineto{\pgfqpoint{0.716603in}{0.878250in}}%
\pgfpathlineto{\pgfqpoint{0.732315in}{0.868730in}}%
\pgfpathlineto{\pgfqpoint{0.744108in}{0.860222in}}%
\pgfpathlineto{\pgfqpoint{0.755860in}{0.851745in}}%
\pgfpathlineto{\pgfqpoint{0.741116in}{0.860685in}}%
\pgfpathlineto{\pgfqpoint{0.728016in}{0.870005in}}%
\pgfpathclose%
\pgfusepath{fill}%
\end{pgfscope}%
\begin{pgfscope}%
\pgfpathrectangle{\pgfqpoint{0.407592in}{0.222904in}}{\pgfqpoint{1.411163in}{1.411163in}}%
\pgfusepath{clip}%
\pgfsetbuttcap%
\pgfsetroundjoin%
\definecolor{currentfill}{rgb}{0.749020,0.772549,0.792157}%
\pgfsetfillcolor{currentfill}%
\pgfsetfillopacity{0.160000}%
\pgfsetlinewidth{0.000000pt}%
\definecolor{currentstroke}{rgb}{0.000000,0.000000,0.000000}%
\pgfsetstrokecolor{currentstroke}%
\pgfsetdash{}{0pt}%
\pgfpathmoveto{\pgfqpoint{0.554333in}{1.014591in}}%
\pgfpathlineto{\pgfqpoint{0.540302in}{1.025170in}}%
\pgfpathlineto{\pgfqpoint{0.526239in}{1.035773in}}%
\pgfpathlineto{\pgfqpoint{0.534210in}{1.021614in}}%
\pgfpathlineto{\pgfqpoint{0.544749in}{1.007673in}}%
\pgfpathlineto{\pgfqpoint{0.558469in}{0.997685in}}%
\pgfpathlineto{\pgfqpoint{0.572151in}{0.987724in}}%
\pgfpathlineto{\pgfqpoint{0.562025in}{1.001050in}}%
\pgfpathlineto{\pgfqpoint{0.554333in}{1.014591in}}%
\pgfpathclose%
\pgfusepath{fill}%
\end{pgfscope}%
\begin{pgfscope}%
\pgfpathrectangle{\pgfqpoint{0.407592in}{0.222904in}}{\pgfqpoint{1.411163in}{1.411163in}}%
\pgfusepath{clip}%
\pgfsetbuttcap%
\pgfsetroundjoin%
\definecolor{currentfill}{rgb}{0.749020,0.772549,0.792157}%
\pgfsetfillcolor{currentfill}%
\pgfsetfillopacity{0.160000}%
\pgfsetlinewidth{0.000000pt}%
\definecolor{currentstroke}{rgb}{0.000000,0.000000,0.000000}%
\pgfsetstrokecolor{currentstroke}%
\pgfsetdash{}{0pt}%
\pgfpathmoveto{\pgfqpoint{1.167767in}{0.744025in}}%
\pgfpathlineto{\pgfqpoint{1.171478in}{0.750415in}}%
\pgfpathlineto{\pgfqpoint{1.175206in}{0.756833in}}%
\pgfpathlineto{\pgfqpoint{1.196228in}{0.759908in}}%
\pgfpathlineto{\pgfqpoint{1.216725in}{0.763549in}}%
\pgfpathlineto{\pgfqpoint{1.211186in}{0.756849in}}%
\pgfpathlineto{\pgfqpoint{1.205671in}{0.750178in}}%
\pgfpathlineto{\pgfqpoint{1.186956in}{0.746842in}}%
\pgfpathlineto{\pgfqpoint{1.167767in}{0.744025in}}%
\pgfpathclose%
\pgfusepath{fill}%
\end{pgfscope}%
\begin{pgfscope}%
\pgfpathrectangle{\pgfqpoint{0.407592in}{0.222904in}}{\pgfqpoint{1.411163in}{1.411163in}}%
\pgfusepath{clip}%
\pgfsetbuttcap%
\pgfsetroundjoin%
\definecolor{currentfill}{rgb}{0.749020,0.772549,0.792157}%
\pgfsetfillcolor{currentfill}%
\pgfsetfillopacity{0.160000}%
\pgfsetlinewidth{0.000000pt}%
\definecolor{currentstroke}{rgb}{0.000000,0.000000,0.000000}%
\pgfsetstrokecolor{currentstroke}%
\pgfsetdash{}{0pt}%
\pgfpathmoveto{\pgfqpoint{1.749686in}{1.124831in}}%
\pgfpathlineto{\pgfqpoint{1.763644in}{1.135716in}}%
\pgfpathlineto{\pgfqpoint{1.777633in}{1.146626in}}%
\pgfpathlineto{\pgfqpoint{1.787171in}{1.163937in}}%
\pgfpathlineto{\pgfqpoint{1.793508in}{1.181291in}}%
\pgfpathlineto{\pgfqpoint{1.779355in}{1.169805in}}%
\pgfpathlineto{\pgfqpoint{1.765228in}{1.158340in}}%
\pgfpathlineto{\pgfqpoint{1.758991in}{1.141560in}}%
\pgfpathlineto{\pgfqpoint{1.749686in}{1.124831in}}%
\pgfpathclose%
\pgfusepath{fill}%
\end{pgfscope}%
\begin{pgfscope}%
\pgfpathrectangle{\pgfqpoint{0.407592in}{0.222904in}}{\pgfqpoint{1.411163in}{1.411163in}}%
\pgfusepath{clip}%
\pgfsetbuttcap%
\pgfsetroundjoin%
\definecolor{currentfill}{rgb}{0.749020,0.772549,0.792157}%
\pgfsetfillcolor{currentfill}%
\pgfsetfillopacity{0.160000}%
\pgfsetlinewidth{0.000000pt}%
\definecolor{currentstroke}{rgb}{0.000000,0.000000,0.000000}%
\pgfsetstrokecolor{currentstroke}%
\pgfsetdash{}{0pt}%
\pgfpathmoveto{\pgfqpoint{1.442518in}{0.880209in}}%
\pgfpathlineto{\pgfqpoint{1.454126in}{0.889095in}}%
\pgfpathlineto{\pgfqpoint{1.465771in}{0.898010in}}%
\pgfpathlineto{\pgfqpoint{1.482584in}{0.908604in}}%
\pgfpathlineto{\pgfqpoint{1.470455in}{0.899390in}}%
\pgfpathlineto{\pgfqpoint{1.458363in}{0.890204in}}%
\pgfusepath{fill}%
\end{pgfscope}%
\begin{pgfscope}%
\pgfpathrectangle{\pgfqpoint{0.407592in}{0.222904in}}{\pgfqpoint{1.411163in}{1.411163in}}%
\pgfusepath{clip}%
\pgfsetbuttcap%
\pgfsetroundjoin%
\definecolor{currentfill}{rgb}{0.749020,0.772549,0.792157}%
\pgfsetfillcolor{currentfill}%
\pgfsetfillopacity{0.160000}%
\pgfsetlinewidth{0.000000pt}%
\definecolor{currentstroke}{rgb}{0.000000,0.000000,0.000000}%
\pgfsetstrokecolor{currentstroke}%
\pgfsetdash{}{0pt}%
\pgfpathmoveto{\pgfqpoint{0.970704in}{0.746235in}}%
\pgfpathlineto{\pgfqpoint{0.964987in}{0.752725in}}%
\pgfpathlineto{\pgfqpoint{0.959245in}{0.759245in}}%
\pgfpathlineto{\pgfqpoint{0.979709in}{0.756286in}}%
\pgfpathlineto{\pgfqpoint{1.000705in}{0.753911in}}%
\pgfpathlineto{\pgfqpoint{1.004636in}{0.747616in}}%
\pgfpathlineto{\pgfqpoint{1.008549in}{0.741349in}}%
\pgfpathlineto{\pgfqpoint{0.989386in}{0.743524in}}%
\pgfpathlineto{\pgfqpoint{0.970704in}{0.746235in}}%
\pgfpathclose%
\pgfusepath{fill}%
\end{pgfscope}%
\begin{pgfscope}%
\pgfpathrectangle{\pgfqpoint{0.407592in}{0.222904in}}{\pgfqpoint{1.411163in}{1.411163in}}%
\pgfusepath{clip}%
\pgfsetbuttcap%
\pgfsetroundjoin%
\definecolor{currentfill}{rgb}{0.749020,0.772549,0.792157}%
\pgfsetfillcolor{currentfill}%
\pgfsetfillopacity{0.160000}%
\pgfsetlinewidth{0.000000pt}%
\definecolor{currentstroke}{rgb}{0.000000,0.000000,0.000000}%
\pgfsetstrokecolor{currentstroke}%
\pgfsetdash{}{0pt}%
\pgfpathmoveto{\pgfqpoint{0.884607in}{0.776574in}}%
\pgfpathlineto{\pgfqpoint{0.875599in}{0.783848in}}%
\pgfpathlineto{\pgfqpoint{0.866553in}{0.791152in}}%
\pgfpathlineto{\pgfqpoint{0.885450in}{0.785614in}}%
\pgfpathlineto{\pgfqpoint{0.905333in}{0.780622in}}%
\pgfpathlineto{\pgfqpoint{0.912833in}{0.773723in}}%
\pgfpathlineto{\pgfqpoint{0.920300in}{0.766854in}}%
\pgfpathlineto{\pgfqpoint{0.902004in}{0.771462in}}%
\pgfpathlineto{\pgfqpoint{0.884607in}{0.776574in}}%
\pgfpathclose%
\pgfusepath{fill}%
\end{pgfscope}%
\begin{pgfscope}%
\pgfpathrectangle{\pgfqpoint{0.407592in}{0.222904in}}{\pgfqpoint{1.411163in}{1.411163in}}%
\pgfusepath{clip}%
\pgfsetbuttcap%
\pgfsetroundjoin%
\definecolor{currentfill}{rgb}{0.749020,0.772549,0.792157}%
\pgfsetfillcolor{currentfill}%
\pgfsetfillopacity{0.160000}%
\pgfsetlinewidth{0.000000pt}%
\definecolor{currentstroke}{rgb}{0.000000,0.000000,0.000000}%
\pgfsetstrokecolor{currentstroke}%
\pgfsetdash{}{0pt}%
\pgfpathmoveto{\pgfqpoint{0.779242in}{0.834877in}}%
\pgfpathlineto{\pgfqpoint{0.767572in}{0.843296in}}%
\pgfpathlineto{\pgfqpoint{0.755860in}{0.851745in}}%
\pgfpathlineto{\pgfqpoint{0.772192in}{0.843230in}}%
\pgfpathlineto{\pgfqpoint{0.790043in}{0.835185in}}%
\pgfpathlineto{\pgfqpoint{0.800608in}{0.827279in}}%
\pgfpathlineto{\pgfqpoint{0.811131in}{0.819403in}}%
\pgfpathlineto{\pgfqpoint{0.794484in}{0.826919in}}%
\pgfpathlineto{\pgfqpoint{0.779242in}{0.834877in}}%
\pgfpathclose%
\pgfusepath{fill}%
\end{pgfscope}%
\begin{pgfscope}%
\pgfpathrectangle{\pgfqpoint{0.407592in}{0.222904in}}{\pgfqpoint{1.411163in}{1.411163in}}%
\pgfusepath{clip}%
\pgfsetbuttcap%
\pgfsetroundjoin%
\definecolor{currentfill}{rgb}{0.749020,0.772549,0.792157}%
\pgfsetfillcolor{currentfill}%
\pgfsetfillopacity{0.160000}%
\pgfsetlinewidth{0.000000pt}%
\definecolor{currentstroke}{rgb}{0.000000,0.000000,0.000000}%
\pgfsetstrokecolor{currentstroke}%
\pgfsetdash{}{0pt}%
\pgfpathmoveto{\pgfqpoint{1.615051in}{1.010449in}}%
\pgfpathlineto{\pgfqpoint{1.628276in}{1.020526in}}%
\pgfpathlineto{\pgfqpoint{1.641536in}{1.030631in}}%
\pgfpathlineto{\pgfqpoint{1.655398in}{1.045230in}}%
\pgfpathlineto{\pgfqpoint{1.666604in}{1.060037in}}%
\pgfpathlineto{\pgfqpoint{1.652867in}{1.049323in}}%
\pgfpathlineto{\pgfqpoint{1.639161in}{1.038634in}}%
\pgfpathlineto{\pgfqpoint{1.628367in}{1.024439in}}%
\pgfpathlineto{\pgfqpoint{1.615051in}{1.010449in}}%
\pgfpathclose%
\pgfusepath{fill}%
\end{pgfscope}%
\begin{pgfscope}%
\pgfpathrectangle{\pgfqpoint{0.407592in}{0.222904in}}{\pgfqpoint{1.411163in}{1.411163in}}%
\pgfusepath{clip}%
\pgfsetbuttcap%
\pgfsetroundjoin%
\definecolor{currentfill}{rgb}{0.749020,0.772549,0.792157}%
\pgfsetfillcolor{currentfill}%
\pgfsetfillopacity{0.160000}%
\pgfsetlinewidth{0.000000pt}%
\definecolor{currentstroke}{rgb}{0.000000,0.000000,0.000000}%
\pgfsetstrokecolor{currentstroke}%
\pgfsetdash{}{0pt}%
\pgfpathmoveto{\pgfqpoint{0.403920in}{1.154973in}}%
\pgfpathlineto{\pgfqpoint{0.389540in}{1.166381in}}%
\pgfpathlineto{\pgfqpoint{0.375133in}{1.177810in}}%
\pgfpathlineto{\pgfqpoint{0.377860in}{1.160458in}}%
\pgfpathlineto{\pgfqpoint{0.383806in}{1.143166in}}%
\pgfpathlineto{\pgfqpoint{0.398199in}{1.132314in}}%
\pgfpathlineto{\pgfqpoint{0.412558in}{1.121487in}}%
\pgfpathlineto{\pgfqpoint{0.406696in}{1.138197in}}%
\pgfpathlineto{\pgfqpoint{0.403920in}{1.154973in}}%
\pgfpathclose%
\pgfusepath{fill}%
\end{pgfscope}%
\begin{pgfscope}%
\pgfpathrectangle{\pgfqpoint{0.407592in}{0.222904in}}{\pgfqpoint{1.411163in}{1.411163in}}%
\pgfusepath{clip}%
\pgfsetbuttcap%
\pgfsetroundjoin%
\definecolor{currentfill}{rgb}{0.749020,0.772549,0.792157}%
\pgfsetfillcolor{currentfill}%
\pgfsetfillopacity{0.160000}%
\pgfsetlinewidth{0.000000pt}%
\definecolor{currentstroke}{rgb}{0.000000,0.000000,0.000000}%
\pgfsetstrokecolor{currentstroke}%
\pgfsetdash{}{0pt}%
\pgfpathmoveto{\pgfqpoint{1.128335in}{0.740005in}}%
\pgfpathlineto{\pgfqpoint{1.130161in}{0.746209in}}%
\pgfpathlineto{\pgfqpoint{1.131996in}{0.752442in}}%
\pgfpathlineto{\pgfqpoint{1.153761in}{0.754340in}}%
\pgfpathlineto{\pgfqpoint{1.175206in}{0.756833in}}%
\pgfpathlineto{\pgfqpoint{1.171478in}{0.750415in}}%
\pgfpathlineto{\pgfqpoint{1.167767in}{0.744025in}}%
\pgfpathlineto{\pgfqpoint{1.148195in}{0.741742in}}%
\pgfpathlineto{\pgfqpoint{1.128335in}{0.740005in}}%
\pgfpathclose%
\pgfusepath{fill}%
\end{pgfscope}%
\begin{pgfscope}%
\pgfpathrectangle{\pgfqpoint{0.407592in}{0.222904in}}{\pgfqpoint{1.411163in}{1.411163in}}%
\pgfusepath{clip}%
\pgfsetbuttcap%
\pgfsetroundjoin%
\definecolor{currentfill}{rgb}{0.749020,0.772549,0.792157}%
\pgfsetfillcolor{currentfill}%
\pgfsetfillopacity{0.160000}%
\pgfsetlinewidth{0.000000pt}%
\definecolor{currentstroke}{rgb}{0.000000,0.000000,0.000000}%
\pgfsetstrokecolor{currentstroke}%
\pgfsetdash{}{0pt}%
\pgfpathmoveto{\pgfqpoint{1.506955in}{0.927119in}}%
\pgfpathlineto{\pgfqpoint{1.519197in}{0.936420in}}%
\pgfpathlineto{\pgfqpoint{1.531478in}{0.945749in}}%
\pgfpathlineto{\pgfqpoint{1.548086in}{0.957919in}}%
\pgfpathlineto{\pgfqpoint{1.562509in}{0.970412in}}%
\pgfpathlineto{\pgfqpoint{1.549462in}{0.960470in}}%
\pgfpathlineto{\pgfqpoint{1.536450in}{0.950555in}}%
\pgfpathlineto{\pgfqpoint{1.522730in}{0.938681in}}%
\pgfpathlineto{\pgfqpoint{1.506955in}{0.927119in}}%
\pgfpathclose%
\pgfusepath{fill}%
\end{pgfscope}%
\begin{pgfscope}%
\pgfpathrectangle{\pgfqpoint{0.407592in}{0.222904in}}{\pgfqpoint{1.411163in}{1.411163in}}%
\pgfusepath{clip}%
\pgfsetbuttcap%
\pgfsetroundjoin%
\definecolor{currentfill}{rgb}{0.749020,0.772549,0.792157}%
\pgfsetfillcolor{currentfill}%
\pgfsetfillopacity{0.160000}%
\pgfsetlinewidth{0.000000pt}%
\definecolor{currentstroke}{rgb}{0.000000,0.000000,0.000000}%
\pgfsetstrokecolor{currentstroke}%
\pgfsetdash{}{0pt}%
\pgfpathmoveto{\pgfqpoint{0.832058in}{0.803742in}}%
\pgfpathlineto{\pgfqpoint{0.821615in}{0.811558in}}%
\pgfpathlineto{\pgfqpoint{0.811131in}{0.819403in}}%
\pgfpathlineto{\pgfqpoint{0.829113in}{0.812368in}}%
\pgfpathlineto{\pgfqpoint{0.848350in}{0.805852in}}%
\pgfpathlineto{\pgfqpoint{0.857470in}{0.798487in}}%
\pgfpathlineto{\pgfqpoint{0.866553in}{0.791152in}}%
\pgfpathlineto{\pgfqpoint{0.848729in}{0.797205in}}%
\pgfpathlineto{\pgfqpoint{0.832058in}{0.803742in}}%
\pgfpathclose%
\pgfusepath{fill}%
\end{pgfscope}%
\begin{pgfscope}%
\pgfpathrectangle{\pgfqpoint{0.407592in}{0.222904in}}{\pgfqpoint{1.411163in}{1.411163in}}%
\pgfusepath{clip}%
\pgfsetbuttcap%
\pgfsetroundjoin%
\definecolor{currentfill}{rgb}{0.749020,0.772549,0.792157}%
\pgfsetfillcolor{currentfill}%
\pgfsetfillopacity{0.160000}%
\pgfsetlinewidth{0.000000pt}%
\definecolor{currentstroke}{rgb}{0.000000,0.000000,0.000000}%
\pgfsetstrokecolor{currentstroke}%
\pgfsetdash{}{0pt}%
\pgfpathmoveto{\pgfqpoint{1.385021in}{0.844902in}}%
\pgfpathlineto{\pgfqpoint{1.395404in}{0.853157in}}%
\pgfpathlineto{\pgfqpoint{1.405825in}{0.861442in}}%
\pgfpathlineto{\pgfqpoint{1.424982in}{0.870606in}}%
\pgfpathlineto{\pgfqpoint{1.442518in}{0.880209in}}%
\pgfpathlineto{\pgfqpoint{1.430949in}{0.871353in}}%
\pgfpathlineto{\pgfqpoint{1.419417in}{0.862526in}}%
\pgfpathlineto{\pgfqpoint{1.402972in}{0.853506in}}%
\pgfpathlineto{\pgfqpoint{1.385021in}{0.844902in}}%
\pgfpathclose%
\pgfusepath{fill}%
\end{pgfscope}%
\begin{pgfscope}%
\pgfpathrectangle{\pgfqpoint{0.407592in}{0.222904in}}{\pgfqpoint{1.411163in}{1.411163in}}%
\pgfusepath{clip}%
\pgfsetbuttcap%
\pgfsetroundjoin%
\definecolor{currentfill}{rgb}{0.749020,0.772549,0.792157}%
\pgfsetfillcolor{currentfill}%
\pgfsetfillopacity{0.160000}%
\pgfsetlinewidth{0.000000pt}%
\definecolor{currentstroke}{rgb}{0.000000,0.000000,0.000000}%
\pgfsetstrokecolor{currentstroke}%
\pgfsetdash{}{0pt}%
\pgfpathmoveto{\pgfqpoint{1.216725in}{0.763549in}}%
\pgfpathlineto{\pgfqpoint{1.222289in}{0.770280in}}%
\pgfpathlineto{\pgfqpoint{1.227878in}{0.777040in}}%
\pgfpathlineto{\pgfqpoint{1.249487in}{0.781578in}}%
\pgfpathlineto{\pgfqpoint{1.270310in}{0.786683in}}%
\pgfpathlineto{\pgfqpoint{1.263018in}{0.779550in}}%
\pgfpathlineto{\pgfqpoint{1.255757in}{0.772447in}}%
\pgfpathlineto{\pgfqpoint{1.236599in}{0.767736in}}%
\pgfpathlineto{\pgfqpoint{1.216725in}{0.763549in}}%
\pgfpathclose%
\pgfusepath{fill}%
\end{pgfscope}%
\begin{pgfscope}%
\pgfpathrectangle{\pgfqpoint{0.407592in}{0.222904in}}{\pgfqpoint{1.411163in}{1.411163in}}%
\pgfusepath{clip}%
\pgfsetbuttcap%
\pgfsetroundjoin%
\definecolor{currentfill}{rgb}{0.749020,0.772549,0.792157}%
\pgfsetfillcolor{currentfill}%
\pgfsetfillopacity{0.160000}%
\pgfsetlinewidth{0.000000pt}%
\definecolor{currentstroke}{rgb}{0.000000,0.000000,0.000000}%
\pgfsetstrokecolor{currentstroke}%
\pgfsetdash{}{0pt}%
\pgfpathmoveto{\pgfqpoint{1.008549in}{0.741349in}}%
\pgfpathlineto{\pgfqpoint{1.004636in}{0.747616in}}%
\pgfpathlineto{\pgfqpoint{1.000705in}{0.753911in}}%
\pgfpathlineto{\pgfqpoint{1.022131in}{0.752134in}}%
\pgfpathlineto{\pgfqpoint{1.043882in}{0.750965in}}%
\pgfpathlineto{\pgfqpoint{1.045919in}{0.744794in}}%
\pgfpathlineto{\pgfqpoint{1.047946in}{0.738652in}}%
\pgfpathlineto{\pgfqpoint{1.028100in}{0.739723in}}%
\pgfpathlineto{\pgfqpoint{1.008549in}{0.741349in}}%
\pgfpathclose%
\pgfusepath{fill}%
\end{pgfscope}%
\begin{pgfscope}%
\pgfpathrectangle{\pgfqpoint{0.407592in}{0.222904in}}{\pgfqpoint{1.411163in}{1.411163in}}%
\pgfusepath{clip}%
\pgfsetbuttcap%
\pgfsetroundjoin%
\definecolor{currentfill}{rgb}{0.749020,0.772549,0.792157}%
\pgfsetfillcolor{currentfill}%
\pgfsetfillopacity{0.160000}%
\pgfsetlinewidth{0.000000pt}%
\definecolor{currentstroke}{rgb}{0.000000,0.000000,0.000000}%
\pgfsetstrokecolor{currentstroke}%
\pgfsetdash{}{0pt}%
\pgfpathmoveto{\pgfqpoint{0.526239in}{1.035773in}}%
\pgfpathlineto{\pgfqpoint{0.512144in}{1.046401in}}%
\pgfpathlineto{\pgfqpoint{0.498016in}{1.057053in}}%
\pgfpathlineto{\pgfqpoint{0.506254in}{1.042282in}}%
\pgfpathlineto{\pgfqpoint{0.517197in}{1.027733in}}%
\pgfpathlineto{\pgfqpoint{0.530992in}{1.017690in}}%
\pgfpathlineto{\pgfqpoint{0.544749in}{1.007673in}}%
\pgfpathlineto{\pgfqpoint{0.534210in}{1.021614in}}%
\pgfpathlineto{\pgfqpoint{0.526239in}{1.035773in}}%
\pgfpathclose%
\pgfusepath{fill}%
\end{pgfscope}%
\begin{pgfscope}%
\pgfpathrectangle{\pgfqpoint{0.407592in}{0.222904in}}{\pgfqpoint{1.411163in}{1.411163in}}%
\pgfusepath{clip}%
\pgfsetbuttcap%
\pgfsetroundjoin%
\definecolor{currentfill}{rgb}{0.749020,0.772549,0.792157}%
\pgfsetfillcolor{currentfill}%
\pgfsetfillopacity{0.160000}%
\pgfsetlinewidth{0.000000pt}%
\definecolor{currentstroke}{rgb}{0.000000,0.000000,0.000000}%
\pgfsetstrokecolor{currentstroke}%
\pgfsetdash{}{0pt}%
\pgfpathmoveto{\pgfqpoint{0.626505in}{0.948151in}}%
\pgfpathlineto{\pgfqpoint{0.612972in}{0.958003in}}%
\pgfpathlineto{\pgfqpoint{0.599402in}{0.967883in}}%
\pgfpathlineto{\pgfqpoint{0.611388in}{0.955450in}}%
\pgfpathlineto{\pgfqpoint{0.625626in}{0.943352in}}%
\pgfpathlineto{\pgfqpoint{0.638566in}{0.934083in}}%
\pgfpathlineto{\pgfqpoint{0.651466in}{0.924842in}}%
\pgfpathlineto{\pgfqpoint{0.637925in}{0.936335in}}%
\pgfpathlineto{\pgfqpoint{0.626505in}{0.948151in}}%
\pgfpathclose%
\pgfusepath{fill}%
\end{pgfscope}%
\begin{pgfscope}%
\pgfpathrectangle{\pgfqpoint{0.407592in}{0.222904in}}{\pgfqpoint{1.411163in}{1.411163in}}%
\pgfusepath{clip}%
\pgfsetbuttcap%
\pgfsetroundjoin%
\definecolor{currentfill}{rgb}{0.749020,0.772549,0.792157}%
\pgfsetfillcolor{currentfill}%
\pgfsetfillopacity{0.160000}%
\pgfsetlinewidth{0.000000pt}%
\definecolor{currentstroke}{rgb}{0.000000,0.000000,0.000000}%
\pgfsetstrokecolor{currentstroke}%
\pgfsetdash{}{0pt}%
\pgfpathmoveto{\pgfqpoint{1.777633in}{1.146626in}}%
\pgfpathlineto{\pgfqpoint{1.791655in}{1.157562in}}%
\pgfpathlineto{\pgfqpoint{1.805709in}{1.168522in}}%
\pgfpathlineto{\pgfqpoint{1.815467in}{1.186405in}}%
\pgfpathlineto{\pgfqpoint{1.821891in}{1.204326in}}%
\pgfpathlineto{\pgfqpoint{1.807686in}{1.192798in}}%
\pgfpathlineto{\pgfqpoint{1.793508in}{1.181291in}}%
\pgfpathlineto{\pgfqpoint{1.787171in}{1.163937in}}%
\pgfpathlineto{\pgfqpoint{1.777633in}{1.146626in}}%
\pgfpathclose%
\pgfusepath{fill}%
\end{pgfscope}%
\begin{pgfscope}%
\pgfpathrectangle{\pgfqpoint{0.407592in}{0.222904in}}{\pgfqpoint{1.411163in}{1.411163in}}%
\pgfusepath{clip}%
\pgfsetbuttcap%
\pgfsetroundjoin%
\definecolor{currentfill}{rgb}{0.749020,0.772549,0.792157}%
\pgfsetfillcolor{currentfill}%
\pgfsetfillopacity{0.160000}%
\pgfsetlinewidth{0.000000pt}%
\definecolor{currentstroke}{rgb}{0.000000,0.000000,0.000000}%
\pgfsetstrokecolor{currentstroke}%
\pgfsetdash{}{0pt}%
\pgfpathmoveto{\pgfqpoint{1.088135in}{0.738201in}}%
\pgfpathlineto{\pgfqpoint{1.088034in}{0.744322in}}%
\pgfpathlineto{\pgfqpoint{1.087933in}{0.750472in}}%
\pgfpathlineto{\pgfqpoint{1.110017in}{0.751151in}}%
\pgfpathlineto{\pgfqpoint{1.131996in}{0.752442in}}%
\pgfpathlineto{\pgfqpoint{1.130161in}{0.746209in}}%
\pgfpathlineto{\pgfqpoint{1.128335in}{0.740005in}}%
\pgfpathlineto{\pgfqpoint{1.108283in}{0.738822in}}%
\pgfpathlineto{\pgfqpoint{1.088135in}{0.738201in}}%
\pgfpathclose%
\pgfusepath{fill}%
\end{pgfscope}%
\begin{pgfscope}%
\pgfpathrectangle{\pgfqpoint{0.407592in}{0.222904in}}{\pgfqpoint{1.411163in}{1.411163in}}%
\pgfusepath{clip}%
\pgfsetbuttcap%
\pgfsetroundjoin%
\definecolor{currentfill}{rgb}{0.749020,0.772549,0.792157}%
\pgfsetfillcolor{currentfill}%
\pgfsetfillopacity{0.160000}%
\pgfsetlinewidth{0.000000pt}%
\definecolor{currentstroke}{rgb}{0.000000,0.000000,0.000000}%
\pgfsetstrokecolor{currentstroke}%
\pgfsetdash{}{0pt}%
\pgfpathmoveto{\pgfqpoint{1.047946in}{0.738652in}}%
\pgfpathlineto{\pgfqpoint{1.045919in}{0.744794in}}%
\pgfpathlineto{\pgfqpoint{1.043882in}{0.750965in}}%
\pgfpathlineto{\pgfqpoint{1.065852in}{0.750410in}}%
\pgfpathlineto{\pgfqpoint{1.087933in}{0.750472in}}%
\pgfpathlineto{\pgfqpoint{1.088034in}{0.744322in}}%
\pgfpathlineto{\pgfqpoint{1.088135in}{0.738201in}}%
\pgfpathlineto{\pgfqpoint{1.067990in}{0.738144in}}%
\pgfpathlineto{\pgfqpoint{1.047946in}{0.738652in}}%
\pgfpathclose%
\pgfusepath{fill}%
\end{pgfscope}%
\begin{pgfscope}%
\pgfpathrectangle{\pgfqpoint{0.407592in}{0.222904in}}{\pgfqpoint{1.411163in}{1.411163in}}%
\pgfusepath{clip}%
\pgfsetbuttcap%
\pgfsetroundjoin%
\definecolor{currentfill}{rgb}{0.749020,0.772549,0.792157}%
\pgfsetfillcolor{currentfill}%
\pgfsetfillopacity{0.160000}%
\pgfsetlinewidth{0.000000pt}%
\definecolor{currentstroke}{rgb}{0.000000,0.000000,0.000000}%
\pgfsetstrokecolor{currentstroke}%
\pgfsetdash{}{0pt}%
\pgfpathmoveto{\pgfqpoint{0.702659in}{0.888170in}}%
\pgfpathlineto{\pgfqpoint{0.689921in}{0.897295in}}%
\pgfpathlineto{\pgfqpoint{0.677143in}{0.906448in}}%
\pgfpathlineto{\pgfqpoint{0.691925in}{0.895932in}}%
\pgfpathlineto{\pgfqpoint{0.708603in}{0.885836in}}%
\pgfpathlineto{\pgfqpoint{0.720480in}{0.877268in}}%
\pgfpathlineto{\pgfqpoint{0.732315in}{0.868730in}}%
\pgfpathlineto{\pgfqpoint{0.716603in}{0.878250in}}%
\pgfpathlineto{\pgfqpoint{0.702659in}{0.888170in}}%
\pgfpathclose%
\pgfusepath{fill}%
\end{pgfscope}%
\begin{pgfscope}%
\pgfpathrectangle{\pgfqpoint{0.407592in}{0.222904in}}{\pgfqpoint{1.411163in}{1.411163in}}%
\pgfusepath{clip}%
\pgfsetbuttcap%
\pgfsetroundjoin%
\definecolor{currentfill}{rgb}{0.749020,0.772549,0.792157}%
\pgfsetfillcolor{currentfill}%
\pgfsetfillopacity{0.160000}%
\pgfsetlinewidth{0.000000pt}%
\definecolor{currentstroke}{rgb}{0.000000,0.000000,0.000000}%
\pgfsetstrokecolor{currentstroke}%
\pgfsetdash{}{0pt}%
\pgfpathmoveto{\pgfqpoint{1.326989in}{0.813739in}}%
\pgfpathlineto{\pgfqpoint{1.335931in}{0.821415in}}%
\pgfpathlineto{\pgfqpoint{1.344910in}{0.829121in}}%
\pgfpathlineto{\pgfqpoint{1.365639in}{0.836759in}}%
\pgfpathlineto{\pgfqpoint{1.385021in}{0.844902in}}%
\pgfpathlineto{\pgfqpoint{1.374675in}{0.836677in}}%
\pgfpathlineto{\pgfqpoint{1.364368in}{0.828482in}}%
\pgfpathlineto{\pgfqpoint{1.346300in}{0.820874in}}%
\pgfpathlineto{\pgfqpoint{1.326989in}{0.813739in}}%
\pgfpathclose%
\pgfusepath{fill}%
\end{pgfscope}%
\begin{pgfscope}%
\pgfpathrectangle{\pgfqpoint{0.407592in}{0.222904in}}{\pgfqpoint{1.411163in}{1.411163in}}%
\pgfusepath{clip}%
\pgfsetbuttcap%
\pgfsetroundjoin%
\definecolor{currentfill}{rgb}{0.749020,0.772549,0.792157}%
\pgfsetfillcolor{currentfill}%
\pgfsetfillopacity{0.160000}%
\pgfsetlinewidth{0.000000pt}%
\definecolor{currentstroke}{rgb}{0.000000,0.000000,0.000000}%
\pgfsetstrokecolor{currentstroke}%
\pgfsetdash{}{0pt}%
\pgfpathmoveto{\pgfqpoint{1.270310in}{0.786683in}}%
\pgfpathlineto{\pgfqpoint{1.277632in}{0.793845in}}%
\pgfpathlineto{\pgfqpoint{1.284986in}{0.801039in}}%
\pgfpathlineto{\pgfqpoint{1.306520in}{0.807116in}}%
\pgfpathlineto{\pgfqpoint{1.326989in}{0.813739in}}%
\pgfpathlineto{\pgfqpoint{1.318081in}{0.806094in}}%
\pgfpathlineto{\pgfqpoint{1.309209in}{0.798479in}}%
\pgfpathlineto{\pgfqpoint{1.290248in}{0.792326in}}%
\pgfpathlineto{\pgfqpoint{1.270310in}{0.786683in}}%
\pgfpathclose%
\pgfusepath{fill}%
\end{pgfscope}%
\begin{pgfscope}%
\pgfpathrectangle{\pgfqpoint{0.407592in}{0.222904in}}{\pgfqpoint{1.411163in}{1.411163in}}%
\pgfusepath{clip}%
\pgfsetbuttcap%
\pgfsetroundjoin%
\definecolor{currentfill}{rgb}{0.749020,0.772549,0.792157}%
\pgfsetfillcolor{currentfill}%
\pgfsetfillopacity{0.160000}%
\pgfsetlinewidth{0.000000pt}%
\definecolor{currentstroke}{rgb}{0.000000,0.000000,0.000000}%
\pgfsetstrokecolor{currentstroke}%
\pgfsetdash{}{0pt}%
\pgfpathmoveto{\pgfqpoint{0.375133in}{1.177810in}}%
\pgfpathlineto{\pgfqpoint{0.360700in}{1.189260in}}%
\pgfpathlineto{\pgfqpoint{0.346239in}{1.200731in}}%
\pgfpathlineto{\pgfqpoint{0.348903in}{1.182812in}}%
\pgfpathlineto{\pgfqpoint{0.354919in}{1.164946in}}%
\pgfpathlineto{\pgfqpoint{0.369379in}{1.154043in}}%
\pgfpathlineto{\pgfqpoint{0.383806in}{1.143166in}}%
\pgfpathlineto{\pgfqpoint{0.377860in}{1.160458in}}%
\pgfpathlineto{\pgfqpoint{0.375133in}{1.177810in}}%
\pgfpathclose%
\pgfusepath{fill}%
\end{pgfscope}%
\begin{pgfscope}%
\pgfpathrectangle{\pgfqpoint{0.407592in}{0.222904in}}{\pgfqpoint{1.411163in}{1.411163in}}%
\pgfusepath{clip}%
\pgfsetbuttcap%
\pgfsetroundjoin%
\definecolor{currentfill}{rgb}{0.749020,0.772549,0.792157}%
\pgfsetfillcolor{currentfill}%
\pgfsetfillopacity{0.160000}%
\pgfsetlinewidth{0.000000pt}%
\definecolor{currentstroke}{rgb}{0.000000,0.000000,0.000000}%
\pgfsetstrokecolor{currentstroke}%
\pgfsetdash{}{0pt}%
\pgfpathmoveto{\pgfqpoint{0.920300in}{0.766854in}}%
\pgfpathlineto{\pgfqpoint{0.912833in}{0.773723in}}%
\pgfpathlineto{\pgfqpoint{0.905333in}{0.780622in}}%
\pgfpathlineto{\pgfqpoint{0.926109in}{0.776200in}}%
\pgfpathlineto{\pgfqpoint{0.947680in}{0.772375in}}%
\pgfpathlineto{\pgfqpoint{0.953476in}{0.765795in}}%
\pgfpathlineto{\pgfqpoint{0.959245in}{0.759245in}}%
\pgfpathlineto{\pgfqpoint{0.939410in}{0.762774in}}%
\pgfpathlineto{\pgfqpoint{0.920300in}{0.766854in}}%
\pgfpathclose%
\pgfusepath{fill}%
\end{pgfscope}%
\begin{pgfscope}%
\pgfpathrectangle{\pgfqpoint{0.407592in}{0.222904in}}{\pgfqpoint{1.411163in}{1.411163in}}%
\pgfusepath{clip}%
\pgfsetbuttcap%
\pgfsetroundjoin%
\definecolor{currentfill}{rgb}{0.749020,0.772549,0.792157}%
\pgfsetfillcolor{currentfill}%
\pgfsetfillopacity{0.160000}%
\pgfsetlinewidth{0.000000pt}%
\definecolor{currentstroke}{rgb}{0.000000,0.000000,0.000000}%
\pgfsetstrokecolor{currentstroke}%
\pgfsetdash{}{0pt}%
\pgfpathmoveto{\pgfqpoint{1.465771in}{0.898010in}}%
\pgfpathlineto{\pgfqpoint{1.477455in}{0.906954in}}%
\pgfpathlineto{\pgfqpoint{1.489178in}{0.915928in}}%
\pgfpathlineto{\pgfqpoint{1.506955in}{0.927119in}}%
\pgfpathlineto{\pgfqpoint{1.494750in}{0.917847in}}%
\pgfpathlineto{\pgfqpoint{1.482584in}{0.908604in}}%
\pgfusepath{fill}%
\end{pgfscope}%
\begin{pgfscope}%
\pgfpathrectangle{\pgfqpoint{0.407592in}{0.222904in}}{\pgfqpoint{1.411163in}{1.411163in}}%
\pgfusepath{clip}%
\pgfsetbuttcap%
\pgfsetroundjoin%
\definecolor{currentfill}{rgb}{0.749020,0.772549,0.792157}%
\pgfsetfillcolor{currentfill}%
\pgfsetfillopacity{0.160000}%
\pgfsetlinewidth{0.000000pt}%
\definecolor{currentstroke}{rgb}{0.000000,0.000000,0.000000}%
\pgfsetstrokecolor{currentstroke}%
\pgfsetdash{}{0pt}%
\pgfpathmoveto{\pgfqpoint{1.641536in}{1.030631in}}%
\pgfpathlineto{\pgfqpoint{1.654833in}{1.040763in}}%
\pgfpathlineto{\pgfqpoint{1.668165in}{1.050923in}}%
\pgfpathlineto{\pgfqpoint{1.682564in}{1.066125in}}%
\pgfpathlineto{\pgfqpoint{1.694172in}{1.081536in}}%
\pgfpathlineto{\pgfqpoint{1.680372in}{1.070774in}}%
\pgfpathlineto{\pgfqpoint{1.666604in}{1.060037in}}%
\pgfpathlineto{\pgfqpoint{1.655398in}{1.045230in}}%
\pgfpathlineto{\pgfqpoint{1.641536in}{1.030631in}}%
\pgfpathclose%
\pgfusepath{fill}%
\end{pgfscope}%
\begin{pgfscope}%
\pgfpathrectangle{\pgfqpoint{0.407592in}{0.222904in}}{\pgfqpoint{1.411163in}{1.411163in}}%
\pgfusepath{clip}%
\pgfsetbuttcap%
\pgfsetroundjoin%
\definecolor{currentfill}{rgb}{0.749020,0.772549,0.792157}%
\pgfsetfillcolor{currentfill}%
\pgfsetfillopacity{0.160000}%
\pgfsetlinewidth{0.000000pt}%
\definecolor{currentstroke}{rgb}{0.000000,0.000000,0.000000}%
\pgfsetstrokecolor{currentstroke}%
\pgfsetdash{}{0pt}%
\pgfpathmoveto{\pgfqpoint{1.805709in}{1.168522in}}%
\pgfpathlineto{\pgfqpoint{1.819795in}{1.179507in}}%
\pgfpathlineto{\pgfqpoint{1.833914in}{1.190518in}}%
\pgfpathlineto{\pgfqpoint{1.843881in}{1.208967in}}%
\pgfpathlineto{\pgfqpoint{1.850378in}{1.227445in}}%
\pgfpathlineto{\pgfqpoint{1.836121in}{1.215875in}}%
\pgfpathlineto{\pgfqpoint{1.821891in}{1.204326in}}%
\pgfpathlineto{\pgfqpoint{1.815467in}{1.186405in}}%
\pgfpathlineto{\pgfqpoint{1.805709in}{1.168522in}}%
\pgfpathclose%
\pgfusepath{fill}%
\end{pgfscope}%
\begin{pgfscope}%
\pgfpathrectangle{\pgfqpoint{0.407592in}{0.222904in}}{\pgfqpoint{1.411163in}{1.411163in}}%
\pgfusepath{clip}%
\pgfsetbuttcap%
\pgfsetroundjoin%
\definecolor{currentfill}{rgb}{0.749020,0.772549,0.792157}%
\pgfsetfillcolor{currentfill}%
\pgfsetfillopacity{0.160000}%
\pgfsetlinewidth{0.000000pt}%
\definecolor{currentstroke}{rgb}{0.000000,0.000000,0.000000}%
\pgfsetstrokecolor{currentstroke}%
\pgfsetdash{}{0pt}%
\pgfpathmoveto{\pgfqpoint{1.175206in}{0.756833in}}%
\pgfpathlineto{\pgfqpoint{1.178951in}{0.763281in}}%
\pgfpathlineto{\pgfqpoint{1.182714in}{0.769759in}}%
\pgfpathlineto{\pgfqpoint{1.205584in}{0.773094in}}%
\pgfpathlineto{\pgfqpoint{1.227878in}{0.777040in}}%
\pgfpathlineto{\pgfqpoint{1.222289in}{0.770280in}}%
\pgfpathlineto{\pgfqpoint{1.216725in}{0.763549in}}%
\pgfpathlineto{\pgfqpoint{1.196228in}{0.759908in}}%
\pgfpathlineto{\pgfqpoint{1.175206in}{0.756833in}}%
\pgfpathclose%
\pgfusepath{fill}%
\end{pgfscope}%
\begin{pgfscope}%
\pgfpathrectangle{\pgfqpoint{0.407592in}{0.222904in}}{\pgfqpoint{1.411163in}{1.411163in}}%
\pgfusepath{clip}%
\pgfsetbuttcap%
\pgfsetroundjoin%
\definecolor{currentfill}{rgb}{0.749020,0.772549,0.792157}%
\pgfsetfillcolor{currentfill}%
\pgfsetfillopacity{0.160000}%
\pgfsetlinewidth{0.000000pt}%
\definecolor{currentstroke}{rgb}{0.000000,0.000000,0.000000}%
\pgfsetstrokecolor{currentstroke}%
\pgfsetdash{}{0pt}%
\pgfpathmoveto{\pgfqpoint{1.531478in}{0.945749in}}%
\pgfpathlineto{\pgfqpoint{1.543797in}{0.955108in}}%
\pgfpathlineto{\pgfqpoint{1.556155in}{0.964496in}}%
\pgfpathlineto{\pgfqpoint{1.573590in}{0.977270in}}%
\pgfpathlineto{\pgfqpoint{1.588709in}{0.990377in}}%
\pgfpathlineto{\pgfqpoint{1.575592in}{0.980381in}}%
\pgfpathlineto{\pgfqpoint{1.562509in}{0.970412in}}%
\pgfpathlineto{\pgfqpoint{1.548086in}{0.957919in}}%
\pgfpathlineto{\pgfqpoint{1.531478in}{0.945749in}}%
\pgfpathclose%
\pgfusepath{fill}%
\end{pgfscope}%
\begin{pgfscope}%
\pgfpathrectangle{\pgfqpoint{0.407592in}{0.222904in}}{\pgfqpoint{1.411163in}{1.411163in}}%
\pgfusepath{clip}%
\pgfsetbuttcap%
\pgfsetroundjoin%
\definecolor{currentfill}{rgb}{0.749020,0.772549,0.792157}%
\pgfsetfillcolor{currentfill}%
\pgfsetfillopacity{0.160000}%
\pgfsetlinewidth{0.000000pt}%
\definecolor{currentstroke}{rgb}{0.000000,0.000000,0.000000}%
\pgfsetstrokecolor{currentstroke}%
\pgfsetdash{}{0pt}%
\pgfpathmoveto{\pgfqpoint{0.755860in}{0.851745in}}%
\pgfpathlineto{\pgfqpoint{0.744108in}{0.860222in}}%
\pgfpathlineto{\pgfqpoint{0.732315in}{0.868730in}}%
\pgfpathlineto{\pgfqpoint{0.749735in}{0.859661in}}%
\pgfpathlineto{\pgfqpoint{0.768792in}{0.851089in}}%
\pgfpathlineto{\pgfqpoint{0.779438in}{0.843121in}}%
\pgfpathlineto{\pgfqpoint{0.790043in}{0.835185in}}%
\pgfpathlineto{\pgfqpoint{0.772192in}{0.843230in}}%
\pgfpathlineto{\pgfqpoint{0.755860in}{0.851745in}}%
\pgfpathclose%
\pgfusepath{fill}%
\end{pgfscope}%
\begin{pgfscope}%
\pgfpathrectangle{\pgfqpoint{0.407592in}{0.222904in}}{\pgfqpoint{1.411163in}{1.411163in}}%
\pgfusepath{clip}%
\pgfsetbuttcap%
\pgfsetroundjoin%
\definecolor{currentfill}{rgb}{0.749020,0.772549,0.792157}%
\pgfsetfillcolor{currentfill}%
\pgfsetfillopacity{0.160000}%
\pgfsetlinewidth{0.000000pt}%
\definecolor{currentstroke}{rgb}{0.000000,0.000000,0.000000}%
\pgfsetstrokecolor{currentstroke}%
\pgfsetdash{}{0pt}%
\pgfpathmoveto{\pgfqpoint{0.498016in}{1.057053in}}%
\pgfpathlineto{\pgfqpoint{0.483856in}{1.067730in}}%
\pgfpathlineto{\pgfqpoint{0.469662in}{1.078431in}}%
\pgfpathlineto{\pgfqpoint{0.478156in}{1.063056in}}%
\pgfpathlineto{\pgfqpoint{0.489491in}{1.047904in}}%
\pgfpathlineto{\pgfqpoint{0.503363in}{1.037805in}}%
\pgfpathlineto{\pgfqpoint{0.517197in}{1.027733in}}%
\pgfpathlineto{\pgfqpoint{0.506254in}{1.042282in}}%
\pgfpathlineto{\pgfqpoint{0.498016in}{1.057053in}}%
\pgfpathclose%
\pgfusepath{fill}%
\end{pgfscope}%
\begin{pgfscope}%
\pgfpathrectangle{\pgfqpoint{0.407592in}{0.222904in}}{\pgfqpoint{1.411163in}{1.411163in}}%
\pgfusepath{clip}%
\pgfsetbuttcap%
\pgfsetroundjoin%
\definecolor{currentfill}{rgb}{0.749020,0.772549,0.792157}%
\pgfsetfillcolor{currentfill}%
\pgfsetfillopacity{0.160000}%
\pgfsetlinewidth{0.000000pt}%
\definecolor{currentstroke}{rgb}{0.000000,0.000000,0.000000}%
\pgfsetstrokecolor{currentstroke}%
\pgfsetdash{}{0pt}%
\pgfpathmoveto{\pgfqpoint{1.405825in}{0.861442in}}%
\pgfpathlineto{\pgfqpoint{1.416284in}{0.869757in}}%
\pgfpathlineto{\pgfqpoint{1.426782in}{0.878103in}}%
\pgfpathlineto{\pgfqpoint{1.447146in}{0.887825in}}%
\pgfpathlineto{\pgfqpoint{1.465771in}{0.898010in}}%
\pgfpathlineto{\pgfqpoint{1.454126in}{0.889095in}}%
\pgfpathlineto{\pgfqpoint{1.442518in}{0.880209in}}%
\pgfpathlineto{\pgfqpoint{1.424982in}{0.870606in}}%
\pgfpathlineto{\pgfqpoint{1.405825in}{0.861442in}}%
\pgfpathclose%
\pgfusepath{fill}%
\end{pgfscope}%
\begin{pgfscope}%
\pgfpathrectangle{\pgfqpoint{0.407592in}{0.222904in}}{\pgfqpoint{1.411163in}{1.411163in}}%
\pgfusepath{clip}%
\pgfsetbuttcap%
\pgfsetroundjoin%
\definecolor{currentfill}{rgb}{0.749020,0.772549,0.792157}%
\pgfsetfillcolor{currentfill}%
\pgfsetfillopacity{0.160000}%
\pgfsetlinewidth{0.000000pt}%
\definecolor{currentstroke}{rgb}{0.000000,0.000000,0.000000}%
\pgfsetstrokecolor{currentstroke}%
\pgfsetdash{}{0pt}%
\pgfpathmoveto{\pgfqpoint{0.866553in}{0.791152in}}%
\pgfpathlineto{\pgfqpoint{0.857470in}{0.798487in}}%
\pgfpathlineto{\pgfqpoint{0.848350in}{0.805852in}}%
\pgfpathlineto{\pgfqpoint{0.868755in}{0.799888in}}%
\pgfpathlineto{\pgfqpoint{0.890235in}{0.794510in}}%
\pgfpathlineto{\pgfqpoint{0.897801in}{0.787551in}}%
\pgfpathlineto{\pgfqpoint{0.905333in}{0.780622in}}%
\pgfpathlineto{\pgfqpoint{0.885450in}{0.785614in}}%
\pgfpathlineto{\pgfqpoint{0.866553in}{0.791152in}}%
\pgfpathclose%
\pgfusepath{fill}%
\end{pgfscope}%
\begin{pgfscope}%
\pgfpathrectangle{\pgfqpoint{0.407592in}{0.222904in}}{\pgfqpoint{1.411163in}{1.411163in}}%
\pgfusepath{clip}%
\pgfsetbuttcap%
\pgfsetroundjoin%
\definecolor{currentfill}{rgb}{0.749020,0.772549,0.792157}%
\pgfsetfillcolor{currentfill}%
\pgfsetfillopacity{0.160000}%
\pgfsetlinewidth{0.000000pt}%
\definecolor{currentstroke}{rgb}{0.000000,0.000000,0.000000}%
\pgfsetstrokecolor{currentstroke}%
\pgfsetdash{}{0pt}%
\pgfpathmoveto{\pgfqpoint{0.599402in}{0.967883in}}%
\pgfpathlineto{\pgfqpoint{0.585795in}{0.977790in}}%
\pgfpathlineto{\pgfqpoint{0.572151in}{0.987724in}}%
\pgfpathlineto{\pgfqpoint{0.584694in}{0.974678in}}%
\pgfpathlineto{\pgfqpoint{0.599623in}{0.961980in}}%
\pgfpathlineto{\pgfqpoint{0.612645in}{0.952651in}}%
\pgfpathlineto{\pgfqpoint{0.625626in}{0.943352in}}%
\pgfpathlineto{\pgfqpoint{0.611388in}{0.955450in}}%
\pgfpathlineto{\pgfqpoint{0.599402in}{0.967883in}}%
\pgfpathclose%
\pgfusepath{fill}%
\end{pgfscope}%
\begin{pgfscope}%
\pgfpathrectangle{\pgfqpoint{0.407592in}{0.222904in}}{\pgfqpoint{1.411163in}{1.411163in}}%
\pgfusepath{clip}%
\pgfsetbuttcap%
\pgfsetroundjoin%
\definecolor{currentfill}{rgb}{0.749020,0.772549,0.792157}%
\pgfsetfillcolor{currentfill}%
\pgfsetfillopacity{0.160000}%
\pgfsetlinewidth{0.000000pt}%
\definecolor{currentstroke}{rgb}{0.000000,0.000000,0.000000}%
\pgfsetstrokecolor{currentstroke}%
\pgfsetdash{}{0pt}%
\pgfpathmoveto{\pgfqpoint{0.959245in}{0.759245in}}%
\pgfpathlineto{\pgfqpoint{0.953476in}{0.765795in}}%
\pgfpathlineto{\pgfqpoint{0.947680in}{0.772375in}}%
\pgfpathlineto{\pgfqpoint{0.969942in}{0.769166in}}%
\pgfpathlineto{\pgfqpoint{0.992788in}{0.766591in}}%
\pgfpathlineto{\pgfqpoint{0.996756in}{0.760236in}}%
\pgfpathlineto{\pgfqpoint{1.000705in}{0.753911in}}%
\pgfpathlineto{\pgfqpoint{0.979709in}{0.756286in}}%
\pgfpathlineto{\pgfqpoint{0.959245in}{0.759245in}}%
\pgfpathclose%
\pgfusepath{fill}%
\end{pgfscope}%
\begin{pgfscope}%
\pgfpathrectangle{\pgfqpoint{0.407592in}{0.222904in}}{\pgfqpoint{1.411163in}{1.411163in}}%
\pgfusepath{clip}%
\pgfsetbuttcap%
\pgfsetroundjoin%
\definecolor{currentfill}{rgb}{0.749020,0.772549,0.792157}%
\pgfsetfillcolor{currentfill}%
\pgfsetfillopacity{0.160000}%
\pgfsetlinewidth{0.000000pt}%
\definecolor{currentstroke}{rgb}{0.000000,0.000000,0.000000}%
\pgfsetstrokecolor{currentstroke}%
\pgfsetdash{}{0pt}%
\pgfpathmoveto{\pgfqpoint{0.346239in}{1.200731in}}%
\pgfpathlineto{\pgfqpoint{0.331751in}{1.212225in}}%
\pgfpathlineto{\pgfqpoint{0.317235in}{1.223739in}}%
\pgfpathlineto{\pgfqpoint{0.319823in}{1.205261in}}%
\pgfpathlineto{\pgfqpoint{0.325897in}{1.186828in}}%
\pgfpathlineto{\pgfqpoint{0.340425in}{1.175874in}}%
\pgfpathlineto{\pgfqpoint{0.354919in}{1.164946in}}%
\pgfpathlineto{\pgfqpoint{0.348903in}{1.182812in}}%
\pgfpathlineto{\pgfqpoint{0.346239in}{1.200731in}}%
\pgfpathclose%
\pgfusepath{fill}%
\end{pgfscope}%
\begin{pgfscope}%
\pgfpathrectangle{\pgfqpoint{0.407592in}{0.222904in}}{\pgfqpoint{1.411163in}{1.411163in}}%
\pgfusepath{clip}%
\pgfsetbuttcap%
\pgfsetroundjoin%
\definecolor{currentfill}{rgb}{0.749020,0.772549,0.792157}%
\pgfsetfillcolor{currentfill}%
\pgfsetfillopacity{0.160000}%
\pgfsetlinewidth{0.000000pt}%
\definecolor{currentstroke}{rgb}{0.000000,0.000000,0.000000}%
\pgfsetstrokecolor{currentstroke}%
\pgfsetdash{}{0pt}%
\pgfpathmoveto{\pgfqpoint{0.811131in}{0.819403in}}%
\pgfpathlineto{\pgfqpoint{0.800608in}{0.827279in}}%
\pgfpathlineto{\pgfqpoint{0.790043in}{0.835185in}}%
\pgfpathlineto{\pgfqpoint{0.809339in}{0.827653in}}%
\pgfpathlineto{\pgfqpoint{0.829995in}{0.820673in}}%
\pgfpathlineto{\pgfqpoint{0.839191in}{0.813247in}}%
\pgfpathlineto{\pgfqpoint{0.848350in}{0.805852in}}%
\pgfpathlineto{\pgfqpoint{0.829113in}{0.812368in}}%
\pgfpathlineto{\pgfqpoint{0.811131in}{0.819403in}}%
\pgfpathclose%
\pgfusepath{fill}%
\end{pgfscope}%
\begin{pgfscope}%
\pgfpathrectangle{\pgfqpoint{0.407592in}{0.222904in}}{\pgfqpoint{1.411163in}{1.411163in}}%
\pgfusepath{clip}%
\pgfsetbuttcap%
\pgfsetroundjoin%
\definecolor{currentfill}{rgb}{0.749020,0.772549,0.792157}%
\pgfsetfillcolor{currentfill}%
\pgfsetfillopacity{0.160000}%
\pgfsetlinewidth{0.000000pt}%
\definecolor{currentstroke}{rgb}{0.000000,0.000000,0.000000}%
\pgfsetstrokecolor{currentstroke}%
\pgfsetdash{}{0pt}%
\pgfpathmoveto{\pgfqpoint{1.833914in}{1.190518in}}%
\pgfpathlineto{\pgfqpoint{1.848065in}{1.201554in}}%
\pgfpathlineto{\pgfqpoint{1.858131in}{1.220283in}}%
\pgfpathlineto{\pgfqpoint{1.864660in}{1.239036in}}%
\pgfpathlineto{\pgfqpoint{1.850378in}{1.227445in}}%
\pgfpathlineto{\pgfqpoint{1.843881in}{1.208967in}}%
\pgfusepath{fill}%
\end{pgfscope}%
\begin{pgfscope}%
\pgfpathrectangle{\pgfqpoint{0.407592in}{0.222904in}}{\pgfqpoint{1.411163in}{1.411163in}}%
\pgfusepath{clip}%
\pgfsetbuttcap%
\pgfsetroundjoin%
\definecolor{currentfill}{rgb}{0.749020,0.772549,0.792157}%
\pgfsetfillcolor{currentfill}%
\pgfsetfillopacity{0.160000}%
\pgfsetlinewidth{0.000000pt}%
\definecolor{currentstroke}{rgb}{0.000000,0.000000,0.000000}%
\pgfsetstrokecolor{currentstroke}%
\pgfsetdash{}{0pt}%
\pgfpathmoveto{\pgfqpoint{0.677143in}{0.906448in}}%
\pgfpathlineto{\pgfqpoint{0.664324in}{0.915631in}}%
\pgfpathlineto{\pgfqpoint{0.651466in}{0.924842in}}%
\pgfpathlineto{\pgfqpoint{0.667083in}{0.913732in}}%
\pgfpathlineto{\pgfqpoint{0.684723in}{0.903063in}}%
\pgfpathlineto{\pgfqpoint{0.696684in}{0.894434in}}%
\pgfpathlineto{\pgfqpoint{0.708603in}{0.885836in}}%
\pgfpathlineto{\pgfqpoint{0.691925in}{0.895932in}}%
\pgfpathlineto{\pgfqpoint{0.677143in}{0.906448in}}%
\pgfpathclose%
\pgfusepath{fill}%
\end{pgfscope}%
\begin{pgfscope}%
\pgfpathrectangle{\pgfqpoint{0.407592in}{0.222904in}}{\pgfqpoint{1.411163in}{1.411163in}}%
\pgfusepath{clip}%
\pgfsetbuttcap%
\pgfsetroundjoin%
\definecolor{currentfill}{rgb}{0.749020,0.772549,0.792157}%
\pgfsetfillcolor{currentfill}%
\pgfsetfillopacity{0.160000}%
\pgfsetlinewidth{0.000000pt}%
\definecolor{currentstroke}{rgb}{0.000000,0.000000,0.000000}%
\pgfsetstrokecolor{currentstroke}%
\pgfsetdash{}{0pt}%
\pgfpathmoveto{\pgfqpoint{1.131996in}{0.752442in}}%
\pgfpathlineto{\pgfqpoint{1.133839in}{0.758705in}}%
\pgfpathlineto{\pgfqpoint{1.135691in}{0.764998in}}%
\pgfpathlineto{\pgfqpoint{1.159378in}{0.767056in}}%
\pgfpathlineto{\pgfqpoint{1.182714in}{0.769759in}}%
\pgfpathlineto{\pgfqpoint{1.178951in}{0.763281in}}%
\pgfpathlineto{\pgfqpoint{1.175206in}{0.756833in}}%
\pgfpathlineto{\pgfqpoint{1.153761in}{0.754340in}}%
\pgfpathlineto{\pgfqpoint{1.131996in}{0.752442in}}%
\pgfpathclose%
\pgfusepath{fill}%
\end{pgfscope}%
\begin{pgfscope}%
\pgfpathrectangle{\pgfqpoint{0.407592in}{0.222904in}}{\pgfqpoint{1.411163in}{1.411163in}}%
\pgfusepath{clip}%
\pgfsetbuttcap%
\pgfsetroundjoin%
\definecolor{currentfill}{rgb}{0.749020,0.772549,0.792157}%
\pgfsetfillcolor{currentfill}%
\pgfsetfillopacity{0.160000}%
\pgfsetlinewidth{0.000000pt}%
\definecolor{currentstroke}{rgb}{0.000000,0.000000,0.000000}%
\pgfsetstrokecolor{currentstroke}%
\pgfsetdash{}{0pt}%
\pgfpathmoveto{\pgfqpoint{1.668165in}{1.050923in}}%
\pgfpathlineto{\pgfqpoint{1.681534in}{1.061110in}}%
\pgfpathlineto{\pgfqpoint{1.694940in}{1.071325in}}%
\pgfpathlineto{\pgfqpoint{1.709867in}{1.087125in}}%
\pgfpathlineto{\pgfqpoint{1.721866in}{1.103134in}}%
\pgfpathlineto{\pgfqpoint{1.708003in}{1.092323in}}%
\pgfpathlineto{\pgfqpoint{1.694172in}{1.081536in}}%
\pgfpathlineto{\pgfqpoint{1.682564in}{1.066125in}}%
\pgfpathlineto{\pgfqpoint{1.668165in}{1.050923in}}%
\pgfpathclose%
\pgfusepath{fill}%
\end{pgfscope}%
\begin{pgfscope}%
\pgfpathrectangle{\pgfqpoint{0.407592in}{0.222904in}}{\pgfqpoint{1.411163in}{1.411163in}}%
\pgfusepath{clip}%
\pgfsetbuttcap%
\pgfsetroundjoin%
\definecolor{currentfill}{rgb}{0.749020,0.772549,0.792157}%
\pgfsetfillcolor{currentfill}%
\pgfsetfillopacity{0.160000}%
\pgfsetlinewidth{0.000000pt}%
\definecolor{currentstroke}{rgb}{0.000000,0.000000,0.000000}%
\pgfsetstrokecolor{currentstroke}%
\pgfsetdash{}{0pt}%
\pgfpathmoveto{\pgfqpoint{1.227878in}{0.777040in}}%
\pgfpathlineto{\pgfqpoint{1.233491in}{0.783831in}}%
\pgfpathlineto{\pgfqpoint{1.239130in}{0.790652in}}%
\pgfpathlineto{\pgfqpoint{1.262487in}{0.795541in}}%
\pgfpathlineto{\pgfqpoint{1.284986in}{0.801039in}}%
\pgfpathlineto{\pgfqpoint{1.277632in}{0.793845in}}%
\pgfpathlineto{\pgfqpoint{1.270310in}{0.786683in}}%
\pgfpathlineto{\pgfqpoint{1.249487in}{0.781578in}}%
\pgfpathlineto{\pgfqpoint{1.227878in}{0.777040in}}%
\pgfpathclose%
\pgfusepath{fill}%
\end{pgfscope}%
\begin{pgfscope}%
\pgfpathrectangle{\pgfqpoint{0.407592in}{0.222904in}}{\pgfqpoint{1.411163in}{1.411163in}}%
\pgfusepath{clip}%
\pgfsetbuttcap%
\pgfsetroundjoin%
\definecolor{currentfill}{rgb}{0.749020,0.772549,0.792157}%
\pgfsetfillcolor{currentfill}%
\pgfsetfillopacity{0.160000}%
\pgfsetlinewidth{0.000000pt}%
\definecolor{currentstroke}{rgb}{0.000000,0.000000,0.000000}%
\pgfsetstrokecolor{currentstroke}%
\pgfsetdash{}{0pt}%
\pgfpathmoveto{\pgfqpoint{1.000705in}{0.753911in}}%
\pgfpathlineto{\pgfqpoint{0.996756in}{0.760236in}}%
\pgfpathlineto{\pgfqpoint{0.992788in}{0.766591in}}%
\pgfpathlineto{\pgfqpoint{1.016105in}{0.764664in}}%
\pgfpathlineto{\pgfqpoint{1.039780in}{0.763396in}}%
\pgfpathlineto{\pgfqpoint{1.041836in}{0.757166in}}%
\pgfpathlineto{\pgfqpoint{1.043882in}{0.750965in}}%
\pgfpathlineto{\pgfqpoint{1.022131in}{0.752134in}}%
\pgfpathlineto{\pgfqpoint{1.000705in}{0.753911in}}%
\pgfpathclose%
\pgfusepath{fill}%
\end{pgfscope}%
\begin{pgfscope}%
\pgfpathrectangle{\pgfqpoint{0.407592in}{0.222904in}}{\pgfqpoint{1.411163in}{1.411163in}}%
\pgfusepath{clip}%
\pgfsetbuttcap%
\pgfsetroundjoin%
\definecolor{currentfill}{rgb}{0.749020,0.772549,0.792157}%
\pgfsetfillcolor{currentfill}%
\pgfsetfillopacity{0.160000}%
\pgfsetlinewidth{0.000000pt}%
\definecolor{currentstroke}{rgb}{0.000000,0.000000,0.000000}%
\pgfsetstrokecolor{currentstroke}%
\pgfsetdash{}{0pt}%
\pgfpathmoveto{\pgfqpoint{1.344910in}{0.829121in}}%
\pgfpathlineto{\pgfqpoint{1.353924in}{0.836857in}}%
\pgfpathlineto{\pgfqpoint{1.362974in}{0.844625in}}%
\pgfpathlineto{\pgfqpoint{1.385126in}{0.852766in}}%
\pgfpathlineto{\pgfqpoint{1.405825in}{0.861442in}}%
\pgfpathlineto{\pgfqpoint{1.395404in}{0.853157in}}%
\pgfpathlineto{\pgfqpoint{1.385021in}{0.844902in}}%
\pgfpathlineto{\pgfqpoint{1.365639in}{0.836759in}}%
\pgfpathlineto{\pgfqpoint{1.344910in}{0.829121in}}%
\pgfpathclose%
\pgfusepath{fill}%
\end{pgfscope}%
\begin{pgfscope}%
\pgfpathrectangle{\pgfqpoint{0.407592in}{0.222904in}}{\pgfqpoint{1.411163in}{1.411163in}}%
\pgfusepath{clip}%
\pgfsetbuttcap%
\pgfsetroundjoin%
\definecolor{currentfill}{rgb}{0.749020,0.772549,0.792157}%
\pgfsetfillcolor{currentfill}%
\pgfsetfillopacity{0.160000}%
\pgfsetlinewidth{0.000000pt}%
\definecolor{currentstroke}{rgb}{0.000000,0.000000,0.000000}%
\pgfsetstrokecolor{currentstroke}%
\pgfsetdash{}{0pt}%
\pgfpathmoveto{\pgfqpoint{1.489178in}{0.915928in}}%
\pgfpathlineto{\pgfqpoint{1.500940in}{0.924931in}}%
\pgfpathlineto{\pgfqpoint{1.512741in}{0.933965in}}%
\pgfpathlineto{\pgfqpoint{1.531478in}{0.945749in}}%
\pgfpathlineto{\pgfqpoint{1.519197in}{0.936420in}}%
\pgfpathlineto{\pgfqpoint{1.506955in}{0.927119in}}%
\pgfusepath{fill}%
\end{pgfscope}%
\begin{pgfscope}%
\pgfpathrectangle{\pgfqpoint{0.407592in}{0.222904in}}{\pgfqpoint{1.411163in}{1.411163in}}%
\pgfusepath{clip}%
\pgfsetbuttcap%
\pgfsetroundjoin%
\definecolor{currentfill}{rgb}{0.749020,0.772549,0.792157}%
\pgfsetfillcolor{currentfill}%
\pgfsetfillopacity{0.160000}%
\pgfsetlinewidth{0.000000pt}%
\definecolor{currentstroke}{rgb}{0.000000,0.000000,0.000000}%
\pgfsetstrokecolor{currentstroke}%
\pgfsetdash{}{0pt}%
\pgfpathmoveto{\pgfqpoint{1.087933in}{0.750472in}}%
\pgfpathlineto{\pgfqpoint{1.087832in}{0.756652in}}%
\pgfpathlineto{\pgfqpoint{1.087730in}{0.762862in}}%
\pgfpathlineto{\pgfqpoint{1.111769in}{0.763598in}}%
\pgfpathlineto{\pgfqpoint{1.135691in}{0.764998in}}%
\pgfpathlineto{\pgfqpoint{1.133839in}{0.758705in}}%
\pgfpathlineto{\pgfqpoint{1.131996in}{0.752442in}}%
\pgfpathlineto{\pgfqpoint{1.110017in}{0.751151in}}%
\pgfpathlineto{\pgfqpoint{1.087933in}{0.750472in}}%
\pgfpathclose%
\pgfusepath{fill}%
\end{pgfscope}%
\begin{pgfscope}%
\pgfpathrectangle{\pgfqpoint{0.407592in}{0.222904in}}{\pgfqpoint{1.411163in}{1.411163in}}%
\pgfusepath{clip}%
\pgfsetbuttcap%
\pgfsetroundjoin%
\definecolor{currentfill}{rgb}{0.749020,0.772549,0.792157}%
\pgfsetfillcolor{currentfill}%
\pgfsetfillopacity{0.160000}%
\pgfsetlinewidth{0.000000pt}%
\definecolor{currentstroke}{rgb}{0.000000,0.000000,0.000000}%
\pgfsetstrokecolor{currentstroke}%
\pgfsetdash{}{0pt}%
\pgfpathmoveto{\pgfqpoint{0.317235in}{1.223739in}}%
\pgfpathlineto{\pgfqpoint{0.302693in}{1.235276in}}%
\pgfpathlineto{\pgfqpoint{0.305237in}{1.216520in}}%
\pgfpathlineto{\pgfqpoint{0.311335in}{1.197808in}}%
\pgfpathlineto{\pgfqpoint{0.325897in}{1.186828in}}%
\pgfpathlineto{\pgfqpoint{0.319823in}{1.205261in}}%
\pgfusepath{fill}%
\end{pgfscope}%
\begin{pgfscope}%
\pgfpathrectangle{\pgfqpoint{0.407592in}{0.222904in}}{\pgfqpoint{1.411163in}{1.411163in}}%
\pgfusepath{clip}%
\pgfsetbuttcap%
\pgfsetroundjoin%
\definecolor{currentfill}{rgb}{0.749020,0.772549,0.792157}%
\pgfsetfillcolor{currentfill}%
\pgfsetfillopacity{0.160000}%
\pgfsetlinewidth{0.000000pt}%
\definecolor{currentstroke}{rgb}{0.000000,0.000000,0.000000}%
\pgfsetstrokecolor{currentstroke}%
\pgfsetdash{}{0pt}%
\pgfpathmoveto{\pgfqpoint{1.284986in}{0.801039in}}%
\pgfpathlineto{\pgfqpoint{1.292371in}{0.808263in}}%
\pgfpathlineto{\pgfqpoint{1.299788in}{0.815517in}}%
\pgfpathlineto{\pgfqpoint{1.322926in}{0.822027in}}%
\pgfpathlineto{\pgfqpoint{1.344910in}{0.829121in}}%
\pgfpathlineto{\pgfqpoint{1.335931in}{0.821415in}}%
\pgfpathlineto{\pgfqpoint{1.326989in}{0.813739in}}%
\pgfpathlineto{\pgfqpoint{1.306520in}{0.807116in}}%
\pgfpathlineto{\pgfqpoint{1.284986in}{0.801039in}}%
\pgfpathclose%
\pgfusepath{fill}%
\end{pgfscope}%
\begin{pgfscope}%
\pgfpathrectangle{\pgfqpoint{0.407592in}{0.222904in}}{\pgfqpoint{1.411163in}{1.411163in}}%
\pgfusepath{clip}%
\pgfsetbuttcap%
\pgfsetroundjoin%
\definecolor{currentfill}{rgb}{0.749020,0.772549,0.792157}%
\pgfsetfillcolor{currentfill}%
\pgfsetfillopacity{0.160000}%
\pgfsetlinewidth{0.000000pt}%
\definecolor{currentstroke}{rgb}{0.000000,0.000000,0.000000}%
\pgfsetstrokecolor{currentstroke}%
\pgfsetdash{}{0pt}%
\pgfpathmoveto{\pgfqpoint{1.043882in}{0.750965in}}%
\pgfpathlineto{\pgfqpoint{1.041836in}{0.757166in}}%
\pgfpathlineto{\pgfqpoint{1.039780in}{0.763396in}}%
\pgfpathlineto{\pgfqpoint{1.063694in}{0.762794in}}%
\pgfpathlineto{\pgfqpoint{1.087730in}{0.762862in}}%
\pgfpathlineto{\pgfqpoint{1.087832in}{0.756652in}}%
\pgfpathlineto{\pgfqpoint{1.087933in}{0.750472in}}%
\pgfpathlineto{\pgfqpoint{1.065852in}{0.750410in}}%
\pgfpathlineto{\pgfqpoint{1.043882in}{0.750965in}}%
\pgfpathclose%
\pgfusepath{fill}%
\end{pgfscope}%
\begin{pgfscope}%
\pgfpathrectangle{\pgfqpoint{0.407592in}{0.222904in}}{\pgfqpoint{1.411163in}{1.411163in}}%
\pgfusepath{clip}%
\pgfsetbuttcap%
\pgfsetroundjoin%
\definecolor{currentfill}{rgb}{0.749020,0.772549,0.792157}%
\pgfsetfillcolor{currentfill}%
\pgfsetfillopacity{0.160000}%
\pgfsetlinewidth{0.000000pt}%
\definecolor{currentstroke}{rgb}{0.000000,0.000000,0.000000}%
\pgfsetstrokecolor{currentstroke}%
\pgfsetdash{}{0pt}%
\pgfpathmoveto{\pgfqpoint{0.469662in}{1.078431in}}%
\pgfpathlineto{\pgfqpoint{0.455436in}{1.089158in}}%
\pgfpathlineto{\pgfqpoint{0.441176in}{1.099909in}}%
\pgfpathlineto{\pgfqpoint{0.449915in}{1.083936in}}%
\pgfpathlineto{\pgfqpoint{0.461632in}{1.068187in}}%
\pgfpathlineto{\pgfqpoint{0.475581in}{1.058032in}}%
\pgfpathlineto{\pgfqpoint{0.489491in}{1.047904in}}%
\pgfpathlineto{\pgfqpoint{0.478156in}{1.063056in}}%
\pgfpathlineto{\pgfqpoint{0.469662in}{1.078431in}}%
\pgfpathclose%
\pgfusepath{fill}%
\end{pgfscope}%
\begin{pgfscope}%
\pgfpathrectangle{\pgfqpoint{0.407592in}{0.222904in}}{\pgfqpoint{1.411163in}{1.411163in}}%
\pgfusepath{clip}%
\pgfsetbuttcap%
\pgfsetroundjoin%
\definecolor{currentfill}{rgb}{0.749020,0.772549,0.792157}%
\pgfsetfillcolor{currentfill}%
\pgfsetfillopacity{0.160000}%
\pgfsetlinewidth{0.000000pt}%
\definecolor{currentstroke}{rgb}{0.000000,0.000000,0.000000}%
\pgfsetstrokecolor{currentstroke}%
\pgfsetdash{}{0pt}%
\pgfpathmoveto{\pgfqpoint{1.556155in}{0.964496in}}%
\pgfpathlineto{\pgfqpoint{1.568551in}{0.973914in}}%
\pgfpathlineto{\pgfqpoint{1.580986in}{0.983361in}}%
\pgfpathlineto{\pgfqpoint{1.599243in}{0.996733in}}%
\pgfpathlineto{\pgfqpoint{1.615051in}{1.010449in}}%
\pgfpathlineto{\pgfqpoint{1.601862in}{1.000399in}}%
\pgfpathlineto{\pgfqpoint{1.588709in}{0.990377in}}%
\pgfpathlineto{\pgfqpoint{1.573590in}{0.977270in}}%
\pgfpathlineto{\pgfqpoint{1.556155in}{0.964496in}}%
\pgfpathclose%
\pgfusepath{fill}%
\end{pgfscope}%
\begin{pgfscope}%
\pgfpathrectangle{\pgfqpoint{0.407592in}{0.222904in}}{\pgfqpoint{1.411163in}{1.411163in}}%
\pgfusepath{clip}%
\pgfsetbuttcap%
\pgfsetroundjoin%
\definecolor{currentfill}{rgb}{0.749020,0.772549,0.792157}%
\pgfsetfillcolor{currentfill}%
\pgfsetfillopacity{0.160000}%
\pgfsetlinewidth{0.000000pt}%
\definecolor{currentstroke}{rgb}{0.000000,0.000000,0.000000}%
\pgfsetstrokecolor{currentstroke}%
\pgfsetdash{}{0pt}%
\pgfpathmoveto{\pgfqpoint{0.905333in}{0.780622in}}%
\pgfpathlineto{\pgfqpoint{0.897801in}{0.787551in}}%
\pgfpathlineto{\pgfqpoint{0.890235in}{0.794510in}}%
\pgfpathlineto{\pgfqpoint{0.912689in}{0.789747in}}%
\pgfpathlineto{\pgfqpoint{0.936010in}{0.785625in}}%
\pgfpathlineto{\pgfqpoint{0.941858in}{0.778985in}}%
\pgfpathlineto{\pgfqpoint{0.947680in}{0.772375in}}%
\pgfpathlineto{\pgfqpoint{0.926109in}{0.776200in}}%
\pgfpathlineto{\pgfqpoint{0.905333in}{0.780622in}}%
\pgfpathclose%
\pgfusepath{fill}%
\end{pgfscope}%
\begin{pgfscope}%
\pgfpathrectangle{\pgfqpoint{0.407592in}{0.222904in}}{\pgfqpoint{1.411163in}{1.411163in}}%
\pgfusepath{clip}%
\pgfsetbuttcap%
\pgfsetroundjoin%
\definecolor{currentfill}{rgb}{0.749020,0.772549,0.792157}%
\pgfsetfillcolor{currentfill}%
\pgfsetfillopacity{0.160000}%
\pgfsetlinewidth{0.000000pt}%
\definecolor{currentstroke}{rgb}{0.000000,0.000000,0.000000}%
\pgfsetstrokecolor{currentstroke}%
\pgfsetdash{}{0pt}%
\pgfpathmoveto{\pgfqpoint{0.572151in}{0.987724in}}%
\pgfpathlineto{\pgfqpoint{0.558469in}{0.997685in}}%
\pgfpathlineto{\pgfqpoint{0.544749in}{1.007673in}}%
\pgfpathlineto{\pgfqpoint{0.557841in}{0.994021in}}%
\pgfpathlineto{\pgfqpoint{0.573455in}{0.980725in}}%
\pgfpathlineto{\pgfqpoint{0.586560in}{0.971338in}}%
\pgfpathlineto{\pgfqpoint{0.599623in}{0.961980in}}%
\pgfpathlineto{\pgfqpoint{0.584694in}{0.974678in}}%
\pgfpathlineto{\pgfqpoint{0.572151in}{0.987724in}}%
\pgfpathclose%
\pgfusepath{fill}%
\end{pgfscope}%
\begin{pgfscope}%
\pgfpathrectangle{\pgfqpoint{0.407592in}{0.222904in}}{\pgfqpoint{1.411163in}{1.411163in}}%
\pgfusepath{clip}%
\pgfsetbuttcap%
\pgfsetroundjoin%
\definecolor{currentfill}{rgb}{0.749020,0.772549,0.792157}%
\pgfsetfillcolor{currentfill}%
\pgfsetfillopacity{0.160000}%
\pgfsetlinewidth{0.000000pt}%
\definecolor{currentstroke}{rgb}{0.000000,0.000000,0.000000}%
\pgfsetstrokecolor{currentstroke}%
\pgfsetdash{}{0pt}%
\pgfpathmoveto{\pgfqpoint{0.732315in}{0.868730in}}%
\pgfpathlineto{\pgfqpoint{0.720480in}{0.877268in}}%
\pgfpathlineto{\pgfqpoint{0.708603in}{0.885836in}}%
\pgfpathlineto{\pgfqpoint{0.727111in}{0.876213in}}%
\pgfpathlineto{\pgfqpoint{0.747376in}{0.867116in}}%
\pgfpathlineto{\pgfqpoint{0.758105in}{0.859087in}}%
\pgfpathlineto{\pgfqpoint{0.768792in}{0.851089in}}%
\pgfpathlineto{\pgfqpoint{0.749735in}{0.859661in}}%
\pgfpathlineto{\pgfqpoint{0.732315in}{0.868730in}}%
\pgfpathclose%
\pgfusepath{fill}%
\end{pgfscope}%
\begin{pgfscope}%
\pgfpathrectangle{\pgfqpoint{0.407592in}{0.222904in}}{\pgfqpoint{1.411163in}{1.411163in}}%
\pgfusepath{clip}%
\pgfsetbuttcap%
\pgfsetroundjoin%
\definecolor{currentfill}{rgb}{0.749020,0.772549,0.792157}%
\pgfsetfillcolor{currentfill}%
\pgfsetfillopacity{0.160000}%
\pgfsetlinewidth{0.000000pt}%
\definecolor{currentstroke}{rgb}{0.000000,0.000000,0.000000}%
\pgfsetstrokecolor{currentstroke}%
\pgfsetdash{}{0pt}%
\pgfpathmoveto{\pgfqpoint{1.426782in}{0.878103in}}%
\pgfpathlineto{\pgfqpoint{1.437319in}{0.886480in}}%
\pgfpathlineto{\pgfqpoint{1.447894in}{0.894888in}}%
\pgfpathlineto{\pgfqpoint{1.469466in}{0.905165in}}%
\pgfpathlineto{\pgfqpoint{1.489178in}{0.915928in}}%
\pgfpathlineto{\pgfqpoint{1.477455in}{0.906954in}}%
\pgfpathlineto{\pgfqpoint{1.465771in}{0.898010in}}%
\pgfpathlineto{\pgfqpoint{1.447146in}{0.887825in}}%
\pgfpathlineto{\pgfqpoint{1.426782in}{0.878103in}}%
\pgfpathclose%
\pgfusepath{fill}%
\end{pgfscope}%
\begin{pgfscope}%
\pgfpathrectangle{\pgfqpoint{0.407592in}{0.222904in}}{\pgfqpoint{1.411163in}{1.411163in}}%
\pgfusepath{clip}%
\pgfsetbuttcap%
\pgfsetroundjoin%
\definecolor{currentfill}{rgb}{0.749020,0.772549,0.792157}%
\pgfsetfillcolor{currentfill}%
\pgfsetfillopacity{0.160000}%
\pgfsetlinewidth{0.000000pt}%
\definecolor{currentstroke}{rgb}{0.000000,0.000000,0.000000}%
\pgfsetstrokecolor{currentstroke}%
\pgfsetdash{}{0pt}%
\pgfpathmoveto{\pgfqpoint{1.694940in}{1.071325in}}%
\pgfpathlineto{\pgfqpoint{1.708382in}{1.081568in}}%
\pgfpathlineto{\pgfqpoint{1.721861in}{1.091839in}}%
\pgfpathlineto{\pgfqpoint{1.737307in}{1.108230in}}%
\pgfpathlineto{\pgfqpoint{1.749686in}{1.124831in}}%
\pgfpathlineto{\pgfqpoint{1.735760in}{1.113970in}}%
\pgfpathlineto{\pgfqpoint{1.721866in}{1.103134in}}%
\pgfpathlineto{\pgfqpoint{1.709867in}{1.087125in}}%
\pgfpathlineto{\pgfqpoint{1.694940in}{1.071325in}}%
\pgfpathclose%
\pgfusepath{fill}%
\end{pgfscope}%
\begin{pgfscope}%
\pgfpathrectangle{\pgfqpoint{0.407592in}{0.222904in}}{\pgfqpoint{1.411163in}{1.411163in}}%
\pgfusepath{clip}%
\pgfsetbuttcap%
\pgfsetroundjoin%
\definecolor{currentfill}{rgb}{0.749020,0.772549,0.792157}%
\pgfsetfillcolor{currentfill}%
\pgfsetfillopacity{0.160000}%
\pgfsetlinewidth{0.000000pt}%
\definecolor{currentstroke}{rgb}{0.000000,0.000000,0.000000}%
\pgfsetstrokecolor{currentstroke}%
\pgfsetdash{}{0pt}%
\pgfpathmoveto{\pgfqpoint{1.182714in}{0.769759in}}%
\pgfpathlineto{\pgfqpoint{1.186494in}{0.776268in}}%
\pgfpathlineto{\pgfqpoint{1.190292in}{0.782806in}}%
\pgfpathlineto{\pgfqpoint{1.215026in}{0.786399in}}%
\pgfpathlineto{\pgfqpoint{1.239130in}{0.790652in}}%
\pgfpathlineto{\pgfqpoint{1.233491in}{0.783831in}}%
\pgfpathlineto{\pgfqpoint{1.227878in}{0.777040in}}%
\pgfpathlineto{\pgfqpoint{1.205584in}{0.773094in}}%
\pgfpathlineto{\pgfqpoint{1.182714in}{0.769759in}}%
\pgfpathclose%
\pgfusepath{fill}%
\end{pgfscope}%
\begin{pgfscope}%
\pgfpathrectangle{\pgfqpoint{0.407592in}{0.222904in}}{\pgfqpoint{1.411163in}{1.411163in}}%
\pgfusepath{clip}%
\pgfsetbuttcap%
\pgfsetroundjoin%
\definecolor{currentfill}{rgb}{0.749020,0.772549,0.792157}%
\pgfsetfillcolor{currentfill}%
\pgfsetfillopacity{0.160000}%
\pgfsetlinewidth{0.000000pt}%
\definecolor{currentstroke}{rgb}{0.000000,0.000000,0.000000}%
\pgfsetstrokecolor{currentstroke}%
\pgfsetdash{}{0pt}%
\pgfpathmoveto{\pgfqpoint{0.651466in}{0.924842in}}%
\pgfpathlineto{\pgfqpoint{0.638566in}{0.934083in}}%
\pgfpathlineto{\pgfqpoint{0.625626in}{0.943352in}}%
\pgfpathlineto{\pgfqpoint{0.642073in}{0.931653in}}%
\pgfpathlineto{\pgfqpoint{0.660673in}{0.920412in}}%
\pgfpathlineto{\pgfqpoint{0.672719in}{0.911722in}}%
\pgfpathlineto{\pgfqpoint{0.684723in}{0.903063in}}%
\pgfpathlineto{\pgfqpoint{0.667083in}{0.913732in}}%
\pgfpathlineto{\pgfqpoint{0.651466in}{0.924842in}}%
\pgfpathclose%
\pgfusepath{fill}%
\end{pgfscope}%
\begin{pgfscope}%
\pgfpathrectangle{\pgfqpoint{0.407592in}{0.222904in}}{\pgfqpoint{1.411163in}{1.411163in}}%
\pgfusepath{clip}%
\pgfsetbuttcap%
\pgfsetroundjoin%
\definecolor{currentfill}{rgb}{0.749020,0.772549,0.792157}%
\pgfsetfillcolor{currentfill}%
\pgfsetfillopacity{0.160000}%
\pgfsetlinewidth{0.000000pt}%
\definecolor{currentstroke}{rgb}{0.000000,0.000000,0.000000}%
\pgfsetstrokecolor{currentstroke}%
\pgfsetdash{}{0pt}%
\pgfpathmoveto{\pgfqpoint{0.848350in}{0.805852in}}%
\pgfpathlineto{\pgfqpoint{0.839191in}{0.813247in}}%
\pgfpathlineto{\pgfqpoint{0.829995in}{0.820673in}}%
\pgfpathlineto{\pgfqpoint{0.851917in}{0.814285in}}%
\pgfpathlineto{\pgfqpoint{0.875004in}{0.808522in}}%
\pgfpathlineto{\pgfqpoint{0.882636in}{0.801501in}}%
\pgfpathlineto{\pgfqpoint{0.890235in}{0.794510in}}%
\pgfpathlineto{\pgfqpoint{0.868755in}{0.799888in}}%
\pgfpathlineto{\pgfqpoint{0.848350in}{0.805852in}}%
\pgfpathclose%
\pgfusepath{fill}%
\end{pgfscope}%
\begin{pgfscope}%
\pgfpathrectangle{\pgfqpoint{0.407592in}{0.222904in}}{\pgfqpoint{1.411163in}{1.411163in}}%
\pgfusepath{clip}%
\pgfsetbuttcap%
\pgfsetroundjoin%
\definecolor{currentfill}{rgb}{0.749020,0.772549,0.792157}%
\pgfsetfillcolor{currentfill}%
\pgfsetfillopacity{0.160000}%
\pgfsetlinewidth{0.000000pt}%
\definecolor{currentstroke}{rgb}{0.000000,0.000000,0.000000}%
\pgfsetstrokecolor{currentstroke}%
\pgfsetdash{}{0pt}%
\pgfpathmoveto{\pgfqpoint{0.790043in}{0.835185in}}%
\pgfpathlineto{\pgfqpoint{0.779438in}{0.843121in}}%
\pgfpathlineto{\pgfqpoint{0.768792in}{0.851089in}}%
\pgfpathlineto{\pgfqpoint{0.789407in}{0.843060in}}%
\pgfpathlineto{\pgfqpoint{0.811487in}{0.835619in}}%
\pgfpathlineto{\pgfqpoint{0.820760in}{0.828130in}}%
\pgfpathlineto{\pgfqpoint{0.829995in}{0.820673in}}%
\pgfpathlineto{\pgfqpoint{0.809339in}{0.827653in}}%
\pgfpathlineto{\pgfqpoint{0.790043in}{0.835185in}}%
\pgfpathclose%
\pgfusepath{fill}%
\end{pgfscope}%
\begin{pgfscope}%
\pgfpathrectangle{\pgfqpoint{0.407592in}{0.222904in}}{\pgfqpoint{1.411163in}{1.411163in}}%
\pgfusepath{clip}%
\pgfsetbuttcap%
\pgfsetroundjoin%
\definecolor{currentfill}{rgb}{0.749020,0.772549,0.792157}%
\pgfsetfillcolor{currentfill}%
\pgfsetfillopacity{0.160000}%
\pgfsetlinewidth{0.000000pt}%
\definecolor{currentstroke}{rgb}{0.000000,0.000000,0.000000}%
\pgfsetstrokecolor{currentstroke}%
\pgfsetdash{}{0pt}%
\pgfpathmoveto{\pgfqpoint{0.441176in}{1.099909in}}%
\pgfpathlineto{\pgfqpoint{0.426884in}{1.110686in}}%
\pgfpathlineto{\pgfqpoint{0.412558in}{1.121487in}}%
\pgfpathlineto{\pgfqpoint{0.421529in}{1.104922in}}%
\pgfpathlineto{\pgfqpoint{0.433618in}{1.088583in}}%
\pgfpathlineto{\pgfqpoint{0.447645in}{1.078371in}}%
\pgfpathlineto{\pgfqpoint{0.461632in}{1.068187in}}%
\pgfpathlineto{\pgfqpoint{0.449915in}{1.083936in}}%
\pgfpathlineto{\pgfqpoint{0.441176in}{1.099909in}}%
\pgfpathclose%
\pgfusepath{fill}%
\end{pgfscope}%
\begin{pgfscope}%
\pgfpathrectangle{\pgfqpoint{0.407592in}{0.222904in}}{\pgfqpoint{1.411163in}{1.411163in}}%
\pgfusepath{clip}%
\pgfsetbuttcap%
\pgfsetroundjoin%
\definecolor{currentfill}{rgb}{0.749020,0.772549,0.792157}%
\pgfsetfillcolor{currentfill}%
\pgfsetfillopacity{0.160000}%
\pgfsetlinewidth{0.000000pt}%
\definecolor{currentstroke}{rgb}{0.000000,0.000000,0.000000}%
\pgfsetstrokecolor{currentstroke}%
\pgfsetdash{}{0pt}%
\pgfpathmoveto{\pgfqpoint{0.947680in}{0.772375in}}%
\pgfpathlineto{\pgfqpoint{0.941858in}{0.778985in}}%
\pgfpathlineto{\pgfqpoint{0.936010in}{0.785625in}}%
\pgfpathlineto{\pgfqpoint{0.960084in}{0.782167in}}%
\pgfpathlineto{\pgfqpoint{0.984796in}{0.779391in}}%
\pgfpathlineto{\pgfqpoint{0.988801in}{0.772976in}}%
\pgfpathlineto{\pgfqpoint{0.992788in}{0.766591in}}%
\pgfpathlineto{\pgfqpoint{0.969942in}{0.769166in}}%
\pgfpathlineto{\pgfqpoint{0.947680in}{0.772375in}}%
\pgfpathclose%
\pgfusepath{fill}%
\end{pgfscope}%
\begin{pgfscope}%
\pgfpathrectangle{\pgfqpoint{0.407592in}{0.222904in}}{\pgfqpoint{1.411163in}{1.411163in}}%
\pgfusepath{clip}%
\pgfsetbuttcap%
\pgfsetroundjoin%
\definecolor{currentfill}{rgb}{0.749020,0.772549,0.792157}%
\pgfsetfillcolor{currentfill}%
\pgfsetfillopacity{0.160000}%
\pgfsetlinewidth{0.000000pt}%
\definecolor{currentstroke}{rgb}{0.000000,0.000000,0.000000}%
\pgfsetstrokecolor{currentstroke}%
\pgfsetdash{}{0pt}%
\pgfpathmoveto{\pgfqpoint{1.512741in}{0.933965in}}%
\pgfpathlineto{\pgfqpoint{1.524581in}{0.943028in}}%
\pgfpathlineto{\pgfqpoint{1.536460in}{0.952122in}}%
\pgfpathlineto{\pgfqpoint{1.556155in}{0.964496in}}%
\pgfpathlineto{\pgfqpoint{1.543797in}{0.955108in}}%
\pgfpathlineto{\pgfqpoint{1.531478in}{0.945749in}}%
\pgfusepath{fill}%
\end{pgfscope}%
\begin{pgfscope}%
\pgfpathrectangle{\pgfqpoint{0.407592in}{0.222904in}}{\pgfqpoint{1.411163in}{1.411163in}}%
\pgfusepath{clip}%
\pgfsetbuttcap%
\pgfsetroundjoin%
\definecolor{currentfill}{rgb}{0.749020,0.772549,0.792157}%
\pgfsetfillcolor{currentfill}%
\pgfsetfillopacity{0.160000}%
\pgfsetlinewidth{0.000000pt}%
\definecolor{currentstroke}{rgb}{0.000000,0.000000,0.000000}%
\pgfsetstrokecolor{currentstroke}%
\pgfsetdash{}{0pt}%
\pgfpathmoveto{\pgfqpoint{1.580986in}{0.983361in}}%
\pgfpathlineto{\pgfqpoint{1.593461in}{0.992838in}}%
\pgfpathlineto{\pgfqpoint{1.605974in}{1.002344in}}%
\pgfpathlineto{\pgfqpoint{1.625047in}{1.016311in}}%
\pgfpathlineto{\pgfqpoint{1.641536in}{1.030631in}}%
\pgfpathlineto{\pgfqpoint{1.628276in}{1.020526in}}%
\pgfpathlineto{\pgfqpoint{1.615051in}{1.010449in}}%
\pgfpathlineto{\pgfqpoint{1.599243in}{0.996733in}}%
\pgfpathlineto{\pgfqpoint{1.580986in}{0.983361in}}%
\pgfpathclose%
\pgfusepath{fill}%
\end{pgfscope}%
\begin{pgfscope}%
\pgfpathrectangle{\pgfqpoint{0.407592in}{0.222904in}}{\pgfqpoint{1.411163in}{1.411163in}}%
\pgfusepath{clip}%
\pgfsetbuttcap%
\pgfsetroundjoin%
\definecolor{currentfill}{rgb}{0.749020,0.772549,0.792157}%
\pgfsetfillcolor{currentfill}%
\pgfsetfillopacity{0.160000}%
\pgfsetlinewidth{0.000000pt}%
\definecolor{currentstroke}{rgb}{0.000000,0.000000,0.000000}%
\pgfsetstrokecolor{currentstroke}%
\pgfsetdash{}{0pt}%
\pgfpathmoveto{\pgfqpoint{1.362974in}{0.844625in}}%
\pgfpathlineto{\pgfqpoint{1.372061in}{0.852424in}}%
\pgfpathlineto{\pgfqpoint{1.381184in}{0.860254in}}%
\pgfpathlineto{\pgfqpoint{1.404763in}{0.868896in}}%
\pgfpathlineto{\pgfqpoint{1.426782in}{0.878103in}}%
\pgfpathlineto{\pgfqpoint{1.416284in}{0.869757in}}%
\pgfpathlineto{\pgfqpoint{1.405825in}{0.861442in}}%
\pgfpathlineto{\pgfqpoint{1.385126in}{0.852766in}}%
\pgfpathlineto{\pgfqpoint{1.362974in}{0.844625in}}%
\pgfpathclose%
\pgfusepath{fill}%
\end{pgfscope}%
\begin{pgfscope}%
\pgfpathrectangle{\pgfqpoint{0.407592in}{0.222904in}}{\pgfqpoint{1.411163in}{1.411163in}}%
\pgfusepath{clip}%
\pgfsetbuttcap%
\pgfsetroundjoin%
\definecolor{currentfill}{rgb}{0.749020,0.772549,0.792157}%
\pgfsetfillcolor{currentfill}%
\pgfsetfillopacity{0.160000}%
\pgfsetlinewidth{0.000000pt}%
\definecolor{currentstroke}{rgb}{0.000000,0.000000,0.000000}%
\pgfsetstrokecolor{currentstroke}%
\pgfsetdash{}{0pt}%
\pgfpathmoveto{\pgfqpoint{1.239130in}{0.790652in}}%
\pgfpathlineto{\pgfqpoint{1.244795in}{0.797504in}}%
\pgfpathlineto{\pgfqpoint{1.250484in}{0.804387in}}%
\pgfpathlineto{\pgfqpoint{1.275602in}{0.809626in}}%
\pgfpathlineto{\pgfqpoint{1.299788in}{0.815517in}}%
\pgfpathlineto{\pgfqpoint{1.292371in}{0.808263in}}%
\pgfpathlineto{\pgfqpoint{1.284986in}{0.801039in}}%
\pgfpathlineto{\pgfqpoint{1.262487in}{0.795541in}}%
\pgfpathlineto{\pgfqpoint{1.239130in}{0.790652in}}%
\pgfpathclose%
\pgfusepath{fill}%
\end{pgfscope}%
\begin{pgfscope}%
\pgfpathrectangle{\pgfqpoint{0.407592in}{0.222904in}}{\pgfqpoint{1.411163in}{1.411163in}}%
\pgfusepath{clip}%
\pgfsetbuttcap%
\pgfsetroundjoin%
\definecolor{currentfill}{rgb}{0.749020,0.772549,0.792157}%
\pgfsetfillcolor{currentfill}%
\pgfsetfillopacity{0.160000}%
\pgfsetlinewidth{0.000000pt}%
\definecolor{currentstroke}{rgb}{0.000000,0.000000,0.000000}%
\pgfsetstrokecolor{currentstroke}%
\pgfsetdash{}{0pt}%
\pgfpathmoveto{\pgfqpoint{1.135691in}{0.764998in}}%
\pgfpathlineto{\pgfqpoint{1.137552in}{0.771321in}}%
\pgfpathlineto{\pgfqpoint{1.139422in}{0.777674in}}%
\pgfpathlineto{\pgfqpoint{1.165049in}{0.779892in}}%
\pgfpathlineto{\pgfqpoint{1.190292in}{0.782806in}}%
\pgfpathlineto{\pgfqpoint{1.186494in}{0.776268in}}%
\pgfpathlineto{\pgfqpoint{1.182714in}{0.769759in}}%
\pgfpathlineto{\pgfqpoint{1.159378in}{0.767056in}}%
\pgfpathlineto{\pgfqpoint{1.135691in}{0.764998in}}%
\pgfpathclose%
\pgfusepath{fill}%
\end{pgfscope}%
\begin{pgfscope}%
\pgfpathrectangle{\pgfqpoint{0.407592in}{0.222904in}}{\pgfqpoint{1.411163in}{1.411163in}}%
\pgfusepath{clip}%
\pgfsetbuttcap%
\pgfsetroundjoin%
\definecolor{currentfill}{rgb}{0.749020,0.772549,0.792157}%
\pgfsetfillcolor{currentfill}%
\pgfsetfillopacity{0.160000}%
\pgfsetlinewidth{0.000000pt}%
\definecolor{currentstroke}{rgb}{0.000000,0.000000,0.000000}%
\pgfsetstrokecolor{currentstroke}%
\pgfsetdash{}{0pt}%
\pgfpathmoveto{\pgfqpoint{0.544749in}{1.007673in}}%
\pgfpathlineto{\pgfqpoint{0.530992in}{1.017690in}}%
\pgfpathlineto{\pgfqpoint{0.517197in}{1.027733in}}%
\pgfpathlineto{\pgfqpoint{0.530828in}{1.013479in}}%
\pgfpathlineto{\pgfqpoint{0.547119in}{0.999591in}}%
\pgfpathlineto{\pgfqpoint{0.560308in}{0.990143in}}%
\pgfpathlineto{\pgfqpoint{0.573455in}{0.980725in}}%
\pgfpathlineto{\pgfqpoint{0.557841in}{0.994021in}}%
\pgfpathlineto{\pgfqpoint{0.544749in}{1.007673in}}%
\pgfpathclose%
\pgfusepath{fill}%
\end{pgfscope}%
\begin{pgfscope}%
\pgfpathrectangle{\pgfqpoint{0.407592in}{0.222904in}}{\pgfqpoint{1.411163in}{1.411163in}}%
\pgfusepath{clip}%
\pgfsetbuttcap%
\pgfsetroundjoin%
\definecolor{currentfill}{rgb}{0.749020,0.772549,0.792157}%
\pgfsetfillcolor{currentfill}%
\pgfsetfillopacity{0.160000}%
\pgfsetlinewidth{0.000000pt}%
\definecolor{currentstroke}{rgb}{0.000000,0.000000,0.000000}%
\pgfsetstrokecolor{currentstroke}%
\pgfsetdash{}{0pt}%
\pgfpathmoveto{\pgfqpoint{1.299788in}{0.815517in}}%
\pgfpathlineto{\pgfqpoint{1.307236in}{0.822803in}}%
\pgfpathlineto{\pgfqpoint{1.314716in}{0.830120in}}%
\pgfpathlineto{\pgfqpoint{1.339469in}{0.837062in}}%
\pgfpathlineto{\pgfqpoint{1.362974in}{0.844625in}}%
\pgfpathlineto{\pgfqpoint{1.353924in}{0.836857in}}%
\pgfpathlineto{\pgfqpoint{1.344910in}{0.829121in}}%
\pgfpathlineto{\pgfqpoint{1.322926in}{0.822027in}}%
\pgfpathlineto{\pgfqpoint{1.299788in}{0.815517in}}%
\pgfpathclose%
\pgfusepath{fill}%
\end{pgfscope}%
\begin{pgfscope}%
\pgfpathrectangle{\pgfqpoint{0.407592in}{0.222904in}}{\pgfqpoint{1.411163in}{1.411163in}}%
\pgfusepath{clip}%
\pgfsetbuttcap%
\pgfsetroundjoin%
\definecolor{currentfill}{rgb}{0.749020,0.772549,0.792157}%
\pgfsetfillcolor{currentfill}%
\pgfsetfillopacity{0.160000}%
\pgfsetlinewidth{0.000000pt}%
\definecolor{currentstroke}{rgb}{0.000000,0.000000,0.000000}%
\pgfsetstrokecolor{currentstroke}%
\pgfsetdash{}{0pt}%
\pgfpathmoveto{\pgfqpoint{1.721861in}{1.091839in}}%
\pgfpathlineto{\pgfqpoint{1.735377in}{1.102138in}}%
\pgfpathlineto{\pgfqpoint{1.748930in}{1.112466in}}%
\pgfpathlineto{\pgfqpoint{1.764885in}{1.129441in}}%
\pgfpathlineto{\pgfqpoint{1.777633in}{1.146626in}}%
\pgfpathlineto{\pgfqpoint{1.763644in}{1.135716in}}%
\pgfpathlineto{\pgfqpoint{1.749686in}{1.124831in}}%
\pgfpathlineto{\pgfqpoint{1.737307in}{1.108230in}}%
\pgfpathlineto{\pgfqpoint{1.721861in}{1.091839in}}%
\pgfpathclose%
\pgfusepath{fill}%
\end{pgfscope}%
\begin{pgfscope}%
\pgfpathrectangle{\pgfqpoint{0.407592in}{0.222904in}}{\pgfqpoint{1.411163in}{1.411163in}}%
\pgfusepath{clip}%
\pgfsetbuttcap%
\pgfsetroundjoin%
\definecolor{currentfill}{rgb}{0.749020,0.772549,0.792157}%
\pgfsetfillcolor{currentfill}%
\pgfsetfillopacity{0.160000}%
\pgfsetlinewidth{0.000000pt}%
\definecolor{currentstroke}{rgb}{0.000000,0.000000,0.000000}%
\pgfsetstrokecolor{currentstroke}%
\pgfsetdash{}{0pt}%
\pgfpathmoveto{\pgfqpoint{0.992788in}{0.766591in}}%
\pgfpathlineto{\pgfqpoint{0.988801in}{0.772976in}}%
\pgfpathlineto{\pgfqpoint{0.984796in}{0.779391in}}%
\pgfpathlineto{\pgfqpoint{1.010022in}{0.777314in}}%
\pgfpathlineto{\pgfqpoint{1.035637in}{0.775947in}}%
\pgfpathlineto{\pgfqpoint{1.037713in}{0.769657in}}%
\pgfpathlineto{\pgfqpoint{1.039780in}{0.763396in}}%
\pgfpathlineto{\pgfqpoint{1.016105in}{0.764664in}}%
\pgfpathlineto{\pgfqpoint{0.992788in}{0.766591in}}%
\pgfpathclose%
\pgfusepath{fill}%
\end{pgfscope}%
\begin{pgfscope}%
\pgfpathrectangle{\pgfqpoint{0.407592in}{0.222904in}}{\pgfqpoint{1.411163in}{1.411163in}}%
\pgfusepath{clip}%
\pgfsetbuttcap%
\pgfsetroundjoin%
\definecolor{currentfill}{rgb}{0.749020,0.772549,0.792157}%
\pgfsetfillcolor{currentfill}%
\pgfsetfillopacity{0.160000}%
\pgfsetlinewidth{0.000000pt}%
\definecolor{currentstroke}{rgb}{0.000000,0.000000,0.000000}%
\pgfsetstrokecolor{currentstroke}%
\pgfsetdash{}{0pt}%
\pgfpathmoveto{\pgfqpoint{1.087730in}{0.762862in}}%
\pgfpathlineto{\pgfqpoint{1.087628in}{0.769101in}}%
\pgfpathlineto{\pgfqpoint{1.087525in}{0.775370in}}%
\pgfpathlineto{\pgfqpoint{1.113537in}{0.776164in}}%
\pgfpathlineto{\pgfqpoint{1.139422in}{0.777674in}}%
\pgfpathlineto{\pgfqpoint{1.137552in}{0.771321in}}%
\pgfpathlineto{\pgfqpoint{1.135691in}{0.764998in}}%
\pgfpathlineto{\pgfqpoint{1.111769in}{0.763598in}}%
\pgfpathlineto{\pgfqpoint{1.087730in}{0.762862in}}%
\pgfpathclose%
\pgfusepath{fill}%
\end{pgfscope}%
\begin{pgfscope}%
\pgfpathrectangle{\pgfqpoint{0.407592in}{0.222904in}}{\pgfqpoint{1.411163in}{1.411163in}}%
\pgfusepath{clip}%
\pgfsetbuttcap%
\pgfsetroundjoin%
\definecolor{currentfill}{rgb}{0.749020,0.772549,0.792157}%
\pgfsetfillcolor{currentfill}%
\pgfsetfillopacity{0.160000}%
\pgfsetlinewidth{0.000000pt}%
\definecolor{currentstroke}{rgb}{0.000000,0.000000,0.000000}%
\pgfsetstrokecolor{currentstroke}%
\pgfsetdash{}{0pt}%
\pgfpathmoveto{\pgfqpoint{1.447894in}{0.894888in}}%
\pgfpathlineto{\pgfqpoint{1.458509in}{0.903327in}}%
\pgfpathlineto{\pgfqpoint{1.469162in}{0.911797in}}%
\pgfpathlineto{\pgfqpoint{1.491942in}{0.922627in}}%
\pgfpathlineto{\pgfqpoint{1.512741in}{0.933965in}}%
\pgfpathlineto{\pgfqpoint{1.500940in}{0.924931in}}%
\pgfpathlineto{\pgfqpoint{1.489178in}{0.915928in}}%
\pgfpathlineto{\pgfqpoint{1.469466in}{0.905165in}}%
\pgfpathlineto{\pgfqpoint{1.447894in}{0.894888in}}%
\pgfpathclose%
\pgfusepath{fill}%
\end{pgfscope}%
\begin{pgfscope}%
\pgfpathrectangle{\pgfqpoint{0.407592in}{0.222904in}}{\pgfqpoint{1.411163in}{1.411163in}}%
\pgfusepath{clip}%
\pgfsetbuttcap%
\pgfsetroundjoin%
\definecolor{currentfill}{rgb}{0.749020,0.772549,0.792157}%
\pgfsetfillcolor{currentfill}%
\pgfsetfillopacity{0.160000}%
\pgfsetlinewidth{0.000000pt}%
\definecolor{currentstroke}{rgb}{0.000000,0.000000,0.000000}%
\pgfsetstrokecolor{currentstroke}%
\pgfsetdash{}{0pt}%
\pgfpathmoveto{\pgfqpoint{1.039780in}{0.763396in}}%
\pgfpathlineto{\pgfqpoint{1.037713in}{0.769657in}}%
\pgfpathlineto{\pgfqpoint{1.035637in}{0.775947in}}%
\pgfpathlineto{\pgfqpoint{1.061515in}{0.775298in}}%
\pgfpathlineto{\pgfqpoint{1.087525in}{0.775370in}}%
\pgfpathlineto{\pgfqpoint{1.087628in}{0.769101in}}%
\pgfpathlineto{\pgfqpoint{1.087730in}{0.762862in}}%
\pgfpathlineto{\pgfqpoint{1.063694in}{0.762794in}}%
\pgfpathlineto{\pgfqpoint{1.039780in}{0.763396in}}%
\pgfpathclose%
\pgfusepath{fill}%
\end{pgfscope}%
\begin{pgfscope}%
\pgfpathrectangle{\pgfqpoint{0.407592in}{0.222904in}}{\pgfqpoint{1.411163in}{1.411163in}}%
\pgfusepath{clip}%
\pgfsetbuttcap%
\pgfsetroundjoin%
\definecolor{currentfill}{rgb}{0.749020,0.772549,0.792157}%
\pgfsetfillcolor{currentfill}%
\pgfsetfillopacity{0.160000}%
\pgfsetlinewidth{0.000000pt}%
\definecolor{currentstroke}{rgb}{0.000000,0.000000,0.000000}%
\pgfsetstrokecolor{currentstroke}%
\pgfsetdash{}{0pt}%
\pgfpathmoveto{\pgfqpoint{0.708603in}{0.885836in}}%
\pgfpathlineto{\pgfqpoint{0.696684in}{0.894434in}}%
\pgfpathlineto{\pgfqpoint{0.684723in}{0.903063in}}%
\pgfpathlineto{\pgfqpoint{0.704319in}{0.892890in}}%
\pgfpathlineto{\pgfqpoint{0.725793in}{0.883268in}}%
\pgfpathlineto{\pgfqpoint{0.736606in}{0.875176in}}%
\pgfpathlineto{\pgfqpoint{0.747376in}{0.867116in}}%
\pgfpathlineto{\pgfqpoint{0.727111in}{0.876213in}}%
\pgfpathlineto{\pgfqpoint{0.708603in}{0.885836in}}%
\pgfpathclose%
\pgfusepath{fill}%
\end{pgfscope}%
\begin{pgfscope}%
\pgfpathrectangle{\pgfqpoint{0.407592in}{0.222904in}}{\pgfqpoint{1.411163in}{1.411163in}}%
\pgfusepath{clip}%
\pgfsetbuttcap%
\pgfsetroundjoin%
\definecolor{currentfill}{rgb}{0.749020,0.772549,0.792157}%
\pgfsetfillcolor{currentfill}%
\pgfsetfillopacity{0.160000}%
\pgfsetlinewidth{0.000000pt}%
\definecolor{currentstroke}{rgb}{0.000000,0.000000,0.000000}%
\pgfsetstrokecolor{currentstroke}%
\pgfsetdash{}{0pt}%
\pgfpathmoveto{\pgfqpoint{0.412558in}{1.121487in}}%
\pgfpathlineto{\pgfqpoint{0.398199in}{1.132314in}}%
\pgfpathlineto{\pgfqpoint{0.383806in}{1.143166in}}%
\pgfpathlineto{\pgfqpoint{0.392998in}{1.126016in}}%
\pgfpathlineto{\pgfqpoint{0.405447in}{1.109093in}}%
\pgfpathlineto{\pgfqpoint{0.419552in}{1.098824in}}%
\pgfpathlineto{\pgfqpoint{0.433618in}{1.088583in}}%
\pgfpathlineto{\pgfqpoint{0.421529in}{1.104922in}}%
\pgfpathlineto{\pgfqpoint{0.412558in}{1.121487in}}%
\pgfpathclose%
\pgfusepath{fill}%
\end{pgfscope}%
\begin{pgfscope}%
\pgfpathrectangle{\pgfqpoint{0.407592in}{0.222904in}}{\pgfqpoint{1.411163in}{1.411163in}}%
\pgfusepath{clip}%
\pgfsetbuttcap%
\pgfsetroundjoin%
\definecolor{currentfill}{rgb}{0.749020,0.772549,0.792157}%
\pgfsetfillcolor{currentfill}%
\pgfsetfillopacity{0.160000}%
\pgfsetlinewidth{0.000000pt}%
\definecolor{currentstroke}{rgb}{0.000000,0.000000,0.000000}%
\pgfsetstrokecolor{currentstroke}%
\pgfsetdash{}{0pt}%
\pgfpathmoveto{\pgfqpoint{0.890235in}{0.794510in}}%
\pgfpathlineto{\pgfqpoint{0.882636in}{0.801501in}}%
\pgfpathlineto{\pgfqpoint{0.875004in}{0.808522in}}%
\pgfpathlineto{\pgfqpoint{0.899148in}{0.803417in}}%
\pgfpathlineto{\pgfqpoint{0.924231in}{0.798997in}}%
\pgfpathlineto{\pgfqpoint{0.930134in}{0.792296in}}%
\pgfpathlineto{\pgfqpoint{0.936010in}{0.785625in}}%
\pgfpathlineto{\pgfqpoint{0.912689in}{0.789747in}}%
\pgfpathlineto{\pgfqpoint{0.890235in}{0.794510in}}%
\pgfpathclose%
\pgfusepath{fill}%
\end{pgfscope}%
\begin{pgfscope}%
\pgfpathrectangle{\pgfqpoint{0.407592in}{0.222904in}}{\pgfqpoint{1.411163in}{1.411163in}}%
\pgfusepath{clip}%
\pgfsetbuttcap%
\pgfsetroundjoin%
\definecolor{currentfill}{rgb}{0.749020,0.772549,0.792157}%
\pgfsetfillcolor{currentfill}%
\pgfsetfillopacity{0.160000}%
\pgfsetlinewidth{0.000000pt}%
\definecolor{currentstroke}{rgb}{0.000000,0.000000,0.000000}%
\pgfsetstrokecolor{currentstroke}%
\pgfsetdash{}{0pt}%
\pgfpathmoveto{\pgfqpoint{0.625626in}{0.943352in}}%
\pgfpathlineto{\pgfqpoint{0.612645in}{0.952651in}}%
\pgfpathlineto{\pgfqpoint{0.599623in}{0.961980in}}%
\pgfpathlineto{\pgfqpoint{0.616895in}{0.949694in}}%
\pgfpathlineto{\pgfqpoint{0.636452in}{0.937885in}}%
\pgfpathlineto{\pgfqpoint{0.648584in}{0.929133in}}%
\pgfpathlineto{\pgfqpoint{0.660673in}{0.920412in}}%
\pgfpathlineto{\pgfqpoint{0.642073in}{0.931653in}}%
\pgfpathlineto{\pgfqpoint{0.625626in}{0.943352in}}%
\pgfpathclose%
\pgfusepath{fill}%
\end{pgfscope}%
\begin{pgfscope}%
\pgfpathrectangle{\pgfqpoint{0.407592in}{0.222904in}}{\pgfqpoint{1.411163in}{1.411163in}}%
\pgfusepath{clip}%
\pgfsetbuttcap%
\pgfsetroundjoin%
\definecolor{currentfill}{rgb}{0.749020,0.772549,0.792157}%
\pgfsetfillcolor{currentfill}%
\pgfsetfillopacity{0.160000}%
\pgfsetlinewidth{0.000000pt}%
\definecolor{currentstroke}{rgb}{0.000000,0.000000,0.000000}%
\pgfsetstrokecolor{currentstroke}%
\pgfsetdash{}{0pt}%
\pgfpathmoveto{\pgfqpoint{1.605974in}{1.002344in}}%
\pgfpathlineto{\pgfqpoint{1.618527in}{1.011881in}}%
\pgfpathlineto{\pgfqpoint{1.631120in}{1.021447in}}%
\pgfpathlineto{\pgfqpoint{1.651003in}{1.036004in}}%
\pgfpathlineto{\pgfqpoint{1.668165in}{1.050923in}}%
\pgfpathlineto{\pgfqpoint{1.654833in}{1.040763in}}%
\pgfpathlineto{\pgfqpoint{1.641536in}{1.030631in}}%
\pgfpathlineto{\pgfqpoint{1.625047in}{1.016311in}}%
\pgfpathlineto{\pgfqpoint{1.605974in}{1.002344in}}%
\pgfpathclose%
\pgfusepath{fill}%
\end{pgfscope}%
\begin{pgfscope}%
\pgfpathrectangle{\pgfqpoint{0.407592in}{0.222904in}}{\pgfqpoint{1.411163in}{1.411163in}}%
\pgfusepath{clip}%
\pgfsetbuttcap%
\pgfsetroundjoin%
\definecolor{currentfill}{rgb}{0.749020,0.772549,0.792157}%
\pgfsetfillcolor{currentfill}%
\pgfsetfillopacity{0.160000}%
\pgfsetlinewidth{0.000000pt}%
\definecolor{currentstroke}{rgb}{0.000000,0.000000,0.000000}%
\pgfsetstrokecolor{currentstroke}%
\pgfsetdash{}{0pt}%
\pgfpathmoveto{\pgfqpoint{1.536460in}{0.952122in}}%
\pgfpathlineto{\pgfqpoint{1.548380in}{0.961246in}}%
\pgfpathlineto{\pgfqpoint{1.560339in}{0.970401in}}%
\pgfpathlineto{\pgfqpoint{1.580986in}{0.983361in}}%
\pgfpathlineto{\pgfqpoint{1.568551in}{0.973914in}}%
\pgfpathlineto{\pgfqpoint{1.556155in}{0.964496in}}%
\pgfusepath{fill}%
\end{pgfscope}%
\begin{pgfscope}%
\pgfpathrectangle{\pgfqpoint{0.407592in}{0.222904in}}{\pgfqpoint{1.411163in}{1.411163in}}%
\pgfusepath{clip}%
\pgfsetbuttcap%
\pgfsetroundjoin%
\definecolor{currentfill}{rgb}{0.749020,0.772549,0.792157}%
\pgfsetfillcolor{currentfill}%
\pgfsetfillopacity{0.160000}%
\pgfsetlinewidth{0.000000pt}%
\definecolor{currentstroke}{rgb}{0.000000,0.000000,0.000000}%
\pgfsetstrokecolor{currentstroke}%
\pgfsetdash{}{0pt}%
\pgfpathmoveto{\pgfqpoint{1.190292in}{0.782806in}}%
\pgfpathlineto{\pgfqpoint{1.194108in}{0.789375in}}%
\pgfpathlineto{\pgfqpoint{1.197941in}{0.795975in}}%
\pgfpathlineto{\pgfqpoint{1.224555in}{0.799828in}}%
\pgfpathlineto{\pgfqpoint{1.250484in}{0.804387in}}%
\pgfpathlineto{\pgfqpoint{1.244795in}{0.797504in}}%
\pgfpathlineto{\pgfqpoint{1.239130in}{0.790652in}}%
\pgfpathlineto{\pgfqpoint{1.215026in}{0.786399in}}%
\pgfpathlineto{\pgfqpoint{1.190292in}{0.782806in}}%
\pgfpathclose%
\pgfusepath{fill}%
\end{pgfscope}%
\begin{pgfscope}%
\pgfpathrectangle{\pgfqpoint{0.407592in}{0.222904in}}{\pgfqpoint{1.411163in}{1.411163in}}%
\pgfusepath{clip}%
\pgfsetbuttcap%
\pgfsetroundjoin%
\definecolor{currentfill}{rgb}{0.749020,0.772549,0.792157}%
\pgfsetfillcolor{currentfill}%
\pgfsetfillopacity{0.160000}%
\pgfsetlinewidth{0.000000pt}%
\definecolor{currentstroke}{rgb}{0.000000,0.000000,0.000000}%
\pgfsetstrokecolor{currentstroke}%
\pgfsetdash{}{0pt}%
\pgfpathmoveto{\pgfqpoint{0.768792in}{0.851089in}}%
\pgfpathlineto{\pgfqpoint{0.758105in}{0.859087in}}%
\pgfpathlineto{\pgfqpoint{0.747376in}{0.867116in}}%
\pgfpathlineto{\pgfqpoint{0.769312in}{0.858592in}}%
\pgfpathlineto{\pgfqpoint{0.792823in}{0.850689in}}%
\pgfpathlineto{\pgfqpoint{0.802175in}{0.843138in}}%
\pgfpathlineto{\pgfqpoint{0.811487in}{0.835619in}}%
\pgfpathlineto{\pgfqpoint{0.789407in}{0.843060in}}%
\pgfpathlineto{\pgfqpoint{0.768792in}{0.851089in}}%
\pgfpathclose%
\pgfusepath{fill}%
\end{pgfscope}%
\begin{pgfscope}%
\pgfpathrectangle{\pgfqpoint{0.407592in}{0.222904in}}{\pgfqpoint{1.411163in}{1.411163in}}%
\pgfusepath{clip}%
\pgfsetbuttcap%
\pgfsetroundjoin%
\definecolor{currentfill}{rgb}{0.749020,0.772549,0.792157}%
\pgfsetfillcolor{currentfill}%
\pgfsetfillopacity{0.160000}%
\pgfsetlinewidth{0.000000pt}%
\definecolor{currentstroke}{rgb}{0.000000,0.000000,0.000000}%
\pgfsetstrokecolor{currentstroke}%
\pgfsetdash{}{0pt}%
\pgfpathmoveto{\pgfqpoint{0.829995in}{0.820673in}}%
\pgfpathlineto{\pgfqpoint{0.820760in}{0.828130in}}%
\pgfpathlineto{\pgfqpoint{0.811487in}{0.835619in}}%
\pgfpathlineto{\pgfqpoint{0.834933in}{0.828805in}}%
\pgfpathlineto{\pgfqpoint{0.859638in}{0.822658in}}%
\pgfpathlineto{\pgfqpoint{0.867338in}{0.815574in}}%
\pgfpathlineto{\pgfqpoint{0.875004in}{0.808522in}}%
\pgfpathlineto{\pgfqpoint{0.851917in}{0.814285in}}%
\pgfpathlineto{\pgfqpoint{0.829995in}{0.820673in}}%
\pgfpathclose%
\pgfusepath{fill}%
\end{pgfscope}%
\begin{pgfscope}%
\pgfpathrectangle{\pgfqpoint{0.407592in}{0.222904in}}{\pgfqpoint{1.411163in}{1.411163in}}%
\pgfusepath{clip}%
\pgfsetbuttcap%
\pgfsetroundjoin%
\definecolor{currentfill}{rgb}{0.749020,0.772549,0.792157}%
\pgfsetfillcolor{currentfill}%
\pgfsetfillopacity{0.160000}%
\pgfsetlinewidth{0.000000pt}%
\definecolor{currentstroke}{rgb}{0.000000,0.000000,0.000000}%
\pgfsetstrokecolor{currentstroke}%
\pgfsetdash{}{0pt}%
\pgfpathmoveto{\pgfqpoint{1.748930in}{1.112466in}}%
\pgfpathlineto{\pgfqpoint{1.762521in}{1.122822in}}%
\pgfpathlineto{\pgfqpoint{1.776149in}{1.133206in}}%
\pgfpathlineto{\pgfqpoint{1.792602in}{1.150760in}}%
\pgfpathlineto{\pgfqpoint{1.805709in}{1.168522in}}%
\pgfpathlineto{\pgfqpoint{1.791655in}{1.157562in}}%
\pgfpathlineto{\pgfqpoint{1.777633in}{1.146626in}}%
\pgfpathlineto{\pgfqpoint{1.764885in}{1.129441in}}%
\pgfpathlineto{\pgfqpoint{1.748930in}{1.112466in}}%
\pgfpathclose%
\pgfusepath{fill}%
\end{pgfscope}%
\begin{pgfscope}%
\pgfpathrectangle{\pgfqpoint{0.407592in}{0.222904in}}{\pgfqpoint{1.411163in}{1.411163in}}%
\pgfusepath{clip}%
\pgfsetbuttcap%
\pgfsetroundjoin%
\definecolor{currentfill}{rgb}{0.749020,0.772549,0.792157}%
\pgfsetfillcolor{currentfill}%
\pgfsetfillopacity{0.160000}%
\pgfsetlinewidth{0.000000pt}%
\definecolor{currentstroke}{rgb}{0.000000,0.000000,0.000000}%
\pgfsetstrokecolor{currentstroke}%
\pgfsetdash{}{0pt}%
\pgfpathmoveto{\pgfqpoint{0.517197in}{1.027733in}}%
\pgfpathlineto{\pgfqpoint{0.503363in}{1.037805in}}%
\pgfpathlineto{\pgfqpoint{0.489491in}{1.047904in}}%
\pgfpathlineto{\pgfqpoint{0.503654in}{1.033053in}}%
\pgfpathlineto{\pgfqpoint{0.520615in}{1.018577in}}%
\pgfpathlineto{\pgfqpoint{0.533889in}{1.009069in}}%
\pgfpathlineto{\pgfqpoint{0.547119in}{0.999591in}}%
\pgfpathlineto{\pgfqpoint{0.530828in}{1.013479in}}%
\pgfpathlineto{\pgfqpoint{0.517197in}{1.027733in}}%
\pgfpathclose%
\pgfusepath{fill}%
\end{pgfscope}%
\begin{pgfscope}%
\pgfpathrectangle{\pgfqpoint{0.407592in}{0.222904in}}{\pgfqpoint{1.411163in}{1.411163in}}%
\pgfusepath{clip}%
\pgfsetbuttcap%
\pgfsetroundjoin%
\definecolor{currentfill}{rgb}{0.749020,0.772549,0.792157}%
\pgfsetfillcolor{currentfill}%
\pgfsetfillopacity{0.160000}%
\pgfsetlinewidth{0.000000pt}%
\definecolor{currentstroke}{rgb}{0.000000,0.000000,0.000000}%
\pgfsetstrokecolor{currentstroke}%
\pgfsetdash{}{0pt}%
\pgfpathmoveto{\pgfqpoint{1.381184in}{0.860254in}}%
\pgfpathlineto{\pgfqpoint{1.390343in}{0.868115in}}%
\pgfpathlineto{\pgfqpoint{1.399540in}{0.876009in}}%
\pgfpathlineto{\pgfqpoint{1.424552in}{0.885152in}}%
\pgfpathlineto{\pgfqpoint{1.447894in}{0.894888in}}%
\pgfpathlineto{\pgfqpoint{1.437319in}{0.886480in}}%
\pgfpathlineto{\pgfqpoint{1.426782in}{0.878103in}}%
\pgfpathlineto{\pgfqpoint{1.404763in}{0.868896in}}%
\pgfpathlineto{\pgfqpoint{1.381184in}{0.860254in}}%
\pgfpathclose%
\pgfusepath{fill}%
\end{pgfscope}%
\begin{pgfscope}%
\pgfpathrectangle{\pgfqpoint{0.407592in}{0.222904in}}{\pgfqpoint{1.411163in}{1.411163in}}%
\pgfusepath{clip}%
\pgfsetbuttcap%
\pgfsetroundjoin%
\definecolor{currentfill}{rgb}{0.749020,0.772549,0.792157}%
\pgfsetfillcolor{currentfill}%
\pgfsetfillopacity{0.160000}%
\pgfsetlinewidth{0.000000pt}%
\definecolor{currentstroke}{rgb}{0.000000,0.000000,0.000000}%
\pgfsetstrokecolor{currentstroke}%
\pgfsetdash{}{0pt}%
\pgfpathmoveto{\pgfqpoint{0.936010in}{0.785625in}}%
\pgfpathlineto{\pgfqpoint{0.930134in}{0.792296in}}%
\pgfpathlineto{\pgfqpoint{0.924231in}{0.798997in}}%
\pgfpathlineto{\pgfqpoint{0.950134in}{0.795289in}}%
\pgfpathlineto{\pgfqpoint{0.976727in}{0.792313in}}%
\pgfpathlineto{\pgfqpoint{0.980771in}{0.785837in}}%
\pgfpathlineto{\pgfqpoint{0.984796in}{0.779391in}}%
\pgfpathlineto{\pgfqpoint{0.960084in}{0.782167in}}%
\pgfpathlineto{\pgfqpoint{0.936010in}{0.785625in}}%
\pgfpathclose%
\pgfusepath{fill}%
\end{pgfscope}%
\begin{pgfscope}%
\pgfpathrectangle{\pgfqpoint{0.407592in}{0.222904in}}{\pgfqpoint{1.411163in}{1.411163in}}%
\pgfusepath{clip}%
\pgfsetbuttcap%
\pgfsetroundjoin%
\definecolor{currentfill}{rgb}{0.749020,0.772549,0.792157}%
\pgfsetfillcolor{currentfill}%
\pgfsetfillopacity{0.160000}%
\pgfsetlinewidth{0.000000pt}%
\definecolor{currentstroke}{rgb}{0.000000,0.000000,0.000000}%
\pgfsetstrokecolor{currentstroke}%
\pgfsetdash{}{0pt}%
\pgfpathmoveto{\pgfqpoint{0.383806in}{1.143166in}}%
\pgfpathlineto{\pgfqpoint{0.369379in}{1.154043in}}%
\pgfpathlineto{\pgfqpoint{0.354919in}{1.164946in}}%
\pgfpathlineto{\pgfqpoint{0.364320in}{1.147219in}}%
\pgfpathlineto{\pgfqpoint{0.377119in}{1.129718in}}%
\pgfpathlineto{\pgfqpoint{0.391303in}{1.119391in}}%
\pgfpathlineto{\pgfqpoint{0.405447in}{1.109093in}}%
\pgfpathlineto{\pgfqpoint{0.392998in}{1.126016in}}%
\pgfpathlineto{\pgfqpoint{0.383806in}{1.143166in}}%
\pgfpathclose%
\pgfusepath{fill}%
\end{pgfscope}%
\begin{pgfscope}%
\pgfpathrectangle{\pgfqpoint{0.407592in}{0.222904in}}{\pgfqpoint{1.411163in}{1.411163in}}%
\pgfusepath{clip}%
\pgfsetbuttcap%
\pgfsetroundjoin%
\definecolor{currentfill}{rgb}{0.749020,0.772549,0.792157}%
\pgfsetfillcolor{currentfill}%
\pgfsetfillopacity{0.160000}%
\pgfsetlinewidth{0.000000pt}%
\definecolor{currentstroke}{rgb}{0.000000,0.000000,0.000000}%
\pgfsetstrokecolor{currentstroke}%
\pgfsetdash{}{0pt}%
\pgfpathmoveto{\pgfqpoint{1.250484in}{0.804387in}}%
\pgfpathlineto{\pgfqpoint{1.256200in}{0.811300in}}%
\pgfpathlineto{\pgfqpoint{1.261941in}{0.818245in}}%
\pgfpathlineto{\pgfqpoint{1.288832in}{0.823836in}}%
\pgfpathlineto{\pgfqpoint{1.314716in}{0.830120in}}%
\pgfpathlineto{\pgfqpoint{1.307236in}{0.822803in}}%
\pgfpathlineto{\pgfqpoint{1.299788in}{0.815517in}}%
\pgfpathlineto{\pgfqpoint{1.275602in}{0.809626in}}%
\pgfpathlineto{\pgfqpoint{1.250484in}{0.804387in}}%
\pgfpathclose%
\pgfusepath{fill}%
\end{pgfscope}%
\begin{pgfscope}%
\pgfpathrectangle{\pgfqpoint{0.407592in}{0.222904in}}{\pgfqpoint{1.411163in}{1.411163in}}%
\pgfusepath{clip}%
\pgfsetbuttcap%
\pgfsetroundjoin%
\definecolor{currentfill}{rgb}{0.749020,0.772549,0.792157}%
\pgfsetfillcolor{currentfill}%
\pgfsetfillopacity{0.160000}%
\pgfsetlinewidth{0.000000pt}%
\definecolor{currentstroke}{rgb}{0.000000,0.000000,0.000000}%
\pgfsetstrokecolor{currentstroke}%
\pgfsetdash{}{0pt}%
\pgfpathmoveto{\pgfqpoint{1.314716in}{0.830120in}}%
\pgfpathlineto{\pgfqpoint{1.322228in}{0.837468in}}%
\pgfpathlineto{\pgfqpoint{1.329773in}{0.844848in}}%
\pgfpathlineto{\pgfqpoint{1.356149in}{0.852223in}}%
\pgfpathlineto{\pgfqpoint{1.381184in}{0.860254in}}%
\pgfpathlineto{\pgfqpoint{1.372061in}{0.852424in}}%
\pgfpathlineto{\pgfqpoint{1.362974in}{0.844625in}}%
\pgfpathlineto{\pgfqpoint{1.339469in}{0.837062in}}%
\pgfpathlineto{\pgfqpoint{1.314716in}{0.830120in}}%
\pgfpathclose%
\pgfusepath{fill}%
\end{pgfscope}%
\begin{pgfscope}%
\pgfpathrectangle{\pgfqpoint{0.407592in}{0.222904in}}{\pgfqpoint{1.411163in}{1.411163in}}%
\pgfusepath{clip}%
\pgfsetbuttcap%
\pgfsetroundjoin%
\definecolor{currentfill}{rgb}{0.749020,0.772549,0.792157}%
\pgfsetfillcolor{currentfill}%
\pgfsetfillopacity{0.160000}%
\pgfsetlinewidth{0.000000pt}%
\definecolor{currentstroke}{rgb}{0.000000,0.000000,0.000000}%
\pgfsetstrokecolor{currentstroke}%
\pgfsetdash{}{0pt}%
\pgfpathmoveto{\pgfqpoint{1.139422in}{0.777674in}}%
\pgfpathlineto{\pgfqpoint{1.141300in}{0.784057in}}%
\pgfpathlineto{\pgfqpoint{1.143188in}{0.790471in}}%
\pgfpathlineto{\pgfqpoint{1.170774in}{0.792850in}}%
\pgfpathlineto{\pgfqpoint{1.197941in}{0.795975in}}%
\pgfpathlineto{\pgfqpoint{1.194108in}{0.789375in}}%
\pgfpathlineto{\pgfqpoint{1.190292in}{0.782806in}}%
\pgfpathlineto{\pgfqpoint{1.165049in}{0.779892in}}%
\pgfpathlineto{\pgfqpoint{1.139422in}{0.777674in}}%
\pgfpathclose%
\pgfusepath{fill}%
\end{pgfscope}%
\begin{pgfscope}%
\pgfpathrectangle{\pgfqpoint{0.407592in}{0.222904in}}{\pgfqpoint{1.411163in}{1.411163in}}%
\pgfusepath{clip}%
\pgfsetbuttcap%
\pgfsetroundjoin%
\definecolor{currentfill}{rgb}{0.749020,0.772549,0.792157}%
\pgfsetfillcolor{currentfill}%
\pgfsetfillopacity{0.160000}%
\pgfsetlinewidth{0.000000pt}%
\definecolor{currentstroke}{rgb}{0.000000,0.000000,0.000000}%
\pgfsetstrokecolor{currentstroke}%
\pgfsetdash{}{0pt}%
\pgfpathmoveto{\pgfqpoint{1.469162in}{0.911797in}}%
\pgfpathlineto{\pgfqpoint{1.479856in}{0.920299in}}%
\pgfpathlineto{\pgfqpoint{1.490589in}{0.928832in}}%
\pgfpathlineto{\pgfqpoint{1.514577in}{0.940213in}}%
\pgfpathlineto{\pgfqpoint{1.536460in}{0.952122in}}%
\pgfpathlineto{\pgfqpoint{1.524581in}{0.943028in}}%
\pgfpathlineto{\pgfqpoint{1.512741in}{0.933965in}}%
\pgfpathlineto{\pgfqpoint{1.491942in}{0.922627in}}%
\pgfpathlineto{\pgfqpoint{1.469162in}{0.911797in}}%
\pgfpathclose%
\pgfusepath{fill}%
\end{pgfscope}%
\begin{pgfscope}%
\pgfpathrectangle{\pgfqpoint{0.407592in}{0.222904in}}{\pgfqpoint{1.411163in}{1.411163in}}%
\pgfusepath{clip}%
\pgfsetbuttcap%
\pgfsetroundjoin%
\definecolor{currentfill}{rgb}{0.749020,0.772549,0.792157}%
\pgfsetfillcolor{currentfill}%
\pgfsetfillopacity{0.160000}%
\pgfsetlinewidth{0.000000pt}%
\definecolor{currentstroke}{rgb}{0.000000,0.000000,0.000000}%
\pgfsetstrokecolor{currentstroke}%
\pgfsetdash{}{0pt}%
\pgfpathmoveto{\pgfqpoint{0.684723in}{0.903063in}}%
\pgfpathlineto{\pgfqpoint{0.672719in}{0.911722in}}%
\pgfpathlineto{\pgfqpoint{0.660673in}{0.920412in}}%
\pgfpathlineto{\pgfqpoint{0.681356in}{0.909691in}}%
\pgfpathlineto{\pgfqpoint{0.704041in}{0.899547in}}%
\pgfpathlineto{\pgfqpoint{0.714938in}{0.891391in}}%
\pgfpathlineto{\pgfqpoint{0.725793in}{0.883268in}}%
\pgfpathlineto{\pgfqpoint{0.704319in}{0.892890in}}%
\pgfpathlineto{\pgfqpoint{0.684723in}{0.903063in}}%
\pgfpathclose%
\pgfusepath{fill}%
\end{pgfscope}%
\begin{pgfscope}%
\pgfpathrectangle{\pgfqpoint{0.407592in}{0.222904in}}{\pgfqpoint{1.411163in}{1.411163in}}%
\pgfusepath{clip}%
\pgfsetbuttcap%
\pgfsetroundjoin%
\definecolor{currentfill}{rgb}{0.749020,0.772549,0.792157}%
\pgfsetfillcolor{currentfill}%
\pgfsetfillopacity{0.160000}%
\pgfsetlinewidth{0.000000pt}%
\definecolor{currentstroke}{rgb}{0.000000,0.000000,0.000000}%
\pgfsetstrokecolor{currentstroke}%
\pgfsetdash{}{0pt}%
\pgfpathmoveto{\pgfqpoint{0.599623in}{0.961980in}}%
\pgfpathlineto{\pgfqpoint{0.586560in}{0.971338in}}%
\pgfpathlineto{\pgfqpoint{0.573455in}{0.980725in}}%
\pgfpathlineto{\pgfqpoint{0.591547in}{0.967857in}}%
\pgfpathlineto{\pgfqpoint{0.612057in}{0.955483in}}%
\pgfpathlineto{\pgfqpoint{0.624276in}{0.946668in}}%
\pgfpathlineto{\pgfqpoint{0.636452in}{0.937885in}}%
\pgfpathlineto{\pgfqpoint{0.616895in}{0.949694in}}%
\pgfpathlineto{\pgfqpoint{0.599623in}{0.961980in}}%
\pgfpathclose%
\pgfusepath{fill}%
\end{pgfscope}%
\begin{pgfscope}%
\pgfpathrectangle{\pgfqpoint{0.407592in}{0.222904in}}{\pgfqpoint{1.411163in}{1.411163in}}%
\pgfusepath{clip}%
\pgfsetbuttcap%
\pgfsetroundjoin%
\definecolor{currentfill}{rgb}{0.749020,0.772549,0.792157}%
\pgfsetfillcolor{currentfill}%
\pgfsetfillopacity{0.160000}%
\pgfsetlinewidth{0.000000pt}%
\definecolor{currentstroke}{rgb}{0.000000,0.000000,0.000000}%
\pgfsetstrokecolor{currentstroke}%
\pgfsetdash{}{0pt}%
\pgfpathmoveto{\pgfqpoint{1.631120in}{1.021447in}}%
\pgfpathlineto{\pgfqpoint{1.643752in}{1.031044in}}%
\pgfpathlineto{\pgfqpoint{1.656425in}{1.040672in}}%
\pgfpathlineto{\pgfqpoint{1.677111in}{1.055814in}}%
\pgfpathlineto{\pgfqpoint{1.694940in}{1.071325in}}%
\pgfpathlineto{\pgfqpoint{1.681534in}{1.061110in}}%
\pgfpathlineto{\pgfqpoint{1.668165in}{1.050923in}}%
\pgfpathlineto{\pgfqpoint{1.651003in}{1.036004in}}%
\pgfpathlineto{\pgfqpoint{1.631120in}{1.021447in}}%
\pgfpathclose%
\pgfusepath{fill}%
\end{pgfscope}%
\begin{pgfscope}%
\pgfpathrectangle{\pgfqpoint{0.407592in}{0.222904in}}{\pgfqpoint{1.411163in}{1.411163in}}%
\pgfusepath{clip}%
\pgfsetbuttcap%
\pgfsetroundjoin%
\definecolor{currentfill}{rgb}{0.749020,0.772549,0.792157}%
\pgfsetfillcolor{currentfill}%
\pgfsetfillopacity{0.160000}%
\pgfsetlinewidth{0.000000pt}%
\definecolor{currentstroke}{rgb}{0.000000,0.000000,0.000000}%
\pgfsetstrokecolor{currentstroke}%
\pgfsetdash{}{0pt}%
\pgfpathmoveto{\pgfqpoint{0.984796in}{0.779391in}}%
\pgfpathlineto{\pgfqpoint{0.980771in}{0.785837in}}%
\pgfpathlineto{\pgfqpoint{0.976727in}{0.792313in}}%
\pgfpathlineto{\pgfqpoint{1.003880in}{0.790085in}}%
\pgfpathlineto{\pgfqpoint{1.031455in}{0.788619in}}%
\pgfpathlineto{\pgfqpoint{1.033551in}{0.782268in}}%
\pgfpathlineto{\pgfqpoint{1.035637in}{0.775947in}}%
\pgfpathlineto{\pgfqpoint{1.010022in}{0.777314in}}%
\pgfpathlineto{\pgfqpoint{0.984796in}{0.779391in}}%
\pgfpathclose%
\pgfusepath{fill}%
\end{pgfscope}%
\begin{pgfscope}%
\pgfpathrectangle{\pgfqpoint{0.407592in}{0.222904in}}{\pgfqpoint{1.411163in}{1.411163in}}%
\pgfusepath{clip}%
\pgfsetbuttcap%
\pgfsetroundjoin%
\definecolor{currentfill}{rgb}{0.749020,0.772549,0.792157}%
\pgfsetfillcolor{currentfill}%
\pgfsetfillopacity{0.160000}%
\pgfsetlinewidth{0.000000pt}%
\definecolor{currentstroke}{rgb}{0.000000,0.000000,0.000000}%
\pgfsetstrokecolor{currentstroke}%
\pgfsetdash{}{0pt}%
\pgfpathmoveto{\pgfqpoint{1.776149in}{1.133206in}}%
\pgfpathlineto{\pgfqpoint{1.789814in}{1.143619in}}%
\pgfpathlineto{\pgfqpoint{1.803517in}{1.154061in}}%
\pgfpathlineto{\pgfqpoint{1.820459in}{1.172186in}}%
\pgfpathlineto{\pgfqpoint{1.833914in}{1.190518in}}%
\pgfpathlineto{\pgfqpoint{1.819795in}{1.179507in}}%
\pgfpathlineto{\pgfqpoint{1.805709in}{1.168522in}}%
\pgfpathlineto{\pgfqpoint{1.792602in}{1.150760in}}%
\pgfpathlineto{\pgfqpoint{1.776149in}{1.133206in}}%
\pgfpathclose%
\pgfusepath{fill}%
\end{pgfscope}%
\begin{pgfscope}%
\pgfpathrectangle{\pgfqpoint{0.407592in}{0.222904in}}{\pgfqpoint{1.411163in}{1.411163in}}%
\pgfusepath{clip}%
\pgfsetbuttcap%
\pgfsetroundjoin%
\definecolor{currentfill}{rgb}{0.749020,0.772549,0.792157}%
\pgfsetfillcolor{currentfill}%
\pgfsetfillopacity{0.160000}%
\pgfsetlinewidth{0.000000pt}%
\definecolor{currentstroke}{rgb}{0.000000,0.000000,0.000000}%
\pgfsetstrokecolor{currentstroke}%
\pgfsetdash{}{0pt}%
\pgfpathmoveto{\pgfqpoint{1.087525in}{0.775370in}}%
\pgfpathlineto{\pgfqpoint{1.087422in}{0.781670in}}%
\pgfpathlineto{\pgfqpoint{1.087318in}{0.788000in}}%
\pgfpathlineto{\pgfqpoint{1.115323in}{0.788852in}}%
\pgfpathlineto{\pgfqpoint{1.143188in}{0.790471in}}%
\pgfpathlineto{\pgfqpoint{1.141300in}{0.784057in}}%
\pgfpathlineto{\pgfqpoint{1.139422in}{0.777674in}}%
\pgfpathlineto{\pgfqpoint{1.113537in}{0.776164in}}%
\pgfpathlineto{\pgfqpoint{1.087525in}{0.775370in}}%
\pgfpathclose%
\pgfusepath{fill}%
\end{pgfscope}%
\begin{pgfscope}%
\pgfpathrectangle{\pgfqpoint{0.407592in}{0.222904in}}{\pgfqpoint{1.411163in}{1.411163in}}%
\pgfusepath{clip}%
\pgfsetbuttcap%
\pgfsetroundjoin%
\definecolor{currentfill}{rgb}{0.749020,0.772549,0.792157}%
\pgfsetfillcolor{currentfill}%
\pgfsetfillopacity{0.160000}%
\pgfsetlinewidth{0.000000pt}%
\definecolor{currentstroke}{rgb}{0.000000,0.000000,0.000000}%
\pgfsetstrokecolor{currentstroke}%
\pgfsetdash{}{0pt}%
\pgfpathmoveto{\pgfqpoint{1.560339in}{0.970401in}}%
\pgfpathlineto{\pgfqpoint{1.572338in}{0.979586in}}%
\pgfpathlineto{\pgfqpoint{1.584377in}{0.988802in}}%
\pgfpathlineto{\pgfqpoint{1.605974in}{1.002344in}}%
\pgfpathlineto{\pgfqpoint{1.593461in}{0.992838in}}%
\pgfpathlineto{\pgfqpoint{1.580986in}{0.983361in}}%
\pgfusepath{fill}%
\end{pgfscope}%
\begin{pgfscope}%
\pgfpathrectangle{\pgfqpoint{0.407592in}{0.222904in}}{\pgfqpoint{1.411163in}{1.411163in}}%
\pgfusepath{clip}%
\pgfsetbuttcap%
\pgfsetroundjoin%
\definecolor{currentfill}{rgb}{0.749020,0.772549,0.792157}%
\pgfsetfillcolor{currentfill}%
\pgfsetfillopacity{0.160000}%
\pgfsetlinewidth{0.000000pt}%
\definecolor{currentstroke}{rgb}{0.000000,0.000000,0.000000}%
\pgfsetstrokecolor{currentstroke}%
\pgfsetdash{}{0pt}%
\pgfpathmoveto{\pgfqpoint{1.035637in}{0.775947in}}%
\pgfpathlineto{\pgfqpoint{1.033551in}{0.782268in}}%
\pgfpathlineto{\pgfqpoint{1.031455in}{0.788619in}}%
\pgfpathlineto{\pgfqpoint{1.059315in}{0.787923in}}%
\pgfpathlineto{\pgfqpoint{1.087318in}{0.788000in}}%
\pgfpathlineto{\pgfqpoint{1.087422in}{0.781670in}}%
\pgfpathlineto{\pgfqpoint{1.087525in}{0.775370in}}%
\pgfpathlineto{\pgfqpoint{1.061515in}{0.775298in}}%
\pgfpathlineto{\pgfqpoint{1.035637in}{0.775947in}}%
\pgfpathclose%
\pgfusepath{fill}%
\end{pgfscope}%
\begin{pgfscope}%
\pgfpathrectangle{\pgfqpoint{0.407592in}{0.222904in}}{\pgfqpoint{1.411163in}{1.411163in}}%
\pgfusepath{clip}%
\pgfsetbuttcap%
\pgfsetroundjoin%
\definecolor{currentfill}{rgb}{0.749020,0.772549,0.792157}%
\pgfsetfillcolor{currentfill}%
\pgfsetfillopacity{0.160000}%
\pgfsetlinewidth{0.000000pt}%
\definecolor{currentstroke}{rgb}{0.000000,0.000000,0.000000}%
\pgfsetstrokecolor{currentstroke}%
\pgfsetdash{}{0pt}%
\pgfpathmoveto{\pgfqpoint{0.875004in}{0.808522in}}%
\pgfpathlineto{\pgfqpoint{0.867338in}{0.815574in}}%
\pgfpathlineto{\pgfqpoint{0.859638in}{0.822658in}}%
\pgfpathlineto{\pgfqpoint{0.885483in}{0.817210in}}%
\pgfpathlineto{\pgfqpoint{0.912344in}{0.812494in}}%
\pgfpathlineto{\pgfqpoint{0.918301in}{0.805730in}}%
\pgfpathlineto{\pgfqpoint{0.924231in}{0.798997in}}%
\pgfpathlineto{\pgfqpoint{0.899148in}{0.803417in}}%
\pgfpathlineto{\pgfqpoint{0.875004in}{0.808522in}}%
\pgfpathclose%
\pgfusepath{fill}%
\end{pgfscope}%
\begin{pgfscope}%
\pgfpathrectangle{\pgfqpoint{0.407592in}{0.222904in}}{\pgfqpoint{1.411163in}{1.411163in}}%
\pgfusepath{clip}%
\pgfsetbuttcap%
\pgfsetroundjoin%
\definecolor{currentfill}{rgb}{0.749020,0.772549,0.792157}%
\pgfsetfillcolor{currentfill}%
\pgfsetfillopacity{0.160000}%
\pgfsetlinewidth{0.000000pt}%
\definecolor{currentstroke}{rgb}{0.000000,0.000000,0.000000}%
\pgfsetstrokecolor{currentstroke}%
\pgfsetdash{}{0pt}%
\pgfpathmoveto{\pgfqpoint{0.489491in}{1.047904in}}%
\pgfpathlineto{\pgfqpoint{0.475581in}{1.058032in}}%
\pgfpathlineto{\pgfqpoint{0.461632in}{1.068187in}}%
\pgfpathlineto{\pgfqpoint{0.476317in}{1.052744in}}%
\pgfpathlineto{\pgfqpoint{0.493941in}{1.037685in}}%
\pgfpathlineto{\pgfqpoint{0.507300in}{1.028116in}}%
\pgfpathlineto{\pgfqpoint{0.520615in}{1.018577in}}%
\pgfpathlineto{\pgfqpoint{0.503654in}{1.033053in}}%
\pgfpathlineto{\pgfqpoint{0.489491in}{1.047904in}}%
\pgfpathclose%
\pgfusepath{fill}%
\end{pgfscope}%
\begin{pgfscope}%
\pgfpathrectangle{\pgfqpoint{0.407592in}{0.222904in}}{\pgfqpoint{1.411163in}{1.411163in}}%
\pgfusepath{clip}%
\pgfsetbuttcap%
\pgfsetroundjoin%
\definecolor{currentfill}{rgb}{0.749020,0.772549,0.792157}%
\pgfsetfillcolor{currentfill}%
\pgfsetfillopacity{0.160000}%
\pgfsetlinewidth{0.000000pt}%
\definecolor{currentstroke}{rgb}{0.000000,0.000000,0.000000}%
\pgfsetstrokecolor{currentstroke}%
\pgfsetdash{}{0pt}%
\pgfpathmoveto{\pgfqpoint{0.354919in}{1.164946in}}%
\pgfpathlineto{\pgfqpoint{0.340425in}{1.175874in}}%
\pgfpathlineto{\pgfqpoint{0.325897in}{1.186828in}}%
\pgfpathlineto{\pgfqpoint{0.335494in}{1.168530in}}%
\pgfpathlineto{\pgfqpoint{0.348632in}{1.150458in}}%
\pgfpathlineto{\pgfqpoint{0.362895in}{1.140073in}}%
\pgfpathlineto{\pgfqpoint{0.377119in}{1.129718in}}%
\pgfpathlineto{\pgfqpoint{0.364320in}{1.147219in}}%
\pgfpathlineto{\pgfqpoint{0.354919in}{1.164946in}}%
\pgfpathclose%
\pgfusepath{fill}%
\end{pgfscope}%
\begin{pgfscope}%
\pgfpathrectangle{\pgfqpoint{0.407592in}{0.222904in}}{\pgfqpoint{1.411163in}{1.411163in}}%
\pgfusepath{clip}%
\pgfsetbuttcap%
\pgfsetroundjoin%
\definecolor{currentfill}{rgb}{0.749020,0.772549,0.792157}%
\pgfsetfillcolor{currentfill}%
\pgfsetfillopacity{0.160000}%
\pgfsetlinewidth{0.000000pt}%
\definecolor{currentstroke}{rgb}{0.000000,0.000000,0.000000}%
\pgfsetstrokecolor{currentstroke}%
\pgfsetdash{}{0pt}%
\pgfpathmoveto{\pgfqpoint{0.747376in}{0.867116in}}%
\pgfpathlineto{\pgfqpoint{0.736606in}{0.875176in}}%
\pgfpathlineto{\pgfqpoint{0.725793in}{0.883268in}}%
\pgfpathlineto{\pgfqpoint{0.749055in}{0.874251in}}%
\pgfpathlineto{\pgfqpoint{0.774003in}{0.865887in}}%
\pgfpathlineto{\pgfqpoint{0.783433in}{0.858272in}}%
\pgfpathlineto{\pgfqpoint{0.792823in}{0.850689in}}%
\pgfpathlineto{\pgfqpoint{0.769312in}{0.858592in}}%
\pgfpathlineto{\pgfqpoint{0.747376in}{0.867116in}}%
\pgfpathclose%
\pgfusepath{fill}%
\end{pgfscope}%
\begin{pgfscope}%
\pgfpathrectangle{\pgfqpoint{0.407592in}{0.222904in}}{\pgfqpoint{1.411163in}{1.411163in}}%
\pgfusepath{clip}%
\pgfsetbuttcap%
\pgfsetroundjoin%
\definecolor{currentfill}{rgb}{0.749020,0.772549,0.792157}%
\pgfsetfillcolor{currentfill}%
\pgfsetfillopacity{0.160000}%
\pgfsetlinewidth{0.000000pt}%
\definecolor{currentstroke}{rgb}{0.000000,0.000000,0.000000}%
\pgfsetstrokecolor{currentstroke}%
\pgfsetdash{}{0pt}%
\pgfpathmoveto{\pgfqpoint{0.811487in}{0.835619in}}%
\pgfpathlineto{\pgfqpoint{0.802175in}{0.843138in}}%
\pgfpathlineto{\pgfqpoint{0.792823in}{0.850689in}}%
\pgfpathlineto{\pgfqpoint{0.817803in}{0.843452in}}%
\pgfpathlineto{\pgfqpoint{0.844134in}{0.836919in}}%
\pgfpathlineto{\pgfqpoint{0.851903in}{0.829773in}}%
\pgfpathlineto{\pgfqpoint{0.859638in}{0.822658in}}%
\pgfpathlineto{\pgfqpoint{0.834933in}{0.828805in}}%
\pgfpathlineto{\pgfqpoint{0.811487in}{0.835619in}}%
\pgfpathclose%
\pgfusepath{fill}%
\end{pgfscope}%
\begin{pgfscope}%
\pgfpathrectangle{\pgfqpoint{0.407592in}{0.222904in}}{\pgfqpoint{1.411163in}{1.411163in}}%
\pgfusepath{clip}%
\pgfsetbuttcap%
\pgfsetroundjoin%
\definecolor{currentfill}{rgb}{0.749020,0.772549,0.792157}%
\pgfsetfillcolor{currentfill}%
\pgfsetfillopacity{0.160000}%
\pgfsetlinewidth{0.000000pt}%
\definecolor{currentstroke}{rgb}{0.000000,0.000000,0.000000}%
\pgfsetstrokecolor{currentstroke}%
\pgfsetdash{}{0pt}%
\pgfpathmoveto{\pgfqpoint{1.197941in}{0.795975in}}%
\pgfpathlineto{\pgfqpoint{1.201793in}{0.802606in}}%
\pgfpathlineto{\pgfqpoint{1.205662in}{0.809268in}}%
\pgfpathlineto{\pgfqpoint{1.234172in}{0.813380in}}%
\pgfpathlineto{\pgfqpoint{1.261941in}{0.818245in}}%
\pgfpathlineto{\pgfqpoint{1.256200in}{0.811300in}}%
\pgfpathlineto{\pgfqpoint{1.250484in}{0.804387in}}%
\pgfpathlineto{\pgfqpoint{1.224555in}{0.799828in}}%
\pgfpathlineto{\pgfqpoint{1.197941in}{0.795975in}}%
\pgfpathclose%
\pgfusepath{fill}%
\end{pgfscope}%
\begin{pgfscope}%
\pgfpathrectangle{\pgfqpoint{0.407592in}{0.222904in}}{\pgfqpoint{1.411163in}{1.411163in}}%
\pgfusepath{clip}%
\pgfsetbuttcap%
\pgfsetroundjoin%
\definecolor{currentfill}{rgb}{0.749020,0.772549,0.792157}%
\pgfsetfillcolor{currentfill}%
\pgfsetfillopacity{0.160000}%
\pgfsetlinewidth{0.000000pt}%
\definecolor{currentstroke}{rgb}{0.000000,0.000000,0.000000}%
\pgfsetstrokecolor{currentstroke}%
\pgfsetdash{}{0pt}%
\pgfpathmoveto{\pgfqpoint{1.803517in}{1.154061in}}%
\pgfpathlineto{\pgfqpoint{1.817258in}{1.164532in}}%
\pgfpathlineto{\pgfqpoint{1.834441in}{1.182940in}}%
\pgfpathlineto{\pgfqpoint{1.848065in}{1.201554in}}%
\pgfpathlineto{\pgfqpoint{1.833914in}{1.190518in}}%
\pgfpathlineto{\pgfqpoint{1.820459in}{1.172186in}}%
\pgfusepath{fill}%
\end{pgfscope}%
\begin{pgfscope}%
\pgfpathrectangle{\pgfqpoint{0.407592in}{0.222904in}}{\pgfqpoint{1.411163in}{1.411163in}}%
\pgfusepath{clip}%
\pgfsetbuttcap%
\pgfsetroundjoin%
\definecolor{currentfill}{rgb}{0.749020,0.772549,0.792157}%
\pgfsetfillcolor{currentfill}%
\pgfsetfillopacity{0.160000}%
\pgfsetlinewidth{0.000000pt}%
\definecolor{currentstroke}{rgb}{0.000000,0.000000,0.000000}%
\pgfsetstrokecolor{currentstroke}%
\pgfsetdash{}{0pt}%
\pgfpathmoveto{\pgfqpoint{1.399540in}{0.876009in}}%
\pgfpathlineto{\pgfqpoint{1.408774in}{0.883934in}}%
\pgfpathlineto{\pgfqpoint{1.418045in}{0.891891in}}%
\pgfpathlineto{\pgfqpoint{1.444495in}{0.901533in}}%
\pgfpathlineto{\pgfqpoint{1.469162in}{0.911797in}}%
\pgfpathlineto{\pgfqpoint{1.458509in}{0.903327in}}%
\pgfpathlineto{\pgfqpoint{1.447894in}{0.894888in}}%
\pgfpathlineto{\pgfqpoint{1.424552in}{0.885152in}}%
\pgfpathlineto{\pgfqpoint{1.399540in}{0.876009in}}%
\pgfpathclose%
\pgfusepath{fill}%
\end{pgfscope}%
\begin{pgfscope}%
\pgfpathrectangle{\pgfqpoint{0.407592in}{0.222904in}}{\pgfqpoint{1.411163in}{1.411163in}}%
\pgfusepath{clip}%
\pgfsetbuttcap%
\pgfsetroundjoin%
\definecolor{currentfill}{rgb}{0.749020,0.772549,0.792157}%
\pgfsetfillcolor{currentfill}%
\pgfsetfillopacity{0.160000}%
\pgfsetlinewidth{0.000000pt}%
\definecolor{currentstroke}{rgb}{0.000000,0.000000,0.000000}%
\pgfsetstrokecolor{currentstroke}%
\pgfsetdash{}{0pt}%
\pgfpathmoveto{\pgfqpoint{1.656425in}{1.040672in}}%
\pgfpathlineto{\pgfqpoint{1.669138in}{1.050329in}}%
\pgfpathlineto{\pgfqpoint{1.681891in}{1.060018in}}%
\pgfpathlineto{\pgfqpoint{1.703375in}{1.075741in}}%
\pgfpathlineto{\pgfqpoint{1.721861in}{1.091839in}}%
\pgfpathlineto{\pgfqpoint{1.708382in}{1.081568in}}%
\pgfpathlineto{\pgfqpoint{1.694940in}{1.071325in}}%
\pgfpathlineto{\pgfqpoint{1.677111in}{1.055814in}}%
\pgfpathlineto{\pgfqpoint{1.656425in}{1.040672in}}%
\pgfpathclose%
\pgfusepath{fill}%
\end{pgfscope}%
\begin{pgfscope}%
\pgfpathrectangle{\pgfqpoint{0.407592in}{0.222904in}}{\pgfqpoint{1.411163in}{1.411163in}}%
\pgfusepath{clip}%
\pgfsetbuttcap%
\pgfsetroundjoin%
\definecolor{currentfill}{rgb}{0.749020,0.772549,0.792157}%
\pgfsetfillcolor{currentfill}%
\pgfsetfillopacity{0.160000}%
\pgfsetlinewidth{0.000000pt}%
\definecolor{currentstroke}{rgb}{0.000000,0.000000,0.000000}%
\pgfsetstrokecolor{currentstroke}%
\pgfsetdash{}{0pt}%
\pgfpathmoveto{\pgfqpoint{1.490589in}{0.928832in}}%
\pgfpathlineto{\pgfqpoint{1.501362in}{0.937397in}}%
\pgfpathlineto{\pgfqpoint{1.512175in}{0.945993in}}%
\pgfpathlineto{\pgfqpoint{1.537372in}{0.957923in}}%
\pgfpathlineto{\pgfqpoint{1.560339in}{0.970401in}}%
\pgfpathlineto{\pgfqpoint{1.548380in}{0.961246in}}%
\pgfpathlineto{\pgfqpoint{1.536460in}{0.952122in}}%
\pgfpathlineto{\pgfqpoint{1.514577in}{0.940213in}}%
\pgfpathlineto{\pgfqpoint{1.490589in}{0.928832in}}%
\pgfpathclose%
\pgfusepath{fill}%
\end{pgfscope}%
\begin{pgfscope}%
\pgfpathrectangle{\pgfqpoint{0.407592in}{0.222904in}}{\pgfqpoint{1.411163in}{1.411163in}}%
\pgfusepath{clip}%
\pgfsetbuttcap%
\pgfsetroundjoin%
\definecolor{currentfill}{rgb}{0.749020,0.772549,0.792157}%
\pgfsetfillcolor{currentfill}%
\pgfsetfillopacity{0.160000}%
\pgfsetlinewidth{0.000000pt}%
\definecolor{currentstroke}{rgb}{0.000000,0.000000,0.000000}%
\pgfsetstrokecolor{currentstroke}%
\pgfsetdash{}{0pt}%
\pgfpathmoveto{\pgfqpoint{0.573455in}{0.980725in}}%
\pgfpathlineto{\pgfqpoint{0.560308in}{0.990143in}}%
\pgfpathlineto{\pgfqpoint{0.547119in}{0.999591in}}%
\pgfpathlineto{\pgfqpoint{0.566026in}{0.986143in}}%
\pgfpathlineto{\pgfqpoint{0.587487in}{0.973208in}}%
\pgfpathlineto{\pgfqpoint{0.599794in}{0.964330in}}%
\pgfpathlineto{\pgfqpoint{0.612057in}{0.955483in}}%
\pgfpathlineto{\pgfqpoint{0.591547in}{0.967857in}}%
\pgfpathlineto{\pgfqpoint{0.573455in}{0.980725in}}%
\pgfpathclose%
\pgfusepath{fill}%
\end{pgfscope}%
\begin{pgfscope}%
\pgfpathrectangle{\pgfqpoint{0.407592in}{0.222904in}}{\pgfqpoint{1.411163in}{1.411163in}}%
\pgfusepath{clip}%
\pgfsetbuttcap%
\pgfsetroundjoin%
\definecolor{currentfill}{rgb}{0.749020,0.772549,0.792157}%
\pgfsetfillcolor{currentfill}%
\pgfsetfillopacity{0.160000}%
\pgfsetlinewidth{0.000000pt}%
\definecolor{currentstroke}{rgb}{0.000000,0.000000,0.000000}%
\pgfsetstrokecolor{currentstroke}%
\pgfsetdash{}{0pt}%
\pgfpathmoveto{\pgfqpoint{0.924231in}{0.798997in}}%
\pgfpathlineto{\pgfqpoint{0.918301in}{0.805730in}}%
\pgfpathlineto{\pgfqpoint{0.912344in}{0.812494in}}%
\pgfpathlineto{\pgfqpoint{0.940090in}{0.808535in}}%
\pgfpathlineto{\pgfqpoint{0.968582in}{0.805358in}}%
\pgfpathlineto{\pgfqpoint{0.972665in}{0.798820in}}%
\pgfpathlineto{\pgfqpoint{0.976727in}{0.792313in}}%
\pgfpathlineto{\pgfqpoint{0.950134in}{0.795289in}}%
\pgfpathlineto{\pgfqpoint{0.924231in}{0.798997in}}%
\pgfpathclose%
\pgfusepath{fill}%
\end{pgfscope}%
\begin{pgfscope}%
\pgfpathrectangle{\pgfqpoint{0.407592in}{0.222904in}}{\pgfqpoint{1.411163in}{1.411163in}}%
\pgfusepath{clip}%
\pgfsetbuttcap%
\pgfsetroundjoin%
\definecolor{currentfill}{rgb}{0.749020,0.772549,0.792157}%
\pgfsetfillcolor{currentfill}%
\pgfsetfillopacity{0.160000}%
\pgfsetlinewidth{0.000000pt}%
\definecolor{currentstroke}{rgb}{0.000000,0.000000,0.000000}%
\pgfsetstrokecolor{currentstroke}%
\pgfsetdash{}{0pt}%
\pgfpathmoveto{\pgfqpoint{0.325897in}{1.186828in}}%
\pgfpathlineto{\pgfqpoint{0.311335in}{1.197808in}}%
\pgfpathlineto{\pgfqpoint{0.321025in}{1.179228in}}%
\pgfpathlineto{\pgfqpoint{0.334328in}{1.160872in}}%
\pgfpathlineto{\pgfqpoint{0.348632in}{1.150458in}}%
\pgfpathlineto{\pgfqpoint{0.335494in}{1.168530in}}%
\pgfusepath{fill}%
\end{pgfscope}%
\begin{pgfscope}%
\pgfpathrectangle{\pgfqpoint{0.407592in}{0.222904in}}{\pgfqpoint{1.411163in}{1.411163in}}%
\pgfusepath{clip}%
\pgfsetbuttcap%
\pgfsetroundjoin%
\definecolor{currentfill}{rgb}{0.749020,0.772549,0.792157}%
\pgfsetfillcolor{currentfill}%
\pgfsetfillopacity{0.160000}%
\pgfsetlinewidth{0.000000pt}%
\definecolor{currentstroke}{rgb}{0.000000,0.000000,0.000000}%
\pgfsetstrokecolor{currentstroke}%
\pgfsetdash{}{0pt}%
\pgfpathmoveto{\pgfqpoint{1.261941in}{0.818245in}}%
\pgfpathlineto{\pgfqpoint{1.267708in}{0.825222in}}%
\pgfpathlineto{\pgfqpoint{1.273502in}{0.832230in}}%
\pgfpathlineto{\pgfqpoint{1.302180in}{0.838171in}}%
\pgfpathlineto{\pgfqpoint{1.329773in}{0.844848in}}%
\pgfpathlineto{\pgfqpoint{1.322228in}{0.837468in}}%
\pgfpathlineto{\pgfqpoint{1.314716in}{0.830120in}}%
\pgfpathlineto{\pgfqpoint{1.288832in}{0.823836in}}%
\pgfpathlineto{\pgfqpoint{1.261941in}{0.818245in}}%
\pgfpathclose%
\pgfusepath{fill}%
\end{pgfscope}%
\begin{pgfscope}%
\pgfpathrectangle{\pgfqpoint{0.407592in}{0.222904in}}{\pgfqpoint{1.411163in}{1.411163in}}%
\pgfusepath{clip}%
\pgfsetbuttcap%
\pgfsetroundjoin%
\definecolor{currentfill}{rgb}{0.749020,0.772549,0.792157}%
\pgfsetfillcolor{currentfill}%
\pgfsetfillopacity{0.160000}%
\pgfsetlinewidth{0.000000pt}%
\definecolor{currentstroke}{rgb}{0.000000,0.000000,0.000000}%
\pgfsetstrokecolor{currentstroke}%
\pgfsetdash{}{0pt}%
\pgfpathmoveto{\pgfqpoint{1.329773in}{0.844848in}}%
\pgfpathlineto{\pgfqpoint{1.337350in}{0.852260in}}%
\pgfpathlineto{\pgfqpoint{1.344960in}{0.859704in}}%
\pgfpathlineto{\pgfqpoint{1.372969in}{0.867510in}}%
\pgfpathlineto{\pgfqpoint{1.399540in}{0.876009in}}%
\pgfpathlineto{\pgfqpoint{1.390343in}{0.868115in}}%
\pgfpathlineto{\pgfqpoint{1.381184in}{0.860254in}}%
\pgfpathlineto{\pgfqpoint{1.356149in}{0.852223in}}%
\pgfpathlineto{\pgfqpoint{1.329773in}{0.844848in}}%
\pgfpathclose%
\pgfusepath{fill}%
\end{pgfscope}%
\begin{pgfscope}%
\pgfpathrectangle{\pgfqpoint{0.407592in}{0.222904in}}{\pgfqpoint{1.411163in}{1.411163in}}%
\pgfusepath{clip}%
\pgfsetbuttcap%
\pgfsetroundjoin%
\definecolor{currentfill}{rgb}{0.749020,0.772549,0.792157}%
\pgfsetfillcolor{currentfill}%
\pgfsetfillopacity{0.160000}%
\pgfsetlinewidth{0.000000pt}%
\definecolor{currentstroke}{rgb}{0.000000,0.000000,0.000000}%
\pgfsetstrokecolor{currentstroke}%
\pgfsetdash{}{0pt}%
\pgfpathmoveto{\pgfqpoint{0.660673in}{0.920412in}}%
\pgfpathlineto{\pgfqpoint{0.648584in}{0.929133in}}%
\pgfpathlineto{\pgfqpoint{0.636452in}{0.937885in}}%
\pgfpathlineto{\pgfqpoint{0.658221in}{0.926618in}}%
\pgfpathlineto{\pgfqpoint{0.682118in}{0.915953in}}%
\pgfpathlineto{\pgfqpoint{0.693101in}{0.907734in}}%
\pgfpathlineto{\pgfqpoint{0.704041in}{0.899547in}}%
\pgfpathlineto{\pgfqpoint{0.681356in}{0.909691in}}%
\pgfpathlineto{\pgfqpoint{0.660673in}{0.920412in}}%
\pgfpathclose%
\pgfusepath{fill}%
\end{pgfscope}%
\begin{pgfscope}%
\pgfpathrectangle{\pgfqpoint{0.407592in}{0.222904in}}{\pgfqpoint{1.411163in}{1.411163in}}%
\pgfusepath{clip}%
\pgfsetbuttcap%
\pgfsetroundjoin%
\definecolor{currentfill}{rgb}{0.749020,0.772549,0.792157}%
\pgfsetfillcolor{currentfill}%
\pgfsetfillopacity{0.160000}%
\pgfsetlinewidth{0.000000pt}%
\definecolor{currentstroke}{rgb}{0.000000,0.000000,0.000000}%
\pgfsetstrokecolor{currentstroke}%
\pgfsetdash{}{0pt}%
\pgfpathmoveto{\pgfqpoint{1.584377in}{0.988802in}}%
\pgfpathlineto{\pgfqpoint{1.596457in}{0.998049in}}%
\pgfpathlineto{\pgfqpoint{1.608577in}{1.007328in}}%
\pgfpathlineto{\pgfqpoint{1.631120in}{1.021447in}}%
\pgfpathlineto{\pgfqpoint{1.618527in}{1.011881in}}%
\pgfpathlineto{\pgfqpoint{1.605974in}{1.002344in}}%
\pgfusepath{fill}%
\end{pgfscope}%
\begin{pgfscope}%
\pgfpathrectangle{\pgfqpoint{0.407592in}{0.222904in}}{\pgfqpoint{1.411163in}{1.411163in}}%
\pgfusepath{clip}%
\pgfsetbuttcap%
\pgfsetroundjoin%
\definecolor{currentfill}{rgb}{0.749020,0.772549,0.792157}%
\pgfsetfillcolor{currentfill}%
\pgfsetfillopacity{0.160000}%
\pgfsetlinewidth{0.000000pt}%
\definecolor{currentstroke}{rgb}{0.000000,0.000000,0.000000}%
\pgfsetstrokecolor{currentstroke}%
\pgfsetdash{}{0pt}%
\pgfpathmoveto{\pgfqpoint{0.461632in}{1.068187in}}%
\pgfpathlineto{\pgfqpoint{0.447645in}{1.078371in}}%
\pgfpathlineto{\pgfqpoint{0.433618in}{1.088583in}}%
\pgfpathlineto{\pgfqpoint{0.448815in}{1.072554in}}%
\pgfpathlineto{\pgfqpoint{0.467095in}{1.056916in}}%
\pgfpathlineto{\pgfqpoint{0.480540in}{1.047285in}}%
\pgfpathlineto{\pgfqpoint{0.493941in}{1.037685in}}%
\pgfpathlineto{\pgfqpoint{0.476317in}{1.052744in}}%
\pgfpathlineto{\pgfqpoint{0.461632in}{1.068187in}}%
\pgfpathclose%
\pgfusepath{fill}%
\end{pgfscope}%
\begin{pgfscope}%
\pgfpathrectangle{\pgfqpoint{0.407592in}{0.222904in}}{\pgfqpoint{1.411163in}{1.411163in}}%
\pgfusepath{clip}%
\pgfsetbuttcap%
\pgfsetroundjoin%
\definecolor{currentfill}{rgb}{0.749020,0.772549,0.792157}%
\pgfsetfillcolor{currentfill}%
\pgfsetfillopacity{0.160000}%
\pgfsetlinewidth{0.000000pt}%
\definecolor{currentstroke}{rgb}{0.000000,0.000000,0.000000}%
\pgfsetstrokecolor{currentstroke}%
\pgfsetdash{}{0pt}%
\pgfpathmoveto{\pgfqpoint{1.143188in}{0.790471in}}%
\pgfpathlineto{\pgfqpoint{1.145085in}{0.796916in}}%
\pgfpathlineto{\pgfqpoint{1.146991in}{0.803392in}}%
\pgfpathlineto{\pgfqpoint{1.176553in}{0.805932in}}%
\pgfpathlineto{\pgfqpoint{1.205662in}{0.809268in}}%
\pgfpathlineto{\pgfqpoint{1.201793in}{0.802606in}}%
\pgfpathlineto{\pgfqpoint{1.197941in}{0.795975in}}%
\pgfpathlineto{\pgfqpoint{1.170774in}{0.792850in}}%
\pgfpathlineto{\pgfqpoint{1.143188in}{0.790471in}}%
\pgfpathclose%
\pgfusepath{fill}%
\end{pgfscope}%
\begin{pgfscope}%
\pgfpathrectangle{\pgfqpoint{0.407592in}{0.222904in}}{\pgfqpoint{1.411163in}{1.411163in}}%
\pgfusepath{clip}%
\pgfsetbuttcap%
\pgfsetroundjoin%
\definecolor{currentfill}{rgb}{0.749020,0.772549,0.792157}%
\pgfsetfillcolor{currentfill}%
\pgfsetfillopacity{0.160000}%
\pgfsetlinewidth{0.000000pt}%
\definecolor{currentstroke}{rgb}{0.000000,0.000000,0.000000}%
\pgfsetstrokecolor{currentstroke}%
\pgfsetdash{}{0pt}%
\pgfpathmoveto{\pgfqpoint{0.976727in}{0.792313in}}%
\pgfpathlineto{\pgfqpoint{0.972665in}{0.798820in}}%
\pgfpathlineto{\pgfqpoint{0.968582in}{0.805358in}}%
\pgfpathlineto{\pgfqpoint{0.997679in}{0.802979in}}%
\pgfpathlineto{\pgfqpoint{1.027232in}{0.801414in}}%
\pgfpathlineto{\pgfqpoint{1.029349in}{0.795001in}}%
\pgfpathlineto{\pgfqpoint{1.031455in}{0.788619in}}%
\pgfpathlineto{\pgfqpoint{1.003880in}{0.790085in}}%
\pgfpathlineto{\pgfqpoint{0.976727in}{0.792313in}}%
\pgfpathclose%
\pgfusepath{fill}%
\end{pgfscope}%
\begin{pgfscope}%
\pgfpathrectangle{\pgfqpoint{0.407592in}{0.222904in}}{\pgfqpoint{1.411163in}{1.411163in}}%
\pgfusepath{clip}%
\pgfsetbuttcap%
\pgfsetroundjoin%
\definecolor{currentfill}{rgb}{0.749020,0.772549,0.792157}%
\pgfsetfillcolor{currentfill}%
\pgfsetfillopacity{0.160000}%
\pgfsetlinewidth{0.000000pt}%
\definecolor{currentstroke}{rgb}{0.000000,0.000000,0.000000}%
\pgfsetstrokecolor{currentstroke}%
\pgfsetdash{}{0pt}%
\pgfpathmoveto{\pgfqpoint{1.087318in}{0.788000in}}%
\pgfpathlineto{\pgfqpoint{1.087213in}{0.794362in}}%
\pgfpathlineto{\pgfqpoint{1.087109in}{0.800754in}}%
\pgfpathlineto{\pgfqpoint{1.117125in}{0.801662in}}%
\pgfpathlineto{\pgfqpoint{1.146991in}{0.803392in}}%
\pgfpathlineto{\pgfqpoint{1.145085in}{0.796916in}}%
\pgfpathlineto{\pgfqpoint{1.143188in}{0.790471in}}%
\pgfpathlineto{\pgfqpoint{1.115323in}{0.788852in}}%
\pgfpathlineto{\pgfqpoint{1.087318in}{0.788000in}}%
\pgfpathclose%
\pgfusepath{fill}%
\end{pgfscope}%
\begin{pgfscope}%
\pgfpathrectangle{\pgfqpoint{0.407592in}{0.222904in}}{\pgfqpoint{1.411163in}{1.411163in}}%
\pgfusepath{clip}%
\pgfsetbuttcap%
\pgfsetroundjoin%
\definecolor{currentfill}{rgb}{0.749020,0.772549,0.792157}%
\pgfsetfillcolor{currentfill}%
\pgfsetfillopacity{0.160000}%
\pgfsetlinewidth{0.000000pt}%
\definecolor{currentstroke}{rgb}{0.000000,0.000000,0.000000}%
\pgfsetstrokecolor{currentstroke}%
\pgfsetdash{}{0pt}%
\pgfpathmoveto{\pgfqpoint{0.859638in}{0.822658in}}%
\pgfpathlineto{\pgfqpoint{0.851903in}{0.829773in}}%
\pgfpathlineto{\pgfqpoint{0.844134in}{0.836919in}}%
\pgfpathlineto{\pgfqpoint{0.871693in}{0.831130in}}%
\pgfpathlineto{\pgfqpoint{0.900346in}{0.826116in}}%
\pgfpathlineto{\pgfqpoint{0.906359in}{0.819289in}}%
\pgfpathlineto{\pgfqpoint{0.912344in}{0.812494in}}%
\pgfpathlineto{\pgfqpoint{0.885483in}{0.817210in}}%
\pgfpathlineto{\pgfqpoint{0.859638in}{0.822658in}}%
\pgfpathclose%
\pgfusepath{fill}%
\end{pgfscope}%
\begin{pgfscope}%
\pgfpathrectangle{\pgfqpoint{0.407592in}{0.222904in}}{\pgfqpoint{1.411163in}{1.411163in}}%
\pgfusepath{clip}%
\pgfsetbuttcap%
\pgfsetroundjoin%
\definecolor{currentfill}{rgb}{0.749020,0.772549,0.792157}%
\pgfsetfillcolor{currentfill}%
\pgfsetfillopacity{0.160000}%
\pgfsetlinewidth{0.000000pt}%
\definecolor{currentstroke}{rgb}{0.000000,0.000000,0.000000}%
\pgfsetstrokecolor{currentstroke}%
\pgfsetdash{}{0pt}%
\pgfpathmoveto{\pgfqpoint{0.725793in}{0.883268in}}%
\pgfpathlineto{\pgfqpoint{0.714938in}{0.891391in}}%
\pgfpathlineto{\pgfqpoint{0.704041in}{0.899547in}}%
\pgfpathlineto{\pgfqpoint{0.728632in}{0.890037in}}%
\pgfpathlineto{\pgfqpoint{0.755023in}{0.881214in}}%
\pgfpathlineto{\pgfqpoint{0.764533in}{0.873534in}}%
\pgfpathlineto{\pgfqpoint{0.774003in}{0.865887in}}%
\pgfpathlineto{\pgfqpoint{0.749055in}{0.874251in}}%
\pgfpathlineto{\pgfqpoint{0.725793in}{0.883268in}}%
\pgfpathclose%
\pgfusepath{fill}%
\end{pgfscope}%
\begin{pgfscope}%
\pgfpathrectangle{\pgfqpoint{0.407592in}{0.222904in}}{\pgfqpoint{1.411163in}{1.411163in}}%
\pgfusepath{clip}%
\pgfsetbuttcap%
\pgfsetroundjoin%
\definecolor{currentfill}{rgb}{0.749020,0.772549,0.792157}%
\pgfsetfillcolor{currentfill}%
\pgfsetfillopacity{0.160000}%
\pgfsetlinewidth{0.000000pt}%
\definecolor{currentstroke}{rgb}{0.000000,0.000000,0.000000}%
\pgfsetstrokecolor{currentstroke}%
\pgfsetdash{}{0pt}%
\pgfpathmoveto{\pgfqpoint{1.031455in}{0.788619in}}%
\pgfpathlineto{\pgfqpoint{1.029349in}{0.795001in}}%
\pgfpathlineto{\pgfqpoint{1.027232in}{0.801414in}}%
\pgfpathlineto{\pgfqpoint{1.057093in}{0.800671in}}%
\pgfpathlineto{\pgfqpoint{1.087109in}{0.800754in}}%
\pgfpathlineto{\pgfqpoint{1.087213in}{0.794362in}}%
\pgfpathlineto{\pgfqpoint{1.087318in}{0.788000in}}%
\pgfpathlineto{\pgfqpoint{1.059315in}{0.787923in}}%
\pgfpathlineto{\pgfqpoint{1.031455in}{0.788619in}}%
\pgfpathclose%
\pgfusepath{fill}%
\end{pgfscope}%
\begin{pgfscope}%
\pgfpathrectangle{\pgfqpoint{0.407592in}{0.222904in}}{\pgfqpoint{1.411163in}{1.411163in}}%
\pgfusepath{clip}%
\pgfsetbuttcap%
\pgfsetroundjoin%
\definecolor{currentfill}{rgb}{0.749020,0.772549,0.792157}%
\pgfsetfillcolor{currentfill}%
\pgfsetfillopacity{0.160000}%
\pgfsetlinewidth{0.000000pt}%
\definecolor{currentstroke}{rgb}{0.000000,0.000000,0.000000}%
\pgfsetstrokecolor{currentstroke}%
\pgfsetdash{}{0pt}%
\pgfpathmoveto{\pgfqpoint{1.681891in}{1.060018in}}%
\pgfpathlineto{\pgfqpoint{1.694685in}{1.069737in}}%
\pgfpathlineto{\pgfqpoint{1.707519in}{1.079488in}}%
\pgfpathlineto{\pgfqpoint{1.729795in}{1.095786in}}%
\pgfpathlineto{\pgfqpoint{1.748930in}{1.112466in}}%
\pgfpathlineto{\pgfqpoint{1.735377in}{1.102138in}}%
\pgfpathlineto{\pgfqpoint{1.721861in}{1.091839in}}%
\pgfpathlineto{\pgfqpoint{1.703375in}{1.075741in}}%
\pgfpathlineto{\pgfqpoint{1.681891in}{1.060018in}}%
\pgfpathclose%
\pgfusepath{fill}%
\end{pgfscope}%
\begin{pgfscope}%
\pgfpathrectangle{\pgfqpoint{0.407592in}{0.222904in}}{\pgfqpoint{1.411163in}{1.411163in}}%
\pgfusepath{clip}%
\pgfsetbuttcap%
\pgfsetroundjoin%
\definecolor{currentfill}{rgb}{0.749020,0.772549,0.792157}%
\pgfsetfillcolor{currentfill}%
\pgfsetfillopacity{0.160000}%
\pgfsetlinewidth{0.000000pt}%
\definecolor{currentstroke}{rgb}{0.000000,0.000000,0.000000}%
\pgfsetstrokecolor{currentstroke}%
\pgfsetdash{}{0pt}%
\pgfpathmoveto{\pgfqpoint{1.418045in}{0.891891in}}%
\pgfpathlineto{\pgfqpoint{1.427353in}{0.899881in}}%
\pgfpathlineto{\pgfqpoint{1.436700in}{0.907902in}}%
\pgfpathlineto{\pgfqpoint{1.464593in}{0.918042in}}%
\pgfpathlineto{\pgfqpoint{1.490589in}{0.928832in}}%
\pgfpathlineto{\pgfqpoint{1.479856in}{0.920299in}}%
\pgfpathlineto{\pgfqpoint{1.469162in}{0.911797in}}%
\pgfpathlineto{\pgfqpoint{1.444495in}{0.901533in}}%
\pgfpathlineto{\pgfqpoint{1.418045in}{0.891891in}}%
\pgfpathclose%
\pgfusepath{fill}%
\end{pgfscope}%
\begin{pgfscope}%
\pgfpathrectangle{\pgfqpoint{0.407592in}{0.222904in}}{\pgfqpoint{1.411163in}{1.411163in}}%
\pgfusepath{clip}%
\pgfsetbuttcap%
\pgfsetroundjoin%
\definecolor{currentfill}{rgb}{0.749020,0.772549,0.792157}%
\pgfsetfillcolor{currentfill}%
\pgfsetfillopacity{0.160000}%
\pgfsetlinewidth{0.000000pt}%
\definecolor{currentstroke}{rgb}{0.000000,0.000000,0.000000}%
\pgfsetstrokecolor{currentstroke}%
\pgfsetdash{}{0pt}%
\pgfpathmoveto{\pgfqpoint{0.792823in}{0.850689in}}%
\pgfpathlineto{\pgfqpoint{0.783433in}{0.858272in}}%
\pgfpathlineto{\pgfqpoint{0.774003in}{0.865887in}}%
\pgfpathlineto{\pgfqpoint{0.800523in}{0.858225in}}%
\pgfpathlineto{\pgfqpoint{0.828493in}{0.851308in}}%
\pgfpathlineto{\pgfqpoint{0.836331in}{0.844098in}}%
\pgfpathlineto{\pgfqpoint{0.844134in}{0.836919in}}%
\pgfpathlineto{\pgfqpoint{0.817803in}{0.843452in}}%
\pgfpathlineto{\pgfqpoint{0.792823in}{0.850689in}}%
\pgfpathclose%
\pgfusepath{fill}%
\end{pgfscope}%
\begin{pgfscope}%
\pgfpathrectangle{\pgfqpoint{0.407592in}{0.222904in}}{\pgfqpoint{1.411163in}{1.411163in}}%
\pgfusepath{clip}%
\pgfsetbuttcap%
\pgfsetroundjoin%
\definecolor{currentfill}{rgb}{0.749020,0.772549,0.792157}%
\pgfsetfillcolor{currentfill}%
\pgfsetfillopacity{0.160000}%
\pgfsetlinewidth{0.000000pt}%
\definecolor{currentstroke}{rgb}{0.000000,0.000000,0.000000}%
\pgfsetstrokecolor{currentstroke}%
\pgfsetdash{}{0pt}%
\pgfpathmoveto{\pgfqpoint{1.512175in}{0.945993in}}%
\pgfpathlineto{\pgfqpoint{1.523029in}{0.954622in}}%
\pgfpathlineto{\pgfqpoint{1.533923in}{0.963284in}}%
\pgfpathlineto{\pgfqpoint{1.560330in}{0.975758in}}%
\pgfpathlineto{\pgfqpoint{1.584377in}{0.988802in}}%
\pgfpathlineto{\pgfqpoint{1.572338in}{0.979586in}}%
\pgfpathlineto{\pgfqpoint{1.560339in}{0.970401in}}%
\pgfpathlineto{\pgfqpoint{1.537372in}{0.957923in}}%
\pgfpathlineto{\pgfqpoint{1.512175in}{0.945993in}}%
\pgfpathclose%
\pgfusepath{fill}%
\end{pgfscope}%
\begin{pgfscope}%
\pgfpathrectangle{\pgfqpoint{0.407592in}{0.222904in}}{\pgfqpoint{1.411163in}{1.411163in}}%
\pgfusepath{clip}%
\pgfsetbuttcap%
\pgfsetroundjoin%
\definecolor{currentfill}{rgb}{0.749020,0.772549,0.792157}%
\pgfsetfillcolor{currentfill}%
\pgfsetfillopacity{0.160000}%
\pgfsetlinewidth{0.000000pt}%
\definecolor{currentstroke}{rgb}{0.000000,0.000000,0.000000}%
\pgfsetstrokecolor{currentstroke}%
\pgfsetdash{}{0pt}%
\pgfpathmoveto{\pgfqpoint{1.205662in}{0.809268in}}%
\pgfpathlineto{\pgfqpoint{1.209550in}{0.815961in}}%
\pgfpathlineto{\pgfqpoint{1.213456in}{0.822686in}}%
\pgfpathlineto{\pgfqpoint{1.243878in}{0.827058in}}%
\pgfpathlineto{\pgfqpoint{1.273502in}{0.832230in}}%
\pgfpathlineto{\pgfqpoint{1.267708in}{0.825222in}}%
\pgfpathlineto{\pgfqpoint{1.261941in}{0.818245in}}%
\pgfpathlineto{\pgfqpoint{1.234172in}{0.813380in}}%
\pgfpathlineto{\pgfqpoint{1.205662in}{0.809268in}}%
\pgfpathclose%
\pgfusepath{fill}%
\end{pgfscope}%
\begin{pgfscope}%
\pgfpathrectangle{\pgfqpoint{0.407592in}{0.222904in}}{\pgfqpoint{1.411163in}{1.411163in}}%
\pgfusepath{clip}%
\pgfsetbuttcap%
\pgfsetroundjoin%
\definecolor{currentfill}{rgb}{0.749020,0.772549,0.792157}%
\pgfsetfillcolor{currentfill}%
\pgfsetfillopacity{0.160000}%
\pgfsetlinewidth{0.000000pt}%
\definecolor{currentstroke}{rgb}{0.000000,0.000000,0.000000}%
\pgfsetstrokecolor{currentstroke}%
\pgfsetdash{}{0pt}%
\pgfpathmoveto{\pgfqpoint{0.547119in}{0.999591in}}%
\pgfpathlineto{\pgfqpoint{0.533889in}{1.009069in}}%
\pgfpathlineto{\pgfqpoint{0.520615in}{1.018577in}}%
\pgfpathlineto{\pgfqpoint{0.540331in}{1.004554in}}%
\pgfpathlineto{\pgfqpoint{0.562740in}{0.991060in}}%
\pgfpathlineto{\pgfqpoint{0.575136in}{0.982118in}}%
\pgfpathlineto{\pgfqpoint{0.587487in}{0.973208in}}%
\pgfpathlineto{\pgfqpoint{0.566026in}{0.986143in}}%
\pgfpathlineto{\pgfqpoint{0.547119in}{0.999591in}}%
\pgfpathclose%
\pgfusepath{fill}%
\end{pgfscope}%
\begin{pgfscope}%
\pgfpathrectangle{\pgfqpoint{0.407592in}{0.222904in}}{\pgfqpoint{1.411163in}{1.411163in}}%
\pgfusepath{clip}%
\pgfsetbuttcap%
\pgfsetroundjoin%
\definecolor{currentfill}{rgb}{0.749020,0.772549,0.792157}%
\pgfsetfillcolor{currentfill}%
\pgfsetfillopacity{0.160000}%
\pgfsetlinewidth{0.000000pt}%
\definecolor{currentstroke}{rgb}{0.000000,0.000000,0.000000}%
\pgfsetstrokecolor{currentstroke}%
\pgfsetdash{}{0pt}%
\pgfpathmoveto{\pgfqpoint{0.433618in}{1.088583in}}%
\pgfpathlineto{\pgfqpoint{0.419552in}{1.098824in}}%
\pgfpathlineto{\pgfqpoint{0.405447in}{1.109093in}}%
\pgfpathlineto{\pgfqpoint{0.421148in}{1.092483in}}%
\pgfpathlineto{\pgfqpoint{0.440076in}{1.076272in}}%
\pgfpathlineto{\pgfqpoint{0.453607in}{1.066578in}}%
\pgfpathlineto{\pgfqpoint{0.467095in}{1.056916in}}%
\pgfpathlineto{\pgfqpoint{0.448815in}{1.072554in}}%
\pgfpathlineto{\pgfqpoint{0.433618in}{1.088583in}}%
\pgfpathclose%
\pgfusepath{fill}%
\end{pgfscope}%
\begin{pgfscope}%
\pgfpathrectangle{\pgfqpoint{0.407592in}{0.222904in}}{\pgfqpoint{1.411163in}{1.411163in}}%
\pgfusepath{clip}%
\pgfsetbuttcap%
\pgfsetroundjoin%
\definecolor{currentfill}{rgb}{0.749020,0.772549,0.792157}%
\pgfsetfillcolor{currentfill}%
\pgfsetfillopacity{0.160000}%
\pgfsetlinewidth{0.000000pt}%
\definecolor{currentstroke}{rgb}{0.000000,0.000000,0.000000}%
\pgfsetstrokecolor{currentstroke}%
\pgfsetdash{}{0pt}%
\pgfpathmoveto{\pgfqpoint{1.608577in}{1.007328in}}%
\pgfpathlineto{\pgfqpoint{1.620739in}{1.016637in}}%
\pgfpathlineto{\pgfqpoint{1.632941in}{1.025978in}}%
\pgfpathlineto{\pgfqpoint{1.656425in}{1.040672in}}%
\pgfpathlineto{\pgfqpoint{1.643752in}{1.031044in}}%
\pgfpathlineto{\pgfqpoint{1.631120in}{1.021447in}}%
\pgfusepath{fill}%
\end{pgfscope}%
\begin{pgfscope}%
\pgfpathrectangle{\pgfqpoint{0.407592in}{0.222904in}}{\pgfqpoint{1.411163in}{1.411163in}}%
\pgfusepath{clip}%
\pgfsetbuttcap%
\pgfsetroundjoin%
\definecolor{currentfill}{rgb}{0.749020,0.772549,0.792157}%
\pgfsetfillcolor{currentfill}%
\pgfsetfillopacity{0.160000}%
\pgfsetlinewidth{0.000000pt}%
\definecolor{currentstroke}{rgb}{0.000000,0.000000,0.000000}%
\pgfsetstrokecolor{currentstroke}%
\pgfsetdash{}{0pt}%
\pgfpathmoveto{\pgfqpoint{0.636452in}{0.937885in}}%
\pgfpathlineto{\pgfqpoint{0.624276in}{0.946668in}}%
\pgfpathlineto{\pgfqpoint{0.612057in}{0.955483in}}%
\pgfpathlineto{\pgfqpoint{0.634912in}{0.943672in}}%
\pgfpathlineto{\pgfqpoint{0.660022in}{0.932489in}}%
\pgfpathlineto{\pgfqpoint{0.671092in}{0.924205in}}%
\pgfpathlineto{\pgfqpoint{0.682118in}{0.915953in}}%
\pgfpathlineto{\pgfqpoint{0.658221in}{0.926618in}}%
\pgfpathlineto{\pgfqpoint{0.636452in}{0.937885in}}%
\pgfpathclose%
\pgfusepath{fill}%
\end{pgfscope}%
\begin{pgfscope}%
\pgfpathrectangle{\pgfqpoint{0.407592in}{0.222904in}}{\pgfqpoint{1.411163in}{1.411163in}}%
\pgfusepath{clip}%
\pgfsetbuttcap%
\pgfsetroundjoin%
\definecolor{currentfill}{rgb}{0.749020,0.772549,0.792157}%
\pgfsetfillcolor{currentfill}%
\pgfsetfillopacity{0.160000}%
\pgfsetlinewidth{0.000000pt}%
\definecolor{currentstroke}{rgb}{0.000000,0.000000,0.000000}%
\pgfsetstrokecolor{currentstroke}%
\pgfsetdash{}{0pt}%
\pgfpathmoveto{\pgfqpoint{1.344960in}{0.859704in}}%
\pgfpathlineto{\pgfqpoint{1.352603in}{0.867180in}}%
\pgfpathlineto{\pgfqpoint{1.360279in}{0.874689in}}%
\pgfpathlineto{\pgfqpoint{1.389930in}{0.882926in}}%
\pgfpathlineto{\pgfqpoint{1.418045in}{0.891891in}}%
\pgfpathlineto{\pgfqpoint{1.408774in}{0.883934in}}%
\pgfpathlineto{\pgfqpoint{1.399540in}{0.876009in}}%
\pgfpathlineto{\pgfqpoint{1.372969in}{0.867510in}}%
\pgfpathlineto{\pgfqpoint{1.344960in}{0.859704in}}%
\pgfpathclose%
\pgfusepath{fill}%
\end{pgfscope}%
\begin{pgfscope}%
\pgfpathrectangle{\pgfqpoint{0.407592in}{0.222904in}}{\pgfqpoint{1.411163in}{1.411163in}}%
\pgfusepath{clip}%
\pgfsetbuttcap%
\pgfsetroundjoin%
\definecolor{currentfill}{rgb}{0.749020,0.772549,0.792157}%
\pgfsetfillcolor{currentfill}%
\pgfsetfillopacity{0.160000}%
\pgfsetlinewidth{0.000000pt}%
\definecolor{currentstroke}{rgb}{0.000000,0.000000,0.000000}%
\pgfsetstrokecolor{currentstroke}%
\pgfsetdash{}{0pt}%
\pgfpathmoveto{\pgfqpoint{1.273502in}{0.832230in}}%
\pgfpathlineto{\pgfqpoint{1.279322in}{0.839270in}}%
\pgfpathlineto{\pgfqpoint{1.285168in}{0.846342in}}%
\pgfpathlineto{\pgfqpoint{1.315646in}{0.852634in}}%
\pgfpathlineto{\pgfqpoint{1.344960in}{0.859704in}}%
\pgfpathlineto{\pgfqpoint{1.337350in}{0.852260in}}%
\pgfpathlineto{\pgfqpoint{1.329773in}{0.844848in}}%
\pgfpathlineto{\pgfqpoint{1.302180in}{0.838171in}}%
\pgfpathlineto{\pgfqpoint{1.273502in}{0.832230in}}%
\pgfpathclose%
\pgfusepath{fill}%
\end{pgfscope}%
\begin{pgfscope}%
\pgfpathrectangle{\pgfqpoint{0.407592in}{0.222904in}}{\pgfqpoint{1.411163in}{1.411163in}}%
\pgfusepath{clip}%
\pgfsetbuttcap%
\pgfsetroundjoin%
\definecolor{currentfill}{rgb}{0.749020,0.772549,0.792157}%
\pgfsetfillcolor{currentfill}%
\pgfsetfillopacity{0.160000}%
\pgfsetlinewidth{0.000000pt}%
\definecolor{currentstroke}{rgb}{0.000000,0.000000,0.000000}%
\pgfsetstrokecolor{currentstroke}%
\pgfsetdash{}{0pt}%
\pgfpathmoveto{\pgfqpoint{0.912344in}{0.812494in}}%
\pgfpathlineto{\pgfqpoint{0.906359in}{0.819289in}}%
\pgfpathlineto{\pgfqpoint{0.900346in}{0.826116in}}%
\pgfpathlineto{\pgfqpoint{0.929950in}{0.821907in}}%
\pgfpathlineto{\pgfqpoint{0.960359in}{0.818528in}}%
\pgfpathlineto{\pgfqpoint{0.964480in}{0.811927in}}%
\pgfpathlineto{\pgfqpoint{0.968582in}{0.805358in}}%
\pgfpathlineto{\pgfqpoint{0.940090in}{0.808535in}}%
\pgfpathlineto{\pgfqpoint{0.912344in}{0.812494in}}%
\pgfpathclose%
\pgfusepath{fill}%
\end{pgfscope}%
\begin{pgfscope}%
\pgfpathrectangle{\pgfqpoint{0.407592in}{0.222904in}}{\pgfqpoint{1.411163in}{1.411163in}}%
\pgfusepath{clip}%
\pgfsetbuttcap%
\pgfsetroundjoin%
\definecolor{currentfill}{rgb}{0.749020,0.772549,0.792157}%
\pgfsetfillcolor{currentfill}%
\pgfsetfillopacity{0.160000}%
\pgfsetlinewidth{0.000000pt}%
\definecolor{currentstroke}{rgb}{0.000000,0.000000,0.000000}%
\pgfsetstrokecolor{currentstroke}%
\pgfsetdash{}{0pt}%
\pgfpathmoveto{\pgfqpoint{1.146991in}{0.803392in}}%
\pgfpathlineto{\pgfqpoint{1.148906in}{0.809899in}}%
\pgfpathlineto{\pgfqpoint{1.150830in}{0.816438in}}%
\pgfpathlineto{\pgfqpoint{1.182387in}{0.819139in}}%
\pgfpathlineto{\pgfqpoint{1.213456in}{0.822686in}}%
\pgfpathlineto{\pgfqpoint{1.209550in}{0.815961in}}%
\pgfpathlineto{\pgfqpoint{1.205662in}{0.809268in}}%
\pgfpathlineto{\pgfqpoint{1.176553in}{0.805932in}}%
\pgfpathlineto{\pgfqpoint{1.146991in}{0.803392in}}%
\pgfpathclose%
\pgfusepath{fill}%
\end{pgfscope}%
\begin{pgfscope}%
\pgfpathrectangle{\pgfqpoint{0.407592in}{0.222904in}}{\pgfqpoint{1.411163in}{1.411163in}}%
\pgfusepath{clip}%
\pgfsetbuttcap%
\pgfsetroundjoin%
\definecolor{currentfill}{rgb}{0.749020,0.772549,0.792157}%
\pgfsetfillcolor{currentfill}%
\pgfsetfillopacity{0.160000}%
\pgfsetlinewidth{0.000000pt}%
\definecolor{currentstroke}{rgb}{0.000000,0.000000,0.000000}%
\pgfsetstrokecolor{currentstroke}%
\pgfsetdash{}{0pt}%
\pgfpathmoveto{\pgfqpoint{1.707519in}{1.079488in}}%
\pgfpathlineto{\pgfqpoint{1.720395in}{1.089269in}}%
\pgfpathlineto{\pgfqpoint{1.733311in}{1.099082in}}%
\pgfpathlineto{\pgfqpoint{1.756371in}{1.115950in}}%
\pgfpathlineto{\pgfqpoint{1.776149in}{1.133206in}}%
\pgfpathlineto{\pgfqpoint{1.762521in}{1.122822in}}%
\pgfpathlineto{\pgfqpoint{1.748930in}{1.112466in}}%
\pgfpathlineto{\pgfqpoint{1.729795in}{1.095786in}}%
\pgfpathlineto{\pgfqpoint{1.707519in}{1.079488in}}%
\pgfpathclose%
\pgfusepath{fill}%
\end{pgfscope}%
\begin{pgfscope}%
\pgfpathrectangle{\pgfqpoint{0.407592in}{0.222904in}}{\pgfqpoint{1.411163in}{1.411163in}}%
\pgfusepath{clip}%
\pgfsetbuttcap%
\pgfsetroundjoin%
\definecolor{currentfill}{rgb}{0.749020,0.772549,0.792157}%
\pgfsetfillcolor{currentfill}%
\pgfsetfillopacity{0.160000}%
\pgfsetlinewidth{0.000000pt}%
\definecolor{currentstroke}{rgb}{0.000000,0.000000,0.000000}%
\pgfsetstrokecolor{currentstroke}%
\pgfsetdash{}{0pt}%
\pgfpathmoveto{\pgfqpoint{0.704041in}{0.899547in}}%
\pgfpathlineto{\pgfqpoint{0.693101in}{0.907734in}}%
\pgfpathlineto{\pgfqpoint{0.682118in}{0.915953in}}%
\pgfpathlineto{\pgfqpoint{0.708043in}{0.905952in}}%
\pgfpathlineto{\pgfqpoint{0.735882in}{0.896670in}}%
\pgfpathlineto{\pgfqpoint{0.745472in}{0.888925in}}%
\pgfpathlineto{\pgfqpoint{0.755023in}{0.881214in}}%
\pgfpathlineto{\pgfqpoint{0.728632in}{0.890037in}}%
\pgfpathlineto{\pgfqpoint{0.704041in}{0.899547in}}%
\pgfpathclose%
\pgfusepath{fill}%
\end{pgfscope}%
\begin{pgfscope}%
\pgfpathrectangle{\pgfqpoint{0.407592in}{0.222904in}}{\pgfqpoint{1.411163in}{1.411163in}}%
\pgfusepath{clip}%
\pgfsetbuttcap%
\pgfsetroundjoin%
\definecolor{currentfill}{rgb}{0.749020,0.772549,0.792157}%
\pgfsetfillcolor{currentfill}%
\pgfsetfillopacity{0.160000}%
\pgfsetlinewidth{0.000000pt}%
\definecolor{currentstroke}{rgb}{0.000000,0.000000,0.000000}%
\pgfsetstrokecolor{currentstroke}%
\pgfsetdash{}{0pt}%
\pgfpathmoveto{\pgfqpoint{0.968582in}{0.805358in}}%
\pgfpathlineto{\pgfqpoint{0.964480in}{0.811927in}}%
\pgfpathlineto{\pgfqpoint{0.960359in}{0.818528in}}%
\pgfpathlineto{\pgfqpoint{0.991417in}{0.815999in}}%
\pgfpathlineto{\pgfqpoint{1.022968in}{0.814334in}}%
\pgfpathlineto{\pgfqpoint{1.025106in}{0.807858in}}%
\pgfpathlineto{\pgfqpoint{1.027232in}{0.801414in}}%
\pgfpathlineto{\pgfqpoint{0.997679in}{0.802979in}}%
\pgfpathlineto{\pgfqpoint{0.968582in}{0.805358in}}%
\pgfpathclose%
\pgfusepath{fill}%
\end{pgfscope}%
\begin{pgfscope}%
\pgfpathrectangle{\pgfqpoint{0.407592in}{0.222904in}}{\pgfqpoint{1.411163in}{1.411163in}}%
\pgfusepath{clip}%
\pgfsetbuttcap%
\pgfsetroundjoin%
\definecolor{currentfill}{rgb}{0.749020,0.772549,0.792157}%
\pgfsetfillcolor{currentfill}%
\pgfsetfillopacity{0.160000}%
\pgfsetlinewidth{0.000000pt}%
\definecolor{currentstroke}{rgb}{0.000000,0.000000,0.000000}%
\pgfsetstrokecolor{currentstroke}%
\pgfsetdash{}{0pt}%
\pgfpathmoveto{\pgfqpoint{1.436700in}{0.907902in}}%
\pgfpathlineto{\pgfqpoint{1.446084in}{0.915957in}}%
\pgfpathlineto{\pgfqpoint{1.455507in}{0.924044in}}%
\pgfpathlineto{\pgfqpoint{1.484848in}{0.934679in}}%
\pgfpathlineto{\pgfqpoint{1.512175in}{0.945993in}}%
\pgfpathlineto{\pgfqpoint{1.501362in}{0.937397in}}%
\pgfpathlineto{\pgfqpoint{1.490589in}{0.928832in}}%
\pgfpathlineto{\pgfqpoint{1.464593in}{0.918042in}}%
\pgfpathlineto{\pgfqpoint{1.436700in}{0.907902in}}%
\pgfpathclose%
\pgfusepath{fill}%
\end{pgfscope}%
\begin{pgfscope}%
\pgfpathrectangle{\pgfqpoint{0.407592in}{0.222904in}}{\pgfqpoint{1.411163in}{1.411163in}}%
\pgfusepath{clip}%
\pgfsetbuttcap%
\pgfsetroundjoin%
\definecolor{currentfill}{rgb}{0.749020,0.772549,0.792157}%
\pgfsetfillcolor{currentfill}%
\pgfsetfillopacity{0.160000}%
\pgfsetlinewidth{0.000000pt}%
\definecolor{currentstroke}{rgb}{0.000000,0.000000,0.000000}%
\pgfsetstrokecolor{currentstroke}%
\pgfsetdash{}{0pt}%
\pgfpathmoveto{\pgfqpoint{0.405447in}{1.109093in}}%
\pgfpathlineto{\pgfqpoint{0.391303in}{1.119391in}}%
\pgfpathlineto{\pgfqpoint{0.377119in}{1.129718in}}%
\pgfpathlineto{\pgfqpoint{0.393314in}{1.112532in}}%
\pgfpathlineto{\pgfqpoint{0.412881in}{1.095753in}}%
\pgfpathlineto{\pgfqpoint{0.426500in}{1.085997in}}%
\pgfpathlineto{\pgfqpoint{0.440076in}{1.076272in}}%
\pgfpathlineto{\pgfqpoint{0.421148in}{1.092483in}}%
\pgfpathlineto{\pgfqpoint{0.405447in}{1.109093in}}%
\pgfpathclose%
\pgfusepath{fill}%
\end{pgfscope}%
\begin{pgfscope}%
\pgfpathrectangle{\pgfqpoint{0.407592in}{0.222904in}}{\pgfqpoint{1.411163in}{1.411163in}}%
\pgfusepath{clip}%
\pgfsetbuttcap%
\pgfsetroundjoin%
\definecolor{currentfill}{rgb}{0.749020,0.772549,0.792157}%
\pgfsetfillcolor{currentfill}%
\pgfsetfillopacity{0.160000}%
\pgfsetlinewidth{0.000000pt}%
\definecolor{currentstroke}{rgb}{0.000000,0.000000,0.000000}%
\pgfsetstrokecolor{currentstroke}%
\pgfsetdash{}{0pt}%
\pgfpathmoveto{\pgfqpoint{0.844134in}{0.836919in}}%
\pgfpathlineto{\pgfqpoint{0.836331in}{0.844098in}}%
\pgfpathlineto{\pgfqpoint{0.828493in}{0.851308in}}%
\pgfpathlineto{\pgfqpoint{0.857778in}{0.845177in}}%
\pgfpathlineto{\pgfqpoint{0.888235in}{0.839866in}}%
\pgfpathlineto{\pgfqpoint{0.894305in}{0.832975in}}%
\pgfpathlineto{\pgfqpoint{0.900346in}{0.826116in}}%
\pgfpathlineto{\pgfqpoint{0.871693in}{0.831130in}}%
\pgfpathlineto{\pgfqpoint{0.844134in}{0.836919in}}%
\pgfpathclose%
\pgfusepath{fill}%
\end{pgfscope}%
\begin{pgfscope}%
\pgfpathrectangle{\pgfqpoint{0.407592in}{0.222904in}}{\pgfqpoint{1.411163in}{1.411163in}}%
\pgfusepath{clip}%
\pgfsetbuttcap%
\pgfsetroundjoin%
\definecolor{currentfill}{rgb}{0.749020,0.772549,0.792157}%
\pgfsetfillcolor{currentfill}%
\pgfsetfillopacity{0.160000}%
\pgfsetlinewidth{0.000000pt}%
\definecolor{currentstroke}{rgb}{0.000000,0.000000,0.000000}%
\pgfsetstrokecolor{currentstroke}%
\pgfsetdash{}{0pt}%
\pgfpathmoveto{\pgfqpoint{1.533923in}{0.963284in}}%
\pgfpathlineto{\pgfqpoint{1.544858in}{0.971978in}}%
\pgfpathlineto{\pgfqpoint{1.555834in}{0.980704in}}%
\pgfpathlineto{\pgfqpoint{1.583451in}{0.993721in}}%
\pgfpathlineto{\pgfqpoint{1.608577in}{1.007328in}}%
\pgfpathlineto{\pgfqpoint{1.596457in}{0.998049in}}%
\pgfpathlineto{\pgfqpoint{1.584377in}{0.988802in}}%
\pgfpathlineto{\pgfqpoint{1.560330in}{0.975758in}}%
\pgfpathlineto{\pgfqpoint{1.533923in}{0.963284in}}%
\pgfpathclose%
\pgfusepath{fill}%
\end{pgfscope}%
\begin{pgfscope}%
\pgfpathrectangle{\pgfqpoint{0.407592in}{0.222904in}}{\pgfqpoint{1.411163in}{1.411163in}}%
\pgfusepath{clip}%
\pgfsetbuttcap%
\pgfsetroundjoin%
\definecolor{currentfill}{rgb}{0.749020,0.772549,0.792157}%
\pgfsetfillcolor{currentfill}%
\pgfsetfillopacity{0.160000}%
\pgfsetlinewidth{0.000000pt}%
\definecolor{currentstroke}{rgb}{0.000000,0.000000,0.000000}%
\pgfsetstrokecolor{currentstroke}%
\pgfsetdash{}{0pt}%
\pgfpathmoveto{\pgfqpoint{0.520615in}{1.018577in}}%
\pgfpathlineto{\pgfqpoint{0.507300in}{1.028116in}}%
\pgfpathlineto{\pgfqpoint{0.493941in}{1.037685in}}%
\pgfpathlineto{\pgfqpoint{0.514461in}{1.023091in}}%
\pgfpathlineto{\pgfqpoint{0.537814in}{1.009042in}}%
\pgfpathlineto{\pgfqpoint{0.550299in}{1.000035in}}%
\pgfpathlineto{\pgfqpoint{0.562740in}{0.991060in}}%
\pgfpathlineto{\pgfqpoint{0.540331in}{1.004554in}}%
\pgfpathlineto{\pgfqpoint{0.520615in}{1.018577in}}%
\pgfpathclose%
\pgfusepath{fill}%
\end{pgfscope}%
\begin{pgfscope}%
\pgfpathrectangle{\pgfqpoint{0.407592in}{0.222904in}}{\pgfqpoint{1.411163in}{1.411163in}}%
\pgfusepath{clip}%
\pgfsetbuttcap%
\pgfsetroundjoin%
\definecolor{currentfill}{rgb}{0.749020,0.772549,0.792157}%
\pgfsetfillcolor{currentfill}%
\pgfsetfillopacity{0.160000}%
\pgfsetlinewidth{0.000000pt}%
\definecolor{currentstroke}{rgb}{0.000000,0.000000,0.000000}%
\pgfsetstrokecolor{currentstroke}%
\pgfsetdash{}{0pt}%
\pgfpathmoveto{\pgfqpoint{1.087109in}{0.800754in}}%
\pgfpathlineto{\pgfqpoint{1.087003in}{0.807177in}}%
\pgfpathlineto{\pgfqpoint{1.086897in}{0.813631in}}%
\pgfpathlineto{\pgfqpoint{1.118946in}{0.814598in}}%
\pgfpathlineto{\pgfqpoint{1.150830in}{0.816438in}}%
\pgfpathlineto{\pgfqpoint{1.148906in}{0.809899in}}%
\pgfpathlineto{\pgfqpoint{1.146991in}{0.803392in}}%
\pgfpathlineto{\pgfqpoint{1.117125in}{0.801662in}}%
\pgfpathlineto{\pgfqpoint{1.087109in}{0.800754in}}%
\pgfpathclose%
\pgfusepath{fill}%
\end{pgfscope}%
\begin{pgfscope}%
\pgfpathrectangle{\pgfqpoint{0.407592in}{0.222904in}}{\pgfqpoint{1.411163in}{1.411163in}}%
\pgfusepath{clip}%
\pgfsetbuttcap%
\pgfsetroundjoin%
\definecolor{currentfill}{rgb}{0.749020,0.772549,0.792157}%
\pgfsetfillcolor{currentfill}%
\pgfsetfillopacity{0.160000}%
\pgfsetlinewidth{0.000000pt}%
\definecolor{currentstroke}{rgb}{0.000000,0.000000,0.000000}%
\pgfsetstrokecolor{currentstroke}%
\pgfsetdash{}{0pt}%
\pgfpathmoveto{\pgfqpoint{1.632941in}{1.025978in}}%
\pgfpathlineto{\pgfqpoint{1.645185in}{1.035351in}}%
\pgfpathlineto{\pgfqpoint{1.657470in}{1.044755in}}%
\pgfpathlineto{\pgfqpoint{1.681891in}{1.060018in}}%
\pgfpathlineto{\pgfqpoint{1.669138in}{1.050329in}}%
\pgfpathlineto{\pgfqpoint{1.656425in}{1.040672in}}%
\pgfusepath{fill}%
\end{pgfscope}%
\begin{pgfscope}%
\pgfpathrectangle{\pgfqpoint{0.407592in}{0.222904in}}{\pgfqpoint{1.411163in}{1.411163in}}%
\pgfusepath{clip}%
\pgfsetbuttcap%
\pgfsetroundjoin%
\definecolor{currentfill}{rgb}{0.749020,0.772549,0.792157}%
\pgfsetfillcolor{currentfill}%
\pgfsetfillopacity{0.160000}%
\pgfsetlinewidth{0.000000pt}%
\definecolor{currentstroke}{rgb}{0.000000,0.000000,0.000000}%
\pgfsetstrokecolor{currentstroke}%
\pgfsetdash{}{0pt}%
\pgfpathmoveto{\pgfqpoint{1.027232in}{0.801414in}}%
\pgfpathlineto{\pgfqpoint{1.025106in}{0.807858in}}%
\pgfpathlineto{\pgfqpoint{1.022968in}{0.814334in}}%
\pgfpathlineto{\pgfqpoint{1.054850in}{0.813543in}}%
\pgfpathlineto{\pgfqpoint{1.086897in}{0.813631in}}%
\pgfpathlineto{\pgfqpoint{1.087003in}{0.807177in}}%
\pgfpathlineto{\pgfqpoint{1.087109in}{0.800754in}}%
\pgfpathlineto{\pgfqpoint{1.057093in}{0.800671in}}%
\pgfpathlineto{\pgfqpoint{1.027232in}{0.801414in}}%
\pgfpathclose%
\pgfusepath{fill}%
\end{pgfscope}%
\begin{pgfscope}%
\pgfpathrectangle{\pgfqpoint{0.407592in}{0.222904in}}{\pgfqpoint{1.411163in}{1.411163in}}%
\pgfusepath{clip}%
\pgfsetbuttcap%
\pgfsetroundjoin%
\definecolor{currentfill}{rgb}{0.749020,0.772549,0.792157}%
\pgfsetfillcolor{currentfill}%
\pgfsetfillopacity{0.160000}%
\pgfsetlinewidth{0.000000pt}%
\definecolor{currentstroke}{rgb}{0.000000,0.000000,0.000000}%
\pgfsetstrokecolor{currentstroke}%
\pgfsetdash{}{0pt}%
\pgfpathmoveto{\pgfqpoint{0.774003in}{0.865887in}}%
\pgfpathlineto{\pgfqpoint{0.764533in}{0.873534in}}%
\pgfpathlineto{\pgfqpoint{0.755023in}{0.881214in}}%
\pgfpathlineto{\pgfqpoint{0.783092in}{0.873128in}}%
\pgfpathlineto{\pgfqpoint{0.812710in}{0.865827in}}%
\pgfpathlineto{\pgfqpoint{0.820619in}{0.858551in}}%
\pgfpathlineto{\pgfqpoint{0.828493in}{0.851308in}}%
\pgfpathlineto{\pgfqpoint{0.800523in}{0.858225in}}%
\pgfpathlineto{\pgfqpoint{0.774003in}{0.865887in}}%
\pgfpathclose%
\pgfusepath{fill}%
\end{pgfscope}%
\begin{pgfscope}%
\pgfpathrectangle{\pgfqpoint{0.407592in}{0.222904in}}{\pgfqpoint{1.411163in}{1.411163in}}%
\pgfusepath{clip}%
\pgfsetbuttcap%
\pgfsetroundjoin%
\definecolor{currentfill}{rgb}{0.749020,0.772549,0.792157}%
\pgfsetfillcolor{currentfill}%
\pgfsetfillopacity{0.160000}%
\pgfsetlinewidth{0.000000pt}%
\definecolor{currentstroke}{rgb}{0.000000,0.000000,0.000000}%
\pgfsetstrokecolor{currentstroke}%
\pgfsetdash{}{0pt}%
\pgfpathmoveto{\pgfqpoint{0.612057in}{0.955483in}}%
\pgfpathlineto{\pgfqpoint{0.599794in}{0.964330in}}%
\pgfpathlineto{\pgfqpoint{0.587487in}{0.973208in}}%
\pgfpathlineto{\pgfqpoint{0.611427in}{0.960856in}}%
\pgfpathlineto{\pgfqpoint{0.637751in}{0.949156in}}%
\pgfpathlineto{\pgfqpoint{0.648909in}{0.940806in}}%
\pgfpathlineto{\pgfqpoint{0.660022in}{0.932489in}}%
\pgfpathlineto{\pgfqpoint{0.634912in}{0.943672in}}%
\pgfpathlineto{\pgfqpoint{0.612057in}{0.955483in}}%
\pgfpathclose%
\pgfusepath{fill}%
\end{pgfscope}%
\begin{pgfscope}%
\pgfpathrectangle{\pgfqpoint{0.407592in}{0.222904in}}{\pgfqpoint{1.411163in}{1.411163in}}%
\pgfusepath{clip}%
\pgfsetbuttcap%
\pgfsetroundjoin%
\definecolor{currentfill}{rgb}{0.749020,0.772549,0.792157}%
\pgfsetfillcolor{currentfill}%
\pgfsetfillopacity{0.160000}%
\pgfsetlinewidth{0.000000pt}%
\definecolor{currentstroke}{rgb}{0.000000,0.000000,0.000000}%
\pgfsetstrokecolor{currentstroke}%
\pgfsetdash{}{0pt}%
\pgfpathmoveto{\pgfqpoint{1.213456in}{0.822686in}}%
\pgfpathlineto{\pgfqpoint{1.217380in}{0.829443in}}%
\pgfpathlineto{\pgfqpoint{1.221324in}{0.836231in}}%
\pgfpathlineto{\pgfqpoint{1.253675in}{0.840864in}}%
\pgfpathlineto{\pgfqpoint{1.285168in}{0.846342in}}%
\pgfpathlineto{\pgfqpoint{1.279322in}{0.839270in}}%
\pgfpathlineto{\pgfqpoint{1.273502in}{0.832230in}}%
\pgfpathlineto{\pgfqpoint{1.243878in}{0.827058in}}%
\pgfpathlineto{\pgfqpoint{1.213456in}{0.822686in}}%
\pgfpathclose%
\pgfusepath{fill}%
\end{pgfscope}%
\begin{pgfscope}%
\pgfpathrectangle{\pgfqpoint{0.407592in}{0.222904in}}{\pgfqpoint{1.411163in}{1.411163in}}%
\pgfusepath{clip}%
\pgfsetbuttcap%
\pgfsetroundjoin%
\definecolor{currentfill}{rgb}{0.749020,0.772549,0.792157}%
\pgfsetfillcolor{currentfill}%
\pgfsetfillopacity{0.160000}%
\pgfsetlinewidth{0.000000pt}%
\definecolor{currentstroke}{rgb}{0.000000,0.000000,0.000000}%
\pgfsetstrokecolor{currentstroke}%
\pgfsetdash{}{0pt}%
\pgfpathmoveto{\pgfqpoint{1.360279in}{0.874689in}}%
\pgfpathlineto{\pgfqpoint{1.367989in}{0.882230in}}%
\pgfpathlineto{\pgfqpoint{1.375732in}{0.889804in}}%
\pgfpathlineto{\pgfqpoint{1.407034in}{0.898472in}}%
\pgfpathlineto{\pgfqpoint{1.436700in}{0.907902in}}%
\pgfpathlineto{\pgfqpoint{1.427353in}{0.899881in}}%
\pgfpathlineto{\pgfqpoint{1.418045in}{0.891891in}}%
\pgfpathlineto{\pgfqpoint{1.389930in}{0.882926in}}%
\pgfpathlineto{\pgfqpoint{1.360279in}{0.874689in}}%
\pgfpathclose%
\pgfusepath{fill}%
\end{pgfscope}%
\begin{pgfscope}%
\pgfpathrectangle{\pgfqpoint{0.407592in}{0.222904in}}{\pgfqpoint{1.411163in}{1.411163in}}%
\pgfusepath{clip}%
\pgfsetbuttcap%
\pgfsetroundjoin%
\definecolor{currentfill}{rgb}{0.749020,0.772549,0.792157}%
\pgfsetfillcolor{currentfill}%
\pgfsetfillopacity{0.160000}%
\pgfsetlinewidth{0.000000pt}%
\definecolor{currentstroke}{rgb}{0.000000,0.000000,0.000000}%
\pgfsetstrokecolor{currentstroke}%
\pgfsetdash{}{0pt}%
\pgfpathmoveto{\pgfqpoint{1.733311in}{1.099082in}}%
\pgfpathlineto{\pgfqpoint{1.746270in}{1.108926in}}%
\pgfpathlineto{\pgfqpoint{1.759269in}{1.118802in}}%
\pgfpathlineto{\pgfqpoint{1.783107in}{1.136235in}}%
\pgfpathlineto{\pgfqpoint{1.803517in}{1.154061in}}%
\pgfpathlineto{\pgfqpoint{1.789814in}{1.143619in}}%
\pgfpathlineto{\pgfqpoint{1.776149in}{1.133206in}}%
\pgfpathlineto{\pgfqpoint{1.756371in}{1.115950in}}%
\pgfpathlineto{\pgfqpoint{1.733311in}{1.099082in}}%
\pgfpathclose%
\pgfusepath{fill}%
\end{pgfscope}%
\begin{pgfscope}%
\pgfpathrectangle{\pgfqpoint{0.407592in}{0.222904in}}{\pgfqpoint{1.411163in}{1.411163in}}%
\pgfusepath{clip}%
\pgfsetbuttcap%
\pgfsetroundjoin%
\definecolor{currentfill}{rgb}{0.749020,0.772549,0.792157}%
\pgfsetfillcolor{currentfill}%
\pgfsetfillopacity{0.160000}%
\pgfsetlinewidth{0.000000pt}%
\definecolor{currentstroke}{rgb}{0.000000,0.000000,0.000000}%
\pgfsetstrokecolor{currentstroke}%
\pgfsetdash{}{0pt}%
\pgfpathmoveto{\pgfqpoint{1.285168in}{0.846342in}}%
\pgfpathlineto{\pgfqpoint{1.291041in}{0.853447in}}%
\pgfpathlineto{\pgfqpoint{1.296941in}{0.860584in}}%
\pgfpathlineto{\pgfqpoint{1.329233in}{0.867227in}}%
\pgfpathlineto{\pgfqpoint{1.360279in}{0.874689in}}%
\pgfpathlineto{\pgfqpoint{1.352603in}{0.867180in}}%
\pgfpathlineto{\pgfqpoint{1.344960in}{0.859704in}}%
\pgfpathlineto{\pgfqpoint{1.315646in}{0.852634in}}%
\pgfpathlineto{\pgfqpoint{1.285168in}{0.846342in}}%
\pgfpathclose%
\pgfusepath{fill}%
\end{pgfscope}%
\begin{pgfscope}%
\pgfpathrectangle{\pgfqpoint{0.407592in}{0.222904in}}{\pgfqpoint{1.411163in}{1.411163in}}%
\pgfusepath{clip}%
\pgfsetbuttcap%
\pgfsetroundjoin%
\definecolor{currentfill}{rgb}{0.749020,0.772549,0.792157}%
\pgfsetfillcolor{currentfill}%
\pgfsetfillopacity{0.160000}%
\pgfsetlinewidth{0.000000pt}%
\definecolor{currentstroke}{rgb}{0.000000,0.000000,0.000000}%
\pgfsetstrokecolor{currentstroke}%
\pgfsetdash{}{0pt}%
\pgfpathmoveto{\pgfqpoint{0.900346in}{0.826116in}}%
\pgfpathlineto{\pgfqpoint{0.894305in}{0.832975in}}%
\pgfpathlineto{\pgfqpoint{0.888235in}{0.839866in}}%
\pgfpathlineto{\pgfqpoint{0.919714in}{0.835406in}}%
\pgfpathlineto{\pgfqpoint{0.952056in}{0.831826in}}%
\pgfpathlineto{\pgfqpoint{0.956217in}{0.825161in}}%
\pgfpathlineto{\pgfqpoint{0.960359in}{0.818528in}}%
\pgfpathlineto{\pgfqpoint{0.929950in}{0.821907in}}%
\pgfpathlineto{\pgfqpoint{0.900346in}{0.826116in}}%
\pgfpathclose%
\pgfusepath{fill}%
\end{pgfscope}%
\begin{pgfscope}%
\pgfpathrectangle{\pgfqpoint{0.407592in}{0.222904in}}{\pgfqpoint{1.411163in}{1.411163in}}%
\pgfusepath{clip}%
\pgfsetbuttcap%
\pgfsetroundjoin%
\definecolor{currentfill}{rgb}{0.749020,0.772549,0.792157}%
\pgfsetfillcolor{currentfill}%
\pgfsetfillopacity{0.160000}%
\pgfsetlinewidth{0.000000pt}%
\definecolor{currentstroke}{rgb}{0.000000,0.000000,0.000000}%
\pgfsetstrokecolor{currentstroke}%
\pgfsetdash{}{0pt}%
\pgfpathmoveto{\pgfqpoint{0.377119in}{1.129718in}}%
\pgfpathlineto{\pgfqpoint{0.362895in}{1.140073in}}%
\pgfpathlineto{\pgfqpoint{0.348632in}{1.150458in}}%
\pgfpathlineto{\pgfqpoint{0.365311in}{1.132703in}}%
\pgfpathlineto{\pgfqpoint{0.385509in}{1.115361in}}%
\pgfpathlineto{\pgfqpoint{0.399217in}{1.105541in}}%
\pgfpathlineto{\pgfqpoint{0.412881in}{1.095753in}}%
\pgfpathlineto{\pgfqpoint{0.393314in}{1.112532in}}%
\pgfpathlineto{\pgfqpoint{0.377119in}{1.129718in}}%
\pgfpathclose%
\pgfusepath{fill}%
\end{pgfscope}%
\begin{pgfscope}%
\pgfpathrectangle{\pgfqpoint{0.407592in}{0.222904in}}{\pgfqpoint{1.411163in}{1.411163in}}%
\pgfusepath{clip}%
\pgfsetbuttcap%
\pgfsetroundjoin%
\definecolor{currentfill}{rgb}{0.749020,0.772549,0.792157}%
\pgfsetfillcolor{currentfill}%
\pgfsetfillopacity{0.160000}%
\pgfsetlinewidth{0.000000pt}%
\definecolor{currentstroke}{rgb}{0.000000,0.000000,0.000000}%
\pgfsetstrokecolor{currentstroke}%
\pgfsetdash{}{0pt}%
\pgfpathmoveto{\pgfqpoint{0.493941in}{1.037685in}}%
\pgfpathlineto{\pgfqpoint{0.480540in}{1.047285in}}%
\pgfpathlineto{\pgfqpoint{0.467095in}{1.056916in}}%
\pgfpathlineto{\pgfqpoint{0.488413in}{1.041755in}}%
\pgfpathlineto{\pgfqpoint{0.512706in}{1.027154in}}%
\pgfpathlineto{\pgfqpoint{0.525283in}{1.018082in}}%
\pgfpathlineto{\pgfqpoint{0.537814in}{1.009042in}}%
\pgfpathlineto{\pgfqpoint{0.514461in}{1.023091in}}%
\pgfpathlineto{\pgfqpoint{0.493941in}{1.037685in}}%
\pgfpathclose%
\pgfusepath{fill}%
\end{pgfscope}%
\begin{pgfscope}%
\pgfpathrectangle{\pgfqpoint{0.407592in}{0.222904in}}{\pgfqpoint{1.411163in}{1.411163in}}%
\pgfusepath{clip}%
\pgfsetbuttcap%
\pgfsetroundjoin%
\definecolor{currentfill}{rgb}{0.749020,0.772549,0.792157}%
\pgfsetfillcolor{currentfill}%
\pgfsetfillopacity{0.160000}%
\pgfsetlinewidth{0.000000pt}%
\definecolor{currentstroke}{rgb}{0.000000,0.000000,0.000000}%
\pgfsetstrokecolor{currentstroke}%
\pgfsetdash{}{0pt}%
\pgfpathmoveto{\pgfqpoint{1.657470in}{1.044755in}}%
\pgfpathlineto{\pgfqpoint{1.669797in}{1.054191in}}%
\pgfpathlineto{\pgfqpoint{1.682166in}{1.063659in}}%
\pgfpathlineto{\pgfqpoint{1.707519in}{1.079488in}}%
\pgfpathlineto{\pgfqpoint{1.694685in}{1.069737in}}%
\pgfpathlineto{\pgfqpoint{1.681891in}{1.060018in}}%
\pgfusepath{fill}%
\end{pgfscope}%
\begin{pgfscope}%
\pgfpathrectangle{\pgfqpoint{0.407592in}{0.222904in}}{\pgfqpoint{1.411163in}{1.411163in}}%
\pgfusepath{clip}%
\pgfsetbuttcap%
\pgfsetroundjoin%
\definecolor{currentfill}{rgb}{0.749020,0.772549,0.792157}%
\pgfsetfillcolor{currentfill}%
\pgfsetfillopacity{0.160000}%
\pgfsetlinewidth{0.000000pt}%
\definecolor{currentstroke}{rgb}{0.000000,0.000000,0.000000}%
\pgfsetstrokecolor{currentstroke}%
\pgfsetdash{}{0pt}%
\pgfpathmoveto{\pgfqpoint{1.555834in}{0.980704in}}%
\pgfpathlineto{\pgfqpoint{1.566852in}{0.989463in}}%
\pgfpathlineto{\pgfqpoint{1.577911in}{0.998256in}}%
\pgfpathlineto{\pgfqpoint{1.606738in}{1.011813in}}%
\pgfpathlineto{\pgfqpoint{1.632941in}{1.025978in}}%
\pgfpathlineto{\pgfqpoint{1.620739in}{1.016637in}}%
\pgfpathlineto{\pgfqpoint{1.608577in}{1.007328in}}%
\pgfpathlineto{\pgfqpoint{1.583451in}{0.993721in}}%
\pgfpathlineto{\pgfqpoint{1.555834in}{0.980704in}}%
\pgfpathclose%
\pgfusepath{fill}%
\end{pgfscope}%
\begin{pgfscope}%
\pgfpathrectangle{\pgfqpoint{0.407592in}{0.222904in}}{\pgfqpoint{1.411163in}{1.411163in}}%
\pgfusepath{clip}%
\pgfsetbuttcap%
\pgfsetroundjoin%
\definecolor{currentfill}{rgb}{0.749020,0.772549,0.792157}%
\pgfsetfillcolor{currentfill}%
\pgfsetfillopacity{0.160000}%
\pgfsetlinewidth{0.000000pt}%
\definecolor{currentstroke}{rgb}{0.000000,0.000000,0.000000}%
\pgfsetstrokecolor{currentstroke}%
\pgfsetdash{}{0pt}%
\pgfpathmoveto{\pgfqpoint{1.150830in}{0.816438in}}%
\pgfpathlineto{\pgfqpoint{1.152764in}{0.823008in}}%
\pgfpathlineto{\pgfqpoint{1.154707in}{0.829610in}}%
\pgfpathlineto{\pgfqpoint{1.188278in}{0.832473in}}%
\pgfpathlineto{\pgfqpoint{1.221324in}{0.836231in}}%
\pgfpathlineto{\pgfqpoint{1.217380in}{0.829443in}}%
\pgfpathlineto{\pgfqpoint{1.213456in}{0.822686in}}%
\pgfpathlineto{\pgfqpoint{1.182387in}{0.819139in}}%
\pgfpathlineto{\pgfqpoint{1.150830in}{0.816438in}}%
\pgfpathclose%
\pgfusepath{fill}%
\end{pgfscope}%
\begin{pgfscope}%
\pgfpathrectangle{\pgfqpoint{0.407592in}{0.222904in}}{\pgfqpoint{1.411163in}{1.411163in}}%
\pgfusepath{clip}%
\pgfsetbuttcap%
\pgfsetroundjoin%
\definecolor{currentfill}{rgb}{0.749020,0.772549,0.792157}%
\pgfsetfillcolor{currentfill}%
\pgfsetfillopacity{0.160000}%
\pgfsetlinewidth{0.000000pt}%
\definecolor{currentstroke}{rgb}{0.000000,0.000000,0.000000}%
\pgfsetstrokecolor{currentstroke}%
\pgfsetdash{}{0pt}%
\pgfpathmoveto{\pgfqpoint{0.682118in}{0.915953in}}%
\pgfpathlineto{\pgfqpoint{0.671092in}{0.924205in}}%
\pgfpathlineto{\pgfqpoint{0.660022in}{0.932489in}}%
\pgfpathlineto{\pgfqpoint{0.687284in}{0.921998in}}%
\pgfpathlineto{\pgfqpoint{0.716577in}{0.912258in}}%
\pgfpathlineto{\pgfqpoint{0.726250in}{0.904448in}}%
\pgfpathlineto{\pgfqpoint{0.735882in}{0.896670in}}%
\pgfpathlineto{\pgfqpoint{0.708043in}{0.905952in}}%
\pgfpathlineto{\pgfqpoint{0.682118in}{0.915953in}}%
\pgfpathclose%
\pgfusepath{fill}%
\end{pgfscope}%
\begin{pgfscope}%
\pgfpathrectangle{\pgfqpoint{0.407592in}{0.222904in}}{\pgfqpoint{1.411163in}{1.411163in}}%
\pgfusepath{clip}%
\pgfsetbuttcap%
\pgfsetroundjoin%
\definecolor{currentfill}{rgb}{0.749020,0.772549,0.792157}%
\pgfsetfillcolor{currentfill}%
\pgfsetfillopacity{0.160000}%
\pgfsetlinewidth{0.000000pt}%
\definecolor{currentstroke}{rgb}{0.000000,0.000000,0.000000}%
\pgfsetstrokecolor{currentstroke}%
\pgfsetdash{}{0pt}%
\pgfpathmoveto{\pgfqpoint{1.455507in}{0.924044in}}%
\pgfpathlineto{\pgfqpoint{1.464968in}{0.932165in}}%
\pgfpathlineto{\pgfqpoint{1.474468in}{0.940318in}}%
\pgfpathlineto{\pgfqpoint{1.505262in}{0.951448in}}%
\pgfpathlineto{\pgfqpoint{1.533923in}{0.963284in}}%
\pgfpathlineto{\pgfqpoint{1.523029in}{0.954622in}}%
\pgfpathlineto{\pgfqpoint{1.512175in}{0.945993in}}%
\pgfpathlineto{\pgfqpoint{1.484848in}{0.934679in}}%
\pgfpathlineto{\pgfqpoint{1.455507in}{0.924044in}}%
\pgfpathclose%
\pgfusepath{fill}%
\end{pgfscope}%
\begin{pgfscope}%
\pgfpathrectangle{\pgfqpoint{0.407592in}{0.222904in}}{\pgfqpoint{1.411163in}{1.411163in}}%
\pgfusepath{clip}%
\pgfsetbuttcap%
\pgfsetroundjoin%
\definecolor{currentfill}{rgb}{0.749020,0.772549,0.792157}%
\pgfsetfillcolor{currentfill}%
\pgfsetfillopacity{0.160000}%
\pgfsetlinewidth{0.000000pt}%
\definecolor{currentstroke}{rgb}{0.000000,0.000000,0.000000}%
\pgfsetstrokecolor{currentstroke}%
\pgfsetdash{}{0pt}%
\pgfpathmoveto{\pgfqpoint{1.759269in}{1.118802in}}%
\pgfpathlineto{\pgfqpoint{1.772311in}{1.128710in}}%
\pgfpathlineto{\pgfqpoint{1.796535in}{1.146423in}}%
\pgfpathlineto{\pgfqpoint{1.817258in}{1.164532in}}%
\pgfpathlineto{\pgfqpoint{1.803517in}{1.154061in}}%
\pgfpathlineto{\pgfqpoint{1.783107in}{1.136235in}}%
\pgfusepath{fill}%
\end{pgfscope}%
\begin{pgfscope}%
\pgfpathrectangle{\pgfqpoint{0.407592in}{0.222904in}}{\pgfqpoint{1.411163in}{1.411163in}}%
\pgfusepath{clip}%
\pgfsetbuttcap%
\pgfsetroundjoin%
\definecolor{currentfill}{rgb}{0.749020,0.772549,0.792157}%
\pgfsetfillcolor{currentfill}%
\pgfsetfillopacity{0.160000}%
\pgfsetlinewidth{0.000000pt}%
\definecolor{currentstroke}{rgb}{0.000000,0.000000,0.000000}%
\pgfsetstrokecolor{currentstroke}%
\pgfsetdash{}{0pt}%
\pgfpathmoveto{\pgfqpoint{0.828493in}{0.851308in}}%
\pgfpathlineto{\pgfqpoint{0.820619in}{0.858551in}}%
\pgfpathlineto{\pgfqpoint{0.812710in}{0.865827in}}%
\pgfpathlineto{\pgfqpoint{0.843734in}{0.859353in}}%
\pgfpathlineto{\pgfqpoint{0.876011in}{0.853744in}}%
\pgfpathlineto{\pgfqpoint{0.882137in}{0.846789in}}%
\pgfpathlineto{\pgfqpoint{0.888235in}{0.839866in}}%
\pgfpathlineto{\pgfqpoint{0.857778in}{0.845177in}}%
\pgfpathlineto{\pgfqpoint{0.828493in}{0.851308in}}%
\pgfpathclose%
\pgfusepath{fill}%
\end{pgfscope}%
\begin{pgfscope}%
\pgfpathrectangle{\pgfqpoint{0.407592in}{0.222904in}}{\pgfqpoint{1.411163in}{1.411163in}}%
\pgfusepath{clip}%
\pgfsetbuttcap%
\pgfsetroundjoin%
\definecolor{currentfill}{rgb}{0.749020,0.772549,0.792157}%
\pgfsetfillcolor{currentfill}%
\pgfsetfillopacity{0.160000}%
\pgfsetlinewidth{0.000000pt}%
\definecolor{currentstroke}{rgb}{0.000000,0.000000,0.000000}%
\pgfsetstrokecolor{currentstroke}%
\pgfsetdash{}{0pt}%
\pgfpathmoveto{\pgfqpoint{0.960359in}{0.818528in}}%
\pgfpathlineto{\pgfqpoint{0.956217in}{0.825161in}}%
\pgfpathlineto{\pgfqpoint{0.952056in}{0.831826in}}%
\pgfpathlineto{\pgfqpoint{0.985095in}{0.829145in}}%
\pgfpathlineto{\pgfqpoint{1.018663in}{0.827381in}}%
\pgfpathlineto{\pgfqpoint{1.020821in}{0.820841in}}%
\pgfpathlineto{\pgfqpoint{1.022968in}{0.814334in}}%
\pgfpathlineto{\pgfqpoint{0.991417in}{0.815999in}}%
\pgfpathlineto{\pgfqpoint{0.960359in}{0.818528in}}%
\pgfpathclose%
\pgfusepath{fill}%
\end{pgfscope}%
\begin{pgfscope}%
\pgfpathrectangle{\pgfqpoint{0.407592in}{0.222904in}}{\pgfqpoint{1.411163in}{1.411163in}}%
\pgfusepath{clip}%
\pgfsetbuttcap%
\pgfsetroundjoin%
\definecolor{currentfill}{rgb}{0.749020,0.772549,0.792157}%
\pgfsetfillcolor{currentfill}%
\pgfsetfillopacity{0.160000}%
\pgfsetlinewidth{0.000000pt}%
\definecolor{currentstroke}{rgb}{0.000000,0.000000,0.000000}%
\pgfsetstrokecolor{currentstroke}%
\pgfsetdash{}{0pt}%
\pgfpathmoveto{\pgfqpoint{0.587487in}{0.973208in}}%
\pgfpathlineto{\pgfqpoint{0.575136in}{0.982118in}}%
\pgfpathlineto{\pgfqpoint{0.562740in}{0.991060in}}%
\pgfpathlineto{\pgfqpoint{0.587763in}{0.978170in}}%
\pgfpathlineto{\pgfqpoint{0.615303in}{0.965956in}}%
\pgfpathlineto{\pgfqpoint{0.626550in}{0.957539in}}%
\pgfpathlineto{\pgfqpoint{0.637751in}{0.949156in}}%
\pgfpathlineto{\pgfqpoint{0.611427in}{0.960856in}}%
\pgfpathlineto{\pgfqpoint{0.587487in}{0.973208in}}%
\pgfpathclose%
\pgfusepath{fill}%
\end{pgfscope}%
\begin{pgfscope}%
\pgfpathrectangle{\pgfqpoint{0.407592in}{0.222904in}}{\pgfqpoint{1.411163in}{1.411163in}}%
\pgfusepath{clip}%
\pgfsetbuttcap%
\pgfsetroundjoin%
\definecolor{currentfill}{rgb}{0.749020,0.772549,0.792157}%
\pgfsetfillcolor{currentfill}%
\pgfsetfillopacity{0.160000}%
\pgfsetlinewidth{0.000000pt}%
\definecolor{currentstroke}{rgb}{0.000000,0.000000,0.000000}%
\pgfsetstrokecolor{currentstroke}%
\pgfsetdash{}{0pt}%
\pgfpathmoveto{\pgfqpoint{0.755023in}{0.881214in}}%
\pgfpathlineto{\pgfqpoint{0.745472in}{0.888925in}}%
\pgfpathlineto{\pgfqpoint{0.735882in}{0.896670in}}%
\pgfpathlineto{\pgfqpoint{0.765509in}{0.888162in}}%
\pgfpathlineto{\pgfqpoint{0.796785in}{0.880477in}}%
\pgfpathlineto{\pgfqpoint{0.804765in}{0.873135in}}%
\pgfpathlineto{\pgfqpoint{0.812710in}{0.865827in}}%
\pgfpathlineto{\pgfqpoint{0.783092in}{0.873128in}}%
\pgfpathlineto{\pgfqpoint{0.755023in}{0.881214in}}%
\pgfpathclose%
\pgfusepath{fill}%
\end{pgfscope}%
\begin{pgfscope}%
\pgfpathrectangle{\pgfqpoint{0.407592in}{0.222904in}}{\pgfqpoint{1.411163in}{1.411163in}}%
\pgfusepath{clip}%
\pgfsetbuttcap%
\pgfsetroundjoin%
\definecolor{currentfill}{rgb}{0.749020,0.772549,0.792157}%
\pgfsetfillcolor{currentfill}%
\pgfsetfillopacity{0.160000}%
\pgfsetlinewidth{0.000000pt}%
\definecolor{currentstroke}{rgb}{0.000000,0.000000,0.000000}%
\pgfsetstrokecolor{currentstroke}%
\pgfsetdash{}{0pt}%
\pgfpathmoveto{\pgfqpoint{1.086897in}{0.813631in}}%
\pgfpathlineto{\pgfqpoint{1.086791in}{0.820118in}}%
\pgfpathlineto{\pgfqpoint{1.086684in}{0.826636in}}%
\pgfpathlineto{\pgfqpoint{1.120784in}{0.827661in}}%
\pgfpathlineto{\pgfqpoint{1.154707in}{0.829610in}}%
\pgfpathlineto{\pgfqpoint{1.152764in}{0.823008in}}%
\pgfpathlineto{\pgfqpoint{1.150830in}{0.816438in}}%
\pgfpathlineto{\pgfqpoint{1.118946in}{0.814598in}}%
\pgfpathlineto{\pgfqpoint{1.086897in}{0.813631in}}%
\pgfpathclose%
\pgfusepath{fill}%
\end{pgfscope}%
\begin{pgfscope}%
\pgfpathrectangle{\pgfqpoint{0.407592in}{0.222904in}}{\pgfqpoint{1.411163in}{1.411163in}}%
\pgfusepath{clip}%
\pgfsetbuttcap%
\pgfsetroundjoin%
\definecolor{currentfill}{rgb}{0.749020,0.772549,0.792157}%
\pgfsetfillcolor{currentfill}%
\pgfsetfillopacity{0.160000}%
\pgfsetlinewidth{0.000000pt}%
\definecolor{currentstroke}{rgb}{0.000000,0.000000,0.000000}%
\pgfsetstrokecolor{currentstroke}%
\pgfsetdash{}{0pt}%
\pgfpathmoveto{\pgfqpoint{0.348632in}{1.150458in}}%
\pgfpathlineto{\pgfqpoint{0.334328in}{1.160872in}}%
\pgfpathlineto{\pgfqpoint{0.351245in}{1.142835in}}%
\pgfpathlineto{\pgfqpoint{0.371757in}{1.125212in}}%
\pgfpathlineto{\pgfqpoint{0.385509in}{1.115361in}}%
\pgfpathlineto{\pgfqpoint{0.365311in}{1.132703in}}%
\pgfusepath{fill}%
\end{pgfscope}%
\begin{pgfscope}%
\pgfpathrectangle{\pgfqpoint{0.407592in}{0.222904in}}{\pgfqpoint{1.411163in}{1.411163in}}%
\pgfusepath{clip}%
\pgfsetbuttcap%
\pgfsetroundjoin%
\definecolor{currentfill}{rgb}{0.749020,0.772549,0.792157}%
\pgfsetfillcolor{currentfill}%
\pgfsetfillopacity{0.160000}%
\pgfsetlinewidth{0.000000pt}%
\definecolor{currentstroke}{rgb}{0.000000,0.000000,0.000000}%
\pgfsetstrokecolor{currentstroke}%
\pgfsetdash{}{0pt}%
\pgfpathmoveto{\pgfqpoint{1.022968in}{0.814334in}}%
\pgfpathlineto{\pgfqpoint{1.020821in}{0.820841in}}%
\pgfpathlineto{\pgfqpoint{1.018663in}{0.827381in}}%
\pgfpathlineto{\pgfqpoint{1.052584in}{0.826543in}}%
\pgfpathlineto{\pgfqpoint{1.086684in}{0.826636in}}%
\pgfpathlineto{\pgfqpoint{1.086791in}{0.820118in}}%
\pgfpathlineto{\pgfqpoint{1.086897in}{0.813631in}}%
\pgfpathlineto{\pgfqpoint{1.054850in}{0.813543in}}%
\pgfpathlineto{\pgfqpoint{1.022968in}{0.814334in}}%
\pgfpathclose%
\pgfusepath{fill}%
\end{pgfscope}%
\begin{pgfscope}%
\pgfpathrectangle{\pgfqpoint{0.407592in}{0.222904in}}{\pgfqpoint{1.411163in}{1.411163in}}%
\pgfusepath{clip}%
\pgfsetbuttcap%
\pgfsetroundjoin%
\definecolor{currentfill}{rgb}{0.749020,0.772549,0.792157}%
\pgfsetfillcolor{currentfill}%
\pgfsetfillopacity{0.160000}%
\pgfsetlinewidth{0.000000pt}%
\definecolor{currentstroke}{rgb}{0.000000,0.000000,0.000000}%
\pgfsetstrokecolor{currentstroke}%
\pgfsetdash{}{0pt}%
\pgfpathmoveto{\pgfqpoint{1.221324in}{0.836231in}}%
\pgfpathlineto{\pgfqpoint{1.225286in}{0.843052in}}%
\pgfpathlineto{\pgfqpoint{1.229266in}{0.849906in}}%
\pgfpathlineto{\pgfqpoint{1.263563in}{0.854799in}}%
\pgfpathlineto{\pgfqpoint{1.296941in}{0.860584in}}%
\pgfpathlineto{\pgfqpoint{1.291041in}{0.853447in}}%
\pgfpathlineto{\pgfqpoint{1.285168in}{0.846342in}}%
\pgfpathlineto{\pgfqpoint{1.253675in}{0.840864in}}%
\pgfpathlineto{\pgfqpoint{1.221324in}{0.836231in}}%
\pgfpathclose%
\pgfusepath{fill}%
\end{pgfscope}%
\begin{pgfscope}%
\pgfpathrectangle{\pgfqpoint{0.407592in}{0.222904in}}{\pgfqpoint{1.411163in}{1.411163in}}%
\pgfusepath{clip}%
\pgfsetbuttcap%
\pgfsetroundjoin%
\definecolor{currentfill}{rgb}{0.749020,0.772549,0.792157}%
\pgfsetfillcolor{currentfill}%
\pgfsetfillopacity{0.160000}%
\pgfsetlinewidth{0.000000pt}%
\definecolor{currentstroke}{rgb}{0.000000,0.000000,0.000000}%
\pgfsetstrokecolor{currentstroke}%
\pgfsetdash{}{0pt}%
\pgfpathmoveto{\pgfqpoint{1.375732in}{0.889804in}}%
\pgfpathlineto{\pgfqpoint{1.383509in}{0.897411in}}%
\pgfpathlineto{\pgfqpoint{1.391320in}{0.905052in}}%
\pgfpathlineto{\pgfqpoint{1.424283in}{0.914149in}}%
\pgfpathlineto{\pgfqpoint{1.455507in}{0.924044in}}%
\pgfpathlineto{\pgfqpoint{1.446084in}{0.915957in}}%
\pgfpathlineto{\pgfqpoint{1.436700in}{0.907902in}}%
\pgfpathlineto{\pgfqpoint{1.407034in}{0.898472in}}%
\pgfpathlineto{\pgfqpoint{1.375732in}{0.889804in}}%
\pgfpathclose%
\pgfusepath{fill}%
\end{pgfscope}%
\begin{pgfscope}%
\pgfpathrectangle{\pgfqpoint{0.407592in}{0.222904in}}{\pgfqpoint{1.411163in}{1.411163in}}%
\pgfusepath{clip}%
\pgfsetbuttcap%
\pgfsetroundjoin%
\definecolor{currentfill}{rgb}{0.749020,0.772549,0.792157}%
\pgfsetfillcolor{currentfill}%
\pgfsetfillopacity{0.160000}%
\pgfsetlinewidth{0.000000pt}%
\definecolor{currentstroke}{rgb}{0.000000,0.000000,0.000000}%
\pgfsetstrokecolor{currentstroke}%
\pgfsetdash{}{0pt}%
\pgfpathmoveto{\pgfqpoint{1.296941in}{0.860584in}}%
\pgfpathlineto{\pgfqpoint{1.302868in}{0.867753in}}%
\pgfpathlineto{\pgfqpoint{1.308823in}{0.874956in}}%
\pgfpathlineto{\pgfqpoint{1.342942in}{0.881950in}}%
\pgfpathlineto{\pgfqpoint{1.375732in}{0.889804in}}%
\pgfpathlineto{\pgfqpoint{1.367989in}{0.882230in}}%
\pgfpathlineto{\pgfqpoint{1.360279in}{0.874689in}}%
\pgfpathlineto{\pgfqpoint{1.329233in}{0.867227in}}%
\pgfpathlineto{\pgfqpoint{1.296941in}{0.860584in}}%
\pgfpathclose%
\pgfusepath{fill}%
\end{pgfscope}%
\begin{pgfscope}%
\pgfpathrectangle{\pgfqpoint{0.407592in}{0.222904in}}{\pgfqpoint{1.411163in}{1.411163in}}%
\pgfusepath{clip}%
\pgfsetbuttcap%
\pgfsetroundjoin%
\definecolor{currentfill}{rgb}{0.749020,0.772549,0.792157}%
\pgfsetfillcolor{currentfill}%
\pgfsetfillopacity{0.160000}%
\pgfsetlinewidth{0.000000pt}%
\definecolor{currentstroke}{rgb}{0.000000,0.000000,0.000000}%
\pgfsetstrokecolor{currentstroke}%
\pgfsetdash{}{0pt}%
\pgfpathmoveto{\pgfqpoint{1.682166in}{1.063659in}}%
\pgfpathlineto{\pgfqpoint{1.694577in}{1.073160in}}%
\pgfpathlineto{\pgfqpoint{1.707030in}{1.082693in}}%
\pgfpathlineto{\pgfqpoint{1.733311in}{1.099082in}}%
\pgfpathlineto{\pgfqpoint{1.720395in}{1.089269in}}%
\pgfpathlineto{\pgfqpoint{1.707519in}{1.079488in}}%
\pgfusepath{fill}%
\end{pgfscope}%
\begin{pgfscope}%
\pgfpathrectangle{\pgfqpoint{0.407592in}{0.222904in}}{\pgfqpoint{1.411163in}{1.411163in}}%
\pgfusepath{clip}%
\pgfsetbuttcap%
\pgfsetroundjoin%
\definecolor{currentfill}{rgb}{0.749020,0.772549,0.792157}%
\pgfsetfillcolor{currentfill}%
\pgfsetfillopacity{0.160000}%
\pgfsetlinewidth{0.000000pt}%
\definecolor{currentstroke}{rgb}{0.000000,0.000000,0.000000}%
\pgfsetstrokecolor{currentstroke}%
\pgfsetdash{}{0pt}%
\pgfpathmoveto{\pgfqpoint{0.467095in}{1.056916in}}%
\pgfpathlineto{\pgfqpoint{0.453607in}{1.066578in}}%
\pgfpathlineto{\pgfqpoint{0.440076in}{1.076272in}}%
\pgfpathlineto{\pgfqpoint{0.462186in}{1.060548in}}%
\pgfpathlineto{\pgfqpoint{0.487416in}{1.045399in}}%
\pgfpathlineto{\pgfqpoint{0.500084in}{1.036260in}}%
\pgfpathlineto{\pgfqpoint{0.512706in}{1.027154in}}%
\pgfpathlineto{\pgfqpoint{0.488413in}{1.041755in}}%
\pgfpathlineto{\pgfqpoint{0.467095in}{1.056916in}}%
\pgfpathclose%
\pgfusepath{fill}%
\end{pgfscope}%
\begin{pgfscope}%
\pgfpathrectangle{\pgfqpoint{0.407592in}{0.222904in}}{\pgfqpoint{1.411163in}{1.411163in}}%
\pgfusepath{clip}%
\pgfsetbuttcap%
\pgfsetroundjoin%
\definecolor{currentfill}{rgb}{0.749020,0.772549,0.792157}%
\pgfsetfillcolor{currentfill}%
\pgfsetfillopacity{0.160000}%
\pgfsetlinewidth{0.000000pt}%
\definecolor{currentstroke}{rgb}{0.000000,0.000000,0.000000}%
\pgfsetstrokecolor{currentstroke}%
\pgfsetdash{}{0pt}%
\pgfpathmoveto{\pgfqpoint{1.577911in}{0.998256in}}%
\pgfpathlineto{\pgfqpoint{1.589012in}{1.007081in}}%
\pgfpathlineto{\pgfqpoint{1.600155in}{1.015940in}}%
\pgfpathlineto{\pgfqpoint{1.630192in}{1.030035in}}%
\pgfpathlineto{\pgfqpoint{1.657470in}{1.044755in}}%
\pgfpathlineto{\pgfqpoint{1.645185in}{1.035351in}}%
\pgfpathlineto{\pgfqpoint{1.632941in}{1.025978in}}%
\pgfpathlineto{\pgfqpoint{1.606738in}{1.011813in}}%
\pgfpathlineto{\pgfqpoint{1.577911in}{0.998256in}}%
\pgfpathclose%
\pgfusepath{fill}%
\end{pgfscope}%
\begin{pgfscope}%
\pgfpathrectangle{\pgfqpoint{0.407592in}{0.222904in}}{\pgfqpoint{1.411163in}{1.411163in}}%
\pgfusepath{clip}%
\pgfsetbuttcap%
\pgfsetroundjoin%
\definecolor{currentfill}{rgb}{0.749020,0.772549,0.792157}%
\pgfsetfillcolor{currentfill}%
\pgfsetfillopacity{0.160000}%
\pgfsetlinewidth{0.000000pt}%
\definecolor{currentstroke}{rgb}{0.000000,0.000000,0.000000}%
\pgfsetstrokecolor{currentstroke}%
\pgfsetdash{}{0pt}%
\pgfpathmoveto{\pgfqpoint{0.888235in}{0.839866in}}%
\pgfpathlineto{\pgfqpoint{0.882137in}{0.846789in}}%
\pgfpathlineto{\pgfqpoint{0.876011in}{0.853744in}}%
\pgfpathlineto{\pgfqpoint{0.909380in}{0.849035in}}%
\pgfpathlineto{\pgfqpoint{0.943672in}{0.845252in}}%
\pgfpathlineto{\pgfqpoint{0.947874in}{0.838523in}}%
\pgfpathlineto{\pgfqpoint{0.952056in}{0.831826in}}%
\pgfpathlineto{\pgfqpoint{0.919714in}{0.835406in}}%
\pgfpathlineto{\pgfqpoint{0.888235in}{0.839866in}}%
\pgfpathclose%
\pgfusepath{fill}%
\end{pgfscope}%
\begin{pgfscope}%
\pgfpathrectangle{\pgfqpoint{0.407592in}{0.222904in}}{\pgfqpoint{1.411163in}{1.411163in}}%
\pgfusepath{clip}%
\pgfsetbuttcap%
\pgfsetroundjoin%
\definecolor{currentfill}{rgb}{0.749020,0.772549,0.792157}%
\pgfsetfillcolor{currentfill}%
\pgfsetfillopacity{0.160000}%
\pgfsetlinewidth{0.000000pt}%
\definecolor{currentstroke}{rgb}{0.000000,0.000000,0.000000}%
\pgfsetstrokecolor{currentstroke}%
\pgfsetdash{}{0pt}%
\pgfpathmoveto{\pgfqpoint{1.474468in}{0.940318in}}%
\pgfpathlineto{\pgfqpoint{1.484007in}{0.948506in}}%
\pgfpathlineto{\pgfqpoint{1.493586in}{0.956726in}}%
\pgfpathlineto{\pgfqpoint{1.525837in}{0.968349in}}%
\pgfpathlineto{\pgfqpoint{1.555834in}{0.980704in}}%
\pgfpathlineto{\pgfqpoint{1.544858in}{0.971978in}}%
\pgfpathlineto{\pgfqpoint{1.533923in}{0.963284in}}%
\pgfpathlineto{\pgfqpoint{1.505262in}{0.951448in}}%
\pgfpathlineto{\pgfqpoint{1.474468in}{0.940318in}}%
\pgfpathclose%
\pgfusepath{fill}%
\end{pgfscope}%
\begin{pgfscope}%
\pgfpathrectangle{\pgfqpoint{0.407592in}{0.222904in}}{\pgfqpoint{1.411163in}{1.411163in}}%
\pgfusepath{clip}%
\pgfsetbuttcap%
\pgfsetroundjoin%
\definecolor{currentfill}{rgb}{0.749020,0.772549,0.792157}%
\pgfsetfillcolor{currentfill}%
\pgfsetfillopacity{0.160000}%
\pgfsetlinewidth{0.000000pt}%
\definecolor{currentstroke}{rgb}{0.000000,0.000000,0.000000}%
\pgfsetstrokecolor{currentstroke}%
\pgfsetdash{}{0pt}%
\pgfpathmoveto{\pgfqpoint{0.660022in}{0.932489in}}%
\pgfpathlineto{\pgfqpoint{0.648909in}{0.940806in}}%
\pgfpathlineto{\pgfqpoint{0.637751in}{0.949156in}}%
\pgfpathlineto{\pgfqpoint{0.666353in}{0.938177in}}%
\pgfpathlineto{\pgfqpoint{0.697107in}{0.927980in}}%
\pgfpathlineto{\pgfqpoint{0.706863in}{0.920103in}}%
\pgfpathlineto{\pgfqpoint{0.716577in}{0.912258in}}%
\pgfpathlineto{\pgfqpoint{0.687284in}{0.921998in}}%
\pgfpathlineto{\pgfqpoint{0.660022in}{0.932489in}}%
\pgfpathclose%
\pgfusepath{fill}%
\end{pgfscope}%
\begin{pgfscope}%
\pgfpathrectangle{\pgfqpoint{0.407592in}{0.222904in}}{\pgfqpoint{1.411163in}{1.411163in}}%
\pgfusepath{clip}%
\pgfsetbuttcap%
\pgfsetroundjoin%
\definecolor{currentfill}{rgb}{0.749020,0.772549,0.792157}%
\pgfsetfillcolor{currentfill}%
\pgfsetfillopacity{0.160000}%
\pgfsetlinewidth{0.000000pt}%
\definecolor{currentstroke}{rgb}{0.000000,0.000000,0.000000}%
\pgfsetstrokecolor{currentstroke}%
\pgfsetdash{}{0pt}%
\pgfpathmoveto{\pgfqpoint{1.154707in}{0.829610in}}%
\pgfpathlineto{\pgfqpoint{1.156660in}{0.836245in}}%
\pgfpathlineto{\pgfqpoint{1.158622in}{0.842912in}}%
\pgfpathlineto{\pgfqpoint{1.194226in}{0.845936in}}%
\pgfpathlineto{\pgfqpoint{1.229266in}{0.849906in}}%
\pgfpathlineto{\pgfqpoint{1.225286in}{0.843052in}}%
\pgfpathlineto{\pgfqpoint{1.221324in}{0.836231in}}%
\pgfpathlineto{\pgfqpoint{1.188278in}{0.832473in}}%
\pgfpathlineto{\pgfqpoint{1.154707in}{0.829610in}}%
\pgfpathclose%
\pgfusepath{fill}%
\end{pgfscope}%
\begin{pgfscope}%
\pgfpathrectangle{\pgfqpoint{0.407592in}{0.222904in}}{\pgfqpoint{1.411163in}{1.411163in}}%
\pgfusepath{clip}%
\pgfsetbuttcap%
\pgfsetroundjoin%
\definecolor{currentfill}{rgb}{0.749020,0.772549,0.792157}%
\pgfsetfillcolor{currentfill}%
\pgfsetfillopacity{0.160000}%
\pgfsetlinewidth{0.000000pt}%
\definecolor{currentstroke}{rgb}{0.000000,0.000000,0.000000}%
\pgfsetstrokecolor{currentstroke}%
\pgfsetdash{}{0pt}%
\pgfpathmoveto{\pgfqpoint{0.562740in}{0.991060in}}%
\pgfpathlineto{\pgfqpoint{0.550299in}{1.000035in}}%
\pgfpathlineto{\pgfqpoint{0.537814in}{1.009042in}}%
\pgfpathlineto{\pgfqpoint{0.563918in}{0.995616in}}%
\pgfpathlineto{\pgfqpoint{0.592676in}{0.982890in}}%
\pgfpathlineto{\pgfqpoint{0.604012in}{0.974406in}}%
\pgfpathlineto{\pgfqpoint{0.615303in}{0.965956in}}%
\pgfpathlineto{\pgfqpoint{0.587763in}{0.978170in}}%
\pgfpathlineto{\pgfqpoint{0.562740in}{0.991060in}}%
\pgfpathclose%
\pgfusepath{fill}%
\end{pgfscope}%
\begin{pgfscope}%
\pgfpathrectangle{\pgfqpoint{0.407592in}{0.222904in}}{\pgfqpoint{1.411163in}{1.411163in}}%
\pgfusepath{clip}%
\pgfsetbuttcap%
\pgfsetroundjoin%
\definecolor{currentfill}{rgb}{0.749020,0.772549,0.792157}%
\pgfsetfillcolor{currentfill}%
\pgfsetfillopacity{0.160000}%
\pgfsetlinewidth{0.000000pt}%
\definecolor{currentstroke}{rgb}{0.000000,0.000000,0.000000}%
\pgfsetstrokecolor{currentstroke}%
\pgfsetdash{}{0pt}%
\pgfpathmoveto{\pgfqpoint{0.812710in}{0.865827in}}%
\pgfpathlineto{\pgfqpoint{0.804765in}{0.873135in}}%
\pgfpathlineto{\pgfqpoint{0.796785in}{0.880477in}}%
\pgfpathlineto{\pgfqpoint{0.829560in}{0.873661in}}%
\pgfpathlineto{\pgfqpoint{0.863671in}{0.867754in}}%
\pgfpathlineto{\pgfqpoint{0.869855in}{0.860733in}}%
\pgfpathlineto{\pgfqpoint{0.876011in}{0.853744in}}%
\pgfpathlineto{\pgfqpoint{0.843734in}{0.859353in}}%
\pgfpathlineto{\pgfqpoint{0.812710in}{0.865827in}}%
\pgfpathclose%
\pgfusepath{fill}%
\end{pgfscope}%
\begin{pgfscope}%
\pgfpathrectangle{\pgfqpoint{0.407592in}{0.222904in}}{\pgfqpoint{1.411163in}{1.411163in}}%
\pgfusepath{clip}%
\pgfsetbuttcap%
\pgfsetroundjoin%
\definecolor{currentfill}{rgb}{0.749020,0.772549,0.792157}%
\pgfsetfillcolor{currentfill}%
\pgfsetfillopacity{0.160000}%
\pgfsetlinewidth{0.000000pt}%
\definecolor{currentstroke}{rgb}{0.000000,0.000000,0.000000}%
\pgfsetstrokecolor{currentstroke}%
\pgfsetdash{}{0pt}%
\pgfpathmoveto{\pgfqpoint{0.735882in}{0.896670in}}%
\pgfpathlineto{\pgfqpoint{0.726250in}{0.904448in}}%
\pgfpathlineto{\pgfqpoint{0.716577in}{0.912258in}}%
\pgfpathlineto{\pgfqpoint{0.747770in}{0.903328in}}%
\pgfpathlineto{\pgfqpoint{0.780715in}{0.895259in}}%
\pgfpathlineto{\pgfqpoint{0.788768in}{0.887851in}}%
\pgfpathlineto{\pgfqpoint{0.796785in}{0.880477in}}%
\pgfpathlineto{\pgfqpoint{0.765509in}{0.888162in}}%
\pgfpathlineto{\pgfqpoint{0.735882in}{0.896670in}}%
\pgfpathclose%
\pgfusepath{fill}%
\end{pgfscope}%
\begin{pgfscope}%
\pgfpathrectangle{\pgfqpoint{0.407592in}{0.222904in}}{\pgfqpoint{1.411163in}{1.411163in}}%
\pgfusepath{clip}%
\pgfsetbuttcap%
\pgfsetroundjoin%
\definecolor{currentfill}{rgb}{0.749020,0.772549,0.792157}%
\pgfsetfillcolor{currentfill}%
\pgfsetfillopacity{0.160000}%
\pgfsetlinewidth{0.000000pt}%
\definecolor{currentstroke}{rgb}{0.000000,0.000000,0.000000}%
\pgfsetstrokecolor{currentstroke}%
\pgfsetdash{}{0pt}%
\pgfpathmoveto{\pgfqpoint{0.952056in}{0.831826in}}%
\pgfpathlineto{\pgfqpoint{0.947874in}{0.838523in}}%
\pgfpathlineto{\pgfqpoint{0.943672in}{0.845252in}}%
\pgfpathlineto{\pgfqpoint{0.978711in}{0.842420in}}%
\pgfpathlineto{\pgfqpoint{1.014314in}{0.840556in}}%
\pgfpathlineto{\pgfqpoint{1.016494in}{0.833952in}}%
\pgfpathlineto{\pgfqpoint{1.018663in}{0.827381in}}%
\pgfpathlineto{\pgfqpoint{0.985095in}{0.829145in}}%
\pgfpathlineto{\pgfqpoint{0.952056in}{0.831826in}}%
\pgfpathclose%
\pgfusepath{fill}%
\end{pgfscope}%
\begin{pgfscope}%
\pgfpathrectangle{\pgfqpoint{0.407592in}{0.222904in}}{\pgfqpoint{1.411163in}{1.411163in}}%
\pgfusepath{clip}%
\pgfsetbuttcap%
\pgfsetroundjoin%
\definecolor{currentfill}{rgb}{0.749020,0.772549,0.792157}%
\pgfsetfillcolor{currentfill}%
\pgfsetfillopacity{0.160000}%
\pgfsetlinewidth{0.000000pt}%
\definecolor{currentstroke}{rgb}{0.000000,0.000000,0.000000}%
\pgfsetstrokecolor{currentstroke}%
\pgfsetdash{}{0pt}%
\pgfpathmoveto{\pgfqpoint{1.086684in}{0.826636in}}%
\pgfpathlineto{\pgfqpoint{1.086576in}{0.833186in}}%
\pgfpathlineto{\pgfqpoint{1.086468in}{0.839769in}}%
\pgfpathlineto{\pgfqpoint{1.122640in}{0.840852in}}%
\pgfpathlineto{\pgfqpoint{1.158622in}{0.842912in}}%
\pgfpathlineto{\pgfqpoint{1.156660in}{0.836245in}}%
\pgfpathlineto{\pgfqpoint{1.154707in}{0.829610in}}%
\pgfpathlineto{\pgfqpoint{1.120784in}{0.827661in}}%
\pgfpathlineto{\pgfqpoint{1.086684in}{0.826636in}}%
\pgfpathclose%
\pgfusepath{fill}%
\end{pgfscope}%
\begin{pgfscope}%
\pgfpathrectangle{\pgfqpoint{0.407592in}{0.222904in}}{\pgfqpoint{1.411163in}{1.411163in}}%
\pgfusepath{clip}%
\pgfsetbuttcap%
\pgfsetroundjoin%
\definecolor{currentfill}{rgb}{0.749020,0.772549,0.792157}%
\pgfsetfillcolor{currentfill}%
\pgfsetfillopacity{0.160000}%
\pgfsetlinewidth{0.000000pt}%
\definecolor{currentstroke}{rgb}{0.000000,0.000000,0.000000}%
\pgfsetstrokecolor{currentstroke}%
\pgfsetdash{}{0pt}%
\pgfpathmoveto{\pgfqpoint{1.707030in}{1.082693in}}%
\pgfpathlineto{\pgfqpoint{1.719526in}{1.092258in}}%
\pgfpathlineto{\pgfqpoint{1.732064in}{1.101857in}}%
\pgfpathlineto{\pgfqpoint{1.759269in}{1.118802in}}%
\pgfpathlineto{\pgfqpoint{1.746270in}{1.108926in}}%
\pgfpathlineto{\pgfqpoint{1.733311in}{1.099082in}}%
\pgfusepath{fill}%
\end{pgfscope}%
\begin{pgfscope}%
\pgfpathrectangle{\pgfqpoint{0.407592in}{0.222904in}}{\pgfqpoint{1.411163in}{1.411163in}}%
\pgfusepath{clip}%
\pgfsetbuttcap%
\pgfsetroundjoin%
\definecolor{currentfill}{rgb}{0.749020,0.772549,0.792157}%
\pgfsetfillcolor{currentfill}%
\pgfsetfillopacity{0.160000}%
\pgfsetlinewidth{0.000000pt}%
\definecolor{currentstroke}{rgb}{0.000000,0.000000,0.000000}%
\pgfsetstrokecolor{currentstroke}%
\pgfsetdash{}{0pt}%
\pgfpathmoveto{\pgfqpoint{1.018663in}{0.827381in}}%
\pgfpathlineto{\pgfqpoint{1.016494in}{0.833952in}}%
\pgfpathlineto{\pgfqpoint{1.014314in}{0.840556in}}%
\pgfpathlineto{\pgfqpoint{1.050296in}{0.839670in}}%
\pgfpathlineto{\pgfqpoint{1.086468in}{0.839769in}}%
\pgfpathlineto{\pgfqpoint{1.086576in}{0.833186in}}%
\pgfpathlineto{\pgfqpoint{1.086684in}{0.826636in}}%
\pgfpathlineto{\pgfqpoint{1.052584in}{0.826543in}}%
\pgfpathlineto{\pgfqpoint{1.018663in}{0.827381in}}%
\pgfpathclose%
\pgfusepath{fill}%
\end{pgfscope}%
\begin{pgfscope}%
\pgfpathrectangle{\pgfqpoint{0.407592in}{0.222904in}}{\pgfqpoint{1.411163in}{1.411163in}}%
\pgfusepath{clip}%
\pgfsetbuttcap%
\pgfsetroundjoin%
\definecolor{currentfill}{rgb}{0.749020,0.772549,0.792157}%
\pgfsetfillcolor{currentfill}%
\pgfsetfillopacity{0.160000}%
\pgfsetlinewidth{0.000000pt}%
\definecolor{currentstroke}{rgb}{0.000000,0.000000,0.000000}%
\pgfsetstrokecolor{currentstroke}%
\pgfsetdash{}{0pt}%
\pgfpathmoveto{\pgfqpoint{1.391320in}{0.905052in}}%
\pgfpathlineto{\pgfqpoint{1.399165in}{0.912726in}}%
\pgfpathlineto{\pgfqpoint{1.407045in}{0.920434in}}%
\pgfpathlineto{\pgfqpoint{1.441679in}{0.929960in}}%
\pgfpathlineto{\pgfqpoint{1.474468in}{0.940318in}}%
\pgfpathlineto{\pgfqpoint{1.464968in}{0.932165in}}%
\pgfpathlineto{\pgfqpoint{1.455507in}{0.924044in}}%
\pgfpathlineto{\pgfqpoint{1.424283in}{0.914149in}}%
\pgfpathlineto{\pgfqpoint{1.391320in}{0.905052in}}%
\pgfpathclose%
\pgfusepath{fill}%
\end{pgfscope}%
\begin{pgfscope}%
\pgfpathrectangle{\pgfqpoint{0.407592in}{0.222904in}}{\pgfqpoint{1.411163in}{1.411163in}}%
\pgfusepath{clip}%
\pgfsetbuttcap%
\pgfsetroundjoin%
\definecolor{currentfill}{rgb}{0.749020,0.772549,0.792157}%
\pgfsetfillcolor{currentfill}%
\pgfsetfillopacity{0.160000}%
\pgfsetlinewidth{0.000000pt}%
\definecolor{currentstroke}{rgb}{0.000000,0.000000,0.000000}%
\pgfsetstrokecolor{currentstroke}%
\pgfsetdash{}{0pt}%
\pgfpathmoveto{\pgfqpoint{0.440076in}{1.076272in}}%
\pgfpathlineto{\pgfqpoint{0.426500in}{1.085997in}}%
\pgfpathlineto{\pgfqpoint{0.412881in}{1.095753in}}%
\pgfpathlineto{\pgfqpoint{0.435778in}{1.079471in}}%
\pgfpathlineto{\pgfqpoint{0.461940in}{1.063777in}}%
\pgfpathlineto{\pgfqpoint{0.474701in}{1.054571in}}%
\pgfpathlineto{\pgfqpoint{0.487416in}{1.045399in}}%
\pgfpathlineto{\pgfqpoint{0.462186in}{1.060548in}}%
\pgfpathlineto{\pgfqpoint{0.440076in}{1.076272in}}%
\pgfpathclose%
\pgfusepath{fill}%
\end{pgfscope}%
\begin{pgfscope}%
\pgfpathrectangle{\pgfqpoint{0.407592in}{0.222904in}}{\pgfqpoint{1.411163in}{1.411163in}}%
\pgfusepath{clip}%
\pgfsetbuttcap%
\pgfsetroundjoin%
\definecolor{currentfill}{rgb}{0.749020,0.772549,0.792157}%
\pgfsetfillcolor{currentfill}%
\pgfsetfillopacity{0.160000}%
\pgfsetlinewidth{0.000000pt}%
\definecolor{currentstroke}{rgb}{0.000000,0.000000,0.000000}%
\pgfsetstrokecolor{currentstroke}%
\pgfsetdash{}{0pt}%
\pgfpathmoveto{\pgfqpoint{1.229266in}{0.849906in}}%
\pgfpathlineto{\pgfqpoint{1.233266in}{0.856793in}}%
\pgfpathlineto{\pgfqpoint{1.237285in}{0.863712in}}%
\pgfpathlineto{\pgfqpoint{1.273544in}{0.868865in}}%
\pgfpathlineto{\pgfqpoint{1.308823in}{0.874956in}}%
\pgfpathlineto{\pgfqpoint{1.302868in}{0.867753in}}%
\pgfpathlineto{\pgfqpoint{1.296941in}{0.860584in}}%
\pgfpathlineto{\pgfqpoint{1.263563in}{0.854799in}}%
\pgfpathlineto{\pgfqpoint{1.229266in}{0.849906in}}%
\pgfpathclose%
\pgfusepath{fill}%
\end{pgfscope}%
\begin{pgfscope}%
\pgfpathrectangle{\pgfqpoint{0.407592in}{0.222904in}}{\pgfqpoint{1.411163in}{1.411163in}}%
\pgfusepath{clip}%
\pgfsetbuttcap%
\pgfsetroundjoin%
\definecolor{currentfill}{rgb}{0.749020,0.772549,0.792157}%
\pgfsetfillcolor{currentfill}%
\pgfsetfillopacity{0.160000}%
\pgfsetlinewidth{0.000000pt}%
\definecolor{currentstroke}{rgb}{0.000000,0.000000,0.000000}%
\pgfsetstrokecolor{currentstroke}%
\pgfsetdash{}{0pt}%
\pgfpathmoveto{\pgfqpoint{1.600155in}{1.015940in}}%
\pgfpathlineto{\pgfqpoint{1.611340in}{1.024833in}}%
\pgfpathlineto{\pgfqpoint{1.622568in}{1.033759in}}%
\pgfpathlineto{\pgfqpoint{1.653815in}{1.048388in}}%
\pgfpathlineto{\pgfqpoint{1.682166in}{1.063659in}}%
\pgfpathlineto{\pgfqpoint{1.669797in}{1.054191in}}%
\pgfpathlineto{\pgfqpoint{1.657470in}{1.044755in}}%
\pgfpathlineto{\pgfqpoint{1.630192in}{1.030035in}}%
\pgfpathlineto{\pgfqpoint{1.600155in}{1.015940in}}%
\pgfpathclose%
\pgfusepath{fill}%
\end{pgfscope}%
\begin{pgfscope}%
\pgfpathrectangle{\pgfqpoint{0.407592in}{0.222904in}}{\pgfqpoint{1.411163in}{1.411163in}}%
\pgfusepath{clip}%
\pgfsetbuttcap%
\pgfsetroundjoin%
\definecolor{currentfill}{rgb}{0.749020,0.772549,0.792157}%
\pgfsetfillcolor{currentfill}%
\pgfsetfillopacity{0.160000}%
\pgfsetlinewidth{0.000000pt}%
\definecolor{currentstroke}{rgb}{0.000000,0.000000,0.000000}%
\pgfsetstrokecolor{currentstroke}%
\pgfsetdash{}{0pt}%
\pgfpathmoveto{\pgfqpoint{1.308823in}{0.874956in}}%
\pgfpathlineto{\pgfqpoint{1.314805in}{0.882192in}}%
\pgfpathlineto{\pgfqpoint{1.320814in}{0.889462in}}%
\pgfpathlineto{\pgfqpoint{1.356774in}{0.896806in}}%
\pgfpathlineto{\pgfqpoint{1.391320in}{0.905052in}}%
\pgfpathlineto{\pgfqpoint{1.383509in}{0.897411in}}%
\pgfpathlineto{\pgfqpoint{1.375732in}{0.889804in}}%
\pgfpathlineto{\pgfqpoint{1.342942in}{0.881950in}}%
\pgfpathlineto{\pgfqpoint{1.308823in}{0.874956in}}%
\pgfpathclose%
\pgfusepath{fill}%
\end{pgfscope}%
\begin{pgfscope}%
\pgfpathrectangle{\pgfqpoint{0.407592in}{0.222904in}}{\pgfqpoint{1.411163in}{1.411163in}}%
\pgfusepath{clip}%
\pgfsetbuttcap%
\pgfsetroundjoin%
\definecolor{currentfill}{rgb}{0.749020,0.772549,0.792157}%
\pgfsetfillcolor{currentfill}%
\pgfsetfillopacity{0.160000}%
\pgfsetlinewidth{0.000000pt}%
\definecolor{currentstroke}{rgb}{0.000000,0.000000,0.000000}%
\pgfsetstrokecolor{currentstroke}%
\pgfsetdash{}{0pt}%
\pgfpathmoveto{\pgfqpoint{1.493586in}{0.956726in}}%
\pgfpathlineto{\pgfqpoint{1.503203in}{0.964981in}}%
\pgfpathlineto{\pgfqpoint{1.512861in}{0.973270in}}%
\pgfpathlineto{\pgfqpoint{1.546575in}{0.985383in}}%
\pgfpathlineto{\pgfqpoint{1.577911in}{0.998256in}}%
\pgfpathlineto{\pgfqpoint{1.566852in}{0.989463in}}%
\pgfpathlineto{\pgfqpoint{1.555834in}{0.980704in}}%
\pgfpathlineto{\pgfqpoint{1.525837in}{0.968349in}}%
\pgfpathlineto{\pgfqpoint{1.493586in}{0.956726in}}%
\pgfpathclose%
\pgfusepath{fill}%
\end{pgfscope}%
\begin{pgfscope}%
\pgfpathrectangle{\pgfqpoint{0.407592in}{0.222904in}}{\pgfqpoint{1.411163in}{1.411163in}}%
\pgfusepath{clip}%
\pgfsetbuttcap%
\pgfsetroundjoin%
\definecolor{currentfill}{rgb}{0.749020,0.772549,0.792157}%
\pgfsetfillcolor{currentfill}%
\pgfsetfillopacity{0.160000}%
\pgfsetlinewidth{0.000000pt}%
\definecolor{currentstroke}{rgb}{0.000000,0.000000,0.000000}%
\pgfsetstrokecolor{currentstroke}%
\pgfsetdash{}{0pt}%
\pgfpathmoveto{\pgfqpoint{0.637751in}{0.949156in}}%
\pgfpathlineto{\pgfqpoint{0.626550in}{0.957539in}}%
\pgfpathlineto{\pgfqpoint{0.615303in}{0.965956in}}%
\pgfpathlineto{\pgfqpoint{0.645249in}{0.954489in}}%
\pgfpathlineto{\pgfqpoint{0.677470in}{0.943837in}}%
\pgfpathlineto{\pgfqpoint{0.687310in}{0.935892in}}%
\pgfpathlineto{\pgfqpoint{0.697107in}{0.927980in}}%
\pgfpathlineto{\pgfqpoint{0.666353in}{0.938177in}}%
\pgfpathlineto{\pgfqpoint{0.637751in}{0.949156in}}%
\pgfpathclose%
\pgfusepath{fill}%
\end{pgfscope}%
\begin{pgfscope}%
\pgfpathrectangle{\pgfqpoint{0.407592in}{0.222904in}}{\pgfqpoint{1.411163in}{1.411163in}}%
\pgfusepath{clip}%
\pgfsetbuttcap%
\pgfsetroundjoin%
\definecolor{currentfill}{rgb}{0.749020,0.772549,0.792157}%
\pgfsetfillcolor{currentfill}%
\pgfsetfillopacity{0.160000}%
\pgfsetlinewidth{0.000000pt}%
\definecolor{currentstroke}{rgb}{0.000000,0.000000,0.000000}%
\pgfsetstrokecolor{currentstroke}%
\pgfsetdash{}{0pt}%
\pgfpathmoveto{\pgfqpoint{0.876011in}{0.853744in}}%
\pgfpathlineto{\pgfqpoint{0.869855in}{0.860733in}}%
\pgfpathlineto{\pgfqpoint{0.863671in}{0.867754in}}%
\pgfpathlineto{\pgfqpoint{0.898947in}{0.862794in}}%
\pgfpathlineto{\pgfqpoint{0.935207in}{0.858810in}}%
\pgfpathlineto{\pgfqpoint{0.939450in}{0.852015in}}%
\pgfpathlineto{\pgfqpoint{0.943672in}{0.845252in}}%
\pgfpathlineto{\pgfqpoint{0.909380in}{0.849035in}}%
\pgfpathlineto{\pgfqpoint{0.876011in}{0.853744in}}%
\pgfpathclose%
\pgfusepath{fill}%
\end{pgfscope}%
\begin{pgfscope}%
\pgfpathrectangle{\pgfqpoint{0.407592in}{0.222904in}}{\pgfqpoint{1.411163in}{1.411163in}}%
\pgfusepath{clip}%
\pgfsetbuttcap%
\pgfsetroundjoin%
\definecolor{currentfill}{rgb}{0.749020,0.772549,0.792157}%
\pgfsetfillcolor{currentfill}%
\pgfsetfillopacity{0.160000}%
\pgfsetlinewidth{0.000000pt}%
\definecolor{currentstroke}{rgb}{0.000000,0.000000,0.000000}%
\pgfsetstrokecolor{currentstroke}%
\pgfsetdash{}{0pt}%
\pgfpathmoveto{\pgfqpoint{0.537814in}{1.009042in}}%
\pgfpathlineto{\pgfqpoint{0.525283in}{1.018082in}}%
\pgfpathlineto{\pgfqpoint{0.512706in}{1.027154in}}%
\pgfpathlineto{\pgfqpoint{0.539892in}{1.013195in}}%
\pgfpathlineto{\pgfqpoint{0.569868in}{0.999959in}}%
\pgfpathlineto{\pgfqpoint{0.581295in}{0.991407in}}%
\pgfpathlineto{\pgfqpoint{0.592676in}{0.982890in}}%
\pgfpathlineto{\pgfqpoint{0.563918in}{0.995616in}}%
\pgfpathlineto{\pgfqpoint{0.537814in}{1.009042in}}%
\pgfpathclose%
\pgfusepath{fill}%
\end{pgfscope}%
\begin{pgfscope}%
\pgfpathrectangle{\pgfqpoint{0.407592in}{0.222904in}}{\pgfqpoint{1.411163in}{1.411163in}}%
\pgfusepath{clip}%
\pgfsetbuttcap%
\pgfsetroundjoin%
\definecolor{currentfill}{rgb}{0.749020,0.772549,0.792157}%
\pgfsetfillcolor{currentfill}%
\pgfsetfillopacity{0.160000}%
\pgfsetlinewidth{0.000000pt}%
\definecolor{currentstroke}{rgb}{0.000000,0.000000,0.000000}%
\pgfsetstrokecolor{currentstroke}%
\pgfsetdash{}{0pt}%
\pgfpathmoveto{\pgfqpoint{1.732064in}{1.101857in}}%
\pgfpathlineto{\pgfqpoint{1.744646in}{1.111488in}}%
\pgfpathlineto{\pgfqpoint{1.772311in}{1.128710in}}%
\pgfpathlineto{\pgfqpoint{1.759269in}{1.118802in}}%
\pgfusepath{fill}%
\end{pgfscope}%
\begin{pgfscope}%
\pgfpathrectangle{\pgfqpoint{0.407592in}{0.222904in}}{\pgfqpoint{1.411163in}{1.411163in}}%
\pgfusepath{clip}%
\pgfsetbuttcap%
\pgfsetroundjoin%
\definecolor{currentfill}{rgb}{0.749020,0.772549,0.792157}%
\pgfsetfillcolor{currentfill}%
\pgfsetfillopacity{0.160000}%
\pgfsetlinewidth{0.000000pt}%
\definecolor{currentstroke}{rgb}{0.000000,0.000000,0.000000}%
\pgfsetstrokecolor{currentstroke}%
\pgfsetdash{}{0pt}%
\pgfpathmoveto{\pgfqpoint{1.158622in}{0.842912in}}%
\pgfpathlineto{\pgfqpoint{1.160594in}{0.849611in}}%
\pgfpathlineto{\pgfqpoint{1.162575in}{0.856344in}}%
\pgfpathlineto{\pgfqpoint{1.200231in}{0.859530in}}%
\pgfpathlineto{\pgfqpoint{1.237285in}{0.863712in}}%
\pgfpathlineto{\pgfqpoint{1.233266in}{0.856793in}}%
\pgfpathlineto{\pgfqpoint{1.229266in}{0.849906in}}%
\pgfpathlineto{\pgfqpoint{1.194226in}{0.845936in}}%
\pgfpathlineto{\pgfqpoint{1.158622in}{0.842912in}}%
\pgfpathclose%
\pgfusepath{fill}%
\end{pgfscope}%
\begin{pgfscope}%
\pgfpathrectangle{\pgfqpoint{0.407592in}{0.222904in}}{\pgfqpoint{1.411163in}{1.411163in}}%
\pgfusepath{clip}%
\pgfsetbuttcap%
\pgfsetroundjoin%
\definecolor{currentfill}{rgb}{0.749020,0.772549,0.792157}%
\pgfsetfillcolor{currentfill}%
\pgfsetfillopacity{0.160000}%
\pgfsetlinewidth{0.000000pt}%
\definecolor{currentstroke}{rgb}{0.000000,0.000000,0.000000}%
\pgfsetstrokecolor{currentstroke}%
\pgfsetdash{}{0pt}%
\pgfpathmoveto{\pgfqpoint{0.716577in}{0.912258in}}%
\pgfpathlineto{\pgfqpoint{0.706863in}{0.920103in}}%
\pgfpathlineto{\pgfqpoint{0.697107in}{0.927980in}}%
\pgfpathlineto{\pgfqpoint{0.729874in}{0.918628in}}%
\pgfpathlineto{\pgfqpoint{0.764500in}{0.910176in}}%
\pgfpathlineto{\pgfqpoint{0.772626in}{0.902701in}}%
\pgfpathlineto{\pgfqpoint{0.780715in}{0.895259in}}%
\pgfpathlineto{\pgfqpoint{0.747770in}{0.903328in}}%
\pgfpathlineto{\pgfqpoint{0.716577in}{0.912258in}}%
\pgfpathclose%
\pgfusepath{fill}%
\end{pgfscope}%
\begin{pgfscope}%
\pgfpathrectangle{\pgfqpoint{0.407592in}{0.222904in}}{\pgfqpoint{1.411163in}{1.411163in}}%
\pgfusepath{clip}%
\pgfsetbuttcap%
\pgfsetroundjoin%
\definecolor{currentfill}{rgb}{0.749020,0.772549,0.792157}%
\pgfsetfillcolor{currentfill}%
\pgfsetfillopacity{0.160000}%
\pgfsetlinewidth{0.000000pt}%
\definecolor{currentstroke}{rgb}{0.000000,0.000000,0.000000}%
\pgfsetstrokecolor{currentstroke}%
\pgfsetdash{}{0pt}%
\pgfpathmoveto{\pgfqpoint{0.796785in}{0.880477in}}%
\pgfpathlineto{\pgfqpoint{0.788768in}{0.887851in}}%
\pgfpathlineto{\pgfqpoint{0.780715in}{0.895259in}}%
\pgfpathlineto{\pgfqpoint{0.815254in}{0.888101in}}%
\pgfpathlineto{\pgfqpoint{0.851214in}{0.881898in}}%
\pgfpathlineto{\pgfqpoint{0.857457in}{0.874809in}}%
\pgfpathlineto{\pgfqpoint{0.863671in}{0.867754in}}%
\pgfpathlineto{\pgfqpoint{0.829560in}{0.873661in}}%
\pgfpathlineto{\pgfqpoint{0.796785in}{0.880477in}}%
\pgfpathclose%
\pgfusepath{fill}%
\end{pgfscope}%
\begin{pgfscope}%
\pgfpathrectangle{\pgfqpoint{0.407592in}{0.222904in}}{\pgfqpoint{1.411163in}{1.411163in}}%
\pgfusepath{clip}%
\pgfsetbuttcap%
\pgfsetroundjoin%
\definecolor{currentfill}{rgb}{0.749020,0.772549,0.792157}%
\pgfsetfillcolor{currentfill}%
\pgfsetfillopacity{0.160000}%
\pgfsetlinewidth{0.000000pt}%
\definecolor{currentstroke}{rgb}{0.000000,0.000000,0.000000}%
\pgfsetstrokecolor{currentstroke}%
\pgfsetdash{}{0pt}%
\pgfpathmoveto{\pgfqpoint{0.412881in}{1.095753in}}%
\pgfpathlineto{\pgfqpoint{0.399217in}{1.105541in}}%
\pgfpathlineto{\pgfqpoint{0.385509in}{1.115361in}}%
\pgfpathlineto{\pgfqpoint{0.409186in}{1.098524in}}%
\pgfpathlineto{\pgfqpoint{0.436278in}{1.082289in}}%
\pgfpathlineto{\pgfqpoint{0.449132in}{1.073016in}}%
\pgfpathlineto{\pgfqpoint{0.461940in}{1.063777in}}%
\pgfpathlineto{\pgfqpoint{0.435778in}{1.079471in}}%
\pgfpathlineto{\pgfqpoint{0.412881in}{1.095753in}}%
\pgfpathclose%
\pgfusepath{fill}%
\end{pgfscope}%
\begin{pgfscope}%
\pgfpathrectangle{\pgfqpoint{0.407592in}{0.222904in}}{\pgfqpoint{1.411163in}{1.411163in}}%
\pgfusepath{clip}%
\pgfsetbuttcap%
\pgfsetroundjoin%
\definecolor{currentfill}{rgb}{0.749020,0.772549,0.792157}%
\pgfsetfillcolor{currentfill}%
\pgfsetfillopacity{0.160000}%
\pgfsetlinewidth{0.000000pt}%
\definecolor{currentstroke}{rgb}{0.000000,0.000000,0.000000}%
\pgfsetstrokecolor{currentstroke}%
\pgfsetdash{}{0pt}%
\pgfpathmoveto{\pgfqpoint{0.943672in}{0.845252in}}%
\pgfpathlineto{\pgfqpoint{0.939450in}{0.852015in}}%
\pgfpathlineto{\pgfqpoint{0.935207in}{0.858810in}}%
\pgfpathlineto{\pgfqpoint{0.972264in}{0.855826in}}%
\pgfpathlineto{\pgfqpoint{1.009923in}{0.853862in}}%
\pgfpathlineto{\pgfqpoint{1.012124in}{0.847192in}}%
\pgfpathlineto{\pgfqpoint{1.014314in}{0.840556in}}%
\pgfpathlineto{\pgfqpoint{0.978711in}{0.842420in}}%
\pgfpathlineto{\pgfqpoint{0.943672in}{0.845252in}}%
\pgfpathclose%
\pgfusepath{fill}%
\end{pgfscope}%
\begin{pgfscope}%
\pgfpathrectangle{\pgfqpoint{0.407592in}{0.222904in}}{\pgfqpoint{1.411163in}{1.411163in}}%
\pgfusepath{clip}%
\pgfsetbuttcap%
\pgfsetroundjoin%
\definecolor{currentfill}{rgb}{0.749020,0.772549,0.792157}%
\pgfsetfillcolor{currentfill}%
\pgfsetfillopacity{0.160000}%
\pgfsetlinewidth{0.000000pt}%
\definecolor{currentstroke}{rgb}{0.000000,0.000000,0.000000}%
\pgfsetstrokecolor{currentstroke}%
\pgfsetdash{}{0pt}%
\pgfpathmoveto{\pgfqpoint{1.622568in}{1.033759in}}%
\pgfpathlineto{\pgfqpoint{1.633838in}{1.042720in}}%
\pgfpathlineto{\pgfqpoint{1.645152in}{1.051715in}}%
\pgfpathlineto{\pgfqpoint{1.677609in}{1.066874in}}%
\pgfpathlineto{\pgfqpoint{1.707030in}{1.082693in}}%
\pgfpathlineto{\pgfqpoint{1.694577in}{1.073160in}}%
\pgfpathlineto{\pgfqpoint{1.682166in}{1.063659in}}%
\pgfpathlineto{\pgfqpoint{1.653815in}{1.048388in}}%
\pgfpathlineto{\pgfqpoint{1.622568in}{1.033759in}}%
\pgfpathclose%
\pgfusepath{fill}%
\end{pgfscope}%
\begin{pgfscope}%
\pgfpathrectangle{\pgfqpoint{0.407592in}{0.222904in}}{\pgfqpoint{1.411163in}{1.411163in}}%
\pgfusepath{clip}%
\pgfsetbuttcap%
\pgfsetroundjoin%
\definecolor{currentfill}{rgb}{0.749020,0.772549,0.792157}%
\pgfsetfillcolor{currentfill}%
\pgfsetfillopacity{0.160000}%
\pgfsetlinewidth{0.000000pt}%
\definecolor{currentstroke}{rgb}{0.000000,0.000000,0.000000}%
\pgfsetstrokecolor{currentstroke}%
\pgfsetdash{}{0pt}%
\pgfpathmoveto{\pgfqpoint{1.407045in}{0.920434in}}%
\pgfpathlineto{\pgfqpoint{1.414960in}{0.928176in}}%
\pgfpathlineto{\pgfqpoint{1.422909in}{0.935952in}}%
\pgfpathlineto{\pgfqpoint{1.459224in}{0.945906in}}%
\pgfpathlineto{\pgfqpoint{1.493586in}{0.956726in}}%
\pgfpathlineto{\pgfqpoint{1.484007in}{0.948506in}}%
\pgfpathlineto{\pgfqpoint{1.474468in}{0.940318in}}%
\pgfpathlineto{\pgfqpoint{1.441679in}{0.929960in}}%
\pgfpathlineto{\pgfqpoint{1.407045in}{0.920434in}}%
\pgfpathclose%
\pgfusepath{fill}%
\end{pgfscope}%
\begin{pgfscope}%
\pgfpathrectangle{\pgfqpoint{0.407592in}{0.222904in}}{\pgfqpoint{1.411163in}{1.411163in}}%
\pgfusepath{clip}%
\pgfsetbuttcap%
\pgfsetroundjoin%
\definecolor{currentfill}{rgb}{0.749020,0.772549,0.792157}%
\pgfsetfillcolor{currentfill}%
\pgfsetfillopacity{0.160000}%
\pgfsetlinewidth{0.000000pt}%
\definecolor{currentstroke}{rgb}{0.000000,0.000000,0.000000}%
\pgfsetstrokecolor{currentstroke}%
\pgfsetdash{}{0pt}%
\pgfpathmoveto{\pgfqpoint{1.086468in}{0.839769in}}%
\pgfpathlineto{\pgfqpoint{1.086360in}{0.846385in}}%
\pgfpathlineto{\pgfqpoint{1.086251in}{0.853033in}}%
\pgfpathlineto{\pgfqpoint{1.124515in}{0.854174in}}%
\pgfpathlineto{\pgfqpoint{1.162575in}{0.856344in}}%
\pgfpathlineto{\pgfqpoint{1.160594in}{0.849611in}}%
\pgfpathlineto{\pgfqpoint{1.158622in}{0.842912in}}%
\pgfpathlineto{\pgfqpoint{1.122640in}{0.840852in}}%
\pgfpathlineto{\pgfqpoint{1.086468in}{0.839769in}}%
\pgfpathclose%
\pgfusepath{fill}%
\end{pgfscope}%
\begin{pgfscope}%
\pgfpathrectangle{\pgfqpoint{0.407592in}{0.222904in}}{\pgfqpoint{1.411163in}{1.411163in}}%
\pgfusepath{clip}%
\pgfsetbuttcap%
\pgfsetroundjoin%
\definecolor{currentfill}{rgb}{0.749020,0.772549,0.792157}%
\pgfsetfillcolor{currentfill}%
\pgfsetfillopacity{0.160000}%
\pgfsetlinewidth{0.000000pt}%
\definecolor{currentstroke}{rgb}{0.000000,0.000000,0.000000}%
\pgfsetstrokecolor{currentstroke}%
\pgfsetdash{}{0pt}%
\pgfpathmoveto{\pgfqpoint{1.014314in}{0.840556in}}%
\pgfpathlineto{\pgfqpoint{1.012124in}{0.847192in}}%
\pgfpathlineto{\pgfqpoint{1.009923in}{0.853862in}}%
\pgfpathlineto{\pgfqpoint{1.047986in}{0.852929in}}%
\pgfpathlineto{\pgfqpoint{1.086251in}{0.853033in}}%
\pgfpathlineto{\pgfqpoint{1.086360in}{0.846385in}}%
\pgfpathlineto{\pgfqpoint{1.086468in}{0.839769in}}%
\pgfpathlineto{\pgfqpoint{1.050296in}{0.839670in}}%
\pgfpathlineto{\pgfqpoint{1.014314in}{0.840556in}}%
\pgfpathclose%
\pgfusepath{fill}%
\end{pgfscope}%
\begin{pgfscope}%
\pgfpathrectangle{\pgfqpoint{0.407592in}{0.222904in}}{\pgfqpoint{1.411163in}{1.411163in}}%
\pgfusepath{clip}%
\pgfsetbuttcap%
\pgfsetroundjoin%
\definecolor{currentfill}{rgb}{0.749020,0.772549,0.792157}%
\pgfsetfillcolor{currentfill}%
\pgfsetfillopacity{0.160000}%
\pgfsetlinewidth{0.000000pt}%
\definecolor{currentstroke}{rgb}{0.000000,0.000000,0.000000}%
\pgfsetstrokecolor{currentstroke}%
\pgfsetdash{}{0pt}%
\pgfpathmoveto{\pgfqpoint{1.237285in}{0.863712in}}%
\pgfpathlineto{\pgfqpoint{1.241324in}{0.870664in}}%
\pgfpathlineto{\pgfqpoint{1.245381in}{0.877651in}}%
\pgfpathlineto{\pgfqpoint{1.283620in}{0.883064in}}%
\pgfpathlineto{\pgfqpoint{1.320814in}{0.889462in}}%
\pgfpathlineto{\pgfqpoint{1.314805in}{0.882192in}}%
\pgfpathlineto{\pgfqpoint{1.308823in}{0.874956in}}%
\pgfpathlineto{\pgfqpoint{1.273544in}{0.868865in}}%
\pgfpathlineto{\pgfqpoint{1.237285in}{0.863712in}}%
\pgfpathclose%
\pgfusepath{fill}%
\end{pgfscope}%
\begin{pgfscope}%
\pgfpathrectangle{\pgfqpoint{0.407592in}{0.222904in}}{\pgfqpoint{1.411163in}{1.411163in}}%
\pgfusepath{clip}%
\pgfsetbuttcap%
\pgfsetroundjoin%
\definecolor{currentfill}{rgb}{0.749020,0.772549,0.792157}%
\pgfsetfillcolor{currentfill}%
\pgfsetfillopacity{0.160000}%
\pgfsetlinewidth{0.000000pt}%
\definecolor{currentstroke}{rgb}{0.000000,0.000000,0.000000}%
\pgfsetstrokecolor{currentstroke}%
\pgfsetdash{}{0pt}%
\pgfpathmoveto{\pgfqpoint{1.512861in}{0.973270in}}%
\pgfpathlineto{\pgfqpoint{1.522558in}{0.981593in}}%
\pgfpathlineto{\pgfqpoint{1.532296in}{0.989951in}}%
\pgfpathlineto{\pgfqpoint{1.567477in}{1.002553in}}%
\pgfpathlineto{\pgfqpoint{1.600155in}{1.015940in}}%
\pgfpathlineto{\pgfqpoint{1.589012in}{1.007081in}}%
\pgfpathlineto{\pgfqpoint{1.577911in}{0.998256in}}%
\pgfpathlineto{\pgfqpoint{1.546575in}{0.985383in}}%
\pgfpathlineto{\pgfqpoint{1.512861in}{0.973270in}}%
\pgfpathclose%
\pgfusepath{fill}%
\end{pgfscope}%
\begin{pgfscope}%
\pgfpathrectangle{\pgfqpoint{0.407592in}{0.222904in}}{\pgfqpoint{1.411163in}{1.411163in}}%
\pgfusepath{clip}%
\pgfsetbuttcap%
\pgfsetroundjoin%
\definecolor{currentfill}{rgb}{0.749020,0.772549,0.792157}%
\pgfsetfillcolor{currentfill}%
\pgfsetfillopacity{0.160000}%
\pgfsetlinewidth{0.000000pt}%
\definecolor{currentstroke}{rgb}{0.000000,0.000000,0.000000}%
\pgfsetstrokecolor{currentstroke}%
\pgfsetdash{}{0pt}%
\pgfpathmoveto{\pgfqpoint{1.320814in}{0.889462in}}%
\pgfpathlineto{\pgfqpoint{1.326852in}{0.896765in}}%
\pgfpathlineto{\pgfqpoint{1.332917in}{0.904103in}}%
\pgfpathlineto{\pgfqpoint{1.370732in}{0.911797in}}%
\pgfpathlineto{\pgfqpoint{1.407045in}{0.920434in}}%
\pgfpathlineto{\pgfqpoint{1.399165in}{0.912726in}}%
\pgfpathlineto{\pgfqpoint{1.391320in}{0.905052in}}%
\pgfpathlineto{\pgfqpoint{1.356774in}{0.896806in}}%
\pgfpathlineto{\pgfqpoint{1.320814in}{0.889462in}}%
\pgfpathclose%
\pgfusepath{fill}%
\end{pgfscope}%
\begin{pgfscope}%
\pgfpathrectangle{\pgfqpoint{0.407592in}{0.222904in}}{\pgfqpoint{1.411163in}{1.411163in}}%
\pgfusepath{clip}%
\pgfsetbuttcap%
\pgfsetroundjoin%
\definecolor{currentfill}{rgb}{0.749020,0.772549,0.792157}%
\pgfsetfillcolor{currentfill}%
\pgfsetfillopacity{0.160000}%
\pgfsetlinewidth{0.000000pt}%
\definecolor{currentstroke}{rgb}{0.000000,0.000000,0.000000}%
\pgfsetstrokecolor{currentstroke}%
\pgfsetdash{}{0pt}%
\pgfpathmoveto{\pgfqpoint{0.615303in}{0.965956in}}%
\pgfpathlineto{\pgfqpoint{0.604012in}{0.974406in}}%
\pgfpathlineto{\pgfqpoint{0.592676in}{0.982890in}}%
\pgfpathlineto{\pgfqpoint{0.623970in}{0.970938in}}%
\pgfpathlineto{\pgfqpoint{0.657663in}{0.959832in}}%
\pgfpathlineto{\pgfqpoint{0.667588in}{0.951817in}}%
\pgfpathlineto{\pgfqpoint{0.677470in}{0.943837in}}%
\pgfpathlineto{\pgfqpoint{0.645249in}{0.954489in}}%
\pgfpathlineto{\pgfqpoint{0.615303in}{0.965956in}}%
\pgfpathclose%
\pgfusepath{fill}%
\end{pgfscope}%
\begin{pgfscope}%
\pgfpathrectangle{\pgfqpoint{0.407592in}{0.222904in}}{\pgfqpoint{1.411163in}{1.411163in}}%
\pgfusepath{clip}%
\pgfsetbuttcap%
\pgfsetroundjoin%
\definecolor{currentfill}{rgb}{0.749020,0.772549,0.792157}%
\pgfsetfillcolor{currentfill}%
\pgfsetfillopacity{0.160000}%
\pgfsetlinewidth{0.000000pt}%
\definecolor{currentstroke}{rgb}{0.000000,0.000000,0.000000}%
\pgfsetstrokecolor{currentstroke}%
\pgfsetdash{}{0pt}%
\pgfpathmoveto{\pgfqpoint{0.512706in}{1.027154in}}%
\pgfpathlineto{\pgfqpoint{0.500084in}{1.036260in}}%
\pgfpathlineto{\pgfqpoint{0.487416in}{1.045399in}}%
\pgfpathlineto{\pgfqpoint{0.515681in}{1.030910in}}%
\pgfpathlineto{\pgfqpoint{0.546875in}{1.017166in}}%
\pgfpathlineto{\pgfqpoint{0.558395in}{1.008545in}}%
\pgfpathlineto{\pgfqpoint{0.569868in}{0.999959in}}%
\pgfpathlineto{\pgfqpoint{0.539892in}{1.013195in}}%
\pgfpathlineto{\pgfqpoint{0.512706in}{1.027154in}}%
\pgfpathclose%
\pgfusepath{fill}%
\end{pgfscope}%
\begin{pgfscope}%
\pgfpathrectangle{\pgfqpoint{0.407592in}{0.222904in}}{\pgfqpoint{1.411163in}{1.411163in}}%
\pgfusepath{clip}%
\pgfsetbuttcap%
\pgfsetroundjoin%
\definecolor{currentfill}{rgb}{0.749020,0.772549,0.792157}%
\pgfsetfillcolor{currentfill}%
\pgfsetfillopacity{0.160000}%
\pgfsetlinewidth{0.000000pt}%
\definecolor{currentstroke}{rgb}{0.000000,0.000000,0.000000}%
\pgfsetstrokecolor{currentstroke}%
\pgfsetdash{}{0pt}%
\pgfpathmoveto{\pgfqpoint{0.385509in}{1.115361in}}%
\pgfpathlineto{\pgfqpoint{0.371757in}{1.125212in}}%
\pgfpathlineto{\pgfqpoint{0.395821in}{1.108101in}}%
\pgfpathlineto{\pgfqpoint{0.423376in}{1.091597in}}%
\pgfpathlineto{\pgfqpoint{0.436278in}{1.082289in}}%
\pgfpathlineto{\pgfqpoint{0.409186in}{1.098524in}}%
\pgfusepath{fill}%
\end{pgfscope}%
\begin{pgfscope}%
\pgfpathrectangle{\pgfqpoint{0.407592in}{0.222904in}}{\pgfqpoint{1.411163in}{1.411163in}}%
\pgfusepath{clip}%
\pgfsetbuttcap%
\pgfsetroundjoin%
\definecolor{currentfill}{rgb}{0.749020,0.772549,0.792157}%
\pgfsetfillcolor{currentfill}%
\pgfsetfillopacity{0.160000}%
\pgfsetlinewidth{0.000000pt}%
\definecolor{currentstroke}{rgb}{0.000000,0.000000,0.000000}%
\pgfsetstrokecolor{currentstroke}%
\pgfsetdash{}{0pt}%
\pgfpathmoveto{\pgfqpoint{0.863671in}{0.867754in}}%
\pgfpathlineto{\pgfqpoint{0.857457in}{0.874809in}}%
\pgfpathlineto{\pgfqpoint{0.851214in}{0.881898in}}%
\pgfpathlineto{\pgfqpoint{0.888413in}{0.876686in}}%
\pgfpathlineto{\pgfqpoint{0.926659in}{0.872500in}}%
\pgfpathlineto{\pgfqpoint{0.930944in}{0.865638in}}%
\pgfpathlineto{\pgfqpoint{0.935207in}{0.858810in}}%
\pgfpathlineto{\pgfqpoint{0.898947in}{0.862794in}}%
\pgfpathlineto{\pgfqpoint{0.863671in}{0.867754in}}%
\pgfpathclose%
\pgfusepath{fill}%
\end{pgfscope}%
\begin{pgfscope}%
\pgfpathrectangle{\pgfqpoint{0.407592in}{0.222904in}}{\pgfqpoint{1.411163in}{1.411163in}}%
\pgfusepath{clip}%
\pgfsetbuttcap%
\pgfsetroundjoin%
\definecolor{currentfill}{rgb}{0.749020,0.772549,0.792157}%
\pgfsetfillcolor{currentfill}%
\pgfsetfillopacity{0.160000}%
\pgfsetlinewidth{0.000000pt}%
\definecolor{currentstroke}{rgb}{0.000000,0.000000,0.000000}%
\pgfsetstrokecolor{currentstroke}%
\pgfsetdash{}{0pt}%
\pgfpathmoveto{\pgfqpoint{1.645152in}{1.051715in}}%
\pgfpathlineto{\pgfqpoint{1.656509in}{1.060744in}}%
\pgfpathlineto{\pgfqpoint{1.667909in}{1.069807in}}%
\pgfpathlineto{\pgfqpoint{1.701576in}{1.085494in}}%
\pgfpathlineto{\pgfqpoint{1.732064in}{1.101857in}}%
\pgfpathlineto{\pgfqpoint{1.719526in}{1.092258in}}%
\pgfpathlineto{\pgfqpoint{1.707030in}{1.082693in}}%
\pgfpathlineto{\pgfqpoint{1.677609in}{1.066874in}}%
\pgfpathlineto{\pgfqpoint{1.645152in}{1.051715in}}%
\pgfpathclose%
\pgfusepath{fill}%
\end{pgfscope}%
\begin{pgfscope}%
\pgfpathrectangle{\pgfqpoint{0.407592in}{0.222904in}}{\pgfqpoint{1.411163in}{1.411163in}}%
\pgfusepath{clip}%
\pgfsetbuttcap%
\pgfsetroundjoin%
\definecolor{currentfill}{rgb}{0.749020,0.772549,0.792157}%
\pgfsetfillcolor{currentfill}%
\pgfsetfillopacity{0.160000}%
\pgfsetlinewidth{0.000000pt}%
\definecolor{currentstroke}{rgb}{0.000000,0.000000,0.000000}%
\pgfsetstrokecolor{currentstroke}%
\pgfsetdash{}{0pt}%
\pgfpathmoveto{\pgfqpoint{0.697107in}{0.927980in}}%
\pgfpathlineto{\pgfqpoint{0.687310in}{0.935892in}}%
\pgfpathlineto{\pgfqpoint{0.677470in}{0.943837in}}%
\pgfpathlineto{\pgfqpoint{0.711820in}{0.934065in}}%
\pgfpathlineto{\pgfqpoint{0.748135in}{0.925230in}}%
\pgfpathlineto{\pgfqpoint{0.756336in}{0.917686in}}%
\pgfpathlineto{\pgfqpoint{0.764500in}{0.910176in}}%
\pgfpathlineto{\pgfqpoint{0.729874in}{0.918628in}}%
\pgfpathlineto{\pgfqpoint{0.697107in}{0.927980in}}%
\pgfpathclose%
\pgfusepath{fill}%
\end{pgfscope}%
\begin{pgfscope}%
\pgfpathrectangle{\pgfqpoint{0.407592in}{0.222904in}}{\pgfqpoint{1.411163in}{1.411163in}}%
\pgfusepath{clip}%
\pgfsetbuttcap%
\pgfsetroundjoin%
\definecolor{currentfill}{rgb}{0.749020,0.772549,0.792157}%
\pgfsetfillcolor{currentfill}%
\pgfsetfillopacity{0.160000}%
\pgfsetlinewidth{0.000000pt}%
\definecolor{currentstroke}{rgb}{0.000000,0.000000,0.000000}%
\pgfsetstrokecolor{currentstroke}%
\pgfsetdash{}{0pt}%
\pgfpathmoveto{\pgfqpoint{1.162575in}{0.856344in}}%
\pgfpathlineto{\pgfqpoint{1.164567in}{0.863109in}}%
\pgfpathlineto{\pgfqpoint{1.166568in}{0.869908in}}%
\pgfpathlineto{\pgfqpoint{1.206296in}{0.873256in}}%
\pgfpathlineto{\pgfqpoint{1.245381in}{0.877651in}}%
\pgfpathlineto{\pgfqpoint{1.241324in}{0.870664in}}%
\pgfpathlineto{\pgfqpoint{1.237285in}{0.863712in}}%
\pgfpathlineto{\pgfqpoint{1.200231in}{0.859530in}}%
\pgfpathlineto{\pgfqpoint{1.162575in}{0.856344in}}%
\pgfpathclose%
\pgfusepath{fill}%
\end{pgfscope}%
\begin{pgfscope}%
\pgfpathrectangle{\pgfqpoint{0.407592in}{0.222904in}}{\pgfqpoint{1.411163in}{1.411163in}}%
\pgfusepath{clip}%
\pgfsetbuttcap%
\pgfsetroundjoin%
\definecolor{currentfill}{rgb}{0.749020,0.772549,0.792157}%
\pgfsetfillcolor{currentfill}%
\pgfsetfillopacity{0.160000}%
\pgfsetlinewidth{0.000000pt}%
\definecolor{currentstroke}{rgb}{0.000000,0.000000,0.000000}%
\pgfsetstrokecolor{currentstroke}%
\pgfsetdash{}{0pt}%
\pgfpathmoveto{\pgfqpoint{0.780715in}{0.895259in}}%
\pgfpathlineto{\pgfqpoint{0.772626in}{0.902701in}}%
\pgfpathlineto{\pgfqpoint{0.764500in}{0.910176in}}%
\pgfpathlineto{\pgfqpoint{0.800814in}{0.902677in}}%
\pgfpathlineto{\pgfqpoint{0.838638in}{0.896175in}}%
\pgfpathlineto{\pgfqpoint{0.844941in}{0.889020in}}%
\pgfpathlineto{\pgfqpoint{0.851214in}{0.881898in}}%
\pgfpathlineto{\pgfqpoint{0.815254in}{0.888101in}}%
\pgfpathlineto{\pgfqpoint{0.780715in}{0.895259in}}%
\pgfpathclose%
\pgfusepath{fill}%
\end{pgfscope}%
\begin{pgfscope}%
\pgfpathrectangle{\pgfqpoint{0.407592in}{0.222904in}}{\pgfqpoint{1.411163in}{1.411163in}}%
\pgfusepath{clip}%
\pgfsetbuttcap%
\pgfsetroundjoin%
\definecolor{currentfill}{rgb}{0.749020,0.772549,0.792157}%
\pgfsetfillcolor{currentfill}%
\pgfsetfillopacity{0.160000}%
\pgfsetlinewidth{0.000000pt}%
\definecolor{currentstroke}{rgb}{0.000000,0.000000,0.000000}%
\pgfsetstrokecolor{currentstroke}%
\pgfsetdash{}{0pt}%
\pgfpathmoveto{\pgfqpoint{1.422909in}{0.935952in}}%
\pgfpathlineto{\pgfqpoint{1.430894in}{0.943762in}}%
\pgfpathlineto{\pgfqpoint{1.438914in}{0.951607in}}%
\pgfpathlineto{\pgfqpoint{1.476919in}{0.961989in}}%
\pgfpathlineto{\pgfqpoint{1.512861in}{0.973270in}}%
\pgfpathlineto{\pgfqpoint{1.503203in}{0.964981in}}%
\pgfpathlineto{\pgfqpoint{1.493586in}{0.956726in}}%
\pgfpathlineto{\pgfqpoint{1.459224in}{0.945906in}}%
\pgfpathlineto{\pgfqpoint{1.422909in}{0.935952in}}%
\pgfpathclose%
\pgfusepath{fill}%
\end{pgfscope}%
\begin{pgfscope}%
\pgfpathrectangle{\pgfqpoint{0.407592in}{0.222904in}}{\pgfqpoint{1.411163in}{1.411163in}}%
\pgfusepath{clip}%
\pgfsetbuttcap%
\pgfsetroundjoin%
\definecolor{currentfill}{rgb}{0.749020,0.772549,0.792157}%
\pgfsetfillcolor{currentfill}%
\pgfsetfillopacity{0.160000}%
\pgfsetlinewidth{0.000000pt}%
\definecolor{currentstroke}{rgb}{0.000000,0.000000,0.000000}%
\pgfsetstrokecolor{currentstroke}%
\pgfsetdash{}{0pt}%
\pgfpathmoveto{\pgfqpoint{0.935207in}{0.858810in}}%
\pgfpathlineto{\pgfqpoint{0.930944in}{0.865638in}}%
\pgfpathlineto{\pgfqpoint{0.926659in}{0.872500in}}%
\pgfpathlineto{\pgfqpoint{0.965753in}{0.869365in}}%
\pgfpathlineto{\pgfqpoint{1.005488in}{0.867300in}}%
\pgfpathlineto{\pgfqpoint{1.007711in}{0.860564in}}%
\pgfpathlineto{\pgfqpoint{1.009923in}{0.853862in}}%
\pgfpathlineto{\pgfqpoint{0.972264in}{0.855826in}}%
\pgfpathlineto{\pgfqpoint{0.935207in}{0.858810in}}%
\pgfpathclose%
\pgfusepath{fill}%
\end{pgfscope}%
\begin{pgfscope}%
\pgfpathrectangle{\pgfqpoint{0.407592in}{0.222904in}}{\pgfqpoint{1.411163in}{1.411163in}}%
\pgfusepath{clip}%
\pgfsetbuttcap%
\pgfsetroundjoin%
\definecolor{currentfill}{rgb}{0.749020,0.772549,0.792157}%
\pgfsetfillcolor{currentfill}%
\pgfsetfillopacity{0.160000}%
\pgfsetlinewidth{0.000000pt}%
\definecolor{currentstroke}{rgb}{0.000000,0.000000,0.000000}%
\pgfsetstrokecolor{currentstroke}%
\pgfsetdash{}{0pt}%
\pgfpathmoveto{\pgfqpoint{1.532296in}{0.989951in}}%
\pgfpathlineto{\pgfqpoint{1.542074in}{0.998343in}}%
\pgfpathlineto{\pgfqpoint{1.551892in}{1.006770in}}%
\pgfpathlineto{\pgfqpoint{1.588546in}{1.019859in}}%
\pgfpathlineto{\pgfqpoint{1.622568in}{1.033759in}}%
\pgfpathlineto{\pgfqpoint{1.611340in}{1.024833in}}%
\pgfpathlineto{\pgfqpoint{1.600155in}{1.015940in}}%
\pgfpathlineto{\pgfqpoint{1.567477in}{1.002553in}}%
\pgfpathlineto{\pgfqpoint{1.532296in}{0.989951in}}%
\pgfpathclose%
\pgfusepath{fill}%
\end{pgfscope}%
\begin{pgfscope}%
\pgfpathrectangle{\pgfqpoint{0.407592in}{0.222904in}}{\pgfqpoint{1.411163in}{1.411163in}}%
\pgfusepath{clip}%
\pgfsetbuttcap%
\pgfsetroundjoin%
\definecolor{currentfill}{rgb}{0.749020,0.772549,0.792157}%
\pgfsetfillcolor{currentfill}%
\pgfsetfillopacity{0.160000}%
\pgfsetlinewidth{0.000000pt}%
\definecolor{currentstroke}{rgb}{0.000000,0.000000,0.000000}%
\pgfsetstrokecolor{currentstroke}%
\pgfsetdash{}{0pt}%
\pgfpathmoveto{\pgfqpoint{1.086251in}{0.853033in}}%
\pgfpathlineto{\pgfqpoint{1.086141in}{0.859714in}}%
\pgfpathlineto{\pgfqpoint{1.086031in}{0.866429in}}%
\pgfpathlineto{\pgfqpoint{1.126408in}{0.867628in}}%
\pgfpathlineto{\pgfqpoint{1.166568in}{0.869908in}}%
\pgfpathlineto{\pgfqpoint{1.164567in}{0.863109in}}%
\pgfpathlineto{\pgfqpoint{1.162575in}{0.856344in}}%
\pgfpathlineto{\pgfqpoint{1.124515in}{0.854174in}}%
\pgfpathlineto{\pgfqpoint{1.086251in}{0.853033in}}%
\pgfpathclose%
\pgfusepath{fill}%
\end{pgfscope}%
\begin{pgfscope}%
\pgfpathrectangle{\pgfqpoint{0.407592in}{0.222904in}}{\pgfqpoint{1.411163in}{1.411163in}}%
\pgfusepath{clip}%
\pgfsetbuttcap%
\pgfsetroundjoin%
\definecolor{currentfill}{rgb}{0.749020,0.772549,0.792157}%
\pgfsetfillcolor{currentfill}%
\pgfsetfillopacity{0.160000}%
\pgfsetlinewidth{0.000000pt}%
\definecolor{currentstroke}{rgb}{0.000000,0.000000,0.000000}%
\pgfsetstrokecolor{currentstroke}%
\pgfsetdash{}{0pt}%
\pgfpathmoveto{\pgfqpoint{1.245381in}{0.877651in}}%
\pgfpathlineto{\pgfqpoint{1.249459in}{0.884670in}}%
\pgfpathlineto{\pgfqpoint{1.253556in}{0.891724in}}%
\pgfpathlineto{\pgfqpoint{1.293792in}{0.897398in}}%
\pgfpathlineto{\pgfqpoint{1.332917in}{0.904103in}}%
\pgfpathlineto{\pgfqpoint{1.326852in}{0.896765in}}%
\pgfpathlineto{\pgfqpoint{1.320814in}{0.889462in}}%
\pgfpathlineto{\pgfqpoint{1.283620in}{0.883064in}}%
\pgfpathlineto{\pgfqpoint{1.245381in}{0.877651in}}%
\pgfpathclose%
\pgfusepath{fill}%
\end{pgfscope}%
\begin{pgfscope}%
\pgfpathrectangle{\pgfqpoint{0.407592in}{0.222904in}}{\pgfqpoint{1.411163in}{1.411163in}}%
\pgfusepath{clip}%
\pgfsetbuttcap%
\pgfsetroundjoin%
\definecolor{currentfill}{rgb}{0.749020,0.772549,0.792157}%
\pgfsetfillcolor{currentfill}%
\pgfsetfillopacity{0.160000}%
\pgfsetlinewidth{0.000000pt}%
\definecolor{currentstroke}{rgb}{0.000000,0.000000,0.000000}%
\pgfsetstrokecolor{currentstroke}%
\pgfsetdash{}{0pt}%
\pgfpathmoveto{\pgfqpoint{0.487416in}{1.045399in}}%
\pgfpathlineto{\pgfqpoint{0.474701in}{1.054571in}}%
\pgfpathlineto{\pgfqpoint{0.461940in}{1.063777in}}%
\pgfpathlineto{\pgfqpoint{0.491283in}{1.048761in}}%
\pgfpathlineto{\pgfqpoint{0.523697in}{1.034512in}}%
\pgfpathlineto{\pgfqpoint{0.535310in}{1.025821in}}%
\pgfpathlineto{\pgfqpoint{0.546875in}{1.017166in}}%
\pgfpathlineto{\pgfqpoint{0.515681in}{1.030910in}}%
\pgfpathlineto{\pgfqpoint{0.487416in}{1.045399in}}%
\pgfpathclose%
\pgfusepath{fill}%
\end{pgfscope}%
\begin{pgfscope}%
\pgfpathrectangle{\pgfqpoint{0.407592in}{0.222904in}}{\pgfqpoint{1.411163in}{1.411163in}}%
\pgfusepath{clip}%
\pgfsetbuttcap%
\pgfsetroundjoin%
\definecolor{currentfill}{rgb}{0.749020,0.772549,0.792157}%
\pgfsetfillcolor{currentfill}%
\pgfsetfillopacity{0.160000}%
\pgfsetlinewidth{0.000000pt}%
\definecolor{currentstroke}{rgb}{0.000000,0.000000,0.000000}%
\pgfsetstrokecolor{currentstroke}%
\pgfsetdash{}{0pt}%
\pgfpathmoveto{\pgfqpoint{1.009923in}{0.853862in}}%
\pgfpathlineto{\pgfqpoint{1.007711in}{0.860564in}}%
\pgfpathlineto{\pgfqpoint{1.005488in}{0.867300in}}%
\pgfpathlineto{\pgfqpoint{1.045652in}{0.866319in}}%
\pgfpathlineto{\pgfqpoint{1.086031in}{0.866429in}}%
\pgfpathlineto{\pgfqpoint{1.086141in}{0.859714in}}%
\pgfpathlineto{\pgfqpoint{1.086251in}{0.853033in}}%
\pgfpathlineto{\pgfqpoint{1.047986in}{0.852929in}}%
\pgfpathlineto{\pgfqpoint{1.009923in}{0.853862in}}%
\pgfpathclose%
\pgfusepath{fill}%
\end{pgfscope}%
\begin{pgfscope}%
\pgfpathrectangle{\pgfqpoint{0.407592in}{0.222904in}}{\pgfqpoint{1.411163in}{1.411163in}}%
\pgfusepath{clip}%
\pgfsetbuttcap%
\pgfsetroundjoin%
\definecolor{currentfill}{rgb}{0.749020,0.772549,0.792157}%
\pgfsetfillcolor{currentfill}%
\pgfsetfillopacity{0.160000}%
\pgfsetlinewidth{0.000000pt}%
\definecolor{currentstroke}{rgb}{0.000000,0.000000,0.000000}%
\pgfsetstrokecolor{currentstroke}%
\pgfsetdash{}{0pt}%
\pgfpathmoveto{\pgfqpoint{1.332917in}{0.904103in}}%
\pgfpathlineto{\pgfqpoint{1.339011in}{0.911474in}}%
\pgfpathlineto{\pgfqpoint{1.345134in}{0.918880in}}%
\pgfpathlineto{\pgfqpoint{1.384817in}{0.926925in}}%
\pgfpathlineto{\pgfqpoint{1.422909in}{0.935952in}}%
\pgfpathlineto{\pgfqpoint{1.414960in}{0.928176in}}%
\pgfpathlineto{\pgfqpoint{1.407045in}{0.920434in}}%
\pgfpathlineto{\pgfqpoint{1.370732in}{0.911797in}}%
\pgfpathlineto{\pgfqpoint{1.332917in}{0.904103in}}%
\pgfpathclose%
\pgfusepath{fill}%
\end{pgfscope}%
\begin{pgfscope}%
\pgfpathrectangle{\pgfqpoint{0.407592in}{0.222904in}}{\pgfqpoint{1.411163in}{1.411163in}}%
\pgfusepath{clip}%
\pgfsetbuttcap%
\pgfsetroundjoin%
\definecolor{currentfill}{rgb}{0.749020,0.772549,0.792157}%
\pgfsetfillcolor{currentfill}%
\pgfsetfillopacity{0.160000}%
\pgfsetlinewidth{0.000000pt}%
\definecolor{currentstroke}{rgb}{0.000000,0.000000,0.000000}%
\pgfsetstrokecolor{currentstroke}%
\pgfsetdash{}{0pt}%
\pgfpathmoveto{\pgfqpoint{0.592676in}{0.982890in}}%
\pgfpathlineto{\pgfqpoint{0.581295in}{0.991407in}}%
\pgfpathlineto{\pgfqpoint{0.569868in}{0.999959in}}%
\pgfpathlineto{\pgfqpoint{0.602512in}{0.987524in}}%
\pgfpathlineto{\pgfqpoint{0.637684in}{0.975965in}}%
\pgfpathlineto{\pgfqpoint{0.647695in}{0.967881in}}%
\pgfpathlineto{\pgfqpoint{0.657663in}{0.959832in}}%
\pgfpathlineto{\pgfqpoint{0.623970in}{0.970938in}}%
\pgfpathlineto{\pgfqpoint{0.592676in}{0.982890in}}%
\pgfpathclose%
\pgfusepath{fill}%
\end{pgfscope}%
\begin{pgfscope}%
\pgfpathrectangle{\pgfqpoint{0.407592in}{0.222904in}}{\pgfqpoint{1.411163in}{1.411163in}}%
\pgfusepath{clip}%
\pgfsetbuttcap%
\pgfsetroundjoin%
\definecolor{currentfill}{rgb}{0.749020,0.772549,0.792157}%
\pgfsetfillcolor{currentfill}%
\pgfsetfillopacity{0.160000}%
\pgfsetlinewidth{0.000000pt}%
\definecolor{currentstroke}{rgb}{0.000000,0.000000,0.000000}%
\pgfsetstrokecolor{currentstroke}%
\pgfsetdash{}{0pt}%
\pgfpathmoveto{\pgfqpoint{1.667909in}{1.069807in}}%
\pgfpathlineto{\pgfqpoint{1.679353in}{1.078906in}}%
\pgfpathlineto{\pgfqpoint{1.713624in}{1.094855in}}%
\pgfpathlineto{\pgfqpoint{1.744646in}{1.111488in}}%
\pgfpathlineto{\pgfqpoint{1.732064in}{1.101857in}}%
\pgfpathlineto{\pgfqpoint{1.701576in}{1.085494in}}%
\pgfusepath{fill}%
\end{pgfscope}%
\begin{pgfscope}%
\pgfpathrectangle{\pgfqpoint{0.407592in}{0.222904in}}{\pgfqpoint{1.411163in}{1.411163in}}%
\pgfusepath{clip}%
\pgfsetbuttcap%
\pgfsetroundjoin%
\definecolor{currentfill}{rgb}{0.749020,0.772549,0.792157}%
\pgfsetfillcolor{currentfill}%
\pgfsetfillopacity{0.160000}%
\pgfsetlinewidth{0.000000pt}%
\definecolor{currentstroke}{rgb}{0.000000,0.000000,0.000000}%
\pgfsetstrokecolor{currentstroke}%
\pgfsetdash{}{0pt}%
\pgfpathmoveto{\pgfqpoint{0.851214in}{0.881898in}}%
\pgfpathlineto{\pgfqpoint{0.844941in}{0.889020in}}%
\pgfpathlineto{\pgfqpoint{0.838638in}{0.896175in}}%
\pgfpathlineto{\pgfqpoint{0.877776in}{0.890713in}}%
\pgfpathlineto{\pgfqpoint{0.918027in}{0.886325in}}%
\pgfpathlineto{\pgfqpoint{0.922354in}{0.879395in}}%
\pgfpathlineto{\pgfqpoint{0.926659in}{0.872500in}}%
\pgfpathlineto{\pgfqpoint{0.888413in}{0.876686in}}%
\pgfpathlineto{\pgfqpoint{0.851214in}{0.881898in}}%
\pgfpathclose%
\pgfusepath{fill}%
\end{pgfscope}%
\begin{pgfscope}%
\pgfpathrectangle{\pgfqpoint{0.407592in}{0.222904in}}{\pgfqpoint{1.411163in}{1.411163in}}%
\pgfusepath{clip}%
\pgfsetbuttcap%
\pgfsetroundjoin%
\definecolor{currentfill}{rgb}{0.749020,0.772549,0.792157}%
\pgfsetfillcolor{currentfill}%
\pgfsetfillopacity{0.160000}%
\pgfsetlinewidth{0.000000pt}%
\definecolor{currentstroke}{rgb}{0.000000,0.000000,0.000000}%
\pgfsetstrokecolor{currentstroke}%
\pgfsetdash{}{0pt}%
\pgfpathmoveto{\pgfqpoint{0.677470in}{0.943837in}}%
\pgfpathlineto{\pgfqpoint{0.667588in}{0.951817in}}%
\pgfpathlineto{\pgfqpoint{0.657663in}{0.959832in}}%
\pgfpathlineto{\pgfqpoint{0.693604in}{0.949639in}}%
\pgfpathlineto{\pgfqpoint{0.731620in}{0.940422in}}%
\pgfpathlineto{\pgfqpoint{0.739897in}{0.932809in}}%
\pgfpathlineto{\pgfqpoint{0.748135in}{0.925230in}}%
\pgfpathlineto{\pgfqpoint{0.711820in}{0.934065in}}%
\pgfpathlineto{\pgfqpoint{0.677470in}{0.943837in}}%
\pgfpathclose%
\pgfusepath{fill}%
\end{pgfscope}%
\begin{pgfscope}%
\pgfpathrectangle{\pgfqpoint{0.407592in}{0.222904in}}{\pgfqpoint{1.411163in}{1.411163in}}%
\pgfusepath{clip}%
\pgfsetbuttcap%
\pgfsetroundjoin%
\definecolor{currentfill}{rgb}{0.749020,0.772549,0.792157}%
\pgfsetfillcolor{currentfill}%
\pgfsetfillopacity{0.160000}%
\pgfsetlinewidth{0.000000pt}%
\definecolor{currentstroke}{rgb}{0.000000,0.000000,0.000000}%
\pgfsetstrokecolor{currentstroke}%
\pgfsetdash{}{0pt}%
\pgfpathmoveto{\pgfqpoint{1.166568in}{0.869908in}}%
\pgfpathlineto{\pgfqpoint{1.168579in}{0.876741in}}%
\pgfpathlineto{\pgfqpoint{1.170600in}{0.883608in}}%
\pgfpathlineto{\pgfqpoint{1.212419in}{0.887118in}}%
\pgfpathlineto{\pgfqpoint{1.253556in}{0.891724in}}%
\pgfpathlineto{\pgfqpoint{1.249459in}{0.884670in}}%
\pgfpathlineto{\pgfqpoint{1.245381in}{0.877651in}}%
\pgfpathlineto{\pgfqpoint{1.206296in}{0.873256in}}%
\pgfpathlineto{\pgfqpoint{1.166568in}{0.869908in}}%
\pgfpathclose%
\pgfusepath{fill}%
\end{pgfscope}%
\begin{pgfscope}%
\pgfpathrectangle{\pgfqpoint{0.407592in}{0.222904in}}{\pgfqpoint{1.411163in}{1.411163in}}%
\pgfusepath{clip}%
\pgfsetbuttcap%
\pgfsetroundjoin%
\definecolor{currentfill}{rgb}{0.749020,0.772549,0.792157}%
\pgfsetfillcolor{currentfill}%
\pgfsetfillopacity{0.160000}%
\pgfsetlinewidth{0.000000pt}%
\definecolor{currentstroke}{rgb}{0.000000,0.000000,0.000000}%
\pgfsetstrokecolor{currentstroke}%
\pgfsetdash{}{0pt}%
\pgfpathmoveto{\pgfqpoint{0.764500in}{0.910176in}}%
\pgfpathlineto{\pgfqpoint{0.756336in}{0.917686in}}%
\pgfpathlineto{\pgfqpoint{0.748135in}{0.925230in}}%
\pgfpathlineto{\pgfqpoint{0.786240in}{0.917389in}}%
\pgfpathlineto{\pgfqpoint{0.825941in}{0.910590in}}%
\pgfpathlineto{\pgfqpoint{0.832305in}{0.903366in}}%
\pgfpathlineto{\pgfqpoint{0.838638in}{0.896175in}}%
\pgfpathlineto{\pgfqpoint{0.800814in}{0.902677in}}%
\pgfpathlineto{\pgfqpoint{0.764500in}{0.910176in}}%
\pgfpathclose%
\pgfusepath{fill}%
\end{pgfscope}%
\begin{pgfscope}%
\pgfpathrectangle{\pgfqpoint{0.407592in}{0.222904in}}{\pgfqpoint{1.411163in}{1.411163in}}%
\pgfusepath{clip}%
\pgfsetbuttcap%
\pgfsetroundjoin%
\definecolor{currentfill}{rgb}{0.749020,0.772549,0.792157}%
\pgfsetfillcolor{currentfill}%
\pgfsetfillopacity{0.160000}%
\pgfsetlinewidth{0.000000pt}%
\definecolor{currentstroke}{rgb}{0.000000,0.000000,0.000000}%
\pgfsetstrokecolor{currentstroke}%
\pgfsetdash{}{0pt}%
\pgfpathmoveto{\pgfqpoint{1.438914in}{0.951607in}}%
\pgfpathlineto{\pgfqpoint{1.446970in}{0.959487in}}%
\pgfpathlineto{\pgfqpoint{1.455062in}{0.967403in}}%
\pgfpathlineto{\pgfqpoint{1.494766in}{0.978210in}}%
\pgfpathlineto{\pgfqpoint{1.532296in}{0.989951in}}%
\pgfpathlineto{\pgfqpoint{1.522558in}{0.981593in}}%
\pgfpathlineto{\pgfqpoint{1.512861in}{0.973270in}}%
\pgfpathlineto{\pgfqpoint{1.476919in}{0.961989in}}%
\pgfpathlineto{\pgfqpoint{1.438914in}{0.951607in}}%
\pgfpathclose%
\pgfusepath{fill}%
\end{pgfscope}%
\begin{pgfscope}%
\pgfpathrectangle{\pgfqpoint{0.407592in}{0.222904in}}{\pgfqpoint{1.411163in}{1.411163in}}%
\pgfusepath{clip}%
\pgfsetbuttcap%
\pgfsetroundjoin%
\definecolor{currentfill}{rgb}{0.749020,0.772549,0.792157}%
\pgfsetfillcolor{currentfill}%
\pgfsetfillopacity{0.160000}%
\pgfsetlinewidth{0.000000pt}%
\definecolor{currentstroke}{rgb}{0.000000,0.000000,0.000000}%
\pgfsetstrokecolor{currentstroke}%
\pgfsetdash{}{0pt}%
\pgfpathmoveto{\pgfqpoint{1.551892in}{1.006770in}}%
\pgfpathlineto{\pgfqpoint{1.561752in}{1.015233in}}%
\pgfpathlineto{\pgfqpoint{1.571653in}{1.023730in}}%
\pgfpathlineto{\pgfqpoint{1.609784in}{1.037305in}}%
\pgfpathlineto{\pgfqpoint{1.645152in}{1.051715in}}%
\pgfpathlineto{\pgfqpoint{1.633838in}{1.042720in}}%
\pgfpathlineto{\pgfqpoint{1.622568in}{1.033759in}}%
\pgfpathlineto{\pgfqpoint{1.588546in}{1.019859in}}%
\pgfpathlineto{\pgfqpoint{1.551892in}{1.006770in}}%
\pgfpathclose%
\pgfusepath{fill}%
\end{pgfscope}%
\begin{pgfscope}%
\pgfpathrectangle{\pgfqpoint{0.407592in}{0.222904in}}{\pgfqpoint{1.411163in}{1.411163in}}%
\pgfusepath{clip}%
\pgfsetbuttcap%
\pgfsetroundjoin%
\definecolor{currentfill}{rgb}{0.749020,0.772549,0.792157}%
\pgfsetfillcolor{currentfill}%
\pgfsetfillopacity{0.160000}%
\pgfsetlinewidth{0.000000pt}%
\definecolor{currentstroke}{rgb}{0.000000,0.000000,0.000000}%
\pgfsetstrokecolor{currentstroke}%
\pgfsetdash{}{0pt}%
\pgfpathmoveto{\pgfqpoint{0.461940in}{1.063777in}}%
\pgfpathlineto{\pgfqpoint{0.449132in}{1.073016in}}%
\pgfpathlineto{\pgfqpoint{0.436278in}{1.082289in}}%
\pgfpathlineto{\pgfqpoint{0.466696in}{1.066750in}}%
\pgfpathlineto{\pgfqpoint{0.500331in}{1.051998in}}%
\pgfpathlineto{\pgfqpoint{0.512038in}{1.043237in}}%
\pgfpathlineto{\pgfqpoint{0.523697in}{1.034512in}}%
\pgfpathlineto{\pgfqpoint{0.491283in}{1.048761in}}%
\pgfpathlineto{\pgfqpoint{0.461940in}{1.063777in}}%
\pgfpathclose%
\pgfusepath{fill}%
\end{pgfscope}%
\begin{pgfscope}%
\pgfpathrectangle{\pgfqpoint{0.407592in}{0.222904in}}{\pgfqpoint{1.411163in}{1.411163in}}%
\pgfusepath{clip}%
\pgfsetbuttcap%
\pgfsetroundjoin%
\definecolor{currentfill}{rgb}{0.749020,0.772549,0.792157}%
\pgfsetfillcolor{currentfill}%
\pgfsetfillopacity{0.160000}%
\pgfsetlinewidth{0.000000pt}%
\definecolor{currentstroke}{rgb}{0.000000,0.000000,0.000000}%
\pgfsetstrokecolor{currentstroke}%
\pgfsetdash{}{0pt}%
\pgfpathmoveto{\pgfqpoint{0.926659in}{0.872500in}}%
\pgfpathlineto{\pgfqpoint{0.922354in}{0.879395in}}%
\pgfpathlineto{\pgfqpoint{0.918027in}{0.886325in}}%
\pgfpathlineto{\pgfqpoint{0.959177in}{0.883038in}}%
\pgfpathlineto{\pgfqpoint{1.001008in}{0.880873in}}%
\pgfpathlineto{\pgfqpoint{1.003254in}{0.874070in}}%
\pgfpathlineto{\pgfqpoint{1.005488in}{0.867300in}}%
\pgfpathlineto{\pgfqpoint{0.965753in}{0.869365in}}%
\pgfpathlineto{\pgfqpoint{0.926659in}{0.872500in}}%
\pgfpathclose%
\pgfusepath{fill}%
\end{pgfscope}%
\begin{pgfscope}%
\pgfpathrectangle{\pgfqpoint{0.407592in}{0.222904in}}{\pgfqpoint{1.411163in}{1.411163in}}%
\pgfusepath{clip}%
\pgfsetbuttcap%
\pgfsetroundjoin%
\definecolor{currentfill}{rgb}{0.749020,0.772549,0.792157}%
\pgfsetfillcolor{currentfill}%
\pgfsetfillopacity{0.160000}%
\pgfsetlinewidth{0.000000pt}%
\definecolor{currentstroke}{rgb}{0.000000,0.000000,0.000000}%
\pgfsetstrokecolor{currentstroke}%
\pgfsetdash{}{0pt}%
\pgfpathmoveto{\pgfqpoint{0.569868in}{0.999959in}}%
\pgfpathlineto{\pgfqpoint{0.558395in}{1.008545in}}%
\pgfpathlineto{\pgfqpoint{0.546875in}{1.017166in}}%
\pgfpathlineto{\pgfqpoint{0.580874in}{1.004249in}}%
\pgfpathlineto{\pgfqpoint{0.617531in}{0.992238in}}%
\pgfpathlineto{\pgfqpoint{0.627630in}{0.984084in}}%
\pgfpathlineto{\pgfqpoint{0.637684in}{0.975965in}}%
\pgfpathlineto{\pgfqpoint{0.602512in}{0.987524in}}%
\pgfpathlineto{\pgfqpoint{0.569868in}{0.999959in}}%
\pgfpathclose%
\pgfusepath{fill}%
\end{pgfscope}%
\begin{pgfscope}%
\pgfpathrectangle{\pgfqpoint{0.407592in}{0.222904in}}{\pgfqpoint{1.411163in}{1.411163in}}%
\pgfusepath{clip}%
\pgfsetbuttcap%
\pgfsetroundjoin%
\definecolor{currentfill}{rgb}{0.749020,0.772549,0.792157}%
\pgfsetfillcolor{currentfill}%
\pgfsetfillopacity{0.160000}%
\pgfsetlinewidth{0.000000pt}%
\definecolor{currentstroke}{rgb}{0.000000,0.000000,0.000000}%
\pgfsetstrokecolor{currentstroke}%
\pgfsetdash{}{0pt}%
\pgfpathmoveto{\pgfqpoint{1.253556in}{0.891724in}}%
\pgfpathlineto{\pgfqpoint{1.257673in}{0.898812in}}%
\pgfpathlineto{\pgfqpoint{1.261810in}{0.905934in}}%
\pgfpathlineto{\pgfqpoint{1.304060in}{0.911869in}}%
\pgfpathlineto{\pgfqpoint{1.345134in}{0.918880in}}%
\pgfpathlineto{\pgfqpoint{1.339011in}{0.911474in}}%
\pgfpathlineto{\pgfqpoint{1.332917in}{0.904103in}}%
\pgfpathlineto{\pgfqpoint{1.293792in}{0.897398in}}%
\pgfpathlineto{\pgfqpoint{1.253556in}{0.891724in}}%
\pgfpathclose%
\pgfusepath{fill}%
\end{pgfscope}%
\begin{pgfscope}%
\pgfpathrectangle{\pgfqpoint{0.407592in}{0.222904in}}{\pgfqpoint{1.411163in}{1.411163in}}%
\pgfusepath{clip}%
\pgfsetbuttcap%
\pgfsetroundjoin%
\definecolor{currentfill}{rgb}{0.749020,0.772549,0.792157}%
\pgfsetfillcolor{currentfill}%
\pgfsetfillopacity{0.160000}%
\pgfsetlinewidth{0.000000pt}%
\definecolor{currentstroke}{rgb}{0.000000,0.000000,0.000000}%
\pgfsetstrokecolor{currentstroke}%
\pgfsetdash{}{0pt}%
\pgfpathmoveto{\pgfqpoint{1.086031in}{0.866429in}}%
\pgfpathlineto{\pgfqpoint{1.085920in}{0.873177in}}%
\pgfpathlineto{\pgfqpoint{1.085809in}{0.879959in}}%
\pgfpathlineto{\pgfqpoint{1.128320in}{0.881216in}}%
\pgfpathlineto{\pgfqpoint{1.170600in}{0.883608in}}%
\pgfpathlineto{\pgfqpoint{1.168579in}{0.876741in}}%
\pgfpathlineto{\pgfqpoint{1.166568in}{0.869908in}}%
\pgfpathlineto{\pgfqpoint{1.126408in}{0.867628in}}%
\pgfpathlineto{\pgfqpoint{1.086031in}{0.866429in}}%
\pgfpathclose%
\pgfusepath{fill}%
\end{pgfscope}%
\begin{pgfscope}%
\pgfpathrectangle{\pgfqpoint{0.407592in}{0.222904in}}{\pgfqpoint{1.411163in}{1.411163in}}%
\pgfusepath{clip}%
\pgfsetbuttcap%
\pgfsetroundjoin%
\definecolor{currentfill}{rgb}{0.749020,0.772549,0.792157}%
\pgfsetfillcolor{currentfill}%
\pgfsetfillopacity{0.160000}%
\pgfsetlinewidth{0.000000pt}%
\definecolor{currentstroke}{rgb}{0.000000,0.000000,0.000000}%
\pgfsetstrokecolor{currentstroke}%
\pgfsetdash{}{0pt}%
\pgfpathmoveto{\pgfqpoint{1.345134in}{0.918880in}}%
\pgfpathlineto{\pgfqpoint{1.351285in}{0.926321in}}%
\pgfpathlineto{\pgfqpoint{1.357464in}{0.933796in}}%
\pgfpathlineto{\pgfqpoint{1.399030in}{0.942190in}}%
\pgfpathlineto{\pgfqpoint{1.438914in}{0.951607in}}%
\pgfpathlineto{\pgfqpoint{1.430894in}{0.943762in}}%
\pgfpathlineto{\pgfqpoint{1.422909in}{0.935952in}}%
\pgfpathlineto{\pgfqpoint{1.384817in}{0.926925in}}%
\pgfpathlineto{\pgfqpoint{1.345134in}{0.918880in}}%
\pgfpathclose%
\pgfusepath{fill}%
\end{pgfscope}%
\begin{pgfscope}%
\pgfpathrectangle{\pgfqpoint{0.407592in}{0.222904in}}{\pgfqpoint{1.411163in}{1.411163in}}%
\pgfusepath{clip}%
\pgfsetbuttcap%
\pgfsetroundjoin%
\definecolor{currentfill}{rgb}{0.749020,0.772549,0.792157}%
\pgfsetfillcolor{currentfill}%
\pgfsetfillopacity{0.160000}%
\pgfsetlinewidth{0.000000pt}%
\definecolor{currentstroke}{rgb}{0.000000,0.000000,0.000000}%
\pgfsetstrokecolor{currentstroke}%
\pgfsetdash{}{0pt}%
\pgfpathmoveto{\pgfqpoint{1.005488in}{0.867300in}}%
\pgfpathlineto{\pgfqpoint{1.003254in}{0.874070in}}%
\pgfpathlineto{\pgfqpoint{1.001008in}{0.880873in}}%
\pgfpathlineto{\pgfqpoint{1.043295in}{0.879845in}}%
\pgfpathlineto{\pgfqpoint{1.085809in}{0.879959in}}%
\pgfpathlineto{\pgfqpoint{1.085920in}{0.873177in}}%
\pgfpathlineto{\pgfqpoint{1.086031in}{0.866429in}}%
\pgfpathlineto{\pgfqpoint{1.045652in}{0.866319in}}%
\pgfpathlineto{\pgfqpoint{1.005488in}{0.867300in}}%
\pgfpathclose%
\pgfusepath{fill}%
\end{pgfscope}%
\begin{pgfscope}%
\pgfpathrectangle{\pgfqpoint{0.407592in}{0.222904in}}{\pgfqpoint{1.411163in}{1.411163in}}%
\pgfusepath{clip}%
\pgfsetbuttcap%
\pgfsetroundjoin%
\definecolor{currentfill}{rgb}{0.749020,0.772549,0.792157}%
\pgfsetfillcolor{currentfill}%
\pgfsetfillopacity{0.160000}%
\pgfsetlinewidth{0.000000pt}%
\definecolor{currentstroke}{rgb}{0.000000,0.000000,0.000000}%
\pgfsetstrokecolor{currentstroke}%
\pgfsetdash{}{0pt}%
\pgfpathmoveto{\pgfqpoint{0.838638in}{0.896175in}}%
\pgfpathlineto{\pgfqpoint{0.832305in}{0.903366in}}%
\pgfpathlineto{\pgfqpoint{0.825941in}{0.910590in}}%
\pgfpathlineto{\pgfqpoint{0.867036in}{0.904877in}}%
\pgfpathlineto{\pgfqpoint{0.909310in}{0.900286in}}%
\pgfpathlineto{\pgfqpoint{0.913679in}{0.893288in}}%
\pgfpathlineto{\pgfqpoint{0.918027in}{0.886325in}}%
\pgfpathlineto{\pgfqpoint{0.877776in}{0.890713in}}%
\pgfpathlineto{\pgfqpoint{0.838638in}{0.896175in}}%
\pgfpathclose%
\pgfusepath{fill}%
\end{pgfscope}%
\begin{pgfscope}%
\pgfpathrectangle{\pgfqpoint{0.407592in}{0.222904in}}{\pgfqpoint{1.411163in}{1.411163in}}%
\pgfusepath{clip}%
\pgfsetbuttcap%
\pgfsetroundjoin%
\definecolor{currentfill}{rgb}{0.749020,0.772549,0.792157}%
\pgfsetfillcolor{currentfill}%
\pgfsetfillopacity{0.160000}%
\pgfsetlinewidth{0.000000pt}%
\definecolor{currentstroke}{rgb}{0.000000,0.000000,0.000000}%
\pgfsetstrokecolor{currentstroke}%
\pgfsetdash{}{0pt}%
\pgfpathmoveto{\pgfqpoint{0.436278in}{1.082289in}}%
\pgfpathlineto{\pgfqpoint{0.423376in}{1.091597in}}%
\pgfpathlineto{\pgfqpoint{0.454331in}{1.075797in}}%
\pgfpathlineto{\pgfqpoint{0.488577in}{1.060795in}}%
\pgfpathlineto{\pgfqpoint{0.500331in}{1.051998in}}%
\pgfpathlineto{\pgfqpoint{0.466696in}{1.066750in}}%
\pgfusepath{fill}%
\end{pgfscope}%
\begin{pgfscope}%
\pgfpathrectangle{\pgfqpoint{0.407592in}{0.222904in}}{\pgfqpoint{1.411163in}{1.411163in}}%
\pgfusepath{clip}%
\pgfsetbuttcap%
\pgfsetroundjoin%
\definecolor{currentfill}{rgb}{0.749020,0.772549,0.792157}%
\pgfsetfillcolor{currentfill}%
\pgfsetfillopacity{0.160000}%
\pgfsetlinewidth{0.000000pt}%
\definecolor{currentstroke}{rgb}{0.000000,0.000000,0.000000}%
\pgfsetstrokecolor{currentstroke}%
\pgfsetdash{}{0pt}%
\pgfpathmoveto{\pgfqpoint{0.657663in}{0.959832in}}%
\pgfpathlineto{\pgfqpoint{0.647695in}{0.967881in}}%
\pgfpathlineto{\pgfqpoint{0.637684in}{0.975965in}}%
\pgfpathlineto{\pgfqpoint{0.675224in}{0.965353in}}%
\pgfpathlineto{\pgfqpoint{0.714953in}{0.955755in}}%
\pgfpathlineto{\pgfqpoint{0.723306in}{0.948071in}}%
\pgfpathlineto{\pgfqpoint{0.731620in}{0.940422in}}%
\pgfpathlineto{\pgfqpoint{0.693604in}{0.949639in}}%
\pgfpathlineto{\pgfqpoint{0.657663in}{0.959832in}}%
\pgfpathclose%
\pgfusepath{fill}%
\end{pgfscope}%
\begin{pgfscope}%
\pgfpathrectangle{\pgfqpoint{0.407592in}{0.222904in}}{\pgfqpoint{1.411163in}{1.411163in}}%
\pgfusepath{clip}%
\pgfsetbuttcap%
\pgfsetroundjoin%
\definecolor{currentfill}{rgb}{0.749020,0.772549,0.792157}%
\pgfsetfillcolor{currentfill}%
\pgfsetfillopacity{0.160000}%
\pgfsetlinewidth{0.000000pt}%
\definecolor{currentstroke}{rgb}{0.000000,0.000000,0.000000}%
\pgfsetstrokecolor{currentstroke}%
\pgfsetdash{}{0pt}%
\pgfpathmoveto{\pgfqpoint{1.455062in}{0.967403in}}%
\pgfpathlineto{\pgfqpoint{1.463190in}{0.975353in}}%
\pgfpathlineto{\pgfqpoint{1.471354in}{0.983339in}}%
\pgfpathlineto{\pgfqpoint{1.512768in}{0.994572in}}%
\pgfpathlineto{\pgfqpoint{1.551892in}{1.006770in}}%
\pgfpathlineto{\pgfqpoint{1.542074in}{0.998343in}}%
\pgfpathlineto{\pgfqpoint{1.532296in}{0.989951in}}%
\pgfpathlineto{\pgfqpoint{1.494766in}{0.978210in}}%
\pgfpathlineto{\pgfqpoint{1.455062in}{0.967403in}}%
\pgfpathclose%
\pgfusepath{fill}%
\end{pgfscope}%
\begin{pgfscope}%
\pgfpathrectangle{\pgfqpoint{0.407592in}{0.222904in}}{\pgfqpoint{1.411163in}{1.411163in}}%
\pgfusepath{clip}%
\pgfsetbuttcap%
\pgfsetroundjoin%
\definecolor{currentfill}{rgb}{0.749020,0.772549,0.792157}%
\pgfsetfillcolor{currentfill}%
\pgfsetfillopacity{0.160000}%
\pgfsetlinewidth{0.000000pt}%
\definecolor{currentstroke}{rgb}{0.000000,0.000000,0.000000}%
\pgfsetstrokecolor{currentstroke}%
\pgfsetdash{}{0pt}%
\pgfpathmoveto{\pgfqpoint{1.571653in}{1.023730in}}%
\pgfpathlineto{\pgfqpoint{1.581596in}{1.032264in}}%
\pgfpathlineto{\pgfqpoint{1.591580in}{1.040833in}}%
\pgfpathlineto{\pgfqpoint{1.631193in}{1.054890in}}%
\pgfpathlineto{\pgfqpoint{1.667909in}{1.069807in}}%
\pgfpathlineto{\pgfqpoint{1.656509in}{1.060744in}}%
\pgfpathlineto{\pgfqpoint{1.645152in}{1.051715in}}%
\pgfpathlineto{\pgfqpoint{1.609784in}{1.037305in}}%
\pgfpathlineto{\pgfqpoint{1.571653in}{1.023730in}}%
\pgfpathclose%
\pgfusepath{fill}%
\end{pgfscope}%
\begin{pgfscope}%
\pgfpathrectangle{\pgfqpoint{0.407592in}{0.222904in}}{\pgfqpoint{1.411163in}{1.411163in}}%
\pgfusepath{clip}%
\pgfsetbuttcap%
\pgfsetroundjoin%
\definecolor{currentfill}{rgb}{0.749020,0.772549,0.792157}%
\pgfsetfillcolor{currentfill}%
\pgfsetfillopacity{0.160000}%
\pgfsetlinewidth{0.000000pt}%
\definecolor{currentstroke}{rgb}{0.000000,0.000000,0.000000}%
\pgfsetstrokecolor{currentstroke}%
\pgfsetdash{}{0pt}%
\pgfpathmoveto{\pgfqpoint{0.748135in}{0.925230in}}%
\pgfpathlineto{\pgfqpoint{0.739897in}{0.932809in}}%
\pgfpathlineto{\pgfqpoint{0.731620in}{0.940422in}}%
\pgfpathlineto{\pgfqpoint{0.771527in}{0.932240in}}%
\pgfpathlineto{\pgfqpoint{0.813123in}{0.925144in}}%
\pgfpathlineto{\pgfqpoint{0.819547in}{0.917850in}}%
\pgfpathlineto{\pgfqpoint{0.825941in}{0.910590in}}%
\pgfpathlineto{\pgfqpoint{0.786240in}{0.917389in}}%
\pgfpathlineto{\pgfqpoint{0.748135in}{0.925230in}}%
\pgfpathclose%
\pgfusepath{fill}%
\end{pgfscope}%
\begin{pgfscope}%
\pgfpathrectangle{\pgfqpoint{0.407592in}{0.222904in}}{\pgfqpoint{1.411163in}{1.411163in}}%
\pgfusepath{clip}%
\pgfsetbuttcap%
\pgfsetroundjoin%
\definecolor{currentfill}{rgb}{0.749020,0.772549,0.792157}%
\pgfsetfillcolor{currentfill}%
\pgfsetfillopacity{0.160000}%
\pgfsetlinewidth{0.000000pt}%
\definecolor{currentstroke}{rgb}{0.000000,0.000000,0.000000}%
\pgfsetstrokecolor{currentstroke}%
\pgfsetdash{}{0pt}%
\pgfpathmoveto{\pgfqpoint{1.170600in}{0.883608in}}%
\pgfpathlineto{\pgfqpoint{1.172631in}{0.890508in}}%
\pgfpathlineto{\pgfqpoint{1.174672in}{0.897443in}}%
\pgfpathlineto{\pgfqpoint{1.218603in}{0.901116in}}%
\pgfpathlineto{\pgfqpoint{1.261810in}{0.905934in}}%
\pgfpathlineto{\pgfqpoint{1.257673in}{0.898812in}}%
\pgfpathlineto{\pgfqpoint{1.253556in}{0.891724in}}%
\pgfpathlineto{\pgfqpoint{1.212419in}{0.887118in}}%
\pgfpathlineto{\pgfqpoint{1.170600in}{0.883608in}}%
\pgfpathclose%
\pgfusepath{fill}%
\end{pgfscope}%
\begin{pgfscope}%
\pgfpathrectangle{\pgfqpoint{0.407592in}{0.222904in}}{\pgfqpoint{1.411163in}{1.411163in}}%
\pgfusepath{clip}%
\pgfsetbuttcap%
\pgfsetroundjoin%
\definecolor{currentfill}{rgb}{0.749020,0.772549,0.792157}%
\pgfsetfillcolor{currentfill}%
\pgfsetfillopacity{0.160000}%
\pgfsetlinewidth{0.000000pt}%
\definecolor{currentstroke}{rgb}{0.000000,0.000000,0.000000}%
\pgfsetstrokecolor{currentstroke}%
\pgfsetdash{}{0pt}%
\pgfpathmoveto{\pgfqpoint{0.546875in}{1.017166in}}%
\pgfpathlineto{\pgfqpoint{0.535310in}{1.025821in}}%
\pgfpathlineto{\pgfqpoint{0.523697in}{1.034512in}}%
\pgfpathlineto{\pgfqpoint{0.559054in}{1.021115in}}%
\pgfpathlineto{\pgfqpoint{0.597202in}{1.008654in}}%
\pgfpathlineto{\pgfqpoint{0.607389in}{1.000428in}}%
\pgfpathlineto{\pgfqpoint{0.617531in}{0.992238in}}%
\pgfpathlineto{\pgfqpoint{0.580874in}{1.004249in}}%
\pgfpathlineto{\pgfqpoint{0.546875in}{1.017166in}}%
\pgfpathclose%
\pgfusepath{fill}%
\end{pgfscope}%
\begin{pgfscope}%
\pgfpathrectangle{\pgfqpoint{0.407592in}{0.222904in}}{\pgfqpoint{1.411163in}{1.411163in}}%
\pgfusepath{clip}%
\pgfsetbuttcap%
\pgfsetroundjoin%
\definecolor{currentfill}{rgb}{0.749020,0.772549,0.792157}%
\pgfsetfillcolor{currentfill}%
\pgfsetfillopacity{0.160000}%
\pgfsetlinewidth{0.000000pt}%
\definecolor{currentstroke}{rgb}{0.000000,0.000000,0.000000}%
\pgfsetstrokecolor{currentstroke}%
\pgfsetdash{}{0pt}%
\pgfpathmoveto{\pgfqpoint{0.918027in}{0.886325in}}%
\pgfpathlineto{\pgfqpoint{0.913679in}{0.893288in}}%
\pgfpathlineto{\pgfqpoint{0.909310in}{0.900286in}}%
\pgfpathlineto{\pgfqpoint{0.952536in}{0.896847in}}%
\pgfpathlineto{\pgfqpoint{0.996484in}{0.894582in}}%
\pgfpathlineto{\pgfqpoint{0.998752in}{0.887710in}}%
\pgfpathlineto{\pgfqpoint{1.001008in}{0.880873in}}%
\pgfpathlineto{\pgfqpoint{0.959177in}{0.883038in}}%
\pgfpathlineto{\pgfqpoint{0.918027in}{0.886325in}}%
\pgfpathclose%
\pgfusepath{fill}%
\end{pgfscope}%
\begin{pgfscope}%
\pgfpathrectangle{\pgfqpoint{0.407592in}{0.222904in}}{\pgfqpoint{1.411163in}{1.411163in}}%
\pgfusepath{clip}%
\pgfsetbuttcap%
\pgfsetroundjoin%
\definecolor{currentfill}{rgb}{0.749020,0.772549,0.792157}%
\pgfsetfillcolor{currentfill}%
\pgfsetfillopacity{0.160000}%
\pgfsetlinewidth{0.000000pt}%
\definecolor{currentstroke}{rgb}{0.000000,0.000000,0.000000}%
\pgfsetstrokecolor{currentstroke}%
\pgfsetdash{}{0pt}%
\pgfpathmoveto{\pgfqpoint{1.357464in}{0.933796in}}%
\pgfpathlineto{\pgfqpoint{1.363673in}{0.941307in}}%
\pgfpathlineto{\pgfqpoint{1.369911in}{0.948853in}}%
\pgfpathlineto{\pgfqpoint{1.413375in}{0.957596in}}%
\pgfpathlineto{\pgfqpoint{1.455062in}{0.967403in}}%
\pgfpathlineto{\pgfqpoint{1.446970in}{0.959487in}}%
\pgfpathlineto{\pgfqpoint{1.438914in}{0.951607in}}%
\pgfpathlineto{\pgfqpoint{1.399030in}{0.942190in}}%
\pgfpathlineto{\pgfqpoint{1.357464in}{0.933796in}}%
\pgfpathclose%
\pgfusepath{fill}%
\end{pgfscope}%
\begin{pgfscope}%
\pgfpathrectangle{\pgfqpoint{0.407592in}{0.222904in}}{\pgfqpoint{1.411163in}{1.411163in}}%
\pgfusepath{clip}%
\pgfsetbuttcap%
\pgfsetroundjoin%
\definecolor{currentfill}{rgb}{0.749020,0.772549,0.792157}%
\pgfsetfillcolor{currentfill}%
\pgfsetfillopacity{0.160000}%
\pgfsetlinewidth{0.000000pt}%
\definecolor{currentstroke}{rgb}{0.000000,0.000000,0.000000}%
\pgfsetstrokecolor{currentstroke}%
\pgfsetdash{}{0pt}%
\pgfpathmoveto{\pgfqpoint{1.261810in}{0.905934in}}%
\pgfpathlineto{\pgfqpoint{1.265967in}{0.913092in}}%
\pgfpathlineto{\pgfqpoint{1.270144in}{0.920284in}}%
\pgfpathlineto{\pgfqpoint{1.314428in}{0.926478in}}%
\pgfpathlineto{\pgfqpoint{1.357464in}{0.933796in}}%
\pgfpathlineto{\pgfqpoint{1.351285in}{0.926321in}}%
\pgfpathlineto{\pgfqpoint{1.345134in}{0.918880in}}%
\pgfpathlineto{\pgfqpoint{1.304060in}{0.911869in}}%
\pgfpathlineto{\pgfqpoint{1.261810in}{0.905934in}}%
\pgfpathclose%
\pgfusepath{fill}%
\end{pgfscope}%
\begin{pgfscope}%
\pgfpathrectangle{\pgfqpoint{0.407592in}{0.222904in}}{\pgfqpoint{1.411163in}{1.411163in}}%
\pgfusepath{clip}%
\pgfsetbuttcap%
\pgfsetroundjoin%
\definecolor{currentfill}{rgb}{0.749020,0.772549,0.792157}%
\pgfsetfillcolor{currentfill}%
\pgfsetfillopacity{0.160000}%
\pgfsetlinewidth{0.000000pt}%
\definecolor{currentstroke}{rgb}{0.000000,0.000000,0.000000}%
\pgfsetstrokecolor{currentstroke}%
\pgfsetdash{}{0pt}%
\pgfpathmoveto{\pgfqpoint{1.085809in}{0.879959in}}%
\pgfpathlineto{\pgfqpoint{1.085697in}{0.886776in}}%
\pgfpathlineto{\pgfqpoint{1.085585in}{0.893626in}}%
\pgfpathlineto{\pgfqpoint{1.130252in}{0.894942in}}%
\pgfpathlineto{\pgfqpoint{1.174672in}{0.897443in}}%
\pgfpathlineto{\pgfqpoint{1.172631in}{0.890508in}}%
\pgfpathlineto{\pgfqpoint{1.170600in}{0.883608in}}%
\pgfpathlineto{\pgfqpoint{1.128320in}{0.881216in}}%
\pgfpathlineto{\pgfqpoint{1.085809in}{0.879959in}}%
\pgfpathclose%
\pgfusepath{fill}%
\end{pgfscope}%
\begin{pgfscope}%
\pgfpathrectangle{\pgfqpoint{0.407592in}{0.222904in}}{\pgfqpoint{1.411163in}{1.411163in}}%
\pgfusepath{clip}%
\pgfsetbuttcap%
\pgfsetroundjoin%
\definecolor{currentfill}{rgb}{0.749020,0.772549,0.792157}%
\pgfsetfillcolor{currentfill}%
\pgfsetfillopacity{0.160000}%
\pgfsetlinewidth{0.000000pt}%
\definecolor{currentstroke}{rgb}{0.000000,0.000000,0.000000}%
\pgfsetstrokecolor{currentstroke}%
\pgfsetdash{}{0pt}%
\pgfpathmoveto{\pgfqpoint{1.001008in}{0.880873in}}%
\pgfpathlineto{\pgfqpoint{0.998752in}{0.887710in}}%
\pgfpathlineto{\pgfqpoint{0.996484in}{0.894582in}}%
\pgfpathlineto{\pgfqpoint{1.040914in}{0.893506in}}%
\pgfpathlineto{\pgfqpoint{1.085585in}{0.893626in}}%
\pgfpathlineto{\pgfqpoint{1.085697in}{0.886776in}}%
\pgfpathlineto{\pgfqpoint{1.085809in}{0.879959in}}%
\pgfpathlineto{\pgfqpoint{1.043295in}{0.879845in}}%
\pgfpathlineto{\pgfqpoint{1.001008in}{0.880873in}}%
\pgfpathclose%
\pgfusepath{fill}%
\end{pgfscope}%
\begin{pgfscope}%
\pgfpathrectangle{\pgfqpoint{0.407592in}{0.222904in}}{\pgfqpoint{1.411163in}{1.411163in}}%
\pgfusepath{clip}%
\pgfsetbuttcap%
\pgfsetroundjoin%
\definecolor{currentfill}{rgb}{0.749020,0.772549,0.792157}%
\pgfsetfillcolor{currentfill}%
\pgfsetfillopacity{0.160000}%
\pgfsetlinewidth{0.000000pt}%
\definecolor{currentstroke}{rgb}{0.000000,0.000000,0.000000}%
\pgfsetstrokecolor{currentstroke}%
\pgfsetdash{}{0pt}%
\pgfpathmoveto{\pgfqpoint{1.591580in}{1.040833in}}%
\pgfpathlineto{\pgfqpoint{1.601606in}{1.049438in}}%
\pgfpathlineto{\pgfqpoint{1.641962in}{1.063736in}}%
\pgfpathlineto{\pgfqpoint{1.679353in}{1.078906in}}%
\pgfpathlineto{\pgfqpoint{1.667909in}{1.069807in}}%
\pgfpathlineto{\pgfqpoint{1.631193in}{1.054890in}}%
\pgfusepath{fill}%
\end{pgfscope}%
\begin{pgfscope}%
\pgfpathrectangle{\pgfqpoint{0.407592in}{0.222904in}}{\pgfqpoint{1.411163in}{1.411163in}}%
\pgfusepath{clip}%
\pgfsetbuttcap%
\pgfsetroundjoin%
\definecolor{currentfill}{rgb}{0.749020,0.772549,0.792157}%
\pgfsetfillcolor{currentfill}%
\pgfsetfillopacity{0.160000}%
\pgfsetlinewidth{0.000000pt}%
\definecolor{currentstroke}{rgb}{0.000000,0.000000,0.000000}%
\pgfsetstrokecolor{currentstroke}%
\pgfsetdash{}{0pt}%
\pgfpathmoveto{\pgfqpoint{0.637684in}{0.975965in}}%
\pgfpathlineto{\pgfqpoint{0.627630in}{0.984084in}}%
\pgfpathlineto{\pgfqpoint{0.617531in}{0.992238in}}%
\pgfpathlineto{\pgfqpoint{0.656679in}{0.981209in}}%
\pgfpathlineto{\pgfqpoint{0.698131in}{0.971230in}}%
\pgfpathlineto{\pgfqpoint{0.706562in}{0.963474in}}%
\pgfpathlineto{\pgfqpoint{0.714953in}{0.955755in}}%
\pgfpathlineto{\pgfqpoint{0.675224in}{0.965353in}}%
\pgfpathlineto{\pgfqpoint{0.637684in}{0.975965in}}%
\pgfpathclose%
\pgfusepath{fill}%
\end{pgfscope}%
\begin{pgfscope}%
\pgfpathrectangle{\pgfqpoint{0.407592in}{0.222904in}}{\pgfqpoint{1.411163in}{1.411163in}}%
\pgfusepath{clip}%
\pgfsetbuttcap%
\pgfsetroundjoin%
\definecolor{currentfill}{rgb}{0.749020,0.772549,0.792157}%
\pgfsetfillcolor{currentfill}%
\pgfsetfillopacity{0.160000}%
\pgfsetlinewidth{0.000000pt}%
\definecolor{currentstroke}{rgb}{0.000000,0.000000,0.000000}%
\pgfsetstrokecolor{currentstroke}%
\pgfsetdash{}{0pt}%
\pgfpathmoveto{\pgfqpoint{0.825941in}{0.910590in}}%
\pgfpathlineto{\pgfqpoint{0.819547in}{0.917850in}}%
\pgfpathlineto{\pgfqpoint{0.813123in}{0.925144in}}%
\pgfpathlineto{\pgfqpoint{0.856191in}{0.919180in}}%
\pgfpathlineto{\pgfqpoint{0.900505in}{0.914386in}}%
\pgfpathlineto{\pgfqpoint{0.904918in}{0.907319in}}%
\pgfpathlineto{\pgfqpoint{0.909310in}{0.900286in}}%
\pgfpathlineto{\pgfqpoint{0.867036in}{0.904877in}}%
\pgfpathlineto{\pgfqpoint{0.825941in}{0.910590in}}%
\pgfpathclose%
\pgfusepath{fill}%
\end{pgfscope}%
\begin{pgfscope}%
\pgfpathrectangle{\pgfqpoint{0.407592in}{0.222904in}}{\pgfqpoint{1.411163in}{1.411163in}}%
\pgfusepath{clip}%
\pgfsetbuttcap%
\pgfsetroundjoin%
\definecolor{currentfill}{rgb}{0.749020,0.772549,0.792157}%
\pgfsetfillcolor{currentfill}%
\pgfsetfillopacity{0.160000}%
\pgfsetlinewidth{0.000000pt}%
\definecolor{currentstroke}{rgb}{0.000000,0.000000,0.000000}%
\pgfsetstrokecolor{currentstroke}%
\pgfsetdash{}{0pt}%
\pgfpathmoveto{\pgfqpoint{1.471354in}{0.983339in}}%
\pgfpathlineto{\pgfqpoint{1.479555in}{0.991361in}}%
\pgfpathlineto{\pgfqpoint{1.487793in}{0.999419in}}%
\pgfpathlineto{\pgfqpoint{1.530926in}{1.011075in}}%
\pgfpathlineto{\pgfqpoint{1.571653in}{1.023730in}}%
\pgfpathlineto{\pgfqpoint{1.561752in}{1.015233in}}%
\pgfpathlineto{\pgfqpoint{1.551892in}{1.006770in}}%
\pgfpathlineto{\pgfqpoint{1.512768in}{0.994572in}}%
\pgfpathlineto{\pgfqpoint{1.471354in}{0.983339in}}%
\pgfpathclose%
\pgfusepath{fill}%
\end{pgfscope}%
\begin{pgfscope}%
\pgfpathrectangle{\pgfqpoint{0.407592in}{0.222904in}}{\pgfqpoint{1.411163in}{1.411163in}}%
\pgfusepath{clip}%
\pgfsetbuttcap%
\pgfsetroundjoin%
\definecolor{currentfill}{rgb}{0.749020,0.772549,0.792157}%
\pgfsetfillcolor{currentfill}%
\pgfsetfillopacity{0.160000}%
\pgfsetlinewidth{0.000000pt}%
\definecolor{currentstroke}{rgb}{0.000000,0.000000,0.000000}%
\pgfsetstrokecolor{currentstroke}%
\pgfsetdash{}{0pt}%
\pgfpathmoveto{\pgfqpoint{0.731620in}{0.940422in}}%
\pgfpathlineto{\pgfqpoint{0.723306in}{0.948071in}}%
\pgfpathlineto{\pgfqpoint{0.714953in}{0.955755in}}%
\pgfpathlineto{\pgfqpoint{0.756676in}{0.947232in}}%
\pgfpathlineto{\pgfqpoint{0.800180in}{0.939838in}}%
\pgfpathlineto{\pgfqpoint{0.806667in}{0.932473in}}%
\pgfpathlineto{\pgfqpoint{0.813123in}{0.925144in}}%
\pgfpathlineto{\pgfqpoint{0.771527in}{0.932240in}}%
\pgfpathlineto{\pgfqpoint{0.731620in}{0.940422in}}%
\pgfpathclose%
\pgfusepath{fill}%
\end{pgfscope}%
\begin{pgfscope}%
\pgfpathrectangle{\pgfqpoint{0.407592in}{0.222904in}}{\pgfqpoint{1.411163in}{1.411163in}}%
\pgfusepath{clip}%
\pgfsetbuttcap%
\pgfsetroundjoin%
\definecolor{currentfill}{rgb}{0.749020,0.772549,0.792157}%
\pgfsetfillcolor{currentfill}%
\pgfsetfillopacity{0.160000}%
\pgfsetlinewidth{0.000000pt}%
\definecolor{currentstroke}{rgb}{0.000000,0.000000,0.000000}%
\pgfsetstrokecolor{currentstroke}%
\pgfsetdash{}{0pt}%
\pgfpathmoveto{\pgfqpoint{0.523697in}{1.034512in}}%
\pgfpathlineto{\pgfqpoint{0.512038in}{1.043237in}}%
\pgfpathlineto{\pgfqpoint{0.500331in}{1.051998in}}%
\pgfpathlineto{\pgfqpoint{0.537050in}{1.038124in}}%
\pgfpathlineto{\pgfqpoint{0.576694in}{1.025215in}}%
\pgfpathlineto{\pgfqpoint{0.586970in}{1.016917in}}%
\pgfpathlineto{\pgfqpoint{0.597202in}{1.008654in}}%
\pgfpathlineto{\pgfqpoint{0.559054in}{1.021115in}}%
\pgfpathlineto{\pgfqpoint{0.523697in}{1.034512in}}%
\pgfpathclose%
\pgfusepath{fill}%
\end{pgfscope}%
\begin{pgfscope}%
\pgfpathrectangle{\pgfqpoint{0.407592in}{0.222904in}}{\pgfqpoint{1.411163in}{1.411163in}}%
\pgfusepath{clip}%
\pgfsetbuttcap%
\pgfsetroundjoin%
\definecolor{currentfill}{rgb}{0.749020,0.772549,0.792157}%
\pgfsetfillcolor{currentfill}%
\pgfsetfillopacity{0.160000}%
\pgfsetlinewidth{0.000000pt}%
\definecolor{currentstroke}{rgb}{0.000000,0.000000,0.000000}%
\pgfsetstrokecolor{currentstroke}%
\pgfsetdash{}{0pt}%
\pgfpathmoveto{\pgfqpoint{1.174672in}{0.897443in}}%
\pgfpathlineto{\pgfqpoint{1.176723in}{0.904413in}}%
\pgfpathlineto{\pgfqpoint{1.178785in}{0.911418in}}%
\pgfpathlineto{\pgfqpoint{1.224849in}{0.915253in}}%
\pgfpathlineto{\pgfqpoint{1.270144in}{0.920284in}}%
\pgfpathlineto{\pgfqpoint{1.265967in}{0.913092in}}%
\pgfpathlineto{\pgfqpoint{1.261810in}{0.905934in}}%
\pgfpathlineto{\pgfqpoint{1.218603in}{0.901116in}}%
\pgfpathlineto{\pgfqpoint{1.174672in}{0.897443in}}%
\pgfpathclose%
\pgfusepath{fill}%
\end{pgfscope}%
\begin{pgfscope}%
\pgfpathrectangle{\pgfqpoint{0.407592in}{0.222904in}}{\pgfqpoint{1.411163in}{1.411163in}}%
\pgfusepath{clip}%
\pgfsetbuttcap%
\pgfsetroundjoin%
\definecolor{currentfill}{rgb}{0.749020,0.772549,0.792157}%
\pgfsetfillcolor{currentfill}%
\pgfsetfillopacity{0.160000}%
\pgfsetlinewidth{0.000000pt}%
\definecolor{currentstroke}{rgb}{0.000000,0.000000,0.000000}%
\pgfsetstrokecolor{currentstroke}%
\pgfsetdash{}{0pt}%
\pgfpathmoveto{\pgfqpoint{1.369911in}{0.948853in}}%
\pgfpathlineto{\pgfqpoint{1.376179in}{0.956434in}}%
\pgfpathlineto{\pgfqpoint{1.382476in}{0.964052in}}%
\pgfpathlineto{\pgfqpoint{1.427851in}{0.973144in}}%
\pgfpathlineto{\pgfqpoint{1.471354in}{0.983339in}}%
\pgfpathlineto{\pgfqpoint{1.463190in}{0.975353in}}%
\pgfpathlineto{\pgfqpoint{1.455062in}{0.967403in}}%
\pgfpathlineto{\pgfqpoint{1.413375in}{0.957596in}}%
\pgfpathlineto{\pgfqpoint{1.369911in}{0.948853in}}%
\pgfpathclose%
\pgfusepath{fill}%
\end{pgfscope}%
\begin{pgfscope}%
\pgfpathrectangle{\pgfqpoint{0.407592in}{0.222904in}}{\pgfqpoint{1.411163in}{1.411163in}}%
\pgfusepath{clip}%
\pgfsetbuttcap%
\pgfsetroundjoin%
\definecolor{currentfill}{rgb}{0.749020,0.772549,0.792157}%
\pgfsetfillcolor{currentfill}%
\pgfsetfillopacity{0.160000}%
\pgfsetlinewidth{0.000000pt}%
\definecolor{currentstroke}{rgb}{0.000000,0.000000,0.000000}%
\pgfsetstrokecolor{currentstroke}%
\pgfsetdash{}{0pt}%
\pgfpathmoveto{\pgfqpoint{0.909310in}{0.900286in}}%
\pgfpathlineto{\pgfqpoint{0.904918in}{0.907319in}}%
\pgfpathlineto{\pgfqpoint{0.900505in}{0.914386in}}%
\pgfpathlineto{\pgfqpoint{0.945828in}{0.910795in}}%
\pgfpathlineto{\pgfqpoint{0.991914in}{0.908430in}}%
\pgfpathlineto{\pgfqpoint{0.994204in}{0.901489in}}%
\pgfpathlineto{\pgfqpoint{0.996484in}{0.894582in}}%
\pgfpathlineto{\pgfqpoint{0.952536in}{0.896847in}}%
\pgfpathlineto{\pgfqpoint{0.909310in}{0.900286in}}%
\pgfpathclose%
\pgfusepath{fill}%
\end{pgfscope}%
\begin{pgfscope}%
\pgfpathrectangle{\pgfqpoint{0.407592in}{0.222904in}}{\pgfqpoint{1.411163in}{1.411163in}}%
\pgfusepath{clip}%
\pgfsetbuttcap%
\pgfsetroundjoin%
\definecolor{currentfill}{rgb}{0.749020,0.772549,0.792157}%
\pgfsetfillcolor{currentfill}%
\pgfsetfillopacity{0.160000}%
\pgfsetlinewidth{0.000000pt}%
\definecolor{currentstroke}{rgb}{0.000000,0.000000,0.000000}%
\pgfsetstrokecolor{currentstroke}%
\pgfsetdash{}{0pt}%
\pgfpathmoveto{\pgfqpoint{1.270144in}{0.920284in}}%
\pgfpathlineto{\pgfqpoint{1.274342in}{0.927511in}}%
\pgfpathlineto{\pgfqpoint{1.278561in}{0.934774in}}%
\pgfpathlineto{\pgfqpoint{1.324895in}{0.941229in}}%
\pgfpathlineto{\pgfqpoint{1.369911in}{0.948853in}}%
\pgfpathlineto{\pgfqpoint{1.363673in}{0.941307in}}%
\pgfpathlineto{\pgfqpoint{1.357464in}{0.933796in}}%
\pgfpathlineto{\pgfqpoint{1.314428in}{0.926478in}}%
\pgfpathlineto{\pgfqpoint{1.270144in}{0.920284in}}%
\pgfpathclose%
\pgfusepath{fill}%
\end{pgfscope}%
\begin{pgfscope}%
\pgfpathrectangle{\pgfqpoint{0.407592in}{0.222904in}}{\pgfqpoint{1.411163in}{1.411163in}}%
\pgfusepath{clip}%
\pgfsetbuttcap%
\pgfsetroundjoin%
\definecolor{currentfill}{rgb}{0.749020,0.772549,0.792157}%
\pgfsetfillcolor{currentfill}%
\pgfsetfillopacity{0.160000}%
\pgfsetlinewidth{0.000000pt}%
\definecolor{currentstroke}{rgb}{0.000000,0.000000,0.000000}%
\pgfsetstrokecolor{currentstroke}%
\pgfsetdash{}{0pt}%
\pgfpathmoveto{\pgfqpoint{1.085585in}{0.893626in}}%
\pgfpathlineto{\pgfqpoint{1.085472in}{0.900511in}}%
\pgfpathlineto{\pgfqpoint{1.085358in}{0.907432in}}%
\pgfpathlineto{\pgfqpoint{1.132203in}{0.908805in}}%
\pgfpathlineto{\pgfqpoint{1.178785in}{0.911418in}}%
\pgfpathlineto{\pgfqpoint{1.176723in}{0.904413in}}%
\pgfpathlineto{\pgfqpoint{1.174672in}{0.897443in}}%
\pgfpathlineto{\pgfqpoint{1.130252in}{0.894942in}}%
\pgfpathlineto{\pgfqpoint{1.085585in}{0.893626in}}%
\pgfpathclose%
\pgfusepath{fill}%
\end{pgfscope}%
\begin{pgfscope}%
\pgfpathrectangle{\pgfqpoint{0.407592in}{0.222904in}}{\pgfqpoint{1.411163in}{1.411163in}}%
\pgfusepath{clip}%
\pgfsetbuttcap%
\pgfsetroundjoin%
\definecolor{currentfill}{rgb}{0.749020,0.772549,0.792157}%
\pgfsetfillcolor{currentfill}%
\pgfsetfillopacity{0.160000}%
\pgfsetlinewidth{0.000000pt}%
\definecolor{currentstroke}{rgb}{0.000000,0.000000,0.000000}%
\pgfsetstrokecolor{currentstroke}%
\pgfsetdash{}{0pt}%
\pgfpathmoveto{\pgfqpoint{0.996484in}{0.894582in}}%
\pgfpathlineto{\pgfqpoint{0.994204in}{0.901489in}}%
\pgfpathlineto{\pgfqpoint{0.991914in}{0.908430in}}%
\pgfpathlineto{\pgfqpoint{1.038509in}{0.907306in}}%
\pgfpathlineto{\pgfqpoint{1.085358in}{0.907432in}}%
\pgfpathlineto{\pgfqpoint{1.085472in}{0.900511in}}%
\pgfpathlineto{\pgfqpoint{1.085585in}{0.893626in}}%
\pgfpathlineto{\pgfqpoint{1.040914in}{0.893506in}}%
\pgfpathlineto{\pgfqpoint{0.996484in}{0.894582in}}%
\pgfpathclose%
\pgfusepath{fill}%
\end{pgfscope}%
\begin{pgfscope}%
\pgfpathrectangle{\pgfqpoint{0.407592in}{0.222904in}}{\pgfqpoint{1.411163in}{1.411163in}}%
\pgfusepath{clip}%
\pgfsetbuttcap%
\pgfsetroundjoin%
\definecolor{currentfill}{rgb}{0.749020,0.772549,0.792157}%
\pgfsetfillcolor{currentfill}%
\pgfsetfillopacity{0.160000}%
\pgfsetlinewidth{0.000000pt}%
\definecolor{currentstroke}{rgb}{0.000000,0.000000,0.000000}%
\pgfsetstrokecolor{currentstroke}%
\pgfsetdash{}{0pt}%
\pgfpathmoveto{\pgfqpoint{0.617531in}{0.992238in}}%
\pgfpathlineto{\pgfqpoint{0.607389in}{1.000428in}}%
\pgfpathlineto{\pgfqpoint{0.597202in}{1.008654in}}%
\pgfpathlineto{\pgfqpoint{0.637966in}{0.997208in}}%
\pgfpathlineto{\pgfqpoint{0.681152in}{0.986849in}}%
\pgfpathlineto{\pgfqpoint{0.689662in}{0.979021in}}%
\pgfpathlineto{\pgfqpoint{0.698131in}{0.971230in}}%
\pgfpathlineto{\pgfqpoint{0.656679in}{0.981209in}}%
\pgfpathlineto{\pgfqpoint{0.617531in}{0.992238in}}%
\pgfpathclose%
\pgfusepath{fill}%
\end{pgfscope}%
\begin{pgfscope}%
\pgfpathrectangle{\pgfqpoint{0.407592in}{0.222904in}}{\pgfqpoint{1.411163in}{1.411163in}}%
\pgfusepath{clip}%
\pgfsetbuttcap%
\pgfsetroundjoin%
\definecolor{currentfill}{rgb}{0.749020,0.772549,0.792157}%
\pgfsetfillcolor{currentfill}%
\pgfsetfillopacity{0.160000}%
\pgfsetlinewidth{0.000000pt}%
\definecolor{currentstroke}{rgb}{0.000000,0.000000,0.000000}%
\pgfsetstrokecolor{currentstroke}%
\pgfsetdash{}{0pt}%
\pgfpathmoveto{\pgfqpoint{0.500331in}{1.051998in}}%
\pgfpathlineto{\pgfqpoint{0.488577in}{1.060795in}}%
\pgfpathlineto{\pgfqpoint{0.525977in}{1.046683in}}%
\pgfpathlineto{\pgfqpoint{0.566372in}{1.033550in}}%
\pgfpathlineto{\pgfqpoint{0.576694in}{1.025215in}}%
\pgfpathlineto{\pgfqpoint{0.537050in}{1.038124in}}%
\pgfusepath{fill}%
\end{pgfscope}%
\begin{pgfscope}%
\pgfpathrectangle{\pgfqpoint{0.407592in}{0.222904in}}{\pgfqpoint{1.411163in}{1.411163in}}%
\pgfusepath{clip}%
\pgfsetbuttcap%
\pgfsetroundjoin%
\definecolor{currentfill}{rgb}{0.749020,0.772549,0.792157}%
\pgfsetfillcolor{currentfill}%
\pgfsetfillopacity{0.160000}%
\pgfsetlinewidth{0.000000pt}%
\definecolor{currentstroke}{rgb}{0.000000,0.000000,0.000000}%
\pgfsetstrokecolor{currentstroke}%
\pgfsetdash{}{0pt}%
\pgfpathmoveto{\pgfqpoint{1.487793in}{0.999419in}}%
\pgfpathlineto{\pgfqpoint{1.496068in}{1.007514in}}%
\pgfpathlineto{\pgfqpoint{1.504380in}{1.015645in}}%
\pgfpathlineto{\pgfqpoint{1.549243in}{1.027723in}}%
\pgfpathlineto{\pgfqpoint{1.591580in}{1.040833in}}%
\pgfpathlineto{\pgfqpoint{1.581596in}{1.032264in}}%
\pgfpathlineto{\pgfqpoint{1.571653in}{1.023730in}}%
\pgfpathlineto{\pgfqpoint{1.530926in}{1.011075in}}%
\pgfpathlineto{\pgfqpoint{1.487793in}{0.999419in}}%
\pgfpathclose%
\pgfusepath{fill}%
\end{pgfscope}%
\begin{pgfscope}%
\pgfpathrectangle{\pgfqpoint{0.407592in}{0.222904in}}{\pgfqpoint{1.411163in}{1.411163in}}%
\pgfusepath{clip}%
\pgfsetbuttcap%
\pgfsetroundjoin%
\definecolor{currentfill}{rgb}{0.749020,0.772549,0.792157}%
\pgfsetfillcolor{currentfill}%
\pgfsetfillopacity{0.160000}%
\pgfsetlinewidth{0.000000pt}%
\definecolor{currentstroke}{rgb}{0.000000,0.000000,0.000000}%
\pgfsetstrokecolor{currentstroke}%
\pgfsetdash{}{0pt}%
\pgfpathmoveto{\pgfqpoint{0.813123in}{0.925144in}}%
\pgfpathlineto{\pgfqpoint{0.806667in}{0.932473in}}%
\pgfpathlineto{\pgfqpoint{0.800180in}{0.939838in}}%
\pgfpathlineto{\pgfqpoint{0.845239in}{0.933623in}}%
\pgfpathlineto{\pgfqpoint{0.891613in}{0.928627in}}%
\pgfpathlineto{\pgfqpoint{0.896071in}{0.921489in}}%
\pgfpathlineto{\pgfqpoint{0.900505in}{0.914386in}}%
\pgfpathlineto{\pgfqpoint{0.856191in}{0.919180in}}%
\pgfpathlineto{\pgfqpoint{0.813123in}{0.925144in}}%
\pgfpathclose%
\pgfusepath{fill}%
\end{pgfscope}%
\begin{pgfscope}%
\pgfpathrectangle{\pgfqpoint{0.407592in}{0.222904in}}{\pgfqpoint{1.411163in}{1.411163in}}%
\pgfusepath{clip}%
\pgfsetbuttcap%
\pgfsetroundjoin%
\definecolor{currentfill}{rgb}{0.749020,0.772549,0.792157}%
\pgfsetfillcolor{currentfill}%
\pgfsetfillopacity{0.160000}%
\pgfsetlinewidth{0.000000pt}%
\definecolor{currentstroke}{rgb}{0.000000,0.000000,0.000000}%
\pgfsetstrokecolor{currentstroke}%
\pgfsetdash{}{0pt}%
\pgfpathmoveto{\pgfqpoint{0.714953in}{0.955755in}}%
\pgfpathlineto{\pgfqpoint{0.706562in}{0.963474in}}%
\pgfpathlineto{\pgfqpoint{0.698131in}{0.971230in}}%
\pgfpathlineto{\pgfqpoint{0.741683in}{0.962366in}}%
\pgfpathlineto{\pgfqpoint{0.787111in}{0.954676in}}%
\pgfpathlineto{\pgfqpoint{0.793661in}{0.947239in}}%
\pgfpathlineto{\pgfqpoint{0.800180in}{0.939838in}}%
\pgfpathlineto{\pgfqpoint{0.756676in}{0.947232in}}%
\pgfpathlineto{\pgfqpoint{0.714953in}{0.955755in}}%
\pgfpathclose%
\pgfusepath{fill}%
\end{pgfscope}%
\begin{pgfscope}%
\pgfpathrectangle{\pgfqpoint{0.407592in}{0.222904in}}{\pgfqpoint{1.411163in}{1.411163in}}%
\pgfusepath{clip}%
\pgfsetbuttcap%
\pgfsetroundjoin%
\definecolor{currentfill}{rgb}{0.749020,0.772549,0.792157}%
\pgfsetfillcolor{currentfill}%
\pgfsetfillopacity{0.160000}%
\pgfsetlinewidth{0.000000pt}%
\definecolor{currentstroke}{rgb}{0.000000,0.000000,0.000000}%
\pgfsetstrokecolor{currentstroke}%
\pgfsetdash{}{0pt}%
\pgfpathmoveto{\pgfqpoint{1.178785in}{0.911418in}}%
\pgfpathlineto{\pgfqpoint{1.180857in}{0.918458in}}%
\pgfpathlineto{\pgfqpoint{1.182939in}{0.925533in}}%
\pgfpathlineto{\pgfqpoint{1.231156in}{0.929530in}}%
\pgfpathlineto{\pgfqpoint{1.278561in}{0.934774in}}%
\pgfpathlineto{\pgfqpoint{1.274342in}{0.927511in}}%
\pgfpathlineto{\pgfqpoint{1.270144in}{0.920284in}}%
\pgfpathlineto{\pgfqpoint{1.224849in}{0.915253in}}%
\pgfpathlineto{\pgfqpoint{1.178785in}{0.911418in}}%
\pgfpathclose%
\pgfusepath{fill}%
\end{pgfscope}%
\begin{pgfscope}%
\pgfpathrectangle{\pgfqpoint{0.407592in}{0.222904in}}{\pgfqpoint{1.411163in}{1.411163in}}%
\pgfusepath{clip}%
\pgfsetbuttcap%
\pgfsetroundjoin%
\definecolor{currentfill}{rgb}{0.749020,0.772549,0.792157}%
\pgfsetfillcolor{currentfill}%
\pgfsetfillopacity{0.160000}%
\pgfsetlinewidth{0.000000pt}%
\definecolor{currentstroke}{rgb}{0.000000,0.000000,0.000000}%
\pgfsetstrokecolor{currentstroke}%
\pgfsetdash{}{0pt}%
\pgfpathmoveto{\pgfqpoint{1.382476in}{0.964052in}}%
\pgfpathlineto{\pgfqpoint{1.388803in}{0.971706in}}%
\pgfpathlineto{\pgfqpoint{1.395161in}{0.979396in}}%
\pgfpathlineto{\pgfqpoint{1.442462in}{0.988836in}}%
\pgfpathlineto{\pgfqpoint{1.487793in}{0.999419in}}%
\pgfpathlineto{\pgfqpoint{1.479555in}{0.991361in}}%
\pgfpathlineto{\pgfqpoint{1.471354in}{0.983339in}}%
\pgfpathlineto{\pgfqpoint{1.427851in}{0.973144in}}%
\pgfpathlineto{\pgfqpoint{1.382476in}{0.964052in}}%
\pgfpathclose%
\pgfusepath{fill}%
\end{pgfscope}%
\begin{pgfscope}%
\pgfpathrectangle{\pgfqpoint{0.407592in}{0.222904in}}{\pgfqpoint{1.411163in}{1.411163in}}%
\pgfusepath{clip}%
\pgfsetbuttcap%
\pgfsetroundjoin%
\definecolor{currentfill}{rgb}{0.749020,0.772549,0.792157}%
\pgfsetfillcolor{currentfill}%
\pgfsetfillopacity{0.160000}%
\pgfsetlinewidth{0.000000pt}%
\definecolor{currentstroke}{rgb}{0.000000,0.000000,0.000000}%
\pgfsetstrokecolor{currentstroke}%
\pgfsetdash{}{0pt}%
\pgfpathmoveto{\pgfqpoint{1.278561in}{0.934774in}}%
\pgfpathlineto{\pgfqpoint{1.282800in}{0.942072in}}%
\pgfpathlineto{\pgfqpoint{1.287060in}{0.949407in}}%
\pgfpathlineto{\pgfqpoint{1.335463in}{0.956122in}}%
\pgfpathlineto{\pgfqpoint{1.382476in}{0.964052in}}%
\pgfpathlineto{\pgfqpoint{1.376179in}{0.956434in}}%
\pgfpathlineto{\pgfqpoint{1.369911in}{0.948853in}}%
\pgfpathlineto{\pgfqpoint{1.324895in}{0.941229in}}%
\pgfpathlineto{\pgfqpoint{1.278561in}{0.934774in}}%
\pgfpathclose%
\pgfusepath{fill}%
\end{pgfscope}%
\begin{pgfscope}%
\pgfpathrectangle{\pgfqpoint{0.407592in}{0.222904in}}{\pgfqpoint{1.411163in}{1.411163in}}%
\pgfusepath{clip}%
\pgfsetbuttcap%
\pgfsetroundjoin%
\definecolor{currentfill}{rgb}{0.749020,0.772549,0.792157}%
\pgfsetfillcolor{currentfill}%
\pgfsetfillopacity{0.160000}%
\pgfsetlinewidth{0.000000pt}%
\definecolor{currentstroke}{rgb}{0.000000,0.000000,0.000000}%
\pgfsetstrokecolor{currentstroke}%
\pgfsetdash{}{0pt}%
\pgfpathmoveto{\pgfqpoint{0.900505in}{0.914386in}}%
\pgfpathlineto{\pgfqpoint{0.896071in}{0.921489in}}%
\pgfpathlineto{\pgfqpoint{0.891613in}{0.928627in}}%
\pgfpathlineto{\pgfqpoint{0.939053in}{0.924884in}}%
\pgfpathlineto{\pgfqpoint{0.987297in}{0.922418in}}%
\pgfpathlineto{\pgfqpoint{0.989611in}{0.915406in}}%
\pgfpathlineto{\pgfqpoint{0.991914in}{0.908430in}}%
\pgfpathlineto{\pgfqpoint{0.945828in}{0.910795in}}%
\pgfpathlineto{\pgfqpoint{0.900505in}{0.914386in}}%
\pgfpathclose%
\pgfusepath{fill}%
\end{pgfscope}%
\begin{pgfscope}%
\pgfpathrectangle{\pgfqpoint{0.407592in}{0.222904in}}{\pgfqpoint{1.411163in}{1.411163in}}%
\pgfusepath{clip}%
\pgfsetbuttcap%
\pgfsetroundjoin%
\definecolor{currentfill}{rgb}{0.749020,0.772549,0.792157}%
\pgfsetfillcolor{currentfill}%
\pgfsetfillopacity{0.160000}%
\pgfsetlinewidth{0.000000pt}%
\definecolor{currentstroke}{rgb}{0.000000,0.000000,0.000000}%
\pgfsetstrokecolor{currentstroke}%
\pgfsetdash{}{0pt}%
\pgfpathmoveto{\pgfqpoint{1.504380in}{1.015645in}}%
\pgfpathlineto{\pgfqpoint{1.512730in}{1.023812in}}%
\pgfpathlineto{\pgfqpoint{1.558462in}{1.036102in}}%
\pgfpathlineto{\pgfqpoint{1.601606in}{1.049438in}}%
\pgfpathlineto{\pgfqpoint{1.591580in}{1.040833in}}%
\pgfpathlineto{\pgfqpoint{1.549243in}{1.027723in}}%
\pgfusepath{fill}%
\end{pgfscope}%
\begin{pgfscope}%
\pgfpathrectangle{\pgfqpoint{0.407592in}{0.222904in}}{\pgfqpoint{1.411163in}{1.411163in}}%
\pgfusepath{clip}%
\pgfsetbuttcap%
\pgfsetroundjoin%
\definecolor{currentfill}{rgb}{0.749020,0.772549,0.792157}%
\pgfsetfillcolor{currentfill}%
\pgfsetfillopacity{0.160000}%
\pgfsetlinewidth{0.000000pt}%
\definecolor{currentstroke}{rgb}{0.000000,0.000000,0.000000}%
\pgfsetstrokecolor{currentstroke}%
\pgfsetdash{}{0pt}%
\pgfpathmoveto{\pgfqpoint{1.085358in}{0.907432in}}%
\pgfpathlineto{\pgfqpoint{1.085244in}{0.914387in}}%
\pgfpathlineto{\pgfqpoint{1.085129in}{0.921378in}}%
\pgfpathlineto{\pgfqpoint{1.134173in}{0.922810in}}%
\pgfpathlineto{\pgfqpoint{1.182939in}{0.925533in}}%
\pgfpathlineto{\pgfqpoint{1.180857in}{0.918458in}}%
\pgfpathlineto{\pgfqpoint{1.178785in}{0.911418in}}%
\pgfpathlineto{\pgfqpoint{1.132203in}{0.908805in}}%
\pgfpathlineto{\pgfqpoint{1.085358in}{0.907432in}}%
\pgfpathclose%
\pgfusepath{fill}%
\end{pgfscope}%
\begin{pgfscope}%
\pgfpathrectangle{\pgfqpoint{0.407592in}{0.222904in}}{\pgfqpoint{1.411163in}{1.411163in}}%
\pgfusepath{clip}%
\pgfsetbuttcap%
\pgfsetroundjoin%
\definecolor{currentfill}{rgb}{0.749020,0.772549,0.792157}%
\pgfsetfillcolor{currentfill}%
\pgfsetfillopacity{0.160000}%
\pgfsetlinewidth{0.000000pt}%
\definecolor{currentstroke}{rgb}{0.000000,0.000000,0.000000}%
\pgfsetstrokecolor{currentstroke}%
\pgfsetdash{}{0pt}%
\pgfpathmoveto{\pgfqpoint{0.991914in}{0.908430in}}%
\pgfpathlineto{\pgfqpoint{0.989611in}{0.915406in}}%
\pgfpathlineto{\pgfqpoint{0.987297in}{0.922418in}}%
\pgfpathlineto{\pgfqpoint{1.036079in}{0.921247in}}%
\pgfpathlineto{\pgfqpoint{1.085129in}{0.921378in}}%
\pgfpathlineto{\pgfqpoint{1.085244in}{0.914387in}}%
\pgfpathlineto{\pgfqpoint{1.085358in}{0.907432in}}%
\pgfpathlineto{\pgfqpoint{1.038509in}{0.907306in}}%
\pgfpathlineto{\pgfqpoint{0.991914in}{0.908430in}}%
\pgfpathclose%
\pgfusepath{fill}%
\end{pgfscope}%
\begin{pgfscope}%
\pgfpathrectangle{\pgfqpoint{0.407592in}{0.222904in}}{\pgfqpoint{1.411163in}{1.411163in}}%
\pgfusepath{clip}%
\pgfsetbuttcap%
\pgfsetroundjoin%
\definecolor{currentfill}{rgb}{0.749020,0.772549,0.792157}%
\pgfsetfillcolor{currentfill}%
\pgfsetfillopacity{0.160000}%
\pgfsetlinewidth{0.000000pt}%
\definecolor{currentstroke}{rgb}{0.000000,0.000000,0.000000}%
\pgfsetstrokecolor{currentstroke}%
\pgfsetdash{}{0pt}%
\pgfpathmoveto{\pgfqpoint{0.597202in}{1.008654in}}%
\pgfpathlineto{\pgfqpoint{0.586970in}{1.016917in}}%
\pgfpathlineto{\pgfqpoint{0.576694in}{1.025215in}}%
\pgfpathlineto{\pgfqpoint{0.619083in}{1.013352in}}%
\pgfpathlineto{\pgfqpoint{0.664014in}{1.002614in}}%
\pgfpathlineto{\pgfqpoint{0.672603in}{0.994713in}}%
\pgfpathlineto{\pgfqpoint{0.681152in}{0.986849in}}%
\pgfpathlineto{\pgfqpoint{0.637966in}{0.997208in}}%
\pgfpathlineto{\pgfqpoint{0.597202in}{1.008654in}}%
\pgfpathclose%
\pgfusepath{fill}%
\end{pgfscope}%
\begin{pgfscope}%
\pgfpathrectangle{\pgfqpoint{0.407592in}{0.222904in}}{\pgfqpoint{1.411163in}{1.411163in}}%
\pgfusepath{clip}%
\pgfsetbuttcap%
\pgfsetroundjoin%
\definecolor{currentfill}{rgb}{0.749020,0.772549,0.792157}%
\pgfsetfillcolor{currentfill}%
\pgfsetfillopacity{0.160000}%
\pgfsetlinewidth{0.000000pt}%
\definecolor{currentstroke}{rgb}{0.000000,0.000000,0.000000}%
\pgfsetstrokecolor{currentstroke}%
\pgfsetdash{}{0pt}%
\pgfpathmoveto{\pgfqpoint{0.800180in}{0.939838in}}%
\pgfpathlineto{\pgfqpoint{0.793661in}{0.947239in}}%
\pgfpathlineto{\pgfqpoint{0.787111in}{0.954676in}}%
\pgfpathlineto{\pgfqpoint{0.834178in}{0.948210in}}%
\pgfpathlineto{\pgfqpoint{0.882632in}{0.943011in}}%
\pgfpathlineto{\pgfqpoint{0.887134in}{0.935801in}}%
\pgfpathlineto{\pgfqpoint{0.891613in}{0.928627in}}%
\pgfpathlineto{\pgfqpoint{0.845239in}{0.933623in}}%
\pgfpathlineto{\pgfqpoint{0.800180in}{0.939838in}}%
\pgfpathclose%
\pgfusepath{fill}%
\end{pgfscope}%
\begin{pgfscope}%
\pgfpathrectangle{\pgfqpoint{0.407592in}{0.222904in}}{\pgfqpoint{1.411163in}{1.411163in}}%
\pgfusepath{clip}%
\pgfsetbuttcap%
\pgfsetroundjoin%
\definecolor{currentfill}{rgb}{0.749020,0.772549,0.792157}%
\pgfsetfillcolor{currentfill}%
\pgfsetfillopacity{0.160000}%
\pgfsetlinewidth{0.000000pt}%
\definecolor{currentstroke}{rgb}{0.000000,0.000000,0.000000}%
\pgfsetstrokecolor{currentstroke}%
\pgfsetdash{}{0pt}%
\pgfpathmoveto{\pgfqpoint{0.698131in}{0.971230in}}%
\pgfpathlineto{\pgfqpoint{0.689662in}{0.979021in}}%
\pgfpathlineto{\pgfqpoint{0.681152in}{0.986849in}}%
\pgfpathlineto{\pgfqpoint{0.726546in}{0.977645in}}%
\pgfpathlineto{\pgfqpoint{0.773915in}{0.969658in}}%
\pgfpathlineto{\pgfqpoint{0.780529in}{0.962149in}}%
\pgfpathlineto{\pgfqpoint{0.787111in}{0.954676in}}%
\pgfpathlineto{\pgfqpoint{0.741683in}{0.962366in}}%
\pgfpathlineto{\pgfqpoint{0.698131in}{0.971230in}}%
\pgfpathclose%
\pgfusepath{fill}%
\end{pgfscope}%
\begin{pgfscope}%
\pgfpathrectangle{\pgfqpoint{0.407592in}{0.222904in}}{\pgfqpoint{1.411163in}{1.411163in}}%
\pgfusepath{clip}%
\pgfsetbuttcap%
\pgfsetroundjoin%
\definecolor{currentfill}{rgb}{0.749020,0.772549,0.792157}%
\pgfsetfillcolor{currentfill}%
\pgfsetfillopacity{0.160000}%
\pgfsetlinewidth{0.000000pt}%
\definecolor{currentstroke}{rgb}{0.000000,0.000000,0.000000}%
\pgfsetstrokecolor{currentstroke}%
\pgfsetdash{}{0pt}%
\pgfpathmoveto{\pgfqpoint{1.182939in}{0.925533in}}%
\pgfpathlineto{\pgfqpoint{1.185032in}{0.932644in}}%
\pgfpathlineto{\pgfqpoint{1.187135in}{0.939791in}}%
\pgfpathlineto{\pgfqpoint{1.237527in}{0.943951in}}%
\pgfpathlineto{\pgfqpoint{1.287060in}{0.949407in}}%
\pgfpathlineto{\pgfqpoint{1.282800in}{0.942072in}}%
\pgfpathlineto{\pgfqpoint{1.278561in}{0.934774in}}%
\pgfpathlineto{\pgfqpoint{1.231156in}{0.929530in}}%
\pgfpathlineto{\pgfqpoint{1.182939in}{0.925533in}}%
\pgfpathclose%
\pgfusepath{fill}%
\end{pgfscope}%
\begin{pgfscope}%
\pgfpathrectangle{\pgfqpoint{0.407592in}{0.222904in}}{\pgfqpoint{1.411163in}{1.411163in}}%
\pgfusepath{clip}%
\pgfsetbuttcap%
\pgfsetroundjoin%
\definecolor{currentfill}{rgb}{0.749020,0.772549,0.792157}%
\pgfsetfillcolor{currentfill}%
\pgfsetfillopacity{0.160000}%
\pgfsetlinewidth{0.000000pt}%
\definecolor{currentstroke}{rgb}{0.000000,0.000000,0.000000}%
\pgfsetstrokecolor{currentstroke}%
\pgfsetdash{}{0pt}%
\pgfpathmoveto{\pgfqpoint{1.395161in}{0.979396in}}%
\pgfpathlineto{\pgfqpoint{1.401548in}{0.987123in}}%
\pgfpathlineto{\pgfqpoint{1.407967in}{0.994887in}}%
\pgfpathlineto{\pgfqpoint{1.457209in}{1.004675in}}%
\pgfpathlineto{\pgfqpoint{1.504380in}{1.015645in}}%
\pgfpathlineto{\pgfqpoint{1.496068in}{1.007514in}}%
\pgfpathlineto{\pgfqpoint{1.487793in}{0.999419in}}%
\pgfpathlineto{\pgfqpoint{1.442462in}{0.988836in}}%
\pgfpathlineto{\pgfqpoint{1.395161in}{0.979396in}}%
\pgfpathclose%
\pgfusepath{fill}%
\end{pgfscope}%
\begin{pgfscope}%
\pgfpathrectangle{\pgfqpoint{0.407592in}{0.222904in}}{\pgfqpoint{1.411163in}{1.411163in}}%
\pgfusepath{clip}%
\pgfsetbuttcap%
\pgfsetroundjoin%
\definecolor{currentfill}{rgb}{0.749020,0.772549,0.792157}%
\pgfsetfillcolor{currentfill}%
\pgfsetfillopacity{0.160000}%
\pgfsetlinewidth{0.000000pt}%
\definecolor{currentstroke}{rgb}{0.000000,0.000000,0.000000}%
\pgfsetstrokecolor{currentstroke}%
\pgfsetdash{}{0pt}%
\pgfpathmoveto{\pgfqpoint{0.576694in}{1.025215in}}%
\pgfpathlineto{\pgfqpoint{0.566372in}{1.033550in}}%
\pgfpathlineto{\pgfqpoint{0.609577in}{1.021480in}}%
\pgfpathlineto{\pgfqpoint{0.655385in}{1.010552in}}%
\pgfpathlineto{\pgfqpoint{0.664014in}{1.002614in}}%
\pgfpathlineto{\pgfqpoint{0.619083in}{1.013352in}}%
\pgfusepath{fill}%
\end{pgfscope}%
\begin{pgfscope}%
\pgfpathrectangle{\pgfqpoint{0.407592in}{0.222904in}}{\pgfqpoint{1.411163in}{1.411163in}}%
\pgfusepath{clip}%
\pgfsetbuttcap%
\pgfsetroundjoin%
\definecolor{currentfill}{rgb}{0.749020,0.772549,0.792157}%
\pgfsetfillcolor{currentfill}%
\pgfsetfillopacity{0.160000}%
\pgfsetlinewidth{0.000000pt}%
\definecolor{currentstroke}{rgb}{0.000000,0.000000,0.000000}%
\pgfsetstrokecolor{currentstroke}%
\pgfsetdash{}{0pt}%
\pgfpathmoveto{\pgfqpoint{1.287060in}{0.949407in}}%
\pgfpathlineto{\pgfqpoint{1.291341in}{0.956777in}}%
\pgfpathlineto{\pgfqpoint{1.295644in}{0.964185in}}%
\pgfpathlineto{\pgfqpoint{1.346134in}{0.971160in}}%
\pgfpathlineto{\pgfqpoint{1.395161in}{0.979396in}}%
\pgfpathlineto{\pgfqpoint{1.388803in}{0.971706in}}%
\pgfpathlineto{\pgfqpoint{1.382476in}{0.964052in}}%
\pgfpathlineto{\pgfqpoint{1.335463in}{0.956122in}}%
\pgfpathlineto{\pgfqpoint{1.287060in}{0.949407in}}%
\pgfpathclose%
\pgfusepath{fill}%
\end{pgfscope}%
\begin{pgfscope}%
\pgfpathrectangle{\pgfqpoint{0.407592in}{0.222904in}}{\pgfqpoint{1.411163in}{1.411163in}}%
\pgfusepath{clip}%
\pgfsetbuttcap%
\pgfsetroundjoin%
\definecolor{currentfill}{rgb}{0.749020,0.772549,0.792157}%
\pgfsetfillcolor{currentfill}%
\pgfsetfillopacity{0.160000}%
\pgfsetlinewidth{0.000000pt}%
\definecolor{currentstroke}{rgb}{0.000000,0.000000,0.000000}%
\pgfsetstrokecolor{currentstroke}%
\pgfsetdash{}{0pt}%
\pgfpathmoveto{\pgfqpoint{0.891613in}{0.928627in}}%
\pgfpathlineto{\pgfqpoint{0.887134in}{0.935801in}}%
\pgfpathlineto{\pgfqpoint{0.882632in}{0.943011in}}%
\pgfpathlineto{\pgfqpoint{0.932208in}{0.939116in}}%
\pgfpathlineto{\pgfqpoint{0.982633in}{0.936550in}}%
\pgfpathlineto{\pgfqpoint{0.984971in}{0.929466in}}%
\pgfpathlineto{\pgfqpoint{0.987297in}{0.922418in}}%
\pgfpathlineto{\pgfqpoint{0.939053in}{0.924884in}}%
\pgfpathlineto{\pgfqpoint{0.891613in}{0.928627in}}%
\pgfpathclose%
\pgfusepath{fill}%
\end{pgfscope}%
\begin{pgfscope}%
\pgfpathrectangle{\pgfqpoint{0.407592in}{0.222904in}}{\pgfqpoint{1.411163in}{1.411163in}}%
\pgfusepath{clip}%
\pgfsetbuttcap%
\pgfsetroundjoin%
\definecolor{currentfill}{rgb}{0.749020,0.772549,0.792157}%
\pgfsetfillcolor{currentfill}%
\pgfsetfillopacity{0.160000}%
\pgfsetlinewidth{0.000000pt}%
\definecolor{currentstroke}{rgb}{0.000000,0.000000,0.000000}%
\pgfsetstrokecolor{currentstroke}%
\pgfsetdash{}{0pt}%
\pgfpathmoveto{\pgfqpoint{1.085129in}{0.921378in}}%
\pgfpathlineto{\pgfqpoint{1.085014in}{0.928404in}}%
\pgfpathlineto{\pgfqpoint{1.084898in}{0.935467in}}%
\pgfpathlineto{\pgfqpoint{1.136164in}{0.936957in}}%
\pgfpathlineto{\pgfqpoint{1.187135in}{0.939791in}}%
\pgfpathlineto{\pgfqpoint{1.185032in}{0.932644in}}%
\pgfpathlineto{\pgfqpoint{1.182939in}{0.925533in}}%
\pgfpathlineto{\pgfqpoint{1.134173in}{0.922810in}}%
\pgfpathlineto{\pgfqpoint{1.085129in}{0.921378in}}%
\pgfpathclose%
\pgfusepath{fill}%
\end{pgfscope}%
\begin{pgfscope}%
\pgfpathrectangle{\pgfqpoint{0.407592in}{0.222904in}}{\pgfqpoint{1.411163in}{1.411163in}}%
\pgfusepath{clip}%
\pgfsetbuttcap%
\pgfsetroundjoin%
\definecolor{currentfill}{rgb}{0.749020,0.772549,0.792157}%
\pgfsetfillcolor{currentfill}%
\pgfsetfillopacity{0.160000}%
\pgfsetlinewidth{0.000000pt}%
\definecolor{currentstroke}{rgb}{0.000000,0.000000,0.000000}%
\pgfsetstrokecolor{currentstroke}%
\pgfsetdash{}{0pt}%
\pgfpathmoveto{\pgfqpoint{0.987297in}{0.922418in}}%
\pgfpathlineto{\pgfqpoint{0.984971in}{0.929466in}}%
\pgfpathlineto{\pgfqpoint{0.982633in}{0.936550in}}%
\pgfpathlineto{\pgfqpoint{1.033625in}{0.935330in}}%
\pgfpathlineto{\pgfqpoint{1.084898in}{0.935467in}}%
\pgfpathlineto{\pgfqpoint{1.085014in}{0.928404in}}%
\pgfpathlineto{\pgfqpoint{1.085129in}{0.921378in}}%
\pgfpathlineto{\pgfqpoint{1.036079in}{0.921247in}}%
\pgfpathlineto{\pgfqpoint{0.987297in}{0.922418in}}%
\pgfpathclose%
\pgfusepath{fill}%
\end{pgfscope}%
\begin{pgfscope}%
\pgfpathrectangle{\pgfqpoint{0.407592in}{0.222904in}}{\pgfqpoint{1.411163in}{1.411163in}}%
\pgfusepath{clip}%
\pgfsetbuttcap%
\pgfsetroundjoin%
\definecolor{currentfill}{rgb}{0.749020,0.772549,0.792157}%
\pgfsetfillcolor{currentfill}%
\pgfsetfillopacity{0.160000}%
\pgfsetlinewidth{0.000000pt}%
\definecolor{currentstroke}{rgb}{0.000000,0.000000,0.000000}%
\pgfsetstrokecolor{currentstroke}%
\pgfsetdash{}{0pt}%
\pgfpathmoveto{\pgfqpoint{0.787111in}{0.954676in}}%
\pgfpathlineto{\pgfqpoint{0.780529in}{0.962149in}}%
\pgfpathlineto{\pgfqpoint{0.773915in}{0.969658in}}%
\pgfpathlineto{\pgfqpoint{0.823008in}{0.962941in}}%
\pgfpathlineto{\pgfqpoint{0.873560in}{0.957540in}}%
\pgfpathlineto{\pgfqpoint{0.878108in}{0.950258in}}%
\pgfpathlineto{\pgfqpoint{0.882632in}{0.943011in}}%
\pgfpathlineto{\pgfqpoint{0.834178in}{0.948210in}}%
\pgfpathlineto{\pgfqpoint{0.787111in}{0.954676in}}%
\pgfpathclose%
\pgfusepath{fill}%
\end{pgfscope}%
\begin{pgfscope}%
\pgfpathrectangle{\pgfqpoint{0.407592in}{0.222904in}}{\pgfqpoint{1.411163in}{1.411163in}}%
\pgfusepath{clip}%
\pgfsetbuttcap%
\pgfsetroundjoin%
\definecolor{currentfill}{rgb}{0.749020,0.772549,0.792157}%
\pgfsetfillcolor{currentfill}%
\pgfsetfillopacity{0.160000}%
\pgfsetlinewidth{0.000000pt}%
\definecolor{currentstroke}{rgb}{0.000000,0.000000,0.000000}%
\pgfsetstrokecolor{currentstroke}%
\pgfsetdash{}{0pt}%
\pgfpathmoveto{\pgfqpoint{1.407967in}{0.994887in}}%
\pgfpathlineto{\pgfqpoint{1.414415in}{1.002688in}}%
\pgfpathlineto{\pgfqpoint{1.464634in}{1.012649in}}%
\pgfpathlineto{\pgfqpoint{1.512730in}{1.023812in}}%
\pgfpathlineto{\pgfqpoint{1.504380in}{1.015645in}}%
\pgfpathlineto{\pgfqpoint{1.457209in}{1.004675in}}%
\pgfusepath{fill}%
\end{pgfscope}%
\begin{pgfscope}%
\pgfpathrectangle{\pgfqpoint{0.407592in}{0.222904in}}{\pgfqpoint{1.411163in}{1.411163in}}%
\pgfusepath{clip}%
\pgfsetbuttcap%
\pgfsetroundjoin%
\definecolor{currentfill}{rgb}{0.749020,0.772549,0.792157}%
\pgfsetfillcolor{currentfill}%
\pgfsetfillopacity{0.160000}%
\pgfsetlinewidth{0.000000pt}%
\definecolor{currentstroke}{rgb}{0.000000,0.000000,0.000000}%
\pgfsetstrokecolor{currentstroke}%
\pgfsetdash{}{0pt}%
\pgfpathmoveto{\pgfqpoint{0.681152in}{0.986849in}}%
\pgfpathlineto{\pgfqpoint{0.672603in}{0.994713in}}%
\pgfpathlineto{\pgfqpoint{0.664014in}{1.002614in}}%
\pgfpathlineto{\pgfqpoint{0.711264in}{0.993071in}}%
\pgfpathlineto{\pgfqpoint{0.760589in}{0.984787in}}%
\pgfpathlineto{\pgfqpoint{0.767268in}{0.977204in}}%
\pgfpathlineto{\pgfqpoint{0.773915in}{0.969658in}}%
\pgfpathlineto{\pgfqpoint{0.726546in}{0.977645in}}%
\pgfpathlineto{\pgfqpoint{0.681152in}{0.986849in}}%
\pgfpathclose%
\pgfusepath{fill}%
\end{pgfscope}%
\begin{pgfscope}%
\pgfpathrectangle{\pgfqpoint{0.407592in}{0.222904in}}{\pgfqpoint{1.411163in}{1.411163in}}%
\pgfusepath{clip}%
\pgfsetbuttcap%
\pgfsetroundjoin%
\definecolor{currentfill}{rgb}{0.749020,0.772549,0.792157}%
\pgfsetfillcolor{currentfill}%
\pgfsetfillopacity{0.160000}%
\pgfsetlinewidth{0.000000pt}%
\definecolor{currentstroke}{rgb}{0.000000,0.000000,0.000000}%
\pgfsetstrokecolor{currentstroke}%
\pgfsetdash{}{0pt}%
\pgfpathmoveto{\pgfqpoint{1.187135in}{0.939791in}}%
\pgfpathlineto{\pgfqpoint{1.189250in}{0.946975in}}%
\pgfpathlineto{\pgfqpoint{1.191375in}{0.954195in}}%
\pgfpathlineto{\pgfqpoint{1.243962in}{0.958517in}}%
\pgfpathlineto{\pgfqpoint{1.295644in}{0.964185in}}%
\pgfpathlineto{\pgfqpoint{1.291341in}{0.956777in}}%
\pgfpathlineto{\pgfqpoint{1.287060in}{0.949407in}}%
\pgfpathlineto{\pgfqpoint{1.237527in}{0.943951in}}%
\pgfpathlineto{\pgfqpoint{1.187135in}{0.939791in}}%
\pgfpathclose%
\pgfusepath{fill}%
\end{pgfscope}%
\begin{pgfscope}%
\pgfpathrectangle{\pgfqpoint{0.407592in}{0.222904in}}{\pgfqpoint{1.411163in}{1.411163in}}%
\pgfusepath{clip}%
\pgfsetbuttcap%
\pgfsetroundjoin%
\definecolor{currentfill}{rgb}{0.749020,0.772549,0.792157}%
\pgfsetfillcolor{currentfill}%
\pgfsetfillopacity{0.160000}%
\pgfsetlinewidth{0.000000pt}%
\definecolor{currentstroke}{rgb}{0.000000,0.000000,0.000000}%
\pgfsetstrokecolor{currentstroke}%
\pgfsetdash{}{0pt}%
\pgfpathmoveto{\pgfqpoint{1.295644in}{0.964185in}}%
\pgfpathlineto{\pgfqpoint{1.299967in}{0.971629in}}%
\pgfpathlineto{\pgfqpoint{1.304313in}{0.979110in}}%
\pgfpathlineto{\pgfqpoint{1.356910in}{0.986345in}}%
\pgfpathlineto{\pgfqpoint{1.407967in}{0.994887in}}%
\pgfpathlineto{\pgfqpoint{1.401548in}{0.987123in}}%
\pgfpathlineto{\pgfqpoint{1.395161in}{0.979396in}}%
\pgfpathlineto{\pgfqpoint{1.346134in}{0.971160in}}%
\pgfpathlineto{\pgfqpoint{1.295644in}{0.964185in}}%
\pgfpathclose%
\pgfusepath{fill}%
\end{pgfscope}%
\begin{pgfscope}%
\pgfpathrectangle{\pgfqpoint{0.407592in}{0.222904in}}{\pgfqpoint{1.411163in}{1.411163in}}%
\pgfusepath{clip}%
\pgfsetbuttcap%
\pgfsetroundjoin%
\definecolor{currentfill}{rgb}{0.749020,0.772549,0.792157}%
\pgfsetfillcolor{currentfill}%
\pgfsetfillopacity{0.160000}%
\pgfsetlinewidth{0.000000pt}%
\definecolor{currentstroke}{rgb}{0.000000,0.000000,0.000000}%
\pgfsetstrokecolor{currentstroke}%
\pgfsetdash{}{0pt}%
\pgfpathmoveto{\pgfqpoint{0.882632in}{0.943011in}}%
\pgfpathlineto{\pgfqpoint{0.878108in}{0.950258in}}%
\pgfpathlineto{\pgfqpoint{0.873560in}{0.957540in}}%
\pgfpathlineto{\pgfqpoint{0.925294in}{0.953493in}}%
\pgfpathlineto{\pgfqpoint{0.977921in}{0.950826in}}%
\pgfpathlineto{\pgfqpoint{0.980283in}{0.943670in}}%
\pgfpathlineto{\pgfqpoint{0.982633in}{0.936550in}}%
\pgfpathlineto{\pgfqpoint{0.932208in}{0.939116in}}%
\pgfpathlineto{\pgfqpoint{0.882632in}{0.943011in}}%
\pgfpathclose%
\pgfusepath{fill}%
\end{pgfscope}%
\begin{pgfscope}%
\pgfpathrectangle{\pgfqpoint{0.407592in}{0.222904in}}{\pgfqpoint{1.411163in}{1.411163in}}%
\pgfusepath{clip}%
\pgfsetbuttcap%
\pgfsetroundjoin%
\definecolor{currentfill}{rgb}{0.749020,0.772549,0.792157}%
\pgfsetfillcolor{currentfill}%
\pgfsetfillopacity{0.160000}%
\pgfsetlinewidth{0.000000pt}%
\definecolor{currentstroke}{rgb}{0.000000,0.000000,0.000000}%
\pgfsetstrokecolor{currentstroke}%
\pgfsetdash{}{0pt}%
\pgfpathmoveto{\pgfqpoint{1.084898in}{0.935467in}}%
\pgfpathlineto{\pgfqpoint{1.084782in}{0.942565in}}%
\pgfpathlineto{\pgfqpoint{1.084665in}{0.949700in}}%
\pgfpathlineto{\pgfqpoint{1.138175in}{0.951249in}}%
\pgfpathlineto{\pgfqpoint{1.191375in}{0.954195in}}%
\pgfpathlineto{\pgfqpoint{1.189250in}{0.946975in}}%
\pgfpathlineto{\pgfqpoint{1.187135in}{0.939791in}}%
\pgfpathlineto{\pgfqpoint{1.136164in}{0.936957in}}%
\pgfpathlineto{\pgfqpoint{1.084898in}{0.935467in}}%
\pgfpathclose%
\pgfusepath{fill}%
\end{pgfscope}%
\begin{pgfscope}%
\pgfpathrectangle{\pgfqpoint{0.407592in}{0.222904in}}{\pgfqpoint{1.411163in}{1.411163in}}%
\pgfusepath{clip}%
\pgfsetbuttcap%
\pgfsetroundjoin%
\definecolor{currentfill}{rgb}{0.749020,0.772549,0.792157}%
\pgfsetfillcolor{currentfill}%
\pgfsetfillopacity{0.160000}%
\pgfsetlinewidth{0.000000pt}%
\definecolor{currentstroke}{rgb}{0.000000,0.000000,0.000000}%
\pgfsetstrokecolor{currentstroke}%
\pgfsetdash{}{0pt}%
\pgfpathmoveto{\pgfqpoint{0.664014in}{1.002614in}}%
\pgfpathlineto{\pgfqpoint{0.655385in}{1.010552in}}%
\pgfpathlineto{\pgfqpoint{0.703568in}{1.000840in}}%
\pgfpathlineto{\pgfqpoint{0.753877in}{0.992408in}}%
\pgfpathlineto{\pgfqpoint{0.760589in}{0.984787in}}%
\pgfpathlineto{\pgfqpoint{0.711264in}{0.993071in}}%
\pgfusepath{fill}%
\end{pgfscope}%
\begin{pgfscope}%
\pgfpathrectangle{\pgfqpoint{0.407592in}{0.222904in}}{\pgfqpoint{1.411163in}{1.411163in}}%
\pgfusepath{clip}%
\pgfsetbuttcap%
\pgfsetroundjoin%
\definecolor{currentfill}{rgb}{0.749020,0.772549,0.792157}%
\pgfsetfillcolor{currentfill}%
\pgfsetfillopacity{0.160000}%
\pgfsetlinewidth{0.000000pt}%
\definecolor{currentstroke}{rgb}{0.000000,0.000000,0.000000}%
\pgfsetstrokecolor{currentstroke}%
\pgfsetdash{}{0pt}%
\pgfpathmoveto{\pgfqpoint{0.982633in}{0.936550in}}%
\pgfpathlineto{\pgfqpoint{0.980283in}{0.943670in}}%
\pgfpathlineto{\pgfqpoint{0.977921in}{0.950826in}}%
\pgfpathlineto{\pgfqpoint{1.031145in}{0.949559in}}%
\pgfpathlineto{\pgfqpoint{1.084665in}{0.949700in}}%
\pgfpathlineto{\pgfqpoint{1.084782in}{0.942565in}}%
\pgfpathlineto{\pgfqpoint{1.084898in}{0.935467in}}%
\pgfpathlineto{\pgfqpoint{1.033625in}{0.935330in}}%
\pgfpathlineto{\pgfqpoint{0.982633in}{0.936550in}}%
\pgfpathclose%
\pgfusepath{fill}%
\end{pgfscope}%
\begin{pgfscope}%
\pgfpathrectangle{\pgfqpoint{0.407592in}{0.222904in}}{\pgfqpoint{1.411163in}{1.411163in}}%
\pgfusepath{clip}%
\pgfsetbuttcap%
\pgfsetroundjoin%
\definecolor{currentfill}{rgb}{0.749020,0.772549,0.792157}%
\pgfsetfillcolor{currentfill}%
\pgfsetfillopacity{0.160000}%
\pgfsetlinewidth{0.000000pt}%
\definecolor{currentstroke}{rgb}{0.000000,0.000000,0.000000}%
\pgfsetstrokecolor{currentstroke}%
\pgfsetdash{}{0pt}%
\pgfpathmoveto{\pgfqpoint{0.773915in}{0.969658in}}%
\pgfpathlineto{\pgfqpoint{0.767268in}{0.977204in}}%
\pgfpathlineto{\pgfqpoint{0.760589in}{0.984787in}}%
\pgfpathlineto{\pgfqpoint{0.811726in}{0.977820in}}%
\pgfpathlineto{\pgfqpoint{0.864397in}{0.972217in}}%
\pgfpathlineto{\pgfqpoint{0.868990in}{0.964860in}}%
\pgfpathlineto{\pgfqpoint{0.873560in}{0.957540in}}%
\pgfpathlineto{\pgfqpoint{0.823008in}{0.962941in}}%
\pgfpathlineto{\pgfqpoint{0.773915in}{0.969658in}}%
\pgfpathclose%
\pgfusepath{fill}%
\end{pgfscope}%
\begin{pgfscope}%
\pgfpathrectangle{\pgfqpoint{0.407592in}{0.222904in}}{\pgfqpoint{1.411163in}{1.411163in}}%
\pgfusepath{clip}%
\pgfsetbuttcap%
\pgfsetroundjoin%
\definecolor{currentfill}{rgb}{0.749020,0.772549,0.792157}%
\pgfsetfillcolor{currentfill}%
\pgfsetfillopacity{0.160000}%
\pgfsetlinewidth{0.000000pt}%
\definecolor{currentstroke}{rgb}{0.000000,0.000000,0.000000}%
\pgfsetstrokecolor{currentstroke}%
\pgfsetdash{}{0pt}%
\pgfpathmoveto{\pgfqpoint{1.304313in}{0.979110in}}%
\pgfpathlineto{\pgfqpoint{1.308680in}{0.986628in}}%
\pgfpathlineto{\pgfqpoint{1.362338in}{0.993994in}}%
\pgfpathlineto{\pgfqpoint{1.414415in}{1.002688in}}%
\pgfpathlineto{\pgfqpoint{1.407967in}{0.994887in}}%
\pgfpathlineto{\pgfqpoint{1.356910in}{0.986345in}}%
\pgfusepath{fill}%
\end{pgfscope}%
\begin{pgfscope}%
\pgfpathrectangle{\pgfqpoint{0.407592in}{0.222904in}}{\pgfqpoint{1.411163in}{1.411163in}}%
\pgfusepath{clip}%
\pgfsetbuttcap%
\pgfsetroundjoin%
\definecolor{currentfill}{rgb}{0.749020,0.772549,0.792157}%
\pgfsetfillcolor{currentfill}%
\pgfsetfillopacity{0.160000}%
\pgfsetlinewidth{0.000000pt}%
\definecolor{currentstroke}{rgb}{0.000000,0.000000,0.000000}%
\pgfsetstrokecolor{currentstroke}%
\pgfsetdash{}{0pt}%
\pgfpathmoveto{\pgfqpoint{1.191375in}{0.954195in}}%
\pgfpathlineto{\pgfqpoint{1.193510in}{0.961451in}}%
\pgfpathlineto{\pgfqpoint{1.195657in}{0.968745in}}%
\pgfpathlineto{\pgfqpoint{1.250462in}{0.973230in}}%
\pgfpathlineto{\pgfqpoint{1.304313in}{0.979110in}}%
\pgfpathlineto{\pgfqpoint{1.299967in}{0.971629in}}%
\pgfpathlineto{\pgfqpoint{1.295644in}{0.964185in}}%
\pgfpathlineto{\pgfqpoint{1.243962in}{0.958517in}}%
\pgfpathlineto{\pgfqpoint{1.191375in}{0.954195in}}%
\pgfpathclose%
\pgfusepath{fill}%
\end{pgfscope}%
\begin{pgfscope}%
\pgfpathrectangle{\pgfqpoint{0.407592in}{0.222904in}}{\pgfqpoint{1.411163in}{1.411163in}}%
\pgfusepath{clip}%
\pgfsetbuttcap%
\pgfsetroundjoin%
\definecolor{currentfill}{rgb}{0.749020,0.772549,0.792157}%
\pgfsetfillcolor{currentfill}%
\pgfsetfillopacity{0.160000}%
\pgfsetlinewidth{0.000000pt}%
\definecolor{currentstroke}{rgb}{0.000000,0.000000,0.000000}%
\pgfsetstrokecolor{currentstroke}%
\pgfsetdash{}{0pt}%
\pgfpathmoveto{\pgfqpoint{0.873560in}{0.957540in}}%
\pgfpathlineto{\pgfqpoint{0.868990in}{0.964860in}}%
\pgfpathlineto{\pgfqpoint{0.864397in}{0.972217in}}%
\pgfpathlineto{\pgfqpoint{0.918309in}{0.968017in}}%
\pgfpathlineto{\pgfqpoint{0.973161in}{0.965250in}}%
\pgfpathlineto{\pgfqpoint{0.975547in}{0.958019in}}%
\pgfpathlineto{\pgfqpoint{0.977921in}{0.950826in}}%
\pgfpathlineto{\pgfqpoint{0.925294in}{0.953493in}}%
\pgfpathlineto{\pgfqpoint{0.873560in}{0.957540in}}%
\pgfpathclose%
\pgfusepath{fill}%
\end{pgfscope}%
\begin{pgfscope}%
\pgfpathrectangle{\pgfqpoint{0.407592in}{0.222904in}}{\pgfqpoint{1.411163in}{1.411163in}}%
\pgfusepath{clip}%
\pgfsetbuttcap%
\pgfsetroundjoin%
\definecolor{currentfill}{rgb}{0.749020,0.772549,0.792157}%
\pgfsetfillcolor{currentfill}%
\pgfsetfillopacity{0.160000}%
\pgfsetlinewidth{0.000000pt}%
\definecolor{currentstroke}{rgb}{0.000000,0.000000,0.000000}%
\pgfsetstrokecolor{currentstroke}%
\pgfsetdash{}{0pt}%
\pgfpathmoveto{\pgfqpoint{1.084665in}{0.949700in}}%
\pgfpathlineto{\pgfqpoint{1.084547in}{0.956872in}}%
\pgfpathlineto{\pgfqpoint{1.084429in}{0.964081in}}%
\pgfpathlineto{\pgfqpoint{1.140207in}{0.965689in}}%
\pgfpathlineto{\pgfqpoint{1.195657in}{0.968745in}}%
\pgfpathlineto{\pgfqpoint{1.193510in}{0.961451in}}%
\pgfpathlineto{\pgfqpoint{1.191375in}{0.954195in}}%
\pgfpathlineto{\pgfqpoint{1.138175in}{0.951249in}}%
\pgfpathlineto{\pgfqpoint{1.084665in}{0.949700in}}%
\pgfpathclose%
\pgfusepath{fill}%
\end{pgfscope}%
\begin{pgfscope}%
\pgfpathrectangle{\pgfqpoint{0.407592in}{0.222904in}}{\pgfqpoint{1.411163in}{1.411163in}}%
\pgfusepath{clip}%
\pgfsetbuttcap%
\pgfsetroundjoin%
\definecolor{currentfill}{rgb}{0.749020,0.772549,0.792157}%
\pgfsetfillcolor{currentfill}%
\pgfsetfillopacity{0.160000}%
\pgfsetlinewidth{0.000000pt}%
\definecolor{currentstroke}{rgb}{0.000000,0.000000,0.000000}%
\pgfsetstrokecolor{currentstroke}%
\pgfsetdash{}{0pt}%
\pgfpathmoveto{\pgfqpoint{0.760589in}{0.984787in}}%
\pgfpathlineto{\pgfqpoint{0.753877in}{0.992408in}}%
\pgfpathlineto{\pgfqpoint{0.806042in}{0.985315in}}%
\pgfpathlineto{\pgfqpoint{0.859780in}{0.979611in}}%
\pgfpathlineto{\pgfqpoint{0.864397in}{0.972217in}}%
\pgfpathlineto{\pgfqpoint{0.811726in}{0.977820in}}%
\pgfusepath{fill}%
\end{pgfscope}%
\begin{pgfscope}%
\pgfpathrectangle{\pgfqpoint{0.407592in}{0.222904in}}{\pgfqpoint{1.411163in}{1.411163in}}%
\pgfusepath{clip}%
\pgfsetbuttcap%
\pgfsetroundjoin%
\definecolor{currentfill}{rgb}{0.749020,0.772549,0.792157}%
\pgfsetfillcolor{currentfill}%
\pgfsetfillopacity{0.160000}%
\pgfsetlinewidth{0.000000pt}%
\definecolor{currentstroke}{rgb}{0.000000,0.000000,0.000000}%
\pgfsetstrokecolor{currentstroke}%
\pgfsetdash{}{0pt}%
\pgfpathmoveto{\pgfqpoint{0.977921in}{0.950826in}}%
\pgfpathlineto{\pgfqpoint{0.975547in}{0.958019in}}%
\pgfpathlineto{\pgfqpoint{0.973161in}{0.965250in}}%
\pgfpathlineto{\pgfqpoint{1.028640in}{0.963935in}}%
\pgfpathlineto{\pgfqpoint{1.084429in}{0.964081in}}%
\pgfpathlineto{\pgfqpoint{1.084547in}{0.956872in}}%
\pgfpathlineto{\pgfqpoint{1.084665in}{0.949700in}}%
\pgfpathlineto{\pgfqpoint{1.031145in}{0.949559in}}%
\pgfpathlineto{\pgfqpoint{0.977921in}{0.950826in}}%
\pgfpathclose%
\pgfusepath{fill}%
\end{pgfscope}%
\begin{pgfscope}%
\pgfpathrectangle{\pgfqpoint{0.407592in}{0.222904in}}{\pgfqpoint{1.411163in}{1.411163in}}%
\pgfusepath{clip}%
\pgfsetbuttcap%
\pgfsetroundjoin%
\definecolor{currentfill}{rgb}{0.749020,0.772549,0.792157}%
\pgfsetfillcolor{currentfill}%
\pgfsetfillopacity{0.160000}%
\pgfsetlinewidth{0.000000pt}%
\definecolor{currentstroke}{rgb}{0.000000,0.000000,0.000000}%
\pgfsetstrokecolor{currentstroke}%
\pgfsetdash{}{0pt}%
\pgfpathmoveto{\pgfqpoint{1.195657in}{0.968745in}}%
\pgfpathlineto{\pgfqpoint{1.197815in}{0.976077in}}%
\pgfpathlineto{\pgfqpoint{1.253736in}{0.980642in}}%
\pgfpathlineto{\pgfqpoint{1.308680in}{0.986628in}}%
\pgfpathlineto{\pgfqpoint{1.304313in}{0.979110in}}%
\pgfpathlineto{\pgfqpoint{1.250462in}{0.973230in}}%
\pgfusepath{fill}%
\end{pgfscope}%
\begin{pgfscope}%
\pgfpathrectangle{\pgfqpoint{0.407592in}{0.222904in}}{\pgfqpoint{1.411163in}{1.411163in}}%
\pgfusepath{clip}%
\pgfsetbuttcap%
\pgfsetroundjoin%
\definecolor{currentfill}{rgb}{0.749020,0.772549,0.792157}%
\pgfsetfillcolor{currentfill}%
\pgfsetfillopacity{0.160000}%
\pgfsetlinewidth{0.000000pt}%
\definecolor{currentstroke}{rgb}{0.000000,0.000000,0.000000}%
\pgfsetstrokecolor{currentstroke}%
\pgfsetdash{}{0pt}%
\pgfpathmoveto{\pgfqpoint{0.864397in}{0.972217in}}%
\pgfpathlineto{\pgfqpoint{0.859780in}{0.979611in}}%
\pgfpathlineto{\pgfqpoint{0.914790in}{0.975335in}}%
\pgfpathlineto{\pgfqpoint{0.970763in}{0.972517in}}%
\pgfpathlineto{\pgfqpoint{0.973161in}{0.965250in}}%
\pgfpathlineto{\pgfqpoint{0.918309in}{0.968017in}}%
\pgfusepath{fill}%
\end{pgfscope}%
\begin{pgfscope}%
\pgfpathrectangle{\pgfqpoint{0.407592in}{0.222904in}}{\pgfqpoint{1.411163in}{1.411163in}}%
\pgfusepath{clip}%
\pgfsetbuttcap%
\pgfsetroundjoin%
\definecolor{currentfill}{rgb}{0.749020,0.772549,0.792157}%
\pgfsetfillcolor{currentfill}%
\pgfsetfillopacity{0.160000}%
\pgfsetlinewidth{0.000000pt}%
\definecolor{currentstroke}{rgb}{0.000000,0.000000,0.000000}%
\pgfsetstrokecolor{currentstroke}%
\pgfsetdash{}{0pt}%
\pgfpathmoveto{\pgfqpoint{1.084429in}{0.964081in}}%
\pgfpathlineto{\pgfqpoint{1.084310in}{0.971328in}}%
\pgfpathlineto{\pgfqpoint{1.141231in}{0.972965in}}%
\pgfpathlineto{\pgfqpoint{1.197815in}{0.976077in}}%
\pgfpathlineto{\pgfqpoint{1.195657in}{0.968745in}}%
\pgfpathlineto{\pgfqpoint{1.140207in}{0.965689in}}%
\pgfusepath{fill}%
\end{pgfscope}%
\begin{pgfscope}%
\pgfpathrectangle{\pgfqpoint{0.407592in}{0.222904in}}{\pgfqpoint{1.411163in}{1.411163in}}%
\pgfusepath{clip}%
\pgfsetbuttcap%
\pgfsetroundjoin%
\definecolor{currentfill}{rgb}{0.749020,0.772549,0.792157}%
\pgfsetfillcolor{currentfill}%
\pgfsetfillopacity{0.160000}%
\pgfsetlinewidth{0.000000pt}%
\definecolor{currentstroke}{rgb}{0.000000,0.000000,0.000000}%
\pgfsetstrokecolor{currentstroke}%
\pgfsetdash{}{0pt}%
\pgfpathmoveto{\pgfqpoint{0.973161in}{0.965250in}}%
\pgfpathlineto{\pgfqpoint{0.970763in}{0.972517in}}%
\pgfpathlineto{\pgfqpoint{1.027378in}{0.971179in}}%
\pgfpathlineto{\pgfqpoint{1.084310in}{0.971328in}}%
\pgfpathlineto{\pgfqpoint{1.084429in}{0.964081in}}%
\pgfpathlineto{\pgfqpoint{1.028640in}{0.963935in}}%
\pgfusepath{fill}%
\end{pgfscope}%
\end{pgfpicture}%
\makeatother%
\endgroup%
}
  \captionsetup{font=small}
  \caption{\textbf{Bures trajectories in the covariance cone.}
  (a) Equal diagonal blocks: invariant hyperbolas.
  (b) Unequal blocks: fixed-\(b\) sections of determinant leaves.
  The displayed trajectories converge to \(r=\beta\).}
  \label{fig:gaussian-w2-cone-geometry}
\end{wrapfigure}
Under special conditions on the conserved spectrum of $J\Sigma$, elimination of hidden blocks becomes possible: Appendix~\ref{app:intrinsic} describes these intrinsic dynamics.

\paragraph{Numerical illustration: the $2\times2$ covariance cone}
For scalar input and output, every covariance has the form
$\Sigma=\begin{psmallmatrix}p&r\\r&q\end{psmallmatrix}\succeq0$.
Take the loss $F(\Sigma):=(r-\beta)^2/2$, with scalar teacher $\beta$.
Writing $(a,b):=((p+q)/2,(p-q)/2)$ identifies $\Sigma\succeq0$ with
the Lorentz cone $a^2\ge r^2+b^2$, $a\ge0$. The Bures flow
\eqref{eq:block-flow-general} becomes
$(\dot a,\dot b,\dot r)=-(2r(r-\beta),0,2a(r-\beta))$,
preserving $b=b_0$ and $\alpha^2:=a^2-r^2-b^2=pq-r^2$.

When $p=q$, set $s:=p$. Then $b=0$ and trajectories follow hyperbolas
$s^2-r^2=\alpha^2$ in $s\ge|r|$ (Figure~\ref{fig:gaussian-w2-cone-geometry}(a)).
For $\alpha>0$, they converge to $(r,s)=(\beta,\sqrt{\beta^2+\alpha^2})$;
the spectral-margin bound is exact for the minimum $s$ on the leaf.
Unequal diagonal blocks give intersections of the planes $b=b_0$ with
determinant leaves in the three-dimensional cone (panel (b)).
If $b_0^2+\alpha^2>0$, the endpoint is
$(r_\infty,a_\infty)=(\beta,\sqrt{\beta^2+b_0^2+\alpha^2})$.

The invariants $2b$ and $-\alpha^2$ are the trace and determinant of
$J\Sigma$, whose eigenvalues are $b_0\pm\sqrt{b_0^2+\alpha^2}$.
Thus $\gap{\Sigma_0}=2\sqrt{b_0^2+\alpha^2}$: the convergence bound
covers both panels, including singular leaves with $b_0\ne0$.
Remark~\ref{rem:degenerate-boundary-leaves} in Appendix~\ref{app:gf}
treats the zero-support-gap boundary. Figure~\ref{fig:gaussian-w2-cone-geometry}
is obtained by integrating $\dot r=-2\sqrt{r^2+b_0^2+\alpha^2}(r-\beta)$.

%%% Local Variables:
%%% mode: LaTeX
%%% TeX-master: "../main"
%%% End:
